% 纲要标题：# 第一卷　古典几何与光滑流形
% 纲要位置：几何学计划纲要.md，第 111 行。
% 正文状态：二十章正文及附录 A—E 初稿；前言、后记及前两编导言待撰写。

% 纲要内容开始
% 
% > 由几何等价和局部模型进入光滑结构、微积分、可积性及对称。古典内容作为研究生的结构性回顾，不按本科知识点逐条分章；不以完成全部古典分类作为流形路线的门槛。
% 
% 纲要内容结束

% !TeX TXS-program:compile = cmd /c build-Book1.bat -NoPause
% !TeX TXS-program:view = txs:///view-pdf-internal --embedded "Book1.pdf"

\documentclass[UTF8,zihao=-4,openany,fontset=fandol]{ctexbook}
\newcommand{\BookTitle}{古典几何与光滑流形}
\newcommand{\BookVolumeName}{第一卷}
\newcommand{\BookNumber}{1}
\newcommand{\BookDraftDate}{2026年10月}
\newcommand{\BookStatus}{写作稿·二十章及附录 A—E 初稿}
% 《几何学》五卷共用的基础排版与语义宏。
% 各卷主文件须先定义 \BookTitle、\BookVolumeName 和 \BookDraftDate。
\newcommand{\BookAuthor}{R. EverStarine}
% 稿件状态由各卷元数据指定；尚未开写的卷沿用框架状态。
\providecommand{\BookStatus}{框架预览·正文未撰写}

\usepackage[a4paper,inner=28mm,outer=24mm,top=27mm,bottom=28mm,
  headheight=16pt,headsep=8mm,footskip=12mm]{geometry}

\usepackage{amsmath}
\usepackage{mathtools}
\usepackage{amsthm}
\usepackage{aliascnt}
\usepackage{unicode-math}
% 数学字母使用默认数学字体内的对应字形；禁止 mathrm 等默认继承正文族。
% 运算符（sin、log、operatorname 等）随 mathrm 使用数学正体。
\unimathsetup{mathrm=sym,mathit=sym,mathbf=sym,mathsf=sym,mathtt=sym}
\setmainfont{texgyrepagella-regular.otf}[
  BoldFont=texgyrepagella-bold.otf,
  ItalicFont=texgyrepagella-italic.otf,
  BoldItalicFont=texgyrepagella-bolditalic.otf
]
\setsansfont{texgyreheros-regular.otf}[
  BoldFont=texgyreheros-bold.otf,
  ItalicFont=texgyreheros-italic.otf,
  BoldItalicFont=texgyreheros-bolditalic.otf
]
\setmonofont{lmmonolt10-regular.otf}[
  BoldFont=lmmonolt10-bold.otf,
  ItalicFont=lmmonolt10-oblique.otf,
  BoldItalicFont=lmmonolt10-boldoblique.otf
]
% 人物注释中的古希腊文、斯拉夫文字、日文异体字与希伯来文原名使用有相应字形的字体。
\newfontfamily\GeometryOriginalNameFont{Times New Roman}
\NewDocumentCommand{\OriginalName}{m}{{\GeometryOriginalNameFont #1}}
\newCJKfontfamily\GeometryOriginalHanFont{SimSun}
\NewDocumentCommand{\OriginalHanName}{m}{{\GeometryOriginalHanFont #1}}
\newfontfamily\GeometryHebrewFont{Times New Roman}[Script=Hebrew]
\TeXXeTstate=1
\NewDocumentCommand{\OriginalHebrewName}{m}{{\GeometryHebrewFont\beginR#1\endR}}
% 行内、行间及标题内公式均使用 unicode-math 默认 Latin Modern Math。
% 公式之外的正文西文仍由 setmainfont 使用 Pagella；黑板粗体单独采用 AMS。
\let\mathbb\relax
\DeclareMathAlphabet{\mathbb}{U}{msb}{m}{n}
% unicode-math 的 Latin Modern Math 将 \varnothing 映为瘦长椭圆；空集单独
% 采用传统 AMS MSBM 的圆形斜杠字形，不改变其余数学符号的字体。
\DeclareSymbolFont{GeometryAMSb}{U}{msb}{m}{n}
\DeclareMathSymbol{\GeometryRoundEmptySet}{\mathord}{GeometryAMSb}{"3F}
% 当前 Latin Modern Math 缺少横截符号 U+22D4；仅此关系符使用 AMSa 的现有字形。
\DeclareSymbolFont{GeometryAMSa}{U}{msa}{m}{n}
\DeclareMathSymbol{\GeometryTransverse}{\mathrel}{GeometryAMSa}{"74}
\AtBeginDocument{%
  \let\varnothing\GeometryRoundEmptySet
  \let\emptyset\GeometryRoundEmptySet
  \let\pitchfork\GeometryTransverse
}

\usepackage{microtype}
\usepackage{xcolor}
\usepackage{graphicx}
\usepackage{booktabs}
\usepackage{longtable}
\usepackage{array}
\usepackage{tabularx}
\usepackage{algorithm}
\usepackage{algpseudocode}
\usepackage{enumitem}
\usepackage{fancyhdr}
% 双面排版补齐奇数页时，空白页不残留上一章的页眉或页码。
\usepackage{emptypage}
\usepackage[most]{tcolorbox}
\usepackage{tikz-cd}
\usepackage{csquotes}
\usepackage[
  backend=biber,
  style=numeric-comp,
  sorting=nyt,
  giveninits=true,
  maxbibnames=99,
  doi=true,
  url=true,
  isbn=true
]{biblatex}
\addbibresource{../Shared/References.bib}
% 骨架没有引文时保留书目入口；写入真实条目后由完整构建自动排入书目。
\newif\ifGeometryHasBibliography
\AtDataInput{\global\GeometryHasBibliographytrue}
\NewDocumentCommand{\PrintBookBibliography}{}{%
  \ifGeometryHasBibliography
    \printbibliography[heading=bibintoc,title={参考文献}]%
  \else
    \chapter*{参考文献}%
    \addcontentsline{toc}{chapter}{参考文献}%
    \noindent 待撰写\par
  \fi
}
% 访问日期保留在书目库供内部核验，成品参考文献不显示 visited on。
\renewbibmacro*{url+urldate}{\usebibmacro{url}}
% biblatex 在书目内重绑定格式宏；在其作用域中将强调统一为正体，
% 同时避免书名中的中文及标点仍继承旧的斜体样式。
\AtBeginBibliography{\let\emph\textup\let\textit\textup}
% 中文条目的责任者连接符：不可套用西文的 “and”，按 GB/T 7714—2015 的著录
% 惯例改用逗号分隔。中文条目在 Shared/References.bib 中以 langid = {chinese}
% 标注（含译本，因其著录信息为中文形式），西文条目保持 “and” 不变。
% 按条目切换须在每条书目项处理时判定；\DeclareDelimFormat 的可选参数不是
% 语言选择器（实测无效），故用 \AtEveryBibitem 配合 \iffieldequalstr。
\AtEveryBibitem{%
  \iffieldequalstr{langid}{chinese}%
    {\DeclareDelimFormat{finalnamedelim}{\addcomma\space}%
     \DeclareDelimFormat{multinamedelim}{\addcomma\space}}%
    {\DeclareDelimFormat{finalnamedelim}{\addspace\bibstring{and}\space}%
     \DeclareDelimFormat{multinamedelim}{\addcomma\space}}%
}
\usepackage{etoolbox}
\usepackage{xparse}
\usepackage[noautomatic]{imakeidx}

\definecolor{GeometryBlue}{HTML}{244A64}
\definecolor{GeometryBlueLight}{HTML}{EEF4F7}
\definecolor{GeometryGray}{HTML}{4F5962}
\definecolor{GeometryGrayLight}{HTML}{F4F5F6}

\ctexset{
  part = {
    % 编题页分两行：第一行“第X编”，与封面“第一卷”同款（大号无衬线灰色）；
    % 第二行编名，与封面“几何学”同款（大号无衬线蓝色粗体），字号保持原样。
    % 第一行与行距由 number／aftername 给出，format 只负责编名的字体，
    % 以免字号被重置或被重复输出。
    format = {\centering\fontsize{30}{38}\selectfont\sffamily\bfseries\color{GeometryBlue}},
    name = {,},
    number = {\centering\normalfont\normalsize{\fontsize{24}{30}\selectfont\sffamily\color{GeometryGray}第\chinese{part}编}},
    aftername = {\par\vspace{0.7cm}},
    beforeskip = 0pt,
    afterskip = 2.5\baselineskip,
    pagestyle = plain
  },
  chapter = {
    format = \raggedright\Huge\bfseries\color{GeometryBlue},
    name = {第,章},
    number = \chinese{chapter},
    aftername = \quad,
    beforeskip = 0pt,
    afterskip = 30pt,
    pagestyle = plain
  },
  section = {
    format = \Large\bfseries\color{GeometryBlue},
    beforeskip = 2.4ex plus .6ex minus .2ex,
    afterskip = 1.1ex plus .2ex
  },
  subsection = {
    format = \large\bfseries\color{GeometryGray},
    beforeskip = 2ex plus .4ex minus .2ex,
    afterskip = .8ex plus .2ex
  }
}

\setlength{\parindent}{2em}
\setlength{\parskip}{0pt}
\linespread{1.18}
% 仅用于两个连续的普通正文段落；列表、公式、定理、标题等结构边界不用。
% 段间额外留白统一由这一长度控制，避免改动全局 \parskip 后扰乱各类环境。
\newlength{\GeometryProseParagraphGap}
\AtBeginDocument{%
  \setlength{\GeometryProseParagraphGap}{0.40\baselineskip}%
}
\NewDocumentCommand{\ProseParagraphBreak}{}{%
  \par\addvspace{\GeometryProseParagraphGap}%
}
% 与代数学保持一致：小节编号并进入目录，便于按具体主题定位。
\setcounter{secnumdepth}{3}
\setcounter{tocdepth}{2}
\allowdisplaybreaks[2]
\emergencystretch=2em
\clubpenalty=10000
\widowpenalty=10000
\displaywidowpenalty=10000

% 固定列表内部段距，避免 flushbottom 将默认的弹性 \parsep 拉大到超过正文段距。
\setlist[itemize]{leftmargin=2.2em,itemsep=.35ex,parsep=0pt,topsep=.6ex}
\setlist[enumerate]{leftmargin=2.5em,itemsep=.35ex,parsep=0pt,topsep=.6ex,
  font=\normalfont\upshape}
\setlist[enumerate,1]{label=\arabic*.,ref=\arabic*}
\setlist[enumerate,2]{label=\arabic*),ref=\arabic*}
\setlist[enumerate,3]{label=\alph*),ref=\alph*}
\setlist[description]{leftmargin=2.5em,labelindent=0pt,itemsep=.35ex,parsep=0pt,topsep=.6ex}

% “下列命题等价”专用列表；普通罗列仍使用 enumerate。
\newlist{equivalences}{enumerate}{2}
\setlist[equivalences,1]{leftmargin=2.8em,itemsep=.35ex,parsep=0pt,topsep=.6ex,
  label=\Roman*.,ref=\Roman*,font=\normalfont\upshape}
\setlist[equivalences,2]{leftmargin=2.8em,itemsep=.35ex,parsep=0pt,topsep=.6ex,
  label=\roman*),ref=\roman*,font=\normalfont\upshape}

\pagestyle{fancy}
\fancyhf{}
\fancyhead[LE]{\small\color{GeometryGray}\thepage\quad\leftmark}
% 奇数页采用页前最后有效的节标记，避免页中开始新节时提前更换页眉。
% 新节或习题的标记从后续页生效；章首页使用 plain，不显示上一章的标记。
\fancyhead[RO]{\small\color{GeometryGray}\TopMark{2e-right}\quad\thepage}
\renewcommand{\headrulewidth}{0.35pt}
% 保留数学标题的原始大小写，避免 book 类的 \MakeUppercase
% 与 unicode-math 命令在页眉标记中发生冲突。
% 偶数页页眉显示“章序＋章标题”（如“第一章 代数对象、映射与基本构造”）；
% 章序与章标题之间以 \quad 分隔，间距与奇数页节眉“节号＋节标题”之间一致；
% 附录章的章序由 \CTEXthechapter 按本卷附录字母（如 A）给出。
\renewcommand{\chaptermark}[1]{\markboth{\CTEXthechapter\quad #1}{}}
\renewcommand{\sectionmark}[1]{\markright{\thesection\quad #1}}
% 章末习题不编号，但进入习题页后应覆盖上一节留下的奇数页页眉；
% 同时写入目录项。PDF 书签由 hyperref 对 \section*{习题} 自动生成，
% 不再另行 \pdfbookmark，避免“习题”下再嵌套一个“习题”。
\newcommand{\BookExercises}{%
  \section*{习题}%
  \markright{习题}%
  \addcontentsline{toc}{section}{习题}%
}
\fancypagestyle{plain}{
  \fancyhf{}
  \fancyfoot[C]{\small\color{GeometryGray}\thepage}
  \renewcommand{\headrulewidth}{0pt}
}

\usepackage[unicode,colorlinks=true,linkcolor=GeometryBlue,
  citecolor=GeometryBlue,urlcolor=GeometryBlue,pdfborder={0 0 0}]{hyperref}
\usepackage{bookmark}
% PDF 书签沿用目录条目中的编号：正文章为“第几章”，附录章为“附录 A”，
% 节与小节分别保留 x.x、x.x.x；无编号的导读、习题和卷末单元仍不加号。
\bookmarksetup{numbered}

% unicode-math 会调整原始 Unicode 希腊字母的类别码。若把这些字母直接写进
% \texorpdfstring 的书签分支，辅助文件回读时类别码不同，LaTeX 会在每一轮都
% 误报标签发生变化。书签分支统一写成下列可展开宏，既保留希腊字母显示，
% 又使写入和回读的标签文本稳定。
\newcommand{\GeometryPdfAlpha}{α}
\newcommand{\GeometryPdfGamma}{Γ}
\newcommand{\GeometryPdfLambda}{λ}
\newcommand{\GeometryPdfMu}{μ}
\newcommand{\GeometryPdfPi}{π}
\newcommand{\GeometryPdfSigma}{σ}
\newcommand{\GeometryPdfTheta}{θ}
\newcommand{\GeometryPdfXi}{ξ}
\newcommand{\GeometryPdfZeta}{ζ}
\newcommand{\GeometryPdfCapitalLambda}{Λ}
\newcommand{\GeometryPdfOmega}{Ω}
\newcommand{\GeometryPdfCapitalSigma}{Σ}
\newcommand{\GeometryPdfDelta}{δ}
\newcommand{\GeometryPdfEpsilon}{ε}
\newcommand{\GeometryPdfIota}{ι}
\newcommand{\GeometryPdfSmallOmega}{ω}
\newcommand{\GeometryPdfEll}{ℓ}

\usepackage[nameinlink,noabbrev]{cleveref}
\newcommand{\crefpairconjunction}{和}
\newcommand{\crefmiddleconjunction}{、}
\newcommand{\creflastconjunction}{和}

% CTeX 2.5.10 的标题前缀宏带一个参数，而 TeX Live 2025 中 cleveref
% 仍按普通前缀宏展开它们，会把节号写成含 \endcsname 的损坏元数据。
% 改为恒等包装后，CTeX 传入的标题编号保持不变，\ref 与 \cref 均可正常工作。
\makeatletter
\def\p@part#1{#1}
\def\p@chapter#1{#1}
\def\p@section#1{#1}
\def\p@subsection#1{#1}
\def\p@subsubsection#1{#1}
\def\p@paragraph#1{#1}
\def\p@subparagraph#1{#1}
\makeatother
% 非数学西文字母统一正体；不改变括号类别、标点或数学字体。
% 2026-09-11：非数学模式的西文字母统一使用 Pagella 正体，保留粗细区别。
% 保留字母级边界，不重分类数字、括号、标点及中文。
% 此设置也覆盖列表序号字母和术语原文；数学模式保持原样。
\newfontfamily\GeometryLetterFont{texgyrepagella-regular.otf}[
  BoldFont=texgyrepagella-bold.otf,
  ItalicFont=texgyrepagella-italic.otf,
  BoldItalicFont=texgyrepagella-bolditalic.otf
]
\ExplSyntaxOn
\bool_new:N \l_geometry_text_letters_bool
\bool_set_true:N \l_geometry_text_letters_bool
\cs_new_protected:Npn \GeometryPreserveMathText
  { \bool_set_false:N \l_geometry_text_letters_bool }
\cs_new_protected:Npn \geometry_letter_group_begin:
  {
    \group_begin:
    \bool_if:NT \l_geometry_text_letters_bool
      { \GeometryLetterFont \upshape }
  }
% 继承 Default 的全部 xeCJK 转移，保留中西文间距及标点处理；
% 不能用只有开关分组的转移覆盖 Boundary 或 CJK 的原处理。
\xeCJK_new_class:n { GeometryLatinLetter }
\seq_map_inline:Nn \g__xeCJK_class_seq
  {
    \str_if_eq:nnF {#1} { GeometryLatinLetter }
      {
        \xeCJK_copy_inter_class_toks:nnnn
          {#1} {GeometryLatinLetter} {#1} {Default}
        \xeCJK_copy_inter_class_toks:nnnn
          {GeometryLatinLetter} {#1} {Default} {#1}
        \xeCJK_app_inter_class_toks:nnn {#1} {GeometryLatinLetter}
          { \geometry_letter_group_begin: }
        \xeCJK_pre_inter_class_toks:nnn {GeometryLatinLetter} {#1}
          { \group_end: }
      }
  }
% ASCII 字母及常用拉丁扩展字母；明确跳过乘号、除号，不重分类标点。
\xeCJKDeclareCharClass { GeometryLatinLetter }
  { "0041->"005A, "0061->"007A, "00C0->"00D6, "00D8->"00F6,
    "00F8->"024F, "1E00->"1EFF }
\ExplSyntaxOff
% 数学组内局部关闭字母切换，连公式中 \text{...} 的字体也保持原样。
% 只追加，不覆盖 amsmath、unicode-math、xeCJK 已有的数学初始化。
\AtBeginDocument{%
  \everymath\expandafter{\the\everymath\GeometryPreserveMathText}%
  \everydisplay\expandafter{\the\everydisplay\GeometryPreserveMathText}%
}


% 定理类环境（定理、定义、命题、引理、推论、断言、猜想、开放问题、注记、
% 习题）的正文：中文用楷体，西文保持正体，数学字体不受影响。
% 注意 \normalfont 只重置西文字族，中文字族仍随默认（宋体），故须显式 \kaishu；
% 不能退回 \itshape——那会把西文一并变斜。
\newtheoremstyle{geometryplain}
  {1.2ex}{1.2ex}{\normalfont\upshape\kaishu}{}
  {\color{GeometryBlue}\bfseries}{.}{0.6em}{}
\newtheoremstyle{geometrydefinition}
  {1.2ex}{1.2ex}{\normalfont\kaishu}{}
  {\color{GeometryBlue}\bfseries}{.}{0.6em}{}
\newtheoremstyle{geometryremark}
  {1.0ex}{1.0ex}{\normalfont\kaishu}{}
  {\color{GeometryGray}\bfseries}{.}{0.6em}{}

\theoremstyle{geometryplain}
\newtheorem{theorem}{定理}[section]
% 正文默认按“章.节.流水号”编号；未分节的附录省略中间的 0。
\renewcommand{\thetheorem}{%
  \ifnum\value{section}=0
    \thechapter.\arabic{theorem}%
  \else
    \thesection.\arabic{theorem}%
  \fi
}
\newaliascnt{lemma}{theorem}
\newtheorem{lemma}[lemma]{引理}
\aliascntresetthe{lemma}
\newaliascnt{proposition}{theorem}
\newtheorem{proposition}[proposition]{命题}
\aliascntresetthe{proposition}
\newaliascnt{corollary}{theorem}
\newtheorem{corollary}[corollary]{推论}
\aliascntresetthe{corollary}
\newaliascnt{claim}{theorem}
\newtheorem{claim}[claim]{断言}
\aliascntresetthe{claim}
\newaliascnt{conjecture}{theorem}
\newtheorem{conjecture}[conjecture]{猜想}
\aliascntresetthe{conjecture}
\newaliascnt{openproblem}{theorem}
\newtheorem{openproblem}[openproblem]{开放问题}
\aliascntresetthe{openproblem}

\theoremstyle{geometrydefinition}
\newaliascnt{definition}{theorem}
\newtheorem{definition}[definition]{定义}
\aliascntresetthe{definition}
\newaliascnt{example}{theorem}
\newtheorem{example}[example]{例}
\aliascntresetthe{example}
% 习题按章单独编号，不占用定理类环境的流水号。
\newtheorem{exercise}{习题}[chapter]

\theoremstyle{geometryremark}
\newaliascnt{remark}{theorem}
\newtheorem{remark}[remark]{注记}
\aliascntresetthe{remark}
\renewcommand{\proofname}{证明}
% 仅在证明紧接定理类陈述时收紧间距，并按陈述结尾是否为行间公式调整。
% 证明标题（证明／解／证明概要）的中文用楷体，
% 西文保持正体；正文不施加楷体，仍为默认宋体。不能退回 \itshape。
\newif\ifGeometryDirectProof
\newif\ifGeometryCoupledProof
\newif\ifGeometryStatementEndsInDisplay
\newskip\GeometryProofTrailingSkip
% 行间公式结束后，TeX 留下空的水平列表；在定理环境结束前辨认这一状态。
% 证明跟在公式之后时不再减去定理尾距，否则公式下方会比上方更紧。
\newcommand{\GeometryRecordStatementTail}{%
  \ifhmode
    \ifnum\lastnodetype=-1\relax
      \global\GeometryStatementEndsInDisplaytrue
    \else
      \global\GeometryStatementEndsInDisplayfalse
    \fi
  \else
    \global\GeometryStatementEndsInDisplayfalse
  \fi
}
\AddToHook{para/begin}{\global\GeometryDirectProoffalse\global\GeometryStatementEndsInDisplayfalse}
\AddToHook{env/theorem/end}{\GeometryRecordStatementTail}
\AddToHook{env/lemma/end}{\GeometryRecordStatementTail}
\AddToHook{env/proposition/end}{\GeometryRecordStatementTail}
\AddToHook{env/corollary/end}{\GeometryRecordStatementTail}
\AddToHook{env/claim/end}{\GeometryRecordStatementTail}
\AddToHook{env/exercise/end}{\GeometryRecordStatementTail}
\AddToHook{env/theorem/after}{\global\GeometryDirectProoftrue}
\AddToHook{env/lemma/after}{\global\GeometryDirectProoftrue}
\AddToHook{env/proposition/after}{\global\GeometryDirectProoftrue}
\AddToHook{env/corollary/after}{\global\GeometryDirectProoftrue}
\AddToHook{env/claim/after}{\global\GeometryDirectProoftrue}
\AddToHook{env/exercise/after}{\global\GeometryDirectProoftrue}
\AddToHook{env/figure/before}{\global\GeometryDirectProoffalse}
\AddToHook{env/table/before}{\global\GeometryDirectProoffalse}
\makeatletter
\renewenvironment{proof}[1][\proofname]{\par
  \pushQED{\qed}%
  \normalfont \topsep6\p@\@plus6\p@\relax
  \GeometryCoupledProoffalse
  \ifGeometryDirectProof
    \GeometryCoupledProoftrue
    \GeometryProofTrailingSkip=\topsep
    \advance\GeometryProofTrailingSkip by \partopsep
    \ifGeometryStatementEndsInDisplay\else
      \kern-5\p@% 只在陈述以普通文字结尾时抵消部分尾距。
    \fi
    \topsep2\p@\partopsep\z@\relax
  \fi
  \global\GeometryDirectProoffalse
  \global\GeometryStatementEndsInDisplayfalse
  \trivlist
  \item[\hskip\labelsep\upshape\kaishu #1\@addpunct{.}]\ignorespaces
}{%
  \popQED
  \ifGeometryCoupledProof\@topsepadd\GeometryProofTrailingSkip\fi
  \endtrivlist\@endpefalse
}
\makeatother
\newenvironment{solution}
  {\begin{proof}[解]}
  {\end{proof}}
\newenvironment{proofsketch}
  {\begin{proof}[证明概要]}
  {\end{proof}}

% 正文选读内容的星号属于编号而不属于标题；第二参数统一接收标签。
\newcommand{\GeometryOptionalMark}{\texorpdfstring{\textsuperscript{*}}{*}}
% 选做习题沿用同一编号星号。可选参数仍只表示习题自身的简短题名，
% 不得再写“选读”或“选做”；星号由本环境自动附在题号右侧。
\NewDocumentEnvironment{optionalexercise}{o}{%
  \begingroup
  \let\GeometrySavedTheExercise\theexercise
  \renewcommand{\theexercise}{\GeometrySavedTheExercise\GeometryOptionalMark}%
  \IfNoValueTF{#1}{\begin{exercise}}{\begin{exercise}[#1]}%
}{%
  \end{exercise}%
  \endgroup
}
% 可选参数只供目录与 PDF 书签使用；正文题名可借此按语义手动换行，
% 又不把换行命令写入目录或书签。
\NewDocumentCommand{\OptionalChapter}{O{}mm}{%
  \begingroup
  \let\GeometrySavedTheChapter\thechapter
  \renewcommand{\thechapter}{\GeometrySavedTheChapter\GeometryOptionalMark}%
  \ctexset{chapter={name={第,章\GeometryOptionalMark},
    number=\chinese{chapter},aftername=\quad}}%
  \ifstrempty{#1}{\chapter{#2}}{\chapter[#1]{#2}}\label{#3}%
  \endgroup
}
\NewDocumentCommand{\OptionalSection}{mm}{%
  \begingroup
  \let\GeometrySavedTheSection\thesection
  \renewcommand{\thesection}{\GeometrySavedTheSection\GeometryOptionalMark}%
  \ctexset{section={number=\thesection}}%
  \section{#1}\label{#2}%
  \endgroup
}
\NewDocumentCommand{\OptionalSubsection}{mm}{%
  \begingroup
  \let\GeometrySavedTheSubsection\thesubsection
  \renewcommand{\thesubsection}{\GeometrySavedTheSubsection\GeometryOptionalMark}%
  \ctexset{subsection={number=\thesubsection}}%
  \subsection{#1}\label{#2}%
  \endgroup
}

% 编题页必须从奇数页开始，题页仅在页脚显示页码，其后保留一张空白页。
\makeatletter
\NewDocumentCommand{\BookPart}{mm}{%
  \cleardoublepage
  \begingroup
  \def\GeometryPartLabel{#2}%
  % \part 会在返回前翻页，故须在 \@endpart 开始处记录题页标签。
  \pretocmd{\@endpart}{\label{\GeometryPartLabel}}{}%
    {\PackageError{geometry}{无法为编题页写入标签}{请检查当前文档类的 \string\@endpart 定义。}}%
  \part{#1}%
  \endgroup
  \thispagestyle{empty}%
  \null
  \clearpage
}
% 各编的“本编导读”是无编号章，在总目录和 PDF 书签中与所属编的各章并列。
% addcontentsline 同时建立目录项和书签，不再另加同名 pdfbookmark。
\NewDocumentCommand{\BookPartGuide}{O{本编导读}}{%
  \chapter*{#1}%
  \addcontentsline{toc}{chapter}{#1}%
}

% “附录”单独排无编号编题页，使用与正文编名相同的字号、字体和颜色。
% 题页从奇数页开始，其后保留一张空白页，再进入第一个附录章。
% startatroot 使根级书签与书后顶层节点分离；phantomsection 须先于
% addcontentsline，目录中的“附录”行才能获得独立锚点而非沿用前一节的
% section*.N。根级“附录”书签连同其子章由 addcontentsline 的 part 行自动
% 建立，不得再手写 \pdfbookmark（实测会产生重复的“附录”节点）。
\NewDocumentCommand{\BookAppendixPart}{m}{%
  \cleardoublepage
  \begingroup
  \def\GeometryPartLabel{#1}%
  % 在 part* 的题页尚未翻走时建立唯一的目录与书签目标。
  \pretocmd{\@endpart}{%
    \bookmarksetup{startatroot}%
    \phantomsection
    \addcontentsline{toc}{part}{附录}%
    \label{\GeometryPartLabel}%
  }{}{\PackageError{geometry}{无法为附录编题页写入目录与标签}{请检查当前文档类的 \string\@endpart 定义。}}%
  \part*{附录}%
  \endgroup
  \thispagestyle{empty}%
  \null
  \clearpage
}
\makeatother

% 列表项的括号与编号整体承担链接，避免只让中间的编号可点击。
\NewDocumentCommand{\itemref}{m}{\hyperref[#1]{(\ref*{#1})}}

% 章、节、小节的引用：整个“第 X 章／第 X 节／第 X 小节”都是跳转链接。
% 链接必须走 hyperref 的标签解析机制（目标名 chapter.N／section.N.M）；
% 早先用 \hyperlink{裸标签名} 指向不存在的目标，链接全部失效。
%
% 章号用汉字（“第五章”）且必须取“被引用那一章”的编号；\thechapter 是当前
% 计数器，在编导读、前言或跨章引用时会给出错误章号（表现为“第零章”）。
% hyperref 下 \ref 受保护、\getrefnumber 在 \edef 中又失真，故改由 zref 记录：
% 各章在 \chapter 之后调用 \BookChapterNumber，把汉字号存为 zref 属性；
% \zref@extract 可展开，能在 \edef 中安全取出。
\usepackage[user,hyperref]{zref}
\makeatletter
\zref@newprop{geometrychapternum}{}
% 在 \chapter 之后调用，参数即该章标签；把本章汉字号存为 zref 属性
\NewDocumentCommand{\BookChapterNumber}{m}{%
  \zref@setcurrent{geometrychapternum}{\chinese{chapter}}%
  \zref@labelbyprops{zch@#1}{geometrychapternum}%
}
% 取出某章标签对应的汉字号（\zref@extract 可展开，能在 \edef 中安全使用）
\NewDocumentCommand{\GeometryChineseChapterOf}{m}{%
  \zref@extract{zch@#1}{geometrychapternum}}
\makeatother
% \ref* 的分组会遮住编号末端的西文字符类别；编号前的间距可由 xeCJK
% 自动产生，编号后的同类间距须在完整链接内部显式补回。
\NewDocumentCommand{\secref}{m}{%
  \hyperref[#1]{第\ref*{#1}\CJKecglue 节}}
\NewDocumentCommand{\subsecref}{m}{%
  \hyperref[#1]{第\ref*{#1}\CJKecglue 小节}}
\NewDocumentCommand{\chapref}{m}{%
  \begingroup
  \xeCJKsetup{CJKecglue={}}%
  \edef\GeometryChapterDisp{第\GeometryChineseChapterOf{#1}章}%
  \hyperref[#1]{\GeometryChapterDisp}%
  \endgroup}

% 算法采用带编号的步骤，并明确标出输入与输出。
\floatname{algorithm}{算法}
\renewcommand{\listalgorithmname}{算法目录}
\algrenewcommand\algorithmicrequire{\textbf{输入：}}
\algrenewcommand\algorithmicensure{\textbf{输出：}}
\numberwithin{algorithm}{chapter}
% 完整中文编号进入题注、各自目录和引用；附录自动使用字母章号。
% 内部 \nobreakspace 与 cleveref 默认的定理引用一致；引用整串置于同一链接，
% 链接后的中文间距由下方 \cref 包装补齐。
\renewcommand{\thealgorithm}{算法\nobreakspace\thechapter-\arabic{algorithm}}
\renewcommand{\thefigure}{图\nobreakspace\thechapter-\arabic{figure}}
\renewcommand{\thetable}{表\nobreakspace\thechapter-\arabic{table}}
\renewcommand{\figurename}{}
\renewcommand{\tablename}{}
\makeatletter
\renewcommand{\fnum@figure}{\thefigure}
\renewcommand{\fnum@table}{\thetable}
\renewcommand{\fnum@algorithm}{\thealgorithm}
\makeatother
\NewDocumentEnvironment{algorithmsteps}{mm}
  {\begin{algorithmic}[1]\Require #1\Ensure #2}
  {\end{algorithmic}}

\newtcolorbox{learninggoals}{
  enhanced,breakable,colback=GeometryBlueLight,colframe=GeometryBlue,
  boxrule=.45pt,arc=1mm,left=1.2mm,right=1.2mm,top=1mm,bottom=1mm,
  title=本章目标,fonttitle=\bfseries
}
\newtcolorbox{partguide}{
  enhanced,breakable,colback=GeometryGrayLight,colframe=GeometryGray,
  boxrule=.35pt,arc=1mm,left=1.2mm,right=1.2mm,top=1mm,bottom=1mm,
  title=本编导读,fonttitle=\bfseries
}
\newtcolorbox{chaptersummary}{
  enhanced,breakable,colback=GeometryBlueLight,colframe=GeometryBlue,
  boxrule=.45pt,arc=1mm,left=1.2mm,right=1.2mm,top=1mm,bottom=1mm,
  title=本章小结,fonttitle=\bfseries,
  % 与阅读提示框一致：仅本框内容为多段散文，须恢复首行缩进。
  before upper={\parindent=2em\parskip=0pt}
}
% 仅用于尚待高阶数学复核的核心断言；成稿前必须处理或移除。
\newtcolorbox{coreverification}{
  enhanced,breakable,colback=yellow!10,colframe=orange!85!black,
  boxrule=.55pt,arc=1mm,left=1.2mm,right=1.2mm,top=1mm,bottom=1mm,
  title=待数学复核,fonttitle=\bfseries
}
% 兼容旧稿的占位命令：只输出普通文字，不另加方框。
\NewDocumentCommand{\GeometryPlaceholder}{m}{待撰写}
\NewDocumentEnvironment{chapterabstract}{}
  {\small\list{}{\rightmargin\leftmargin
    \listparindent=2\ccwd\itemindent=\listparindent}%
    \item\relax\color{GeometryGray}}
  {\endlist}
% 引诗与短引文统一居中，中文使用定理类环境的楷体；章摘要已在上方独立定义，
% 不继承这里的字形与对齐方式。
\RenewDocumentEnvironment{quote}{}
  {\list{}{\rightmargin\leftmargin}\item\relax\centering\kaishu}
  {\endlist}
% 引诗环境保留为语义清楚的兼容入口，版式与 quote 完全一致。
\NewDocumentEnvironment{poemquote}{}
  {\begin{quote}}
  {\end{quote}}
% 引诗末尾的作者或出处在引诗区域内右对齐。
\NewDocumentCommand{\quoteattribution}{m}{%
  \par\smallskip
  {\raggedleft #1\par}%
}

\makeindex[name=chinese,title=中文索引,columns=2,intoc,options=-q]
\makeindex[name=foreign,title=外文索引,columns=2,intoc,options=-q]
\makeindex[name=symbols,title=符号索引,columns=2,intoc,options=-q]
% 索引采用稍小字号与自然栏底，字母组间仅留少量空白；末页平衡不得越过正文高度。
\indexsetup{othercode={\small\raggedcolumns\maxbalancingoverflow=0pt
  \renewcommand{\indexspace}{\par\addvspace{0.2\baselineskip}}}}

% 有条目时读入 makeindex 生成的真实索引；尚无条目时仍保留书后结构与目录入口。
% 一键脚本会在 .idx 为空时删除旧 .ind，故这里不会误排上一次构建的陈旧索引。
\NewDocumentCommand{\PrintBookIndex}{m m}{%
  \IfFileExists{#1.ind}{%
    % 双栏输出可能丢失页眉标记；在索引范围内直接绑定标题并完成分页。
    \cleardoublepage
    \begingroup
      \fancyhead[LE]{\small\color{GeometryGray}\thepage\quad #2}%
      \fancyhead[RO]{\small\color{GeometryGray}#2\quad\thepage}%
      \printindex[#1]%
      \clearpage
    \endgroup
  }{%
    \chapter*{#2}%
    \addcontentsline{toc}{chapter}{#2}%
    \noindent 待撰写\par
  }%
}

% 所有术语在每卷首次正式定义或引入时都显示完整原文，并登记中外文索引。
% 星号形式与缺失原文均视为错误，防止旧的隐藏原文规则被重新带回源码。
% 定理类正文处于楷体；Fandol 楷体没有粗体字形，直接 \textbf 会退回普通楷体。
% 术语主称因此局部恢复正文的宋体字族，再使用与普通正文一致的粗体。
\NewDocumentCommand{\term}{s O{} m o}{%
  \IfBooleanT{#1}{%
    \PackageError{geometry}{\string\term* 已取消}{请改用 \string\term，并保留末尾的原文注解。}%
  }%
  {\songti\textbf{\textup{#3}}}%
  \IfNoValueTF{#4}{%
    \PackageError{geometry}{术语缺少原文注解}{每个术语首次正式定义或引入时都必须写成 \string\term[排序键]{中文主称}[原文]。}%
  }{%
    \textnormal{(\textup{#4})}%
  }%
  \ifstrempty{#2}{\index[chinese]{#3}}{\index[chinese]{#2@#3}}%
  \IfNoValueF{#4}{\index[foreign]{#4}}%
}
\NewDocumentCommand{\termidx}{O{}m}{%
  \ifstrempty{#1}{\index[chinese]{#2}}{\index[chinese]{#1@#2}}%
}
% 别称与主称一样用 \termidx 登记实际页码，不生成“见某某”的转引条目。
\NewDocumentCommand{\foreignidx}{O{}m}{%
  \ifstrempty{#1}{\index[foreign]{#2}}{\index[foreign]{#1@#2}}%
}
\newcounter{geometrySymbolIndexOrder}
\NewDocumentCommand{\symidx}{m}{%
  \stepcounter{geometrySymbolIndexOrder}%
  \edef\GeometrySymbolSortKey{%
    \number\numexpr10000+\value{geometrySymbolIndexOrder}\relax}%
  \index[symbols]{\GeometrySymbolSortKey@#1}%
}

% 一般强调与术语粗体分离。
\NewDocumentCommand{\keypoint}{m}{\CJKunderdot{#1}}
\NewDocumentCommand{\foreignkeypoint}{m}{\underline{\textup{#1}}}


\crefname{chapter}{章}{章}
\crefname{part}{编}{编}
\crefname{section}{节}{节}
\crefname{subsection}{小节}{小节}
\crefname{theorem}{定理}{定理}
\crefname{lemma}{引理}{引理}
\crefname{proposition}{命题}{命题}
\crefname{corollary}{推论}{推论}
\crefname{claim}{断言}{断言}
\crefname{conjecture}{猜想}{猜想}
\crefname{openproblem}{开放问题}{开放问题}
\crefname{definition}{定义}{定义}
\crefname{example}{例}{例}
\crefname{remark}{注记}{注记}
\crefname{exercise}{习题}{习题}
\crefname{equation}{式}{式}
\crefname{figure}{}{}
\crefname{table}{}{}
\crefname{algorithm}{}{}
\crefformat{figure}{#2#1#3}
\Crefformat{figure}{#2#1#3}
\crefformat{table}{#2#1#3}
\Crefformat{table}{#2#1#3}
\crefformat{algorithm}{#2#1#3}
\Crefformat{algorithm}{#2#1#3}

% cleveref 会在自己的分组中排出“定义 3.1.5”一类完整链接，类别名与编号之间
% 使用 \nobreakspace。分组结束后，xeCJK 无法再从末端识别出西文数字，
% 因而不会在紧随其后的汉字前插入 \CJKecglue。这里只在下一个源码记号确为
% CJK 正文字时补回同一种间距；空格、控制序列以及中英文标点均不处理，
% 从而不在“定义3.1.5，”或“定义3.1.5.”前制造额外空白。
\NewCommandCopy{\GeometryOriginalCref}{\cref}
\NewCommandCopy{\GeometryOriginalCrefCapital}{\Cref}
\ExplSyntaxOn
\cs_new_protected:Npn \geometry_cref_maybe_cjk_glue:
  { \peek_after:Nw \geometry_cref_maybe_cjk_glue_aux: }
\cs_new_protected:Npn \geometry_cref_maybe_cjk_glue_aux:
  {
    \token_if_space:NF \l_peek_token
      {
        \token_if_cs:NF \l_peek_token
          {
            \xeCJK_if_same_class:NNT \l_peek_token 中
              { \CJKecglue }
          }
      }
  }
\RenewDocumentCommand{\cref}{s m}
  {
    \IfBooleanTF{#1}
      { \GeometryOriginalCref*{#2} }
      { \GeometryOriginalCref{#2} }
    \geometry_cref_maybe_cjk_glue:
  }
\RenewDocumentCommand{\Cref}{s m}
  {
    \IfBooleanTF{#1}
      { \GeometryOriginalCrefCapital*{#2} }
      { \GeometryOriginalCrefCapital{#2} }
    \geometry_cref_maybe_cjk_glue:
  }
% 附录标签沿用字母编号；名称与编号整体进入同一链接，并与 \cref
% 采用相同的链接后中西文间距规则。
\NewDocumentCommand{\appref}{m}
  {
    \hyperref[#1]{附录\CJKecglue\ref*{#1}}
    \geometry_cref_maybe_cjk_glue:
  }
\ExplSyntaxOff

\newcommand{\MakeBookTitle}[1][封面]{%
  % 书名页在 PDF 侧栏中作为顶级入口；卷级文件可传入更具辨识性的书签名。
  \bookmarksetup{startatroot}%
  \pdfbookmark[0]{#1}{GeometryCover}%
  \begin{titlepage}
    \centering
    \vspace*{2.2cm}
    {\fontsize{36}{46}\selectfont\sffamily 几何学\par}
    \vspace{0.7cm}
    {\Large\sffamily\color{GeometryGray}\BookVolumeName\par}
    \vspace{1.0cm}
    {\fontsize{24}{32}\selectfont\bfseries\color{GeometryBlue}\BookTitle\par}
    \vfill
    {\large \BookAuthor\ 编写\par}
    \vspace{0.5cm}
    {\normalsize\BookDraftDate\par}
    \vspace{0.7cm}
    {\large\sffamily\bfseries\color{GeometryGray}\BookStatus\par}
    \vspace*{1.8cm}
  \end{titlepage}
}

% 目录也应有独立的顶级 PDF 书签。使用公共命令避免五卷分别遗漏。
\newcommand{\MakeBookContents}{%
  \cleardoublepage
  \bookmarksetup{startatroot}%
  \pdfbookmark[0]{目录}{GeometryContents}%
  \tableofcontents
}

% 图、表、算法目录不进入总目录，也不建立 PDF 书签；各自页面使用原有章级标题。
% 只另起一页而不补单双页。
\makeatletter
\newcommand{\GeometryListStart}{%
  \clearpage
}
% 各列表命令自带 \chapter*，局部消除其中的再次换页，保留原有标题与列表排版。
\newcommand{\GeometryListOnCurrentPage}[1]{%
  \begingroup
    \let\clearpage\relax
    \let\cleardoublepage\relax
    #1%
  \endgroup
}
\newcommand{\MakeBookListOfFigures}{%
  \GeometryListStart
  \GeometryListOnCurrentPage{\listoffigures}%
}
\newcommand{\MakeBookListOfTables}{%
  \begingroup
    \renewcommand{\listtablename}{表目录}%
    \GeometryListStart
    \GeometryListOnCurrentPage{\listoftables}%
  \endgroup
}
\newcommand{\MakeBookListOfAlgorithms}{%
  \begingroup
    \GeometryListStart
    % float.sty 在 listof 内重设编号栏；在标题回调中再加宽它。
    \let\GeometryOriginalFloatListHead\float@listhead
    \renewcommand{\float@listhead}[1]{%
      \renewcommand*\l@algorithm{\@dottedtocline{1}{1.5em}{5em}}%
      \GeometryOriginalFloatListHead{##1}%
    }%
    \GeometryListOnCurrentPage{\listofalgorithms}%
  \endgroup
}
\makeatother

% 无编号章（序言、目录、编导读、后记、参考文献、索引等）同样写入页眉标记，
% 使后续偶数页页眉显示该单元标题。置于 hyperref 之后，保证保留 hyperref 的章锚点。
\makeatletter
\let\Geometry@orig@schapter\@schapter
\def\@schapter#1{%
  \markboth{#1}{#1}%
  \Geometry@orig@schapter{#1}%
}
% 目录按层级着色：编为黑色、小节为深灰；章与节保持默认链接蓝。
\pretocmd{\l@part}{\hypersetup{linkcolor=black}}{}{%
  \PackageWarning{geometry}{无法为目录编条目着色}}
\pretocmd{\l@chapter}{\hypersetup{linkcolor=GeometryBlue}}{}{%
  \PackageWarning{geometry}{无法为目录章条目着色}}
\pretocmd{\l@section}{\hypersetup{linkcolor=GeometryBlue}}{}{%
  \PackageWarning{geometry}{无法为目录节条目着色}}
\pretocmd{\l@subsection}{\hypersetup{linkcolor=GeometryGray}}{}{%
  \PackageWarning{geometry}{无法为目录小节条目着色}}
% 编的目录条目与 PDF 书签须用纯文本“第X编”。ctex 的 \CTEX@part@tocline
% 取 \CTEXthepart，而该宏带上编题页的颜色与字体设置，会把格式代码写进目录，
% 并使书签出现 “=GeometryGray第一编” 之类乱码。此覆盖须置于 \ctexset 之后，
% 否则会被 ctex 重新定义；\CTEXifname 是 ctex 的内部开关，此处沿用。
\def\CTEX@part@tocline#1#2{%
  \CTEXifname
    {第\chinese{part}编\hspace{1em}}%
    {}%
  #2}
\makeatother

\hypersetup{pdftitle={几何学·\BookVolumeName：\BookTitle},pdfauthor={\BookAuthor},pdfsubject={第一卷写作稿：二十章正文及附录 A—E 初稿；前言、后记及前两编导言待撰写}}

\begin{document}
\frontmatter
% !TeX root = ../Book1.tex
\MakeBookTitle[\BookTitle]
\clearpage
\thispagestyle{empty}
\null
\clearpage
% !TeX root = ../Book1.tex
\chapter*{序言}
\addcontentsline{toc}{chapter}{序言}
同一个圆，可以由距离方程描述，也可以看成一条闭合曲线、一个拓扑空间或一个群。每一种描述都保留某些性质，并提出不同的问题：哪些变换保持距离？沿圆周走一圈留下什么信息？局部参数怎样拼成整体？《几何学》从这些具体问题出发，逐步建立研究几何结构、空间与映射的方法。本卷讨论古典几何与光滑流形，是全书五卷中的第一卷。

\ProseParagraphBreak

古典几何中的坐标、变换和分类，为后来的抽象构造提供了可以反复检验的例子。圆锥曲线的标准形取决于允许仿射变换、等距变换还是射影变换；曲面的局部度量可以写在一张坐标图上，曲面整体的拓扑却要通过粘合来认识。进入流形以后，这些问题仍在：换坐标时什么量保持意义，局部给出的数据何时能相容，怎样由相容的数据构造整体对象。本卷的安排围绕这些问题展开。

\ProseParagraphBreak

本书面向正在首次或第二次系统学习几何的数学系本科生和研究生。开始阅读需要有限维线性代数和微积分的基本训练：能够使用基、线性映射、对偶、行列式与内积，理解导数、积分、极限和连续性。曲线曲面与流形部分还要用到多元微积分；有关微分方程、拓扑和群作用的准备，在实际使用处交代。已经学过某个定理却忘记了条件的读者，可以先看该处的回顾和引用，再决定是否回读其证明。

\ProseParagraphBreak

前两编先讨论几何对象怎样分类、空间怎样构造。第一编从仿射与 Euclid 结构进入射影几何和凸性，把坐标化简与不变量联系起来；第二编建立拓扑、基本群和覆盖，继而通过切割与粘合研究紧曲面。这里的分类同时包含两项工作：证明给出的模型能够容纳所有对象，并找出区分不同模型的依据。这样的区分也将用于后面的曲线基本定理、平坦结构和齐性空间。

\ProseParagraphBreak

第三编研究曲线与曲面的微分几何。曲率、基本形式和联络逐步把局部计算联系起来，曲面的整体公式则把曲率积分接到拓扑。球面与双曲模型让读者比较不同度量下的距离、测地线和群作用。第四编将这些构造放到一般光滑流形上，依次处理切空间、向量丛、流、微分形式与联络。第五编研究连续对称性，借助 Lie 群、齐性空间和主丛，把换标架、群作用与平行移动纳入共同的几何描述；最后两章继续讨论紧群的结构与表示。

\ProseParagraphBreak

读者可以按自己的目的选择起点。希望从熟悉的图形逐步学习者，可以依次阅读前三编，再进入流形。已经具备线性代数和多元微积分基础、希望直接学习光滑流形者，可以从\chapref{g1:ch:11}开始；所需的开集、连续映射、同胚、子空间与积空间，可在\chapref{g1:ch:04}查阅。逆函数、隐函数与常秩定理在\chapref{g1:ch:11}重新建立，单位分解的光滑构造在\chapref{g1:ch:12}展开。读到平坦联络与单值时，再补读\chapref{g1:ch:05}的基本群和覆盖理论。古典分类与曲面分类可以另行阅读。

\ProseParagraphBreak

各编导读说明本编的问题和所需基础，各章的摘要与学习目标便于确认当前任务。正文中的例子既用于理解定义，也用于检验条件；同一个环面、射影空间或球面再次出现时，可以比较这一次新增了哪些结构。习题配有解答，阅读时宜先独立计算或构造，再用解答检查关键步骤。对于较长证明，先辨认待构造的对象及主要障碍，再跟随局部计算回到完整结论，通常比逐式记忆更有帮助。

\ProseParagraphBreak

章、节及习题中的星号标明可按学习目的选读的内容，附录提供补充专题。附录 A 讨论公理化几何与几何史的方法；附录 B、E 延伸离散、凸性与反射群问题；附录 C 讨论计算与图示，附录 D 继续研究几何结构。它们可以在相应正文之后选读。各章的阅读说明给出适合对照学习的教材，具体引文则标明所用结果或证明的出处；专题专著和原始论文用于继续追问某一问题。卷末参考文献与索引供检索，阅读时以当前问题所需的资料为先。

\ProseParagraphBreak

后续各卷分别讨论代数拓扑与微分拓扑、黎曼几何与几何分析、复几何与辛几何、代数几何。本卷的流形、丛、微分形式和联络会在其中反复使用，所需的拓扑、分析或代数工具则随问题继续建立。当前二十章、附录及书前书后文字均为初稿，尚未经过正式审核。

\MakeBookContents
% 图、表目录收录已写正文中的实际对象。
\MakeBookListOfFigures
\MakeBookListOfTables
% \MakeBookListOfAlgorithms


\mainmatter
% !TeX root = ../Book1.tex
% 纲要标题：# 第一编　几何结构、变换与分类
% 纲要位置：几何学计划纲要.md，第 115 行。
% 正文状态：未撰写；本文件目前是框架。

% 纲要内容开始
% 
% 纲要内容结束

\BookPart{几何结构、变换与分类}{g1:part:01}
\BookPartGuide
待撰写

% !TeX root = ../../Book1.tex
% 纲要标题：## 第1章　仿射、Euclid 结构与变换群
% 纲要位置：几何学计划纲要.md，第 117 行。
% 正文状态：初稿；尚未正式审核。

% 纲要内容开始
% 
% > 本章目标与证明范围：以“选择什么结构、允许什么变换、分类保留什么不变量”为统一问题。完整处理仿射构造、实二次型和 Euclid 等距变换；群作用在实际分类中建立，不另作一章名词准备。
% 
% 纲要内容结束

\chapter{仿射、Euclid 结构与变换群}
\label{g1:ch:01}
\BookChapterNumber{g1:ch:01}

\begin{chapterabstract}
	几何分类取决于指定的结构与允许的变换。本章从点与位移的区别建立仿射空间，说明坐标、重心与子空间相交的内在含义；加入内积后，证明正交投影、自伴谱分解、实对称型的惯性分类，并用内积矩阵的行列式计算体积。随后从距离保持推出 Euclid 等距变换的仿射形式，分解其正交部分与不可消去的平移，证明反射生成并给出低维分类。最后以轨道和稳定子群统一这些分类，比较 Euclid 、相似与仿射等价，考察退化参数处不变量的变化。
\end{chapterabstract}

\begin{learninggoals}
	\begin{itemize}
		\item 区分点与向量，计算仿射组合、标架坐标及仿射子空间的相交与包络。
		\item 说明正交投影和谱分解的构造，区分正交相似、一般相似与合同，求实对称型的惯性，并由内积矩阵计算体积。
		\item 从距离条件求出等距变换的正交部分和平移部分，判断不动点、最小位移及低维运动类型。
		\item 用轨道、稳定子群和完整分类参数比较不同等价关系，并判断连续变化何时跨越退化对象。
	\end{itemize}
\end{learninggoals}

把一个三角形在桌面上平移或旋转，边长没有改变；把它画到另一个比例的地图上，边长改变而比例没有改变；把图形沿一个方向拉长，连角度也可能改变。但只要三个顶点始终不共线，它们仍然确定同一种仿射标架，线段的中点也仍被送到中点。这三个操作分别提醒我们：长度、长度比和仿射组合是不同层次的信息。说两个图形“相同”之前，必须说明我们希望保存其中哪些信息。

\ProseParagraphBreak

本章将这一问题落实到可计算的构造中。先不选长度，只研究点之间的位移；再在方向空间上指定内积，使长度、正交和体积获得意义。二次型提供了观察“同一数据在不同坐标中怎样表达”的具体机会，而 Euclid 运动则让我们看到，仅由距离也能恢复仿射结构。旋转中心、螺旋轴或滑移方向会从同一个线性方程中出现，群作用随后把这些计算概括为分类问题。

\ProseParagraphBreak

全文限于有限维实空间。所需基础是向量空间、基、线性映射、矩阵、行列式与秩的基本运算；关键的几何用法和谱定理本身在本章证明。另用有限维空间中闭有界集的紧性，以及连续实函数在非空紧集上取得最大、最小值这一基本分析结论。例如在单位球面上取函数值趋近上确界的点列，紧性给出收敛子列，连续性便使极限点取得该上确界；这一短论证正是后文极值法的依据。

\ProseParagraphBreak

仿射构造可对照 Gallier 的正式第二版\cite[第2.1--2.8, 2.10--2.11节]{Gallier2011Geometric}；内积、等距变换与矩阵标准形可分别参看同书\cite[第6.1--6.4, 8.1--8.2, 8.4--8.7, 12.3--12.4节]{Gallier2011Geometric}。

% !TeX root = ../../Book1.tex
% 纲要标题：### 1.1 仿射空间的内在构造
% 纲要位置：几何学计划纲要.md，第 121 行。
% 正文状态：初稿；尚未正式审核。

% 纲要内容开始
% 
% * 平移作用、点差与仿射组合
% * 仿射标架、坐标变换与重心坐标
% * 子空间的相交、平行与维数公式
% 
% 纲要内容结束

\section{仿射空间的内在构造}
\label{g1:sec:0101}

在一张地图上改变坐标原点，不会改变两个地点之间的位移，却会同时改变两个地点的坐标。因此，“地点”与“位移”应当是两类对象。向量空间天然带有零向量，仿射空间则把原点的选择从几何结构中分离出来。这样做以后，中点、直线和平行关系仍可定义，而两个点直接相加通常没有意义。

\subsection{平移作用、点差与仿射组合}
\label{g1:subsec:0101:01}

\begin{definition}[仿射空间]\label{g1:def:affine-space}
	设 \(V\) 是有限维实向量空间。给非空集合 \(A\) 指定运算 \(A\times V\rightarrow A\)，记为 \((p,v)\mapsto p+v\)。若该运算满足
	\[
		p+0=p,\qquad (p+u)+v=p+(u+v),
	\]
	且对每对 \(p,q\in A\)，恰有一个 \(v\in V\) 满足 \(p+v=q\)，则称这个带有运算的集合为以 \(V\) 为方向空间的\term[fangshekongjian]{仿射空间}[affine space]。满足上述等式的唯一向量称为\term[diancha]{点差}[difference vector]，记为 \(q-p\)；定义 \(\dim A=\dim V\)。
\end{definition}

映射 \(\tau_v:p\mapsto p+v\) 称为\term[pingyi]{平移}[translation]。上述条件说，每一个点都能通过唯一的平移向量移到另一个点。特别地，\(q=p+(q-p)\)，而点差满足
\[
	(r-q)+(q-p)=r-p,\qquad p-q=-(q-p).
\]
例如，第一个等式由 \(p+((q-p)+(r-q))=q+(r-q)=r\) 及位移的唯一性得到。这里的加法发生在方向空间中，不能把它误读成三个点的普通加减。

\begin{example}[方程的解集]\label{g1:exm:affine-solutions}
	设 \(L:V\rightarrow W\) 为实线性映射，方程 \(Lx=b\) 有一个解 \(p\)。其解集恰为 \(p+\ker L\)：若 \(x\) 也是解，则 \(L(x-p)=0\)；若 \(u\in\ker L\)，则 \(L(p+u)=b\)。它是以 \(\ker L\) 为方向空间的仿射空间。若 \(b\ne0\)，这个解集不含零向量，不能当作 \(V\) 的向量子空间。例如 \(x+y+z=1\) 是仿射平面，两个解之和一般已不是解。
\end{example}

选定 \(o\in A\) 后，\(p\mapsto p-o\) 将 \(A\) 与 \(V\) 一一对应。若改选 \(o'=o+a\)，则同一点的向量表示变为 \(p-o'=(p-o)-a\)。这说明可以借助原点计算，但要检查计算结果是否依赖它。特别值得保留的运算是系数和为 \(1\) 的加权平均。

\begin{proposition}[仿射组合]\label{g1:prop:affine-combination}
	设 \(A\) 是实仿射空间，\(p_0,\ldots,p_m\in A\)， \(\lambda_0,\ldots,\lambda_m\in\mathbb R\) 且 \(\sum_i\lambda_i=1\)。则
	\[
		o+\sum_{i=0}^m\lambda_i(p_i-o)
	\]
	与 \(o\in A\) 的选择无关，称为这些点的\term[fangshezuhe]{仿射组合}[affine combination]，简记为 \(\sum_i\lambda_i p_i\)。
\end{proposition}
\begin{proof}
	以 \(o'=o+a\) 代替 \(o\)，新表达式为
	\[
		o+a+\sum_i\lambda_i\bigl((p_i-o)-a\bigr)
		=o+\sum_i\lambda_i(p_i-o)+(1-\sum_i\lambda_i)a.
	\]
	最后一项为零，因而得到同一点。
\end{proof}

系数和条件正是消去原点影响的条件。特别地，当 \(p\ne q\) 时，\((1-t)p+tq=p+t(q-p)\) 描述穿过 \(p,q\) 的直线；当 \(0\le t\le1\) 时描述线段，当 \(t=1/2\) 时给出中点。若 \(p=q\)，这些组合都退化为同一点。若系数和为零，则 \(\sum_i\lambda_i(p_i-o)\) 也与原点无关，但它给出的是向量而不是点。这一差别解释了为什么 \(q-p\) 合法，而 \(p+q\) 不能作为不依赖原点的点运算。

\begin{definition}[仿射映射]\label{g1:def:affine-map}
	设 \(A,B\) 分别以实向量空间 \(V,W\) 为方向空间。对于映射 \(f:A\rightarrow B\)，若存在实线性映射 \(L:V\rightarrow W\)，使
	\[
		f(p+v)=f(p)+Lv\qquad(p\in A,\ v\in V),
	\]
	则称这个映射为\term[fangsheyingshe]{仿射映射}[affine map]，称 \(L\) 为 \(f\) 的\term[xianxingbufen]{线性部分}[linear part]。
\end{definition}

线性部分唯一，因为 \(Lv=f(p+v)-f(p)\)。选定原点后，仿射映射就是 \(x\mapsto Lx+b\)。如果 \(L\) 可逆，逆映射是 \(y\mapsto L^{-1}(y-b)\)，此时称 \(f\) 为\term[fangshetonggou]{仿射同构}[affine isomorphism]；反过来，若仿射映射是双射，则其线性部分也是双射。

\begin{proposition}\label{g1:prop:affine-preservation}
	实仿射空间间的映射 \(f:A\rightarrow B\) 是仿射映射，当且仅当它保持一切有限仿射组合。
\end{proposition}
\begin{proof}
	若 \(f\) 的线性部分为 \(L\)，在组合的计算中以 \(o\) 与 \(f(o)\) 为原点，便有
	\[
		f\left(o+\sum_i\lambda_i(p_i-o)\right)
		=f(o)+\sum_i\lambda_i\bigl(f(p_i)-f(o)\bigr).
	\]
	这正是保持仿射组合的等式。

	反过来，固定 \(o\)，定义 \(L(v)=f(o+v)-f(o)\)。从组合 \((1-t)o+t(o+v)\) 得到 \(L(tv)=tL(v)\)，其中 \(t\) 为任意实数。再对三点 \(o+u,o+v,o\) 使用系数 \(1,1,-1\)，得到 \(L(u+v)=L(u)+L(v)\)。故 \(L\) 线性，且 \(f(p+v)=f(o)+L(p-o+v)=f(p)+Lv\)。
\end{proof}

\subsection{仿射标架、坐标变换与重心坐标}
\label{g1:subsec:0101:02}

要给一个点写出数值坐标，需要选原点与方向基；要用若干几何点描述它，则可以用一个有序点组同时指定原点与方向基。

\begin{definition}[标架与仿射无关]\label{g1:def:affine-frame}
	在 \(n\) 维实仿射空间 \(A\) 中，若 \(p_1-p_0,\ldots,p_m-p_0\) 线性无关，就称 \(p_0,\ldots,p_m\) \term[fangshewuguan]{仿射无关}[affinely independent]。当 \(m=n\) 时，有序点组 \((p_0,\ldots,p_n)\) 称为\term[fangshebiaojia]{仿射标架}[affine frame]。
\end{definition}

仿射无关不依赖于把哪个点放在首位。事实上，它等价于：若 \(\sum_i c_i=0\) 且 \(\sum_i c_i(p_i-o)=0\)，则所有 \(c_i\) 为零。取 \(o=p_0\) 就得到原定义；这个对称表述不再区别各点。

\begin{proposition}[坐标及其变换]\label{g1:prop:affine-coordinates}
	设 \(A\) 为 \(n\) 维实仿射空间。给定仿射标架 \((p_0,\ldots,p_n)\)，每个 \(p\in A\) 唯一写成
	\[
		p=p_0+\sum_{i=1}^n x_i(p_i-p_0)
		=\sum_{i=0}^n\lambda_i p_i,\qquad
		\lambda_0=1-\sum_{i=1}^n x_i,\quad\lambda_i=x_i\ (i\ge1).
	\]
	称 \((x_1,\ldots,x_n)\) 为\term[fangshezuobiao]{仿射坐标}[affine coordinates]，\((\lambda_0,\ldots,\lambda_n)\) 为相对于此标架的\term[zhongxinzuobiao]{重心坐标}[barycentric coordinates]。若新原点相对于旧标架的坐标为 \(a\)，新方向基在旧方向基中的列矩阵为 \(S\)，则新旧坐标满足
	\[
		x=a+Sx',\qquad x'=S^{-1}(x-a).
	\]
\end{proposition}
\begin{proof}
	基展开给出 \(p-p_0\) 的唯一系数 \(x_i\)，而系数和为 \(1\) 唯一确定 \(\lambda_0\)。设旧方向基的组合记为 \(Ex=\sum_i x_i(p_i-p_0)\)，则新原点为 \(p_0+Ea\)，新方向基对应线性映射 \(ES\)。于是 \(p=p_0+Ex=p_0+Ea+ESx'\)。由于 \(E\) 单射，得到 \(x=a+Sx'\)；新方向基仍是基，所以 \(S\) 可逆。
\end{proof}

\begin{example}[三角形中的位置]\label{g1:exm:barycentric-triangle}
	取 \(p_0=(0,0),p_1=(3,0),p_2=(0,2)\)。点 \(p=(1,1)\) 的重心坐标为 \((1/6,1/3,1/2)\)，因为后两个系数分别由 \(3\lambda_1=1\)、 \(2\lambda_2=1\) 决定，再由系数和求出第一个。三个系数全非负表示点处在闭三角形内，全为正表示点处在内部；在全非负的条件下，某个系数为零表示点处在对应的边上。例如 \(q=(2,1)\) 的坐标为 \((-1/6,2/3,1/2)\)，它越过了与 \(p_0\) 相对的边。
\end{example}

\begin{figure}[htbp]
	\centering
	\input{Figures/Chapter01-Barycentric}
	\caption[重心坐标与三角形]{重心坐标与三角形。直线 \(p_1p_2\) 上 \(\lambda_0=0\)，点 \(p\) 与顶点 \(p_0\) 在同侧，点 \(q\) 在另一侧。}
	\label{g1:fig:barycentric}
\end{figure}

图中的符号判断可以直接计算，而无须凭图形猜测：此例 \(\lambda_0=1-x/3-y/2\)，所以直线两侧的符号由一次函数决定。一般情形下，系数全非负的仿射组合构成这组仿射无关点张成的\term[dan-chun-xing]{单纯形}[simplex]。这种位置关系在任何仿射坐标中都成立，因为仿射同构及其逆映射都保持这些组合，而不要求保持长度或角度。

\begin{proposition}[由标架确定映射]\label{g1:prop:frame-map}
	设 \((p_0,\ldots,p_n)\) 为实仿射空间 \(A\) 的标架，\(B\) 为另一个实仿射空间。任给 \(q_0,\ldots,q_n\in B\)，存在唯一仿射映射 \(f:A\rightarrow B\) 满足 \(f(p_i)=q_i\)。若 \(B\) 也为 \(n\) 维，则 \(f\) 是同构当且仅当这些 \(q_i\) 构成标架。
\end{proposition}
\begin{proof}
	定义线性映射 \(L\) 将基向量 \(p_i-p_0\) 送到 \(q_i-q_0\)，再置 \(f(p)=q_0+L(p-p_0)\)。它是满足指定条件的仿射映射。任何满足条件的映射都必须有这一线性部分及同一 \(f(p_0)\)，故唯一。两空间同维时，\(L\) 可逆恰好要求其基向量的像构成基。
\end{proof}

\subsection{子空间的相交、平行与维数公式}
\label{g1:subsec:0101:03}

在向量空间中，子空间一定交于零向量；在仿射空间中，两条直线可能平行，也可能既不相交又不平行。为了区别这些情形，除了方向空间，还要记录两个空间之间的位移。

\begin{definition}[仿射子空间]\label{g1:def:affine-subspace}
	设 \(A\) 为以 \(V\) 为方向空间的实仿射空间，\(p\in A\)，\(U\subseteq V\) 为向量子空间。集合 \(F=p+U\) 称为\term[fangshezikongjian]{仿射子空间}[affine subspace]，其方向空间为 \(U\)，维数为 \(\dim U\)。若两个仿射子空间的方向空间之一包含于另一个，本章就称它们\term[pingxing]{平行}[parallel]；对同维子空间，这等价于两个方向空间相等。允许重合的空间互相平行。
\end{definition}

不同文献对异维空间的“平行”用语不完全一致，后文的判断将直接写出方向空间的包含或相等。对于 \(F=p+U\)，选其中任何 \(r\) 为基点仍有 \(F=r+U\)，因为 \(r-p\in U\)。同时 \(F-F=U\)，所以方向空间由点集本身唯一决定。

\begin{proposition}[相交与仿射包络]\label{g1:prop:affine-intersection}
	在有限维实仿射空间 \(A\) 中，设 \(F=p+U\)、\(G=q+W\) 为非空仿射子空间。则
	\[
		F\cap G\ne\varnothing\quad\Longleftrightarrow\quad q-p\in U+W.
	\]
	若 \(r\in F\cap G\)，则 \(F\cap G=r+(U\cap W)\)。包含 \(F\cup G\) 的最小仿射子空间，称为其\term[fangshebaoluo]{仿射包络}[affine hull]，为
	\begin{equation}\label{g1:eq:affine-join}
		\operatorname{aff}(F\cup G)=p+\bigl(U+W+\mathbb R(q-p)\bigr).
	\end{equation}
	特别地，令 \(d=\dim(U\cap W)\)，则
	\[
		\dim\operatorname{aff}(F\cup G)
		=\begin{cases}
			\dim F+\dim G-d,&F\cap G\ne\varnothing,\\
			\dim F+\dim G-d+1,&F\cap G=\varnothing.
		\end{cases}
	\]
	在第一种情形中，\(d=\dim(F\cap G)\)。
\end{proposition}
\begin{proof}
	交点的方程为 \(p+u=q+w\)，也就是 \(q-p=u-w\)。这个方程有解，当且仅当 \(q-p\in U+W\)。若有交点 \(r\)，则写成 \(F=r+U\)、\(G=r+W\)，一个点属于二者等价于它与 \(r\) 的差同时属于 \(U,W\)，给出交集公式。

	式\eqref{g1:eq:affine-join}右边包含 \(F,G\)。任何包含二者的仿射子空间都含 \(p\)，其方向空间必须包含 \(U,W\) 及 \(q-p\)，因而也包含右边的方向空间。这证明了最小性。

	最后回顾线性空间的维数计算。先取 \(U\cap W\) 的一组基，分别扩充成 \(U,W\) 的基；把交集基只写一次，再合并两个扩充部分，就得到 \(U+W\) 的基。张成性来自两个扩充基；若合并后有线性关系，则 \(U\) 的扩充部分之和落在 \(W\)，因而落在交集，这迫使该部分系数为零，再依次得到其余系数为零。所以 \(\dim(U+W)=\dim U+\dim W-d\)。交点存在时 \(q-p\) 已在这个和中；交点不存在时再加入它恰增加一维。
\end{proof}

\begin{example}[平行线与异面直线]\label{g1:exm:skew-lines}
	在 \(\mathbb R^3\) 中，令
	\[
		F=\{(s,0,0):s\in\mathbb R\},\qquad
		G=\{(0,t,1):t\in\mathbb R\}.
	\]
	它们的方向分别为 \(\mathbb Re_1,\mathbb Re_2\)，基点之差为 \(e_3\)，不在 \(\mathbb Re_1+\mathbb Re_2\) 中，故不相交；方向也不平行。两条这样的直线称为\term[yimianzhixian]{异面直线}[skew lines]。其仿射包络维数为 \(1+1-0+1=3\)。若改成 \(G'=\{(t,0,1):t\in\mathbb R\}\)，方向交集为一维，包络维数变为 \(2\)，这才是两条不同的平行线。
\end{example}

对任意非空点集 \(S\subseteq A\)，选定 \(p_0\in S\) 后，同一最小性论证给出
\[
	\operatorname{aff}(S)=p_0+\operatorname{span}\{p-p_0:p\in S\}.
\]
这里的 \(\operatorname{span}\) 只取有限线性组合。展开右边后可见，它也恰好是 \(S\) 中有限多个点的一切仿射组合。因此有限点组 \(p_0,\ldots,p_m\) 的仿射包络维数就是列向量 \(p_1-p_0,\ldots,p_m-p_0\) 的秩，是否仿射无关可以转化为一次矩阵秩计算。

% !TeX root = ../../Book1.tex
% 纲要标题：### 1.2 度量、配对与二次型
% 纲要位置：几何学计划纲要.md，第 127 行。
% 正文状态：初稿；尚未正式审核。

% 纲要内容开始
% 
% * 内积、对偶与正交投影
% * 自伴谱分解与正交标准形
% * 对称双线性型的合同分类与惯性
% * Gram 行列式、体积与退化的检测
% 
% 纲要内容结束

\section{度量、配对与二次型}
\label{g1:sec:0102}

仿射结构能说明三个点是否共线、两条直线是否平行，却不能决定一条线段有多长。把平面沿一个方向拉伸，上一节讨论的仿射关系仍然成立，圆却会变成椭圆。要分辨这样的变化，就需要在方向向量空间上再指定一种结构，用来测量长度和夹角。内积承担这一任务；更一般的对称双线性型则让我们比较正负方向、识别退化，并描述二次方程。

\ProseParagraphBreak

本节始终在实有限维向量空间上讨论。线性映射可以用矩阵表示，但同一个矩阵也可能表示双线性型；二者换基时的公式不同。先弄清对象，再讨论它的标准形，是后面几何分类中避免混淆的关键。

\subsection{内积、对偶与正交投影}
\label{g1:subsec:0102:01}

设 \(V\) 是实向量空间。我们希望长度的平方是一个二次表达式，并且非零向量的长度严格为正。为此，先规定两个向量之间的配对。

\begin{definition}[内积与欧氏结构]
	\label{g1:def:inner-product}
	实向量空间 \(V\) 上的\term[neiji]{内积}[inner product]是映射
	\[
		\langle\,\cdot\,,\,\cdot\,\rangle\colon V\times V\longrightarrow\mathbb R,
	\]
	满足对两个变量分别线性、\(\langle u,v\rangle=\langle v,u\rangle\)，并且对每个非零 \(v\in V\) 有 \(\langle v,v\rangle>0\)。最后一个条件称为\term[zhengdingxing]{正定性}[positive definiteness]。记
	\[
		\lVert v\rVert\coloneqq\sqrt{\langle v,v\rangle},
	\]
	称为内积给出的\term[fanshu]{范数}[norm]。带有内积的实有限维向量空间称为\term[oushixiangliangkongjian]{欧氏向量空间}[Euclidean vector space]。若仿射空间 \(A\) 的方向空间 \(V\) 带有内积，则称 \(A\) 带有\term[oushijiegou]{欧氏结构}[Euclidean structure]，两点的距离定义为
	\[
		d(p,q)\coloneqq\lVert q-p\rVert.
	\]
	\symidx{\(\langle u,v\rangle\)：内积}
	\symidx{\(\lVert v\rVert\)：内积范数}
\end{definition}

这里的正定性有两项直接作用。它保证长度为零的向量只有零向量，也保证长度满足三角不等式。后一项可以从一个简单的非负平方看出。若 \(v\ne0\)，则
\[
	0\leq\left\lVert u-\frac{\langle u,v\rangle}{\lVert v\rVert^2}v\right\rVert^2
	=\lVert u\rVert^2-\frac{\langle u,v\rangle^2}{\lVert v\rVert^2}.
\]
所以 \(\lvert\langle u,v\rangle\rvert\leq\lVert u\rVert\cdot\lVert v\rVert\)，\(v=0\) 时也成立。代入
\[
	\lVert u+v\rVert^2=\lVert u\rVert^2+2\langle u,v\rangle+\lVert v\rVert^2
	\leq(\lVert u\rVert+\lVert v\rVert)^2
\]
便得到三角不等式。结合上一节的点差相加公式，\(d\) 确实满足距离的三角不等式。对于两个非零向量，上述配对不等式还保证
\[
	\cos\theta=\frac{\langle u,v\rangle}{\lVert u\rVert\cdot\lVert v\rVert}
\]
确定唯一的 \(\theta\in[0,\pi]\)，这就是它们的\term[jiajiao]{夹角}[angle]。

\begin{proposition}[内积的恢复]
	\label{g1:prop:polarization}
	在实内积空间 \(V\) 上，长度确定内积，具体有
	\begin{equation}
		\label{g1:eq:inner-polarization}
		\langle u,v\rangle
		=\frac{\lVert u+v\rVert^2-\lVert u\rVert^2-\lVert v\rVert^2}{2}
		=\frac{\lVert u+v\rVert^2-\lVert u-v\rVert^2}{4}.
	\end{equation}
\end{proposition}

\begin{proof}
	由双线性和对称性展开 \(\langle u+v,u+v\rangle\) 与 \(\langle u-v,u-v\rangle\)，交叉项分别是 \(2\langle u,v\rangle\) 与 \(-2\langle u,v\rangle\)。整理即得两式。
\end{proof}

公式\eqref{g1:eq:inner-polarization}称为内积的\term[jihuahengdengshi]{极化恒等式}[polarization identity]。它说明保持所有长度的线性映射也保持所有内积。这里的“线性”暂时不可省去，因为展开时要知道映射怎样处理 \(u+v\)；下一节将解释，保持所有点间距离的映射为什么会恢复出仿射结构。

\ProseParagraphBreak

内积还使向量和线性函数联系起来。记 \(V^*\) 为 \(V\) 的\term[duioukongjian]{对偶空间}[dual space]，即所有线性映射 \(\alpha\colon V\rightarrow\mathbb R\) 组成的向量空间。\(\alpha(v)\) 是不需要内积就能定义的\term[duioupeidui]{对偶配对}[duality pairing]。给定内积以后，每个向量 \(v\) 又产生一个线性函数
\[
	v^\flat\colon w\longmapsto\langle v,w\rangle.
\]
映射 \(\flat\colon V\rightarrow V^*\) 是线性的；若 \(v^\flat=0\)，取 \(w=v\) 得 \(v=0\)，故它单射。给定 \(V\) 的一组基，逐个取坐标的线性函数构成 \(V^*\) 的一组基：任意线性函数由它在原基上的值唯一确定。因此 \(\dim V^*=\dim V\)，上述单射也是满射。于是每个 \(\alpha\in V^*\) 恰有一个向量 \(\alpha^\sharp\) 满足
\[
	\alpha(w)=\langle\alpha^\sharp,w\rangle\qquad(w\in V).
\]
\symidx{\(V^*\)：对偶空间}
\symidx{\(v^\flat,\alpha^\sharp\)：由内积给出的向量与对偶向量的对应}
对偶空间本身不依赖内积，但这项对应依赖内积，不能把它当作一个无需选择的相等关系。

\begin{example}[同一个线性函数对应不同的向量]
	\label{g1:exm:metric-dual}
	在 \(\mathbb R^2\) 上取
	\[
		\langle x,y\rangle=x_1y_1+4x_2y_2,\qquad
		\alpha(x)=x_1+2x_2.
	\]
	由 \(\alpha(x)=\langle(1,1/2),x\rangle\) 得 \(\alpha^\sharp=(1,1/2)\)。若改用通常的点积，则同一个 \(\alpha\) 对应 \((1,2)\)。更一般地，若基 \(b_1,\ldots,b_n\) 下的内积矩阵为 \(G=(\langle b_i,b_j\rangle)\)，向量 \(v\) 的坐标列为 \(c\)，则 \(v^\flat\) 的系数列是 \(Gc\)，因为 \(\langle v,w\rangle=c^{\mathsf T}Gx\)。把向量坐标与线性函数的系数直接等同，只有在 \(G=I\) 时才正确。
\end{example}

在内积为零时，称两个向量\term[zhengjiao]{正交}[orthogonal]，记 \(u\perp v\)。一组基 \(e_1,\ldots,e_n\) 若满足 \(\langle e_i,e_j\rangle=\delta_{ij}\)，便称为\term[biaozhunzhengjiaoji]{标准正交基}[orthonormal basis]。标准正交基把长度公式化成平方和；它的存在应当从内积公理推出，而不应预先藏在“欧氏”这个名称中。

\begin{proposition}[保持逐步张成空间的正交化]
	\label{g1:prop:orthonormalization}
	设 \(v_1,\ldots,v_k\) 是实有限维内积空间 \(V\) 中的线性无关向量。则存在标准正交向量组 \(e_1,\ldots,e_k\)，使得对每个 \(1\leq j\leq k\)，都有
	\[
		\operatorname{span}(e_1,\ldots,e_j)
		=\operatorname{span}(v_1,\ldots,v_j).
	\]
	特别地，\(V\) 有标准正交基，任意子空间的标准正交基都能扩充为 \(V\) 的标准正交基。
\end{proposition}

\begin{proof}
	构造的想法是从 \(v_j\) 中减去已经处理的方向，只保留新增加的方向。令
	\[
		w_1=v_1,\qquad e_1=\frac{w_1}{\lVert w_1\rVert}.
	\]
	若已经构造了前 \(j-1\) 个向量，定义
	\begin{equation}
		\label{g1:eq:orthogonalization}
		w_j=v_j-\sum_{i=1}^{j-1}\langle v_j,e_i\rangle e_i,
		\qquad e_j=\frac{w_j}{\lVert w_j\rVert}.
	\end{equation}
	对每个 \(\ell<j\)，标准正交性给出
	\[
		\langle w_j,e_\ell\rangle
		=\langle v_j,e_\ell\rangle
		-\sum_{i<j}\langle v_j,e_i\rangle\delta_{i\ell}=0.
	\]
	若 \(w_j=0\)，则 \(v_j\) 属于前 \(j-1\) 个 \(e_i\) 的张成空间，也就是前 \(j-1\) 个 \(v_i\) 的张成空间，与线性无关矛盾。因此除以 \(\lVert w_j\rVert\) 合法。式\eqref{g1:eq:orthogonalization}又说明 \(e_j\) 可由前 \(j\) 个 \(v_i\) 表出，而 \(v_j\) 可由前 \(j\) 个 \(e_i\) 表出，所以逐步张成空间相同。归纳完成构造。

	取 \(k=\dim V\) 便得到标准正交基。若从某个子空间的标准正交基出发，先把它扩充为 \(V\) 的一组普通基，再作上述构造，原有标准正交向量不变，因此得到所需扩充。
\end{proof}

这一递推构造称为\term[GramSchmidtzhengjiaohua]{Gram--Schmidt 正交化}[Gram--Schmidt orthogonalization]。\footnote{Gram（约尔根·佩德森·格拉姆，Jørgen Pedersen Gram，1850--1916），丹麦数学家，研究正交化、概率与数值方法；Schmidt（埃尔哈德·施密特，Erhard Schmidt，1876--1959），德国数学家，研究积分方程与无穷维内积空间。}输入是一组线性无关向量，输出是一组保留逐步张成空间的标准正交向量。它也说明正定性为何不可缺少：若允许非零向量的平方长度为零，构造就可能在归一化时停止。

\ProseParagraphBreak

现在考虑一个几何问题：给定向量 \(x\)，如何在子空间 \(W\) 中找到离它最近的向量？正交化提示我们，应当减去 \(x\) 在 \(W\) 的各个标准正交方向上的分量。记
\[
	W^\perp\coloneqq\{\,z\in V:\langle z,w\rangle=0\text{ 对所有 }w\in W\,\},
\]
称为 \(W\) 的\term[zhengjiaobu]{正交补}[orthogonal complement]。
\symidx{\(W^\perp\)：正交补}

\begin{proposition}[正交投影与最小距离]
	\label{g1:prop:projection}
	设 \(V\) 是实有限维内积空间，\(W\subseteq V\) 是线性子空间。每个 \(x\in V\) 都唯一地写成
	\[
		x=p+r,\qquad p\in W,\quad r\in W^\perp.
	\]
	映射 \(P_Wx=p\) 是线性的，满足 \(P_W^2=P_W\)、\(\operatorname{im}P_W=W\)、\(\ker P_W=W^\perp\)。而且对每个 \(w\in W\)，
	\begin{equation}
		\label{g1:eq:projection-distance}
		\lVert x-w\rVert^2
		=\lVert x-P_Wx\rVert^2+\lVert P_Wx-w\rVert^2.
	\end{equation}
	因此 \(P_Wx\) 是 \(W\) 中离 \(x\) 最近的唯一向量。
\end{proposition}

\begin{proof}
	取 \(W\) 的标准正交基 \(e_1,\ldots,e_k\)，令
	\[
		p=\sum_{i=1}^k\langle x,e_i\rangle e_i,\qquad r=x-p.
	\]
	逐项计算得 \(\langle r,e_j\rangle=0\)，从而 \(r\in W^\perp\)。若另有 \(x=p'+r'\)，其中 \(p'\in W\)、\(r'\in W^\perp\)，则 \(p-p'=r'-r\) 同时属于 \(W\) 和 \(W^\perp\)。这个向量与自己正交，正定性迫使它等于零。这证明唯一性，也说明构造结果不依赖所取标准正交基。

	上面的有限求和公式给出线性性。对 \(w\in W\) 有 \(P_Ww=w\)，对 \(r\in W^\perp\) 有 \(P_Wr=0\)，故幂等性、像与核的断言都成立。最后，\(x-w=(x-p)+(p-w)\) 是正交的两项之和，平方范数中的交叉项为零，得到式\eqref{g1:eq:projection-distance}。第二项非负，且恰在 \(w=p\) 时为零，故最小值在 \(p\) 处唯一取得。
\end{proof}

\(P_W\) 称为到 \(W\) 的\term[zhengjiaotouying]{正交投影}[orthogonal projection]。由分解可见
\[
	V=W\oplus W^\perp,\qquad \dim W^\perp=\dim V-\dim W,
	\qquad (W^\perp)^\perp=W.
\]
最后一式可由包含 \(W\subseteq(W^\perp)^\perp\) 和相同维数得到。\cref{g1:fig:orthogonal-projection}把距离恒等式画成一个直角三角形：换一个候选点 \(w\)，新增的平方距离正是沿子空间移动所付出的 \(\lVert P_Wx-w\rVert^2\)。
\symidx{\(P_W\)：到子空间 \(W\) 的正交投影}

\begin{figure}[htbp]
	\centering
	\input{Figures/OrthogonalProjection.tex}
	\caption[正交投影与最小距离]{正交投影与最小距离。图中以一维子空间示意，\(x-P_Wx\) 垂直于整个 \(W\)，因而也垂直于从 \(P_Wx\) 到任意 \(w\in W\) 的位移。}
	\label{g1:fig:orthogonal-projection}
\end{figure}

若 \(F=a+W\) 是欧氏仿射空间中的非空仿射子空间，最近点则是
\[
	\pi_F(x)=a+P_W(x-a).
\]
这是把原点临时移到 \(a\) 后应用\cref{g1:prop:projection}。若改用 \(a'=a+w_0\)，其中 \(w_0\in W\)，由 \(P_Ww_0=w_0\) 可得同一结果，故该公式不依赖基点选择。它刻画的最近点满足 \(x-\pi_F(x)\perp W\)，而非一定与位置向量 \(\pi_F(x)\) 正交。

\begin{example}[到一个仿射平面的距离]
	\label{g1:exm:plane-projection}
	在带通常点积的 \(\mathbb R^3\) 中，令 \(F=\{\,(x_1,x_2,x_3):x_1+x_2+x_3=1\,\}\)。其方向空间是 \(W=\{\,v:v_1+v_2+v_3=0\,\}\)，法向量 \(n=(1,1,1)\) 与 \(W\) 正交。因此原点的最近点必须形如 \(cn\)，再由 \(cn\in F\) 得 \(3c=1\)。最近点是 \((1/3,1/3,1/3)\)，距离为 \(1/\sqrt3\)。更一般地，若 \(n\ne0\)，到超平面 \(\langle n,y\rangle=b\) 的最近点为
	\[
		x-\frac{\langle n,x\rangle-b}{\lVert n\rVert^2}n,
	\]
	因为这个点在超平面上，而它与 \(x\) 的差沿着法向方向。
\end{example}

\subsection{自伴谱分解与正交标准形}
\label{g1:subsec:0102:02}

在标准正交坐标下，内积本身是平方和；若另有一个二次表达式，通常还会含有交叉项。例如 \(2x^2+2xy+2y^2\) 的等值线沿原来的坐标轴并不对称。是否可以只旋转或反射坐标轴，就把这些交叉项消去？答案来自自伴算子的特征方向。

\ProseParagraphBreak

给定实有限维内积空间上的线性算子 \(T\)，固定 \(v\) 后，\(u\mapsto\langle Tu,v\rangle\) 是线性函数。由本节已经证明的向量与对偶向量的对应，存在唯一向量 \(T^*v\)，使
\[
	\langle Tu,v\rangle=\langle u,T^*v\rangle\qquad(u,v\in V).
\]
唯一性同时说明 \(v\mapsto T^*v\) 线性。称 \(T^*\) 为 \(T\) 的\term[banshuisuanzi]{伴随算子}[adjoint operator]。在标准正交基下，它的矩阵是 \(T\) 的矩阵的转置；从定义连续使用两次配对等式还得到 \((AB)^*=B^*A^*\)、\((T^*)^*=T\)。例如 \(\langle ABu,v\rangle=\langle Bu,A^*v\rangle=\langle u,B^*A^*v\rangle\)，这说明乘积伴随为何颠倒次序。
\symidx{\(T^*\)：内积伴随算子}

\begin{definition}[自伴算子]
	\label{g1:def:self-adjoint}
	实内积空间 \(V\) 上的线性算子 \(T\colon V\rightarrow V\) 称为\term[zibansuanzi]{自伴算子}[self-adjoint operator]，若
	\[
		\langle Tu,v\rangle=\langle u,Tv\rangle\qquad(u,v\in V).
	\]
	若存在非零向量 \(v\) 与实数 \(\lambda\) 满足 \(Tv=\lambda v\)，则称 \(\lambda\) 为 \(T\) 的\term[tezhengzhi]{特征值}[eigenvalue]，\(v\) 为相应\term[tezhengxiangliang]{特征向量}[eigenvector]，并记\term[tezhengzikongjian]{特征子空间}[eigenspace]为
	\[
		E_\lambda(T)\coloneqq\ker(T-\lambda I).
	\]
	\symidx{\(E_\lambda(T)\)：特征子空间}
\end{definition}

若内积矩阵为 \(G\)，算子矩阵为 \(A\)，自伴条件写成 \(A^{\mathsf T}G=GA\)。只有在标准正交基下，即 \(G=I\) 时，它才简化为 \(A^{\mathsf T}=A\)。因此“矩阵对称”是自伴性在特定类型坐标中的表达。

\begin{theorem}[实自伴谱定理]
	\label{g1:thm:spectral}
	设 \(V\) 是 \(n\) 维实内积空间，\(T\colon V\rightarrow V\) 是自伴线性算子。则 \(V\) 有由 \(T\) 的特征向量组成的标准正交基。在这样的基下，\(T\) 的矩阵为
	\[
		\operatorname{diag}(\lambda_1,\ldots,\lambda_n),\qquad\lambda_i\in\mathbb R.
	\]
	不同特征值的特征子空间互相正交，\(V\) 是这些特征子空间的正交直和。特征值连同各自重数所组成的多重集由 \(T\) 唯一确定。
\end{theorem}

\begin{proof}
	关键是先找到一个实特征方向，再证明它的正交补仍由 \(T\) 保持，这样就能降低维数。若 \(n=0\)，结论由空基成立。设 \(n\geq1\)。由\cref{g1:prop:orthonormalization}，可以先取任意标准正交坐标。在这些坐标下，单位球面 \(\{\,x:\lVert x\rVert=1\,\}\) 是 \(\mathbb R^n\) 中的闭有界集合，因而紧；连续实函数在非空紧集上取得最大值。这里使用的是有限维实分析中的紧性与极值定理，有限维和正定性共同保证这个球面是通常的紧球面。

	函数 \(f(x)=\langle Tx,x\rangle\) 连续，所以在某个单位向量 \(u\) 处取得球面最大值，记为 \(\lambda=\langle Tu,u\rangle\)。对任意 \(v\perp u\)，曲线
	\[
		\gamma(t)=\frac{u+tv}{\sqrt{1+t^2\lVert v\rVert^2}}
	\]
	始终在单位球面上，并经过 \(u\)。因此 \(f(\gamma(t))\) 在 \(t=0\) 有极大值，其导数为零。利用自伴性直接展开：
	\[
		f(\gamma(t))
		=\frac{\lambda+2t\langle Tu,v\rangle+t^2\langle Tv,v\rangle}
		{1+t^2\lVert v\rVert^2}.
	\]
	在 \(t=0\) 求导得 \(2\langle Tu,v\rangle=0\)。所以 \(Tu\) 与所有 \(v\in u^\perp\) 正交。由正交分解，\(Tu\in\mathbb Ru\)；再与 \(u\) 配对得到 \(Tu=\lambda u\)。

	若 \(v\in u^\perp\)，则
	\[
		\langle Tv,u\rangle=\langle v,Tu\rangle
		=\lambda\langle v,u\rangle=0.
	\]
	故 \(u^\perp\) 在 \(T\) 下不变，限制 \(T|_{u^\perp}\) 仍满足自伴条件。对维数 \(n-1\) 应用归纳假设，得到 \(u^\perp\) 的标准正交特征基，与 \(u\) 合起来就是所求的基。

	若 \(Tu=\lambda u\)、\(Tv=\mu v\) 且 \(\lambda\ne\mu\)，则
	\[
		\lambda\langle u,v\rangle
		=\langle Tu,v\rangle
		=\langle u,Tv\rangle
		=\mu\langle u,v\rangle,
	\]
	所以 \(u\perp v\)。已有特征基又说明这些特征子空间张成整个 \(V\)。

	最后说明唯一性的含义。若换基矩阵为可逆矩阵 \(C\)，同一算子的两个矩阵为 \(A\) 与 \(C^{-1}AC\)，于是
	\[
		\det(tI-C^{-1}AC)=\det\bigl(C^{-1}(tI-A)C\bigr)=\det(tI-A).
	\]
	在所得特征基下，这个多项式等于 \(\prod_{i=1}^n(t-\lambda_i)\)。实一次因子的根和重数唯一，因而特征值的多重集不依赖所选特征基。
\end{proof}

这个结果称为实有限维的\term[pudingli]{谱定理}[spectral theorem]。“谱”在这里就是全部特征值。定理不声称特征基唯一：每个重特征值的特征子空间内，可以重新选择任意标准正交基；即使特征值互异，单位特征向量仍可改变符号。真正不依赖这些选择的是特征子空间及其正交投影，算子可以写成
\[
	T=\sum_{\lambda}\lambda P_{E_\lambda(T)}.
\]
若 \(x=\sum_i x_i e_i\)，则
\[
	\langle Tx,x\rangle=\sum_i\lambda_i x_i^2.
\]
因此 \(\langle Tx,x\rangle>0\) 对所有 \(x\ne0\) 成立，当且仅当全部 \(\lambda_i>0\)；非负性则对应全部 \(\lambda_i\geq0\)。

\begin{corollary}[实对称矩阵的正交对角化]
	\label{g1:cor:orthogonal-diagonalization}
	对每个实对称 \(n\times n\) 矩阵 \(A\)，存在满足 \(Q^{\mathsf T}Q=I\) 的实矩阵 \(Q\)，使
	\[
		Q^{\mathsf T}AQ=\operatorname{diag}(\lambda_1,\ldots,\lambda_n).
	\]
	两个同阶实对称矩阵能由这样的正交换基相互转化，当且仅当它们的特征值多重集相同。
\end{corollary}

\begin{proof}
	把 \(A\) 看作通常点积空间 \(\mathbb R^n\) 上的算子，它因 \(A^{\mathsf T}=A\) 而自伴。把\cref{g1:thm:spectral}所得标准正交特征基作为 \(Q\) 的各列，便有 \(Q^{\mathsf T}Q=I\) 及所示公式。正交换基也是相似变换，因此保持特征值多重集。反过来，若两个矩阵具有相同特征值多重集，调整特征基顺序，使两者的对角矩阵相同，复合这两个正交换基便得到所求变换。
\end{proof}

\begin{example}[椭圆的轴从哪里来]
	\label{g1:exm:ellipse-spectral}
	对称矩阵
	\[
		A=\begin{pmatrix}2&1\\1&2\end{pmatrix}
	\]
	分别在单位向量 \(e_+=(1,1)/\sqrt2\)、\(e_-=(1,-1)/\sqrt2\) 上取特征值 \(3\)、\(1\)。若 \((x,y)=u e_++v e_-\)，则
	\[
		2x^2+2xy+2y^2=3u^2+v^2.
	\]
	所以等值线 \(2x^2+2xy+2y^2=1\) 的两条主轴沿 \(e_+\)、\(e_-\)，对应半轴长分别为 \(1/\sqrt3\)、\(1\)。较大的特征值对应较短的半轴，因为沿该方向走同样的距离，二次表达式增长得更快。这里换成 \((u,v)\) 没有改变长度，只改变了标准正交坐标的方向。
\end{example}

\subsection{对称双线性型的合同分类与惯性}
\label{g1:subsec:0102:03}

正定性适合测量距离，但许多二次方程具有正负号，例如 \(x^2-y^2=1\)。这时仍有对称双线性配对，却不能把其对角值的平方根当作长度。我们因此把对称性与正定性分开。

\begin{definition}[对称双线性型与二次型]
	\label{g1:def:symmetric-form}
	设 \(V\) 是实向量空间。映射 \(B\colon V\times V\rightarrow\mathbb R\) 若对每个变量分别线性，并满足 \(B(u,v)=B(v,u)\)，称为\term[duichengshuangxianxingxing]{对称双线性型}[symmetric bilinear form]。其对应的\term[ercixing]{二次型}[quadratic form]是
	\[
		q(v)\coloneqq B(v,v).
	\]
	其\term[genji]{根基}[radical]定义为
	\[
		\operatorname{rad}B\coloneqq\{\,v\in V:B(v,w)=0\text{ 对所有 }w\in V\,\}.
	\]
	若 \(\operatorname{rad}B=\{0\}\)，称 \(B\) \term[feituihua]{非退化}[nondegenerate]；否则称其退化。
	\symidx{\(\operatorname{rad}B\)：双线性型的根基}
\end{definition}

实数域上，\(q\) 确定 \(B\)，因为
\[
	B(u,v)=\frac{q(u+v)-q(u)-q(v)}2.
\]
这里与式\eqref{g1:eq:inner-polarization}有相同的展开过程，完全不需要正定性。若取基 \(b_1,\ldots,b_n\)，令 \(A_{ij}=B(b_i,b_j)\)，则 \(A\) 对称，且坐标列为 \(x,y\) 的两个向量满足
\[
	B(u,v)=x^{\mathsf T}Ay,\qquad q(u)=x^{\mathsf T}Ax.
\]
这也说明二次多项式中的交叉项系数是相应非对角矩阵元的两倍。

\ProseParagraphBreak

若新基的各向量在旧基下组成列矩阵 \(C\)，则同一向量的旧坐标与新坐标满足 \(x=Cx'\)。因此同一个双线性型在新基下的矩阵是
\[
	A'=C^{\mathsf T}AC.
\]
这称为\term[hetong]{合同}[congruence]变换。相比之下，同一个线性算子的矩阵变为 \(C^{-1}AC\)，称为\term[xiangsi]{相似}[similarity]变换。\cref{g1:tab:similarity-congruence}列出区别。只在 \(C\) 正交、即 \(C^{-1}=C^{\mathsf T}\) 时，两种公式才重合。

\begin{table}[htbp]
	\centering
	\caption{相似变换与合同变换}
	\label{g1:tab:similarity-congruence}
	\begin{tabularx}{\linewidth}{@{}lXX@{}}
		\toprule
		& 线性算子 & 对称双线性型 \\
		\midrule
		作用对象 & 一个输入向量与一个输出向量 & 两个输入向量与一个实数值 \\
		换基公式 & \(C^{-1}AC\) & \(C^{\mathsf T}AC\) \\
		一般换基所保留 & 特征多项式、特征值 & 正负方向的最大维数、根基维数 \\
		\bottomrule
	\end{tabularx}
\end{table}

为了运用刚证得的谱定理，可以在 \(V\) 上临时选择一个辅助内积。由向量与对偶空间的对应，对每个 \(u\) 恰有一个 \(Tu\) 使
\[
	B(u,v)=\langle Tu,v\rangle\qquad(v\in V).
\]
配对的线性性和表示的唯一性说明 \(T\) 是线性算子，\(B\) 的对称性说明 \(T\) 自伴。谱定理于是把 \(B\) 化成实对角形。辅助内积只是获得某组基的工具；接下来的唯一性证明会把最终分类量刻画成 \(B\) 自己的性质。

\begin{theorem}[实对称双线性型的惯性定理]
	\label{g1:thm:inertia}
	设 \(V\) 是 \(n\) 维实向量空间，\(B\) 是 \(V\) 上的对称双线性型。则存在一组基，使 \(B\) 的矩阵为
	\[
		\operatorname{diag}(I_p,-I_q,0_r),\qquad p+q+r=n.
	\]
	三元组 \((p,q,r)\) 唯一。其中 \(p\) 是使 \(B|_{U\times U}\) 正定的子空间 \(U\) 的最大维数，\(q\) 是使 \(-B|_{U\times U}\) 正定的子空间 \(U\) 的最大维数，\(r=\dim\operatorname{rad}B\)。两个同维实对称双线性型合同，当且仅当这三个整数相同。
\end{theorem}

\begin{proof}
	先证明标准形存在。任取辅助内积，将 \(B\) 写成 \(\langle Tu,v\rangle\)。由\cref{g1:thm:spectral}，存在标准正交特征基 \(e_1,\ldots,e_n\)，其中
	\[
		B(e_i,e_j)=\lambda_i\delta_{ij}.
	\]
	对每个非零 \(\lambda_i\)，把 \(e_i\) 改为 \(e_i/\sqrt{\lvert\lambda_i\rvert}\)，对应的对角元就变成 \(1\) 或 \(-1\)；零特征值对应的向量保持不变。把正、负、零三组重新排序便得到所示标准形。这次缩放通常不是正交换基，所以它可以改变特征值的绝对大小。

	再证明计数与基无关。设所得三组向量张成的空间分别为 \(V_+\)、\(V_-\)、\(V_0\)。显然 \(B\) 在 \(V_+\) 上正定，故正定子空间的最大维数至少为 \(p\)。若 \(U\) 是任意正定子空间，则
	\[
		U\cap(V_-\oplus V_0)=\{0\},
	\]
	因为右边空间中的每个向量都满足 \(B(v,v)\leq0\)，而 \(U\) 中每个非零向量满足 \(B(v,v)>0\)。将 \(U\) 沿 \(V_-\oplus V_0\) 投影到 \(V_+\)，上述交集条件说明该线性映射单射，故 \(\dim U\leq\dim V_+=p\)。因此 \(p\) 恰等于所有正定子空间的最大维数。对 \(-B\) 应用同一论证得负定方向的最大维数为 \(q\)。

	在标准形坐标中，\(B(v,w)=0\) 对所有 \(w\) 成立，当且仅当 \(v\) 的前 \(p+q\) 个坐标都为零。因此 \(\operatorname{rad}B=V_0\)，其维数是 \(r\)。这些描述只涉及 \(B\)，没有涉及所选基或辅助内积，故三元组唯一。合同变换是同一配对的换基，必保持这三个数；反过来，三元组相同的两个型都能化为同一个标准矩阵，复合相应换基就得到二者的合同关系。
\end{proof}

三元组 \((p,q,r)\) 称为 \(B\) 的\term[guanxing]{惯性}[inertia]；\(p\)、\(q\) 分别称\term[zhengguanxingzhishu]{正惯性指数}[positive index of inertia]与\term[fuguanxingzhishu]{负惯性指数}[negative index of inertia]。\(p+q\) 是 \(B\) 的秩，\(r\) 是零化度。

\begin{example}[零值方向不等于退化方向]
	\label{g1:exm:isotropic-radical}
	在 \(\mathbb R^3\) 上考虑
	\[
		q(x,y,z)=x^2+4xy+y^2=(x+2y)^2-3y^2.
	\]
	取新坐标 \(u=x+2y\)、\(v=\sqrt3\,y\)、\(w=z\)，得到 \(q=u^2-v^2\)，惯性为 \((1,1,1)\)，根基为 \(z\) 轴。另一方面，在 \((u,v,w)\) 坐标中，向量 \((1,1,0)\) 满足 \(q=0\)，却不在根基中，因为它与 \((1,0,0)\) 的配对为 \(1\)。非零而满足 \(q(v)=0\) 的向量称为\term[mixiangxiangliang]{迷向向量}[isotropic vector]。迷向性只检查向量与自己的配对，根基则要求它与所有向量的配对都为零。这一差别也说明，不定二次型不能直接给出距离。
\end{example}

\begin{example}[相同惯性仍可有不同轴长]
	\label{g1:exm:congruence-vs-orthogonal}
	矩阵 \(I_2\) 与 \(\operatorname{diag}(1,4)\) 都有惯性 \((2,0,0)\)，并且
	\[
		\begin{pmatrix}1&0\\0&1/2\end{pmatrix}^{\!\mathsf T}
		\begin{pmatrix}1&0\\0&4\end{pmatrix}
		\begin{pmatrix}1&0\\0&1/2\end{pmatrix}=I_2.
	\]
	它们因而合同，却不能由正交换基相互转化，因为特征值多重集分别为 \(\{1,1\}\) 和 \(\{1,4\}\)。几何上，一个坐标方向的缩放可以把椭圆变成圆，保持长度的坐标变换却不能做到。允许的变换越多，分类留下的信息就可能越少。
\end{example}

\subsection{Gram 行列式、体积与退化的检测}
\label{g1:subsec:0102:04}

测量一条边只需要它的长度；测量一个由多条边张成的平行四边形或平行多面体，还需要知道这些边之间是否接近平行。把所有两两内积放进一个矩阵，正好同时保存长度与夹角的信息。

\begin{definition}[Gram 矩阵]
	\label{g1:def:gram-matrix}
	设 \(v_1,\ldots,v_k\) 是实内积空间中的向量，它们的\term[Gramjuzhen]{Gram 矩阵}[Gram matrix]为
	\[
		\Gamma(v_1,\ldots,v_k)\coloneqq
		\bigl(\langle v_i,v_j\rangle\bigr)_{1\leq i,j\leq k}.
	\]
	其行列式称为\term[Gramhanglieshi]{Gram 行列式}[Gram determinant]。
	\symidx{\(\Gamma(v_1,\ldots,v_k)\)：Gram 矩阵}
\end{definition}

对任意列向量 \(c=(c_1,\ldots,c_k)^{\mathsf T}\)，双线性给出
\begin{equation}
	\label{g1:eq:gram-quadratic}
	c^{\mathsf T}\Gamma c
	=\left\lVert\sum_{i=1}^k c_i v_i\right\rVert^2.
\end{equation}
因此 \(\Gamma\) 总是\term[banzhengding]{半正定的}[positive semidefinite]，即对应二次型处处非负；只有当这些向量线性无关时，它才正定。这个等式已经预示：Gram 行列式为零应当对应边向量的线性相关，而它的正平方根应当测量所张成的体积。

\ProseParagraphBreak

先明确体积的含义。对线性无关的 \(v_1,\ldots,v_k\)，记 \(W=\operatorname{span}(v_1,\ldots,v_k)\)。取 \(W\) 的标准正交基，把这些向量写成一个 \(k\times k\) 坐标矩阵 \(C\) 的各列。它们张成的平行多面体为
\[
	\mathcal P(v_1,\ldots,v_k)
	=\left\{\,\sum_{i=1}^k t_iv_i:0\leq t_i\leq1\,\right\},
\]
其 \(k\) 维欧氏体积定义为 \(\lvert\det C\rvert\)。这个定义采用了单位正交立方体体积为 \(1\) 的归一化：有向体积对各边线性、交换两边变号，且在所选的标准正交基上的值为 \(1\)，按各列的坐标展开就得到行列式公式。取绝对值后不再记录定向。

\ProseParagraphBreak

若换一个标准正交基，坐标矩阵变为 \(Q^{\mathsf T}C\)，其中 \(Q^{\mathsf T}Q=I\)。取行列式可得 \((\det Q)^2=1\)，所以绝对行列式不变。体积因此不依赖所选标准正交坐标。若 \(v_1,\ldots,v_k\) 线性相关，我们把它们的 \(k\) 维体积定义为零：此时多面体实际上落在低于 \(k\) 维的空间内。这个约定不否认它可能仍有非零的较低维体积。\(k=0\) 时零维平行多面体 \(\{0\}\) 的体积约定为 \(1\)。

\begin{proposition}[Gram 行列式与体积]
	\label{g1:prop:gram-volume}
	设 \(V\) 是 \(n\) 维实内积空间，\(v_1,\ldots,v_k\in V\)。则
	\[
		\operatorname{Vol}_k\bigl(\mathcal P(v_1,\ldots,v_k)\bigr)^2
		=\det\Gamma(v_1,\ldots,v_k).
	\]
	Gram 矩阵的秩等于 \(\operatorname{span}(v_1,\ldots,v_k)\) 的维数。特别地，当 \(k\geq1\) 时，Gram 行列式为正当且仅当各向量线性无关，为零当且仅当各向量线性相关；若 \(k>n\)，它必为零。
\end{proposition}

\begin{proof}
	令 \(L\colon\mathbb R^k\rightarrow V\) 为 \(L(c)=\sum_i c_iv_i\)。若 \(L(c)=0\)，则 \(\Gamma c\) 的第 \(j\) 个分量 \(\langle v_j,L(c)\rangle\) 为零，故 \(\ker L\subseteq\ker\Gamma\)。反之，若 \(\Gamma c=0\)，式\eqref{g1:eq:gram-quadratic}给出 \(\lVert L(c)\rVert^2=0\)，由正定性知 \(L(c)=0\)。因此两者的核相同。对同一个定义域应用秩与零化度公式，得到
	\[
		\operatorname{rank}\Gamma=\operatorname{rank}L
		=\dim\operatorname{span}(v_1,\ldots,v_k).
	\]
	若各向量线性相关，\(\Gamma\) 不可逆，行列式和所定义的 \(k\) 维体积都为零。

	若各向量线性无关，在它们张成的 \(k\) 维空间中取标准正交基，记坐标矩阵为 \(C\)。逐项计算有 \(\Gamma=C^{\mathsf T}C\)，所以
	\[
		\det\Gamma=\det(C^{\mathsf T})\det C=(\det C)^2
		=\operatorname{Vol}_k(\mathcal P)^2>0.
	\]
	线性无关时行列式为正，与线性相关时的零值一起给出了退化判据。\(k=0\) 时空矩阵的行列式约定为 \(1\)，与体积约定相容。
\end{proof}

体积公式还可以从“底乘高”理解。对线性无关的边向量作式\eqref{g1:eq:orthogonalization}的正交化。记
\[
	h_j=\lVert w_j\rVert
	=d\bigl(v_j,\operatorname{span}(v_1,\ldots,v_{j-1})\bigr),
\]
其中第一步的张成空间是 \(\{0\}\)。由\cref{g1:prop:projection}，\(w_j\) 正是减去正交投影后的剩余向量。在所得标准正交基中，第 \(j\) 列只在前 \(j\) 行有非零坐标，对角元为 \(h_j\)。因此坐标矩阵上三角，得到
\[
	\operatorname{Vol}_k(\mathcal P)=h_1h_2\cdots h_k.
\]
每增加一条边，体积乘上它离已有边所张成空间的高度。任一高度为零都意味着并未增加一个新维度。

\begin{example}[三维空间中的面积与体积]
	\label{g1:exm:gram-volume}
	在通常的 \(\mathbb R^3\) 中取
	\[
		v_1=(1,1,0),\qquad v_2=(1,0,1),\qquad v_3=(0,1,1).
	\]
	前两条边的 Gram 矩阵为
	\[
		\begin{pmatrix}2&1\\1&2\end{pmatrix},
	\]
	故平行四边形面积为 \(\sqrt3\)。同一数也可由底乘高得到：底长是 \(\sqrt2\)，而
	\[
		\left\lVert v_2-\frac{\langle v_2,v_1\rangle}{\lVert v_1\rVert^2}v_1\right\rVert^2
		=2-\frac12=\frac32.
	\]
	三条边的 Gram 矩阵为
	\[
		\begin{pmatrix}2&1&1\\1&2&1\\1&1&2\end{pmatrix},
	\]
	其行列式为 \(4\)，所以平行六面体体积为 \(2\)。各边长都是 \(\sqrt2\)，但体积不是 \((\sqrt2)^3\)，因为三条边并不两两正交。
\end{example}

\begin{example}[长度不消失而面积退化]
	\label{g1:exm:gram-degeneration}
	在 \(\mathbb R^2\) 中令 \(v_1=(1,0)\)、\(v_2(t)=(1,t)\)。它们的长度在 \(t\rightarrow0\) 时都趋于 \(1\)，但
	\[
		\Gamma(t)=\begin{pmatrix}1&1\\1&1+t^2\end{pmatrix},
		\qquad \det\Gamma(t)=t^2,
	\]
	所以面积为 \(\lvert t\rvert\)，趋于零。退化发生在两个方向逐渐重合，而不是某条边变短。对任意两个非零向量，Gram 公式也写成
	\[
		\operatorname{Area}^2
		=\lVert v_1\rVert^2\lVert v_2\rVert^2-\langle v_1,v_2\rangle^2
		=\lVert v_1\rVert^2\lVert v_2\rVert^2\sin^2\theta,
	\]
	这正是熟悉的“边长乘边长再乘夹角正弦”。
\end{example}

最后要保留正定性这个条件。若把 Gram 矩阵中的内积换成任意对称双线性型，行列式为零便未必表示线性相关。例如 \(B((x,y),(x',y'))=xx'-yy'\) 非退化，而非零向量 \(v=(1,1)\) 所组成的一元向量组线性无关，其配对矩阵却是 \((0)\)。\cref{g1:prop:gram-volume}证明中从平方长度为零推出向量为零的一步，此时已经失效。

% !TeX root = ../../Book1.tex
% 纲要标题：### 1.3 Euclid 等距变换的结构定理
% 纲要位置：几何学计划纲要.md，第 134 行。
% 正文状态：初稿；尚未正式审核。

% 纲要内容开始
% 
% * 由距离保持恢复仿射性
% 
%   先用中点及极化恢复线性结构，再分离正交部分和平移部分。
% 
% * 正交算子的实标准分块
% * 不动点、螺旋运动与平移分量
% * 反射生成与低维分类
% 
% 纲要内容结束

\section{Euclid 等距变换的结构定理}
\label{g1:sec:0103}

一张刚性平面图形可以平移、旋转，也可以翻转后与原图形重合。这些动作都保持两点间的距离，但只从“距离保持”出发，是否还可能有形状更复杂的变换？关键是用距离唯一确定中点：保持距离便保持中点，连续性再把中点关系推广为任意仿射组合。得到仿射形式后，正交部分控制旋转与反射，平移部分则由不动方向决定哪些分量可以通过换原点消去。

\subsection{由距离保持恢复仿射性}
\label{g1:subsec:0103:01}

在一个有限维实仿射空间的平移向量空间上给定正定内积，就得到了\term[oushikongjian]{Euclid 空间}[Euclidean space]。其距离定义为 \(d(p,q)=\lVert q-p\rVert\)。选取原点和标准正交方向基后，可以将它写成 \(\mathbb R^n\)，内积成为通常的点积；更换原点不会改变点差。以下先在这种坐标中论证，再把结果解释为不依赖原点的几何结构。

\begin{definition}
	\label{g1:def:isometry}
	设 \(E\) 是有限维 Euclid 空间。双射 \(f\colon E\rightarrow E\) 若对所有 \(p,q\in E\) 满足
	\[
		d(f(p),f(q))=d(p,q),
	\]
	就称为 \(E\) 的\term[dengjubianhuan]{等距变换}[isometry]。若暂不要求双射，只要求上述等式，称 \(f\) 为距离保持映射。
\end{definition}

距离保持首先给出连续性：两像点之间的距离恰好等于原距离。更关键的是，它还给出了中点。对任意 \(a,b,z\in\mathbb R^n\)，展开内积可得
\[
	\lVert z-a\rVert^2+\lVert z-b\rVert^2
	=2\left\lVert z-\frac{a+b}{2}\right\rVert^2
	+\frac12\lVert a-b\rVert^2.
\]
因此，同时满足
\[
	\lVert z-a\rVert=\lVert z-b\rVert=\frac12\lVert a-b\rVert
\]
的点只有 \((a+b)/2\)。这里既要求两段等长，也要求其长度恰好是端点距离的一半；当 \(a\ne b\) 时，仅有“到两端等距”会得到整个中垂超平面，不能确定中点。

\ProseParagraphBreak

若实有限维内积空间上的线性算子 \(Q\) 对所有向量都满足 \(\langle Qu,Qv\rangle=\langle u,v\rangle\)，就称其为\term[zhengjiaosuanzi]{正交算子}[orthogonal operator]。由伴随的定义，这等价于 \(Q^*Q=I\)。这样的算子单射，因定义域和值域同维而可逆，故 \(Q^*=Q^{-1}\)。在标准正交基下，其矩阵满足 \(Q^{\mathsf T}Q=I\)。复合与求逆都保持这一性质，所有这样的 \(n\) 阶矩阵组成\term[zhengjiaoqun]{正交群}[orthogonal group]，记为 \(O(n)\)。
\symidx{\(O(n)\)：实正交群}

\begin{theorem}[等距变换的仿射形式]
	\label{g1:thm:isometry-affine}
	设 \(n\geq 1\)，且 \(f\colon\mathbb R^n\rightarrow\mathbb R^n\) 对所有 \(x,y\) 都满足 \(\lVert f(x)-f(y)\rVert=\lVert x-y\rVert\)。那么存在唯一的向量 \(t\in\mathbb R^n\) 和唯一的正交线性算子 \(Q\)，使
	\[
		f(x)=Qx+t.
	\]
	特别地，\(f\) 自动为双射。反之，每个这种形式的映射都是等距变换。
\end{theorem}

\begin{proof}
	先证明中点保持，再利用连续性把整数、有理数的运算推广到实数。若 \(m=(x+y)/2\)，则距离保持说明 \(f(m)\) 到 \(f(x),f(y)\) 的距离均为 \(\lVert f(x)-f(y)\rVert/2\)。由刚才证明的中点唯一性，
	\[
		f\left(\frac{x+y}{2}\right)=\frac{f(x)+f(y)}2.
	\]
	这一步只涉及已知像点，不需要预先知道 \(f\) 满射。

	令 \(t=f(0)\)，\(g(x)=f(x)-t\)，则 \(g(0)=0\)，且 \(g\) 同样保持中点与距离。取端点为 \(0,2x\)，得到 \(g(2x)=2g(x)\)，从而
	\[
		g(x+y)=2g\left(\frac{x+y}{2}\right)=g(x)+g(y).
	\]
	由加法性可得 \(g(-x)=-g(x)\)、\(g(mx)=mg(x)\) 对整数 \(m\) 成立；若 \(r=m/k\) 为有理数且 \(k>0\)，则 \(kg(rx)=g(krx)=mg(x)\)，所以 \(g(rx)=rg(x)\)。对于任意实数 \(a\)，取有理数列 \(r_j\rightarrow a\)。距离保持给出
	\[
		\lVert g(r_jx)-g(ax)\rVert=|r_j-a|\lVert x\rVert\longrightarrow 0,
	\]
	因而 \(g(ax)=ag(x)\)。这就证明了 \(g\) 为实线性算子，记为 \(Q\)。

	最后由极化公式恢复内积。将\cref{g1:prop:polarization}中的平方范数展开用于 \(x-y\)，得到
	\[
		\langle Qx,Qy\rangle
		=\frac12\bigl(\lVert Qx\rVert^2+\lVert Qy\rVert^2
		-\lVert Q(x-y)\rVert^2\bigr)
		=\langle x,y\rangle.
	\]
	所以 \(Q\) 正交。特别地，\(Qx=0\) 蕴含 \(x=0\)，有限维同维线性映射的秩与核维数公式便说明 \(Q\) 可逆；于是 \(f\) 的逆为 \(y\mapsto Q^{-1}(y-t)\)。唯一性来自 \(t=f(0)\) 和 \(Qx=f(x)-f(0)\)。反向验证只需计算 \(\lVert Qx+t-(Qy+t)\rVert=\lVert Q(x-y)\rVert=\lVert x-y\rVert\)。
\end{proof}

按照\cref{g1:def:affine-map}，上一定理正是说：Euclid 空间中只用距离提出的保持条件，已经强制映射具有仿射形式。有限维和定义域、值域同维在由单射推出满射时起作用。例如 \(x\mapsto(x,0)\) 从 \(\mathbb R^n\) 到 \(\mathbb R^{n+1}\) 保持距离，却不能覆盖整个值域。

\subsection{正交算子的实标准分块}
\label{g1:subsec:0103:02}

平移部分暂时放在一边。正交算子满足 \(Q^*Q=I\)，其中 \(Q^*\) 是内积伴随；在标准正交基下，这就是矩阵转置满足 \(Q^{\mathsf T}Q=I\)。它通常并不自伴：平面上一个转角既不是 \(0\) 也不是 \(\pi\)（模 \(2\pi\)）的旋转，就没有非零实特征向量。因此不能直接把\cref{g1:thm:spectral}应用于 \(Q\)。我们将对它的自伴部分应用该定理，余下的部分恰好给出旋转平面。

\begin{theorem}[正交算子的实分块]
	\label{g1:thm:orthogonal-blocks}
	设 \(V\) 是 \(n\) 维实内积空间，\(Q\colon V\rightarrow V\) 为正交算子。存在一组标准正交基，使其矩阵成为
	\[
		\operatorname{diag}\bigl(I_p,-I_q,R_{\theta_1},\ldots,R_{\theta_k}\bigr),
		\qquad p+q+2k=n,
	\]
	其中
	\begin{equation}
		\label{g1:eq:rotation-block}
		R_\theta=
		\begin{pmatrix}\cos\theta&-\sin\theta\\
			\sin\theta&\cos\theta\end{pmatrix},
		\qquad 0<\theta<\pi.
	\end{equation}
	整数 \(p,q\) 和角度的无序多重集 \(\{\theta_1,\ldots,\theta_k\}\) 由 \(Q\) 唯一确定；各个旋转平面的选择未必唯一。特别地，\(\det Q=(-1)^q\)。
\end{theorem}

\begin{proof}
	取自伴算子
	\[
		S=\frac{Q+Q^*}{2}.
	\]
	由\cref{g1:thm:spectral}，有限维实内积空间中的自伴算子有正交特征分解，故 \(V\) 是 \(S\) 的各特征空间 \(E_\lambda\) 的正交直和。因为 \(Q^*=Q^{-1}\)，所以 \(QS=SQ\)，每个 \(E_\lambda\) 均被 \(Q\) 保持。若 \(v\in E_\lambda\)，则
	\[
		\lambda\lVert v\rVert^2=\langle Sv,v\rangle
		=\langle Qv,v\rangle,
	\]
	由内积不等式得 \(|\lambda|\leq 1\)。对于端点，
	\[
		\lVert Qv-v\rVert^2=2(1-\lambda)\lVert v\rVert^2,
		\qquad
		\lVert Qv+v\rVert^2=2(1+\lambda)\lVert v\rVert^2,
	\]
	故 \(Q|_{E_1}=I\)，\(Q|_{E_{-1}}=-I\)。

	现在设 \(-1<\lambda<1\)，并记 \(s=\sqrt{1-\lambda^2}>0\)。在 \(E_\lambda\) 上有 \(Q+Q^*=2\lambda I\)；左乘 \(Q\) 得 \(Q^2-2\lambda Q+I=0\)。因此，去掉 \(Q\) 的标量部分并重新归一化，得到
	\[
		J=\frac{Q-\lambda I}{s},\qquad J^2=-I,
		\qquad J^*=-J,
		\qquad J^*J=I.
	\]
	这个构造的作用是把一般角度旋转还原成直角旋转：\(J^2=-I\) 表明连续做两次得到反向，而 \(J^*=-J\) 给出 \(\langle e,Je\rangle=0\)。选取任意单位向量 \(e\in E_\lambda\)，则 \(e,Je\) 是正交单位向量，其张成的平面 \(P\) 被 \(J\) 保持，并且
	\[
		Qe=\lambda e+sJe,\qquad Q(Je)=-se+\lambda Je.
	\]
	取唯一的 \(\theta\in(0,\pi)\) 使 \(\cos\theta=\lambda\)、\(\sin\theta=s\)，便得到\eqref{g1:eq:rotation-block}。

	在这个平面的正交补中可以继续拆分。若 \(w\in E_\lambda\cap P^\perp\) 且 \(z\in P\)，则 \(\langle Jw,z\rangle=-\langle w,Jz\rangle=0\)，所以 \(E_\lambda\cap P^\perp\) 也被 \(J\) 保持。在这个子空间中重复选取单位向量及其 \(J\) 像，每次使维数降低二，直至空间耗尽。若余下一维，任取非零向量仍会产生与它正交的非零 \(J\) 像，矛盾。因此 \(\dim E_\lambda\) 必为偶数，产生 \(\dim E_\lambda/2\) 个相同角度的旋转块。

	最后，\(p=\dim E_1\)、\(q=\dim E_{-1}\)，其余角度及重数由 \(S\) 的特征值和特征空间维数确定，故与基的选择无关。每个旋转块行列式为 \(1\)，其余一维块为 \(1\) 或 \(-1\)，于是得到行列式公式。
\end{proof}

这里把 \(\theta=0\) 归入两个 \(+1\) 块，把 \(\theta=\pi\) 归入两个 \(-1\) 块，因而标准形的记法不会把端点重复计算。一个平面的基向量顺序反过来时，\(R_\theta\) 会变为 \(R_{-\theta}\)；定理允许任意正交基，所以可以统一取 \(0<\theta<\pi\)。若平面已有固定定向，并且只准使用与之相容的基，旋转角的正负就仍须保留。

\begin{definition}
	\label{g1:def:direct-isometry}
	对于有限维 Euclid 空间的等距变换 \(f(x)=Qx+t\)，若 \(\det Q=1\)，称其\term[baoxiangdengjubianhuan]{保向}[orientation-preserving]；若 \(\det Q=-1\)，称其\term[fanxiangdengjubianhuan]{反向}[orientation-reversing]。
\end{definition}

这一称呼可以直接用有序基理解：在两组有序基之间，换基矩阵行列式为正时将它们视为同一定向，为负时视为相反定向。正交算子的行列式只能是 \(\pm1\)，分别保持或交换这两类基；平移不改变差向量，因而不影响定向。

\subsection{不动点、螺旋运动与平移分量}
\label{g1:subsec:0103:03}

写出 \(f(x)=Qx+t\) 后，接着考察换原点怎样改变平移部分。纯平移的向量在换原点时保持不变。若新原点的旧坐标是 \(c\)，则新坐标为 \(y=x-c\)，且
\begin{equation}
	\label{g1:eq:isometry-origin-change}
	f(c+y)-c=Qy+t-(I-Q)c.
\end{equation}
能消去哪些平移分量，取决于 \(I-Q\) 的像；特别是 \(Q\) 固定的方向，不能由 \((I-Q)c\) 产生。下面把这一障碍与可消去的部分分开。

\begin{theorem}[等距变换的平移分解]
	\label{g1:thm:isometry-normal}
	设 \(f(x)=Qx+t\) 是 \(\mathbb R^n\) 的等距变换，令
	\[
		F=\ker(I-Q),\qquad t=a+b,\qquad a\in F,\quad b\in F^\perp
	\]
	为 \(t\) 的正交分解。存在唯一的 \(c\in F^\perp\) 满足 \((I-Q)c=b\)，且对所有 \(u\in F\)、\(v\in F^\perp\)，
	\begin{equation}
		\label{g1:eq:isometry-normal}
		f(c+u+v)=c+u+a+Qv.
	\end{equation}
	变换 \(f\) 有不动点当且仅当 \(a=0\)；这时不动点集恰为 \(c+F\)。对所有 \(x\in\mathbb R^n\)，有
	\[
		\lVert f(x)-x\rVert\geq\lVert a\rVert,
	\]
	且等号恰在仿射子空间 \(c+F\) 上成立。该子空间被 \(f\) 保持，\(f\) 在其上就是沿 \(a\) 的平移。
\end{theorem}

\begin{proof}
	由\cref{g1:prop:projection}，有限维内积空间中的向量有相对于 \(F\) 的唯一正交分解。若 \(v\in F^\perp\)、\(u\in F\)，则
	\[
		\langle Qv,u\rangle=\langle v,Q^*u\rangle
		=\langle v,u\rangle=0,
	\]
	其中 \(Q^*u=Q^{-1}u=u\)，故 \(F^\perp\) 被 \(Q\) 和 \(I-Q\) 保持。在 \(F^\perp\) 上，\((I-Q)v=0\) 蕴含 \(v\in F\cap F^\perp=\{0\}\)，所以这个子空间上的线性算子 \(I-Q\) 单射，也就可逆。于是存在唯一 \(c\in F^\perp\) 使 \((I-Q)c=b\)。将 \(t=a+b\) 代入\eqref{g1:eq:isometry-origin-change}，再用 \(Qu=u\)，得到\eqref{g1:eq:isometry-normal}。

	若 \(x=c+u+v\)，位移分成两个正交部分：
	\[
		f(x)-x=a+(Q-I)v,
		\qquad
		\lVert f(x)-x\rVert^2=\lVert a\rVert^2+\lVert(Q-I)v\rVert^2.
	\]
	由于 \(I-Q\) 在 \(F^\perp\) 上可逆，第二项为零当且仅当 \(v=0\)。这既给出最小位移及其取得位置，也说明存在零位移点当且仅当 \(a=0\)，此时所有零位移点为 \(c+F\)。在 \(v=0\) 时，\eqref{g1:eq:isometry-normal}还直接给出 \(f(c+u)=c+u+a\)。
\end{proof}

向量 \(a\) 是无法通过平移原点消去的部分：由\eqref{g1:eq:isometry-origin-change}，原点改变只给 \(t\) 添加 \(F^\perp\) 中的向量，所以它在 \(F\) 上的投影不变。虽然 \(c\) 的坐标依赖起初的原点，\(c+F\) 却可以直接从 \(f\) 识别为“位移最小的点的集合”。这说明分解中的几何对象比最初写下的 \(t\) 更有意义。

\begin{example}[滑移反射及其最小位移]
	\label{g1:exm:glide-reflection}
	考虑平面变换 \(f(x,y)=(x+h,-y)\)，其中 \(h\ne0\)。它先关于 \(x\) 轴反射，再沿该轴平移 \((h,0)\)，称为\term[huayifanshe]{滑移反射}[glide reflection]。这两步可以交换次序，因为反射不改变轴方向的向量。对点 \((x,y)\)，位移为 \((h,-2y)\)，长度为 \(\sqrt{h^2+4y^2}\)，所以它没有不动点；轴上的位移最小，最小值为 \(|h|\)。\cref{g1:fig:glide-displacement}把沿轴的平移和垂直于轴的翻转分别画出。若 \(h=0\)，变换退化为普通轴反射，本文不再将其列入非平凡滑移反射。
\end{example}

\begin{figure}[htbp]
	\centering
	\input{Figures/Chapter01-GlideReflection.tex}
	\caption[滑移反射的两部分]{滑移反射的两部分。先由 \(p\) 反射到 \(r(p)\)，再沿轴平移到 \(f(p)\)；虚线表示垂直于轴的位移分量，蓝色斜箭头表示总位移。轴上每一点只发生长度为 \(|h|\) 的平移。}
	\label{g1:fig:glide-displacement}
\end{figure}

在三维空间中，若 \(Q\) 固定一条直线，并在其垂直平面上旋转角 \(\theta\)，正规形为
\[
	(x,y,z)\longmapsto
	\bigl(x\cos\theta-y\sin\theta,
	x\sin\theta+y\cos\theta,z+h\bigr).
\]
当 \(\theta\not\equiv0\pmod{2\pi}\) 且 \(h\ne0\) 时，称为\term[luoxuanyundong]{螺旋运动}[screw motion]。每次作用绕同一轴转动相同角度，同时沿轴前进相同距离；旋转与轴向平移可交换。距轴为 \(r\) 的点在垂直平面内移动的弦长平方为 \(2r^2(1-\cos\theta)\)，所以它的位移长度为
\[
	\sqrt{h^2+4r^2\sin^2(\theta/2)}.
\]
最小位移为 \(|h|\)，恰在轴上取得。\(h=0\) 时是绕轴旋转；\(\theta\equiv0\) 时是纯平移或恒等变换，此时不能由变换本身选出唯一的轴。

\begin{example}[从坐标公式找螺旋轴]
	\label{g1:exm:screw-axis}
	设
	\[
		f(x,y,z)=(-y+2,x+1,z+3).
	\]
	其正交部分在 \(xy\) 平面旋转 \(\pi/2\)，固定方向为 \(F=\mathbb R(0,0,1)\)。平移向量分为 \(t=(0,0,3)+(2,1,0)\)，故 \(a=(0,0,3)\)。为找出能够消去横向平移的新原点，求解 \((I-Q)c=(2,1,0)\) 且 \(c\in F^\perp\)，即
	\[
		c_x+c_y=2,\qquad -c_x+c_y=1,\qquad c_z=0.
	\]
	得到 \(c=(1/2,3/2,0)\)。在 \((u,v,w)=(x-1/2,y-3/2,z)\) 坐标中，变换成为 \((u,v,w)\mapsto(-v,u,w+3)\)。因此螺旋轴是 \(\{(1/2,3/2,z):z\in\mathbb R\}\)，每次沿轴前进 \(3\)，没有不动点。连续作用四次时旋转恢复原状，故 \(f^4(x,y,z)=(x,y,z+12)\)。
\end{example}

\subsection{反射生成与低维分类}
\label{g1:subsec:0103:04}

旋转和平移看起来与镜面反射差别很大，但它们都可以由有限次反射合成。先考虑过原点的超平面。若 \(u\ne0\)，关于 \(u^\perp\) 的\term[chaopingmianfanshe]{超平面反射}[hyperplane reflection]为
\[
	r_u(x)=x-2\frac{\langle x,u\rangle}{\langle u,u\rangle}u.
\]
向量沿 \(u^\perp\) 的分量保持不变，沿 \(u\) 的分量反号，因此 \(r_u\) 正交，\(r_u^2=I\)，\(\det r_u=-1\)。一般仿射超平面 \(c+u^\perp\) 的反射则是 \(x\mapsto c+r_u(x-c)\)。

\begin{theorem}[反射生成]
	\label{g1:thm:reflection-generation}
	\(n\) 维实内积空间的每个正交算子都是至多 \(n\) 个过原点的超平面反射的复合；恒等算子允许用零个反射表示。进一步，\(n\) 维 Euclid 空间的每个等距变换都是至多 \(n+1\) 个仿射超平面反射的复合。
\end{theorem}

\begin{proof}
	先对维数证明线性结论。零维情形只有恒等算子。设 \(n>0\)，取单位向量 \(e\)。若 \(Qe=e\)，则 \(e^\perp\) 被 \(Q\) 保持；对这一 \(n-1\) 维空间使用归纳假设，将 \(Q|_{e^\perp}\) 写成至多 \(n-1\) 个反射。每个反射的法向量都在 \(e^\perp\)，用同一法向量在整个空间定义反射，便在 \(e\) 方向上保持恒等，且其复合正是 \(Q\)。

	若 \(Qe\ne e\)，我们先用一次反射把 \(Qe\) 送回 \(e\)，之后就能沿用上一步的降维。令 \(u=Qe-e\)，则
	\[
		\lVert u\rVert^2=2(1-\langle Qe,e\rangle),
		\qquad
		\langle Qe,u\rangle=1-\langle Qe,e\rangle,
	\]
	所以 \(r_u(Qe)=Qe-u=e\)。算子 \(r_uQ\) 固定 \(e\)，按上一段是至多 \(n-1\) 个反射的复合。由 \(Q=r_u(r_uQ)\)，总数至多为 \(n\)。

	对于仿射结论，先用\cref{g1:thm:isometry-normal}把变换写成 \(c+y\mapsto c+Qy+a\)，其中 \(a\in F=\ker(I-Q)\)。若 \(a=0\)，只须将线性结论中的反射超平面全部平移到经过 \(c\)，使用至多 \(n\) 个反射。若 \(a\ne0\)，令 \(e=a/\lVert a\rVert\)，在中心为 \(c\) 的坐标里，关于超平面
	\[
		H_s=\{y:\langle y,e\rangle=s\}
	\]
	的反射记为 \(\rho_s\)，其公式为 \(\rho_s(y)=y+2(s-\langle y,e\rangle)e\)。直接代入得 \(\rho_{\lVert a\rVert/2}\rho_0(y)=y+a\)，因此沿 \(a\) 的平移由两个平行超平面反射组成。另一方面，只须在 \(F^\perp\) 上分解 \(Q\)，再将反射扩充为在 \(F\) 上恒等，故需要至多 \(n-\dim F\) 个反射。因 \(a\ne0\) 保证 \(\dim F\geq1\)，总数不超过 \(n-\dim F+2\leq n+1\)。
\end{proof}

下面的低维分类来自两个已经证明的约束：正交块的维数相加必须等于空间维数，而不能消去的平移只能位于 \(+1\) 特征空间。保向与反向的情形都由这两个约束确定。以下“取坐标”均指更换原点和标准正交基；同一类中的角度与位移仍可能给出不同变换，其参数关系将在\secref{g1:sec:0104}中进一步讨论。

\ProseParagraphBreak

在一维，\(Q\) 只有 \(+1\) 与 \(-1\)。保向变换为 \(x\mapsto x+t\)，其中 \(t=0\) 是恒等，\(t\ne0\) 是非零平移。反向变换为 \(x\mapsto-x+t\)，有唯一不动点 \(t/2\)，即关于该点的反射。

\ProseParagraphBreak

在二维，保向正交算子或者为 \(I\)，或者是非恒等旋转。如果 \(Q=I\)，便得到恒等或非零平移；否则 \(F=0\)，所有平移都可消去，得到绕唯一中心的旋转，角度可以包含 \(\pi\)。反向正交算子只能有一个 \(+1\) 块与一个 \(-1\) 块，所以 \(F\) 是一条直线。不可消去的平移为零时得到关于一条直线的反射，非零时得到沿该直线的滑移反射。这五种互不重叠的情形列于\cref{g1:tab:plane-isometries}。

\begin{table}[htbp]
	\centering
	\caption{平面等距变换的类型}
	\label{g1:tab:plane-isometries}
	\begin{tabularx}{\linewidth}{@{}l l X@{}}
		\toprule
		类型 & 正规形及参数 & 不动点集 \\
		\midrule
		恒等 & \((x,y)\mapsto(x,y)\) & 整个平面 \\
		非零平移 & \((x,y)\mapsto(x+h,y),\ h>0\) & 空集 \\
		旋转 & \(v\mapsto R_\theta v,\ 0<\theta\leq\pi\) & 唯一中心 \\
		直线反射 & \((x,y)\mapsto(x,-y)\) & 反射轴 \\
		滑移反射 & \((x,y)\mapsto(x+h,-y),\ h>0\) & 空集 \\
		\bottomrule
	\end{tabularx}
\end{table}

二维表格允许翻转坐标轴方向，所以把非零平移长度取为正，把旋转角归入 \((0,\pi]\)。若坚持使用事先固定的平面定向，则旋转角应记作模 \(2\pi\) 的带方向角，不能仅凭该表消去正负。

\ProseParagraphBreak

在三维，保向而不恒等的正交算子必有形式 \(\operatorname{diag}(R_\theta,1)\)，其中允许 \(\theta=\pi\)。确实，若有旋转块，余下的一维块必须为 \(+1\)；若全部为 \(\pm1\) 块，行列式为正迫使 \(-1\) 块的数目为零或二，后一种也就是转角 \(\pi\) 的情形。因此保向等距变换只有恒等、非零平移、绕轴旋转和非平凡螺旋运动。后两种在 \(0<\theta\leq\pi\) 时有唯一的轴，即最小位移点集；\(h=0\) 与 \(h\ne0\) 将它们分开。

\ProseParagraphBreak

三维反向正交算子则可写成 \(\operatorname{diag}(R_\theta,-1)\)，其中 \(0\leq\theta\leq\pi\)。当 \(\theta=0\) 时，\(F\) 是二维平面，得到平面反射，或者平面反射后沿该平面的非零平移；后一种也称滑移反射。当 \(0<\theta\leq\pi\) 时，\(F=0\)，变换有唯一不动点，全部平移都能消去。此时它由绕一条轴的旋转和关于过中心、垂直于该轴的平面反射组成，称为\term[xuanzhuanfanshe]{旋转反射}[rotatory reflection]。这两步在坐标中分别作用于 \((x,y)\) 与 \(z\)，故可交换次序。当 \(\theta=\pi\) 时，正规形成为 \(-I\)，即\term[zhongxinduichen]{中心对称}[central inversion]；此时中心唯一，却没有唯一的旋转轴，任意过中心的直线都可用于上述分解。

\begin{table}[htbp]
	\centering
	\caption{三维等距变换的类型}
	\label{g1:tab:space-isometries}
	\begin{tabularx}{\linewidth}{@{}l l X@{}}
		\toprule
		类型 & 正规形及参数 & 不动点集 \\
		\midrule
		恒等 & \(v\mapsto v\) & 整个空间 \\
		非零平移 & \((x,y,z)\mapsto(x,y,z+h),\ h\ne0\) & 空集 \\
		绕轴旋转 & \((v,z)\mapsto(R_\theta v,z)\) & 旋转轴 \\
		螺旋运动 & \((v,z)\mapsto(R_\theta v,z+h),\ h\ne0\) & 空集 \\
		平面反射 & \((x,y,z)\mapsto(x,y,-z)\) & 反射平面 \\
		滑移反射 & \((x,y,z)\mapsto(x+h,y,-z),\ h>0\) & 空集 \\
		旋转反射 & \((v,z)\mapsto(R_\theta v,-z),\ 0<\theta<\pi\) & 唯一中心 \\
		中心对称 & \(v\mapsto-v\) & 唯一中心 \\
		\bottomrule
	\end{tabularx}
\end{table}

\cref{g1:tab:space-isometries}中的绕轴旋转与螺旋运动取 \(0<\theta\leq\pi\)，这两行的 \(v\) 是二维向量；恒等与中心对称两行的 \(v\) 是三维向量。表中将旋转反射的端点 \(0\) 和 \(\pi\) 分别归入平面反射与中心对称，将螺旋运动的 \(h=0\) 归入绕轴旋转，从而各行表示互不重叠的类型。纯平移的最小位移点集是整个空间，滑移反射的是反射平面，螺旋运动的是螺旋轴；这些集合比仅仅问“有没有不动点”提供了更细的几何信息。

% !TeX root = ../../Book1.tex
% 纲要标题：### 1.4 群作用如何组织几何分类
% 纲要位置：几何学计划纲要.md，第 144 行。
% 正文状态：初稿；尚未正式审核。

% 纲要内容开始
% 
% * 轨道、稳定子群与等价关系
% * 共轭、坐标选择和分类参数
% * Euclid 、仿射与相似分类的比较
% * 矩阵族退化时不变量的变化
% 
% ---
% 
% 纲要内容结束

\section{群作用如何组织几何分类}
\label{g1:sec:0104}

一个椭圆经刚性移动仍有原来的长短轴，经统一缩放会改变轴长，经一般仿射变换却可以成为圆。因此，分类之前必须先回答：哪些改变算作同一个对象？前面已经出现了仿射变换、正交变换和等距变换，现在把这些具体变换组织起来，就能说明何谓分类参数，以及参数为什么随允许的变换而改变。

\subsection{轨道、稳定子群与等价关系}
\label{g1:subsec:0104:01}

把可逆变换按复合相乘时，复合满足结合律，有恒等变换，每个变换也有逆变换。抽象地，一个集合及其乘法若满足结合律、存在单位元并且每个元素都有逆元，就称为\term[qun]{群}[group]。一个含单位元且对乘法、取逆封闭的子集称为\term[ziqun]{子群}[subgroup]。下面保留这些运算性质，而把变换实际作用的对象另记为 \(X\)。

\begin{definition}[群作用]\label{g1:def:group-action}
	设 \(G\) 为群，单位元为 \(e\)，\(X\) 为集合。 \(G\) 在 \(X\) 上的\term[qunzuoyong]{群作用}[group action]是运算 \((g,x)\mapsto g\cdot x\)，满足 \(e\cdot x=x\) 与 \((gh)\cdot x=g\cdot(h\cdot x)\)。定义
	\[
		G\cdot x\coloneqq\{g\cdot x:g\in G\},\qquad
		G_x\coloneqq\{g\in G:g\cdot x=x\},
	\]
	分别称为 \(x\) 的\term[guidao]{轨道}[orbit]与\term[wendingziqun]{稳定子群}[stabilizer subgroup]。
	\symidx{\(G\cdot x\)：点 \(x\) 的轨道}
	\symidx{\(G_x\)：点 \(x\) 的稳定子群}
\end{definition}

群在这里首先可以理解为一组可逆变换：乘法是复合，单位元是不改变对象的变换，逆元撤销原变换。稳定子群中的变换允许移动周围空间，但必须把指定对象留在原处；例如，一个球心固定的球仍容许许多旋转。

\begin{proposition}[轨道与稳定子群]\label{g1:prop:orbit-stabilizer}
	对任意群 \(G\) 在集合 \(X\) 上的作用，关系 \(x\sim y\Longleftrightarrow y\in G\cdot x\) 是等价关系，\(G_x\) 是 \(G\) 的子群。记 \(gG_x=\{gh:h\in G_x\}\)，称为子群 \(G_x\) 的\term[zuopeiji]{左陪集}[left coset]。左陪集集合 \(G/G_x=\{gG_x:g\in G\}\) 与轨道 \(G\cdot x\) 之间有双射
	\[
		gG_x\longmapsto g\cdot x.
	\]
	若 \(y=h\cdot x\)，则 \(G_y=hG_xh^{-1}\)。
\end{proposition}
\begin{proof}
	单位元给出自反性，逆元给出对称性，乘法给出传递性。若 \(g,h\) 都固定 \(x\)，则 \(gh\) 固定 \(x\)，而 \(g^{-1}\) 也固定 \(x\)，故 \(G_x\) 为子群。两个左陪集 \(gG_x,hG_x\) 相等，当且仅当 \(h^{-1}g\in G_x\)，这又等价于 \(g\cdot x=h\cdot x\)。所以所给映射既良定义又单射；满射由轨道定义得到。最后，\(k\) 固定 \(h\cdot x\) 当且仅当 \(h^{-1}kh\) 固定 \(x\)，这就是稳定子群的共轭公式。
\end{proof}

上述双射使用的是陪集集合 \(G/G_x\)；当 \(G_x\) 正规时，这个集合才带有自然的商群结构。若 \(G\) 有限，每个陪集的元素数等于 \(|G_x|\)，因而 \(|G\cdot x|=|G|/|G_x|\)。几何变换群通常无限，上面的集合双射仍然成立。

\begin{example}[点与有序标架]\label{g1:exm:frame-orbits}
	在 \(\mathbb R^n\) 中，所有可逆仿射映射组成\term[fangshequn]{仿射群}[affine group]，记为 \(\operatorname{Aff}(n,\mathbb R)\)。用 \((P,b)\) 表示 \(x\mapsto Px+b\)，其中 \(P\in\operatorname{GL}(n,\mathbb R)\)，则
	\[
		(P,b)(S,c)=(PS,Pc+b),\qquad
		(P,b)^{-1}=(P^{-1},-P^{-1}b).
	\]
	这些公式同时验证了复合与求逆仍在其中。这里 \(\operatorname{GL}(n,\mathbb R)\) 表示所有可逆 \(n\) 阶实矩阵组成的\term[yibanxianxingqun]{一般线性群}[general linear group]。仿射群在点集上的作用只有一个轨道，因为平移可以连接任意两点；原点的稳定子群为 \(\{(P,0)\}\)。它在有序仿射标架上的作用也只有一个轨道，但一个标架的稳定子群只含单位元：由\cref{g1:prop:frame-map}，标架的像唯一确定仿射映射。
\end{example}

分类需要\term[bubianliang]{不变量}[invariant]，即在同一轨道上取相同值的量。如果某个不变量还满足“取值相同就处于同一轨道”，便称它为\term[wanquanbubianliang]{完全不变量}[complete invariant]。秩是二次型合同变换的不变量，但实数域上仅知道秩不够： \(x_1^2+x_2^2\) 与 \(x_1^2-x_2^2\) 都有秩 \(2\)，惯性却不同。\cref{g1:thm:inertia}中的三个惯性数合在一起才给出完整分类。

\subsection{共轭、坐标选择和分类参数}
\label{g1:subsec:0104:02}

几何变换本身也可以作为待分类对象。若用可逆仿射映射 \(h(x)=Px+c\) 把旧坐标送到新坐标，计算新坐标中的作用时便先用 \(h^{-1}\) 还原旧坐标，再施行 \(f\)，最后用 \(h\) 返回新坐标。因此，\(f(x)=Qx+b\) 在新坐标中变成 \(hfh^{-1}\)。直接代入逆映射得
\begin{equation}\label{g1:eq:affine-conjugation}
	hfh^{-1}(y)=PQP^{-1}y+Pb+(I-PQP^{-1})c.
\end{equation}
这称为\term[gonge]{共轭}[conjugation]。取 \(P\) 为正交矩阵，便得到等距坐标改变下的公式。共轭允许换观察位置和方向，但保留迭代关系，因为 \((hfh^{-1})^k=hf^kh^{-1}\)。

\begin{proposition}[等距变换的共轭参数]\label{g1:prop:isometry-conjugacy}
	设 \(f_1,f_2\) 是同一 \(n\) 维实 Euclid 空间的等距变换。分别选取原点，将它们按\cref{g1:thm:isometry-normal}写成 \(f_i(x)=Q_ix+t_i\)，其中 \(t_i\in F_i\coloneqq\ker(I-Q_i)\)。则二者在 Euclid 等距变换群中共轭，当且仅当 \(Q_1,Q_2\) 正交共轭且 \(\lVert t_1\rVert=\lVert t_2\rVert\)。因此，正交标准分块中的 \(+1,-1\) 重数、\(0<\theta<\pi\) 的旋转角多重集，连同非负数 \(\lVert t_i\rVert\)，构成完全不变量。
\end{proposition}
\begin{proof}
	分别选择原点，等价于先把两个映射分别作一次平移共轭；如果变换后的映射可以相互共轭，前后复合这些平移及其逆就得到原映射间的共轭，反之亦然。因此可以比较各自的正规形。

	若 \(h(x)=Px+c\) 为共轭它们的等距变换，则 \(P\) 正交，且由\eqref{g1:eq:affine-conjugation}有 \(Q_2=PQ_1P^{-1}\)。所以 \(P(F_1)=F_2\)， \(Pt_1\in F_2\)。另一方面，\((I-Q_2)c\) 属于 \(F_2^\perp\)：对 \(v\in F_2\)， \(\langle(I-Q_2)c,v\rangle=\langle c,(I-Q_2^*)v\rangle=0\)。比较平移部分的 \(F_2\) 分量便得 \(t_2=Pt_1\)，从而长度相同。

	反过来，设正交映射 \(P_0\) 满足 \(P_0Q_1P_0^{-1}=Q_2\)。向量 \(P_0t_1\) 与 \(t_2\) 都在 \(F_2\) 且等长。若它们为零，取 \(U=I\)；否则把各自除以公共长度，再分别扩充为 \(F_2\) 的标准正交基，把前一组基映到后一组，得到 \(F_2\) 上的正交映射 \(U\)。令它在 \(F_2^\perp\) 上为恒等。由于 \(Q_2\) 在 \(F_2\) 上为恒等且保持其正交补， \(UQ_2=Q_2U\)。于是 \(P=UP_0\) 仍共轭两个正交部分，并满足 \(Pt_1=t_2\)。取 \(h(x)=Px\) 即得 \(hf_1h^{-1}=f_2\)。正交部分的参数由\cref{g1:thm:orthogonal-blocks}确定。
\end{proof}

这里允许所有等距坐标改变，含改变定向的变换。例如在一个二维旋转平面中翻转一个基向量，会把 \(R_\theta\) 变为 \(R_{-\theta}\)，所以参数可选在 \((0,\pi)\)。若只允许保定向的共轭，分类还可能需要额外的方向信息。

\ProseParagraphBreak

还要区分共轭与双线性型的坐标变换。同一个线性算子的输入与输出都换基，若旧坐标 \(x=Sx'\)，矩阵变成 \(S^{-1}AS\)；同一个二次型 \(x^{\mathsf T}Ax\) 改用新坐标则变成 \(x'^{\mathsf T}S^{\mathsf T}ASx'\)。前者是矩阵相似，后者是合同。矩阵相似是线性算子的换基；下一小节的几何相似变换则指统一缩放再作等距变换，名称相近而所讨论的对象不同。对正交矩阵，由 \(S^{-1}=S^{\mathsf T}\) 可知，两种公式在这个特殊条件下重合。

\subsection{Euclid 、仿射与相似分类的比较}
\label{g1:subsec:0104:03}

把正交条件加到\cref{g1:exm:frame-orbits}的仿射变换中，就得到\term[oushiqun]{Euclid 群}[Euclidean group] \(E(n)\)。它的元素为 \(x\mapsto Qx+b\)，由\cref{g1:thm:isometry-affine}也恰好是全部保持距离的自映射。进一步允许统一比例缩放，得到\term[xiangsiqun]{相似群}[similarity group] \(\operatorname{Sim}(n)\)，元素为 \(x\mapsto rQx+b\)，其中 \(r>0\)、\(Q\) 正交。这些映射把所有距离乘以同一个 \(r\)；反过来，具有这一距离性质的映射除去比例和平移后为等距映射。

\begin{table}[htbp]
	\centering\small
	\caption{三种变换所保留的信息}\label{g1:tab:transformation-groups}
	\begin{tabularx}{\linewidth}{@{}l l X@{}}
		\toprule
		变换 & 坐标形式 & 必定保留的几何量或关系 \\
		\midrule
		Euclid 等距变换 & \(Qx+b\) & 距离、角度、仿射组合与平行关系 \\
		相似变换 & \(rQx+b\) & 角度、非零长度之比、仿射组合与平行关系 \\
		可逆仿射变换 & \(Px+b\) & 仿射组合、子空间维数、相交与平行关系 \\
		\bottomrule
	\end{tabularx}
\end{table}

表中的包含关系为 \(E(n)\subseteq\operatorname{Sim}(n)\subseteq\operatorname{Aff}(n,\mathbb R)\)。允许的群越大，同一轨道可能包含越多对象，因而分类中能保留的信息越少。角度的保持可由内积公式核对：相似变换把内积与两个长度的乘积同时乘以 \(r^2\)，所以它们的比值不变。

\begin{example}[非退化椭球的三种分类]\label{g1:exm:ellipsoid-classification}
	在 \(\mathbb R^n\)、\(n\ge1\) 中，设 \(A\) 为实对称正定矩阵，考虑点集
	\[
		\mathcal E(A,c)=\{x:(x-c)^{\mathsf T}A(x-c)=1\}.
	\]
	\cref{g1:thm:spectral}将它化为 \(\sum_i\lambda_i y_i^2=1\)，其中 \(\lambda_i>0\)，半轴长为 \(a_i=\lambda_i^{-1/2}\)。在 Euclid 等距变换下，计入重数的无序半轴长多重集 \(\{a_1,\ldots,a_n\}\) 是完全不变量；在相似变换下，只剩它们的比例，即有序排列后归一化为最短半轴为 \(1\) 的长度组；在可逆仿射变换下，所有这样的椭球属于同一类。
\end{example}
\begin{proof}
	先证明椭球点集唯一确定中心和对称矩阵。正定性及谱分解说明该点集非空而有界。若 \(c,d\) 都是它的对称中心，两次中心对称的复合就是平移 \(2(d-c)\)。点集若在非零平移下不变，任取一点并反复平移便得到无界序列，与有界性矛盾。因此中心唯一。中心移到零后，每个单位向量 \(u\) 上的正射线与椭球相交于 \(u/\sqrt{u^{\mathsf T}Au}\)。所以点集确定 \(u^{\mathsf T}Au\)，由齐次性确定所有二次型值，再由极化公式确定对称矩阵 \(A\)。

	等距映射必须把唯一中心映到唯一中心，故中心化后两个矩阵通过正交合同联系，由谱定理的唯一性，特征值及半轴长的多重集一致。反过来，若半轴长组相同，匹配相应的标准正交特征基，再恢复中心的平移，就给出所求等距映射。相似变换则先把全部半轴长乘以公共正数 \(r\)，再作这种匹配，故判据恰为比例一致。最后，在正交特征坐标中置 \(z_i=\sqrt{\lambda_i}y_i\)，即可通过可逆线性变换把任意椭球送到单位球面。
\end{proof}

例如 \(x^2/4+y^2=1\) 与 \(x^2/16+y^2/4=1\) 相似而不等距，因为半轴由 \((2,1)\) 变为 \((4,2)\)。它们又都与圆仿射等价，但都不与圆相似，因为长短轴之比为 \(2\)。这也说明“能化成同一个标准方程”必须连同所用的坐标变换类别一起陈述。

\subsection{矩阵族退化时不变量的变化}
\label{g1:subsec:0104:04}

分类一个固定对象之后，还可以问：数据连续变化时，分类参数怎样变化？矩阵每个元素都很接近，并不保证它们处于同一个轨道。以下把矩阵族的连续性理解为各矩阵元的连续性；所用估计都在有限维空间内。

\begin{proposition}[秩与惯性的局部变化]\label{g1:prop:rank-inertia-continuity}
	设 \(A_0\) 为实矩阵。充分接近 \(A_0\) 的同型实矩阵 \(A\) 均满足 \(\operatorname{rank}A\ge\operatorname{rank}A_0\)。若矩阵皆为实对称 \(n\) 阶矩阵，记正、负惯性指数为 \(p(A),q(A)\)，则在 \(A_0\) 的某个邻域内
	\[
		p(A)\ge p(A_0),\qquad q(A)\ge q(A_0).
	\]
	所以在秩固定的连续实对称矩阵族中，惯性局部不变；特别地，非退化型的惯性局部不变。
\end{proposition}
\begin{proof}
	若 \(r=\operatorname{rank}A_0>0\)，取一个非零 \(r\) 阶子式。其行列式是矩阵元的多项式，故在充分小的邻域内仍非零，这证明秩的结论；\(r=0\) 时结论直接成立。

	由\cref{g1:thm:inertia}，有维数 \(p(A_0)\) 的子空间 \(U_+\)，使 \(u^{\mathsf T}A_0u\) 在 \(U_+\setminus\{0\}\) 上为正。若 \(p(A_0)>0\)，它在 \(U_+\) 的单位球面上的最小值 \(m_+\) 为正，因为该球面紧且型连续。对 \(D=A-A_0\)，用矩阵元的 Euclid 范数 \(\lVert D\rVert_{\mathrm F}=(\sum_{ij}d_{ij}^2)^{1/2}\) 有
	\[
		|u^{\mathsf T}Du|
		\le\left(\sum_{ij}d_{ij}^2\right)^{1/2}
		\left(\sum_{ij}u_i^2u_j^2\right)^{1/2}
		=\lVert D\rVert_{\mathrm F}\,\lVert u\rVert^2.
	\]
	所以当 \(\lVert D\rVert_{\mathrm F}<m_+/2\) 时，\(A\) 在同一个 \(U_+\) 上仍正定，从而 \(p(A)\ge p(A_0)\)；\(p(A_0)=0\) 时无需估计。在负定子空间上对 \(-A\) 作相同论证，得到负指数不等式。若同时有 \(\operatorname{rank}A=\operatorname{rank}A_0\)，利用 \(p(A)+q(A)=\operatorname{rank}A\)，两项不等式只能同时取等号。
\end{proof}

序列收敛时，极限处的秩可以降低；邻近矩阵的秩则至少等于极限矩阵的秩。这一性质称为秩的下半连续性。在一个参数区间上，秩固定就使惯性处处局部不变，因而在整个区间上不变：否则，取一个惯性值的原像以及其他所有惯性值原像的并；局部常值性说明二者都是开集，它们将该区间分成互不相交的两个非空开集，违背区间的连通性。

\begin{example}[越过退化型]\label{g1:exm:degenerating-form}
	实对称矩阵族 \(A_t=\operatorname{diag}(1,t)\) 对应 \(x^2+ty^2\)。当 \(t>0\) 时惯性为 \((2,0,0)\)，当 \(t=0\) 时为 \((1,0,1)\)，当 \(t<0\) 时为 \((1,1,0)\)。每个正参数都可用 \(z=\sqrt t\,y\) 化为 \(x^2+z^2\)，但这个变换在 \(t=0\) 时不可逆。因此，同一个非退化合同类的极限可以落入另一个轨道；分类所用的坐标不能保证跨过退化参数仍有效。
\end{example}

\begin{example}[特征值不变而相似类型改变]\label{g1:exm:shear-degeneration}
	考虑线性变换族
	\[
		N_t=\begin{pmatrix}1&t\\0&1\end{pmatrix}.
	\]
	当 \(t\ne0\) 时，取 \(P_t=\operatorname{diag}(t,1)\)，便有 \(N_t=P_tN_1P_t^{-1}\)，所以这些变换都相似。但 \(N_0=I\) 不与 \(N_1\) 相似：\(\operatorname{rank}(N_t-I)\) 在 \(t\ne0\) 时为 \(1\)，在 \(t=0\) 时为 \(0\)，而共轭保持该秩。全部特征值始终为 \(1\)，所以特征值多重集对一般线性算子并不是完全不变量。几何上，非零剪切只逐点固定一条直线，极限恒等变换却逐点固定整个平面。
\end{example}

谱定理中的自伴条件于是有了另一个可见作用：它排除了上述剪切型的现象，使特征值多重集足以完成正交分类。对于变化的几何对象，应分别观察连续数值参数、离散分类数据和达到退化点时失效的构造；不能仅凭某一个熟悉的不变量没有变化，就判断对象的类型没有变化。


\begin{chaptersummary}
	仿射空间把点与向量分开：点差是向量，系数和为 \(1\) 的组合才自然给出点。标架使这些运算可计算，但交集、方向空间与仿射包络不依赖坐标。加入内积后，正交补产生投影，自伴算子产生正交谱分解；一般对称型经合同变换只保留正、负、零惯性数，而 Gram 行列式同时给出体积的平方与线性依赖的判据。

	\ProseParagraphBreak

	距离保持足以推出 \(f(x)=Qx+b\)。正交部分分解为固定方向、反向方向与旋转平面；平移中垂直于固定空间的分量可以通过换原点消去，沿固定空间的分量则决定最小位移。低维的旋转、反射、滑移与螺旋运动是这一分解的具体形式。分类还须指定变换群：同一对象在 Euclid 、相似与仿射作用下有不同不变量；连续族达到退化点时，秩和惯性等离散信息可能改变，原先的标准化坐标也可能失效。
\end{chaptersummary}

\BookExercises

\begin{exercise}\label{g1:ex:prescribed-affine-map}
	在 \(\mathbb R^2\) 中，求将 \((0,0),(1,0),(0,1)\) 分别送到 \((1,1),(3,1),(2,4)\) 的仿射映射。求 \(p=(2,1)\) 的像和它相对于像标架的重心坐标。该映射是否为等距变换或相似变换？
\end{exercise}
\begin{solution}
	原点的像确定平移 \(b=(1,1)\)，两个基向量的像差确定矩阵的两列，故
	\[
		f(x,y)=\begin{pmatrix}2&1\\0&3\end{pmatrix}\binom{x}{y}+\binom11,
		\qquad f(2,1)=(6,4).
	\]
	其行列式为 \(6\ne0\)，所以像点组仍为标架。点 \(p\) 在原标架下的重心坐标为 \((-2,2,1)\)，因为这些坐标的和为 \(1\)，后两个系数是其仿射坐标。由仿射组合的保持性，同一重心坐标给出 \(f(p)\)，直接代入也得到 \(-2(1,1)+2(3,1)+(2,4)=(6,4)\)。原来两个垂直单位向量的像分别为 \((2,0),(1,3)\)，内积为 \(2\ne0\)，所以该映射既不等距也不相似。坐标与重心关系的保持并不要求角度保持。
\end{solution}

\begin{exercise}\label{g1:ex:parallelogram-midpoints}
	在实仿射空间中，设 \(p,q,r\) 仿射无关。证明恰有一个点 \(s\)，使 \(p,r\) 的中点等于 \(q,s\) 的中点，写出其仿射表达式，并说明 \(p,q,r,s\) 按此顺序是一个平行四边形。若不要求这三个点仿射无关，唯一性与平行四边形的非退化性中哪项仍成立？
\end{exercise}
\begin{solution}
	以 \(p\) 为原点，设 \(u=q-p,v=r-p,w=s-p\)。中点等式为 \(v/2=(u+w)/2\)，故唯一可能为 \(w=v-u\)，且它确实满足等式。所以 \(s=p-q+r\)，此式的系数和为 \(1\)，是合法仿射组合。对边向量分别满足 \(q-p=r-s=u\) 与 \(r-q=s-p=v-u\)。若 \(p,q,r\) 仿射无关，则 \(u,v\) 无关，\(u,v-u\) 也无关，四边形因此处于一个平面且不退化。无关性去掉后，方程仍唯一确定 \(s\)，但可能得到共线甚至重合的点，故非退化性不再成立。
\end{solution}

\begin{exercise}\label{g1:ex:affine-intersection-parameter}
	在 \(\mathbb R^3\) 中，令
	\[
		F=\{(s,0,0):s\in\mathbb R\},\qquad
		G_t=\{(0,u,t):u\in\mathbb R\}.
	\]
	分别求 \(F\cap G_t\) 与 \(\operatorname{aff}(F\cup G_t)\)，并说明只知道两条直线方向的夹角，为何不能判断它们是否相交。
\end{exercise}
\begin{solution}
	方程 \((s,0,0)=(0,u,t)\) 强制 \(s=u=t=0\)。所以 \(t=0\) 时交集为原点，\(t\ne0\) 时交集为空。包络的方向为 \(\operatorname{span}(e_1,e_2,te_3)\)，故 \(t=0\) 时包络为 \(z=0\) 平面， \(t\ne0\) 时为整个 \(\mathbb R^3\)。两条直线的方向始终垂直，发生变化的是基点差是否属于方向空间之和，夹角没有记录这一信息。
\end{solution}

\begin{exercise}\label{g1:ex:distance-configurations}
	设 \(p_0,\ldots,p_m\) 与 \(q_0,\ldots,q_m\) 是 \(\mathbb R^n\) 中两组有序点，满足所有两两距离 \(\lVert p_i-p_j\rVert=\lVert q_i-q_j\rVert\)。证明存在 Euclid 等距变换把每个 \(p_i\) 送到 \(q_i\)。若第一组点的仿射包络为 \(\mathbb R^n\)，证明此变换唯一。
\end{exercise}
\begin{solution}
	令 \(v_i=p_i-p_0,w_i=q_i-q_0\)。极化公式给出
	\[
		\langle v_i,v_j\rangle
		=\tfrac12\bigl(\lVert v_i\rVert^2+\lVert v_j\rVert^2-\lVert v_i-v_j\rVert^2\bigr)
		=\langle w_i,w_j\rangle.
	\]
	尝试定义 \(L(\sum_i a_iv_i)=\sum_i a_iw_i\)。若 \(\sum_i a_iv_i=0\)，由上述内积相等可知 \(\lVert\sum_i a_iw_i\rVert^2=\sum_{ij}a_ia_j\langle v_i,v_j\rangle=0\)，所以对应像也为零，定义因此与表示无关。同一计算表明 \(L\) 保内积，并把 \(\operatorname{span}(v_i)\) 双射到 \(\operatorname{span}(w_i)\)。取该子空间的标准正交基及其像，分别补成全空间的标准正交基，将补上的基向量逐一匹配，便将 \(L\) 扩充为正交算子 \(Q\)。所求映射为 \(x\mapsto q_0+Q(x-p_0)\)。若原点组仿射张成全空间，则这些 \(v_i\) 张成 \(\mathbb R^n\)，已唯一确定 \(Q\)，而 \(p_0\) 的像又唯一确定平移，故等距变换唯一。
\end{solution}

% 供章末习题整合；题目由本节概念独立设计。

\begin{exercise}
	\label{g1:ex:weighted-projection}
	在 \(\mathbb R^3\) 上取内积
	\[
		\langle x,y\rangle=x_1y_1+4x_2y_2+2x_3y_3.
	\]
	令 \(x=(2,0,1)\)。求 \(x\) 在直线 \(W=\mathbb R(1,1,0)\) 上的正交投影及到 \(W\) 的距离，再求它在仿射平面 \(F=\{\,y:y_1+y_2+y_3=1\,\}\) 上的最近点及到 \(F\) 的距离。说明为何不能直接把通常点积下的法向量 \((1,1,1)\) 用于第二问。
\end{exercise}

\begin{solution}
	令 \(v=(1,1,0)\)。由 \(\langle x,v\rangle=2\)、\(\lVert v\rVert^2=5\)，投影为
	\[
		P_Wx=\frac25v=(2/5,2/5,0).
	\]
	投影距离公式给出
	\[
		d(x,W)^2=\lVert x\rVert^2-\lVert P_Wx\rVert^2
		=6-\frac45=\frac{26}{5}.
	\]

	对平面，令 \(\alpha(y)=y_1+y_2+y_3\)。需要的是满足 \(\alpha(y)=\langle n,y\rangle\) 的向量，故 \(n=\alpha^\sharp=(1,1/4,1/2)\)，且 \(\lVert n\rVert^2=7/4\)。平面方程即 \(\langle n,y\rangle=1\)，而 \(\langle n,x\rangle=3\)。由\cref{g1:exm:plane-projection}的公式，最近点为
	\[
		x-\frac{3-1}{7/4}n
		=\left(\frac67,-\frac27,\frac37\right),
	\]
	其坐标和等于 \(1\)。距离平方为
	\[
		\left(\frac87\right)^2\frac74=\frac{16}{7}.
	\]
	故两问距离分别为 \(\sqrt{26/5}\) 与 \(4/\sqrt7\)。向量 \((1,1,1)\) 在指定内积下对应的线性函数是 \(y_1+4y_2+2y_3\)，并非平面的定义函数；它甚至不与平面方向向量 \((1,-1,0)\) 正交，二者内积为 \(-3\)。
\end{solution}

\begin{exercise}
	\label{g1:ex:simultaneous-spectral}
	设 \(S,T\) 是同一个实有限维内积空间上的自伴算子。证明：存在同时由 \(S\)、\(T\) 的特征向量组成的标准正交基，当且仅当 \(ST=TS\)。再用
	\[
		S=\begin{pmatrix}1&0\\0&-1\end{pmatrix},\qquad
		T=\begin{pmatrix}0&1\\1&0\end{pmatrix}
	\]
	说明两个算子各自可以正交对角化，并不保证它们可以在同一组基中对角化。
\end{exercise}

\begin{solution}
	若有共同特征基，则两个矩阵在这组基下都为对角矩阵，故可交换。

	反过来，设 \(ST=TS\)。由\cref{g1:thm:spectral}，空间是 \(S\) 的各特征子空间 \(E_\lambda(S)\) 的正交直和。若 \(v\in E_\lambda(S)\)，则
	\[
		S(Tv)=T(Sv)=\lambda Tv,
	\]
	所以 \(Tv\in E_\lambda(S)\)。因此可以将 \(T\) 限制到每个 \(E_\lambda(S)\)，限制仍自伴：两边内积只是原内积的限制，自伴等式不变。再次应用谱定理，在每个 \(E_\lambda(S)\) 中取 \(T\) 的标准正交特征基。空间之间互相正交，合并这些基得到整个空间的标准正交基，其中每个向量又都是 \(S\) 的特征向量。

	题给矩阵均为实对称矩阵，但
	\[
		ST=\begin{pmatrix}0&1\\-1&0\end{pmatrix},\qquad
		TS=\begin{pmatrix}0&-1\\1&0\end{pmatrix}.
	\]
	两者不同，故不存在共同标准正交特征基。
\end{solution}

\begin{exercise}
	\label{g1:ex:inertia-family}
	对实参数 \(t\)，求二次型
	\[
		q_t(x,y,z)=x^2+2txy+y^2+z^2
	\]
	的惯性、根基与非退化参数范围。给出一组对所有 \(t\) 都适用的正交坐标，并说明惯性在哪些参数处发生变化。
\end{exercise}

\begin{solution}
	取 \(u=(x+y)/\sqrt2\)、\(v=(x-y)/\sqrt2\)、\(w=z\)，这是与 \(t\) 无关的正交换基，给出
	\[
		q_t=(1+t)u^2+(1-t)v^2+w^2.
	\]
	当 \(-1<t<1\) 时三个系数都正，惯性为 \((3,0,0)\)。当 \(t>1\) 或 \(t<-1\) 时恰有一个系数为负，惯性为 \((2,1,0)\)。当 \(t=1\) 或 \(t=-1\) 时恰有一个系数为零，其余两个为正，惯性为 \((2,0,1)\)。

	在 \(t\ne\pm1\) 时根基为 \(\{0\}\)。在 \(t=1\) 时根基是原坐标中的 \(\mathbb R(1,-1,0)\)；在 \(t=-1\) 时根基是 \(\mathbb R(1,1,0)\)。因此非退化参数恰为 \(\mathbb R\setminus\{-1,1\}\)。惯性在这两个参数处改变，对应的是一个特征值经过零；在三个互不相连的非退化参数区间内，它保持不变。
\end{solution}

\begin{exercise}
	\label{g1:ex:equidistant-four-points}
	设 Euclid 空间中四点 \(p_0,p_1,p_2,p_3\) 两两距离都是 \(1\)。不用预先画出正四面体，证明它们不能同处于一个仿射平面，并在 \(\mathbb R^3\) 中给出一组满足条件的坐标。提示：令 \(v_i=p_i-p_0\)，由距离恢复 \(\langle v_i,v_j\rangle\)。
\end{exercise}

\begin{solution}
	由 \(\lVert v_i\rVert=1\)，且 \(\lVert v_i-v_j\rVert=1\)，展开平方范数得
	\[
		\langle v_i,v_j\rangle
		=\frac{\lVert v_i\rVert^2+\lVert v_j\rVert^2-\lVert v_i-v_j\rVert^2}{2}
		=\frac12\qquad(i\ne j).
	\]
	故三条向量的 Gram 矩阵为
	\[
		\Gamma=\begin{pmatrix}1&1/2&1/2\\1/2&1&1/2\\1/2&1/2&1\end{pmatrix},
		\qquad \det\Gamma=\frac12>0.
	\]
	由\cref{g1:prop:gram-volume}，它们线性无关，四点的仿射张成空间因而是三维的，不能包含于一个仿射平面。

	在 \(\mathbb R^3\) 中可以取
	\[
		p_0=(0,0,0),\quad p_1=(1,0,0),\quad
		p_2=(1/2,\sqrt3/2,0),\quad
		p_3=(1/2,1/(2\sqrt3),\sqrt{2/3}).
	\]
	后三个位置向量的长度平方分别为 \(1\)，两两内积都是 \(1/2\)：其中 \(\langle p_2,p_3\rangle=1/4+1/4=1/2\)。于是它们与 \(p_0\) 的距离为 \(1\)，其余两两距离平方为 \(1+1-2(1/2)=1\)。这验证了全部六条距离条件。
\end{solution}

\begin{exercise}
	\label{g1:ex:positive-square-root}
	设 \(T\) 是实有限维内积空间 \(V\) 上的自伴算子，且 \(\langle Tv,v\rangle>0\) 对所有 \(v\ne0\) 成立。证明存在唯一的自伴正定算子 \(R\) 满足 \(R^2=T\)。说明若删除对 \(R\) 的正定要求，唯一性为何失败；再说明内积 \(B(u,v)=\langle Tu,v\rangle\) 与原内积之间有什么关系。
\end{exercise}

\begin{solution}
	由谱定理，\(V\) 是 \(T\) 的各特征子空间 \(E_\lambda(T)\) 的正交直和，且全部特征值 \(\lambda>0\)。在 \(E_\lambda(T)\) 上定义 \(R=\sqrt\lambda I\)，再按直和延拓。这个算子在标准正交特征基下是正对角矩阵，所以自伴、正定，且逐个特征子空间上都有 \(R^2=T\)。

	若 \(S\) 是另一个自伴正定平方根，则 \(ST=SS^2=S^3=S^2S=TS\)，故 \(S\) 保持每个 \(E_\lambda(T)\)。限制到这个子空间，\(S\) 仍自伴、正定，因而可由谱定理正交对角化。它的每个特征值 \(\mu\) 都满足 \(\mu>0\) 和 \(\mu^2=\lambda\)，所以只能等于 \(\sqrt\lambda\)。于是 \(S\) 在该子空间上等于 \(\sqrt\lambda I\)，即等于 \(R\)。各特征子空间张成 \(V\)，故 \(S=R\)。

	若 \(V\ne\{0\}\)，删除正定条件后，\(R\) 和 \(-R\) 是两个不同的自伴平方根；还可以在不同特征方向上选择不同正负号。最后，由 \(R\) 自伴及 \(R^2=T\)，
	\[
		B(u,v)=\langle R^2u,v\rangle=\langle Ru,Rv\rangle.
	\]
	因此可逆线性映射 \(R\) 把带内积 \(B\) 的 \(V\) 等距地映到带原内积的 \(V\)。它解释了为什么任意正定型都能经线性坐标伸缩变成平方和；该伸缩相对于原内积一般不保持长度。
\end{solution}


% 供第一章习题整合；以下均不承担正文证明责任。
\begin{exercise}
	\label{g1:ex:isometry-frame}
	设 \(f\colon\mathbb R^2\rightarrow\mathbb R^2\) 保持距离，并且
	\[
		f(0,0)=(1,2),\qquad f(1,0)=(1,3),\qquad f(0,1)=(0,2).
	\]
	求 \(f\) 的公式，判断它的类型，并求不动点。
\end{exercise}
\begin{solution}
	由\cref{g1:thm:isometry-affine}，\(f(x)=Qx+t\)，其中 \(t=(1,2)\)。减去 \(f(0,0)\) 后，两列分别为 \(Q(1,0)=(0,1)\)、\(Q(0,1)=(-1,0)\)，故
	\[
		f(x,y)=(-y+1,x+2).
	\]
	它的正交部分为旋转 \(\pi/2\)，固定方向空间为零，因而有唯一不动点。直接解 \(x=-y+1\)、\(y=x+2\)，得到 \((-1/2,3/2)\)。以此为原点便是逆时针 \(\pi/2\) 旋转。这里三个点的像决定公式，是因为其两个差向量构成一组基，不能只凭“给了三个条件”来判断唯一性。
\end{solution}

\begin{exercise}
	\label{g1:ex:space-glide}
	设 \(f\colon\mathbb R^3\rightarrow\mathbb R^3\) 为
	\[
		f(x,y,z)=(y+1,x-1,z+2).
	\]
	求其最小位移及取得最小位移的点集，说明它的几何类型，并求 \(f^2\)。
\end{exercise}
\begin{solution}
	正交部分为 \(Q(x,y,z)=(y,x,z)\)，所以
	\[
		F=\{(x,y,z):x=y\},\qquad F^\perp=\mathbb R(1,-1,0).
	\]
	平移向量分解为 \(t=(0,0,2)+(1,-1,0)\)，因此 \(a=(0,0,2)\)。取 \(c=(1/2,-1/2,0)\)，有 \((I-Q)c=(1,-1,0)\)。由\cref{g1:thm:isometry-normal}，最小位移为 \(2\)，在平面 \(c+F=\{x-y=1\}\) 上且仅在其上取得。\(Q\) 在 \(F\) 上恒等、在其法向上反号，故 \(f\) 是关于平面 \(x-y=1\) 的反射，再沿该平面平移 \((0,0,2)\)，属于反向的滑移反射。螺旋运动的非恒等正交部分保向，并且最小位移点集是一条直线，两者在此已可区分。最后直接计算
	\[
		f^2(x,y,z)=f(y+1,x-1,z+2)=(x,y,z+4).
	\]
\end{solution}

\begin{exercise}
	\label{g1:ex:two-reflections}
	在已固定标准定向的 Euclid 平面中，记 \(r_\alpha\) 为关于过原点且方向角为 \(\alpha\) 的直线的反射。
	\begin{enumerate}
		\item 求 \(r_\beta r_\alpha\)，并说明复合次序对结果的影响。
		\item 记 \(s_a\) 为关于直线 \(x=a\) 的反射。求 \(s_b s_a\)，并求 \(a\ne b\) 时的位移长度。
	\end{enumerate}
\end{exercise}
\begin{solution}
	1. 关于 \(x\) 轴的反射矩阵为 \(D=\operatorname{diag}(1,-1)\)。先旋转 \(-\alpha\) 使反射轴成为 \(x\) 轴，再反射，最后旋转回去，得到
	\[
		r_\alpha=R_\alpha D R_{-\alpha}
		=\begin{pmatrix}
			\cos2\alpha&\sin2\alpha\\
			\sin2\alpha&-\cos2\alpha
		\end{pmatrix}.
	\]
	由 \(DR_\gamma=R_{-\gamma}D\)，
	\[
		r_\beta r_\alpha
		=R_\beta D R_{\alpha-\beta}D R_{-\alpha}
		=R_{2(\beta-\alpha)}.
	\]
	这里右侧映射先作用，所以先关于方向角 \(\alpha\) 的直线反射，再关于方向角 \(\beta\) 的直线反射。交换次序会把转角改为 \(2(\alpha-\beta)\)，即得到逆旋转。反射轴的方向角只在模 \(\pi\) 意义下确定，但双倍角在模 \(2\pi\) 意义下不受此选择影响。

	2. 由于 \(s_a(x,y)=(2a-x,y)\)，有
	\[
		s_b s_a(x,y)=(2b-2a+x,y).
	\]
	因此复合为沿两条直线的公共法向平移 \(2(b-a)\)，位移长度为 \(2|b-a|\)，等于两条平行线距离的两倍。
\end{solution}

\begin{exercise}
	\label{g1:ex:minimal-reflections}
	设 \(Q\) 是 \(n\) 维实内积空间的正交算子。证明：把 \(Q\) 表示为过原点的超平面反射的复合，所需反射的最少个数等于 \(\operatorname{rank}(I-Q)\)。
\end{exercise}
\begin{solution}
	先证下界。任意同维线性算子 \(A,B\) 满足
	\[
		I-AB=(I-A)+A(I-B),
	\]
	故 \(\operatorname{rank}(I-AB)\leq\operatorname{rank}(I-A)+\operatorname{rank}(I-B)\)。若 \(r\) 为一个超平面反射，则 \(I-r\) 的像恰为法向直线，秩为 \(1\)。反复使用上述不等式，若 \(Q=r_1\cdots r_m\)，便得到 \(\operatorname{rank}(I-Q)\leq m\)。

	再证该下界能够达到。令 \(F=\ker(I-Q)\)，并设 \(d=\dim F^\perp\)。秩与核维数公式给出 \(d=\operatorname{rank}(I-Q)\)。空间 \(F^\perp\) 被 \(Q\) 保持，限制在其上的算子仍正交；由\cref{g1:thm:reflection-generation}，它是至多 \(d\) 个 \(F^\perp\) 内的超平面反射的复合。将每个反射按同一法向量扩充到整个空间，便在 \(F\) 上恒等，复合得到 \(Q\)。所需个数至多 \(d\)，与下界合起来恰为 \(d\)。\(d=0\) 时 \(Q=I\)，与零次反射的约定一致。
\end{solution}

\begin{exercise}
	\label{g1:ex:isometry-bounded-orbit}
	设 \(f\) 是 \(\mathbb R^n\) 的等距变换。对一点 \(x\)，其正向轨道指集合 \(\{f^k(x):k=0,1,2,\ldots\}\)。
	\begin{enumerate}
		\item 证明：\(f\) 有一条有界正向轨道，当且仅当 \(f\) 有不动点；在此情形，每条轨道都有界。
		\item 证明：存在正整数 \(m\) 使 \(f^m=\operatorname{id}\)，当且仅当 \(f\) 的不可消去平移分量为零，且正交部分各旋转块的角度除以 \(2\pi\) 都是有理数。
		\item 给出有界轨道不蕴含变换有限阶的具体例子。
	\end{enumerate}
\end{exercise}
\begin{solution}
	1. 按\cref{g1:thm:isometry-normal}，写 \(x=c+u+v\)，其中 \(u\in F\)、\(v\in F^\perp\)，则归纳得到
	\[
		f^k(x)=c+u+ka+Q^kv.
	\]
	其在 \(F\) 上相对于 \(c\) 的投影为 \(u+ka\)。若 \(a\ne0\)，该投影长度随 \(k\) 无界，所以轨道不可能有界。若 \(a=0\)，则 \(\lVert f^k(x)-(c+u)\rVert=\lVert v\rVert\)，每条轨道都有界；同一定理又说明 \(a=0\) 当且仅当 \(f\) 有不动点。

	2. 若 \(f^m=\operatorname{id}\)，则每条轨道都是有限集，由第1问知 \(a=0\)。以不动点为原点后 \(f=Q\)，且 \(Q^m=I\)。对一个旋转块，\(R_\theta^m=R_{m\theta}\)，故 \(m\theta\) 是 \(2\pi\) 的整数倍，\(\theta/(2\pi)\) 为有理数。反之，若 \(a=0\) 且全部旋转角满足此条件，取这些有理数分母的公倍数，再乘以 \(2\)，得到正整数 \(m\)。每个旋转块的 \(m\) 次方为 \(I\)，\(-1\) 块的 \(m\) 次方也为 \(1\)，而 \(+1\) 块始终不变，因此 \(Q^m=I\)，进而 \(f^m=\operatorname{id}\)。没有旋转块时直接取 \(m=2\) 即可。

	3. 平面上绕原点旋转角 \(2\pi\sqrt2\) 的变换，保持每个点到原点的距离，所以轨道都有界。若某个正整数 \(m\) 使其 \(m\) 次方为恒等，就会有 \(m\sqrt2\in\mathbb Z\)，与 \(\sqrt2\) 无理矛盾。因此它不是有限阶变换。
\end{solution}


\begin{exercise}\label{g1:ex:ellipsoid-stabilizer}
	设 \(A=\operatorname{diag}(\lambda_1,\ldots,\lambda_n)\) 为正定实对称矩阵。证明保持椭球 \(x^{\mathsf T}Ax=1\) 的 Euclid 等距变换恰为满足 \(QA=AQ\) 的正交变换。若不同特征值的重数为 \(m_1,\ldots,m_k\)，描述其稳定子群；比较所有半轴长互异与所有半轴长相等的情形。
\end{exercise}
\begin{solution}
	由\cref{g1:exm:ellipsoid-classification}的中心唯一性，保持点集的等距变换必须固定原点，所以为 \(x\mapsto Qx\)。该证明还说明中心化椭球的点集唯一决定二次型，故保持点集等价于 \(Q^{\mathsf T}AQ=A\)。由于 \(Q^{-1}=Q^{\mathsf T}\)，这等价于 \(AQ=QA\)。交换关系说明 \(Q\) 保持每个特征空间，反之在各特征空间上独立选择正交变换便与 \(A\) 交换。因此稳定子群同构于 \(O(m_1)\times\cdots\times O(m_k)\)。半轴互异时每个空间一维，只能独立翻转各轴，群共有 \(2^n\) 个元素；全部半轴相等时椭球为球面，稳定子群为整个 \(O(n)\)。
\end{solution}

\begin{exercise}\label{g1:ex:affine-versus-euclidean}
	在 \(\mathbb R^2\) 中考虑非零平移 \(\tau_v(x)=x+v\)。证明任意两个非零平移在仿射群中共轭，而在 Euclid 群中共轭当且仅当平移向量等长。再说明面积保持并不足以推出一个线性变换是等距变换。
\end{exercise}
\begin{solution}
	代入\eqref{g1:eq:affine-conjugation}中 \(Q=I,b=v\)，得到 \(h\tau_vh^{-1}=\tau_{Pv}\)。对任意非零 \(v,w\)，分别将它们扩充成一组基，将两组基匹配便得到可逆 \(P\) 满足 \(Pv=w\)，故仿射共轭成立。当 \(P\) 必须正交时，必要条件为 \(\lVert v\rVert=\lVert w\rVert\)；等长时匹配以归一化向量为首向量的标准正交基，又给出充分性。最后，剪切 \((x,y)\mapsto(x+y,y)\) 行列式为 \(1\)，保持平行四边形面积，却把单位向量 \((0,1)\) 变为长度 \(\sqrt2\) 的向量，故不等距。
\end{solution}

\begin{exercise}\label{g1:ex:orbit-closure}
	令 \(D=\operatorname{diag}(1,-1)\)。在实对称 \(2\) 阶矩阵空间中，构造 \(S_t\in\operatorname{GL}(2,\mathbb R)\)、\(t>0\)，使 \(S_t^{\mathsf T}DS_t\) 当 \(t\rightarrow0\) 时分别趋于 \(\operatorname{diag}(1,0)\)、\(\operatorname{diag}(0,-1)\) 与零矩阵。证明这同一个合同轨道中的矩阵序列不可能趋于单位矩阵。
\end{exercise}
\begin{solution}
	分别取 \(S_t=\operatorname{diag}(1,t)\)、\(\operatorname{diag}(t,1)\) 与 \(tI\)，所得矩阵为 \(\operatorname{diag}(1,-t^2)\)、\(\operatorname{diag}(t^2,-1)\) 与 \(t^2D\)，极限如题。三个极限的惯性都与 \(D\) 不同，故轨道不包含它们。若存在轨道内矩阵 \(B_j\rightarrow I\)，则当 \(j\) 充分大时 \(\lVert B_j-I\rVert_{\mathrm F}<1/2\)，从而对非零 \(x\) 有 \(x^{\mathsf T}B_jx\ge(1-1/2)\lVert x\rVert^2>0\)。这与轨道中每个矩阵的负惯性指数为 \(1\) 矛盾。因此轨道闭包增加的是某些退化型，不能任意增加别的非退化惯性类型。
\end{solution}

% !TeX root = ../../Book1.tex
% 纲要标题：## 第2章　射影空间、二次曲面与对偶几何
% 纲要位置：几何学计划纲要.md，第 153 行。
% 正文状态：初稿；尚未正式审核。

% 纲要内容开始
% 
% > 本章目标与证明范围：把射影构造、二次对象的分类和极性合成一条连续主线。必须有标准形的存在与唯一性论证，也有由对偶产生的新命题；不把圆锥曲线仅当作公式表。
% 
% 纲要内容结束

\chapter{射影空间、二次曲面与对偶几何}
\label{g1:ch:02}
\BookChapterNumber{g1:ch:02}

\begin{chapterabstract}
	射影空间把仿射位置与无穷远方向放在同一个线性代数模型中。齐次坐标由一维子空间产生，标架固定坐标自由度，交比刻画射影线上的有序四点。二次对象的分类则把上一章的谱定理与惯性定理用于非齐次方程，再经齐次化统一处理实、复射影二次曲面。非退化二次型使点与超平面对应，并把圆锥曲线的六点共线定理转为六条切线的共点定理。最后，通过乘积、二次单项式、子空间与关联条件的坐标模型，看到射影几何如何把一族对象本身组织成新的几何空间。
\end{chapterabstract}

\begin{learninggoals}
	\begin{itemize}
		\item 在齐次坐标与仿射图之间转换，计算子空间的交与张成空间，并用标架和交比确定射影变换。
		\item 对实二次多项式判断中心是否存在，写出 Euclid 、仿射标准形，并说明标准形中哪些数据唯一。
		\item 区分实、复射影二次方程的分类，求曲面所含线性子空间的最大维数，并识别实点集不能保留的信息。
		\item 从配对构造极性，计算切超平面与对偶二次曲面，利用参数与对偶证明圆锥曲线上的构造定理。
		\item 用齐次坐标描述乘积、二次映射和子空间族，检查定义不依赖代表，并从方程恢复原来的几何对象。
	\end{itemize}
\end{learninggoals}

在实平面中，圆 \(x^2+y^2=1\)、双曲线 \(x^2-y^2=1\) 与抛物线 \(y=x^2\) 的外形差别很大。若只允许平移和旋转，这些差别确实无法消除；若还允许改变“哪些点在无穷远处”，情况会怎样？把三个方程分别改写成
\[
	X_1^2+X_2^2=X_0^2,\qquad
	X_1^2-X_2^2=X_0^2,\qquad
	X_1^2=X_0X_2,
\]
并在 \(X_0\ne0\) 时令 \(x=X_1/X_0,y=X_2/X_0\)，原方程便重新出现。而当 \(X_0=0\) 时，三个方程分别没有实射影点、有两个点、有一个点。这提示我们：仿射图中看到的不同形状，有一部分来自二次方程，有一部分来自所选的无穷远直线。

\ProseParagraphBreak

齐次坐标将非零向量的共同缩放归为同一点，使仿射图与无穷远直线可以一起参与坐标变换。二次方程的标准形由此区分方程本身的分类数据与仿射图的选择。随后，二次型的配对 \(B(v,w)\) 会提供一个新的构造：固定一个点的代表 \(v\)，令 \(B(v,w)=0\)，便得到一个超平面。当这个构造可逆时，点的命题可以转成超平面的命题；在平面中，共线与共点由此相互对应。

\ProseParagraphBreak

本章实际使用的线性代数包括基、秩与核、对偶空间，以及实对称矩阵的两种标准形。其中\cref{g1:thm:spectral}要求实内积与正交换基，保留特征值；\cref{g1:thm:inertia}允许一般实换基，保留正、负、零三个维数。两种结果的区别会在分类中反复出现，所需条件将在调用处恢复。切线与切超平面由二次方程的一次变化直接计算。射影空间、标架、交比与对偶的补充阅读可参看 Gallier 的正式第二版\cite[第5.2--5.9, 5.11节]{Gallier2011Geometric}。

% !TeX root = ../../Book1.tex
% 纲要标题：### 2.1 从仿射图到射影空间
% 纲要位置：几何学计划纲要.md，第 157 行。
% 正文状态：初稿；尚未正式审核。

% 纲要内容开始
% 
% * 一维子空间、齐次坐标与标准仿射图
% * 射影子空间的关联与维数
% * 射影标架及射影变换的唯一确定
% * 射影线上的交比与置换对称
% 
% 纲要内容结束

\section{从仿射图到射影空间}
\label{g1:sec:0201}

在仿射平面中，两条不同直线的关系要分成相交与平行两种。可是从透视图看，平行铁轨似乎也交于远方。射影空间把这种现象变成一个精确的构造：把方向也作为点加入，使平面内任意两条不同直线都有唯一交点。这些方向点将由向量的一维子空间来定义。

\subsection{一维子空间、齐次坐标与标准仿射图}
\label{g1:subsec:0201:01}

本章的基域 \(K\) 为 \(\mathbb R\) 或 \(\mathbb C\)。上一章的仿射组合、基与维数计算也可对复向量空间逐字进行；长度、正负号与正交分类则仍须注明是实 Euclid 结构，不能随之自动推广。

\begin{definition}[射影空间]\label{g1:def:projective-space}
	设 \(K=\mathbb R\) 或 \(\mathbb C\)，\(V\) 是 \(n+1\) 维 \(K\) 向量空间，\(n\geq0\)。以 \(V\) 的一维向量子空间为元素的集合
	\[
		\mathbb P(V)\coloneqq\{Kv:v\in V\setminus\{0\}\}
	\]
	称为 \(V\) 的\term[sheyingkongjian]{射影空间}[projective space]，其射影维数定义为 \(n\)。记 \([v]=Kv\)，并记 \(\mathbb P^n(K)=\mathbb P(K^{n+1})\)。选定 \(V\) 的基后，若 \(v=(X_0,\ldots,X_n)\ne0\)，则 \([X_0:\cdots:X_n]\) 称为该点的\term[qicizuobiao]{齐次坐标}[homogeneous coordinates]。
\end{definition}
\symidx{\(\mathbb P(V)\)：向量空间的一维子空间所成的射影空间}

两个非零向量表示同一点，当且仅当相差一个非零倍数：
\[
	[X_0:\cdots:X_n]=[\lambda X_0:\cdots:\lambda X_n]\qquad(\lambda\in K^\times).
\]
零向量不表示任何射影点。在实数情形，\(v\) 与 \(-v\) 也表示同一点，所以这里取的是过原点的整条直线，并非带方向的射线。\(\mathbb P^0(K)\) 只有一个点，而 \(\mathbb P^1(K)\) 已有无穷多个点。

\begin{proposition}[标准仿射图]\label{g1:prop:projective-charts}
	设 \(K=\mathbb R\) 或 \(\mathbb C\)。对 \(0\leq i\leq n\)，令
	\[
		U_i=\{[X_0:\cdots:X_n]:X_i\ne0\}.
	\]
	映射 \(U_i\rightarrow K^n\) 将一点送到除去第 \(i\) 项后的各个比值 \(X_j/X_i\)，是双射。这些集合覆盖 \(\mathbb P^n(K)\)。在 \(U_i\cap U_k\) 上，若 \(x_j=X_j/X_i\)，并补记 \(x_i=1\)，则第 \(k\) 张图的坐标为 \(y_j=x_j/x_k\)，其中 \(j\ne k\)。
\end{proposition}
\begin{proof}
	同乘非零数不改变任何比值，故映射良定义。反过来，把第 \(i\) 项取为 \(1\)，其余项依次取为所给的 \(n\) 个数，便得到唯一射影点。每个非零向量至少有一个非零分量，所以各图覆盖整个空间。在重叠处 \(x_k\ne0\)，直接相除给出 \(X_j/X_k=(X_j/X_i)/(X_k/X_i)\)，即所述换图公式。
\end{proof}

称 \(U_i\) 为\term[fangshetu]{仿射图}[affine chart]，因为选定归一化 \(X_i=1\) 后它具有通常的仿射坐标。不同图的换图映射一般不是仿射映射。例如在射影线上，\(t=X_0/X_1\) 与 \(s=X_1/X_0\) 满足 \(s=1/t\)。图的用途是让我们在一块区域内使用熟悉的坐标，并明确知道走到另一块区域时如何转换。

\begin{example}[一条射影线中的无穷远点]\label{g1:exm:projective-line}
	在 \(\mathbb P^1(K)\) 中，\(U_1\) 的点唯一写成 \([t:1]\)，余下只有 \([1:0]\)，记作 \(\infty\)。因此作为集合有 \(\mathbb P^1(K)=K\cup\{\infty\}\)。这里的记号不允许直接参与普通实数或复数运算；例如涉及 \(\infty\) 的分式应改用齐次坐标计算。在实数情形，\cref{g1:fig:projective-line-chart}用穿过原点的直线表示这些点。
\end{example}

\begin{figure}[htbp]
	\centering
	\input{Figures/Chapter02-ProjectiveLine}
	\caption[射影线与一张仿射图]{射影线与一张仿射图。与 \(X_1=1\) 相交的每条过原点直线给出一个有限坐标；水平直线 \(X_1=0\) 给出唯一遗漏的射影点 \([1:0]\)。}
	\label{g1:fig:projective-line-chart}
\end{figure}

对 \(\mathbb P^n(K)\)，其中 \(n\geq1\)，选取 \(U_0\)，便把 \(K^n\) 放入射影空间：
\[
	(x_1,\ldots,x_n)\longmapsto[1:x_1:\cdots:x_n].
\]
补集 \(H_\infty=\{X_0=0\}\cong\mathbb P^{n-1}(K)\) 称为这张图的\term[wuqiongyuanchaopingmian]{无穷远超平面}[hyperplane at infinity]。它取决于仿射图的选择，并不是射影空间中预先指定的特殊部分。当 \(n=1\) 时它是一个点；当 \(n=2\) 时是一条射影直线。

\subsection{射影子空间的关联与维数}
\label{g1:subsec:0201:02}

由线性方程定义子空间的办法仍然适用，不过现在方程必须齐次，才能使“是否满足方程”不依赖代表向量。

\begin{definition}[射影子空间]\label{g1:def:projective-subspace}
	设 \(K=\mathbb R\) 或 \(\mathbb C\)，\(V\) 为有限维 \(K\) 向量空间。若 \(W\subseteq V\) 为非零向量子空间，称 \(\mathbb P(W)\subseteq\mathbb P(V)\) 为\term[sheyingzikongjian]{射影子空间}[projective subspace]，维数为 \(\dim_K W-1\)。维数为 \(1\) 的称为\term[sheyingzhixian]{射影直线}[projective line]；余维为 \(1\) 的称为\term[sheyingchaopingmian]{射影超平面}[projective hyperplane]。为统一相交公式，也把空集视为空射影子空间，约定 \(\mathbb P(0)=\varnothing\)，\(\dim\varnothing=-1\)。
\end{definition}

\begin{proposition}[交与张成空间]\label{g1:prop:projective-dimension}
	设 \(K=\mathbb R\) 或 \(\mathbb C\)，\(V\) 为有限维 \(K\) 向量空间，\(U,W\subseteq V\) 为向量子空间，记 \(E=\mathbb P(U)\)、\(F=\mathbb P(W)\)，并约定空射影子空间的维数为 \(-1\)。则 \(\mathbb P(U)\cap\mathbb P(W)=\mathbb P(U\cap W)\)，包含二者的最小射影子空间为 \(\mathbb P(U+W)\)，称为它们的张成空间。若用 \(\langle E,F\rangle\) 表示两个射影子空间的张成空间，则
	\[
		\dim\langle E,F\rangle+\dim(E\cap F)=\dim E+\dim F.
	\]
	特别地，在 \(\mathbb P^n(K)\) 中，\(\dim E+\dim F\geq n\) 蕴含 \(E\cap F\ne\varnothing\)。
\end{proposition}
\begin{proof}
	一个一维空间同时包含于 \(U,W\)，等价于它包含于 \(U\cap W\)，故得交集公式。任何包含 \(\mathbb P(U),\mathbb P(W)\) 的射影子空间，其对应向量空间包含 \(U,W\)，因而包含 \(U+W\)；后者本身也满足包含要求，故为最小者。上一章证明\cref{g1:prop:affine-intersection}时已经通过扩充交集的基得到 \(\dim(U+W)+\dim(U\cap W)=\dim U+\dim W\)。四个维数各减 \(1\)，即得到此处公式，零空间也符合约定。若交集为空，公式给出 \(\dim\langle E,F\rangle=\dim E+\dim F+1\leq n\)，与最后的假设矛盾。
\end{proof}

所以射影平面中两条不同直线恰交于一点：必有交点，而若交空间维数至少为 \(1\)，两条直线就相同。但射影三维空间仍有不相交的直线。例如 \(\mathbb P(\operatorname{span}(e_0,e_1))\) 与 \(\mathbb P(\operatorname{span}(e_2,e_3))\) 不交。加入无穷远点消除了平面内平行线的例外，没有消除高维空间中的异面直线。

\begin{proposition}[仿射子空间的射影补全]\label{g1:prop:affine-projective-completion}
	设 \(K=\mathbb R\) 或 \(\mathbb C\)，\(n\geq1\)，\(F=p+L\subseteq K^n\) 为非空仿射子空间，\(\dim L=d\)。用 \(x\mapsto[1:x]\) 把 \(K^n\) 识别为 \(U_0=\{X_0\ne0\}\subseteq\mathbb P^n(K)\)，并记 \(H_\infty=\{X_0=0\}\)。在此嵌入下，包含 \(F\) 的最小射影子空间为
	\[
		\widehat F=\mathbb P\bigl(K(1,p)+(0,L)\bigr).
	\]
	它具有维数 \(d\)，与 \(U_0\) 的交为 \(F\)，与 \(H_\infty\) 的交为 \(\mathbb P(0,L)\)。
\end{proposition}
\begin{proof}
	所写向量空间包含所有 \((1,p+u)=(1,p)+(0,u)\)，其中 \(u\in L\)。反过来，任何包含这些向量所张成各条直线的向量空间都含 \((1,p)\)，再作差便含全部 \((0,u)\)，故包含所写空间。由于 \((1,p)\notin(0,L)\)，其向量维数为 \(d+1\)。一般向量为 \((a,ap+u)\)：若 \(a\ne0\)，归一化得到 \([1:p+u/a]\)；若 \(a=0\)，得到 \([0:u]\)。这分别给出两种交集。
\end{proof}

\begin{example}[平行线的公共方向点]\label{g1:exm:parallel-infinity}
	在实仿射平面中，两条直线 \(y=0\)、\(y=1\) 的射影补全分别为 \(X_2=0\)、\(X_2=X_0\)。联立可得唯一交点 \([0:1:0]\)，它代表水平方向。一般地，不同平行直线 \(p+Ku\)、\(q+Ku\) 的射影补全都含 \([0:u]\)，而有限部分不交，故交点恰是该方向点。
\end{example}

\subsection{射影标架及射影变换的唯一确定}
\label{g1:subsec:0201:03}

上一章的仿射标架由 \(n+1\) 个点构成，它们唯一确定仿射映射。在射影空间中，每个代表向量还允许独立缩放，所以只指定 \(n+1\) 个独立的点不够；还需要一个点来协调这些缩放。

\begin{definition}[射影变换与射影标架]\label{g1:def:projective-frame}
	设 \(K=\mathbb R\) 或 \(\mathbb C\)，\(V,W\) 是维数相同的有限维 \(K\) 向量空间。可逆线性映射 \(T:V\rightarrow W\) 诱导的双射 \([v]\mapsto[Tv]\) 称为\term[sheyingbianhuan]{射影变换}[projective transformation]，记为 \(\mathbb P(T)\)。在 \(\mathbb P^n(K)\) 中，设 \(n\geq1\)。若一个有序点组的 \(n+2\) 个点 \((p_0,\ldots,p_n,p_*)\) 中，任意 \(n+1\) 个点的代表向量都线性无关，就称这个点组为\term[sheyingbiaojia]{射影标架}[projective frame]。
\end{definition}

这里的射影变换由可逆线性映射诱导。复共轭也保持直线，却涉及基域的自同构，属于不同的变换类别。射影变换把 \(\mathbb P(U)\) 送到 \(\mathbb P(TU)\)，因此保持子空间维数、相交与包含关系。

\begin{theorem}[标架决定射影变换]\label{g1:thm:projective-frame}
	设 \(K=\mathbb R\) 或 \(\mathbb C\)，\(n\geq1\)。任意两个 \(n\) 维射影空间的有序射影标架之间，存在唯一按顺序匹配各点的射影变换。
\end{theorem}
\begin{proof}
	先固定基 \(e_0,\ldots,e_n\)，称 \(([e_0],\ldots,[e_n],[e_0+\cdots+e_n])\) 为标准标架。给定标架，任选 \(p_i\) 的非零代表 \(v_i\)；它们组成基。将 \(p_*\) 的代表写成 \(w=\sum_i a_i v_i\)。每个 \(a_i\ne0\)，否则 \(w\) 与除去 \(v_i\) 的其余 \(n\) 个向量线性相关，违反标架条件。因此线性映射 \(T(e_i)=a_i v_i\) 可逆，且 \(T(\sum_i e_i)=w\)，把标准标架送到给定标架。对另一个标架作同样构造，再将相应映射与前一个映射的逆复合，就得到存在性。

	为证唯一性，只需证明固定标准标架各点的射影变换恒等。其线性代表 \(S\) 固定各坐标直线，故 \(Se_i=c_i e_i\)，其中 \(c_i\ne0\)。又 \([\sum_i c_i e_i]=[\sum_i e_i]\)，由基展开唯一性知所有 \(c_i\) 相同。于是 \(S=cI\)，诱导恒等射影变换。
\end{proof}

同一证明还说明，两个可逆线性映射诱导同一射影变换，当且仅当相差非零标量。因此 \(\mathbb P^n(K)\) 上的射影变换群记为
\[
	\operatorname{PGL}(n+1,K)\coloneqq\operatorname{GL}(n+1,K)/\{cI:c\in K^\times\},
\]
称为\term[sheyingyibanxianxingqun]{射影一般线性群}[projective general linear group]。商掉标量矩阵，把诱导同一射影变换的线性映射归为一类。

\begin{example}[平面标架与分式坐标]\label{g1:exm:projective-frame-map}
	在 \(\mathbb P^2(K)\) 中，取前三点 \([e_0],[e_1],[e_2]\)，再取 \([2:3:1]\)。三个系数非零，所以构成标架。从标准标架到它的变换由 \(\operatorname{diag}(2,3,1)\) 表示。在源、目标都采用 \(X_0\ne0\) 的图时，一般矩阵 \(M=(m_{ij})\) 所诱导的坐标公式为
	\[
		x_i'=
		\frac{m_{i0}+\sum_{j=1}^n m_{ij}x_j}{m_{00}+\sum_{j=1}^n m_{0j}x_j},\qquad i=1,\ldots,n,
	\]
	仅在分母非零的部分成立。分母为零时，射影变换仍有定义，只是像落到目标图的无穷远超平面。例如 \([X_0:X_1:X_2]\mapsto[X_1:X_0:X_2]\) 在 \(x_1\ne0\) 处写成 \((x_1,x_2)\mapsto(1/x_1,x_2/x_1)\)，而 \(x_1=0\) 上的有限点被送到无穷远处。
\end{example}

射影线上的标架就是三个两两不同的点。因此三个指定的不同点可以唯一移到 \(\infty,0,1\)；第四个点在这个归一化之后留下的坐标，就是我们下一步要构造的不变量。

\subsection{射影线上的交比与置换对称}
\label{g1:subsec:0201:04}

线段长之比依赖仿射图，换成分式坐标后通常改变。用四个点形成的适当比值却能把这些变化全部抵消。行列式定义能同时处理有限点与无穷远点。

\begin{definition}[交比]\label{g1:def:cross-ratio}
	设 \(K=\mathbb R\) 或 \(\mathbb C\)，\(A,B,C,D\in\mathbb P^1(K)\) 两两不同，分别取非零代表 \(a,b,c,d\in K^2\)。记 \(\Delta(u,v)=\det(u,v)\)，定义\term[jiaobi]{交比}[cross-ratio]为
	\begin{equation}\label{g1:eq:cross-ratio}
		(A,B;C,D)\coloneqq
		\frac{\Delta(a,c)\Delta(b,d)}{\Delta(a,d)\Delta(b,c)}.
	\end{equation}
\end{definition}

两点不同恰好意味着其代表线性无关，故上述分母非零。每个代表向量的缩放因子在分子分母各出现一次，因而抵消；改变 \(K^2\) 的基又使每个行列式乘以同一非零数，总商仍不变。因此定义也适用于任意二维向量空间的射影线。采用有限坐标 \(A=[a:1]\) 等时，公式成为
\[
	(A,B;C,D)=\frac{(a-c)(b-d)}{(a-d)(b-c)}.
\]
不同书籍有不同的排列约定，比较公式时应先检查哪三个点被归一化为 \(\infty,0,1\)。本书约定满足
\begin{equation}\label{g1:eq:cross-ratio-normalization}
	(\infty,0;1,t)=t.
\end{equation}

\begin{proposition}[有序四点的完全不变量]\label{g1:prop:cross-ratio-orbits}
	设 \(K=\mathbb R\) 或 \(\mathbb C\)，\((A,B,C,D)\) 与 \((A',B',C',D')\) 为 \(\mathbb P^1(K)\) 中两组各自两两不同的有序四点。\(\mathbb P^1(K)\) 的射影自同构保持交比；两组四点可由这样的自同构逐点匹配，当且仅当交比相等，且这个自同构若存在则唯一。交比的取值恰为 \(K\setminus\{0,1\}\)。
\end{proposition}
\begin{proof}
	若变换由 \(T\in\operatorname{GL}(2,K)\) 诱导，则 \(\Delta(Tu,Tv)=\det(T)\Delta(u,v)\)，四个因子抵消，证明不变性。由\cref{g1:thm:projective-frame}，存在唯一变换把前三点变为 \(\infty,0,1\)。第四点不同于前三点，故变为唯一的 \(t\in K\setminus\{0,1\}\)；直接以 \((1,0),(0,1),(1,1),(t,1)\) 代入定义，得到交比等于 \(t\)。于是交比相等的两组四点有同一归一化形式，复合归一化变换便逐点匹配；唯一性仍由前三点决定。反过来，任取 \(t\ne0,1\)，这组标准四点实现所需值。
\end{proof}

\begin{proposition}[置换产生的六种表达式]\label{g1:prop:cross-ratio-permutations}
	设 \(K=\mathbb R\) 或 \(\mathbb C\)，\(A,B,C,D\in\mathbb P^1(K)\) 两两不同，记 \(\lambda=(A,B;C,D)\)。任意重排四点得到的值，都属于下表的六种表达式；表中每种表达式都由所列重排实现，但数值可能重复。


\begin{table}[htbp]
	\centering
	\caption{交比与四点的次序}
	\label{g1:tab:cross-ratio-permutations}
	\begin{tabular}{cccc}
		\toprule
		排列&交比&排列&交比\\
		\midrule
		\((A,B;C,D)\)&\(\lambda\)&\((A,C;D,B)\)&\(1/(1-\lambda)\)\\
		\((A,B;D,C)\)&\(1/\lambda\)&\((A,D;B,C)\)&\((\lambda-1)/\lambda\)\\
		\((A,C;B,D)\)&\(1-\lambda\)&\((A,D;C,B)\)&\(\lambda/(\lambda-1)\)\\
		\bottomrule
	\end{tabular}
\end{table}
\end{proposition}

\begin{proof}
	由不变性可把四点设成 \(\infty,0,1,\lambda\)，代入行列式公式即逐项得到表中各值。还需说明其余排列不会增加新值。交换第一、二个点或第三、四个点，都把当前交比 \(t\) 变为 \(1/t\)；交换第二、三个点则把它变为 \(1-t\)，后者也可在归一化四点上直接计算。任何排列都可通过相邻交换得到：依次把所需的首项移到第一位，再把第二项移到第二位，继续即可。表中的六个分式对 \(t\mapsto1/t\) 与 \(t\mapsto1-t\) 封闭，所以没有其他值。
\end{proof}

\begin{example}[调和四点与次序的影响]\label{g1:exm:harmonic-quadruple}
	交比为 \(-1\) 的有序四点称为\term[tiaohe-sidian]{调和四点}[harmonic quadruple]。例如 \((\infty,0;1,-1)=-1\)，两个有限点 \(1,-1\) 关于 \(0\) 对称。一般地，在实仿射线上 \((\infty,b;c,d)=(b-d)/(b-c)\)，所以等于 \(-1\) 恰好意味着 \(b=(c+d)/2\)。“中点”本身依赖所选无穷远点，而四点构成的射影关系不依赖图。对于 \(\lambda=-1\)，\cref{g1:tab:cross-ratio-permutations}的六式只有 \(-1,2,1/2\) 三个不同值，说明交换次序可以保留同一无序点集，却改变交比。
\end{example}

于是，标架把位置选择中的自由度固定下来，交比记录仍未消除的参数。这个“先归一化，再读出不变量”的步骤将在圆锥曲线的参数与六点构造中再次使用。

% !TeX root = ../../Book1.tex
% 纲要标题：### 2.2 仿射和 Euclid 二次对象的分类
% 纲要位置：几何学计划纲要.md，第 164 行。
% 正文状态：初稿；尚未正式审核。

% 纲要内容开始
% 
% * 二次项、线性项与中心判据
% * 平移消去线性项及其障碍
% * 二次多项式的 Euclid 与仿射标准形
% * 退化参数、直纹及实点集
% 
% 纲要内容结束

\section{仿射和 Euclid 二次对象的分类}
\label{g1:sec:0202}

平面上的椭圆、双曲线和抛物线都由二次方程给出，但它们的差别不能只从二次项的正负号读出。例如 \(x^2+y^2=1\) 与 \(x^2+y^2=-1\) 的二次项相同，后一个方程却没有实解；\(x^2-y=0\) 中的线性项则决定了一条抛物线。要分类这些对象，必须同时处理二次项、线性项和常数项，并说明坐标变换是否需要保持长度。

\subsection{二次项、线性项与中心判据}
\label{g1:subsec:0202:01}

在 \(\mathbb R^n\) 上，每个次数不超过二的实多项式都唯一写成
\begin{equation}
	\label{g1:eq:affine-quadratic-polynomial}
	F(x)=x^{\mathsf T}Ax+2b^{\mathsf T}x+c,
	\qquad A=A^{\mathsf T}.
\end{equation}
交叉项均分给 \(a_{ij}\) 与 \(a_{ji}\)，这就确定了对称矩阵 \(A\)；在线性项前写 \(2\)，是为了使配方时的公式较简洁。当 \(A\ne0\) 时，称 \(F=0\) 为一个\term[fangsheerciqumian]{仿射二次曲面}[affine quadric]的方程。这里“曲面”泛指任意维数的二次对象；在仿射平面中也称为二次曲线。

\ProseParagraphBreak

本节把方程本身保留下来：两个方程在允许的坐标变换下等价，是指相应多项式经过该变换后相差一个非零常数倍。仅把实解看成点集会丢掉信息。例如 \(x^2+y^2+1=0\) 与 \(x^2+1=0\) 都给出空集，但二次项的秩不同，不能经可逆仿射变换互相转化。为看清常数倍究竟消去了什么，我们先分类固定的多项式，再把整体非零倍数也纳入等价关系。

\begin{definition}[二次多项式的中心]
	\label{g1:def:quadric-center}
	对次数不超过二的实多项式 \(F\colon\mathbb R^n\rightarrow\mathbb R\)，若点 \(a\) 满足
	\[
		F(a+y)=F(a-y)\qquad(y\in\mathbb R^n),
	\]
	则称 \(a\) 为 \(F\) 的\term[erciduoxiangshizhongxin]{中心}[center of a quadratic polynomial]。中心存在时称该多项式有中心。
\end{definition}

这个条件使每条等值集都关于 \(a\) 中心对称，因而当然使零点集中心对称。反过来，仅检查某一个点集的对称性未必能恢复多项式条件：空集关于每一点都对称，而 \(x^2+y^2+1\) 作为多项式只有原点这个中心。

\begin{proposition}[中心判据]
	\label{g1:prop:quadric-center}
	设 \(n\) 为非负整数，\(A\) 为实对称 \(n\) 阶矩阵，\(b,a\in\mathbb R^n\)、\(c\in\mathbb R\)，并将实多项式 \(F(x)=x^{\mathsf T}Ax+2b^{\mathsf T}x+c\) 写成式\eqref{g1:eq:affine-quadratic-polynomial}的形式。点 \(a\) 是 \(F\) 的中心，当且仅当
	\[
		Aa=-b.
	\]
	中心存在当且仅当 \(b\in\operatorname{im}A\)，等价地，\(b\) 与 \(\ker A\) 正交。若 \(a_0\) 是一个中心，则全部中心构成仿射子空间 \(a_0+\ker A\)，并且 \(F\) 在这个子空间上取同一个值。
\end{proposition}

\begin{proof}
	平移原点后直接展开，得到
	\begin{equation}
		\label{g1:eq:quadratic-translation}
		F(a+y)=y^{\mathsf T}Ay+2(Aa+b)^{\mathsf T}y+F(a).
	\end{equation}
	把 \(y\) 换成 \(-y\)，两式相减为 \(4(Aa+b)^{\mathsf T}y\)。它对所有 \(y\) 为零，当且仅当 \(Aa+b=0\)。这个线性方程有解恰好意味着 \(-b\in\operatorname{im}A\)，且有解时全部解为 \(a_0+\ker A\)。

	还需说明正交条件。因为 \(A\) 对称，若 \(v\in\ker A\)，则 \(\langle Ax,v\rangle=\langle x,Av\rangle=0\)。所以 \(\operatorname{im}A\subseteq(\ker A)^\perp\)；两边维数都是 \(n-\dim\ker A\)，故相等。最后在式\eqref{g1:eq:quadratic-translation}中取中心 \(a_0\) 和 \(y\in\ker A\)，二次项、线性项都为零，得到 \(F(a_0+y)=F(a_0)\)。
\end{proof}

\begin{example}[唯一中心、一族中心与没有中心]
	\label{g1:exm:quadric-centers}
	\(F(x,y)=(x-1)^2+4(y+2)^2-1\) 的矩阵 \(A\) 可逆，中心唯一，为 \((1,-2)\)。在 \(\mathbb R^3\) 中，\(G(x,y,z)=(x-1)^2+y^2-1\) 的中心是直线 \(\{(1,0,z):z\in\mathbb R\}\)，因为沿 \(z\) 方向平移不改变多项式。对 \(H(x,y)=x^2-y\)，有 \(\ker A=\mathbb R(0,1)\)，而线性项在这个方向上不为零，所以没有中心。抛物线的顶点存在，却不是中心；顶点的作用不能由“把原点移到那里就消去全部线性项”来解释。
\end{example}

\subsection{平移消去线性项及其障碍}
\label{g1:subsec:0202:02}

平移只能消去落在二次项像空间中的线性项；剩余分量由核记录。记
\[
	\mathbb R^n=\operatorname{im}A\oplus\ker A,
	\qquad b=b_1+b_0,
	\quad b_1\in\operatorname{im}A,\quad b_0\in\ker A.
\]
这里用的是通常点积下的正交分解；它的存在与唯一性由\cref{g1:prop:projection}保证。限制 \(A|_{\operatorname{im}A}\) 是可逆的：若像空间中的向量也落在核中，它同时属于两个正交子空间，便只能为零；有限维下单射即满射。因此存在唯一 \(a\in\operatorname{im}A\) 满足 \(Aa=-b_1\)。由式\eqref{g1:eq:quadratic-translation}，
\begin{equation}
	\label{g1:eq:residual-linear-term}
	F(a+y)=y^{\mathsf T}Ay+2b_0^{\mathsf T}y+\kappa,
	\qquad \kappa=F(a).
\end{equation}
平移对 \(b\) 的影响只是加上 \(Aa\)，始终属于 \(\operatorname{im}A\)。所以投影到 \(\ker A\) 的分量 \(b_0\) 是任何平移都无法改变的障碍。

\ProseParagraphBreak

若 \(b_0=0\)，我们已经到达中心，常数 \(\kappa\) 是中心上的值。若 \(b_0\ne0\)，令 \(\beta=\lVert b_0\rVert>0\)，在核中选第一条单位坐标轴为 \(b_0/\beta\)。沿这条轴的坐标记为 \(z_1\)，剩余线性项就是 \(2\beta z_1\)。再把 \(z_1\) 平移 \(-\kappa/(2\beta)\)，就能消去常数。这个方向没有二次项，因此此次平移不会重新制造线性项或交叉项。这正是有中心与无中心两类标准形不同的原因。

\begin{example}[一次完整配方]
	\label{g1:exm:parabolic-completion}
	考虑
	\[
		F(x,y,z)=x^2+2xy+y^2+2x+2y-4z+3.
	\]
	先取正交坐标 \(u=(x+y)/\sqrt2\)、\(v=(x-y)/\sqrt2\)、\(w=z\)，得到
	\[
		F=2u^2+2\sqrt2\,u-4w+3
		=2(u+1/\sqrt2)^2-4w+2.
	\]
	令 \(U=u+1/\sqrt2\)、\(W=w-1/2\)，则方程为 \(U^2=2W\)，而 \(v\) 完全自由。它是一条抛物线沿 \(v\) 方向平移得到的柱面。这里二次项的核由 \(v,w\) 两个方向张成，但只有 \(w\) 方向上的线性项不能消去；“有核”与“没有中心”不是同一句话。
\end{example}

\subsection{二次多项式的 Euclid 与仿射标准形}
\label{g1:subsec:0202:03}

上一章已经分别给出两种对角化。\cref{g1:thm:spectral}在实有限维内积空间中用标准正交特征基消去交叉项，保留特征值；\cref{g1:thm:inertia}允许各方向分别缩放，把非零系数归一为 \(1\) 或 \(-1\)，保留正、负和零的个数。现在加上刚才的平移，就得到完整的非齐次标准形。

\begin{theorem}[二次多项式的 Euclid 标准形]
	\label{g1:thm:euclidean-quadric-normal}
	设 \(F(x)=x^{\mathsf T}Ax+2b^{\mathsf T}x+c\) 是 \(\mathbb R^n\) 上次数不超过二的实多项式，\(A\) 对称，秩为 \(k\)。通过 Euclid 坐标变换 \(x=Qy+a\)，其中 \(Q^{\mathsf T}Q=I\)，它可化为下列标准形，且恰属于其中一类：
	\begin{enumerate}
		\item 有中心时，
		\[
			F=\lambda_1y_1^2+\cdots+\lambda_ky_k^2+\kappa,
			\qquad \lambda_i\ne0,
		\]
		其余 \(n-k\) 个变量不出现。
		\item 没有中心时，\(k<n\)，且
		\[
			F=\lambda_1y_1^2+\cdots+\lambda_ky_k^2+2\beta z_1,
			\qquad \lambda_i\ne0,\quad\beta>0,
		\]
		其中 \(z_1\) 是核中的一个坐标，其余 \(n-k-1\) 个变量不出现。
	\end{enumerate}
	在第一类中，非零特征值多重集 \(\{\lambda_i\}\) 和中心值 \(\kappa\) 是完全不变量；在第二类中，完全不变量是同一个特征值多重集及 \(\beta=\lVert P_{\ker A}b\rVert\)。这里完全不变量的意思是：固定 \(n\) 后，两个多项式能由 Euclid 坐标变换相互转化，当且仅当所属类别和相应数据相同。
\end{theorem}

\begin{proof}
	由\cref{g1:cor:orthogonal-diagonalization}，可以在 \(\operatorname{im}A\) 中选择对应全部非零特征值的标准正交基，并在 \(\ker A\) 中选择标准正交基。配合式\eqref{g1:eq:residual-linear-term}及其后的核内平移，得到两种所列标准形。若 \(b_0\ne0\)，核非零，所以 \(k<n\)。这证明存在性。

	对于唯一性，任意 Euclid 坐标变换使二次项矩阵变成 \(Q^{\mathsf T}AQ\)，故特征值多重集不变。中心存在与否由定义保持，且有中心时多项式在中心上的值不变。无中心时，平移把 \(b\) 改为 \(b+Aa\)，不改变它在核中的投影；正交换基把核及这个投影同时作等距变换，故投影的长度 \(\beta\) 不变。反过来，若这些数据相同，重排非零特征轴后，两个多项式具有完全相同的标准表达式；把各自的坐标变换复合，就得到它们之间的 Euclid 变换。
\end{proof}

若分类的是方程 \(F=0\)，则还允许把 \(F\) 整体乘以 \(\mu\ne0\)。在有中心的情形，数据变成 \((\{\mu\lambda_i\},\mu\kappa)\)；在无中心的情形，数据变成 \((\{\mu\lambda_i\},|\mu|\beta)\)。这是 Euclid 方程分类中唯一需要再作的认同，不能任意改变每一个 \(\lambda_i\) 的大小。例如椭圆 \(x^2/a^2+y^2/b^2=1\) 的半轴长仍被 Euclid 等价保留。

\begin{theorem}[二次多项式的仿射标准形]
	\label{g1:thm:affine-quadric-normal}
	设 \(n\) 为非负整数，\(A\) 为实对称 \(n\) 阶矩阵，\(b\in\mathbb R^n\)、\(c\in\mathbb R\)，\(F\colon\mathbb R^n\rightarrow\mathbb R\) 为式\eqref{g1:eq:affine-quadratic-polynomial}所示的多项式 \(F(x)=x^{\mathsf T}Ax+2b^{\mathsf T}x+c\)，\(A\) 的惯性为 \((p,q,r)\)，其中 \(p+q+r=n\)。允许任意可逆仿射坐标变换后，有如下标准形与唯一性。
	\begin{enumerate}
		\item 有中心时，
		\[
			F=u_1^2+\cdots+u_p^2-v_1^2-\cdots-v_q^2+\kappa,
		\]
		其余 \(r\) 个变量不出现；\((p,q,r,\kappa)\) 是完全不变量。
		\item 没有中心时，\(r\ge1\)，且
		\[
			F=u_1^2+\cdots+u_p^2-v_1^2-\cdots-v_q^2+2z_1,
		\]
		其余 \(r-1\) 个变量不出现；\((p,q,r)\) 是完全不变量。
	\end{enumerate}
	若进一步只分类方程的非零常数倍，则第一类可写为
	\begin{equation}
		\label{g1:eq:affine-centered-quadric}
		u_1^2+\cdots+u_p^2-v_1^2-\cdots-v_q^2=\varepsilon,
		\qquad\varepsilon\in\{-1,0,1\},
	\end{equation}
	唯一的重复是 \((p,q,\varepsilon)\) 与 \((q,p,-\varepsilon)\)；第二类可写为
	\begin{equation}
		\label{g1:eq:affine-parabolic-quadric}
		u_1^2+\cdots+u_p^2-v_1^2-\cdots-v_q^2=2z_1,
	\end{equation}
	唯一的重复是交换 \(p,q\)。各式中的自由变量数仍由 \(n\) 决定。
\end{theorem}

\begin{proof}
	从 Euclid 标准形出发，对每个非零特征方向用 \(u_i=\sqrt{|\lambda_i|}\,y_i\) 缩放，即可把二次项化成正负平方和。这是\cref{g1:thm:inertia}中的同一个缩放。无中心时再用 \(z'_1=\beta z_1\)，便把线性项系数化为 \(2\)。所有缩放因相应系数非零而可逆。

	一般坐标变换 \(x=Cy+a\) 使二次项矩阵变成 \(C^{\mathsf T}AC\)。惯性定理说明 \((p,q,r)\) 不变。中心性质与中心值都由多项式本身的等式定义，也不因坐标改变而变。故所列数据是必要条件；数据相同则标准表达式相同，逆用坐标变换即得充分性。

	最后考虑方程的倍数。有中心时，若 \(\kappa\ne0\)，先乘以正数 \(1/|\kappa|\)，再缩放二次项各坐标，便得到右端 \(\varepsilon=-\operatorname{sign}\kappa\)；若 \(\kappa=0\)，右端已为零。正倍数不改变惯性，负倍数同时交换 \(p,q\) 并改变 \(\varepsilon\) 的符号，所以恰有陈述中的认同。没有中心时，常数已经消去，任意非零倍数都可由重新缩放线性坐标及二次坐标吸收；负倍数只额外交换正、负方向。固定多项式的唯一性说明不存在其他认同。
\end{proof}

上述结论也覆盖 \(A=0\) 的边界：非零线性多项式都没有中心，经仿射变换成为 \(2z_1\)；常数多项式有每一点作为中心，值 \(\kappa\) 保持不变。它们的零点集分别为超平面、空集或整个空间，但不属于本节要求 \(A\ne0\) 的二次曲面。

\begin{example}[同一仿射类中的不同 Euclid 形状]
	\label{g1:exm:ellipses-affine-euclidean}
	方程 \(x^2+4y^2=1\) 在替换 \(u=x\)、\(v=2y\) 后成为单位圆，因此与圆仿射等价。若它与单位圆 Euclid 等价，则有中心方程的分类数据应满足
	\[
		(\{1,4\},-1)=(\{\mu,\mu\},-\mu)
	\]
	至多相差特征值的次序。常数项迫使 \(\mu=1\)，而特征值不相同，矛盾。两种分类中允许的变换不同，留下的几何信息也不同。
\end{example}

\subsection{退化参数、直纹及实点集}
\label{g1:subsec:0202:04}

标准形可以直接给出实点是否存在、哪些方向完全自由，以及直线能否落在曲面中的判据。对式\eqref{g1:eq:affine-centered-quadric}，只看出现的 \(p+q\) 个变量；每一个解都可以再带上任意 \(r\) 个自由坐标。\cref{g1:tab:affine-real-quadric-types}列出实点的区别，其中假定 \(p+q>0\)。

\begin{table}[htbp]
	\centering
	\caption{有中心标准形的实点集}
	\label{g1:tab:affine-real-quadric-types}
	\begin{tabularx}{\linewidth}{@{}l l X@{}}
		\toprule
		二次项 & 右端 & 出现变量中的实解 \\
		\midrule
		\(p>0,\ q=0\) & \(\varepsilon=1\) & 单位球面；\(p=1\) 时为两点 \\
		\(p>0,\ q=0\) & \(\varepsilon=0\) & 仅原点 \\
		\(p>0,\ q=0\) & \(\varepsilon=-1\) & 空集 \\
		\(p,q>0\) & 任意 \(\varepsilon\) & 均非空；\(\varepsilon=0\) 时为过原点的二次锥 \\
		\bottomrule
	\end{tabularx}
\end{table}

半负定的情形把方程乘以 \(-1\) 即归入前三行。表中前三行来自平方和非负，最后一行则可分别用一个正方向和一个负方向取得任意实数。右端为零时，若 \(x\) 是解，则 \(tx\) 对任意实数 \(t\) 都是解，因而出现直线。右端不为零时仍可能有直线，但它们通常不经过中心。

\begin{proposition}[曲面中的仿射直线]
	\label{g1:prop:affine-quadric-line}
	设 \(A\) 为实对称 \(n\) 阶矩阵，\(b,a,v\in\mathbb R^n\)、\(c\in\mathbb R\)，且 \(v\ne0\)，并令 \(F(x)=x^{\mathsf T}Ax+2b^{\mathsf T}x+c\)。整条直线 \(a+\mathbb Rv\) 包含在 \(F=0\) 中，当且仅当
	\[
		F(a)=0,\qquad (Aa+b)^{\mathsf T}v=0,
		\qquad v^{\mathsf T}Av=0.
	\]
\end{proposition}

\begin{proof}
	把直线参数代入方程：
	\[
		F(a+tv)=F(a)+2t(Aa+b)^{\mathsf T}v+t^2v^{\mathsf T}Av.
	\]
	一元二次多项式对全部实数 \(t\) 为零，当且仅当三个系数都为零。这给出两方向的蕴含。
\end{proof}

\begin{example}[两种抛物面的直线]
	\label{g1:exm:paraboloid-lines}
	在 \(z=x^2+y^2\) 上，直线方向 \((a,b,c)\) 必须满足 \(a^2+b^2=0\)，故 \(a=b=0\)；一次项条件再迫使 \(c=0\)，与方向非零矛盾，所以曲面不包含实直线。相反，\(z=xy\) 对每个固定的 \(s\) 都包含直线
	\[
		(x,y,z)=(s,t,st)\qquad(t\in\mathbb R),
	\]
	也对每个固定的 \(t\) 包含 \((s,t,st)\) 随 \(s\) 变化的直线。这两个已给出的直线族中，各恰有一条直线通过每个曲面点。由连续变化的一族直线扫出的曲面称为\term[zhiwenqumian]{直纹曲面}[ruled surface]；这两族直线称为两个\term[zhiwenzu]{直纹族}[rulings]。\cref{g1:fig:quadric-rulings}画出其中有限段，下一节将补全这两个族，并证明它们包含该曲面射影化后的全部直线。
\end{example}

\begin{figure}[htbp]
	\centering
	\input{Figures/Chapter02-QuadricRulings.tex}
	\caption[双曲抛物面的两个直纹族]{双曲抛物面 \(z=xy\) 的两个直纹族。蓝线固定 \(x\)，橙线固定 \(y\)，标出的点为 \(P=(1/2,1/2,1/4)\)。图中只画 \(-1\le x,y\le1\) 的部分，空间位置采用平行投影示意。两族线在曲面上相交，同族线互不相交。}
	\label{g1:fig:quadric-rulings}
\end{figure}

无中心标准形\eqref{g1:eq:affine-parabolic-quadric}可直接解出 \(z_1\)，因此总有实点，且是其他 \(n-1\) 个变量的图像；未出现的变量给出柱面的自由方向。有中心标准形中的自由变量也给出这种方向。例如 \(x^2+y^2=1\) 在 \(\mathbb R^3\) 中是一张圆柱面，虽然它的二次项矩阵有核，但其每一个实点附近都保持通常的圆柱形状。“矩阵有核”不能自动理解为实点集到处出现尖点或交叉。

\begin{example}[参数变化要看哪一项]
	\label{g1:exm:affine-quadric-degeneration}
	族 \(x^2+y^2=t\) 的二次项矩阵始终为 \(I_2\)，但当 \(t>0\)、\(t=0\)、\(t<0\) 时，实点集分别为圆、一个点、空集。改变的量是中心值 \(-t\)。另一族
	\[
		x^2+y^2-tz^2=1
	\]
	在 \(t<0\) 时为椭球面，\(t=0\) 时为圆柱面，\(t>0\) 时为单叶双曲面；后一名称对应截面圆半径 \(\sqrt{1+tz^2}\)，它对每个实数 \(z\) 都存在。这次发生变化的是二次项的秩与惯性。\cref{g1:prop:rank-inertia-continuity}说明，在连续实对称矩阵族中，秩固定时惯性局部不变；本例正是在 \(t=0\) 丢掉一个非零特征值，才从一个惯性类型走到另一个类型。
\end{example}

为了把二次项、线性项和常数项放进同一个对称矩阵，下一节引入一个齐次坐标。这会使上述两种看似不同的退化都能用一个齐次二次型描述，也会说明椭圆、抛物线和双曲线为何可能属于同一个射影类型。

% !TeX root = ../../Book1.tex
% 纲要标题：### 2.3 射影二次曲面与秩
% 纲要位置：几何学计划纲要.md，第 171 行。
% 正文状态：初稿；尚未正式审核。

% 纲要内容开始
% 
% * 齐次化及无穷远截面
% * 基域、秩和射影等价
% * 极大线性子空间及直纹族
% * 复化与实形式的不同分类信息
% 
% 纲要内容结束

\section{射影二次曲面与秩}
\label{g1:sec:0203}

仿射二次方程的三类系数在平移时互相混合，因此上一节的分类要分别讨论有无中心。齐次坐标提供了另一种看法：用新增坐标补齐次数，把所有项放进同一个二次型。这样可以直接使用合同分类，但代价是我们不再固定哪一个超平面被称作无穷远。椭圆、抛物线和双曲线的区别，恰好会在这个选择中重新出现。

\subsection{齐次化及无穷远截面}
\label{g1:subsec:0203:01}

设 \(F(x)=x^{\mathsf T}Ax+2b^{\mathsf T}x+c\) 是 \(K^n\) 上次数恰为二的多项式，其中 \(K=\mathbb R\) 或 \(\mathbb C\)，\(A\) 对称且非零。新增坐标 \(X_0\)，把原来的坐标写为 \(X=(X_1,\ldots,X_n)^{\mathsf T}\)，定义它的\term[qicihua]{齐次化}[homogenization]为
\begin{equation}
	\label{g1:eq:quadric-homogenization}
	\begin{aligned}
		F^{\mathrm h}(X_0,X)
		&=X^{\mathsf T}AX+2X_0b^{\mathsf T}X+cX_0^2\\
		&=\begin{pmatrix}X_0&X^{\mathsf T}\end{pmatrix}
		M\begin{pmatrix}X_0\\X\end{pmatrix},\\
		M&=\begin{pmatrix}c&b^{\mathsf T}\\b&A\end{pmatrix}.
	\end{aligned}
\end{equation}
它的每一项总次数都是二，所以向量整体乘以 \(t\ne0\) 时，函数值乘以 \(t^2\)；零值条件因此能在射影空间上定义。

\begin{definition}[射影二次曲面]
	\label{g1:def:projective-quadric}
	设 \(V\) 是有限维 \(K\) 向量空间，\(K=\mathbb R\) 或 \(\mathbb C\)，\(q\) 是 \(V\) 上非零二次型。由齐次方程 \(q=0\) 定义的\term[sheyingerciqumian]{射影二次曲面}[projective quadric]记为
	\[
		Q(q)=\{\,[v]\in\mathbb P(V):q(v)=0\,\}.
	\]
	讨论方程等价时，保留 \(q\) 到整体非零倍数的信息。其秩是对应对称矩阵的秩；矩阵可逆时称该二次曲面方程非退化，否则称其退化。
\end{definition}

在标准仿射图 \(X_0\ne0\) 中，用 \(x_i=X_i/X_0\)，有
\[
	F^{\mathrm h}(X_0,X)=X_0^2F(X/X_0).
\]
因此这个图中的零点恰好是原来的 \(F=0\)。在无穷远超平面 \(X_0=0\) 上，剩下的方程则是 \(X^{\mathsf T}AX=0\)。它记录二次项的零值方向，称为相对于这个仿射图的\term[wuqiongyuanjiemian]{无穷远截面}[section at infinity]。若把次数小于二的多项式也机械地补到二次，就可能多出整个无穷远超平面，例如线性方程 \(x_1=0\) 被写成 \(X_0X_1=0\)；这正是这里要求原方程次数恰为二的原因。

\begin{example}[三种实二次曲线的无穷远点]
	\label{g1:exm:conics-at-infinity}
	仿射平面中的三个方程齐次化后成为
	\[
		\begin{array}{c|c|c}
			\text{仿射方程}&\text{齐次方程}&\text{实无穷远点}\\ \hline
			x_1^2+x_2^2=1&X_1^2+X_2^2-X_0^2=0&\text{无}\\
			x_1^2-x_2^2=1&X_1^2-X_2^2-X_0^2=0&[0:1:1],\ [0:1:-1]\\
			x_1^2=x_2&X_1^2-X_0X_2=0&[0:0:1]
		\end{array}
	\]
	最后一行在无穷远超平面上的限制为平方 \(X_1^2\)，因此只给出一个实无穷远点。中间一行有两个不同的零值方向；第一行的正定二次项没有非零实零值方向。这里说的无穷远，始终相对于已经选定的 \(X_0=0\)。
\end{example}

在实数域上，还应区分齐次化与通常拓扑中的闭包。例如 \(x_1^2+1=0\) 在 \(\mathbb R^2\) 中没有实点，而齐次方程 \(X_1^2+X_0^2=0\) 在 \(\mathbb P^2(\mathbb R)\) 中有点 \([0:0:1]\)。所以齐次化的实零点集不必等于原实零点集的拓扑闭包。本节依据的是多项式和齐次坐标的计算，不把实点的逼近性质暗中放进“补入无穷远点”的说法中。

\ProseParagraphBreak

仿射坐标变换 \(x=Cy+a\) 在齐次坐标中对应可逆矩阵
\[
	T=\begin{pmatrix}1&0\\a&C\end{pmatrix}.
\]
它把 \((Y_0,Y)\) 送到 \((Y_0,aY_0+CY)\)，所以保持 \(X_0=0\)，并使 \(M\) 变为 \(T^{\mathsf T}MT\)。若允许任意可逆的 \((n+1)\) 阶矩阵，就不必再保持这个超平面，得到范围更大的射影等价关系。

\begin{proposition}[齐次化的秩]
	\label{g1:prop:homogenized-rank}
	设 \(F\) 是 \(\mathbb R^n\) 上次数恰为二的实多项式，二次项矩阵 \(A\) 的惯性为 \((p,q,r)\)，秩为 \(k=p+q\)。记齐次化矩阵为 \(M\)。
	\begin{enumerate}
		\item 若 \(F\) 有中心，中心值为 \(\kappa\)，则 \(M\) 合同于 \(\operatorname{diag}(\kappa,A)\)，其秩为 \(k\) 或 \(k+1\)，分别对应 \(\kappa=0\) 或 \(\kappa\ne0\)。
		\item 若 \(F\) 没有中心，则 \(\operatorname{rank}M=k+2\)，惯性为 \((p+1,q+1,r-1)\)。
	\end{enumerate}
\end{proposition}

\begin{proof}
	有中心时，把该中心移到原点，式\eqref{g1:eq:quadratic-translation}给出 \(F=y^{\mathsf T}Ay+\kappa\)，其齐次化矩阵正是 \(\operatorname{diag}(\kappa,A)\)。

	无中心时，\cref{g1:thm:euclidean-quadric-normal}给出 \(F=\sum_{i=1}^k\lambda_i y_i^2+2\beta z_1\)，其中 \(\beta>0\)。齐次化后，\(z_1\) 与新增坐标 \(Y_0\) 组成矩阵块
	\[
		\begin{pmatrix}0&\beta\\\beta&0\end{pmatrix},
	\]
	它在方向 \((1,1)\)、\((1,-1)\) 上分别有特征值 \(\beta\)、\(-\beta\)。所以它贡献一个正方向、一个负方向，原核中的其余 \(r-1\) 个方向仍为零。由此同时得到秩与惯性公式。
\end{proof}

这说明抛物线的二次项虽有核，齐次化的二次型却可以非退化；也说明一个有中心方程可能因中心值变成零而使齐次化的秩下降。对上一节的族 \(x^2+y^2=t\)，齐次化矩阵为 \(\operatorname{diag}(-t,1,1)\)，正好在 \(t=0\) 退化。

\subsection{基域、秩和射影等价}
\label{g1:subsec:0203:02}

对于 \(K=\mathbb R\) 或 \(\mathbb C\)，二次型 \(q\) 都对应唯一的对称双线性型
\[
	B(u,v)=\frac{q(u+v)-q(u)-q(v)}2.
\]
这与\cref{g1:def:symmetric-form}的实数情形相同。复数情形仍然对两个变量分别线性，矩阵公式是 \(u^{\mathsf T}Av\)，不含复共轭；它并不是用来测长度的复内积。根基仍定义为与所有向量配对均为零的向量组成的子空间。

\ProseParagraphBreak

在实数域上，惯性定理已经证明正负号无法通过实可逆换基消去；在复数域上，\(i^2=-1\) 却使 \(-x^2=(ix)^2\)。因此复数域上的分类应当更粗。为避免把实自伴谱定理误用到复对称矩阵上，我们直接给出所需的复数证明。

\begin{lemma}[复对称型的合同标准形]
	\label{g1:lem:complex-symmetric-normal}
	设 \(V\) 是 \(N\) 维复向量空间，\(B\) 是其上的对称双线性型，秩为 \(s\)。存在一组复基，使其矩阵为 \(\operatorname{diag}(I_s,0_{N-s})\)。两个同维复对称双线性型合同，当且仅当它们的秩相同。
\end{lemma}

\begin{proof}
	若 \(B=0\)，任意基都满足结论。若 \(B\ne0\)，则存在 \(v\) 使 \(B(v,v)\ne0\)：否则由极化公式，所有 \(B(u,v)\) 都为零。非零复数都有平方根，所以缩放 \(v\) 可得到 \(e\) 满足 \(B(e,e)=1\)。对任意 \(w\in V\)，分解
	\[
		w=B(w,e)e+\bigl(w-B(w,e)e\bigr)
	\]
	的第二项与 \(e\) 正交。并且 \(\mathbb Ce\cap e^\perp=0\)，因为 \(B(te,e)=t\)。这里 \(e^\perp\) 表示关于 \(B\) 的正交空间，故 \(V=\mathbb Ce\oplus e^\perp\)，在这个分解下矩阵为一个 \(1\) 与较低维对称型的直和。

	继续对后者归纳，直到剩下零型，就得到对角元只含 \(1\) 与 \(0\) 的矩阵。乘以可逆矩阵不改变秩，所以其中 \(1\) 的个数必为 \(s\)。这既给出数目唯一性，也说明秩相同的两个型都合同于同一个标准矩阵；反向蕴含由秩在合同变换下不变得到。
\end{proof}

\begin{theorem}[实、复射影二次方程的分类]
	\label{g1:thm:projective-quadric-classification}
	设 \(K=\mathbb R\) 或 \(\mathbb C\)，\(q,q'\) 是 \(N\) 维 \(K\) 向量空间上的非零二次型。称其射影方程等价，是指存在 \(T\in\operatorname{GL}(N,K)\) 及 \(\mu\in K\setminus\{0\}\)，使
	\[
		q'(v)=\mu q(Tv)\qquad(v\in K^N).
	\]
	在复数域上，等价的充要条件是二者秩相同。在实数域上，若惯性分别为 \((p,q_-,r)\) 与 \((p',q'_-,r')\)，等价的充要条件是
	\[
		r=r',\qquad \{p,q_-\}=\{p',q'_-\},
	\]
	其中花括号表示允许交换两个惯性指数的无序二元组。
\end{theorem}

\begin{proof}
	矩阵变化为 \(A' =\mu T^{\mathsf T}AT\)，故任何基域上秩都保持不变。复数情形由\cref{g1:lem:complex-symmetric-normal}，相同秩已足以给出合同，甚至不必再乘一个标量。

	实数情形由\cref{g1:thm:inertia}，合同保持 \((p,q_-,r)\)。乘以正数不改变这三个数；乘以负数则交换正、负惯性指数，保持根基维数。故所列条件必要。若条件成立，两组惯性相同，或把其中一个型乘以 \(-1\) 后相同；再用惯性定理的标准形便得到所需合同，故条件也充分。
\end{proof}

这里用 \(q_-\) 表示负惯性指数，是为了与二次型 \(q\) 区别；下面没有这种冲突时仍写惯性为 \((p,q,r)\)。定理分类的是方程及其基域，不只是肉眼所见的实点形状。例如在 \(\mathbb P^2(\mathbb R)\) 中，方程 \(X_0^2=0\) 与 \(X_0^2+X_1^2=0\) 的实零点集分别是一条直线与一个点；它们的秩分别为 \(1,2\)。即使两个方程的实零点集偶然相同，也不能因此认定它们按上述定义等价。

\begin{example}[椭圆、双曲线和抛物线的射影类型]
	\label{g1:exm:one-projective-conic-type}
	\cref{g1:exm:conics-at-infinity}中的三个齐次型都具有秩 \(3\)，正、负惯性指数的无序二元组均为 \(\{2,1\}\)，所以它们的实射影方程等价。第三个式子中，\(-X_0X_2\) 在 \(X_0+X_2\) 与 \(X_0-X_2\) 两个方向上给出相反符号，再加上 \(X_1^2\) 就能直接核对惯性。它们作为仿射曲线仍然不同，因为给出等价的射影变换不必保持已选的无穷远直线。区分“曲线”与“曲线加上一条指定的无穷远直线”，就能同时理解这两个结论。
\end{example}

\subsection{极大线性子空间及直纹族}
\label{g1:subsec:0203:03}

一个射影点 \([v]\) 在 \(Q(q)\) 上，只要求 \(q(v)=0\)。若希望整条射影直线、甚至更高维的射影子空间都在曲面上，就必须要求它的全部向量都满足这个条件。由\cref{g1:prop:projective-dimension}，\(k\) 维射影子空间来自 \(k+1\) 维向量子空间；因此以下问题可以完全化成双线性型问题。

\begin{definition}[全迷向子空间]
	\label{g1:def:totally-isotropic}
	设 \(V\) 是实或复向量空间，\(B\) 是对称双线性型，\(q(v)=B(v,v)\)。子空间 \(U\subseteq V\) 若满足 \(q|_U=0\)，称为\term[quanmixiangzikongjian]{全迷向子空间}[totally isotropic subspace]。等价地，\(B(u,w)=0\) 对所有 \(u,w\in U\) 成立；等价性由极化公式得到。不能真包含于更大全迷向子空间中的 \(U\) 称为极大，而可能达到的最大维数还需证明。
\end{definition}

注意，全迷向条件要求子空间内任意两个向量之间的配对也为零。两个分别迷向的向量未必张成全迷向平面：对 \(x^2-y^2\)，\((1,1)\) 与 \((1,-1)\) 都迷向，但它们的和为 \((2,0)\)，二次型值为 \(4\)。

\begin{theorem}[二次曲面中的最大线性维数]
	\label{g1:thm:quadric-linear-dimension}
	设 \(V\) 是 \(N\) 维实或复向量空间，\(q\ne0\) 是二次型，对应型为 \(B\)，根基维数为 \(r\)。
	\begin{enumerate}
		\item 若基域为实数，\(B\) 的惯性为 \((p,q_-,r)\)，则全迷向子空间的最大维数为
		\[
			r+\min(p,q_-).
		\]
		\item 若基域为复数，\(B\) 的秩为 \(s=N-r\)，则该最大维数为
		\[
			r+\lfloor s/2\rfloor.
		\]
	\end{enumerate}
	每个按包含关系极大的全迷向子空间都达到相应最大维数。若该数为 \(d>0\)，则 \(Q(q)\) 中最大射影线性子空间的维数为 \(d-1\)；若 \(d=0\)，则 \(Q(q)\) 没有基域内的点。
\end{theorem}

\begin{proof}
	先处理根基。由标准形可取直和 \(V=R\oplus W\)，其中 \(R=\operatorname{rad}B\)，\(B|_W\) 非退化。任意全迷向子空间 \(U\) 投影到 \(W\) 的像 \(\overline U\) 仍全迷向，而且
	\[
		\dim U=\dim(U\cap R)+\dim\overline U
		\le r+\dim\overline U.
	\]
	反过来，\(W\) 中的全迷向子空间与 \(R\) 直和仍全迷向。所以最大维数等于 \(r\) 加上非退化部分中的最大维数。

	在实非退化情形，标准形为 \(\lVert x\rVert^2-\lVert y\rVert^2\)，其中 \(x\in\mathbb R^p\)、\(y\in\mathbb R^{q_-}\)。对全迷向子空间，投影到任一组坐标都是单射：例如 \(x=0\) 时，零值条件给出 \(\lVert y\rVert^2=0\)，迫使 \(y=0\)。所以其维数同时不超过 \(p\) 与 \(q_-\)。取正、负标准基 \(e_i,f_j\)，向量
	\[
		e_1+f_1,\ldots,e_m+f_m,
		\qquad m=\min(p,q_-),
	\]
	彼此配对均为零，且线性无关，故上界可达到。

	在复非退化的 \(s\) 维空间中，若 \(U\) 全迷向，则 \(U\subseteq U^\perp\)。这里 \(\dim U^\perp=s-\dim U\)：非退化性使 \(w\mapsto B(w,-)\) 成为到对偶空间的同构，而任意 \(U\) 上的线性函数都可以先扩充一组基，再延拓到全空间；因此 \(w\mapsto B(w,-)|_U\) 满射到 \(U^*\)，其核就是 \(U^\perp\)。秩与零化度公式给出所需维数。于是 \(2\dim U\le s\)。在\cref{g1:lem:complex-symmetric-normal}的标准基中，
	\[
		e_1+ie_2,\ e_3+ie_4,\ldots,
		e_{2m-1}+ie_{2m},\qquad m=\lfloor s/2\rfloor,
	\]
	构成线性无关的全迷向向量组，故复数情形的上界也可达到。

	最后证明极大子空间不可能停在更小的维数。极大全迷向子空间必包含 \(R\)，否则加上根基中的一个新方向仍然全迷向；此时它等于 \(R\oplus(U\cap W)\)。于是只需在非退化部分中证明结论，仍把其中的全迷向子空间记为 \(U\)。若 \(u\in U\) 非零，非退化性给出 \(v\) 满足 \(B(u,v)=1\)。以
	\[
		v'=v-\frac{B(v,v)}2u
	\]
	替代 \(v\)，就有 \(B(v',v')=0\)、\(B(u,v')=1\)。平面 \(H=\operatorname{span}(u,v')\) 的配对矩阵为 \(\left(\begin{smallmatrix}0&1\\1&0\end{smallmatrix}\right)\)，因而非退化。任意 \(w\) 都可按
	\[
		w=B(w,v')u+B(w,u)v'
		+\bigl(w-B(w,v')u-B(w,u)v'\bigr)
	\]
	分成 \(H\) 与 \(H^\perp\) 两部分，且交集为零。\(H^\perp\) 上的限制仍非退化：若其中一个向量与整个 \(H^\perp\) 正交，它也与 \(H\) 正交，便与整个空间正交，故只能为零。若 \(w\in U\)，则 \(B(w,u)=0\)，所以 \(U=Ku\oplus(U\cap H^\perp)\)。

	重复这一步，若 \(\dim U=k\)，便剥离出 \(k\) 个这样的二维平面，余下非退化空间 \(W_0\) 与它们正交。在实数域中每剥离一次，正、负惯性指数各减一，所以若 \(k<\min(p,q_-)\)，\(W_0\) 仍有一正一负两个方向；它们可组成新的迷向向量。在复数域中，若 \(k<\lfloor s/2\rfloor\)，则 \(\dim W_0\ge2\)，也可由标准形的前两个方向组成 \(e_1+ie_2\)。这些新向量都与原 \(U\) 正交，使 \(U\) 真正扩大，违背极大性。这也覆盖 \(U=0\) 而上界为正的情形。

	最后，非零全迷向子空间 \(U\) 与曲面内的射影子空间 \(\mathbb P(U)\) 恰好对应，其维数减少一。若不存在非零全迷向子空间，连一个迷向向量都不存在，曲面便没有相应基域的射影点。
\end{proof}

定理中的两个维数计数揭示了实、复几何的差别。例如实球面齐次化后的型为 \(X_1^2+X_2^2+X_3^2-X_0^2\)，惯性为 \((3,1,0)\)。全迷向子空间最大只有一维，所以球面射影化后没有实射影直线。允许复坐标以后，同一个矩阵秩为 \(4\)，最大全迷向维数变为二，因而出现复射影直线。

\ProseParagraphBreak

对于惯性 \((2,2,0)\) 的实型，最大射影线性子空间就是直线。下面的模型不仅表明直线存在，而且把它们全部写出。这个表示也将在\secref{g1:sec:0206}中成为一个射影映射的例子。

\begin{proposition}[二次曲面的两个直纹族]
	\label{g1:prop:quadric-two-rulings}
	设 \(K=\mathbb R\) 或 \(\mathbb C\)，在 \(\mathbb P^3(K)\) 中考虑
	\[
		Q:\ X_0X_3-X_1X_2=0.
	\]
	每个点唯一写为
	\begin{equation}
		\label{g1:eq:quadric-rank-one-param}
		[X_0:X_1:X_2:X_3]
		=[su:sv:tu:tv],
		\qquad [s:t],[u:v]\in\mathbb P^1(K).
	\end{equation}
	固定 \([s:t]\) 而让 \([u:v]\) 变化，或固定 \([u:v]\) 而让 \([s:t]\) 变化，各得到一个直纹族。这两族包含 \(Q\) 上的全部射影直线；同族不同直线互不相交，异族任意两条直线恰交于一点。过 \(Q\) 上的每一点恰有两条直线，每族各一条。
\end{proposition}

\begin{proof}
	把四个坐标排成矩阵
	\[
		Z=\begin{pmatrix}X_0&X_1\\X_2&X_3\end{pmatrix}.
	\]
	方程是 \(\det Z=0\)，而射影坐标要求 \(Z\ne0\)，所以 \(Z\) 的秩恰为一。选一个非零列，其他列都是它的倍数，便可写成 \(Z=a b^{\mathsf T}\)，其中非零列向量 \(a=(s,t)^{\mathsf T}\)、\(b=(u,v)^{\mathsf T}\)。它的列空间决定 \([a]\)，行空间决定 \([b]\)，故这两个射影点都唯一。分别改变 \(a,b\) 的代表，只把 \(Z\) 整体乘以非零数，所以式\eqref{g1:eq:quadric-rank-one-param}良定义。

	固定非零 \(a\) 后，映射 \(b\mapsto ab^{\mathsf T}\) 线性且单射，因此其像是二维向量子空间，射影化后正是一条直线；固定 \(b\) 同理。分解的唯一性说明同族不同直线没有公共点，而分别固定 \([a]\)、\([b]\) 的异族两条线只在 \([ab^{\mathsf T}]\) 相交。

	还须排除第三种直线。设 \(L\subseteq Q\) 为任意射影直线，取它对应二维向量空间的一组基 \(ab^{\mathsf T}\)、\(cd^{\mathsf T}\)。对所有 \(\tau\in K\)，矩阵 \(ab^{\mathsf T}+\tau cd^{\mathsf T}\) 的行列式为零。按两列展开，得到
	\[
		\det\bigl(ab^{\mathsf T}+\tau cd^{\mathsf T}\bigr)
		=\tau\det(a,c)\det(b,d).
	\]
	取 \(\tau=1\)，可知 \(a,c\) 成比例或 \(b,d\) 成比例；二者不能同时成比例，否则两基向量相关。在第一种情形，整条 \(L\) 属于固定 \([a]\) 的族；在第二种情形，属于固定 \([b]\) 的族。全部断言随之成立。
\end{proof}

在图 \(X_0\ne0\) 中令 \(x=X_1/X_0\)、\(y=X_2/X_0\)、\(z=X_3/X_0\)，方程成为 \(z=xy\)，两个直纹族正是上一节画出的固定 \(x\) 与固定 \(y\) 的直线族。不过射影曲面还包含无穷远处的两条线：把 \(X_0=0\) 代入，得到 \(X_1X_2=0\)，即 \(X_0=X_1=0\) 与 \(X_0=X_2=0\)。它们也各属于一个直纹族，交于 \([0:0:0:1]\)，因此射影描述补全了仿射图没有显示的直线。

\subsection{复化与实形式的不同分类信息}
\label{g1:subsec:0203:04}

一个实二次方程可以保持系数不变而允许复数坐标，这称为把方程\term[fuhua]{复化}[complexification]。矩阵的秩不变，因为某阶子式是否非零不因把实数看成复数而改变；惯性的正负号却不再是复换基的不变量。乘以 \(i\) 就能把一个正平方改成负平方，因而所有相同秩的复二次型归入同一类。

\ProseParagraphBreak

为说明丢掉了哪些信息，\cref{g1:tab:real-quadric-forms}比较 \(\mathbb P^3(\mathbb R)\) 中的三种满秩方程。它们的复化都具有秩 \(4\)，所以由\cref{g1:thm:projective-quadric-classification}彼此复射影等价。

\begin{table}[htbp]
	\centering
	\caption{同一个复二次曲面类型的不同实形式}
	\label{g1:tab:real-quadric-forms}
	\begin{tabularx}{\linewidth}{@{}l c X@{}}
		\toprule
		实齐次方程 & 惯性 & 实几何 \\
		\midrule
		\(X_0^2+X_1^2+X_2^2+X_3^2=0\) & \((4,0,0)\) & 没有实点 \\
		\(X_1^2+X_2^2+X_3^2-X_0^2=0\) & \((3,1,0)\) & 有实点，没有实直线 \\
		\(X_0X_3-X_1X_2=0\) & \((2,2,0)\) & 有两个实直纹族 \\
		\bottomrule
	\end{tabularx}
\end{table}

第一行中，正定性排除了非零实零值向量；第二行在 \(X_0=1\) 中就是单位球面，维数定理已经说明它不含实直线；第三行则由两个符号相反的乘积各贡献一个正、一个负方向，且直纹由\cref{g1:prop:quadric-two-rulings}完全描述。这三种实方程可称为该复射影类型的不同\term[shixingshi]{实形式}[real forms]：它们的复化等价，但实数域上的方程并不等价。这里的名称只用于这项具体比较，不预设一般实形式理论。

\ProseParagraphBreak

由实到复并非简单地给图上每一点再加一个虚部。实点的存在、实直线的存在和实惯性都可能在复化后不再可辨；保存下来的秩则仍然控制根基维数及最大复线性子空间维数。后面讨论极性时，非退化性还会决定一个点能否唯一对应一个超平面。因此，在说“二次曲面由秩分类”以前，必须先说明基域；在说“这个曲面有直纹”以前，也必须说明这些直线在哪个基域上定义。

% !TeX root = ../../Book1.tex
% 纲要标题：### 2.4 射影对偶与极性
% 纲要位置：几何学计划纲要.md，第 178 行。
% 正文状态：初稿；尚未正式审核。

% 纲要内容开始
% 
% * 超平面参数空间与双重对偶
% * 非退化配对产生的极点、极超平面
% * 切空间、切超平面与对偶二次曲面
% * 退化配对中极性的失效
% 
% 纲要内容结束

\section{射影对偶与极性}
\label{g1:sec:0204}

平面上一点写成齐次坐标 \([x_0:x_1:x_2]\)，一条直线也由一组非全零系数 \([a_0:a_1:a_2]\) 决定，其方程是 \(a_0x_0+a_1x_1+a_2x_2=0\)。点的坐标与直线的系数虽然形式相似，却属于不同的向量空间。先保留这个区别，再利用二次型建立二者的对应，就能说明为什么一条共线定理会产生另一条共点定理，也能把二次曲面的切超平面组织成新的二次曲面。

\subsection{超平面参数空间与双重对偶}
\label{g1:subsec:0204:01}

设 \(V\) 是 \(K=\mathbb R\) 或 \(\mathbb C\) 上的 \(n+1\) 维向量空间，\(n\ge1\)。回顾对偶空间 \(V^*\) 的元素是线性函数 \(\alpha:V\rightarrow K\)。每个非零 \(\alpha\) 给出射影超平面 \(H_\alpha=\mathbb P(\ker\alpha)\)。两个非零函数给出同一超平面，当且仅当相差非零倍数：若核相同，取 \(v\) 使 \(\alpha(v)=1\)，分解 \(x=\alpha(x)v+(x-\alpha(x)v)\)，即得 \(\beta(x)=\beta(v)\alpha(x)\)。反之，乘非零常数不改变核。每个余维一子空间都是某个非零函数的核，因为可以取它的一组基，补上一个向量，并令线性函数在前者上为零、在补上的向量上为 \(1\)。

\begin{definition}[射影对偶空间]\label{g1:def:projective-dual}
	对于 \(K=\mathbb R\) 或 \(\mathbb C\) 上有限维向量空间 \(V\)，称 \(\mathbb P(V^*)\) 为 \(\mathbb P(V)\) 的\term[sheyingduioukongjian]{射影对偶空间}[dual projective space]。其中一点 \([\alpha]\) 参数化原空间的超平面 \(H_\alpha\)。点 \([v]\) 与超平面 \([\alpha]\) 的关联条件是 \(\alpha(v)=0\)。
\end{definition}

这个条件对 \(v,\alpha\) 的非零倍数都不变。在选定的基及其对偶基中，它是 \(\sum_i a_ix_i=0\)；但把点坐标列直接当成超平面系数列，需要额外选择，不能仅凭两列有同样长度就认作自然对应。真正不需要选基的是到双重对偶的求值映射
\[
	\iota:V\longrightarrow V^{**},\qquad
	\iota(v)(\alpha)=\alpha(v).
\]
若 \(v\ne0\)，把 \(v\) 扩充成基便可构造 \(\alpha(v)=1\) 的函数，所以 \(\iota\) 单射。取一组基可知 \(\dim V^*=\dim V\)，再用一次得到 \(\dim V^{**}=\dim V\)，故 \(\iota\) 是同构。因此双重射影对偶自然回到原空间。

\begin{proposition}[零化子与关联反转]\label{g1:prop:annihilator-duality}
	设 \(V\) 是 \(K=\mathbb R\) 或 \(\mathbb C\) 上的有限维向量空间，\(U,W\subseteq V\) 为子空间。定义 \(U\) 的\term[linghuazi]{零化子}[annihilator]
	\[
		U^\circ\coloneqq\{\,\alpha\in V^*: \alpha(u)=0\text{ 对一切 }u\in U\,\}.
	\]
	经由求值同构把 \(V^{**}\) 认作 \(V\)，有
	\[
		\begin{gathered}
			\dim U^\circ=\dim V-\dim U,\qquad (U^\circ)^\circ=U,\\
			(U+W)^\circ=U^\circ\cap W^\circ,\qquad
			(U\cap W)^\circ=U^\circ+W^\circ.
		\end{gathered}
	\]
	因此取零化子反转包含关系，也将射影子空间的张成与相交互换。
\end{proposition}
\begin{proof}
	把 \(U\) 的基 \(u_1,\ldots,u_k\) 补成 \(V\) 的基 \(u_1,\ldots,u_m\)。其对偶基中最后 \(m-k\) 个函数构成 \(U^\circ\) 的基，得到维数公式。每个 \(u\in U\) 被 \(U^\circ\) 中的所有函数零化，故 \(U\subseteq(U^\circ)^\circ\)；两边维数均为 \(k\)，于是相等。

	一个函数在 \(U+W\) 上为零，当且仅当它同时在 \(U,W\) 上为零，得到第一条运算公式。再在 \(V^*\) 中对 \(U^\circ,W^\circ\) 使用这条公式，得 \((U^\circ+W^\circ)^\circ=U\cap W\)；取第二次零化子便得到另一条。最后，\(U\subseteq W\) 直接推出 \(W^\circ\subseteq U^\circ\)。
\end{proof}
\symidx{\(U^\circ\)：子空间的零化子}

若 \(\dim\mathbb P(V)=n\)，则一个 \(k\) 维射影子空间对应一个 \(n-k-1\) 维对偶子空间。采用 \(\mathbb P(0)=\varnothing\)、维数 \(-1\) 的约定时，整个空间与空集也互相对应。在射影平面中，点与直线互换；两点的连线与两条对偶直线的交点互换。三个点共线，等价于它们的代表向量张成空间维数至多为 \(2\)，也等价于三个对偶直线具有公共点。这些结论来自零化子公式，不需要图形的位置猜测。

\subsection{非退化配对产生的极点、极超平面}
\label{g1:subsec:0204:02}

射影对偶把原空间与另一个空间联系起来。若还想在原空间内部给每一点指定一个超平面，就需要一个从 \(V\) 到 \(V^*\) 的同构。对称双线性型恰好提供这样的映射，而且它的对称性将变成几何上的互反关系。

\begin{definition}[二次型的极性]\label{g1:def:quadric-polarity}
	设 \(V\) 是 \(K=\mathbb R\) 或 \(\mathbb C\) 上维数至少为 \(2\) 的向量空间， \(B\) 是非退化对称双线性型。记
	\[
		b:V\longrightarrow V^*,\quad b(v)=B(v,-),\qquad
		U^{\perp_B}=\{\,w\in V:B(u,w)=0\text{ 对一切 }u\in U\,\}.
	\]
	点 \(p=[v]\) 的\term[jichaopingmian]{极超平面}[polar hyperplane]是 \(p^{\perp_B}=\mathbb P((Kv)^{\perp_B})\)；超平面 \(H_\alpha\) 的\term[jidian]{极点}[pole]是 \([b^{-1}\alpha]\)。这一反转关联的对应称为由 \(B\) 给出的\term[jixing]{极性}[polarity]。在射影平面中，极超平面通常称为\term[jixian]{极线}[polar line]。
\end{definition}

非退化的含义仍是 \(B(v,w)=0\) 对所有 \(w\) 成立时必有 \(v=0\)，与\cref{g1:def:symmetric-form}的实情形相同。它保证 \(b\) 单射；有限维且两端同维又保证 \(b\) 满射，所以上述逆映射确实存在。把 \(B\) 乘以非零常数不改变极性。这里只使用双线性型；在复数域上不作复共轭，不能把它换成带复共轭的内积而仍沿用相同公式。

\begin{proposition}[极性的互反性]\label{g1:prop:polarity}
	设 \(V\) 是 \(K=\mathbb R\) 或 \(\mathbb C\) 上 \(n+1\) 维向量空间，\(n\ge1\)， \(B\) 为非退化对称双线性型。对每个子空间 \(U\subseteq V\)，有
	\[
		\dim U^{\perp_B}=n+1-\dim U,\qquad
		(U^{\perp_B})^{\perp_B}=U.
	\]
	对点 \(p=[v]\)、\(q=[w]\)，有
	\[
		p\in q^{\perp_B}\quad\Longleftrightarrow\quad q\in p^{\perp_B}.
	\]
	点 \(p\) 落在自己的极超平面上，当且仅当 \(B(v,v)=0\)。
\end{proposition}
\begin{proof}
	对称性给出 \(U^{\perp_B}=b^{-1}(U^\circ)\)，故维数公式由\cref{g1:prop:annihilator-duality}得到。又有 \(U\subseteq(U^{\perp_B})^{\perp_B}\)，因为 \(B(u,w)=B(w,u)\)；两边维数相同，故相等。点的互反性正是 \(B(w,v)=0\) 与 \(B(v,w)=0\) 的等价，最后令 \(w=v\) 即得自关联条件。
\end{proof}

因此二次曲面 \(Q=\{\,[v]:B(v,v)=0\,\}\) 可以从极性中恢复为所有自关联点。反过来，给定二次方程及其对称配对便得到极性。这里仍应记住前两节对方程与点集的区别：例如正定实型的射影零点集为空，但它给出的极性存在，并且绝不由“空集”这一个集合唯一决定。

\begin{example}[圆的极线]\label{g1:exm:circle-polar}
	在 \(\mathbb P^2(\mathbb R)\) 中考虑 \(X^2+Y^2-Z^2=0\)。点 \([a:b:c]\) 的极线是 \(aX+bY-cZ=0\)。在仿射图 \(Z=1\) 中，点 \((2,0)\) 的极线是 \(x=1/2\)，圆心 \([0:0:1]\) 的极线则是无穷远直线 \(Z=0\)。无穷远点 \([1:0:0]\) 的极线为直径所在的直线 \(X=0\)。这些计算说明极线不要求通过原来的点；只有圆上的点才有这一性质。
\end{example}

\subsection{切空间、切超平面与对偶二次曲面}
\label{g1:subsec:0204:03}

为什么极线与圆锥曲线的切线有关？沿着一条过曲面上点的直线代入二次方程，常数项已经为零，一次项描述离开曲面的第一阶变化；令一次项为零，恰好得到极超平面的方程。下面把这个观察写成不依赖图像的局部计算。

\begin{proposition}[二次曲面的线性化]\label{g1:prop:quadric-tangent}
	设 \(K=\mathbb R\) 或 \(\mathbb C\)，\(\dim V=n+1\)、\(n\ge2\)， \(B\) 为对称双线性型，\(Q\subseteq\mathbb P(V)\) 由 \(q(x)=B(x,x)=0\) 定义。设 \(p=[v]\in Q\)，且 \(B(v,-)\ne0\)。则 \(Q\) 在 \(p\) 处的一阶方程给出唯一的射影超平面
	\[
		\mathrm{T}_pQ=\mathbb P\{\,w\in V:B(v,w)=0\,\}.
	\]
	选取 \(\alpha\in V^*\) 满足 \(\alpha(v)=1\)，在仿射图 \(\alpha=1\) 内，切向量组成的空间是
	\[
		\{\,u\in\ker\alpha:B(v,u)=0\,\},
	\]
	其 \(K\)-维数为 \(n-1\)，并通过商映射同构于 \(\ker B(v,-)/Kv\)。
\end{proposition}
\begin{proof}
	仿射图中以 \(v\) 为基点，每一点写成 \(v+u\)，其中 \(u\in\ker\alpha\)。代入得
	\[
		q(v+u)=2B(v,u)+q(u),
	\]
	因为 \(q(v)=0\)。因此一阶线性部分为 \(u\mapsto2B(v,u)\)。它在 \(\ker\alpha\) 上非零：否则 \(B(v,-)\) 是 \(\alpha\) 的倍数，把 \(v\) 代入得倍数为 \(B(v,v)=0\)，与假设矛盾。所以其核在 \(n\) 维方向空间 \(\ker\alpha\) 中有余维一。

	一阶仿射平面由 \(v+u\)、\(\alpha(u)=B(v,u)=0\) 组成。其射影张成空间为 \(\mathbb P(\ker B(v,-))\)：对任何 \(w\in\ker B(v,-)\)，分解 \(w=\alpha(w)v+(w-\alpha(w)v)\)，后项同时属于两个核。同一分解也说明商映射将切向量空间一一对应到 \(\ker B(v,-)/Kv\)。换另一个仿射图时仍得到同一个齐次超平面；把 \(v\) 换成非零倍数也不改变其核。
\end{proof}

这里将方程的一阶部分非零的点称为该二次方程的\term[guanghuadian]{光滑点}[smooth point]，将 \(\mathrm{T}_pQ\) 称为\term[qiechaopingmian]{切超平面}[tangent hyperplane]。上述仿射图中的方向空间是\term[qiekongjian]{切空间}[tangent space]的当前坐标描述；当 \(n=2\) 时，切超平面就是\term[qiexian]{切线}[tangent line]。我们在这里用多项式的一阶项定义这些对象，不需要预先调用隐函数定理。若一条过 \(p\) 的直线由 \(v,w\) 张成，则
\[
	q(v+tw)=2tB(v,w)+t^2q(w).
\]
它的方向落在切超平面中，恰好表示一次项消失；若二次项也消失，整条直线就在 \(Q\) 上。因此在含有直纹的二次曲面上，“切”不排除整条直线包含在曲面内。

\begin{theorem}[非退化二次曲面的对偶]\label{g1:thm:dual-quadric}
	设 \(K=\mathbb R\) 或 \(\mathbb C\)，\(n\ge2\)， \(A\) 为可逆的 \((n+1)\times(n+1)\) 对称矩阵， \(Q=\{\,[x]\in\mathbb P^n(K):x^{\mathsf T}Ax=0\,\}\)。以超平面系数 \([a]\) 为对偶坐标，所有 \(Q\) 的切超平面组成\term[duiouerciqumian]{对偶二次曲面}[dual quadric]
	\[
		Q^*=\{\,[a]\in\mathbb P((K^{n+1})^*):a^{\mathsf T}A^{-1}a=0\,\}.
	\]
	映射 \([x]\mapsto[Ax]\) 在 \(Q\) 与 \(Q^*\) 间给出双射；对 \(Q^*\) 再作同样的对偶构造，返回 \(Q\)。
\end{theorem}
\begin{proof}
	对 \([x]\in Q\)，可逆性保证 \(Ax\ne0\)，所以\cref{g1:prop:quadric-tangent}适用，其切超平面系数为 \(a=Ax\)。由对称性，
	\[
		a^{\mathsf T}A^{-1}a=x^{\mathsf T}Ax=0.
	\]
	反过来，若 \(a^{\mathsf T}A^{-1}a=0\)，令 \(x=A^{-1}a\)，则 \(x\ne0\)，且 \(x^{\mathsf T}Ax=0\)；在这个点的切超平面恰为 \([a]\)。两个映射互为逆，且将 \(A\) 换为 \(A^{-1}\) 再作一次便返回原方程。若实零点集为空，上述两个集合同时为空，这个集合等式和矩阵公式仍然成立。
\end{proof}

\begin{example}[切线的参数方程]\label{g1:exm:dual-conic}
	对圆锥曲线 \(XZ-Y^2=0\)，可取
	\[
		A=\begin{pmatrix}0&0&1/2\\0&-1&0\\1/2&0&0\end{pmatrix}.
	\]
	在点 \([s^2:st:t^2]\) 处的切线为 \(t^2X-2stY+s^2Z=0\)。所以切线系数 \([a:b:c]=[t^2:-2st:s^2]\) 满足 \(b^2-4ac=0\)。这也是由 \(a^{\mathsf T}A^{-1}a=0\) 得到的方程。一族切线之所以能成为对偶平面中的曲线，是因为每条线由一个系数点表示；两边的“点”分别参数化不同种类的几何对象。
\end{example}

\subsection{退化配对中极性的失效}
\label{g1:subsec:0204:04}

退化时不应只在非退化公式中删掉一个逆矩阵符号。若 \(R=\operatorname{rad}B\ne0\)，映射 \(b:v\mapsto B(v,-)\) 的核为 \(R\)，像为 \(R^\circ\)：像包含于 \(R^\circ\) 是对称性的结果，两者维数都等于 \(\dim V-\dim R\)。因此 \([v]\in\mathbb P(R)\) 对应零函数，根本没有极超平面；在其余点，\([v]\mapsto[B(v,-)]\) 仍有定义，但 \(v\) 与 \(v+r\) 给出相同函数，通常不再单射。得到的每个超平面还都包含 \(\mathbb P(R)\)，故也不能遍历全部超平面。

\begin{proposition}[退化型的双重正交补]\label{g1:prop:degenerate-polarity}
	设 \(V\) 是 \(K=\mathbb R\) 或 \(\mathbb C\) 上的有限维向量空间， \(B\) 为对称双线性型，\(R=\operatorname{rad}B\)。则对每个子空间 \(U\subseteq V\)，
	\[
		(U^{\perp_B})^{\perp_B}=U+R.
	\]
	在商空间 \(V/R\) 上，公式 \(\overline B(v+R,w+R)=B(v,w)\) 定义非退化对称双线性型。
\end{proposition}
\begin{proof}
	改变代表向量只会增加根基中的向量，所有新增配对项为零，所以 \(\overline B\) 良定义。若 \(\overline B(v+R,w+R)=0\) 对所有 \(w\) 成立，则 \(v\in R\)，证明非退化性。记商映射为 \(\pi:V\rightarrow V/R\)。由定义， \(U^{\perp_B}=\pi^{-1}((\pi U)^{\perp_{\overline B}})\)。在商空间使用\cref{g1:prop:polarity}的双重正交补结论，再取逆像，得到 \((U^{\perp_B})^{\perp_B}=\pi^{-1}(\pi U)=U+R\)。商空间维数为 \(0\) 或 \(1\) 时，所用向量空间公式仍由同一维数证明成立。
\end{proof}

\begin{example}[锥顶与重合的切平面]\label{g1:exm:singular-polarity}
	在 \(\mathbb P^3(\mathbb R)\) 中，令 \(q(X,Y,Z,W)=X^2+Y^2-Z^2\)。根基为 \(\mathbb R(0,0,0,1)\)，锥顶 \([0:0:0:1]\) 处一次项全为零，没有由一阶方程确定的切平面。对每个 \(w\in\mathbb R\)，点 \([1:0:1:w]\) 在二次曲面上，其切平面都为 \(X-Z=0\)。这些点沿一条母线变化，切平面却不变，展示了对偶映射丢失的自由度。
\end{example}

更极端地，在射影平面中，方程 \(X^2=0\) 的实点集和复点集都是直线 \(X=0\)，但作为这个二次方程，它在每个点的一阶部分都为零。若把它先改成一次方程 \(X=0\)，当然得到一条光滑直线；这已经更换了定义方程所保留的信息。所以本节对退化对象的线性化始终针对给定的二次方程，不能仅从裸点集判断其一阶结构。

% !TeX root = ../../Book1.tex
% 纲要标题：### 2.5 圆锥曲线上的构造与定理
% 纲要位置：几何学计划纲要.md，第 185 行。
% 正文状态：初稿；尚未正式审核。

% 纲要内容开始
% 
% * 由投影得到有理参数及其逆
% * 圆锥曲线上的交比
% * Pascal 定理的坐标证明
% * Brianchon 定理与对偶推导
% 
% 纲要内容结束

\section{圆锥曲线上的构造与定理}
\label{g1:sec:0205}

射影分类告诉我们哪些圆锥曲线可以通过换坐标互相变换，但还没有说明怎样在一条曲线上移动和构造点。一个可操作的办法是固定曲线上一点，让过它的直线转动，用直线与曲线的另一个交点作参数。这样，圆锥曲线上的计算转化为射影直线上的计算；交比因此有了新的适用对象，而六个点的关联关系也能压缩为含三个参数的多项式恒等式。

\subsection{由投影得到有理参数及其逆}
\label{g1:subsec:0205:01}

本节讨论 \(K=\mathbb R\) 或 \(\mathbb C\) 上的非退化圆锥曲线，即射影平面中由可逆对称矩阵给出的二次方程。实数域上还须要求存在实点；例如 \(X^2+Y^2+Z^2=0\) 没有实射影点，不能从 \(\mathbb P^1(\mathbb R)\) 参数化一个空集。由\cref{g1:thm:projective-quadric-classification}，复数域上的非退化三元型都等价于 \(X^2+Y^2+Z^2\)，它含有点 \([1:\mathrm i:0]\)，所以在复数域上非空条件自动满足。下面仍把一个曲线点作为构造的输入，从它直接找出参数坐标。

\begin{proposition}[圆锥曲线的参数化]\label{g1:prop:conic-parametrization}
	设 \(K=\mathbb R\) 或 \(\mathbb C\)，\(C\subseteq\mathbb P^2(K)\) 是非退化圆锥曲线，且有一个 \(K\)-点。则存在射影坐标使 \(C\) 的方程为 \(XZ-Y^2=0\)。在这些坐标中，映射
	\begin{equation}\label{g1:eq:conic-parametrization}
		\nu:\mathbb P^1(K)\longrightarrow C,
		\qquad [s:t]\longmapsto[s^2:st:t^2]
	\end{equation}
	是双射。逆映射在 \(Z\ne0\) 上为 \([X:Y:Z]\mapsto[Y:Z]\)，在 \(X\ne0\) 上为 \([X:Y:Z]\mapsto[X:Y]\)，二者在交集上一致。
\end{proposition}
\begin{proof}
	先说明坐标为何存在。写 \(q(v)=B(v,v)\)，取非零 \(v\) 满足 \(q(v)=0\)。非退化性使我们能选 \(w\) 满足 \(B(v,w)=1\)。令 \(w'=w-\frac12q(w)v\)，则 \(q(w')=0\)、\(B(v,w')=1\)。在 \(H=\operatorname{span}(v,w')\) 上，配对矩阵为 \(\left(\begin{smallmatrix}0&1\\1&0\end{smallmatrix}\right)\)，可逆。因此每个向量可唯一减去 \(v,w'\) 的线性组合而变成与 \(H\) 正交的向量：两个待定系数由这一个可逆矩阵确定。于是 \(V=H\oplus H^{\perp_B}\)，且后一空间为一维。取其非零生成元 \(u\)，必有 \(a=q(u)\ne0\)，否则 \(u\) 与整个 \(V\) 正交，违反非退化性。故
	\[
		q(\alpha v+\beta w'+\gamma u)=2\alpha\beta+a\gamma^2.
	\]
	置 \(\alpha=X\)、\(\beta=-aZ/2\)、\(\gamma=Y\)，零点方程就变成 \(XZ-Y^2=0\)。

	接着核对参数映射。\(s,t\) 不全为零时三个二次单项式不全为零，把 \((s,t)\) 乘 \(\lambda\ne0\) 只将像乘 \(\lambda^2\)，故良定义，且像满足曲线方程。若 \(Z\ne0\)，方程给出
	\[
		[X:Y:Z]=[(Y/Z)^2:Y/Z:1]=\nu([Y:Z]).
	\]
	若 \(Z=0\)，必有 \(Y=0\)，该点为 \([1:0:0]=\nu([1:0])\)。在 \(X\ne0\) 的图上同样得到逆公式 \([X:Y]\)。两图覆盖 \(C\)，因为 \(X=Z=0\) 会迫使所有坐标为零；交集上 \([X:Y]=[Y:Z]\) 正是 \(XZ=Y^2\)。这证明满射和逆映射的唯一性。
\end{proof}

在 \(t\ne0\) 的参数图上置 \(r=s/t\)，便有 \(\nu(r)=[r^2:r:1]\)；缺少的一个点记作 \(\nu(\infty)=[1:0:0]\)。这种用参数的有理函数表示坐标的方式称为\term[youlicanshuhua]{有理参数化}[rational parametrization]。两个参数图各自有公式，在交集相容，共同覆盖整条曲线。

\begin{figure}[htbp]
	\centering
	\input{Figures/Chapter02-ConicProjection}
	\caption[圆上的投影参数]{圆上的投影参数。固定 \(P=(-1,0)\)，过 \(P\) 的直线由斜率 \(m\) 标记；第二交点给出圆的参数化。切线 \(x=-1\) 对应参数直线上的无穷远点，并在延拓后对应 \(P\) 本身。}
	\label{g1:fig:conic-projection}
\end{figure}

这个双射正来自投影。对标准曲线固定 \(O=\nu(\infty)\)，过 \(O\) 的直线写成 \(tY-sZ=0\)，由 \([s:t]\in\mathbb P^1(K)\) 参数化。若 \(t\ne0\)，该直线与曲线相交于 \(O\) 和 \(\nu([s:t])\)；若 \(t=0\)，它成为 \(Z=0\)，即 \(O\) 处的切线。所以“另一个交点”的构造原先漏掉了 \(O\)，补上切线所对应的参数后，才得到整条曲线上的双射。

\begin{example}[单位圆的有理公式]\label{g1:exm:circle-parametrization}
	在实平面取 \(P=(-1,0)\)，直线 \(y=m(x+1)\) 与 \(x^2+y^2=1\) 相交。代入得到 \((x+1)((1+m^2)x+m^2-1)=0\)。把已知根 \(x=-1\) 提出后，另一个根为
	\[
		x=\frac{1-m^2}{1+m^2},\qquad y=\frac{2m}{1+m^2}.
	\]
	反过来，圆上除 \(P\) 外的点满足 \(m=y/(x+1)\)。参数 \(m=\infty\) 对应竖直切线，并补回点 \(P\)，如\cref{g1:fig:conic-projection}所示。这比三角函数参数更适合只涉及直线求交和代数计算的问题；两种参数表达同一个圆，却保留不同的计算便利。
\end{example}

\subsection{圆锥曲线上的交比}
\label{g1:subsec:0205:02}

先观察一条曲线上的四个点怎样变成射影直线上的四个点。过固定点 \(P\) 的全部直线称为以 \(P\) 为中心的\term[xianshu]{线束}[pencil of lines]。若 \(P=[p]\)，这些直线系数构成二维子空间 \((Kp)^\circ\) 的射影化，所以线束本身就是一条射影直线。这说明可以对四条线定义交比，但还须证明结果与投影中心无关。

\begin{proposition}[曲线交比的选择无关性]\label{g1:prop:conic-cross-ratio}
	设 \(K=\mathbb R\) 或 \(\mathbb C\)，\(C\subseteq\mathbb P^2(K)\) 是有 \(K\)-点的非退化圆锥曲线。通过\cref{g1:prop:conic-parametrization}的参数化，把四个两两不同的点 \(A,B,C_0,D\) 拉回 \(\mathbb P^1(K)\)，所得交比与所选标准射影坐标无关。它也等于从曲线上任一点 \(P\) 投影得到的四条线的交比，其中 \(P\) 自己的像解释为 \(P\) 处切线。射影变换保持这个数。
\end{proposition}
\begin{proof}
	关键是证明投影给出的参数变化为射影变换，而不只是集合上的双射。在标准曲线上令 \(P=\nu([a:b])\)。过 \(P\) 与 \(\nu([s:t])\) 的直线系数可以取为
	\begin{equation}\label{g1:eq:conic-pencil-map}
		[bt:-(at+bs):as].
	\end{equation}
	代入两个点的坐标，分别得到零；当 \([s:t]=[a:b]\) 时，式子变为 \([b^2:-2ab:a^2]\)，正是\cref{g1:exm:dual-conic}中的切线系数。这里从 \((s,t)\) 到三个系数的映射是线性的，且核为零：若 \(a,b\) 都非零，由第一和第三项即得 \(s=t=0\)；若其中一个为零，则另一个非零，由中间项和剩余的一端同样得到 \(s=t=0\)。其像位于 \(P\) 对应的二维零化子，因而恰好是整个零化子。故从参数直线到线束的对应是一个射影同构，保持\cref{g1:def:cross-ratio}所定义的交比。

	现在改变平面射影坐标，或用射影变换 \(S\) 将 \(C\) 送到另一条非退化曲线。选取 \(S\) 的一个可逆线性代表，仍记为 \(S\)。它将过 \(P\) 的线束线性地送到过 \(SP\) 的线束：若直线系数为 \(\alpha\)，变换后的系数为 \(\alpha\circ S^{-1}\)。由上一段，两个线束各自与其参数直线射影同构。把这三个同构复合，得到两个参数之间的射影变换。根据\cref{g1:prop:cross-ratio-orbits}的交比不变性，所得数与坐标和中心都无关，也被 \(S\) 保持。
\end{proof}

这个数称为\term[yuanzhuiquxianjiaobi]{圆锥曲线交比}[cross-ratio on a conic]，记为 \((A,B;C_0,D)_C\)。它仍然依赖四点的次序，置换规律与射影直线完全相同。若一个投影中心位于曲线外部，投影通常会把一条割线的两个交点合并，便不再给出上述射影同构；命题中特意要求中心在曲线上。

\begin{example}[计算与对称性]\label{g1:exm:conic-cross-ratio}
	在 \(XZ-Y^2=0\) 上取参数为 \(\infty,0,1,-1\) 的四点。按本章约定，\((\infty,0;1,-1)=-1\)，所以这四点在圆锥曲线上的交比也为 \(-1\)。无论从曲线上的哪一点投影，只要将中心点自身解释为切线，所得四条线都有同样的有序交比。换成参数 \(\infty,0,1,2\) 时交比为 \(2\)，故不存在保持曲线、并将前一有序四点组逐点送到后一组的射影变换。
\end{example}

\subsection{Pascal 定理的坐标证明}
\label{g1:subsec:0205:03}

把曲线上的三个点规范化，可以把六点构造的自由参数减到三个。所需的射影直线变换可以通过二次单项式提升到整个射影平面。若
\[
	G=\begin{pmatrix}a&b\\c&d\end{pmatrix}\in\operatorname{GL}_2(K),
\]
则由 \((s,t)\mapsto(as+bt,cs+dt)\) 得到平面线性变换
\[
	\widetilde G=
	\begin{pmatrix}
		a^2&2ab&b^2\\ ac&ad+bc&bd\\ c^2&2cd&d^2
	\end{pmatrix},\qquad
	\widetilde G\begin{pmatrix}s^2\\st\\t^2\end{pmatrix}
	=\begin{pmatrix}(as+bt)^2\\(as+bt)(cs+dt)\\(cs+dt)^2\end{pmatrix}.
\]
代入两次线性变换表明 \(\widetilde{G_1G_2}=\widetilde G_1\widetilde G_2\)，故 \(\widetilde G\) 的逆为 \(\widetilde{G^{-1}}\)。它诱导保持标准圆锥曲线的平面射影变换。由\cref{g1:thm:projective-frame}在射影直线上的情形，任意有序三点可送到 \(\infty,0,1\)，因而曲线上的三点也能通过这种平面变换规范化。

下面的\term[Pascaldingli]{Pascal 定理}[Pascal's theorem]描述六个曲线点所产生的三个交点。\footnote{Blaise Pascal（1623--1662），法国数学家、物理学家与哲学家，对射影几何和概率论有重要贡献。人物资料见University of St Andrews \href{https://mathshistory.st-andrews.ac.uk/Biographies/Pascal/}{MacTutor 传记}。}“六边形”仅表示六个点的循环次序，不要求它在某张实仿射图中凸，也不要求各边只在顶点相交。

\begin{theorem}[Pascal]\label{g1:thm:pascal}
	设 \(K=\mathbb R\) 或 \(\mathbb C\)，\(A,B,C,D,E,F\) 是 \(\mathbb P^2(K)\) 中一条非退化圆锥曲线上的六个两两不同的点。则
	\[
		U=AB\cap DE,\qquad V=BC\cap EF,\qquad W=CD\cap FA
	\]
	共线。这里 \(AB\) 表示穿过 \(A,B\) 的射影直线，交点允许位于所选仿射图的无穷远处。
\end{theorem}
\begin{proof}
	证明的想法是先用射影变换消去三个参数，再直接写出经过三个交点的直线。由\cref{g1:prop:conic-parametrization}及上面的提升构造，可以设曲线为 \(XZ-Y^2=0\)，且六个点的参数依次为
	\[
		\infty,\quad0,\quad1,\quad d,\quad e,\quad f.
	\]
	六点不同保证 \(d,e,f\in K\setminus\{0,1\}\) 且两两不同。这一步没有舍弃任何构形：唯一被送到参数无穷远处的点就是 \(A\)。

	先核对所需交点都有定义。任一平面直线 \(aX+bY+cZ=0\) 在参数曲线上的逆像满足 \(as^2+bst+ct^2=0\)。若直线系数不全为零，这个二元齐次二次式最多有两个不同的射影零点：在 \(t\ne0\) 的图上是次数至多为二的一元方程；若 \([1:0]\) 也是零点，则 \(a=0\)，有限图中次数至多为一。所以曲线上没有三个不同点共线，各对需要相交的边也不重合。

	参数为 \(r,s\) 的两个有限点的连线方程是 \(X-(r+s)Y+rsZ=0\)，参数为 \(\infty,r\) 的连线为 \(Y-rZ=0\)。分别求交得到
	\[
		\begin{aligned}
			U&=[-de:0:1],\\
			V&=[ef:ef:e+f-1],\\
			W&=[(1+d)f-d:f:1].
		\end{aligned}
	\]
	这些齐次坐标都非零；特别地，即使 \(e+f-1=0\)，第二个点仍由 \(ef\ne0\) 保证有效。 \(U,V\) 不同，因为前者的 \(Y\) 坐标为零而后者不为零。它们的连线是
	\begin{equation}\label{g1:eq:pascal-line}
		-fX+(de+df-d+f)Y-defZ=0.
	\end{equation}
	将 \(U,V\) 的坐标代入可分别核对为零；对 \(W\)，左边为
	\[
		-f\bigl((1+d)f-d\bigr)+(de+df-d+f)f-def=0,
	\]
	各项展开即两两抵消。由于 \(f\ne0\)，这确是一条直线，故三个交点共线。再用最初射影变换的逆变换恢复原构形，结论仍成立。
\end{proof}

\begin{figure}[htbp]
	\centering
	\input{Figures/Chapter02-Pascal}
	\caption[Pascal 的六点构形]{Pascal 的六点构形。三组同色的对边分别交于 \(U,V,W\)，红色虚线是它们的公共直线。六点的循环次序允许产生自交边。}
	\label{g1:fig:pascal}
\end{figure}

\cref{g1:fig:pascal}取单位圆上的
\[
	\begin{aligned}
		A&=(-1,0),&\qquad B&=(0,-1),\\
		C&=(0,1),& D&=(-3/5,-4/5),\\
		E&=(-3/5,4/5),& F&=(1,0).
	\end{aligned}
\]
这些点由\cref{g1:exm:circle-parametrization}的参数 \(\infty,-1,1,-2,2,0\) 给出。直线求交得到 \(U=(-3/5,-2/5)\)、\(V=(0,1/2)\)、\(W=(-1/3,0)\)，三点满足 \(y=3x/2+1/2\)。图中六边形虽有自交，关联关系仍然清楚。

\begin{example}[一个交点在无穷远处]\label{g1:exm:pascal-infinity}
	在定理的标准坐标中取 \(d=2\)、\(e=3\)、\(f=-2\)。则
	\[
		U=[-6:0:1],\qquad V=[-6:-6:0],\qquad W=[-8:-2:1].
	\]
	三个点都满足 \(X-Y+6Z=0\)。在仿射图 \(Z=1\) 中，\(BC\) 与 \(EF\) 是平行线，其交点 \(V\) 位于无穷远直线；因此仿射图像中的平行情形已经包含在同一个定理和证明里。
\end{example}

六点两两不同的条件使每条割线都有唯一含义。若相邻点重合，相应的边须改用切线，并要求剩下的交点仍有定义。

\subsection{Brianchon 定理与对偶推导}
\label{g1:subsec:0205:04}

\cref{g1:thm:pascal}涉及曲线上的点、两点连线、两线交点与共线。在对偶平面中，曲线点变成原平面中的切线，连线与交点互换，共线则变成共点。由此得到\term[Brianchondingli]{Brianchon 定理}[Brianchon's theorem]。\footnote{Charles-Julien Brianchon（1783--1864），法国数学家，研究射影几何与圆锥曲线。人物资料见University of St Andrews \href{https://mathshistory.st-andrews.ac.uk/Biographies/Brianchon/}{MacTutor 传记}。}下面逐一写出点、直线与关联关系的对偶对应。

\begin{table}[htbp]
	\centering
	\caption{Pascal 与 Brianchon 构形的对应}
	\label{g1:tab:pascal-brianchon}
	\begin{tabularx}{\textwidth}{@{}XXX@{}}
		\toprule
		Pascal 构形中的对象 & 对偶运算 & Brianchon 构形中的对象 \\
		\midrule
		圆锥曲线上的六点 & 点表示原平面的直线 & 六条不同的曲线切线 \\
		相邻两点的连线 & 张成与相交互换 & 相邻两条切线的交点 \\
		三组对边的交点共线 & 共线与共点互换 & 三条对角线通过公共点 \\
		\bottomrule
	\end{tabularx}
\end{table}

\cref{g1:tab:pascal-brianchon}列出的运算都来自零化子；证明还需要确认六条切线确实落在同一条对偶圆锥曲线上。

\begin{theorem}[Brianchon]\label{g1:thm:brianchon}
	设 \(K=\mathbb R\) 或 \(\mathbb C\)，\(\ell_1,\ldots,\ell_6\) 是 \(\mathbb P^2(K)\) 中一条非退化圆锥曲线的六条两两不同的切线。按循环次序置
	\[
		A=\ell_1\cap\ell_2,\quad B=\ell_2\cap\ell_3,\quad C=\ell_3\cap\ell_4,
		\quad D=\ell_4\cap\ell_5,\quad E=\ell_5\cap\ell_6,\quad F=\ell_6\cap\ell_1.
	\]
	则三条对角线 \(AD,BE,CF\) 共点。
\end{theorem}
\begin{proof}
	把每条 \(\ell_i\) 看作对偶平面中的点 \(L_i\)。由\cref{g1:thm:dual-quadric}，这六个不同的点位于一条非退化对偶圆锥曲线上。该曲线上没有三个不同点共线，这已在\cref{g1:thm:pascal}的证明中由参数二次式说明；由\cref{g1:prop:annihilator-duality}，原平面中的三条不同切线因而不共点。于是所列相邻交点存在，且 \(A\ne D\)、\(B\ne E\)、\(C\ne F\)，各对角线都有定义。

	在对偶平面对 \(L_1,\ldots,L_6\) 使用\cref{g1:thm:pascal}，得到
	\[
		U^*=L_1L_2\cap L_4L_5,\qquad
		V^*=L_2L_3\cap L_5L_6,\qquad
		W^*=L_3L_4\cap L_6L_1
	\]
	共线。直线 \(L_1L_2\) 是原平面中点 \(A\) 的对偶直线，\(L_4L_5\) 是原平面中点 \(D\) 的对偶直线；两条对偶直线的交点 \(U^*\) 因而对应原直线 \(AD\)。同理 \(V^*,W^*\) 分别对应 \(BE,CF\)。包含 \(U^*,V^*,W^*\) 的对偶直线对应原平面的一个点，关联反转保证该点同时属于 \(AD,BE,CF\)，即得结论。
\end{proof}

例如，在实平面取一条圆和一个外切六边形，三条对角线即使看上去没有明显对称性，仍然共点。定理不要求六条切线围成有界的凸区域，公共点也可能位于无穷远处；这些情形与前面允许割线平行的情形相对应。这说明对偶真正保留的是关联及其由线性方程表达的逻辑关系，长度、角度和“内外”都不是推导所需的数据。

% !TeX root = ../../Book1.tex
% 纲要标题：### 2.6 射影映射与参数空间
% 纲要位置：几何学计划纲要.md，第 192 行。
% 正文状态：初稿；尚未正式审核。

% 纲要内容开始
% 
% * Segre 与二次 Veronese 映射的构造
% * Grassmann 空间与 Plücker 坐标
% * 点与超平面的关联簇及其投影
% 
% ---
% 
% 纲要内容结束

\section{射影映射与参数空间}
\label{g1:sec:0206}

圆锥曲线的参数化已经展示过一种变化：把一个点的坐标写成参数的二次式，便把一条射影线放进了另一个射影空间。还可以反过来问：一对射影点、一条直线或一个超平面，能否也用一个射影点来记录？合适的坐标会把“来自同一个对象”的相容条件写成方程，从而让整族几何对象本身成为可以研究的几何集合。本节始终取 \(K=\mathbb R\) 或 \(\mathbb C\)，只用多项式、矩阵秩与坐标图完成这些构造。

\subsection{Segre 与二次 Veronese 映射的构造}
\label{g1:subsec:0206:01}

先考虑两个独立的射影点 \([x_0:\cdots:x_m]\) 和 \([y_0:\cdots:y_n]\)。直接把坐标并排写成 \([x_0:\cdots:x_m:y_0:\cdots:y_n]\) 并不合法：改变第一组的代表而保持第二组不变，通常会改变所得射影点。若改用所有乘积 \(x_i y_j\)，两组代表分别乘以 \(\lambda\) 和 \(\mu\) 时，全部新坐标都乘以同一个 \(\lambda\mu\)。这正是齐次坐标允许的变化。

\begin{definition}[Segre 映射]\label{g1:def:segre-map}
	设 \(K=\mathbb R\) 或 \(\mathbb C\)，\(m,n\ge0\)。映射
	\[
		\sigma_{m,n}:\mathbb P^m(K)\times\mathbb P^n(K)
		\longrightarrow\mathbb P^{(m+1)(n+1)-1}(K),\qquad
		([x],[y])\longmapsto[x_i y_j]_{0\le i\le m,\,0\le j\le n}
	\]
	称为\term[Segreyingshe]{Segre 映射}[Segre map]，目标坐标先按 \(i\)、再按 \(j\) 排列。\footnote{Corrado Segre（1863--1924），意大利数学家，研究高维射影几何、代数曲面和直纹几何。}
\end{definition}

\begin{proposition}[乘积坐标与秩一矩阵]\label{g1:prop:segre-rank-one}
	设 \(K=\mathbb R\) 或 \(\mathbb C\)，\(m,n\ge0\)。Segre 映射良定义且为单射。把目标点的坐标排成 \((m+1)\times(n+1)\) 矩阵 \(Z=(z_{ij})\)，则其像恰由方程
	\[
		z_{ij}z_{k\ell}-z_{i\ell}z_{kj}=0
		\qquad(0\le i,k\le m,\ 0\le j,\ell\le n)
	\]
	给出，也就是所有非零秩一矩阵的射影类。在 \(z_{ab}\ne0\) 的部分，原来两个点可由第 \(b\) 列和第 \(a\) 行恢复。
\end{proposition}
\begin{proof}
	选取非零坐标 \(x_a,y_b\)，则 \(x_a y_b\ne0\)，所以输出坐标不全为零；代表的缩放只产生共同倍数，因而映射良定义。其矩阵是列向量与行向量的乘积 \(Z=xy^{\mathsf T}\)，每个二阶子式均为零。

	反过来，设 \(Z\ne0\) 且全部二阶子式为零，选定 \(z_{ab}\ne0\)。含第 \(a,i\) 行和第 \(b,j\) 列的子式给出
	\[
		z_{ij}=\frac{z_{ib}z_{aj}}{z_{ab}}.
	\]
	因此取 \(x_i=z_{ib}\)、\(y_j=z_{aj}/z_{ab}\)，便有 \(Z=xy^{\mathsf T}\)。这里 \(x_a\ne0\)、\(y_b=1\)，两向量均非零，证明了像的描述。对这样的分解，第 \(b\) 列与 \(x\) 成比例，第 \(a\) 行与 \(y^{\mathsf T}\) 成比例，故它们的射影类唯一，映射为单射。更具体地，在这张坐标图上，逆映射使用的仿射坐标为
	\[
		\frac{x_i}{x_a}=\frac{z_{ib}}{z_{ab}},\qquad
		\frac{y_j}{y_b}=\frac{z_{aj}}{z_{ab}}.
	\]
	分母在所选图上不为零，故这些公式也说明不同图中的恢复结果相容。
\end{proof}

\begin{example}[乘积中的两个方向]\label{g1:exm:segre-fibers}
	取 \(m=n=1\)，则
	\[
		([s:t],[u:v])\longmapsto[su:sv:tu:tv]
	\]
	将 \(\mathbb P^1(K)\times\mathbb P^1(K)\) 一一对应到 \(z_{00}z_{11}-z_{01}z_{10}=0\) 的二次曲面。固定 \([s:t]\) 而改变 \([u:v]\)，像是由 \((s,0,t,0)\)、\((0,s,0,t)\) 张成的射影直线；这两个向量线性无关。固定第二个因子同样得到一条直线。同一族中两条不同直线不相交，因为单射性不允许第一因子或第二因子同时取两个不同值；两族各取一条则恰交于相应有序对的像。\cref{g1:prop:quadric-two-rulings}中的两族直纹，在这里变成了乘积的两个独立方向。
\end{example}

还有另一种坐标构造。若希望把一个点的所有二次表达式同时作为坐标，就应取 \(x_i x_j\) 而不是取两组互不相关的变量。由此得到二次 Veronese 映射。\footnote{Giuseppe Veronese（1854--1917），意大利数学家，研究高维射影几何与几何基础。}

\begin{definition}[二次 Veronese 映射]\label{g1:def:veronese-map}
	设 \(K=\mathbb R\) 或 \(\mathbb C\)，\(n\ge1\)，\(N=\binom{n+2}{2}-1\)。映射
	\[
		\nu_2:\mathbb P^n(K)\longrightarrow\mathbb P^N(K),\qquad
		[x_0:\cdots:x_n]\longmapsto[x_i x_j]_{0\le i\le j\le n}
	\]
	称为\term[Veroneseyingshe]{二次 Veronese 映射}[quadratic Veronese map]。目标坐标按指标 \((i,j)\) 排列，将 \(z_{ij}\) 对称地补成矩阵 \(Z=(z_{ij})\)。
\end{definition}

\begin{proposition}[二次单项式坐标]\label{g1:prop:veronese-image}
	设 \(K=\mathbb R\) 或 \(\mathbb C\)，\(n\ge1\)。二次 Veronese 映射良定义且为单射，其像恰为非零对称秩一矩阵的射影类，亦即对称矩阵 \(Z\) 的所有二阶子式同时为零的射影点集。在 \(z_{aa}\ne0\) 的部分，逆映射由 \(x_i/x_a=z_{ia}/z_{aa}\) 给出。
\end{proposition}
\begin{proof}
	非零向量至少有一个非零坐标 \(x_a\)，故 \(x_a^2\ne0\)；缩放 \(x\) 使全部二次单项式乘以相同的平方倍数。这证明良定义，且像矩阵 \(xx^{\mathsf T}\) 为对称秩一矩阵。

	设对称矩阵 \(Z\ne0\) 的所有二阶子式均为零。必有某个对角元 \(z_{aa}\ne0\)：否则对任意 \(i,j\)，主子式 \(z_{ii}z_{jj}-z_{ij}^2=0\) 会强迫 \(z_{ij}=0\)，与 \(Z\ne0\) 矛盾。令 \(x_i=z_{ia}/z_{aa}\)，于是 \(x_a=1\)，而子式关系给出
	\[
		\frac{z_{ij}}{z_{aa}}
		=\frac{z_{ia}}{z_{aa}}\frac{z_{ja}}{z_{aa}}=x_i x_j.
	\]
	所以 \([Z]=[xx^{\mathsf T}]\)，并且上述比值唯一恢复 \([x]\)。这同时证明了像的描述、单射性及逆坐标公式。
\end{proof}

在实数情形，上一证明没有要求 \(z_{aa}\) 有实平方根。因为讨论的是矩阵的射影类，可以先把整个矩阵除以 \(z_{aa}\)；这也解释了为什么负半定的秩一矩阵仍表示像中的点。当 \(n=1\) 时，\(\nu_2([s:t])=[s^2:st:t^2]\)，正是前一节的圆锥曲线参数化。当 \(n=2\) 时，得到 \(\mathbb P^5(K)\) 中的 Veronese 曲面。

\ProseParagraphBreak

这两种映射还有一个直接用途：使原来的高次方程成为新空间的线性方程。例如任意二次齐次式 \(q(x)=\sum_{i\le j}c_{ij}x_i x_j\) 都满足
\[
	q(x)=0\quad\Longleftrightarrow\quad
	\sum_{i\le j}c_{ij}z_{ij}=0\quad\text{在 }[z]=\nu_2([x])\text{ 上成立}.
\]
右边是一个超平面与 Veronese 像的相交条件。同样，双线性关系 \(\sum a_{ij}x_i y_j=0\) 经 Segre 映射成为超平面截面。目标空间中的子式关系刻画了哪些射影点属于这些映射的像。这些模型的代数几何推广可参看\cite[第I章, 第2, Exercises节 2.12--2.15]{Hartshorne1977}。

\subsection{Grassmann 空间与 Plücker 坐标}
\label{g1:subsec:0206:02}

射影空间以一维向量子空间为点。若把“维数一”换成“维数 \(k\)”，同样得到一个值得研究的对象集合。困难在于：一个子空间有许多组基，把基向量的坐标直接抄下来会记录过多的选择。行列式恰能把换基的影响压缩成一个共同倍数。

\begin{definition}[Grassmann 空间]\label{g1:def:grassmannian}
	设 \(K=\mathbb R\) 或 \(\mathbb C\)，\(V\) 为 \(n\) 维 \(K\)-向量空间，\(1\le k\le n\)。所有 \(k\) 维向量子空间组成的集合
	\[
		\operatorname{Gr}(k,V)=\{\,W\subseteq V:W\text{ 为向量子空间且 }\dim_K W=k\,\}
	\]
	称为\term[Grassmannkongjian]{Grassmann 空间}[Grassmannian]，当 \(V=K^n\) 时也记为 \(\operatorname{Gr}(k,n;K)\)。\footnote{Hermann Günter Grassmann（1809--1877），德国数学家，发展了向量运算和外代数，并以坐标研究不同维数的子空间。}
	\symidx{\(\operatorname{Gr}(k,V)\)：Grassmann 空间}
\end{definition}

这里的 \(k,n\) 都是向量空间维数；\(W\) 对应 \(\mathbb P(V)\) 中的 \((k-1)\) 维射影子空间。特别地，\(\operatorname{Gr}(1,V)=\mathbb P(V)\)，而 \(\operatorname{Gr}(2,4;K)\) 参数化射影三维空间中的所有直线。先把后一情形具体写出。选取二维子空间的有序基 \(u,v\)，将它们写成矩阵的两行，并定义六个子式
\[
	p_{ij}=u_i v_j-u_j v_i\qquad(1\le i<j\le4).
\]
它们称为该直线的\term[Pluckerzuobiao]{Plücker 坐标}[Plücker coordinates]。\footnote{Julius Plücker（1801--1868），德国数学家和物理学家，研究解析几何、直线几何及光谱。}矩阵的秩为 \(2\)，故至少一个二阶子式非零。若将有序基替换为 \((u',v')^{\mathsf T}=S(u,v)^{\mathsf T}\)，其中 \(S\in\operatorname{GL}_2(K)\)，每个子式都乘以 \(\det S\)，所以射影点 \([p_{12}:p_{13}:p_{14}:p_{23}:p_{24}:p_{34}]\) 与基的选择无关。

\begin{theorem}[射影三维空间中直线的坐标]\label{g1:thm:plucker-gr24}
	设 \(K=\mathbb R\) 或 \(\mathbb C\)。Plücker 坐标给出 \(\operatorname{Gr}(2,4;K)\) 到 \(\mathbb P^5(K)\) 的单射，其像恰为二次曲面
	\begin{equation}\label{g1:eq:plucker-gr24}
		p_{12}p_{34}-p_{13}p_{24}+p_{14}p_{23}=0.
	\end{equation}
	在 \(p_{12}\ne0\) 的部分，每个子空间有唯一的行基矩阵
	\[
		\begin{pmatrix}
			1&0&a&b\\
			0&1&c&d
		\end{pmatrix},
	\]
	其中 \(a,b,c,d\in K\) 可任意取值，相应射影坐标为 \([1:c:d:-a:-b:ad-bc]\)。
\end{theorem}
\begin{proof}
	先说明关系式从何而来。将 \(p_{ij}=u_i v_j-u_j v_i\) 代入左边，展开得
	\[
		\begin{aligned}
			&(u_1v_2-u_2v_1)(u_3v_4-u_4v_3)\\
			&\quad-(u_1v_3-u_3v_1)(u_2v_4-u_4v_2)\\
			&\quad+(u_1v_4-u_4v_1)(u_2v_3-u_3v_2)=0;
		\end{aligned}
	\]
	每个含两个 \(u\) 坐标和两个 \(v\) 坐标的单项式出现两次且符号相反。这证明必要性。

	若 \(p_{12}\ne0\)，基矩阵的前两列组成可逆矩阵 \(C\)。将整矩阵左乘 \(C^{-1}\)，得到所述行基矩阵，而且前两坐标为 \((1,0)\)、\((0,1)\) 的两行在该子空间中唯一。逐项计算子式即得 \((1,c,d,-a,-b,ad-bc)\)。反过来，对满足关系式且 \(p_{12}\ne0\) 的射影点，将全部坐标除以 \(p_{12}\)，记归一化后的坐标为 \(\widetilde p_{ij}\)。取
	\[
		a=-\widetilde p_{23},\quad b=-\widetilde p_{24},\quad
		c=\widetilde p_{13},\quad d=\widetilde p_{14}.
	\]
	对应行基矩阵的前五个子式就是指定坐标；第六个由
	\[
		ad-bc=\widetilde p_{13}\widetilde p_{24}
		-\widetilde p_{14}\widetilde p_{23}=\widetilde p_{34}
	\]
	保证一致。这给出了存在性，也说明相同坐标只能恢复同一个子空间。

	任意非零六元组都有某个 \(p_{ij}\ne0\)。交换坐标轴可使这对指标成为 \(1,2\)。这里交换时同时使用 \(p_{ji}=-p_{ij}\)；关系式的左边在交换任意相邻两个指标时只改变符号，故它的零条件保持不变。于是上述构造适用于每一张图，并在图的重叠处由恢复出的子空间的唯一性相容，证明全局结论。
\end{proof}

\begin{example}[由六个数恢复一条直线]\label{g1:exm:plucker-reconstruction}
	在 \(K=\mathbb R\) 或 \(\mathbb C\) 上，六元组 \([1:2:3:4:5:-2]\) 满足 \(-2-10+12=0\)。按\cref{g1:thm:plucker-gr24}，其二维子空间有行基
	\[
		u=(1,0,-4,-5),\qquad v=(0,1,2,3),
	\]
	故对应直线为 \([s:t:-4s+2t:-5s+3t]\)，其中 \([s:t]\in\mathbb P^1(K)\)。若把最后一个坐标改成 \(0\)，关系式左边变成 \(2\)，所得射影点便不代表任何直线。所以六个齐次坐标之间的二次关系是实质约束。
\end{example}

一般的 \(\operatorname{Gr}(k,n;K)\) 也有同样可操作的坐标。将一组基写成秩为 \(k\) 的 \(k\times n\) 矩阵 \(M\)，选择一个非零的 \(k\) 阶子式。若它来自前 \(k\) 列，左乘这 \(k\) 列组成的矩阵的逆，即得唯一形式
\[
	M=[I_k\mid Z],\qquad Z\in\operatorname{Mat}_{k\times(n-k)}(K).
\]
唯一性的理由是：若 \(S[I_k\mid Z]=[I_k\mid Z']\)，比较前 \(k\) 列便有 \(S=I_k\)，继而 \(Z=Z'\)。任意 \(Z\) 都给出秩为 \(k\) 的矩阵。因此，这一部分中的子空间可用 \(k(n-k)\) 个自由参数描述。其他非零子式给出其他图；在重叠处，变换是先选出新的可逆子矩阵 \(C\)，再作 \(C^{-1}M\)。逆矩阵的余子式公式说明，新参数是旧参数的有理函数，分母就是被假设非零的 \(\det C\)。

\ProseParagraphBreak

所有 \(k\) 阶子式在换基时仍同乘 \(\det S\)，故又给出一个射影点，落在 \(\mathbb P^{\binom nk-1}(K)\) 中。这些坐标确实能恢复子空间：在选定子式归一化为 \(1\) 后， \([I_k\mid Z]\) 中的每个 \(Z_{ij}\) 都是用相应外部列替换第 \(i\) 个单位列所得子式，至多差一个由列排序产生的符号。因此所得到的射影映射为单射。一般维数时子式之间还有一组二次关系；式\eqref{g1:eq:plucker-gr24}给出了其中最小的非平凡模型。这里的参数数目由坐标图中的自由参数直接给出。

\subsection{点与超平面的关联簇及其投影}
\label{g1:subsec:0206:03}

对偶空间把一个超平面视为一个点；若同时保留原空间中的点，就能把“点属于超平面”本身作为一个几何对象。设 \(V=K^{n+1}\)、\(n\ge1\)。对偶空间 \(V^*\) 是 \(V\) 上的线性函数全体，非零函数 \(\alpha\) 的核给出超平面 \(\mathbb P(\ker\alpha)\)。定义
\[
	\mathcal I_n=
	\{\,([x],[\alpha])\in\mathbb P(V)\times\mathbb P(V^*):\alpha(x)=0\,\},
\]
称为点与超平面的\term[guanliancu]{关联簇}[incidence variety]。这里先将“簇”理解为这组方程给出的几何集合，关于代数簇的结构与一般性质将在第五卷中系统建立。在对偶基下，若 \(\alpha\) 的坐标为 \((a_0,\ldots,a_n)\)，其方程就是
\[
	\sum_{i=0}^n a_i x_i=0.
\]
分别缩放 \(a\) 和 \(x\) 不改变零条件，所以定义与代表选择无关。

\begin{proposition}[关联集合的两个投影]\label{g1:prop:incidence-fibers}
	设 \(K=\mathbb R\) 或 \(\mathbb C\)，\(V=K^{n+1}\)、\(n\ge1\)，并记 \(\mathcal I_n=\{([x],[\alpha])\in\mathbb P(V)\times\mathbb P(V^*):\alpha(x)=0\}\)。从 \(\mathcal I_n\) 到 \(\mathbb P(V)\)、\(\mathbb P(V^*)\) 的坐标投影都为满射。固定 \([x]\) 后，其上方的集合是 \(\mathbb P(\{\alpha:\alpha(x)=0\})\)，固定 \([\alpha]\) 后，其上方的集合是 \(\mathbb P(\ker\alpha)\)；二者都同构于 \(\mathbb P^{n-1}(K)\)。经 Segre 映射，\(\mathcal I_n\) 恰成为 Segre 像与超平面 \(\sum_i z_{ii}=0\) 的交集。
\end{proposition}
\begin{proof}
	固定非零向量 \(x\)，线性映射 \(V^*\rightarrow K\)、\(\alpha\mapsto\alpha(x)\) 是满射：将 \(x\) 扩为一组基，令线性函数在 \(x\) 上取值 \(1\)、在其余基向量上取值 \(0\) 即可。其核因而为 \(n\) 维，非零核向量的射影类就是所有经过 \([x]\) 的超平面。固定非零 \(\alpha\) 时，\(\alpha:V\rightarrow K\) 也满射，所以 \(\ker\alpha\) 同样为 \(n\) 维。由于 \(n\ge1\)，这些核都含非零向量，故两个投影都满射。最后，Segre 坐标取 \(z_{ij}=x_i a_j\) 时，关联方程恰为 \(\sum_i z_{ii}=0\)。\cref{g1:prop:segre-rank-one}给出的单射性和逆映射保证此描述没有额外或遗漏的点。
\end{proof}

分别记这两个投影为 \(\pi_1([x],[\alpha])=[x]\)、\(\pi_2([x],[\alpha])=[\alpha]\)。

\begin{figure}[htbp]
	\centering
	\input{Figures/Chapter02-Incidence.tex}
	\caption[点与超平面的关联]{点与超平面的关联。左边一个点的上方记录所有经过它的超平面；右边一个超平面的上方记录其全部点。}
	\label{g1:fig:incidence-projections}
\end{figure}

\cref{g1:fig:incidence-projections}把同一个条件的两种读法并排呈现。当 \(n=2\) 时，固定点得到过该点的直线束，固定直线得到这条直线上的点列；当 \(n=1\) 时，两边的纤维都是一个点，因为射影直线上的超平面本身就是点。这里映射的\term[xianwei]{纤维}[fiber]指一个指定点的逆像。

\ProseParagraphBreak

这些方程也解释了关联为什么是射影不变量。若 \(g\in\operatorname{GL}(V)\)，同时变换
\[
	[x]\longmapsto[gx],\qquad [\alpha]\longmapsto[\alpha\circ g^{-1}],
\]
则 \((\alpha\circ g^{-1})(gx)=\alpha(x)\)。因此点属于超平面的条件在改变射影坐标后保持不变。这些模型提出了三个共同问题：哪些条件描述的是一个子空间而非某组基，哪些参数变化只改变描述而不改变对象，哪些局部坐标能覆盖整个对象族。本节的秩一分解、子式恢复和关联方程已经分别给出了可直接计算的答案。后续研究更一般的曲线、曲面和参数空间时，这三个问题仍然存在。


\begin{chaptersummary}
	射影空间的点是一维向量子空间，仿射图是把一个非零坐标归一为 \(1\) 后得到的坐标区域。子空间的关联因此还原为线性代数，标架则把代表向量的相对比例一并固定。射影线上三个点可以完全归一化，第四个点留下交比；排列四点会改变这个值，故不变量必须与有序或无序对象的选择相对应。

	\ProseParagraphBreak

	二次多项式先分离能由平移消去的线性项与核中的障碍，再依允许的变换选择 Euclid 或仿射标准形。齐次化将二次项、线性项和常数项放进一个矩阵，但同时必须保留仿射图，才能恢复原来的仿射对象。实射影分类需要惯性，复射影分类只需秩；实零点集相同并不保证二次方程等价。曲面中的线性子空间由全迷向子空间控制，两个直纹族则把这套线性代数变成可见的几何。

	\ProseParagraphBreak

	对偶空间本来就把超平面参数化，非退化配对进一步为每个点指定一个极超平面。在二次曲面上，这个超平面恰是切超平面；配对退化时，对应可能失去定义或唯一性。圆锥曲线的有理参数把六点构造化成可检查的坐标等式，对偶再给出六条切线的共点结论。乘积、二次单项式、子空间坐标与关联方程展示了同一种构造方法：先记录几何对象所决定的坐标，证明选择无关，再说明哪些方程刻画像，并在局部坐标中恢复原对象。
\end{chaptersummary}

\BookExercises

\begin{exercise}\label{g1:ex:projective-completion}
	在 \(\mathbb R^2\) 中，将直线 \(2x-y=1\) 与 \(2x-y=3\) 补全到 \(\mathbb P^2(\mathbb R)\)。求其交点，并在以 \(X_1\ne0\) 为条件的图中写出交点坐标。解释为什么它在这张图中不是无穷远点。
\end{exercise}
\begin{solution}
	采用 \(x=X_1/X_0,y=X_2/X_0\)，补全方程为 \(2X_1-X_2=X_0\)、\(2X_1-X_2=3X_0\)。相减得 \(X_0=0\)，继而 \(X_2=2X_1\)，故唯一交点为 \([0:1:2]\)。在图 \(X_1\ne0\) 中取坐标 \((X_0/X_1,X_2/X_1)\)，交点坐标为 \((0,2)\)。原图的无穷远条件为 \(X_0=0\)，新图的无穷远条件则是 \(X_1=0\)。同一点可在一张图之外而在另一张图之内，所以“无穷远”描述的是相对于图的位置。
\end{solution}

\begin{exercise}\label{g1:ex:projective-frame}
	求把 \([e_0],[e_1],[e_2],[1:1:1]\) 依次送到 \([1:1:0],[0:1:1],[1:0:1],[4:3:5]\) 的射影变换，并求它把 \([1:2:-1]\) 送到哪里。
\end{exercise}
\begin{solution}
	取前三个像的代表 \(v_0=(1,1,0),v_1=(0,1,1),v_2=(1,0,1)\)。其列矩阵行列式为 \(2\)，故组成基。由 \(a_0v_0+a_1v_1+a_2v_2=(4,3,5)\)，解得 \((a_0,a_1,a_2)=(1,2,3)\)，均非零，故所给像也构成标架。依\cref{g1:thm:projective-frame}的构造，变换的矩阵可取
	\[
		M=\begin{pmatrix}1&0&3\\1&2&0\\0&2&3\end{pmatrix}.
	\]
	它将最后一个代表送到 \((4,3,5)\)，且 \(\det M=12\ne0\)。计算 \(M(1,2,-1)^{\mathsf T}=(-2,5,1)^{\mathsf T}\)，故所求像为 \([-2:5:1]\)。若只把 \(v_i\) 原样作为列，前三个像虽然正确，第四个像却不正确；列的相对比例正是第四点所确定的数据。
\end{solution}

\begin{exercise}\label{g1:ex:cross-ratio-test}
	在实射影线上，判断有序四点 \((\infty,0,1,3)\) 与 \((2,0,1,3/2)\) 是否射影等价。若等价，给出唯一变换。再判断把第二组的最后一点改成 \(4\) 后是否仍然等价。
\end{exercise}
\begin{solution}
	把前三点 \(\infty,0,1\) 送到 \(2,0,1\) 的分式为 \(f(t)=2t/(t+1)\)，齐次矩阵为 \(\begin{pmatrix}2&0\\1&1\end{pmatrix}\)，行列式非零。它满足 \(f(3)=3/2\)，所以第一组与原第二组等价；由前三点构成标架知变换唯一。改成 \(4\) 后，同一个唯一候选变换不能匹配第四点，故不等价。也可直接检查：\((2,0;1,3/2)=3\)，而 \((2,0;1,4)=-2\)，与原交比 \(3\) 不同。
\end{solution}

\begin{exercise}\label{g1:ex:affine-projective-group}
	设 \(K=\mathbb R\) 或 \(\mathbb C\)，\(n\geq1\)。证明 \(\mathbb P^n(K)\) 的射影变换保持超平面 \(X_0=0\) 当且仅当它在图 \(X_0\ne0\) 上的限制是仿射同构。说明每个仿射同构都有唯一这样的射影扩充。
\end{exercise}
\begin{solution}
	将线性代表分块写为 \(M=\begin{pmatrix}a&\ell\\b&S\end{pmatrix}\)。若保持 \(X_0=0\)，则对所有 \(x\in K^n\) 都有 \(\ell x=0\)，故 \(\ell=0\)。由于 \(M\) 可逆，\(a\ne0\)、\(S\) 可逆；除以 \(a\) 后，矩阵形如 \(\begin{pmatrix}1&0\\b'&S'\end{pmatrix}\)，在仿射图中诱导 \(x\mapsto b'+S'x\)。

	反过来，任给仿射同构 \(x\mapsto b'+S'x\)，所写块矩阵可逆，并保持 \(X_0=0\)，给出扩充。若已有射影变换在整个图上的限制为该仿射同构，那么图被双射到自身，其补集也被保持，所以其代表可归一化为同一种块矩阵。仿射函数的常数项与线性部分唯一，故扩充唯一。这里先说明限制在整个图上存在；仅在图中某个小区域看到分式表达式，并不能直接套用结论。
\end{solution}

\begin{exercise}
	\label{g1:ex:quadric-center-axes}
	在通常 Euclid 空间 \(\mathbb R^3\) 中考虑
	\[
		2x^2+2xy+2y^2+4x-2y+z^2+4=0.
	\]
	求中心、正交标准形和各半轴长，并求其齐次化矩阵的秩与惯性。
\end{exercise}
\begin{solution}
	二次项矩阵和线性系数列为
	\[
		A=\begin{pmatrix}2&1&0\\1&2&0\\0&0&1\end{pmatrix},
		\qquad b=(2,-1,0)^{\mathsf T}.
	\]
	由\cref{g1:prop:quadric-center}，中心满足 \(2x+y=-2\)、\(x+2y=1\)、\(z=0\)，所以唯一中心为 \(a=(-5/3,4/3,0)\)。在中心处，利用 \(Aa=-b\)，有
	\[
		F(a)=a^{\mathsf T}Aa+2b^{\mathsf T}a+4
		=b^{\mathsf T}a+4=-\frac23.
	\]
	矩阵 \(A\) 在单位方向 \((1,1,0)/\sqrt2\)、\((1,-1,0)/\sqrt2\)、\((0,0,1)\) 上的特征值分别为 \(3,1,1\)。因此令
	\[
		U=\frac{x+y+1/3}{\sqrt2},\qquad
		V=\frac{x-y+3}{\sqrt2},\qquad W=z,
	\]
	得到 \(3U^2+V^2+W^2=2/3\)。这是椭球面，半轴长分别为 \(\sqrt2/3\)、\(\sqrt{2/3}\)、\(\sqrt{2/3}\)，后两条轴可以在它们张成的平面内作任意正交选择。

	齐次化矩阵合同于 \(\operatorname{diag}(-2/3,A)\)，因为平移到中心已经消去了全部线性项。故其秩为 \(4\)，惯性为 \((3,1,0)\)。这项结论同时说明该射影曲面没有实射影直线。
\end{solution}

\begin{exercise}
	\label{g1:ex:paraboloid-parameter}
	对 \(t\in\mathbb R\)，在 \(\mathbb R^3\) 中令
	\[
		F_t(x,y,z)=x^2+z^2+2ty-1.
	\]
	\begin{enumerate}
		\item 求 \(F_t\) 的中心集、其零点集的形状，以及齐次化的秩与惯性。
		\item 对 \(s,t\ne0\)，证明方程 \(F_s=0\) 与 \(F_t=0\) 总是仿射等价，而 Euclid 等价当且仅当 \(|s|=|t|\)。
	\end{enumerate}
\end{exercise}
\begin{solution}
	1. 二次项矩阵为 \(\operatorname{diag}(1,0,1)\)，核为 \(y\) 轴，线性项的系数列是 \((0,t,0)\)。当 \(t=0\) 时，中心是整条 \(y\) 轴，零点集为圆柱面 \(x^2+z^2=1\)，中心值为 \(-1\)。由\cref{g1:prop:homogenized-rank}，齐次化的秩为 \(3\)，惯性为 \((2,1,1)\)。当 \(t\ne0\) 时没有中心。令 \(Y=y-1/(2t)\)，方程成为 \(x^2+z^2+2tY=0\)，所以是以 \(y\) 方向为轴的抛物面；其齐次化的秩为 \(4\)，惯性为 \((3,1,0)\)。

	2. 对 \(s,t\ne0\)，变换 \((x,y,z)\mapsto(x,(t/s)y,z)\) 使 \(F_s\) 变为 \(F_t\)，故方程仿射等价。Euclid 标准形中的非零特征值为 \(\{1,1\}\)，剩余线性项长度为 \(|t|\)。若两个方程 Euclid 等价，还可以有整体倍数 \(\mu\ne0\)；二次项的特征值必须满足 \(\{1,1\}=\{\mu,\mu\}\)，故 \(\mu=1\)，再由\cref{g1:thm:euclidean-quadric-normal}得到 \(|s|=|t|\)。反过来，\(s=t\) 时取恒等变换，\(s=-t\) 时取 \(y\mapsto-y\) 的反射，均给出 Euclid 等价。
\end{solution}

\begin{exercise}
	\label{g1:ex:quadric-locus-information}
	在 \(\mathbb R^3\) 中令 \(F=x^2+y^2\)、\(G=x^2+4y^2\)。证明它们的零点集相同，但作为方程不 Euclid 等价；再给出一个保持这个共同零点集、却不把 \(G\) 变成自身非零倍数的 Euclid 变换。
\end{exercise}
\begin{solution}
	因为每一项都非负，两个方程都等价于 \(x=y=0\)，故实零点集都是 \(z\) 轴。两者有中心，中心值都为零；二次项的非零特征值多重集分别为 \(\{1,1\}\) 和 \(\{1,4\}\)。方程 Euclid 等价要求存在非零 \(\mu\)，使 \(\{\mu,\mu\}=\{1,4\}\)，这是不可能的。因此相同实点集不能恢复 Euclid 方程类别。

	取绕 \(z\) 轴的正交变换
	\[
		(x,y,z)\longmapsto
		\left(\frac{x-y}{\sqrt2},\frac{x+y}{\sqrt2},z\right).
	\]
	它固定整条 \(z\) 轴，而变换后的 \(G\) 为
	\[
		\frac{(x-y)^2}{2}+2(x+y)^2
		=\frac52x^2+3xy+\frac52y^2.
	\]
	新出现的交叉项表明它不是 \(G\) 的非零倍数。这说明点集的对称群可以严格大于指定二次方程的对称群。
\end{solution}

\begin{exercise}
	\label{g1:ex:isotropic-graph-orthogonal}
	在 \(\mathbb R^5=\mathbb R^2\oplus\mathbb R^2\oplus\mathbb R\) 上考虑
	\[
		q(u,v,z)=\lVert u\rVert^2-\lVert v\rVert^2.
	\]
	证明它的每个极大全迷向子空间都唯一写成
	\[
		U_T=\{\,(u,Tu,z):u\in\mathbb R^2,\ z\in\mathbb R\,\},
		\qquad T\in O(2).
	\]
	据此说明对应射影二次曲面的极大射影平面都经过哪一个点。
\end{exercise}
\begin{solution}
	该型的惯性为 \((2,2,1)\)，根基为 \(R=\{(0,0,z)\}\)。由\cref{g1:thm:quadric-linear-dimension}，极大全迷向子空间 \(U\) 维数为 \(3\)，并包含 \(R\)。将它投影到前四个坐标，得到二维全迷向子空间 \(W\)，且 \(U=W\oplus R\)。

	从 \(W\) 到第一个 \(\mathbb R^2\) 的投影单射：若 \(u=0\)，则 \(q=0\) 迫使 \(v=0\)。由于两边都是二维，这个投影也是满射，故恰有一个线性映射 \(T\) 使 \(W=\{(u,Tu)\}\)。全迷向条件给出 \(\lVert Tu\rVert^2=\lVert u\rVert^2\) 对所有 \(u\) 成立。用实内积的极化公式可得 \(\langle Tu,Tw\rangle=\langle u,w\rangle\)，即 \(T^{\mathsf T}T=I\)。因此 \(T\in O(2)\)，且其唯一性来自图像对每个 \(u\) 的唯一取值。

	反过来，任意正交 \(T\) 都使 \(q\) 在 \(U_T\) 上恒为零；\(U_T\) 维数为 \(3\)，已经达到最大值，故极大。射影化后得到全部极大射影平面，它们都包含根基对应的点 \([0:0:0:0:1]\)。
\end{solution}

\begin{exercise}
	\label{g1:ex:quadric-cone-family}
	在 \(\mathbb P^3(\mathbb R)\) 中考虑参数族
	\[
		Q_t:\ X_0X_3-X_1^2-tX_2^2=0.
	\]
	\begin{enumerate}
		\item 按 \(t\) 的符号求秩、惯性和最大射影线性子空间维数。
		\item 当 \(t=0\) 时，证明曲面中的全部射影直线恰为
		\[
			L_{[a:b]}=\mathbb P\!\left(
			\operatorname{span}\bigl((a^2,ab,0,b^2),(0,0,1,0)\bigr)\right),
			\quad [a:b]\in\mathbb P^1(\mathbb R).
		\]
	\end{enumerate}
\end{exercise}
\begin{solution}
	1. 乘积 \(X_0X_3\) 贡献一个正方向、一个负方向；\(-X_1^2\) 贡献一个负方向。因此 \(t>0\) 时惯性为 \((1,3,0)\)，秩为 \(4\)，最大全迷向维数为 \(1\)，曲面内最大射影线性子空间只是点，没有实直线。\(t<0\) 时惯性为 \((2,2,0)\)，秩为 \(4\)，最大射影线性子空间维数为 \(1\)；由射影分类及\cref{g1:prop:quadric-two-rulings}，有两个直纹族。\(t=0\) 时惯性为 \((1,2,1)\)，秩为 \(3\)，最大射影线性维数仍为 \(1\)。

	2. 在 \(t=0\) 时，根基由 \(e_2=(0,0,1,0)\) 张成。任意曲面内直线对应一个二维全迷向子空间，已经达到最大维数，所以包含根基。模去根基后，它由 \((X_0,X_1,X_3)\ne0\) 且 \(X_0X_3=X_1^2\) 的一个射影点决定。若 \(X_0\ne0\)，把该点写为 \([1:u:u^2]\)，对应 \([a:b]=[1:u]\)；若 \(X_0=0\)，则 \(X_1=0\)、\(X_3\ne0\)，对应 \([a:b]=[0:1]\)。这给出全部且唯一的参数点。

	对每个 \([a:b]\)，向量 \((a^2,ab,0,b^2)\) 非零、迷向，且不在根基中；它与根基向量的所有配对为零。因此所列二维空间全迷向，其射影化确为曲面内直线。全部直线都经过顶点 \([0:0:1:0]\)，与满秩分裂曲面的两个直纹族有不同的相交关系。
\end{solution}

\begin{exercise}
	\label{g1:ex:complex-line-real-sphere}
	考虑复射影二次曲面
	\[
		Q:\ X_1^2+X_2^2+X_3^2-X_0^2=0
		\quad\text{于 }\mathbb P^3(\mathbb C).
	\]
	证明 \(u=(0,1,i,0)\)、\(v=(1,0,0,1)\) 张成的二维复子空间给出 \(Q\) 中的一条复射影直线，求它上面的全部实射影点，并解释这为何不与实球面没有实直线矛盾。
\end{exercise}
\begin{solution}
	分别计算得到 \(q(u)=1+i^2=0\)、\(q(v)=1-1=0\)，且对应双线性型满足 \(B(u,v)=0\)。所以对任意复数 \(a,b\)，
	\[
		q(au+bv)=a^2q(u)+2abB(u,v)+b^2q(v)=0.
	\]
	\(u,v\) 线性无关，故其张成空间二维，射影化是一条复直线。该直线的点写成 \([b:a:ia:b]\)。若 \(b\ne0\)，归一化为 \([1:w:iw:1]\)，它有实代表当且仅当 \(w\) 与 \(iw\) 同时为实数，故 \(w=0\)。若 \(b=0\)，唯一射影点为 \([0:1:i:0]\)，没有实代表。因此直线只有一个实点 \([1:0:0:1]\)。

	实射影直线须由二维实向量子空间射影化而来，包含无穷多个实点；本例的二维空间是复空间，不能据此推出存在二维实全迷向子空间。实型的惯性为 \((3,1,0)\)，仍由\cref{g1:thm:quadric-linear-dimension}排除实射影直线。
\end{solution}

\begin{exercise}\label{g1:ex:polarity-coordinate-change}
	设 \(K=\mathbb R\) 或 \(\mathbb C\)，\(A\) 为可逆对称矩阵。在旧坐标中，二次曲面为 \(x^{\mathsf T}Ax=0\)。改用 \(x=Sy\)，其中 \(S\) 可逆。求新坐标中的二次型矩阵、极超平面系数和对偶二次型矩阵，并解释为什么对点和超平面的坐标不能使用同一个变换公式。
\end{exercise}
\begin{solution}
	直接代入得到新矩阵 \(A'=S^{\mathsf T}AS\)。旧超平面 \(a^{\mathsf T}x=0\) 变成 \((S^{\mathsf T}a)^{\mathsf T}y=0\)，故新系数为 \(a'=S^{\mathsf T}a\)。若旧点为 \(x\)，新点为 \(y=S^{-1}x\)，其极超平面的新系数为
	\[
		A'y=S^{\mathsf T}ASS^{-1}x=S^{\mathsf T}Ax,
	\]
	恰与先在旧坐标取极超平面再变换一致。对偶方程的新矩阵是
	\[
		(A')^{-1}=S^{-1}A^{-1}S^{-\mathsf T}.
	\]
	这是因为旧系数 \(a=S^{-\mathsf T}a'\)，将其代入旧对偶方程便得到同一矩阵。点的坐标用 \(S^{-1}\)，超平面系数用 \(S^{\mathsf T}\)；它们必须保持求值 \(a^{\mathsf T}x=(a')^{\mathsf T}y\)，不能因都是坐标列而使用相同规则。
\end{solution}

\begin{exercise}\label{g1:ex:ellipse-tangents}
	在实仿射平面中考虑椭圆 \(x^2/9+y^2/4=1\)。用对偶二次曲面的方法，求从点 \((6,0)\) 引出的全部实切线及其切点。再解释为什么从原点引不出实切线。
\end{exercise}
\begin{solution}
	齐次方程为 \(4X^2+9Y^2-36Z^2=0\)。由\cref{g1:thm:dual-quadric}，直线系数 \([a:b:c]\) 对应切线的充要条件是 \(9a^2+4b^2-c^2=0\)。过 \([6:0:1]\) 的条件为 \(c=-6a\)，于是 \(4b^2=27a^2\)。这里 \(a\ne0\)，否则三个系数全零。按比例取 \(a=2\)，得到两条直线
	\[
		2x+3\sqrt3\,y=12,\qquad 2x-3\sqrt3\,y=12.
	\]
	矩阵逆像给出切点的齐次代表
	\[
		\operatorname{diag}(1/4,1/9,-1/36)
		\begin{pmatrix}2\\ \pm3\sqrt3\\ -12\end{pmatrix}
		=\begin{pmatrix}1/2\\ \pm\sqrt3/3\\ 1/3\end{pmatrix},
	\]
	故切点为 \((3/2,\pm\sqrt3)\)，分别代入椭圆和相应直线均成立。过原点则 \(c=0\)，对偶方程成为 \(9a^2+4b^2=0\)。在实数域上只能有 \(a=b=0\)，不是合法直线系数，故没有实切线。
\end{solution}

\begin{exercise}\label{g1:ex:cone-dual-fibers}
	在 \(\mathbb P^3(K)\)、\(K=\mathbb R\) 或 \(\mathbb C\) 中，令 \(Q\) 由 \(XZ-Y^2=0\) 定义。求全部光滑点处切平面组成的对偶点集，并求每一个切平面对应的全部切点。
\end{exercise}
\begin{solution}
	配对的根基为 \(K(0,0,0,1)\)，因此唯一非光滑点是顶点 \(O=[0:0:0:1]\)。在其余点 \([X:Y:Z:W]\)，切平面系数可取 \([Z:-2Y:X:0]\)。故若对偶坐标写成 \([a:b:c:d]\)，像满足
	\[
		d=0,\qquad b^2-4ac=0.
	\]
	反过来，任一满足这些方程的射影点都有 \((a,b,c)\ne(0,0,0)\)，点 \([2c:-b:2a:w]\) 对每个 \(w\in K\) 都位于 \(Q\setminus\{O\}\)，且切平面正是 \([a:b:c:0]\)。若一个点给出同一切平面，存在 \(\lambda\ne0\) 使 \((Z,-2Y,X)=\lambda(a,b,c)\)，缩放其代表向量后便是上述形式。因此每个切平面的全部切点是一条经过顶点的射影直线去掉顶点。对偶像是超平面 \(d=0\) 内的一条非退化圆锥曲线；不同母线被区分，而同一母线上各光滑点被合并。
\end{solution}

\begin{exercise}\label{g1:ex:five-conic-points}
	在 \(K=\mathbb R\) 或 \(\mathbb C\) 上，取标准圆锥曲线 \(XZ-Y^2=0\) 的参数 \(\infty,0,1,-1,2\) 所对应的五个点。证明：在这五点处都为零的齐次二次多项式必为 \(XZ-Y^2\) 的常数倍。再说明同样结论对该曲线上任意五个不同点成立。
\end{exercise}
\begin{solution}
	一般齐次二次式写成 \(F=aX^2+bXY+cXZ+dY^2+eYZ+fZ^2\)。代入参数得
	\[
		F(s^2,st,t^2)=as^4+bs^3t+(c+d)s^2t^2+est^3+ft^4.
	\]
	参数 \(\infty\) 使 \(a=0\)。在有限图中得到次数至多为三的多项式 \(br^3+(c+d)r^2+er+f\)，它在 \(0,1,-1,2\) 四个不同数处为零，故各系数全零。于是 \(b=e=f=0\)、\(d=-c\)，即 \(F=c(XZ-Y^2)\)。这里用了非零一元多项式的根数不超过次数；可由每个根都提出一个一次因子反复得到。

	任意五个不同参数对应一个二元齐次四次式的五个不同射影零点。若无穷远不是其中之一，在有限图中是次数至多四的多项式有五个不同根，故恒为零；若包含无穷远，则最高次系数为零，有限图中的次数至多三，而另有四个不同根，仍恒为零。比较五个系数同样得到 \(F=c(XZ-Y^2)\)。对任意有点非退化圆锥曲线，再通过\cref{g1:prop:conic-parametrization}的平面换坐标，就得到相应的唯一性。
\end{solution}

\begin{exercise}\label{g1:ex:pascal-converse}
	在 \(\mathbb P^2(K)\)、\(K=\mathbb R\) 或 \(\mathbb C\) 中，令 \(A,B,C,D,E\) 为 \(XZ-Y^2=0\) 上参数依次为 \(\infty,0,1,2,3\) 的点。任取与这五点都不同的点 \(F\)。证明 \(AB\cap DE\)、\(BC\cap EF\)、\(CD\cap FA\) 共线，当且仅当 \(F\) 也在该圆锥曲线上。
\end{exercise}
\begin{solution}
	写 \(F=[x:y:z]\)。前五点的坐标为 \([1:0:0],[0:0:1],[1:1:1],[4:2:1],[9:3:1]\)。前一对边的交点为 \(U=[-6:0:1]\)。直线 \(EF\) 的系数为 \([3z-y:x-9z:9y-3x]\)，与 \(BC:X-Y=0\) 相交得到
	\[
		V=[3x-9y:3x-9y:x-y-6z].
	\]
	其坐标同时为零只可能有 \(x=3y\)、\(y=3z\)，即 \(F=E\)，已排除。又由 \(FA:zY-yZ=0\)、\(CD:X-3Y+2Z=0\) 得到
	\[
		W=[3y-2z:y:z].
	\]
	该坐标同时为零只可能 \(y=z=0\)，即 \(F=A\)，也已排除。因为 \(U\notin BC\) 而 \(V\in BC\)，二者不同。三个非零向量所代表的射影点共线，恰当且仅当它们的三阶行列式为零。展开得到
	\[
		\det\begin{pmatrix}
			-6&3x-9y&3y-2z\\
			0&3x-9y&y\\
			1&x-y-6z&z
		\end{pmatrix}
		=-12(xz-y^2).
	\]
	在实数域和复数域中 \(-12\ne0\)，故共线恰等价于 \(xz-y^2=0\)，这就是所求条件。
\end{solution}

\begin{exercise}\label{g1:ex:brianchon-computation}
	在实射影平面中，按参数 \(\infty,0,1,2,3,-2\) 的次序，取圆锥曲线 \(XZ-Y^2=0\) 上六点的切线。求相邻切线交点所组成六边形的三条对角线及其公共点，允许顶点位于无穷远处。
\end{exercise}
\begin{solution}
	按\cref{g1:exm:dual-conic}，六条切线依次为
	\[
		\begin{aligned}
			&Z=0,&&X=0,\\
			&X-2Y+Z=0,&&X-4Y+4Z=0,\\
			&X-6Y+9Z=0,&&X+4Y+4Z=0.
		\end{aligned}
	\]
	逐次求相邻交点，按\cref{g1:thm:brianchon}的记号得到
	\[
		\begin{aligned}
			A&=[0:1:0],&\qquad B&=[0:1:2],\\
			C&=[4:3:2],& D&=[12:5:2],\\
			E&=[-12:1:2],& F&=[4:-1:0].
		\end{aligned}
	\]
	于是三条对角线分别为
	\[
		AD:X-6Z=0,\qquad BE:2Y-Z=0,\qquad CF:X+4Y-8Z=0.
	\]
	前两条交于 \([12:1:2]\)，代入第三条亦为零。因此公共点在仿射图 \(Z=1\) 中为 \((6,1/2)\)。两个顶点 \(A,F\) 位于无穷远处，而三条对角线与公共点都仍有明确的射影含义。
\end{solution}

\begin{exercise}\label{g1:ex:segre-diagonal}
	设 \(K=\mathbb R\) 或 \(\mathbb C\)。在 Segre 映射 \(\sigma_{1,1}([s:t],[u:v])=[su:sv:tu:tv]\) 下，求乘积中对角线 \(\Delta=\{([s:t],[s:t])\}\) 的像。证明它位于一个射影平面内，并在该平面内为非退化圆锥曲线。说明这条曲线与 Segre 曲面两族直纹的相交情况。
\end{exercise}
\begin{solution}
	对角线的像为 \([s^2:st:st:t^2]\)，故位于平面 \(z_{01}=z_{10}\)。以 \([A:B:C]=[z_{00}:z_{01}:z_{11}]\) 作为该平面坐标，曲面方程化为 \(AC-B^2=0\)。其对称矩阵为
	\[
		\begin{pmatrix}0&0&1/2\\0&-1&0\\1/2&0&0\end{pmatrix},
	\]
	行列式为 \(1/4\ne0\)，故是非退化圆锥曲线。若 \(A\ne0\)，令 \([s:t]=[A:B]\) 即得到指定点，因为 \(AC=B^2\)；若 \(A=0\)，则 \(B=0\)、\(C\ne0\)，由 \([s:t]=[0:1]\) 得到。所以像恰为整条圆锥曲线。固定第一因子 \([s:t]=p\) 的直纹与对角线恰交于 \((p,p)\)，固定第二因子的直纹也一样；Segre 映射的单射性保证相交点唯一。
\end{solution}

\begin{exercise}\label{g1:ex:veronese-hyperplane}
	对 \(K=\mathbb R\) 或 \(\mathbb C\)，将 \(\mathbb P^2(K)\) 的二次 Veronese 映射写成
	\[
		[x:y:z]\longmapsto[A:B:C:D:E:F]
		=[x^2:xy:xz:y^2:yz:z^2].
	\]
	求超平面 \(A+D-F=0\) 在原平面中的逆像。证明直线 \(z=0\) 的像是一条平面圆锥曲线，并比较它与该超平面在实数和复数情形的交点。
\end{exercise}
\begin{solution}
	超平面的逆像为 \(x^2+y^2-z^2=0\)。直线 \(z=0\) 的像是 \([x^2:xy:0:y^2:0:0]\)，位于平面 \(C=E=F=0\) 内，并满足 \(AD-B^2=0\)。反向恢复 \([x:y]\) 与二次 Veronese 映射的逆公式相同，所以像恰为该圆锥曲线。与指定超平面相交还要求 \(A+D=0\)，等价于 \(x^2+y^2=0\)。在 \(\mathbb R\) 上只有零向量解，因而没有射影交点；在 \(\mathbb C\) 上有两个不同解 \([x:y]=[1:i]\)、\([1:-i]\)，像分别是 \([1:i:0:-1:0:0]\)、\([1:-i:0:-1:0:0]\)。改变基域保留了方程，却改变了可见的交点集。
\end{solution}

\begin{exercise}\label{g1:ex:plucker-intersection}
	设 \(K=\mathbb R\) 或 \(\mathbb C\)，射影三维空间中两条直线的 Plücker 坐标分别为 \([p_{ij}]\)、\([q_{ij}]\)。证明它们相交（允许重合）当且仅当
	\[
		p_{12}q_{34}-p_{13}q_{24}+p_{14}q_{23}
		+p_{23}q_{14}-p_{24}q_{13}+p_{34}q_{12}=0.
	\]
	用此条件判断由 \((1,0,0,0),(0,1,0,0)\) 张成的直线与由 \((0,0,1,0),(0,0,0,1)\) 张成的直线是否相交。
\end{exercise}
\begin{solution}
	分别选取两条直线对应的二维向量子空间 \(U,W\) 的基 \(u,v\) 和 \(a,b\)。两条射影直线相交等价于 \(U\cap W\ne\{0\}\)。由向量子空间的维数公式，这又等价于 \(\dim(U+W)<4\)，即以 \(u,v,a,b\) 为四行的矩阵行列式为零。沿前两行展开此行列式：选定的两列为 \(1,2\)、\(1,3\)、\(1,4\)、\(2,3\)、\(2,4\)、\(3,4\) 时，余下两列的子式分别是 \(q_{34},q_{24},q_{23},q_{14},q_{13},q_{12}\)，相应符号依次为 \(+,-,+,+,-,+\)。因此展开式正是题中表达式。换基只使其乘以非零数，故零条件与选基无关。所给两条直线分别只有 \(p_{12}=1\)、\(q_{34}=1\) 非零，表达式等于 \(1\)，所以不相交。这也直接说明射影三维空间中的两条直线未必相交。
\end{solution}

\begin{exercise}\label{g1:ex:incidence-chart}
	设 \(K=\mathbb R\) 或 \(\mathbb C\)，在 \(\mathcal I_2\) 中考虑 \(x_0\ne0\)、\(a_2\ne0\) 的部分。分别将两组坐标归一化为 \(x_0=a_2=1\)，求一组无约束参数。解释两个投影如何在这些参数中表示，并说明这张图为什么没有覆盖整个 \(\mathcal I_2\)。
\end{exercise}
\begin{solution}
	写 \([x]=[1:s:t]\)、\([a]=[b:c:1]\)，关联方程为 \(b+cs+t=0\)，故 \(b=-cs-t\)。于是所求部分与 \(K^3\) 一一对应：
	\[
		(s,t,c)\longmapsto([1:s:t],[-cs-t:c:1]).
	\]
	反向由 \(s=x_1/x_0\)、\(t=x_2/x_0\)、\(c=a_1/a_2\) 唯一恢复参数。第一投影取 \((s,t,c)\mapsto[1:s:t]\)，固定点后 \(c\) 可任意取值；第二投影取 \((s,t,c)\mapsto[-cs-t:c:1]\)。若固定右边为 \([b:c:1]\)，则 \(s\) 可任意取值而 \(t=-b-cs\)。这两个自由的一参数集合分别只是相应完整射影直线的一张仿射图。例如 \(([0:1:0],[1:0:0])\) 属于 \(\mathcal I_2\)，但不满足两个非零条件，因而不在本图中。
\end{solution}

% !TeX root = ../../Book1.tex
% 纲要标题：## 第3章　凸性、多面体与有限维几何方法
% 纲要位置：几何学计划纲要.md，第 200 行。
% 正文状态：初稿；尚未正式审核。

% 纲要内容开始
% 
% > 本章目标与证明范围：以分离、有限支撑与对偶为主线，完整证明所选有限维分离定理、Carathéodory、Radon 和 Helly 定理。把群作用移至第1章，避免与凸几何并列拼接。
% 
% 纲要内容结束

\chapter{凸性、多面体与有限维几何方法}
\label{g1:ch:03}
\BookChapterNumber{g1:ch:03}

\begin{chapterabstract}
	凸性要求两点之间的整条线段留在集合内，却能把许多整体问题转成线性不等式。最近点投影产生支撑与分离超平面；仿射依赖限制凸组合所需的点数，并控制有限凸集族的相交；多面体则把这些方法收束为有限个顶点与有限个半空间之间的转换。面与法锥记录边界结构，极体把包含关系反转。最后，距离公式、环面作用的线性模型和格点的单项式映射展示这些方法如何进入后续几何。
\end{chapterabstract}

\begin{learninggoals}
	\begin{itemize}
		\item 区分内部与相对内部，证明闭凸集的最近点存在且唯一，并从投影不等式构造支撑及分离超平面。
		\item 说明弱分离、严格分离和有正间隔的分离各需要什么条件，识别闭性与紧性不能互换的情形。
		\item 利用仿射依赖完整证明 Carathéodory、Radon 和 Helly 定理，并用例子检验维数、凸性与有限族条件。
		\item 在有界情形证明半空间表示与顶点表示等价，计算面、法锥及极体，理解面对应为何反转包含关系。
		\item 用支撑函数计算距离，从整数权重计算矩映射的像，从格点关系写出射影方程。
	\end{itemize}
\end{learninggoals}

平面上的三角形
\[
	T=\{(x,y):x\ge0,\ y\ge0,\ x+y\le1\}
\]
既由三个不等式描述，又是三个点 \((0,0),(1,0),(0,1)\) 的凸包。后一种说法可以直接核验：若 \((x,y)\in T\)，则
\[
	(x,y)=(1-x-y)(0,0)+x(1,0)+y(0,1),
\]
三个系数非负且和为 \(1\)。不等式告诉我们哪些位置允许进入，顶点告诉我们如何生成所有允许的位置。更复杂的集合是否也有这两种描述？若一个点不在集合里，能否给出一条线性不等式作为证据？如果约束很多，又能否只看少量约束就判断它们不能同时满足？

\ProseParagraphBreak

本章逐步回答这些问题。先不要求集合由有限个条件给定，而从线段封闭这一性质出发。对非空闭凸集，集合外一点的最近点给出一个法向；这个法向产生的不等式把点与整个集合分开。接着利用有限维线性代数，把“点或条件太多”转化为依赖关系。最后回到有限个半空间的交，证明有界时可以恢复有限个顶点，并追踪边界各部分如何在对偶构造中交换角色。

\ProseParagraphBreak

所需的代数基础是仿射组合、秩、核与对偶空间。\cref{g1:prop:affine-combination}说明系数和为 \(1\) 的组合不依赖原点；本章再要求系数非负，从而保留线段内的位置。分析方面使用有限维欧氏空间中“闭且有界等价于紧”、紧集上连续实函数取得最值，以及有界序列有收敛子列；应用这些事实时会指出闭性与有界性从何而来。本章先在 \(\mathbb R^n\) 中证明所需结论，不要求泛函分析中的一般分离定理。凸集与分离的补充阅读可参看\cite[Sections 2.1, 2.5--2.6]{BoydVandenberghe2004}和\cite[Sections 3.1--3.2, 3.4, 7.1--7.2]{Gallier2011Geometric}；多面体的有限表示与面结构可参看\cite[Sections 1.1, 2.1--2.4]{Ziegler1995}。

% !TeX root = ../../Book1.tex
% 纲要标题：### 3.1 凸集、支撑与分离
% 纲要位置：几何学计划纲要.md，第 204 行。
% 正文状态：初稿；尚未正式审核。

% 纲要内容开始
% 
% * 凸包、相对内部和仿射包络
% * 最近点投影及支撑超平面
% * 凸集的弱分离、严格分离与强分离
% * 闭凸集的支撑函数表示
% 
% 纲要内容结束

\section{凸集、支撑与分离}
\label{g1:sec:0301}

在平面上围住一组点，可以用曲线描出边界，也可以逐个指定包含这些点的半平面。前一种描述关注边界在哪里，后一种描述关注各个方向上能走多远。对于凸集，两种描述能够互相恢复。证明这一点的关键是：若一个点在闭凸集之外，能否找到一个半空间，把整个集合留在一侧而把该点排除在外？本节从最近点出发构造这样的半空间。除特别说明外，环境为带标准内积的 \(\mathbb R^n\)，\(n\geq1\)，范数记为 \(\lVert\cdot\rVert\)。

\subsection{凸包、相对内部和仿射包络}
\label{g1:subsec:0301:01}

\cref{g1:prop:affine-combination}说明，系数和为 \(1\) 的仿射组合不依赖原点。若进一步要求系数非负，就不会沿着两点的连线越过端点。这一限制把仿射几何中的“张成”变成了凸几何中的“围住”。

\begin{definition}[凸集与凸包]\label{g1:def:convex-hull}
	设集合 \(C\subseteq\mathbb R^n\)。若对任意 \(x,y\in C\) 和 \(t\in[0,1]\)，都有 \((1-t)x+ty\in C\)，则称其为\term[tuji]{凸集}[convex set]。有限个点的仿射组合 \(\sum_{i=1}^m\lambda_i x_i\) 若满足 \(\lambda_i\geq0\)、\(\sum_i\lambda_i=1\)，称为\term[tuzuhe]{凸组合}[convex combination]。对任意 \(S\subseteq\mathbb R^n\)，其\term[tubao]{凸包}[convex hull] \(\operatorname{conv}S\) 是所有包含 \(S\) 的凸集之交。
	\symidx{\(\operatorname{conv}S\)：凸包}
\end{definition}

按定义，空集和单点集都是凸集；\(\operatorname{conv}\varnothing=\varnothing\)。任意一族凸集的交仍为凸集，因为交集中任意两点之间的线段属于这族集合的每一个成员。这里不要求集合闭：开球和闭球都凸，开球并不包含自己的边界。

\begin{proposition}[凸包的有限组合表示]\label{g1:prop:convex-finite-combinations}
	设 \(S\subseteq\mathbb R^n\) 非空，则
	\[
		\operatorname{conv}S=\left\{\,\sum_{i=1}^m\lambda_i x_i:m\geq1,\ x_i\in S,\ \lambda_i\geq0,\ \sum_{i=1}^m\lambda_i=1\,\right\}.
	\]
	此外，凸集的闭包仍为凸集，仿射映射把凸集映为凸集。
\end{proposition}
\begin{proof}
	先用点数归纳说明凸集包含其点的一切有限凸组合。一个点的情形成立。若 \(\lambda_m=1\)，组合等于 \(x_m\)；否则把前 \(m-1\) 项除以 \(1-\lambda_m>0\)，先得到归纳假设所允许的一个点，再与 \(x_m\) 作两点凸组合。因而上式右边包含于每个包含 \(S\) 的凸集中。

	反过来，右边包含 \(S\)，而两个有限凸组合再按权重 \(1-t,t\) 组合，只需把原来两组系数分别乘以 \(1-t,t\)，仍得到非负且总和为 \(1\) 的有限组系数。因此右边本身凸，由凸包的定义得到相等。

	若 \(x,y\in\overline C\)，取 \(x_j,y_j\in C\) 分别趋于 \(x,y\)。对固定 \(t\in[0,1]\)，点 \((1-t)x_j+ty_j\in C\) 趋于 \((1-t)x+ty\)，故极限在 \(\overline C\) 中。最后，\cref{g1:prop:affine-preservation}给出 \(f((1-t)x+ty)=(1-t)f(x)+tf(y)\)，因而 \(f(C)\) 凸。
\end{proof}

“有限”不是对集合 \(S\) 的限制，而是对每一个具体表示的限制；所需点数暂时可以依赖于所表示的点。另一个需要区分的概念是仿射包络。\secref{g1:sec:0101}末尾已说明，对非空 \(S\) 及 \(a\in S\)，
\[
	\operatorname{aff}S=a+\operatorname{span}\{x-a:x\in S\}.
\]
它允许系数为负的一切有限仿射组合，因而 \(\operatorname{conv}S\subseteq\operatorname{aff}S\)，且两者有相同的仿射包络。记 \(\dim C=\dim\operatorname{aff}C\)，此记号在这里用于非空凸集。

\begin{definition}[相对内部与相对边界]\label{g1:def:relative-interior}
	设非空集合 \(C\subseteq\mathbb R^n\) 的仿射包络为 \(F=a+L\)。它的\term[xiangduineibu]{相对内部}[relative interior]是它在 \(F\) 中的内部，记为
	\[
		\operatorname{ri}C=\{\,x\in C:\text{存在 }r>0\text{ 使 }\{x+v:v\in L,\ \lVert v\rVert<r\}\subseteq C\,\}.
	\]
	它的\term[xiangduibianjie]{相对边界}[relative boundary]为 \(\operatorname{rbd}C=\overline C\setminus\operatorname{ri}C\)。这里仿射子空间 \(F\) 在有限维 Euclid 空间中闭，故相对闭包与环境空间中的闭包相同。对空集约定 \(\operatorname{ri}\varnothing=\operatorname{rbd}\varnothing=\varnothing\)。
	\symidx{\(\operatorname{ri}C,\operatorname{rbd}C\)：相对内部与相对边界}
\end{definition}

\begin{example}[一条线段的两种内部]\label{g1:exm:relative-segment}
	令 \(C=\{(t,0):0\leq t\leq1\}\subseteq\mathbb R^2\)。平面中的任何小圆盘都不能放进 \(C\)，故 \(\operatorname{int}C=\varnothing\)，整个 \(C\) 都是环境空间中的边界。它的仿射包络却只是 \(x\) 轴，于是 \(\operatorname{ri}C=\{(t,0):0<t<1\}\)，相对边界只有两个端点。若 \(C=\{p\}\)，仿射包络是零维空间 \(\{p\}\)，相对内部等于 \(C\)，相对边界为空。所谓“内部”必须说明在哪个空间中计算。
\end{example}

\begin{proposition}[相对内部的存在与闭包]\label{g1:prop:relative-interior}
	设 \(C\subseteq\mathbb R^n\) 为非空凸集，\(F=\operatorname{aff}C\)。则 \(\operatorname{ri}C\ne\varnothing\)。若 \(a\in\operatorname{ri}C\)、\(b\in\overline C\) 和 \(0\leq t<1\)，则 \((1-t)a+tb\in\operatorname{ri}C\)。并且
	\[
		\overline{\operatorname{ri}C}=\overline C,\qquad
		\operatorname{ri}(\overline C)=\operatorname{ri}C.
	\]
\end{proposition}
\begin{proof}
	令 \(d=\dim F\)。当 \(d=0\) 时，\(C\) 为单点，全部结论成立。若 \(d\geq1\)，从 \(C\) 中逐次选择点，直至得到张成 \(F\) 的仿射无关组 \(p_0,\ldots,p_d\)：只要已有点的仿射包络不是 \(F\)，就还可从 \(C\) 选到包络之外的点，每次使维数增加一。由\cref{g1:prop:affine-coordinates}，相对于这组点的重心坐标 \(\lambda_i\) 是 \(F\) 上的连续仿射函数。重心 \(a=(p_0+\cdots+p_d)/(d+1)\) 的各坐标均为正，因此在 \(F\) 中的某个小球上各 \(\lambda_i\) 仍为正。这个小球包含于这些点的凸包，再包含于 \(C\)，故 \(a\in\operatorname{ri}C\)。

	现任取 \(a\in\operatorname{ri}C\)，选 \(r>0\) 使相对球 \(B_F(a,r)\subseteq C\)。固定 \(t<1\)，取 \(b_j\in C\) 趋于 \(b\)。由凸性，
	\[
		B_F((1-t)a+tb_j,(1-t)r)\subseteq C.
	\]
	当 \(j\) 充分大时，球心与 \((1-t)a+tb\) 的距离小于 \((1-t)r/2\)，故以后一个点为中心、半径为 \((1-t)r/2\) 的相对球仍包含于 \(C\)。这证明线段上除 \(b\) 外的各点都在相对内部。令 \(t\) 趋于 \(1\)，可知每个 \(b\in\overline C\) 都属于 \(\overline{\operatorname{ri}C}\)；反向包含来自 \(\operatorname{ri}C\subseteq C\)。

	由于 \(F\) 闭且 \(C\subseteq\overline C\subseteq F\)，\(\overline C\) 的仿射包络仍为 \(F\)，并有 \(\operatorname{ri}C\subseteq\operatorname{ri}(\overline C)\)。若 \(x\in\operatorname{ri}(\overline C)\)，固定一个 \(a\in\operatorname{ri}C\)。\(x=a\) 时无须再证；否则可取充分小的 \(\varepsilon>0\)，使 \(b=x+\varepsilon(x-a)\in\overline C\)。此时
	\[
		x=\frac{\varepsilon}{1+\varepsilon}a+\frac{1}{1+\varepsilon}b,
	\]
	右边是刚刚讨论过的线段内部点，故 \(x\in\operatorname{ri}C\)，得到最后一个等式。
\end{proof}

特别地，具有非空普通内部的凸集满足 \(\operatorname{int}(\overline C)=\operatorname{int}C\)。对这样的凸集，取闭包可以补入原来缺失的边界点，内部则保持不变。

\subsection{最近点投影及支撑超平面}
\label{g1:subsec:0301:02}

\cref{g1:prop:projection}给出子空间中的唯一最近点：把 \(x\) 分成沿子空间的分量和正交分量。对于一般凸集，最近点处不一定有一条可以向两边自由移动的直线，所以“与所有允许方向正交”需要改成“沿任何允许方向移动都不会缩短距离”。这会把等式条件变成一个内积不等式。

\begin{theorem}[闭凸集的最近点投影]\label{g1:thm:convex-projection}
	设 \(C\subseteq\mathbb R^n\) 为非空闭凸集。每个 \(x\in\mathbb R^n\) 都有唯一最近点 \(P_Cx\in C\)，满足
	\[
		\lVert x-P_Cx\rVert=d(x,C)\coloneqq\inf_{y\in C}\lVert x-y\rVert.
	\]
	而且，对 \(p\in C\)，有等价刻画
	\begin{equation}\label{g1:eq:convex-projection-variation}
		p=P_Cx\quad\Longleftrightarrow\quad
		\langle x-p,y-p\rangle\leq0\quad\text{对所有 }y\in C.
	\end{equation}
	因此对所有 \(y\in C\)，
	\begin{equation}\label{g1:eq:convex-projection-distance}
		\lVert x-y\rVert^2\geq\lVert x-P_Cx\rVert^2+\lVert y-P_Cx\rVert^2.
	\end{equation}
\end{theorem}
\begin{proof}
	先证明最小值能取到。取 \(c_0\in C\)，则 \(d(x,C)\leq\lVert x-c_0\rVert<\infty\)。按下确界定义可取 \(y_j\in C\)，使 \(\lVert x-y_j\rVert\to d(x,C)\)。这些距离有界，所以 \((y_j)\) 有界。有限维实分析中的有界序列子列定理保证它有收敛子列；因 \(C\) 闭，其极限 \(p\) 仍在 \(C\)，范数连续性给出 \(\lVert x-p\rVert=d(x,C)\)。这里没有要求整个 \(C\) 有界，只需要控制这一列候选点。

	若 \(p\) 是最近点，任取 \(y\in C\)。由凸性，\(p+t(y-p)\in C\) 对 \(0\leq t\leq1\) 成立，于是
	\[
		0\leq\lVert x-p-t(y-p)\rVert^2-\lVert x-p\rVert^2
		=-2t\langle x-p,y-p\rangle+t^2\lVert y-p\rVert^2.
	\]
	对 \(t>0\) 除以 \(t\)，再令 \(t\to0^+\)，得到所需内积不等式。反过来，若 \(p\in C\) 满足该不等式，则
	\[
		\lVert x-y\rVert^2
		=\lVert x-p\rVert^2+\lVert y-p\rVert^2-2\langle x-p,y-p\rangle
		\geq\lVert x-p\rVert^2+\lVert y-p\rVert^2.
	\]
	它首先说明 \(p\) 是最近点；若 \(y\) 也是最近点，左右最短距离相等，便迫使 \(\lVert y-p\rVert=0\)。这同时证明唯一性和最后一个不等式。
\end{proof}

映射 \(P_C\) 称为到 \(C\) 的\term[zuijindiantouying]{最近点投影}[nearest-point projection]。如果 \(C=W\) 是子空间，将 \(y=p+w\) 和 \(y=p-w\) 同时代入式\eqref{g1:eq:convex-projection-variation}，就得到 \(\langle x-p,w\rangle=0\)，恢复了第一章的正交投影。对于一般凸集，只能沿集合内部的方向移动，因此只能得到不等式。
\symidx{\(P_Cx\)：闭凸集的最近点投影}

\begin{example}[球上的投影及两项条件的作用]\label{g1:exm:ball-projection}
	对 \(C=\{y:\lVert y\rVert\leq r\}\)，其中 \(r>0\)，有
	\[
		P_Cx=\begin{cases}x,&\lVert x\rVert\leq r,\\[2pt]
		r x/\lVert x\rVert,&\lVert x\rVert>r.
		\end{cases}
	\]
	第二种情形中，令 \(p=rx/\lVert x\rVert\)，则 \(x-p\) 是 \(p\) 的正倍数，而对每个 \(y\in C\)，\(\langle p,y-p\rangle\leq r\lVert y\rVert-r^2\leq0\)，所以可用投影刻画验证该式。这个投影通常不是线性的。若把闭球改成开球，外点的距离下确界仍为 \(\lVert x\rVert-r\)，却没有最近点，说明闭性用于保证存在。若取闭集 \(C=\{-1,1\}\subseteq\mathbb R\)，则 \(0\) 有两个最近点，说明缺少凸性会破坏唯一性。
\end{example}

最近点还给出一个方向：从集合指向外点的向量 \(x-P_Cx\)。投影不等式说，整个集合都位于过 \(P_Cx\)、垂直于这个向量的超平面的一侧，如\cref{g1:fig:convex-separation}所示。

\begin{theorem}[点与闭凸集的分离]\label{g1:thm:point-closed-separation}
	设 \(C\subseteq\mathbb R^n\) 为非空闭凸集，\(x\notin C\)。令 \(p=P_Cx\)、\(u=(x-p)/\lVert x-p\rVert\)，则 \(\lVert u\rVert=1\)，且
	\begin{equation}\label{g1:eq:point-closed-separation}
		\langle u,y\rangle\leq\langle u,p\rangle
		<\langle u,x\rangle=\langle u,p\rangle+d(x,C)
		\quad(y\in C).
	\end{equation}
	特别地，存在 \(\alpha\in\mathbb R\) 和 \(\varepsilon>0\)，使 \(\langle u,y\rangle\leq\alpha-\varepsilon\) 对所有 \(y\in C\) 成立，而 \(\langle u,x\rangle\geq\alpha+\varepsilon\)。
\end{theorem}
\begin{proof}
	由 \(x\notin C\) 和最近点存在可知 \(x\ne p\)，故 \(d(x,C)=\lVert x-p\rVert>0\)。将式\eqref{g1:eq:convex-projection-variation}除以此正数，就得到 \(\langle u,y-p\rangle\leq0\)。同时 \(\langle u,x-p\rangle=\lVert x-p\rVert\)，得所示不等式。取 \(\alpha=(\langle u,p\rangle+\langle u,x\rangle)/2\)、\(\varepsilon=d(x,C)/2\) 即得最后一项。
\end{proof}

\begin{figure}[htbp]
	\centering
	\input{Figures/Chapter03-Separation}
	\caption[最近点、支撑超平面与分离超平面]{最近点、支撑超平面与分离超平面。实线过最近点 \(p\)，虚线位于 \(p\) 与外点 \(x\) 之间；\(C\) 中所有方向 \(y-p\) 与 \(x-p\) 的内积都非正。}
	\label{g1:fig:convex-separation}
\end{figure}

\begin{definition}[支撑超平面]\label{g1:def:supporting-hyperplane}
	设 \(C\subseteq\mathbb R^n\) 非空，\(p\in\overline C\)。若 \(u\ne0\)，且 \(\langle u,y\rangle\leq\langle u,p\rangle\) 对每个 \(y\in C\) 成立，则称
	\[
		H=\{y:\langle u,y\rangle=\langle u,p\rangle\}
	\]
	为 \(C\) 在 \(p\) 处的\term[zhichengchaopingmian]{支撑超平面}[supporting hyperplane]，向量 \(u\) 指向所选的外侧。
\end{definition}

定义允许整个 \(C\) 都在超平面上，因而也适用于低维集合。过最近点的超平面已经给出了某些支撑超平面，但一般边界点未必事先被表示成某个外点的投影。可以让外点趋近指定边界点，利用单位法向量的紧性留下一个极限方向。

\begin{theorem}[边界点的支撑超平面]\label{g1:thm:supporting-hyperplane}
	设 \(C\subseteq\mathbb R^n\) 为非空闭凸集。对每个 \(p\in\partial C\)，都存在 \(u\in\mathbb R^n\)、\(\lVert u\rVert=1\)，使
	\[
		\langle u,y-p\rangle\leq0\qquad(y\in C).
	\]
	若在 \(F=\operatorname{aff}C=a+L\) 中讨论，\(\dim F\geq1\)，则每个 \(p\in\operatorname{rbd}C\) 都有法向量 \(u\in L\setminus\{0\}\) 的相对支撑超平面；这个超平面不会包含整个 \(C\)。
\end{theorem}
\begin{proof}
	由 \(p\in\partial C\)，可选 \(x_j\notin C\) 且 \(x_j\to p\)。令 \(p_j=P_Cx_j\)。因为 \(p\in C\)，有
	\[
		\lVert p_j-p\rVert\leq\lVert p_j-x_j\rVert+\lVert x_j-p\rVert
		=d(x_j,C)+\lVert x_j-p\rVert\leq2\lVert x_j-p\rVert\to0.
	\]
	取单位向量 \(u_j=(x_j-p_j)/\lVert x_j-p_j\rVert\)。单位球面紧，故可取子列使 \(u_j\to u\)，其中 \(\lVert u\rVert=1\)。对任意固定 \(y\in C\)，投影刻画给出 \(\langle u_j,y-p_j\rangle\leq0\)，沿该子列取极限得 \(\langle u,y-p\rangle\leq0\)。同一个收敛子列对所有这些不等式都适用，因为选取它只用到了 \((u_j)\)，不依赖于 \(y\)。

	对相对边界，在 \(F\) 内选择趋于 \(p\) 的外点，并以 \(L\) 的标准正交坐标重复上述证明，所得单位法向量属于 \(L\)。若相应超平面包含整个 \(C\)，则 \(\langle u,y-p\rangle=0\) 对所有 \(y\in C\) 成立。由于 \(p\in C\) 且这些差向量张成 \(L\)，便有 \(u\perp L\)，与 \(0\ne u\in L\) 矛盾。
\end{proof}

普通边界与相对边界在这里给出不同信息。例如，线段的内点在平面中也有支撑直线，但这条直线只能包含整条线段；在其一维仿射包络内，它并不是边界点，也就不需要有支撑超平面。零维非空凸集没有相对边界，故相对版本在零维时没有待讨论的点。

\subsection{凸集的弱分离、严格分离与强分离}
\label{g1:subsec:0301:03}

分离两个集合时，“在两侧”可能只要求不越过同一个超平面，也可能要求两边各自留下间隙。先明确这些量词，才能判断闭性、开性或紧性分别解决什么问题。下列用语中的法向量总是非零向量，否则常数函数会给出毫无几何内容的等式。

\begin{definition}[分离的三种条件]\label{g1:def:convex-separation}
	设 \(A,B\subseteq\mathbb R^n\) 非空。
	\begin{enumerate}
		\item 若存在 \(u\ne0\)、\(\alpha\in\mathbb R\)，使 \(\langle u,a\rangle\leq\alpha\leq\langle u,b\rangle\) 对所有 \(a\in A,b\in B\) 成立，称 \(A,B\) 被该超平面\term[ruofenli]{弱分离}[weakly separated]。
		\item 若上述两个不等号对每个 \(a,b\) 都严格成立，称为\term[yangefenli]{严格分离}[strict separation]。
		\item 若存在 \(u\ne0\)、\(\alpha\in\mathbb R\)、\(\varepsilon>0\)，使 \(\langle u,a\rangle\leq\alpha-\varepsilon\) 且 \(\langle u,b\rangle\geq\alpha+\varepsilon\) 对所有 \(a,b\) 成立，称为\term[qiangfenli]{强分离}[strong separation]。
	\end{enumerate}
\end{definition}

强分离等价于存在 \(u\ne0\) 使 \(\sup_{a\in A}\langle u,a\rangle<\inf_{b\in B}\langle u,b\rangle\)：两个端值之间取中点和半间距即可。严格分离只要求逐点严格，两个端值仍可能相同。文献对“严格分离”和“强分离”的命名并不完全统一，本书按上面的量词区分；比较定理时应比较不等式本身。分离与支撑的基本观点可参阅 \cite[第2章，第2.5节]{BoydVandenberghe2004}。

\begin{theorem}[一个集合开时的弱分离]\label{g1:thm:convex-weak-separation}
	设 \(A,B\subseteq\mathbb R^n\) 为非空凸集，且至少一个为开集。则 \(A\cap B=\varnothing\) 当且仅当它们可以弱分离。若 \(A\) 开，分离不等式在 \(A\) 一侧必严格；若 \(B\) 开，在 \(B\) 一侧必严格。
\end{theorem}
\begin{proof}
	先设 \(A\cap B=\varnothing\)，考虑差集
	\[
		D=B-A=\{b-a:a\in A,\ b\in B\}.
	\]
	分别对两个分量取凸组合可知 \(D\) 凸。若 \(B\) 开，\(D=\bigcup_{a\in A}(B-a)\) 是开集；若 \(A\) 开，改用 \(D=\bigcup_{b\in B}(b-A)\) 得到同一结论。\(D\) 非空，且 \(0\notin D\)。由\cref{g1:prop:relative-interior}在满维开凸集上的结论，\(\operatorname{int}\overline D=D\)。因此 \(0\) 要么在闭凸集 \(\overline D\) 之外，要么在它的边界。

	若 \(0\notin\overline D\)，应用\cref{g1:thm:point-closed-separation}并把法向量反向，得到某个 \(u\ne0\)，使 \(\langle u,d\rangle>0\) 对所有 \(d\in D\) 成立。若 \(0\in\partial\overline D\)，则对 \(\overline D\) 在 \(0\) 处应用\cref{g1:thm:supporting-hyperplane}，也把法向量反向，得到 \(\langle u,d\rangle\geq0\)。两种情形都有
	\[
		\langle u,a\rangle\leq\langle u,b\rangle\qquad(a\in A,\ b\in B).
	\]
	固定 \(b_0\in B\) 可知左边的上确界有限，固定 \(a_0\in A\) 可知右边的下确界有限。故可取实数 \(\alpha\) 介于这两个端值之间，得到弱分离。

	现设已有弱分离，且 \(A\) 开。若某个 \(a\in A\) 满足 \(\langle u,a\rangle=\alpha\)，则充分小的 \(t>0\) 使 \(a+tu\in A\)，但 \(\langle u,a+tu\rangle=\alpha+t\lVert u\rVert^2>\alpha\)，矛盾。因此 \(A\) 一侧全部严格。若 \(B\) 开，使用 \(b-tu\) 得到 \(B\) 一侧全部严格。至少一侧严格，便不可能有同一点同时属于 \(A,B\)，这也证明了反向命题。
\end{proof}

这里的开性不仅用于构造证明，还保证弱分离能反推出不相交。若去掉开性，\(A=B=\{0\}\) 就可以被任何过原点的超平面弱分离，两者却相交。

\begin{theorem}[紧凸集与闭凸集的强分离]\label{g1:thm:compact-closed-separation}
	设 \(A,B\subseteq\mathbb R^n\) 为非空凸集，\(A\) 紧、\(B\) 闭，且 \(A\cap B=\varnothing\)。则它们能够强分离；更具体地，存在单位向量 \(u\)，使
	\[
		\inf_{b\in B}\langle u,b\rangle-\sup_{a\in A}\langle u,a\rangle
		=d(A,B)>0,
	\]
	其中 \(d(A,B)=\inf\{\lVert b-a\rVert:a\in A,b\in B\}\)。
\end{theorem}
\begin{proof}
	仍令 \(D=B-A\)，它非空且凸。先说明此时 \(D\) 闭。若 \(d_j=b_j-a_j\to d\)，紧性允许取子列使 \(a_j\to a\in A\)。于是 \(b_j=d_j+a_j\to d+a\)，由 \(B\) 闭可知 \(d+a\in B\)，所以 \(d\in D\)。又因 \(A,B\) 不交，\(0\notin D\)。

	取 \(q=P_D0\)，则 \(q\ne0\)。把投影刻画中的外点取为 \(0\)，得到 \(\langle -q,d-q\rangle\leq0\)，即
	\[
		\langle q,d\rangle\geq\lVert q\rVert^2\qquad(d\in D).
	\]
	令 \(u=q/\lVert q\rVert\)，便有 \(\langle u,b\rangle-\langle u,a\rangle\geq\lVert q\rVert\) 对所有 \(a,b\) 成立。因为 \(q\in D\)，可写成 \(q=b_0-a_0\)，此处取得等号。固定 \(b_0\) 说明 \(a_0\) 取得 \(\langle u,a\rangle\) 的上确界，固定 \(a_0\) 说明 \(b_0\) 取得 \(\langle u,b\rangle\) 的下确界。两者之差就是 \(\lVert q\rVert=d(0,D)=d(A,B)>0\)，在两个端值中间取超平面即得强分离。
\end{proof}

\begin{example}[闭而不紧时没有统一间距]\label{g1:exm:zero-distance-separation}
	在 \(\mathbb R^2\) 中令
	\[
		A=\{(x,y):y\leq0\},\qquad
		B=\{(x,y):x>0,\ y\geq1/x\}.
	\]
	两者闭、凸且不交。\(B\) 的闭性可直接检查：若 \((x_j,y_j)\in B\) 收敛到有限点 \((x,y)\)，则 \(x\geq0\)；若 \(x=0\)，\(y_j\geq1/x_j\to\infty\)，与收敛矛盾。因此 \(x>0\)，取极限得到 \(y\geq1/x\)。它的凸性来自 \(f(x)=1/x\) 在 \(x>0\) 上满足 \(f''(x)=2/x^3>0\)，从而 \(f((1-t)x_1+tx_2)\leq(1-t)f(x_1)+tf(x_2)\)。这项一元分析结论也可由导数单调及割线斜率比较得到。

	点 \((j,0)\in A\) 与 \((j,1/j)\in B\) 的距离趋于零，故 \(d(A,B)=0\)。它们由 \(y=0\) 弱分离，但连逐点严格分离也不可能：若 \(u=(u_1,u_2)\ne0\) 使 \(\sup_A\langle u,a\rangle\) 有限，则 \(A\) 在水平方向无界迫使 \(u_1=0\)，向下无界又迫使 \(u_2\geq0\)，故 \(u_2>0\)。此时 \(\sup_A\langle u,a\rangle=0=\inf_B\langle u,b\rangle\)，而 \(A\) 上确实取到零值；严格分离要求超平面常数同时大于 \(0\) 且不大于 \(0\)，无法成立。
\end{example}

距离为零本身只排除强分离，不必排除逐点严格分离。例如 \(\{y<0\}\) 与 \(\{y>0\}\) 被 \(y=0\) 严格分离，距离却仍为零。相反，若有强分离不等式，则由内积不等式得到 \(\lVert b-a\rVert\geq2\varepsilon/\lVert u\rVert\)，从而距离严格为正。

\subsection{闭凸集的支撑函数表示}
\label{g1:subsec:0301:04}

一个超平面可以排除一个外点；把所有方向的信息放在一起，就能恢复整个闭凸集。对给定方向 \(u\)，只需记录 \(\langle u,x\rangle\) 在集合上最大的可能值，未必要记录取得它的是哪个点。无界集合在某些方向上可能没有有限上界，因此定义中要允许 \(+\infty\)。

\begin{definition}[支撑函数]\label{g1:def:support-function}
	设 \(C\subseteq\mathbb R^n\) 非空。它的\term[zhichenghanshu]{支撑函数}[support function]为
	\[
		h_C:\mathbb R^n\longrightarrow\mathbb R\cup\{+\infty\},\qquad
		h_C(u)\coloneqq\sup_{x\in C}\langle x,u\rangle.
	\]
	这里 \(h_C(0)=0\)；由于 \(C\) 非空，上确界不会是 \(-\infty\)。
	\symidx{\(h_C(u)\)：支撑函数}
\end{definition}

\begin{proposition}[闭凸集的支撑函数表示]\label{g1:prop:support-function-representation}
	设 \(C\subseteq\mathbb R^n\) 为非空闭凸集，则
	\begin{equation}\label{g1:eq:support-halfspaces}
		C=\bigcap_{u\in\mathbb R^n}\{x:\langle x,u\rangle\leq h_C(u)\}.
	\end{equation}
	其中 \(h_C(u)=+\infty\) 时对应的集合理解为整个 \(\mathbb R^n\)，故真正起作用的是具有有限支撑值的方向。对于非空闭凸集 \(C,D\)，有
	\[
		C\subseteq D\quad\Longleftrightarrow\quad h_C(u)\leq h_D(u)\text{ 对所有 }u\in\mathbb R^n.
	\]
\end{proposition}
\begin{proof}
	由上确界定义，\(C\) 中每个点都满足式\eqref{g1:eq:support-halfspaces}右边全部不等式。若 \(x\notin C\)，用\cref{g1:thm:point-closed-separation}取 \(p=P_Cx\) 及单位向量 \(u=(x-p)/\lVert x-p\rVert\)。因为 \(p\in C\)，该定理给出的 \(\langle y,u\rangle\leq\langle p,u\rangle\) 说明
	\[
		h_C(u)=\langle p,u\rangle<\langle x,u\rangle.
	\]
	所以这个方向的半空间已经排除了 \(x\)，证明反向包含。

	若 \(C\subseteq D\)，对同一线性函数在较小集合上取上确界，直接得到 \(h_C\leq h_D\)。反过来，若各方向上支撑函数有序，则每个 \(x\in C\) 都满足 \(\langle x,u\rangle\leq h_C(u)\leq h_D(u)\)。对 \(D\) 使用刚证得的表示式可知 \(x\in D\)。
\end{proof}

\begin{proposition}[支撑函数保留的信息]\label{g1:prop:support-function-basic}
	设 \(S\subseteq\mathbb R^n\) 非空，记 \(K=\overline{\operatorname{conv}S}\)。则 \(h_S=h_K\)。对任意 \(u,v\in\mathbb R^n\)、\(t>0\)，有
	\[
		h_S(tu)=t h_S(u),\qquad h_S(u+v)\leq h_S(u)+h_S(v).
	\]
	若 \(S\) 紧，则 \(h_S\) 处处有限且每个方向的上确界都能取到。
\end{proposition}
\begin{proof}
	令 \(M=h_S(u)\)。若 \(M<\infty\)，则 \(\langle x,u\rangle\leq M\) 对 \(S\) 中的点成立；按非负权重相加，这个不等式对凸组合仍成立；再由连续性，它对凸组合的极限也成立。故 \(h_K(u)\leq h_S(u)\)，反向不等式由 \(S\subseteq K\) 得到。若 \(h_S(u)=+\infty\)，同一包含直接使 \(h_K(u)=+\infty\)。

	对正数 \(t\)，乘以 \(t\) 保持上界与上确界，得齐次等式。对每个 \(x\in S\)，有 \(\langle x,u+v\rangle\leq h_S(u)+h_S(v)\)，取上确界得次可加性；右边为 \(+\infty\) 时不等式同样有意义。最后，紧集有界，且线性函数连续，所以它在非空紧集上取得有限最大值。
\end{proof}

因而支撑函数不区分一个集合与其闭凸包。例如，圆周和闭圆盘具有同一个支撑函数；开圆盘与闭圆盘也一样。式\eqref{g1:eq:support-halfspaces}中的闭性与凸性各有作用，删掉它们就不能据支撑函数恢复原来的集合。

\begin{example}[球、线段与半空间的支撑函数]\label{g1:exm:support-functions}
	对闭球 \(C=\{x:\lVert x-a\rVert\leq r\}\)，其中 \(r\geq0\)，内积不等式及沿 \(u\) 方向取点给出
	\[
		h_C(u)=\langle a,u\rangle+r\lVert u\rVert.
	\]
	当 \(u\ne0\) 时，上确界在 \(a+ru/\lVert u\rVert\) 取得；\(u=0\) 时每点都取得零值。对线段 \(C=[a,b]\)，线性函数在线段上的值是端点值的凸组合，故 \(h_C(u)=\max\{\langle a,u\rangle,\langle b,u\rangle\}\)。这也说明有些方向上的最大值会在整条线段上取得。

	对半空间 \(H=\{x:x_1\leq c\}\subseteq\mathbb R^n\)，若 \(u=t e_1\)、\(t\geq0\)，则 \(h_H(u)=tc\)；其余方向上 \(h_H(u)=+\infty\)。事实上，若 \(u\) 有非零的其余坐标，沿对应的无约束方向就可使配对任意大；若 \(u_1<0\)，令 \(x_1\to-\infty\) 也可做到。式\eqref{g1:eq:support-halfspaces}虽列出所有方向，此例只有 \(e_1\) 的正倍方向提供同一个实质约束。
\end{example}

支撑函数保存的这些方向信息还能够测量距离：\cref{g1:prop:distance-support-formula}将利用本节的投影与分离，把点到闭凸集的距离写成各方向支撑差值的最大值。

% !TeX root = ../../Book1.tex
% 纲要标题：### 3.2 有限维性的组合后果
% 纲要位置：几何学计划纲要.md，第 211 行。
% 正文状态：初稿；尚未正式审核。

% 纲要内容开始
% 
% * 仿射依赖与 Carathéodory 定理
% * Radon 分割的线性代数证明
% * Helly 定理的归纳论证
% * 去掉维数或凸性条件后的失败
% 
% 纲要内容结束

\section{有限维性的组合后果}
\label{g1:sec:0302}

在平面上，用十个点的加权平均表示一个位置时，是否真的需要十个点？若许多凸区域每三个都有公共点，是否就能找到同时属于所有区域的点？两个问题看似一个关于表示、一个关于相交，却都受仿射维数控制。本节始终在有限维实空间中讨论；关键是把“点太多”转化为仿射依赖，再把依赖关系中的正负系数转化为凸组合。

\subsection{仿射依赖与 Carathéodory 定理}
\label{g1:subsec:0302:01}

按\secref{g1:sec:0301}的凸包定义，若 \(x\in\operatorname{conv}S\)，就有某个有限表示 \(x=\sum_{i=1}^m\lambda_i p_i\)，其中 \(p_i\in S\)、\(\lambda_i\ge0\)、\(\sum_i\lambda_i=1\)。定义没有限制 \(m\)；即使 \(S\) 很大，也要等写出具体表示之后才知道使用了多少点。Carathéodory 定理给出一个只依赖维数的统一上界。\footnote{Constantin Carathéodory（1873--1950），希腊数学家，在变分法、测度论与函数论中作出重要贡献。}

\begin{theorem}[Carathéodory 定理]\label{g1:thm:caratheodory}
	\termidx[Caratheodorydingli]{Carathéodory 定理}\foreignidx{Carathéodory's theorem}设 \(S\subseteq\mathbb R^n\) 非空，\(d=\dim\operatorname{aff}(S)\)。则每个 \(x\in\operatorname{conv}S\) 都可写成 \(S\) 中至多 \(d+1\) 个仿射无关点的凸组合。特别地，至多需要 \(n+1\) 个点。
\end{theorem}
\begin{proof}
	取 \(x=\sum_{i=1}^m\lambda_i p_i\) 的一个有限凸组合表示，删去零系数并合并重复点，使各 \(\lambda_i>0\)。如果这些点已仿射无关，证明就结束了，因为 \(d\) 维仿射空间中至多有 \(d+1\) 个仿射无关点。下面说明：只要它们仿射相关，就能减少至少一个点，同时保持表示的点 \(x\) 不变。

	回顾\cref{g1:def:affine-frame}之后的对称判别：仿射相关意味着存在不全为零的实数 \(c_i\)，满足
	\[
		\sum_{i=1}^m c_i=0,\qquad \sum_{i=1}^m c_i p_i=0.
	\]
	第二个等式是向量等式；若将所有 \(p_i\) 平移同一向量，第一个等式保证它不变。因此，把系数由 \(\lambda_i\) 改成 \(\lambda_i-tc_i\)，既不会改变系数和，也不会改变所得点。剩下的障碍只有非负性：怎样让某个系数降到零，又不让别的系数先变负？

	由于 \(\sum_i c_i=0\) 且 \((c_i)\ne0\)，其中既有正数也有负数。取
	\[
		t=\min_{c_i>0}\frac{\lambda_i}{c_i}>0,\qquad \mu_i=\lambda_i-tc_i.
	\]
	当 \(c_i>0\) 时，最小值的选择保证 \(\mu_i\ge0\)，而至少一个这样的 \(\mu_i=0\)；当 \(c_i\le0\) 时，\(\mu_i\ge\lambda_i>0\)。又有
	\[
		\sum_i\mu_i=1,\qquad \sum_i\mu_i p_i=x.
	\]
	删去零系数后，得到使用更少点的凸组合。只要剩余点仍仿射相关，就重复此操作。点数每次严格下降，故有限次后停止；此时剩余点仿射无关，最多为 \(d+1\) 个。
\end{proof}

沿仿射依赖移动全部权重，既保持点的位置和权重总和，又能用最小比值守住非负条件。这里用的是 \(S\) 的仿射维数。例如，一批点虽然写成 \(\mathbb R^{100}\) 中的坐标，但若全在同一个仿射平面内，凸包中的每个点仍只需三个点来表示。

\begin{example}[沿依赖关系减少权重]\label{g1:exm:convex-coefficient-reduction}
	令 \(p_1=(0,0),p_2=(2,0),p_3=(2,2),p_4=(0,2)\)，并取
	\[
		x=\frac18p_1+\frac38p_2+\frac14p_3+\frac14p_4=\left(\frac54,1\right).
	\]
	关系 \(p_1-p_2+p_3-p_4=0\) 的系数和为零。按证明取 \(t=\min\{1/8,1/4\}=1/8\)，便得到
	\[
		x=\frac12p_2+\frac18p_3+\frac38p_4.
	\]
	新系数非负且和为 \(1\)，故 \(x\) 落在由后三个顶点构成的三角形内。定理只保证存在一个短表示，并不保证它唯一，也不保证每次消去都会得到使用点数最少的表示。例如正方形中心可用一条对角线的两个端点表示。
\end{example}

\subsection{Radon 分割的线性代数证明}
\label{g1:subsec:0302:02}

上面的证明把仿射依赖用于调整已有权重。同一依赖式还可以直接产生两份凸组合：将它分成正、负两边，两边经同一个正数归一化，就会成为表示同一点的两个凸组合。这给出 Radon 定理。\footnote{Johann Radon（1887--1956），奥地利数学家，研究测度论、变分法和微分几何，并提出由积分数据恢复函数的 Radon 变换。}

\begin{theorem}[Radon 分割]\label{g1:thm:radon}
	\termidx[Radondingli]{Radon 定理}\foreignidx{Radon's theorem}设 \(p_1,\ldots,p_{d+2}\) 为 \(\mathbb R^d\) 中的点，允许重复。则指标集可分成两个非空、不相交的集合 \(I,J\)，使
	\[
		I\cup J=\{1,\ldots,d+2\},\qquad
		\operatorname{conv}\{p_i:i\in I\}\cap\operatorname{conv}\{p_j:j\in J\}\ne\varnothing.
	\]
	这样的指标分组称为\term[Radonfenge]{Radon 分割}[Radon partition]。
\end{theorem}
\begin{proof}
	把每个点提升为 \(\mathbb R^{d+1}\) 中的向量 \((p_i,1)\)。有 \(d+2\) 个向量而空间维数为 \(d+1\)，所以存在不全为零的实系数 \(c_i\)，使
	\[
		\sum_i c_i(p_i,1)=0.
	\]
	最后一个坐标给出 \(\sum_i c_i=0\)，其余坐标给出 \(\sum_i c_i p_i=0\)。令 \(I_+=\{i:c_i>0\}\)、\(I_-=\{i:c_i<0\}\)，它们都非空，且
	\[
		a\coloneqq\sum_{i\in I_+}c_i=-\sum_{j\in I_-}c_j>0.
	\]
	把正、负项分别移到等号两边并除以 \(a\)，得到
	\[
		z\coloneqq\sum_{i\in I_+}\frac{c_i}{a}p_i
		=\sum_{j\in I_-}\frac{-c_j}{a}p_j.
	\]
	两边系数都非负且和为 \(1\)，故 \(z\) 同时属于两组点的凸包。最后把所有零系数指标任意分配到两组，例如全部加入 \(I_+\)，令 \(I=\{i:c_i\ge0\}\)、\(J=I_-\)。加入点不会缩小凸包，因此交点 \(z\) 仍然保留，且 \(I,J\) 构成所要求的分割。
\end{proof}

这里分割的是指标。若两个指标对应同一点，仍可把它们放到不同组；这也解释了为什么定理允许重复点。提升后的向量 \((p_i,1)\) 的最后一个坐标专门用来记录系数和条件，与\secref{g1:sec:0101}中仿射组合不依赖原点的条件完全对应。

\begin{example}[平面四点的两种分割]\label{g1:exm:radon-plane}
	若四点依次为凸四边形的顶点，两条对角线相交，对角顶点各成一组，得到两段凸包的交点。若一点位于另三点所成三角形的内部，就把内部点单独作为一组，另三个点作为一组。此时交集是内部点本身，不能要求两组各有两个点。\cref{g1:fig:radon}画出了这两种情形；共线、重合等退化情形由同一线性代数证明处理，不需要另凭图形分类。
\end{example}

\begin{figure}[htbp]
	\centering
	\input{Figures/Chapter03-Radon}
	\caption[平面四点的 Radon 分割]{平面四点的 Radon 分割。左图两组各有两个点，公共点 \(z\) 是对角线交点；右图一组只有 \(p_4=z\)，另一组的凸包是三角形。颜色表示分组。}
	\label{g1:fig:radon}
\end{figure}

\subsection{Helly 定理的归纳论证}
\label{g1:subsec:0302:03}

直线上的有限个非空区间若每两个相交，就有公共点。到了平面，每两个凸集相交已不够：三角形的三条边两两相交，却没有共同交点。正确的检查数量随维数增加为 \(d+1\)，这就是 Helly 定理。\footnote{Eduard Helly（1884--1943），奥地利数学家，研究凸集、积分方程与泛函分析，并在有界线性泛函的延拓理论中作出早期贡献。}

\begin{theorem}[有限族的 Helly 定理]\label{g1:thm:helly}
	\termidx[Hellydingli]{Helly 定理}\foreignidx{Helly's theorem}设 \(C_1,\ldots,C_N\) 是 \(\mathbb R^d\) 中的有限个凸集，\(N\ge1\)。若其中每个至多含 \(d+1\) 个集合的非空子族都有非空交集，则
	\[
		\bigcap_{i=1}^N C_i\ne\varnothing.
	\]
	不要求这些凸集闭、有界或有非空内部。
\end{theorem}
\begin{proof}
	当 \(N\le d+1\) 时，全族本身就是假设中的子族。第一个需要证明的情形是 \(N=d+2\)：每次删去一个集合，剩下的 \(d+1\) 个集合有公共点。我们先利用这些点和 Radon 分割找到全族的交点，再把任意较大的族压缩为少一个集合的族。

	设 \(N=d+2\)。对每个 \(i\)，选取
	\[
		p_i\in\bigcap_{j\ne i}C_j.
	\]
	点 \(p_i\) 只可能不属于 \(C_i\)。对这 \(d+2\) 个点应用\cref{g1:thm:radon}，得到分割 \(I,J\) 及两凸包的公共点 \(z\)。任意固定一个指标 \(k\)：若 \(k\in I\)，则 \(J\) 中每个指标都不等于 \(k\)，故所有 \(p_j\)（\(j\in J\)）都在 \(C_k\) 中。由 \(C_k\) 凸，整个 \(\operatorname{conv}\{p_j:j\in J\}\) 都包含在 \(C_k\) 中，特别是 \(z\in C_k\)。若 \(k\in J\)，改用 \(I\) 中各点，同样得到 \(z\in C_k\)。因此 \(z\) 属于全部 \(d+2\) 个集合。

	现在对 \(N\) 归纳。设 \(N>d+2\)，且结论对少于 \(N\) 个集合成立。把末尾两个集合合并为
	\[
		D=C_{N-1}\cap C_N,
	\]
	考虑由 \(C_1,\ldots,C_{N-2},D\) 组成的 \(N-1\) 个凸集。这样合并不改变全族的交集；需要检查的是，新族是否仍满足小子族相交条件。

	任取新族中至多 \(d+1\) 个集合。若不含 \(D\)，原假设直接给出非空交集。若含 \(D\)，把它重新展开为 \(C_{N-1},C_N\)，就得到原族中至多 \(d+2\) 个集合。展开后至多有 \(d+1\) 个时，交集由原假设非空；恰有 \(d+2\) 个时，这个子族的每个至多 \(d+1\) 元子族仍满足原假设，所以由已经证明的 \(d+2\) 情形，它的全交也非空。于是新族满足归纳假设，其全交非空，而这正是原族的全交。
\end{proof}

证明中凸性的作用很具体：所有 \(p_j\) 都在 \(C_k\) 中，才能把这些点的凸组合 \(z\) 也放进 \(C_k\)。维数则先控制产生仿射依赖所需的点数，再决定要检查的子族大小。对于 \(N\ge d+1\)，只检查恰有 \(d+1\) 个集合的子族已经足够，因为较小子族可以扩充到这个大小；定理采用“至多”的表述，是为了同时覆盖 \(N<d+1\) 的情形。

\ProseParagraphBreak

若所有集合都在一个 \(r\) 维仿射子空间中，就可通过仿射坐标把该子空间看作 \(\mathbb R^r\)，检查数量相应降为 \(r+1\)。反过来说，若有限个凸约束无共同解，必能从中挑出至多 \(d+1\) 个已经互相矛盾的约束。这一逆否表述给出了规模受维数控制的不可行子族。相关定理及进一步推广可参阅\cite[第3.2, 3.4节]{Gallier2011Geometric}。

\subsection{去掉维数或凸性条件后的失败}
\label{g1:subsec:0302:04}

这三个定理中的数量界都是最优的。先看一个同时检验这些界的模型：取 \(d\ge1\)，令 \(v_0,\ldots,v_d\) 为 \(\mathbb R^d\) 中的仿射无关点，记其单纯形为 \(\Delta\)。由\cref{g1:prop:affine-coordinates}，单纯形中的每个点都有唯一的非负重心坐标 \((\lambda_0,\ldots,\lambda_d)\)，且这些坐标之和为 \(1\)。

\begin{example}[单纯形给出的最优界]\label{g1:exm:convex-dimension-sharp}
	单纯形的重心 \(b=(v_0+\cdots+v_d)/(d+1)\) 的全部重心坐标都为正。由坐标唯一性，\(b\) 不能由少于 \(d+1\) 个顶点作凸组合得到，所以 Carathéodory 定理中的 \(d+1\) 不能统一缩小。

	同样，这 \(d+1\) 个顶点不可能有 Radon 分割：若两个不相交顶点组的凸包有公共点，把等号两边的凸组合相减，就得到一个非零、系数和为零的仿射依赖关系，与仿射无关矛盾。因此保证分割所需的 \(d+2\) 个点也不能统一减少。

	最后，对每个 \(i\) 令
	\[
		F_i=\operatorname{conv}\{v_j:j\ne i\}.
	\]
	它在重心坐标中由 \(\lambda_i=0\) 描述。任取 \(d\) 个这样的面，设未取的下标为 \(k\)，顶点 \(v_k\) 就属于所取各面的交；但全部 \(d+1\) 个面相交会要求每个 \(\lambda_i=0\)，与系数和为 \(1\) 矛盾。因此 Helly 定理不能把检查数量从 \(d+1\) 降到 \(d\)。
\end{example}

维数上界也不能彻底取消。考虑所有仅有有限个非零坐标的实数列所成的向量空间，以 \(e_1,e_2,\ldots\) 表示标准坐标向量。对任意 \(m\)，点 \((e_1+\cdots+e_m)/m\) 都属于 \(\{e_i:i\ge1\}\) 的凸包，但要从这些向量作凸组合表示它，前 \(m\) 个坐标迫使 \(e_1,\ldots,e_m\) 全部出现。因此不存在与所在有限维子空间的维数无关的统一点数界。对任意事先指定的正整数 \(k\)，再把一个 \(k\) 维单纯形的 \(k+1\) 个面放入这个空间，由上一例，它们的每个至多含 \(k\) 个集合的子族均相交，全族却不相交。因此 Helly 定理也没有与维数无关的统一检查数量。

\begin{example}[非凸集合不能只检查小子族]\label{g1:exm:helly-nonconvex}
	在直线 \(\mathbb R\) 上取
	\[
		C_1=\{0,1\},\qquad C_2=\{1,2\},\qquad C_3=\{0,2\}.
	\]
	任意两个集合的交分别为 \(\{1\},\{0\},\{2\}\)，而三个集合没有公共点。它们都是闭且有界的集合，问题出在凸性。例如 \(0,2\in C_3\)，但其中点 \(1\notin C_3\)；Helly 证明中把若干已有点换成凸组合的步骤便失去了依据。
\end{example}

\begin{example}[无限族还需要额外条件]\label{g1:exm:helly-infinite-family}
	在 \(\mathbb R\) 中令 \(C_k=[k,\infty)\)，其中 \(k=1,2,\ldots\)。每个有限子族的交都是非空闭射线，但全族的交为空，因为一个实数不能大于等于所有正整数。这些集合甚至都是闭凸集，所以“有限族”不能直接删除。

	若只补上有界性，取 \(D_k=(0,1/k)\) 仍有同样的有限相交而全交为空现象。此时每个集合都凸且有界，却在端点处不闭。有限 Helly 定理所做的是把有限次检查的数量降下来；把有限相交再推进到无限相交，是另外一个需要紧性等条件的问题。
\end{example}

% !TeX root = ../../Book1.tex
% 纲要标题：### 3.3 多面体的两种表示与对偶
% 纲要位置：几何学计划纲要.md，第 218 行。
% 正文状态：初稿；尚未正式审核。

% 纲要内容开始
% 
% * 半空间表示与顶点表示
% * 有界多面体的等价描述
% * 面、法锥与面格
% * 极体及包含关系的反转
% * 单纯形、立方体和交叉多面体的对照
% 
% 纲要内容结束

\section{多面体的两种表示与对偶}
\label{g1:sec:0303}

要把一个三角形交给别人，可以给出它的三个顶点，也可以给出三条边所在的直线，并指出每条线的哪一侧应当保留。前者便于生成其中的点，后者便于判断一个给定点是否在里面。维数增加以后，这两种描述仍然等价，但“看见所有边”已不能代替证明。本节先说明有限的不等式为何产生有限顶点，再把面及其法向方向组织起来，最后用极体把顶点与不等式相互转换。

\subsection{半空间表示与顶点表示}
\label{g1:subsec:0303:01}

以下使用 \(\mathbb R^n\) 的标准内积。\secref{g1:sec:0301}中的闭半空间具有形式 \(\{x:\langle a,x\rangle\le b\}\)，其中 \(a\ne0\) 指向被排除的一侧。有限个这样的限制给出一种容易核验的描述。

\begin{definition}[两种有限描述]
	\label{g1:def:polyhedron-polytope}
	\(\mathbb R^n\) 中有限个闭半空间的交称为\term[duomianti]{多面体}[polyhedron]。写成
	\[
		P=\{x\in\mathbb R^n:\langle a_i,x\rangle\le b_i,\ 1\le i\le m\}
	\]
	的描述称为\term[bankongjianbiaoshi]{半空间表示}[H-representation]。有限个点的凸包称为\term[duobaoxing]{多胞形}[polytope]，写成
	\[
		P=\operatorname{conv}\{v_1,\ldots,v_N\}
	\]
	的描述称为\term[dingdianbiaoshi]{顶点表示}[V-representation]。字母 \(H\) 与 \(V\) 分别来自半空间与顶点的英文。允许表示中出现冗余不等式或冗余点；因此顶点表示中列出的点未必全是真正的顶点。
\end{definition}

空集作为退化情形，约定 \(\operatorname{conv}\varnothing=\varnothing\)；全空间可看作没有不等式时的交。在矩阵记法中写 \(Ax\le b\)，不等号逐坐标理解；零行只表示恒真或矛盾条件。等式 \(\langle a,x\rangle=b\) 可写成一对相反的不等式，所以线段、平面内的三角形等低维对象也在讨论范围内。

\begin{example}[一个三角形的两份数据]
	\label{g1:exm:triangle-hv}
	平面中的集合
	\[
		T=\{(x_1,x_2):x_1\ge0,\ x_2\ge0,\ x_1+x_2\le1\}
	\]
	等于 \(\operatorname{conv}\{(0,0),(1,0),(0,1)\}\)。事实上，任一点 \((x_1,x_2)\in T\) 都有分解
	\[
		(x_1,x_2)=(1-x_1-x_2)(0,0)+x_1(1,0)+x_2(0,1),
	\]
	三个系数非负且和为一；反过来，这种凸组合必满足三条不等式。增加不等式 \(x_1+x_2\le2\) 或增加生成点 \((1/3,1/3)\) 都不改变集合，却使数据变得冗余。
\end{example}

“顶点”需要一个与图画无关的定义。一个凸集中的点若处在某条非退化线段的内部，就可以继续分解；顶点则是这种分解停止的位置。

\begin{definition}[极点与顶点]
	\label{g1:def:extreme-vertex}
	设 \(C\subseteq\mathbb R^n\) 为凸集。点 \(v\in C\) 称为其\term[jidian]{极点}[extreme point]，若每当
	\[
		v=(1-t)x+ty,\qquad x,y\in C,\quad 0<t<1,
	\]
	都有 \(x=y=v\)。多面体或多胞形的极点称为\term[dingdian]{顶点}[vertex]，全体顶点记为 \(\operatorname{Vert}(P)\)。
\end{definition}

\symidx{\(\operatorname{Vert}(P)\)：多面体或多胞形的顶点集}
在有限点集的凸包中，极点必属于原来的点集：若一个点不在原名单中，却是名单中点的凸组合，那么任选一个正系数项，把其余项合成第二个点，就得到非平凡线段分解。因此有限点集提供的是顶点的候选名单，仍需剔除其中能够由其他点生成的点。相反，半空间表示直接给出限制，顶点藏在若干限制同时取等号的位置。

\subsection{有界多面体的等价描述}
\label{g1:subsec:0303:02}

有界性是两种有限描述等价的关键条件。例如第一象限由两个半空间给出，却不可能是有限点集的凸包，因为有限点集的凸包必有界。下面证明有界时不再需要另外的生成方向。

\begin{theorem}[有界多面体的两种表示]
	\label{g1:thm:polytope-hv}
	设 \(P\subseteq\mathbb R^n\)。则 \(P\) 是有界多面体，当且仅当 \(P\) 是多胞形。若 \(P\ne\varnothing\)，其顶点集有限且非空，并且
	\[
		P=\operatorname{conv}\operatorname{Vert}(P).
	\]
	这些结论允许 \(P\) 的维数小于 \(n\)。
\end{theorem}

\begin{proof}
	先从有界的半空间表示出发。空集单独按上述约定处理；若 \(P\) 是一个点，结论也已成立。一般把环境限制到仿射包络 \(L=\operatorname{aff}P\)，选取其中的原点与正交坐标，便可把 \(P\) 看成 \(\mathbb R^d\) 中的满维集合，其中 \(d=\dim P>0\)。原不等式限制到 \(L\) 后仍是有限个仿射不等式。以下写
	\[
		P=\{x:\langle a_i,x\rangle\le b_i,\ 1\le i\le m\}\subseteq\mathbb R^d.
	\]
	对 \(x\in P\)，把恰好取等号的指标组成的集合记为
	\[
		I(x)=\{i:\langle a_i,x\rangle=b_i\},
	\]
	称这些不等式在 \(x\) 处活动。令 \(r(x)\) 为活动法向量 \(a_i\)、\(i\in I(x)\) 的线性张成空间的维数。

	首先证明：\(x\) 是顶点，当且仅当 \(r(x)=d\)。若 \(r(x)<d\)，可取非零向量 \(u\) 与全部活动法向量正交。活动不等式在 \(x\pm\varepsilon u\) 处仍为等式，其他有限条不等式在 \(x\) 处严格成立，故当 \(\varepsilon>0\) 足够小时，它们在这两个点也成立。于是 \(x\) 是 \(P\) 中一条非退化线段的中点，不是顶点。反过来，若 \(r(x)=d\) 且 \(x=(1-t)y+tz\)、\(0<t<1\)，则每条在 \(x\) 处活动的不等式都必须在 \(y,z\) 处同时取等号：两个不超过 \(b_i\) 的数以正权重平均成 \(b_i\)，只能都等于 \(b_i\)。所以所有活动法向量都垂直于 \(y-x,z-x\)。它们张成全空间，迫使 \(y=z=x\)。

	每个顶点可由 \(d\) 条线性无关的活动等式唯一确定，而这样的指标组合只有有限个，故顶点总数有限。还需证明这些顶点确实生成 \(P\)，而不是只存在有限个候选点。若 \(x\) 尚非顶点，沿上面选取的方向 \(u\) 考察
	\[
		J=\{t\in\mathbb R:x+tu\in P\}.
	\]
	由闭性、凸性和有界性，\(J\) 是有限闭区间 \([\alpha,\beta]\)；前面的微小扰动又说明 \(\alpha<0<\beta\)。两个端点 \(x_-=x+\alpha u\)、\(x_+=x+\beta u\) 满足
	\[
		x=\frac{\beta}{\beta-\alpha}x_-+\frac{-\alpha}{\beta-\alpha}x_+.
	\]
	在每个端点都至少有一条原先不活动的不等式变成等式，否则还能稍微延长区间。新活动法向量不在原活动法向量的张成空间中：若在其中，它与 \(u\) 的配对为零，沿整条直线的值不变，就不可能从严格不等式变为等式。因此 \(r(x_-),r(x_+)>r(x)\)。对两个端点重复分解，活动法向量的秩每次增加，最多经过 \(d-r(x)\) 层便都到达顶点。把各层凸组合代回，即得 \(x\) 是有限个顶点的凸组合。这也证明非空有界多面体确实有顶点。

	现在从有限点集的凸包出发，设 \(P=\operatorname{conv}\{v_1,\ldots,v_N\}\)。它是标准单纯形在连续映射 \((\lambda_i)\mapsto\sum_i\lambda_iv_i\) 下的像，因此紧，尤其闭且有界。仍在其仿射包络中处理；零维情形已知。由\cref{g1:prop:relative-interior}，可选相对内点并把它移到原点，于是在 \(\mathbb R^d\) 中有
	\[
		rB\subseteq P\subseteq RB\qquad(0<r\le R),
	\]
	其中 \(B\) 是闭单位球。为从生成点造出有限个限制，先收集所有满足统一上界的线性测量：
	\[
		Q=\{y:\langle y,v_i\rangle\le1,\ 1\le i\le N\}.
	\]
	因为线性函数在凸组合上的值也是相应的凸组合，上式等价于 \(\langle y,z\rangle\le1\) 对一切 \(z\in P\) 成立。若 \(y\in Q\) 且 \(y\ne0\)，取 \(z=ry/\lVert y\rVert\in P\)，便得 \(r\lVert y\rVert\le1\)；故 \(Q\) 有界。另一方面，\(\lVert y\rVert\le1/R\) 时所有限制都成立，所以 \(Q\) 还包含原点的一个邻域。

	\(Q\) 已经是有界半空间交。应用刚刚证完的方向，可写 \(Q=\operatorname{conv}\{w_1,\ldots,w_M\}\)。我们断言
	\[
		P=\{x:\langle w_j,x\rangle\le1,\ 1\le j\le M\}.
	\]
	左边包含于右边由 \(w_j\in Q\) 立即得到。若 \(x\notin P\)，\cref{g1:thm:point-closed-separation}给出非零 \(u\)，使
	\[
		h=\max_{z\in P}\langle u,z\rangle<\langle u,x\rangle.
	\]
	最大值由紧性存在，而且 \(rB\subseteq P\) 保证 \(h\ge r\lVert u\rVert>0\)。因此 \(y=u/h\) 属于 \(Q\)，但 \(\langle y,x\rangle>1\)。若所有 \(\langle w_j,x\rangle\le1\)，把 \(y\) 写成 \(w_j\) 的凸组合便会得到相反结论。故至少一个 \(w_j\) 排除了 \(x\)，断言得证。

	最后恢复平移，并把仿射包络 \(L\) 的有限个线性方程各写成两条不等式，就得到原 \(\mathbb R^n\) 中的半空间表示。第一部分随即还说明 \(P\) 等于其全部顶点的凸包。
\end{proof}

证明中沿直线移到端点的步骤，正是“还没有被足够多独立限制固定，就继续移动”的几何含义。有界性保证直线与集合相交的区间有两个有限端点；若集合无界，这一步可能永远沿某个方向走下去。反向证明中的 \(Q\) 则把点 \(v_i\) 变成线性不等式的系数；\subsecref{g1:subsec:0303:04}将把这一构造正式称作极体。

\subsection{面、法锥与面格}
\label{g1:subsec:0303:03}

顶点只记住零维的边界部分。正方体还有边、四边形侧面以及它们的包含关系；若只用“某条不等式取等号”来定义这些对象，就可能把一个特定表示中的冗余数据误当成几何本身。下面先用线段给出不依赖表示的定义，再证明它对多胞形等价于支撑超平面的描述。

\begin{definition}[面与暴露面]
	\label{g1:def:polytope-face}
	设 \(P\subseteq\mathbb R^n\) 为非空多胞形。凸子集 \(F\subseteq P\) 称为 \(P\) 的\term[mian]{面}[face]，若一条端点在 \(P\) 中的线段有内点落在 \(F\)，则两个端点也都在 \(F\)。允许 \(F=\varnothing\) 和 \(F=P\)，其余称为非空真面。对 \(u\in\mathbb R^n\)，集合
	\[
		F_P(u)=\{x\in P:\langle u,x\rangle=h_P(u)\},\qquad h_P(u)=\max_{z\in P}\langle u,z\rangle,
	\]
	称为一个\term[baolumian]{暴露面}[exposed face]；允许 \(u=0\)，此时 \(F_P(0)=P\)。若 \(d=\dim P\)，其 \(d-1\) 维面称为\term[lengmian]{棱面}[facet]。零维面恰为单个顶点，一维面称为棱。
\end{definition}

\symidx{\(F_P(u)\)：线性函数 \(\langle u,-\rangle\) 在 \(P\) 上的最大值集}
这里的 \(h_P\) 就是\secref{g1:sec:0301}中定义的支撑函数。暴露面的线段性质来自一个简单事实：两个不大于最大值的数以正权重平均成最大值，二者都必须已是最大值。因此每个暴露面都是面。对于一般紧凸集，反向结论会失败，见\cref{g1:exr:nonexposed-point}；有限生成性正是下面证明的有效条件。

\begin{theorem}[多胞形的面]
	\label{g1:thm:polytope-faces}
	设 \(P\subseteq\mathbb R^n\) 是非空多胞形，顶点集为 \(V\)。则对每个面 \(F\)，有
	\[
		F=\operatorname{conv}(V\cap F).
	\]
	特别地，每个面仍是多胞形，面的总数有限。每个非空面都是暴露面；空面单独按约定加入。若在 \(P\) 的仿射包络中给出半空间表示，则每个非空面都是若干活动等式在 \(P\) 中的共同解集。
\end{theorem}

\begin{proof}
	若 \(x\in F\)，由\cref{g1:thm:polytope-hv}，可把 \(x\) 写成顶点的凸组合。每个系数为正的顶点都必须属于 \(F\)：若其系数小于一，先把这一项与其余项分开，便由面的线段性质得到结论；若系数为一，则该顶点就是 \(x\)。所以 \(F\subseteq\operatorname{conv}(V\cap F)\)，反向包含来自 \(F\) 的凸性。每个面由 \(V\) 的某个子集确定，而有限集只有有限个子集，故面的总数有限。

	现设 \(F\ne\varnothing\)，把其中的全部顶点写成 \(v_1,\ldots,v_s\)，取它们的平均点 \(x=s^{-1}\sum_jv_j\)。在仿射包络的坐标中写 \(P=\{z:\langle a_i,z\rangle\le b_i\}\)，仍令 \(I(x)\) 表示在 \(x\) 处取等号的指标，并置
	\[
		G=\{z\in P:\langle a_i,z\rangle=b_i\text{ 对每个 }i\in I(x)\}.
	\]
	对 \(i\in I(x)\)，所有 \(v_j\) 的函数值平均为 \(b_i\)，每项又不超过 \(b_i\)，故每项都等于 \(b_i\)。这说明 \(F\subseteq G\)。反过来，若 \(y\in G\)，取 \(z=x+\varepsilon(x-y)\)。活动等式在 \(z\) 处仍成立，其余有限条不等式在 \(x\) 处有严格余量，因此足够小的 \(\varepsilon>0\) 保证 \(z\in P\)。现在
	\[
		x=\frac{\varepsilon}{1+\varepsilon}y+\frac1{1+\varepsilon}z\in F,
	\]
	面的线段性质给出 \(y\in F\)。于是 \(F=G\)。

	令 \(u=\sum_{i\in I(x)}a_i\)、\(c=\sum_{i\in I(x)}b_i\)。对 \(z\in P\)，每个数 \(b_i-\langle a_i,z\rangle\) 都非负，其和为零当且仅当它们全为零。所以
	\[
		F=\{z\in P:\langle u,z\rangle=c\}=F_P(u).
	\]
	若 \(I(x)\) 为空，等式条件为空，得到 \(F=P\) 和 \(u=0\)。在低维情况下，把仿射包络中的线性函数延拓到环境空间，最大值集保持不变，故同样得到环境空间中的暴露面。
\end{proof}

由证明还可看出，\(F\) 的顶点正是 \(P\) 中落在 \(F\) 的顶点。一个方向是环境中的极点在子集中仍为极点；反方向若 \(v\) 是 \(F\) 的极点而在 \(P\) 中可以非平凡分解，面的线段性质会迫使这次分解完全发生在 \(F\) 内，产生矛盾。因而“先取面，再找顶点”和“先找顶点，再保留属于面的部分”给出同一份数据。

\begin{definition}[面上的外法锥]
	\label{g1:def:normal-cone}
	设 \(F\) 是非空多胞形 \(P\subseteq\mathbb R^n\) 的非空面。它的\term[fazhui]{法锥}[normal cone]定义为
	\[
		N_P(F)=\{u\in\mathbb R^n:\langle u,z-x\rangle\le0\text{ 对所有 }z\in P,\ x\in F\}.
	\]
	这是外法向约定，等价于 \(F\subseteq F_P(u)\)。它由闭半空间相交得到，对非负线性组合封闭，因而是闭凸锥；这里“锥”只表示乘以非负数后仍在集合内。
\end{definition}

\symidx{\(N_P(F)\)：面 \(F\) 的外法锥}
法锥收集的是能够在整个 \(F\) 上达到最大值的方向，而不要求最大值集恰好等于 \(F\)。这一点可以从正方形看清：对 \(P=[-1,1]^2\)，右边 \(E=\{1\}\times[-1,1]\) 的法锥为 \(\{(a,0):a\ge0\}\)，右上顶点 \(v=(1,1)\) 的法锥为 \(\{(a,b):a,b\ge0\}\)。若 \(a,b>0\)，最大值只在 \(v\) 取得；若 \(a>0,b=0\)，整条右边都达到最大值。一个较小的面允许更多方向同时在它上面取最大值。

\begin{proposition}[法锥反转面包含]
	\label{g1:prop:normal-cone-order}
	若 \(F,G\) 是非空多胞形 \(P\subseteq\mathbb R^n\) 的两个非空面，则
	\[
		F\subseteq G\quad\Longleftrightarrow\quad N_P(G)\subseteq N_P(F).
	\]
\end{proposition}

\begin{proof}
	若 \(F\subseteq G\)，在 \(G\) 的每一点达到最大值，当然也在 \(F\) 的每一点达到最大值，故得到正向包含。反过来，由\cref{g1:thm:polytope-faces}选取 \(u\) 使 \(G=F_P(u)\)。此时 \(u\in N_P(G)\subseteq N_P(F)\)，所以 \(F\subseteq F_P(u)=G\)。
\end{proof}

对于低维 \(P\)，垂直于其仿射包络方向空间的所有向量都属于每个法锥，因为这些线性函数在 \(P\) 上恒定。若只想保留 \(P\) 内部的几何方向，可以把法锥投影到该方向空间；本节需要环境法锥时，始终采用上述定义。

\begin{definition}[面格]
	\label{g1:def:face-lattice}
	多胞形 \(P\) 的全部面连同空面，按包含关系排序，称为它的\term[miange]{面格}[face lattice]，记作 \(\mathcal F(P)\)。
\end{definition}

\symidx{\(\mathcal F(P)\)：多胞形的面格}
“格”在这里表示任意两个面都有最大公共下面与最小公共上面。前者就是交集 \(F\cap G\)：交集仍凸，且每个面的线段性质同时保留下来。后者是所有包含 \(F\cup G\) 的面的交，记为 \(F\vee G\)；这样的面至少有 \(P\)，交集仍是面，故确实是所求最小上面。有限性由\cref{g1:thm:polytope-faces}保证。注意 \(F\vee G\) 一般不等于 \(\operatorname{conv}(F\cup G)\)：正方形的两个相对顶点的凸包是对角线，而包含两者的最小面是整个正方形。因为这样的面包含中心，把中心写成任意一对关于中心对称的点的中点，面的线段性质就迫使正方形内的全部点都属于它。

\subsection{极体及包含关系的反转}
\label{g1:subsec:0303:04}

在两种表示的等价证明中，每个点 \(x\) 产生了一个条件 \(\langle y,x\rangle\le1\)。当 \(x\) 遍历原来的集合时，满足全部条件的 \(y\) 又组成一个凸集。这个构造把“有哪些点”改写成“允许哪些线性测量”，也解释了为什么原来的点越多，得到的集合反而越小。

\begin{definition}[极集与极体]
	\label{g1:def:polar-body}
	对集合 \(K\subseteq\mathbb R^n\)，它相对于原点和标准内积的\term[jiji]{极集}[polar set]定义为
	\[
		K^\circ=\{y\in\mathbb R^n:\langle y,x\rangle\le1\text{ 对所有 }x\in K\}.
	\]
	若 \(K\) 是紧凸集且原点在其内部，通常称 \(K^\circ\) 为 \(K\) 的\term[jiti]{极体}[polar body]。
\end{definition}

\symidx{\(K^\circ\)：集合 \(K\) 相对于原点的极集}
不论 \(K\) 是否闭或凸，\(K^\circ\) 都是含原点的闭凸集，因为它由形如 \(\langle y,x\rangle\le1\) 的闭半空间相交得到。并且直接由量词可知
\begin{equation}
	\label{g1:eq:polar-inclusion}
	K\subseteq L\quad\Longrightarrow\quad L^\circ\subseteq K^\circ.
\end{equation}
若 \(K=\operatorname{conv}\{v_i\}\)，只要对所有 \(v_i\) 检验不等式，就等于对全部 \(K\) 检验。因此顶点表示立即变成极集的半空间表示。

\begin{theorem}[双极与有界性]
	\label{g1:thm:bipolar}
	设 \(K\subseteq\mathbb R^n\) 为含原点的非空闭凸集，则
	\[
		(K^\circ)^\circ=K.
	\]
	若还存在 \(0<r\le R\) 使 \(rB\subseteq K\subseteq RB\)，则
	\[
		\frac1R B\subseteq K^\circ\subseteq\frac1r B.
	\]
	特别地，含原点为内点的多胞形的极体仍是含原点为内点的多胞形。
\end{theorem}

\begin{proof}
	对 \(x\in K\)、\(y\in K^\circ\)，定义给出 \(\langle x,y\rangle\le1\)，故 \(K\subseteq(K^\circ)^\circ\)。若 \(x\notin K\)，由\cref{g1:thm:point-closed-separation}，存在 \(u\ne0\) 和实数 \(a\)，使
	\[
		\langle u,z\rangle\le a<\langle u,x\rangle\qquad(z\in K).
	\]
	因为 \(0\in K\)，必有 \(a\ge0\)。若 \(a>0\)，取 \(y=u/a\)，就有 \(y\in K^\circ\) 且 \(\langle y,x\rangle>1\)。若 \(a=0\)，则 \(\langle u,x\rangle>0\)，取足够大的 \(t>0\)，使 \(t\langle u,x\rangle>1\)；此时 \(y=tu\) 仍满足对 \(K\) 的全部配对不超过零，也属于 \(K^\circ\)。两种情形都证明 \(x\notin(K^\circ)^\circ\)，从而得到双极公式。

	若 \(\lVert y\rVert\le1/R\)，对 \(x\in K\) 由 Cauchy--Schwarz 不等式有 \(\langle y,x\rangle\le\lVert y\rVert\lVert x\rVert\le1\)，故 \(y\in K^\circ\)。反过来，若 \(y\in K^\circ\) 且 \(y\ne0\)，把 \(ry/\lVert y\rVert\in K\) 代入定义，得到 \(r\lVert y\rVert\le1\)。这就证明两个球的夹逼。

	若 \(K\) 是多胞形，其有限顶点给出 \(K^\circ\) 的有限半空间表示；球的夹逼又保证它有界并含原点为内点。\cref{g1:thm:polytope-hv}遂说明它也是多胞形。
\end{proof}

原点条件有两个不同作用，不宜混在一起。双极公式只要求 \(0\in K\)，无需原点是内点；极体保持有界却需要原点周围有足够多方向。例如 \([0,1]^\circ=(-\infty,1]\)，再次取极集仍得到 \([0,1]\)，但第一次取极已变成无界集合。若连 \(0\in K\) 都删除，\([1,2]^\circ=(-\infty,1/2]\)，其双极为 \([0,2]\)，原来的集合已经不能恢复。低维集合也有同样的有界性问题：\([-1,1]\times\{0\}\) 在平面中的极集是条带 \(\{|y_1|\le1\}\)，沿 \(y_2\) 方向无界；若在其一维仿射包络内取极，则得到有界区间。因而取极之前必须说明环境空间和原点。

\begin{example}[正方形与菱形]
	\label{g1:exm:square-diamond-polar}
	令 \(C=[-1,1]^2\)。给定 \(y=(y_1,y_2)\)，分别选择 \(x_i\) 与 \(y_i\) 同号，就得到
	\[
		h_C(y)=\max_{x\in C}(y_1x_1+y_2x_2)=|y_1|+|y_2|.
	\]
	因此
	\[
		C^\circ=\{y:|y_1|+|y_2|\le1\}=\operatorname{conv}\{\pm e_1,\pm e_2\}.
	\]
	右边 \(x_1=1\) 上所有点的共同等式 \(\langle y,x\rangle=1\) 迫使 \(y=(1,0)\)，所以一条边对应极体的一个顶点。顶点 \((1,1)\) 则给出极体中的一条边 \(y_1+y_2=1\)。\cref{g1:fig:square-polar}把这两种对应同时画出。
\end{example}

\begin{figure}[htbp]
	\centering
	\input{Figures/Chapter03-Polar}
	\caption[正方形与极体的面对应]{正方形与极体的面对应。蓝边 \(E\) 对应蓝顶点 \(e\)；红顶点 \(v\) 对应红边 \(D\)。包含关系 \(v\in E\) 变为 \(e\in D\)。两幅图使用各自的坐标平面。}
	\label{g1:fig:square-polar}
\end{figure}

这种对应不仅存在于平面图形。原点位于内部时，每个支撑值都可以归一化成一，从而把一个面的外法向方向截成极体上的面。

\begin{proposition}[极体的面对应]
	\label{g1:prop:polar-face-duality}
	设 \(P\subseteq\mathbb R^d\) 为满维多胞形，且 \(0\in\operatorname{int}P\)。对非空真面 \(F\)，定义
	\[
		F^*=\{y\in P^\circ:\langle y,x\rangle=1\text{ 对所有 }x\in F\}.
	\]
	它是 \(P^\circ\) 的非空真面，且 \((F^*)^*=F\)。连同约定 \(P^*=\varnothing\)、\(\varnothing^*=P^\circ\)，这给出两个面格之间反转包含关系的双射。对非空真面还有
	\[
		\dim F+\dim F^*=d-1.
	\]
	特别地，顶点对应棱面，棱面对应顶点。
\end{proposition}

\begin{proof}
	选取线性函数 \(\langle u,-\rangle\) 暴露出 \(F\)，即 \(F=F_P(u)\)。因为 \(F\) 为真面，\(u\ne0\)；又因原点在内部，\(h_P(u)>0\)。于是 \(y_0=u/h_P(u)\) 属于 \(P^\circ\)，在 \(F\) 上配对为一，在其他顶点上配对严格小于一。这首先证明 \(F^*\ne\varnothing\)。取 \(F\) 全部顶点的平均点 \(x_0\)。对 \(y\in P^\circ\)，\(\langle y,x_0\rangle=1\) 当且仅当在 \(F\) 的每个顶点都配对为一，也就当且仅当在整个 \(F\) 上配对为一。因此
	\[
		F^*=\{y\in P^\circ:\langle y,x_0\rangle=1\}
	\]
	是一个暴露面。它不含 \(0\)，所以是一个真面。

	由双极公式，\((P^\circ)^\circ=P\)。按定义，\(F\subseteq(F^*)^*\)；若 \(x\in P\setminus F\)，刚才的 \(y_0\) 满足 \(\langle y_0,x\rangle<1\)，又有 \(y_0\in F^*\)，故 \(x\notin(F^*)^*\)。于是 \((F^*)^*=F\)。包含关系的反转直接来自条件“对所有 \(x\in F\)”；再对反向映射应用同一结论，得到反转包含的双射。空面与全体的约定也与包含关系相容。

	最后核对维数。设 \(k=\dim F\)，\(W\) 为 \(F\) 的仿射包络的方向空间，则 \(\dim W=k\)。因 \(F\) 位于支撑超平面 \(\langle y_0,x\rangle=1\) 中，\(x_0\notin W\)。所有对 \(F\) 配对为一的向量组成仿射空间
	\[
		Y=\{y:\langle y,x_0\rangle=1,\ \langle y,w\rangle=0\ (w\in W)\}.
	\]
	这些是 \(k+1\) 个独立线性条件，所以 \(\dim Y=d-k-1\)。在 \(y_0\) 附近、沿 \(Y\) 作足够小的扰动，对 \(F\) 的顶点仍配对为一；对有限个不在 \(F\) 中的顶点，由严格余量，配对仍小于一。因此 \(F^*\) 包含 \(y_0\) 在 \(Y\) 中的一个相对邻域，又完全位于 \(Y\) 内，故 \(\dim F^*=d-k-1\)。
\end{proof}

法锥与极体的面现在也可以直接比较：对非空真面 \(F\) 及 \(x_0\in F\)，有
\[
	F^*=N_P(F)\cap\{y:\langle y,x_0\rangle=1\},\qquad N_P(F)=\{ty:t\ge0,\ y\in F^*\}.
\]
第一个等式把最大值归一化；第二个等式使用了非零外法向量的支撑值严格为正。因而一个面先决定一整束外法向方向，再由极体从每条射线上选取一个代表。原点改变时，归一化的位置也会改变，极体不是单纯由图形的仿射外观决定的。

\subsection{单纯形、立方体和交叉多面体的对照}
\label{g1:subsec:0303:05}

取 \(d\ge2\)。记标准单纯形为 \(\Delta_d=\operatorname{conv}\{0,e_1,\ldots,e_d\}\)，将其重心移到原点后记为 \(S_d\)；另记
\[
	C_d=[-1,1]^d,\qquad O_d=\operatorname{conv}\{\pm e_1,\ldots,\pm e_d\}.
\]
其中 \(\Delta_d\) 是\secref{g1:sec:0101}已经定义的单纯形：它由 \(d+1\) 个仿射无关点生成，维数为 \(d\)。\(C_d\) 是 \(d\) 维立方体，\(O_d\) 称为\term[jiaochaduomianti]{交叉多面体}[cross-polytope]。平移单纯形只是为了满足取有界极体所需的原点条件，不改变其面及包含关系。

\begin{proposition}[三个模型的面数与对偶]
	\label{g1:prop:basic-polytope-counts}
	设 \(d\ge2\)，并以 \(f_k(P)\) 表示 \(P\) 的 \(k\) 维面的数目。则对 \(0\le k\le d-1\)，\cref{g1:tab:basic-polytopes}中的公式成立。其中 \(C_d^\circ=O_d\)、\(O_d^\circ=C_d\)，而 \(S_d^\circ\) 也是一个 \(d\) 维单纯形。
\end{proposition}

\begin{table}[htbp]
	\centering
	\caption[三个基本多胞形的计数与对偶]{三个基本多胞形的计数与对偶。\(0\le k\le d-1\)；每个多胞形另有一个空面和一个全体面。单纯形的极体记为 \(S'_d\)，不表示它与原单纯形逐点相同。}
	\label{g1:tab:basic-polytopes}
	\begin{tabular}{@{}lcccc@{}}
		\toprule
		对象 & 顶点数 & 棱面数 & \(f_k\) & 极体 \\
		\midrule
		单纯形 \(S_d\) & \(d+1\) & \(d+1\) & \(\binom{d+1}{k+1}\) & \(S'_d\) \\
		立方体 \(C_d\) & \(2^d\) & \(2d\) & \(2^{d-k}\binom dk\) & \(O_d\) \\
		交叉多面体 \(O_d\) & \(2d\) & \(2^d\) & \(2^{k+1}\binom d{k+1}\) & \(C_d\) \\
		\bottomrule
	\end{tabular}
\end{table}

\begin{proof}
	单纯形的每个非空面由一部分顶点生成；反过来，任意顶点子集都生成一个面。为说明后一点，使用单纯形上的重心坐标 \(\lambda_0,\ldots,\lambda_d\)：它们是仿射函数，非负且和为一。给定非空真子集 \(I\)，仿射函数 \(\sum_{i\in I}\lambda_i\) 的最大值为一，恰好在其余重心坐标全为零时取得，也就是选定顶点的凸包。仿射函数与其线性部分的最大值集相同，所以这确是一个暴露面。选出 \(k+1\) 个顶点便得到一个 \(k\) 维面，数目为 \(\binom{d+1}{k+1}\)。

	对于立方体，线性函数 \(\sum_i a_ix_i\) 取最大值时，\(a_i>0\) 迫使 \(x_i=1\)，\(a_i<0\) 迫使 \(x_i=-1\)，只有 \(a_i=0\) 的坐标可以自由变化。故 \(k\) 维面由 \(k\) 个自由坐标和其余坐标的 \(\pm1\) 选择确定，共有 \(\binom dk2^{d-k}\) 个。最大值本身是 \(\sum_i|a_i|\)，从而
	\[
		C_d^\circ=\{y:\textstyle\sum_i|y_i|\le1\}=O_d.
	\]
	最后一个等号可以直接核验：\(\sum_i|y_i|\le1\) 时，以 \(|y_i|\) 为系数使用相应符号的 \(e_i\)，剩余权重放在 \(0\)；而 \(0=(e_1+(-e_1))/2\) 本已在这些顶点的凸包中。双极公式给出 \(O_d^\circ=C_d\)。

	对于交叉多面体，非零线性函数 \(\langle a,-\rangle\) 的最大值为 \(M=\max_i|a_i|>0\)，达到最大值的顶点正是 \(\operatorname{sgn}(a_i)e_i\)、\(|a_i|=M\)。它们每个坐标方向至多选一个，且仿射无关。要得到 \(k\) 维面，应选 \(k+1\) 个坐标方向，每个方向任选正负号，共有 \(\binom d{k+1}2^{k+1}\) 种；每种确实能实现，只需让所选坐标的 \(a_i\) 等于相应的 \(\pm1\)，其余为零。

	还需确认平移单纯形的极体也是单纯形，而不只核对面数。置 \(b=(1/(d+1),\ldots,1/(d+1))\)，则 \(S_d=\Delta_d-b\)，故
	\[
		S_d=\{x:-(d+1)x_i\le1\ (1\le i\le d),\ (d+1)\textstyle\sum_i x_i\le1\}.
	\]
	记 \(q_i=-(d+1)e_i\) 和 \(q_0=(d+1)(1,\ldots,1)\)。这些 \(d+1\) 个点仿射无关，且平均为零：若 \(\sum_{i=1}^d c_i(q_i-q_0)=0\)，逐坐标得到 \(c_j=-\sum_i c_i\)，求和即迫使全部系数为零。因此 \(Q=\operatorname{conv}\{q_0,\ldots,q_d\}\) 是含原点的单纯形；上式恰好说明 \(Q^\circ=S_d\)。对闭凸含原点的 \(Q\) 应用双极公式，便有 \(S_d^\circ=Q\)。
\end{proof}

例如三维立方体有 \(8\) 个顶点、\(12\) 条棱和 \(6\) 个棱面，而三维交叉多面体即正八面体，分别有 \(6,12,8\) 个相应的面。两份数目倒序对应，反映的是 \(\dim F+\dim F^*=2\)，不是偶然的数字对称。在一维，单纯形、立方体和交叉多面体都只是线段；在零维，非空多胞形只有一个点，需要在零维环境中单独理解，不能把表中的棱面公式机械延伸过去。

有限表示、面格与极性还可以继续用于研究多胞形的组合结构。按这些主题延伸阅读，可查阅 \cite[Sections 1.1, 2.1--2.4]{Ziegler1995}。

% !TeX root = ../../Book1.tex
% 纲要标题：### 3.4 凸方法在后续几何中的位置
% 纲要位置：几何学计划纲要.md，第 226 行。
% 正文状态：初稿；尚未正式审核。

% 纲要内容开始
% 
% 把支撑函数、对偶锥和格多面体分别接到第三卷距离比较、第四卷矩映射以及第五卷环面簇；这里不调用后续分类定理。
% 
% ---
% 
% 纲要内容结束

\section{凸方法在后续几何中的位置}
\label{g1:sec:0304}

一个凸集可以由它包含的点来描述，也可以由线性函数在它上面的取值界来描述。\cref{g1:prop:support-function-representation}把这两种描述联系起来：只要 \(C\subseteq\mathbb R^n\) 非空、闭且凸，全部不等式 \(\langle u,x\rangle\le h_C(u)\) 就恰好刻画 \(C\)。这一事实不仅能回答“点是否属于集合”，还可以测量点离集合多远。以下仍使用标准欧氏内积与范数，先把这个测量写成一个精确等式。

\begin{proposition}[距离的支撑函数表示]\label{g1:prop:distance-support-formula}
	设 \(C\subseteq\mathbb R^n\) 是非空闭凸集，\(h_C(u)=\sup_{y\in C}\langle u,y\rangle\)。对每个 \(x\in\mathbb R^n\)，有
	\begin{equation}\label{g1:eq:distance-support-formula}
		d(x,C)=\max_{\lVert u\rVert\le1}\bigl(\langle u,x\rangle-h_C(u)\bigr).
	\end{equation}
	当 \(h_C(u)=+\infty\) 时，括号中的值约定为 \(-\infty\)。若 \(x\notin C\)，令 \(p=P_Cx\) 为最近点，则 \(u=(x-p)/\lVert x-p\rVert\) 取得最大值；若 \(x\in C\)，可取 \(u=0\)。
\end{proposition}
\begin{proof}
	对任意 \(y\in C\) 及 \(\lVert u\rVert\le1\)，Cauchy--Schwarz 不等式给出
	\[
		\langle u,x\rangle-h_C(u)\le\langle u,x-y\rangle\le\lVert x-y\rVert.
	\]
	对 \(y\) 取下确界，便知右边的每个候选值都不超过 \(d(x,C)\)。若 \(x\in C\)，取 \(u=0\) 得到 \(0\)，所以等号成立。

	若 \(x\notin C\)，\cref{g1:thm:convex-projection}保证最近点 \(p\) 存在且唯一，并且对所有 \(y\in C\) 有 \(\langle x-p,y-p\rangle\le0\)。令 \(u=(x-p)/\lVert x-p\rVert\)，这条不等式说明 \(h_C(u)=\langle u,p\rangle\)：右侧是所有取值的上界，且在 \(p\in C\) 处达到。因此
	\[
		\langle u,x\rangle-h_C(u)=\langle u,x-p\rangle=\lVert x-p\rVert=d(x,C).
	\]
	这样还直接找到了最大值的取得方向，不需要先假定支撑函数在整个单位球上有限或连续。
\end{proof}

对满足 \(h_C(u)<+\infty\) 的单位向量 \(u\)，超平面 \(\langle u,y\rangle=h_C(u)\) 给出沿该方向的界，\(\langle u,x\rangle-h_C(u)\) 是点 \(x\) 到这个超平面的有向距离。一般情况下上确界未必在 \(C\) 中达到，因此这个超平面未必接触集合；但上面由最近点选出的方向确实给出在 \(p\) 处接触的支撑超平面。对集合外的点，这个方向测得的距离恰好就是到集合的距离。特别地，\(d(\,\cdot\,,C)\) 是凸函数：把每个有限值方向对应的仿射函数在 \((1-t)x+ty\) 处的值展开，再取上确界，就得到
\[
	d((1-t)x+ty,C)\le(1-t)d(x,C)+t\,d(y,C)\qquad(0\le t\le1).
\]
这里取上确界的方向集合与 \(x,y\) 无关，因此每个方向上的仿射等式都受同一个右侧控制。

\begin{example}[球和半空间的距离]\label{g1:exm:support-distance-models}
	对闭球 \(C=\{y:\lVert y-c\rVert\le r\}\)，其中 \(r\ge0\)，有 \(h_C(u)=\langle u,c\rangle+r\lVert u\rVert\)。将它代入\eqref{g1:eq:distance-support-formula}，先沿 \(x-c\) 的方向取 \(u\)，再在 \(0\le\lVert u\rVert\le1\) 中优化，得到
	\[
		d(x,C)=\max\{0,\lVert x-c\rVert-r\}.
	\]
	若 \(x=c\)，直接取 \(u=0\) 即可。对半空间 \(H=\{y:\langle a,y\rangle\le b\}\)，\(a\ne0\)，沿单位法向 \(a/\lVert a\rVert\) 的超出量则给出
	\[
		d(x,H)=\frac{\max\{0,\langle a,x\rangle-b\}}{\lVert a\rVert}.
	\]
	第二式也可直接验证：当 \(x\notin H\) 时，从 \(x\) 减去 \((\langle a,x\rangle-b)a/\lVert a\rVert^2\)，就到达边界上的最近点。
\end{example}

这些计算说明，距离虽然由对所有点求下确界定义，却可以由比较函数从另一侧控制。比如在 \(x_0\notin C\) 处选取上面得到的单位向量 \(u_0\)，仿射函数
\[
	\ell(x)=\langle u_0,x\rangle-h_C(u_0)
\]
满足 \(\ell(x)\le d(x,C)\)，并在 \(x_0\) 处相等。这是一种从下方接触距离函数的比较。第三卷第7章第3、4节将处理曲空间中的距离与割迹，那里的“支撑函数”还指在一点接触目标函数的局部比较函数，不应直接等同于本章定义在方向空间上的 \(h_C\)。比较时究竟从上方还是下方接触、在哪个邻域光滑，以及曲率对二阶导数的限制，都需要另行说明；这里的欧氏公式提供的是一个可以完全算清的模型。

\ProseParagraphBreak

再看同时满足一组齐次不等式的方向。设 \(V\) 是有限维实向量空间。包含 \(0\) 且对非负线性组合封闭的集合 \(C\subseteq V\) 称为\term[tuzhui]{凸锥}[convex cone]；它的\term[duiouzhui]{对偶锥}[dual cone]定义在对偶空间中：
\[
	C^\vee=\{\,\varphi\in V^*:\varphi(v)\ge0\ \text{对所有 }v\in C\,\}.
\]
本节采用非负号约定。选定内积后可以用向量表示线性函数，但 \(v\) 与 \(\varphi\) 原本属于不同空间：一个是允许的方向，另一个是检测该方向的线性量。由定义，\(C^\vee\) 总是闭凸锥。它不是把极体定义中的常数 \(1\) 原样保留的另一种记号；对锥而言，\(\varphi(tv)\le1\) 对所有 \(t\ge0\) 成立，当且仅当 \(\varphi(v)\le0\)，所以前节的极集满足 \(C^\circ=-C^\vee\)。

\begin{proposition}[闭凸锥的双对偶]\label{g1:prop:dual-cone-bipolar}
	设 \(V\) 是有限维实向量空间，\(C\subseteq V\) 是闭凸锥。在自然识别 \(V^{**}=V\) 下，有 \(C^{\vee\vee}=C\)。若
	\[
		C=\left\{\,\sum_{j=1}^m t_jv_j:t_j\ge0\,\right\},
	\]
	则无论是否预先知道此锥闭，都有
	\[
		C^\vee=\{\,\varphi:\varphi(v_j)\ge0,\ j=1,\ldots,m\,\}.
	\]
\end{proposition}
\begin{proof}
	第一式的包含 \(C\subseteq C^{\vee\vee}\) 直接来自定义。若 \(x\notin C\)，由\cref{g1:thm:point-closed-separation}，存在非零线性函数 \(u\)，使 \(u(x)>\sup_{y\in C}u(y)\)，并且右侧有限。因为 \(0\in C\)，该上确界不小于 \(0\)；又因 \(ty\in C\) 对所有 \(t\ge0\) 成立，若某个 \(u(y)>0\)，上确界就会是 \(+\infty\)。故 \(u(y)\le0\) 对所有 \(y\in C\) 成立，且上确界等于 \(0\)。于是 \(-u\in C^\vee\)，而 \((-u)(x)<0\)，排除了 \(x\in C^{\vee\vee}\) 的可能。

	对第二式，\(\varphi\in C^\vee\) 显然在每个生成向量上非负；反之，若这些有限个值都非负，则 \(\varphi(\sum_jt_jv_j)=\sum_jt_j\varphi(v_j)\ge0\)。这正是从全部方向的检验缩减为有限次检验的理由。
\end{proof}

这个定义也解释了为何对偶锥出现在矩映射中。下面的计算只用复数、偏导数和群作用。取整数向量 \(\alpha_1,\ldots,\alpha_N\in\mathbb Z^d\)，记 \(T^d=(\mathbb R/2\pi\mathbb Z)^d\)。在 \(\mathbb C^N\) 上定义作用
\[
	t\cdot(z_1,\ldots,z_N)=
	\bigl(e^{-i\langle\alpha_1,t\rangle}z_1,\ldots,e^{-i\langle\alpha_N,t\rangle}z_N\bigr).
\]
整数条件保证把 \(t\) 的任一坐标加上 \(2\pi\) 不改变作用。定义
\begin{equation}\label{g1:eq:linear-moment-map}
	\mu(z)=\frac12\sum_{j=1}^N|z_j|^2\alpha_j.
\end{equation}
这是此作用的一个\term[juyingshe]{矩映射}[moment map]。要理解名称背后的要求，可把 \(z_j=x_j+iy_j\) 写成实坐标。在两个实切向量 \((a,b)\)、\((c,e)\) 上定义交替配对
\[
	\omega_0((a,b),(c,e))=\sum_j(a_je_j-b_jc_j).
\]
它就是第四卷要使用的标准辛形式 \(\sum_j\mathrm{d}x_j\wedge \mathrm{d}y_j\)，当前只需把它看成明确给定的双线性型。对方向 \(\xi\in\mathbb R^d\)，令 \(\beta_j=\langle\alpha_j,\xi\rangle\)。沿一参数作用 \(s\xi\) 求导，得到速度向量
\[
	X_\xi(z)=\bigl((\beta_jy_j,-\beta_jx_j)\bigr)_{j=1}^N.
\]
而函数 \(H_\xi(z)=\langle\mu(z),\xi\rangle=\frac12\sum_j\beta_j(x_j^2+y_j^2)\) 在方向 \(Y=(c,e)\) 上的导数满足
\[
	\mathrm{d}H_\xi(Y)=\sum_j\beta_j(x_jc_j+y_je_j)=\omega_0(X_\xi,Y).
\]
此外，\(\mu\) 在上述作用下不变，因为每个 \(|z_j|\) 不变。这两个条件是环面作用的矩映射条件在本例中的具体形式。负号放在作用的指数中，是为了与本书 \(\iota_X\omega=\mathrm{d}H\) 的约定相容；若保持 \(\omega_0\) 不变而把作用改成正向旋转，相应的 \(\mu\) 就要变号。一般定义与这些线性例子可参看\cite[Section 5.6]{Cannas2006Symplectic}。

\begin{example}[从权重得到锥与多胞形]\label{g1:exm:moment-cone-polytope}
	对\eqref{g1:eq:linear-moment-map}，非负系数 \(|z_j|^2/2\) 可以独立取遍所有非负实数，故
	\[
		\mu(\mathbb C^N)=\operatorname{cone}(\alpha_1,\ldots,\alpha_N).
	\]
	对任意 \(u\in\mathbb R^d\)，函数 \(\langle u,\mu(z)\rangle\) 处处非负，当且仅当 \(\langle u,\alpha_j\rangle\ge0\) 对每个 \(j\) 成立：必要性只需令其余坐标为零。因此对偶锥恰好记录了哪些线性组合的矩映射分量不会取负值。

	若限制在球面 \(S=\{z:\sum_j|z_j|^2=2\}\) 上，系数之和变为 \(1\)，从而
	\[
		\mu(S)=\operatorname{conv}\{\alpha_1,\ldots,\alpha_N\}.
	\]
	反向包含也有显式验证：给定 \(t_j\ge0\)、\(\sum_jt_j=1\)，取 \(z_j=\sqrt{2t_j}\) 即可。这里是在计算同一个映射限制到球面后的像，并未把这个奇数维球面当作辛流形。\cref{g1:thm:caratheodory}还说明，像中的每一点都可由至多 \(d+1\) 个非零坐标实现；它并不要求原来的每个原像都只有这么多个非零坐标。
\end{example}

第四卷第10章第1节将把这套定义推广到流形上的群作用，第12章第2、3节再讨论紧连通情形的凸性和从多面体构造几何对象的问题。线性模型中像的凸性直接来自可独立调节的非负系数；一般流形并没有这样一套全局坐标，因此需要新的整体论证。

\ProseParagraphBreak

把实向量空间中的整数点也保留下来，会出现另一类信息。称 \(\mathbb R^d\) 中由一组实基 \(e_1,\ldots,e_d\) 的整数线性组合构成的集合为\term[ge]{格}[lattice]，记为 \(M\)。以格点为顶点的多胞形称为\term[geduobaoxing]{格多胞形}[lattice polytope]，也称格多面体。选择格基后可写 \(M=\mathbb Z^d\)。保持格的线性坐标变换必须是整数矩阵，且逆矩阵仍为整数矩阵；由行列式乘法可知这等价于行列式为 \(1\) 或 \(-1\)。因此一般实可逆线性变换不再是允许的坐标变换。

对 \(m=(m_1,\ldots,m_d)\in M\)，在 \((\mathbb C^*)^d\) 上定义
\[
	\chi^m(z)=z_1^{m_1}\cdots z_d^{m_d}.
\]
所有坐标非零，所以负指数也有意义。这些函数满足 \(\chi^{m+m'}=\chi^m\chi^{m'}\)，把格的加法转成函数的乘法。若 \(P\) 是格多胞形，有限集合 \(A=P\cap M=\{m_0,\ldots,m_r\}\) 就给出映射
\begin{equation}\label{g1:eq:lattice-monomial-map}
	\Phi_A:(\mathbb C^*)^d\longrightarrow\mathbb P^r(\mathbb C),\qquad
	z\longmapsto[\chi^{m_0}(z):\cdots:\chi^{m_r}(z)].
\end{equation}
格点集有限，是因为有界集合中每个整数坐标只有有限种可能；各单项式均非零，故这确实定义了射影点。把全部格点平移同一个 \(a\in M\)，所有坐标都乘上共同的非零因子 \(\chi^a(z)\)，映射不变。若 \(m_i+m_j=m_k+m_l\)，像上便有齐次方程 \(Z_iZ_j=Z_kZ_l\)。这说明多面体中的整数关系如何转成几何方程。

\begin{example}[正方形给出的射影二次曲面]\label{g1:exm:lattice-square-segre}
	取 \(P=[0,1]^2\subseteq\mathbb R^2\)，按 \((0,0),(1,0),(0,1),(1,1)\) 排列格点，则
	\[
		\Phi_A(z,w)=[1:z:w:zw],\qquad Z_0Z_3=Z_1Z_2.
	\]
	不但像包含在这个二次曲面中，它的闭包还等于整个曲面。为证明这一点，把第二章的乘积坐标写成
	\[
		([s_0:s_1],[t_0:t_1])\longmapsto
		[s_0t_0:s_1t_0:s_0t_1:s_1t_1].
	\]
	由\cref{g1:prop:segre-rank-one}，其像恰为上述方程的非零齐次解：坐标排成的 \(2\times2\) 矩阵行列式为零，故可分解为两个非零向量的乘积。两个因子的全部坐标均非零时，分别除以 \(s_0,t_0\) 就得到 \(\Phi_A(z,w)\)。任意一对因子都可以用坐标全非零的因子逼近，因而曲面上的每一点都是这个像的极限；另一方面，齐次多项式的零点集是闭集，极限仍满足方程。这里的闭包结论已经用通常的复射影拓扑证明，不需要预先掌握代数簇理论。
\end{example}

同样，区间 \([0,2]\) 的三个格点给出 \(z\mapsto[1:z:z^2]\)，补齐端点后得到 \(Z_0Z_2=Z_1^2\)，正是第二章的二次射影参数曲线。区间 \([0,1]\) 与 \([0,2]\) 作为实凸集可以伸缩对应，内部格点数却不同，于是给出的射影坐标数和嵌入次数也不同。这就是第五卷第28章需要同时记录凸性与整数结构的原因：第1节由对偶格与锥构造仿射模型，第3节再把格点与线丛的截面联系起来。当前的单项式计算给出了可以直接观察的起点。



\begin{chaptersummary}
	凸包由有限凸组合生成，相对内部则在集合真正占据的仿射空间内衡量内部。对非空闭凸集，最近点的存在来自有限维紧性，唯一性来自距离平方沿线段的严格凸性；投影不等式进一步给出点与闭凸集的严格分离。边界支撑、开凸集的弱分离以及紧集与闭集之间的正间隔分离各有自己的假设。支撑函数允许取 \(+\infty\)，其全部有限不等式仍能完整恢复一个非空闭凸集。

	\ProseParagraphBreak

	有限维性的组合后果来自同一条仿射依赖式。调整已有权重并让一个系数降到零，得到 Carathéodory 定理；把正负系数分开归一化，得到 Radon 分割；再把分割用于“每次少一个条件”的公共点，得到有限族的 Helly 定理。单纯形说明这些数量界一般不能降低；无限族的相交还需要额外的紧性条件。

	\ProseParagraphBreak

	有界多面体恰是有限点集的凸包。活动约束决定顶点和面，线性函数的最大值记录暴露的边界，法锥记录哪些方向看见同一个面。原点在内部时，极体仍为满维多胞形，取极两次恢复原物，面对应反转包含关系并交换维数。单纯形、立方体与交叉多面体分别展示了重心坐标、独立坐标约束和绝对值约束所产生的结构。距离、矩映射和格点构造随后使用的，正是这些不等式、非负组合与对偶关系。
\end{chaptersummary}

\BookExercises
以下各题均附完整解答。分离题先核对闭性、凸性与紧性条件；多面体题同时保留几何维数和所在空间的维数，计算极体前先确认原点的位置。


\begin{exercise}\label{g1:ex:relative-triangle}
	在 \(\mathbb R^3\) 中令 \(p_0=(1,0,0)\)、\(p_1=(0,1,0)\)、\(p_2=(0,0,1)\)，并令 \(C=\operatorname{conv}\{p_0,p_1,p_2\}\)。求 \(\operatorname{aff}C\)、\(\operatorname{int}C\)、\(\operatorname{ri}C\) 与 \(\operatorname{rbd}C\)。再证明：若 \(f:F\rightarrow G\) 为有限维实仿射空间间的仿射同构，\(S\subseteq F\) 非空，则 \(f(\operatorname{ri}S)=\operatorname{ri}(f(S))\)，其中两边的相对内部都在相应仿射包络中计算。
\end{exercise}
\begin{solution}
	三点仿射无关，因为 \(p_1-p_0=(-1,1,0)\)、\(p_2-p_0=(-1,0,1)\) 线性无关。故
	\[
		\operatorname{aff}C=\{(x,y,z):x+y+z=1\},\qquad
		C=\{(x,y,z):x+y+z=1,\ x,y,z\geq0\}.
	\]
	它的普通内部为空，因为三维开球不可能包含于这个平面。若三个坐标都为正，在平面内充分小的邻域中它们仍为正，所以这些点属于相对内部。若某个坐标为零，可在平面内作任意小的扰动使该坐标为负，同时调整另一个坐标保持总和为 \(1\)，因此该点不属于相对内部。于是
	\[
		\operatorname{ri}C=\{(x,y,z):x+y+z=1,\ x,y,z>0\},
	\]
	而相对边界由 \(C\) 中至少一个坐标为零的点组成，即三条闭边的并。

	对于一般仿射同构 \(f\)，它与逆映射都保持仿射组合，故 \(f(\operatorname{aff}S)=\operatorname{aff}(f(S))\)。在有限维坐标中，\(f\) 及其逆映射都是线性映射加平移，因而连续；因此这一限制是两个仿射包络之间的同胚。若 \(x\in\operatorname{ri}S\)，它在 \(\operatorname{aff}S\) 中有一个包含于 \(S\) 的开邻域，其像便是 \(f(x)\) 在像仿射包络中包含于 \(f(S)\) 的开邻域。反过来对 \(f^{-1}\) 使用同一论证，得到相等。这个结论只用到仿射同构，未要求 \(S\) 凸。
\end{solution}

\begin{exercise}\label{g1:ex:halfspace-projection}
	设 \(a\in\mathbb R^n\setminus\{0\}\)、\(b\in\mathbb R\)，令 \(H=\{y:\langle a,y\rangle\leq b\}\)。证明
	\[
		P_Hx=x-\frac{\max\{\langle a,x\rangle-b,0\}}{\lVert a\rVert^2}a.
	\]
	再求到闭正象限 \(Q=\{y:y_i\geq0\ (1\leq i\leq n)\}\) 的投影和距离，并说明 \(P_Q\) 为什么一般不是线性映射。
\end{exercise}
\begin{solution}
	若 \(\langle a,x\rangle\leq b\)，则 \(x\in H\)，最近点就是 \(x\)，公式成立。若 \(\langle a,x\rangle>b\)，置 \(t=(\langle a,x\rangle-b)/\lVert a\rVert^2>0\)、\(p=x-ta\)。有 \(\langle a,p\rangle=b\)，故 \(p\in H\)。任取 \(y\in H\)，
	\[
		\langle x-p,y-p\rangle=t(\langle a,y\rangle-b)\leq0.
	\]
	由\cref{g1:thm:convex-projection}，\(p=P_Hx\)。这项计算也给出 \(d(x,H)=\max\{\langle a,x\rangle-b,0\}/\lVert a\rVert\)。

	对正象限，距离平方为 \(\sum_i(x_i-y_i)^2\)，每个坐标的约束互相独立。若 \(x_i\geq0\)，取 \(y_i=x_i\) 使该项为零；若 \(x_i<0\)，对 \(y_i\geq0\) 有 \((x_i-y_i)^2\geq x_i^2\)，且仅在 \(y_i=0\) 时取等号。因此
	\[
		(P_Qx)_i=\max\{x_i,0\},\qquad
		d(x,Q)=\left(\sum_{i:x_i<0}x_i^2\right)^{1/2}.
	\]
	例如 \(P_Q(e_1)=e_1\) 而 \(P_Q(-e_1)=0\)，不满足线性映射必须具有的 \(P_Q(-e_1)=-P_Q(e_1)\)。
\end{solution}

\begin{exercise}\label{g1:ex:nonexpansive-projection}
	设 \(C\subseteq\mathbb R^n\) 非空、闭且凸。证明对所有 \(x,y\in\mathbb R^n\)，
	\[
		\lVert P_Cx-P_Cy\rVert^2\leq\langle P_Cx-P_Cy,x-y\rangle,
		\qquad\lVert P_Cx-P_Cy\rVert\leq\lVert x-y\rVert.
	\]
	再证明，对任意非空集合 \(S\subseteq\mathbb R^n\)，距离函数都满足 \(\lvert d(x,S)-d(y,S)\rvert\leq\lVert x-y\rVert\)；这一步不要求 \(S\) 闭或凸。
\end{exercise}
\begin{solution}
	置 \(p=P_Cx\)、\(q=P_Cy\)。对两次投影分别使用式\eqref{g1:eq:convex-projection-variation}，得到 \(\langle x-p,q-p\rangle\leq0\) 与 \(\langle y-q,p-q\rangle\leq0\)。将第一式反号后与第二式比较，有
	\[
		0\leq\langle x-p,p-q\rangle-\langle y-q,p-q\rangle
		=\langle x-y,p-q\rangle-\lVert p-q\rVert^2,
	\]
	即第一项结论。再用内积不等式，\(\lVert p-q\rVert^2\leq\lVert p-q\rVert\cdot\lVert x-y\rVert\)。当 \(p=q\) 时第二项结论直接成立；否则除以正数 \(\lVert p-q\rVert\) 即可。这说明输入点略有移动时，最近点的移动不会更大。

	对于一般非空 \(S\)，任给 \(\varepsilon>0\)，可取 \(s\in S\) 使 \(\lVert y-s\rVert<d(y,S)+\varepsilon\)。三角不等式给出
	\[
		d(x,S)\leq\lVert x-s\rVert\leq\lVert x-y\rVert+\lVert y-s\rVert
		<\lVert x-y\rVert+d(y,S)+\varepsilon.
	\]
	令 \(\varepsilon\to0^+\)，再交换 \(x,y\)，得到绝对值不等式。选择近似最近点而不是要求真正最近点，正是能够去掉闭性假设的原因。
\end{solution}

\begin{exercise}\label{g1:ex:nonclosed-convex-image}
	取\cref{g1:exm:zero-distance-separation}中的闭凸集 \(B=\{(x,y):x>0,\ y\geq1/x\}\)，令 \(\pi:\mathbb R^2\rightarrow\mathbb R\) 为 \(\pi(x,y)=x\)。求 \(\pi(B)\) 及其闭包的支撑函数，并说明这给“线性映射保持闭凸集”与“支撑函数确定所有凸集”各提供了什么反例。
\end{exercise}
\begin{solution}
	每个 \((x,y)\in B\) 都满足 \(x>0\)；反过来，每个 \(x>0\) 都有 \((x,1/x)\in B\)。因此 \(\pi(B)=(0,\infty)\)，是凸集却不是闭集，它的闭包为 \([0,\infty)\)。这说明线性映射虽然保持凸性，但不一定保持闭性，甚至正交投影到坐标轴也不能保证闭性。

	把 \(u\in\mathbb R\) 看作一维方向，直接取上确界得
	\[
		h_{(0,\infty)}(u)=h_{[0,\infty)}(u)
		=\begin{cases}0,&u\leq0,\\+\infty,&u>0.\end{cases}
	\]
	当 \(u<0\) 时，第一个集合只在 \(x\to0^+\) 时趋近上确界，第二个集合则在 \(x=0\) 取得上确界；支撑函数记录的都是同一个数 \(0\)。所以它并不记录某个边界点是否被删除，无法仅凭这项函数区分两个不同的凸集。闭凸集表示定理中的闭性不能省去。
\end{solution}

\begin{exercise}\label{g1:ex:dense-support-directions}
	设 \(C\subseteq\mathbb R^n\) 为非空紧凸集，\(E\) 为单位球面中的稠密子集。证明
	\[
		C=\bigcap_{u\in E}\{x:\langle x,u\rangle\leq h_C(u)\}.
	\]
	证明闭单位圆盘不能只用有限个方向的支撑半平面表示。最后给出一个非空闭凸集和一个单位圆周中的稠密方向集，说明第一个结论在没有紧性时可能失败。
\end{exercise}
\begin{solution}
	由于 \(C\) 紧，可取 \(R\geq0\) 使 \(\lVert c\rVert\leq R\) 对每个 \(c\in C\) 成立。于是
	\[
		\langle c,u\rangle\leq\langle c,v\rangle+R\lVert u-v\rVert
		\leq h_C(v)+R\lVert u-v\rVert.
	\]
	对 \(c\) 取上确界，再交换 \(u,v\)，得 \(\lvert h_C(u)-h_C(v)\rvert\leq R\lVert u-v\rVert\)，故支撑函数连续。若 \(x\) 满足 \(E\) 中所有方向的支撑不等式，对任意单位向量 \(u\) 选择 \(u_j\in E\) 且 \(u_j\to u\)，取极限得到 \(\langle x,u\rangle\leq h_C(u)\)。由正齐次性，不等式对所有非零方向成立，零方向也成立。\cref{g1:prop:support-function-representation}因此给出 \(x\in C\)。反向包含由支撑函数定义成立。

	对平面闭单位圆盘，单位方向 \(u\) 的支撑值为 \(1\)。给定有限个单位方向 \(u_1,\ldots,u_m\)，选择不同于它们的一个单位向量 \(z\)。则 \(\langle z,u_i\rangle<1\) 对所有 \(i\) 成立。因为只有有限个严格不等式，可以取充分接近 \(1\) 的 \(t>1\)，使 \(\langle tz,u_i\rangle\leq1\) 对每个 \(i\) 仍成立；例如令 \(M=\max_i\langle z,u_i\rangle<1\)，若 \(M>0\) 就取 \(1<t<1/M\)，若 \(M\leq0\) 则任意 \(t>1\) 均可。这样 \(tz\) 属于所有给定半平面的交，却在圆盘之外。没有约束的空方向组更不能表示圆盘。

	最后取 \(C=\mathbb R\times\{0\}\)，令 \(E\) 为单位圆周去掉 \((0,1)\) 和 \((0,-1)\) 后的集合。\(E\) 仍稠密，但每个 \(u\in E\) 的第一坐标非零，沿直线 \(C\) 可以使 \(\langle x,u\rangle\) 任意大，所以 \(h_C(u)=+\infty\)。这些方向的半空间没有提供任何约束，其交为整个 \(\mathbb R^2\)。只有被删去的两个法向方向能检测该直线的竖直位置，这说明紧性保证的有限性与连续性在第一个论证中确有作用。
\end{solution}

\begin{exercise}\label{g1:ex:convex-two-moments}
	给定实数 \(t_1,\ldots,t_m\) 与非负权重 \(a_1,\ldots,a_m\)，其中 \(\sum_i a_i=1\)。证明可以只保留至多三个原来的数，并重新选择非负且和为 \(1\) 的权重，使 \(\sum_i a_i t_i\) 与 \(\sum_i a_i t_i^2\) 同时不变。对原数 \(0,1,2,3\) 各取权重 \(1/4\)，给出三点表示，并证明不能只使用两个原数。
\end{exercise}
\begin{solution}
	把 \(t_i\) 提升为平面点 \(q_i=(t_i,t_i^2)\)。所要保留的两个数恰好是 \(\sum_i a_i q_i\) 的两个坐标。由\cref{g1:thm:caratheodory}，这个凸组合可由至多三个 \(q_i\) 表示；投影到两个坐标，便同时保留两个指定的量。

	给定数据的两个量为 \(3/2\) 与 \(7/2\)。在 \(0,1,3\) 上依次取权重 \(1/6,1/2,1/3\)，有
	\[
		0\cdot\frac16+1\cdot\frac12+3\cdot\frac13=\frac32,\qquad
		0^2\cdot\frac16+1^2\cdot\frac12+3^2\cdot\frac13=\frac72.
	\]
	若只用两个不同的原数 \(u<v\)，并要求加权平均为 \(\mu=3/2\)，则两个权重唯一确定为 \((v-\mu)/(v-u)\)、\((\mu-u)/(v-u)\)。它们非负要求 \(u\le3/2\le v\)，故只有四种候选：\((0,2),(0,3),(1,2),(1,3)\)。对应的平方加权平均由恒等式 \(t^2=(u+v)t-uv\)（在 \(t=u,v\) 处成立）分别得到
	\[
		3,\quad \frac92,\quad \frac52,\quad 3.
	\]
	没有一个等于 \(7/2\)。单点表示的平均又只能是某个原数，不可能等于 \(3/2\)，所以三个点确实必需。
\end{solution}

\begin{exercise}\label{g1:ex:radon-explicit}
	令 \(p_1=(0,0),p_2=(3,0),p_3=(2,2),p_4=(0,1)\)。求这四点的全部仿射依赖关系，写出一个 Radon 分割及其公共点，并判断分割是否唯一，交换两组不计。
\end{exercise}
\begin{solution}
	关系 \(\sum_i c_i p_i=0\)、\(\sum_i c_i=0\) 等价于
	\[
		3c_2+2c_3=0,\qquad 2c_3+c_4=0,\qquad c_1+c_2+c_3+c_4=0.
	\]
	以 \(c_3=3s\) 参数化，得到 \((c_1,c_2,c_3,c_4)=s(5,-2,3,-6)\)。取 \(s>0\)，正、负指标给出 \(\{1,3\}\) 与 \(\{2,4\}\)，归一化为
	\[
		\frac58p_1+\frac38p_3=\frac14p_2+\frac34p_4=\left(\frac34,\frac34\right).
	\]
	这就是两条对角线的交点。任何 Radon 分割若给出两凸组合相等，相减便是一个非零仿射依赖关系。这里所有这样的关系都是上述向量的非零倍数，且四个系数都非零；因此两组只能是它的正指标与负指标，分割在交换两组的意义下唯一。
\end{solution}

\begin{exercise}\label{g1:ex:helly-on-line}
	在 \(\mathbb R^n\) 中给定有限个闭球 \(B_i=\{x:\|x-a_i\|\le r_i\}\)，其中 \(r_i\ge0\)，以及直线 \(L=\{p+tv:t\in\mathbb R\}\)，其中 \(\|v\|=1\)。证明：若任意至多两个球与 \(L\) 有共同点，则所有球与 \(L\) 有共同点。写出每个 \(B_i\cap L\) 所对应的参数区间。进一步，在平面中取圆心 \((0,0),(2,0)\)、半径均为 \(\sqrt2\) 的两个闭圆盘，求它们与直线 \(y=1\) 的共同点。
\end{exercise}
\begin{solution}
	记 \(\alpha_i=\langle a_i-p,v\rangle\)，并将 \(a_i-p=\alpha_i v+w_i\)，其中 \(w_i\perp v\)。由正交性，
	\[
		\|p+tv-a_i\|^2=(t-\alpha_i)^2+\|w_i\|^2.
	\]
	若 \(r_i<\|w_i\|\)，交集为空；否则参数构成闭区间
	\[
		I_i=\left[\alpha_i-\sqrt{r_i^2-\|w_i\|^2},\ \alpha_i+\sqrt{r_i^2-\|w_i\|^2}\right].
	\]
	假设保证这些区间非空且两两相交。由\cref{g1:thm:helly}的一维情形，所有 \(I_i\) 有公共参数，对应所有球与 \(L\) 的公共点。也可直接取所有左端点的最大值 \(a\) 和所有右端点的最小值 \(b\)：取得这两个端点的区间相交，故 \(a\le b\)，任取 \(t\in[a,b]\) 即可。

	在给定平面例子中，取 \(p=(0,1),v=(1,0)\)。两圆心到直线的距离都是 \(1\)，区间半长度为 \(\sqrt{2-1}=1\)，故两个参数区间为 \([-1,1]\) 与 \([1,3]\)。唯一公共参数是 \(t=1\)，得到共同点 \((1,1)\)。这里检查两个球就够，是因为未知公共点被限制在一维直线上。
\end{solution}

\begin{exercise}\label{g1:ex:helly-disk-cover}
	设 \(S\) 是平面中的有限非空点集。若 \(S\) 中每个至多含三个点的子集都能被某个半径为 \(1\) 的闭圆盘包含，证明整个 \(S\) 也能被一个这样的圆盘包含。不同子集所用的圆盘允许不同。说明把“三个”改成“两个”后结论不成立。
\end{exercise}
\begin{solution}
	未知量是圆盘中心。对每个 \(p\in S\)，令 \(C_p=\{c\in\mathbb R^2:\|c-p\|\le1\}\)。一个中心 \(c\) 给出的单位闭圆盘包含 \(p\)，当且仅当 \(c\in C_p\)。每个 \(C_p\) 是凸集：若 \(c_1,c_2\in C_p\) 且 \(0\le t\le1\)，由三角不等式，
	\[
		\|(1-t)c_1+tc_2-p\|\le(1-t)\|c_1-p\|+t\|c_2-p\|\le1.
	\]
	题设恰好说任意至多三个 \(C_p\) 有非空交集。由\cref{g1:thm:helly}的二维情形，全部 \(C_p\) 有公共中心，此中心的单位闭圆盘包含整个 \(S\)。

	反例取 \(S=\{(-1,0),(1,0),(0,\sqrt3)\}\)。三条边的长度均为 \(2\)，每对点可被以其中点为中心的单位闭圆盘包含。若一个中心 \((u,v)\) 到前两点的距离都不超过 \(1\)，则
	\[
		(u+1)^2+v^2\le1,\qquad (u-1)^2+v^2\le1.
	\]
	相加得 \(u^2+v^2\le0\)，所以中心只能是 \((0,0)\)，但它到第三点的距离为 \(\sqrt3>1\)。因此三个点不能同处于一个单位闭圆盘中。
\end{solution}

\begin{exercise}
	\label{g1:exr:clipped-square}
	设 \(P=[-1,1]^2\cap\{(x_1,x_2):x_1+x_2\le1\}\)。求其顶点与棱面，求顶点 \((1,0)\) 的外法锥，并给出 \(P^\circ\) 的顶点表示。
\end{exercise}

\begin{solution}
	原正方形的右上角被直线 \(x_1+x_2=1\) 截去。新直线分别与右边、上边交于 \((1,0)\)、\((0,1)\)，所以候选顶点为
	\[
		(-1,-1),\quad (1,-1),\quad (1,0),\quad (0,1),\quad (-1,1).
	\]
	在每个列出的点都有两条法向量线性无关的限制取等号，故由\cref{g1:thm:polytope-hv}证明中的顶点判别，它们确是顶点。任何其他点至多有一条独立限制取等号，不能成为顶点。五个棱面就是上述相邻顶点连接而成的五条边，其所在直线依次为 \(x_2=-1\)、\(x_1=1\)、\(x_1+x_2=1\)、\(x_2=1\)、\(x_1=-1\)。

	记 \(v=(1,0)\)。若 \(u=(a,b)\in N_P(v)\)，把相邻顶点 \((1,-1)\) 和 \((0,1)\) 代入法锥定义，得到
	\[
		-b\le0,\qquad -a+b\le0,
	\]
	也就是 \(a\ge b\ge0\)。反过来，此时对任意 \(x\in P\)，有
	\[
		\langle u,x-v\rangle=(a-b)(x_1-1)+b(x_1+x_2-1)\le0.
	\]
	所以
	\[
		N_P(v)=\{(a,b):a\ge b\ge0\}=\{s(1,0)+t(1,1):s,t\ge0\}.
	\]
	这说明这个角的外法锥由两条相邻边的外法向生成。

	令 \(Q=\operatorname{conv}\{e_1,-e_1,e_2,-e_2,(1,1)\}\)。由极集的定义，\(Q^\circ\) 的不等式恰好是 \(P\) 的五条限制，所以 \(Q^\circ=P\)。\(Q\) 是闭凸且含原点的集合，双极公式给出
	\[
		P^\circ=Q=\operatorname{conv}\{(1,0),(-1,0),(0,1),(0,-1),(1,1)\}.
	\]
	这些点全都是顶点：\((1,1)\) 唯一最大化 \(y_1+y_2\)，\((1,0)\) 唯一最大化 \(2y_1-y_2\)，\((0,1)\) 唯一最大化 \(-y_1+2y_2\)，而 \((-1,0)\)、\((0,-1)\) 分别唯一最大化 \(-y_1\)、\(-y_2\)。
\end{solution}

\begin{exercise}
	\label{g1:exr:segment-ambient-polar}
	设 \(S=[-1,1]\times\{0\}\subseteq\mathbb R^2\)。写出其半空间表示和全部面，求 \(N_S(S)\)、\(N_S(\{(1,0)\})\) 及环境极集 \(S^\circ\)。再在直线 \(L=\mathbb R\times\{0\}\) 中重复这些计算，说明差异来自哪里。
\end{exercise}

\begin{solution}
	半空间表示为
	\[
		-1\le x_1\le1,\qquad x_2\le0,\qquad -x_2\le0.
	\]
	全部面是 \(\varnothing\)、两个端点单点集以及 \(S\) 本身。事实上，若一个面含有线段的某个内点，面的性质迫使两个端点都在其中，凸性再迫使它等于 \(S\)。

	若线性函数 \(ax_1+bx_2\) 在全部 \(S\) 上达到最大值，它沿 \(x_1\) 方向必须恒定，故 \(a=0\)，而 \(b\) 任意。因此
	\[
		N_S(S)=\{(0,b):b\in\mathbb R\}.
	\]
	若只要求在 \((1,0)\) 取最大值，则条件为 \(a(x_1-1)\le0\) 对 \(-1\le x_1\le1\) 成立，等价于 \(a\ge0\)。所以
	\[
		N_S(\{(1,0)\})=\{(a,b):a\ge0,\ b\in\mathbb R\}.
	\]
	极集要求 \(\sup_{|x_1|\le1}ax_1=|a|\le1\)，故
	\[
		S^\circ=\{(a,b):|a|\le1,\ b\in\mathbb R\},
	\]
	是无界条带。在 \(L\) 中只允许法向量也属于 \(L\)，上述三个集合分别变为 \(\{0\}\)、非负半轴和 \([-1,1]\times\{0\}\)。面本身没有改变；消失的是所有垂直于 \(L\) 的方向。它们在 \(S\) 上的配对恒为零，环境极集无法用不等式约束这些分量。
\end{solution}

\begin{exercise}
	\label{g1:exr:polar-transformation}
	设 \(K\subseteq\mathbb R^n\) 为非空紧凸集，\(T\) 为可逆线性变换，\(a\in\mathbb R^n\)。证明
	\[
		(TK)^\circ=T^{-\mathsf T}K^\circ,\qquad (K+a)^\circ=\{y:h_K(y)+\langle y,a\rangle\le1\}.
	\]
	计算矩形 \([-2,2]\times[-3,3]\) 的极体；再以区间为例说明平移后的极集通常不是原极集的相同平移。
\end{exercise}

\begin{solution}
	由内积与转置的关系，\(y\in(TK)^\circ\) 当且仅当
	\[
		\langle y,Tx\rangle=\langle T^{\mathsf T}y,x\rangle\le1\qquad(x\in K),
	\]
	即 \(T^{\mathsf T}y\in K^\circ\)。乘以逆变换即得第一式。第二式来自
	\[
		\max_{x\in K}\langle y,x+a\rangle=h_K(y)+\langle y,a\rangle.
	\]
	矩形是正方形经 \(T=\operatorname{diag}(2,3)\) 的像，所以其极体为
	\[
		\{(y_1,y_2):2|y_1|+3|y_2|\le1\},
	\]
	顶点是 \((\pm1/2,0)\)、\((0,\pm1/3)\)。长轴方向在取极以后变短，原因是原集合沿该方向越远，相应线性函数允许的系数就越小。

	取 \(K=[-1,1]\)、\(a=1/2\)。则 \(K+a=[-1/2,3/2]\)，正数 \(y\) 要满足 \((3/2)y\le1\)，负数 \(y\) 要满足 \((-1/2)y\le1\)，从而
	\[
		(K+a)^\circ=[-2,2/3].
	\]
	而 \(K^\circ+a=[-1/2,3/2]\)，二者不同。这也说明，即使平移前后原点都在内部，极体仍然依赖所选原点。
\end{solution}

\begin{exercise}
	\label{g1:exr:polar-intersection-hull}
	设 \(K,L\subseteq\mathbb R^n\) 为紧凸集，且 \(0\) 同时在二者内部。证明
	\[
		\bigl(\operatorname{conv}(K\cup L)\bigr)^\circ=K^\circ\cap L^\circ,\qquad (K\cap L)^\circ=\operatorname{conv}(K^\circ\cup L^\circ).
	\]
	用 \(K=[-2,2]\times[-1,1]\)、\(L=[-1,1]\times[-2,2]\) 核验第二式。
\end{exercise}

\begin{solution}
	一个线性函数在 \(\operatorname{conv}(K\cup L)\) 上不超过一，当且仅当它在 \(K\) 和 \(L\) 上分别不超过一，故第一式直接成立。

	证明第二式时需要确保再次取极不会丢失原集合。由\cref{g1:thm:bipolar}的球夹逼，\(K^\circ,L^\circ\) 仍为紧凸集，并含原点为内点。令 \(Q=\operatorname{conv}(K^\circ\cup L^\circ)\)。因为两个集合各自凸，任意混合凸组合都可按来自哪个集合归组，故
	\[
		Q=\{(1-t)y+tz:0\le t\le1,\ y\in K^\circ,\ z\in L^\circ\}.
	\]
	这是紧集 \([0,1]\times K^\circ\times L^\circ\) 在连续映射下的像，所以 \(Q\) 紧，特别是闭。对第一式应用双极公式，有
	\[
		Q^\circ=(K^\circ)^\circ\cap(L^\circ)^\circ=K\cap L.
	\]
	由于 \(Q\) 闭、凸且含原点，再取一次极便得 \((K\cap L)^\circ=Q\)。

	例中 \(K\cap L=[-1,1]^2\)，其极体是 \(D=\{|y_1|+|y_2|\le1\}\)。另一方面，
	\[
		K^\circ=\{2|y_1|+|y_2|\le1\},\qquad L^\circ=\{|y_1|+2|y_2|\le1\}.
	\]
	二者都在 \(D\) 内；\(K^\circ\) 含 \(\pm e_2\)，\(L^\circ\) 含 \(\pm e_1\)，所以它们的凸包包含 \(D\) 的全部顶点，也就包含 \(D\)。两向包含给出所需等式。
\end{solution}

\begin{exercise}
	\label{g1:exr:cube-face-join}
	在 \(C_3=[-1,1]^3\) 中，令 \(v=(1,1,1)\)、\(w=(-1,1,1)\)。求面格中的 \(\{v\}\vee\{w\}\) 与 \(\{v\}\vee\{-v\}\)，并在极体 \(O_3\) 中计算对应面的交，核对包含关系反转。
\end{exercise}

\begin{solution}
	包含 \(v,w\) 的面必须包含线段 \([v,w]\)。该线段本身是立方体的棱：它是线性函数 \(x_2+x_3\) 的最大值集。因此
	\[
		\{v\}\vee\{w\}=[-1,1]\times\{1\}\times\{1\}.
	\]
	若一个面含有 \(v,-v\)，则由凸性含有原点。对任意 \(z\in C_3\)，还有 \(-z\in C_3\)，而 \(0=(z+(-z))/2\)。面的线段性质迫使 \(z\) 属于该面，所以这个面只能是整个 \(C_3\)。故 \(\{v\}\vee\{-v\}=C_3\)，尽管两个点的凸包只是一条体对角线。

	由极面对偶的定义，
	\[
		\{v\}^*=\operatorname{conv}\{e_1,e_2,e_3\},\qquad \{w\}^*=\operatorname{conv}\{-e_1,e_2,e_3\}.
	\]
	第一个三角形中 \(y_1\ge0\)，第二个中 \(y_1\le0\)，所以交集要求 \(y_1=0\)；余下条件是 \(y_2,y_3\ge0\)、\(y_2+y_3=1\)。因此交集恰为 \(\operatorname{conv}\{e_2,e_3\}\)，也是原棱的极对偶面。另一方面，\(\{-v\}^*=\operatorname{conv}\{-e_1,-e_2,-e_3\}\)。这个三角形上的坐标和为 \(-1\)，而 \(\{v\}^*\) 上为 \(1\)，二者交集为空，正好对应 \(C_3^*=\varnothing\)。
\end{solution}

\begin{exercise}
	\label{g1:exr:nonexposed-point}
	令 \(B=\{(x_1,x_2):x_1^2+x_2^2\le1\}\)、\(p=(2,0)\)，并置 \(K=\operatorname{conv}(B\cup\{p\})\)。证明 \(q=(1/2,\sqrt3/2)\) 是 \(K\) 的极点，但 \(\{q\}\) 不是暴露面。指出多胞形的面定理为何不能应用于这个例子。
\end{exercise}

\begin{solution}
	注意 \(\lVert q\rVert=1\) 且 \(\langle q,p\rangle=1\)。在单位圆盘上，Cauchy--Schwarz 不等式给出 \(\langle q,x\rangle\le1\)，等号仅在 \(x=q\) 时成立。因此对 \(K\)，线性函数 \(\langle q,-\rangle\) 的最大值为一，最大值集恰为线段 \([q,p]\)：把 \(K\) 中的点写成 \((1-t)b+tp\)，\(b\in B\)，若 \(t<1\) 且配对为一，就必须有 \(b=q\)。

	假设 \(q=(1-t)x+ty\)，其中 \(x,y\in K\)、\(0<t<1\)。对两边与 \(q\) 配对，得到两个不超过一的数以正权重平均为一，故 \(x,y\in[q,p]\)。而 \(q\) 是这条线段的端点，所以 \(x=y=q\)。这证明 \(q\) 是极点。

	若 \(\{q\}\) 是暴露面，就存在非零 \(u\) 使 \(q\) 成为 \(\langle u,-\rangle\) 在 \(K\) 上的唯一最大值点，因而也是它在 \(B\) 上的最大值点。单位圆盘上线性函数的唯一最大值点是 \(u/\lVert u\rVert\)，所以 \(u=tq\)、\(t>0\)。但这时
	\[
		\langle u,p\rangle=t\langle q,p\rangle=t=\langle u,q\rangle,
	\]
	与唯一性矛盾。因此 \(\{q\}\) 虽然是一个面，却不是暴露面。

	这里的 \(K\) 并非多胞形。例如单位圆左半边每个点 \(r\) 都满足 \(r_1<0\)，函数 \(\langle r,-\rangle\) 在 \(B\) 上唯一最大化于 \(r\)，而 \(\langle r,p\rangle=2r_1<1\)，故在 \(K\) 上仍唯一最大化于 \(r\)。于是 \(K\) 有无穷多个极点，不能由有限点集生成。有限性假设在面定理中确实不可省略。
\end{solution}

\begin{exercise}\label{g1:ex:orthant-distance}
	令 \(C=\mathbb R_{\ge0}^n\)，并定义 \((x_+)_i=\max\{x_i,0\}\)、\((x_-)_i=\max\{-x_i,0\}\)。求 \(P_Cx\)、\(d(x,C)\) 及 \(h_C(u)\)，并在 \(x\notin C\) 时求\eqref{g1:eq:distance-support-formula}的一个最大化方向。
\end{exercise}
\begin{solution}
	平方距离 \(\sum_i(x_i-y_i)^2\) 在约束 \(y_i\ge0\) 下可以逐坐标最小化。当 \(x_i\ge0\) 时取 \(y_i=x_i\)，否则取 \(y_i=0\)，故 \(P_Cx=x_+\)、\(d(x,C)=\lVert x_-\rVert\)。若某个 \(u_i>0\)，令 \(y_i\to+\infty\) 而其余坐标为零，得到 \(h_C(u)=+\infty\)；若所有 \(u_i\le0\)，则 \(\langle u,y\rangle\le0\)，且在 \(y=0\) 处达到零，所以 \(h_C(u)=0\)。当 \(x_-\ne0\) 时，取 \(u=-x_-/\lVert x_-\rVert\)。它是单位向量且各坐标非正，于是
	\[
		\langle u,x\rangle-h_C(u)=\frac{\langle-x_-,x_+-x_-\rangle}{\lVert x_-\rVert}
		=\lVert x_-\rVert.
	\]
	这里 \(\langle x_-,x_+\rangle=0\)，因为每一坐标不可能同时有非零的正部与负部。
\end{solution}

\begin{exercise}\label{g1:ex:cone-lattice-generators}
	在 \(\mathbb R^2\) 中令 \(C=\operatorname{cone}\{(1,0),(1,2)\}\)。
	\begin{enumerate}
		\item 写出 \(C\) 与 \(C^\vee\) 的半空间表示，并给出 \(C^\vee\) 的两个生成向量。
		\item 证明 \(C\cap\mathbb Z^2\) 中每一点都是 \((1,0),(1,1),(1,2)\) 的非负整数线性组合；说明只用原来的两个生成向量为何不够。
	\end{enumerate}
\end{exercise}
\begin{solution}
	1. 写 \((x,y)=s(1,0)+t(1,2)\)，得到 \(y=2t\)、\(x=s+t\)，所以 \(C=\{(x,y):y\ge0,\ 2x-y\ge0\}\)。由生成向量上的有限检验，
	\[
		C^\vee=\{(a,b):a\ge0,\ a+2b\ge0\}.
	\]
	每个这样的 \((a,b)\) 可以写成 \((a/2)(2,-1)+(b+a/2)(0,1)\)，两个系数都非负；反之这两个生成向量均满足不等式。因此 \(C^\vee=\operatorname{cone}\{(2,-1),(0,1)\}\)。

	2. 若 \((x,y)\in C\cap\mathbb Z^2\)，将 \(y\) 除以 \(2\)，写成 \(y=2q+r\)，其中 \(q\ge0\)、\(r\in\{0,1\}\)。由 \(x\ge y/2\) 及 \(x\) 为整数，得 \(x\ge q+r\)，从而
	\[
		(x,y)=(x-q-r)(1,0)+r(1,1)+q(1,2).
	\]
	相反，三个向量的非负整数线性组合显然都属于 \(C\cap\mathbb Z^2\)。原来的两个向量只能产生第二坐标为偶数的整数线性组合，不能产生 \((1,1)\)。因此描述同一个实锥所需的向量，与描述其全部格点加法所需的向量，可以不同。
\end{solution}

\begin{exercise}\label{g1:ex:moment-square-fiber}
	在\eqref{g1:eq:linear-moment-map}中取四个权重 \((0,0),(1,0),(0,1),(1,1)\)，并限制在 \(\sum_{j=0}^3|z_j|^2=2\) 上。证明像为 \([0,1]^2\)，且对给定的 \((a,b)\in[0,1]^2\)，所有可能的四个数 \(t_j=|z_j|^2/2\) 恰为
	\[
		(t_0,t_1,t_2,t_3)=(1-a-b+s,\ a-s,\ b-s,\ s),
	\]
	其中 \(\max\{0,a+b-1\}\le s\le\min\{a,b\}\)。解释为何“凸组合的系数”通常不能唯一确定原像。
\end{exercise}
\begin{solution}
	约束为 \(t_j\ge0\)、\(\sum_jt_j=1\)，而映射值为 \((t_1+t_3,t_2+t_3)\)。令 \(s=t_3\)，解这三个线性等式就得到题中的四元组；非负条件恰好给出所列区间。当 \(0\le a,b\le1\) 时，\(0\le\min\{a,b\}\)，且 \(a+b-1\le a,b\)，所以区间非空。对任意其中的 \(s\)，取 \(z_j=\sqrt{2t_j}\) 就得到所需原像。反过来，任何原像都产生这样的系数，且 \(a=t_1+t_3\)、\(b=t_2+t_3\) 均属于 \([0,1]\)，故像恰为正方形。

	在内部点，例如 \(a=b=1/2\) 处，参数 \(s\) 可以在 \([0,1/2]\) 中变化，已给出不同的模长分配。即使四个 \(t_j\) 都固定，每个非零 \(z_j\) 的复相位仍可改变。因此系数只记录原像的一部分信息，不能把像中的凸组合表示当作原像的唯一坐标。
\end{solution}


% !TeX root = ../Book1.tex
% 纲要标题：# 第二编　拓扑构造、覆盖与曲面分类
% 纲要位置：几何学计划纲要.md，第 232 行。
% 正文状态：未撰写；本文件目前是框架。

% 纲要内容开始
% 
% 纲要内容结束

\BookPart{拓扑构造、覆盖与曲面分类}{g1:part:02}
\BookPartGuide
待撰写

% !TeX root = ../../Book1.tex
% 纲要标题：## 第4章　拓扑空间、商构造与局部—整体工具
% 纲要位置：几何学计划纲要.md，第 234 行。
% 正文状态：初稿；尚未正式审核。

% 纲要内容开始
% 
% > 本章目标与证明范围：为后续流形和商空间建立足够的一般拓扑。基本定义在同一构造中连续出现；较深部分集中在分离性、局部有限性、度量化和单位分解，而非平均分配给每个名词。
% 
% 纲要内容结束

\chapter{拓扑空间、商构造与局部—整体工具}
\label{g1:ch:04}
\BookChapterNumber{g1:ch:04}

\begin{chapterabstract}
	拓扑记录邻近关系与连续变化，使空间不必先以某个 Euclid 空间中的子集出现。子空间、积与商分别处理限制、同时记录和识别；紧性与连通性帮助辨认这些构造保留了什么。为了把局部 Euclid 模型组织成可操作的整体，还需要分离性、可数性和局部有限性。本章证明局部紧第二可数 Hausdorff 空间的仿紧性、正则第二可数空间的度量化，以及仿紧 Hausdorff 空间上连续单位分解的存在，并用环面、实射影空间与两个非分离商模型核验这些条件。
\end{chapterabstract}

\begin{learninggoals}
	\begin{itemize}
		\item 用开集、拓扑基和连续映射描述子空间、积空间及商空间，逐项检查映射的下降与泛性质。
		\item 使用紧性证明最值与同胚判据，区分连通、道路连通和局部道路连通，判断分支在映射下如何变化。
		\item 从局部紧性与第二可数性构造紧穷竭和局部有限细化，说明 Hausdorff、正则与正规条件在证明中的不同作用。
		\item 构造分离闭集的连续函数，并由可数函数族建立与原拓扑相容的距离。
		\item 从覆盖收缩与截断函数构造单位分解，粘合连续函数和正定型，同时识别不能直接加权的条件。
		\item 完整核验环面与实射影空间的商模型，利用双原点和稠密轨道商说明局部形式不保证整体分离性。
	\end{itemize}
\end{learninggoals}

把线段 \([0,1]\) 的两个端点识别成一个点，会得到圆周。这个说法在图形上很直观，但如果只把 \(0\) 与 \(1\) 规定为同一个等价类，还没有说明新空间中的“附近”是什么意思。特别是粘合点附近的一小段圆弧，在原线段中会分成靠近两端的两小段。映射
\[
	q:[0,1]\longrightarrow S^1,\qquad q(t)=(\cos2\pi t,\sin2\pi t)
\]
提示了合适的判断：圆周中的一个集合应当开，当且仅当它在原线段中的原像开。这样既保留了参数的连续变化，也给识别点规定了正确的邻域。

\ProseParagraphBreak

同样的任务会在几何中反复出现。一个圆柱可以由长方形的一对边识别得到，第二章中的实射影空间可以由球面的对径点识别得到；局部坐标又要求我们把空间的一小块同 Euclid 开集相比较。但是识别过程有时会让不同点无法用不交邻域分开，局部公式也可能在无限多个区域聚集的地方失去连续性。分离性控制点能否由邻域区分，局部有限性控制无穷多份局部数据怎样在每点附近相加；紧性与可数性使这些控制能够在全空间上组织起来。

\ProseParagraphBreak

本章从开集出发建立这些工具。所需基础是集合的并、交、原像与等价关系，实数的完备性、闭有界区间的紧性、连续实函数的基本性质，以及有限维线性代数。Euclid 空间中的开球与极限会作为具体模型；一般空间中的紧性以开覆盖定义，序列刻画则在具备相应条件时使用。关于基础构造与紧性的进一步阅读，可参看 Lee 的正式第二版\cite[第2章, 第19--37页; 第3章, 第49--80页; 第4章, 第86--117页]{Lee2011Topological}，其中分离性与可数性、商空间、局部紧性和仿紧性按相应主题展开。本章所用的度量化和单位分解均在正文证明，光滑截断函数则在光滑流形的微积分工具中继续构造。

% !TeX root = ../../Book1.tex
% 纲要标题：### 4.1 拓扑与映射的基本构造
% 纲要位置：几何学计划纲要.md，第 238 行。
% 正文状态：初稿完成；尚未正式审核。

% 纲要内容开始
% 
% * 度量拓扑、拓扑基及连续性
% * 子空间、积空间与映射的泛性质
% * 商拓扑、粘合及商映射的检验
% * 同胚不变量与商空间的反例
% 
% 纲要内容结束

\section{拓扑与映射的基本构造}
\label{g1:sec:0401}

在平面中判断一条曲线是否连续，可以测量点的距离；把曲线的两个端点识别后，距离公式却未必还合用。我们仍希望知道哪些点足够接近、原来的连续映射何时能通过识别后的空间。开集提供了一套能同时处理这两种情形的语言。本节始终把构造空间与检验映射放在一起：子空间和积空间适合检验映入的映射，商空间适合检验从它映出的映射。

\subsection{度量拓扑、拓扑基及连续性}
\label{g1:subsec:0401:01}

Euclid 开球有两个与具体半径无关的性质：开集的任意并仍开，有限个开集的交仍开。第二个条件只能要求有限交，例如实直线上的开区间 \((-1/m,1/m)\) 的交是单点集 \(\{0\}\)。我们把这两项性质保留下来，允许点附近的形状由问题决定。

\begin{definition}[拓扑空间]\label{g1:def:topology}
	设 \(X\) 为集合。若其子集族 \(\mathcal T\) 含 \(\varnothing,X\)，对任意并和有限交封闭，则称 \(\mathcal T\) 为 \(X\) 上的一个\term[tuopu]{拓扑}[topology]，称 \((X,\mathcal T)\) 为\term[tuopukongjian]{拓扑空间}[topological space]。\(\mathcal T\) 中的集合称为\term[kaiji]{开集}[open set]，开集的补集称为\term[biji]{闭集}[closed set]。若 \(N\subseteq X\) 包含一个含 \(x\) 的开集，则称 \(N\) 为 \(x\) 的\term[linyu]{邻域}[neighborhood]。
\end{definition}

邻域本身不必开；“开邻域”才同时要求开性。由补集运算，任意闭集的交和有限闭集的并仍闭。\(\varnothing\) 与 \(X\) 既开又闭，因此“开”与“闭”并非互相排斥的两种分类。一个集合也可以既不开又不闭，例如通常实直线中的 \([0,1)\)。

对 \(A\subseteq X\)，所有包含于 \(A\) 的开集之并称为它的\term[neibu]{内部}[interior]，记为 \(\operatorname{Int}_X A\)；所有包含 \(A\) 的闭集之交称为它的\term[bibao]{闭包}[closure]，记为 \(\overline A^{\,X}\)，环境明确时省去上标。它们分别是包含于 \(A\) 的最大开集和包含 \(A\) 的最小闭集。集合 \(\overline A\setminus\operatorname{Int}_X A\) 称为 \(A\) 的\term[bianjie]{边界}[boundary]，记为 \(\partial_X A\)。这里的边界由所在空间决定；同一子集放入较小的环境，边界可能改变。

\begin{proposition}[闭包的邻域判据]\label{g1:prop:closure-neighborhood}
	设 \(A\) 为拓扑空间 \(X\) 的子集。点 \(x\) 属于 \(\overline A\)，当且仅当 \(x\) 的每个开邻域都与 \(A\) 相交。此外，
	\[
		X\setminus\overline A=\operatorname{Int}_X(X\setminus A).
	\]
\end{proposition}
\begin{proof}
	若 \(x\notin\overline A\)，则开集 \(X\setminus\overline A\) 是一个不交 \(A\) 的邻域。反之，若存在含 \(x\) 的开集 \(U\) 与 \(A\) 不交，则闭集 \(X\setminus U\) 包含 \(A\)，所以 \(\overline A\subseteq X\setminus U\)，特别是 \(x\notin\overline A\)。这个判据同时说明，\(x\notin\overline A\) 等价于它属于 \(X\setminus A\) 中的某个开集，给出所写等式。
\end{proof}

若 \(\overline A=X\)，称 \(A\) 在 \(X\) 中\term[choumi]{稠密}[dense]。由刚才的判据，这等价于 \(A\) 与每个非空开集相交；有理数在通常实直线中稠密，就是这一条件的熟悉例子。

度量是产生拓扑的一种重要方式，但不是拓扑定义的一部分。

\begin{definition}[度量及其拓扑]\label{g1:def:metric-topology}
	设 \(X\) 为集合。若函数 \(d:X\times X\to[0,\infty)\) 对所有 \(x,y,z\in X\) 都满足
	\[
		d(x,y)=0\Longleftrightarrow x=y,\qquad d(x,y)=d(y,x),\qquad d(x,z)\le d(x,y)+d(y,z),
	\]
	则称这个函数为\term[duliang]{度量}[metric]。记 \(B_d(x,r)=\{y:d(x,y)<r\}\)，其中 \(r>0\)。规定 \(U\subseteq X\) 开，当且仅当每个 \(x\in U\) 都有 \(r>0\) 使 \(B_d(x,r)\subseteq U\)；所得拓扑称为 \(d\) 的\term[duliangtuopu]{度量拓扑}[metric topology]。
\end{definition}

这确实给出拓扑：并集中的点已经落在某一个成员中，可沿用该成员提供的球；有限交中的点可取有限个正半径的最小值。每个球本身也开，因为若 \(y\in B_d(x,r)\)，则三角不等式给出
\[
	B_d\bigl(y,r-d(x,y)\bigr)\subseteq B_d(x,r).
\]
前三章使用的 Euclid 距离就是这个定义的例子。另一个例子是离散度量：不同点的距离都取 \(1\)。此时每个单点都是半径 \(1/2\) 的球，所有子集都开；这一拓扑称为\term[lisantuopu]{离散拓扑}[discrete topology]。在另一端，只有 \(\varnothing,X\) 开的拓扑称为\term[pingfantuopu]{平凡拓扑}[indiscrete topology]。当 \(X\) 至少有两个点时，平凡拓扑不能来自度量，因为度量总能用足够小的球区分不同点。

\begin{definition}[拓扑基]\label{g1:def:topological-basis}
	设 \(X\) 为拓扑空间，\(\mathcal B\) 是它的一族开集。若对任意开集 \(U\) 和 \(x\in U\)，都存在 \(B\in\mathcal B\) 使 \(x\in B\subseteq U\)，则称这族开集为一个\term[tuopuji]{拓扑基}[topological basis]。等价地，每个开集都是若干基元素的并。
\end{definition}

基把“检查所有开集”缩减为“检查一批足够小的开集”。度量空间的全体开球构成基；实直线中有理端点的开区间也构成基，因为每个开区间内部的点都能夹在两个合适的有理数之间。

\begin{proposition}[从集合族产生拓扑]\label{g1:prop:basis-test}
	设 \(\mathcal B\) 为集合 \(X\) 的子集族。它是某个拓扑的基，当且仅当它覆盖 \(X\)，并且对任意 \(B_1,B_2\in\mathcal B\) 及 \(x\in B_1\cap B_2\)，存在 \(B_3\in\mathcal B\) 使 \(x\in B_3\subseteq B_1\cap B_2\)。此时全体基元素的并组成唯一的相应拓扑。
\end{proposition}
\begin{proof}
	若已经是拓扑基，覆盖性质来自开集 \(X\)，交集条件来自开集 \(B_1\cap B_2\)。反过来，把所有基元素的并组成的集合族记为 \(\mathcal T\)。空并是 \(\varnothing\)，覆盖条件给出 \(X\in\mathcal T\)，任意并显然仍在 \(\mathcal T\) 中。若 \(x\in U\cap V\)，其中 \(U,V\in\mathcal T\)，可先找含 \(x\) 的基元素 \(B_1\subseteq U\)、\(B_2\subseteq V\)，再由假设找到 \(B_3\subseteq U\cap V\)。因此 \(U\cap V\) 也是基元素的并。有限交由归纳得到，故 \(\mathcal T\) 为拓扑，且 \(\mathcal B\) 是它的基。任何以 \(\mathcal B\) 为基的拓扑都必须由这些并组成，所以唯一。
\end{proof}

\begin{definition}[连续映射]\label{g1:def:continuous-map}
	设 \(f:X\to Y\) 为拓扑空间之间的映射。若每个开集 \(V\subseteq Y\) 的原像 \(f^{-1}(V)\) 都在 \(X\) 中开，则称其为\term[lianxuyingshe]{连续映射}[continuous map]。
\end{definition}

这个定义也可只在目标的一组基上检验，因为原像保持任意并。它等价于闭集的原像闭；还等价于逐点条件：对每个 \(x\in X\) 及 \(f(x)\) 的开邻域 \(V\)，存在 \(x\) 的开邻域 \(U\)，使 \(f(U)\subseteq V\)。确实，连续时可取 \(U=f^{-1}(V)\)；反之，各点提供的 \(U\) 把 \(f^{-1}(V)\) 表成开集的并。在度量空间中把 \(V\) 取为 \(\varepsilon\)-球，并在 \(U\) 中取一个 \(\delta\)-球，就恢复了分析学中的 \(\varepsilon\)--\(\delta\) 条件。

在一般拓扑空间中，给定序列 \((x_m)\)。如果 \(x\) 的每个开邻域都包含该序列从某一项起的全部项，就称该序列\term[shoulian]{收敛}[converge]于 \(x\)，记为 \(x_m\to x\)。连续映射保持这种收敛：若 \(x_m\to x\)，则对 \(f(x)\) 的每个开邻域 \(V\)，连续性使 \(f^{-1}(V)\) 成为 \(x\) 的开邻域，因而最终 \(x_m\in f^{-1}(V)\)，即最终 \(f(x_m)\in V\)。所以 \(f(x_m)\to f(x)\)。

恒等映射连续，两个连续映射的复合连续，因为 \((g\circ f)^{-1}(V)=f^{-1}(g^{-1}(V))\)。但一般连续映射不把开集映成开集：常值映射 \(\mathbb R\to\mathbb R\) 已是反例。把每个开集都映成开集的映射称为\term[kaiyingshe]{开映射}[open map]；类似定义\term[biyingshe]{闭映射}[closed map]。连续性检验的是原像，开映射、闭映射检验的是像，三个条件不能混用。

\begin{definition}[同胚]\label{g1:def:homeomorphism}
	双射 \(f:X\to Y\) 若与逆映射 \(f^{-1}:Y\to X\) 都连续，则称为\term[tongpei]{同胚}[homeomorphism]。存在同胚的两个空间称为同胚的，记为 \(X\cong Y\)。
\end{definition}

对连续双射，逆映射连续等价于原映射开，也等价于原映射闭。这不是额外定理：逆映射对于集合的原像正是原映射对于该集合的像。因而证明一个显式双射是同胚时，既可以写出逆映射并检查连续性，也可以证明开性或闭性。

\subsection{子空间、积空间与映射的泛性质}
\label{g1:subsec:0401:02}

圆周 \(S^1=\{(x,y)\in\mathbb R^2:x^2+y^2=1\}\) 上的开集通常不是平面开集：一小段圆弧在平面内没有面积，却应当是圆周上的开集。若把对象放在一个已有空间中，就应把环境的开集限制到这个对象上。

设 \(A\subseteq X\)。集合族 \(\{A\cap U:U\subseteq X\text{ 开}\}\) 在 \(A\) 上组成拓扑，称为\term[zikongjiantuopu]{子空间拓扑}[subspace topology]。公理直接来自交集与并、交运算的相容性。这个定义还说明：子空间 \(A\) 中的闭集恰为 \(A\cap F\)，其中 \(F\) 在 \(X\) 中闭。例如 \([0,1)\) 在子空间 \([0,1]\) 中开，因为它等于 \([0,1]\cap(-1,1)\)；它并不在 \(\mathbb R\) 中开。

\begin{proposition}[子空间的映入判据]\label{g1:prop:subspace-universal}
	设 \(A\subseteq X\) 带子空间拓扑，\(i:A\hookrightarrow X\) 为包含映射。对任意拓扑空间 \(Z\) 及集合映射 \(f:Z\to A\)，\(f\) 连续当且仅当 \(i\circ f\) 连续。
\end{proposition}
\begin{proof}
	包含映射连续，因为 \(i^{-1}(U)=A\cap U\)。所以 \(f\) 连续时复合连续。反过来，若 \(i\circ f\) 连续，则对 \(A\) 中开集 \(A\cap U\)，有 \(f^{-1}(A\cap U)=(i\circ f)^{-1}(U)\)，故原像开。
\end{proof}

这一性质把子空间拓扑与映射的目标联系起来：给定连续映射 \(g:Z\to X\)，若它的像包含于 \(A\)，就存在唯一的连续映射 \(f:Z\to A\) 使下图交换。
\[
	\begin{tikzcd}[column sep=large,row sep=large]
		& A \arrow[d,hook,"i"] \\
		Z \arrow[ur,dashed,"f"] \arrow[r,"g"'] & X
	\end{tikzcd}
\]
这里“唯一”来自集合映射的取值已经被 \(g\) 决定；拓扑保证这个唯一的映射连续。这类用存在唯一映射来刻画构造的性质称为\term[fanxingzhi]{泛性质}[universal property]。

积空间处理的是把几个量同时记录下来。对拓扑空间 \(X,Y\)，全体开矩形 \(U\times V\) 满足\cref{g1:prop:basis-test}，由它们生成的拓扑称为 \(X\times Y\) 的\term[jituopu]{积拓扑}[product topology]。有限个空间的积同样以 \(\prod_{i=1}^mU_i\) 为基。对任意指标集 \(I\)，积 \(\prod_{i\in I}X_i\) 的基本开集则规定为
\begin{equation}\label{g1:eq:product-basic-open}
	\bigcap_{i\in F}\pi_i^{-1}(U_i),\qquad F\subseteq I\text{ 有限},\quad U_i\subseteq X_i\text{ 开},
\end{equation}
其中 \(\pi_i\) 是第 \(i\) 个坐标投影。也就是说，一次只限制有限个坐标，其余坐标任取。交两个这样的集合仍只限制有限个坐标，所以确实给出一组基。

\begin{proposition}[积空间的映入判据]\label{g1:prop:product-universal}
	设 \((X_i)_{i\in I}\) 为拓扑空间族，积空间取上述积拓扑。映射 \(f:Z\to\prod_iX_i\) 连续，当且仅当每个坐标映射 \(f_i=\pi_i\circ f\) 连续。因此给定连续映射族 \(f_i:Z\to X_i\)，存在唯一的连续映射 \(f\) 满足 \(\pi_i\circ f=f_i\)。
\end{proposition}
\begin{proof}
	投影 \(\pi_i\) 连续，因为 \(\pi_i^{-1}(U_i)\) 是基本开集。故 \(f\) 连续时各坐标连续。反过来，对\eqref{g1:eq:product-basic-open}中的基本开集，其原像为有限交 \(\bigcap_{i\in F}f_i^{-1}(U_i)\)，所以开。存在唯一的集合映射由 \(f(z)=(f_i(z))_i\) 给出，刚才的判据保证连续。
\end{proof}

有限积的图示为
\[
	\begin{tikzcd}[column sep=large,row sep=large]
		& Z \arrow[dl,"f_1"'] \arrow[d,dashed,"f"] \arrow[dr,"f_2"] & \\
		X_1 & X_1\times X_2 \arrow[l,"\pi_1"] \arrow[r,"\pi_2"'] & X_2 .
	\end{tikzcd}
\]
如果在无限积中允许每个坐标都受一个开集限制，所得的是通常更细的盒拓扑，逐坐标连续就不再够用。例如 \(t\mapsto(t,t,\ldots)\) 从 \(\mathbb R\) 到 \(\mathbb R^{\mathbb N}\) 的所有坐标都连续，而集合 \(\prod_{n\ge1}(-1/n,1/n)\) 的原像是不开的 \(\{0\}\)。这个集合在盒拓扑中开，在积拓扑中并不开：任何含零序列的基本开集只限制有限个坐标，不能保证其余坐标满足所写不等式。

有限个度量空间 \((X_i,d_i)\) 的积拓扑由度量
\[
	D(x,y)=\max_{1\le i\le m}d_i(x_i,y_i)
\]
产生。三角不等式逐坐标成立，再取最大值即可；\(D\)-球是相同半径的开球之积，而每个含 \(x\) 的基本开矩形都包含某个这样的球。因此有限积既能用坐标也能用距离来理解。连续映射在子集上的限制仍连续，这由包含映射连续和复合连续立即得到。

\subsection{商拓扑、粘合及商映射的检验}
\label{g1:subsec:0401:03}

现在把问题倒过来。若 \(q:X\to Y\) 已经规定哪些点要合并，如何给 \(Y\) 选拓扑？目标是使“合并后的一块区域开”能够在原来的空间中检查。由于这块区域可能包含多个代表，必须同时检查它的全部原像，而不能只挑一份代表。

\begin{definition}[商拓扑与商映射]\label{g1:def:quotient-map}
	设 \(q:X\to Y\) 为从拓扑空间到集合的满射。规定
	\[
		U\subseteq Y\text{ 开}\quad\Longleftrightarrow\quad q^{-1}(U)\subseteq X\text{ 开},
	\]
	得到的拓扑称为由 \(q\) 给出的\term[shangtuopu]{商拓扑}[quotient topology]。若 \(Y\) 已有拓扑，且满射 \(q\) 满足这个等价关系，则称 \(q\) 为\term[shangyingshe]{商映射}[quotient map]。当 \(Y=X/{\sim}\) 是等价类的集合且 \(q(x)=[x]\) 时，所得空间称为\term[shangkongjian]{商空间}[quotient space]。
\end{definition}

原像保持任意并和有限交，故上述定义确实满足拓扑公理。商映射必连续；进一步说，商拓扑是使 \(q\) 连续的最细拓扑，因为任何这样的拓扑中的开集都必须有开原像。若 \(A\subseteq X\) 是一些完整等价类的并，即 \(q^{-1}(q(A))=A\)，称它为\term[baoheji]{饱和集}[saturated set]。商空间的开集与原空间的饱和开集一一对应。

\begin{proposition}[映射从商空间下降]\label{g1:prop:quotient-descent}
	设 \(q:X\to Y\) 为商映射，\(Z\) 为拓扑空间。集合映射 \(h:Y\to Z\) 连续，当且仅当 \(h\circ q\) 连续。特别地，若连续映射 \(f:X\to Z\) 在 \(q\) 的每个纤维上取常值，则存在唯一连续映射 \(\bar f:Y\to Z\)，使 \(f=\bar f\circ q\)。
\end{proposition}
\begin{proof}
	一个方向来自复合连续。另一个方向中，对 \(Z\) 的开集 \(V\)，等式
	\[
		q^{-1}\bigl(h^{-1}(V)\bigr)=(h\circ q)^{-1}(V)
	\]
	说明左侧开，再由商映射的判据得到 \(h^{-1}(V)\) 开。若 \(f\) 在纤维上常值，就可定义 \(\bar f(y)=f(x)\)，其中任选 \(q(x)=y\)；纤维上的常值性使选择无关，满射性保证定义处处成立且唯一，连续性由刚证的判据得到。
\end{proof}

下降的图示为
\[
	\begin{tikzcd}[column sep=large,row sep=large]
		X \arrow[r,"f"] \arrow[d,two heads,"q"'] & Z \\
		Y \arrow[ur,dashed,"\bar f"'] &
	\end{tikzcd}
\]
它把构造分成两个实质条件：纤维常值保证映射良定义，商拓扑保证连续。只写一个代表元公式而不检验第一项，就还没有得到商空间上的映射。

\begin{proposition}[商映射的实用判据]\label{g1:prop:quotient-tests}
	连续满射 \(q:X\to Y\) 若为开映射或闭映射，则为商映射。若 \(q\) 是商映射，\(U\subseteq Y\) 开，则限制 \(q:q^{-1}(U)\to U\) 仍为商映射。
\end{proposition}
\begin{proof}
	只需补证“原像开蕴含目标开”。若 \(q\) 开且 \(q^{-1}(V)\) 开，满射性给出 \(V=q(q^{-1}(V))\) 开。若 \(q\) 闭，则 \(q^{-1}(Y\setminus V)\) 闭，其像 \(Y\setminus V\) 闭，所以 \(V\) 开。

	对限制的结论，\(q^{-1}(U)\) 在 \(X\) 中开。若 \(V\subseteq U\) 的原像在 \(q^{-1}(U)\) 中开，它也在 \(X\) 中开，故 \(V\) 在 \(Y\) 中开，因而在 \(U\) 中开；相反方向由限制连续得到。
\end{proof}

开映射和闭映射只是充分条件。实际工作中也可以直接找一组饱和开集来描述商拓扑。尤其要注意，连续满射只能保证“目标开蕴含原像开”，还缺反方向。例如把带离散拓扑的 \(\mathbb R\) 恒等映到带通常拓扑的 \(\mathbb R\)，是连续满射，却不是商映射。

粘合是商构造的一个常见来源。给定空间 \(X,Y\)，先取它们的\term[bujiaobing]{不交并}[disjoint union] \(X\sqcup Y\)：点保留来自哪一份的标签，一个集合开当且仅当它与每一份的交开。给定 \(A\subseteq X\) 及连续映射 \(a:A\to Y\)，再取由关系 \(x\sim a(x)\)（\(x\in A\)）生成的等价关系，得到粘合空间 \(X\cup_aY\)。这里“生成”包括传递性所迫使的所有识别；若 \(a\) 非单射，\(A\) 中不同的点也可能合并。根据下降判据，给定连续映射 \(f_X:X\to Z\)、\(f_Y:Y\to Z\)，它们恰在
\[
	f_X|_A=f_Y\circ a
\]
时粘成唯一连续映射 \(X\cup_aY\to Z\)。这条等式就是粘合处必须满足的相容条件。

另一个方向是：一个空间已经由若干子空间覆盖，能否从各片上的连续性推出整体连续性？需要对覆盖加条件。

\begin{proposition}[粘合引理]\label{g1:prop:gluing-lemma}
	设 \(X=\bigcup_{\alpha\in I}A_\alpha\)，且每个 \(A_\alpha\) 带子空间拓扑。假设或者所有 \(A_\alpha\) 都在 \(X\) 中开，或者 \(I\) 有限且所有 \(A_\alpha\) 都在 \(X\) 中闭。若连续映射 \(f_\alpha:A_\alpha\to Y\) 在交集上两两相同，则它们定义唯一连续映射 \(f:X\to Y\)。
\end{proposition}
\begin{proof}
	交集上的相等性使 \(f(x)=f_\alpha(x)\) 与所选覆盖成员无关，覆盖性保证处处有定义，也给出唯一性。开覆盖情形，对开集 \(V\subseteq Y\)，每个 \(f_\alpha^{-1}(V)\) 在 \(A_\alpha\) 中开，因 \(A_\alpha\) 在 \(X\) 中开，它也在 \(X\) 中开；这些集合的并是 \(f^{-1}(V)\)，所以开。有限闭覆盖情形改用闭集 \(F\subseteq Y\)：\(f_\alpha^{-1}(F)\) 在 \(A_\alpha\) 中闭，因 \(A_\alpha\) 在 \(X\) 中闭，它也在 \(X\) 中闭；有限并 \(f^{-1}(F)\) 仍闭，故 \(f\) 连续。
\end{proof}

例如 \(|x|\) 由两个闭半轴上的公式 \(-x\)、\(x\) 粘成，它们在交点 \(0\) 处一致。有限闭覆盖的“有限”不能删去，因为闭集的任意并不一定闭；习题给出各片均有非空区间的具体反例。后面局部有限覆盖所起的作用之一，正是恢复这种闭集并的控制。

\subsection{同胚不变量与商空间的反例}
\label{g1:subsec:0401:04}

同胚只保留由开集表达的结构。距离、角度、面积不在其中，例如 \(x\mapsto2x\) 是 \(\mathbb R^n\) 的同胚，却改变长度。与此相对，能否把两个不同的点放进不相交的开邻域，是只涉及开集的问题。

\begin{definition}[Hausdorff 空间]\label{g1:def:hausdorff}
	设 \(X\) 为拓扑空间。若任意不同的点 \(x,y\in X\) 都有不相交的开邻域 \(U\ni x\)、\(V\ni y\)，则称它为\term[Hausdorffkongjian]{Hausdorff 空间}[Hausdorff space]。\footnote{Felix Hausdorff（1868--1942），德国数学家，对一般拓扑、集合论及维数理论作出重要贡献。}
\end{definition}

度量空间都是 Hausdorff 的：若 \(d(x,y)=r>0\)，两个半径 \(r/3\) 的球若相交，三角不等式就给出 \(r<2r/3\)，矛盾。子空间继承 Hausdorff 性，因为可把两个不交的环境开邻域与子空间相交；任意积也继承这一性质，因为两个不同的点必在某一坐标不同，拉回该坐标中分离它们的两个开集即可。

若 \(h:X\to Y\) 为同胚且 \(X\) 为 Hausdorff 空间，则 \(Y\) 中不同两点的原像能被不交开集分离，而这两个开集在 \(h\) 下的像仍开，因此 \(Y\) 也为 Hausdorff 空间。像这样在同胚下保持的性质称为\term[tongpeibubianliang]{同胚不变量}[topological invariant]。后面引入的紧性与连通性也将给出不变量；它们有用，是因为一旦发现不变量不同，就能排除所有可能的同胚，而不必逐一尝试公式。

\begin{example}[连续双射不一定是同胚]\label{g1:exm:half-open-circle}
	给 \([0,1)\) 通常的子空间拓扑，给 \(S^1\subseteq\mathbb C\cong\mathbb R^2\) 子空间拓扑。映射
	\[
		f:[0,1)\longrightarrow S^1,\qquad f(t)=e^{2\pi it}
	\]
	连续且为双射，但不是同胚。事实上，\([0,1/4)\) 在 \([0,1)\) 中开，而它的像包含 \(1\)，却不含任何角度为负且足够接近 \(0\) 的圆周点。每个 \(1\) 的圆周开邻域都包含这样的点，所以该像不开。因而 \(f^{-1}\) 不连续。
\end{example}

这个例子的两个空间都来自 Euclid 空间，问题仍会发生。连续双射允许把原来相距很远的端部压到相邻位置，逆映射因此可能把一个很小的邻域撕开。\cref{g1:thm:compact-bijection}将给出一个常用补救：若源空间紧、目标空间 Hausdorff，连续双射就必为同胚。

\begin{example}[合并点可以破坏分离性]\label{g1:exm:two-point-quotient}
	在通常实直线中保留原点为一个等价类，把所有非零点合并为另一个等价类。商集合记为 \(\{a,b\}\)，其中 \(a=[0]\)。单点集 \(\{b\}\) 的原像是开集 \(\mathbb R\setminus\{0\}\)，而 \(\{a\}\) 的原像 \(\{0\}\) 不开，所以商拓扑恰为
	\[
		\{\varnothing,\{b\},\{a,b\}\}.
	\]
	含 \(a\) 的唯一开邻域是整个空间，因此 \(a,b\) 不能被不交开邻域分离。源空间为 Hausdorff 空间，商空间却不是。
\end{example}

所以“把某些点看成同一点”只确定了商构造的输入，并不自动保证结果是通常想象中的几何空间。后面处理环面与实射影空间时，我们会分别核验等价类、连续性与分离性；双原点和有理平移的商则显示这些核验分别可能在哪里失败。

% !TeX root = ../../Book1.tex
% 纲要标题：### 4.2 紧性、连通性与局部条件
% 纲要位置：几何学计划纲要.md，第 245 行。
% 正文状态：初稿完成；未送审。

% 纲要内容开始
% 
% * 紧性与有限交性质
% * 局部紧性与相对紧邻域
% * 连通、道路连通与局部道路连通
% * 分支及其在映射下的行为
% 
% 纲要内容结束

\section{紧性、连通性与局部条件}
\label{g1:sec:0402}

在第三章，单位球面的紧性使一列法向量留下极限方向；有限 Helly 定理却不能直接处理无穷多个约束，因为可行点可能不断逃向远处。本节把这些现象放到一般拓扑空间中：紧性让任意开覆盖缩减为有限覆盖，连通性则排除把空间分成两个互不接触的开部分。两者解决不同问题。局部紧性只要求在每一点附近保留紧性；局部道路连通性则保证足够小的邻域里能够用道路连接点。它们将把后文的局部坐标和整体构造联系起来。

\subsection{紧性与有限交性质}
\label{g1:subsec:0402:01}

设 \(K\subseteq X\)。一个由 \(X\) 的开集组成的族 \(\{U_\alpha\}_{\alpha\in A}\) 若满足 \(K\subseteq\bigcup_{\alpha\in A}U_\alpha\)，就称为 \(K\) 的开覆盖。这里的指标集不必可数，覆盖中的集合也不必彼此不同。

\begin{definition}[紧性与有限交性质]\label{g1:def:compact}
	若拓扑空间 \(X\) 的每个开覆盖都有有限子覆盖，则称 \(X\) 为\term[jinkongjian]{紧空间}[compact space]。若子集 \(K\subseteq X\) 在子空间拓扑下是紧空间，就称其为\term[jinji]{紧集}[compact set]；等价地，\(X\) 中每个覆盖 \(K\) 的开集族都有有限子族覆盖 \(K\)。

	若集合族 \(\{F_\alpha\}_{\alpha\in A}\) 中每个有限子族的交都非空，就称它具有\term[youxianjiaoxingzhi]{有限交性质}[finite intersection property]；空子族的交约定为整个所在空间。
\end{definition}

上述子集定义的等价性来自子空间开集的形式 \(K\cap U\)：把子空间中的覆盖写成这种形式，有限子覆盖的选取与环境空间中的覆盖完全对应。空集是紧集；有限集也紧，因为给每个点选一个覆盖它的开集便已足够。

\ProseParagraphBreak

本章的“紧”只表示开覆盖性质，不暗含 Hausdorff 条件；需要分离性时会在陈述中另列。部分文献把同一开覆盖性质称为\term[nijin]{拟紧}[quasi-compact]，并把“紧”保留给拟紧且 Hausdorff 的空间。读到这类资料时须先转换约定。后文的流形本身要求 Hausdorff，届时两种用法在流形的紧子集上没有差别。

\begin{theorem}[紧性的闭集刻画]\label{g1:thm:compact-fip}
	拓扑空间 \(X\) 紧，当且仅当 \(X\) 的每个具有有限交性质的闭集族都有非空交集。
\end{theorem}
\begin{proof}
	设 \(X\) 紧，而闭集族 \(\{F_\alpha\}\) 具有有限交性质。若全交为空，则补集族 \(\{X\setminus F_\alpha\}\) 是开覆盖。由紧性取有限子覆盖，取补集便得到某个有限子族的交为空，与假设矛盾。

	反之，设每个具有有限交性质的闭集族都有公共点。若开覆盖 \(\{U_\alpha\}\) 没有有限子覆盖，则对任意有限个指标 \(\alpha_1,\ldots,\alpha_m\)，都有
	\[
		\bigcap_{j=1}^m(X\setminus U_{\alpha_j})
		=X\setminus\bigcup_{j=1}^mU_{\alpha_j}\ne\varnothing.
	\]
	因此补集闭族具有有限交性质，应该有公共点，但这个点不属于任何 \(U_\alpha\)，与开覆盖矛盾。若 \(X=\varnothing\)，空覆盖已有有限子覆盖，而约定排除了在空空间中具有有限交性质的族，两种刻画仍然一致。
\end{proof}

这一表述把“有限次约束可解”与“全部约束同时可解”接上了。闭性不能删除：\((0,1/n)\) 是紧空间 \([0,1]\) 的子集，任意有限个相交非空，全交却为空，因为这些子集没有保留趋近端点的极限。紧性也不能删除：\([n,\infty)\) 是 \(\mathbb R\) 中的闭集，有限交非空，全交为空。这正是\cref{g1:exm:helly-infinite-family}中的两个失败机制。

\begin{proposition}[闭子集与有限并]\label{g1:prop:compact-closed-subset}
	紧空间的闭子集是紧集。任意拓扑空间中，有限个紧子集的并仍紧。
\end{proposition}
\begin{proof}
	若 \(F\) 闭于紧空间 \(X\)，给定 \(F\) 的一个环境开覆盖，把 \(X\setminus F\) 加入后便覆盖 \(X\)。取有限子覆盖，再删去 \(X\setminus F\)，得到原覆盖中有限个集合覆盖 \(F\)。

	若 \(K_1,\ldots,K_m\) 都紧，给定它们的并的开覆盖，对每个 \(K_j\) 各取一个有限子覆盖，再合并这 \(m\) 个有限族即可。
\end{proof}

\begin{proposition}[紧集的连续像]\label{g1:prop:compact-image}
	若 \(f:X\to Y\) 连续且 \(K\subseteq X\) 紧，则 \(f(K)\) 是 \(Y\) 的紧子集；这里不要求 \(X,Y\) 为 Hausdorff 空间。
\end{proposition}
\begin{proof}
	给定 \(f(K)\) 的开覆盖 \(\{V_\alpha\}\)，由连续性，\(\{f^{-1}(V_\alpha)\}\) 是覆盖 \(K\) 的开集族。取有限子覆盖后，对应的有限个 \(V_\alpha\) 就覆盖 \(f(K)\)。
\end{proof}

\begin{proposition}[紧集在 Hausdorff 空间中闭]\label{g1:prop:compact-hausdorff-closed}
	若 \(Y\) 为 Hausdorff 空间，\(K\subseteq Y\) 紧，则 \(K\) 闭于 \(Y\)。
\end{proposition}
\begin{proof}
	只需对每个 \(y\in Y\setminus K\) 找到与 \(K\) 不交的开邻域。若 \(K=\varnothing\)，没有需要证明的内容。否则，对每个 \(x\in K\)，Hausdorff 条件给出开集 \(U_x,V_x\)，分别含 \(x,y\)，且 \(U_x\cap V_x=\varnothing\)。从覆盖 \(K\) 的 \(\{U_x\}_{x\in K}\) 中取有限子覆盖 \(U_{x_1},\ldots,U_{x_m}\)，则
	\[
		V=V_{x_1}\cap\cdots\cap V_{x_m}
	\]
	是 \(y\) 的开邻域，并与覆盖 \(K\) 的每个 \(U_{x_j}\) 不交。因此 \(V\cap K=\varnothing\)，证明 \(Y\setminus K\) 开。
\end{proof}

连续实函数在非空紧空间上取得最大值和最小值。事实上，若 \(X\) 非空且紧、\(f:X\rightarrow\mathbb R\) 连续，由\cref{g1:prop:compact-image}，\(f(X)\) 是非空紧子集；由上一命题它在 \(\mathbb R\) 中闭。开区间族 \(\{(-n,n)\}_{n\ge1}\) 覆盖 \(f(X)\)，紧性给出有限子覆盖，所以 \(f(X)\) 还有界。实数完备性给出其上、下确界，它们都是 \(f(X)\) 的极限点或已在其中的点，闭性保证二者属于 \(f(X)\)。因此确有原像取得这两个最值。

这里 Hausdorff 条件先分开一对点，紧性再把所有这些局部分离合并为一次有限交。若没有分离性，结论会失败：两点集 \(\{a,b\}\) 取拓扑 \(\{\varnothing,\{a\},\{a,b\}\}\)，单点 \(\{a\}\) 紧，却不闭，因为其补集 \(\{b\}\) 不开。

\begin{theorem}[由紧性判定同胚]\label{g1:thm:compact-bijection}
	若 \(X\) 紧、\(Y\) 为 Hausdorff 空间，且 \(f:X\to Y\) 是连续双射，则 \(f\) 是同胚。
\end{theorem}
\begin{proof}
	双射已经给出了集合映射 \(f^{-1}:Y\to X\)，只需证明它连续。对任意闭集 \(F\subseteq X\)，\cref{g1:prop:compact-closed-subset}给出 \(F\) 紧；其连续像 \(f(F)\) 紧，由 \(Y\) 为 Hausdorff 及\cref{g1:prop:compact-hausdorff-closed}知 \(f(F)\) 闭。注意这里的 \(f(F)\) 正是映射 \(f^{-1}\) 对 \(F\) 的原像。因此逆映射对闭集取原像仍为闭集，也就连续。
\end{proof}

这一判据常用于参数化和粘合空间：连续性与双射容易直接核对，逆映射却不容易写出。紧性和 Hausdorff 条件让逆映射的连续性自动跟上。条件确有作用：连续双射 \(t\mapsto(\cos2\pi t,\sin2\pi t)\) 从 \([0,1)\) 到圆周不是同胚，原因是在圆周上趋近 \((1,0)\) 的点既可对应趋近 \(0\) 的参数，也可对应趋近 \(1\) 的参数；定义域不紧。若定义域取两点离散空间、值域取同一集合上的平凡拓扑，则恒等映射是从紧空间出发的连续双射，逆映射却不连续；此时值域不是 Hausdorff。

\begin{proposition}[有限积的紧性]\label{g1:prop:finite-product-compact}
	有限个紧拓扑空间的积在积拓扑下紧，不要求各因子为 Hausdorff。
\end{proposition}
\begin{proof}
	只需证明两个紧空间 \(X,Y\) 的积紧，再归纳。若任一因子为空，积为空。设两者非空，给定 \(X\times Y\) 的开覆盖 \(\mathcal U\)。积拓扑的开矩形是基：对每个 \((x,y)\)，可选开矩形 \(A_{x,y}\times B_{x,y}\) 含 \((x,y)\)，并使它包含在某个 \(U_{x,y}\in\mathcal U\) 中。

	先固定 \(x\)。开集族 \(\{B_{x,y}\}_{y\in Y}\) 覆盖 \(Y\)，由 \(Y\) 紧，可选 \(y_1,\ldots,y_m\) 使相应的 \(B_{x,y_j}\) 覆盖 \(Y\)。令
	\[
		A_x=\bigcap_{j=1}^m A_{x,y_j}.
	\]
	这是 \(x\) 的开邻域，而且整个 \(A_x\times Y\) 都被有限个 \(U_{x,y_j}\) 覆盖：给定 \((a,b)\)，先选 \(j\) 使 \(b\in B_{x,y_j}\)，而 \(a\in A_x\subseteq A_{x,y_j}\)。

	再让 \(x\) 变化。由 \(X\) 紧，从覆盖 \(\{A_x\}_{x\in X}\) 中选有限子覆盖 \(A_{x_1},\ldots,A_{x_r}\)。对每个 \(x_i\) 保留上一步的有限个 \(U_{x_i,y_j}\)，把它们合并便覆盖 \(X\times Y\)。有限个因子的结论逐次应用这一证明得到；空积约定为单点空间，也紧。
\end{proof}

有限积证明两次使用紧性：先沿一条“纤维”找到有限个覆盖片，再把这个有限选择保持到纤维附近，最后沿另一因子再取有限次。它只处理有限个因子，没有据此证明任意无限积的紧性。在 Euclid 空间中，本书已使用的基本分析结论“闭且有界当且仅当紧”与这里的开覆盖定义相符；但一般拓扑空间没有范数，也没有可供替代定义的“有界”概念。

\subsection{局部紧性与相对紧邻域}
\label{g1:subsec:0402:02}

\(\mathbb R^n\) 本身不紧，任一点却都藏在一个闭球的内部，而闭球紧。后文常只需在一个给定区域内处理数据，希望先把关注的紧集包进一个稍大的开集，再使这个开集的闭包仍然紧且不碰到区域外部。这样既留下调整数据的空间，又保留紧集上的有限性。

\begin{definition}[局部紧性]\label{g1:def:locally-compact}
	设 \(X\) 为拓扑空间。若每个 \(x\in X\) 都有紧邻域，即存在紧集 \(N\subseteq X\) 和开集 \(O\)，使 \(x\in O\subseteq N\)，则称该空间是\term[jubujin]{局部紧}[locally compact]的。若子集 \(A\subseteq X\) 的闭包 \(\overline A\) 紧，则称 \(A\) 为\term[xiangduijin]{相对紧}[relatively compact]的。
\end{definition}

紧邻域本身未必开。定义只先给出一个紧邻域，而几何构造往往需要在任意指定的小邻域内选择这样的对象。对 Hausdorff 空间，这一步可以做到；证明先处理紧空间内部的缩小。

\begin{lemma}[紧 Hausdorff 空间中的邻域缩小]\label{g1:lem:compact-hausdorff-shrink}
	若 \(X\) 为紧 Hausdorff 空间，\(G\subseteq X\) 开且 \(x\in G\)，则存在开集 \(V\)，使
	\[
		x\in V\subseteq\overline V\subseteq G.
	\]
\end{lemma}
\begin{proof}
	令 \(F=X\setminus G\)。它是紧空间的闭子集，因而紧。若 \(F=\varnothing\)，取 \(V=X\)。否则，对每个 \(y\in F\)，由 Hausdorff 条件选不交开集 \(P_y,Q_y\)，分别含 \(x,y\)。有限个 \(Q_{y_1},\ldots,Q_{y_m}\) 覆盖 \(F\)。取
	\[
		V=\bigcap_{j=1}^mP_{y_j},\qquad Q=\bigcup_{j=1}^mQ_{y_j}.
	\]
	则 \(V\) 开且含 \(x\)，而 \(V\subseteq X\setminus Q\)。后者是闭集，所以 \(\overline V\subseteq X\setminus Q\subseteq X\setminus F=G\)。
\end{proof}

\begin{lemma}[局部紧 Hausdorff 空间中的邻域缩小]\label{g1:lem:local-compact-shrink}
	设 \(X\) 为局部紧 Hausdorff 空间，\(x\in U\subseteq X\)，其中 \(U\) 开。则存在开集 \(V\)，使 \(x\in V\)、\(\overline V\) 紧，且 \(\overline V\subseteq U\)。
\end{lemma}
\begin{proof}
	取 \(x\) 的紧邻域 \(N\) 及开集 \(O\)，满足 \(x\in O\subseteq N\)。由 Hausdorff 条件和\cref{g1:prop:compact-hausdorff-closed}，\(N\) 闭于 \(X\)；它作为子空间仍为 Hausdorff。于是 \(N\) 是紧 Hausdorff 空间，\(G=O\cap U\) 是 \(N\) 中含 \(x\) 的开集。

	在 \(N\) 中应用上一引理，得到 \(N\) 中的开集 \(W\)，满足 \(x\in W\subseteq\overline W^{\,N}\subseteq G\)。写 \(W=N\cap H\)，其中 \(H\) 开于 \(X\)。由于 \(W\subseteq G\subseteq O\subseteq N\)，有
	\[
		W=H\cap O.
	\]
	因此 \(W\) 实际上也开于 \(X\)。又因 \(N\) 闭，\(X\) 中的闭包 \(\overline W^{\,X}\) 包含在 \(N\) 内，恰等于 \(\overline W^{\,N}\)。它是紧空间 \(N\) 的闭子集，因而紧，并包含在 \(G\subseteq U\) 中。取 \(V=W\) 即可。
\end{proof}

这说明在 Hausdorff 条件下，局部紧性等价于“每点都有由相对紧开集组成的邻域基”。这里的闭包必须在原空间 \(X\) 中取。只说某个较小子空间内的闭包紧，不能自动控制它在 \(X\) 中还会增加哪些边界点。

\begin{proposition}[紧集的相对紧开邻域]\label{g1:prop:compact-neighborhood-shrink}
	设 \(X\) 为局部紧 Hausdorff 空间，\(K\subseteq X\) 紧，且 \(U\subseteq X\) 是含 \(K\) 的开集。则存在开集 \(V\)，使
	\[
		K\subseteq V\subseteq\overline V\subseteq U,\qquad \overline V\text{ 紧}.
	\]
\end{proposition}
\begin{proof}
	若 \(K=\varnothing\)，取 \(V=\varnothing\)。否则，对每个 \(x\in K\)，由\cref{g1:lem:local-compact-shrink}选开邻域 \(V_x\)，使 \(\overline{V_x}\) 紧且包含在 \(U\) 中。紧性给出有限个 \(V_{x_1},\ldots,V_{x_m}\) 覆盖 \(K\)。令 \(V=\bigcup_{j=1}^mV_{x_j}\)。有限并与闭包相容，所以
	\[
		\overline V=\bigcup_{j=1}^m\overline{V_{x_j}}\subseteq U.
	\]
	右边的有限个紧集之并紧，故 \(\overline V\) 紧。
\end{proof}

这里必须先取有限子覆盖再取并。无限多个相对紧开集的并未必相对紧，例如 \(\bigcup_{m\ge1}(-m,m)=\mathbb R\)。对后续构造，命题给出了两层余地：\(K\) 在 \(V\) 的内部，而 \(\overline V\) 仍留在 \(U\) 内。若一个连续函数最终被构造成在 \(X\setminus V\) 上为零，那么其非零点集的闭包就包含在紧集 \(\overline V\) 中；这正是控制支撑所需的几何条件。如何构造这种函数，还需下一节的分离工具，不能仅由集合包含关系推出。

\ProseParagraphBreak

Euclid 空间以及它的开子集局部紧：在给定点与给定开邻域中选一个足够小的闭球即可。更一般地，若一个 Hausdorff 空间每点附近都同胚于 \(\mathbb R^n\) 的开子集，就可以把小闭球经逆坐标映射送回来；连续像保持紧性，而开球的像提供内部邻域，因此这种空间局部紧。这一论证将直接适用于后文的拓扑流形。相比之下，\(\mathbb Q\) 取实直线诱导的拓扑时并不局部紧，相关极限缺失将在习题中具体说明。

\subsection{连通、道路连通与局部道路连通}
\label{g1:subsec:0402:03}

在直线中，从一个点连续移动到另一点必须经过中间位置；但一般空间的连通性先只排除“分成两块”，并不承诺存在可以行走的参数曲线。这两个要求需要分别定义。

\begin{definition}[连通与道路连通]\label{g1:def:connected-path-connected}
	若拓扑空间 \(X\) 不能写成两个非空、不交开集的并，则称它\term[liantong]{连通}[connected]。这样的并分解若存在，称为 \(X\) 的一个分离。等价地，\(X\) 中同时开且闭的子集只有 \(\varnothing,X\)。子集的连通性总指其子空间拓扑。

	从 \(x\in X\) 到 \(y\in X\) 的一条\term[daolu]{道路}[path]是连续映射 \(\gamma:[0,1]\to X\)，满足 \(\gamma(0)=x\)、\(\gamma(1)=y\)。若任意两点间都有道路，则称 \(X\)\term[daoluliantong]{道路连通}[path-connected]。
\end{definition}

空空间按上述定义既连通又道路连通。以后谈某点所在的分支时，空间当然含有该点。开闭子集刻画是因为分离 \(X=U\sqcup V\) 同时使 \(U\) 的补集 \(V\) 开；反过来，非空真开闭子集与其补集就给出分离。

\begin{lemma}[区间的连通性]\label{g1:lem:interval-connected}
	实直线的每个区间都连通。
\end{lemma}
\begin{proof}
	若区间 \(I=A\sqcup B\) 是分离，选 \(a\in A,b\in B\)，必要时交换两者，使 \(a<b\)。因为 \(I\) 是区间，\([a,b]\subseteq I\)。令 \(c=\sup(A\cap[a,b])\)。集合 \(A\) 在 \(I\) 中既开又闭，而上确界性质表明 \(c\) 属于 \(A\cap[a,b]\) 的闭包，故 \(c\in A\)。特别地，\(c\ne b\)，因为 \(b\in B\)。于是 \(c<b\)，由 \(A\) 在 \(I\) 中开，可在 \(c\) 的右边找到仍属于 \(A\cap[a,b]\) 的点，与 \(c\) 为上确界矛盾。
\end{proof}

\begin{proposition}[连续像、闭包与相交的连通集]\label{g1:prop:connected-image-closure}
	拓扑空间中有下列结论。
	\begin{enumerate}
		\item 连通空间在连续映射下的像连通。
		\item 若 \(A\) 连通且 \(A\subseteq B\subseteq\overline A\)，则 \(B\) 连通；特别是 \(\overline A\) 连通。
		\item 若一族连通子集 \(\{C_\alpha\}\) 有公共点，则它们的并连通。
	\end{enumerate}
\end{proposition}
\begin{proof}
	若 \(f(X)\) 有分离，把两个相对开部分取原像，就会得到 \(X\) 的两个非空、不交开集，它们的并为 \(X\)，与 \(X\) 连通矛盾。这证明第一项。

	对第二项，\(A=\varnothing\) 时，\(B=\varnothing\)，结论成立。否则，假设 \(B=U\sqcup V\) 是分离。把分离限制到 \(A\)，连通性迫使 \(A\) 全在其中一侧，不妨 \(A\subseteq U\)。取 \(v\in V\)。由 \(V\) 在 \(B\) 中开，可写 \(V=B\cap O\)，其中 \(O\) 开于环境空间且含 \(v\)。但 \(v\in B\subseteq\overline A\)，由闭包的邻域刻画，\(O\cap A\ne\varnothing\)；又因 \(A\subseteq B\)，这意味着 \(V\cap A\ne\varnothing\)，矛盾。

	对第三项，固定公共点 \(p\)。若 \(\bigcup_\alpha C_\alpha=U\sqcup V\) 是分离，不妨 \(p\in U\)。每个 \(C_\alpha\) 都连通且遇到 \(U\)，因而不能再遇到 \(V\)，于是所有 \(C_\alpha\subseteq U\)，迫使 \(V=\varnothing\)，矛盾。
\end{proof}

\begin{corollary}\label{g1:cor:path-connected-connected}
	道路连通空间连通，且道路连通空间的连续像仍道路连通。
\end{corollary}
\begin{proof}
	若道路连通空间 \(X\) 有分离 \(X=U\sqcup V\)，取两侧的点及连接它们的道路 \(\gamma\)。由\cref{g1:lem:interval-connected}和上一命题，\(\gamma([0,1])\) 连通，却同时遇到分离的两侧，矛盾。若 \(f:X\to Y\) 连续，给定像中的两点，选它们在 \(X\) 中的原像，再把连接原像的道路与 \(f\) 复合，就得到像中的道路。
\end{proof}

例如 Euclid 空间中的凸集道路连通，因为 \(\gamma(t)=(1-t)x+ty\) 始终位于该集合中。连通性允许的形状更广，它不要求两点间的线段留在集合里，甚至不保证存在任何道路。下面的例子把这一区别落实到一个平面紧集。

\begin{example}[拓扑正弦曲线]\label{g1:exm:topologist-sine}
	令
	\[
		G=\{(x,\sin(1/x)):0<x\le1\},\qquad I_0=\{0\}\times[-1,1],\qquad S=G\cup I_0.
	\]
	集合 \(S\) 称为\term[tuopuzhengxianquxian]{拓扑正弦曲线}[topologist's sine curve]。集合 \(G\) 是区间 \((0,1]\) 在连续映射 \(x\mapsto(x,\sin(1/x))\) 下的像，所以连通且道路连通。其在 \(\mathbb R^2\) 中的闭包恰为 \(S\)。一方面，正横坐标处图像的极限仍由连续函数给出，横坐标为零时纵坐标只能落在 \([-1,1]\)。另一方面，对任意 \(y\in[-1,1]\)，选 \(\theta\) 使 \(\sin\theta=y\)，则对充分大的整数 \(n\)，
	\[
		x_n=(2\pi n+\theta)^{-1}>0,\qquad (x_n,\sin(1/x_n))=(x_n,y)\longrightarrow(0,y).
	\]
	因此 \(I_0\subseteq\overline G\)。由\cref{g1:prop:connected-image-closure}，\(S=\overline G\) 连通；它闭且有界，故也是平面中的紧集。

	不过，没有道路能从 \(I_0\) 到达 \(G\)。假设存在这样的道路，写为 \(\gamma(t)=(a(t),b(t))\)，取方向使 \(a(0)=0\)、\(a(1)>0\)。零点集 \(a^{-1}(\{0\})\) 是 \([0,1]\) 的非空闭子集，因而有最大值 \(t_0<1\)。于是 \(a(t)>0\) 对 \(t_0<t\le1\) 成立，在这段上必有
	\[
		b(t)=\sin(1/a(t)).
	\]
	由 \(b\) 在 \(t_0\) 处连续，可取 \(0<\delta<1-t_0\)，使 \(|b(t)-b(t_0)|<1/2\) 对 \(t_0\le t\le t_0+\delta\) 成立。然而 \(a(t_0)=0\)、\(a(t_0+\delta)>0\)，由介值性质，\(a\) 在该参数区间上取遍 \([0,a(t_0+\delta)]\)。可在其中选两个正数 \(r_+,r_-\)，满足 \(\sin(1/r_+)=1\)、\(\sin(1/r_-)=-1\)，于是 \(b\) 在该区间内既取 \(1\) 又取 \(-1\)。两个相差 \(2\) 的值不可能都距 \(b(t_0)\) 小于 \(1/2\)，矛盾。
\end{example}

\begin{figure}[htbp]
	\centering
	\input{Figures/Chapter04-SineCurve}
	\caption[拓扑正弦曲线]{拓扑正弦曲线 \(S=G\cup I_0\)。曲线 \(G\) 不断靠近竖线段 \(I_0\)，使整个集合连通；沿一条道路离开竖线段却会迫使纵坐标在任意短的参数时间内反复走遍 \([-1,1]\)，与连续性矛盾。左侧密集振荡仅画出有限部分。}
	\label{g1:fig:topologist-sine}
\end{figure}

图中的竖线段不是与曲线“隔着一个正距离”的另一块：它的每个点都是曲线的极限点。这解释了为什么集合没有开集分离；道路的障碍则来自趋近竖线段时无限振荡的纵坐标。连通与道路连通分别检测了这两件不同的事。

\begin{definition}[局部道路连通]\label{g1:def:locally-path-connected}
	若拓扑空间 \(X\) 的每个点 \(x\)，在每个给定开邻域 \(U\) 中，都有一个道路连通开邻域 \(V\) 满足 \(x\in V\subseteq U\)，则称 \(X\)\term[jubudaoluliantong]{局部道路连通}[locally path-connected]。
\end{definition}

定义要求能够任意缩小，并非仅要求每点落在某个道路连通开集中。Euclid 开集局部道路连通，因为充分小的开球凸；局部同胚于 Euclid 开集的空间也局部道路连通，只需把小球送回原空间。拓扑正弦曲线在竖线段上的任何点都没有道路连通开邻域：每个这样的开邻域都遇到 \(G\)，而上一例已证明不存在从竖线段到 \(G\) 的道路。

\subsection{分支及其在映射下的行为}
\label{g1:subsec:0402:04}

一个空间未必整体连通，此时可以按“属于同一个连通部分”把点分组。这里的“部分”应是最大的连通子集，而不能只是任意选出的一个连通邻域；否则分组会依赖选择。

\begin{definition}[连通分支与道路分支]\label{g1:def:components}
	拓扑空间 \(X\) 的极大非空连通子集称为\term[liantongfenzhi]{连通分支}[connected component]。对 \(x\in X\)，所有可用道路从 \(x\) 到达的点组成的集合称为 \(x\) 的\term[daolufenzhi]{道路分支}[path component]。
\end{definition}

反向参数 \(t\mapsto1-t\) 可以反转道路；把两条道路分别重参数化到 \([0,1/2]\)、\([1/2,1]\) 并在公共端点相接，便可拼成一条道路。拼接的连续性由闭集粘合\cref{g1:prop:gluing-lemma}保证，因为两段在 \(t=1/2\) 取值相同。再加上常值道路，“可以用道路相连”就是等价关系。每个等价类本身道路连通：连接其中两点的道路经过的任一点，都能通过取道路的一段而与这两点相连，所以整条道路都留在同一个等价类中。因此道路分支也就是极大道路连通子集。

\begin{proposition}[分支的存在与基本性质]\label{g1:prop:components}
	设 \(X\) 为拓扑空间。
	\begin{enumerate}
		\item 对每个 \(x\in X\)，包含 \(x\) 的所有连通子集之并 \(C_x\) 是唯一包含 \(x\) 的连通分支；各连通分支构成 \(X\) 的一个划分。
		\item 每个连通分支闭于 \(X\)，但不一定开。每个道路分支包含在唯一的连通分支中。
		\item 若 \(f:X\to Y\) 连续，则一个连通分支的像包含在 \(Y\) 的某个连通分支中，一个道路分支的像包含在 \(Y\) 的某个道路分支中。
	\end{enumerate}
\end{proposition}
\begin{proof}
	第一项中，所有参加并集的连通子集都含 \(x\)，所以由\cref{g1:prop:connected-image-closure}，它们的并 \(C_x\) 连通。任何包含 \(x\) 的连通子集都已在此并中，因此 \(C_x\) 极大且唯一。如果两个这样的集合相交，它们的并仍连通，极大性便使两者相等；每个点至少属于自己的 \(C_x\)，所以这些分支构成划分。

	第二项中，\(\overline{C_x}\) 由同一命题仍连通，且包含 \(C_x\)，极大性给出 \(\overline{C_x}=C_x\)，故 \(C_x\) 闭。道路分支道路连通，因而连通，所以它包含在唯一的 \(C_x\) 中。连通分支未必开，例如 \(\mathbb Q\) 中每个连通分支都是单点，却没有单点是开集，证明见习题。

	第三项中，连通分支的连续像连通，因而被包含在其中任一点所属的连通分支中。对于道路分支，用道路与 \(f\) 复合即可说明其像道路连通，故同样被包含在某个道路分支中。
\end{proof}

此处结论是“包含在”，不能改成“等于”。例如常值映射把任意连通空间送到 \(\mathbb R\) 的一个点，而 \(\mathbb R\) 的连通分支是整个直线。甚至映射本身满射，也未必把每个分支满射到目标分支，习题将给出例子。

\begin{theorem}[局部道路连通时两种分支一致]\label{g1:thm:locally-path-components}
	若 \(X\) 局部道路连通，则 \(X\) 的每个道路分支都是开集，并且道路分支与连通分支相同。因此局部道路连通空间连通，当且仅当它道路连通。
\end{theorem}
\begin{proof}
	设 \(P\) 为一个道路分支，\(x\in P\)。局部道路连通性给出 \(x\) 的道路连通开邻域 \(V\)。\(V\) 中每点都能与 \(x\) 用道路连接，故 \(V\subseteq P\)。因此 \(P\) 中每点都有包含在 \(P\) 内的开邻域，\(P\) 开。

	设 \(C\) 为包含 \(P\) 的连通分支。若 \(C\setminus P\ne\varnothing\)，因为其他道路分支也开，\(X\setminus P\) 是开集；于是 \(P=C\cap P\) 与 \(C\setminus P=C\cap(X\setminus P)\) 是 \(C\) 的两个非空、不交相对开集，其并为 \(C\)，与 \(C\) 连通矛盾。因此 \(C=P\)。对非空连通空间，唯一的连通分支就是整个空间，所以结论随之成立；空空间也符合约定。
\end{proof}

拓扑正弦曲线说明局部条件的作用：它只有一个连通分支 \(S\)，却有两个道路分支 \(G,I_0\)，其中 \(I_0\) 不是开集。流形的局部 Euclid 性质则排除了这种现象。在流形上讨论某一连通分支时，也就在讨论一个开且道路连通的部分；这将允许我们把沿道路的局部构造在同一分支内传递。有关连通性、紧性和局部紧性的进一步阅读，可参阅\cite[第4章, Connectedness, 第86页; Compactness, 第94页; Local Compactness, 第104页]{Lee2011Topological}。

% !TeX root = ../../Book1.tex
% 纲要标题：### 4.3 流形所需的分离与可数性
% 纲要位置：几何学计划纲要.md，第 252 行。
% 正文状态：初稿；尚未正式审核。

% 纲要内容开始
% 
% * Hausdorff、正则与正规空间
% * 第二可数和可数穷竭
% * 局部紧、第二可数 Hausdorff 空间的仿紧性
% * Urysohn 引理与正则第二可数空间的度量化
% 
% 纲要内容结束

\section{流形所需的分离与可数性}
\label{g1:sec:0403}

Euclid 空间有两种看似不同的便利：可以把一个闭集与外面的点隔开，也可以用可数个小球描述所有开集。前一种便利允许我们在邻域内部留出余地，后一种便利使全空间上的选择能够按序组织。它们结合起来，将产生紧集穷竭、局部有限覆盖和与原拓扑相容的距离。本节把这些结论的假设逐项分开；后面定义流形时采用 Hausdorff 和第二可数条件，正是为了能够使用这些整体构造。

\subsection{Hausdorff、正则与正规空间}
\label{g1:subsec:0403:01}

\cref{g1:def:hausdorff}中的 Hausdorff 条件要求不同的两个点有不交开邻域。但在构造函数时，我们常要让一个函数在整块闭集上取零，而在另一块闭集上取一；只分离两个点还不够。下面的条件依次扩大需要分离的对象。

\begin{definition}[分离条件]\label{g1:def:regular-normal}
	设 \(X\) 为拓扑空间。若每个单点集都闭，称 \(X\) 满足\term[t1kongjian]{\(T_1\) 条件}[\(T_1\) separation axiom]。本书采用以下约定。
	\begin{enumerate}
		\item 若 \(X\) 满足 \(T_1\) 条件，且对每个闭集 \(F\subseteq X\) 及每个 \(x\notin F\)，都存在不交开集 \(U,V\)，使 \(x\in U\)、\(F\subseteq V\)，则称它为\term[zhengzekongjian]{正则空间}[regular space]。
		\item 若 \(X\) 满足 \(T_1\) 条件，且任意两个不交闭集 \(A,B\subseteq X\) 都有不交开邻域，则称它为\term[zhengguikongjian]{正规空间}[normal space]。
	\end{enumerate}
\end{definition}

一些文献把上述分离性质与 \(T_1\) 条件分开命名，因此使用外部结果时需要核对约定。在本书的约定下，
\[
	\text{正规}\ \Longrightarrow\ \text{正则}\ \Longrightarrow\ \text{Hausdorff}\ \Longrightarrow\ T_1.
\]
前两个箭头分别把一个闭集取成单点、把两个闭集都取成单点。最后一个箭头来自：若 \(y\ne x\)，Hausdorff 条件给出一个含 \(y\) 而不含 \(x\) 的开邻域；这些邻域之并是 \(X\setminus\{x\}\)。

\begin{lemma}[分离与邻域缩小]\label{g1:lem:regular-normal-shrink}
	设 \(X\) 满足 \(T_1\) 条件。
	\begin{enumerate}
		\item \(X\) 正则，当且仅当对每个 \(x\in G\)（\(G\) 开），存在开集 \(V\)，使 \(x\in V\subseteq\overline V\subseteq G\)。
		\item \(X\) 正规，当且仅当对每个闭集 \(A\) 和包含它的开集 \(G\)，存在开集 \(V\)，使 \(A\subseteq V\subseteq\overline V\subseteq G\)。
	\end{enumerate}
\end{lemma}
\begin{proof}
	先看第二项。若 \(X\) 正规，对不交闭集 \(A\) 与 \(X\setminus G\) 取不交开邻域 \(V,W\)。因为 \(X\setminus W\) 闭且包含 \(V\)，有 \(\overline V\subseteq X\setminus W\subseteq G\)。反之，对不交闭集 \(A,B\)，把缩小性质用于 \(A\subseteq X\setminus B\)。所得 \(V\) 与 \(X\setminus\overline V\) 是分别包含 \(A,B\) 的不交开集。第一项的论证相同，只需把 \(A\) 换成单点 \(\{x\}\)。
\end{proof}

把闭包也留在原邻域内的缩小方式，为取极限、控制支集与零延拓留出了余地。\cref{g1:lem:local-compact-shrink}已经说明，局部紧 Hausdorff 空间具有这种点的邻域缩小性质，而且可以要求闭包紧。因此局部紧 Hausdorff 空间正则；此处尚未要求它第二可数。

\begin{example}[度量空间中的分离函数]\label{g1:exm:metric-normal}
	每个度量空间 \((X,d)\) 都正规。单点集闭；若 \(A,B\) 是非空不交闭集，令
	\[
		d(x,A)=\inf_{a\in A}d(x,a),\qquad
		h(x)=\frac{d(x,A)}{d(x,A)+d(x,B)}.
	\]
	三角不等式给出 \(|d(x,A)-d(y,A)|\le d(x,y)\)，故两个距离函数连续。又因 \(A,B\) 闭且不交，它们的距离函数不可能在同一点同时为零，所以分母处处为正。于是 \(h\) 连续，在 \(A\) 上为零，在 \(B\) 上为一；开集 \(\{h<1/3\}\) 与 \(\{h>2/3\}\) 分离 \(A,B\)。若有一个闭集为空，直接取空集和全空间即可。
\end{example}

这里没有要求两个闭集之间的距离为正。例如平面中 \(A=\{(t,0):t\ge1\}\) 与 \(B=\{(t,1/t):t\ge1\}\) 都闭，二者距离的下确界却为零。上例只需要每个固定点处的分母为正，并不需要全空间上一致的正下界。

\subsection{第二可数和可数穷竭}
\label{g1:subsec:0403:02}

为了在不同地方选择局部数据，我们希望这些选择能够列成一张可数清单。局部与整体的可数性具有不同的量词：“每一点有可数个基本邻域”允许不同的点各自使用一张清单；“全空间有可数个基本开集”要求一张清单适用于所有点。

\begin{definition}[可数性条件]\label{g1:def:countability}
	设 \(X\) 为拓扑空间。若每一点 \(x\) 都有可数邻域基，即存在含 \(x\) 的开集 \(U_1,U_2,\ldots\)，使每个含 \(x\) 的开集都包含某个 \(U_n\)，则称该空间\term[diyikeshu]{第一可数}[first countable]。若 \(X\) 的整个拓扑有一个可数基，则称它\term[dierkeshu]{第二可数}[second countable]。若 \(X\) 有可数稠密子集，称它\term[kefenkongjian]{可分}[separable]。若每个开覆盖都有可数子覆盖，称它具有\term[lindelofxingzhi]{Lindelöf 性质}[Lindelöf property]。\footnote{Ernst Leonard Lindelöf（1870--1946），芬兰数学家，研究函数论与微分方程。}
\end{definition}

可数在这里包括有限。第二可数立即蕴含第一可数：只取全空间的基中那些包含给定点的成员。每个度量空间都第一可数，因为 \(B(x,1/n)\) 构成 \(x\) 的邻域基；但一个度量空间不一定第二可数。若一个离散空间不可数，其每个单点都是开集，任何拓扑基都必须包含所有单点。

\begin{proposition}[第二可数的两项后果]\label{g1:prop:second-countable-lindelof}
	第二可数空间是可分空间，并具有 Lindelöf 性质。第二可数性遗传给任意子空间；Lindelöf 性质遗传给闭子空间。
\end{proposition}
\begin{proof}
	设 \(\{B_n\}\) 是可数基，忽略空的基元素。在每个 \(B_n\) 中取一点 \(x_n\)，所得可数点集与每个非空开集相交，因为每个这样的开集都包含某个 \(B_n\)。所以它稠密。

	给定开覆盖 \(\mathcal U\)，对每个能够包含在某个覆盖成员中的 \(B_n\)，选定一个这样的 \(U_n\in\mathcal U\)。所得至多可数个 \(U_n\) 覆盖 \(X\)：若 \(x\in U\in\mathcal U\)，基的性质给出 \(x\in B_n\subseteq U\)，于是选定的 \(U_n\) 也包含 \(x\)。

	若 \(Y\subseteq X\)，则 \(\{B_n\cap Y\}\) 是子空间的可数基。最后设 \(F\) 是 Lindelöf 空间 \(X\) 的闭子集，且 \(\{F\cap U_\alpha\}\) 是 \(F\) 的相对开覆盖，其中 \(U_\alpha\) 在 \(X\) 中开。把 \(X\setminus F\) 加入 \(\{U_\alpha\}\) 得到 \(X\) 的开覆盖，取可数子覆盖再与 \(F\) 相交即可。
\end{proof}

在 Euclid 空间中，以有理坐标点为中心、正有理数为半径的开球构成可数基。其关键不只是中心稠密，还在于可以把半径取得足够小：给定 \(x\in U\)，先取 \(B(x,\varepsilon)\subseteq U\)，再取靠近 \(x\) 的有理点 \(q\) 和满足 \(\lVert x-q\rVert<r<\varepsilon-\lVert x-q\rVert\) 的正有理数 \(r\)，便有 \(x\in B(q,r)\subseteq U\)。

\begin{definition}[紧穷竭]\label{g1:def:compact-exhaustion}
	设 \(X\) 为拓扑空间。一列紧集 \((K_n)_{n\ge1}\) 若满足
	\[
		K_n\subseteq\operatorname{int}K_{n+1},\qquad
		X=\bigcup_{n\ge1}K_n,
	\]
	称为 \(X\) 的一个\term[jinqiongjie]{紧穷竭}[compact exhaustion]。内部和闭包均在 \(X\) 中计算。
\end{definition}

\cref{g1:def:locally-compact}中的相对紧是指子集在环境空间中的闭包紧。下面会先得到由相对紧开集组成的可数基，再把这些开集分批装进紧穷竭。

\begin{theorem}[可数相对紧基与紧穷竭]\label{g1:thm:compact-exhaustion}
	局部紧、第二可数的 Hausdorff 空间 \(X\) 有可数相对紧开基，并有紧穷竭。对这样的穷竭，每个紧子集 \(C\subseteq X\) 都包含于某个 \(\operatorname{int}K_N\)。
\end{theorem}
\begin{proof}
	空空间取所有 \(K_n=\varnothing\)。以下设 \(X\ne\varnothing\)。从一个可数开基出发，只保留相对紧的基元素。它们仍是基：给定 \(x\in U\) 且 \(U\) 开，\cref{g1:lem:local-compact-shrink}给出开集 \(V\)，满足 \(x\in V\)、\(\overline V\) 紧且 \(\overline V\subseteq U\)。再选原来的基元素 \(B\) 使 \(x\in B\subseteq V\)，则 \(\overline B\) 是紧集 \(\overline V\) 的闭子集，因而紧。这里闭包包含于 \(\overline V\) 的事实也保证了 \(B\subseteq U\)。

	把保留下来的非空基元素列成 \(B_1,B_2,\ldots\)；若只有有限个，可以重复列出。令 \(K_1=\overline B_1\)。已构造 \(K_n\) 后，紧集 \(K_n\cup\overline B_{n+1}\) 可由有限个基元素 \(B_{j_1},\ldots,B_{j_s}\) 覆盖。定义
	\[
		K_{n+1}=\overline B_{j_1}\cup\cdots\cup\overline B_{j_s}.
	\]
	这是有限个紧集之并，仍紧，而且
	\[
		K_n\cup\overline B_{n+1}
		\subseteq B_{j_1}\cup\cdots\cup B_{j_s}
		\subseteq\operatorname{int}K_{n+1}.
	\]
	因此所需的内部包含关系成立，各 \(B_n\) 也都被某个 \(K_n\) 包含，故这些紧集穷尽 \(X\)。

	最后，由 \(K_n\subseteq\operatorname{int}K_{n+1}\)，开集 \(\operatorname{int}K_n\) 覆盖 \(X\)。它们递增，所以对紧集 \(C\) 取有限子覆盖后，最大的下标 \(N\) 已满足 \(C\subseteq\operatorname{int}K_N\)。
\end{proof}

紧穷竭除了满足 \(X=\bigcup_nK_n\)，还通过相邻紧集的内部包含关系留出余地。内部包含关系使 \(K_n\) 与 \(X\setminus K_{n+1}\) 之间隔着一层开邻域；下一小节将利用这层余地，让越来越远处选择的开集不再影响一个固定点附近。

\subsection{局部紧、第二可数 Hausdorff 空间的仿紧性}
\label{g1:subsec:0403:03}

在一个非紧空间上，开覆盖通常不能缩减为有限覆盖。然而粘合函数时，全空间上只有有限项并非必需：只要每个点附近实际上只有有限项出现，求和的连续性便仍是有限和的连续性。这就是局部有限条件的作用。

\begin{definition}[细化、局部有限与仿紧]\label{g1:def:locally-finite-refinement}
	设 \(\mathcal U\) 是 \(X\) 的覆盖，\(\mathcal V\) 是另一个覆盖。若每个 \(V\in\mathcal V\) 都包含于某个 \(U\in\mathcal U\)，则称新覆盖为原覆盖的\term[xihua]{细化}[refinement]。对于一族子集 \(\{A_i\}_{i\in I}\)，若每个 \(x\in X\) 都有一个只与有限个 \(A_i\) 相交的开邻域，则称这族集合是\term[jubuyouxian]{局部有限}[locally finite]的。若 \(X\) 的每个开覆盖都有局部有限开细化，称 \(X\) 为\term[fangjinkongjian]{仿紧空间}[paracompact space]。本书的仿紧性定义本身不包含 Hausdorff 条件。
\end{definition}

细化可以增加覆盖成员的数目；它要求新开集更贴合原覆盖，并不要求新覆盖是原覆盖的子族。局部有限也不同于“每个点只属于有限个成员”：后一条件只控制点本身，前一条件还控制它附近的所有点。例如 \(\{(1/(n+1),1/n):n\ge1\}\) 在 \(\mathbb R\) 中每点至多属于一个成员，却在 \(0\) 附近不是局部有限的。

\begin{lemma}[局部有限与闭包]\label{g1:lem:locally-finite-closures}
	若 \(\{A_i\}_{i\in I}\) 是拓扑空间 \(X\) 中的局部有限族，则 \(\{\overline A_i\}_{i\in I}\) 仍局部有限，并且
	\[
		\overline{\bigcup_{i\in I}A_i}=\bigcup_{i\in I}\overline A_i.
	\]
	特别地，局部有限闭集族的并是闭集；其任意子族的并也闭。
\end{lemma}
\begin{proof}
	若开集 \(N\) 与 \(\overline A_i\) 相交，取交点后利用闭包定义可知 \(N\) 也与 \(A_i\) 相交。因此一个只与有限个 \(A_i\) 相交的开邻域，也只与对应的有限个闭包相交。

	为证明右边的并闭，取 \(x\notin\bigcup_i\overline A_i\)，并取只与有限个闭包 \(\overline A_{i_1},\ldots,\overline A_{i_s}\) 相交的开邻域 \(N\)。开集 \(N\setminus\bigcup_{k=1}^s\overline A_{i_k}\) 包含 \(x\)，且与所有闭包都不交。故右边是包含 \(\bigcup_iA_i\) 的闭集，包含其闭包。反向包含逐项由 \(A_i\subseteq\bigcup_jA_j\) 得到。子族仍局部有限，最后一句随即成立。
\end{proof}

\begin{theorem}[由紧穷竭得到局部有限细化]\label{g1:thm:lch-paracompact}
	设 \(X\) 为局部紧、第二可数的 Hausdorff 空间。对 \(X\) 的每个开覆盖 \(\mathcal U\)，存在可数的局部有限开覆盖 \(\{V_j\}\)，使每个 \(\overline V_j\) 都紧，并包含于某个 \(U\in\mathcal U\)。特别地，\(X\) 仿紧。
\end{theorem}
\begin{proof}
	由\cref{g1:thm:compact-exhaustion}取紧穷竭 \((K_n)\)，并约定 \(K_{-1}=K_0=\varnothing\)。把 \(X\) 分成相互允许重叠的紧壳层
	\[
		S_n=K_n\setminus\operatorname{int}K_{n-1},
		\qquad
		G_n=\operatorname{int}K_{n+1}\setminus K_{n-2}.
	\]
	每个 \(S_n\) 紧，每个 \(G_n\) 开，而且 \(S_n\subseteq G_n\)：外侧包含来自 \(K_n\subseteq\operatorname{int}K_{n+1}\)，内侧避开 \(K_{n-2}\) 则来自 \(K_{n-2}\subseteq\operatorname{int}K_{n-1}\)。这些 \(S_n\) 覆盖 \(X\)，因为每个点第一次进入某个 \(K_n\) 时就属于 \(S_n\)。

	对每个 \(x\in S_n\)，选原覆盖成员 \(U_x\) 包含 \(x\)，再由\cref{g1:lem:local-compact-shrink}在开集 \(G_n\cap U_x\) 中缩小邻域，得到
	\[
		x\in V_x,\qquad
		\overline V_x\text{ 紧},\qquad
		\overline V_x\subseteq G_n\cap U_x.
	\]
	利用 \(S_n\) 紧，为这一层保留有限个 \(V_{n,1},\ldots,V_{n,r_n}\)。空壳层不需选择成员。所有层的选择合起来是可数开细化，并覆盖 \(X\)。

	剩下要证明的是局部有限；每层有限还不足以直接给出这个结论。给定 \(x\in X\)，选 \(m\) 使 \(x\in\operatorname{int}K_m\)。当 \(n\ge m+2\) 时，\(K_m\subseteq K_{n-2}\)，而这一层的全部 \(V_{n,k}\) 都位于 \(G_n\subseteq X\setminus K_{n-2}\)，故它们与 \(\operatorname{int}K_m\) 不交。于是该邻域只可能遇到前 \(m+1\) 层中的有限个成员，这正是局部有限。
\end{proof}

在这个构造中，\(S_n\) 是需要覆盖的区域，\(G_n\) 是允许选择邻域的稍大区域。让第 \(n\) 层避开 \(K_{n-2}\)，而不只要求它留在 \(K_{n+1}\) 以内，保证了远处的开集不会重新伸回固定点附近。

\begin{figure}[htbp]
	\centering
	\input{Figures/Chapter04-Exhaustion.tex}
	\caption{蓝色壳层 \(S_n\) 位于灰色区域 \(G_n\) 内，内外各留一层余地。曲线仅表示包含关系，并未假定空间嵌入平面。}
	\label{g1:fig:compact-exhaustion-shells}
\end{figure}

\begin{theorem}[仿紧 Hausdorff 空间的正规性]\label{g1:thm:paracompact-normal}
	每个仿紧 Hausdorff 空间都是正规空间。
\end{theorem}
\begin{proof}
	先证明正则性，不能在这一步预先使用闭包缩小。给定闭集 \(F\) 和 \(x\notin F\)，对每个 \(y\in F\)，由 Hausdorff 条件取分别含 \(x,y\) 的不交开集 \(P_y,Q_y\)。于是 \(x\notin\overline{Q_y}\)。对开覆盖
	\[
		\{X\setminus F\}\cup\{Q_y:y\in F\}
	\]
	取局部有限开细化 \(\{W_i\}_{i\in I}\)，并令 \(I_F=\{i:W_i\cap F\ne\varnothing\}\)。若 \(i\in I_F\)，则 \(W_i\) 不可能包含于 \(X\setminus F\)，所以必包含于某个 \(Q_y\)，从而 \(x\notin\overline W_i\)。由\cref{g1:lem:locally-finite-closures}，
	\[
		H=\bigcup_{i\in I_F}\overline W_i
	\]
	是一个不含 \(x\) 的闭集。于是 \(X\setminus H\) 与 \(\bigcup_{i\in I_F}W_i\) 是分别包含 \(x,F\) 的不交开集。Hausdorff 条件已保证 \(T_1\)，故 \(X\) 正则。

	现在设 \(A,B\) 是不交闭集。由刚证的正则性及\cref{g1:lem:regular-normal-shrink}，对每个 \(a\in A\) 取开集 \(O_a\)，满足 \(a\in O_a\)、\(\overline O_a\subseteq X\setminus B\)。对覆盖
	\[
		\{X\setminus A\}\cup\{O_a:a\in A\}
	\]
	再取局部有限开细化 \(\{V_j\}_{j\in J}\)。保留那些与 \(A\) 相交的成员，其下标集记为 \(J_A\)。每个这样的 \(V_j\) 必包含于某个 \(O_a\)，所以它的闭包避开 \(B\)。局部有限性保证 \(F_A=\bigcup_{j\in J_A}\overline V_j\) 闭，且仍避开 \(B\)。因此 \(\bigcup_{j\in J_A}V_j\) 与 \(X\setminus F_A\) 是 \(A,B\) 的不交开邻域，得到正规性。
\end{proof}

紧空间的每个开覆盖都有有限子覆盖，而有限族总是局部有限的。因此紧 Hausdorff 空间正规；这里紧性提供有限子覆盖，Hausdorff 条件负责最初的点分离。对流形将使用的局部紧、第二可数 Hausdorff 空间，上面两条定理则把“局部紧”转化成了整个空间上的正规性。

\subsection{Urysohn 引理与正则第二可数空间的度量化}
\label{g1:subsec:0403:04}

如果拓扑已经来自距离，\cref{g1:exm:metric-normal}用距离函数构造了取值在 \([0,1]\) 中的连续分离函数。现在反过来：先仅用开集与闭集构造足够多的这样的函数，再把函数值的差组合成距离。需要解决两个问题：如何在没有距离时构造连续函数，以及如何把全空间的分离任务缩减到可数个函数。

\begin{lemma}[正则 Lindelöf 空间的正规性]\label{g1:lem:regular-lindelof-normal}
	每个正则且具有 Lindelöf 性质的空间都是正规空间。
\end{lemma}
\begin{proof}
	设 \(A,B\) 为不交闭集。由正则性，为 \(A\) 的各点取闭包避开 \(B\) 的开邻域。闭子空间 \(A\) 仍有 Lindelöf 性质，故可取可数个这样的开集 \(U_n\) 覆盖 \(A\)。同样可取开集 \(V_n\) 覆盖 \(B\)，使 \(\overline V_n\cap A=\varnothing\)。若某一族有限或为空，可补空集写成数列。令
	\[
		U=\bigcup_{n\ge1}\left(U_n\setminus\bigcup_{i=1}^n\overline V_i\right),
		\qquad
		V=\bigcup_{n\ge1}\left(V_n\setminus\bigcup_{i=1}^n\overline U_i\right).
	\]
	括号中只减去有限个闭集，因而 \(U,V\) 都开。因为所有 \(\overline V_i\) 避开 \(A\)，\(U\) 仍包含 \(A\)；同理 \(V\) 包含 \(B\)。若某点同时落在左边第 \(n\) 项与右边第 \(m\) 项中，当 \(m\le n\) 时，左边已减去 \(\overline V_m\)，不可能属于 \(V_m\)；当 \(n\le m\) 时，右边已减去 \(\overline U_n\)，同样矛盾。因此 \(U\cap V=\varnothing\)。
\end{proof}

这个交错删去有限个闭集的构造解释了可数性的用途：第 \(n\) 个开集只需避开对方的前 \(n\) 个闭包，而任意一对下标终究有一个先后关系。若直接从每个 \(U_n\) 中减去对方所有闭包，剩下的集合未必开。

\begin{lemma}[Urysohn 引理]\label{g1:lem:urysohn}
	设 \(X\) 为正规空间，\(A,B\subseteq X\) 为不交闭集。则存在连续函数 \(f:X\to[0,1]\)，使 \(f|_A=0\)、\(f|_B=1\)。
\end{lemma}

这一结果称为\term[urysohnyinli]{Urysohn 引理}[Urysohn lemma]。\footnote{Pavel Samuilovich Urysohn（1898--1924），苏联数学家，研究一般拓扑与维数理论。}构造从安排函数的“低于某一高度的区域”入手。为保证函数连续，不同高度之间需要留出闭包包含关系。

\begin{proof}
	令 \(D=\{k/2^m:0\le k\le2^m,\ m\ge0\}\subseteq[0,1]\) 为二进有理数集。我们要为每个 \(r\in D\) 构造开集 \(U_r\)，使
	\[
		A\subseteq U_0,\qquad U_1=X\setminus B,\qquad
		\overline U_r\subseteq U_s\quad(r<s,\ r,s\in D).
	\]
	先由\cref{g1:lem:regular-normal-shrink}，为闭集 \(A\subseteq X\setminus B\) 取开邻域 \(U_0\)，满足 \(\overline U_0\subseteq U_1=X\setminus B\)。假设分母为 \(2^m\) 的各开集已经构造。对相邻的 \(r=k/2^m\) 与 \(s=(k+1)/2^m\)，在闭集 \(\overline U_r\subseteq U_s\) 之间再缩小一次，得到
	\[
		\overline U_r\subseteq U_{(r+s)/2}
		\subseteq\overline U_{(r+s)/2}\subseteq U_s.
	\]
	每一层只加入这些中间开集，旧的开集保持不变。任意两个二进有理数都出现在某个共同的有限层中，逐个使用相邻包含关系，便得到全部所需的嵌套关系。

	定义
	\[
		f(x)=\inf\bigl(\{r\in D:x\in U_r\}\cup\{1\}\bigr).
	\]
	加入 \(\{1\}\) 使这个集合总非空，所以 \(f\) 在每一点都有定义并属于 \([0,1]\)。若 \(x\in A\)，则 \(x\in U_0\)，故 \(f(x)=0\)；若 \(x\in B\)，则它不属于任何 \(U_r\)，故 \(f(x)=1\)。

	下面分别检验函数值低于和高于给定高度的集合开。对 \(0<a\le1\)，下确界的性质给出
	\[
		\{f<a\}=\bigcup_{\substack{r\in D\\r<a}}U_r.
	\]
	对 \(0\le a<1\)，则有
	\[
		\{f>a\}=\bigcup_{\substack{r\in D\\a<r<1}}(X\setminus\overline U_r).
	\]
	为核对第二式，若 \(x\notin\overline U_r\)，则它不属于任何 \(U_t\)（\(t<r\)），所以 \(f(x)\ge r>a\)。反之，若 \(f(x)>a\)，取二进有理数 \(a<r<f(x)\)。若 \(x\in\overline U_r\)，则对每个二进有理数 \(s>r\) 都有 \(x\in U_s\)，从而 \(f(x)\le r\)，矛盾。故 \(x\notin\overline U_r\)。这证明第二式。上述两个集合都开；当 \(a\) 在其余范围时，它们是空集或全空间。实数开区间的原像由这两类集合相交得到，故 \(f\) 连续。
\end{proof}

若映射 \(F:X\to Y\) 是 \(X\) 到其像 \(F(X)\) 的同胚，其中像赋予 \(Y\) 的子空间拓扑，就称 \(F\) 为\term[tuopuqianru]{拓扑嵌入}[topological embedding]。下面将用可数个连续函数把原空间实现为一个可数积空间的子空间。

\begin{theorem}[Urysohn 度量化定理]\label{g1:thm:urysohn-metrization}
	设 \(X\) 为正则、第二可数的拓扑空间。则存在一列连续函数 \(f_n:X\to[0,1]\)，使
	\begin{equation}\label{g1:eq:urysohn-metric}
		d(x,y)=\sum_{n=1}^{\infty}2^{-n}|f_n(x)-f_n(y)|
	\end{equation}
	成为 \(X\) 上的度量，并且它所诱导的拓扑恰好等于原拓扑。映射 \(x\mapsto(f_n(x))_{n\ge1}\) 将 \(X\) 拓扑嵌入带积拓扑的 \([0,1]^{\mathbb N}\)。
\end{theorem}

若一个空间的拓扑可由某个距离诱导，就称这个空间\term[kelianghua]{可度量化}[metrizable]；并不预先指定距离，也不保证这个距离完备。上述结果称为\term[urysohndulianghua]{Urysohn 度量化定理}[Urysohn metrization theorem]。这里正则性含 \(T_1\) 条件。定理比“局部紧、第二可数 Hausdorff 空间可度量化”更一般，因为它不要求局部紧。

\begin{proof}
	空空间的结论按空映射成立。设 \(X\ne\varnothing\)，取可数开基 \(\{B_i\}\)。由\cref{g1:prop:second-countable-lindelof,g1:lem:regular-lindelof-normal}，\(X\) 正规。对每一对满足 \(\overline B_i\subseteq B_j\) 的下标 \((i,j)\)，把\cref{g1:lem:urysohn}用于不交闭集 \(X\setminus B_j\) 与 \(\overline B_i\)，得到连续函数 \(f_{ij}:X\to[0,1]\)，满足
	\[
		f_{ij}=1\text{ 于 }\overline B_i,\qquad
		f_{ij}=0\text{ 于 }X\setminus B_j.
	\]
	这些下标对至多可数。列出全部所得函数，并在必要时补上常值零函数，记成 \((f_n)_{n\ge1}\)。

	这些函数具有比“分开两个点”更强的性质：若 \(x\in O\) 且 \(O\) 开，可以找到一个 \(f_k\)，使 \(f_k(x)=1\)，并在 \(X\setminus O\) 上为零。事实上，先取基元素 \(x\in B_j\subseteq O\)，再由正则性选 \(x\in V\subseteq\overline V\subseteq B_j\)，最后取 \(x\in B_i\subseteq V\)。于是 \(\overline B_i\subseteq B_j\)，所对应的 \(f_{ij}\) 就满足要求。特别地，对 \(x\ne y\)，因 \(T_1\) 条件使 \(X\setminus\{y\}\) 开，这一性质给出某个 \(f_k(x)=1\)、\(f_k(y)=0\)。

	现在用\eqref{g1:eq:urysohn-metric}定义 \(d\)。每项不超过 \(2^{-n}\)，故级数收敛。对称性与非负性逐项成立；刚证的点分离性质保证 \(x\ne y\) 时 \(d(x,y)>0\)。每个函数值差满足实数的三角不等式，把前 \(N\) 项相加再令 \(N\to\infty\)，得到 \(d(x,z)\le d(x,y)+d(y,z)\)。所以 \(d\) 是度量。

	还必须证明原来的开集与度量开集相同，单有点分离并不能保证这一点。给定 \(x\in X\) 与 \(\varepsilon>0\)，取 \(N\) 使 \(\sum_{n>N}2^{-n}<\varepsilon/2\)。由前 \(N\) 个函数的连续性，存在原拓扑中的开邻域 \(V\ni x\)，使对 \(y\in V\) 及 \(1\le n\le N\)，都有 \(|f_n(x)-f_n(y)|<\varepsilon/2\)。于是前 \(N\) 项之和小于 \(\varepsilon/2\)，尾项之和也小于 \(\varepsilon/2\)，故 \(V\subseteq B_d(x,\varepsilon)\)。这说明从原拓扑到度量拓扑的恒等映射连续。

	反过来，给定原拓扑的开邻域 \(O\ni x\)，取上述满足 \(f_k(x)=1\) 且在 \(X\setminus O\) 上为零的函数。如果 \(d(x,y)<2^{-k-1}\)，则
	\[
		2^{-k}|f_k(x)-f_k(y)|\le d(x,y)<2^{-k-1},
	\]
	所以 \(f_k(y)>1/2\)，从而 \(y\in O\)。因此 \(B_d(x,2^{-k-1})\subseteq O\)，另一方向的恒等映射也连续。

	最后，令 \(F(x)=(f_n(x))\)。它由点分离性质单射，而在立方体 \([0,1]^{\mathbb N}\) 上，
	\[
		D(s,t)=\sum_{n\ge1}2^{-n}|s_n-t_n|
	\]
	是一个度量，且 \(D(F(x),F(y))=d(x,y)\)。这个度量给出积拓扑：其开球包含由有限个坐标的小变化所确定的基本开邻域，理由正是上面的尾项估计；反过来，对规定有限个坐标的邻域，取球半径小于各坐标容许变化量乘以对应 \(2^{-n}\) 的最小值，整个球便包含于该邻域。因此 \(F\) 是 \(X\) 到其像的同胚，即所称的拓扑嵌入。
\end{proof}

可以把 \((f_n(x))\) 看作一列连续测量值。正则性与 Urysohn 引理保证：无论给点 \(x\) 指定多小的开邻域，总有一项测量在 \(x\) 取一、在邻域外取零。第二可数性把所有这种局部区分压缩为可数个测量；权重 \(2^{-n}\) 则让后面无穷多项的总影响可以统一控制。

\begin{corollary}[流形所用的拓扑结论]\label{g1:cor:lch-second-countable}
	每个局部紧、第二可数的 Hausdorff 空间都可度量化、仿紧且正规，并具有紧穷竭。
\end{corollary}
\begin{proof}
	局部紧 Hausdorff 空间由邻域缩小性质正则，因此可应用\cref{g1:thm:urysohn-metrization}。紧穷竭、仿紧性与正规性分别由\cref{g1:thm:compact-exhaustion,g1:thm:lch-paracompact,g1:thm:paracompact-normal}得到。
\end{proof}

这个推论给出与开集相容的点间距离，式\eqref{g1:eq:urysohn-metric}中的距离未必完备。Riemann 度量则在每个切空间上指定随点光滑变化的内积；下一节的连续单位分解将先说明这类局部数据怎样加权粘合。

% !TeX root = ../../Book1.tex
% 纲要标题：### 4.4 局部有限覆盖与单位分解
% 纲要位置：几何学计划纲要.md，第 259 行。
% 正文状态：初稿；尚未正式审核。

% 纲要内容开始
% 
% * 覆盖的细化与闭集收缩
% * 连续截断函数的构造
% * 从局部有限覆盖构造单位分解
% * 连续函数与正定双线性型的整体粘合
% 
% 纲要内容结束

\section{局部有限覆盖与单位分解}
\label{g1:sec:0404}

在一条直线的左侧使用一种公式、右侧使用另一种公式，如果两式在交叠处相同，就可以直接粘合。若它们各有用途，却不完全相同，又怎样平稳地过渡？一个自然想法是给两式分配非负权重，并使权重之和为 \(1\)。对于许多张局部坐标图，同样的想法要求我们同时解决三个问题：每个权重只在允许的区域起作用，各点附近只出现有限项，而且任一点至少有一个权重为正。本节用覆盖的收缩与连续截断函数逐项实现这些要求。

\subsection{覆盖的细化与闭集收缩}
\label{g1:subsec:0404:01}

\cref{g1:def:locally-finite-refinement}中的局部有限性说，每一点都有一个邻域，只与给定族中的有限多个集合相交。它比“每个点只属于有限多个集合”更强：前者控制了点附近的变化，后者只控制某一点上的计数。在 \(\mathbb R\) 中，令 \(c_n=2^{-n}\)、\(r_n=2^{-n-2}\)，闭区间 \([c_n-r_n,c_n+r_n]\) 两两不交，故每点至多属于一个，但原点的每个邻域都遇到无穷多个。这种聚集会使看似逐点有意义的和失去连续性。

已由\cref{g1:lem:locally-finite-closures}证明，局部有限族取闭包后仍局部有限，且局部有限闭集族的并为闭集。这里闭包结论中的“邻域”尤其有用：一个开集若不交某集合，也不交它的闭包。这使我们可以在同一个邻域内同时控制函数的非零区域与这些区域的边界。

\begin{lemma}[带闭包控制的细化与收缩]\label{g1:lem:cover-shrinking}
	设 \(X\) 为仿紧 Hausdorff 空间。
	\begin{enumerate}
		\item 对任意开覆盖 \(\{U_\alpha\}_{\alpha\in A}\)，存在局部有限开覆盖 \(\{V_j\}_{j\in J}\) 及指标选择 \(\alpha(j)\)，使 \(\overline{V_j}\subseteq U_{\alpha(j)}\)。
		\item 对任意局部有限开覆盖 \(\{V_j\}_{j\in J}\)，存在同指标闭集族 \(\{F_j\}_{j\in J}\)，满足 \(F_j\subseteq V_j\)、\(\bigcup_jF_j=X\)。这族闭集仍局部有限，允许某些 \(F_j\) 为空。
	\end{enumerate}
\end{lemma}
\begin{proof}
	由\cref{g1:thm:paracompact-normal}，\(X\) 正规，因而也正则。对每个 \(x\in X\)，先选包含它的一个 \(U_{\alpha_x}\)，再用正则性选开邻域 \(W_x\)，使
	\[
		x\in W_x\subseteq\overline{W_x}\subseteq U_{\alpha_x}.
	\]
	仿紧性给出开覆盖 \(\{W_x\}\) 的一个局部有限开细化 \(\{V_j\}\)。对每个 \(j\) 选 \(x(j)\) 使 \(V_j\subseteq W_{x(j)}\)，便有 \(\overline{V_j}\subseteq\overline{W_{x(j)}}\subseteq U_{\alpha_{x(j)}}\)，证明第一项。

	现从第二项给定的局部有限开覆盖 \(\{V_j\}\) 出发，把第一项用于它，得到局部有限开覆盖 \(\{W_k\}_{k\in K}\)，并为每个 \(k\) 选 \(j(k)\)，使 \(\overline{W_k}\subseteq V_{j(k)}\)。按被分配到的原指标合并，令
	\[
		F_j=\bigcup_{k:\,j(k)=j}\overline{W_k}.
	\]
	闭包族局部有限，故其中任何子族的并都闭；于是每个 \(F_j\) 是闭集，且 \(F_j\subseteq V_j\)。因为 \(W_k\) 覆盖 \(X\)，这些 \(F_j\) 也覆盖 \(X\)。最后，任何只遇到有限多个 \(V_j\) 的邻域，也只会遇到有限多个 \(F_j\)，故闭集族局部有限。
\end{proof}

第二项称为原覆盖的一个\term[bisousuo]{闭收缩}[closed shrinking]。这里既要 \(F_j\) 留在 \(V_j\) 内，也要所有 \(F_j\) 仍覆盖空间。任意把每个开集缩小并不能保证后一点；证明通过另一个仍覆盖空间的细化，再将闭包按原指标合并，保住了覆盖性。

\subsection{连续截断函数的构造}
\label{g1:subsec:0404:02}

覆盖缩小以后，希望给闭集 \(F_j\) 上的点权重 \(1\)，并在 \(V_j\) 外将权重关掉。只规定函数在外面为零还不够：它的非零点可能一直逼近不允许进入的边界，因此需要控制非零区域的闭包。

\begin{definition}[支集与截断函数]\label{g1:def:support-cutoff}
	连续函数 \(f:X\rightarrow\mathbb R\) 的\term[zhiji]{支集}[support]定义为
	\[
		\operatorname{supp}f=\overline{\{x\in X:f(x)\ne0\}}.
	\]
	给定 \(A\subseteq U\subseteq X\)，若连续函数 \(\rho:X\rightarrow[0,1]\) 满足 \(\rho=1\) 于 \(A\)，且 \(\operatorname{supp}\rho\subseteq U\)，则称它为在 \(A\) 上等于 \(1\)、支集位于 \(U\) 内的\term[jieduanhanshu]{截断函数}[cutoff function]。
	\symidx{\(\operatorname{supp}f\)：函数的支集}
\end{definition}

例如 \(f(x)=\max\{1-|x|,0\}\) 的非零点集是 \((-1,1)\)，支集却是 \([-1,1]\)。所以“在 \((-1,1)\) 外等于零”并不能推出“支集包含在 \((-1,1)\) 内”。构造截断函数时，需要在闭集与允许区域的边界之间再留一层空间。

\begin{proposition}[闭集截断与紧支集截断]\label{g1:prop:continuous-cutoff}
	以下两种情形均存在连续函数 \(\rho:X\to[0,1]\)，使 \(\rho|_A=1\)、\(\operatorname{supp}\rho\subseteq U\)；第二种情形以 \(K\) 代替 \(A\)，并可使支集紧。
	\begin{enumerate}
		\item \(X\) 正规，\(A\) 为闭集，\(U\) 为包含 \(A\) 的开集。
		\item \(X\) 局部紧且 Hausdorff，\(K\subseteq U\)，其中 \(K\) 紧、\(U\) 开。此时可以令 \(\rho=1\) 于 \(K\)，并使 \(\operatorname{supp}\rho\) 为包含在 \(U\) 内的紧集。
	\end{enumerate}
\end{proposition}
\begin{proof}
	第一种情形中，由正规性的闭邻域缩小形式，取开集 \(W\)，使
	\[
		A\subseteq W\subseteq\overline W\subseteq U.
	\]
	对不交闭集 \(A\) 与 \(X\setminus W\) 应用\cref{g1:lem:urysohn}，得到连续 \(\rho:X\rightarrow[0,1]\)，在 \(A\) 上为 \(1\)，在 \(X\setminus W\) 上为 \(0\)。于是非零点都在 \(W\) 内，支集包含于 \(\overline W\)，进而包含于 \(U\)。

	第二种情形若 \(K=\varnothing\)，取零函数即可。否则两次使用\cref{g1:prop:compact-neighborhood-shrink}，得到开集 \(V,W\)，使
	\[
		K\subseteq V\subseteq\overline V\subseteq W\subseteq\overline W\subseteq U,
	\]
	且 \(\overline V,\overline W\) 紧。令 \(C=\overline W\)。这是紧 Hausdorff 空间，又因每个开覆盖有有限子覆盖而仿紧，故由\cref{g1:thm:paracompact-normal}正规。在 \(C\) 中对不交闭集 \(K\) 与 \(C\setminus V\) 使用 Urysohn 引理，得到 \(\rho_C\)，分别取值 \(1\) 与 \(0\)。把它在 \(X\setminus V\) 上补成零。两个定义域 \(C\) 与 \(X\setminus V\) 是覆盖 \(X\) 的闭集，函数在交集上相同，故\cref{g1:prop:gluing-lemma}保证得到的 \(\rho\) 连续。其支集包含于紧集 \(\overline V\subseteq U\)，因此本身也紧。
\end{proof}

这个构造还解决了一种常用的延拓问题：若函数已在紧集的一个开邻域上定义，就可把它截成全空间上的函数，同时保持紧集上的值。

\begin{corollary}[截断后的零延拓]\label{g1:cor:cutoff-extension}
	设 \(X\) 为拓扑空间，\(U\subseteq X\) 开，\(f:U\rightarrow\mathbb R^r\) 连续。若连续 \(\rho:X\rightarrow[0,1]\) 满足 \(\operatorname{supp}\rho\subseteq U\)，则
	\[
		\widetilde f(x)=
		\begin{cases}
			\rho(x)f(x),&x\in U,\\
			0,&x\notin U
		\end{cases}
	\]
	为 \(X\) 上的连续函数。若 \(\rho=1\) 于 \(A\subseteq U\)，则 \(\widetilde f=f\) 于 \(A\)。
\end{corollary}
\begin{proof}
	在开集 \(U\) 上，连续性来自连续函数的乘法。若 \(x\notin U\)，则 \(x\notin\operatorname{supp}\rho\)，所以它有一个与支集不交的开邻域；在此邻域中 \(\widetilde f\) 恒为零。这同时处理了 \(U\) 的边界，并不要求 \(f\) 在接近边界时有界。在 \(A\) 上，等式由 \(\rho=1\) 立即得到。
\end{proof}

支集控制使补零位置附近有一块函数恒零的区域，因而截断后的函数在开邻域边缘也连续。

\subsection{从局部有限覆盖构造单位分解}
\label{g1:subsec:0404:03}

有了截断函数，最后一步是归一化。有限个非负数只要不全为零，就可以除以它们的和，得到总和为 \(1\) 的权重。局部有限性允许把同一操作逐点用于一个可能无限的函数族，并保证所得函数连续。

\begin{definition}[从属于开覆盖的单位分解]\label{g1:def:partition-of-unity}
	设 \(\{U_\alpha\}_{\alpha\in A}\) 是拓扑空间 \(X\) 的开覆盖。对于一族连续函数 \(\{\varphi_\alpha:X\rightarrow[0,1]\}_{\alpha\in A}\)，若其支集族局部有限，且
	\[
		\operatorname{supp}\varphi_\alpha\subseteq U_\alpha,\qquad
		\sum_{\alpha\in A}\varphi_\alpha(x)=1\quad(x\in X),
	\]
	则称这族函数为从属于此覆盖的\term[danweifenjie]{单位分解}[partition of unity]。允许某些函数恒为零。局部有限性使每一点附近的和实际上只有有限个可能非零的项。
\end{definition}

\begin{theorem}[连续单位分解的存在]\label{g1:thm:partition-of-unity}
	仿紧 Hausdorff 空间的每个开覆盖都有从属于它的连续单位分解。特别地，由\cref{g1:thm:lch-paracompact}，局部紧、第二可数 Hausdorff 空间满足这一结论。
\end{theorem}
\begin{proof}
	记原覆盖为 \(\{U_\alpha\}\)。由\cref{g1:lem:cover-shrinking}，先取局部有限开覆盖 \(\{V_j\}\)，为各 \(j\) 选择 \(\alpha(j)\)，使 \(\overline{V_j}\subseteq U_{\alpha(j)}\)；再取覆盖 \(X\) 的闭收缩 \(F_j\subseteq V_j\)。空间正规，故\cref{g1:prop:continuous-cutoff}给出连续函数 \(\rho_j:X\rightarrow[0,1]\)，满足
	\[
		\rho_j=1\text{ 于 }F_j,\qquad \operatorname{supp}\rho_j\subseteq V_j.
	\]
	支集族局部有限，因为 \(V_j\) 族局部有限。所以
	\[
		s(x)=\sum_j\rho_j(x)
	\]
	在每一点附近都是有限个连续函数之和，从而连续。每个点属于某个 \(F_j\)，因此至少一个 \(\rho_j(x)=1\)，得到 \(s(x)\ge1\)。现在令 \(\psi_j=\rho_j/s\)。分母连续且不为零，故各 \(\psi_j\) 连续、非负，且 \(\sum_j\psi_j=1\)；它们与 \(\rho_j\) 有相同的非零点集及支集。

	这些函数目前从属于细化覆盖，还需按原覆盖的指标合并。定义
	\[
		\varphi_\alpha=\sum_{j:\,\alpha(j)=\alpha}\psi_j.
	\]
	局部有限性保证每个这样的和连续，也保证合并后仍局部有限：对一个只遇到有限多个 \(\operatorname{supp}\psi_j\) 的邻域，相应的原指标也只有有限多个。更具体地，非负性给出
	\[
		\{\varphi_\alpha\ne0\}=\bigcup_{j:\,\alpha(j)=\alpha}\{\psi_j\ne0\}.
	\]
	局部有限闭集族 \(\{\operatorname{supp}\psi_j\}\) 的子族并为闭集，且包含于 \(U_\alpha\)，故
	\[
		\operatorname{supp}\varphi_\alpha
		\subseteq\bigcup_{j:\,\alpha(j)=\alpha}\operatorname{supp}\psi_j
		\subseteq U_\alpha.
	\]
	最后，各点只有有限个非零项，重排与归组不改变总和，所以 \(\sum_\alpha\varphi_\alpha=1\)。每一项非负且不超过总和 \(1\)，全部要求得证。
\end{proof}

证明中三项条件各承担一个任务：闭收缩保证归一化前的和处处为正，截断保证函数的支集落在允许区域内，局部有限保证求和与合并后的连续性。即使原覆盖十分冗余，分解也不必让每个指标都有非零函数；留下哪些局部数据由细化和构造决定。

\begin{example}[直线上的帐篷函数]\label{g1:exm:hat-partition}
	对每个整数 \(k\)，令
	\[
		\varphi_k(x)=\max\{1-|x-k|,0\},\qquad U_k=(k-3/2,k+3/2).
	\]
	支集为 \([k-1,k+1]\subseteq U_k\)，这族区间局部有限。若 \(x\in[k,k+1]\)，可能非零的函数只有 \(\varphi_k(x)=k+1-x\) 和 \(\varphi_{k+1}(x)=x-k\)，它们的和为 \(1\)；在整数点其余函数也都为零。因此 \(\{\varphi_k\}\) 已经是一组单位分解，无须再归一化。\cref{g1:fig:hat-partition}画出三个相邻权重。它们连续但在峰顶有折角；后续构造光滑对象时，还要把截断函数换成光滑函数。
\end{example}

\begin{figure}[htbp]
	\centering
	\input{Figures/Chapter04-Partition}
	\caption[直线上的单位分解]{直线上的单位分解。蓝、红、绿三条帐篷函数分别以 \(0,1,2\) 为中心；阴影区间 \([0,1]\) 内，前两个权重相加为 \(1\)。其余整数处还有同样平移得到的函数。}
	\label{g1:fig:hat-partition}
\end{figure}

\begin{example}[逐点有限仍可能破坏连续性]\label{g1:exm:pointfinite-sum}
	沿用本节开头的 \(c_n=2^{-n}\)、\(r_n=2^{-n-2}\)，令
	\[
		f_n(x)=\max\{1-|x-c_n|/r_n,0\}.
	\]
	支集是两两不交的闭区间，因此对每个 \(x\)，和 \(f(x)=\sum_nf_n(x)\) 至多有一项非零。然而 \(f(0)=0\)，而 \(f(c_n)=1\) 且 \(c_n\to0\)，所以 \(f\) 在 \(0\) 处不连续。失败的原因并非每一点上的求和项太多，而是无穷多个非零区域聚集到同一点附近。局部有限排除的正是这种情形。
\end{example}

\subsection{连续函数与正定双线性型的整体粘合}
\label{g1:subsec:0404:04}

把单位分解看成随位置变化的凸组合权重，就能接上第三章的凸性。如果局部数据属于同一个向量空间，加权平均有明确意义；如果它们还满足对凸组合封闭的条件，这个条件也能保留下来。

\begin{proposition}[连续局部函数的加权粘合]\label{g1:prop:weighted-gluing}
	设 \(\{U_\alpha\}\) 是拓扑空间 \(X\) 的开覆盖，\(\{\varphi_\alpha\}\) 是从属于它的连续单位分解，且 \(f_\alpha:U_\alpha\rightarrow\mathbb R^r\) 连续。则约定每项在 \(U_\alpha\) 外为零后，
	\[
		f=\sum_\alpha\varphi_\alpha f_\alpha
	\]
	为 \(X\) 上的连续函数。在每个 \(x\)，它属于所有满足 \(\varphi_\alpha(x)>0\) 的 \(f_\alpha(x)\) 的凸包。若局部函数在交叠处相同，则 \(f\) 等于它们通常粘合得到的函数；若所有局部值都落在一个固定凸集 \(C\subseteq\mathbb R^r\) 中，则 \(f(X)\subseteq C\)。
\end{proposition}
\begin{proof}
	由\cref{g1:cor:cutoff-extension}，每项的零延拓连续；这里该推论只使用支集包含，不要求权重在某个集合上等于 \(1\)。这些项的支集包含于 \(\operatorname{supp}\varphi_\alpha\)，因而局部有限，求和后连续。固定 \(x\)，只有有限个正权重，其和为 \(1\)，所以所得值就是所列局部值的一个有限凸组合。若这些值相同，组合仍等于共同值；若它们都属于凸集，\cref{g1:prop:convex-finite-combinations}保证组合仍在该凸集中。
\end{proof}

两份局部数据不相同时，加权粘合在交叠处给出它们的过渡值。若数据在交叠处本来一致，\cref{g1:prop:gluing-lemma}已足够；单位分解更重要的用途，是把未必一致、但允许作凸组合的局部选择合成一个整体选择。

\begin{proposition}[正定双线性型的粘合]\label{g1:prop:positive-form-gluing}
	设 \(X\) 为拓扑空间，\(\{U_\alpha\}_{\alpha\in A}\) 为其开覆盖，\(\{\varphi_\alpha:X\to[0,1]\}_{\alpha\in A}\) 为从属于此覆盖的连续单位分解，即支集族局部有限、\(\operatorname{supp}\varphi_\alpha\subseteq U_\alpha\)，且 \(\sum_{\alpha\in A}\varphi_\alpha(x)=1\)。设 \(r\) 为正整数，\(\operatorname{Sym}_r(\mathbb R)\) 表示实对称 \(r\times r\) 矩阵空间。对每个 \(\alpha\)，给定连续矩阵函数 \(G_\alpha:U_\alpha\rightarrow\operatorname{Sym}_r(\mathbb R)\)，且每个 \(G_\alpha(x)\) 正定。则逐矩阵元作零延拓求和得到
	\[
		G(x)=\sum_\alpha\varphi_\alpha(x)G_\alpha(x),
	\]
	它是连续的正定对称矩阵函数。其所定义的内积在连续换基下按通常的合同变换相容。
\end{proposition}
\begin{proof}
	连续性由上一命题逐矩阵元应用得到，对称性由转置与有限求和交换得到。固定 \(x\) 及非零向量 \(v\in\mathbb R^r\)，有
	\[
		v^{\mathsf T}G(x)v=\sum_{\alpha:\,\varphi_\alpha(x)>0}
		\varphi_\alpha(x)\,v^{\mathsf T}G_\alpha(x)v>0,
	\]
	因为至少一个权重为正，而列出的每个二次型值都严格为正。若换基矩阵为可逆矩阵 \(A(x)\)，各局部矩阵变为 \(A^{\mathsf T}G_\alpha A\)，则
	\[
		\sum_\alpha\varphi_\alpha A^{\mathsf T}G_\alpha A
		=A^{\mathsf T}\Bigl(\sum_\alpha\varphi_\alpha G_\alpha\Bigr)A.
	\]
	因此先换基再求和与先求和再换基给出同一双线性型。
\end{proof}

这里的正定双线性型是一个点处向量空间上的内积，后续讨论流形时会把这种数据称为度量。它与\cref{g1:thm:urysohn-metrization}中的“两点之间的距离函数”属于不同层次：前者的两个输入是同一空间内的向量，后者的两个输入是空间中的点。由点态内积经过曲线长度构造距离，还要另作定义与证明。当前计算说明，只要局部向量坐标按换基关系相容，正定性就允许加权粘合；在第12章讨论向量丛时，局部基的语言会得到系统定义。

最后，允许求凸组合这一条件不能省去。例如两份取值于圆周的局部函数若分别恒为 \((1,0)\) 与 \((-1,0)\)，在权重各为 \(1/2\) 的位置，它们的平均值为原点，已经离开圆周。类似地，正定型可以求和，但一般方程的解、坐标映射或几何约束未必在加权后保留。单位分解的粘合因而适用于能够在同一向量空间相加、且约束对凸组合稳定的局部数据。

% !TeX root = ../../Book1.tex
% 纲要标题：### 4.5 商模型的完整核验
% 纲要位置：几何学计划纲要.md，第 266 行。
% 正文状态：初稿完成；尚未正式审核。

% 纲要内容开始
% 
% * 环面的积模型与边识别模型
% * 实射影空间的商拓扑和仿射图
% * 双原点与稠密轨道商的非分离现象
% 
% ---
% 
% 纲要内容结束

\section{商模型的完整核验}
\label{g1:sec:0405}

把正方形的对边接起来，或者把球面上的对径点合并，图形会提示一个新的空间。要把提示变成证明，必须说明：哪些点被识别，诱导映射是否良定义且双射，两个模型的拓扑是否相同。紧性与 Hausdorff 性常能完成最后一步，但它们也需要实际检验。本节先完成两个常用的几何模型，再考察两种看似温和却会破坏分离性的识别。

\subsection{环面的积模型与边识别模型}
\label{g1:subsec:0405:01}

记 \(S^1=\{z\in\mathbb C:|z|=1\}\)，带来自 \(\mathbb C\cong\mathbb R^2\) 的子空间拓扑。它的两个独立角参数组成\term[huanmian]{环面}[torus]
\[
	T^2=S^1\times S^1.
\]
这里先用积空间定义环面，因而不依赖在三维空间里画出的具体形状。角参数取在闭区间 \([0,1]\) 时，\(0\) 与 \(1\) 表示同一个圆周点，所以两个参数的端点都需要识别。

\begin{theorem}[环面的方形商模型]\label{g1:thm:torus-square}
	记 \(S^1=\{z\in\mathbb C:|z|=1\}\) 带 Euclid 子空间拓扑，\(T^2=S^1\times S^1\) 带积拓扑。设 \(Q=[0,1]^2\)，令 \(\sim\) 为由
	\[
		(0,t)\sim(1,t),\qquad (s,0)\sim(s,1)\qquad(0\le s,t\le1)
	\]
	生成的等价关系。则映射
	\[
		F:Q\longrightarrow T^2,\qquad F(s,t)=(e^{2\pi is},e^{2\pi it})
	\]
	诱导同胚 \(\bar F:Q/{\sim}\to T^2\)。
\end{theorem}
\begin{proof}
	先确定 \(F\) 的纤维，而不只检查题中写出的两种识别。对于 \(s,s'\in[0,1]\)，有 \(e^{2\pi is}=e^{2\pi is'}\) 当且仅当 \(s-s'\in\mathbb Z\)，也就是 \(s=s'\) 或 \(\{s,s'\}=\{0,1\}\)。第二坐标同理。所以两个点的像相等，恰好意味着可以依次改变第一坐标的端点、第二坐标的端点而从一个点到另一个点。这正是所给关系生成的等价关系。特别地，内部点的等价类是单点，边上非角点的等价类含两个点，四个角点属于同一个等价类。

	两个坐标函数连续，故由积空间的映入判据，\(F\) 连续。每个圆周点都有角参数，因而 \(F\) 满射。商下降命题于是给出连续映射 \(\bar F\)，上面对纤维的计算说明它是双射。

	最后核验逆映射连续。闭区间紧，故由\cref{g1:prop:finite-product-compact}，\(Q\) 紧；由\cref{g1:prop:compact-image}，其商空间也是紧的。圆周是度量空间的子空间，故为 Hausdorff 空间，两个圆周的积仍为 Hausdorff 空间。因此\cref{g1:thm:compact-bijection}适用，连续双射 \(\bar F\) 为同胚。
\end{proof}

\begin{figure}[htbp]
	\centering
	\input{Figures/Chapter04-TorusSquare}
	\caption[环面的对边识别]{环面的对边识别。同字母、同参数方向的边彼此对应，四个角合并为一个点。阴影给出该点一个小邻域在正方形中的四块原像；识别以后它们接成一块。}
	\label{g1:fig:torus-square}
\end{figure}

对边识别的方向必须写清。例如竖边上的 \((0,t)\) 与 \((1,t)\) 使用相同的 \(t\)，所以两侧箭头都向上；若改成 \((0,t)\sim(1,1-t)\)，就已经更改等价关系，不能沿用上面的 \(F\)。图中的箭头只是记录这个参数对应，并不代替纤维计算。

边或角在正方形里看起来特殊，但识别后没有成为局部的缺口。这个事实可以直接在积模型中看出。对任意 \(c\in S^1\)，取开半圆
\[
	A_c=\{z\in S^1:\operatorname{Re}(z/c)>0\}.
\]
映射 \(u\mapsto ce^{2\pi iu}\) 把 \((-1/4,1/4)\) 同胚地映到 \(A_c\)。其连续逆映射为
\[
	z\longmapsto\frac1{2\pi}\arctan\frac{\operatorname{Im}(z/c)}{\operatorname{Re}(z/c)};
\]
正实部保证这里选定的反正切分支唯一。于是 \((c,d)\in T^2\) 的开邻域 \(A_c\times A_d\) 同胚于一个开矩形。通过 \(\bar F\) 拉回，处于一条边中间的点附近通常需要两块半矩形拼合，角点附近需要四块角区域拼合，而得到的局部空间都是一个开矩形。这也解释了为什么不能只在正方形中取某一个代表的单侧邻域来决定商拓扑。

\subsection{实射影空间的商拓扑和仿射图}
\label{g1:subsec:0405:02}

\cref{g1:def:projective-space}把实射影空间定义为 \(\mathbb R^{n+1}\) 中过原点的一维线性子空间的集合，并用非零向量的比例类 \([v]\) 表示点。\cref{g1:prop:projective-charts}进一步给出了 \(v_i\ne0\) 时的归一化坐标。现在要给这个集合补上拓扑，并证明已有坐标确实是同胚坐标，而不只是集合之间的双射。

对 \(n\ge0\)，令 \(E=\mathbb R^{n+1}\setminus\{0\}\)，给 \(E\) Euclid 子空间拓扑。记
\[
	p:E\longrightarrow\mathbb RP^n=\mathbb P^n(\mathbb R),\qquad p(v)=[v],
\]
并给 \(\mathbb RP^n\) 由 \(p\) 决定的商拓扑；这里 \(v\sim w\) 的含义是 \(w=\lambda v\)，其中 \(\lambda\in\mathbb R\setminus\{0\}\)。每条这样的直线与单位球面 \(S^n\) 恰交于两个对径点，所以也应能从球面构造同一个空间。

\begin{theorem}[实射影空间的两个商模型]\label{g1:thm:real-projective-quotient}
	设 \(n\ge0\)，\(\mathbb RP^n=(\mathbb R^{n+1}\setminus\{0\})/\mathbb R^\times\) 带非零实数倍关系的商拓扑，\(S^n\) 为 \(\mathbb R^{n+1}\) 的单位球面。把 \(S^n\) 上的 \(u,-u\) 识别，所得商空间 \(S^n/\{\pm1\}\) 与上述 \(\mathbb RP^n\) 自然同胚。空间 \(\mathbb RP^n\) 紧且 Hausdorff，并且在单位代表上由
	\begin{equation}\label{g1:eq:projective-quotient-metric}
		\delta([u],[v])=\min\{\lVert u-v\rVert,\lVert u+v\rVert\},\qquad u,v\in S^n,
	\end{equation}
	给出一个产生其商拓扑的度量。
\end{theorem}
\begin{proof}
	先构造互逆连续映射，不预先假设任何商模型的分离性。记球面上的商映射为 \(q:S^n\to S^n/\{\pm1\}\)，包含映射为 \(i:S^n\hookrightarrow E\)。\(p\circ i\) 在对径类上常值，故下降为连续映射
	\[
		a:S^n/\{\pm1\}\longrightarrow\mathbb RP^n,\qquad a(\{u,-u\})=[u].
	\]
	反向利用连续的归一化映射 \(r:E\to S^n\)、\(r(v)=v/\lVert v\rVert\)。当 \(v\) 乘以非零实数时，\(r(v)\) 只可能保持不变或变成对径点，故 \(q\circ r\) 在 \(p\) 的纤维上常值，也下降为连续映射 \(b:\mathbb RP^n\to S^n/\{\pm1\}\)。因为 \(p(r(v))=p(v)\)、\(r(u)=u\)，两个复合都恒等，因而 \(a,b\) 为互逆同胚。

	再在球面商上检查度量。式\eqref{g1:eq:projective-quotient-metric}在任一代表变号时都不变，故良定义；它对称、非负，且为零恰好意味着两个对径类相同。为证三角不等式，取符号 \(\varepsilon,\eta\in\{1,-1\}\) 使
	\[
		\delta([u],[v])=\lVert u-\varepsilon v\rVert,\qquad
		\delta([v],[w])=\lVert v-\eta w\rVert.
	\]
	于是
	\[
		\begin{aligned}
			\delta([u],[w])&\le\lVert u-\varepsilon\eta w\rVert\\
			&\le\lVert u-\varepsilon v\rVert+\lVert\varepsilon v-\varepsilon\eta w\rVert\\
			&=\delta([u],[v])+\delta([v],[w]).
		\end{aligned}
	\]
	所以 \(\delta\) 为度量。

	最后比较拓扑。任一 \(\delta\)-球的原像满足
	\[
		q^{-1}\bigl(B_\delta([u],\rho)\bigr)
		=B_{S^n}(u,\rho)\cup B_{S^n}(-u,\rho),
	\]
	故在球面中开，说明度量开集都是商拓扑开集。反过来，若 \(U\) 在商拓扑中开且 \([u]\in U\)，则 \(q^{-1}(U)\) 开并含 \(u\)，所以某个球面球 \(B_{S^n}(u,\rho)\) 包含于这个原像。任意满足 \(\delta([u],[v])<\rho\) 的类，都有一个单位代表位于该球内，所以 \(B_\delta([u],\rho)\subseteq U\)。这证明 \(U\) 在度量拓扑中开。两种拓扑相同，因而商空间为 Hausdorff 空间。球面是紧立方体 \([-1,1]^{n+1}\) 中的闭集，故由\cref{g1:prop:compact-closed-subset}紧；它的连续商像由\cref{g1:prop:compact-image}仍紧，完成证明。
\end{proof}

这里紧性来自球面，而不是来自非紧空间 \(E\)；Hausdorff 性来自显式度量，而不是来自“商掉一个对称”的直觉。两种论证承担不同任务。这个度量也有直接含义：比较两条无向直线时，允许把一个单位代表变号，再取两种距离中较小的一个。

\begin{proposition}[实射影空间的仿射图]\label{g1:prop:real-projective-topological-charts}
	设 \(n\ge1\)，\(0\le i\le n\)。集合 \(U_i=\{[X_0:\cdots:X_n]:X_i\ne0\}\) 在 \(\mathbb RP^n\) 中开，且映射
	\[
		\varphi_i:U_i\longrightarrow\mathbb R^n,\qquad
		[X_0:\cdots:X_n]\longmapsto
		\left(\frac{X_0}{X_i},\ldots,\widehat{\frac{X_i}{X_i}},\ldots,\frac{X_n}{X_i}\right)
	\]
	为同胚；帽号表示略去该项。这 \(n+1\) 张图覆盖 \(\mathbb RP^n\)，在重叠部分的换图映射由相应坐标相除给出。
\end{proposition}
\begin{proof}
	原像 \(p^{-1}(U_i)=\{X\in E:X_i\ne0\}\) 开，所以 \(U_i\) 在商拓扑中开。第二章已经说明 \(\varphi_i\) 是双射：每个类唯一地归一化为 \(X_i=1\)，其余 \(n\) 个坐标任意。

	为核验连续性，令 \(E_i=p^{-1}(U_i)\)。由\cref{g1:prop:quotient-tests}，限制 \(p_i:E_i\to U_i\) 仍为商映射。复合 \(\varphi_i\circ p_i\) 是分母不为零的坐标商，故连续，商下降判据于是给出 \(\varphi_i\) 连续。反过来，将 \((x_0,\ldots,\widehat{x_i},\ldots,x_n)\) 在第 \(i\) 位插入 \(1\)，得到连续映射 \(s_i:\mathbb R^n\to E_i\)，并有 \(\varphi_i^{-1}=p_i\circ s_i\)。故逆映射也连续。

	每个非零向量至少有一个非零坐标，所以这些图覆盖全空间。在第 \(i\) 张图中把 \(x_i=1\) 补回，重叠条件为 \(x_j\ne0\)；转到第 \(j\) 张图，就是把所有坐标除以 \(x_j\)，再略去第 \(j\) 位。因此新坐标中第 \(i\) 位为 \(1/x_j\)，其他保留位为 \(x_k/x_j\)。分母非零同时保证换图公式连续。
\end{proof}

例如在 \(\mathbb RP^2\) 的 \(U_0\) 中写点为 \([1:x:y]\)，转入 \(U_1\) 要求 \(x\ne0\)，新坐标为 \((1/x,y/x)\)。直线 \(x=0\) 上的点没有失去意义，它们只是落在 \(U_1\) 之外。这个例子把第二章的代数换图公式与当前的拓扑构造接上了：坐标图是开子空间上的同胚，换图则发生在真正的开重叠区域中。

\subsection{双原点与稠密轨道商的非分离现象}
\label{g1:subsec:0405:03}

前两个模型的识别都留下了足够多的开集来区分点。下面把两份直线几乎全部粘在一起，只保留两个不同的原点。问题不在于点数太多，而在于两个未被识别的点附近，已经没有足够独立的邻域。

\begin{example}[双原点直线]\label{g1:exm:double-origin}
	取不交并 \(\mathbb R_a\sqcup\mathbb R_b\)，把 \(x_a\) 与 \(x_b\) 在所有 \(x\ne0\) 处识别，保留 \(0_a,0_b\) 不同。所得商空间 \(D\) 称为\term[shuangyuandianzhixian]{双原点直线}[line with two origins]。它的每个点都有同胚于开区间的开邻域，每个单点闭，但它不是 Hausdorff 空间。
\end{example}
\begin{proof}
	记商映射为 \(q\)。先证明 \(q\) 开，这样才能准确判断两份直线的像。若 \(O=O_a\sqcup O_b\) 在不交并中开，则它的饱和集在第 \(a\) 份直线中对应
	\[
		O_a\cup(O_b\setminus\{0\}),
	\]
	在第 \(b\) 份中对应 \(O_b\cup(O_a\setminus\{0\})\)，两者都是通常开集。因此 \(q^{-1}(q(O))\) 开，由商拓扑定义，\(q(O)\) 开。

	包含每一份直线再取商，得到连续单射 \(j_a:\mathbb R\to D\)、\(j_b:\mathbb R\to D\)。它们都是到其像的开映射，故给出
	\[
		\mathbb R\cong D\setminus\{0_b\},\qquad \mathbb R\cong D\setminus\{0_a\},
	\]
	而这两个像在 \(D\) 中开。因此每个点确实有通常直线的局部模型。任一单点的原像或者是一份直线的一个原点，或者是两份直线中同坐标的两个点；这些原像在不交并中闭。商映射对闭集的判据便说明每个单点在 \(D\) 中闭。

	但若 \(U\ni0_a\)、\(V\ni0_b\) 都开，则 \(q^{-1}(U)\) 在第 \(a\) 份中包含某个 \((-\varepsilon,\varepsilon)\)，\(q^{-1}(V)\) 在第 \(b\) 份中包含某个 \((-\eta,\eta)\)。取 \(0<|x|<\min\{\varepsilon,\eta\}\)，其公共等价类同时属于 \(U,V\)。所以这两个点没有不交开邻域，\(D\) 非 Hausdorff。
\end{proof}

\begin{figure}[htbp]
	\centering
	\input{Figures/Chapter04-DoubleOrigin}
	\caption[双原点直线的两个邻域]{双原点直线的两个邻域。上下两份中坐标相同的非零点被识别，两个原点仍不同；各自围住原点的区间都包含共同的非零点，故像中的邻域必相交。}
	\label{g1:fig:double-origin}
\end{figure}

从序列也能看到同一障碍。由每一份直线到 \(D\) 的连续映射，非零点序列 \(1/m\) 在 \(D\) 中同时收敛到 \(0_a\) 与 \(0_b\)。局部坐标各自没有问题，却不能排除全局出现两个极限。这说明后来定义流形时，Hausdorff 条件不能从“每一点附近像 Euclid 空间”中省略。

用取值于 Hausdorff 空间的连续映射也无法区分这两个原点。若 \(f:D\to Y\) 连续且 \(Y\) 为 Hausdorff 空间，而 \(f(0_a)\ne f(0_b)\)，则取这两个像的不交开邻域，其原像会把 \(0_a,0_b\) 分离，矛盾。因此 \(f(0_a)=f(0_b)\)。具体地，令 \(r:D\to\mathbb R\) 忘去原点的标签，它由两份恒等坐标映射下降得到，故连续。每个这样的 \(f\) 都唯一写成 \(f=g\circ r\)，其中 \(g=f\circ j_a\) 连续。商空间仍然保存两个点，到 Hausdorff 空间的连续映射却看不出这个差别。

另一种失败来自稠密的等价类。不同等价类在集合意义下完全分离，却散布在彼此的每个小区间里。

\begin{example}[有理平移的轨道商]\label{g1:exm:rational-orbit-quotient}
	在通常实直线上定义 \(x\sim y\) 当且仅当 \(x-y\in\mathbb Q\)，记商空间为 \(\mathbb R/\mathbb Q\)。它至少有两个点，但拓扑只有空集与全空间，所以不是 Hausdorff 空间，且每个单点都不闭。
\end{example}
\begin{proof}
	等价类是 \(x+\mathbb Q\)。有理数的稠密性说明每个类都与每个非空开区间相交。商集合至少有两个元素，因为 \(0\) 与 \(\sqrt2\) 之差不是有理数。

	设 \(W\) 为商空间的非空开集，则 \(U=q^{-1}(W)\) 是非空的饱和开集。取开区间 \((a-\varepsilon,a+\varepsilon)\subseteq U\)。对任意 \(x\in\mathbb R\)，有理数的稠密性给出 \(r\in\mathbb Q\)，使 \(x+r\) 落在这个区间中。由于 \(U\) 饱和，而 \(x\sim x+r\)，得到 \(x\in U\)。因此 \(U=\mathbb R\)，满射性给出 \(W=\mathbb R/\mathbb Q\)。这证明拓扑平凡。这个空间有不止一个点，所以任何单点都不等于全空间，也非空，因而既不开也不闭。
\end{proof}

这个商映射甚至是开映射：对 \(\mathbb R\) 中的开集 \(O\)，
\[
	q^{-1}(q(O))=\bigcup_{r\in\mathbb Q}(O+r)
\]
是开集，故 \(q(O)\) 开。于是“源空间很好，商映射连续、满且开”仍然不够保证目标的分离性。在 \(\mathbb R/\mathbb Q\) 中，每个点的闭包都是整个空间；任何到 Hausdorff 空间的连续映射都为常值，因为若有两个不同的像，不交开邻域的原像将给出两个非空不交开集，而这里根本不存在这样的开集。

两种反例的差别值得保留。双原点直线的每个单点闭，普通点还有熟悉的局部坐标，失败集中在两个原点不能被开集分开；有理平移商连单点闭性也失去了，每个非空开集都扩张成全空间。研究一个新的商空间时，先计算饱和开集，往往就能看出这种差别。关于积、商与粘合的系统介绍，可参阅\cite[第3章, Product Spaces, 第60页; Quotient Spaces, 第65页; Adjunction Spaces, 第73页]{Lee2011Topological}。


\begin{chaptersummary}
	拓扑以开集表达邻近关系，连续性由开集的原像检验。子空间和积空间由映入其中的映射刻画，商空间则由从中出发的映射刻画。商映射要求原像的开性能够反映商集的开性；连续满射本身并不足够，连续双射也未必是同胚。紧到 Hausdorff 的连续双射提供了一个经常可用的同胚判据。

	\ProseParagraphBreak

	紧性把任意覆盖缩成有限覆盖，有限交性质是它的闭集表达。局部紧性在 Hausdorff 条件下提供闭包仍受控制的小邻域。连通性排除分离，道路连通性要求实际路径，局部道路连通性则使两种分支一致并成为开集。分离公理控制点与闭集能否分开，可数性控制需要多少邻域数据；两者不能互相代替。

	\ProseParagraphBreak

	局部紧第二可数 Hausdorff 空间可以用紧集逐层穷尽。对紧壳层作有限选择，便能构造局部有限细化；仿紧 Hausdorff 空间因而具有足够的正规性。Urysohn 引理把闭集分离转成连续函数，正则第二可数空间再用可数个这样的函数编码成一个相容距离。覆盖收缩与截断函数随后给出正的局部有限和，归一化得到单位分解。它能粘合可作凸组合的局部数据，支集控制与局部有限性保证补零和求和的连续性。

	\ProseParagraphBreak

	环面的边识别模型与积模型同胚，实射影空间的球面商、比例商与仿射图描述相容。双原点空间虽局部像直线，却不能分离两个原点；稠密有理平移轨道的商甚至只有两个开集。这些模型将局部坐标与整体分离性分开：前者由各点的邻域描述，后者还取决于商拓扑留下了哪些开集。
\end{chaptersummary}

\BookExercises
以下各题均给出完整解答。处理商空间时先确定拓扑和映射方向；处理紧性、分离与粘合时，注明每次使用的是开覆盖、闭性、局部有限性还是可数性。

\begin{exercise}\label{g1:ex:truncated-metric}
	设 \((X,d)\) 为度量空间，令 \(\delta(x,y)=\min\{1,d(x,y)\}\)。证明 \(\delta\) 也是度量，并且与 \(d\) 产生相同拓扑。由此说明，一个具体度量下的有界性不是拓扑本身决定的性质。
\end{exercise}
\begin{solution}
	正定性与对称性直接继承。对非负数 \(a,b\)，若 \(a+b<1\)，则 \(\min\{1,a+b\}=a+b=\min\{1,a\}+\min\{1,b\}\)；若 \(a+b\ge1\)，右侧至少为 \(1\)。因此总有
	\[
		\min\{1,a+b\}\le\min\{1,a\}+\min\{1,b\}.
	\]
	结合 \(d\) 的三角不等式，就得到 \(\delta\) 的三角不等式。对 \(0<r<1\)，两种度量的 \(r\)-球完全相同。任一点的每个度量邻域都包含这种小半径的球，所以两者的开集相同。在 \(\mathbb R\) 上取 \(d(x,y)=|x-y|\)，其整个空间无界，而 \(\delta\le1\) 使整个空间有界；两者的拓扑却相同。
\end{solution}

\begin{exercise}\label{g1:ex:continuous-graph}
	设 \(f:X\to Y\) 连续，\(\Gamma_f=\{(x,f(x)):x\in X\}\subseteq X\times Y\) 带子空间拓扑。证明 \(x\mapsto(x,f(x))\) 给出 \(X\) 与 \(\Gamma_f\) 的同胚。若 \(Y\) 为 Hausdorff 空间，再证明 \(\Gamma_f\) 在 \(X\times Y\) 中闭。
\end{exercise}
\begin{solution}
	映射 \(s(x)=(x,f(x))\) 的两个坐标连续，故作为到积空间的映射连续；由子空间的映入判据，它到 \(\Gamma_f\) 仍连续。它是双射，逆映射为第一坐标投影的限制，因而连续。

	若 \((a,b)\notin\Gamma_f\)，则 \(b\ne f(a)\)。取 \(Y\) 中不交开集 \(U\ni f(a)\)、\(V\ni b\)。开矩形 \(f^{-1}(U)\times V\) 含 \((a,b)\)，而其中若有图像点 \((x,f(x))\)，就会有 \(f(x)\in U\cap V\)，矛盾。因此图像的补集中每一点都有包含于补集的开邻域，补集开，图像闭。
\end{solution}

\begin{exercise}\label{g1:ex:absolute-value-quotient}
	证明 \(q:\mathbb R\to[0,\infty)\)、\(q(x)=|x|\) 是开商映射。对任意拓扑空间 \(Z\)，证明连续映射 \(f:\mathbb R\to Z\) 能写成 \(f=g\circ q\)，其中 \(g:[0,\infty)\to Z\) 连续，当且仅当 \(f(x)=f(-x)\) 对所有 \(x\) 成立。
\end{exercise}
\begin{solution}
	映射 \(q\) 连续且满。若 \(O\subseteq\mathbb R\) 开，考虑 \(t\in q(O)\)。当 \(t>0\) 时，取 \(x\in O\) 满足 \(|x|=t\)，再取足够小的开区间围住 \(x\)，既包含于 \(O\) 又不跨过零，其绝对值像包含 \(t\) 的一个开区间。当 \(t=0\) 时，\(0\in O\)，故某个 \((-\varepsilon,\varepsilon)\subseteq O\)，从而 \([0,\varepsilon)\subseteq q(O)\)。这两种情形都给出 \([0,\infty)\) 中的开邻域，所以 \(q(O)\) 开；由开连续满射判据，\(q\) 为商映射。

	若 \(f=g\circ q\)，则 \(q(x)=q(-x)\) 立即给出偶性。反之，纤维在 \(t>0\) 时为 \(\{t,-t\}\)，在 \(0\) 处为单点，偶性恰好说明 \(f\) 在每个纤维上常值。由商下降命题得到唯一连续 \(g\)，具体为 \(g(t)=f(t)\)。
\end{solution}

\begin{exercise}\label{g1:ex:infinite-closed-pasting}
	在 \(X=[0,1]\) 中取 \(A_n=\{0\}\cup[1/n,1]\)，\(n\ge1\)。证明这些集合组成闭覆盖。定义 \(f(0)=0\)，\(f(x)=1\)（\(x>0\)）。证明每个限制 \(f|_{A_n}:A_n\to\mathbb R\) 连续，但 \(f\) 不连续，并指出有限闭覆盖证明中失效的一步。
\end{exercise}
\begin{solution}
	每个 \(A_n\) 是两个闭集的有限并，故闭；任意 \(x>0\) 对足够大的 \(n\) 满足 \(1/n\le x\)，所以它们覆盖 \(X\)。在 \(A_n\) 中，\(\{0\}=A_n\cap(-1/(2n),1/(2n))\) 是开集，其补集 \([1/n,1]\) 也开，因为它等于 \(A_n\cap(1/(2n),2)\)。函数在这两个不交开集上分别常值，故限制连续。

	然而 \(f^{-1}((-1/2,1/2))=\{0\}\) 不在 \(X\) 中开，所以 \(f\) 不连续。若改用闭集原像看，\(\{1\}\) 在 \(\mathbb R\) 中闭，而
	\[
		f^{-1}(\{1\})=\bigcup_{n\ge1}[1/n,1]=(0,1]
	\]
	在 \(X\) 中不闭。有限闭覆盖的证明需要把闭集原像作有限并；这里恰恰变成了不能保持闭性的无限并。
\end{solution}

\begin{exercise}[无限 Helly 定理]\label{g1:exr:infinite-helly-compact}
	设 \(\{C_\alpha\}_{\alpha\in A}\) 是 \(\mathbb R^d\) 中的非空族闭凸集，其中至少一个集合紧。若任意至多 \(d+1\) 个集合都有公共点，证明 \(\bigcap_{\alpha\in A}C_\alpha\ne\varnothing\)。明确说明有限维性与紧性分别在哪一步使用。
\end{exercise}
\begin{solution}
	固定紧集 \(C_{\alpha_0}\)。它由单个集合的相交条件非空。令 \(F_\alpha=C_\alpha\cap C_{\alpha_0}\)，则每个 \(F_\alpha\) 都是 \(C_{\alpha_0}\) 中的闭集。任取有限个指标 \(\alpha_1,\ldots,\alpha_m\)，对有限族
	\[
		C_{\alpha_0},C_{\alpha_1},\ldots,C_{\alpha_m}
	\]
	应用\cref{g1:thm:helly}。其中每个至多含 \(d+1\) 个集合的子族均有公共点，故整个有限族的交非空。这意味着 \(F_{\alpha_1}\cap\cdots\cap F_{\alpha_m}\ne\varnothing\)，所以 \(\{F_\alpha\}\) 具有有限交性质。由紧空间 \(C_{\alpha_0}\) 上的闭集刻画，
	\[
		\varnothing\ne\bigcap_{\alpha\in A}F_\alpha
		=\bigcap_{\alpha\in A}C_\alpha.
	\]
	有限维性用来把每个有限子族的相交检查缩减为至多 \(d+1\) 个集合的检查；紧性再把所有有限次相交合并为无限次相交。集合的闭性保证 \(F_\alpha\) 是紧空间中的闭集，也是最后一步的必要输入。
\end{solution}

\begin{exercise}\label{g1:exr:compact-closed-projection}
	设 \(X\) 为任意拓扑空间，\(K\) 为紧空间。证明投影 \(\pi:X\times K\to X\) 把闭集送到闭集。再在 \(\mathbb R^2\) 中给出一个闭集，说明若第二因子不紧，投影的像可以不闭。
\end{exercise}
\begin{solution}
	设 \(F\subseteq X\times K\) 闭，取 \(x\notin\pi(F)\)。这意味着整条纤维 \(\{x\}\times K\) 都在开集 \((X\times K)\setminus F\) 内。对每个 \(y\in K\)，取开矩形 \(U_y\times V_y\) 含 \((x,y)\)，并包含在这个补集中。由 \(K\) 紧，有限个 \(V_{y_1},\ldots,V_{y_m}\) 覆盖 \(K\)。于是
	\[
		U=U_{y_1}\cap\cdots\cap U_{y_m}
	\]
	是 \(x\) 的开邻域，且 \(U\times K\) 与 \(F\) 不交，从而 \(U\cap\pi(F)=\varnothing\)。因此 \(X\setminus\pi(F)\) 开，\(\pi(F)\) 闭。若 \(K=\varnothing\)，像为空，同样闭。

	取 \(F=\{(x,y)\in\mathbb R^2:xy=1\}\)。它是连续函数 \((x,y)\mapsto xy\) 对闭集 \(\{1\}\) 的原像，故闭；它在第一坐标上的投影为 \(\mathbb R\setminus\{0\}\)，却不闭。趋近零的横坐标必须配上越来越大的纵坐标，紧因子所禁止的正是这种逃离。
\end{solution}

\begin{exercise}\label{g1:exr:rationals-local-properties}
	令 \(\mathbb Q\) 带有实直线诱导的拓扑。证明它的每个连通分支、每个道路分支都是单点，但它既不局部道路连通，也不局部紧。
\end{exercise}
\begin{solution}
	若一个子集 \(A\subseteq\mathbb Q\) 含有不同的点 \(p<q\)，在二者之间取无理数 \(\alpha\)。则
	\[
		A=(A\cap(-\infty,\alpha))\sqcup(A\cap(\alpha,\infty))
	\]
	把 \(A\) 分成两个非空、不交相对开集，因为 \(\alpha\notin\mathbb Q\)，没有遗漏点。因此有理数空间的非空连通子集只能是单点，各连通分支也都是单点。道路连通蕴含连通，所以道路分支同样是单点。任何单点都不在 \(\mathbb Q\) 中开，因为每个非空实区间都含有其他有理数；于是没有点具有道路连通开邻域，局部道路连通性失败。

	假设某个 \(q\in\mathbb Q\) 有紧邻域 \(N\subseteq\mathbb Q\)。可选 \(\varepsilon>0\)，使 \((q-\varepsilon,q+\varepsilon)\cap\mathbb Q\subseteq N\)。包含映射 \(\mathbb Q\hookrightarrow\mathbb R\) 连续，故 \(N\) 在 \(\mathbb R\) 中也紧；实直线为 Hausdorff 空间，所以 \(N\) 闭于 \(\mathbb R\)。但在上述区间中取无理数 \(\beta\)，再取区间内趋近 \(\beta\) 的有理数列。该列全在 \(N\) 中，闭性迫使 \(\beta\in N\subseteq\mathbb Q\)，矛盾。因此 \(\mathbb Q\) 不局部紧。
\end{solution}

\begin{exercise}\label{g1:exr:polygonal-connected-open}
	设 \(U\subseteq\mathbb R^n\) 是非空连通开集。证明任意两点都可由完全位于 \(U\) 内的有限折线相连。这里有限折线是有限条依次相接的直线段，不要求它没有自交。
\end{exercise}
\begin{solution}
	固定 \(p\in U\)，令 \(A\) 为能够通过有限折线从 \(p\) 到达的点集，允许只含一个点的退化折线，所以 \(p\in A\)。若 \(x\in A\)，因 \(U\) 开，可选开球 \(B(x,r)\subseteq U\)。对任意 \(y\in B(x,r)\)，把线段 \([x,y]\subseteq B(x,r)\) 接在原折线后，就得到从 \(p\) 到 \(y\) 的有限折线。因此 \(B(x,r)\subseteq A\)，\(A\) 开于 \(U\)。

	若 \(x\in U\setminus A\)，同样选 \(B(x,r)\subseteq U\)。假如该球含 \(y\in A\)，就可以把从 \(p\) 到 \(y\) 的折线接上线段 \([y,x]\)，得到 \(x\in A\)，矛盾。因此 \(B(x,r)\subseteq U\setminus A\)，补集也开于 \(U\)。由 \(U\) 连通，非空开闭子集 \(A\) 必为整个 \(U\)。取任意终点即可得到所需折线。对折线各段分别作线性参数化，再用有限次闭集粘合，也得到通常意义的连续道路。
\end{solution}

\begin{exercise}\label{g1:exr:maps-components}
	\begin{enumerate}
		\item 给实数集合分别赋予离散拓扑与通常拓扑。证明恒等集合映射从前者到后者是连续满射，但不会把一个连通分支满射到目标的连通分支。
		\item 设 \(X\) 为非空连通空间，\(D\) 为离散空间。证明任何连续映射 \(f:X\to D\) 都为常值映射。
		\item 设 \(X\) 局部道路连通。证明 \(f:X\to D\) 连续当且仅当 \(f\) 在每个连通分支上为常值。说明若删除局部条件，反向结论可能失败。
	\end{enumerate}
\end{exercise}
\begin{solution}
	1. 离散定义域的每个子集都开，因此任意开集的原像都开，恒等集合映射连续；它显然满射。离散空间中含有至少两个点的子集，可由其中一个单点及其补集分离，所以每个连通分支是单点。通常实直线连通，其唯一分支为整条直线。故每个源分支的像只是目标分支中的一个点。

	2. 若 \(f\) 取至少两个不同值，取其中一个值 \(d\)。离散空间的 \(\{d\}\) 及其补集都开，所以 \(f^{-1}(\{d\})\) 是 \(X\) 的非空真开闭子集，与连通性矛盾。

	3. 连续时，把第二问用于每个连通分支即可。反之，由\cref{g1:thm:locally-path-components}，\(X\) 的各连通分支开。若 \(f\) 在每个分支上为常值，则对任意 \(E\subseteq D\)，原像 \(f^{-1}(E)\) 是若干分支的并，因而开；这正是 \(f\) 连续。

	若删去局部条件，令 \(X=\mathbb Q\) 带通常子空间拓扑，\(D=\{0,1\}\) 离散，并定义 \(f(0)=1\)、对 \(q\ne0\) 令 \(f(q)=0\)。\cref{g1:exr:rationals-local-properties}证明各连通分支均为单点，故 \(f\) 在每个分支上为常值；但 \(f^{-1}(\{1\})=\{0\}\) 不开，\(f\) 不连续。
\end{solution}

\begin{exercise}\label{g1:ex:lower-limit-topology}
	在 \(\mathbb R\) 上，以所有半开区间 \([a,b)\)（\(a<b\)）为基定义\term[xiaxiantuopu]{下限拓扑}[lower limit topology]。
	\begin{enumerate}
		\item 证明它是 Hausdorff、正则、第一可数且可分的空间。
		\item 证明它不第二可数。再证明每个可分度量空间都第二可数，并据此判断下限拓扑能否由某个距离诱导。
	\end{enumerate}
\end{exercise}
\begin{solution}
	1. 这些区间覆盖 \(\mathbb R\)，两个区间的非空交仍是同类区间，所以确实构成拓扑基。若 \(x<y\)，开集 \([x,y)\) 与 \([y,y+1)\) 分别包含两点且不交，故空间 Hausdorff。每个基元素 \([a,b)\) 也闭，因为其补集 \((-\infty,a)\cup[b,\infty)\) 是同类基元素之并。给定 \(x\in U\) 且 \(U\) 开，选 \(x\in[a,b)\subseteq U\)，便有 \(x\in[a,b)=\overline{[a,b)}\subseteq U\)。结合 \(T_1\) 条件与\cref{g1:lem:regular-normal-shrink}，得到正则性。对每个 \(x\)，区间 \([x,x+1/n)\) 构成可数邻域基，故空间第一可数。每个非空基元素都包含有理数，因此 \(\mathbb Q\) 稠密，空间可分。

	2. 设 \(\mathcal B\) 为任意拓扑基。对每个实数 \(x\)，取 \(B_x\in\mathcal B\) 满足 \(x\in B_x\subseteq[x,x+1)\)。若 \(x<y\)，则 \(B_y\) 不含 \(x\)，而 \(B_x\) 含 \(x\)，所以 \(B_x\ne B_y\)。于是 \(\mathcal B\) 至少与 \(\mathbb R\) 一样大，不可能可数。

	另一方面，设度量空间 \((X,d)\) 有可数稠密集 \(D\)。所有以 \(D\) 中的点为中心、正有理数为半径的球组成可数族。给定 \(x\in U\) 且 \(U\) 开，取 \(B(x,\varepsilon)\subseteq U\)，再取 \(q\in D\) 使 \(d(x,q)<\varepsilon/3\)，并取正有理数 \(r\) 满足
	\[
		d(x,q)<r<\varepsilon-d(x,q).
	\]
	三角不等式给出 \(x\in B(q,r)\subseteq B(x,\varepsilon)\subseteq U\)。故这些球是可数基。若下限拓扑可度量化，其已证的可分性就会推出第二可数性，与前面的结论矛盾。因此它不可度量化。这也说明 Urysohn 度量化定理中的第二可数不能简单替换成“第一可数且可分”。
\end{solution}

\begin{exercise}\label{g1:ex:open-euclidean-exhaustion}
	设 \(\Omega\subseteq\mathbb R^n\) 为非空开集。若 \(\Omega\ne\mathbb R^n\)，定义
	\[
		K_m=\{x\in\mathbb R^n:\lVert x\rVert\le m,\ d(x,\mathbb R^n\setminus\Omega)\ge1/m\},
		\qquad m\ge1.
	\]
	证明这是 \(\Omega\) 的紧穷竭，说明其中两条不等式各自排除了什么。若 \(\Omega=\mathbb R^n\)，给出相应的简单选择。
\end{exercise}
\begin{solution}
	到非空集合的距离函数满足 Lipschitz 不等式，故连续。因此 \(K_m\) 是两个闭集的交，并且有界，由 Euclid 空间的闭有界紧性定理可知它紧。若 \(x\notin\Omega\)，则其到补集的距离为零，所以 \(K_m\subseteq\Omega\)。

	对 \(x\in K_m\)，有
	\[
		\lVert x\rVert\le m<m+1,\qquad
		d(x,\mathbb R^n\setminus\Omega)\ge1/m>1/(m+1).
	\]
	两条严格不等式和连续性说明 \(x\) 有一个包含于 \(K_{m+1}\) 的 Euclid 开邻域，故 \(K_m\subseteq\operatorname{int}_\Omega K_{m+1}\)。若 \(x\in\Omega\)，因为 \(\Omega\) 开，其到补集的距离为某个正数；取 \(m\) 大于 \(\lVert x\rVert\) 并使 \(1/m\) 不超过这一正数，即得 \(x\in K_m\)。所以它们穷尽 \(\Omega\)。

	条件 \(\lVert x\rVert\le m\) 控制向无穷远逃逸；条件 \(d(x,\mathbb R^n\setminus\Omega)\ge1/m\) 控制向边界靠近。一般开集可能同时有这两类非紧性。当 \(\Omega=\mathbb R^n\) 时没有边界需要避开，直接取 \(K_m=\{x:\lVert x\rVert\le m\}\) 即可。
\end{solution}

\begin{exercise}\label{g1:ex:uncountable-discrete}
	设 \(X\) 是不可数集合，赋予离散拓扑。证明它局部紧、正规、仿紧且可度量化，但既不第二可数，也不具有 Lindelöf 性质，并且不能写成可数个紧子集之并。说明这些事实与\cref{g1:cor:lch-second-countable}并不矛盾。
\end{exercise}
\begin{solution}
	每个单点都是紧开邻域，故 \(X\) 局部紧。任意子集同时开且闭，所以任意两个不交闭集本身就是不交开邻域，得到正规性。对任意开覆盖，所有单点组成一个开细化；每个点的单点邻域只与这一族中的一个成员相交，所以细化局部有限，空间仿紧。离散距离 \(d(x,y)=1\)（\(x\ne y\)）及 \(d(x,x)=0\) 诱导离散拓扑，故空间可度量化。

	每个拓扑基都必须含所有单点，因而不可能可数。单点开覆盖没有可数子覆盖，所以 Lindelöf 性质也不成立。若 \(K\subseteq X\) 紧，则其单点开覆盖有有限子覆盖，故 \(K\) 必有限。因此可数个紧子集之并至多可数，不可能等于 \(X\)。推论中的条件是获得全部结论的一组充分条件，并没有宣称可度量化或仿紧必然推出第二可数；本例正好区分了这些方向。
\end{solution}

\begin{exercise}\label{g1:ex:compatible-metric-incomplete}
	在 \(X=(0,1)\) 上取 \(f_n(x)=x^n\)，并定义
	\[
		d(x,y)=\sum_{n\ge1}2^{-n}|x^n-y^n|.
	\]
	求此距离的闭式，证明它诱导 \((0,1)\) 的通常拓扑，并证明序列 \(x_k=1/(k+1)\) 关于 \(d\) 为 Cauchy 列却不在 \(X\) 中收敛。由此说明度量化结论没有自动给出哪一项额外性质。
\end{exercise}
\begin{solution}
	若 \(x\ge y\)，则每项 \(x^n-y^n\ge0\)，应用几何级数求和公式得到
	\[
		d(x,y)=\sum_{n\ge1}(x/2)^n-\sum_{n\ge1}(y/2)^n
		=\frac{x}{2-x}-\frac{y}{2-y}.
	\]
	由对称性，一般有 \(d(x,y)=|g(x)-g(y)|\)，其中 \(g(x)=x/(2-x)\)。函数 \(g\) 将 \((0,1)\) 连续严格递增地映满 \((0,1)\)，反函数为 \(g^{-1}(t)=2t/(1+t)\)，也连续。因此 \(d\) 是通常距离经同胚 \(g\) 拉回所得的度量，诱导通常拓扑。

	对所给序列，\(g(x_k)=1/(2k+1)\to0\)，所以对任意 \(\varepsilon>0\)，充分大的 \(k,l\) 满足
	\[
		d(x_k,x_l)=\left|\frac1{2k+1}-\frac1{2l+1}\right|<\varepsilon,
	\]
	即它是 Cauchy 列。若它在 \(X\) 中关于 \(d\) 收敛于 \(x\)，则 \(g(x_k)\to g(x)\)，故 \(g(x)=0\)，这与 \(x\in(0,1)\) 矛盾。因此这个相容度量不完备。空间的可度量化只保证存在诱导其拓扑的距离；对给出的某个距离是否每个 Cauchy 列都收敛，还需要另外证明。
\end{solution}

\begin{exercise}\label{g1:ex:interval-partition-width}
	给定 \(a>0\)，对每个 \(k\in\mathbb Z\) 令 \(U_k=(k-a,k+a)\)。证明存在连续函数族 \(\{\varphi_k\}\)，使支集族局部有限、\(\operatorname{supp}\varphi_k\subseteq U_k\) 且 \(\sum_k\varphi_k=1\)，当且仅当 \(a>1/2\)。在可行情形给出一个公式。
\end{exercise}
\begin{solution}
	若 \(a\le1/2\)，点 \(x=1/2\) 不在任何 \(U_k\) 中，所以每个 \(\varphi_k(1/2)\) 都为零，不可能相加为 \(1\)。这里甚至连覆盖条件都未满足。

	若 \(a>1/2\)，选择 \(b\) 满足 \(1/2<b<a\)，并定义
	\[
		\rho_k(x)=\max\{b-|x-k|,0\},\qquad
		\varphi_k(x)=\frac{\rho_k(x)}{\sum_{j\in\mathbb Z}\rho_j(x)}.
	\]
	每个 \(\rho_k\) 的支集为 \([k-b,k+b]\subseteq U_k\)。任一有界区间只与有限多个这样的支集相交，故该族局部有限；分母因而连续。每个实数到某个整数的距离不超过 \(1/2<b\)，所以至少一个 \(\rho_j(x)>0\)，分母处处为正。归一化后各函数连续、非负、总和为 \(1\)，支集仍为原来的闭区间。中间选择 \(b<a\)，使支集连同端点也留在指定开区间内。
\end{solution}

\begin{exercise}\label{g1:ex:explicit-cutoff-extension}
	设 \(U=(0,2)\)、\(K=[1/2,3/2]\)，且 \(f(x)=1/(x(2-x))\) 定义在 \(U\) 上。构造 \(\rho:\mathbb R\rightarrow[0,1]\)，使其在 \(K\) 上为 \(1\)、支集包含在 \([1/4,7/4]\) 中。写出 \(\rho f\) 的零延拓，证明它连续，并解释为何 \(f\) 本身不能连续延拓到 \([0,2]\) 上的实值函数。
\end{exercise}
\begin{solution}
	可取分段线性函数
	\[
		\rho(x)=
		\begin{cases}
			0,&x\le1/4,\\
			4x-1,&1/4\le x\le1/2,\\
			1,&1/2\le x\le3/2,\\
			7-4x,&3/2\le x\le7/4,\\
			0,&x\ge7/4.
		\end{cases}
	\]
	相邻表达式在接点有相同取值，故 \(\rho\) 连续，且支集恰为 \([1/4,7/4]\)。令
	\[
		F(x)=
		\begin{cases}
			\rho(x)/(x(2-x)),&0<x<2,\\
			0,&x\le0\text{ 或 }x\ge2.
		\end{cases}
	\]
	在 \(U\) 内，分母不为零，因此 \(F\) 连续。在 \(0\) 的邻域 \((-1/8,1/8)\) 和 \(2\) 的邻域 \((15/8,17/8)\) 内，它恒为零，所以两端的补零也是连续的。它在 \(K\) 上与 \(f\) 相同。反之，\(f(x)\to+\infty\) 当 \(x\to0^+\) 或 \(x\to2^-\)，故不可能为这两个端点指定有限实数值而得到连续延拓。截断保留了指定紧集上的信息，同时去除了边界处的发散。
\end{solution}

\begin{exercise}\label{g1:ex:partition-approximation}
	设 \(f:X\rightarrow\mathbb R^r\) 连续，\(\{U_\alpha\}\) 有从属连续单位分解 \(\{\varphi_\alpha\}\)，且连续函数 \(g_\alpha:U_\alpha\rightarrow\mathbb R^r\) 满足 \(\lVert g_\alpha(x)-f(x)\rVert<\varepsilon\)，其中 \(\varepsilon>0\) 固定。证明加权粘合 \(g=\sum_\alpha\varphi_\alpha g_\alpha\) 仍满足 \(\lVert g-f\rVert<\varepsilon\)。若还已知 \(\lVert f(x)\rVert=1\) 且 \(\varepsilon<1\)，证明 \(h=g/\lVert g\rVert\) 连续且
	\[
		\lVert h(x)-f(x)\rVert<2\varepsilon.
	\]
	说明此处把平均值重新投到单位球面为什么合法。
\end{exercise}
\begin{solution}
	在每个 \(x\)，正权重只有有限个，且其和为 \(1\)。由三角不等式，
	\[
		\lVert g(x)-f(x)\rVert
		\le\sum_{\alpha:\,\varphi_\alpha(x)>0}\varphi_\alpha(x)
		\lVert g_\alpha(x)-f(x)\rVert<\varepsilon.
	\]
	严格不等式成立，是因为这是有限个严格小于 \(\varepsilon\) 的数的凸组合。若 \(\lVert f(x)\rVert=1\)，反三角不等式给出 \(\lVert g(x)\rVert\ge1-\lVert g(x)-f(x)\rVert>1-\varepsilon>0\)。因此归一化的分母从不为零，\(h\) 连续并取值于单位球面。又有
	\[
		\lVert h-g\rVert=|1-\lVert g\rVert|
		\le\lVert g-f\rVert,
	\]
	从而 \(\lVert h-f\rVert\le\lVert h-g\rVert+\lVert g-f\rVert<2\varepsilon\)。一般球面值函数的平均可能为零；这里已有共同的单位长度目标 \(f\)，误差小于 \(1\) 保证平均值始终避开原点，才允许作归一化。
\end{solution}

\begin{exercise}\label{g1:ex:averaging-positive-forms}
	令 \(G_0=\operatorname{diag}(1,4)\)、\(G_1=\operatorname{diag}(4,1)\)，并对 \(t\in[0,1]\) 取 \(G(t)=(1-t)G_0+tG_1\)。求它的特征值与行列式，证明对所有 \(v\in\mathbb R^2\) 有
	\[
		\lVert v\rVert^2\le v^{\mathsf T}G(t)v\le4\lVert v\rVert^2.
	\]
	再比较 \(G(1/2)^{-1}\) 与 \((G_0^{-1}+G_1^{-1})/2\)。
\end{exercise}
\begin{solution}
	直接计算
	\[
		G(t)=\operatorname{diag}(1+3t,4-3t),\qquad
		\det G(t)=4+9t-9t^2.
	\]
	两个对角元就是特征值，均在 \([1,4]\) 内。因此对 \(v=(v_1,v_2)\)，分别估计两个平方项，就得到题设上下界，尤其 \(G(t)\) 始终正定。

	当 \(t=1/2\) 时，\(G(t)=(5/2)I\)，故其逆为 \((2/5)I\)。另一方面，\(G_0^{-1}=\operatorname{diag}(1,1/4)\)、\(G_1^{-1}=\operatorname{diag}(1/4,1)\)，平均为 \((5/8)I\)，与前者不同。正定型构成对凸组合封闭的集合，但求逆不是线性运算，不能与加权平均任意交换。
\end{solution}

\begin{exercise}\label{g1:ex:closed-cover-partition-failure}
	实线由两个闭集 \(A=(-\infty,0]\)、\(B=[0,\infty)\) 覆盖。证明不存在连续函数 \(\varphi_A,\varphi_B:\mathbb R\rightarrow[0,1]\)，同时满足
	\[
		\operatorname{supp}\varphi_A\subseteq A,\qquad
		\operatorname{supp}\varphi_B\subseteq B,\qquad
		\varphi_A+\varphi_B=1.
	\]
	这与有限闭覆盖上的连续函数粘合引理是否矛盾？
\end{exercise}
\begin{solution}
	支集条件迫使 \(\varphi_A(x)=0\) 于 \(x>0\)。令 \(x\to0^+\)，连续性给出 \(\varphi_A(0)=0\)。同理，从 \(x<0\) 一侧得到 \(\varphi_B(0)=0\)。于是二者在 \(0\) 处的和为零，矛盾。

	闭集粘合引理以局部函数在交集上一致为前提；此时直接按已有共同值定义函数即可，无须寻找过渡权重。单位分解则要求每个权重在指定区域外消失，并在重叠处合起来保持总和为一。这个闭覆盖的重叠只有一个边界点，连续性迫使两个权重在那里同时为零，故不能承担该任务。
\end{solution}

\begin{exercise}\label{g1:ex:cylinder-quotient}
	在 \(Q=[0,1]^2\) 中只规定 \((0,t)\sim(1,t)\)，不再识别水平边。证明所得商空间同胚于 \(S^1\times[0,1]\)。写出两条水平边在这个同胚下的像，并说明本题中四个角形成几个等价类。
\end{exercise}
\begin{solution}
	取 \(F(s,t)=(e^{2\pi is},t)\)。它连续、满射，且 \(F(s,t)=F(s',t')\) 等价于 \(t=t'\) 以及 \(s=s'\) 或 \(\{s,s'\}=\{0,1\}\)，恰好给出题中的等价类。由商下降得到连续双射 \(\bar F:Q/{\sim}\to S^1\times[0,1]\)。源空间为紧正方形的连续像，故紧；目标是两个度量空间的有限积，故 Hausdorff。由紧到 Hausdorff 的连续双射定理，\(\bar F\) 为同胚。

	下边映到 \(S^1\times\{0\}\)，上边映到 \(S^1\times\{1\}\)，它们没有被彼此识别。四个角形成两个类：\((0,0)\sim(1,0)\) 与 \((0,1)\sim(1,1)\)。这与环面中四角合并成一个类不同。
\end{solution}

\begin{exercise}\label{g1:ex:real-projective-circle}
	证明
	\[
		\Phi:\mathbb RP^1\longrightarrow S^1,\qquad
		[x:y]\longmapsto
		\left(\frac{x^2-y^2}{x^2+y^2},\frac{2xy}{x^2+y^2}\right)
	\]
	为同胚。在仿射图 \([1:t]\) 中写出公式，并确定漏去的射影点和圆周点。
\end{exercise}
\begin{solution}
	两个分子分母都在 \((x,y)\) 同时乘以非零实数时乘以同一平方因子，故公式良定义；恒等式 \((x^2-y^2)^2+(2xy)^2=(x^2+y^2)^2\) 保证像在圆周上。原非零向量空间上的公式连续，所以由商下降，\(\Phi\) 连续。

	取单位代表 \(z=x+iy\in S^1\)，该公式就是 \(z\mapsto z^2\)。任一单位复数可写为 \(e^{i\theta}\)，因而有平方根 \(e^{i\theta/2}\)；两个平方根恰相差负号。于是每个像点恰对应一个对径类，\(\Phi\) 为双射。由正文，\(\mathbb RP^1\) 紧，圆周 Hausdorff，所以 \(\Phi\) 为同胚。

	在 \([1:t]\) 图中，
	\[
		t\longmapsto\left(\frac{1-t^2}{1+t^2},\frac{2t}{1+t^2}\right).
	\]
	漏去的射影点为 \([0:1]\)，其像为 \((-1,0)\)。它同时是 \(t\to+\infty\) 与 \(t\to-\infty\) 时的极限，这正反映实射影直线只添了一个无穷远点。
\end{solution}

\begin{exercise}\label{g1:ex:projective-projection-matrix}
	设 \(n\ge0\)，\(\operatorname{Sym}_{n+1}(\mathbb R)\) 带矩阵条目给出的 Euclid 拓扑。证明
	\[
		J:\mathbb RP^n\longrightarrow\operatorname{Sym}_{n+1}(\mathbb R),\qquad
		[v]\longmapsto\frac{vv^{\mathsf T}}{v^{\mathsf T}v}
	\]
	给出到其像的同胚，并证明像恰为满足 \(P^2=P\)、\(\operatorname{tr}P=1\) 的实对称矩阵集合。
\end{exercise}
\begin{solution}
	代表向量乘以 \(\lambda\ne0\) 时，分子分母均乘以 \(\lambda^2\)，故 \(J\) 良定义。非零向量空间上的矩阵公式连续，所以 \(J\) 连续。矩阵 \(J([v])\) 的像空间为 \(\mathbb Rv\)：它把任意 \(w\) 送到 \(v(v^{\mathsf T}w)/(v^{\mathsf T}v)\)，且把 \(v\) 固定。因此若两个这样的矩阵相等，其像空间相同，两个射影点就相同，\(J\) 单射。源空间紧，矩阵空间 Hausdorff，其子空间 \(J(\mathbb RP^n)\) 也 Hausdorff，所以连续双射 \(\mathbb RP^n\to J(\mathbb RP^n)\) 为同胚。

	直接相乘得 \(J([v])^2=J([v])\)，对角线求和得迹为 \(1\)，且矩阵对称。反过来，若实对称矩阵 \(P\) 满足这两个条件，\cref{g1:thm:spectral}给出标准正交特征基；幂等性迫使每个特征值满足 \(\lambda^2=\lambda\)，故只能是 \(0\) 或 \(1\)。迹为 \(1\) 说明恰有一个特征值为 \(1\)。取其单位特征向量 \(u\)，则 \(P\) 在 \(u\) 上恒等、在 \(u^\perp\) 上为零，故 \(P=uu^{\mathsf T}=J([u])\)。这正是到一维子空间的正交投影矩阵。
\end{solution}

\begin{exercise}\label{g1:ex:double-origin-compact}
	在正文的双原点直线 \(D\) 中，令 \(j_a:\mathbb R\to D\) 为第 \(a\) 份直线的像映射，并取 \(K=j_a([-1,1])\)。证明 \(K\) 紧而不闭，并计算 \(\overline K\)。解释这个例子与“紧子集在 Hausdorff 空间中闭”的关系。
\end{exercise}
\begin{solution}
	闭区间紧，\(j_a\) 连续，故 \(K\) 紧。它含 \(0_a\) 而不含 \(0_b\)。任一 \(0_b\) 的开邻域都包含足够小的非零点，而这些点属于 \(K\)，所以由闭包的邻域判据，\(0_b\in\overline K\setminus K\)。因此 \(K\) 不闭。

	除 \(0_b\) 外，任何不在 \(K\) 中的点有一个非零坐标 \(x\) 且 \(|x|>1\)。在普通实直线上取足够小的开区间 \(I\ni x\)，使它与 \([-1,1]\) 不交且不含零。其在 \(D\) 中的像为开集，与 \(K\) 不交，故该点不在 \(\overline K\)。于是 \(\overline K=K\cup\{0_b\}\)。定理要求环境空间 Hausdorff，而 \(D\) 不满足这一条件；即使它每个单点闭且处处有直线的局部模型，也不能删去该假设。
\end{solution}

% !TeX root = ../../Book1.tex
% 纲要标题：## 第5章　基本群、覆盖与单值
% 纲要位置：几何学计划纲要.md，第 274 行。
% 正文状态：初稿；尚未正式审核。

% 纲要内容开始
% 
% > 本章目标与证明范围：连续证明道路提升、同伦提升、万有覆盖及覆盖分类。把“计算基本群”与“由基本群重建覆盖”连成双向过程，明确局部道路连通和半局部单连通条件。
% 
% 纲要内容结束

\chapter{基本群、覆盖与单值}
\label{g1:ch:05}
\BookChapterNumber{g1:ch:05}

\begin{chapterabstract}
	基本群把环路的连续变形变成群运算，覆盖则把同一底空间上的不同延拓结果放到不同的纤维点上。本章先从定端点同伦、圆周绕数和空间粘合计算基本群，再依次证明道路提升、同伦提升、万有覆盖及覆盖分类。局部道路连通与半局部单连通分别控制邻域中的道路和小环路的消除；在这些条件下，子群与纤维上的单值作用可以反过来重建覆盖。环面、曲面和图的例子将贯穿这两个方向。
\end{chapterabstract}

\begin{learninggoals}
	\begin{itemize}
		\item 区分道路与道路的同伦类，验证拼接、反路和基点变化，并用形变收缩与绕数计算基本群。
		\item 通过道路细分和方形网格证明 van Kampen 定理，解释粘合为何产生生成元和关系。
		\item 从覆盖片逐段构造提升，证明唯一性、同伦提升和基本群提升判据。
		\item 给道路类集合定义拓扑，证明万有覆盖与子群分类，并逐步指出两个局部假设的作用。
		\item 将右单值作用、陪集、覆盖变换相互对应，利用置换重建覆盖，计算其连通分支、叶数和对称。
	\end{itemize}
\end{learninggoals}

在一个小圆弧上，圆周上的点可以用一个连续变化的实角度表示。但如果绕圆周走完一周，角度增加了 \(2\pi\)，终点回到原处，实角度却没有回到原值。这里出现了两个相关的问题：一条环路在空间中能否缩掉？沿环路连续延拓某个局部选择后，为什么可能回到另一个值？前者由基本群记录，后者由覆盖与单值作用记录。

\ProseParagraphBreak

指数映射 \(t\mapsto e^{2\pi it}\) 把实直线投到圆周，每个充分小的圆弧上方都有一列彼此分开的区间。指定起始角度后，圆周上的道路可以唯一地在实直线上继续走；一周的环路会从一个整数走到下一个整数。这个模型会帮助我们理解一般覆盖的局部结构。局部逆沿道路的延续给出提升，要求它随参数连续变化又把同伦纳入讨论。把同终点道路按可忽略的环路关系合并后，能否得到新的覆盖，取决于这些道路类集合能否获得由局部逆描述的拓扑。

\ProseParagraphBreak

本章使用上一章的子空间、积、商、紧性与局部道路连通。群的同态、子群、正规子群与商群只需基本定义；自由群、自由积、群展示、陪集作用及正规化子将在使用时说明。道路的拼接始终按先走左边、再走右边的顺序，相关的单值采用右作用，以便把公式与实际行走保持一致。进一步阅读可参见 Hatcher 的正式著作\cite[第1.1--1.3节; 第1.A节]{Hatcher2002}。

% !TeX root = ../../Book1.tex
% 纲要标题：### 5.1 道路、基点与同伦运算
% 纲要位置：几何学计划纲要.md，第 278 行。
% 正文状态：初稿；尚未正式审核。

% 纲要内容开始
% 
% * 定端点同伦和道路复合
% * 基本群的函子性与基点变化
% * 形变收缩与基本群计算
% 
% 纲要内容结束

\section{道路、基点与同伦运算}
\label{g1:sec:0501}

在平面中绕原点走一圈，若允许穿过原点，就可以把这一圈缩成一个点；在去掉原点的平面里，同样的收缩却不可能完成。我们需要记录的不是道路的精确形状或速度，而是它在指定空间中能怎样连续变形。先固定起点与终点，再把可以这样变形的道路归为一类，拼接道路便产生了基本群的运算。

\subsection{定端点同伦和道路复合}
\label{g1:subsec:0501:01}

沿用\cref{g1:def:connected-path-connected}，道路是连续映射 \(\alpha:I\rightarrow X\)，其中 \(I=[0,1]\)。参数 \(t\) 记录行进时间，\(\alpha(0),\alpha(1)\) 是起终点。同一条几何曲线可以有不同走法，甚至可以往返或中途停顿；这些信息在判断能否变形时需要通过定义处理。

\begin{definition}[定端点同伦]\label{g1:def:path-homotopy}
	记 \(I=[0,1]\)。设 \(X\) 为拓扑空间，道路 \(\alpha,\beta:I\rightarrow X\) 都从 \(x\) 到 \(y\)。若存在连续映射 \(H:I\times I\rightarrow X\)，使
	\[
		H(t,0)=\alpha(t),\quad H(t,1)=\beta(t),\quad
		H(0,u)=x,\quad H(1,u)=y,
	\]
	则称它们\term[duandingdiantonglun]{定端点同伦}[homotopic relative to endpoints]，记为 \(\alpha\simeq\beta\;(\operatorname{rel}\partial I)\)。这里第二个参数 \(u\) 表示变形时间。若一条道路的起点与终点相同，则称它为以该点为基点的\term[huanlu]{环路}[loop]。
\end{definition}

下面用 \(\simeq\;(\operatorname{rel}\partial I)\) 时始终固定两端，不能让端点在变形中绕着空间移动。恒定同伦 \(H(t,u)=\alpha(t)\) 给出自反性，\(u\mapsto1-u\) 给出对称性；两个同伦各占变形时间的一半，再由\cref{g1:prop:gluing-lemma}粘合，给出传递性。因此定端点同伦是等价关系，其等价类记为 \([\alpha]\)。

\begin{definition}[拼接与反路]\label{g1:def:path-concatenation}
	设 \(\alpha\) 从 \(x\) 到 \(y\)，\(\beta\) 从 \(y\) 到 \(z\)。定义它们的\term[daolufuhe]{道路复合}[path concatenation]为
	\[
		(\alpha*\beta)(t)=
		\begin{cases}
			\alpha(2t),&0\le t\le1/2,\\
			\beta(2t-1),&1/2\le t\le1.
		\end{cases}
	\]
	本章始终约定 \(\alpha*\beta\) 先走 \(\alpha\)，后走 \(\beta\)。反向道路记为 \(\bar\alpha(t)=\alpha(1-t)\)，点 \(x\) 处的常值道路记为 \(e_x\)。
\end{definition}

复合在接点处连续，因为两段都取值 \(y\)；有限闭覆盖粘合引理保证整个道路连续。这里的星号是道路的先后拼接，连续映射的函数复合 \(g\circ f\) 则仍表示先 \(f\) 后 \(g\)。把两种约定分清，后面的群作用才不会发生方向颠倒。

\begin{proposition}[道路类的运算]\label{g1:prop:path-classes}
	记 \(I=[0,1]\)。在任意拓扑空间 \(X\) 中，道路复合在定端点同伦类上良定义，并满足结合律；常值道路类是单位，反路类是逆。具体地，若 \(\alpha,\beta\) 都从 \(x\) 到 \(y\)，则
	\[
		\alpha\simeq\beta\;(\operatorname{rel}\partial I)
		\quad\Longleftrightarrow\quad
		\alpha*\bar\beta\simeq e_x\;(\operatorname{rel}\partial I).
	\]
\end{proposition}
\begin{proof}
	若 \(H\) 把 \(\alpha\) 变到 \(\alpha'\)，\(K\) 把 \(\beta\) 变到 \(\beta'\)，则在 \(t\le1/2\) 使用 \(H(2t,u)\)，在 \(t\ge1/2\) 使用 \(K(2t-1,u)\)。两式在接缝恒取共同端点，因此粘成从 \(\alpha*\beta\) 到 \(\alpha'*\beta'\) 的定端点同伦。

	先说明改变速度为何不影响道路类。若 \(\rho:I\rightarrow I\) 连续且 \(\rho(0)=0,\rho(1)=1\)，则
	\[
		H(t,u)=\alpha\bigl((1-u)t+u\rho(t)\bigr)
	\]
	把 \(\alpha\) 定端点同伦到 \(\alpha\circ\rho\)。对三条依次相接的道路，先令 \(P\) 在三个等长时间段依次走完它们。道路 \((\alpha*\beta)*\gamma\) 与 \(\alpha*(\beta*\gamma)\) 都是 \(P\) 的分段线性重参数化：它们依次分配的时间分别是 \((1/4,1/4,1/2)\) 与 \((1/2,1/4,1/4)\)。上述同伦使两者的类都等于 \([P]\)，得到结合律。类似地，\(\alpha*e_y\) 与 \(e_x*\alpha\) 分别由 \(\rho(t)=\min\{2t,1\}\) 与 \(\rho(t)=\max\{2t-1,0\}\) 得到，所以常值类为单位。

	最后，沿道路出去再原路返回可以逐渐缩短。公式
	\[
		H(t,u)=\alpha\bigl((1-u)\min\{2t,2-2t\}\bigr)
	\]
	在 \(u=0\) 给出 \(\alpha*\bar\alpha\)，在 \(u=1\) 给出 \(e_x\)，两端始终为 \(x\)。应用于 \(\bar\alpha\) 又得到另一侧逆关系。于是所有上述恒等式在道路类上成立。若 \([\alpha]=[\beta]\)，显然 \([\alpha*\bar\beta]=[e_x]\)；反之，在此等式右边拼上 \([\beta]\)，用结合律和逆关系消去 \([\bar\beta*\beta]\)，就得到 \([\alpha]=[\beta]\)。
\end{proof}

这些运算允许起终点各不相同，只有可接续的道路才能相乘。若只保留从某个固定点出发并回到该点的道路，所有乘积都有定义，就得到一个群。

\begin{definition}[基本群]\label{g1:def:fundamental-group}
	设 \(X\) 为拓扑空间，\(x_0\in X\)。以 \(x_0\) 为基点的环路的定端点同伦类，在乘法
	\[
		[\alpha][\beta]=[\alpha*\beta]
	\]
	下组成的群称为\term[jibenqun]{基本群}[fundamental group]，记为 \(\pi_1(X,x_0)\)。单位为 \([e_{x_0}]\)，逆元为 \([\bar\alpha]\)。与常值环路定端点同伦的环路称为\term[linglun]{零伦}[null-homotopic]，即它在基本群中代表单位元；后文说“定基点零伦”时特别强调收缩中固定基点。
\end{definition}

\begin{definition}[单连通空间]\label{g1:def:simply-connected}
	非空拓扑空间 \(X\) 若道路连通，且在每个点 \(x\) 都有 \(\pi_1(X,x)=\{1\}\)，称为\term[danliantong]{单连通}[simply connected]。因此单连通不仅要求没有不能收缩的环路，也要求任意两点可以用道路连接。
\end{definition}

由\cref{g1:prop:path-classes}，单连通等价于：空间非空、道路连通，而且同一起终点之间的任意两条道路定端点同伦。下一个小节还会说明，道路连通时只需在一个基点核验基本群平凡。所谓“没有洞”只能作为某些图形的提示，不能代替这个定义。

\subsection{基本群的函子性与基点变化}
\label{g1:subsec:0501:02}

连续映射把一条道路送到另一条道路，也把一整个连续变形送到另一个连续变形。因此基本群能够记录空间之间的映射，而不只是分别给每个空间贴上一个群的名称。

\begin{proposition}[诱导同态]\label{g1:prop:homotopy-functor}
	设 \(f:X\rightarrow Y\) 连续，\(x_0\in X\)。公式
	\[
		f_*:\pi_1(X,x_0)\longrightarrow\pi_1(Y,f(x_0)),\qquad
		f_*[\alpha]=[f\circ\alpha]
	\]
	给出群同态。对可复合的连续映射 \(f,g\)，有 \((g\circ f)_*=g_*\circ f_*\)，且 \((\operatorname{id}_X)_*=\operatorname{id}\)。若 \(f,g:X\rightarrow Y\) 有同伦 \(H\)，且 \(H(x_0,u)\) 恒定，则 \(f_*=g_*\)。
\end{proposition}
\begin{proof}
	若两条道路由 \(K\) 定端点同伦，则 \(f\circ K\) 是像道路之间的定端点同伦，所以公式良定义。等式 \(f\circ(\alpha*\beta)=(f\circ\alpha)*(f\circ\beta)\) 说明它保乘法。复合与恒等映射的结论直接逐个道路类计算即可。最后，\((t,u)\mapsto H(\alpha(t),u)\) 给出 \(f\circ\alpha\) 与 \(g\circ\alpha\) 的同伦；由于基点在整个变形中固定，这个同伦固定两端，得到所需等式。
\end{proof}

这些与映射复合相容的性质称为基本群的\term[hanzixing]{函子性}[functoriality]。特别是同胚诱导群同构。但两个空间的基本群同构并不自动给出空间同胚，后面的形变收缩会提供许多这样的例子。

\begin{proposition}[沿道路改变基点]\label{g1:prop:basepoint-change}
	设 \(c\) 为拓扑空间 \(X\) 中从 \(x_0\) 到 \(x_1\) 的道路。则
	\[
		c_\#:\pi_1(X,x_0)\longrightarrow\pi_1(X,x_1),\qquad
		c_\#[\gamma]=[\bar c*\gamma*c]
	\]
	为同构，其逆是 \(\bar c_\#\)。它只依赖 \(c\) 的定端点同伦类，一般会依赖这条道路的选择。
\end{proposition}
\begin{proof}
	公式先从新基点 \(x_1\) 沿反路回到 \(x_0\)，走完 \(\gamma\)，再沿 \(c\) 返回，所以确实是新基点处的环路。道路类运算说明它良定义。计算两个像的乘积时，中间的 \(c*\bar c\) 可以缩成常值道路，得到
	\[
		[\bar c*\gamma*c][\bar c*\eta*c]=[\bar c*\gamma*\eta*c].
	\]
	因此它是同态。先沿 \(c\) 改基点再沿 \(\bar c\) 改回来，出现的往返道路都可消去，所以这两个同态互逆。

	若 \(d\) 是另一条从 \(x_0\) 到 \(x_1\) 的道路，令 \(k=[\bar d*c]\in\pi_1(X,x_1)\)。直接消去相邻的往返段得到
	\[
		c_\#[\gamma]=k^{-1}\,d_\#[\gamma]\,k.
	\]
	故两种基点变化相差一个共轭同构；这并不保证它们作为具体同态相等。
\end{proof}

这也解释了为何覆盖分类中“选定一个纤维点”与“不选纤维点”给出不同的分类数据：前者会确定一个具体子群，后者通常只确定它的共轭类。

\begin{proposition}[基点移动的同伦公式]\label{g1:prop:moving-basepoint}
	记 \(I=[0,1]\)，设 \(X,Y\) 为拓扑空间，\(H:X\times I\rightarrow Y\) 是连续映射 \(f,g:X\rightarrow Y\) 之间的同伦，即 \(H(x,0)=f(x)\)、\(H(x,1)=g(x)\)。固定 \(x_0\in X\)，令 \(c(u)=H(x_0,u)\)。则
	\[
		g_*=c_\#\circ f_*:\pi_1(X,x_0)\longrightarrow\pi_1(Y,g(x_0)).
	\]
\end{proposition}
\begin{proof}
	对以 \(x_0\) 为基点的环路 \(\alpha\)，考虑方形映射 \(K(t,u)=H(\alpha(t),u)\)。沿方形边界从 \((0,0)\) 依次走下、右、上、左四边，其像为
	\[
		(f\circ\alpha)*c*\overline{(g\circ\alpha)}*\bar c.
	\]
	这个边界环路在方形中可以固定 \((0,0)\) 收缩：若边界参数化为 \(b(t)\)，取 \((1-v)b(t)\) 即可。与 \(K\) 复合后，像环路也可收缩。由道路类的逆与消去律，
	\[
		[f\circ\alpha]\,[c]=[c]\,[g\circ\alpha],
	\]
	此式在可接续道路类的意义下成立。左边再拼 \([\bar c]\)，便得到 \([g\circ\alpha]=[\bar c*(f\circ\alpha)*c]\)。
\end{proof}

若忽略 \(c\)，就会把不同基点处的群同态误写成相等。尤其是一般同伦不一定固定基点；使用同伦不变性前，应先看基点是否真的留在原处。

\subsection{形变收缩与基本群计算}
\label{g1:subsec:0501:03}

许多空间虽不与较简单的空间同胚，却能连续地缩到后者上，而且整个过程中保留后者不动。这使基本群计算可以先去掉不影响环路的多余方向。圆环的径向厚度就是一个例子；圆周自身绕行的次数却不能用这种方式消去。

\begin{definition}[形变收缩]\label{g1:def:deformation-retract}
	记 \(I=[0,1]\)。设 \(X\) 为拓扑空间，\(A\subseteq X\) 带子空间拓扑，\(i:A\hookrightarrow X\) 为包含映射。连续映射 \(r:X\rightarrow A\) 若满足 \(r|_A=\operatorname{id}_A\)，称为\term[suohui]{缩回}[retraction]。若还存在连续 \(H:X\times I\rightarrow X\)，满足
	\[
		H(x,0)=x,\qquad H(x,1)=i(r(x)),\qquad H(a,u)=a\quad(a\in A),
	\]
	则称 \(A\) 是 \(X\) 的\term[xingbiansuosuo]{形变收缩核}[deformation retract]。本书此处要求同伦逐点固定 \(A\)，也称强形变收缩。
\end{definition}

\begin{proposition}[形变收缩保持基本群]\label{g1:prop:deformation-retract}
	设 \(X\) 为拓扑空间，\(A\subseteq X\) 是强形变收缩核，\(i:A\hookrightarrow X\) 为包含映射，\(r:X\to A\) 为该形变收缩中的缩回，\(a_0\in A\)。则 \(i_*:\pi_1(A,a_0)\rightarrow\pi_1(X,a_0)\) 与 \(r_*\) 互为逆同构。
\end{proposition}
\begin{proof}
	因为 \(r\circ i=\operatorname{id}_A\)，函子性给出 \(r_*i_*=\operatorname{id}\)。因为 \(i\circ r\) 与 \(\operatorname{id}_X\) 的同伦固定 \(a_0\)，\cref{g1:prop:homotopy-functor}又给出 \(i_*r_*=\operatorname{id}\)。这两条等式同时给出单射与满射。
\end{proof}

凸集 \(C\subseteq\mathbb R^n\) 中选一点 \(a\)，公式 \(H(x,u)=(1-u)x+ua\) 把它形变收缩到 \(a\)，所以非空凸集单连通。对穿孔空间 \(\mathbb R^n\setminus\{0\}\)，则可取
\[
	r(x)=\frac{x}{\lVert x\rVert},\qquad
	H(x,u)=\left(1-u+\frac{u}{\lVert x\rVert}\right)x.
\]
系数始终正，故过程不经过被去掉的原点；当 \(\lVert x\rVert=1\) 时所有点不动。因此它形变收缩到 \(S^{n-1}\)。尤其穿孔平面的计算化成了圆周的计算。

\begin{lemma}[紧空间上的一致细分尺度]\label{g1:lem:lebesgue-number}
	设 \((K,d)\) 为紧度量空间，\(\mathcal U\) 是它的开覆盖。则存在 \(\delta>0\)，使 \(K\) 中每个直径小于 \(\delta\) 的非空子集都包含在某个 \(U\in\mathcal U\) 中。
\end{lemma}

这样的 \(\delta\) 称为覆盖的\term[lebesgueshu]{Lebesgue 数}[Lebesgue number]。\footnote{Henri Léon Lebesgue（1875--1941），法国数学家，建立了测度与积分理论的重要基础。}它把“每个点附近可以处理”变成“细分得足够小时，每一小块都可以处理”，是后面把同伦方形切成有限网格的依据。

\begin{proof}
	空空间任取 \(\delta>0\)。否则，对每个 \(x\in K\) 选 \(r_x>0\) 与覆盖成员 \(U_x\)，使 \(B(x,2r_x)\subseteq U_x\)。由紧性，有限个 \(B(x_j,r_{x_j})\) 覆盖 \(K\)。取 \(\delta=\min_jr_{x_j}>0\)。若非空集合 \(A\) 的直径小于 \(\delta\)，选 \(a\in A\) 和满足 \(a\in B(x_j,r_{x_j})\) 的 \(j\)。对每个 \(b\in A\)，
	\[
		d(b,x_j)\le d(b,a)+d(a,x_j)<\delta+r_{x_j}\le2r_{x_j},
	\]
	故 \(A\subseteq B(x_j,2r_{x_j})\subseteq U_{x_j}\)。
\end{proof}

下面直接从连续角度计算圆周的基本群。局部角度可以选择，但走完一圈后不一定回到原角度；这个整周差正是要记录的整数。写
\[
	E:\mathbb R\longrightarrow S^1,\qquad E(s)=e^{2\pi is}.
\]
对每个 \(z\in S^1\)，其开半圆 \(A_z=\{w:\operatorname{Re}(w/z)>0\}\) 有连续角度
\[
	\theta_z(w)=\frac{1}{2\pi}\arctan\frac{\operatorname{Im}(w/z)}{\operatorname{Re}(w/z)}
	\in(-1/4,1/4),\qquad E(\theta_z(w))=w/z.
\]
这是\secref{g1:sec:0405}中给圆周作局部坐标时已经使用的显式公式。

\begin{lemma}[沿圆周道路连续选择角度]\label{g1:lem:circle-angle-lift}
	记 \(I=[0,1]\)，\(E:\mathbb R\to S^1\) 为 \(E(t)=e^{2\pi it}\)。设 \(\alpha:I\rightarrow S^1\) 连续，\(a\in\mathbb R\) 满足 \(E(a)=\alpha(0)\)。则存在唯一连续 \(\widetilde\alpha:I\rightarrow\mathbb R\)，满足 \(\widetilde\alpha(0)=a\)、\(E\circ\widetilde\alpha=\alpha\)。
\end{lemma}
\begin{proof}
	开集 \(\alpha^{-1}(A_z)\) 覆盖紧区间。由\cref{g1:lem:lebesgue-number}，可把 \(I\) 分成有限个闭小区间，使每一段的像落在某个 \(A_z\) 中。在第一段上，选 \(b\) 使 \(E(b)=z\)，用 \(b+\theta_z(\alpha(t))\) 表示角度，再加上唯一适当整数使它在起点等于 \(a\)。在下一段重复这个操作，把起点值接到前一段的终点值。有限闭集粘合保证所得函数连续，并满足公式。

	若另有 \(\widetilde\alpha'\) 满足要求，则它们的差连续且处处为整数，因为 \(E(s)=E(t)\) 当且仅当 \(s-t\in\mathbb Z\)。连续整数值函数在连通区间上恒定：每个整数的原像都开且闭，若有两个取值便会分离区间。起点之差为零，故处处为零。
\end{proof}

\begin{theorem}[圆周的基本群]\label{g1:thm:circle-fundamental}
	设 \(\alpha:[0,1]\to S^1\) 为以 \(1\) 为基点的环路，令 \(\widetilde\alpha:[0,1]\to\mathbb R\) 为满足 \(e^{2\pi i\widetilde\alpha(t)}=\alpha(t)\)、\(\widetilde\alpha(0)=0\) 的唯一连续实值函数。定义
	\[
		w(\alpha)=\widetilde\alpha(1)\in\mathbb Z.
	\]
	这个整数只依赖环路的定端点同伦类，并给出同构
	\[
		w:\pi_1(S^1,1)\longrightarrow\mathbb Z.
	\]
	它称为环路的\term[raoshu]{绕数}[winding number]；\(\alpha_n(t)=e^{2\pi int}\) 对应整数 \(n\)。
\end{theorem}
\begin{proof}
	因为 \(\alpha(1)=1\)，角度终值必为整数。先证明同伦不改变这个整数。若两条基点为 \(1\) 的环路 \(\alpha,\beta\) 满足 \(|\beta(t)-\alpha(t)|<1\) 对所有 \(t\) 成立，则 \(\beta(t)/\alpha(t)\) 落在 \(1\) 的开半圆内，可以写
	\[
		d(t)=\theta_1(\beta(t)/\alpha(t)),\qquad d(0)=d(1)=0.
	\]
	故 \(\widetilde\alpha+d\) 是 \(\beta\) 从零开始的角度，由唯一性等于 \(\widetilde\beta\)，两者终值相同。若 \(H(t,u)\) 为定端点同伦，将 \(S^1\) 用半径 \(1/3\) 的圆周开球覆盖，把它们拉回到紧方形，再取一个 Lebesgue 数 \(\delta>0\)。方形内距离小于 \(\delta\) 的两点落在同一个原像中，故其像的距离小于 \(2/3\)。于是把 \(u\) 区间细分成长度小于 \(\delta\) 的小段，相邻两条环路在所有 \(t\) 处的距离便都小于 \(1\)。逐段应用刚才的论证，得到两端绕数相同。

	拼接 \(\alpha*\beta\) 时，前半段角度走完 \(\widetilde\alpha\)，后半段应走 \(\widetilde\beta+w(\alpha)\)，才能在接点连续。因此 \(w(\alpha*\beta)=w(\alpha)+w(\beta)\)，所以 \(w\) 是群同态。每个 \(n\in\mathbb Z\) 由 \(\alpha_n\) 实现，故它满射。

	若 \(w(\alpha)=0\)，则 \(\widetilde\alpha\) 的两端都为零，公式
	\[
		H(t,u)=E\bigl((1-u)\widetilde\alpha(t)\bigr)
	\]
	把 \(\alpha\) 定端点收缩为常值道路。故同态的核只有单位，得到单射。更具体地，任一绕数为 \(n\) 的环路可由角度的线性插值
	\[
		E\bigl((1-u)\widetilde\alpha(t)+unt\bigr)
	\]
	定端点变到标准环路 \(\alpha_n\)，这直接说明整数已经记录了全部环路类。
\end{proof}

绕数不是某段参数上的角度，而是沿整个环路连续记录后的总增量。因此不存在连续函数 \(L:S^1\rightarrow\mathbb R\) 满足 \(E\circ L=\operatorname{id}_{S^1}\)：若存在，把它复合到标准一周环路，所得角度的起终值必相同；但角度唯一性又要求终值比起值大 \(1\)。连续选择局部角度没有问题，整体单值选择却受到绕行的阻碍。

\begin{proposition}[积空间的基本群]\label{g1:prop:product-fundamental}
	设 \(X,Y\) 为拓扑空间，\(x_0\in X,y_0\in Y\)。坐标投影诱导同构
	\[
		\pi_1(X\times Y,(x_0,y_0))
		\longrightarrow\pi_1(X,x_0)\times\pi_1(Y,y_0),\qquad
		[(\alpha,\beta)]\longmapsto([\alpha],[\beta]).
	\]
\end{proposition}
\begin{proof}
	由\cref{g1:prop:product-universal}，到积的道路连续当且仅当两条坐标道路连续；到积的同伦也恰好等于两份同时给出的坐标同伦。因此两条积道路同伦当且仅当各坐标分别同伦，所写映射良定义且单射。任意两条坐标环路可以同时参数化为 \(t\mapsto(\alpha(t),\beta(t))\)，所以满射。道路拼接逐坐标进行，故它保群运算。
\end{proof}

于是由圆周绕数与\cref{g1:thm:torus-square}的环面模型，
\[
	\pi_1(\mathbb R^2\setminus\{0\},(1,0))\cong\mathbb Z,\qquad
	\pi_1(T^2,(1,1))\cong\mathbb Z^2.
\]
前一个同构使用径向形变收缩；后一个把两个圆周方向的绕数分别记录下来。这些计算还没有说明：空间由几块拼合而成时，怎样根据各块的基本群求出整体的基本群。下一节将处理这种需要增加生成元与关系的情形。

% !TeX root = ../../Book1.tex
% 纲要标题：### 5.2 空间粘合与 van Kampen 定理
% 纲要位置：几何学计划纲要.md，第 284 行。
% 正文状态：初稿完成；尚未正式审核。

% 纲要内容开始
% 
% * 自由积及带关系群的最低准备
% * 细分道路和开覆盖的证明
% * 多重交集、基点选择与推广范围
% * 图、曲面和去点空间的群展示
% 
% 纲要内容结束

\section{空间粘合与 van Kampen 定理}
\label{g1:sec:0502}

圆周上的绕数记录绕行次数，但两个圆周在一点相接以后，仅记录各绕了几圈便不够了：先走左圈再走右圈，与先右后左未必能够互相变形。另一方面，若在这两个圆周上再粘一张圆盘，沿圆盘边界走出的某个词就会变成可以收缩的闭路。基本群因而同时记录“可以怎样走”和“哪些走法被填上的区域消去”。本节把这两件事分别写成生成元与关系，再证明它们确实由空间的开覆盖决定。

\subsection{自由积及带关系群的最低准备}
\label{g1:subsec:0502:01}

我们只补充当前计算所需的群论。群的单位元记为 \(1\)，群乘法的书写顺序沿用道路约定：\([\alpha][\beta]=[\alpha*\beta]\) 表示先走 \(\alpha\)，再走 \(\beta\)。一个词是有限个字母排成的序列；空词也允许，长度为零。

\begin{definition}[自由群]\label{g1:def:free-group}
	给定符号集合 \(S\)，对每个 \(s\in S\) 增加一个形式逆符号 \(s^{-1}\)。删去词中相邻的 \(ss^{-1}\) 或 \(s^{-1}s\)，直到不能再删，所得词称为约化词。所有约化词组成的\term[ziyouqun]{自由群}[free group]记为 \(F(S)\)：乘法是先连接两个词再约化，单位元为空词，逆元是倒转字母顺序并逐字取逆。若 \(S=\{a_1,\ldots,a_r\}\)，也记 \(F_r\)。
\end{definition}

这个定义需要排除约化顺序带来的歧义。每次删除都使长度下降。若两次可作的删除位置不交，先后次序可交换；若位置相交，局部词只能是 \(ss^{-1}s\) 或 \(s^{-1}ss^{-1}\)，两种删除得到同一个单字母。对词长归纳：任意两条约化过程的第一步或者相同，或者经过上述交换抵达同一较短词，再用归纳假设，最终约化词一致。因此连接后约化是良定义的，而且把三个词一起连接时，无论先约化哪一段，最后结果相同，这就证明了结合律。逆元公式通过逐对消去验证。

任意指定一族元素 \(h_s\in H\)，其中 \(H\) 为群，都恰好确定一个同态 \(F(S)\to H\)：将 \(s\) 替换为 \(h_s\)，将 \(s^{-1}\) 替换为 \(h_s^{-1}\)，再按顺序相乘。删除相邻逆字母不改变乘积，故映射良定义；同态性来自连接词，唯一性来自每个词都由这些字母组成。这就是“自由”的精确含义：除群公理迫使成立的等式外，没有另加关系。特别地，\(F(\{a\})\cong\mathbb Z\)，而在 \(F(\{a,b\})\) 中，\(ab\) 与 \(ba\) 是不同的约化词。

空间的不同部分往往已经有各自的基本群，我们希望先保留每部分原有的关系，再把它们放在一起，而不提前要求不同部分的元素交换。

\begin{definition}[自由积]\label{g1:def:free-product}
	设 \((G_\alpha)_{\alpha\in A}\) 为一族群。它们的\term[ziyouji]{自由积}[free product] \(\mathop{*}_{\alpha\in A}G_\alpha\) 由下列约化词组成：每个字母是某个 \(G_\alpha\) 的非单位元，相邻字母来自不同的群。乘法是连接词，再把同一群中相邻的两个字母按该群的乘法合并，并删去出现的单位元。不同群中的字母保留所属群的标签。
\end{definition}

约化的唯一性可用与自由群相同的长度归纳证明。两处不相交的合并可以交换；相交时涉及同一群中的三个连续元素，两种合并最终一致正是该群的结合律，乘积为单位元时继续删除即可。因此这确实给出一个群，每个 \(G_\alpha\) 通过单字母词嵌入其中。对任一群 \(H\) 及同态 \(f_\alpha:G_\alpha\to H\)，公式
\[
	f(g_1\cdots g_m)=f_{\alpha_1}(g_1)\cdots f_{\alpha_m}(g_m)
\]
与约化相容，给出唯一的同态 \(f:\mathop{*}_\alpha G_\alpha\to H\)。两群情形的泛性质写成下图，其中两条斜箭头是给定的同态，虚线箭头是唯一延拓：
\[
	\begin{tikzcd}[column sep=large]
		G \arrow[r,"j_G"] \arrow[dr,"f_G"'] & G*H \arrow[d,dashed,"f"] & H \arrow[l,"j_H"'] \arrow[dl,"f_H"]\\
		& K &
	\end{tikzcd}
\]
例如 \(\mathbb Z*\mathbb Z=F_2\)，与直积 \(\mathbb Z\times\mathbb Z\) 不同：直积中两坐标的元素交换，自由积中一般不交换。

接着说明怎样加入关系。子群 \(N\subseteq G\) 称为\term[zhengguiziqun]{正规子群}[normal subgroup]，若对所有 \(g\in G\) 有 \(gNg^{-1}=N\)。这个条件保证陪集乘法 \((gN)(hN)=ghN\) 不依赖代表，从而定义\term[shangqun]{商群}[quotient group] \(G/N\)。给定 \(R\subseteq G\)，要使每个 \(r\in R\) 都成为单位元，必须同时消去它们的逆、乘积和所有共轭。因此包含 \(R\) 的最小正规子群为
\[
	\langle\!\langle R\rangle\!\rangle
	=\left\{\prod_{j=1}^{m}g_jr_j^{\varepsilon_j}g_j^{-1}:
	m\ge0,\ r_j\in R,\ \varepsilon_j=\pm1,\ g_j\in G\right\}.
\]
右边对乘法、逆及共轭封闭，且任何包含 \(R\) 的正规子群都包含右边，因而确实具有所说的最小性。它称为 \(R\) 的\term[zhengguibibao]{正规闭包}[normal closure]。

\begin{definition}[群展示]\label{g1:def:group-presentation}
	若 \(R\) 是自由群 \(F(S)\) 中的一组词，则
	\[
		\langle S\mid R\rangle=F(S)/\langle\!\langle R\rangle\!\rangle
	\]
	称为由生成元 \(S\) 和关系词 \(R\) 给出的\term[qunzhanshi]{群展示}[group presentation]。写关系 \(u=v\) 等价于把 \(uv^{-1}\) 加入关系词。
\end{definition}

由最小性可知，一个映射 \(S\to H\) 恰好在所有关系词映成 \(1\) 时延拓为 \(\langle S\mid R\rangle\to H\) 的同态。这个判据比“按关系化简到某种样子”更可靠，因为一般的群展示未必提供容易执行的唯一约化算法。例如
\[
	\langle a\mid a^m=1\rangle\cong\mathbb Z/m\mathbb Z\quad(m\ge1),\qquad
	\langle a,b\mid aba^{-1}b^{-1}=1\rangle\cong\mathbb Z^2.
\]
第二个同构可以直接验证：关系使 \(a,b\) 交换，所以每个词都能写成 \(a^pb^q\)；映射 \(a\mapsto(1,0)\)、\(b\mapsto(0,1)\) 给出同态到 \(\mathbb Z^2\)，并证明不同的整数对不会再被识别。这里既证明了这些表达能够覆盖所有元素，也证明了没有额外关系。

\subsection{细分道路和开覆盖的证明}
\label{g1:subsec:0502:02}

设 \(X=U\cup V\)。一条闭路未必全部落在 \(U\) 或 \(V\) 中，但紧参数区间可以切成有限段，每段落在其中一个开集。切出的段一般不是闭路；我们必须从公共基点走到各分点，再沿原路段前进，最后返回基点。不同接回方式之间的差别恰好发生在交集内，这便说明为什么群的关系应来自 \(U\cap V\)。

\begin{theorem}[van Kampen 定理]\label{g1:thm:van-kampen}
	设拓扑空间 \(X=U\cup V\)，其中 \(U,V\) 都开且道路连通，\(U\cap V\) 也道路连通，并取 \(x_0\in U\cap V\)。记包含映射诱导的同态为
	\[
		i_U:\pi_1(U\cap V,x_0)\to\pi_1(U,x_0),\qquad
		i_V:\pi_1(U\cap V,x_0)\to\pi_1(V,x_0).
	\]
	则两包含映射诱导满同态
	\[
		\Phi:\pi_1(U,x_0)*\pi_1(V,x_0)\longrightarrow\pi_1(X,x_0),
	\]
	其核恰为所有 \(i_U(\omega)i_V(\omega)^{-1}\) 的正规闭包。因此
	\begin{equation}\label{g1:eq:van-kampen-quotient}
		\pi_1(X,x_0)\cong
		\frac{\pi_1(U,x_0)*\pi_1(V,x_0)}
		{\langle\!\langle i_U(\omega)i_V(\omega)^{-1}:\omega\in\pi_1(U\cap V,x_0)\rangle\!\rangle}.
	\end{equation}
\end{theorem}

这就是\term[vankampendingli]{van Kampen 定理}[van Kampen theorem]。\footnote{Egbert Rudolf van Kampen（埃格伯特·鲁道夫·范坎彭，1908--1942），生于比利时，在荷兰求学，后在美国任教，研究拓扑与拓扑群。} 其中的关系只要求把交集中的同一条闭路在两边识别，并不声称 \(i_U,i_V\) 单射。例如交集中的圆圈可能已经在 \(U\) 内围成一个圆盘，此时它在 \(V\) 中的像也会被关系消去。式\eqref{g1:eq:van-kampen-quotient}有时称为沿交集群的合并自由积，但计算时仍应保留两条具体同态。

\begin{proof}
	记式\eqref{g1:eq:van-kampen-quotient}右边的群为 \(P\)，两因子到 \(P\) 的同态为 \(j_U,j_V\)。交集中的同一闭路在 \(X\) 中当然给出同一元素，故 \(\Phi\) 下降为同态 \(\bar\Phi:P\to\pi_1(X,x_0)\)。我们将对闭路构造反向映射，并逐一说明分段、连接道路和同伦都不会改变结果。

	先取闭路 \(\gamma:[0,1]\to X\)。开覆盖 \(\{\gamma^{-1}(U),\gamma^{-1}(V)\}\) 由\cref{g1:lem:lebesgue-number}有 Lebesgue 数，所以存在分点 \(0=t_0<\cdots<t_m=1\)，使每段 \(\gamma_k=\gamma|_{[t_{k-1},t_k]}\) 的像落在一个标记开集 \(A_k\in\{U,V\}\) 中；各段均线性重新参数化为 \([0,1]\)。在分点 \(y_k=\gamma(t_k)\) 处选道路 \(c_k:x_0\to y_k\)：若相邻标记不同，在 \(U\cap V\) 中选择；若相同，在该标记开集中选择。取 \(c_0,c_m\) 为常值道路。道路连通的假设保证这些选择存在。令
	\[
		W(\gamma)=\prod_{k=1}^{m}j_{A_k}\bigl([c_{k-1}*\gamma_k*\bar c_k]\bigr)\in P.
	\]
	每个括号内都是 \(A_k\) 中以 \(x_0\) 为基点的闭路。把它们在 \(X\) 中连接，连续出现的 \(\bar c_k*c_k\) 可相对端点缩去，因此 \(\bar\Phi(W(\gamma))=[\gamma]\)。

	若将 \(c_k\) 改为另一条允许的 \(d_k\)，记 \(h_k=[c_k*\bar d_k]\)。在相邻标记相同时，\(h_k\) 属于同一个因子；标记不同时，两条连接道路都在交集中，故 \(h_k\) 在两个因子中的像由 \(P\) 的关系识别。替换后第 \(k\) 项变为
	\[
		h_{k-1}^{-1}\,j_{A_k}([c_{k-1}*\gamma_k*\bar c_k])\,h_k,
	\]
	这里将 \(h_k\) 也记作其在 \(P\) 中的像。相邻的 \(h_kh_k^{-1}\) 依次消去，且 \(h_0=h_m=1\)，故 \(W(\gamma)\) 不变。

	在一段内部增加分点时，只需在同一个标记开集中增加一条连接道路；插入的反路与正路相消，所以值仍不变。如果某一段同时位于 \(U,V\)，要改变它的标记，可以先把两端连接道路都改选在 \(U\cap V\) 中。这在原标记和新标记下都允许，且刚刚已证明不影响乘积；随后该段对应的闭路全在交集中，换标记正是 \(P\) 中的一条关系。任意两种分段都有共同细分，细分后的标记可逐段改变。因此 \(W(\gamma)\) 不依赖这些选择。

	还须证明相对端点同伦不改变 \(W\)。设 \(H:[0,1]^2\to X\) 是 \(\gamma_0\) 到 \(\gamma_1\) 的同伦，且 \(H(0,t)=H(1,t)=x_0\)。对方形上的拉回开覆盖应用 Lebesgue 数引理，取足够细的矩形网格，使每个闭小矩形的像全在 \(U\) 或 \(V\) 中，并给它相应的标记。对每个网格顶点 \(v\)，选道路 \(c_v:x_0\to H(v)\)：若相邻小矩形出现两种标记，就在交集中选；否则在唯一的标记开集中选。左右边界顶点的 \(c_v\) 都取常值。对从顶点 \(v\) 到 \(w\) 的有向网格边 \(e\)，记
	\[
		E(e)=j_A([c_v*(H|_e)*\bar c_w])\in P,
	\]
	其中 \(A\) 是任一相邻小矩形的标记。若有两种标记，该闭路全在交集中，故两种取法给出同一元素；反向边的元素为 \(E(e)^{-1}\)。

	对任一小矩形，把下边、右边、上边、左边分别按向右、向上、向右、向上定向，得到元素 \(b,r,t,\ell\)。将矩形边界沿矩形内的直线缩到左下角，再复合 \(H\)，便给出该标记开集内保持基点的收缩；其基点通过 \(c_v\) 接到 \(x_0\) 不改变可收缩性。消去中间连接道路便得
	\[
		brt^{-1}\ell^{-1}=1,\qquad\text{即 }br=\ell t.
	\]
	这条等式允许把网格路径的“下边后右边”换成“左边后上边”。逐格从最下一行向上移动：同一行中相邻矩形的公共竖边一正一反，替换时相消；扫完整行后，下面的水平边乘积换成上面的水平边乘积，外侧只留下左右边界的竖边。后者经 \(H\) 都是常值，元素为 \(1\)。逐行应用，得到整个方形底边与顶边的乘积相等，即 \(W(\gamma_0)=W(\gamma_1)\)。\cref{g1:fig:van-kampen-grid}显示的正是这个局部替换，而非交换不相邻的群元素。

	因此 \(W\) 给出映射 \(\Psi:\pi_1(X,x_0)\to P\)。连接两条闭路时，把两套分段接在一起，并在接合处取常值连接道路，就有 \(W(\alpha*\beta)=W(\alpha)W(\beta)\)，故 \(\Psi\) 为同态。已经证明 \(\bar\Phi\Psi\) 恒等。若闭路全在 \(U\) 中，取只有一段的分割便有 \(\Psi\bar\Phi(j_U([\gamma]))=j_U([\gamma])\)，\(V\) 同理；两因子的像生成 \(P\)，所以 \(\Psi\bar\Phi\) 也恒等。于是 \(\bar\Phi\) 是同构，原同态的核恰为所述正规闭包。
\end{proof}

\begin{figure}[htbp]
	\centering
	\input{Figures/Chapter05-VanKampenGrid}
	\caption[同伦方格中的道路替换]{同伦方形被细分，使每个小格的像落在一个开集中。右图把一个小格的下边后右边换为左边后上边；所有边通过连接道路接到同一基点，因而得到群等式 \(br=\ell t\)。}
	\label{g1:fig:van-kampen-grid}
\end{figure}

这个证明也给出泛性质。若 \(f_U:\pi_1(U,x_0)\to K\) 与 \(f_V:\pi_1(V,x_0)\to K\) 满足 \(f_Ui_U=f_Vi_V\)，自由积上的延拓会消去所有关系词，故唯一下降为 \(\pi_1(X,x_0)\to K\)。因此计算的关键是两边的群、交集的群和两条包含同态，而不只是三个群的抽象同构类型。

\subsection{多重交集、基点选择与推广范围}
\label{g1:subsec:0502:03}

有限图常需要同时使用好几个开集。下面给出一个条件稍强、却足以直接使用同一网格证明的多开集版本。它允许任意多个开集，但所有开集共用一个基点。

\begin{corollary}[共同基点的多开集版本]\label{g1:cor:many-set-van-kampen}
	设 \(X=\bigcup_{\alpha\in A}U_\alpha\)，每个 \(U_\alpha\) 都开且包含同一点 \(x_0\)，并且任意非空有限个 \(U_\alpha\) 的交集都道路连通。则包含映射诱导同构
	\[
		\pi_1(X,x_0)\cong
		\left(\mathop{*}_{\alpha\in A}\pi_1(U_\alpha,x_0)\right)\big/N,
	\]
	其中 \(N\) 是所有成对交集闭路在两个因子中所给元素之差
	\(i_{\alpha\beta}(\omega)i_{\beta\alpha}(\omega)^{-1}\) 的正规闭包。
\end{corollary}
\begin{proof}
	对一条道路或一个同伦方形，拉回的开覆盖仍有 Lebesgue 数，故只需要有限段或有限个小矩形。段的交点处，在两个标记的交集中选择连接道路；网格顶点处，在所有相邻小矩形标记的交集中选择连接道路，最多涉及四个开集。改变某段标记时，若它的端点还邻接其他标记，就把端点连接道路选在所有这些标记的交集中，所需交集仍是有限交。假设保证上述选择都能完成。

	其余论证逐项与双开集情形相同：连接道路改变所造成的因子在成对交集关系下相消；共同细分比较不同分段；每条公共网格边在其两侧因子中表示同一元素；逐格等式 \(br=\ell t\) 证明同伦不变。由此构造与自然同态互逆的同态，得到同构。开覆盖可以无限，并不要求 \(X\) 紧；用到紧性的只是道路和同伦的参数空间。
\end{proof}

这里要求所有有限交道路连通，是为了能直接在普通四边网格的顶点选择连接道路。更精细的版本只要求成对与三重交道路连通，并把网格错开，使每个顶点至多接触三个标记；本节的应用使用上述已证明的充分条件。一般没有公共基点、或者交集不道路连通时，需要保留多个基点之间的道路信息，不能直接套用同一个商群公式。

最小的警示就是圆周。取 \(U=S^1\setminus\{-1\}\)、\(V=S^1\setminus\{1\}\)，它们都是可缩开弧，交集却有两个道路分支。若忽略交集的道路连通条件，两个平凡群的自由积会错误地给出平凡群；而\cref{g1:thm:circle-fundamental}已经证明圆周的基本群是 \(\mathbb Z\)。从一条开弧走进另一条开弧后再返回，需要知道经过交集的哪一个分支。只有一个基点的交集基本群没有记录这个信息。

改变基点也不是把记号无条件删去。由\cref{g1:prop:basepoint-change}，道路 \(c:x_0\to x_1\) 给出
\[
	\pi_1(X,x_0)\longrightarrow\pi_1(X,x_1),\qquad
	[\gamma]\longmapsto[\bar c*\gamma*c].
\]
若分别给各个开集选了基点，计算包含同态前必须先指定把它们接到公共基点的道路。更改连接道路通常会把相应元素作共轭；这不改变“加入某条闭路的全部共轭作为关系”的正规闭包，却可能改变具体生成元写出的词。因而群的同构类型可以不依赖某些选择，展示的字母与几何闭路之间的对应仍须写清。

还须注意开性。证明中先把参数区间或方形的开覆盖细分；若把空间任意拆成两个闭子集，拉回覆盖不再提供这种细分保证。下面的图和圆盘附着计算都会明确取开邻域，再通过形变收缩把这些邻域化为易算的对象。

\subsection{图、曲面和去点空间的群展示}
\label{g1:subsec:0502:04}

先看只有边而没有面的空间。这里的有限图是把有限多个闭区间的端点附着到有限个顶点所得的商空间，允许一条边两端相同，也允许两顶点之间有多条边；连通指顶点之间可由有限边路连接。有限图带上述商拓扑，等价地可在每条边赋长度 \(1\) 后用最短边路长度定义度量：边内部的小球是区间，顶点的小球是有限条短边段接成的星形，这些正好给出商拓扑的局部开集。图因而是紧 Hausdorff 空间。

一个不含回路的连通子图称为树。若有限连通图中还有回路，删去该回路的一条边的内部不会破坏顶点间的连通；有限次删除以后得到包含所有顶点的树 \(T\)，称为\term[shengchengshu]{生成树}[spanning tree]。树中任意两顶点间恰有一条不重复顶点的边路：存在性来自连通，若有两条不同的简单边路，分岔后第一次重逢就产生回路。有限树若不止一个顶点，任一最长简单边路的两个端点都只有一条邻边，否则可延长或产生回路。逐次删除这种末端边便知，含 \(v\) 个顶点的树有 \(v-1\) 条边。

树还能形变收缩到任一指定顶点。每次选择一个不同于该顶点的末端顶点，将其唯一邻边按区间坐标线性缩到另一端，其他边保持不动；这些映射在共同端点处相容，故由有限闭集粘合连续。有限次组合便给出保持指定顶点不动的收缩。这是下面去掉树状部分时所需的实际同伦。

\begin{proposition}[有限连通图的基本群]\label{g1:prop:graph-fundamental}
	设 \(G\) 是有 \(v\) 个顶点、\(e\) 条边的有限连通图，\(T\) 是一棵生成树，\(x_0\) 是顶点。则
	\[
		\pi_1(G,x_0)\cong F_{e-v+1}.
	\]
	给每条不在 \(T\) 中的边选定方向，从 \(x_0\) 沿树走到该边起点，走过该边，再沿树返回 \(x_0\)，所得闭路给出一组自由生成元。
\end{proposition}
\begin{proof}
	若没有非树边，\(G=T\) 可缩，结论为平凡群 \(F_0\)。否则对每条非树边 \(d\)，令 \(U_d\) 包含整棵树、整条边 \(d\)，以及其他非树边两端参数长度小于 \(1/3\) 的短边段。每个顶点附近的所有短边段都被包括，所以 \(U_d\) 在图中开；这些集合覆盖 \(G\)，并都包含 \(x_0\)。

	把额外短边段各自线性缩到端点，\(U_d\) 形变收缩到 \(T\cup d\)。树中连接 \(d\) 两端的唯一路径与 \(d\) 组成一个圆周；其余边是附着在这个圆周上的有限树枝，依次收缩末端边就把 \(T\cup d\) 收缩到该圆周。若 \(x_0\) 不在圆周上，收缩时它沿树路移到圆周，再用基点变化公式，得到的生成元正是命题中说明的闭路。因此由\cref{g1:thm:circle-fundamental}及形变收缩不改变基本群，\(\pi_1(U_d,x_0)\cong\mathbb Z\)。

	两个或更多不同 \(U_d\) 的交集只含 \(T\) 及各非树边的端点短边段。缩去这些短边段，再收缩树，交集便可缩，且上述收缩保持 \(x_0\) 不动。因此所有有限交道路连通，成对交集基本群平凡。由\cref{g1:cor:many-set-van-kampen}，\(\pi_1(G,x_0)\) 是每条非树边所贡献的 \(\mathbb Z\) 的自由积。非树边数为 \(e-(v-1)\)，结论成立。
\end{proof}

例如两顶点之间的三条边组成的图有基本群 \(F_2\)。选中间一条边为树，另外两条边各给出一个生成元。虽然在纸面图上它围出两个有界区域，这些区域并不属于图本身，所以没有提供把闭路填满的同伦。往图中加入一张面，才会真正增加关系。

\begin{figure}[htbp]
	\centering
	\input{Figures/Chapter05-GraphAndCell}
	\caption[图的生成元与附着面的关系]{左图中间的蓝边是一棵生成树，其余两边分别贡献自由生成元；纸面上围出的区域不属于图。右图的方形内部则是实际加入的圆盘，按箭头识别对边后，其边界词给出关系 \(aba^{-1}b^{-1}=1\)。}
	\label{g1:fig:graph-and-cell}
\end{figure}

\begin{proposition}[附着一个圆盘]\label{g1:prop:attach-disk}
	设 \(X\) 是道路连通的紧 Hausdorff 空间，\(f:S^1\to X\) 连续，\(x_0=f(1)\)。令
	\[
		Y=X\cup_fD^2=(X\sqcup D^2)/(z\sim f(z)\text{ 对所有 }z\in S^1),
	\]
	即只把圆盘边界按 \(f\) 附着到 \(X\)。记边界闭路为 \(\beta(t)=f(e^{2\pi it})\)。则包含映射给出满同态，且
	\[
		\pi_1(Y,x_0)\cong\pi_1(X,x_0)/\langle\!\langle[\beta]\rangle\!\rangle.
	\]
\end{proposition}
\begin{proof}
	先交代这个商空间及所用收缩的连续性。圆盘内部没有识别，其像是开子空间。对 \(X\) 中的开集 \(O\)，集合
	\[
		O\ \cup\ \{ru:1-\varepsilon<r\le1,\ u\in S^1,\ f(u)\in O\}
	\]
	在商空间中给出 \(O\) 的开邻域；这里边界点已按 \(f\) 识别。若 \(X\) 中两点不同，取不交开邻域 \(O_1,O_2\)，上述邻域也不交。若一个点在圆盘内部，则取半径较大的薄边界环带把它与 \(X\) 分开；两个圆盘内部点用普通小圆盘分开。因此 \(Y\) 为 Hausdorff 空间，且由紧空间 \(X\sqcup D^2\) 的连续像而紧。相同邻域描述也表明 \(X\) 到其像为同胚。

	记圆盘中心的像为 \(o\)，取开集 \(U=Y\setminus\{o\}\) 与 \(V=\{z\in D^2:|z|<2/3\}\) 的像。它们覆盖 \(Y\)，\(V\) 为开圆盘，而交集是 \(0<|z|<2/3\) 的环带。\(U\) 形变收缩到 \(X\)：保持 \(X\) 不动，把每个 \(ru\)（\(r>0\)）送到 \(((1-t)r+t)u\)，终点 \(u\) 按附着关系就是 \(f(u)\)。在边界上公式始终给出 \(f(u)\)，所以与 \(X\) 上的恒等映射相容。为核验同伦一起连续，可注意商映射 \(q:X\sqcup D^2\to Y\) 与 \(\operatorname{id}_{[0,1]}\) 的积是紧空间到 Hausdorff 空间的连续满射，故为闭商映射；限制到开集 \(U\times[0,1]\) 的原像后仍为商映射，上述相容的连续公式遂下降为连续同伦。

	在交集中取基点 \(y_0=1/2\)，其基本群由逆时针圆周 \(t\mapsto\tfrac12e^{2\pi it}\) 生成。径向收缩把这个闭路送到 \(\beta\)；把 \(x_0\) 与 \(y_0\) 沿半径连接并使用基点变化，就得到交集群到 \(\pi_1(U)\) 的像由 \([\beta]\) 生成。另一边 \(V\) 可缩，包含同态把这个生成元送到单位元。由\cref{g1:thm:van-kampen}，商去的正是 \([\beta]\) 的正规闭包。最后沿同一半径把基点换回 \(x_0\)，便得到陈述中的包含同态及其核。
\end{proof}

为何要取正规闭包，而不只删掉整数次的 \(\beta\)？从 \(x_0\) 先沿另一条闭路 \(\eta\) 绕一段，再走圆盘边界，最后沿 \(\bar\eta\) 返回，也能利用同一张圆盘收缩。因此 \([\eta][\beta][\eta]^{-1}\) 必须同时成为单位元。有限次附着圆盘时可逐次应用上面的命题；每一步仍是紧 Hausdorff 空间，最终加入所有附着边界词的关系。

\begin{example}[三个曲面模型的群]\label{g1:exm:surface-presentations}
	环面 \(T^2\) 的基本群为 \(\mathbb Z^2\)，实射影平面 \(\mathbb RP^2\) 的基本群为 \(\mathbb Z/2\mathbb Z\)，从环面去掉一点所得空间的基本群为 \(F_2\)。
\end{example}
\begin{proof}
	在\cref{g1:thm:torus-square}的方形模型中，边界识别后成为有一个顶点、两条环边 \(a,b\) 的图。从左下角出发，逆时针沿方形边界走，依次读出
	\[
		aba^{-1}b^{-1}.
	\]
	方形的内部是一张附着到这个图上的圆盘，故附着命题与图的计算给出
	\[
		\pi_1(T^2)\cong\langle a,b\mid aba^{-1}b^{-1}=1\rangle\cong\mathbb Z^2.
	\]
	这也与\cref{g1:prop:product-fundamental}和 \(T^2=S^1\times S^1\) 的计算相合。边界词记录了具体的附着方向，不是仅凭图形有两个“洞”猜出两个交换的生成元。

	实射影平面可由闭圆盘边界上的对径点识别得到。确切地，将 \(D^2\) 同胚地送到 \(S^2\) 的上半球，再送到对径商 \(S^2/\{\pm1\}\)：每个对径类都有上半球代表，上半球内部的代表唯一，赤道上的两个对径代表恰好被识别。于是得到从圆盘边界对径商到 \(\mathbb RP^2\) 的连续双射；源空间紧，目标由\cref{g1:thm:real-projective-quotient}为 Hausdorff，故是同胚。

	边界自身的商仍是圆周，其坐标可写为 \(e^{i\theta}\mapsto e^{2i\theta}\)。所以绕圆盘边界一周在商圆周上走两周，附着词为 \(a^2\)。得到
	\[
		\pi_1(\mathbb RP^2)\cong\langle a\mid a^2=1\rangle\cong\mathbb Z/2\mathbb Z.
	\]
	最后从环面方形模型的内部去掉中心点，剩下的方形沿从中心发出的射线形变收缩到边界，边界点保持不动，所以同伦与对边识别相容。也可把方形看作圆盘后直接用附着命题证明中的 \(U\) 的径向收缩。所得图是两个圆周在一点相接，故去点环面的基本群为 \(F_2\)。去掉任意其他一点得到相同结论，因为积模型中的旋转 \((z,w)\mapsto(cz,dw)\) 可将任一点移到指定点。
\end{proof}

这些计算针对已经明确构造的曲面，没有使用曲面分类定理。更一般地，若一个多边形模型的边界识别后是一顶点图，先给图中的环边选方向，再沿多边形边界读出词 \(w\)，上述命题便给出 \(\langle a_1,\ldots,a_r\mid w=1\rangle\)。若要断言某个任意给定的曲面必属于某类多边形模型，还需要下一章的分类论证；群展示本身不承担这一任务。

\begin{proposition}[高维球面的单连通性]\label{g1:prop:sphere-simply-connected}
	对每个整数 \(k\ge2\)，球面 \(S^k\) 单连通。
\end{proposition}
\begin{proof}
	当 \(k\ge2\) 时，\(S^k\) 去掉北极或南极得到两个开集，各由球极投影同胚于 \(\mathbb R^k\)：写点为 \((x,t)\in\mathbb R^k\times\mathbb R\)，从北极投影为 \(x/(1-t)\)，逆映射为
	\[
		u\longmapsto\left(\frac{2u}{1+\lVert u\rVert^2},
		\frac{\lVert u\rVert^2-1}{1+\lVert u\rVert^2}\right).
	\]
	南极情形将末坐标变号。这两个开集可缩，交集由
	\[
		(v,t)\longmapsto(\sqrt{1-t^2}\,v,t)
	\]
	同胚于 \(S^{k-1}\times(-1,1)\)，因而道路连通。这里 \(S^{k-1}\) 道路连通可直接看出：两个不对径的单位向量用线段归一化连接，对径时先经一个不与它们共线的单位向量。两个开集道路连通且交集非空，故球面道路连通；再由 van Kampen 定理，两个平凡群的自由积再取商仍平凡，所以 \(\pi_1(S^k)=1\)。
\end{proof}

\begin{proposition}[欧氏空间去掉有限个点]\label{g1:prop:punctured-euclidean}
	设 \(P\subseteq\mathbb R^n\) 是含 \(m\) 个点的有限集合。若 \(n=2\)，则 \(\mathbb R^2\setminus P\) 道路连通，基本群同构于 \(F_m\)，可选从基点出发绕各缺失点一周再返回的闭路作为自由生成元。若 \(n\ge3\)，则 \(\mathbb R^n\setminus P\) 单连通。
\end{proposition}
\begin{proof}
	先说明单点情形。平移把缺失点移到原点，公式
	\[
		H(x,t)=\left(1-t+\frac{t}{\lVert x\rVert}\right)x
	\]
	把 \(\mathbb R^n\setminus\{0\}\) 形变收缩到 \(S^{n-1}\)。当 \(n=2\) 时得到 \(\mathbb Z\)；当 \(n\ge3\) 时，\(n-1\ge2\)，由\cref{g1:prop:sphere-simply-connected}得到单连通性。

	对一般 \(m\) 归纳，\(m=0\) 是可缩的欧氏空间。为使 \(P\) 中各点的第一坐标互不相同，考虑 \(v(t)=(1,t,\ldots,t^{n-1})\)：对任意两不同点之差 \(d\)，内积 \(\langle d,v(t)\rangle\) 是非零多项式，只有有限个实根。取 \(t\) 避开有限个点差所对应的全部实根，再以 \(v(t)\) 的方向为第一坐标轴，就得到所需坐标。取相邻的两个第一坐标之间的数 \(a<b\)，令开条带 \(a<x_1<b\) 不含缺失点，并使 \(x_1<b\) 的开半空间只含最左边的一个缺失点。于是
	\[
		U=(\{x_1<b\}\setminus P),\qquad
		V=(\{x_1>a\}\setminus P)
	\]
	覆盖补空间，交集是凸开条带。每个开半空间同胚于 \(\mathbb R^n\)，例如第一坐标可用 \(\log(b-x_1)\) 或 \(\log(x_1-a)\) 替换，其余坐标不变。因此 \(U\) 是一个点的补空间，\(V\) 对应 \(m-1\) 个点的补空间；它们由归纳假设道路连通，交集可缩。van Kampen 定理给出两基本群的自由积。平面情形为 \(\mathbb Z*F_{m-1}=F_m\)，高维情形为平凡群。

	取交集中的公共基点，并逐次沿条带及半空间中的道路连接到各个小圆周，便得到陈述中的自由生成元；道路的改变会改变所选基及其共轭表达，不改变自由群的同构类型。这个归纳同时证明了补空间道路连通，因为两个道路连通开集有非空交集。
\end{proof}

去点空间说明维数会改变绕行障碍：在平面中，一个缺失点留下一个自由生成元；在至少三维时，闭路可以借额外方向绕过该点并收缩。它也提醒我们不要把“各点绕数都为零”当成平面多点补空间中闭路可收缩的判据，自由群还保留走法的次序。相关的一般方法可参阅\cite[Section 1.2]{Hatcher2002}。
% 纲要条目排为编号小节标题，并进入目录与书签；说明段仍保留于纲要注释中。

% !TeX root = ../../Book1.tex
% 纲要标题：### 5.3 覆盖映射的提升理论
% 纲要位置：几何学计划纲要.md，第 291 行。
% 正文状态：初稿完成，未正式审核。

% 纲要内容开始
% 
% * 覆盖的局部乘积结构
% * 道路提升及唯一性
% * 同伦提升与提升判据
% * 甲板变换及其对纤维的作用
% 
% 纲要内容结束

\section{覆盖映射的提升理论}
\label{g1:sec:0503}

沿圆周连续选择一个实数角度时，可以从任意初始角度出发，并一直跟随道路；但绕圆一周以后，角度增加了 \(2\pi\)，没有回到原值。覆盖映射把这个现象推广到一般空间：底空间上的一个小邻域在上方分成若干互不相交的副本，局部选择一个副本以后，连续移动就决定了下一步该去哪里。真正需要证明的是这种局部选择能否沿整条道路延续、能否随参数连续变化，以及沿不同道路得到的选择何时一致。

\subsection{覆盖的局部乘积结构}
\label{g1:subsec:0503:01}

设想把若干张相同地图叠在一起，每张都投影到同一个区域。这里地图只是局部模型；走到区域外以后，各张地图可以重新连接，所以全空间未必是不交的若干份底空间。

\begin{definition}[覆盖映射]\label{g1:def:covering-map}
	设 \(E,X\) 为拓扑空间。连续满射 \(p:E\to X\) 称为\term[fugaiyingshe]{覆盖映射}[covering map]，如果每个 \(x\in X\) 都有开邻域 \(U\)，使
	\[
		p^{-1}(U)=\bigsqcup_{\lambda\in\Lambda}V_\lambda,
		\qquad p|_{V_\lambda}:V_\lambda\longrightarrow U
	\]
	为同胚，且各 \(V_\lambda\) 都在 \(E\) 中开。这样的 \(U\) 称为\term[junyunfugailinyu]{均匀覆盖邻域}[evenly covered neighborhood]，各 \(V_\lambda\) 称为该邻域上方的\term[fugaipian]{覆盖片}[sheet]。空间 \(X\) 称为底空间，\(E\) 称为覆盖空间；集合 \(p^{-1}(x)\) 称为点 \(x\) 上的纤维，与\secref{g1:sec:0206}中“一个指定点的逆像”的用法相同。
\end{definition}

给指标集 \(\Lambda\) 离散拓扑，映射 \(V_\lambda\ni e\mapsto(p(e),\lambda)\) 将 \(p^{-1}(U)\) 同胚地识别为 \(U\times\Lambda\)，并把 \(p\) 化为第一投影。这就是覆盖的局部乘积结构。反过来，这种与投影相容的局部乘积表示也给出定义中的分解。覆盖片是相对于所选 \(U\) 而言的，不能未经证明就把某一片延伸为整个 \(E\) 的一个分支。若连续映射在每个源点附近都限制为到目标某个开集的同胚，称它为\term[jubutongpei]{局部同胚}[local homeomorphism]；这里邻域选在源空间中，定义本身不要求满射。

\begin{proposition}[覆盖的局部性质]\label{g1:prop:covering-local-properties}
	覆盖映射 \(p:E\to X\) 是局部同胚、开映射和商映射，每个纤维带子空间拓扑时都是离散空间。若 \(X\) 局部道路连通，则 \(E\) 也局部道路连通。
\end{proposition}
\begin{proof}
	每个 \(e\in E\) 都落在某张覆盖片上，而这张片到相应开邻域的限制是同胚，故 \(p\) 为局部同胚。若 \(O\subseteq E\) 开，将 \(O\) 与所有覆盖片相交，每一交集的像都在相应 \(U\) 中开，从而在 \(X\) 中开；这些像的并就是 \(p(O)\)，所以 \(p\) 开。连续开满射是商映射：若 \(p^{-1}(A)\) 开，则 \(A=p(p^{-1}(A))\) 开；另一方向由连续性得到。

	在 \(x\) 上方，每张覆盖片恰与 \(p^{-1}(x)\) 相交于一点，所以纤维的每个单点都相对开。最后，若 \(X\) 局部道路连通，给定 \(e\) 的开邻域 \(O\)，先在一张含 \(e\) 的覆盖片 \(V\) 中看 \(O\cap V\)，再在其像内取含 \(p(e)\) 的道路连通开邻域。该邻域经 \((p|_V)^{-1}\) 的像是包含在 \(O\) 内的道路连通开邻域，证明了 \(E\) 的局部道路连通性。
\end{proof}

若纤维有 \(n\) 个点，常称覆盖在该点为 \(n\) 叶。均匀覆盖邻域内各点的纤维与同一个 \(\Lambda\) 一一对应，所以纤维基数局部不变。特别地，底空间连通时该基数在整个底空间上相同：与一个给定纤维等势的点组成开闭子集。允许无限叶，也允许覆盖空间不连通。

\begin{example}[圆周的两个覆盖]\label{g1:exm:circle-coverings}
	映射
	\[
		q:\mathbb R\longrightarrow S^1,\qquad q(t)=e^{2\pi it}
	\]
	是覆盖。确实，取 \(a\in\mathbb R\) 和 \(0<\varepsilon<1/4\)，令 \(U=q((a-\varepsilon,a+\varepsilon))\)。这是圆周开弧，并且
	\[
		q^{-1}(U)=\bigsqcup_{k\in\mathbb Z}(a-\varepsilon+k,a+\varepsilon+k).
	\]
	每个限制都是同胚：在旋转后的开弧 \(e^{-2\pi ia}U\) 上实部为正，故角度可由连续函数 \(\arctan(\operatorname{Im}z/\operatorname{Re}z)\) 恢复，再加上 \(a+k\)。每个纤维是一份整数集。

	对整数 \(n\ge1\)，映射
	\[
		p_n:S^1\longrightarrow S^1,\qquad p_n(z)=z^n
	\]
	也是覆盖。上面开弧 \(U\) 的原像由 \(n\) 条不交开弧组成，它们分别由
	\[
		z=e^{2\pi i(t+j)/n},\qquad
		a-\varepsilon<t<a+\varepsilon,\quad j=0,\ldots,n-1
	\]
	参数化；\(p_n\) 将每条开弧同胚地映到 \(U\)。所以它是 \(n\) 叶覆盖。
\end{example}

局部同胚满射还不足以保证覆盖。例如 \((0,2)\to S^1\)、\(t\mapsto e^{2\pi it}\) 处处局部同胚且满射，但 \(1\) 上只有一个原像 \(1\)，而其附近的非 \(1\) 点有两个原像，不可能有均匀覆盖邻域。另一方面，\(\mathbb C\to\mathbb C\)、\(z\mapsto z^n\) 在 \(n\ge2\) 时也不是覆盖：原点的每个邻域中都有两个不同的非零点具有相同 \(n\) 次方，故原点附近连局部单射都不成立。覆盖要求的正是这种分支现象不出现。

\subsection{道路提升及唯一性}
\label{g1:subsec:0503:02}

给定连续映射 \(f:Y\to X\)，寻找连续映射 \(\widetilde f:Y\to E\) 满足 \(p\circ\widetilde f=f\)，称为将 \(f\)\term[tisheng]{提升}[lift]到覆盖空间。它所要求的交换关系是
\[
	\begin{tikzcd}[column sep=large,row sep=large]
		& E\arrow[d,"p"]\\
		Y\arrow[ur,dashed,"\widetilde f"]\arrow[r,"f"']&X.
	\end{tikzcd}
\]
提升如果存在，通常还需指定一个点的取值。先证明这个取值会带来多强的唯一性。

\begin{lemma}[连通空间上的提升唯一性]\label{g1:lem:lift-uniqueness}
	设 \(p:E\to X\) 为覆盖，\(Y\) 连通，且连续映射 \(f_1,f_2:Y\to E\) 满足 \(p\circ f_1=p\circ f_2\)。若它们在某一点相等，则处处相等。
\end{lemma}
\begin{proof}
	令 \(A=\{y:f_1(y)=f_2(y)\}\)。对任一点 \(y\)，选 \(p(f_1(y))\) 的均匀覆盖邻域 \(U\)，并令 \(V_1,V_2\) 为分别包含 \(f_1(y),f_2(y)\) 的覆盖片。在开邻域 \(f_1^{-1}(V_1)\cap f_2^{-1}(V_2)\) 内，若 \(V_1=V_2\)，两映射由该片上 \(p\) 的单射性必相等；若 \(V_1\ne V_2\)，它们则因覆盖片不交而处处不等。因此 \(A\) 及其补集都开。由于 \(Y\) 连通且 \(A\ne\varnothing\)，有 \(A=Y\)。
\end{proof}

这里没有假设 \(E\) 为 Hausdorff 空间。相等集的闭性来自覆盖片的分离结构，而不是一般的“连续映射相等集为闭集”结论；后者通常需要目标的 Hausdorff 条件。

\begin{theorem}[道路提升]\label{g1:thm:path-lifting}
	设 \(p:E\to X\) 为覆盖，\(\alpha:[0,1]\to X\) 为道路，且 \(e_0\in E\) 满足 \(p(e_0)=\alpha(0)\)。则存在唯一道路 \(\widetilde\alpha:[0,1]\to E\)，满足
	\[
		p\circ\widetilde\alpha=\alpha,\qquad \widetilde\alpha(0)=e_0.
	\]
\end{theorem}
\begin{proof}
	唯一性由\cref{g1:lem:lift-uniqueness}及区间的连通性得到。构造时，要把道路分成有限段，使每段都落在一个均匀覆盖邻域内，然后由已确定的接点选择下一张片。

	把 \(X\) 的均匀覆盖邻域拉回到 \([0,1]\)，得到该紧度量空间的开覆盖。\cref{g1:lem:lebesgue-number}给出 \(\delta>0\)，使直径小于 \(\delta\) 的子集包含在一个覆盖成员内。取分割 \(0=t_0<t_1<\cdots<t_m=1\)，每段长度小于 \(\delta\)。于是可选均匀覆盖邻域 \(U_j\)，使 \(\alpha([t_{j-1},t_j])\subseteq U_j\)。

	在第一段，取包含 \(e_0\) 的覆盖片 \(V_1\)，定义 \(\widetilde\alpha=(p|_{V_1})^{-1}\circ\alpha\)。若前 \(j-1\) 段已定义，就取 \(U_j\) 上方包含已知端点 \(\widetilde\alpha(t_{j-1})\) 的那一片 \(V_j\)，并在第 \(j\) 段定义 \(\widetilde\alpha=(p|_{V_j})^{-1}\circ\alpha\)。两个相邻公式在接点上相同，故由有限闭集粘合\cref{g1:prop:gluing-lemma}得到连续道路。每段都满足投影等式，第一段又给出指定初值，所以所构造的道路即为所求。
\end{proof}

连续性禁止道路在同一均匀覆盖邻域的不同片之间突然跳跃；有限分割则保证每一次局部选择都能传递到下一段。对反向道路与拼接道路，唯一性立即给出相容性：提升 \(\bar\alpha\) 时若从 \(\widetilde\alpha(1)\) 出发，就得到反向道路 \(\overline{\widetilde\alpha}\)；提升 \(\alpha*\beta\) 时，先提升 \(\alpha\)，再从其终点提升 \(\beta\)，所得拼接就是整条提升。

\begin{figure}[htbp]
	\centering
	\input{Figures/Chapter05-PathLift}
	\caption[圆周环路的非闭提升]{在覆盖 \(q(t)=e^{2\pi it}\) 下，圆周上从 \(1\) 正向绕行一周的环路，从上方 \(0\) 出发提升为线段 \(t\mapsto t\)，终点为 \(1\)。上下同写数字 \(1\) 的点属于不同空间；上方的 \(0,1\) 则属于同一纤维。虚线仅表示这两个端点的投影。}
	\label{g1:fig:path-lift}
\end{figure}

环路的提升不一定闭合，这不是提升定理的缺陷，反而是覆盖保存的整体信息。例如对 \(q\)，从 \(a\) 出发提升 \(t\mapsto e^{2\pi i(a+kt)}\)，得到 \(t\mapsto a+kt\)；整数 \(k\) 直接记录提升终点相对于起点的位移。下面将证明，终点只依赖底空间道路的定端点同伦类。

\subsection{同伦提升与提升判据}
\label{g1:subsec:0503:03}

逐条提升一族道路还不够：若参数稍变而提升突然换片，就不能得到连续同伦。关键在于，一条固定道路只需要有限个局部公式；把参数限制在足够小的邻域内，这些公式可以一起使用。下述版本允许参数空间完全一般，不要求它连通、局部道路连通或紧。

\begin{theorem}[覆盖的同伦提升]\label{g1:thm:homotopy-lifting}
	设 \(p:E\to X\) 为覆盖，\(Y\) 为任意拓扑空间，\(H:Y\times[0,1]\to X\) 连续。若连续映射 \(h_0:Y\to E\) 满足 \(p(h_0(y))=H(y,0)\)，则存在唯一连续映射 \(\widetilde H:Y\times[0,1]\to E\)，满足
	\[
		p\circ\widetilde H=H,\qquad \widetilde H(y,0)=h_0(y).
	\]
\end{theorem}
\begin{proof}
	对每个 \(y\)，由道路提升定理，提升 \(t\mapsto H(y,t)\)，并指定起点 \(h_0(y)\)。逐点定义的这些提升给出唯一可能的函数 \(\widetilde H\)。问题只剩联合连续性，而不只是每条竖直道路的连续性。

	固定 \(y_0\in Y\)。如道路提升证明那样，选有限分割 \(0=t_0<\cdots<t_m=1\)，使每段 \(H(\{y_0\}\times[t_{j-1},t_j])\) 落在均匀覆盖邻域 \(U_j\) 中。对每个 \(j\)，存在 \(y_0\) 的开邻域 \(N_j\)，满足
	\[
		H(N_j\times[t_{j-1},t_j])\subseteq U_j.
	\]
	为说明这一点，对每个 \(t\in[t_{j-1},t_j]\)，用 \(H^{-1}(U_j)\) 的开放性选一个含 \((y_0,t)\) 的乘积开邻域，再由参数闭区间的紧性选有限个覆盖它；取这些邻域在 \(Y\) 方向的有限交即可。再取所有 \(N_j\) 的交，得到可同时处理全部时间段的开邻域 \(N\)。

	令 \(V_1\) 为 \(U_1\) 上方包含 \(h_0(y_0)\) 的片。由 \(h_0\) 连续，缩小 \(N\) 后有 \(h_0(N)\subseteq V_1\)，所以在 \(N\times[t_0,t_1]\) 上可用连续公式 \((p|_{V_1})^{-1}\circ H\)。这个公式在 \(t_1\) 处给出连续的端点函数。令 \(V_2\) 为 \(U_2\) 上方包含该函数在 \(y_0\) 处取值的片，再次缩小 \(N\)，使全部这些端点都在 \(V_2\) 内，然后在下一段用 \((p|_{V_2})^{-1}\circ H\)。照此继续有限次；此前的公式限制到更小的 \(N\) 上仍有效。

	最终在某个含 \(y_0\) 的开邻域 \(N\) 上，得到有限段连续公式，它们在每个接点对全部 \(y\in N\) 都一致。有限闭集粘合引理说明它们合成 \(N\times[0,1]\) 上的连续映射。对每个固定 \(y\)，该映射都是具有指定初值的道路提升，所以由唯一性，它恰等于先前逐点定义的 \(\widetilde H\)。这样的 \(N\) 覆盖 \(Y\)，故 \(\widetilde H\) 联合连续。任意另一个满足条件的提升在每条竖直道路上相等，因而全局相等。
\end{proof}

\begin{corollary}[定端点同伦的提升]\label{g1:cor:lifted-path-homotopy}
	设 \(p:E\to X\) 为覆盖，\(\alpha,\beta\) 是从 \(x_0\) 到 \(x_1\) 的定端点同伦道路。若两者的提升都从同一点 \(e_0\in p^{-1}(x_0)\) 出发，则它们终点相同，并且在 \(E\) 中定端点同伦。
\end{corollary}
\begin{proof}
	设 \(H(s,t)\) 为该同伦，其中 \(s\) 是同伦参数、\(t\) 是道路参数，故 \(H(s,0)=x_0\)、\(H(s,1)=x_1\)。用初始提升 \(s\mapsto e_0\) 提升 \(H\)。每条道路 \(\widetilde H(s,-)\) 从 \(e_0\) 出发，而终点函数 \(s\mapsto\widetilde H(s,1)\) 连续且取值于离散纤维 \(p^{-1}(x_1)\)。区间连通，连续像也连通，离散空间的连通子集至多含一点，因此该终点函数恒定。于是提升同伦的两端都固定，且由道路提升唯一性，其边界道路正是指定的两条提升。
\end{proof}

\begin{proposition}[覆盖诱导的基本群单射]\label{g1:prop:covering-pi1-injective}
	设 \(p:E\to X\) 为覆盖，\(e_0\in E\)，\(x_0=p(e_0)\)。则
	\[
		p_*:\pi_1(E,e_0)\longrightarrow\pi_1(X,x_0)
	\]
	为单射。
\end{proposition}
\begin{proof}
	由\cref{g1:prop:homotopy-functor}，\(p_*\) 是群同态，只需证明其核平凡。设 \(E\) 中基于 \(e_0\) 的环路 \(\gamma\) 满足 \(p\circ\gamma\) 定端点同伦于 \(x_0\) 处的常值道路。将这一同伦从起点 \(e_0\) 提升。由\cref{g1:cor:lifted-path-homotopy}，所得同伦固定两端；第一条边界提升是 \(\gamma\)，另一条边界提升是 \(e_0\) 处的常值道路，因为后者已是常值底道路的一个提升。因此 \([\gamma]=1\)，单射性成立。
\end{proof}

\begin{corollary}[环路提升的闭合判据]\label{g1:cor:closed-lift-criterion}
	设 \(p:(E,e_0)\to(X,x_0)\) 为带指定点的覆盖。基于 \(x_0\) 的环路 \(\alpha\) 从 \(e_0\) 出发的提升闭合，当且仅当
	\[
		[\alpha]\in p_*\pi_1(E,e_0).
	\]
\end{corollary}
\begin{proof}
	若提升 \(\widetilde\alpha\) 闭合，则它代表 \(\pi_1(E,e_0)\) 的元素，且 \(p_*[\widetilde\alpha]=[\alpha]\)。反过来，若 \([\alpha]=p_*[\gamma]\)，其中 \(\gamma\) 是 \(E\) 中基于 \(e_0\) 的环路，则 \(\alpha\) 与 \(p\circ\gamma\) 定端点同伦。从 \(e_0\) 出发的两条提升由\cref{g1:cor:lifted-path-homotopy}有相同终点，而后一提升就是环路 \(\gamma\)，所以前一提升也闭合。
\end{proof}

对于圆周的 \(n\) 叶覆盖 \(p_n\)，基本群在绕数同构 \(\pi_1(S^1,1)\cong\mathbb Z\) 下由 \(k\mapsto nk\) 给出：生成环路 \(t\mapsto e^{2\pi it}\) 与 \(p_n\) 复合后绕行 \(n\) 周。因此，一个圆周环路从 \(1\) 出发的提升闭合，恰好当它的绕数被 \(n\) 整除。这里使用的绕数同构见\cref{g1:thm:circle-fundamental}。

\ProseParagraphBreak

现在考虑整个映射 \(f:Y\to X\) 的提升。可以从 \(y_0\) 沿道路走到每个 \(y\)，提升其像，并用提升终点定义 \(\widetilde f(y)\)。可能出现的障碍是：沿两条不同道路抵达同一点时，得到的终点不同。把一条道路与另一条的反路拼起来，这一障碍便变成环路提升能否闭合的问题。

\begin{theorem}[映射的提升判据]\label{g1:thm:lifting-criterion}
	设 \(p:(E,e_0)\to(X,x_0)\) 为覆盖，\(Y\) 为道路连通且局部道路连通的空间，\(y_0\in Y\)，连续映射 \(f:Y\to X\) 满足 \(f(y_0)=x_0\)。存在连续提升 \(\widetilde f:Y\to E\) 满足 \(\widetilde f(y_0)=e_0\)，当且仅当
	\[
		f_*\pi_1(Y,y_0)\subseteq p_*\pi_1(E,e_0).
	\]
	若提升存在，则它唯一。
\end{theorem}
\begin{proof}
	必要性来自函子性：若 \(f=p\circ\widetilde f\)，则 \(f_*=p_*\circ\widetilde f_*\)，故左边的像包含于右边。唯一性来自\cref{g1:lem:lift-uniqueness}，因为道路连通蕴含连通。

	设子群包含关系成立。对每个 \(y\)，取道路 \(\alpha_y\) 从 \(y_0\) 到 \(y\)，把 \(f\circ\alpha_y\) 从 \(e_0\) 提升，并以其终点定义 \(\widetilde f(y)\)。下面分别检查选择无关性与连续性。

	若 \(\beta_y\) 是另一选择，则 \(\alpha_y*\bar\beta_y\) 为基于 \(y_0\) 的环路。假设保证其像的同伦类属于 \(p_*\pi_1(E,e_0)\)，故由\cref{g1:cor:closed-lift-criterion}，它从 \(e_0\) 出发的提升闭合。该提升先走 \(f\circ\alpha_y\) 的提升，再提升 \(f\circ\bar\beta_y\)，最终回到 \(e_0\)。反转第二段，就得到从 \(e_0\) 出发的 \(f\circ\beta_y\) 的提升，且终点正是第一段的终点。道路提升唯一性于是说明两种选择给出相同的 \(\widetilde f(y)\)。取常值道路还说明 \(\widetilde f(y_0)=e_0\)，而构造直接给出 \(p\circ\widetilde f=f\)。

	固定 \(y\)，取 \(f(y)\) 的均匀覆盖邻域 \(U\)。由 \(Y\) 局部道路连通，可在开集 \(f^{-1}(U)\) 内取含 \(y\) 的道路连通开邻域 \(V\)。令 \(W\) 为 \(U\) 上方含 \(\widetilde f(y)\) 的片。对任意 \(z\in V\)，选完全位于 \(V\) 的道路 \(\delta\) 从 \(y\) 到 \(z\)，并用 \(\alpha_y*\delta\) 作为定义 \(\widetilde f(z)\) 的道路。其后一段的像位于 \(U\)，故从 \(\widetilde f(y)\) 出发的提升为 \((p|_W)^{-1}\circ f\circ\delta\)。因此
	\[
		\widetilde f|_V=(p|_W)^{-1}\circ f|_V,
	\]
	右端连续。这对每个 \(y\) 都成立，所以 \(\widetilde f\) 连续。
\end{proof}

道路连通性用于从一个基点到达所有点，局部道路连通性则用于最后的连续性证明：它让附近各点的比较道路始终留在一个均匀覆盖邻域的原像内。两者在证明中承担不同任务。特别地，若 \(Y\) 道路连通、局部道路连通且单连通，子群条件自动成立，任何这样的 \(f\) 在指定初值以后都能唯一提升。

\subsection{甲板变换及其对纤维的作用}
\label{g1:subsec:0503:04}

在覆盖 \(q:\mathbb R\to S^1\) 中，整数平移改变了实数角度，却不改变底空间的点。一般覆盖也可能有这样的对称性：它在上方交换点，而投影后完全看不出变化。

\begin{definition}[甲板变换]\label{g1:def:deck-transformation}
	设 \(p:E\to X\) 为覆盖。满足 \(p\circ h=p\) 的同胚 \(h:E\to E\) 称为 \(p\) 的\term[jiabanbianhuan]{甲板变换}[deck transformation]，也称覆盖变换。所有甲板变换在函数通常复合
	\[
		(hk)(e)=h(k(e))
	\]
	下组成群，记为 \(\operatorname{Deck}(p)\)。
\end{definition}

确实，恒等映射满足投影等式，两个这样的同胚复合仍满足等式；从 \(p\circ h=p\) 右复合 \(h^{-1}\)，又得 \(p\circ h^{-1}=p\)，所以逆映射也在其中。每个甲板变换保持每个纤维，因而在纤维上给出置换。

\begin{proposition}[甲板变换由一点确定]\label{g1:prop:deck-determined}
	设 \(p:E\to X\) 为覆盖，且 \(E\) 非空连通。若两个甲板变换在某一点相等，则它们相同。特别地，非恒等甲板变换没有不动点；对任意 \(e_0\in E\)，映射
	\[
		\operatorname{Deck}(p)\longrightarrow p^{-1}(p(e_0)),
		\qquad h\longmapsto h(e_0)
	\]
	为单射。
\end{proposition}
\begin{proof}
	任意两个甲板变换都是映射 \(p:E\to X\) 的提升，因此可直接使用\cref{g1:lem:lift-uniqueness}。取其中一个为恒等映射，就得到不动点结论；再固定 \(e_0\)，就得到所写映射的单射性。
\end{proof}

群在集合上的作用称为自由的，是指只有单位元能固定一个点。因此在连通覆盖空间上，甲板变换群的作用自由；但自由不等于传递，不能据此断言同一纤维中的任意两点都能由某个甲板变换互换。后一性质将在下一节与正规子群联系起来。若底空间局部道路连通，则\cref{g1:prop:covering-local-properties}与\cref{g1:thm:locally-path-components}还说明：覆盖空间的连通分支就是道路分支；这里的唯一性命题本身则只需连通。

\begin{example}[圆周覆盖的甲板变换群]\label{g1:exm:circle-deck-groups}
	对于 \(q:\mathbb R\to S^1\)，整数平移 \(T_k(t)=t+k\) 都是甲板变换。反过来，若 \(h\) 是甲板变换，则 \(q(h(t))=q(t)\) 意味着 \(h(t)-t\in\mathbb Z\)。函数 \(t\mapsto h(t)-t\) 连续，\(\mathbb R\) 连通而 \(\mathbb Z\) 离散，故它恒为某个整数 \(k\)。所以全部甲板变换恰为 \(T_k\)，且
	\[
		T_k\circ T_\ell=T_{k+\ell},\qquad \operatorname{Deck}(q)\cong\mathbb Z.
	\]

	对于 \(p_n:S^1\to S^1\)，令 \(\zeta=e^{2\pi i/n}\)。旋转 \(R_j(z)=\zeta^jz\) 是甲板变换。反过来，投影等式给出 \((h(z)/z)^n=1\)，所以连续函数 \(h(z)/z\) 取值于有限离散集 \(\{1,\zeta,\ldots,\zeta^{n-1}\}\)。圆周道路连通，故该函数恒定；于是 \(h=R_j\)，且
	\[
		\operatorname{Deck}(p_n)\cong\mathbb Z/n\mathbb Z.
	\]
	这两个例子的甲板变换都在每个纤维上传递，但这一事实来自上述具体计算。
\end{example}

还有一种纤维置换来自在底空间走环路并取提升终点。它与甲板变换的来源不同，却彼此相容：若 \(\widetilde\alpha_e\) 是从 \(e\) 出发的道路提升，那么 \(h\circ\widetilde\alpha_e\) 是从 \(h(e)\) 出发的同一道路的提升，因而
\[
	h\bigl(\widetilde\alpha_e(1)\bigr)
	=\widetilde\alpha_{h(e)}(1).
\]
这一等式只用了提升唯一性。到\secref{g1:sec:0505}中，我们会将取提升终点的操作组织成基本群在纤维上的右作用，并用它重建覆盖。

% !TeX root = ../../Book1.tex
% 纲要标题：### 5.4 万有覆盖与分类定理
% 纲要位置：几何学计划纲要.md，第 298 行。
% 正文状态：初稿；尚未正式审核。

% 纲要内容开始
% 
% * 道路类构造和拓扑的定义
% * 局部假设保证覆盖性质
% * 带基点连通覆盖与子群的对应
% * 不带基点的共轭类及正规覆盖
% * 局部假设失败时分类论证在哪里中断
% 
% 纲要内容结束

\section{万有覆盖与分类定理}
\label{g1:sec:0504}

沿圆周行走时，只记录终点会丢掉已经绕过的圈数；把道路提到实线上，这段历史便保留在终点的整数偏移中。一般空间也可以尝试保留从某个固定点出发的全部道路历史。不过，连续变形不应产生新的历史，所以真正要记录的是定端点同伦类。这样得到的集合是否仍能在每一点附近看起来像原空间，正是本节首先要解决的问题。解决以后，覆盖的不同之处就可以解释为：允许把哪些环路的历史遗忘掉。

\subsection{道路类构造和拓扑的定义}
\label{g1:subsec:0504:01}

局部收缩条件有两种不同的范围。邻域中的环路能够在整个空间中缩掉，并不表示它能够在这个邻域内部缩掉。

\begin{definition}[半局部单连通]\label{g1:def:semilocally-simply-connected}
	设 \(X\) 为拓扑空间。如果每个 \(x\in X\) 都有开邻域 \(U\)，使每条以 \(x\) 为基点、完全位于 \(U\) 的环路在 \(X\) 中定基点零伦，即包含映射诱导的同态 \(\pi_1(U,x)\to\pi_1(X,x)\) 为零，就称该空间为\term[banjubudanliantong]{半局部单连通}[semilocally simply connected]。若每一点都有由单连通开邻域组成的邻域基，则称 \(X\)\term[jubudanliantong]{局部单连通}[locally simply connected]。
\end{definition}

局部单连通蕴含局部道路连通与半局部单连通；本节的分类定理只要求后两项。若 \(X\) 局部道路连通且半局部单连通，则每个开邻域 \(O\ni x\) 都包含一个道路连通开邻域 \(U\ni x\)，使 \(U\) 中的全部环路在 \(X\) 中零伦。事实上，先取半局部条件中的邻域 \(W\)，再在 \(W\cap O\) 内取道路连通开邻域 \(U\)。以 \(x\) 为基点的结论成立；以其他 \(y\in U\) 为基点的环路，可以沿 \(U\) 内一条从 \(x\) 到 \(y\) 的道路搬回 \(x\)，由\cref{g1:prop:basepoint-change}得到同样结论。下文把这样的道路连通开集 \(U\) 简称为好邻域；它不必本身单连通。

固定 \(x_0\in X\)。令 \(\widetilde X\) 为所有从 \(x_0\) 出发的道路的定端点同伦类所成的集合，定义
\[
	\widetilde x_0=[e_{x_0}],\qquad p:\widetilde X\longrightarrow X,\qquad p([\alpha])=\alpha(1).
\]
这里 \(e_{x_0}\) 是常值道路，符号 \([\alpha]\) 包含终点信息：不同终点的道路不可能代表同一个点。如何描述 \([\alpha]\) 附近的点？合理的选择是保留已有道路，只允许在它的末端追加一小段道路。若 \(U\) 是包含 \(\alpha(1)\) 的好邻域，令
\[
	B([\alpha],U)=\{[\alpha*\beta]:\beta(0)=\alpha(1),\ \beta(I)\subseteq U\}.
\]
沿用全章约定，\(\alpha*\beta\) 表示先走 \(\alpha\)，再走 \(\beta\)。改变 \(\alpha\) 的同伦代表不改变上述集合，因为道路拼接与定端点同伦相容，见\cref{g1:prop:path-classes}。

\begin{lemma}[道路类的基本邻域]\label{g1:lem:path-class-basis}
	设 \(X\) 道路连通、局部道路连通且半局部单连通，固定 \(x_0\in X\)。令 \(\widetilde X\) 为从 \(x_0\) 出发的道路的定端点同伦类集合，\(p([\alpha])=\alpha(1)\)。对含 \(\alpha(1)\) 的道路连通开集 \(U\)，要求其中所有基于 \(\alpha(1)\) 的环路在 \(X\) 中零伦，并记 \(B([\alpha],U)=\{[\alpha*\beta]:\beta:[0,1]\to U,\ \beta(0)=\alpha(1)\}\)。所有 \(B([\alpha],U)\) 构成 \(\widetilde X\) 的一个拓扑基；在此拓扑下，限制映射 \(p:B([\alpha],U)\to U\) 是同胚。
\end{lemma}
\begin{proof}
	每个 \([\alpha]\) 都属于这样的集合，因为可以追加常值道路。若 \([\delta]\in B([\alpha],U)\cap B([\gamma],V)\)，取包含 \(\delta(1)\) 的好邻域 \(W\subseteq U\cap V\)。由 \([\delta]=[\alpha*\beta]\)、\(\beta(I)\subseteq U\)，在 \(\delta\) 后追加一条位于 \(W\) 的道路，仍然得到 \(B([\alpha],U)\) 中的点；对 \(V\) 同理。所以
	\[
		B([\delta],W)\subseteq B([\alpha],U)\cap B([\gamma],V),
	\]
	这就验证了拓扑基的交集条件。

	限制映射满射，因为 \(U\) 道路连通。若 \(\beta,\beta'\) 从 \(\alpha(1)\) 出发并在 \(U\) 内到达同一个点，则 \(\beta*\bar\beta'\) 是 \(U\) 中的环路，在 \(X\) 中零伦。因此 \(\alpha*\beta\) 与 \(\alpha*\beta'\) 定端点同伦，限制映射也单射。这是好邻域条件的第一次关键使用：追加道路的终点已经决定了追加后的历史。

	对 \(X\) 中开集 \(O\) 及任意 \([\delta]\in p^{-1}(O)\)，在 \(O\) 内取 \(\delta(1)\) 的好邻域 \(W\)，则 \(B([\delta],W)\subseteq p^{-1}(O)\)。故 \(p\) 连续。另一方面，\(p\) 将每个基元素 \(B([\delta],W)\) 映为开集 \(W\)，因而是开映射。限制到一个开基元素后，仍是开映射；所以刚才得到的连续双射 \(B([\alpha],U)\to U\) 是同胚。
\end{proof}

这个拓扑记录道路的同伦历史：一个基本邻域保留已有道路，允许在末端追加位于 \(U\) 内的一小段道路。

\subsection{局部假设保证覆盖性质}
\label{g1:subsec:0504:02}

终点映射在每个基本邻域上已经是同胚。要使它成为单连通覆盖，还要证明同一底邻域上方互不相同的基本邻域彼此不交，并把总空间的环路转回道路类来计算。

\begin{definition}[万有覆盖]\label{g1:def:universal-cover}
	覆盖映射 \(p:\widetilde X\to X\) 的总空间若单连通，即非空、道路连通且基本群平凡，则称它为 \(X\) 的\term[wanyoufugai]{万有覆盖}[universal covering]。
\end{definition}

\begin{theorem}[道路类构造万有覆盖]\label{g1:thm:path-class-universal-cover}
	设 \(X\) 道路连通、局部道路连通且半局部单连通，\(x_0\in X\)。令 \(\widetilde X\) 为从 \(x_0\) 出发的道路的定端点同伦类集合。称道路连通开集 \(U\) 为好邻域，是指其中的环路在 \(X\) 中均为零伦；以 \(B([\alpha],U)=\{[\alpha*\beta]:\beta:[0,1]\to U,\ \beta(0)=\alpha(1)\}\) 为基赋予 \(\widetilde X\) 拓扑，其中 \(U\) 遍历含 \(\alpha(1)\) 的好邻域。记 \(\widetilde x_0=[e_{x_0}]\)、\(p([\alpha])=\alpha(1)\)。则 \(p:(\widetilde X,\widetilde x_0)\to(X,x_0)\) 是万有覆盖，每个好邻域都被 \(p\) 均匀覆盖。
\end{theorem}
\begin{proof}
	先固定好邻域 \(U\) 及 \(u\in U\)。在 \(p^{-1}(u)\) 中的每个道路类 \([\alpha]\) 上取基本邻域 \(B([\alpha],U)\)。这些集合覆盖 \(p^{-1}(U)\)：若 \([\eta]\) 的终点是 \(z\in U\)，取 \(U\) 内从 \(u\) 到 \(z\) 的道路 \(\beta\)，则 \([\eta*\bar\beta]\in p^{-1}(u)\)，且 \([\eta]\in B([\eta*\bar\beta],U)\)。

	不同的 \([\alpha],[\alpha']\in p^{-1}(u)\) 给出互不相交的基本邻域。否则可写 \([\alpha*\beta]=[\alpha'*\beta']\)，其中 \(\beta,\beta'\) 在 \(U\) 内由 \(u\) 到达同一点。环路 \(\beta*\bar\beta'\) 在 \(X\) 中零伦，故消去末端道路后得到 \([\alpha]=[\alpha']\)。结合\cref{g1:lem:path-class-basis}中的局部同胚，得到
	\[
		p^{-1}(U)=\coprod_{[\alpha]\in p^{-1}(u)}B([\alpha],U).
	\]
	道路连通性还保证 \(p\) 满射，所以这确实是覆盖映射。

	下面证明总空间道路连通。对任意道路 \(\alpha\) 定义初段 \(\alpha_t(s)=\alpha(ts)\)，则
	\[
		\widetilde\alpha(t)=[\alpha_t],\qquad 0\le t\le1
	\]
	从 \(\widetilde x_0\) 走到 \([\alpha]\)。它连续：在任意 \(t_0\) 附近，取好邻域 \(U\ni\alpha(t_0)\)，再把参数区间缩小，使相应的道路小段都在 \(U\) 内。无论 \(t\ge t_0\) 还是 \(t\le t_0\)，\([\alpha_t]\) 都可由 \([\alpha_{t_0}]\) 追加这段道路或其反路得到。因此在此区间内 \(\widetilde\alpha\) 等于局部逆 \((p|_{B([\alpha_{t_0}],U)})^{-1}\circ\alpha\)，从而连续。

	最后设 \(\lambda\) 是 \(\widetilde X\) 中以 \(\widetilde x_0\) 为基点的环路，令 \(\alpha=p\circ\lambda\)。上面的初段构造给出 \(\alpha\) 从 \(\widetilde x_0\) 出发的提升；由\cref{g1:thm:path-lifting}的唯一性，它就是 \(\lambda\)。因此 \([\alpha]=\widetilde\alpha(1)=\widetilde x_0\)，即 \(\alpha\) 在 \(X\) 中零伦。由\cref{g1:prop:covering-pi1-injective}，覆盖诱导的 \(p_*\) 是单射，所以 \([\lambda]=1\)。这证明了 \(\widetilde X\) 单连通。
\end{proof}

这个证明显示了两个局部条件的不同职责。局部道路连通使好邻域足够多，能够恢复 \(X\) 的原拓扑；半局部单连通则使每片基本邻域不会因绕小圈而在同一个终点上产生多个历史。

\begin{corollary}[万有覆盖向任意连通覆盖的映射]\label{g1:cor:universal-cover-map}
	设 \(X\) 为局部道路连通拓扑空间，\(x_0\in X\)，\(p:(\widetilde X,\widetilde x_0)\to(X,x_0)\) 是万有覆盖（其道路类构造见\cref{g1:thm:path-class-universal-cover}），\(q:(E,e_0)\to(X,x_0)\) 是连通覆盖。存在唯一连续映射 \(F:\widetilde X\to E\)，使 \(qF=p\)、\(F(\widetilde x_0)=e_0\)；这个 \(F\) 本身也是满射覆盖映射。
	\[
		\begin{tikzcd}
			(\widetilde X,\widetilde x_0)\arrow[rr,dashed,"F"]\arrow[dr,"p"']&&(E,e_0)\arrow[dl,"q"]\\
			&(X,x_0)&
		\end{tikzcd}
	\]
\end{corollary}
\begin{proof}
	覆盖的局部同胚性质使 \(E\) 与 \(\widetilde X\) 局部道路连通；因此连通的 \(E\) 道路连通，见\cref{g1:thm:locally-path-components}。因为 \(\pi_1(\widetilde X)=1\)，\cref{g1:thm:lifting-criterion}的子群包含条件自动满足，给出所需唯一映射。任取 \(e\in E\)，沿 \(E\) 中从 \(e_0\) 到 \(e\) 的道路投影，再提升到 \(\widetilde X\)。由提升唯一性，所得末端映到 \(e\)，所以 \(F\) 满射。

	在 \(X\) 中取一个道路连通开邻域 \(U\)，使它同时被 \(p,q\) 均匀覆盖。\(p\) 的每一片道路连通，故其在 \(F\) 下的像落在 \(q\) 的一片中，且在这两片上 \(F=(q|_{\text{该片}})^{-1}\circ p\)，是同胚。对 \(q\) 的任一片，其逆像便是若干这样的 \(p\) 的片的不交并。这说明 \(F\) 是覆盖映射。
\end{proof}

\subsection{带基点连通覆盖与子群的对应}
\label{g1:subsec:0504:03}

令 \(G=\pi_1(X,x_0)\)。前节已经说明，一个带基点覆盖 \(q:(E,e_0)\to(X,x_0)\) 决定子群 \(q_*\pi_1(E,e_0)\le G\)：它恰好收集那些从 \(e_0\) 出发后仍回到 \(e_0\) 的环路类。现在反过来，从指定子群 \(H\le G\) 出发构造覆盖。两条道路 \(\alpha,\beta\) 应被认为具有相同历史，当且仅当它们终点相同，而且相差的环路 \(\alpha*\bar\beta\) 属于 \(H\)。也就是说，\(H\) 指定了允许遗忘的环路。

\begin{theorem}[子群对应的覆盖]\label{g1:thm:subgroup-covering}
	设 \(X\) 道路连通、局部道路连通且半局部单连通，\(x_0\in X\)，\(H\le G=\pi_1(X,x_0)\)。在从 \(x_0\) 出发的道路集合上规定
	\[
		\alpha\sim_H\beta\quad\Longleftrightarrow\quad
		\alpha(1)=\beta(1)\ \text{且}\ [\alpha*\bar\beta]\in H.
	\]
	以 \([\alpha]_H\) 记等价类，令 \(E_H\) 为这些等价类的集合，\(e_H=[e_{x_0}]_H\)。以
	\[
		B_H([\alpha]_H,U)=\{[\alpha*\beta]_H:\beta:[0,1]\to U,\ \beta(0)=\alpha(1)\}
	\]
	为基赋予拓扑，其中 \(U\) 遍历含 \(\alpha(1)\) 的道路连通开邻域，且 \(U\) 中所有基于 \(\alpha(1)\) 的环路均在 \(X\) 中零伦。则终点映射 \(p_H:(E_H,e_H)\to(X,x_0)\) 是连通覆盖，且
	\[
		(p_H)_*\pi_1(E_H,e_H)=H,\qquad p_H^{-1}(x_0)\cong H\backslash G=\{Hg:g\in G\}.
	\]
	尤其覆盖的叶数为子群指数 \([G:H]\)，允许此数无限。
\end{theorem}
\begin{proof}
	关系自反、对称分别来自单位元与取逆。若 \(\alpha\sim_H\beta\)、\(\beta\sim_H\gamma\)，则
	\[
		[\alpha*\bar\gamma]=[\alpha*\bar\beta]\,[\beta*\bar\gamma]\in H,
	\]
	其中消去 \(\bar\beta*\beta\) 用的是道路类的拼接与逆关系。因此这是等价关系；特别地，同伦的道路必等价。若 \([\alpha]_H=[\alpha']_H\)，追加同一条 \(\beta\) 不改变等价性，因为 \([\alpha*\beta*\bar\beta*\bar\alpha']=[\alpha*\bar\alpha']\)。于是基元素的定义不依赖代表。

	基的交集条件与\cref{g1:lem:path-class-basis}相同：若 \([\delta]_H\) 在两个基元素中，在两好邻域交内取更小好邻域 \(W\)，则 \(B_H([\delta]_H,W)\) 同时包含在二者之内。每个 \(B_H([\alpha]_H,U)\to U\) 满射。若两个追加道路终点相同，它们在 \(U\) 中相差的环路在 \(X\) 中为单位元，因而属于 \(H\)，故此限制映射单射。它的连续性与开性也由同一基论证得到，所以是同胚。

	固定 \(u\in U\)，这些基元素在 \([\alpha]_H\) 遍历 \(p_H^{-1}(u)\) 时覆盖 \(p_H^{-1}(U)\)。若其中两片相交，则有 \([\alpha*\beta]_H=[\alpha'*\beta']_H\)，其中 \(\beta,\beta'\) 都在 \(U\) 内由 \(u\) 到达同一点。由于 \(\beta*\bar\beta'\) 在 \(X\) 中零伦，等价条件化为 \([\alpha*\bar\alpha']\in H\)，即两个索引本来就相同。因此这些片互不相交，\(p_H\) 是覆盖。初段道路 \(t\mapsto[\alpha_t]_H\) 的连续性由局部逆证明，故 \(E_H\) 道路连通。

	以 \(x_0\) 为基点的环路 \(\gamma\) 从 \(e_H\) 出发提升后终点为 \([\gamma]_H\)，它等于 \(e_H\) 当且仅当 \([\gamma]\in H\)。\cref{g1:cor:closed-lift-criterion}告诉我们，闭提升的环路类恰为 \((p_H)_*\pi_1(E_H,e_H)\)，所以该子群正是 \(H\)。纤维内的道路全是环路；若其类为 \(g,k\in G\)，等价条件为 \(gk^{-1}\in H\)，等价于 \(Hg=Hk\)。这给出所述纤维对应。其他纤维与此纤维之间可沿一条道路提升得到双射，故叶数恒为 \([G:H]\)。
\end{proof}

\begin{theorem}[带基点的连通覆盖分类]\label{g1:thm:based-covering-classification}
	设 \(X\) 道路连通、局部道路连通且半局部单连通，\(x_0\in X\)。带基点连通覆盖 \(q:(E,e_0)\to(X,x_0)\) 的等价类与 \(G=\pi_1(X,x_0)\) 的子群一一对应；这里等价指保持基点且与投影交换的同胚。对应把 \((E,e_0)\) 送到 \(q_*\pi_1(E,e_0)\)，逆对应把子群 \(H\le G\) 送到满足 \(q_*\pi_1(E,e_0)=H\) 的唯一带基点连通覆盖等价类。
\end{theorem}
\begin{proof}
	给定覆盖，令 \(H=q_*\pi_1(E,e_0)\)。总空间 \(E\) 局部道路连通，且连通，因而道路连通。定义 \(F:E_H\to E\)：把 \([\alpha]_H\) 送到 \(\alpha\) 从 \(e_0\) 出发的提升终点。需要证明，此终点恰好只忘掉 \(H\) 中的历史。

	若 \(\alpha,\beta\) 有相同终点，则其提升终点相同，当且仅当 \(\alpha*\bar\beta\) 有从 \(e_0\) 出发的闭提升。一方向可以直接把两条提升拼接；反方向把闭提升分为前后两段，反转后半段，再用 \(\beta\) 的提升唯一性。由闭提升判据，这又等价于 \([\alpha*\bar\beta]\in H\)。所以 \(F\) 良定义且单射。任意 \(e\in E\) 可沿 \(E\) 中一条从 \(e_0\) 到 \(e\) 的道路到达，取其投影就得 \(F\) 的一个原像，因此 \(F\) 满射。

	对 \([\alpha]_H\)，取包含 \(\alpha(1)\) 的好邻域 \(U\)，并使它被 \(q\) 均匀覆盖。在 \(B_H([\alpha]_H,U)\) 上，\(F\) 就是 \(p_H\) 后接经过 \(F([\alpha]_H)\) 的那片的局部逆。因此 \(F\) 局部为同胚，进而是同胚，且 \(qF=p_H\)、\(F(e_H)=e_0\)。

	如果两个带基点覆盖等价，等价同胚诱导基本群同构，故它们在 \(G\) 中的像子群相同；反过来，若像子群相同，上述构造说明它们都与同一个 \(E_H\) 等价。最后，\cref{g1:thm:subgroup-covering}保证每个子群确实出现，完成一一对应的证明。
\end{proof}

\(H=\{1\}\) 时，上述构造就是万有覆盖；\(H=G\) 时，所有同终点道路都等价，得到恒等覆盖 \(X\to X\)。子群越大，允许遗忘的历史越多，覆盖中的点就越少。万有覆盖的“万有”对应\cref{g1:cor:universal-cover-map}所述的映射性质。

\subsection{不带基点的共轭类及正规覆盖}
\label{g1:subsec:0504:04}

同一个覆盖若换了纤维中的起点，能够闭合的环路集合也可能改变。为看清改变方式，暂将从 \(e\in q^{-1}(x_0)\) 提升代表 \(g\in G\) 的环路所得终点记为 \(e\cdot g\)。同伦提升保证它与代表无关，先走后走的约定给出
\[
	(e\cdot g)\cdot k=e\cdot(gk),\qquad e\cdot1=e.
\]
这就是纤维上的右作用。连通覆盖的总空间道路连通，所以此作用传递：任意两个纤维点可由道路相连，其投影是所需环路。对 \(E_H\)，纤维对应 \(H\backslash G\)，且右作用为 \((Hg)\cdot k=H(gk)\)。于是
\[
	\{k:(Hg)\cdot k=Hg\}=g^{-1}Hg.
\]
左边是以 \(Hg\) 为起点时有闭提升的环路类，因而也是该基点处基本群在 \(G\) 中的像。

\begin{proposition}[不带基点的分类]\label{g1:prop:unbased-covering-classification}
	设 \(X\) 道路连通、局部道路连通且半局部单连通，\(x_0\in X\)。不指定总空间基点的连通覆盖等价类，与 \(\pi_1(X,x_0)\) 的子群共轭类一一对应。
\end{proposition}
\begin{proof}
	在同一个连通覆盖内换基点，刚才的计算说明像子群由 \(H\) 变成 \(g^{-1}Hg\)，且所有这样的共轭子群都可得到。若两个不带基点的覆盖等价，任选第一个覆盖的一个纤维点，并在第二个中选它的像，便得到相同的像子群；与第二个覆盖原来任选的基点比较，只相差共轭。反过来，若两个覆盖的像子群共轭，可以在其中一个覆盖内更换纤维基点，使两个像子群相同，再应用\cref{g1:thm:based-covering-classification}得到等价。
\end{proof}

\begin{definition}[正规覆盖]\label{g1:def:normal-covering}
	连通覆盖 \(q:E\to X\) 若其覆盖变换群在每个纤维上传递，称为\term[zhengguifugai]{正规覆盖}[regular covering]；英文文献也使用 normal covering。\foreignidx{normal covering}在道路连通底空间上，只需在一个纤维上传递：沿连接两个底点的道路提升，可把覆盖变换与两纤维间的对应交换。
\end{definition}

正规性允许在覆盖内部把任一张“同样的片”搬到另一张。这个要求比“各纤维点数相同”强得多，后者对这里所有覆盖都成立。

\begin{theorem}[正规化子与覆盖变换]\label{g1:thm:normal-covering-deck}
	设 \(X\) 道路连通、局部道路连通且半局部单连通，\(x_0\in X\)，\(G=\pi_1(X,x_0)\)，\(H\le G\)。令 \(p_H:(E_H,e_H)\to(X,x_0)\) 为满足 \((p_H)_*\pi_1(E_H,e_H)=H\) 的带基点连通覆盖。记
	\[
		N_G(H)=\{n\in G:nHn^{-1}=H\}.
	\]
	这个子群称为 \(H\) 在 \(G\) 中的\term[zhengguihuazi]{正规化子}[normalizer]。则有群同构
	\[
		\operatorname{Deck}(E_H/X)\cong N_G(H)/H.
	\]
	覆盖 \(E_H\to X\) 正规当且仅当 \(H\) 是 \(G\) 的正规子群；此时覆盖变换群为 \(G/H\)。特别地，万有覆盖的覆盖变换群同构于 \(G\)，在道路类模型上由
	\[
		L_g[\alpha]=[\gamma*\alpha],\qquad [\gamma]=g
	\]
	给出左作用，并且 \(L_g\circ L_k=L_{gk}\)。
\end{theorem}
\begin{proof}
	先构造变换，再说明没有遗漏。若 \(n\in N_G(H)\)，选代表环路 \(\nu\)，令 \(L_n[\alpha]_H=[\nu*\alpha]_H\)。若 \([\alpha]_H=[\beta]_H\)，则
	\[
		[\nu*\alpha*\bar\beta*\bar\nu]
		=n[\alpha*\bar\beta]n^{-1}\in H,
	\]
	所以公式良定义，也不依赖 \(\nu\) 的同伦代表。它把 \(B_H([\alpha]_H,U)\) 同胚地送到 \(B_H([\nu*\alpha]_H,U)\)，逆为 \(L_{n^{-1}}\)，并保持投影，故是覆盖变换。拼接顺序直接给出 \(L_m\circ L_n=L_{mn}\)，这是到覆盖变换群的同态。\(L_n\) 为恒等当且仅当 \(n\in H\)：一方向代入基点 \(e_H\)；另一方向用 \([\nu*\alpha*\bar\alpha]=n\in H\)。因此核为 \(H\)。

	任一覆盖变换 \(T\) 都与纤维的右作用交换，因为 \(T\) 将一条道路的提升变成从像点出发的同一道路的提升，而提升唯一。写 \(T(e_H)=Hn\)。由 \(T\) 可逆，一个 \(g\) 固定 \(e_H\) 当且仅当它固定 \(T(e_H)\)。前者的稳定子是 \(H\)，后者是 \(n^{-1}Hn\)，所以 \(n^{-1}Hn=H\)，即 \(n\in N_G(H)\)。\(T\) 与 \(L_n\) 在 \(e_H\) 相等，由\cref{g1:prop:deck-determined}，它们在整个连通覆盖上相等。故上述同态满射，商去核得到群同构。

	基点 \(e_H\) 的覆盖变换轨道恰为 \(\{Hn:n\in N_G(H)\}\)。它等于全部纤维 \(H\backslash G\)，当且仅当 \(N_G(H)=G\)：若 \(Hg=Hn\)，则 \(g\in Hn\subseteq N_G(H)\)。而 \(N_G(H)=G\) 正是 \(H\trianglelefteq G\)。正规性与正规子群条件的等价因而成立。最后令 \(H=\{1\}\)，就得到万有覆盖的结论。
\end{proof}

这里的左右两种作用承担不同任务。覆盖变换群在道路前面追加环路，所以是左作用；纤维的路径提升在道路后面追加环路，所以是右作用。在 \(H\backslash G\) 上，前者写为 \(Hg\mapsto Hng\)（\(n\in N_G(H)\)），后者写为 \(Hg\mapsto Hgk\)（\(k\in G\)）。左右乘相互交换，因而这两类纤维置换彼此相容。

也可以把 \(E_H\) 写成 \(H\backslash\widetilde X\)。事实上，\([\alpha]_H=[\beta]_H\) 等价于存在 \(h\in H\) 使 \([\alpha]=L_h[\beta]\)：取 \(h=[\alpha*\bar\beta]\) 即可。因此自然满射 \(\widetilde X\to E_H\) 的纤维正是 \(H\) 的左轨道。它把基本邻域 \(B([\alpha],U)\) 映为 \(B_H([\alpha]_H,U)\)，是连续开满射，故 \(E_H\) 的拓扑正是此轨道商拓扑。这些基本邻域就是\cref{g1:thm:subgroup-covering}中构造出的覆盖片。

\subsection{局部假设失败时分类论证在哪里中断}
\label{g1:subsec:0504:05}

下面先看半局部单连通失败的空间：它仍然道路连通、局部道路连通，但根本没有单连通覆盖。

\begin{example}[缩小圆周汇聚在同一点]\label{g1:exm:hawaiian-earring}
	在平面中令
	\[
		C_n=\{(x,y):(x-1/n)^2+y^2=1/n^2\},\qquad
		Z=\bigcup_{n\ge1}C_n,\qquad o=(0,0).
	\]
	赋予 \(Z\) 平面的子空间拓扑。这个空间称为\term[xiaweiyierhuan]{Hawaii 耳环}[Hawaiian earring]。所有圆周相交于 \(o\)，所以 \(Z\) 道路连通。非 \(o\) 点附近有普通开弧邻域；\(o\) 附近的 \(Z\cap B(o,\varepsilon)\) 在每个圆周上截得包含 \(o\) 的一段弧或整个圆周，这些集合都能沿弧连到 \(o\)，所以也道路连通。故 \(Z\) 局部道路连通。
	\begin{figure}[htbp]
		\centering
		\input{Figures/Chapter05-Earring}
		\caption{Hawaii 耳环中的若干圆周。图中只画出有限多个；完整空间含全部 \(C_n\)，它们都通过 \(o\)，且直径趋于零。任意小的 \(o\) 的邻域仍包含完整的小圆周。}
		\label{g1:fig:hawaiian-earring}
	\end{figure}

	固定 \(n\)，定义 \(r_n:Z\to C_n\)，在 \(C_n\) 上恒等，在其他圆周上恒取 \(o\)。它连续：在 \(o\) 处有 \(\|r_n(z)\|\le\|z\|\)；在非 \(o\) 点附近，可以避开所有其他圆周，或避开 \(C_n\)，因为足够小的圆周全都挤在 \(o\) 附近，而其余只有有限个。因而 \(r_n\) 是连续缩回。绕 \(C_n\) 一周的环路若在 \(Z\) 中零伦，复合 \(r_n\) 就会使它在 \(C_n\) 中零伦，这与\cref{g1:thm:circle-fundamental}矛盾。

	然而 \(o\) 的每个邻域都包含充分大的 \(n\) 对应的整个 \(C_n\)，因为 \(C_n\subseteq\overline B(o,2/n)\)。所以任意小邻域都有在 \(Z\) 中不能缩掉的环路，\(Z\) 不半局部单连通。

	若存在单连通覆盖 \(q:Y\to Z\)，取 \(o\) 的均匀覆盖邻域 \(U\) 及其上面任意一片。充分小的整个 \(C_n\) 位于 \(U\)；绕它一周的环路可经该片的局部逆提升为 \(Y\) 中的环路。它在单连通的 \(Y\) 中零伦，投影后便在 \(Z\) 中零伦，矛盾。因此 \(Z\) 不存在万有覆盖。这个论证只使用单个圆周的基本群，不需要描述 \(\pi_1(Z,o)\) 的全部结构。
\end{example}

在道路类构造中，失败发生在一个明确位置：如果邻域 \(U\) 中还有在 \(X\) 内非零伦的环路 \(\beta\)，则 \([\alpha]\) 与 \([\alpha*\beta]\) 有相同投影，却未必相同。因此 \(B([\alpha],U)\to U\) 可能不单射，也就不能作为覆盖的一片。

局部道路连通的职责则不同。\cref{g1:exm:topologist-sine}中的拓扑正弦曲线 \(S\) 在竖直线段上的道路分支不是开集，故没有道路连通开邻域基。取它的锥，可以保留局部道路连通性的障碍，同时使整个空间可缩：
\[
	C(S)=\{((1-t)x,(1-t)y,t):(x,y)\in S,\ 0\le t\le1\}\subseteq\mathbb R^3.
\]
顶点为 \(v=(0,0,1)\)。每个点都能沿直线连到 \(v\)，且直线同伦 \((z,s)\mapsto(1-s)z+sv\) 始终留在锥中，因此 \(C(S)\) 可缩，特别是道路连通、单连通并半局部单连通。不过在开集 \(V=\{t<1/2\}\cap C(S)\) 内，坐标变换
\[
	((1-t)x,(1-t)y,t)\longmapsto((x,y),t)
\]
给出 \(V\cong S\times[0,1/2)\)。点 \(((0,0),0)\) 的每个相对开邻域都遇到正弦图像上的点；若它有包含于 \(V\) 的道路连通开邻域，投影到 \(S\) 就会给出从竖直线段到正弦图像的道路，与\cref{g1:exm:topologist-sine}矛盾。因此这个单连通空间仍不局部道路连通。在它上面，恒等映射当然是万有覆盖，但上述“在任意给定邻域内取好邻域”的拓扑基论证不能照搬。

在道路类构造中，半局部单连通保证局部单射，局部道路连通保证好邻域构成原拓扑的基。锥的例子同时说明，特定空间即使缺少后一个条件，也可能有万有覆盖。关于本节道路类模型与覆盖分类的进一步阅读，可参见\cite[第1.3节]{Hatcher2002}。

% !TeX root = ../../Book1.tex
% 纲要标题：### 5.5 群作用、单值与覆盖的重建
% 纲要位置：几何学计划纲要.md，第 306 行。
% 正文状态：初稿；尚未正式审核。

% 纲要内容开始
% 
% * 适当不连续作用构造商覆盖
% * 纤维上的置换单值
% * 由单值数据重建覆盖
% * 环面、曲面及图覆盖的计算
% 
% ---
% 
% 纲要内容结束

\section{群作用、单值与覆盖的重建}
\label{g1:sec:0505}

沿底空间的道路行走时，覆盖中的起点一旦选定，整条提升就被确定了。如果底空间中的道路回到原点，提升却可能到达同一纤维的另一个点。于是环路给出纤维的一次置换。上一节用子群描述一个连通覆盖；本节把纤维的所有点一起保留，用这种置换记录整个覆盖，也允许覆盖不连通。

\subsection{适当不连续作用构造商覆盖}
\label{g1:subsec:0505:01}

第四章已经提醒，群轨道的商不一定是良好的几何空间：有理平移的商甚至只有平凡拓扑。要得到覆盖，所需条件相当具体：每一点应有一个小邻域，它的不同平移副本互不相交。这样，一个小邻域在取商时只被复制和识别，不会在内部折叠。

\begin{definition}[用于覆盖的适当不连续作用]\label{g1:def:properly-discontinuous-covering}
	设群 \(\Gamma\) 通过同胚左作用于拓扑空间 \(E\)。本节称该作用\term[shidangbulianxuzuoyong]{适当不连续}[properly discontinuous]，如果对每个 \(e\in E\)，都存在开邻域 \(V\ni e\)，使
	\[
		gV\cap V=\varnothing\qquad(g\in\Gamma,\ g\ne1).
	\]
	此约定包含自由性：若 \(ge=e\)，上述邻域立即迫使 \(g=1\)。有的文献允许非平凡点稳定子群而使用同一术语，应用其结论时需要核对定义。
\end{definition}

\begin{proposition}[作用的商是覆盖]\label{g1:prop:action-quotient-cover}
	设 \(\Gamma\) 在非空拓扑空间 \(E\) 上有上述作用，商集 \(E/\Gamma\) 取商拓扑。则轨道映射 \(q:E\rightarrow E/\Gamma\) 为覆盖。若 \(E\) 连通，则这个覆盖的甲板变换群恰为给定的 \(\Gamma\)；若 \(E\) 还单连通，则它是一个万有覆盖。
\end{proposition}
\begin{proof}
	对任意开集 \(O\subseteq E\)，有
	\[
		q^{-1}(q(O))=\bigcup_{g\in\Gamma}gO,
	\]
	右边开，故 \(q\) 为开映射。取定义中的 \(V\)。它的各个平移 \(gV\) 两两不交：若 \(gV\cap hV\ne\varnothing\)，则 \(h^{-1}gV\cap V\ne\varnothing\)，从而 \(g=h\)。由于 \(q|_V\) 连续、单射、开，且映满 \(q(V)\)，它为同胚；每个 \(q|_{gV}\) 同样为同胚。因此
	\[
		q^{-1}(q(V))=\bigsqcup_{g\in\Gamma}gV
	\]
	正是覆盖的局部形式。

	每个 \(g\) 保持轨道，故为甲板变换。反过来，若 \(E\) 连通且 \(T\) 是任意甲板变换，固定 \(e\in E\)，由同纤维就是同轨道，可找唯一 \(g\) 使 \(T(e)=ge\)。映射 \(T\) 与 \(g\) 都是 \(q\) 到自身的提升，由\cref{g1:prop:deck-determined}，在一点相等便处处相等。所以没有额外的甲板变换。单连通时，按万有覆盖的定义即得最后结论。
\end{proof}

整数群在 \(\mathbb R\) 上的平移满足此条件，因为长度小于 \(1\) 的开区间与其非零整数平移不交。由 \(t\mapsto e^{2\pi it}\) 可识别 \(\mathbb R/\mathbb Z\cong S^1\)，得到圆周的万有覆盖。反之，\(\mathbb Q\) 的平移在任何开区间中都有非零的小平移与自身相交，恰好失去这一条件。反射 \(x\mapsto-x\) 的作用在原点有稳定子群，其商映射 \(\mathbb R\rightarrow[0,\infty)\) 在原点附近也不是覆盖：任一含原点的开区间都包含两个被识别的点，不能成为一张同胚覆盖片。

\begin{proposition}[万有覆盖上的基本群作用]\label{g1:prop:universal-deck-discontinuous}
	设 \(X\) 非空、道路连通、局部道路连通且半局部单连通，\(\widetilde X\) 是从基点 \(x_0\) 出发的道路类万有覆盖，\(G=\pi_1(X,x_0)\)。由
	\[
		L_g[\alpha]=[\gamma*\alpha]\qquad(g=[\gamma])
	\]
	给出的 \(G\) 左作用适当不连续，而且在每个纤维上自由且传递。
\end{proposition}
\begin{proof}
	这些映射及其群乘法由\cref{g1:thm:normal-covering-deck}给出。若道路 \(\alpha,\beta\) 终点相同，取 \(g=[\beta*\bar\alpha]\)，便有 \(L_g[\alpha]=[\beta]\)，故作用在纤维上传递；自由性来自甲板变换由一点决定。

	取底空间一个连通的均匀覆盖邻域 \(U\)，并取上方一张片 \(V\)。每个 \(L_g\) 把 \(V\) 映到 \(U\) 上方的一张片，因为它保持投影，且片可由道路提升从其中一点得到。若 \(L_gV\) 与 \(V\) 相交，两片必相同。此时投影在 \(V\) 上单射，迫使 \(L_g|_V=\operatorname{id}_V\)，再由甲板变换的确定性得到 \(g=1\)。因此 \(V\) 的不同群平移两两不交。
\end{proof}

这里的“均匀覆盖邻域”指其原像分解成同胚覆盖片的开邻域；连通的选择由底空间局部道路连通保证。它与前一节道路类构造中的好邻域是同一类局部工具。

\subsection{纤维上的置换单值}
\label{g1:subsec:0505:02}

现在固定任意覆盖 \(p:E\rightarrow X\) 和基点 \(x_0\in X\)，令 \(S=p^{-1}(x_0)\)。这一纤维是离散空间，但其点不需要预先排列成编号。我们把“从哪一点出发，沿环路以后到哪一点”作为置换的定义。

\begin{proposition}[单值的右作用]\label{g1:prop:monodromy-action}
	设 \(p:E\rightarrow X\) 为覆盖，\(x_0\in X\)，\(G=\pi_1(X,x_0)\)。对 \(s\in S=p^{-1}(x_0)\) 与 \(g=[\gamma]\in G\)，定义
	\[
		s\cdot g=\widetilde\gamma_s(1),
	\]
	其中 \(\widetilde\gamma_s\) 是从 \(s\) 出发的唯一道路提升。则此公式良定义，并构成 \(G\) 在 \(S\) 上的右作用。这称为覆盖的\term[danzhizuoyong]{单值作用}[monodromy action]。
\end{proposition}
\begin{proof}
	\cref{g1:cor:lifted-path-homotopy}保证同伦的环路从同一纤维点提升后有相同终点，所以结果只依赖 \(g\)。常值道路的提升常值，故 \(s\cdot1=s\)。先走 \(\gamma\) 再走 \(\eta\) 时，可以先提升 \(\gamma\)，再从其终点提升 \(\eta\)；两段拼接是 \(\gamma*\eta\) 的提升，由唯一性就是全部提升。因此
	\[
		(s\cdot g)\cdot h=s\cdot(gh).
	\]
	沿反路提升把终点送回起点，所以每个 \(g\) 给出的映射均为置换，逆置换由 \(g^{-1}\) 给出。
\end{proof}

本书先走左边道路的乘法约定，自然产生右作用。若把 \(s\mapsto s\cdot g\) 记作置换 \(M_g\)，并按通常函数复合记乘法，则 \(M_{gh}=M_h\circ M_g\)。若希望得到按通常复合的群同态 \(G\rightarrow\operatorname{Sym}(S)\)，应取 \(g\mapsto M_{g^{-1}}\)。后面直接使用右作用等式，避免在两个约定之间来回转换。

\begin{proposition}[轨道、稳定子群与覆盖的连通性]\label{g1:prop:monodromy-components}
	设 \(X\) 非空、道路连通且局部道路连通，\(p:E\rightarrow X\) 为覆盖，\(S=p^{-1}(x_0)\)。
	\begin{enumerate}
		\item \(S\) 的单值轨道与 \(E\) 的连通分支一一对应。因此 \(E\) 连通当且仅当单值作用传递。
		\item 对 \(s\in S\)，其稳定子群是 \(p_*\pi_1(E,s)\subseteq G\)。
		\item 设该子群为 \(H\)。轨道 \(sG\) 与陪集集合 \(H\backslash G=\{Hg:g\in G\}\) 右等变双射，对应为 \(Hg\mapsto s\cdot g\)。
	\end{enumerate}
\end{proposition}
\begin{proof}
	由覆盖的局部形式，\(E\) 局部道路连通，故由\cref{g1:thm:locally-path-components}，连通分支就是道路分支。每个分支都遇到 \(S\)：从任一点的投影沿道路到 \(x_0\)，提升后就在同一分支中到达 \(S\)。两点 \(s,t\in S\) 能在 \(E\) 中用道路相连，当且仅当某个底空间环路的提升从 \(s\) 到 \(t\)，这正是二者属于同一单值轨道，得到第一项。

	等式 \(s\cdot g=s\) 意味着对应提升是闭路；由\cref{g1:cor:closed-lift-criterion}，这恰好等价于 \(g\in p_*\pi_1(E,s)\)，得到第二项。最后，
	\[
		s\cdot g=s\cdot h
		\quad\Longleftrightarrow\quad gh^{-1}\in H
		\quad\Longleftrightarrow\quad Hg=Hh.
	\]
	故陪集对应良定义且单射，满射由轨道的定义得到；右乘的相容性也直接来自作用等式。
\end{proof}

例如 \(z\mapsto z^n\)（\(n\ge1\)）的纤维可标为 \(s_j=e^{2\pi ij/n}\)，\(j\in\mathbb Z/n\mathbb Z\)。一周正向环路从 \(s_j\) 提升为 \(e^{2\pi i(t+j)/n}\)，终点是 \(s_{j+1}\)。所以基本群的生成元在纤维上作一个 \(n\)-循环，其稳定子群为 \(n\mathbb Z\)。一个局部选出的 \(n\) 次根沿圆周继续一周以后会改变分支，这就是单值作用所记录的现象。

\subsection{由单值数据重建覆盖}
\label{g1:subsec:0505:03}

反过来，若只给一个集合 \(S\)，以及底空间环路如何置换 \(S\) 的规则，能否找回覆盖？万有覆盖已经把每条道路的行走历史保存为道路类；把历史改变时纤维编号应如何改变一并规定，就能作出所需商空间。

\begin{theorem}[由右群作用重建覆盖]\label{g1:thm:monodromy-reconstruction}
	设 \(X\) 非空、道路连通、局部道路连通且半局部单连通，固定 \(x_0\)，令 \(G=\pi_1(X,x_0)\)，\(\widetilde p:\widetilde X\rightarrow X\) 为道路类万有覆盖。对任意非空右 \(G\) 集合 \(S\)，赋予 \(S\) 离散拓扑，并在 \(S\times\widetilde X\) 上定义左作用
	\[
		g\cdot(s,u)=(s\cdot g^{-1},L_g u).
	\]
	则轨道商
	\[
		E_S=(S\times\widetilde X)/G,\qquad p_S([s,u])=\widetilde p(u)
	\]
	是 \(X\) 上的覆盖。其基点纤维自然等于 \(S\)，单值作用就是原给定右作用。每个非空覆盖均由此得到，两个这样的覆盖同构当且仅当其纤维右 \(G\) 集合同构。
\end{theorem}
\begin{proof}
	先核验群作用方向。连续施以 \(h\) 再施以 \(g\)，第一坐标变为 \(s\cdot h^{-1}g^{-1}=s\cdot(gh)^{-1}\)，第二坐标变为 \(L_{gh}u\)，所以确为左作用。每个作用映射为同胚，且保持 \(\widetilde p(u)\)，因此 \(p_S\) 良定义且连续。商映射 \(q:S\times\widetilde X\rightarrow E_S\) 开，证明与\cref{g1:prop:action-quotient-cover}相同。

	取连通的均匀覆盖邻域 \(U\subseteq X\) 和万有覆盖的一张片 \(V\)。由\cref{g1:prop:universal-deck-discontinuous}，其他各片恰为 \(L_gV\)，而且指标 \(g\) 唯一。对每个 \(s\in S\)，令 \(W_s=q(\{s\}\times V)\)。这些集合开，并且 \(q\) 将 \(\{s\}\times V\) 同胚地映到 \(W_s\)：若其中两点被识别，则某个 \(L_gV\) 与 \(V\) 相交，迫使 \(g=1\)。同样的论证说明不同 \(s\) 给出的 \(W_s\) 不交。若 \(u=L_gv\)、\(v\in V\)，轨道关系给出
	\begin{equation}\label{g1:eq:associated-cover-relation}
		[s,L_gv]=[s\cdot g,v],
	\end{equation}
	所以这些 \(W_s\) 覆盖 \(p_S^{-1}(U)\)，且每个 \(p_S|_{W_s}\) 与 \(\widetilde p|_V\) 对应，为到 \(U\) 的同胚。这证明覆盖性质。

	记万有覆盖基点为 \(\widetilde x_0=[e_{x_0}]\)。纤维点写成 \([s,\widetilde x_0]\) 的方式唯一，故基点纤维与 \(S\) 相同。从此点提升代表 \(g\) 的环路 \(\gamma\)，可以使用道路 \(t\mapsto[s,\widetilde\gamma(t)]\)，其中 \(\widetilde\gamma\) 从 \(\widetilde x_0\) 出发。其终点为
	\[
		[s,L_g\widetilde x_0]=[s\cdot g,\widetilde x_0],
	\]
	所以单值恰为指定的右作用。

	再设 \(p:E\rightarrow X\) 是任意覆盖，\(S=p^{-1}(x_0)\) 取其单值作用。对每个 \(s\)，由\cref{g1:thm:lifting-criterion}，\(\widetilde p\) 有唯一连续提升 \(F_s:\widetilde X\rightarrow E\) 满足 \(F_s(\widetilde x_0)=s\)：定义域道路连通、局部道路连通且单连通，因此基本群包含条件自动成立。具体地，\(F_s([\alpha])\) 是 \(\alpha\) 从 \(s\) 出发的提升终点。映射 \(F(s,u)=F_s(u)\) 在离散积上连续。先从 \(s\cdot g^{-1}\) 提升 \(\gamma\)，会回到 \(s\)，然后提升 \(\alpha\)，故
	\[
		F(s\cdot g^{-1},L_g[\alpha])=F(s,[\alpha]).
	\]
	它因而下降为连续映射 \(\bar F:E_S\rightarrow E\)，并保持到底空间的投影。

	给定 \(x\in X\) 及一条从 \(x_0\) 到 \(x\) 的道路 \(\alpha\)，沿 \(\alpha\) 提升给出纤维 \(S\rightarrow p^{-1}(x)\) 的双射，其逆由反路提升给出。另一方面，\(E_S\) 中 \(x\) 上方的点也唯一写成 \([s,[\alpha]]\)。所以 \(\bar F\) 在每个纤维上双射，因而全局双射。它在足够小的共同均匀覆盖邻域上逐片等于两张同胚片之间的坐标对应，故为局部同胚；双射局部同胚是同胚。因此找回了原覆盖。

	最后，覆盖同构限制到基点纤维后与道路提升相容，故是右 \(G\) 集合的同构。反过来，右等变双射 \(\phi:S\rightarrow S'\) 使 \((s,u)\mapsto(\phi(s),u)\) 与轨道关系相容，下降为 \(E_S\rightarrow E_{S'}\) 的同构，逆映射由 \(\phi^{-1}\) 给出。
\end{proof}

若允许空总空间，空集合也对应一个空的局部覆盖；本章覆盖按满射定义，所以定理明确使用非空 \(S\)。在连通情形，传递右 \(G\) 集合写成 \(H\backslash G\)，便恢复上一节的子群分类。这里的重建还保留了不同轨道，因此同时处理不连通覆盖。

\begin{corollary}[甲板变换与等变置换]\label{g1:cor:deck-centralizer}
	在上一定理的假设下，任一覆盖 \(E\rightarrow X\) 的甲板变换群，等于其纤维 \(S\) 上与全部单值置换交换的置换群。
\end{corollary}
\begin{proof}
	甲板变换将从 \(s\) 出发的道路提升送到从 \(T(s)\) 出发的提升，因此 \(T(s\cdot g)=T(s)\cdot g\)。反过来，每个满足此等式的纤维置换由定理的重建公式扩张为甲板变换。扩张唯一：每一点都是某条从基点纤维出发的提升道路的终点，而该终点的像已由纤维置换决定。限制与扩张保持函数复合，所以给出群同构。
\end{proof}

单值作用和甲板变换作用因此不是同一件事。前者由底空间环路推动纤维点，后者是整个覆盖保持投影的对称；后者必须与前者交换。

\subsection{环面、曲面及图覆盖的计算}
\label{g1:subsec:0505:04}

先看环面。由\cref{g1:prop:product-fundamental}，\(\pi_1(T^2)\cong\mathbb Z^2\)，而逐坐标指数映射
\[
	P:\mathbb R^2\longrightarrow T^2,\qquad P(u,v)=(e^{2\pi iu},e^{2\pi iv})
\]
为万有覆盖：两个圆周覆盖片的积给出覆盖局部形式，\(\mathbb R^2\) 凸而单连通。甲板群是整数平移。对子群 \(H\le\mathbb Z^2\)，得到连通覆盖 \(\mathbb R^2/H\rightarrow T^2\)，纤维为 \(H\backslash\mathbb Z^2\)。因为底群交换，这些连通覆盖全部正规。

\begin{example}[一个六叶环面覆盖]\label{g1:exm:torus-six-cover}
	取 \(H=\langle(2,0),(1,3)\rangle\)。线性映射 \(A(u,v)=(2u+v,3v)\) 将 \(\mathbb Z^2\) 映到 \(H\)，并诱导同胚 \(\mathbb R^2/\mathbb Z^2\rightarrow\mathbb R^2/H\)。所以相应覆盖可写成
	\[
		f:T^2\longrightarrow T^2,\qquad f(z,w)=(z^2w,w^3).
	\]
	每个陪集都有唯一代表 \((r,s)\)，其中 \(r=0,1\)，\(s=0,1,2\)：先用 \((1,3)\) 把第二坐标约化到这三个值，再用 \((2,0)\) 约化第一坐标；若两个代表之差在 \(H\)，第二坐标差迫使 \((1,3)\) 的系数为零，再由第一坐标得到两代表相同。因此叶数为 \(6\)。

	令 \(e_1=(1,0),e_2=(0,1)\)。商群中 \(2e_1=0\)、\(e_1+3e_2=0\)，所以 \(e_1=-3e_2\)，全群由 \(e_2\) 生成且有六个元素，即甲板群为 \(\mathbb Z/6\mathbb Z\)。沿两个圆周生成环路的单值分别是陪集加 \(e_1,e_2\)，在这一循环编号下分别为加 \(3\)、加 \(1\)。这同时给出了覆盖公式、叶数、单值与甲板群。
\end{example}

对由标准边字
\[
	[a_1,b_1]\cdots[a_g,b_g],\qquad [a,b]=aba^{-1}b^{-1},
\]
得到的可定向多边形曲面模型 \(\Sigma_g\)，先对其 \(2g\) 条边组成的一点并图使用\cref{g1:prop:graph-fundamental}，再由\cref{g1:prop:attach-disk}沿该边字附着圆盘，就得到以这些生成元和该关系为展示的基本群（\(g\ge1\)）。构造 \(n\) 叶覆盖时，可在 \(S=\{1,\ldots,n\}\) 上指定各生成元的右置换；唯一需要检查的关系是按该边字依次作用后恒等。这保证自由群作用通过所展示的商群下降，再由\cref{g1:thm:monodromy-reconstruction}得到覆盖。生成置换的轨道决定覆盖的连通分支。这里使用的是明确的多边形模型及其关系，不需要预先知道所有紧曲面的分类。

例如 \(g\ge1\) 时，令 \(a_1\) 作 \(n\)-循环，其余 \(a_i,b_i\) 全部恒等。每个交换子作用都恒等，所以关系满足；\(n\)-循环使作用传递，故得到连通 \(n\) 叶覆盖。单值通过满同态 \(\pi_1(\Sigma_g)\rightarrow\mathbb Z/n\mathbb Z\) 的正则作用给出，稳定子群就是该同态的核，因而此覆盖正规，甲板群为 \(\mathbb Z/n\mathbb Z\)。

\begin{example}[图覆盖中单值与甲板群的差别]\label{g1:exm:three-sheet-graph}
	设 \(B\) 为两个圆周 \(a,b\) 在一点相接的图，\(\pi_1(B)=F(a,b)\)。在三个纤维点上规定右置换
	\[
		s\cdot a=(12)(s),\qquad s\cdot b=(23)(s).
	\]
	自由群没有关系，故这给出一个右作用：一个字从左到右逐个作用即可。构造覆盖图时，在顶点 \(1,2,3\) 之间，为每个 \(s\) 放一条从 \(s\) 到 \(s\cdot a\) 的有向 \(a\)-边，再放一条从 \(s\) 到 \(s\cdot b\) 的有向 \(b\)-边。每条边按参数同胚地映到基图相应边的内部。

	在每个顶点附近，各有一条 \(a\)-边进入、一条 \(a\)-边离开，以及对应的两条 \(b\)-半边；环边也分别贡献进入、离开的半边。因此顶点的小星形邻域同胚地映到基图基点的小星形邻域，边内部也显然满足覆盖条件，得到三叶图覆盖。两次换位生成传递作用，所以覆盖连通。图有三个顶点和六条边，由\cref{g1:prop:graph-fundamental}，其基本群为秩 \(6-3+1=4\) 的自由群。

	单值置换生成 \(S_3\)，但甲板群平凡。事实上，若置换 \(T\) 同时与 \((12)\)、\((23)\) 交换，它必须分别保留二者唯一的不动点 \(3,1\)，于是也固定 \(2\)，所以 \(T=\operatorname{id}\)。由\cref{g1:cor:deck-centralizer}，覆盖没有非平凡甲板变换。覆盖连通只说明单值传递，并不保证甲板变换能把任意纤维点送到另一个。
\end{example}

\begin{figure}[htbp]
	\centering
	\input{Figures/Chapter05-GraphCover}
	\caption{二圈图的三叶覆盖。蓝色边标记 \(a\)，红色边标记 \(b\)；同一对顶点间的两条反向边是两条不同的边。每个顶点都有各标记的一入一出，纤维置换分别为 \((12)\) 与 \((23)\)。}
	\label{g1:fig:three-sheet-graph}
\end{figure}

这些例子显示两种计算可以相互检查。先从图形或粘合求基本群，再用它在纤维上的作用建立覆盖；反过来，提升终点、叶数和甲板变换又能检验原来的群与关系。关于覆盖、群作用和图的进一步阅读，可参阅\cite[Section 1.3; Section 1.A]{Hatcher2002}。


\begin{chaptersummary}
	定端点同伦允许改变道路的形状与速度，但保留起终点。环路类的拼接给出基本群，连续映射诱导群同态；沿道路改变基点得到同构，不同选择之间相差内自同构。形变收缩保留基本群，圆周角度的提升把基本群识别为整数绕数，积空间的基本群则逐坐标计算。

	\ProseParagraphBreak

	van Kampen 定理将粘合转成自由积与关系。细分道路给出生成元，细分同伦方形说明全部关系由交集中的环路产生。有限连通图的生成元来自生成树之外的边，附着圆盘将边界词的正规闭包消去；这种计算同时解释了图、去点空间和曲面模型中环路为何保留或消失。

	\ProseParagraphBreak

	覆盖的均匀局部结构使道路逐段唯一提升，同伦提升保证终点只依赖道路类。覆盖在基本群上诱导单射，一个映射能否提升由其基本群像是否落入覆盖子群判断。覆盖变换由一个点的像决定，需要总空间连通；它反映整个覆盖的对称，而不能仅从纤维点数推断。

	\ProseParagraphBreak

	道路类构造把“到达同一点的不同历史”保留下来。局部道路连通提供适用的道路邻域，半局部单连通保证小环路在全空间内零伦，从而让终点映射在基本邻域上单射。带基点连通覆盖对应基本群子群，忘掉基点对应共轭类，正规覆盖对应正规子群；一般覆盖变换群为 \(N_G(H)/H\)，正规时才是 \(G/H\)。

	\ProseParagraphBreak

	沿环路提升使基本群右作用于纤维，轨道给出覆盖的连通分支，稳定子给出相应子群。反过来，把这个右作用与万有覆盖的左覆盖变换作用配对取商，就能重建覆盖。覆盖变换群由与全部单值置换交换的置换组成；环面和图的计算说明，单值传递、覆盖正规以及覆盖变换群传递必须分别判断。
\end{chaptersummary}

\BookExercises
以下各题均给出完整解答。计算基本群时注明基点及所用粘合条件；构造覆盖时检查局部拓扑，并区分单值作用、子群与覆盖变换。

\begin{exercise}\label{g1:ex:loop-disk-extension}
	设 \(f:S^1\rightarrow X\) 连续，\(x_0=f(1)\)，并令 \(\alpha(t)=f(e^{2\pi it})\)。证明 \([\alpha]=1\in\pi_1(X,x_0)\)，当且仅当 \(f\) 可以连续延拓为 \(F:D^2\rightarrow X\)，其中 \(D^2=\{z\in\mathbb C:|z|\le1\}\)。
\end{exercise}
\begin{solution}
	若 \(H(t,u)\) 是把 \(\alpha\) 固定端点收缩到 \(x_0\) 的同伦，定义
	\[
		F(re^{2\pi it})=H(t,1-r)\qquad(0\le r\le1).
	\]
	当 \(r=0\) 时右边恒为 \(x_0\)，当 \(t=0,1\) 时两式也都为 \(x_0\)，所以选择极角无关。为验证圆心处的连续性而不依赖极坐标直觉，使用连续满射 \(q:[0,1]^2\rightarrow D^2\)、\(q(t,r)=re^{2\pi it}\)。定义域紧、目标 Hausdorff，所以 \(q\) 闭，从而是商映射。连续函数 \((t,r)\mapsto H(t,1-r)\) 在它的每个纤维上常值，故由商下降得到连续 \(F\)。边界 \(r=1\) 处的值正是 \(f\)。

	反之，若 \(F\) 存在，取
	\[
		H(t,u)=F\bigl((1-u)e^{2\pi it}+u\bigr).
	\]
	单位圆盘凸，括号中的点始终在圆盘内。\(u=0\) 时得到 \(\alpha\)，\(u=1\) 时恒为 \(F(1)=x_0\)；\(t=0,1\) 时括号中都为 \(1\)，所以两端固定。这就是所需同伦。
\end{solution}

\begin{exercise}\label{g1:ex:no-disk-boundary-retraction}
	证明不存在连续映射 \(r:D^2\rightarrow S^1\) 满足 \(r|_{S^1}=\operatorname{id}_{S^1}\)。为什么这不妨碍穿孔圆盘向边界作径向形变收缩？
\end{exercise}
\begin{solution}
	若存在，包含映射 \(i:S^1\hookrightarrow D^2\) 满足 \(r\circ i=\operatorname{id}\)，所以 \(r_*i_*=\operatorname{id}_{\mathbb Z}\)。然而圆盘凸，基本群平凡，故 \(i_*\) 的像为零，复合 \(r_*i_*\) 也只能是零同态，与 \(\mathbb Z\) 的恒等同态不同。因此缩回不存在。

	穿孔圆盘 \(\{0<|z|\le1\}\) 上 \(r(z)=z/|z|\) 连续，且
	\[
		H(z,u)=\left(1-u+\frac{u}{|z|}\right)z
	\]
	的模为 \((1-u)|z|+u\in(0,1]\)，故始终留在穿孔圆盘内，并固定边界。这个公式在 \(z=0\) 无定义，不能连续补上；正是删去圆心改变了基本群。
\end{solution}

\begin{exercise}\label{g1:ex:topological-group-pi1-abelian}
	设 \(K\) 是拓扑群，单位元为 \(e\)；这里群乘法 \(K\times K\rightarrow K\) 和取逆都连续。证明 \(\pi_1(K,e)\) 为交换群，不要求 \(K\) 本身交换。
\end{exercise}
\begin{solution}
	取两个基于 \(e\) 的环路 \(\alpha,\beta\)，考虑方形映射 \(F(s,t)=\alpha(s)\beta(t)\)。沿下边再右边从 \((0,0)\) 到 \((1,1)\)，其像是 \(\alpha*\beta\)；沿左边再上边，其像是 \(\beta*\alpha\)，因为 \(\alpha,\beta\) 在端点都取 \(e\)。这两条方形内的折线路径可以定端点线性插值，插值仍留在凸方形中。与 \(F\) 复合，就给出 \(\alpha*\beta\simeq\beta*\alpha\) 的定端点同伦。所以任意两个基本群元素交换。整个证明没有交换 \(K\) 中两个任意元素，只比较了方形边界上因端点取单位元而简化的乘积。
\end{solution}

\begin{exercise}\label{g1:ex:torus-induced-matrix}
	对映射 \(f:T^2\rightarrow T^2\)、\(f(z,w)=(z^2w^{-1},z^3w^2)\)，在两个坐标圆周环路给出的 \(\pi_1(T^2,(1,1))\cong\mathbb Z^2\) 下求 \(f_*\)。证明 \(f\) 不是同胚，也不与常值映射同伦。
\end{exercise}
\begin{solution}
	第一生成环路 \(a(t)=(e^{2\pi it},1)\) 的像绕数为 \((2,3)\)，第二生成环路 \(b(t)=(1,e^{2\pi it})\) 的像绕数为 \((-1,2)\)。所以
	\[
		f_*\binom{m}{n}
		=\begin{pmatrix}2&-1\\3&2\end{pmatrix}\binom{m}{n}
		=\binom{2m-n}{3m+2n}.
	\]
	它不满射：若像为 \((1,0)\)，由 \(n=2m-1\) 得 \(7m=2\)，没有整数解。同胚必须诱导群同构，故 \(f\) 不是同胚。若 \(f\) 与常值映射同伦，基点可能在同伦中移动，但\cref{g1:prop:moving-basepoint}说明相应诱导同态只相差一个基点变化同构。常值映射的诱导同态为零，于是 \(f_*\) 也须为零，与 \(f_*(1,0)=(2,3)\) 矛盾。
\end{solution}

\begin{exercise}\label{g1:ex:radial-winding}
	在 \(\mathbb C\setminus\{0\}\) 中考虑从 \(3\) 出发并回到 \(3\) 的环路
	\[
		\gamma(t)=(2+\cos2\pi t)e^{4\pi it}.
	\]
	求它经径向缩回到圆周后的绕数，写出与两周标准圆周环路之间的自由同伦（中间各曲线仍为环路，但基点允许移动），并说明该同伦为什么不是原基点 \(3\) 处的定端点同伦。
\end{exercise}
\begin{solution}
	因为 \(2+\cos2\pi t\ge1\)，\(\gamma\) 从不为零，其径向缩回为 \(e^{4\pi it}\)，连续角度为 \(2t\)，所以绕数为 \(2\)。公式
	\[
		H(t,u)=\bigl((1-u)(2+\cos2\pi t)+u\bigr)e^{4\pi it}
	\]
	给出连续同伦，每个 \(u\) 处都是环路，且系数始终正。在 \(t=0,1\) 处，它的值为 \(3-2u\)，从 \(3\) 移动到 \(1\)，所以端点虽彼此相同却没有固定不动。若要在基点 \(3\) 处比较基本群元素，应把圆周环路与从 \(3\) 到 \(1\) 的径向道路及其反路拼接，使用\cref{g1:prop:basepoint-change}，不能把不同基点的两个环路直接认作同一基本群中的元素。
\end{solution}

\begin{exercise}\label{g1:ex:free-product-orders}
	设 \(G=\langle a,b\mid a^2=1,\ b^3=1\rangle\)。证明 \(G\cong(\mathbb Z/2\mathbb Z)*(\mathbb Z/3\mathbb Z)\)，证明 \(ab\) 有无限阶，并判断 \(aba^{-1}b^{-1}\) 是否为单位元。
\end{exercise}
\begin{solution}
	把 \(a\) 送到第一循环因子的生成元，把 \(b\) 送到第二循环因子的生成元，两个关系都成立，故群展示的泛性质给出到自由积的同态。反方向，由 \(a^2=1,b^3=1\)，两个循环群分别同态地映入 \(G\)，自由积的泛性质给出反向同态。两个复合在 \(a,b\) 上恒等，故在它们生成的群上恒等，得到所述同构。

	在自由积中，正整数 \(k\) 对应的词 \((ab)^k\) 有 \(2k\) 个非单位字母，相邻字母分属两个不同因子，故已经约化且不是空词。因此 \((ab)^k\ne1\)，\(ab\) 有无限阶。又 \(a^{-1}=a\)、\(b^{-1}=b^2\)，所以交换子为 \(abab^2\)。它同样是长度 \(4\) 的约化词，故不是单位元。两个生成元各自有有限阶，并不使自由积有限，也不使它们交换。
\end{solution}

\begin{exercise}\label{g1:ex:cube-graph}
	设 \(G\) 为立方体的顶点与棱组成的图，不包括任何面。计算 \(\pi_1(G)\)。若删去任一条棱的内部，所得图 \(G'\) 的基本群是什么？若改为把整张立方体表面加入 \(G\)，基本群又是什么？说明后两种操作为何不能混为一谈。
\end{exercise}
\begin{solution}
	立方体有 \(8\) 个顶点和 \(12\) 条棱，且图连通，故\cref{g1:prop:graph-fundamental}给出 \(\pi_1(G)\cong F_{12-8+1}=F_5\)。删去任一条棱的内部后，其两端仍可沿包含它的一个正方形面的其余三条棱连接；原先经过所删棱的任何边路都可用这三条棱替换，因此图仍连通。它有 \(11\) 条边、\(8\) 个顶点，故 \(\pi_1(G')\cong F_4\)。

	加入所有面所得空间同胚于 \(S^2\)。若立方体以原点为中心，边长为 \(2\)，由球面到其边界的同胚为
	\[
		u\longmapsto\frac{u}{\lVert u\rVert_\infty},\qquad
		\lVert u\rVert_\infty=\max\{|u_1|,|u_2|,|u_3|\},
	\]
	逆映射为 \(x\mapsto x/\lVert x\rVert\)，两个分母在各自空间中都不为零。正文已通过南北极开覆盖证明 \(\pi_1(S^2)=1\)，所以立方体表面的基本群平凡。删边减少可走的路线；添面保留原来的路线，却提供新的同伦把某些环路消去。这两种空间变化的机制不同，基本群的计算也不同。
\end{solution}

\begin{exercise}\label{g1:ex:degree-disk-attachment}
	对整数 \(m\)，将一个圆盘沿 \(f_m:S^1\to S^1,\ z\mapsto z^m\) 附着到一个圆周上，所得空间记为 \(X_m\)；当 \(m=0\) 时，\(f_0\) 为常值映射。计算 \(\pi_1(X_m)\)，并说明附着圆盘的包含关系为何不保证原圆周的基本群仍嵌入。
\end{exercise}
\begin{solution}
	取圆周上的基点 \(1\)，正向生成元为 \(a\)。沿圆盘边界走一周时，像环路为 \(t\mapsto e^{2\pi imt}\)，其绕数为 \(m\)，故代表 \(a^m\)。附着命题给出
	\[
		\pi_1(X_m)\cong\langle a\mid a^m=1\rangle
		\cong
		\begin{cases}
			\mathbb Z,&m=0,\\
			\mathbb Z/|m|\mathbb Z,&m\ne0.
		\end{cases}
	\]
	当 \(m=0\) 时，关系词已经是单位元，未添加新关系；当 \(m\ne0\) 时，包含映射在基本群上是 \(\mathbb Z\to\mathbb Z/|m|\mathbb Z\)，核为 \(m\mathbb Z\)，不是单射。空间中的点并未因为附着而把原圆周的不同点额外识别，但新增的圆盘允许原圆周中的环路在更大的空间内变形；点之间的映射是否单射，与同伦类之间的映射是否单射，是不同的问题。
\end{solution}

\begin{exercise}\label{g1:ex:puncture-commutator}
	取 \(X=\mathbb R^2\setminus\{p,q\}\)，按正文的去点空间构造选择自由生成元 \(a,b\)，分别绕 \(p,q\) 一周。证明环路类 \(aba^{-1}b^{-1}\) 非平凡，但它在两个包含映射
	\[
		X\hookrightarrow\mathbb R^2\setminus\{p\},\qquad
		X\hookrightarrow\mathbb R^2\setminus\{q\}
	\]
	诱导的同态下都变为单位元。说明这给“用每个缺失点的绕数判断是否可收缩”带来什么限制。
\end{exercise}
\begin{solution}
	\(\pi_1(X)\cong F(\{a,b\})\)，其中 \(aba^{-1}b^{-1}\) 没有任何相邻的逆字母，因而是长度 \(4\) 的非空约化词，不是单位元。

	在只去掉 \(p\) 的空间中，绕 \(q\) 的小圆周可以在不遇到 \(p\) 的小圆盘内收缩；接到基点的道路仅产生共轭，所以 \(b\) 的像为单位元。\(a\) 的像是该单点补空间的正向生成元，整个词的像遂为 \(a\cdot1\cdot a^{-1}\cdot1=1\)。另一个包含映射下 \(a\) 的像为单位元、\(b\) 为生成元，同样得到 \(1\)。

	因此这条环路对每个缺失点的整数绕数都为零，却不能在同时避开两点的空间中收缩。两个绕数只记录两个字母各自的指数和，会丢失先后次序。它们构成到 \(\mathbb Z^2\) 的同态，但不是对非交换自由群元素的完全判别。
\end{solution}

\begin{exercise}\label{g1:ex:torus-slope-attachment}
	取整数 \(p,q\)，沿环路 \(t\mapsto(e^{2\pi ipt},e^{2\pi iqt})\) 向环面 \(T^2\) 附着一张圆盘，所得空间为 \(Y_{p,q}\)。计算其基本群。对 \((p,q)\ne(0,0)\)，令 \(d=\gcd(|p|,|q|)\)，证明结果为 \(\mathbb Z\oplus\mathbb Z/d\mathbb Z\)。
\end{exercise}
\begin{solution}
	环面的生成元 \(a,b\) 交换，所给环路在积空间基本群中的坐标为 \((p,q)\)，亦即元素 \(a^pb^q\)。附着命题给出
	\[
		\pi_1(Y_{p,q})\cong
		\langle a,b\mid aba^{-1}b^{-1}=1,\ a^pb^q=1\rangle
		\cong\mathbb Z^2/\mathbb Z(p,q).
	\]
	若 \(p=q=0\)，所商子群为零，结果为 \(\mathbb Z^2\)。

	若 \((p,q)\ne(0,0)\)，取整数 \(r,s\) 使 \(rp+sq=d\)。这是最大公因数的整系数表示，可由辗转相除逐步回代求得。矩阵
	\[
		A=\begin{pmatrix}r&s\\-q/d&p/d\end{pmatrix}
	\]
	的行列式为 \(1\)，其逆矩阵仍有整数条目，故给出 \(\mathbb Z^2\) 的自同构；同时 \(A(p,q)^{\mathsf T}=(d,0)^{\mathsf T}\)。于是
	\[
		\mathbb Z^2/\mathbb Z(p,q)
		\cong\mathbb Z^2/\mathbb Z(d,0)
		\cong(\mathbb Z/d\mathbb Z)\oplus\mathbb Z.
	\]
	特别地，若 \(p,q\) 互素，结果是无限循环群；若 \(d>1\)，附着关系会同时留下一个自由方向和一个有限阶方向。这里一般并未声称附着后的空间仍为曲面。
\end{solution}

\begin{exercise}\label{g1:ex:pullback-covering}
	设 \(p:E\to X\) 为覆盖，\(f:Y\to X\) 连续。给集合
	\[
		f^*E=\{(y,e)\in Y\times E:f(y)=p(e)\}
	\]
	积空间中的子空间拓扑，定义 \(r(y,e)=y\)。证明 \(r:f^*E\to Y\) 是覆盖。进一步，若连续映射 \(a:Z\to Y\)、\(b:Z\to E\) 满足 \(f\circ a=p\circ b\)，证明存在唯一连续映射 \(c:Z\to f^*E\)，其两个坐标分别为 \(a,b\)。
\end{exercise}
\begin{solution}
	连续性由积投影的连续性得到；对任意 \(y\in Y\)，覆盖 \(p\) 的满射性给出 \(e\) 满足 \(p(e)=f(y)\)，故 \(r\) 满射。固定 \(y_0\)，取 \(f(y_0)\) 的均匀覆盖邻域 \(U\)，写 \(p^{-1}(U)=\bigsqcup_\lambda V_\lambda\)，并令 \(s_\lambda=(p|_{V_\lambda})^{-1}\)。开集 \(W=f^{-1}(U)\) 含 \(y_0\)，且
	\[
		r^{-1}(W)=\bigsqcup_\lambda W_\lambda,\qquad
		W_\lambda=\{(y,s_\lambda(f(y))):y\in W\}.
	\]
	各 \(W_\lambda=f^*E\cap(W\times V_\lambda)\) 在 \(f^*E\) 中开，彼此不交。投影 \(r|_{W_\lambda}\) 的逆为连续映射 \(y\mapsto(y,s_\lambda(f(y)))\)，故它是同胚。因此 \(W\) 是 \(r\) 的均匀覆盖邻域，证明了覆盖性质。

	所求 \(c\) 只能是 \(c(z)=(a(z),b(z))\)，所以唯一。等式 \(f(a(z))=p(b(z))\) 保证其值在 \(f^*E\) 中；两个坐标连续，先给出到 \(Y\times E\) 的连续映射，再由子空间的泛性质得到到 \(f^*E\) 的连续性。这个构造把一个已有覆盖沿任意连续映射搬到新的底空间上，通常称为覆盖的拉回。
\end{solution}

\begin{exercise}\label{g1:ex:finite-cover-composition}
	设 \(p:Y\to X\)、\(q:Z\to Y\) 均为覆盖，并且 \(p\) 的每个纤维有限。证明 \(p\circ q:Z\to X\) 是覆盖，指出有限性用在何处。将结论用于 \(p_n:S^1\to S^1\)、\(p_n(z)=z^n\)（\(n\ge1\)）与 \(q:\mathbb R\to S^1\)、\(q(t)=e^{2\pi it}\)，求复合映射在 \(1\) 上的纤维及其覆盖变换群。
\end{exercise}
\begin{solution}
	固定 \(x\in X\)。取 \(p\) 的均匀覆盖邻域 \(U\)，写 \(p^{-1}(U)=V_1\sqcup\cdots\sqcup V_m\)。片数有限，因为每片恰含纤维 \(p^{-1}(x)\) 的一个点。令 \(y_i=(p|_{V_i})^{-1}(x)\)。对每个 \(i\)，取 \(q\) 的均匀覆盖邻域 \(W_i\)，并缩小使 \(y_i\in W_i\subseteq V_i\)；均匀覆盖邻域的开子集仍均匀覆盖，只需限制各片。令
	\[
		U'=\bigcap_{i=1}^m p(W_i).
	\]
	各 \(p(W_i)\) 都开且含 \(x\)，有限交保证 \(U'\) 仍是开邻域。对每个 \(i\)，令 \(V_i'=(p|_{V_i})^{-1}(U')\)，则 \(V_i'\subseteq W_i\)。将 \(q\) 在 \(W_i\) 上的各片限制到 \(V_i'\)，每个限制都同胚地映到 \(V_i'\)，再经 \(p\) 同胚地映到 \(U'\)。对所有 \(i\) 及全部片取并，恰好得到 \((p\circ q)^{-1}(U')\) 的不交开分解，故复合是覆盖。有限性正用于同时缩小所有 \(y_i\) 上方的邻域，使它们投影后仍有一个共同的开邻域。

	本题复合为 \(r(t)=e^{2\pi int}\)，故 \(r^{-1}(1)=\frac1n\mathbb Z\)。平移 \(t\mapsto t+k/n\) 都是覆盖变换。若 \(h\) 为任意覆盖变换，则 \(n(h(t)-t)\) 是连续的整数值函数；实直线连通，故它恒为某个 \(k\in\mathbb Z\)。所以覆盖变换恰为上述平移，群同构于加法群 \(\mathbb Z\)。这里纤维的间距改变了，但群的抽象同构类型没有改变。
\end{solution}

\begin{exercise}\label{g1:ex:continuous-circle-root}
	设 \(Y\) 道路连通且局部道路连通，\(f:Y\to S^1\) 连续，\(f(y_0)=1\)，\(n\ge1\)。证明存在连续函数 \(g:Y\to S^1\)，使 \(g(y_0)=1\) 且 \(g(y)^n=f(y)\)，当且仅当每条基于 \(y_0\) 的环路 \(\alpha\) 都使 \(f\circ\alpha\) 的绕数属于 \(n\mathbb Z\)。据此判定 \(f(z)=z^m\)（\(z\in S^1,m\in\mathbb Z\)）何时有连续 \(n\) 次根，并列出全部这样的根。
\end{exercise}
\begin{solution}
	等式 \(g^n=f\) 正是把 \(f\) 提升到覆盖 \(p_n:S^1\to S^1\)。由\cref{g1:thm:lifting-criterion}，指定 \(g(y_0)=1\) 的提升存在，当且仅当 \(f_*\pi_1(Y,y_0)\subseteq(p_n)_*\pi_1(S^1,1)\)。在绕数同构下，右端是 \(n\mathbb Z\)，所以这等价于题述条件；指定初值以后根还唯一。

	对 \(f(z)=z^m\)，基本群同态为乘以 \(m\)，像为 \(m\mathbb Z\)。故存在这样的根，当且仅当 \(n\mid m\)。当 \(m=n\ell\) 时，指定初值的根是 \(g(z)=z^\ell\)，所有根为
	\[
		g_j(z)=e^{2\pi ij/n}z^\ell,\qquad j=0,\ldots,n-1.
	\]
	这些函数均满足要求中的幂等式。任意根在 \(z=1\) 处必须取某个 \(n\) 次单位根；除以上述相应常数以后，就成为初值为 \(1\) 的根，由提升唯一性等于 \(z^\ell\)。这也证明列表无遗漏。例如恒等映射 \(S^1\to S^1\) 没有连续平方根，尽管每个点都各自有两个平方根。
\end{solution}

\begin{exercise}\label{g1:ex:interval-cover-trivial}
	设 \(p:E\to[0,1]\) 为覆盖，\(F=p^{-1}(0)\) 带离散拓扑。证明存在同胚
	\[
		\Phi:F\times[0,1]\longrightarrow E,\qquad
		p(\Phi(e,t))=t,\quad \Phi(e,0)=e.
	\]
	因此，若 \(E\) 连通，则 \(p\) 本身是同胚。
\end{exercise}
\begin{solution}
	把 \(H:F\times[0,1]\to[0,1]\) 定义为 \(H(e,t)=t\)，并以包含映射 \(e\mapsto e\) 作为初始提升。由一般参数空间的同伦提升\cref{g1:thm:homotopy-lifting}，得到连续的 \(\Phi\)，满足所写等式。

	固定 \(t\)，对任意 \(a\in p^{-1}(t)\)，从 \(a\) 提升底空间道路 \(s\mapsto(1-s)t\)，终点落在 \(F\) 中，记为 \(e\)。反转这条提升，再由道路提升唯一性，它与 \(u\mapsto\Phi(e,u)\) 在 \(0\le u\le t\) 上相同，所以 \(\Phi(e,t)=a\)。这证明满射性；若 \(\Phi(e,t)=\Phi(e',t)\)，沿同一反向道路提升的唯一性又给出 \(e=e'\)。不同的 \(t\) 有不同投影，故 \(\Phi\) 是双射。

	还须证明逆连续，不能只停在连续双射。固定 \((e,t)\)，在 \(t\) 的一个均匀覆盖邻域内取相对开区间 \(J\ni t\)。道路 \(u\mapsto\Phi(e,u)\) 在 \(J\) 上连续，其像在 \(p^{-1}(J)\) 内；由于 \(J\) 连通，像只能留在一张覆盖片 \(V\) 中。该映射每点的投影都是 \(u\)，所以它就是 \((p|_V)^{-1}\)，且 \(\Phi(\{e\}\times J)=V\)。集合 \(\{e\}\times J\) 在离散积空间中开，\(V\) 在 \(E\) 中开，限制 \(\Phi|_{\{e\}\times J}\) 为同胚。这样的开邻域覆盖整个源空间，故 \(\Phi\) 为开映射，其逆连续。

	若 \(F\) 至少有两个点，则 \(F\times[0,1]\) 可分成一份 \(\{e\}\times[0,1]\) 与其余各份的并，两边都非空开，因而不连通。故 \(E\) 连通时 \(F\) 只有一个点，所构造的同胚就说明 \(p:E\to[0,1]\) 是同胚。
\end{solution}

\begin{exercise}\label{g1:ex:disconnected-deck-fixed-point}
	设 \(X\) 为非空连通空间，\(E=X\times\{1,2,3\}\)，第二个因子取离散拓扑，\(p:E\to X\) 为第一投影。求 \(\operatorname{Deck}(p)\)，并构造有不动点的非恒等覆盖变换，说明\cref{g1:prop:deck-determined}中覆盖空间连通性的作用。
\end{exercise}
\begin{solution}
	每个置换 \(\sigma\in S_3\) 定义覆盖变换 \(h_\sigma(x,j)=(x,\sigma(j))\)，且通常的函数复合满足 \(h_\sigma\circ h_\tau=h_{\sigma\circ\tau}\)。

	反过来，若 \(h\) 为覆盖变换，投影等式迫使 \(h(x,j)=(x,\sigma_x(j))\)，其中对每个 \(x\)，映射 \(\sigma_x\) 为集合 \(\{1,2,3\}\) 的置换。固定 \(j\)，由 \(h\) 连续，函数 \(x\mapsto\sigma_x(j)\) 是连通空间到离散空间的连续映射，所以恒定。对 \(j=1,2,3\) 分别应用这一点，得到整个置换 \(\sigma_x\) 与 \(x\) 无关。于是全部覆盖变换恰为 \(h_\sigma\)，故 \(\operatorname{Deck}(p)\cong S_3\)。

	取交换 \(2,3\) 而固定 \(1\) 的置换，则相应覆盖变换不是恒等，却固定 \(X\times\{1\}\) 中每个点。一个点的像只决定它所在连通分支上的提升，不能决定其他分支上如何交换覆盖片；这就是连通性缺失时唯一性结论失效的原因。
\end{solution}

\begin{exercise}\label{g1:ex:annulus-cover-classification}
	令 \(A=\mathbb C\setminus\{0\}\)，基点取 \(1\)。利用径向形变收缩和本节分类，列出 \(A\) 的全部连通覆盖，并写出各自的覆盖变换群。证明可用以下模型表示它们：
	\[
		p_m:A\longrightarrow A,\quad p_m(re^{i\theta})=re^{im\theta}\quad(m\ge1),
	\]
	以及 \(p_0:\mathbb R\times(0,\infty)\to A\)，\(p_0(t,r)=re^{2\pi it}\)。
\end{exercise}
\begin{solution}
	同伦 \(F(z,s)=((1-s)+s/|z|)z\) 把 \(A\) 强形变收缩到 \(S^1\)，在圆周上始终固定。由\cref{g1:prop:deformation-retract,g1:thm:circle-fundamental}，\(\pi_1(A,1)\cong\mathbb Z\)，其中正向绕原点一周对应 \(1\)。\(A\) 是平面的开子集，每一点有不含原点的凸小圆盘，所以满足分类定理的全部条件。

	\(\mathbb Z\) 的子群恰为 \(m\mathbb Z\)，\(m\ge0\)：非零子群中取最小正整数 \(m\)，用整数除法可知每个元素都是 \(m\) 的倍数；零子群写成 \(0\mathbb Z=\{0\}\)。由于 \(\mathbb Z\) 交换，子群共轭不再作额外识别。

	对 \(m\ge1\)，在一个角度长度小于 \(2\pi\) 的扇形邻域中选定连续辐角，\(p_m\) 有 \(m\) 个局部逆，辐角分别除以 \(m\) 后加上 \(2\pi j/m\)，\(j=0,\ldots,m-1\)。这些逆的像互不相交并穷尽原像，所以 \(p_m\) 是 \(m\) 叶覆盖。其总空间仍是道路连通的 \(A\)；正向一周环路在投影后绕 \(m\) 周，故像子群为 \(m\mathbb Z\)。覆盖变换恰为 \(z\mapsto e^{2\pi ij/m}z\)，其群为 \(\mathbb Z/m\mathbb Z\)。这些变换已给出 \(m\) 个，且由分类的覆盖变换群公式没有其他变换；\(m=1\) 给出恒等覆盖。

	对 \(p_0\)，同一辐角邻域的原像分成整数平移的一列条带，每片投影同胚。总空间 \(\mathbb R\times(0,\infty)\) 凸，故单连通，它对应零子群，是万有覆盖。覆盖变换为 \((t,r)\mapsto(t+k,r)\)，\(k\in\mathbb Z\)。所有子群都已得到，因此这些模型穷尽全部连通覆盖。
\end{solution}

\begin{exercise}\label{g1:ex:intermediate-covering}
	设 \(X\) 道路连通、局部道路连通且半局部单连通，\(G=\pi_1(X,x_0)\)，\(H\le K\le G\)。证明自然映射
	\[
		r_{H,K}:E_H\longrightarrow E_K,\qquad [\alpha]_H\longmapsto[\alpha]_K
	\]
	是保持基点的覆盖，叶数为 \([K:H]\)，且
	\[
		\operatorname{Deck}(E_H/E_K)\cong N_K(H)/H.
	\]
	反过来，若存在保持基点且满足 \(p_KF=p_H\) 的连续映射 \(F:E_H\to E_K\)，证明 \(H\le K\)，并证明它必为上述映射。由此解释，带基点的连通中间覆盖对应包含 \(H\) 的子群。
\end{exercise}
\begin{solution}
	因为 \(H\subseteq K\)，\(H\) 等价的道路也 \(K\) 等价，故 \(r_{H,K}\) 良定义；任意 \([\alpha]_K\) 都有原像 \([\alpha]_H\)，所以满射。在好邻域 \(U\) 上，两空间的基本片都经终点映射同胚到 \(U\)，而 \(r_{H,K}\) 在每一片上的表达就是这两个局部同胚的复合。每个 \(E_K\) 的片的原像是若干 \(E_H\) 的片的不交并，所以 \(r_{H,K}\) 是覆盖。基点 \(e_K\) 的原像由满足 \([\alpha]\in K\) 的环路给出，等价关系是 \(Hg=Hk\)，故纤维为 \(H\backslash K\)，叶数为 \([K:H]\)。

	由覆盖基本群的单射性，\((p_K)_*\) 把 \(\pi_1(E_K,e_K)\) 同构到 \(K\)。等式 \(p_Kr_{H,K}=p_H\) 说明 \((r_{H,K})_*\pi_1(E_H,e_H)\) 在这个识别下恰为 \(H\)。空间 \(E_K\) 继承局部道路连通和半局部单连通：在一个覆盖片中选取来自 \(X\) 的好邻域，片中的环路投影后在 \(X\) 中为零，再由 \((p_K)_*\) 单射知它在 \(E_K\) 中为零。因此可将\cref{g1:thm:normal-covering-deck}用于底空间 \(E_K\)，得到覆盖变换群 \(N_K(H)/H\)。

	反过来，若 \(F\) 保持基点并满足投影交换，则
	\[
		H=(p_H)_*\pi_1(E_H,e_H)
		=(p_K)_*F_*\pi_1(E_H,e_H)\subseteq K.
	\]
	对道路 \(\alpha\)，从 \(e_H\) 到 \([\alpha]_H\) 的初段道路经 \(F\) 变成 \(\alpha\) 从 \(e_K\) 出发的提升。提升唯一性给出 \(F([\alpha]_H)=[\alpha]_K\)，所以 \(F=r_{H,K}\)。任意连通中间覆盖先用分类定理写成 \(E_K\)，以上两方向便说明它恰对应一个 \(H\le K\le G\)。
\end{solution}

\begin{exercise}\label{g1:ex:a4-nonregular-cover}
	设 \(X\) 满足本节分类定理的条件，并给定 \(\pi_1(X,x_0)\cong A_4\)，其中 \(A_4\) 是四个字母上的偶置换群。令 \(H=\langle(123)\rangle\)。证明对应覆盖 \(E_H\to X\) 有四叶，却只有恒等覆盖变换；再证明它不存在位于恒等覆盖与自身之间的非平凡连通中间覆盖。
\end{exercise}
\begin{solution}
	\(|A_4|=12\)、\(|H|=3\)，所以叶数为 \([A_4:H]=4\)。\(H\) 的所有元素共同固定的字母恰为 \(4\)。如果 \(n\in A_4\) 正规化 \(H\)，则 \(n\) 必把这个共同不动点集合送到自身，故 \(n(4)=4\)。固定 \(4\) 的偶置换恰为 \(1,(123),(132)\)，即 \(H\)。因此 \(N_{A_4}(H)=H\)，覆盖变换群为 \(H/H=\{1\}\)。纤维有四点而覆盖变换群轨道只有单点，故该覆盖不正规。

	若 \(H\) 真包含于一个真子群 \(K<A_4\)，由有限群的阶数整除关系，\(|K|\) 只能为 \(6\)。此时 \(K\) 指数为 \(2\)，所以正规，商群为二阶群。任一三循环的阶为 \(3\)，它在二阶商群中的像只能为单位元。但三循环生成 \(A_4\)：\(A_4\) 的非单位元素是三循环或两个不交对换之积，而对互异的 \(a,b,c,d\)，有
	\[
		(ab)(cd)=(acb)(acd)
	\]
	（置换复合按通常函数复合从右向左作用）。所以该商映射把整个 \(A_4\) 都送到单位元，与商群为二阶群矛盾。于是不存在这样的 \(K\)。由\cref{g1:ex:intermediate-covering}，也就不存在非平凡连通中间覆盖。若原中间分解未指定基点，先选择一个 \(E_H\) 的基点，再取它在中间空间的像，即可归入此论证。
\end{solution}

\begin{exercise}\label{g1:ex:finite-cover-compactness}
	设 \(p:E\to X\) 是覆盖，每个纤维均有限，且 \(X\) 紧。证明 \(E\) 紧；此处不假设 Hausdorff。若再假设 \(E\) 紧而 \(X\) Hausdorff，证明每个纤维必有限。举例说明，紧底空间的无限叶覆盖可以不紧。
\end{exercise}
\begin{solution}
	给定 \(E\) 的开覆盖 \(\mathcal W\)。对每个 \(x\in X\)，取均匀覆盖邻域 \(U\)，其片数就是有限数 \(|p^{-1}(x)|\)。在每片的纤维点 \(e\) 处选 \(W_e\in\mathcal W\)，然后把 \(U\) 缩小到开邻域 \(V_x\)，使每片在 \(V_x\) 上的部分都包含在对应 \(W_e\) 中。具体说，把各 \(W_e\) 与其所在片相交后投影，得到有限个含 \(x\) 的开集，再取它们与 \(U\) 的交即可。因此 \(p^{-1}(V_x)\) 被有限个 \(W_e\) 覆盖。

	由 \(X\) 紧，从 \(\{V_x\}\) 中选有限子覆盖，再合并相应有限多个 \(W_e\)，便得到 \(E\) 的有限子覆盖。整个论证只用开覆盖，不用分离性。

	若 \(E\) 紧、\(X\) Hausdorff，则 \(\{x\}\) 闭，纤维 \(p^{-1}(x)\) 是紧空间 \(E\) 的闭子集，因而紧。覆盖的局部片说明纤维是离散空间；离散空间由全部单点构成的开覆盖有有限子覆盖，当且仅当它有限。最后，\(\mathbb R\to S^1\)、\(t\mapsto e^{2\pi it}\) 是紧圆周的无限叶覆盖，总空间 \(\mathbb R\) 不紧。这说明不能删去第一问的有限纤维条件。
\end{solution}

\begin{exercise}\label{g1:ex:earring-cone-semilocal}
	取\cref{g1:exm:hawaiian-earring}的Hawaii 耳环 \(Z\subseteq\mathbb R^2\)，并在 \(\mathbb R^3\) 中定义它的锥
	\[
		C(Z)=\{((1-t)x,(1-t)y,t):(x,y)\in Z,\ 0\le t\le1\}.
	\]
	证明 \(C(Z)\) 可缩，因而半局部单连通；证明它也局部道路连通，但在底部公共点 \(o'=(0,0,0)\) 不局部单连通。说明这个例子怎样区分“在邻域内缩掉”和“在整个空间内缩掉”。
\end{exercise}
\begin{solution}
	记顶点为 \(v=(0,0,1)\)。直线同伦 \(F(z,s)=(1-s)z+sv\) 连续，且始终沿同一条锥的母线移动，故从恒等映射变到常值映射。因此 \(C(Z)\) 可缩。可缩空间中的全部环路在全空间内零伦，所以每点的半局部单连通条件可直接取全空间作为邻域。

	在 \(t<1\) 的部分，坐标除以 \(1-t\) 给出与 \(Z\times[0,1)\) 的同胚。\(Z\) 局部道路连通，区间也局部道路连通；两个道路连通开邻域的积仍道路连通，所以此处局部道路连通。顶点附近可以用截顶锥 \(\{t>a\}\cap C(Z)\) 作为道路连通邻域基。它们都是相对开集，每个点能沿母线连到顶点；又因 \(Z\) 有界，当 \(a\to1\) 时截顶锥的直径趋于零，确实组成顶点的邻域基。

	令 \(V=\{t<1/2\}\cap C(Z)\cong Z\times[0,1/2)\)。若 \(o'\) 有包含于 \(V\) 的单连通开邻域 \(U\)，则 \(U\) 在底部 \(Z\times\{0\}\) 中包含 \(o\) 的一个邻域，因而含有充分小的整个圆周 \(C_n\times\{0\}\)。绕此圆周一周的环路若在 \(U\) 中零伦，也在 \(V\) 中零伦；再复合投影 \(V\to Z\) 及缩回 \(r_n:Z\to C_n\)，便会得到绕 \(C_n\) 一周在圆周内零伦的矛盾。因此 \(U\) 不存在，\(C(Z)\) 不局部单连通。

	这些小环路确实能在整个锥中经过顶点缩掉，却不能在避开顶点的足够小邻域内缩掉。本节覆盖构造只要求前一种零伦性，因此局部单连通比实际所需条件更强。
\end{solution}

\begin{exercise}\label{g1:ex:two-orbit-circle-monodromy}
	在五元集合 \(S=\{1,2,3,4,5\}\) 上，让 \(\pi_1(S^1)\cong\mathbb Z\) 的正生成元作用为 \((12)(345)\)。描述相应覆盖的全部连通分支、每个分支的叶数及整个覆盖的覆盖变换群。
\end{exercise}
\begin{solution}
	作用有两个轨道 \(\{1,2\}\)、\(\{3,4,5\}\)，所以覆盖有两个连通分支，分别为二叶、三叶圆周覆盖，同构模型是 \(z\mapsto z^2\) 与 \(z\mapsto z^3\) 的不交并。与单值交换的纤维置换必须把一个轨道等变地送到等基数的轨道，因此不能交换这两个分支。

	在一个长度为 \(m\) 的循环轨道上，与该循环交换的置换在各点的取值由一个点的像决定，且该像可以任取轨道中的点；这些置换恰为该循环的各次幂。因此两个轨道分别贡献 \(\mathbb Z/2\mathbb Z\)、\(\mathbb Z/3\mathbb Z\)，整个覆盖变换群为它们的直积，同构于 \(\mathbb Z/6\mathbb Z\)。这个六元群不在五点纤维上传递，因为它分别保持两个轨道。
\end{solution}

\begin{exercise}[自由子群的秩公式]\label{g1:ex:finite-index-free-rank}
	设 \(B\) 为有限连通图，有 \(v\) 个顶点、\(e\) 条边，\(r=e-v+1\)。设 \(\widehat B\rightarrow B\) 是连通 \(d\) 叶覆盖，其中 \(1\le d<\infty\)。证明
	\[
		\pi_1(\widehat B)\cong F_{\,1+d(r-1)}.
	\]
	据此证明秩 \(r\) 的自由群中每个指数 \(d\) 的子群都是秩 \(1+d(r-1)\) 的自由群。
\end{exercise}
\begin{solution}
	底空间的每个顶点有 \(d\) 个原像，共得到 \(dv\) 个顶点。把底空间的边按参数看作区间：从起点纤维中的每一点唯一提升该参数道路，得到 \(d\) 条边；沿反向道路的唯一性保证它们在内部互不相交并穷尽边内部的原像。顶点的小星形覆盖邻域说明这些提升边按原来的关联关系连接，因此覆盖空间本身是有 \(dv\) 个顶点、\(de\) 条边的有限图。由图基本群公式，其秩为
	\[
		de-dv+1=1+d(e-v)=1+d(r-1).
	\]
	特别地，取由 \(r\) 个圆周在一点相接的图，它的基本群为 \(F_r\)。有限图局部同胚于区间或有限星形，因而局部道路连通且半局部单连通。子群分类给每个指数 \(d\) 的子群一个连通 \(d\) 叶覆盖；覆盖的基本群单射到 \(F_r\)，像正是该子群。结合刚算出的秩，得到结论。这是一个从覆盖空间的边数、顶点数读出群生成元数目的例子。
\end{solution}

\begin{exercise}\label{g1:ex:torus-monodromy-relation}
	试图在三元纤维上让环面的两个标准生成元分别作用为 \((12)\)、\((23)\)。证明这样的覆盖不存在。若改成 \((123)\)、\((132)\)，求所得覆盖的连通性、叶数、对应的 \(\mathbb Z^2\) 子群及覆盖变换群。
\end{exercise}
\begin{solution}
	环面基本群的两个生成元交换，所以其右置换也必须交换。然而先施以 \((12)\) 再施以 \((23)\)，把 \(1\) 送到 \(3\)；反序则把 \(1\) 送到 \(2\)。两者不交换，故第一组数据不满足环面的关系，不能来自覆盖。

	第二组中，两置换互为逆，故交换。把纤维写成 \(\mathbb Z/3\mathbb Z\)，两生成元分别加 \(1\)、减 \(1\)。一般元素 \((m,n)\) 加 \(m-n\)，因此作用传递，覆盖连通且为三叶。基点稳定子群为
	\[
		H=\{(m,n)\in\mathbb Z^2:m-n\equiv0\pmod3\}
		=\langle(3,0),(1,1)\rangle.
	\]
	所有子群在 \(\mathbb Z^2\) 中都正规，故覆盖变换群为 \(\mathbb Z^2/H\cong\mathbb Z/3\mathbb Z\)。与二圈图相比，环面多出的交换关系实质地限制了允许的置换。
\end{solution}

\begin{exercise}\label{g1:ex:translation-subgroup-discrete}
	设 \(\Gamma\le(\mathbb R^n,+)\) 为加法子群，按平移作用于 \(\mathbb R^n\)。证明轨道映射 \(\mathbb R^n\rightarrow\mathbb R^n/\Gamma\) 是覆盖，当且仅当存在 \(\varepsilon>0\)，使 \(B(0,\varepsilon)\cap\Gamma=\{0\}\)。不要预先使用离散子群的基分类。
\end{exercise}
\begin{solution}
	若存在这样的 \(\varepsilon\)，取任一点 \(x\) 的开球 \(V=B(x,\varepsilon/3)\)。若 \(V+\gamma\) 与 \(V\) 相交，则 \(\gamma\) 是两个 \(V\) 中点的差，故 \(|\gamma|<2\varepsilon/3<\varepsilon\)，只能为零。于是作用满足本节的适当不连续条件，由\cref{g1:prop:action-quotient-cover}，轨道映射为覆盖。

	反之，若轨道映射为覆盖，取包含原点的一张覆盖片 \(V\)，其上投影单射。选 \(B(0,\varepsilon)\subseteq V\)。若有非零 \(\gamma\in B(0,\varepsilon)\cap\Gamma\)，则 \(0,\gamma\in V\) 是两个不同点，却属于同一轨道，违反单射性。因此该交只有零。这个条件正是子群在零点附近离散；通过平移，所有子群点也都孤立。
\end{solution}

\begin{exercise}\label{g1:ex:projective-double-cover}
	设 \(n\ge2\)。证明球面对径商 \(S^n\rightarrow\mathbb RP^n\) 是万有二叶覆盖，求 \(\pi_1(\mathbb RP^n)\)，并分类 \(\mathbb RP^n\) 的全部连通覆盖。说明 \(n=1\) 时为何不能沿用同一基本群结论。
\end{exercise}
\begin{solution}
	对径作用是自由的；在每个 \(u\in S^n\) 处取足够小的球面邻域，使它与其对径邻域不交，故由作用的商覆盖判据得到二叶覆盖。\cref{g1:thm:real-projective-quotient}说明这里的商拓扑正是此前定义的射影空间拓扑。由\cref{g1:prop:sphere-simply-connected}，\(n\ge2\) 时 \(S^n\) 单连通，所以此覆盖为万有覆盖。

	连通性与两点纤维给出覆盖变换群恰为 \(\{1,-1\}\)，即 \(\mathbb Z/2\mathbb Z\)。万有覆盖的覆盖变换群同构于底空间基本群，故 \(\pi_1(\mathbb RP^n)\cong\mathbb Z/2\mathbb Z\)。射影空间的仿射图保证局部道路连通和半局部单连通，因此子群分类适用；该二元群只有平凡子群和全群，故连通覆盖只有万有二叶覆盖与恒等覆盖两种同构类。

	当 \(n=1\) 时，\(S^1\) 的基本群为 \(\mathbb Z\)，并不单连通；对径二叶覆盖仍存在，却不再是万有覆盖。由第四章实射影直线与圆周的同胚，\(\pi_1(\mathbb RP^1)\cong\mathbb Z\)，不能从一个非万有覆盖的覆盖变换群直接读出底空间基本群。
\end{solution}

% !TeX root = ../../Book1.tex
% 纲要标题：## 第6章　曲面的拓扑分类
% 纲要位置：几何学计划纲要.md，第 315 行。
% 正文状态：初稿；尚未正式审核。

% 纲要内容开始
% 
% > 本章目标与证明范围：把三角剖分、标准模型的存在和分类不变量的唯一性作为三个独立证明阶段。主范围为紧拓扑曲面，带边界情形同时处理；不把高维流形三角剖分问题混入本章结论。
% 
% 纲要内容结束

\chapter{曲面的拓扑分类}
\label{g1:ch:06}
\BookChapterNumber{g1:ch:06}

\begin{chapterabstract}
	一张紧曲面可以由有限个三角形拼成，继而化为球面上添若干手柄或交叉帽、再挖去有限个圆盘的模型。本章分别证明这两个存在性结论，再用可定向性、Euler 示性数和边界分支数说明模型的唯一性。三角剖分从纯拓扑坐标出发，先建立平面 Jordan--Schoenflies 定理与简单弧的直化；多边形约化则通过具体的切割粘合完成。最后讨论曲线的邻域、切开后的曲面与裤分解，为后续几何结构提供拓扑模型。
\end{chapterabstract}

\begin{learninggoals}
	\begin{itemize}
		\item 分清内点、边界点和闭曲面的含义，用绕数定义局部定向，并说明定向如何沿坐标图相容。
		\item 从平面分离与延拓证明出发，整理有限圆盘覆盖，构造紧拓扑曲面的有限三角剖分。
		\item 用生成树把剖分化为边配对多边形，说明每一步切割、重粘与封盘为什么保留所需拓扑信息。
		\item 抽取手柄和交叉帽，处理两者混合的情形，得到带边界的全部标准模型。
		\item 通过定向覆盖、Euler 示性数和边界分支数证明分类唯一性，并以基本群展示作独立核验。
		\item 判断曲线的一侧或两侧、分离或非分离，计算切开后的类型并构造可定向负 Euler 示性数曲面的裤分解。
	\end{itemize}
\end{learninggoals}

上一章已经算过环面和若干多边形商空间的基本群。但列出几个模型，并不等于知道所有曲面都在名单之中。即使每一点附近都是圆盘，把局部圆盘接成整体时仍可能产生复杂的环路；而一个基本群计算正确，也尚未给出从原曲面到某个模型的同胚。本章要回答的正是：紧曲面的名单能否列全，名单中的参数又能否唯一确定？

\ProseParagraphBreak

证明分成三个阶段。首先把任意紧拓扑曲面化为有限的三角形拼接；其次用切割和粘合整理这些有限材料，得到标准模型；最后找出同胚必须保持的量，区分模型。第一阶段需要认真对待连续曲线，因为拓扑坐标并没有导数，两条曲线的交点也未必有限。第二阶段则要追踪边的方向和粘合参数，不能只对边字作代数变形。第三阶段使用的 Euler 示性数，要先证明不依赖所选剖分，才能承担分类不变量的作用。

\ProseParagraphBreak

本章的曲面均为 Hausdorff、第二可数的二维拓扑流形，允许有边界；分类定理讨论紧曲面，并先逐个处理连通分支。用到上一章的圆周基本群时，会重新指出绕数的作用；用到覆盖分类、有限图的基本群和附着圆盘公式时，也会标出相应结论。平面拓扑与三角剖分部分的证明依据见 Thomassen 的论文\cite[Sections 2--4]{Thomassen1992}。希望从另一种教材组织方式复习紧曲面，可参见 Lee 的著作\cite[Chapter 6]{Lee2011Topological}。

% !TeX root = ../../Book1.tex
% 纲要标题：### 6.1 曲面与局部组合结构
% 纲要位置：几何学计划纲要.md，第 319 行。
% 正文状态：初稿；尚未正式审核。

% 纲要内容开始
% 
% * 内点、边界点与连通分支
% * 局部圆盘的相容缩小
% * 定向的局部定义
% 
% 纲要内容结束

\section{曲面与局部组合结构}
\label{g1:sec:0601}

球面、环面和圆盘在大部分点附近看起来都像一小块平面，差别在于这些小块怎样接成整体。圆盘的边缘还需要另一种局部模型：只能向一侧延伸的半平面。本节先把这两种点的区别、局部邻域的控制和定向说清楚。这里没有切向量，也没有坐标变换的导数；判断方向将使用上一章已经建立的绕数。

\subsection{内点、边界点与连通分支}
\label{g1:subsec:0601:01}

\begin{definition}[拓扑曲面]\label{g1:def:topological-surface}
	设 \(M\) 为 Hausdorff、第二可数的拓扑空间。若每个点都有一个开邻域 \(U\)，同胚于闭半平面
	\[
		\mathbb H^2=\{(x,y)\in\mathbb R^2:y\ge0\}
	\]
	中的相对开集，则称 \(M\) 为允许带边界的\term[tuopuqumian]{拓扑曲面}[topological surface]。这样的同胚 \(\varphi:U\to\varphi(U)\subseteq\mathbb H^2\) 称为\term[zuobiaotu]{坐标图}[coordinate chart]。当 \(\varphi(p)\) 的第二坐标大于零时，称 \(p\) 为内点；等于零时，称 \(p\) 为曲面的边界点。所有边界点组成的集合记为 \(\partial M\)，内点集记为 \(M^\circ\)。
\end{definition}

定义暂时用某张图区分两类点。必须证明它不随换图改变，否则 \(\partial M\) 就不是空间本身的性质。平面圆盘去掉中心后有不能缩掉的环路，半圆盘去掉直边中心后却没有；这提供了一个不依赖坐标名称的判断。

\begin{proposition}[边界的内在性]\label{g1:prop:boundary-invariance}
	在允许带边界的拓扑曲面中，一个点不可能在一张坐标图中为内点、在另一张图中为边界点。任一曲面同胚 \(f:M\to N\) 因而满足 \(f(\partial M)=\partial N\)。内点集开，边界集闭；\(\partial M\) 在子空间拓扑下是一维无边界拓扑流形。
\end{proposition}
\begin{proof}
	假设 \(p\) 同时有两种图。先在内点图中取以 \(p\) 为中心的坐标开圆盘 \(B\)，使它包含于两图的共同定义域；再在边界图中取以 \(p\) 为直边中心的坐标开半圆盘 \(D\subset B\)；最后取内点图中更小的同心圆盘 \(B'\subset D\)。选 \(z\in B'\setminus\{p\}\) 为共同基点。包含映射给出
	\[
		\pi_1(B'\setminus\{p\},z)\longrightarrow
		\pi_1(D\setminus\{p\},z)\longrightarrow
		\pi_1(B\setminus\{p\},z).
	\]
	两端的群都是 \(\mathbb Z\)，而复合为同构：同心穿孔圆盘沿半径收缩到小圆周，包含映射把绕一周送到绕一周，所用计算是\cref{g1:thm:circle-fundamental}。中间的群却平凡，因为穿孔半圆盘沿半径形变收缩到一条闭半圆弧，后者同胚于闭区间而可缩。这与复合同态为同构矛盾。

	同胚将坐标图带到坐标图，故保持两种点。内点有完全由内点组成的小圆盘邻域，所以 \(M^\circ\) 开，\(\partial M\) 闭。在边界图中，刚证的内在性说明边界恰对应直边上的一个开区间；这些限制图给出 \(\partial M\) 的一维无边界流形结构。
\end{proof}

这里的 \(\partial M\) 是曲面的内在边界，不是任意嵌入下的子集边界。例如圆盘 \(D^2\) 的曲面边界是 \(S^1\)，球面 \(S^2\) 的曲面边界为空；但把 \(S^2\) 放在 \(\mathbb R^3\) 中，它作为子集的拓扑边界是整个 \(S^2\)。第四章的 \(\partial_X A\) 仍指后一种依赖环境的边界。紧且边界为空的曲面称为\term[biqumian]{闭曲面}[closed surface]；“闭”在这里同时包含紧性，不能只理解为某个环境中的闭子集。

\begin{proposition}[分支与边界圆周]\label{g1:prop:boundary-circles}
	允许带边界的拓扑曲面局部道路连通，其连通分支就是道路分支，并且开。紧曲面只有有限个连通分支，其边界是有限多个彼此不交的圆周的并。
\end{proposition}
\begin{proof}
	小圆盘与小半圆盘都道路连通，且可以任意缩小，故曲面局部道路连通。由\cref{g1:thm:locally-path-components}，两种分支相同并且开。若曲面紧，用这些开分支覆盖它，只能有有限多个分支。边界闭，因而在紧曲面中紧；上命题又说明它是一维无边界流形。

	下面补上所需的一维事实，不预设一维流形分类。设 \(C\) 为紧的一维无边界流形。在每一点选取一段嵌入的闭区间，使该点落在其中部；由紧性，有限条闭区间的内部覆盖 \(C\)。把这些区间的所有端点组成有限集 \(P\)。每条区间去掉 \(P\) 后只有有限个开小段。因此 \(C\setminus P\) 只有有限个连通分支：它由有限个这样的连通小段覆盖。

	任取一个分支 \(J\) 和其中一点，选择覆盖该点的一个区间内部 \(U\)。由于 \(U\) 的两个端点都已删去，\(U\setminus P\) 在 \(C\setminus P\) 中既开又闭，故 \(J\) 完全落在其中一段内。于是 \(J\) 本身就是一个开区间，它的两端在 \(P\) 中。这样 \(C\) 成为有限图；每个顶点的区间邻域恰有左右两个半边，所以顶点的度数为二。从任一边出发沿唯一的下一条边继续，有限性迫使回到起点，且不能遗下同一连通分支中的其他边。因此每个连通分支是一个有限循环图，同胚于 \(S^1\)。应用于 \(C=\partial M\) 即得结论；空边界对应零个圆周。
\end{proof}

这个结论允许我们先分类连通紧曲面，再取有限不交并。以后用 \(b\) 记边界圆周数；它是曲面同胚必须保留的非负整数。

\subsection{局部圆盘的相容缩小}
\label{g1:subsec:0601:02}

一张坐标图不能覆盖整个环面，但在已有坐标图里，可以预留一圈余地后再动手修改曲线。这种“较小集合的闭包仍在较大坐标图中”的控制，会防止局部切割碰到图的边缘。此时相容指拓扑包含和覆盖关系，不预设坐标变换可微。

\begin{lemma}[带余量的坐标圆盘]\label{g1:lem:nested-surface-charts}
	设 \(M\) 为允许带边界的拓扑曲面，\(p\in M\)，\(O\) 为 \(p\) 的开邻域。可以在同一坐标图中选取三个相似的闭圆盘或闭半圆盘 \(K_1,K_2,K_3\)，使
	\[
		p\in\operatorname{Int}_M K_1,\qquad
		K_1\subset\operatorname{Int}_M K_2,\qquad
		K_2\subset\operatorname{Int}_M K_3,\qquad K_3\subset O.
	\]
	半圆盘的直边位于 \(\partial M\)。若 \(M\) 紧，可以选有限组这样的集合，使最小各片的相对内部覆盖 \(M\)。
\end{lemma}
\begin{proof}
	先限制坐标图到 \(O\) 内。在内点处，将坐标中心平移到原点，选闭球完全落在坐标像内，再取三个严格递增的小半径。在边界点处，对半平面取以直边上的坐标点为中心的三个闭半球。其直边中部在 \(M\) 中仍是相对内部点，因为半球本来就只需要包含半平面内的邻域。

	这些闭圆盘或半圆盘的坐标原像紧；\(M\) 为 Hausdorff，所以它们在 \(M\) 中闭，闭包关系与坐标中的关系一致。这说明每一级都保留了所需余量。对所有 \(p\) 先作选择，最小各片的内部成为开覆盖，紧性再给出有限子覆盖。
\end{proof}

每一点都有紧邻域，所以曲面局部紧。由\cref{g1:cor:lch-second-countable}，Hausdorff、第二可数与局部紧性保证存在相容距离；本章在需要控制网格大小时固定其中任意一个。距离的数值不是新的曲面不变量，它只帮助把“充分小”写成统一的有限条件。特别是对紧曲面的有限坐标覆盖，\cref{g1:lem:lebesgue-number}提供正数 \(\delta\)，使直径小于 \(\delta\) 的集合完全落入某一坐标图。

\ProseParagraphBreak

还需注意：两张图中的圆看起来都很规则，它们在重叠处却可能无限次相交。坐标变换目前只是同胚，不能用“稍微移动就横截”代替证明。下一节将在平面中建立简单曲线和简单弧的具体控制，再把有限圆盘覆盖整理成有限曲线网。有限覆盖本身尚不等于有限三角剖分。

\subsection{定向的局部定义}
\label{g1:subsec:0601:03}

在欧氏平面中，“逆时针”有确定意义。换到另一张拓扑坐标图后，一个小圆可能变成非常不规则的简单闭曲线，但它围绕像点的绕数仍有意义。对不经过 \(a\in\mathbb R^2\) 的环路 \(\gamma\)，定义
\[
	w(\gamma,a)=w\left(t\longmapsto
	\frac{\gamma(t)-a}{\lVert\gamma(t)-a\rVert}\right).
\]
右边是\cref{g1:thm:circle-fundamental}的圆周绕数。起点不必为 \(1\)：把角度提升的终值减去初值即可，改变初始角度不会改变这个整数。移动基点的环路同伦仍保持总角度增量，亦可由\cref{g1:prop:moving-basepoint}和 \(\mathbb Z\) 的交换性得到。

\begin{proposition}[平面换图的局部方向符号]\label{g1:prop:orientation-local-sign}
	设 \(U,V\subseteq\mathbb R^2\) 为开集，\(f:U\to V\) 为同胚。对 \(x\in U\)，取以 \(x\) 为中心、闭包包含于 \(U\) 的小圆，按逆时针参数化为 \(\gamma_x\)。整数
	\[
		\varepsilon_f(x)=w(f\circ\gamma_x,f(x))
	\]
	与小圆半径无关，只能等于 \(1\) 或 \(-1\)，且关于 \(x\) 局部常值。对可复合的同胚，有
	\[
		\varepsilon_{g\circ f}(x)=\varepsilon_g(f(x))\,\varepsilon_f(x),
		\qquad \varepsilon_{\operatorname{id}}=1.
	\]
\end{proposition}
\begin{proof}
	同心小圆可沿半径在 \(U\setminus\{x\}\) 中变形到彼此，因此它们的像有相同绕数。为证明这个数的绝对值为一，取以 \(x\) 为中心的开圆盘 \(B\)，再取以 \(f(x)\) 为中心的圆盘 \(D\)，使 \(\overline D\subset f(B)\)，最后取以 \(x\) 为中心的圆盘 \(B'\)，使 \(f(B')\subset D\)。穿孔空间上的同态组成
	\[
		\pi_1(B'\setminus\{x\})\xrightarrow{\,f_*\,}
		\pi_1(D\setminus\{f(x)\})\longrightarrow
		\pi_1(f(B)\setminus\{f(x)\}).
	\]
	这里各群都是无限循环群：前两者由径向收缩得到，第三者由同胚 \(f\) 与 \(B\setminus\{x\}\) 得到。复合同态是同构，因为它就是 \(f\) 复合同心穿孔圆盘的包含同构。把两个同态写成整数乘法，其乘数的乘积为 \(1\) 或 \(-1\)，所以两个乘数都为 \(1\) 或 \(-1\)。第一个乘数正是所定义的绕数。基点可沿相应小区域中的道路对齐；循环群的内自同构平凡，不影响这些整数。

	证明局部常值时，固定一条围住 \(x\) 的坐标圆 \(C\)，其闭圆盘包含于 \(U\)。对圆盘内部任何 \(y\)，\(C\) 都代表绕 \(y\) 一周的类：以 \(y\) 为中心向一个充分小的圆作径向收缩，所有径向线段留在这个凸圆盘中。因此
	\[
		\varepsilon_f(y)=w(f\circ C,f(y)).
	\]
	当 \(y\) 沿圆盘内部的道路移动时，\(f(y)\) 始终避开 \(f(C)\)，上式的归一化方向连续变化，给出圆周环路的同伦。于是整数不变，证明局部常值。

	最后在 \(x,f(x),g(f(x))\) 处逐级取充分小的圆盘。穿孔小邻域中的基本群均由正向小圆指定生成元；\(f_*\) 和 \(g_*\) 分别是乘以相应符号，复合便乘以两者之积。恒等映射显然把正向圆保留为正向圆。
\end{proof}

若换图涉及边界，先在半圆盘的内点部分使用上述符号。这个部分连通，所以符号恒定；把相同符号赋给其直边点。缩小半圆盘不会改变符号，乘法公式也通过内点的公式延续到边界。这里不把半圆盘边界上的点当作平面内点来计算绕数。

\begin{definition}[局部定向与曲面的定向]\label{g1:def:surface-orientation}
	在固定点 \(p\)，把一张坐标图 \(\varphi\) 和一个符号 \(s_\varphi\in\{1,-1\}\) 配在一起；若另一对 \((\psi,s_\psi)\) 满足 \(s_\psi=\varepsilon_{\psi\circ\varphi^{-1}}(\varphi(p))s_\varphi\)，便把这两对视为等价。符号的乘法公式保证这是等价关系，每个点恰有两个等价类，称为该点的两个局部定向。它们记录同一方向在不同坐标中的写法。
	设 \(M\) 为允许带边界的拓扑曲面。在每个点的坐标图中标记一个符号 \(s_\varphi\in\{1,-1\}\)，其中符号在图内局部常值；若在任意两图交叠处都满足
	\[
		s_\psi(p)=\varepsilon_{\psi\circ\varphi^{-1}}(\varphi(p))\,s_\varphi(p),
	\]
	并将满足此关系的局部描述视为同一选择，则称这些相容选择为 \(M\) 的一个\term[dingxiang]{定向}[orientation]。存在定向的曲面称为\term[kedingxiang]{可定向}[orientable]，否则称为不可定向。直观上，符号 \(1\) 选定图中的逆时针方向，符号 \(-1\) 选定相反方向。
\end{definition}

只需在覆盖 \(M\) 的一组坐标图上给出相容符号：乘法公式保证它们唯一地确定所有其他图中的符号。这个定义不要求选择一个全局坐标系；它要求的是各处的局部方向能接在一起。对连通可定向曲面，定向恰有两个，因为两个定向在每点的符号比值为局部常值的 \(\pm1\)，在连通空间上只能恒定。

\begin{example}[平面方向与环面的定向]\label{g1:exm:torus-local-orientation}
	平移与旋转把正向小圆送到正向小圆，故方向符号为 \(1\)；反射 \((x,y)\mapsto(x,-y)\) 的符号为 \(-1\)。在环面模型 \(\mathbb R^2/\mathbb Z^2\) 中，充分小的投影圆盘作为坐标图，重叠处的坐标变换为整数平移。因此所有这些图均选符号 \(1\) 就得到环面的定向。将全部符号同时取反，得到另一个定向。
\end{example}

对平面中带边界的区域，正定向规定沿边界行走时区域留在左侧。若取半平面 \(y\ge0\) 的正定向，直边上的正向就是 \(x\) 增大的方向。定向相容的换图保留这个选择：把一段直边与内侧一条简单弧合成小半圆盘的边界，边界的行走方向由内点的正向小环路决定，换图符号为正时这个方向不变。下一节的简单弧直化会给出这个边界描述的圆盘坐标。反转曲面定向时，各边界圆周的诱导方向也同时反转。

% !TeX root = ../../Book1.tex
% 纲要标题：### 6.2 紧曲面三角剖分的证明
% 纲要位置：几何学计划纲要.md，第 325 行。
% 正文状态：初稿；尚未正式审核。

% 纲要内容开始
% 
% * 平面 Jordan–Schoenflies 定理与简单弧直化
% 
%   在局部平面图中先建立 Jordan–Schoenflies 的平面 Jordan 曲线版本，以及圆盘内适当嵌入简单弧的直化与延拓引理，再处理有限曲线网。三角剖分的主范围为紧、Hausdorff、第二可数的拓扑二维流形，边界允许存在；平面拓扑引理为本章明确的证明任务，不预设曲面已有光滑图册，也不调用高维三角剖分结论。
% 
% * 交叠边界的局部有限交叉
% * 由曲线网得到有限胞腔分解
% * 圆盘胞腔的三角化
% * 同痕后的公共细分与 Euler 示性数不变性
% 
% 纲要内容结束

\section{紧曲面三角剖分的证明}
\label{g1:sec:0602}

形成有限三角剖分，需要把坐标图中的曲线整理成有限的边与顶点。简单闭曲线可能处处没有切线，圆盘边界也可能无限相交；平面中的分离与直化结论使我们能够逐个调整边界，同时保留圆盘覆盖。平面部分采用有限图与逐次加密的构造；这条证明思路见\cite[第2--4节]{Thomassen1992}。

\subsection{平面 Jordan–Schoenflies 定理与简单弧直化}
\label{g1:subsec:0602:01}

这里的简单弧是 \([0,1]\) 的嵌入像，\term[Jordanquxian]{Jordan 曲线}[Jordan curve]是 \(S^1\) 的嵌入像。\footnote{Camille Jordan（1838--1922），法国数学家，研究群论与拓扑。}嵌入的参数映射可以非常不规则；整个论证只用连续性和单射性。先给出后面所需的平面分离与有限图结论。

\begin{lemma}[有限折线与平面图]\label{g1:lem:polygonal-plane-tools}
	下列事实成立。
	\begin{enumerate}
		\item 平面连通开集中的任意两点可以用一条简单折线连接。
		\item 简单闭折线把平面分成两个连通开集，其中一个有界；二者的边界都是这条折线。闭折线内部的一条简单折线弧若恰在两个端点与边界相交，就把内部切成两个区域，其边界分别由该弧与原边界的两段组成。
		\item 任意有限图在平面中的嵌入都可以重新画成边为折线的嵌入。完全二部图 \(K_{3,3}\) 没有平面嵌入；把一个平面图中的互不相交的有限树各收缩成一点，仍得到平面图。
		\item 若一个有限平面图连通、没有环边、没有割点，且至少有三个顶点（若原图有环边，先将环边细分，再检验这些条件），则它可以从任一圈开始，逐次加入两端在旧图上、内部与旧图不交的简单弧得到。若各边为折线，它的每个面以一个简单圈为边界。
	\end{enumerate}
\end{lemma}

\begin{proof}
	对于第一项，固定起点，把能用有限折线到达的点组成集合。小圆盘内可以添一条线段，所以这个集合及其补集在给定开集中都开；连通性说明它等于整个开集。有限条线段的交点和重叠端点只有有限个，把它们作为顶点得到有限图，删去回路便得到简单折线。

	对于第二项，先给出折线两侧的具体含义。在每条边的内部取窄矩形，在每个顶点取足够小的圆盘，使不相邻的这些小邻域不相交；顶点圆盘内只有两条相接的半径方向，可以沿这两条边的两侧分别绕过顶点。沿整个多边形走一周，得到两条互不相交的窄通路。平面补集中的任一点可以沿折线走向边界，取第一次接触以前的一小段，就到达这两条通路之一；若接触顶点，改在顶点圆盘中的相应扇形内行走。因此补集至多有两个分支。

	再证明两侧不能相通。从不在边界上的点向无穷远发一条射线，先让射线避开顶点且不平行于边，记录它与边界的交点数模 \(2\)。改变射线方向经过一个顶点时，两次交点同时出现或消失，或者一次交点换到相邻边，故奇偶性不变。沿不碰边界的短线段移动起点也不改变奇偶性：在扫过的窄四边形中，每段边界从边界进入还必须离开，交点变化为偶数。有限折线连接把这一结论推广到同一分支中的任意两点。跨过一条边的内部时奇偶性恰好改变，因此确有两个分支。足够大的圆外连通，故只有一个分支无界；窄通路同时说明边界上的每一点都在两分支的闭包中。对于内部弧，把它与边界的两段分别组成简单闭折线，重复奇偶性论证；弧的两侧不同，而每一侧都沿对应边界窄通路连通，恰好得到所说的两个区域。

	第三项的重新画法不要求原边可微。先给环边添一个顶点，使每条边有不同的端点。在每个顶点周围取互不相交的小圆盘，只截到与该顶点相接的边。对每条边取最后离开起点圆盘以后、随后第一次进入终点圆盘以前的闭子弧；这些子弧两两不交，因而有两两不交的开邻域。用第一项在各邻域内改画简单折线，并在顶点圆盘内用互不相交的径向线段接到顶点。\(K_{3,3}\) 包含一个六边圈和连接三对相对顶点的三条弦。在折线画法中，至少两条弦落在该圈的同一侧；它们的四个端点沿圈交错，与第二项的切分结论矛盾。有限折线树的细邻域可以由顶点小圆盘和边的细矩形组成；从叶边开始依次添加矩形和端点圆盘，每次仍是一个闭盘。在该闭盘内把树及其附着边改画成一颗星，就实现树的收缩。这也说明若一个图经过收缩互不相交的树能得到 \(K_{3,3}\)，原图便不能平面嵌入。

	最后证明第四项。已有一个没有割点的真子图 \(H\)。若剩余某条边的两端都在 \(H\) 中，可直接加入。否则，取 \(G\setminus H\) 的一个连通分支；它必须接到 \(H\) 的至少两个不同顶点，因为只有一个接点会成为 \(G\) 的割点。在该分支中选连接两个接点的最短路径，便得到所需的新弧。加入它仍没有割点，有限次便遍历全部边。从初始圈出发，每加入一条弧，它都位于某个旧面的内部；第二项说明它把该面分成两个以简单圈为边界的面。于是归纳得到最后的结论。
\end{proof}

下面两步把折线结论推广到任意连续弧和闭曲线。关键区别是：任意简单弧仍然不能隔断平面，而闭曲线确实会隔断。

\begin{lemma}[简单弧不分离平面]\label{g1:lem:arc-nonseparation}
	若 \(A\subset\mathbb R^2\) 是简单弧，则 \(\mathbb R^2\setminus A\) 道路连通。
\end{lemma}

\begin{proof}
	固定 \(x,y\notin A\)，取 \(d>0\)，使它们到 \(A\) 的距离都大于 \(4d\)。按参数把 \(A\) 分成有限段 \(A_1,\ldots,A_k\)，每段直径小于 \(d\)。不相邻的闭段互不相交，所以其间距离的最小值为正；取 \(0<\delta<d\) 小于这些距离。若只有一两段，直接任取这样的 \(\delta\)。在每段上取更细的点列，使相邻点间的子弧直径小于 \(\delta/8\)，并在每个点处画边长 \(\delta/2\)、边平行于坐标轴的正方形。记第 \(i\) 段所用正方形边界的并为 \(G_i\)，把有限交点添作顶点。

	相邻正方形的边界至少有两个公共点；由\cref{g1:lem:polygonal-plane-tools}最后一项，或直接删去任一顶点检查，\(G_i\) 没有割点。相邻的 \(G_i,G_{i+1}\) 包含同一个端点正方形；不相邻的图不交。每一对相邻图都包含在以公共弧端点为中心、半径 \(3d\) 的圆盘内，所以 \(x,y\) 位于这对图的无界面。

	需要一个有限图的观察：这种链状排列的图若把一点围在有界面内，必有一对相邻图已经把它围住。证明如下。由前一引理，整个并没有割点，各有界面的边界都是圈。若观察不成立，取围住该点的圈 \(C\)，使它所用图的最大下标 \(j\) 最小，再使 \(C\) 使用 \(G_j\setminus G_{j-1}\) 的边数最少。若 \(C\) 只在 \(G_{j-1}\cup G_j\) 内，已经矛盾；否则 \(C\) 在 \(G_{j-1}\) 上至少有两段。\(G_{j-1}\) 没有割点，故可取一条内部避开 \(C\) 的路径连接其中两段：在 \(G_{j-1}\) 中取连接两段的最短路径，并截去首次接触之前和末次接触之后的部分即可。它把 \(C\) 切成两个圈；点不在新路径上，必在其中一个圈的内部。该圈或者降低最大下标，或者减少所数的边，均与选择矛盾。

	因此 \(x,y\) 都位于 \(G=\bigcup_iG_i\) 的无界面。这个面是连通开集，可用简单折线连接两点。每一段 \(A_i\) 都被相应正方形覆盖，而每个正方形内部与 \(G\) 的无界面不交；故所得折线避开 \(A\)。
\end{proof}

\begin{theorem}[Jordan 曲线分离定理]\label{g1:thm:jordan-separation}
	任意 Jordan 曲线 \(C\subset\mathbb R^2\) 的补集恰有两个分支：有界分支 \(D\) 和无界分支 \(E\)，且 \(\partial D=\partial E=C\)。两分支都是道路连通的开集。
\end{theorem}

\begin{proof}
	先证至少有两个分支。取两条竖直支撑线，使 \(C\) 位于它们之间，在每条支撑线上取 \(C\) 的最高点 \(p_1,p_2\)。曲线分成从 \(p_1\) 到 \(p_2\) 的两段。在两支撑线之间再取一条竖线，这两段与该线都相交；沿该线找到两个相邻的、分别属于两段的交点 \(a,b\)，使开线段 \((a,b)\) 避开 \(C\)。这是可行的，因为两组交点是互不相交的紧集。另从 \(p_1\) 向上、经曲线上方的水平线、再降到 \(p_2\)，画一条内部避开 \(C\) 的折线 \(L\)。若线段 \((a,b)\) 的中点可以在补集中到达无穷远，它就可以接到 \(L\) 的内部；截取首次接触，所得连接弧的内部同时避开 \(C\)、\([a,b]\) 和 \(L\)。于是 \(C\)、\([a,b]\)、\(L\) 及这条连接弧组成一个 \(K_{3,3}\) 的细分：两类顶点分别是 \(p_1,p_2\) 与连接弧在 \([a,b]\) 上的端点，以及 \(a,b\) 与连接弧在 \(L\) 上的端点。这与\cref{g1:lem:polygonal-plane-tools}矛盾。因此补集不连通。

	对补集的任意分支 \(U\)，考察从 \(U\) 内出发、恰在终点接触 \(C\) 的折线。这些折线所能到达的端点在 \(C\) 中稠密。事实上，任取一个开子弧 \(J\subset C\)，\(C\setminus J\) 是简单闭弧；由\cref{g1:lem:arc-nonseparation}，可用避开该闭弧的折线连接 \(U\) 中的一点和另一个分支中的一点。它第一次离开 \(U\) 时只能碰到 \(J\)，这就给出一个所需的可达点。子弧可以任意缩小，稠密性随之成立。

	假设补集有三个分支 \(U_1,U_2,U_3\)，在曲线上选三段互不相交的短闭子弧 \(J_1,J_2,J_3\)。每个 \(U_i\) 中可以画一棵有限的折线树，分别接到三个 \(J_j\)，且除这些端点外避开 \(C\)：先用上一段选可达点，再在连通开集 \(U_i\) 内连接三条到达折线，删除多余回路。三棵树互不相交。把各树及每个 \(J_j\) 上连接三个接触点的最短子弧各收缩为一点，得到 \(K_{3,3}\)，仍与有限图引理矛盾。所以恰有两个分支。

	足够大的圆外是连通的，因此恰有一个无界分支。补集是开集，平面的连通开集道路连通。边界只能落在 \(C\) 上，而上述可达点在 \(C\) 上的稠密性说明 \(C\) 上每一点都属于每个分支的边界。
\end{proof}

\begin{lemma}[横切弧与有限弦网]\label{g1:lem:jordan-crosscuts}
	设 \(C\) 是 Jordan 曲线，\(D\) 是其有界补分支。若简单折线弧 \(P\) 的内部在 \(D\) 中，且两个端点是 \(C\) 上的不同点，那么 \(P\) 把 \(D\) 切成两个区域，边界分别由 \(P\) 与 \(C\) 的两段组成。设 \(G\) 为由 \(C\) 和 \(D\) 中有限条折线弧组成的连通有限图，将环边先作细分，并假设细分后的图没有割点。则 \(G\) 的每个有界面都以 Jordan 曲线为边界；在某个面中添加一条上述弧，恰好把该面分成两个。
\end{lemma}

\begin{proof}
	记 \(C\) 的两段为 \(A_1,A_2\)，\(C_i=P\cup A_i\) 的有界内部为 \(D_i\)。原曲线的无界外部 \(E\) 是避开 \(C_i\) 的连通无界集，故在 \(C_i\) 的外部。\(A_2\) 的内点属于 \(\overline E\) 又不在 \(C_1\) 上，因此也在 \(C_1\) 的外部；同理 \(A_1\) 的内点在 \(C_2\) 外部。于是 \(D_1\) 不碰 \(C\)，且不在 \(E\) 中，故 \(D_1\subset D\)，同理 \(D_2\subset D\)。

	\(D_1\) 连通且不碰 \(C_2\)，所以整体位于 \(C_2\) 的一侧。若 \(D_1\subset D_2\)，闭包包含关系将使 \(A_1\) 的内点落在 \(\overline{D_2}\)，与上一段矛盾。因此 \(D_1\cap D_2=\varnothing\)。在 \(P\) 的任一内点处取足够小的圆盘，使 \(P\) 在其中是一条折线直径；若恰为折点，先用扇形坐标把两段边变直。在这个小盘中，\(C_1,C_2\) 都只有同一段 \(P\)。Jordan 定理说明每个 \(D_i\) 都接近该点；两侧半盘连通，因此其中一侧属于 \(D_i\)。两内部互不相交，说明它们恰好占据相反的两侧。

	所以 \(D_1\cup P^\circ\cup D_2\) 在 \(D\) 中是开集。它也是 \(D\) 中的闭集，因为它等于 \((\overline{D_1}\cup\overline{D_2})\cap D\)。\(D\) 连通且这个集合非空，故它等于 \(D\)。这就证明 \(D\setminus P=D_1\sqcup D_2\)，并且各自边界正是 \(P\cup A_i\)。有限图部分再按\cref{g1:lem:polygonal-plane-tools}的逐次加弧构造归纳；这里每次只调用刚证明的横切弧结论。
\end{proof}

把 \(D\cup C\) 实现为闭圆盘，需要将边界上的对应延拓到内部。Jordan--Schoenflies 定理的构造同时细化原像和像中的网格，使两侧相邻点的控制在极限中都保留下来，从而得到连续且互逆的映射；\cref{g1:fig:double-mesh}给出构造的示意。\footnote{Arthur Sch\"onflies（1853--1928），德国数学家，研究拓扑与晶体群。}

\input{Figures/Chapter06-DoubleMesh}

\begin{theorem}[Jordan--Schoenflies 定理]\label{g1:thm:jordan-schoenflies}
	任意两条平面 Jordan 曲线之间的同胚都可以延拓为整个平面的同胚。特别地，每条 Jordan 曲线及其有界补分支的并同胚于闭圆盘，而且给定的边界同胚可以一并延拓。
\end{theorem}

\begin{proof}
	先把给定边界同胚写成 \(f:C\to\partial Q\)，其中 \(Q\) 是正方形；若这个情形成立，对两条曲线分别构造到正方形的延拓再复合即可。记 \(D\) 为 \(C\) 的内部。我们构造一列有限弦网 \(G_n\subset\overline D\)、\(G'_n\subset Q\)，始终包含边界，内部边为折线。每个图没有割点，其删去边界后的部分连通；两图具有相同的顶点、边、面关系，边界对应始终是 \(f\)。在这个阶段只固定顶点的对应，不固定内部边上每一点的参数。

	从 \(D\) 可沿折线到达的边界点中选可数稠密集 \(A\)，这是\cref{g1:thm:jordan-separation}证明中得到的。取两个这样的点，用 \(D\) 中的简单折线连接，得到初始弦；在正方形内取对应弦。于是得到 \(G_0,G'_0\)。以后依次把 \(A\) 的各点加入顶点集：从新点沿到达折线走进 \(D\)，再连到旧图内部，截到第一次碰到旧图为止。它在某个旧面内成为一条新弦。在对应的正方形面内也可画出相应弦；该面内部连通，端点位于折线边或已有可达的边界点，先取端点处的小段，再用连通开集中的简单折线相接即可。\cref{g1:lem:jordan-crosscuts}保证两面以相同方式各分为两面。

	同时必须加密网格。令
	\[
		K_n=\{x\in\overline D:\operatorname{dist}(x,C)\ge 1/n\}.
	\]
	在 \(D\) 内加入有限个小正方形的边及其中的横、竖网格，使 \(K_n\) 被这些正方形的内部覆盖，且网格的每格直径小于 \(1/n\)。所有这些边与旧图的内部折线只有有限个交点，并与 \(C\) 不交。必要时添加有限条内部折线，使新图没有割点且内部连通。这个补充可以如下完成：先连接各分支；若有割点 \(v\)，在 \(v\) 两侧的分支中各取靠近 \(v\) 的点，在 \(D\setminus\{v\}\) 中用折线连接。连通开集删去一点仍连通，因为一条经过该点的折线可以在一个小圆盘内绕过它。把新交点作为顶点后，这条连接把割点处分开的块合并，使块数至少减少一；按有限图的块数归纳即可结束。每次新增弧都可拆成有限次在旧面内加弦，因而可以在 \(Q\) 中逐条复制其组合关系。

	接着反过来，在 \(Q\) 的内部加入足够细的有限网格，使到边界距离至少 \(1/n\) 的部分被直径小于 \(1/n\) 的网格包围；再把这些弦逐条复制回 \(D\)。新端点若落在已有边的内部，就在对应边上选一个新点，保持已有各点的次序；若在边界上，则选取 \(f(A)\) 中的点。可在边界附近略弯折网格末端以做到这一点，内部网格不受影响。仍用前一段的加弦过程补足连通性。后加的弦只细分已有面，不会破坏先前的网格。因此，对于原像与像两侧各自的每个内部紧集，以及每个 \(\varepsilon>0\)，从某一步起，紧集上每一点都位于一个由图中边组成的简单小圈内部，而这个小圈连同内部的直径小于 \(\varepsilon\)。若点恰在网格边或顶点上，就取环绕它的两行、三行小格的外框；把网格取得更细即可达到同一要求。

	把所有顶点及整个边界组成集合 \(F\)，相应集合记为 \(F'\)。构造给出双射 \(f_0:F\to F'\)，且两集合分别在 \(\overline D,Q\) 中稠密。它保持有限弦网中每个圈的内外关系：这在每次加弦时由横切弧引理保持。现在证明上述对应延拓为同胚。

	取 \(x\in D\)。在某个图中选围住 \(x\) 的内部圈 \(L\)，其对应圈 \(L'\) 也在 \(Q\) 内部。趋近 \(x\) 的 \(F\) 中点最终在 \(L\) 内，像点因而留在 \(L'\) 围出的紧区域中。任取其像点的两个聚点 \(y,z\)。若 \(y\ne z\)，利用像侧越来越细的网格，在 \(y\) 附近取两个互不相交的嵌套网格圈，它们都把 \(y\) 与 \(z\) 分开。对应的两个原像圈也互不相交，至多一个经过 \(x\)。取另一个圈 \(M\)；趋近 \(x\) 的点最终都在 \(M\) 的同一侧，像点也都在对应圈的同一侧，不可能同时趋近位于相反侧且不在圈上的 \(y,z\)。故像点只有一个极限，定义它为 \(h(x)\)。同样的论证对反向顶点对应成立，得到 \(k:Q^\circ\to D\)。这里聚点不会落在边界：环绕任一内部点的内部小圈及其对应圈已经把聚点限制在内部的紧区域。常值顶点序列也是允许的趋近序列，所以延拓在每个已有顶点上与 \(f_0\) 一致。

	这些延拓连续。给定 \(h(x)\) 的一个小邻域，取像侧两个互不相交的嵌套内部小圈，都把 \(h(x)\) 围在其中，并令闭内部落在该邻域；至少一个对应原像圈不经过 \(x\)，选取这一个。它必须围住 \(x\)，否则稠密顶点的对应违反内外关系。这个原像圈的内部便是 \(x\) 的一个邻域，整个邻域的像都落在给定小邻域中。同理 \(k\) 连续。两个复合映射在稠密顶点集上都是恒等映射，连续性给出 \(k\circ h=\mathrm{id}\)、\(h\circ k=\mathrm{id}\)。

	最后检查边界。设 \(x_m\in D\) 趋近 \(a\in C\)，而某个像子列趋近 \(b\in Q\)。若 \(b\in Q^\circ\)，内部圈及其原像紧区域迫使 \(x_m\) 留在 \(D\) 的一个紧子集，矛盾。故 \(b\in\partial Q\)。若 \(b\ne f(a)\)，在 \(C\) 上选两个 \(A\) 中的点 \(a_1,a_2\)，把 \(a\) 与 \(f^{-1}(b)\) 分开。它们最终都已加入图；因图的内部连通，可以在图内找到一条从 \(a_1\) 到 \(a_2\)、其余部分均在内部的简单弦。该弦及其对应弦把两边各分成两块。靠近 \(a\) 的点位于与 \(a\) 相邻的块，其像位于与 \(f(a)\) 相邻的块，后者的闭包不含 \(b\)，矛盾。因此 \(b=f(a)\)。紧性说明整个像序列都趋近 \(f(a)\)，从而 \(h\cup f\) 连续。它是紧空间到 Hausdorff 空间的连续双射，故为闭盘同胚。

	还要延拓到外部。先选各曲线内部的一点作为原点，对两平面作反演 \(x\mapsto x/\|x\|^2\)。每条曲线的外部连同无穷远点变成一条 Jordan 曲线的闭内部，原来的无穷远点对应原点。已经证明的闭盘延拓给出反演后两内部之间的同胚。可使它把原点映到原点：先把目标闭盘映到正方形，在正方形中从给定内部点与目标内部点分别向四角作扇形三角剖分，逐三角形的仿射映射给出固定边界、移动该点的同胚，再共轭回去。于是反演回去得到外部同胚，与内部同胚在原边界上相合。两部分的粘合连续，其逆也同样连续，故得到全平面同胚。
\end{proof}

\begin{lemma}[简单弧的直化与延拓]\label{g1:lem:proper-arc-straightening}
	平面简单弧的每个内点都有一个开矩形型邻域，使该邻域与原弧的交集在某个坐标下恰为水平直径。若 \(B\) 是闭 Jordan 盘，\(A\subset B\) 是简单弧，且 \(A\cap\partial B\) 恰为其两个端点，则 \(B\setminus A\) 有两个分支，闭包都是闭 Jordan 盘。给定把这两个端点送到另一闭盘中一条直径端点的边界同胚，可以延拓为把 \(A\) 送到该直径的闭盘同胚。
\end{lemma}

\begin{proof}
	先证局部结论。在所给内点 \(p\) 两侧各选一个短子弧。每个短子弧内都有从原弧补集沿线段可达的点：取靠近该子弧、但不在原弧上的点，再在原弧上选取离它最近的点，连接这两个点的开线段不碰原弧；起点取得足够近，最近点仍在指定短子弧内。补集由\cref{g1:lem:arc-nonseparation}连通，所以可用其中的折线连接两条到达线段，删去回路，得到一条与原弧只在所选两端点相交的简单弧。它与经过 \(p\) 的原子弧组成 Jordan 曲线。\cref{g1:thm:jordan-schoenflies}把这条曲线变成圆；再在圆上的对应点附近取普通矩形坐标，并缩小矩形以避开原弧剩下的紧集，就得到所求直化。

	对闭盘中的弧，将它分别与盘边界的两段组成 Jordan 曲线。重复\cref{g1:lem:jordan-crosscuts}的两个有界内部的论证：两内部不交且包含在 \(B\) 中；本引理刚证的局部直化使 \(A\) 每个内点附近都有两个半矩形，两个内部各占一侧；它们与 \(A^\circ\) 的并在 \(B^\circ\) 中同时开闭。因此它们恰为 \(B\setminus A\) 的两块，闭包只沿 \(A\) 相交。任取从 \(A\) 到目标直径、与给定两个端点相容的同胚。它与给定的边界同胚在两个子盘的边界上分别组成同胚；对两子盘应用\cref{g1:thm:jordan-schoenflies}并沿 \(A\) 粘合，便得到整个闭盘的同胚。
\end{proof}

这里有一个以后会反复用到的直接推论：位于一个平面图中的闭 Jordan 盘，其边界有双侧环带邻域。用全平面同胚把边界变成单位圆，再取 \(1-\varepsilon<r<1+\varepsilon\) 的逆像即可。有限条在一点相交、除此之外互不相交的弧也可以同时直化为有限条径向线段：先把两条组成的弧直化，再在各半盘内依次用上一引理直化其余弧；每次缩小中心邻域以避开各弧已经截去的紧尾部。因弧数有限，此过程结束后得到一个顶点小圆盘和互不相交的边矩形。因此，有限嵌入图具有由顶点盘、边带组成的细邻域。

\subsection{交叠边界的局部有限交叉}
\label{g1:subsec:0602:02}

现在回到曲面。紧性允许选有限个坐标圆盘，但不能保证它们的边界相交有限次；我们要改变边界，仍保留原来较小圆盘的覆盖性质。下面的论证只在一个坐标图内进行，因此同时适用于曲面上的有限条曲线。

\begin{lemma}[边界的有限化]\label{g1:lem:disk-boundary-regularization}
	设曲面 \(M\) 中有有限条简单曲线组成的嵌入有限图 \(G\)，以及一个坐标图内的三个闭 Jordan 盘
	\[
		B_0\subset\operatorname{int}B_1\subset B_1\subset\operatorname{int}B_2.
	\]
	存在闭 Jordan 盘 \(B\)，使 \(B_0\subset\operatorname{int}B\subset B\subset\operatorname{int}B_2\)，且 \(\partial B\) 与 \(G\) 只有有限个交点。可以进一步使交点均不为 \(G\) 的顶点，并且在局部矩形中两条边分别为横、竖直径。
\end{lemma}

\begin{proof}
	先把 \(B_2\) 当作平面中的圆盘。在 \(G\cap\operatorname{int}B_2\) 中，称两端在 \(\partial B_2\)、内部在盘内的一段边为穿入段。虽然穿入段可能无限多，碰到 \(B_1\) 的穿入段只有有限多：每条这样的段都从 \(\partial B_2\) 走到与它正距离的 \(B_1\)，而一条边的参数映射一致连续，所以它占据的参数区间长度有统一正下界；这些参数区间互不相交，因此只能有有限多个。先在 \(G\) 的有限个顶点处用小盘截断，也把这些顶点附近的有限段一并列入。

	把这些有限段与 \(\partial B_2\) 合在一起，必要时在盘内添有限条弧，使所得有限图没有割点。\cref{g1:lem:proper-arc-straightening}及\cref{g1:thm:jordan-schoenflies}允许固定 \(\partial B_2\)，把这个有限图改画成折线图：按逐次加弧的次序，第一次延拓边界同胚，以后在当前弧所在的两个子盘内分别延拓，再粘合。每一步保持已经画好的边不动。

	在这个坐标变化以后，未列入的穿入段仍不碰新的 \(B_1\)。沿 \(\partial B_1\) 取有限个小正方形，闭正方形都避开 \(B_0\) 及这些未列入段，并使它们覆盖 \(\partial B_1\)。为何能同时避开全部未列入段？在 \(\partial B_1\) 的每一点，若它们有点列趋近该点，则原有限边中的一个子列有参数极限；该极限落在 \(B_1\) 上，所处穿入段也应列入，且一个小参数邻域属于这同一穿入段，矛盾。必要时先把 \(B_1\) 略缩到仍包含 \(B_0\) 的内层盘，便可使闭方形与所有未列入段分离。

	这些正方形连同它们围住的 \(B_0\) 构成有限折线区域；取围住 \(B_0\) 的外边界，并在有限个自接触顶点处用很小的折线绕开，得到 \(B_0\) 与 \(\partial B_2\) 之间的一条简单闭折线。它只可能碰到已列入的有限折线段，所以交点有限。最后在每个交点附近作小的折线改动，避开顶点、重叠和相切：沿每段取不平行于所遇旧边的小线段，并把端点避开有限条直线即可。改动仍在所选小方形内，故既不破坏包含关系，也不会碰到未列入的段。把新边界映回原坐标便得结论。
\end{proof}

局部横交的意思只是两条弧在一个小矩形中呈十字交叉，并没有涉及导数。对于紧无边界曲面，先取有限个坐标盘 \(B_{0,i}\) 的内部覆盖曲面，并为每个盘预留两层较大的坐标盘。逐个应用\cref{g1:lem:disk-boundary-regularization}，已选边界组成旧图，便得到仍覆盖曲面的有限个闭盘 \(B_i\)，其边界两两有限横交，且没有三条边界交于一点。避开有限个已有交点就能满足最后一项。

\subsection{由曲线网得到有限胞腔分解}
\label{g1:subsec:0602:03}

上述圆盘边界形成有限嵌入图 \(G\)。在每条没有交点的圆周上任添一个顶点，就可以把它也看成有限条边。\(M\setminus G\) 的每个分支都位于某个坐标盘的内部：选该分支中的一点和包含它的盘；在不越过该盘边界的情况下，连通分支不可能从盘内走到盘外。这里使用的是盘边界在本地坐标中已经证明的分离性质。

一个分支的边界不一定已经是一个简单圈：同一顶点或同一边可能在边界行走中出现多次，分支也可能有若干个洞。沿图取细邻域，可以把重复的边界行走拆开，使余下区域具有互不相交的边界圆周。

\begin{lemma}[有限曲线网的圆盘化]\label{g1:lem:finite-network-disk-cells}
	设紧无边界曲面 \(M\) 上的有限嵌入图 \(G\) 的每个补分支都包含在某个平面坐标图内。则可以添加有限条简单弧，并对所得曲线网作有限细分，使 \(M\) 被分成有限个闭 Jordan 盘；不同盘只沿整条边或顶点相交。
\end{lemma}

\begin{proof}
	按\cref{g1:lem:proper-arc-straightening}后的构造，为每个顶点取互不相交的小盘，为每条截短的边取细矩形，各矩形只在两个端部接到相应顶点盘。边数和顶点数有限，紧的不相交部分可以用互不相交的邻域隔开，因此各次缩小能够同时完成。所得闭邻域 \(N\) 的边界由有限段弧拼成有限条互不相交的圆周。顶点盘与边矩形本身已经是有限个盘形块。

	\(M\setminus\operatorname{int}N\) 的每个分支 \(R\) 是一个紧平面区域，其边界是上述圆周中的有限条。这样的分支只有有限多个：每个分支都接到 \(\partial N\) 的至少一条圆周，而一条圆周的外侧局部属于唯一分支。若有分支没有边界，它在 \(M\) 中同时开闭；在所考虑的连通分支内这会迫使 \(G\) 为空，可先添一个坐标小圆周排除该情形。每个 \(R\) 的边界恰由一个外圈和有限个内圈组成。由 Jordan 分离定理，任一边界圈的一侧属于 \(R\)，另一侧不属于；在最外圈内反复排除内圈的闭内部，就得到 \(R\)，不存在额外遗漏区域，否则其边界还应包含另一条 \(\partial R\) 圈。

	若 \(R\) 有内圈，利用其内部的折线连通性，从一个内圈接一条简单弧到外圈，只在两端碰边界。取该弧的窄矩形带并切开，相当于把这两个边界圈合并成一个圈，内圈数减少一。弧与边界相接处使用半矩形坐标，故切开后仍是带边界的紧曲面；在平面中也可直接把矩形带从区域移去，它的两条长边取代原弧。有限次以后只剩一个边界圈，由\cref{g1:thm:jordan-schoenflies}所得区域是闭盘。把先前取出的矩形带作为单独的闭盘形块加回来，便得到有限盘形分解。最后将各端点和交点统一添为顶点，使两盘共有的边都成为完整边。
\end{proof}

\subsection{圆盘胞腔的三角化}
\label{g1:subsec:0602:04}

盘形分解已将曲面分成有限个闭圆盘，各块沿完整的边或顶点相交。在每个圆盘中添中心并连至边界分点，可以把它分成三角形；公共边上使用同一组分点，使这些三角形能够拼成全曲面的剖分。

\begin{definition}[有限三角剖分]\label{g1:def:surface-triangulation}
	曲面 \(M\) 的一个\term[sanjiaopoufen]{有限三角剖分}[finite triangulation]是有限单纯复形 \(K\) 及同胚 \(|K|\to M\)。这里 \(K\) 由有限个顶点、边和三角形组成，每个单形的面仍在 \(K\) 中，两个单形的交集是各自的一个公共面或空集；\(|K|\) 表示按这些面粘合后的拓扑空间。若 \(M\) 有边界，还要求 \(\partial M\) 是一些边和顶点的并。
\end{definition}

\begin{theorem}[紧拓扑曲面的有限三角剖分]\label{g1:thm:surface-triangulation}
	每个紧、Hausdorff、第二可数的二维拓扑流形都有有限三角剖分；允许存在边界。对于给定的、与边界相容的有限嵌入图，可以预先指定有限个点，并使该图和这些点成为剖分的子复形；这里与边界相容指图与边界的交集是它的若干整条边及顶点。尤其可以保留有限条两两不相交的内嵌简单闭曲线和适当嵌入简单弧。
\end{theorem}

\begin{proof}
	先考虑无边界情形。有限圆盘覆盖及其边界正则化给出\cref{g1:lem:finite-network-disk-cells}所需曲线网，因而得到有限闭盘分解。若预先指定有限条曲线与弧，就从这些曲线组成的图开始添加圆盘边界；在交点或预先指定点处分边，后续步骤保持它们作为图的一部分。给每个盘的边界添足至少三个顶点，选一个内部点。把该盘同胚到正多边形，在内部点与边界各顶点之间连线，再映回来，就把该盘分成有限个三角形。盘与盘之间的边界划分事先统一，所以这些三角形可以沿相同的边粘合。

	这一步可能只得到一种容许重边的三角形粘合，例如两条不同的边具有相同端点。作一次重心细分便可消除此问题：给每个原顶点、原边和原盘分别赋予不同的新顶点，每条严格包含链“顶点属于边，边属于盘”对应一个小三角形。为避免一个盘边界中同一边重复导致重复链，使用上一引理的闭 Jordan 盘分解，其边界已经单射，且共有部分均为完整的边或顶点。于是一个包含链决定唯一小三角形，交集恰为共同的子链，真正得到有限单纯复形。同胚在每个闭盘上连续并相容，有限闭集粘合给出全局连续双射；紧性和 Hausdorff 性使它成为同胚。

	若 \(M\) 有边界，先把两个副本沿边界逐点粘合，得到双倍曲面 \(\widehat M\)。\cref{g1:prop:boundary-invariance}保证这里的边界是内在定义的；边界点的半平面图把两个半圆盘粘成整圆盘，因此 \(\widehat M\) 是无边界曲面。它紧且第二可数；Hausdorff 性也可直接在图中验证：不同内点用各自副本的邻域分离，两个不同边界点先在 \(M\) 中取闭包不交的小半圆盘，再把对应的两半邻域粘合。由\cref{g1:prop:boundary-circles}，\(\partial M\) 是有限个圆周的不交并，在双倍中成为有限条简单闭曲线。

	对 \(\widehat M\) 应用无边界情形，并把这些边界圆周以及预先给定弧的两个副本加入初始曲线网，指定点也保留。所得剖分以 \(\partial M\) 为子复形；每个小三角形内部不越过这条子复形，所以选取原副本一侧的三角形及其面，就得到 \(M\) 的有限三角剖分。若 \(M\) 不连通，局部连通性使各分支开，紧性说明分支只有有限个，分别构造再取不交并即可。
\end{proof}

紧性在这里提供有限覆盖、有限边界分支与最终的有限剖分。第二可数性是本书对曲面的统一假设；在紧 Hausdorff、局部 Euclid 的情形中，它也可由有限图册推出：每个坐标图有可数基，把有限个图的基合并即可。

\subsection{同痕后的公共细分与 Euler 示性数不变性}
\label{g1:subsec:0602:05}

两个任意拓扑剖分可能有无限交叉：一条边可以在另一条边附近来回摆动无穷多次，而两边仍各自是简单弧。先用环境同痕调整一份剖分，使两套边网只有有限横交，再逐个细分这些交点，就能寻找共同的三角剖分。

\begin{definition}[环境同痕]\label{g1:def:ambient-isotopy}
	空间 \(M\) 上的\term[huanjingtonghen]{环境同痕}[ambient isotopy]是一族同胚 \(H_t:M\to M\)，其中 \(0\le t\le1\)，\(H_0=\mathrm{id}_M\)，并且 \((x,t)\mapsto(H_t(x),t)\) 是 \(M\times[0,1]\) 的同胚。它连续移动整个空间，所以同时控制被移动图形与其周围的部分。对本节的紧 Hausdorff 曲面，只要 \((x,t)\mapsto H_t(x)\) 连续且每个 \(H_t\) 是同胚，上述乘积映射便是紧空间到 Hausdorff 空间的连续双射，逆映射自动连续。
\end{definition}

\begin{lemma}[有限图的小环境同痕]\label{g1:lem:graph-small-isotopy}
	设 \(K\) 是紧曲面 \(M\) 的一个有限三角剖分，\(G\subset M\) 是另一个与边界相容的有限嵌入图，即 \(G\cap\partial M\) 是 \(G\) 的若干整条边与顶点。存在从恒等映射开始的环境同痕 \(H_t:M\to M\)，使 \(H_1(G)\) 的每条边在 \(K\) 的三角形坐标中由有限条线段组成。可以把同痕限制在 \(G\) 的任意给定邻域内；有边界时保持 \(\partial M\) 作为集合不变，并可保持预先共有的边界顶点不动。
\end{lemma}

\begin{proof}
	取 \(G\) 的互不相交顶点小盘及边的窄矩形，使用\cref{g1:lem:proper-arc-straightening}后的有限图邻域构造，并使这些邻域都落在指定邻域中。先在每个顶点盘内把相交的有限条初始边改成径向星形，再把截短边作为各矩形的中线处理。由于 \(K\) 有限，在 \(|K|\) 中充分细分以后，可以在每个盘、矩形内画出一颗折线星或一条简单折线中线，与原来的端点及循环次序一致：沿原弧取有限条小三角形链，连接相邻小邻域中的点，删除回路；不同边使用互不相交的窄矩形，故不会产生新交叉。顶点附近则在一张三角形扇形的有限平面图中按给定循环次序选择短线段，再用小折线连接到边矩形。

	旧星形与新星形之间的对应可以延拓为固定顶点盘边界的同胚。逐条应用闭盘适当弧的延拓引理，把一个扇区分成两个，再在各剩余子盘内延拓即可；只有有限个扇区。各边矩形中，固定矩形边界，将旧中线送到新折线的同胚同样由\cref{g1:lem:proper-arc-straightening}给出。按“先顶点盘、后截短矩形”的次序复合这些同胚，它们在各自支撑边界上恒等，故与外部的恒等映射粘合。

	每个固定闭盘边界的同胚都同痕于恒等映射。把盘识别为单位圆盘，若该同胚为 \(u\)，定义
	\[
		U_t(x)=
		\begin{cases}
			t\,u(x/t),&\|x\|\le t,\quad t>0,\\
			x,&\|x\|\ge t,
		\end{cases}
		\qquad U_0(x)=x.
	\]
	在 \(\|x\|=t\) 两式相等；每个 \(U_t\) 的逆由同一公式将 \(u\) 换成 \(u^{-1}\) 得到，且 \(\|U_t(x)-x\|\le2t\) 保证在 \(t=0\) 连续。这就把有限次盘内修改逐一提升为环境同痕。边界半盘的修改保持其直径上的点或指定端点，反射成整盘后使用相同构造，再限制回半盘即可。
\end{proof}

\begin{theorem}[同痕后的公共细分]\label{g1:thm:common-subdivision-isotopy}
	设 \(K,L\) 是紧拓扑曲面 \(M\) 的两个有限三角剖分。存在环境同痕 \(H_t\)，使 \(K\) 与 \(H_1(L)\) 有一个有限公共细分。有边界时，可以保持边界为集合。
\end{theorem}

\begin{proof}
	对 \(L\) 的边网应用\cref{g1:lem:graph-small-isotopy}，并在有限个小盘内再作小的折线调整，使它避开 \(K\) 的顶点、没有重叠边段，且与 \(K\) 的边只作有限次横交；边界上允许共有边段，把所有端点统一细分即可。两个有限折线网的并是有限嵌入图。每个补分支都在 \(K\) 的一个三角形及 \(H_1(L)\) 的一个三角形之内，但可能有洞或重复的边界行走；按\cref{g1:lem:finite-network-disk-cells}的顶点盘、边带及连接内圈的构造，把它进一步切成闭盘，再在每个盘内三角化。取出的顶点盘和边带也分别按两个原剖分的边切开，因而每个新小三角形都包含在两个剖分各自的某个单形内。统一共有边上的分点，最后作重心细分，即得有限公共细分。
\end{proof}

\begin{proposition}[Euler 示性数与剖分无关]\label{g1:prop:euler-triangulation-independent}
	若紧曲面 \(M\) 的有限三角剖分有 \(V\) 个顶点、\(E\) 条边、\(F\) 个三角形，则整数
	\[
	\chi(M)=V-E+F
	\]
	不依赖所选三角剖分。更一般地，任何有限闭盘分解的顶点数减边数加盘数也等于 \(\chi(M)\)。这个数称为\term[Eulershu]{Euler 示性数}[Euler characteristic]。\footnote{Leonhard Euler（1707--1783），瑞士数学家，在分析与几何等领域作出广泛贡献。}
\end{proposition}

\begin{proof}
	先记住三种局部计数。给一条边加一个顶点时，\(V,E\) 各增加 \(1\)；用一条新弧把一个盘分成两个盘时，\(E,F\) 各增加 \(1\)；在一个盘内部加一个顶点并向其边界的 \(r\) 个顶点连线时，\(V,E,F\) 分别增加 \(1,r,r-1\)。三种变化都保持 \(V-E+F\)。

	任意有限细分也保持这个数。先在各旧边上加入所有新顶点，再处理每个旧盘内部的有限连通边网。把它连同盘边界画在平面中，选择一个生成树；树有 \(V'-1\) 条边，且补集只有一个面，因为可从叶边开始，用小圆盘绕过叶边，并逐条删去叶边。每增加一条非树边，它与旧图中连接两端的一条简单路径组成 Jordan 曲线。这条新边的两侧因该曲线而不能再相通；它内部有矩形邻域，所以只把所在的一个面分成两个，其他面不变。因此有界面数为 \(E'-(V'-1)\)，即这个闭盘分解满足 \(V'-E'+F'=1\)。边界上若有 \(b\) 个顶点，就恰有 \(b\) 条边，它们在交替和中抵消，故盘内新增的顶点、边与面取代原来一个盘时仍贡献 \(1\)。对各旧盘求和，便证明细分前后计数相同，也同时证明盘形分解与其三角化的计数相同。

	给两个三角剖分应用\cref{g1:thm:common-subdivision-isotopy}。环境同痕只移动单形，不改变它们的数目；同痕后的两个剖分具有公共细分，而细分保持计数，故两者的 \(V-E+F\) 相等。
\end{proof}

Euler 示性数的 \(V-E+F\) 计数以盘形面为单位。环面上有一条圆周，其补集是一个环带，这个环带就不是一个可以按“一个面”计数的闭盘。后文使用多边形或带分解计算 Euler 示性数时，会先把各块细分为本节允许的盘形块。

% !TeX root = ../../Book1.tex
% 纲要标题：### 6.3 由剖分到边识别多边形
% 纲要位置：几何学计划纲要.md，第 336 行。
% 正文状态：初稿；尚未正式审核。

% 纲要内容开始
% 
% * 生成树切割和单个多边形
% * 边配对及边界记录
% * 切割粘合操作的合法性
% 
% 纲要内容结束

\section{由剖分到边识别多边形}
\label{g1:sec:0603}

三角剖分把曲面化成有限资料，但三角形的个数通常与曲面的复杂程度相差很远。一个球面可以有成千上万个三角形，却不应需要成千上万个分类参数。本节先沿一棵对偶生成树把三角形展开成单个多边形，再把原来的顶点树吸收到一张圆盘中。留下的有限带条才记录曲面上不能靠改变分割消去的粘合。

\subsection{生成树切割和单个多边形}
\label{g1:subsec:0603:01}

固定紧连通曲面 \(M\) 的有限三角剖分，其存在性由\cref{g1:thm:surface-triangulation}保证。必要时作一次细分，以便在每个顶点周围看到由三角形组成的扇：内点的扇环绕一整圈，边界点的扇只绕半圈。原剖分的顶点与边组成一个有限图，称为其一维骨架。另放一个顶点在每个三角形内部；对每条非边界边，把它两侧的三角形中心经边的中点连接，得到\term[duioutu]{对偶图}[dual graph]。这些连线实际画在三角形中，只有在指定的边中点才穿过原骨架。

\begin{lemma}[对偶连通性]\label{g1:lem:surface-dual-connected}
	紧连通曲面的任一有限三角剖分，其对偶图连通。若曲面无边界，且 \(T\) 为一维骨架的一棵生成树，则删除与 \(T\) 的边相交的对偶边以后，对偶图仍连通。
\end{lemma}
\begin{proof}
	把经非边界边相邻的三角形归入同一类。同一内边两侧当然同类。在一个顶点周围，三角形按圆周或区间顺次排列，相邻者共享一条经过该顶点的非边界边，所以全部同类。于是各类三角形的并不但闭，而且在顶点、边内点和面内点处都有包含于该并的小邻域，因而也开。曲面连通，只有一类，证明第一项。

	设曲面无边界，删除指定对偶边后却出现两个以上分支。取其中一个分支的全部三角形为集合 \(\mathcal A\)。将一侧在 \(\mathcal A\)、另一侧不在 \(\mathcal A\) 的原边收集成子图 \(L\)。原对偶图连通，故 \(L\) 非空；而未被删除的对偶边不能连接两个新分支，故 \(L\subseteq T\)。

	在每个原顶点周围沿一圈看相邻三角形，进入 \(\mathcal A\) 与离开 \(\mathcal A\) 的次数相同。因此 \(L\) 在每个顶点的度数都是偶数。另一方面，树的任意非空边子图是森林，其中一个含边分支的最长简单边路有度数为一的端点。这与所有度数为偶数矛盾，故删边后的对偶图仍连通。
\end{proof}

\begin{theorem}[单个多边形表示]\label{g1:thm:surface-polygon}
	每个紧连通拓扑曲面都同胚于一张闭多边形按部分边的成对识别所得的商。闭曲面的每条边都参与且只参与一次配对；带边界时，未配对边的像组成曲面的边界。每一对边的识别均可在选定边参数后取为线性识别。
\end{theorem}
\begin{proof}
	选择对偶图的一棵生成树 \(T^*\)。先把各三角形当作彼此独立的副本，只沿 \(T^*\) 所穿过的原边粘合，不作其余任何识别。根三角形是一张圆盘。按树离根的距离逐个加入其他三角形，每次新三角形只沿一条完整边接到已有圆盘的一条边上；顶点和其他边仍保留各自副本。这样的粘合仍是圆盘：两张三角形或多边形沿边拼成的区域可把共同边直化，再在两侧使用半圆盘坐标；亦可把抽象的边界顶点依序放在一个凸多边形上，对角线连接后逐三角形定义分片线性同胚。归纳得到一张闭多边形 \(P\)。

	未穿过 \(T^*\) 的每条内边现在在 \(\partial P\) 上出现两次，按原边参数重新配对；原边界边只出现一次，不配对。顶点识别由这些边识别自动产生，因为同一原顶点周围的三角形扇是连通的。于是原三角形到 \(M\) 的映射下降为 \(P/{\sim}\to M\) 的连续双射。源空间紧，目标 Hausdorff，故由\cref{g1:thm:compact-bijection}为同胚。面内点不参与剩余识别，配对边内点有两张半圆盘，未配对边内点只有一张半圆盘；在顶点处，原有的圆形或区间形扇被恢复。因此未配对边的像恰是 \(\partial M\)。
\end{proof}

这一步只用了对偶树。一张单面多边形的各个顶点仍可能代表曲面中的不同点。下一步同时使用原骨架的一棵树，目的是把顶点附近的树状部分并成一张圆盘，而不是把一条任意闭路压成一点。

\begin{proposition}[一个顶点圆盘与有限带条]\label{g1:prop:one-vertex-band-model}
	设 \(M\) 为紧连通闭拓扑曲面。存在一张闭圆盘 \(D\)、有限条矩形带条和一张面圆盘 \(F\)，使 \(M\) 可按如下方式构成：每条带条的两个短边分别粘到 \(\partial D\) 上互不相交的小区间，长边保持为边界；\(D\) 与全部带条组成的曲面 \(R\) 恰有一个边界圆周，再沿此圆周粘上 \(F\)。若没有带条，则 \(M\cong S^2\)。
\end{proposition}
\begin{proof}
	在一维骨架中选生成树 \(T\)。由上引理，可在不穿过 \(T\) 的对偶图中选生成树 \(T^*\)，故两棵嵌入的树不相交。先围绕各原顶点取互不相交的小多边形圆盘，沿每条原边取窄矩形，在每个三角形中央留下一个小圆盘。这些选择可在有限三角形的重心细分中同时完成：顶点小片、边的窄条和面中央小片只沿各自边界相接，覆盖整张曲面。

	把属于 \(T\) 的边条加入各顶点小圆盘。每加一条树边，它连接的是此前不同的两个树分支，所以只是用一个矩形把两张圆盘接起来；所得仍是圆盘。按树归纳，所有顶点片及树边条合成一张圆盘 \(D\)。同样，把 \(T^*\) 所穿边处的横向窄条与各面中央圆盘合并，得到另一张圆盘 \(F\)。这两种合并不冲突，因为 \(T^*\) 没有穿过 \(T\)。

	尚余的区域恰对应既不在 \(T\) 中、也不被 \(T^*\) 穿过的原边。沿每条这样的边保留矩形带条，两个短边都接到 \(D\) 上，两个长边都接到面圆盘 \(F\) 的边界上。有限个多边形角可以在各原三角形内重新参数化成直角，不改变粘合。故 \(D\) 加全部带条的补部就是 \(\operatorname{Int}F\)，其边界为 \(\partial F\) 这一条圆周，得到所需描述。

	无带条时是两张圆盘沿整个边界粘合。任一边界粘合同胚 \(h:S^1\to S^1\) 可径向延拓为 \(ru\mapsto r h(u)\) 的圆盘同胚，故可把粘合改成恒等粘合；两盘分别放在上下半球便得到 \(S^2\)。
\end{proof}

在这个构造中，树并未被直接商成一个点。我们把树的一个盘状邻域当作新的顶点圆盘，所有改变都在已有曲面内进行。这也解释了为何带条的扭转仍须记录：树可以吸收，带条的两端以何种方向接到圆盘却影响整张曲面的定向与拓扑。

\subsection{边配对及边界记录}
\label{g1:subsec:0603:02}

给每对待识别边规定同一名称和箭头，沿多边形边界选定一个行走方向。顺着箭头读 \(a\)，逆着箭头读 \(a^{-1}\)，便得到循环的\term[bianzi]{边字}[boundary word]。闭曲面中每个字母恰出现两次，符号可能相同，也可能相反。未配对边另作标记，不与任何其他边识别。边字记录的是边的参数配对；顶点的等价类也必须由箭头逐项追踪，不能任意规定所有角为同一点。

\begin{example}[两种四边形]\label{g1:ex:polygon-vertex-counts}
	边字 \(aba^{-1}b^{-1}\) 的四个角属于同一等价类，给出环面。边字 \(aa^{-1}bb^{-1}\) 却有三个顶点等价类：依次记四角为 \(v_0,v_1,v_2,v_3\)，\(a\) 配对只令 \(v_0=v_2\)，\(b\) 配对也只令 \(v_2=v_0\)，而 \(v_1,v_3\) 各自保留。这时一维骨架是一棵有两条边的树，顶点盘加带模型没有剩余带条，所以得到球面。
\end{example}

两字在相应群关系中的表现很不同，但这里先按边识别构造空间，不能先对字作群的自由约化，再宣称所得商空间不变。允许的空间操作将在下个小节说明。

\begin{proposition}[边界领圈]\label{g1:prop:boundary-collar}
	设 \(M\) 为紧拓扑曲面。它的边界有一个\term[lingquan]{领圈}[collar]：存在嵌入
	\[
		c:\partial M\times[0,1]\longrightarrow M,\qquad c(x,0)=x,
	\]
	使 \(c(\partial M\times[0,1))\) 是 \(\partial M\) 的开邻域。不同边界分支的领圈可以取为互不相交。
\end{proposition}
\begin{proof}
	取含边界子复形的有限三角剖分。沿一条边界边的内部，在邻接三角形中取一条窄矩形；在边界顶点处，邻接三角形形成半圆盘状的有限扇。在这个扇中以一条靠近顶点的折线连接相邻两条边上的截点，折线与边界上的两小段围成一个圆盘。把它参数化成矩形，使下边对应原边界小段，左右两边接到上述边条，顶边对应内侧折线。这样的参数化存在：给四段边界指定对应，利用平面圆盘的边界延拓；也可在此有限扇中逐个三角形线性构造。

	于是每条边界圆周旁得到有限个矩形首尾相接，每个矩形的下边都在原边界上。相邻矩形在横截区间上的粘合固定下端和上端，所以其区间参数是递增同胚。依次改变各矩形的参数，把除最后一次之外的粘合都改为恒等；最后一条缝上的粘合是某个递增同胚 \(h:[0,1]\to[0,1]\)。它固定两端，且
	\[
		h_s(t)=(1-s)t+s h(t)\qquad(0\le s\le1)
	\]
	对每个 \(s\) 仍为递增同胚。从这个首尾粘合的矩形到标准环带的映射
	\[
		[s,t]\longmapsto(e^{2\pi i s},h_s(t))
	\]
	在 \((1,t)\sim(0,h(t))\) 处相容，是连续双射，因紧到 Hausdorff 而为同胚。这表明沿边界没有横向反转的可能：反转会把属于原边界的下端送到内侧上端。

	所得闭环带的外边是原边界，去掉内侧边后是相对开邻域；各块中的矩形坐标也直接核验这一点。边界只有有限个分支，且相互为不交紧集，先缩小所选扇片和边条即可使全部领圈不交。
\end{proof}

\begin{proposition}[封住边界]\label{g1:prop:cap-boundary}
	设 \(M\) 为有 \(b\) 个边界圆周的紧连通拓扑曲面。沿每个边界圆周用任意同胚粘入一张闭圆盘，所得空间 \(\widehat M\) 是紧连通闭曲面，且 \(M\) 同胚于 \(\widehat M\) 去掉这 \(b\) 张圆盘的内部。所得同胚类型不依赖封口同胚的选择。
\end{proposition}
\begin{proof}
	用上命题的领圈，在每条粘合缝附近使用圆周参数与法向参数：原曲面一侧的法向参数非负，所添圆盘薄环带一侧的参数非正。拼合以后成为圆周乘一个开区间，所以缝上各点都有平面圆盘邻域。新圆盘中心和原曲面内点都有原来的圆盘图，因而没有边界点。

	空间由有限个紧空间粘合而紧，并且连通。它是 Hausdorff：不同缝上的点可取不交薄领圈；同一缝的不同点取不交圆弧乘薄区间；不在缝上的点可在各自原片内取避开缝的邻域，或先把薄领圈缩小后分开。这些邻域也是商拓扑中的开集。有限组来自原曲面和添入圆盘的可数坐标基给出第二可数性，因此所得确为闭拓扑曲面。

	若改变某个封口同胚，两种选择之差为圆周同胚。把它径向延拓到所添圆盘，原曲面上保持恒等，即得到两个商空间之间的同胚。这也保持原曲面的副本，故不仅抽象空间同胚，原曲面及封口位置也相容。
\end{proof}

所以有边界情形可以先封盘化为闭曲面，再追踪取下了多少张盘。但“同样数目的盘可以移到同样的位置”仍需证明，下一节恢复边界时会补上这个位置无关性，而不把它藏在图形直觉中。

\subsection{切割粘合操作的合法性}
\label{g1:subsec:0603:03}

边字只是表示方法，下面的变化应当被理解为同一块材料重新切开和接合。若没有把边界参数与剩余识别一并带走，即使变换后的群展示同构，也不能据此断言两张曲面同胚。

\begin{lemma}[圆盘的延拓与相对边界同痕]\label{g1:lem:disk-extension}
	任意圆周同胚 \(h:S^1\to S^1\) 都延拓为圆盘同胚。若圆盘同胚 \(f:D^2\to D^2\) 在边界上恒等，则存在从恒等映射到 \(f\) 的连续同胚族，且族中每个同胚都在边界上恒等。
\end{lemma}
\begin{proof}
	第一项用 \(H(ru)=r h(u)\)，其中 \(u\in S^1,0<r\le1\)，并令 \(H(0)=0\)。它与以 \(h^{-1}\) 定义的同类公式互逆；在原点处由像点的模等于 \(r\) 得到连续性。

	第二项对 \(0<t\le1\) 定义
	\[
		F_t(x)=
		\begin{cases}
			t f(x/t),&\lVert x\rVert\le t,\\
			x,&t\le\lVert x\rVert\le1,
		\end{cases}
		\qquad F_0(x)=x.
	\]
	在 \(\lVert x\rVert=t\) 处两式一致，因为 \(f\) 固定单位圆。每个 \(F_t\) 的逆由相同公式中的 \(f^{-1}\) 给出，故为同胚。\(t>0\) 时分片粘合保证联合连续；当 \(t\to0\) 时，所有移动只发生在半径 \(t\) 的圆盘内，且 \(\lVert F_t(x)-x\rVert\le2t\)，故在 \(t=0\) 也联合连续。最后 \(F_1=f\)，并始终固定单位圆。
\end{proof}

上述连续同胚族也称为\term[tonghen]{同痕}[isotopy]；把它在圆盘外接上恒等映射，就得到\cref{g1:def:ambient-isotopy}的环境同痕。它比一般同伦严格：每一时刻都必须可逆。本章用它调整切割位置时，曲面同胚类型始终保留。

\begin{proposition}[多边形表示的基本变化]\label{g1:prop:polygon-moves}
	对已表示拓扑曲面的边配对多边形，下列操作保持所得曲面的同胚类型：
	\begin{enumerate}
		\item 同时重命名一对边、反向选择该边的箭头，或改变边界字的起读位置；也可反向走完整个多边形边界，此时次序和所有指数同时反转。
		\item 在一对边上加入相应的分点，或删去这种成对分点。
		\item 沿多边形内一条连接不同边界点的适当嵌入简单弧切成两张圆盘，给两份切缝加入相应识别；随后可选择其他一对完整边将两盘重新合为一盘，所有其余配对按原参数保留。
	\end{enumerate}
	最后一种操作中的两张盘允许先各自换用圆盘坐标或翻面，但必须同时带走边界箭头及其参数。
\end{proposition}
\begin{proof}
	前两项只改变同一商映射的参数名称或分段，等价关系未变。第三项由\cref{g1:lem:proper-arc-straightening}将切弧化为圆盘直径，因此切开后确是两盘，按原参数重新接上切缝便逐点找回原盘。若改先粘另一对边，抽象的两盘沿一段完整边拼合仍是圆盘；其边界各段继承原识别。无论先作哪次粘合，最终对两盘不交并所加的等价关系完全相同，所以两种分步商都连续双射到同一个原曲面。紧到 Hausdorff 的判据将这些双射提升为同胚。

	圆盘坐标的改变由上引理将边界参数延拓到内部；翻面也是圆盘同胚，只是改变边界行走方向。因此只要同步改变箭头及边界次序，最后的等价关系仍相同。若把箭头保留而单独翻面，则可能改变识别，这不属于所述合法操作。
\end{proof}

一个常用的特例是消去边字中的相邻 \(aa^{-1}\)。这两边在商图中形成一条从其余图伸出的末端边；围住这条边及其末端顶点的小邻域是圆盘，把它并入相邻面就去掉了这个“树枝”。若还有其他边，得到余下的边字；若整字只有 \(aa^{-1}\)，则商空间由两个盘组成，是球面，不能把空字误读为空空间。同号的 \(aa\) 不属于这种消去，它记录了不同的粘合。

\ProseParagraphBreak

下一节将反复使用这些圆盘与带条操作抽出手柄和交叉帽。证明的终止性由带条数或边对数的严格减少保证；Euler 数稍后用于区分最终模型，不能提前拿来替代存在性中的切割论证。

% !TeX root = ../../Book1.tex
% 纲要标题：### 6.4 标准模型的约化
% 纲要位置：几何学计划纲要.md，第 342 行。
% 正文状态：初稿；尚未正式审核。

% 纲要内容开始
% 
% * 手柄和交叉帽的抽取
% * 可定向标准字
% * 不可定向标准字
% * 连通和与边界分支的恢复
% 
% 纲要内容结束

\section{标准模型的约化}
\label{g1:sec:0604}

上一节已把紧曲面的局部拼接信息压缩到有限条边上，但边数少并不意味着字已经容易辨认。例如，同一条边的两个字母相隔很远，可能隐藏着一个交叉帽；两个方向相反的边对互相穿插，则可能隐藏着一个手柄。本节用实际的切割和粘合把这些部件逐次取出。我们先处理闭曲面，最后把封住的边界圆盘恢复。这里要得到的是标准模型的存在；不同参数是否可能给出同胚曲面，由下一节另行判别。

\subsection{手柄和交叉帽的抽取}
\label{g1:subsec:0604:01}

先把多边形换成一种更容易局部操作的表示。\cref{g1:prop:one-vertex-band-model}把闭连通曲面写成一个顶点圆盘 \(D\)、有限条矩形带和一个面圆盘 \(F\) 的并：每条带的两条短边粘到 \(\partial D\) 上互不相交的小闭区间，且
\[
	B=D\cup\{\text{全部矩形带}\},\qquad \partial B\cong S^1,
\]
面圆盘 \(F\) 恰好封住这条边界。带的长边和顶点盘边界上没有被占用的弧组成 \(\partial B\)。这比抽象边字多保留了一点宽度，因而可以直接看见切下的对象是不是曲面、其边界有几条。

给 \(D\) 选一个定向。沿一条带把这个定向从一个粘合区间传到另一个区间，若与 \(D\) 原来的定向一致，称这条带无扭；若相反，称它有扭。此处说的扭转只记录两端的粘合方式，不需要把带嵌入三维空间再数半圈。顶点盘加一条有扭带是\term[mobiusdai]{Möbius 带}[Möbius band]\footnote{August Ferdinand Möbius（1790--1868），德国数学家、天文学家，研究射影几何、重心坐标及单侧曲面。}，可写成
\[
	\mathcal M=([0,1]\times[-1,1])/((0,t)\sim(1,-t)).
\]
顶点盘加一条无扭带则是环带。若两条带的四个粘合区间沿 \(\partial D\) 依次为 \(a,b,a,b\)，称这两个边对交错；这只说明端点在顶点盘周围的循环次序。

\begin{definition}[连通和与标准模型]\label{g1:def:surface-standard-models}
	从两个连通闭曲面 \(M,N\) 各去掉一个嵌入闭圆盘的内部，再沿两条剩下的边界圆周作同胚粘合，所得曲面称为它们的\term[liantonghe]{连通和}[connected sum]，记为 \(M\#N\)。本节先对所选圆盘及粘合实施这个构造，随后说明标准模型中这些选择为何不影响同胚类型。记
	\[
		\Sigma_{0,0}=S^2,\qquad \Sigma_{g,0}=\underbrace{T^2\#\cdots\#T^2}_{g\text{ 个}}\ (g\ge1),\qquad N_{k,0}=\underbrace{\mathbb RP^2\#\cdots\#\mathbb RP^2}_{k\text{ 个}}\ (k\ge1).
	\]
	再从相应闭模型去掉 \(b\) 个互不相交嵌入闭圆盘的内部，所得模型记为 \(\Sigma_{g,b}\) 和 \(N_{k,b}\)，其中 \(b\ge0\)。从一个闭曲面接上单孔环面，相应部件称为\term[shoubing]{手柄}[handle]；接上 Möbius 带，相应部件称为\term[jiaochamao]{交叉帽}[crosscap]。后者也就是与 \(\mathbb RP^2\) 作连通和，因为 \(\mathbb RP^2\) 去掉一个面内圆盘后正是 Möbius 带。
\end{definition}

连通和的两条边界各有环带邻域，把它们粘合后，接缝处的两个半圆盘组成一个圆盘，所以构造确实给出无边界曲面。\cref{g1:exm:surface-presentations}已分别从方形和圆盘对径识别得到环面与实射影平面；取其顶点及边的小带状邻域，就得到这里的单孔环面与 Möbius 带模型。下面把这个识别和余下部分同时说清楚。

\begin{lemma}[带条的抽取]\label{g1:lem:band-extraction}
	设闭连通曲面 \(M\) 由一个顶点闭盘 \(D\)、\(m\) 条矩形带及一个面闭盘 \(F\) 组成：各带的两个短边粘到 \(\partial D\) 上互不相交的区间，\(D\) 与全部带的并只有一个边界圆周，再由 \(F\) 封住。给 \(D\) 选定向，沿带传递此定向后，若与另一端 \(D\) 的定向一致，称该带无扭，否则称其有扭。两条分别标记为 \(a,b\) 的带称为交错，是指它们的粘合区间沿 \(\partial D\) 的循环次序为 \(a,b,a,b\)。若存在一条有扭带，则 \(M\cong\mathbb RP^2\#M'\)，其中 \(M'\) 有 \(m-1\) 条带的一顶点一面表示。若有两条交错的无扭带，则 \(M\cong T^2\#M'\)，其中 \(M'\) 有 \(m-2\) 条带的一顶点一面表示。
\end{lemma}
\begin{proof}
	第一种情形令 \(N\) 为 \(D\) 加选中的一条有扭带，第二种情形令 \(N\) 为 \(D\) 加选中的两条无扭带。未选带的粘合区间仍在 \(\partial N\) 上，暂时不将它们粘上。

	先核对 \(N\) 的类型。带条模型由粘合区间在顶点盘边界上的循环次序、每条带的两端方向决定：保持这些数据的边界同胚，可以在每个间隙弧上延拓，再延拓到顶点盘；带上按矩形坐标延拓。这些同胚在粘合区间上相容，故给出整个带面的同胚。上面的矩形 Möbius 模型恰好具有单条有扭带的粘合数据；环面方形模型中两条生成边的小邻域则具有两条交错无扭带的粘合数据。沿这个邻域的边界，从 \(a\) 的一侧出发，依次经过 \(a,b,a^{-1},b^{-1}\) 四段；其边界是一条圆周，补上一张盘就是方形的对边识别。因此第二种 \(N\) 是单孔环面。两种情形的 \(\partial N\) 都只有一条圆周。

	现在用一个新圆盘 \(D'\) 代替 \(N\)，将所有未选带原封不动地粘到 \(\partial D'\) 上对应的区间，再保留原来的面圆盘 \(F\)。换盘时边界上的粘合区间、间隙弧及其次序全部保持，因此新带面边界仍是原来供 \(F\) 封住的一条圆周。得到的 \(M'\) 是闭连通曲面，且正好少了一条或两条带。这里连通性也可以直接从模型读出：新顶点盘与每条余带相接，面盘沿整个边界相接。

	给 \(N\) 的边界补一个盘，所得闭曲面分别为 \(\mathbb RP^2\) 或 \(T^2\)。把这个补盘及 \(M'\) 中的 \(D'\) 各自去掉，再按先前的边界对应粘合，恰好恢复 \(M\)。这正是所声称的连通和分解，同时给出了余下带条的具体表示。
\end{proof}

上述证明没有假定“去掉一个手柄后仍应有某个亏格”。余下曲面是通过换盘实际造出来的，带条数确实下降；这将保证约化在有限步后停止。若读者习惯看边字，交叉帽的抽取还可写成下面的切粘等式，它稍后也用来把混合模型化为统一形式。

\begin{lemma}[一个同向边对的切粘]\label{g1:lem:crosscap-word-move}
	设一个边配对多边形的循环边字为 \(aUaV\)，字母 \(a\) 不在 \(U,V\) 中出现，其余每个边标号仍按原来的方式配对。则通过一次切割和一次重粘，可将它改为
	\begin{equation}\label{g1:eq:crosscap-word-move}
		aUaV\ \sim\ xxU^{-1}V.
	\end{equation}
	其中 \(x\) 为新边标号，\(U^{-1}\) 表示倒转 \(U\) 的次序并反转每个字母的指数；符号 \(\sim\) 表示配对所得曲面同胚，而非只表示基本群同构。
\end{lemma}
\begin{proof}
	在多边形内部取一条对角线 \(d\)，连接两条 \(a\) 边的起点。沿它切开，得到边字分别为
	\[
		aUd^{-1},\qquad daV
	\]
	的两个盘；若 \(U\) 为空，可把新角点稍作分离，或者直接保留原来相邻的 \(aa\)。把第一张盘反向读作 \(dU^{-1}a^{-1}\)，再沿两张盘中标为 \(a\) 的边粘合。粘合两盘的一整条边仍得到一个盘；剩余外边界按次序为 \(ddU^{-1}V\)。旧 \(a\) 变成内部接缝，新 \(d\) 的两份取代它成为配对边，重标为 \(x\) 即得所需字。旧边的配对始终未变，故\cref{g1:prop:polygon-moves}保证这是商曲面的同胚。图\ref{g1:fig:crosscap-word-move}标出实际切线与两张盘的边字。
\end{proof}

\begin{figure}[htbp]
	\centering
	\input{Figures/Chapter06-CrosscapMove}
	\caption[同向边对的切粘]{蓝色对角线连接两条 \(a\) 边的起点。先切出 \(aUd^{-1}\) 与 \(daV\)，将前一盘反向读边，再沿旧 \(a\) 边粘合，新的外边界为 \(ddU^{-1}V\)。图中 \(U,V\) 可以分别代表多条边。}
	\label{g1:fig:crosscap-word-move}
\end{figure}

\subsection{可定向标准字}
\label{g1:subsec:0604:02}

为什么可定向性恰好让抽取只产生手柄？所需信息已经在带条的两端，只要把局部定向沿带延拓一次，就能判别。

\begin{lemma}[带条与定向]\label{g1:lem:band-orientation}
	设闭连通曲面 \(M\) 由一个顶点闭盘 \(D\)、有限条矩形带及一个面闭盘组成：带的短边粘到 \(\partial D\) 上互不相交的区间，顶点盘与带的并只有一条边界，再由面盘封住。给 \(D\) 选定向；若沿一条带传递此定向后与另一端 \(D\) 的定向一致，就称该带无扭。则 \(M\) 可定向，当且仅当每条带都无扭。
\end{lemma}
\begin{proof}
	若曲面有定向，将它限制到顶点盘和一条带，沿带延拓必须返回同一个局部定向，因此该带无扭。反过来，先给顶点盘定向，无扭条件恰好保证每条矩形带可选一个定向，使其两端同时与顶点盘相容。各带内部互不相交，于是这些局部定向给整个带面 \(B\) 定向。\(B\) 的边界只有一条圆周，给面圆盘选择使接缝两侧相容的定向，即延拓到闭曲面。这里使用的是沿边界粘合时两侧诱导的边界方向相反；一条连通边界上相容性的符号不会跳变，所以一处相容便沿整条边界相容。
\end{proof}

\begin{proposition}[可定向标准字]\label{g1:prop:orientable-standard-word}
	对非负整数 \(g\)，记 \(\Sigma_{g,0}\) 为球面添 \(g\) 个手柄所得的曲面，\(\Sigma_{0,0}=S^2\)。每个非空紧连通、可定向、无边界拓扑曲面同胚于某个 \(\Sigma_{g,0}\)，其中 \(g\ge0\)。当 \(g\ge1\) 时，它有一顶点一面边字
	\[
		a_1b_1a_1^{-1}b_1^{-1}\cdots a_gb_ga_g^{-1}b_g^{-1}.
	\]
	当 \(g=0\) 时模型为 \(S^2\)，可直接看成两张圆盘沿整条边界粘合。
\end{proposition}
\begin{proof}
	用\cref{g1:prop:one-vertex-band-model}取得有限带条模型；若没有带，则顶点盘与面盘沿边界粘合，所得为球面。否则所有带均无扭。我们先证明必有两条带的端点交错。

	假设没有交错的边对。在圆周上的有限配对区间中，选一对，使其某一侧的开弧所含端点数最小。若这条弧非空，则弧内任一端点的配对端点也在弧内，否则它与所选边对交错。于是弧内又有一对，其一侧所含端点更少，矛盾。因此有一条带的两个粘合区间相邻，中间没有其他带端。该带无扭，所以它的一条长边与这段不含端点的顶点盘边界弧构成 \(\partial B\) 的一整条圆周。它不经过任何其他带；该带另一条长边还属于另一条边界圆周。这与 \(\partial B\) 只有一个分支矛盾。

	故存在交错的两条无扭带，\cref{g1:lem:band-extraction}给出 \(M\cong T^2\#M'\)，并把带数减少二。余曲面仍可定向：沿切下部件的边界将其余部分的定向延拓到新顶点盘即可。反复进行，每步减少两条带，有限步后没有带，余下为球面。于是得到有限个环面的连通和；与球面作连通和只是把一个盘取出后换回另一个盘，不改变同胚类型。

	最后说明边字而不只命名模型。单孔环面的边界在其顶点与两条核心边附近读出 \(aba^{-1}b^{-1}\)。把若干单孔模型通过一个连接各顶点的小圆盘串联，沿该盘外围依次读各部件，就得到上述交换子串。等价地，在两个面盘内部各去一小圆盘，用一条弧把新边界接到原多边形的顶点，沿弧切开后，连通和的面字是两者边字的串接；这条弧的两份作为内部接缝相互消去。所有端点汇到同一个顶点，最后补上一张面盘。这个构造也直接说明串接所得对象就是所取的环面连通和。
\end{proof}

注意，算法停止的理由是带条数这个非负整数严格下降，不是预先知道某个 Euler 示性数必须等于 \(2-2g\)。后一个等式将在下一节从模型计算出来，承担的是区分模型的职责。

\subsection{不可定向标准字}
\label{g1:subsec:0604:03}

若模型中有扭带，先取出一个交叉帽，余下部分可能仍不可定向，也可能变得可定向。因此直接抽取所得通常是手柄与交叉帽的混合。真正需要的新步骤，是在已有一个交叉帽时把一个手柄换成两个额外的交叉帽。

\begin{lemma}[一个手柄与一个交叉帽]\label{g1:lem:handle-crosscap-trade}
	存在同胚
	\[
		T^2\#\mathbb RP^2\cong\mathbb RP^2\#\mathbb RP^2\#\mathbb RP^2.
	\]
\end{lemma}
\begin{proof}
	左边的边字可循环读作 \(ccaba^{-1}b^{-1}\)。在下面的等式链中，每个带“切粘”的箭头均应用\cref{g1:eq:crosscap-word-move}或其逆操作，每个带“循环”的箭头只改变读边的起点；\(x,y,z\) 每次都是新的边标号：
	\begin{align*}
		ccaba^{-1}b^{-1}
		&\ \overset{\text{切粘}}{\sim}\ cb^{-1}a^{-1}ca^{-1}b^{-1}\\
		&\ \overset{\text{循环}}{\sim}\ a^{-1}ca^{-1}b^{-1}cb^{-1}\\
		&\ \overset{\text{切粘}}{\sim}\ xxc^{-1}b^{-1}cb^{-1}\\
		&\ \overset{\text{循环}}{\sim}\ b^{-1}cb^{-1}xxc^{-1}\\
		&\ \overset{\text{切粘}}{\sim}\ yyc^{-1}xxc^{-1}\\
		&\ \overset{\text{循环}}{\sim}\ c^{-1}xxc^{-1}yy\\
		&\ \overset{\text{切粘}}{\sim}\ zzx^{-1}x^{-1}yy.
	\end{align*}
	为核对第一步，可在基本公式中取 \(U=b^{-1}a^{-1}\)、\(V=a^{-1}b^{-1}\)；则 \(U^{-1}V=aba^{-1}b^{-1}\)，所以使用逆操作恰好得到第一行右端。其后三次分别抽取同向的 \(a^{-1}\)、\(b^{-1}\)、\(c^{-1}\) 边对。最后把 \(x^{-1}\) 重标为 \(u\)，边字成为 \(zzuuyy\)。这是三个实射影平面模型的字串接，因而给出所需同胚。每一步都作用于盘及边的实际识别，不是从群展示相同推测曲面相同。
\end{proof}

\begin{proposition}[不可定向标准字]\label{g1:prop:nonorientable-standard-word}
	对正整数 \(k\)，记 \(N_{k,0}\) 为 \(k\) 个 \(\mathbb RP^2\) 的连通和。每个紧连通、不可定向、无边界拓扑曲面同胚于某个 \(N_{k,0}\)，其中 \(k\ge1\)，并有一顶点一面边字
	\[
		c_1c_1c_2c_2\cdots c_kc_k.
	\]
\end{proposition}
\begin{proof}
	取得一顶点一面带条模型。由\cref{g1:lem:band-orientation}，至少有一条有扭带，先按\cref{g1:lem:band-extraction}取出一个交叉帽。以后只要余模型还有有扭带，就再取出一个；若余模型的带均无扭，则按可定向情形的证明取出手柄。每次带条数至少减一，所以终止后得到
	\[
		\underbrace{T^2\#\cdots\#T^2}_{h\text{ 个}}\#\underbrace{\mathbb RP^2\#\cdots\#\mathbb RP^2}_{r\text{ 个}},\qquad r\ge1.
	\]
	连通和中的部件可以重新排列、重新加括号：取一个有若干圆孔的球面，把每个去盘部件分别接到一条边界上，便同时表示这些次序与括号。交换两个孔可在平面圆盘里交换两小圆盘，再保持其外侧不动；边界参数的保向改变可在环带内延拓。标准的环面、实射影平面各有反射对称，它在指定的小面盘上反转边界方向，因此粘合方向的另一选择也给出同胚模型。于是可以每次把一个手柄与已有的一个交叉帽放在一起，应用\cref{g1:lem:handle-crosscap-trade}。每用一次，手柄数减一、交叉帽数增二；最终得到 \(N_{2h+r,0}\)。

	实射影平面的一个面盘边字为 \(cc\)，将这些部件按前一命题末尾的串接构造连接，所得标准字为 \(c_1^2\cdots c_k^2\)。每个部件中顶点均相同，串接处又识别这些顶点，故这是一个顶点、\(k\) 条边和一张面的表示。
\end{proof}

\begin{example}[Klein 瓶的两个边字]\label{g1:exm:klein-two-crosscaps}
	通常将方形按 \(aba^{-1}b\) 配对得到的曲面称为\term[kleinping]{Klein 瓶}[Klein bottle]\footnote{Felix Christian Klein（1849--1925），德国数学家，以变换群组织几何的 Erlangen 纲领及函数论研究著称。}。循环读为 \(ba^{-1}ba\)，在\cref{g1:eq:crosscap-word-move}中取 \(U=a^{-1}\)、\(V=a\)，立即得到 \(xxaa\)。因此 Klein 瓶同胚于 \(N_{2,0}\)。这个计算说明：是否已经出现相邻的平方，不是判别有无交叉帽的条件；真正的数据是同一边对的方向和全部粘合次序。
\end{example}

\subsection{连通和与边界分支的恢复}
\label{g1:subsec:0604:04}

对带边界曲面，先封盘再约化很方便，但最后还必须说明：得到闭模型之后，原先补入的盘可能位置很不规则，能否把它们移到标准的互不相交小盘？下面给出所需的有限配置结论。

\begin{lemma}[有限个圆盘的位置]\label{g1:lem:disk-configuration}
	设 \(M\) 是连通拓扑曲面，\(D_1,\ldots,D_b\) 与 \(E_1,\ldots,E_b\) 是其内部的两组两两不相交的嵌入闭圆盘，且各圆盘边界在曲面内有双侧环带邻域。则存在 \(M\) 的自同胚，将第一组圆盘作为无序集合送到第二组。若需要保持圆盘的编号，也可按任意指定的对应实现。
\end{lemma}
\begin{proof}
	先说明两个局部操作。一个有双侧环带的闭圆盘连同稍厚的外环带仍是闭圆盘。在这个较大盘中，将原盘的边界参数化为某个同心圆，再作分段线性的径向伸缩，可把原盘缩到内部任意小的圆盘；伸缩在较大盘的边界固定，因而可在外面接恒等映射。这里用\cref{g1:lem:disk-extension}把边界参数延拓到盘内，外侧环带则按乘积参数延拓。

	其次，在一个坐标圆盘内，可以把一个足够小的圆盘移到邻近位置，同时固定坐标圆盘外部。取半径 \(r<R\) 的同心圆盘，小盘包含待移动的盘，大盘仍在坐标域内。令 \(\rho(x)=\max\{0,\min\{1,(R-|x|)/(R-r)\}\}\)，则 \(\rho\) 在小盘上等于一，在大盘外等于零，并且 \(|\rho(x)-\rho(y)|\le L|x-y|\)，其中 \(L=1/(R-r)\)。若向量 \(v\) 足够短，使 \(L|v|<1\)，则映射 \(x\mapsto x+\rho(x)v\) 在小盘上作平移、在大盘外保持不动。由
	\[
		|x+\rho(x)v-y-\rho(y)v|\ge(1-L|v|)|x-y|
	\]
	知其单射。对每个 \(z\)，方程 \(x=z-\rho(x)v\) 的右端为压缩映射，逐次迭代收敛到唯一解：相邻迭代差按同一小于一的比例衰减，故为 Cauchy 列，极限满足方程。于是该映射满射，逆映射由上式连续，确为所需自同胚。把平移分成足够短的若干步，就可沿位于有限坐标盘中的道路移动小盘。

	还需保证已经放好的盘不会挡住道路。连通曲面去掉有限个互不相交内盘的内部后仍连通：所得空间是有边界的局部道路连通空间，若有多个连通分支，每条新边界圆周完全落在其中一个分支；把所去的盘沿该圆周补回，只能与这个分支相接，不会连接不同分支。局部半圆盘及所补盘的邻域使这些分支补盘后仍给出原曲面的不交开分解，与连通性矛盾。因此余下曲面道路连通，利用边界环带可将道路略微推入其内部。

	现逐个安置圆盘。目标盘之外有非空开集：有限个互不相交的内闭盘不可能覆盖整个连通曲面。先将全部原盘缩得足够小，使该开集内有一个与它们不交的坐标小盘，再在这里选 \(b\) 个暂存小盘。将原盘逐个送到这些暂存位置，移动时保持其他当前圆盘不动；此阶段目标盘只标记最终位置，允许移动操作经过它们。随后逐个沿避开其他当前圆盘的道路移到指定目标盘内部，再在目标盘及其外环带中作径向扩张，填满目标盘。已经放好的目标盘保持不动。每一步只有有限个盘需要避开，前段保证可选择道路；每条道路的紧性又保证只需有限个局部移动。复合所得自同胚即实现所给对应。
\end{proof}

暂存位置的作用是避开一种实际困难：原盘可能横跨几个目标盘，直接逐个扩张到目标位置会撞上尚未搬走的盘。先全部移到目标区域之外，再逐个放入，便把这个障碍消除了。对于 \(b=0\)，结论就是恒等映射。

\begin{theorem}[紧曲面标准模型的存在]\label{g1:thm:surface-standard-existence}
	记 \(\Sigma_{g,b}\) 为球面添 \(g\) 个手柄后挖去 \(b\) 个互不相交闭圆盘的内部所得的曲面，\(N_{k,b}\) 为 \(k\) 个 \(\mathbb RP^2\) 的连通和后作同样的 \(b\) 次挖孔，其中 \(g,b\ge0\)、\(k\ge1\) 均为整数。每个非空紧连通、Hausdorff、第二可数的拓扑曲面 \(M\)，边界允许非空，都同胚于某个
	\[
		\Sigma_{g,b}\quad(g,b\ge0),\qquad\text{或}\qquad N_{k,b}\quad(k\ge1,\ b\ge0).
	\]
	其中 \(b\) 为 \(M\) 的边界分支数；可定向的 \(M\) 得到第一类，不可定向的 \(M\) 得到第二类。
\end{theorem}
\begin{proof}
	由\cref{g1:prop:cap-boundary}，\(\partial M\) 是有限条圆周的并，逐一封盘得到闭连通曲面 \(\widehat M\)。若 \(M\) 可定向，其定向在每个补盘上按边界相容的方式延拓；若 \(\widehat M\) 可定向，其限制当然给 \(M\) 定向。因此封盘不改变可定向性。

	由前面的两种标准字结论，\(\widehat M\) 同胚于 \(\Sigma_{g,0}\) 或 \(N_{k,0}\)。补入的 \(b\) 张盘在这个同胚下仍是两两不交的闭盘，接缝两侧的环带使它们满足\cref{g1:lem:disk-configuration}。用该引理把这些盘移到标准模型中指定的 \(b\) 个小盘，再去掉它们的内部，便得到 \(M\cong\Sigma_{g,b}\) 或 \(M\cong N_{k,b}\)。去盘后每张盘贡献一条边界且没有别的边界，所以参数 \(b\) 正好为原边界分支数。
\end{proof}

圆盘配置引理也补足了连通和构造的选择问题：改变取盘位置可由曲面自同胚实现；两个同向的圆周粘合参数之差，可以沿环带作同伦延拓。具体地，将保向圆周同胚提升为严格递增的 \(f:\mathbb R\to\mathbb R\)，满足 \(f(t+1)=f(t)+1\)，则 \((1-s)f(t)+st\) 仍严格递增，给出到恒等映射的同伦；把参数 \(s\) 取为环带的横坐标即可延拓。标准模型还有反射对称，允许反转某一粘合边界的方向。各部件接到带孔球面上的描述说明重新加括号和交换部件也不改变结果。因此 \(\Sigma_{g,b}\)、\(N_{k,b}\) 在固定参数下确有明确的同胚类型。

例如，\(\Sigma_{0,1}\) 是圆盘，\(\Sigma_{0,2}\) 是环带，\(N_{1,1}\) 是 Möbius 带；\(\Sigma_{0,3}\) 是三孔球面。若从 \(T^2\#T^2\#\mathbb RP^2\) 去掉三个内盘，连续使用两次手柄与交叉帽的换位得到 \(N_{5,3}\)。这些参数是由实际切粘读出的。下一节将证明它们不能再被改成另一组参数，从而把本节的模型存在性提升为分类定理。
% 纲要条目排为编号小节标题，并进入目录与 PDF 书签；说明段仍保留于纲要注释中。

% !TeX root = ../../Book1.tex
% 纲要标题：### 6.5 分类不变量与唯一性
% 纲要位置：几何学计划纲要.md，第 349 行。
% 正文状态：初稿；尚未正式审核。

% 纲要内容开始
% 
% * 定向覆盖及可定向性
% * Euler 示性数与亏格
% * 标准模型的非同胚性
% * 基本群展示的独立核验
% 
% 纲要内容结束

\section{分类不变量与唯一性}
\label{g1:sec:0605}

多边形约化说明每个紧连通曲面都能变成一个标准模型，但它尚未排除同一个曲面有两种不同的标准答案。例如，一个手柄能否经复杂的剪贴变成两个交叉帽？单看 Euler 示性数无法排除这种可能，因为两者都会使球面的示性数减少 \(2\)。我们还需要记录局部定向绕行一周后是否反转，并把边界分支数与示性数一起使用。

\subsection{定向覆盖及可定向性}
\label{g1:subsec:0605:01}

回忆\cref{g1:def:surface-orientation}：在每一点有两个局部定向，坐标之间的局部符号决定哪些选择彼此相容。可定向性要求这些局部选择能在整个曲面上连续接起来。即使曲面不可定向，我们也可以把两种选择都保留为两个不同的点，得到第五章所讨论的二叶覆盖。

\begin{theorem}[定向二叶覆盖]\label{g1:thm:orientation-cover}
	设 \(S\) 是非空连通曲面，允许有边界。令
	\[
		\widehat S=\{(x,o):x\in S,\ o\text{ 是 }x\text{ 处的一个局部定向}\},
		\qquad p(x,o)=x.
	\]
	在每个坐标邻域上，把与坐标正定向相容和与之相反的选择分别作为两张开片，则这些开片给出 \(\widehat S\) 的拓扑，使 \(p:\widehat S\rightarrow S\) 成为二叶覆盖，称为\term[dingxiangfugai]{定向覆盖}[orientation covering]。曲面 \(\widehat S\) 自然可定向，映射 \(\tau(x,o)=(x,-o)\) 是反转定向的甲板变换。并且：
	\begin{enumerate}
		\item \(S\) 可定向时，\(\widehat S\) 是两个分别同胚于 \(S\) 的连通分支；
		\item \(S\) 不可定向时，\(\widehat S\) 连通，甲板群为 \(\{1,\tau\}\)。
	\end{enumerate}
\end{theorem}
\begin{proof}
	设 \(\varphi_i:U_i\rightarrow V_i\) 为坐标。用符号 \(\varepsilon=\pm1\) 表示相对于 \(\varphi_i\) 的局部定向。若 \(\varphi_j\circ\varphi_i^{-1}\) 在点 \(x\) 的符号为 \(\sigma_{ji}(x)\)，则同一点的两种表示满足
	\[
		(x,\varepsilon)_i=(x,\sigma_{ji}(x)\varepsilon)_j.
	\]
	由\cref{g1:prop:orientation-local-sign}，\(\sigma_{ji}\) 局部常值，并满足复合的乘法规则。因此三张图重叠时识别相容，且任意两张上述开片的交集在各自坐标中开。这就给出了拓扑；每个 \(p^{-1}(U_i)\) 恰分成两张片，分别同胚于 \(U_i\)，故 \(p\) 为二叶覆盖。不同底点可用 \(S\) 的 Hausdorff 性分离，同一底点上方的两点可用两张片分离；从 \(S\) 的可数坐标基提起两份基又得第二可数性，所以总空间确为曲面。

	在 \((x,o)\) 附近，选与 \(o\) 相容的底空间坐标，将它经 \(p\) 提到所在片上。这些提升坐标之间保持定向，因而定义 \(\widehat S\) 的定向。交换 \(o\) 与 \(-o\) 则反转这一定向，并且保持投影，所以得到所述甲板变换 \(\tau\)。

	若 \(S\) 已选定全局定向 \(o_x\)，则 \(\{(x,o_x)\}\) 与 \(\{(x,-o_x)\}\) 是两个不交开闭子集，投影分别为同胚；连通性保证它们就是两个分支。反过来，若 \(\widehat S\) 不连通，由\cref{g1:prop:monodromy-components}，其分支对应二点纤维上的单值轨道。曲面连通且局部道路连通，因而道路连通，满足该命题的条件。二点集上的不传递作用只能有两个单点轨道，于是每个分支恰含每个纤维中的一个点；投影在这个分支上为双射的局部同胚，故为同胚。取其逆便得到连续的局部定向选择，说明 \(S\) 可定向。不可定向时总空间因此连通，而甲板变换由一个纤维点的像决定；二点纤维最多容许两个这样的变换，已经构造出的 \(1,\tau\) 就是全部。
\end{proof}

固定 \(x_0\) 的一个局部定向 \(o_0\)，沿环路 \(\gamma\) 提升到 \(\widehat S\)。终点或者是 \((x_0,o_0)\)，或者是 \((x_0,-o_0)\)。记相应符号为
\[
	w_S([\gamma])\in\{+1,-1\}.
\]
这称为\term[dingxiangtezheng]{定向特征}[orientation character]。第五章的同伦提升说明它只依赖环路类；按先走左边的约定拼接环路，符号相乘，故
\[
	w_S:\pi_1(S,x_0)\longrightarrow\{+1,-1\}
\]
为群同态。两种起始局部定向给出的符号相同。曲面可定向当且仅当 \(w_S\) 恒为 \(+1\)；不可定向时，定向覆盖对应子群 \(\ker w_S\)。这里曲面的小圆盘或小半圆盘可缩，故它局部道路连通且半局部单连通，确实满足\cref{g1:thm:based-covering-classification}的条件。

\ProseParagraphBreak

这个符号有直接的几何读法。在一个交叉帽的带条中绕中心走一周，横向方向反转，因而得到 \(-1\)；在一个手柄的两条标准环路上都得到 \(+1\)。边界分支的绕行也给出 \(+1\)：由\cref{g1:prop:boundary-collar}，边界附近有一个半环形领圈，它本身可定向。这个观察稍后会确定定向覆盖的边界数。

\subsection{Euler 示性数与亏格}
\label{g1:subsec:0605:02}

\cref{g1:prop:euler-triangulation-independent}已经说明，三角剖分中的数 \(V-E+F\) 不随剖分改变。为了计算标准模型，可以先使用边数少得多的多边形分解，再把它细分为三角形：分一条边会同时增加一个顶点和一条边，在一个有 \(m\) 条边的面内添中心并连到边界顶点，则顶点、边、面分别增加 \(1,m,m-1\)，都保持 \(V-E+F\)。若粘合后的边或顶点有重复，先在边上增加分点，并在面内取一圈较小的多边形，把各个角分开后逐片三角化；这些同样是上述保持交错和的细分。因此这样的有限多边形分解也可用于计算 \(\chi\)。

\begin{proposition}[标准曲面的 Euler 示性数]\label{g1:prop:surface-euler-formulas}
	对\cref{g1:def:surface-standard-models}中的标准模型，有
	\[
		\chi(\Sigma_{g,b})=2-2g-b\quad(g,b\ge0),\qquad
		\chi(N_{k,b})=2-k-b\quad(k\ge1,\ b\ge0).
	\]
	这里 \(\Sigma_{0,0}=S^2\)。沿边界封上一个圆盘使 \(\chi\) 增加 \(1\)；沿两个曲面的圆盘孔口作连通和时，
	\[
		\chi(S\mathbin{\#}T)=\chi(S)+\chi(T)-2.
	\]
\end{proposition}
\begin{proof}
	球面可取四面体边界的三角剖分，故 \(\chi(S^2)=4-6+4=2\)。当 \(g\ge1\) 时，可定向标准字给出一个顶点、\(2g\) 条边和一个面，因此 \(\chi(\Sigma_{g,0})=1-2g+1\)。不可定向标准字同样给出一个顶点、\(k\) 条边和一个面，因此 \(\chi(N_{k,0})=1-k+1\)。

	封盘时，先在边界圆周上取 \(m\ge3\) 个顶点；在加入的圆盘内取一个中心，连出 \(m\) 条边并得到 \(m\) 个三角形。已有边界的顶点和边不重复计数，所以新增加的交错和是 \(1-m+m=1\)。封 \(b\) 个盘便增加 \(b\)，给出带边界的公式。

	连通和先从两曲面各去一个开圆盘，所以交错和总共减 \(2\)。把两条孔口边界用相同数目的顶点与边细分后粘合，若边界各有 \(m\) 个顶点和 \(m\) 条边，粘合使总顶点数与总边数都减少 \(m\)，面数不变。因此粘合本身不改变交错和，得到最后的等式。
\end{proof}

\begin{proposition}[有限覆盖与 Euler 示性数]\label{g1:prop:surface-finite-cover-euler}
	设 \(S\) 是紧连通曲面，\(p:T\rightarrow S\) 为 \(d\) 叶覆盖，其中 \(1\le d<\infty\)。则 \(T\) 是紧曲面，允许不连通，并且
	\[
		\chi(T)=d\,\chi(S),
	\]
	其中不连通空间的 \(\chi\) 是各连通分支之和。
\end{proposition}
\begin{proof}
	覆盖的局部同胚性质给出曲面坐标；Hausdorff 性与第二可数性也可分别由底空间及有限张覆盖片取得。每个点选一个均匀覆盖的坐标邻域 \(U\)，再在其中选较小的邻域 \(V\)，使闭包 \(\overline V\) 紧且包含于 \(U\)。紧性给出有限个这样的 \(V\) 覆盖 \(S\)。每个 \(p^{-1}(\overline V)\) 是 \(d\) 个与 \(\overline V\) 同胚的紧集之并，有限个这样的原像覆盖 \(T\)，故 \(T\) 紧。

	对 \(S\) 的一个有限三角剖分作充分多次细分，使每个闭三角形都包含在某个均匀覆盖邻域中。这可以在有限复形的度量下使用\cref{g1:lem:lebesgue-number}，因为重心细分后三角形直径趋于零。每个顶点、开边和开三角形都恰有 \(d\) 个提升；一张三角形的各个面在提升中仍按原来的包含关系相接，给出 \(T\) 的三角剖分。于是三种数分别乘以 \(d\)，交错和也乘以 \(d\)。
\end{proof}

若 \(S=N_{k,b}\)，其定向覆盖连通且可定向，边界数为 \(2b\)：每条边界的单值恒等，所以边界的原像是两个分别一叶覆盖它的圆周。设总空间为 \(\Sigma_{h,2b}\)，则
\[
	2-2h-2b=2(2-k-b),\qquad h=k-1.
\]
所以实射影平面的定向覆盖是球面，Klein 瓶的定向覆盖是环面，Möbius 带的定向覆盖是环带。它们是同一个构造的三个低维数值例子，而非三个互不相关的特殊技巧。

\subsection{标准模型的非同胚性}
\label{g1:subsec:0605:03}

\begin{theorem}[紧连通曲面的分类与参数唯一性]\label{g1:thm:surface-classification-unique}
	每个非空紧连通曲面，允许有边界，恰属于下列同胚类型之一：
	\[
		\Sigma_{g,b}\quad(g,b\ge0),\qquad
		N_{k,b}\quad(k\ge1,\ b\ge0).
	\]
	两种类型彼此不同；在同一类型中，参数也唯一。等价地，边界分支数、可定向性和 Euler 示性数这三个量共同确定其同胚类型。
\end{theorem}
\begin{proof}
	存在性是\cref{g1:thm:surface-standard-existence}的结论。由\cref{g1:prop:boundary-invariance}，同胚必须把边界送到边界，因而保持边界分支数 \(b\)。同胚还把局部定向及其相容关系带到目标曲面，所以保持可定向性。由\cref{g1:lem:band-orientation}，标准的手柄粘合可选相容的定向，而交叉帽内有定向反转环路，故可定向模型与不可定向模型不可能同胚。最后，由三角剖分无关性，Euler 示性数也是同胚不变量。在可定向类型中，
	\[
		g=\frac{2-b-\chi(S)}2;
	\]
	在不可定向类型中，
	\[
		k=2-b-\chi(S).
	\]
	这两个公式分别唯一恢复剩下的参数。反过来，三个量相同的两曲面属于同一个标准模型，故彼此同胚。
\end{proof}

可定向模型的参数 \(g\) 称为\term[kuige]{亏格}[genus]。对不可定向模型，文献也把交叉帽数 \(k\) 称为不可定向亏格；本书在有混淆可能时直接写“交叉帽数”。这两个参数在 Euler 公式中的系数不同，不能在未说明可定向性的情况下使用同一个公式。

\begin{example}[三个不变量各自保留了什么]\label{g1:exm:surface-invariants}
	圆盘和 Möbius 带都有一条边界，但前者可定向、\(\chi=1\)，后者不可定向、\(\chi=0\)。环面和 Klein 瓶都是无边界且 \(\chi=0\)，由可定向性区分。圆盘和球面都可定向，却因边界数及 \(\chi\) 不同而不等价。可定向曲面若有一条边界，则其 Euler 示性数只能取 \(1,-1,-3,\ldots\)；例如“可定向、一条边界、\(\chi=0\)”不对应任何紧连通曲面。
\end{example}

\subsection{基本群展示的独立核验}
\label{g1:subsec:0605:04}

分类的唯一性没有使用基本群的完整结构。因此，重新从标准模型计算基本群可以核对剪贴时的生成元和关系，同时看清基本群究竟保留或遗忘了哪些几何信息。下面的交换子约定仍为 \([a,b]=aba^{-1}b^{-1}\)。

\begin{proposition}[曲面群展示]\label{g1:prop:surface-group-presentations}
	为标准曲面选定基点，并用道路把各边界环路接到基点，则可选生成元使
	\begin{align*}
		\pi_1(\Sigma_{g,b})
		&\cong\left\langle a_1,b_1,\ldots,a_g,b_g,d_1,\ldots,d_b
		\ \middle|\ [a_1,b_1]\cdots[a_g,b_g]d_1\cdots d_b=1\right\rangle,\\
		\pi_1(N_{k,b})
		&\cong\left\langle c_1,\ldots,c_k,d_1,\ldots,d_b
		\ \middle|\ c_1^2\cdots c_k^2d_1\cdots d_b=1\right\rangle.
	\end{align*}
	当 \(b=0\) 时删去全部 \(d_j\)；\(\Sigma_{0,0}=S^2\) 的群为平凡群。特别地，当 \(b\ge1\) 时，
	\[
		\pi_1(\Sigma_{g,b})\cong F_{2g+b-1},\qquad
		\pi_1(N_{k,b})\cong F_{k+b-1}.
	\]
\end{proposition}
\begin{proof}
	无边界且不为球面的情形，标准字的边界识别后是一点并的有限圆周图，其基本群由\cref{g1:prop:graph-fundamental}自由生成。再把多边形内部作为一个圆盘附着，\cref{g1:prop:attach-disk}恰把边界词的正规闭包商掉，所以分别得到交换子乘积或平方乘积的关系。球面基本群平凡见\cref{g1:prop:sphere-simply-connected}。

	对于带边界的模型，把被挖去的 \(b\) 个圆盘安排在标准多边形内部，从每条新边界到外边界选互不相交、内部在面内的弧，逐条沿弧切开。这样的弧可先把各小圆盘沿树连接，再取树的一组互不交的细带中心线；平面有孔圆盘沿这些弧切开后成为一个圆盘。若闭模型为球面，则先把球面的一张圆盘当作面，另一张当作外侧填充来做同样的构造。切口的两份方向相反，连接弧因而在附着词中成对出现；把这些弧和必要的顶点连接边选入生成树后，它们不贡献自由生成元。剩下的边是各手柄或交叉帽生成元，以及各边界圆周的 \(d_j\)。圆盘的边界依次读出标准字和 \(d_1,\ldots,d_b\)；通过选择绕行方向和连接弧次序，可以使关系具有陈述中的形式。再次应用图的基本群和附着圆盘命题，得到所写展示。若某个边界的方向或接入道路换了，字会逆转或共轭，但得到的是同一群的另一组生成元。

	当 \(b\ge1\) 时，关系中 \(d_b\) 只出现一次且指数为 \(1\)，所以它被唯一表示为其余生成元的一个词。把该词代回即可消去 \(d_b\) 与唯一的关系，留下自由群。更明确地，从其余字母的自由群到所展示群有自然满同态；反向把 \(d_b\) 送到关系所规定的逆词、其他字母保持不变，关系便成立，得到互逆同态。因此并未额外引入隐藏的关系。
\end{proof}

若只想用较容易计算的群不变量核验，可以令所有生成元交换。对任意群 \(G\)，把全部交换子生成的正规子群商掉，所得 \(G_{\mathrm{ab}}\) 称为\term[jiaohuanhua]{交换化}[abelianization]；每个到交换群的同态都唯一地经过这个商。由刚才的展示可直接得到
\[
	\pi_1(\Sigma_{g,0})_{\mathrm{ab}}\cong\mathbb Z^{2g},\qquad
	\pi_1(N_{k,0})_{\mathrm{ab}}\cong\mathbb Z^{k-1}\oplus\mathbb Z/2\mathbb Z.
\]
第一个等式因为每个交换子自动消失。第二个等式的唯一加法关系为 \(2(c_1+\cdots+c_k)=0\)：用 \(u=c_1+\cdots+c_k,c_2,\ldots,c_k\) 代替原基，逆变换为 \(c_1=u-c_2-\cdots-c_k\)，这是整数可逆换基，于是只有 \(u\) 受到 \(2u=0\) 的限制。闭可定向曲面因此没有这里的二阶扭元，而每个闭不可定向模型都有一个。

\ProseParagraphBreak

还有一个不必处理整数扭元的核验。到加法群 \(\mathbb Z/2\mathbb Z\) 的同态由生成元的像决定；在这个群中，交换子和平方都自动为零。所以对闭曲面，同态分别有 \(2^{2g}\) 个和 \(2^k\) 个，两个指数都等于 \(2-\chi\)。有边界时，唯一关系变成 \(d_1+\cdots+d_b=0\)，同态数则为 \(2^{1-\chi}\)。这说明模 \(2\) 的计数能检验 Euler 公式，却不能单独分辨可定向性。

\begin{example}[同一个自由群不足以确定曲面]\label{g1:exm:surface-group-not-classify}
	一次穿孔环面 \(\Sigma_{1,1}\)、三边界球面 \(\Sigma_{0,3}\) 和一次穿孔 Klein 瓶 \(N_{2,1}\) 的基本群都是 \(F_2\)。前两者的边界数不同，第一与第三者的可定向性不同，因此三个曲面互不同胚。若只记抽象基本群，边界环路对应的共轭类以及定向特征都可能被遗忘；这正是分类定理仍然保留边界与定向数据的原因。
\end{example}
% 纲要条目排为编号小节标题，并进入目录与 PDF 书签；说明段仍保留于纲要注释中。

% !TeX root = ../../Book1.tex
% 纲要标题：### 6.6 曲面上的本质曲线
% 纲要位置：几何学计划纲要.md，第 356 行。
% 正文状态：初稿；尚未正式审核。

% 纲要内容开始
% 
% * 分离和非分离曲线
% * 裤分解的存在与计数
% 
% ---
% 
% 纲要内容结束

\section{曲面上的本质曲线}
\label{g1:sec:0606}

分类给出了整张曲面的名字，但许多几何问题还要求知道怎样把它切成更小的部分。在球面上，任意简单闭曲线都会把曲面分开；在环面上，沿一条经线切开却只得到一个连通的环带。差别来自曲线在整张曲面中的位置，而不能仅凭曲线本身都是圆周来判断。先弄清曲线附近的两种带状模型，切开后的边界数与 Euler 示性数便能把这个问题化成可计算的分类问题。

\subsection{分离和非分离曲线}
\label{g1:subsec:0606:01}

本节的简单闭曲线指嵌入 \(c:S^1\hookrightarrow\operatorname{int}S\) 的像，常仍记为 \(c\)；它总在曲面的内部。紧性与 Hausdorff 性保证，连续单射 \(S^1\rightarrow S\) 已经是到其像的同胚。边界分支附近的平行曲线也可看成这样的内部曲线，但边界分支本身不作为切口。

\begin{lemma}[简单闭曲线的带状邻域]\label{g1:lem:curve-band-neighborhood}
	设 \(S\) 为允许带边界的拓扑曲面，\(c:S^1\hookrightarrow\operatorname{int}S\) 为嵌入，并仍以 \(c\) 表示其像。则 \(c\) 有一个闭邻域 \(B\subseteq\operatorname{int}S\)，使 \((B,c)\) 同胚于下列两种由带与中心圆周组成的模型之一：
	\[
		\bigl(S^1\times[-1,1],\,S^1\times\{0\}\bigr),
		\qquad
		\left(\frac{[0,1]\times[-1,1]}{(0,t)\sim(1,-t)},\,\{t=0\}/\sim\right).
	\]
	两种情形分别称 \(c\) 为\term[liangcequxian]{两侧曲线}[two-sided curve]和\term[yicequxian]{一侧曲线}[one-sided curve]。环路 \(c\) 的定向特征分别为 \(+1\) 和 \(-1\)，故两种类型与邻域的选择无关。
\end{lemma}
\begin{proof}
	先说明从局部直化怎样得到整个带。由\cref{g1:lem:proper-arc-straightening}，曲线上每一点都有矩形坐标，曲线在其中恰为一条水平中线。沿圆周参数选有限个依次重叠的这样的矩形，并在相邻矩形的交叠弧内部选一个横截小区间。可以把这些横截区间选得两两不交：所选中心点有限且互异，先取不交的小邻域，再在其中直化。两相邻横截区间之间的曲线段是紧的，把中线两边的宽度充分缩小，就能取一个以这两个区间为端截面的小矩形带；这一步可逐个穿过有限直化图，在交叠处沿已经选定的横截区间接续。非相邻曲线段的紧子段彼此不交，故还可缩小宽度，使不同小带除共同端截面外不再相交。

	先不接最后一个端截面，这串小矩形接成一个长矩形，中心线是一条弧。各次端截面的参数都是区间同胚，可递推改变下一矩形的横向参数，使前面所有接缝成为相同参数的识别；只剩最后一对端截面的识别 \(t\mapsto h(t)\)，其中 \(h\) 是闭区间同胚且 \(h(0)=0\)。区间同胚严格单调，所以 \(h\) 或保持两端，或交换两端。

	若 \(h\) 递增，则 \((1-u)h(t)+ut\) 对每个 \(u\in[0,1]\) 都是保持两端与 \(0\) 不动的区间同胚；沿长矩形纵向使用这族重参数，把末端识别直化为 \(t\mapsto t\)，得到环带。若 \(h\) 递减，则先复合 \(t\mapsto-t\)，用同一方法把递增的部分直化，末端识别便成为 \(t\mapsto-t\)，得到 Möbius 带。所有选择都在 \(\operatorname{int}S\) 内，且小带的并包含曲线的开邻域，故所构造的 \(B\) 符合要求。

	沿环带中心绕行时，纵向和横向的局部定向一起回到原值；沿 Möbius 带绕行时纵向不变而横向反转。因此两种带的定向特征分别为 \(+1,-1\)。定向特征只依赖原曲面中的环路类，所以同一曲线不能同时具有两种类型的带状邻域。
\end{proof}

特别地，可定向曲面内的简单闭曲线都是两侧的。不可定向曲面中则两种曲线都可能出现；“曲面不可定向”并不表示它的每一条闭曲线都反转定向。

\begin{definition}[沿曲线切开]\label{g1:def:cut-surface}
	设 \(S\) 是紧曲面，\(c\subseteq\operatorname{int}S\) 是简单闭曲线，\(B\subseteq\operatorname{int}S\) 是以 \(c\) 为中心圆周的闭环带或闭 Möbius 带邻域。把
	\[
		S_c\coloneqq S\setminus\operatorname{int}B
	\]
	称为沿 \(c\) 切开所得的曲面。这里把切口稍稍张开成带的宽度；两侧情形增加两条边界，一侧情形增加一条边界。它也有不依赖带宽的描述：在每个把 \(c\) 直化为水平中线的矩形中，将中线分别作为上、下两个闭半矩形的边界，原来的非中线点保持不变；在图的交叠处按原坐标变换接回同一侧。坐标变换将每个足够小的侧半矩形送到一侧，所以这些局部识别相容，定义了一个确定的切开空间。在带的横向参数中把 \([0,1]\) 线性伸缩到 \([\varepsilon,1]\)，两侧分别进行，一侧情形按端点反转相容地进行，便把这个空间同胚到删去较窄开带的模型。因此任一带模型给出相同的切开曲面类型。
\end{definition}

曲面 \(S_c\) 紧，边界由原边界和 \(B\) 的边界组成。在 \(S_c\) 上重新装回相应的带，就恢复原曲面；两侧情形也可用边界领圈直接把两条新边界逐点粘合。新出现的边界圆周属于切开后曲面的内在边界。

\begin{proposition}[切开的计数]\label{g1:prop:cut-surface-euler}
	设 \(S\) 是紧连通曲面，边界分支数为 \(b\)。沿内部简单闭曲线 \(c\) 切开后，\(S_c\) 至多有两个连通分支，并且
	\[
		\chi(S_c)=\chi(S).
	\]
	若 \(c\) 两侧，则 \(S_c\) 的边界分支总数为 \(b+2\)；若 \(c\) 一侧，则为 \(b+1\)，且 \(S_c\) 连通。
\end{proposition}
\begin{proof}
	在每个带的横向坐标中，把 \(t\ne0\) 沿同号方向推到相应的外边缘，公式可写为 \(t\mapsto \operatorname{sgn}(t)((1-u)|t|+u)\)。带外保持不动，便把 \(S\setminus c\) 形变收缩到 \(S_c\)。因此二者的连通分支一一对应。\(S_c\) 的每个分支必须接触一条新边界；否则该分支作为曲面的一个分支，在 \(S\) 中也是开闭子集，与 \(S\) 连通矛盾。每条新边界是一个连通圆周，只能属于一个分支，所以两侧时至多两个分支，一侧时只有一个。

	为计算交错和，对 \(S_c\) 和 \(B\) 分别三角剖分，再沿公共的边界圆周细分到相同的顶点。粘合每个边界圆周时，重复计算的顶点数和边数相同，因此 \(\chi(S)=\chi(S_c)+\chi(B)\)。环带和 Möbius 带的 Euler 示性数都是 \(0\)，由\cref{g1:prop:surface-euler-formulas}可得所需等式。边界总数直接来自两个带的边界分别有两条和一条圆周。
\end{proof}

若 \(S\setminus c\) 不连通，则称 \(c\) 为\term[fenliquxian]{分离曲线}[separating curve]；否则称为\term[feifenliquxian]{非分离曲线}[nonseparating curve]。一侧曲线总是非分离的。两侧曲线的判据可用一条横穿它的道路表达。

\begin{proposition}[两侧曲线的分离判据]\label{g1:prop:curve-separation-test}
	设 \(S\) 是紧连通曲面，\(c\) 是其内部的两侧简单闭曲线。则 \(c\) 非分离，当且仅当存在一条简单闭曲线 \(\beta\)，它与 \(c\) 恰交于一点，并在该点的矩形坐标中横穿 \(c\)。若 \(c\) 分离，则每一条只与 \(c\) 有有限个横穿交点的环路，交点数都为偶数。
\end{proposition}
\begin{proof}
	若 \(S\setminus c\) 连通，切开后的两条新边界属于同一分支。选在带坐标中对应的两个边界点 \(x_+,x_-\)，在切开曲面中用一条内部不碰边界的简单弧连接它们。这里可先用道路连通性取道路，再用有限个圆盘或半圆盘坐标把它替换为有限折线，并在三角剖分中消去重复边路，得到简单弧；端点附近用边界领圈保证其余部分都在内部。把这条弧与带中从 \(x_-\) 到 \(x_+\) 的一条横向区间相接，得到简单闭曲线 \(\beta\)，只在带的中心横穿 \(c\) 一次。

	若 \(c\) 分离，由两侧邻域及前述分支论证，补集恰有两个分支，分别邻接 \(c\) 的两侧。一条环路每横穿一次 \(c\) 就从一个分支进入另一个；从起点出发最后仍回到同一分支，所以交点数为偶数。因而不可能有上述只横穿一次的 \(\beta\)。若原起点在交点上，先沿环路略移起点到补集中即可。
\end{proof}

\begin{corollary}[可定向曲面切开后的亏格]\label{g1:cor:orientable-cut-types}
	对非负整数 \(g,b\)，记 \(\Sigma_{g,b}\) 为球面添 \(g\) 个手柄后挖去 \(b\) 个互不相交闭圆盘的内部所得的曲面，\(\Sigma_{0,0}=S^2\)。设 \(c\subseteq\operatorname{int}\Sigma_{g,b}\) 是简单闭曲线。若 \(c\) 非分离，则切开曲面为 \(\Sigma_{g-1,b+2}\)，特别地 \(g\ge1\)。若 \(c\) 分离，则切开曲面为
	\[
		\Sigma_{g_1,b_1+1}\sqcup\Sigma_{g_2,b_2+1},
		\qquad g_1+g_2=g,\quad b_1+b_2=b.
	\]
	这里 \(b_i\) 只数落在第 \(i\) 个分支中的原边界，另加的 \(1\) 是切口边界。
\end{corollary}
\begin{proof}
	曲面及切开后的各分支可定向，分类定理适用。非分离时，设新亏格为 \(h\)，由边界数和 Euler 等式得到
	\[
		2-2h-(b+2)=2-2g-b,
	\]
	故 \(h=g-1\)。分离时，两条新边界分别属于两个分支；将两个分支的 Euler 公式相加，与原曲面比较，即得 \(g_1+g_2=g\)。
\end{proof}

需要排除的曲线有两类直观的冗余：围住一个圆盘的曲线只切下一个盘，紧贴已有边界的曲线只切下一个环带。后一种称为\term[zhoubianquxian]{周边曲线}[peripheral curve]，也称边界平行曲线，精确含义是曲线与某条原边界共同围出一个嵌入的环带。在不可定向曲面上，还要注意围住 Möbius 带的两侧曲线。

\begin{definition}[本节的本质曲线约定]\label{g1:def:essential-surface-curve}
	对紧曲面内部的简单闭曲线 \(c\)，若它不围出嵌入圆盘，不是周边曲线，并且不作为一个嵌入 Möbius 带的边界，则称其为\term[benzhiquxian]{本质曲线}[essential curve]。最后一个排除项只在不可定向曲面上可能发生；本书讨论不可定向曲面的曲线系统时采用这一约定。若两条位于内部且互不相交的简单闭曲线共同围出一个嵌入环带，称它们平行。
\end{definition}

这个定义直接针对切割所得的曲面，而不把“本质”仅理解为基本群中的非单位元。例如 Klein 瓶可由两个 Möbius 带沿边界粘合，中间的接缝是两侧分离曲线；其群元素可以非平凡，却已被上述约定排除。在可定向曲面中没有 Möbius 带子曲面，定义只保留前两个排除项。由\cref{g1:cor:orientable-cut-types}，一条分离曲线不是本质曲线，恰好是在某一侧切出圆盘或只有一条原边界的环带；一条非分离曲线则必为本质曲线，因为圆盘或周边环带的内部会成为补集的一个单独分支。

\subsection{裤分解的存在与计数}
\label{g1:subsec:0606:02}

一个三边界球面 \(\Sigma_{0,3}\) 称为\term[kuxingqumian]{裤形曲面}[pair of pants]。它可画成从一个闭圆盘挖去两个不交开圆盘，三条边界分别看作腰口和两条裤脚；这些名字不赋予边界不同的数学地位。它的 Euler 示性数为 \(-1\)，而且其中每条内部简单闭曲线都围盘或与边界平行：由上面的切开公式，曲线必分离，原来的三条边界在两侧分配，必有一侧至多包含一条原边界。

\begin{definition}[裤分解]\label{g1:def:pants-decomposition}
	紧可定向曲面 \(S\) 的一个\term[kufenjie]{裤分解}[pants decomposition]是一组有限的、两两不交且两两不平行的本质简单闭曲线，使沿这些曲线全部切开后，每个连通分支都是裤形曲面。空曲线组也允许，例如裤形曲面自身已有这样的分解。
\end{definition}

\begin{theorem}[裤分解的存在与计数]\label{g1:thm:pants-decomposition}
	对非负整数 \(g,b\)，记 \(\Sigma_{g,b}\) 为球面添 \(g\) 个手柄后挖去 \(b\) 个互不相交闭圆盘的内部所得的曲面，\(\Sigma_{0,0}=S^2\)。设 \(S\cong\Sigma_{g,b}\) 为紧连通可定向曲面，且 \(\chi(S)=2-2g-b<0\)。则 \(S\) 有裤分解。每个裤分解都恰有
	\[
		3g-3+b
	\]
	条切割曲线，切开后恰有
	\[
		2g-2+b
	\]
	个裤形曲面。
\end{theorem}
\begin{proof}
	先构造分解。在标准的 \(g\) 个手柄上，逐次选取一条手柄经线。其补集连通：在圆盘加手柄的模型中，去掉这条经线只把该手柄打开成两个边界孔，其他部分仍接在同一张曲面上。也可画出同一手柄的纬线，它只横穿经线一次，再用\cref{g1:prop:curve-separation-test}判断。由\cref{g1:cor:orientable-cut-types}，每切一次，亏格减少 \(1\)，边界数增加 \(2\)。重复 \(g\) 次，得到有 \(n=b+2g\) 条边界的球面。假设 \(\chi<0\) 正好给出 \(n\ge3\)。

	若 \(n=3\)，已经是一条裤子。若 \(n\ge4\)，在球面上用一条曲线围住恰好两个边界孔，切下一个三边界球面；另一部分有 \(n-1\) 条边界。这样的曲线可以在平面有孔圆盘模型中构造：用一条避开其余孔的简单弧连接所选两孔，取这两个孔连同该弧的一个细邻域，其外边界就是所需曲线。反复进行，直到剩下三个孔，得到 \(n-2\) 个裤形曲面。所有新切口都在此前切开曲面的内部，所以放回原曲面时，各曲线仍两两不交。

	还须验证这些切口在原曲面中本质且互不平行。假如某一条围盘，或它与原边界之间围出环带，或两条切口共同围出环带，就把所围出的圆盘或环带再沿落在其中的其余切口全部切开。切开后的每个部分都是最终分解中的裤形曲面，且至少有一个。因此 Euler 示性数应是某个正整数的负数；但切开不改变 Euler 示性数，圆盘为 \(1\)，环带为 \(0\)，矛盾。原曲面可定向，也不会含 Möbius 带。因此上述构造确实给出裤分解。

	最后证明数目不依赖构造。设裤子数为 \(P\)，切口数为 \(C\)。切开不改变 Euler 示性数，每条裤子贡献 \(-1\)，故
	\[
		-P=2-2g-b,\qquad P=2g-2+b.
	\]
	切开后共有 \(3P\) 条边界，每条原边界被数一次，每条切口产生两条边界，所以
	\[
		3P=b+2C.
	\]
	代入 \(P\) 就得 \(C=3g-3+b\)。
\end{proof}

\begin{figure}[htbp]
	\centering
	\input{Figures/Chapter06-PantsTree}
	\caption[五边界球面的裤分解]{从五边界球面依次切下两条裤子，最后留下第三条。右图记录三个裤形分支的邻接关系：两个内部接缝各对应一条切割曲线，五个外端分别对应原来的五条边界。图记录拓扑粘合，不预设长度或角度。}
	\label{g1:fig:pants-tree}
\end{figure}

例如，闭的亏格 \(2\) 曲面需要 \(3\) 条切口，得到 \(2\) 条裤子；一次穿孔环面需要 \(1\) 条切口，得到 \(1\) 条裤子；五边界球面需要 \(2\) 条切口，得到 \(3\) 条裤子，如\cref{g1:fig:pants-tree}。恢复原曲面时，同一条裤子的两条边界可以彼此粘合，切口两侧因而可能属于同一块裤形曲面。

\ProseParagraphBreak

条件 \(\chi<0\) 来自每块裤形曲面的负 Euler 示性数。球面、圆盘、环带和闭环面分别有 Euler 示性数 \(2,1,0,0\)，沿闭曲线切开仍保持该数，因而不可能变成非空的裤形曲面之并。闭环面可以沿一条经线切成环带，但环带不是裤子。另一方面，裤形曲面本身满足 \(\chi=-1\)，其分解使用 \(0\) 条切口，符合公式，给出允许空曲线组的基本例子。

\ProseParagraphBreak

裤分解为后面的双曲几何提供了一个清楚的粘合问题：先给每条裤子配置适当的几何结构，再把对应边界按相容的数据接回去。拓扑分解记录块数与连接圆周；选定双曲几何结构后，边界长度与粘合扭转才成为新的数据。双曲裤子的存在与唯一性、边界长度的作用以及粘合中的扭转参数，都需要后卷的度量与测地线理论。这里已经确定的是分成多少块、沿哪些圆周连接，以及这些切口怎样反映曲面的亏格和边界。
% 纲要条目排为编号小节标题，并进入目录与 PDF 书签；说明段仍保留于纲要注释中。


\begin{chaptersummary}
	内点和边界点由局部拓扑区分；紧曲面只有有限个连通分支，边界是有限个圆周的不交并。平面同胚对小环路绕数的作用给出方向符号，符号在换图下相乘，由此定义局部定向、曲面定向与定向二叶覆盖。

	\ProseParagraphBreak

	Jordan 分离定理给出闭曲线的两侧，Schoenflies 延拓进一步把两侧的闭区域变成圆盘。简单弧直化和有限图邻域随后使圆盘边界能够逐个正则化，形成有限盘形分解并三角化。不同拓扑剖分的边可能无限相交，因此先通过环境同痕整理，再取公共细分；细分的局部计数证明 Euler 数与剖分无关。

	\ProseParagraphBreak

	对偶生成树把三角形粘成一个多边形；在闭曲面中再配合原骨架的生成树，可以得到一个顶点圆盘、有限条带和一个补圆盘。扭带抽出交叉帽，交错的无扭带抽出手柄；每次带数减少，约化因而终止。一个手柄与一个交叉帽可改写为三个交叉帽，这使不可定向情形归于统一模型。边界先封盘、最后恢复，圆盘搬移保证挖孔位置不会另添分类参数。

	\ProseParagraphBreak

	紧连通曲面恰为 \(\Sigma_{g,b}\) 或 \(N_{k,b}\)，其中前者可定向，后者不可定向。它们的 Euler 示性数分别为 \(2-2g-b\) 和 \(2-k-b\)；结合边界分支数与可定向性，即可唯一恢复参数。基本群展示检验这些模型中的环路和关系，但带边界曲面的自由基本群单独不足以完成分类。

	\ProseParagraphBreak

	简单闭曲线有环带或 Möbius 带邻域，沿曲线切开保持 Euler 示性数，改变边界数，并可能改变连通性与可定向性。负 Euler 示性数的可定向紧曲面可以分成裤子：\(\Sigma_{g,b}\) 的裤子数为 \(2g-2+b\)，切口数为 \(3g-3+b\)。这些是拓扑分解；在它们上面选择长度、扭转和双曲度量，还需要后续的几何论证。
\end{chaptersummary}

\BookExercises
以下各题均附完整解答。辨认曲面时先说明紧性、连通性、边界数与可定向性；使用 Euler 数前确认所用分解的面是圆盘；切割问题还需说明切口两侧怎样接回。


\begin{exercise}\label{g1:ex:surface-intrinsic-boundary}
	分别将闭圆盘 \(D^2\)、圆环 \(A=\{z\in\mathbb C:1\le |z|\le2\}\) 视为平面子集，将 \(S^2\) 视为三维空间子集。求各自的曲面边界及环境中的子集边界。说明开圆盘虽在平面中有圆周作为子集边界，却没有曲面边界；再说明本章为什么不把它称为闭曲面。
\end{exercise}
\begin{solution}
	闭圆盘的曲面边界与平面子集边界均为 \(S^1\)。圆环的两种边界均为 \(\{|z|=1\}\sqcup\{|z|=2\}\)。球面的每点都有平面圆盘图，所以曲面边界为空；它在 \(\mathbb R^3\) 中闭且内部为空，因此其环境边界是整个球面。

	开圆盘的每个点都有足够小的平面开圆盘邻域，所以没有曲面边界。它在 \(\mathbb R^2\) 中的闭包是闭圆盘，减去自身得到环境边界 \(S^1\)，但这些点不属于这个曲面。开圆盘不紧，例如递增开集 \(\{|z|<1-1/n\}\)、\(n\ge2\)，覆盖它却没有有限子覆盖。闭曲面要求紧且无曲面边界，因此开圆盘不满足前一条件。
\end{solution}

\begin{exercise}\label{g1:ex:wedge-disks-not-surface}
	把两张闭圆盘的中心识别为同一个点 \(p\)，其他点不识别，所得商空间为 \(X\)。证明 \(X\) 紧且 Hausdorff，但不是允许带边界的拓扑曲面。比较把两张圆盘沿整个边界逐点粘合的结果。
\end{exercise}
\begin{solution}
	有限个紧圆盘的不交并紧，故其商紧。不同于 \(p\) 的点可在原圆盘中取避开中心的相对开邻域；同一圆盘中的两点可按平面分开，不同圆盘中避开中心的邻域天然不交。把 \(p\) 与另一点分开时，在两张圆盘的中心各取小圆盘的并，另一个点取避开它们的邻域。这些集合均饱和，所以它们的像在商中开，证明 Hausdorff。

	\(X\setminus\{p\}\) 恰有来自两张圆盘的两个不交开闭部分。每个 \(p\) 的邻域都遇到两者，所以删去 \(p\) 后都不连通。但无论圆盘内点还是半圆盘直边中点，都有一个删去中心后仍道路连通的坐标邻域。这一必要局部性质失败，因此 \(X\) 不是曲面。

	沿整个边界粘合的情形不同。把两张圆盘分别映到上下闭半球，公共边界按同一赤道参数识别，便得到到 \(S^2\) 的连续双射；紧到 Hausdorff 的判据使它为同胚。赤道附近两个半圆盘共同组成圆盘邻域，并未产生前面的局部分叉。
\end{solution}

\begin{exercise}\label{g1:ex:orientation-pullback-cover}
	设 \(p:E\to M\) 是覆盖，\(M\) 是已定向的拓扑曲面，且 \(E\) 为 Hausdorff、第二可数。证明覆盖片使 \(E\) 成为拓扑曲面，并给出一个使 \(p\) 局部保持定向的唯一提升定向。证明任一甲板变换保持此定向。
\end{exercise}
\begin{solution}
	把 \(M\) 的坐标图缩到均匀覆盖邻域内；每张覆盖片经 \(p\) 同胚地映到该邻域，再接底空间坐标便得到 \(E\) 的圆盘或半圆盘图。题设已经给出两项整体拓扑条件，故 \(E\) 为曲面。

	在每张这样的提升图上采用底图的定向符号。两张提升图交叠时，坐标变换就是对应底图的坐标变换，所以相容关系成立。若要求 \(p\) 局部保持定向，任一点的符号就必须是这个符号，因此唯一。甲板变换 \(T\) 满足 \(pT=p\)，在上述坐标中，它只是从一张片移到另一张片而保持坐标值不变，局部方向符号为 \(1\)，故保持提升定向。
\end{solution}

\begin{exercise}\label{g1:ex:oscillating-edges}
	在平面中考虑两条简单弧
	\[
		A=\{(x,0):0\le x\le1\},\qquad B=\{(x,f(x)):0\le x\le1\},
	\]
	其中 \(f(0)=0\)，\(f(x)=x^2\sin(1/x)\) 对 \(x>0\)。证明它们有无限多个交点，却都是拓扑线段。说明为什么把这两条弧原样保留为边的两个有限剖分不可能有有限公共细分。最后给出一个局部环境同痕，把 \(B\) 的一个内点附近的短段直化。
\end{exercise}
\begin{solution}
	投影到第一坐标是每条弧到 \([0,1]\) 的连续双射，其逆分别是 \(x\mapsto(x,0)\)、\(x\mapsto(x,f(x))\)，因而两条弧都是拓扑线段。交点包含 \((1/(k\pi),0)\)，其中 \(k\ge1\)，这些点两两不同并趋近原点；它们没有构成公共线段，因为 \(f\) 不在任何正长度区间上恒为零。

	在一个有限单纯复形中，两条作为子复形的边链的交集仍为子复形。如果交集不含一条边，它就只能是有限个顶点。这里交集有无限多个孤立点且没有弧段，故不可能成为有限公共细分中的交集。这指出了“叠放两剖分即可细分”所缺少的有限性。

	取内点 \(x_0\in(0,1)\) 及一个小区间 \(J\) 包含 \(x_0\)。在竖直方向上用映射 \((x,y)\mapsto(x,y-tf(x))\) 即可把该段图像连续移到水平线；若要让修改有局部支撑，先在一条包含图像的窄矩形带中使用\cref{g1:lem:proper-arc-straightening}，得到固定矩形边界、把中线图像移到指定直径的同胚，再用\cref{g1:lem:graph-small-isotopy}证明中的圆盘公式把它连到恒等映射。这样修改只发生在指定带内，且每一时刻都是整个邻域的同胚。
\end{solution}

\begin{exercise}\label{g1:ex:double-euler}
	设 \(M\) 是紧曲面，边界非空，\(\widehat M\) 是沿边界粘合两个副本得到的双倍。直接用与边界相容的三角剖分证明 \(\chi(\widehat M)=2\chi(M)\)。分别对闭圆盘和圆环核对所得结果。
\end{exercise}
\begin{solution}
	记 \(M\) 的顶点、边、三角形数为 \(V,E,F\)，边界中的顶点数和边数为 \(V_\partial,E_\partial\)。两份副本粘合后，边界顶点和边界边各成对识别，三角形不识别，因此
	\[
		\chi(\widehat M)=(2V-V_\partial)-(2E-E_\partial)+2F
		=2\chi(M)-(V_\partial-E_\partial).
	\]
	每个边界分支是有限圈，顶点数等于边数，所以括号为零。双倍粘合得到的有限三角形分解若暂有重边，重心细分不改变交替和，故上述计数仍给出 \(\chi(\widehat M)\)。

	闭盘可取一个三角形，故 \(\chi(D^2)=3-3+1=1\)，双倍是两个盘沿圆周相接的球面，其 Euler 数为 \(2\)。圆环取内、外两圈各 \(m\ge3\) 个顶点，连接对应顶点，并把每个四边形沿一条对角线分开，得到 \(2m\) 个顶点、\(4m\) 条边和 \(2m\) 个三角形，交替和为 \(0\)。圆环双倍是环面，Euler 数仍为 \(0\)。
\end{solution}

\begin{exercise}\label{g1:ex:disk-two-marked-points}
	设 \(B,B'\) 是闭 Jordan 盘，\(f:\partial B\to\partial B'\) 是同胚，\(p\in B^\circ\)、\(p'\in B'^\circ\)。证明 \(f\) 有一个把 \(p\) 送到 \(p'\) 的闭盘同胚延拓。再证明，固定闭盘边界的任意同胚都与恒等映射经固定边界的环境同痕相连。
\end{exercise}
\begin{solution}
	先由\cref{g1:thm:jordan-schoenflies}延拓为 \(F:B\to B'\)，记 \(q=F(p)\)。把 \(B'\) 同胚到一个正方形，把 \(q,p'\) 的像分别向四个角连线，得到两个具有相同组合关系的扇形三角剖分。在对应三角形上取仿射同胚，它们在公共半径上相合且在正方形边界上恒等，因而粘成一个固定边界、把 \(q\) 的像送到 \(p'\) 的像的同胚。共轭回 \(B'\)，再与 \(F\) 复合即得所求。

	第二问把盘识别为单位圆盘，使用正文的缩放公式 \(U_t(x)=tu(x/t)\)（\(\|x\|\le t\)）以及 \(U_t(x)=x\)（\(\|x\|\ge t\)）。每个 \(U_t\) 固定边界，逆由 \(u^{-1}\) 的同一公式给出；\(\|U_t(x)-x\|\le2t\) 保证 \(t\to0\) 的连续性。于是得到固定边界的同痕。注意这个结论要求原同胚在边界上逐点恒等；反射圆盘的同胚不满足这一条件。
\end{solution}

\begin{exercise}\label{g1:ex:relative-triangulation-euler}
	设紧曲面 \(M\) 上嵌入了一个有限图 \(G\)，图与边界相容，且所有补区域都是开圆盘，每条边的两侧允许属于同一补区域。说明为什么能够选一个以 \(G\) 为子复形的有限三角剖分。若某补区域的边界行走重复经过顶点或边，为什么不能直接给这条边界行走加一个中心点并声称得到单纯复形？用环面的标准一顶点、两条环边、一张方形面的分解计算 Euler 数，并说明与真正三角剖分的联系。
\end{exercise}
\begin{solution}
	由\cref{g1:thm:surface-triangulation}的相对形式，先把 \(G\) 加入有限曲线网，再作顶点盘、边矩形和补区域的圆盘化，即可保留 \(G\) 并得到有限三角剖分。添加的顶点只把原边分成有限段，所以 \(G\) 成为一个子复形。

	如果边界行走把同一顶点访问两次，则从中心连出的两条半径可能有相同端点而是不同的边；如果同一边出现两次，同一个三角形的不同边也可能在商空间中识别。单纯复形由顶点集合唯一确定单形，因而这种未经处理的扇形画法只给出三角形粘合描述。先取顶点盘与边矩形，再细分这些盘形块，才能得到符合定义的三角剖分。

	环面的方形分解有 \(V=1,E=2,F=1\)，所以交替和为 \(0\)。为把这个计数与正式定义联系起来，可在方形上取 \(m\times m\) 方格并沿同向对角线分割，其中 \(m\ge3\)，然后识别相对边。商中有 \(m^2\) 个顶点、\(3m^2\) 条边和 \(2m^2\) 个三角形，确为单纯复形，交替和仍为 \(m^2-3m^2+2m^2=0\)。粗分解虽然便于计算，细分则说明这个数确实是环面的 \(\chi\)。
\end{solution}

\begin{exercise}\label{g1:ex:tetrahedron-two-trees}
	用四面体边界作为球面剖分，顶点为 \(0,1,2,3\)。选原骨架的生成树 \(T=\{01,02,03\}\)。写出不穿 \(T\) 的对偶生成树，并指出顶点圆盘加带条模型中剩多少条带条。
\end{exercise}
\begin{solution}
	四个面分别为 \(012,013,023,123\)。未在 \(T\) 中的边是 \(12,13,23\)；它们的对偶边分别连接 \(123\) 与 \(012,013,023\)。因此这三条对偶边是一棵以面 \(123\) 为中心的星形生成树。六条原边已经分别被 \(T\) 或对偶树使用，没有剩余带条。于是两个树的盘状邻域分别形成顶点盘与面盘，沿唯一圆周粘合，得到球面。整个构造直接给出了球面模型，而没有先用 Euler 数判别。
\end{solution}

\begin{exercise}\label{g1:ex:two-letter-polygons}
	对边字 \(aabb\) 和 \(abab^{-1}\)，分别追踪四个角的等价类，计算相应胞腔分解的 \(v-e+f\)。再与 \(aa^{-1}bb^{-1}\) 比较。这里尚不要求只凭这些整数判断前两张曲面的同胚类型。
\end{exercise}
\begin{solution}
	依边界顺序写角为 \(v_0,v_1,v_2,v_3\)。对 \(aabb\)，两条 \(a\) 同向配对给出 \(v_0=v_1=v_2\)，两条 \(b\) 的配对再令 \(v_2=v_3=v_0\)，所以只有一个顶点。对 \(abab^{-1}\)，\(a\) 的同向配对给出 \(v_0=v_2,v_1=v_3\)，而 \(b\) 的反向配对给出 \(v_1=v_0,v_2=v_3\)，同样只有一个顶点。两者均有两条商边、一张面，所以 \(v-e+f=1-2+1=0\)。

	对于 \(aa^{-1}bb^{-1}\)，正文已追踪出三个顶点，仍有两条边和一张面，故为 \(3-2+1=2\)。这说明“一个四边形、两对边”不足以确定 Euler 数，顶点等价关系不可省略。前两张曲面最终都属于不可定向的两交叉帽模型，但这一同胚还须靠后面的合法切粘或分类证明，而不是由相同整数单独断言。
\end{solution}

\begin{exercise}\label{g1:ex:capping-parameter}
	设 \(M\) 为紧曲面，\(C\) 为其一个边界圆周，\(h_0,h_1:S^1\to C\) 为两个同胚。显式构造两个封口空间 \(M\cup_{h_0}D^2\) 与 \(M\cup_{h_1}D^2\) 之间在 \(M\) 上恒等的同胚，并说明为何即使 \(h_0,h_1\) 在圆周上的方向相反也成立。
\end{exercise}
\begin{solution}
	在圆盘边界上取 \(k=h_1^{-1}h_0\)，径向延拓为 \(K(ru)=r k(u)\)、\(K(0)=0\)。在不交并 \(M\sqcup D^2\) 上，令映射在 \(M\) 上恒等、在圆盘上为 \(K\)。源空间中 \(u\sim h_0(u)\)，而目标空间中 \(K(u)=h_1^{-1}h_0(u)\sim h_0(u)\)，故它尊重识别，下降为连续映射。用 \(K^{-1}\) 得到逆映射，所以是同胚。

	径向延拓只要求 \(k\) 是圆周同胚，不要求保向。如果两个封口参数方向相反，\(K\) 就反转所添圆盘的方向，但并不改变未预先定向的封口曲面。若另要求封口后延续一个给定的曲面定向，则应在新圆盘上选择相应方向；这是附加的定向资料问题。
\end{solution}

\begin{exercise}
	用合法切粘识别边字 \(abab\) 与 \(aba b^{-1}\) 给出的闭曲面；说明为何不能只数它们都有两种字母，就认定是同一个模型。
\end{exercise}
\begin{solution}
	对 \(abab\)，在\cref{g1:eq:crosscap-word-move}取 \(U=b\)、\(V=b\)，得到 \(xxb^{-1}b\)。相邻反向边 \(b^{-1}b\) 可按\cref{g1:prop:polygon-moves}消去，余字 \(xx\) 给出 \(\mathbb RP^2=N_{1,0}\)。这里消去后仍有 \(xx\)，不涉及把球面的退化双边形误看成空空间的问题。

	对 \(abab^{-1}\)，同样取 \(U=b\)、\(V=b^{-1}\)，得到 \(xxb^{-1}b^{-1}\)。将 \(b^{-1}\) 重标为 \(y\)，便是 \(xxyy\)，故为 \(N_{2,0}\)，也就是 Klein 瓶。两份 \(b\) 的方向不同，改变了端点识别与可抽出的带条数据。边字中的字母数还可能因原来有多个顶点或相邻可消边而冗余，必须连同全部配对规则一起使用。
\end{solution}

\begin{exercise}
	对有边界曲面也可在内部各取一个盘作连通和。证明
	\[
		\Sigma_{g,b}\#\Sigma_{h,c}\cong\Sigma_{g+h,b+c},\qquad
		\Sigma_{g,b}\#N_{k,c}\cong N_{2g+k,b+c},\qquad
		N_{k,b}\#N_{\ell,c}\cong N_{k+\ell,b+c}.
	\]
	其中 \(g,h,b,c\ge0\)，\(k,\ell\ge1\)。说明这些构造不需要先使用下一节的 Euler 数公式。
\end{exercise}
\begin{solution}
	先将原有的 \(b+c\) 条边界全部封盘，再在避开这些盘的位置作连通和。由于取盘及粘合都可与所封的盘保持不交，最后去掉原补盘就恢复带边界的连通和；\cref{g1:lem:disk-configuration}保证这些余孔的位置不影响模型。因此只需计算闭模型，最后保留 \(b+c\) 个孔。

	两种可定向模型的边字都是交换子串，将字串接就有 \(g+h\) 个交换子，给出 \(\Sigma_{g+h,0}\)。两种不可定向模型的平方字串接，得到 \(k+\ell\) 个平方，给出 \(N_{k+\ell,0}\)。混合情形先得到 \(g\) 个手柄与 \(k\) 个交叉帽；因为 \(k\ge1\)，可连续 \(g\) 次使用\cref{g1:lem:handle-crosscap-trade}，每次把一个手柄及一个已有交叉帽换成三个交叉帽，最终有 \(2g+k\) 个交叉帽。整个论证逐步构造同胚，没有借助同胚不变量倒推模型。
\end{solution}

\begin{exercise}
	在圆盘上粘四条无扭带，八个粘合区间沿边界依次为 \(a,b,c,d,d,c,b,a\)。证明带面有五条边界；沿全部五条边界分别补盘后得到球面。由此说明“四条无扭带应当是两个手柄”的推断漏掉了什么条件。
\end{exercise}
\begin{solution}
	先只粘最内侧的 \(d\) 带。它的两个端区间之间没有其他端点，因此它的一条长边与这段盘边界形成一个新圆周，另一侧仍承载 \(a,b,c\) 的全部端点。再粘 \(c\) 带时，其两个端点在承载剩余端点的同一条边界上相邻，又分出一个没有其他端点的新圆周。\(b\)、\(a\) 两带同理。每加一条带，边界数增加一；从圆盘的一条边界出发，最终有五条。

	把每次分出的那条新圆周立即封盘。该带、所封圆盘和盘边界之间的小区域合在一起，只是把原盘沿一个边界弧向外扩大：沿横穿带的一条弧切开，可看成两盘沿整条边粘合，因而仍是盘。四次如此消去四条带，最后余下一个盘；再封住它唯一的边界，得到两盘沿圆周粘合的球面。

	抽取手柄的命题要求余带面只有一条边界，并且两条无扭带的端点交错。本题没有交错边对，也没有单边界条件，不能把带数直接除以二当作手柄数。
\end{solution}

\begin{exercise}
	设 \(M\) 是闭连通曲面，\(p_1,\ldots,p_b\) 是互异点，\(b\ge1\)。证明 \(M\setminus\{p_1,\ldots,p_b\}\) 形变收缩到一个从 \(M\) 去掉 \(b\) 个内开盘所得的紧曲面，但二者并不同胚。
\end{exercise}
\begin{solution}
	取互不相交的闭坐标盘 \(D_i\)，在每个盘的单位圆盘坐标中令 \(p_i\) 对应原点。对 \(0<r\le1\)，定义
	\[
		H_s(ru)=((1-s)r+s)u,\qquad u\in S^1,\quad 0\le s\le1.
	\]
	在各 \(D_i\setminus\{p_i\}\) 上使用此公式，在全部盘外保持不动。公式在 \(r=1\) 上固定，所以与盘外恒等映射相容，给出连续形变收缩，其终像为 \(M\setminus\bigcup_i\operatorname{int}D_i\)。

	这个终像是闭曲面 \(M\) 的闭子集，因而紧。去点空间不紧：若它紧，由 \(M\) 为 Hausdorff，它应为 \(M\) 的闭子集；但在任一被去点的坐标盘中，非中心点趋近中心，说明该去点空间不闭。因此两者不能同胚。这里的形变收缩只保证同伦类型相同，不保留紧性或流形的边界。
\end{solution}

\begin{exercise}\label{g1:ex:surface-cover-genus}
	\begin{enumerate}
		\item 设 \(p:T\rightarrow\Sigma_{g,0}\) 是连通 \(d\) 叶覆盖。证明 \(T\) 可定向，并计算其亏格。当 \(g=0\) 时，说明这限制了哪些 \(d\)。
		\item 一个连通三叶覆盖以一次穿孔环面为底空间，若边界环路在三点纤维上的单值是一个三循环，确定总空间的同胚类型。
	\end{enumerate}
\end{exercise}
\begin{solution}
	(1) 将底空间的定向通过覆盖片提起，坐标变换仍保持定向，故 \(T\) 可定向。有限覆盖的总空间紧，且底空间无边界，覆盖的局部模型也没有边界。若 \(T=\Sigma_{h,0}\)，则 \(2-2h=d(2-2g)\)，所以 \(h=1+d(g-1)\)。当 \(g=0\) 时，\(h=1-d\ge0\) 迫使 \(d=1\)，这与球面单连通、没有非平凡连通覆盖相符。

	(2) 底空间可定向，故总空间可定向；由 \(\chi(\Sigma_{1,1})=-1\)，得到 \(\chi(T)=-3\)。边界上方的连通分支对应边界环路在纤维上的轨道，一个三循环只有一个轨道，所以 \(T\) 只有一条边界。于是 \(2-2h-1=-3\)，得 \(h=2\)，即 \(T\cong\Sigma_{2,1}\)。这个单值条件可以实现：在 \(F_2=\langle a,b\rangle\) 上令 \(a,b\) 分别作换位 \((12),(23)\)，它们生成传递作用，而交换子作三循环；根据正文的展示，边界环路是交换子的逆，也仍作三循环。覆盖重建定理因而确实给出一个这样的连通覆盖。
\end{solution}

\begin{exercise}\label{g1:ex:surface-invariant-recovery}
	对下列数据，确定所有符合条件的紧连通曲面，或证明不存在：
	\begin{enumerate}
		\item 有两条边界，Euler 示性数为 \(-4\)，未说明可定向性。
		\item 无边界且可定向，Euler 示性数为 \(-3\)。
		\item 有一条边界，定向覆盖连通且同胚于 \(\Sigma_{3,2}\)。
	\end{enumerate}
\end{exercise}
\begin{solution}
	(1) 可定向时，\(2-2g-2=-4\)，故 \(g=2\)，得到 \(\Sigma_{2,2}\)。不可定向时，\(2-k-2=-4\)，故 \(k=4\)，得到 \(N_{4,2}\)。二者都存在，并因可定向性不同而不同胚。

	(2) 闭可定向曲面的示性数 \(2-2g\) 为偶数，不能为 \(-3\)，故不存在。

	(3) 定向覆盖连通说明底曲面不可定向。覆盖总空间的示性数为 \(2-6-2=-6\)，二叶性给出底空间示性数为 \(-3\)。代入 \(2-k-1=-3\)，得 \(k=4\)，所以底曲面为 \(N_{4,1}\)；其定向覆盖公式也恰给出 \(\Sigma_{3,2}\)。
\end{solution}

\begin{exercise}\label{g1:ex:surface-abelianization-test}
	设 \(G=\langle a,b,c\mid a^2b^2c^2=1\rangle\)。
	\begin{enumerate}
		\item 计算 \(G_{\mathrm{ab}}\) 及从 \(G\) 到 \(\mathbb Z/2\mathbb Z\) 的同态数。
		\item 找出一个与 \(G\) 的交换化有相同自由秩、但交换化不同构的闭可定向曲面群。
		\item 在 \(G\) 的七个非平凡同态 \(G\rightarrow\mathbb Z/2\mathbb Z\) 中，哪个核给出定向覆盖？该覆盖是什么曲面？
	\end{enumerate}
\end{exercise}
\begin{solution}
	(1) 交换化中只有关系 \(2a+2b+2c=0\)。以 \(u=a+b+c,b,c\) 为整数基，关系变成 \(2u=0\)，故 \(G_{\mathrm{ab}}\cong\mathbb Z^2\oplus\mathbb Z/2\mathbb Z\)。到 \(\mathbb Z/2\mathbb Z\) 时，原关系对三个生成元的任意像都成立，所以有 \(2^3=8\) 个同态。

	(2) 环面群的交换化是 \(\mathbb Z^2\)，自由秩同为 \(2\)，但没有非零二阶元素，故与上式不同构。

	(3) \(G\) 是 \(N_{3,0}\) 的群，三个标准交叉帽环路都反转定向。因此定向特征在加法记法下把 \(a,b,c\) 都送到 \(1\in\mathbb Z/2\mathbb Z\)。它非平凡，核的指数为 \(2\)，对应连通定向二叶覆盖。由定向覆盖公式，总空间为 \(\Sigma_{2,0}\)。其他六个非平凡同态也给出二叶覆盖，但不对应这个核；不能仅凭覆盖的叶数把它们认作定向覆盖。
\end{solution}

\begin{exercise}\label{g1:ex:capping-surface-group}
	设 \(S=\Sigma_{g,b}\) 或 \(N_{k,b}\)，且 \(b\ge1\)。沿其中一条边界封上圆盘。说明包含同态在基本群上的核，并以一次穿孔环面与 Möbius 带为例，分别写出封盘前后的群。
\end{exercise}
\begin{solution}
	将被封边界的环路用连接道路带到基点，记其类为 \(d\)。封盘正是沿该边界附着一个圆盘，故由\cref{g1:prop:attach-disk}，包含同态满射，核为 \(\langle\!\langle d\rangle\!\rangle\)。改变连接道路只把 \(d\) 共轭，改变方向只把它取逆，都不改变这个正规闭包。

	一次穿孔环面的群为 \(F_2=\langle a,b\rangle\)，边界类可取 \([a,b]^{-1}\)；封盘后加入 \([a,b]=1\)，得到 \(\mathbb Z^2\)，即闭环面群。Möbius 带的群为 \(\langle c\rangle\cong\mathbb Z\)，其边界类为 \(c^{-2}\)；封盘后加入 \(c^2=1\)，得到 \(\mathbb Z/2\mathbb Z\)，即实射影平面群。第二个例子还表明，边界圆周的一个生成元不必映成曲面群的原始生成元。
\end{solution}

\begin{exercise}\label{g1:ex:genus-two-cuts}
	在闭可定向亏格 \(2\) 曲面上，分别构造一条非分离简单闭曲线和一条分离的本质简单闭曲线，确定两种切开曲面的类型，并说明两条曲线不可能被曲面的某个自同胚互相送到对方。
\end{exercise}
\begin{solution}
	把曲面看成两个环面的连通和。在其中一个手柄上取经线，同一手柄的纬线只横穿它一次，所以经线非分离；切开后为 \(\Sigma_{1,2}\)。再取连通和接缝附近的中心圆周，它把两个手柄分开，切开后为两个 \(\Sigma_{1,1}\)。这两个分支都不是圆盘或环带，故接缝是本质曲线。自同胚给出相应补集之间的同胚，必须保持补集连通分支数；前者补集连通，后者有两个分支，所以不可能互相送到对方。
\end{solution}

\begin{exercise}\label{g1:ex:klein-curve-types}
	把 Klein 瓶写成两个 Möbius 带沿边界的并，选其中一条中心曲线 \(a\)，再选两带的公共边界 \(c\)。
	\begin{enumerate}
		\item 判断 \(a,c\) 的一侧或两侧、分离或非分离，并确定切开后的曲面。
		\item 用展示 \(\langle u,v\mid u^2v^2=1\rangle\) 证明 \(c\) 在基本群中非平凡，却不是本节约定的本质曲线。
	\end{enumerate}
\end{exercise}
\begin{solution}
	(1) \(a\) 是一侧曲线，所以非分离；切开第一条 Möbius 带的中心后，余下部分为一条边界领圈，接到另一条 Möbius 带上不改变后者的类型，因此总的切开曲面仍为 Möbius 带。也可由 \(\chi=0\)、一条新边界及仍含另一条一侧曲线，恢复 \(N_{1,1}\)。公共边界 \(c\) 有来自两条带各一半的环形邻域，故为两侧曲线，切开后是两条 Möbius 带，所以分离。

	(2) 取 \(u,v\) 为两条带的中心生成元，\(c\) 的类是 \(u^2\) 或其逆、共轭。定义到加法群 \(\mathbb Z\) 的同态 \(u\mapsto1,v\mapsto-1\)，则关系 \(u^2v^2\) 映为 \(2-2=0\)，同态良定义，而 \(u^2\) 映为 \(2\ne0\)。所以 \(c\) 非平凡。它却围出嵌入 Möbius 带，按本节约定被排除在本质曲线之外。这说明不可定向情形必须读清“本质”的具体定义。
\end{solution}

\begin{exercise}\label{g1:ex:pants-dual-graph}
	设 \(\Sigma_{g,b}\) 满足 \(\chi<0\)，给定一个裤分解。作一个有限图：每条裤子给一个顶点，每条切口给一条连接其两侧裤子的边，允许环边；原边界只作标记，不计作边。证明这个图连通，其基本群为 \(F_g\)。分别描述一次穿孔环面和闭亏格 \(2\) 曲面可能出现的图。
\end{exercise}
\begin{solution}
	原曲面连通，若裤子的邻接图分成两部分且没有连边，则把每部分的裤子按其内部接缝粘合后，得到原曲面的两个不交非空开闭部分，矛盾。因此图连通。顶点数 \(P=2g-2+b\)，边数 \(C=3g-3+b\)，由有限图的基本群公式，
	\[
		\pi_1(\text{邻接图})\cong F_{C-P+1}=F_g.
	\]
	一次穿孔环面只有一条裤子和一条切口，所以图是一个顶点加一条环边，剩下一条裤子边界作原边界标记。闭亏格 \(2\) 曲面有两个顶点、三条边，可以取三条边都连接这两个顶点的图：把两条裤子的三对边界分别粘合即可。也可以各在一个顶点添一条环边，再用第三条边连接两个顶点，这对应先把每条裤子的一对边界自粘成一次穿孔环面，再沿剩余边界粘合。两图都连通且边数减顶点数加一为 \(2\)，但它们的边与曲面内分离曲线的对应有所不同。
\end{solution}

\begin{exercise}\label{g1:ex:one-sided-cut-orientation}
	设 \(c\) 是 \(N_{k,b}\) 内的一侧简单闭曲线。证明：
	\begin{enumerate}
		\item 若切开后不可定向，则切开曲面为 \(N_{k-1,b+1}\)，并且 \(k\ge2\)。
		\item 若切开后可定向，则 \(k\) 为奇数，切开曲面为 \(\Sigma_{(k-1)/2,b+1}\)。
		\item 分别给出切开后不可定向、可定向的具体例子，说明一侧曲线不能只凭“删去一个交叉帽”这句话判断余下曲面的定向性。
	\end{enumerate}
\end{exercise}
\begin{solution}
	一侧切开后连通，边界数为 \(b+1\)，Euler 示性数仍为 \(2-k-b\)。若所得曲面为 \(N_{r,b+1}\)，则 \(2-r-(b+1)=2-k-b\)，所以 \(r=k-1\)，不可定向要求 \(r\ge1\)，即 \(k\ge2\)。若所得曲面为 \(\Sigma_{h,b+1}\)，则 \(2-2h-(b+1)=2-k-b\)，故 \(2h=k-1\)，得到第二项。

	第三项的不可定向例子是 Klein 瓶中一个交叉帽的中心曲线：切开后为 Möbius 带。可定向例子可取实射影平面中的射影直线，它是 Möbius 带的中心，切开后为圆盘。还有亏格较高的例子：把一次穿孔环面与 Möbius 带沿边界粘合，由\cref{g1:lem:handle-crosscap-trade}得到 \(T^2\#\mathbb RP^2\cong N_{3,0}\)；沿这条 Möbius 带的中心切开，余下的薄领圈只把一次穿孔环面的边界向外移，所得仍为 \(\Sigma_{1,1}\)。同样在 \(N_{3,0}\) 的三个交叉帽标准模型中，沿一个标准交叉帽中心切开却得到 \(N_{2,1}\)。可见余下曲面是否可定向取决于这条曲线的具体位置。
\end{solution}


% !TeX root = ../Book1.tex
% 纲要标题：# 第三编　曲线曲面与非 Euclid 几何
% 纲要位置：几何学计划纲要.md，第 363 行。
% 正文状态：导读初稿；本编各章状态见各章及审核进度。

% 纲要内容开始
% 
% 纲要内容结束

\BookPart{曲线曲面与非 Euclid 几何}{g1:part:03}
\BookPartGuide

一条曲线在空间中弯曲，一张曲面上的行者测量距离与角度，这两类经验引出不同的几何信息。曲线的曲率和挠率描述它在环境空间中的变化；曲面的第一基本形式记录沿面运动时能够测到的长度，第二基本形式则记录曲面在空间中的弯曲。把这些数据与参数选择分开，才能讨论哪些量属于形状本身，以及怎样从数据反过来恢复形状。

\ProseParagraphBreak

第七章由弧长与活动标架提取曲线的微分不变量，再由常微分方程重建曲线；闭合条件、转角和总曲率把局部数据与整体要求联系起来。第八章将重建问题推广到两个参数，在小参数邻域内，相容方程使沿不同方向积分得到同一个曲面，也将 Gauss 曲率写成第一基本形式的内在量。第九章通过平行移动和边界转角，将这个逐点曲率的积分与 Euler 示性数联系起来；第十章再利用角和与面积的关系构造常曲率多边形，研究它们的粘合及模型商。局部微分数据由此逐步连接到整张曲面的拓扑与几何结构。

\ProseParagraphBreak

所需分析基础是有限维微积分、Taylor 展开、连续函数的积分及基本收敛概念；曲线重建实际使用的线性常微分方程将在第七章给出证明。线性代数使用内积、正交基、行列式与线性方程。涉及闭曲线转角时会回顾第五章的圆周角度提升，涉及简单闭曲线的内部时会调用第六章的 Jordan 分离定理；具体调用处会回顾所用结论与记号。

% !TeX root = ../../Book1.tex
% 纲要标题：## 第7章　曲线的局部不变量与整体几何
% 纲要位置：几何学计划纲要.md，第 365 行。
% 正文状态：初稿；尚未正式审核。

% 纲要内容开始
% 
% > 本章目标与证明范围：以由微分不变量重建曲线为起点，再讨论闭合条件和整体曲率约束。Frenet 标架的适用范围、平行法标架和总曲率分别说明，避免把局部分类当成闭曲线分类。
% 
% 纲要内容结束

\chapter{曲线的局部不变量与整体几何}
\label{g1:ch:07}
\BookChapterNumber{g1:ch:07}

\begin{chapterabstract}
	曲线的弧长消除行走速度的影响，曲率与挠率记录单位长度内切向量和法方向的变化。本章先建立 Frenet 方程，说明零曲率点为何需要其他法标架，再完整证明由微分不变量重建曲线的存在与刚性唯一性。重建后的曲线是否闭合还需满足方向与位移的相合条件。对于闭曲线，切向映射把局部弯曲转成整体转角与总曲率，由此证明平面转角定理、凸性判据及空间曲线的 Fenchel 不等式，并说明结型所带来的进一步约束。
\end{chapterabstract}

\begin{learninggoals}
	\begin{itemize}
		\item 区分参数映射与轨迹，证明弧长公式与弧长重参数化，并正确处理方向和重复绕行。
		\item 推导曲率、挠率和 Frenet 方程，说明适用正则性及正曲率条件，构造能穿过零曲率点的平行法标架。
		\item 从连续系数线性方程的积分级数建立存在唯一性，利用正交性保持从给定不变量重建曲线。
		\item 分清保向刚性运动、空间反射和参数反向，检验周期数据何时真正给出闭曲线。
		\item 用切向角提升计算旋转次数，证明简单闭平面曲线的转角定理，并辨析曲率同号、凸性和简单性的关系。
		\item 将总曲率解释为球面切向曲线的长度，证明其下界与等号条件，读清非平凡结的总曲率限制所需的假设。
	\end{itemize}
\end{learninggoals}

把一根细线弯成圆，再弯成向上盘旋的螺线，我们可以在每一点描述它朝哪个方向走、向哪边弯。但这些描述还没有回答一个反问题：若只给出沿线的弯曲和扭转数据，是否足以重新摆出整根线？如果能够恢复，起点和朝向以外还会不会有别的自由？曲线基本定理将把这个问题化为一个标架方程和一次积分。

\ProseParagraphBreak

局部恢复以后，还需要问末端能否接回起点。圆柱螺线有恒定的曲率与挠率，却一直向上延伸；另一类平面曲线可以让切向量恰转一周，位置仍留下一个缺口。用 Frenet 方程重建正曲率曲线时，闭合还要求位置与 Frenet 标架在一周后相合。曲线真正闭合以后，其切向量还会受到额外限制：简单平面曲线的净转角只有一周，而空间闭曲线无论怎样扭动，总弯曲量至少为 \(2\pi\)。这些结论会把微分计算与整体几何接起来。

\ProseParagraphBreak

本章使用有限维微积分、内积与行列式。线性常微分方程的所需版本在正文证明；使用圆周的角度提升和 Jordan 分离定理时，将明确回引前两章的相应结果。参数化、曲率和经典平面曲线的进一步阅读可参见 Pressley 第二版的相应各节\cite[第1.1--1.4, 2.1--2.3, 3.1节]{Pressley2010}。关于非平凡结的严格总曲率下界，本章陈述结论并引用原始论文中的证明位置。

% !TeX root = ../../Book1.tex
% 纲要标题：### 7.1 参数、弧长与标架
% 纲要位置：几何学计划纲要.md，第 369 行。
% 正文状态：初稿；尚未正式审核。

% 纲要内容开始
% 
% * 正则参数化与重参数化
% * 弧长参数及方向选择
% * 曲率、挠率与 Frenet 方程
% * 零曲率点及替代法标架
% 
% 纲要内容结束

\section{参数、弧长与标架}
\label{g1:sec:0701}

同一条圆周可以匀速走过，也可以时快时慢地走过；速度和加速度会随走法改变，圆周的大小却没有改变。要从导数中提取形状信息，首先应消除这种速度选择。弧长参数使每单位参数恰好走过单位长度，随后切向量的变化才直接衡量曲线弯曲得多快。空间曲线还可能离开原来的弯曲平面，因此需要同时记录法方向怎样转动。

\subsection{正则参数化与重参数化}
\label{g1:subsec:0701:01}

本章主要讨论 \(\mathbb R^2\) 与 \(\mathbb R^3\) 中的曲线，使用标准 Euclid 内积；三维叉积按右手定向规定，即 \(\langle u\times v,w\rangle=\det(u,v,w)\)。当参数是一般的 \(t\) 时用上点表示求导；以弧长 \(s\) 为参数后，用撇号表示对 \(s\) 求导。这样做是为了在同一个公式中看清“每单位时间的变化”和“每单位长度的变化”。

\begin{definition}[正则曲线与重参数化]\label{g1:def:regular-curve}
	设 \(I,J\subset\mathbb R\) 是具有非空内部的区间，\(r\ge1\) 为整数。一个 \(C^r\) 映射 \(\alpha:I\to\mathbb R^n\) 称为 \(C^r\) 参数曲线；若 \(\dot\alpha(t)\ne0\) 对每个 \(t\in I\) 成立，就称其为\term[zhengzequxian]{正则曲线}[regular curve]。端点处的正则性按单侧导数理解，所说 \(C^r\) 均允许在端点附近作 \(C^r\) 延拓。若 \(h:J\to I\) 是导数处处非零的 \(C^r\) 双射，则 \(\beta=\alpha\circ h\) 称为 \(\alpha\) 的\term[chongcanshuhua]{重参数化}[reparametrization]；\(h'>0\) 时保方向，\(h'<0\) 时反方向。
\end{definition}

连续且不为零的 \(h'\) 在区间上符号恒定，所以这两种情况已经穷尽。逆函数连续可微，导数为 \(1/(h'\circ h^{-1})\)；对这个等式逐次求导便得到 \(C^r\) 正则性，因而定义中的 \(h\) 确实是微分同胚。链式法则给出 \(\dot\beta(u)=\dot\alpha(h(u))h'(u)\)，说明重参数化保持正则性。

\ProseParagraphBreak

正则性是参数映射的性质，不能只看轨迹集合。例如 \(t\mapsto(t^3,0)\) 的像是直线，导数却在 \(0\) 为零；它与 \(s\mapsto(s,0)\) 有相同轨迹，但 \(t\mapsto t^3\) 在 \(0\) 不满足这里的重参数化条件。另一方面，正则曲线可以自交：\(\alpha(t)=(\sin t,\sin2t)\) 的导数 \((\cos t,2\cos2t)\) 从不同时为零，而 \(\alpha(0)=\alpha(\pi)\)。速度不消失只保证局部能顺畅经过，不能保证不同时间不会经过同一点。

\begin{definition}[闭曲线与简单性]\label{g1:def:regular-closed-curve}
	设 \(P>0\)。一个 \(P\) 周期的 \(C^r\) 正则映射 \(\alpha:\mathbb R\to\mathbb R^n\) 称为 \(C^r\) 正则闭曲线，也视为定义在 \(\mathbb R/P\mathbb Z\) 上的映射。若 \(\alpha(t)=\alpha(u)\) 当且仅当 \(t-u\in P\mathbb Z\)，则称它为\term[jiandanbiquxian]{简单闭曲线}[simple closed curve]。不满足这个单射条件的闭曲线仍可能正则，甚至可能把一个简单轨迹重复走过多次。
\end{definition}

第六章中的 Jordan 曲线只要求圆周的连续嵌入；这里的简单正则闭曲线额外具有导数。若只给出 \([0,P]\) 上的公式，要求它成为 \(C^r\) 闭曲线，就必须检查端点处直到 \(r\) 阶的导数相合。仅有 \(\alpha(0)=\alpha(P)\) 还可能留下一个角点。

\subsection{弧长参数及方向选择}
\label{g1:subsec:0701:02}

内接折线的长度上确界给出曲线长度的几何定义。对连续曲线 \(\alpha:[a,b]\to\mathbb R^n\)，令
\[
	\ell(\alpha)=\sup_{a=t_0<\cdots<t_m=b}\sum_{i=1}^m\lVert\alpha(t_i)-\alpha(t_{i-1})\rVert.
\]
这个上确界可以为无穷；它把每次沿曲线走过的部分都计入，重复走圆周两次就会得到两倍周长。

\begin{proposition}[长度公式与弧长参数]\label{g1:prop:arc-length-param}
	设 \(\alpha:I\to\mathbb R^n\) 是 \(C^1\) 曲线。对每个紧子区间 \([a,b]\subset I\)，有
	\[
		\ell(\alpha|_{[a,b]})=\int_a^b\lVert\dot\alpha(t)\rVert\,\mathrm dt.
	\]
	若进一步 \(\alpha\) 为 \(C^r\) 正则曲线，\(r\ge1\)，固定 \(t_0\in I\)，定义
	\[
		s(t)=\int_{t_0}^t v(u)\,\mathrm du,\qquad v(t)=\lVert\dot\alpha(t)\rVert.
	\]
	则 \(s:I\to s(I)\) 是保方向的 \(C^r\) 微分同胚，\(\gamma(s)=\alpha(t(s))\) 满足 \(\lVert\gamma'(s)\rVert=1\)。它称为\term[huchangcanshu]{弧长参数}[arc-length parameter]下的曲线。两个描述同一条参数曲线的单位速率参数之间，只相差 \(s\mapsto s+c\) 或 \(s\mapsto -s+c\)。
\end{proposition}
\begin{proof}
	对任意分割，微积分基本定理及积分的三角不等式给出
	\[
		\sum_i\lVert\alpha(t_i)-\alpha(t_{i-1})\rVert
		=\sum_i\left\lVert\int_{t_{i-1}}^{t_i}\dot\alpha(t)\,\mathrm dt\right\rVert
		\le\int_a^b\lVert\dot\alpha(t)\rVert\,\mathrm dt.
	\]
	为了得到反向不等式，利用 \(\dot\alpha\) 在紧区间上的一致连续性。给定 \(\varepsilon>0\)，分割足够细时，在每个小区间内都有 \(\lVert\dot\alpha(t)-\dot\alpha(t_{i-1})\rVert<\varepsilon\)，于是
	\[
		\lVert\alpha(t_i)-\alpha(t_{i-1})\rVert
		\ge\lVert\dot\alpha(t_{i-1})\rVert(t_i-t_{i-1})-\varepsilon(t_i-t_{i-1}).
	\]
	求和并使分割网格趋于零，右侧的 Riemann 和趋于积分，再令 \(\varepsilon\to0\)，即得长度公式。

	正则时 \(v>0\)，所以 \(s'=v>0\)，\(s\) 严格递增。局部逆函数拼成整个像区间上的逆，且 \(\mathrm dt/\mathrm ds=1/v(t(s))\)。速度 \(v\) 属于 \(C^{r-1}\)，故 \(s\) 及其逆为 \(C^r\)，链式法则给出 \(\gamma'=\dot\alpha/v\)，其范数为 \(1\)。若另一单位速率表示是 \(\gamma(h(u))\)，则 \(1=\lVert\gamma'(h(u))h'(u)\rVert=|h'(u)|\)，连续性说明 \(h'\) 恒为 \(1\) 或恒为 \(-1\)，积分就得到结论。
\end{proof}

弧长反函数 \(t(s)\) 的存在由正速度保证。例如椭圆 \((a\cos t,b\sin t)\) 的弧长积分一般不容易显式反解。实际计算可以保留原来的参数，把每次弧长求导换成
\[
	\frac{\mathrm d}{\mathrm ds}=\frac1v\frac{\mathrm d}{\mathrm dt}.
\]
若曲线以 \(P\) 为周期，令 \(L=\int_0^P v(t)\,\mathrm dt\)，则 \(s(t+P)=s(t)+L\)，所以重参数化后成为 \(L\) 周期的单位速率闭曲线。这里改变的是行走速度，方向及走过的次数均保留。

\begin{example}[圆周与重复行走]\label{g1:exm:circle-speed}
	对半径 \(R>0\) 的圆 \(\alpha(t)=R(\cos t,\sin t)\)，有 \(v=R\)，以 \(\alpha(0)\) 为起点时 \(s=Rt\)，单位速率公式为 \(\gamma(s)=R(\cos(s/R),\sin(s/R))\)。一次周长是 \(2\pi R\)。若原参数在 \([0,4\pi]\) 上运行，总长度为 \(4\pi R\)；弧长化不会把两周自动变成一周。因此“形状相同”“参数等价”和“走过相同次数”是需要分别说明的事情。
\end{example}

\subsection{曲率、挠率与 Frenet 方程}
\label{g1:subsec:0701:03}

以下先用弧长。单位切向量 \(T=\gamma'\) 满足 \(\langle T,T\rangle=1\)，求导得到 \(\langle T',T\rangle=0\)：转向只能发生在法方向。单位长度内方向改变的大小，自然由 \(\lVert T'\rVert\) 衡量。

\begin{definition}[曲率与 Frenet 标架]\label{g1:def:frenet-frame}
	设 \(\gamma:I\to\mathbb R^3\) 为实区间 \(I\) 上的 \(C^2\) 单位速率曲线，令 \(T=\gamma'\)，定义其\term[qulv]{曲率}[curvature]为 \(\kappa(s)=\lVert T'(s)\rVert\)。在 \(\kappa>0\) 的部分，定义主法向量 \(N=T'/\kappa\) 与副法向量 \(B=T\times N\)。右手正交标架 \((T,N,B)\) 称为\term[Frenetbiaojia]{Frenet 标架}[Frenet frame]。\footnote{Frenet（Jean Frédéric Frenet，1816--1900）是法国数学家，以空间曲线的活动标架与相应微分公式著称。}若 \(\gamma\) 为 \(C^3\)，则 \(N,B\) 为 \(C^1\)，定义\term[naolv]{挠率}[torsion]为 \(\tau=-\langle B',N\rangle\)。
	\symidx{\(\kappa\)：空间曲线的非负曲率}
	\symidx{\(\tau\)：空间曲线的挠率}
\end{definition}

由 \(T,N\) 张成并经过曲线点的平面称为\term[mijiepingmian]{密切平面}[osculating plane]。Taylor 展开
\[
	\gamma(s+h)=\gamma(s)+hT(s)+\tfrac12h^2\kappa(s)N(s)+o(h^2)
\]
说明，在二阶近似中曲线沿着这张平面弯曲。副法向量 \(B\) 是这张平面的单位法向量，所以挠率衡量的是密切平面沿曲线的转动，正负号则依赖空间定向。

\begin{proposition}[Frenet 方程]\label{g1:prop:frenet-equations}
	设 \(\gamma:I\to\mathbb R^3\) 是 \(C^3\) 单位速率曲线，其曲率 \(\kappa\) 在整个实区间 \(I\) 上为正；以 \((T,N,B)\) 记其右手 Frenet 标架，以 \(\tau\) 记挠率。则
	\begin{equation}\label{g1:eq:frenet-equations}
		T'=\kappa N,\qquad N'=-\kappa T+\tau B,\qquad B'=-\tau N.
	\end{equation}
	这些等式也称为 Frenet--Serret 方程（Frenet--Serret equations）。\footnote{Serret（Joseph-Alfred Serret，1819--1885）是法国数学家，研究微分几何与分析；他独立给出了空间曲线的相应标架公式。}
\end{proposition}
\begin{proof}
	第一式就是 \(N\) 的定义。由 \(\langle B,B\rangle=1\) 得 \(B'\perp B\)；再对 \(\langle B,T\rangle=0\) 求导，有 \(\langle B',T\rangle=-\langle B,T'\rangle=-\kappa\langle B,N\rangle=0\)。因此 \(B'\) 只可能沿 \(N\)，其系数按定义是 \(-\tau\)。同样，\(N'\perp N\)，而
	\[
		\langle N',T\rangle=-\langle N,T'\rangle=-\kappa,\qquad
		\langle N',B\rangle=-\langle N,B'\rangle=\tau.
	\]
	把 \(N'\) 在正交基 \((T,N,B)\) 中展开就得到第二式。
\end{proof}

\begin{definition}[平面的有向曲率]\label{g1:def:signed-plane-curvature}
	在具有标准定向的平面 \(\mathbb R^2\) 中，令 \(J(x,y)=(-y,x)\)，即逆时针旋转 \(\pi/2\)。设 \(\gamma:I\to\mathbb R^2\) 为以弧长 \(s\) 为参数的 \(C^2\) 单位速率曲线，令 \(T=\gamma'\)。则 \((T,JT)\) 始终为正向正交标架；唯一的函数 \(k\) 使
	\[
		T'=k\,JT,\qquad k=\det(T,T')
	\]
	成立，称为\term[youxiangqulv]{有向曲率}[signed curvature]。曲线作为空间曲线时的非负曲率为 \(\kappa=|k|\)。左转时 \(k>0\)，右转时 \(k<0\)，该定义也允许 \(k=0\)。
	\symidx{\(k\)：定向平面曲线的有向曲率}
\end{definition}

在切向量落入一个小圆弧的区间内，可选连续角度 \(\theta\) 使 \(T=(\cos\theta,\sin\theta)\)。这时 \(T'=\theta'JT\)，所以 \(k=\theta'\)。角度换成 \(\theta+2m\pi\) 不改变导数；全局怎样选角度、闭合后角度能增加多少，将在\secref{g1:sec:0703}通过圆周覆盖来说明。

\begin{proposition}[一般参数下的计算公式]\label{g1:prop:curvature-torsion-formulas}
	设 \(\alpha:I\to\mathbb R^3\) 是 \(C^2\) 正则曲线，速度为 \(v=\lVert\dot\alpha\rVert\)。其非负曲率为
	\[
		\kappa=\frac{\lVert\dot\alpha\times\ddot\alpha\rVert}{v^3}.
	\]
	若 \(\alpha\) 为 \(C^3\) 且 \(\dot\alpha\times\ddot\alpha\ne0\)，则
	\[
		\tau=\frac{\det(\dot\alpha,\ddot\alpha,\dddot\alpha)}
		{\lVert\dot\alpha\times\ddot\alpha\rVert^2}.
	\]
	对 \(C^2\) 正则平面曲线，按参数递增的方向，其有向曲率为
	\[
		k=\frac{\det(\dot\alpha,\ddot\alpha)}{v^3}.
	\]
\end{proposition}
\begin{proof}
	写 \(T=\dot\alpha/v\)，则
	\[
		T'=\frac{\ddot\alpha}{v^2}-\frac{\dot v\,\dot\alpha}{v^3},\qquad
		\dot v=\frac{\langle\dot\alpha,\ddot\alpha\rangle}{v}.
	\]
	展开范数平方并使用 \(\lVert u\times w\rVert^2=\lVert u\rVert^2\lVert w\rVert^2-\langle u,w\rangle^2\)，就得曲率公式，这一计算在 \(\kappa=0\) 时也有效。平面的行列式计算同样消去平行于 \(\dot\alpha\) 的第二项，得到 \(k\) 的公式。

	在 \(\kappa>0\) 时，\(\dot\alpha=vT\)、\(\ddot\alpha=\dot vT+v^2\kappa N\)。对第二式按 \(t\) 求导，利用 Frenet 方程，\(\dddot\alpha\) 的 \(B\) 分量是 \(v^3\kappa\tau\)。因此
	\[
		\det(\dot\alpha,\ddot\alpha,\dddot\alpha)=v^6\kappa^2\tau,\qquad
		\lVert\dot\alpha\times\ddot\alpha\rVert^2=v^6\kappa^2,
	\]
	相除就得到挠率公式。分母不为零正是所需条件。
\end{proof}

\begin{example}[圆柱螺线]\label{g1:exm:circular-helix-invariants}
	设 \(a>0\)、\(b\in\mathbb R\)，考虑 \(\alpha(t)=(a\cos t,a\sin t,bt)\)。记 \(c=\sqrt{a^2+b^2}\)，则 \(v=c\)，并且
	\[
		\dot\alpha\times\ddot\alpha=(ab\sin t,-ab\cos t,a^2),\qquad
		\det(\dot\alpha,\ddot\alpha,\dddot\alpha)=a^2b.
	\]
	所以 \(\kappa=a/c^2\)、\(\tau=b/c^2\)。当 \(b=0\) 时是平面圆，挠率为零；当 \(b\ne0\) 时每转一周高度增加 \(2\pi b\)。曲率、挠率虽然都是常数，曲线却不闭合。下一节会反过来从给定的两个常数恢复这个模型。
\end{example}

反转行走方向和反转空间定向有不同影响。对 \(\widetilde\gamma(s)=\gamma(-s)\)，有 \(\widetilde T=-T\)、\(\widetilde N=N\)、\(\widetilde B=-B\)（右侧在 \(-s\) 处取值），由定义可得 \(\widetilde\kappa=\kappa\)、\(\widetilde\tau=\tau\)。平面的有向曲率却变成 \(\widetilde k=-k\)。若作空间等距变换 \(x\mapsto Qx+a\)，正交矩阵 \(Q\) 保持叉积的模，行列式则乘以 \(\det Q\)，故 \(\kappa\) 不变、\(\tau\) 乘以 \(\det Q\)；平面有向曲率也乘以平面变换的行列式。刚性变换具有上述形式，见\cref{g1:thm:isometry-affine}。

\subsection{零曲率点及替代法标架}
\label{g1:subsec:0701:04}

Frenet 标架在 \(\kappa=0\) 时不能用 \(N=T'/\kappa\) 定义。光滑曲线仍可经过这种点，也可具有连续的其他法向选择。例如 \(\alpha(t)=(t,t^3,0)\) 处处正则，而
\[
	k(t)=\frac{6t}{(1+9t^4)^{3/2}}.
\]
取 \(v=\sqrt{1+9t^4}\)，在 \(t\ne0\) 时，\(N=\operatorname{sgn}(t)(-3t^2,1,0)/v\)、\(B=\operatorname{sgn}(t)(0,0,1)\)，左右极限相反。给 \(N(0)\) 取任意值都不能使它连续。相反，平面法向 \(JT=(-3t^2,1,0)/v\) 在拐点处完全连续；让曲率带符号，就能避免将法向跟随弯曲侧翻转。

\ProseParagraphBreak

\cref{g1:fig:inflection-frame}对比了这两种法向选择：蓝色左法向连续变化，Frenet 主法向在左侧取相反方向。

\input{Figures/Chapter07-InflectionFrame}

在三维中没有固定的平面旋转 \(J\) 可以选出唯一法向，但可以先在一个点选取一对正交法向，再规定它们沿曲线怎样继续。要求法向量的导数只沿切向，意味着法向量只作维持垂直所必需的转动，不额外在法平面内旋转。为构造这样的场，我们会使用\cref{g1:lem:linear-frame-ode}的连续系数线性常微分方程存在唯一性；该引理及其级数证明见下一节开头。

\begin{theorem}[平行法标架]\label{g1:thm:parallel-normal-frame}
	设 \(\gamma:I\to\mathbb R^3\) 为实区间 \(I\) 上的 \(C^2\) 单位速率曲线，令 \(T=\gamma'\)，固定 \(s_0\in I\) 和一组右手正交基 \((T(s_0),E_1^0,E_2^0)\)。存在唯一的 \(C^1\) 正交标架 \((T,E_1,E_2)\)，取给定初值并满足 \(E_i'\) 始终平行于 \(T\)。称它为\term[pingxingfabiaojia]{平行法标架}[parallel normal frame]。若 \(k_i=\langle T',E_i\rangle\)，则
	\begin{equation}\label{g1:eq:parallel-frame}
		T'=k_1E_1+k_2E_2,\qquad E_1'=-k_1T,\qquad E_2'=-k_2T,\qquad
		\kappa=\sqrt{k_1^2+k_2^2}.
	\end{equation}
	另一组同向平行法标架与这一组之间相差一个与 \(s\) 无关的法平面旋转。
\end{theorem}
\begin{proof}
	如果 \(E\perp T\) 且 \(E'\) 平行于 \(T\)，求导 \(\langle E,T\rangle=0\) 可得必须满足的方程
	\[
		E'=-\langle T',E\rangle T.
	\]
	它是未知向量 \(E\) 的线性方程，矩阵为 \(-T(T')^{\mathsf T}\)，系数连续。由上述线性方程引理，分别取初值 \(E_1^0,E_2^0\)，得到整个 \(I\) 上的唯一 \(C^1\) 解。还要确认解确实留在法平面中。直接计算
	\[
		\frac{\mathrm d}{\mathrm ds}\langle E_i,T\rangle
		=-\langle T',E_i\rangle\langle T,T\rangle+\langle E_i,T'\rangle=0.
	\]
	初值为零，所以始终垂直。继而 \((\langle E_i,E_j\rangle)'=0\)，因为 \(E_i',E_j'\) 都沿 \(T\)。因此长度与相互正交保持，标架的行列式连续且只能取 \(\pm1\)，初始右手性也保持。展开 \(T'\perp T\) 便得到全部方程及 \(\kappa\) 的表达式。

	若 \((\widetilde E_1,\widetilde E_2)\) 是另一组同向平行法标架，计算 \((\langle\widetilde E_i,E_j\rangle)'=0\)，两项都因法向与切向垂直而消失。所以变换矩阵常值，并且正交、行列式为 \(1\)，即一个固定的二维旋转。反过来，固定旋转后的标架仍满足同样的平行条件。
\end{proof}

平行法标架在 \(\kappa=0\) 时仍然有定义，此时 \(k_1=k_2=0\)，两个法向量的导数同时为零。在刚才的平面拐点例中，取 \(E_1=JT\)、\(E_2=(0,0,1)\)，便有 \(k_1=k\)、\(k_2=0\)。这两个带符号的分量可连续穿过零曲率点，使法向量保持连续。

\ProseParagraphBreak

在 \(\kappa>0\) 的区间，若写
\[
	N=\cos\theta\,E_1+\sin\theta\,E_2,\qquad
	B=-\sin\theta\,E_1+\cos\theta\,E_2,
\]
则 \(k_1=\kappa\cos\theta\)、\(k_2=\kappa\sin\theta\)。对第一式求导并代入\eqref{g1:eq:parallel-frame}，得到 \(N'=-\kappa T+\theta'B\)；与 Frenet 方程比较可知 \(\tau=\theta'\)。挠率因而是主法向相对于平行法标架的转角速度。在零曲率点，\((k_1,k_2)=(0,0)\) 没有确定的辐角，这正是不能一般地继续定义 Frenet 挠率的原因。

\ProseParagraphBreak

这里的存在唯一性是在一个参数区间上给定初值而言。若曲线闭合，沿一周继续平行法标架后，两个法向量未必回到起始值，可能相差一个固定旋转。初值唯一性确定沿途标架，一周后的法向旋转则记录闭合时的整体相合条件。

% !TeX root = ../../Book1.tex
% 纲要标题：### 7.2 曲线基本定理
% 纲要位置：几何学计划纲要.md，第 376 行。
% 正文状态：初稿；尚未正式审核。

% 纲要内容开始
% 
% * 标架常微分方程的存在与正交性保持
% * 从标架积分得到曲线
% * 刚性变换意义下的唯一性
% * 给定不变量与曲线闭合的额外条件
% 
% 纲要内容结束

\section{曲线基本定理}
\label{g1:sec:0702}

曲率与挠率是从曲线求导得到的两个函数。反过来，若只知道每一点弯曲多少、密切平面如何转动，能否恢复曲线？Frenet 方程提供了一种办法：先恢复三个随动的单位向量，再把切向量积分成位置。这里有两个需要分别回答的问题：方程是否有解，解出来的三个向量是否始终构成正交标架。Frenet 方程的斜对称系数将使内积导数为零，把初值处的正交性传到整个区间。

\subsection{标架常微分方程的存在与正交性保持}
\label{g1:subsec:0702:01}

下面以连续函数的积分、函数级数的一致收敛和微积分基本定理证明所需的线性常微分方程结论。以下区间均有非空内部；在区间包含的端点处，导数按单侧理解。

\begin{lemma}[连续系数线性方程]
	\label{g1:lem:linear-frame-ode}
	设 \(I\subseteq\mathbb R\) 是具有非空内部的区间，\(s_0\in I\)，\(A\colon I\rightarrow M_m(\mathbb R)\) 连续，端点处的导数按单侧理解。对任意 \(u_0\in\mathbb R^m\)，存在唯一的 \(C^1\) 函数 \(u\colon I\rightarrow\mathbb R^m\)，满足
	\[
		u'(s)=A(s)u(s),\qquad u(s_0)=u_0.
	\]
	特别地，连续矩阵系数方程 \(F'=FA\) 或 \(F'=AF\) 也有指定矩阵初值的唯一全区间解。
\end{lemma}

\begin{proof}
	把求导方程改写为积分方程
	\[
		u(s)=u_0+\int_{s_0}^s A(t)u(t)\,\mathrm{d}t.
	\]
	从常向量开始逐次积分，令
	\[
		u_0(s)=u_0,\qquad u_{j+1}(s)=\int_{s_0}^s A(t)u_j(t)\,\mathrm{d}t.
	\]
	这里 \(u_j\) 是将要相加的第 \(j\) 项，而不是第 \(j\) 次近似解。取包含 \(s_0\) 的任意紧子区间 \(J\subseteq I\)，令 \(M=\max_{t\in J}\lVert A(t)\rVert\)，范数取 Euclid 范数诱导的算子范数。反复积分并归纳可得
	\[
		\lVert u_j(s)\rVert\leq\lVert u_0\rVert\frac{M^j|s-s_0|^j}{j!},\qquad s\in J.
	\]
	当 \(s<s_0\) 时将积分反向后取绝对值，同一估计仍成立。于是级数 \(u=\sum_{j=0}^{\infty}u_j\) 在 \(J\) 上一致收敛。系数 \(A\) 有界，故可以把求和移入积分，得到上面的积分方程；被积函数连续，微积分基本定理便给出 \(u\in C^1\) 及 \(u'=Au\)。

	为证唯一性，设两解之差为 \(w\)，则 \(w(s_0)=0\) 且 \(w(s)=\int_{s_0}^s A(t)w(t)\,\mathrm{d}t\)。记 \(K=\max_J\lVert w\rVert\)，反复代入同一积分式，任意 \(j\geq0\) 均有
	\[
		\lVert w(s)\rVert\leq K\frac{M^j|s-s_0|^j}{j!}.
	\]
	右端随 \(j\to\infty\) 趋于零，故 \(w=0\)。

	每个 \(s\in I\) 都与 \(s_0\) 同处某个这样的紧区间。上述积分项在所有区间上由同一公式定义，因此得到 \(I\) 上的同一个解；无需假定 \(I\) 有界，也不要求 \(A\) 在整个 \(I\) 上有统一上界。矩阵方程 \(F'=AF\) 对每一列应用已证结论；\(F'=FA\) 则先转置，再对每一列应用它。
\end{proof}

上述阶乘估计解释了线性方程为何不会在有限的内部参数值处突然失去解：只要系数在这一点附近连续，紧区间上的整串积分仍然收敛。端点若不属于 \(I\)，则可能出现发散。例如 \(u'=u/(1-s)\) 在 \(I=(-\infty,1)\) 上的解可能在 \(s\to1\) 时发散。

\ProseParagraphBreak

把右手正交标架写成以向量为列的矩阵 \(F=(T,N,B)\)。由\cref{g1:prop:frenet-equations}，我们要解的是
\begin{equation}
	\label{g1:eq:frenet-matrix}
	F'=FA,\qquad
	A(s)=
	\begin{pmatrix}
		0&-\kappa(s)&0\\
		\kappa(s)&0&-\tau(s)\\
		0&\tau(s)&0
	\end{pmatrix}.
\end{equation}
第一列说 \(T'=\kappa N\)，第二列说 \(N'=-\kappa T+\tau B\)，第三列说 \(B'=-\tau N\)。矩阵 \(A\) 的各列记录导向量在当前标架中的系数，因此乘法写在右边。

\ProseParagraphBreak

行列式为 \(1\) 的正交矩阵组成\term[teshuzhengjiaoqun]{特殊正交群}[special orthogonal group]，这里记
\[
	\operatorname{SO}(3)\coloneqq\{Q\in M_3(\mathbb R):Q^{\mathsf T}Q=I,\ \det Q=1\}.
\]
它的列向量正好是右手正交标架。\symidx{\(\operatorname{SO}(3)\)：三维特殊正交群}

\begin{proposition}[斜对称系数保持标架]
	\label{g1:prop:skew-frame-orthogonality}
	设 \(I\subseteq\mathbb R\) 为具有非空内部的区间，\(s_0\in I\)，\(A\colon I\rightarrow M_3(\mathbb R)\) 连续且 \(A^{\mathsf T}=-A\)。设 \(F:I\to M_3(\mathbb R)\) 是方程 \(F'=FA\) 的 \(C^1\) 解，端点导数按单侧理解。若其初值 \(F(s_0)\) 属于 \(\operatorname{SO}(3)\)，则 \(F(s)\in \operatorname{SO}(3)\) 对每个 \(s\in I\) 成立。
\end{proposition}

\begin{proof}
	\cref{g1:lem:linear-frame-ode}给出全区间的 \(C^1\) 解。直接求导可得
	\[
		(FF^{\mathsf T})'=FAF^{\mathsf T}+FA^{\mathsf T}F^{\mathsf T}=0.
	\]
	因此 \(FF^{\mathsf T}=I\)。这使 \(F\) 可逆，且 \(F^{-1}=F^{\mathsf T}\)，从而也有 \(F^{\mathsf T}F=I\)。正交矩阵的行列式只能为 \(1\) 或 \(-1\)；它随 \(s\) 连续，而区间连通，所以只能恒等于初值 \(1\)。
\end{proof}

斜对称性恰好使长度与内积的变化抵消，因而把初值的一次正交条件传播到整个区间。

\subsection{从标架积分得到曲线}
\label{g1:subsec:0702:02}

现在可以把“弯曲和扭转决定形状”写成精确结论。条件 \(\kappa>0\) 保证 \(N=T'/\kappa\) 有意义；正则性条件则保证所得曲线确实具有本章用三阶导数定义的经典挠率。

\begin{theorem}[空间曲线基本定理：存在性]
	\label{g1:thm:space-curve-fundamental}
	设 \(I\subseteq\mathbb R\) 为具有非空内部的区间，\(\kappa\in C^1(I)\) 处处为正，\(\tau\in C^0(I)\)，端点处的导数按单侧理解。给定 \(s_0\in I\)、点 \(p_0\in\mathbb R^3\) 和右手正交标架 \((T_0,N_0,B_0)\)，存在唯一的 \(C^3\) 单位速率曲线 \(\gamma\colon I\rightarrow\mathbb R^3\)，使其曲率、挠率分别为 \(\kappa,\tau\)，并满足
	\[
		\gamma(s_0)=p_0,\qquad
		(T,N,B)(s_0)=(T_0,N_0,B_0).
	\]
	这一结论称为\term[quxianjibendingli]{曲线基本定理}[fundamental theorem of space curves]。
\end{theorem}

\begin{proof}
	用给定的 \(\kappa,\tau\) 组成\eqref{g1:eq:frenet-matrix}中的连续斜对称矩阵，解出初值 \(F(s_0)=(T_0,N_0,B_0)\) 的方程。由\cref{g1:prop:skew-frame-orthogonality}，三个列向量始终是右手正交单位向量。定义
	\begin{equation}
		\label{g1:eq:curve-frame-integral}
		\gamma(s)=p_0+\int_{s_0}^s T(t)\,\mathrm{d}t.
	\end{equation}
	于是 \(\gamma'=T\)，速率恒为 \(1\)。又因 \(T'=\kappa N\)，且 \(\kappa,N\) 都是 \(C^1\) 函数，故 \(T\in C^2\)，从而 \(\gamma\in C^3\)。这一步用到了 \(\kappa\in C^1\)，没有用到 \(\tau\) 可微。

	曲线的曲率为 \(\lVert\gamma''\rVert=\lVert\kappa N\rVert=\kappa\)。因 \(\kappa>0\)，其主法向量恰为构造中的 \(N\)，副法向量为 \(T\times N=B\)。最后 \(B'=-\tau N\)，所以其挠率 \(-\langle B',N\rangle\) 就是给定的 \(\tau\)。位置及标架的初值也由构造满足。

	任何满足条件的曲线，其 Frenet 标架都解同一矩阵初值问题，故由\cref{g1:lem:linear-frame-ode}与 \(F\) 相同。切向量既相同，位置初值也相同，积分式\eqref{g1:eq:curve-frame-integral}便使曲线相同。
\end{proof}

若只给连续的正函数 \(\kappa\) 和连续的 \(\tau\)，矩阵方程仍可解，积分仍给出 \(C^2\) 曲线，其 \(N,B\) 为 \(C^1\)。要得到 \(\gamma\in C^3\) 并使用涉及 \(\gamma'''\) 的经典挠率公式，还需加强数据的正则性。相反，若数据都光滑，反复对方程求导可知标架和曲线都光滑。

\begin{example}[常曲率和常挠率]
	\label{g1:exm:constant-curvature-helix}
	给定常数 \(\kappa_0>0\)、\(\tau_0\in\mathbb R\)，令 \(\omega=(\kappa_0^2+\tau_0^2)^{1/2}\)。曲线
	\[
		\gamma(s)=\left(
			\frac{\kappa_0}{\omega^2}\cos(\omega s),
			\frac{\kappa_0}{\omega^2}\sin(\omega s),
			\frac{\tau_0}{\omega}s
		\right)
	\]
	以单位速率运动，因为
	\[
		T=\left(-\frac{\kappa_0}{\omega}\sin(\omega s),
		\frac{\kappa_0}{\omega}\cos(\omega s),\frac{\tau_0}{\omega}\right).
	\]
	求导后得到 \(T'=\kappa_0N\)，其中
	\[
		N=(-\cos(\omega s),-\sin(\omega s),0),\qquad
		B=\left(\frac{\tau_0}{\omega}\sin(\omega s),
		-\frac{\tau_0}{\omega}\cos(\omega s),\frac{\kappa_0}{\omega}\right).
	\]
	直接计算 \(B'=-\tau_0N\)，所以曲率、挠率确为 \(\kappa_0,\tau_0\)。当 \(\tau_0=0\) 时得到半径 \(1/\kappa_0\) 的圆；当 \(\tau_0\ne0\) 时，第三坐标严格单调，得到不闭合的圆柱螺线。因而“曲率、挠率都是常数”并不蕴含闭合。
\end{example}

平面情形还能少解一个方程。由\cref{g1:def:signed-plane-curvature}，有向曲率 \(k\) 测量切向量向左转的速率。它可以为负，也可以经过零；只要改用转角，便不用除以曲率。

\begin{proposition}[由有向曲率重建平面曲线]
	\label{g1:prop:plane-curve-reconstruction}
	设 \(I\subseteq\mathbb R\) 为具有非空内部的区间，\(k\colon I\rightarrow\mathbb R\) 连续，端点处的导数按单侧理解。给定 \(s_0\in I\)、\(p_0\in\mathbb R^2\) 、\(\theta_0\in\mathbb R\) 及初始单位切向量 \((\cos\theta_0,\sin\theta_0)\)，令
	\[
		\theta(s)=\theta_0+\int_{s_0}^s k(t)\,\mathrm{d}t,\qquad
		\gamma(s)=p_0+\int_{s_0}^s(\cos\theta(t),\sin\theta(t))\,\mathrm{d}t.
	\]
	那么 \(\gamma\) 是具有这些初值和有向曲率 \(k\) 的唯一 \(C^2\) 单位速率平面曲线。
\end{proposition}

\begin{proof}
	定义式给出 \(\gamma'=T=(\cos\theta,\sin\theta)\)，而
	\[
		T'=k(-\sin\theta,\cos\theta)=kJT,\qquad J(x,y)=(-y,x).
	\]
	故 \(\det(T,T')=k\)，曲线为 \(C^2\) 且速率为 \(1\)。反之，任意具有有向曲率 \(k\) 的单位速率曲线都满足 \(\langle T,T'\rangle=0\) 以及 \(\det(T,T')=k\)，所以 \(T'=kJT\)。这是连续系数线性方程，指定初始切向量后由\cref{g1:lem:linear-frame-ode}唯一；再指定起点便唯一确定其积分。
\end{proof}

\subsection{刚性变换意义下的唯一性}
\label{g1:subsec:0702:03}

存在性定理指定了位置和整个标架，因而得到逐点相等的唯一性。通常说“不变量决定曲线”，并没有指定曲线在房间里的起点和朝向，此时应把平移与旋转造成的差别除去。\cref{g1:thm:isometry-affine}已经证明，Euclid 等距变换恰有形式 \(x\mapsto Qx+b\)，其中 \(Q\) 正交；保定向时还要求 \(\det Q=1\)。

\begin{proposition}[曲线的刚性唯一性]
	\label{g1:prop:curve-rigid-uniqueness}
	设 \(I\subseteq\mathbb R\) 为具有非空内部的区间，\(\gamma_1,\gamma_2\colon I\rightarrow\mathbb R^3\) 是其上的 \(C^3\) 单位速率曲线，曲率处处为正。若它们在同一参数 \(s\) 下有相同的曲率 \(\kappa(s)\) 和挠率 \(\tau(s)\)，则存在唯一的保定向 Euclid 等距变换 \(x\mapsto Qx+b\)，使
	\[
		\gamma_2(s)=Q\gamma_1(s)+b\qquad(s\in I).
	\]
	平面 \(C^2\) 单位速率曲线具有相同有向曲率时，也由唯一的保定向平面等距变换联系，不要求曲率非零。
\end{proposition}

\begin{proof}
	固定 \(s_0\in I\)，令两曲线的右手 Frenet 标架为 \(F_1,F_2\)，取
	\[
		Q=F_2(s_0)F_1(s_0)^{\mathsf T},\qquad
		b=\gamma_2(s_0)-Q\gamma_1(s_0).
	\]
	则 \(Q\in \operatorname{SO}(3)\)。因为数据相同，\(QF_1\) 和 \(F_2\) 满足同一方程 \(F'=FA\) 及同一初值，故相等。其第一列给出 \(\gamma_2'=Q\gamma_1'\)，积分即得所需等式。

	若另一个保定向刚性变换也逐参数匹配两曲线，对等式求一、二阶导数，得到它必须把 \(T_1(s_0)\)、\(N_1(s_0)\) 分别送到 \(T_2(s_0)\)、\(N_2(s_0)\)。保定向性又强制它把 \(B_1=T_1\times N_1\) 送到 \(B_2\)，因此正交部分必为 \(Q\)，平移部分再由起点确定。

	平面上，取唯一把初始单位切向量 \(T_1(s_0)\) 送到 \(T_2(s_0)\) 的旋转 \(Q\)，同样指定平移。旋转与 \(J\) 交换，故 \(QT_1\) 和 \(T_2\) 满足同一个方程 \(T'=kJT\) 及初值，于是曲线重合。唯一性只需初始切向量和起点。
\end{proof}

这里比较的是同一弧长参数下的函数。若先移动弧长原点，数据也要作相同平移；若逆转参数方向，也要先按\secref{g1:sec:0701}的变换规则重写数据。

\ProseParagraphBreak

还需分清参数反向与环境反射。若 \(\widetilde\gamma=Q\gamma+b\)，其中 \(Q\in O(3)\)，则 \(\widetilde T=QT\)、\(\widetilde N=QN\)，但
\[
	\widetilde B=\widetilde T\times\widetilde N=(\det Q)QB.
\]
这个叉积变换式可由 \(\langle u\times v,w\rangle=\det(u,v,w)\) 及正交矩阵的行列式直接验证。因此
\[
	\widetilde\kappa=\kappa,\qquad
	\widetilde\tau=-\langle\widetilde B',\widetilde N\rangle=(\det Q)\tau.
\]
环境反射把挠率变号；空间曲线的参数反向则不把挠率变号。若允许全部 Euclid 等距变换来比较曲线，相同曲率且挠率相反的一对曲线也可以互为镜像，挠率的整体改号对应镜像。

\subsection{给定不变量与曲线闭合的额外条件}
\label{g1:subsec:0702:04}

闭合要求走过一周的总位移为零，并要求参数接点处的切向量与法向量相容。标架方程的初值唯一性可以把这两项条件写成一周后的数据。

\begin{proposition}[周期数据的闭合判据]
	\label{g1:prop:periodic-curve-closure}
	设 \(L>0\)，\(\kappa\in C^1(\mathbb R)\) 为处处正的 \(L\) 周期函数，\(\tau\in C^0(\mathbb R)\) 也为 \(L\) 周期函数。令 \(A=\begin{psmallmatrix}0&-\kappa&0\\\kappa&0&-\tau\\0&\tau&0\end{psmallmatrix}\)，取任意 \(F(0)\in\operatorname{SO}(3)\)，以 \(F:\mathbb R\to M_3(\mathbb R)\) 解 \(F'=FA\)，即\eqref{g1:eq:frenet-matrix}。以 \(T\) 记 \(F\) 的第一列，给定 \(p_0\in\mathbb R^3\)，并令 \(\gamma(s)=p_0+\int_0^sT(t)\,\mathrm{d}t\)。则 \(\gamma\) 为 \(L\) 周期曲线，当且仅当
	\begin{equation}
		\label{g1:eq:space-closure-conditions}
		F(L)=F(0),\qquad \int_0^L T(s)\,\mathrm{d}s=0.
	\end{equation}
	此时它给出 \(\mathbb R/L\mathbb Z\) 上的 \(C^3\) 正则参数曲线，但不一定单射。条件是否成立与初始位置和右手标架的选择无关。
\end{proposition}

\begin{proof}
	若 \(\gamma\) 为 \(L\) 周期，则 \(\gamma'\)、\(\gamma''\) 也为 \(L\) 周期。由 \(T=\gamma'\)、\(N=\gamma''/\kappa\) 以及 \(B=T\times N\)，可知 \(F(L)=F(0)\)；积分条件就是 \(\gamma(L)-\gamma(0)=0\)。

	反过来，由系数的周期性，函数 \(F(s+L)\) 和 \(F(s)\) 解同一矩阵方程。它们在 \(s=0\) 的值相同，唯一性便给出 \(F(s+L)=F(s)\) 对所有 \(s\) 成立。特别地 \(T\) 周期，故函数 \(\int_s^{s+L}T(t)\,\mathrm{d}t\) 的导数为 \(T(s+L)-T(s)=0\)，它恒等于 \(\int_0^L T(t)\,\mathrm{d}t=0\)。于是 \(\gamma(s+L)=\gamma(s)\)。

	周期的 \(C^3\) 函数可通过商映射 \(\mathbb R\to\mathbb R/L\mathbb Z\) 下降：在每个长度小于 \(L\) 的坐标弧内用原函数表示，重叠处只差整数倍 \(L\) 的参数平移，因而各阶导数相容；速率仍为 \(1\)。改变初始位置只平移曲线，改变右手标架则使所有 \(F,T\) 同时左乘一个固定的 \(Q\in \operatorname{SO}(3)\)。这些操作都不改变\eqref{g1:eq:space-closure-conditions}中的两个等式是否成立。
\end{proof}

为看清周期性实际给出的内容，令 \(Q=F(L)F(0)^{\mathsf T}\)。即使 \(Q\ne I\)，方程唯一性仍然给出 \(F(s+L)=QF(s)\)。进一步积分得到
\[
	\gamma(s+L)=Q\gamma(s)+\bigl(\gamma(L)-Q\gamma(0)\bigr).
\]
所以周期数据保证，相邻两段长度 \(L\) 的曲线之间都相差同一个刚性运动；只有满足闭合条件时，这种重复才成为逐参数的闭合。

\ProseParagraphBreak

平面上，若 \(k\) 是连续的 \(L\) 周期函数，则\cref{g1:prop:plane-curve-reconstruction}把闭合条件化为
\begin{equation}
	\label{g1:eq:plane-closure-conditions}
	\int_0^L k(s)\,\mathrm{d}s\in2\pi\mathbb Z,\qquad
	\int_0^L \cos\theta(s)\,\mathrm{d}s=0,\qquad
	\int_0^L \sin\theta(s)\,\mathrm{d}s=0.
\end{equation}
第一个条件是切向量回到初始方向，后两个条件是位置回到起点。其充分性来自 \(T(s+L)=T(s)\) 与位移为零；必要性则由闭曲线的切向量周期和积分式得到。总转角为整周不能代替位移条件。

\begin{example}[转满一周却没有闭合]
	\label{g1:exm:turning-with-displacement}
	取 \(0<\varepsilon<1\)，先以转角 \(\theta\) 为参数，定义
	\[
		r(\theta)=\left(
			\sin\theta+\varepsilon\left(\frac{\theta}{2}+\frac{\sin2\theta}{4}\right),
			1-\cos\theta+\frac{\varepsilon}{2}\sin^2\theta
		\right).
	\]
	其导数为
	\[
		r'(\theta)=(1+\varepsilon\cos\theta)(\cos\theta,\sin\theta),
	\]
	速率严格为正。弧长为 \(s=\theta+\varepsilon\sin\theta\)，这是严格递增且映满 \(\mathbb R\) 的光滑函数；令其反函数为 \(\theta(s)\)。于是 \(\gamma(s)=r(\theta(s))\) 是单位速率曲线，有向曲率
	\[
		k(s)=\frac{1}{1+\varepsilon\cos\theta(s)}>0.
	\]
	由 \(s(\theta+2\pi)=s(\theta)+2\pi\)，得 \(\theta(s+2\pi)=\theta(s)+2\pi\)，故 \(k\) 以 \(L=2\pi\) 为周期，切向量也恰转一周。但是
	\[
		\gamma(2\pi)-\gamma(0)=r(2\pi)-r(0)=(\pi\varepsilon,0)\ne0.
	\]
	因此它满足标架接合条件却不满足位移条件。把它放入 \(\mathbb R^3\) 的水平面，便同时得到 \(\kappa=k>0,\tau=0\) 的空间反例。
\end{example}

\begin{figure}[htbp]
	\centering
	\input{Figures/Chapter07-Closure}
	\caption[切向量接合与位置闭合]{切向量接合与位置闭合。\(\varepsilon=\tfrac12\) 时，曲线转角从 \(0\) 增至 \(2\pi\)，两端的切向量相同，但端点相差 \((\pi/2,0)\)。箭头只表示相应端点的切向方向。}
	\label{g1:fig:closure-displacement}
\end{figure}

正面的例子是 \(k(s)=1/R\)、\(L=2\pi R\)，取初始切向量为 \((1,0)\)，积分得到
\[
	\gamma(s)=\bigl(R\sin(s/R),R(1-\cos(s/R))\bigr).
\]
三项闭合条件都成立，曲线恰为圆。若改用周期 \(4\pi R\)，同一公式仍满足条件，却把圆走了两遍。简单性还要求一个周期内只经过每个点一次。

% !TeX root = ../../Book1.tex
% 纲要标题：### 7.3 平面闭曲线的绕转
% 纲要位置：几何学计划纲要.md，第 383 行。
% 正文状态：初稿；尚未正式审核。

% 纲要内容开始
% 
% * 切向映射的次数
% * 转角定理与简单闭曲线
% * 曲率符号、凸性及反例
% 
% 纲要内容结束

\section{平面闭曲线的绕转}
\label{g1:sec:0703}

沿一条平面闭曲线走一周，位置回到原点，切向量也回到原方向；但连续记录的切向角可能多转了几周。这个整周差把局部的有向曲率同闭曲线的整体形状联系起来。简单性将这个整数限制为 \(1\) 或 \(-1\)，曲率同号则使切向角单调，从而把转角与凸性联系起来。

\subsection{切向映射的次数}
\label{g1:subsec:0703:01}

设 \(\gamma:\mathbb R/L\mathbb Z\rightarrow\mathbb R^2\) 是 \(C^1\) 正则闭曲线；用弧长参数时，\(L\) 是总长，\(T(s)=\gamma'(s)\in S^1\)。闭合在这里包含参数接点的导数相容，因而 \(T(0)=T(L)\)。\cref{g1:lem:circle-angle-lift}说明，给定初始角度后，可以沿区间 \([0,L]\) 唯一地选择连续函数 \(\theta\)，使
\[
	T(s)=(\cos\theta(s),\sin\theta(s)).
\]
前章把圆周角度按整周归一为 \(e^{2\pi iu}\)；本节使用弧度，所以这里的 \(\theta\) 是那里角度提升的 \(2\pi\) 倍。等式 \(T(0)=T(L)\) 只保证 \(\theta(L)-\theta(0)\in2\pi\mathbb Z\)，不要求这个差为零。

\begin{definition}[旋转次数]\label{g1:def:rotation-index}
	设 \(P>0\)，\(\gamma:\mathbb R/P\mathbb Z\to\mathbb R^2\) 为 \(C^1\) 正则闭曲线，令 \(T(t)=\gamma'(t)/\|\gamma'(t)\|\)，并取 \(T\) 在 \([0,P]\) 上的连续角度提升 \(\theta\)。单位切向映射的\term[cishu]{次数}[degree]称为曲线的\term[xuanzhuancishu]{旋转次数}[rotation index]，记为
	\[
		\operatorname{rot}(\gamma)\coloneqq\deg T
		=\frac{\theta(P)-\theta(0)}{2\pi}\in\mathbb Z.
	\]
	若采用弧长参数，则以总长 \(L\) 替代 \(P\)。这里次数就是\cref{g1:thm:circle-fundamental}中圆周映射沿定义域正向一周的绕数；任意两个角度提升相差常数 \(2\pi m\)，故右边不依赖初始角度的选择。
\end{definition}

切向映射在方向圆周上走，原曲线在位置平面上走。旋转次数记录前者，不能与原曲线绕某个指定点的绕数混同：后者需要先选一个不在曲线上的点，而前者不需要这样一个点。

\begin{proposition}[有向曲率的整周积分]\label{g1:prop:total-turning}
	设 \(\gamma:\mathbb R/L\mathbb Z\rightarrow\mathbb R^2\) 是 \(C^2\) 弧长参数正则闭曲线，\(k\) 是按 \(T'=kJT\) 定义的有向曲率，其中 \(J(x,y)=(-y,x)\)。则其连续切向角提升是 \(C^1\) 的，且
	\begin{equation}\label{g1:eq:total-signed-curvature}
		k(s)=\theta'(s),\qquad \int_0^L k(s)\,\mathrm{d}s=2\pi\operatorname{rot}(\gamma).
	\end{equation}
	对一般正则参数 \(t\in[0,P]\)，同一公式写为
	\[
		2\pi\operatorname{rot}(\gamma)
		=\int_0^P\frac{\det(\dot\gamma(t),\ddot\gamma(t))}{|\dot\gamma(t)|^2}\,\mathrm{d}t.
	\]
\end{proposition}
\begin{proof}
	在任一小开弧上，圆周有光滑角度坐标。连续提升在这样的坐标区间内与该局部角度相差固定的 \(2\pi\) 整数倍，所以它具有与 \(T\) 相同的 \(C^1\) 正则性。微分给出
	\[
		T'(s)=\theta'(s)(-\sin\theta(s),\cos\theta(s))=\theta'(s)JT(s).
	\]
	与有向曲率的定义比较得到 \(k=\theta'\)，然后应用微积分基本定理。一般参数下，\cref{g1:prop:arc-length-param}给出 \(\mathrm{d}s=|\dot\gamma|\,\mathrm{d}t\)，而有向曲率公式为 \(k=\det(\dot\gamma,\ddot\gamma)/|\dot\gamma|^3\)；将这两个等式相乘即得最后一式。
\end{proof}

\begin{proposition}[方向与连续变形]\label{g1:prop:rotation-invariance}
	旋转次数在保向正则重参数化和保向平面刚体运动下不变；反向重参数化或平面反射各使它变号。若 \(\gamma_u:\mathbb R/P\mathbb Z\rightarrow\mathbb R^2\) 是一族闭曲线，\((t,u)\mapsto\gamma_u(t)\) 与 \((t,u)\mapsto\partial_t\gamma_u(t)\) 均连续，且后者处处非零，\(u\in[0,1]\)，则 \(\operatorname{rot}(\gamma_u)\) 与 \(u\) 无关。
\end{proposition}
\begin{proof}
	保向重参数化只按原次序走过切向映射；反向重参数化使切向量成为原切向量的负向并逆序走过。取负向只给切向角加上 \(\pi\)，逆序则把总增量变号。保向平面刚体运动给角度加一个常数，反射把角度变为某个常数减去原角度，这证明前两句。

	对连续变形，单位切向量
	\[
		T_u(t)=\frac{\partial_t\gamma_u(t)}{|\partial_t\gamma_u(t)|}
	\]
	构成圆周映射的连续同伦。把平面识别为复平面，考虑 \(T_u(t)/T_u(0)\)，就得到基点始终为 \(1\) 的环路同伦。\cref{g1:thm:circle-fundamental}的同伦不变性使它们的绕数相同，而除以 \(T_u(0)\) 只减去常数角度。
\end{proof}

命题最后一类变形称为\term[zhengzetonglun]{正则同伦}[regular homotopy]：在变形中允许曲线穿过自身，但切向量不能消失。旋转次数因而是正则同伦的不变量。

\subsection{转角定理与简单闭曲线}
\label{g1:subsec:0703:02}

简单闭曲线要求 \(\gamma:\mathbb R/L\mathbb Z\rightarrow\mathbb R^2\) 单射。\cref{g1:thm:jordan-separation}给出其有界内部 \(\Omega\) 和无界外部。在曲线每一点，以切线为横轴，正则性使一小段曲线成为 \(C^1\) 函数的图像；把坐标矩形取得足够小，还能避开其余曲线。图像分出的两侧恰分别属于内部和外部，因为两者的边界都是整条曲线。内部处于行进方向的左侧还是右侧，是随参数局部不变的选择，故在整个参数圆周上一致。称内部位于左侧的行进方向为正向；反转参数即可在两种方向间转换。

\begin{theorem}[转角定理]\label{g1:thm:turning-tangent}
	设 \(\gamma:\mathbb R/L\mathbb Z\rightarrow\mathbb R^2\) 是 \(C^1\) 正则简单闭曲线。正向参数下 \(\operatorname{rot}(\gamma)=1\)，反向参数下 \(\operatorname{rot}(\gamma)=-1\)。特别地，若 \(\gamma\) 为 \(C^2\) 弧长参数曲线，则
	\[
		\int_0^L k(s)\,\mathrm{d}s=
		\begin{cases}
			2\pi,&\text{正向},\\
			-2\pi,&\text{反向}.
		\end{cases}
	\]
\end{theorem}

简单曲线可以有许多凹陷，切向角也可以来回变化。用两个点合拢时的弦方向表示相切方向，就能在参数三角形上统一记录这些变化。\cref{g1:fig:chord-triangle}显示了三条边各自承担的角度变化。

\begin{figure}[htbp]
	\centering
	\input{Figures/Chapter07-ChordTriangle}
	\caption[弦方向的参数三角形]{左图以曲线最低点为基点，此点发出的所有弦均指向闭上半平面。右图的两条直角边分别记录基点弦方向及其负向，各有角增量 \(\pi\)；斜边逆向记录切向映射。图中曲线仅为示意，证明不要求凸性。}
	\label{g1:fig:chord-triangle}
\end{figure}

\begin{proof}
	先取弧长参数；\(C^1\) 正则性足以使用\cref{g1:prop:arc-length-param}，旋转次数不因这种保向重参数化而改变。旋转、平移平面并移动参数起点，使 \(\gamma(0)=\gamma(L)=0\) 是曲线第二坐标的最小点，整条曲线位于 \(y\geq0\)。该点的切向量水平。正向时 \(T(0)=(1,0)\)：最低点下方属于无界外部，局部内部在切线的上侧，正向要求它位于左方。

	在闭三角形
	\[
		\Delta=\{(s,t):0\leq s\leq t\leq L\}
	\]
	上定义一个圆周值映射。对 \(s<t\)、\((s,t)\ne(0,L)\)，令
	\[
		F(s,t)=\frac{\gamma(t)-\gamma(s)}{|\gamma(t)-\gamma(s)|}.
	\]
	简单性保证分母不为零。在对角线上令 \(F(s,s)=T(s)\)，在余下角点令 \(F(0,L)=-T(0)\)。这些赋值确实连续。对角线附近，
	\[
		\frac{\gamma(t)-\gamma(s)}{t-s}
		=\frac1{t-s}\int_s^t T(v)\,\mathrm{d}v
	\]
	一致趋近相应的单位切向量；这是紧区间上 \(T\) 一致连续的直接结果。在角点 \((0,L)\) 附近，设 \(a=s\)、\(b=L-t\)，由两端相容的导数 \(T(0)=T(L)\)，
	\[
		\gamma(L-b)-\gamma(a)=-(a+b)T(0)+o(a+b).
	\]
	归一化后恰趋向 \(-T(0)\)，也包括 \(a=0\) 或 \(b=0\) 的情形。

	令 \(A=(0,0)\)、\(B=(0,L)\)、\(C=(L,L)\)。沿 \(A\to B\) 走，\(F(0,t)\) 是从原点发出的弦方向，始终在闭上半圆。因此其连续角度可直接取为 \([0,\pi]\) 中的值，从 \(0\) 变到 \(\pi\)，总增量为 \(\pi\)。沿 \(B\to C\) 走，
	\[
		F(s,L)=-F(0,s),
	\]
	故在前一角度上加 \(\pi\)，得到从 \(\pi\) 到 \(2\pi\) 的提升，总增量仍为 \(\pi\)。沿 \(C\to A\) 的对角线走，方向映射为逆向的 \(T\)，总增量是 \(-2\pi\operatorname{rot}(\gamma)\)。

	另一方面，边界环路 \(F|_{\partial\Delta}\) 的总增量必须为零。具体地，把三角形边界上的点 \(x\) 沿 \((1-u)x+uA\) 缩到 \(A\)，再复合 \(F\)，就得到固定基点 \(F(A)\) 的零伦；\cref{g1:thm:circle-fundamental}说明它的绕数为零。三段角度相接，因而
	\[
		0=\pi+\pi-2\pi\operatorname{rot}(\gamma).
	\]
	这给出正向的结论。反向时旋转次数变号；\(C^2\) 情形的曲率积分再由\eqref{g1:eq:total-signed-curvature}得到。
\end{proof}

简单性在证明中有明确的用途：它保证任意两个不同参数点之间的弦方向都有定义。自交点会在三角形内部产生未定义点，边界映射就未必能延拓到整个三角形，绕数为零的论证也随之失效。

\subsection{曲率符号、凸性及反例}
\label{g1:subsec:0703:03}

若 \(k\geq0\)，切向角不减；再要求它只转一周，就能同时控制绕行次数并判断曲线的凸性。若一条简单闭曲线的像是某个有非空内部的紧凸集的边界，就称它为\term[tuquxian]{凸曲线}[convex curve]；这里允许边界包含线段。

\begin{theorem}[曲率符号与凸性]\label{g1:thm:curvature-convexity}
	设 \(L>0\)，\(\gamma:\mathbb R/L\mathbb Z\rightarrow\mathbb R^2\) 是 \(C^2\) 弧长参数正则闭曲线，平面取标准定向，以 \(k\) 记沿参数递增方向的有向曲率。
	\begin{enumerate}
		\item 若 \(k\geq0\) 且 \(\operatorname{rot}(\gamma)=1\)，则 \(\gamma\) 自动为简单闭曲线，且按正向恰好描出一个紧凸集的边界一周。
		\item 若已知 \(\gamma\) 是正向简单闭曲线，则它为凸曲线，当且仅当 \(k\geq0\)。
	\end{enumerate}
	相应地，反向时把 \(k\geq0\)、旋转次数 \(1\) 改为 \(k\leq0\)、旋转次数 \(-1\)。
\end{theorem}
\begin{proof}
	先假设 \(k\geq0\)、旋转次数为 \(1\)。切向角提升满足 \(\theta'=k\geq0\)、\(\theta(s+L)=\theta(s)+2\pi\)。任意选一个起点，平移和旋转后可设它的参数为 \(0\)、位置为 \(0\)、切向角为 \(0\)。记 \(\gamma(s)=(x(s),y(s))\)，则
	\[
		y'(s)=\sin\theta(s),\qquad 0\leq\theta(s)\leq2\pi,\qquad y(0)=y(L)=0.
	\]
	在切向角从 \(0\) 增至 \(\pi\) 的阶段，\(y\) 不减；从 \(\pi\) 增至 \(2\pi\) 的阶段，\(y\) 不增。因此 \(y(s)\geq0\)：整条曲线在所选起点的有向切线左侧。

	还须排除中途回到起点，不能仅由半平面包含关系断言简单性。若 \(0<\theta(s)\leq\pi\)，连续角度在此前经过 \((0,\pi)\) 中的一段，故 \(y(s)=\int_0^s\sin\theta(v)\,\mathrm{d}v>0\)。若 \(\pi\leq\theta(s)<2\pi\)，则
	\[
		y(s)=-\int_s^L\sin\theta(v)\,\mathrm{d}v>0.
	\]
	余下的 \(\theta(s)=0\) 情形，由单调性知 \([0,s]\) 上 \(T=(1,0)\)，所以 \(\gamma(s)=(s,0)\ne0\)；\(\theta(s)=2\pi\) 时，同理 \(\gamma(s)=(-(L-s),0)\ne0\)，这里均取 \(0<s<L\)。起点任意，所以两个不同参数类不可能具有同一像点，\(\gamma\) 为简单闭曲线。

	设其有界内部为 \(\Omega\)，\(K=\overline\Omega\)。转角定理表明当前参数为正向，故局部内部在左侧。前两段已经证明每一点的闭左半平面
	\[
		H_s=\{p:\langle p-\gamma(s),JT(s)\rangle\geq0\}
	\]
	包含整条曲线，也包含 \(K\)。后一断言可这样看：线性函数 \(\langle p-\gamma(s),JT(s)\rangle\) 在紧集 \(K\) 上取到最小值，若它为负且不在边界取到，就会在开集 \(\Omega\) 内有局部最小值，这与其非零线性梯度矛盾。

	事实上 \(K=\bigcap_sH_s\)。为证反向包含，取 \(p\notin K\)，在 \(K\) 中取距离 \(p\) 最近的点 \(q\)。最近点必在边界上，写 \(q=\gamma(s)\)；对 \(|p-\gamma(v)|^2\) 求导可知 \(p-q\perp T(s)\)。线段 \((q,p]\) 不碰 \(K\)，否则有更近的点。若 \(p-q\) 是左法向，线段靠近 \(q\) 的一小段就位于局部内部，与此矛盾。因此 \(p-q\) 是右法向，\(p\notin H_s\)。所求等式成立，而闭半平面的交是凸集。这证明第1项。

	现在假设曲线已简单并取正向。若 \(k\geq0\)，转角定理给出旋转次数 \(1\)，第1项适用。反之，设 \(K=\overline\Omega\) 凸。先说明每条切线都是支撑线。把任一边界点移到原点，使正向切向量为 \((1,0)\)，则附近边界写成 \(y=f(x)\)，其中 \(f(0)=f'(0)=0\)，局部 \(K\) 位于图像上方。如果某个 \(p=(p_x,p_y)\in K\) 满足 \(p_y<0\)，凸性使 \(tp\in K\) 对 \(0<t<1\) 成立。但 \(f(tp_x)=o(t)\)，故充分小时 \(tp_y<f(tp_x)\)，矛盾。因此 \(K\) 在切线的闭上半平面内，特别地 \(f(x)\geq0\)。于是 \(f''(0)\geq0\)。在这一坐标和方向下有向曲率为 \(f''(0)/(1+f'(0)^2)^{3/2}\)，所以 \(k\geq0\)。边界点任意，第2项得证。
\end{proof}

在证明中，允许 \(k=0\) 至关重要：它使 \(\theta\) 可以保持不变，对应曲线中的直线段。若正向简单闭曲线处处 \(k>0\)，则每条支撑切线只接触曲线一个点。事实上，对固定起点，切向角严格增加，证明中的法向高度在其余参数处都严格为正。因此此时所围凸集严格凸，即边界不包含非退化线段。逆命题却不成立：严格凸边界的个别点仍可以曲率为零，章末的四次曲线给出例子。

\begin{example}[正曲率仍可能重复绕行]\label{g1:exm:multiple-circle-turning}
	对整数 \(m\geq2\)，考虑 \(\gamma(t)=(\cos mt,\sin mt)\)，\(t\in\mathbb R/2\pi\mathbb Z\)。它正则，速率为 \(m\)，有向曲率恒为 \(1\)，但同一圆周被走了 \(m\) 遍。其总长为 \(2\pi m\)，旋转次数为 \(m\)。所以在尚未知道简单性时，不能把 \(k>0\) 直接当作简单凸曲线的判据。
\end{example}

\begin{example}[旋转一次仍可能有自交]\label{g1:exm:rotation-one-self-intersection}
	令
	\[
		\gamma(t)=(\cos t,\sin t+2\sin3t),\qquad t\in\mathbb R/2\pi\mathbb Z.
	\]
	当 \(\sin t\ne0\) 时，第一坐标导数 \(-\sin t\) 非零；在 \(t=0,\pi\) 处，第二坐标导数分别为 \(7,-7\)，故曲线处处正则。对 \(0<t<\pi\)，切向量的第一坐标为负，其角度可以连续取在 \((\pi/2,3\pi/2)\)，并从 \(\pi/2\) 走到 \(3\pi/2\)；对 \(\pi<t<2\pi\)，第一坐标为正，角度继续从 \(3\pi/2\) 走到 \(5\pi/2\)。两段合计增量 \(2\pi\)，所以旋转次数为 \(1\)。

	然而
	\[
		\sin t+2\sin3t=\sin t\,(8\cos^2t-1).
	\]
	选 \(t_0\in(0,\pi/2)\) 满足 \(\cos t_0=1/\sqrt8\)，就有 \(\gamma(t_0)=\gamma(2\pi-t_0)\)，因而曲线不简单。它的曲率也不是同号：\(\det(\dot\gamma,\ddot\gamma)\) 在 \(t=0\) 为 \(7\)，在 \(t=\pi/2\) 为 \(-17\)。这个例子把凸性判据的两个条件分开了：整周数为 \(1\) 只控制净转角，曲率符号才控制转向是否来回反复。
\end{example}

本节积分的是有向曲率，左转和右转会抵消。下一节改看 \(\int|k|\,\mathrm{d}s\) 及其空间曲线版本：即使净转角很小，实际弯曲的总量仍可能很大。
% 纲要条目排为编号小节标题，并进入目录与书签；说明段仍保留于纲要注释中。

% !TeX root = ../../Book1.tex
% 纲要标题：### 7.4 总曲率与空间曲线
% 纲要位置：几何学计划纲要.md，第 389 行。
% 正文状态：初稿；尚未正式审核。

% 纲要内容开始
% 
% * 球面切向曲线和总曲率
% * Fenchel 不等式的几何证明
% * 等号条件与平面凸曲线
% * 结型施加的进一步限制：准确陈述，不将结理论预设为先修
% 
% ---
% 
% 纲要内容结束

\section{总曲率与空间曲线}
\label{g1:sec:0704}

平面简单闭曲线的切向量一共转过一周，但空间曲线没有一个固定平面供我们比较转角的正负。仍然可以把每个单位切向量的起点移到原点，观察它的终点在单位球面上走了多远。闭合条件将迫使这条球面曲线有足够的长度；当长度恰好达到最小值时，空间曲线反而必须落回一个平面。这个论证只需要单位切向量，不要求处处有 Frenet 法向量，因而也能处理零曲率点。

\subsection{球面切向曲线和总曲率}
\label{g1:subsec:0704:01}

\begin{definition}[总曲率与切向示性曲线]\label{g1:def:total-curvature}
	设 \(\gamma:[0,L]\to\mathbb R^3\) 是 \(C^2\) 单位速率曲线，其中 \(L>0\)，允许曲线不闭合。其单位切向量 \(T(s)=\gamma'(s)\) 给出一条 \(C^1\) 球面曲线
	\[
		T:[0,L]\longrightarrow S^2,\qquad
		S^2=\{x\in\mathbb R^3:\|x\|=1\},
	\]
	称为\term[qiexiangshixingquxian]{切向示性曲线}[tangent indicatrix]。定义 \(\gamma\) 的\term[zongqulv]{总曲率}[total curvature]为
	\begin{equation}\label{g1:eq:total-curvature}
		\operatorname{TC}(\gamma)\coloneqq\int_0^L\kappa(s)\,\mathrm{d}s
		=\int_0^L\|T'(s)\|\,\mathrm{d}s.
	\end{equation}
	\symidx{\(\operatorname{TC}(\gamma)\)：曲线的总曲率}
	对于任意紧参数区间上的 \(C^2\) 正则曲线，先按\cref{g1:prop:arc-length-param}换成弧长参数，或等价地在原参数 \(t\) 下使用 \(\int\kappa(t)\|\dot\gamma(t)\|\,\mathrm{d}t\)。
\end{definition}

右端正是切向示性曲线的长度，计入重复经过同一球面弧的次数。它并不是集合 \(T([0,L])\) 的“去重后长度”。原曲线的一段直线使 \(T\) 停在一点，不给总曲率增加长度；原曲线绕同一圆走两遍，则切向示性曲线也走两遍。由于 \(\kappa=\|T'\|\ge0\)，总曲率没有正负抵消。平面有向曲率 \(k\) 与这里的量满足 \(\kappa=|k|\)，因此 \(\int k\,\mathrm{d}s\) 和 \(\operatorname{TC}(\gamma)\) 应当区分。

\begin{example}[尺度与重复参数]\label{g1:exm:total-curvature-circle}
	半径为 \(R>0\) 的圆走一周时，\(\kappa=1/R\)，长度为 \(2\pi R\)，故总曲率是 \(2\pi\)。更一般地，对任意 \(a>0\)，伸缩后的曲线 \(a\gamma\) 的弧长元变成 \(a\,\mathrm{d}s\)，曲率变成 \(\kappa/a\)，所以总曲率不变。若 \(\gamma_m(t)=(R\cos mt,R\sin mt,0)\)，其中 \(m\) 是非零整数且 \(0\le t\le2\pi\)，则总曲率为 \(2\pi|m|\)。同一个圆形像集并未确定总曲率；走过它的次数仍然是参数曲线的一部分。
\end{example}

闭曲线还满足一个比 \(T(0)=T(L)\) 更有用的约束：
\begin{equation}\label{g1:eq:tangent-balance}
	\int_0^L T(s)\,\mathrm{d}s=\gamma(L)-\gamma(0)=0.
\end{equation}
它表示各方向的位移必须抵消。例如，若存在非零向量 \(a\)，使 \(a\cdot T(s)>0\) 对所有 \(s\) 成立，那么整条曲线始终向 \(a\) 方向前进，不可能闭合。稍后还需要边界情形：若 \(a\cdot T(s)\ge0\) 处处成立，则连续性和\eqref{g1:eq:tangent-balance}迫使 \(a\cdot T(s)=0\) 处处成立，此时 \(\gamma\) 落在一个垂直于 \(a\) 的平面内。

\subsection{Fenchel 不等式的几何证明}
\label{g1:subsec:0704:02}

我们把球面上两点 \(p,q\) 的夹角记为 \(d_S(p,q)=\arccos(p\cdot q)\in[0,\pi]\)。下面证明，在球面上绕行不能比这段角距离更短。这也是下一小节判断等号的工具。

\begin{lemma}[球面弧长与端点夹角]\label{g1:lem:spherical-arc-distance}
	设 \(c:[a,b]\to S^2\) 是 \(C^1\) 曲线，允许 \(c'\) 为零。则
	\[
		d_S(c(a),c(b))\le\int_a^b\|c'(t)\|\,\mathrm{d}t.
	\]
	若 \(c(a)=p\)、\(c(b)=-p\)，且右端等于 \(\pi\)，则存在一个与 \(p\) 垂直的单位向量 \(v\)，使
	\[
		c(t)=p\cos u(t)+v\sin u(t),\qquad
		u(t)=\int_a^t\|c'(r)\|\,\mathrm{d}r.
	\]
	因而该弧沿一条半大圆单调前进，可以停顿，但不能折返。
\end{lemma}

\begin{proof}
	记 \(p=c(a)\)。由于 \(c\cdot c'=0\)，将 \(p\) 分解到 \(c(t)\) 方向及其垂直平面，得到
	\[
		|p\cdot c'(t)|\le\sqrt{1-(p\cdot c(t))^2}\,\|c'(t)\|.
	\]
	直接对 \(\arccos(p\cdot c)\) 求导会在 \(c=\pm p\) 处遇到分母为零。为保留这些端点，先取 \(0<\varepsilon<1\)，令
	\[
		\theta_\varepsilon(t)=\arccos\bigl((1-\varepsilon)p\cdot c(t)\bigr).
	\]
	此时分母严格为正，而且
	\[
		|\theta_\varepsilon'(t)|
		\le
		\frac{(1-\varepsilon)\sqrt{1-(p\cdot c(t))^2}}
		{\sqrt{1-(1-\varepsilon)^2(p\cdot c(t))^2}}\,
		\|c'(t)\|
		\le\|c'(t)\|.
	\]
	积分并令 \(\varepsilon\to0\)，便得端点夹角的估计。相同论证适用于曲线的任意子区间。

	现在设端点为 \(p,-p\)，总长恰为 \(\pi\)。令 \(\theta(t)=\arccos(p\cdot c(t))\)。分别从两端使用刚证的估计，得到
	\[
		\theta(t)\le u(t),\qquad
		\pi-\theta(t)\le\pi-u(t),
	\]
	所以 \(\theta(t)=u(t)\)。在 \(0<u(t)<\pi\) 的区间上，定义
	\[
		v(t)=\frac{c(t)-p\cos u(t)}{\sin u(t)}.
	\]
	它是 \(C^1\) 单位向量，且 \(v(t)\perp p\)。微分上式所给的 \(c=p\cos u+v\sin u\)，利用 \(v'\perp p,v\)，可得
	\[
		\|c'\|^2=(u')^2+\sin^2u\,\|v'\|^2.
	\]
	定义给出 \(u'=\|c'\|\)，故 \(v'=0\)。因为 \(u\) 单调，\(0<u<\pi\) 的原像是一个区间，因而这里的 \(v\) 是同一个常向量。剩余端段上 \(u=0\) 或 \(u=\pi\)，对应 \(c=p\) 或 \(c=-p\)，公式仍然成立。
\end{proof}

\begin{lemma}[短球面闭曲线的半球约束]\label{g1:lem:short-spherical-loop}
	设 \(c\) 是一条 \(C^1\) 球面闭曲线，将它按经过顺序分成长度相等的两段，分点为 \(p,q\)。若总长度小于 \(2\pi\)，则
	\[
		(p+q)\cdot c(t)>0\qquad\text{对所有 }t.
	\]
	若总长度等于 \(2\pi\) 且 \(q\ne-p\)，则上式改为非负不等式，并且在分点 \(p\) 处严格为正。
\end{lemma}

\begin{proof}
	累计长度是连续函数，所以可以取到半长分点；长度为零时任取分点即可。任一点 \(x=c(t)\) 位于两段之一。从该段的两端走到 \(x\)，由\cref{g1:lem:spherical-arc-distance}有
	\[
		\alpha+\beta\le\frac{\operatorname{length}(c)}2,\qquad
		\alpha=d_S(p,x),\quad\beta=d_S(q,x).
	\]
	而
	\[
		(p+q)\cdot x=\cos\alpha+\cos\beta
		=2\cos\frac{\alpha+\beta}{2}\cos\frac{\alpha-\beta}{2}.
	\]
	若总长小于 \(2\pi\)，则 \(\alpha+\beta<\pi\)，两个余弦都严格为正。若总长等于 \(2\pi\)，则两者都非负；在 \(x=p\) 处，该内积为 \(1+p\cdot q>0\)，最后一个严格不等式正用了 \(q\ne-p\)。
\end{proof}

这个引理把“转得太少”变成了“所有方向都偏向同一侧”。其中 \(p+q\) 指向夹角的中间方向；两段都不足半个大圆时，它们就不能绕到这个方向的背后。

\begin{theorem}[Fenchel 不等式]\label{g1:thm:fenchel}
	设 \(\gamma:\mathbb R/L\mathbb Z\to\mathbb R^3\) 是 \(C^2\) 正则闭曲线，按弧长参数化，\(L>0\)。则
	\[
		\operatorname{TC}(\gamma)=\int_0^L\kappa(s)\,\mathrm{d}s\ge2\pi.
	\]
	这里不假设 \(\gamma\) 简单，也不假设 \(\kappa>0\)。
\end{theorem}

\begin{proof}
	若总曲率小于 \(2\pi\)，切向示性曲线 \(T\) 的长度也小于 \(2\pi\)。对它应用\cref{g1:lem:short-spherical-loop}，取得两个半长分点 \(p,q\)，使 \((p+q)\cdot T(s)>0\) 处处成立。于是
	\[
		(p+q)\cdot\int_0^L T(s)\,\mathrm{d}s>0,
	\]
	这与闭合约束\eqref{g1:eq:tangent-balance}矛盾。
\end{proof}

这一结论通常称为\term[fenchelbudengshi]{Fenchel 不等式}[Fenchel's inequality]，包括等号条件的整体陈述也称为 Fenchel 定理。\footnote{Fenchel，芬切尔，Werner Fenchel（1905--1988），生于德国、后在丹麦工作的数学家，研究凸性理论及非欧几何中的离散群。}原曲线即使在空间里不断改变弯曲平面，也不能把闭合所需的总转向减少到一周以下。文献中的另一种证明从内接多边形的外角和出发，见\cite[Theorem 3.4]{Milnor1950}；这里的球面证明则使零曲率段的作用直接可见。

\subsection{等号条件与平面凸曲线}
\label{g1:subsec:0704:03}

不等式的证明很短，但等号还要回答三个不同的问题：切向量是否在固定平面内，它沿这个平面的单位圆是否折返，以及原曲线是否可能重叠或自交。以下论证依次排除这些可能。

\begin{theorem}[Fenchel 不等式的等号条件]\label{g1:thm:fenchel-equality}
	设 \(\gamma:\mathbb R/L\mathbb Z\to\mathbb R^3\) 是 \(C^2\) 单位速率闭曲线，\(L>0\)。则 \(\operatorname{TC}(\gamma)=2\pi\) 当且仅当 \(\gamma\) 简单，其像是某个仿射平面内具有非空内部的紧凸集的边界，并且参数将该边界走一周。允许曲率在点或区间上为零。
\end{theorem}

\begin{proof}
	先设总曲率为 \(2\pi\)，记累计曲率
	\[
		u(s)=\int_0^s\kappa(r)\,\mathrm{d}r,\qquad 0\le s\le L.
	\]
	把 \(u\) 延拓为满足 \(u(s+L)=u(s)+2\pi\) 的函数。任取 \(s<t\le s+L\)，使 \(u(t)-u(s)=\pi\)。这两个时刻把切向示性曲线分成长度均为 \(\pi\) 的两段。令 \(p=T(s)\)、\(q=T(t)\)。如果 \(q\ne-p\)，\cref{g1:lem:short-spherical-loop}给出
	\[
		(p+q)\cdot T(r)\ge0\quad\text{处处成立},
		\qquad (p+q)\cdot T(s)>0.
	\]
	连续性使严格正值保持在一个参数小区间上，因此其积分严格为正，仍与\eqref{g1:eq:tangent-balance}矛盾。所以任何这样的一对半长分点都互为对径点。

	取 \(p=T(0)\)，并取 \(s_1,s_2,s_3\)，使 \(u(s_j)=j\pi/2\)。第一半段从 \(p\) 到 \(-p\)，长度为 \(\pi\)，故\cref{g1:lem:spherical-arc-distance}给出一个单位向量 \(v\perp p\)，使
	\[
		T(s)=p\cos u(s)+v\sin u(s)\qquad(0\le s\le s_2).
	\]
	特别地 \(T(s_1)=v\)。第二半段也必须沿某条半大圆前进，可以写成
	\[
		T(s)=-p\cos(u(s)-\pi)+w\sin(u(s)-\pi)\qquad(s_2\le s\le L),
	\]
	其中 \(w\perp p\) 是单位向量。由于 \(s_1,s_3\) 又是一对半长分点，\(w=T(s_3)=-T(s_1)=-v\)。因此两段实际合成同一个大圆，而且
	\begin{equation}\label{g1:eq:fenchel-equality-tangent}
		T(s)=p\cos u(s)+v\sin u(s)\qquad(0\le s\le L).
	\end{equation}

	积分表明 \(\gamma(s)-\gamma(0)\) 始终属于平面 \(\operatorname{span}\{p,v\}\)。在以 \(p,v\) 为正向正交基的这个平面里，\(u(s)\) 正是一个连续切角；它单调不减，总增量为 \(2\pi\)。\cref{g1:thm:curvature-convexity}的充分性方向适用于闭正则曲线，并不预设简单性：切角单调增加一周时，所有切线都是同侧支撑线，而且曲线没有重复点。因此 \(\gamma\) 简单，并且沿一个紧凸集的边界走一周。

	反过来，设 \(\gamma\) 简单且是平面紧凸集的边界。选择使凸集位于左侧的方向，由\cref{g1:thm:curvature-convexity}，其有向曲率 \(k\ge0\)；由\cref{g1:thm:turning-tangent}，\(\int_0^L k\,\mathrm{d}s=2\pi\)。于是 \(\operatorname{TC}(\gamma)=\int_0^L|k|\,\mathrm{d}s=2\pi\)。反向行走时 \(k\) 改变符号，总曲率不变。
\end{proof}

\input{Figures/Chapter07-TotalCurvature.tex}

证明中的 \(u\) 是累计曲率，未必是合法的正则重参数：在直线段上 \(u'=0\)，甚至可能一段时间不变化。我们只把它代入\eqref{g1:eq:fenchel-equality-tangent}，没有除以 \(\kappa\)，也没有要求 \(u\) 可逆。\cref{g1:fig:total-curvature}所画的圆周运动可以在个别方向上停顿；因此等号条件是“平面凸曲线”，不是“圆”，也不是“处处严格正曲率的平面曲线”。

\begin{corollary}[非平面性的严格代价]\label{g1:cor:nonplanar-total-curvature}
	任何像不包含在一个仿射平面内的 \(C^2\) 正则闭空间曲线，都有 \(\operatorname{TC}(\gamma)>2\pi\)。任何自交的 \(C^2\) 正则闭空间曲线也有同样的严格不等式。
\end{corollary}

\begin{proof}
	\cref{g1:thm:fenchel}给出非严格下界，而\cref{g1:thm:fenchel-equality}排除了这两种情形中的等号。
\end{proof}

\subsection{结型施加的进一步限制：准确陈述，不将结理论预设为先修}
\label{g1:subsec:0704:04}

非平面性取决于当前摆放方式。一个平面圆被轻微拉出平面后，总曲率已经严格大于 \(2\pi\)，但仍然可以完整地拉回原来的圆，不必让曲线穿过自身。结型问的是后一类变形能否做到，其限制比“目前不在平面内”更强。

\begin{definition}[结与平凡结]\label{g1:def:knot-unknot}
	本小节的\term[jie]{结}[knot]指 \(\mathbb R^3\) 中 \(C^2\) 正则嵌入圆周的像。两个结 \(K_0,K_1\) 具有相同\term[jiexing]{结型}[knot type]，是指存在 \(\mathbb R^3\) 的环境同痕 \(H_t\)，使 \(H_1(K_0)=K_1\)。这里沿用\cref{g1:def:ambient-isotopy}的定义：\(H_0=\mathrm{id}\)，且 \((x,t)\mapsto(H_t(x),t)\) 是 \(\mathbb R^3\times[0,1]\) 的同胚。与平面圆有相同结型的结称为\term[pingfanjie]{平凡结}[unknot]，否则称为非平凡结。
\end{definition}

环境同痕移动的是整个周围空间，每一时刻都保持可逆，因而曲线上不同的点不能在变形途中穿过彼此。一般的道路同伦没有这个限制；在 \(\mathbb R^3\) 中任意闭路都可以连续收缩，并不能据此说所有结都平凡。结型只保留嵌入的拓扑性质，不要求变形保持长度或曲率。

\begin{theorem}[Fáry--Milnor 定理]\label{g1:thm:fary-milnor}
	设 \(\gamma:S^1\to\mathbb R^3\) 是 \(C^2\) 正则嵌入，且它的像是非平凡结。则
	\[
		\operatorname{TC}(\gamma)>4\pi.
	\]
	等价地，若一个 \(C^2\) 正则嵌入圆周的总曲率不超过 \(4\pi\)，它必为平凡结。
\end{theorem}

这就是\term[farymilnordingli]{Fáry--Milnor 定理}[Fáry--Milnor theorem]。\footnote{Fáry，法里，匈牙利姓名为 Fáry István（1922--1984），匈牙利出生的数学家，研究几何与代数拓扑，证明了平面图的直线嵌入定理。Milnor，米尔诺，John Willard Milnor（1931年生），美国数学家，在微分拓扑、代数 \(K\) 理论和动力系统中有重要工作。}本章引用这一结论，不展开证明。所用光滑版本可由\cite[Theorem 2.2; Corollary 4.6]{Milnor1950}直接得到：前一处将原文以多边形定义的总曲率与这里的积分相等同，后一处给出严格的 \(4\pi\) 下界。其证明还需要把有限总曲率的结在环境同痕下替换为适当的内接多边形，并研究各投影方向的高度极值数；这些结论在该文的 Sections 3--4 中建立。

\begin{example}[很大的总曲率并不保证打结]\label{g1:exm:unknot-large-curvature}
	考虑平面极坐标曲线
	\[
		\gamma(t)=r(t)(\cos t,\sin t,0),\qquad
		r(t)=1+\tfrac12\cos4t,\qquad 0\le t\le2\pi.
	\]
	因为 \(r\ge1/2\)，每个极角只有一个点，而且 \(\|\dot\gamma(t)\|^2=r'(t)^2+r(t)^2>0\)，所以这是 \(C^\infty\) 正则简单闭曲线。它是平凡结：在柱坐标 \((\rho,\phi,z)\) 中，
	\[
		H_a(\rho,\phi,z)=
		\left(\frac{\rho}{1+(a/2)\cos4\phi},\phi,z\right),
		\qquad 0\le a\le1,
	\]
	给出把它移到单位圆的环境同痕。轴上令 \(H_a(0,0,z)=(0,0,z)\)；分母一致介于 \(1/2\) 与 \(3/2\) 之间，因此映射及逆映射在轴上也连续。

	它的总曲率却大于 \(4\pi\)。把平面视为复平面，有 \(\dot\gamma(t)=e^{it}(r'(t)+ir(t))\)，所以可取切角
	\[
		\theta(t)=t+\frac\pi2-\arctan\frac{r'(t)}{r(t)}.
	\]
	在 \(t_j=(2j+1)\pi/8\) 处，\(r=1\)、\(r'=-2(-1)^j\)，故
	\[
		\theta(t_j)=t_j+\frac\pi2+(-1)^j A,\qquad A=\arctan2.
	\]
	把 \(j=0,\ldots,8\) 的相邻切角差取绝对值并求和，八项按奇偶配成四对，每对为 \(4A\)。由于总变差至少等于任何分割上的变差和，
	\[
		\operatorname{TC}(\gamma)=\int_{t_0}^{t_0+2\pi}|\theta'(t)|\,\mathrm{d}t
		\ge16\arctan2>4\pi.
	\]
	所以 Fáry--Milnor 定理给出的是非平凡结的必要曲率代价，不能反过来凭一个大总曲率数值断定曲线打结。
\end{example}
% 纲要条目排为编号小节标题，并进入目录与书签；说明段仍保留于纲要注释中。


\begin{chaptersummary}
	弧长把行走速度从几何数据中消去。平面有向曲率 \(k\) 是切向角的导数，空间非负曲率 \(\kappa\) 是单位切向量的变化率；在 \(\kappa>0\) 时，挠率与 Frenet 方程继续记录密切平面的转动。零曲率不等于曲线失去正则性，平行法标架可以在这样的点继续，而 Frenet 主法向一般不能连续延拓。

	\ProseParagraphBreak

	斜对称的标架方程保持正交性，积分其切向量便恢复曲线。给定同一弧长参数下的正曲率与挠率，曲线在保向刚性运动意义下唯一；平面则使用允许变号或为零的有向曲率。空间反射使挠率变号，参数反向不使挠率变号。周期数据只保证重复的刚性关系，真正闭合还要求标架和总位移同时相合。

	\ProseParagraphBreak

	平面闭曲线的有向曲率积分等于 \(2\pi\) 倍的旋转次数。简单性使全部弦方向能在参数三角形上连续组织起来，从而迫使转角为正或负一周。对正向简单闭曲线，曲率非负恰好刻画凸性；在尚未知道简单性时，非负曲率还需与旋转次数为 \(1\) 合用。严格凸允许个别零曲率点，正曲率也不能排除多次绕行。

	\ProseParagraphBreak

	总曲率计算切向量在单位球面上走过的总长度，没有有向转角的正负抵消。闭合使所有切向量的弧长平均为零，由此得到 Fenchel 不等式 \(\int\kappa\,\mathrm ds\ge2\pi\)，等号恰对应一周的平面凸简单闭曲线。非平凡嵌入结的总曲率还必须严格大于 \(4\pi\)；较大的总曲率本身并不保证打结。局部不变量、闭合与嵌入类型因而承担不同的几何信息。
\end{chaptersummary}

\BookExercises

\begin{exercise}\label{g1:ex:curve-cusp-regularity}
	\begin{enumerate}
		\item 说明 \(t\mapsto(t^3,0)\) 在原点不正则，但其轨迹可以正则参数化。
		\item 证明半立方抛物线 \(C=\{(x,y):y^2=x^3,\ x\ge0\}\) 不能在原点附近用经过原点的正则曲线参数化。这里要求原点对应参数区间的一个内点。
	\end{enumerate}
\end{exercise}
\begin{solution}
	(1) 原参数的导数是 \((3t^2,0)\)，在 \(0\) 为零；轨迹却是整个横轴，由 \(s\mapsto(s,0)\) 正则参数化。替换 \(s=t^3\) 不是原点附近的微分同胚，所以这并不与正则性在允许的重参数化下保持相矛盾。

	(2) 假设有 \(C^1\) 正则曲线 \(\beta(u)=(x(u),y(u))\) 在内点 \(u=0\) 经过原点，且像在 \(C\) 中。由于 \(x(u)\ge0=x(0)\)，内点极小值给出 \(x'(0)=0\)，故 \(x(u)=o(|u|)\)。曲线上 \(|y(u)|=x(u)^{3/2}\)，于是
	\[
		\frac{|y(u)|}{|u|}=\frac{x(u)}{|u|}\sqrt{x(u)}\longrightarrow0.
	\]
	因此 \(y'(0)=0\)，与 \(\beta'(0)\ne0\) 矛盾。这个障碍来自轨迹在原点的形状；仅换一个更快或更慢的参数不能解决。
\end{solution}

\begin{exercise}\label{g1:ex:logarithmic-spiral}
	设 \(a>0\)，\(\alpha(t)=e^{at}(\cos t,\sin t)\)。求速度、有向曲率及从 \(t=0\) 起算的弧长 \(s(t)\)，把曲率改写成 \(s\) 的函数，并计算参数从 \(0\) 到 \(T>0\) 时的总曲率。
\end{exercise}
\begin{solution}
	令 \(c=\sqrt{1+a^2}\)。以径向单位向量 \(e_r=(\cos t,\sin t)\) 及 \(e_\theta=(-\sin t,\cos t)\) 为基，
	\[
		\dot\alpha=e^{at}(ae_r+e_\theta),\qquad
		\ddot\alpha=e^{at}((a^2-1)e_r+2ae_\theta).
	\]
	所以速度为 \(ce^{at}\)，行列式为 \(e^{2at}(a^2+1)\)，得到
	\[
		k(t)=\frac{e^{-at}}c,\qquad
		s(t)=\frac ca(e^{at}-1),\qquad
		k(s)=\frac1{as+c}.
	\]
	弧长参数的像为 \((-c/a,\infty)\)，分母在这个区间内为正。曲率没有负号，故总曲率为
	\[
		\int_0^T\kappa(t)\lVert\dot\alpha(t)\rVert\,\mathrm dt
		=\int_0^T\frac{e^{-at}}c\,ce^{at}\,\mathrm dt=T.
	\]
	这也可以从 \(\dot\alpha\) 的方向看出：向量 \(ae_r+e_\theta\) 与 \(e_r\) 成一个固定角，所以切向角随 \(t\) 增加同样的量。半径增长会使局部曲率变小，却不减少每一整周的总转角。
\end{solution}

\begin{exercise}\label{g1:ex:quartic-zero-curvature}
	考虑空间中的平面曲线 \(\alpha(t)=(t,t^4,0)\)。求其曲率，说明 Frenet 主法向在 \(t\ne0\) 时的表达式能否连续延拓到 \(0\)，并与正文中的 \(y=x^3\) 比较。能否因此把一般曲线在零曲率点的主法向仍定义为 \(T'/\kappa\)？
\end{exercise}
\begin{solution}
	速度为 \(v=\sqrt{1+16t^6}\)，有向曲率为
	\[
		k(t)=\frac{12t^2}{(1+16t^6)^{3/2}}\ge0.
	\]
	所以 \(\kappa=k\)，仅在原点为零。在 \(t\ne0\) 时，
	\[
		N(t)=\frac{(-4t^3,1,0)}{\sqrt{1+16t^6}},\qquad B(t)=(0,0,1).
	\]
	它们可以连续延拓为 \(N(0)=(0,1,0)\)、\(B(0)=(0,0,1)\)。与 \(y=x^3\) 的区别在于这里 \(k\) 在两侧均为正，没有翻转弯曲侧。

	然而 \(T'/\kappa\) 在原点仍是 \(0/0\)，原公式没有定义。此例只能说明某条具体曲线有可用的延拓，不能保证一般曲线都有这种延拓，也不能使所有零曲率点获得统一的主法向定义。平行法标架的存在则不需要检查这类比值的极限。
\end{solution}

\begin{exercise}\label{g1:ex:rotating-normal-frame}
	沿单位圆 \(\gamma(s)=(\cos s,\sin s,0)\)，令
	\[
		E_1=(-\cos s,-\sin s,0),\qquad E_2=(0,0,1).
	\]
	对 \(C^1\) 函数 \(\phi\)，令 \(\widetilde E_1=\cos\phi\,E_1+\sin\phi\,E_2\)、\(\widetilde E_2=-\sin\phi\,E_1+\cos\phi\,E_2\)。证明原法标架平行，计算新标架的导数，并判断它何时平行。解释为什么新标架中出现的法平面旋转速度不能直接叫作曲线的挠率。
\end{exercise}
\begin{solution}
	这里 \(T=(-\sin s,\cos s,0)\)，\(T'=E_1\)、\(E_1'=-T\)、\(E_2'=0\)，所以原标架满足平行法标架方程。求导得到
	\[
		\widetilde E_1'=-\cos\phi\,T+\phi'\widetilde E_2,\qquad
		\widetilde E_2'=\sin\phi\,T-\phi'\widetilde E_1.
	\]
	两个导数都平行于 \(T\) 当且仅当 \(\phi'=0\)，也就是 \(\phi\) 在区间上常值。圆的 Frenet 副法向 \(B=(0,0,1)\) 恒定，挠率始终为零；而 \(\phi'\) 是人为旋转所选法标架产生的系数，可以任意指定。只有选定正文规定的 Frenet 标架后，相应法向转动系数才是那里的 \(\tau\)。
\end{solution}

\begin{exercise}\label{g1:ex:curvature-tangent-estimate}
	设 \(\gamma:[a,b]\to\mathbb R^n\) 是 \(C^2\) 单位速率曲线，\(\lVert\gamma''(s)\rVert\le K\)。对 \(s_0,s_0+h\in[a,b]\)，证明
	\[
		\lVert T(s_0+h)-T(s_0)\rVert\le K|h|,\qquad
		\lVert\gamma(s_0+h)-\gamma(s_0)-hT(s_0)\rVert\le\tfrac12Kh^2.
	\]
	由此给出两端弦长的上下界，并说明 \(K=0\) 时曲线是什么。
\end{exercise}
\begin{solution}
	由 \(T'=\gamma''\) 和微积分基本定理，
	\[
		T(s_0+u)-T(s_0)=\int_0^u\gamma''(s_0+v)\,\mathrm dv,
	\]
	取范数便得第一个估计。再积分一次：
	\[
		\gamma(s_0+h)-\gamma(s_0)-hT(s_0)=\int_0^h(T(s_0+u)-T(s_0))\,\mathrm du.
	\]
	当 \(h<0\) 时反转积分端点，同样得到范数不超过 \(\int_0^{|h|}Ku\,\mathrm du=Kh^2/2\)。三角不等式及反三角不等式于是给出
	\[
		\max\{0,|h|-\tfrac12Kh^2\}\le\lVert\gamma(s_0+h)-\gamma(s_0)\rVert\le|h|.
	\]
	上界直接使用单位速率和长度公式，比分解为切线近似得到的 \(|h|+Kh^2/2\) 更强。若 \(K=0\)，则 \(T\) 恒定，\(\gamma(s)=\gamma(s_0)+(s-s_0)T(s_0)\)，是直线的单位速率参数化。
\end{solution}

\begin{exercise}
	\label{g1:ex:curve-data-scaling}
	设 \(\gamma\colon I\rightarrow\mathbb R^3\) 是 \(C^3\) 单位速率曲线，曲率 \(\kappa>0\)，挠率为 \(\tau\)。给定 \(a>0\)，定义 \(\widetilde\gamma(s)=a\gamma(s/a)\)，其中 \(s\in aI\)。求其曲率和挠率，并说明为什么相似变换通常不保留曲率数据。
\end{exercise}
\begin{solution}
	求导得 \(\widetilde\gamma'(s)=\gamma'(s/a)=T(s/a)\)，所以它仍以单位速率参数化。又有
	\[
		\widetilde T'(s)=\frac{\kappa(s/a)}aN(s/a),
	\]
	故 \(\widetilde\kappa(s)=\kappa(s/a)/a\)，\(\widetilde N(s)=N(s/a)\)，\(\widetilde B(s)=B(s/a)\)。因此
	\[
		\widetilde B'(s)=-\frac{\tau(s/a)}a\widetilde N(s),\qquad
		\widetilde\tau(s)=\frac{\tau(s/a)}a.
	\]
	长度扩大 \(a\) 倍，单位弧长上的弯曲与扭转相应缩小 \(a\) 倍，并且参数本身也被缩放。刚性唯一性要求同一弧长参数下数据相同，所以不能把允许缩放的相似等价直接代入该结论。
\end{solution}

\begin{exercise}
	\label{g1:ex:zero-torsion-planarity}
	设 \(\gamma\colon I\rightarrow\mathbb R^3\) 是连通区间上的 \(C^3\) 单位速率曲线，且 \(\kappa>0\)。证明 \(\tau\equiv0\) 当且仅当曲线包含在某个仿射平面中。说明哪里用到了 \(\kappa>0\)。
\end{exercise}
\begin{solution}
	若 \(\tau=0\)，则 Frenet 方程给出 \(B'=0\)，故 \(B=B_0\) 为固定单位向量。对任意固定 \(s_0\)，有
	\[
		\frac{\mathrm{d}}{\mathrm{d}s}\langle\gamma(s)-\gamma(s_0),B_0\rangle
		=\langle T(s),B(s)\rangle=0.
	\]
	起点处内积为零，所以曲线位于法向量为 \(B_0\) 且过 \(\gamma(s_0)\) 的平面。

	反之，若曲线位于法向量为固定单位向量 \(n\) 的仿射平面，求导得 \(T\perp n\) 和 \(T'=\kappa N\perp n\)。因为 \(\kappa>0\)，所以 \(N\perp n\)，从而 \(B=T\times N\) 只能为 \(n\) 或 \(-n\)。\(B\) 连续且区间连通，因此符号不变，\(B'=0\)，即 \(\tau=0\)。正曲率使 Frenet 法向量可定义，并使从 \(T'\perp n\) 推出 \(N\perp n\) 合法；曲率为零的直线虽平面，却没有本章的 Frenet 挠率。
\end{solution}

\begin{exercise}
	\label{g1:ex:half-period-curvature}
	设 \(L>0\)，连续函数 \(k\colon\mathbb R\rightarrow\mathbb R\) 满足
	\[
		k(s+L/2)=k(s),\qquad \int_0^{L/2}k(s)\,\mathrm{d}s=\pi.
	\]
	证明由 \(k\) 重建的任意单位速率平面曲线均以 \(L\) 为周期。给出一个非常数、处处正的具体 \(k\)，并说明为何无需显式算出位置积分。
\end{exercise}
\begin{solution}
	因为 \(k\) 以 \(L/2\) 为周期，函数 \(\int_s^{s+L/2}k(t)\,\mathrm{d}t\) 的导数为 \(k(s+L/2)-k(s)=0\)，故其恒等于 \(\pi\)。于是转角满足 \(\theta(s+L/2)=\theta(s)+\pi\)，从而 \(T(s+L/2)=-T(s)\)。特别地 \(T\) 以 \(L\) 为周期，而且
	\[
		\int_0^L T(s)\,\mathrm{d}s
		=\int_0^{L/2}\bigl(T(s)+T(s+L/2)\bigr)\,\mathrm{d}s=0.
	\]
	由\eqref{g1:eq:plane-closure-conditions}，曲线以 \(L\) 为周期。

	例如取
	\[
		k(s)=\frac{2\pi}{L}\left(1+\frac12\cos\frac{4\pi s}{L}\right).
	\]
	它处处大于零，具有半周期 \(L/2\)，而余弦项在半周期上的积分为零。两半段切向量逐点相反，已经直接证明总位移抵消，所以不需要求出位置积分的初等表达式。
\end{solution}

\begin{exercise}
	\label{g1:ex:helix-mirror}
	设 \(\gamma\) 是\cref{g1:exm:constant-curvature-helix}中 \(\tau_0\ne0\) 的单位速率螺线，\(R(x,y,z)=(x,y,-z)\)。证明不存在 \(Q\in \operatorname{SO}(3)\)、\(b\in\mathbb R^3\) 和常数 \(c\in\mathbb R\)，使 \(R\gamma(s)=Q\gamma(s+c)+b\) 对所有 \(s\) 成立。若把右端改为 \(Q\gamma(c-s)+b\)，结论是否改变？
\end{exercise}
\begin{solution}
	反射 \(R\) 的行列式为 \(-1\)，所以 \(R\gamma\) 的曲率仍为 \(\kappa_0\)，挠率为 \(-\tau_0\)。右端 \(Q\gamma(s+c)+b\) 的挠率为 \(\tau_0\)：平移弧长原点不改变常函数，且保向等距变换不改变挠率。若两条参数曲线相同，它们的挠率必须相同，于是 \(-\tau_0=\tau_0\)，与假设矛盾。

	参数反向也不改变空间曲线的挠率。直接核对时，\(\widehat T(s)=-T(c-s)\)、\(\widehat N(s)=N(c-s)\)、\(\widehat B(s)=-B(c-s)\)，故
	\[
		\widehat B'(s)=B'(c-s)=-\tau_0\widehat N(s).
	\]
	所以右端改成 \(Q\gamma(c-s)+b\) 后，挠率仍为 \(\tau_0\)，同样不可能。这里正是环境反射和参数反向的区别。
\end{solution}

\begin{exercise}\label{g1:ex:affine-rotation}
	设 \(\gamma:\mathbb R/P\mathbb Z\rightarrow\mathbb R^2\) 是 \(C^1\) 正则闭曲线，\(A\in GL(2,\mathbb R)\)。证明
	\[
		\operatorname{rot}(A\gamma)=\operatorname{sgn}(\det A)\operatorname{rot}(\gamma).
	\]
	再对 \(a,b>0\) 的椭圆 \(\eta(t)=(a\cos t,b\sin t)\) 算出有向曲率，并求其总有向曲率。
\end{exercise}
\begin{solution}
	对 \(A\) 的两个列向量作 Gram--Schmidt 正交化，可写 \(A=QR\)，其中 \(Q\) 正交，\(R\) 为对角元均为正的上三角矩阵。令 \(R_u=(1-u)R+uI\)，其对角元仍均为正，故始终可逆。于是 \(QR_u\gamma\) 是一族正则闭曲线：导数 \(QR_u\dot\gamma\) 不会为零，且连续依赖 \((t,u)\)。由\cref{g1:prop:rotation-invariance}，\(\operatorname{rot}(A\gamma)=\operatorname{rot}(Q\gamma)\)。二维正交变换按行列式符号分为旋转与反射后旋转，又 \(\det R>0\)，所以得到所求公式。

	椭圆的速率为 \((a^2\sin^2t+b^2\cos^2t)^{1/2}\)，且 \(\det(\dot\eta,\ddot\eta)=ab\)，因而
	\[
		k(t)=\frac{ab}{(a^2\sin^2t+b^2\cos^2t)^{3/2}}>0.
	\]
	它是单位圆经 \(\operatorname{diag}(a,b)\) 的像，旋转次数为 \(1\)，故总有向曲率为 \(2\pi\)。这里无需先用椭圆积分求出弧长：\eqref{g1:eq:total-signed-curvature}计算的是整周角度增量。
\end{solution}

\begin{exercise}\label{g1:ex:figure-eight-rotation}
	对 \(\gamma(t)=(\sin t,\sin2t)\)，\(t\in\mathbb R/2\pi\mathbb Z\)，证明它处处正则、在原点具有横向自交，并求旋转次数。说明为何不能通过切向量始终非零的连续变形，把它变成只走一周的圆。
\end{exercise}
\begin{solution}
	导数为 \((\cos t,2\cos2t)\)。若第一分量为零，则 \(t=\pi/2\) 或 \(3\pi/2\)，第二分量为 \(-2\)，因此导数从不为零。参数 \(0,\pi\) 都映到原点，对应切向量为 \((1,2)\)、\((-1,2)\)，其行列式为 \(4\ne0\)，故两支在原点横向相交。

	单位切向量满足 \(T(2\pi-t)=T(t)\)。这意味着 \(t\in[\pi,2\pi]\) 时，切向映射恰好把 \(t\in[0,\pi]\) 时的道路反向走一遍；两段的角度增量互为相反数。所以 \(\operatorname{rot}(\gamma)=0\)。圆的一周参数无论选哪种方向，旋转次数都是 \(1\) 或 \(-1\)，与零不同。\cref{g1:prop:rotation-invariance}排除了题设的正则变形；若存在某种连续变形，在这一过程中必有切向量为零或不再满足该命题的连续正则性条件之处。
\end{solution}

\begin{exercise}\label{g1:ex:superellipse-zero-curvature}
	考虑闭曲线
	\[
		x^4+y^4=1.
	\]
	证明它有光滑正则的正向简单闭参数，所围区域严格凸，而曲率在四个坐标轴交点处为零。因此严格凸性不蕴含曲率处处为正。
\end{exercise}
\begin{solution}
	取
	\[
		r(t)=(\cos^4t+\sin^4t)^{-1/4},\qquad
		\gamma(t)=r(t)(\cos t,\sin t).
	\]
	分母处处为正，所以参数光滑。又 \(\det(\gamma,\dot\gamma)=r(t)^2>0\)，故 \(\dot\gamma\ne0\)；每条从原点发出的射线恰遇曲线一次，所以它在参数圆周上单射，并按极角递增的正向走一周。

	函数 \(x\mapsto x^4\) 严格凸，因为其导数 \(4x^3\) 严格递增。因此 \(F(x,y)=x^4+y^4\) 严格凸：两点不同，至少一个坐标不同，该坐标的凸性不等式便严格。若 \(p\ne q\) 且 \(F(p),F(q)\leq1\)，那么对 \(0<t<1\)，
	\[
		F((1-t)p+tq)<(1-t)F(p)+tF(q)\leq1.
	\]
	所以 \(K=\{F\leq1\}\) 的任意两个不同点之间的开线段都在其内部，\(K\) 严格凸。

	为计算曲率，沿正向弧长参数写 \(T=J\nabla F/|\nabla F|\)。在曲线上 \(\nabla F=(4x^3,4y^3)\ne0\)，而 \(JT=-\nabla F/|\nabla F|\)。两次微分 \(F(\gamma(s))=1\)，得
	\[
		0=T^{\mathsf T}(\mathrm{D}^2F)T+\langle\nabla F,kJT\rangle,
		\qquad
		k=\frac{T^{\mathsf T}(\mathrm{D}^2F)T}{|\nabla F|}.
	\]
	代入 \(\mathrm{D}^2F=\operatorname{diag}(12x^2,12y^2)\) 及 \(x^4+y^4=1\)，得到
	\[
		k=\frac{3x^2y^2}{(x^6+y^6)^{3/2}}.
	\]
	它在 \((\pm1,0)\)、\((0,\pm1)\) 为零，其余点均为正。
\end{solution}

\begin{exercise}\label{g1:ex:smooth-convex-flat-segment}
	构造一条 \(C^\infty\) 正则简单凸闭曲线，使其边界含非退化直线段。可先构造非负光滑周期函数 \(k\)，再积分切向角；须同时验证闭合，而不只验证切向量闭合。
\end{exercise}
\begin{solution}
	在 \([0,1]\) 上令
	\[
		k(s)=
		\begin{cases}
			c\exp\!\left(-\dfrac1{(s-\frac14)(\frac34-s)}\right),&\frac14<s<\frac34,\\
			0,&\text{其余情形},
		\end{cases}
	\]
	并按周期 \(1\) 延拓。常数 \(c>0\) 选成 \(\int_0^1k(s)\,\mathrm{d}s=\pi\)。这个选择存在，因为未乘 \(c\) 的积分有限且为正；端点处各阶导数均为零，故周期延拓为 \(C^\infty\) 函数。端点光滑性也可直接由 \(u^{-N}e^{-1/u}\to0\)（\(u\downarrow0\)，任意 \(N\)）逐次检查。

	令
	\[
		\theta(s)=\int_0^s k(v)\,\mathrm{d}v,\qquad
		T(s)=(\cos\theta(s),\sin\theta(s)),\qquad
		\gamma(s)=\int_0^sT(v)\,\mathrm{d}v.
	\]
	周期性给出 \(\theta(s+1)=\theta(s)+\pi\)，所以 \(T(s+1)=-T(s)\)。因此
	\[
		\gamma(2)=\int_0^1T(v)\,\mathrm{d}v+\int_1^2T(v)\,\mathrm{d}v=0,
	\]
	而 \(T\) 及其各阶导数为 \(2\) 周期，\(\gamma\) 便定义了参数圆周 \(\mathbb R/2\mathbb Z\) 上的光滑闭曲线。其速率为 \(1\)，有向曲率就是所给 \(k\geq0\)，切向角在两周期内增加 \(2\pi\)。由\cref{g1:thm:curvature-convexity}，它自动简单且为凸边界。在 \(0\leq s\leq1/4\) 上，\(\theta=0\)，故 \(\gamma(s)=(s,0)\)，提供了所需的非退化直线段。
\end{solution}

\begin{exercise}\label{g1:ex:helix-total-curvature}
	设 \(a>0\)、\(b\ne0\)，考虑螺旋线的一段
	\[
		\gamma(t)=(a\cos t,a\sin t,bt),\qquad 0\le t\le2\pi.
	\]
	求它的长度、曲率与总曲率，并说明为什么它的单位切向量虽然首尾相同，总曲率仍可小于 \(2\pi\)。指出 Fenchel 不等式证明中究竟哪一步不再适用。
\end{exercise}
\begin{solution}
	由 \(\|\dot\gamma(t)\|=\sqrt{a^2+b^2}\)，得到长度 \(L=2\pi\sqrt{a^2+b^2}\)。单位切向量是
	\[
		T(t)=\frac{(-a\sin t,a\cos t,b)}{\sqrt{a^2+b^2}},
	\]
	所以 \(\|\mathrm{d}T/\mathrm{d}t\|=a/\sqrt{a^2+b^2}\)，再除以速率得到 \(\kappa=a/(a^2+b^2)\)。于是
	\[
		\operatorname{TC}(\gamma)=\frac{2\pi a}{\sqrt{a^2+b^2}}<2\pi.
	\]
	切向示性曲线是高度 \(b/\sqrt{a^2+b^2}\) 的纬圆，确实闭合，而且包含在一个开半球内。原曲线的端点却相差 \((0,0,2\pi b)\)，因而
	\[
		\int_0^L T(s)\,\mathrm{d}s=(0,0,2\pi b)\ne0.
	\]
	证明中的位移抵消条件失效。切向量首尾相同只能保证接缝处的一阶数据相容，不能代替位置闭合。
\end{solution}

\begin{exercise}\label{g1:ex:bending-energy-lower-bound}
	设 \(\gamma\) 是长度为 \(L>0\) 的 \(C^2\) 正则闭空间曲线，定义 \(E(\gamma)=\int_0^L\kappa^2\,\mathrm{d}s\)。
	\begin{enumerate}
		\item 若 \(\kappa\le K\)，其中 \(K>0\)，证明 \(L\ge2\pi/K\)，并描述等号情形。
		\item 证明 \(E(\gamma)\ge4\pi^2/L\)，并证明等号当且仅当曲线为一个圆，走一周。
		\item 若曲线是非平凡结，利用\cref{g1:thm:fary-milnor}给出更强的严格能量下界。
	\end{enumerate}
\end{exercise}
\begin{solution}
	1. Fenchel 不等式给出 \(2\pi\le\int\kappa\,\mathrm{d}s\le KL\)，从而得到长度下界。若 \(L=2\pi/K\)，两个不等式都为等号。连续函数 \(K-\kappa\ge0\) 的积分为零，所以 \(\kappa\equiv K\)；同时等号定理说明曲线是平面凸曲线一周。在选定正向后，切角满足 \(\theta'=K\)，积分 \(T=(\cos\theta,\sin\theta)\) 得半径 \(1/K\) 的圆。该圆确实达到等号。

	2. 对函数 \(\kappa\) 与常函数 \(1\) 使用 Cauchy--Schwarz 不等式，
	\[
		4\pi^2\le\left(\int_0^L\kappa\,\mathrm{d}s\right)^2
		\le L\int_0^L\kappa^2\,\mathrm{d}s.
	\]
	若最终为等号，则总曲率为 \(2\pi\)，并且 \(\int_0^L(\kappa-2\pi/L)^2\,\mathrm{d}s=E-4\pi^2/L=0\)。被积函数连续非负，故处处为零，即 \(\kappa\equiv2\pi/L\)。第一问的积分论证说明曲线是半径 \(L/(2\pi)\) 的圆一周。反之直接代入可得 \(E=4\pi^2/L\)。

	3. 非平凡结满足 \(\operatorname{TC}(\gamma)>4\pi\)，故同一不等式给出 \(E(\gamma)>16\pi^2/L\)。严格性来自总曲率的严格下界，不需要另外讨论 Cauchy--Schwarz 能否取等号。
\end{solution}

\begin{exercise}\label{g1:ex:small-wavy-circle}
	令 \(0<|a|<\sqrt3/4\)，并设
	\[
		\gamma_a(t)=(\cos t,\sin t,a\sin2t),\qquad t\in\mathbb R/2\pi\mathbb Z.
	\]
	证明这是一条不包含在任何仿射平面内的正则嵌入闭曲线，且
	\[
		2\pi<\operatorname{TC}(\gamma_a)<4\pi.
	\]
	另给一个显式环境同痕将它移到平面圆，从而不用 Fáry--Milnor 定理也能判定其结型。
\end{exercise}
\begin{solution}
	前两坐标已经区分了圆周参数，且 \(\|\dot\gamma_a(t)\|\ge1\)，所以正则性与单射性成立。圆周紧而 \(\mathbb R^3\) 为 Hausdorff 空间，由\cref{g1:thm:compact-bijection}可知它与其像同胚。若整条曲线属于平面 \(Ax+By+Cz=D\)，就有
	\[
		A\cos t+B\sin t+Ca\sin2t=D\qquad\text{对所有 }t.
	\]
	依次乘以 \(1,\cos t,\sin t,\sin2t\) 并在 \([0,2\pi]\) 上积分，得到 \(D=A=B=Ca=0\)。因为 \(a\ne0\)，这迫使 \(A=B=C=0\)，不能定义平面。因此\cref{g1:cor:nonplanar-total-curvature}给出严格下界。

	记 \(v=\|\dot\gamma_a\|\)。对 \(T=\dot\gamma_a/v\) 求导，它的导数是 \(\ddot\gamma_a/v\) 在 \(T^\perp\) 上的正交投影，因此
	\[
		\left\|\frac{\mathrm{d}T}{\mathrm{d}t}\right\|
		\le\frac{\|\ddot\gamma_a(t)\|}{\|\dot\gamma_a(t)\|}
		\le\sqrt{1+16a^2}.
	\]
	积分得到 \(\operatorname{TC}(\gamma_a)\le2\pi\sqrt{1+16a^2}<4\pi\)。

	最后定义 \(H_u(x,y,z)=(x,y,z-2ua\,xy)\)，其中 \(0\le u\le1\)。其逆为 \((x,y,z)\mapsto(x,y,z+2ua\,xy)\)，二者都连续依赖 \(u\)，而 \(H_0=\mathrm{id}\)。在曲线上 \(2xy=\sin2t\)，故 \(H_1\circ\gamma_a(t)=(\cos t,\sin t,0)\)，给出了所求环境同痕。
\end{solution}

\begin{exercise}\label{g1:ex:spherical-small-circle}
	一条 \(C^2\) 单位速率曲线落在半径为 \(R>0\) 的球面上时，由 \(\|\gamma(s)\|^2=R^2\) 可得 \(\kappa(s)\ge1/R\)。有人把这个结论与 Fenchel 不等式合并，声称球面上的任意正则闭曲线都长至少 \(2\pi R\)。解释这一推理为什么不成立，并以纬圆给出反例。
\end{exercise}
\begin{solution}
	将球心移到原点，两次微分给出 \(\gamma\cdot T=0\) 以及 \(\gamma\cdot\gamma''=-1\)。由 Cauchy--Schwarz，\(1\le R\|\gamma''\|=R\kappa\)。积分得到的是 \(\operatorname{TC}(\gamma)\ge L/R\)，它与 \(\operatorname{TC}(\gamma)\ge2\pi\) 都是总曲率的下界，不能比较右端的 \(L/R\) 和 \(2\pi\)。

	取 \(0<\alpha<\pi/2\)，令
	\[
		\gamma(t)=
		(R\sin\alpha\cos t,R\sin\alpha\sin t,R\cos\alpha).
	\]
	这是球面上一个半径 \(R\sin\alpha\) 的平面圆，长度 \(2\pi R\sin\alpha<2\pi R\)，曲率 \(1/(R\sin\alpha)\)，总曲率仍为 \(2\pi\)。\(\alpha\to0\) 时长度趋于零，曲率却相应增大。
\end{solution}


% !TeX root = ../../Book1.tex
% 纲要标题：## 第8章　曲面的基本形式、相容条件与局部实现
% 纲要位置：几何学计划纲要.md，第 398 行。
% 正文状态：初稿；尚未正式审核。

% 纲要内容开始
% 
% > 本章目标与证明范围：完整建立从浸入到两基本形式、再由相容数据恢复浸入的双向理论。第一基本形式不再独立为浅章；Gauss–Codazzi 的必要性和局部存在性分别展开。
% 
% 纲要内容结束

\chapter{曲面的基本形式、相容条件与局部实现}
\label{g1:ch:08}
\BookChapterNumber{g1:ch:08}

\begin{chapterabstract}
	一张曲面同时携带沿面测量的度量与在空间中弯曲的法向数据。本章从正则参数、切平面与坐标变换建立两种基本形式，推导长度、面积、微分算子和主曲率的计算方法，再以混合偏导相容性联系内在曲率与外在弯曲。反向构造从光滑正定度量及对称第二基本形式出发：相容方程使标架系统可积分，闭的候选微分再恢复位置，从而得到曲面基本定理的局部存在与唯一性。最后用球面、圆柱、旋转曲面及一阶形变说明数据、法向、局部实现和整体刚性各自的作用。
\end{chapterabstract}

\begin{learninggoals}
	\begin{itemize}
		\item 从参数映射检查浸入条件，求两种基本形式，并在换参数或换法向后正确转换数据。
		\item 利用第一基本形式计算长度、面积、梯度、散度及本书符号约定下的 Laplace 算子，辨别局部等距与环境刚性运动。
		\item 求形算子、主曲率、平均曲率与法曲率，判断脐点及渐近方向，并说明点态方向与局部积分曲线的区别。
		\item 从切向投影推导结构方程，证明曲率由度量决定，并分清两组相容条件各自排除的障碍。
		\item 通过矩形上的标架积分和位置积分重建曲面，核验秩、法向、两种形式及初值唯一性。
		\item 用明确算例区分局部存在、整体单值、整体单射与一阶形变。
	\end{itemize}
\end{learninggoals}

在平面纸片上画一个小三角形，再把纸片卷成圆柱的一部分，它的三条边沿纸测量时仍可保持原来的长度，但空间中的三个顶点及其切平面已经改变。这个例子提出两个问题：哪些数据描述纸片内部的测量，哪些数据描述它在空间中的摆放？反过来，如果只给出这两类数据，是否总能找到一个曲面把它们实现出来？

\ProseParagraphBreak

第一类数据来自两个参数切向量的内积，即第一基本形式；第二类来自法向量随位置的改变，即第二基本形式。后者不是任意的三个函数。一个位置函数 \(X(u,v)\) 同时沿两个方向变化，混合偏导必须相等，因此长度和弯曲数据要满足相容方程。曲线的重建只需沿一条参数线解常微分方程；曲面的重建则还要证明，先横走后竖走与相反次序不会给出冲突的标架。

\ProseParagraphBreak

本章将这一步直接落实为积分证明：先归一化待求标架，再建立矩形上的双变量系统，最后对两个候选切向量积分得到位置。正定性用于保证得到的是浸入；相容性用于消除混合偏导的冲突；指定初值用于说明唯一性。绕环路后的标架与位移另作检查，因为每个小邻域能够实现，并不保证带孔区域上存在单值实现。

\ProseParagraphBreak

所需基础是多元微积分中的链式法则、偏导交换、Taylor 展开、积分换元与分部积分，以及内积和矩阵运算。局部逆映射的所需版本、二维自伴算子的主方向、局部积分曲线和双变量系统的存在性均在本章证明。第七章的长度公式和线性方程引理将在使用处回顾。参数曲面、两种基本形式和曲率的延伸阅读可查 Pressley 第二版的明确位置\cite[第4.1--4.5, 6.1--6.2, 7.1--7.3, 8.1--8.2, 10.1--10.2节]{Pressley2010}。

% !TeX root = ../../Book1.tex
% 纲要标题：### 8.1 正则曲面和度量数据
% 纲要位置：几何学计划纲要.md，第 402 行。
% 正文状态：初稿；尚未正式审核。

% 纲要内容开始
% 
% * 浸入参数、图像模型和切平面
% * 第一基本形式的坐标变换
% * 长度、面积和局部等距
% * 梯度、散度与曲面上的微分算子
% 
% 纲要内容结束

\section{正则曲面和度量数据}
\label{g1:sec:0801}

把一张纸弯成圆柱的一部分，纸上两条相交曲线的夹角和沿纸测量的长度可以保持不变，纸在空间中的法方向却已经改变。要把这两种现象分开，首先应弄清参数怎样记录曲面，以及哪些计算只需两个切向量的内积。本章均在光滑即 \(C^\infty\) 的范围内讨论；这使后面交换偏导与求解标架系统时无需反复改变正则性条件，并不表示每一条单独的公式都需要无限次可微。

\subsection{浸入参数、图像模型和切平面}
\label{g1:subsec:0801:01}

曲线的正则性要求速度不为零，曲面的对应要求是两个参数方向不能塌缩成一条直线。以下记 \(u^1=u,u^2=v\)，\(\partial_i=\partial/\partial u^i\)，\(X_i=\partial_iX\)。

\begin{definition}[正则参数曲面]\label{g1:def:surface-immersion}
	设 \(U\subseteq\mathbb R^2\) 开。光滑映射 \(X:U\rightarrow\mathbb R^3\) 若处处满足 \(\operatorname{rank}\mathrm{d}X=2\)，称为一个\term[zhengzecanshuqumian]{正则参数曲面}[regular parametrized surface]，也称参数域到 \(\mathbb R^3\) 的\term[jinru]{浸入}[immersion]。等价地，\(X_1,X_2\) 线性无关，或 \(X_1\times X_2\ne0\)。在参数点 \(q\in U\) 处的\term[qiepingmian]{切平面}[tangent plane]作为向量空间定义为
	\[
		\mathrm{T}_qX\coloneqq \mathrm{d}X_q(\mathbb R^2)=\operatorname{span}\{X_1(q),X_2(q)\}.
	\]
	经过曲面点的仿射切平面是 \(X(q)+\mathrm{T}_qX\)。若两个参数点具有同一空间像，记号 \(\mathrm{T}_qX\) 仍记录所选参数点，而不擅自把两个分支的切平面等同。
\end{definition}

最直接的模型是图像 \(X(u,v)=(u,v,f(u,v))\)：前两个分量已经保证两个偏导线性无关。反过来，每个浸入在足够小的参数邻域内都能换成这样的模型。为免把“选适当坐标”隐藏为一句话，先说明这里需要的逆映射事实。

\begin{lemma}[局部逆映射的一个版本]\label{g1:lem:surface-local-inverse}
	设 \(f:V\rightarrow\mathbb R^m\) 光滑，\(V\subseteq\mathbb R^m\) 开，且 \(\mathrm{D}f(a)\) 可逆。则 \(a\) 有开邻域 \(V_0\)，使 \(f|_{V_0}\) 是到某个开集的光滑双射，逆映射也光滑。
\end{lemma}

\begin{proof}
	平移并左乘 \(\mathrm{D}f(a)^{-1}\)，可设 \(a=0,f(0)=0,\mathrm{D}f(0)=I\)。取闭球 \(\overline B_r\subset V\)，使其中 \(\|\mathrm{D}f-I\|\le q<1\)。沿线段积分导数得到
	\[
		\|f(x)-f(z)-(x-z)\|\le q\|x-z\|,\qquad
		\|f(x)-f(z)\|\ge(1-q)\|x-z\|.
	\]
	对 \(\|y\|<(1-q)r\)，映射 \(\Phi_y(x)=y+x-f(x)\) 将闭球映入自身，且距离至多缩短为原来的 \(q\) 倍。从球内任一点迭代，相邻两项的距离至多为 \(q^j\) 倍初始距离，故迭代列为 Cauchy 列；闭球完备，极限 \(x\) 满足 \(f(x)=y\)。同一缩短估计证明解唯一。令 \(V_0=B_r\cap f^{-1}(B_{(1-q)r})\)，便得所需双射；上面的第二个不等式使逆映射连续且局部 Lipschitz。

	写 \(x=f^{-1}(y)\) 及 \(x+h=f^{-1}(y+k)\)。由 \(k=\mathrm{D}f(x)h+o(\|h\|)\) 和 \(\|h\|\le(1-q)^{-1}\|k\|\)，有
	\[
		h=\mathrm{D}f(x)^{-1}k+o(\|k\|).
	\]
	因此逆映射可微，导数为 \(\mathrm{D}f(f^{-1}(y))^{-1}\)，且连续。对这个表达式反复求导，依次得到全部高阶连续导数。恢复最初的平移和可逆线性变换即得结论。
\end{proof}

\begin{proposition}[浸入的局部图像]\label{g1:prop:surface-local-graph}
	设 \(X:U\rightarrow\mathbb R^3\) 为光滑浸入，\(q\in U\)。选择空间中适当的两个坐标后，\(q\) 有邻域 \(U_0\)，使 \(X(U_0)\) 是这两个坐标上的光滑函数图像；\(X|_{U_0}\) 是到其像的同胚。其切平面恰为所有通过 \(q\) 的光滑参数曲线在空间中的速度所组成的集合。特别地，可先用刚性运动把该点移到原点、切平面移成水平面，再得到满足 \(f(0)=0,\mathrm{d}f(0)=0\) 的图像。
\end{proposition}

\begin{proof}
	矩阵 \(\mathrm{d}X_q\) 的秩为二，故某个二阶子式不为零。设对应前两个空间坐标，令 \(\pi(x,y,z)=(x,y)\)。由前一引理，\(\pi\circ X\) 在较小邻域内有光滑逆 \(\psi\)。于是
	\[
		X\circ\psi(x,y)=(x,y,f(x,y)).
	\]
	图像的逆坐标是连续的投影 \(\pi\)，所以原参数映射在该邻域内也是同胚。参数曲线 \(c(t)\) 满足 \(c(0)=q\) 时，\((X\circ c)'(0)=\mathrm{d}X_q(c'(0))\) 属于上述张成空间；反之，每个 \(\mathrm{d}X_q(w)\) 都由短直线 \(c(t)=q+tw\) 实现。若先将切平面移为水平面，投影 \(\pi\) 限制在该切平面上就是线性同构，所以上述逆映射论证仍适用；图像的水平切平面又使两个一阶偏导都为零。
\end{proof}

若子集 \(S\subseteq\mathbb R^3\) 在每一点附近都恰好是一个这样的图像，称为\term[zhengzequmian]{正则曲面}[regular surface]，此时也称嵌入曲面。它的图像参数之间的过渡映射光滑，因为可写为一张图像的投影与另一张参数映射的复合。浸入只保证每个参数分支局部如此，不保证整个像是嵌入曲面。例如
\[
	X(u,v)=(\sin u,\sin2u,v),\qquad (u,v)\in\mathbb R^2
\]
处处为浸入：\((\cos u,2\cos2u)\) 不会同时为零。但 \(u=0\) 与 \(u=\pi\) 的分支在同一竖直线上相交，两个切平面不同。本章的局部计算始终在选定分支上进行。

\subsection{第一基本形式的坐标变换}
\label{g1:subsec:0801:02}

参数方向 \(a=(a^1,a^2)\) 本身的欧氏长度通常不是曲面上相应速度的长度。实际速度是 \(a^1X_1+a^2X_2\)，其平方范数把所需数据集中为一个二次型。

\begin{definition}[第一基本形式]\label{g1:def:first-fundamental-form}
	设 \(X:U\rightarrow\mathbb R^3\) 为光滑浸入。定义
	\[
		g_{ij}=\langle X_i,X_j\rangle,\qquad
		E=g_{11},\quad F=g_{12}=g_{21},\quad G=g_{22},\quad D=EG-F^2.
	\]
	参数方向上的内积 \(g(a,b)=\sum_{i,j=1}^2g_{ij}a^ib^j\) 称为曲面的\term[diyijibenxingshi]{第一基本形式}[first fundamental form]，记
	\[
		I=E\,\mathrm{d}u^2+2F\,\mathrm{d}u\,\mathrm{d}v+G\,\mathrm{d}v^2.
	\]
	这里 \(\mathrm{d}u^2\) 表示 \(\mathrm{d}u\otimes \mathrm{d}u\) 在同一向量上求值所得的平方，混合项同理；它不是对参数函数取二阶导数。第一基本形式在每点正定，故 \(E>0,D>0\)。
	\symidx{\(g_{ij}\)：第一基本形式的系数}
	\symidx{\(I\)：第一基本形式}
\end{definition}

\begin{proposition}[度量的换元公式]\label{g1:prop:surface-metric-change}
	设 \(X:U\rightarrow\mathbb R^3\) 为光滑浸入，\(\phi:\widetilde U\rightarrow U\) 为光滑微分同胚，即光滑且有光滑逆的双射。置 \(\widetilde X=X\circ\phi\)，以 \(J=(\partial u^i/\partial\widetilde u^a)\) 记换元导数矩阵。则
	\begin{equation}\label{g1:eq:surface-metric-change}
		\widetilde g=J^{\mathsf T}(g\circ\phi)J,\qquad
		\det\widetilde g=(\det J)^2(D\circ\phi).
	\end{equation}
	特别地，\(g\) 的正定性以及相应切向量之间的长度和夹角不依赖坐标。
\end{proposition}

\begin{proof}
	链式法则给出 \(\widetilde X_a=\sum_iX_i\,\partial_a\phi^i\)。两边取内积即为第一式，再取行列式即为第二式。若同一参数方向在两坐标中的分量满足 \(a=J\widetilde a\)，则 \(\widetilde a^{\mathsf T}\widetilde g\widetilde b=a^{\mathsf T}gb\)，这也直接说明所有内积计算保持一致。
\end{proof}

因此，度量不是某张矩阵本身，而是按这一规律相互转换的一族内积。即使暂时没有空间中的曲面，参数域上的任意光滑正定对称矩阵场也定义一个\term[limanduliang]{黎曼度量}[Riemannian metric]；本章重建问题正是询问这样的 \(g\) 与另一个对称数据 \(h\) 能否同时来自浸入。此名称所指的内在度量与环境空间中的欧氏距离应区分。\footnote{Riemann，黎曼，德文全名 Georg Friedrich Bernhard Riemann（1826--1866），德国数学家，在复分析、几何及分析学中作出基础性贡献。}

\begin{example}[坐标网格可以不正交]\label{g1:exm:oblique-plane-metric}
	平面参数 \(X(u,v)=(u+v,v,0)\) 给出 \(E=1,F=1,G=2\)，所以 \(I=\mathrm{d}u^2+2\,\mathrm{d}u\,\mathrm{d}v+2\,\mathrm{d}v^2\)。两个参数轴在空间中的夹角余弦为 \(1/\sqrt2\)，而不是零。另一方面，参数 \(X(r,\theta)=(r\cos\theta,r\sin\theta,0)\) 在 \(r>0\) 时给出 \(I=\mathrm{d}r^2+r^2\mathrm{d}\theta^2\)。系数变化可能仅来自坐标；这两张矩阵都描述平面的欧氏度量。
\end{example}

\subsection{长度、面积和局部等距}
\label{g1:subsec:0801:03}

度量的用途先从测量开始。若 \(c:[a,b]\rightarrow U\) 分段 \(C^1\)，空间曲线的速度为 \(\mathrm{d}X(c')\)，由\cref{g1:prop:arc-length-param}中的长度积分公式，
\begin{equation}\label{g1:eq:surface-length}
	L_g(c)=\int_a^b\sqrt{\sum_{i,j}g_{ij}(c(t))\dot c^i(t)\dot c^j(t)}\,\mathrm{d}t.
\end{equation}
两非零方向 \(a,b\) 的夹角余弦为 \(g(a,b)/\sqrt{g(a,a)g(b,b)}\)。这里距离沿曲面测量；将两点在空间中直线相连所得到的弦长一般不等于沿面的最短路程。

\begin{definition}[面积密度]\label{g1:def:surface-area-density}
	对光滑浸入 \(X:U\rightarrow\mathbb R^3\)，其\term[mianjimidu]{面积密度}[area density]及参数区域 \(A\subset U\) 的面积定义为
	\[
		\mathrm{d}A_g=\sqrt D\,\mathrm{d}u\,\mathrm{d}v,\qquad
		\operatorname{Area}_X(A)=\int_A\sqrt D\,\mathrm{d}u\,\mathrm{d}v,
	\]
	其中 \(A\) 可取紧包含于 \(U\) 的分片光滑有界区域。若不同参数部分覆盖同一空间部分，此积分计入覆盖次数。
\end{definition}

这个因子的几何来源是微分 \(\mathrm{d}X\) 把参数单位小正方形送成由 \(X_1,X_2\) 张成的平行四边形，而
\[
	\|X_1\times X_2\|^2=\|X_1\|^2\|X_2\|^2-\langle X_1,X_2\rangle^2=D.
\]
面积定义使用微分给出的局部伸缩；并不声称随意选取的空间多面体逼近，其面积都自动收敛到这个值。由\eqref{g1:eq:surface-metric-change}，\(\sqrt{\det\widetilde g}=\sqrt{D\circ\phi}\,|\det J|\)，结合多元积分换元公式便知面积不依赖参数。换元公式在这里的具体用途是：可微双射的导数行列式绝对值恰补偿参数区域面积的改变。

\begin{definition}[局部等距]\label{g1:def:surface-local-isometry}
	设 \(U,V\subseteq\mathbb R^2\) 带光滑正定度量 \(g,\widetilde g\)。光滑映射 \(\phi:U\rightarrow V\) 若在每点附近为微分同胚，且
	\[
		g_q(a,b)=\widetilde g_{\phi(q)}(\mathrm{d}\phi_q(a),\mathrm{d}\phi_q(b))
	\]
	对所有 \(q,a,b\) 成立，称为\term[jubudengju]{局部等距}[local isometry]。对嵌入曲面的映射，在局部参数中采用同一定义。若映射在整体上是微分同胚并满足此式，则称为两度量曲面的等距。
\end{definition}

\begin{proposition}[长度刻画局部等距]\label{g1:prop:local-isometry-length}
	设 \(\phi:(U,g)\rightarrow(V,\widetilde g)\) 为光滑局部微分同胚。则 \(\phi\) 是局部等距，当且仅当每条分段 \(C^1\) 参数曲线 \(c\) 都满足 \(L_{\widetilde g}(\phi\circ c)=L_g(c)\)。局部等距还保持夹角及在单射邻域内计算的面积。
\end{proposition}

\begin{proof}
	保持内积直接给出速度模相同，积分即得长度相同；夹角公式和度量行列式的换元式给出另两项。反过来，固定 \(q\) 与方向 \(a\)，取短线段 \(c(t)=q+ta\)，\(0\le t\le\varepsilon\)。将长度相等除以 \(\varepsilon\)，令 \(\varepsilon\to0^+\)，连续性给出 \(g_q(a,a)=\widetilde g_{\phi(q)}(\mathrm{d}\phi_q(a),\mathrm{d}\phi_q(a))\)。利用
	\[
		2g(a,b)=g(a+b,a+b)-g(a,a)-g(b,b)
	\]
	恢复全部内积，故得局部等距。
\end{proof}

这个定义比较的是切向量及沿面测量，不要求 \(\phi\) 是环境空间某个刚性运动的限制。\subsecref{g1:subsec:0805:02}会用圆柱把这一区别算出来。局部等距也不保证整体单射，因而不能将局部保持长度直接理解为全局最短距离和拓扑都不改变。

\subsection{梯度、散度与曲面上的微分算子}
\label{g1:subsec:0801:04}

已知一个温度函数 \(f\)，沿方向 \(a\) 的变化率 \(\mathrm{d}f(a)=\sum_i(\partial_i f)a^i\) 不需要度量；但问“沿哪个单位方向增长最快”，就必须知道什么叫单位长度。反过来，描述切向流的局部流出量时，还必须使用正确的面积密度。这两点分别导出梯度与散度。

\begin{definition}[梯度]\label{g1:def:surface-gradient}
	在带光滑正定度量 \(g\) 的平面开集 \(U\) 上，光滑函数 \(f\) 的\term[tidu]{梯度}[gradient]是唯一满足 \(g(\operatorname{grad}_g f,a)=\mathrm{d}f(a)\) 的向量场。若 \(g^{-1}=(g^{ij})\)，则
	\begin{equation}\label{g1:eq:surface-gradient}
		\operatorname{grad}_g f=\sum_{i,j}g^{ij}(\partial_j f)\partial_i.
	\end{equation}
	在浸入曲面上，右端的 \(\partial_i\) 经 \(\mathrm{d}X\) 对应到 \(X_i\)。
\end{definition}

求解线性方程 \(g_{ij}(\operatorname{grad}f)^j=\partial_i f\) 就得到公式。它说明直接把两个偏导数当作梯度分量，只有在度量矩阵为单位阵的坐标中才正确。若梯度非零，Cauchy--Schwarz 不等式给出 \(\mathrm{d}f(a)\le\|\operatorname{grad}f\|_g\)（\(\|a\|_g=1\)），等号在梯度方向达到；若 \(f\) 是环境函数 \(\widehat f\) 的限制，曲面梯度恰为欧氏梯度 \(\nabla\widehat f\) 的切向正交投影，因为两者与任何切向量的内积都等于该方向导数。

\begin{definition}[散度与 Laplace 算子]\label{g1:def:surface-divergence-laplacian}
	设 \(g\) 为 \(U\subseteq\mathbb R^2\) 上的光滑正定度量，\(W=\sum_iW^i\partial_i\) 为光滑向量场。定义其\term[sandu]{散度}[divergence]为
	\begin{equation}\label{g1:eq:surface-divergence}
		\operatorname{div}_gW=\frac1{\sqrt D}\sum_i\partial_i(\sqrt D\,W^i).
	\end{equation}
	函数上的\term[LaplaceBeltramisuanzi]{Laplace--Beltrami 算子}[Laplace--Beltrami operator]采用本书的非负号约定：
	\begin{equation}\label{g1:eq:surface-laplacian}
		\Delta_g f=-\operatorname{div}_g\operatorname{grad}_g f
		=-\frac1{\sqrt D}\sum_{i,j}\partial_i\bigl(\sqrt D\,g^{ij}\partial_jf\bigr).
	\end{equation}
	\symidx{\(\Delta_g\)：非负号约定的 Laplace--Beltrami 算子}
\end{definition}

上述名称分别纪念 Laplace 和 Beltrami。\footnote{Laplace，拉普拉斯，法文全名 Pierre-Simon Laplace（1749--1827），法国数学家和天文学家，研究概率、势函数与天体力学。Beltrami，贝尔特拉米，意大利文名 Eugenio Beltrami（1835--1900），意大利数学家，研究微分几何、非欧几何模型及微分算子。}为看清散度为何需要 \(\sqrt D\)，取紧支于 \(U\) 内的光滑函数 \(\eta\)，记作 \(\eta\in C_c^\infty(U)\)。通常的分部积分没有边界项，因此
\begin{equation}\label{g1:eq:surface-divergence-weak}
	\int_U\eta\,\operatorname{div}_gW\,\mathrm{d}A_g
	=-\int_U \mathrm{d}\eta(W)\,\mathrm{d}A_g.
\end{equation}
右边只涉及方向导数和已证明与坐标无关的面积，因而不依赖坐标。这一等式还唯一确定连续函数 \(\operatorname{div}_gW\)：若两个候选的差在一点为正，则在一个小圆盘上仍为正，在该圆盘取非负且不恒零的光滑函数作 \(\eta\)，便与积分为零矛盾。这样的函数可以显式取为圆盘内的 \(\exp(-1/(r^2-\|x-a\|^2))\)，在圆盘外取零；负值情形同理。于是散度公式确实定义了一个内在的标量，而不是坐标的偶然组合。

\begin{proposition}[非负号约定与局部等距]\label{g1:prop:surface-laplacian-energy}
	设 \(g\) 为平面开集 \(U\) 上的光滑正定度量，\(f,\eta\in C_c^\infty(U)\)。则
	\[
		\int_U\eta\Delta_g f\,\mathrm{d}A_g
		=\int_U g(\operatorname{grad}_g\eta,\operatorname{grad}_gf)\,\mathrm{d}A_g,
		\qquad
		\int_Uf\Delta_g f\,\mathrm{d}A_g\ge0.
	\]
	若 \(\phi:(U,g)\rightarrow(V,\widetilde g)\) 为光滑局部等距，则每个光滑函数 \(f\) 满足 \(\Delta_g(f\circ\phi)=(\Delta_{\widetilde g}f)\circ\phi\)。
\end{proposition}

\begin{proof}
	第一式在\eqref{g1:eq:surface-divergence-weak}中令 \(W=\operatorname{grad}f\) 并使用梯度的定义；令 \(\eta=f\) 得到第二式。对局部等距，定义直接给出 \(\mathrm{d}\phi(\operatorname{grad}_g(f\circ\phi))=(\operatorname{grad}_{\widetilde g}f)\circ\phi\)。在 \(\phi\) 单射的小邻域中，用面积换元和\eqref{g1:eq:surface-divergence-weak}可知相对应向量场的散度也对应。将两项结合即得算子等式；等式是逐点的，所以这些邻域足够覆盖整个定义域。
\end{proof}

“非负”说的是积分二次型，不是说 \(\Delta_g f\) 对每个 \(f\) 都逐点非负。在平面的极坐标中，\(g=\operatorname{diag}(1,r^2)\)、\(\sqrt D=r\)，故
\[
	\operatorname{grad}f=f_r\partial_r+\frac{f_\theta}{r^2}\partial_\theta,\quad
	\operatorname{div}W=\frac1r\partial_r(rW^r)+\partial_\theta W^\theta,\quad
	\Delta f=-f_{rr}-\frac1rf_r-\frac1{r^2}f_{\theta\theta}.
\]
例如 \(f=r^2\) 时 \(\Delta f=-4\)。这些式子只在 \(r>0\) 的极坐标区域使用；原点处公式的分母反映坐标失效，不能据此断言平面的几何有奇点。

% !TeX root = ../../Book1.tex
% 纲要标题：### 8.2 法向几何与第二基本形式
% 纲要位置：几何学计划纲要.md，第 409 行。
% 正文状态：初稿；尚未正式审核。

% 纲要内容开始
% 
% * Gauss 映射和形算子
% * 主曲率、法曲率与 Euler 公式
% * 平均曲率、Gauss 曲率及脐点
% * 曲率线、渐近线及其退化
% 
% 纲要内容结束

\section{法向几何与第二基本形式}
\label{g1:sec:0802}

一张纸卷成圆柱后，小范围内的长度和角度没有改变，切平面却随着位置转动。第一基本形式记录前一种信息；要描述后一种变化，就应比较邻近点的法向量。曲线只有一个切方向，曲面在一点有整整一圈方向，因此需要一个线性算子同时记录沿各方向前进时法向量怎样变化。以下所有曲面参数均为光滑浸入 \(X:U\to\mathbb R^3\)，其中 \(U\subset\mathbb R^2\) 为开集；讨论在参数邻域中进行，自交处的不同分支分别处理。

\subsection{Gauss 映射和形算子}
\label{g1:subsec:0802:01}

令 \(X_i=\partial X/\partial u^i\)。上一节的正则性保证 \(X_1\times X_2\ne0\)，因而在每张参数片上可以选择光滑单位法向
\[
	n=\frac{X_1\times X_2}{\sqrt D},\qquad
	D=EG-F^2,\quad E=\langle X_1,X_1\rangle,\quad F=\langle X_1,X_2\rangle,\quad G=\langle X_2,X_2\rangle.
\]
也可以选相反的 \(n\)。以下先固定这一选择；曲面若不能作整体的法向选择，这些构造仍可在每张小参数片上使用。

\begin{definition}[Gauss 映射与形算子]\label{g1:def:shape-operator}
	设 \(X:U\to\mathbb R^3\) 为光滑浸入，\(n:U\to S^2\) 为一个光滑单位法向场。把每一点送到其单位法向的映射 \(n\) 称为\term[Gaussyingshe]{Gauss 映射}[Gauss map]。\footnote{Gauss（高斯，Carl Friedrich Gauss，1777--1855）是德国数学家，在数论、微分几何与测地学等领域作出重要贡献。}对每个参数点 \(q\in U\)，令 \(p=X(q)\)，并在这一参数分支上定义\term[xingsuanzi]{形算子}[shape operator]
	\[
		S_q:\mathrm{T}_qX\longrightarrow \mathrm{T}_qX,\qquad
		S_q\!\left(\sum_i a^iX_i\right)=-\sum_i a^i n_i.
	\]
	它也记为 \(S=-\mathrm dn\)，其中用 \(\mathrm dX\) 把参数切向量与曲面切向量对应起来。
	\symidx{\(S\)：曲面的形算子 \(S=-\mathrm dn\)}
\end{definition}

这个定义中有两个需要核对的地方。首先，\(X_1,X_2\) 线性无关，所以系数 \(a^i\) 唯一；其次，对 \(\langle n,n\rangle=1\) 求导得到 \(\langle n_i,n\rangle=0\)，故 \(n_i\) 确实落在 \(\mathrm{T}_qX=n^\perp\) 中。Gauss 映射的微分本来取值于 \(\mathrm{T}_{n(q)}S^2\)，这里使用了这两个切平面作为 \(\mathbb R^3\) 线性子空间恰好都是 \(n^\perp\) 的事实。

\ProseParagraphBreak

负号使“曲线的法向加速度”稍后直接等于形算子的二次型。它也决定了球面的符号：半径 \(R\) 的球面若取外向法向 \(n(q)=p/R\)，则 \(\mathrm dn(v)=v/R\)，所以 \(S=-I/R\)。不要把别处采用 \(S=\mathrm dn\) 的主曲率公式直接移入这里。

\begin{definition}[第二基本形式]\label{g1:def:second-fundamental-form}
	在带选定单位法向的光滑浸入曲面上，对每个参数点 \(q\) 定义
	\[
		\mathrm{II}_q(v,w)\coloneqq\langle S_qv,w\rangle,\qquad v,w\in \mathrm{T}_qX.
	\]
	这个双线性型称为\term[dierjibenxingshi]{第二基本形式}[second fundamental form]。它的对称性由下面的计算给出。在坐标中记 \(h_{ij}=\mathrm{II}(X_i,X_j)\)，于是
	\[
		\mathrm{II}=h_{11}(\mathrm du^1)^2+2h_{12}\,\mathrm du^1\mathrm du^2+h_{22}(\mathrm du^2)^2.
	\]
	\symidx{\(\mathrm{II}\)：第二基本形式}
	\symidx{\(h_{ij}\)：第二基本形式的坐标系数}
\end{definition}

\begin{proposition}[形算子的自伴性与坐标表示]\label{g1:prop:shape-self-adjoint}
	设 \(X:U\to\mathbb R^3\) 为光滑浸入，并选定光滑单位法向 \(n\)。则
	\begin{equation}\label{g1:eq:second-form-coefficients}
		h_{ij}=-\langle n_i,X_j\rangle=\langle X_{ij},n\rangle.
	\end{equation}
	第二基本形式对称，形算子关于第一基本形式自伴，即 \(\langle Sv,w\rangle=\langle v,Sw\rangle\)。以列向量表示系数时，形算子在基 \(X_1,X_2\) 下的矩阵为
	\begin{equation}\label{g1:eq:shape-matrix}
		[S]=g^{-1}h,\qquad g=(g_{ij}),\quad h=(h_{ij}).
	\end{equation}
	若作光滑坐标变换 \(u=u(v)\)，令 \(A=(\partial u^i/\partial v^a)\)，并保留同一个几何法向场，则 \(\widetilde h=A^{\mathsf T}hA\)。因此 \(\mathrm{II}\) 不依赖参数的选择。
\end{proposition}
\begin{proof}
	对 \(\langle n,X_j\rangle=0\) 关于 \(u^i\) 求导，得到 \(\langle n_i,X_j\rangle+\langle n,X_{ij}\rangle=0\)。混合偏导相等给出 \(h_{ij}=h_{ji}\)，再由双线性性得到自伴性。若 \(S(X_j)=\sum_i S^i{}_jX_i\)，则
	\[
		\sum_i g_{ki}S^i{}_j=\langle X_k,S(X_j)\rangle=h_{kj}.
	\]
	即 \(g[S]=h\)，而 \(g\) 正定可逆，故得\eqref{g1:eq:shape-matrix}。

	对 \(\widetilde X(v)=X(u(v))\) 两次求导，有
	\[
		\widetilde X_{ab}=\sum_{i,j}X_{ij}A^i{}_aA^j{}_b+\sum_i X_i\frac{\partial^2u^i}{\partial v^a\partial v^b}.
	\]
	与 \(n\) 作内积后，第二项消失，第一项正是 \((A^{\mathsf T}hA)_{ab}\)。变换后的切向量系数代入这个式子，双线性型的值保持不变。
\end{proof}

这里的自伴性由内积 \(g\) 衡量。一般参数的 \(X_1,X_2\) 不一定正交，矩阵 \(g^{-1}h\) 因而未必等于自己的转置；直接把它当作标准内积下的对称矩阵，会混淆两个不同的内积。另一方面，若换参数后重新机械地采用叉积法向，则反向坐标变换会同时把 \(n\) 变成 \(-n\)，此时还要把 \(\mathrm{II}\) 整体变号。

\begin{example}[图像上的二阶数据]\label{g1:exm:graph-second-form}
	对图像 \(X(u,v)=(u,v,f(u,v))\)，取向上的单位法向
	\[
		n=\frac{(-f_u,-f_v,1)}W,\qquad W=\sqrt{1+f_u^2+f_v^2}.
	\]
	直接求导得到
	\[
		g=\begin{pmatrix}1+f_u^2&f_uf_v\\f_uf_v&1+f_v^2\end{pmatrix},
		\qquad
		h=\frac1W\begin{pmatrix}f_{uu}&f_{uv}\\f_{uv}&f_{vv}\end{pmatrix}.
	\]
	若 \(\mathrm df(0)=0\)，原点的切平面水平，\(g(0)=I\)，所以形算子就是 \(f\) 在原点的二阶偏导矩阵。此时
	\[
		f(u,v)=f(0)+\tfrac12\bigl(h_{11}u^2+2h_{12}uv+h_{22}v^2\bigr)+o(u^2+v^2).
	\]
	第二基本形式给出了曲面离开切平面的二阶高度；它没有包含三阶及更高阶的信息。
\end{example}

\subsection{主曲率、法曲率与 Euler 公式}
\label{g1:subsec:0802:02}

沿同一切方向，可以画出许多不同的曲面曲线，它们的空间加速度当然未必相同。值得提取的是加速度沿曲面法向的分量，因为它只由这一方向和曲面的二阶数据决定。

\begin{proposition}[曲面曲线的法向加速度]\label{g1:prop:surface-normal-acceleration}
	设 \(X:U\to\mathbb R^3\) 为光滑浸入，选定单位法向 \(n\)，并设 \(\gamma(t)=X(u(t))\) 是其上的 \(C^2\) 正则曲线。则
	\[
		\langle\ddot\gamma,n\rangle=\mathrm{II}(\dot\gamma,\dot\gamma).
	\]
	曲线在一点沿切方向 \(v\ne0\) 的单位速率法向加速度为
	\begin{equation}\label{g1:eq:normal-curvature-quotient}
		k_n(v)\coloneqq\frac{\mathrm{II}(v,v)}{\langle v,v\rangle}.
	\end{equation}
	它称为该方向的\term[faqulv]{法曲率}[normal curvature]，与经过该点且具有这一切方向的曲线选择无关，也与把 \(v\) 乘以任意非零实数无关。
\end{proposition}
\begin{proof}
	链式法则给出
	\[
		\ddot\gamma=\sum_iX_i\ddot u^i+\sum_{i,j}X_{ij}\dot u^i\dot u^j.
	\]
	第一项与 \(n\) 正交，第二项的法向分量由\eqref{g1:eq:second-form-coefficients}等于 \(\sum_{i,j}h_{ij}\dot u^i\dot u^j\)。按\cref{g1:prop:arc-length-param}把这条正则曲线改用弧长 \(s\)，令 \(v_0=\lVert\dot\gamma\rVert\)，则
	\[
		\gamma''=\frac{\ddot\gamma}{v_0^2}-\frac{\dot v_0}{v_0^3}\dot\gamma.
	\]
	第二项仍为切向，所以 \(\langle\gamma'',n\rangle=\mathrm{II}(\dot\gamma,\dot\gamma)/v_0^2\)。二次型的分子分母在向量缩放时乘上同一个平方因子，得到其余断言。
\end{proof}

法曲率可以为负，它不是空间曲线的非负曲率 \(\kappa\)。对单位速率曲线，\(\kappa=\lVert\gamma''\rVert\)，而 \(k_n=\langle\gamma'',n\rangle\)，故 \(|k_n|\le\kappa\)；曲线还可能在切平面内转弯。若只想观察曲面沿方向 \(v\) 怎样弯曲，就用经过 \(p\)、由 \(v,n\) 张成的平面去截曲面，得到\term[fajiexian]{法截线}[normal section]。借助\cref{g1:prop:surface-local-graph}，先把切平面写成 \(z=0\)，再把 \(v\) 作为 \(x\) 方向，曲面成为 \(z=f(x,y)\)、\(\mathrm df(p)=0\)，而法截线为 \((x,0,f(x,0))\)，因此它在 \(p\) 确实正则且沿 \(v\)。其单位速率加速度在 \(p\) 沿 \(n\)，有向大小就是 \(k_n(v)\)。

\begin{theorem}[主曲率与 Euler 公式]\label{g1:thm:euler-normal-curvature}
	在带选定单位法向的光滑浸入曲面的每一点，形算子有一组正交单位特征向量 \(e_1,e_2\)。相应实特征值 \(k_1,k_2\) 称为该点的\term[zhuqulv]{主曲率}[principal curvatures]，其特征直线称为\term[zhufangxiang]{主方向}[principal directions]。若 \(v=\cos\theta\,e_1+\sin\theta\,e_2\)，则
	\begin{equation}\label{g1:eq:euler-normal-curvature}
		k_n(v)=k_1\cos^2\theta+k_2\sin^2\theta.
	\end{equation}
	这个等式称为\term[Eulerfaqulvgongshi]{Euler 法曲率公式}[Euler's formula for normal curvature]。法曲率在所有单位方向上的最小值和最大值恰为两个主曲率。
\end{theorem}
\begin{proof}
	这里可以直接证明所需的二维自伴算子结论。函数 \(q(v)=\langle Sv,v\rangle\) 在切平面的单位圆上连续，故有最大值点 \(e_1\)。取垂直于 \(e_1\) 的单位向量 \(e_2\)，对 \(q(\cos t\,e_1+\sin t\,e_2)\) 在 \(t=0\) 求导，极值条件给出
	\[
		0=\langle Se_2,e_1\rangle+\langle Se_1,e_2\rangle=2\langle Se_1,e_2\rangle.
	\]
	故 \(Se_1=k_1e_1\)。自伴性又给出 \(\langle Se_2,e_1\rangle=0\)，所以 \(Se_2=k_2e_2\)。把 \(v\) 按这组正交单位基展开，交叉项消失，便得\eqref{g1:eq:euler-normal-curvature}。右端是 \(k_1,k_2\) 的非负加权平均，在 \(e_1,e_2\) 处取到两个端值，因此它们就是最大值与最小值。
\end{proof}

主方向首先是一条无定向的直线：\(e_i\) 和 \(-e_i\) 给出同一个方向。若 \(k_1\ne k_2\)，恰有两条彼此正交的主方向；若 \(k_1=k_2\)，则 \(S=k_1I\)，每条方向都是主方向。不能在后一种情况下还声称曲面指定了唯一的一对主方向。

\begin{example}[一张鞍面上的不同法截线]\label{g1:exm:saddle-normal-curvature}
	对图像 \(z=x^2-\tfrac12y^2\)，原点取向上法向，有 \(S=\operatorname{diag}(2,-1)\)，两个坐标方向就是主方向。沿单位方向 \(v=(\cos\theta,\sin\theta,0)\)，法曲率为 \(2\cos^2\theta-\sin^2\theta\)。它在两个坐标方向上分别为 \(2,-1\)，在 \(\tan^2\theta=2\) 的两条无向直线上为零。\cref{g1:fig:normal-sections}画出三种法截线；零法曲率方向在此例中对应直线，但一般只能说明法向加速度在该点为零。
\end{example}

\input{Figures/Chapter08-NormalSections.tex}

\subsection{平均曲率、Gauss 曲率及脐点}
\label{g1:subsec:0802:03}

两个主曲率交换编号时，和与积都不改变。把它们用于刻画曲面，可以避免不断选择和排列主方向。和反映两个正交方向的平均弯曲，积则区分同向弯曲、反向弯曲与二阶退化。

\begin{definition}[平均曲率、Gauss 曲率与脐点]\label{g1:def:surface-curvatures}
	在带选定单位法向的光滑浸入曲面上，定义\term[pingjunqulv]{平均曲率}[mean curvature]与\term[Gaussqulv]{Gauss 曲率}[Gaussian curvature]为
	\[
		H\coloneqq\tfrac12\operatorname{tr}S=\tfrac12(k_1+k_2),\qquad K\coloneqq\det S=k_1k_2.
	\]
	若某点 \(k_1=k_2\)，则称为\term[qidian]{脐点}[umbilic point]，等价地，该点 \(S=HI\)。若 \(S=0\)，称其为\term[pingdian]{平点}[planar point]；平点是两个主曲率都为零的脐点。
	\symidx{\(H\)：平均曲率}
	\symidx{\(K\)：曲面的 Gauss 曲率}
\end{definition}

这里的 \(1/2\) 属于平均曲率的约定，有些文献把主曲率之和直接叫作平均曲率。我们后续的公式均使用上面的平均值。Gauss 曲率目前是通过曲面在三维空间中的法向变化定义的；它是否也能仅由第一基本形式恢复，是下一节要证明的问题。

\begin{proposition}[曲率的计算与法向选择]\label{g1:prop:surface-curvature-formulas}
	设 \(X:U\to\mathbb R^3\) 是光滑浸入，选定单位法向，并使用第一、第二基本形式的系数 \(E,F,G,h_{ij}\)。则
	\begin{equation}\label{g1:eq:surface-HK-formulas}
		K=\frac{h_{11}h_{22}-h_{12}^2}{D},\qquad
		H=\frac{Gh_{11}-2Fh_{12}+Eh_{22}}{2D},\qquad D=EG-F^2.
	\end{equation}
	有 \(H^2-K\ge0\)，两个主曲率是 \(H\pm\sqrt{H^2-K}\)，脐点恰由 \(H^2=K\) 刻画。把法向改为 \(-n\) 时，\(S,\mathrm{II},H\) 同时变号，两个主曲率各自变号，\(K\) 与主方向不变。
\end{proposition}
\begin{proof}
	由\eqref{g1:eq:shape-matrix}，
	\[
		[S]=\frac1D
		\begin{pmatrix}
			Gh_{11}-Fh_{12}&Gh_{12}-Fh_{22}\\
			Eh_{12}-Fh_{11}&Eh_{22}-Fh_{12}
		\end{pmatrix}.
	\]
	取行列式和迹即得\eqref{g1:eq:surface-HK-formulas}。又
	\[
		H^2-K=\tfrac14(k_1-k_2)^2,\qquad
		\det(\lambda I-S)=\lambda^2-2H\lambda+K,
	\]
	所以平方根为实数，二次方程的根给出主曲率，等根恰为脐点。法向翻转后 \(S=-\mathrm dn\) 变号，其余各项由定义得到；若主曲率原先按大小排序，翻转后还须交换顺序。
\end{proof}

例如，外向法向的半径 \(R\) 球面处处有 \(H=-1/R\)、\(K=1/R^2\)，且每一点都是脐点。平面取任一固定单位法向则有 \(S=0\)。与此不同，图像 \(z=x^3-3xy^2\) 只在原点有二阶偏导全部为零，所以原点是平点；在其余位置，二阶偏导矩阵的行列式为 \(-36(x^2+y^2)\)，由图像公式及\eqref{g1:eq:surface-HK-formulas}知 \(K<0\)。因此“一点的第二基本形式为零”不意味着附近的曲面是一张平面。

\ProseParagraphBreak

平均曲率还真的是各方向法曲率的角平均。由 Euler 公式，
\[
	\frac1{2\pi}\int_0^{2\pi}k_n(\cos\theta\,e_1+\sin\theta\,e_2)\,\mathrm d\theta
	=\frac{k_1+k_2}2=H.
\]
平均值为零时，各方向的弯曲可能恰好抵消，不能据此断言每个方向都不弯曲。例如上一小节的鞍面若改为 \(z=x^2-y^2\)，则原点 \(H=0\) 而 \(K=-4\)。

\subsection{曲率线、渐近线及其退化}
\label{g1:subsec:0802:04}

一点上的方向能否沿曲面连成曲线，是另一个问题。先规定要沿哪些方向行走，再讨论方向能否光滑选取以及微分方程是否有解；这两步在脐点或曲率变号处都可能发生变化。

\begin{definition}[曲率线与渐近线]\label{g1:def:principal-asymptotic-curves}
	在带选定单位法向的光滑浸入曲面上，若一条光滑正则曲线的切向量处处为形算子的特征向量，称其为\term[qulvxian]{曲率线}[line of curvature]。若非零切向量 \(v\) 满足 \(\mathrm{II}(v,v)=0\)，称它张成的无向直线为\term[jianjinfangxiang]{渐近方向}[asymptotic direction]；切向量处处为渐近方向的光滑正则曲线称为\term[jianjinxian]{渐近线}[asymptotic curve]。
\end{definition}

法向翻转不改变特征方向，也不改变 \(\mathrm{II}(v,v)=0\) 的解，因此这两类曲线可以不依赖整体法向选择来讨论。曲率线要求法向变化 \(\mathrm dn(v)\) 与行进方向平行；渐近线要求曲线的法向加速度为零。两者是不同条件。

\begin{lemma}[局部积分曲线]\label{g1:lem:surface-field-integral-curve}
	设 \(V:U\to\mathbb R^2\) 是开集 \(U\subset\mathbb R^2\) 上的光滑向量场，\(u_0\in U\)。则在某个 \((-\varepsilon,\varepsilon)\) 上，初值问题
	\[
		\dot u(t)=V(u(t)),\qquad u(0)=u_0
	\]
	有唯一的局部光滑解。若 \(V(u_0)\ne0\)，缩小区间后该解正则。
\end{lemma}
\begin{proof}
	取闭球 \(\overline B(u_0,r)\subset U\)。由于 \(V\) 及其一阶导数在闭球上有界，存在常数 \(M,L\ge0\)，使 \(\lVert V\rVert\le M\) 且 \(\lVert V(x)-V(y)\rVert\le L\lVert x-y\rVert\)。选 \(a>0\) 足够小，使 \(aM<r\)、\(q=aL<1\)。在所有取值于此闭球的连续路径 \(u:[-a,a]\to\overline B(u_0,r)\) 上定义
	\[
		(\mathcal Tu)(t)=u_0+\int_0^t V(u(s))\,\mathrm ds.
	\]
	估计 \(aM<r\) 保证 \(\mathcal T\) 仍取值于闭球，而一致范数满足
	\[
		\lVert\mathcal Tu-\mathcal Tw\rVert_\infty\le q\lVert u-w\rVert_\infty.
	\]
	从常路径 \(u_0\) 迭代 \(u_{m+1}=\mathcal Tu_m\)，相邻两次的差至多按 \(q^m\) 衰减。几何级数收敛说明 \(u_m\) 一致收敛；连续路径的一致极限仍连续且取值于闭球。把极限送入积分，得到 \(u=\mathcal Tu\)，所以 \(u\) 为 \(C^1\) 且满足方程。若 \(u,w\) 都是这个积分方程的解，同一估计给出 \(\lVert u-w\rVert_\infty\le q\lVert u-w\rVert_\infty\)，因而两者相同。任意另一局部解在足够短的公共区间内也留在闭球中，这就给出局部唯一性。由 \(\dot u=V(u)\) 逐次求导，\(u\) 为光滑；最后用 \(V(u_0)\ne0\) 及连续性保证附近速度不消失。
\end{proof}

\begin{proposition}[主方向与渐近方向的局部曲线]\label{g1:prop:principal-asymptotic-lines}
	设 \(X:U\to\mathbb R^3\) 为光滑浸入。
	\begin{enumerate}
		\item 每个非脐点附近有两族局部曲率线；经过给定点、沿指定主方向的局部曲率线唯一到正则重参数化为止。
		\item 每个 \(K<0\) 的点附近有两族局部渐近线；经过给定点、沿指定渐近方向的局部渐近线唯一到正则重参数化为止。
		\item 在一点，若 \(K>0\)，没有渐近方向；若 \(K<0\)，有两条不同的渐近方向；若 \(K=0\) 且 \(S\ne0\)，恰有一条渐近方向，为 \(\ker S\)；若 \(S=0\)，所有方向都是渐近方向。
	\end{enumerate}
\end{proposition}
\begin{proof}
	在非脐点附近，\(H^2-K>0\)。所以
	\[
		k_+=H+\sqrt{H^2-K},\qquad k_-=H-\sqrt{H^2-K}
	\]
	是光滑的不同特征值。相应的线性算子
	\[
		P_+=\frac{S-k_-I}{k_+-k_-},\qquad P_-=\frac{S-k_+I}{k_--k_+}
	\]
	分别投影到两条特征直线上。先取一个光滑切向量场 \(W\)，使 \(P_+W\) 在指定点不为零；例如把该点的一个特征向量按坐标基用常系数延伸。缩小邻域后 \(P_+W\) 始终不为零，归一化得到光滑单位特征向量场 \(e_+\)。同理得到 \(e_-\)，并可按需要调整其符号。把 \(e_+\) 在 \(X_1,X_2\) 下的系数看作 \(U\) 上的非零光滑向量场，应用\cref{g1:lem:surface-field-integral-curve}并复合 \(X\)，即得所需曲率线。另一主方向同样处理。任意沿同一方向的正则曲线，缩短后可调整方向并改用弧长，使其导数等于所选单位场；初值唯一性就给出无参数轨迹的唯一性。

	当 \(K<0\) 时，\(k_+>0>k_-\)。在上面选出的光滑正交单位场中，二次型的零方向由
	\[
		k_+a^2+k_-b^2=0
	\]
	确定，分别由两个光滑非零场
	\[
		\sqrt{-k_-}\,e_++\sqrt{k_+}\,e_-,
		\qquad
		\sqrt{-k_-}\,e_+-\sqrt{k_+}\,e_-
	\]
	张成。归一化后再次应用局部积分曲线引理，得到存在与唯一性。

	最后，在任一点用主方向正交基把 \(\mathrm{II}\) 写成 \(k_1a^2+k_2b^2\)。若两系数同号，此式对非零向量不为零；若异号，得到上面的两条直线；若恰有一个系数为零，解是一条坐标轴；若两者都为零，每个向量都是解。这证明了第三项。
\end{proof}

这份分类只需解一个二元二次方程，可整理为\cref{g1:tab:asymptotic-directions}。其中 \(K=0\) 的两行必须分开：仅凭 \(K\) 的值，无法判断零方向是一条还是所有方向。

\begin{table}[htbp]
	\centering
	\caption{第二基本形式与渐近方向}
	\label{g1:tab:asymptotic-directions}
	\begin{tabular}{lll}
		\toprule
		条件 & 主方向基下的二次型 & 渐近方向 \\
		\midrule
		\(K>0\) & 两个非零系数同号 & 无 \\
		\(K<0\) & 两个系数异号 & 两条不同直线 \\
		\(K=0,\ S\ne0\) & 恰有一个系数为零 & 一条直线 \\
		\(S=0\) & 零二次型 & 所有方向 \\
		\bottomrule
	\end{tabular}
\end{table}

在一整片 \(K=0,\ S\ne0\) 的区域上，零特征值与另一个特征值仍然分离，上面的投影构造同样给出一条光滑方向场及其积分曲线。但若只知道某一点 \(K=0\)，不能据此断言有渐近线穿过它。例如 \(z=x^2+y^4\) 在原点的形算子为 \(\operatorname{diag}(2,0)\)，只有 \(y\) 轴方向渐近；在 \(y\ne0\) 处却有 \(K>0\)，不存在渐近方向。若有正则渐近线通过原点，其速度必须沿 \(y\) 轴且不为零，于是附近离开原点后 \(y\ne0\)，矛盾。因此原点的渐近方向不能延续成穿过该点的渐近线。

\ProseParagraphBreak

脐点处的困难则是主方向不再分离。球面的每条正则曲线都是曲率线，经过一点和一条切向的曲线远不唯一。即使避开脐点，在一张较大的曲面上也不能从局部单位特征向量的选择自动得到全局连续选择：沿闭路延伸后，向量可能反向，而它张成的直线仍相同。曲率线和渐近线的定义使用无向直线，正是为了把这种方向选择与曲线本身区别开。

% !TeX root = ../../Book1.tex
% 纲要标题：### 8.3 结构方程与内在曲率
% 纲要位置：几何学计划纲要.md，第 416 行。
% 正文状态：初稿；尚未正式审核。

% 纲要内容开始
% 
% * Gauss–Weingarten 方程
% * 协变导数与 Christoffel 系数
% * Gauss 方程与绝妙定理
% * Codazzi 方程及混合偏导相容性
% 
% 纲要内容结束

\section{结构方程与内在曲率}
\label{g1:sec:0803}

第一基本形式记录切平面中的长度，第二基本形式记录曲面向法向弯曲的程度。它们看起来各管一件事，但不能任意指定：如果它们来自同一个位置函数 \(X(u,v)\)，先对 \(u\) 求导再对 \(v\) 求导，必须与相反的顺序得到相同结果。这一简单要求将同时给出 Gauss 方程和 Codazzi 方程。Gauss 方程还将把主曲率的乘积写成第一基本形式及其导数的表达式，将法向弯曲与内在度量联系起来。

\ProseParagraphBreak

本节取光滑浸入 \(X\colon U\rightarrow\mathbb R^3\)，其中 \(U\subseteq\mathbb R^2\) 开，记 \(u^1=u\)、\(u^2=v\)、\(X_i=\partial_iX\)。沿用 \(g_{ij}=\langle X_i,X_j\rangle\)、\(D=\det(g_{ij})>0\) 与选定的单位法向 \(n\)，并记 \(h_{ij}=\langle X_{ij},n\rangle\)。向量场及其导数都按参数点定义；若像有自交，不把不同参数点上的切向量混为一谈。指标 \(i,j,k,\ell,m\) 均取 \(1,2\)；在同一项中重复出现的一上一下指标表示对这两个值求和。这个简写只省去求和号，不改变矩阵的乘法顺序。

\subsection{Gauss–Weingarten 方程}
\label{g1:subsec:0803:01}

前一节已经把 \(X_{ij}\) 的法向分量写成 \(h_{ij}n\)。剩下的切向分量也需要记录，因为坐标向量 \(X_1,X_2\) 一般随位置改变，既不保持长度，也不保持夹角。令 \((g^{ij})\) 为 \((g_{ij})\) 的逆矩阵，并写 \(h_i{}^k=g^{k\ell}h_{i\ell}\)。

\begin{proposition}[Gauss–Weingarten 方程]
	\label{g1:prop:gauss-weingarten}
	设 \(U\subseteq\mathbb R^2\) 开，\(X:U\rightarrow\mathbb R^3\) 是光滑浸入，\(n\) 为光滑单位法向。记 \(X_i=\partial_iX\)、\(g_{ij}=\langle X_i,X_j\rangle\)、\(h_{ij}=\langle\partial_i\partial_jX,n\rangle\)、\((g^{ij})=(g_{ij})^{-1}\)，并令 \(h_i{}^k=\sum_\ell g^{k\ell}h_{i\ell}\)。重复指标取 \(1,2\) 并求和。存在唯一的光滑函数 \(\Gamma^k_{ij}\)，使
	\begin{equation}
		\label{g1:eq:gauss-weingarten}
		\partial_iX_j=\Gamma^k_{ij}X_k+h_{ij}n,\qquad
		\partial_i n=-h_i{}^kX_k.
	\end{equation}
	其中 \(\Gamma^k_{ij}=\Gamma^k_{ji}\)，并且
	\begin{equation}
		\label{g1:eq:christoffel-metric}
		\Gamma^k_{ij}
		=\frac12g^{k\ell}\bigl(\partial_i g_{j\ell}+\partial_j g_{i\ell}-\partial_\ell g_{ij}\bigr).
	\end{equation}
	因此 \(X_{ij}\) 的切向分解系数完全由第一基本形式决定。
\end{proposition}

\begin{proof}
	浸入条件保证 \(X_1,X_2,n\) 线性无关，所以向量 \(X_{ij}\) 有唯一的上述分解。混合偏导相等给出 \(\Gamma^k_{ij}=\Gamma^k_{ji}\)。对 \(g_{j\ell}=\langle X_j,X_\ell\rangle\) 求导，得到
	\[
		\partial_i g_{j\ell}=\langle X_{ij},X_\ell\rangle+\langle X_j,X_{i\ell}\rangle.
	\]
	把交换 \(i,j,\ell\) 所得的三个式子按前两个相加、第三个相减，交叉项抵消，故
	\[
		2\langle X_{ij},X_\ell\rangle
		=\partial_i g_{j\ell}+\partial_j g_{i\ell}-\partial_\ell g_{ij}
		=2\Gamma^k_{ij}g_{k\ell}.
	\]
	乘逆矩阵便得\eqref{g1:eq:christoffel-metric}。

	由 \(\langle n,n\rangle=1\) 可知 \(\partial_i n\) 是切向量；对 \(\langle n,X_j\rangle=0\) 求导，又有
	\[
		\langle\partial_i n,X_j\rangle=-\langle n,X_{ij}\rangle=-h_{ij}.
	\]
	若 \(\partial_i n=a^kX_k\)，则 \(a^kg_{kj}=-h_{ij}\)，即 \(a^k=-g^{k\ell}h_{i\ell}\)。
\end{proof}

\eqref{g1:eq:gauss-weingarten}中的第一式称为 Gauss 公式，第二式称为 Weingarten 公式；两式合称\term[gaussweingartenfangcheng]{Gauss–Weingarten 方程}[Gauss--Weingarten equations]。第二式正是形算子 \(S=-\mathrm{d}n\) 的坐标表达；负号与\cref{g1:def:shape-operator}的约定一致。\footnote{Weingarten，Julius Weingarten（1836--1910），德国数学家，研究曲面等距变形与曲率线。}

\ProseParagraphBreak

这里的标架 \((X_1,X_2,n)\) 通常不是正交单位标架。因此 \(\Gamma^k_{ij}\) 无须像\eqref{g1:eq:frenet-matrix}中的系数矩阵那样斜对称。其作用是补偿坐标基自身的变化；这个变化有些来自曲面的弯曲，有些仅来自选了弯曲的坐标线。

\subsection{协变导数与 Christoffel 系数}
\label{g1:subsec:0803:02}

设 \(V=V^iX_i\) 是沿曲面的光滑切向量场。把它当作取值于 \(\mathbb R^3\) 的函数求导，\(\partial_iV\) 未必仍在切平面内。例如球面上的切向量沿曲线变化时，普通导数可能含有径向分量。如果我们只记录曲面自身的切向变化，就应把普通导数投影回当前切平面。

\begin{definition}
	\label{g1:def:surface-covariant-derivative}
	设 \(U\subseteq\mathbb R^2\) 开，\(X:U\rightarrow\mathbb R^3\) 为光滑浸入，\(X_i=\partial_iX\)，\(V=V^iX_i\)、\(W=W^iX_i\) 为沿 \(X\) 的光滑切向量场，\(f\in C^\infty(U)\)，\(\Gamma^k_{ij}\) 为诱导度量的 Christoffel 系数。定义它沿坐标方向的\term[xiebiandaoshu]{协变导数}[covariant derivative]为
	\[
		\nabla_iV\coloneqq(\partial_iV)^\top
		=\bigl(\partial_iV^k+\Gamma^k_{ij}V^j\bigr)X_k,
	\]
	其中 \({}^\top\) 表示向切平面的正交投影。若 \(W=W^iX_i\)，则定义 \(\nabla_WV=W^i\nabla_iV\)。函数的方向导数记为 \(W(f)=W^i\partial_i f\)。\symidx{\(\nabla_WV\)：切向量场的协变导数}
\end{definition}

上式中的 \(\Gamma^k_{ij}\) 称为\term[christoffelxishu]{Christoffel 系数}[Christoffel symbols]，也称 Christoffel 符号。\footnote{Christoffel，Elwin Bruno Christoffel（1829--1900），德国数学家，研究二次微分形式、张量演算与共形映射。}在这里“协变”所强调的是：重参数化后，计算仍给出同一个切向量，而非要求各个数 \(\Gamma^k_{ij}\) 保持不变。定义中的正交投影及方向导数都是几何操作，已经说明其不依赖坐标；公式中新增的 \(\Gamma^k_{ij}V^j\) 恰好补上了基向量求导的部分。\symidx{\(\Gamma^k_{ij}\)：Christoffel 系数}

\begin{example}[平面也可以有非零的 Christoffel 系数]
	\label{g1:exm:polar-christoffel}
	在平面的极坐标参数 \(X(r,\theta)=(r\cos\theta,r\sin\theta,0)\) 中，取 \(r>0\) 及一个适当的角度区间，有 \(g_{rr}=1\)、\(g_{r\theta}=0\)、\(g_{\theta\theta}=r^2\)。代入\eqref{g1:eq:christoffel-metric}得到仅有的非零系数
	\[
		\Gamma^r_{\theta\theta}=-r,\qquad
		\Gamma^\theta_{r\theta}=\Gamma^\theta_{\theta r}=\frac1r.
	\]
	例如 \(\partial_\theta X_\theta=-rX_r\) 只是平面内的径向加速度，并不意味着平面发生了法向弯曲。在直角坐标下，同一平面的全部 Christoffel 系数都为零。
\end{example}

这种现象说明 \(\Gamma^k_{ij}\) 不是一个张量的分量：如果张量在某组基下全部分量为零，换基后仍应全部为零。更直接地，设新坐标为 \(y^a\)，旧坐标为 \(u^i=u^i(y)\)，对 \(\widetilde X(y)=X(u(y))\) 求两次导数，并比较切向分量，得到
\[
	\widetilde\Gamma^c_{ab}
	=\frac{\partial y^c}{\partial u^k}
	\left(
		\Gamma^k_{ij}\frac{\partial u^i}{\partial y^a}\frac{\partial u^j}{\partial y^b}
		+\frac{\partial^2u^k}{\partial y^a\partial y^b}
	\right).
\]
最后的二阶导数项不是逐点线性换基产生的项。协变导数整体仍然与坐标无关，因为向量分量 \(V^i\) 的导数也发生变化，恰好抵消这项。

\ProseParagraphBreak

为了比较先沿两个方向求导的结果，需要注意：任意两个向量场的方向导数未必可交换。定义它们的\term[xiangliangchangkuohao]{向量场括号}[Lie bracket]为
\[
	[W,V]=\bigl(W^i\partial_iV^k-V^i\partial_iW^k\bigr)X_k.
\]
对每个光滑函数 \(f\)，直接展开可得 \([W,V](f)=W(V(f))-V(W(f))\)，其中二阶偏导项成对抵消。这也给出了括号与坐标无关的含义。这一运算也称为 Lie 括号。\footnote{Lie，Marius Sophus Lie（1842--1899），挪威数学家，创立连续变换群理论，并研究其与微分方程的联系。}特别地，坐标向量场满足 \([X_1,X_2]=0\)；这并不是说作为 \(\mathbb R^3\) 值函数的 \(X_1,X_2\) 是常向量。

\begin{proposition}[切向求导的基本性质]
	\label{g1:prop:surface-connection-properties}
	设 \(U\subseteq\mathbb R^2\) 开，\(X:U\to\mathbb R^3\) 为光滑浸入，\(\nabla_WV\) 由普通方向导数沿 \(W\) 求导后向切平面正交投影得到。此协变导数对向量场可加，且对任意光滑函数 \(f\) 及切向量场 \(V,W,Z\)，有
	\begin{align*}
		\nabla_{fW}V&=f\nabla_WV,&
		\nabla_W(fV)&=W(f)V+f\nabla_WV,\\
		W\langle V,Z\rangle&=\langle\nabla_WV,Z\rangle+\langle V,\nabla_WZ\rangle,&
		\nabla_WV-\nabla_VW&=[W,V].
	\end{align*}
	后两条分别称为\term[duliangxiangrong]{度量相容}[metric compatibility]和\term[wunao]{无挠性}[torsion-freeness]。更一般地，在开集 \(U\subseteq\mathbb R^2\) 上给定任意光滑正定矩阵场 \(g=(g_{ij})\)，令 \((g^{ij})=(g_{ij})^{-1}\)，按 Christoffel 公式\eqref{g1:eq:christoffel-metric}，即 \(\Gamma^k_{ij}=\tfrac12g^{k\ell}(\partial_i g_{j\ell}+\partial_j g_{i\ell}-\partial_\ell g_{ij})\)，及 \(\nabla_WV=W^i(\partial_iV^k+\Gamma^k_{ij}V^j)\partial_k\) 定义的运算也满足这两条性质，而且它是满足这些性质与前两条求导规则的唯一运算。
\end{proposition}

\begin{proof}
	前两条由普通求导的线性及乘法法则得到。对于诱导度量，普通内积求导后出现的法向分量与 \(V,Z\) 垂直，所以投影不改变结果，得到度量相容。再展开 \(\nabla_WV-\nabla_VW\)，由于 \(\Gamma^k_{ij}=\Gamma^k_{ji}\)，所有含 \(\Gamma\) 的项相消，余项就是括号。

	对于任意给定的正定 \(g\)，公式\eqref{g1:eq:christoffel-metric}仍给出 \(\Gamma^k_{ij}=\Gamma^k_{ji}\)，所以无挠性的证明不变。另有
	\[
		\Gamma^k_{ij}g_{k\ell}+\Gamma^k_{i\ell}g_{jk}
		=\partial_i g_{j\ell},
	\]
	因为将右式中的两个 \(\Gamma\) 换成定义后，其余偏导项互相抵消。这正是对坐标向量场的度量相容，乘以向量场分量并用乘法法则即得一般情形。

	最后，设另一运算满足相同规则，写其坐标系数为 \(C^k_{ij}\)。无挠性给出 \(C^k_{ij}=C^k_{ji}\)，度量相容给出 \(\partial_i g_{j\ell}=C^k_{ij}g_{k\ell}+C^k_{i\ell}g_{jk}\)。如\cref{g1:prop:gauss-weingarten}的证明那样将三个式子相加减，就强制得到 \(C^k_{ij}=\Gamma^k_{ij}\)。
\end{proof}

这一运算也称为度量的\term[lianluo]{联络}[connection]。在只给定 \(g\) 而尚未构造曲面时，我们把坐标基写作 \(\partial_1,\partial_2\)，完全可以照同一公式求导。重参数化后，度量相容、无挠性和两条求导规则都保持；上面的唯一性说明，各套坐标得到的是同一个运算。因此，“把普通导数投影回去”是诱导度量的直观构造，“由 \(g\) 的系数求出 \(\Gamma\)”则使这个构造不再依赖于周围的三维空间。

\subsection{Gauss 方程与绝妙定理}
\label{g1:subsec:0803:03}

普通二阶偏导可交换，协变二阶导数却是在每次求导后都投影一次，两次中间使用的切平面不同，结果可能不同。若方向本身也是变化的向量场，还必须扣除它们的括号造成的差别。扣除这一差别后，结果才能只依赖给定点上的方向。

\begin{definition}
	\label{g1:def:surface-curvature-tensor}
	设 \((M,g)\) 为光滑度量曲面，\(\nabla\) 为 \(g\) 的无挠相容协变导数。对任意光滑切向量场 \(V,W,Z\)，定义
	\[
		R(V,W)Z\coloneqq\nabla_V\nabla_WZ-\nabla_W\nabla_VZ-\nabla_{[V,W]}Z.
	\]
	这个运算称为\term[qulvzhangliang]{曲率张量}[curvature tensor]。本书固定上述符号顺序。\symidx{\(R(V,W)Z\)：曲率张量}
\end{definition}

\begin{proposition}
	\label{g1:prop:curvature-tensorial}
	设 \(M\) 为光滑曲面，\(\nabla\) 为其切丛上的光滑联络，\(V,W,Z\) 为光滑切向量场。令 \(R(V,W)Z=\nabla_V\nabla_WZ-\nabla_W\nabla_VZ-\nabla_{[V,W]}Z\)，此处不要求联络无挠或与度量相容。则 \(R(V,W)Z\) 对三个向量场均为逐点线性：任意光滑函数 \(f\) 均满足
	\begin{align*}
		R(fV,W)Z&=fR(V,W)Z,\\
		R(V,fW)Z&=fR(V,W)Z,\\
		R(V,W)(fZ)&=fR(V,W)Z.
	\end{align*}
	因此在点 \(p\) 的值只依赖 \(V(p),W(p),Z(p)\)，不依赖这些向量在邻域中的延拓。
\end{proposition}

\begin{proof}
	先由括号的坐标式得到 \([fV,W]=f[V,W]-W(f)V\)。代入定义，第二项求导产生的 \(-W(f)\nabla_VZ\) 正好与括号项中的 \(W(f)\nabla_VZ\) 抵消，证明第一式。由 \(R(W,V)=-R(V,W)\) 得第二式。

	对第三式，把两次乘法法则代入定义。含 \(\nabla_VZ,\nabla_WZ\) 的一阶导数项成对抵消，剩余的额外项只有
	\[
		\bigl(V(W(f))-W(V(f))-[V,W](f)\bigr)Z=0.
	\]
	故第三式成立。将三个向量场各自在坐标基中展开，所得值是各自分量在 \(p\) 的乘积乘以 \(R(X_i,X_j)X_k\) 在 \(p\) 的值，所以确实只需要三个输入向量的点值。
\end{proof}

这里“张量”的具体含义是：在每一点，这个多线性映射都把三个切向量映为一个切向量。

\ProseParagraphBreak

由于 \([X_1,X_2]=0\)，计算坐标分量只须展开两次求导：
\begin{equation}
	\label{g1:eq:surface-curvature-components}
	\begin{aligned}
		R(X_1,X_2)X_j&=R^k{}_{j12}X_k,\\
		R^k{}_{j12}
		&=\partial_1\Gamma^k_{2j}-\partial_2\Gamma^k_{1j}
		+\Gamma^k_{1m}\Gamma^m_{2j}-\Gamma^k_{2m}\Gamma^m_{1j}.
	\end{aligned}
\end{equation}
例如 \(\nabla_1(\Gamma^m_{2j}X_m)\) 同时对系数及基向量求导，才得到偏导项和两个系数的乘积。遗漏后一部分，就会把坐标基变化误当作曲率。

\begin{theorem}[Gauss 方程]
	\label{g1:thm:gauss-equation}
	设 \(U\subseteq\mathbb R^2\) 开，\(X:U\rightarrow\mathbb R^3\) 是光滑浸入，取光滑单位法向 \(n\)，形算子为 \(S=-\mathrm{d}n\)，第二基本形式为 \(h(V,W)=\langle SV,W\rangle\)。以 \(\nabla\) 记诱导度量的无挠相容协变导数，定义 \(R(V,W)Z=\nabla_V\nabla_WZ-\nabla_W\nabla_VZ-\nabla_{[V,W]}Z\)，其中 \(V,W,Z\) 为光滑切向量场。令 \(X_i=\partial_iX\)、\(g_{ij}=\langle X_i,X_j\rangle\)、\((g^{ij})=(g_{ij})^{-1}\)、\(h_{ij}=h(X_i,X_j)\)、\(h_i{}^k=g^{k\ell}h_{i\ell}\)，并以 \(R(X_1,X_2)X_j=R^k{}_{j12}X_k\) 定义曲率分量。则曲率满足
	\begin{equation}
		\label{g1:eq:gauss-equation}
		R(V,W)Z=h(W,Z)SV-h(V,Z)SW.
	\end{equation}
	特别地，
	\[
		R^k{}_{j12}=h_{2j}h_1{}^k-h_{1j}h_2{}^k,\qquad
		\langle R(X_1,X_2)X_2,X_1\rangle=h_{11}h_{22}-h_{12}^2.
	\]
\end{theorem}

\begin{proof}
	记 \(\mathrm{D}_V\) 为对 \(\mathbb R^3\) 值函数逐分量作方向导数。由定义及\cref{g1:prop:gauss-weingarten}，任意切向量场满足
	\[
		\mathrm{D}_VZ=\nabla_VZ+h(V,Z)n,\qquad \mathrm{D}_Vn=-SV.
	\]
	Euclid 空间中的逐分量求导满足
	\[
		\mathrm{D}_V\mathrm{D}_WZ-\mathrm{D}_W\mathrm{D}_VZ-\mathrm{D}_{[V,W]}Z=0,
	\]
	因为对每个分量这正是括号的定义。将前两个分解式代入，并只取切向分量：
	\[
		0=\nabla_V\nabla_WZ-h(W,Z)SV
		-\nabla_W\nabla_VZ+h(V,Z)SW-\nabla_{[V,W]}Z.
	\]
	移项便得 Gauss 方程。最后取 \(V=X_1,W=Z=X_2\)，再与 \(X_1\) 作内积，使用 \(h(X_i,X_j)=h_{ij}\)，即可得到坐标结论。
\end{proof}

上面的等式称为\term[gaussfangcheng]{Gauss 方程}[Gauss equation]。左边只由度量求导得到，右边却是法向弯曲的数据。由\cref{g1:prop:surface-curvature-formulas}，\(K=\det(h_{ij})/\det(g_{ij})\)，因此立即得到
\begin{equation}
	\label{g1:eq:intrinsic-gauss-curvature}
	K=\frac{\langle R(X_1,X_2)X_2,X_1\rangle}{D}
	=\frac{g_{1k}R^k{}_{212}}{D}.
\end{equation}
其右端只使用 \(g_{ij}\)、它们的逆矩阵和一二阶偏导数。

\begin{theorem}[Gauss 绝妙定理]
	\label{g1:thm:theorema-egregium}
	光滑曲面的 Gauss 曲率由第一基本形式决定。若两个光滑正则曲面片局部等距，则对应点的 Gauss 曲率相等；此结论不要求局部等距来自三维 Euclid 空间的刚性变换，也不要求保持所选法向。
\end{theorem}

\begin{proof}
	设局部等距在参数域中由微分同胚 \(\phi\) 表示。用 \(\phi\) 重参数化第二个曲面片后，它与第一个曲面片的 \(g_{ij}\) 完全相同，因而\eqref{g1:eq:christoffel-metric}给出相同的 \(\Gamma^k_{ij}\)，\eqref{g1:eq:surface-curvature-components}给出相同的 \(R^k{}_{j12}\)。再由\eqref{g1:eq:intrinsic-gauss-curvature}，两者的 \(K\) 相同。反转法向虽使 \(h,S\) 变号，却不改变 \(\det S\)，所以结论也不依赖该选择。
\end{proof}

这一定理称为\term[juemiaodingli]{绝妙定理}[theorema egregium]。内在度量确定两主曲率的乘积；平面与圆柱可以局部等距而主曲率不同，说明单个主曲率还记录着外在弯曲。另一方面，标准球面有 \(K=1/R^2\)，平面有 \(K=0\)，所以球面上任何开片都不能通过保持长度的局部等距完全铺到平面上。

\ProseParagraphBreak

下一节将从给定的度量与对称形式出发，尚无曲面可供求导。为此要明确：二维度量的曲率也只有一个独立的数，因而一个标量 Gauss 方程足以表达全部曲率条件。

\begin{proposition}[二维曲率只有一个独立分量]
	\label{g1:prop:surface-curvature-two-dim}
	设 \(g\) 是开集 \(U\subseteq\mathbb R^2\) 上的光滑正定度量，令 \((g^{ij})=(g_{ij})^{-1}\)，按\eqref{g1:eq:christoffel-metric}，即 \(\Gamma^k_{ij}=\tfrac12g^{k\ell}(\partial_i g_{j\ell}+\partial_jg_{i\ell}-\partial_\ell g_{ij})\)，定义 \(\nabla_i\partial_j=\Gamma^k_{ij}\partial_k\)，并令 \(R(V,W)Z=\nabla_V\nabla_WZ-\nabla_W\nabla_VZ-\nabla_{[V,W]}Z\)。令
	\[
		K_g\coloneqq\frac{g(R(\partial_1,\partial_2)\partial_2,\partial_1)}
		{\det(g_{ij})}.
	\]
	则 \(K_g\) 不依赖坐标，且任意切向量 \(V,W,Z\) 都满足
	\begin{equation}
		\label{g1:eq:two-dimensional-curvature}
		R(V,W)Z=K_g\bigl(g(W,Z)V-g(V,Z)W\bigr).
	\end{equation}
	若另给 \(U\) 上的一个光滑对称双线性形式 \(h\)，令 \(S=g^{-1}h\)，则标量等式 \(K_g=\det S\) 等价于完整的 Gauss 方程\eqref{g1:eq:gauss-equation}。
\end{proposition}

\begin{proof}
	固定一点，令 \(C(V,W,Z,T)=g(R(V,W)Z,T)\)。定义给出 \(C\) 在前两个变量中反对称。度量相容又给出
	\[
		g(R(V,W)Z,T)+g(Z,R(V,W)T)=0.
	\]
	要验证这式，只须对 \(g(Z,T)\) 应用恒等式 \(VW-WV-[V,W]=0\)，每次用度量相容展开；含一次协变导数的交叉内积抵消，剩余的正是这两项。因此 \(C\) 在最后两个变量中也反对称。

	取该点的任意正交单位基 \(e_1,e_2\)。二维空间的交替双线性形式只有一个独立分量：若 \(V=ae_1+be_2\)、\(W=ce_1+de_2\)，则前两个位置按双线性与反对称性展开后只剩系数 \(ad-bc\)。对最后两个位置作同样展开，整个 \(C\) 就由
	\[
		k=C(e_1,e_2,e_2,e_1)
	\]
	决定。具体地，
	\[
		C(V,W,Z,T)=k\bigl(g(W,Z)g(V,T)-g(V,Z)g(W,T)\bigr).
	\]
	取 \(V=\partial_1,W=Z=\partial_2,T=\partial_1\)，括号内是 \(g_{11}g_{22}-g_{12}^2\)，故 \(k=K_g\)，同时得到\eqref{g1:eq:two-dimensional-curvature}。右式也表明，在任意另一组基上用同一比值求出的仍是 \(k\)，所以与坐标无关。

	最后，\(S=g^{-1}h\) 关于 \(g\) 自伴。在一点选取它的正交单位特征基，设特征值为 \(a,b\)。此时直接在基向量上计算可得
	\[
		h(W,Z)SV-h(V,Z)SW
		=ab\bigl(g(W,Z)V-g(V,Z)W\bigr).
	\]
	两边对 \(V,W,Z\) 多线性，因此对任意向量成立。又 \(ab=\det S\)，与已证公式比较便得到等价性。
\end{proof}

这里的 \(K_g\) 可以在尚未证明度量能够浸入 \(\mathbb R^3\) 时定义。若确有这样的浸入，Gauss 方程才把它与外在定义的 \(K=\det S\) 认作同一个函数。这一区别会在给定数据的实现问题中起作用。

\begin{example}[共形坐标中的内在曲率]
	\label{g1:exm:conformal-curvature}
	设在参数域上给定度量
	\[
		g=e^{2\phi(u,v)}(\mathrm du^2+\mathrm dv^2),
	\]
	其中 \(\phi\) 光滑。两坐标方向垂直，且长度按同一倍数缩放。由\eqref{g1:eq:christoffel-metric}，
	\[
		\begin{array}{lll}
			\Gamma^1_{11}=\phi_u,&\Gamma^1_{12}=\Gamma^1_{21}=\phi_v,&\Gamma^1_{22}=-\phi_u,\\
			\Gamma^2_{11}=-\phi_v,&\Gamma^2_{12}=\Gamma^2_{21}=\phi_u,&\Gamma^2_{22}=\phi_v.
		\end{array}
	\]
	计算 \(R^1{}_{212}\) 时，两个乘积和恰好抵消：
	\[
		R^1{}_{212}
		=-\phi_{uu}-\phi_{vv}
		+\bigl(-\phi_u^2+\phi_v^2\bigr)
		-\bigl(\phi_v^2-\phi_u^2\bigr)
		=-\phi_{uu}-\phi_{vv}.
	\]
	由于 \(g_{12}=0\)、\(g_{11}=e^{2\phi}\)、\(D=e^{4\phi}\)，得到
	\begin{equation}
		\label{g1:eq:conformal-gauss-curvature}
		K_g=-e^{-2\phi}(\phi_{uu}+\phi_{vv}).
	\end{equation}
	例如 \(\phi=\log2-\log(1+u^2+v^2)\) 时，\(\phi_{uu}+\phi_{vv}=-4/(1+u^2+v^2)^2\)，故 \(K_g=1\)。求出这个数不需要先构造法向或主曲率。
\end{example}

\subsection{Codazzi 方程及混合偏导相容性}
\label{g1:subsec:0803:04}

Gauss 方程比较的是二阶求导的切向部分。法向部分也必须相等，而其中包含的是 \(h_{ij}\) 的一阶变化。直接比较 \(\partial_1h_{2j}\) 与 \(\partial_2h_{1j}\) 还不够，因为 \(h\) 的两个输入基向量也随位置改变，需要把这两种变化扣除。

\begin{definition}
	\label{g1:def:covariant-h-derivative}
	设 \((M,g)\) 为光滑度量曲面，\(\nabla\) 为其无挠相容协变导数，\(h\) 为 \(M\) 上的光滑对称双线性形式。对光滑切向量场 \(V,W,Z\)，定义
	\[
		(\nabla_Vh)(W,Z)
		\coloneqq V\bigl(h(W,Z)\bigr)-h(\nabla_VW,Z)-h(W,\nabla_VZ).
	\]
	在局部坐标中，令 \(h_{jk}=h(\partial_j,\partial_k)\)、\(\nabla_i\partial_j=\Gamma^m_{ij}\partial_m\)，其坐标分量记为
	\[
		\nabla_i h_{jk}\coloneqq\partial_i h_{jk}
		-\Gamma^m_{ij}h_{mk}-\Gamma^m_{ik}h_{jm}.
	\]
\end{definition}

这个定义中的两个减项分别对应 \(h\) 的两个输入。利用协变导数的乘法法则可以检查：将 \(W\) 或 \(Z\) 换成 \(fW\) 或 \(fZ\)，求导产生的 \(V(f)\) 项被相应减项抵消；对 \(V\) 则直接逐点线性。因此 \(\nabla h\) 的值同样不依赖输入向量场的延拓。

\begin{theorem}[Codazzi 方程]
	\label{g1:thm:codazzi-equation}
	设 \(U\subseteq\mathbb R^2\) 开，\(X:U\rightarrow\mathbb R^3\) 是光滑浸入，\(n\) 为光滑单位法向，\(h\) 为相应第二基本形式，\(S=-\mathrm{d}n\) 为形算子，\(\nabla\) 为诱导度量的无挠相容协变导数，则
	\[
		(\nabla_Vh)(W,Z)=(\nabla_Wh)(V,Z)
	\]
	对任意光滑切向量场 \(V,W,Z\) 成立。在局部坐标中，令 \(h_{ij}=h(\partial_i,\partial_j)\)、\(\nabla_i\partial_j=\Gamma^m_{ij}\partial_m\)，坐标表达为两个方程
	\begin{equation}
		\label{g1:eq:codazzi-components}
		\nabla_1h_{2j}-\nabla_2h_{1j}
		=\partial_1h_{2j}-\partial_2h_{1j}
		+\Gamma^m_{2j}h_{1m}-\Gamma^m_{1j}h_{2m}=0,
		\qquad j=1,2.
	\end{equation}
	若定义 \((\nabla_VS)W=\nabla_V(SW)-S(\nabla_VW)\)，则同一条件也可写成 \((\nabla_VS)W=(\nabla_WS)V\)。
\end{theorem}

\begin{proof}
	对 Gauss–Weingarten 方程求一次偏导。由于 \(\partial_i n\) 与 \(n\) 垂直，
	\[
		\langle\partial_1\partial_2X_j,n\rangle
		=\Gamma^m_{2j}h_{1m}+\partial_1h_{2j},\qquad
		\langle\partial_2\partial_1X_j,n\rangle
		=\Gamma^m_{1j}h_{2m}+\partial_2h_{1j}.
	\]
	混合偏导相等立即给出\eqref{g1:eq:codazzi-components}。它与协变导数的表达相同，是因为展开 \(\nabla_1h_{2j}-\nabla_2h_{1j}\) 后，项 \(-\Gamma^m_{12}h_{mj}\) 与 \(+\Gamma^m_{21}h_{mj}\) 相消。其余恰为上式。前面已经证明 \(\nabla h\) 对三个输入逐点线性，所以在坐标基上验证便给出一般向量场的结论。

	最后，度量相容以及 \(h(W,Z)=g(SW,Z)\) 给出
	\[
		(\nabla_Vh)(W,Z)
		=g\bigl(\nabla_V(SW)-S(\nabla_VW),Z\bigr).
	\]
	因度量非退化，对所有 \(Z\) 的内积相等当且仅当对应向量相等，这就得到形算子形式。
\end{proof}

这些关系称为\term[codazzifangcheng]{Codazzi 方程}[Codazzi equations]。\footnote{Codazzi，Delfino Codazzi（1824--1873），意大利数学家，研究曲面基本形式的相容关系与曲线坐标。}它们控制第二基本形式沿不同方向的变化，Gauss 方程则控制第二基本形式与度量曲率的对应。把两种分量放回同一次计算，可直接看见它们的共同来源：
\begin{equation}
	\label{g1:eq:mixed-surface-derivatives}
	\begin{aligned}
		0&=\partial_1\partial_2X_j-\partial_2\partial_1X_j\\
		&=\bigl(R^k{}_{j12}-h_{2j}h_1{}^k+h_{1j}h_2{}^k\bigr)X_k
		+\bigl(\nabla_1h_{2j}-\nabla_2h_{1j}\bigr)n.
	\end{aligned}
\end{equation}
这里切向与法向线性无关，所以两类系数必须分别消失。对于已经存在的曲面，这给出了必要条件；下一节将反过来证明，在足够小的参数邻域内，这些条件也足以从 \(g\) 与 \(h\) 恢复曲面。

% !TeX root = ../../Book1.tex
% 纲要标题：### 8.4 曲面基本定理的连续证明
% 纲要位置：几何学计划纲要.md，第 423 行。
% 正文状态：初稿；尚未正式审核。

% 纲要内容开始
% 
% * 相容数据的矩阵组织
% 
%   把第一、第二基本形式转写成带正交约束的标架系统，确认正定性及法向约定。
% 
% * Gauss–Codazzi 与平坦性
% 
%   计算结构方程，逐项识别给定条件恰好消除混合偏导障碍。
% 
% * 相容一阶系统的局部解
% 
%   建立本章所需的积分引理；不能以尚未证明的一般 Frobenius 定理替代。
% 
% * 从标架到闭的向量值一形式
% 
%   以第一基本形式组装待积分的微分，证明其外微分为零。
% 
% * 浸入的构造和秩检验
% 
%   积分得到位置向量，并用正定性证明微分秩为二。
% 
% * 两基本形式的恢复
% 
%   逐项核对所得浸入的度量和第二基本形式。
% 
% * 局部唯一性与整体单值
% 
%   比较两个标架初值；局部相差刚性变换，非单连通区域另有运输障碍。
% 
% 纲要内容结束

\section{曲面基本定理的连续证明}
\label{g1:sec:0804}

上一节由已经存在的曲面导出了相容条件。本节反过来：给定长度数据 \(g_{ij}\) 和弯曲数据 \(h_{ij}\)，若它们满足那些条件，是否确有一个曲面实现它们？第七章恢复曲线时只需沿一个参数积分标架；曲面有两个参数，沿不同方向走到同一点必须得到相同结果。因而这里真正增加的困难是两组微分方程能否同时成立。我们先把这个困难归结为一个矩阵恒等式，再用两次普通积分解决它。

\subsection{相容数据的矩阵组织}
\label{g1:subsec:0804:01}

设 \(U\subseteq\mathbb R^2\) 为开集，坐标记为 \(u^1=u,\ u^2=v\)，给定光滑函数 \(g_{ij}=g_{ji}\) 和 \(h_{ij}=h_{ji}\)，其中矩阵 \(g=(g_{ij})\) 处处正定。这里的“光滑”指 \(C^\infty\)。记
\[
	g=\begin{pmatrix}E&F\\F&G\end{pmatrix},\qquad D=EG-F^2>0,
	\qquad h_i{}^k=\sum_{\ell=1}^2g^{k\ell}h_{i\ell}.
\]
由 \(g\) 按\eqref{g1:eq:christoffel-metric}计算 \(\Gamma^k_{ij}\)。暂时还没有曲面，因而 \(\Gamma^k_{ij}\) 只是由给定函数计算出的函数。\cref{g1:prop:surface-connection-properties}的证明只用了 \(g\) 的正定性和这个公式，故其无挠条件 \(\Gamma^k_{ij}=\Gamma^k_{ji}\) 与度量相容等式在这里仍然成立。

\ProseParagraphBreak

假如曲面已经存在，它的三个列向量 \(\mathcal F=(X_u,X_v,n)\) 应满足 Gauss–Weingarten 方程。照\eqref{g1:eq:gauss-weingarten}把各列的系数排列起来，就得到待求系统
\begin{equation}
	\label{g1:eq:surface-coordinate-frame-system}
	\partial_i\mathcal F=\mathcal F A_i,\qquad
	A_i=
	\begin{pmatrix}
		\Gamma^1_{i1}&\Gamma^1_{i2}&-h_i{}^1\\
		\Gamma^2_{i1}&\Gamma^2_{i2}&-h_i{}^2\\
		h_{i1}&h_{i2}&0
	\end{pmatrix},\quad i=1,2.
\end{equation}
乘法仍在右边：例如第三列规定 \(\partial_i n=-h_i{}^1X_u-h_i{}^2X_v\)。与曲线的 Frenet 标架不同，\(X_u,X_v\) 通常既不单位也不垂直。因此这里要保持的不是 \(\mathcal F^{\mathsf T}\mathcal F=I\)，而是
\[
	\mathcal F^{\mathsf T}\mathcal F=\mathcal G,\qquad
	\mathcal G=\begin{pmatrix}E&F&0\\F&G&0\\0&0&1\end{pmatrix}.
\]
度量相容性和 \(h_{ij}=h_{ji}\) 给出
\begin{equation}
	\label{g1:eq:surface-matrix-metric-compatibility}
	\partial_i\mathcal G=A_i^{\mathsf T}\mathcal G+\mathcal G A_i.
\end{equation}
左上角的等式就是 \(\partial_i g_{jk}=\Gamma^\ell_{ij}g_{\ell k}+\Gamma^\ell_{ik}g_{j\ell}\)，右上角则是 \(h_{ij}-g_{j\ell}h_i{}^\ell=0\)，其余各块由转置或 \(0=0\) 得到。本节沿用上一节的求和约定：同一项中重复的一上一下指标从 \(1\) 到 \(2\) 求和。

\ProseParagraphBreak

为了直接使用正交标架的积分办法，先把给定长度拆出来。取
\begin{equation}
	\label{g1:eq:surface-metric-factor}
	P=
	\begin{pmatrix}
		\sqrt E&F/\sqrt E&0\\
		0&\sqrt D/\sqrt E&0\\
		0&0&1
	\end{pmatrix}.
\end{equation}
由于 \(E,D>0\)，它光滑、可逆，且 \(P^{\mathsf T}P=\mathcal G\)、\(\det P=\sqrt D>0\)。若令 \(\mathcal F=WP\)，所求三个向量就写成
\[
	a_1=\sqrt E\,e_1,\qquad
	a_2=\frac F{\sqrt E}e_1+\frac{\sqrt D}{\sqrt E}e_2,\qquad
	n=e_3,
\]
其中 \(W=(e_1,e_2,e_3)\) 是右手正交标架。这正是按给定长度和夹角拼出两个切向量；正平方根保证 \(a_1\times a_2=\sqrt D\,n\)，同时固定法向约定。

\ProseParagraphBreak

将 \(\mathcal F=WP\) 代入\eqref{g1:eq:surface-coordinate-frame-system}，可知只需解
\begin{equation}
	\label{g1:eq:surface-orthonormal-frame-system}
	\partial_iW=W\Omega_i,\qquad
	\Omega_i=(PA_i-\partial_iP)P^{-1}.
\end{equation}
这些系数确实斜对称。事实上，由\eqref{g1:eq:surface-matrix-metric-compatibility}，
\[
	P^{\mathsf T}(\Omega_i^{\mathsf T}+\Omega_i)P
	=A_i^{\mathsf T}\mathcal G+\mathcal G A_i-\partial_i(P^{\mathsf T}P)=0.
\]
因 \(P\) 可逆，故 \(\Omega_i^{\mathsf T}=-\Omega_i\)。至此，给定的非正交长度约束已变成一个正交标架系统。

\subsection{Gauss–Codazzi 与平坦性}
\label{g1:subsec:0804:02}

若可逆矩阵函数 \(\mathcal F\) 满足两组方程，交换混合偏导得到
\[
	0=\partial_1\partial_2\mathcal F-\partial_2\partial_1\mathcal F
	=\mathcal F\bigl(\partial_1A_2-\partial_2A_1+A_1A_2-A_2A_1\bigr).
\]
因此括号内的矩阵必须为零。对于一个矩阵系统，条件
\begin{equation}
	\label{g1:eq:surface-flatness}
	\mathcal C\coloneqq\partial_1A_2-\partial_2A_1+[A_1,A_2]=0,
	\qquad [A_1,A_2]=A_1A_2-A_2A_1
\end{equation}
称为该系统的\term[pingtanxing]{平坦性}[flatness]。这个名称指三个标架向量的整体运输方程，不是说给定曲面一定有 \(K=0\)；曲面的内在曲率和法向弯曲可以在 \(\mathcal C\) 中恰好抵消。

\begin{proposition}[相容条件与矩阵平坦性]
	\label{g1:prop:surface-data-flatness}
	设 \(g\) 是平面开集 \(U\) 上的光滑正定对称矩阵，\(h\) 是光滑对称矩阵，各矩阵均为 \(2\times2\)。令 \(D=\det g\)、\((g^{ij})=g^{-1}\)、\(h_i{}^k=g^{k\ell}h_{i\ell}\)，\(\Gamma\) 由 \(g\) 的 Christoffel 公式确定。对 \(i=1,2\)，令 \(A_i=\begin{psmallmatrix}\Gamma^1_{i1}&\Gamma^1_{i2}&-h_i{}^1\\\Gamma^2_{i1}&\Gamma^2_{i2}&-h_i{}^2\\h_{i1}&h_{i2}&0\end{psmallmatrix}\)，即\eqref{g1:eq:surface-coordinate-frame-system}中的系数矩阵，并令 \(\mathcal C=\partial_1A_2-\partial_2A_1+A_1A_2-A_2A_1\)。则 \(\mathcal C=0\) 等价于 Gauss、Codazzi 两组分量等式
	\begin{align}
		\label{g1:eq:surface-data-gauss}
		R^k{}_{j12}&=h_1{}^kh_{2j}-h_2{}^kh_{1j},&&j,k=1,2,\\
		\label{g1:eq:surface-data-codazzi}
		\nabla_1h_{2j}&=\nabla_2h_{1j},&&j=1,2.
	\end{align}
	其中 \(R\) 是 \(\nabla_i\partial_j=\Gamma^k_{ij}\partial_k\) 的曲率，约定 \(R(V,W)Z=\nabla_V\nabla_WZ-\nabla_W\nabla_VZ-\nabla_{[V,W]}Z\)，并以 \(R(\partial_1,\partial_2)\partial_j=R^k{}_{j12}\partial_k\) 定义分量，按\eqref{g1:eq:surface-curvature-components}计算；\(\nabla_i h_{jk}=\partial_i h_{jk}-\Gamma^\ell_{ij}h_{\ell k}-\Gamma^\ell_{ik}h_{j\ell}\)。等价地，第一组条件可以换成 \(K_g=(h_{11}h_{22}-h_{12}^2)/D\)，其中 \(K_g=\langle R(\partial_u,\partial_v)\partial_v,\partial_u\rangle_g/D\) 完全由 \(g\) 决定。
\end{proposition}

\begin{proof}
	按前两行与最后一行分块。矩阵 \(\mathcal C\) 左上角的第 \(k,j\) 项为
	\[
		\partial_1\Gamma^k_{2j}-\partial_2\Gamma^k_{1j}
		+\Gamma^k_{1\ell}\Gamma^\ell_{2j}-\Gamma^k_{2\ell}\Gamma^\ell_{1j}
		-h_1{}^kh_{2j}+h_2{}^kh_{1j}.
	\]
	前四项正是 \(R^k{}_{j12}\)，故这一块消失恰好等价于 Gauss 方程。左下角的第 \(j\) 项为
	\[
		\partial_1h_{2j}-\partial_2h_{1j}
		+h_{1\ell}\Gamma^\ell_{2j}-h_{2\ell}\Gamma^\ell_{1j}
		=\nabla_1h_{2j}-\nabla_2h_{1j}.
	\]
	最后一步中另外两项 \(-\Gamma^\ell_{12}h_{\ell j}\) 与 \(+\Gamma^\ell_{21}h_{\ell j}\) 因无挠条件抵消。右下角是 \(-h_{1\ell}h_2{}^\ell+h_{2\ell}h_1{}^\ell=0\)，因为 \(g^{-1}\) 对称。

	还需检查右上角，不能仅凭矩阵对称性猜测它。交换\eqref{g1:eq:surface-matrix-metric-compatibility}的混合偏导并代回同一等式，展开后得到
	\[
		\mathcal C^{\mathsf T}\mathcal G+\mathcal G\mathcal C=0.
	\]
	当左上角和左下角已为零时，写 \(\mathcal C=\left(\begin{smallmatrix}0&b\\0&0\end{smallmatrix}\right)\)，上式的右上角便是 \(gb=0\)。正定矩阵 \(g\) 可逆，所以 \(b=0\)。这证明两组几何条件足以使整个 \(\mathcal C\) 消失；反向则直接读取上述两块。

	最后，\cref{g1:prop:surface-curvature-two-dim}已对任意光滑正定 \(g\) 证明，二维曲率张量由 \(\langle R(\partial_u,\partial_v)\partial_v,\partial_u\rangle_g\) 这一分量确定。Gauss 方程右边也有相同的交替性和度量斜对称性；其这一分量为 \(h_{11}h_{22}-h_{12}^2\)。故完整的 Gauss 分量条件与所述标量条件等价。
\end{proof}

换成正交标架不会增加新的相容条件。用 \(\partial_iP^{-1}=-P^{-1}(\partial_iP)P^{-1}\) 展开\eqref{g1:eq:surface-orthonormal-frame-system}，可以把二阶导数相互抵消，得到
\begin{equation}
	\label{g1:eq:surface-flatness-change-frame}
	\partial_1\Omega_2-\partial_2\Omega_1+[\Omega_1,\Omega_2]
	=P\mathcal CP^{-1}.
\end{equation}
这个换标架公式也可不逐项展开地核对：对任意矩阵函数 \(W\)，先写
\[
	\partial_i(WP)-(WP)A_i=(\partial_iW-W\Omega_i)P,
\]
再先后应用 \(i=1,2\) 两个运算并相减；左边留下 \(-(WP)\mathcal C\)，右边留下 \(-W(\partial_1\Omega_2-\partial_2\Omega_1+[\Omega_1,\Omega_2])P\)，即得上式。因此 Gauss–Codazzi 条件也恰好使正交标架系统平坦。

\subsection{相容一阶系统的局部解}
\label{g1:subsec:0804:03}

接下来解决分析上的问题：两个方向各有一条普通微分方程，如何知道依次积分的结果会同时满足它们？可以先沿一条水平线解，再沿每条竖线解。竖向方程由构造满足；平坦性则保证先前满足的横向方程不会在竖向移动时丢失。

\begin{lemma}[矩形上的相容矩阵系统]
	\label{g1:lem:rectangular-matrix-system}
	设 \(V=I\times J\subseteq\mathbb R^2\) 为开矩形，\((u_0,v_0)\in V\)，\(\Omega_1,\Omega_2\colon V\rightarrow M_3(\mathbb R)\) 光滑，且
	\[
		\partial_u\Omega_2-\partial_v\Omega_1+[\Omega_1,\Omega_2]=0.
	\]
	对任意矩阵 \(W_0\in M_3(\mathbb R)\)，存在唯一光滑函数 \(W\colon V\rightarrow M_3(\mathbb R)\)，使
	\[
		W_u=W\Omega_1,\qquad W_v=W\Omega_2,\qquad W(u_0,v_0)=W_0.
	\]
	若两系数斜对称且 \(W_0\in \operatorname{SO}(3)\)，则 \(W(u,v)\in \operatorname{SO}(3)\)。
\end{lemma}

\begin{proof}
	由\cref{g1:lem:linear-frame-ode}，先在整个 \(I\) 上解
	\[
		H'(u)=H(u)\Omega_1(u,v_0),\qquad H(u_0)=W_0.
	\]
	系数光滑，对方程反复求导可知 \(H\) 光滑。随后对每个固定的 \(u\)，在整个 \(J\) 上解
	\[
		W_v(u,v)=W(u,v)\Omega_2(u,v),\qquad W(u,v_0)=H(u).
	\]
	这立即给出 \(v\) 方向的方程，但还需证明所得 \(W\) 对 \(u\) 光滑，才能检查另一个方向。

	这一点可从前引理的迭代积分证明直接补出。在任意紧子矩形上，取包含 \(v_0\) 的紧 \(v\) 区间，写
	\[
		W=\sum_{n=0}^\infty Y_n,\qquad Y_0(u,v)=H(u),\qquad
		Y_{n+1}(u,v)=\int_{v_0}^vY_n(u,t)\Omega_2(u,t)\,\mathrm{d}t.
	\]
	第 \(n\) 项是含 \(n\) 个系数因子的有序积分。给定非负整数 \(r\)，设 \(M_r\geq1\) 控制该紧矩形上 \(H,\Omega_2\) 直到 \(r\) 阶的所有 \(u\) 导数，\(T\) 控制 \(|v-v_0|\)。对这 \(n+1\) 个因子使用乘积求导法则，其系数之和为 \((n+1)^r\)，所以
	\[
		\|\partial_u^r Y_n(u,v)\|
		\leq (n+1)^r M_r^{\,n+1}\frac{T^n}{n!}.
	\]
	右边对 \(n\) 可求和，故各阶 \(u\) 导数级数一致收敛，可以逐项求导。再用 \(W_v=W\Omega_2\) 对 \(v\) 及混合导数逐阶求导，即知 \(W\) 光滑。向 \(v_0\) 左侧积分时，只需将有序积分反向后取绝对值，估计不变。紧子矩形任意，因而光滑性在整个 \(V\) 上成立。

	定义尚待消除的误差 \(Z=W_u-W\Omega_1\)。利用已经成立的竖向方程求导：
	\begin{align*}
		Z_v
		&=(W\Omega_2)_u-W\Omega_2\Omega_1-W\partial_v\Omega_1\\
		&=Z\Omega_2+W\bigl(\partial_u\Omega_2-\partial_v\Omega_1+[\Omega_1,\Omega_2]\bigr)
		=Z\Omega_2.
	\end{align*}
	在 \(v=v_0\) 上，\(Z=H'-H\Omega_1=0\)。固定 \(u\) 后，\cref{g1:lem:linear-frame-ode}的唯一性使这个零初值方程只有零解，故 \(Z=0\)，横向方程也成立。

	任意另一个满足双系统的解，限制到 \(v=v_0\) 后必须等于 \(H\)，再沿每条竖线必须等于 \(W\)，故解唯一。若系数斜对称，\cref{g1:prop:skew-frame-orthogonality}先沿水平线、再沿竖线保持 \(\operatorname{SO}(3)\)，得到最后的结论。
\end{proof}

\cref{g1:fig:surface-rectangle-integration}画出了构造中选定的折线路径。若改成先竖后横，同样的证明仍给出满足双系统、具有相同初值的解，故由唯一性得到同一个标架。这里的关键不是随意约定一条路径，而是证明路径上的解满足全部两个方向的方程。

\begin{figure}[htbp]
	\centering
	\input{Figures/Chapter08-RectangleIntegration}
	\caption[矩形上的标架积分]{先沿水平线积分，再沿竖线积分。平坦性使横向误差 \(Z\) 始终为零，故交换这两个积分方向也得到同一终点标架。}
	\label{g1:fig:surface-rectangle-integration}
\end{figure}

\subsection{从标架到闭的向量值一形式}
\label{g1:subsec:0804:04}

现在假定给定数据满足 Gauss–Codazzi 条件，在任意包含基点 \(q_0=(u_0,v_0)\) 的矩形 \(V\subseteq U\) 上，用\cref{g1:lem:rectangular-matrix-system}解出初值 \(W_0\in \operatorname{SO}(3)\) 的正交标架 \(W\)，再令
\[
	\mathcal F=WP=(a_1,a_2,n).
\]
由\eqref{g1:eq:surface-orthonormal-frame-system}的定义，
\[
	\partial_i\mathcal F
	=W\Omega_iP+W\partial_iP
	=WPA_i=\mathcal F A_i.
\]
所以 \(a_1,a_2\) 已有了候选切向量应有的导数关系。还欠位置函数 \(X\)：必须证明这两个向量确实能同时写成 \(X_u,X_v\)。

\ProseParagraphBreak

把候选微分写作
\[
	\alpha=a_1\,\mathrm{d}u+a_2\,\mathrm{d}v.
\]
这是一个\term[xiangliangzhiyixingshi]{向量值一形式}[vector-valued 1-form]；这里，它将参数平面中的位移 \((\xi^1,\xi^2)\) 送到向量 \(a_1\xi^1+a_2\xi^2\)。若它的每个实分量都满足混合偏导条件，就称它为闭的；该条件写作
\[
	\mathrm{d}\alpha=(\partial_u a_2-\partial_v a_1)\,\mathrm{d}u\wedge \mathrm{d}v=0.
\]
这里 \(\mathrm{d}u\wedge \mathrm{d}v\) 只记录两个方向交换时变号；不需要预先掌握微分形式的一般理论。由 \(\partial_i\mathcal F=\mathcal F A_i\) 的前两列，
\[
	\partial_u a_2-\partial_v a_1
	=(\Gamma^1_{12}-\Gamma^1_{21})a_1
	+(\Gamma^2_{12}-\Gamma^2_{21})a_2+(h_{12}-h_{21})n=0.
\]
这一步说明了两种对称性的不同职责：Gauss–Codazzi 条件先使标架能够积分，而 \(\Gamma^k_{12}=\Gamma^k_{21}\)、\(h_{12}=h_{21}\) 再使位置的候选微分闭合。两步缺一不可。

\subsection{浸入的构造和秩检验}
\label{g1:subsec:0804:05}

在矩形上，把闭的向量值一形式积分成位置无需调用另一个一般存在定理。指定 \(x_0\in\mathbb R^3\)，直接沿刚才的折线定义
\begin{equation}
	\label{g1:eq:surface-position-rectangle-integral}
	X(u,v)=x_0+\int_{u_0}^u a_1(s,v_0)\,\mathrm{d}s+\int_{v_0}^v a_2(u,t)\,\mathrm{d}t.
\end{equation}
对 \(v\) 求导得到 \(X_v=a_2\)。对 \(u\) 求导，并用 \(\partial_u a_2=\partial_v a_1\)，得到
\[
	X_u=a_1(u,v_0)+\int_{v_0}^v\partial_u a_2(u,t)\,\mathrm{d}t
	=a_1(u,v_0)+a_1(u,v)-a_1(u,v_0)=a_1(u,v).
\]
于是 \(\mathrm{d}X=\alpha\)，且 \(X(q_0)=x_0\)。积分中各函数光滑，故 \(X\) 光滑。

\ProseParagraphBreak

有了位置函数还不能直接宣布它是正则曲面，必须检查两个切向量处处独立。由 \(W^{\mathsf T}W=I\) 及 \(P^{\mathsf T}P=\mathcal G\)，对任意 \(\xi=(\xi^1,\xi^2)\in\mathbb R^2\)，
\[
	\|\mathrm{d}X(\xi)\|^2
	=\|a_1\xi^1+a_2\xi^2\|^2
	=g_{ij}\xi^i\xi^j.
\]
当 \(\xi\ne0\) 时，正定性使右端严格为正，因此 \(\mathrm{d}X\) 的核为零、秩为 \(2\)。这正是最初坚持 \(g\) 正定的原因；若它退化，即使所有偏导相容，也不能保证得到浸入。

\subsection{两基本形式的恢复}
\label{g1:subsec:0804:06}

构造出了浸入，还须核对它实现的是原先的数据。第一基本形式由 \(\mathcal F^{\mathsf T}\mathcal F=\mathcal G\) 立即得到：
\[
	\langle X_i,X_j\rangle=\langle a_i,a_j\rangle=g_{ij}.
\]
从\eqref{g1:eq:surface-metric-factor}的列向量关系得到
\[
	X_u\times X_v=a_1\times a_2=\sqrt D\,n,
\]
故标架第三列确实等于按参数顺序选择的单位法向 \((X_u\times X_v)/\|X_u\times X_v\|\)。取 \(\partial_i\mathcal F=\mathcal F A_i\) 的第 \(j\) 列，
\[
	X_{ij}=\Gamma^1_{ij}X_u+\Gamma^2_{ij}X_v+h_{ij}n.
\]
与 \(n\) 作内积即得 \(\langle X_{ij},n\rangle=h_{ij}\)，所以第二基本形式也恢复了。特别地，第二基本形式的符号并非由事后任取法向决定；构造中的 \(P\) 和 \(W\in \operatorname{SO}(3)\) 已使它与给定参数方向一致。

\begin{example}[平直度量配上非零弯曲]
	\label{g1:exm:fundamental-cylinder-data}
	取 \(g=\mathrm{d}u^2+\mathrm{d}v^2\)、\(h=\lambda\,\mathrm{d}u^2\)，其中 \(\lambda\ne0\) 为常数。此时 \(\Gamma=0\)、\(K_g=0\)、\(\det h=0\)，Codazzi 条件也因系数常数而成立。取 \(P=I\)、\(W(0,0)=I\)，系统成为
	\[
		A_1=\begin{pmatrix}0&0&-\lambda\\0&0&0\\\lambda&0&0\end{pmatrix},
		\qquad A_2=0.
	\]
	于是
	\[
		a_1=(\cos\lambda u,0,\sin\lambda u),\quad
		a_2=(0,1,0),\quad n=(-\sin\lambda u,0,\cos\lambda u).
	\]
	以 \(X(0,0)=0\) 积分，得到
	\[
		X(u,v)=\left(\frac{\sin\lambda u}{\lambda},\,v,\,
		\frac{1-\cos\lambda u}{\lambda}\right).
	\]
	这是半径 \(1/|\lambda|\) 的圆柱面参数化。平直的长度数据允许一个切向方向发生弯曲；Gauss 方程只要求两个主曲率的乘积为零，并不要求每个主曲率都为零。
\end{example}

\subsection{局部唯一性与整体单值}
\label{g1:subsec:0804:07}

现在把构造和唯一性合并成完整结论。初始标架所规定的是两个切向量按顺序正交归一化后的方向，而不是任意指定 \(X_u,X_v\) 的长度；后者已由 \(g\) 给定。

\begin{theorem}[曲面基本定理]
	\label{g1:thm:surface-fundamental}
	设 \(U\subseteq\mathbb R^2\) 为开集，\(g=(g_{ij})\) 为光滑正定对称矩阵，\(h=(h_{ij})\) 为光滑对称矩阵，各矩阵均为 \(2\times2\)。令 \((g^{ij})=g^{-1}\)，\(\Gamma\) 为 \(g\) 的 Christoffel 系数，\(\nabla_i\partial_j=\Gamma^k_{ij}\partial_k\)，\(R\) 为该联络的曲率，约定 \(R(V,W)Z=\nabla_V\nabla_WZ-\nabla_W\nabla_VZ-\nabla_{[V,W]}Z\)。以 \(R(\partial_1,\partial_2)\partial_j=R^k{}_{j12}\partial_k\) 定义分量，令 \(h_i{}^k=g^{k\ell}h_{i\ell}\)、\(\nabla_i h_{jk}=\partial_i h_{jk}-\Gamma^\ell_{ij}h_{\ell k}-\Gamma^\ell_{ik}h_{j\ell}\)，要求 \(R^k{}_{j12}=h_1{}^kh_{2j}-h_2{}^kh_{1j}\) 及 \(\nabla_1h_{2j}=\nabla_2h_{1j}\) 对 \(j,k=1,2\) 成立，即\eqref{g1:eq:surface-data-gauss}、\eqref{g1:eq:surface-data-codazzi}。令 \(E=g_{11}\)、\(F=g_{12}\)、\(D=\det g\)，并令 \(P=\begin{psmallmatrix}\sqrt E&F/\sqrt E&0\\0&\sqrt D/\sqrt E&0\\0&0&1\end{psmallmatrix}\)。给定 \(q_0=(u_0,v_0)\in U\)、\(x_0\in\mathbb R^3\) 及右手正交标架矩阵 \(W_0\in\operatorname{SO}(3)\)，则在每个含 \(q_0\) 的开矩形 \(V\subseteq U\) 上，存在唯一的光滑浸入 \(X\colon V\rightarrow\mathbb R^3\)，其第一、第二基本形式分别为 \(g,h\)，且
	\[
		X(q_0)=x_0,\qquad
		(X_u,X_v,n)(q_0)=W_0P(q_0),\qquad
		n=\frac{X_u\times X_v}{\|X_u\times X_v\|}.
	\]
	这里 \(P\) 由\eqref{g1:eq:surface-metric-factor}确定。若不指定这些初值，同一矩形上的任意两个这样的浸入满足
	\[
		\widetilde X=QX+b,\qquad Q\in \operatorname{SO}(3),\quad b\in\mathbb R^3.
	\]
	这一结论称为\term[qumianjibendingli]{曲面基本定理}[fundamental theorem of surface theory]。
\end{theorem}

\begin{proof}
	由\cref{g1:prop:surface-data-flatness}及\eqref{g1:eq:surface-flatness-change-frame}，\(\Omega_1,\Omega_2\) 满足平坦性；它们的斜对称性已由\eqref{g1:eq:surface-matrix-metric-compatibility}证明。\cref{g1:lem:rectangular-matrix-system}给出初值为 \(W_0\) 的 \(\operatorname{SO}(3)\) 值解。于是前面三小节的构造依次给出闭的候选微分、\eqref{g1:eq:surface-position-rectangle-integral}中的光滑浸入，以及其基本形式和法向的逐项核验，这证明存在性。

	对于任意另一个具有相同数据的浸入 \(\widetilde X\)，\cref{g1:prop:gauss-weingarten}使其坐标标架 \(\widetilde{\mathcal F}=(\widetilde X_u,\widetilde X_v,\widetilde n)\) 满足相同的 \(\partial_i\widetilde{\mathcal F}=\widetilde{\mathcal F}A_i\)。因 \(g\) 相同，\(\widetilde W=\widetilde{\mathcal F}P^{-1}\) 满足同一正交标架系统。若初值也相同，矩形积分引理的唯一性给出 \(\widetilde W=W\)，从而 \(\mathrm{d}\widetilde X=\mathrm{d}X\)。两者之差的两个偏导都为零，沿水平、竖直线段积分可知其在矩形上为常数；基点位置相同便使这个常数为零。

	若不指定初值，令 \(Q=\widetilde W(q_0)W(q_0)^{\mathsf T}\in \operatorname{SO}(3)\)。矩阵 \(QW\) 满足与 \(\widetilde W\) 相同的系统和初值，因此 \(\widetilde W=QW\)。于是 \(\mathrm{d}\widetilde X=Q\,\mathrm{d}X\)，积分得到 \(\widetilde X=QX+b\)，其中 \(b=\widetilde X(q_0)-QX(q_0)\)。
\end{proof}

这里的“唯一”保留了参数点之间的对应；若允许重新参数化，就必须同时按该坐标变换拉回 \(g,h\)。此外，保向刚性变换的限定与法向约定有关：若用一个行列式为 \(-1\) 的正交矩阵反射曲面，按参数叉积重新选择的法向会多出一个负号，第二基本形式随之变号。

\ProseParagraphBreak

对一般开集，定理保证每点附近的局部实现。若沿一条道路继续这些局部解，相交处可用同一位置和标架初值保证它们一致；每一段有限道路只需有限多个小矩形。但绕环路回来时，标架可能已经转动，或者标架虽回到原值，位置却发生平移。这与\cref{g1:prop:periodic-curve-closure}中“标架回到初值”和“位移积分为零”两个条件的区别相同。以下两个例子直接显示这两种障碍，而不把局部存在性误当成任意区域上的单值存在性。

\begin{example}[转动障碍与位移障碍]
	\label{g1:exm:surface-monodromy-obstructions}
	在环带 \(A=\{(x,y):1<\sqrt{x^2+y^2}<2\}\) 上，写 \(r=\sqrt{x^2+y^2}\)。角度 \(\theta\) 不是整个环带上的单值实函数，但
	\[
		\mathrm{d}r=\frac{x\,\mathrm{d}x+y\,\mathrm{d}y}{r},\qquad
		\mathrm{d}\theta=\frac{-y\,\mathrm{d}x+x\,\mathrm{d}y}{r^2}
	\]
	都是单值的一形式；局部选定角度后，它们就是通常的坐标微分。

	先取 \(g=\mathrm{d}r^2+c^2r^2\mathrm{d}\theta^2\)、\(h=0\)，其中 \(c>0\)。在展开角度的矩形 \((1,2)\times\mathbb R\) 上，
	\[
		\widetilde X(r,\theta)=(r\cos c\theta,r\sin c\theta,0)
	\]
	实现这两组数据，因而各局部坐标中的 Gauss–Codazzi 条件都成立。它的单位径向、角向标架绕一周后旋转 \(2\pi c\)。若 \(c\notin\mathbb Z\)，这个标架不能在 \(\theta\) 与 \(\theta+2\pi\) 处匹配。假如环带上有实现数据的单值浸入，把它写到展开角度的矩形上，再用曲面基本定理的唯一性，它必与 \(\widetilde X\) 相差固定保向刚性变换；但固定刚性变换不能消除标架的不匹配，故这样的浸入不存在。

	再取 \(g=\mathrm{d}r^2+\mathrm{d}\theta^2\)、\(h=0\)。在展开角度的矩形上，\(\widetilde X(r,\theta)=(r,\theta,0)\) 实现数据，沿径向、角向的标架都是常向量，已经没有转动障碍。然而候选位置微分为 \(\alpha=(\mathrm{d}r,\mathrm{d}\theta,0)\)，沿任一正向圆周积分得到
	\[
		\oint\alpha=(0,2\pi,0)\ne0.
	\]
	一个单值位置函数的微分沿环路积分必须等于终点减起点，即零；这里却存在非零位移。也可直接由 \(\widetilde X(r,\theta+2\pi)-\widetilde X(r,\theta)=(0,2\pi,0)\) 和同样的唯一性论证排除整体浸入。
\end{example}

这两个例子中的数据处处光滑并且满足相容方程，障碍只在绕环带一周后出现。因而 Gauss–Codazzi 方程精确回答了局部能否实现的问题；在带孔区域上还要另查标架运输和环路位移，不能仅检查点上的曲率等式。

% !TeX root = ../../Book1.tex
% 纲要标题：### 8.5 模型、形变与刚性的界线
% 纲要位置：几何学计划纲要.md，第 453 行。
% 正文状态：初稿；尚未正式审核。

% 纲要内容开始
% 
% * 球面、圆柱和旋转曲面的核验
% * 平面与圆柱的局部等距
% * 一阶、局部和整体刚性问题的区分
% 
% ---
% 
% 纲要内容结束

\section{模型、形变与刚性的界线}
\label{g1:sec:0805}

曲面基本定理需要同时给定 \(I\) 和 \(II\)。如果只给 \(I\)，内在曲率虽然已经确定，两个主曲率及曲面在空间中的形状却还可能变化。本节把这种自由算在几个熟悉模型上，再区分三个常被混在一起的问题：一个变动在一阶上是否保持度量，是否存在真实的等距形变，以及局部构造能否在整张曲面上成立。

\subsection{球面、圆柱和旋转曲面的核验}
\label{g1:subsec:0805:01}

\begin{example}[球面与法向的符号]\label{g1:exm:sphere-two-forms}
	半径 \(R>0\) 的球面去掉两极后，可在经度的局部区间中用
	\[
		X(s,\theta)=\bigl(R\sin(s/R)\cos\theta,\ R\sin(s/R)\sin\theta,\ R\cos(s/R)\bigr),
		\quad 0<s<\pi R
	\]
	参数化。直接求导得到
	\[
		I=\mathrm{d}s^2+R^2\sin^2(s/R)\,\mathrm{d}\theta^2,\qquad
		n=\frac{X}{R}=\frac{X_s\times X_\theta}{\|X_s\times X_\theta\|}.
	\]
	因为 \(\mathrm{d}n=\mathrm{d}X/R\)，本书约定 \(S=-\mathrm{d}n\) 给出 \(S=-R^{-1}\operatorname{id}\)，所以
	\[
		II=-R^{-1}I,\qquad k_1=k_2=-R^{-1},\qquad H=-R^{-1},\qquad K=R^{-2}.
	\]
	因此标准球面仍有正的 Gauss 曲率；外向法向下主曲率为负，与 \(S=-\mathrm{d}n\) 一致。反向选择法向时 \(II,H,k_1,k_2\) 改号，\(I,K\) 不变。参数在两极退化，因为 \(X_\theta=0\)，但球面本身在那里可以改用图像坐标，没有几何奇点。
\end{example}

\begin{example}[圆柱的两个方向]\label{g1:exm:cylinder-two-forms}
	取
	\[
		C(u,v)=\bigl(R\cos(u/R),R\sin(u/R),v\bigr),\qquad R>0.
	\]
	有 \(C_u=(-\sin(u/R),\cos(u/R),0)\)、\(C_v=(0,0,1)\)，两者单位且正交。外向法向 \(n=(\cos(u/R),\sin(u/R),0)\) 等于 \(C_u\times C_v\)，而
	\[
		n_u=R^{-1}C_u,\qquad n_v=0.
	\]
	故
	\[
		I=\mathrm{d}u^2+\mathrm{d}v^2,\qquad II=-R^{-1}\mathrm{d}u^2,\qquad
		k_1=-R^{-1},\quad k_2=0,\quad H=-\frac1{2R},\quad K=0.
	\]
	圆周方向发生弯曲，母线方向没有法向弯曲。这里有一个零主曲率并不表示整张曲面是平面，也不表示参数不正则。
\end{example}

这两例都可由旋转曲面的统一公式找回。取光滑单位速率平面曲线 \((r(s),z(s))\)，假设 \(r>0\) 且 \(r'^2+z'^2=1\)，绕 \(z\) 轴旋转：
\[
	X(s,\theta)=(r(s)\cos\theta,r(s)\sin\theta,z(s)).
\]
弧长条件只作用于母线；\(\theta\) 方向的长度因子仍为 \(r\)。

\begin{proposition}[旋转曲面的基本形式]\label{g1:prop:revolution-two-forms}
	在上述光滑单位速率母线、\(r>0\) 及局部角参数条件下，选择法向
	\[
		n=(-z'\cos\theta,-z'\sin\theta,r'),
	\]
	则
	\[
		I=\mathrm{d}s^2+r^2\mathrm{d}\theta^2,\qquad
		II=(r'z''-z'r'')\,\mathrm{d}s^2+rz'\,\mathrm{d}\theta^2.
	\]
	子午线和纬线方向为主方向，主曲率及 Gauss 曲率为
	\begin{equation}\label{g1:eq:revolution-curvatures}
		k_s=r'z''-z'r'',\qquad k_\theta=\frac{z'}r,\qquad
		K=-\frac{r''}r,\qquad H=\frac12\left(r'z''-z'r''+\frac{z'}r\right).
	\end{equation}
	在脐点，这两方向仍是主方向，但不再是仅有的主方向。
\end{proposition}

\begin{proof}
	计算 \(X_s,X_\theta\) 得到 \(\|X_s\|=1\)、\(\|X_\theta\|=r\)、\(\langle X_s,X_\theta\rangle=0\)，且 \(X_s\times X_\theta=rn\)。二阶导数与 \(n\) 的内积分别是
	\[
		\langle X_{ss},n\rangle=-z'r''+r'z'',\qquad
		\langle X_{s\theta},n\rangle=0,\qquad
		\langle X_{\theta\theta},n\rangle=rz'.
	\]
	因此 \(I,II\) 都对角，形算子的两个特征值是相应对角项之比。最后，由单位速率条件求导得 \(r'r''+z'z''=0\)，从而
	\[
		z'(r'z''-z'r'')=-r'^2r''-z'^2r''=-r'',
	\]
	即得 \(K=-r''/r\)。此推导没有除以 \(z'\)，故 \(z'=0\) 的点也包括在内。
\end{proof}

令 \(r=R\sin(s/R),z=R\cos(s/R)\)，就得到外向球面的两个负主曲率。令 \(r=R,z=s\)，则上面指定的法向朝向圆柱轴线，因而 \(k_\theta=1/R\)；它与前一例的负号差别完全来自法向，不能直接混用。对圆锥的去顶点部分，取 \(r=as,z=bs\)，其中 \(s>0,a>0,a^2+b^2=1\)，可得 \(K=0\)。\(r=0\) 的顶点不属于命题的适用区域，零曲率公式不能被延伸为顶点正则性的证明。

\begin{example}[一个非平面的负曲率模型]\label{g1:exm:catenoid-curvatures}
	取 \(r(s)=\sqrt{a^2+s^2}\)、\(z(s)=a\,\operatorname{arsinh}(s/a)\)，其中 \(a>0\)。这里 \(\operatorname{arsinh}\) 是双曲正弦的反函数，导数为 \(1/\sqrt{1+t^2}\)。计算得
	\[
		r'=\frac{s}{r},\quad z'=\frac ar,\quad
		r''=\frac{a^2}{r^3},\quad z''=-\frac{as}{r^3}.
	\]
	因此母线确实单位速率，而旋转曲面的 \(k_s=-a/r^2\)、\(k_\theta=a/r^2\)，故 \(H=0\)、\(K=-a^2/r^4<0\)。这是\term[xuanlianmian]{悬链面}[catenoid]的一个弧长参数表示。平均曲率为零表示两个带符号主曲率相互抵消，绝不是两个主曲率都为零。
\end{example}

\subsection{平面与圆柱的局部等距}
\label{g1:subsec:0805:02}

平面 \(P(u,v)=(u,v,0)\) 与圆柱参数 \(C(u,v)\) 的第一基本形式都为 \(\mathrm{d}u^2+\mathrm{d}v^2\)。在宽度小于 \(2\pi R\) 的开条带上，\(C\) 单射，故对应 \(P(u,v)\mapsto C(u,v)\) 是局部曲面之间的等距。由\cref{g1:prop:local-isometry-length}，任意纸上曲线的长度以及相交角度都保持；这正是“无拉伸卷起”的数学含义。

\ProseParagraphBreak

它不可能由空间刚性运动得到：平面的第二基本形式恒为零，而圆柱的第二基本形式不为零；刚性运动保持第二基本形式（同时推送所选法向），反射至多因法向约定带来整体改号，也不能把零变成非零。第一基本形式的相同因此不足以触发上一节的曲面唯一性结论。

\begin{example}[从平面连续卷起]\label{g1:exm:plane-cylinder-bending}
	在固定有界矩形 \(U=(-a,a)\times(-b,b)\) 上定义
	\[
		X_t(u,v)=
		\begin{cases}
			\bigl(\sin(tu)/t,\ v,\ (1-\cos(tu))/t\bigr),&t\ne0,\\
			(u,v,0),&t=0.
		\end{cases}
	\]
	函数 \(\sin(tu)/t=\int_0^u\cos(t\xi)\,\mathrm{d}\xi\)、\((1-\cos(tu))/t=\int_0^u\sin(t\xi)\,\mathrm{d}\xi\) 表明 \(X_t\) 对 \((t,u,v)\) 光滑，并无真实的 \(t=0\) 奇点。因为
	\[
		(X_t)_u=(\cos(tu),0,\sin(tu)),\qquad (X_t)_v=(0,1,0),
	\]
	全部 \(t\) 对应的第一基本形式都为 \(\mathrm{d}u^2+\mathrm{d}v^2\)。取 \(n_t=(-\sin(tu),0,\cos(tu))\)，则 \(II_t=t\,\mathrm{d}u^2\)。当 \(|t|a<\pi/2\) 时第一坐标对 \(u\) 严格递增，故每个 \(X_t\) 还是嵌入。这个例子既给出真实的局部等距形变，又清楚显示被改变的数据是第二基本形式。
\end{example}

\cref{g1:fig:plane-cylinder-bending}画出同一参数网格在 \(t=0\) 与 \(t=1\) 时的样子。网格线在图纸上的投影长度可以不同；被保持的是曲面自身的 \(I\)，并非观看图像时的平面投影距离。

\input{Figures/Chapter08-PlaneCylinder}

如果在整个参数平面上使用 \(C\)，则 \(C(u+2\pi R,v)=C(u,v)\)。它仍逐点保持度量，却不再单射，因而不是平面到整张圆柱的等距。更强地，平面与整张圆柱之间不存在任何整体微分同胚：平面可收缩，而圆柱投影到圆周的映射与固定高度的圆周包含互为收缩关系，绕轴一周的道路不可能收缩成点。后一断言可直接由\cref{g1:thm:circle-fundamental}中 \(\pi_1(S^1)\cong\mathbb Z\) 及投影、包含的复合为恒等得到。这个拓扑障碍与局部度量计算并不矛盾。

\subsection{一阶、局部和整体刚性问题的区分}
\label{g1:subsec:0805:03}

“刚性”必须说明保持什么、在哪个区域比较以及是否允许改变参数。曲面基本定理的唯一性固定了参数域上的两种形式：若两个浸入的 \(I,II\) 相同，才比较它们是否由一个固定的空间刚性运动联系。若给定一个参数对应 \(\phi\)，应先将第二个浸入换成 \(\widetilde X\circ\phi\)，再逐项比较两种形式。

\begin{definition}[无穷小等距变动]\label{g1:def:infinitesimal-bending}
	设 \(X:U\rightarrow\mathbb R^3\) 为光滑浸入，\(V:U\rightarrow\mathbb R^3\) 为沿 \(X\) 的光滑向量场。若
	\begin{equation}\label{g1:eq:infinitesimal-bending}
		\langle\partial_iV,X_j\rangle+\langle X_i,\partial_jV\rangle=0,\qquad i,j=1,2,
	\end{equation}
	称 \(V\) 为一个\term[wuqiongxiaodengjubiandong]{无穷小等距变动}[infinitesimal bending]。若 \(V=a+\Omega X\)，其中 \(a\in\mathbb R^3\) 固定、\(\Omega\) 为固定斜对称矩阵，称其为刚性运动产生的变动。
\end{definition}

几何上 \(V(q)\) 是曲面点最初的移动速度。若光滑族 \(X_t\) 满足 \(X_0=X\)，则对 \(g_{ij}(t)=\langle(X_t)_i,(X_t)_j\rangle\) 求 \(t\) 导数，正好得到\eqref{g1:eq:infinitesimal-bending}，其中 \(V=\partial_tX_t|_{t=0}\)。因此真正的等距形变一定产生无穷小等距变动。斜对称矩阵给出
\[
	\langle\Omega X_i,X_j\rangle+\langle X_i,\Omega X_j\rangle=0,
\]
所以刚性运动产生的变动总满足这个线性方程。

\begin{example}[一阶条件并不足够]\label{g1:exm:nonintegrable-plane-bending}
	对平面 \(X(u,v)=(u,v,0)\)，任意光滑函数 \(f\) 都给出 \(V=f(u,v)(0,0,1)\)，并满足\eqref{g1:eq:infinitesimal-bending}，因为 \(\partial_iV\) 总与平面垂直。它不一定来自真正的等距形变。取 \(f=u^2+v^2\)，假设有对 \((t,u,v)\) 光滑的浸入族 \(X_t\)，使 \(X_0=X\)、\(\partial_tX_t|_0=V\)，且 \(I_t=\mathrm{d}u^2+\mathrm{d}v^2\)。取由参数定向决定的法向 \(n_t\)。因为 \(X_{ij}|_{t=0}=0\)，
	\[
		h_{ij}(0)=0,\qquad
		\left.\partial_t h_{ij}(t)\right|_0
		=\langle\partial_i\partial_jV,(0,0,1)\rangle=f_{ij}.
	\]
	因此在任一固定参数点有 \(h(t)=2tI_2+O(t^2)\)，从而 \(\det h(t)=4t^2+O(t^3)\)。可是固定的平面度量具有 \(K=0\)，由\cref{g1:thm:gauss-equation}必须有 \(\det h(t)=K\det g=0\)，矛盾。这里一阶约束全满足，二阶相容性却阻止了真实形变。
\end{example}

相比之下，\cref{g1:exm:plane-cylinder-bending}在 \(t=0\) 的变动为 \(V=(u^2/2)(0,0,1)\)，它确实来自整个等距族；其二阶导数矩阵的行列式为零。这个比较只说明上述必要障碍在后一例中消失，不能反过来把“行列式为零”未经证明地宣布为任意形变都可积分的充分条件。

\begin{proposition}[两种形式相同的整体唯一性]\label{g1:prop:surface-global-uniqueness}
	设 \(U\subseteq\mathbb R^2\) 为连通开集，\(X,\widetilde X:U\rightarrow\mathbb R^3\) 为光滑浸入，均取参数定向的单位法向。若两者在同一参数下具有相同的第一、第二基本形式，则存在一个固定 \(Q\in \operatorname{SO}(3)\) 和 \(a\in\mathbb R^3\)，使 \(\widetilde X=QX+a\) 在整个 \(U\) 上成立。这里不要求 \(U\) 单连通。
\end{proposition}

\begin{proof}
	在一点由两组标架与位置确定 \(Q,a\)，再将 \(\widetilde X\) 替换为 \(Q^{-1}(\widetilde X-a)\)。设两者的位置、两个一阶偏导及法向全部相同的点集为 \(A\)。这个集合非空，由连续性为闭集。每个 \(q\in A\) 都有一个矩形邻域；\cref{g1:thm:surface-fundamental}中指定位置及标架初值的唯一性使两浸入在该矩形内相同，所以 \(A\) 也开。连通性给出 \(A=U\)。只比较已经存在的两个整体解时，不会发生“沿不同道路构造出不同解”的问题，因此无需用单连通性排除运输障碍。
\end{proof}

由此可见，局部存在、整体存在与整体唯一性有不同的条件。相容性给出小邻域上的实现；绕闭路运输能否回到同一标架和同一位置，决定能否将局部实现接成单值的整体实现；两个已经存在的实现若保持全部数据，其唯一性又可由连通性传播。本章没有把仅保持第一基本形式的闭曲面刚性问题归结为这一命题。


\begin{chaptersummary}
	浸入要求两个参数方向在空间中仍然独立，它只保证局部图像模型。第一基本形式是切向量的内积，按坐标换元公式转换；长度、面积、梯度、散度与 Laplace 算子由它确定。参数矩阵的样子可以改变，而内积及其测量含义不改变。局部等距保持这套数据，却可能改变曲面在空间中的形状。

	\ProseParagraphBreak

	固定法向以后，形算子 \(S=-\mathrm{d}n\) 和第二基本形式记录法向弯曲。自伴性给出正交主方向，Euler 公式把任意方向的法曲率写成两个主曲率的加权平均。平均曲率与 Gauss 曲率分别是半迹与行列式；法向反转使前者改号，却不改变后者。脐点、零主曲率和渐近方向的退化须逐项区分，不能从一个点的二阶数据直接推出邻域的形状。

	\ProseParagraphBreak

	结构方程把坐标标架的导数分为切向与法向两部分。协变导数记录切向变化，其系数由度量决定；混合偏导相容性同时给出 Gauss 与 Codazzi 方程。Gauss 方程使两个主曲率的乘积成为内在量，这解释了球面不能局部等距铺平，也解释了平面和圆柱虽有相同的零 Gauss 曲率，仍可有不同的主曲率。

	\ProseParagraphBreak

	反向构造先把给定度量分解为标架的长度约束，再把相容条件转为矩阵系统的平坦性。两次普通微分方程积分、误差的零初值唯一性，以及闭向量值一形式的显式积分，共同得到局部浸入；正定性落实秩为二，列向量的内积和导数恢复全部数据。在同一参数下，两种形式共同决定曲面到保向刚性运动为止。整体单值还可能遇到转动或位移障碍，整体单射和仅保持度量的刚性则是进一步的问题。
\end{chaptersummary}

\BookExercises
以下习题均附完整解答。计算前先固定参数和法向；使用存在或唯一性结论时，区分给定的是第一基本形式、两种基本形式，还是只有它们的一阶变化。

\begin{exercise}\label{g1:ex:oblique-metric-measurements}
	对 \(X(u,v)=(u+v,v,0)\)，求参数单位正方形的面积、两条对角线的长度及函数 \(f(u,v)=u\) 的曲面梯度。说明参数分量 \((1,0)\) 为什么不是该梯度。
\end{exercise}
\begin{solution}
	度量矩阵及其逆为
	\[
		g=\begin{pmatrix}1&1\\1&2\end{pmatrix},\qquad
		g^{-1}=\begin{pmatrix}2&-1\\-1&1\end{pmatrix},\qquad D=1.
	\]
	故单位正方形的像为面积 \(1\) 的平行四边形。两条对角线的参数方向分别是 \((1,1)\) 与 \((1,-1)\)，在参数长度为 \(1\) 的区间上走过，平方速度分别为 \(5\) 与 \(1\)，所以长度为 \(\sqrt5\) 与 \(1\)。

	由\eqref{g1:eq:surface-gradient}，\(\operatorname{grad}_gf=2\partial_u-\partial_v\)，空间向量为 \(2X_u-X_v=(1,-1,0)\)。空间平面上 \(f=x-y\)，确实有这个 Euclid 梯度。如果误用 \(\partial_u\)，则它与 \(\partial_v\) 的内积是 \(g_{12}=1\)，而 \(\mathrm{d}f(\partial_v)=0\)，不满足梯度的定义。
\end{solution}

\begin{exercise}\label{g1:ex:graph-area-gradient}
	设 \(X(u,v)=(u,v,f(u,v))\) 为光滑图像，令 \(w=(f_u,f_v)^{\mathsf T}\)。
	\begin{enumerate}
		\item 证明 \(g=I_2+ww^{\mathsf T}\)，并求 \(D\)、\(g^{-1}\) 与高度函数 \(f\) 的曲面梯度。
		\item 对 \(f=(u^2+v^2)/2\)，求参数圆盘 \(u^2+v^2\le R^2\) 对应曲面片的面积。
	\end{enumerate}
\end{exercise}
\begin{solution}
	1. 由 \(X_u=(1,0,f_u),X_v=(0,1,f_v)\)，立即得到所给矩阵。展开二阶行列式，有 \(D=1+|w|^2\)。直接相乘可验
	\[
		(I_2+ww^{\mathsf T})\left(I_2-\frac{ww^{\mathsf T}}{1+|w|^2}\right)=I_2,
	\]
	所以 \(g^{-1}=I_2-ww^{\mathsf T}/(1+|w|^2)\)。于是梯度的参数分量为 \(w/(1+|w|^2)\)，对应空间向量
	\[
		\operatorname{grad}f=\frac{(f_u,f_v,f_u^2+f_v^2)}{1+f_u^2+f_v^2}.
	\]
	它也等于 \((0,0,1)-\langle(0,0,1),n\rangle n\)，其中 \(n=(-f_u,-f_v,1)/\sqrt{1+|w|^2}\)，与投影解释相符。

	2. 此时 \(D=1+u^2+v^2\)，用平面极坐标积分得
	\[
		\operatorname{Area}=2\pi\int_0^Rr\sqrt{1+r^2}\,\mathrm{d}r
		=\frac{2\pi}{3}\bigl((1+R^2)^{3/2}-1\bigr).
	\]
	被积函数中的一个 \(r\) 来自参数平面的极坐标换元，\(\sqrt{1+r^2}\) 才是图像的面积伸缩因子；两者不可遗漏或混同。
\end{solution}

\begin{exercise}\label{g1:ex:conformal-laplacian}
	在平面开集 \(U\) 上取 \(g=e^{2\phi}(\mathrm{d}u^2+\mathrm{d}v^2)\)，其中 \(\phi\) 光滑。求 \(\operatorname{grad}_gf\) 和 \(\Delta_gf\)，并证明对任意 \(f\in C_c^\infty(U)\)，积分 \(\int_U|\operatorname{grad}_gf|_g^2\,\mathrm{d}A_g\) 与 \(\phi\) 无关。由此判断，在这个度量下满足 \(\Delta_gf=0\) 的函数是否与 Euclid 度量下相同。
\end{exercise}
\begin{solution}
	这里 \(g^{ij}=e^{-2\phi}\delta_{ij}\)，\(\sqrt D=e^{2\phi}\)。因此
	\[
		\operatorname{grad}_gf=e^{-2\phi}(f_u\partial_u+f_v\partial_v),\qquad
		\Delta_gf=-e^{-2\phi}(f_{uu}+f_{vv}).
	\]
	散度中的乘积 \(\sqrt D\,g^{ij}\) 恰为 \(\delta_{ij}\)，这就是没有出现 \(\phi\) 的导数的原因。又
	\[
		|\operatorname{grad}_gf|_g^2=e^{-2\phi}(f_u^2+f_v^2),
		\qquad
		\int_U|\operatorname{grad}_gf|_g^2\,\mathrm{d}A_g
		=\int_U(f_u^2+f_v^2)\,\mathrm{d}u\,\mathrm{d}v.
	\]
	因 \(e^{-2\phi}>0\)，\(\Delta_gf=0\) 等价于 \(f_{uu}+f_{vv}=0\)，解的集合确实相同；不过算子本身和度量长度通常不同，不能因此把两个度量判为等距。
\end{solution}

\begin{exercise}\label{g1:ex:area-angle-not-isometry}
	在两个标准 Euclid 平面之间，分别考虑 \(\phi(u,v)=(2u,v/2)\) 和 \(\psi(u,v)=(2u,2v)\)。判断它们是否保持面积、是否保持夹角、是否局部等距，并为每个否定结论给出具体的方向或区域。
\end{exercise}
\begin{solution}
	\(\phi\) 的导数行列式为 \(1\)，故由面积换元保持每个区域的面积。但 \((1,0)\) 的长度从 \(1\) 变为 \(2\)，所以不是局部等距。原来正交的方向 \(a=(1,1)\)、\(b=(1,-1)\) 映为 \((2,1/2)\)、\((2,-1/2)\)，内积为 \(15/4\ne0\)，故也不保持夹角。

	\(\psi\) 把内积乘以 \(4\)，向量长度乘以 \(2\)，夹角公式中的因子抵消，故保持夹角；但单位正方形面积变成 \(4\)，单位线段长度变成 \(2\)，所以既不保持面积，也不是局部等距。单独保持面积或单独保持夹角都不足以代替第一基本形式的完整条件。
\end{solution}

\begin{exercise}\label{g1:ex:shape-oblique-basis}
	在某曲面的一点，以坐标切基 \(X_1,X_2\) 表示的两基本形式矩阵为
	\[
		g=\begin{pmatrix}2&1\\1&1\end{pmatrix},\qquad h=\begin{pmatrix}2&0\\0&-1\end{pmatrix}.
	\]
	求形算子矩阵、主曲率及一组正交单位主向量，并求沿 \(X_1,X_2\) 的法曲率。形算子矩阵不是通常的对称矩阵，为什么不违反其自伴性？
\end{exercise}
\begin{solution}
	因为 \(\det g=1\)，
	\[
		[S]=g^{-1}h=\begin{pmatrix}1&-1\\-1&2\end{pmatrix}
		\begin{pmatrix}2&0\\0&-1\end{pmatrix}
		=\begin{pmatrix}2&1\\-2&-2\end{pmatrix}.
	\]
	迹为 \(0\)，行列式为 \(-2\)，所以主曲率为 \(k_\pm=\pm\sqrt2\)。对应的特征向量可取
	\[
		w_\pm=X_1+(k_\pm-2)X_2.
	\]
	用给定的 \(g\) 而不是把系数当作直角坐标计算，得到
	\[
		\langle w_+,w_-\rangle
		=2+(k_+-2)+(k_--2)+(k_+-2)(k_--2)=0,
	\]
	且 \(\lVert w_\pm\rVert^2=2+2(k_\pm-2)+(k_\pm-2)^2=4-2k_\pm>0\)。于是 \(e_\pm=w_\pm/\sqrt{4-2k_\pm}\) 是所求正交单位主向量。

	两个坐标方向的法曲率分别为 \(h_{11}/g_{11}=1\) 和 \(h_{22}/g_{22}=-1\)，均不同于主曲率，说明坐标方向并非主方向。自伴性在这个基下应写成 \([S]^{\mathsf T}g=g[S]\)，两边都等于 \(h\)；只有 \(g=I\) 的正交单位基才把这个条件化成矩阵自身对称。
\end{solution}

\begin{exercise}\label{g1:ex:graph-mean-curvature}
	设 \(X(u,v)=(u,v,f(u,v))\) 为光滑图像，取向上法向。推导
	\[
		H=\frac{(1+f_v^2)f_{uu}-2f_uf_vf_{uv}+(1+f_u^2)f_{vv}}
		{2(1+f_u^2+f_v^2)^{3/2}},
		\qquad
		K=\frac{f_{uu}f_{vv}-f_{uv}^2}{(1+f_u^2+f_v^2)^2}.
	\]
	对 \(f(u,v)=u^2-v^2\) 求 \(H,K\)，说明一个点的 \(H=0\) 既不能直接推出整个邻域的 \(H=0\)，也不能推出曲面在该点附近为平面。
\end{exercise}
\begin{solution}
	由\cref{g1:exm:graph-second-form}，令 \(W=\sqrt{1+f_u^2+f_v^2}\)，则 \(E=1+f_u^2\)、\(F=f_uf_v\)、\(G=1+f_v^2\)、\(D=W^2\)，并有 \(h_{ij}=f_{ij}/W\)。分别代入\eqref{g1:eq:surface-HK-formulas}，分子的第二基本形式贡献一个或两个 \(W^{-1}\)，便得到题设的 \(H,K\) 公式。

	对 \(f=u^2-v^2\)，有 \(f_u=2u\)、\(f_v=-2v\)、\(f_{uu}=2\)、\(f_{vv}=-2\)、\(f_{uv}=0\)。因此
	\[
		H=\frac{4(v^2-u^2)}{(1+4u^2+4v^2)^{3/2}},\qquad
		K=\frac{-4}{(1+4u^2+4v^2)^2}.
	\]
	原点 \(H=0\)，但任意邻域都含 \(v^2\ne u^2\) 的点，故 \(H\) 在该邻域不恒为零；而 \(K<0\) 已说明此曲面不是平面。这些计算分别排除了两个推断，原点的平均曲率为零只表示两个主曲率在该点相互抵消。
\end{solution}

\begin{exercise}\label{g1:ex:helicoid-asymptotic-curves}
	先证明曲面内的任意正则直线段都是渐近线。再设 \(a>0\)，考虑
	\[
		X(u,v)=(u\cos v,u\sin v,av),\qquad u>0,\quad v\in\mathbb R.
	\]
	取 \(X_u\times X_v\) 确定的法向，求两基本形式及 \(H,K\)。证明两族坐标曲线都是渐近线，并说明其中一族不是直线，从而否定上述结论的逆命题。
\end{exercise}
\begin{solution}
	直线段改用弧长后加速度为零，由\cref{g1:prop:surface-normal-acceleration}知 \(\mathrm{II}(\gamma',\gamma')=0\)，所以是渐近线。

	对题设曲面，记 \(W=\sqrt{u^2+a^2}\)。有
	\[
		X_u=(\cos v,\sin v,0),\qquad X_v=(-u\sin v,u\cos v,a),\qquad
		n=\frac{(a\sin v,-a\cos v,u)}W.
	\]
	由 \(E=1,F=0,G=W^2\) 可知浸入处处正则。又
	\[
		X_{uu}=0,\qquad X_{uv}=(-\sin v,\cos v,0),\qquad
		X_{vv}=(-u\cos v,-u\sin v,0),
	\]
	与 \(n\) 作内积得到 \(h_{11}=h_{22}=0\)、\(h_{12}=-a/W\)。所以
	\[
		\mathrm I=(\mathrm du)^2+(u^2+a^2)(\mathrm dv)^2,\qquad
		\mathrm{II}=-\frac{2a}{W}\,\mathrm du\,\mathrm dv,\qquad
		H=0,\quad K=-\frac{a^2}{(u^2+a^2)^2}.
	\]
	因 \(a/W\ne0\)，渐近方向恰由 \(\mathrm du\,\mathrm dv=0\) 给出，故 \(u\) 固定与 \(v\) 固定的两族坐标曲线都是渐近线。固定 \(v\) 得到直线；固定 \(u=u_0>0\) 则得到圆柱螺线 \((u_0\cos v,u_0\sin v,av)\)，其曲率为 \(u_0/(u_0^2+a^2)>0\)，不是直线。这与渐近条件并不矛盾：它的加速度沿曲面的切平面，而不是沿曲面法向。
\end{solution}

\begin{exercise}\label{g1:ex:gauss-map-area-factor}
	设光滑浸入 \(X:U\to\mathbb R^3\) 取单位法向 \(n=(X_1\times X_2)/\sqrt D\)。证明
	\[
		n_1\times n_2=K(X_1\times X_2),\qquad
		\lVert n_1\times n_2\rVert=|K|\sqrt D.
	\]
	解释后一式中 \(|K|\) 的面积意义，并用图像 \(z=x^2\) 说明：\(K\) 恒为零也不能推出 Gauss 映射为常值。
\end{exercise}
\begin{solution}
	以 \(X_1,X_2\) 为基，\(n_i=-S(X_i)\)。二维线性变换对有向平行四边形的面积乘以行列式；在这里直接按两个系数展开叉积，也得到
	\[
		n_1\times n_2=(-S(X_1))\times(-S(X_2))
		=\det(-S)(X_1\times X_2)=K(X_1\times X_2).
	\]
	两边取范数，并用 \(\lVert X_1\times X_2\rVert=\sqrt D\)，即得第二式。曲面原来的面积密度为 \(\sqrt D\)，Gauss 映射在单位球面上扫过的参数面积密度为 \(\lVert n_1\times n_2\rVert\)；两者之比是 \(|K|\)。这个说法是局部的面积变化率，若不同参数点映到同一法向，参数面积积分仍会计入各次覆盖。

	对 \(X(u,v)=(u,v,u^2)\)，有
	\[
		n(u,v)=\frac{(-2u,0,1)}{\sqrt{1+4u^2}},\qquad
		h=\frac1{\sqrt{1+4u^2}}\begin{pmatrix}2&0\\0&0\end{pmatrix}.
	\]
	因 \(\det h=0\)，处处 \(K=0\)，但 \(n\) 随 \(u\) 变化。这里 Gauss 映射的微分秩为 \(1\)，所以它没有扫出二维面积，却仍扫过一条球面曲线。行列式为零只表示面积塌缩，并不表示整个微分为零。
\end{solution}


\begin{exercise}\label{g1:ex:polar-parallel-vector}
	在去掉原点的平面上，令 \(e_r=(\cos\theta,\sin\theta,0)\)、\(e_\theta=(-\sin\theta,\cos\theta,0)\)。沿圆周 \(r=r_0>0\)，写切向量场 \(V(\theta)=a(\theta)e_r+b(\theta)e_\theta\)。利用本章的协变导数求出 \(\nabla_{\mathrm d/\mathrm d\theta}V=0\) 的所有解，并解释为什么分量改变而几何向量可以不变。这里沿参数曲线的协变导数指普通导数向平面切空间的投影。
\end{exercise}
\begin{solution}
	平面切空间恒定，普通导数已经是切向量。由 \(\mathrm de_r/\mathrm d\theta=e_\theta\)、\(\mathrm de_\theta/\mathrm d\theta=-e_r\)，
	\[
		\nabla_{\mathrm d/\mathrm d\theta}V
		=(\dot a-b)e_r+(\dot b+a)e_\theta.
	\]
	所以方程等价于 \(\dot a=b\)、\(\dot b=-a\)，其全部解为
	\[
		a=A\cos\theta+B\sin\theta,\qquad
		b=-A\sin\theta+B\cos\theta,
	\]
	其中 \(A,B\) 为常数。可直接代入验证，而任意初值唯一确定 \(A,B\)，也可由\cref{g1:lem:linear-frame-ode}的唯一性确认没有其他解。代回原式，得到 \(V=(A,B,0)\)，确为固定向量。变化的 \(a,b\) 只是相对于旋转基的分量；如果错误地只令 \(\dot a=\dot b=0\)，通常得到的是随圆周一起转动的向量。
\end{solution}

\begin{exercise}\label{g1:ex:two-conformal-metrics}
	\begin{enumerate}
		\item 在 \(\mathbb R^2\) 上考虑 \(g=e^{2u}(\mathrm du^2+\mathrm dv^2)\)。求它的曲率，并验证 \(X(u,v)=(e^u\cos v,e^u\sin v,0)\) 实现这个度量。这个参数化是否是到其像的整体一一对应？
		\item 在上半平面 \(v>0\) 上考虑 \(g=v^{-2}(\mathrm du^2+\mathrm dv^2)\)。求它的曲率，判断它能否局部等距于某个 Euclid 平面开集。
	\end{enumerate}
\end{exercise}
\begin{solution}
	(1) 此时 \(\phi=u\)，所以\eqref{g1:eq:conformal-gauss-curvature}给出 \(K_g=0\)。直接求导，
	\[
		X_u=e^u(\cos v,\sin v,0),\qquad
		X_v=e^u(-\sin v,\cos v,0),
	\]
	二者正交且长度均为 \(e^u\)，所以 \(X\) 是浸入并诱导所给度量。然而 \(X(u,v+2\pi)=X(u,v)\)，它不是整体一一对应。这是局部度量完全实现而参数仍有整体重复的例子。

	(2) 此时 \(\phi=-\log v\)，有 \(\phi_{uu}=0\)、\(\phi_{vv}=v^{-2}\)，故 \(K_g=-v^2\cdot v^{-2}=-1\)。Euclid 平面开集的曲率为零。若存在局部等距，重参数化后度量一致，由\cref{g1:prop:surface-curvature-two-dim}及曲率的坐标公式，曲率也应一致，矛盾。这里仅在参数域上计算给定度量，不需要先假定它在三维空间中存在整体浸入。
\end{solution}

\begin{exercise}\label{g1:ex:gauss-codazzi-independent}
	在参数域 \(U=\{(u,v):v>0\}\) 上取平坦度量 \(g=\mathrm du^2+\mathrm dv^2\)。
	\begin{enumerate}
		\item 令 \(h=v\,\mathrm du^2\)，证明它满足 Gauss 方程，但不满足 Codazzi 方程。
		\item 令 \(h=\mathrm du^2+\mathrm dv^2\)，证明它满足 Codazzi 方程，但不满足 Gauss 方程。
	\end{enumerate}
	由此判断两组数据能否作为某个曲面片的两基本形式。
\end{exercise}
\begin{solution}
	度量系数恒定，所以全部 \(\Gamma^k_{ij}=0\)，从而 \(R=0\)、\(K_g=0\)，且 \(\nabla_i h_{jk}=\partial_i h_{jk}\)。

	(1) 矩阵 \(h=\operatorname{diag}(v,0)\) 的行列式为零，故 \(K_g=\det(g^{-1}h)=0\)。由\cref{g1:prop:surface-curvature-two-dim}，这就是完整的 Gauss 方程。然而
	\[
		\nabla_1h_{21}=0,\qquad \nabla_2h_{11}=1,
	\]
	不满足 Codazzi 方程。

	(2) 矩阵 \(h=I\) 恒定，所有协变导数均为零，故 Codazzi 方程成立。但 \(\det(g^{-1}h)=1\ne0=K_g\)，Gauss 方程不成立。两组数据分别违反一个必要条件，因此均不能在任何非空开子域上作为光滑曲面片的两基本形式。
\end{solution}

\begin{exercise}\label{g1:ex:totally-umbilic-surface}
	设 \(U\subseteq\mathbb R^2\) 是连通开集，\(X\colon U\rightarrow\mathbb R^3\) 是光滑浸入，取光滑单位法向 \(n\)。假设每一点都是脐点，即 \(S=\lambda I\)，其中 \(\lambda\) 为光滑函数。证明 \(\lambda\) 为常数，并证明 \(X(U)\) 包含在某个平面或某个球面中。说明只在一个点为脐点为什么不足以得到这个结论。
\end{exercise}
\begin{solution}
	由 \(\nabla_V(\lambda W)=V(\lambda)W+\lambda\nabla_VW\)，有 \((\nabla_VS)W=V(\lambda)W\)。Codazzi 方程的形算子形式因此化为
	\[
		V(\lambda)W=W(\lambda)V.
	\]
	取 \(V=X_1,W=X_2\)，二者线性无关，故 \(\partial_1\lambda=\partial_2\lambda=0\)。函数 \(\lambda\) 在每个小圆盘内为常数，因而是局部常值函数；连通性使它在整个 \(U\) 上为同一个常数。

	若 \(\lambda=0\)，则 \(\mathrm{d}n=-S\,\mathrm{d}X=0\)，同理 \(n=n_0\) 为常向量。又 \(\partial_i\langle X,n_0\rangle=\langle X_i,n_0\rangle=0\)，所以 \(\langle X,n_0\rangle=c\)，像包含在平面内。

	若 \(\lambda\ne0\)，则
	\[
		\partial_i\left(X+\frac1\lambda n\right)
		=X_i-\frac1\lambda S X_i=0.
	\]
	故 \(X+n/\lambda=a\) 为常向量，且 \(\lVert X-a\rVert=1/|\lambda|\)，像包含在以 \(a\) 为心的球面上。这里的结论是像的包含关系，并不要求参数化整体单射。

	若只在一点有 \(S=\lambda I\)，便不能在邻域内对这个恒等式求导，Codazzi 方程也不能推出邻域中的 \(\lambda\) 为常数。例如图像 \(X(u,v)=(u,v,u^3+v^3)\) 在原点的第一导数、第二导数使 \(S(0)=0\)，原点为脐点；当 \(u=v\ne0\) 时其 Gauss 曲率为
	\[
		K(u,u)=\frac{36u^2}{(1+18u^4)^2}>0,
	\]
	而原点的曲率为零。任何包含原点的这类邻域都既不是平面片，也不可能包含在一个固定球面内，因为球面的曲率是固定的正数。
\end{solution}

\begin{exercise}
	\label{g1:ex:constant-surface-data}
	在 \(\mathbb R^2\) 上给定
	\[
		g=\mathrm{d}u^2+\mathrm{d}v^2,\qquad h=a\,\mathrm{d}u^2+2b\,\mathrm{d}u\,\mathrm{d}v+c\,\mathrm{d}v^2,
	\]
	其中 \(a,b,c\) 是实常数。确定这些数据可局部实现的充要条件，并在条件成立时给出一个定义于整个 \(\mathbb R^2\) 的实现。
\end{exercise}

\begin{solution}
	度量系数常数使 \(\Gamma=0\)、\(K_g=0\)，常数 \(h\) 自动满足 Codazzi 条件。Gauss 条件要求 \(ac-b^2=0\)，故这是必要条件。以下构造同时说明充分性。

	若 \(a=b=c=0\)，取 \(X(u,v)=(u,v,0)\)。否则对称矩阵 \(h\) 的两个特征值为 \(\lambda\ne0\) 和 \(0\)。选择单位特征向量 \(t=(t_1,t_2)\) 对应 \(\lambda\)，再取 \(q=(-t_2,t_1)\)；令 \(\xi=t_1u+t_2v\)、\(\eta=-t_2u+t_1v\)。这是保向正交坐标变换，并且 \(g=\mathrm{d}\xi^2+\mathrm{d}\eta^2\)、\(h=\lambda\,\mathrm{d}\xi^2\)。定义
	\[
		X(u,v)=\left(\frac{\sin(\lambda\xi)}{\lambda},\,\eta,\,
		\frac{1-\cos(\lambda\xi)}{\lambda}\right).
	\]
	有 \(X_\xi=(\cos\lambda\xi,0,\sin\lambda\xi)\)、\(X_\eta=(0,1,0)\)，所以第一基本形式正确；法向为 \(n=(-\sin\lambda\xi,0,\cos\lambda\xi)\)，而 \(X_{\xi\xi}=\lambda n\)、\(X_{\xi\eta}=X_{\eta\eta}=0\)，所以第二基本形式也正确。变换 \((u,v)\mapsto(\xi,\eta)\) 保向，故不会额外改变所定法向。所得浸入在整个平面上成立，但当 \(\lambda\ne0\) 时会周期性覆盖圆柱，不能据此称为嵌入。
\end{solution}

\begin{exercise}
	\label{g1:ex:simply-connected-surface-overlap}
	在整个 \(\mathbb R^2\) 上考虑数据 \(g=e^{2u}(\mathrm{d}u^2+\mathrm{d}v^2)\)、\(h=0\)。构造一个实现它们的浸入，并证明没有实现同样数据、采用参数叉积法向的单射浸入。解释这与曲面基本定理是否矛盾。
\end{exercise}

\begin{solution}
	取 \(X(u,v)=(e^u\cos v,e^u\sin v,0)\)。其偏导为
	\begin{align*}
		X_u&=e^u(\cos v,\sin v,0),\\
		X_v&=e^u(-\sin v,\cos v,0).
	\end{align*}
	内积矩阵是 \(e^{2u}I\)，叉积为 \((0,0,e^{2u})\ne0\)，所以它是所需浸入，单位法向是常向量 \((0,0,1)\)，第二基本形式为零。

	设 \(\widetilde X\) 是另一个实现。\(\mathbb R^2=\mathbb R\times\mathbb R\) 本身就是开矩形，故\cref{g1:thm:surface-fundamental}的唯一性在整个域上适用，得到 \(\widetilde X=QX+b\)。但 \(X(u,v+2\pi)=X(u,v)\)，所以 \(\widetilde X\) 同样不是单射。基本定理保证的是微分秩为二的浸入；局部点可分辨并不排除远处参数点映到同一点。这里没有整体单值障碍，只有全局单射性的失败。
\end{solution}

\begin{exercise}
	\label{g1:ex:annulus-cylinder-period}
	在环带 \(A=\{1<r<2\}\) 上给定 \(g=\mathrm{d}r^2+\mathrm{d}\theta^2\)、\(h=\lambda\,\mathrm{d}\theta^2\)，其中 \(\lambda\in\mathbb R\) 为常数，\(\mathrm{d}r,\mathrm{d}\theta\) 按\cref{g1:exm:surface-monodromy-obstructions}理解。确定哪些 \(\lambda\) 允许数据由环带上的单值浸入实现；在可以实现的情形中，判断哪些允许单射实现。
\end{exercise}

\begin{solution}
	局部坐标 \((r,\theta)\) 中 \(g,h\) 的系数都是常数，\(\Gamma=0\)，且 \(K_g=\det h=0\)，故所有 \(\lambda\) 都满足局部相容条件。在展开角度的矩形 \((1,2)\times\mathbb R\) 上，若 \(\lambda\ne0\)，取
	\[
		\widetilde X(r,\theta)=
		\left(r,\frac{\sin\lambda\theta}{\lambda},
		\frac{1-\cos\lambda\theta}{\lambda}\right).
	\]
	其偏导为 \(\widetilde X_r=(1,0,0)\)、\(\widetilde X_\theta=(0,\cos\lambda\theta,\sin\lambda\theta)\)，参数法向为 \(n=(0,-\sin\lambda\theta,\cos\lambda\theta)\)，且 \(\widetilde X_{\theta\theta}=\lambda n\)，因此实现给定数据。它在 \(\theta\mapsto\theta+2\pi\) 下为周期函数，当且仅当 \(\lambda\in\mathbb Z\)。对于 \(\lambda=0\)，实现为 \(\widetilde X=(r,\theta,0)\)，有非零的平移周期，不能下降到环带。

	任意其他实现提升到展开角度的矩形后，均由\cref{g1:thm:surface-fundamental}与上述实现相差固定保向刚性变换，故不能改变周期条件。于是恰好在 \(\lambda\in\mathbb Z\setminus\{0\}\) 时存在环带上的浸入。当 \(|\lambda|=1\) 时，第一坐标分辨 \(r\)，后两坐标分辨模 \(2\pi\) 的 \(\theta\)，所以实现单射。当 \(|\lambda|\geq2\) 时，\(\theta\) 与 \(\theta+2\pi/|\lambda|\) 在环带上是不同点，却有相同像；刚性变换也不能消除此重叠。因此单射实现恰对应 \(\lambda=\pm1\)。
\end{solution}

\begin{exercise}
	\label{g1:ex:surface-frame-choice}
	设矩形上的数据满足曲面基本定理的条件，\(P,\Omega_i,W\) 如正文构造。另取任意光滑函数 \(O\colon V\rightarrow \operatorname{SO}(3)\)，令 \(\widehat P=OP\)。求相应的 \(\widehat\Omega_i=(\widehat P A_i-\partial_i\widehat P)\widehat P^{-1}\)，并说明在正确对应初值后，用 \(\widehat P\) 代替 \(P\) 得到的是同一个位置函数。
\end{exercise}

\begin{solution}
	由 \((OP)^{-1}=P^{-1}O^{-1}\)，
	\[
		\widehat\Omega_i=O\Omega_iO^{-1}-(\partial_iO)O^{-1}.
	\]
	不能删去第二项；它记录所选正交坐标本身随参数旋转。令 \(\widehat W=WO^{-1}\)，直接求导得
	\[
		\partial_i\widehat W
		=W\Omega_iO^{-1}-WO^{-1}(\partial_iO)O^{-1}
		=\widehat W\widehat\Omega_i.
	\]
	因此对应的初值应取 \(\widehat W(q_0)=W(q_0)O(q_0)^{-1}\)。矩阵 \(\widehat\Omega_i\) 仍斜对称，因为正交共轭保持斜对称性，且对 \(OO^{\mathsf T}=I\) 求导可知 \((\partial_iO)O^{-1}\) 斜对称。最关键的是
	\[
		\widehat W\widehat P=WO^{-1}OP=WP=\mathcal F.
	\]
	故候选切向量与法向完全相同，以同一位置初值积分后，位置函数也完全相同。这里允许 \(O\) 随参数变化，说明正交归一化是计算选择，并不在给定的两基本形式之外增加几何数据。
\end{solution}

\begin{exercise}\label{g1:ex:torus-curvature-zones}
	设 \(a>b>0\)，将单位速率圆形母线
	\[
		r(s)=a+b\cos(s/b),\qquad z(s)=b\sin(s/b)
	\]
	绕 \(z\) 轴旋转。采用\cref{g1:prop:revolution-two-forms}指定的法向，求两个主曲率及 Gauss 曲率，说明曲率在哪些部位为正、为负或为零，并判断有无脐点。
\end{exercise}
\begin{solution}
	写 \(\eta=s/b\)，则 \(r'=-\sin\eta,z'=\cos\eta,r''=-\cos\eta/b,z''=-\sin\eta/b\)，故
	\[
		k_s=\frac1b,\qquad k_\theta=\frac{\cos\eta}{a+b\cos\eta},\qquad
		K=\frac{\cos\eta}{b(a+b\cos\eta)}.
	\]
	因为 \(a>b\)，所有分母正且参数始终正则。远离旋转轴的一侧 \(\cos\eta>0\) 有正曲率，靠近轴的一侧 \(\cos\eta<0\) 有负曲率，最上与最下两条纬线 \(\cos\eta=0\) 有零曲率。若有脐点，则 \(1/b=\cos\eta/(a+b\cos\eta)\)，交叉相乘得 \(a=0\)，与假设矛盾。因此没有脐点。改变法向会使两主曲率同时改号，不改变这些结论。
\end{solution}

\begin{exercise}\label{g1:ex:parallel-surface}
	设 \(X:U\rightarrow\mathbb R^3\) 是光滑浸入，选光滑单位法向 \(n\)，令 \(X_t=X+tn\)，其中 \(t\) 是常数。在 \(X_t\) 仍为浸入的区域，沿用向量 \(n\) 作单位法向，不另外强制它与 \(X_t\) 的参数定向一致。
	\begin{enumerate}
		\item 用 \(X\) 的主曲率写出 \(X_t\) 正则的充要条件及其主曲率。
		\item 求 \(H_t,K_t\)，并解释球面和圆柱在哪个平移量发生退化。
	\end{enumerate}
\end{exercise}
\begin{solution}
	1. 对每个参数方向 \(w\)，
	\[
		\mathrm{d}X_t(w)=\mathrm{d}X(w)+t\,\mathrm{d}n(w)=(\operatorname{id}-tS)\mathrm{d}X(w).
	\]
	在 \(S\) 的正交主方向下，这一线性映射的两个因子为 \(1-tk_1\)、\(1-tk_2\)。因此 \(X_t\) 正则当且仅当二者都不为零；此时新切平面与旧切平面作为向量空间相同，\(n\) 确实仍为单位法向。由 \(S_t \mathrm{d}X_t(w)=-\mathrm{d}n(w)=S \mathrm{d}X(w)\)，得
	\[
		S_t=S(\operatorname{id}-tS)^{-1},\qquad
		k_i(t)=\frac{k_i}{1-tk_i}.
	\]

	2. 记 \(P=(1-tk_1)(1-tk_2)=1-2tH+t^2K\)。相加和相乘给出
	\[
		H_t=\frac{H-tK}{P},\qquad K_t=\frac K P.
	\]
	外向法向的半径 \(R\) 球面有 \(k_1=k_2=-1/R\)，在 \(t=-R\) 时全部塌缩为球心；圆柱有 \(k_1=-1/R,k_2=0\)，也在 \(t=-R\) 时退化，此时圆周方向塌缩而母线方向保留，像落到轴线上。退化值处不能再使用新曲率公式，因为 \(X_t\) 已不满足浸入条件。
\end{solution}

\begin{exercise}\label{g1:ex:plane-infinitesimal-solutions}
	在开矩形 \(U\subseteq\mathbb R^2\) 上取 \(X(u,v)=(u,v,0)\)。求所有无穷小等距变动 \(V=(A,B,C)\)，并确定其中哪些来自空间刚性运动。
\end{exercise}
\begin{solution}
	将 \(X_u=(1,0,0),X_v=(0,1,0)\) 代入\eqref{g1:eq:infinitesimal-bending}，得到
	\[
		A_u=0,\qquad B_v=0,\qquad A_v+B_u=0.
	\]
	矩形内可先写 \(A=a(v),B=b(u)\)。第三式为 \(a'(v)=-b'(u)\)，对任意 \(u,v\) 成立，两边必须是同一个常数 \(c\)。所以全部解为
	\[
		V(u,v)=(cv+a_0,\,-cu+b_0,\,f(u,v)),
	\]
	其中 \(f\) 是任意光滑函数。若 \(V=d+\Omega X\)，斜对称矩阵使第三分量必为 \(d_3+\Omega_{31}u+\Omega_{32}v\)，故 \(f\) 必须为仿射函数。反之任意仿射 \(f=\alpha u+\beta v+\gamma\) 都由
	\[
		d=(a_0,b_0,\gamma),\qquad
		\Omega=\begin{pmatrix}0&c&-\alpha\\-c&0&-\beta\\\alpha&\beta&0\end{pmatrix}
	\]
	实现。非仿射的 \(f\) 因而给出非刚体的一阶变动；这些变动中有些可来自实际等距族，有些不能，正文两个例子已经显示这种差别。
\end{solution}

% !TeX root = ../../Book1.tex
% 纲要标题：## 第9章　曲面联络、测地线与 Gauss–Bonnet
% 纲要位置：几何学计划纲要.md，第 461 行。
% 正文状态：初稿；尚未正式审核。

% 纲要内容开始
% 
% > 本章目标与证明范围：从切向投影建立联络，以平行移动和转角说明曲率的整体意义。二维 Gauss–Bonnet 包括边界和角点的完整证明，不借用后置高维示性类理论。
% 
% 纲要内容结束

\chapter{曲面联络、测地线与 Gauss–Bonnet}
\label{g1:ch:09}
\BookChapterNumber{g1:ch:09}

\begin{chapterabstract}
	沿曲面移动切向量，需要比较不同切平面中的方向。协变导数把这种比较写成平行移动方程，联络一形式则记录移动标架自身的转动。本章建立测地线的初值理论、长度与能量的第一变分以及径向测地线的局部最短性，再把边界切向转角分解为测地曲率与联络积分。局部曲率公式经有限圆盘分解、公共边抵消和顶点计数，成为含边界及角点的 Gauss--Bonnet 定理。最后讨论曲率符号、亏格、向量场零点指数和改变度量时的总曲率不变性。
\end{chapterabstract}

\begin{learninggoals}
	\begin{itemize}
		\item 从切向投影和度量相容性理解沿曲线的协变导数，计算平行移动及绕环路后的旋转。
		\item 证明测地线初值解的存在唯一性与光滑依赖，区分长度、能量、参数选择和端点条件。
		\item 用第一变分与 Gauss 引理说明测地线为何在足够短的范围内最短，并明确局部结论的范围。
		\item 固定联络、曲率、边界方向和外角的符号，完整证明局部及整体 Gauss--Bonnet 公式。
		\item 将总曲率与 Euler 示性数、孤立零点指数联系起来，计算保角改变度量后的内部与边界补偿。
	\end{itemize}
\end{learninggoals}

在球面上沿三个大圆弧走出一个三角形，即使每条边都没有在曲面内主动转弯，最后的方向仍可能与出发时不同。这个现象同时提出了两个问题：沿曲面“保持方向”应如何定义？一条闭合边界的转角又怎样记录它围住的几何？第八章已把曲率写成度量的局部表达式，本章要把它与沿路比较方向、围成区域以及整张曲面的拓扑联系起来。

\ProseParagraphBreak

我们采用 \(\omega(V)=\langle\nabla_Ve_1,e_2\rangle\) 的联络约定，所以曲率关系是 \(\mathrm{d}\omega=-K\,\mathrm{d}A\)。边界按区域在左的方向行进，内角 \(\beta\) 对应外角 \(\alpha=\pi-\beta\)。先在一张坐标图中的圆盘上，借平面转角定理与 Green 公式证明曲率积分公式；随后真正构造分片光滑的有限圆盘分解，把公共边和新增角逐项消去。这样得到的 Gauss--Bonnet 定理包括非定向曲面、光滑边界和有限角点。\footnote{Bonnet，Pierre Ossian Bonnet（1819--1892），法国数学家，研究曲面微分几何、测地曲率与曲率积分公式。}

\ProseParagraphBreak

阅读所需的主要旧结论是第五章圆周角度提升与绕数、第六章曲面分类和 Euler 计数、第七章平面转角定理，以及第八章的度量、协变导数和曲率。使用时会回引具体结果并重述关键条件。常微分方程的光滑依赖、测地线的变分论证与二维 Green 公式所需版本在本章补出；一般流形的微分形式理论和高维示性类理论不是前提。相关教材阅读位置为 Pressley 的平行移动部分\cite[第7.4节]{Pressley2010}、测地线部分\cite[第9.1--9.5节]{Pressley2010}和 Gauss--Bonnet、环路平行移动与向量场零点部分\cite[第13.1--13.4, 13.6--13.7节]{Pressley2010}。

% !TeX root = ../../Book1.tex
% 纲要标题：### 9.1 协变微分和曲面上的移动标架
% 纲要位置：几何学计划纲要.md，第 465 行。
% 正文状态：初稿；尚未正式审核。

% 纲要内容开始
% 
% * 切向投影与度量相容性
% * 沿曲线的平行移动
% * 正交标架的联络一形式
% * 曲率二形式与局部 holonomy
% 
% 纲要内容结束

\section{协变微分和曲面上的移动标架}
\label{g1:sec:0901}

在平面内移动一个箭头，可以要求它在固定直角坐标中的分量不变。曲面上的箭头却必须始终留在当前位置的切平面内；切平面一变，环境空间中的常向量通常就不再合适。第八章的切向投影给出了另一种要求：允许箭头为适应切平面而向法向改变，但不让它在切平面内自行转动。沿路保持这种状态，就是本节要建立的平行移动。

\subsection{切向投影与度量相容性}
\label{g1:subsec:0901:01}

先说明本章的曲面范围。第六章的拓扑曲面默认 Hausdorff、第二可数。给这样的曲面选定坐标图，要求重叠处的坐标变换光滑且具有光滑逆映射，就得到本章所说的\term[guanghuaqumian]{光滑曲面}[smooth surface]。坐标中的切向量按坐标变换的微分换基；光滑正定矩阵 \((g_{ij})\) 若按第八章\eqref{g1:eq:surface-metric-change}在重叠处变换，就在各切平面上给出同一个内积 \(g\)。这样的 \((M,g)\) 称为\term[limanduliangqumian]{黎曼度量曲面}[Riemannian surface]。这里的英文指带黎曼度量的实二维曲面，不是作为复一维对象的 Riemann surface。

\ProseParagraphBreak

正则曲面在 \(\mathbb R^3\) 中的诱导度量是主要例子。不过，平面上按平移识别得到的环面也有相容的欧氏度量；讨论它的内在几何时，不必先把这张带度量的曲面放进三维空间。有光滑边界时用半平面坐标，要求度量系数光滑延拓到坐标域的邻域。后面还会考虑曲面内具有有限角点的区域，边界各光滑段正则，角点内角在 \((0,2\pi)\) 内，因而不包括尖点。

\ProseParagraphBreak

以下指标 \(i,j,k\) 取 \(1,2\)，重复指标求和。回顾\cref{g1:prop:surface-connection-properties}：给定度量后，唯一同时满足度量相容和无挠性的协变导数为
\[
	\nabla_VW=V^i\bigl(\partial_iW^k+\Gamma^k_{ij}W^j\bigr)\partial_k,
	\qquad
	\Gamma^k_{ij}=\frac12g^{k\ell}(\partial_i g_{j\ell}+\partial_j g_{i\ell}-\partial_\ell g_{ij}).
\]
这里 \((g^{ij})=(g_{ij})^{-1}\)。第八章已经证明这套公式在换坐标后给出同一个运算；对于诱导度量，它就是 \(\nabla_VW=(\mathrm{d}W(V))^\top\)。因此本章下面的坐标计算也适用于只给定度量的曲面。

\begin{definition}
	\label{g1:def:surface-along-derivative}
	设 \(\gamma:I\rightarrow M\) 是光滑曲线，\(W(t)\in \mathrm{T}_{\gamma(t)}M\) 是沿曲线的光滑切向量场。在坐标中写 \(\gamma(t)=(u^1(t),u^2(t))\)、\(W(t)=W^k(t)\partial_k|_{\gamma(t)}\)，定义\term[yanquxianxiebiandaoshu]{沿曲线的协变导数}[covariant derivative along a curve]为
	\begin{equation}
		\label{g1:eq:along-covariant-derivative}
		\mathrm{D}_tW\coloneqq\left(\frac{\mathrm{d}W^k}{\mathrm{d}t}+\Gamma^k_{ij}(\gamma(t))\dot u^iW^j\right)\partial_k.
	\end{equation}
	若曲面已浸入 \(\mathbb R^3\)，此式就是 \(\mathrm{D}_tW=(\mathrm{d}W/\mathrm{d}t)^\top\)。\symidx{\(\mathrm{D}_tW\)：沿曲线的协变导数}
\end{definition}

沿曲线的向量场只给每个参数值配置一个向量。即使曲线经过同一点两次，这两次的向量也可以不同，因此不必把 \(W\) 延拓成曲面邻域内的单值向量场。换坐标时，对 \(W\) 的分量应用链式法则，再代入 Christoffel 系数的变换式，二阶坐标导数项相消，故 \(\mathrm{D}_tW\) 仍是同一个切向量。

\begin{proposition}
	\label{g1:prop:along-derivative-properties}
	在光滑度量曲面上，沿光滑曲线的协变导数满足乘法法则及
	\[
		\frac{\mathrm{d}}{\mathrm{d}t}g(V,W)=g(\mathrm{D}_tV,W)+g(V,\mathrm{D}_tW).
	\]
	对光滑双参数映射 \(F(s,t)\)，有 \(\mathrm{D}_s\partial_tF=\mathrm{D}_t\partial_sF\)。若曲面已定向，以 \(J\) 表示每个切平面内的正向直角旋转，则 \(\mathrm{D}_t(JW)=J(\mathrm{D}_tW)\)。
\end{proposition}

\begin{proof}
	乘法法则直接来自\eqref{g1:eq:along-covariant-derivative}。内积求导时，用第八章已证的 \(\partial_i g_{j\ell}=\Gamma^k_{ij}g_{k\ell}+\Gamma^k_{i\ell}g_{jk}\) 将度量系数的导数分到两个向量上，即得第一式。混合导数的第 \(k\) 个分量分别为
	\[
		\partial_s\partial_tF^k+\Gamma^k_{ij}(F)\partial_sF^i\partial_tF^j,
		\qquad
		\partial_t\partial_sF^k+\Gamma^k_{ij}(F)\partial_tF^i\partial_sF^j.
	\]
	偏导可交换，且 \(\Gamma^k_{ij}=\Gamma^k_{ji}\)，所以两式相等；\(F\) 不必是浸入。

	最后局部选正向正交单位标架 \((e_1,e_2)\)。由内积求导，存在标量 \(a(t)\)，使 \(\mathrm{D}_te_1=ae_2\)、\(\mathrm{D}_te_2=-ae_1\)。将 \(W=xe_1+ye_2\)、\(JW=-ye_1+xe_2\) 分别求导，便得 \(\mathrm{D}_t(JW)=J(\mathrm{D}_tW)\)。
\end{proof}

这里的 \(J\) 仅用度量与定向定义，满足 \(J^2=-\mathrm{id}\)；在取单位法向 \(n\) 的空间曲面上，\(JW=n\times W\)。以后边界的符号都以这个正向直角旋转为准。

\subsection{沿曲线的平行移动}
\label{g1:subsec:0901:02}

\begin{definition}
	沿分片光滑曲线 \(\gamma\) 的连续切向量场 \(W\)，若在每个光滑段上满足 \(\mathrm{D}_tW=0\)，就称为沿 \(\gamma\) 的\term[pingxingxiangliangchang]{平行向量场}[parallel vector field]。从给定初向量出发求出终向量的操作称为\term[pingxingyidong]{平行移动}[parallel transport]，记为 \(P_\gamma\)。\symidx{\(P_\gamma\)：沿曲线的平行移动}
\end{definition}

\begin{theorem}
	\label{g1:thm:surface-parallel-transport}
	设 \((M,g)\) 是光滑度量曲面，\(\gamma:[a,b]\rightarrow M\) 分片光滑。给定 \(w\in \mathrm{T}_{\gamma(a)}M\)，存在唯一沿 \(\gamma\) 的平行场，初值为 \(w\)。所得 \(P_\gamma:\mathrm{T}_{\gamma(a)}M\rightarrow \mathrm{T}_{\gamma(b)}M\) 是线性等距同构；在定向曲面上还保持定向。它在保向重参数化下不变，并满足
	\[
		P_{\gamma_1*\gamma_2}=P_{\gamma_2}\circ P_{\gamma_1},
		\qquad P_{\bar\gamma}=P_\gamma^{-1}.
	\]
	这里沿用第五章的次序，\(\gamma_1*\gamma_2\) 先走 \(\gamma_1\) 再走 \(\gamma_2\)，\(\bar\gamma\) 表示反向道路。
\end{theorem}

\begin{proof}
	在一张坐标图中，平行条件是线性方程
	\[
		\dot W^k=-\Gamma^k_{ij}(\gamma(t))\dot u^iW^j.
	\]
	每个光滑闭段上的系数连续，故\cref{g1:lem:linear-frame-ode}给出指定初值的唯一解。把曲线的坐标邻域拉回参数区间，利用紧区间的有限细分，使每段落入一张图。在接点把所得几何向量换成下一张图的分量作为初值；坐标无关性与唯一性保证相邻解相容，也保证改变细分不会改变结果。

	方程线性，所以解依初值线性。对两个平行场，由上命题得 \(\frac{\mathrm{d}}{\mathrm{d}t}g(V,W)=0\)，于是长度与夹角保持。反向参数化时 \(\mathrm{D}_\tau W=t'(\tau)\mathrm{D}_tW\)，故平行性不变；唯一性给出反向和拼接公式，也给出可逆性。在定向情形，把一个正向单位标架平行移动，两向量始终正交且单位长；其方向符号连续而不可能经过零，故一直保持正向。
\end{proof}

\begin{example}
	\label{g1:exm:parallel-polar-frame}
	在欧氏平面沿圆 \(\gamma(t)=(r\cos t,r\sin t)\) 运动，取 \(e_r=(\cos t,\sin t)\)、\(e_\theta=(-\sin t,\cos t)\)。若 \(W=a e_r+b e_\theta\)，则
	\[
		\mathrm{D}_tW=(a'-b)e_r+(b'+a)e_\theta.
	\]
	由 \(W(0)=(1,0)\) 出发的平行场满足 \(a=\cos t\)、\(b=-\sin t\)，所以 \(W(t)=(1,0)\) 在环境平面内确实不变。反而是分量固定的 \(e_r\) 本身不平行：它一直朝向圆外，因而在转动。分量变化并不等于向量发生了内在转动，必须同时计算标架的变化。
\end{example}

\subsection{正交标架的联络一形式}
\label{g1:subsec:0901:03}

在一张定向坐标图中，先对 \(\partial_1,\partial_2\) 作 Gram--Schmidt 正交化，得到光滑正向正交标架 \((e_1,e_2)\)。内积为常数意味着 \(\nabla_Ve_1\) 沿 \(e_2\)，而 \(\nabla_Ve_2\) 沿 \(-e_1\)。因此所有方向上的标架转动，只需用一个对 \(V\) 线性的数来记录。

\begin{definition}
	\label{g1:def:surface-connection-form}
	对定向光滑度量曲面上的局部正向正交标架，定义\term[lianluoyixingshi]{联络一形式}[connection one-form]为 \(\omega(V)=g(\nabla_Ve_1,e_2)\)。于是
	\begin{equation}
		\label{g1:eq:surface-frame-connection}
		\nabla_Ve_1=\omega(V)e_2,\qquad \nabla_Ve_2=-\omega(V)e_1.
	\end{equation}
\end{definition}

这里“一形式”表示在每点给定一个随点光滑变化的线性函数。在坐标中它就是 \(A\,\mathrm{d}u+B\,\mathrm{d}v\)，其中 \(\mathrm{d}u(V)=V^u\)。两一形式的外积定义为 \((\alpha\wedge\beta)(V,W)=\alpha(V)\beta(W)-\alpha(W)\beta(V)\)，因而 \(\mathrm{d}u\wedge\mathrm{d}v\) 记录两个向量的有向坐标面积。令 \(\vartheta^i(V)=g(V,e_i)\)，则
\[
	\mathrm{d}A=\vartheta^1\wedge\vartheta^2=\sqrt{\det(g_{ij})}\,\mathrm{d}u\wedge\mathrm{d}v
\]
是正向面积形式。计算非负面积时取其对应的面积元；反向坐标使用 Jacobian 行列式的绝对值。\symidx{\(\mathrm{d}A\)：曲面的面积元}

对一形式定义
\[
	\mathrm{d}(A\,\mathrm{d}u+B\,\mathrm{d}v)
	=(\partial_uB-\partial_vA)\,\mathrm{d}u\wedge\mathrm{d}v.
\]
对函数 \(f\) 则有 \(\mathrm{d}f=f_u\,\mathrm{d}u+f_v\,\mathrm{d}v\)，混合偏导相等给出 \(\mathrm{d}(\mathrm{d}f)=0\)。这些是本章使用的\term[waiweifen]{外微分}[exterior derivative]。换坐标时应用链式法则，二阶坐标导数项两两相消；更具几何意味的同一公式是
\[
	\mathrm{d}\alpha(V,W)=V(\alpha(W))-W(\alpha(V))-\alpha([V,W]).
\]
将 \(V,W\) 的坐标分量展开可验证此式，也直接看出其与坐标无关。我们在这里需要的只是函数、一形式和面积形式，不以掌握高维微分形式为前提。

\begin{proposition}
	\label{g1:prop:connection-frame-change}
	若另一正向正交标架在邻域内写成
	\[
		\widetilde e_1=\cos\varphi\,e_1+\sin\varphi\,e_2,\qquad
		\widetilde e_2=-\sin\varphi\,e_1+\cos\varphi\,e_2,
	\]
	则相应联络形式满足 \(\widetilde\omega=\omega+\mathrm{d}\varphi\)。沿曲线写单位向量 \(W=\cos\theta\,e_1+\sin\theta\,e_2\)，有
	\begin{equation}
		\label{g1:eq:parallel-angle}
		\mathrm{D}_tW=(\theta'+\omega(\dot\gamma))JW.
	\end{equation}
	特别地，平行条件为 \(\theta'=-\omega(\dot\gamma)\)。
\end{proposition}

\begin{proof}
	将两个标架变换式求协变导数，用\eqref{g1:eq:surface-frame-connection}合并正弦、余弦项，得到 \(\nabla_V\widetilde e_1=(V\varphi+\omega(V))\widetilde e_2\)，即第一式。对 \(W\) 作相同展开，得到第二式。沿参数区间的连续角度由\cref{g1:lem:circle-angle-lift}选取；只要 \(W\) 光滑，局部反三角函数又说明提升角度光滑。
\end{proof}

联络形式依赖标架；标架多转了 \(\varphi\)，所记录的转动便多出 \(\mathrm{d}\varphi\)。但 \(\mathrm{d}\widetilde\omega=\mathrm{d}\omega\)。在环带上换架角度未必有全域单值的实函数表示，上述关系仍在每个小邻域成立，外微分不依赖标架的结论因而仍然成立。

\subsection{曲率二形式与局部 holonomy}
\label{g1:subsec:0901:04}

\begin{proposition}
	\label{g1:prop:connection-curvature-form}
	在定向光滑度量曲面上，任意局部正向正交标架的联络形式均满足
	\begin{equation}
		\label{g1:eq:surface-curvature-form}
		\mathrm{d}\omega=-K\,\mathrm{d}A.
	\end{equation}
	这里 \(K\) 是第八章定义的内在 Gauss 曲率，曲率张量约定为 \(R(V,W)=\nabla_V\nabla_W-\nabla_W\nabla_V-\nabla_{[V,W]}\)。
\end{proposition}

\begin{proof}
	对 \(e_1\) 应用曲率定义与\eqref{g1:eq:surface-frame-connection}，两次求导中沿 \(e_1\) 的乘积项抵消，得到
	\[
		R(V,W)e_1=\bigl(V\omega(W)-W\omega(V)-\omega([V,W])\bigr)e_2
		=\mathrm{d}\omega(V,W)e_2.
	\]
	另一方面，\cref{g1:prop:surface-curvature-two-dim}给出 \(R(e_1,e_2)e_1=-Ke_2\)。所以 \(\mathrm{d}\omega(e_1,e_2)=-K\)。二维交替双线性形式由在一组基上的这个值决定，而 \(\mathrm{d}A(e_1,e_2)=1\)，故得结论。
\end{proof}

联络的曲率矩阵是 \(\begin{psmallmatrix}0&-\mathrm{d}\omega\\\mathrm{d}\omega&0\end{psmallmatrix}\)；它记录了先后两个方向求导的差。负号 \(\mathrm{d}\omega=-K\,\mathrm{d}A\) 同时来自本书的曲率次序和 \(\omega=g(\nabla e_1,e_2)\) 的选择。改变其中任何一项约定时，都必须重新核对这个符号。

\begin{lemma}[Green 公式]
	\label{g1:lem:green-formula}
	设 \(D\) 是一张正向坐标图内的紧区域，边界由有限条分片 \(C^2\) 正则简单闭曲线组成，彼此不交；角点内角在 \((0,2\pi)\) 内。若一形式 \(\alpha\) 在 \(D\) 的邻域上为 \(C^1\)，则
	\[
		\int_{\partial D}\alpha=\int_D\mathrm{d}\alpha.
	\]
	边界方向取区域在左的方向，故有孔时，外边界逆时针，孔的边界顺时针。
\end{lemma}

这一公式称为 Green 公式，也就是二维的 Stokes 公式。\footnote{Green，格林，George Green（1793--1841），英国数学家，研究位势理论，提出 Green 函数及相关积分公式。Stokes，斯托克斯，George Gabriel Stokes（1819--1903），生于爱尔兰的英国数学家、物理学家，研究流体力学与光学。}下面用平面微积分证明本章所需的版本。

\begin{proof}
	先设边界光滑。在边界附近选坐标 \((u,t)\)，使区域为 \(t\ge0\)、边界为 \(t=0\)，并保持定向；这由边界正则性及逆函数定理得到。例如边界原为 \(v=f(u)\)，取 \(t=v-f(u)\)。若 \(\alpha=P\,\mathrm{d}u+Q\,\mathrm{d}t\) 的支撑紧含于此坐标邻域，则积分基本定理与累次积分给出
	\[
		\int_{t\ge0}(\partial_uQ-\partial_tP)\,\mathrm{d}u\,\mathrm{d}t
		=\int P(u,0)\,\mathrm{d}u.
	\]
	第一项沿 \(u\) 积分为零，第二项只留下 \(t=0\) 的值；右边正是把区域放在左侧的边界积分。若支撑在内部，同样在整个平面上积分，两个偏导积分都为零。

	把 \(D\) 用有限个内部邻域和这样的边界邻域覆盖，再用支撑在其中的光滑函数 \(\rho_j\) 分解 \(1\)。这里不需调用曲面上的一般分解定理：在坐标平面内取有限个更小圆盘覆盖紧集，以 \(\exp(-1/(1-|x|^2))\) 在单位圆盘内、在外取零并经平移缩放得到非负光滑函数，最后除以它们在紧集附近处处为正的和即可。对 \(\rho_j\alpha\) 应用局部计算并相加，由 \(\sum\rho_j=1\)、\(\sum\mathrm{d}\rho_j=0\) 得到原式。换坐标的积分公式保证所有局部积分相容。

	有角点时，先在每个角点附近用一条长度为 \(O(\varepsilon)\) 的光滑小弧切去角点，并把新接点圆滑化。两侧边界有非零夹角，可在其切线张成的扇形内完成此构造；所删区域面积为 \(O(\varepsilon^2)\)，所改边界的总长度为 \(O(\varepsilon)\)。\(\alpha\) 及 \(\mathrm{d}\alpha\) 在紧邻域上有界，所以新旧边界积分之差趋于零，面积积分之差也趋于零。对光滑截角区域的公式取极限，便得分片光滑情形；这一论证也适用于多个边界分支。
\end{proof}

对任意紧曲面区域，面积积分也可以先作有限坐标分解再相加。在曲面坐标中取上述非负光滑截断函数，支撑紧含于坐标邻域后延为零，再除以有限和，得到坐标覆盖的有限单位分解。以局部表达计算 \(\int\rho_j f\,\mathrm{d}A\)；换坐标的 Jacobian 公式及 \(\sum\rho_j=1\) 保证结果与分解无关。面积元在不可定向曲面上仍有意义，因为它使用绝对面积；带方向的二形式积分则需要定向。

\begin{proposition}
	\label{g1:prop:surface-local-holonomy}
	设 \(D\) 是上述坐标图内由简单分片光滑闭曲线围成的圆盘区域，\(\gamma=\partial D\) 取正向。沿 \(\gamma\) 的平行移动回到起点后是一个平面旋转，其旋转角 \(\beta\) 满足
	\begin{equation}
		\label{g1:eq:surface-holonomy-integral}
		\beta\equiv\int_D K\,\mathrm{d}A\pmod{2\pi}.
	\end{equation}
\end{proposition}

\begin{proof}
	在整张坐标图上将坐标基正交化，得到单值标架。绕行一周，标架回到原值。平行单位向量相对标架的连续角度由\eqref{g1:eq:parallel-angle}满足 \(\theta(b)-\theta(a)=-\int_{\partial D}\omega\)。再用 Green 公式和\eqref{g1:eq:surface-curvature-form}，得到角度增量 \(\int_DK\,\mathrm{d}A\)。终点向量只记录该数模 \(2\pi\) 的值，这便是所述旋转角。
\end{proof}

绕闭路返回同一切平面后所得的变换也称为该闭路的 holonomy。\foreignidx{holonomy}它通常不只取决于闭路的同伦类：同一个坐标圆盘内的闭路都可缩，但围起不同曲率积分的闭路可能给出不同旋转。对缩到点 \(p\) 的小圆盘，选取趋于零的角度代表后，有
\[
	\frac{\beta}{\operatorname{Area}(D)}
	=\frac{1}{\operatorname{Area}(D)}\int_DK\,\mathrm{d}A\longrightarrow K(p),
\]
因为曲率的连续性控制了平均值与 \(K(p)\) 的差。这给出了曲率的一种直接测量方式。

\begin{example}
	\label{g1:exm:spherical-latitude-holonomy}
	半径 \(R\) 的球面用余纬 \(\theta\in(0,\pi)\)、经度 \(\varphi\) 表示，度量为 \(R^2\mathrm{d}\theta^2+R^2\sin^2\theta\,\mathrm{d}\varphi^2\)。取 \(e_1=R^{-1}\partial_\theta\)、\(e_2=(R\sin\theta)^{-1}\partial_\varphi\)，由 Christoffel 公式求得
	\[
		\omega=\cos\theta\,\mathrm{d}\varphi,\qquad
		\mathrm{d}\omega=-\sin\theta\,\mathrm{d}\theta\wedge\mathrm{d}\varphi=-R^{-2}\mathrm{d}A.
	\]
	沿余纬 \(\theta_0\) 处的纬线增加经度一周，平行向量相对这个周期标架的角度变化为 \(-2\pi\cos\theta_0\)。而北球冠的曲率积分是 \(2\pi(1-\cos\theta_0)\)，两者相差 \(2\pi\)，表示同一个旋转。这里的经纬标架在北极无定义，不能将它的 \(\omega\) 直接用于整个球冠的 Green 公式；换用在球冠上光滑的坐标标架，角度代表就与球冠积分一致。赤道绕行的旋转为恒等，但北半球的曲率积分为 \(2\pi\)，这也说明模 \(2\pi\) 不能随意删去。
\end{example}

% !TeX root = ../../Book1.tex
% 纲要标题：### 9.2 测地线与变分
% 纲要位置：几何学计划纲要.md，第 472 行。
% 正文状态：初稿；尚未正式审核。

% 纲要内容开始
% 
% * 测地线方程的存在唯一性
% * 弧长、能量与端点条件
% * 测地曲率及第一变分
% * 测地坐标和径向局部极小性
% 
% 纲要内容结束

\section{测地线与变分}
\label{g1:sec:0902}

在平面上，匀速直线同时具有两种特征：速度沿运动保持不变，较短的直线段又是两端之间的最短路。曲面上不宜直接要求空间中的速度不变，因为这个向量可能立即离开切平面。上一节的平行移动给出了替代条件：让速度沿曲线平行。它定义的是一个微分方程；这一方程与长度极小之间的关系还需要证明，而且最短性首先只在短程成立。

\ProseParagraphBreak

以下讨论的曲面 \(M\) 都带光滑正定度量：可以取正则曲面及其诱导度量，也可以取上一节所说的二维光滑度量曲面。局部计算取坐标 \((u^1,u^2)\)，记 \(g_{ij}=\langle\partial_i,\partial_j\rangle\)。讨论初值与邻域时，所取点均为曲面的内点；边界会在变分的端点条件中单独出现。

\subsection{测地线方程的存在唯一性}
\label{g1:subsec:0902:01}

\begin{definition}
	\label{g1:def:geodesic}
	光滑度量曲面 \(M\) 上的光滑曲线 \(\gamma:I\to M\)，若满足
	\[
		\mathrm{D}_t\dot\gamma=0,
	\]
	就称为\term[cedixian]{测地线}[geodesic]。这里 \(\mathrm{D}_t\) 是\cref{g1:def:surface-along-derivative}定义的沿曲线协变导数，参数 \(t\) 称为仿射参数。常值曲线也包括在内。
\end{definition}

由\cref{g1:prop:along-derivative-properties}的度量相容性，测地线满足
\[
	\frac{\mathrm{d}}{\mathrm{d}t}|\dot\gamma|^2=2\langle \mathrm{D}_t\dot\gamma,\dot\gamma\rangle=0.
\]
所以每条非常值测地线都以恒定的正速度行进，改用弧长后仍是测地线。若 \(\gamma\) 是非常值测地线，重参数曲线 \(\eta(s)=\gamma(\phi(s))\) 满足
\[
	\mathrm{D}_s\dot\eta=\phi''(s)\dot\gamma(\phi(s)),
\]
因而保持测地线方程的正则重参数化恰为 \(\phi(s)=as+b\)、\(a\ne0\)。同一条轨迹以忽快忽慢的速度走过，便可能不再满足上述带参数的定义。

\ProseParagraphBreak

把沿曲线求导的坐标式代入定义，测地线方程成为
\begin{equation}
	\label{g1:eq:geodesic-equation}
	\ddot u^k+\Gamma^k_{ij}(u)\dot u^i\dot u^j=0,\qquad k=1,2.
\end{equation}
这里对 \(i,j\) 求和。\(\Gamma^k_{ij}\) 由\eqref{g1:eq:christoffel-metric}计算，故这一方程只依赖于度量。在嵌入曲面上，它也表示空间加速度全部沿法向；一般并不表示空间加速度为零。

\begin{lemma}[初值的光滑依赖]
	\label{g1:lem:smooth-ode-flow}
	设 \(F:\Omega\to\mathbb R^m\) 是开集 \(\Omega\subset\mathbb R^m\) 上的光滑映射，\(z_0\in\Omega\)。则存在 \(a>0\) 及 \(z_0\) 的邻域 \(B\)，使
	\[
		\partial_t\Phi(t,z)=F(\Phi(t,z)),\qquad \Phi(0,z)=z
	\]
	对 \((t,z)\in(-a,a)\times B\) 有唯一局部解，且 \(\Phi\) 对时间与初值联合光滑。若一条解已经存在于紧时间区间，则该区间上的解也对其附近的初值光滑依赖。
\end{lemma}
\begin{proof}
	\cref{g1:lem:surface-field-integral-curve}用积分迭代证明了二维向量场的局部存在唯一性。该证明只使用 Euclid 范数与导数界，在 \(\mathbb R^m\) 中同样成立。具体地，取 \(\overline B(z_0,2r)\subset\Omega\)，在此球上以 \(M,L\) 分别控制 \(|F|,\|\mathrm{D}F\|\)。令初值 \(z\in B(z_0,r)\)，并取 \(aM<r\)、\(q=aL<1\)，则
	\[
		(\mathcal T_zx)(t)=z+\int_0^tF(x(s))\,\mathrm{d}s
	\]
	把取值于大闭球的连续路径映入自身，且一致距离至多缩小为 \(q\) 倍。相邻迭代的差构成收敛几何级数，极限就是唯一解。比较两个初值得到
	\[
		\|\Phi(\cdot,z+h)-\Phi(\cdot,z)\|_\infty\le\frac{|h|}{1-q}.
	\]
	这也先证明了初值依赖的连续性。

	为证明可微，令矩阵 \(A(t,z)\) 解积分方程
	\[
		A(t,z)=I+\int_0^t \mathrm{D}F(\Phi(s,z))A(s,z)\,\mathrm{d}s.
	\]
	同一个缩短估计给出其唯一连续解，并由比较积分方程得到它对 \(z\) 连续。记 \(\delta_h=\Phi(\cdot,z+h)-\Phi(\cdot,z)\)。\(\mathrm{D}F\) 在较小闭球上一致连续，Taylor 余项与上面的 Lipschitz 估计遂给出
	\[
		\|\delta_h-A(\cdot,z)h\|_\infty
		\le q\|\delta_h-A(\cdot,z)h\|_\infty+o(|h|).
	\]
	把第一项移到左边，就有 \(\mathrm{D}_z\Phi=A\)。高阶导数可以用相同办法逐阶得到：对初值求 \(k\) 阶导后，未知的最高阶导数 \(B_k\) 总满足
	\[
		B_k(t)=B_k(0)+\int_0^t\bigl(\mathrm{D}F(\Phi(s))B_k(s)+C_k(s)\bigr)\,\mathrm{d}s,
	\]
	其中 \(C_k\) 只由 \(F\) 的至多 \(k\) 阶导数和已经求得的低阶导数组成，\(k\ge2\) 时 \(B_k(0)=0\)。归纳时，对前一阶方程取差商，Taylor 余项趋零，再用 \(1-q>0\) 的同一估计，便证明该方程的解确为下一阶导数而且连续。时间导数则由 \(\partial_t\Phi=F(\Phi)\) 得到，混合导数也随之连续，故 \(\Phi\) 联合光滑。

	最后，已有解在紧时间区间上的像是紧集，可沿它选取有限个上述局部时间步。逐步缩小初值邻域，使各步的终点仍落在下一步的允许邻域；有限次光滑映射的复合给出整个区间上的光滑依赖，唯一性保证相邻步骤相接。
\end{proof}

\begin{proposition}[测地线初值问题]
	\label{g1:prop:geodesic-initial-values}
	设 \(M\) 是光滑度量曲面，\(p\) 为内点，\(v\in \mathrm{T}_pM\)。则存在以 \(0\) 为内点的区间及其上的唯一局部测地线 \(\gamma_v\)，使 \(\gamma_v(0)=p\)、\(\dot\gamma_v(0)=v\)。在共同存在的紧时间区间上，解对起点与初速度光滑依赖。若 \(v\ne0\)，其速度恒为 \(|v|\)；若 \(v=0\)，则为常值曲线。
\end{proposition}
\begin{proof}
	置 \(w^k=\dot u^k\)，将\eqref{g1:eq:geodesic-equation}改写成四维一阶系统
	\[
		\dot u^k=w^k,\qquad \dot w^k=-\Gamma^k_{ij}(u)w^iw^j.
	\]
	右端光滑，前一引理给出全部初值结论。换坐标后所得曲线仍满足几何定义 \(\mathrm{D}_t\dot\gamma=0\)，所以局部唯一性使各坐标中的解相符。速度性质已由度量相容性证明；零初速度时，常值曲线是一个解，唯一性排除了其他解。
\end{proof}

初值问题给定的是“一点及一个速度”，并不保证任意两点间有唯一测地线，也不保证解能够走任意长时间。例如去掉原点的 Euclid 平面中，指向原点的直线在有限时间撞上被删去的点；而球面的南北两极之间有许多不同的半大圆。

\begin{example}[球面的测地线]
	\label{g1:exm:sphere-geodesics}
	半径 \(R\) 的球面取外向单位法向 \(n=\gamma/R\)。若 \(\gamma\) 是速度 \(c>0\) 的测地线，则空间加速度平行于 \(\gamma\)。两次求导 \(\langle\gamma,\gamma\rangle=R^2\)，有 \(\langle\ddot\gamma,\gamma\rangle=-c^2\)，所以
	\[
		\ddot\gamma=-\frac{c^2}{R^2}\gamma,\qquad
		\gamma(t)=p\cos\frac{ct}{R}+\frac{R}{c}v\sin\frac{ct}{R},
	\]
	其中 \(|p|=R\)、\(\langle p,v\rangle=0\)、\(|v|=c\)。这个公式也可直接核验初值与方程，唯一性说明它给出全部非常值测地线。因此球面的测地线轨迹是大圆；绕一整圈仍满足方程，却显然不是起终点之间的最短路。
\end{example}

\subsection{弧长、能量与端点条件}
\label{g1:subsec:0902:02}

比较曲线时，长度天然不在乎走得快慢，微分方程却记录参数。引入能量可以把参数也纳入比较，使随后的求导不出现速度范数的分母。对于固定区间 \([a,b]\)、\(a<b\) 上的分片 \(C^1\) 曲线，定义
\[
	L(\gamma)=\int_a^b|\dot\gamma(t)|\,\mathrm{d}t,\qquad
	E(\gamma)=\frac12\int_a^b|\dot\gamma(t)|^2\,\mathrm{d}t.
\]
这里的 \(E\) 称为曲线的\term[nengliang]{能量}[energy]。\symidx{\(E(\gamma)\)：曲线能量}长度是各坐标段中的度量长度；对于嵌入曲面，\cref{g1:prop:arc-length-param}的长度公式说明它与空间曲线的折线定义一致。

\begin{proposition}[长度与能量]
	\label{g1:prop:length-energy}
	设 \((M,g)\) 为光滑度量曲面，\(a<b\)，\(\gamma:[a,b]\to M\) 为分片 \(C^1\) 曲线。令 \(L(\gamma)=\int_a^b|\dot\gamma|_g\,\mathrm dt\)、\(E(\gamma)=\tfrac12\int_a^b|\dot\gamma|_g^2\,\mathrm dt\)，则
	\[
		L(\gamma)^2\le2(b-a)E(\gamma).
	\]
	等号成立当且仅当速度范数几乎处处为常数；对 \(C^1\) 曲线，这等价于处处等速。若一条曲线在给定端点的全部竞争曲线中长度最小且等速，则它在同一区间上也使能量最小。
\end{proposition}
\begin{proof}
	对函数 \(|\dot\gamma|\) 与常函数 \(1\) 使用 Cauchy--Schwarz 不等式，便得到所述估计。等号条件是两个函数几乎处处成比例。若 \(\gamma_0\) 等速且最短，对任意同端点的竞争曲线 \(\eta\)，则
	\[
		E(\eta)\ge\frac{L(\eta)^2}{2(b-a)}\ge\frac{L(\gamma_0)^2}{2(b-a)}=E(\gamma_0).
	\]
\end{proof}

对正则曲线，也可把其轨迹改为等速行进而不改长度。具体地，令 \(\ell(t)=\int_a^t|\dot\gamma(\tau)|\,\mathrm{d}\tau\)，它严格递增，按\cref{g1:prop:arc-length-param}中正速度积分取逆的构造，再把 \([0,L]\) 线性缩放到 \([a,b]\)，便得到速度 \(L/(b-a)\) 的参数。该构造只使用正的速度函数，对这里的度量速度同样成立。因此在正则曲线的比较中，能量最小同时要求轨迹最短和参数等速；长度最小则不强制后一条件。

\ProseParagraphBreak

为了精确表达“小幅移动一条曲线”，取光滑映射
\[
	F:(-\varepsilon,\varepsilon)\times[a,b]\to M,\qquad F(0,t)=\gamma(t).
\]
称 \(F\) 为 \(\gamma\) 的一个\term[bianfen]{变分}[variation]，称 \(V(t)=\partial_sF(0,t)\) 为变分向量场。固定端点是要求 \(F(s,a)=p\)、\(F(s,b)=q\)，它给出 \(V(a)=V(b)=0\)。若只要求端点分别落在给定的光滑曲线 \(C_0,C_1\) 上，则 \(V(a)\in \mathrm{T}_pC_0\)、\(V(b)\in \mathrm{T}_qC_1\)，这些向量通常不为零。若端点可以在曲面内任意移动，相应的端点向量更没有限制。这三种条件不能混用，下一小节中的边界项会记录它们的差别。

\subsection{测地曲率及第一变分}
\label{g1:subsec:0902:03}

空间曲率把曲线所有弯曲都算在内。对于曲面内部的最短路，真正应受约束的是曲线在切平面内向侧方偏转的程度。以下先取一个曲面定向，记 \(J\) 为切平面内的正向直角旋转；嵌入曲面中，\(JW=n\times W\)。

\begin{definition}
	\label{g1:def:geodesic-curvature}
	设 \((M,g)\) 为定向光滑度量曲面，\(\mathrm D\) 为其无挠相容联络沿曲线的导数，\(J\) 为各切平面的正向直角旋转。设 \(\gamma\) 是 \(M\) 上的 \(C^2\) 正则曲线，\(s\) 为沿行进方向的弧长，\(T=\mathrm{d}\gamma/\mathrm{d}s\)。定义有向\term[cediqulv]{测地曲率}[geodesic curvature]为
	\[
		k_g=\langle \mathrm{D}_sT,JT\rangle.
	\]
	由于 \(|T|=1\)，度量相容性给出 \(\mathrm{D}_sT\perp T\)，故 \(\mathrm{D}_sT=k_gJT\)。\symidx{\(k_g\)：有向测地曲率}
\end{definition}

于是正则曲线的弧长参数化为测地线，当且仅当 \(k_g=0\)。反转曲面定向会使 \(k_g\) 变号；保持曲面定向而反转曲线行进方向也使 \(k_g\) 变号。若 \(\gamma\) 是嵌入曲面上的单位速曲线，\cref{g1:prop:gauss-weingarten}还给出空间加速度的分解
\begin{equation}
	\label{g1:eq:geodesic-normal-curvature}
	\gamma''=k_gJT+\mathrm{II}(T,T)n,
	\qquad \kappa^2=k_g^2+\mathrm{II}(T,T)^2.
\end{equation}
测地线只消去这里的切向项，通常仍有法曲率。例如球面大圆的测地曲率为零，空间曲率却是 \(1/R\)。

\begin{proposition}[第一变分]
	\label{g1:prop:first-variation}
	设 \((M,g)\) 为光滑度量曲面，\(a<b\)、\(\varepsilon>0\)，\(F:(-\varepsilon,\varepsilon)\times[a,b]\to M\) 是光滑变分，\(\gamma(t)=F(0,t)\)、\(V(t)=\partial_sF(0,t)\)。令 \(E(c)=\tfrac12\int_a^b|\dot c|_g^2\,\mathrm dt\)、\(L(c)=\int_a^b|\dot c|_g\,\mathrm dt\)，\(\mathrm D\) 取 \(g\) 的无挠相容联络。则
	\begin{equation}
		\label{g1:eq:energy-first-variation}
		\left.\frac{\mathrm{d}}{\mathrm{d}s}E(F(s,\cdot))\right|_{s=0}
		=\bigl[\langle V,\dot\gamma\rangle\bigr]_a^b-\int_a^b\langle V,\mathrm{D}_t\dot\gamma\rangle\,\mathrm{d}t.
	\end{equation}
	若 \(\gamma\) 正则，令 \(T=\dot\gamma/|\dot\gamma|\)，则
	\begin{equation}
		\label{g1:eq:length-first-variation}
		\left.\frac{\mathrm{d}}{\mathrm{d}s}L(F(s,\cdot))\right|_{s=0}
		=\bigl[\langle V,T\rangle\bigr]_a^b-\int_a^b\langle V,\mathrm{D}_tT\rangle\,\mathrm{d}t.
	\end{equation}
	若 \(M\) 已定向，以 \(J\) 记正向直角旋转，\(k_g=\langle\mathrm D_\ell T,JT\rangle\) 记测地曲率，把最后一项改用弧长 \(\ell\) 积分，即为 \(-\int_\gamma k_g\langle V,JT\rangle\,\mathrm{d}\ell\)。
\end{proposition}
\begin{proof}
	记 \(W=\partial_tF\)。由\cref{g1:prop:along-derivative-properties}，度量相容性及混合协变导数相等分别给出
	\[
		\partial_s\frac12\langle W,W\rangle=\langle \mathrm{D}_sW,W\rangle,
		\qquad \mathrm{D}_s\partial_tF=\mathrm{D}_t\partial_sF.
	\]
	后一个等式的坐标理由是 \(\partial_s\partial_tu^k=\partial_t\partial_su^k\) 及 \(\Gamma^k_{ij}=\Gamma^k_{ji}\)，并不要求曲率为零。对紧区间积分可在积分号内求导，故能量导数是 \(\int_a^b\langle \mathrm{D}_tV,\dot\gamma\rangle\,\mathrm{d}t\)。再由内积的求导公式作分部积分，得到\eqref{g1:eq:energy-first-variation}。

	若 \(\gamma\) 正则，其速度在紧区间上有正的下界，缩小变分参数区间后附近曲线也正则。因此可对速度范数直接求导，得到长度导数
	\[
		\int_a^b\left\langle \mathrm{D}_tV,\frac{\dot\gamma}{|\dot\gamma|}\right\rangle\,\mathrm{d}t.
	\]
	再次分部积分即得\eqref{g1:eq:length-first-variation}。最后 \(\mathrm{D}_tT=|\dot\gamma|\mathrm{D}_\ell T\)，而 \(\mathrm{d}\ell=|\dot\gamma|\mathrm{d}t\)，代入测地曲率定义便得弧长形式。
\end{proof}

所谓驻定，是指对所有允许变分的一阶导数均为零，不是指附近所有曲线都更长。固定端点时，\eqref{g1:eq:energy-first-variation}立即说明测地线使能量驻定。反过来，如果能量对所有固定端点变分驻定，则 \(\mathrm{D}_t\dot\gamma=0\)：若某个内点处该向量非零，取其附近一段落在同一坐标图中的小区间，并选非负光滑函数 \(\chi\)，支撑在此小区间内且在该点为正。令 \(V=\chi \mathrm{D}_t\dot\gamma\)，在坐标中用 \(u_s(t)=u(t)+sV^i(t)\) 实现这个变分，并在区间外保持原曲线，则能量导数为
\[
	-\int_a^b\chi|\mathrm{D}_t\dot\gamma|^2\,\mathrm{d}t<0,
\]
产生矛盾。所需的 \(\chi\) 可由区间内部的函数 \(\exp(-1/(\delta^2-(t-t_0)^2))\) 在区间外补零得到。对长度公式作相同论证，正则曲线的固定端点长度驻定等价于 \(\mathrm{D}_tT=0\)，即轨迹按弧长走时为测地线；它本身未必等速。

\ProseParagraphBreak

若端点分别可在 \(C_0,C_1\) 上移动，先用支撑在内部的变分仍得到测地线条件，再保留边界项：驻定还要求终点速度分别垂直于 \(\mathrm{T}_pC_0,\mathrm{T}_qC_1\)。任意端点切向量都可实现：先在相应端点曲线上取具有该速度的短曲线，再在端点附近的坐标图中用光滑截断函数把这一移动延伸到变分，远处保持不动。因而上述垂直条件也是必要的；反过来，测地线条件和垂直条件使全部项消失，所以它们充分。如果一个端点可在曲面内任意移动，长度的边界项便不可能对所有向量都为零，因为 \(T\) 是非零切向量。这说明最短路问题必须先说清楚端点允许怎样动。

\subsection{测地坐标和径向局部极小性}
\label{g1:subsec:0902:04}

第一变分只给出了极小的必要条件；要证明短测地线确实最短，还需要一种能同时比较许多竞争曲线的坐标。直觉是从同一点向各个方向发出测地线，用初速度标记它们的终点。这一标记是否会重合，必须先在局部解决。

\begin{definition}
	\label{g1:def:surface-exponential}
	设 \((M,g)\) 为光滑度量曲面，\(p\) 为其内点。对 \(v\in\mathrm T_pM\)，以 \(\gamma_v\) 记满足 \(\gamma_v(0)=p\)、\(\dot\gamma_v(0)=v\) 的测地线。对使 \(\gamma_v\) 存在于 \([0,1]\) 的 \(v\in \mathrm{T}_pM\)，定义\term[zhishuyingshe]{指数映射}[exponential map]
	\[
		\exp_p(v)=\gamma_v(1).
	\]
	这个名称表示由初速度生成测地线的映射，并不是坐标逐项取通常的指数函数。\symidx{\(\exp_p\)：曲面的指数映射}
\end{definition}

\begin{proposition}[正常邻域]
	\label{g1:prop:normal-neighborhood}
	设 \(M\) 为光滑度量曲面，\(p\) 为内点。则 \(\exp_p\) 在 \(\mathrm{T}_pM\) 中零向量的某个开邻域上光滑，并满足 \(\mathrm{d}(\exp_p)_0=\operatorname{id}_{\mathrm{T}_pM}\)。存在 \(\rho>0\)，使它在 \(B_{2\rho}(0)\subset \mathrm{T}_pM\) 上为到开邻域的光滑微分同胚，其中球由 \(p\) 处的内积定义。
\end{proposition}
\begin{proof}
	零初速度给出常值测地线，在任意紧时间区间上都存在。由\cref{g1:lem:smooth-ode-flow}的有限时间结论，附近初速度的解也存在于包含 \([0,1]\) 的时间区间，并且光滑依赖初值，故 \(\exp_p\) 在零附近定义且光滑。

	初值唯一性给出缩放关系 \(\gamma_{tv}(s)=\gamma_v(ts)\)，只要两边有定义。因此对于固定 \(v\) 及足够小的 \(t\)，有 \(\exp_p(tv)=\gamma_v(t)\)，求 \(t=0\) 处的导数便得 \(\mathrm{d}(\exp_p)_0(v)=v\)。在 \(p\) 附近取一张坐标图后，\cref{g1:lem:surface-local-inverse}适用于这个二维光滑映射，给出零附近的光滑逆。最后在允许的邻域中选取足够小的 \(B_{2\rho}(0)\)。
\end{proof}

把 \(\mathrm{T}_pM\) 的一组正交单位基写作 \(e_1,e_2\)，用
\[
	(x^1,x^2)\longmapsto\exp_p(x^1e_1+x^2e_2)
\]
得到的坐标称为\term[cedizuobiao]{测地坐标}[geodesic normal coordinates]，其像称为\term[zhengchanglinyu]{正常邻域}[normal neighborhood]。这些坐标的优点是在中心点恢复 Euclid 度量的一阶形态。确实，\(\mathrm{d}(\exp_p)_0=\operatorname{id}\) 给出 \(g_{ij}(p)=\delta_{ij}\)；直线坐标 \(x(t)=tv\) 对应测地线，代入方程得 \(\Gamma^k_{ij}(p)v^iv^j=0\)。对所有 \(v\) 成立，加上 \(\Gamma^k_{ij}=\Gamma^k_{ji}\)，推出 \(\Gamma^k_{ij}(p)=0\)。度量相容的坐标式又给出 \(\partial_k g_{ij}(p)=0\)。这些消失只发生在中心点，不能据此认定整个邻域平坦。

\begin{lemma}[Gauss 引理]
	\label{g1:lem:gauss-lemma}
	设 \(p\) 是光滑度量曲面的内点，\(\exp_p\) 在零点的星形开邻域 \(B\subset \mathrm{T}_pM\) 上光滑且其中所有径向测地线均定义至时间 \(1\)。则对 \(v\in B\) 及 \(w\in \mathrm{T}_pM\)，有
	\[
		\left\langle\mathrm{d}(\exp_p)_v(v),\mathrm{d}(\exp_p)_v(w)\right\rangle_{\exp_p(v)}
		=\langle v,w\rangle_p.
	\]
	特别地，指数映射保持径向长度，并把与径向垂直的向量映为与径向测地线垂直的向量。
\end{lemma}
\begin{proof}
	\(v=0\) 时两端均为零。其余情况取小 \(s\)，令 \(F(t,s)=\exp_p(t(v+sw))\)、\(0\le t\le1\)。星形条件保证这些线段留在定义域内。记 \(A=\partial_tF\)、\(B_1=\partial_sF\)。每条 \(t\) 曲线为测地线，故 \(\mathrm{D}_tA=0\)，且速度平方为 \(|v+sw|_p^2\)。利用混合协变导数相等，得到
	\[
		\partial_t\langle A,B_1\rangle
		=\langle A,\mathrm{D}_sA\rangle
		=\frac12\partial_s|A|^2
		=\langle v+sw,w\rangle_p.
	\]
	在 \(t=0\) 处，\(F(0,s)=p\)，所以 \(B_1(0,s)=0\)。积分到 \(t=1\) 并取 \(s=0\)，左边正是所求的内积。令 \(w=v\) 给出径向长度保持，令 \(w\perp v\) 给出径向正交性。
\end{proof}

这一结论称为\term[gaussyinli]{Gauss 引理}[Gauss lemma]。它没有说指数映射保持所有长度。事实上，取 \(e(\theta)=e_1\cos\theta+e_2\sin\theta\)，在 \(r>0\) 处使用极坐标
\[
	(r,\theta)\longmapsto\exp_p(re(\theta)),
\]
Gauss 引理只保证径向速度长度为 \(1\)、径向与角向正交。度量因而具有形式
\begin{equation}
	\label{g1:eq:geodesic-polar-metric}
	I=\mathrm{d}r^2+a(r,\theta)^2\,\mathrm{d}\theta^2,\qquad a(r,\theta)>0.
\end{equation}
其中 \(a(r,\theta)=|\mathrm{d}(\exp_p)_{re(\theta)}(r e'(\theta))|\)。光滑性给出 \(a/r\to1\)，但一般 \(a\ne r\)。例如球面中心为 \(p\) 时，\cref{g1:exm:sphere-geodesics}给出 \(a=R\sin(r/R)\)，在 \(0<r<\pi R\) 内有定义，到了对跖点则所有径向方向相遇。极坐标在 \(r=0\) 的角度没有定义，光滑的测地直角坐标却没有这个缺陷。

\input{Figures/Chapter09-NormalNeighborhood.tex}

\begin{theorem}[径向测地线的最短性]
	\label{g1:thm:radial-minimizing}
	设 \((M,g)\) 是光滑度量曲面，\(p\) 为内点，并取 \(\rho>0\)，使 \(\exp_p\) 在 \(B_{2\rho}(0)\) 上为微分同胚。令 \(V=\exp_p(B_\rho(0))\)。对 \(q=\exp_p(v)\in V\)，曲线 \(t\mapsto\exp_p(tv)\)、\(0\le t\le1\)，在 \(M\) 内所有连接 \(p,q\) 且参数化在 \([0,1]\) 上的分片 \(C^1\) 曲线中长度最小，长度为 \(|v|\)。当 \(q\ne p\) 时，达到该长度的正则分片光滑曲线具有同一径向轨迹，并沿同一方向行进；在共同区间 \([0,1]\) 上以恒速参数化后，最短参数曲线唯一。
\end{theorem}
\begin{proof}
	记 \(r(x)=|\exp_p^{-1}(x)|\)。先比较留在较大正常邻域中的曲线。Gauss 引理把非中心点处任意速度 \(W\) 正交分解为径向部分与角向部分，因而
	\[
		|W|^2=\bigl(\mathrm{d}r(W)\bigr)^2+|W_{\mathrm{ang}}|^2
		\ge\bigl(\mathrm{d}r(W)\bigr)^2.
	\]
	这也可以直接从\eqref{g1:eq:geodesic-polar-metric}读出。设竞争曲线 \(\eta:[a,b]\to V\) 从 \(p\) 到 \(q\ne p\)。为避开中心处的极坐标问题，对任意 \(0<\varepsilon<r(q)\)，取最后一个满足 \(r(\eta(t_\varepsilon))=\varepsilon\) 的时刻。连续性及终点条件保证它存在，且随后 \(r>\varepsilon\)。在这一段上逐段积分，得到
	\[
		L(\eta)\ge\int_{t_\varepsilon}^b|(r\circ\eta)'(t)|\,\mathrm{d}t
		\ge r(q)-\varepsilon.
	\]
	令 \(\varepsilon\downarrow0\)，便有 \(L(\eta)\ge r(q)\)。径向曲线的速度恒为 \(|v|\)，所以恰好达到这个界。

	现在允许竞争曲线离开 \(V\)。由于 \(\exp_p\) 在更大的球 \(B_{2\rho}\) 上是微分同胚，\(\overline V=\exp_p(\overline B_\rho)\) 为紧集，其边界就是 \(\exp_p(\partial B_\rho)\)。任何从 \(p\) 出发又离开 \(V\) 的连续曲线，必有一个首先到达该边界的时刻；此前的弧段仍处于正常坐标内。刚才的估计对这一弧段给出长度至少为 \(\rho\)，而 \(r(q)<\rho\)。因此离开 \(V\) 的曲线不可能达到径向曲线的长度。

	最后说明等号。先注意，位于 \(V\) 的任意弧段的长度至少为两端的径向坐标差的绝对值：不经过 \(p\) 时直接积分 \(|(r\circ\eta)'|\)；经过 \(p\) 时，取第一次及最后一次到达 \(p\) 的时刻，分别把两端弧段按上述从中心出发的估计处理，便得到更强的下界 \(r(\eta(a))+r(\eta(b))\)。因此，若总长度恰为 \(r(q)\)，在曲线不经过 \(p\) 的任意一小段上，径向增量的积分不能严格小于长度：否则把这段的严格差值与前后两段的径向估计相加，会使总长度严格大于 \(r(q)\)。于是该段必须同时满足 \(W_{\mathrm{ang}}=0\) 及 \((r\circ\eta)'\ge0\)。正则性使离开 \(p\) 后的径向速度为正，角方向保持常数；它也不可能返回 \(p\)，否则已经走过的正长度不能被后续消去。因此轨迹为从 \(p\) 沿同一方向到 \(q\) 的径向段。其恒速参数由端点、参数区间和长度唯一确定。
\end{proof}

最短性是以初始点为中心、在足够小范围内成立的结论。正常邻域中的任意两点未必能直接用这一中心点的径向线连接；要讨论它们之间的唯一最短线，须重新以其中一点作上述构造，并核对另一点落在允许范围内。球面大圆可以继续延伸，超过半圈后却不再最短。初值方程控制测地线的延伸，Gauss 引理则给出短段的长度比较；延伸范围与最短范围由此成为两个问题。

% !TeX root = ../../Book1.tex
% 纲要标题：### 9.3 局部 Gauss–Bonnet 公式
% 纲要位置：几何学计划纲要.md，第 479 行。
% 正文状态：初稿；尚未正式审核。

% 纲要内容开始
% 
% * 切向转角与联络积分
% * 分片光滑边界及外角
% * 圆盘区域上的公式与定向符号
% * 测地三角形的角亏
% 
% 纲要内容结束

\section{局部 Gauss–Bonnet 公式}
\label{g1:sec:0903}

平面上的正向简单闭曲线，其切向量净转过一周。曲面上的切向量却分属不同的切平面，必须先说明用什么方向作参照。联络给出比较方向的方法，测地曲率记录相对于平行移动的转向；把它沿边界积分，通常不再得到 \(2\pi\)。本节将说明，恰好缺少的量是区域内部的 Gauss 曲率积分。边界若有角点，还须计入转弯时发生的方向跳跃。

\subsection{切向转角与联络积分}
\label{g1:subsec:0903:01}

设曲面已经定向，\(J\) 表示切平面内正向旋转 \(\pi/2\)。取正向正交单位标架 \((e_1,e_2)\)，所以 \(Je_1=e_2\)、\(Je_2=-e_1\)。回忆\eqref{g1:eq:surface-frame-connection}中的约定：联络一形式 \(\omega\) 满足
\[
	\nabla_Ve_1=\omega(V)e_2,\qquad \nabla_Ve_2=-\omega(V)e_1.
\]
沿一段按弧长参数化的 \(C^2\) 正则曲线 \(\gamma\)，单位切向量 \(T=\gamma'\) 可写成
\[
	T(s)=\cos\theta(s)e_1+\sin\theta(s)e_2.
\]
这里标架在 \(\gamma(s)\) 处取值。\cref{g1:lem:circle-angle-lift}允许在参数区间上连续选择角度；因切向量是 \(C^1\) 的，局部光滑辐角使这个提升也是 \(C^1\) 的。角度不是曲面本身提供的绝对方向，而是切向量相对于所选标架的方向。

\begin{proposition}[切向转角公式]
	\label{g1:prop:tangent-connection-angle}
	设 \(\gamma\) 是有向光滑曲面上的 \(C^2\) 弧长参数正则曲线，\((e_1,e_2)\) 是其邻域中的光滑正向正交单位标架，\(\omega=\langle\nabla e_1,e_2\rangle\)。若 \(T=\cos\theta\,e_1+\sin\theta\,e_2\)，则
	\begin{equation}
		\label{g1:eq:tangent-connection-angle}
		k_g=\theta'+\omega(T),\qquad
		\int_\gamma k_g\,\mathrm{d}s=\theta(b)-\theta(a)+\int_\gamma\omega.
	\end{equation}
\end{proposition}
\begin{proof}
	按\cref{g1:prop:along-derivative-properties}的乘法法则沿曲线求协变导数，得到
	\begin{align*}
		\mathrm{D}_sT
		&=-\theta'\sin\theta\,e_1+\theta'\cos\theta\,e_2
		+\cos\theta\,\omega(T)e_2-\sin\theta\,\omega(T)e_1\\
		&=\bigl(\theta'+\omega(T)\bigr)JT.
	\end{align*}
	\cref{g1:def:geodesic-curvature}规定 \(k_g=\langle \mathrm{D}_sT,JT\rangle\)，故得第一式。对 \(s\) 积分，并用一形式的线积分定义 \(\int_\gamma\omega=\int_a^b\omega(T)\,\mathrm{d}s\)，便得第二式。
\end{proof}

这个等式把两种转动分开：\(\theta'\) 是切向量相对于标架的转动，\(\omega(T)\) 是标架自身相对于平行移动的转动。若沿曲线选择平行标架，第二项为零；若选择曲面上一组便于计算的坐标标架，第二项一般不能省略。尤其在曲线闭合时，沿它平行移动所得的标架未必首尾一致，因此不能同时把它当作周期标架，又任意令 \(\omega(T)=0\)。

\ProseParagraphBreak

换一组正向正交标架不会改变 \(k_g\)。例如局部令 \(\widetilde e_1=\cos\varphi\,e_1+\sin\varphi\,e_2\)、\(\widetilde e_2=J\widetilde e_1\)，则新角度是 \(\widetilde\theta=\theta-\varphi\)，新联络形式是 \(\widetilde\omega=\omega+\mathrm{d}\varphi\)。因而 \(\widetilde\theta'+\widetilde\omega(T)=\theta'+\omega(T)\)。计算中的角度和联络积分各自依赖参照，二者之和才是曲线的测地曲率积分。

\subsection{分片光滑边界及外角}
\label{g1:subsec:0903:02}

三角形的边可以是光滑的，但顶点处没有唯一切向量。若只把三条直线边的曲率积分相加，结果是零，整周转角便丢失了。因此要把沿边的连续转向和顶点处的瞬时转向一起记录。

\begin{definition}
	\label{g1:def:boundary-exterior-angle}
	设紧曲面区域 \(D\) 的边界由有限条 \(C^2\) 正则弧首尾相接而成，各弧在端点有非零单侧切向量，弧的内部互不相交。按区域位于左侧的方向走过边界，在角点 \(p_j\) 记到达和离开的单位切向量为 \(T_j^-\)、\(T_j^+\)。若区域在该点占据的切向扇形角为 \(\beta_j\in(0,2\pi)\)，称 \(\beta_j\) 为\term[neijiao]{内角}[interior angle]，并称
	\[
		\alpha_j\coloneqq\pi-\beta_j\in(-\pi,\pi)
	\]
	为该点的有向\term[waijiao]{外角}[exterior angle]。它是从 \(T_j^-\) 转到 \(T_j^+\) 的有向短转角；向左转取正，向右转取负。允许 \(\beta_j=\pi\)，此时只是人为分段点，外角为零。
\end{definition}

内角是由指向来路的 \(-T_j^-\) 与指向去路的 \(T_j^+\) 围成、朝向区域内部的角，不是两个前进切向量之间较小的夹角。\cref{g1:fig:boundary-exterior-angles}同时画出凸角与凹角：凹角的内角大于 \(\pi\)，外角便是负数。这里排除了内角为 \(0\) 或 \(2\pi\) 的尖点；在这种退化情形，两个前进切向量可能反向，不能直接使用唯一短转角。

\begin{figure}[htbp]
	\centering
	\input{Figures/Chapter09-ExteriorAngles}
	\caption[边界外角的符号]{区域位于行进方向左侧。实线箭头表示先到达顶点、再沿下一边离开；虚线是到达方向的延长线。左图的内角小于 \(\pi\)，外角为正；右图的内角大于 \(\pi\)，外角为负。浅色扇形只示意顶点附近的内部方向。}
	\label{g1:fig:boundary-exterior-angles}
\end{figure}

将边界各段的切向角连续提升，并在每个角点接上转角 \(\alpha_j\)，便得到包含角点跳跃的总转角。这个总量应是多少，不能只从曲线闭合推出：闭合只使它成为 \(2\pi\) 的整数倍。下面说明，单图圆盘的正向简单边界使这个整数恰好为 \(1\)。

\begin{lemma}[分片边界的转角]
	\label{g1:lem:piecewise-turning}
	设有向光滑曲面上的紧区域 \(D\) 同胚于闭圆盘，并且包含在一张正向坐标图中。假设 \(\partial D\) 是正向分片 \(C^2\) 正则简单闭曲线，各角点内角均属于 \((0,2\pi)\)。将该图的坐标向量按顺序正交单位化，得到正向标架 \((e_1,e_2)\)。若 \(\Delta\theta_i\) 是第 \(i\) 条边的切向量相对于该标架的角增量，\(\alpha_j\) 为各外角，则
	\begin{equation}
		\label{g1:eq:piecewise-total-turning}
		\sum_i\Delta\theta_i+\sum_j\alpha_j=2\pi.
	\end{equation}
\end{lemma}
\begin{proof}
	先在坐标平面中证明带角点的转角定理。\cref{g1:thm:turning-tangent}已经证明 \(C^1\) 正则简单闭平面曲线的正向总转角为 \(2\pi\)；现在只需把角点的贡献补回来。在某个角点，两个前进单位切向量不反向，取它们的角平分方向为横轴，则两个向量在该横轴上的分量都为正。于是足够短的两侧曲线合在一起可以写成连续图像 \(y=f(x)\)，其左右单侧导数存在，且 \(f\) 在两侧为 \(C^2\)。取互不相交的小邻域处理全部角点，并使这些邻域避开其余曲线。

	在一个这样的邻域中，把 \(f\) 在角点附近换成光滑近似，并用光滑截断函数使它在邻域边缘与原函数完全相同。这个操作可以具体实施为：先在稍大的区间上延拓 \(f\)，用非负光滑积分核作卷积，再在角点附近取卷积结果、远离角点取原函数，在两者之间以截断函数相接。卷积在小尺度下一致趋于连续函数 \(f\)，所以改动后的图像仍留在所选邻域。它始终是关于 \(x\) 的图像，故不产生自交，也不会使前进切向量为零。所得闭曲线为 \(C^2\) 正则简单曲线，并沿用原曲线的正向。

	图像段的切向角可取为 \(\arctan f'(x)\)，值始终在 \((-\pi/2,\pi/2)\) 内；因此，无论中间怎样光滑化，这一小段的总转角恰为其终端与初端的切向角之差。让修改区间缩向原角点，两端角度趋向单侧切向角，差趋向该角点的有向短转角。未修改各段的角增量则趋向原边的角增量。将光滑曲线的转角定理取极限，得到坐标平面中“沿边角增量加外角等于 \(2\pi\)”。内角为 \(\pi\) 的分段点不改变这一论证。

	还需比较坐标平面的角和曲面度量测得的角。记坐标为 \((u,v)\)，度量系数为 \(E,F,G\)，\(D_g=EG-F^2>0\)。正交单位化给出
	\[
		\partial_u=\sqrt E\,e_1,\qquad
		\partial_v=\frac{F}{\sqrt E}\,e_1+\sqrt{\frac{D_g}{E}}\,e_2.
	\]
	因此，坐标速度 \(w=(\dot u,\dot v)^{\mathsf T}\) 在正交标架中的分量是
	\[
		Pw,\qquad
		P=\begin{pmatrix}\sqrt E&F/\sqrt E\\0&\sqrt{D_g/E}\end{pmatrix}.
	\]
	对 \(0\leq t\leq1\)，矩阵 \(P_t=(1-t)I+tP\) 都是对角元为正的上三角矩阵，故可逆且保向。沿每条边将 \(P_tw\) 归一化，在每个角点将左右极限用有向短圆弧接起来，就得到方向圆周上的连续闭路。这个闭路随 \(t\) 连续变化：非反向的两个端点经过可逆线性变换后仍不反向，故短圆弧不会突然跳到另一边。可给每个角点插入一段固定长度的辅助参数区间，使这里的闭路具有同一个参数圆周。

	\cref{g1:thm:circle-fundamental}的绕数同伦不变性说明，这个方向闭路在 \(t=0\) 与 \(t=1\) 时有相同总角增量。\(t=0\) 是刚证明的平面分片边界，总增量为 \(2\pi\)；\(t=1\) 正是曲面切向量相对于 \((e_1,e_2)\) 的方向及度量外角。这就得到\eqref{g1:eq:piecewise-total-turning}。
\end{proof}

平面角和曲面角一般不相等；上面保持的是绕一周的整数，而不是每个顶点的数值。这也解释了为何不能直接把坐标图画出的三角形内角，当成曲面上的内角代入公式。

\subsection{圆盘区域上的公式与定向符号}
\label{g1:subsec:0903:03}

现在已有两个可以相接的等式：切向转角公式把边界的测地曲率写成角增量与联络积分之和，曲率二形式则把联络的外微分写成 \(-K\,\mathrm{d}A\)。正号、负号都由这两个约定决定，不另作猜测。

\begin{theorem}[局部 Gauss–Bonnet 公式]
	\label{g1:thm:local-gauss-bonnet}
	设 \(M\) 是带光滑正定度量的有向光滑曲面，\(D\subset M\) 是包含在一张正向坐标图中的紧闭圆盘区域。假设其边界是有限条 \(C^2\) 正则弧组成的简单闭曲线，所有内角 \(\beta_j\) 均在 \((0,2\pi)\) 内。边界按 \(D\) 位于左侧定向，\(k_g=\langle \mathrm{D}_sT,JT\rangle\)，\(\alpha_j=\pi-\beta_j\)，则
	\begin{equation}
		\label{g1:eq:local-gauss-bonnet}
		\int_D K\,\mathrm{d}A
		+\int_{\partial D}k_g\,\mathrm{d}s
		+\sum_j\alpha_j=2\pi.
	\end{equation}
	边界积分表示各光滑边上积分之和；没有角点时，最后一项为空和。
\end{theorem}
\begin{proof}
	在给定坐标图上将坐标向量正交单位化，得到覆盖整个 \(D\) 及其邻域的正向标架，记其联络形式为 \(\omega\)。对每条边使用\eqref{g1:eq:tangent-connection-angle}，加上所有外角，再用\cref{g1:lem:piecewise-turning}，得到
	\[
		\int_{\partial D}k_g\,\mathrm{d}s+\sum_j\alpha_j
		=2\pi+\int_{\partial D}\omega.
	\]
	由\cref{g1:lem:green-formula}，当前分片 \(C^2\) 边界及正向参数满足二维积分公式的条件，故
	\[
		\int_{\partial D}\omega=\int_D\mathrm{d}\omega
		=-\int_D K\,\mathrm{d}A,
	\]
	最后一步是\cref{g1:prop:connection-curvature-form}的曲率形式公式。代回即得结论。
\end{proof}

这就是带边界和角点的\term[jubugaussbonnetgongshi]{局部 Gauss–Bonnet 公式}[local Gauss--Bonnet formula]。曲面可以嵌入三维欧氏空间，也可以只给定二维度量；证明仅用了度量的协变导数及其曲率。区域是拓扑圆盘，并不自动意味着它被当前某张坐标图覆盖。本定理明确保留单图假设；\secref{g1:sec:0904}将通过分解与粘合去掉它，并处理有孔的区域。

\ProseParagraphBreak

公式中的 \(\mathrm{d}A\) 在正向坐标下就是正的面积测度。若保持曲面定向，却把边界反向走，\(k_g\) 和有向转角都变号；以 \(-\partial D\) 表示这个反向边界，正确的式子是
\[
	-\int_D K\,\mathrm{d}A+\int_{-\partial D}k_g\,\mathrm{d}s+\sum_j(-\alpha_j)=-2\pi.
\]
这里仍按正面积测量 \(D\)，所以面积积分前也须带负号。若将曲面定向与边界方向一起反转，则 \(J\) 与 \(T\) 同时变号，\(\mathrm{D}_sT\) 保持不变，故 \(\langle \mathrm{D}_sT,JT\rangle\) 不变；区域内角与 \(K\) 也不变，原公式便仍成立。正向边界上的 \(JT\) 是指向区域内部的单位余法向，因而此时 \(k_g\) 也可理解为边界加速度的向内分量。

\begin{example}[平面圆盘与多边形]
	\label{g1:exm:planar-gauss-bonnet}
	在平面圆盘上，\(K=0\)，半径 \(r\) 的正向圆周满足 \(k_g=1/r\)、长度为 \(2\pi r\)，因此边界积分为 \(2\pi\)。对有 \(m\) 个顶点的平面简单多边形，各边是直线，\(k_g=0\)，所以公式给出
	\[
		\sum_{j=1}^m(\pi-\beta_j)=2\pi,
		\qquad \sum_{j=1}^m\beta_j=(m-2)\pi.
	\]
	这不要求多边形凸；凹顶点必须采用大于 \(\pi\) 的内部角，外角相应取负。
\end{example}

\subsection{测地三角形的角亏}
\label{g1:subsec:0903:04}

平面三角形的边没有沿平面的弯曲，因此全部转向都集中在顶点。在曲面上，测地线正好承担同样的角色：其单位切向量平行，故 \(\mathrm{D}_sT=0\)，从而 \(k_g=0\)。把三条测地线弧组成的简单边界代入局部公式，就能直接读出三角形的角和。

\begin{corollary}[测地三角形的角和]
	\label{g1:cor:geodesic-triangle-defect}
	设有向光滑曲面上一块单图闭圆盘区域 \(D\) 的边界由三条正则测地线弧组成，弧的内部互不相交，三个内角 \(\beta_1,\beta_2,\beta_3\) 均属于 \((0,2\pi)\)。则
	\begin{equation}
		\label{g1:eq:geodesic-triangle-defect}
		\beta_1+\beta_2+\beta_3-\pi=\int_D K\,\mathrm{d}A.
	\end{equation}
	特别地，若 \(K\equiv c\) 在 \(D\) 上为常数，则角和为 \(\pi+c\operatorname{Area}(D)\)。
\end{corollary}
\begin{proof}
	沿每条测地线取弧长参数后，测地线方程给出 \(\mathrm{D}_sT=0\)，所以三条边的测地曲率积分全为零。由\eqref{g1:eq:local-gauss-bonnet}，
	\[
		\int_D K\,\mathrm{d}A+3\pi-(\beta_1+\beta_2+\beta_3)=2\pi.
	\]
	移项即得结论。常曲率情形把常数提出面积积分即可。
\end{proof}

称 \(\delta=\pi-(\beta_1+\beta_2+\beta_3)\) 为三角形的\term[jiaokui]{角亏}[angular defect]；因此本约定下 \(\delta=-\int_DK\,\mathrm{d}A\)。正曲率给出角盈，负曲率给出正的角亏；若曲率变号，则由整个区域的积分决定，不能只看某个顶点的曲率。

\begin{example}[球面上的三个直角]
	\label{g1:exm:spherical-octant}
	在半径 \(R\) 的球面上，取三个点 \(A=(R,0,0)\)、\(B=(0,R,0)\)、\(C=(0,0,R)\)，用坐标平面中的四分之一大圆弧连接，围成第一卦限中的球面三角形。\cref{g1:exm:sphere-geodesics}已经给出大圆的测地性：用弧长参数化后，其空间加速度是 \(-\gamma/R^2\)，沿球面法向，故切向投影为零。每个顶点处两条边的切向量互相垂直，内角都是 \(\pi/2\)。

	这块闭三角形包含在开半球 \(x+y+z>0\) 中；沿法向 \((1,1,1)\) 的正交投影把该半球表示为球面的一个图像，故满足单图条件。\cref{g1:prop:surface-curvature-formulas}给出球面 \(K=1/R^2\)，于是
	\[
		\frac{\operatorname{Area}(D)}{R^2}
		=\frac{3\pi}{2}-\pi=\frac\pi2,
		\qquad \operatorname{Area}(D)=\frac{\pi R^2}{2}.
	\]
	这一结果也等于球面面积 \(4\pi R^2\) 的八分之一。三条边各自在曲面内都没有弯曲，三个直角仍能闭合，差别由区域内部的正曲率承担。
\end{example}

同样的代入还适用于 \(m\) 条测地线围成的单图圆盘：\(\sum_j\beta_j=(m-2)\pi+\int_DK\,\mathrm{d}A\)。但角和公式本身不保证给定三点之间能找到唯一的测地边，也不保证任意三条测地线围出圆盘。存在性、边的选择以及区域的拓扑条件仍须分别核对。

% !TeX root = ../../Book1.tex
% 纲要标题：### 9.4 整体公式的粘合证明
% 纲要位置：几何学计划纲要.md，第 486 行。
% 正文状态：初稿；尚未正式审核。

% 纲要内容开始
% 
% * 适应边界的曲面分解
% * 公共边项抵消及顶点贡献
% * Euler 示性数的出现
% * 闭曲面与带边界曲面的统一公式
% 
% 纲要内容结束

\section{整体公式的粘合证明}
\label{g1:sec:0904}

一个圆盘的局部公式右端是 \(2\pi\)，若把整张曲面切成 \(F\) 个圆盘，直接相加便得到 \(2\pi F\)。但切割也制造了原来没有的边和角：公共边的曲率会抵消，新增角的贡献却必须逐点数清。最后留下的整数并不是圆盘数，而是第六章的 \(V-E+F\)。这一过程还需要先解决一个正则性问题：第六章给出的是拓扑剖分，拓扑边没有可供积分的曲率，不能原样代入本章公式。

\subsection{适应边界的曲面分解}
\label{g1:subsec:0904:01}

下面允许 \(M\) 有分片 \(C^2\) 边界，角点有限，每个角点的内角属于 \((0,2\pi)\)。这里曲面和度量在各局部图中光滑；边界和角点表示在这样的光滑局部环境内截取一个区域，因而内角、边的切向量和协变加速度都有意义。内角为 \(\pi\) 的分段点也允许保留。

\begin{lemma}[坐标圆盘分解]\label{g1:lem:smooth-coordinate-disk-decomposition}
	设 \(M\) 是紧、Hausdorff、第二可数的光滑度量曲面，允许有分片 \(C^2\) 边界及有限个内角属于 \((0,2\pi)\) 的角点。则存在有限个闭圆盘 \(D_1,\ldots,D_F\)，满足：各 \(D_j\) 位于一张光滑坐标图中，边界由有限条 \(C^2\) 正则弧组成，各角的内角属于 \((0,2\pi)\)；圆盘的内部两两不交，合起来覆盖 \(M\)；两个圆盘的公共部分是若干条完整的边及顶点；\(\partial M\) 由分解的边和顶点组成，原有角点均为顶点。
\end{lemma}
\begin{proof}
	我们沿用\cref{g1:lem:finite-network-disk-cells}的构造思路：先造有限曲线网，再取顶点盘、边带，最后把平面区域中的孔连接起来。需要补充的是，所有新边都可以选成分片光滑的，并且每一块仍留在一张坐标图中。

	先在各点附近选带余量的小坐标片。在内点处用圆盘，在光滑边界点处用半圆盘；角点附近，在环境坐标中用小圆周截取两条边界弧围成的扇形区域。后一构造确实给出闭圆盘：两条正则弧从角点出发时，沿各弧的坐标半径在足够短的一段上严格增加，故每条小圆周恰与每侧相交一次；两侧的不同切向射线保证区域有正角。把最小各片的相对内部取成有限覆盖，再为每片保留稍大的同类片，其半径可以在一个开区间内变化而仍包含原来最小的片。这样的余量与\cref{g1:lem:nested-surface-charts}中的三层坐标片起同样作用。

	逐个选择这些较大片的半径，使新增的圆周弧与已经选定的边及 \(\partial M\) 有限次横交，并避开已有顶点。这个选择只需要如下一个实变量事实。若 \(f\in C^1([a,b])\)，则临界值集 \(f(\{f'=0\})\) 可以用总长度任意小的开区间覆盖。事实上，对任意 \(\varepsilon>0\)，在包含 \(\{f'=0\}\) 的开集中令 \(|f'|<\varepsilon\)；把该开集写成不交区间之并，平均值定理说明各像区间的长度不超过相应区间长度的 \(\varepsilon\) 倍，总和不超过 \(\varepsilon(b-a)\)。端点另用任意短的区间覆盖。因此有限个这类临界值集与有限个指定值，不能填满一个开区间。

	对每条已有光滑弧，令 \(f\) 为沿弧的坐标半径平方。只在所选的闭坐标环带上使用此函数；可以把径向函数在更大圆盘外取为常数，光滑延拓后再复合弧的参数，所以弧反复进出坐标图也不会妨碍上述论证。选择一个同时不是各弧临界值、也不是已有顶点半径平方的值。相交处沿弧的半径导数非零，正是两条曲线的切线不同；交点因而孤立。交点集又是紧集，故有限。避开已有顶点也排除了三条新旧边在同一点相交的情形。初始曲线网取为 \(\partial M\)，并把原有角点标出；最后给没有交点的闭曲线添上几个顶点，就得到一个有限的分片光滑嵌入图 \(G\)。

	每个 \(M\setminus G\) 的连通分支都留在某个已选坐标片内。确实，选分支中一点及包含它的坐标片内部；该片的相对边界已在 \(G\) 中，所以在不碰 \(G\) 的连通分支内，不能从片内走到片外。它的闭包也位于这张坐标图中，因为所选闭片与图的边缘之间仍有余量。不过补分支未必是盘，同一条边还可能在其边界行走中出现两次；下面的加宽步骤就是为了消除这种现象。

	在各顶点取两两不交的小坐标盘；边界顶点取相应半盘或角区域。半径取得足够小，每条关联边只以一个短弧进入顶点盘，并横交它的圆弧边界。把余下的边进一步有限分段，使每个闭段在一张坐标图中都是一条函数图像 \(v=f(u)\)；这是正则性与紧性给出的有限选择。在这样的图中，\(|v-f(u)|\le\delta\) 给出一条细带，位于 \(\partial M\) 上的边只取曲面内侧的半条带。可以令带的两端接到相邻顶点盘的边界弧上，并在接点附近调整侧边使其横交圆弧。有限条互不相交的截短边及其他不关联部分都是紧集，逐次缩小带宽即可使不同带的内部不交，且每个顶点盘、每条带都包含在一张坐标图中。顶点盘和带的并记为 \(N\)。它是 \(G\) 的闭邻域，包含 \(\partial M\) 的相对邻域，边界由有限条分片光滑 Jordan 曲线组成。所有接点都采用横交弧，故不会出现零角尖点。

	余下的每个闭区域 \(R\subset M\setminus\operatorname{int}_M N\) 位于原来一个补分支的坐标图中，其边界为有限条互不相交的分片光滑圈。区域只有有限多个：每个非空区域都接到 \(\partial N\) 的某条边界圈，而一条圈的外侧只邻接一个区域。\cref{g1:lem:finite-network-disk-cells}的平面部分说明，每个连通的 \(R\) 是一个外圈内去掉有限个内圈所围闭盘后的区域。若有内圈，在其一个非顶点处与外圈的非顶点处之间画一条内部位于 \(R^\circ\) 的简单分片光滑弧。具体地，先在两端取横交边界的小线段，再用 \(R^\circ\) 中的有限折线连接两段内端，删去折线回路，并在不碰边界的小圆盘内磨圆折线拐角。取这条弧的一条窄闭带，使其短边沿两条边界弧贴合。

	将窄带作为一个独立盘形块取出，剩下区域的两个边界圈合为一个，内圈数减少一；此处保留带的两个不同侧边，不把它们压回同一条割线。故余下区域始终具有真正的、互不相交的 Jordan 边界。有限次以后，每块余下区域只有一个边界圈，由\cref{g1:thm:jordan-schoenflies}是闭圆盘。它们和取出的窄带、顶点盘、原有边带一起给出有限闭盘分解，新增块都仍在所操作的原坐标图中。最后把所有接点统一作为顶点，边界上的分点也同时加入相邻两块。这样公共部分由完整边和顶点组成，原有边界和角点均保留下来。所有边由上述光滑弧构成，角可以选成严格介于 \(0\) 和 \(2\pi\) 之间，得到所需分解。
\end{proof}

这里没有要求分解的边是测地线。普通光滑边已经足够，因为它们的曲率积分会在粘合时抵消；若强求测地三角剖分，反而会额外引入尚未需要的存在性问题。闭盘也不必是三角形，下面按边和角的实际出现次数计数即可。

\subsection{公共边项抵消及顶点贡献}
\label{g1:subsec:0904:02}

先假设 \(M\) 可定向并选定定向，每块 \(D_j\) 的边界都按区域在左的方向行进。在一条内部公共边上，两块给出的行进方向相反。同一个曲线加速度在反向参数化后不变，而左侧的单位切向法向变号，因此两侧测地曲率互为相反数。两次积分使用相同的弧长测度，恰好抵消。

\ProseParagraphBreak

这个抵消其实不需要整张曲面可定向。对每块圆盘分别选一个局部定向，令其边界仍把本块放在左边。写 \(\nu_j\) 为沿边、垂直于边且指向 \(D_j\) 内侧的单位切向量，则正向行进时 \(\nu_j=JT\)，因而
\[
	k_{g,j}=\langle\nabla_TT,\nu_j\rangle.
\]
反向行进将 \(T\) 换成 \(-T\)，但 \(\nabla_{-T}(-T)=\nabla_TT\)；只要固定内侧，右端便不变。公共边两侧的 \(\nu_j\) 相反，所以无论各面选择的定向能否拼成整体定向，两项都抵消。沿原有边界，圆盘内侧与 \(M\) 内侧一致，留下的恰是 \(M\) 的边界曲率积分。于是
\begin{equation}\label{g1:eq:gb-edge-cancellation}
	\sum_{j=1}^F\int_{\partial D_j}k_{g,j}\,\mathrm{d}s
	=\int_{\partial M}k_g\,\mathrm{d}s.
\end{equation}
右端沿各条光滑边求和，不把角点另算成曲率函数的值。

\ProseParagraphBreak

再看角。用 \(V_{\mathrm i},V_{\mathrm b}\) 分别表示内部与边界顶点数，用 \(E_{\mathrm i},E_{\mathrm b}\) 表示内部与边界边数。每次一个圆盘在某顶点占据一个角，都算作一次角的出现，记其内角为 \(\beta_{j,v}\)。在内点，所有小角在切平面内围满一周，总和为 \(2\pi\)；在一个光滑边界点，它们铺满内侧半平面，总和为 \(\pi\)；在原有角点 \(v\)，总和等于 \(M\) 的内角 \(\beta_v\)，而不是 \(\pi\)。\cref{g1:fig:gb-vertex-angles}画的是顶点处带内积的切平面，各扇区就是这些角的切向模型。

\input{Figures/Chapter09-VertexAngles.tex}

为统一记号，给每个边界顶点置 \(\alpha_v=\pi-\beta_v\)，在新增的光滑分点处置 \(\alpha_v=0\)。则
\begin{equation}\label{g1:eq:gb-interior-angle-count}
	\sum_{j,v}\beta_{j,v}
	=2\pi V_{\mathrm i}+\pi V_{\mathrm b}
	-\sum_{v\in\partial M}\alpha_v.
\end{equation}
这里的角度是度量角，选择正交坐标后才与图中的 Euclid 角一致；角能相加是同一个切平面中的事实，不要求坐标图把曲面等距铺平。

\subsection{Euler 示性数的出现}
\label{g1:subsec:0904:03}

每个闭圆盘的边数等于角的出现次数。内部边出现两次，边界边出现一次，故总出现次数为 \(2E_{\mathrm i}+E_{\mathrm b}\)。把每个面角换成外角 \(\pi-\beta_{j,v}\)，再代入\eqref{g1:eq:gb-interior-angle-count}，得到
\begin{align*}
	\sum_{j,v}(\pi-\beta_{j,v})
	&=\pi(2E_{\mathrm i}+E_{\mathrm b})-2\pi V_{\mathrm i}-\pi V_{\mathrm b}
	+\sum_{v\in\partial M}\alpha_v\\
	&=2\pi(E_{\mathrm i}-V_{\mathrm i})+\sum_{v\in\partial M}\alpha_v.
\end{align*}
第二步用到 \(E_{\mathrm b}=V_{\mathrm b}\)：\cref{g1:prop:boundary-circles}说明紧曲面的边界是有限个圆周，每条圆周分成多少条边，就有多少个分点。这也解释了为何边界顶点不能与内部顶点一起按 \(2\pi\) 计角。

\ProseParagraphBreak

现在对每块 \(D_j\) 使用\cref{g1:thm:local-gauss-bonnet}。分解引理保证它是落在单张图中的闭圆盘，边界具有所要求的正则性和非零角，因此每块的右端都是 \(2\pi\)。对各块相加，面积积分可加、公共边按\eqref{g1:eq:gb-edge-cancellation}抵消，故
\[
	\int_MK\,\mathrm{d}A+\int_{\partial M}k_g\,\mathrm{d}s
	+\sum_{v\in\partial M}\alpha_v
	=2\pi(F-E_{\mathrm i}+V_{\mathrm i}).
\]
由于 \(E_{\mathrm b}=V_{\mathrm b}\)，右端括号等于
\[
	F-(E_{\mathrm i}+E_{\mathrm b})+(V_{\mathrm i}+V_{\mathrm b})=V-E+F=\chi(M).
\]
最后一步正是\cref{g1:prop:euler-triangulation-independent}：它不仅证明三角剖分的交错和不变，也证明任何有限闭盘分解的交错和都等于同一个 Euler 示性数。这里已经构造出真正的闭盘分解，因而满足该命题的条件。

\ProseParagraphBreak

可以把这段计数理解为两笔账。曲率积分记下各小片内部与边缘的几何；面片外角的总和则记下这些小片在顶点附近如何拼接。把人为切出的边和角消去后，几何账中留下总曲率，组合账中留下 \(\chi(M)\)。

\subsection{闭曲面与带边界曲面的统一公式}
\label{g1:subsec:0904:04}

\begin{theorem}[Gauss--Bonnet]\label{g1:thm:global-gauss-bonnet}
	设 \(M\) 为紧、Hausdorff、第二可数的光滑度量曲面，允许不连通、不可定向以及分片 \(C^2\) 边界。边界由有限条正则弧组成，角点的内角 \(\beta_v\) 属于 \((0,2\pi)\)。度量在边界及角点的局部光滑环境中光滑延拓。令 \(K\) 为 Gauss 曲率，\(\mathrm{d}A\) 为度量面积测度，沿每条边界弧取指向 \(M\) 内侧的单位切向法向 \(\nu\)，并定义
	\[
		k_g=\langle\nabla_TT,\nu\rangle,\qquad \alpha_v=\pi-\beta_v.
	\]
	这里 \(T\) 是沿该弧任一方向的单位切向量。则
	\begin{equation}\label{g1:eq:global-gauss-bonnet}
		\int_MK\,\mathrm{d}A+\int_{\partial M}k_g\,\mathrm{d}s+\sum_v\alpha_v
		=2\pi\chi(M).
	\end{equation}
	当 \(M\) 可定向时，边界按区域在左的方向行进，上述 \(k_g\) 就是 \(\langle\nabla_TT,JT\rangle\)。不连通时，\(\chi(M)\) 为各连通分支的 Euler 示性数之和。
\end{theorem}
\begin{proof}
	对每个连通分支使用\cref{g1:lem:smooth-coordinate-disk-decomposition}，按前两小节的边项抵消和顶点计数相加，即得\eqref{g1:eq:global-gauss-bonnet}。局部面积测度与 \(K\) 不因反转定向而改变；内侧测地曲率的抵消已经逐边证明，故不需要全局定向。最后，紧曲面只有有限个连通分支，这是\cref{g1:prop:boundary-circles}的另一项结论；将有限个公式相加即可。
\end{proof}

若边界光滑，角项为空；若 \(M\) 是闭曲面，边界积分也为空，公式成为
\begin{equation}\label{g1:eq:closed-gauss-bonnet}
	\int_MK\,\mathrm{d}A=2\pi\chi(M).
\end{equation}
另一方面，若 \(M\) 本身是一个圆盘，即使它不落在任何一张坐标图内，也有 \(\chi(M)=1\)，所以仍得到右端为 \(2\pi\) 的带边界公式。单图条件是局部证明所用的条件，经粘合后便不再是整体结论的限制。

\begin{example}\label{g1:exm:flat-mobius-gauss-bonnet}
	取带 Euclid 度量的矩形 \([0,L]\times[-a,a]\)，其中 \(L,a>0\)，按 \((0,y)\sim(L,-y)\) 识别两条竖边。因为识别映射在邻近条带上由刚性变换 \((x,y)\mapsto(x+L,-y)\) 给出，度量光滑粘合成一条 Möbius 带。这里的度量由矩形粘合定义，不需要把它作为 \(\mathbb R^3\) 中的曲面实现。

	局部图都是 Euclid 图，所以 \(K=0\)。上下两条水平边接成同一条光滑边界，各局部段都是直线，故 \(k_g=0\)，识别处也没有角点。\cref{g1:prop:surface-euler-formulas}给出 \(\chi(N_{1,1})=2-1-1=0\)，于是公式两边都是零。这个例子中不存在连续的全局定向，但“指向带内侧”的边界法向仍然存在，正好对应定理中采用的表述。
\end{example}

有边界时不能只看面积曲率。平面圆盘的 \(K\) 处处为零而 \(\chi=1\)，缺少的 \(2\pi\) 全由边界转动提供；平面环带的内圈则带来负的边界测地曲率，抵消外圈的贡献。角项承担同样的补偿作用，所以把光滑边界换成多边形边界后，公式的右端仍是同一个拓扑数。

% !TeX root = ../../Book1.tex
% 纲要标题：### 9.5 曲率积分的几何后果
% 纲要位置：几何学计划纲要.md，第 493 行。
% 正文状态：初稿；尚未正式审核。

% 纲要内容开始
% 
% * 亏格与曲率符号的限制
% * 总曲率与向量场零点的联系
% * 度量变化下的积分不变量
% 
% ---
% 
% 纲要内容结束

\section{曲率积分的几何后果}
\label{g1:sec:0905}

整体公式把一个可以逐点变化的函数，积分成了只能取离散值的拓扑量。因而它既能排除某些曲率分布，也能解释为什么改变度量时，总有一部分局部变化必须互相补偿。本节先看曲率符号，再把同一个积分理解为向量场绕零点的转动，最后计算保角改变度量时这种补偿怎样发生。

\subsection{亏格与曲率符号的限制}
\label{g1:subsec:0905:01}

本小节的曲面均紧、连通且无边界。第六章的\cref{g1:prop:surface-euler-formulas}给出：可定向亏格为 \(g\) 的曲面 \(\Sigma_g\) 满足 \(\chi=2-2g\)，不可定向标准曲面 \(N_k\) 满足 \(\chi=2-k\)，其中 \(k\ge1\)。因此\eqref{g1:eq:closed-gauss-bonnet}立即成为
\begin{equation}
	\label{g1:eq:total-curvature-genus}
	\int_{\Sigma_g}K\,\mathrm{d}A=4\pi(1-g),\qquad
	\int_{N_k}K\,\mathrm{d}A=2\pi(2-k).
\end{equation}
这里限制的是整张曲面的总积分。它并没有说一个点的曲率可以由亏格决定，也没有说总积分为零就意味着处处平坦。

\begin{proposition}
	\label{g1:prop:curvature-sign-topology}
	紧连通闭度量曲面满足下列限制。
	\begin{enumerate}
		\item 若 \(K\ge0\)，则 \(\chi\ge0\)；若还有一点 \(K>0\)，则 \(\chi>0\)。
		\item 若 \(K\le0\)，则 \(\chi\le0\)；若还有一点 \(K<0\)，则 \(\chi<0\)。
		\item 若 \(\chi=0\)，且 \(K\) 在全曲面上不变号，则 \(K\equiv0\)。
	\end{enumerate}
\end{proposition}
\begin{proof}
	积分保持非负性和非正性，故前两项的弱不等式来自 Gauss--Bonnet。若连续函数 \(K\) 在一点为正，则在某个小坐标圆盘上有 \(K\ge c>0\)。度量正定使这个圆盘的面积为正，所以在 \(K\ge0\) 的前提下，全曲率积分严格为正。负号情形同理。最后，若 \(K\) 不变号而又不恒为零，前两项就使积分严格非零，与 \(\chi=0\) 矛盾。
\end{proof}

结合\cref{g1:thm:surface-classification-unique}，这些限制可写成下表。表中列的是可能的拓扑类型，并不单凭这个必要条件断言某种指定曲率的度量一定存在。
\begin{center}
	\begin{tabular}{lll}
		\toprule
		曲率条件 & 可定向情形 & 不可定向情形\\
		\midrule
		\(K\ge0\)，且某处 \(K>0\) & 球面 \(\Sigma_0\) & 射影平面 \(N_1\)\\
		\(K\equiv0\) & 环面 \(\Sigma_1\) & Klein 瓶 \(N_2\)\\
		\(K\le0\)，且某处 \(K<0\) & \(\Sigma_g\)，\(g\ge2\) & \(N_k\)，\(k\ge3\)\\
		\bottomrule
	\end{tabular}
\end{center}

若 \(\Sigma_g\) 的曲率为常数 \(c<0\)，则还得到面积必须满足
\[
	\operatorname{Area}(\Sigma_g)=\frac{4\pi(g-1)}{-c}.
\]
所以固定曲率为 \(-1\) 后，闭可定向曲面的面积只能是 \(4\pi(g-1)\)。这个结论的前提仍是已经给定这样的度量；其构造属于后面的几何内容。

\begin{example}
	\label{g1:exm:two-torus-curvatures}
	平坦环面 \(\mathbb R^2/(a\mathbb Z\times b\mathbb Z)\) 的各坐标图通过平移相接，欧氏度量因而可以粘合。局部 Christoffel 系数全为零，故 \(K=0\)，总曲率当然为零。

	另一种熟悉的环面是三维空间中的旋转曲面
	\[
		x(u,v)=((R+r\cos v)\cos u,(R+r\cos v)\sin u,r\sin v),\qquad R>r>0.
	\]
	两个参数都以 \(2\pi\) 为周期。直接计算第一基本形式与形算子的行列式，得到
	\[
		\mathrm{d}A=r(R+r\cos v)\,\mathrm{d}u\,\mathrm{d}v,
		\qquad K=\frac{\cos v}{r(R+r\cos v)}.
	\]
	这个计算也可从\cref{g1:prop:revolution-two-forms}代入生成圆读出。外侧 \(\cos v>0\) 的曲率为正，内侧为负，且
	\[
		\int K\,\mathrm{d}A=\int_0^{2\pi}\int_0^{2\pi}\cos v\,\mathrm{d}v\,\mathrm{d}u=0.
	\]
	两种度量都符合环面的 Euler 示性数，但后者并不平坦。拓扑相同，也不意味着这两个度量局部等距。
\end{example}

\subsection{总曲率与向量场零点的联系}
\label{g1:subsec:0905:02}

一个处处非零的切向量场经单位化后，可以充当移动标架的第一根轴。若这样的轴只能在删掉某些点以后选取，就要问它绕着这些点转了几圈。第五章的圆周绕数正好能记录这个局部现象。

\begin{definition}
	\label{g1:def:surface-vector-index}
	设 \(V\) 是光滑切向量场，\(p\) 是其孤立零点。在 \(p\) 附近选一个定向坐标圆盘和正向光滑正交单位标架 \((e_1,e_2)\)，写 \(V=a e_1+b e_2\)。沿包围 \(p\)、内部没有其他零点的小正向坐标圆周 \(C\)，定义 \(V\) 在 \(p\) 的\term[xiangliangchangzhishu]{向量场指数}[index of a vector field]为映射
	\[
		C\longrightarrow S^1,\qquad q\longmapsto\frac{(a(q),b(q))}{\sqrt{a(q)^2+b(q)^2}}
	\]
	的绕数，记为 \(\operatorname{ind}_p(V)\)。等价地，若沿一周连续选择角度 \(\theta\)，则
	\[
		\operatorname{ind}_p(V)=\frac{\theta(\text{终点})-\theta(\text{起点})}{2\pi}.
	\]
	角度提升及所得整数的同伦不变性见\cref{g1:lem:circle-angle-lift,g1:thm:circle-fundamental}。\symidx{\(\operatorname{ind}_p(V)\)：切向量场在孤立零点的指数}
\end{definition}

这个整数与定义中的辅助选择无关。两条足够小的圆周在无零点的穿孔圆盘中可相互变形，单位化向量给出相应的圆周映射同伦。换正向正交标架时，新旧标架之间的旋转定义在整个圆盘上，其限制到边界的绕数为零；向量的角度减去这个旋转角，不改变一周的总增量。若改用另一张小坐标圆盘，可先在二者交中取更小的圆盘作比较。反转局部定向则同时反转边界行进方向和记录向量角度的方向，两个负号抵消。因此在不可定向曲面上，这个局部指数也有意义。

\begin{proposition}
	\label{g1:prop:nondegenerate-vector-index}
	在上述标架中置 \(F=(a,b)\)，并以零点为坐标原点。如果 \(A=\mathrm{d}F_0\) 可逆，则
	\[
		\operatorname{ind}_p(V)=\operatorname{sgn}\det A.
	\]
\end{proposition}
\begin{proof}
	由可逆性，存在 \(c>0\) 使 \(|Ax|\ge c|x|\)。Taylor 展开给出 \(F(x)=Ax+r(x)\)，其中 \(|r(x)|/|x|\to0\)。在足够小的圆周上，\(Ax+t r(x)\) 对 \(0\le t\le1\) 始终非零，所以 \(F/|F|\) 与 \(Ax/|Ax|\) 同伦。

	对 \(A\) 的列作 Gram--Schmidt 正交化，可写 \(A=QR\)，其中 \(R\) 上三角且对角元为正，\(Q\) 正交。矩阵 \((1-t)R+tI\) 始终可逆，故只需算 \(Q\) 的绕数。若 \(\det Q=1\)，它是旋转，绕数为 \(1\)；若 \(\det Q=-1\)，它是旋转与反射的复合，绕数为 \(-1\)。这正是 \(\operatorname{sgn}\det A\)。
\end{proof}

例如平面场 \((x,y)\) 与 \((-x,-y)\) 的指数都为 \(1\)：向外或向内改变的是箭头方向，不改变绕原点一周时箭头也转一周的事实。鞍点场 \((x,-y)\) 的指数是 \(-1\)。这一符号差别比“是不是零点”本身提供了更多信息。

\begin{theorem}[Poincar\'e--Hopf]
	\label{g1:thm:surface-poincare-hopf}
	设 \(M\) 为紧光滑闭曲面，\(V\) 为只有孤立零点的光滑切向量场。则零点只有有限个，且
	\begin{equation}
		\label{g1:eq:surface-index-sum}
		\sum_{V(p)=0}\operatorname{ind}_p(V)=\chi(M).
	\end{equation}
	允许 \(V\) 没有零点，此时左端为空和 \(0\)。
\end{theorem}

这是 Poincar\'e--Hopf 定理的二维情形。\footnote{Poincar\'e，庞加莱，Jules Henri Poincar\'e（1854--1912），法国数学家，对拓扑学、微分方程与动力系统作出奠基性贡献。}\footnote{Hopf，霍普夫，Heinz Hopf（1894--1971），德国出生的瑞士数学家，研究拓扑学、向量场及同伦不变量。}我们用刚证明的曲率公式给出证明，不预先调用高维版本。
\begin{proof}
	零点集闭且位于紧曲面中。每个零点都能被一个只含该零点的邻域隔开；若零点无限，紧性会给出一个聚点，连续性又使该聚点成为非孤立零点，矛盾。只需逐连通分支证明。

	先设 \(M\) 可定向。删去每个零点附近互不相交的坐标小开圆盘 \(D_p(\varepsilon)\)，所得曲面记为 \(M_\varepsilon\)。在它上面，
	\[
		E_1=\frac{V}{|V|},\qquad E_2=JE_1
	\]
	是一组全局正交单位标架，其联络形式 \(\Omega\) 满足 \(\mathrm{d}\Omega=-K\,\mathrm{d}A\)。Green 公式在这里也能逐块相加：用\cref{g1:lem:smooth-coordinate-disk-decomposition}分解 \(M_\varepsilon\)，对每块应用\cref{g1:lem:green-formula}，公共边上一形式的积分因方向相反而抵消。因此
	\[
		\int_{M_\varepsilon}K\,\mathrm{d}A=-\int_{\partial M_\varepsilon}\Omega.
	\]
	在 \(p\) 附近取一组能延到圆盘中心的正向光滑标架，联络形式记为 \(\omega_p\)。沿小圆周写 \(E_1=\cos\theta\,e_1+\sin\theta\,e_2\)。\cref{g1:prop:connection-frame-change}给出 \(\Omega=\omega_p+\mathrm{d}\theta\)。\(\theta\) 可以沿一周提升，不要求在穿孔圆盘上成为单值函数。按小圆盘的正向积分，便有
	\[
		\int_{\partial D_p(\varepsilon)}\Omega
		=\int_{\partial D_p(\varepsilon)}\omega_p+2\pi\operatorname{ind}_p(V).
	\]
	洞的边界在 \(M_\varepsilon\) 中取相反方向，故
	\[
		\int_{M_\varepsilon}K\,\mathrm{d}A
		=\sum_p\left(\int_{\partial D_p(\varepsilon)}\omega_p+2\pi\operatorname{ind}_p(V)\right).
	\]
	\(\omega_p\) 在整个小圆盘上光滑，其系数有界，圆周长度趋于零，因而前一项趋于零。删掉的圆盘面积也趋于零，而 \(K\) 有界。令 \(\varepsilon\to0\)，得到 \(\int_M K\,\mathrm{d}A=2\pi\sum_p\operatorname{ind}_p(V)\)，再用闭曲面的 Gauss--Bonnet 公式即得结论。如果没有零点，就直接在整个闭曲面上对 \(\Omega\) 积分，得到总曲率为零。

	若 \(M\) 不可定向，取\cref{g1:thm:orientation-cover}中的定向二重覆盖 \(\pi:\widetilde M\to M\)。以局部坐标的提升作为光滑坐标，并赋予拉回度量，\(\pi\) 就是光滑局部等距。场 \(V\) 唯一提升为 \(\widetilde V\)，每个零点提升成两个零点。局部微分同胚只是在同一指数定义中换坐标；即使它反转定向，前面解释的两个负号仍然抵消，所以两个指数都等于原指数。可定向情形给出 \(2\sum_p\operatorname{ind}_p(V)=\chi(\widetilde M)\)，再由\cref{g1:prop:surface-finite-cover-euler}知 \(\chi(\widetilde M)=2\chi(M)\)，除以 \(2\) 即得所求。
\end{proof}

\begin{example}
	\label{g1:exm:sphere-vector-indices}
	单位球面上高度函数 \(h(p)=p_3\) 的梯度是
	\[
		V(p)=e_3-p_3p.
	\]
	因为对每个切向量 \(w\) 都有 \(\langle V,w\rangle=w_3=\mathrm{d}h(w)\)，这个表达式确实给出曲面梯度。它仅在南、北极为零。北极处沿切方向求导得到 \(-I\)，南极处得到 \(I\)；二者在二维中行列式都为正，所以两个指数均为 \(1\)，总和为球面的 \(\chi=2\)。特别地，球面不可能有处处非零的光滑切向量场。平坦环面上的常向量场则能穿过平移识别而下降为处处非零的场，与 \(\chi=0\) 相符。
\end{example}

指数并不是给每个零点机械地计数一次。退化零点可以有更大的正指数、负指数或零指数；定理约束的是带指数的总和。它的证明还指出了联系所在：删去零点后标架存在，但标架绕这些洞的转角，恰好补回了整张曲面的曲率积分。

\subsection{度量变化下的积分不变量}
\label{g1:subsec:0905:03}

固定一张光滑闭曲面 \(M\)，任取两个光滑度量 \(g_0,g_1\)。虽然曲率和面积测度都依赖度量，Gauss--Bonnet 却给出
\[
	\int_MK_{g_0}\,\mathrm{d}A_{g_0}=\int_MK_{g_1}\,\mathrm{d}A_{g_1}=2\pi\chi(M).
\]
这里不需要先证明两种度量之间有什么特殊几何关系。事实上 \(g_t=(1-t)g_0+tg_1\) 对 \(0\le t\le1\) 始终正定，提供了一个直接的光滑变化；沿整个变化，总积分都固定。有边界时，固定的则是曲率积分、边界测地曲率积分与外角和三者之和。

为看清局部补偿，考虑保角改变度量 \(\widetilde g=e^{2f}g\)，其中 \(f\) 是光滑实函数。第八章在平面保角坐标中算过曲率；下面给出相对于任意曲面度量的公式。仍采用第八章约定 \(\Delta_g=-\operatorname{div}_g\operatorname{grad}_g\)。

\begin{proposition}
	\label{g1:prop:conformal-gauss-bonnet-change}
	在定向光滑度量曲面上，若 \(\widetilde g=e^{2f}g\)，则
	\begin{align}
		\mathrm{d}\widetilde A&=e^{2f}\mathrm{d}A,\label{g1:eq:conformal-area-change}\\
		\widetilde K\,\mathrm{d}\widetilde A&=(K+\Delta_g f)\,\mathrm{d}A.\label{g1:eq:conformal-curvature-change}
	\end{align}
	对区域的一段光滑边界，令 \(T\) 为原度量的正向单位切向量，\(\nu=JT\) 为内侧单位切向法向，\(\partial_\nu f=\mathrm{d}f(\nu)\)。则
	\begin{equation}
		\label{g1:eq:conformal-boundary-change}
		\widetilde k_g\,\mathrm{d}\widetilde s=(k_g-\partial_\nu f)\,\mathrm{d}s.
	\end{equation}
	保角变化保持每个角点的内角及外角。
\end{proposition}
\begin{proof}
	二维度量矩阵乘以 \(e^{2f}\) 后，其行列式乘以 \(e^{4f}\)，开方即得面积公式。为求曲率，先比较协变导数。定义
	\[
		\widetilde\nabla_VW=\nabla_VW+V(f)W+W(f)V-g(V,W)\operatorname{grad}_g f.
	\]
	右端新增项关于 \(V,W\) 对称，故仍无挠。对 \(\widetilde g(W,Z)=e^{2f}g(W,Z)\) 求导并展开，新增项在两个内积中的贡献为 \(2V(f)e^{2f}g(W,Z)\)，其余交叉项互相抵消，正好给出度量相容性。因此由\cref{g1:prop:surface-connection-properties}的唯一性，这就是新度量的协变导数。

	取原度量的正向正交单位标架 \((e_1,e_2)\)，以 \((\vartheta^1,\vartheta^2)\) 表示其对偶一形式。新标架为 \(\widetilde e_i=e^{-f}e_i\)。把它代入刚得到的公式，可得
	\[
		\widetilde\omega(V)=\omega(V)+e_1(f)g(V,e_2)-e_2(f)g(V,e_1).
	\]
	记 \(*\vartheta^1=\vartheta^2\)、\(*\vartheta^2=-\vartheta^1\)，线性延拓到一形式，则上述等式就是 \(\widetilde\omega=\omega+*\mathrm{d}f\)。这个星号这里只表示一形式在定向正交标架中的直角旋转；换正向正交标架不改变它的定义。

	还需要算出这个一形式的外微分。置 \(Z=\operatorname{grad}_g f=Z^u\partial_u+Z^v\partial_v\)，在正向坐标中写 \(h=\sqrt{\det(g_{ij})}\)。由定义
	\[
		*\mathrm{d}f=h(Z^u\mathrm{d}v-Z^v\mathrm{d}u),\qquad
		\mathrm{d}(*\mathrm{d}f)=\bigl(\partial_u(hZ^u)+\partial_v(hZ^v)\bigr)\mathrm{d}u\wedge\mathrm{d}v.
	\]
	第八章的散度公式表明，右端为 \((\operatorname{div}_g Z)\mathrm{d}A=-(\Delta_g f)\mathrm{d}A\)。因此
	\[
		-\widetilde K\,\mathrm{d}\widetilde A
		=\mathrm{d}\widetilde\omega
		=\mathrm{d}\omega+\mathrm{d}(*\mathrm{d}f)
		=-(K+\Delta_g f)\mathrm{d}A,
	\]
	即得曲率公式。

	沿边界，新旧标架中切向量的角度 \(\theta\) 相同，且 \(\mathrm{d}\widetilde s=e^f\mathrm{d}s\)。用\cref{g1:prop:tangent-connection-angle}将曲率乘弧长写成 \(\mathrm{d}\theta+\omega\)，得到
	\[
		\widetilde k_g\,\mathrm{d}\widetilde s
		=k_g\,\mathrm{d}s+(*\mathrm{d}f)(T)\,\mathrm{d}s
		=(k_g-\partial_\nu f)\,\mathrm{d}s.
	\]
	最后，度量乘以同一个正数不改变两条射线夹角的余弦，故角点的角也保持不变。
\end{proof}

曲率面积形式新增的 \(\Delta_g f\) 与边界新增项恰好抵消。具体地，由刚才的局部计算和逐块 Green 公式，
\[
	\int_M\Delta_g f\,\mathrm{d}A
	=-\int_{\partial M}*\mathrm{d}f
	=\int_{\partial M}\partial_\nu f\,\mathrm{d}s.
\]
这里使用的是内侧法向；若换成外侧，最后一项前面要加负号。把此式与\eqref{g1:eq:conformal-boundary-change}相加，总变化为零。闭曲面没有边界，便直接得到 \(\int_M\Delta_g f\,\mathrm{d}A=0\)。不可定向时，这些曲率和内侧法向公式在局部反转定向后仍不变，可以提升到定向二重覆盖再除以 \(2\)，所以总积分的结论照样成立。

\begin{example}
	\label{g1:exm:conformal-flat-torus}
	在 \(u,v\) 都以 \(2\pi\) 为周期的平坦环面上，取 \(f(u,v)=\varepsilon\cos u\)，其中 \(\varepsilon>0\)。新度量为 \(e^{2\varepsilon\cos u}(\mathrm{d}u^2+\mathrm{d}v^2)\)。因 \(K=0\)、\(\Delta_g f=-f_{uu}-f_{vv}=\varepsilon\cos u\)，有
	\[
		\widetilde K=\varepsilon e^{-2\varepsilon\cos u}\cos u,\qquad
		\widetilde K\,\mathrm{d}\widetilde A=\varepsilon\cos u\,\mathrm{d}u\,\mathrm{d}v.
	\]
	局部曲率已经出现正负变化，总积分却仍为零。这里曲率的补偿可以直接由周期积分看见；一般度量变化中，保证这种补偿的就是整体 Gauss--Bonnet 公式。
\end{example}

本章由切向投影出发，经联络形式、边界转角和有限分解，得到总曲率的拓扑不变性。向量场指数提供了同一整数的另一种计算方式。这些联系都已在二维中完成，不需要预先知道一般流形的特征类理论；后面的推广应当解释这里已经看清的现象，而不是替代这些具体证明。


\begin{chaptersummary}
	协变导数比较沿曲线移动的切向量；平行移动由线性常微分方程给出，并保持内积。在正向正交标架中，平行向量的角度满足 \(\theta'=-\omega(\dot\gamma)\)，而联络形式满足 \(\mathrm{d}\omega=-K\,\mathrm{d}A\)。因此绕一个坐标圆盘的平行移动，其旋转角等于内部曲率积分模 \(2\pi\) 的值。

	\ProseParagraphBreak

	测地线的速度沿自身平行。初值方程给出局部唯一解和光滑依赖，第一变分说明固定端点下的驻值条件；能量与长度之间还需要核对参数和恒速条件。指数映射与 Gauss 引理把径向和角向分开，在足够小的正常邻域内证明径向测地线最短。能继续延伸并不意味着能一直保持最短。

	\ProseParagraphBreak

	正向边界满足 \(k_g=\theta'+\omega(T)\)，角点另贡献 \(\alpha=\pi-\beta\)。圆盘的总切向转角为 \(2\pi\)，与 Green 公式相合便得到局部 Gauss--Bonnet。对有限坐标圆盘分解相加，内部边的测地曲率抵消，内部与边界顶点分别按完整角和内侧角计数，最终留下 \(2\pi(V-E+F)\)。内侧法向的表述使整体公式也适用于不可定向曲面。

	\ProseParagraphBreak

	闭曲面的总曲率等于 \(2\pi\chi\)，从而限制曲率符号与亏格。沿向量场的单位方向选标架，则零点附近的转动给出指数，总指数同样等于 \(\chi\)。改变度量可以改变每一点的曲率；在保角变化中，新增的曲率面积项为 \(\Delta_g f\,\mathrm{d}A\)，其积分由边界项补偿。总公式保持不变，局部几何仍可以发生显著变化。
\end{chaptersummary}

\BookExercises

\begin{exercise}\label{g1:ex:warped-annulus-holonomy}
	在环带 \(I\times(\mathbb R/2\pi\mathbb Z)\) 上给定度量 \(\mathrm{d}u^2+f(u)^2\mathrm{d}v^2\)，其中 \(I\) 为开区间，\(f>0\) 光滑，定向为 \(\mathrm{d}u\wedge\mathrm{d}v\)。
	\begin{enumerate}
		\item 对标架 \(e_1=\partial_u\)、\(e_2=f^{-1}\partial_v\)，求联络形式、Gauss 曲率及沿 \(u=u_0\) 增加 \(v\) 一周的平行移动旋转角。
		\item 取 \(I=(0,\infty)\)、\(f(u)=cu\)、\(c>0\)，说明曲率处处为零仍可能有非恒等的环路平行移动，并指出这为何不与局部曲率积分公式矛盾。
	\end{enumerate}
\end{exercise}
\begin{solution}
	1. 非零 Christoffel 系数为 \(\Gamma^u_{vv}=-ff'\)、\(\Gamma^v_{uv}=\Gamma^v_{vu}=f'/f\)。所以 \(\nabla_{\partial_u}e_1=0\)、\(\nabla_{\partial_v}e_1=f'e_2\)，给出 \(\omega=f'\mathrm{d}v\)。又有 \(\mathrm{d}A=f\,\mathrm{d}u\wedge\mathrm{d}v\)，故
	\[
		\mathrm{d}\omega=f''\mathrm{d}u\wedge\mathrm{d}v=-K\,\mathrm{d}A,
		\qquad K=-\frac{f''}{f}.
	\]
	这个标架沿 \(v\) 是周期的，平行向量相对它的角增量为 \(-\int\omega=-2\pi f'(u_0)\)。因此返回后的旋转角为该数模 \(2\pi\) 的值。

	2. 此时 \(K=0\)，旋转角却是 \(-2\pi c\pmod{2\pi}\)，当 \(c\notin\mathbb Z\) 时并非恒等。纬圈在环带内不围成圆盘，不能直接使用\cref{g1:prop:surface-local-holonomy}。把 \(u=0\) 添进去也不能随意补救：若 \(c\ne1\)，小圆周周长与径向距离之比是 \(2\pi c\)，不是光滑平面点应有的极限 \(2\pi\)，新添点具有锥形奇性。局部平坦性只保证足够小的可缩环路平行移动为恒等，不消除这种绕孔的整体效应。
\end{solution}

\begin{exercise}\label{g1:ex:sphere-octant-transport}
	在外向定向的单位球面上，依次沿短大圆弧从 \(N=(0,0,1)\) 到 \(A=(1,0,0)\)，再到 \(B=(0,1,0)\)，最后回到 \(N\)。把初始切向量 \((1,0,0)\in \mathrm{T}_NS^2\) 平行移动一周，直接求终点向量，并与所围球面八分之一的曲率积分比较。
\end{exercise}
\begin{solution}
	第一段可写 \(\gamma_1(t)=(\sin t,0,\cos t)\)，\(0\le t\le\pi/2\)。其单位切向量 \((\cos t,0,-\sin t)\) 的导数为 \(-\gamma_1(t)\)，纯属法向，所以它沿大圆平行。初始向量正是该切向量，到 \(A\) 时成为 \((0,0,-1)\)。

	第二段为赤道弧 \((\cos t,\sin t,0)\)。常向量 \((0,0,-1)\) 处处与球面法向垂直，环境导数为零，所以它平行地到达 \(B\)。第三段取 \(\gamma_3(t)=(0,\cos t,\sin t)\)，其单位切向量在 \(B\) 为 \((0,0,1)\)。待移动向量是该切向量的负值，故沿弧保持为 \((0,\sin t,-\cos t)\)，到 \(N\) 时成为 \((0,1,0)\)。北极的正向直角旋转把 \((1,0,0)\) 送到 \((0,1,0)\)，所以净旋转角为 \(\pi/2\)。

	该路径把第一卦限球面区域置于左侧；球面对坐标反射对称，八个卦限面积相同，总面积为 \(4\pi\)，故此块面积为 \(\pi/2\)。单位球面 \(K=1\)，曲率积分也为 \(\pi/2\)，与直接平行移动计算一致。该闭盘可放入一张避开其他球面点的坐标图，因此局部环路平行移动公式的条件满足；角点只分段改变行进方向，并不额外转动被移动的向量。
\end{solution}

\begin{exercise}\label{g1:ex:flat-rectangle-frame}
	设定向光滑度量曲面的一个坐标域为开矩形 \(U=(a,b)\times(c,d)\)，且 \(K\equiv0\)。不用一般的 Poincar\'e 引理，证明 \(U\) 上存在光滑正交单位标架，使两根标架向量沿任意曲线均平行。再证明任意两组这样的正向标架相差一个常旋转。
\end{exercise}
\begin{solution}
	把坐标基正交单位化，设所得联络形式为 \(\omega=A\mathrm{d}u+B\mathrm{d}v\)。曲率为零给出 \(B_u=A_v\)。固定 \((u_0,v_0)\in U\)，令
	\[
		h(u,v)=\int_{u_0}^uA(s,v_0)\,\mathrm{d}s+\int_{v_0}^vB(u,t)\,\mathrm{d}t.
	\]
	则 \(h_v=B\)，且
	\[
		h_u=A(u,v_0)+\int_{v_0}^v B_u(u,t)\,\mathrm{d}t
		=A(u,v_0)+\int_{v_0}^v A_v(u,t)\,\mathrm{d}t=A(u,v).
	\]
	所以 \(\mathrm{d}h=\omega\)。把原标架旋转角 \(-h\)，\cref{g1:prop:connection-frame-change}使新联络形式成为 \(\omega-\mathrm{d}h=0\)。两根新标架向量的协变导数因而在每个方向上都为零，沿任何曲线都平行。

	若 \((e_1,e_2)\) 与 \((f_1,f_2)\) 都有此性质，则对每个切向量 \(X\)，度量相容性给出 \(X\langle e_i,f_j\rangle=0\)。矩形连通，故四个内积都是常数。换基矩阵因此为常正交矩阵，两组标架均正向又使其行列式为 \(1\)，所以是常旋转。
\end{solution}

\begin{exercise}
	\label{g1:ex:cylinder-geodesics}
	在半径 \(R>0\) 的圆柱面上使用局部参数
	\[
		X(u,z)=(R\cos(u/R),R\sin(u/R),z).
	\]
	求出全部仿射参数测地线，判定哪些为闭曲线，并说明为什么闭合的一整圈不能是同端点之间的最短路。
\end{exercise}
\begin{solution}
	有 \(\langle X_u,X_u\rangle=\langle X_z,X_z\rangle=1\)、\(\langle X_u,X_z\rangle=0\)，所以度量为 \(\mathrm{d}u^2+\mathrm{d}z^2\)，全部 Christoffel 系数为零。测地线方程成为 \(\ddot u=\ddot z=0\)，因此
	\[
		\gamma(t)=\left(R\cos\frac{at+b}{R},R\sin\frac{at+b}{R},ct+d\right).
	\]
	这些公式在全部实数时间上有定义，并与每个局部初值相符，由初值唯一性给出全部测地线。\(a=0,c\ne0\) 时为竖直母线；\(a\ne0,c=0\) 时为水平圆；\(a,c\) 都非零时为圆柱螺线；\(a=c=0\) 时为常值曲线。非常值曲线闭合当且仅当 \(c=0\)，此时可取周期 \(2\pi R/|a|\)。走完一整圈的长度为 \(2\pi R>0\)，同端点之间的常值曲线长度为零，所以整圈不是最短路；这与其每个足够短的弧段为最短路并不冲突。
\end{solution}

\begin{exercise}
	\label{g1:ex:latitude-geodesic-curvature}
	半径 \(R>0\) 的球面取外向定向。对固定 \(0<\theta<\pi\)，按 \(\phi\) 增大的方向走纬线
	\[
		\gamma_\theta(\phi)=R(\sin\theta\cos\phi,\sin\theta\sin\phi,\cos\theta),\qquad 0\le\phi\le2\pi.
	\]
	计算其测地曲率，并用第一变分重新得到 \(\frac{\mathrm{d}}{\mathrm{d}\theta}L(\gamma_\theta)\)。这里闭曲线变分两端相接，端点向量不为零，需说明边界项为何仍相消。
\end{exercise}
\begin{solution}
	单位切向量和左侧单位向量分别为
	\begin{align*}
		T&=(-\sin\phi,\cos\phi,0),\\
		JT=n\times T&=(-\cos\theta\cos\phi,-\cos\theta\sin\phi,\sin\theta).
	\end{align*}
	由于 \(\mathrm{d}\ell=R\sin\theta\,\mathrm{d}\phi\)，空间中有
	\[
		\frac{\mathrm{d}T}{\mathrm{d}\ell}=\frac{(-\cos\phi,-\sin\phi,0)}{R\sin\theta}.
	\]
	与 \(JT\) 内积时法向投影不影响结果，故 \(k_g=\cot\theta/R\)。特别地，赤道 \(\theta=\pi/2\) 才是纬线族中的测地线。

	让 \(\theta\) 增大所给出的变分向量为 \(V=\partial_\theta\gamma_\theta=-RJT\)。在 \(\phi=0,2\pi\) 两端，\(V\) 和 \(T\) 都相同，故边界项 \([\langle V,T\rangle]_0^{2\pi}\) 相消。由\eqref{g1:eq:length-first-variation}，
	\[
		\frac{\mathrm{d}L}{\mathrm{d}\theta}
		=-\int_{\gamma_\theta}k_g\langle V,JT\rangle\,\mathrm{d}\ell
		=\frac{\cos\theta}{\sin\theta}\,(2\pi R\sin\theta)
		=2\pi R\cos\theta.
	\]
	这与直接对 \(L=2\pi R\sin\theta\) 求导一致。北半球 \(\theta<\pi/2\) 中，增大 \(\theta\) 是向纬线的右侧移动，长度增加；符号同时核对了 \(JT\) 的方向。
\end{solution}

\begin{exercise}
	\label{g1:ex:revolution-geodesic-integral}
	设旋转曲面由 \(X(s,\theta)=(r(s)\cos\theta,r(s)\sin\theta,z(s))\) 给出，其中 \(r>0\)、\(r'^2+z'^2=1\)。证明其仿射参数测地线满足 \(r(s)^2\dot\theta=\text{常数}\)。判定参数曲线 \(\theta=\theta_0\) 和 \(s=s_0\) 何时在适当参数下为测地线，并解释守恒量与旋转对称的关系。
\end{exercise}
\begin{solution}
	度量为 \(\mathrm{d}s^2+r(s)^2\mathrm{d}\theta^2\)。由\eqref{g1:eq:christoffel-metric}，非零的 Christoffel 系数至多为
	\[
		\Gamma^s_{\theta\theta}=-rr',\qquad
		\Gamma^\theta_{s\theta}=\Gamma^\theta_{\theta s}=\frac{r'}r.
	\]
	测地线方程于是为
	\[
		\ddot s-rr'\dot\theta^2=0,\qquad
		\ddot\theta+2\frac{r'}r\dot s\dot\theta=0.
	\]
	第二式乘 \(r^2\)，正是 \(\frac{\mathrm{d}}{\mathrm{d}t}(r(s)^2\dot\theta)=0\)。当 \(\theta=\theta_0\) 时，只需令 \(s=at+b\)，所以每条子午线都是测地线轨迹。对于纬线 \(s=s_0\)，第一式成为 \(r(s_0)r'(s_0)\dot\theta^2=0\)。非常值纬线满足 \(\dot\theta\ne0\)，故它为测地线当且仅当 \(r'(s_0)=0\)；此时第二式要求 \(\theta\) 匀速变化。

	旋转把 \(\theta\) 变为 \(\theta+\alpha\)，而度量系数完全不变。守恒量可写成 \(\langle\dot\gamma,\partial_\theta\rangle=r^2\dot\theta\)，记录速度在旋转方向上的内积分量。这里的守恒不是仅凭对称直觉宣布的：上面对测地线方程的求导已经直接证明它。若曲线为单位速，且与子午线方向的有向角为 \(\beta\)，则 \(\sin\beta=r\dot\theta\)，守恒式也写为 \(r\sin\beta=\text{常数}\)。
\end{solution}

\begin{exercise}
	\label{g1:ex:length-energy-parametrization}
	在 Euclid 平面中取 \(\gamma(t)=(t+t^2,0)\)、\(0\le t\le1\)。证明它在固定端点下使长度驻定，但不使能量驻定；给出一个固定端点变分，使能量的一阶导数不为零。再证明在连接 \((0,0)\)、\((2,0)\) 且定义于 \([0,1]\) 的分片 \(C^1\) 曲线中，能量最小值为 \(2\)，并求出达到最小值的曲线。
\end{exercise}
\begin{solution}
	曲线速度 \(\dot\gamma=(1+2t,0)\) 始终向右，所以单位切向量 \(T=(1,0)\) 恒定，\(\mathrm{D}_tT=0\)。固定端点时，长度第一变分的边界项和积分项都为零。但 \(\ddot\gamma=(2,0)\ne0\)，所以能量不驻定。更直接地，取
	\[
		F(s,t)=(t+t^2+s\,t(1-t),0),\qquad V=(t(1-t),0).
	\]
	端点保持不变，而
	\[
		\left.\frac{\mathrm{d}E}{\mathrm{d}s}\right|_{s=0}
		=-2\int_0^1t(1-t)\,\mathrm{d}t=-\frac13.
	\]
	这条曲线的长度确为 \(2\)，但能量为 \(\frac12\int_0^1(1+2t)^2\,\mathrm{d}t=13/6\)。

	对任意竞争曲线 \(\eta\)，有 \(L(\eta)\ge2\)，故\cref{g1:prop:length-energy}给出 \(E(\eta)\ge L(\eta)^2/2\ge2\)。等号要求沿两端之间的线段单向行进且等速。也可直接对向量 \(\dot\eta\) 使用积分的 Cauchy--Schwarz 不等式：
	\[
		4=\left|\int_0^1\dot\eta(t)\,\mathrm{d}t\right|^2
		\le\int_0^1|\dot\eta(t)|^2\,\mathrm{d}t=2E(\eta).
	\]
	此处等号要求 \(\dot\eta\) 几乎处处为同一向量，端点条件强制该向量是 \((2,0)\)。积分后唯一得到 \(\eta(t)=(2t,0)\)，其能量恰为 \(2\)。
\end{solution}

\begin{exercise}
	\label{g1:ex:concave-boundary-angles}
	在定向平面中，按顺序连接
	\[
		(0,0),\ (2,0),\ (2,1),\ (1,1),\ (1,2),\ (0,2),\ (0,0).
	\]
	证明这是一条正向简单闭多边形边界，求六个内角和有向外角，并用局部 Gauss–Bonnet 公式核验。若把每个外角误取为两个前进切向量之间的非负夹角，会得到什么错误结果？
\end{exercise}
\begin{solution}
	该多边形围成集合 \(D=([0,2]\times[0,1])\cup([0,1]\times[0,2])\)，各条边除相邻端点外互不相交。逐边看去，\(D\) 都在左侧，故方向为正。顶点 \((1,1)\) 处，前进方向由向左变为向上，是向右转 \(\pi/2\)；这里内部占据另外三个象限，内角为 \(3\pi/2\)，外角为 \(-\pi/2\)。其余五个顶点都向左转 \(\pi/2\)，内角与外角均为 \(\pi/2\)。

	平面 \(K=0\)，每条直线边的 \(k_g=0\)，而外角之和为 \(5\pi/2-\pi/2=2\pi\)，正好满足公式；内角之和为 \(5\pi/2+3\pi/2=4\pi=(6-2)\pi\)。若六处都只取非负夹角 \(\pi/2\)，就会算得总转角 \(3\pi\)，甚至不再是 \(2\pi\) 的整数倍。错误在于抹掉了凹角处向右转的负号。
\end{solution}

\begin{exercise}
	\label{g1:ex:spherical-cap-boundary}
	半径 \(R>0\) 的球面取外向定向。固定 \(0<\varphi_0<\pi\)，令 \(D\) 是余纬 \(0\leq\varphi\leq\varphi_0\) 的北球冠，其边界为
	\[
		\gamma(\theta)=R(\sin\varphi_0\cos\theta,\sin\varphi_0\sin\theta,\cos\varphi_0),
		\qquad 0\leq\theta\leq2\pi.
	\]
	验证参数增大的方向是球冠的正向边界，计算其测地曲率、边界积分与球冠面积，核验 Gauss–Bonnet 公式。解释 \(\varphi_0>\pi/2\) 时边界项为何为负。
\end{exercise}
\begin{solution}
	沿边界取单位切向量和子午方向单位向量
	\begin{align*}
		T&=(-\sin\theta,\cos\theta,0),\\
		e_\varphi&=(\cos\varphi_0\cos\theta,\cos\varphi_0\sin\theta,-\sin\varphi_0).
	\end{align*}
	有 \(e_\varphi\times T=n=\gamma/R\)，所以 \(JT=-e_\varphi\)，指向余纬减小的球冠内部；边界方向确为正。由 \(\mathrm{d}s=R\sin\varphi_0\,\mathrm{d}\theta\)，空间导数为
	\[
		\frac{\mathrm{d}T}{\mathrm{d}s}
		=-\frac{(\cos\theta,\sin\theta,0)}{R\sin\varphi_0}.
	\]
	与切向量 \(JT\) 作内积时，法向分量不作贡献，因此
	\[
		k_g=\left\langle\frac{\mathrm{d}T}{\mathrm{d}s},-e_\varphi\right\rangle
		=\frac{\cot\varphi_0}{R},\qquad
		\int_{\partial D}k_g\,\mathrm{d}s=2\pi\cos\varphi_0.
	\]
	球面面积元是 \(R^2\sin\varphi\,\mathrm{d}\varphi\,\mathrm{d}\theta\)，所以
	\[
		\operatorname{Area}(D)=2\pi R^2(1-\cos\varphi_0),\qquad
		\int_DK\,\mathrm{d}A=2\pi(1-\cos\varphi_0).
	\]
	两积分相加为 \(2\pi\)。球冠避开南极，包含在删去南极的球面坐标图中，且无角点，故这里适用局部公式。这个图就是\cref{g1:prop:sphere-simply-connected}证明中的球极投影：从南极投向赤道平面，半径 \(R\) 时写为 \((x,y,z)\mapsto(Rx/(R+z),Ry/(R+z))\)；那里给出的有理逆映射经缩放后适用，正反映射的分母均非零，因而也是光滑坐标图。

	当 \(\varphi_0>\pi/2\) 时，球冠超过半球，曲率面积积分大于 \(2\pi\)，故边界积分必须为负。直接计算也显示 \(\cot\varphi_0<0\)：边界加速度的切向分量指向余纬增大的一侧，即球冠外侧。负号描述的是在曲面内的弯向，与这条空间圆周的非负曲率不是同一件事。
\end{solution}

\begin{exercise}
	\label{g1:ex:nonpositive-geodesic-bigon}
	设 \(D\) 是有向光滑曲面上一张坐标图内的紧闭圆盘区域，\(K\leq0\) 在 \(D\) 上成立。
	\begin{enumerate}
		\item 若 \(\partial D\) 由 \(m\) 条测地线弧组成，各内角 \(\beta_j\in(0,2\pi)\)，证明 \(\sum_j\beta_j\leq(m-2)\pi\)，并说明何时取等号。
		\item 证明两条测地线弧不能以两个不同公共端点围成这样的圆盘，并在两个端点具有非零且非整周的内角。证明 \(D\) 的整条光滑边界也不可能是闭测地线。
	\end{enumerate}
\end{exercise}
\begin{solution}
	1. 每条测地边上 \(k_g=0\)，由\cref{g1:thm:local-gauss-bonnet}，
	\[
		\sum_{j=1}^m\beta_j=(m-2)\pi+\int_DK\,\mathrm{d}A\leq(m-2)\pi.
	\]
	等号成立当且仅当 \(\int_DK\,\mathrm{d}A=0\)。由于 \(K\) 连续且非正，若内点处 \(K<0\)，某个正面积小邻域的积分就严格为负；故等号迫使所有内点 \(K=0\)，再由连续性得到 \(K\) 在闭区域上恒为零。反之 \(K\equiv0\) 显然给出等号。

	2. 两条边的情形是 \(m=2\)，第一问会给出 \(\beta_1+\beta_2\leq0\)，与两个内角都为正矛盾。若整条边界是光滑闭测地线，则没有角点、\(k_g=0\)，公式给出 \(\int_DK\,\mathrm{d}A=2\pi\)，同样矛盾。这些结论只排除题设的圆盘边界；例如平圆柱上的闭测地线不围圆盘，并不受此排除。
\end{solution}

\begin{exercise}
	\label{g1:ex:conformal-boundary-gauss-bonnet}
	设 \(D\subset\mathbb R^2\) 是具有 \(C^2\) 正则简单闭边界的紧圆盘区域，\(\varphi\) 在其邻域光滑。在该邻域上给定度量 \(g=e^{2\varphi}(\mathrm{d}x^2+\mathrm{d}y^2)\)。记 \(s_0\)、\(t\)、\(k_0\) 分别为正向边界的 Euclid 弧长、单位切向量和有向曲率，\(\nu=Jt\) 为 Euclid 向内单位法向。
	\begin{enumerate}
		\item 用标架 \(e_1=e^{-\varphi}\partial_x\)、\(e_2=e^{-\varphi}\partial_y\) 求联络一形式，并证明
		\[
			k_g\,\mathrm{d}s_g=(k_0-\partial_\nu\varphi)\,\mathrm{d}s_0.
		\]
		\item 用 \(K\,\mathrm{d}A=-(\varphi_{xx}+\varphi_{yy})\,\mathrm{d}x\,\mathrm{d}y\) 与平面 Green 公式直接证明 \(\int_DK\,\mathrm{d}A+\int_{\partial D}k_g\,\mathrm{d}s_g=2\pi\)，说明度量改变时哪两项互相补偿。
	\end{enumerate}
\end{exercise}
\begin{solution}
	1. 将 \(g_{11}=g_{22}=e^{2\varphi}\)、\(g_{12}=0\) 代入\eqref{g1:eq:christoffel-metric}，得到
	\begin{align*}
		\nabla_{\partial_x}e_1&=-\varphi_y e_2,&
		\nabla_{\partial_y}e_1&=\varphi_x e_2.
	\end{align*}
	例如 \(\nabla_{\partial_x}\partial_x=\varphi_x\partial_x-\varphi_y\partial_y\)，对系数 \(e^{-\varphi}\) 求导后，\(\partial_x\) 分量恰相消；另一式同理。因此
	\[
		\omega=-\varphi_y\,\mathrm{d}x+\varphi_x\,\mathrm{d}y.
	\]
	写 \(t=(\cos\theta,\sin\theta)\)，则度量单位切向量为 \(T=e^{-\varphi}t=\cos\theta\,e_1+\sin\theta\,e_2\)，所以两种度量下相对于相应标架的角度相同，而 \(\mathrm{d}s_g=e^\varphi\mathrm{d}s_0\)。由\eqref{g1:eq:tangent-connection-angle}，
	\[
		k_g\,\mathrm{d}s_g=\mathrm{d}\theta+\omega(t)\,\mathrm{d}s_0.
	\]
	Euclid 曲率给出 \(\mathrm{d}\theta=k_0\,\mathrm{d}s_0\)，且
	\[
		\omega(t)=-\varphi_y\cos\theta+\varphi_x\sin\theta
		=-\partial_\nu\varphi,
	\]
	因为 \(\nu=(-\sin\theta,\cos\theta)\)。所求边界公式成立。

	2. 向外单位法向是 \(-\nu\)。将\cref{g1:lem:green-formula}用于一形式 \(-\varphi_y\,\mathrm{d}x+\varphi_x\,\mathrm{d}y\)，可得
	\[
		\int_D(\varphi_{xx}+\varphi_{yy})\,\mathrm{d}x\,\mathrm{d}y
		=-\int_{\partial D}\partial_\nu\varphi\,\mathrm{d}s_0.
	\]
	因而面积项为 \(\int_{\partial D}\partial_\nu\varphi\,\mathrm{d}s_0\)，与第一问的边界项中的法向导数部分恰好抵消。剩下
	\[
		\int_DK\,\mathrm{d}A+\int_{\partial D}k_g\,\mathrm{d}s_g
		=\int_{\partial D}k_0\,\mathrm{d}s_0=2\pi,
	\]
	最后一式是\cref{g1:thm:turning-tangent}。共形度量改变了内部的曲率面积积分，也改变了边界的测地曲率积分；边界法向导数具体说明了两者如何补偿。
\end{solution}

\begin{exercise}\label{g1:ex:gb-finite-cover}
	设 \(M\) 是紧连通光滑度量曲面，允许有分片光滑边界及有限个非零内角，\(p:\widetilde M\to M\) 是有限 \(m\) 叶光滑覆盖。给 \(\widetilde M\) 配备拉回度量 \(p^*g\)。不用三角形提升计数，利用 Gauss--Bonnet 公式证明 \(\chi(\widetilde M)=m\chi(M)\)，并说明不可定向情形为何不改变证明。
\end{exercise}
\begin{solution}
	有限叶覆盖的总空间紧，可用\cref{g1:prop:surface-finite-cover-euler}证明的第一段构造直接看出：用有限个闭包位于均匀覆盖坐标图内的小邻域覆盖 \(M\)，每个小闭包的原像是 \(m\) 个紧副本，其有限并覆盖 \(\widetilde M\)。下面只使用这一紧性结论，不使用该命题的 Euler 计数结论。

	在均匀覆盖的每一片上，\(p\) 是局部等距，故保留 \(K\)、面积测度、边界的内侧单位法向以及内角。取有限个这样的图，在它们内把积分区域分成互不重叠的可积小块；每块有 \(m\) 个等距副本，所以
	\[
		\int_{\widetilde M}\widetilde K\,\mathrm{d}\widetilde A=m\int_MK\,\mathrm{d}A,
		\qquad
		\int_{\partial\widetilde M}\widetilde k_g\,\mathrm{d}\widetilde s=m\int_{\partial M}k_g\,\mathrm{d}s.
	\]
	每个角点也恰有 \(m\) 个原像，且其外角相同，故外角和同样乘以 \(m\)。两边分别应用\cref{g1:thm:global-gauss-bonnet}，得到 \(2\pi\chi(\widetilde M)=2\pi m\chi(M)\)。面积测度、曲率和内侧测地曲率都不要求全局定向，局部等距也可以反转定向，所以证明同时适用于不可定向曲面。
\end{solution}

\begin{exercise}\label{g1:ex:gb-cutting}
	设 \(M\) 是紧光滑度量曲面，边界光滑。对下列两种切割，用 Gauss--Bonnet 公式证明相应的 Euler 示性数变化。
	\begin{enumerate}
		\item 沿内部一条光滑嵌入闭曲线 \(C\) 切开，所得曲面为 \(M_C\)。证明 \(\chi(M_C)=\chi(M)\)，并同时处理 \(C\) 为单侧曲线的情形。
		\item 沿一条光滑嵌入弧 \(A\) 切开，假设其两个不同端点在 \(\partial M\) 上，内部在 \(M^\circ\) 中，且两端都与边界横交。所得带角点曲面为 \(M_A\)。证明 \(\chi(M_A)=\chi(M)+1\)。
	\end{enumerate}
\end{exercise}
\begin{solution}
	1. 切开后每个 \(C\) 的小弧都有两个边界副本，二者的内侧单位法向互为相反数。协变加速度相同，故两侧的 \(k_g\) 相反，积分抵消。若 \(C\) 为双侧曲线，这两个副本分别接成两条边界；若为单侧曲线，它们接成一条二叶覆盖 \(C\) 的边界。后一情形仍可先沿 \(C\) 分成有限小弧，再将两侧积分配对，总和依然为零。切割不改变面积积分，也不产生角点，原有边界的积分保持不变。比较 \(M_C\) 与 \(M\) 的 Gauss--Bonnet 公式，即得 \(\chi(M_C)=\chi(M)\)。

	2. 弧内部的两侧积分仍然抵消，面积积分和原有边界各段也不变。在弧的一个端点处，原来的内角为 \(\pi\)，切割将其分为 \(\theta\) 与 \(\pi-\theta\)，其中横交保证 \(0<\theta<\pi\)。这两个新角的外角和是
	\[
		(\pi-\theta)+\bigl(\pi-(\pi-\theta)\bigr)=\pi.
	\]
	两个端点共增加 \(2\pi\)。因此 \(M_A\) 的 Gauss--Bonnet 左端比 \(M\) 增加 \(2\pi\)，右端也增加 \(2\pi\)，故 \(\chi(M_A)=\chi(M)+1\)。无论切割后是否连通，这个计算均成立。
\end{solution}

\begin{exercise}\label{g1:ex:gb-cone-points}
	设 \(M\) 是紧拓扑曲面，除有限个内部点 \(q_1,\ldots,q_r\) 外具有光滑度量，边界分片光滑且与这些点不交。假设每个 \(q_i\) 的穿孔邻域等距于
	\[
		0<\rho<\rho_i,\qquad \theta\in\mathbb R/2\pi\mathbb Z,
		\qquad g=\mathrm{d}\rho^2+a_i^2\rho^2\,\mathrm{d}\theta^2,
	\]
	其中 \(a_i>0\)。称 \(\Theta_i=2\pi a_i\) 为该点的锥角。证明
	\[
		\int_{M\setminus\{q_1,\ldots,q_r\}}K\,\mathrm{d}A
		+\int_{\partial M}k_g\,\mathrm{d}s+\sum_{v\in\partial M}\alpha_v
		+\sum_{i=1}^r(2\pi-\Theta_i)=2\pi\chi(M).
	\]
	其中边界角点仍假定内角属于 \((0,2\pi)\)。
\end{exercise}
\begin{solution}
	在每个锥点附近，局部改用角变量 \(\varphi=a_i\theta\) 就把度量写成平面极坐标度量，所以穿孔邻域的 \(K=0\)。去掉互不相交的小开圆盘 \(\rho<\varepsilon\)，得到光滑度量曲面 \(M_\varepsilon\)。由\cref{g1:prop:surface-euler-formulas}中封盘增加 \(1\) 的结论，\(\chi(M_\varepsilon)=\chi(M)-r\)。

	新边界 \(\rho=\varepsilon\) 的单位切向量可取 \(T=(a_i\varepsilon)^{-1}\partial_\theta\)。由度量系数直接计算 \(\Gamma^\rho_{\theta\theta}=-a_i^2\rho\)，所以
	\[
		\nabla_TT=-\frac1\varepsilon\partial_\rho.
	\]
	\(M_\varepsilon\) 的内侧在 \(\rho\) 增大的方向，故 \(\nu=\partial_\rho\)，\(k_g=-1/\varepsilon\)。圆周长度为 \(2\pi a_i\varepsilon\)，其边界积分就是 \(-2\pi a_i=-\Theta_i\)。这些新边界没有角点；面积曲率积分在删去的穿孔邻域内为零，因此它不依赖足够小的 \(\varepsilon\)。对 \(M_\varepsilon\) 应用\cref{g1:thm:global-gauss-bonnet}，得到
	\[
		\int K\,\mathrm{d}A+\int_{\partial M}k_g\,\mathrm{d}s
		+\sum_{v\in\partial M}\alpha_v-\sum_{i=1}^r\Theta_i
		=2\pi\bigl(\chi(M)-r\bigr).
	\]
	移项即为所求。这里的 \(2\pi-\Theta_i\) 是集中在锥点处的角亏；该点不必具有通常意义下的光滑曲率函数。
\end{solution}

\begin{exercise}\label{g1:ex:negative-curvature-area-bounds}
	设 \(M\) 为亏格 \(g\ge2\) 的紧连通可定向闭度量曲面，且 \(-b\le K\le-a<0\)，其中 \(0<a\le b\)。证明
	\[
		\frac{4\pi(g-1)}b\le\operatorname{Area}(M)\le\frac{4\pi(g-1)}a,
	\]
	并给出两端分别取等号的充要条件。
\end{exercise}
\begin{solution}
	置 \(A=\operatorname{Area}(M)\)。对 \(a\le-K\le b\) 积分，并用\eqref{g1:eq:total-curvature-genus}，有
	\[
		aA\le-\int_MK\,\mathrm{d}A=4\pi(g-1)\le bA.
	\]
	除以正数 \(a,b\) 即得面积界。下界取等号当且仅当 \(\int_M(b+K)\,\mathrm{d}A=0\)。被积函数连续且非负，若某点为正则其小邻域积分为正，所以等价于 \(K\equiv-b\)。同理，上界取等号当且仅当 \(K\equiv-a\)。这里的等号结论是针对给定曲面度量的判定，不另外断言任意面积都能实现。
\end{solution}

\begin{exercise}\label{g1:ex:degenerate-vector-indices}
	把平面的点和向量分别用复数 \(z=x+\mathrm{i}y\) 表示；这只是实数对的简写。
	\begin{enumerate}
		\item 对整数 \(m\ge1\)，求向量场 \(V(z)=z^m\) 与 \(W(z)=\overline z^{,m}\) 在原点的指数。
		\item 把球面写成两张复坐标图，重叠处用 \(w=1/z\) 换坐标。证明在 \(z\) 图上给定常速度 \(\dot z=1\)，可光滑延拓为球面上的切向量场，恰有一个零点，并求其指数。
	\end{enumerate}
\end{exercise}
\begin{solution}
	1. 沿正向圆周 \(z=re^{\mathrm{i}\theta}\)，两个单位化向量分别为 \(e^{\mathrm{i}m\theta}\) 与 \(e^{-\mathrm{i}m\theta}\)，故指数为 \(m\) 与 \(-m\)。这也表明退化零点的指数并不限于 \(\pm1\)。

	2. 先说明这套坐标确实描述标准光滑球面。可取 \(z=(X+\mathrm{i}Y)/(1-Z)\) 为去掉北极的立体投影坐标，另一图取 \(w=(X-\mathrm{i}Y)/(1+Z)\)；在单位球面重叠处，由 \(X^2+Y^2=1-Z^2\) 得 \(zw=1\)。两图的坐标函数和逆映射都光滑。

	切向量按坐标映射的微分换分量，所以在另一张图上
	\[
		\dot w=-z^{-2}\dot z=-w^2.
	\]
	右端是 \(\Re w,\Im w\) 的二次多项式，在 \(w=0\) 光滑且为零；在其余位置非零。因此所得球面场恰有北极一个零点。围绕 \(w=0\) 一周，\(-w^2\) 的角度增加 \(4\pi\)，指数为 \(2\)。使用坐标基而非正交标架不改变此计算：把坐标基正交化所产生的正对角上三角换基矩阵可经 \((1-t)A+tI\) 在可逆矩阵中变到恒等，因而单位化向量的绕数不变。总指数为 \(2=\chi(S^2)\)，与\cref{g1:thm:surface-poincare-hopf}相符；该定理并不要求球面场至少有两个零点。
\end{solution}

\begin{exercise}\label{g1:ex:conformal-torus-sign}
	设 \(f(u,v)\) 为两个变量均以 \(2\pi\) 为周期的光滑实函数。在环面上赋予度量 \(\widetilde g=e^{2f}(\mathrm{d}u^2+\mathrm{d}v^2)\)。证明：若其 Gauss 曲率在全曲面上不变号，则 \(f\) 为常数。因此任何非常数 \(f\) 都使新度量的曲率同时取正值和负值。
\end{exercise}
\begin{solution}
	环面 \(\chi=0\)，由\cref{g1:prop:curvature-sign-topology}，不变号的曲率必恒为零。\eqref{g1:eq:conformal-curvature-change}又给出
	\[
		0=\widetilde K e^{2f}=-(f_{uu}+f_{vv}).
	\]
	在周期矩形 \([0,2\pi]^2\) 上乘以 \(f\) 并分部积分；相对两边的函数值及其一阶导数相等，边界项成对抵消，故
	\[
		0=-\int_0^{2\pi}\int_0^{2\pi}f(f_{uu}+f_{vv})\,\mathrm{d}u\,\mathrm{d}v
		=\int_0^{2\pi}\int_0^{2\pi}(f_u^2+f_v^2)\,\mathrm{d}u\,\mathrm{d}v.
	\]
	连续非负函数的积分为零，使 \(f_u=f_v=0\)，所以 \(f\) 在连通平面上为常数。反之常数 \(f\) 只把平坦度量按常数缩放，曲率仍为零。若 \(f\) 非常数，曲率便不能非负，也不能非正，因而必兼有正负值。
\end{solution}

% !TeX root = ../../Book1.tex
% 纲要标题：## 第10章　非 Euclid 模型、离散作用与几何结构
% 纲要位置：几何学计划纲要.md，第 501 行。
% 正文状态：初稿；尚未正式审核。

% 纲要内容开始
% 
% > 本章目标与证明范围：保留球面与双曲几何的独立深入模块。先计算模型，再研究离散群商与局部齐性结构；一般统一化是第四卷的后续任务，不从模型存在直接推出。
% 
% 纲要内容结束

\chapter{非 Euclid 模型、离散作用与几何结构}
\label{g1:ch:10}
\BookChapterNumber{g1:ch:10}

\begin{chapterabstract}
	常曲率曲面既可以用具体公式计算，也可以作为局部几何的标准模型。本章构造双曲面、圆盘和上半平面之间的等距映射，求距离、测地线、等距群及球面与双曲三角恒等式，并把角和与面积相联系。随后从等距变换的分类研究离散作用，证明自由固有作用的商结构，按边和顶点的实际条件粘合基本多边形，构造闭双曲曲面、尖点与漏斗。最后沿道路延续几何坐标，建立展开映射与单值表示，并证明完备模型结构与模型商之间的关系。
\end{chapterabstract}

\begin{learninggoals}
	\begin{itemize}
		\item 在三种双曲模型之间明确转换点、切向量、度量和测地线，区分 Euclid 图形与模型中的长度、角度。
		\item 推导球面和双曲三角公式，利用曲率归一化计算面积，判定正则测地多边形的角条件。
		\item 用固定点、矩阵迹和标准型区分椭圆、抛物与双曲变换，判断离散性、自由性和固有性的作用。
		\item 检查边配对的长度匹配、顶点链接和总角，构造具有光滑双曲度量的闭曲面及其典型非紧端。
		\item 从几何图册构造展开映射，核对单值表示的乘法顺序与等变性，说明完备性如何保证整体商的实现。
	\end{itemize}
\end{learninggoals}

在 Euclid 平面里，三角形的三个角之和总是 \(\pi\)。球面上的短大圆三角形却有更大的角和；第九章说明，多出来的角恰好记录它围住的曲率积分。那么，角和小于 \(\pi\) 的几何能否同样明确地计算？如果把许多这样的图形按边粘合，怎样保证接缝和顶点附近仍是同一种光滑几何？这些问题把局部的测量公式与整张曲面的构造联系起来。

\ProseParagraphBreak

双曲几何有几种外观差别很大的模型。上半平面中，靠近实轴的一小段 Euclid 距离可能对应很长的双曲距离；圆盘中的测地线通常画成圆弧；双曲面则把距离计算转为一个带负号的双线性型。证明模型间的等距，才能在最便于眼前计算的模型中工作。它们的理想边界也需要与通常的度量边界分清：图中看得见的边缘，并不是走有限长度便能抵达的缺口。

\ProseParagraphBreak

研究商空间时，局部模型仍不够。一个群可以离散却有固定点，配对的边也可以等长却在顶点留下锥形奇点。我们逐项核对这些条件，再将问题反过来：如果只知道每一小块都具有同样的几何，能否把整张曲面放回一个标准模型中？展开映射先回答万有覆盖上的延续问题；单值表示记录绕路后的坐标变化；完备性则负责排除有限长度处中断的情形。这里的结构单值与第九章切向量的环路平行移动变换有不同的对象和作用，使用时会加以区分。

\ProseParagraphBreak

本章使用有限维线性代数、多元微积分与复数运算，复数只作为平面坐标及分式变换的记法。第五章的覆盖与道路拼接、第六章的曲面边配对、第九章的测地邻域与 Gauss--Bonnet 公式在调用处恢复所需条件。进一步阅读可查 Ratcliffe 第三版：球面、双曲模型及三角计算见\cite[\S\S~2.3--2.5, 3.1--3.5, 4.3--4.6]{Ratcliffe2019}；离散作用、基本区域和曲面粘合见\cite[\S\S~4.7, 5.2--5.3, 6.6, 9.2, 9.7--9.8]{Ratcliffe2019}；展开与完备性见\cite[\S\S~8.3--8.5]{Ratcliffe2019}。本章构造具体的常曲率结构，并证明给定完备结构的商描述；一般统一化仍属于第四卷的后续内容。

% !TeX root = ../../Book1.tex
% 纲要标题：### 10.1 常曲率模型的度量几何
% 纲要位置：几何学计划纲要.md，第 505 行。
% 正文状态：初稿；尚未正式审核。

% 纲要内容开始
% 
% * 双曲面、圆盘和上半平面
% * 模型间等距映射的构造
% * 距离、测地线及理想边界
% * Möbius 变换与等距群
% * 球面与双曲三角恒等式
% * 角和、面积及曲率归一化
% * 测地多边形和角条件
% 
% 纲要内容结束

\section{常曲率模型的度量几何}
\label{g1:sec:1001}

同样是圆盘，在通常的欧氏度量下，走到边缘只需有限长度；若越靠近边缘，度量把长度放大得越快，边缘就可能变得无限遥远。双曲平面的圆盘模型正是这样的例子。为了看清距离、测地线和对称变换，我们还会把同一几何放在双曲面与上半平面中表达。三个模型各有便于计算的地方，彼此间的显式等距映射使这些计算能够相互转用。

\subsection{双曲面、圆盘和上半平面}
\label{g1:subsec:1001:01}

先取单位球面 \(S^2\subset\mathbb R^3\) 及其诱导度量，它的曲率为 \(1\)。我们寻找一个负曲率的对应模型。把通常内积中的一项改成负号，定义 \(\mathbb R^3\) 上的双线性型
\[
	\langle x,y\rangle_L\coloneqq-x_0y_0+x_1y_1+x_2y_2.
\]
这是\term[lorentzxing]{Lorentz 型}[Lorentzian form]，不是正定内积；例如 \(\langle(1,0,0),(1,0,0)\rangle_L=-1\)。我们并不把它在整个 \(\mathbb R^3\) 上当作长度，而只将它限制到下面曲面的切平面。\footnote{Lorentz，中文常译洛伦兹；Hendrik Antoon Lorentz（1853--1928），荷兰理论物理学家，研究电磁理论、电子理论与运动物体的变换。}

\begin{definition}
	\label{g1:def:hyperbolic-models}
	\term[shuangqumianmoxing]{双曲面模型}[hyperboloid model]是
	\[
		\mathbb H_L^2\coloneqq\{x:\langle x,x\rangle_L=-1,\ x_0>0\},
		\qquad g_L(v,w)=\langle v,w\rangle_L.
	\]
	另外定义\term[yuanpanmoxing]{圆盘模型}[disc model]与\term[shangbanpingmianmoxing]{上半平面模型}[upper half-plane model]：
	\begin{align*}
		\mathbb D&=\{w=u+iv:|w|<1\},&
		g_{\mathbb D}&=\frac{4(\mathrm{d}u^2+\mathrm{d}v^2)}{(1-|w|^2)^2},\\
		\mathbb H&=\{z=x+iy:y>0\},&
		g_{\mathbb H}&=\frac{\mathrm{d}x^2+\mathrm{d}y^2}{y^2}.
	\end{align*}
	复数记号 \(|\mathrm{d}z|^2\) 只简写实二次型 \(\mathrm{d}x^2+\mathrm{d}y^2\)。后两个模型采用平面的标准定向。
\end{definition}

双曲面上的度量确实正定。写 \(x=(x_0,x')\)，则 \(x_0=\sqrt{1+|x'|^2}\)，所以它首先是一个光滑图像。切向量 \(v=(v_0,v')\) 满足 \(x_0v_0=x'\cdot v'\)，从而
\[
	\langle v,v\rangle_L=|v'|^2-\frac{(x'\cdot v')^2}{x_0^2}
	\geq\frac{|v'|^2}{x_0^2}>0\qquad(v\ne0).
\]
这里最后一个严格不等式来自：若 \(v'=0\)，切向条件也使 \(v_0=0\)。因此这个不定型在每个切平面上都给出通常意义的正定长度与角度。它与双曲面作为欧氏空间曲面时的诱导度量不同，不能用欧氏空间中该图像的主曲率代替它的内在曲率。

\ProseParagraphBreak

在上半平面上，\(g_{\mathbb H}=e^{2\phi}(\mathrm{d}x^2+\mathrm{d}y^2)\)，其中 \(\phi=-\log y\)。回忆\eqref{g1:eq:conformal-gauss-curvature}，共形度量的曲率为 \(-e^{-2\phi}(\phi_{xx}+\phi_{yy})\)。这里 \(\phi_{xx}=0\)、\(\phi_{yy}=1/y^2\)，故 \(K=-1\)。圆盘与双曲面是否也给出相同几何，需要下一小节的等距计算，而不能只凭外形或名称判断。

\subsection{模型间等距映射的构造}
\label{g1:subsec:1001:02}

从点 \((-1,0,0)\) 向平面 \(x_0=0\) 投影，会把双曲面上半支送到单位圆盘。这一选择中的 \(-1\) 使投影分母始终为正，也让无穷远处趋向圆周。上半平面到圆盘则用一个把实轴变成圆周的分式映射。

\begin{proposition}[三个模型的等距]
	\label{g1:prop:hyperbolic-model-isometries}
	以下映射均为光滑等距微分同胚：
	\begin{align*}
		P:\mathbb H_L^2&\longrightarrow\mathbb D,&
		P(x)&=\frac{x_1+ix_2}{x_0+1},\\
		C:\mathbb H&\longrightarrow\mathbb D,&
		C(z)&=\frac{z-i}{z+i}.
	\end{align*}
	双曲面的定向取为使 \(P\) 保向的定向；\(C\) 也保持定向。因此三个模型都具有常曲率 \(-1\)，统称为\term[shuangqupingmian]{双曲平面}[hyperbolic plane]。
\end{proposition}
\begin{proof}
	先由 \(x_1^2+x_2^2=x_0^2-1\)，得到 \(|P(x)|^2=(x_0-1)/(x_0+1)<1\)。反解为
	\[
		P^{-1}(u+iv)=\frac{(1+u^2+v^2,\,2u,\,2v)}{1-u^2-v^2}.
	\]
	这给出了互逆的光滑映射。记 \(q=1-u^2-v^2\)、\(a=u\,\mathrm{d}u+v\,\mathrm{d}v\)。对逆映射求导，有
	\[
		\mathrm{d}x_0=\frac{4a}{q^2},\qquad
		\mathrm{d}(x_1,x_2)=\frac{2(\mathrm{d}u,\mathrm{d}v)}q+\frac{4(u,v)a}{q^2}.
	\]
	代入 Lorentz 型，含 \(a^2\) 的系数为 \(16/q^3+16(u^2+v^2-1)/q^4=0\)，剩下 \(4(\mathrm{d}u^2+\mathrm{d}v^2)/q^2\)，正是圆盘度量。

	对 \(C\)，直接算得
	\[
		1-|C(z)|^2=\frac{4y}{|z+i|^2},\qquad
		C'(z)=\frac{2i}{(z+i)^2},\qquad C^{-1}(w)=i\frac{1+w}{1-w}.
	\]
	这些式子证明 \(C\) 恰把上半平面双射到圆盘，且正反映射光滑。复导数乘法把实长度放大 \(|C'(z)|\) 倍，故
	\[
		C^*g_{\mathbb D}
		=\frac{4|C'(z)|^2|\mathrm{d}z|^2}{(1-|C(z)|^2)^2}
		=\frac{|\mathrm{d}z|^2}{y^2}.
	\]
	非零复数乘法的实行列式为其模平方，所以 \(C\) 保向。最后，等距保持内在曲率，这是\cref{g1:thm:theorema-egregium}及其度量表达\eqref{g1:eq:intrinsic-gauss-curvature}的结论，故三个模型的曲率都是 \(-1\)。
\end{proof}

这些坐标图的用途不同：双曲面的线性代数适合距离与三角公式，圆盘适合展示对称性，上半平面的分式变换适合研究群作用。圆盘和上半平面的度量都是欧氏度量乘以一个正函数，因而保留欧氏角度；它们并不保留欧氏长度。

\subsection{距离、测地线及理想边界}
\label{g1:subsec:1001:03}

在连通度量曲面上，用全部连接 \(p,q\) 的分片 \(C^1\) 曲线定义长度距离
\[
	d(p,q)\coloneqq\inf_{\gamma:p\to q}L(\gamma).
\]
三角不等式来自道路拼接，反向参数给出对称性。局部坐标中的正定度量在小闭坐标球上有正的下界，故走出这个球需要正长度，这也保证不同点的距离不为零。前章的径向极小定理只在正常邻域内比较短曲线；对双曲平面，可以利用模型把比较推广到任意两点。

\begin{theorem}[双曲距离与测地线]
	\label{g1:thm:hyperbolic-distance-geodesics}
	双曲面模型中，若 \(p\in\mathbb H_L^2\)、\(v\in \mathrm{T}_p\mathbb H_L^2\) 且 \(\langle v,v\rangle_L=1\)，则所有单位速测地线由
	\begin{equation}
		\label{g1:eq:hyperboloid-geodesic}
		\gamma(t)=p\cosh t+v\sinh t,\qquad t\in\mathbb R
	\end{equation}
	给出。任意两个不同点之间有唯一测地线段，并且它在全部分片 \(C^1\) 连接曲线中最短。三个模型中的距离分别满足
	\begin{align}
		\cosh d(p,q)&=-\langle p,q\rangle_L,
		\label{g1:eq:hyperboloid-distance}\\
		\cosh d(w_1,w_2)&=1+\frac{2|w_1-w_2|^2}{(1-|w_1|^2)(1-|w_2|^2)},
		\label{g1:eq:disc-distance}\\
		\cosh d(z_1,z_2)&=1+\frac{|z_1-z_2|^2}{2\operatorname{Im}z_1\operatorname{Im}z_2}.
		\label{g1:eq:half-plane-distance}
	\end{align}
	双曲平面单连通，其长度距离完备，每个闭有界球紧；每条测地线也都能在整个实数轴上延拓。
\end{theorem}
\begin{proof}
	先说明如何把任意点移到 \(o=(1,0,0)\)。空间坐标 \((x_1,x_2)\) 的旋转保持 Lorentz 型，而矩阵
	\[
		B_r=\begin{pmatrix}\cosh r&\sinh r&0\\\sinh r&\cosh r&0\\0&0&1\end{pmatrix}
	\]
	也保持该型，并把 \(o\) 送到 \((\cosh r,\sinh r,0)\)。任意点都能写成
	\[
		x=(\cosh r,\sinh r\cos\theta,\sinh r\sin\theta),
	\]
	所以这些变换在双曲面上传递；它们保持上半支，因为变换后的第零坐标不可能为零，且上半支连通。其微分保持切平面上的度量，故都是等距。

	双曲面的协变导数可以由 Lorentz 正交投影计算：沿 \(\gamma\) 的切向量场 \(W\) 满足
	\[
		\mathrm{D}_tW=\dot W+\langle\dot W,\gamma\rangle_L\gamma.
	\]
	右边确与 \(\gamma\) 正交；逐项求导内积给出度量相容，交换坐标导数给出无挠性，因此\cref{g1:prop:surface-connection-properties}的唯一性说明它就是度量联络。对单位速曲线，\(\langle\ddot\gamma,\gamma\rangle_L=-1\)，故测地线方程等价于 \(\ddot\gamma=\gamma\)。解这个线性方程即得\eqref{g1:eq:hyperboloid-geodesic}；代入可核验它始终位于上半支且速度为 \(1\)。前章\cref{g1:prop:geodesic-initial-values}的唯一性说明没有遗漏同初值的其他测地线。

	以 \(o\) 为中心，上面的径向参数给出
	\begin{equation}
		\label{g1:eq:hyperbolic-polar-metric}
		g_L=\mathrm{d}r^2+\sinh^2r\,\mathrm{d}\theta^2,\qquad r>0.
	\end{equation}
	这是直接微分 \((\cosh r,\sinh r\cos\theta,\sinh r\sin\theta)\) 所得。任何从 \(o\) 到径向坐标 \(r_0\) 的曲线，其速度不小于径向速度的绝对值。为避开 \(r=0\) 的极坐标奇点，取 \(0<\varepsilon<r_0\)，只比较曲线最后一次离开 \(r=\varepsilon\) 后的一段，得到长度至少为 \(r_0-\varepsilon\)；让 \(\varepsilon\) 趋零，便得下界 \(r_0\)。固定 \(\theta\) 的径向测地段达到该下界。

	将任意起点 \(p\) 等距移到 \(o\)，可知两点间距离是 \(\operatorname{arcosh}(-\langle p,q\rangle_L)\)。当 \(p\ne q\) 时记它为 \(a>0\)，则到达 \(q\) 的单位初速度必须是
	\[
		v=\frac{q-p\cosh a}{\sinh a}.
	\]
	该向量与 \(p\) 正交且 Lorentz 长度为 \(1\)，于是相应测地段存在且唯一。双曲面测地线的轨迹是过原点的二维平面与上半支之交，且每条单位速参数线都不重复。将 \(P^{-1}\) 代入 Lorentz 内积，整理即得\eqref{g1:eq:disc-distance}；再用
	\[
		|C(z_1)-C(z_2)|^2=\frac{4|z_1-z_2|^2}{|z_1+i|^2|z_2+i|^2}
	\]
	及上一小节的 \(1-|C(z)|^2\) 公式，得到\eqref{g1:eq:half-plane-distance}。

	由距离公式，\(o\) 的半径 \(R\) 闭球恰是双曲面上 \(x_0\leq\cosh R\) 的部分，它在 \(\mathbb R^3\) 中闭且有界，故紧。距离公式又说明坐标收敛蕴含距离收敛；反过来，将极限点移到 \(o\)，由 \(x_0=\cosh r\)、\(|x'|=\sinh r\) 可知距离收敛蕴含坐标收敛。因此上述紧性也是长度距离的紧性。任意 Cauchy 序列有界，落在某个这样的紧球内，有收敛子列，再由 Cauchy 条件整个序列收敛，故度量完备。公式\eqref{g1:eq:hyperboloid-geodesic}另已直接给出测地线的全时延拓；这里没有把两种完备性未经证明地互换。最后，圆盘凸且可沿直线缩到原点，故单连通，模型间的微分同胚将这一性质传给双曲面和上半平面。
\end{proof}

圆盘中的测地线是直径或与单位圆周正交的圆弧。具体地，双曲面测地线所在平面可写为 \(a_0x_0+a_1x_1+a_2x_2=0\)，投影后成为
\[
	a_0(1+u^2+v^2)+2a_1u+2a_2v=0.
\]
当 \(a_0=0\) 时是过原点的直线；否则圆心为 \((-a_1/a_0,-a_2/a_0)\)，半径平方等于圆心到原点距离平方减去 \(1\)，这正是与单位圆正交的条件。每条相交圆弧都对应一个这样的测地平面。把 \(w=C(z)\) 代入，方程化成 \(A(x^2+y^2)+Bx+C_0=0\)，所以在上半平面中，它们是以实轴上点为圆心的半圆，或垂直实轴的直线。通过两个不同内点的这类弧唯一，正与定理的唯一性一致。

\begin{figure}[htbp]
	\centering
	\input{Figures/Chapter10-HyperbolicModels}
	\caption[双曲平面的测地线与理想边界]{圆盘中的测地线为直径及正交于圆周的圆弧；上半平面中则为垂直实轴的直线及圆心在实轴上的半圆。灰色边界用于表示理想端点，不属于模型本身。图中的两组曲线分别示意各模型中的形状，不标示逐条对应。}
	\label{g1:fig:hyperbolic-model-geodesics}
\end{figure}

圆盘的圆周称为其\term[lixiangbianjie]{理想边界}[ideal boundary]，记为 \(\partial_\infty\mathbb H^2\)。上半平面的相应边界是 \(\mathbb R\cup\{\infty\}\)，其中一点 \(\infty\) 由 \(C\) 送到圆周上的 \(1\)。理想端点只描述测地线向哪个方向无限延伸，不是度量完备化中缺少的点。事实上
\[
	d_{\mathbb D}(0,w)=2\operatorname{artanh}|w|
	=\log\frac{1+|w|}{1-|w|}\longrightarrow\infty\qquad(|w|\to1).
\]
双曲面上，同一个端点由向量 \(p+v\) 的正比例类表示，因为\eqref{g1:eq:hyperboloid-geodesic}在 \(t\to+\infty\) 时沿 \(e^t(p+v)/2\) 增长，而 \(\langle p+v,p+v\rangle_L=0\)。不同理想端点恰好确定一条完整测地线；在上半平面中，两有限端点决定一个半圆，有一个端点为 \(\infty\) 时则决定一条竖线。

\ProseParagraphBreak

球面的距离提供了重要对照。\cref{g1:exm:sphere-geodesics}已证明单位球面的测地线是大圆，而且可全时延拓。固定 \(p\in S^2\)，写 \(q=p\cos r+u\sin r\)，其中 \(u\perp p\)、\(|u|=1\)、\(0<r<\pi\)。直接微分给出 \(\mathrm{d}r^2+\sin^2r\,\mathrm{d}\theta^2\)，同样的径向长度比较得到
\begin{equation}
	\label{g1:eq:spherical-distance}
	d_{S^2}(p,q)=\arccos(p\cdot q)\in[0,\pi].
\end{equation}
到对径点时可先截到 \(r=\pi-\varepsilon\) 再取极限，故同一公式仍成立。非对径的不同点之间，较短大圆弧是唯一最短测地段；对径点之间所有半大圆都最短。距离公式还说明它与球面原有拓扑相同，因此球面的紧性保证度量完备以及闭球紧。球面单连通已由\cref{g1:prop:sphere-simply-connected}证明；欧氏平面的相应三种性质则分别来自凸性、欧氏 Cauchy 序列收敛和闭有界集紧性。

\subsection{Möbius 变换与等距群}
\label{g1:subsec:1001:04}

已经知道模型之间怎样等距，现在要找同一模型的全部等距。仅写出一族保持度量的公式还不够，还须证明没有遗漏。一个实用的办法是先证明：等距只要在一点具有相同的位置和微分，便处处相同；随后用显式变换实现所有可能的初始数据。

\begin{proposition}[等距的初始数据]
	\label{g1:prop:isometry-initial-data}
	设 \(M,N\) 是光滑度量曲面，\(M\) 连通，\(F_1,F_2:M\to N\) 为局部等距。若在某一点 \(p\) 有 \(F_1(p)=F_2(p)\) 且 \(\mathrm{d}(F_1)_p=\mathrm{d}(F_2)_p\)，则 \(F_1=F_2\)。特别地，两个局部等距若在非空开集上相同，就在其连通公共定义域上相同。
\end{proposition}
\begin{proof}
	局部等距保持联络：在可逆邻域上，把目标联络拉回源曲面，所得运算满足度量相容、无挠性及求导法则，故由\cref{g1:prop:surface-connection-properties}的唯一性，等于源联络。因此它把带原仿射参数的测地线映为测地线。\cref{g1:prop:geodesic-initial-values}的初值唯一性给出
	\[
		F_i(\exp_qv)=\exp_{F_i(q)}\bigl(\mathrm{d}(F_i)_qv\bigr)
	\]
	对充分小的 \(v\) 成立。

	令 \(S\) 为两映射的位置与微分都相同的点集。假设使它非空。若 \(q\in S\)，由\cref{g1:prop:normal-neighborhood}，\(\exp_q\) 在零附近给出一个开邻域坐标，上式使两映射在整个邻域相同，其微分自然也相同，故 \(S\) 开。若 \(q_j\in S\) 且 \(q_j\to q\)，连续性先给出 \(F_1(q)=F_2(q)\)；在该公共像点附近选一张坐标图，再用两映射的一阶导数连续性，得到 \(\mathrm{d}(F_1)_q=\mathrm{d}(F_2)_q\)，故 \(S\) 闭。连通性迫使 \(S=M\)。
\end{proof}

以下的等距指保持度量的光滑微分同胚，组成等距群 \(\operatorname{Isom}\)。取实矩阵 \(A=\begin{psmallmatrix}a&b\\c&d\end{psmallmatrix}\) 且 \(ad-bc=1\)，定义\term[mobiusbianhuan]{Möbius 变换}[Möbius transformation]
\[
	f_A(z)=\frac{az+b}{cz+d}.
\]
分母在上半平面不为零。其在实轴及 \(\infty\) 上的取值按分式极限解释，例如 \(c\ne0\) 时 \(f_A(\infty)=a/c\)。

\begin{theorem}[双曲平面的等距群]
	\label{g1:thm:hyperbolic-isometry-group}
	上半平面的全部保向等距恰为上述实 Möbius 变换，且
	\[
		\operatorname{Isom}^+(\mathbb H)
		\cong\operatorname{PSL}(2,\mathbb R)
		=\operatorname{SL}(2,\mathbb R)/\{I,-I\}.
	\]
	它在正向正交单位标架上传递，且实现任意给定的源标架与目标标架的等距唯一。全部反向等距都是保向等距与 \(\sigma(z)=-\overline z\) 的复合。圆盘模型的全部保向等距可写成
	\[
		w\longmapsto e^{i\theta}\frac{w-a}{1-\overline a w},\qquad |a|<1,\quad\theta\in\mathbb R.
	\]
\end{theorem}
\begin{proof}
	直接计算有
	\[
		\operatorname{Im}f_A(z)=\frac{y}{|cz+d|^2},\qquad
		f_A'(z)=\frac1{(cz+d)^2}.
	\]
	于是 \(|\mathrm{d}f_A|^2/(\operatorname{Im}f_A)^2=|\mathrm{d}z|^2/y^2\)。逆矩阵给出逆映射，非零复导数保持定向，故它们都是保向等距。代入分式可得 \(f_A\circ f_B=f_{AB}\)。若 \(f_A(z)=z\) 恒成立，比较多项式 \(az+b=cz^2+dz\) 的系数，得到 \(b=c=0\)、\(a=d\)，再由行列式为 \(1\) 得 \(A=\pm I\)，这解释了商群中的两个矩阵。

	对点 \(x+iy\)，映射 \(z\mapsto x+yz\) 将 \(i\) 送到它，对应矩阵
	\[
		\begin{pmatrix}\sqrt y&x/\sqrt y\\0&1/\sqrt y\end{pmatrix}.
	\]
	而矩阵 \(K_t=\begin{psmallmatrix}\cos t&\sin t\\-\sin t&\cos t\end{psmallmatrix}\) 对应的变换固定 \(i\)，其在 \(i\) 处的微分为复数乘法 \(e^{2it}\)，可以实现任意角度的旋转。组合这两类变换及其逆，就能将任意一点的正向单位标架送到任意另一标架。

	现在给定任意保向等距 \(F\)，在一个点选正向单位标架，并从上述变换中选 \(f_A\) 与 \(F\) 的像点和微分一致。\cref{g1:prop:isometry-initial-data}迫使二者处处相同，也证明给定标架数据后的唯一性。反射 \(\sigma\) 保持 \(y\) 与欧氏长度，故是反向等距；任意反向等距与它复合后保向，因而没有其他类型。

	圆盘中的 \(h_a(w)=(w-a)/(1-\overline a w)\) 满足
	\[
		1-|h_a(w)|^2=\frac{(1-|a|^2)(1-|w|^2)}{|1-\overline a w|^2},\qquad
		h_a'(w)=\frac{1-|a|^2}{(1-\overline a w)^2}.
	\]
	所以它是把 \(a\) 送到 \(0\) 的保向等距；其逆为 \(h_{-a}\)。旋转 \(w\mapsto e^{i\theta}w\) 也保持度量，并在 \(0\) 处实现所有正向正交微分。给定保向等距 \(F\)，取 \(a=F^{-1}(0)\)，则 \(F\circ h_a^{-1}\) 固定原点，用初始数据命题知它必为某个旋转，得到所列公式。
\end{proof}

等长的有序双曲线段可以由唯一保向等距相互匹配。确切地，若 \(p\ne q\)、\(p'\ne q'\) 且 \(d(p,q)=d(p',q')\)，取各自唯一最短段的初始单位切向量 \(v,v'\)，将 \((p,v,Jv)\) 送到 \((p',v',Jv')\) 的等距便把整个线段及终点一起送过去；反过来，匹配有序端点的保向等距必须具有这些初始数据。这个结论以后将用于基本多边形的边配对。

\ProseParagraphBreak

球面与欧氏平面也有同样的标架描述。把三个模型的结论统一写出，可以直接用于下一节的边配对及局部几何图册。

\begin{proposition}[模型等距的初始数据]
	\label{g1:prop:model-isometry-data}
	令 \(X\) 为欧氏平面、单位球面或曲率 \(-1\) 的双曲平面。给定 \(p,q\in X\) 及线性内积等距 \(L:\mathrm{T}_pX\to \mathrm{T}_qX\)，存在唯一全局等距 \(F:X\to X\)，使 \(F(p)=q\)、\(\mathrm{d}F_p=L\)。这里允许 \(L\) 保向或反向。因而三个模型的等距群都在同向单位标架上传递，并可用反向等距匹配相反定向的标架。定义在非空连通开集 \(U\subset X\) 上的任意局部等距 \(f:U\to X\)，都有唯一的全局模型等距延拓。
\end{proposition}
\begin{proof}
	双曲情形，保向初值由\cref{g1:thm:hyperbolic-isometry-group}的标架传递性实现；反向初值先与模型反射的微分复合，化成保向初值，再把反射复合回来即可。球面情形，选取 \(\mathrm{T}_pS^2\) 的正交基 \((e_1,e_2)\)。\((p,e_1,e_2)\) 与 \((q,Le_1,Le_2)\) 都是 \(\mathbb R^3\) 的正交基，唯一将前者送到后者的正交线性变换限制到球面，就实现了所需初值；当 \(L\) 保向时矩阵行列式为 \(1\)，反向时为 \(-1\)。欧氏平面则直接取 \(F(x)=q+L(x-p)\)。

	唯一性由\cref{g1:prop:isometry-initial-data}给出。对于局部等距 \(f\)，取 \(p\in U\)，以 \(q=f(p)\)、\(L=\mathrm{d}f_p\) 构造全局等距 \(F\)；同一唯一性命题应用于连通域 \(U\)，便得 \(F|_U=f\)，任何另一全局延拓也具有相同初始数据。
\end{proof}

这还给出了另外两个模型的全部等距：球面保向等距群是 \(\operatorname{SO}(3)\)，全等距群是 \(\operatorname{O}(3)\)；欧氏平面的全等距为 \(x\mapsto Qx+b\)，其中 \(Q\in\operatorname{O}(2)\)，保向时 \(\det Q=1\)。

\subsection{球面与双曲三角恒等式}
\label{g1:subsec:1001:05}

三角公式比较的是边长与角，不依赖画出的曲线看起来有多弯。在双曲平面中，取三个不共测地线的点，以唯一测地段连接；球面中则取包含在一个开半球内的凸测地三角形，三边都用较短大圆弧。记顶点为 \(A,B,C\)，对应内角为 \(\alpha,\beta,\gamma\)，其对边长分别为 \(a,b,c\)。这些条件使三角形非退化，所有角属于 \((0,\pi)\)；球面的边长还属于 \((0,\pi)\)。

\ProseParagraphBreak

双曲三角形确实围成凸圆盘，可以用一个辅助投影看清。令 \(Q(x)=(x_1/x_0,x_2/x_0)\)，则 \(Q\) 把双曲面微分同胚地送到单位圆盘，逆映射为 \(u\mapsto(1,u)/\sqrt{1-|u|^2}\)。测地线所在的过原点平面在该投影下成为直线，故三条测地边恰对应普通三角形的三条线段。这称为\term[kleinmoxing]{Klein 模型}[Klein model]；它使凸性清楚，却一般不保欧氏角度。球面开半球向其切向平面的中心投影同样把较短大圆段变为直线段，所以本节所取的球面三角形也有清楚的凸内部。

\begin{theorem}[常曲率三角公式]
	\label{g1:thm:constant-curvature-trigonometry}
	对上述非退化双曲三角形，曲率为 \(-1\) 时有
	\begin{align}
		\cosh a&=\cosh b\cosh c-\sinh b\sinh c\cos\alpha,
		\label{g1:eq:hyperbolic-cosine-law}\\
		\frac{\sin\alpha}{\sinh a}&=\frac{\sin\beta}{\sinh b}=\frac{\sin\gamma}{\sinh c},
		\label{g1:eq:hyperbolic-sine-law}\\
		\cos\alpha&=-\cos\beta\cos\gamma+\sin\beta\sin\gamma\cosh a.
		\label{g1:eq:hyperbolic-angle-cosine-law}
	\end{align}
	对上述单位球面上的凸三角形，有
	\begin{align}
		\cos a&=\cos b\cos c+\sin b\sin c\cos\alpha,
		\label{g1:eq:spherical-cosine-law}\\
		\frac{\sin\alpha}{\sin a}&=\frac{\sin\beta}{\sin b}=\frac{\sin\gamma}{\sin c},
		\label{g1:eq:spherical-sine-law}\\
		\cos\alpha&=-\cos\beta\cos\gamma+\sin\beta\sin\gamma\cos a.
		\label{g1:eq:spherical-angle-cosine-law}
	\end{align}
	各余弦公式也可按 \((a,\alpha),(b,\beta),(c,\gamma)\) 循环置换。
\end{theorem}
\begin{proof}
	先看双曲面。令 \(u,v\in \mathrm{T}_A\mathbb H_L^2\) 分别为从 \(A\) 指向 \(B,C\) 的单位切向量，故 \(\langle u,v\rangle_L=\cos\alpha\)。测地线公式给出
	\[
		B=A\cosh c+u\sinh c,\qquad
		C=A\cosh b+v\sinh b.
	\]
	取负 Lorentz 内积，利用 \(\langle A,A\rangle_L=-1\)、\(u,v\perp_L A\)，立刻得到\eqref{g1:eq:hyperbolic-cosine-law}。球面情形改用 \(B=A\cos c+u\sin c\)、\(C=A\cos b+v\sin b\)，并取通常内积，就得到\eqref{g1:eq:spherical-cosine-law}。这两个计算解释了公式中双曲余弦与普通余弦、以及中间正负号的区别。

	为证明正弦公式，取 \(\mathcal D=|\det(A,B,C)|\)。在双曲情形，扩充 \(A\) 为 Lorentz 正交单位基 \((A,e_1,e_2)\)，其行列式绝对值为 \(1\)：对基矩阵的关系 \(F^{\mathsf T}\operatorname{diag}(-1,1,1)F=\operatorname{diag}(-1,1,1)\) 取行列式即可看出。上述表达式遂给出
	\[
		\mathcal D=\sinh b\sinh c\sin\alpha
		=\sinh c\sinh a\sin\beta
		=\sinh a\sinh b\sin\gamma.
	\]
	三个角都在 \((0,\pi)\)，故正弦为正，没有额外绝对值符号。除以 \(\sinh a\sinh b\sinh c\) 得所求结论。球面用通常正交基作同样的行列式计算，把各 \(\sinh\) 换成 \(\sin\)，也得到正弦公式。

	最后证明角形式的余弦公式，以免在后面使用直角三角形时暗加另一条恒等式。双曲情形记 \(x=\cosh a\)、\(y=\cosh b\)、\(z=\cosh c\)。将三个位置向量的 Lorentz Gram 矩阵取行列式，得到
	\[
		\mathcal D^2=1+2xyz-x^2-y^2-z^2.
	\]
	边余弦公式给出 \(\cos\beta=(xz-y)/(\sinh a\sinh c)\)、\(\cos\gamma=(xy-z)/(\sinh a\sinh b)\)，行列式公式给出对应的正弦。于是 \(-\cos\beta\cos\gamma+x\sin\beta\sin\gamma\) 的共同分子为
	\[
		-(xz-y)(xy-z)+x\mathcal D^2=(x^2-1)(yz-x),
	\]
	共同分母为 \((x^2-1)\sinh b\sinh c\)。约去后是 \(\cos\alpha\)，证明\eqref{g1:eq:hyperbolic-angle-cosine-law}。

	球面情形令 \(x=\cos a\)、\(y=\cos b\)、\(z=\cos c\)，通常 Gram 行列式仍给出同一个 \(\mathcal D^2\) 表达式；这时相应分子为
	\[
		-(y-xz)(z-xy)+x\mathcal D^2=(1-x^2)(x-yz),
	\]
	分母为 \((1-x^2)\sin b\sin c\)，约去后仍为 \(\cos\alpha\)。
\end{proof}

例如双曲直角三角形的直角对边为 \(c\) 时，\(\cosh c=\cosh a\cosh b\)；球面对应式为 \(\cos c=\cos a\cos b\)。这些是相应的勾股关系。若边长都很小，把 \(\cosh t=1+t^2/2+O(t^4)\)、\(\sinh t=t+O(t^3)\) 代入双曲余弦公式，最低阶便恢复 \(a^2=b^2+c^2-2bc\cos\alpha\)。局部接近欧氏几何，并不意味着这些公式对较大三角形仍可用欧氏公式替代。

\subsection{角和、面积及曲率归一化}
\label{g1:subsec:1001:06}

三角公式通过线性代数计算边和角，面积则由前章的曲率积分直接给出。这里必须使用测地边，因为一般曲线边还带有测地曲率积分。

\begin{proposition}[测地多边形的面积]
	\label{g1:prop:geodesic-polygon-area}
	设常曲率 \(K\) 的光滑度量曲面中，紧闭圆盘区域 \(D\) 的边界由 \(n\) 条测地线弧组成，边只在相邻端点相交，各内角 \(\alpha_j\in(0,2\pi)\)。则
	\begin{equation}
		\label{g1:eq:constant-curvature-polygon-area}
		K\operatorname{Area}(D)=\sum_{j=1}^n\alpha_j-(n-2)\pi.
	\end{equation}
	尤其在曲率 \(-1\) 的双曲平面中，面积为 \((n-2)\pi-\sum_j\alpha_j\)；单位球面上面积为 \(\sum_j\alpha_j-(n-2)\pi\)。
\end{proposition}
\begin{proof}
	\cref{g1:thm:global-gauss-bonnet}允许圆盘不在单张坐标图内。这里圆盘的 Euler 示性数为 \(1\)，每条测地边的 \(k_g=0\)，外角是 \(\pi-\alpha_j\)，故
	\[
		K\operatorname{Area}(D)+\sum_{j=1}^n(\pi-\alpha_j)=2\pi.
	\]
	整理即得公式。全部假设中的紧性、圆盘拓扑以及非退化角点，正是所引公式在这里所需的条件。
\end{proof}

非退化双曲三角形有正面积，所以角和严格小于 \(\pi\)；单位球面的上述凸三角形角和严格大于 \(\pi\)。这与\cref{g1:cor:geodesic-triangle-defect}相符。反过来，仅知道三角形某些角度满足一个数值不等式，还不能在未经构造时断言所需边长与边界都已经存在。

\begin{example}[理想三角形]
	\label{g1:exm:ideal-triangle-area}
	在上半平面取理想顶点 \(0,1,\infty\)。它的三边是 \(x=0\)、\(x=1\) 与直径为 \([0,1]\) 的半圆，内部为 \(0<x<1\)、\(y>\sqrt{x(1-x)}\)。面积元是 \(y^{-2}\,\mathrm{d}x\,\mathrm{d}y\)，故
	\[
		\operatorname{Area}(D)
		=\int_0^1\int_{\sqrt{x(1-x)}}^\infty\frac{\mathrm{d}y\,\mathrm{d}x}{y^2}
		=\int_0^1\frac{\mathrm{d}x}{\sqrt{x(1-x)}}=\pi.
	\]
	最后可令 \(x=(1+\cos t)/2\) 积分。这个区域非紧，其顶点也不在双曲平面内，因此不能不加解释地把角度 \(0\) 代入紧圆盘定理；这里的直接积分给出了依据。任意三个不同且按正向排列的理想端点都可由实分式变换送到 \(0,1,\infty\)：若三者有限，先用 \((z-a)(b-c)/((z-c)(b-a))\) 把它们送到 \(0,1,\infty\)，涉及 \(\infty\) 时取相应极限，选择保持边界循环次序的排列保证它保持上半平面；矩阵再归一化为行列式 \(1\)。所以每个理想三角形的面积都是 \(\pi\)。
\end{example}

曲率的归一化会改变面积公式中的尺度。若把度量乘以常数 \(R^2>0\)，长度乘以 \(R\)，面积乘以 \(R^2\)，角度不变。\eqref{g1:eq:christoffel-metric}中度量与逆度量的常数因子相消，所以联络不变；\eqref{g1:eq:intrinsic-gauss-curvature}的分子乘以 \(R^2\)、分母乘以 \(R^4\)，故曲率乘以 \(R^{-2}\)。因此半径 \(R\) 的球面与曲率 \(-1/R^2\) 的双曲平面，三角公式中的每个边长都应写成 \(a/R,b/R,c/R\)，而面积公式分别成为
\[
	\operatorname{Area}_{+}(D)=R^2\left(\sum_j\alpha_j-(n-2)\pi\right),\qquad
	\operatorname{Area}_{-}(D)=R^2\left((n-2)\pi-\sum_j\alpha_j\right).
\]
两式中的角度都以弧度计。选定曲率 \(-1\) 相当于选定长度单位，并非说任意常负曲率度量的数字都天然等于 \(-1\)。

\subsection{测地多边形和角条件}
\label{g1:subsec:1001:07}

为了随后把多边形的边两两粘合，需要能同时控制边长与顶点内角。欧氏正 \(n\) 边形缩放后角度不变；双曲正 \(n\) 边形则会随大小增加而使内角减小。这一变化恰好提供了曲面粘合所需的角度。

\begin{proposition}[正则双曲多边形]
	\label{g1:prop:regular-hyperbolic-polygons}
	对整数 \(n\geq3\) 及任意
	\[
		0<\alpha<\frac{(n-2)\pi}{n},
	\]
	曲率 \(-1\) 的双曲平面中存在一个凸正则测地 \(n\) 边形，其边长相等、内角均为 \(\alpha\)。以中心到顶点的距离 \(R>0\) 作参数时，内角恰为
	\begin{equation}
		\label{g1:eq:regular-hyperbolic-polygon-angle}
		\alpha(R)=2\arctan\frac{\cot(\pi/n)}{\cosh R}.
	\end{equation}
	它连续严格递减，从 \((n-2)\pi/n\) 趋向 \(0\)。特别地，对每个 \(g\geq2\)，存在内角均为 \(\pi/(2g)\) 的正则 \(4g\) 边形，其外接半径与面积满足
	\[
		\cosh R=\cot^2\frac{\pi}{4g},\qquad
		\operatorname{Area}=4\pi(g-1).
	\]
\end{proposition}
\begin{proof}
	在双曲面的径向坐标中取中心 \(O=o\)，顶点为
	\[
		V_j=(\cosh R,\sinh R\cos(2\pi j/n),\sinh R\sin(2\pi j/n)),
		\quad 0\leq j<n.
	\]
	用唯一测地段连接相邻顶点。Klein 投影 \(Q\) 把它变成欧氏圆周半径 \(\tanh R\) 上的正 \(n\) 边形，故这些边围成凸圆盘；空间坐标的旋转与反射是双曲等距，使所有边长及内角分别相等。

	取一条边的双曲中点 \(M\)，以及它的一个端点 \(V\)。关于该边两端顶点的径向角平分线反射，会交换两个顶点，因而由测地段唯一性也交换整条边的两半。固定轴从中心经过 \(M\)，其微分在轴向为 \(+1\)、在横向为 \(-1\)，而边在 \(M\) 处的切向量被取负，所以 \(OM\) 与该边垂直。关于 \(OV\) 的反射则交换 \(V\) 的两条邻边，故平分内角。于是 \(OMV\) 是一个直角三角形，其在 \(O\) 处的角为 \(\pi/n\)，在 \(V\) 处的角为 \(\alpha/2\)，直角所对边长为 \(R\)，如\cref{g1:fig:regular-hyperbolic-polygon}所示。

	将\eqref{g1:eq:hyperbolic-angle-cosine-law}用于直角，得到
	\[
		0=-\cos\frac\pi n\cos\frac\alpha2
		+\sin\frac\pi n\sin\frac\alpha2\cosh R.
	\]
	所有角都在 \((0,\pi/2)\)，故可约去正因子，得到\eqref{g1:eq:regular-hyperbolic-polygon-angle}。\(\cosh R\) 对 \(R>0\) 严格增加，所以 \(\alpha\) 连续严格减少；当 \(R\to0\) 时其极限为 \(2(\pi/2-\pi/n)\)，当 \(R\to\infty\) 时趋于零。介值定理给出每个允许内角的存在性，严格单调性也给出半径唯一性。

	最后取 \(n=4g\)、\(\alpha=\pi/(2g)=2\pi/n\)。当且仅当 \(g>1\) 时，目标角严格小于上述上限，故有正半径解。代入角公式得 \(\cosh R=\cot^2(\pi/n)\)。再用\cref{g1:prop:geodesic-polygon-area}，面积是 \((4g-2)\pi-4g\cdot\pi/(2g)=4\pi(g-1)\)。
\end{proof}

\begin{figure}[htbp]
	\centering
	\input{Figures/Chapter10-RegularPolygon}
	\caption[正则双曲多边形的直角三角形]{圆盘模型中的正则八边形。中心 \(O\)、边中点 \(M\) 与顶点 \(V\) 构成直角三角形，\(OV=R\)，中心角为 \(\pi/n\)，顶点角为 \(\alpha/2\)。虚线圆周是理想边界；测地边向圆盘中心一侧弯曲。}
	\label{g1:fig:regular-hyperbolic-polygon}
\end{figure}

角条件在粘合时有直接含义。若若干多边形顶点被识别成一个点，它们围出的角必须合计 \(2\pi\)，才可能成为普通光滑曲面点；否则会留下圆锥型的角亏或角盈。上面的 \(4g\) 个角之和正是 \(2\pi\)，而边长相等允许按保向等距匹配有序端点。下一节还会逐项检查边两侧的光滑接合、顶点邻域以及整体商空间；这里的多边形构造提供了所需的具体长度与角度数据。

% !TeX root = ../../Book1.tex
% 纲要标题：### 10.2 离散群与曲面商
% 纲要位置：几何学计划纲要.md，第 515 行。
% 正文状态：初稿；尚未正式审核。

% 纲要内容开始
% 
% * 椭圆、抛物和双曲变换
% * 自由固有作用及基本区域
% * 基本多边形的边配对与粘合
% * 闭曲面、尖点与漏斗的模型
% 
% 纲要内容结束

\section{离散群与曲面商}
\label{g1:sec:1002}

双曲平面中每一点的局部几何都相同，但把不同的点按等距变换识别以后，可以得到有闭测地线的圆柱、面积有限的尖点端，乃至任意亏格不小于二的闭曲面。这里有两个不同的构造问题：给定一个等距变换群，商是否仍是光滑曲面；给定一张多边形和边配对，接缝及顶点附近的度量是否光滑。前一个问题由作用的自由性和固有性控制，后一个问题必须逐条检查边的等距配对和顶点周围的总角。

\subsection{椭圆、抛物和双曲变换}
\label{g1:subsec:1002:01}

在上半平面 \(\mathbb H=\{z=x+iy:y>0\}\) 上，度量为 \(g_{\mathbb H}=|\mathrm{d}z|^2/y^2\)。\cref{g1:thm:hyperbolic-isometry-group}已经说明，其保向等距变换恰为
\[
	z\longmapsto\frac{az+b}{cz+d},\qquad
	A=\begin{pmatrix}a&b\\c&d\end{pmatrix}\in\operatorname{SL}(2,\mathbb R),
\]
其中 \(A\) 与 \(-A\) 给出同一个变换，因此群为 \(\operatorname{PSL}(2,\mathbb R)\)。这些变换也作用于理想边界 \(\mathbb R\cup\{\infty\}\)；那里不是双曲平面的一部分，边界不动点与平面内不动点必须区别。

\begin{theorem}[保向等距变换的分类]\label{g1:thm:psl-element-classification}
	设 \(h\in\operatorname{PSL}(2,\mathbb R)\) 不是恒等变换，取代表矩阵 \(A\in\operatorname{SL}(2,\mathbb R)\)，记 \(\tau=(\operatorname{tr}A)^2\)。则有且仅有下列一种情形。
	\begin{enumerate}
		\item \(\tau<4\)：\(h\) 在 \(\mathbb H\) 中有唯一不动点，在理想边界没有不动点；移到圆盘模型并把该点移到原点后，\(h\) 为非恒等旋转 \(w\mapsto e^{i\theta}w\)。称 \(h\) 为\term[tuoyuanbianhuan]{椭圆变换}[elliptic isometry]。
		\item \(\tau=4\)：\(h\) 在 \(\mathbb H\) 中没有不动点，在理想边界有唯一不动点；经保向等距共轭后，\(h\) 为 \(z\mapsto z+1\) 或 \(z\mapsto z-1\)。称 \(h\) 为\term[paowubianhuan]{抛物变换}[parabolic isometry]。
		\item \(\tau>4\)：\(h\) 在 \(\mathbb H\) 中没有不动点，在理想边界有两个不同不动点；经保向等距共轭后，\(h\) 为 \(z\mapsto e^\ell z\)，其中 \(\ell>0\) 且
		\[
			|\operatorname{tr}A|=2\cosh(\ell/2).
		\]
		称 \(h\) 为\term[shuangqubianhuan]{双曲变换}[hyperbolic isometry]。
	\end{enumerate}
\end{theorem}
\begin{proof}
	更换代表 \(A\mapsto-A\) 不改变 \(\tau\)，共轭也不改变它。有限不动点满足
	\[
		cz^2+(d-a)z-b=0,
	\]
	其判别式为 \((a+d)^2-4(ad-bc)=\tau-4\)。若 \(c=0\)，还要把无穷远点计入；统一地说，理想边界不动点就是实矩阵 \(A\) 的实特征直线，所以判别式同样区分零个、一个和两个实不动点。

	若 \(\tau<4\)，两个非实根互为共轭，恰有一个在 \(\mathbb H\) 中。设它为 \(p=x+iy\)，变换 \(z\mapsto(z-x)/y\) 将其送到 \(i\)。一个实 Möbius 变换固定 \(i\) 当且仅当其代表可写成
	\[
		R_\theta=\begin{pmatrix}
			\cos(\theta/2)&\sin(\theta/2)\\
			-\sin(\theta/2)&\cos(\theta/2)
		\end{pmatrix}.
	\]
	用圆盘坐标 \(w=(z-i)/(z+i)\) 代入，所得变换为 \(w\mapsto e^{i\theta}w\)。非恒等性排除 \(\theta\equiv0\pmod{2\pi}\)，也给出平面内不动点的唯一性。

	若 \(\tau=4\)，非恒等性说明 \(A\ne\pm I\)，故只有一条实特征直线。将它送到 \(\infty\)，变换便写成 \(z\mapsto a^2z+ab\)，其中代表矩阵的对角元为 \(a,a^{-1}\)。迹平方等于四迫使 \(a=\pm1\)，于是它是非零平移 \(z\mapsto z+t\)。再用正缩放共轭，将 \(|t|\) 变成 \(1\)。正、负两个方向都保留在陈述中，因为保向共轭不能随意把这两个标准型混为同一个。

	若 \(\tau>4\)，将两个理想不动点送到 \(0,\infty\)，得到 \(z\mapsto\lambda z\)，其中 \(\lambda>0\) 且 \(\lambda\ne1\)。若 \(\lambda<1\)，再用 \(z\mapsto-1/z\) 交换两端，便可令 \(\lambda=e^\ell>1\)。代表矩阵可取 \(\operatorname{diag}(e^{\ell/2},e^{-\ell/2})\)，给出迹的公式。三个标准型的全部不动点均已列出，证明分类。
\end{proof}

迹平方判别的是三种类型，不是所有共轭参数。例如圆盘中旋转 \(\theta\) 与旋转 \(-\theta\) 具有相同的迹平方，但保向共轭还保留旋转方向。下一步要构造商时，更重要的是一个变换把点移动了多远。

\begin{definition}[平移长度]\label{g1:def:hyperbolic-translation-length}
	设 \(\mathbb H\) 为曲率 \(-1\) 的双曲平面，\(d_{\mathbb H}\) 为其长度距离，\(h:\mathbb H\to\mathbb H\) 为等距变换。它的\term[pingyichangdu]{平移长度}[translation length]定义为
	\[
		L(h)=\inf_{p\in\mathbb H}d_{\mathbb H}(p,hp).
	\]
\end{definition}

对椭圆变换，\(L(h)=0\)，且在唯一内点不动点处达到。对抛物变换 \(z\mapsto z+t\)，\cref{g1:thm:hyperbolic-distance-geodesics}中的距离公式给出
\[
	\cosh d_{\mathbb H}(z,z+t)=1+\frac{t^2}{2y^2},
\]
所以 \(L(h)=0\)，但没有任何平面内点达到这个下确界。对双曲标准型，距离公式则给出
\begin{equation}\label{g1:eq:hyperbolic-displacement}
	\cosh d_{\mathbb H}(z,e^\ell z)
	=\cosh\ell+2\sinh^2(\ell/2)\frac{x^2}{y^2}.
\end{equation}
因此 \(L(h)=\ell\)，恰在虚轴上达到。一般双曲变换的两个理想不动点之间的测地线称为它的\term[zhou]{轴}[axis]；沿轴的作用就是长度为 \(\ell\) 的平移。下确界是否达到，已经把抛物与双曲两种无内点不动点的行为区别开来。

\subsection{自由固有作用及基本区域}
\label{g1:subsec:1002:02}

为了使商附近仍是一张普通小圆盘，需要排除两种障碍：一个群元素固定该点，或者无限多个不同群元素在附近反复堆积。前者由自由性排除，后者由下面的紧集条件控制。

\begin{definition}[离散子群与固有作用]\label{g1:def:discrete-proper-isometric-action}
	给 \(\operatorname{PSL}(2,\mathbb R)\) 赋予矩阵群模去 \(\{\pm I\}\) 的拓扑。若子群 \(\Gamma\) 有一个单位元邻域只包含其单位元，就称它为\term[lisanziqun]{离散子群}[discrete subgroup]。对群在 \(\mathbb H\) 上的作用，若 \(\gamma p=p\) 蕴含 \(\gamma=1\)，就称它为\term[ziyouzuoyong]{自由作用}[free action]；若对每个紧集 \(K\subset\mathbb H\)，集合
	\[
		\{\gamma\in\Gamma:\gamma K\cap K\ne\varnothing\}
	\]
	有限，就称它为\term[guyouzuoyong]{固有作用}[proper action]。
\end{definition}

这里把 \(\Gamma\) 当作离散群，使用的是适合离散作用的紧集定义。第五章\cref{g1:def:properly-discontinuous-covering}中的“适当不连续”还要求每点都有与所有非平凡平移不交的小邻域，因而已经包含自由性。本节的固有性允许有限稳定子；要恢复第五章的条件，需要同时假设自由。

\begin{proposition}[离散性与固有性]\label{g1:prop:discrete-hyperbolic-action}
	子群 \(\Gamma\subset\operatorname{PSL}(2,\mathbb R)\) 离散，当且仅当其在 \(\mathbb H\) 上的作用固有。离散子群中的椭圆元素都有有限阶；其作用自由，当且仅当它没有非单位有限阶元素。
\end{proposition}
\begin{proof}
	先说明离散子群与任意紧的群元素集合只有有限交。若紧集内有无限个不同元素 \(\gamma_n\)，由矩阵空间的紧性可取收敛子列；于是不同的 \(\gamma_n^{-1}\gamma_m\) 可以任意接近单位元，与单位元孤立矛盾。这里并未假定子群先验闭；同一个论证也说明其在矩阵群中的任何收敛序列最终恒定，所以它确实闭。

	给 \(p=x+iy\) 选连续依赖于 \(p\) 的等距变换
	\[
		u_p(z)=yz+x,\qquad
		u_p=\left[\begin{pmatrix}\sqrt y&x/\sqrt y\\0&1/\sqrt y\end{pmatrix}\right],
	\]
	它将 \(i\) 送到 \(p\)。若 \(h p=q\)，则 \(u_q^{-1}hu_p\) 固定 \(i\)，故属于前面写出的旋转群 \(\{R_\theta\}\)，这是一个紧群。若 \(p,q\in K\)，那么 \(h\) 位于紧集
	\[
		\{u_qR_\theta u_p^{-1}:p,q\in K,\ 0\le\theta\le2\pi\}
	\]
	之中。离散子群与该紧集只有有限交，故作用固有。反过来，若有不同的 \(\gamma_n\) 趋于单位元，则 \(\gamma_ni\to i\)；包含这些点和 \(i\) 的一个小紧邻域 \(K\) 将满足 \(\gamma_nK\cap K\ne\varnothing\)，与固有性矛盾。

	椭圆元素的所有幂都位于一个共轭旋转群中，后者紧，所以离散性迫使其幂只有有限多个，即该元素为有限阶。分类定理表明，非恒等抛物或双曲变换都有无限阶。因此没有非单位有限阶元素等价于没有椭圆元素，又等价于没有平面内不动点，即作用自由。
\end{proof}

例如绕同一点旋转 \(2\pi/m\) 生成的有限群离散且作用固有，却不自由；这个点在商中有锥角 \(2\pi/m\)。绕同一点旋转一个与 \(2\pi\) 之比为无理数的角所生成的群则不离散，因为它有无限多个幂落在同一个紧旋转群中。离散性不会自动排除所有不动点。

\begin{theorem}[自由固有等距作用的商]\label{g1:thm:hyperbolic-quotient}
	设 \(\Gamma\subset\operatorname{PSL}(2,\mathbb R)\) 自由且固有地作用于 \(\mathbb H\)。则轨道商 \(M=\mathbb H/\Gamma\) 是 Hausdorff、第二可数的可定向光滑曲面，且有唯一度量使投影 \(q:\mathbb H\to M\) 为局部等距。该度量曲率恒为 \(-1\)，长度距离完备；\(q\) 是万有覆盖，覆盖变换群恰为 \(\Gamma\)。
\end{theorem}
\begin{proof}
	固定 \(p\in\mathbb H\)，取包含它的紧邻域 \(K\)。只有有限个 \(\gamma\) 满足 \(\gamma K\cap K\ne\varnothing\)。自由性给出 \(\gamma p\ne p\)；对有限个这样的非单位元素，同时缩小 \(p\) 的邻域 \(U\)，即可使 \(\gamma U\cap U=\varnothing\)。其余元素由 \(U\subset K\) 已经排除。因此第五章的强适当不连续条件成立。\cref{g1:prop:action-quotient-cover}的证明给出 \(q\) 为开映射，以及
	\[
		q^{-1}(q(U))=\bigsqcup_{\gamma\in\Gamma}\gamma U,
	\]
	每张片都同胚映到 \(q(U)\)。上半平面凸且可缩，特别是单连通；该命题还说明这是万有覆盖，覆盖变换群正是 \(\Gamma\)。

	仍须核验商的分离性。若 \(p,r\) 不在同一轨道，取它们的紧邻域之并 \(K\)。可能把第一个邻域送到第二个的群元素只有有限个；又因 \(\gamma p\ne r\)，可以同时缩小成 \(U\ni p,V\ni r\)，使所有 \(\gamma U\) 都不交 \(V\)。于是 \(q(U),q(V)\) 是不交开邻域，故商 Hausdorff。\(\mathbb H\) 的可数开基经开映射 \(q\) 取像，仍是商的可数开基，得到第二可数性。

	用 \(q|_U\) 的逆把 \(q(U)\) 当作上半平面小片。两个这样的图在每个足够小的交叠部分上相差唯一一个 \(\gamma\in\Gamma\)，因而转换是光滑保向等距。各片的度量便相容地下降，唯一性也由局部等距要求决定；局部曲率保持为 \(-1\)。

	最后证明完备性，而不由“局部等距”直接推出。距离公式说明 \(d_{\mathbb H}(i,x+iy)\le R\) 等价于
	\[
		x^2+(y-\cosh R)^2\le\sinh^2R.
	\]
	这个 Euclid 闭圆盘的最低高度为 \(e^{-R}>0\)，所以双曲闭球紧。固有性于是保证轨道与每个闭有界球只有有限交，故
	\begin{equation}\label{g1:eq:hyperbolic-quotient-distance}
		d_M(qx,qy)=\min_{\gamma\in\Gamma}d_{\mathbb H}(x,\gamma y).
	\end{equation}
	具体地，右端的最小值存在；不同轨道之间它为正，三角不等式由等距不变性得到。投影任意连接 \(x\) 与 \(\gamma y\) 的测地段，说明商长度距离不超过它。反过来，商中的每条分片光滑连接曲线从 \(x\) 提升后终于某个 \(\gamma y\)，提升保持长度，故其长度至少为右端。这样就证明了等式。

	对 \(M\) 中任一 Cauchy 序列，取子列 \(p_n\)，使 \(d_M(p_n,p_{n+1})<2^{-n}\)。从任意提升 \(x_1\) 开始，逐次用\eqref{g1:eq:hyperbolic-quotient-distance}选择提升 \(x_{n+1}\)，满足 \(d_{\mathbb H}(x_n,x_{n+1})<2^{-n}\)。这些提升构成 Cauchy 序列，在完备的 \(\mathbb H\) 中收敛；其投影给出 \(p_n\) 的极限。原序列因 Cauchy 性也收敛，故商完备。
\end{proof}

上述证明把拓扑和度量的责任分开了：局部不交邻域给出覆盖，固有性还保证不同轨道能被分离，而等距性才使局部度量在商上相容。

\begin{definition}[Dirichlet 区域]\label{g1:def:dirichlet-domain}
	设 \(\Gamma\subset\operatorname{PSL}(2,\mathbb R)\) 的作用自由且固有，固定 \(p\in\mathbb H\)。集合
	\[
		D_p=\{x\in\mathbb H:d_{\mathbb H}(x,p)\le d_{\mathbb H}(x,\gamma p)\ \text{对所有 }\gamma\in\Gamma\}
	\]
	称为以 \(p\) 为中心的\term[Dirichletquyu]{Dirichlet 区域}[Dirichlet domain]。它把“离哪一个轨道点最近”转化成一组半平面条件。\footnote{Johann Peter Gustav Lejeune Dirichlet（1805--1859），德国数学家，对数论、Fourier 级数与势理论作出重要贡献。}
\end{definition}

\begin{proposition}[Dirichlet 区域的铺盖]\label{g1:prop:dirichlet-domain}
	设 \(\Gamma\subseteq\operatorname{PSL}(2,\mathbb R)\) 以标准双曲等距自由且固有地作用于 \(\mathbb H\)，\(p\in\mathbb H\)，\(d\) 为双曲距离。令 \(D_p=\{x\in\mathbb H:d(x,p)\le d(x,\gamma p)\text{ 对所有 }\gamma\in\Gamma\}\)。则 \(D_p\) 是有非空内部的闭测地凸集，所有 \(\gamma D_p\) 覆盖 \(\mathbb H\)，不同平移的内部不交；这些平移局部有限。\(D_p\) 在每个有界部分由有限条测地边界限制。若 \(\mathbb H/\Gamma\) 紧，则 \(D_p\) 是紧的有限边测地多边形。
\end{proposition}
\begin{proof}
	两个不同点 \(p,q\) 的等距轨迹是一条测地线，\(d(x,p)\le d(x,q)\) 选出包含 \(p\) 的闭半平面。在双曲面模型中可以直接看出这一点：距离公式把等距条件写成 Lorentz 内积的线性等式 \(\langle X,P-Q\rangle_L=0\)，而 \(P-Q\) 为空间样向量；该平面截双曲面得到测地线。每一侧测地凸，因为长为 \(L>0\) 的测地段在双曲面上的表达是
	\[
		X(t)=\frac{\sinh((1-t)L)}{\sinh L}X(0)
		+\frac{\sinh(tL)}{\sinh L}X(1),\qquad 0\le t\le1,
	\]
	两系数均非负，保持同侧的线性不等式。因此 \(D_p\) 是闭测地凸集。固有性和自由性给出 \(\inf_{\gamma\ne1}d(p,\gamma p)>0\)，取该下界一半以内的小球，就包含在 \(D_p\) 的内部。

	对任意 \(x\)，轨道 \(\Gamma p\) 中有一个距 \(x\) 最近的点，存在性与商距离的最小值相同。若最近点为 \(\gamma p\)，则 \(\gamma^{-1}x\in D_p\)，所以各平移覆盖。若 \(x\) 同时位于 \(\gamma D_p\) 和 \(\eta D_p\) 的内部，且 \(\gamma\ne\eta\)，第一个内部条件给出 \(d(x,\gamma p)<d(x,\eta p)\)，第二个给出反向严格不等式，矛盾。严格性来自等距轨迹是一条边界测地线，不能包含内点。

	若 \(x\in\overline B(p,R)\cap\gamma D_p\)，最近点条件说明 \(d(x,\gamma p)\le d(x,p)\le R\)，于是 \(d(p,\gamma p)\le2R\)。这样的 \(\gamma\) 只有有限个，得到平移的局部有限性。同样，在 \(B(p,R)\) 内能限制 \(D_p\) 的半平面只可能来自 \(d(p,\gamma p)\le2R\) 的元素，因为更远的轨道点在此球中总比 \(p\) 远。因此每个有界部分只需有限条边界。

	若商紧，递增开集 \(q(B(p,R))\) 覆盖它，紧性给出某个 \(R\) 已经覆盖全商。每个 \(x\in D_p\) 都与 \(B(p,R)\) 中某点等价，故它距某个轨道点小于 \(R\)，最近点条件进一步给出 \(d(x,p)<R\)。于是 \(D_p\) 闭且有界，从而紧；它全部位于一个稍大的球内，所需边界只有有限条，得到有限测地多边形。
\end{proof}

满足覆盖且内部不重叠的这样的闭区域称为群作用的\term[jibenquyu]{基本区域}[fundamental domain]。边界上可以有多个代表，必须按群元素识别，不能要求每条轨道与闭区域恰交一次。非紧商的 Dirichlet 区域也未必只有有限条边，上命题只保证有界范围内的有限性。

\subsection{基本多边形的边配对与粘合}
\label{g1:subsec:1002:03}

反过来，可以先画一张双曲多边形，再沿成对的边接合。第六章\cref{g1:thm:surface-polygon}及\cref{g1:prop:orientable-standard-word}解释了边字如何记录拓扑粘合；现在还要使这些粘合保留度量。边内点的邻域由两张半圆盘组成，顶点的邻域则由若干个角扇组成，两处需要不同的检查。

\begin{proposition}[双曲多边形的光滑粘合]\label{g1:prop:hyperbolic-polygon-gluing}
	设 \(P\subset\mathbb H\) 是紧凸测地多边形，将不同的边两两配对，每对等长，并用保向双曲等距变换按相反的边界行进方向识别。对每个顶点等价类，将其各角的一条小圆弧按边配对接起；假设所得空间是一条圆周，且该类内的角度总和为 \(2\pi\)。则商 \(P/{\sim}\) 是紧、可定向、无边界的光滑曲面，\(P\) 的双曲度量在其上光滑粘合，曲率恒为 \(-1\)。
\end{proposition}
\begin{proof}
	等长测地边的端点可以按要求配对：先用\cref{g1:prop:model-isometry-data}将一条边的起点及单位切向量送到另一条边的终点及反向单位切向量；测地线的初值唯一性使整条等长边随之对应。保向性保证原来的内侧被送到另一条边的外侧，因此两侧窄半盘在识别后确实拼成一个双曲小圆盘，度量在接缝处光滑。

	这里的商也是 Hausdorff。等价关系在 \(P\times P\) 中由对角线、有限个闭边配对映射的图以及有限个顶点等价对组成，故为闭集。若 \(C\subset P\) 闭，它的饱和集是此闭关系与 \(C\times P\) 的交在第二因子的投影，因紧性而闭。两个不同等价类是 \(P\) 中不交的有限集，取不交开邻域 \(U,V\) 包含它们；从 \(U\) 中去掉所有与 \(P\setminus U\) 等价的点，得到包含原等价类的饱和开集，\(V\) 同理。两者下降为商中不交开邻域。商紧性来自 \(P\) 的紧性。

	在一个顶点类处，将各原顶点附近截成同样小半径的角扇。边等距保留距相应顶点的距离，故各扇可以用同一个径向变量 \(r\) 粘合。双曲面极坐标
	\[
		(r,\theta)\longmapsto(\cosh r,\sinh r\cos\theta,\sinh r\sin\theta)
	\]
	直接给出度量 \(\mathrm{d}r^2+\sinh^2r\,\mathrm{d}\theta^2\)。把各扇按圆弧连接的次序排好，角变量逐段累计；小圆弧接成一条圆周意味着绕行只形成一个扇环，总角 \(2\pi\) 则使这个扇环恰好拼成普通双曲小圆盘。末缝处 \(\theta=0\) 与 \(2\pi\) 的极坐标映射以及所有导数相容，中心也由上述双曲面坐标给出通常的光滑结构。因此顶点没有锥形奇点。

	面内、边内和顶点处的这些局部图在交叠处相差双曲等距变换，形成相容的光滑图册和度量。各边都已配对，故没有边界；配对反转边界行进方向而保持环境定向，所以局部定向也相容。紧性允许从局部图取有限覆盖，每张图有可数基，故商第二可数。这就得到所述光滑度量曲面。
\end{proof}

顶点附近的小圆弧拼成的圆周常称为该顶点的\term[lianjie]{链接}[link]。链接控制局部邻域是不是一张盘，总角则控制这张盘上的度量是否光滑。若链接仍是一条圆周但总角为 \(\Theta\ne2\pi\)，小半径圆周的长度就是 \(\Theta\sinh r\)，除以半径后趋于 \(\Theta\)，不能是光滑曲面中应有的 \(2\pi\)。因此检查边长相等还不够。

\ProseParagraphBreak

这一命题构造的是有限多边形的度量商。若另把边配对变换放进 \(\operatorname{PSL}(2,\mathbb R)\) 并取它们生成的群，还需要证明其离散性及原多边形的平移是否铺盖整个平面；这些整体结论不是一句“角和为 \(2\pi\)”就能直接替代的。下一节将用展开映射说明完备模型结构的整体实现。

\begin{theorem}[闭双曲曲面的多边形构造]\label{g1:thm:closed-hyperbolic-polygon}
	设 \(g\ge2\) 为整数，\(\Sigma_{g,0}\) 为可定向亏格 \(g\) 的闭曲面。存在一张与 \(\Sigma_{g,0}\) 同胚的紧光滑曲面，带曲率恒为 \(-1\) 的度量。它可以由曲率 \(-1\) 双曲平面中一个内角均为 \(\pi/(2g)\) 的凸正则测地 \(4g\) 边形按边字
	\[
		a_1b_1a_1^{-1}b_1^{-1}\cdots a_gb_ga_g^{-1}b_g^{-1}
	\]
	将同字母边按相反边界方向，用保向双曲等距两两配对得到，面积为 \(4\pi(g-1)\)。
\end{theorem}
\begin{proof}
	\cref{g1:prop:regular-hyperbolic-polygons}保证：给定 \(n\ge3\) 及 \(0<\alpha<(n-2)\pi/n\)，存在内角为 \(\alpha\) 的正则双曲 \(n\) 边形。对 \(n=4g\)、\(\alpha=\pi/(2g)\)，所需严格不等式恰由 \(g\ge2\) 保证。所有边等长，故按上命题可实施指定的等距配对。

	还要检查顶点，不能只因共有 \(4g\) 个角就默认它们拼在同一点。依次记顶点为 \(v_0,\ldots,v_{4g-1}\)，第 \(j\) 条边从 \(v_j\) 指向 \(v_{j+1}\)，下标模 \(4g\)。在第 \(k\) 个四边块中，\(a\) 配对给出
	\[
		v_{4k}\sim v_{4k+3},\qquad v_{4k+1}\sim v_{4k+2},
	\]
	\(b\) 配对给出 \(v_{4k+1}\sim v_{4k+4}\) 及 \(v_{4k+2}\sim v_{4k+3}\)。这些关系把该块的五个连续端点连在一起，因而所有顶点属于一个等价类。

	再从一个角的小圆弧穿过配对边进入下一个角，所得链接的后继顺序为
	\[
		4k+1\longmapsto4k+2\longmapsto4k+3\longmapsto4k
		\longmapsto4(k-1)+1
	\]
	（所有下标模 \(4g\)）。它先走过一个四边块的四个角，再进入前一个四边块，最终恰走过全部 \(4g\) 个角一次才返回，故链接是单一圆周。总角为 \(4g\cdot\pi/(2g)=2\pi\)，上命题给出光滑双曲度量。

	这个边字的拓扑类型由\cref{g1:prop:orientable-standard-word}中标准字与手柄连通和的对应给出，正是 \(\Sigma_{g,0}\)。最后，\cref{g1:prop:geodesic-polygon-area}给出
	\[
		\operatorname{Area}(P)=(4g-2)\pi-4g\frac{\pi}{2g}=4\pi(g-1).
	\]
	边及顶点的识别不改变面积，故商也有这个面积。
\end{proof}

对 \(g=1\)，所需四边形的四个角都为 \(\pi/2\)，角和达到平面四边形的角和 \(2\pi\)，双曲面积公式将给出零面积，所以这种非退化双曲四边形不存在。这与第九章 Gauss--Bonnet 排除闭环面上的负常曲率度量相符。

\subsection{闭曲面、尖点与漏斗的模型}
\label{g1:subsec:1002:04}

刚才得到的闭曲面紧，长度距离因而完备：局部双曲图中的度量在小闭片上与 Euclid 度量双边可比，长度距离给出原拓扑；任一 Cauchy 序列都有由紧性提供的收敛子列，进而整个序列收敛。它们的面积有限，但非紧曲面也可能有有限面积的端部。两种最简单的非紧端部来自抛物和双曲循环群。

\begin{example}[尖点]\label{g1:exm:hyperbolic-cusp}
	令 \(T(z)=z+1\)。群 \(\langle T\rangle\) 离散且自由，整个商 \(\mathbb H/\langle T\rangle\) 是完备双曲圆柱。固定 \(Y>0\)，取其中由 \(y\ge Y\) 给出的端部
	\[
		C_Y=\{y\ge Y\}/\langle T\rangle.
	\]
	这是\term[jiandian]{尖点}[cusp]的标准模型。其一条基本条带为 \(0\le x\le1,y\ge Y\)，两条竖边配对。令 \(r=\log(y/Y)\)，则
	\begin{equation}\label{g1:eq:hyperbolic-cusp-metric}
		g=\mathrm{d}r^2+Y^{-2}e^{-2r}\,\mathrm{d}x^2,
		\qquad r\ge0,\quad x\in\mathbb R/\mathbb Z.
	\end{equation}
	截圈 \(r=R\) 的长度为 \(Y^{-1}e^{-R}\)，向尖端越来越短；但是任何向 \(r\to\infty\) 行进的曲线长度至少为其 \(r\) 的总变化，所以尖端处在无限距离处。另一方面，
	\[
		\operatorname{Area}(C_Y)=\int_0^1\int_Y^\infty\frac{\mathrm{d}y\,\mathrm{d}x}{y^2}=\frac1Y.
	\]
	因此无限距离与有限面积可以同时出现。这里 \(C_Y\) 是带截口的端部；整个完备平移商还包含 \(0<y<Y\) 的另一端，其面积无限。不能把有限面积端部和整个圆柱的面积混为一谈。
\end{example}

\begin{proposition}[双曲圆柱的 Fermi 坐标]\label{g1:prop:hyperbolic-cylinder-fermi}
	设 \(\ell>0\)，令 \(H(z)=e^\ell z\)。完备双曲圆柱 \(\mathbb H/\langle H\rangle\) 有全局坐标
	\[
		r\in\mathbb R,\qquad t\in\mathbb R/\ell\mathbb Z,
		\qquad g=\mathrm{d}r^2+\cosh^2r\,\mathrm{d}t^2.
	\]
	\(r=0\) 是长度 \(\ell\) 的闭测地线，\(|r|\) 是距它的距离；\(r\ge0\) 与 \(r\le0\) 分别给出一个\term[loudou]{漏斗}[funnel]，各漏斗的面积无限。
\end{proposition}
\begin{proof}
	在 \(\mathbb H\) 中直接置
	\[
		x=e^t\tanh r,\qquad y=e^t\operatorname{sech}r.
	\]
	其逆为 \(t=\log\sqrt{x^2+y^2}\)、\(r=\operatorname{arsinh}(x/y)\)，所以这是全局坐标。求微分后交叉项抵消，得到
	\[
		\mathrm{d}x^2+\mathrm{d}y^2
		=e^{2t}\bigl(\mathrm{d}t^2+\operatorname{sech}^2r\,\mathrm{d}r^2\bigr).
	\]
	除以 \(y^2=e^{2t}\operatorname{sech}^2r\) 即得所述度量。缩放 \(z\mapsto e^\ell z\) 恰对应 \(t\mapsto t+\ell\)，故坐标下降到圆柱。群的离散性与自由性由缩放标准型直接得到，完备性由\cref{g1:thm:hyperbolic-quotient}保证。

	\(r=0\) 是虚轴的像，因而为闭测地线，其长度就是 \(t\) 的周期 \(\ell\)。任一从 \(r\) 到零的曲线长度至少为 \(|r|\)，固定 \(t\) 的径向段达到它，因此 \(|r|\) 是到轴的距离。在 \(r\ge0\) 一侧，截圈长度为 \(\ell\cosh r\)，面积元为 \(\cosh r\,\mathrm{d}r\,\mathrm{d}t\)，所以
	\[
		\operatorname{Area}(0\le r\le R)=\ell\sinh R\longrightarrow\infty.
	\]
	另一侧由 \(r\mapsto-r\) 对称得到相同结论。
\end{proof}

上述以测地线为基线、以法向距离 \(r\) 为横坐标的坐标称为\term[Fermizuobiao]{Fermi 坐标}[Fermi coordinates]。\footnote{Enrico Fermi（1901--1954），意大利出生的美国物理学家，研究量子理论、统计力学与核物理。}尖点截圈的长度指数衰减，漏斗截圈的长度按 \(\cosh r\) 增长，这是两种端部面积行为不同的直接原因。\cref{g1:fig:cusp-funnel-domains}把两个循环群的基本区域画在同一个上半平面模型中。

\input{Figures/Chapter10-CuspFunnel.tex}

闭曲面、尖点端和漏斗都具有相同的局部曲率 \(-1\)，它们的紧性、面积和环路长度却不同。模型与粘合给出了这些具体度量；若只知道一张抽象曲面处处局部等距于双曲平面，还不能省略展开及单值的论证就把它写成一个全局商。下一节将把这一环节独立证明。

% !TeX root = ../../Book1.tex
% 纲要标题：### 10.3 局部齐性结构及其单值
% 纲要位置：几何学计划纲要.md，第 522 行。
% 正文状态：初稿；尚未正式审核。

% 纲要内容开始
% 
% * 几何图册与过渡群
% * 展开映射的构造
% * 单值表示及等变性
% * 完备性与模型商的关系
% 
% ---
% 
% 纲要内容结束

\section{局部齐性结构及其单值}
\label{g1:sec:1003}

把一张平纸卷成圆柱，每一小片仍可摊回平面，但沿圆柱绕一圈以后，展开的位置会相差一个平移。这个差别不来自小片的曲率，而来自不同小片怎样拼接。前一节从模型中的群作用构造曲面，本节反过来从曲面上的局部模型出发，追踪一张图沿道路延续时如何变化。所得展开映射与单值表示将说明：什么时候局部模型确实能够拼成整个模型的商，什么时候只能得到一个不完备的局部结构。

\subsection{几何图册与过渡群}
\label{g1:subsec:1003:01}

本节的曲面均为非空、连通、无边界的 Hausdorff、第二可数曲面，沿用\cref{g1:def:topological-surface}的分离与可数条件。无边界的要求有实际作用：模型取整个 \(S^2,\mathbb R^2,\mathbb H^2\) 的开集，半圆盘边缘不在这类图册中。曲面局部道路连通，连通性因而保证道路连通；坐标小圆盘又保证局部单连通，能够使用第五章的万有覆盖。

\ProseParagraphBreak

先取连通光滑模型曲面 \(X\)，以及通过光滑微分同胚作用在 \(X\) 上的群 \(G\)。若只有单位元在整个 \(X\) 上逐点不动，就称该作用有效。本节还要求一个更强的局部确定性质：
\begin{equation}
	\label{g1:eq:model-local-determination}
	g|_O=h|_O,\quad O\subset X\text{ 非空开}
	\quad\Longrightarrow\quad g=h.
\end{equation}
有效性本身不足以推出这一点。一般光滑微分同胚可以在一个开集上为恒等，在远处才发生变化；此时从局部换图无法唯一恢复群元素。对于连通模型的等距群，\cref{g1:prop:isometry-initial-data}说明等距由一点及该点的导数决定，故满足\eqref{g1:eq:model-local-determination}。通常还要求 \(G\) 在 \(X\) 上传递，表示模型各点没有区别；下述延续构造实际使用的是局部确定性质。

\begin{definition}
	\label{g1:def:geometry-atlas}
	设 \(X\) 为连通光滑曲面，\(G\le\operatorname{Diff}(X)\)，并满足局部确定性：若 \(g,h\in G\) 在 \(X\) 的某个非空开集上一致，则 \(g=h\)。设 \(M\) 为无边界拓扑曲面，\(\{U_i\}\) 为它的开覆盖。\(M\) 的一个\term[jihetuce]{几何图册}[geometric atlas]，或 \((G,X)\) 图册，是覆盖 \(M\) 的图
	\[
		\varphi_i:U_i\longrightarrow\varphi_i(U_i)\subset X,
	\]
	其中各 \(\varphi_i\) 为到开集的同胚，而且在任意交叠 \(U_i\cap U_j\) 的每个连通分支 \(C\) 上，存在 \(g_{ij,C}\in G\)，使
	\[
		\varphi_i|_C=g_{ij,C}\circ\varphi_j|_C.
	\]
	两个图册若其并仍满足此条件，就称为相容。一个相容图册类称为 \(M\) 的\term[jubuqixingjiegou]{局部齐性结构}[locally homogeneous structure]或 \((G,X)\) 结构。
\end{definition}

这些图自动给出光滑结构，因为换图映射是光滑微分同胚的限制。交叠处的群元素由\eqref{g1:eq:model-local-determination}唯一确定；在三图的共同小邻域中，
\[
	g_{ik}=g_{ij}g_{jk}.
\]
这里群乘法是函数复合：\((gh)(x)=g(h(x))\)。只在每个交叠分支上指定元素，是为了允许交叠不连通；没有理由要求几个分支使用同一个元素。

\ProseParagraphBreak

若 \(X\) 带度量 \(g_X\)、\(G\le\operatorname{Isom}(X)\)，各片上的 \(\varphi_i^*g_X\) 在交叠处一致，故定义了 \(M\) 上的光滑度量。这是一个具体的粘合验证：\(g_{ij}\) 保持 \(g_X\)，所以
\[
	\varphi_i^*g_X=\varphi_j^*(g_{ij}^*g_X)=\varphi_j^*g_X.
\]
由此得到的曲面每点附近都与模型等距。若过渡群包含非等距变换，例如平面的伸缩，图册仍有意义，却不能用这个公式粘合 Euclid 度量。

常曲率度量是否真的给出这种图册，也需要一个桥梁。这个问题只涉及小邻域，可以直接使用第九章已经建立的测地极坐标。

\begin{proposition}[常曲率的局部模型]
	\label{g1:prop:constant-curvature-local-model}
	设 \(M\) 是无边界光滑度量曲面，Gauss 曲率恒为 \(\kappa\in\{1,0,-1\}\)。则每点附近分别局部等距于单位球面、Euclid 平面或曲率为 \(-1\) 的双曲平面。这些局部等距给出 \((\operatorname{Isom}(X),X)\) 图册；若 \(M\) 已定向，可将过渡群限制为保向等距群。
\end{proposition}
\begin{proof}
	固定 \(p\in M\)。由\cref{g1:prop:normal-neighborhood,g1:lem:gauss-lemma}，中心之外的测地极坐标具有度量
	\[
		\mathrm{d}r^2+a(r,\theta)^2\,\mathrm{d}\theta^2,\qquad a>0.
	\]
	局部取与 \((\partial_r,\partial_\theta)\) 一致的定向。令 \(e_1=\partial_r\)、\(e_2=a^{-1}\partial_\theta\)。度量的 Christoffel 公式给出 \(\nabla_{\partial_r}e_1=0\)、\(\nabla_{\partial_\theta}e_1=a_r e_2\)，所以本书约定的联络形式为 \(\omega=a_r\,\mathrm{d}\theta\)。由\cref{g1:prop:connection-curvature-form}，
	\[
		a_{rr}\,\mathrm{d}r\wedge\mathrm{d}\theta
		=\mathrm{d}\omega=-K a\,\mathrm{d}r\wedge\mathrm{d}\theta,
		\qquad a_{rr}+\kappa a=0.
	\]

	中心处的初值不能只从极坐标公式形式地代入。取 \(\mathrm{T}_pM\) 中的单位向量 \(e(\theta)=(\cos\theta,\sin\theta)\)。第九章的构造实际给出
	\[
		a(r,\theta)=r\,b(r,\theta),\qquad
		b(r,\theta)=\left|\mathrm{d}(\exp_p)_{re(\theta)}(e'(\theta))\right|.
	\]
	指数映射在零附近光滑，且其零点微分为恒等；因而被取范数的向量在零附近不为零，\(b\) 对 \(r,\theta\) 光滑且 \(b(0,\theta)=1\)。所以 \(a\) 从 \(r\ge0\) 一侧光滑延到零，并满足 \(a(0,\theta)=0\)、\(a_r(0,\theta)=1\)。由\cref{g1:lem:smooth-ode-flow}所含的一阶系统唯一性，固定 \(\theta\) 后该标量方程的解只能是
	\[
		a(r,\theta)=S_\kappa(r),\qquad
		S_1(r)=\sin r,\quad S_0(r)=r,\quad S_{-1}(r)=\sinh r.
	\]
	这些函数直接满足方程与初值，且不依赖 \(\theta\)。

	对应模型的曲率同为 \(\kappa\)，对它作上述计算得到相同极坐标度量。选模型点 \(o\) 及线性等距 \(A:\mathrm{T}_pM\to \mathrm{T}_oX\)，则
	\[
		F=\exp_o\circ A\circ\exp_p^{-1}
	\]
	在充分小的正常邻域内光滑。极坐标计算说明它在去掉中心后局部等距；在中心处，两个指数映射的微分均为恒等，故 \(\mathrm{d}F_p=A\) 也等距。这证明局部模型存在，并且没有留下极坐标在中心处的光滑性缺口。

	最后验证过渡图属于整个模型的等距群。在一个连通交叠分支上，过渡映射 \(f\) 是模型开集之间的局部等距。由\cref{g1:prop:model-isometry-data}的初值实现结论，存在全局模型等距 \(h\)，使它在某点的像及导数与 \(f\) 相同。局部等距保持协变导数：把联络经它搬运，仍度量相容且无挠，由\cref{g1:prop:surface-connection-properties}的唯一性就是原联络；因而它保持短测地线。测地线初值唯一性使 \(f,h\) 在该点附近相等。具有相同值及导数的点组成相对闭集，又由同样的短测地线论证为开集；连通性使它等于整个交叠分支。所以过渡映射确是 \(h\) 的限制。若曲面定向，最初选择保向 \(A\)，便只需保向标架与保向 \(h\)。
\end{proof}
\subsection{展开映射的构造}
\label{g1:subsec:1003:02}

固定基点 \(p_0\in M\)，并在其附近选一张初始图 \(\varphi_0\)。两张局部映射若在指定点的某个邻域内相等，称它们在该点有相同的\term[yingsheya]{映射芽}[germ of a map]。这里保留的是附近整张小图，而不只是点的像；点的像相同一般无法确定换图的群元素。我们只延续形如 \(a\circ\varphi_i\)、\(a\in G\) 的局部映射。

\begin{lemma}[图的沿路延续]
	\label{g1:lem:chart-continuation}
	设 \(X\) 为连通光滑曲面，\(G\le\operatorname{Diff}(X)\)，且任意 \(g,h\in G\) 若在某个非空开集上一致就有 \(g=h\)。设无边界曲面 \(M\) 带 \((G,X)\) 结构，\(\alpha:[0,1]\to M\) 连续，并在 \(\alpha(0)\) 给定一个相容的图芽。只允许延续为局部形如 \(a\circ\varphi_i\)、\(a\in G\) 的图芽，其中 \(\varphi_i\) 为结构图册中的图。该图芽可沿 \(\alpha\) 延续，且每个时刻的芽唯一。终点芽只依赖于 \(\alpha\) 的定端点同伦类。
\end{lemma}
\begin{proof}
	先构造延续。用图域覆盖道路像，将其拉回到紧区间；\cref{g1:lem:lebesgue-number}给出有限分割，使每一小段都包含在一张图中。在第一段按给定图芽调整这张图。若前段使用 \(a_i\varphi_i\)，下一段的接点附近满足 \(\varphi_i=g_{ij}\varphi_j\)，就取
	\[
		a_j=a_i g_{ij}.
	\]
	于是 \(a_i\varphi_i=a_j\varphi_j\) 在接点附近相等，两段确实接续。每个段内使用同一局部公式，就得到了沿整条道路的芽。

	为什么不受分割及中途换图影响？取任意两套延续并作公共细分；在每个足够短的参数区间内，把两套公式都换写在同一张图 \(\varphi\) 下，即 \(a\varphi\) 与 \(b\varphi\)。若它们在一个时刻的芽相同，就在非空开集上相等，\eqref{g1:eq:model-local-determination}强制 \(a=b\)，所以在该小段内始终相同。把这一结论沿有限分割推进，全部芽相同。这也证明在道路中插入反向再返回的小段不会改变延续结果。

	再处理同伦，而不是直接把“同伦”当成延续唯一性的另一种说法。设 \(H:[0,1]^2\to M\) 是定端点同伦；用图域的原像覆盖参数方形，取足够细的有限矩形网格，使每个闭小矩形的像落在一张图中。在一张图内，给定起始芽 \(a\varphi\) 后，沿任何道路到同一点都给出这个同一公式的芽。因此，沿一个小矩形的下边再右边，或沿左边再上边，产生相同的终点芽。

	从方形下边的网格路径开始，每次把某个小矩形的两条边替换为另两条边，逐格扫过一行，再扫过全部行，即可把下边连同右侧边替换为左侧边连同上边。每次替换保留延续芽；内部边若先后反向经过，则按上一段消去。左右侧在 \(H\) 下均为常值道路，因为同伦固定端点，它们不改变芽。因此下边与上边的延续有相同终点芽。这里有限网格的用途正是把任意连续同伦转成有限次、可核验的换路操作。
\end{proof}

现在把可能有不同延续结果的道路分开保存。由\cref{g1:thm:path-class-universal-cover}，曲面的万有覆盖可写为
\[
	\widetilde M=\{[\alpha]:\alpha(0)=p_0\},\qquad
	p([\alpha])=\alpha(1),
\]
其中同伦固定两端。记 \(\varphi_\alpha\) 为 \(\varphi_0\) 沿 \(\alpha\) 延续所得的终点图芽，定义
\begin{equation}
	\label{g1:eq:developing-definition}
	\operatorname{dev}([\alpha])=\varphi_\alpha(\alpha(1)).
\end{equation}
前一引理说明它良定义。这一映射称为\term[zhankaiyingshe]{展开映射}[developing map]。\symidx{\(\operatorname{dev}\)：几何结构的展开映射}

\ProseParagraphBreak

它为何是局部微分同胚，也可直接从第五章的道路类拓扑看出。在 \(\alpha(1)\) 取足够小的坐标圆盘 \(U\)，使终点芽由 \(f=a\varphi_i\) 在 \(U\) 上表示。\cref{g1:lem:path-class-basis}给出覆盖片
\[
	B([\alpha],U)=\{[\alpha*\beta]:\beta\text{ 在 }U\text{ 内从 }\alpha(1)\text{ 出发}\}.
\]
在这一片上，继续沿 \(\beta\) 不改变公式 \(f\)，所以 \(\operatorname{dev}=f\circ p\)。覆盖空间的光滑图由 \(p\) 拉回，\(p\) 在这一片上为微分同胚，\(f\) 也是微分同胚，因而展开映射局部光滑可逆。若过渡群保持 \(g_X\)，同一公式还说明 \(\operatorname{dev}\) 对拉回度量是局部等距。这里的“局部”不可省去：即使源空间单连通，展开映射仍可能重叠，也可能只覆盖模型的一部分。

\subsection{单值表示及等变性}
\label{g1:subsec:1003:03}

设 \(\gamma\) 是以 \(p_0\) 为基点的环路。其终点仍为 \(p_0\)，但延续后的图芽可能变成另一张图。该芽与初始图相容，故存在唯一 \(a_\gamma\in G\)，使
\[
	\varphi_\gamma=a_\gamma\circ\varphi_0
\]
在基点附近成立。图芽的同伦不变性说明 \(a_\gamma\) 只依赖 \([\gamma]\)。

\begin{theorem}[展开与单值]
	\label{g1:thm:developing-map}
	设 \(X\) 为连通光滑曲面，\(G\le\operatorname{Diff}(X)\)，且任意 \(g,h\in G\) 在某个非空开集上一致就有 \(g=h\)，即满足\eqref{g1:eq:model-local-determination}。设 \(M\) 为连通无边界曲面，带 \((G,X)\) 结构。固定 \(p_0\in M\) 及其相容初始图芽 \(\varphi_0\)，将从 \(p_0\) 出发的道路按定端点同伦取类，得到万有覆盖 \(\widetilde M\)。以 \(\varphi_\alpha\) 记沿 \(\alpha\) 延续的终点图芽，定义 \(\operatorname{dev}([\alpha])=\varphi_\alpha(\alpha(1))\)。对基点环路 \(\gamma\)，以唯一的 \(a_\gamma\in G\) 满足 \(\varphi_\gamma=a_\gamma\circ\varphi_0\)。则 \(\operatorname{dev}:\widetilde M\to X\) 为局部微分同胚，并给出群同态
	\[
		\rho:\pi_1(M,p_0)\longrightarrow G,\qquad \rho([\gamma])=a_\gamma.
	\]
	对道路类万有覆盖上的左覆盖变换作用 \(L_g[\alpha]=[\gamma*\alpha]\)、\(g=[\gamma]\)，有
	\begin{equation}
		\label{g1:eq:developing-equivariance}
		\operatorname{dev}\circ L_g=\rho(g)\circ\operatorname{dev}.
	\end{equation}
	若初始图芽改为 \(h\circ\varphi_0\)、\(h\in G\)，所得数据变为
	\[
		\operatorname{dev}'=h\circ\operatorname{dev},\qquad
		\rho'(g)=h\rho(g)h^{-1}.
	\]
\end{theorem}
\begin{proof}
	局部微分同胚性质已在构造后证明。必须注意乘法的方向：本书 \(\gamma*\eta\) 先走 \(\gamma\)，再走 \(\eta\)。从 \(\varphi_0\) 沿 \(\gamma\) 走完以后，初始图芽已变成 \(\rho([\gamma])\varphi_0\)。把一个固定的群元素左乘到整套延续公式上，仍是一套延续；故继续沿 \(\eta\) 会得到
	\[
		\rho([\gamma])\rho([\eta])\varphi_0.
	\]
	终点芽的唯一性因此给出 \(\rho([\gamma][\eta])=\rho([\gamma])\rho([\eta])\)。常值环路对应单位元，所以 \(\rho\) 确为同态。

	对于任意从 \(p_0\) 出发的道路 \(\alpha\)，完全相同的延续比较给出
	\[
		\varphi_{\gamma*\alpha}=\rho([\gamma])\varphi_\alpha.
	\]
	在终点评值，即得\eqref{g1:eq:developing-equivariance}。这与\cref{g1:prop:universal-deck-discontinuous}的左作用一致：\(L_g\circ L_k=L_{gk}\)。第五章在覆盖纤维上另有右单值作用，那是“在末端接环路”；这里使用“在前端接环路”的覆盖变换作用，不把两者的乘法混写。

	最后，将初始图改为 \(h\varphi_0\)，所有沿路延续都变为 \(h\varphi_\alpha\)，所以 \(\operatorname{dev}'=h\operatorname{dev}\)。绕 \(g\) 后的芽为 \(h\rho(g)\varphi_0=(h\rho(g)h^{-1})(h\varphi_0)\)，唯一性便给出所述共轭公式。
\end{proof}

同态 \(\rho\) 称为结构的\term[danzhibiaoshi]{单值表示}[holonomy representation]，其像称为结构的单值群。\symidx{\(\rho\)：几何结构的单值表示}它与展开映射一起记录图册延续，而不是仅用一份群表示代替整个几何结构。即使 \(\rho\) 相同，展开映射的像和重叠方式仍可能不同。

\ProseParagraphBreak

第九章的联络的环路平行移动变换 作用在一个切平面上，通过平行移动比较向量；本节的结构单值 则是模型 \(X\) 上的整体变换。两者不能混为同一个旋转角。例如球面本身单连通，标准球面结构的 \(\rho\) 平凡，绕小圆盘的平行移动却可能有非零的曲率转角。反之，下面圆柱的结构单值包含非零平移，但其平坦联络有全局平行标架，联络的环路平行移动变换 恒为恒等。

\begin{example}[平坦圆柱的展开]
	\label{g1:exm:flat-cylinder-development}
	设 \(L>0\)，在 \(\mathbb R^2\) 上令 \(T(x,y)=(x+L,y)\)，并取平坦圆柱 \(M=\mathbb R^2/\langle T\rangle\)。小圆盘投影的逆给出 Euclid 几何图册。万有覆盖是 \(\mathbb R^2\)，选与原坐标相符的初始图，便有
	\[
		\operatorname{dev}(x,y)=(x,y),\qquad
		\rho(n)(x,y)=(x+nL,y),\quad n\in\mathbb Z.
	\]
	绕一圈的延续改变整个位置图，却不改变常向量场 \(\partial_x,\partial_y\)，所以这里能直接看出环路平行移动与结构单值的区别。
\end{example}

\begin{example}[不完备锥面的单值]
	\label{g1:exm:incomplete-holonomy}
	取 \(a>0\)，在 \(M=(0,\infty)\times(\mathbb R/2\pi\mathbb Z)\) 上赋予
	\[
		g=\mathrm{d}r^2+a^2r^2\,\mathrm{d}\theta^2.
	\]
	在每个足够窄的角区间内，\((r,\theta)\mapsto(r\cos(a\theta),r\sin(a\theta))\) 为到 Euclid 平面开集的局部等距；换角度分支只相差一个旋转，所以它给出 Euclid 结构。在万有覆盖 \((0,\infty)\times\mathbb R\) 上，
	\[
		\operatorname{dev}(r,t)=(r\cos(at),r\sin(at)),\qquad
		\rho(n)=R_{2\pi an}.
	\]
	其中 \(R_\beta\) 是平面旋转。展开像是穿孔平面，既不是全部 \(\mathbb R^2\)，展开映射也不单射。沿固定角度取 \(r_j=1/j\)，径向长度估计给出 \(d(r_j,r_k)\le|1/j-1/k|\)，故这是 Cauchy 列，却没有曲面内的极限。

	当 \(a=1\) 时，\(\rho\) 平凡，所以不忠实；当 \(a\) 为无理数时，\(\rho\) 忠实却不离散。后一结论可由抽屉原理直接看出：把 \(0,a,\ldots,Na\) 的小数部分放入 \(N\) 个等长区间，存在 \(1\le n\le N\)，使 \(na\) 到某个整数的距离小于或等于 \(1/N\)。无理性使这些非零小旋转不可能最终固定为同一个群元素，故有非恒等旋转趋于恒等。单值表示作为同态始终存在；离散性与忠实性则需要额外条件。
\end{example}

\subsection{完备性与模型商的关系}
\label{g1:subsec:1003:04}

以后取 \(X=S^2,\mathbb R^2\) 或 \(\mathbb H^2\)，分别带单位球面、Euclid 或曲率为 \(-1\) 的双曲度量，并令 \(G\le\operatorname{Isom}(X)\)。这些模型连通、单连通且度量完备；双曲情形见\cref{g1:thm:hyperbolic-distance-geodesics}，球面与平面的对应性质也在模型部分说明。我们将证明：若曲面自身没有有限距离的缺口，则展开不再重叠或漏掉模型点。

\ProseParagraphBreak

先明确这里的完备。对连通光滑度量曲面 \(N\)，令
\[
	d_N(p,q)=\inf_\alpha L(\alpha),
\]
其中 \(\alpha\) 遍历连接两点的分片 \(C^1\) 曲线。连通曲面上任意两点都能用这样的曲线连接：能够从指定点经有限条坐标直线段到达的点构成非空开闭集，故是全曲面。因而该下确界有限。若每个 \(d_N\)-Cauchy 列都有极限，就称该度量完备。本节不预先把它与测地线能否无限延伸等同。

\begin{proposition}[长度距离与紧性]
	\label{g1:prop:compact-metric-complete}
	连通光滑度量曲面的长度距离确为距离，并诱导原有拓扑。若该曲面紧，则它度量完备。
\end{proposition}
\begin{proof}
	对称性和三角不等式由反转、拼接曲线得到。关键是局部正定性。以一点 \(p\) 为坐标原点，在图内取闭圆盘 \(\overline B_{2r}\)。光滑正定矩阵在这个紧集上有一致上下界：存在 \(0<c\le C\)，使
	\[
		c|v|\le\sqrt{g_x(v,v)}\le C|v|.
	\]
	若 \(x\in B_r\)，坐标直线段给出 \(d_N(p,x)\le C|x|\)。反向估计时，留在 \(B_{2r}\) 的任意连接曲线长度至少为 \(c|x|\)；离开 \(B_{2r}\) 的曲线仅走到坐标半径 \(2r\) 就已长至少 \(2cr\)，仍不小于 \(c|x|\)。故 \(d_N(p,x)\ge c|x|\)。同理，任何从 \(p\) 到 \(B_r\) 外的曲线长度至少为 \(cr\)。这既说明不同点的距离为正，也说明度量小球与坐标邻域互相包含，拓扑相同。

	若曲面紧，它在该距离下是紧度量空间。每个序列都有收敛子列：若没有聚点，则各点都有只含有限多个序列项的邻域，有限子覆盖会迫使整个序列只有有限多个项，矛盾；取聚点的逐次缩小距离球便得到收敛子列。Cauchy 列一旦有收敛子列，三角不等式使整个序列收敛。因此紧曲面度量完备。
\end{proof}

\begin{proposition}[覆盖与完备性]
	\label{g1:prop:complete-pullback-cover}
	设 \((N,g_N)\) 为连通光滑度量曲面，\(p:Y\to N\) 为连通拓扑覆盖。在 \(Y\) 上取使 \(p\) 成为局部微分同胚的拉回光滑结构及拉回度量 \(g_Y=p^*g_N\)。则 \(Y\) 度量完备当且仅当 \(N\) 度量完备。
\end{proposition}
\begin{proof}
	覆盖局部等距，故保留分片光滑曲线的长度，且 \(d_N(p(y),p(z))\le d_Y(y,z)\)。先设 \(N\) 完备，\((y_j)\) 为 \(Y\) 中的 Cauchy 列。投影列收敛于某点 \(x\)。取 \(x\) 的均匀覆盖邻域 \(U\)，再选 \(r>0\) 使 \(B_N(x,3r)\subset U\)，这可由前一命题的拓扑相容性保证。

	在尾部有 \(p(y_j)\in B_N(x,r)\)，且 \(d_Y(y_j,y_k)<r\)。固定其中一个 \(y_{j_0}\)。对每个后续 \(j\)，按下确界定义取从 \(y_{j_0}\) 到 \(y_j\) 的曲线，长度小于 \(r\)。其投影始终在 \(B_N(x,2r)\subset U\) 中；曲线连通，故留在 \(p^{-1}(U)\) 的同一张覆盖片上。于是全部尾项属于同一片。该片的逆映射把 \(p(y_j)\to x\) 提为 \(y_j\to y\)，并由距离与拓扑相容性得到距离收敛。这证明 \(Y\) 完备。

	反过来，设 \(Y\) 完备，\((x_j)\) 是 \(N\) 中的 Cauchy 列。选子列 \((x_{j_k})\)，使相邻距离小于 \(2^{-k-2}\)，并以长度小于 \(2^{-k-1}\) 的分片光滑曲线相连。任取 \(x_{j_1}\) 的一个原像，依次使用\cref{g1:thm:path-lifting}提升这些曲线。局部逆光滑，故提升仍分片光滑且长度不变。所得端点 \(y_k\) 满足
	\[
		d_Y(y_k,y_\ell)\le\sum_{i=k}^{\ell-1}2^{-i-1}\longrightarrow0
		\quad(k,\ell\to\infty).
	\]
	它们收敛于 \(y\in Y\)，投影子列遂收敛于 \(p(y)\)。原列 Cauchy，故也收敛于 \(p(y)\)。
\end{proof}

这个命题已经处理了一个容易遗漏的环节：曲面完备以后，它的万有覆盖配拉回度量也完备，不能只因每张覆盖片局部等距就略去这一证明。

\begin{lemma}[完备局部等距]
	\label{g1:lem:complete-local-isometry}
	设 \(f:Y\to X\) 是连通无边界光滑度量曲面间的局部等距，且 \(Y\) 度量完备。则 \(f\) 为满射覆盖。
\end{lemma}
\begin{proof}
	先沿道路提升，再由统一的模型小球构造覆盖片。设 \(c:[0,1]\to X\) 为分片 \(C^1\) 道路，\(f(y_0)=c(0)\)。局部逆从 \(y_0\) 给出短时间提升；有同一起点的两条提升在相等处可用同一个局部逆继续，故唯一。把所有连续相接的局部提升合并，得到一个最大时间区间。若它在 \(b\le1\) 以前停止，则对 \(s,t\uparrow b\)，
	\[
		d_Y(\widetilde c(s),\widetilde c(t))
		\le L(\widetilde c|_{[s,t]})=L(c|_{[s,t]})\longrightarrow0.
	\]
	这里尾段长度趋零来自原道路分片 \(C^1\)，并不要求它是测地线。完备性给出极限 \(y_b\)，连续性使 \(f(y_b)=c(b)\)；在 \(y_b\) 处使用局部逆就能继续提升，并覆盖 \(b\) 本身。因此提升存在到时间 \(1\)。由于 \(X\) 连通，任一点可经分片光滑道路从 \(f(y_0)\) 到达，提升终点说明 \(f\) 满射。

	固定 \(x\in X\)。由\cref{g1:prop:normal-neighborhood,g1:thm:radial-minimizing}，取一个小正常球 \(B\)；其中每个 \(z\) 都有从 \(x\) 出发的径向测地线 \(c_z\)，而且 \(c_z\) 随 \(z\) 光滑变化。对每个 \(y\in f^{-1}(x)\)，令 \(s_y(z)\) 为 \(c_z\) 从 \(y\) 出发的提升终点。刚才的延拓保证它对所有 \(z\in B\) 定义，且 \(f\circ s_y=\operatorname{id}_B\)。

	还须确认 \(s_y\) 光滑，才能说其像是一张开覆盖片。固定 \(z_0\)，沿已得的紧提升曲线选有限个局部逆，作时间分割，使每段留在相应局部逆的像中。\(c_z\) 对 \(z\) 连续且时间区间紧，故当 \(z\) 足够接近 \(z_0\) 时，这一组有限局部逆仍能逐段使用。每段的提升由光滑局部逆与 \(c_z\) 复合给出，在接点的选片由前段终点确定；所以最终的 \(s_y(z)\) 光滑。对 \(f\circ s_y=\operatorname{id}\) 求导，其导数可逆，故 \(s_y\) 在每点附近为局部逆；从而 \(s_y(B)\) 开，且 \(f:s_y(B)\to B\) 为微分同胚。

	这些像两两不交。若 \(s_y(z)=s_{y'}(z)\)，反向提升同一条径向线，唯一性迫使 \(y=y'\)。它们也穷尽 \(f^{-1}(B)\)：从任意 \(u\in f^{-1}(z)\) 反向提升 \(c_z\)，到达某个 \(y\in f^{-1}(x)\)，再反过来就有 \(u=s_y(z)\)。因此
	\[
		f^{-1}(B)=\bigsqcup_{y\in f^{-1}(x)}s_y(B)
	\]
	是所需的均匀覆盖分解。
\end{proof}

\begin{theorem}[完备结构的模型商]
	\label{g1:thm:complete-model-quotient}
	取 \(X\) 为带标准度量的单位球面 \(S^2\)、Euclid 平面 \(\mathbb R^2\) 或曲率 \(-1\) 的双曲平面 \(\mathbb H^2\)，\(G\le\operatorname{Isom}(X)\)。设连通无边界曲面 \(M\) 带有 \((G,X)\) 结构，并取由图册拉回的度量。固定 \(p_0\in M\) 及其相容初始图芽，以 \(\operatorname{dev}:\widetilde M\to X\)、\(\rho:\pi_1(M,p_0)\to G\) 记相应的展开映射及单值表示。若 \(M\) 度量完备，则其展开映射
	\[
		\operatorname{dev}:\widetilde M\longrightarrow X
	\]
	为等距微分同胚，单值表示 \(\rho\) 忠实，其像 \(\Gamma=\rho(\pi_1(M,p_0))\) 是自由固有作用的离散等距群，且
	\[
		M\cong X/\Gamma
	\]
	为保持局部几何结构的等距同构。反过来，完备模型上任意自由固有等距群作用的曲面商，带商度量后完备，并具有相应的模型结构。
\end{theorem}
\begin{proof}
	由\cref{g1:prop:complete-pullback-cover}，\(\widetilde M\) 的拉回度量完备；展开映射是局部等距，故前一引理使它成为到 \(X\) 的满射覆盖。因为 \(X\) 单连通，这个连通覆盖只能有一叶。具体地，若两个原像 \(y_1,y_2\) 映到同一点，取连接它们的道路；其投影为零伦环路，与常值道路定端点同伦。\cref{g1:cor:lifted-path-homotopy}使二者从 \(y_1\) 提升后有相同终点，所以 \(y_2=y_1\)。因此展开映射为微分同胚。它及其逆都保持曲线长度，取下确界后即为全局等距。

	等变性把左覆盖变换作用搬到 \(X\) 上，变成 \(\rho(g)\) 的作用。若 \(\rho(g)=1\)，展开映射单射迫使 \(L_g=1\)，而万有覆盖的左作用自由，故 \(g=1\)。这证明忠实性和 \(\Gamma\) 作用的自由性。覆盖变换作用的轨道正是覆盖纤维，故等变微分同胚下降为 \(M\cong X/\Gamma\)，且图册与度量吻合。

	再核实固有性与离散性，而不把它们自动附在“群商”之后。这里对离散群的固有作用，使用紧集判据：每个紧集 \(K\subset X\) 只有有限个 \(\gamma\in\Gamma\) 满足 \(\gamma K\cap K\ne\varnothing\)。若覆盖变换作用违反这一点，便有不同的覆盖变换 \(L_{g_j}\) 及 \(u_j,v_j\) 落在同一紧集内，满足 \(L_{g_j}u_j=v_j\)。取收敛子列，设 \(u_j\to u\)、\(v_j\to v\)。投影相同使 \(p(u)=p(v)\)。在该投影点取连通均匀覆盖邻域；充分靠后的 \(u_j,v_j\) 分别落在包含 \(u,v\) 的两张固定覆盖片中。覆盖变换若把第一片送到第二片，便由投影在片上的单射性唯一确定。因此这些 \(L_{g_j}\) 最终相同，与选择矛盾。等距识别遂给出 \(\Gamma\) 的固有性。由\cref{g1:prop:universal-deck-discontinuous}，又可取 \(x\) 的邻域 \(V\)，使 \(\gamma x\notin V\) 对所有非恒等 \(\gamma\) 成立。在等距群的紧开拓扑中，条件 \(h(x)\in V\) 是恒等元的开邻域，它与 \(\Gamma\) 只交于单位元，所以 \(\Gamma\) 离散。

	反向设 \(\Gamma\) 在完备 \(X\) 上自由且满足紧集判据。每点可取闭包紧的坐标邻域；仅有有限个 \(\gamma\) 可能使该邻域与其像相交。自由性使这些非恒等 \(\gamma\) 都不固定中心点，逐个缩小邻域即可排除全部有限个相交，于是得到与所有非恒等平移像不交的小邻域。\cref{g1:prop:action-quotient-cover}因此使 \(X\to X/\Gamma\) 成为覆盖，小片的逆图给出 \((G,X)\) 结构，等距条件使模型度量下降。商空间为 Hausdorff：对不同轨道的两点，在各自紧邻域间只有有限个可能相交的群平移，缩小两邻域排除它们，便得不交的商邻域。投影为开映射，将 \(X\) 的可数基投影即可得到商的可数基。所以商确为所要求的光滑曲面。最后将\cref{g1:prop:complete-pullback-cover}用于这个局部等距覆盖，由 \(X\) 完备得到商完备。
\end{proof}

在上述商表示中，曲面 \(M\) 紧恰好等价于 \(\Gamma\) 余紧，因为余紧要求商 \(X/\Gamma\) 紧。平坦圆柱是完备但不紧的例子，其平移群也不余紧。特别地，紧的无边界模型曲面自动满足完备性，因而其万有覆盖就是相应模型。例如前一节已经用多边形的小邻域验证了闭双曲曲面的局部结构；紧性加上本定理，才进一步把它识别为 \(\mathbb H^2/\Gamma\)。这一步既不把所有拓扑曲面都预先宣布为常曲率商，也不从单值表示的存在直接推出离散性。缺少完备性时，\cref{g1:exm:incomplete-holonomy}展示的漏点、重叠及非离散单值都可能发生。
% 纲要条目排为编号小节标题，并进入目录与书签；说明段仍保留于纲要注释中。


\begin{chaptersummary}
	双曲面、圆盘与上半平面是同一个曲率为 \(-1\) 的几何的不同坐标表达。双曲面上的双线性型给出距离，上半平面的实系数 Möbius 变换给出保向等距群，圆盘便于观察理想边界。球面与双曲三角公式的差别体现曲率符号；角和与面积的关系则由第九章的 Gauss--Bonnet 公式统一起来。

	\ProseParagraphBreak

	椭圆、抛物与双曲等距变换有不同的固定点和位移行为。离散性保证双曲等距群作用的固有性，自由性排除商空间中的稳定子奇性。基本多边形的粘合还须检查每个顶点的链接与总角：角拼满一周，局部才恢复光滑双曲圆盘。尖点的横截周长向无穷远收缩，漏斗的横截周长则增大，二者的面积行为因而不同。

	\ProseParagraphBreak

	几何图册通过指定变换群限制坐标变化。沿道路调整相邻坐标，可在万有覆盖上得到展开映射；同伦不变性给出单值表示，等变关系把两者联系起来。不完备时，展开映射可能不单射，单值群也可能不离散。对完备的常曲率模型结构，有限长路径的提升不会在中途失去终点，展开映射最终成为整个单连通模型的等距识别，曲面因此由自由固有等距作用取商得到。
\end{chaptersummary}

\BookExercises

\begin{exercise}
	\label{g1:ex:hyperbolic-metric-balls}
	在曲率 \(-1\) 的上半平面中，固定 \(z_0=x_0+iy_0\) 与 \(r>0\)。
	\begin{enumerate}
		\item 证明以 \(z_0\) 为双曲中心、半径为 \(r\) 的闭球，在 Euclid 坐标中为
		\[
			(x-x_0)^2+(y-y_0\cosh r)^2\leq y_0^2\sinh^2r.
		\]
		说明为什么它的 Euclid 中心一般不是 \(z_0\)。
		\item 求这个双曲球的面积与边界长度，并比较它们在 \(r\to0\) 及 \(r\to\infty\) 时的增长。求圆盘模型中双曲中心为 \(0\) 的同半径球的 Euclid 半径。
	\end{enumerate}
\end{exercise}
\begin{solution}
	1. 由\eqref{g1:eq:half-plane-distance}及 \(\cosh\) 在非负半轴上严格递增，\(d(z,z_0)\leq r\) 等价于
	\[
		(x-x_0)^2+(y-y_0)^2\leq2yy_0(\cosh r-1).
	\]
	把含 \(y\) 的项配方，即得所求圆盘。它的最低点为 \(y_0(\cosh r-\sinh r)=y_0e^{-r}>0\)，所以整个 Euclid 闭圆盘都留在上半平面。Euclid 中心的高度是 \(y_0\cosh r\)，大于 \(y_0\)：模型中的长度放大因子 \(1/y\) 随位置改变，度量对称性不等于坐标图中的 Euclid 对称性。

	2. 双曲等距可以把中心送到 \(o\)。由\eqref{g1:eq:hyperbolic-polar-metric}，面积元为 \(\sinh\rho\,\mathrm{d}\rho\,\mathrm{d}\theta\)，半径 \(r\) 圆周上的长度元为 \(\sinh r\,\mathrm{d}\theta\)，故
	\[
		A(r)=2\pi(\cosh r-1),\qquad L(r)=2\pi\sinh r.
	\]
	当 \(r\to0\) 时，\(A(r)=\pi r^2+O(r^4)\)、\(L(r)=2\pi r+O(r^3)\)，最低阶与 Euclid 圆相同。当 \(r\to\infty\) 时，\(A(r)/e^r\to\pi\)、\(L(r)/e^r\to\pi\)，两者都指数增长。圆盘模型中 \(d(0,w)=2\operatorname{artanh}|w|\)，所以 Euclid 半径为 \(\tanh(r/2)\)。它始终小于 \(1\)，尽管双曲半径可以任意大。
\end{solution}

\begin{exercise}
	\label{g1:ex:klein-metric-angles}
	令 \(Q:\mathbb H_L^2\to\mathbb D\) 为 \(Q(x_0,x_1,x_2)=(x_1/x_0,x_2/x_0)\)，在其像中用实向量 \(u=(u_1,u_2)\) 作坐标。
	\begin{enumerate}
		\item 由逆映射计算 Klein 模型的度量，证明
		\[
			g_K=\frac{(1-|u|^2)|\mathrm{d}u|^2+(u\cdot\mathrm{d}u)^2}{(1-|u|^2)^2}.
		\]
		\item 在 \(u=(a,0)\)、\(0<a<1\) 处，求向量 \(v=(1,1)\)、\(w=(1,-1)\) 的度量夹角，说明直线弦模型为什么不能直接用 Euclid 量角器量双曲角。
	\end{enumerate}
\end{exercise}
\begin{solution}
	1. 记 \(q=1-|u|^2>0\)、\(b=u\cdot\mathrm{d}u\)。由 \(Q^{-1}(u)=(1,u)/\sqrt q\)，得
	\[
		\mathrm{d}x_0=q^{-3/2}b,\qquad
		\mathrm{d}x'=q^{-1/2}\mathrm{d}u+q^{-3/2}ub.
	\]
	代入 \(-\mathrm{d}x_0^2+|\mathrm{d}x'|^2\)，得到
	\[
		q^{-1}|\mathrm{d}u|^2+2q^{-2}b^2+q^{-3}(|u|^2-1)b^2
		=q^{-1}|\mathrm{d}u|^2+q^{-2}b^2,
	\]
	即所求公式。分子对非零切向量为正，亦可直接看出这是正定度量。

	2. 在该点，度量矩阵为 \(\operatorname{diag}((1-a^2)^{-2},(1-a^2)^{-1})\)。因此
	\[
		g_K(v,w)=\frac{a^2}{(1-a^2)^2},\qquad
		g_K(v,v)=g_K(w,w)=\frac{2-a^2}{(1-a^2)^2}.
	\]
	夹角 \(\vartheta\) 满足 \(\cos\vartheta=a^2/(2-a^2)>0\)，故为锐角；两向量的 Euclid 内积却是零。Klein 模型把测地线变直的同时，使不同方向的长度缩放不一致，因此不保 Euclid 角。它与共形圆盘模型各自简化了不同的几何信息。
\end{solution}

\begin{exercise}
	\label{g1:ex:hyperbolic-reflection-distance}
	在上半平面中考虑测地线 \(\ell=\{iy:y>0\}\)。
	\begin{enumerate}
		\item 对点 \(z=x+iy\)，求它到 \(\ell\) 的距离及唯一最近点。证明连接点与最近点的测地段垂直于 \(\ell\)。
		\item 令 \(\sigma(z)=-\overline z\)。证明 \(d(z,\sigma(z))=2d(z,\ell)\)。对圆心为 \(c\in\mathbb R\)、半径为 \(r>0\) 的测地半圆，证明
		\[
			\tau(z)=c+\frac{r^2}{\overline z-c}
		\]
		是逐点固定该半圆的反向等距，并说明同样的两倍距离结论也成立。
	\end{enumerate}
\end{exercise}
\begin{solution}
	1. 对 \(t>0\)，距离公式化为
	\[
		\cosh d(z,it)=\frac{x^2+y^2+t^2}{2yt}
		=\frac1{2y}\left(t+\frac{x^2+y^2}{t}\right).
	\]
	正数的算术几何平均不等式给出唯一极小点 \(t_0=\sqrt{x^2+y^2}\)，因而最近点为 \(it_0\)，距离满足
	\[
		\cosh d(z,\ell)=\frac{\sqrt{x^2+y^2}}y,
		\qquad d(z,\ell)=\operatorname{arsinh}\frac{|x|}y.
	\]
	若 \(x\ne0\)，连接 \(z\) 与 \(it_0\) 的测地线是以 \(0\) 为圆心、半径 \(t_0\) 的半圆；在 \(it_0\) 处其切线水平，与竖直的 \(\ell\) 正交。模型共形，故 Euclid 正交也是度量正交。若 \(x=0\)，点已在 \(\ell\) 上，距离为零，此时没有需要指定切向方向的非退化垂线段。

	2. 直接代入距离公式，
	\[
		\cosh d(z,\sigma(z))=1+\frac{2x^2}{y^2}
		=\cosh\left(2\operatorname{arsinh}\frac{|x|}y\right).
	\]
	两边的距离均非负，\(\cosh\) 在非负半轴上单射，故得两倍关系。

	对于 \(\tau\)，有
	\[
		\operatorname{Im}\tau(z)=\frac{r^2y}{|z-c|^2},\qquad
		|\mathrm{d}\tau|=\frac{r^2}{|z-c|^2}|\mathrm{d}z|.
	\]
	所以它保持双曲度量；共轭使方向反转，且直接代入得 \(\tau^2=\operatorname{id}\)。固定点条件为 \(|z-c|^2=r^2\)，恰是该测地半圆。也可取实分式变换
	\[
		f(z)=\frac{z-(c-r)}{(c+r)-z}.
	\]
	其矩阵行列式为 \(2r>0\)，除以 \(\sqrt{2r}\) 后即属 \(\operatorname{SL}(2,\mathbb R)\)，所以它是保向等距，并把半圆的两个理想端点送到 \(0,\infty\)，从而把该半圆送到 \(\ell\)。共轭 \(f\tau f^{-1}\) 是固定 \(\ell\) 的反向等距。在 \(i\) 处它固定沿 \(\ell\) 的单位向量，而反转其垂直方向，故与 \(\sigma\) 有相同初始数据；\cref{g1:prop:isometry-initial-data}使二者相同。由等距保持到集合的距离，第一问及两倍距离关系遂传回原半圆。
\end{solution}

\begin{exercise}
	\label{g1:ex:equilateral-curvature-triangles}
	分别研究曲率 \(-1\) 的双曲平面与单位球面的等边测地三角形。球面三角形要求非退化、凸且包含在一个开半球内。
	\begin{enumerate}
		\item 证明双曲情形对每个边长 \(a>0\) 都有这样的三角形，且内角 \(\alpha\) 满足 \(\cos\alpha=\cosh a/(1+\cosh a)\)。求内角和面积在 \(a\to0\)、\(a\to\infty\) 时的极限。
		\item 证明球面情形的允许边长恰为 \(0<a<2\pi/3\)，且 \(\cos\alpha=\cos a/(1+\cos a)\)。求面积的变化范围，并解释上端极限的几何形状。
	\end{enumerate}
\end{exercise}
\begin{solution}
	1. 在双曲面的半径 \(R\) 圆周上取径向角相差 \(2\pi/3\) 的三个点。相邻点距离的 Lorentz 公式给出
	\[
		\cosh a=\cosh^2R-\sinh^2R\cos\frac{2\pi}3
		=1+\frac32\sinh^2R.
	\]
	对每个 \(a>0\)，取 \(R=\operatorname{arsinh}\sqrt{2(\cosh a-1)/3}>0\) 即得等边三角形。旋转对称使三个内角相等；也可从三个边余弦公式看出其余弦相等。将 \(a=b=c\) 代入\eqref{g1:eq:hyperbolic-cosine-law}，约去 \(\cosh a-1>0\)，得到所求角公式。

	当 \(a\) 从零增向无穷时，\(\cosh a/(1+\cosh a)\) 从 \(1/2\) 严格增向 \(1\)，所以 \(\alpha\) 从 \(\pi/3\) 严格减向零。面积 \(\pi-3\alpha\) 因而从零严格增向 \(\pi\)，每个有限边长对应的面积仍小于 \(\pi\)。后一个极限与理想三角形的面积相符，理想顶点自身则不是有限边长三角形的顶点。

	2. 若三个球面顶点为单位向量 \(A,B,C\)，等边条件使两两内积均为 \(\cos a\)，于是
	\[
		|A+B+C|^2=3+6\cos a\geq0.
	\]
	因较短边有 \(0<a<\pi\)，这迫使 \(a\leq2\pi/3\)。等号时 \(A+B+C=0\)，三个向量共面，三个顶点落在同一个大圆上，不是所需非退化三角形，故必须严格小于该上限。

	反过来，给定 \(0<a<2\pi/3\)，令 \(h=\sqrt{(1+2\cos a)/3}\in(0,1)\)，在高度 \(h\) 的纬圆上取
	\[
		V_j=\left(\sqrt{1-h^2}\cos\frac{2\pi j}3,
		\sqrt{1-h^2}\sin\frac{2\pi j}3,h\right),\qquad j=0,1,2.
	\]
	两两内积为 \((3h^2-1)/2=\cos a\)，所以较短大圆边长都是 \(a\)。三点都在北半球，较短大圆弧也在其中；向平面 \(z=1\) 的中心投影显示三条边围成凸三角形。球面余弦公式给出
	\[
		\cos\alpha=\frac{\cos a-\cos^2a}{1-\cos^2a}
		=\frac{\cos a}{1+\cos a}.
	\]
	随着 \(a\) 从零增到 \(2\pi/3\)，\(\alpha\) 从 \(\pi/3\) 严格增到 \(\pi\)，面积 \(3\alpha-\pi\) 从零严格增到 \(2\pi\)。上端时 \(h\to0\)，三条较短大圆弧趋向拼成整个赤道，凸区域趋向一个半球；顶点角趋于 \(\pi\)，所以这个极限已经不再是非退化球面三角形。
\end{solution}

\begin{exercise}\label{g1:ex:psl-parameter-family}
	设 \(s,t\in\mathbb R\)，考察代表矩阵
	\[
		A_{s,t}=\begin{pmatrix}1&s\\t&1+st\end{pmatrix}
	\]
	给出的上半平面等距变换。
	\begin{enumerate}
		\item 用 \(u=st\) 判别其类型，并单独处理恒等变换。
		\item 椭圆情形求出平面内不动点；双曲情形求出平移长度。
		\item 在抛物情形分别讨论 \(u=0\) 与 \(u=-4\)，求出理想不动点。
	\end{enumerate}
\end{exercise}
\begin{solution}
	1. 矩阵的行列式为 \(1\)，迹为 \(2+u\)，故
	\[
		(\operatorname{tr}A_{s,t})^2-4=u(u+4).
	\]
	当 \(-4<u<0\) 时为椭圆变换；当 \(u>0\) 或 \(u<-4\) 时为双曲变换。当 \(u=0\) 或 \(u=-4\) 时，除恒等变换外均为抛物变换。恒等只发生在 \(s=t=0\)：因为此矩阵左上角为 \(1\)，它不可能等于 \(-I\)，而等于 \(I\) 恰要求 \(s=t=0\)。

	2. 不动点方程为 \(tz^2+uz-s=0\)。椭圆情形 \(t\ne0\)，取虚部为正的根，得到
	\[
		z=-\frac s2+i\frac{\sqrt{-u(u+4)}}{2|t|}.
	\]
	双曲情形由\cref{g1:thm:psl-element-classification}及平移长度的计算得到
	\[
		L(A_{s,t})=2\operatorname{arcosh}\frac{|2+u|}{2}.
	\]
	这里绝对值不可删去，因为 \(u<-4\) 时所选代表矩阵的迹为负，而 \(A\) 与 \(-A\) 表示同一等距变换。

	3. 若 \(u=0\) 且 \(t=0,s\ne0\)，变换为 \(z\mapsto z+s\)，唯一理想不动点是 \(\infty\)。若 \(u=0\) 且 \(s=0,t\ne0\)，变换为 \(z\mapsto z/(tz+1)\)，唯一理想不动点是 \(0\)。若 \(u=-4\)，则 \(s,t\) 均非零，不动点方程具有二重根 \(z=2/t=-s/2\)。
\end{solution}

\begin{exercise}\label{g1:ex:cyclic-dirichlet-domains}
	以 \(i\in\mathbb H\) 为中心，分别求下列群的 Dirichlet 区域：\(\langle z\mapsto z+1\rangle\) 与 \(\langle z\mapsto e^\ell z\rangle\)，其中 \(\ell>0\)。要求从距离不等式导出全部边界，并解释为何只需比较生成元与其逆元，而不必另添所有高次幂给出的边界。
\end{exercise}
\begin{solution}
	对平移群，轨道点为 \(n+i\)，其中 \(n\in\mathbb Z\)。距离公式给出
	\[
		d(z,i)\le d(z,n+i)
		\quad\Longleftrightarrow\quad x^2\le(x-n)^2
		\quad\Longleftrightarrow\quad 2nx\le n^2.
	\]
	对所有正整数 \(n\)，条件为 \(x\le n/2\)，最强者是 \(n=1\)；对所有负整数，最强者是 \(n=-1\) 给出的 \(x\ge-1/2\)。所以区域为 \(-1/2\le x\le1/2,y>0\)，两条边是竖直测地线。

	对缩放群，轨道点为 \(ai\)，其中 \(a=e^{n\ell}\)。由
	\[
		\cosh d(z,ai)=\frac{x^2+y^2+a^2}{2ay}
	\]
	得到
	\[
		d(z,i)\le d(z,ai)
		\quad\Longleftrightarrow\quad(a-1)|z|^2\le a(a-1).
	\]
	当 \(n>0\) 时，这要求 \(|z|^2\le e^{n\ell}\)，最强者是 \(n=1\)；当 \(n<0\) 时，要求 \(|z|^2\ge e^{n\ell}\)，最强者是 \(n=-1\)。故区域为
	\[
		e^{-\ell/2}\le|z|\le e^{\ell/2},\qquad y>0.
	\]
	两条边是以实轴原点为中心的半圆测地线，由缩放 \(e^\ell\) 配对。两种区域都只有两条边，却都不是紧集；有限边数本身不能推出商紧。
\end{solution}

\begin{exercise}\label{g1:ex:hyperbolic-polygon-cone-angle}
	固定 \(g\ge2\)，取任意 \(0<\alpha<(4g-2)\pi/(4g)\)，将内角为 \(\alpha\) 的正则双曲 \(4g\) 边形按标准边字 \(a_1b_1a_1^{-1}b_1^{-1}\cdots a_gb_ga_g^{-1}b_g^{-1}\) 配对。
	\begin{enumerate}
		\item 说明所得空间的拓扑类型不随 \(\alpha\) 改变，并求唯一可能的度量奇点的总角 \(\Theta\)。
		\item 求面积，并说明何时度量在该点光滑；当 \(\Theta=4\pi\) 时，绕行一周的方向又返回原方向，这是否足以保证光滑？
		\item 验证面积、角亏与 Euler 示性数满足
		\[
			-\operatorname{Area}+(2\pi-\Theta)=2\pi(2-2g).
		\]
	\end{enumerate}
\end{exercise}
\begin{solution}
	1. 边字与端点配对完全相同；\cref{g1:thm:closed-hyperbolic-polygon}证明中的顶点识别及链接循环不依赖 \(\alpha\)。因此所有顶点仍接成同一个点，链接仍是一条圆周，拓扑类型仍为 \(\Sigma_{g,0}\)。沿边内点的粘合始终是光滑等距粘合，唯一可能的度量奇点来自这个顶点，其总角为 \(\Theta=4g\alpha\)。

	2. 由测地多边形面积公式，
	\[
		\operatorname{Area}=(4g-2)\pi-4g\alpha.
	\]
	顶点附近由总角为 \(\Theta\) 的双曲扇环组成，其半径 \(r\) 圆周长度为 \(\Theta\sinh r\)。光滑度量必须使圆周长度除以 \(r\) 趋于 \(2\pi\)，故必要条件是 \(\Theta=2\pi\)，即 \(\alpha=\pi/(2g)\)；\cref{g1:prop:hyperbolic-polygon-gluing}说明这个条件也充分。若 \(\Theta=4\pi\)，相应角度在旋转群中虽与零角等价，但小圆周长度与半径之比趋于 \(4\pi\)，中心仍是锥形奇点。方向返回不意味着局部度量只铺满了一张普通圆盘。

	3. 代入前两项得到
	\[
		-\bigl((4g-2)\pi-4g\alpha\bigr)+(2\pi-4g\alpha)
		=2\pi(2-2g).
	\]
	这里面积曲率为 \(-1\)，角亏 \(2\pi-\Theta\) 则集中在顶点；两项随 \(\alpha\) 的变化相互抵消，而拓扑不变。
\end{solution}

\begin{exercise}\label{g1:ex:cusp-funnel-boundary-curvature}
	分别在尖点和漏斗中截取以下紧环带：
	\[
		C_R:\quad 0\le r\le R,\quad x\in\mathbb R/\mathbb Z,
		\quad g=\mathrm{d}r^2+Y^{-2}e^{-2r}\,\mathrm{d}x^2,
	\]
	\[
		F_R:\quad 0\le r\le R,\quad t\in\mathbb R/\ell\mathbb Z,
		\quad g=\mathrm{d}r^2+\cosh^2r\,\mathrm{d}t^2,
	\]
	其中 \(R,Y,\ell>0\)。按指向环带内侧的法向计算两条边界的测地曲率及其积分，直接验证两种环带的 Gauss--Bonnet 公式。
\end{exercise}
\begin{solution}
	先统一计算 \(\mathrm{d}r^2+f(r)^2\,\mathrm{d}u^2\) 的平行圈。单位切向量为 \(T=f^{-1}\partial_u\)，由 \(\Gamma^r_{uu}=-ff'\) 得
	\[
		\nabla_TT=-\frac{f'}f\partial_r.
	\]
	下边界 \(r=0\) 的内侧法向为 \(+\partial_r\)，上边界 \(r=R\) 的内侧法向为 \(-\partial_r\)，所以两者的测地曲率分别为 \(-f'(0)/f(0)\) 与 \(f'(R)/f(R)\)。

	尖点中 \(f(r)=Y^{-1}e^{-r}\)，故下、上边界测地曲率分别为 \(+1,-1\)，长度分别为 \(Y^{-1},Y^{-1}e^{-R}\)。总边界积分为 \(Y^{-1}(1-e^{-R})\)，恰等于环带面积。由于 \(K=-1\)，有
	\[
		\int_{C_R}K\,\mathrm{d}A+\int_{\partial C_R}k_g\,\mathrm{d}s
		=-Y^{-1}(1-e^{-R})+Y^{-1}(1-e^{-R})=0.
	\]

	漏斗中 \(f(r)=\cosh r\)，\(u=t\) 的周期为 \(\ell\)。下边界测地曲率为零；上边界为 \(\tanh R\)，长度为 \(\ell\cosh R\)，总边界积分为 \(\ell\sinh R\)。其面积也是 \(\ell\sinh R\)，所以
	\[
		\int_{F_R}K\,\mathrm{d}A+\int_{\partial F_R}k_g\,\mathrm{d}s
		=-\ell\sinh R+\ell\sinh R=0.
	\]
	两种环带的边界均无角点，Euler 示性数均为零，正与\cref{g1:thm:global-gauss-bonnet}相符。
\end{solution}

\begin{exercise}
	\label{g1:ex:strip-cylinder-development}
	设 \(L>0\)，分别取 Euclid 平面和开条带 \(S=\mathbb R\times(0,1)\)，都按 \(T(x,y)=(x+L,y)\) 取商。写出两者的展开映射和单值表示，说明它们的单值表示相同，但一个完备、另一个不完备。这是否说明单值表示不能单独决定一个 \((G,X)\) 结构？
\end{exercise}
\begin{solution}
	平面与开条带都单连通，整数平移的作用自由且固有，所以它们分别是商曲面的万有覆盖。以标准坐标作为初始图，展开映射分别为
	\[
		\operatorname{dev}_1:\mathbb R^2\to\mathbb R^2,\quad
		\operatorname{dev}_2:S\hookrightarrow\mathbb R^2,
	\]
	两者都是包含公式 \((x,y)\mapsto(x,y)\)，而
	\[
		\rho_1(n)(x,y)=\rho_2(n)(x,y)=(x+nL,y).
	\]
	完整平面上的商由\cref{g1:prop:complete-pullback-cover}完备。条带上的商中，点列 \([(0,1/j)]\)、\(j\ge2\)，为 Cauchy 列，因为竖直线段给出相邻点之间的距离至多为相应高度差。若它在商中有极限，高度函数 \(y\) 的连续性会使极限的高度为零，但商内全部高度属于 \((0,1)\)，矛盾。

	两个展开像分别为整个平面和真开条带，任一 Euclid 刚性运动都不可能把前者变为后者。因此它们不可能通过整体改变初始图得到同一展开数据。单值表示相同只记录了同样的绕行变换，未记录模型中究竟覆盖了哪些位置；它确实不能单独决定结构。
\end{solution}

\begin{exercise}
	\label{g1:ex:spiral-similarity-torus}
	设 \(a>0\)、\(b\in\mathbb R\)，令 \(\Lambda\subset\mathbb R^2\) 由 \((a,b)\)、\((0,2\pi)\) 生成。证明
	\[
		F(u,v)=e^u(\cos v,\sin v)
	\]
	在环面 \(M=\mathbb R^2/\Lambda\) 上诱导平面保向相似几何结构，并求展开映射及两个基环路的单值。为什么这个紧曲面不违反\cref{g1:thm:complete-model-quotient}？
\end{exercise}
\begin{solution}
	两个格向量线性无关；线性坐标变换把该格化成 \(\mathbb Z^2\)，故其作用自由固有，商为紧环面，\(\mathbb R^2\) 是万有覆盖。\(F\) 的两个偏导为
	\[
		F_u=e^u(\cos v,\sin v),\qquad F_v=e^u(-\sin v,\cos v),
	\]
	行列式为 \(e^{2u}>0\)，所以 \(F\) 局部微分同胚。若在足够小的商邻域中选取一个覆盖片，再限制 \(F\) 为单射，就得到一张平面图。两个覆盖片相差格平移 \((ma,mb+2\pi n)\)，而
	\[
		F(u+ma,v+mb+2\pi n)=e^{ma}R_{mb}F(u,v).
	\]
	故图的过渡映射属于平面保向相似群，给出所需结构。选初始图与 \(F\) 一致，展开映射就是 \(F\)，单值为
	\[
		\rho(m,n)(x)=e^{ma}R_{mb}x.
	\]
	第一基环路的单值是伸缩 \(e^a\) 复合旋转 \(b\)，第二基环路的单值为恒等，故表示的核是 \(\{0\}\times\mathbb Z\)。

	完备模型商定理要求过渡群保持模型度量。这里 \(e^a\ne1\)，第一单值是非等距的相似变换，因而 Euclid 度量不能经这些模型图下降。事实上 \(F^*g_{\mathrm{Euc}}=e^{2u}(\mathrm{d}u^2+\mathrm{d}v^2)\)，按 \((u,v)\mapsto(u+a,v+b)\) 平移会乘上 \(e^{2a}\)。环面当然可以另取平坦度量 \(\mathrm{d}u^2+\mathrm{d}v^2\)，但它不是这份相似图册由 Euclid 模型诱导的度量。紧性与另一份度量的完备性不能补上定理缺失的等距假设。
\end{solution}

\begin{exercise}
	\label{g1:ex:klein-geometric-holonomy}
	在 Euclid 平面上固定 \(L,H>0\)，令
	\[
		A(x,y)=(x+L,y),\qquad B(x,y)=(-x,y+H).
	\]
	证明它们生成的群自由固有地作用于平面，其商是紧的不可定向曲面。写出展开与单值，并说明为什么反射成分不会破坏完备模型商定理。
\end{exercise}
\begin{solution}
	直接计算 \(BAB^{-1}=A^{-1}\)，所以任意群元素可写成 \(A^mB^n\)。这一表示唯一：它在 \(y\) 坐标上移动 \(nH\)，确定 \(n\) 后，\(x\) 方向的平移又确定 \(m\)。其作用为
	\[
		A^mB^n(x,y)=((-1)^n x+mL,y+nH).
	\]
	若有不动点，第二坐标迫使 \(n=0\)，第一坐标继而迫使 \(m=0\)，所以作用自由。若紧集 \(K\) 包含在 \([-C,C]^2\)，而 \(A^mB^nK\cap K\ne\varnothing\)，则 \(|n|H\le2C\)、\(|m|L\le2C\)，只有有限种整数选择，故作用固有。每个轨道都有代表落在紧矩形 \([0,L]\times[0,H]\) 中：先用 \(B\) 的幂调整高度，再用 \(A\) 的幂调整横坐标。因此商紧。这是 Klein 瓶的边配对模型。

	若商可定向，其定向拉回到连通平面，必须是平面的两种定向之一，并被所有覆盖变换保持。但 \(\det(\mathrm{d}B)=-1\)，与此矛盾，所以商不可定向。万有覆盖为平面，展开映射可取恒等，\(\rho\) 把两个基环路分别送到 \(A,B\)，是忠实表示，关系为 \(\rho(b)\rho(a)\rho(b)^{-1}=\rho(a)^{-1}\)。两种变换均保持 Euclid 度量；完备模型商定理允许整个等距群，未要求保向。故商仍完备，不可定向性只表示单值群没有落在保向子群内。
\end{solution}

\begin{exercise}
	\label{g1:ex:surjective-development-not-covering}
	令 \(U=\mathbb R\times(0,4\pi)\)，并取 Euclid 模型中的局部映射
	\[
		F(u,v)=e^u(\cos v,\sin v).
	\]
	把 \(F\) 的单射限制作为 \(U\) 的 Euclid 几何图册。求其展开映射、单值及诱导度量；证明展开映射满到其像 \(\mathbb R^2\setminus\{0\}\)，却不是到该像的覆盖。说明局部等距为什么没有自动保证道路提升到最后。
\end{exercise}
\begin{solution}
	与前题相同的偏导计算说明 \(F\) 局部可逆；其各单射限制在重叠处完全相等，故过渡映射可取单位元。\(U\) 是凸开集，因而单连通，初始图选为 \(F\) 的限制后，展开映射就是 \(F\)，基本群和单值表示均平凡。诱导度量为
	\[
		g=e^{2u}(\mathrm{d}u^2+\mathrm{d}v^2).
	\]
	任意非零点可写成 \(r(\cos\theta,\sin\theta)\)，其中可选 \(0<\theta<4\pi\)，取 \(u=\log r\) 即得原像，所以 \(F\) 满到穿孔平面。

	正实轴上的点只有一个原像角度 \(v=2\pi\)；角度 \(0<\theta<2\pi\) 的非正实轴点却有 \(v=\theta,\theta+2\pi\) 两个原像。覆盖在连通的小邻域上具有固定的纤维基数，见\cref{g1:prop:covering-local-properties}后面的说明，因此正实轴点不可能有均匀覆盖邻域，\(F\) 不是覆盖。

	例如从 \((u,v)=(0,\pi)\) 出发，提升单位圆上按角度从 \(\pi\) 递减到 \(0\) 的道路，唯一局部提升是 \((0,v(t))\)，它在终点趋向 \((0,0)\)，已离开 \(U\)。相应长度仅为 \(\pi\)，提升却不能到达终点。缺口并非来自目标点不存在，而来自源空间在有限长度处不完备；度量完备正是\cref{g1:lem:complete-local-isometry}中用来填补这种提升缺口的条件。
\end{solution}


% !TeX root = ../Book1.tex
% 纲要标题：# 第四编　光滑流形的微积分与联络
% 纲要位置：几何学计划纲要.md，第 531 行。
% 正文状态：编导读初稿；各章状态分别维护。

% 纲要内容开始
% 
% 纲要内容结束

\BookPart{光滑流形的微积分与联络}{g1:part:04}
\BookPartGuide

怎样在没有统一坐标的空间上求导、积分，并解出随时间变化的轨道？本编从光滑相容的局部坐标出发，把 Euclid 微积分中的构造转到流形上。这里可以作为流形学习的起点。需要的基础是有限维线性代数中的基、对偶、秩与行列式，多元微积分中的全微分、链式法则与积分，以及开集、连续性、同胚、子空间和乘积等基本拓扑概念。前三编提供的曲面经验可以帮助理解这些构造，开始本编则无需先完成古典分类。

\ProseParagraphBreak

\chapref{g1:ch:11}首先回答哪些空间和映射可以接受光滑计算：图册确定结构，切空间记录一阶方向，微分的秩决定局部标准形。逆函数、隐函数和常秩定理在 Euclid 空间中重建，再通过坐标应用于流形。\chapref{g1:ch:12}把切空间推广为随点变化的向量空间，用过渡函数组织丛和张量，并用单位分解粘合局部数据。Möbius 线丛与丛度量的对照值得反复体会：局部基未必能拼成整体基，局部正定内积却能经加权相加组成整体内积。

\ProseParagraphBreak

逐点指定的方向还要通过积分才能形成运动。\chapref{g1:ch:13}从常微分方程建立局部流，以 Lie 括号记录两种运动交换次序后的差异，并由 Frobenius 定理判断固定秩分布何时具有积分叶。局部流的存在时间、轨道的完备性和叶自身的拓扑在这里逐一出现；环面上的无理斜率运动说明，局部坐标中的直线可以在整体上形成稠密的叶。

\ProseParagraphBreak

\chapref{g1:ch:14}把换元和边界积分写成微分形式的语言。外微分与同伦公式给出局部原形式，定向和单位分解使积分成为整体运算，Stokes 定理则把内部的微分与边界上的积分联系起来。穿孔平面的角形式保留着一个可直接计算的整体障碍：沿闭路的非零周期阻止局部原函数拼成单值函数。这也解释了为什么后续上同调理论要记录闭形式对恰当形式的差别。

\ProseParagraphBreak

\chapref{g1:ch:15}用联络规定沿道路比较不同纤维中向量的方法。平行移动由常微分方程求出，曲率记录小环路运输的首项，平坦局部模型再把零曲率转成平行标架。较大环路留下的单值信息需要基本群与覆盖；直接从本编开始的读者，可在阅读\secref{g1:sec:1505}时回看\chapref{g1:ch:05}的道路同伦和万有覆盖。至此，求导、运动、积分和运输都有了可在局部计算、又能比较整体结果的表达方式，随后便可把它们用于 Lie 群及主丛。

% !TeX root = ../../Book1.tex
% 纲要标题：## 第11章　光滑流形、切丛与局部标准形
% 纲要位置：几何学计划纲要.md，第 533 行。
% 正文状态：初稿；尚未正式审核。

% 纲要内容开始
% 
% > 本章目标与证明范围：流形路线从本章独立进入。按“坐标相容—切映射—秩条件—局部标准形—子流形”的顺序证明，不先使用尚未定义的切丛；必要多元微分定理就地重建。
% 
% 纲要内容结束

\chapter{光滑流形、切丛与局部标准形}
\label{g1:ch:11}
\BookChapterNumber{g1:ch:11}

\begin{chapterabstract}
	流形上的微分计算由光滑换坐标相互联系，像上的拓扑则决定局部参数怎样组成整体。本章从光滑相容的图册构造球面、射影空间与带边界流形，以局部函数芽的导子定义切空间，并证明内点处与曲线速度的等价，建立切丛、余切空间和微分的链式法则。随后用压缩估计完整证明定量逆函数定理，推导参数隐函数定理与常秩定理，得到浸入、淹没和正则水平集的局部标准形。最后区分单射浸入与嵌入，证明限制映射的光滑性与固有嵌入判据，并明确边界版本所需的附加条件。
\end{chapterabstract}

\begin{learninggoals}
	\begin{itemize}
		\item 用明确的坐标图、换坐标和拓扑条件构造光滑流形，区分同胚、光滑双射与微分同胚。
		\item 从导子恢复切空间的坐标基，计算微分与余切拉回，并说明边界处完整切空间与可实现曲线速度的区别。
		\item 由导数估计构造局部逆，给出逆映射的正则性证明，并在有参数的方程中明确存在唯一性的局部范围。
		\item 用常秩坐标化简光滑映射，辨认浸入、淹没与正则水平集，计算水平集的切空间。
		\item 检查像空间的拓扑，判断单射浸入何时是嵌入；在带边界问题中分别核对微分秩与边界相交条件。
	\end{itemize}
\end{learninggoals}

一个圆周可以由方程 \(x^2+y^2=1\) 描述，也可以在小段上用角度参数化。第一种写法把它放在平面里，第二种写法只需一个实变量。若只关心圆周上的光滑函数和运动速度，究竟哪些量依赖平面，哪些量依赖角度的选择，哪些量属于圆周本身？同样的问题出现在球面、射影空间和更高维约束集合上。光滑流形把“每一小块可以使用 Euclid 微积分”组织成一个与坐标选择相容的整体结构。

\ProseParagraphBreak

光滑映射、可逆微分与像上的拓扑描述着不同的性质。映射 \(t\mapsto t^3\) 光滑且双射，逆却在零点不可微；一条参数曲线可以处处有非零速度、没有重复经过同一点，而其像的邻域仍从很远的参数处迎来另一支。要分辨这些现象，需要同时记录一阶微分、局部坐标与像上的拓扑。逆函数定理和常秩定理负责局部，嵌入的定义则保留不可由微分秩代替的整体要求。

\ProseParagraphBreak

本章可以直接开始学习，无须先通读前三编。所需基础是有限维实线性代数中的基、对偶、秩和行列式，以及多元微积分中的偏导数、全微分、链式法则和一元微积分基本定理。拓扑方面只需理解开集、连续映射、同胚、子空间与乘积；Hausdorff、第二可数和商拓扑在实际构造处作简短回顾。逆函数、隐函数和常秩三个定理在本章从 Euclid 微积分重建。没有给流形选定度量时，切向量没有预先指定的长度，余切向量也不能自动认作切向量。

\ProseParagraphBreak

可对照 Lee 第二版的第1--2章学习光滑结构与映射，第3--5章学习切空间、常秩、嵌入和子流形，第11章的 Covectors、The Differential of a Function、Pullbacks of Covector Fields 三部分学习余切与拉回\cite[第1--5章及第11章，第272--286页]{Lee2012Smooth}。这里以局部函数芽定义切空间；一般向量丛和整体构造将在\chapref{g1:ch:12}展开。带边界情形随定义和结果就近讨论，可以在第一次阅读时先理解内点证明，再对照边界处增加的条件。

% !TeX root = ../../Book1.tex
% 纲要标题：### 11.1 光滑结构与典型构造
% 纲要位置：几何学计划纲要.md，第 537 行。
% 正文状态：初稿；尚未正式审核。

% 纲要内容开始
% 
% * 图册、极大图册及光滑映射
% * 球面、射影空间、乘积与开子流形
% * 带边界流形和边界上的光滑性
% * 光滑结构与底层拓扑的区别
% 
% 纲要内容结束

\section{光滑结构与典型构造}
\label{g1:sec:1101}

在圆周的一小段上，可以用横坐标、纵坐标或适当选择的角度表示点。但横坐标在圆周的最右端不能区分邻近的上下两点，角度也不能在整条圆周上连续地取单值。流形的定义允许使用多张局部坐标图，同时要求重叠处的换坐标足够光滑。这样，在一张图中求出的导数才有可能与另一张图中的计算表达同一个对象。

\subsection{图册、极大图册及光滑映射}
\label{g1:subsec:1101:01}

先回顾本节真正需要的拓扑条件。空间是 Hausdorff 的，指任意两个不同点有不相交的开邻域；它是第二可数的，指存在可数个开集，使每个开集都能写成其中某些集合的并。前一个条件使极限不会同时落在两个不同点上，后一个条件排除无法由可数个局部坐标控制的过大空间。开子集总取给定拓扑中的开集，不能在选择坐标时另换一套拓扑。

\begin{definition}
	\label{g1:def:smooth-manifold}
	设 \(n\geq0\)，\(M\) 为 Hausdorff、第二可数拓扑空间。若 \(U\subset M\) 是开集，\(\varphi:U\to\varphi(U)\subset\mathbb R^n\) 是到欧氏开集的同胚，称 \((U,\varphi)\) 为一张\term[zuobiaotu]{坐标图}[coordinate chart]。图域覆盖 \(M\) 的坐标图族称为\term[tuce]{图册}[atlas]。两张图 \((U,\varphi)\)、\((V,\psi)\) 称为\term[guanghuaxiangrong]{光滑相容的}[smoothly compatible]，若
	\[
		\psi\circ\varphi^{-1}:\varphi(U\cap V)\longrightarrow\psi(U\cap V)
	\]
	及其逆映射都为 \(C^\infty\) 映射；交集为空时条件自动成立。两两光滑相容的图册称为\term[guanghuatuce]{光滑图册}[smooth atlas]。不能再加入新的相容坐标图的光滑图册称为\term[jidaguanghuatuce]{极大光滑图册}[maximal smooth atlas]，也称 \(M\) 的一个\term[guanghuajiegou]{光滑结构}[smooth structure]。配备这一结构的 \(M\) 称为 \(n\) 维\term[guanghualiuxing]{光滑流形}[smooth manifold]。
\end{definition}

本节先讨论无边界流形；带边界的定义在\subsecref{g1:subsec:1101:03}给出。约定 \(\mathbb R^0\) 是单点空间，因此零维流形是可数离散空间。坐标图写成 \(\varphi=(x^1,\ldots,x^n)\) 时，\(x^i\) 是定义在 \(U\) 上的实值函数，上标是坐标编号，并非幂。换坐标同时改变所有坐标函数，点 \(p\in M\) 本身不变。

\begin{proposition}
	\label{g1:prop:atlas-maximal-extension}
	设 \(\mathcal A\) 是 Hausdorff、第二可数空间 \(M\) 上的一个光滑图册。所有与 \(\mathcal A\) 中每张图相容的坐标图组成唯一包含 \(\mathcal A\) 的极大光滑图册。两个光滑图册决定同一光滑结构，当且仅当它们的并仍是光滑图册。
\end{proposition}

\begin{proof}
	记所述图族为 \(\overline{\mathcal A}\)。要证明它是图册，关键是其中新加入的两张图也相容。设 \((V,\psi)\)、\((W,\eta)\in\overline{\mathcal A}\)，在 \(p\in V\cap W\) 处选一张原图 \((U,\varphi)\in\mathcal A\) 覆盖 \(p\)。在 \(U\cap V\cap W\) 上，
	\[
		\eta\circ\psi^{-1}=(\eta\circ\varphi^{-1})\circ(\varphi\circ\psi^{-1}).
	\]
	右端两个映射都光滑，逆变换也一样。光滑性是局部性质，故两张新图在整个交叠上相容。原图已经覆盖 \(M\)，所以 \(\overline{\mathcal A}\) 是光滑图册。

	与 \(\overline{\mathcal A}\) 相容的图必与 \(\mathcal A\) 相容，因而已经包含在其中，这证明极大性。任一包含 \(\mathcal A\) 的极大图册只能由与 \(\mathcal A\) 相容的图组成，故包含于 \(\overline{\mathcal A}\)，由极大性二者相等。最后，两个图册的并若光滑，其极大扩张同时包含二者，唯一性说明两者决定同一结构；反向结论来自同一极大图册中的图两两相容。
\end{proof}

因此，给出几张足够覆盖空间的坐标图就能指定光滑结构，不必逐张列出极大图册。限制坐标图到更小开集、把欧氏坐标再作一个光滑可逆变换，都不会改变这个结构。

\ProseParagraphBreak

下面用到的欧氏微分只需要一个基本事实。映射 \(f\) 在 \(a\) 可微，意味着有线性映射 \(\mathrm{D}f(a)\) 使 \(f(a+h)=f(a)+\mathrm{D}f(a)h+o(|h|)\)。把这个式子代入 \(g\) 在 \(f(a)\) 的一阶展开，由 \(f(a+h)-f(a)=O(|h|)\) 得到
\[
	\mathrm{D}(g\circ f)(a)=\mathrm{D}g(f(a))\,\mathrm{D}f(a).
\]
这就是链式法则。各阶偏导数连续的含义也在欧氏开集上理解。逆函数定理比链式法则强得多，将在\secref{g1:sec:1103}另行证明。

\begin{definition}
	\label{g1:def:smooth-map}
	设 \(M,N\) 分别为 \(m,n\) 维光滑流形。连续映射 \(F:M\to N\) 称为\term[guanghuayingshe]{光滑映射}[smooth map]，若对每个 \(p\in M\)，可选 \(p\) 附近的图 \((U,\varphi)\) 和 \(F(p)\) 附近的图 \((V,\psi)\)，其中 \(F(U)\subset V\)，使
	\[
		\psi\circ F\circ\varphi^{-1}:\varphi(U)\longrightarrow\psi(V)
	\]
	为欧氏空间之间的 \(C^\infty\) 映射。若 \(F\) 双射，且 \(F,F^{-1}\) 都光滑，则称 \(F\) 为\term[weifentongpei]{微分同胚}[diffeomorphism]。
\end{definition}

要求 \(F(U)\subset V\) 只是为了让坐标表达有意义；连续性允许围绕任意指定的 \(p\) 缩小 \(U\) 来满足它。图中的三条箭头说明计算发生在哪里：
\[
	\begin{tikzcd}[column sep=large,row sep=large]
		U\subset M \arrow[r,"F"] \arrow[d,"\varphi"'] & V\subset N \arrow[d,"\psi"] \\
		\varphi(U)\subset\mathbb R^m \arrow[r,"\psi\circ F\circ\varphi^{-1}"'] & \psi(V)\subset\mathbb R^n.
	\end{tikzcd}
\]

\begin{proposition}
	\label{g1:prop:smooth-coordinate-independence}
	光滑映射的定义与所选相容坐标图无关。光滑映射的复合光滑，恒等映射光滑。如果非空的 \(m\) 维与 \(n\) 维光滑流形微分同胚，则 \(m=n\)。
\end{proposition}

\begin{proof}
	换成源端坐标 \(\widetilde\varphi\) 与目标坐标 \(\widetilde\psi\) 后，在共同定义的邻域中，新表达为
	\[
		\widetilde\psi\circ F\circ\widetilde\varphi^{-1}
		=(\widetilde\psi\circ\psi^{-1})\circ(\psi\circ F\circ\varphi^{-1})\circ(\varphi\circ\widetilde\varphi^{-1}).
	\]
	两侧换坐标与中间表达都光滑，故复合光滑。复合映射的结论由同样的坐标计算得到，恒等映射的坐标表达就是换坐标。

	若 \(F\) 是微分同胚，在对应点附近用坐标写 \(f\) 与 \(g=f^{-1}\)，令 \(A=\mathrm{D}f(a)\)、\(B=\mathrm{D}g(f(a))\)。由链式法则有 \(BA=I_m\)、\(AB=I_n\)，故 \(A\) 同时单射、满射，有限维线性代数给出 \(m=n\)。这里证明的是光滑可逆映射保持维数，并未把它当作一般同胚保持维数的拓扑证明。
\end{proof}

\subsection{球面、射影空间、乘积与开子流形}
\label{g1:subsec:1101:02}

给一个集合列出参数式，还没有给出流形。需要先有 Hausdorff、第二可数拓扑，再检查参数的逆是否连续、图域是否开、换坐标是否光滑。下面的构造把这些步骤逐一落实，也说明抽象定义包含了熟悉的几何对象。

\begin{example}[球面的坐标]
	\label{g1:exm:sphere-smooth-atlas}
	令 \(S^n=\{x\in\mathbb R^{n+1}:|x|=1\}\)，取欧氏子空间拓扑。对 \(i=1,\ldots,n+1\) 和 \(\varepsilon\in\{1,-1\}\)，令
	\[
		U_i^\varepsilon=\{x\in S^n:\varepsilon x_i>0\},\qquad
		\varphi_i^\varepsilon(x)=(x_1,\ldots,\widehat{x_i},\ldots,x_{n+1}),
	\]
	其中帽号表示删掉该项。像为开单位球 \(B^n\)，逆映射在第 \(i\) 位插入 \(\varepsilon\sqrt{1-|u|^2}\)。所以每张图是同胚，图域覆盖球面。换坐标只是在一个点的坐标中插入这一光滑平方根，再删去另一项，故在交叠区域光滑；逆变换同理。球面是欧氏空间的子空间，Hausdorff 性直接继承，可数欧氏开球与球面的交给出可数基。因此这些图使 \(S^n\) 成为 \(n\) 维光滑流形。

	坐标在 \(|u|=1\) 时失效，是这一张图到了边缘，球面本身没有在那里出现边界。换一张不删去零坐标的图即可继续。若 \(n\geq1\)，球面不能由一张欧氏开集坐标图覆盖：球面紧，而 \(\mathbb R^n\) 中非空开集不可能紧。后者若紧便闭，因 \(\mathbb R^n\) 连通只能等于整个 \(\mathbb R^n\)，又与其非紧性矛盾。
\end{example}

\begin{example}[实射影空间]
	\label{g1:exm:projective-smooth-atlas}
	实射影空间 \(\mathbb{RP}^n\) 的点是 \(\mathbb R^{n+1}\) 中过原点的一维线性子空间，写为 \([x_0:\cdots:x_n]\)。非零向量乘以同一非零实数表示同一点。每条直线与球面相交于一对对径点，因而给它球面商 \(S^n/(x\sim-x)\) 的拓扑。商映射 \(q:S^n\to\mathbb{RP}^n\) 是开映射，因为对球面开集 \(O\)，有 \(q^{-1}(q(O))=O\cup(-O)\)。

	令 \(U_i=\{[x]:x_i\neq0\}\)，定义
	\[
		\varphi_i([x])=(x_0/x_i,\ldots,\widehat{x_i/x_i},\ldots,x_n/x_i)\in\mathbb R^n.
	\]
	这个式子与代表向量无关。\(U_i\) 开，坐标函数连续；逆映射把 \(1\) 插入第 \(i\) 位，再取该向量张成的直线，也连续。若要通过球面商直接检查逆的连续性，只需把这个插入后的向量除以其长度，再应用 \(q\)。换到 \(U_j\) 的坐标，就是把各分量除以第 \(j\) 个齐次分量；交叠条件恰好保证这个分母非零，故换坐标及其逆都光滑。

	还需核对商拓扑的分离性质，不能从球面的 Hausdorff 性直接跳到商空间。定义
	\[
		P([x])=\frac{xx^{\mathsf T}}{x^{\mathsf T}x}\in\operatorname{Sym}_{n+1}(\mathbb R),
	\]
	右端是到直线 \(\mathbb Rx\) 的正交投影矩阵。它与代表无关，并由球面上的连续映射下降，故连续。两个不同直线对应不同投影矩阵，在矩阵空间取不相交的小球，其原像就分离这两个射影点。因此 \(\mathbb{RP}^n\) 是 Hausdorff 的。有限张图中分别取有理中心、半径的坐标球，其逆像合起来是可数基，第二可数性也成立。于是得到 \(n\) 维光滑结构。球面商同时说明 \(\mathbb{RP}^n\) 紧。
\end{example}

射影坐标的分母为零只表示需要换图。与球面相同，坐标表达中的无穷远不必是空间的缺失点；在另一张图中，它可以有完全普通的有限坐标。

\begin{proposition}
	\label{g1:prop:basic-smooth-manifolds}
	设 \(M,N\) 分别为 \(m,n\) 维无边界光滑流形。任意开集 \(W\subset M\) 由坐标图的限制得到 \(m\) 维光滑结构，包含映射 \(W\hookrightarrow M\) 光滑。乘积拓扑下的 \(M\times N\) 具有由乘积坐标定义的 \(m+n\) 维光滑结构，两坐标投影光滑。映射 \(F:L\to M\times N\) 光滑，当且仅当两个分量映射都光滑。
\end{proposition}

\begin{proof}
	Hausdorff 性与第二可数性传给开子空间。将 \(M\) 的图域与 \(W\) 相交，像仍是欧氏开集，换坐标是原换坐标的限制，包含映射的表达是恒等映射。这给出开集的结论，称这样的 \(W\) 为\term[kaiziliuxing]{开子流形}[open submanifold]。

	乘积中的两点若不同，至少一个分量不同，在该分量中选不相交开邻域就能分离两点；两个可数基的所有乘积组成可数基。对 \(M,N\) 的图 \((U,\varphi)\)、\((V,\psi)\)，取 \((U\times V,\varphi\times\psi)\)。它是到 \(\mathbb R^{m+n}\) 开集的同胚，换坐标为两个光滑换坐标的乘积。投影在这些图中是普通线性投影；一个欧氏向量值函数光滑当且仅当其所有分量光滑，从而也得到关于 \(F\) 的断言。
\end{proof}

例如，\(S^1\times S^1\) 是二维光滑环面，不必先把它画成 \(\mathbb R^3\) 中的曲面。映射 \(t\mapsto((\cos t,\sin t),(\cos\alpha t,\sin\alpha t))\) 对每个实数 \(\alpha\) 都光滑，因为两分量的局部坐标表达光滑。这是先构造空间、再在其上计算的一个直接例子。

\subsection{带边界流形和边界上的光滑性}
\label{g1:subsec:1101:03}

闭圆盘上的微积分需要包括圆周，但在圆周附近可用的点只占平面邻域的一侧。局部模型因此从开球改为半空间中的相对开集。边界点也需要完整的一阶导数，不能只取沿边界走动时能够看见的方向。

\begin{definition}
	\label{g1:def:manifold-with-boundary}
	对 \(n\geq1\)，记
	\[
		\mathbb H^n=\{(x^1,\ldots,x^n)\in\mathbb R^n:x^n\geq0\}.
	\]
	若 \(A\subset\mathbb H^n\) 相对开，称 \(f:A\to\mathbb R^k\) 光滑，是指每个 \(a\in A\) 附近都有欧氏开邻域 \(O\) 和光滑函数 \(\widetilde f:O\to\mathbb R^k\)，使 \(\widetilde f=f\) 于 \(O\cap A\)。到另一半空间的映射也用这个定义，不要求延拓在半空间外仍取值于目标半空间。

	一个 Hausdorff、第二可数空间，若有到 \(\mathbb H^n\) 相对开集的坐标图覆盖，且各换坐标及其逆按上述延拓意义光滑，则由该图册的极大扩张得到一个 \(n\) 维\term[daibianjieguanghualiuxing]{带边界光滑流形}[smooth manifold with boundary]。在某张图中最后坐标为零的点称为边界点，组成 \(\partial M\)；其余点组成内部 \(\operatorname{Int}M\)。零维流形另约定边界为空。流形之间的光滑映射仍按局部坐标表达定义。
\end{definition}

这个定义需要两项确认：不同延拓是否给出同一导数，不同坐标图是否把同一点都判为边界点。两者均由局部计算解决。

\begin{proposition}
	\label{g1:prop:boundary-smooth-derivative}
	设 \(A\subset\mathbb H^n\) 相对开，\(f:A\to\mathbb R^k\) 按局部延拓意义光滑。任意两个局部光滑延拓在 \(a\in A\) 给出的所有阶偏导数一致，因而 \(\mathrm{D}f(a):\mathbb R^n\to\mathbb R^k\) 有确定含义。半空间相对开集之间的光滑映射满足链式法则。
\end{proposition}

\begin{proof}
	两延拓在共同邻域的半空间一侧相等，故在其中的欧氏开集 \(x^n>0\) 上各阶偏导数都相等。每个边界点可由这类内点逼近；偏导数连续，于是在边界也相等。对复合映射，先取内、外两个映射的局部延拓，再缩小源端的欧氏邻域，使内映射的延拓落入外延拓的定义域。复合延拓满足普通链式法则，由导数的唯一性得到所需公式。
\end{proof}

\begin{proposition}
	\label{g1:prop:boundary-chart-invariance}
	带边界光滑流形的边界点判定与坐标图无关。\(\operatorname{Int}M\) 是开集并带有 \(n\) 维无边界光滑结构；\(\partial M\) 是闭集，若非空则具有由边界坐标诱导的 \(n-1\) 维无边界光滑结构。带边界流形之间的微分同胚把边界映到边界。
\end{proposition}

\begin{proof}
	设 \(F:A\to B\) 是两个半空间相对开集之间的光滑坐标变换，逆映射记为 \(G\)。假设 \(a\in A\) 满足 \(a^n=0\)，但 \(b=F(a)\) 是 \(B\) 的内点。于是 \(b\) 有一个完全包含于 \(B\) 的欧氏开球。这个球上的 \(G\) 取值于 \(\mathbb H^n\)，故最后分量 \(G^n\geq0\)，且 \(G^n(b)=0\)。可微函数在内点取得极小值时各方向导数为零，所以 \(\mathrm{D}G(b)\) 的最后一行全为零。另一方面，\(F\circ G\) 是恒等映射，由\cref{g1:prop:boundary-smooth-derivative}得 \(\mathrm{D}F(a)\mathrm{D}G(b)=I_n\)，矛盾。对逆变换同样推理，边界与内部均保持。

	内部在任一图中由 \(x^n>0\) 描述，因而开；这些图限制后就是欧氏开集坐标。边界是其补集，故闭。在边界上删去恒为零的最后坐标，图像是 \(\mathbb R^{n-1}\) 的开集，换坐标及其逆由原光滑延拓限制到 \(x^n=0\) 得到，仍光滑。Hausdorff 性与第二可数性传给子空间，因此确实得到无边界光滑流形。一般微分同胚在局部也成为上述 \(F\)，所以同样保持边界。
\end{proof}

\begin{example}[闭球与边界条件]
	\label{g1:exm:closed-ball-boundary}
	闭球 \(\overline B^n\) 的内点已有欧氏坐标。在边界点附近先作正交坐标变换，把该点移到北极 \((0,\ldots,0,1)\)。写 \(u\in\mathbb R^{n-1}\)，取足够小的 \(u\) 与 \(t\geq0\)，令
	\[
		\Phi(u,t)=\bigl(u,\sqrt{1-|u|^2}-t\bigr).
	\]
	球内靠近北极的点恰由这类参数给出；逆映射为 \(u=(x_1,\ldots,x_{n-1})\)、\(t=\sqrt{1-|u|^2}-x_n\)。在北极的一个小邻域内，球的不等式等价于 \(t\geq0\)，另一侧球面远离此邻域。两式都有欧氏光滑延拓，故这些边界图与内部图的转换光滑。由此 \(\overline B^n\) 是带边界流形，且 \(\partial\overline B^n=S^{n-1}\)，诱导结构正是前面的球面结构。

	光滑映射本身可以不保持边界。映射 \(t\mapsto t+1\) 从 \([0,\infty)\) 到自身光滑，却把边界点送到内点；映射 \(s\mapsto s^2\) 从 \(\mathbb R\) 到 \([0,\infty)\) 光滑，却把内点送到边界。保持边界是微分同胚的结论，不能从“光滑”两字推出。
\end{example}

在 \(\partial M\) 的坐标中只保留前 \(n-1\) 个分量，但边界处 \(M\) 上函数的微分仍有 \(n\) 个分量。后一事实来自函数在半空间一侧的变化，而不是任选延拓带来的额外自由。\secref{g1:sec:1102}将据此区分 \(\mathrm{T}_pM\)、\(\mathrm{T}_p\partial M\) 和单侧曲线的速度。

\ProseParagraphBreak

一个带边界流形与一个无边界流形的乘积仍有半空间坐标，只需把有边界的那个坐标移到最后。两个带边界流形的自然乘积图却可能以 \(\mathbb H^m\times\mathbb H^n\) 为模型，例如象限的原点同时遇到两条边。这样的图带有角，不能未经说明就当作本节的半空间光滑图；若要讨论带角结构或重新选择平滑方式，需要另给定义。

\subsection{光滑结构与底层拓扑的区别}
\label{g1:subsec:1101:04}

坐标图首先必须是同胚，但同胚只控制连续性，不控制函数在坐标中的导数。为了看清两层结构的区别，不必立即寻找复杂流形；同一条实直线就足够。

\begin{example}[运输来的直线坐标]
	\label{g1:exm:transported-smooth-line}
	在具有通常拓扑的集合 \(\mathbb R\) 上，分别用全局坐标 \(x(t)=t\) 与 \(u(t)=t^3\) 生成两个极大光滑图册，所得流形记为 \(M_x\)、\(M_u\)。\(u\) 是同胚，所以单张图确实定义一个光滑结构。但两张图不相容：\(u\circ x^{-1}(s)=s^3\) 光滑，其逆 \(x\circ u^{-1}(s)=\sqrt[3]{s}\) 在零点不可微。因此这是同一个拓扑空间上的两个不同光滑结构。

	集合的恒等映射 \(I:M_x\to M_u\) 的坐标表达为 \(s^3\)，它光滑；反向恒等映射的表达为 \(\sqrt[3]{s}\)，不光滑。然而映射 \(F:M_u\to M_x\)、\(F(t)=t^3\) 是微分同胚，因为
	\[
		x\circ F\circ u^{-1}(s)=s,
	\]
	其逆的坐标表达也为恒等映射。故“同一集合上的结构不同”与“两个光滑流形不微分同胚”是不同判断，这个例子只证明前者。
\end{example}

同一个实值函数是否光滑也取决于源端结构。例如 \(t\mapsto t\) 在 \(M_x\) 上光滑，在 \(M_u\) 上的表达却是 \(u\mapsto\sqrt[3]{u}\)。相反，\(t\mapsto t^3\) 是 \(M_u\) 的坐标函数，当然光滑。写“这个函数光滑”时，通常已固定的图册正是这句话的必要背景。

\ProseParagraphBreak

一般的光滑双射也不必是微分同胚：通常光滑直线上的 \(t\mapsto t^3\) 已给出反例。其障碍在零点的一阶导数为零，而不在连续性或单射性。后面的逆函数定理将说明，导数可逆足以局部修复这个障碍；它仍然不会把局部可逆自动提升为整体单射。

\ProseParagraphBreak

本节通过具体图册构造了光滑空间，并检验了换坐标。对于一个仅仅给定拓扑的局部欧氏空间，定义本身并不提供制造相容光滑图册的方法，也不保证不同选择必然等价。当前讨论所需的是明确说明选了哪种结构，并在这一结构中计算；后文的切空间和微分都建立在这一选择之上。
% 纲要条目排为编号小节标题，并进入目录与书签；说明段仍保留于纲要注释中。

% !TeX root = ../../Book1.tex
% 纲要标题：### 11.2 切空间和微分
% 纲要位置：几何学计划纲要.md，第 544 行。
% 正文状态：初稿；尚未正式审核。

% 纲要内容开始
% 
% * 曲线速度与导子定义
% * 切空间的坐标基与曲线速度
% * 微分、链式法则与切丛图册
% * 余切空间和拉回
% 
% 纲要内容结束

\section{切空间和微分}
\label{g1:sec:1102}

在 Euclid 空间里，曲线经过一点时的速度是一个向量。换成流形以后，两个附近的点一般不能相减，因而“位置差除以时间差”没有直接意义。但沿曲线观察一个光滑函数，它仍是一个普通的实变量函数，可以求导。把所有这些观测的变化率同时记下来，就能定义一个不依赖外部空间的速度。这个想法既连接曲线与函数，也会说明微分为什么把切向量向前送、把线性函数向后拉。

\ProseParagraphBreak

本节取光滑 \(n\) 维流形 \(M\)，必要时允许边界；光滑性与边界图沿用\cref{g1:def:smooth-manifold,g1:def:manifold-with-boundary}。边界模型始终是 \(\mathbb H^n=\{x^n\ge0\}\)，光滑函数在半空间上的定义是局部延到 \(\mathbb R^n\) 后光滑。边界处的切空间仍将是 \(n\) 维向量空间；能在流形内实际走出的速度，则需要按时间是双侧还是单侧另行辨认。

\subsection{曲线速度与导子定义}
\label{g1:subsec:1102:01}

先看内点 \(p\in M\)。取两条光滑曲线 \(\gamma,\eta:(-\varepsilon,\varepsilon)\to M\)，满足 \(\gamma(0)=\eta(0)=p\)。在坐标图 \(x=(x^1,\ldots,x^n)\) 中，可以比较
\[
	(x\circ\gamma)'(0),\qquad (x\circ\eta)'(0)\in\mathbb R^n.
\]
若二者相等，称两条曲线在 \(p\) 有相同的速度。这个条件不依赖图的选择：换到 \(y\) 坐标，Euclid 链式法则给出
\[
	(y\circ\gamma)'(0)
	=\mathrm{D}(y\circ x^{-1})_{x(p)}(x\circ\gamma)'(0),
\]
对 \(\eta\) 也有同式；逆坐标变换给出反向结论。速度等价只比较一阶信息，不要求两条曲线在附近重合。为了给速度等价类定义加法，下面将把沿曲线求导写成作用于函数的线性算子。

\ProseParagraphBreak

固定 \(p\)。在 \(p\) 的某个开邻域上定义的两个光滑实函数，如果在更小的 \(p\) 邻域内相等，就把它们视为同一个\term[hanshuya]{函数芽}[germ of a function]，记作 \([f]_p\)。全体这些芽组成实代数 \(C^\infty_p(M)\)：加法与乘法在两个定义域的交上进行，换代表不改变结果。求值 \(f(p)\) 也只依赖芽。采用芽，是为了只保留当前点附近的信息；一个速度不应因远处怎样延长函数而改变。

\begin{definition}
	\label{g1:def:tangent-derivation}
	设 \(M\) 为光滑流形，可带边界，\(p\in M\)。一个实线性映射
	\[
		v:C^\infty_p(M)\longrightarrow\mathbb R
	\]
	若对所有函数芽满足
	\[
		v(fg)=f(p)v(g)+g(p)v(f),
	\]
	就称为 \(p\) 处的\term[daozi]{导子}[derivation]，或一个\term[qiexiangliang]{切向量}[tangent vector]。这里简写 \(v(f)=v([f]_p)\)。所有这样的导子按逐项相加与数乘组成实向量空间，称为 \(p\) 处的\term[qiekongjian]{切空间}[tangent space]，记为 \(\mathrm{T}_pM\)。\symidx{\(\mathrm{T}_pM\)：点处的切空间}
\end{definition}

导子的乘法法则来自沿曲线求导。若 \(\gamma(0)=p\)，则
\[
	v_\gamma(f)=\left.\frac{\mathrm{d}}{\mathrm{d}t}f(\gamma(t))\right|_{t=0}
\]
确为导子。代表函数在 \(p\) 附近相等时，复合在 \(t=0\) 附近相等，所以它作用于函数芽是良定义的；线性和乘法法则来自一元微积分。相同坐标速度给出相同导子，因为 \(f\circ\gamma=(f\circ x^{-1})\circ(x\circ\gamma)\)，可以使用链式法则。

\ProseParagraphBreak

导子自动消去常数。将乘法法则用于 \(1\cdot1\)，得到 \(v(1)=2v(1)\)，所以 \(v(1)=0\)，再由实线性得 \(v(c)=0\)。它也消去两个在 \(p\) 为零的函数芽的乘积。这提示导子只能看到一阶变化；下一个小节将把任意导子写成坐标方向导数的线性组合。

\subsection{切空间的坐标基与曲线速度}
\label{g1:subsec:1102:02}

设 \(x:U\to x(U)\) 是 \(p\) 附近的坐标图，记 \(a=x(p)\)。对于函数芽 \(f\)，定义
\[
	\left.\frac{\partial}{\partial x^i}\right|_p(f)
	=\frac{\partial(f\circ x^{-1})}{\partial x^i}(a).
\]
若 \(p\) 在边界，右边使用光滑延拓的偏导数。这一值与延拓无关，见\cref{g1:prop:boundary-smooth-derivative}：两个延拓在半空间内侧相等，其各阶偏导在内侧相等，再由连续性取极限到边界。偏导数的线性与乘法法则说明上述算子确为切向量，称为坐标导数。

\begin{proposition}[切空间的坐标基]
	\label{g1:prop:derivation-coordinate-basis}
	设 \(M\) 为光滑 \(n\) 维流形，可带边界，\(x\) 为 \(p\) 附近的坐标图。则
	\[
		\left.\frac{\partial}{\partial x^1}\right|_p,\ldots,
		\left.\frac{\partial}{\partial x^n}\right|_p
	\]
	是 \(\mathrm{T}_pM\) 的一组基。具体地，每个 \(v\in \mathrm{T}_pM\) 唯一写为
	\begin{equation}
		\label{g1:eq:tangent-coordinate-basis}
		v=\sum_{i=1}^n v(x^i)\left.\frac{\partial}{\partial x^i}\right|_p.
	\end{equation}
	因此内点和边界点的切空间均为 \(n\) 维。
\end{proposition}
\begin{proof}
	核心是将一个函数的一阶部分从余项中分离出来，而不要求导子具有任何预先假定的连续性。设 \(F=f\circ x^{-1}\)；边界情形先在 \(a\) 附近选光滑延拓。缩小到 Euclid 小球后，沿连接 \(a,z\) 的线段使用微积分基本定理，得到
	\[
		F(z)-F(a)=\sum_{i=1}^n(z^i-a^i)H_i(z),\qquad
		H_i(z)=\int_0^1\partial_iF(a+t(z-a))\,\mathrm{d}t.
	\]
	在较小邻域内，积分中的函数及任意阶导数连续，参数 \(t\) 处于紧区间，所以可以逐阶在积分号内求导；故 \(H_i\) 光滑，且 \(H_i(a)=\partial_iF(a)\)。边界时把此恒等式限制回半空间即可，因而它给出函数芽的等式
	\[
		f-f(p)=\sum_i(x^i-a^i)(H_i\circ x).
	\]
	导子作用于右端时，乘法法则中的 \(x^i(p)-a^i\) 为零，于是
	\[
		v(f)=\sum_i v(x^i)\,H_i(a)
		=\sum_i v(x^i)\frac{\partial(f\circ x^{-1})}{\partial x^i}(a).
	\]
	这证明了张成性。又有 \(\left.\frac{\partial}{\partial x^i}\right|_p(x^j)=\delta_i^j\)，对任何线性关系逐个作用于 \(x^j\)，便得到全部系数为零，所以这些向量线性无关，表示也唯一。
\end{proof}

公式中的 \(v(x^i)\) 是同一个切向量在 \(x\) 坐标下的分量；坐标函数本身只是附近的函数，不要求延长到整个流形。若换用坐标 \(y\)，则对每个 \(y^a\) 使用上述公式，得到
\begin{equation}
	\label{g1:eq:tangent-coordinate-change}
	v_y^a=\sum_i\frac{\partial y^a}{\partial x^i}(p)v_x^i,
	\qquad
	\left.\frac{\partial}{\partial x^i}\right|_p
	=\sum_a\frac{\partial y^a}{\partial x^i}(p)
	\left.\frac{\partial}{\partial y^a}\right|_p.
\end{equation}
这里 \(\partial y^a/\partial x^i\) 是 \(y\circ x^{-1}\) 的偏导。变换矩阵可逆，是因为正反坐标变换的导数由链式法则互逆；边界处同样成立。

\begin{theorem}[曲线速度与导子]
	\label{g1:thm:tangent-curve-derivation}
	设 \(p\) 是光滑流形 \(M\) 的内点。考虑满足 \(\gamma(0)=p\) 的光滑曲线 \(\gamma:(-\varepsilon,\varepsilon)\to M\)，其中 \(\varepsilon>0\)。若两条这样的曲线在某张、从而任意一张 \(p\) 附近坐标图中的 \(0\) 时刻导数相同，就称它们速度等价。所得等价类集合与 \(\mathrm{T}_pM\) 一一对应，对应把 \(\gamma\) 送到导子 \(v_\gamma(f)=(f\circ\gamma)'(0)\)，其中 \(f\) 是 \(p\) 附近的光滑函数芽。在任意坐标图中，
	\[
		v_\gamma=\sum_i(x^i\circ\gamma)'(0)
		\left.\frac{\partial}{\partial x^i}\right|_p.
	\]
\end{theorem}
\begin{proof}
	曲线给出导子及其坐标公式，已由链式法则说明。若两个导子相等，把它们作用于各坐标函数便得到相同坐标速度；反过来，相同坐标速度由坐标基公式给出相同导子，所以对应良定义且单射。任取 \(v=\sum_i v^i\partial_{x^i}|_p\)。因为 \(p\) 是内点，\(a=x(p)\) 的一个完整 Euclid 小球在图像内，故足够小的双侧时间上
	\[
		\gamma(t)=x^{-1}(a+t(v^1,\ldots,v^n))
	\]
	有定义，其坐标速度为指定的 \((v^1,\ldots,v^n)\)，所以 \(v_\gamma=v\)，这证明满射。
\end{proof}

现在可以把速度等价类相加：先把它们看成导子相加，再由定理选一条代表曲线。坐标里，这只是把速度列向量相加，但不需要把流形中的两个点相加。切向量也不是一条完整曲线；许多形状不同、加速度不同的曲线可代表同一个切向量。

\begin{example}[极坐标中的同一速度]
	\label{g1:exm:tangent-polar-coordinates}
	在平面原点之外的一个角扇区中，\(x=r\cos\theta\)、\(y=r\sin\theta\)。链式法则给出
	\[
		\partial_r=\cos\theta\,\partial_x+\sin\theta\,\partial_y,\qquad
		\partial_\theta=-r\sin\theta\,\partial_x+r\cos\theta\,\partial_y.
	\]
	所以同一向量 \(v=A\partial_r+B\partial_\theta\) 的直角坐标分量为
	\[
		(v^x,v^y)=(A\cos\theta-rB\sin\theta,\ A\sin\theta+rB\cos\theta).
	\]
	例如在 \((x,y)=(2,0)\) 处，\(\partial_\theta=2\partial_y\)。符号 \(\partial_\theta\) 指每单位角度变化的速度，而不是已经归一化的单位向量。这里没有改变切向量，只改变了用来记录它的坐标基。
\end{example}

\begin{proposition}[边界处的曲线速度]
	\label{g1:prop:boundary-velocities}
	设 \(M\) 为光滑 \(n\) 维带边界流形，\(n\ge1\)，\(p\in\partial M\)，边界坐标为 \(x\)，其中 \(x^n(p)=0\)。
	\begin{enumerate}
		\item 经过 \(p\) 的双侧光滑曲线 \(\gamma:(-\varepsilon,\varepsilon)\to M\) 所能给出的速度，恰为 \(v^n=0\) 的子空间。它自然识别为 \(\mathrm{T}_p(\partial M)\)。
		\item 从 \(p\) 出发的单侧光滑曲线 \(\gamma:[0,\varepsilon)\to M\) 所能给出的右速度，恰为 \(v^n\ge0\) 的闭半空间，称为\term[neixiangqiezhui]{内向切锥}[inward tangent cone]，记为 \(\mathrm{T}_p^+M\)。
	\end{enumerate}
	这两个集合都与边界图的选择无关；完整切空间 \(\mathrm{T}_pM\) 则允许任意 \(v^n\in\mathbb R\)。
\end{proposition}
\begin{proof}
	双侧曲线的函数 \(h(t)=x^n(\gamma(t))\) 非负且 \(h(0)=0\)，所以 \(0\) 为内点极小，\(h'(0)=0\)。反过来，任意 \(v^n=0\) 的坐标向量都由短曲线 \(x^{-1}(a+tv)\) 表示，它始终留在边界。单侧时，\(h(t)/t\ge0\) 对 \(t>0\) 成立，取右极限得到 \(h'_+(0)\ge0\)；任意 \(v^n\ge0\) 又由同一线性公式在 \(t\ge0\) 时实现。这就准确给出了两类速度。

	\cref{g1:prop:boundary-chart-invariance}已经说明 \(\partial M\) 自身为 \(n-1\) 维光滑流形。其导子 \(w\in \mathrm{T}_p(\partial M)\) 可作用于限制函数 \(f|_{\partial M}\)，从而给出 \(\mathrm{T}_pM\) 中的导子 \(f\mapsto w(f|_{\partial M})\)。在边界坐标中，这个线性映射把边界坐标基送到 \(\partial_{x^1},\ldots,\partial_{x^{n-1}}\)，所以单射，像恰为 \(v^n=0\)。两个集合都由实际曲线的速度定义，而曲线及其导子与坐标无关，故子空间和内向锥的描述也与坐标选择无关。
\end{proof}

\input{Figures/Chapter11-BoundaryVelocities.tex}

例如半直线 \([0,\infty)\) 在 \(0\) 处的切空间是 \(\mathbb R\)，坐标导子 \(\partial_x|_0\) 及 \(-\partial_x|_0\) 都合法；双侧曲线在 \(0\) 处的速度却只能为零，单侧速度则只能非负。因此若把“曲线速度”不加时间方向说明地当成边界切空间的定义，就会分别得到维数太小的空间或不对取负封闭的锥。导子定义保留了做线性代数所需的完整空间，曲线再标记其中哪些方向能够向前实际走出。

\subsection{微分、链式法则与切丛图册}
\label{g1:subsec:1102:03}

现在设 \(F:M\to N\) 是光滑映射，\(p\in M\)。若在 \(N\) 上观察函数 \(f\)，沿 \(F\) 的像曲线观察它，就是在 \(M\) 上观察 \(f\circ F\)。所以一个输入速度 \(v\) 应送到算子 \(f\mapsto v(f\circ F)\)。这个公式只用了复合，不需要先选择坐标。

\begin{definition}
	\label{g1:def:manifold-differential}
	设 \(F:M\to N\) 为光滑映射，流形可带边界。它在 \(p\) 处的\term[weifen]{微分}[differential]或切映射为
	\[
		\mathrm{d}F_p:\mathrm{T}_pM\longrightarrow \mathrm{T}_{F(p)}N,\qquad
		(\mathrm{d}F_pv)(f)=v(f\circ F).
	\]
\end{definition}

若两个函数在 \(F(p)\) 附近相等，其复合在 \(p\) 附近相等，故此公式对芽良定义。复合保留加法与乘法，而 \(f\circ F\) 在 \(p\) 的值为 \(f(F(p))\)，所以 \(\mathrm{d}F_pv\) 确实满足目标点的导子乘法法则。它对 \(v\) 线性。若 \(v=v_\gamma\) 是曲线速度，则
\[
	\mathrm{d}F_pv_\gamma=v_{F\circ\gamma},
\]
因为对任意函数 \(f\)，两边均为 \(\left.\frac{\mathrm{d}}{\mathrm{d}t}f(F(\gamma(t)))\right|_0\)。边界处的单侧曲线也有相同结论。

\begin{proposition}[坐标微分与链式法则]
	\label{g1:prop:differential-chain-rule}
	设 \(F:M^n\to N^m\)、\(G:N\to P\) 为光滑映射，可涉及带边界流形。取 \(p\) 与 \(F(p)\) 附近的坐标 \(x,y\)，记 \(F^a=y^a\circ F\circ x^{-1}\)。则
	\begin{equation}
		\label{g1:eq:manifold-differential-matrix}
		\mathrm{d}F_p\left(\sum_i v^i\partial_{x^i}|_p\right)
		=\sum_a\left(\sum_i\frac{\partial F^a}{\partial x^i}(x(p))v^i\right)
		\partial_{y^a}|_{F(p)}.
	\end{equation}
	并且
	\[
		\mathrm{d}(G\circ F)_p=\mathrm{d}G_{F(p)}\circ\mathrm{d}F_p,
		\qquad \mathrm{d}(\operatorname{id}_M)_p=\operatorname{id}_{\mathrm{T}_pM}.
	\]
	微分同胚的微分为线性同构，其逆为相应逆映射的微分。
\end{proposition}
\begin{proof}
	目标切向量的第 \(a\) 个分量是它作用于坐标函数 \(y^a\) 的结果，即 \(v(y^a\circ F)\)。用\eqref{g1:eq:tangent-coordinate-basis}展开，便得坐标公式。边界处的偏导数由光滑延拓定义且唯一，Euclid 链式法则延到边界的版本已在\cref{g1:prop:boundary-smooth-derivative}证明，所以同式仍成立。

	链式法则还可直接在导子层面核对。对 \(G(F(p))\) 附近的任意函数芽 \(f\)，有
	\[
		\bigl(\mathrm{d}(G\circ F)_p v\bigr)(f)
		=v(f\circ G\circ F)
		=\bigl(\mathrm{d}G_{F(p)}(\mathrm{d}F_pv)\bigr)(f).
	\]
	恒等映射的公式立即由定义得到。若 \(F\) 为微分同胚，把 \(F^{-1}\circ F\) 与 \(F\circ F^{-1}\) 的微分分别展开，便得到两个方向的互逆关系。
\end{proof}

这说明微分的矩阵就是多元微积分中的 Jacobian 矩阵，但线性映射本身不依赖坐标。它的秩记为 \(\operatorname{rank}\mathrm{d}F_p\)，换坐标只在矩阵两侧乘可逆矩阵，故秩不变。映射本身双射并不足以使微分可逆，例如 \(F:\mathbb R\to\mathbb R\)、\(F(t)=t^3\) 在零点的微分为零；它的逆虽连续，却在零点不可微。

\ProseParagraphBreak

各点的切空间都已经定义，现在才把它们组织为一个空间。令
\[
	\mathrm{T}M=\bigsqcup_{p\in M}\mathrm{T}_pM,\qquad \pi:\mathrm{T}M\longrightarrow M,\quad \pi(v_p)=p.
\]
不交并符号强调不同基点的向量属于不同的纤维，即使它们的坐标分量恰好相同，也不是同一个元素。对坐标图 \(x:U\to x(U)\)，定义双射
\[
	\Phi_x:\pi^{-1}(U)\longrightarrow x(U)\times\mathbb R^n,\qquad
	\Phi_x(v_p)=(x(p),v(x^1),\ldots,v(x^n)).
\]
前 \(n\) 个数记录向量位于哪里，后 \(n\) 个数记录它在该点坐标基下的分量。

\begin{theorem}[切丛的光滑图册]
	\label{g1:thm:tangent-bundle-atlas}
	设 \(M\) 为 Hausdorff、第二可数的光滑 \(n\) 维流形，可带边界。令 \(\mathrm{T}M=\coprod_{p\in M}\mathrm{T}_pM\)，并令 \(\pi(v_p)=p\)。对每张坐标图 \(x=(x^1,\ldots,x^n):U\to x(U)\)，在 \(\pi^{-1}(U)\) 上定义 \(\Phi_x(v_p)=(x(p),v_p(x^1),\ldots,v_p(x^n))\)。这些图唯一确定 \(\mathrm{T}M\) 上的拓扑与光滑结构，使其成为 Hausdorff、第二可数的 \(2n\) 维光滑流形，称为 \(M\) 的\term[qiecong]{切丛}[tangent bundle]。投影 \(\pi\) 光滑。若 \(M\) 带边界，则
	\[
		\partial(\mathrm{T}M)=\pi^{-1}(\partial M).
	\]
\end{theorem}
\begin{proof}
	首先计算过渡映射。设 \(h=y\circ x^{-1}\)，将切向量分量看作列向量，由\eqref{g1:eq:tangent-coordinate-change}可得
	\begin{equation}
		\label{g1:eq:tangent-bundle-transition}
		\Phi_y\circ\Phi_x^{-1}(z,v)=(h(z),\mathrm{D}h(z)v).
	\end{equation}
	它在交叠像上光滑，逆映射为 \((w,\eta)\mapsto(h^{-1}(w),\mathrm{D}(h^{-1})(w)\eta)\)，因而是微分同胚。边界情形中，这些公式由 \(h\) 的局部光滑延拓给出；延拓的偏导唯一，正反导数仍互逆。模型 \(x(U)\times\mathbb R^n\subset\mathbb H^n\times\mathbb R^n\) 在把非负坐标 \(z^n\) 重排到最后一位后，就是 \(\mathbb H^{2n}\) 的相对开集。

	定义集合 \(A\subset \mathrm{T}M\) 为开，当且仅当每个 \(\Phi_x(A\cap\pi^{-1}(U))\) 在相应模型中开。过渡映射为同胚，保证一张图中的开集在另一张图内仍为开；因此每个 \(\pi^{-1}(U)\) 开，每个 \(\Phi_x\) 是同胚，这些图组成拓扑图册。按相同公式，\(\pi\) 在局部为 \((z,v)\mapsto z\)，故连续且光滑。

	下面补足流形定义中的两个整体条件。若 \(v_p,w_q\) 的基点不同，利用 \(M\) 的 Hausdorff 性取不交基点邻域，其在 \(\pi\) 下的逆像分离两向量。若 \(p=q\) 但向量不同，在同一张 \(\Phi_x\) 图中，它们对应不同 Euclid 点，可取不交开邻域；这些图内开集也是 \(\mathrm{T}M\) 的开集。所以 \(\mathrm{T}M\) 为 Hausdorff。

	由 \(M\) 第二可数，可以从任意坐标覆盖中选出可数子覆盖：对每个包含在某图域内的可数基元素，选一张包含它的图；每个点附近都有这样的基元素，所选图遂覆盖 \(M\)。记这些图为 \(U_1,U_2,\ldots\)。每个 \(x_j(U_j)\times\mathbb R^n\) 都有可数基，例如取有理中心、半径的 Euclid 球再与模型相交；把这些基通过 \(\Phi_{x_j}^{-1}\) 拉回，再对 \(j\) 取并，得到 \(\mathrm{T}M\) 的可数基。故 \(\mathrm{T}M\) 第二可数。

	各图的维数为 \(2n\)，光滑过渡已经验证。边界图中的非负坐标是基点坐标 \(z^n\)，与纤维分量 \(v^n\) 无关，因此边界恰为基点落在 \(\partial M\) 上的全部向量。最后，任何使所列图为坐标图的拓扑都必须以其图内开集为基，故拓扑唯一；光滑结构由同一相容图册唯一确定。
\end{proof}

\(\partial(\mathrm{T}M)\) 与 \(\mathrm{T}(\partial M)\) 不相同。前者包含边界点处全部 \(n\) 维切向量，维数为 \(2n-1\)；后者只包含沿边界的向量，在边界坐标中还满足 \(v^n=0\)，维数为 \(2n-2\)。例如半直线的切丛是 \([0,\infty)\times\mathbb R\)，其边界是一整条纤维；边界点本身作为零维流形，切空间却只有零向量。

\ProseParagraphBreak

切丛还把逐点微分组成一个光滑映射。对光滑 \(F:M\to N\)，定义
\[
	\mathrm{T}F:\mathrm{T}M\longrightarrow \mathrm{T}N,\qquad \mathrm{T}F(v_p)=\mathrm{d}F_p(v_p).
\]
由\eqref{g1:eq:manifold-differential-matrix}，它在切丛坐标中的表达为
\[
	(z,v)\longmapsto(\widehat F(z),\mathrm{D}\widehat F(z)v),
	\qquad \widehat F=y\circ F\circ x^{-1},
\]
所以确实光滑。并有 \(\pi_N\circ \mathrm{T}F=F\circ\pi_M\)、\(\mathrm{T}(G\circ F)=\mathrm{T}G\circ \mathrm{T}F\)。零向量也给出光滑映射 \(0_M:M\to \mathrm{T}M\)，其局部公式为 \(z\mapsto(z,0)\)，称为零截面。这里“丛”所记录的基本事实，是在每张图上 \(\mathrm{T}M\) 看起来像“基点邻域乘一个向量空间”，而换图时向量分量按随基点变化的可逆线性映射变换。

\subsection{余切空间和拉回}
\label{g1:subsec:1102:04}

切向量记录变化的方向和速度；要从一个速度读出一个数，还需要线性函数。例如一个光滑实函数 \(f\) 在点 \(p\) 的变化率，就是 \(v\mapsto v(f)\)。这类线性观测应放在切空间的对偶中。

\begin{definition}
	\label{g1:def:cotangent-space}
	设 \(M\) 为光滑流形，可带边界，\(p\in M\)。实对偶空间
	\[
		\mathrm{T}_p^*M=\operatorname{Hom}_{\mathbb R}(\mathrm{T}_pM,\mathbb R)
	\]
	称为 \(p\) 处的\term[yuqiekongjian]{余切空间}[cotangent space]，其中的元素称为\term[xiebianxiangliang]{协变向量}[covector]。对 \(p\) 附近的光滑实函数 \(f\)，定义 \(\mathrm{d}f_p(v)=v(f)\)。\symidx{\(\mathrm{T}_p^*M\)：点处的余切空间}
\end{definition}

这里的 \(\mathrm{d}f_p\) 与把 \(f:M\to\mathbb R\) 看作映射时的微分相符：利用 \(\partial_t|_{f(p)}\) 将 \(\mathrm{T}_{f(p)}\mathbb R\) 识别为 \(\mathbb R\)，前一小节的微分就给出这个实值线性函数。对坐标函数有
\[
	\mathrm{d}x^i_p\left(\partial_{x^j}|_p\right)=\delta^i_j,
\]
所以 \(\mathrm{d}x^1_p,\ldots,\mathrm{d}x^n_p\) 恰为坐标切基的对偶基。每个协变向量唯一写为
\[
	\alpha=\sum_i\alpha_i\,\mathrm{d}x^i_p,\qquad
	\alpha_i=\alpha(\partial_{x^i}|_p),
	\qquad
	\mathrm{d}f_p=\sum_i\partial_i(f\circ x^{-1})(x(p))\,\mathrm{d}x^i_p.
\]
因而每个协变向量都是某个局部函数的微分：取 \(f=\sum_i\alpha_i x^i\) 即可。这里的系数 \(\alpha_i\) 是指定点处的实数，不是把任意随点变化的一形式都断言为某个函数的微分。

\ProseParagraphBreak

切空间与余切空间同维，但在没有额外结构时，并没有指定它们之间的自然识别。梯度需要先有内积，才能由 \(\langle\operatorname{grad}f,v\rangle=\mathrm{d}f(v)\) 求出一个向量；微分 \(\mathrm{d}f\) 本身无需度量。例如平面上的 \(\mathrm{d}x\) 固定不变，用 Euclid 度量时它对应 \(\partial_x\)，用 \(2\,\mathrm{d}x^2+\mathrm{d}y^2\) 时则对应 \(\frac12\partial_x\)。

\begin{proposition}[余切拉回]
	\label{g1:prop:covector-pullback}
	设 \(F:M\to N\)、\(G:N\to P\) 为光滑映射，流形可带边界；取 \(p\in M\)，并令 \(f\) 为 \(F(p)\) 附近的光滑实函数。其微分的对偶映射为
	\[
		(\mathrm{d}F_p)^*:\mathrm{T}_{F(p)}^*N\longrightarrow \mathrm{T}_p^*M,\qquad
		(\mathrm{d}F_p)^*\alpha=\alpha\circ\mathrm{d}F_p,
	\]
	称为协变向量的\term[lahui]{拉回}[pullback]。它满足
	\[
		(\mathrm{d}F_p)^*(\mathrm{d}f_{F(p)})=\mathrm{d}(f\circ F)_p,
	\]
	对复合映射还有
	\[
		\bigl(\mathrm{d}(G\circ F)_p\bigr)^*
		=(\mathrm{d}F_p)^*\circ(\mathrm{d}G_{F(p)})^*.
	\]
	若 \(\mathrm{d}F_p\) 的矩阵为 \(J\)，则对偶基下的拉回矩阵为 \(J^{\mathsf T}\)。
\end{proposition}
\begin{proof}
	线性映射与线性函数的复合仍为线性函数，故拉回良定义。对任意 \(v\in \mathrm{T}_pM\)，
	\[
		\bigl((\mathrm{d}F_p)^*\mathrm{d}f_{F(p)}\bigr)(v)
		=(\mathrm{d}F_pv)(f)=v(f\circ F)
		=\mathrm{d}(f\circ F)_p(v),
	\]
	得到第一式。复合公式由微分的链式法则及函数复合的结合性得到：先将 \(v\) 经 \(\mathrm{d}F_p\)、\(\mathrm{d}G_{F(p)}\) 向前送，再由目标协变向量读数，等价于先后反向拉回这两个线性观测。最后，列向量 \(v\) 的像为 \(Jv\)，协变向量列系数 \(\alpha\) 对它的读数为 \(\alpha^{\mathsf T}Jv=(J^{\mathsf T}\alpha)^{\mathsf T}v\)，所以矩阵是转置。
\end{proof}

拉回总能定义，不要求 \(F\) 可逆。反之，要把源点的一个协变向量向前搬到目标点，通常没有唯一办法：若 \(\mathrm{d}F_p\) 压缩了方向或漏掉了目标方向，便可能不相容或有多种延长。因此“向量向前、协变向量向后”来自定义中的复合次序，而不是单纯的记号偏好。

\ProseParagraphBreak

也可将余切空间合成 \(\mathrm{T}^*M=\bigsqcup_p\mathrm{T}_p^*M\)，称为余切丛。将 \(\alpha_p=\sum_i a_i\,\mathrm{d}x^i_p\) 记录为 \((x(p),a)\)。若 \(y=h(x)\)，则 \(\mathrm{d}y=\mathrm{D}h\,\mathrm{d}x\) 表示各坐标微分的线性关系，因此同一个协变向量的系数满足 \(a_x=\mathrm{D}h^{\mathsf T}a_y\)，过渡映射为
\[
	(z,a)\longmapsto\bigl(h(z),\mathrm{D}h(z)^{-{\mathsf T}}a\bigr).
\]
它光滑可逆：矩阵求逆可由余子式除以非零行列式表示。按\cref{g1:thm:tangent-bundle-atlas}的同一构造赋予拓扑，图内过渡为同胚；基点不同时用基空间分离，同基点时用一张图分离；可数坐标覆盖和局部可数基给出第二可数性。于是 \(\mathrm{T}^*M\) 也为 \(2n\) 维光滑流形，带边界时其边界同样是基点落在 \(\partial M\) 的全部余切向量。切向量与协变向量的两种换基公式由此分别构成切丛与余切丛。

\ProseParagraphBreak

如果给 \(M\) 的每一点选定协变向量 \(\alpha_p\)，且在坐标表达 \(\alpha=\sum_i a_i(x)\,\mathrm{d}x^i\) 中各系数光滑，就称它为光滑一形式。上述换基公式说明这与坐标无关。对于 \(N\) 上的一形式 \(\alpha\)，逐点定义
\[
	(F^*\alpha)_p=(\mathrm{d}F_p)^*\alpha_{F(p)}.
\]
在坐标中它是
\[
	F^*\left(\sum_a b_a(y)\,\mathrm{d}y^a\right)
	=\sum_i\left(\sum_a b_a(\widehat F(x))
	\frac{\partial\widehat F^a}{\partial x^i}(x)\right)\mathrm{d}x^i,
\]
所以仍光滑，并满足 \((G\circ F)^*=F^*\circ G^*\)、\(F^*(\mathrm{d}f)=\mathrm{d}(f\circ F)\)。这里仅使用一形式的线性含义；外积和外微分将另行定义。
% 纲要条目排为编号小节标题，并进入目录与书签；说明段仍保留于纲要注释中。

% !TeX root = ../../Book1.tex
% 纲要标题：### 11.3 局部逆与秩定理
% 纲要位置：几何学计划纲要.md，第 551 行。
% 正文状态：初稿；尚未正式审核。

% 纲要内容开始
% 
% * 逆函数定理的定量局部证明
% * 隐函数定理与参数依赖
% * 常秩定理的坐标约化
% * 秩不恒定时的反例
% 
% 纲要内容结束

\section{局部逆与秩定理}
\label{g1:sec:1103}

微分把光滑映射在一点附近的变化近似为线性映射。如果这个线性映射可逆，原映射是否也能在附近反解？如果微分只保留一部分方向，能否用坐标把这部分方向同被压缩的方向分开？这两个问题分别由逆函数定理与常秩定理回答。本节从有限维微积分和线性代数重新证明所需结论，具体说明邻域如何选择、反解为何光滑，以及秩的条件用在什么地方。

\subsection{逆函数定理的定量局部证明}
\label{g1:subsec:1103:01}

先把问题放在 \(\mathbb R^n\) 中。Euclid 范数记为 \(|\cdot|\)，矩阵使用对应的算子范数 \(\|A\|=\sup_{|v|=1}|Av|\)。记 \(B_r(a)\) 为开球，\(\overline B_r(a)\) 为闭球；省略中心时指中心为零。对 Euclid 映射写 \(\mathrm{D}f(x)\) 表示导数矩阵，它表示上一节中微分 \(\mathrm{d}f_x\) 在标准坐标基下的矩阵，并不是一种新的求导运算。

\ProseParagraphBreak

如果已经通过平移与可逆线性变换化成 \(f(0)=0\)、\(\mathrm{D}f(0)=I\)，那么 \(f(x)\) 在原点附近接近 \(x\)。将 \(f(x)=y\) 改写为
\[
	x=y-\bigl(f(x)-x\bigr)
\]
就给出一个逐步修正的办法：先猜一个 \(x\)，再减去它偏离目标的误差。我们需要的不是这一步误差必然很小，而是两次猜测经过修正后彼此更接近。下面的引理把这个思路变成收敛证明。

\begin{lemma}[压缩映射]
	\label{g1:lem:euclidean-contraction}
	设 \(K=\overline B_r(a)\subset\mathbb R^n\) 是非空闭球，\(T:K\to K\) 满足
	\[
		|T(x)-T(x')|\leq q|x-x'|,\qquad 0\leq q<1.
	\]
	则 \(T\) 有唯一不动点 \(x_*\)。从任意 \(x_0\in K\) 出发，递推 \(x_{j+1}=T(x_j)\)，均收敛到 \(x_*\)；当 \(0<q<1\) 时有
	\[
		|x_j-x_*|\leq\frac{q^j}{1-q}|x_1-x_0|.
	\]
\end{lemma}
\begin{proof}
	若 \(q=0\)，映射 \(T\) 为常值，第二次迭代起就不再变化，结论立即成立。以下设 \(0<q<1\)。由归纳，
	\[
		|x_{j+1}-x_j|\leq q^j|x_1-x_0|.
	\]
	对 \(k>j\)，三角不等式和几何级数给出
	\[
		|x_k-x_j|\leq\sum_{\ell=j}^{k-1}|x_{\ell+1}-x_\ell|
		\leq\frac{q^j}{1-q}|x_1-x_0|.
	\]
	因此各坐标都是实数 Cauchy 序列，逐坐标收敛给出 \(\mathbb R^n\) 中的极限 \(x_*\)。闭球的闭性保证 \(x_*\in K\)，而压缩不等式蕴含 \(T\) 连续，所以取极限得 \(x_*=T(x_*)\)。让上式中的 \(k\to\infty\)，也得到误差估计。若还有不动点 \(z_*\)，则
	\[
		|x_*-z_*|=|T(x_*)-T(z_*)|\leq q|x_*-z_*|,
	\]
	因 \(q<1\)，只能有 \(x_*=z_*\)。
\end{proof}

闭球在这里承担两个不同职责：迭代不能跑出定义域，极限也不能掉到定义域之外。严格不等式 \(q<1\) 则使累计误差成为收敛的几何级数。仅有“距离不增加”不能给出同一结论，例如闭球上的恒等映射有许多不动点。

\begin{theorem}[定量局部逆]
	\label{g1:thm:quantitative-inverse-function}
	设 \(n\geq1\)，\(f:U\to\mathbb R^n\) 为 \(C^k\) 映射，其中 \(U\subset\mathbb R^n\) 开，\(1\leq k\leq\infty\)。假设 \(\overline B_r\subset U\)、\(r>0\)，且
	\[
		f(0)=0,\qquad \mathrm{D}f(0)=I,\qquad
		\sup_{x\in\overline B_r}\|\mathrm{D}f(x)-I\|\leq q<1,\quad q\geq0.
	\]
	则 \(f\) 在 \(\overline B_r\) 上单射，且对 \(x,x'\in\overline B_r\)，
	\begin{equation}
		\label{g1:eq:quantitative-inverse-bounds}
		(1-q)|x-x'|\leq|f(x)-f(x')|\leq(1+q)|x-x'|.
	\end{equation}
	像 \(\Omega=f(B_r)\) 是开集，并满足
	\[
		B_{(1-q)r}\subset\Omega.
	\]
	映射 \(f:B_r\to\Omega\) 是 \(C^k\) 微分同胚。其逆映射 \(g\) 满足
	\begin{equation}
		\label{g1:eq:inverse-derivative}
		\mathrm{D}g(y)=\bigl(\mathrm{D}f(g(y))\bigr)^{-1},\qquad
		|g(y)-g(y')|\leq\frac{|y-y'|}{1-q}.
	\end{equation}
	特别地，对 \(|y|<(1-q)r\)，方程 \(f(x)=y\) 在 \(B_r\) 内有唯一解，而且 \(|x|\leq|y|/(1-q)\)。
\end{theorem}
\begin{proof}
	令 \(E(x)=f(x)-x\)。闭球是凸的，所以任意两点间的线段仍在闭球中。对函数 \(t\mapsto E(x'+t(x-x'))\) 逐分量应用一元微积分基本定理，得到
	\[
		E(x)-E(x')=\int_0^1\mathrm{D}E(x'+t(x-x'))(x-x')\,\mathrm{d}t.
	\]
	因此 \(|E(x)-E(x')|\leq q|x-x'|\)。在 \(f(x)-f(x')=(x-x')+(E(x)-E(x'))\) 中分别使用三角不等式及其反向形式，便得\eqref{g1:eq:quantitative-inverse-bounds}，特别地 \(f\) 单射。

	固定 \(|y|\leq(1-q)r\)，令 \(T_y(x)=y-E(x)\)。因为 \(E(0)=0\)，对 \(x\in\overline B_r\) 有
	\[
		|T_y(x)|\leq|y|+q|x|\leq r.
	\]
	所以它把闭球映入自身，且压缩因子至多为 \(q\)。\cref{g1:lem:euclidean-contraction}给出唯一不动点 \(x\)，不动点方程正是 \(f(x)=y\)。再由\eqref{g1:eq:quantitative-inverse-bounds}取 \(x'=0\)，得到 \(|x|\leq|y|/(1-q)\)。当目标满足严格不等式时，解在开球内，从而证明 \(B_{(1-q)r}\subset f(B_r)\)。

	还须证明整个像 \(\Omega\) 开，而不能只知道它包含原点附近的一个球。任取 \(x_0\in B_r\)，选 \(\rho>0\) 使 \(\overline B_\rho(x_0)\subset B_r\)。对
	\[
		|y-f(x_0)|<(1-q)\rho
	\]
	仍使用同一个 \(T_y\)。此时
	\[
		|T_y(x)-x_0|
		\leq|y-f(x_0)|+|E(x)-E(x_0)|
		\leq|y-f(x_0)|+q\rho<\rho.
	\]
	它把这个较小闭球映入自身，故每个这样的 \(y\) 都在 \(f(B_r)\) 中。这就证明了 \(\Omega\) 开。\eqref{g1:eq:quantitative-inverse-bounds}又立即给出逆映射的 Lipschitz 估计，所以 \(g\) 首先是连续的。

	对任意 \(x\in B_r\) 与向量 \(v\)，导数条件给出 \(|\mathrm{D}f(x)v|\geq(1-q)|v|\)。故 \(\mathrm{D}f(x)\) 单射；由于定义域和值域维数同为 \(n\)，它可逆，并且 \(\|\mathrm{D}f(x)^{-1}\|\leq(1-q)^{-1}\)。固定 \(y=f(x)\)，对充分小的 \(h\) 令 \(\xi=g(y+h)-x\)。逆映射的 Lipschitz 估计给出 \(|\xi|\leq|h|/(1-q)\)。利用 \(f\) 在 \(x\) 的可微性，
	\[
		h=\mathrm{D}f(x)\xi+o(|\xi|),
		\qquad
		\xi=\mathrm{D}f(x)^{-1}h+o(|h|).
	\]
	这正是 \(g\) 在 \(y\) 可微及\eqref{g1:eq:inverse-derivative}的含义。

	矩阵求逆在可逆矩阵集合上光滑：每个矩阵元由伴随矩阵的多项式除以非零行列式给出。因 \(\mathrm{D}f\) 与 \(g\) 连续，导数公式先使 \(\mathrm{D}g\) 连续，故 \(g\in C^1\)。若 \(f\in C^k\)，且已经知道 \(g\in C^j\)、\(j<k\)，则 \(\mathrm{D}f\in C^{k-1}\)，同一公式使 \(\mathrm{D}g\in C^j\)，于是 \(g\in C^{j+1}\)。逐阶归纳得到 \(C^k\)；当 \(k=\infty\) 时，对每个有限阶作此归纳即可。
\end{proof}

这里的定量信息不仅是逆映射存在。\((1-q)r\) 给出一个必定包含在像中的目标半径，\((1-q)^{-1}\) 控制反解对目标误差的放大。实际像可以更大；定理只承诺由当前导数界能够保证的部分。

\begin{theorem}[逆函数定理]
	\label{g1:thm:inverse-function}
	设 \(f:U\subset\mathbb R^n\to\mathbb R^n\) 为 \(C^k\) 映射，\(U\) 开，\(1\leq k\leq\infty\)。若 \(\mathrm{D}f(a)\) 可逆，则存在 \(a\) 与 \(f(a)\) 的开邻域 \(U_0,V_0\)，使 \(f:U_0\to V_0\) 为 \(C^k\) 微分同胚。
	
	相应地，若 \(M,N\) 是同维无边界光滑流形，\(F:M\to N\) 光滑，且 \(\mathrm{d}F_p:\mathrm{T}_pM\to \mathrm{T}_{F(p)}N\) 为线性同构，则 \(F\) 在 \(p\) 的某个开邻域上为到 \(F(p)\) 开邻域的微分同胚。
\end{theorem}
\begin{proof}
	设 \(A=\mathrm{D}f(a)\)。对 \(h\) 接近零，定义
	\[
		\widetilde f(h)=A^{-1}\bigl(f(a+h)-f(a)\bigr).
	\]
	则 \(\widetilde f(0)=0\)、\(\mathrm{D}\widetilde f(0)=I\)。导数的连续性允许选取 \(r>0\) 及 \(q<1\)，例如 \(q=1/2\)，使闭球包含在定义域中，且 \(\|A^{-1}\mathrm{D}f(a+h)-I\|\leq q\) 在闭球上成立。应用\cref{g1:thm:quantitative-inverse-function}，再还原平移与线性变换，便得第一部分。

	流形情形，在 \(p\) 与 \(F(p)\) 附近选坐标图。\cref{g1:prop:differential-chain-rule}说明，坐标表达的微分是在 \(\mathrm{d}F_p\) 两侧复合坐标微分同构，故也是可逆矩阵。将 Euclid 结论应用于这一坐标表达，再用坐标图送回流形，得到所需邻域及光滑逆。零维情形的坐标域是单点，结论直接成立。
\end{proof}

一般情形仍有明确尺度。若归一化后的导数界在 \(\overline B_r(a)\) 上为 \(q\)，则
\[
	f(a)+A B_{(1-q)r}\subset f(B_r(a)).
\]
特别地，像包含以 \(f(a)\) 为中心、半径 \((1-q)r/\|A^{-1}\|\) 的 Euclid 球；在这个像上的逆映射满足 Lipschitz 常数不超过 \(\|A^{-1}\|/(1-q)\)。这里先乘 \(A^{-1}\) 的作用，是把原映射在该点不同方向的线性伸缩统一成恒等变换。

\begin{example}
	\label{g1:exm:local-global-inverse}
	映射 \(f(s,t)=(e^s\cos t,e^s\sin t)\) 的导数行列式为 \(e^{2s}>0\)，因此在每一点附近都可光滑反解。但 \(f(s,t+2\pi)=f(s,t)\)，故它不是全局单射。逆函数定理确定的是一个足够小的邻域，并没有把不同邻域的逆自动拼成全局逆。
\end{example}

\subsection{隐函数定理与参数依赖}
\label{g1:subsec:1103:02}

方程 \(F(x,y)=0\) 往往已经给出一个几何对象，却没有把它写成函数图像。若 \(y\) 方向的变化能独立控制全部方程，就有希望把 \(y\) 解成 \(x\) 的函数。为把这个问题转成逆函数问题，保留 \(x\) 不动，把新的另一组坐标取为方程值 \(F\)；然后令这组新坐标等于零。

\begin{theorem}[带参数的隐函数定理]
	\label{g1:thm:implicit-function-parameters}
	设 \(\Omega\subset\mathbb R^p\times\mathbb R^n\times\mathbb R^m\) 开，\(m\geq1\)，\(p,n\geq0\)，且 \(F:\Omega\to\mathbb R^m\) 为 \(C^k\) 映射，\(1\leq k\leq\infty\)。将变量写成 \((\lambda,x,y)\)，并假设
	\[
		F(\lambda_0,x_0,y_0)=0,\qquad \mathrm{D}_yF(\lambda_0,x_0,y_0)\ \text{可逆}.
	\]
	则存在三个点各自的开邻域 \(P,U,V\)，使 \(P\times U\times V\subset\Omega\)，并存在唯一 \(C^k\) 映射 \(\varphi:P\times U\to V\)，满足
	\begin{equation}
		\label{g1:eq:implicit-graph}
		F(\lambda,x,y)=0
		\quad\Longleftrightarrow\quad
		y=\varphi(\lambda,x)
		\qquad((\lambda,x,y)\in P\times U\times V).
	\end{equation}
	缩小这些邻域后，\(\mathrm{D}_yF\) 在其中可逆，而且
	\begin{equation}
		\label{g1:eq:implicit-derivative}
		\mathrm{D}_{(\lambda,x)}\varphi
		=-\bigl(\mathrm{D}_yF(\lambda,x,\varphi)\bigr)^{-1}
		\mathrm{D}_{(\lambda,x)}F(\lambda,x,\varphi).
	\end{equation}
\end{theorem}
\begin{proof}
	合并记 \(z=(\lambda,x)\)，并构造
	\[
		H(z,y)=(z,F(z,y)).
	\]
	其在基点的导数矩阵为
	\[
		\mathrm{D}H=\begin{pmatrix}I&0\\\mathrm{D}_zF&\mathrm{D}_yF\end{pmatrix}.
	\]
	这个矩阵可逆：若给定输出增量 \((a,b)\)，先取输入的 \(z\) 增量为 \(a\)，再取 \(y\) 增量为 \((\mathrm{D}_yF)^{-1}(b-\mathrm{D}_zF\,a)\)，就得到唯一解。由\cref{g1:thm:inverse-function}，\(H\) 在基点附近是微分同胚。因为第一组坐标始终保持为 \(z\)，它的逆必写成
	\[
		H^{-1}(z,w)=(z,\psi(z,w)).
	\]
	令 \(\varphi(z)=\psi(z,0)\)，便得到光滑候选解。

	需要确保存在性与唯一性在同一个乘积邻域上成立。先在 \(H\) 的可逆邻域内取基点的一个小乘积盒 \(Z_1\times V\)。目标可逆邻域是开集，因此缩小 \(Z_1\) 后，所有 \((z,0)\) 都落在目标可逆邻域中；再由 \(\psi(z_0,0)=y_0\) 及连续性，进一步缩小 \(Z_1\)，使 \(\psi(z,0)\in V\)。最后在这组 \(z\) 邻域中取乘积 \(P\times U\)。于是每个 \((\lambda,x)\in P\times U\) 都有解 \(\varphi(\lambda,x)\in V\)；任何另一个 \(y\in V\) 若使 \(F=0\)，就与该解经 \(H\) 具有同一像 \((\lambda,x,0)\)，由 \(H\) 在所取邻域中的单射性，二者相等。这证明\eqref{g1:eq:implicit-graph}。

	行列式连续且基点处非零，故还可把乘积邻域缩小到 \(\mathrm{D}_yF\) 始终可逆。对恒等式 \(F(\lambda,x,\varphi(\lambda,x))=0\) 使用链式法则，得到
	\[
		\mathrm{D}_{(\lambda,x)}F+\mathrm{D}_yF\,\mathrm{D}_{(\lambda,x)}\varphi=0,
	\]
	左乘 \((\mathrm{D}_yF)^{-1}\) 便得导数公式。
\end{proof}

参数 \(\lambda\) 在证明中也是变量。把它并入 \(z\)，使我们一次得到对参数与自变量的联合光滑性，而不只知道“每个固定参数分别有一个解”。唯一性也有明确范围：它只排除所选 \(V\) 内的第二个解，不排除其他地方还有另一条分支。

\begin{example}
	\label{g1:exm:parameter-circle-implicit}
	令 \(F(\lambda,x,y)=x^2+y^2-\lambda\)。在 \((\lambda,x,y)=(1,0,1)\) 处有 \(\mathrm{D}_yF=2\)，所以在附近存在唯一接近 \(1\) 的解
	\[
		y=\varphi(\lambda,x)=\sqrt{\lambda-x^2}.
	\]
	导数公式给出 \(\partial_\lambda\varphi=1/(2\varphi)\)、\(\partial_x\varphi=-x/\varphi\)，与直接求导一致。附近另有接近 \(-1\) 的解，但它不在当前的 \(V\) 中。到了 \(\lambda=x^2\) 且 \(y=0\) 时，\(\mathrm{D}_yF=0\)，这条光滑反解不能照原坐标方式继续；当 \(\lambda<x^2\) 时更没有实数解。对参数的局部光滑依赖并不承诺参数任意变化后仍有同一条光滑分支。
\end{example}

隐函数定理也没有指定哪些字母天生是未知量。若另一组 \(m\) 个变量对应的偏导数块可逆，就可以把它们作为待解变量；余下变量成为自变量或参数。选择可逆子块的办法，正是常秩定理的起点。

\subsection{常秩定理的坐标约化}
\label{g1:subsec:1103:03}

现在允许源空间与目标空间维数不同。一个秩为 \(r\) 的线性映射，经过两边换基可写成 \((u_1,\ldots,u_m)\mapsto(u_1,\ldots,u_r,0,\ldots,0)\)。非线性映射未必等于其线性近似，但如果微分的秩在附近始终为 \(r\)，就能用光滑坐标变化得到同样的局部形式。证明会先把 \(r\) 个输出直接用作源坐标，再说明其余输出不再含新的独立变化。

\begin{theorem}[常秩定理]
	\label{g1:thm:constant-rank}
	设 \(f:U\subset\mathbb R^m\to\mathbb R^n\) 为光滑映射，\(U\) 开，且 \(\operatorname{rank}\mathrm{D}f(x)=r\) 在 \(a\) 的某个邻域内恒成立，其中 \(0\leq r\leq\min(m,n)\)。则存在 \(a\) 的开邻域 \(U_0\)、\(f(a)\) 的开邻域 \(V_0\)，以及光滑坐标变换
	\[
		\phi:U_0\longrightarrow\widetilde U,\qquad
		\psi:V_0\longrightarrow\widetilde V,
	\]
	使 \(f(U_0)\subset V_0\)、\(\phi(a)=0\)、\(\psi(f(a))=0\)，并且
	\begin{equation}
		\label{g1:eq:constant-rank-normal-form}
		(\psi\circ f\circ\phi^{-1})(u_1,\ldots,u_m)
		=(u_1,\ldots,u_r,0,\ldots,0).
	\end{equation}
	最后一组零共有 \(n-r\) 个。
\end{theorem}
\begin{proof}
	先设 \(r>0\)。平移源与目标，使 \(a=0\)、\(f(a)=0\)。秩为 \(r\) 意味着 \(\mathrm{D}f(0)\) 有一个 \(r\) 阶子式非零；重排源、目标坐标后，可设左上角的 \(r\times r\) 子块可逆。记 \(x=(x',x'')\)，其中 \(x'\in\mathbb R^r\)、\(x''\in\mathbb R^{m-r}\)，并定义
	\[
		H(x)=\bigl(f_1(x),\ldots,f_r(x),x''\bigr).
	\]
	它的导数为分块上三角矩阵，左上角正是所选可逆块，右下角是 \(I_{m-r}\)，故 \(\mathrm{D}H(0)\) 可逆。逆函数定理给出 \(H\) 的局部光滑逆。把新源坐标记为 \((u,v)=H(x)\)，则
	\[
		(f\circ H^{-1})(u,v)=(u,G(u,v)).
	\]
	这一步的含义是把已经知道彼此独立的 \(r\) 个输出变成前 \(r\) 个输入坐标。

	链式法则以及 \(\mathrm{D}H^{-1}\) 可逆，保证改换源坐标后秩仍为 \(r\)。另一方面，
	\[
		\mathrm{D}(f\circ H^{-1})
		=\begin{pmatrix}I_r&0\\\mathrm{D}_uG&\mathrm{D}_vG\end{pmatrix}.
	\]
	从下方各行减去 \(\mathrm{D}_uG\) 给出的前 \(r\) 行线性组合，得到 \(\begin{psmallmatrix}I_r&0\\0&\mathrm{D}_vG\end{psmallmatrix}\)，其秩是 \(r+\operatorname{rank}\mathrm{D}_vG\)。因此常秩假设迫使 \(\mathrm{D}_vG=0\) 在整个所取邻域内成立；只知道基点处的秩为 \(r\)，无法得到这一步。

	缩小到以零为中心的乘积球 \(B'\times B''\)，使整个乘积都位于新坐标域中。对固定 \(u\)，线段 \(t\mapsto tv\) 始终在 \(B''\) 中，故微积分基本定理给出
	\[
		G(u,v)-G(u,0)=\int_0^1\mathrm{D}_vG(u,tv)v\,\mathrm{d}t=0.
	\]
	令 \(g(u)=G(u,0)\)，于是 \(f\circ H^{-1}(u,v)=(u,g(u))\)。最后在目标取坐标变化
	\[
		K(y',y'')=\bigl(y',y''-g(y')\bigr),\qquad
		K^{-1}(z',z'')=\bigl(z',z''+g(z')\bigr).
	\]
	这两个公式明确互逆且光滑。为同时安排定义域，可取
	\[
		V_0=\{(y',y''):y'\in B',\ y''-g(y')\in B_\varepsilon^{\,n-r}\},
	\]
	其中 \(\varepsilon>0\)；它是包含 \(f(0)\) 的开邻域，并包含 \(f(H^{-1}(B'\times B''))\)。于是令 \(\phi=H\)、\(\psi=K\)，并限制到这些邻域，就得到标准形。

	当 \(r=m\) 时，变量 \(v\) 不存在，以上 \(G\) 本来就只依赖 \(u\)；当 \(r=n\) 时，没有余下输出 \(G\)，目标变换可以取恒等，以上步骤仍成立。若 \(r=0\)，则 \(\mathrm{D}f=0\) 在一个小凸球上恒成立，沿任意线段积分可知 \(f\) 在该球上常值；将这个常值移到目标原点，便得到全部分量为零的标准形。涉及零维 Euclid 空间时，也正是这些没有相应变量的情形。
\end{proof}

上述证明中的两个坐标变化有不同任务：源坐标 \(H\) 选出独立输出，目标坐标 \(K\) 再减掉其余输出对它们的依赖。不是直接把 \(\mathrm{D}f(a)\) 化成标准矩阵后就停止；决定非线性映射本身能否化简的，是邻域中处处成立的 \(\mathrm{D}_vG=0\)。

\begin{example}
	\label{g1:exm:constant-rank-curved-fibres}
	取 \(f(x,y)=(x+y^2,(x+y^2)^2)\)。令 \(u=x+y^2\)，其导数矩阵为
	\[
		\mathrm{D}f(x,y)=\begin{pmatrix}1&2y\\2u&4uy\end{pmatrix},
	\]
	秩处处为 \(1\)。选择
	\[
		H(x,y)=(x+y^2,y),\qquad K(a,b)=(a,b-a^2),
	\]
	得到 \(H^{-1}(u,v)=(u-v^2,v)\) 及 \(K\circ f\circ H^{-1}(u,v)=(u,0)\)。原来的同值曲线 \(x+y^2=c\) 被源坐标拉直成 \(u=c\)，原来的像抛物线 \(b=a^2\) 被目标坐标拉直成第二坐标为零的直线。\cref{g1:fig:constant-rank-coordinates}将这两个操作分别画出。
\end{example}

\begin{figure}[htbp]
	\centering
	\input{Figures/Chapter11-ConstantRankCoordinates}
	\caption[常秩坐标如何同时整理原像与像]{例中的映射 \(f(x,y)=(x+y^2,(x+y^2)^2)\)。上排是原坐标，下排是经 \(H(x,y)=(x+y^2,y)\) 与 \(K(a,b)=(a,b-a^2)\) 改换后的坐标。同色源曲线映到同色目标点；新坐标使各条原像成为竖直线段，使像成为水平直线。图只显示选取的坐标范围。}
	\label{g1:fig:constant-rank-coordinates}
\end{figure}

\begin{corollary}[流形上的常秩标准形]
	\label{g1:cor:manifold-constant-rank}
	设 \(F:M^m\to N^n\) 是无边界光滑流形间的光滑映射，取 \(p\in M\)，且 \(\operatorname{rank}\mathrm{d}F_q=r\) 在 \(p\) 的某个邻域内恒定。则存在 \(p\) 与 \(F(p)\) 附近的光滑坐标，分别以这两点为原点，使 \(F\) 的坐标表达为 \((u^1,\ldots,u^m)\mapsto(u^1,\ldots,u^r,0,\ldots,0)\)，其中末尾有 \(n-r\) 个零，即\eqref{g1:eq:constant-rank-normal-form}。这里不要求坐标变化保持原先给定的度量、长度或角度。
\end{corollary}
\begin{proof}
	先选任意两张坐标图。由\cref{g1:def:manifold-differential,g1:prop:differential-chain-rule}，坐标表达的导数在两侧与坐标微分的可逆线性映射复合，所以它的秩恰为 \(\mathrm{d}F\) 的秩。Euclid 常秩定理给出源、目标坐标域中的进一步微分同胚；将它们与原坐标图复合，仍是光滑坐标图，就得到结论。构造只使用光滑可逆性，未施加任何度量限制。
\end{proof}

\subsection{秩不恒定时的反例}
\label{g1:subsec:1103:04}

常秩假设控制的是一个邻域，不是单点。单点的导数只记录一阶信息，不能独自判定附近原像的形状。不过微分矩阵的连续性确实提供一个有用的单侧结论。

\begin{proposition}
	\label{g1:prop:rank-lower-bound}
	对光滑映射 \(F:M^m\to N^n\)，若 \(\operatorname{rank}\mathrm{d}F_p\geq r\)，则在 \(p\) 的某个邻域中处处有 \(\operatorname{rank}\mathrm{d}F_q\geq r\)。特别地，若在 \(p\) 处秩达到最大可能值 \(\min(m,n)\)，则它在附近自动恒定。
\end{proposition}
\begin{proof}
	在坐标中，秩至少为 \(r>0\) 意味着某个 \(r\) 阶子式非零。该子式是导数矩阵元的多项式，随点连续，所以在足够小邻域中仍非零，秩遂仍至少为 \(r\)。当 \(r=0\) 时结论本身没有限制。最后秩始终不超过 \(\min(m,n)\)，因而达到这一上界时，邻域中的上下界相同。
\end{proof}

\begin{example}
	\label{g1:exm:cubic-rank-change}
	映射 \(f(t)=t^3\) 在原点的导数为零，在其余点的秩为 \(1\)。它仍是 \(\mathbb R\) 到自身的连续双射，逆函数为实立方根；但
	\[
		\frac{f^{-1}(h)-f^{-1}(0)}h=|h|^{-2/3}\longrightarrow\infty
		\qquad(h\to0),
	\]
	所以逆函数在零处不可微。导数可逆是局部光滑可逆的必要条件：若光滑逆存在，对两个复合恒等式求微分，便得到互逆的导数矩阵。它却不是连续可逆的必要条件，立方函数正好区分了这两件事。
\end{example}

在原点把 \(f\) 的秩记为零，并不能将立方函数在邻域中改成常值映射。任意可逆坐标变化都保持每一点的微分秩，而这个例子在任意原点邻域内都有秩为 \(1\) 的点。

\begin{example}
	\label{g1:exm:rank-changing-product}
	考虑 \(F(x,y)=(x,xy)\)，则
	\[
		\mathrm{D}F(x,y)=\begin{pmatrix}1&0\\y&x\end{pmatrix}.
	\]
	它在 \(x\ne0\) 时秩为 \(2\)，在整条 \(x=0\) 的直线上秩为 \(1\)。同值集合也发生明显变化：
	\[
		F^{-1}(a,b)=
		\begin{cases}
			\{(a,b/a)\},&a\ne0,\\
			\{(0,y):y\in\mathbb R\},&(a,b)=(0,0),\\
			\varnothing,&a=0,\ b\ne0.
		\end{cases}
	\]
	尤其在原点附近，整条短竖线被压到一点，旁边的点却各自局部可逆。若源邻域取为 \(|x|<\varepsilon,\ |y|<\delta\)，其像为
	\[
		\{(a,b):0<|a|<\varepsilon,\ |b|<\delta|a|\}\cup\{(0,0)\}.
	\]
	这个像在原点附近既不是开邻域，也不能由一段光滑坐标直线描述。常秩定理证明中的余分量在这里是 \(G(u,v)=uv\)，其 \(\mathrm{D}_vG=u\) 只在 \(u=0\) 上消失，不能因此删掉全部 \(v\) 依赖。
\end{example}

由此可见，若希望一个方程的解集或映射的原像具有稳定的局部维数，必须检查相应邻域中的正则条件。下一节将把满秩情形分别解释为浸入与淹没，并将这些坐标形式用于子流形；本节的结果只处理无边界邻域，半空间边界还需要相应的边界条件。


% !TeX root = ../../Book1.tex
% 纲要标题：### 11.4 浸入、淹没与嵌入
% 纲要位置：几何学计划纲要.md，第 558 行。
% 正文状态：初稿；尚未正式审核。

% 纲要内容开始
% 
% * 浸入和淹没的标准形
% * 正则水平集定理
% * 嵌入、初始子流形结构及限制映射
% * 单射浸入未必嵌入
% * 带边界版本中切空间与边界条件
% 
% ---
% 
% 纲要内容结束

\section{浸入、淹没与嵌入}
\label{g1:sec:1104}

第八章用秩为二的参数映射描述曲面，现在可以把同一个问题放到任意维数的流形上。一张参数片是否丢失切向方向，由微分是否单射判断；一个方程组能否在附近独立控制所有输出，由微分是否满射判断。但即使每张参数片都很规整，远处的参数仍可能回到眼前，这就需要另加拓扑条件。常秩定理先给出局部图像，嵌入的定义再保证这些局部图像与整个像的拓扑相符。本节前四个小节先讨论无边界流形，最后专门说明边界处哪些条件需要补充。

\subsection{浸入和淹没的标准形}
\label{g1:subsec:1104:01}

\begin{definition}[浸入与淹没]\label{g1:def:immersion-submersion}
	设 \(F:M^m\to N^n\) 是光滑映射。若 \(\mathrm{d}F_p:\mathrm{T}_pM\to \mathrm{T}_{F(p)}N\) 单射，称 \(F\) 在 \(p\) 处为浸入；若它满射，称 \(F\) 在 \(p\) 处为\term[yanmo]{淹没}[submersion]。处处满足相应条件时，分别称 \(F\) 为浸入或淹没。
\end{definition}

这里的浸入推广了\cref{g1:def:surface-immersion}中的曲面参数：其定义域维数为二，所以满列秩正是两个参数方向不被压在一起。一般浸入必须有 \(m\le n\)，一般淹没必须有 \(m\ge n\)。淹没也不要求映射在集合意义下满射，例如 \(t\mapsto e^t\) 是 \(\mathbb R\to\mathbb R\) 的淹没，而像只是正半轴。

\begin{theorem}[浸入与淹没的局部标准形]\label{g1:thm:immersion-submersion-normal-form}
	设 \(M^m,N^n\) 是无边界光滑流形，\(F:M\to N\) 光滑，\(p\in M\)。
	\begin{enumerate}
		\item 若 \(\mathrm{d}F_p\) 单射，则存在分别以 \(p,F(p)\) 为原点的源、目标坐标，使 \(F\) 在源的一个邻域中写成
		\[
			(u^1,\ldots,u^m)\longmapsto(u^1,\ldots,u^m,0,\ldots,0).
		\]
		\item 若 \(\mathrm{d}F_p\) 满射，则可选这样的坐标，使 \(F\) 写成
		\[
			(u^1,\ldots,u^m)\longmapsto(u^1,\ldots,u^n).
		\]
	\end{enumerate}
\end{theorem}
\begin{proof}
	单射情形，坐标 Jacobian 在 \(p\) 处有一个非零的 \(m\) 阶子式，连续性使它在某邻域内仍非零。秩又不能超过源维数 \(m\)，故该邻域的秩恒为 \(m\)。满射情形同理：一个非零的 \(n\) 阶子式保证邻域内的秩恒为 \(n\)。现在分别对秩 \(m\) 和秩 \(n\) 使用\cref{g1:cor:manifold-constant-rank}，就得到两种标准形。这里先证明邻域内秩恒定，才能调用常秩结论；仅知道一点的非极大秩，通常不能这样做。
\end{proof}

从标准形可以读出两项常用结论。浸入限制到足够小的源邻域后，是到其像的同胚：在坐标中只需投影到前 \(m\) 个坐标就得到连续逆映射。这说的是一个指定的参数分支，并没有断言其他参数分支不会进入同一个目标邻域。

\ProseParagraphBreak

无边界流形之间的淹没是开映射。事实上，把任一源开集中的点放进包含在该开集内的坐标小盒，投影标准形把这个小盒送到目标中的开盒，所有这些像的并就是原开集的像。此外，若 \(q=F(p)\)，在投影标准形中取被忘掉的坐标恒为零，就得到一个通过 \(p\) 的光滑局部截面 \(s\)，满足 \(F\circ s=\operatorname{id}\)。因此淹没允许在目标点附近光滑选择一个原像，但未必能把这些局部选择拼成全局截面。

\subsection{正则水平集定理}
\label{g1:subsec:1104:02}

方程 \(F(p)=y\) 应当减少多少自由度，取决于它实际提供了多少个独立约束。若目标维数为 \(n\)，微分满射正是所有 \(n\) 个约束在一阶上独立。淹没标准形将这个判断直接变成坐标中固定前 \(n\) 个数。

\begin{definition}[正则值]\label{g1:def:regular-value}
	设 \(F:M\to N\) 光滑。若 \(\mathrm{d}F_p\) 满射，称 \(p\) 为 \(F\) 的\term[zhengzedian]{正则点}[regular point]；否则称为\term[linjiedian]{临界点}[critical point]。若 \(y\in N\) 的每个原像都是正则点，称 \(y\) 为\term[zhengzezhi]{正则值}[regular value]；否则称为\term[linjiezhi]{临界值}[critical value]。没有原像的点也按此定义算作正则值。
\end{definition}

\begin{theorem}[正则水平集]\label{g1:thm:regular-level-set}
	设 \(F:M^m\to N^n\) 是无边界光滑流形之间的光滑映射，\(y\in N\) 是正则值。若 \(S=F^{-1}(y)\) 非空，则 \(S\) 在子空间拓扑下有自然的 \((m-n)\) 维光滑流形结构，包含映射 \(i:S\hookrightarrow M\) 是浸入，且
	\begin{equation}\label{g1:eq:regular-level-tangent}
		\mathrm{d}i_p(\mathrm{T}_pS)=\ker\mathrm{d}F_p\qquad(p\in S).
	\end{equation}
	通过单射 \(\mathrm{d}i_p\) 把 \(\mathrm{T}_pS\) 视为 \(\mathrm{T}_pM\) 的子空间时，可简写为 \(\mathrm{T}_pS=\ker\mathrm{d}F_p\)。
\end{theorem}
\begin{proof}
	对每个 \(p\in S\) 使用淹没标准形，将 \(y\) 的目标坐标取为零。在所得源坐标 \((u^1,\ldots,u^m)\) 中，水平集恰为
	\[
		u^1=\cdots=u^n=0.
	\]
	因此后 \(m-n\) 个坐标给出 \(S\) 的局部坐标，而且这些坐标与 \(S\) 的子空间拓扑一致。两张这样的坐标之间，转换映射是在环境坐标转换中先代入前 \(n\) 个坐标为零，再取所需坐标分量，故光滑；逆转换同理。这构成光滑图册。Hausdorff 性和第二可数性从环境的子空间继承，所以确实得到流形。

	在这些坐标中，包含映射的微分把后 \(m-n\) 个坐标方向单射送入 \(\mathrm{T}_pM\)。另一方面，\(\mathrm{d}F_p\) 是前 \(n\) 个方向的投影，其核恰由后 \(m-n\) 个方向张成。于是得到\eqref{g1:eq:regular-level-tangent}。坐标改变通过\cref{g1:prop:differential-chain-rule}的微分链式法则相容，所以这个子空间描述不依赖所选坐标。
\end{proof}

下一小节将通过包含映射的性质进一步说明这个自然结构为何唯一。这里并非随意给同一集合指定一种流形拓扑。只需在水平集各点满足正则值的条件，不要求 \(F\) 在整个 \(M\) 上都是淹没。

\begin{example}\label{g1:exm:regular-sphere-and-singular-level}
	设 \(F:\mathbb R^{k+1}\to\mathbb R\)，\(F(x)=\|x\|^2\)，其中 \(k\ge1\)。因为
	\[
		\mathrm{d}F_x(v)=2\langle x,v\rangle,
	\]
	每个 \(c>0\) 都是正则值，球面 \(F^{-1}(c)\) 的维数为 \(k\)，并且 \(\mathrm{T}_xF^{-1}(c)=x^\perp\)。负数也为正则值，但水平集为空。零为临界值，其水平集只有一个点；该集合仍可以是零维流形，却不具有正则水平集定理所给的预期维数 \(k\)。

	另取 \(G(x,y)=xy\)。它的零水平集是两条坐标轴的并，在原点形成四条分支，不是原点附近的一维流形。两个例子说明：临界水平集有时仍规整，有时出现交叉，仅由“不是正则值”不能断言所有水平集具有同一种奇异性。
\end{example}

\subsection{嵌入、初始子流形结构及限制映射}
\label{g1:subsec:1104:03}

第四章的拓扑嵌入要求映射与其像之间互为连续逆映射，像必须使用环境的子空间拓扑。加入这个条件以后，浸入的局部参数片才会与整个像的局部邻域一致。

\begin{definition}[光滑嵌入]\label{g1:def:smooth-embedding}
	无边界光滑流形之间的光滑映射 \(f:M\to N\)，若既是浸入，又是 \(M\) 到 \(f(M)\) 的同胚，其中 \(f(M)\) 取 \(N\) 的子空间拓扑，则称为\term[guanghuaqianru]{光滑嵌入}[smooth embedding]，简称嵌入。
\end{definition}

\begin{definition}[嵌入子流形]\label{g1:def:embedded-submanifold}
	设 \(N^n\) 为无边界光滑流形。子集 \(S\subset N\) 若对每个 \(p\in S\) 都有环境坐标 \((V,y)\)，使
	\[
		y(S\cap V)=y(V)\cap\bigl(\mathbb R^k\times\{0\}^{n-k}\bigr),
	\]
	则称为 \(N\) 的 \(k\) 维\term[qianruziliuxing]{嵌入子流形}[embedded submanifold]。给 \(S\) 子空间拓扑，并用这些切片的前 \(k\) 个坐标作为图册。
\end{definition}

坐标切片之间的转换正如正则水平集证明中那样光滑，所以这个定义确实给出 \(S\) 的光滑结构。下述等价关系说明，这个集合式定义与前面的映射式定义对应。

\begin{proposition}[嵌入与坐标切片]\label{g1:prop:embedding-slice-characterization}
	设 \(M^k,N^n\) 是无边界光滑流形。若 \(f:M\to N\) 为光滑嵌入，则其像 \(S=f(M)\) 为嵌入子流形，且 \(f:M\to S\) 是微分同胚。反之，每个嵌入子流形的包含映射都是光滑嵌入。
\end{proposition}
\begin{proof}
	取 \(p\in M\)，浸入标准形给出源邻域 \(U\) 及 \(f(p)\) 附近的坐标，使 \(f|_U\) 的像是一块坐标切片。必要时缩小坐标盒，可以使这块切片完整覆盖目标小盒中的前 \(k\) 个坐标。关键是排除 \(M\setminus U\) 的点也映入眼前的小盒：因为 \(f\) 是到子空间 \(S\) 的同胚，\(f(U)\) 在 \(S\) 中开，故存在环境开集 \(W\subset N\)，满足 \(S\cap W=f(U)\)。把目标坐标盒进一步缩到 \(W\) 内，盒内的整个 \(S\) 就恰为该切片，而不只是其中一个参数分支。这样的盒覆盖 \(S\)，给出嵌入子流形结构。在对应坐标中，\(f:M\to S\) 是恒等坐标映射，故连同逆映射都光滑。

	反过来，包含映射在适配坐标中为 \(u\mapsto(u,0)\)，所以是浸入。\(S\) 已取环境的子空间拓扑，包含映射自然是到其像的同胚，故为光滑嵌入。
\end{proof}

\begin{proposition}[嵌入子流形的初始性质]\label{g1:prop:embedded-initial-property}
	设 \(S\subset N\) 为无边界光滑流形 \(N\) 的嵌入子流形，\(i:S\hookrightarrow N\) 为包含映射。对任意光滑流形 \(P\)（允许有边界）及任意集合映射 \(h:P\to S\)，有
	\[
		h:P\to S\text{ 光滑}\quad\Longleftrightarrow\quad i\circ h:P\to N\text{ 光滑}.
	\]
	特别地，光滑映射 \(H:P\to N\) 若像包含在 \(S\) 中，就唯一地给出光滑映射 \(P\to S\)。在固定子集 \(S\subset N\) 上，使其成为嵌入子流形的光滑结构唯一。
\end{proposition}
\begin{proof}
	正向由光滑映射的复合得到。反向先处理一个容易遗漏的拓扑问题：若 \(i\circ h\) 光滑，则它连续；由于 \(S\) 取子空间拓扑，\(h\) 也连续。于是对 \(h(p)\) 附近的适配图 \(V\)，集合 \(h^{-1}(S\cap V)\) 在 \(P\) 中开。在这片上，\(h\) 的坐标表示就是 \(i\circ h\) 的坐标表示的前 \(\dim S\) 个分量，所以光滑。即使 \(p\) 在 \(\partial P\)，这些分量仍由环境中的光滑延拓限制得到，论证不变。局部结论合起来给出 \(h\) 光滑。

	像包含在 \(S\) 的映射只有一种集合意义下的限制值域方式，上述等价立即证明其光滑性。若同一个子集带有两种嵌入子流形结构，两种包含映射都光滑。对两种结构之间的恒等集合映射分别使用已证等价，得到恒等映射在两个方向都光滑，故两种极大光滑图册相同。
\end{proof}

这里的“初始”说的是通过包含映射检测光滑性的性质，不是另给 \(S\) 加一种比子空间拓扑更细的拓扑。例如光滑函数 \(f:N\to\mathbb R\) 限制到 \(S\) 后光滑；若 \(f(S)\subset T\)，其中 \(T\) 也是嵌入子流形，则 \(f|_S:S\to T\) 光滑。正则水平集由环境坐标切片构造，因而属于这一情形。

\ProseParagraphBreak

包含映射的微分通常被省略在记号里：对嵌入子流形写 \(\mathrm{T}_pS\subset \mathrm{T}_pN\)，实际是用单射 \(\mathrm{d}i_p\) 作识别。这不是把两个定义不同的导子集合逐项当作相同集合，而是利用包含映射，在前者与后者的一个子空间之间建立自然的线性同构。

\subsection{单射浸入未必嵌入}
\label{g1:subsec:1104:04}

单射解决了“一个像点有几个原像”，却没有解决“附近像点的原像是否也在附近”。下面的平面曲线已经能把两者区别开，不需要依靠高维或稠密轨道的例子。

\begin{example}\label{g1:exm:injective-immersion-not-embedding}
	令
	\[
		\gamma:(-\pi,\pi)\longrightarrow\mathbb R^2,
		\qquad\gamma(t)=(\sin t,\sin2t).
	\]
	它是单射浸入，但不是嵌入。先检查浸入：\(\gamma'(t)=(\cos t,2\cos2t)\)，若第一分量为零，则 \(t=\pm\pi/2\)，第二分量为 \(-2\ne0\)，故导数从不为零。

	再检查单射。\(\sin t\) 的正负区分正、负参数，且在本区间内只在 \(t=0\) 为零。若 \(s,t\in(0,\pi)\) 且 \(\sin s=\sin t\)，则 \(s=t\) 或 \(s=\pi-t\)；后一情形的 \(\sin2s=-\sin2t\)，两者相等只能有 \(t=\pi/2\)，此时仍为 \(s=t\)。负参数的情形用 \(\gamma(-t)=-\gamma(t)\) 化为正参数，因而 \(\gamma\) 单射。

	然而取 \(t_j\to\pi\)，便有 \(\gamma(t_j)\to(0,0)=\gamma(0)\)，而 \(t_j\) 不趋于零。因此从像的子空间拓扑返回参数的逆映射不连续。\cref{g1:fig:injective-immersion-topology}中，蓝色短弧来自零附近的参数，另一方向逼近交点的曲线来自接近参数区间端点的部分。浸入标准形只能控制前一支。

	这个像甚至是闭集：把同一公式延拓到紧区间 \([-\pi,\pi]\)，两个端点都映到已有的 \(\gamma(0)\)，所以延拓前后的像相同，是紧集。可见“单射浸入且像闭”仍不能替代嵌入条件。
\end{example}

\input{Figures/Chapter11-InjectiveImmersion.tex}

还可以在这个例子中具体看出，嵌入子流形的初始性质不能自动移植给任意单射浸入像。给集合 \(S=\gamma((-\pi,\pi))\) 从参数区间搬来流形结构，使 \(\gamma\) 成为微分同胚；此时包含 \(S\to\mathbb R^2\) 是浸入。对充分小的 \(\varepsilon>0\)，考虑值在 \(S\) 中的曲线
\[
	h(s)=(-\sin s,\sin2s),\qquad |s|<\varepsilon.
\]
作为到 \(\mathbb R^2\) 的映射，它光滑；但它在 \(S\) 的参数坐标为
\[
	\gamma^{-1}(h(s))=
	\begin{cases}
		\pi+s,&s<0,\\
		0,&s=0,\\
		-\pi+s,&s>0,
	\end{cases}
\]
在零处不连续，故 \(h:(-\varepsilon,\varepsilon)\to S\) 不是光滑映射。问题恰是这份流形结构的拓扑不同于 \(S\) 的子空间拓扑。

\begin{definition}[固有映射]\label{g1:def:proper-map-manifolds}
	光滑流形之间的连续映射 \(f:M\to N\) 若对每个紧集 \(K\subset N\)，原像 \(f^{-1}(K)\) 都紧，则称为\term[guyouyingshe]{固有映射}[proper map]。
\end{definition}

固有性控制的是“落在一个固定紧目标区域内的所有参数”。上一例中整个像紧，而其原像是非紧区间 \((-\pi,\pi)\)，所以 \(\gamma\) 不固有。这一条件能排除远处参数悄悄回到紧区域的现象。

\begin{proposition}[固有单射浸入]\label{g1:prop:proper-injective-immersion}
	设 \(M,N\) 为无边界光滑流形，\(f:M\to N\) 为单射浸入。若 \(f\) 固有，则 \(f\) 是嵌入，且像闭于 \(N\)。特别地，若 \(M\) 紧，则每个单射浸入 \(M\to N\) 都是闭像嵌入。反之，闭像嵌入是固有映射。
\end{proposition}
\begin{proof}
	先证明一个连续固有映射 \(f:M\to N\) 是闭映射。取闭集 \(A\subset M\) 及 \(y\in\overline{f(A)}\)。在 \(N\) 的坐标图内取一个紧邻域 \(K\ni y\)，使 \(y\in\operatorname{int}K\)。固有性使 \(f^{-1}(K)\) 紧，故其闭子集 \(A\cap f^{-1}(K)\) 紧。它的连续像
	\[
		f(A\cap f^{-1}(K))=f(A)\cap K
	\]
	也紧，并由\cref{g1:prop:compact-hausdorff-closed}闭于 \(N\)。回忆这里闭性的理由：对紧集外一点，用 Hausdorff 性把它与紧集的每一点用开邻域分开，再以有限子覆盖替代紧集一侧的全部邻域，取外点一侧有限个邻域的交，便得到与紧集不交的开邻域。因为 \(y\) 位于 \(K\) 的内部且是 \(f(A)\) 的闭包点，它也是 \(f(A)\cap K\) 的闭包点，因而属于该闭集。于是 \(f(A)\) 闭。

	现在 \(f\) 连续且单射，并将闭集送到闭集，所以 \(f:M\to f(M)\) 的逆映射连续。它又是浸入，故为嵌入；取 \(A=M\) 还得到像闭。若 \(M\) 紧，每个目标紧集在 Hausdorff 空间 \(N\) 中闭，其原像为 \(M\) 的闭子集，因而紧，所以 \(f\) 自动固有。此时也可直接使用\cref{g1:thm:compact-bijection}判断到像的逆连续。

	反之，若 \(f\) 是闭像嵌入，则 \(K\cap f(M)\) 对每个紧 \(K\subset N\) 都是紧集；由同胚 \(M\cong f(M)\)，其原像 \(f^{-1}(K)\) 紧，故 \(f\) 固有。
\end{proof}

并非所有嵌入都固有。例如开区间 \((0,1)\hookrightarrow\mathbb R\) 是嵌入，但紧集 \([0,1]\) 的原像不紧。固有单射浸入给出的是一种保证闭像嵌入的充分判据，并没有把“嵌入”的含义改成“闭子集”。

\subsection{带边界版本中切空间与边界条件}
\label{g1:subsec:1104:05}

对带边界流形，浸入、淹没仍按完整微分的单射、满射定义，嵌入仍要求同时是浸入和到子空间像的同胚。边界坐标采用 \(\mathbb H^m=\{x^m\ge0\}\)。\cref{g1:prop:derivation-coordinate-basis}及\cref{g1:prop:boundary-velocities}说明，即使 \(p\in\partial M\)，完整切空间 \(\mathrm{T}_pM\) 仍是 \(m\) 维向量空间；向内半空间只是其中能够由向内单侧曲线实现的速度集合，不能代替整个切空间。若以 \(\rho=x^m\) 为局部边界定义函数，则
\[
	\mathrm{T}_p(\partial M)=\ker\mathrm{d}\rho_p\subset \mathrm{T}_pM.
\]
边界上各种满秩条件因此仍按完整切空间计算，但由此得到的 Euclid 局部变换未必保留半空间。

\begin{proposition}[边界处的局部逆]\label{g1:prop:boundary-local-inverse}
	设 \(M^m,N^m\) 为带边界光滑流形，\(m\ge1\)，\(F:M\to N\) 光滑。若 \(p\in\partial M\)、\(F(p)\in\partial N\)，在 \(p\) 的某邻域内有 \(F(\partial M)\subset\partial N\)，并且 \(\mathrm{d}F_p\) 可逆，则 \(F\) 限制到 \(p\) 的某个相对开邻域后，是到 \(F(p)\) 的一个相对开邻域的微分同胚。
\end{proposition}
\begin{proof}
	取两边的半空间坐标，把 \(p,F(p)\) 都置于零。按边界光滑映射的定义，坐标表示 \(G\) 可光滑延拓到 \(\mathbb R^m\) 的一个开邻域。保边界条件给出 \(G^m(x',0)=0\)，所以由一元微积分基本定理，
	\[
		G^m(x',x^m)=x^m a(x',x^m),\qquad
		a(x',x^m)=\int_0^1\partial_mG^m(x',t x^m)\,\mathrm{d}t.
	\]
	\(\mathrm{d}G_0\) 的最后一行在前 \(m-1\) 列为零；可逆性因而使 \(a(0)\ne0\)。又因为 \(G\) 把 \(x^m\ge0\) 的点送到目标半空间，一侧导数给出 \(a(0)\ge0\)，故 \(a(0)>0\)。缩小邻域使 \(a>0\)，则 \(G^m\) 与 \(x^m\) 同号。对 Euclid 延拓使用\cref{g1:thm:inverse-function}，得到开邻域间的光滑逆；同号性质说明该逆也保留半空间。把这对互逆映射限制到半空间，便得到所需边界局部逆。
\end{proof}

保边界条件不能只检查在一个点成立。例如
\[
	G:\mathbb H^2\to\mathbb H^2,\qquad G(x,y)=(x,y+x^2)
\]
在原点的微分是恒等映射，也把原点送到原点，但它把附近的边界点送到抛物线 \(v=u^2\)。其像在 \(v\ge u^2\) 中，不包含原点在目标半平面中的任何相对开邻域，所以不是边界局部微分同胚。Euclid 延拓的逆 \((u,v)\mapsto(u,v-u^2)\) 光滑，却未必取值于源半平面。

\ProseParagraphBreak

对水平集，正确的附加条件要求约束在边界切向方向中仍然独立。这样，约束不会把唯一的向内方向全部吃掉。若带边界子流形 \(S\subset M\) 满足 \(\partial S=S\cap\partial M\)，且在这些点满足
\[
	\mathrm{T}_pS+\mathrm{T}_p(\partial M)=\mathrm{T}_pM,
\]
本书用\term[neatziliuxing]{与边界相容的横截嵌入子流形}[neat submanifold]描述同时满足这两个条件的子流形。这里包含映射要求为嵌入；横截等式意味着 \(\mathrm{T}_pS\) 中还保留至少一个不沿环境边界的方向。

\begin{theorem}[带边界的正则水平集]\label{g1:thm:boundary-regular-level-set}
	设 \(M^m\) 为带边界光滑流形，\(N^n\) 为无边界光滑流形，\(F:M\to N\) 光滑。若 \(y\in N\) 同时为 \(F\) 及 \(F|_{\partial M}\) 的正则值，则非空水平集 \(S=F^{-1}(y)\) 自然成为 \((m-n)\) 维带边界嵌入子流形，且
	\[
		\partial S=S\cap\partial M,\qquad \mathrm{T}_pS=\ker\mathrm{d}F_p.
	\]
	它是与边界相容的横截嵌入子流形；在边界点还满足
	\[
		\mathrm{T}_p(\partial S)=\ker\mathrm{d}F_p\cap \mathrm{T}_p(\partial M).
	\]
\end{theorem}
\begin{proof}
	内部点使用\cref{g1:thm:regular-level-set}的局部证明。取 \(p\in S\cap\partial M\)，用半空间坐标将 \(p\) 置于零，用目标坐标将 \(y\) 置于零，记映射的坐标分量为 \(f^1,\ldots,f^n\)。\(F|_{\partial M}\) 在 \(p\) 处满秩，意味着 Jacobian 的前 \(m-1\) 列中有一个非零的 \(n\) 阶子式。重排前 \(m-1\) 个坐标，可令该子式来自前 \(n\) 列。特别地，若水平集碰到边界，必有 \(m-n\ge1\)。

	对 \(f\) 的 Euclid 光滑延拓构造
	\[
		\Phi(x)=\bigl(f^1(x),\ldots,f^n(x),x^{n+1},\ldots,x^{m-1},x^m\bigr).
	\]
	\(\mathrm{d}\Phi_0\) 可逆，故由\cref{g1:thm:inverse-function}，它在足够小的 Euclid 开邻域间为微分同胚。与任意使用常秩定理不同，这里特意把最后一个坐标原样保留：\(\Phi^m(x)=x^m\)。因此 \(\Phi\) 及其逆都保留半空间和边界。以 \(u=\Phi(x)\) 为新的边界坐标，水平集成为
	\begin{equation}\label{g1:eq:boundary-level-slice}
		u^1=\cdots=u^n=0,\qquad u^m\ge0.
	\end{equation}
	后 \(m-n\) 个坐标给出水平集的半空间图，其边界恰是 \(u^m=0\)，也就是环境边界与水平集的交。图册转换由环境光滑转换的限制得到，子空间拓扑、Hausdorff 性及第二可数性同无边界情形，因此得到带边界嵌入子流形。

	在这些坐标中，\(\mathrm{T}_pS\) 由后 \(m-n\) 个完整坐标方向张成，等于 \(\ker\mathrm{d}F_p\)；其边界切空间再去掉最后一个方向，正是 \(\ker\mathrm{d}F_p\cap \mathrm{T}_p(\partial M)\)。最后一个方向 \(\partial/\partial u^m\) 属于 \(\mathrm{T}_pS\) 却不属于 \(\mathrm{T}_p(\partial M)\)，而后者余维为一，故两者之和为整个 \(\mathrm{T}_pM\)。这就同时验证了边界交等式与横截条件。
\end{proof}

上述半空间切片还说明，与边界相容的横截水平集仍具有初始性质：一个映射若值在 \(S\) 中，其是否光滑可以通过组合包含 \(S\hookrightarrow M\) 来检验。证明仍先用子空间拓扑得到连续性，再取\eqref{g1:eq:boundary-level-slice}中后 \(m-n\) 个坐标；这些坐标有环境光滑延拓，因此同样适用于带边界的源。

\ProseParagraphBreak

只检查完整切空间上的 \(\mathrm{d}F\) 满射还不够。例如 \(F:\mathbb H^m\to\mathbb R\)，\(F(x)=x^m\)，处处满秩，但 \(F^{-1}(0)=\partial M\) 是一张无边界的 \((m-1)\) 维流形；其自身边界为空，而与 \(\partial M\) 的交却是整个水平集，故不是 与边界相容的横截水平集。更明显地，\(F(x,y)=y+x^2\) 在 \(\mathbb H^2\) 上也处处满秩，但零水平集只有原点，连预期的一维局部形状都没有。两例中 \(F|_{\partial M}\) 都在相关点失去满秩性，正是附加条件所要排除的现象。


\begin{chaptersummary}
	光滑结构由两两相容的坐标图确定，坐标图的极大扩张消除了任意多列或少列相容图的区别。球面和射影空间可以直接用图册构造；同一个底层拓扑却不能独自决定哪些函数光滑。边界图使用半空间中的局部延拓，延拓函数可以不同，但边界点处各阶导数相同。

	\ProseParagraphBreak

	切向量是作用于局部函数芽的一阶导子，坐标导数给出它的基。内点处每个切向量都能由双侧曲线实现；边界处双侧速度只给出边界切空间，单侧速度给出内向锥，完整切空间仍为 \(n\) 维。切丛把不同点的切空间通过微分换坐标组织在一起；光滑映射的微分向前送切向量，其对偶反向拉回余切向量。

	\ProseParagraphBreak

	逆函数定理以可逆线性近似控制非线性余项，通过压缩迭代产生局部逆。隐函数定理把待求变量与参数一起放入这一构造；常秩定理再利用秩在邻域中不变，消去剩余坐标的变化，得到局部投影或坐标包含。满秩约束因此给出光滑水平集，切空间正是约束微分的核。

	\ProseParagraphBreak

	浸入只控制局部切方向，嵌入还要求源拓扑与像的子空间拓扑相符；单射和闭像各自都不足以代替这一条件，固有性则提供实用的充分条件。边界问题还必须辨认边界是否对应、约束是否沿边界仍满秩。Euclid 延拓的局部可逆，不会自动保证它的逆留在所需的半空间一侧。
\end{chaptersummary}

\BookExercises


\begin{exercise}
	\label{g1:ex:stereographic-smooth-charts}
	设 \(n\geq1\)，\(N=(0,\ldots,0,1)\)、\(S=-N\) 是 \(S^n\) 的两极。对 \(x=(x',x_{n+1})\) 定义
	\[
		\sigma_N(x)=\frac{x'}{1-x_{n+1}},\qquad
		\sigma_S(x)=\frac{x'}{1+x_{n+1}}.
	\]
	求这两个映射的逆，证明它们给出与\cref{g1:exm:sphere-smooth-atlas}相同的光滑结构，并求换坐标 \(\sigma_S\circ\sigma_N^{-1}\)。说明换坐标在零点没有定义的几何原因。
\end{exercise}

\begin{solution}
	解 \(u=x'/(1-x_{n+1})\) 与 \(|x'|^2+x_{n+1}^2=1\)，得到
	\[
		\sigma_N^{-1}(u)=\left(\frac{2u}{1+|u|^2},\frac{|u|^2-1}{1+|u|^2}\right),\qquad
		\sigma_S^{-1}(v)=\left(\frac{2v}{1+|v|^2},\frac{1-|v|^2}{1+|v|^2}\right).
	\]
	这些式子对任意 \(u,v\in\mathbb R^n\) 有意义，代回即得到恒等式，且与原映射均连续。因此它们分别是 \(S^n\setminus\{N\}\)、\(S^n\setminus\{S\}\) 到 \(\mathbb R^n\) 的坐标图。与正文半球图转换时，一向由上述有理式后删除一个分量得到，另一向由半球图的光滑平方根表达代入 \(\sigma_N\) 或 \(\sigma_S\) 得到，分母在相应交叠处非零，故两种图册相容。

	在 \(u\neq0\) 时，把逆式代入可得
	\[
		(\sigma_S\circ\sigma_N^{-1})(u)=\frac{u}{|u|^2}.
	\]
	该映射的逆仍是自身。\(u=0\) 在北极投影图中代表南极，而南极不属于另一张图的图域，故这里的未定义来自图域没有交叠，并非球面的奇点。
\end{solution}

\begin{exercise}
	\label{g1:ex:projective-line-circle}
	证明映射
	\[
		F:\mathbb{RP}^1\longrightarrow S^1,\qquad
		F([x:y])=\left(\frac{x^2-y^2}{x^2+y^2},\frac{2xy}{x^2+y^2}\right)
	\]
	是微分同胚。要求直接给出局部逆的坐标公式，不只利用紧空间到 Hausdorff 空间的连续双射性质。
\end{exercise}

\begin{solution}
	两个分式在同时缩放 \((x,y)\) 后不变，且其平方和为 \(1\)，所以 \(F\) 良定义。在图 \(x\neq0\) 中写 \(t=y/x\)，有
	\[
		F([1:t])=\left(\frac{1-t^2}{1+t^2},\frac{2t}{1+t^2}\right).
	\]
	它的像为圆周除去 \((-1,0)\)，逆坐标为 \(t=b/(1+a)\)，其中 \((a,b)\in S^1\)、\(a\neq-1\)。在图 \(y\neq0\) 中写 \(s=x/y\)，得到
	\[
		F([s:1])=\left(\frac{s^2-1}{s^2+1},\frac{2s}{s^2+1}\right),\qquad s=\frac{b}{1-a}\quad(a\neq1).
	\]
	两组逆式覆盖整个圆周。在同时有定义处，\(st=b^2/((1-a)(1+a))=1\)，所以它们给出同一射影点。两组公式还证明满射与单射；正文的球面图与射影图中各式都光滑，因此 \(F\) 及其逆光滑。
\end{solution}

\begin{exercise}
	\label{g1:ex:transported-line-powers}
	对正奇数 \(r\)，在通常拓扑的 \(\mathbb R\) 上用全局坐标 \(u_r(t)=t^r\) 定义光滑流形 \(M_r\)。对两个正奇数 \(r,s\)，判断集合恒等映射 \(I_{r,s}:M_r\to M_s\) 何时光滑、何时为微分同胚。证明无论 \(r,s\) 如何，\(M_r\) 与 \(M_s\) 总微分同胚。
\end{exercise}

\begin{solution}
	恒等映射的坐标表达为 \(u\mapsto\operatorname{sgn}(u)|u|^{s/r}\)。若 \(s/r\) 是整数，由 \(r,s\) 都为奇数知它是正奇数，表达式就是光滑多项式 \(u^{s/r}\)。若 \(\alpha=s/r\) 不是整数，取 \(k=\lfloor\alpha\rfloor+1\)，在正半轴上第 \(k\) 阶导数等于非零常数乘 \(u^{\alpha-k}\)，当 \(u\downarrow0\) 时无界，因此不可能在零点 \(C^k\)，更不光滑。故 \(I_{r,s}\) 光滑当且仅当 \(r\) 整除 \(s\)；正反两个恒等映射同时光滑当且仅当 \(r=s\)。

	不过映射 \(F=u_s^{-1}\circ u_r:M_r\to M_s\) 始终是微分同胚，因为 \(u_s\circ F\circ u_r^{-1}\) 为恒等映射，逆也一样。具体集合公式为 \(F(t)=\operatorname{sgn}(t)|t|^{r/s}\)。这个公式在通常坐标中可能不光滑，并不妨碍它在指定的源、目标结构之间成为微分同胚。
\end{solution}

\begin{exercise}
	\label{g1:ex:boundary-even-extension}
	设 \(f:[0,\infty)\to\mathbb R\) 按局部延拓意义光滑。证明偶延拓 \(E_f(t)=f(|t|)\) 在 \(\mathbb R\) 上光滑，当且仅当 \(f^{(2j+1)}(0)=0\) 对所有 \(j\geq0\) 成立。再构造一个在负半轴不恒为零的光滑函数 \(h:\mathbb R\to\mathbb R\)，使 \(h\) 在非负半轴恒为零，解释为什么边界处各阶导数的唯一性不意味着延拓函数唯一。
\end{exercise}

\begin{solution}
	如果 \(E_f\) 光滑，由它是偶函数可得 \(E_f^{(k)}(0)=(-1)^kE_f^{(k)}(0)\)，故所有奇阶导数为零；在右侧这些导数就是 \(f\) 的边界导数。

	反过来，假设所有奇阶边界导数为零。在正、负半轴上，\(E_f\) 的第 \(k\) 阶导数分别是 \(f^{(k)}(t)\)、\((-1)^kf^{(k)}(-t)\)。当 \(k\) 为偶数时两边极限相同，为奇数时两边极限均为零，故每一阶候选导数均可连续粘合。为确认它们确为导数，依次用微积分基本定理：若第 \(k+1\) 阶候选函数连续，则其从 \(0\) 到 \(t\) 的积分在两侧都等于第 \(k\) 阶候选函数减去其零点值，因而第 \(k\) 阶候选函数在零点也可微。对任意有限阶重复这一论证，得到 \(E_f\in C^\infty\)。例如 \(f(t)=t^2\) 的偶延拓光滑，\(f(t)=t\) 的偶延拓则是 \(|t|\)，在零点不可微。

	取 \(h(t)=e^{-1/t^2}\) 当 \(t<0\)，\(h(t)=0\) 当 \(t\geq0\)。每次在负半轴求导，都得到一个关于 \(1/t\) 的多项式乘 \(e^{-1/t^2}\)。由 \(u^Ne^{-u^2}\to0\)（\(u\to\infty\)）对每个 \(N\) 成立，这些导数以及它们除以 \(t\) 后的极限均为零，归纳可知零点的各阶导数存在且为零。因此 \(h\) 光滑而不恒为零。对任何一个已取的光滑延拓 \(\widetilde f\)，\(\widetilde f+c h\) 在共同定义的邻域上都是另一个延拓；所有边界导数相同，负半轴上的值却可不同。
\end{solution}
\begin{exercise}
	\label{g1:ex:point-derivation-orders}
	对实直线零点处的光滑函数芽，固定 \(a,b,c\in\mathbb R\)，令
	\[
		L(f)=a f(0)+b f'(0)+c f''(0).
	\]
	求出使 \(L\) 成为零点处导子的全部系数组合，并说明本题为什么体现了导子只记录一阶变化。
\end{exercise}
\begin{solution}
	导子必消去常数，对 \(f=1\) 得 \(a=0\)。再取 \(f(x)=g(x)=x\)，因为 \(f(0)=g(0)=0\)，乘法法则要求 \(L(x^2)=0\)。左边等于 \(2c\)，所以 \(c=0\)。系数 \(b\) 则不受限制：普通求导的乘法法则说明 \(L(f)=b f'(0)\) 对任意 \(b\in\mathbb R\) 都是导子。

	所以全部解是 \((a,b,c)=(0,b,0)\)。零阶求值项被 \(L(1)=0\) 排除，二阶项则被两个在该点为零的函数相乘排除。二阶导数的乘积公式含有额外交叉项 \((fg)''(0)=f''(0)g(0)+2f'(0)g'(0)+f(0)g''(0)\)，因此不满足点导子的乘法法则。
\end{solution}

\begin{exercise}
	\label{g1:ex:sphere-derivation-tangent}
	单位球面 \(S^2\subset\mathbb R^3\) 取标准光滑结构，\(\iota:S^2\to\mathbb R^3\) 为包含映射。用本节的导子或曲线定义证明 \(\mathrm{d}\iota_p\) 单射，且
	\[
		\mathrm{d}\iota_p(\mathrm{T}_pS^2)=\{w\in\mathbb R^3:\langle p,w\rangle=0\}.
	\]
	在这个识别下求对径映射 \(A(p)=-p\) 的微分，以及高度函数 \(h(p)=p_3\) 的微分；找出 \(\mathrm{d}h_p=0\) 的点。
\end{exercise}
\begin{solution}
	记环境坐标限制到球面的函数为 \(x_1,x_2,x_3\)。切向量 \(v\in \mathrm{T}_pS^2\) 的像为 \(w=(v(x_1),v(x_2),v(x_3))\)。对函数恒等式 \(\sum_i x_i^2=1\) 作用导子，得 \(2\sum_i p_iw_i=0\)，故像包含在正交平面内。

	选 \(i\) 使 \(p_i\ne0\)。在 \(p\) 的小邻域内，另外两个环境坐标是球面坐标，遗漏的一项由
	\[
		x_i=\operatorname{sgn}(p_i)\sqrt{1-\sum_{j\ne i}x_j^2}
	\]
	恢复。若 \(\mathrm{d}\iota_p(v)=0\)，则 \(v\) 作用于这两个局部坐标均为零，坐标基公式迫使 \(v=0\)，所以微分单射。任意 \(w\perp p\)、\(w\ne0\) 又可由球面曲线
	\[
		\gamma(t)=p\cos(|w|t)+\frac{w}{|w|}\sin(|w|t)
	\]
	实现：由 \(|p|=1\)、\(p\perp w\) 可知 \(|\gamma(t)|=1\)，且 \(\gamma(0)=p\)、\(\gamma'(0)=w\)。\(w=0\) 用常值曲线实现。因此所述平面恰为全部像。

	对径映射把代表曲线 \(\gamma\) 送到 \(-\gamma\)，所以 \(\mathrm{d}A_p(w)=-w\in \mathrm{T}_{-p}S^2\)。虽然两个切平面在环境中是同一个线性子空间，它们作为切丛的纤维仍有不同基点。高度微分为 \(\mathrm{d}h_p(w)=w_3\)。它在全部 \(w\perp p\) 上为零，当且仅当 \(e_3=(0,0,1)\) 与 \(p\) 平行，故恰在南北两极 \((0,0,\pm1)\) 为零。
\end{solution}

\begin{exercise}
	\label{g1:ex:circle-tangent-trivialization}
	圆周 \(S^1=\{(x,y):x^2+y^2=1\}\) 取标准角坐标。对 \(p=(x,y)\) 及 \(s\in\mathbb R\)，令 \(v_{p,s}\) 为曲线 \(t\mapsto R_{st}p\) 在零时刻的速度，其中 \(R_\theta\) 是平面旋转。证明
	\[
		\Psi:S^1\times\mathbb R\longrightarrow \mathrm{T}S^1,\qquad
		\Psi(p,s)=v_{p,s}
	\]
	是微分同胚。再将 \(S^1\) 记成复数单位圆，求映射 \(F_k(z)=z^k\)、\(k\in\mathbb Z\) 在这一表示下的切映射。
\end{exercise}
\begin{solution}
	在局部角坐标 \(p=(\cos\theta,\sin\theta)\) 中，曲线 \(R_{st}p\) 的角度为 \(\theta+st\)，所以 \(v_{p,s}=s\partial_\theta|_p\)。切空间一维，故每个 \(v\in \mathrm{T}_pS^1\) 都有且仅有一个这样的 \(s\)。因此 \(\Psi\) 双射。角度换成 \(\theta+2\pi m\) 不改变 \(\partial_\theta\) 或系数 \(s\)，所以这些局部描述在交叠处一致。切丛图中，\(\Psi\) 的表达为 \((\theta,s)\mapsto(\theta,s)\)，逆映射也有同样的坐标表达，故二者均光滑。这证明 \(\mathrm{T}S^1\) 可以在全局写成 \(S^1\times\mathbb R\)。

	环境中的相应速度是 \(s(-y,x)\)，这也给出系数的直接恢复式 \(s=\langle w,(-y,x)\rangle\)。这个内积只是本题在平面模型中写逆映射的方便方法，切丛本身的定义没有使用它。

	由于 \(F_k(e^{i\theta})=e^{ik\theta}\)，局部角度微分把速度 \(s\) 送到 \(ks\)。因此在上述乘积表示下，
	\[
		\mathrm{T}F_k(p,s)=(p^k,ks).
	\]
	\(k=0\) 时，切映射把全部纤维向量送到点 \(1\) 上的零向量；\(k\ne0\) 时，各纤维上的线性映射可逆，但 \(F_k\) 是否全局单射仍取决于 \(k\)，只有 \(k=\pm1\) 时是微分同胚。
\end{solution}

\begin{exercise}
	\label{g1:ex:halfplane-square-differential}
	设 \(M=N=\mathbb H^2=\{(x,y):y\ge0\}\)，令 \(F(x,y)=(x,y^2)\)，目标坐标记为 \((X,Y)\)。计算 \(\mathrm{d}F\)、\(\mathrm{T}F\) 及 \(F^*(\mathrm{d}Y)\)。在边界点说明完整切空间、内向切锥和边界切空间各自怎样被微分映射，并判定 \(F\) 是否为微分同胚。
\end{exercise}
\begin{solution}
	多项式公式给出光滑延拓，故 \(F\) 光滑。对 \(a\partial_x+b\partial_y\)，有
	\[
		\mathrm{d}F_{(x,y)}(a\partial_x+b\partial_y)
		=a\partial_X+2yb\partial_Y.
	\]
	在切丛坐标中，
	\[
		\mathrm{T}F(x,y;a,b)=(x,y^2;a,2yb),\qquad
		F^*(\mathrm{d}Y)=\mathrm{d}(y^2)=2y\,\mathrm{d}y.
	\]
	边界点 \(p=(x,0)\) 的完整切空间由任意 \((a,b)\in\mathbb R^2\) 组成，微分的像为 \(\mathbb R\partial_X\)，核为 \(\mathbb R\partial_y\)。内向锥 \(b\ge0\) 也被送入目标边界切空间，且其像仍是整条 \(\mathbb R\partial_X\)；源边界切空间 \(b=0\) 则被同构地送到目标边界切空间。这里向内的非零速度可以在微分下消失，因为 \(y^2\) 在边界只有二阶变化。

	目标点的 \(\mathrm{d}Y\) 是非零协变向量，但拉回在边界为零；这与微分漏掉目标法向方向完全一致。又有 \(F^{-1}(X,Y)=(X,\sqrt Y)\)，所以 \(F\) 是同胚，却不是微分同胚。若逆映射光滑，链式法则会使边界处的微分可逆，与刚才的非零核矛盾。最后，\(\mathrm{T}M\) 的边界条件是 \(y=0\)，并不要求 \(b=0\)；额外要求 \(b=0\) 才得到 \(\mathrm{T}(\partial M)\)。
\end{solution}
\begin{exercise}
	\label{g1:ex:global-small-perturbation}
	设 \(h:\mathbb R^n\to\mathbb R^n\) 光滑，且 \(\sup_x\|\mathrm{D}h(x)\|\leq q<1\)，其中 \(q\geq0\)。
	\begin{enumerate}
		\item 证明 \(f(x)=x+h(x)\) 是全局光滑微分同胚，并且其逆映射的 Lipschitz 常数不超过 \((1-q)^{-1}\)。须说明对任意目标点如何选取包含迭代的闭球，不能只引用局部逆函数定理声称满射。
		\item 对 \(F_\varepsilon(x,y)=(x+\varepsilon\sin y,y+\varepsilon\sin x)\)，求它成为 \(\mathbb R^2\) 上全局光滑微分同胚的全部实参数 \(\varepsilon\)。
	\end{enumerate}
\end{exercise}
\begin{solution}
	1. 沿两点间的直线段积分 \(\mathrm{D}h\)，得到 \(|h(x)-h(x')|\leq q|x-x'|\)，因此
	\[
		|f(x)-f(x')|\geq(1-q)|x-x'|.
	\]
	这先给出全局单射和逆映射一旦存在时的 Lipschitz 控制。固定任意目标 \(b\in\mathbb R^n\)，令 \(T_b(x)=b-h(x)\)。选取
	\[
		R>\frac{|b-h(0)|}{1-q},
	\]
	则对 \(|x|\leq R\)，
	\[
		|T_b(x)|\leq|b-h(0)|+|h(x)-h(0)|
		\leq|b-h(0)|+qR<R.
	\]
	所以 \(T_b\) 把闭球 \(\overline B_R\) 映入自身，压缩因子不超过 \(q\)。\cref{g1:lem:euclidean-contraction}给出唯一不动点，即 \(f(x)=b\) 的解。由于 \(b\) 任意，\(f\) 满射，故为双射。

	此外 \(|\mathrm{D}f(x)v|\geq(1-q)|v|\)，所以 \(\mathrm{D}f(x)\) 处处可逆。逆函数定理在每一点给出光滑局部逆，这些局部逆都等于已经构造的全局集合逆，因而全局逆在每一点附近光滑，也就整体光滑。第一段的不等式给出所求全局 Lipschitz 界。

	2. 写 \(h_\varepsilon(x,y)=(\varepsilon\sin y,\varepsilon\sin x)\)，则
	\[
		\mathrm{D}h_\varepsilon=
		\begin{pmatrix}0&\varepsilon\cos y\\\varepsilon\cos x&0\end{pmatrix},
		\qquad \|\mathrm{D}h_\varepsilon\|\leq|\varepsilon|.
	\]
	这个范数界也可由矩阵 \(\mathrm{D}h_\varepsilon^{\mathsf T}\mathrm{D}h_\varepsilon\) 的两个对角元直接看出。故 \(|\varepsilon|<1\) 时，第一问证明它是全局光滑微分同胚。

	若 \(|\varepsilon|\geq1\)，可选 \(t\in[0,\pi]\) 满足 \(\cos t=1/\varepsilon\)。在 \((t,t)\) 处，
	\[
		\mathrm{D}F_\varepsilon(t,t)=\begin{pmatrix}1&1\\1&1\end{pmatrix}
	\]
	不可逆。若存在光滑逆，对复合恒等式求导会迫使这一矩阵可逆，矛盾。因此全部允许参数恰为 \(-1<\varepsilon<1\)。
\end{solution}

\begin{exercise}
	\label{g1:ex:cubic-implicit-parameters}
	考虑方程 \(y^3-\lambda y-x=0\)，其中 \((\lambda,x)\) 为参数，\(y\) 为未知量。
	\begin{enumerate}
		\item 证明在 \((\lambda,x,y)=(1,0,1)\) 附近，存在唯一接近 \(1\) 的光滑解 \(y=\varphi(\lambda,x)\)，并求该点的 \(\partial_x\varphi,\partial_\lambda\varphi,\partial_x^2\varphi\)。
		\item 列出 \((\lambda,x)=(1,0)\) 时的全部实数解，说明它们为何不违背上一问的唯一性。
		\item 找出所有满足方程但 \(\mathrm{D}_yF=0\) 的点，并写出其在参数平面中的关系。说明在这些点，方程是否仍能把 \(x\) 光滑地解成 \((\lambda,y)\) 的函数。
	\end{enumerate}
\end{exercise}
\begin{solution}
	1. 令 \(F(\lambda,x,y)=y^3-\lambda y-x\)，则在指定点有 \(F=0\)、\(\mathrm{D}_yF=3y^2-\lambda=2\ne0\)。\cref{g1:thm:implicit-function-parameters}给出所述局部光滑解。对方程求导，得到
	\[
		(3\varphi^2-\lambda)\partial_x\varphi=1,\qquad
		(3\varphi^2-\lambda)\partial_\lambda\varphi=\varphi.
	\]
	因此在 \((1,0)\) 处，两项一阶导数均为 \(1/2\)。对第一式再对 \(x\) 求导，
	\[
		6\varphi(\partial_x\varphi)^2
		+(3\varphi^2-\lambda)\partial_x^2\varphi=0,
	\]
	代入 \(\varphi=1\) 得 \(\partial_x^2\varphi=-3/4\)。

	2. 此时方程为 \(y^3-y=y(y-1)(y+1)=0\)，共有 \(y=-1,0,1\) 三个解。隐函数定理中的 \(V\) 是 \(1\) 的一个小邻域，可以排除另外两点；它断言在所选 \(V\) 内唯一，并未断言整条实轴上唯一。

	3. 联立 \(3y^2-\lambda=0\) 与 \(F=0\)，得
	\[
		\lambda=3y^2,\qquad x=-2y^3.
	\]
	故参数满足 \(27x^2=4\lambda^3\)、\(\lambda\geq0\)。反过来，这条参数曲线上的每一点都可用某个实数 \(y\) 如此表示，所以没有遗漏。这里以 \(y\) 为待解变量的偏导数失去可逆性，但 \(\mathrm{D}_xF=-1\) 始终非零；事实上方程全局等价于
	\[
		x=y^3-\lambda y,
	\]
	它对 \((\lambda,y)\) 是光滑函数。一个特定的未知量选择失效，不等于同一方程在所有坐标选择下都失去光滑图像表示。
\end{solution}

\begin{exercise}
	\label{g1:ex:rank-two-explicit-coordinates}
	设 \(F:\mathbb R^3\to\mathbb R^3\) 为
	\[
		F(x,y,z)=\bigl(x+y,(x+y)^2,ze^{x+y}\bigr).
	\]
	证明它处处秩为 \(2\)，给出显式的全局光滑坐标变化 \(\phi,\psi\)，使 \(\psi\circ F\circ\phi^{-1}(u,v,w)=(u,v,0)\)。再求各个目标点的原像以及 \(\ker\mathrm{d}F_{(x,y,z)}\)。
\end{exercise}
\begin{solution}
	记 \(s=x+y\)。导数矩阵为
	\[
		\mathrm{D}F=\begin{pmatrix}
			1&1&0\\
			2s&2s&0\\
			ze^s&ze^s&e^s
		\end{pmatrix}.
	\]
	第二行是第一行的 \(2s\) 倍，而第一、第三行因第三列的 \(e^s\ne0\) 线性无关，所以秩恰为 \(2\)。

	取
	\[
		\phi(x,y,z)=(x+y,ze^{x+y},y),\qquad
		\psi(a,b,c)=(a,c,b-a^2).
	\]
	它们的逆分别为
	\[
		\phi^{-1}(u,v,w)=(u-w,w,ve^{-u}),\qquad
		\psi^{-1}(\eta_1,\eta_2,\eta_3)
		=(\eta_1,\eta_3+\eta_1^2,\eta_2).
	\]
	所有公式在整个 \(\mathbb R^3\) 上定义且光滑，故这确是全局坐标变化。直接复合得 \(F\circ\phi^{-1}(u,v,w)=(u,u^2,v)\)，再用 \(\psi\) 就得到 \((u,v,0)\)。

	因此目标点 \((a,b,c)\) 若满足 \(b\ne a^2\)，原像为空；若 \(b=a^2\)，则
	\[
		F^{-1}(a,a^2,c)=\{(a-t,t,ce^{-a}):t\in\mathbb R\}.
	\]
	最后对向量 \((\xi,\eta,\zeta)\) 解 \(\mathrm{D}F(\xi,\eta,\zeta)^{\mathsf T}=0\)：第一行给出 \(\xi+\eta=0\)，第三行随后给出 \(\zeta=0\)。所以核为 \(\operatorname{span}\{(1,-1,0)\}\)，与这些原像直线的方向一致。一般常秩定理只保证局部坐标；本题之所以得到全局坐标，是因为上面确实写出了全局定义的互逆公式。
\end{solution}

\begin{exercise}
	\label{g1:ex:singular-level-velocities}
	令 \(f_+(x,y)=x^2+y^2\)、\(f_-(x,y)=x^2-y^2\)，并记 \(Z_\pm=f_\pm^{-1}(0)\)。它们在原点的微分都是零。
	\begin{enumerate}
		\item 求 \(Z_+,Z_-\)，并分别求所有在区间 \((-\varepsilon,\varepsilon)\) 上定义、满足 \(\gamma(0)=0\) 且 \(\gamma(t)\in Z_\pm\) 的 \(C^1\) 曲线所能有的速度 \(\gamma'(0)\)。
		\item 解释为什么不能仅凭 \(\mathrm{d}f_\pm|_0=0\)，就把 \(\ker\mathrm{d}f_\pm|_0=\mathbb R^2\) 当作相应零集的切空间。证明 \(Z_-\) 在原点附近不能由光滑坐标变化变成任何线性坐标子空间的一部分。
	\end{enumerate}
\end{exercise}
\begin{solution}
	1. 因平方和非负，\(Z_+=\{(0,0)\}\)，任何满足条件的曲线都是常值曲线，所以速度集合是 \(\{0\}\)。另一方面，
	\[
		Z_-=\{y=x\}\cup\{y=-x\}.
	\]
	若 \(\gamma(t)=(x(t),y(t))\) 的速度为 \((a,b)\)，则 \(x(t)=at+o(t)\)、\(y(t)=bt+o(t)\)。将 \(x(t)^2-y(t)^2=0\) 除以 \(t^2\) 后取极限，得到 \(a^2-b^2=0\)，即 \(b=a\) 或 \(b=-a\)。反过来，每个满足此条件的向量都由直线曲线 \(\gamma(t)=t(a,b)\) 实现。因此速度集合恰为这两条直线的并。

	2. 链式法则只说明所有实际曲线速度都落在微分的核中，不能在没有正则条件时反推每个核向量都来自零集中的曲线。这里两个微分的核同为整个 \(\mathbb R^2\)，但第一个速度集合只有零，第二个速度集合是两条直线的并，都不是整个平面。第二个集合甚至不是向量空间：它含有 \((1,1)\) 与 \((1,-1)\)，却不含二者之和 \((2,0)\)。

	若存在原点附近的光滑坐标变化，把 \(Z_-\) 变成某个线性坐标子空间在开邻域中的部分，那么该部分的曲线速度集合就是那个线性子空间。对坐标变化及其逆分别使用链式法则，原来的速度集合必是此子空间经一个可逆线性映射的逆像，也必为线性子空间，与上一段矛盾。这证明不能用光滑坐标把交叉零集拉直。两个例子都在原点之外具有秩 \(1\)，在原点处变为秩 \(0\)，正好显示常秩与正则性条件的作用。
\end{solution}
\begin{exercise}\label{g1:ex:orthogonal-group-level-set}
	把 \(n\times n\) 实矩阵空间看作 \(\mathbb R^{n^2}\)，对称矩阵空间记为 \(\operatorname{Sym}_n(\mathbb R)\)。设 \(n\ge1\)，令
	\[
		F:\operatorname{Mat}_n(\mathbb R)\to\operatorname{Sym}_n(\mathbb R),
		\qquad F(A)=A^{\mathsf T}A.
	\]
	证明单位矩阵 \(I\) 是正则值，求 \(\operatorname{O}(n)=F^{-1}(I)\) 的维数及每点切空间，并说明 \(\operatorname{SO}(n)=\{A\in\operatorname{O}(n):\det A=1\}\) 也是同维数的紧嵌入子流形。
\end{exercise}
\begin{solution}
	对任意矩阵 \(H\)，展开
	\[
		F(A+tH)=A^{\mathsf T}A+t(A^{\mathsf T}H+H^{\mathsf T}A)+t^2H^{\mathsf T}H,
	\]
	所以 \(\mathrm{d}F_A(H)=A^{\mathsf T}H+H^{\mathsf T}A\)。若 \(A^{\mathsf T}A=I\)，给定任意对称矩阵 \(B\)，取 \(H=AB/2\)，就有 \(\mathrm{d}F_A(H)=B\)。因此 \(I\) 为正则值。对称矩阵由 \(n\) 个对角元及 \(n(n-1)/2\) 个上三角元决定，维数为 \(n(n+1)/2\)，故
	\[
		\dim\operatorname{O}(n)=n^2-\frac{n(n+1)}2=\frac{n(n-1)}2.
	\]
	正则水平集的切空间公式给出
	\[
		\mathrm{T}_A\operatorname{O}(n)=\{H:A^{\mathsf T}H+H^{\mathsf T}A=0\}
		=\{A\Omega:\Omega^{\mathsf T}=-\Omega\}.
	\]
	最后一个等式由 \(\Omega=A^{\mathsf T}H\) 得到。

	\(\operatorname{O}(n)\) 是多项式映射 \(F\) 的闭水平集；每列长度为一，所以所有矩阵元的绝对值都不超过一，因而它闭且有界，故紧。其上 \((\det A)^2=1\)，所以行列式只取 \(\pm1\)。集合 \(\det A=1\) 在 \(\operatorname{O}(n)\) 中既开又闭：例如它也是连续函数行列式的正值原像。因此 \(\operatorname{SO}(n)\) 是同维数的开子流形，也是紧集的闭子集。其包含到矩阵空间的映射为两个嵌入的复合，仍为嵌入。
\end{solution}

\begin{exercise}\label{g1:ex:transverse-level-intersection}
	设 \(M^m\) 为无边界光滑流形，\(F:M\to\mathbb R^a\)、\(G:M\to\mathbb R^b\) 光滑，零分别是二者的正则值。记 \(S=F^{-1}(0)\)、\(T=G^{-1}(0)\)。证明：若对每个 \(p\in S\cap T\) 都有 \(\mathrm{T}_pS+\mathrm{T}_pT=\mathrm{T}_pM\)，则非空交集 \(S\cap T\) 是 \((m-a-b)\) 维嵌入子流形，并且
	\[
		\mathrm{T}_p(S\cap T)=\mathrm{T}_pS\cap \mathrm{T}_pT.
	\]
	再用 \(F(x,y,z)=z\)、\(G(x,y,z)=z-xy\) 说明，两组约束各自正则并不足以保证这个结论。
\end{exercise}
\begin{solution}
	先证明所需的线性代数事实。设 \(A:V\to\mathbb R^a\)、\(B:V\to\mathbb R^b\) 都满射。联合映射 \((A,B)\) 满射，当且仅当 \(B|_{\ker A}\) 满射：一个方向取联合目标 \((0,w)\) 即可；反方向先选 \(v\) 使 \(Av=u\)，再在 \(\ker A\) 中找 \(v'\) 使 \(Bv'=w-Bv\)，就有 \((A,B)(v+v')=(u,w)\)。而 \(B|_{\ker A}\) 满射又等价于 \(\ker A+\ker B=V\)。若它满射，对任意 \(v\) 找 \(v'\in\ker A\) 满足 \(Bv'=Bv\)，则 \(v-v'\in\ker B\)；反之，将一个满足 \(Bv=w\) 的向量分解成两核之和，就得到 \(w\) 在 \(\ker A\) 中的原像。

	在交点处取 \(A=\mathrm{d}F_p\)、\(B=\mathrm{d}G_p\)。由正则水平集定理，\(\mathrm{T}_pS=\ker A\)、\(\mathrm{T}_pT=\ker B\)，题设因而使联合映射 \(\mathrm{d}(F,G)_p\) 满射。于是零是 \((F,G):M\to\mathbb R^{a+b}\) 的正则值，其水平集恰为 \(S\cap T\)，得到维数结论及
	\[
		\mathrm{T}_p(S\cap T)=\ker\mathrm{d}(F,G)_p
		=\ker\mathrm{d}F_p\cap\ker\mathrm{d}G_p.
	\]

	对所给反例，两个函数对 \(z\) 的偏导都为一，所以处处为淹没。但交集是 \(z=0,xy=0\)，即平面内两条相交的坐标轴。在原点附近去掉原点得到四条分支，不能是一维子流形；而在原点两个切平面都为 \(\{v_z=0\}\)，它们的和没有填满 \(\mathbb R^3\)。这正是题设横截条件的失败。
\end{solution}

\begin{exercise}\label{g1:ex:graph-proper-embedding}
	设 \(M,N\) 为无边界光滑流形，\(f:M\to N\) 光滑。证明图映射 \(j_f:M\to M\times N\)，\(j_f(p)=(p,f(p))\)，是固有嵌入，并求图像的切空间。再给出一个例子，说明图像嵌入固有不意味着原映射 \(f\) 固有。
\end{exercise}
\begin{solution}
	图映射单射，且
	\[
		\mathrm{d}(j_f)_p(v)=(v,\mathrm{d}f_p(v)),
	\]
	因第一分量为 \(v\) 而单射，故 \(j_f\) 为浸入。图像到 \(M\) 的逆映射是第一因子投影的限制，因此连续，\(j_f\) 是嵌入。其切空间经包含识别为
	\[
		\mathrm{T}_{(p,f(p))}\operatorname{graph}(f)
		=\{(v,\mathrm{d}f_p(v)):v\in \mathrm{T}_pM\}.
	\]

	还需说明图像闭。若 \((p,q)\) 不在图像中，则 \(q\ne f(p)\)；由 \(N\) 的 Hausdorff 性，取不交开邻域 \(V\ni q,W\ni f(p)\)，再由 \(f\) 连续取 \(U\ni p\) 使 \(f(U)\subset W\)。开集 \(U\times V\) 与图像不交，所以图像的补集开。由\cref{g1:prop:proper-injective-immersion}的闭像嵌入判据，\(j_f\) 固有。

	例如取 \(f:\mathbb R\to\mathbb R\) 为零映射。目标紧集 \(\{0\}\) 的原像是非紧的 \(\mathbb R\)，所以 \(f\) 不固有，但图映射 \(t\mapsto(t,0)\) 是固有嵌入。二者的目标空间不同；图映射的第一坐标额外记录了参数是否逃向无穷远。
\end{solution}

\begin{exercise}\label{g1:ex:ball-boundary-regular-level}
	设 \(m\ge2\)，给闭单位球 \(M=\{x\in\mathbb R^m:\|x\|\le1\}\) 以标准带边界光滑结构，令 \(F:M\to\mathbb R\)，\(F(x)=x^1\)。分别讨论 \(|c|<1\)、\(|c|=1\) 与 \(|c|>1\) 时 \(c\) 对 \(F\) 及 \(F|_{\partial M}\) 的正则性，描述其水平集，并在 \(|c|<1\) 时核验边界公式及横截性。
\end{exercise}
\begin{solution}
	闭球的标准边界结构可由局部边界函数 \(\rho(x)=1-\|x\|^2\) 看出：在球面点处至少一个偏导不为零，使用其余坐标加 \(\rho\)，由逆函数定理得到 \(\rho\ge0\) 的半空间图。因而 \(\mathrm{T}_xM\) 在所有点都可识别为 \(\mathbb R^m\)，而在边界上 \(\mathrm{T}_x\partial M=x^\perp\)。

	完整微分 \(\mathrm{d}F_x(v)=v^1\) 处处满射，所以每个 \(c\in\mathbb R\) 都是 \(F\) 的正则值。边界限制的微分则为 \(v\mapsto v^1\) 在 \(x^\perp\) 上的限制。它不满射恰好意味着 \(x^\perp\subset(e_1)^\perp\)，等价于 \(x=\pm e_1\)。因此边界限制仅有两个临界值 \(\pm1\)。

	若 \(|c|<1\)，水平集为
	\[
		S_c=\{(c,x^2,\ldots,x^m):(x^2)^2+\cdots+(x^m)^2\le1-c^2\},
	\]
	是一张 \((m-1)\) 维闭球；其边界恰为 \(S_c\cap S^{m-1}\)。在此交集上，\(x\ne\pm e_1\)，向量 \(w=x-ce_1\) 满足 \(w^1=0\)，所以 \(w\in \mathrm{T}_xS_c\)，但
	\[
		\langle x,w\rangle=1-c^2>0,
	\]
	故 \(w\notin \mathrm{T}_x\partial M\)。由于 \(\mathrm{T}_x\partial M\) 余维为一，这说明 \(\mathrm{T}_xS_c+\mathrm{T}_x\partial M=\mathrm{T}_xM\)，即横截。其边界切空间也直接等于 \(\{v:v^1=0,\langle x,v\rangle=0\}\)，符合带边界水平集定理。

	若 \(c=\pm1\)，水平集分别只有 \(\pm e_1\)；尽管完整微分仍满射，边界限制已在该点失去满秩，水平集因而不是预期的 \((m-1)\) 维带边界片。若 \(|c|>1\)，水平集为空，对两个映射都按空原像约定成为正则值。
\end{solution}

% !TeX root = ../../Book1.tex
% 纲要标题：## 第12章　向量丛、张量与整体构造
% 纲要位置：几何学计划纲要.md，第 568 行。
% 正文状态：初稿；尚未正式审核。

% 纲要内容开始
% 
% > 本章目标与证明范围：由切丛推广到一般向量丛，证明局部表示和整体对象之间的对应。多线性代数围绕几何运算重建，单位分解承担明确的构造任务。
% 
% 纲要内容结束

\chapter{向量丛、张量与整体构造}
\label{g1:ch:12}
\BookChapterNumber{g1:ch:12}

\begin{chapterabstract}
	切丛的各个纤维都是向量空间，但在不同坐标图中需要用随点变化的可逆矩阵互相识别。本章由这一结构定义一般向量丛，证明余循环粘合、常秩同态的核与像以及拉回、商丛的构造。随后以泛性质建立张量积，以行列式配对固定外幂和楔积的系数，区分体积形式与密度。在整体构造方面，完整证明光滑单位分解、闭集上的截面延拓、丛度量的存在与短正合丛列的分裂，最后把这些工具用于张量场、自然拉回和几类结构群约化。
\end{chapterabstract}

\begin{learninggoals}
	\begin{itemize}
		\item 从局部标架写出过渡矩阵及余循环条件，说明它们如何决定整体向量丛，并辨认局部平凡与整体平凡的区别。
		\item 通过泛性质与基展开解释张量、对称幂、外幂和收缩，计算楔积的系数、分次符号与换标架公式。
		\item 构造支集局部有限的光滑单位分解，明确闭集延拓及其相对版本所需的相容条件。
		\item 用单位分解构造丛度量，再用正交补分裂短正合丛列，同时辨认这些构造中依赖选择的部分。
		\item 把局部分量解释为光滑张量场，区分拉回丛与源流形上的张量，并用适配标架描述度量、定向及其他逐点几何结构。
	\end{itemize}
\end{learninggoals}

上一章的切空间在每张坐标图中都有一组基，但从一张图走进另一张图时，这组基会改变。过渡矩阵把一组基中的分量转换到另一组基中；这些局部识别的整体相容方式决定了能否选出一组全局基。Möbius 线丛将展示最小的例子：每个小圆弧上的纤维都是一条普通实直线，绕圆周一周后的符号变化，却阻止了处处非零的整体截面。局部模型相同，并不足以决定整体对象平凡。

\ProseParagraphBreak

另一些对象又确实总能从局部构造出来。例如，在每张局部标架中都能声明基向量正交归一；虽然这些局部内积未必一致，却可以用适当权重相加，得到整体丛度量。为什么同样的操作不能给 Möbius 线丛制造非零截面？关键不只在于权重是否光滑，还在于加权相加能否保持所需条件。正定性可以保持，非零性、非退化性和某些代数恒等式却未必可以。

\ProseParagraphBreak

本章沿用\chapref{g1:ch:11}的光滑流形、切丛和微分定义，默认底流形有限维、Hausdorff 且第二可数；正文讨论的实光滑向量丛均为有限秩，相关构造允许底流形带边界。读者只需熟悉有限维线性代数中的基、对偶、商空间和行列式。张量积的泛性质、对称与交替运算在这里由多线性映射重建。单位分解所需的紧穷竭与局部有限细化，也会恢复实际使用的构造。

\ProseParagraphBreak

阅读时可以把每一项构造分为两步：先检查同一纤维上的线性代数，再检查它随底点是否光滑，以及换标架后是否代表同一个对象。Lee 第二版的第10章可用于对照向量丛，第12--13章讨论张量与 Riemann 度量，第2章的 Partitions of Unity 一节对应光滑权重的构造；第14--15章及第16章的 Densities 一节可补充交替形式、定向与密度的阅读\cite[第40--47、249--271、304--399、427--433页]{Lee2012Smooth}。本章只建立形式的逐点代数及光滑性，外微分与积分留在\chapref{g1:ch:14}。

% !TeX root = ../../Book1.tex
% 纲要标题：### 12.1 向量丛及其粘合
% 纲要位置：几何学计划纲要.md，第 572 行。
% 正文状态：初稿；正文、证明、例子与配套习题解答已撰写，尚未正式审核。

% 纲要内容开始
% 
% * 局部平凡化、过渡函数和余循环
% * 丛同态、子丛及常秩条件
% * 拉回丛、直和与商丛
% * 标架丛与全局标架
% 
% 纲要内容结束

\section{向量丛及其粘合}
\label{g1:sec:1201}

切丛把每点的切空间组织成一个光滑对象。它的局部坐标同时记住基点与向量分量，而换坐标时，分量按随基点光滑变化的可逆矩阵变换。许多几何对象也有这种形式：曲面每点的法向量、线性方程在每点的解空间，以及随后要讨论的张量，都不是固定向量空间中的一组数，而是一族随基点变化的向量空间。向量丛保留了切丛构造中的这一共同部分。

\ProseParagraphBreak

本节讨论有限秩实光滑向量丛，底空间为 Hausdorff、第二可数的光滑流形，允许有边界；边界附近的光滑性沿用局部延拓到 Euclid 开集的定义。我们先弄清局部数据怎样粘合，再讨论如何逐纤维取子空间、商空间，以及如何更换底空间。关键问题始终是：逐点成立的线性代数，何时能随基点光滑地进行。

\subsection{局部平凡化、过渡函数和余循环}
\label{g1:subsec:1201:01}

\cref{g1:thm:tangent-bundle-atlas}已经证明，切丛在坐标域 \(U\) 上可写成 \(U\times\mathbb R^n\)，并验证了总空间的 Hausdorff 性与第二可数性。这种乘积表示只在基点的一个邻域上给定；换图时可以出现非恒定矩阵，因此它还没有给出整个切丛与 \(M\times\mathbb R^n\) 的同构。

\begin{definition}
	\label{g1:def:vector-bundle}
	设 \(M\) 为 Hausdorff、第二可数的光滑 \(n\) 维流形，可带边界，\(r\geq0\) 为整数。\(M\) 上的秩 \(r\) 的实\term[xiangliangcong]{向量丛}[vector bundle]由光滑流形 \(E\)、光滑满射 \(\pi:E\to M\) 及每个集合 \(E_p=\pi^{-1}(p)\) 上的 \(r\) 维实向量空间结构组成，并要求每个 \(p\in M\) 都有开邻域 \(U\) 和微分同胚
	\[
		\phi_U:\pi^{-1}(U)\longrightarrow U\times\mathbb R^r,
		\qquad \phi_U(v_p)=(p,[v_p]_U),
	\]
	使 \([\cdot]_U:E_q\to\mathbb R^r\) 对每个 \(q\in U\) 都是线性同构。这样的 \(\phi_U\) 称为\term[jubupingfanhua]{局部平凡化}[local trivialization]。\(E\)、\(M\)、\(E_p\) 分别称为总空间、底空间和 \(p\) 上的\term[xianwei]{纤维}[fiber]；\(E|_U=\pi^{-1}(U)\) 表示限制丛。
\end{definition}

这里总空间同样要求 Hausdorff、第二可数。局部模型说明 \(\dim E=n+r\)；若底空间有边界，模型为 \(\mathbb H^n\times\mathbb R^r\)，非负坐标来自底空间，所以 \(\partial E=\pi^{-1}(\partial M)\)。秩零的丛每个纤维只有零向量，其总空间可由 \(\pi\) 识别为 \(M\)。本文固定一个秩；若底空间不连通，也可在不同连通分支分别规定秩。

\begin{definition}
	\label{g1:def:bundle-section}
	设 \(\pi:E\to M\) 为有限秩实光滑向量丛，\(U\subset M\) 开。满足 \(\pi\circ s=\operatorname{id}_U\) 的光滑映射 \(s:U\to E\) 称为 \(E\) 在 \(U\) 上的\term[jiemian]{截面}[section]，全体光滑截面记为 \(\Gamma(U,E)\)，并简记 \(\Gamma(E)=\Gamma(M,E)\)。若秩为 \(r\)，一组截面 \(e_1,\ldots,e_r\) 在每个 \(p\in U\) 都组成 \(E_p\) 的有序基，就称为 \(U\) 上的\term[jububiaojia]{局部标架}[local frame]。
\end{definition}

局部平凡化把标准基向量送回纤维，便给出局部标架。反过来，给定局部标架 \(e=(e_1,\ldots,e_r)\)，记
\[
	e(p)v=\sum_{a=1}^r v^a e_a(p),\qquad v\in\mathbb R^r.
\]
在任一已有平凡化中，\(e_a\) 的分量组成一个光滑可逆矩阵 \(A(p)\)，因而 \((p,v)\mapsto e(p)v\) 及其逆都光滑。逆的光滑性来自 \(A^{-1}=\operatorname{adj}(A)/\det A\)，这里行列式处处非零。于是局部标架与局部平凡化确实等价，而不仅是逐点选基。

\ProseParagraphBreak

在开覆盖 \(\{U_i\}\) 上选标架 \(e_i\)，其中下标 \(i\) 标记开集，标架内部的向量仍用 \(e_{i,a}\) 表示。在交集上存在唯一光滑函数
\[
	g_{ij}:U_i\cap U_j\longrightarrow\operatorname{GL}(r,\mathbb R)
\]
使
\begin{equation}
	\label{g1:eq:bundle-frame-transition}
	e_j=e_i g_{ij},\qquad v_i=g_{ij}v_j.
\end{equation}
第二式中 \(v_i,v_j\) 是同一个纤维向量在两组标架下的列分量，因为 \(e_i v_i=e_jv_j\)。因此标架变换与分量变换在这个下标约定下都出现 \(g_{ij}\)，但前一式把第 \(j\) 组基写成第 \(i\) 组基的线性组合，后一式则把第 \(j\) 组分量转换成第 \(i\) 组分量。

\ProseParagraphBreak

这些\term[guoduhanshu]{过渡函数}[transition function]满足
\begin{equation}
	\label{g1:eq:bundle-cocycle}
	g_{ii}=I_r,\qquad g_{ji}=g_{ij}^{-1},\qquad
	g_{ij}g_{jk}=g_{ik}.
\end{equation}
最后一个等式在三重交集上成立，称为\term[yuxunhuantiaojian]{余循环条件}[cocycle condition]：先由第 \(k\) 组分量换到第 \(j\) 组，再换到第 \(i\) 组，必须与直接转换相同。它由 \(e_k=e_jg_{jk}=e_i g_{ij}g_{jk}\) 及基的线性无关性得到。对截面 \(s\)，局部表达 \(s=e_i s_i\) 中的列函数 \(s_i:U_i\to\mathbb R^r\) 光滑，且满足 \(s_i=g_{ij}s_j\)；反之，满足这些关系的光滑列函数唯一确定一个光滑截面。纤维中的加法与数乘由此定义了截面的逐点运算，并使 \(\Gamma(E)\) 成为 \(C^\infty(M)\) 上的模。

\begin{example}
	\label{g1:exm:tangent-bundle-transition-functions}
	对光滑流形 \(M\)，在坐标图 \(x_i:U_i\to x_i(U_i)\) 上取切标架 \(e_i=(\partial_{x_i^1},\ldots,\partial_{x_i^n})\)。由\eqref{g1:eq:tangent-bundle-transition}，
	\[
		g_{ij}(p)=\mathrm{D}(x_i\circ x_j^{-1})(x_j(p)).
	\]
	链式法则正好给出 \(g_{ij}g_{jk}=g_{ik}\)，所以切丛的过渡函数是坐标变换的 Jacobian 矩阵。一般向量丛的矩阵只须光滑、可逆并满足余循环条件，无须来自某个底空间坐标变换的导数。乘积丛 \(M\times\mathbb R^r\) 则有处处相同的标准标架，过渡函数都可取为 \(I_r\)。
\end{example}

余循环条件足以把一批局部乘积真正粘成向量丛。证明中既要保证同一个向量不会因走不同的换图路径而得到不同代表，也要保证粘合后的空间满足流形的整体拓扑条件。

\begin{theorem}[向量丛的粘合]
	\label{g1:thm:bundle-cocycle-gluing}
	设 \(M\) 为 Hausdorff、第二可数的光滑流形，可带边界；\(\{U_i\}_{i\in I}\) 为开覆盖，\(r\geq0\) 为固定整数。若光滑映射 \(g_{ij}:U_i\cap U_j\to\operatorname{GL}(r,\mathbb R)\) 在相应的二重、三重交集上满足 \(g_{ii}=I_r\)、\(g_{ji}=g_{ij}^{-1}\)、\(g_{ij}g_{jk}=g_{ik}\)，即\eqref{g1:eq:bundle-cocycle}，则它们确定一个秩 \(r\) 的实光滑向量丛，其局部平凡化按 \((p,v_j)\mapsto(p,g_{ij}(p)v_j)\) 转换。该丛在保持这些局部识别的丛同构意义下唯一。
\end{theorem}
\begin{proof}
	先把各局部乘积视为互不相交的副本，令
	\[
		\mathcal E=\coprod_{i\in I}(U_i\times\mathbb R^r),\qquad
		(i,p,v)\sim(j,q,w)
		\quad\Longleftrightarrow\quad
		p=q\ \text{且}\ v=g_{ij}(p)w.
	\]
	自反性来自 \(g_{ii}=I_r\)，对称性来自 \(g_{ji}=g_{ij}^{-1}\)。若 \(v=g_{ij}(p)w\)、\(w=g_{jk}(p)z\)，则余循环条件给出 \(v=g_{ik}(p)z\)，所以关系也传递。令 \(E=\mathcal E/\!\sim\)，赋予商拓扑，并记商映射为 \(Q\)。由于等价关系不改变基点，\(\pi([i,p,v])=p\) 良定义且连续。

	商映射 \(Q\) 是开映射。事实上，对 \(\mathcal E\) 中的开集 \(O\)，它的饱和集在第 \(j\) 个副本中的部分为
	\[
		\bigcup_{i\in I}\left\{(p,w):p\in U_i\cap U_j,\ (p,g_{ij}(p)w)\in O_i\right\},
	\]
	其中 \(O_i=O\cap(U_i\times\mathbb R^r)\)。每项都是光滑同胚 \((p,w)\mapsto(p,g_{ij}(p)w)\) 下的开原像，故此并为开。商拓扑的定义于是说明 \(Q(O)\) 开。

	对每个 \(i\)，映射
	\[
		\tau_i:U_i\times\mathbb R^r\longrightarrow\pi^{-1}(U_i),
		\qquad\tau_i(p,v)=[i,p,v]
	\]
	是双射：每个基点在 \(U_i\) 的等价类都可换到第 \(i\) 个副本，而在同一副本中相等的两个代表必有相同分量。它连续且开，故为同胚；其像 \(\pi^{-1}(U_i)\) 也是 \(E\) 中的开集。在交集上有
	\[
		\tau_i^{-1}\tau_j(p,w)=(p,g_{ij}(p)w).
	\]
	再将每个 \(U_i\) 限制到底空间坐标域，便得到 \(E\) 的光滑坐标图；上述公式以及底空间换图的光滑性保证它们相容。在边界附近，矩阵元有局部光滑延拓，乘法公式同样给出半空间模型之间的光滑映射。

	还需核验 Hausdorff 性。两个点若有不同基点，用底空间中不交开邻域的逆像将它们分离；若基点相同，就把两点放进同一个 \(\tau_i\) 的像内，利用 \(U_i\times\mathbb R^r\) 的 Hausdorff 性分离。该像是 \(E\) 的开集，所以所得邻域也是总空间的开邻域。

	第二可数性不能由任意多个局部图直接推出，不过底空间的第二可数性提供可数子覆盖。具体地，对底空间可数基中每个包含在某个 \(U_i\) 内的基元素，选择一个包含它的 \(U_i\)；这些被选开的集合构成可数覆盖。每个对应乘积 \(U_i\times\mathbb R^r\) 有可数基，把它们通过 \(\tau_i\) 送到 \(E\) 再取并，就得到总空间的可数基。

	在纤维 \(E_p\) 中，选任一个 \(p\in U_i\)，用 \(\tau_i(p,\cdot)\) 运输 \(\mathbb R^r\) 的向量空间运算。换到另一副本时使用的是线性同构 \(g_{ij}(p)\)，故运算与所选副本无关。投影在局部是 \((p,v)\mapsto p\)，光滑且满射，\(\tau_i^{-1}\) 则是所需局部平凡化。

	最后，若另一个丛 \(E'\) 有同样局部识别 \(\tau_i'\)，则 \([i,p,v]\mapsto\tau_i'(p,v)\) 由相同过渡关系而良定义。它在每个局部乘积上为恒等坐标映射，故是光滑丛同构。任何保持局部识别的映射都必须使用这个公式，因此唯一。秩零时取 \(\operatorname{GL}(0,\mathbb R)=\{I_0\}\)，所有论证仍成立。
\end{proof}

这里的唯一性保留了指定的局部识别；改变标架当然会改变矩阵。若在 \(U_i\) 上以 \(\widetilde v_i=h_i v_i\) 改换分量，其中 \(h_i:U_i\to\operatorname{GL}(r,\mathbb R)\) 光滑，则
\begin{equation}
	\label{g1:eq:bundle-cocycle-change}
	\widetilde e_i=e_i h_i^{-1},\qquad
	\widetilde g_{ij}=h_i g_{ij}h_j^{-1}.
\end{equation}
后一个等式由 \(\widetilde v_i=h_i g_{ij}h_j^{-1}\widetilde v_j\) 得出。反过来，两组余循环若满足此式，\([i,p,v]\mapsto[i,p,h_i(p)v]\) 就给出所粘丛之间的同构：把 \(v=g_{ij}w\) 代入，得到 \(h_i v=\widetilde g_{ij}h_jw\)，恰好保证映射与代表选择无关。不同覆盖上的表示可先限制到两覆盖的交集，再作同一比较。因此过渡矩阵是丛在选定标架中的数据，并不是无需改换的绝对记号。

\begin{example}[Möbius 线丛]
	\label{g1:exm:mobius-line-bundle}
	令 \(M=S^1=\{e^{\mathrm{i}t}:t\in\mathbb R\}\)，把 \(\mathbb R\times\mathbb R\) 按
	\[
		(t+2\pi,v)\sim(t,-v)
	\]
	及其迭代关系取商，记所得空间为 \(L\)，投影为 \([t,v]\mapsto e^{\mathrm{i}t}\)。每个足够短的圆弧 \(U_i\) 上可选光滑角度提升 \(t_i:U_i\to\mathbb R\)，从而有局部表示 \((p,v)\mapsto[t_i(p),v]\)。若在交叠上 \(t_j-t_i=2\pi k_{ij}\)，则 \(k_{ij}\) 为局部常值整数，且
	\[
		[t_j(p),v]=[t_i(p),(-1)^{k_{ij}(p)}v].
	\]
	所以过渡函数为 \(g_{ij}=(-1)^{k_{ij}}\)，满足余循环条件。\cref{g1:thm:bundle-cocycle-gluing}给出秩一光滑丛结构；这些局部表示也是所写商空间的图，因为任一长度小于 \(2\pi\) 的提升圆弧及其 \(2\pi\) 平移副本，经商映射的饱和集仍开。这个丛称为 Möbius 线丛。

	它的光滑截面恰对应满足
	\begin{equation}
		\label{g1:eq:mobius-section-antiperiodic}
		f(t+2\pi)=-f(t)
	\end{equation}
	的光滑函数 \(f:\mathbb R\to\mathbb R\)，截面公式为 \(s(e^{\mathrm{i}t})=[t,f(t)]\)。因为 \([t+2\pi,f(t+2\pi)]=[t,-f(t+2\pi)]\)，此条件正是代表选择无关性；光滑性可在各提升圆弧上核验。若截面处处非零，\(f(0)\) 与 \(f(2\pi)=-f(0)\) 异号，中间值定理迫使某处 \(f\) 为零，矛盾。故 \(L\) 没有处处非零的光滑截面，甚至没有这样的连续截面。乘积线丛却有常值分量为 \(1\) 的截面，任何丛同构都会保持非零性，因此 \(L\) 与 \(S^1\times\mathbb R\) 不同构。
\end{example}

\begin{figure}[htbp]
	\centering
	\input{Figures/Chapter12-MobiusSection}
	\caption[Möbius 线丛中的截面与端点粘合]{矩形只画出线丛纤维的有限一段。左右竖边按相反方向粘合，\((2\pi,v)\) 对应 \((0,-v)\)。蓝线为截面分量 \(f(t)=\cos(t/2)\)，它的两端经过粘合相合，而在中间穿过零。红色虚线的分量恒为正，左右端点却不能按粘合关系相合，因而不是圆周上的截面。}
	\label{g1:fig:mobius-section}
\end{figure}

这说明局部平凡并未排除全局扭转。扭转并不是某个纤维内部的缺陷：每条纤维都是一维实向量空间；障碍出现在沿圆周比较各条纤维的基时。

\subsection{丛同态、子丛及常秩条件}
\label{g1:subsec:1201:02}

对同一个底空间上的两族向量空间，可以逐点施加线性映射。不过逐点线性只保证每个纤维内的行为，并不保证映射随基点连续或光滑。局部标架把后一个要求化为矩阵元的光滑性。

\begin{definition}
	\label{g1:def:bundle-morphism}
	设 \(E\to M\)、\(F\to M\) 为有限秩实光滑向量丛。覆盖恒等映射的\term[congtongtai]{丛同态}[bundle morphism]是光滑映射 \(\Phi:E\to F\)，使 \(\pi_F\Phi=\pi_E\)，且每个限制 \(\Phi_p:E_p\to F_p\) 都线性。如果它有同样性质的光滑逆，称为\term[congtonggou]{丛同构}[bundle isomorphism]。若 \(E\to N\)、\(F\to M\) 的底空间不同，覆盖光滑映射 \(f:N\to M\) 的丛同态则要求 \(\pi_F\Phi=f\pi_E\)，并要求 \(E_q\to F_{f(q)}\) 线性。
\end{definition}

\begin{proposition}[丛同态的矩阵]
	\label{g1:prop:bundle-morphism-matrices}
	设 \(E,F\) 为 Hausdorff、第二可数光滑流形 \(M\) 上的秩 \(r,s\) 的实光滑向量丛，允许 \(M\) 有边界。在共同开覆盖上取局部标架，其过渡函数为 \(g_{ij}^E,g_{ij}^F\)。一个覆盖恒等映射的逐纤维线性映射 \(\Phi:E\to F\) 为光滑丛同态，当且仅当它的局部矩阵 \(A_i:U_i\to\operatorname{Mat}_{s\times r}(\mathbb R)\) 光滑；这些矩阵满足
	\begin{equation}
		\label{g1:eq:bundle-morphism-transition}
		A_i g_{ij}^E=g_{ij}^F A_j.
	\end{equation}
	反之，任何满足此关系的光滑矩阵族都确定唯一丛同态。特别地，逐纤维为线性同构的光滑丛同态，其逆自动光滑。
\end{proposition}
\begin{proof}
	局部矩阵由 \(\Phi(e_i v)=f_i A_i v\) 定义。若 \(\Phi\) 光滑，则其作用于每个光滑基截面 \(e_{i,a}\) 后的分量光滑，故 \(A_i\) 各列光滑。反过来，光滑矩阵给出的局部映射 \((p,v)\mapsto(p,A_i(p)v)\) 光滑。交叠上，先换源分量再映射得到 \(A_i g_{ij}^E v_j\)，先映射再换目标分量得到 \(g_{ij}^F A_jv_j\)，二者对所有 \(v_j\) 相同恰好等价于\eqref{g1:eq:bundle-morphism-transition}。因此相容矩阵族唯一粘成光滑映射。最后，若每个 \(A_i(p)\) 都可逆，其逆矩阵光滑且局部逆彼此相容，从而构成光滑丛逆。
\end{proof}

例如映射的微分 \(\mathrm{T}F:\mathrm{T}M\to \mathrm{T}N\) 覆盖 \(F:M\to N\)，局部矩阵为坐标表达的导数，见\cref{g1:prop:differential-chain-rule}。底空间不同的情形只须在目标变换矩阵中代入 \(F(p)\)，矩阵判据及其证明不变。另一个常用例子是同底丛同态作用于截面：\(s\mapsto\Phi\circ s\) 把 \(\Gamma(E)\) 送入 \(\Gamma(F)\)，并与光滑函数的逐点乘法相容。

\begin{definition}
	\label{g1:def:vector-subbundle}
	设 \(E\to M\) 为秩 \(r\) 的实光滑向量丛，\(0\leq k\leq r\)。若每个 \(p\) 给定一个 \(k\) 维线性子空间 \(F_p\subset E_p\)，且每点附近存在 \(k\) 个光滑截面 \(s_1,\ldots,s_k\)，在每个邻近点都线性无关并张成相应的 \(F_p\)，则称 \(F=\coprod_p F_p\subset E\) 为秩 \(k\) 的光滑\term[zicong]{子丛}[subbundle]。
\end{definition}

定义中的光滑局部张成条件不可删去。例如在 \(\mathbb R\times\mathbb R^2\) 中，令 \(F_t=\mathbb R(1,0)\) 当 \(t\leq0\)，令 \(F_t=\mathbb R(0,1)\) 当 \(t>0\)。各个子空间的维数都为一，但若零附近有一个处处非零的连续张成截面，它在零点的值必须既落在横轴又落在纵轴，因而只能为零，矛盾。恒定维数没有替代连续变化，更没有替代光滑变化。

\begin{proposition}[子丛的适配标架]
	\label{g1:prop:subbundle-adapted-frame}
	设 \(E\to M\) 是 Hausdorff、第二可数光滑流形上的秩 \(r\) 实光滑向量丛，允许有边界。若 \(s_1,\ldots,s_k\) 在点 \(p\) 处线性无关，则在 \(p\) 的某个邻域上可将这些光滑截面补成 \(E\) 的局部标架。特别地，秩 \(k\) 子丛 \(F\subset E\) 带有唯一的光滑向量丛结构，使其局部标架给出的坐标为光滑平凡化；其拓扑是 \(E\) 的子空间拓扑，并且存在局部标架使
	\[
		E|_U\simeq U\times\mathbb R^r,\qquad
		F|_U\simeq U\times(\mathbb R^k\times\{0\}).
	\]
\end{proposition}
\begin{proof}
	先在 \(p\) 附近固定一个 \(E\) 的平凡化，把 \(s_a\) 写成光滑列向量。在 \(p\) 处用常向量 \(b_{k+1},\ldots,b_r\) 将其补为一组基，并把这些常向量在该平凡化下视为光滑截面。所成 \(r\times r\) 矩阵的行列式在 \(p\) 非零，由连续性，在某邻域仍非零。因此这些截面在整个邻域组成标架。\(k=0\) 时直接选任意标架，\(k=r\) 时无须补充。

	对子丛应用这个结论，前 \(k\) 个坐标向量正好张成 \(F\)。因此环境平凡化限制到 \(F\) 后为到 \(U\times\mathbb R^k\) 的同胚，局部拓扑正是子空间拓扑。两组这种标架的变换矩阵把 \(\mathbb R^k\times\{0\}\) 映到自身，其左上 \(k\times k\) 块为光滑可逆矩阵，故这些限制图光滑相容。Hausdorff 性与第二可数性由子空间继承，局部图给出维数 \(\dim M+k\) 及向量丛结构。唯一性由规定的局部图立即得到。边界点处仅用了光滑矩阵及非零行列式，所得局部模型仍为 \(\mathbb H^n\times\mathbb R^k\)。
\end{proof}

现在考虑一个丛同态的核与像。每条纤维内它们都是线性子空间；要成为子丛，维数先要在附近不变，而常秩条件还会给出光滑的局部基。\cref{g1:thm:constant-rank}中用过“选非零子式，再消去其余行”的办法；此处只需要对一族矩阵作消元，不需要对底空间改换常秩标准坐标。

\begin{theorem}[常秩丛同态]
	\label{g1:thm:constant-rank-bundle-morphism}
	设 \(E,F\) 为 Hausdorff、第二可数光滑流形 \(M\) 上的秩 \(r,s\) 实光滑向量丛，允许 \(M\) 有边界；\(\Phi:E\to F\) 为覆盖恒等映射的光滑丛同态。若 \(\operatorname{rank}\Phi_p=k\) 对所有 \(p\in M\) 恒成立，则
	\[
		\ker\Phi=\coprod_{p\in M}\ker\Phi_p,\qquad
		\operatorname{im}\Phi=\coprod_{p\in M}\operatorname{im}\Phi_p
	\]
	分别是 \(E,F\) 的秩 \(r-k,k\) 的光滑子丛。若秩只要求局部恒定，则此结论按各秩恒定的开集分别成立。
\end{theorem}
\begin{proof}
	固定一点并取局部标架，把 \(\Phi\) 表示为光滑 \(s\times r\) 矩阵 \(A(p)\)。先设 \(k>0\)。在基点选择一个非零 \(k\) 阶子式，重排源与目标的基后，将对应子块放到左上角。缩小邻域使该块一直可逆，写成
	\[
		A=\begin{pmatrix}B&C\\D&H\end{pmatrix},\qquad B\in\operatorname{GL}(k,\mathbb R).
	\]
	左乘可逆矩阵作消元，得到
	\[
		\begin{pmatrix}I&0\\-DB^{-1}&I\end{pmatrix}A
		=\begin{pmatrix}B&C\\0&H-DB^{-1}C\end{pmatrix}.
	\]
	再把右上块通过列消元变为零，可见右边的秩为 \(k+\operatorname{rank}(H-DB^{-1}C)\)。常秩条件遂迫使 \(H=DB^{-1}C\) 在整个邻域上成立。

	于是 \(A(x,y)=0\) 等价于 \(Bx+Cy=0\)，亦即 \(x=-B^{-1}Cy\)。所以矩阵
	\[
		\begin{pmatrix}-B^{-1}C\\I_{r-k}\end{pmatrix}
	\]
	的 \(r-k\) 个列向量给出核的光滑局部基。从下方的单位矩阵可见，这些列向量线性无关。另一方面，\(A\) 的前 \(k\) 列组成 \(\begin{psmallmatrix}B\\D\end{psmallmatrix}\)，因 \(B\) 可逆而线性无关；其余列为这些列乘以 \(B^{-1}C\)，故前 \(k\) 列正好给出像的光滑局部基。两族列向量都在先前选定标架下表示光滑截面，因而核与像满足子丛定义。

	若 \(k=0\)，矩阵处处为零，核为 \(E\)，像为零子丛。若 \(k=r\) 或 \(k=s\)，上述相应的空块与空基直接略去，结论仍成立。边界点处，\(B^{-1}\) 的矩阵元是光滑函数除以非零的 \(\det B\)，所以在半空间图中仍可光滑延拓；所有等式只在实际底空间邻域内使用，不要求延拓后的矩阵在外侧仍保持秩 \(k\)。
\end{proof}

\begin{example}
	\label{g1:exm:variable-rank-bundle-map}
	在 \(\mathbb R\) 上的两个乘积线丛之间，令 \(\Phi(t,v)=(t,tv)\)。当 \(t\ne0\) 时，它的核为零、像为整条纤维；在 \(t=0\) 时，核为整条纤维、像为零。故核与像都不能在零附近成为向量丛。即使把底空间改成 \([0,\infty)\)，同样的问题仍在边界点零发生。失败原因是秩的变化，不是边界使矩阵消元失效。
\end{example}

特别地，一个逐纤维单射的光滑丛同态把源丛同构到目标的像子丛。像是子丛由常秩定理得到；把同态看成到这个像的映射后，它逐纤维可逆，\cref{g1:prop:bundle-morphism-matrices}保证逆光滑。这允许我们把一个抽象的单射丛同态与其实现的子丛明确对应起来。
\subsection{拉回丛、直和与商丛}
\label{g1:subsec:1201:03}

设 \(f:N\to M\) 光滑，\(E\to M\) 是向量丛。如果希望在 \(N\) 的每点使用 \(f(q)\) 处的纤维，就应把新纤维定义为 \(E_{f(q)}\)。不同的 \(q\) 即使具有相同的像，也必须保留为不同的基点，因此不能只取 \(E\) 中这些纤维的并，而要把 \(q\) 一同记下来。

\begin{proposition}[拉回丛]
	\label{g1:prop:pullback-bundle}
	设 \(M,N\) 为 Hausdorff、第二可数的光滑流形，均可带边界；\(f:N\to M\) 光滑，\(\pi:E\to M\) 为秩 \(r\) 的实光滑向量丛。集合
	\[
		f^*E=\{(q,v)\in N\times E:f(q)=\pi(v)\},\qquad
		\pi'(q,v)=q
	\]
	带有自然的秩 \(r\) 实光滑向量丛结构，其拓扑是 \(N\times E\) 的子空间拓扑，称为 \(E\) 沿 \(f\) 的\term[lahuicong]{拉回丛}[pullback bundle]。映射 \(\widetilde f:f^*E\to E\)，\(\widetilde f(q,v)=v\)，是覆盖 \(f\) 的光滑丛同态。

	它具有如下泛性质：对任意 \(N\) 上的有限秩实光滑向量丛 \(H\)，以及覆盖 \(f\) 的光滑丛同态 \(\Psi:H\to E\)，存在唯一覆盖 \(\operatorname{id}_N\) 的光滑丛同态 \(\widehat\Psi:H\to f^*E\)，使 \(\widetilde f\widehat\Psi=\Psi\)。
\end{proposition}
\begin{proof}
	若 \(e_i\) 是 \(E\) 在 \(U_i\) 上的局部标架，定义
	\[
		\tau_i':f^{-1}(U_i)\times\mathbb R^r\longrightarrow(\pi')^{-1}(f^{-1}(U_i)),\qquad
		(q,a)\longmapsto(q,e_i(f(q))a).
	\]
	它是双射，逆为 \((q,v)\mapsto(q,[v]_i)\)。两个公式对所给子空间拓扑都连续，且定义域与像分别在相应局部范围内开，所以 \(\tau_i'\) 是局部平凡化的拓扑候选。它们的过渡映射为
	\[
		(q,a_j)\longmapsto(q,g_{ij}(f(q))a_j).
	\]
	这些矩阵光滑、可逆，且满足余循环条件。由\cref{g1:thm:bundle-cocycle-gluing}，得到唯一相容的光滑向量丛结构；其拓扑与上述子空间拓扑相同，因为两者在覆盖总空间的每个局部乘积上都由同一同胚给定。Hausdorff 性与第二可数性也可直接由 \(N\times E\) 的子空间性质得到。边界来自 \(N\) 的半空间坐标，纤维仍为完整的 \(\mathbb R^r\)。

	映射 \(\widetilde f\) 在这些图中写成 \((q,a)\mapsto(f(q),a)\)，故光滑，并且逐纤维线性。对泛性质，唯一可能的公式是
	\[
		\widehat\Psi(u)=\bigl(\pi_H(u),\Psi(u)\bigr).
	\]
	因为 \(\pi\Psi=f\pi_H\)，这个有序对确实落在 \(f^*E\) 中。局部矩阵中，它与 \(\Psi\) 使用相同的纤维矩阵，只是底映射变为恒等，所以光滑且逐纤维线性。此公式同时给出 \(\widetilde f\widehat\Psi=\Psi\) 与 \(\pi'\widehat\Psi=\pi_H\)，而这两个等式决定有序对的两个分量，故唯一。
\end{proof}

这两个交换关系可集中画成
\[
	\begin{tikzcd}[column sep=large,row sep=large]
		H \arrow[rr,bend left=16,"\Psi"] \arrow[dr,"\pi_H"'] \arrow[r,dashed,"\widehat\Psi"] & f^*E \arrow[r,"\widetilde f"] \arrow[d,"\pi'"'] & E \arrow[d,"\pi"] \\
		& N \arrow[r,"f"'] & M.
	\end{tikzcd}
\]
虚线箭头不增加一种任意选择；它只是把 \(\Psi(u)\) 与其新基点记录在一起。特别地，对 \(E\) 的截面 \(s\)，
\[
	(f^*s)(q)=(q,s(f(q)))
\]
是 \(f^*E\) 的光滑截面。对流形映射的微分，这一构造把覆盖 \(f\) 的 \(\mathrm{T}f:\mathrm{T}N\to \mathrm{T}M\) 改写成同底丛同态 \(\mathrm{T}N\to f^*\mathrm{T}M\)，其在 \(q\) 处仍是 \(\mathrm{d}f_q\)。纤维 \((f^*\mathrm{T}M)_q\) 是 \(\mathrm{T}_{f(q)}M\)，通常不是 \(\mathrm{T}_qN\)。

\ProseParagraphBreak

如果再有光滑 \(h:P\to N\)，则
\[
	h^*(f^*E)\longrightarrow(f\circ h)^*E,
	\qquad(p,(h(p),v))\longmapsto(p,v)
\]
是自然丛同构，其逆把 \((p,v)\) 送到 \((p,(h(p),v))\)。在从 \(E\) 拉回的局部平凡化中，两个公式的纤维分量均为恒等，故光滑。类似地，\(\operatorname{id}_M^*E\simeq E\)。这些明确同构说明，连续两次拉回与沿复合映射拉回给出相同的几何对象。

\ProseParagraphBreak

线性代数的直和则保留同一基点上的两种向量。对 \(M\) 上的秩 \(r,s\) 向量丛 \(E,F\)，定义\term[zhihecong]{直和丛}[direct sum bundle]
\[
	E\oplus F=\{(u,v)\in E\times F:\pi_E(u)=\pi_F(v)\},
	\qquad(E\oplus F)_p=E_p\oplus F_p.
\]
在共同平凡化下，它为 \(U_i\times(\mathbb R^r\oplus\mathbb R^s)\)，过渡矩阵为
\[
	\begin{pmatrix}g_{ij}^E&0\\0&g_{ij}^F\end{pmatrix}.
\]
分块乘法验证余循环条件，粘合定理给出光滑丛结构。局部表示与在 \(E\times F\) 中的子空间拓扑相容，证明与拉回丛的局部表示相同：从局部乘积到这个集合的映射及其逆都由两组平凡化给出。逐纤维投影和包含在局部都是固定线性映射，所以是光滑丛同态。这里的总空间不是整个 \(E\times F\)：直和要求两个向量属于同一基点，因而其维数为 \(\dim M+r+s\)。

\ProseParagraphBreak

商丛的目标是在每条纤维中忽略一个指定子空间的分量。虽然计算时可补选基来表示剩余分量，但商本身不应依赖所补的基，也不应要求预先选出全局补空间。

\begin{proposition}[商丛]
	\label{g1:prop:quotient-vector-bundle}
	设 \(M\) 为 Hausdorff、第二可数的光滑流形，可带边界；\(E\to M\) 为秩 \(r\) 实光滑向量丛，\(F\subset E\) 为秩 \(k\) 光滑子丛。集合
	\[
		E/F=\coprod_{p\in M}E_p/F_p
	\]
	有自然的秩 \(r-k\) 实光滑向量丛结构，其总空间拓扑为映射 \(q:E\to E/F\)，\(v_p\mapsto[v_p]\)，确定的商拓扑。\(q\) 是光滑满射丛同态，其逐纤维核为 \(F\)。如果 \(H\to M\) 是实光滑向量丛，\(\Phi:E\to H\) 是同底光滑丛同态，且 \(\Phi|_F=0\)，则存在唯一光滑丛同态 \(\overline\Phi:E/F\to H\) 使 \(\Phi=\overline\Phi q\)。
\end{proposition}
\begin{proof}
	由\cref{g1:prop:subbundle-adapted-frame}，局部可取 \(E\) 的标架 \(e_i\)，其前 \(k\) 个向量组成 \(F\) 的标架。把分量写成 \((a,b)\in\mathbb R^k\oplus\mathbb R^{r-k}\)。在每条纤维中，两个向量代表同一个商类，当且仅当它们的 \(b\) 分量相同。因此
	\[
		[e_i(p)(a,b)]\longmapsto(p,b)
	\]
	是商集合在 \(U_i\) 上到 \(U_i\times\mathbb R^{r-k}\) 的双射。

	在交叠上，两组适配标架的转换保持前 \(k\) 维子空间，故其矩阵必为
	\[
		g_{ij}=\begin{pmatrix}B_{ij}&C_{ij}\\0&D_{ij}\end{pmatrix}.
	\]
	其限制到 \(F_p\) 可逆，所以 \(B_{ij}\) 可逆；再由 \(\det g_{ij}=\det B_{ij}\det D_{ij}\ne0\)，知 \(D_{ij}\) 也可逆。商分量转换恰为 \(b_i=D_{ij}b_j\)。原余循环等式的右下块给出 \(D_{ij}D_{jk}=D_{ik}\)，于是粘合定理赋予商集合光滑向量丛结构。换用另一组适配标架时，同样的分块计算给出光滑可逆过渡，所以所得结构与补基选择无关。

	按这些局部图，\(q\) 的表达是 \((p,a,b)\mapsto(p,b)\)，故连续、光滑，并且开。开性可在每个 \(U_i\) 上验证，因为 Euclid 纤维投影把开集送到开集。具体地，对总空间开集 \(O\)，\(q(O)\) 在商丛限制到 \(U_i\) 中的部分，正是 \(O\cap E|_{U_i}\) 的局部投影像，因而开。连续开满射总是商映射：如果 \(q^{-1}(A)\) 开，则 \(A=q(q^{-1}(A))\) 开。因此这里的拓扑确实等于原集合商的商拓扑。核、满射性与秩都由同一局部公式得到。

	最后，逐纤维定义 \(\overline\Phi([v])=\Phi(v)\)。若换代表，差向量落在 \(F\)，被 \(\Phi\) 送到零，所以定义无关选择。在适配标架与 \(H\) 的任意局部标架下，\(\Phi\) 的前 \(k\) 列为零，矩阵写成 \((0,A_i)\)。\(\overline\Phi\) 的局部矩阵就是光滑矩阵 \(A_i\)，故它光滑。由于 \(q\) 满射，因子映射也唯一。
\end{proof}

对常秩丛同态 \(\Phi:E\to F\)，现在得到纤维中的第一同构定理的光滑版本：
\begin{equation}
	\label{g1:eq:bundle-first-isomorphism}
	E/\ker\Phi\ \simeq\ \operatorname{im}\Phi,
	\qquad [v_p]\longmapsto\Phi_p(v_p).
\end{equation}
核与像由\cref{g1:thm:constant-rank-bundle-morphism}成为子丛；商丛泛性质给出这个光滑映射。线性代数说明它在每条纤维上双射，而\cref{g1:prop:bundle-morphism-matrices}使其逆也光滑。因此这不是只在每个点分别存在的一族同构，而是总空间上的丛同构。

\begin{definition}
	\label{g1:def:short-exact-bundle-sequence}
	设 \(F,E,Q\) 是同一光滑流形上的有限秩实光滑向量丛，\(\iota:F\to E\)、\(p:E\to Q\) 为同底光滑丛同态。如果对每个底点 \(x\)，\(\iota_x\) 单射、\(p_x\) 满射且 \(\operatorname{im}\iota_x=\ker p_x\)，则称
	\[
		0\longrightarrow F\xrightarrow{\ \iota\ }E\xrightarrow{\ p\ }Q\longrightarrow0
	\]
	为向量丛的\term[duanzhengheconglie]{短正合列}[short exact sequence]。
\end{definition}

在这样的列中，\(\iota\) 将 \(F\) 识别为 \(E\) 的像子丛，而\eqref{g1:eq:bundle-first-isomorphism}把 \(Q\) 识别为 \(E/\iota(F)\)。局部补全 \(\iota(F)\) 的标架后，商分量可用补出的基向量提升，因而总能局部选择光滑右逆 \(Q|_U\to E|_U\)。不同邻域补出的子空间却未必彼此相同；把这些局部选择组织成整体选择，需要随后单位分解所提供的工具。即使整体选择存在，它通常也不由商丛本身唯一决定。

\subsection{标架丛与全局标架}
\label{g1:subsec:1201:04}

局部标架使向量有了分量，但一个纤维里并没有优先选定的一组基。若把所有可能的基都保留下来，换基矩阵就从计算中的辅助量变成这个新空间上的作用。这样做也能准确表达“能否在整个底空间上光滑选基”。

\begin{definition}
	\label{g1:def:frame-bundle}
	设 \(E\to M\) 为秩 \(r\) 的实光滑向量丛。定义
	\[
		\operatorname{Fr}(E)_p=\{u:\mathbb R^r\to E_p\mid u\text{ 是线性同构}\},\qquad
		\operatorname{Fr}(E)=\coprod_{p\in M}\operatorname{Fr}(E)_p.
	\]
	每个 \(u\) 等价于有序基 \((u(\varepsilon_1),\ldots,u(\varepsilon_r))\)。对任一局部标架 \(e_i\)，将 \(u\in\operatorname{Fr}(E)_p\) 唯一写成 \(u=e_i(p)A_i\)，其中 \(A_i\in\operatorname{GL}(r,\mathbb R)\)，并规定 \(u\mapsto(p,A_i)\) 为局部平凡化。若 \(e_j=e_i g_{ij}\)，则 \(A_i=g_{ij}A_j\)；用这些平凡化指定拓扑与光滑结构。所得空间称为 \(E\) 的\term[biaojiacong]{标架丛}[frame bundle]，其投影为 \(u\mapsto p\)。
\end{definition}

在 \(U_i\) 上选定参考标架 \(e_i\) 后，任意 \(u\in\operatorname{Fr}(E)_p\) 都唯一写成
\[
	u=e_i(p)A_i,\qquad A_i\in\operatorname{GL}(r,\mathbb R).
\]
因此得到 \(\operatorname{Fr}(E)|_{U_i}\simeq U_i\times\operatorname{GL}(r,\mathbb R)\)。若参考标架满足 \(e_j=e_i g_{ij}\)，则同一 \(u\) 的矩阵满足
\begin{equation}
	\label{g1:eq:frame-bundle-transition}
	A_i=g_{ij}A_j.
\end{equation}
这些过渡映射光滑可逆，并由余循环条件相容。它们赋予 \(\operatorname{Fr}(E)\) 唯一的拓扑与光滑结构，理由可直接沿用粘合证明的局部步骤：局部乘积中的开集通过过渡同胚相容，故组成拓扑基；不同基点用底空间的开邻域分离，同一基点的不同标架在同一局部乘积中分离；从底空间取可数平凡化子覆盖，再取各乘积的可数基，得到第二可数性。这里 \(\operatorname{GL}(r,\mathbb R)\) 是矩阵空间中行列式非零的开集，因此维数为 \(r^2\)，总空间维数为 \(\dim M+r^2\)。若底空间有边界，其边界仍是底空间边界的逆像。标架纤维是可逆线性映射的集合，一般不是向量空间，所以这个构造没有把 \(\operatorname{Fr}(E)\) 再当作向量丛。

\ProseParagraphBreak

群 \(\operatorname{GL}(r,\mathbb R)\) 在标架丛上有自然右作用
\begin{equation}
	\label{g1:eq:frame-bundle-right-action}
	u\cdot B=u\circ B.
\end{equation}
确为右作用，因为 \((u\cdot B)\cdot C=u\circ(BC)=u\cdot(BC)\)。局部表示为 \((p,A)\cdot B=(p,AB)\)，故光滑。作用在每条纤维内自由且传递：\(uB=u\) 蕴含 \(B=I_r\)；给定同一纤维内的 \(u,v\)，唯一使 \(uB=v\) 的矩阵为 \(B=u^{-1}\circ v\)。因此任意两组基都可通过唯一的矩阵相互转换，却没有哪一组基天然是单位元；只有选定参考标架后，纤维才被识别为矩阵群。这是以后讨论主丛时要保留的结构。

\begin{proposition}[全局标架与平凡丛]
	\label{g1:prop:global-frame-triviality}
	设 \(M\) 为 Hausdorff、第二可数光滑流形，可带边界，\(E\to M\) 为秩 \(r\) 实光滑向量丛。下列条件等价：
	\begin{equivalences}
		\item \(E\) 与乘积向量丛 \(M\times\mathbb R^r\) 光滑同构。
		\item 存在 \(r\) 个全局光滑截面，在每个点组成纤维的有序基。
		\item 标架丛投影 \(\operatorname{Fr}(E)\to M\) 有光滑截面。
	\end{equivalences}
	满足这些条件时，称 \(E\) 为\term[pingfancong]{平凡丛}[trivial bundle]。
\end{proposition}
\begin{proof}
	若给定丛同构 \(\Theta:M\times\mathbb R^r\to E\)，令 \(s_a(p)=\Theta(p,\varepsilon_a)\)。这些截面光滑，且由 \(\Theta\) 的逐纤维线性同构性构成基，故第一项推出第二项。

	给定第二项中的全局标架，定义 \(\sigma(p):\mathbb R^r\to E_p\) 为 \(\sigma(p)v=\sum_a v^a s_a(p)\)。在任一参考局部标架下，\(s_a\) 的分量组成一个光滑可逆矩阵，因此 \(\sigma\) 是标架丛的光滑截面。反过来，给定这样的 \(\sigma\)，令 \(s_a(p)=\sigma(p)\varepsilon_a\)，同一局部矩阵计算给出光滑全局标架，故第二、三项等价。

	最后，给定 \(\sigma\)，令 \(\Theta(p,v)=\sigma(p)v\)。它逐纤维为线性同构，局部表达为 \((p,v)\mapsto(p,A(p)v)\)，其中 \(A\) 光滑且可逆，所以 \(\Theta\) 与其逆 \((p,w)\mapsto(p,A(p)^{-1}w)\) 都光滑。这证明第一项。秩零时取空标架，所有构造与结论仍成立。
\end{proof}

Möbius 线丛没有处处非零的截面，因而连秩一所需的全局标架也没有。另一方面，切丛有全局标架并不要求流形具有一张全局坐标图。例如圆周上的切向量 \((-\sin t,\cos t)\) 随 \(e^{\mathrm{i}t}\) 光滑且处处非零，给出 \(\mathrm{T}S^1\) 的全局标架；圆周仍不能由一张 Euclid 开区间坐标图覆盖。标架提供的是每个点的线性坐标，流形坐标还须把邻域中的不同基点同时参数化，这两种要求应当区分。
% 纲要条目排为编号小节标题，并进入目录与书签；说明段仍保留于纲要注释中。

% !TeX root = ../../Book1.tex
% 纲要标题：### 12.2 几何中的多线性代数
% 纲要位置：几何学计划纲要.md，第 579 行。
% 正文状态：初稿；尚未正式审核。

% 纲要内容开始
% 
% * 张量积的泛性质与收缩
% * 对偶、对称幂及外幂
% * 交替性、楔积与分次符号
% * 行列式线、密度与体积
% * 不同标架下的分量及不变量
% 
% 纲要内容结束

\section{几何中的多线性代数}
\label{g1:sec:1202}

在一点处，方向属于切空间，测量方向的线性函数属于余切空间。内积却要同时接收两个方向，有向面积也要同时接收两个方向，而且交换这两个方向时会变号。要把这些量沿流形逐点组织起来，先须知道它们在一个向量空间里是什么，以及换基时怎样保持同一个对象。本节暂时固定有限维实向量空间；随点变化及光滑性将在\secref{g1:sec:1204}处理。

\ProseParagraphBreak

所需的线性代数基础是基、对偶、线性映射和行列式。回顾\cref{g1:def:cotangent-space,g1:prop:covector-pullback}：\(V^*=\operatorname{Hom}_{\mathbb R}(V,\mathbb R)\) 中的元素把向量读成实数；线性映射 \(A:V\to W\) 的对偶为 \(A^*\alpha=\alpha\circ A\)，方向从 \(W^*\) 返回 \(V^*\)。下面把一次线性观测推广到多次线性观测，同时保留不依赖基的定义。

\subsection{张量积的泛性质与收缩}
\label{g1:subsec:1202:01}

设 \(V,W,U\) 为实向量空间。对映射 \(b:V\times W\to U\)，若固定任一输入后，它对另一输入线性，就称该映射为双线性的。内积就是这样的映射，但它一般不是乘积向量空间 \(V\times W\) 上的线性映射：例如同时将两个输入加倍，输出会乘以 \(4\)。我们希望构造一个新的向量空间，使每个双线性运算都能写成“先作同一个规定的运算，再作一个普通线性映射”。

\begin{theorem}[张量积的泛性质]
	\label{g1:thm:tensor-product-universal}
	对有限维实向量空间 \(V,W\)，存在实向量空间 \(V\otimes W\) 及双线性映射
	\[
		\tau:V\times W\longrightarrow V\otimes W,\qquad (v,w)\longmapsto v\otimes w,
	\]
	使得对任意实向量空间 \(U\) 及双线性映射 \(b:V\times W\to U\)，存在唯一线性映射 \(\widetilde b:V\otimes W\to U\) 满足 \(\widetilde b(v\otimes w)=b(v,w)\)。这样的二元组 \((V\otimes W,\tau)\) 在保持规定映射的唯一线性同构意义下唯一。若 \(e_1,\ldots,e_m\) 和 \(f_1,\ldots,f_n\) 分别为 \(V,W\) 的基，则
	\[
		\{e_i\otimes f_j:1\le i\le m,\ 1\le j\le n\}
	\]
	为 \(V\otimes W\) 的基，因而 \(\dim(V\otimes W)=mn\)。
\end{theorem}

\[
	\begin{tikzcd}[column sep=large,row sep=large]
		V\times W \arrow[r,"\tau"] \arrow[dr,"b"'] & V\otimes W \arrow[d,dashed,"\exists!\,\widetilde b"] \\
		& U
	\end{tikzcd}
\]
图中的 \(\tau,b\) 是双线性映射，虚线箭头是线性映射。交换性只说从 \((v,w)\) 出发的两条路径给出同一个值；唯一性则说任何双线性运算都恰好对应一个这样的线性运算。

\begin{proof}
	先用一组基制造所需空间，再证明结果不依赖这次选择。取一个以符号 \(E_{ij}\) 为基的 \(mn\) 维实向量空间 \(T\)，定义
	\[
		\tau\left(\sum_i v^ie_i,\sum_j w^jf_j\right)=\sum_{i,j}v^iw^jE_{ij}.
	\]
	每个系数对两组变量分别线性，故 \(\tau\) 双线性。给定 \(b\)，规定 \(\widetilde b(E_{ij})=b(e_i,f_j)\)，唯一地定义一个线性映射。由双线性，
	\[
		b(v,w)=\sum_{i,j}v^iw^j b(e_i,f_j)=\widetilde b(\tau(v,w)).
	\]
	反之，任何满足这个等式的线性映射都必须在 \(E_{ij}=\tau(e_i,f_j)\) 上取上述值，故唯一。把 \(T\) 记为 \(V\otimes W\)，把 \(\tau(v,w)\) 记为 \(v\otimes w\)，便得到构造及所述基。

	若另一个二元组 \((T',\tau')\) 具有同样的性质，分别取测试映射 \(\tau'\) 和 \(\tau\)，得到唯一线性映射 \(F:T\to T'\)、\(G:T'\to T\)，满足 \(F\tau=\tau'\)、\(G\tau'=\tau\)。于是 \(GF\tau=\tau\)，而恒等映射也具有这个性质；唯一性给出 \(GF=\operatorname{id}_T\)。同理 \(FG=\operatorname{id}_{T'}\)。因此换一组基构造的空间与原空间有一个由规定元素决定的唯一同构。这正是构造不依赖基的含义，而不是说坐标列表本身不改变。零维情形中 \(T=\{0\}\)，双线性映射也只能为零，上述结论仍成立。
\end{proof}

称 \(V\otimes W\) 为 \(V,W\) 的\term[zhangliangji]{张量积}[tensor product]。元素 \(v\otimes w\) 称为纯张量；一般元素是纯张量的有限和，未必本身就是一个纯张量。双线性意味着
\[
	\begin{aligned}
		(v+v')\otimes w&=v\otimes w+v'\otimes w,\\
		v\otimes(w+w')&=v\otimes w+v\otimes w',\\
		(cv)\otimes w&=v\otimes(cw)=c(v\otimes w).
	\end{aligned}
\]
这些关系不允许把一切和都合成一个纯张量。例如在两个二维空间中，\(e_1\otimes f_1+e_2\otimes f_2\) 若等于 \(v\otimes w\)，其系数矩阵将既是单位矩阵，又是列向量 \((v^i)\) 与行向量 \((w^j)\) 的乘积；后者的列都在同一条直线上，不可能等于单位矩阵。

\ProseParagraphBreak

泛性质也规定了张量积上的映射。若 \(A:V\to V'\)、\(B:W\to W'\) 线性，则 \((v,w)\mapsto Av\otimes Bw\) 双线性，因而唯一给出 \(A\otimes B\)，满足 \((A\otimes B)(v\otimes w)=Av\otimes Bw\)。复合和恒等式在纯张量上成立，便在整个空间成立，因为纯张量张成该空间。同样，对三个或更多输入可按基构造多线性映射的通用空间；\((V\otimes W)\otimes U\) 与 \(V\otimes(W\otimes U)\) 都有这一性质，规定 \((v\otimes w)\otimes u\mapsto v\otimes(w\otimes u)\) 的同构使我们能够省略括号。记 \(V^{\otimes k}\) 为 \(k\) 个因子的张量积，并约定 \(V^{\otimes0}=\mathbb R\)。

\begin{proposition}[线性映射的张量表示]
	\label{g1:prop:tensor-hom-duality}
	设 \(V,W\) 为有限维实向量空间。映射
	\[
		\Phi:W\otimes V^*\longrightarrow\operatorname{Hom}_{\mathbb R}(V,W),\qquad
		\Phi(w\otimes\alpha)(v)=\alpha(v)w
	\]
	是无需选基的线性同构。若 \(e_1,\ldots,e_n\) 为 \(V\) 的任意一组基，\(e^1,\ldots,e^n\) 为对偶基，则
	\[
		\Phi^{-1}(L)=\sum_i L(e_i)\otimes e^i.
	\]
	此外，求值映射 \(V\to V^{**}\)、\(v\mapsto[\alpha\mapsto\alpha(v)]\) 为自然线性同构。
\end{proposition}
\begin{proof}
	\((w,\alpha)\mapsto[v\mapsto\alpha(v)w]\) 是双线性的，故由泛性质得到 \(\Phi\)。将给定 \(L\) 对应的张量作用于 \(v=\sum_i e^i(v)e_i\)，得到
	\[
		\Phi\left(\sum_iL(e_i)\otimes e^i\right)(v)=\sum_i e^i(v)L(e_i)=L(v).
	\]
	反过来，对纯张量 \(w\otimes\alpha\)，先经 \(\Phi\) 再用所列逆公式，得到
	\[
		\sum_i\alpha(e_i)w\otimes e^i
		=w\otimes\left(\sum_i\alpha(e_i)e^i\right)=w\otimes\alpha.
	\]
	由线性可知两个映射互逆。虽然逆公式用了基，它是一个无基映射的逆，所以换基所得张量相同。最后，非零向量 \(v\) 至少有一个非零坐标，相应的对偶基元素在 \(v\) 上不为零，故求值映射 \(V\to V^{**}\) 单射；两空间同维，因而是同构。它由求值定义，也没有基的选择。
\end{proof}

一般地，定义
\[
	T^r_s(V)=V^{\otimes r}\otimes(V^*)^{\otimes s},\qquad r,s\ge0,
\]
称其元素为 \((r,s)\) 型\term[zhangliang]{张量}[tensor]。\symidx{\(T^r_s(V)\)：\((r,s)\) 型张量空间}它自然作用于 \(r\) 个协变向量和 \(s\) 个向量：纯张量 \(v_1\otimes\cdots\otimes v_r\otimes\alpha_1\otimes\cdots\otimes\alpha_s\) 对输入 \((\xi_1,\ldots,\xi_r;w_1,\ldots,w_s)\) 的值是
\[
	\prod_{a=1}^r\xi_a(v_a)\prod_{b=1}^s\alpha_b(w_b).
\]
这不只是举出一种作用方式，而是给出了张量空间与相应多线性函数空间的同构：选基后，任一多线性函数都由它在各输入基上的值唯一决定，逐项展开便得到张量基的唯一系数表示。这里把向量视为对协变向量的线性函数，使用的正是 \(V\simeq V^{**}\)。因此上标 \(r\) 数的是向量因子，实际输入却是 \(r\) 个协变向量；不要把两者混淆。

\begin{proposition}[收缩与迹]
	\label{g1:prop:tensor-contraction}
	设 \(V\) 为有限维实向量空间，\(r,s\ge1\)。指定一个向量因子的位置 \(a\) 和一个对偶因子的位置 \(b\)，在纯张量上将它们求值、删除这两个因子，并保留其余因子的先后次序，得到唯一线性映射
	\[
		C^a_b:T^r_s(V)\longrightarrow T^{r-1}_{s-1}(V).
	\]
	这一运算不依赖基，称为\term[shousuo]{收缩}[contraction]。特别地，经 \(V\otimes V^*\simeq\operatorname{End}(V)\) 识别，\(C^1_1\) 给出线性映射 \(L\) 的迹：
	\[
		\operatorname{tr}L=\sum_i e^i(L(e_i)).
	\]
	右端在任意基及其对偶基中均取同一值。
\end{proposition}
\begin{proof}
	具体规定为
	\[
		C^a_b(v_1\otimes\cdots\otimes v_r\otimes\alpha_1\otimes\cdots\otimes\alpha_s)
		=\alpha_b(v_a)\,v_1\otimes\cdots\widehat{v_a}\cdots\otimes v_r
		\otimes\alpha_1\otimes\cdots\widehat{\alpha_b}\cdots\otimes\alpha_s,
	\]
	帽号表示省略。右端对原来的每个输入分别线性，多因子张量积的泛性质便给出唯一线性映射。这个定义只使用内在的求值 \(\alpha_b(v_a)\)，没有使用基。对于 \(r=s=1\)，\cref{g1:prop:tensor-hom-duality}给出 \(L=\Phi(\sum_iL(e_i)\otimes e^i)\)，收缩后正是所列求和。因此矩阵对角元之和是一个无基运算的坐标表达，换基不改变它。
\end{proof}

收缩为什么只配对一个向量因子和一个对偶因子？因为 \(\alpha(v)\) 已有定义，而两个向量之间并没有预定的实数值。若另选内积，当然可以先用内积把一个向量变成协变向量再收缩，但这时结果包含了所选内积的信息。例如一个双线性形式属于 \(T^0_2(V)\)；仅凭这个张量本身没有规定的“迹”，要用内积的逆张量把它的两个指标配对，才能谈相对于该内积的迹。

\subsection{对偶、对称幂及外幂}
\label{g1:subsec:1202:02}

有些多线性观测不区分输入的排列，例如内积；另一些观测在交换两个输入时改变符号，例如二维有向面积。与其在每次计算时附加相同条件，不如把这些条件直接加入张量积的构造。这里的商空间只是在说：若所有允许的观测都必须把某种差别读成零，就把这种差别规定为零。

\begin{definition}
	\label{g1:def:symmetric-exterior-powers}
	设 \(V\) 为有限维实向量空间，\(k\ge2\)。其 \(k\) 次\term[duichengmi]{对称幂}[symmetric power]为
	\[
		\operatorname{Sym}^kV=V^{\otimes k}/R_{\mathrm{sym}},
	\]
	其中 \(R_{\mathrm{sym}}\) 由所有纯张量与其任意排列之差张成。纯张量的像记为 \(v_1\odot\cdots\odot v_k\)。其 \(k\) 次\term[waimi]{外幂}[exterior power]为
	\[
		\Lambda^kV=V^{\otimes k}/R_{\mathrm{alt}},
	\]
	其中 \(R_{\mathrm{alt}}\) 由任意两个位置输入相同的纯张量张成；纯张量的像记为 \(v_1\wedge\cdots\wedge v_k\)。约定 \(\operatorname{Sym}^0V=\Lambda^0V=\mathbb R\)，\(\operatorname{Sym}^1V=\Lambda^1V=V\)。
\end{definition}

外幂的定义立刻给出重复输入为零。将两个相同位置的输入都换成 \(u+v\) 并展开，两个重复项各自为零，余下两项相加为零。因此交换任意两个位置会变号。这也说明：一个多线性映射在重复输入时为零，当且仅当交换任意两个输入时变号；反向推导使用 \(a=-a\Rightarrow a=0\)，这里实数域的条件是有用的。这样的映射称为\term[jiaotixing]{交替的}[alternating]。

\ProseParagraphBreak

商的定义给出相应的泛性质：对称 \(k\) 线性映射唯一地通过 \(\operatorname{Sym}^kV\) 分解，交替 \(k\) 线性映射唯一地通过 \(\Lambda^kV\) 分解。理由是张量积已经把多线性映射变成线性映射；对称或交替条件恰好说这个线性映射消去相应关系子空间，因而能且只能下降到商。特别地，\((\Lambda^kV)^*\) 就是所有交替 \(k\) 线性实函数的空间。

\begin{proposition}[对称幂与外幂的基]
	\label{g1:prop:symmetric-exterior-bases}
	设 \(V\) 为 \(n\) 维实向量空间，基为 \(e_1,\ldots,e_n\)。对 \(k\ge1\)，
	\[
		e_1^{\odot a_1}\odot\cdots\odot e_n^{\odot a_n},
		\qquad a_i\ge0,\quad \sum_i a_i=k
	\]
	是 \(\operatorname{Sym}^kV\) 的一组基，而
	\[
		e_{i_1}\wedge\cdots\wedge e_{i_k},\qquad i_1<\cdots<i_k
	\]
	是 \(\Lambda^kV\) 的一组基。因此，当 \(n\ge1\) 时，
	\[
		\dim\operatorname{Sym}^kV=\binom{n+k-1}{k},\qquad
		\dim\Lambda^kV=\binom nk,
	\]
	其中 \(k>n\) 时外幂为零；零次幂均为一维，零维空间的正次幂均为零。
\end{proposition}
\begin{proof}
	展开各输入的基表示，并在对称幂中重新排列因子，得到第一组元素张成。为证独立，把每个 \(e_i\) 对应于形式变量 \(X_i\)，规定
	\[
		(v_1,\ldots,v_k)\longmapsto
		\prod_{\ell=1}^k\left(\sum_i e^i(v_\ell)X_i\right).
	\]
	它是取值于 \(k\) 次齐次多项式空间的对称多线性映射，故下降到 \(\operatorname{Sym}^kV\)，并把所列元素送到各不相同的单项式 \(X_1^{a_1}\cdots X_n^{a_n}\)。形式多项式的系数表示唯一，因此这些像线性无关，原元素也线性无关。非负整数解 \(a_1+\cdots+a_n=k\) 的数量是 \(\binom{n+k-1}{n-1}\)：把 \(k\) 个记号排成一行，再用 \(n-1\) 个隔板划分为各组，允许空组，即得到这个计数。

	在外幂中也先展开基表示；重复指标项为零，不重复项可按指标递增排列，代价只是排列的符号，所以第二组元素张成。对每个严格递增指标组 \(I=(i_1,\ldots,i_k)\)，定义
	\[
		D_I(v_1,\ldots,v_k)=\det\bigl[e^{i_a}(v_b)\bigr]_{a,b=1}^k.
	\]
	行列式对各列线性，重复列时为零，故它下降为 \(\Lambda^kV\) 上的线性函数。在所列的递增基元素 \(e_{j_1}\wedge\cdots\wedge e_{j_k}\) 上，它恰在 \(I=J\) 时取 \(1\)，其余取 \(0\)。逐个应用这些函数，就知所列元素独立，其数量为 \(\binom nk\)。若 \(k>n\)，每个展开项必有重复指标，所以全部为零。
\end{proof}

现在把外幂与交替形式准确地对应起来。对 \(\alpha_i\in V^*\) 和 \(v_j\in V\)，规定
\begin{equation}
	\label{g1:eq:exterior-determinant-pairing}
	\left\langle\alpha_1\wedge\cdots\wedge\alpha_k,\,
	v_1\wedge\cdots\wedge v_k\right\rangle
	=\det[\alpha_i(v_j)].
\end{equation}
行列式对两组输入分别多线性并交替，故由两次泛性质，这给出 \(\Lambda^k(V^*)\) 与 \(\Lambda^kV\) 之间的双线性配对。对偶基的递增楔积与向量基的递增楔积在此配对下互为对偶，因而得到自然同构
\[
	\Lambda^k(V^*)\simeq(\Lambda^kV)^*
	\simeq\{\text{\(V\) 上的交替 \(k\) 线性实函数}\}.
\]
这里“自然”具体指不需选基，而且与线性映射相容：若 \(A:V\to W\)，令 \(\Lambda^kA(v_1\wedge\cdots\wedge v_k)=Av_1\wedge\cdots\wedge Av_k\)，则配对的两边都给出同一个行列式 \(\det[\alpha_i(Av_j)]\)，所以 \((\Lambda^kA)^*\) 对应于 \(\Lambda^k(A^*)\)。以后简写 \(\Lambda^kV^*=\Lambda^k(V^*)\)，称其中元素为 \(V\) 上的 \(k\) 次形式。

\ProseParagraphBreak

对称幂有相似的对偶识别，但归一化不同。规定
\[
	\left\langle\alpha_1\odot\cdots\odot\alpha_k,\,
	v_1\odot\cdots\odot v_k\right\rangle
	=\frac1{k!}\sum_{\sigma\in S_k}\prod_{i=1}^k\alpha_i(v_{\sigma(i)}).
\]
这里 \(S_k\) 是 \(k\) 个位置的全排列集合。右端对两组输入都对称多线性，因而给出 \(\operatorname{Sym}^k(V^*)\to(\operatorname{Sym}^kV)^*\)。在上一命题的单项式基上，两个多重指标不同时配对为零；都为 \((a_1,\ldots,a_n)\) 时配对为 \(a_1!\cdots a_n!/k!\)，是非零实数。因此该映射也是无基的同构，并把 \(\operatorname{Sym}^k(V^*)\) 识别为对称 \(k\) 线性实函数。它与拉回相容，因为将各 \(v_j\) 换成 \(Av_j\)，等价于将各 \(\alpha_i\) 换成 \(A^*\alpha_i\)。例如
\[
	(\alpha\odot\beta)(u,v)=\frac{\alpha(u)\beta(v)+\alpha(v)\beta(u)}2,\qquad
	(\alpha\wedge\beta)(u,v)=\alpha(u)\beta(v)-\alpha(v)\beta(u).
\]
这两个公式中的系数差别应记住；本书的有向面积形式采用第二个公式，没有额外的 \(1/2\)。

\subsection{交替性、楔积与分次符号}
\label{g1:subsec:1202:03}

将两个外幂中的代表张量直接首尾相接，便得到乘法
\[
	\Lambda^pV\times\Lambda^qV\longrightarrow\Lambda^{p+q}V.
\]
它确实不依赖代表：若一个因子改变了一个重复输入生成的关系项，相接后的张量仍有重复输入，在目标商空间中为零。于是乘法由纯元素上的
\[
	(v_1\wedge\cdots\wedge v_p)\wedge(w_1\wedge\cdots\wedge w_q)
	=v_1\wedge\cdots\wedge v_p\wedge w_1\wedge\cdots\wedge w_q
\]
决定，称为\term[xieji]{楔积}[wedge product]。在 \(V^*\) 上也有同一构造。把它视为交替形式的乘法时，必须说明它对一般输入怎样求值。

\ProseParagraphBreak

先对任意 \(k\) 线性实函数 \(T\) 定义\term[jiaotihua]{交替化}[alternation]
\[
	\operatorname{Alt}_kT(v_1,\ldots,v_k)
	=\frac1{k!}\sum_{\sigma\in S_k}\operatorname{sgn}(\sigma)
	T(v_{\sigma(1)},\ldots,v_{\sigma(k)}).
\]
交换两个输入会将每项对应到一个符号相反的排列，所以结果交替；若 \(T\) 已交替，每一项连同符号都等于 \(T\)，故 \(\operatorname{Alt}_kT=T\)。于是 \(\operatorname{Alt}_k\) 是投影到交替形式空间的线性算子。由行列式展开，\(\alpha_1\wedge\cdots\wedge\alpha_k\) 作为形式等于 \(k!\operatorname{Alt}_k(\alpha_1\otimes\cdots\otimes\alpha_k)\)，不能把这里的阶乘删去。

\begin{proposition}[楔积的代数性质]
	\label{g1:prop:wedge-algebra}
	设 \(V\) 为有限维实向量空间，\(\alpha\in\Lambda^pV^*\)、\(\beta\in\Lambda^qV^*\)、\(\gamma\in\Lambda^rV^*\)。在行列式配对\eqref{g1:eq:exterior-determinant-pairing}的约定下，
	\begin{equation}
		\label{g1:eq:wedge-shuffle}
		\begin{split}
			(\alpha\wedge\beta)(v_1,\ldots,v_{p+q})
			&=\sum_{\sigma\in\operatorname{Sh}(p,q)}\operatorname{sgn}(\sigma)
			\alpha(v_{\sigma(1)},\ldots,v_{\sigma(p)})
			\beta(v_{\sigma(p+1)},\ldots,v_{\sigma(p+q)})\\
			&=\frac1{p!\,q!}\sum_{\sigma\in S_{p+q}}\operatorname{sgn}(\sigma)
			\alpha(v_{\sigma(1)},\ldots,v_{\sigma(p)})
			\beta(v_{\sigma(p+1)},\ldots,v_{\sigma(p+q)}).
		\end{split}
	\end{equation}
	其中 \(\operatorname{Sh}(p,q)\) 由满足 \(\sigma(1)<\cdots<\sigma(p)\)、\(\sigma(p+1)<\cdots<\sigma(p+q)\) 的排列组成。楔积双线性，并满足
	\[
		(\alpha\wedge\beta)\wedge\gamma=\alpha\wedge(\beta\wedge\gamma),\qquad
		\alpha\wedge\beta=(-1)^{pq}\beta\wedge\alpha.
	\]
	零次形式是实数，\(1\) 是乘法单位。对任意有限维实向量空间 \(W\) 及线性映射 \(A:V\to W\)，按
	\[
		(A^*\eta)(v_1,\ldots,v_k)=\eta(Av_1,\ldots,Av_k)
	\]
	定义形式的拉回，则 \(A^*(\eta\wedge\zeta)=A^*\eta\wedge A^*\zeta\)，且复合的拉回按相反顺序复合。
\end{proposition}
\begin{proof}
	先说明公式与已经定义的乘法一致。由于一形式的楔积张成各次外幂，只须取 \(\alpha=\alpha_1\wedge\cdots\wedge\alpha_p\)、\(\beta=\beta_1\wedge\cdots\wedge\beta_q\)。相接所得形式在 \(v_1,\ldots,v_{p+q}\) 上的值，是以 \(\alpha_i,\beta_j\) 为行的 \((p+q)\) 阶行列式。按哪些 \(p\) 列被分配给前 \(p\) 行将行列式展开归组，每组选列对应唯一的 \((p,q)\) 洗牌排列；组内的两个行列式正是 \(\alpha\) 和 \(\beta\) 的值，分组选列的符号就是该排列的符号。这得到第一式。任意全排列又唯一分为一个洗牌排列和两个块内排列。块内排列的符号会分别被 \(\alpha,\beta\) 的交替性抵消，所以一个洗牌项在全排列和中重复 \(p!q!\) 次，得到第二式。

	结合性也能直接从相接代表看出：两种加括号方式都给出同一个首尾相接的张量。用形式公式核对，展开任一边时，都是将所有输入分为内部递增、大小分别为 \(p,q,r\) 的三个有序块；每种分法出现一次，系数都是将原输入排列成三个块的符号。因此公式中的归一化与结合性一致。交换前后两个块时，前块的每一个位置都须越过后块的每一个位置，共 \(pq\) 次交换，故产生 \((-1)^{pq}\)，这就给出所列交换律。零次时没有输入，所有公式退化为实数乘法。

	拉回仍为交替多线性函数。将洗牌公式的每个 \(v_i\) 换成 \(Av_i\)，两边分别就是 \(A^*(\eta\wedge\zeta)\) 和 \(A^*\eta\wedge A^*\zeta\)。若再有线性映射 \(B:W\to Z\)，则两种拉回都把输入送成 \(BAv_i\)，所以 \((B\circ A)^*=A^*\circ B^*\)。
\end{proof}

这里的交换律称为分次交换律：符号由两个次数的乘积决定。特别地，奇数次形式满足 \(\alpha\wedge\alpha=0\)，因为它等于自身的相反数；偶数次形式没有这一结论。例如在四维空间中，
\[
	(e^1\wedge e^2+e^3\wedge e^4)^2
	=2e^1\wedge e^2\wedge e^3\wedge e^4\ne0.
\]
右边的两项相加，是因为交换两个二次块经过 \(2\cdot2\) 次换位，符号为正。由此也可见，“交替”指的是一个形式内部交换向量输入的规律，不能理解为任意两个形式相乘都要变号。

\begin{definition}
	\label{g1:def:interior-product}
	设 \(V\) 为实向量空间，\(k\ge1\) 为整数，\(v\in V\)，\(\alpha\in\Lambda^kV^*\)。定义\term[charu]{插入}[interior product]
	\[
		(\iota_v\alpha)(v_1,\ldots,v_{k-1})=\alpha(v,v_1,\ldots,v_{k-1}).
	\]
	它属于 \(\Lambda^{k-1}V^*\)，并约定 \(\iota_v c=0\) 对零次形式 \(c\in\mathbb R\) 成立。
\end{definition}

插入固定了一个输入，所以次数降低 \(1\)。输入固定在最左端是本书的约定；把它放在末端会产生与次数有关的符号，必须相应修改下面的公式。

\begin{proposition}[插入的乘法公式]
	\label{g1:prop:interior-product-rule}
	设 \(V\) 为有限维实向量空间，\(v\in V\)，\(\alpha\in\Lambda^pV^*\)、\(\beta\in\Lambda^qV^*\)。则
	\[
		\iota_v(\alpha\wedge\beta)
		=(\iota_v\alpha)\wedge\beta+(-1)^p\alpha\wedge(\iota_v\beta).
	\]
	因此 \(\iota_v\) 是外代数上次数为 \(-1\) 的分次导子。另有 \(\iota_v^2=0\)；若 \(A:V\to W\) 线性，则 \(\iota_v(A^*\eta)=A^*(\iota_{Av}\eta)\)。
\end{proposition}
\begin{proof}
	先取一形式 \(\eta_1,\ldots,\eta_k\)。把 \(\eta_1\wedge\cdots\wedge\eta_k\) 在 \((v,v_1,\ldots,v_{k-1})\) 上的行列式沿第一列展开，得到
	\[
		\iota_v(\eta_1\wedge\cdots\wedge\eta_k)
		=\sum_{j=1}^k(-1)^{j-1}\eta_j(v)
		\,\eta_1\wedge\cdots\widehat{\eta_j}\cdots\wedge\eta_k.
	\]
	若前 \(p\) 个因子来自 \(\alpha\)，后 \(q\) 个来自 \(\beta\)，前 \(p\) 项合成 \((\iota_v\alpha)\wedge\beta\)；第 \(p+j\) 项的符号为 \((-1)^{p+j-1}=(-1)^p(-1)^{j-1}\)，后 \(q\) 项因而合成第二项。纯楔积张成所有形式，按双线性便得到一般公式。\(p=0\) 或 \(q=0\) 时公式就是插入对常数乘法的线性性质。两次插入同一向量得到 \(\alpha(v,v,\ldots)=0\)，故 \(\iota_v^2=0\)；最后一式的两边在任意输入上都等于 \(\eta(Av,Av_1,\ldots)\)。
\end{proof}

\subsection{行列式线、密度与体积}
\label{g1:subsec:1202:04}

对 \(n\) 维空间，最高次外幂只有一维，却保存着换基的行列式。若把一组有序基写作行向量 \(e=(e_1,\ldots,e_n)\)，换基约定为
\[
	e'=eA,\qquad e'_j=\sum_i e_i A^i{}_j,\qquad A\in\operatorname{GL}(n,\mathbb R),
\]
那么按多线性展开并消去重复项，有
\[
	e'_1\wedge\cdots\wedge e'_n=(\det A)e_1\wedge\cdots\wedge e_n.
\]
行列式的符号记录定向，绝对值记录平行多面体的体积缩放。要把这两个作用区分开，就会得到不同的一维空间。

\begin{definition}
	\label{g1:def:determinant-density-line}
	设 \(V\) 为 \(n\) 维实向量空间。其\term[hanglieshixian]{行列式线}[determinant line]为 \(\det V=\Lambda^nV\)。最高次形式组成的是另一条线 \(\det(V^*)=\Lambda^n(V^*)\)，经行列式配对自然同构于 \((\det V)^*\)。

	\(V\) 上的一个\term[midu]{密度}[density]是有序基集合上的实值函数 \(\mu\)，满足对每组有序基 \(e\) 和每个可逆实矩阵 \(A\)，
	\[
		\mu(eA)=|\det A|\,\mu(e).
	\]
	全体密度按逐点加法与数乘构成实向量空间，记为 \(\mathcal D(V)\)，也记为 \(|\Lambda^nV^*|\)。若 \(\mu(e)>0\) 对一组、从而对每组有序基成立，称 \(\mu\) 为正密度。
\end{definition}

密度线的记号 \(|\Lambda^nV^*|\) 是一种惯用记号，其中的元素仍可取负值，不能把它理解为只有非负元素的集合。固定一组基 \(e\) 后，任意指定实数 \(c\) 都唯一地定义密度：若 \(f=eA\)，令 \(\mu(f)=|\det A|c\)。两个连续换基的行列式绝对值相乘，故这个定义满足规定关系。因此 \(\mu\mapsto\mu(e)\) 是 \(\mathcal D(V)\to\mathbb R\) 的线性同构，密度线确为一维；正密度只构成其中一个正半轴，不是线性子空间。

\ProseParagraphBreak

最高次形式 \(\omega\in\Lambda^nV^*\) 给出密度 \(|\omega|\)，定义为 \(|\omega|(e)=|\omega(e_1,\ldots,e_n)|\)。这是因为 \(\omega(eA)=\det A\,\omega(e)\)。但是 \(\omega\mapsto|\omega|\) 不是线性同构：\(\omega\) 和 \(-\omega\) 给出同一密度，而两个相反形式之和为零。对 \(n\ge1\)，定向可由有序基按正行列式换基划成的两类中选择一类来指定。选定定向以后，令 \(s(e)=1\) 对正向基成立，\(s(e)=-1\) 对负向基成立，则
\[
	\omega\longmapsto\bigl[e\longmapsto s(e)\omega(e)\bigr]
\]
才是最高次形式线到密度线的线性同构。正密度与所选定向一起唯一确定一个在正向基上取正值的最高次形式；仅有正密度时，两个可能的定向仍未选定。

\ProseParagraphBreak

这些对象的换基行为可以集中写成下表。设 \(e^1,\ldots,e^n\) 为 \(e\) 的对偶基，\(\varepsilon=e^1\wedge\cdots\wedge e^n\)，\(\varepsilon'=e'^1\wedge\cdots\wedge e'^n\)。
\[
	\begin{array}{c|c|c}
		\text{空间} & \text{由基产生的生成元} & \text{换到 }e'=eA\\ \hline
		\det V & e_1\wedge\cdots\wedge e_n & (\det A)\,e_1\wedge\cdots\wedge e_n\\
		\det(V^*) & \varepsilon & \varepsilon'=(\det A)^{-1}\varepsilon\\
		\mathcal D(V) & |\varepsilon| & |\varepsilon'|=|\det A|^{-1}|\varepsilon|
	\end{array}
\]
生成元的变换与同一对象的系数变换互为倒数。例如 \(\mu=m|\varepsilon|=m'|\varepsilon'|\)，便有 \(m'=|\det A|m\)，且 \(m=\mu(e)\)。这一点在把密度逐点粘成丛时会再次使用。

\ProseParagraphBreak

若 \(L:V\to W\) 是同维实向量空间之间的线性同构，则 \((L^*\mu)(e)=\mu(Le)\) 定义了密度的线性拉回，并与复合相容，因为 \(L(eA)=(Le)A\)。这里要求 \(L\) 可逆，才能直接把有序基送成有序基。对不可逆的同维线性映射，可以另用绝对行列式约定得到零密度，但不能据此把密度对光滑映射的拉回与普通微分形式的拉回不加区别地处理；行列式绝对值在变号处有不同的光滑性问题。

\begin{proposition}[内积的体积密度]
	\label{g1:prop:gram-volume-density}
	设 \(V\) 为有限维实向量空间，\(g\) 为正定内积。对有序基 \(v=(v_1,\ldots,v_n)\)，公式
	\[
		\mu_g(v_1,\ldots,v_n)=\sqrt{\det[g(v_i,v_j)]_{i,j=1}^n}
	\]
	定义一个正密度。若另给定向，则存在唯一在每组正向正交单位基上取值 \(1\) 的最高次形式 \(\operatorname{vol}_g\)，并有 \(|\operatorname{vol}_g|=\mu_g\)。对任意 \(n\) 个向量，同一平方根公式在它们线性相关时取零，在线性无关时就是相应平行多面体的体积。
\end{proposition}
\begin{proof}
	先说明根号内为正。逐个从基向量中减去沿先前向量的正交投影，再除以自身的正长度，得到正交单位基 \(u_1,\ldots,u_n\)；每一步的余向量不为零，因为原向量组线性无关。将 \(v\) 写成 \(uB\)，其内积矩阵就是 \(B^{\mathsf T}B\)，行列式为 \((\det B)^2>0\)。换基 \(v'=vA\) 后，矩阵变成 \(A^{\mathsf T}[g(v_i,v_j)]A\)，所以平方根乘以 \(|\det A|\)。这正是密度的变换规则，并证明了正性。

	选一组正向正交单位基 \(u\)，定义 \(\operatorname{vol}_g\) 在它上取值 \(1\)。其他正向正交单位基的变换矩阵 \(Q\) 满足 \(Q^{\mathsf T}Q=I\) 及 \(\det Q>0\)，因而 \(\det Q=1\)；所以定义不依赖所选正向正交单位基。最高次形式由在一组基上的值唯一决定，故唯一。对 \(v=uB\)，有 \(\operatorname{vol}_g(v)=\det B\)，而 \(\mu_g(v)=|\det B|\)，得到绝对值关系。若给定的 \(n\) 个向量线性相关，其坐标矩阵 \(B\) 奇异，Gram 行列式为零；否则 \(|\det B|\) 是正交坐标中平行多面体的体积。这也解释了该公式的几何意义。零维时采用空行列式等于 \(1\) 的约定。
\end{proof}

例如反射一组正交单位基中的一个向量，最高次体积形式的值由 \(1\) 变为 \(-1\)，体积密度的值仍为 \(1\)。在没有整体定向的几何问题里，长度和体积仍应是可测量的正量；密度保留的正是这一部分信息。

\subsection{不同标架下的分量及不变量}
\label{g1:subsec:1202:05}

前面的对象都先于基而存在，实际计算却需要系数。本小节把这两件事接起来。固定 \(n\) 维实向量空间的一组基 \(e\)，也称一个标架，并始终用 \(e'=eA\) 表示新基；矩阵的第 \(j\) 列是第 \(j\) 个新基向量在旧基中的分量。若 \(v\) 的两组分量列向量为 \([v]_e,[v]_{e'}\)，则 \(e[v]_e=eA[v]_{e'}\)，所以
\[
	[v]_{e'}=A^{-1}[v]_e.
\]
对偶基必须满足 \(e'^i(e'_j)=\delta^i_j\)，故 \(e'^i=\sum_a(A^{-1})^i{}_a e^a\)。但同一协变向量的系数是它在各基向量上的值，因而满足
\[
	[\alpha]_{e'}=A^{\mathsf T}[\alpha]_e.
\]
对偶基的变化和协变向量系数的变化是两件互逆的事。

\begin{proposition}[张量的分量变换]
	\label{g1:prop:tensor-component-change}
	设 \(V\) 为 \(n\) 维实向量空间，\(T\in T^r_s(V)\)，有序基满足 \(e'=eA\)。在各自的张量基中写
	\[
		T=\sum T^{i_1\cdots i_r}{}_{j_1\cdots j_s}
		\,e_{i_1}\otimes\cdots\otimes e_{i_r}\otimes e^{j_1}\otimes\cdots\otimes e^{j_s},
	\]
	以及对应的带撇表达，则
	\begin{equation}
		\label{g1:eq:tensor-change-components}
		T'^{i_1\cdots i_r}{}_{j_1\cdots j_s}
		=\sum_{\substack{a_1,\ldots,a_r\\ b_1,\ldots,b_s}}
		(A^{-1})^{i_1}{}_{a_1}\cdots(A^{-1})^{i_r}{}_{a_r}
		A^{b_1}{}_{j_1}\cdots A^{b_s}{}_{j_s}
		T^{a_1\cdots a_r}{}_{b_1\cdots b_s}.
	\end{equation}
	收缩的分量由一对上下指标取相同值后求和得到；这与换基相容。尤其所有输入都与张量配对之后得到的实数，与标架无关。
\end{proposition}
\begin{proof}
	以多线性函数的观点看，原系数恰为
	\[
		T^{i_1\cdots i_r}{}_{j_1\cdots j_s}
		=T(e^{i_1},\ldots,e^{i_r};e_{j_1},\ldots,e_{j_s}).
	\]
	新系数把这些输入全部换成新基及新对偶基。分别代入 \(e'_j=\sum_b e_bA^b{}_j\) 和 \(e'^i=\sum_a(A^{-1})^i{}_a e^a\)，再逐项使用多线性，即得到\eqref{g1:eq:tensor-change-components}。

	对纯张量收缩时，\(e^j(e_i)=\delta_i^j\)，故分量规则就是指定的上下指标相等后求和。若先换基再收缩，被收缩的一对指标带来的因子满足
	\[
		\sum_k(A^{-1})^k{}_a A^b{}_k=\delta^b_a,
	\]
	两矩阵相消，其余因子恰是低两阶张量的换基公式。因此先收缩或先换基给出同一张量。所有输入的分量也按相应规则变换，完全配对时每个因子都有逆因子与之抵消；所得数就是原来无基定义的多线性函数值。
\end{proof}

例如线性映射 \(L:V\to V\) 的矩阵满足 \(L'=A^{-1}LA\)，而双线性形式 \(g\) 的矩阵满足 \(G'=A^{\mathsf T}GA\)。它们的两种变换公式不同，是因为 \(L\) 有一个向量因子和一个对偶因子，\(g\) 有两个对偶因子。对任意 \(v,w\)，
\[
	[v]_{e'}^{\mathsf T}G'[w]_{e'}
	=[v]_e^{\mathsf T}G[w]_e;
\]
这是同一个数 \(g(v,w)\)。又如
\[
	L=\begin{pmatrix}2&0\\0&3\end{pmatrix},\qquad
	A=\begin{pmatrix}1&1\\0&1\end{pmatrix}
	\quad\Longrightarrow\quad
	L'=\begin{pmatrix}2&-1\\0&3\end{pmatrix}.
\]
新矩阵不再对角，但迹仍为 \(5\)，其线性映射也仍是原来那个映射。矩阵形状的改变与几何对象的改变须分开判断。

\begin{example}[极坐标中的度量与面积]
	\label{g1:exm:tensor-polar-components}
	在平面原点之外的一个角扇区上，取 \(x=r\cos\theta\)、\(y=r\sin\theta\)，其中 \(r>0\)。\cref{g1:exm:tangent-polar-coordinates}已经给出切基的变换；固定其中一点，写成矩阵就是
	\[
		(\partial_r,\partial_\theta)=(\partial_x,\partial_y)A,\qquad
		A=\begin{pmatrix}\cos\theta&-r\sin\theta\\ \sin\theta&r\cos\theta\end{pmatrix}.
	\]
	此处每个向量都在同一个切空间中。Euclid 内积在直角坐标基中的矩阵为 \(I\)，所以在极坐标基中为
	\[
		A^{\mathsf T}A=\begin{pmatrix}1&0\\0&r^2\end{pmatrix},\qquad
		g=\mathrm{d}r\otimes\mathrm{d}r+r^2\,\mathrm{d}\theta\otimes\mathrm{d}\theta.
	\]
	因此 \(a\partial_r+b\partial_\theta\) 的长度平方为 \(a^2+r^2b^2\)。出现 \(r^2\) 是因为 \(\partial_\theta\) 表示每单位角度的速度，其长度为 \(r\)，并不表示 Euclid 内积换成了另一种内积。

	对偶变换也可直接求导：
	\[
		\mathrm{d}x=\cos\theta\,\mathrm{d}r-r\sin\theta\,\mathrm{d}\theta,\qquad
		\mathrm{d}y=\sin\theta\,\mathrm{d}r+r\cos\theta\,\mathrm{d}\theta.
	\]
	使用楔积的双线性和一次形式交换变号，得到
	\[
		\mathrm{d}x\wedge\mathrm{d}y
		=r\,\mathrm{d}r\wedge\mathrm{d}\theta,\qquad
		\mu_g=|\mathrm{d}x\wedge\mathrm{d}y|
		=r\,|\mathrm{d}r\wedge\mathrm{d}\theta|.
	\]
	度量的 Gram 行列式为 \(r^2\)，平方根也给出密度系数 \(r\)，与行列式换基结果一致。原点处 \(A\) 不可逆，所以这些公式是在极坐标有效的区域内比较两组标架；不能用退化的极坐标基描述原点处的张量。
\end{example}
% 纲要条目排为编号小节标题，并进入目录与书签；说明段仍保留于纲要注释中。

% !TeX root = ../../Book1.tex
% 纲要标题：### 12.3 光滑单位分解与延拓
% 纲要位置：几何学计划纲要.md，第 587 行。
% 正文状态：初稿；尚未正式审核。

% 纲要内容开始
% 
% * 光滑截断函数和局部有限细化
% * 单位分解的完整构造
% * 截面延拓与相对版本
% * 丛度量、正交补和短正合丛列的分裂
% 
% 纲要内容结束

\section{光滑单位分解与延拓}
\label{g1:sec:1203}

在每张局部标架中，把纤维坐标的通常内积搬回来，就得到一个局部丛度量。可是两张标架之间的过渡矩阵通常不是正交矩阵，两个局部内积因而未必相同，不能直接粘合。截面延拓也有相似的困难：局部公式可以很好地描述一块区域，却可能在接近定义域边界时发散。我们需要一族光滑权重，让每个局部对象只在自己的定义域内起作用，并使这些权重在各点的总和为 \(1\)。

\ProseParagraphBreak

本节中的流形均为有限维、Hausdorff 且第二可数的光滑流形，允许有边界；向量丛均为有限秩实光滑向量丛。边界处的光滑性沿用\cref{g1:def:manifold-with-boundary}的局部延拓定义。构造首先在 Euclid 空间及半空间中完成，再借助局部有限的坐标覆盖搬到流形上；局部有限性保证以下无穷和在每一点附近其实都是有限和。

\subsection{光滑截断函数和局部有限细化}
\label{g1:subsec:1203:01}

一个函数在开集 \(U\) 外等于零，并不表示它的支集包含于 \(U\)。按照\cref{g1:def:support-cutoff}，支集还包括非零点的所有极限点：
\[
	\operatorname{supp}f=\overline{\{p:f(p)\ne0\}}.
\]
例如 \(f(t)=\max\{1-|t|,0\}\) 在 \((-1,1)\) 外为零，支集却是 \([-1,1]\)。我们要把支集也留在允许的开集内部，并把折角换成任意阶可微的过渡。

\begin{lemma}[光滑截断函数]\label{g1:lem:smooth-bump-function}
	函数
	\[
		\eta(t)=
		\begin{cases}
			e^{-1/t},&t>0,\\
			0,&t\le0
		\end{cases}
	\]
	在 \(\mathbb R\) 上光滑，且 \(\eta^{(k)}(0)=0\) 对每个非负整数 \(k\) 成立。由它可构造光滑函数 \(\vartheta:\mathbb R\to[0,1]\)，满足 \(\vartheta=0\) 于 \((-\infty,0]\)、\(\vartheta=1\) 于 \([1,\infty)\)。因此，若 \(0<r<s<R\)，在 Euclid 空间或半空间 \(H^m=\{x^m\ge0\}\) 内，存在光滑函数，等于 \(1\) 于半径 \(r\) 的闭球或相对闭半球，支集包含于半径 \(s\) 的相对闭球，因而紧含于半径 \(R\) 的相对开球。
\end{lemma}
\begin{proof}
	首先要检查拼接点 \(0\)，仅知道正半轴上的公式光滑是不够的。对 \(t>0\) 反复求导得到
	\[
		\eta^{(k)}(t)=e^{-1/t}P_k(1/t),
	\]
	其中 \(P_k\) 是多项式。确实 \(P_0=1\)，写 \(u=1/t\)，下一次求导给出 \(P_{k+1}(u)=u^2(P_k(u)-P_k'(u))\)。对任意非负整数 \(N\)，有 \(u^N e^{-u}\to0\)（\(u\to+\infty\)）：例如从指数级数取第 \(N+1\) 项，便有 \(u^N e^{-u}\le (N+1)!/u\)。

	把每个 \(e^{-1/t}P_k(1/t)\) 在 \(t\le0\) 补成零，所得函数在 \(0\) 连续。而它在 \(0\) 的右侧差商又是 \(e^{-1/t}P_k(1/t)/t\)，仍趋于零；左侧差商恒为零。因此这个延拓的导数在 \(0\) 为零，在两侧则是下一阶公式。逐阶归纳即得 \(\eta\in C^\infty(\mathbb R)\)，且所有阶导数在 \(0\) 消失。

	令
	\[
		\vartheta(t)=\frac{\eta(t)}{\eta(t)+\eta(1-t)}.
	\]
	分母处处为正：若第一项为零，则 \(t\le0\)，此时 \(1-t>0\)。其余性质由非负性及 \(\eta\) 的零点集直接得到。对中心 \(a\) 定义
	\[
		b_{a,r,s}(x)=
		\vartheta\!\left(\frac{s^2-|x-a|^2}{s^2-r^2}\right).
	\]
	它在 \(|x-a|\le r\) 时为 \(1\)，在 \(|x-a|\ge s\) 时为零，支集包含于闭球 \(\overline B_s(a)\)。限制到 \(H^m\) 后仍是光滑函数，因为原来的公式就是其 Euclid 光滑延拓；支集按半空间的相对拓扑计算，包含于 \(\overline B_s(a)\cap H^m\)。
\end{proof}

若坐标图中半径 \(R\) 的相对闭球仍留在坐标域内，就可把上述函数搬到流形，并在坐标域外补零。补零为何光滑？支集已留在半径 \(s<R\) 的紧集内，每个坐标域外的点都有一个与该支集不交的邻域，所得函数在这个邻域恒为零。边界坐标图使用相对球，论证完全相同。零维流形的每个单点都是开集，其示性函数就是所需的光滑截断函数。

\ProseParagraphBreak

许多局部函数相加时，另一个要求是控制它们在同一邻域内出现的数目。回忆\cref{g1:def:locally-finite-refinement}：集合族 \(\{C_i\}\) 局部有限，是指每一点都有一个开邻域，只与有限个 \(C_i\) 相交；逐点有限则只要求每一点属于有限个 \(C_i\)。后一个条件不能控制邻近点发生的变化。

\begin{example}\label{g1:exm:pointwise-finite-smooth-sum}
	在 \(\mathbb R\) 中令 \(c_n=2^{-n}\)、\(r_n=2^{-n-2}\)，并用前述构造选光滑函数 \(b_n\)，使 \(0\le b_n\le1\)、\(b_n(c_n)=1\)，且
	\[
		\operatorname{supp}b_n\subseteq[c_n-r_n,c_n+r_n].
	\]
	这些闭区间两两不交，所以每一点最多落在一个支集中，\(\sum_n b_n(x)\) 在每一点都只有至多一个非零项。然而这个和在 \(0\) 等于零，在每个 \(c_n\) 处等于 \(1\)，因而在 \(0\) 不连续。问题正是任何原点邻域都会遇到无穷多个支集。\cref{g1:fig:locally-finite-bumps}把这种聚集与沿直线向远处移动的支集作了对照。
\end{example}

\input{Figures/Chapter12-LocallyFiniteBumps.tex}

下面的细化既给出覆盖，又为每个截断函数预留一圈支集边缘。记号 \(V\Subset U\) 表示 \(\overline V\) 紧且 \(\overline V\subseteq U\)，闭包一律在 \(M\) 中计算。

\begin{lemma}[相对紧坐标细化]\label{g1:lem:precompact-coordinate-refinement}
	设 \(M\) 为有限维、Hausdorff 且第二可数的光滑流形，允许有边界，\(\{O_\alpha\}_{\alpha\in I}\) 为任意开覆盖。存在至多可数族坐标邻域对 \(V_i\Subset U_i\) 及指标选择 \(\alpha(i)\)，使
	\[
		M=\bigcup_iV_i,\qquad
		\overline{U_i}\text{ 紧},\qquad
		\overline{U_i}\subseteq O_{\alpha(i)},
	\]
	且 \(\{U_i\}\) 局部有限。每对 \(V_i,U_i\) 可在同一坐标图中选为同心的相对开球；存在光滑函数 \(b_i:M\to[0,1]\)，满足 \(b_i=1\) 于 \(\overline{V_i}\)、\(\operatorname{supp}b_i\Subset U_i\)。
\end{lemma}
\begin{proof}
	\cref{g1:thm:compact-exhaustion,g1:thm:lch-paracompact}分别给出了紧穷竭与壳层细化。这里恢复构造，并在每一层加入所需的两层坐标球。这样从第十一章开始阅读的读者也能看到第二可数性与局部紧性各用在何处。

	先构造紧集列 \(K_n\subseteq\operatorname{int}K_{n+1}\)，使其并为 \(M\)。流形的坐标球或半球可以缩小为紧闭球或紧闭半球，所以每一点的每个邻域都包含一个相对紧开邻域。从一个可数拓扑基中仅保留相对紧的成员，仍得到基：在上述较小邻域中再选基元素，其闭包就是紧闭包的闭子集。把非空成员列为 \(B_1,B_2,\ldots\)，有限时允许重复。令 \(K_1=\overline B_1\)。若 \(K_n\) 已选，由紧性用有限个 \(B_j\) 覆盖 \(K_n\cup\overline B_{n+1}\)，令 \(K_{n+1}\) 为这些 \(B_j\) 的闭包之并。于是
	\[
		K_n\cup\overline B_{n+1}\subseteq\operatorname{int}K_{n+1},
		\qquad
		M=\bigcup_{n\ge1}K_n=\bigcup_{n\ge1}\operatorname{int}K_n.
	\]
	各 \(K_n\) 紧且在 Hausdorff 空间中闭；后一事实可由外点与紧集各点的分离邻域、再取有限子覆盖得到。空流形无需作任何选择。

	约定 \(K_{-1}=K_0=\varnothing\)，定义需要覆盖的紧壳层和允许选择邻域的较大开区域
	\[
		S_n=K_n\setminus\operatorname{int}K_{n-1},
		\qquad
		G_n=\operatorname{int}K_{n+1}\setminus K_{n-2}.
	\]
	有 \(S_n\subseteq G_n\)：外侧由内部包含关系保证，内侧则因为 \(K_{n-2}\subseteq\operatorname{int}K_{n-1}\)。每一点首次进入某个 \(K_n\) 时属于 \(S_n\)，所以各壳层覆盖 \(M\)。

	对每个 \(p\in S_n\)，选择含它的 \(O_{\alpha(p)}\)，在开邻域 \(G_n\cap O_{\alpha(p)}\) 中取一张坐标图。选以 \(p\) 的坐标为中心的三层相对球，使内层 \(V_p\) 的闭包在中层 \(U_p\) 内，中层的闭包又紧含于该坐标图；特别地，\(\overline{U_p}\subseteq G_n\cap O_{\alpha(p)}\)。由壳层的紧性，保留有限个内层球 \(V_{n,1},\ldots,V_{n,k_n}\) 覆盖 \(S_n\)，并同时保留其外层球 \(U_{n,j}\)。所有层合起来至多可数，内层球覆盖 \(M\)。

	还须检查局部有限性；仅仅每层有限并不够。给定 \(p\)，取 \(m\) 使 \(p\in\operatorname{int}K_m\)。当 \(n\ge m+2\) 时，这一层的全部 \(\overline{U_{n,j}}\) 位于 \(G_n\subseteq M\setminus K_{n-2}\)，而 \(K_m\subseteq K_{n-2}\)，故它们与 \(\operatorname{int}K_m\) 不交。这个固定邻域只能遇到前 \(m+1\) 层的有限个外层球，局部有限性得证。最后，在每对内外球之间选一个中间半径，由\cref{g1:lem:smooth-bump-function}构造 \(b_i\)，再补零到 \(M\)。它的支集位于中间闭球内，因而紧含于 \(U_i\)。
\end{proof}

\subsection{单位分解的完整构造}
\label{g1:subsec:1203:02}

现在已有非负函数 \(b_i\)，每一点至少落在一个 \(b_i=1\) 的区域内，各点附近又只出现有限项。只需除以这些函数的总和，就能得到总和为 \(1\) 的权重。这里不能省略支集控制：局部对象只在 \(O_\alpha\) 上定义时，权重的非零点及其极限点都应留在 \(O_\alpha\) 内。

\begin{definition}[光滑单位分解]\label{g1:def:smooth-partition-unity}
	设 \(\{O_\alpha\}_{\alpha\in I}\) 是光滑流形 \(M\) 的开覆盖。一族光滑函数 \(\varphi_\alpha:M\to[0,1]\) 若满足
	\[
		\operatorname{supp}\varphi_\alpha\subseteq O_\alpha,\qquad
		\{\operatorname{supp}\varphi_\alpha\}_{\alpha\in I}\text{ 局部有限},\qquad
		\sum_{\alpha\in I}\varphi_\alpha(p)=1,
	\]
	则称为从属于这个覆盖的\term[guanghuadanweifenjie]{光滑单位分解}[smooth partition of unity]。允许某些 \(\varphi_\alpha\) 恒为零。这是\cref{g1:def:partition-of-unity}中连续单位分解的光滑版本。
\end{definition}

\begin{theorem}[光滑单位分解的存在]\label{g1:thm:smooth-partition-unity}
	有限维、Hausdorff 且第二可数的光滑流形 \(M\)，包括带边界流形，其每个开覆盖都有从属于该覆盖的光滑单位分解。构造可以先得到一个从属于可数细化、每项紧支且支集族局部有限的光滑单位分解，再按原覆盖指标归组；归组后的各项不一定紧支。若 \(M\) 紧，可以只使用有限个非零项。
\end{theorem}
\begin{proof}
	给定覆盖 \(\{O_\alpha\}\)，用\cref{g1:lem:precompact-coordinate-refinement}选取 \(V_i,U_i,b_i\)。由于 \(\operatorname{supp}b_i\subseteq U_i\)，支集族局部有限。因此
	\[
		b=\sum_i b_i
	\]
	在每一点的某个邻域内都是有限个光滑函数之和，从而光滑。内层 \(V_i\) 覆盖 \(M\)，故每个点处至少有一个 \(b_i=1\)，得到 \(b\ge1\)。定义
	\[
		\psi_i=\frac{b_i}{b}.
	\]
	这些函数光滑、非负、总和为 \(1\)，且 \(\psi_i\) 与 \(b_i\) 有相同的非零点集，所以
	\[
		\operatorname{supp}\psi_i=\operatorname{supp}b_i\Subset U_i\subseteq O_{\alpha(i)}.
	\]
	这给出了细化上的单位分解。

	若需要与原覆盖同指标，令
	\[
		\varphi_\alpha=\sum_{\alpha(i)=\alpha}\psi_i.
	\]
	每个和局部只有有限项，因此光滑。说明支集包含关系时，要防止无限并增加边界点。局部有限闭集族的任意子族之并仍闭：在并集外的一个点附近，只遇到有限个闭集，去掉这有限个闭集后仍留下一个开邻域。将这点用于 \(\{\operatorname{supp}\psi_i\}\)，得到
	\[
		\operatorname{supp}\varphi_\alpha
		\subseteq\bigcup_{\alpha(i)=\alpha}\operatorname{supp}\psi_i
		\subseteq O_\alpha.
	\]
	前一个包含成立，是因为右侧是包含所有非零点的闭集。一个只与有限个 \(\operatorname{supp}\psi_i\) 相交的邻域，也只可能遇到相应有限个 \(\operatorname{supp}\varphi_\alpha\)，所以归组后支集族仍局部有限。各点处重排的只是有限和，故 \(\sum_\alpha\varphi_\alpha=1\)。非负性又保证各项不超过 \(1\)。若某个原指标没有分配到成员，就取零函数。

	当 \(M\) 紧时，先从覆盖 \(M\) 的内层球 \(V_i\) 中选有限子覆盖，只保留相应的 \(b_i\)，再作同样的归一化即可。一般情形下，按原指标归组会把无穷多个紧支集并起来，这个并虽然闭，却未必紧。
\end{proof}

局部有限还保证可以逐项求任意阶导数，因为在计算点的一个固定邻域内，本来就只有有限个函数。这里既不需要估计导数级数，也不涉及交换无限求和与极限；\cref{g1:exm:pointwise-finite-smooth-sum}说明仅有逐点有限时，连连续性都可能失败。

\begin{corollary}[闭集的光滑截断]\label{g1:cor:smooth-cutoff-closed}
	设 \(M\) 为有限维、Hausdorff 且第二可数的光滑流形，允许有边界，\(A\subseteq M\) 闭，\(U\subseteq M\) 开且 \(A\subseteq U\)。存在光滑函数 \(\chi:M\to[0,1]\)，使
	\[
		\chi=1\text{ 于 }A\text{ 的某个开邻域},\qquad
		\operatorname{supp}\chi\subseteq U.
	\]
	若 \(A\) 紧，还可要求 \(\operatorname{supp}\chi\) 紧。
\end{corollary}
\begin{proof}
	对开覆盖 \(\{U,M\setminus A\}\) 取单位分解 \(\chi+\rho=1\)，其中 \(\operatorname{supp}\chi\subseteq U\)、\(\operatorname{supp}\rho\subseteq M\setminus A\)。开集 \(M\setminus\operatorname{supp}\rho\) 包含 \(A\)，在这里 \(\rho=0\)，所以 \(\chi=1\)。

	若 \(A\) 紧，在 \(U\) 中给每个 \(p\in A\) 选相对紧坐标邻域，再取有限个覆盖 \(A\)。它们的并 \(W\) 满足 \(A\subseteq W\) 且 \(\overline W\) 紧、\(\overline W\subseteq U\)。把前面的构造用于 \(A\subseteq W\)，所得支集闭且包含于紧集 \(\overline W\)，因而紧。若 \(A=\varnothing\)，直接取 \(\chi=0\)。
\end{proof}

“支集包含于 \(U\)”与“支集紧”是两件事。即使 \(U\) 非紧，只要 \(A\) 紧，仍可选紧支截断；反过来，若闭集 \(A\) 非紧，任何在 \(A\) 上等于 \(1\) 的函数都不可能紧支，因为其支集包含 \(A\)。例如 \(A=\mathbb Z\subset U=\mathbb R\)，截断函数存在，却不能紧支。

\subsection{截面延拓与相对版本}
\label{g1:subsec:1203:03}

回忆\cref{g1:def:bundle-section}，丛 \(E\to M\) 的截面 \(s\) 在每一点 \(p\) 取值于对应纤维 \(E_p\)，光滑性可在局部标架中逐分量判断。截面的支集同样定义为
\[
	\operatorname{supp}s=\overline{\{p:s(p)\ne0_p\}}.
\]
对开集 \(U\) 上的光滑截面 \(s\)，若光滑函数 \(\chi\) 满足 \(\operatorname{supp}\chi\subseteq U\)，就可以令 \(\chi s\) 在 \(U\) 外等于零。它在 \(U\) 内逐分量光滑，在 \(U\) 外的每个点附近恒为零，因此是整体光滑截面。这里不要求 \(s\) 在接近 \(U\) 的边界时有界。

\begin{definition}[闭子集上的光滑截面]\label{g1:def:smooth-section-closed-set}
	设 \(E\to M\) 为光滑向量丛，\(A\subseteq M\) 为闭集。给定一个赋值 \(a(p)\in E_p\)（\(p\in A\)）。局部光滑延拓的条件是：每个 \(p\in A\) 都有 \(M\) 中的开邻域 \(O_p\) 和光滑截面 \(s_p\in\Gamma(E|_{O_p})\)，满足
	\[
		s_p(q)=a(q)\qquad(q\in A\cap O_p).
	\]
	若满足这一条件，就称这一赋值为 \(A\) 上的光滑截面。
\end{definition}

这个定义允许 \(A\) 不是子流形，只要求所给纤维值在各点附近可以光滑延拓。不同 \(s_p\) 在 \(A\) 外可以不同；当前数据没有指定法向导数或更高阶导数。

\begin{theorem}[截面延拓]\label{g1:thm:section-extension}
	设 \(M\) 为有限维、Hausdorff 且第二可数的光滑流形，允许有边界，\(E\to M\) 为有限秩实光滑向量丛，\(A\subseteq M\) 闭，\(U\supseteq A\) 开。
	\begin{enumerate}
		\item 若 \(s\in\Gamma(E|_U)\)，则存在整体光滑截面 \(\widetilde s\)，在 \(A\) 的某开邻域与 \(s\) 相同，且 \(\operatorname{supp}\widetilde s\subseteq U\)。
		\item 若 \(a\) 为\cref{g1:def:smooth-section-closed-set}意义下 \(A\) 上的光滑截面，则存在整体光滑截面 \(\widetilde a\)，满足 \(\widetilde a|_A=a\) 且 \(\operatorname{supp}\widetilde a\subseteq U\)。
	\end{enumerate}
	若 \(A\) 紧，两项都可要求所得截面紧支。
\end{theorem}
\begin{proof}
	第一项中，由\cref{g1:cor:smooth-cutoff-closed}选取在 \(A\) 附近等于 \(1\)、支集包含于 \(U\) 的 \(\chi\)。令
	\[
		\widetilde s(p)=
		\begin{cases}
			\chi(p)s(p),&p\in U,\\
			0_p,&p\notin U.
		\end{cases}
	\]
	前面的零延拓论证说明其光滑，支集包含于 \(\operatorname{supp}\chi\)，并且在 \(\chi=1\) 的那个邻域与 \(s\) 相同。紧性结论也由 \(\chi\) 的紧支性得到。

	第二项先粘合各个局部延拓，再控制支集。为每个 \(p\in A\) 取定义中的 \(O_p,s_p\)，并加入开集 \(O_*=M\setminus A\) 及其零截面 \(s_*=0\)，得到 \(M\) 的开覆盖。取从属单位分解 \(\{\varphi_p\}\cup\{\varphi_*\}\)。每个乘积 \(\varphi_p s_p\) 都可在 \(O_p\) 外光滑补零，支集仍包含于 \(\operatorname{supp}\varphi_p\)。因此
	\[
		t=\sum_{p\in A}\varphi_p s_p+\varphi_*s_*
	\]
	是整体光滑截面：每一点附近只需求有限项相加。若 \(q\in A\)，则 \(\varphi_*(q)=0\)，而每个在 \(q\) 处权重非零的 \(s_p\) 都满足 \(s_p(q)=a(q)\)，从而
	\[
		t(q)=\left(\sum_{p\in A}\varphi_p(q)\right)a(q)=a(q).
	\]
	最后再用第一项中的截断函数，取 \(\widetilde a=\chi t\)，即可控制支集并保持 \(A\) 上的值；\(A\) 紧时选紧支的 \(\chi\)。
\end{proof}

第一项保留的是一个邻域上的整条截面，第二项保留的是闭集上指定的值；不能从第二项自动推断它同时保留各个局部延拓的全部导数。闭性也有实际作用。取平凡线丛 \(\mathbb R\times\mathbb R\)，在非闭集 \(A=(0,1)\) 上给定 \(a(x)=1/x\)。它甚至已在开集 \(A\) 上光滑，却不可能延拓为 \(\mathbb R\) 上的连续截面，因为 \(x\downarrow0\) 时发散。

\ProseParagraphBreak

有时已经有一个整体截面 \(s_0\)，只希望在某闭集附近把它换成指定截面，并在另一闭集附近保持原样。两块区域相遇时需要相容；仅在相遇点取值相同，未必足以要求双方都在邻域内保持不变。

\begin{proposition}[相对截面延拓]\label{g1:prop:relative-section-extension}
	设 \(M\) 为 Hausdorff、第二可数的有限维光滑流形，允许带边界，\(E\to M\) 为有限秩实光滑向量丛，\(A,B\subseteq M\) 为闭集，\(s_0\in\Gamma(E)\)。这些流形与丛条件保证可以应用\cref{g1:thm:section-extension}。
	\begin{enumerate}
		\item 设 \(U\supseteq A\) 开，\(s\in\Gamma(E|_U)\)。若有开集 \(W\supseteq A\cap B\)，使 \(s=s_0\) 于 \(U\cap W\)，则存在 \(\widetilde s\in\Gamma(E)\)，在 \(A\) 的某开邻域等于 \(s\)，在 \(B\) 的某开邻域等于 \(s_0\)，且 \(\operatorname{supp}(\widetilde s-s_0)\subseteq U\)。
		\item 设 \(a\) 是闭集 \(A\) 上的光滑截面。若有开集 \(W\supseteq B\)，使 \(a=s_0|_A\) 于 \(A\cap W\)，则存在整体延拓 \(\widetilde a|_A=a\)，并且 \(\widetilde a=s_0\) 于 \(B\) 的某开邻域。
	\end{enumerate}
	若 \(A\) 紧，还可在两项中要求对 \(s_0\) 的改变量紧支。
\end{proposition}
\begin{proof}
	第一项中，闭集 \(B\setminus W\) 与 \(A\) 不交，所以
	\[
		O=U\setminus(B\setminus W)
	\]
	是包含 \(A\) 的开集。选在 \(A\) 附近为 \(1\)、支集包含于 \(O\) 的截断函数 \(\chi\)，把 \(\chi(s-s_0)\) 从 \(U\) 补零，令
	\[
		\widetilde s=s_0+\chi(s-s_0).
	\]
	在 \(A\) 附近有 \(\chi=1\)，故 \(\widetilde s=s\)。在 \(W\) 上改变量为零，因为 \(U\cap W\) 内两截面相等，\(U\) 外则已补零。在 \(M\setminus\operatorname{supp}\chi\) 上改变量也为零。这两个开集的并包含 \(B\)：属于 \(B\cap W\) 的点在前者中，属于 \(B\setminus W\) 的点不在支集内，因而在后者中。支集控制由构造得到。

	第二项中，先用\cref{g1:thm:section-extension}取 \(a\) 的一个整体延拓 \(t\)。闭集 \(C=A\setminus W\) 与 \(B\) 不交，故可选在 \(C\) 附近为 \(1\)、支集包含于 \(M\setminus B\) 的光滑截断函数 \(\chi\)。令
	\[
		\widetilde a=s_0+\chi(t-s_0).
	\]
	在 \(A\setminus W\) 上 \(\chi=1\)，在 \(A\cap W\) 上则有 \(t=a=s_0\)，所以两部分都给出 \(\widetilde a=a\)。开集 \(M\setminus\operatorname{supp}\chi\) 包含 \(B\)，其上 \(\widetilde a=s_0\)。若 \(A\) 紧，第一项中的 \(A\) 与第二项中的 \(C\) 都紧，可选紧支的 \(\chi\)，于是改变量紧支。
\end{proof}

例如在平凡线丛上取 \(A=[0,1]\)、\(B=\{0\}\)、\(a(x)=x\)、\(s_0=0\)。虽然 \(a\) 与 \(s_0\) 在 \(A\cap B\) 上相等，却不能找一个同时延拓 \(a\) 并在 \(0\) 附近恒为零的截面。因此相对版本中的邻域相容条件不能只改成相交点上的相等。当 \(A,B\) 不交时，第一项的相容条件自动满足，可取 \(W=\varnothing\)。

\subsection{丛度量、正交补和短正合丛列的分裂}
\label{g1:subsec:1203:04}

截面的局部公式相加前，必须先确认它们在同一点属于同一个向量空间。丛的局部标架本身不提供这种相加规则：不同标架中的同一列数字，通常表示不同的向量。单位分解粘合的是已解释为纤维对象的局部截面、双线性型或丛同态，再通过标架检验其光滑性。

\begin{definition}[丛度量]\label{g1:def:bundle-metric}
	设 \(E\to M\) 为有限秩实光滑向量丛。给每个纤维 \(E_p\) 指定正定对称双线性型 \(h_p\)，若在任一光滑局部标架 \(e_1,\ldots,e_r\) 中，所有函数 \(h(e_a,e_b)\) 都光滑，则称 \(h\) 为 \(E\) 上的\term[congduliang]{丛度量}[bundle metric]。
\end{definition}

只需在一组覆盖 \(M\) 的局部标架中检查光滑性。事实上，按照\cref{g1:def:bundle-section}的行标架约定 \(e_j=e_i g_{ij}\)，度量矩阵满足
\[
	H_j=g_{ij}^{\mathsf T}H_i g_{ij}.
\]
右侧是光滑矩阵的乘积，因而光滑。这里说的是各纤维内的内积；在尚未另作定义之前，它不是底空间 \(M\) 上两点间的距离。

\begin{theorem}[丛度量的存在]\label{g1:thm:bundle-metric-existence}
	设 \(M\) 为有限维、Hausdorff 且第二可数的光滑流形，允许有边界，\(E\to M\) 为有限秩实光滑向量丛，则 \(E\) 存在光滑丛度量。
\end{theorem}
\begin{proof}
	取局部标架覆盖 \(\{O_\alpha\}\)。在每个 \(O_\alpha\) 上，把该标架声明为正交归一，得到局部度量
	\[
		h_\alpha\!\left(\sum_a v^a e_{\alpha,a},\sum_a w^a e_{\alpha,a}\right)
		=\sum_a v^a w^a.
	\]
	它在别的标架中仍有光滑矩阵，因为过渡矩阵光滑。取从属于 \(\{O_\alpha\}\) 的光滑单位分解，定义
	\[
		h=\sum_\alpha\varphi_\alpha h_\alpha,
	\]
	每个加权局部双线性型在定义域外补零。其支集控制保证补零光滑；局部有限性保证每一矩阵分量在各点附近都是有限个光滑函数之和。

	所得 \(h_p\) 对称且双线性。若 \(v\in E_p\) 非零，则每个有定义的 \(h_{\alpha,p}(v,v)>0\)，而权重非负且至少一个为正，故
	\[
		h_p(v,v)=\sum_\alpha\varphi_\alpha(p)h_{\alpha,p}(v,v)>0.
	\]
	这给出正定性，完成构造。
\end{proof}

正定性能够保留，是因为正定对称双线性型构成凸集。给定两个丛度量 \(h_0,h_1\)，对 \(0\le t\le1\)，\((1-t)h_0+th_1\) 仍是丛度量。存在性并未选出一个唯一度量，局部标架与单位分解的选择通常都会改变结果。

\begin{proposition}[正交补子丛]\label{g1:prop:orthogonal-complement-bundle}
	设 \(E\to M\) 为有限秩实光滑向量丛，\(h\) 为光滑丛度量，\(F\subseteq E\) 为秩 \(k\) 的光滑子丛。则逐纤维的正交补
	\[
		F^\perp_p=\{v\in E_p:h_p(v,w)=0\text{ 对所有 }w\in F_p\}
	\]
	构成光滑子丛，并有光滑直和分解 \(E=F\oplus F^\perp\)。到 \(F\) 的正交投影是光滑丛同态，\(h\) 在两个子丛上的限制仍为光滑丛度量。
\end{proposition}
\begin{proof}
	每个纤维中的线性代数已给出向量空间直和；要补足的是正交补和投影如何随底点光滑变化。由\cref{g1:prop:subbundle-adapted-frame}，在一点附近取环境丛标架 \(e_1,\ldots,e_r\)，使前 \(k\) 项张成 \(F\)。令
	\[
		A_{ab}=h(e_a,e_b)\quad(1\le a,b\le k).
	\]
	矩阵 \(A\) 光滑且正定，故可逆，\(A^{-1}=\operatorname{adj}(A)/\det A\) 也是光滑矩阵。对纤维向量 \(v\) 定义
	\[
		P(v)=\sum_{a,b=1}^k e_a(A^{-1})_{ab}\,h(e_b,v).
	\]
	对每个 \(c\le k\)，有
	\[
		h(e_c,P(v))=\sum_{a,b=1}^k A_{ca}(A^{-1})_{ab}h(e_b,v)
		=h(e_c,v).
	\]
	所以 \(v-P(v)\) 与 \(F\) 正交；若 \(v\in F\)，则公式又给出 \(P(v)=v\)。这正是纤维上的唯一正交投影，因而不同局部标架得到的公式自动相同，可以粘合为整体光滑丛同态。

	对 \(j=k+1,\ldots,r\) 令 \(f_j=e_j-P(e_j)\)。它们光滑且落在 \(F^\perp\) 中；若其线性组合为零，在商空间 \(E_p/F_p\) 中取像，得到 \(e_{k+1},\ldots,e_r\) 的对应线性组合为零，从而全部系数为零。它们共有 \(r-k\) 个，正好构成 \(F^\perp\) 的局部标架，证明了子丛的光滑性。映射 \(F\oplus F^\perp\to E\) 的逆为 \(v\mapsto(Pv,v-Pv)\)，两者都光滑。度量限制的矩阵光滑，正定性逐纤维保持；\(k=0\) 或 \(k=r\) 时按零子丛解释同样成立。
\end{proof}

有了正交补，就能把逐纤维的短正合列拆成整体直和。回忆\cref{g1:def:short-exact-bundle-sequence}，短正合要求第一张箭头在每个纤维单射、第二张箭头满射，且前者的像等于后者的核；这些条件本身还没有指定一个与核互补的子丛。

\begin{theorem}[短正合丛列的光滑分裂]\label{g1:thm:exact-bundle-splitting}
	设 \(M\) 为有限维、Hausdorff 且第二可数的光滑流形，允许有边界，且
	\[
		0\longrightarrow E'\xrightarrow{\ i\ }E\xrightarrow{\ q\ }E''\longrightarrow0
	\]
	是 \(M\) 上有限秩实光滑向量丛的短正合列。则存在光滑丛同态 \(j:E''\to E\)，使 \(qj=\operatorname{id}_{E''}\)，并由
	\[
		E'\oplus E''\longrightarrow E,\qquad (u,w)\longmapsto i(u)+j(w)
	\]
	得到光滑丛同构。这样的分裂一般不唯一。
\end{theorem}
\begin{proof}
	由\cref{g1:thm:constant-rank-bundle-morphism}，\(F=i(E')=\ker q\) 是光滑子丛；这里 \(i,q\) 的纤维秩由逐纤维单射、满射保证恒定。为 \(E\) 选丛度量，令 \(C=F^\perp\)。限制 \(q|_C:C\to E''\) 逐纤维为线性同构：其核是 \(C_p\cap F_p=0\)，任意 \(q_p(v)\) 又等于分解 \(v=v_F+v_C\) 后的 \(q_p(v_C)\)，故它满射。

	\cref{g1:prop:bundle-morphism-matrices}说明，逐纤维可逆的光滑丛同态具有光滑逆；其理由是在局部标架中取可逆光滑矩阵的逆，仍由伴随矩阵除以不为零的行列式给出。于是取 \(j=(q|_C)^{-1}\)，再包含到 \(E\) 中，就有 \(qj=\operatorname{id}\)。对任意 \(v\in E\)，
	\[
		v=(v-jqv)+jqv,\qquad v-jqv\in F,
	\]
	所以所述直和同构的逆为
	\[
		v\longmapsto\bigl(i^{-1}(v-jqv),qv\bigr).
	\]
	这里 \(i:E'\to F\) 也是光滑丛同构，因此该逆光滑。

	最后说明选择的自由度。若 \(j_0,j_1\) 都是分裂，则 \(q(j_1-j_0)=0\)，故其差唯一写成 \(i\circ a\)，其中 \(a:E''\to E'\) 为光滑丛同态；反过来，任意这样的 \(a\) 都使 \(j_0+i\circ a\) 成为分裂。例如平凡丛列 \(\mathbb R\to\mathbb R^2\to\mathbb R\) 中，第一箭头包含为第一坐标、第二箭头投影到第二坐标，则每个光滑函数 \(a\) 都给出分裂 \(t\mapsto(a(p)t,t)\)。因此短正合列本身通常没有指定一个分裂。
\end{proof}

这些构造也说明单位分解可以保持什么。若 \(L:E\to F\) 是固定的丛同态，局部截面都满足同一个线性条件 \(L(s_\alpha)=t\)，则
\[
	L\!\left(\sum_\alpha\varphi_\alpha s_\alpha\right)
	=\sum_\alpha\varphi_\alpha L(s_\alpha)
	=\left(\sum_\alpha\varphi_\alpha\right)t=t.
\]
正定性也由非负权重保持。但两个处处非零的局部截面可以相互抵消，例如平凡线丛中的 \(1\) 与 \(-1\) 以相同权重平均就得到零。逐向量平均两组正交归一标架，甚至未必得到一组标架，更不能自动保持正交归一；非退化交替形式与几乎复结构的代数条件也不能直接由平均保持，\secref{g1:sec:1204}将给出具体例子；含导数的条件还须另作检验。因此“局部存在”能否推出“整体存在”，仍须检验所需条件是否被这个加权构造保留。

% 纲要条目排为编号小节标题，并进入目录与书签；说明段仍保留于纲要注释中。

% !TeX root = ../../Book1.tex
% 纲要标题：### 12.4 张量场与几何结构
% 纲要位置：几何学计划纲要.md，第 594 行。
% 正文状态：初稿；尚未正式审核。

% 纲要内容开始
% 
% * 度量、微分形式和场的光滑性
% * 自然丛运算与拉回
% * 结构群约化的模型及其存在问题
% 
% ---
% 
% 纲要内容结束

\section{张量场与几何结构}
\label{g1:sec:1204}

前面已经分别建立了两种构造：过渡函数把一族向量空间粘成向量丛，单位分解把相容的局部要求组织成整体截面。现在把它们用于几何。度量给每个切空间一个内积，微分形式给每个切空间一个交替的多线性函数；局部矩阵按换标架规律相互转换，并随基点光滑变化，从而给出整体几何对象。哪些变化规律来自代数运算，哪些几何条件可以通过粘合保持，是本节要分开的两个问题。

\subsection{度量、微分形式和场的光滑性}
\label{g1:subsec:1204:01}

先说明怎样把\secref{g1:sec:1202}的逐点运算变成一个光滑空间。设 \(E,F\) 是同一光滑流形 \(M\) 上的有限秩实光滑向量丛。直和已经构造；对偶、张量积和外幂也可以在每个纤维中进行，但要确认它们的换标架仍满足余循环条件。

\begin{proposition}[由纤维运算构造向量丛]
	\label{g1:prop:tensor-bundle-constructions}
	设 \(M\) 为 Hausdorff、第二可数光滑流形，可带边界，\(E,F\to M\) 为有限秩实光滑向量丛。纤维分别为 \(E_p^*\)、\(E_p\otimes F_p\)、\(\operatorname{Hom}(E_p,F_p)\)、\(\operatorname{Sym}^kE_p\)、\(\Lambda^kE_p\) 的不交并，都有与局部标架相容的自然光滑向量丛结构。逐点的求值、收缩、对称化、楔积以及线性映射的复合，均给出相应光滑丛映射。
\end{proposition}

\begin{proof}
	取共同细化的平凡化覆盖，设 \(E,F\) 的向量分量过渡分别为 \(v_i=g_{ij}v_j\)、\(w_i=h_{ij}w_j\)。对偶、张量积和同态的分量过渡依次为
	\[
		\alpha_i=g_{ij}^{-\mathsf T}\alpha_j,\qquad
		z_i=(g_{ij}\otimes h_{ij})z_j,\qquad
		A_i=h_{ij}A_jg_{ij}^{-1}.
	\]
	对称幂和外幂使用由 \(g_{ij}\) 诱导的 \(\operatorname{Sym}^kg_{ij}\)、\(\Lambda^kg_{ij}\)。这些矩阵元由原矩阵元的多项式与可逆矩阵的逆组成，故光滑。每一种代数构造保持恒等映射与复合，因此新的过渡函数也满足余循环条件。例如 \((g_{ij}\otimes h_{ij})(g_{jk}\otimes h_{jk})=g_{ik}\otimes h_{ik}\)。\cref{g1:thm:bundle-cocycle-gluing}便给出相应向量丛。

	如果改变原平凡化，新、旧标架之间的同构也诱导上述各空间之间的同构，所以纤维上的恒等识别给出同一个光滑结构。在这些图中，求值与收缩是坐标分量的有限次乘法和求和，对称化是有限次线性组合，楔积是带符号的有限次乘法和求和，复合是矩阵乘法，故都光滑。它们的定义不依赖基，因而在重叠图中相容，确实成为整体丛映射。
\end{proof}

\begin{definition}
	\label{g1:def:smooth-tensor-field}
	设 \(M\) 为光滑流形，可带边界，\(r,s\geq0\)。记
	\[
		T^r_sM=(\mathrm{T}M)^{\otimes r}\otimes(\mathrm{T}^*M)^{\otimes s}.
	\]
	它的光滑截面称为 \((r,s)\) 型\term[zhangliangchang]{张量场}[tensor field]。零次张量积约定为平凡实线丛，因此 \((0,0)\) 型张量场就是光滑实函数。\(\mathrm{T}M\) 的光滑截面称为\term[xiangliangchang]{向量场}[vector field]，其空间记为 \(\mathfrak X(M)\)；\(\Lambda^k\mathrm{T}^*M\) 的光滑截面称为光滑 \(k\) 次\term[weifenxingshi]{微分形式}[differential form]，其空间记为 \(\Omega^k(M)\)。\symidx{\(T^r_sM\)：张量丛}\symidx{\(\Omega^k(M)\)：光滑微分形式空间}
\end{definition}

设局部切标架为 \(e_1,\ldots,e_n\)，对偶标架为 \(\varepsilon^1,\ldots,\varepsilon^n\)。一个张量场写成
\[
	T=\sum T^{i_1\cdots i_r}_{j_1\cdots j_s}\,
	e_{i_1}\otimes\cdots\otimes e_{i_r}\otimes
	\varepsilon^{j_1}\otimes\cdots\otimes\varepsilon^{j_s}.
\]
它光滑当且仅当所有分量函数光滑，这是向量丛截面的局部判据。\cref{g1:prop:tensor-component-change}说明换标架后分量乘上相应过渡矩阵及其逆，因而光滑性不依赖选择。这些分量不是各自独立的几何不变量，只有连同换标架规律理解，才能恢复同一个 \(T\)。

\ProseParagraphBreak

张量场还可以通过它如何作用于其他场来识别。一个 \((r,s)\) 型张量在每一点输入 \(r\) 个协变向量和 \(s\) 个切向量；逐点求值得到光滑函数。这里的多线性不仅是对实数线性，而且是对光滑函数线性。例如对一个协变二阶张量，\(T(fX,Y)=fT(X,Y)\)，不应额外出现 \(f\) 的导数。

\begin{proposition}[张量场的多线性判据]
	\label{g1:prop:tensor-field-multilinearity}
	在 Hausdorff、第二可数光滑流形 \(M\)（可带边界）上，光滑 \((r,s)\) 型张量场与 \(C^\infty(M)\) 多线性映射
	\[
		\underbrace{\Omega^1(M)\times\cdots\times\Omega^1(M)}_{r\text{ 个}}
		\times\underbrace{\mathfrak X(M)\times\cdots\times\mathfrak X(M)}_{s\text{ 个}}
		\longrightarrow C^\infty(M)
	\]
	一一对应，对应为逐点求值。当 \(r=s=0\) 时，右端只表示指定一个光滑函数。
\end{proposition}

\begin{proof}
	张量场给出的映射具有所述多线性，坐标公式还说明输出光滑。反向设给定多线性映射 \(A\)，首先证明它在点 \(p\) 的值只看各输入场在 \(p\) 的取值。

	固定其他输入，若某个输入场 \(Z\) 在 \(p\) 的邻域内恒为零，取一个在 \(p\) 附近为一、支集包含于该邻域的光滑截断函数 \(\chi\)，其存在由\cref{g1:cor:smooth-cutoff-closed}保证。于是 \(\chi Z=0\)，由多线性，\(\chi A(\ldots,Z,\ldots)=0\)，故输出在 \(p\) 为零。因此局部相同的输入可互相替换。

	在 \(p\) 附近取一个局部标架，将它乘以支集包含于图域且在更小邻域为一的截断函数，再在图域外延零，就得到在较小邻域等于该标架的整体场；对偶标架同样处理。任一输入场 \(Z\) 在该邻域可展开为 \(\sum z^i e_i\)。再用一个支集包含于此邻域的截断函数把系数延成整体光滑函数，由刚才的局部性与 \(C^\infty(M)\) 多线性，在点 \(p\) 得
	\[
		A(\ldots,Z,\ldots)(p)=\sum_i z^i(p)A(\ldots,e_i,\ldots)(p),
	\]
	其中右端的 \(e_i\) 表示已经延拓的场。对每个输入重复，便证明输出只依赖于各输入在 \(p\) 的值，并给出一个有限维多线性函数，即张量 \(T_p\)。这里允许把任意点处向量作为输入，是因为局部常系数标架组合也能用同样方法延拓。

	在固定的较小图域内，\(T\) 的每个分量都是 \(A\) 作用于若干固定整体标架场后得到的光滑函数，故 \(T\) 光滑。它的逐点求值正是 \(A\)，且各点全部输入值都可由整体场实现，所以 \(T\) 唯一。
\end{proof}

这个判据也帮助排除看起来像张量的表达。在坐标区间上，令 \(B(f\partial_x,g\partial_x)=f g'\)。它对实数双线性，却有 \(B(X,hY)=hB(X,Y)+X(h)g\)，一般不对光滑函数线性，所以不是二阶张量。额外出现的导数表明它还在观察输入场在附近如何变化，而不只观察一个点的值。

\begin{definition}
	\label{g1:def:riemannian-metric-tensor}
	光滑流形 \(M\) 上的\term[limanduliang]{Riemann 度量}[Riemannian metric]是一个光滑截面 \(g\in\Gamma(\operatorname{Sym}^2\mathrm{T}^*M)\)，使每个 \(g_p\) 都是 \(\mathrm{T}_pM\) 上的正定内积。这里将对称张量与对称多线性形式作\secref{g1:sec:1202}中的自然识别。
\end{definition}

\cref{g1:thm:bundle-metric-existence}应用于 \(\mathrm{T}M\)，保证每个本书意义下的光滑流形都有 Riemann 度量。度量输入同一点的两个向量，不是直接输入两个流形点；由它构造曲线长度与距离，还需另作定义。在标架中记 \(G=(g(e_i,e_j))\)，换标架 \(e'=eA\) 时
\[
	G'=A^{\mathsf T}GA.
\]
因此一个正定矩阵函数只有按这个规律在交叠处相容，才表示整体度量。

\begin{proposition}[度量的对偶识别]
	\label{g1:prop:metric-duality-smooth}
	设 \(g\) 为光滑流形 \(M\)（可带边界）上的 Riemann 度量。映射
	\[
		g^\flat:\mathrm{T}M\longrightarrow \mathrm{T}^*M,\qquad
		v\longmapsto g(v,\cdot)
	\]
	是光滑丛同构，记其逆为 \(g^\sharp\)。对任意 \(f\in C^\infty(M)\)，向量场 \(\operatorname{grad}_g f=g^\sharp(\mathrm{d}f)\) 是唯一满足 \(g(\operatorname{grad}_g f,X)=\mathrm{d}f(X)\) 的光滑向量场。
\end{proposition}

\begin{proof}
	在局部标架中，\(g^\flat\) 的矩阵就是 \(G\)。正定性使 \(G\) 可逆，逆矩阵各项由余子式除以非零行列式表示，所以 \(G^{-1}\) 光滑。这给出光滑的逐纤维逆，且求值定义保证各图中的构造相容。令 \(\operatorname{grad}_g f=g^\sharp(\mathrm{d}f)\) 就得到等式；非退化性保证唯一。坐标表达为 \(\operatorname{grad}_g f=\sum_{i,j}g^{ij}\partial_jf\,\partial_i\)，其中 \((g^{ij})=(g_{ij})^{-1}\)。
\end{proof}

对微分形式，若 \(x^1,\ldots,x^n\) 是坐标，则每个 \(k\) 次形式唯一写成
\[
	\alpha=\sum_{i_1<\cdots<i_k}a_{i_1\cdots i_k},
	\mathrm{d}x^{i_1}\wedge\cdots\wedge\mathrm{d}x^{i_k},
\]
光滑性等价于这些系数光滑。\(k>n\) 时只有零形式，\(\Omega^0(M)=C^\infty(M)\)。逐点楔积定义 \(\Omega^p(M)\times\Omega^q(M)\to\Omega^{p+q}(M)\)，逐点插入向量场定义 \(\iota_X\alpha\in\Omega^{k-1}(M)\)。它们光滑，并满足\cref{g1:prop:wedge-algebra,g1:prop:interior-product-rule}的分次公式，因为每一点上都满足同样的代数恒等式。外微分将由\chapref{g1:ch:14}另行构造；当前的一形式 \(\mathrm{d}f\) 已有第十一章的定义。

\begin{example}[极坐标中的度量与面积]
	\label{g1:exm:polar-tensors-density}
	在 \(r>0\) 的一个角扇区中，\(x=r\cos\theta\)、\(y=r\sin\theta\)。Euclid 度量和标准二形式的坐标表达分别为
	\[
		g=\mathrm{d}r\otimes\mathrm{d}r+r^2\mathrm{d}\theta\otimes\mathrm{d}\theta,
		\mathrm{d}x\wedge\mathrm{d}y=r\,\mathrm{d}r\wedge\mathrm{d}\theta.
	\]
	第一式由 \(\mathrm{d}x=\cos\theta\,\mathrm{d}r-r\sin\theta\,\mathrm{d}\theta\)、\(\mathrm{d}y=\sin\theta\,\mathrm{d}r+r\cos\theta\,\mathrm{d}\theta\) 展开后交叉项抵消得到；第二式由楔积的反交换关系得到。\(g\) 的矩阵为 \(\operatorname{diag}(1,r^2)\)，所以体积密度为 \(r|\mathrm{d}r\wedge\mathrm{d}\theta|\)。若把角度坐标反向，二形式的相应系数改变符号，密度使用绝对行列式，因而仍为正。不同坐标中的系数虽不同，度量、二形式与密度各自仍是同一个整体对象。
\end{example}

\subsection{自然丛运算与拉回}
\label{g1:subsec:1204:02}

“在 \(N\) 上有一个对象，把它拉回到 \(M\)”有两种不同含义。只改变基空间，得到的是拉回丛 \(F^*E\) 上的数据；如果还要让它作用于 \(M\) 自己的切向量，就需要使用 \(\mathrm{d}F\)。前一步保留目标空间的纤维，后一步把源切向量作为输入，二者由微分联系。

\begin{proposition}[丛运算与基空间拉回]
	\label{g1:prop:bundle-operations-pullback}
	设 \(F:M\to N\) 为光滑映射，两流形可带边界，\(E,H\to N\) 为有限秩实光滑向量丛。存在自然光滑丛同构
	\[
		F^*(E\otimes H)\cong F^*E\otimes F^*H,\qquad
		F^*(E^*)\cong(F^*E)^*,\qquad
		F^*(\Lambda^kE)\cong\Lambda^k(F^*E),
	\]
	对直和、对称幂和同态丛也有相同性质。每个光滑截面 \(s\) 都给出拉回截面 \((F^*s)(p)=(p,s(F(p)))\)，且这些同构与各代数运算相容。
\end{proposition}

\begin{proof}
	在 \(p\) 处，拉回丛的纤维就是带基点标记的 \(E_{F(p)}\)。因此第一项同构把 \((p,e\otimes h)\) 送到 \((p,e)\otimes(p,h)\)，其余各项由同一个纤维的对偶、外幂等恒等识别给出。它们与线性换标架相容。在一张 \(E,H\) 的平凡化图拉回后的图中，上述映射在纤维分量上都是固定线性同构，所以光滑，逆也光滑。截面的局部系数从 \(s^a(y)\) 变为 \(s^a(F(x))\)，故光滑；逐点定义保证它与所列运算相容。
\end{proof}

例如，\(N\) 上的向量场 \(Y\) 总能给出 \(F^*\mathrm{T}N\) 的截面 \(p\mapsto(p,Y_{F(p)})\)，但这还是位于目标切空间的向量。为把 \(M\) 的向量送到那里，可以用覆盖恒等基映射的丛同态
\[
	\widehat{\mathrm{d}F}:\mathrm{T}M\longrightarrow F^*\mathrm{T}N,
	\qquad v_p\longmapsto(p,\mathrm{d}F_pv_p).
\]
\cref{g1:prop:differential-chain-rule}及第十一章的切丛坐标表明它光滑。其对象关系是
\[
	\begin{tikzcd}[column sep=large,row sep=large]
		\mathrm{T}M \arrow[r,"\widehat{\mathrm{d}F}"] \arrow[d] & F^*\mathrm{T}N \arrow[r] \arrow[d] & \mathrm{T}N \arrow[d] \\
		M \arrow[r,"\operatorname{id}_M"'] & M \arrow[r,"F"'] & N.
	\end{tikzcd}
\]
右上箭头忘记拉回纤维的基点标记，三条竖箭头均为丛投影。

\begin{proposition}[协变张量的拉回]
	\label{g1:prop:covariant-tensor-pullback}
	设 \(F:M\to N\)、\(G:N\to P\) 为光滑映射，各流形可带边界，\(T\) 是 \(N\) 上的光滑协变 \(k\) 阶张量。公式
	\[
		(F^*T)_p(v_1,\ldots,v_k)
		=T_{F(p)}(\mathrm{d}F_pv_1,\ldots,\mathrm{d}F_pv_k)
	\]
	定义 \(M\) 上的光滑协变 \(k\) 阶张量。拉回保持张量积、对称性、交替性及楔积，且对 \(P\) 上的同类张量满足 \((G\circ F)^*=F^*\circ G^*\)。当 \(T=g\) 为 Riemann 度量时，\(F^*g\) 半正定；它正定当且仅当 \(F\) 为浸入。
\end{proposition}

\begin{proof}
	先把 \(T\) 拉到 \(F^*((\mathrm{T}^*N)^{\otimes k})\)，再在每个输入中复合 \(\widehat{\mathrm{d}F}\)，便得到所给逐点公式。局部展开时，其系数由 \(T\) 的系数与 \(F\) 复合，再乘上 \(F\) 的一阶偏导数并有限求和，因此光滑。交换或交替输入不改变这种复合，故对称性和交替性保留。张量积与楔积的相容性由各自的多线性公式逐项得到；复合公式来自微分的链式法则。

	度量情形中，\((F^*g)_p(v,v)=g_{F(p)}(\mathrm{d}F_pv,\mathrm{d}F_pv)\geq0\)，且等于零当且仅当 \(\mathrm{d}F_pv=0\)。所以它在每点正定，恰好等价于每点微分单射。
\end{proof}

例如抛物线参数 \(F(t)=(t,t^2)\) 将 Euclid 度量拉回为 \((1+4t^2)\,\mathrm{d}t\otimes\mathrm{d}t\)，处处正定。投影 \(P(x,y)=x\) 却把直线度量拉回为 \(\mathrm{d}x\otimes\mathrm{d}x\)，在 \(\partial_y\) 方向为零。这两个计算分别体现浸入保留全部输入方向、淹没可以忘掉方向。

\ProseParagraphBreak

对向量场不能直接照抄同一公式。若要从 \(N\) 上的 \(Y\) 得到 \(M\) 上的 \(X\)，至少要解 \(\mathrm{d}F_pX_p=Y_{F(p)}\)。对投影 \(P:\mathbb R^2\to\mathbb R\) 与 \(Y=\partial_t\)，\(X=\partial_x+b(x,y)\partial_y\) 都是解，没有指定唯一选择；对包含 \(I(t)=(t,0)\) 与目标向量场 \(\partial_y\)，则没有解。相比之下，协变张量只需让目标观测复合已有的 \(\mathrm{d}F\)，所以总能拉回。

\begin{proposition}[可逆微分下的混合张量]
	\label{g1:prop:mixed-tensor-pullback}
	设 \(F:M\to N\) 为光滑映射，两流形可带边界，并假设每个 \(\mathrm{d}F_p\) 都是线性同构。则 \(N\) 上任意光滑 \((r,s)\) 型张量场都可按
	\[
		(F^*T)_p=
		\bigl((\mathrm{d}F_p)^{-1}\bigr)^{\otimes r}
		\otimes\bigl((\mathrm{d}F_p)^*\bigr)^{\otimes s}
		\bigl(T_{F(p)}\bigr)
	\]
	拉回为 \(M\) 上的同型张量场。该操作与张量积和收缩相容；两个满足条件的映射复合时，拉回仍满足反向复合公式。特别地，微分同胚满足这些条件。
\end{proposition}

\begin{proof}
	\(\mathrm{d}F\) 的局部矩阵光滑且可逆，故逆矩阵光滑，公式因而定义光滑张量场。收缩中一次切向量与一次协变向量的配对变成 \((\alpha\circ\mathrm{d}F)((\mathrm{d}F)^{-1}v)=\alpha(v)\)，所以收缩不变。张量积相容性由定义直接得到，复合公式则使用链式法则以及复合线性同构的逆按反向顺序复合。
\end{proof}

密度还需要单独辨析。密度丛本身作为向量丛，可以沿任意 \(F\) 拉回；但若要把目标的顶维密度变成源上的同维密度，公式中会出现 \(|\det \mathrm{D}F|\)。在局部微分同胚下，行列式处处非零，绝对值仍光滑，所以这一步没有问题。对一般光滑映射则可能失去光滑性：\(F(t)=t^2\) 会把 \(|\mathrm{d}y|\) 按该逐点规则变成 \(2|t|\,|\mathrm{d}t|\)，它在零点不光滑。\cref{g1:prop:covariant-tensor-pullback}适用于交替形式，不能因此不加条件地套到密度上。

\subsection{结构群约化的模型及其存在问题}
\label{g1:subsec:1204:03}

一个任意标架允许全部可逆线性换基。如果选定内积，就可以只使用正交单位标架，其间换基矩阵只能是正交矩阵。这个观察把“额外几何结构”与“允许怎样换标架”联系起来。这里先用局部标架给出定义；一般主丛的语言将在\chapref{g1:ch:18}系统展开。

\begin{definition}
	\label{g1:def:frame-structure-reduction}
	设 \(E\to M\) 为秩 \(r\geq1\) 的实光滑向量丛，\(H\subset\operatorname{GL}(r,\mathbb R)\) 为矩阵子群。一份 \(H\)-\term[jiegouqunyuehua]{结构群约化}[reduction of structure group]由覆盖上的局部光滑标架 \(e_i\) 给出，要求所有过渡矩阵 \(e_j=e_i g_{ij}\) 都取值于 \(H\)。在同一基点处把 \(e_i(p)\) 右乘全部 \(H\) 元素得到的标架子集，因过渡条件而与 \(i\) 无关；记为 \(P_{H,p}\)。定义同一约化的局部标架，允许在各自图中再右乘光滑的 \(H\) 值函数，并允许细化覆盖。
\end{definition}

局部标架之间的矩阵约定与\cref{g1:def:frame-bundle}一致：标架是同构 \(\mathbb R^r\to E_p\)，右作用为复合同构。在 \(P_{H,p}\) 上，\(H\) 的右作用仍自由且传递。几何结构由这组适配标架确定；限制允许的换基群，正是要求额外结构在换标架后保持不变。

\begin{proposition}[度量与正交标架]
	\label{g1:prop:metric-orthogonal-reduction}
	对 Hausdorff、第二可数光滑流形 \(M\)（可带边界）上的秩 \(r\) 实光滑向量丛 \(E\to M\)，给定丛度量等价于给定 \(\operatorname{O}(r)\) 结构群约化。这样的约化总存在。
\end{proposition}

\begin{proof}
	若给定度量，在任一局部标架 \(e_1,\ldots,e_r\) 上逐步定义
	\[
		u_k=e_k-\sum_{j<k}g(e_k,v_j)v_j,
		\qquad v_k=\frac{u_k}{\sqrt{g(u_k,u_k)}}.
	\]
	先前的 \(v_j\) 张成 \(e_1,\ldots,e_{k-1}\) 的空间，原标架线性无关，所以 \(u_k\neq0\)。分母因正定性处处为正，各步光滑，最终得到局部正交单位标架。两个这样的标架之间的矩阵 \(A\) 满足 \(A^{\mathsf T}A=I\)，所以属于 \(\operatorname{O}(r)\)。这些标架的 \(\operatorname{O}(r)\) 轨道就是全部正交单位标架。

	反向给定这样的约化，在每个所选标架中规定度量矩阵为 \(I\)。因为过渡矩阵正交，不同图中的规定一致，得到整体光滑正定度量。两种构造互逆。最后由\cref{g1:thm:bundle-metric-existence}，度量总存在，故约化也总存在。
\end{proof}

\begin{proposition}[定向与体积形式的标架]
	\label{g1:prop:orientation-volume-reductions}
	设 \(E\to M\) 为秩 \(r\geq1\) 的实光滑向量丛。连续选择每个纤维的定向，且在局部标架中选择一致，等价于 \(\operatorname{GL}^+(r,\mathbb R)\) 约化，其中上标 \(+\) 表示行列式为正。一个处处非零的光滑截面 \(\nu\in\Gamma(\det E^*)\) 等价于 \(\operatorname{SL}(r,\mathbb R)\) 约化。若底流形 Hausdorff、第二可数，可带边界，则这种 \(\nu\) 存在当且仅当 \(E\) 可定向。
\end{proposition}

\begin{proof}
	给定定向，局部取正向标架；两正向标架之间行列式为正。反向由正行列式的过渡保证各图的定向一致，所以第一项等价成立。

	给定非零 \(\nu\)，在一个局部标架上设 \(a=\nu(e_1,\ldots,e_r)\)，它是处处非零光滑函数。把 \(e_1\) 换成 \(e_1/a\)，其余不动，即得满足 \(\nu(e_1,\ldots,e_r)=1\) 的标架。这些标架之间的行列式为一，给出 \(\operatorname{SL}\) 约化。反向在约化标架中规定 \(\nu(e_1,\ldots,e_r)=1\)，过渡矩阵行列式为一使规定相容，故得到非零光滑顶次形式，且两种构造互逆。

	非零顶次形式将使其取值为正的标架定义为正向，因而给出定向。若已有定向，先由\cref{g1:thm:bundle-metric-existence}取度量，再取正向正交单位标架；把它们的顶次对偶楔积定义为 \(\nu\)。两这样的标架之间矩阵既正交又有正行列式，故行列式为一，顶次形式在交叠处相同。这就构造出所需 \(\nu\)。
\end{proof}

度量总给出正密度，而不必给出带符号的非零顶次形式。差别在于，密度的转换使用绝对行列式，反向标架不造成符号冲突。因而“有处处正的体积密度”不能用于证明可定向。

\begin{proposition}[非退化交替形式的标架]
	\label{g1:prop:symplectic-frame-reduction}
	设 \(E\to M\) 为秩 \(2m\) 的实光滑向量丛，\(m\geq1\)。给定光滑、逐纤维非退化的交替双线性形式 \(\omega\)，等价于过渡矩阵取值于
	\[
		\operatorname{Sp}(2m,\mathbb R)
		=\{A\in\operatorname{GL}(2m,\mathbb R):A^{\mathsf T}J_0A=J_0\},
		\qquad J_0=\begin{pmatrix}0&I_m\\-I_m&0\end{pmatrix}
	\]
	的结构群约化。这样的 \(\omega\) 总给出非零顶次形式 \(\omega^m/m!\)，因此使 \(E\) 可定向。
\end{proposition}

\begin{proof}
	需要证明适配标架可以光滑选择，而不只是在每个点分别选出一个基。在一点附近先取一个非零局部截面 \(u_1\)。非退化性保证在该点某个标架向量 \(w\) 满足 \(\omega(u_1,w)\neq0\)，缩小邻域后这个函数仍非零，令 \(v_1=w/\omega(u_1,w)\)，就有 \(\omega(u_1,v_1)=1\)。

	对任意局部向量 \(z\)，令
	\[
		z'=z+\omega(v_1,z)u_1-\omega(u_1,z)v_1.
	\]
	直接代入得 \(\omega(u_1,z')=\omega(v_1,z')=0\)。因此这个光滑投影的像为两者的共同交替正交空间 \(W\)，核为 \(\operatorname{span}\{u_1,v_1\}\)，秩恒为 \(2m-2\)。由\cref{g1:thm:constant-rank-bundle-morphism}，\(W\) 是光滑子丛。\(\omega|_W\) 非退化：若 \(z\in W\) 与 \(W\) 中所有向量配对为零，它又与 \(u_1,v_1\) 配对为零，因直和分解而与整个纤维配对为零，所以 \(z=0\)。在 \(W\) 上重复，得到光滑标架 \(u_1,\ldots,u_m,v_1,\ldots,v_m\)，其中 \(\omega(u_i,v_j)=\delta_{ij}\)，同类向量间配对为零。

	在这些标架中 \(\omega\) 的矩阵为 \(J_0\)，故过渡矩阵满足所列条件。反向，在这样的约化标架中规定矩阵为 \(J_0\)，便粘合成光滑非退化交替形式。最后在对偶标架中写 \(\omega=\sum_i u^i\wedge v^i\)。各个二次项两两交换，各自平方为零，故
	\[
		\frac{\omega^m}{m!}=u^1\wedge v^1\wedge\cdots\wedge u^m\wedge v^m\neq0.
	\]
	由\cref{g1:prop:orientation-volume-reductions}得到定向。
\end{proof}

另一个常见的逐点结构，是实向量丛上的光滑丛自同态 \(J\)，满足 \(J^2=-I\)。规定 \((a+ib)v=av+bJv\)，就把每个实纤维变成复向量空间。在切丛上，这样的 \(J\) 称为\term[jihufujiegou]{几乎复结构}[almost complex structure]；此处只要求代数恒等式与光滑性，尚没有要求特殊坐标的存在。

\begin{proposition}[复线性标架]
	\label{g1:prop:complex-frame-reduction}
	设 \(\pi:E\to M\) 为有限秩实光滑向量丛，底流形 Hausdorff、第二可数，可带边界。若有覆盖 \(\mathrm{id}_M\) 的光滑实线性丛自同态 \(J:E\to E\)，即 \(\pi\circ J=\pi\) 且每个 \(J_p:E_p\to E_p\) 实线性，并满足 \(J_p^2=-\mathrm{id}_{E_p}\)，则秩为偶数 \(2m\)，且存在局部光滑复标架；相应实标架的过渡属于 \(\operatorname{GL}(m,\mathbb C)\) 在 \(\operatorname{GL}(2m,\mathbb R)\) 中的标准实表示。反向，这样的约化定义唯一的 \(J\)。该结构总使 \(E\) 可定向。
\end{proposition}

\begin{proof}
	每个纤维通过上述复数乘法成为有限维复向量空间，实维数因此为偶数。为光滑地选基，先取任意丛度量 \(g\)，定义
	\[
		h(v,w)=g(v,w)+g(Jv,Jw).
	\]
	这是光滑正定度量，且 \(h(Jv,Jw)=h(v,w)\)。一个局部非零截面 \(u_1\) 与 \(Ju_1\) 实线性无关：若 \(Ju_1=a u_1\)，再作用 \(J\) 就得 \(-u_1=a^2u_1\)，矛盾。它们张成一个 \(J\) 不变的秩二子丛，其 \(h\) 正交补也是 \(J\) 不变的，因为 \(h(Jv,w)=-h(v,Jw)\)。由\cref{g1:prop:orthogonal-complement-bundle}正交补光滑，在其上重复选择，得到 \(u_1,Ju_1,\ldots,u_m,Ju_m\) 的局部光滑实标架，等价于复标架 \(u_1,\ldots,u_m\)。

	两个复标架之间是可逆复矩阵，视为实矩阵时与标准复乘法相交换。反向若过渡矩阵都复线性，在每张标架中规定 \(J\) 为乘以 \(i\)，交叠处便相容，给出所需 \(J\)。两项构造互逆。

	若复矩阵为 \(A+iB\)，按实部、虚部分组的实矩阵为 \(\begin{psmallmatrix}A&-B\\B&A\end{psmallmatrix}\)。在复数上改变基，把实部与虚部组合为 \(v+iw\)、\(v-iw\)，该矩阵变成对角块 \(A+iB\)、\(A-iB\)，所以它的实行列式为 \(|\det_{\mathbb C}(A+iB)|^2>0\)。所有复标架因而给出一致实定向。
\end{proof}

这里的存在性已经出现了整体限制。令 \(L\to S^1\) 为\cref{g1:exm:mobius-line-bundle}的 Möbius 实线丛，并取 \(E=L\oplus\underline{\mathbb R}\)。沿圆周绕一周，\(E\) 的标架过渡可表示为 \(\operatorname{diag}(-1,1)\)，所以其行列式线丛 \(\det E\) 与 \(L\) 相同。若 \(E\) 可定向，取度量并选正向单位体积向量，就会得到 \(\det E\) 的处处非零截面；但 \(L\) 没有这样的截面。因此 \(E\) 不可定向。它虽然秩为二，仍不能有上面两类结构中的任何一种。这说明偶数秩只是纤维上的必要条件，不能代替过渡函数的整体相容。

\ProseParagraphBreak

单位分解不能消除这些障碍。在二维向量空间上，\(\omega\) 与 \(-\omega\) 都可非退化，但等权平均为零；\(J\) 与 \(-J\) 都满足平方为 \(-I\)，平均后却不再满足。相反，正定双线性型的非负加权和，只要权重和为一，就仍正定。这正是丛度量总能粘合、而其他结构不能照搬同一证明的原因。即使某种逐点结构已存在，它是否满足进一步的微分条件，仍须另行研究。
% 纲要条目排为编号小节标题，并进入目录与书签；说明段仍保留于纲要注释中。


\begin{chaptersummary}
	向量丛把局部乘积通过光滑可逆矩阵粘合；三重交叠上的余循环条件保证识别相容。一个整体标架等价于整体平凡化，局部标架则总存在。丛同态的核与像要成为固定秩光滑子丛，需要相应的常秩条件；拉回改变底空间，商丛则逐纤维消去指定子丛，两者都必须同时构造局部平凡化。

	\ProseParagraphBreak

	张量积将多线性运算改写为线性运算，收缩来自向量与协变向量的求值。外幂的行列式配对确定了楔积的归一化，交换次数为 \(p,q\) 的形式会产生 \((-1)^{pq}\)。分量随标架改变，完整配对得到的数却不变。最高次形式记录符号，密度使用绝对行列式；度量总产生正密度，产生相应的体积形式还需要定向。

	\ProseParagraphBreak

	光滑单位分解由紧穷竭、局部有限的坐标细化、截断函数和归一化构成。局部有限性使各点附近的求和真正成为有限和；支集控制使局部对象能够光滑补零。它们共同给出闭集上的截面延拓和整体丛度量。正交补进一步分裂短正合丛列，但度量与分裂通常依赖选择，相对延拓也须保留明确的邻域相容条件。

	\ProseParagraphBreak

	张量场是在张量丛中的光滑截面，光滑分量与局部标架无关。任意光滑映射都能拉回向量丛，也能利用微分拉回协变张量；把目标向量变为源向量则一般缺少所需的逆。结构群约化用适配标架组织额外几何结构。正交约化总存在，定向、非退化交替形式和几乎复结构的存在却受整体相容性限制，不能仅凭各纤维有相同的线性模型就断定。
\end{chaptersummary}

\BookExercises

\begin{exercise}\label{g1:ex:mobius-sum-double-pullback}
	令 \(L\to S^1\) 为关系 \((t+2\pi,v)\sim(t,-v)\) 定义的 Möbius 线丛，令 \(f:S^1\to S^1\)，\(f(e^{\mathrm{i}t})=e^{2\mathrm{i}t}\)。
	\begin{enumerate}
		\item 构造 \(L\oplus L\) 的全局标架，并证明它是平凡秩二实向量丛。
		\item 构造 \(f^*L\) 的处处非零光滑截面，并解释这与 \(L\) 非平凡为何不矛盾。
	\end{enumerate}
\end{exercise}
\begin{solution}
	1. 记
	\[
		R(\theta)=\begin{pmatrix}\cos\theta&-\sin\theta\\\sin\theta&\cos\theta\end{pmatrix}.
	\]
	在角度参数 \(t\) 下，用 \(R(t/2)\) 的两列分别作为两个截面的纤维分量，即
	\[
		s_1(e^{\mathrm{i}t})=([t,\cos(t/2)],[t,\sin(t/2)]),\qquad
		s_2(e^{\mathrm{i}t})=([t,-\sin(t/2)],[t,\cos(t/2)]).
	\]
	把 \(t\) 换成 \(t+2\pi\) 后，四个分量都变号；同时两条线丛的粘合也都变号，所以这两组有序对不变。它们因此在圆周上良定义，并且在每个提升圆弧上由光滑函数表示。两截面的分量矩阵为 \(R(t/2)\)，行列式恒为一，故逐点线性无关。\cref{g1:prop:global-frame-triviality}给出平凡化
	\[
		S^1\times\mathbb R^2\longrightarrow L\oplus L,
		\qquad(p,(a,b))\longmapsto a s_1(p)+b s_2(p).
	\]
	其局部逆通过乘 \(R(-t/2)\) 求出，因而也光滑。这个构造用了两条截面，而不只是秩二丛的一条非零截面。

	2. 对 \(f^*L\)，定义
	\[
		s(e^{\mathrm{i}t})=(e^{\mathrm{i}t},[2t,1]).
	\]
	右边第二分量位于 \(e^{2\mathrm{i}t}=f(e^{\mathrm{i}t})\) 上，故确实属于拉回纤维。把 \(t\) 换成 \(t+2\pi\) 时，第二分量变为 \([2t+4\pi,1]=[2t,1]\)，因为绕目标圆周两次产生两次变号。于是截面良定义，局部光滑且处处非零，故给出拉回线丛的全局标架。\(L\oplus L\) 改变了纤维秩，\(f^*L\) 改变了沿底空间环路所经历的粘合；两个结论都没有在原线丛 \(L\) 上构造非零全局截面，因此与其非平凡性相容。
\end{solution}

\begin{exercise}\label{g1:ex:smooth-idempotent-subbundles}
	设 \(M\) 为连通的 Hausdorff、第二可数光滑流形，可带边界，\(E\to M\) 为有限秩实光滑向量丛。若同底光滑丛同态 \(P:E\to E\) 满足 \(P^2=P\)，证明 \(\ker P\) 与 \(\operatorname{im}P\) 是光滑子丛，且加法给出自然丛同构
	\[
		\ker P\oplus\operatorname{im}P\simeq E.
	\]
	再对 \(S^2\times\mathbb R^3\) 上的 \(P_p(v)=v-\langle v,p\rangle p\) 求核、像和相应商丛，并写出这个分解的两个投影。
\end{exercise}
\begin{solution}
	对任一纤维向量 \(v\)，有
	\[
		v=(v-Pv)+Pv,\qquad P(v-Pv)=0.
	\]
	而若 \(w\in\ker P\cap\operatorname{im}P\)，写成 \(w=Pu\) 后有 \(w=P^2u=Pw=0\)。所以逐纤维有直和分解 \(E_p=\ker P_p\oplus\operatorname{im}P_p\)。在适应该分解的一组线性基下，\(P_p\) 的矩阵是 \(\operatorname{diag}(0,I_k)\)，故 \(\operatorname{tr}P_p=\operatorname{rank}P_p\)。迹不依赖基，因为 \(\operatorname{tr}(AB)=\operatorname{tr}(BA)\)，所以局部矩阵的迹给出一个整体光滑函数。

	这个函数只取整数值，故局部常值：在一点附近把它与该点值的差控制在 \(1/2\) 以内，整数性就迫使两者相同。底空间连通，因而秩在整个 \(M\) 上恒定。由\cref{g1:thm:constant-rank-bundle-morphism}，核与像都是光滑子丛。加法映射逐纤维同构，且逆明确为
	\[
		E\longrightarrow\ker P\oplus\operatorname{im}P,\qquad
		v\longmapsto(v-Pv,Pv).
	\]
	这些映射在子丛局部标架中光滑，故确为光滑丛同构。这里投影 \(P\) 本身已经指定了补丛，所以不需要另选内积。

	对题给球面例，\(|p|=1\) 蕴含
	\[
		P_p^2=P_p,\qquad
		\ker P_p=\mathbb Rp,\qquad
		\operatorname{im}P_p=p^\perp.
	\]
	核丛由全局非零截面 \(p\mapsto p\) 张成，故是平凡线丛。像丛就是 \(\mathrm{T}S^2\) 在环境 \(\mathbb R^3\) 中的实现：对球面内曲线求导 \(\langle\gamma,\gamma\rangle=1\)，得到速度垂直于 \(p\)；反之，对 \(0\ne v\perp p\)，曲线
	\[
		\gamma(s)=\cos(s|v|)p+\frac{\sin(s|v|)}{|v|}v
	\]
	位于球面且满足 \(\gamma(0)=p\)、\(\gamma'(0)=v\)，零速度由常曲线实现。球面的局部图是一个坐标关于另两个坐标的光滑图像，所以其包含映射微分单射；上述速度计算遂识别了整个切空间。于是
	\[
		(S^2\times\mathbb R^3)/\ker P\simeq \mathrm{T}S^2,
		\qquad [v]\longmapsto v-\langle v,p\rangle p.
	\]
	所需两个投影分别为 \((I-P)_p(v)=\langle v,p\rangle p\) 与 \(P_p(v)=v-\langle v,p\rangle p\)，并给出 \((S^2\times\mathbb R)\oplus \mathrm{T}S^2\simeq S^2\times\mathbb R^3\)。这里切丛与一条平凡线丛的直和是平凡丛，并未据此断言切丛本身平凡。
\end{solution}

\begin{exercise}\label{g1:ex:latitude-pullback-bundle}
	对固定 \(\alpha\in[0,\pi]\)，定义
	\[
		f_\alpha:S^1\longrightarrow S^2,\qquad
		f_\alpha(e^{\mathrm{i}t})=(\sin\alpha\cos t,\sin\alpha\sin t,\cos\alpha).
	\]
	将球面切向量识别为环境中垂直于基点的向量。构造拉回丛 \(f_\alpha^*\mathrm{T}S^2\) 的全局标架；用它写出微分给出的同底丛同态 \(\mathrm{T}S^1\to f_\alpha^*\mathrm{T}S^2\)，并分别讨论 \(0<\alpha<\pi\) 与 \(\alpha=0,\pi\) 时的核、像和商丛。
\end{exercise}
\begin{solution}
	令 \(p(t)=f_\alpha(e^{\mathrm{i}t})\)，定义环境向量
	\[
		u_1(t)=(-\sin t,\cos t,0),\qquad
		u_2(t)=(-\cos\alpha\cos t,-\cos\alpha\sin t,\sin\alpha).
	\]
	直接计算得到 \(\langle p,u_1\rangle=\langle p,u_2\rangle=0\)、\(|u_1|=|u_2|=1\)、\(\langle u_1,u_2\rangle=0\)。因此两向量在 \(\mathrm{T}_{p(t)}S^2\) 中组成基。它们都是 \(2\pi\) 周期的光滑函数，所以
	\[
		e_a(e^{\mathrm{i}t})=(e^{\mathrm{i}t},u_a(t)),\qquad a=1,2,
	\]
	是拉回丛的全局光滑标架。这个公式在 \(\alpha=0,\pi\) 时仍成立：虽然 \(p(t)\) 此时为常点，所选基可以在这一固定平面内随 \(t\) 旋转。

	在源圆周上取全局切标架 \(\tau(t)=(-\sin t,\cos t)\)。微分作用由参数曲线求导得到
	\[
		\mathrm{d}(f_\alpha)_{e^{\mathrm{i}t}}(a\tau(t))
		=a\sin\alpha\,u_1(t).
	\]
	因此改写成到拉回丛的同底丛同态后，其在标架 \(\tau\) 与 \((e_1,e_2)\) 下的矩阵为列 \((\sin\alpha,0)^{\mathsf T}\)。若 \(0<\alpha<\pi\)，这个映射单射，核为零子丛，像为 \(e_1\) 张成的平凡线丛，而目标除以像的商丛由 \([e_2]\) 给出全局标架，亦为平凡线丛。具体商坐标为 \([a e_1+b e_2]\mapsto b\)。

	若 \(\alpha=0\) 或 \(\pi\)，微分为零，核是整个 \(\mathrm{T}S^1\)，像是零子丛；目标除以像仍为秩二的平凡拉回丛。拉回纤维的维数始终为二，微分的秩却可以为一或零。因而拉回丛不能被定义为微分的像，也不等同于源流形的切丛。
\end{solution}

\begin{exercise}\label{g1:ex:quotient-complement-shears}
	在 \(\mathbb R\) 上令 \(E=\mathbb R\times\mathbb R^2\)，\(F_t=\mathbb R(1,0)\)，并把商映射写成 \(q(t,a,b)=(t,b)\)。求出 \(q\) 的全部光滑线性右逆及其像补丛。再考虑丛自同构
	\[
		T_c(t,a,b)=(t,a+c(t)b,b),\qquad c\in C^\infty(\mathbb R).
	\]
	证明每个 \(T_c\) 在 \(F\) 上为恒等，且在 \(E/F\) 上诱导恒等；是否存在一个补丛在所有这些自同构下都不变？
\end{exercise}
\begin{solution}
	若 \(j:E/F\to E\) 是光滑丛同态且 \(qj=\operatorname{id}\)，其在每条纤维上的线性表达必为
	\[
		j_h(t,b)=(t,h(t)b,b).
	\]
	第一分量的系数 \(h(t)\) 由 \(j(t,1)\) 得到，所以光滑；反过来，任意光滑 \(h\) 都给出这样的右逆。因此全部右逆由光滑函数 \(h\) 参数化，其像为
	\[
		C_h|_t=\mathbb R(h(t),1).
	\]
	它由一个处处非零的光滑截面张成，且每个 \((a,b)\) 有唯一分解
	\[
		(a,b)=(a-h(t)b,0)+(h(t)b,b),
	\]
	所以 \(E=F\oplus C_h\)。反之，任一光滑补丛 \(C\) 上的 \(q|_C\) 逐纤维可逆，由\cref{g1:prop:bundle-morphism-matrices}有光滑逆，故也来自某个 \(h\)。

	矩阵 \(\begin{psmallmatrix}1&c(t)\\0&1\end{psmallmatrix}\) 的逆由把 \(c\) 换成 \(-c\) 得到，因此 \(T_c\) 是光滑丛自同构。若 \(b=0\)，它保持 \((a,0)\) 不变，故固定 \(F\) 的每个向量；第二纤维坐标 \(b\) 不变，故诱导商上的恒等。同时
	\[
		T_c\circ j_h=j_{h+c},\qquad T_c(C_h)=C_{h+c}.
	\]
	取 \(c\equiv1\) 就有 \(C_{h+1}\ne C_h\)：若两条直线相同，将生成向量的第二分量都归一为一，就会得到 \(h+1=h\)，矛盾。因此不存在在所有这些自同构下不变的补丛。这个例子中每个补丛都能整体光滑选出，但子丛与商映射并没有指定一个对全部保留这些数据的自同构不变的选择。
\end{solution}

\begin{exercise}
	\label{g1:exr:1202-tensor-rank}
	设 \(V,W\) 为有限维实向量空间，\(T\in W\otimes V^*\)，其对应的线性映射为 \(L=\Phi(T)\)。证明：把 \(T\) 写成纯张量之和所需的最少项数，等于 \(\operatorname{rank}L\)，约定零张量允许空和。由此判断在 \(V=W=\mathbb R^2\) 中，\(e_1\otimes e^1+e_2\otimes e^2\) 能否写成一个纯张量。再给出一个非零而收缩为零的 \(V\otimes V^*\) 中的元素，解释“迹为零”为什么比“张量为零”弱。
\end{exercise}
\begin{solution}
	若 \(T=\sum_{j=1}^k w_j\otimes\alpha_j\)，则
	\[
		L(v)=\sum_{j=1}^k\alpha_j(v)w_j,
	\]
	故 \(\operatorname{im}L\subseteq\operatorname{span}\{w_1,\ldots,w_k\}\)，得到 \(\operatorname{rank}L\le k\)。反过来，设 \(\operatorname{rank}L=r>0\)，取像空间的一组基 \(w_1,\ldots,w_r\)。对每个 \(v\)，唯一写 \(L(v)=\sum_{j=1}^r\alpha_j(v)w_j\)。由 \(L\) 线性及系数表示唯一，\(\alpha_j(v+v')=\alpha_j(v)+\alpha_j(v')\)、\(\alpha_j(cv)=c\alpha_j(v)\)，所以 \(\alpha_j\in V^*\)。于是 \(L=\Phi(\sum_jw_j\otimes\alpha_j)\)，由 \(\Phi\) 单射得到 \(T=\sum_jw_j\otimes\alpha_j\)。若 \(r=0\)，则 \(L=0\)，同样由同构知 \(T=0\)。最少项数因此恰为 \(r\)。

	给出的二维张量对应恒等映射，秩为 \(2\)，故不能写成一个纯张量。另一方面，\(S=e_1\otimes e^2\) 将 \(e_2\) 送到 \(e_1\)，所以非零；但其收缩是 \(e^2(e_1)=0\)。收缩只保留一个实数，是从整个张量空间到 \(\mathbb R\) 的线性观测，会消去许多非零元素；这一例中，它消去的是一个非零的非对角分量。
\end{solution}

\begin{exercise}
	\label{g1:exr:1202-quadratic-symmetric}
	设 \(V\) 为有限维实向量空间，\(B\in\operatorname{Sym}^2(V^*)\) 按本节约定视为对称双线性形式，令 \(q(v)=B(v,v)\)。证明 \(q\) 唯一决定 \(B\)，并写出恢复公式。反过来，在一组基的坐标中给定任意实齐次二次多项式，证明它来自唯一这样的 \(B\)。对
	\[
		q(xe_1+ye_2)=3x^2+4xy+2y^2
	\]
	写出 \(B\) 的矩阵及它在 \(\operatorname{Sym}^2(V^*)\) 中的表达，并计算 \(B(e_1+e_2,e_1-e_2)\)。
\end{exercise}
\begin{solution}
	双线性与对称性给出
	\[
		q(u+v)=B(u,u)+2B(u,v)+B(v,v),
	\]
	所以
	\[
		B(u,v)=\frac12\bigl(q(u+v)-q(u)-q(v)\bigr).
	\]
	这先证明唯一性。若给定齐次二次多项式
	\[
		q(x)=\sum_i a_i(x^i)^2+\sum_{i<j}b_{ij}x^ix^j,
	\]
	取矩阵 \(G\) 的对角元为 \(a_i\)，非对角元 \(G_{ij}=G_{ji}=b_{ij}/2\)，则 \(B(u,v)=[u]^{\mathsf T}G[v]\) 是对称双线性形式，并满足 \(B(v,v)=q(v)\)，证明存在。恢复公式只使用向量加法与实数值，因此换基并不改变它定义的对象。

	本题中
	\[
		G=\begin{pmatrix}3&2\\2&2\end{pmatrix},\qquad
		B=3e^1\odot e^1+4e^1\odot e^2+2e^2\odot e^2.
	\]
	中间项的系数为 \(4\)，因为 \((e^1\odot e^2)(u,v)=\tfrac12(e^1(u)e^2(v)+e^1(v)e^2(u))\)，对应矩阵的两个非对角元都是 \(1/2\)。所以系数 \(4\) 正好产生两个矩阵元 \(2\)。最后，
	\[
		B(e_1+e_2,e_1-e_2)
		=(1,1)\begin{pmatrix}3&2\\2&2\end{pmatrix}\binom{1}{-1}=1.
	\]
\end{solution}

\begin{exercise}
	\label{g1:exr:1202-wedge-insertion}
	设 \(V=\mathbb R^4\)，基与对偶基分别为 \(e_i,e^i\)，令
	\[
		\omega=e^1\wedge e^2+e^3\wedge e^4.
	\]
	计算 \(\omega\wedge\omega\) 以及 \(\iota_v\omega\)，其中 \(v=\sum_iv^ie_i\)。证明 \(v\mapsto\iota_v\omega\) 是 \(V\to V^*\) 的线性同构，且 \(\omega\) 不能写成两个一形式的楔积。最后计算 \(\iota_{e_1}(\omega\wedge\omega)\)，同时用插入的乘法公式核对符号与系数。
\end{exercise}
\begin{solution}
	两个平方项含重复的一形式，均为零。两个交叉项交换的是两个二次块，符号为 \((-1)^4=1\)，故
	\[
		\omega\wedge\omega=2e^1\wedge e^2\wedge e^3\wedge e^4.
	\]
	按插入固定第一位置的约定，
	\[
		\iota_v\omega=v^1e^2-v^2e^1+v^3e^4-v^4e^3.
	\]
	若此一形式为零，对偶基的独立性依次给出 \(v^2=v^1=v^4=v^3=0\)，所以该线性映射单射。两个空间都为四维，故为同构；也可直接把任意一形式 \(\sum_i a_ie^i\) 的原像写成 \(a_2e_1-a_1e_2+a_4e_3-a_3e_4\)。

	若 \(\omega=\alpha\wedge\beta\)，那么
	\[
		\omega\wedge\omega=\alpha\wedge\beta\wedge\alpha\wedge\beta=0,
	\]
	因为移项后出现两个相同的一形式。这与已计算出的非零最高次形式矛盾，所以 \(\omega\) 不是纯楔积。外幂由纯楔积张成，并不意味着每个元素都是纯楔积。

	由最高次形式的表达直接插入，得到
	\[
		\iota_{e_1}(\omega\wedge\omega)=2e^2\wedge e^3\wedge e^4.
	\]
	另一方面，\(\iota_{e_1}\omega=e^2\)，且 \(\deg\omega=2\)，故乘法公式给出
	\[
		\iota_{e_1}(\omega\wedge\omega)
		=e^2\wedge\omega+\omega\wedge e^2
		=2e^2\wedge\omega
		=2e^2\wedge e^3\wedge e^4.
	\]
	第二步交换一次与二次形式，符号为 \((-1)^{1\cdot2}=1\)，正好得到系数 \(2\)。
\end{solution}

\begin{exercise}
	\label{g1:exr:1202-density-orientation}
	在 \(V=\mathbb R^2\) 的标准基 \(e=(e_1,e_2)\) 中，取内积矩阵及新基
	\[
		G=\begin{pmatrix}2&1\\1&2\end{pmatrix},\qquad
		e'=eA,\qquad A=\begin{pmatrix}1&1\\0&-2\end{pmatrix}.
	\]
	证明 \(G\) 正定，计算新基中的内积矩阵，以及 \(\mu_g(e)\)、\(\mu_g(e')\)。取使标准基正向的定向，写出体积形式 \(\omega\) 和体积密度 \(\mu_g\) 在新旧两组对偶基中的表达。若把定向反过来，哪些对象发生改变？
\end{exercise}
\begin{solution}
	对非零 \((x,y)\)，有
	\[
		(x,y)G\binom{x}{y}=2x^2+2xy+2y^2=(x+y)^2+x^2+y^2>0,
	\]
	故 \(G\) 正定。直接相乘得
	\[
		G'=A^{\mathsf T}GA=\begin{pmatrix}2&0\\0&6\end{pmatrix}.
	\]
	因此 \(\mu_g(e)=\sqrt{\det G}=\sqrt3\)，而 \(\mu_g(e')=\sqrt{12}=2\sqrt3\)，也等于 \(|\det A|\mu_g(e)\)。

	记 \(\varepsilon=e^1\wedge e^2\)、\(\varepsilon'=e'^1\wedge e'^2\)。由于 \(\det A=-2\)，有 \(\varepsilon'=-\tfrac12\varepsilon\)。标准基虽不必是正交单位基，但正向体积形式在它上的值必须为正且绝对值等于 \(\sqrt3\)，所以
	\[
		\omega=\sqrt3\,\varepsilon=-2\sqrt3\,\varepsilon',\qquad
		\mu_g=\sqrt3\,|\varepsilon|=2\sqrt3\,|\varepsilon'|.
	\]
	新基为负向基，故体积形式的系数为负，而体积密度的系数仍正。反转定向会把 \(\omega\) 变为 \(-\omega\)，不改变 \(g\)、\(\mu_g\)、两组固定基或它们的变换矩阵。体积密度因此不能单独判断所选的是哪一个定向。
\end{solution}

\begin{exercise}
	\label{g1:exr:1202-polar-rotation}
	在平面一个 \(r>0\) 的极坐标扇区内，固定一点。令 \(g\) 为 Euclid 内积，\(J\) 为正向旋转 \(\pi/2\) 的线性映射，即 \(J\partial_x=\partial_y\)、\(J\partial_y=-\partial_x\)。求 \(J\) 在基 \((\partial_r,\partial_\theta)\) 中的矩阵与 \((1,1)\) 型张量表达。由这些表达验证 \(J^2=-I\)、\(\operatorname{tr}J=0\) 和 \(g(Ju,Jv)=g(u,v)\)。再求双线性形式 \(\beta(u,v)=g(Ju,v)\)，说明它为何是二次交替形式，并写成 \(\mathrm{d}r,\mathrm{d}\theta\) 的楔积。
\end{exercise}
\begin{solution}
	直角坐标中的矩阵为 \(J_0=\begin{psmallmatrix}0&-1\\1&0\end{psmallmatrix}\)。极坐标换基矩阵可写成
	\[
		A=\begin{pmatrix}\cos\theta&-\sin\theta\\\sin\theta&\cos\theta\end{pmatrix}
		\begin{pmatrix}1&0\\0&r\end{pmatrix}.
	\]
	旋转矩阵与 \(J_0\) 交换，所以
	\[
		J'=A^{-1}J_0A=\begin{pmatrix}0&-r\\1/r&0\end{pmatrix}.
	\]
	等价地，\(J\partial_r=r^{-1}\partial_\theta\)、\(J\partial_\theta=-r\partial_r\)，故
	\[
		J=\frac1r\,\partial_\theta\otimes\mathrm{d}r
		-r\,\partial_r\otimes\mathrm{d}\theta.
	\]
	矩阵平方为 \(-I\)，对角元和为零。极坐标内积矩阵是 \(G'=\operatorname{diag}(1,r^2)\)，相乘得 \((J')^{\mathsf T}G'J'=G'\)，从而验证内积保持不变。也可令 \(u=a\partial_r+b\partial_\theta\)、\(v=c\partial_r+d\partial_\theta\)，直接算出
	\[
		g(Ju,Jv)=(-rb)(-rd)+r^2(a/r)(c/r)=r^2bd+ac=g(u,v).
	\]

	同样，
	\[
		\beta(u,v)=g(-rb\partial_r+(a/r)\partial_\theta,\,
		c\partial_r+d\partial_\theta)=r(ad-bc).
	\]
	交换 \(u,v\) 时此值变号，重复输入时为零，因此它是交替二次形式，并且
	\[
		\beta=r\,\mathrm{d}r\wedge\mathrm{d}\theta
		=\mathrm{d}x\wedge\mathrm{d}y.
	\]
	这里 \(J\) 的矩阵含有 \(r\) 与 \(1/r\)，却仍表示固定的直角旋转；内积的矩阵同时发生相应变化，配对之后得到的有向面积形式与直角坐标中的对象一致。
\end{solution}

\begin{exercise}\label{g1:ex:smooth-partition-on-line}
	令 \(\eta\) 为\cref{g1:lem:smooth-bump-function}中的函数，\(b(t)=\eta(1-t^2)\)。对 \(k\in\mathbb Z\) 定义
	\[
		H(x)=\sum_{k\in\mathbb Z}b(x-k),\qquad
		\psi_k(x)=\frac{b(x-k)}{H(x)}.
	\]
	证明这些公式给出从属于开覆盖 \(O_k=(k-\tfrac32,k+\tfrac32)\) 的光滑单位分解，并证明
	\[
		H(x+1)=H(x),\qquad
		\psi_k(x)=\psi_0(x-k).
	\]
	求在整数点与半整数点处非零的权重，并证明对每个整数 \(j\ge1\)，有 \(\sum_k\psi_k^{(j)}(x)=0\)。
\end{exercise}
\begin{solution}
	函数 \(b\) 光滑，在 \((-1,1)\) 上为正，在外部为零，故 \(\operatorname{supp}b=[-1,1]\)。因此 \(\operatorname{supp}b(\,\cdot-k)=[k-1,k+1]\subset O_k\)。给定 \(x\)，邻域 \((x-\tfrac14,x+\tfrac14)\) 只能与指标落在有限区间 \([x-\tfrac54,x+\tfrac54]\) 内的这些闭区间相交，故支集族局部有限，\(H\) 光滑。每个实数与某个整数的距离不超过 \(1/2\)，相应项严格为正，所以 \(H>0\)。

	归一化后，\(\psi_k\) 光滑、非负、总和为 \(1\)，且其支集与对应的 \(b(x-k)\) 相同，全部从属条件成立。换指标 \(\ell=k-1\) 得
	\[
		H(x+1)=\sum_k b(x+1-k)=\sum_\ell b(x-\ell)=H(x).
	\]
	于是 \(H(x-k)=H(x)\)，给出 \(\psi_k(x)=b(x-k)/H(x-k)=\psi_0(x-k)\)。

	在整数 \(m\) 处，只有 \(b(m-m)=b(0)>0\)，相邻项的自变量为 \(\pm1\)，也为零，故 \(\psi_m(m)=1\)，其余项为零。在 \(m+\tfrac12\) 处，只有指标 \(m,m+1\) 的两项非零；由于 \(b\) 为偶函数，两项相等，权重各为 \(1/2\)。最后，在任意计算点的固定邻域内，恒等式 \(\sum_k\psi_k=1\) 都是有限和恒等式，逐项求 \(j\) 阶导数即得所求。若把覆盖误写为 \((k-1,k+1)\)，支集就不包含于对应开集，不能得到本书定义中的从属关系。
\end{solution}

\begin{exercise}\label{g1:ex:sections-on-closed-discrete-set}
	设 \(M\) 为有限维、Hausdorff 且第二可数的光滑流形，允许有边界，\(E\to M\) 为有限秩实光滑向量丛，\(D\subseteq M\) 是闭离散子集，即 \(D\) 闭，且每个 \(p\in D\) 有一个只与 \(D\) 交于 \(p\) 的开邻域。任意指定向量 \(v_p\in E_p\)。证明存在整体光滑截面 \(s\)，使 \(s(p)=v_p\) 对所有 \(p\in D\) 成立。进一步证明，可要求 \(s\) 紧支，当且仅当集合 \(\{p\in D:v_p\ne0\}\) 有限。
\end{exercise}
\begin{solution}
	在每个 \(p\in D\) 附近选择局部标架，并把邻域缩小到只含一个 \(D\) 中的点。用 \(v_p\) 在该点的分量作常系数线性组合，得到局部光滑截面 \(s_p\)，满足 \(s_p(p)=v_p\)。因为这个邻域与 \(D\) 只交于 \(p\)，它已延拓该邻域内的全部指定值。故 \(p\mapsto v_p\) 正是\cref{g1:def:smooth-section-closed-set}意义下闭集 \(D\) 上的光滑截面，\cref{g1:thm:section-extension}给出整体延拓。

	若整体截面 \(s\) 紧支，令 \(K=\operatorname{supp}s\)。由于 \(D\) 闭，\(D\cap K\) 紧；它又是离散空间，而紧离散空间只能有限，因为由所有单点组成的开覆盖必须有有限子覆盖。因此非零数据的指标集包含于有限集 \(D\cap K\)，也有限。

	反之，设只有 \(p_1,\ldots,p_N\) 的指定值非零。为每个 \(p_j\) 取前面的局部截面 \(s_j\)，其定义域 \(U_j\) 只与 \(D\) 交于 \(p_j\)。对紧闭集 \(\{p_j\}\subset U_j\) 选紧支光滑截断函数 \(\chi_j\)，在 \(p_j\) 附近为 \(1\)，支集包含于 \(U_j\)。将 \(\chi_j s_j\) 补零到 \(M\)，然后取有限和
	\[
		s=\sum_{j=1}^N\chi_j s_j.
	\]
	该截面在 \(p_j\) 取值 \(v_{p_j}\)，在 \(D\) 的其余点为零，支集包含于有限个紧集之并，所以紧。没有非零数据时取零截面即可。
\end{solution}

\begin{exercise}\label{g1:ex:relative-lift-of-section}
	设 \(M\) 为有限维、Hausdorff 且第二可数的光滑流形，允许有边界，\(q:E\to F\) 是有限秩实光滑向量丛之间覆盖 \(\operatorname{id}_M\) 的逐纤维满射同态。证明每个 \(t\in\Gamma(F)\) 都有整体光滑提升 \(s\in\Gamma(E)\)，即 \(q\circ s=t\)。若 \(A\subseteq M\) 闭，且已在 \(A\) 的开邻域 \(U\) 上给定提升 \(s_U\)，证明可选整体提升与 \(s_U\) 在 \(A\) 附近相同。说明两条整体提升之差属于哪个截面空间。
\end{exercise}
\begin{solution}
	由\cref{g1:thm:constant-rank-bundle-morphism}，\(K=\ker q\) 为光滑子丛，且
	\[
		0\longrightarrow K\longrightarrow E\xrightarrow{\ q\ }F\longrightarrow0
	\]
	为短正合丛列。\cref{g1:thm:exact-bundle-splitting}给出光滑右逆 \(j:F\to E\)。于是 \(s_0=j\circ t\) 光滑，且 \(q\circ s_0=t\)。

	在 \(U\) 上，\(d=s_U-s_0|_U\) 满足 \(q(d)=0\)，故逐点落在 \(K\) 中，并且是 \(K\) 的光滑截面。最后这一点可在适配 \(K\) 的局部标架中看出：后面的分量恒为零，前面的分量仍光滑。用\cref{g1:thm:section-extension}把 \(d\) 延拓为 \(K\) 的整体截面 \(\widetilde d\)，保持 \(A\) 附近的值，令
	\[
		s=s_0+\widetilde d.
	\]
	则 \(q(s)=q(s_0)=t\)，且在 \(A\) 附近 \(s=s_U\)。若 \(A\) 紧，还可使改变量 \(\widetilde d\) 紧支。

	若 \(s_1,s_2\) 都提升 \(t\)，则 \(q(s_1-s_2)=0\)，其差属于 \(\Gamma(K)\)；反之，任意给一个提升加上 \(\Gamma(K)\) 中的截面，仍为提升。因此一旦选定一个提升，所有提升恰好由它加上核子丛的任意整体截面得到。
\end{solution}

\begin{exercise}\label{g1:ex:metric-dependent-splitting}
	设 \(M\) 为光滑流形，\(a\in C^\infty(M)\)。在平凡丛 \(E=M\times\mathbb R^2\) 上定义
	\[
		h_{a,p}\bigl((u,v),(u',v')\bigr)
		=(u+a(p)v)(u'+a(p)v')+vv'.
	\]
	证明 \(h_a\) 是丛度量，求第一坐标子丛 \(L=M\times(\mathbb R\times\{0\})\) 的正交补、到 \(L\) 的正交投影，以及由该正交补给出的投影 \(q:E\to M\times\mathbb R\)、\(q(p,u,v)=(p,v)\) 的右逆。进一步证明，\(q\) 的每个光滑丛右逆都可由某个这样的 \(h_a\) 得到。
\end{exercise}
\begin{solution}
	定义显然给出对称双线性型，其矩阵
	\[
		H_a=\begin{pmatrix}1&a\\a&1+a^2\end{pmatrix}
	\]
	光滑。对非零 \((u,v)\)，平方和 \((u+av)^2+v^2\) 严格为正，所以 \(h_a\) 为丛度量。

	向量 \((u,v)\) 与 \((1,0)\) 正交，当且仅当 \(u+av=0\)。因此 \(L^\perp\) 有整体标架 \((-a,1)\)。分解
	\[
		(u,v)=(u+av,0)+v(-a,1)
	\]
	给出正交投影 \(P_a(p,u,v)=(p,u+a(p)v,0)\)。由于 \(q(p,-a(p)t,t)=(p,t)\)，对应的丛右逆为
	\[
		j_a(p,t)=(p,-a(p)t,t).
	\]

	任意逐纤维线性的右逆 \(j\) 都由其在 \(t=1\) 处的值决定；右逆条件使第二坐标必须为 \(1\)，所以 \(j(p,t)=(p,c(p)t,t)\)。丛同态的光滑性等价于 \(c\) 光滑。取 \(a=-c\)，便有 \(j=j_a\)。可见同一条短正合丛列的不同分裂都能在本例中通过改变度量实现；丛列本身并未规定其中哪一个。
\end{solution}

\begin{exercise}
	\label{g1:ex:tensor-field-changing-frame}
	在 \(\mathbb R^2\) 上取 \(X=\partial_x+y\partial_y\)、\(\alpha=\mathrm{d}x+\mathrm{d}y\)，令 \(L=X\otimes\alpha\)，把它视为同态场 \(L(w)=\alpha(w)X\)。换用全局坐标 \(u=x\)、\(v=e^xy\)，求 \(X,\alpha,L\) 的新分量，并在两套坐标中核对 \(\operatorname{tr}L=\alpha(X)\)。
\end{exercise}
\begin{solution}
	坐标变化的逆为 \(x=u\)、\(y=ve^{-u}\)，两个方向都光滑。由 \(X(u)=1\)、\(X(v)=e^xy+ye^x=2v\)，得
	\[
		X=\partial_u+2v\partial_v.
	\]
	又由 \(\mathrm{d}x=\mathrm{d}u\)、\(\mathrm{d}y=e^{-u}\mathrm{d}v-ve^{-u}\mathrm{d}u\)，得
	\[
		\alpha=(1-ve^{-u})\mathrm{d}u+e^{-u}\mathrm{d}v.
	\]
	向量与协变向量按不同方向变换，合起来给出同态的矩阵。在旧、新坐标基中分别为
	\[
		[L]_{x,y}=\begin{pmatrix}1&1\\y&y\end{pmatrix},\qquad
		[L]_{u,v}=\begin{pmatrix}1-ve^{-u}&e^{-u}\\2v(1-ve^{-u})&2ve^{-u}\end{pmatrix}.
	\]
	所以旧坐标中迹为 \(1+y\)，新坐标中为 \(1+ve^{-u}\)，正是同一函数。直接求值也有 \(\alpha(X)=1+y=(1-ve^{-u})+2ve^{-u}\)。矩阵分量随标架改变，收缩时换基因子相消，因而迹保持不变。
\end{solution}

\begin{exercise}
	\label{g1:ex:exponential-polar-pullbacks}
	设 \(F:\mathbb R^2\to\mathbb R^2\setminus\{0\}\)，\(F(u,v)=(e^u\cos v,e^u\sin v)\)。在目标上取 Euclid 度量 \(g=\mathrm{d}x\otimes\mathrm{d}x+\mathrm{d}y\otimes\mathrm{d}y\)、一形式 \(\alpha=x\,\mathrm{d}y-y\,\mathrm{d}x\)、二形式 \(\omega=\mathrm{d}x\wedge\mathrm{d}y\)。求三者的拉回，并求 \(F^*g\) 下函数 \(u\) 的梯度。说明即使 \(F\) 不是整体单射，这些构造仍然有意义。
\end{exercise}
\begin{solution}
	有
	\[
		F^*(\mathrm{d}x)=e^u\cos v\,\mathrm{d}u-e^u\sin v\,\mathrm{d}v,
		\qquad
		F^*(\mathrm{d}y)=e^u\sin v\,\mathrm{d}u+e^u\cos v\,\mathrm{d}v.
	\]
	平方的对称交叉项相消，故
	\[
		F^*g=e^{2u}(\mathrm{d}u\otimes\mathrm{d}u+\mathrm{d}v\otimes\mathrm{d}v).
	\]
	把函数 \(x,y\) 同时与 \(F\) 复合，再代入上式中的一形式，得到 \(F^*\alpha=e^{2u}\mathrm{d}v\)。按楔积展开，则 \(F^*\omega=e^{2u}\mathrm{d}u\wedge\mathrm{d}v\)。该度量的逆矩阵为 \(e^{-2u}I\)，所以 \(\operatorname{grad}_{F^*g}u=e^{-2u}\partial_u\)。

	\(F(u,v+2\pi)=F(u,v)\)，因此 \(F\) 不单射；但 \(\det \mathrm{D}F=e^{2u}>0\)，它在每点附近为局部微分同胚。协变张量的拉回本来就不要求单射，拉回度量的正定性只要求微分单射，故本题没有矛盾。若在目标取 \(h(x,y)=\log\sqrt{x^2+y^2}\)，则 \(h\circ F=u\)，并有
	\[
		\mathrm{d}F(\operatorname{grad}_{F^*g}u)
		=e^{-2u}(x,y)=\operatorname{grad}_g h\big|_{F(u,v)},
	\]
	也可直接核对这个局部等距识别下的梯度关系。
\end{solution}

\begin{exercise}
	\label{g1:ex:oriented-plane-complex-structure}
	设 \(E\to M\) 是可定向的秩二实光滑向量丛，固定定向和丛度量 \(g\)。在每张正向正交单位标架 \((e_1,e_2)\) 中规定 \(Je_1=e_2\)、\(Je_2=-e_1\)。证明这定义整体光滑 \(J\)，且 \(J^2=-I\)。若 \(\nu\) 是由定向和 \(g\) 决定的单位体积形式，证明 \(\nu(v,w)=g(Jv,w)\)。最后将度量改为 \(\widetilde g=e^{2f}g\)，其中 \(f\) 光滑，比较两者对应的 \(J\) 与体积形式。
\end{exercise}
\begin{solution}
	两张正向正交单位标架之间的变换矩阵在 \(\operatorname{SO}(2)\) 中。任一这样的矩阵都可写为 \(\begin{psmallmatrix}a&-b\\b&a\end{psmallmatrix}\)，其中 \(a^2+b^2=1\)，它与 \(J_0=\begin{psmallmatrix}0&-1\\1&0\end{psmallmatrix}\) 交换。因此各图中规定的 \(J\) 在交叠处一致，且局部矩阵恒为 \(J_0\)，所以 \(J\) 整体光滑并满足 \(J^2=-I\)。

	对 \(v=ae_1+be_2\)、\(w=ce_1+de_2\)，有 \(g(Jv,w)=ad-bc\)。同一正向单位标架中，\(\nu=e^1\wedge e^2\)，故 \(\nu(v,w)=ad-bc\)，得到所需恒等式。这也说明 \(J\) 由度量与定向唯一确定：对每个 \(v\)，非退化性使满足 \(g(Jv,\cdot)=\nu(v,\cdot)\) 的向量唯一。

	对新度量，\((e^{-f}e_1,e^{-f}e_2)\) 为正向正交单位标架。在它上面规定的四分之一转动仍是原来的 \(J\)，因为两个基向量作了相同的正缩放。对应的对偶标架各乘 \(e^f\)，所以 \(\widetilde\nu=e^{2f}\nu\)。因此正向共形缩放不改变 \(J\)，但改变面积的尺度。
\end{solution}

\begin{exercise}
	\label{g1:ex:prescribed-positive-density}
	设 \(M\) 为 \(n\geq1\) 维 Hausdorff、第二可数光滑流形，可带边界，\(\mu\) 是一个光滑正密度。
	\begin{enumerate}
		\item 证明存在 Riemann 度量 \(g\) 的体积密度恰为 \(\mu\)，不要求 \(M\) 可定向。
		\item 若两个度量 \(g_0,g_1\) 的体积密度均为 \(\mu\)，构造连接它们且始终保持体积密度 \(\mu\) 的光滑度量族。说明普通凸组合为何未必保持体积密度。
	\end{enumerate}
\end{exercise}
\begin{solution}
	1. 先由\cref{g1:thm:bundle-metric-existence}给 \(\mathrm{T}M\) 取一个度量 \(h\)，记其正体积密度为 \(\mu_h\)。两个正密度的比 \(\rho=\mu/\mu_h\) 是定义在 \(M\) 上的正光滑函数：在任一局部标架中是两个正光滑系数的比，而二者换标架时都乘以同一绝对行列式，所以比值与标架无关。若令 \(g=\rho^{2/n}h\)，其局部矩阵行列式为 \(\rho^2\det H\)，故 \(\mu_g=\rho\mu_h=\mu\)。构造只使用密度的正性与比值，不需要选择定向。

	2. 先令 \(k_t=(1-t)g_0+tg_1\)、\(0\leq t\leq1\)。它对每个非零向量均取正值，故是一族光滑度量，但体积密度可能变化。再按第一问逐个归一化：
	\[
		g_t=\left(\frac{\mu}{\mu_{k_t}}\right)^{2/n}k_t.
	\]
	\(\mu_{k_t}\) 的系数是正行列式的平方根，联合光滑依赖于 \((t,p)\)；因此上式给出联合光滑度量族，且每个 \(g_t\) 的体积密度均为 \(\mu\)。端点处比值为一，故确实连接给定 \(g_0,g_1\)。

	普通凸组合不保体积的例子已可取在平凡二维丛中：\(G_0=\operatorname{diag}(2,1/2)\)、\(G_1=\operatorname{diag}(1/2,2)\) 的行列式都为一，但平均矩阵为 \((5/4)I\)，其体积密度系数为 \(5/4\)，不再是一。正定性对平均稳定，固定行列式却不是线性条件，因而需要额外归一化。
\end{solution}

% !TeX root = ../../Book1.tex
% 纲要标题：## 第13章　向量场、流与 Frobenius 可积性
% 纲要位置：几何学计划纲要.md，第 602 行。
% 正文状态：初稿；尚未正式审核。

% 纲要内容开始
% 
% > 本章目标与证明范围：由常微分方程和交换流进入分布的可积性，连续证明局部流、流盒与 Frobenius。叶的整体结构和叶空间不是可积性定理的自动结论。
% 
% 纲要内容结束

\chapter{向量场、流与 Frobenius 可积性}
\label{g1:ch:13}
\BookChapterNumber{g1:ch:13}

\begin{chapterabstract}
	向量场指定瞬时速度，局部流则把这些速度组织为随时间变化的局部微分同胚。本章从初值问题的存在唯一性出发，完整证明对初值和参数的光滑依赖，建立最大局部流、延拓判据、紧支完备性与流盒定理。Lie 括号记录不同方向的运动是否相容，张量场的 Lie 导数则描述沿流搬运后的变化。由交换化标架和交换流，我们连续证明常秩分布的 Frobenius 定理，并构造具有自身流形拓扑的最大连通叶。环面上的线性叶、返回映射、沿叶道路的横截面映射芽、非 Hausdorff 叶空间及接触分布，说明局部可积性能够保证什么，以及哪些整体问题仍需另外研究。
\end{chapterabstract}

\begin{learninggoals}
	\begin{itemize}
		\item 从积分方程构造初值解，用误差估计和变分方程证明对初值、时间与参数的联合光滑依赖。
		\item 明确最大流的定义域，使用逃逸判据判断完备性，并用横截片与逆函数定理构造流盒。
		\item 计算 Lie 括号与张量 Lie 导数，说明括号为零与局部流交换之间的等价关系，辨认自然性所需的映射条件。
		\item 将对合分布的标架交换化，利用共同小时间的流构造适应坐标，完整证明 Frobenius 定理及最大连通叶的存在唯一性。
		\item 区分叶片、最大叶、返回映射与沿叶道路的横截面映射芽，用具体例子判断叶是否嵌入、叶空间是否 Hausdorff，以及分布为什么不可积。
	\end{itemize}
\end{learninggoals}

上一章允许我们在各点光滑地指定一个切向量、一组切平面或一个张量。本章追问这些逐点数据怎样组织成运动和子流形。一个向量场给出速度，并不直接给出所有时间的运动；一个光滑平面分布在每点都很整齐，也未必能够拼成处处与之相切的曲面。存在、延拓与相容性是三个不同的问题，需要分别处理。

\ProseParagraphBreak

例如，平面上的 \(\partial_x\) 与 \(x\partial_y\) 各自都能产生所有时间的运动，但先水平移动会改变随后竖直移动的速度，交换次序后终点就不同。另一方面，环面上的常斜率方向在每个小邻域都能变成水平线，其完整轨道却可能是闭圆，也可能在整个环面中稠密。前一个例子解释为什么要引入 Lie 括号，后一个例子说明为什么局部流盒和局部积分叶还不能替代整体分析。

\ProseParagraphBreak

本章默认流形有限维、Hausdorff、第二可数，向量场与分布均光滑；主要定理的流形无边界，边界上的双侧流条件在\secref{g1:sec:1302}另行说明。分布均为固定秩子丛，不把维数跳跃的情形并入 Frobenius 定理。所用前置工具是\chapref{g1:ch:11}的切向量、微分、逆函数与隐函数定理，以及\chapref{g1:ch:12}的局部标架、张量拉回、插入和截断函数。常微分方程方面只预期读者熟悉积分和求导；解空间的完备性、误差估计及高阶光滑依赖会在本章实际用到的位置证明。

\ProseParagraphBreak

阅读时要留意两类定义域：流的时间区间随初值变化，叶上的邻近关系也可能不同于环境流形中的邻近关系。复合流需要共同的存在时间，构造最大叶则需要给出叶自身的拓扑。Lee 第二版附录 D 中的 Existence, Uniqueness, and Smoothness 一节可用于对照初值理论；第8章的 Lie Brackets 及第9章的前六节对应流、流盒和 Lie 导数；第19章的 Distributions and Involutivity、The Frobenius Theorem、Foliations 三节对应可积性，第22章的 Contact Structures 一节补充接触模型\cite[附录D, 第663--671页; 第8章, 第185--188页; 第9章, 第206--238页; 第19章, 第491--504页; 第22章, 第581--584页]{Lee2012Smooth}。本章使用微分形式的逐点代数，外微分与 Cartan 公式的证明接在\chapref{g1:ch:14}。

% !TeX root = ../../Book1.tex
% 纲要标题：### 13.1 常微分方程与局部流
% 纲要位置：几何学计划纲要.md，第 606 行。
% 正文状态：初稿；尚未正式审核。

% 纲要内容开始
% 
% * 局部存在唯一性及最大解区间
% * 初值与参数的光滑依赖
% * 流的局部群性质
% * 时间依赖向量场与演化映射
% 
% 纲要内容结束

\section{常微分方程与局部流}
\label{g1:sec:1301}

向量场在每个点指定一个瞬时速度。若让一个点始终按所在位置的速度运动，就得到积分曲线；若同时移动一片邻域中的所有点，就希望得到一族光滑变换。后一种说法比“每个初值各有一条解”要求更多：不同初值的解必须能在共同时间区间上比较，终点要光滑地依赖起点，而且先运动一段时间再继续运动，应与一次运动相应的总时间一致。这些要求将把微分方程的解组织成局部流。

\ProseParagraphBreak

本节主体中的流形均为有限维、Hausdorff、第二可数的无边界光滑流形，向量场均光滑。\cref{g1:def:smooth-tensor-field}已把向量场定义为切丛的光滑截面，因此在局部坐标中，其分量是光滑函数。我们先在 Euclid 空间重建所需的初值理论，再说明换图与延拓为何相容。边界点的双侧运动还需要向量场切于边界的条件，留到下一节单独处理。

\subsection{局部存在唯一性及最大解区间}
\label{g1:subsec:1301:01}

\begin{definition}
	\label{g1:def:integral-curve}
	设 \(M\) 为无边界光滑流形，\(X\in\mathfrak X(M)\)。若开区间 \(I\subset\mathbb R\) 上的光滑曲线 \(\gamma:I\to M\) 满足
	\[
		\gamma'(t)=X_{\gamma(t)}\qquad(t\in I),
	\]
	就称 \(\gamma\) 为 \(X\) 的\term[jifenquxian]{积分曲线}[integral curve]。给定 \(t_0\in I\) 和 \(p\in M\) 的条件 \(\gamma(t_0)=p\) 称为初值。常值曲线也包括在定义中。
\end{definition}

在坐标 \(x=(x^1,\ldots,x^n)\) 中写 \(X=\sum_i a^i\partial_{x^i}\)，积分曲线条件成为
\[
	\frac{\mathrm{d}x^i(\gamma(t))}{\mathrm{d}t}=a^i(\gamma(t)).
\]
链式法则说明这组方程与所用坐标无关。为同时处理参数与随后出现的时间依赖情形，先考虑右端 \(F(t,z,\lambda)\)，其中 \(z\in\mathbb R^n\) 是未知位置，\(\lambda\in\mathbb R^\ell\) 是固定参数。没有参数时取 \(\ell=0\)。

\begin{theorem}[局部初值问题]
	\label{g1:thm:ode-local-existence-uniqueness}
	设 \(\Omega\subset\mathbb R\times\mathbb R^n\times\mathbb R^\ell\) 开，\(F:\Omega\to\mathbb R^n\) 光滑，且 \((t_0,z_0,\lambda_0)\in\Omega\)。则存在 \(a>0\) 和 \(t_0,z_0,\lambda_0\) 的开邻域 \(T,B,\Lambda\)，使对所有 \((\tau,y,\lambda)\in T\times B\times\Lambda\)，初值问题
	\begin{equation}
		\label{g1:eq:parameter-ode-initial-value}
		\frac{\mathrm{d}z}{\mathrm{d}t}=F(t,z,\lambda),\qquad z(\tau)=y
	\end{equation}
	在 \(|t-\tau|<a\) 上有解。任意两条同参数、同初值的解，在包含初始时刻的公共区间上相同。每条解对时间光滑。
\end{theorem}
\begin{proof}
	平移符号后仍保留基点记号。由 \(\Omega\) 开，可取 \(\delta,r,\rho>0\)，使紧乘积
	\[
		K=[t_0-2\delta,t_0+2\delta]\times\overline B_{2r}(z_0)\times\overline B_\rho(\lambda_0)
	\]
	包含在 \(\Omega\) 中。若没有参数，省略最后因子。连续函数在紧集上有界，故存在 \(C,L\geq0\) 使 \(|F|\leq C\)、\(\|\mathrm{D}_zF\|\leq L\) 在 \(K\) 上成立。对固定时间及参数，沿大球中的线段积分导数，得到
	\[
		|F(t,z,\lambda)-F(t,w,\lambda)|\leq L|z-w|.
	\]
	取 \(T=(t_0-\delta,t_0+\delta)\)、\(B=B_r(z_0)\)，再取稍小的参数球 \(\Lambda\subset B_\rho(\lambda_0)\)。选择 \(0<a<\delta\)，使 \(aC<r\)、\(q=aL<1\)。即使 \(C\) 或 \(L\) 为零，也可以这样选择。

	改用相对时间 \(s=t-\tau\)。在所有取值于 \(\overline B_{2r}(z_0)\) 的连续路径 \(u:[-a,a]\to\mathbb R^n\) 上，定义
	\[
		(\mathcal T u)(s)=y+\int_0^s F(\tau+v,u(v),\lambda)\,\mathrm{d}v.
	\]
	积分公式把速度累积成位置；它的不动点正是待求解。由 \(|y-z_0|<r\) 与 \(aC<r\)，\(\mathcal Tu\) 仍取值于大闭球。用一致距离 \(\|u-w\|_\infty=\sup_{|s|\leq a}|u(s)-w(s)|\)，有
	\[
		\|\mathcal Tu-\mathcal Tw\|_\infty\leq q\|u-w\|_\infty.
	\]
	这个估计在正、负时间上都成立，因为反向积分的长度仍不超过 \(a\)。

	\cref{g1:lem:euclidean-contraction}曾在有限维闭球上使用过相邻迭代差的几何级数。此处对象是连续路径，因此补明所需完备性：一致 Cauchy 的路径列在每个时间点都有 Euclid 极限；一致 Cauchy 条件使逐点极限也是一致极限，而连续函数的一致极限连续，闭球的闭性又保证极限仍在闭球内。于是从常路径 \(u_0(s)=y\) 出发，令 \(u_{m+1}=\mathcal Tu_m\)，有
	\[
		\|u_{m+1}-u_m\|_\infty\leq q^m\|u_1-u_0\|_\infty.
	\]
	几何级数说明迭代一致收敛到某条连续路径 \(u\)。\(\mathcal T\) 对一致距离连续，所以 \(u=\mathcal Tu\)。微积分基本定理使 \(u\) 为 \(C^1\)，且满足方程与初值。由 \(u'=F(\tau+s,u,\lambda)\) 逐阶求导，又得 \(u\) 光滑。

	若两条路径都是上述积分方程的不动点，压缩估计给出 \(\|u-w\|_\infty\leq q\|u-w\|_\infty\)，因而相同。对任意另一个通过同一初值的解，连续性保证它在足够短的公共区间内也留在所取球中，所以得到局部唯一性。再对两解公共区间中的任意相遇时刻作同样处理，它们相等的时间集合是开集；由连续性又是闭集，且含初始时刻。公共区间连通，故两解在整个公共区间相等。
\end{proof}

共同时间 \(a\) 是这一结论的重要部分：它适用于附近一整族初值与参数。它来自对 \(F\) 和 \(\mathrm{D}_zF\) 的共同界，而不是先对每个初值选一个时间后再假定这些时间有正的下界。常数 \(C\) 控制路径是否会离开允许的坐标范围，\(L\) 则控制不同路径之间的差距。

\ProseParagraphBreak

将这个定理用于向量场的局部分量，便得到流形上的局部积分曲线。两张图中得到的曲线在重叠时满足同一个坐标方程，所以唯一性使它们一致。仍须分清局部解与最大解：某张图不够大时可以换图继续，但能否一直继续到所有实数时间，是另一个问题。

\begin{theorem}[最大积分曲线]
	\label{g1:thm:maximal-integral-curve}
	设 \(M\) 为有限维 Hausdorff、第二可数的无边界光滑流形，\(X\in\mathfrak X(M)\)。对每个 \(p\in M\)，存在唯一开区间
	\[
		I_p=(\alpha(p),\beta(p)),\qquad -\infty\leq\alpha(p)<0<\beta(p)\leq\infty,
	\]
	及积分曲线 \(\gamma_p:I_p\to M\)，满足 \(\gamma_p(0)=p\)，并使任何具有该初值的积分曲线都是它的限制。称其为通过 \(p\) 的\term[zuidajifenquxian]{最大积分曲线}[maximal integral curve]。若 \(X_p=0\)，则 \(I_p=\mathbb R\)，且 \(\gamma_p\) 为常值曲线。
\end{theorem}
\begin{proof}
	局部初值定理先给出至少一条过 \(p\) 的积分曲线。若两条这样的曲线在某时刻相遇，选相遇点的坐标图，由连续性可将两条曲线都限制到该图中；Euclid 唯一性表明它们在该时刻附近相同。因此两曲线相等的时间集合在其公共区间中开。\(M\) 的 Hausdorff 性保证对角集合闭，所以两曲线相等的时间集合也闭。它包含零，公共区间连通，故两曲线在全部公共区间相同。

	取所有过 \(p\)、初始时刻为零的积分曲线定义区间之并，记为 \(I_p\)。它是含零的开区间，因为所有被并区间都含零。对 \(t\in I_p\)，任选一条在 \(t\) 有定义的解，令 \(\gamma_p(t)\) 等于它在该处的值。前段唯一性说明此定义与选择无关；在每个局部区间它就是一条已有解，故整体光滑且满足方程。它包含所有其他同初值解，也就不能再作严格延拓。若还有另一条具有相同最大性质的解，二者互为限制，所以定义域与曲线都相同。

	当 \(X_p=0\) 时，常值曲线 \(\gamma(t)=p\) 在整个 \(\mathbb R\) 上满足方程。最大性质与唯一性给出最后一项。零维流形上的向量场处处为零，也由这一情形覆盖。
\end{proof}

同样的并集证明适用于固定参数的非自治 Euclid 方程：给定初始时刻 \(\tau\)，所有解区间的并是包含 \(\tau\) 的最大开区间，图像始终要求留在方程的定义域中。“最大”指不能再延拓，并不表示区间一定无界，也不表示不同初值具有相同区间。

\begin{example}
	\label{g1:exm:quadratic-maximal-interval}
	在 \(\mathbb R\) 上取 \(X=x^2\partial_x\)。初值 \(x(0)=p\) 的解为
	\[
		x(t)=\frac{p}{1-tp},\qquad
		I_p=\begin{cases}
			(-\infty,1/p),&p>0,\\
			\mathbb R,&p=0,\\
			(1/p,\infty),&p<0.
		\end{cases}
	\]
	直接求导验证方程，唯一性说明它给出初值解。有限端点处函数绝对值趋于无穷，故不可能以实数值连续延拓。虽然分式在极点另一侧也有意义，那一侧与初始时刻不在同一个定义区间，不能一并作为这条初值解。作为另一种有限时间终止，开区间 \((0,1)\) 上的常向量场 \(\partial_x\) 的解为 \(p+t\)，最大区间是 \((-p,1-p)\)；其坐标没有趋向无穷，却在有限时间离开了底流形。
\end{example}

\subsection{初值与参数的光滑依赖}
\label{g1:subsec:1301:02}

在\cref{g1:lem:smooth-ode-flow}中，测地线初值问题已用过自主方程的光滑依赖。这里需要同时处理初始时刻、初值、参数及其高阶导数。不能从“每条轨道对时间光滑”直接推出这些结论；两者分别比较同一条轨道上的位置与不同轨道之间的变化。

\begin{lemma}[积分不等式]
	\label{g1:lem:gronwall-estimate}
	设 \(u:[0,T]\to[0,\infty)\) 连续，\(a,b,L\geq0\)，且
	\[
		u(t)\leq a+bt+L\int_0^t u(s)\,\mathrm{d}s.
	\]
	则当 \(L>0\) 时，
	\[
		u(t)\leq a e^{Lt}+\frac bL(e^{Lt}-1);
	\]
	当 \(L=0\) 时，\(u(t)\leq a+bt\)。
\end{lemma}
\begin{proof}
	令 \(v(t)=a+bt+L\int_0^t u(s)\,\mathrm{d}s\)，则 \(u\leq v\)、\(v(0)=a\)，而 \(v'=b+Lu\leq b+Lv\)。因此
	\[
		\frac{\mathrm{d}}{\mathrm{d}t}(e^{-Lt}v(t))\leq b e^{-Lt}.
	\]
	从零积分到 \(t\)，并乘回 \(e^{Lt}\)，即得 \(L>0\) 时的结论。\(L=0\) 时原假设本身就是所需估计。
\end{proof}

这一估计常称为\term[gronwallguji]{Gronwall 估计}[Gronwall estimate]。\footnote{Thomas Hakon Grönwall（1877—1932），瑞典数学家，研究分析不等式及数学在物理化学中的应用。}它把“误差由原有误差加累计误差控制”变成明确的指数界；例如 \(a=b=0\) 时必有 \(u=0\)。负时间的估计可先置 \(v=-t\)，再对反向时间区间应用同一引理。

\begin{theorem}[初值与参数的光滑依赖]
	\label{g1:thm:ode-smooth-dependence}
	设 \(\Omega\subset\mathbb R\times\mathbb R^n\times\mathbb R^\ell\) 开，\(F:\Omega\to\mathbb R^n\) 光滑，\((t_0,z_0,\lambda_0)\in\Omega\)。存在 \(a>0\) 及 \(t_0,z_0,\lambda_0\) 的开邻域 \(T,B,\Lambda\)，使初值问题 \(z'=F(t,z,\lambda)\)、\(z(\tau)=y\) 对所有 \((\tau,y,\lambda)\in T\times B\times\Lambda\) 都在 \(|t-\tau|<a\) 上有唯一解，其存在唯一性见\cref{g1:thm:ode-local-existence-uniqueness}，且相对时间表达
	\[
		u(s;\tau,y,\lambda)=z(\tau+s;\tau,y,\lambda)
	\]
	在 \((-a,a)\times T\times B\times\Lambda\) 上对全部变量联合光滑。等价地，绝对时间表达 \(z(t;\tau,y,\lambda)\) 在 \(|t-\tau|<a\) 的相应开集上光滑。若一条解已存在于某个紧时间区间，则初始数据充分接近时也能在该区间及其一个开邻域上求解，并对时间、初始时刻、初值和参数联合光滑。
\end{theorem}
\begin{proof}
	先证明自主 Euclid 方程 \(z'=G(z)\) 的初值依赖，再把时间与参数并入状态。以下 \(G\) 为 Euclid 开集上的任意光滑映射，局部解记为 \(\psi(s,y)\)。在前一定理使用的较小初值球与共同闭时间区间上，记 \(q=a\sup\|\mathrm{D}G\|<1\)。两个积分方程相减，得到
	\[
		\|\psi(\cdot,y+h)-\psi(\cdot,y)\|_\infty
		\leq |h|+q\|\psi(\cdot,y+h)-\psi(\cdot,y)\|_\infty,
	\]
	故其差不超过 \(|h|/(1-q)\)。同一局部范围中速度有界，又给出对时间的共同连续性。因此 \(\psi\) 已经对 \((s,y)\) 联合连续。

	为了找出初值导数，先独立解矩阵积分方程
	\[
		A(s,y)=I+\int_0^s \mathrm{D}G(\psi(v,y))A(v,y)\,\mathrm{d}v.
	\]
	在连续矩阵函数的一致范数中，右边关于 \(A\) 仍是压缩因子不超过 \(q\) 的仿射映射。连续矩阵函数空间的完备性与连续路径情形相同，故相邻迭代的几何级数给出唯一解，并且
	\[
		\|A(\cdot,y)\|_\infty\leq\frac1{1-q}.
	\]
	比较 \(y,y'\) 所对应的两个方程，得到
	\[
		(1-q)\|A(\cdot,y')-A(\cdot,y)\|_\infty
		\leq a\|\mathrm{D}G(\psi(\cdot,y'))-\mathrm{D}G(\psi(\cdot,y))\|_\infty\,
		\|A(\cdot,y)\|_\infty.
	\]
	右端随 \(y'\to y\) 趋零，所以 \(A\) 对初值连续；其积分方程又给出对时间的联合连续性。

	现在证明 \(A\) 确实是导数，而不只是一条形式变分方程的解。记 \(\delta_h=\psi(\cdot,y+h)-\psi(\cdot,y)\)。\(\mathrm{D}G\) 在所用紧凸球上一致连续，故存在 \(\omega(\varepsilon)\to0\)，使沿线段积分得到的一阶余项满足
	\[
		|G(z+w)-G(z)-\mathrm{D}G(z)w|\leq\omega(|w|)|w|
	\]
	对该球中的 \(z,z+w\) 一致成立。已知 \(\|\delta_h\|_\infty\leq C_0|h|\)，其中 \(C_0=(1-q)^{-1}\)。将两解之差的积分方程减去 \(A(s,y)h\) 的方程，便有
	\[
		\|\delta_h-A(\cdot,y)h\|_\infty
		\leq q\|\delta_h-A(\cdot,y)h\|_\infty
		+a\omega(C_0|h|)C_0|h|.
	\]
	移项并除以 \(|h|\)，右端趋零。因此 \(\mathrm{D}_y\psi=A\)，而且这个导数连续。时间导数 \(\partial_s\psi=G(\psi)\) 同样联合连续。把一个联合增量拆成时间增量与初值增量，分别使用一元积分公式与上述导数估计，即得 \(\psi\) 对 \((s,y)\) 联合为 \(C^1\)。

	在继续求高阶导数以前，说明如何把任一已得的局部正则性沿一段已有解传播。设一条解的定义域含紧区间 \([c,d]\)。其像紧，可用有限个初值邻域 \(V_i\) 覆盖，每个 \(V_i\) 都有共同短时间 \((-a_i,a_i)\) 上的局部解映射。取时间步长小于 \(\min_i a_i/2\)，在给定区间内选有限分点。每到一个分点，在包含该解位置的某个 \(V_i\) 中使用对应局部映射到达下一个分点；唯一性保证所得轨道与原解相同。对附近初值，前一步的终点随初值连续，故可缩小初值邻域，使其仍落在下一步所需的 \(V_i\) 中。有限步只需有限次缩小。于是各步终点和每一步内的任意时间位置都由有限次局部解映射复合得到。若局部映射为 \(C^k\)，这些复合也为 \(C^k\)；在分点附近，各公式存在重叠时间区间并由唯一性一致，所以联合给出光滑至 \(k\) 阶的映射。初始时刻在区间内部时，向前、向后分别作同样构造；端点还可先略微延长，再覆盖所需紧区间。

	下面对所有光滑自主 Euclid 系统同时归纳正则阶数。已证任意这样的系统有 \(C^1\) 初值依赖。假设已知任意光滑系统的解映射为 \(C^k\)，其中 \(k\geq1\)。考虑有限维增广系统
	\[
		z'=G(z),\qquad A'=\mathrm{D}G(z)A,
		\qquad (z(0),A(0))=(y,I).
	\]
	它的右端是 \((z,A)\) 的光滑函数，因此归纳假设适用。由唯一性，其解恰为已经构造的 \((\psi(s,y),\mathrm{D}_y\psi(s,y))\)。在原共同时间区间中的任一紧子区间，这条增广解已存在：\(z\) 是原解，\(A\) 是上面矩阵积分方程的解。利用前段的有限步传播，增广解在该区间上对初值为 \(C^k\)；初值映射 \(y\mapsto(y,I)\) 光滑，故 \(\mathrm{D}_y\psi\) 对 \((s,y)\) 为 \(C^k\)。另一方面，归纳假设已使 \(\psi\) 为 \(C^k\)，所以 \(\partial_s\psi=G(\psi)\) 也为 \(C^k\)。全部一阶偏导都为 \(C^k\)，便得到 \(\psi\in C^{k+1}\)。归纳对每个有限阶成立，因而 \(\psi\) 为 \(C^\infty\)。这里每一阶的局部计算都传播到原来同一个解区间，不是将无穷多个逐渐缩短的时间区间取交。

	例如上述归纳的第二阶给出，对初值方向 \(h,k\)，令 \(B=\mathrm{D}_y^2\psi[h,k]\)，则
	\[
		B'=\mathrm{D}G(\psi)B+\mathrm{D}^2G(\psi)[Ah,Ak],\qquad B(0)=0.
	\]
	最高阶未知量仍线性出现，而新的非齐次项由已经控制的低阶变化构成。增广系统的归纳将这个机制用于所有阶，且每一步真正使用了已经证明的初值依赖。

	最后回到 \(F(t,z,\lambda)\)。在开集 \(\Omega\subset\mathbb R^{1+n+\ell}\) 上定义自主系统
	\[
		(a,z,\lambda)'=\bigl(1,F(a,z,\lambda),0\bigr).
	\]
	从 \((\tau,y,\lambda)\) 出发，其第一分量是 \(\tau+s\)，最后分量保持 \(\lambda\)，中间分量由唯一性正是 \(u(s;\tau,y,\lambda)\)。自主系统的光滑初值依赖便给出所需全部参数的联合光滑性；再作光滑换元 \(s=t-\tau\)，得到绝对时间表达。沿已有紧时间区间的结论，也由前述有限步传播应用于这个增广系统得到。
\end{proof}

把已证明的光滑性代入积分方程，现在可以合法地对初值与参数求导。对相对时间解 \(u=u(s;\tau,y,\lambda)\)，记 \(P=\mathrm{D}_yu\)、\(Q=\mathrm{D}_\lambda u\)，则
\begin{equation}
	\label{g1:eq:ode-variational-equation}
	\begin{aligned}
		P'&=\mathrm{D}_zF(\tau+s,u,\lambda)P,&\qquad P(0)&=I_n,\\
		Q'&=\mathrm{D}_zF(\tau+s,u,\lambda)Q+\mathrm{D}_\lambda F(\tau+s,u,\lambda),& Q(0)&=0.
	\end{aligned}
\end{equation}
第一行测量初值扰动怎样传播，第二行还加入方程本身改变所产生的源项。它们称为相应的\term[bianfenfangcheng]{变分方程}[variational equation]。若 \(\|\mathrm{D}_zF\|\leq L\)，\cref{g1:lem:gronwall-estimate}给出 \(\|P(s)\|\leq e^{L|s|}\)；若另有 \(\|\mathrm{D}_\lambda F\|\leq C\)，则 \(L>0\) 时 \(\|Q(s)\|\leq C(e^{L|s|}-1)/L\)，\(L=0\) 时为 \(C|s|\)。这些界适用于解仍处于所取共同范围内的时间。

\ProseParagraphBreak

光滑依赖也没有承诺最大存在时间是光滑函数。\cref{g1:exm:quadratic-maximal-interval}中，初值为正时未来端点是 \(1/p\)，初值为零或负时未来端点为 \(+\infty\)。真正适于微分的对象是所有“时间与初值”的允许组合所成的开集及其上的解映射，下面就把它作为流的定义域。
\subsection{流的局部群性质}
\label{g1:subsec:1301:03}

将各个初值的最大积分曲线放在一起，定义
\[
	\mathcal D_X=\{(t,p)\in\mathbb R\times M:t\in I_p\},\qquad
	\Phi(t,p)=\gamma_p(t),\qquad \Phi_t(p)=\Phi(t,p).
\]
集合 \(\mathcal D_X\) 的每个竖直截面都是包含零的开区间，但这还不能单凭集合形式推出 \(\mathcal D_X\) 在乘积空间中开。一个初值能运动到某个时刻时，附近初值也能运动到相近时刻，正是还需用光滑依赖证明的部分。

\begin{theorem}[最大局部流]
	\label{g1:thm:maximal-local-flow}
	设 \(M\) 为有限维 Hausdorff、第二可数的无边界光滑流形，\(X\in\mathfrak X(M)\)。对每个 \(p\in M\)，令 \(\gamma_p:I_p\to M\) 为满足 \(\gamma_p(0)=p\) 的最大积分曲线，定义 \(\mathcal D_X=\{(t,p):p\in M,\ t\in I_p\}\) 及 \(\Phi(t,p)=\gamma_p(t)\)，并记 \(\Phi_t(p)=\Phi(t,p)\)。集合 \(\mathcal D_X\) 是包含 \(\{0\}\times M\) 的开集，映射 \(\Phi:\mathcal D_X\to M\) 光滑，且
	\[
		\Phi_0(p)=p,\qquad \partial_t\Phi_t(p)=X_{\Phi_t(p)}.
	\]
	若 \(s\in I_p\)，则
	\begin{equation}
		\label{g1:eq:maximal-flow-time-shift}
		I_{\Phi_s(p)}=I_p-s,
		\qquad \Phi_t(\Phi_s(p))=\Phi_{t+s}(p)\quad(t+s\in I_p).
	\end{equation}
	对于固定 \(t\)，集合 \(D_t=\{p:(t,p)\in\mathcal D_X\}\) 开，且
	\[
		\Phi_t:D_t\longrightarrow D_{-t}
	\]
	为微分同胚，逆为 \(\Phi_{-t}\)。这个由最大积分曲线组成的映射称为 \(X\) 的\term[zuidajubuliu]{最大局部流}[maximal local flow]，并且它唯一由 \(X\) 决定。
\end{theorem}
\begin{proof}
	先证开性与光滑性。Euclid 光滑依赖定理应用于坐标中的向量场分量，给出每点附近的开集 \(V\subset M\)、正数 \(a\) 和光滑局部解映射 \(\psi:(-a,a)\times V\to M\)，其中 \(\psi(0,q)=q\)。坐标唯一性及最大积分曲线定理说明，这个局部映射是 \(\Phi\) 的限制。

	固定 \((t_0,p_0)\in\mathcal D_X\)。在 \(I_{p_0}\) 内取一个紧区间 \([c,d]\)，使零与 \(t_0\) 都属于其内部。轨道段 \(\gamma_{p_0}([c,d])\) 紧，可被有限个上述 \(V_i\) 覆盖，相应共同时间记为 \(a_i>0\)。把 \([0,d]\) 与 \([c,0]\) 分别划成步长小于 \(\min_i a_i/2\) 的有限时间段。在每个分点，选择包含该轨道点的某个 \(V_i\)，并用该处局部映射走到下一分点。由唯一性，固定轨道的这些局部解总与 \(\gamma_{p_0}\) 相同。

	现在让起点 \(p\) 在 \(p_0\) 附近变化。第一步终点是 \(p\) 的光滑函数；因第二步所需的 \(V_i\) 开，可缩小起点邻域，使第一步终点仍落入其中。依次进行，有限步得到一个共同开邻域 \(U\ni p_0\)，使所有这些局部解都能按相同分点复合。时间段内部的位置也有公式 \(\psi_i(t-t_j,p_j(p))\)，其中 \(p_j(p)\) 是前面有限步得到的光滑终点。步长严格小于各自允许时间，故这些表达在相邻分点附近都有时间重叠；重叠时是同初值的积分曲线，唯一性使其一致。于是它们粘成 \((c,d)\times U\) 上的光滑解映射。

	这条拼成的曲线包含初始时刻零，所以它是每个 \(p\in U\) 的最大积分曲线的限制。因此 \((c,d)\times U\subset\mathcal D_X\)，且 \(\Phi\) 在这个开集上光滑。它包含原先任取的 \((t_0,p_0)\)，从而证明整个定义域开、整个解映射光滑。定义域包含零时刻，初值与速度等式则直接来自各条积分曲线。

	再证时间平移关系。固定 \(s\in I_p\)，令 \(q=\Phi_s(p)\)。曲线 \(t\mapsto\gamma_p(s+t)\) 在 \(I_p-s\) 上是过 \(q\) 的积分曲线，故最大性质给出
	\[
		I_p-s\subset I_q,\qquad \gamma_q(t)=\gamma_p(s+t)\quad(t\in I_p-s).
	\]
	特别地，\(-s\in I_q\) 且 \(\gamma_q(-s)=p\)。因此 \(u\mapsto\gamma_q(u-s)\) 在 \(s+I_q\) 上是过 \(p\)、初始时刻为零的积分曲线，另一方向的最大性质给出 \(s+I_q\subset I_p\)。两次包含说明 \(I_q=I_p-s\)，唯一性同时给出所述复合公式。这里使用了方程自主：平移时间后，速度仍由同一个 \(X\) 决定。

	由于 \(\mathcal D_X\) 开，固定时间切片 \(D_t\) 开。如果 \(p\in D_t\)，时间平移关系给出 \(q=\Phi_t(p)\in D_{-t}\) 与 \(\Phi_{-t}(q)=p\)。反过来从 \(D_{-t}\) 出发也得到 \(\Phi_t\Phi_{-t}=\operatorname{id}\)。两个映射都是光滑 \(\Phi\) 的固定时间限制，所以互为微分同胚。最后，任何满足同一速度方程与初值的局部运动，在每条轨道上都必须限制自 \(\gamma_p\)，故最大局部流唯一。
\end{proof}

公式\eqref{g1:eq:maximal-flow-time-shift}说明时间加法如何作用于运动，但它并不保证每个 \(\Phi_t\) 都定义在整个 \(M\) 上。对 \(X=x^2\partial_x\)，最大定义域与流为
\[
	\mathcal D_X=\{(t,p):1-tp>0\},\qquad
	\Phi_t(p)=\frac p{1-tp}.
\]
例如 \(t>0\) 时，\(D_t=(-\infty,1/t)\)，而像为 \(D_{-t}=(-1/t,\infty)\)。逆映射 \(q\mapsto q/(1+tq)\) 正好在后一个开区间上定义。把分式的所有非极点都列入定义域，会错误地加入不属于初值轨道的另一段时间。

\ProseParagraphBreak

局部群公式还说明流保持产生它的向量场。固定 \((t,p)\in\mathcal D_X\)，对足够小的 \(s\) 有 \(\Phi_t(\Phi_s(p))=\Phi_{t+s}(p)\)。在 \(s=0\) 求导，得到
\begin{equation}
	\label{g1:eq:flow-preserves-generator}
	\mathrm{d}(\Phi_t)_p(X_p)=X_{\Phi_t(p)}.
\end{equation}
左边先取起点的速度再用微分送到终点，右边直接取终点的速度，两者一致。这个等式在后面比较两个不同向量场的流时会再次出现。

\begin{example}
	\label{g1:exm:rotation-local-flow}
	平面向量场 \(X=-y\partial_x+x\partial_y\) 的流为
	\[
		\Phi_t\binom{x}{y}
		=\begin{pmatrix}\cos t&-\sin t\\\sin t&\cos t\end{pmatrix}\binom{x}{y}.
	\]
	求导验证速度方程，\(t=0\) 时为恒等，因此唯一性说明这就是其最大流；公式对全部实数时间有定义。原点是静止轨道，其余轨道为圆。两条轨道若在同一时刻相遇，局部唯一性与逆映射会迫使它们有相同起点。圆周轨道却可以在不同时刻回到同一点，所以轨道参数化未必单射。
\end{example}

\subsection{时间依赖向量场与演化映射}
\label{g1:subsec:1301:04}

若速度规律会随绝对时间改变，只知道经过了多久通常不够，还必须知道从哪个时刻开始。设 \(J\subset\mathbb R\) 是开区间。一个光滑时间依赖向量场是光滑映射 \(X:J\times M\to \mathrm{T}M\)，满足 \(X(t,p)\in \mathrm{T}_pM\)，也写作 \(X_t(p)\)。它的初值问题是
\[
	\gamma'(t)=X_t(\gamma(t)),\qquad \gamma(s)=p.
\]
这里 \(s\) 是起始时刻，\(t\) 是到达时刻。

\ProseParagraphBreak

把时间本身当成新的位置坐标，便能把这个问题转为自主方程。在无边界流形 \(J\times M\) 上定义
\[
	\widehat X_{(t,p)}=\partial_t+X_t(p).
\]
第一分量速度恒为一，因此从 \((s,p)\) 出发、经历时间 \(h\) 后，第一分量必为 \(s+h\)。剩下的分量正好解原来的时间依赖方程。这一构造保留了绝对时间，而不是把它误当成固定参数。

\begin{proposition}[演化映射]
	\label{g1:prop:time-dependent-evolution}
	设 \(M\) 为有限维 Hausdorff、第二可数的无边界光滑流形，\(J\subset\mathbb R\) 为开区间，\(X:J\times M\to \mathrm{T}M\) 为光滑时间依赖向量场。对每个 \((s,p)\in J\times M\)，存在包含 \(s\) 的唯一最大开时间区间 \(I_{s,p}\subset J\) 及初值解。写其在 \(t\) 时刻的值为 \(\Phi_{t,s}(p)\)，则
	\[
		\mathcal E_X=\{(t,s,p)\in J\times J\times M:t\in I_{s,p}\}
	\]
	为开集，\((t,s,p)\mapsto\Phi_{t,s}(p)\) 光滑，且
	\begin{equation}
		\label{g1:eq:evolution-initial-value}
		\Phi_{s,s}(p)=p,\qquad
		\partial_t\Phi_{t,s}(p)=X_t(\Phi_{t,s}(p)).
	\end{equation}
	若 \(r\in I_{s,p}\)，则
	\begin{equation}
		\label{g1:eq:evolution-composition}
		I_{r,\Phi_{r,s}(p)}=I_{s,p},\qquad
		\Phi_{t,r}(\Phi_{r,s}(p))=\Phi_{t,s}(p)\quad(t\in I_{s,p}).
	\end{equation}
	固定 \(t,s\) 后，\(\Phi_{t,s}\) 在其开定义域上为微分同胚，逆为 \(\Phi_{s,t}\)。此外，在 \(\mathcal E_X\) 上有
	\begin{equation}
		\label{g1:eq:evolution-initial-time-derivative}
		\partial_s\Phi_{t,s}(p)
		=-\mathrm{d}(\Phi_{t,s})_p(X_s(p)).
	\end{equation}
\end{proposition}
\begin{proof}
	对 \(\widehat X\) 使用\cref{g1:thm:maximal-local-flow}，记其最大流为 \(\widehat\Phi_h\)，定义域为 \(\mathcal D_{\widehat X}\)。第一分量满足 \(a'=1\)、\(a(0)=s\)，所以
	\[
		\widehat\Phi_h(s,p)=\bigl(s+h,\Phi_{s+h,s}(p)\bigr).
	\]
	反过来，任何时间依赖方程的解 \(t\mapsto\gamma(t)\)，都给出增广系统的解 \(h\mapsto(s+h,\gamma(s+h))\)。因此增广系统的存在、唯一及最大性，恰好给出原方程的存在、唯一及最大区间；它只能在 \(s+h\in J\) 时存在。

	定义域可以写成
	\[
		\mathcal E_X=\{(t,s,p):(t-s,(s,p))\in\mathcal D_{\widehat X}\}.
	\]
	右边是开集在连续映射 \((t,s,p)\mapsto(t-s,(s,p))\) 下的原像，所以开。\(\Phi_{t,s}(p)\) 是将 \(\widehat\Phi_{t-s}(s,p)\) 投到第二分量得到的光滑映射，故联合光滑；初值和到达时间导数也由增广方程直接得到。

	若先从绝对时间 \(s\) 走到 \(r\)，增广系统经历的时间是 \(r-s\)；再走到 \(t\)，经历的是 \(t-r\)。对 \(\widehat\Phi\) 应用最大时间平移关系，得到同一轨道的最大绝对时间区间不变，并且
	\[
		\widehat\Phi_{t-r}\bigl(\widehat\Phi_{r-s}(s,p)\bigr)
		=\widehat\Phi_{t-s}(s,p).
	\]
	取第二分量就得到\eqref{g1:eq:evolution-composition}。令 \(t=s\) 给出从 \(s\) 到 \(r\) 再返回的恒等映射；交换两个时刻得到另一侧的逆。因此 \(\Phi_{t,s}\) 与 \(\Phi_{s,t}\) 在各自开定义域上互逆，且均光滑。

	最后固定允许的 \((t,s,p)\)。开性保证对足够小的 \(h\)，复合公式
	\[
		\Phi_{t,s+h}\bigl(\Phi_{s+h,s}(p)\bigr)=\Phi_{t,s}(p)
	\]
	有定义。对 \(h\) 在零处求导，链式法则给出
	\[
		\partial_s\Phi_{t,s}(p)+\mathrm{d}(\Phi_{t,s})_p(X_s(p))=0,
	\]
	即所需初始时刻导数公式。
\end{proof}

\begin{example}
	\label{g1:exm:time-dependent-translation}
	在实直线上取 \(X_t=t\partial_x\)。积分得到
	\[
		\Phi_{t,s}(x)=x+\frac{t^2-s^2}{2}.
	\]
	它依赖起始与到达两个时刻，并满足 \(\Phi_{t,r}\Phi_{r,s}=\Phi_{t,s}\)。若只保留 \(\Psi_t=\Phi_{t,0}\)，则
	\[
		\Psi_t(\Psi_s(x))=x+\frac{t^2+s^2}{2},\qquad
		\Psi_{t+s}(x)=x+\frac{(t+s)^2}{2},
	\]
	一般不同。这不是唯一性失效，而是第二次使用 \(\Psi_t\) 又把绝对时间从零开始计算了。对自主向量场，速度不随绝对时间变化，才有 \(\Phi_{t,s}=\Phi_{t-s}\) 这样的单参数表达。
\end{example}

至此，局部求解、对初值求导和移动时间起点都具有明确的定义域。下一节将判断何时这些区间能延伸到整个实轴，并利用局部流把非零向量场附近的坐标整理成沿一条坐标轴的平移。
% 纲要条目排为编号小节标题，并进入目录与书签；说明段仍保留于纲要注释中。
% !TeX root = ../../Book1.tex
% 纲要标题：### 13.2 完备性和流盒
% 纲要位置：几何学计划纲要.md，第 613 行。
% 正文状态：初稿；尚未正式审核。

% 纲要内容开始
% 
% * 延拓判据与逃逸现象
% * 紧支向量场的完备性
% * 流盒定理的证明
% * 由截断构造具有指定局部行为的整体流
% 
% 纲要内容结束

\section{完备性和流盒}
\label{g1:sec:1302}

局部存在定理给每个初始点一小段运动，却没有承诺可以一直走下去。即使向量场处处光滑，速度也可能越来越大，使轨道在有限时间内逃向无穷远；反过来，只要轨道不断返回一个固定紧集，局部解就能重新启动。本节先把这种延拓机制说清，再用流建立适配坐标，最后用上一章的截断函数把指定的局部运动装进一个整体流中。除边界说明外，流形均无边界，并具有有限维、Hausdorff 和第二可数性质。

\subsection{延拓判据与逃逸现象}
\label{g1:subsec:1302:01}

设 \(X\) 为光滑向量场，过 \(p\) 的最大积分曲线记为 \(\gamma_p:(a_p,b_p)\to M\)，其中 \(0\in(a_p,b_p)\)。\cref{g1:thm:maximal-integral-curve}中的“最大”是按定义区间包含关系说的：任何具有同一初值的解都是它的限制。若 \(b_p<\infty\)，端点不属于定义域；需要判断的是能否在接近这个时刻时找到另一个可用的初始点。

\begin{lemma}[紧集上的统一存在时间]
	\label{g1:lem:uniform-flow-time-compact}
	设 \(X\) 是无边界光滑流形 \(M\) 上的光滑向量场，\(K\subset M\) 紧。存在 \(\varepsilon>0\) 和包含 \(K\) 的开集 \(V\)，使最大流 \(\Phi\) 在 \((-\varepsilon,\varepsilon)\times V\) 上有定义。
\end{lemma}
\begin{proof}
	由\cref{g1:thm:maximal-local-flow}，流的定义域是包含 \(\{0\}\times M\) 的开集。每个 \(q\in K\) 因而有乘积邻域 \((-\varepsilon_q,\varepsilon_q)\times V_q\) 留在这个定义域内。用有限个 \(V_{q_1},\ldots,V_{q_N}\) 覆盖 \(K\)，取 \(\varepsilon=\min_j\varepsilon_{q_j}>0\) 及 \(V=\bigcup_jV_{q_j}\)，即得结论。空紧集时任取一个正数即可。
\end{proof}

\begin{theorem}[有限端点的逃逸判据]
	\label{g1:thm:flow-escape-criterion}
	设 \(X\) 是无边界光滑流形 \(M\) 上的光滑向量场，\(\gamma:(a,b)\to M\) 是最大积分曲线。如果 \(b<\infty\)，则对每个紧集 \(K\subset M\)，存在 \(t_K<b\)，使 \(\gamma(t)\notin K\) 对所有 \(t_K<t<b\) 成立。有限左端点具有相应结论。因此，若有 \(t_j\uparrow b<\infty\) 使 \(\gamma(t_j)\) 都落在一个固定紧集内，则原解一定能够越过 \(b\) 延拓。
\end{theorem}
\begin{proof}
	若结论失败，就有趋向 \(b\) 的时刻 \(t_j\)，使 \(\gamma(t_j)\in K\)。取上一引理给出的统一存在时间 \(\varepsilon\)，并取 \(j\) 使 \(b-t_j<\varepsilon\)。从 \(\gamma(t_j)\) 出发的曲线
	\[
		t\longmapsto\Phi_{t-t_j}(\gamma(t_j)),\qquad |t-t_j|<\varepsilon,
	\]
	与原曲线在 \(t_j\) 取同一值。局部唯一性及沿公共区间逐点延续，说明两者在交叠区间上相同，所以可以粘合。新曲线的定义域包含 \((a,t_j+\varepsilon)\)，而 \(t_j+\varepsilon>b\)，与最大性矛盾。把时间反向，等价于把 \(X\) 换为 \(-X\)，便得左端点结论。
\end{proof}

这个判据不要求先证明 \(\gamma(t)\) 在端点处有极限。紧性在这里提供的是统一的重新启动时间；只要足够晚的某个初值仍在紧集里，已经足以越过旧端点。也不能把“离开一个坐标域”误认为“最大解结束”：坐标域可能很小，曲线离开后仍能在另一张图中继续。

\begin{example}[有限时间逃逸]
	\label{g1:exm:finite-time-flow-escape}
	在 \(\mathbb R\) 上取 \(X=x^2\partial_x\)。对初值 \(x(0)=x_0\neq0\)，分离变量得到
	\[
		x(t)=\frac{x_0}{1-tx_0}.
	\]
	若 \(x_0>0\)，最大区间为 \((-\infty,1/x_0)\)，并且 \(t\uparrow1/x_0\) 时 \(x(t)\to+\infty\)。若 \(x_0<0\)，最大区间为 \((1/x_0,+\infty)\)，有限逃逸发生在反向时间；\(x_0=0\) 则给出整体常值解。向量场处处光滑，失效的是轨道在有限时间内留在任何固定紧集中的可能性。

	另取开流形 \(M=(0,1)\) 和 \(X=\partial_x\)，解为 \(x(t)=x_0+t\)，最大区间是 \((-x_0,1-x_0)\)。坐标数值始终有界，但在 \((0,1)\) 中，每个紧集都与两个缺失端点保持正距离。因此这仍是逃离流形内所有紧集的情形，不是逃逸判据的反例。
\end{example}

\subsection{紧支向量场的完备性}
\label{g1:subsec:1302:02}

\begin{definition}
	\label{g1:def:complete-vector-field}
	无边界光滑流形上的光滑向量场 \(X\) 称为\term[wanbeixiangliangchang]{完备向量场}[complete vector field]，如果它的每条最大积分曲线都定义在整个 \(\mathbb R\) 上。此时最大流称为\term[zhengtiliu]{整体流}[global flow]，是光滑映射 \(\Phi:\mathbb R\times M\to M\)。
\end{definition}

由\cref{g1:thm:maximal-local-flow}，整体流满足 \(\Phi_0=\operatorname{id}\)、\(\Phi_{s+t}=\Phi_s\circ\Phi_t\)，且 \(\Phi_t^{-1}=\Phi_{-t}\)。因此每个固定时间给出整个流形的微分同胚。完备性是向量场的性质，此处并没有给 \(M\) 选定距离，更没有假设度量完备。

\begin{theorem}[紧支向量场的完备性]
	\label{g1:thm:compact-support-complete}
	设 \(M\) 为无边界光滑流形，\(X\) 是光滑向量场。如果
	\[
		\operatorname{supp}X=\overline{\{p:X_p\neq0\}}
	\]
	为紧集，则 \(X\) 完备。特别地，紧无边界光滑流形上的每个光滑向量场都完备。
\end{theorem}
\begin{proof}
	记 \(K=\operatorname{supp}X\)。若一条积分曲线在某个时刻经过零点 \(q\)，常曲线 \(t\mapsto q\) 是具有同一初值的整体解，由唯一性，原最大解就是这条常曲线。因此非恒定积分曲线不经过任何零点，尤其始终留在 \(K\) 中。

	若某条非恒定最大解有有限右端点，它在接近端点时仍留在紧集 \(K\)，违反\cref{g1:thm:flow-escape-criterion}。左端点同理。常值解本来就定义在整个实轴上，所以全部解都整体存在。若 \(M\) 紧，\(\operatorname{supp}X\) 是 \(M\) 的闭子集，因而紧，得到最后一项。
\end{proof}

紧支是充分条件，并非必要条件。例如 \(\partial_x\) 在 \(\mathbb R\) 上的流为 \(x\mapsto x+t\)，\(x\partial_x\) 的流为 \(x\mapsto e^t x\)，两者都完备而不紧支。第二个例子的速度还无界；真正要排除的是有限时间内无法延拓，而不是任意大的速度。

\begin{proposition}[边界上的切向条件]
	\label{g1:prop:boundary-tangent-flow}
	设 \(M\) 为带边界光滑流形，\(X\) 为光滑向量场，并满足 \(X_p\in \mathrm{T}_p\partial M\) 对每个 \(p\in\partial M\) 成立。则 \(X\) 有保持边界和内点集合的双侧局部流；若 \(X\) 紧支，则该流整体存在。反之，若一个双侧局部流的每个时间映射为流形的局部微分同胚，它的生成向量场必切于边界。
\end{proposition}
\begin{proof}
	在边界图中记点为 \((y,r)\)，其中 \(r\geq0\)，写 \(X=(a(y,r),b(y,r))\)。光滑性允许将分量局部延拓到欧氏开集；切向条件为 \(b(y,0)=0\)。由微积分基本定理，缩小为乘积邻域后有
	\[
		b(y,r)=r\,c(y,r),\qquad c(y,r)=\int_0^1\partial_r b(y,\theta r)\,\mathrm{d}\theta,
	\]
	其中 \(c\) 光滑。欧氏延拓方程的解在此邻域内满足
	\[
		r(t)=r(0)\exp\!\left(\int_0^t c(y(s),r(s))\,\mathrm{d}s\right).
	\]
	这可以由线性标量方程 \(r'=c(y(t),r(t))r\) 的唯一性验证。因此 \(r(0)=0\) 时一直在边界，\(r(0)>0\) 时保持为正，正负小时间都不会越出半空间。欧氏局部流遂限制为所需局部流，并且不同延拓的限制由唯一性相同。紧集上的有限覆盖和逃逸判据的证明完全适用，所以紧支时整体存在。

	反向，流形的局部微分同胚保持边界；故对边界点 \(p\)，曲线 \(t\mapsto\Phi_t(p)\) 位于 \(\partial M\)，其在零时刻的速度属于 \(\mathrm{T}_p\partial M\)。这正是生成向量场的值。
\end{proof}

例如闭区间 \([0,1]\) 上的 \(\partial_x\) 虽然光滑且紧支，却不切于边界，也不可能产生双侧整体流。这说明上一完备性定理中的无边界条件不能直接删掉。后面的流盒和 Frobenius 定理仍在无边界流形上陈述；边界适配坐标需要另外核对切向或横截条件。

\subsection{流盒定理的证明}
\label{g1:subsec:1302:03}

在 \(X_p\neq0\) 的点附近，积分曲线能否作为一族坐标线？只沿一条曲线计时还不够：我们还要用 \(n-1\) 个参数区分邻近曲线。办法是在 \(p\) 处放一张与 \(X_p\) 横截的小截面，再从截面上每一点出发沿流运动。时间与截面参数合起来，正好给出 \(n\) 个候选坐标。下面给出\term[liuhe]{流盒定理}[flow-box theorem]。

\begin{theorem}[流盒定理]
	\label{g1:thm:flow-box}
	设 \(M\) 为 \(n\geq1\) 维无边界光滑流形，\(X\) 是光滑向量场，\(X_p\neq0\)。则存在包含 \(p\) 的坐标邻域 \(U\)，其坐标为 \((s,z^2,\ldots,z^n)\)，取值于乘积开集 \((-\varepsilon,\varepsilon)\times B\)，使
	\[
		X=\frac{\partial}{\partial s}\quad\text{于 }U.
	\]
	在这些坐标中，局部流为 \((s,z)\mapsto(s+t,z)\)，只要相应线段留在这个乘积内。
\end{theorem}
\begin{proof}
	先选 \(p\) 附近的坐标并作一个可逆线性变换，使 \(p\) 对应原点，\(X_p\) 在该点为第一坐标方向。令 \(\sigma(z)\) 为局部超平面 \(x^1=0\) 的参数化，\(\sigma(0)=p\)。由流的定义域开放性，在足够小的乘积邻域上可定义
	\[
		F(s,z)=\Phi_s(\sigma(z)).
	\]
	\cref{g1:thm:maximal-local-flow}保证 \(F\) 光滑。对时间求导得到 \(\partial_sF(0,0)=X_p\)，对 \(z\) 求导则得到 \(\mathrm{d}\sigma_0\) 的 \(n-1\) 个切向量。这些向量与 \(X_p\) 合起来组成 \(\mathrm{T}_pM\) 的基，因此 \(\mathrm{d}F_{(0,0)}\) 可逆。

	由\cref{g1:thm:inverse-function}，\(F\) 在 \((0,0)\) 的某邻域上是微分同胚；在该邻域内再取一个小乘积 \((-\varepsilon,\varepsilon)\times B\)，其像 \(U\) 仍为开集。将 \(F^{-1}\) 作为新坐标。因为 \(s\mapsto F(s,z)\) 本来就是积分曲线，处处有
	\[
		\mathrm{d}F_{(s,z)}(\partial_s)=X_{F(s,z)}.
	\]
	所以新坐标中的向量场恰为 \(\partial_s\)，不是仅在原点处相等。固定 \(z\) 的直线满足新坐标下的微分方程，唯一性给出所述平移公式。\(n=1\) 时截面只有一个点，同一论证仍成立。
\end{proof}

\input{Figures/Chapter13-FlowBox.tex}

流盒坐标把时间保留下来，因此比只把轨道画成直线更强：若只改变轨道的形状而不校正速度，得到的可能只是 \(a(s,z)\partial_s\)。另一方面，这个定理完全是局部的。一条整体轨道可以离开流盒再回来，在同一个流盒中留下几段彼此分开的轨迹，不能把整条最大轨道都认作一根局部坐标线。

\begin{example}[径向场的流盒]
	\label{g1:exm:radial-flow-box}
	在半平面 \(x>0\) 上取 \(X=x\partial_x+y\partial_y\)，其流为 \(\Phi_t(x,y)=(e^t x,e^t y)\)。选截面 \(\sigma(z)=(1,z)\)，得到
	\[
		F(s,z)=(e^s,e^s z),\qquad s=\log x,\qquad z=\frac yx.
	\]
	直接计算 \(X(s)=1\)、\(X(z)=0\)，所以 \(X=\partial_s\)。量 \(z\) 标记不同射线，量 \(s\) 沿每条射线计时。若将同一个场看作整个平面上的场，原点处 \(X=0\)，那里不可能有这种流盒：微分同胚的微分把零向量仍送成零，不能送成非零的坐标向量。
\end{example}

\subsection{由截断构造具有指定局部行为的整体流}
\label{g1:subsec:1302:04}

局部向量场不一定完备，但许多几何构造只需它在一小块区域内保持原来的行为。上一章的截断函数正好用于限制支集；本节的完备性定理则把这个支集条件转成全时间存在。应当同时记住：在一块区域内保留了向量场，只能保证轨道留在该区域期间的运动不变。

\begin{proposition}[局部向量场的完备延拓]
	\label{g1:prop:complete-local-extension}
	设 \(M\) 为无边界光滑流形，\(U\subset M\) 开，\(K\subset U\) 紧，\(X\) 是 \(U\) 上的光滑向量场。存在 \(M\) 上的紧支完备光滑向量场 \(Y\)，在 \(K\) 的某开邻域 \(W\subset U\) 与 \(X\) 相同，且 \(\operatorname{supp}Y\subset U\)。\(Y\) 的整体流在 \(M\setminus\operatorname{supp}Y\) 上恒为恒等。对初值位于 \(W\) 的 \(X\) 积分曲线，其包含初始时刻且留在 \(W\) 的区间内，\(Y\) 的积分曲线与它相同。
\end{proposition}
\begin{proof}
	由\cref{g1:cor:smooth-cutoff-closed}，选 \(\chi:M\to[0,1]\)，在 \(K\) 的某开邻域 \(W\) 为一，并且 \(\operatorname{supp}\chi\) 紧、包含于 \(U\)。定义
	\[
		Y_p=
		\begin{cases}
			\chi(p)X_p,&p\in U,\\
			0_p,&p\notin U.
		\end{cases}
	\]
	\cref{g1:thm:section-extension}的零延拓论证说明 \(Y\) 光滑：\(U\) 外每一点都有避开 \(\operatorname{supp}\chi\) 的邻域，在那里整个表达恒为零。它的支集包含于紧集 \(\operatorname{supp}\chi\)，故由\cref{g1:thm:compact-support-complete}完备。

	支集外的点是零点，故由唯一性得到恒定轨道。在 \(W\) 中，\(X=Y\)，所以任何始终留在 \(W\) 内的 \(X\) 解也满足 \(Y\) 的方程；同一初值的唯一性使两条解在所述区间相同。
\end{proof}

这一结论还能同时控制邻近初值。若一条 \(X\) 轨道的紧时间段 \(\{\Phi_t(p):|t|\leq T\}\) 包含于 \(W\)，则存在 \(p\) 的邻域 \(V\)，使所有 \(q\in V\) 的 \(X\) 轨道在 \([-T,T]\) 上有定义且留在 \(W\)。理由是 \([-T,T]\times\{p\}\) 为流定义域中映到 \(W\) 的紧集，利用流的连续性与有限个乘积邻域，就可找出共同的 \(V\)。因此两个流在整个 \([-T,T]\times V\) 上相同，而非仅沿最初选定的那一条曲线相同。

\ProseParagraphBreak

例如对上一例的径向场，给定一个不含原点的小闭圆盘及包含它的开集，可将该场截断成整个平面上的紧支场。新流在小圆盘附近仍具有同样的径向运动，在远处则固定各点。反过来，对 \(x^2\partial_x\) 作紧支截断后得到整体流，原来的有限时间发散不会保留：轨道进入截断改变向量场的区域时，微分方程已经不同。

\begin{proposition}[共同紧支的时间依赖场]
	\label{g1:prop:compact-time-dependent-evolution}
	设 \(M\) 为无边界光滑流形，\(J\subset\mathbb R\) 为开区间，\(X_t\) 关于 \((t,p)\) 光滑。如果存在固定紧集 \(K\subset M\)，使 \(\operatorname{supp}X_t\subset K\) 对每个 \(t\in J\) 成立，则演化映射 \(\Phi_{t,s}\) 对任意 \(s,t\in J\) 及所有初值都有定义，并为整体微分同胚，逆为 \(\Phi_{s,t}\)。
\end{proposition}
\begin{proof}
	固定一个初始时刻 \(s\)。若初值在 \(K\) 外，该点对每个 \(X_t\) 都是零点，故其唯一解恒定。由同样的唯一性，初值在 \(K\) 中的解不能到达 \(K\) 外的任何一点；否则从那一时刻向前后比较常值解，会使初值也在 \(K\) 外。

	若解在 \(J\) 内有有限最大端点 \(b\)，选一个包含 \(s,b\) 的紧时间区间 \([c,d]\subset J\)。对增维场 \(\partial_t+X_t\)，该轨道在接近端点时留在紧集 \([c,d]\times K\) 中；\cref{g1:thm:flow-escape-criterion}排除了这样的端点。左端点同理，所以原解在整个 \(J\) 上存在。光滑性、复合律与逆映射公式由\cref{g1:prop:time-dependent-evolution}得到。
\end{proof}

共同紧支的假设控制的是整段时间内的轨道活动范围，不能只从“每个固定时刻的向量场分别紧支”就省去对支集如何随时间变化的检查。

% 纲要条目排为编号小节标题，并进入目录与书签；说明段仍保留于纲要注释中。

% !TeX root = ../../Book1.tex
% 纲要标题：### 13.3 Lie 括号与无穷小变换
% 纲要位置：几何学计划纲要.md，第 620 行。
% 正文状态：初稿；尚未正式审核。

% 纲要内容开始
% 
% * 导子交换子与坐标表达
% * 流的交换和 Lie 括号
% * 张量场的 Lie 导数
% * 自然性及分次 Cartan 运算的准备
% 
% 纲要内容结束

\section{Lie 括号与无穷小变换}
\label{g1:sec:1303}

沿两个方向先后运动，交换次序后是否到达同一点？如果两个方向是固定的平面平移，答案是肯定的；如果第二个方向的速度随位置变化，第一步运动就可能改变第二步的效果。这个差别不能只看一点处的两个向量，还要看附近的向量场如何变化。本节把它写成向量场的 Lie 括号，并用流解释这一运算的几何含义。

\ProseParagraphBreak

本节的流形均为有限维、Hausdorff、第二可数、无边界的光滑流形，向量场和张量场均光滑。局部向量场可在其定义开集上应用同样的结果。所需的分析结论是\cref{g1:thm:maximal-local-flow}给出的光滑局部流、流的复合规律及唯一性；不同流的复合总在共同有定义的区域内讨论，不假定向量场完备。

\subsection{导子交换子与坐标表达}
\label{g1:subsec:1303:01}

向量场 \(X\) 作用于一个光滑函数 \(f\)，得到函数 \(Xf\)，其在 \(p\) 的值为 \(X_p(f)\)。由\cref{g1:prop:derivation-coordinate-basis}，在坐标中
\[
	X=\sum_iX^i\partial_i,\qquad Xf=\sum_iX^i\partial_i f,
\]
所以 \(Xf\) 仍光滑，并满足 \(X(fg)=X(f)g+fX(g)\)。这些运算也适用于局部函数：某点附近的函数相等，所得导数在该点附近也相等。于是可以把 \(X\) 看作作用于光滑函数的一个一阶微分算子。

\ProseParagraphBreak

复合算子 \(XY\) 指先作 \(Y\)，再作 \(X\)。它通常包含二阶导数，因而未必还是一个向量场；但是在 \(XY-YX\) 中，二阶部分恰好抵消。

\begin{definition}
	\label{g1:def:lie-bracket}
	设 \(M\) 为无边界光滑流形，\(X,Y\in\mathfrak X(M)\)。定义它们的\term[likuohao]{Lie 括号}[Lie bracket]在局部光滑函数上的作用为
	\[
		[X,Y]f=X(Yf)-Y(Xf).
	\]
	右端是两个算子的交换子。下述计算说明，它确实定义一个光滑向量场，而不只是一个一般微分算子。
\end{definition}

为检查乘法法则，先展开
\[
	XY(fg)=XY(f)\,g+Y(f)X(g)+X(f)Y(g)+f\,XY(g).
\]
交换 \(X,Y\) 后，中间两个交叉项仍是同一对项，相减后得到
\[
	[X,Y](fg)=[X,Y](f)\,g+f[X,Y](g).
\]
因此在每个点，\([X,Y]\) 是函数芽上的导子。按照\cref{g1:prop:derivation-coordinate-basis}，它对应一个切向量，其坐标分量是它作用于坐标函数的值。由 \(Y(x^j)=Y^j\)、\(X(x^j)=X^j\)，得到
\begin{equation}
	\label{g1:eq:lie-bracket-coordinates}
	[X,Y]=\sum_j\left(\sum_iX^i\partial_iY^j-Y^i\partial_iX^j\right)\partial_j.
\end{equation}
系数光滑，故它是光滑向量场。这个坐标表达由无坐标的导子定义推出，所以换坐标时仍表示同一个向量场。

\begin{proposition}[Lie 括号的基本性质]
	\label{g1:prop:lie-bracket-properties}
	设 \(M\) 为无边界光滑流形，\(X,Y,Z\in\mathfrak X(M)\)，\(f,g\in C^\infty(M)\)。Lie 括号对实数双线性，并满足
	\[
		[X,Y]=-[Y,X],\qquad
		[X,[Y,Z]]+[Y,[Z,X]]+[Z,[X,Y]]=0.
	\]
	后一个等式称为 Jacobi 恒等式。此外，
	\begin{equation}
		\label{g1:eq:lie-bracket-product}
		[X,fY]=f[X,Y]+X(f)Y,\qquad
		[fX,gY]=fg[X,Y]+fX(g)Y-gY(f)X.
	\end{equation}
	所有这些等式也适用于同一开集上的局部向量场与函数。
\end{proposition}
\begin{proof}
	实双线性和反对称性直接来自算子复合的分配律。为证明 Jacobi 恒等式，把复合算子作用于任意局部光滑函数，利用复合的结合性展开：
	\[
		\begin{aligned}
			[X,[Y,Z]]&=XYZ-XZY-YZX+ZYX,\\
			[Y,[Z,X]]&=YZX-YXZ-ZXY+XZY,\\
			[Z,[X,Y]]&=ZXY-ZYX-XYZ+YXZ.
		\end{aligned}
	\]
	相加后每个三重复合恰好出现一次正号、一次负号，故总和为零。向量场由它对局部函数的作用确定，所以这是向量场的恒等式。

	对任意局部函数 \(h\)，有
	\[
		[X,fY]h=X(fYh)-fY(Xh)
		=X(f)Yh+f[X,Y]h.
	\]
	这给出第一条乘法公式。再由反对称性得到 \([fX,Y]=f[X,Y]-Y(f)X\)，将第一式用于 \([fX,gY]\)，即得第二式。
\end{proof}

因此 \(\mathfrak X(M)\) 是一个实 Lie 代数，即一个实向量空间，其上配有实双线性括号，且该括号满足反对称性与 Jacobi 恒等式。但括号一般不是对 \(C^\infty(M)\) 双线性的。\eqref{g1:eq:lie-bracket-product}中的函数导数正是区别。比如在平面上，\(X=\partial_x\)、\(Y=x\partial_y\) 满足 \([X,Y]=\partial_y\)。虽然 \(Y\) 在原点为零，括号在原点却不为零；将 \(Y\) 换成恒零向量场，它在原点的值没有改变，括号却变成零。这说明 Lie 括号不是只依赖 \(X_p,Y_p\) 的逐点双线性张量。

\ProseParagraphBreak

坐标场满足 \([\partial_i,\partial_j]=0\)，但任意局部标架没有这一性质。例如 \(e_1=\partial_x\)、\(e_2=e^x\partial_y\) 处处独立，却有 \([e_1,e_2]=e_2\)。已有局部标架并不保证能找到彼此交换的标架，后者还需要一项微分条件；下一节讨论分布时将用到这一区别。

\subsection{流的交换和 Lie 括号}
\label{g1:subsec:1303:02}

令 \(\Phi=\Phi^X\) 为 \(X\) 的最大局部流。回顾\cref{g1:thm:maximal-local-flow}：它定义在含 \(\{0\}\times M\) 的开集 \(\mathcal D_X\subset\mathbb R\times M\) 上，\(\Phi_0=\operatorname{id}\)，并在各项有定义时满足 \(\Phi_{t+s}=\Phi_t\circ\Phi_s\)。固定 \(t\) 后，\(\Phi_t\) 是其开定义域与像之间的微分同胚，逆为 \(\Phi_{-t}\)。这些性质允许我们把不同点的向量带到同一个切空间中比较。

\ProseParagraphBreak

直接求 \(Y_{\Phi_t(p)}-Y_p\) 没有内在意义，因为两个向量的基点不同。应先用逆流的微分将第一个向量带回 \(\mathrm{T}_pM\)，定义
\begin{equation}
	\label{g1:eq:flow-vector-pullback}
	\bigl((\Phi_t)^*Y\bigr)_p
	=\mathrm{d}(\Phi_{-t})_{\Phi_t(p)}Y_{\Phi_t(p)}
	=\bigl(\mathrm{d}(\Phi_t)_p\bigr)^{-1}Y_{\Phi_t(p)}.
\end{equation}
这是\cref{g1:prop:mixed-tensor-pullback}的向量场拉回。对函数则是 \((\Phi_t)^*f=f\circ\Phi_t\)。拉回先移动观察点，再将该点的几何量搬回原点，所以其变化率同时记录向量场和搬运方式的变化。

\begin{proposition}[沿流的拉回导数]
	\label{g1:prop:flow-pullback-derivative}
	设 \(M\) 为无边界光滑流形，\(X,Y\in\mathfrak X(M)\)，\(\Phi\) 为 \(X\) 的局部流。则在共同有定义的区域内，
	\[
		\left.\frac{\mathrm{d}}{\mathrm{d}t}\right|_{t=0}(\Phi_t)^*Y=[X,Y],
		\qquad
		\frac{\mathrm{d}}{\mathrm{d}t}(\Phi_t)^*Y
		=(\Phi_t)^*[X,Y]=[X,(\Phi_t)^*Y].
	\]
	对任意光滑函数 \(f\)，还有
	\[
		\frac{\mathrm{d}}{\mathrm{d}t}(\Phi_t)^*f=(\Phi_t)^*(Xf).
	\]
	特别地，\([X,Y]=0\) 当且仅当 \(Y\) 被 \(X\) 的局部流保持，即 \((\Phi_t)^*Y=Y\)；等价地，
	\[
		\mathrm{d}(\Phi_t)_pY_p=Y_{\Phi_t(p)}
	\]
	在流的定义域内成立。
\end{proposition}
\begin{proof}
	先在一点附近选坐标，把 \(X,Y\) 的分量视为列向量，流的坐标表达记为 \(\phi_t(x)\)。由\cref{g1:thm:ode-smooth-dependence}的光滑依赖及流方程，
	\[
		\phi_0(x)=x,\qquad
		\left.\partial_t\phi_t(x)\right|_0=X(x),\qquad
		\left.\partial_t\mathrm{D}\phi_t(x)\right|_0=\mathrm{D}X(x).
	\]
	记 \(A(t)=\mathrm{D}\phi_t(x)\)。对 \(A(t)^{-1}A(t)=I\) 求导，并用 \(A(0)=I\)，得到 \((A^{-1})'(0)=-\mathrm{D}X(x)\)。因此\eqref{g1:eq:flow-vector-pullback}的导数为
	\[
		\left.\frac{\mathrm{d}}{\mathrm{d}t}\right|_0
		\bigl(\mathrm{D}\phi_t(x)^{-1}Y(\phi_t(x))\bigr)
		=-\mathrm{D}X(x)Y(x)+\mathrm{D}Y(x)X(x).
	\]
	第 \(j\) 个分量正是\eqref{g1:eq:lie-bracket-coordinates}的系数，证明了零时刻公式。

	固定一个有定义的 \((t,p)\)。由流定义域的开放性，在 \(p\) 的适当邻域及 \(h=0\) 附近，可同时使用
	\[
		(\Phi_{t+h})^*=(\Phi_t)^*(\Phi_h)^*
		=(\Phi_h)^*(\Phi_t)^*.
	\]
	第一种分解对 \(h\) 求导，得到 \((\Phi_t)^*[X,Y]\)；第二种分解把 \((\Phi_t)^*Y\) 视为局部向量场，使用已经证明的零时刻公式，得到 \([X,(\Phi_t)^*Y]\)。函数公式则直接由沿积分曲线的链式法则得到。

	若 \([X,Y]=0\)，固定 \(p\) 后，上述拉回在包含零的时间定义区间内导数为零。它始终位于固定的向量空间 \(\mathrm{T}_pM\)，按任意基逐分量应用一元微积分，便知它恒等于零时刻的 \(Y_p\)。反之，拉回保持不变时在零时刻求导，得到括号为零。最后一个等式由乘以可逆的 \(\mathrm{d}(\Phi_t)_p\) 得到。
\end{proof}

公式中的负号来自逆微分的求导，不能只沿轨道对 \(Y\) 的分量求导而略去它。另一个直接的检验是 \(X\) 自己被其流保持：对 \(\Phi_t(\Phi_s(p))=\Phi_{t+s}(p)\) 在 \(s=0\) 求导，得到 \(\mathrm{d}(\Phi_t)_pX_p=X_{\Phi_t(p)}\)，与 \([X,X]=0\) 一致。

\begin{theorem}[交换的局部流]
	\label{g1:thm:commuting-local-flows}
	设 \(M\) 为无边界光滑流形，\(X,Y\in\mathfrak X(M)\)，局部流分别为 \(\Phi,\Psi\)。以下两项等价：
	\begin{enumerate}
		\item \([X,Y]=0\)。
		\item 对每个 \(p\in M\)，存在开邻域 \(U\) 和 \(\varepsilon>0\)，使对所有 \(q\in U\) 及 \(|s|,|t|<\varepsilon\)，两种复合均有定义且
		\[
			\Phi_t(\Psi_s(q))=\Psi_s(\Phi_t(q)).
		\]
	\end{enumerate}
	这里不要求两向量场完备，结论是共同定义域内的局部交换。
\end{theorem}
\begin{proof}
	假设括号为零。两个复合映射的定义条件各自在 \(\mathbb R^2\times M\) 中构成开集，并且都包含 \((0,0,p)\)：这是因为各流定义域开，且中间点随参数连续变化。取这两个开集交中的一个乘积邻域 \((-\varepsilon,\varepsilon)^2\times U\)，便得到陈述中统一的共同定义域。

	固定 \(q\in U\) 及 \(|t|<\varepsilon\)。由于 \(X\) 的流保持 \(Y\)，曲线 \(c(s)=\Phi_t(\Psi_s(q))\) 满足
	\[
		c'(s)=\mathrm{d}(\Phi_t)_{\Psi_s(q)}Y_{\Psi_s(q)}
		=Y_{c(s)},\qquad c(0)=\Phi_t(q).
	\]
	曲线 \(s\mapsto\Psi_s(\Phi_t(q))\) 也满足同一个初值问题。积分曲线唯一性使它们在 \((-\varepsilon,\varepsilon)\) 上相同，故两种流交换。

	反过来，假设第二项成立。将交换式对 \(s\) 在零处求导，得到
	\[
		\mathrm{d}(\Phi_t)_qY_q=Y_{\Phi_t(q)}.
	\]
	所以 \((\Phi_t)^*Y=Y\) 在相应邻域和小时间上成立。再对 \(t\) 在零处求导，由\cref{g1:prop:flow-pullback-derivative}得到 \([X,Y]_q=0\)。每个点都有这样的邻域，故括号处处为零。
\end{proof}

\begin{example}[运动次序的差别]
	\label{g1:exm:noncommuting-plane-flows}
	在 \(\mathbb R^2\) 上取 \(X=\partial_x\)、\(Y=x\partial_y\)。它们的流在所有时间有定义：
	\[
		\Phi_t(x,y)=(x+t,y),\qquad
		\Psi_s(x,y)=(x,y+sx).
	\]
	两种次序分别得到
	\[
		\Phi_t\Psi_s(x,y)=(x+t,y+sx),\qquad
		\Psi_s\Phi_t(x,y)=(x+t,y+s(x+t)).
	\]
	第二个终点比第一个多出竖直位移 \(st\)，如\cref{g1:fig:flow-order}所示。原因是先沿 \(X\) 移动会增大随后沿 \(Y\) 运动的速度；而先沿 \(Y\) 运动时，这个增大尚未发生。括号 \([X,Y]=\partial_y\) 正好指出这个差别的方向。沿 \(X\) 的拉回也能直接看到它：\((\Phi_t)^*Y=(x+t)\partial_y\)，对 \(t\) 求导得到 \(\partial_y\)。
\end{example}

\input{Figures/Chapter13-FlowOrder.tex}

\subsection{张量场的 Lie 导数}
\label{g1:subsec:1303:03}

流不仅可以搬运向量，也可以搬运度量、形式以及其他张量。\cref{g1:prop:mixed-tensor-pullback}已经说明，微分同胚拉回一个 \((r,s)\) 型张量时，对每个向量因子使用逆微分，对每个对偶因子使用微分的对偶。因此对 \(X\) 的流 \(\Phi\)，
\[
	\bigl((\Phi_t)^*T\bigr)_p
	=\left(\bigl(\mathrm{d}(\Phi_t)_p\bigr)^{-1}\right)^{\otimes r}
	\otimes\left(\bigl(\mathrm{d}(\Phi_t)_p\bigr)^*\right)^{\otimes s}
	\bigl(T_{\Phi_t(p)}\bigr).
\]
这个表达始终位于 \(T^r_s(\mathrm{T}_pM)\)，可以在固定纤维中对时间求导。

\begin{definition}
	\label{g1:def:lie-derivative-tensor}
	设 \(M\) 为无边界光滑流形，\(X\in\mathfrak X(M)\)，\(T\) 为光滑 \((r,s)\) 型张量场。定义 \(T\) 沿 \(X\) 的\term[lidaoshu]{Lie 导数}[Lie derivative]为
	\[
		\mathcal L_XT=\left.\frac{\mathrm{d}}{\mathrm{d}t}\right|_{t=0}(\Phi_t^X)^*T.
	\]
	导数在各点固定的张量纤维中取。局部流对时间和初值光滑，张量拉回的局部公式也光滑，因此这一定义给出同型的光滑张量场。
\end{definition}

在函数上有 \(\mathcal L_Xf=Xf\)，在向量场上有 \(\mathcal L_XY=[X,Y]\)。所以 Lie 导数扩展了同一个沿流比较的想法，无需另选内积或联络；它比较的方式由 \(X\) 的流本身确定。

\begin{proposition}[张量 Lie 导数的计算]
	\label{g1:prop:lie-derivative-tensor-rules}
	设 \(M\) 为无边界光滑流形，\(X\in\mathfrak X(M)\)，\(T,S\) 为光滑张量场，\(f\) 为光滑函数。Lie 导数实线性，并满足
	\[
		\mathcal L_X(T\otimes S)=(\mathcal L_XT)\otimes S+T\otimes(\mathcal L_XS),\qquad
		\mathcal L_X(fT)=X(f)T+f\mathcal L_XT.
	\]
	对任意指定位置的收缩 \(C\)，有 \(\mathcal L_X(CT)=C(\mathcal L_XT)\)。在局部坐标中，
	\begin{equation}
		\label{g1:eq:lie-derivative-tensor-coordinates}
		\begin{split}
			(\mathcal L_XT)^{i_1\cdots i_r}{}_{j_1\cdots j_s}
			={}&\sum_kX^k\partial_kT^{i_1\cdots i_r}{}_{j_1\cdots j_s}\\
			&-\sum_{a=1}^r\sum_k(\partial_kX^{i_a})
			T^{i_1\cdots i_{a-1}k i_{a+1}\cdots i_r}{}_{j_1\cdots j_s}\\
			&+\sum_{b=1}^s\sum_k(\partial_{j_b}X^k)
			T^{i_1\cdots i_r}{}_{j_1\cdots j_{b-1}k j_{b+1}\cdots j_s}.
		\end{split}
	\end{equation}
	此外，在流的共同定义域内，
	\[
		\frac{\mathrm{d}}{\mathrm{d}t}(\Phi_t^X)^*T
		=(\Phi_t^X)^*(\mathcal L_XT)
		=\mathcal L_X\bigl((\Phi_t^X)^*T\bigr).
	\]
	因此 \(\mathcal L_XT=0\) 当且仅当 \(T\) 被 \(X\) 的局部流保持。
\end{proposition}
\begin{proof}
	拉回保持张量积与收缩，见\cref{g1:prop:mixed-tensor-pullback}。对这些恒等式在 \(t=0\) 求导，用通常的乘积求导法则，便得到前两条乘法公式及收缩相容性。它也保持因子的排列，因此对称张量和交替张量的 Lie 导数仍保持各自的对称性。

	为得到坐标公式，先计算坐标基及对偶基的导数。向量场情形已知
	\[
		\mathcal L_X\partial_i=[X,\partial_i]=-\sum_k(\partial_iX^k)\partial_k.
	\]
	对任意一形式 \(\alpha\) 与向量场 \(Y\)，拉回和求值相容：
	\[
		(\Phi_t)^*(\alpha(Y))=((\Phi_t)^*\alpha)((\Phi_t)^*Y).
	\]
	在零时刻求导，得到
	\begin{equation}
		\label{g1:eq:lie-one-form-evaluation}
		(\mathcal L_X\alpha)(Y)=X(\alpha(Y))-\alpha([X,Y]).
	\end{equation}
	令 \(\alpha=\mathrm{d}x^j\)、\(Y=\partial_i\)，由于 \(\mathrm{d}x^j(\partial_i)=\delta_i^j\) 为常数，便有
	\[
		\mathcal L_X(\mathrm{d}x^j)=\sum_i(\partial_iX^j)\,\mathrm{d}x^i.
	\]
	将 \(T\) 在坐标张量基中展开，分别对系数、每个向量基因子及每个对偶基因子求导，就得到\eqref{g1:eq:lie-derivative-tensor-coordinates}。上指标的负号来自向量基的导数，下指标的正号来自对偶基的导数。

	最后，与\cref{g1:prop:flow-pullback-derivative}的证明一样，在一个共同定义的邻域中分别用
	\[
		(\Phi_{t+h})^*=(\Phi_t)^*(\Phi_h)^*
		=(\Phi_h)^*(\Phi_t)^*
	\]
	对 \(h\) 在零处求导，得到一般时间的两个等式。若 \(\mathcal L_XT=0\)，固定 \(p\) 后，\(\bigl((\Phi_t)^*T\bigr)_p\) 在固定有限维张量空间中导数为零，故等于 \(T_p\)；反向只需在 \(t=0\) 求导。
\end{proof}

同样对一般求值恒等式求导，可以得到一个无需坐标的计算方法。若 \(T\) 为协变 \(k\) 阶张量，则
\begin{equation}
	\label{g1:eq:lie-covariant-evaluation}
	(\mathcal L_XT)(Y_1,\ldots,Y_k)
	=X\bigl(T(Y_1,\ldots,Y_k)\bigr)
	-\sum_{a=1}^kT(Y_1,\ldots,[X,Y_a],\ldots,Y_k).
\end{equation}
右端第一项对整个实值函数求导，后面的各项扣除输入向量场随流变化的影响。对混合张量，还要对每个协变向量输入扣除对应的 \(\mathcal L_X\alpha\) 项；这由相同的求值乘积法则得到。只保留第一个方向导数而省略输入的修正，一般不会得到张量运算。

\begin{example}[膨胀下的张量变化]
	\label{g1:exm:lie-dilation-tensors}
	在 \(\mathbb R^n\) 上令 \(X=\sum_i x^i\partial_i\)，其流为 \(\Phi_t(x)=e^t x\)。常向量场和坐标一形式的拉回分别满足
	\[
		(\Phi_t)^*\partial_i=e^{-t}\partial_i,\qquad
		(\Phi_t)^*(\mathrm{d}x^i)=e^t\,\mathrm{d}x^i.
	\]
	所以 \(\mathcal L_X\partial_i=-\partial_i\)，而 \(\mathcal L_X(\mathrm{d}x^i)=\mathrm{d}x^i\)。Euclid 内积与标准最高次形式因而满足
	\[
		\mathcal L_X\left(\sum_i\mathrm{d}x^i\otimes\mathrm{d}x^i\right)
		=2\sum_i\mathrm{d}x^i\otimes\mathrm{d}x^i,\qquad
		\mathcal L_X(\mathrm{d}x^1\wedge\cdots\wedge\mathrm{d}x^n)
		=n\,\mathrm{d}x^1\wedge\cdots\wedge\mathrm{d}x^n.
	\]
	这些张量的坐标系数都是常数，Lie 导数却不为零，因为膨胀会改变向量输入的尺度。相反，对任意常数倍平移场，流的微分是恒等映射，这些常系数张量的 Lie 导数为零。
\end{example}

\subsection{自然性及分次 Cartan 运算的准备}
\label{g1:subsec:1303:04}

讨论映射与 Lie 括号的相容性时，不能沿任意光滑映射直接“推出一个向量场”。微分确实把单个切向量向前送，但同一目标点可能有多个原像，得到的向量未必一致；也可能目标方向根本不在微分的像中。需要明确要求两边的场彼此对应。

\begin{definition}
	\label{g1:def:related-vector-fields}
	设 \(F:M\to N\) 为无边界光滑流形之间的光滑映射。若 \(X\in\mathfrak X(M)\)、\(\overline X\in\mathfrak X(N)\) 满足
	\[
		\mathrm{d}F_pX_p=\overline X_{F(p)}\qquad(p\in M),
	\]
	称它们为 \(F\)-\term[xiangguanxiangliangchang]{相关向量场}[related vector fields]。
\end{definition}

由\cref{g1:def:manifold-differential}，相关条件等价于对目标上的每个局部光滑函数 \(f\) 有
\[
	X(f\circ F)=(\overline Xf)\circ F.
\]
反向也成立，因为将这一等式用于目标坐标函数，就能逐分量确定两个目标切向量相等。因此相关条件可以用函数检验，无需给 \(F\) 额外的可逆性。

\begin{proposition}[Lie 导数的自然性]
	\label{g1:prop:lie-derivative-natural}
	设 \(F:M\to N\) 为无边界光滑流形之间的光滑映射，\(X,\overline X\) 为 \(F\)-相关向量场。
	\begin{enumerate}
		\item 若 \(Y,\overline Y\) 也 \(F\)-相关，则 \([X,Y]\) 与 \([\overline X,\overline Y]\) \(F\)-相关。
		\item 记 \(X,\overline X\) 的局部流为 \(\Phi,\overline\Phi\)，则在每个初始点附近的充分小共同时间内，
		\[
			F\circ\Phi_t=\overline\Phi_t\circ F.
		\]
		\item 对 \(N\) 上任意光滑协变张量 \(T\)，有
		\[
			\mathcal L_X(F^*T)=F^*(\mathcal L_{\overline X}T).
		\]
		若 \(F\) 为局部微分同胚，同式也适用于任意混合张量。对任意微分形式 \(\alpha\)，无须可逆性即有
		\[
			\iota_X(F^*\alpha)=F^*(\iota_{\overline X}\alpha).
		\]
	\end{enumerate}
\end{proposition}
\begin{proof}
	对任意目标局部函数 \(f\)，连续使用相关条件可得
	\[
		\begin{aligned}
			[X,Y](f\circ F)
			&=X\bigl((\overline Yf)\circ F\bigr)
			-Y\bigl((\overline Xf)\circ F\bigr)\\
			&=(\overline X\overline Yf-\overline Y\overline Xf)\circ F
			=([\overline X,\overline Y]f)\circ F.
		\end{aligned}
	\]
	由函数检验条件得到第一项。这里既没有使用单射性，也没有使用满射性。

	对固定 \(p\)，曲线 \(t\mapsto F(\Phi_t(p))\) 的导数为
	\[
		\mathrm{d}F_{\Phi_t(p)}X_{\Phi_t(p)}
		=\overline X_{F(\Phi_t(p))},
	\]
	初值为 \(F(p)\)。它与 \(t\mapsto\overline\Phi_t(F(p))\) 满足同一初值问题，故在共同时间内相同。局部流的开定义域和连续性允许在 \(p\) 的一个邻域中同时选取充分小的时间区间，得到第二项的局部映射恒等式。

	协变张量的拉回始终有定义，并与复合相容，见\cref{g1:prop:covariant-tensor-pullback}。将第二项拉回，得到
	\[
		\Phi_t^*(F^*T)=F^*(\overline\Phi_t^*T).
	\]
	对 \(t\) 在零处求导，得到第三项的 Lie 导数公式。右边的 \(F^*\) 由固定的微分 \(\mathrm{d}F\) 作用于各输入，局部表达是固定系数与张量分量的有限次乘法和求和，所以求导可以穿过它。若 \(F\) 为局部微分同胚，各微分可逆，\cref{g1:prop:mixed-tensor-pullback}使同一个论证也适用于混合张量。

	最后，对任意 \(v_1,\ldots,v_{k-1}\in \mathrm{T}_pM\)，有
	\[
		\begin{aligned}
			\bigl(\iota_XF^*\alpha\bigr)_p(v_1,\ldots,v_{k-1})
			&=\alpha_{F(p)}(\mathrm{d}F_pX_p,\mathrm{d}F_pv_1,\ldots,\mathrm{d}F_pv_{k-1})\\
			&=\alpha_{F(p)}(\overline X_{F(p)},\mathrm{d}F_pv_1,\ldots,\mathrm{d}F_pv_{k-1}),
		\end{aligned}
	\]
	恰为右端的值。零次形式时两边均为零。
\end{proof}

若 \(F\) 是微分同胚，可以定义目标向量场
\[
	(F_*X)_q=\mathrm{d}F_{F^{-1}(q)}X_{F^{-1}(q)}.
\]
此时相关性自动成立，命题给出 \(F_*[X,Y]=[F_*X,F_*Y]\)。但对于投影 \(P(x,y)=x\)，向量场 \(y\partial_x\) 在不同原像 \((x,y)\) 上会送出不同的目标速度，因此没有与它相关的目标向量场。自然性命题的假设不能由“映射光滑”代替。

\ProseParagraphBreak

最后整理微分形式上的两个运算。\(\mathcal L_X\) 保持形式次数，\(\iota_X\) 将次数降低 \(1\)，并在函数上取零。\cref{g1:prop:interior-product-rule}已经证明
\[
	\iota_X(\alpha\wedge\beta)
	=(\iota_X\alpha)\wedge\beta+(-1)^p\alpha\wedge(\iota_X\beta),
	\qquad \alpha\in\Omega^p(M).
\]
符号来自将被插入的向量越过前面 \(p\) 个输入位置。把这类带次数的运算放在一起时，合适的交换子是
\[
	[A,B]_{\mathrm{gr}}=A\circ B-(-1)^{ab}B\circ A,
\]
其中 \(a,b\) 是各自改变形式次数的数。这称为分次交换子；两个次数为 \(-1\) 的运算之间对应的是反交换和，而不是普通差。

\begin{proposition}[插入与 Lie 导数]
	\label{g1:prop:cartan-insertion-identities}
	设 \(M\) 为无边界光滑流形，\(X,Y\in\mathfrak X(M)\)，\(\alpha,\beta\) 为微分形式。则
	\[
		\mathcal L_X(\alpha\wedge\beta)
		=(\mathcal L_X\alpha)\wedge\beta+\alpha\wedge(\mathcal L_X\beta),
	\]
	并且
	\begin{align*}
		\iota_X\iota_Y+\iota_Y\iota_X&=0,\\
		\mathcal L_X\iota_Y-\iota_Y\mathcal L_X&=\iota_{[X,Y]},\\
		\mathcal L_X\mathcal L_Y-\mathcal L_Y\mathcal L_X&=\mathcal L_{[X,Y]}.
	\end{align*}
	最后一式也适用于任意混合张量场。
\end{proposition}
\begin{proof}
	拉回保持楔积，故对 \(\Phi_t^*(\alpha\wedge\beta)=\Phi_t^*\alpha\wedge\Phi_t^*\beta\) 在零时刻求导，得到第一式。Lie 导数的次数是零，因此第二项前没有额外的分次符号。

	对至少二次的形式 \(\alpha\)，先插入 \(Y\)，再插入 \(X\)，得到
	\[
		(\iota_X\iota_Y\alpha)(v_1,\ldots)=\alpha(Y,X,v_1,\ldots).
	\]
	反向次序给出 \(\alpha(X,Y,v_1,\ldots)\)，由交替性二者相加为零。次数为零或一时，两次插入都给出零，因此第一条交换关系对所有次数成立。

	取 \(X\) 的局部流 \(\Phi\)。\cref{g1:prop:lie-derivative-natural}的插入相容性用于微分同胚 \(\Phi_t\)，其中源向量场是 \(\Phi_t^*Y\)，得到
	\[
		\Phi_t^*(\iota_Y\alpha)
		=\iota_{\Phi_t^*Y}(\Phi_t^*\alpha).
	\]
	插入对向量场和形式各自线性，且逐点为双线性运算。对右端用乘积求导法则，并用 \(\left.\partial_t\right|_0\Phi_t^*Y=[X,Y]\)，便得
	\[
		\mathcal L_X(\iota_Y\alpha)
		=\iota_{[X,Y]}\alpha+\iota_Y(\mathcal L_X\alpha),
	\]
	这就是第二条交换关系。

	为证明最后一条，记 \(D=\mathcal L_X\mathcal L_Y-\mathcal L_Y\mathcal L_X-\mathcal L_{[X,Y]}\)。两个次数为零的张量积导子的交换子仍是张量积导子：将它作用于 \(T\otimes S\)，展开后交叉项 \((\mathcal L_XT)\otimes(\mathcal L_YS)\) 与 \((\mathcal L_YT)\otimes(\mathcal L_XS)\) 分别抵消，余下的正是两端各自求导之和。因此 \(D\) 是张量积导子，并且与收缩相容。

	在函数上，\(Df=XYf-YXf-[X,Y]f=0\)。在向量场 \(Z\) 上，由 Jacobi 恒等式，
	\[
		DZ=[X,[Y,Z]]-[Y,[X,Z]]-[[X,Y],Z]=0.
	\]
	对任意一形式 \(\eta\) 与局部向量场 \(Z\)，求值是张量积之后的收缩，所以
	\[
		0=D(\eta(Z))=(D\eta)(Z)+\eta(DZ)=(D\eta)(Z).
	\]
	取坐标向量场逐个检验，得到 \(D\eta=0\)。任意张量在局部都是光滑函数乘若干向量场与一形式的张量积之和，导子法则遂给出 \(DT=0\)。这证明了所有混合张量上的最后一式，形式情形包含在其中。
\end{proof}

这些等式把运动、插入和交替性联系起来。在\secref{g1:sec:1401}中，将构造把形式次数提高 \(1\) 的外微分 \(\mathrm{d}\)，并证明 Cartan 公式\footnote{Cartan，Élie Joseph Cartan（1869--1951），法国数学家，研究 Lie 群、Lie 代数与微分几何，发展了活动标架方法。}
\[
	\mathcal L_X=\mathrm{d}\iota_X+\iota_X\mathrm{d}.
\]
这里的加号与两个奇次数运算的分次交换子相符；定义外微分以后，将给出该公式及其证明。
% 纲要条目排为编号小节标题，并进入目录与书签；说明段仍保留于纲要注释中。

% !TeX root = ../../Book1.tex
% 纲要标题：### 13.4 Frobenius 定理与积分叶
% 纲要位置：几何学计划纲要.md，第 627 行。
% 正文状态：初稿；尚未正式审核。

% 纲要内容开始
% 
% * 分布、局部标架与积分流形
% * 对合条件的必要性
% * 对合分布的交换化标架
% * 适应坐标和局部积分叶的构造
% * 最大连通积分叶及唯一性
% 
% 纲要内容结束

\section{Frobenius 定理与积分叶}
\label{g1:sec:1304}

给定一条非零向量场，可以沿其积分曲线组织局部坐标。如果在每一点指定的是一个 \(r\) 维切平面，能否找到 \(r\) 维子流形，使它在各点的切空间恰好等于所指定的平面？这里的困难不只是分别沿 \(r\) 个方向求积分曲线：先沿一个方向、再沿另一个方向移动，可能产生离开原平面的额外方向。Lie 括号正好测量这种不相容性。

\ProseParagraphBreak

本节始终在有限维、Hausdorff、第二可数且无边界的光滑流形上讨论常秩光滑分布。证明先把对合分布的局部标架换成两两交换的标架，再用交换流构造坐标。局部坐标片的粘合会给出最大连通积分叶；叶片的粘合还将赋予叶自身的拓扑，并保证 Hausdorff 性和第二可数性。

\subsection{分布、局部标架与积分流形}
\label{g1:subsec:1304:01}

\begin{definition}[光滑分布]\label{g1:def:smooth-distribution}
	设 \(M^n\) 为无边界光滑流形，\(0\le r\le n\)。切丛 \(\mathrm{T}M\) 的秩 \(r\) 光滑向量子丛 \(D\subseteq \mathrm{T}M\) 称为一个秩 \(r\) 的\term[guanghuafenbu]{光滑分布}[smooth distribution]。等价地，在每个 \(p\in M\) 的某邻域 \(U\) 上，存在光滑向量场 \(Z_1,\ldots,Z_r\)，使它们在每一点线性无关，且
	\[
		D_q=\operatorname{span}\{Z_1(q),\ldots,Z_r(q)\}\qquad(q\in U).
	\]
	这样的向量场组称为 \(D\) 的局部标架。记 \(\Gamma(U,D)\) 为 \(U\) 上所有取值于 \(D\) 的光滑向量场。
\end{definition}

这个等价描述正是\cref{g1:def:vector-subbundle}对 \(\mathrm{T}M\) 的应用。固定维数还不够，局部标架的光滑性也属于条件。另一方面，“取值于 \(D\) 的光滑向量场”确实是子丛的光滑截面：由\cref{g1:prop:subbundle-adapted-frame}把局部标架补成 \(\mathrm{T}M\) 的标架后，后面的分量恒为零，前面的分量仍光滑。因而任意 \(X\in\Gamma(U,D)\) 都唯一写为
\[
	X=\sum_{i=1}^r f^iZ_i
\]
且各 \(f^i\) 光滑。

\begin{definition}[积分流形]\label{g1:def:integral-manifold}
	设 \(D\subseteq \mathrm{T}M\) 为秩 \(r\) 光滑分布。一个 \(r\) 维无边界光滑流形 \(N\)，连同单射浸入 \(\iota:N\to M\)，若满足
	\[
		\mathrm{d}\iota_q(\mathrm{T}_qN)=D_{\iota(q)}\qquad(q\in N),
	\]
	则称为 \(D\) 的一个\term[jifenliuxing]{积分流形}[integral manifold]。把 \(N\) 与其像识别时，像使用从 \(N\) 搬来的流形拓扑和光滑结构，不预先要求它们等于子空间拓扑及嵌入子流形结构。若每个 \(p\in M\) 都位于某个积分流形的像内，则称 \(D\) 可积。
\end{definition}

这里单射使一个像点只有一个参数点，浸入则保证局部参数确实给出一块 \(r\) 维曲面。然而，\cref{g1:exm:injective-immersion-not-embedding}已经说明，单射浸入未必是嵌入：可能有环境空间中相近的点，在参数流形中相距很远。因此积分流形的拓扑可以比像的子空间拓扑更细。后面构造最大叶时，会明确选出这个拓扑。

\begin{example}\label{g1:exm:submersion-integral-levels}
	若 \(F:M^n\to N^{n-r}\) 是光滑淹没，则 \(D=\ker\mathrm{d}F\) 为秩 \(r\) 光滑分布。其子丛性质来自\cref{g1:thm:constant-rank-bundle-morphism}：把 \(\mathrm{d}F\) 看成 \(\mathrm{T}M\to F^*\mathrm{T}N\) 的满射丛同态，其核秩恒为 \(r\)。对每个非空水平集 \(F^{-1}(y)\)，\cref{g1:thm:regular-level-set}给出嵌入子流形结构，且
	\[
		\mathrm{T}_pF^{-1}(y)=\ker\mathrm{d}F_p=D_p.
	\]
	因此水平集及其连通分支都是积分流形。在这个例子中，积分流形由一些保持不变的函数值刻画；\term[frobeniusdingli]{Frobenius 定理}[Frobenius theorem]\footnote{Frobenius，Ferdinand Georg Frobenius（1849--1917），德国数学家，研究群的表示与特征标以及微分方程。}将证明，对合条件正好保证局部总能找到这样的横向坐标函数。
\end{example}

秩零分布的连通积分流形是单点；秩 \(n\) 的分布只能是 \(\mathrm{T}M\)，其局部积分流形是开集。以下证明也包含这两个情形，空的标架或横向坐标按通常的零维约定理解。

\subsection{对合条件的必要性}
\label{g1:subsec:1304:02}

若两条向量场都沿同一个积分流形运动，它们的无穷小交换也应当沿这个流形。为把这个直观判断写成坐标无关的条件，我们使用\cref{g1:def:lie-bracket}中的括号
\[
	[X,Y]f=X(Yf)-Y(Xf).
\]

\begin{definition}[对合分布]\label{g1:def:involutive-distribution}
	设 \(D\subseteq \mathrm{T}M\) 为光滑分布。若对任意开集 \(U\subseteq M\) 及任意 \(X,Y\in\Gamma(U,D)\)，都有
	\[
		[X,Y]\in\Gamma(U,D),
	\]
	则称 \(D\) 为\term[duihefenbu]{对合分布}[involutive distribution]。
\end{definition}

检查对合性时，只需检查一组局部标架的两两括号。事实上，由\cref{g1:prop:lie-bracket-properties}的乘积法则，
\[
	\left[\sum_i f^iZ_i,\sum_j g^jZ_j\right]
	=\sum_{i,j}f^ig^j[Z_i,Z_j]
	+\sum_{i,j}f^iZ_i(g^j)Z_j
	-\sum_{i,j}g^jZ_j(f^i)Z_i.
\]
后两项自动取值于 \(D\)，第一项则由标架括号条件控制。对合不要求原标架彼此交换。例如 \(\mathrm{T}\mathbb R^2\) 的标架 \(\partial_x,e^x\partial_y\) 满足 \([\partial_x,e^x\partial_y]=e^x\partial_y\ne0\)，但括号仍在该分布中。秩一分布自动对合，因为 \([fX,gX]=(fXg-gXf)X\)。

\begin{proposition}[可积性的必要条件]\label{g1:prop:involutivity-necessary}
	设 \(M\) 为无边界光滑流形，\(D\subseteq \mathrm{T}M\) 为常秩光滑分布。若每个点都有通过它的积分流形，则 \(D\) 对合。更具体地，若 \(\iota:N\to M\) 是一个积分流形，\(X,Y\) 为某开集 \(U\) 上取值于 \(D\) 的光滑向量场，则 \([X,Y]\) 在 \(U\cap\iota(N)\) 上也取值于 \(D\)。
\end{proposition}
\begin{proof}
	在 \(\iota^{-1}(U)\) 上，因 \(\mathrm{d}\iota_q:\mathrm{T}_qN\to D_{\iota(q)}\) 为线性同构，可唯一定义向量场 \(\widetilde X,\widetilde Y\)，使
	\[
		\mathrm{d}\iota\,\widetilde X=X\circ\iota,\qquad
		\mathrm{d}\iota\,\widetilde Y=Y\circ\iota.
	\]
	它们光滑。具体地，在局部标架中，\(\mathrm{d}\iota\) 给出到拉回分布 \(\iota^*D\) 的光滑可逆矩阵；其逆由伴随矩阵除以非零行列式给出，故光滑，正如\cref{g1:prop:bundle-morphism-matrices}所述。

	相关向量场的括号仍相关，见\cref{g1:prop:lie-derivative-natural}。当前所用的计算很短：对 \(f\in C^\infty(U)\)，有
	\[
		\begin{aligned}
			[\widetilde X,\widetilde Y](f\circ\iota)
			&=\widetilde X\bigl((Yf)\circ\iota\bigr)-\widetilde Y\bigl((Xf)\circ\iota\bigr)\\
			&=([X,Y]f)\circ\iota.
		\end{aligned}
	\]
	以导子判断切向量，得到
	\[
		[X,Y]_{\iota(q)}
		=\mathrm{d}\iota_q\bigl([\widetilde X,\widetilde Y]_q\bigr)
		\in \mathrm{d}\iota_q(\mathrm{T}_qN)=D_{\iota(q)}.
	\]
	若每个点都有这样的积分流形，就在每一点得到所需的包含关系，故 \([X,Y]\) 为 \(D\) 的光滑截面。
\end{proof}

浸入微分把 \(\mathrm{T}N\) 同构到沿像的 \(D\)，因而能将局部的分布截面提升到积分流形；括号自然性随后把它们的括号送回 \(D\)。

\subsection{对合分布的交换化标架}
\label{g1:subsec:1304:03}

对合条件只说 \([Z_i,Z_j]\) 能用原标架线性表示，而交换流需要的是括号等于零。关键在于先选择若干坐标方向，使它们在 \(D\) 上能够独立测量所有分量，再把这些分量归一化。这样新标架的括号既必须留在 \(D\)，又失去了在这些方向上的全部分量，只能为零。

\begin{lemma}[交换化分布标架]\label{g1:lem:commuting-distribution-frame}
	设 \(D\) 为无边界光滑流形 \(M^n\) 上的秩 \(r\) 对合光滑分布。对每个 \(p\in M\)，存在以 \(p\) 为原点的坐标
	\[
		(x^1,\ldots,x^r,y^1,\ldots,y^{n-r})
	\]
	及 \(D\) 的光滑局部标架 \(X_1,\ldots,X_r\)，满足
	\begin{equation}\label{g1:eq:normalized-distribution-frame}
		X_i=\frac{\partial}{\partial x^i}
		+\sum_{\alpha=1}^{n-r}a_i^\alpha(x,y)\frac{\partial}{\partial y^\alpha},
		\qquad [X_i,X_j]=0.
	\end{equation}
\end{lemma}
\begin{proof}
	先取任意局部标架 \(Z_1,\ldots,Z_r\) 和环境坐标 \(z^1,\ldots,z^n\)。在 \(p\) 处，系数矩阵 \((Z_jz^a)_{a,j}\) 的秩为 \(r\)，所以有一个非零的 \(r\) 阶子式。重排坐标后，可令它由前 \(r\) 行给出，记这部分坐标为 \(x\)，其余为 \(y\)。缩小邻域使
	\[
		B_{aj}=Z_jx^a\qquad(1\le a,j\le r)
	\]
	始终可逆，并定义
	\[
		X_i=\sum_{j=1}^r(B^{-1})_{ji}Z_j.
	\]
	逆矩阵光滑，新向量场仍为 \(D\) 的标架；并且
	\[
		X_ix^a=\sum_jB_{aj}(B^{-1})_{ji}=\delta_i^a.
	\]
	这就给出\eqref{g1:eq:normalized-distribution-frame}中的第一式。

	投影到前 \(r\) 个坐标的微分 \(\mathrm{d}x:D_q\to\mathbb R^r\) 在整个邻域中为同构，因为它把 \(X_i(q)\) 送到第 \(i\) 个标准基向量。另一方面，
	\[
		[X_i,X_j]x^a=X_i(\delta_j^a)-X_j(\delta_i^a)=0.
	\]
	所以括号的全部 \(x\) 分量为零。同时，对合性使它属于 \(D_q\)。一个属于 \(D_q\) 且被 \(\mathrm{d}x\) 送到零的向量只能为零，故 \([X_i,X_j]=0\)。
\end{proof}

这一步只换了分布的标架，还没有找到使 \(X_i\) 本身成为坐标向量的坐标。后一任务由流完成。若直接展开刚才消失的横向分量，还可把对合条件写为
\[
	\partial_{x^i}a_j^\alpha-\partial_{x^j}a_i^\alpha
	+\sum_{\beta=1}^{n-r}\left(
		a_i^\beta\partial_{y^\beta}a_j^\alpha
		-a_j^\beta\partial_{y^\beta}a_i^\alpha
	\right)=0.
\]
这些是系数之间的相容方程；不必另外求解这组偏微分方程，交换的局部流会直接给出积分片。

\subsection{适应坐标和局部积分叶的构造}
\label{g1:subsec:1304:04}

交换化标架给出沿各个允许方向的交换流。把这些流依次作用在一张横向小截面上，流时间与截面坐标便组成候选坐标；逆函数定理将证明它们确实覆盖一个邻域。

\begin{theorem}[Frobenius 定理]\label{g1:thm:frobenius}
	设 \(M^n\) 为有限维、Hausdorff、第二可数且无边界的光滑流形，\(D\subseteq \mathrm{T}M\) 为秩 \(r\) 光滑分布。下列条件等价：
	\begin{enumerate}
		\item \(D\) 对合；
		\item 每个点附近都有坐标 \((u^1,\ldots,u^r,v^1,\ldots,v^{n-r})\)，使
		\[
			D=\operatorname{span}\left\{\frac{\partial}{\partial u^1},\ldots,\frac{\partial}{\partial u^r}\right\};
		\]
		\item 每个点都位于某个 \(r\) 维积分流形内。
	\end{enumerate}
	第二项中的坐标域可取为乘积 \(B\times T\)，其中 \(B\subseteq\mathbb R^r\)、\(T\subseteq\mathbb R^{n-r}\) 为开球或开盒；固定 \(v=c\) 的切片给出局部嵌入积分流形。
\end{theorem}
\begin{proof}
	设 \(D\) 对合。固定 \(p\)，用\cref{g1:lem:commuting-distribution-frame}取得局部坐标 \((x,y)\) 和两两交换的标架 \(X_1,\ldots,X_r\)。令 \(\Phi_i^t\) 为 \(X_i\) 的局部流。\cref{g1:thm:maximal-local-flow,g1:thm:ode-smooth-dependence}保证它们在零时间附近存在，且对时间与初值联合光滑；不需要这些向量场完备。由\cref{g1:thm:commuting-local-flows}，在足够小的共同邻域内，各局部流两两交换。

	取横截片 \(\sigma(y)=(0,y)\)，定义
	\begin{equation}\label{g1:eq:frobenius-flow-coordinates}
		\Psi(t,y)=\Phi_r^{t^r}\circ\cdots\circ\Phi_1^{t^1}(\sigma(y)).
	\end{equation}
	这时必须使一个参数盒内的全部复合都有定义。因为只有有限条流及交换时所需的有限种复合，零参数时所有中间点都在 \(p\) 附近；各流的定义域开放，故可同时缩小时间盒和横截片，使这些复合及其中间点都留在上述共同邻域。于是 \(\Psi\) 是一个开参数盒上的光滑映射。

	利用交换性，把第 \(i\) 条流移到最外层，便有
	\[
		\Psi(t+s e_i,y)=\Phi_i^s(\Psi(t,y))
	\]
	对充分小的 \(s\) 成立。在 \(s=0\) 求导得到
	\begin{equation}\label{g1:eq:frobenius-pushed-coordinate-fields}
		\mathrm{d}\Psi_{(t,y)}\!\left(\frac{\partial}{\partial t^i}\right)
		=X_i(\Psi(t,y)).
	\end{equation}
	在 \((0,0)\) 处，时间方向的导数为 \(X_i(p)\)，横截方向的导数为原坐标的 \(\partial_{y^\alpha}|_p\)。因此按 \((x,y)\) 排列的导数矩阵为
	\[
		\mathrm{D}\Psi(0,0)=
		\begin{pmatrix}
			I_r&0\\
			A&I_{n-r}
		\end{pmatrix},
		\qquad A_{\alpha i}=a_i^\alpha(p),
	\]
	它可逆。\cref{g1:thm:inverse-function}给出 \(\Psi\) 在原点附近的光滑逆；再把参数域缩小为乘积 \(B\times T\)，其像仍为开集 \(U\)，且 \(\Psi:B\times T\to U\) 为微分同胚。

	取 \((u,v)=\Psi^{-1}\) 为 \(U\) 上的新坐标。\eqref{g1:eq:frobenius-pushed-coordinate-fields}说明原来的 \(X_i\) 正好成为新坐标向量 \(\partial_{u^i}\)，从而得到第二项。秩零时没有流，直接使用原坐标；满秩时横截片是单点，同一个可逆矩阵论证仍成立。

	第二项推出第三项：固定 \(v=c\)，映射 \(u\mapsto\Psi(u,c)\) 是一块嵌入参数片，其微分像正是 \(D\)，故为积分流形。第三项推出第一项已经由\cref{g1:prop:involutivity-necessary}证明。
\end{proof}

在这样的适应坐标中，把 \(U\) 称为一个\term[yezhuangzuobiaohe]{叶状坐标盒}[foliation box]；对 \(c\in T\)，把
\[
	P_c=\Psi(B\times\{c\})
\]
称为该盒的一张\term[yepian]{叶片}[plaque]。叶片是连通的局部嵌入积分流形。这里的“叶片”只指一个盒中的局部切片，后面所说的“叶”则是最大连通积分流形，不宜混用。\cref{g1:fig:frobenius-plaques}中，横向坐标 \(v\) 区分不同叶片，沿 \(D\) 的运动只改变 \(u\)。

\input{Figures/Chapter13-FrobeniusPlaques.tex}

\begin{lemma}[局部积分叶片的唯一性]\label{g1:lem:local-integral-plaque}
	设 \(M\) 为无边界光滑流形，\(D\) 为其上的秩 \(r\) 光滑分布。设开集 \(U\subset M\) 具有乘积坐标 \((u,v)\in B\times T\)，其中 \(B,T\) 为开球，且 \(D|_U=\operatorname{span}\{\partial_{u^1},\ldots,\partial_{u^r}\}\)。对 \(c\in T\)，以 \(P_c\) 记切片 \(\{(u,c):u\in B\}\) 在 \(U\) 中的原像。若 \(\iota:N\to M\) 是连通积分流形且 \(\iota(N)\subseteq U\)，则 \(v\circ\iota\) 为常值 \(c\)，且
	\[
		\iota:N\longrightarrow P_c
	\]
	是到 \(P_c\) 某开子集的微分同胚。因此，任何两个通过同一点的积分流形，在分别缩小其参数邻域后，像在该点附近一致。
\end{lemma}
\begin{proof}
	由于 \(\mathrm{d}\iota(\mathrm{T}N)=D\)，而每个横向坐标 \(v^\alpha\) 的微分在 \(D\) 上为零，故
	\[
		\mathrm{d}(v^\alpha\circ\iota)=0.
	\]
	光滑函数微分为零时局部常值：在源坐标小球内沿线段使用一元微积分，函数值不变。连通空间上的局部常值函数为常值，因为一个取值的原像及其补集都是开集。因此 \(v\circ\iota=c\)。

	前 \(r\) 个坐标的微分在 \(D\) 上为同构，所以
	\[
		\mathrm{d}(u\circ\iota)_q
		=\mathrm{d}u_{\iota(q)}\circ\mathrm{d}\iota_q:\mathrm{T}_qN\to\mathbb R^r
	\]
	可逆。逆函数定理表明 \(\iota:N\to P_c\) 是局部微分同胚，因此为开映射。它又单射，故是到其像的同胚，局部光滑逆粘合成整体光滑逆，像在 \(P_c\) 中开。

	最后，对两条通过同一点 \(p\) 的积分流形，先各取连通参数邻域，使像落在同一个适应坐标盒中。两者的横向常值都等于 \(v(p)\)，像都是同一 \(P_{v(p)}\) 中含 \(p\) 的开集。取这两个开集的交，再限制各参数邻域，即得到相同的局部像。
\end{proof}

若一条积分流形没有整个落在 \(U\) 内，只能把引理用于其原像 \(\iota^{-1}(U)\) 的各个连通分支。不同分支可以落在同一个坐标盒的不同叶片上。局部唯一性因此不表示一个完整的积分叶与 \(U\) 只能相交一张叶片，也不保证完整积分叶是嵌入的。

\subsection{最大连通积分叶及唯一性}
\label{g1:subsec:1304:05}

沿相交的叶片继续前进，可以把经过一点的所有局部积分片连在一起。首先需要确定哪些点属于同一片最大对象，然后给这个点集一个使各叶片成为开坐标片的拓扑。这个拓扑下，远行后重新靠近原点的部分不必成为原来那一张局部坐标片中的邻近点。

\begin{theorem}[最大连通积分叶]\label{g1:thm:maximal-connected-leaf}
	设 \(M\) 为有限维、Hausdorff、第二可数且无边界的光滑流形，\(D\subseteq \mathrm{T}M\) 为秩 \(r\) 对合光滑分布。对每个 \(p\in M\)，存在唯一的最大连通积分流形 \(L_p\)，其像包含 \(p\)。更精确地，\(L_p\) 是 \(M\) 的一个子集，带有 Hausdorff、第二可数的 \(r\) 维光滑流形结构，包含映射 \(i_p:L_p\hookrightarrow M\) 为浸入，且任何像包含 \(p\) 的连通积分流形 \(\iota:N\to M\) 都唯一分解为
	\[
		N\xrightarrow{\ \widehat\iota\ }L_p\xrightarrow{\ i_p\ }M,
	\]
	其中 \(\widehat\iota\) 是开嵌入。不同的 \(L_p\) 要么相同，要么不交，故它们将 \(M\) 分割成最大连通积分流形。这些最大积分流形称为 \(D\) 的\term[ye]{叶}[leaf]。
\end{theorem}
\begin{proof}
	先取至多可数个叶状坐标盒 \(U_1,U_2,\ldots\) 覆盖 \(M\)。第二可数性保证可以这样取：在一个可数拓扑基中，对每个能包含于某个给定坐标盒的基元素选一个这样的盒；每点附近都有这样的基元素，所以选出的至多可数个盒仍覆盖 \(M\)。记第 \(j\) 个盒的坐标为 \((u_j,v_j)\in B_j\times T_j\)，其叶片记为 \(P_{j,c}\)。

	如果存在有限叶片链
	\[
		P_0,P_1,\ldots,P_k,\qquad
		P_{\ell-1}\cap P_\ell\ne\varnothing,
	\]
	使两点分别位于首尾叶片，就称这两点可由叶片相连。每点属于某叶片，反转链给出对称性，拼接链给出传递性，所以这是等价关系。令 \(L_p\) 为 \(p\) 的等价类。若一个叶片与 \(L_p\) 相交，便可加到相应的有限链末尾，故整个叶片都包含于 \(L_p\)。

	先说明相交叶片的拓扑如何衔接。设 \(P,Q\) 为所选坐标盒中的两张叶片，且 \(z\in P\cap Q\)。在 \(P\) 内取 \(P\cap U_Q\) 中含 \(z\) 的一个小连通开邻域，其中 \(U_Q\) 是含 \(Q\) 的坐标盒。这个邻域是积分流形，像落在 \(U_Q\) 中，由\cref{g1:lem:local-integral-plaque}，它位于通过 \(z\) 的唯一叶片 \(Q\) 内。因此 \(P\cap Q\) 在 \(P\) 中开；交换二者，同样在 \(Q\) 中开。两张叶片坐标之间的转换，是环境坐标转换限制到这些开集后取前 \(r\) 个分量，因而光滑，其逆也光滑。

	给 \(L_p\) 如下拓扑：\(W\subseteq L_p\) 开，当且仅当对每张包含于 \(L_p\) 的上述叶片 \(P\)，交集 \(W\cap P\) 在 \(P\) 中开。刚才的交集性质说明每张 \(P\) 在这个拓扑中开，而且它的诱导拓扑恰好是原来的 Euclid 片拓扑。例如 \(P\) 中任一开集 \(V\) 与其他叶片 \(Q\) 的交，在开集 \(P\cap Q\) 内经光滑坐标转换仍为开集，因此 \(V\) 也是 \(L_p\) 的开集。各叶片坐标遂组成光滑相容图册。包含映射在每个片上就是坐标切片的包含，所以连续、光滑且微分像为 \(D\)。任意两点可在 \(M\) 中用不交开邻域分开，取这些邻域与 \(L_p\) 的交，仍为叶拓扑中的开集，故 \(L_p\) Hausdorff。

	接着验证第二可数性。这不能由“\(L_p\) 是第二可数空间 \(M\) 的子集”推出，因为刚定义的叶拓扑可能比子空间拓扑更细。我们证明这条叶实际只用到可数张叶片。固定一张叶片 \(P\) 以及一个坐标盒 \(U_j\)。集合 \(P\cap U_j\) 是 \(P\) 的开子集，至多有可数个连通分支：在 \(P\) 的 Euclid 坐标中，每个连通分支都是开集，因为每点的小球连通；不同非空分支各含一个不同的有理坐标点，故只有可数个。秩零时开子集至多为单点，结论同样成立。每个分支 \(C\) 是连通积分流形且位于 \(U_j\) 内，由局部唯一性，\(v_j|_C\) 常值，故 \(C\) 只遇到 \(U_j\) 中的一张叶片。因此对固定的 \(P,U_j\)，与 \(P\) 相交的 \(U_j\) 叶片至多可数。

	坐标盒本来只有可数个，故与固定 \(P\) 相交的全部所选叶片至多可数。现在从一张含 \(p\) 的叶片开始：第一步加入所有与它相交的叶片，第二步加入所有与第一步叶片相交的叶片，依此继续。每一步得到可数集合，可数步之并仍可数，并且恰好包含所有能用有限链从 \(p\) 到达的叶片。它们覆盖 \(L_p\)。每张叶片取一个可数 Euclid 拓扑基，所有这些基的并给出 \(L_p\) 的可数基，因为叶片是开坐标片。至此，\(L_p\) 确实是通常意义下的 Hausdorff、第二可数光滑流形。

	每张叶片同胚于开球或开盒，所以可以在片内用光滑路径连接两点。沿任一有限叶片链选择相邻交点并连接这些路径，得到叶内的连续分片光滑路径。因此 \(L_p\) 路径连通，尤其连通。

	最后证明最大性与唯一性。设 \(\iota:N\to M\) 为连通积分流形，且像包含 \(p\)。对任意 \(q\in N\)，选一个连通参数邻域，其像落在某个 \(U_j\) 内。局部唯一性说明这个像落在一张叶片中，所以 \(N\) 中映入同一叶片等价类的点构成开集；这些开集两两不交并覆盖 \(N\)。由连通性只有一个非空类，且它是 \(L_p\)。故 \(\iota\) 在集合意义下唯一地分解为 \(\widehat\iota:N\to L_p\)。

	在刚才的每个小参数邻域内，\(\widehat\iota\) 的像位于一张开叶片中，而\cref{g1:lem:local-integral-plaque}使它成为到该叶片开子集的微分同胚。因此 \(\widehat\iota\) 是局部微分同胚，特别地连续且开。它由 \(\iota\) 的单射性而单射，故为开嵌入。这证明 \(L_p\) 包含经过 \(p\) 的每个连通积分流形的像。任何另一个最大连通积分流形必须具有相同的像，两个结构之间的恒等映射又由上述分解在两个方向光滑，所以最大对象连同流形结构都唯一。最后，等价类本来相同或不交，故叶构成 \(M\) 的分割。
\end{proof}

这个由最大叶组成的分割，连同上述局部乘积坐标，称为 \(M\) 上的一个 \(r\) 维\term[yezhuangjiegou]{光滑叶状结构}[smooth foliation]，常记为 \(\mathcal F\)。各叶在每个坐标盒中分解为若干叶片；叶片之间的切向空间相容，合起来正是原来的分布 \(D\)。

\ProseParagraphBreak

这个证明同时给出了叶状坐标盒的内在用途：它们既提供叶的局部图，又用可数传播控制叶的整体拓扑。对任一固定盒，一条叶可能经过可数多张不同叶片，而不是只有一张；叶中的邻近由单张叶片上的坐标决定，环境中的再次靠近则未必意味着叶内的靠近。

\ProseParagraphBreak

Frobenius 定理由常秩、光滑性与对合条件得到局部适应坐标，并得到具有自身流形结构的最大连通叶。叶是否嵌入、压缩各叶所得商空间是否 Hausdorff，则取决于叶离开局部坐标盒后的整体行为。下一节将在环面流和横截面返回的模型中计算这些行为。

% 纲要条目排为编号小节标题，并进入目录与书签；说明段仍保留于纲要注释中。

% !TeX root = ../../Book1.tex
% 纲要标题：### 13.5 叶状结构的整体问题与不可积模型
% 纲要位置：几何学计划纲要.md，第 635 行。
% 正文状态：初稿；尚未正式审核。

% 纲要内容开始
% 
% * 环面上的有理和无理斜率叶
% * 横截面、返回映射及沿叶映射芽
% * 非 Hausdorff 叶空间
% * 标准接触分布的不可积性
% 
% ---
% 
% 纲要内容结束

\section{叶状结构的整体问题与不可积模型}
\label{g1:sec:1305}

Frobenius 定理给每个点附近一叠坐标切片，最大叶定理则说明怎样沿这些切片继续前进。继续以后会发生什么，仍须另外研究：一条叶可能闭合，也可能反复进入同一个坐标盒，甚至在整个流形中稠密；沿叶绕一圈后，邻近叶的横向位置也可能改变。本节用几个能够直接计算的模型分辨这些现象。全节的流形均无边界，分布均为常秩光滑分布。

\subsection{环面上的有理和无理斜率叶}
\label{g1:subsec:1305:01}

把二维环面写为 \(T^2=\mathbb R^2/\mathbb Z^2\)，即分别把两个坐标模 \(1\) 识别。局部坐标来自小矩形的商像，换坐标只是整数平移，所以常向量场
\[
	X_\alpha=\partial_x+\alpha\partial_y,\qquad \alpha\in\mathbb R,
\]
下降为环面上的处处非零光滑向量场。它的流为
\[
	\Phi_t([x,y])=[x+t,y+\alpha t].
\]
这给出一个处处秩一的分布 \(D_\alpha=\mathbb RX_\alpha\)。秩一分布总对合，因为 \([fX_\alpha,gX_\alpha]=(fX_\alpha g-gX_\alpha f)X_\alpha\)；由\cref{g1:thm:maximal-connected-leaf}，其叶就是这些流轨道，带上各自的积分流形拓扑。

\begin{proposition}[环面上的线性叶]
	\label{g1:prop:linear-torus-leaves}
	设 \(T^2=\mathbb R^2/\mathbb Z^2\)，\(\alpha\in\mathbb R\)，\(X_\alpha=\partial_x+\alpha\partial_y\)，\(D_\alpha=\mathbb R X_\alpha\)，其流为 \(\Phi_t([x,y])=[x+t,y+\alpha t]\)。若 \(\alpha=p/q\)，其中 \(p\in\mathbb Z\)、\(q\in\mathbb Z_{\geq1}\) 且 \(\gcd(p,q)=1\)，则每条 \(D_\alpha\) 叶为嵌入圆周，最小正流周期为 \(q\)。若 \(\alpha\) 无理，则每条叶为 \(\mathbb R\) 的单射浸入像，在整个环面中稠密，且不是嵌入子流形。
\end{proposition}
\begin{proof}
	轨道在时刻 \(t\) 返回初始点，当且仅当 \(t\in\mathbb Z\) 且 \(\alpha t\in\mathbb Z\)。有理情形中这等价于 \(t\in q\mathbb Z\)。于是轨道参数化下降为单射浸入 \(\mathbb R/q\mathbb Z\to T^2\)。定义域紧而目标 Hausdorff，所以它到像为同胚，结合浸入性即为光滑嵌入。

	无理情形没有非零返回时间，故轨道参数化 \(\mathbb R\to T^2\) 单射，其微分是处处非零的 \(X_\alpha\)，因此是浸入。先证明整数倍 \(\{n\alpha\bmod1:n\in\mathbb Z\}\) 在圆周稠密。将 \(0,\alpha,\ldots,N\alpha\) 的小数部分放入 \(N\) 个等长区间，至少两点落在同一段。两者之差取适当正负号，就得到某个非零整数 \(k\) 及整数 \(m\)，使
	\[
		0<\delta=k\alpha-m<\frac1N.
	\]
	各数 \(0,\delta,\ldots,\lfloor1/\delta\rfloor\delta\) 都来自整数倍 \(\alpha\) 模 \(1\)，并且与 \([0,1]\) 中任一点的距离不超过 \(\delta\)。令 \(N\) 任意大，便得稠密性。对环面中任意小开矩形，先固定其中一个第一坐标 \(u\)。轨道在时刻 \(t=u-x_0+n\) 的第一坐标均为 \(u\) 模 \(1\)，第二坐标为
	\[
		y_0+\alpha(u-x_0)+n\alpha\pmod1,
	\]
	这些值稠密，所以某个时刻落入给定矩形。这证明整条叶稠密。

	还须辨认叶的拓扑。由同一逼近论证，可以选 \(|n_j|\to\infty\)，使 \(n_j\alpha\bmod1\to0\)；若这些整数的绝对值有界，则其到整数的非零距离有正的最小值，与任意精度逼近矛盾。于是 \(\Phi_{n_j}(p)\to p\) 在环面中成立，但参数 \(n_j\) 不趋于零。故参数化的逆在子空间像上不连续，它不是嵌入。另一方面，参数化到最大叶是双射局部微分同胚，因为在每个流盒中都是时间坐标的平移，故与叶自身的光滑结构微分同胚。叶本身为 \(\mathbb R\)，稠密性发生在它到环面的浸入像上。
\end{proof}

\input{Figures/Chapter13-TorusLeaves.tex}

局部流盒看不出斜率是否有理：两种情形都能把向量场写成 \(\partial_s\)。区别来自绕过环面后坐标的整数识别。无理叶也没有在一个流盒中占据整块二维区域；它在盒内出现许多互不相连的局部线段，这些线段的并可以稠密，而每一段仍只有一维。

\subsection{横截面、返回映射及沿叶映射芽}
\label{g1:subsec:1305:02}

一条轨道离开小流盒后怎样回来，可以通过它与横截面的交点来记录。附近轨道也给出相应交点，因而沿流运动在横截面之间诱导局部映射。

\begin{definition}
	\label{g1:def:foliation-transversal}
	设 \(D\subset \mathrm{T}M\) 是秩 \(r\) 的可积分布，\(\dim M=n\)。一个 \((n-r)\) 维局部嵌入子流形 \(S\) 若满足
	\[
		\mathrm{T}_pM=D_p\oplus \mathrm{T}_pS\qquad(p\in S),
	\]
	称为这族叶的\term[hengjiemian]{横截面}[transversal]。对于处处非零向量场产生的秩一分布，这表示向量场不切于 \(S\)。
\end{definition}

在适应坐标 \((x,y)\) 中，叶片由 \(y\) 固定给出，横截面可以取 \(x=0\)。由互补条件，横向投影 \((x,y)\mapsto y\) 限制到任意横截面时，其微分可逆；\cref{g1:thm:inverse-function}说明它局部为微分同胚。所以在一个足够小的叶状坐标盒中，附近叶片与横截面的交点可用横向坐标唯一标记。

\begin{proposition}[横截面间的局部返回]
	\label{g1:prop:local-return-map}
	设 \(X\) 为无边界光滑流形 \(M\) 上的光滑向量场，\(\Phi\) 为它的最大局部流，\(\Phi_T(p)=q\)，其中 \(T>0\)，且 \(X\) 沿这段轨道不为零。在 \(p,q\) 附近分别取横截面 \(S,S'\)。存在 \(S\) 中 \(p\) 的邻域 \(V\) 及光滑函数 \(\tau:V\to\mathbb R\)，满足 \(\tau(p)=T\)，使
	\[
		P(z)=\Phi_{\tau(z)}(z)\in S'
	\]
	为到 \(S'\) 中 \(q\) 的一个邻域的微分同胚。这里 \(\tau(z)\) 是接近指定时间 \(T\) 的交会时刻，未必是第一次正返回时间。
\end{proposition}
\begin{proof}
	在 \(q\) 附近将 \(S'\) 写成 \(h=0\)，其中 \(\mathrm{d}h_q\neq0\)。横截性等价于 \(\mathrm{d}h_q(X_q)\neq0\)。\cref{g1:thm:maximal-local-flow}保证 \((t,z)\mapsto\Phi_t(z)\) 在 \((T,p)\) 附近光滑，于是方程
	\[
		H(t,z)=h(\Phi_t(z))=0
	\]
	满足 \(H(T,p)=0\)、\(\partial_tH(T,p)=\mathrm{d}h_q(X_q)\neq0\)。参数隐函数定理给出唯一接近 \(T\) 的光滑解 \(t=\tau(z)\)。这里可在 \(S\) 的局部坐标中直接应用\cref{g1:thm:implicit-function-parameters}。

	沿同一轨道反向运动，用 \(S\) 的局部定义函数作相同构造，就得到 \(S'\) 中 \(q\) 附近到 \(S\) 的光滑交会映射 \(Q\)，时间接近 \(-T\)。适当缩小邻域后，\(Q(P(z))=z\)：原来的 \(z\) 已是接近 \(-T\) 时刻的交点，而隐函数定理保证此交点唯一。同理 \(P(Q(w))=w\)。故两映射互逆。
\end{proof}

若 \(q=p\) 且 \(S'=S\)，就得到闭轨道附近的\term[fanhuiyingshe]{返回映射}[return map]。一般轨道可能不再遇到选定截面，也可能在所选 \(T\) 之前遇到截面；存在一个局部交会映射，不等于全体点具有共同定义的首次返回映射。在环面线性流中，截面 \(S=\{x=0\}\) 则确实是整体横截面，所有点的首次正返回时间都是 \(1\)，返回映射为
\[
	P([y])=[y+\alpha].
\]
于是有理斜率给出有限次迭代后的返回，无理斜率给出无周期的圆周旋转。

\ProseParagraphBreak

高维叶没有统一的“沿叶前进时间”，但仍可指定一条叶内道路，并看附近叶怎样跟随它穿过一串坐标盒。设 \(\gamma:[0,1]\to L\) 是一条叶内道路，端点为 \(p,q\)，在两端固定横截面 \(S_p,S_q\)。取有限分割，使每小段道路都落在一张叶状坐标盒内。沿每一段保持横向坐标不变，抵达下一张小横截面；在交叠处使用光滑换坐标，再作下一段。有限次复合得到 \(p\) 附近到 \(q\) 附近的局部微分同胚。

\begin{definition}
	\label{g1:def:leaf-holonomy}
	设 \(D\) 为无边界光滑流形 \(M\) 上的常秩光滑可积分布，\(L\) 为 \(D\) 的一条最大连通积分叶，\(\gamma:[0,1]\to L\) 为叶内道路，\(p=\gamma(0)\)、\(q=\gamma(1)\)，并在两端固定局部横截面 \(S_p,S_q\)。沿覆盖道路的有限个适应坐标盒，在每一张盒中按相同横向坐标作叶片对应，再依次复合，得到 \(S_p\) 中 \(p\) 附近到 \(S_q\) 中 \(q\) 附近的局部微分同胚。两个这样的局部映射若在 \(p\) 的某个更小邻域相同，就称它们在 \(p\) 有相同的映射芽。上述复合在 \(p\) 的芽记为 \(h_\gamma\)，称为\term[yeholonomy]{沿叶道路的横截面映射芽}[leaf holonomy]。当 \(\gamma\) 为基于 \(p\) 的环路时，在两端采用同一横截面，所得 \(h_\gamma\) 是固定 \(p\) 的横截面返回映射芽；所有这样的芽在复合下组成该叶在 \(p\) 的横截面返回映射芽群。
\end{definition}

采用芽的原因是，改变坐标盒大小会改变代表映射的定义域，却不应改变靠近初始点的横向效应。这里移动的是邻近叶的交点，并没有选择联络，也没有沿曲线搬运切向量；它与\chapref{g1:ch:09}的环路平行移动变换是不同的构造。

\begin{proposition}[沿叶横截面映射芽的相容性]
	\label{g1:prop:leaf-holonomy-homotopy}
	设 \(D\) 为无边界光滑流形上的常秩光滑可积分布，固定道路两端的局部横截面。沿叶道路的横截面映射芽不依赖所选坐标盒与分割，只依赖道路在该叶中相对端点的同伦类。若 \(\gamma*\eta\) 表示先走 \(\gamma\) 后走 \(\eta\)，则
	\[
		h_{\gamma*\eta}=h_\eta\circ h_\gamma,\qquad
		h_{\gamma^{-1}}=h_\gamma^{-1}.
	\]
	改变基点处横截面，用两截面间沿叶片的局部对应识别后，环路产生的芽按共轭变换。
\end{proposition}
\begin{proof}
	先看一张小坐标盒。在适应坐标中，叶片是 \(y=\text{常数}\)，各横截面经投影到 \(y\) 坐标都是局部微分同胚。因此两横截面间的对应是“投影到 \(y\)，再用另一截面的投影逆”，只由两截面决定。插入一张中间截面时，两个投影互相消去，复合不变。

	在两张适应图交叠处，第二张图的横向坐标对第一张图的叶方向导数为零，因为两图都将这些方向识别为 \(D\)。缩小到交叠内的乘积小盒后，它因此只依赖第一张图的横向坐标；这个横向变化的微分可逆，因为整个换坐标微分可逆且有相应的块三角形式。故同一小段道路用任一张图计算，得到相同的横截面对应芽。

	对一条固定道路，将两种有限分割取公共加细，再把每段缩小到上述交叠小盒中。插入中间截面和改变坐标均不改变对应芽，有限次复合仍相同。这证明了定义与选择无关。拼接和逆向的公式也由依次复合与逐段求逆得到；所有比较只需在有限个代表映射的共同小定义域内进行。

	再设 \(H:[0,1]^2\to L\) 是固定端点的叶内同伦。这里连续性使用叶自身的拓扑。每个点都有叶中开邻域，其像落在一张叶片内；这些邻域的原像覆盖紧参数方形。有限子覆盖和 Lebesgue 数给出足够细的网格，使每个小格的像都落在同一张叶片内。为说明所用紧性步骤，可对有限开覆盖选每点一个仍留在某成员中的小球，取有限个半径减半的球覆盖方形；这些半径的正最小值就给出一个可用的细分尺度。

	每个小格边界的返回映射芽 是恒等芽，因为在同一张叶片内，沿边的所有横截面对应都由同一个横向坐标决定，绕一圈后投影及其逆全部抵消。把网格逐格拼起来，公共边的两个反向对应也互相抵消，得到方形外边界的返回映射芽为恒等芽。两条竖边因端点固定而给出恒等芽，所以上下两条道路的横截面映射芽相同。最后，若横截面通过局部对应 \(c\) 改变，沿环路前后分别插入 \(c^{-1}\) 与 \(c\)，所得芽为 \(c\circ h_\gamma\circ c^{-1}\)，即共轭。
\end{proof}

\begin{example}[收缩的横向返回]
	\label{g1:exm:contracting-leaf-holonomy}
	在圆柱 \((\mathbb R/\mathbb Z)\times\mathbb R\) 上取
	\[
		X=\partial_\theta-a r\partial_r,\qquad a>0.
	\]
	其整体流为 \(\Phi_t(\theta,r)=(\theta+t,e^{-at}r)\)。中间圆周 \(r=0\) 是闭叶，截面 \(\theta=0\) 的返回映射为 \(P(r)=e^{-a}r\)。沿闭叶正向一圈的横截面返回映射芽因而是这个收缩映射在零点的芽；绕 \(k\) 圈得到 \(r\mapsto e^{-ak}r\)。它不是恒等芽，反映的是邻近叶的横向收缩。

	相比之下，无理斜率环面的每条叶自身是 \(\mathbb R\)，所有叶内环路均可缩，所以由上一命题，横截面返回映射芽群平凡。截面上的无理旋转返回把一个交点送到另一个交点，并非沿该叶环路返回同一个叶内基点。这说明返回运动复杂与沿环路的横截面返回映射芽非平凡不是同一件事。
\end{example}

\subsection{非 Hausdorff 叶空间}
\label{g1:subsec:1305:03}

把每条叶压成一点后，叶的聚集行为就转化为商空间的邻域关系。开集能否将不同叶分开，将决定这个商能否作为通常意义下的流形继续研究。

\begin{definition}
	\label{g1:def:leaf-space}
	设 \(\mathcal F\) 为光滑流形 \(M\) 上的光滑叶状结构。将同一叶中的所有点识别，得到的集合记为 \(M/\mathcal F\)，称为\term[yekongjian]{叶空间}[leaf space]，以 \(q:M\to M/\mathcal F\) 表示将每点送到其所在叶的映射，并赋予商拓扑。即 \(V\subset M/\mathcal F\) 开，当且仅当\(q^{-1}(V)\) 在 \(M\) 中开；这样的逆像是若干整条叶的并，也称饱和开集。
\end{definition}

这个拓扑描述的是能否用开区域区分整条叶，与上一节每条叶自身的流形拓扑不同。在有理斜率环面上，叶空间确实是圆周。令 \(\alpha=p/q\) 为既约分数，定义
\[
	f:T^2\longrightarrow\mathbb R/\mathbb Z,\qquad f([x,y])=[qy-px].
\]
它良定义、满射，且沿 \(X_\alpha\) 的导数为零。若两点的 \(f\) 值相同，取实代表后有 \(q\Delta y-p\Delta x=k\in\mathbb Z\)。互素性给出整数 \(m,n\) 满足 \(qn-pm=k\)，从而 \(q(\Delta y-n)=p(\Delta x-m)\)。这恰好说明两点在同一条斜率 \(p/q\) 的轨道上。因此纤维就是叶，\(f\) 诱导叶空间到圆周的连续双射。叶空间是紧环面的连续像，故紧；圆周 Hausdorff，这个连续双射就是同胚。

\begin{proposition}[无理斜率的叶空间]
	\label{g1:prop:irrational-leaf-quotient}
	设 \(T^2=\mathbb R^2/\mathbb Z^2\)，\(\alpha\in\mathbb R\setminus\mathbb Q\)，\(\mathcal F_\alpha\) 为分布 \(\operatorname{span}\{\partial_x+\alpha\partial_y\}\) 的积分叶所成的叶状结构。叶空间 \(T^2/\mathcal F_\alpha\)不止一个点，但只有空集和全空间两个开集。特别地，它不是 Hausdorff 空间。
\end{proposition}
\begin{proof}
	每条叶与横截面 \(x=0\) 的交点集为一个整数旋转轨道，至多可数；圆周上的点不可数，所以不止一条叶。若 \(U\subset T^2\) 是非空饱和开集，每条叶都稠密，因此每条叶都与 \(U\) 相交。饱和性又要求整条相交的叶包含于 \(U\)，所以 \(U=T^2\)。按商拓扑的定义，这恰好给出所述两个开集。不同叶所对应的两点没有不交邻域，故非 Hausdorff。
\end{proof}

叶空间的分离失败并不只发生在稠密叶中。下面所有叶都是闭的嵌入子流形，叶空间仍然无法分离某两个点。

\begin{example}[分裂的零层]
	\label{g1:exm:punctured-plane-leaf-space}
	令 \(M=\mathbb R^2\setminus\{(0,0)\}\)，\(D=\mathbb R\partial_x\)。当 \(y\neq0\) 时，直线 \(\{(x,y):x\in\mathbb R\}\) 是一条叶；当 \(y=0\) 时，缺失的原点将该直线分成两条叶
	\[
		L_- =\{(x,0):x<0\},\qquad L_+=\{(x,0):x>0\}.
	\]
	这些叶都在 \(M\) 中闭，且以其通常参数为嵌入子流形。例如 \(L_+\) 的环境闭包可能增加的唯一有限点是原点，而原点已被删除。

	但 \(L_-\) 与 \(L_+\) 无法被不交的饱和开集分开。若饱和开集 \(U_-\) 包含 \(L_-\)，它包含 \((-1,0)\) 的一个小开矩形，因此对所有充分小的 \(y\neq0\)，\(U_-\) 都与高度 \(y\) 的水平叶相交，进而包含整条该叶。包含 \(L_+\) 的饱和开集 \(U_+\) 同样包含所有充分小非零高度的整条水平叶。所以 \(U_-\cap U_+\neq\varnothing\)。叶空间因此非 Hausdorff，尽管每个叶类是闭集、商空间的每个单点也是闭集。
\end{example}

局部横向坐标仍然存在；问题是绕过缺口继续这些坐标时，零高度对应了两条不同叶，而邻近非零高度各只对应一条叶。构造一个整体流形作为叶空间，需要超出局部可积性的额外条件。

\subsection{标准接触分布的不可积性}
\label{g1:subsec:1305:04}

前面的例子都有积分叶，只是叶的整体组织方式不同。现在看另一种障碍：一族光滑平面甚至不能在任何一点附近拼成与它们相切的曲面。在 \(\mathbb R^3\) 中定义
\[
	\alpha=\mathrm{d}z-x\,\mathrm{d}y,\qquad
	D=\ker\alpha
	=\operatorname{span}\{A,B\},\qquad
	A=\partial_x,\quad B=\partial_y+x\partial_z.
\]
因为 \(\alpha(\partial_z)=1\)，这确为常秩二的光滑分布。\(A\) 允许沿 \(x\) 方向移动，\(B\) 允许沿 \(y\) 方向移动，但后者同时按当前位置 \(x\) 改变高度。

\begin{proposition}[标准接触平面的不可积性]
	\label{g1:prop:contact-model-nonintegrable}
	在 \(\mathbb R^3\) 的标准坐标 \((x,y,z)\) 中，令 \(\alpha=\mathrm{d}z-x\,\mathrm{d}y\)，\(A=\partial_x\)，\(B=\partial_y+x\partial_z\)，\(D=\ker\alpha=\operatorname{span}\{A,B\}\)。分布 \(D\) 在每一点都不对合，因而不存在任何二维积分流形。具体地，\([A,B]=\partial_z\notin D\)。
\end{proposition}
\begin{proof}
	由括号坐标公式，只有 \(B\) 的 \(\partial_z\) 分量 \(x\) 沿 \(A=\partial_x\) 求导产生非零项，故 \([A,B]=\partial_z\)。而 \(\alpha(\partial_z)=1\)，所以该向量不属于核。若存在二维积分流形，则其上任一点都必须满足\cref{g1:prop:involutivity-necessary}的必要条件，与此矛盾。

	还可直接看出曲面方程的矛盾。若有积分曲面，投影 \((x,y,z)\mapsto(x,y)\) 在其切空间上可逆，因为它把 \(A,B\) 送为平面坐标基。逆函数定理遂将曲面局部写成 \(z=f(x,y)\)。切向条件 \(\alpha|_{\mathrm{T}S}=0\) 要求 \(f_x=0\)、\(f_y=x\)，从而 \(f_{xy}=0\) 而 \(f_{yx}=1\)，违反光滑函数混合偏导相等。这把括号障碍解释成了相切方程的不相容。
\end{proof}

为什么这个模型称为接触分布？在三维中，可以直接用括号给出不依赖外微分的定义。设 \(D\) 为秩二分布，对 \(u,v\in D_p\)，选局部 \(D\) 截面 \(U,V\) 延拓它们，令
\[
	\mathcal B_p(u,v)=[U,V]_p\pmod{D_p}\in \mathrm{T}_pM/D_p.
\]
这个值不依赖延拓。若 \(U\) 改变为 \(U+\sum_i f_iE_i\)，其中 \(E_i\) 为 \(D\) 的局部标架且 \(f_i(p)=0\)，则括号的变化在 \(p\) 为 \(-\sum_iV(f_i)(p)E_i(p)\)，属于 \(D_p\)；改变 \(V\) 同理。因此 \(\mathcal B_p\) 是良定义的交替双线性映射。若它在每点非零，称 \(D\) 为三维流形上的\term[jiechufenbu]{接触分布}[contact distribution]。因 \(D_p\) 为二维、商空间为一维，非零也等价于这个交替配对非退化。上述模型中 \(\mathcal B_p(A_p,B_p)=[\partial_z]\neq0\)，所以处处满足定义。

\ProseParagraphBreak

接触条件比“某处不对合”更强，它要求在每点都有这种非退化的横向括号效应。反过来，对合条件正是 \(\mathcal B_p=0\) 对所有 \(p\) 成立。这一比较把可积与接触两端放在同一个逐点量中：前者的允许方向在交换后不产生新方向，后者的两个允许方向交换后能产生缺少的第三个方向。\chapref{g1:ch:14}建立外微分后，三维接触条件还可写成 \(\alpha\wedge\mathrm{d}\alpha\neq0\)。

% 纲要条目排为编号小节标题，并进入目录与书签；说明段仍保留于纲要注释中。


\begin{chaptersummary}
	局部初值问题由积分方程的压缩迭代求解，唯一性使坐标解可以粘合为最大积分曲线。光滑依赖先由统一余项估计证明一阶导数，再用变分系统逐阶提升；最大流的定义域因而是时间与初值空间中的开集。自主流满足局部时间加法律，时间依赖方程则使用记录起始与到达时刻的演化映射 \(\Phi_{t,s}\)。

	\ProseParagraphBreak

	有限时间不能继续的最大轨道，最终会离开每个紧集；紧支向量场因此完备。非零向量场在局部流盒中成为一条坐标轴上的平移，零点则构成直接障碍。截断可以保留指定紧集附近的局部向量场并获得完整的整体流，但流的相同性只在轨道仍处于保留区域时成立。带边界流形上的双侧流还要求向量场切于边界。

	\ProseParagraphBreak

	Lie 括号 \([X,Y]=XY-YX\) 是导子交换子，其为零等价于充分小共同时间内的局部流交换。Lie 导数定义为沿流拉回后的变化率，在向量场上为 Lie 括号，在混合张量上与张量积及收缩相容。向量场的相关性保证括号和协变张量 Lie 导数的自然性；混合张量拉回还需要微分可逆。插入与 Lie 导数的分次恒等式为下一章的外微分运算准备符号。

	\ProseParagraphBreak

	对合性使归一化的分布标架两两交换，交换流随即构造适应坐标，这给出 Frobenius 定理的充分性；括号自然性给出必要性。局部叶片沿有限链粘合为最大连通叶，叶片的可数传播保证叶自身第二可数。叶可能在环境流形中稠密，叶空间可能不满足 Hausdorff 分离，沿叶环路的横截面返回映射芽也可能非平凡。接触模型 \(\ker(\mathrm{d}z-x\,\mathrm{d}y)\) 则因括号离开原分布而不可积。
\end{chaptersummary}

\BookExercises

\begin{exercise}\label{g1:ex:ode-regularity-counterexamples}
	在实直线上分别考虑 \(x'=\sqrt{|x|}\) 与 \(x'=|x|\)。
	\begin{enumerate}
		\item 对第一个方程，构造无穷多条满足 \(x(0)=0\) 的 \(C^1\) 解，说明仅要求右端连续不足以保证唯一性。
		\item 对第二个方程，求全部初值解，直接证明唯一性，并判断固定 \(t\ne0\) 时解对初值在零点是否可微。
	\end{enumerate}
\end{exercise}
\begin{solution}
	1. 对任意 \(a\geq0\)，令
	\[
		x_a(t)=\begin{cases}0,&t\leq a,\\(t-a)^2/4,&t>a.\end{cases}
	\]
	在拼接点 \(a\) 两侧，函数值与一阶导数都趋于零，所以 \(x_a\) 为 \(C^1\)。当 \(t>a\) 时，\(x_a'(t)=(t-a)/2=\sqrt{|x_a(t)|}\)；其余时间两边都为零。因此它在整个实轴上解方程，且满足初值。不同 \(a\) 给出不同解，常值零解也满足条件。特别地，\(a=0\) 的解与零解在任意零点邻域内不同，所以局部唯一性已经失败。

	2. 对初值 \(x(0)=p\)，直接验证
	\[
		\Phi_t(p)=\begin{cases}p e^t,&p\geq0,\\p e^{-t},&p<0\end{cases}
	\]
	满足方程，且对全部实数时间有定义。虽然 \(|x|\) 在零处不可微，但它满足 \(\bigl||x|-|y|\bigr|\leq|x-y|\)。若 \(u,v\) 是同初值解，在任一正向公共紧时间段内，由积分方程得
	\[
		|u(t)-v(t)|\leq\int_0^t|u(s)-v(s)|\,\mathrm{d}s.
	\]
	\cref{g1:lem:gronwall-estimate}取 \(a=b=0\)、\(L=1\)，给出两解相同；负时间改用反向参数即可。故初值解唯一。

	但当 \(t\ne0\) 固定时，\(p\mapsto\Phi_t(p)\) 在零点的右导数为 \(e^t\)，左导数为 \(e^{-t}\)，二者不同，所以不可微。这个例子将两个问题分开：局部 Lipschitz 的误差估计足以给出唯一性，初值的光滑依赖还需要右端相应的光滑性。它们不能由右端连续这一条件同时得到。
\end{solution}

\begin{exercise}\label{g1:ex:parameter-quadratic-flow}
	对实参数 \(\lambda\)，考虑 \(x'=\lambda x^2\)、\(x(0)=y\)。求最大解及其关于 \((t,y,\lambda)\) 的定义域，计算 \(P=\partial_y x\)、\(Q=\partial_\lambda x\)，并验证它们的变分方程。特别说明公式在 \(\lambda=0\) 或 \(y=0\) 附近的光滑性，以及 \(y=0\) 时 \(P\) 是否为零。
\end{exercise}
\begin{solution}
	直接代入可得
	\[
		x(t;y,\lambda)=\frac{y}{1-\lambda yt}.
	\]
	若 \(\lambda y>0\)，最大区间为 \((-\infty,1/(\lambda y))\)；若 \(\lambda y<0\)，为 \((1/(\lambda y),\infty)\)；若 \(\lambda y=0\)，为整个实轴。当 \(\lambda y\ne0\) 时，分母在所述有限端点为零，且 \(y\ne0\)，所以函数绝对值趋于无穷，不能延拓。只取包含初始时刻零的分母非零连通分支，联合定义域因此为开集
	\[
		\mathcal D=\{(t,y,\lambda)\in\mathbb R^3:1-\lambda yt>0\}.
	\]
	该域上的分式联合光滑，包括 \(\lambda=0\) 和 \(y=0\) 的所有允许点；此处不需要除以 \(\lambda\) 或 \(y\)。

	令 \(d=1-\lambda yt\)，普通求导给出
	\[
		P=d^{-2},\qquad Q=t y^2d^{-2}.
	\]
	这里 \(d\) 只是分母的缩写，并非微分符号。对时间求导，得到
	\[
		P'=2\lambda y d^{-3}=2\lambda xP,
	\]
	以及
	\[
		Q'=y^2d^{-2}+2\lambda t y^3d^{-3}
		=x^2+2\lambda xQ.
	\]
	同时 \(P(0)=1\)、\(Q(0)=0\)，正好是\eqref{g1:eq:ode-variational-equation}应用于 \(F(x,\lambda)=\lambda x^2\) 的结果。

	当 \(y=0\) 时，轨道为常值零，但 \(P=1\)，并不为零。常值轨道只说明沿这条轨道的时间导数为零；改变初始位置会换到另一条轨道，其一阶变化仍可非零。此时 \(Q=0\)，因为保持初值为零而只改变 \(\lambda\)，仍得到同一条零轨道。
\end{solution}

\begin{exercise}\label{g1:ex:saddle-flow-open-restriction}
	在 \(\mathbb R^2\) 上取 \(X=x\partial_x-y\partial_y\)。求其最大流，并证明 \(xy\) 沿轨道不变。再把向量场限制到开子流形 \(U=\{(x,y):x+y>0\}\)，求每个初值的最大区间，解释为何原来对所有时间存在的流限制后可能在有限时间终止。
\end{exercise}
\begin{solution}
	两个标量方程分别为 \(x'=x\)、\(y'=-y\)，所以
	\[
		\Phi_t(x,y)=(e^t x,e^{-t}y),\qquad t\in\mathbb R.
	\]
	公式满足初值并在全部时间上有意义，唯一性说明它就是最大流。沿轨道的坐标乘积为 \((e^t x)(e^{-t}y)=xy\)，也可由 \((xy)'=x'y+xy'=xy-xy=0\) 验证不变性。

	对初值 \((x,y)\in U\)，限制后的轨道只能保留在 \(U\) 中的那段时间。条件为
	\[
		e^t x+e^{-t}y>0\quad\Longleftrightarrow\quad e^{2t}x+y>0.
	\]
	右侧作为时间条件给出一个包含零的开区间，故无须再从多个分支中选择。具体地，
	\[
		I_{(x,y)}^U=
		\begin{cases}
			\mathbb R,&x\geq0,\ y\geq0,\ x+y>0,\\
			\bigl(\tfrac12\log(-y/x),\infty\bigr),&x>0,\ y<0,\ x+y>0,\\
			\bigl(-\infty,\tfrac12\log(y/(-x))\bigr),&x<0,\ y>0,\ x+y>0.
		\end{cases}
	\]
	这些情形覆盖 \(U\)，因为其中不可能两个坐标都非正。第二、三种情形的有限端点处，环境轨道趋向直线 \(x+y=0\) 上的非零点，而该点不属于 \(U\)。若存在值仍在 \(U\) 中的延拓，复合包含映射后它也会成为平面方程的延拓；平面解的唯一性迫使端点值等于上述不在 \(U\) 中的点，矛盾。因此所列区间最大。这里终止并非速度变得不光滑，也非环境轨道爆破，而是底空间删去了它将到达的位置。
\end{solution}

\begin{exercise}\label{g1:ex:affine-time-dependent-evolution}
	设 \(J\subset\mathbb R\) 是开区间，\(a,b\in C^\infty(J)\)。求方程 \(x'=a(t)x+b(t)\) 从时间 \(s\) 到时间 \(t\) 的演化映射，证明它对所有 \(s,t\in J\) 及全部实初值都有定义，并验证二参数合成公式及初始时刻导数公式。
\end{exercise}
\begin{solution}
	记
	\[
		E(t,s)=\exp\left(\int_s^t a(r)\,\mathrm{d}r\right).
	\]
	它始终为正，且 \(E(t,r)E(r,s)=E(t,s)\)。所求映射为
	\[
		\Phi_{t,s}(x)=E(t,s)x+\int_s^t E(t,r)b(r)\,\mathrm{d}r.
	\]
	对 \(t\) 求导，利用 \(\partial_tE(t,s)=a(t)E(t,s)\)，得
	\[
		\partial_t\Phi_{t,s}(x)=a(t)\Phi_{t,s}(x)+b(t),
	\]
	而 \(\Phi_{s,s}(x)=x\)，所以这确实是初值解。因为 \(J\) 是区间，每两个时刻间的闭线段都包含于 \(J\)，相应光滑函数的积分有限；公式对全部 \(s,t\in J\) 有定义。所有联合光滑性也可从积分换元到固定区间 \([0,1]\) 后的公式看出。

	对任意 \(s,r,t\in J\)，代入复合表达得
	\[
		\begin{aligned}
		\Phi_{t,r}(\Phi_{r,s}(x))
		&=E(t,r)E(r,s)x+\int_s^r E(t,r)E(r,v)b(v)\,\mathrm{d}v
		+\int_r^t E(t,v)b(v)\,\mathrm{d}v\\
		&=E(t,s)x+\int_s^t E(t,v)b(v)\,\mathrm{d}v
		=\Phi_{t,s}(x).
		\end{aligned}
	\]
	定向积分的可加性保证这组等式不要求三个时刻按大小排列。取 \(t=s\) 便得到 \(\Phi_{s,r}\Phi_{r,s}=\operatorname{id}\)，所以各演化映射的逆为反向演化。

	最后 \(\partial_sE(t,s)=-a(s)E(t,s)\)，而积分项对下限 \(s\) 的导数为 \(-E(t,s)b(s)\)。因此
	\[
		\partial_s\Phi_{t,s}(x)
		=-E(t,s)\bigl(a(s)x+b(s)\bigr).
	\]
	又有 \(\partial_x\Phi_{t,s}(x)=E(t,s)\)，故该式恰为
	\[
		\partial_s\Phi_{t,s}(x)
		=-\mathrm{d}(\Phi_{t,s})_x\bigl(X_s(x)\bigr),\qquad
		X_s(x)=\bigl(a(s)x+b(s)\bigr)\partial_x.
	\]
	若把原问题的时间范围固定为 \(J\)，每条解的最大绝对时间区间就是 \(J\)；是否能把系数和方程继续定义到 \(J\) 外，是另一个延拓问题。
\end{solution}

\begin{exercise}\label{g1:ex:linear-growth-completeness}
	设 \(X\) 是 \(\mathbb R^n\) 上的光滑向量场，\(n\geq1\)。若存在常数 \(A,B\geq0\)，使
	\[
		|X(x)|\leq A+B|x|\qquad(x\in\mathbb R^n),
	\]
	证明 \(X\) 完备。证明中须给出轨道在任意有限时间内的明确范数界，并分别处理正向和反向时间。
\end{exercise}
\begin{solution}
	设初值为 \(p\) 的最大积分曲线为 \(\gamma:(a,b)\to\mathbb R^n\)，其中 \(a<0<b\)。对 \(0\leq t<b\)，由积分方程及增长条件得到
	\[
		|\gamma(t)|\leq|p|+\int_0^t|X(\gamma(s))|\,\mathrm{d}s
		\leq|p|+At+B\int_0^t|\gamma(s)|\,\mathrm{d}s.
	\]
	在任意包含于 \([0,b)\) 的紧时间区间上应用\cref{g1:lem:gronwall-estimate}，得
	\[
		|\gamma(t)|\leq
		\begin{cases}
			|p|e^{Bt}+\dfrac AB(e^{Bt}-1),&B>0,\\
			|p|+At,&B=0.
		\end{cases}
	\]
	因为所取紧时间区间可以包含任意指定的 \(t<b\)，这个估计对全部 \(0\leq t<b\) 成立。

	若 \(b<\infty\)，把右端的 \(t\) 换成 \(b\)，便得到一个与 \(t\) 无关的有限上界 \(R\)。于是 \(\gamma([0,b))\) 留在 Euclid 闭球 \(\overline B_R(0)\) 内。该球紧，违反\cref{g1:thm:flow-escape-criterion}关于有限最大右端点必须逃离每个紧集的结论。因此 \(b=+\infty\)。

	对负向时间，令 \(\eta(u)=\gamma(-u)\)。它满足 \(\eta'=-X(\eta)\)、\(\eta(0)=p\)，且 \(-X\) 仍满足同一增长界。若 \(a> -\infty\)，则 \(\eta\) 的最大右端点为 \(-a<\infty\)：否则将其延拓反向就会延拓 \(\gamma\)。把前两段用于 \(-X\)，排除这个有限端点，得到 \(a=-\infty\)。正负时间的估计合并为
	\[
		|\gamma(t)|\leq
		\begin{cases}
			|p|e^{B|t|}+\dfrac AB(e^{B|t|}-1),&B>0,\\
			|p|+A|t|,&B=0,
		\end{cases}
		\qquad t\in\mathbb R.
	\]
	每个初值的最大区间因而都是 \(\mathbb R\)，故 \(X\) 完备。这里 Euclid 空间的闭有界球紧，使范数估计能直接用于逃逸判据；若只在一个真开子集上给出同样的数值界，还须检查轨道是否会接近该开子集缺失的点。
\end{solution}

\begin{exercise}\label{g1:ex:sum-of-complete-fields}
	在 \(\mathbb R^2\) 上令
	\[
		X=y^2\partial_x,\qquad Y=x^2\partial_y.
	\]
	分别求 \(X,Y\) 的整体流，再证明 \(X+Y\) 不完备。对后一个结论，求出初值 \((1,1)\) 的最大积分曲线及最大区间。
\end{exercise}
\begin{solution}
	沿 \(X\) 的积分曲线，\(y'=0\)，所以 \(y(t)=y_0\)；随后 \(x'=y_0^2\) 给出
	\[
		\Phi_t^X(x_0,y_0)=(x_0+t y_0^2,y_0).
	\]
	同理，沿 \(Y\) 的第一坐标保持不变，故
	\[
		\Phi_t^Y(x_0,y_0)=(x_0,y_0+t x_0^2).
	\]
	两个公式对全部初值和全部实数时间都有意义，并满足各自的方程与初值。因此它们就是最大流，\(X,Y\) 都完备。

	和向量场的方程则为
	\[
		x'=y^2,\qquad y'=x^2.
	\]
	曲线
	\[
		\gamma(t)=\left(\frac1{1-t},\frac1{1-t}\right),\qquad t\in(-\infty,1),
	\]
	直接满足这组方程及初值 \(\gamma(0)=(1,1)\)。由\cref{g1:thm:maximal-integral-curve}的唯一性，它是该初值最大解的限制。其定义域已向左延至负无穷；当 \(t\uparrow1\) 时两个坐标都趋于正无穷，不能以 \(\mathbb R^2\) 中的点为值连续延拓，所以最大右端点正是 \(1\)。故这条曲线的最大区间为 \((-\infty,1)\)，\(X+Y\) 不完备。

	在两个单独方程中，决定速度的另一坐标始终固定；相加以后，两个坐标会同时增大，并进一步增大彼此的速度。这个例子说明完备向量场一般不对加法封闭，不能仅凭两个独立运动都能在全部实数时刻进行，就断定由它们的速度之和决定的运动也能如此。
\end{solution}

\begin{exercise}\label{g1:ex:global-flow-box-periodic-obstruction}
	在 \(\mathbb R^2\) 上取 \(X=\partial_x+y\partial_y\)。证明
	\[
		s=x,\qquad z=ye^{-x}
	\]
	给出一组全局坐标，并将 \(X\) 化为 \(\partial_s\)。再对圆环
	\[
		\mathcal A=\{(x,y):1<x^2+y^2<4\}
	\]
	上的转动场 \(R=-y\partial_x+x\partial_y\)，证明不存在满足 \(R(s)=1\) 的全局光滑实函数 \(s:\mathcal A\to\mathbb R\)。说明这与流盒定理的局部结论为何相容。
\end{exercise}
\begin{solution}
	映射 \(H:\mathbb R^2\to\mathbb R^2\)，\(H(x,y)=(x,ye^{-x})\)，有明确的逆
	\[
		H^{-1}(s,z)=(s,e^s z).
	\]
	两个映射都光滑，所以它们是全局微分同胚，\((s,z)\) 确为全局坐标。直接计算得
	\[
		X(s)=1,\qquad
		X(z)=\partial_x(ye^{-x})+y\partial_y(ye^{-x})=-ye^{-x}+ye^{-x}=0.
	\]
	由微分的坐标表达，\(H_*X=\partial_s\)。原坐标中的方程也可直接积分为
	\[
		\Phi_t^X(x,y)=(x+t,e^t y).
	\]
	因此 \(H(\Phi_t^X(x,y))=(s+t,z)\)，新坐标中的运动在所有时间上就是平移。

	圆环上的转动场具有流
	\[
		\Phi_t^R(x,y)=(x\cos t-y\sin t,x\sin t+y\cos t).
	\]
	它保持 \(x^2+y^2\)，因而始终留在圆环中，且 \(\Phi_{2\pi}^R(p)=p\)。若存在题设全局实函数 \(s\)，对任意 \(p\in\mathcal A\) 沿轨道使用链式法则，得
	\[
		\frac{\mathrm{d}}{\mathrm{d}t}s(\Phi_t^R(p))
		=R(s)(\Phi_t^R(p))=1.
	\]
	从零积分到 \(2\pi\)，得到 \(s(p)-s(p)=2\pi\)，矛盾。所以这样的全局实值时间坐标不存在。

	\(R\) 在圆环上处处非零，故\cref{g1:thm:flow-box}仍保证每个点附近有实值局部时间坐标。沿一个足够短的角度分支可以令 \(s\) 等于角度，但绕一整圈以后实角增加 \(2\pi\)，无法成为单值实函数。流盒只承诺局部坐标，不排除轨道离开该坐标域后再次返回同一点。
\end{solution}

\begin{exercise}\label{g1:ex:positive-time-change-completeness}
	设 \(M\) 是有限维 Hausdorff、第二可数的无边界光滑流形，\(X\in\mathfrak X(M)\)，\(f:M\to(0,\infty)\) 光滑，令 \(Y=fX\)。
	\begin{enumerate}
		\item 对 \(X\) 的最大积分曲线 \(\gamma:I\to M\)，\(\gamma(0)=p\)，用积分给出相应 \(Y\) 轨道的换时公式，并证明两条最大轨道具有相同的像及相同运动方向。
		\item 取 \(M=\mathbb R\)、\(X=\partial_x\)、\(f(x)=1+x^2\)，求 \(Y\) 的最大积分曲线，说明正光滑倍数可以破坏完备性。
		\item 若 \(M\) 紧，证明 \(X,Y\) 都完备，并用换时公式解释此时为何不会发生有限时间压缩到无限原时间的现象。
	\end{enumerate}
\end{exercise}
\begin{solution}
	1. 记 \(I=(a,b)\)，其中 \(a<0<b\)。以 \(t\) 表示 \(X\) 的时间，以 \(s\) 表示 \(Y\) 的时间，定义
	\[
		s(t)=\int_0^t\frac{\mathrm{d}u}{f(\gamma(u))}.
	\]
	因为 \(f>0\)，有 \(s'(t)=1/f(\gamma(t))>0\)，所以 \(s\) 严格递增。其像为开区间
	\[
		J=(A,B),\qquad
		A=\lim_{t\downarrow a}s(t),\quad B=\lim_{t\uparrow b}s(t),
	\]
	其中端点允许为无穷。\(s:I\to J\) 双射，导数处处非零；一维逆函数定理使其逆 \(t=t(s)\) 光滑。令
	\[
		\eta(s)=\gamma(t(s)),\qquad s\in J.
	\]
	链式法则与逆函数导数给出
	\[
		\eta'(s)=X_{\gamma(t(s))}\,t'(s)
		=f(\eta(s))X_{\eta(s)}=Y_{\eta(s)},\qquad \eta(0)=p.
	\]
	因此它是 \(Y\) 的积分曲线，并且因为 \(t:J\to I\) 严格递增且满射，\(\eta(J)=\gamma(I)\)，方向也保持不变。

	还须证明 \(\eta\) 最大，才能比较两条最大轨道。若它可延拓为较大开区间 \(\widehat J\supsetneq J\) 上的解 \(\widehat\eta\)，定义反向换时
	\[
		\widehat t(s)=\int_0^s f(\widehat\eta(v))\,\mathrm{d}v.
	\]
	该函数仍严格递增且有光滑逆。在 \(J\) 上，它与原来的 \(t(s)\) 有相同导数和零处值，所以相同，进而 \(\widehat t(J)=I\)。由于 \(\widehat J\) 严格扩大，严格递增映射的像 \(\widehat t(\widehat J)\) 也严格包含 \(I\)。用 \(\widehat t\) 的逆重新参数化 \(\widehat\eta\)，链式法则得到 \(X\) 的一条解，并在 \(I\) 上等于 \(\gamma\)，这与 \(\gamma\) 最大矛盾。因此 \(J\) 正是 \(Y\) 的最大时间区间。若 \(X_p=0\)，两条曲线均为常值，公式退化为 \(s=t/f(p)\)，同样成立。

	2. \(X\) 的积分曲线为 \(\gamma(t)=x_0+t\)，对全部实数时间存在。换时公式给出
	\[
		s(t)=\int_0^t\frac{\mathrm{d}u}{1+(x_0+u)^2}
		=\arctan(x_0+t)-\arctan x_0.
	\]
	令 \(\theta_0=\arctan x_0\in(-\pi/2,\pi/2)\)，则
	\[
		\eta(s)=\tan(s+\theta_0),\qquad
		J=(-\pi/2-\theta_0,\pi/2-\theta_0).
	\]
	在两个有限端点，轨道分别趋于负、正无穷，因此不能延拓。其像仍是整个实轴，与 \(X\) 的轨道像相同；变化的是沿这条轨道走到无穷远所需的时间。尤其初值 \(x_0=0\) 时得到 \(\eta(s)=\tan s\)，而 \(X\) 完备、\(Y=(1+x^2)\partial_x\) 不完备。

	3. 若 \(M\) 紧且无边界，\cref{g1:thm:compact-support-complete}直接说明两个光滑向量场都完备。也可以在先知道 \(X\) 完备后，从换时公式恢复 \(Y\) 的结论。紧性和正性给出常数 \(0<c\leq C<\infty\)，使 \(c\leq f\leq C\)。因此
	\[
		\frac1C\leq s'(t)\leq\frac1c.
	\]
	对 \(t\geq0\) 积分得到 \(t/C\leq s(t)\leq t/c\)，对 \(t\leq0\) 则得到 \(t/c\leq s(t)\leq t/C\)。所以 \(t\to+\infty\) 时 \(s(t)\to+\infty\)，\(t\to-\infty\) 时 \(s(t)\to-\infty\)，其像为整个实轴，\(Y\) 也完备。这个上下界排除了上一问中速度倍数沿轨道无限增大、把无限原时间压成有限新时间的情况。
\end{solution}

\begin{exercise}
	\label{g1:exr:1303-rescaling-frame}
	在 \(\mathbb R^2\) 上令 \(X=\partial_x\)、\(Y=e^x\partial_y\)。求所有满足 \([X,fY]=0\) 的光滑函数 \(f\)，并说明哪些这样的 \(f\) 使 \(X,fY\) 仍处处构成标架。给出一个处处正的选择，使这两个向量场成为坐标向量场，并写出它们的流。
\end{exercise}
\begin{solution}
	首先 \([X,Y]=e^x\partial_y=Y\)。由括号的乘法公式，
	\[
		[X,fY]=f[X,Y]+X(f)Y=(f+\partial_xf)Y.
	\]
	因为 \(Y\) 处处不为零，所求条件等价于 \(\partial_xf+f=0\)。这又等价于 \(\partial_x(e^xf)=0\)，所以
	\[
		f(x,y)=e^{-x}h(y),
	\]
	其中 \(h\) 为任意光滑实函数。反向代入即能验证所有这些函数满足要求。它们使 \(X,fY\) 构成标架，当且仅当 \(h(y)\) 处处不为零；此时 \(fY=h(y)\partial_y\)，仍与 \(\partial_x\) 独立。

	取 \(h\equiv1\)，则 \(f=e^{-x}>0\)，且 \(fY=\partial_y\)。两场正是原坐标的坐标场，流分别为
	\[
		\Phi_t(x,y)=(x+t,y),\qquad \Psi_s(x,y)=(x,y+s).
	\]
	它们在所有时间交换。原来的 \(X,Y\) 虽已构成标架，却并不交换；乘以非恒定函数能够改变括号，所以不能把括号看成对函数双线性的逐点运算。
\end{solution}

\begin{exercise}
	\label{g1:exr:1303-affine-flows}
	在 \(\mathbb R\) 上令 \(X=x\partial_x\)、\(Y=\partial_x\)。求 Lie 括号、两场的流及 \((\Phi_t^X)^*Y\)。验证
	\[
		\Phi_t^X\circ\Phi_s^Y=\Phi_{e^ts}^Y\circ\Phi_t^X,
	\]
	并说明它与两个流交换有什么区别。用拉回的导数再次核对 Lie 括号的符号。
\end{exercise}
\begin{solution}
	由一维坐标公式，
	\[
		[X,Y]=(x\cdot0-1\cdot1)\partial_x=-Y.
	\]
	解两个积分曲线方程，分别得到
	\[
		\Phi_t^X(x)=e^tx,\qquad \Phi_s^Y(x)=x+s.
	\]
	二者都在所有实时间有定义。由于 \(\mathrm{d}\Phi_t^X\) 是乘以 \(e^t\)，而 \(Y\) 的目标分量恒为 \(1\)，故
	\[
		(\Phi_t^X)^*Y=e^{-t}Y.
	\]
	这说明膨胀后的单位平移向量被带回原切空间时，要缩小 \(e^{-t}\) 倍。

	复合计算给出
	\[
		\Phi_t^X(\Phi_s^Y(x))=e^tx+e^ts
		=\Phi_{e^ts}^Y(\Phi_t^X(x)).
	\]
	交换次序时，平移的时间参数同时变成了 \(e^ts\)，因此这不是使用相同参数 \(s,t\) 的交换关系。实际上
	\[
		\Phi_s^Y(\Phi_t^X(x))=e^tx+s,
	\]
	当 \(s\ne0,t\ne0\) 时与前式不同。最后，
	\[
		\left.\frac{\mathrm{d}}{\mathrm{d}t}\right|_0(\Phi_t^X)^*Y
		=\left.\frac{\mathrm{d}}{\mathrm{d}t}\right|_0e^{-t}Y=-Y,
	\]
	与所算括号一致；若误把微分而非逆微分用于向量场拉回，就会得到相反符号。
\end{solution}

\begin{exercise}
	\label{g1:exr:1303-function-times-lie}
	设 \(M\) 为无边界光滑流形，\(f\in C^\infty(M)\)、\(X\in\mathfrak X(M)\)、\(\alpha\in\Omega^k(M)\)，其中 \(k\ge1\)。仅用协变张量 Lie 导数的求值公式、Lie 括号的乘法公式与楔积，证明
	\[
		\mathcal L_{fX}\alpha=f\mathcal L_X\alpha+\mathrm{d}f\wedge\iota_X\alpha.
	\]
	这里 \(\mathrm{d}f\) 只是函数的微分。另写出零次形式的相应结论，并用一个一维例子说明，Lie 导数一般不对下标中的向量场满足 \(C^\infty(M)\) 线性。
\end{exercise}
\begin{solution}
	对任意局部向量场 \(Y_1,\ldots,Y_k\)，有
	\[
		[fX,Y_a]=f[X,Y_a]-Y_a(f)X.
	\]
	将它代入\eqref{g1:eq:lie-covariant-evaluation}，并利用 \(\alpha\) 对函数多线性，得到
	\[
		\begin{split}
			(\mathcal L_{fX}\alpha)(Y_1,\ldots,Y_k)
			={}&f(\mathcal L_X\alpha)(Y_1,\ldots,Y_k)\\
			&+\sum_{a=1}^kY_a(f)\,
			\alpha(Y_1,\ldots,Y_{a-1},X,Y_{a+1},\ldots,Y_k).
		\end{split}
	\]
	最后一项中，把 \(X\) 从第 \(a\) 个位置移到第一位置，需要 \(a-1\) 次交换，因此该和等于
	\[
		\sum_{a=1}^k(-1)^{a-1}\mathrm{d}f(Y_a)
		(\iota_X\alpha)(Y_1,\ldots,\widehat{Y_a},\ldots,Y_k).
	\]
	按一次形式与 \(k-1\) 次形式的洗牌公式，这恰为 \((\mathrm{d}f\wedge\iota_X\alpha)(Y_1,\ldots,Y_k)\)。所有输入上取值相同，就证明了张量等式。若 \(\alpha=h\) 为函数，则直接由定义有 \(\mathcal L_{fX}h=fX(h)=f\mathcal L_Xh\)。

	在 \(\mathbb R\) 上取 \(X=\partial_x\)、\(f=x\)、\(\alpha=\mathrm{d}x\)。此时 \(\mathcal L_X\alpha=0\)，但 \(\iota_X\alpha=1\)，故
	\[
		\mathcal L_{x\partial_x}(\mathrm{d}x)=\mathrm{d}x
		\ne x\mathcal L_{\partial_x}(\mathrm{d}x)=0.
	\]
	附加项记录了速度缩放函数 \(f\) 在空间中的变化；它不能由在一点只乘上 \(f(p)\) 来替代。
\end{solution}

\begin{exercise}
	\label{g1:exr:1303-projection-related}
	设 \(P:\mathbb R^2\to\mathbb R\)、\(P(x,y)=x\)，并用 \(u\) 表示目标坐标。给定 \(X=a(x,y)\partial_x+b(x,y)\partial_y\)，证明存在与 \(X\) 为 \(P\)-相关的目标向量场，当且仅当 \(a(x,y)=a_0(x)\) 只依赖于 \(x\)，并说明目标场是否唯一。若 \(X\) 与 \(Y=c_0(x)\partial_x+d(x,y)\partial_y\) 都满足此条件，求其括号在目标上的相关场。最后对目标一形式 \(\alpha=h(u)\,\mathrm{d}u\) 直接计算 \(\mathcal L_X(P^*\alpha)\)，检验自然性公式，并用 \(X=y\partial_x\) 展示相关条件失效时发生什么。
\end{exercise}
\begin{solution}
	微分满足 \(\mathrm{d}P(\partial_x)=\partial_u\)、\(\mathrm{d}P(\partial_y)=0\)。若目标场为 \(\overline X=A(u)\partial_u\)，相关条件就是
	\[
		a(x,y)=A(x)\qquad\text{对所有 }(x,y).
	\]
	它等价于 \(a\) 只依赖 \(x\)。若条件成立，可取 \(a_0(x)=a(x,0)\)，这个函数光滑，且目标场必须为 \(\overline X=a_0(u)\partial_u\)。由于 \(P\) 满射，每个目标点的系数都被上述公式确定，所以唯一；竖直系数 \(b\) 则不影响相关条件。

	对另一个满足条件的场 \(Y\)，括号的水平系数是
	\[
		X(c_0)-Y(a_0)=a_0c_0'-c_0a_0'.
	\]
	所以它与 \((a_0c_0'-c_0a_0')\partial_u=[\overline X,\overline Y]\) 相关，竖直分量无须为零。

	不先假设相关条件，直接有 \(P^*\alpha=h(x)\,\mathrm{d}x\)，而
	\[
		\mathcal L_X(\mathrm{d}x)=(\partial_xa)\,\mathrm{d}x+(\partial_ya)\,\mathrm{d}y.
	\]
	故
	\[
		\mathcal L_X(P^*\alpha)
		=(a h'+h\,\partial_xa)\,\mathrm{d}x
		+h\,\partial_ya\,\mathrm{d}y.
	\]
	若 \(a=a_0(x)\)，第二项消失，第一项为 \((a_0h'+ha_0')\,\mathrm{d}x\)，恰等于 \(P^*(\mathcal L_{\overline X}\alpha)\)。若取 \(X=y\partial_x\) 与 \(\alpha=\mathrm{d}u\)，则
	\[
		\mathcal L_X(P^*\alpha)=\mathcal L_{y\partial_x}(\mathrm{d}x)=\mathrm{d}y.
	\]
	任何目标一形式的拉回都只含 \(\mathrm{d}x\) 项，不可能等于 \(\mathrm{d}y\)。此时 \(X\) 在同一投影纤维上的水平速度随 \(y\) 改变，没有可用的相关目标场。
\end{solution}

\begin{exercise}\label{g1:ex:frobenius-explicit-graph}
	在 \(\mathbb R^3\) 上取
	\[
		X=\partial_x+y\partial_z,\qquad
		Y=\partial_y+x\partial_z,\qquad
		D=\operatorname{span}\{X,Y\}.
	\]
	计算 Lie 括号和两条流，并从横截线 \((0,0,c)\) 出发构造适应坐标。求 \(D\) 的全部最大连通叶。
\end{exercise}
\begin{solution}
	两向量场的 \(x,y\) 分量分别为 \((1,0)\)、\((0,1)\)，故处处线性无关，\(D\) 为光滑秩二分布。括号的前两分量为零，最后分量为
	\[
		X(x)-Y(y)=1-1=0,
	\]
	所以 \([X,Y]=0\)。积分方程直接给出
	\[
		\Phi_X^t(x,y,z)=(x+t,y,z+yt),\qquad
		\Phi_Y^s(x,y,z)=(x,y+s,z+xs).
	\]
	两者先后复合的结果均为 \((x+t,y+s,z+yt+xs+ts)\)，故它们在所有时间交换。

	从横截线出发，得到
	\[
		\Psi(t,s,c)=\Phi_Y^s\Phi_X^t(0,0,c)=(t,s,c+ts).
	\]
	它是整体微分同胚，逆为
	\[
		(x,y,z)\longmapsto(t,s,c)=(x,y,z-xy).
	\]
	其时间方向微分分别为 \(X,Y\)，所以这些是适应坐标。固定 \(c\) 得到连通积分曲面
	\[
		L_c=\{(x,y,z):z=xy+c\}.
	\]
	函数 \(F=z-xy\) 满足 \(XF=YF=0\)。在任何连通积分流形上，\(F\) 的微分都为零，因而为常值；所以经过 \(L_c\) 中一点的连通积分流形只能落在 \(L_c\) 内。而 \(L_c\) 自身已是连通积分流形，故它就是该点的最大叶。不同 \(c\in\mathbb R\) 给出全部叶。
\end{solution}

\begin{exercise}\label{g1:ex:submersion-fiber-leaves}
	设 \(F:M^m\to N^n\) 为有限维、Hausdorff、第二可数无边界光滑流形之间的淹没。证明 \(D=\ker\mathrm{d}F\) 对合，且其最大连通叶恰为各纤维 \(F^{-1}(y)\) 的连通分支。具体讨论 \(M=\mathbb R^2\setminus\{(0,0)\}\)、\(F(x,y)=x\) 的全部叶，说明为何不能总把整个纤维称为一条叶。
\end{exercise}
\begin{solution}
	\cref{g1:exm:submersion-integral-levels}已说明 \(D\) 为秩 \(m-n\) 光滑分布，非空纤维为切空间等于 \(D\) 的嵌入子流形。也可以直接检查对合性：若 \(X,Y\) 取值于 \(D\)，对目标上的任意局部光滑函数 \(h\)，有
	\[
		X(h\circ F)=Y(h\circ F)=0,
		\qquad [X,Y](h\circ F)=0.
	\]
	在目标坐标函数上应用此式，就得到 \(\mathrm{d}F([X,Y])=0\)，故 \([X,Y]\) 取值于 \(D\)。

	取一条最大连通叶 \(L\)。沿叶有 \(\mathrm{d}(F|_L)=0\)，所以 \(F|_L\) 局部常值，进而由连通性为常值 \(y\)。包含映射连续，故其连通像落在 \(F^{-1}(y)\) 的某个连通分支 \(C\) 中。纤维是流形，局部坐标球连通，故其连通分支开；因此 \(C\) 本身也是连通积分流形。最大性给出 \(C\subseteq L\)，结合原来的包含便得 \(L=C\)，且流形结构由最大叶的唯一性一致。

	在所给平面例子中，\(\mathrm{d}F=\mathrm{d}x\) 处处满射，\(D=\mathbb R\partial_y\)。对 \(c\ne0\)，纤维 \(\{(c,y):y\in\mathbb R\}\) 连通，整条竖直线是一条叶。零纤维则为
	\[
		\{(0,y):y>0\}\ \cup\ \{(0,y):y<0\},
	\]
	它有两个连通分支，分别是一条叶。这里微分的秩没有发生变化，分成两条叶的原因是纤维不连通。
\end{solution}

\begin{exercise}\label{g1:ex:submersion-pullback-distribution}
	设 \(f:M^m\to N^n\) 为有限维、Hausdorff、第二可数无边界光滑流形之间的淹没，\(D\subseteq \mathrm{T}N\) 为秩 \(r\) 对合光滑分布。逐点定义
	\[
		E_p=\{v\in \mathrm{T}_pM:\mathrm{d}f_p(v)\in D_{f(p)}\}.
	\]
	证明 \(E\) 是秩 \(m-n+r\) 的对合光滑分布。再令 \(f:\mathbb R^2\to\mathbb R^2\)、\(f(x,y)=(x^2,y)\)，目标分布为 \(D=\mathbb R\partial_{y}\)，说明省去淹没条件后，这个逐点定义未必给出常秩分布。
\end{exercise}
\begin{solution}
	在 \(f(p)\) 附近为 \(D\) 取适应坐标 \((u^1,\ldots,u^r,v^1,\ldots,v^{n-r})\)。在其原像中定义
	\[
		H=(v^1,\ldots,v^{n-r})\circ f.
	\]
	横向坐标投影的微分满射，\(\mathrm{d}f\) 也满射，因此 \(\mathrm{d}H\) 满射。由适应坐标中 \(D=\ker\mathrm{d}v\)，有
	\[
		E=\ker\mathrm{d}H.
	\]
	所以在这个邻域内，\(E\) 为秩 \(m-(n-r)\) 光滑子丛；这些局部描述都等于题设的同一个逐点子空间，故给出整体光滑分布。

	若 \(X,Y\) 为取值于 \(E\) 的局部光滑向量场，则 \(XH^\alpha=YH^\alpha=0\)，从而
	\[
		[X,Y]H^\alpha=X(YH^\alpha)-Y(XH^\alpha)=0.
	\]
	故 \([X,Y]\in\ker\mathrm{d}H=E\)。这是局部条件，因此 \(E\) 对合。

	最后的例子中，
	\[
		\mathrm{d}f_{(x,y)}(a,b)=(2xa,b).
	\]
	其像属于竖直线，当且仅当 \(2xa=0\)。于是 \(x\ne0\) 时 \(E_{(x,y)}=\mathbb R\partial_y\)，\(x=0\) 时 \(E_{(0,y)}=\mathrm{T}_{(0,y)}\mathbb R^2\)。维数从一变为二，因此不是本节意义下的常秩光滑分布。
\end{solution}

\begin{exercise}\label{g1:ex:maps-tangent-to-leaves}
	设 \(D\) 为有限维、Hausdorff、第二可数无边界光滑流形 \(M\) 上的常秩对合光滑分布，\(P\) 为任意维数的连通无边界光滑流形，\(h:P\to M\) 光滑。固定 \(p_0\in P\)，令 \(L\) 为通过 \(h(p_0)\) 的最大叶。证明下列条件等价：
	\[
		\mathrm{d}h_q(\mathrm{T}_qP)\subseteq D_{h(q)}\quad(q\in P);
		\qquad
		h=i_L\circ\widehat h\text{，其中 }\widehat h:P\to L\text{ 光滑}.
	\]
	证明光滑因子若存在则唯一。再证明，两点位于同一最大叶，当且仅当它们可以由一条有限分段光滑且每段速度属于 \(D\) 的曲线连接。
\end{exercise}
\begin{solution}
	若光滑因子存在，由链式法则，
	\[
		\mathrm{d}h=\mathrm{d}i_L\circ\mathrm{d}\widehat h,
	\]
	而 \(\mathrm{d}i_L(\mathrm{T}L)=D\)，故得到所需包含。因 \(i_L\) 是集合包含，因子至多一个。

	反过来假设微分像包含于 \(D\)。对每个 \(q\in P\)，选连通的小参数邻域 \(V\)，使 \(h(V)\) 落在一个叶状坐标盒中，记其坐标为 \((u,v)\)。于是
	\[
		\mathrm{d}(v\circ h)=0
	\]
	于 \(V\)，故 \(v\circ h\) 常值，\(h(V)\) 落在单张叶片中。这说明 \(P\) 中映入同一个叶的点组成开集。不同叶给出的这些开集两两不交并覆盖 \(P\)，所以由连通性全部像都在通过 \(h(p_0)\) 的 \(L\) 内。在每个 \(V\) 上，把 \(h\) 看成到该开叶片的映射，其坐标为光滑的 \(u\circ h\)，所以所定义的唯一 \(\widehat h:P\to L\) 光滑。这里没有要求 \(\dim P=\operatorname{rank}D\)，也没有要求 \(h\) 为浸入。

	若两点位于同一叶，最大叶构造给出连接它们的有限叶片链。在各片的球或盒坐标中用线段连接相邻交点，再映回 \(M\)，便得到有限分段光滑曲线；每段速度都属于该叶片的切空间 \(D\)。

	反过来，设 \(\gamma:[0,1]\to M\) 是这样的连续分段光滑曲线。对任意参数值，在它附近取一个相对开小区间，使曲线像落在某叶状坐标盒中。在每个光滑段上，横向坐标满足
	\[
		\frac{\mathrm{d}}{\mathrm{d}t}\bigl(v\circ\gamma\bigr)=0,
	\]
	故为常值；经过分段点时，两侧常值又由连续性相同。因此在这个小区间内，像落在一张叶片中。曲线所处的最大叶随参数局部不变，而 \([0,1]\) 连通，故整条曲线位于同一叶，两个端点也在同一叶。
\end{solution}

\begin{exercise}
	\label{g1:exr:1305-leaf-constant-functions}
	在 \(T^2=\mathbb R^2/\mathbb Z^2\) 上考虑向量场 \(X_\alpha=\partial_x+\alpha\partial_y\) 的叶。证明：若 \(\alpha\) 无理，则每个在各叶上为常数的连续函数 \(H:T^2\to\mathbb R\) 都在整个环面上为常数。若 \(\alpha=p/q\)，其中 \(p\in\mathbb Z\)、\(q\ge1\) 互素，则每个在各叶上为常数的光滑函数 \(H:T^2\to\mathbb R\) 都能唯一写成
	\[
		H=h\circ f,\qquad
		f:T^2\longrightarrow\mathbb R/\mathbb Z,\qquad
		f([x,y])=[qy-px],
	\]
	其中 \(h:\mathbb R/\mathbb Z\to\mathbb R\) 光滑。说明为什么由叶常性得到的集合映射 \(h\) 确实光滑，而不只是良定义。
\end{exercise}
\begin{solution}
	无理情形中，\cref{g1:prop:linear-torus-leaves}说明每条叶都稠密。固定一条叶 \(L\)，设 \(H|_L=c\)。由于 \(H\) 连续，集合 \(H^{-1}(\{c\})\) 闭，又包含 \(L\)，因而包含其闭包 \(T^2\)。所以 \(H\equiv c\)。这里连续性是关键：仅在每条叶上分别规定一个任意常数，并不自动得到连续函数。

	现在设 \(\alpha=p/q\)。将 \((x,y)\) 换成 \((x+m,y+n)\) 时，\(qy-px\) 增加整数 \(qn-pm\)，所以 \(f\) 良定义且光滑。它满射：对任意实代表 \(u\)，点 \([0,u/q]\) 映到 \([u]\)。沿流轨道，
	\[
		f([x+t,y+(p/q)t])=[qy-px],
	\]
	所以每条叶包含于一个纤维。

	还要证明一个纤维只有一条叶。若两点取实代表后的差为 \((\Delta x,\Delta y)\)，且它们的 \(f\) 值相同，则
	\[
		q\Delta y-p\Delta x=k\in\mathbb Z.
	\]
	因 \(p,q\) 互素，可以取整数 \(m,n\) 使 \(qn-pm=k\)。于是
	\[
		q(\Delta y-n)=p(\Delta x-m).
	\]
	令 \(t=\Delta x-m\)，便有 \((\Delta x-m,\Delta y-n)=(t,(p/q)t)\)，即两点在环面上位于同一流轨道。纤维因此恰为叶。

	由叶常性可以规定 \(h(u)=H(z)\)，其中 \(f(z)=u\)。纤维只有一条叶保证它不依赖 \(z\)；满射性保证它有定义且唯一。为证明光滑性，在圆周任一点附近取小开弧 \(U\)，选一个光滑实数角坐标 \(\ell:U\to\mathbb R\)，满足 \([\ell(u)]=u\)。定义局部截面
	\[
		\sigma:U\longrightarrow T^2,\qquad
		\sigma(u)=[0,\ell(u)/q].
	\]
	它光滑且 \(f\circ\sigma=\operatorname{id}_U\)，所以 \(h|_U=H\circ\sigma\) 光滑。每个圆周点都有这样的开弧，故 \(h\) 整体光滑。不同局部截面落在同一纤维上，得到的函数一致，无须使所选实数角坐标全局一致。反过来，每个光滑 \(h\) 的复合 \(h\circ f\) 都光滑并在各叶上为常数。
\end{solution}

\begin{exercise}
	\label{g1:exr:1305-variable-cylinder-holonomy}
	设 \(a:\mathbb R\to\mathbb R\) 为周期 \(1\) 的光滑函数，在圆柱 \((\mathbb R/\mathbb Z)\times\mathbb R\) 上取
	\[
		X=\partial_\theta-a(\theta)r\partial_r.
	\]
	求截面 \(\theta=0\) 的首次正返回映射，并证明中间闭叶 \(r=0\) 的 横截面返回映射芽群平凡，当且仅当 \(\int_0^1a(\theta)\,\mathrm{d}\theta=0\)。在积分为零时，构造一个整体光滑第一积分，并说明它的各个水平集是什么。
\end{exercise}
\begin{solution}
	把初始角度暂时提升为实数 \(\theta_0\)。流方程为 \(\theta'=1\)、\(r'=-a(\theta)r\)，故
	\[
		\theta(t)=\theta_0+t,\qquad
		r(t)=r_0\exp\left(-\int_0^t a(\theta_0+u)\,\mathrm{d}u\right).
	\]
	周期性保证更换 \(\theta_0\) 的整数提升不改变圆柱上的曲线。该公式在所有实时间有定义；它满足初值与速度方程，因此由唯一性就是整体流。

	从 \(\theta=0\) 出发，角度第一次在正时间回到零模 \(1\) 是 \(t=1\)，所以令
	\[
		A=\int_0^1a(u)\,\mathrm{d}u,
	\]
	便有返回映射 \(P(r)=e^{-A}r\)。中间叶 \(r=0\) 是周期为 \(1\) 的圆周，沿正向一圈产生 \(P\) 在零点的芽。圆周上的任意基点环路可逐段选取连续实数角度，提升的终点与起点相差整数 \(k\)；将提升与相同端点的直线作线性同伦，便知该环路与绕 \(k\) 圈同伦。因此由\cref{g1:prop:leaf-holonomy-homotopy}，全部横截面返回映射芽恰由
	\[
		P^k(r)=e^{-kA}r,\qquad k\in\mathbb Z,
	\]
	给出。若 \(A=0\)，这些映射都是恒等映射。若 \(A\ne0\)，则 \(e^{-A}\ne1\)，每个零点邻域都含非零 \(r\)，且 \(e^{-A}r\ne r\)；故连生成元的芽也不是恒等芽。于是中间闭叶的 横截面返回映射芽群平凡恰好等价于 \(A=0\)。

	在 \(A=0\) 时，定义
	\[
		B(\theta)=\int_0^\theta a(u)\,\mathrm{d}u.
	\]
	由周期性，\(B(\theta+1)-B(\theta)=A=0\)，所以 \(B\) 是周期 \(1\) 的光滑函数。于是
	\[
		I([\theta],r)=r e^{B(\theta)}
	\]
	是圆柱上良定义的整体光滑函数，并且
	\[
		X(I)=r e^{B(\theta)}a(\theta)-a(\theta)r e^{B(\theta)}=0.
	\]
	故 \(I\) 为第一积分；其 \(\partial_r I=e^{B(\theta)}>0\)，微分处处非零。每个水平集 \(I=c\) 为图
	\[
		r=c e^{-B(\theta)},\qquad [\theta]\in\mathbb R/\mathbb Z,
	\]
	是一条嵌入圆周。向量场在其中处处非零且角速度为 \(1\)，一圈遍历整个水平集，所以这些水平集恰为各条叶。若 \(A\ne0\)，同一实轴公式满足 \(B(\theta+1)=B(\theta)+A\)，因而 \(r e^{B(\theta)}\) 不能直接下降到圆柱上的函数。
\end{solution}

\begin{exercise}
	\label{g1:exr:1305-triple-origin}
	在
	\[
		M=\mathbb R^2\setminus\{(-1,0),(1,0)\}
	\]
	上考虑水平分布 \(D=\mathbb R\partial_x\)。描述其最大叶，并证明叶空间同胚于具有三个不同原点的实直线。请明确给出这个空间的拓扑，直接检验它与叶空间的商拓扑相同，并证明它满足 \(T_1\) 分离公理，却不是 Hausdorff 空间。
\end{exercise}
\begin{solution}
	当 \(y=c\ne0\) 时，整条水平直线 \(L_c=\{(x,c):x\in\mathbb R\}\) 是一条叶。当 \(y=0\) 时，两个删除点把水平直线分成
	\[
		L_-=\{(x,0):x<-1\},\quad
		L_0=\{(x,0):-1<x<1\},\quad
		L_+=\{(x,0):x>1\}.
	\]
	这三条开区间分别是叶。确实，沿 \(D\) 的任何曲线都满足 \(y'=0\)，所以不能改变高度；在零高度也不能穿越删除点。另一方面，每个所列连通水平区间内的点都能沿水平线段连接，因此这些正是最大连通叶。

	定义集合
	\[
		T=(\mathbb R\setminus\{0\})\sqcup\{0_-,0_0,0_+\}.
	\]
	对非零点使用不含零的通常开区间作为邻域基；对三个原点分别规定基本邻域
	\[
		U_i(\varepsilon)=\{0_i\}\cup
		\bigl((-\varepsilon,\varepsilon)\setminus\{0\}\bigr),
		\qquad i\in\{-,0,+\},\quad\varepsilon>0.
	\]
	这些集合构成拓扑基：同一个原点的两个基本邻域可取较小半径，不同原点的基本邻域之交是非零点的通常开集，其余交集也按普通开区间处理。这个拓扑使每个 \((\mathbb R\setminus\{0\})\cup\{0_i\}\) 都以自然对应同胚于实直线，而三个零点在整体集合中彼此不同。

	定义 \(Q:M\to T\)：在 \(y\ne0\) 时令 \(Q(x,y)=y\)，在 \(L_i\) 上令 \(Q(x,0)=0_i\)。它满射，且纤维恰为上述各条叶。先证连续性。非零开区间 \(J\) 的逆像是开水平带 \(\mathbb R\times J\)。另一方面，
	\[
		Q^{-1}(U_i(\varepsilon))
		=\{(x,y):0<|y|<\varepsilon\}\cup L_i.
	\]
	其中非零高度的点都有留在该集合中的小矩形。若 \((x,0)\in L_i\)，因 \(x\) 位于相应开区间内部，可以选一个小水平区间全部留在该区间，再配上充分小的竖直区间，所得矩形也包含于上式。因此此逆像在 \(M\) 中开，\(Q\) 连续。

	还需证明这正是商拓扑。设 \(A\subset T\)，且 \(Q^{-1}(A)\) 在 \(M\) 中开。若非零点 \(c\in A\)，取 \((0,c)\in Q^{-1}(A)\)，其一个小开矩形仍在该逆像中，竖直投影给出一个包含 \(c\)、不含零且包含于 \(A\) 的开区间。若 \(0_i\in A\)，在 \(L_i\) 中分别选点 \((-2,0)\)、\((0,0)\) 或 \((2,0)\)。逆像的开放性给出该点的一个小矩形；它含所有充分小非零高度上的某些点，故这些非零高度全部属于 \(A\)。连同 \(0_i\)，便得到 \(U_i(\varepsilon)\subset A\)。所以 \(A\) 在规定拓扑中开。结合连续性，已经证明
	\[
		A\subset T\text{ 开}\quad\Longleftrightarrow\quad Q^{-1}(A)\subset M\text{ 开}.
	\]

	记真正的叶空间为 \(M/\mathcal F\)，商映射为 \(\pi\)。由于 \(Q\) 的纤维正是叶，它诱导双射 \(\overline Q:M/\mathcal F\to T\)，且 \(Q=\overline Q\circ\pi\)。两边的开集都由在 \(M\) 中的逆像是否开来判断，刚证的等价便说明 \(\overline Q\) 及其逆连续。因此它是所求同胚。

	最后检验分离性。对于原点 \(0_i\)，其补集中的其他原点各有不含 \(0_i\) 的基本邻域，所有非零点也有不含原点的普通邻域，所以 \(\{0_i\}\) 闭。对于非零点 \(c\)，在每个原点处取半径小于 \(|c|\) 的基本邻域，就能避开 \(c\)；其余非零点可用普通开区间避开 \(c\)，故 \(\{c\}\) 也闭。每个单点均闭，空间因此满足 \(T_1\) 公理。然而任意两个不同原点的邻域都包含所有充分小的非零实数，它们必相交。这两个点没有不交邻域，所以 \(T\) 不是 Hausdorff 空间。
\end{solution}

\begin{exercise}
	\label{g1:exr:1305-horizontal-lift}
	在 \(\mathbb R^3\) 上取接触模型
	\[
		D=\ker(\mathrm{d}z-x\,\mathrm{d}y).
	\]
	给定一个周期为 \(T>0\) 的光滑平面曲线 \(t\mapsto(x(t),y(t))\) 及初始高度 \(z_0\)，证明存在唯一从高度 \(z_0\) 出发、投影为该曲线且始终切于 \(D\) 的提升。证明提升为周期 \(T\) 的闭曲线，当且仅当
	\[
		\int_0^T x(t)y'(t)\,\mathrm{d}t=0.
	\]
	对 \(x(t)=R\cos t\)、\(y(t)=R\sin t\)、\(0\le t\le2\pi\)、\(R>0\)，直接求出提升及一周后的高度变化，不使用外微分或 Stokes 定理。
\end{exercise}
\begin{solution}
	提升曲线 \(\widetilde\gamma(t)=(x(t),y(t),z(t))\) 始终切于 \(D\)，等价于其速度满足
	\[
		z'(t)-x(t)y'(t)=0.
	\]
	因此所求曲线必须且可以取为
	\[
		z(t)=z_0+\int_0^t x(u)y'(u)\,\mathrm{d}u.
	\]
	它光滑，满足初值与切向条件；任何另一个满足条件的高度函数与它之差的导数为零，且初值为零，所以两者相同。

	由于 \(x(t)y'(t)\) 周期为 \(T\)，对任意 \(t\) 有
	\[
		z(t+T)-z(t)=\int_t^{t+T}x(u)y'(u)\,\mathrm{d}u
		=\int_0^T x(u)y'(u)\,\mathrm{d}u.
	\]
	所以右端为零时 \(z\) 也周期为 \(T\)，提升确为光滑闭曲线；若提升闭合，则令 \(t=0\) 就得到右端为零。这个条件同时保证周期光滑性，而不只是把两端的位置勉强接在一起。

	对给定圆周，\(y'(t)=R\cos t\)，故
	\[
		z'(t)=R^2\cos^2t,\qquad
		z(t)=z_0+R^2\left(\frac t2+\frac{\sin2t}{4}\right).
	\]
	一周后的高度变化为
	\[
		z(2\pi)-z(0)=R^2\int_0^{2\pi}\cos^2t\,\mathrm{d}t=\pi R^2.
	\]
	这一非零变化与初始高度无关，所以平面投影已经闭合，水平提升仍没有闭合。切向条件把竖直速度与当前位置 \(x\) 及水平速度 \(y'\) 联系起来，绕行一周可以积累净高度位移；这里的数值完全由一元积分计算得到。
\end{solution}


% !TeX root = ../../Book1.tex
% 纲要标题：## 第14章　微分形式、积分与 Stokes 定理
% 纲要位置：几何学计划纲要.md，第 644 行。
% 正文状态：初稿；尚未正式审核。

% 纲要内容开始
% 
% > 本章目标与证明范围：把外微分、同伦算子、定向和积分组织为完整的流形微积分。闭形式的局部原函数与整体周期障碍贯穿全章，为第二卷上同调比较准备链级输入。
% 
% 纲要内容结束

\chapter{微分形式、积分与 Stokes 定理}
\label{g1:ch:14}
\BookChapterNumber{g1:ch:14}

\begin{chapterabstract}
	外微分把函数的微分推广到所有次数的微分形式，其平方为零，使闭形式与原形式的问题自然出现。本章从坐标构造及唯一性出发，证明外微分的自然性与 Cartan 公式，再以显式区间积分建立同伦公式和 Poincaré 引理。定向与密度解决两种不同的积分需求；单位分解把局部积分粘成整体，并保证坐标与分解的选择不影响结果。外法向优先的约定确定边界符号，半空间中的微积分基本定理随后给出流形及带角区域的 Stokes 定理。最后对光滑奇异单形积分，证明积分与边界相容，用周期解释局部原形式的整体障碍，并提出后卷将证明的次数公式与 de Rham 比较问题。
\end{chapterabstract}

\begin{learninggoals}
	\begin{itemize}
		\item 构造外微分，证明其坐标无关性、唯一性、平方为零以及与拉回的相容性，并用 Cartan 公式联系插入和李导数。
		\item 写出同伦算子，逐项核对区间积分的符号，证明星形区域与流形局部的 Poincaré 引理，区分零次与正次情形。
		\item 区分定向线、顶次形式和密度，构造紧支积分并证明其与坐标及单位分解无关，正确使用定向换元公式。
		\item 从外法向优先约定推导边界方向，完整证明 Stokes 定理及带角版本，说明 Green、Gauss 公式中的方向和边界项。
		\item 对光滑奇异链积分，证明积分与链边界相容，利用周期判断恰当性的障碍，明确次数公式与 de Rham 定理尚需的后续理论。
	\end{itemize}
\end{learninggoals}

在一元微积分中，求导和积分由端点差联系起来。在平面上，沿闭曲线的环流又能写成区域内部某种导数的积分。这些公式的被积对象维数不同，边界方向也各有说法，却都在比较“区域内部的变化”与“沿边界的累积”。微分形式提供适合拉回和定向的被积对象，外微分提供共同的导数，Stokes 定理则把这种联系写成一个维数无关的等式。

\ProseParagraphBreak

另一个问题同样贯穿本章：知道一个对象在每个小邻域都有原函数，为什么还可能找不到整体原函数？穿孔平面的角形式最能显出差别。沿短弧可以连续选择角度，绕一圈后角度却增加 \(2\pi\)。局部微分计算没有出错，真正未被它消除的是整体绕行。我们会先用同伦公式制造局部原形式，再用非零周期证明某些局部原形式无法拼合。

\ProseParagraphBreak

本章默认流形有限维、Hausdorff、第二可数，除明确要求无边界之处外允许边界；光滑性沿用\chapref{g1:ch:11}的局部延拓定义。需要回忆的前置内容是\chapref{g1:ch:12}的交替张量、楔积、插入、拉回、密度和光滑单位分解，以及\secref{g1:sec:1303}的李导数。每次调用这些工具时会指出具体结论和约定。积分部分调用欧氏多元换元、Fubini 定理与一元微积分基本定理，其假设在\secref{g1:sec:1403}列明；这里不要求读者已经学习过一般流形积分。

\ProseParagraphBreak

本章主线中的构造和证明在正文完成。最后一节的次数积分公式依纲要先作准确陈述，一般次数的定义、整数性和公式的证明留在第二卷第16.1节；de Rham 比较的完整证明留在第二卷第7.3节。它们不作为前面论证的依据。延伸阅读可按 Lee 第二版的主题查找：Chapter 14 的 Exterior Derivatives，Chapter 15 的 Orientations of Manifolds，Chapter 16 的 Integration of Differential Forms、Stokes's Theorem、Manifolds with Corners 与 Densities，以及 Chapter 17 的 Homotopy Invariance\cite[Chapter 14, pp. 362--372; Chapter 15, pp. 380--387; Chapter 16, pp. 402--433; Chapter 17, pp. 443--447]{Lee2012Smooth}。

% !TeX root = ../../Book1.tex
% 纲要标题：### 14.1 外微分与自然运算
% 纲要位置：几何学计划纲要.md，第 648 行。
% 正文状态：初稿已撰写，尚未经正式审核。

% 纲要内容开始
% 
% * 拉回、楔积及坐标无关性
% * 外微分的构造与唯一性
% * 内乘和 Cartan 公式
% * 外微分与拉回、Lie 导数的相容
% 
% 纲要内容结束

\section{外微分与自然运算}
\label{g1:sec:1401}

第12章把微分形式定义为交替协变张量场，第13章又介绍了沿向量场的 Lie 导数。它们回答了两个不同的问题：形式怎样接受切向量作为输入，形式怎样随一个流变化。本节增加一种不需要先选向量场的微分运算。它把 \(k\) 形式变成 \(k+1\) 形式，函数的普通微分是它的零次情形；它与拉回相容，因而可以在不同参数空间之间传递计算。后面讨论闭形式、局部原函数和积分时，反复使用的正是这些性质。

\ProseParagraphBreak

本节的流形均为有限维、Hausdorff 且第二可数的实光滑流形，允许有边界；没有特别说明时，形式和映射均光滑。边界图仍采用 \(\mathbb H^n=\{x_n\geq0\}\)，光滑性指局部可延拓到 Euclid 开集。涉及流的解释时会另行说明向量场的条件。

\subsection{拉回、楔积及坐标无关性}
\label{g1:subsec:1401:01}

回忆\cref{g1:def:smooth-tensor-field}，\(M\) 上的 \(k\) 形式是 \(\Lambda^k\mathrm{T}^*M\) 的光滑截面，空间记作 \(\Omega^k(M)\)。约定 \(\Omega^0(M)=C^\infty(M)\)，且在 \(k<0\) 或 \(k>\dim M\) 时取零空间。若 \((x^1,\ldots,x^n)\) 为局部坐标，记
\[
	\mathrm{d}x^I=\mathrm{d}x^{i_1}\wedge\cdots\wedge\mathrm{d}x^{i_k},\qquad I=(i_1<\cdots<i_k).
\]
任意 \(k\) 形式唯一表示为
\[
	\omega=\sum_{i_1<\cdots<i_k}\omega_{i_1\cdots i_k}\,\mathrm{d}x^{i_1}\wedge\cdots\wedge\mathrm{d}x^{i_k}
	=\sum_I\omega_I\,\mathrm{d}x^I,
\]
其中 \(\omega_I=\omega(\partial_{i_1},\ldots,\partial_{i_k})\)。这里仅对递增指标求和，因而没有额外的 \(1/k!\)。若改为对全部有序指标求和，并令系数交替，才需要除以 \(k!\)。零次情形把空楔积理解为 \(1\)，便是一个函数自身。

\ProseParagraphBreak

楔积在每个点由交替形式的代数运算定义。\cref{g1:prop:wedge-algebra}已证明：若 \(\alpha\in\Omega^p(M)\)、\(\beta\in\Omega^q(M)\)，则
\[
	\alpha\wedge\beta=(-1)^{pq}\beta\wedge\alpha,
\]
并且楔积结合、对函数的乘法为零次楔积。例如 \(\mathrm{d}y\wedge\mathrm{d}x=-\mathrm{d}x\wedge\mathrm{d}y\)，而一个二形式与另一个二形式交换时没有负号。奇数次形式的自楔积为零；偶数次形式则未必如此，例如在 \(\mathbb R^4\) 中，\((\mathrm{d}x^1\wedge\mathrm{d}x^2+\mathrm{d}x^3\wedge\mathrm{d}x^4)^2=2\,\mathrm{d}x^1\wedge\cdots\wedge\mathrm{d}x^4\)。这些符号来自交换输入的次序，与坐标的具体选择无关。

\ProseParagraphBreak

设 \(F:M\to N\) 光滑。由\cref{g1:prop:covariant-tensor-pullback}，形式的拉回定义为
\[
	(F^*\omega)_p(v_1,\ldots,v_k)
	=\omega_{F(p)}\bigl(\mathrm{d}F_p(v_1),\ldots,\mathrm{d}F_p(v_k)\bigr).
\]
先把输入的切向量送到目标点，再在那里求值，所得自然是定义域上的形式。这个定义只需要 \(\mathrm{d}F_p\)，不需要它可逆；因此嵌入、投影和有临界点的映射都可以拉回形式。对函数，定义为 \(F^*f=f\circ F\)。

\ProseParagraphBreak

在目标坐标 \((y^1,\ldots,y^m)\) 中写 \(\omega=\sum_I\omega_I\,\mathrm{d}y^I\)，并记 \(F^a=y^a\circ F\)。函数微分的链式法则给出 \(F^*(\mathrm{d}y^a)=\mathrm{d}F^a\)。逐点的楔积与线性拉回相容，于是
\begin{equation}
	\label{g1:eq:form-pullback-coordinates}
	F^*\omega=\sum_{i_1<\cdots<i_k}(\omega_{i_1\cdots i_k}\circ F)\,\mathrm{d}F^{i_1}\wedge\cdots\wedge\mathrm{d}F^{i_k}.
\end{equation}
同样由逐点定义可得
\[
	F^*(\alpha\wedge\beta)=F^*\alpha\wedge F^*\beta,\qquad
	(G\circ F)^*=F^*\circ G^*,\qquad (\operatorname{id}_M)^*=\operatorname{id}.
\]
第一式也可直接对楔积的洗牌求和公式逐项代入 \(\mathrm{d}F_p(v_i)\) 验证；第二式就是 \(\mathrm{d}(G\circ F)=\mathrm{d}G\circ\mathrm{d}F\)。因而这些公式在边界点同样成立。此处的 \(\mathrm{d}F^a\) 仍只是已经定义的函数微分，尚未用到本节将构造的高次外微分。

\begin{example}
	令 \(F:\mathbb R^2\to\mathbb R^2\)、\(F(u,v)=(u^2,v)\)，目标坐标记作 \((x,y)\)。则
	\[
		F^*(\mathrm{d}x\wedge\mathrm{d}y)=\mathrm{d}(u^2)\wedge\mathrm{d}v=2u\,\mathrm{d}u\wedge\mathrm{d}v.
	\]
	它在 \(u=0\) 处为零，恰好反映该点的微分无法提供两个线性无关的目标切向量。拉回仍处处光滑。更一般地，若 \(\operatorname{rank}\mathrm{d}F_p<k\)，任何 \(k\) 形式在 \(p\) 处的拉回都为零，因为它的 \(k\) 个输入像线性相关。特别地，把二形式拉回一条曲线，得到的是零二形式。
\end{example}

“坐标无关”并不意味着坐标系数不变。比如 \(y=x^2\) 在 \(x>0\) 上给出 \(\mathrm{d}y=2x\,\mathrm{d}x\)，同一个一形式的系数就可能随坐标改变。它意味着：把系数和基同时按链式法则变换后，对同一组切向量求值得到相同的数。下面构造外微分时，也必须检查这一点，不能只在一组坐标里写下公式便结束。

\subsection{外微分的构造与唯一性}
\label{g1:subsec:1401:02}

函数的微分记录一阶变化，但直接对一个形式的各个系数求导，会同时出现一个导数指标和原来的交替指标。将新指标也作交替处理，就得到一个多一次的形式。楔积负责这一交替化；代价是丢掉部分二阶信息，好处是这个运算不需要度量或联络，而且连续做两次总为零。

\begin{theorem}[外微分]\label{g1:thm:exterior-derivative}
	设 \(M\) 为有限维、Hausdorff 且第二可数的光滑流形，允许有边界。存在唯一的一族实线性映射
	\[
		\mathrm{d}:\Omega^k(M)\longrightarrow\Omega^{k+1}(M),\qquad k\geq0,
	\]
	满足以下性质：
	\begin{enumerate}
		\item 对 \(f\in C^\infty(M)\)，\(\mathrm{d}f\) 是函数的普通微分，即 \((\mathrm{d}f)_p(v)=v(f)\)。
		\item 若 \(\alpha\in\Omega^p(M)\)、\(\beta\in\Omega^q(M)\)，则
		\[
			\mathrm{d}(\alpha\wedge\beta)=\mathrm{d}\alpha\wedge\beta+(-1)^p\alpha\wedge\mathrm{d}\beta.
		\]
		\item 对任意形式 \(\omega\)，有 \(\mathrm{d}(\mathrm{d}\omega)=0\)。
	\end{enumerate}
	这个运算称为\term[waiweifen]{外微分}[exterior derivative]。它与限制到开集相容，并且在任意坐标域上有
	\begin{equation}
		\label{g1:eq:exterior-coordinate-formula}
		\mathrm{d}\omega
		=\sum_I\mathrm{d}\omega_I\wedge\mathrm{d}x^I
		=\sum_I\sum_{j=1}^n\frac{\partial\omega_I}{\partial x^j}\,\mathrm{d}x^j\wedge\mathrm{d}x^I,
		\qquad \omega=\sum_I\omega_I\,\mathrm{d}x^I.
	\end{equation}
\end{theorem}
\begin{proof}
	先在一个坐标域 \(U\) 中用\eqref{g1:eq:exterior-coordinate-formula}定义算子 \(\mathrm{D}_x\)。边界图中对系数任取局部光滑延拓后求偏导；\cref{g1:prop:boundary-smooth-derivative}说明任意两种延拓的所有阶偏导在半空间上相同，因此定义不依赖延拓。系数仍光滑，且零次公式就是通常的 \(\mathrm{d}f\)。

	分次乘积法则可以在单项式上验证，再由双线性推广。令 \(|I|=p\)，则
	\[
		\mathrm{D}_x\bigl((a\,\mathrm{d}x^I)\wedge(b\,\mathrm{d}x^J)\bigr)
		=(b\,\mathrm{d}a+a\,\mathrm{d}b)\wedge\mathrm{d}x^I\wedge\mathrm{d}x^J.
	\]
	另一方面，所要证明的右端是
	\[
		b\,\mathrm{d}a\wedge\mathrm{d}x^I\wedge\mathrm{d}x^J
		+(-1)^p a\,\mathrm{d}x^I\wedge\mathrm{d}b\wedge\mathrm{d}x^J.
	\]
	把第二项中的 \(\mathrm{d}b\) 移到 \(p\) 个一形式之前，产生另一个 \((-1)^p\)，两式因此相等。即使指标重复、相应楔积为零，等式仍然成立。

	对单项式再作一次 \(\mathrm{D}_x\)，得到
	\[
		\mathrm{D}_x^2(a\,\mathrm{d}x^I)
		=\sum_{j,\ell}\frac{\partial^2a}{\partial x^\ell\partial x^j}\,\mathrm{d}x^\ell\wedge\mathrm{d}x^j\wedge\mathrm{d}x^I=0.
	\]
	当 \(j=\ell\) 时，楔积为零；当 \(j\ne\ell\) 时，交换 \(j,\ell\) 的两项系数相等、楔积互为相反数。混合偏导可交换，在边界处也可由任一光滑延拓验证。因此 \(\mathrm{D}_x^2=0\)。

	下面先在坐标域内证明唯一性。若 \(\delta\) 是该域上一族满足三条性质的算子，则
	\[
		\delta(\mathrm{d}x^i)=\delta(\delta x^i)=0.
	\]
	反复使用乘积法则可得 \(\delta(\mathrm{d}x^I)=0\)，从而
	\[
		\delta\Bigl(\sum_I a_I\,\mathrm{d}x^I\Bigr)=\sum_I\mathrm{d}a_I\wedge\mathrm{d}x^I=\mathrm{D}_x\Bigl(\sum_I a_I\,\mathrm{d}x^I\Bigr).
	\]
	这一步只用了函数微分和已经验证的代数性质。

	若 \(x,y\) 两组坐标重叠，\(\mathrm{D}_x,\mathrm{D}_y\) 都在重叠开集上满足上述三条性质，并且在函数上都是同一个普通微分。刚才的局部唯一性于是给出 \(\mathrm{D}_x=\mathrm{D}_y\)。所以各坐标算子可粘合为 \(M\) 上的 \(\mathrm{d}\)，并且显然与限制相容。至此已证明存在性和坐标无关性；这里没有预先使用外微分与任意映射拉回相容的结论。

	最后说明，整体唯一性不需要额外假定候选算子是局部的。设 \(D\) 是整体形式上满足三条性质的另一族算子。若整体形式 \(\eta\) 在 \(p\) 的开邻域 \(U\) 上为零，用\cref{g1:cor:smooth-cutoff-closed}选取 \(\chi\)，使它在 \(p\) 附近等于 \(1\)、支集包含于 \(U\)。于是 \(\chi\eta=0\)，且
	\[
		0=D(\chi\eta)=\mathrm{d}\chi\wedge\eta+\chi D\eta.
	\]
	在 \(\chi=1\) 的邻域，\(\mathrm{d}\chi=0\)，故 \(D\eta=0\)。因此 \(D\eta\) 在一点处的值仅由 \(\eta\) 在该点附近的值决定。

	对只在开集上定义的形式，在每一点附近先乘上述截断函数再向外补零，便得到一个整体光滑延拓；把 \(D\) 作用于延拓后限制回来。刚证明的局部性保证结果在该点附近不依赖延拓；各点的结果遂粘合为局部算子。它仍满足三条性质，因为在一点附近可选共同的整体代表，作两次 \(D\) 时也可使用第一次结果的同一代表。局部唯一性给出这个算子等于坐标构造的 \(\mathrm{d}\)，因而原来的整体 \(D\) 也等于 \(\mathrm{d}\)。
\end{proof}

\(\mathrm{d}^2=0\) 并不表示函数的全部二阶导数为零。例如 \(f=x^2\) 有 \(\partial_x^2f=2\)，但
\[
	\mathrm{d}(\mathrm{d}f)=\mathrm{d}(2x\,\mathrm{d}x)=2\,\mathrm{d}x\wedge\mathrm{d}x=0.
\]
消失的是二阶变化的交替部分。外微分也不是把函数变成梯度向量场：\(\mathrm{d}f\) 是一形式，它的定义不需要度量；将它转成梯度，还要用度量识别切空间与余切空间。

\begin{example}
	在平面上令 \(\alpha=P\,\mathrm{d}x+Q\,\mathrm{d}y\)。按坐标公式计算，
	\[
		\mathrm{d}\alpha
		=\mathrm{d}P\wedge\mathrm{d}x+\mathrm{d}Q\wedge\mathrm{d}y
		=(\partial_xQ-\partial_yP)\,\mathrm{d}x\wedge\mathrm{d}y.
	\]
	只有交叉偏导留下来，且 \(\partial_yP\) 项的负号来自 \(\mathrm{d}y\wedge\mathrm{d}x=-\mathrm{d}x\wedge\mathrm{d}y\)。例如 \(y\,\mathrm{d}x\) 的外微分为 \(-\mathrm{d}x\wedge\mathrm{d}y\)，而 \(x\,\mathrm{d}x=\mathrm{d}(x^2/2)\) 的外微分为零。
\end{example}

\begin{definition}[闭形式与恰当形式]\label{g1:def:closed-exact-form}
	设 \(M\) 为本节约定的光滑流形，\(\omega\in\Omega^k(M)\)。若 \(\mathrm{d}\omega=0\)，称 \(\omega\) 为\term[bixingshi]{闭形式}[closed form]。若存在 \(\eta\in\Omega^{k-1}(M)\) 使 \(\omega=\mathrm{d}\eta\)，称 \(\omega\) 为\term[qiadangxingshi]{恰当形式}[exact form]，并称 \(\eta\) 为 \(\omega\) 的一个\term[yuanxingshi]{原形式}[primitive]。按照 \(\Omega^{-1}(M)=0\) 的约定，零次恰当形式只有零函数。
\end{definition}

由 \(\mathrm{d}^2=0\)，恰当形式一定闭。零次闭形式恰为局部常值函数：在小的坐标球或半球内，沿线段对函数求导，\(\mathrm{d}f=0\) 使函数值不变；反向显然成立。至于高次闭形式是否有原形式，以及原形式能否从局部拼成整体，尚未由上述代数恒等式回答。还要留意，\(n\) 维流形上的每个 \(n\) 形式都闭，因为 \(\Omega^{n+1}(M)=0\)；这也不自动给出它的原形式。

\subsection{内乘和 Cartan 公式}
\label{g1:subsec:1401:03}

给形式增加一个交替输入的是外微分；给它预先指定一个输入，则是第12章的\termidx[charu]{插入}运算。若 \(X\) 为光滑向量场，定义 \(\iota_X\omega\) 在每点由 \(X_p\) 插入第一输入得到，即
\[
	(\iota_X\omega)(Y_1,\ldots,Y_{k-1})=\omega(X,Y_1,\ldots,Y_{k-1}),\qquad k\geq1,
\]
并约定 \(\iota_Xf=0\)。这一运算也称内乘；“第一输入”的约定不能省略，改用最后输入会引入依赖次数的符号。对局部单项式有
\[
	\iota_X(a\,\mathrm{d}x^{i_1}\wedge\cdots\wedge\mathrm{d}x^{i_k})
	=a\sum_{b=1}^k(-1)^{b-1}X^{i_b}\,\mathrm{d}x^{i_1}\wedge\cdots\wedge\widehat{\mathrm{d}x^{i_b}}\wedge\cdots\wedge\mathrm{d}x^{i_k},
\]
其中帽号表示删去该因子。这可由行列式沿第一个输入展开得到，也表明插入后的形式仍光滑。

\ProseParagraphBreak

\cref{g1:prop:interior-product-rule}的逐点代数证明给出
\[
	\iota_X(\alpha\wedge\beta)=(\iota_X\alpha)\wedge\beta+(-1)^p\alpha\wedge\iota_X\beta,
	\qquad \alpha\in\Omega^p(M),
\]
以及 \(\iota_X^2=0\)。为记住乘积法则的符号，可以把插入看成依次寻找它所作用的一形式因子：若这一因子在第二组中，便要先越过第一组的 \(p\) 个一形式。与外微分不同，插入对函数系数是线性的，\(\iota_X(f\omega)=f\iota_X\omega\)；外微分则多出 \(\mathrm{d}f\wedge\omega\)。

\ProseParagraphBreak

再回忆 Lie 导数的约定。在无边界流形上，\cref{g1:def:lie-derivative-tensor}定义
\[
	\mathcal L_X\omega=\left.\frac{\mathrm{d}}{\mathrm{d}t}\right|_{t=0}(\Phi_t^X)^*\omega.
\]
这里使用正向流的拉回。\cref{g1:prop:lie-derivative-tensor-rules,g1:prop:cartan-insertion-identities}已证明 \(\mathcal L_X\) 为次数零的楔积导子，而且 \(\mathcal L_Xf=Xf\)。对函数微分还可直接计算
\[
	\mathcal L_X(\mathrm{d}f)
	=\left.\frac{\mathrm{d}}{\mathrm{d}t}\right|_{0}\mathrm{d}(f\circ\Phi_t^X)
	=\mathrm{d}(Xf).
\]
第一个等式只用普通函数的链式法则；第二个等式在坐标中交换光滑函数的时间偏导与空间偏导。因此这一计算没有预用高次外微分与拉回相容的结论。

\ProseParagraphBreak

在有边界的流形上，一般向量场未必产生保持流形的双侧流，因而不能直接照搬上面的流定义。此时可以在边界图中把 \(X\) 和形式的系数局部光滑延拓到 Euclid 开集，在那里求 Lie 导数，再限制回来。具体而言，\cref{g1:prop:lie-derivative-tensor-rules}的坐标公式对形式写成
\begin{equation}
	\label{g1:eq:boundary-form-lie-derivative}
	\mathcal L_X\omega
	=\sum_I X(\omega_I)\,\mathrm{d}x^I
	+\sum_I\sum_{b=1}^k\omega_I\,\mathrm{d}x^{i_1}\wedge\cdots\wedge\mathrm{d}X^{i_b}\wedge\cdots\wedge\mathrm{d}x^{i_k}.
\end{equation}
它只用系数的一阶导数，故由\cref{g1:prop:boundary-smooth-derivative}不依赖延拓。两张边界图得到的结果在内部相同，而内部点在每个边界邻域中稠密，所以结果在边界处也由连续性相同。这就定义了任意光滑 \(X\) 在边界上的 \(\mathcal L_X\)。上述乘积法则及 \(\mathcal L_X\mathrm{d}f=\mathrm{d}(Xf)\) 均由局部延拓继续成立。若 \(X\) 切于边界，\cref{g1:prop:boundary-tangent-flow}给出保持 \(M\) 的双侧局部流，此定义便与流拉回导数一致；对一般 \(X\)，本节只使用刚定义的局部微分运算。

\begin{theorem}[Cartan 公式]\label{g1:thm:cartan-formula}
	设 \(M\) 为有限维、Hausdorff 且第二可数的光滑流形，允许有边界，\(X\) 为光滑向量场。有边界时，Lie 导数按\eqref{g1:eq:boundary-form-lie-derivative}定义。对每个 \(\omega\in\Omega^k(M)\)，有
	\[
		\mathcal L_X\omega=\mathrm{d}(\iota_X\omega)+\iota_X(\mathrm{d}\omega).
	\]
\end{theorem}
\begin{proof}
	令 \(C_X=\mathrm{d}\iota_X+\iota_X\mathrm{d}\)，先证明它为次数零的楔积导子。若 \(\alpha\) 的次数为 \(p\)，分别展开两项的分次乘积法则，得到
	\begin{align*}
		C_X(\alpha\wedge\beta)
		={}&(C_X\alpha)\wedge\beta+\alpha\wedge C_X\beta\\
		&+\bigl((-1)^{p-1}+(-1)^p\bigr)(\iota_X\alpha)\wedge\mathrm{d}\beta\\
		&+\bigl((-1)^p+(-1)^{p+1}\bigr)\mathrm{d}\alpha\wedge\iota_X\beta.
	\end{align*}
	最后两行恰好抵消，故所需的导子法则成立；\(p=0\) 时含 \(\iota_X\alpha\) 的项本来就是零。

	对函数 \(f\)，有 \(C_Xf=\iota_X\mathrm{d}f=Xf=\mathcal L_Xf\)。对函数微分，由 \(\mathrm{d}^2=0\)，
	\[
		C_X(\mathrm{d}f)=\mathrm{d}(\iota_X\mathrm{d}f)=\mathrm{d}(Xf)=\mathcal L_X(\mathrm{d}f).
	\]
	任意形式在坐标域中都是函数系数乘坐标一形式的楔积之和。\(C_X\) 与 \(\mathcal L_X\) 都实线性、都满足次数零的导子法则，又在函数和每个 \(\mathrm{d}x^i\) 上相同，所以在全部形式上相同。边界处所需的性质已由局部延拓证明，因而同一个论证也适用。
\end{proof}

公式的两项有不同的作用顺序：\(\mathrm{d}\iota_X\omega\) 先指定一个输入，再看剩余系数如何变化；\(\iota_X\mathrm{d}\omega\) 则先对全部输入作交替微分，再代入 \(X\)。两项相加恰好包含了 Lie 导数所需的全部变化，包括随流移动的输入方向。比如在 \(\mathbb R^2\) 中取 \(X=x\partial_x\)、\(\omega=\mathrm{d}x\wedge\mathrm{d}y\)，则 \(\iota_X\omega=x\,\mathrm{d}y\)、\(\mathrm{d}\omega=0\)，所以 \(\mathcal L_X\omega=\mathrm{d}x\wedge\mathrm{d}y\)。这与流 \((x,y)\mapsto(e^t x,y)\) 对该二形式的拉回为 \(e^t\omega\) 一致。

\begin{proposition}[一形式的微分]\label{g1:prop:one-form-exterior-evaluation}
	设 \(M\) 为有限维、Hausdorff 且第二可数的光滑流形，允许有边界，\(\alpha\in\Omega^1(M)\)，\(X,Y\) 为光滑向量场。则
	\begin{equation}
		\label{g1:eq:exterior-one-form-evaluation}
		\mathrm{d}\alpha(X,Y)=X\bigl(\alpha(Y)\bigr)-Y\bigl(\alpha(X)\bigr)-\alpha([X,Y]).
	\end{equation}
	有边界时，括号以坐标公式定义，等价于把向量场局部光滑延拓后取括号再限制。
\end{proposition}
\begin{proof}
	Cartan 公式给出
	\[
		(\mathcal L_X\alpha)(Y)=Y\bigl(\alpha(X)\bigr)+\mathrm{d}\alpha(X,Y).
	\]
	另一方面，\eqref{g1:eq:lie-one-form-evaluation}给出 \((\mathcal L_X\alpha)(Y)=X(\alpha(Y))-\alpha([X,Y])\)。两式相减即得结论。边界处的括号只涉及系数的一阶导数，因此与延拓无关；前面的两式也可由局部延拓后限制得到。
\end{proof}

括号修正项说明，不能只把外微分理解为“对系数交叉求导”。在坐标基 \(\partial_i,\partial_j\) 中，括号为零，交叉求导公式确实成立；换成不对易的向量场作为输入，便必须补上最后一项。它也使右侧对 \(X,Y\) 的函数倍乘仍为双线性，符合二形式在每一点的性质。

\begin{example}\label{g1:exm:contact-form-exterior-test}
	第13章的接触例子使用 \(\mathbb R^3\) 上的一形式 \(\alpha=\mathrm{d}z-x\,\mathrm{d}y\)。现在可直接计算
	\[
		\mathrm{d}\alpha=-\mathrm{d}x\wedge\mathrm{d}y,\qquad
		\alpha\wedge\mathrm{d}\alpha=-\mathrm{d}x\wedge\mathrm{d}y\wedge\mathrm{d}z.
	\]
	令 \(A=\partial_x\)、\(B=\partial_y+x\partial_z\)。它们张成 \(\ker\alpha\)，且 \(\mathrm{d}\alpha(A,B)=-1=-\alpha([A,B])\)，与\cref{g1:prop:contact-model-nonintegrable}的括号计算相符。

	更一般地，在无边界流形上，若一形式 \(\alpha\) 处处非零，则 \(\alpha:\mathrm{T}M\to M\times\mathbb R\) 是秩为 \(1\) 的丛同态，\cref{g1:thm:constant-rank-bundle-morphism}保证 \(D=\ker\alpha\) 为秩 \(n-1\) 的光滑分布。对取值于 \(D\) 的局部向量场 \(U,V\)，\eqref{g1:eq:exterior-one-form-evaluation}化为
	\[
		\mathrm{d}\alpha(U,V)=-\alpha([U,V]).
	\]
	当 \(n\geq2\) 时，\(\alpha\wedge\mathrm{d}\alpha=0\) 当且仅当 \(\mathrm{d}\alpha\) 在 \(D\times D\) 上为零。事实上，局部选 \(Z\) 使 \(\alpha(Z)=1\)，则
	\[
		(\alpha\wedge\mathrm{d}\alpha)(Z,u,v)=\mathrm{d}\alpha(u,v),\qquad u,v\in D.
	\]
	反向把每个向量分解成 \(Z\) 方向与 \(D\) 方向；三个输入全在 \(D\) 中时楔积自动为零，恰有一个在 \(Z\) 方向时由假设为零，至少两个在 \(Z\) 方向时由交替性为零。这证明等价。由此和\cref{g1:thm:frobenius}，\(D\) 局部可积当且仅当 \(\alpha\wedge\mathrm{d}\alpha=0\)。在三维情形，\(\alpha\wedge\mathrm{d}\alpha\) 处处非零则恰好表达第13章定义的接触条件：括号在 \(\mathrm{T}M/D\) 中的交替配对处处非零。若改用同一分布的另一局部定义形式 \(f\alpha\)，其中 \(f\) 处处非零，则
	\[
		(f\alpha)\wedge\mathrm{d}(f\alpha)=f^2\alpha\wedge\mathrm{d}\alpha,
	\]
	所以零与非零的判断不受这个选择影响。
\end{example}

\subsection{外微分与拉回、Lie 导数的相容}
\label{g1:subsec:1401:04}

外微分的构造已经证明换坐标相容；现在可以进一步问，两个不同流形之间的光滑映射是否也与它相容。这个问题比换坐标更强，因为映射可以降维，也可以在某些点秩下降。

\begin{proposition}[拉回与外微分]\label{g1:prop:exterior-pullback}
	设 \(M,N\) 为有限维、Hausdorff 且第二可数的光滑流形，允许有边界，\(F:M\to N\) 为光滑映射，不要求 \(F\) 的微分具有常秩。对任意 \(\omega\in\Omega^k(N)\)，有
	\[
		\mathrm{d}_M(F^*\omega)=F^*(\mathrm{d}_N\omega).
	\]
	两侧的下标只用于指出所在流形，通常省去。
\end{proposition}
\begin{proof}
	在目标坐标域中写 \(\omega=\sum_I\omega_I\,\mathrm{d}y^I\)，在它的原像内记 \(F^a=y^a\circ F\)。由\eqref{g1:eq:form-pullback-coordinates}和外微分的乘积法则，
	\[
		\mathrm{d}_M(F^*\omega)
		=\sum_I\mathrm{d}(\omega_I\circ F)\wedge\mathrm{d}F^{i_1}\wedge\cdots\wedge\mathrm{d}F^{i_k}.
	\]
	对各个 \(\mathrm{d}F^{i_a}\) 再求外微分的项全部为零，因为 \(\mathrm{d}^2F^{i_a}=0\)。函数微分的链式法则给出
	\[
		\mathrm{d}(\omega_I\circ F)=\sum_{b=1}^{\dim N}\bigl((\partial_b\omega_I)\circ F\bigr)\,\mathrm{d}F^b.
	\]
	代回便得到
	\[
		\mathrm{d}_M(F^*\omega)
		=\sum_{I,b}\bigl((\partial_b\omega_I)\circ F\bigr)\,\mathrm{d}F^b\wedge\mathrm{d}F^{i_1}\wedge\cdots\wedge\mathrm{d}F^{i_k}
		=F^*(\mathrm{d}_N\omega).
	\]
	这些局部等式覆盖 \(M\)，所以结论整体成立。边界图中的偏导和链式法则已由\cref{g1:prop:boundary-smooth-derivative}保证与延拓无关，故上述证明原样适用；即使 \(F\) 的像完全位于目标边界，也没有使用其像附近的内部点来替代计算。
\end{proof}

拉回因而把闭形式送到闭形式，把恰当形式送到恰当形式：若 \(\omega=\mathrm{d}\eta\)，则 \(F^*\omega=\mathrm{d}(F^*\eta)\)。这与映射是否单射、满射都无关。反过来却不能由 \(F^*\omega\) 闭或恰当推出 \(\omega\) 具有相应性质，例如常值映射把所有正次形式都拉回为零。对边界包含映射 \(j:\partial M\hookrightarrow M\)，同一命题特别给出 \(\mathrm{d}_{\partial M}(j^*\omega)=j^*(\mathrm{d}_M\omega)\)；这只是形式与微分的相容式，此时还没有用到积分定理。

\begin{proposition}[Lie 导数与外微分]\label{g1:prop:exterior-lie-commute}
	设 \(M\) 为有限维、Hausdorff 且第二可数的光滑流形，允许有边界，\(X\) 为光滑向量场；有边界时沿用\eqref{g1:eq:boundary-form-lie-derivative}的定义。则
	\[
		[\mathcal L_X,\mathrm{d}]=0,
		\qquad\text{即}\qquad \mathcal L_X\mathrm{d}\omega=\mathrm{d}\mathcal L_X\omega
	\]
	对任意形式 \(\omega\) 成立。
\end{proposition}
\begin{proof}
	Cartan 公式与 \(\mathrm{d}^2=0\) 给出
	\[
		\mathrm{d}\mathcal L_X
		=\mathrm{d}(\mathrm{d}\iota_X+\iota_X\mathrm{d})
		=\mathrm{d}\iota_X\mathrm{d},\qquad
		\mathcal L_X\mathrm{d}
		=(\mathrm{d}\iota_X+\iota_X\mathrm{d})\mathrm{d}
		=\mathrm{d}\iota_X\mathrm{d}.
	\]
	两者相同。这里 \(\mathcal L_X\) 的次数为零，故括号是通常的算子交换子，没有附加的次数符号。
\end{proof}

若有保持 \(M\) 的局部流，也可从 \((\Phi_t^X)^*\mathrm{d}\omega=\mathrm{d}((\Phi_t^X)^*\omega)\) 对时间求导得到同一结论；刚才的代数证明还适用于边界处只有局部延拓定义的 Lie 导数。由此，Lie 导数保持闭形式和恰当形式。而对一个正次闭形式，Cartan 公式给出更强的事实：
\[
	\mathrm{d}\omega=0\quad\Longrightarrow\quad \mathcal L_X\omega=\mathrm{d}(\iota_X\omega).
\]
所以闭形式沿一个流的无穷小变化本身总是恰当的。这里说的是变化量 \(\mathcal L_X\omega\)，并没有断言原来的 \(\omega\) 恰当。

\ProseParagraphBreak

最后把三种运算放回同一映射中。若 \(F:M\to N\) 光滑，\(X\) 与 \(Y\) 满足 \(\mathrm{d}F_p(X_p)=Y_{F(p)}\)，则逐点代入第一输入可得
\[
	\iota_XF^*\omega=F^*(\iota_Y\omega).
\]
利用外微分与拉回相容，再用 Cartan 公式，便有
\[
	\mathcal L_XF^*\omega
	=(\mathrm{d}\iota_X+\iota_X\mathrm{d})F^*\omega
	=F^*(\mathrm{d}\iota_Y+\iota_Y\mathrm{d})\omega
	=F^*\mathcal L_Y\omega.
\]
这重新得到第13章对相关向量场证明的 Lie 导数自然性，并使边界处的局部版本同样清楚。外微分无需选择向量场，因而它与任意光滑映射拉回相容；插入和 Lie 导数各自用到了一个向量场，比较两个流形上的运算时就必须同时满足上述相关条件。
% 纲要条目排为编号小节标题，并进入目录与书签；说明段仍保留于纲要注释中。

% !TeX root = ../../Book1.tex
% 纲要标题：### 14.2 同伦公式与 Poincaré 引理
% 纲要位置：几何学计划纲要.md，第 655 行。
% 正文状态：初稿；尚未正式审核。

% 纲要内容开始
% 
% * 区间积分和显式同伦算子
% * 星形区域的 Poincaré 引理
% * 闭形式的局部恰当性
% * 圆周与穿孔空间上的周期障碍
% 
% 纲要内容结束

\section{同伦公式与 Poincaré 引理}
\label{g1:sec:1402}

上一节已经证明 \(\mathrm{d}^2=0\)，所以一个形式若是另一个形式的外微分，它自己的外微分必为零。逆向问题更有内容：若 \(\mathrm{d}\omega=0\)，能否找到 \(\eta\) 使 \(\omega=\mathrm{d}\eta\)？一元微积分用沿区间积分构造原函数；多维情形也可以沿一族道路积分，但还要记录道路上的移动方向怎样插入形式。这个构造既会给出局部原形式，也会说明整体原形式可能遇到什么障碍。

\ProseParagraphBreak

本节采用\cref{g1:def:closed-exact-form}的术语：\(\mathrm{d}\omega=0\) 表示闭，\(\omega=\mathrm{d}\eta\) 表示恰当。流形仍为有限维、Hausdorff、第二可数，允许边界，边界处的光滑性采用\cref{g1:def:manifold-with-boundary}的局部延拓约定。所用的外微分性质来自\cref{g1:thm:exterior-derivative,g1:prop:exterior-pullback}，积分则先只是一元参数积分。

\subsection{区间积分和显式同伦算子}
\label{g1:subsec:1402:01}

令 \(I=[0,1]\)，\(i_t:M\to I\times M\) 为 \(i_t(p)=(t,p)\)。在 \(I\times M\) 上说一个形式光滑，是指其乘积坐标 \((t,x)\) 中的系数在每点附近都可延到欧氏开集上光滑；这包括时间端点及 \(M\) 的边界同时出现的点。外微分按相同的坐标公式计算。坐标相容性与拉回相容性在 \((0,1)\times\operatorname{Int}M\) 上已经成立，各系数及导数连续，因而延至端点和边界。这样就能在整个参数区间上使用同一套公式。

\ProseParagraphBreak

固定 \(k\ge1\)。任意 \(\omega\in\Omega^k(I\times M)\) 都唯一分解为
\begin{equation}
	\label{g1:eq:interval-form-splitting}
	\omega=\mathrm{d}t\wedge\alpha_t+\beta_t,
\end{equation}
其中 \(\alpha_t\in\Omega^{k-1}(M)\)、\(\beta_t\in\Omega^k(M)\) 是光滑参数族；右边将它们视为不含 \(\mathrm{d}t\) 的乘积形式。这里始终将 \(\mathrm{d}t\) 放在最前面。事实上，坐标外幂基中的每一项要么含一个 \(\mathrm{d}t\)，要么不含；含有的那一项可用交替性移到最前面。两部分也有不依赖坐标的描述：
\[
	\alpha_t=i_t^*(\iota_{\partial_t}\omega),\qquad
	\beta_t=i_t^*\omega.
\]
\cref{g1:prop:interior-product-rule}中的插入固定第一位置，恰好保证这里没有额外符号。零次形式只有 \(\beta_t=f_t\)，没有 \(\alpha_t\) 部分。

\ProseParagraphBreak

现在把时间方向上的分量累加起来。约定 \(\Omega^{-1}(M)=\{0\}\)，定义次数降低 \(1\) 的\term[tonglunsuanzi]{同伦算子}[homotopy operator]
\begin{equation}
	\label{g1:eq:homotopy-operator}
	K:\Omega^k(I\times M)\longrightarrow\Omega^{k-1}(M),\qquad
	K\omega=\int_0^1\alpha_t\,\mathrm{d}t
	=\int_0^1 i_t^*(\iota_{\partial_t}\omega)\,\mathrm{d}t
	\quad(k\ge1),
\end{equation}
并在 \(k=0\) 时规定 \(Kf=0\)。这个积分的具体含义是：固定 \(p\) 和 \(v_1,\ldots,v_{k-1}\in \mathrm{T}_pM\)，令
\[
	(K\omega)_p(v_1,\ldots,v_{k-1})
	=\int_0^1(\alpha_t)_p(v_1,\ldots,v_{k-1})\,\mathrm{d}t.
\]
被积函数是实值函数，全部输入始终属于固定的 \(\mathrm{T}_pM\)；因此不需要在不同切空间之间作额外识别。

\ProseParagraphBreak

这个逐点构造确为光滑形式。若 \(\alpha_t=\sum_J a_J(t,x)\,\mathrm{d}x^J\)，它的系数是 \(\int_0^1a_J(t,x)\,\mathrm{d}t\)。在任意固定点的小邻域内，各阶偏导数在紧时间区间上连续；差商对 \(t\) 一致收敛到相应偏导数，所以可逐阶将对 \(x\) 的微分移入积分。若该点在边界，先用紧区间的有限延拓片覆盖，再用只依赖 \(t\) 的一元光滑截断函数归一化、拼合系数的局部延拓；所需截断函数由\cref{g1:lem:smooth-bump-function}给出。有限次拼合后得到共同的欧氏 \(x\) 邻域上的延拓，故积分在边界也光滑。换 \(M\) 的坐标时，对固定 \(x\) 的形式分量施加一个与 \(t\) 无关的线性变换，它与积分交换，所以这些局部系数表示同一个形式。

\begin{theorem}[同伦公式]
	\label{g1:thm:homotopy-formula}
	设 \(M\) 为光滑流形，可带边界，\(I=[0,1]\)，参数积上的光滑形式采用上述局部延拓约定。对每个 \(k\ge0\) 及 \(\omega\in\Omega^k(I\times M)\)，有
	\begin{equation}
		\label{g1:eq:interval-homotopy-formula}
		\mathrm{d}K\omega+K\mathrm{d}\omega=i_1^*\omega-i_0^*\omega.
	\end{equation}
	若 \(N\) 也是光滑流形，可带边界，\(H:I\times M\to N\) 是光滑同伦，\(F_t(p)=H(t,p)\)，令 \(K_H=K\circ H^*\)，则
	\begin{equation}
		\label{g1:eq:map-homotopy-formula}
		F_1^*\eta-F_0^*\eta=\mathrm{d}K_H\eta+K_H\mathrm{d}\eta
		\qquad(\eta\in\Omega^k(N)).
	\end{equation}
	这里光滑同伦要求参数与基点联合光滑，包括参数端点处的局部光滑延拓。
\end{theorem}
\begin{proof}
	为分开两种微分，暂用 \(\mathrm{d}_M\) 表示只对 \(M\) 变量作外微分。对不含 \(\mathrm{d}t\) 的参数族 \(\beta_t\)，坐标公式给出
	\[
		\mathrm{d}\beta_t=\mathrm{d}t\wedge\partial_t\beta_t+\mathrm{d}_M\beta_t.
	\]
	对分解\eqref{g1:eq:interval-form-splitting}中的另一项，外微分的分次乘法法则给出
	\[
		\mathrm{d}(\mathrm{d}t\wedge\alpha_t)
		=-\mathrm{d}t\wedge\mathrm{d}\alpha_t
		=-\mathrm{d}t\wedge\mathrm{d}_M\alpha_t.
	\]
	第二步消去了含两个 \(\mathrm{d}t\) 的项。因此
	\begin{equation}
		\label{g1:eq:interval-exterior-derivative}
		\mathrm{d}\omega
		=\mathrm{d}t\wedge(\partial_t\beta_t-\mathrm{d}_M\alpha_t)
		+\mathrm{d}_M\beta_t.
	\end{equation}
	这里的负号来自一次形式 \(\mathrm{d}t\) 的次数，正是稍后抵消所需的符号。

	对 \(k\ge1\)，将上述分解代入 \(K\)，再使用微分与参数积分的相容性，得到
	\[
		\begin{aligned}
			\mathrm{d}_MK\omega+K\mathrm{d}\omega
			&=\int_0^1\mathrm{d}_M\alpha_t\,\mathrm{d}t
			+\int_0^1(\partial_t\beta_t-\mathrm{d}_M\alpha_t)\,\mathrm{d}t\\
			&=\beta_1-\beta_0
			=i_1^*\omega-i_0^*\omega.
		\end{aligned}
	\]
	最后一步只是对每个系数使用一元微积分基本定理。若 \(k=0\)，写 \(\omega=f(t,p)\)，则 \(Kf=0\)，而
	\[
		\mathrm{d}f=(\partial_tf)\,\mathrm{d}t+\mathrm{d}_Mf,\qquad
		K\mathrm{d}f=\int_0^1\partial_tf\,\mathrm{d}t=f(1,\cdot)-f(0,\cdot).
	\]
	所以零次情形同样满足\eqref{g1:eq:interval-homotopy-formula}。

	对一般同伦，把已经证明的公式用于 \(H^*\eta\)。由\cref{g1:prop:exterior-pullback}，\(\mathrm{d}(H^*\eta)=H^*(\mathrm{d}\eta)\)；又有 \(i_t^*H^*\eta=F_t^*\eta\)，故
	\[
		F_1^*\eta-F_0^*\eta
		=\mathrm{d}(KH^*\eta)+KH^*(\mathrm{d}\eta),
	\]
	即所需结论。参数积存在边界交会处时，上述等式先在内点的乘积坐标中成立，再由系数及导数的连续性成立于整个参数积。
\end{proof}

一般同伦的算子也有直接的求值公式。对 \(k\ge1\)，
\begin{equation}
	\label{g1:eq:homotopy-map-operator}
	\begin{split}
		(K_H\eta)_p(v_1,\ldots,v_{k-1})
		=\int_0^1\eta_{H(t,p)}
		\bigl(\partial_tH(t,p),\,
		\mathrm{d}(F_t)_p v_1,\ldots,\mathrm{d}(F_t)_p v_{k-1}\bigr)\,\mathrm{d}t.
	\end{split}
\end{equation}
它把移动方向放在第一个输入，其余输入随 \(F_t\) 的微分送到当前点。这个公式不要求 \(F_t\) 是微分同胚，也不要求同伦来自某个流；在收缩同伦的起点，微分甚至可以为零。

\ProseParagraphBreak

若 \(\eta\) 是正次数闭形式，同伦公式化为 \(F_1^*\eta-F_0^*\eta=\mathrm{d}(K_H\eta)\)。因此同伦期间拉回的变化有一个明确的原形式。若 \(\eta\) 是闭函数，则 \(K_H\eta=0\) 且 \(\mathrm{d}\eta=0\)，公式给出 \(F_1^*\eta=F_0^*\eta\)。次数的区别在收缩到一点时尤其重要。

\subsection{星形区域的 Poincaré 引理}
\label{g1:subsec:1402:02}

设 \(U\subset\mathbb R^n\) 是含原点的开集，并且每条从原点到 \(x\in U\) 的线段都在 \(U\) 中，即
\[
	tx\in U\qquad(x\in U,\ 0\le t\le1).
\]
称 \(U\) 关于原点星形。这一条件正好保证 \(H(t,x)=tx\) 是在 \(U\) 内定义的光滑同伦，连接常值映射 \(c_0(x)=0\) 与恒等映射。若星形中心是 \(a\)，先用平移 \(x\mapsto x-a\) 将它化到原点即可。

\begin{theorem}[Poincaré 引理]
	\label{g1:thm:poincare-lemma}
	设 \(U\subset\mathbb R^n\) 为关于原点星形的开集。对 \(k\ge1\)，定义
	\begin{equation}
		\label{g1:eq:radial-homotopy-operator}
		(R\omega)_x(v_1,\ldots,v_{k-1})
		=\int_0^1 t^{k-1}\omega_{tx}(x,v_1,\ldots,v_{k-1})\,\mathrm{d}t,
	\end{equation}
	其中以通常的平移识别各欧氏切空间，并在零次形式上规定 \(R=0\)。则 \(R\omega\) 是光滑形式，且在正次数上满足
	\[
		\mathrm{d}R\omega+R\mathrm{d}\omega=\omega.
	\]
	所以每个正次数闭形式都恰当，可以取原形式 \(R\omega\)。对光滑函数 \(f\)，相应公式为
	\[
		R\mathrm{d}f=f-f(0);
	\]
	特别地，\(\mathrm{d}f=0\) 时 \(f\) 在 \(U\) 上为常数。
\end{theorem}
\begin{proof}
	对径向同伦 \(H(t,x)=tx\)，有
	\[
		\partial_tH(t,x)=x,\qquad \mathrm{d}(F_t)_x(v)=tv.
	\]
	代入\eqref{g1:eq:homotopy-map-operator}，由 \(\omega\) 对各输入多线性，便得到 \(K_H=R\) 及系数 \(t^{k-1}\)：除了第一个移动方向 \(x\)，其余 \(k-1\) 个输入各产生一个 \(t\)。在 \(k\ge1\) 时这个幂的指数非负，被积式在 \(t=0\) 也光滑；光滑参数积分已保证 \(R\omega\) 光滑。

	同伦公式给出 \(\mathrm{d}R+R\mathrm{d}=\operatorname{id}^*-c_0^*\)。若形式次数为正，常值映射的微分为零，所以 \(c_0^*\omega=0\)，得到所列等式。若又有 \(\mathrm{d}\omega=0\)，右端便是 \(\mathrm{d}(R\omega)\)。零次时，常值映射的拉回是常函数 \(f(0)\)，而 \(Rf=0\)，所以 \(R\mathrm{d}f=f-f(0)\)。闭函数的结论随即得到。
\end{proof}

\foreignidx{Poincaré lemma}
这里的径向公式只对正次数写成 \(t^{k-1}\) 的形式；不能在 \(k=0\) 时机械地写入 \(t^{-1}\)。零次闭形式是常数，而按照 \(\Omega^{-1}=0\) 的约定，零次恰当形式只有零函数。因此“闭则恰当”必须保留正次数条件。

\begin{example}[算出一个原形式]
	\label{g1:exm:radial-primitive-computation}
	在 \(\mathbb R^2\) 上考虑
	\[
		\omega=(2xy+y^2)\,\mathrm{d}x+(x^2+2xy)\,\mathrm{d}y.
	\]
	两个交叉偏导都等于 \(2x+2y\)，所以 \(\mathrm{d}\omega=0\)。径向公式对一次形式给出
	\[
		\begin{aligned}
			R\omega(x,y)
			&=\int_0^1\bigl((2t^2xy+t^2y^2)x
			+(t^2x^2+2t^2xy)y\bigr)\,\mathrm{d}t\\
			&=x^2y+xy^2.
		\end{aligned}
	\]
	于是 \(\omega=\mathrm{d}(x^2y+xy^2)\)，而不必先猜测原函数。对常二次形式 \(\eta=\mathrm{d}x\wedge\mathrm{d}y\)，同一个算子给出
	\[
		R\eta=\frac12(x\,\mathrm{d}y-y\,\mathrm{d}x),\qquad
		\mathrm{d}(R\eta)=\mathrm{d}x\wedge\mathrm{d}y.
	\]
	系数 \(1/2\) 来自 \(\int_0^1t\,\mathrm{d}t\)，也与对右端两项分别求外微分时产生的两份面积形式相符。
\end{example}

星形条件保证了这一个具体构造在整个区域中有定义。换一个星形中心，得到的原形式可能改变，但其外微分相同，故两者之差为闭形式。对一次形式，这意味着两个原函数在连通区域上只差常数；对更高次数，原形式的自由度仍是一个低一次的闭形式。

\subsection{闭形式的局部恰当性}
\label{g1:subsec:1402:03}

流形未必有一个整体星形坐标域，但每一点附近都有足够小的球形坐标域。边界点附近则用半球；从边界中心到半球内各点的线段仍留在半球内。于是径向证明同样能在这里使用。

\begin{corollary}[闭形式的局部原形式]
	\label{g1:cor:local-exactness}
	设 \(M\) 为光滑 \(n\) 维流形，可带边界，\(\omega\in\Omega^k(M)\) 满足 \(\mathrm{d}\omega=0\)。
	\begin{enumerate}
		\item 若 \(k\ge1\)，则对每个 \(p\in M\)，存在相对开邻域 \(U\) 及光滑 \((k-1)\) 次形式 \(\eta\in\Omega^{k-1}(U)\)，使 \(\omega|_U=\mathrm{d}\eta\)。若 \(p\) 在边界，\(\eta\) 光滑至边界。
		\item 若 \(k=0\)，则 \(\omega\) 局部为常数；在每个连通分支上为常数。
	\end{enumerate}
\end{corollary}
\begin{proof}
	先设 \(p\) 为内点。取以 \(p\) 为坐标原点的图，再缩小，使图像是一个开球 \(B_\varepsilon(0)\)。记这张缩小后的图为 \(x:U\to B_\varepsilon(0)\)。坐标形式 \(\widehat\omega=(x^{-1})^*(\omega|_U)\) 闭，因为外微分与拉回交换。开球星形，Poincaré 引理给出 \(\widehat\omega=\mathrm{d}(R\widehat\omega)\)。再拉回，令 \(\eta=x^*(R\widehat\omega)\)，便有 \(\mathrm{d}\eta=\omega|_U\)。

	若 \(p\) 在边界，取半空间坐标，再沿边界方向平移，使 \(x(p)=0\)。缩小后的图像可取
	\[
		B_\varepsilon(0)\cap\mathbb H^n,\qquad
		\mathbb H^n=\{x^n\ge0\}.
	\]
	这是半空间中的相对开集，且对 \(0\le t\le1\) 有 \(tx\) 仍在该半球中。径向同伦在参数积上光滑至所有边界，故\cref{g1:thm:homotopy-formula}适用；其移动方向和空间微分仍分别为 \(x\) 与 \(tI\)，所以相同的径向公式给出光滑至边界的原形式。按前段拉回即可。若 \(k>n\)，形式本来为零，也可直接取零原形式。

	零次情形在这些球或半球上满足 \(f(x)-f(0)=R\mathrm{d}f=0\)，所以局部常值。一个局部常值函数的任一非空取值纤维是开集，其补集也是其他取值纤维的并，因而开。连通分支不能分成两个不交的非空开集，所以函数在其中为常数。
\end{proof}

局部结论没有自动给出整体原形式。若 \(\omega=\mathrm{d}u_i\) 在各开集 \(U_i\) 上成立，则在交叠处 \(\mathrm{d}(u_i-u_j)=0\)。对一次形式，差 \(u_i-u_j\) 在交叠的每个连通分支上为常数，但不同分支的常数可能不同。要把所有 \(u_i\) 拼成一个函数，还必须使这些常数同时相容。

\ProseParagraphBreak

即便用单位分解 \(\{\rho_i\}\) 把局部原形式加权，也会出现新的项：
\[
	\mathrm{d}\left(\sum_i\rho_i\eta_i\right)
	=\omega+\sum_i\mathrm{d}\rho_i\wedge\eta_i,
	\qquad \mathrm{d}\eta_i=\omega|_{U_i}.
\]
这里各 \(\rho_i\eta_i\) 按从属支集延零，求和局部有限，因此公式合法；问题在于最后那一项一般不为零。单位分解保证了光滑拼合，却不能单凭这一步保证外微分方程仍然成立。下面的角形式将给出一个可以直接计算的整体障碍。

\subsection{圆周与穿孔空间上的周期障碍}
\label{g1:subsec:1402:04}

先明确本小节所需的积分。若 \(\gamma:[a,b]\to M\) 为光滑路径，\(\lambda\) 为一形式，则其拉回必可写成
\[
	\gamma^*\lambda=g(t)\,\mathrm{d}t,\qquad
	g(t)=\lambda_{\gamma(t)}(\gamma'(t)).
\]
定义
\[
	\int_\gamma\lambda=\int_a^b g(t)\,\mathrm{d}t.
\]
这只是实函数的一元积分。对分片光滑路径分段相加；对保持方向的光滑重参数化，链式法则和一元换元公式保证积分不变，反向参数化使其变号。若 \(\gamma(a)=\gamma(b)\)，称这个值为沿该闭路径的周期。

\begin{lemma}[恰当一形式的路径积分]
	\label{g1:lem:exact-path-period}
	设 \(M\) 为光滑流形，可带边界，\(u\in C^\infty(M)\)，\(\gamma:[a,b]\to M\) 为分片光滑路径。则
	\[
		\int_\gamma\mathrm{d}u=u(\gamma(b))-u(\gamma(a)).
	\]
	特别地，恰当一形式沿每条闭路径的周期为零。
\end{lemma}
\begin{proof}
	在每个光滑片段上，由链式法则，
	\[
		(\mathrm{d}u)_{\gamma(t)}(\gamma'(t))
		=\frac{\mathrm{d}}{\mathrm{d}t}(u\circ\gamma)(t).
	\]
	微积分基本定理给出该段的端点差。将有限个片段相加，中间端点值两两抵消，得到整条路径的端点差。闭路径的两个端点相同，所以差为零。
\end{proof}

\begin{example}[角形式的非零周期]
	\label{g1:exm:angular-period}
	在穿孔平面 \(P=\mathbb R^2\setminus\{0\}\) 上定义
	\[
		\vartheta=\frac{-y\,\mathrm{d}x+x\,\mathrm{d}y}{x^2+y^2}.
	\]
	它处处光滑。令 \(r^2=x^2+y^2\)，则
	\[
		\partial_x\left(\frac{x}{r^2}\right)
		=\frac{y^2-x^2}{r^4}
		=\partial_y\left(\frac{-y}{r^2}\right),
	\]
	所以一形式的坐标外微分公式给出 \(\mathrm{d}\vartheta=0\)。取半径 \(R>0\) 的圆周路径
	\[
		\gamma_R(t)=(R\cos t,R\sin t),\qquad 0\le t\le2\pi.
	\]
	直接代入可得
	\[
		\gamma_R^*\vartheta
		=\frac{R^2\sin^2t+R^2\cos^2t}{R^2}\,\mathrm{d}t
		=\mathrm{d}t,\qquad
		\int_{\gamma_R}\vartheta=2\pi.
	\]
	若 \(\vartheta=\mathrm{d}u\) 有整体原函数，上一引理就要求这个周期为零，矛盾。因此 \(\vartheta\) 闭而不恰当。

	在一个能选取极角 \(\theta\) 的角扇区中，代入 \(x=r\cos\theta\)、\(y=r\sin\theta\) 则得到 \(\vartheta=\mathrm{d}\theta\)。例如在 \(x>0\) 的半平面可以取 \(\theta=\arctan(y/x)\)。局部角函数存在，绕原点一圈却增加 \(2\pi\)，不能成为穿孔平面上单值的实函数。这里的局部表达不与非恰当性冲突。
\end{example}

\input{Figures/Chapter14-RadialAndPeriod.tex}

\cref{g1:fig:radial-and-period}把两种构造放在一起：星形区域中径向同伦始终留在定义域内，而穿孔平面中不能把终点设为被删除的原点。穿孔平面上的非恰当性由上面的周期计算证明，不能仅凭一个选定的径向构造失效便作判断。

\ProseParagraphBreak

将角形式限制到单位圆周也得到一个反例。以下取 \(S^1=\{(x,y):x^2+y^2=1\}\)，令 \(j:S^1\hookrightarrow P\) 为包含映射，并记
\[
	\sigma=j^*\vartheta,\qquad
	q:\mathbb R\longrightarrow S^1,\quad q(t)=(\cos t,\sin t).
\]
于是 \(q^*\sigma=\mathrm{d}t\)，沿正向一周的周期为 \(2\pi\)，故 \(\sigma\) 在圆周上不恰当。圆周上的每个一形式都闭，因为二次形式空间为零；所以在一维流形上，闭性本身也不能消除绕圈后的整体障碍。

\begin{proposition}[圆周上的周期判据]
	\label{g1:prop:circle-exactness-period}
	设 \(\lambda\) 为单位圆周 \(S^1\) 上的光滑一形式，写
	\[
		q^*\lambda=a(t)\,\mathrm{d}t,\qquad
		\mathcal P(\lambda)=\int_0^{2\pi}a(t)\,\mathrm{d}t.
	\]
	则 \(a\) 是光滑的 \(2\pi\) 周期函数，且 \(\lambda\) 恰当当且仅当 \(\mathcal P(\lambda)=0\)。更一般地，每个这样的形式都可写成
	\[
		\lambda=c\sigma+\mathrm{d}u,\qquad
		c=\frac{\mathcal P(\lambda)}{2\pi},
	\]
	其中 \(c\) 唯一，光滑函数 \(u:S^1\to\mathbb R\) 在相差常数的意义下唯一。
\end{proposition}
\begin{proof}
	由于 \(q(t+2\pi)=q(t)\) 且 \(q'(t+2\pi)=q'(t)\)，系数 \(a(t)=\lambda_{q(t)}(q'(t))\) 光滑并具有所述周期。若 \(\lambda=\mathrm{d}u\)，沿 \(q|_{[0,2\pi]}\) 使用\cref{g1:lem:exact-path-period}，得到 \(\mathcal P(\lambda)=0\)。

	反过来，设周期为零，令
	\[
		A(t)=\int_0^t a(s)\,\mathrm{d}s.
	\]
	周期性给出 \(A(t+2\pi)-A(t)=\mathcal P(\lambda)=0\)，所以 \(A\) 也为 \(2\pi\) 周期函数。于是可以定义 \(u(q(t))=A(t)\)：两个实数给出同一圆周点时相差 \(2\pi\) 的整数倍，故这个值与代表无关。

	还须检查下降后的函数光滑。\(q'(t)\) 从不为零，由\cref{g1:thm:inverse-function}，\(q\) 在每点附近到一个圆周小弧是微分同胚；在这个小弧上，\(u\) 就是光滑函数 \(A\) 与光滑局部角坐标的复合。因此 \(u\) 光滑。又有
	\[
		q^*(\mathrm{d}u)=\mathrm{d}(u\circ q)=\mathrm{d}A
		=a(t)\,\mathrm{d}t=q^*\lambda.
	\]
	\(q\) 满射且局部微分同胚，每个圆周切向量都能在某个实数点由其微分得到，所以拉回相等蕴含 \(\mathrm{d}u=\lambda\)。

	对一般 \(\lambda\)，取 \(c=\mathcal P(\lambda)/(2\pi)\)，则 \(\lambda-c\sigma\) 的周期为零，刚才的结论给出所需 \(u\)。若存在另一种分解，取周期便知常系数同为 \(c\)；随后两个原函数之差的微分为零。圆周连通，局部常值结论使它们只差常数。
\end{proof}

例如 \(\sigma\) 没有整体原函数，但 \(x\sigma\) 有：在角参数中，它的拉回为 \(\cos t\,\mathrm{d}t=\mathrm{d}(\sin t)\)，所以 \(x\sigma=\mathrm{d}(y|_{S^1})\)。而 \((1+x)\sigma\) 的周期仍为 \(2\pi\)，不恰当。这里决定能否取整体原函数的是一周累积的周期，而不是形式在某个局部坐标中的系数是否为零。
% 纲要条目排为编号小节标题，并进入目录与书签；说明段仍保留于纲要注释中。

% !TeX root = ../../Book1.tex
% 纲要标题：### 14.3 定向、密度与积分
% 纲要位置：几何学计划纲要.md，第 662 行。
% 正文状态：初稿；尚未正式审核。

% 纲要内容开始
% 
% * 定向线和相容图册
% * 密度及不可定向流形的积分
% * 紧支顶次形式的积分
% * 坐标选择和单位分解无关性
% * 定向映射的换元公式
% 
% 纲要内容结束

\section{定向、密度与积分}
\label{g1:sec:1403}

一个局部公式 \(a(x)\,\mathrm{d}x^1\wedge\cdots\wedge\mathrm{d}x^n\) 很像多元积分中的被积表达式，但换坐标时它带的是行列式，而通常的体积换元带绝对行列式。两者之间的符号差，正是定向要处理的问题。密度保留体积缩放而不保留方向，因此在不可定向流形上也能积分。本节先分清这两种对象，再用单位分解把局部积分拼成整体积分。

\ProseParagraphBreak

本节流形均有限维、Hausdorff、第二可数，可以带边界。积分先对光滑紧支对象定义；非紧流形上的非紧支积分还要另查收敛性。分析方面使用一元微积分基本定理、绝对可积函数的 Fubini 定理，以及下面的欧氏换元定理：若 \(F:U\to V\) 是欧氏开集之间的 \(C^1\) 微分同胚，\(h\) 在 \(V\) 上 Lebesgue 可积，则 \((h\circ F)|\det \mathrm{D}F|\) 在 \(U\) 上可积，且
\begin{equation}\label{g1:eq:euclidean-integration-input}
	\int_U h(F(x))|\det \mathrm{D}F_x|\,\mathrm{d}\lambda(x)=\int_V h(y)\,\mathrm{d}\lambda(y).
\end{equation}
这里 \(\lambda\) 是欧氏体积测度，\(\mathrm{d}\lambda\) 没有定向符号。这些是所调用的多元分析基础，分析背景的阅读位置为\cite[Appendix C, Multiple Integrals, pp. 649--655]{Lee2012Smooth}。在半空间图中，对内部应用这些结果；边界超平面是零测集，因而不影响顶维积分。紧支光滑系数有界，其支集位于有限个有界坐标片内，保证下文所用积分绝对收敛。

\subsection{定向线和相容图册}
\label{g1:subsec:1403:01}

回顾\cref{g1:def:determinant-density-line}：当 \(n\geq1\) 时，有序基之间按换基行列式为正划分为两个定向类。在流形上逐点任选一个类还不够，选择必须在局部保持一致。沿小邻域连续移动一个正向标架时，不能突然把它改称反向。

\begin{definition}
	\label{g1:def:manifold-orientation}
	设 \(M\) 为 \(n\geq1\) 维光滑流形。若在每个 \(\mathrm{T}_pM\) 上选定一个定向，并且每点附近存在在各点都为正向的光滑标架，则称这个选择为 \(M\) 的\term[liuxingdingxiang]{定向}[orientation]。允许这样的选择时称 \(M\) 可定向，连同一个选择时称为定向流形。零维流形的定向约定为在每个点选一个符号 \(\varepsilon_p\in\{1,-1\}\)；它用于给有向端点的积分赋号。
\end{definition}

在坐标图 \(x\) 上，令 \(\varepsilon_x(p)\) 表示坐标标架 \((\partial_{x^1},\ldots,\partial_{x^n})\) 是否正向：正向时为 \(1\)，反向时为 \(-1\)。局部一致性等价于 \(\varepsilon_x\) 局部常值。我们通常把图域取为连通的，于是它是一个常数。在两图重叠处，链式法则给出
\begin{equation}\label{g1:eq:orientation-chart-sign}
	\varepsilon_y=\varepsilon_x\operatorname{sgn}\det\frac{\partial x}{\partial y}.
\end{equation}
在无边界情形，必要时反射一个坐标，便可使每张图都正向；于是定向等价于过渡映射的 Jacobian 行列式均为正的相容图册。对于边界模型 \(x^n\geq0\)，保留带符号的表达更方便。尤其在一维区间右端点，只能用向内增大的半直线坐标，不能既要求它正向又要求其取值非负；\(\varepsilon_x\) 会准确记录这个符号。

\ProseParagraphBreak

还可以把符号的变化组织成一个丛。在每张图上取一个实线生成元 \(o_x\)，在交叠处规定
\[
	o_y=\operatorname{sgn}\det\frac{\partial x}{\partial y}\,o_x.
\]
行列式乘法保证三重交叠的余循环条件，所以\cref{g1:thm:bundle-cocycle-gluing}的粘合构造给出一条实光滑线丛，称为\term[dingxiangxian]{定向线丛}[orientation line bundle]，记为 \(\mathfrak o_M\)。它只记录行列式的符号，与使用整个行列式的行列式线丛不同。定向等价于在此线丛中连续选择 \(\pm o_x\) 的一个：表达为 \(\varepsilon_xo_x\)，\eqref{g1:eq:orientation-chart-sign}恰好保证它在换图后不变。零维时取平凡实线丛，在各点选择其正或负单位元。

\begin{proposition}[定向与顶次形式]
	\label{g1:prop:orientation-nonvanishing-form}
	设 \(M\) 为上述光滑 \(n\) 维流形。\(M\) 可定向当且仅当存在处处非零的光滑 \(n\) 形式。给定定向后，可以选择在每个正向有序基上取正值的顶次形式；零维时这表示函数在每点的符号与所选定向一致。
\end{proposition}
\begin{proof}
	对 \(n\geq1\)，若 \(\nu\) 处处非零，就把满足 \(\nu_p(e_1,\ldots,e_n)>0\) 的基规定为正向。局部坐标中 \(\nu=a(x)\,\mathrm{d}x^1\wedge\cdots\wedge\mathrm{d}x^n\)，系数连续且不为零，其符号局部不变，所以确实给出定向。

	反向，\cref{g1:thm:bundle-metric-existence}在 \(\mathrm{T}M\) 上给出光滑正定度量。每张图内把正向标架作光滑正交归一化，正向正交标架之间的矩阵属于 \(\operatorname{SO}(n)\)，行列式为一。因此各正向正交标架的顶次对偶楔积在交叠处相同，粘成所需 \(\nu\)。这也是\cref{g1:prop:orientation-volume-reductions}用于切丛的具体情形。零维时取函数 \(\nu(p)=\varepsilon_p\)；零维流形离散，故它光滑且处处非零。
\end{proof}

例如开 Möbius 带可写成 \(\mathbb R\times(-1,1)\) 在 \((t,u)\sim(t+1,-u)\) 下的商。若它有非零顶次形式，拉回到条带必为 \(a(t,u)\,\mathrm{d}t\wedge\mathrm{d}u\)，且不变性要求
\[
	a(t+1,-u)=-a(t,u).
\]
在中线上，\(a(t+1,0)=-a(t,0)\)，连续性迫使它在一圈中经过零，矛盾。这里每个小矩形都能定向，阻碍来自绕过整个带后的符号变化。

\subsection{密度及不可定向流形的积分}
\label{g1:subsec:1403:02}

\cref{g1:def:determinant-density-line}把密度定义为有序基上的实值函数，换基时乘 \(|\det A|\)。逐点使用这一构造得到密度线丛 \(\mathcal D(M)=|\Lambda^n\mathrm{T}^*M|\)。在坐标图中，光滑密度写为
\[
	\mu=m_x(x)\,|\mathrm{d}x^1\cdots\mathrm{d}x^n|,
\]
换成 \(y\) 坐标后，系数满足
\begin{equation}\label{g1:eq:density-coordinate-change}
	m_y(y)=m_x(x(y))\left|\det\frac{\partial x}{\partial y}\right|.
\end{equation}
系数可以为负，密度线仍是实线性空间；正密度才要求 \(m_x>0\)。若一条光滑密度的支集紧且包含于图域 \(U\)，定义它在该图中的积分为
\[
	I_x(\mu)=\int_{x(U)}m_x(x)\,\mathrm{d}\lambda(x).
\]
由\eqref{g1:eq:density-coordinate-change}和欧氏换元公式，这个数在另一张包含支集的图中不变。这里恰好需要绝对行列式，所以不必指定两图是否同向。

\ProseParagraphBreak

若支集不能放在一张图里，选从属于坐标覆盖的非负光滑单位分解 \(\{\rho_i\}\)，并取 \(\operatorname{supp}\rho_i\subset U_i\)。设 \(\mu\) 紧支，则 \(\rho_i\mu\) 的支集是 \(U_i\) 中的紧集。局部有限性还使只有有限个 \(\rho_i\mu\) 不为零：用有限个局部有限邻域覆盖 \(\operatorname{supp}\mu\)，它们总共只遇到有限个支集。于是可以提出整体积分的定义
\begin{equation}\label{g1:eq:density-integral-definition}
	\int_M\mu=\sum_i I_{x_i}(\rho_i\mu).
\end{equation}
这项和式与单位分解无关的证明，将与顶次形式一起在\subsecref{g1:subsec:1403:04}给出；目前求和的各项和有限性都已经有定义。

\begin{example}[Möbius 带上的面积]
	\label{g1:exm:mobius-density-integral}
	在闭条带 \(\mathbb R\times[-1,1]\) 的上述商上，\(|\mathrm{d}t\,\mathrm{d}u|\) 被变换 \((t,u)\mapsto(t+1,-u)\) 保持，因此下降为正光滑密度。用 \(0\leq t\leq1\)、\(-1\leq u\leq1\) 作基本区域，其积分为 \(2\)。计算时可以把基本区域划成落在商坐标图中的小矩形；公共边以及识别的两条竖边都是二维零测集，不会被重复计入面积。带边界的这个 Möbius 带仍不可定向，因为前面的中线论证不变，却有良定义的面积。
\end{example}

密度解决了测量体积的问题，却不自行产生一个与外微分相容的带符号边界公式。后面的 Stokes 定理仍对定向流形上的微分形式陈述，不能因密度总能积分就删去定向条件。

\subsection{紧支顶次形式的积分}
\label{g1:subsec:1403:03}

现在给 \(M\) 选定定向。若 \(\omega\in\Omega_c^n(M)\) 的支集在一个连通图域 \(U\) 中，并写成
\[
	\omega=a_x(x)\,\mathrm{d}x^1\wedge\cdots\wedge\mathrm{d}x^n,
\]
定义
\begin{equation}\label{g1:eq:oriented-chart-integral}
	I_x^{\mathrm{or}}(\omega)=\varepsilon_x\int_{x(U)}a_x(x)\,\mathrm{d}\lambda(x).
\end{equation}
在正向图中 \(\varepsilon_x=1\)，得到熟悉的公式；反向图中要额外乘 \(-1\)。这不是对被积系数取绝对值，形式的积分可以为负，也可以因正负部分抵消而为零。

\ProseParagraphBreak

对一般紧支顶次形式，以从属于坐标覆盖的单位分解规定
\begin{equation}\label{g1:eq:top-form-integral-definition}
	\int_M\omega=\sum_i I_{x_i}^{\mathrm{or}}(\rho_i\omega).
\end{equation}
和式仍只有有限个非零项。零维流形上的紧集只有有限个点，此时 \(0\) 形式就是函数，定义专门写为
\[
	\int_M f=\sum_{p\in M}\varepsilon_p f(p),
\]
其中紧支性保证求和有限。特别地，一个正向点的积分为函数值，负向点的积分为其相反数。这将使区间的两个边界端点分别贡献 \(f(b)\) 与 \(-f(a)\)。

\ProseParagraphBreak

定向也给出顶次形式到密度的线性对应。令 \(S_o\omega\) 在 \(x\) 图中的密度系数为 \(\varepsilon_xa_x\)。因为
\[
	a_y=a_x\det\frac{\partial x}{\partial y},\qquad
	\varepsilon_y=\varepsilon_x\operatorname{sgn}\det\frac{\partial x}{\partial y},
\]
乘积 \(\varepsilon_ya_y\) 恰按绝对行列式变换，所以这些局部系数粘成密度。这与非线性的 \(|\omega|\) 不同：\(S_o\omega\) 保留相对于给定定向的正负号，\(|\omega|\) 只保留绝对大小。由局部定义，\(\int_M\omega=\int_M S_o\omega\)。反转定向时 \(S_o\) 变号，故形式积分也变号，密度积分本身则不变。

\subsection{坐标选择和单位分解无关性}
\label{g1:subsec:1403:04}

\begin{theorem}[流形积分的良定义性]
	\label{g1:thm:integral-well-defined}
	设 \(M\) 为有限维、Hausdorff、第二可数光滑流形，允许边界。\eqref{g1:eq:density-integral-definition}对每个紧支光滑密度定义一个与坐标覆盖和从属单位分解无关的实数。若另给定向，则\eqref{g1:eq:top-form-integral-definition}对每个紧支光滑顶次形式也有同样性质。两种积分均线性；非负且不恒为零的光滑密度积分严格为正。
\end{theorem}
\begin{proof}
	密度在一张图中的积分已经由欧氏换元证明与图无关。形式也如此：在两个图的交中，其系数相差有符号行列式，图的定向符号满足\eqref{g1:eq:orientation-chart-sign}。两者相乘得到绝对行列式，故
	\[
		\varepsilon_y\int a_y\,\mathrm{d}\lambda(y)
		=\varepsilon_x\int a_x(x(y))\left|\det\frac{\partial x}{\partial y}\right|\mathrm{d}\lambda(y)
		=\varepsilon_x\int a_x\,\mathrm{d}\lambda(x).
	\]
	若交集不连通，对各连通分支使用相同公式，或直接把局部常值符号保留在被积函数中。边界图的换元在内部进行，边界零测不改变结果。

	现取两套从属于坐标覆盖的单位分解 \(\{\rho_i\}\)、\(\{\sigma_j\}\)，图分别记为 \(x_i,y_j\)。令 \(u\) 表示待积的密度或顶次形式。只有有限个支集与 \(\operatorname{supp}u\) 相交，因此下面实为有限次拆分：
	\[
		\sum_i I_{x_i}(\rho_i u)
		=\sum_{i,j}I_{x_i}(\rho_i\sigma_j u)
		=\sum_{i,j}I_{y_j}(\rho_i\sigma_j u)
		=\sum_j I_{y_j}(\sigma_j u).
	\]
	形式情形中的 \(I\) 理解为 \(I^{\mathrm{or}}\)。中间每个对象的支集紧且包含在 \(U_i\cap V_j\) 中，所以第一步的单图换元确实适用。最外两次等式分别使用 \(\sum_j\sigma_j=1\)、\(\sum_i\rho_i=1\)，并利用欧氏积分的线性。由此同时排除了覆盖、坐标与单位分解的选择依赖。

	线性可以在同一套覆盖和单位分解中逐项验证。若 \(\mu\geq0\) 且某点处严格为正，连续性给出某张图内的一块区域，其系数有正下界。该点处至少一个 \(\rho_i>0\)，缩小区域后 \(\rho_i\mu\) 的系数仍有正下界，因此这一项积分正，其余项非负。边界点附近可取正体积的半球内部，零维时则是一项正数。这证明严格正性。
\end{proof}

证明中的共同细化 \(\rho_i\sigma_j\) 是关键：不必在两套坐标之间直接比较整个和式，只需把它们拆成同时落在两张图中的小块。局部有限性让这一分解在紧支集上化为有限运算，避免了任何未经说明的无穷级数换序。

\begin{proposition}[度量的体积积分]
	\label{g1:prop:metric-volume-integration}
	设 \(g\) 为 \(M^n\) 上的光滑正定度量。其正体积密度在坐标中为
	\[
		\mu_g=\sqrt{\det(g_{ij})}\,|\mathrm{d}x^1\cdots\mathrm{d}x^n|.
	\]
	若 \(M\) 定向，则相应的正体积形式为
	\[
		\operatorname{vol}_g=\varepsilon_x\sqrt{\det(g_{ij})}\,\mathrm{d}x^1\wedge\cdots\wedge\mathrm{d}x^n,
	\]
	并且对每个紧支光滑函数 \(f\)，有 \(\int_M f\operatorname{vol}_g=\int_M f\mu_g\)。
\end{proposition}
\begin{proof}
	\cref{g1:prop:gram-volume-density}在每个切空间给出 \(\mu_g(e)=\sqrt{\det[g(e_i,e_j)]}\)，取坐标标架就得到密度公式；正定性使行列式正，其正平方根随点光滑。把坐标标架的定向符号乘入顶次形式，便使其在正向标架上取正值，且换图时相容。最后 \(S_o(f\operatorname{vol}_g)\) 的局部系数为 \(f\sqrt{\det(g_{ij})}\)，正是 \(f\mu_g\)，所以积分相同。
\end{proof}

例如平面的极坐标在 \(r>0\) 的角度分支上给出 \(g=\mathrm{d}r^2+r^2\mathrm{d}\theta^2\)，所以体积密度为 \(r|\mathrm{d}r\,\mathrm{d}\theta|\)。\(r=0\) 不是这张坐标图的一点，不能把极坐标表达在原点的退化误认为欧氏度量退化。对半径为 \(R\) 的圆盘积分得到 \(\int_0^{2\pi}\int_0^R r\,\mathrm{d}r\,\mathrm{d}\theta=\pi R^2\)；被忽略的中心及切开的射线面积为零。

\subsection{定向映射的换元公式}
\label{g1:subsec:1403:05}

对任意光滑映射都能拉回微分形式，但积分的简单换元公式首先适用于微分同胚。非单射映射可能重复覆盖区域，微分降秩也可能压缩维数；这些情况不能仅凭拉回有定义就沿用同一个积分等式。

\begin{theorem}[定向换元]
	\label{g1:thm:oriented-change-variables}
	设 \(F:M^n\to N^n\) 为光滑流形之间的微分同胚，允许边界。
	\begin{enumerate}
		\item 对 \(N\) 上的紧支光滑密度 \(\mu\)，其微分同胚拉回满足 \(\int_M F^*\mu=\int_N\mu\)。
		\item 若两流形定向，定义 \(s_F(p)=1\) 当 \(\mathrm{d}F_p\) 保定向，反向时为 \(-1\)。则 \(s_F\) 局部常值，且对任意 \(\omega\in\Omega_c^n(N)\)，
		\[
			\int_M F^*\omega=\int_N(s_F\circ F^{-1})\omega.
		\]
		特别地，整体保定向时积分不变，整体反定向时积分变号。若 \(M\) 连通，符号只有上述两种整体可能。
	\end{enumerate}
	零维时，用源点和像点的定向符号之比定义 \(s_F\)，同一公式成立。
\end{theorem}
\begin{proof}
	微分同胚将紧集的原像变为紧集，所以两种拉回仍紧支。对目标支集选有限个所需坐标片及从属单位分解，再拉回到源；它们的拉回仍构成从属单位分解。故只需证明单图的结果。

	在源图 \(x\) 和目标图 \(y\) 中，记 \(G=y\circ F\circ x^{-1}\)。若 \(\mu=m(y)|\mathrm{d}y^1\cdots\mathrm{d}y^n|\)，则其拉回系数为 \(m(G(x))|\det \mathrm{D}G_x|\)，欧氏换元直接给出第一项。若 \(\omega=a(y)\,\mathrm{d}y^1\wedge\cdots\wedge\mathrm{d}y^n\)，则拉回系数为 \(a(G(x))\det \mathrm{D}G_x\)。源、目标图的定向符号分别为 \(\varepsilon_x,\varepsilon_y\)，有
	\[
		s_F=\varepsilon_x\varepsilon_y\operatorname{sgn}\det \mathrm{D}G,
	\]
	所以 \(\varepsilon_x\det \mathrm{D}G=s_F\varepsilon_y|\det \mathrm{D}G|\)。代入单图积分，再用欧氏换元，便得到第二项；逐片相加完成证明。符号的局部常值性来自非零行列式的连续性，在连通流形上因此恒定。零维公式就是把有限项求和按双射重新编号。
\end{proof}

以反射 \(F(x^1,x^2)=(-x^1,x^2)\) 为例，在标准定向的平面上，\(F^*(a\,\mathrm{d}x^1\wedge\mathrm{d}x^2)=-a(-x^1,x^2)\,\mathrm{d}x^1\wedge\mathrm{d}x^2\)，故紧支形式的积分变号；密度的拉回则没有这个负号。又如圆周映射 \(e^{\mathrm{i}\theta}\mapsto e^{2\mathrm{i}\theta}\) 不是微分同胚，它把角形式拉回为两倍，积分也变成两倍。后者涉及重复覆盖次数，将在\secref{g1:sec:1405}进一步说明，不能当作当前换元定理的一个无条件推论。

% 纲要条目排为编号小节标题，并进入目录与书签；说明段仍保留于纲要注释中。

% !TeX root = ../../Book1.tex
% 纲要标题：### 14.4 边界定向与 Stokes 的证明
% 纲要位置：几何学计划纲要.md，第 670 行。
% 正文状态：初稿；尚未正式审核。

% 纲要内容开始
% 
% * 外法向优先的边界约定
% * 半空间中的局部公式
% * 通过单位分解粘合
% * 带角区域及逐项边界符号
% * Green、Gauss 公式作为特例
% 
% 纲要内容结束

\section{边界定向与 Stokes 的证明}
\label{g1:sec:1404}

在一元微积分中，导数的积分只留下区间两个端点的函数值，而且右端点取正号、左端点取负号。Stokes 定理把这个关系推广到微分形式：区域内部的外微分积分，等于边界上原形式的积分。这里“边界”必须带有由区域定向确定的方向，否则等式连符号都不能确定。

\ProseParagraphBreak

第九章已经在曲面坐标中证明了所需的 Green 公式，并据此把联络形式的边界积分换成曲率积分，见\cref{g1:lem:green-formula,g1:thm:local-gauss-bonnet}。本节从同一个积分基本定理出发，先处理任意维半空间，再用第十二章的光滑单位分解得到整体公式。所有形式都光滑至边界；非紧流形上的形式要求紧支，支集可以接触实际边界。

\subsection{外法向优先的边界约定}
\label{g1:subsec:1404:01}

设 \(M^n\) 为有向带边界光滑流形，\(n\ge1\)，包含映射记为 \(j:\partial M\hookrightarrow M\)。\cref{g1:prop:boundary-chart-invariance}说明，\(\partial M\) 本身是 \(n-1\) 维无边界光滑流形；其完整切空间由边界坐标中最后一个分量为零的向量组成。若在边界图中把 \(M\) 写成 \(x^n\ge0\)，则指向外侧的横截向量满足 \(\mathrm{d}x^n(\nu)<0\)。这个描述只用边界两侧的区分，不需要长度或正交关系。

\begin{definition}[边界定向]\label{g1:def:boundary-orientation}
	设 \(M^n\) 为有向带边界光滑流形，\(n\ge1\)。在 \(p\in\partial M\) 取一个外向横截向量 \(\nu\)，即在边界定义坐标 \(\rho\ge0\) 下满足 \(\mathrm{d}\rho_p(\nu)<0\)。边界取\term[waifaxiangyouxian]{外法向优先}[outward-normal-first]的定向：当 \(n\ge2\) 时，基 \((v_1,\ldots,v_{n-1})\) 在 \(\mathrm{T}_p\partial M\) 中为正向，当且仅当
	\[
		(\nu,v_1,\ldots,v_{n-1})
	\]
	在 \(\mathrm{T}_pM\) 中为正向。等价地，若 \(\Omega_p\) 是表示 \(M\) 正定向的非零顶次协变向量，则 \((\iota_\nu\Omega_p)|_{\mathrm{T}_p\partial M}\) 表示边界的正定向。后一表述也适用于 \(n=1\)：它给边界点指定一个正负号。
\end{definition}

名称中的“法向”只指横截方向；若另有黎曼度量，才可以特选外向单位法向。先核对这个约定不依赖选择。两种边界定义坐标 \(\rho,\rho'\) 在边界上都为零，局部可写成
\[
	\rho'=\rho\,a,\qquad
	a(x',\rho)=\int_0^1\partial_\rho\rho'(x',t\rho)\,\mathrm{d}t.
\]
边界坐标变换可逆，且把内侧映到内侧，所以 \(a(p)>0\)：其边界处的法向导数不为零，又不能为负。这也可由\cref{g1:prop:boundary-smooth-derivative}的链式法则和一侧非负性直接判断。因此 \(\mathrm{d}\rho'_p=a(p)\mathrm{d}\rho_p\)，外向条件与坐标无关。

\ProseParagraphBreak

若 \(\nu,\nu'\) 都外向，则在一维商空间 \(\mathrm{T}_pM/\mathrm{T}_p\partial M\) 中，它们位于同一条半直线上，故
\[
	\nu'=c\nu+w,\qquad c>0,\quad w\in \mathrm{T}_p\partial M.
\]
把第一列由 \(\nu\) 换为 \(\nu'\)，行列式只乘以正数 \(c\)，因为 \(w\) 是其余切向列的线性组合。等价地，插入 \(w\) 后限制到边界切空间的顶次形式为零。因此边界定向与外向横截向量及正顶次形式的选择无关，且随点光滑变化。

在标准定向半空间
\[
	H^n=\{x^n\ge0\},\qquad
	\Omega=\mathrm{d}x^1\wedge\cdots\wedge\mathrm{d}x^n
\]
中，可以取 \(\nu=-\partial_{x^n}\)。于是
\begin{equation}\label{g1:eq:half-space-boundary-orientation}
	\left.\iota_{-\partial_{x^n}}\Omega\right|_{\partial H^n}
	=(-1)^n\,\mathrm{d}x^1\wedge\cdots\wedge\mathrm{d}x^{n-1}.
\end{equation}
所以边界坐标 \((x^1,\ldots,x^{n-1})\) 的通常定向，与诱导边界定向相差 \((-1)^n\)。这个符号来自把最后一个坐标向量移到第一位，再乘上外向向量的负号。

\ProseParagraphBreak

当 \(n=1\) 时，标准半直线的边界点 \(0\) 取负号；对按递增参数定向的区间 \([a,b]\)，左端点取负、右端点取正。按照上一节的有向零维积分约定，
\[
	\int_{\partial[a,b]}f=f(b)-f(a).
\]
零维流形另约定边界为空，不再应用“降低一维”的基向量描述。此时没有需要讨论的负一次形式；若记 \(\Omega^{-1}=0\)，Stokes 等式只表示 \(0=0\)。

\subsection{半空间中的局部公式}
\label{g1:subsec:1404:02}

高维公式的局部计算只多了一个排列符号。把一般 \(n-1\) 次形式按“缺少哪个坐标微分”分组，并在第 \(i\) 项预先放入 \((-1)^{i-1}\)，可以让外微分后的各项都具有相同正号。以下用帽子表示删去相应因子。

\begin{lemma}[半空间中的 Stokes 公式]\label{g1:lem:half-space-stokes}
	设 \(n\ge1\)，\(H^n=\{x^n\ge0\}\) 取标准定向，边界取外法向优先定向。若 \(\alpha\) 是 \(H^n\) 上光滑至边界的紧支 \(n-1\) 次形式，则
	\[
		\int_{H^n}\mathrm{d}\alpha=\int_{\partial H^n}j^*\alpha.
	\]
	在无边界模型 \(\mathbb R^n\) 上，任意紧支光滑 \(n-1\) 次形式都满足 \(\int_{\mathbb R^n}\mathrm{d}\alpha=0\)。紧支集允许与 \(x^n=0\) 相交。
\end{lemma}
\begin{proof}
	唯一写成
	\[
		\alpha=\sum_{i=1}^n(-1)^{i-1}a_i(x)\,
		\mathrm{d}x^1\wedge\cdots\wedge\widehat{\mathrm{d}x^i}\wedge\cdots\wedge\mathrm{d}x^n.
	\]
	每个 \(a_i\) 光滑且紧支。由\cref{g1:thm:exterior-derivative}的坐标公式，求外微分时，只有 \(\partial_i a_i\,\mathrm{d}x^i\) 那一项不因重复微分而消失。把 \(\mathrm{d}x^i\) 移到第 \(i\) 位又产生 \((-1)^{i-1}\)，两个符号抵消，故
	\begin{equation}\label{g1:eq:stokes-local-divergence}
		\mathrm{d}\alpha=\left(\sum_{i=1}^n\partial_i a_i\right)
		\mathrm{d}x^1\wedge\cdots\wedge\mathrm{d}x^n.
	\end{equation}

	取 \(R\) 足够大，使全部系数支集位于 \((-R,R)^{n-1}\times[0,R)\)。连续系数及其导数在这个紧盒上可积，累次积分的 Fubini 公式允许先沿任一坐标积分。对 \(i<n\)，沿 \(x^i\) 的积分区间两端，\(a_i\) 都为零，所以
	\[
		\int_{H^n}\partial_i a_i\,\mathrm{d}x=0.
	\]
	最后一个坐标则从零开始，积分基本定理给出
	\[
		\int_{H^n}\partial_n a_n\,\mathrm{d}x
		=-\int_{\mathbb R^{n-1}}a_n(x',0)\,\mathrm{d}x'.
	\]
	因此体积分只留下这一个负的边界值。

	在边界上拉回时，含 \(\mathrm{d}x^n\) 的项全部消失，故
	\[
		j^*\alpha=(-1)^{n-1}a_n(x',0)\,
		\mathrm{d}x^1\wedge\cdots\wedge\mathrm{d}x^{n-1}.
	\]
	但边界的定向符号还要乘上\eqref{g1:eq:half-space-boundary-orientation}中的 \((-1)^n\)。于是
	\[
		\int_{\partial H^n}j^*\alpha
		=(-1)^n(-1)^{n-1}\int_{\mathbb R^{n-1}}a_n(x',0)\,\mathrm{d}x'
		=-\int_{\mathbb R^{n-1}}a_n(x',0)\,\mathrm{d}x',
	\]
	恰与体积分相同。\(n=1\) 时，右侧按零维积分解释为 \(-a_1(0)\)，同样成立。

	在 \(\mathbb R^n\) 上，每个坐标都在完整实轴上积分，紧支性使各方向两端的系数为零，故所有偏导积分都为零。
\end{proof}

这里没有把 \(\alpha\) 越过实际边界 \(x^n=0\) 补成零。那样补零通常不光滑，而且会错误地抹掉边界项。紧支性只保证在足够远的辅助截断边缘处系数为零，真正的边界值 \(a_n(x',0)\) 完整保留。

\subsection{通过单位分解粘合}
\label{g1:subsec:1404:03}

\begin{theorem}[Stokes 定理]\label{g1:thm:stokes}
	设 \(M^n\) 为有限维、Hausdorff、第二可数的有向光滑流形，允许有边界，\(n\ge1\)。给 \(\partial M\) 外法向优先定向，记包含映射为 \(j\)。对每个紧支光滑 \(n-1\) 次形式 \(\omega\in\Omega_c^{n-1}(M)\)，有
	\begin{equation}\label{g1:eq:stokes}
		\int_M\mathrm{d}\omega=\int_{\partial M}j^*\omega.
	\end{equation}
	此式称为\term[stokesdingli]{Stokes 定理}[Stokes theorem]。不要求 \(M\) 紧；若 \(M\) 紧，则任意光滑 \(n-1\) 次形式自动满足支集条件。
\end{theorem}
\begin{proof}
	令 \(K=\operatorname{supp}\omega\)。在 \(M\setminus K\) 的每个点附近，\(\omega\) 恒为零，故 \(\operatorname{supp}\mathrm{d}\omega\subseteq K\)。由于边界是闭集，\(j^*\omega\) 的支集包含于紧集 \(K\cap\partial M\)。因此两边的积分均在上一节已定义的紧支积分范围内。

	用内部坐标图和边界坐标图覆盖 \(M\)，再由\cref{g1:thm:smooth-partition-unity}取从属光滑单位分解 \(\{\rho_i\}\)。局部有限的支集族只会有有限个成员与紧集 \(K\) 相交：给 \(K\) 中每一点选一个只遇到有限个支集的邻域，再从这些邻域取有限子覆盖即可。因此
	\[
		\omega=\sum_i\omega_i,\qquad \omega_i=\rho_i\omega
	\]
	实际上只有有限个非零项；每个 \(\omega_i\) 的支集是位于对应图域 \(U_i\) 内的紧集。

	将 \(\omega_i\) 写入坐标。其坐标支集紧含于相对开图域，故可越过这个图域的人为边缘补零到整个 \(\mathbb R^n\) 或 \(H^n\)，而保持光滑；补零没有越过 \(H^n\) 的实际边界。由\cref{g1:prop:exterior-pullback}，坐标拉回与外微分交换。对补零形式应用\cref{g1:lem:half-space-stokes}，再用\cref{g1:thm:oriented-change-variables}把积分搬回，便得
	\[
		\int_M\mathrm{d}\omega_i=\int_{\partial M}j^*\omega_i.
	\]
	这里任一图都可以带有定向符号 \(\varepsilon_i=\pm1\)：若坐标参数化反向，体积分和它诱导的边界积分同时乘以 \(-1\)，所以局部等式仍成立。特别是在一维右端点附近，无须强行要求到标准半直线的坐标正向。

	最后把有限个等式相加。注意局部公式作用于 \(\mathrm{d}(\rho_i\omega)\)，而不是直接作用于 \(\rho_i\mathrm{d}\omega\)；乘积法则产生的项正好消去：
	\[
		\begin{aligned}
			\sum_i\mathrm{d}(\rho_i\omega)
			&=\sum_i\mathrm{d}\rho_i\wedge\omega+\sum_i\rho_i\,\mathrm{d}\omega\\
			&=\mathrm{d}\!\left(\sum_i\rho_i\right)\wedge\omega+\mathrm{d}\omega
			=\mathrm{d}\omega.
		\end{aligned}
	\]
	这些等式可先按局部有限和理解，而乘上 \(\omega\) 后只有前述有限个非零项。边界上的和同样为 \(j^*\omega\)。积分的线性及\cref{g1:thm:integral-well-defined}的坐标、单位分解无关性遂给出\eqref{g1:eq:stokes}。
\end{proof}

若 \(M\) 无边界，结论是 \(\int_M\mathrm{d}\omega=0\)；若 \(M=[a,b]\)，结论就是积分基本定理。非紧情形的紧支要求控制了“远处的边界”：仅仅知道 \(\mathrm{d}\omega\) 的积分收敛，并不能保证无穷远没有留下额外项，习题中会给出一维例子。

\subsection{带角区域及逐项边界符号}
\label{g1:subsec:1404:04}

矩形和单形在角点处不具有半空间局部模型。这里先明确所讨论的对象，再把局部公式推广到正交象限；角棱不贡献新的积分项，但邻接的每张面都必须保留。

\begin{definition}[带角流形]\label{g1:def:manifold-with-corners}
	一个 \(n\) 维\term[daijiaoliuxing]{带角流形}[manifold with corners]是 Hausdorff、第二可数空间，带有到
	\[
		Q_k^n=[0,\infty)^k\times\mathbb R^{n-k},
		\qquad 0\le k\le n
	\]
	的相对开集的坐标图；换坐标及其逆都要求局部延拓为欧氏空间中的光滑映射。函数、微分形式和定向也按这种局部延拓定义。在一点附近，只保留在那里取零的受限坐标，其个数称为该点的深度；深度一部分是光滑的余维一边界，深度至少二的部分称为角棱。
\end{definition}

这个深度与坐标无关。像第十一章半空间的论证一样，局部延拓的导数由内部稠密性唯一确定，互逆变换的导数互逆。在深度 \(k\) 的点，单侧曲线的速度构成锥
\[
	C_p=\{v:v^1,\ldots,v^k\ge0\},
\]
其最大线性子空间 \(C_p\cap(-C_p)\) 的维数为 \(n-k\)。坐标变换的微分把这个速度锥同构到目标的速度锥，因此保持 \(k\)。深度一部分在小坐标邻域内分成 \(k\) 张局部面 \(x^i=0\) 的内部；换坐标保持深度，并把这些连通的局部面分支一一对应。每张面的定向先在深度一点由外法向优先确定，再连续延伸到它的角棱。

\ProseParagraphBreak

顶次积分沿用上一节的坐标定义。正交象限中边界在 \(n\) 维测度下为零，不影响体积分；对每张余维一面，角棱在该面内的 \(n-1\) 维测度也为零，因为它们在面坐标中是一个或多个坐标取零的集合。因此边界积分按每张局部面分别计算，同一角点邻接的不同面各计一次，而不把整个拓扑边界误认为一张光滑流形。

\ProseParagraphBreak

这些面上的积分与坐标选择相容：在深度一部分应用有向换元公式即可。若某积分的支集接触角棱，可以先去掉面坐标中宽度为 \(\varepsilon\) 的角棱条带，再令 \(\varepsilon\downarrow0\)。在所用紧坐标片内，形式系数与换元 Jacobian 有界，条带体积趋于零，遗漏积分的绝对值也趋于零。于是普通换元公式延续到接触角棱的积分。用有限坐标分解求和时，两种分解可通过权重乘积作共同细化，逐面利用换元和权重总和为 \(1\)，证明边界积分不依赖分解。以下的 \(\int_{\partial M}\omega\) 就表示这种按局部面分支计算的积分；面上的形式均指包含拉回。

\begin{theorem}[带角流形上的 Stokes 公式]\label{g1:thm:stokes-corners}
	设 \(M^n\) 为有向、Hausdorff、第二可数的光滑带角流形，\(n\ge1\)，\(\omega\) 为光滑至各面的紧支 \(n-1\) 次形式。各余维一面取外法向优先定向，边界积分按局部面分支计算，则
	\[
		\int_M\mathrm{d}\omega=\int_{\partial M}\omega.
	\]
	余维至少二的角棱不增加独立积分项。
\end{theorem}
\begin{proof}
	先在标准定向的 \(Q_k^n\) 上取紧支形式
	\[
		\alpha=\sum_{i=1}^n(-1)^{i-1}a_i\,
		\mathrm{d}x^1\wedge\cdots\wedge\widehat{\mathrm{d}x^i}\wedge\cdots\wedge\mathrm{d}x^n.
	\]
	外微分仍由\eqref{g1:eq:stokes-local-divergence}给出。对不受限坐标 \(i>k\)，偏导积分为零；对 \(i\le k\)，沿 \([0,\infty)\) 使用积分基本定理，得到
	\[
		\int_{Q_k^n}\mathrm{d}\alpha
		=-\sum_{i=1}^k\int_{F_i}a_i|_{x^i=0}\,\mathrm{d}x^{\widehat i},
		\qquad F_i=Q_k^n\cap\{x^i=0\}.
	\]
	这里右侧暂按其余坐标的递增次序使用通常正测度。面 \(F_i\) 的外向向量为 \(-\partial_{x^i}\)，所以相对于这组面坐标，诱导定向的符号为 \((-1)^i\)。拉回 \(\alpha\) 后只留下 \((-1)^{i-1}a_i\,\mathrm{d}x^{\widehat i}\)，两者相乘正好是 \(-a_i\,\mathrm{d}x^{\widehat i}\)。因此上式就是各有向面的积分之和。各坐标面的相交部分测度为零，不会在这个计算之外再留下角点项。

	现在令 \(K=\operatorname{supp}\omega\)。紧集附近只需有限个坐标截断函数。具体地，在每点的带角坐标中取两层相对球，把\cref{g1:lem:smooth-bump-function}中的欧氏截断函数限制到正交象限，得到支集紧含于外层图域、在内层为 \(1\) 的光滑函数。选有限个内层覆盖 \(K\)，记对应函数为 \(b_1,\ldots,b_N\)。其和 \(b\) 在 \(K\) 的开邻域 \(W\) 上为正，令 \(\rho_i=b_i/b\) 于 \(W\)。于是 \(\sum_i\rho_i=1\) 于 \(W\)。

	在 \(W\) 上取 \(\omega_i=\rho_i\omega\)，再在 \(M\setminus W\) 补零。因为 \(K\subset W\)，每个补零位置附近 \(\omega\) 本来就为零，故 \(\omega_i\) 光滑；其支集位于 \(K\cap\operatorname{supp}b_i\)，是相应图域内的紧集。对各 \(\omega_i\) 使用刚证的象限公式，按图的定向符号搬回，并把有限个等式相加。由 \(\sum_i\omega_i=\omega\)，或等价地由 \(\sum_i\mathrm{d}\rho_i=0\) 于 \(W\)，得到所需整体等式。
\end{proof}

如果一个区域由有限个这样的有向块拼成，并且相邻块沿完整面光滑贴合，就对每块应用公式再相加。公共面的外向横截方向对两块相反，而环境定向相同，所以两次诱导面定向相反，内部面项抵消。这里先按面在各块中出现的次数计入，再作抵消；不能一开始就把同一个几何面只算一次。

\ProseParagraphBreak

对按 \((x^1,\ldots,x^n)\) 定向的矩形 \(R=\prod_{i=1}^n[a_i,b_i]\)，记 \(R_{\widehat i}\) 为删去第 \(i\) 个坐标的矩形，取其剩余坐标递增次序的定向；令 \(\delta_i^+,\delta_i^-\) 分别插入 \(x^i=b_i,a_i\)。上面的向量次序给出：上侧面的外向向量为 \(+\partial_{x^i}\)，符号为 \((-1)^{i-1}\)；下侧面的外向向量为 \(-\partial_{x^i}\)，符号为 \((-1)^i\)。于是
\begin{equation}\label{g1:eq:rectangle-boundary-signs}
	\int_{\partial R}\omega
	=\sum_{i=1}^n(-1)^{i-1}\left(
		\int_{R_{\widehat i}}(\delta_i^+)^*\omega
		-\int_{R_{\widehat i}}(\delta_i^-)^*\omega
	\right).
\end{equation}
若把 \(\omega\) 写成前面的 \(\sum_i(-1)^{i-1}a_i\,\mathrm{d}x^{\widehat i}\)，两组排列符号抵消，留下的正是每个系数“上端值减下端值”。这也从逐坐标的积分基本定理直接验证了矩形公式。

\begin{lemma}[有序单形的边界定向]\label{g1:lem:simplex-boundary-orientation}
	设 \(n\ge1\)。给非退化有序仿射单形 \([v_0,\ldots,v_n]\) 以有序边向量 \((v_1-v_0,\ldots,v_n-v_0)\) 所确定的定向。删去第 \(i\) 个顶点后，面按剩余顶点的原次序定向，则该面的诱导边界定向等于此面定向的 \((-1)^i\) 倍。特别地，对标准单形
	\[
		\Delta^n=\left\{t\in\mathbb R^n:t^j\ge0,\ \sum_{j=1}^nt^j\le1\right\}
		=[0,e_1,\ldots,e_n],
	\]
	令 \(\delta_i:\Delta^{n-1}\to\Delta^n\) 为按该顶点次序参数化第 \(i\) 面的仿射映射。对光滑至边界的 \(n-1\) 次形式 \(\alpha\)，有
	\begin{equation}\label{g1:eq:simplex-stokes}
		\int_{\Delta^n}\mathrm{d}\alpha
		=\sum_{i=0}^n(-1)^i\int_{\Delta^{n-1}}\delta_i^*\alpha.
	\end{equation}
\end{lemma}
\begin{proof}
	先对标准单形计算。对 \(i\ge1\)，第 \(i\) 面为 \(t^i=0\)，删去 \(e_i\) 后的有序边向量为
	\[
		(e_1,\ldots,\widehat e_i,\ldots,e_n).
	\]
	取外向横截向量 \(-e_i\)，则
	\[
		\det(-e_i,e_1,\ldots,\widehat e_i,\ldots,e_n)=(-1)^i.
	\]
	因此该面的边界定向正好是有序面定向的 \((-1)^i\) 倍。

	第零面为 \(\sum_jt^j=1\)，其有序顶点为 \(e_1,\ldots,e_n\)，故有序切向基为
	\[
		(e_2-e_1,\ldots,e_n-e_1).
	\]
	这里可取外向横截向量 \(e_1\)，因为沿它移动会增大 \(\sum_jt^j\)。把第一列加到后面各列，可得
	\[
		\det(e_1,e_2-e_1,\ldots,e_n-e_1)=1.
	\]
	所以第零面的符号为 \(+1=(-1)^0\)。这一步使用横截向量而非必须垂直于面的单位法向，正体现了边界定向的选择无关性。

	一般有序仿射单形由映射 \(t\mapsto v_0+\sum_jt^j(v_j-v_0)\) 从标准单形得到；按题定的定向，它保向，并把各有序面与外向半侧对应，故保持上述符号。标准单形在每一点可用独立的活动面定义函数作为部分坐标，因而是带角流形；对它应用\cref{g1:thm:stokes-corners}即得\eqref{g1:eq:simplex-stokes}。当 \(n=1\) 时，这给出 \([v_1]-[v_0]\)；零维单形则另约定边界为空。
\end{proof}

\subsection{Green、Gauss 公式作为特例}
\label{g1:subsec:1404:05}

先看平面。设 \(D\subset\mathbb R^2\) 为紧的有向光滑带边界区域，或可按上一小节分块处理的带角区域，取标准面积定向；\(P,Q\) 在其邻域光滑。令
\[
	\alpha=P\,\mathrm{d}x+Q\,\mathrm{d}y,\qquad
	\mathrm{d}\alpha=(\partial_xQ-\partial_yP)\,\mathrm{d}x\wedge\mathrm{d}y.
\]
Stokes 公式立即给出
\begin{equation}\label{g1:eq:green-from-stokes}
	\int_{\partial D}P\,\mathrm{d}x+Q\,\mathrm{d}y
	=\int_D(\partial_xQ-\partial_yP)\,\mathrm{d}x\,\mathrm{d}y.
\end{equation}
现在解释边界方向。若 \(\nu\) 为外向单位法向，\(J\) 为平面的正向四分之一转动，则正向单位切向为 \(T=J\nu\)，因为 \((\nu,J\nu)\) 是正向基。行进方向的左侧为 \(JT=-\nu\)，恰好指向区域内部。因此正向边界使区域在左，外圈逆时针，孔洞的边界顺时针，如\cref{g1:fig:boundary-outward-first}。这与第九章的约定一致；\cref{g1:lem:green-formula}还单独证明了该章使用的分片 \(C^2\) 边界、\(C^1\) 一形式版本，并用截角极限处理内角在 \((0,2\pi)\) 内的角点。

\input{Figures/Chapter14-BoundaryOrientation.tex}

再看三维的散度公式。设 \(\Omega\subset\mathbb R^3\) 为紧的光滑或带角区域，取标准体积定向，\(V=P\partial_x+Q\partial_y+R\partial_z\) 在其邻域光滑。将向量场插入体积形式，得到二形式
\[
	\beta=\iota_V(\mathrm{d}x\wedge\mathrm{d}y\wedge\mathrm{d}z)
	=P\,\mathrm{d}y\wedge\mathrm{d}z
	-Q\,\mathrm{d}x\wedge\mathrm{d}z
	+R\,\mathrm{d}x\wedge\mathrm{d}y,
\]
以及
\[
	\mathrm{d}\beta=(\partial_xP+\partial_yQ+\partial_zR)\,
	\mathrm{d}x\wedge\mathrm{d}y\wedge\mathrm{d}z
	=(\operatorname{div}V)\,\mathrm{d}x\wedge\mathrm{d}y\wedge\mathrm{d}z.
\]
在每张光滑边界面上取外向单位法向 \(\nu\)。由外法向优先定向，其面积形式为 \(\mathrm{d}A=(\iota_\nu\mathrm{d}V)|_{\mathrm{T}\partial\Omega}\)，这里 \(\mathrm{d}V=\mathrm{d}x\wedge\mathrm{d}y\wedge\mathrm{d}z\)。把 \(V\) 分解为切向部分与 \(\langle V,\nu\rangle\nu\)，切向部分与两个面切向量线性相关，插入三次交替形式后为零，所以
\[
	\beta|_{\mathrm{T}\partial\Omega}=\langle V,\nu\rangle\,\mathrm{d}A.
\]
由\cref{g1:prop:metric-volume-integration}，这就是通常的通量积分。应用 Stokes 公式，得到
\begin{equation}\label{g1:eq:gauss-divergence-stokes}
	\int_\Omega\operatorname{div}V\,\mathrm{d}V
	=\int_{\partial\Omega}\langle V,\nu\rangle\,\mathrm{d}A.
\end{equation}
带角时逐面使用各自的外向单位法向，角棱上不必定义唯一法向，因为它们在面积积分中测度为零。

\ProseParagraphBreak

最后看通常称为曲面旋度公式的三维 Stokes 公式。设 \(S\subset\mathbb R^3\) 为紧的定向嵌入曲面，可有光滑边界或上述意义的角点。令 \(N\) 为与曲面定向相容的单位法向，\(V=(P,Q,R)\) 在 \(S\) 的邻域光滑。取一形式
\[
	\alpha=P\,\mathrm{d}x+Q\,\mathrm{d}y+R\,\mathrm{d}z.
\]
直接展开外微分可得
\[
	\begin{aligned}
		\mathrm{d}\alpha
		&=(\partial_xQ-\partial_yP)\,\mathrm{d}x\wedge\mathrm{d}y
		+(\partial_xR-\partial_zP)\,\mathrm{d}x\wedge\mathrm{d}z
		+(\partial_yR-\partial_zQ)\,\mathrm{d}y\wedge\mathrm{d}z\\
		&=\iota_{\operatorname{curl}V}(\mathrm{d}x\wedge\mathrm{d}y\wedge\mathrm{d}z),
	\end{aligned}
\]
其中
\[
	\operatorname{curl}V
	=(\partial_yR-\partial_zQ,\ \partial_zP-\partial_xR,\ \partial_xQ-\partial_yP).
\]
记 \(i:S\hookrightarrow\mathbb R^3\)。拉回与外微分交换，且同样的法向分解给出
\[
	\mathrm{d}(i^*\alpha)=i^*(\mathrm{d}\alpha)
	=\langle\operatorname{curl}V,N\rangle\,\mathrm{d}A.
\]
在曲面边界上，令 \(\nu\in \mathrm{T}S\) 为指向曲面外侧、垂直于边界的单位向量；它是曲面内的外向余法向，与空间法向 \(N\) 不同。正向边界切向 \(T\) 满足 \(\nu\times T=N\)，等价地 \(T=N\times\nu\)。一形式沿边界的值为 \(\alpha(T)=\langle V,T\rangle\)，故在 \(S\) 上使用 Stokes 公式得到
\begin{equation}\label{g1:eq:curl-surface-stokes}
	\int_S\langle\operatorname{curl}V,N\rangle\,\mathrm{d}A
	=\int_{\partial S}\langle V,T\rangle\,\mathrm{d}s.
\end{equation}
这三个经典公式使用了不同次数的形式，却都来自同一条\eqref{g1:eq:stokes}。向量、法向和面积元把欧氏内积带入了表达式，而抽象的 Stokes 定理本身只需要光滑结构、定向和支集条件。

% 纲要条目排为编号小节标题，并进入目录与书签；说明段仍保留于纲要注释中。

% !TeX root = ../../Book1.tex
% 纲要标题：### 14.5 积分、次数与周期
% 纲要位置：几何学计划纲要.md，第 678 行。
% 正文状态：初稿；尚未正式审核。

% 纲要内容开始
% 
% * 对光滑奇异单形积分
% * 边界、外微分与链映射
% * 次数积分公式的前瞻陈述
% 
%   对同维闭连通定向光滑流形之间的光滑映射 $f:M\to N$，对任意光滑顶次形式 $\eta$ 准确陈述 $\int_M f^*\eta=\deg(f)\int_N\eta$。本节只说明公式的几何含义，不把次数的整数性和一般定义作为已经证明的结果；正则值、同调与积分三种定义的完整比较在 II.16.1 完成。此前 Stokes 与积分构造不依赖本公式。
% 
% * 由周期提出 de Rham 比较问题
% 
% ---
% 
% 纲要内容结束

\section{积分、次数与周期}
\label{g1:sec:1405}

Stokes 定理可以用一句相容关系概括：先对形式求外微分再积分，等于先取积分区域的边界再积分。要把这句话带到拓扑中，需要一种能相加、取边界并允许退化的“积分区域”。光滑奇异单形承担这个作用。它不要求自己的像嵌入流形，因而能记录重复覆盖和相反方向，而不只记录一个点集。

\ProseParagraphBreak

本节的流形仍有限维、Hausdorff、第二可数；光滑形式不要求紧支，因为单形本身紧。我们只建立所需的链、上链与积分映射，并证明其相容性。把这些构造与一般奇异上同调比较，还需要第二卷的代数拓扑工具。次数积分公式依纲要先准确陈述，其一般定义与证明不参与本章此前的论证。

\subsection{对光滑奇异单形积分}
\label{g1:subsec:1405:01}

对 \(k\geq1\)，取标准单形
\[
	\Delta^k=\left\{(t^1,\ldots,t^k)\in\mathbb R^k:t^i\geq0,\ \sum_i t^i\leq1\right\},
\]
其有序顶点为 \(v_0=0,v_1=e_1,\ldots,v_k=e_k\)，定向为 \(\mathrm{d}t^1\wedge\cdots\wedge\mathrm{d}t^k\)。\(\Delta^0\) 是一个正向点。单形有边和角，但\cref{g1:thm:stokes-corners}已经涵盖这种积分区域。

\begin{definition}
	\label{g1:def:smooth-singular-simplex}
	流形 \(M\) 中的一个\term[guanghuaqiyidanxing]{光滑奇异单形}[smooth singular simplex]是映射 \(\sigma:\Delta^k\to M\)，其局部坐标分量在单形的每一点附近可光滑延拓到 \(\mathbb R^k\) 的开集。若 \(M\) 带边界，延拓只要求坐标分量光滑，不要求延拓后的所有点仍落在目标半空间。对 \(\omega\in\Omega^k(M)\)，定义
	\[
		\int_\sigma\omega=\int_{\Delta^k}\sigma^*\omega.
	\]
	当 \(k=0\) 时，右端为函数在 \(\sigma(\Delta^0)\) 的值。
\end{definition}

局部延拓使边和顶点处的导数有意义，其值由单形内部的导数连续延伸而唯一确定。拉回因此不依赖延拓选择。\(\sigma^*\omega\) 是紧单形上的光滑顶次形式，可以用上一节的带角积分；有限坐标覆盖保证其系数有界，所以积分有限。

\ProseParagraphBreak

“奇异”允许 \(\sigma\) 不单射，也允许微分降秩。如果 \(\operatorname{rank}\mathrm{d}\sigma<k\) 处处成立，则任意 \(k\) 个像向量都线性相关，交替性使 \(\sigma^*\omega=0\)，所以积分为零。另一方面，相同的像点集可以有不同积分：圆周绕行两次与绕行一次的参数化具有相同像，但前者对角形式的积分是后者的两倍。

\begin{example}[同一条圆周的不同参数化]
	\label{g1:exm:circle-simplex-integral}
	在 \(\mathbb R^2\setminus\{0\}\) 上取角形式
	\[
		\alpha=\frac{-y\,\mathrm{d}x+x\,\mathrm{d}y}{x^2+y^2},
	\]
	并对整数 \(m\) 取光滑一单形 \(\sigma_m(t)=(\cos 2\pi mt,\sin 2\pi mt)\)，\(0\leq t\leq1\)。直接拉回得 \(\sigma_m^*\alpha=2\pi m\,\mathrm{d}t\)，故 \(\int_{\sigma_m}\alpha=2\pi m\)。负整数表示反向绕行，零表示常值单形。这个计算只用参数化和一元积分，还没有调用一般映射次数。
\end{example}

\subsection{边界、外微分与链映射}
\label{g1:subsec:1405:02}

令 \(\delta_i:\Delta^{k-1}\to\Delta^k\) 为仿射面映射，其有序顶点依次送到 \(v_0,\ldots,\widehat v_i,\ldots,v_k\)。其中帽号表示省略。\cref{g1:lem:simplex-boundary-orientation}已证明，第 \(i\) 个面的边界定向相对于这一顶点顺序有系数 \((-1)^i\)。因此自然定义
\[
	\partial\sigma=\sum_{i=0}^k(-1)^i\sigma\circ\delta_i.
\]
右端是映射的形式和，而不是若干子集的并。对一单形，它是终点减起点；对二单形，它是 \([v_1,v_2]-[v_0,v_2]+[v_0,v_1]\)，也就是三角形的有向边界。

\begin{definition}
	\label{g1:def:smooth-singular-chains}
	记 \(C_k^\infty(M;\mathbb R)\) 为以全部光滑奇异 \(k\) 单形为基的自由实向量空间，其元素是有限形式和 \(c=\sum_j a_j\sigma_j\)，称为光滑奇异 \(k\) 链。将 \(\partial\) 线性延拓到链，并令 \(\partial:C_0^\infty\to0\)。令
	\[
		\int_c\omega=\sum_j a_j\int_{\sigma_j}\omega.
	\]
	实系数允许线性组合，但每条链只能有有限个非零系数。
\end{definition}

\begin{lemma}[边界的抵消]
	\label{g1:lem:smooth-chain-boundary-square}
	上述边界算子满足 \(\partial^2=0\)。
\end{lemma}
\begin{proof}
	对单个 \(k\) 单形，当 \(k\geq2\) 时，每个省略两个原顶点 \(v_i,v_j\)、\(i<j\) 的面出现两次。先省略 \(v_i\) 后，原来的 \(v_j\) 在剩余序列中编号为 \(j-1\)，所以这次的系数为 \((-1)^i(-1)^{j-1}\)。先省略 \(v_j\) 再省略 \(v_i\)，系数为 \((-1)^j(-1)^i\)。两个复合仿射映射在全部顶点相同，因而是同一映射；两个系数相反，恰好抵消。\(k=1\) 时边界已是零链，而零链的边界按定义为零。线性性给出一般链的结论。
\end{proof}

由此可以定义圈与边界：\(Z_k^\infty=\ker\partial\) 中的链称为圈，\(B_k^\infty=\operatorname{im}\partial\) 中的链称为边界，且 \(B_k^\infty\subset Z_k^\infty\)。商空间 \(H_k^\infty(M;\mathbb R)=Z_k^\infty/B_k^\infty\) 记录圈在允许加上边界后的等价类。这里给出的只是光滑奇异链模型；与所有连续奇异单形构成的模型是否相同，要由光滑逼近另行证明。

\begin{theorem}[积分与边界相容]
	\label{g1:thm:integration-chain-map}
	设 \(M\) 为光滑流形，\(k\geq1\)，\(c\in C_k^\infty(M;\mathbb R)\)，\(\omega\in\Omega^{k-1}(M)\)。则
	\begin{equation}\label{g1:eq:chain-stokes}
		\int_c\mathrm{d}\omega=\int_{\partial c}\omega.
	\end{equation}
	令 \(C^k_\infty(M;\mathbb R)=\operatorname{Hom}_{\mathbb R}(C_k^\infty(M;\mathbb R),\mathbb R)\)，上链微分定义为 \((\delta u)(c)=u(\partial c)\)。积分映射
	\[
		\mathcal I:\Omega^k(M)\longrightarrow C^k_\infty(M;\mathbb R),\qquad
		(\mathcal I\eta)(c)=\int_c\eta
	\]
	满足 \(\delta\mathcal I=\mathcal I\mathrm{d}\)。
\end{theorem}
\begin{proof}
	由线性性，只需对一个单形 \(\sigma\) 证明。外微分与拉回相容，故
	\[
		\int_\sigma\mathrm{d}\omega
		=\int_{\Delta^k}\mathrm{d}(\sigma^*\omega)
		=\int_{\partial\Delta^k}\sigma^*\omega
		=\sum_{i=0}^k(-1)^i\int_{\Delta^{k-1}}(\sigma\circ\delta_i)^*\omega.
	\]
	中间等式使用\cref{g1:thm:stokes-corners}，末式使用各面的诱导定向。\(\sigma\) 的像无需嵌入：Stokes 是在参数单形上应用的。\(k=1\) 时，边界零维积分给出终点值减起点值，与同一公式一致。对链求和即得\eqref{g1:eq:chain-stokes}。

	再对每个 \((k+1)\) 链 \(c\) 计算：
	\[
		(\delta\mathcal I\eta)(c)
		=(\mathcal I\eta)(\partial c)
		=\int_{\partial c}\eta
		=\int_c\mathrm{d}\eta
		=(\mathcal I\mathrm{d}\eta)(c).
	\]
	这就证明了上链映射的恒等式。\(\delta^2=0\) 则来自 \(\partial^2=0\)，故两边确实都是复形。
\end{proof}

这个关系可以画成交换图：
\[
	\begin{tikzcd}[column sep=large,row sep=large]
		\Omega^k(M) \arrow[r,"\mathrm{d}"] \arrow[d,"\mathcal I"'] & \Omega^{k+1}(M) \arrow[d,"\mathcal I"] \\
		C^k_\infty(M;\mathbb R) \arrow[r,"\delta"'] & C^{k+1}_\infty(M;\mathbb R).
	\end{tikzcd}
\]
上边在形式上求导，下边在输入链上取边界；竖边通过积分将它们联系起来。通常把这种次数保持且与微分交换的映射称为链映射，在此处更确切地说是上链复形之间的映射。

\subsection{次数积分公式的前瞻陈述}
\label{g1:subsec:1405:03}

圆周绕行的例子提示：若一个映射反复覆盖目标，积分可能把这些覆盖按方向加起来。一般光滑映射的纤维却可能很复杂，甚至含无限多个点，不能对任意目标点直接“数原像”来定义次数。下面先记录后续理论的准确结论。

\begin{theorem}[次数积分公式]
	\label{g1:thm:degree-integral-preview}
	设 \(M,N\) 为同维、闭、连通、定向光滑流形，\(f:M\to N\) 光滑。存在一个由 \(f\) 及两流形所选定向决定的整数 \(\deg(f)\)，使对每个 \(\eta\in\Omega^{\dim N}(N)\)，都有
	\begin{equation}\label{g1:eq:degree-integral-preview}
		\int_M f^*\eta=\deg(f)\int_N\eta.
	\end{equation}
\end{theorem}

此处“闭”指紧且无边界，所以两边积分自动有限。这个定理在本节是前瞻陈述：次数的一般定义、整数性及公式的证明安排在第二卷第16.1节“度的不同定义”，届时将比较正则值处的局部符号、同调定义与积分定义。不能把一个顶次形式的积分比直接宣布为已经定义好的整数，也不能在目标积分为零时用比值取代\eqref{g1:eq:degree-integral-preview}。此前积分构造、Stokes 定理与\cref{g1:thm:integration-chain-map}都未使用这个结论。

\ProseParagraphBreak

几何上，保定向的覆盖贡献正号，反向的覆盖贡献负号，次数记录其代数总数。我们可以在明确模型中验证积分倍数，而不借此证明一般定理。把圆周写为 \(\mathbb R/2\pi\mathbb Z\)，取 \(f_m([\theta])=[m\theta]\)，其中 \(m\in\mathbb Z\)。任何一形式都可写成 \(a(\theta)\,\mathrm{d}\theta\)，系数 \(a\) 为 \(2\pi\) 周期光滑函数。若 \(m>0\)，则
\[
	\int_{S^1}f_m^*(a\,\mathrm{d}\theta)
	=\int_0^{2\pi}m a(m\theta)\,\mathrm{d}\theta
	=\int_0^{2\pi m}a(u)\,\mathrm{d}u
	=m\int_0^{2\pi}a(u)\,\mathrm{d}u.
\]
负整数时，同样的有向一元换元给出负号及 \(|m|\) 次覆盖；\(m=0\) 时映射常值，拉回一形式为零。这个完整的直接计算说明，在该模型中积分倍数确实为 \(m\)。一般次数理论将把它识别为 \(\deg(f_m)\)。

\ProseParagraphBreak

对于保定向或反定向微分同胚，\cref{g1:thm:oriented-change-variables}也已经分别得到倍数 \(1\) 或 \(-1\)。这些模型解释公式，却没有处理任意映射的临界点或证明某个整数对所有顶次形式都适用。进一步阅读可定位到\cite[Chapter 17, Degree Theory, pp. 457--459]{Lee2012Smooth}。

\subsection{由周期提出 de Rham 比较问题}
\label{g1:subsec:1405:04}

对闭 \(k\) 形式 \(\omega\) 和光滑奇异 \(k\) 圈 \(z\)，数 \(\int_z\omega\) 称为 \(\omega\) 在该圈上的\term[zhouqi]{周期}[period]。它测量的是沿一个闭合的有向链累积的值，不必来自某个单一嵌入子流形。角形式的 \(2\pi\) 就是一个周期；它阻止该形式具有全局单值原函数。

\begin{proposition}[周期的两种不变性]
	\label{g1:prop:period-pairing}
	若 \(\mathrm{d}\omega=0\)、\(\partial z=0\)，则 \(\int_z\omega\) 在把 \(\omega\) 改为 \(\omega+\mathrm{d}\beta\)，或把 \(z\) 改为 \(z+\partial b\) 后均不变。特别地，恰当形式在所有圈上的周期为零；闭形式在所有边界上的积分为零。
\end{proposition}
\begin{proof}
	由\eqref{g1:eq:chain-stokes}，
	\[
		\int_z\mathrm{d}\beta=\int_{\partial z}\beta=0,
		\qquad
		\int_{\partial b}\omega=\int_b\mathrm{d}\omega=0.
	\]
	对第一条式子，\(k=0\) 时约定负次形式为零，不存在非零恰当零形式；第二条仍适用。这就覆盖所有次数。
\end{proof}

高维的穿孔空间也有相同障碍。在 \(\mathbb R^n\setminus\{0\}\)、\(n\geq2\) 上，记 \(r^2=\sum_i(x^i)^2\)，定义
\begin{equation}\label{g1:eq:radial-period-form}
	\eta=\frac1{r^n}\sum_{i=1}^n(-1)^{i-1}x^i\,
	\mathrm{d}x^1\wedge\cdots\wedge\widehat{\mathrm{d}x^i}\wedge\cdots\wedge\mathrm{d}x^n.
\end{equation}
分母只在被删除的原点为零，所以它在所述流形上光滑。外微分的坐标公式给出
\[
	\mathrm{d}\eta
	=\left(\sum_i\partial_i(x^i r^{-n})\right)\mathrm{d}x^1\wedge\cdots\wedge\mathrm{d}x^n
	=\left(nr^{-n}-nr^{-n-2}\sum_i(x^i)^2\right)\mathrm{d}x^1\wedge\cdots\wedge\mathrm{d}x^n=0.
\]
在单位球面 \(S^{n-1}\) 上，径向向量就是外向单位法向。按外法向优先的定向，\(\eta\) 的限制正是球面的正体积形式，所以其积分严格为正。若 \(\eta=\mathrm{d}\beta\)，拉回到球面再用 Stokes 定理，由 \(\partial S^{n-1}=\varnothing\) 得到该积分为零，矛盾。因此 \(\eta\) 闭而不恰当。对 \(n=3\)，这给出
\[
	\eta=\frac{x\,\mathrm{d}y\wedge\mathrm{d}z-y\,\mathrm{d}x\wedge\mathrm{d}z+z\,\mathrm{d}x\wedge\mathrm{d}y}{(x^2+y^2+z^2)^{3/2}},
	\qquad \int_{S^2}\eta=4\pi.
\]
其中球面积可由球坐标直接算出：\(\int_0^{2\pi}\int_0^\pi\sin\theta\,\mathrm{d}\theta\,\mathrm{d}\varphi=4\pi\)，坐标缺失的两极和切开的经线不影响面积积分。这里不能用穿孔空间中的 Stokes 定理将球面当成单位球的边界，因为原点不在该空间中。原点虽然只有一个，却使那个填充球不再是允许的积分区域。

\ProseParagraphBreak

为组织这些障碍，定义
\begin{equation}\label{g1:eq:de-rham-cohomology-definition}
	H_{\mathrm{dR}}^k(M)=
	\frac{\ker(\mathrm{d}:\Omega^k(M)\to\Omega^{k+1}(M))}
	{\operatorname{im}(\mathrm{d}:\Omega^{k-1}(M)\to\Omega^k(M))},
\end{equation}
称为 \(M\) 的\term[derhamshangtongdiao]{de Rham 上同调}[de Rham cohomology]。\footnote{de Rham（德拉姆），Georges de Rham（1903--1990），瑞士数学家，研究微分拓扑与代数拓扑，建立了微分形式与拓扑上同调之间的基本比较。}其中 \(\Omega^{-1}(M)=0\)；分母确实包含于分子，因为 \(\mathrm{d}^2=0\)。恰当差不会改变周期，所以可以在这个商空间上继续讨论周期。\cref{g1:prop:period-pairing}给出双线性配对
\[
	H_{\mathrm{dR}}^k(M)\times H_k^\infty(M;\mathbb R)\longrightarrow\mathbb R,
	\qquad ([\omega],[z])\longmapsto\int_z\omega.
\]
类似地，令 \(H^k_\infty(M;\mathbb R)=\ker(\delta:C^k_\infty\to C^{k+1}_\infty)/\operatorname{im}(\delta:C^{k-1}_\infty\to C^k_\infty)\)，称为光滑奇异上同调；负次上链取零。\cref{g1:thm:integration-chain-map}使 \(\mathcal I\) 下降为从 \(H_{\mathrm{dR}}^k(M)\) 到这个商空间的线性映射：闭形式给闭上链，恰当形式给上链边界，因此商上的映射良定义。

\ProseParagraphBreak

现在才出现比较问题：若一个闭形式在所有圈上的周期都为零，它是否恰当？拓扑上允许的每个实上同调类，是否都来自某个闭微分形式？本章已经证明构造这条比较映射所需的 Stokes 相容性，但其单射性、满射性，以及光滑奇异链与连续奇异链的比较，尚需另外证明。第二卷第7.3节“积分与 de Rham 定理”将通过相对光滑逼近、局部 Poincaré 引理与整体比较完成这些步骤。进一步阅读的位置为\cite[Chapter 18, Smooth Singular Homology, pp. 473--479; The de Rham Theorem, pp. 480--486]{Lee2012Smooth}。

\ProseParagraphBreak

已有的同伦公式说明了另一项可立即得到的事实。若 \(F_0,F_1:M\to N\) 光滑同伦，对闭形式 \(\omega\) 有
\[
	F_1^*\omega-F_0^*\omega=\mathrm{d}(K_H\omega).
\]
所以两映射诱导同一个 de Rham 上同调拉回。这里使用的是\cref{g1:thm:homotopy-formula}的显式原函数，而不是尚未证明的拓扑比较定理。局部收缩能制造原函数，整体绕行可能留下周期，积分则将这两种现象放在同一个公式中。

% 纲要条目排为编号小节标题，并进入目录与书签；说明段仍保留于纲要注释中。


\begin{chaptersummary}
	外微分由局部公式 \(\mathrm{d}(\sum_I a_I\,\mathrm{d}x^I)=\sum_I\mathrm{d}a_I\wedge\mathrm{d}x^I\) 构造。函数微分、分次乘积法则和平方为零唯一确定它，局部唯一性保证换坐标相容。外微分与任意光滑映射拉回交换；插入和李导数的自然性还须比较相关向量场。Cartan 公式 \(\mathcal L_X=\mathrm{d}\iota_X+\iota_X\mathrm{d}\) 将这三种运算联系起来。

	\ProseParagraphBreak

	在区间乘积上把 \(\mathrm{d}t\) 放在最前，积分其系数得到同伦算子，满足 \(\mathrm{d}K+K\mathrm{d}=i_1^*-i_0^*\)。径向收缩因此给出星形区域正次闭形式的显式原形式，也给出流形上的局部恰当性。零次闭形式只是局部常值，并非都恰当。角形式的非零周期说明，局部结论不能直接提升为整体原形式的存在性。

	\ProseParagraphBreak

	顶次形式按行列式换元，密度按绝对行列式换元。指定定向后可以对前者积分，不可定向流形仍可对后者积分。紧支性与局部有限的单位分解使积分构造只涉及有限个局部项；共同细化和欧氏换元证明其独立于所有选择。反转定向改变形式积分的符号，却不改变密度积分。

	\ProseParagraphBreak

	边界采用外法向优先的定向。半空间 \(x^n\geq0\) 的外向方向是 \(-\partial_n\)，相应的边界符号由这一约定实际算出。局部的一元积分、单位分解以及各面符号给出 \(\int_M\mathrm{d}\omega=\int_{\partial M}\omega\)，并延伸到带角区域。低维角层不另加顶维边界项；单形第 \(i\) 个面的有序顶点方向与边界方向相差 \((-1)^i\)。

	\ProseParagraphBreak

	对光滑奇异单形拉回再积分，不要求单形的像嵌入。Stokes 定理给出 \(\delta\mathcal I=\mathcal I\mathrm{d}\)，所以积分诱导形式上同调与光滑奇异上同调之间的比较映射。闭形式的周期只依赖形式模恰当形式的类及圈模边界的类；非零周期构成整体恰当性的障碍。比较映射是否为同构，以及一般映射怎样获得整数次数，仍是后续要证明的问题。
\end{chaptersummary}

\BookExercises

\begin{exercise}\label{g1:ex:exterior-pullback-graph}
	令 \(F:\mathbb R^2\to\mathbb R^3\)、\(F(u,v)=(u,v,uv)\)，并在 \(\mathbb R^3\) 上取一形式
	\[
		\omega=z\,\mathrm{d}x+x\,\mathrm{d}y+y\,\mathrm{d}z.
	\]
	分别计算 \(\mathrm{d}\omega\)、\(F^*\omega\)、\(\mathrm{d}(F^*\omega)\) 和 \(F^*(\mathrm{d}\omega)\)，直接核对外微分与拉回相容。说明计算 \(F^*\omega\) 时为何不能只把 \((x,y,z)\) 换成 \((u,v,uv)\)，却保留 \(\mathrm{d}x,\mathrm{d}y,\mathrm{d}z\) 不变。
\end{exercise}
\begin{solution}
	先对目标上的形式求外微分，得到
	\[
		\mathrm{d}\omega=\mathrm{d}z\wedge\mathrm{d}x+\mathrm{d}x\wedge\mathrm{d}y+\mathrm{d}y\wedge\mathrm{d}z.
	\]
	在定义域上，\(F^*\mathrm{d}x=\mathrm{d}u\)、\(F^*\mathrm{d}y=\mathrm{d}v\)、\(F^*\mathrm{d}z=\mathrm{d}(uv)=v\,\mathrm{d}u+u\,\mathrm{d}v\)，所以
	\[
		F^*\omega=uv\,\mathrm{d}u+u\,\mathrm{d}v+v(v\,\mathrm{d}u+u\,\mathrm{d}v)
		=(uv+v^2)\,\mathrm{d}u+(u+uv)\,\mathrm{d}v.
	\]
	按平面一形式的公式，
	\[
		\mathrm{d}(F^*\omega)
		=\bigl(\partial_u(u+uv)-\partial_v(uv+v^2)\bigr)\,\mathrm{d}u\wedge\mathrm{d}v
		=(1-u-v)\,\mathrm{d}u\wedge\mathrm{d}v.
	\]
	另一方面，逐项拉回 \(\mathrm{d}\omega\) 得
	\begin{align*}
		F^*(\mathrm{d}\omega)
		={}&(v\,\mathrm{d}u+u\,\mathrm{d}v)\wedge\mathrm{d}u
		+\mathrm{d}u\wedge\mathrm{d}v
		+\mathrm{d}v\wedge(v\,\mathrm{d}u+u\,\mathrm{d}v)\\
		={}&(-u+1-v)\,\mathrm{d}u\wedge\mathrm{d}v,
	\end{align*}
	与前式相同。

	拉回后的形式应当接受 \(\mathbb R^2\) 上的切向量作为输入。\(\mathrm{d}x,\mathrm{d}y,\mathrm{d}z\) 原来属于目标空间的余切空间，必须通过 \(\mathrm{d}F\) 一并拉回，才能在定义域的切向量上求值。只替换系数会遗漏这些输入的变化；本题中尤其会漏掉 \(\mathrm{d}(uv)\) 同时包含 \(\mathrm{d}u\) 与 \(\mathrm{d}v\) 的事实。
\end{solution}

\begin{exercise}\label{g1:ex:exterior-polar-coframe}
	在 \(U=(0,\infty)\times(-\pi,\pi)\) 上取极坐标映射 \(F(r,\theta)=(r\cos\theta,r\sin\theta)\)。令目标上的 \(\alpha=-y\,\mathrm{d}x+x\,\mathrm{d}y\)，并在 \(U\) 上令
	\[
		\eta^1=\mathrm{d}r,\qquad \eta^2=r\,\mathrm{d}\theta,\qquad
		E_1=\partial_r,\qquad E_2=r^{-1}\partial_\theta.
	\]
	计算 \(F^*\alpha\) 与 \(\mathrm{d}(F^*\alpha)\)，并用 \(F^*(\mathrm{d}\alpha)\) 核对。再计算 \(\mathrm{d}\eta^1\)、\(\mathrm{d}\eta^2\) 和 \([E_1,E_2]\)，用\eqref{g1:eq:exterior-one-form-evaluation}核对 \(\mathrm{d}\eta^2(E_1,E_2)\)。解释为何不能把 \(\eta^1,\eta^2\) 都当成坐标微分，直接令它们的外微分为零。
\end{exercise}
\begin{solution}
	普通链式法则给出
	\[
		F^*\mathrm{d}x=\cos\theta\,\mathrm{d}r-r\sin\theta\,\mathrm{d}\theta,\qquad
		F^*\mathrm{d}y=\sin\theta\,\mathrm{d}r+r\cos\theta\,\mathrm{d}\theta.
	\]
	代入可得 \(F^*\alpha=r^2\,\mathrm{d}\theta\)，故
	\[
		\mathrm{d}(F^*\alpha)=2r\,\mathrm{d}r\wedge\mathrm{d}\theta.
	\]
	在目标上，\(\mathrm{d}\alpha=-\mathrm{d}y\wedge\mathrm{d}x+\mathrm{d}x\wedge\mathrm{d}y=2\,\mathrm{d}x\wedge\mathrm{d}y\)。将前面的两个拉回一形式作楔积，得到 \(F^*(\mathrm{d}x\wedge\mathrm{d}y)=r\,\mathrm{d}r\wedge\mathrm{d}\theta\)，因此结果一致。

	对所给余标架，
	\[
		\mathrm{d}\eta^1=0,\qquad
		\mathrm{d}\eta^2=\mathrm{d}r\wedge\mathrm{d}\theta
		=\frac1r\eta^1\wedge\eta^2.
	\]
	向量场的括号为
	\[
		[E_1,E_2]=[\partial_r,r^{-1}\partial_\theta]
		=-r^{-2}\partial_\theta=-\frac1r E_2.
	\]
	由于 \(\eta^2(E_1)=0\)、\(\eta^2(E_2)=1\)，这两个求值函数均为常数，故一形式的微分公式给出
	\[
		\mathrm{d}\eta^2(E_1,E_2)
		=E_1(1)-E_2(0)-\eta^2\left(-\frac1rE_2\right)=\frac1r,
	\]
	正好与上面的二形式求值相符。

	\(\eta^1,\eta^2\) 的确在每点构成余切空间的一组基，但余标架不一定由一组坐标函数的微分组成。这里若 \(\eta^2\) 在某开邻域等于 \(\mathrm{d}q\)，就应有 \(\mathrm{d}\eta^2=\mathrm{d}^2q=0\)，与计算矛盾。坐标公式中 \(\mathrm{d}(\mathrm{d}x^i)=0\) 来自这些一形式是函数微分，不能推广为任意局部标架都“无需微分”。
\end{solution}

\begin{exercise}\label{g1:ex:homogeneous-closed-form-primitive}
	在 \(\mathbb R^n\) 上令 \(E=\sum_{j=1}^n x^j\partial_{x^j}\)，并设 \(\omega\) 是次数 \(k\geq1\) 的微分形式，其坐标系数都是次数 \(m\geq0\) 的齐次多项式。证明
	\[
		\mathcal L_E\omega=(m+k)\omega.
	\]
	若再假设 \(\mathrm{d}\omega=0\)，证明 \((m+k)^{-1}\iota_E\omega\) 是一个原形式。用这个公式给出 \(\mathrm{d}x\wedge\mathrm{d}y\) 在 \(\mathbb R^2\) 上的一个原形式，并对 \(\beta=x\,\mathrm{d}y\) 检查闭性条件为何不可省略。
\end{exercise}
\begin{solution}
	\(E\) 的流为 \(\Phi_t(x)=e^t x\)，对一切实数 \(t\) 有定义。若 \(a\) 为次数 \(m\) 的齐次多项式，则 \(a(e^t x)=e^{mt}a(x)\)，而 \(\Phi_t^*\mathrm{d}x^j=e^t\mathrm{d}x^j\)。一个 \(k\) 次楔积因而再贡献 \(e^{kt}\)，所以逐项得到
	\[
		\Phi_t^*\omega=e^{(m+k)t}\omega.
	\]
	在 \(t=0\) 求导，即得 \(\mathcal L_E\omega=(m+k)\omega\)。这里 \(k\) 的贡献来自一形式基的拉回；仅对系数使用齐次关系会漏掉它。

	若 \(\mathrm{d}\omega=0\)，Cartan 公式给出
	\[
		(m+k)\omega=\mathrm{d}(\iota_E\omega)+\iota_E(\mathrm{d}\omega)
		=\mathrm{d}(\iota_E\omega).
	\]
	由于 \(m\geq0\)、\(k\geq1\)，分母 \(m+k\) 非零，故 \(\eta=(m+k)^{-1}\iota_E\omega\) 满足 \(\mathrm{d}\eta=\omega\)。

	对于 \(\omega=\mathrm{d}x\wedge\mathrm{d}y\)，有 \(m=0\)、\(k=2\)，且 \(\iota_E\omega=x\,\mathrm{d}y-y\,\mathrm{d}x\)。因此可取
	\[
		\eta=\frac12(x\,\mathrm{d}y-y\,\mathrm{d}x),\qquad
		\mathrm{d}\eta=\mathrm{d}x\wedge\mathrm{d}y.
	\]
	对 \(\beta=x\,\mathrm{d}y\)，虽然 \(m=k=1\)，却有 \(\mathrm{d}\beta=\mathrm{d}x\wedge\mathrm{d}y\ne0\)。若忽略闭性，得到的候选函数是 \(\tfrac12\iota_E\beta=xy/2\)，其微分为
	\[
		\mathrm{d}(xy/2)=\frac12(y\,\mathrm{d}x+x\,\mathrm{d}y)\ne\beta.
	\]
	差异正是 Cartan 公式中被错误删去的 \(\iota_E\mathrm{d}\beta\) 项。事实上非闭形式根本不可能恰当，因为 \(\mathrm{d}^2=0\)。
\end{solution}

\begin{exercise}\label{g1:ex:lie-function-multiple-area}
	设 \(M\) 为无边界光滑流形，\(X\) 为光滑向量场，\(f\in C^\infty(M)\)。证明对任意形式 \(\omega\)，
	\[
		\mathcal L_{fX}\omega=f\mathcal L_X\omega+\mathrm{d}f\wedge\iota_X\omega.
	\]
	再在 \(\mathbb R^2\) 上令 \(X=-y\partial_x+x\partial_y\)、\(\mu=\mathrm{d}x\wedge\mathrm{d}y\)，证明 \(fX\) 的局部流保持 \(\mu\) 当且仅当 \(Xf=0\)，并把这个条件解释为 \(f\) 在每个以原点为中心的圆周上为常数。分别判断 \(f=x\) 与 \(f=x^2+y^2\)。
\end{exercise}
\begin{solution}
	由于 \(\iota_{fX}\omega=f\iota_X\omega\)，Cartan 公式与乘积法则给出
	\begin{align*}
		\mathcal L_{fX}\omega
		&=\mathrm{d}(f\iota_X\omega)+f\iota_X\mathrm{d}\omega\\
		&=\mathrm{d}f\wedge\iota_X\omega+f\bigl(\mathrm{d}\iota_X\omega+\iota_X\mathrm{d}\omega\bigr)\\
		&=\mathrm{d}f\wedge\iota_X\omega+f\mathcal L_X\omega.
	\end{align*}
	零次形式的插入为零，这个公式也包含 \(\mathcal L_{fX}g=fXg\)。正次形式则一般多出第一项，所以李导数对向量场的函数倍乘通常不是线性的。

	对平面的转动场，
	\[
		\iota_X\mu=-x\,\mathrm{d}x-y\,\mathrm{d}y
		=-\mathrm{d}\bigl((x^2+y^2)/2\bigr).
	\]
	又因 \(\mathrm{d}\mu=0\)，Cartan 公式给出 \(\mathcal L_X\mu=0\)。于是
	\begin{align*}
		\mathcal L_{fX}\mu
		&=(\partial_xf\,\mathrm{d}x+\partial_yf\,\mathrm{d}y)\wedge(-x\,\mathrm{d}x-y\,\mathrm{d}y)\\
		&=(-y\partial_xf+x\partial_yf)\,\mathrm{d}x\wedge\mathrm{d}y
		=(Xf)\mu.
	\end{align*}
	\cref{g1:prop:lie-derivative-tensor-rules}已证明，一个张量被局部流保持当且仅当其李导数为零。因此所求条件恰为 \(Xf=0\)，不需要预先假设 \(fX\) 完备。

	固定 \(r>0\)，沿圆周 \(c(\theta)=(r\cos\theta,r\sin\theta)\) 有 \(c'(\theta)=X_{c(\theta)}\)，故
	\[
		\frac{\mathrm{d}}{\mathrm{d}\theta}f(c(\theta))=(Xf)(c(\theta)).
	\]
	所以 \(Xf=0\) 当且仅当 \(f\) 在每个这样的圆周上为常数；原点处 \(X=0\)，条件自动成立。若 \(f=x\)，则 \(Xf=-y\)，并非处处为零，局部流不保持 \(\mu\)。若 \(f=x^2+y^2\)，则 \(Xf=-2xy+2xy=0\)，局部流保持 \(\mu\)。后一个向量场的转动速度随半径变化，但同一圆周上速度的倍乘因子保持不变。
\end{solution}

\begin{exercise}\label{g1:ex:path-homotopy-integral}
	设 \(\alpha\) 是光滑流形 \(M\) 上的一形式，\(H:[0,1]^2\to M\) 光滑，记 \(\gamma_s(t)=H(s,t)\)。写
	\[
		H^*\alpha=A(s,t)\,\mathrm{d}s+B(s,t)\,\mathrm{d}t,\qquad
		I(s)=\int_{\gamma_s}\alpha=\int_0^1B(s,t)\,\mathrm{d}t.
	\]
	不用 Stokes 定理，证明
	\[
		I'(s)=A(s,1)-A(s,0)+\int_0^1(\mathrm{d}\alpha)_{H(s,t)}
		\bigl(\partial_sH(s,t),\partial_tH(s,t)\bigr)\,\mathrm{d}t.
	\]
	由此证明：若 \(\alpha\) 闭，且同伦固定路径的两个端点，或者每条 \(\gamma_s\) 都是闭路，则 \(I(s)\) 与 \(s\) 无关。最后取 \(\mathbb R^2\) 上的 \(\alpha=y\,\mathrm{d}x\) 和 \(H(s,t)=(t,st(1-t))\)，计算 \(I(s)\)，说明闭性条件不能去掉。
\end{exercise}
\begin{solution}
	参数正方形上的光滑性采用局部光滑延拓的约定，所以系数及其偏导在边界处也有确定意义。外微分的坐标公式给出
	\[
		\mathrm{d}(H^*\alpha)
		=\mathrm{d}A\wedge\mathrm{d}s+\mathrm{d}B\wedge\mathrm{d}t
		=(\partial_sB-\partial_tA)\,\mathrm{d}s\wedge\mathrm{d}t.
	\]
	这里的负号来自 \(\mathrm{d}t\wedge\mathrm{d}s=-\mathrm{d}s\wedge\mathrm{d}t\)。另一方面，外微分与拉回相容，将 \(\mathrm{d}(H^*\alpha)=H^*(\mathrm{d}\alpha)\) 在 \((\partial_s,\partial_t)\) 上求值，得到
	\[
		\partial_sB-\partial_tA
		=(\mathrm{d}\alpha)_{H(s,t)}
		\bigl(\partial_sH(s,t),\partial_tH(s,t)\bigr).
	\]
	光滑系数在紧参数区间上允许对积分求导，因此
	\[
		I'(s)=\int_0^1\partial_sB(s,t)\,\mathrm{d}t
		=\int_0^1\partial_tA(s,t)\,\mathrm{d}t
		+\int_0^1(\mathrm{d}\alpha)_{H(s,t)}
		\bigl(\partial_sH(s,t),\partial_tH(s,t)\bigr)\,\mathrm{d}t.
	\]
	对第一项使用一元微积分基本定理，便得到所求等式。它在 \(0<s<1\) 上成立，并由系数连续性延至端点的一侧导数。

	若 \(\mathrm{d}\alpha=0\)，只剩端点差。固定端点时，\(H(s,0)\) 与 \(H(s,1)\) 各自不随 \(s\) 变化，故 \(\partial_sH(s,0)=\partial_sH(s,1)=0\)，从而 \(A(s,0)=A(s,1)=0\)。若每条路径都是闭路，则 \(H(s,0)=H(s,1)\)；对 \(s\) 求导还得到 \(\partial_sH(s,0)=\partial_sH(s,1)\)。又因为
	\[
		A(s,t)=\alpha_{H(s,t)}\bigl(\partial_sH(s,t)\bigr),
	\]
	这时同样有 \(A(s,0)=A(s,1)\)，即使闭路的共同起终点随 \(s\) 移动也不影响结论。两种情形均给出 \(I'(s)=0\)，所以 \(I\) 在整个参数区间上为常数。

	对所给反例，所有路径都从 \((0,0)\) 到 \((1,0)\)，端点固定。但
	\[
		\mathrm{d}\alpha=\mathrm{d}y\wedge\mathrm{d}x\ne0,\qquad
		H^*\alpha=st(1-t)\,\mathrm{d}t,
	\]
	所以
	\[
		I(s)=s\int_0^1t(1-t)\,\mathrm{d}t=\frac{s}{6},
	\]
	确实随 \(s\) 改变。此时 \(A=0\)，而
	\[
		(\mathrm{d}\alpha)(\partial_sH,\partial_tH)
		=(\mathrm{d}y\wedge\mathrm{d}x)\bigl((0,t(1-t)),(1,s(1-2t))\bigr)=t(1-t).
	\]
	变化量正是公式中的外微分积分项；端点固定本身不足以保证积分不变。
\end{solution}

\begin{exercise}\label{g1:ex:compactly-supported-primitive-line}
	设 \(f\in C^\infty_c(\mathbb R)\)。证明一形式 \(f(x)\,\mathrm{d}x\) 总是恰当的，但它存在紧支撑光滑原函数的充要条件是
	\[
		\int_{-\infty}^{+\infty}f(x)\,\mathrm{d}x=0.
	\]
	解释这为什么不与 Poincaré 引理矛盾。
\end{exercise}
\begin{solution}
	取 \(A<B\) 使 \(\operatorname{supp}f\subset[A,B]\)，定义
	\[
		F(x)=\int_{-\infty}^x f(s)\,\mathrm{d}s.
	\]
	这里的积分实际上只涉及一个有界区间。微积分基本定理给出 \(F'=f\)，再逐次求导可知 \(F\) 光滑，所以 \(\mathrm{d}F=f(x)\,\mathrm{d}x\)。当 \(x<A\) 时 \(F(x)=0\)；当 \(x>B\) 时，\(F(x)\) 等于总积分。因此总积分为零时，\(\operatorname{supp}F\subset[A,B]\)，它就是一个紧支撑原函数。

	反过来，若 \(G\) 是紧支撑原函数，则 \((G-F)'=0\)，所以 \(G-F\) 在连通的实线上为常数。在足够靠左的区间上，\(G\) 与 \(F\) 都为零，故此常数为零。在足够靠右的区间上，\(G=0\)，于是 \(F\) 的常值，即 \(f\) 的总积分，也必须为零。

	例如，取一个非负且不恒为零的光滑紧支撑函数 \(f\)，其总积分为正，便得到没有紧支撑原函数的恰当一形式。Poincaré 引理保证某个光滑原形式存在，并没有保证可以同时控制它的支撑；本题增加的支撑条件确实带来了新的障碍。
\end{solution}

\begin{exercise}\label{g1:ex:angle-branches-gluing}
	在穿孔平面 \(P=\mathbb R^2\setminus\{0\}\) 上，令
	\[
		U_-=P\setminus\{(x,0):x<0\},\qquad
		U_+=P\setminus\{(x,0):x>0\}.
	\]
	在 \(U_-\) 上选取极角 \(\theta_-\in(-\pi,\pi)\)，在 \(U_+\) 上选取极角 \(\theta_+\in(0,2\pi)\)。证明二者均为\cref{g1:exm:angular-period}中角形式 \(\vartheta\) 的光滑原函数。求出 \(\theta_+-\theta_-\) 在交集的每个连通分支上的值，并证明不存在常数 \(c_-,c_+\) 使 \(\theta_-+c_-\) 与 \(\theta_++c_+\) 在整个交集上相等。
\end{exercise}
\begin{solution}
	映射 \((r,\theta)\mapsto(r\cos\theta,r\sin\theta)\) 在 \(r>0\) 时的导数行列式为 \(r\ne0\)。分别把 \(\theta\) 限制到题给的两个区间，映射各自双射到 \(U_-\)、\(U_+\)；逆函数定理说明其局部逆光滑，而局部逆由唯一性彼此相容。因此 \(r\) 与两个极角函数都是相应区域上的光滑函数。由
	\[
		\mathrm{d}x=\cos\theta\,\mathrm{d}r-r\sin\theta\,\mathrm{d}\theta,\qquad
		\mathrm{d}y=\sin\theta\,\mathrm{d}r+r\cos\theta\,\mathrm{d}\theta
	\]
	可得
	\[
		\frac{-y\,\mathrm{d}x+x\,\mathrm{d}y}{x^2+y^2}=\mathrm{d}\theta,
	\]
	所以两个函数确实都是局部原函数。

	交集 \(U_-\cap U_+=\{y\ne0\}\) 有上半平面与下半平面两个连通分支。在上半平面，两个角都位于 \((0,\pi)\)，因此 \(\theta_+=\theta_-\)；在下半平面，\(\theta_-\in(-\pi,0)\)、\(\theta_+\in(\pi,2\pi)\)，因此 \(\theta_+=\theta_-+2\pi\)。若添加常数后可以拼接，上半平面要求 \(c_-=c_+\)，下半平面却要求 \(c_-=c_++2\pi\)，二者矛盾。

	这里每个局部区域本身都连通，所以其中任意两个原函数只相差一个常数。因而失败之处已经覆盖了对这两个局部原函数的一切可能修正；问题来自交集的两个分支提出了互不相容的常数条件。
\end{solution}

\begin{exercise}\label{g1:ex:punctured-plane-closed-forms}
	设 \(\eta\) 是 \(P=\mathbb R^2\setminus\{0\}\) 上的光滑闭一形式，令 \(\gamma(t)=(\cos t,\sin t)\)、\(0\le t\le2\pi\)。证明存在光滑函数 \(f:P\to\mathbb R\) 和唯一实数 \(c\)，使
	\[
		\eta=c\,\vartheta+\mathrm{d}f,\qquad
		c=\frac1{2\pi}\int_\gamma\eta.
	\]
	其中 \(\vartheta=(-y\,\mathrm{d}x+x\,\mathrm{d}y)/(x^2+y^2)\)，而 \(f\) 在相差常数的意义下唯一。可以使用径向投影 \(\rho(x)=x/\lVert x\rVert\) 及同伦
	\[
		H(t,x)=\left(1-t+\frac{t}{\lVert x\rVert}\right)x.
	\]
\end{exercise}
\begin{solution}
	记 \(j:S^1\to P\) 为包含映射，并令 \(\sigma=j^*\vartheta\)。因为 \(x\ne0\) 且 \(0\le t\le1\)，同伦公式中的系数 \(1-t+t/\lVert x\rVert\) 始终严格为正。因此 \(H(t,x)\ne0\)，这确实是 \(P\) 内的光滑同伦，端点为 \(H_0=\operatorname{id}_P\)、\(H_1=j\circ\rho\)。

	由 \(\mathrm{d}\eta=0\) 和\cref{g1:thm:homotopy-formula}，
	\[
		\rho^*j^*\eta-\eta=\mathrm{d}(K_H\eta).
	\]
	另一方面，\cref{g1:prop:circle-exactness-period}给出
	\[
		j^*\eta=c\,\sigma+\mathrm{d}h,\qquad
		c=\frac1{2\pi}\int_\gamma\eta
	\]
	以及某个 \(h\in C^\infty(S^1)\)。为计算 \(\rho^*\sigma\)，写 \(r=\sqrt{x^2+y^2}\)、\(u=x/r\)、\(v=y/r\)。由于圆周上的 \(\sigma\) 为 \(u\,\mathrm{d}v-v\,\mathrm{d}u\)，有
	\[
		\rho^*\sigma
		=\frac{x}{r}\left(\frac{\mathrm{d}y}{r}-\frac{y\,\mathrm{d}r}{r^2}\right)
		-\frac{y}{r}\left(\frac{\mathrm{d}x}{r}-\frac{x\,\mathrm{d}r}{r^2}\right)
		=\vartheta.
	\]
	结合以上各式并使用外微分与拉回相容，得到
	\[
		\eta=c\,\vartheta+\mathrm{d}\bigl(h\circ\rho-K_H\eta\bigr).
	\]
	所以取 \(f=h\circ\rho-K_H\eta\) 即可，且它光滑。

	若另有 \(\eta=c'\vartheta+\mathrm{d}f'\)，沿 \(\gamma\) 作一元路径积分。\cref{g1:lem:exact-path-period}说明恰当项的积分为零，而 \(\int_\gamma\vartheta=2\pi\)，故 \(c'=c\)。随后 \(\mathrm{d}(f'-f)=0\)，所以 \(f'-f\) 局部常值。穿孔平面连通：任意两点可以先沿径向线段移动到单位圆，再用圆弧连接。因此该差在整个 \(P\) 上是一个常数。特别地，这也证明 \(\eta\) 恰当当且仅当它沿单位圆的周期为零。
\end{solution}

\begin{exercise}\label{g1:ex:density-critical-map}
	考虑 \(F:\mathbb R\to\mathbb R\)、\(F(x)=x^2\)，目标坐标记作 \(y\)。尝试按绝对 Jacobian 规则把目标密度 \(\mu=|\mathrm{d}y|\) 变成源上的密度，并与微分形式 \(F^*(\mathrm{d}y)\) 比较光滑性。再把 \(F\) 限制为 \([-1,1]\to[0,1]\)，两个区间均取坐标增大的定向，分别计算候选密度与拉回一形式在源上的积分，并与目标积分比较。说明这些现象与“任意光滑映射都可拉回密度线丛”是否矛盾，以及本节换元定理为何要求微分同胚。
\end{exercise}
\begin{solution}
	绝对 Jacobian 给出的候选表达为
	\[
		\nu=|F'(x)|\,|\mathrm{d}x|=2|x|\,|\mathrm{d}x|.
	\]
	它的系数连续，但在 \(0\) 处的左导数为 \(-2\)、右导数为 \(2\)，所以不是光滑密度。相反，普通微分形式的拉回为
	\[
		F^*(\mathrm{d}y)=\mathrm{d}(x^2)=2x\,\mathrm{d}x,
	\]
	其系数处处光滑。关键区别在于绝对值；形式保留行列式的符号。

	虽然 \(\nu\) 不属于本节定义的光滑密度空间，其连续系数仍可按通常的一元积分计算。于是
	\[
		\int_{-1}^1 2|x|\,\mathrm{d}x=2,
		\qquad \int_{[0,1]}|\mathrm{d}y|=1,
	\]
	而形式积分为
	\[
		\int_{[-1,1]}F^*(\mathrm{d}y)=\int_{-1}^1 2x\,\mathrm{d}x=0,
		\qquad \int_{[0,1]}\mathrm{d}y=1.
	\]
	每个 \(y\in(0,1)\) 有两个原像。密度计算把两支都作正贡献，故得到两倍；形式计算中正、负半轴分别贡献 \(1\) 与 \(-1\)，恰好抵消。这些值都是直接的一元积分结果，不需要覆盖次数的积分公式。

	沿任意 \(F\) 拉回密度线丛，得到的是一个位于源上的线丛 \(F^*\mathcal D(\mathbb R)\)，其在 \(x\) 处的纤维是目标点 \(x^2\) 处的密度线。目标密度确实给出该拉回丛的光滑截面；这一步不含绝对 Jacobian。若还要借助 \(\mathrm{d}F\) 把它变成源自身的密度丛 \(\mathcal D(\mathbb R)\) 中的截面，才会出现 \(|F'(x)|\)，并在本题中失去光滑性。因此两件事并不矛盾。

	微分同胚使 Jacobian 处处非零，其绝对值因而光滑，又避免同一目标区域被多次计入。一般光滑映射不能同时自动保证这两点。本题 \(F\) 既在 \(0\) 处降秩，又不是单射，所以\cref{g1:thm:oriented-change-variables}的条件不成立。局部微分同胚足以保证密度的局部拉回光滑，但若没有整体单射性，仍不能直接断言源、目标积分相等。
\end{solution}

\begin{exercise}\label{g1:ex:conformal-density-sphere-area}
	设 \(M\) 为 \(n\geq1\) 维光滑流形，\(g\) 为光滑正定度量，\(u\in C^\infty(M)\)。令 \(\widetilde g=e^{2u}g\)，证明其体积密度满足 \(\mu_{\widetilde g}=e^{nu}\mu_g\)；若另给定向，相应的正体积形式有怎样的变化？再为半径 \(R>0\) 的欧氏球面 \(S_R^2\subset\mathbb R^3\) 写出球坐标中的面积密度，具体积分求出总面积，并核对它与单位球面积之间的缩放关系。
\end{exercise}
\begin{solution}
	在任一坐标图中，\(\widetilde g\) 的矩阵为 \(e^{2u}(g_{ij})\)，所以
	\[
		\det(\widetilde g_{ij})=e^{2nu}\det(g_{ij}),\qquad
		\sqrt{\det(\widetilde g_{ij})}=e^{nu}\sqrt{\det(g_{ij})}.
	\]
	由\cref{g1:prop:metric-volume-integration}的坐标公式，\(\mu_{\widetilde g}=e^{nu}\mu_g\)。因为指定的定向未改变，正体积形式中的坐标定向符号也未改变，故同样有
	\[
		\operatorname{vol}_{\widetilde g}=e^{nu}\operatorname{vol}_g.
	\]
	长度乘以 \(e^u\)，而 \(n\) 维体积同时包含 \(n\) 个方向的缩放，因此指数为 \(nu\)，并非仅为 \(u\) 或 \(2u\)。

	对球面使用参数
	\[
		q(\theta,\varphi)
		=(R\sin\theta\cos\varphi,R\sin\theta\sin\varphi,R\cos\theta),
		\qquad 0<\theta<\pi,\quad 0<\varphi<2\pi.
	\]
	两个坐标切向量的内积为
	\[
		\langle q_\theta,q_\theta\rangle=R^2,\qquad
		\langle q_\theta,q_\varphi\rangle=0,\qquad
		\langle q_\varphi,q_\varphi\rangle=R^2\sin^2\theta.
	\]
	所以诱导度量及面积密度的拉回为
	\[
		q^*g_R=R^2\,\mathrm{d}\theta^2+R^2\sin^2\theta\,\mathrm{d}\varphi^2,
		\qquad q^*\mu_{g_R}=R^2\sin\theta\,|\mathrm{d}\theta\,\mathrm{d}\varphi|.
	\]
	这里 \(\sin\theta>0\)，取平方根时没有额外的符号。

	这张参数图遗漏两极和一条子午线。两极在局部图中是有限点；被遗漏的子午线包含于一个大圆，大圆在每点附近可用光滑坐标表示为一条坐标轴。由紧性选有限张这样的图，再用欧氏积分中直线的零测性，可知这些集合在二维坐标中均为零测集。用有限单位分解计算球面积分时，删去它们不改变任何一项，余下部分可以在上述参数图中换元并合并积分。于是
	\[
		\int_{S_R^2}\mu_{g_R}
		=\int_0^{2\pi}\int_0^\pi R^2\sin\theta\,\mathrm{d}\theta\,\mathrm{d}\varphi
		=R^2(2)(2\pi)=4\pi R^2.
	\]
	参数表达在两极的退化是球坐标失效，不能解释为球面度量在那里退化。

	最后，缩放映射 \(s_R:S_1^2\to S_R^2\)、\(s_R(p)=Rp\) 是微分同胚，其微分把每个切向量乘以 \(R\)，所以 \(s_R^*g_R=R^2g_1\)。在前面的共形公式中取 \(n=2\)、\(u=\log R\)，便有 \(s_R^*\mu_{g_R}=R^2\mu_{g_1}\)。密度换元因此给出总面积乘以 \(R^2\)，与直接积分一致。
\end{solution}

\begin{exercise}\label{g1:ex:unit-integral-supported-form}
	设 \(M\) 为有限维、Hausdorff、第二可数的定向光滑 \(n\) 维流形，允许有边界，\(U\subseteq M\) 为非空开集。证明存在 \(\omega\in\Omega_c^n(M)\)，使
	\[
		\operatorname{supp}\omega\subseteq U,\qquad \int_M\omega=1.
	\]
	要求同时处理 \(n=0\)。若不假设 \(M\) 可定向，证明仍存在支集包含于 \(U\)、积分等于 \(1\) 的非负光滑紧支密度。
\end{exercise}
\begin{solution}
	先设 \(n\geq1\)。取 \(p\in U\)，由\cref{g1:cor:smooth-cutoff-closed}对紧集 \(\{p\}\subset U\) 的结论，存在 \(\chi\in C^\infty(M)\)，满足
	\[
		0\leq\chi\leq1,\qquad \chi=1\text{ 于 }p\text{ 的某邻域},\qquad
		\operatorname{supp}\chi\text{ 紧且包含于 }U.
	\]
	由\cref{g1:thm:bundle-metric-existence}给切丛选一个光滑正定度量 \(g\)。给定定向后，\cref{g1:prop:metric-volume-integration}给出正体积形式 \(\operatorname{vol}_g\)。令
	\[
		\omega_0=\chi\operatorname{vol}_g,\qquad c=\int_M\omega_0=\int_M\chi\mu_g.
	\]
	\(\omega_0\) 光滑紧支，支集包含于 \(\operatorname{supp}\chi\)。密度 \(\chi\mu_g\) 非负且在 \(p\) 附近为正，故由\cref{g1:thm:integral-well-defined}的严格正性，\(0<c<\infty\)。即使 \(p\) 位于边界，它的开邻域也含有半空间内部的一块正体积区域，所以严格正性仍成立。于是 \(\omega=c^{-1}\omega_0\) 具有所需支集且积分为 \(1\)。

	当 \(n=0\) 时，\(M\) 为离散流形。选 \(p\in U\)，令 \(\varepsilon_p\in\{1,-1\}\) 是该点的定向符号，定义函数
	\[
		\omega(q)=\begin{cases}\varepsilon_p,&q=p,\\0,&q\ne p.\end{cases}
	\]
	离散流形上的函数都是光滑的，单点支集紧且包含于 \(U\)。按零维积分的定义，\(\int_M\omega=\varepsilon_p\omega(p)=\varepsilon_p^2=1\)。若把负向点的函数值误取成 \(1\)，积分就会是 \(-1\)，因此这一符号不能省略。

	不定向的情形只需保留密度构造：当 \(n\geq1\) 时，取上面的 \(g,\chi\)，再令
	\[
		\mu=\frac{\chi\mu_g}{\int_M\chi\mu_g}.
	\]
	分母有限且严格为正，所得密度非负、光滑、紧支并满足要求。零维时在 \(p\) 处取密度值 \(1\)，其余点取 \(0\)，积分同样为 \(1\)，且不需要定向符号。
\end{solution}

\begin{exercise}\label{g1:ex:mobius-double-cover-integral}
	令
	\[
		B=(\mathbb R\times[-1,1])/\bigl((t,u)\sim(t+1,-u)\bigr),\qquad
		C=(\mathbb R/2\mathbb Z)\times[-1,1].
	\]
	自然映射 \(p:C\to B\)、\(p([t]_2,u)=[t,u]\) 是一个二重覆盖，其非平凡覆盖变换为 \(\tau([t]_2,u)=([t+1]_2,-u)\)。设 \(a\) 是条带上的光滑实值函数，满足 \(a(t+1,-u)=a(t,u)\)，令 \(\mu\) 为 \(a(t,u)|\mathrm{d}t\,\mathrm{d}u|\) 在 \(B\) 上下降的密度。通过基本区域中的具体积分证明
	\[
		\int_C p^*\mu=2\int_B\mu,
	\]
	这里的密度拉回按覆盖的局部微分同胚定义。给 \(C\) 选取 \(\mathrm{d}t\wedge\mathrm{d}u\) 指定的定向，说明与 \(p^*\mu\) 对应的二形式 \(\omega=a(t,u)\,\mathrm{d}t\wedge\mathrm{d}u\) 在 \(a\) 不恒为零时不能下降到 \(B\)。这是否意味着 \(B\) 上不存在非零二形式？请给出明确例子说明。
\end{exercise}
\begin{solution}
	把给定恒等式应用两次，可得 \(a(t+2,u)=a(t,u)\)，所以该系数在圆柱上有定义。覆盖在每个小矩形上由相同的 \((t,u)\) 坐标给出局部微分同胚，拉回密度正是 \(a(t,u)|\mathrm{d}t\,\mathrm{d}u|\)。两个商空间都紧，故这些光滑密度均紧支，积分有限。

	分别取 \([0,1]\times[-1,1]\) 与 \([0,2]\times[-1,1]\) 作 \(B,C\) 的基本区域。可以把它们拆成商坐标图中的小矩形；识别的竖边和其余公共边均为二维零测集，所以积分分别为
	\[
		\int_B\mu=\int_0^1\int_{-1}^1a(t,u)\,\mathrm{d}u\,\mathrm{d}t,
		\qquad
		\int_Cp^*\mu=\int_0^2\int_{-1}^1a(t,u)\,\mathrm{d}u\,\mathrm{d}t.
	\]
	把后式分成 \(0\leq t\leq1\) 与 \(1\leq t\leq2\)。在第二块中先令 \(t=s+1\)，再令 \(u=-v\)，便得
	\begin{align*}
		\int_1^2\int_{-1}^1a(t,u)\,\mathrm{d}u\,\mathrm{d}t
		&=\int_0^1\int_{-1}^1a(s+1,-v)\,\mathrm{d}v\,\mathrm{d}s\\
		&=\int_0^1\int_{-1}^1a(s,v)\,\mathrm{d}v\,\mathrm{d}s
		=\int_B\mu.
	\end{align*}
	故得到所求的两倍关系。这对可正可负的 \(a\) 均成立，证明只用了基本区域中的一元换元；特别地 \(a=1\) 时，带的面积为 \(2\)，圆柱面积为 \(4\)。

	虽然 \(\tau^*(p^*\mu)=p^*\mu\)，对应的顶次形式却满足
	\[
		\tau^*\omega=a(t+1,-u)\,\mathrm{d}(t+1)\wedge\mathrm{d}(-u)
		=-a(t,u)\,\mathrm{d}t\wedge\mathrm{d}u=-\omega.
	\]
	若有 \(B\) 上的二形式 \(\eta\) 使 \(p^*\eta=\omega\)，由 \(p\circ\tau=p\) 应有 \(\tau^*\omega=\omega\)。两式合起来迫使 \(\omega=0\)，与 \(a\) 不恒为零矛盾。这里不能把圆柱上依靠定向建立的“密度对应顶次形式”不加修改地传到不可定向的带上：\(\tau\) 反转圆柱定向，恰好改变对应形式的符号。

	这不排除 \(B\) 上存在非零但有零点的二形式。例如圆柱上的
	\[
		\beta=\sin(\pi t)\,\mathrm{d}t\wedge\mathrm{d}u
	\]
	有周期 \(2\)，并且
	\[
		\tau^*\beta=\sin(\pi(t+1))\,\mathrm{d}t\wedge(-\mathrm{d}u)=\beta.
	\]
	因此它能够下降。具体地，在 \(B\) 的每个小图上任选一支局部逆 \(s\) 并定义 \(s^*\beta\)；另一支为 \(\tau\circ s\)，其拉回仍为 \(s^*\beta\)，故这些局部二形式粘合为整体光滑二形式。它在 \(t\) 非整数的点非零，在 \(t\) 为整数的点为零。不可定向性排除的是处处非零的顶次形式，而不是所有非零顶次形式。
\end{solution}

\begin{exercise}\label{g1:ex:stokes-noncompact-end}
	给半直线 \(M=[0,\infty)\) 取定向 \(\mathrm{d}t\)。设 \(f\) 在 \(M\) 上光滑到边界，且 \(\int_0^\infty|f'(t)|\,\mathrm{d}t<\infty\)。证明 \(f_\infty=\lim_{t\to\infty}f(t)\) 存在，并计算 \(\int_M\mathrm{d}f\) 与 \(\int_{\partial M}f\)。判断两者相等的充要条件，再用 \(f(t)=1-e^{-t}\) 说明：把 Stokes 定理的紧支假设改成“左端的积分绝对收敛”，并不能保持结论。
\end{exercise}
\begin{solution}
	由一元微积分基本定理，对每个 \(R>0\)，有
	\[
		f(R)=f(0)+\int_0^R f'(t)\,\mathrm{d}t.
	\]
	右边的积分因绝对可积而在 \(R\to\infty\) 时收敛，所以
	\[
		f_\infty=f(0)+\int_0^\infty f'(t)\,\mathrm{d}t,\qquad
		\int_M\mathrm{d}f=f_\infty-f(0).
	\]
	半直线只有边界点 \(0\)，其外向向量为 \(-\partial_t\)，故由外法向优先约定，\(0\) 带负定向。因此 \(\int_{\partial M}f=-f(0)\)，两个积分相等当且仅当 \(f_\infty=0\)。紧支函数在充分大的 \(t\) 上为零，自然满足这个条件。

	对 \(f(t)=1-e^{-t}\)，有 \(f(0)=0\)、\(f_\infty=1\) 以及 \(f'(t)=e^{-t}\)，于是
	\[
		\int_M\mathrm{d}f=1,\qquad \int_{\partial M}f=0.
	\]
	差值来自 \(f\) 在趋向无穷远时留下的极限值。若先在紧区间 \([0,R]\) 使用 Stokes 定理，等式始终是 \(\int_0^R f'=f(R)-f(0)\)；令 \(R\to\infty\) 时，不能直接丢去右端的 \(f(R)\)。
\end{solution}

\begin{exercise}\label{g1:ex:stokes-annulus-orientation}
	设 \(0<r<R\)，给环域
	\[
		A=\{(x,y)\in\mathbb R^2:r^2\leq x^2+y^2\leq R^2\}
	\]
	取平面标准定向，并令 \(\alpha=(x^2+y^2)(x\,\mathrm{d}y-y\,\mathrm{d}x)\)。分别沿两条边界计算 \(\int_{\partial A}\alpha\)，再直接计算 \(\int_A\mathrm{d}\alpha\)，核对 Stokes 公式。说明内圆的积分为什么取负号。
\end{exercise}
\begin{solution}
	沿半径为 \(a\) 的逆时针参数曲线 \((x,y)=(a\cos\theta,a\sin\theta)\)，有 \(x\,\mathrm{d}y-y\,\mathrm{d}x=a^2\,\mathrm{d}\theta\)，故 \(\alpha=a^4\,\mathrm{d}\theta\)。外圆的外向量指向半径增大方向，边界正向为逆时针，贡献 \(2\pi R^4\)。在内圆上，环域的外向量指向圆心，外法向优先约定给出顺时针方向，因此贡献 \(-2\pi r^4\)。于是
	\[
		\int_{\partial A}\alpha=2\pi(R^4-r^4).
	\]

	另一方面，令 \(h=x^2+y^2\)，利用外微分乘积法则得到
	\begin{align*}
		\mathrm{d}\alpha
		&=(2x\,\mathrm{d}x+2y\,\mathrm{d}y)\wedge(x\,\mathrm{d}y-y\,\mathrm{d}x)
		+h\,\mathrm{d}(x\,\mathrm{d}y-y\,\mathrm{d}x)\\
		&=2h\,\mathrm{d}x\wedge\mathrm{d}y+2h\,\mathrm{d}x\wedge\mathrm{d}y
		=4(x^2+y^2)\,\mathrm{d}x\wedge\mathrm{d}y.
	\end{align*}
	使用极坐标换元，面积因子为 \(s\)，所以
	\[
		\int_A\mathrm{d}\alpha
		=\int_0^{2\pi}\int_r^R4s^3\,\mathrm{d}s\,\mathrm{d}\theta
		=2\pi(R^4-r^4).
	\]
	两边相等。内圆的负号由环域的外侧在圆孔中这一事实决定；若把两条圆周都取逆时针，它们便不再同时具有环域诱导的边界定向。
\end{solution}

\begin{exercise}\label{g1:ex:stokes-simplex-volume}
	给标准单形 \(\Delta^n=\{t\in\mathbb R^n:t^i\geq0,\ \sum_i t^i\leq1\}\) 取标准定向，记其欧氏体积为 \(V_n\)，并约定 \(V_0=1\)。对 \(n\geq1\)，令
	\[
		\alpha=\sum_{i=1}^n(-1)^{i-1}t^i\,\mathrm{d}t^1\wedge\cdots\wedge\widehat{\mathrm{d}t^i}\wedge\cdots\wedge\mathrm{d}t^n.
	\]
	利用 Stokes 公式及\cref{g1:lem:simplex-boundary-orientation}证明 \(nV_n=V_{n-1}\)，从而求出 \(V_n\)。特别说明：在斜面 \(\sum_i t^i=1\) 上使用有序顶点参数时，积分为何恰好是 \(V_{n-1}\)，而不是再乘一个斜面的欧氏面积因子。
\end{exercise}
\begin{solution}
	逐项求外微分可得 \(\mathrm{d}\alpha=n\,\mathrm{d}t^1\wedge\cdots\wedge\mathrm{d}t^n\)。在坐标面 \(t^i=0\) 上，第 \(i\) 项的系数为零，其余各项包含拉回后为零的 \(\mathrm{d}t^i\)，所以 \(\alpha\) 在这些面上的拉回全部为零。

	剩下的是与顶点 \(0\) 相对的面，其有序顶点为 \(e_1,\ldots,e_n\)，在诱导边界中取正号。令 \(s=(s^1,\ldots,s^{n-1})\in\Delta^{n-1}\)，该面的有序仿射参数为
	\[
		\delta_0(s)=\left(1-\sum_{j=1}^{n-1}s^j,s^1,\ldots,s^{n-1}\right).
	\]
	记 \(t=\delta_0(s)\)。形式 \(\alpha\) 是向量 \(t=\sum_i t^i\partial_{t^i}\) 与标准体积形式的缩并，因此
	\[
		(\delta_0^*\alpha)(\partial_{s^1},\ldots,\partial_{s^{n-1}})
		=\det(t,e_2-e_1,\ldots,e_n-e_1).
	\]
	因为 \(t=e_1+\sum_{j=1}^{n-1}s^j(e_{j+1}-e_1)\)，从行列式第一列减去其余列的相应线性组合，便得到 \(\det(e_1,e_2-e_1,\ldots,e_n-e_1)=1\)。所以
	\[
		\delta_0^*\alpha=\mathrm{d}s^1\wedge\cdots\wedge\mathrm{d}s^{n-1},\qquad
		\int_{\Delta^n}\mathrm{d}\alpha
		=\int_{\Delta^{n-1}}\delta_0^*\alpha=V_{n-1}.
	\]
	当 \(n=1\) 时，同一计算表示在正向端点 \(1\) 上评价零形式 \(\alpha=t^1\)，结果为 \(1=V_0\)。于是对所有 \(n\geq1\)，都有 \(nV_n=V_{n-1}\)，归纳得到 \(V_n=1/n!\)。

	这里积分的对象是 \(\alpha\) 这个微分形式，其拉回已经同时处理系数和切向量的变化。它在斜面上的取值并不是单位欧氏面积形式；再次乘一个斜面面积因子，会把参数变化重复计算。
\end{solution}

\begin{exercise}\label{g1:ex:stokes-positive-volume-obstruction}
	设 \(M\) 是非空、紧致、定向且无边界的 \(n\geq1\) 维光滑流形，\(\mu\) 是处处为正的光滑 \(n\) 形式。证明不存在 \((n-1)\) 形式 \(\beta\) 使 \(\mu=\mathrm{d}\beta\)。据此说明标准球面 \(S^2\) 的正面积形式不是恰当形式。再说明为什么这一结论与欧氏空间中的恒等式 \(\mathrm{d}(x^1\,\mathrm{d}x^2\wedge\cdots\wedge\mathrm{d}x^n)=\mathrm{d}x^1\wedge\cdots\wedge\mathrm{d}x^n\) 不矛盾。
\end{exercise}
\begin{solution}
	先说明 \(\int_M\mu>0\)。在任一点 \(p\) 附近选正向坐标，则 \(\mu=a(x)\,\mathrm{d}x^1\wedge\cdots\wedge\mathrm{d}x^n\)，其中 \(a(p)>0\)。连续性给出一个相对紧的小坐标球，使 \(a\) 在其闭包上有正下界。取非负、非零且支集位于这个球中的光滑函数 \(\chi\)，并令 \(0\leq\chi\leq1\)。于是 \(\int_M\chi\mu>0\)。而 \((1-\chi)\mu\) 非负，其积分非负，故 \(\int_M\mu>0\)；这里也可直接调用\cref{g1:thm:integral-well-defined}的正性结论。

	若 \(\mu=\mathrm{d}\beta\)，由于 \(M\) 紧，\(\beta\) 自动紧支，Stokes 定理给出
	\[
		\int_M\mu=\int_M\mathrm{d}\beta
		=\int_{\partial M}\beta=0,
	\]
	与正性矛盾。标准球面满足这些假设，所以它的正面积形式不是恰当形式。

	欧氏空间 \(\mathbb R^n\) 非紧，上述原形式也不紧支，因此不能对整个 \(\mathbb R^n\) 直接使用这一紧支版本的 Stokes 定理。若限制到一个紧长方体，区域便有边界，其正体积恰好由边界积分给出。紧致与无边界在本题的矛盾中各有作用。
\end{solution}

\begin{exercise}\label{g1:ex:triangle-chain-stokes}
	在标准二单形
	\[
		\Delta^2=\{(x,y):x\ge0,\ y\ge0,\ x+y\le1\}
	\]
	上取定向 \(\mathrm{d}x\wedge\mathrm{d}y\) 和一形式 \(\omega=x\,\mathrm{d}y\)。不直接调用 Stokes 定理，分别计算 \(\int_{\Delta^2}\mathrm{d}\omega\) 和三个有向边界积分，验证两者相等。再写出三条有向边的端点符号，说明边界的边界为何抵消。
\end{exercise}
\begin{solution}
	因为 \(\mathrm{d}\omega=\mathrm{d}x\wedge\mathrm{d}y\)，直接用二重积分得
	\[
		\int_{\Delta^2}\mathrm{d}\omega
		=\int_0^1\int_0^{1-x}1\,\mathrm{d}y\,\mathrm{d}x
		=\int_0^1(1-x)\,\mathrm{d}x=\frac12.
	\]
	记 \(v_0=(0,0)\)、\(v_1=(1,0)\)、\(v_2=(0,1)\)。依照\cref{g1:lem:simplex-boundary-orientation}，有向边界为
	\[
		[v_1,v_2]-[v_0,v_2]+[v_0,v_1].
	\]
	分别用参数 \(0\le t\le1\) 写出三个按顶点次序定向的面映射：
	\[
		\delta_0(t)=(1-t,t),\qquad
		\delta_1(t)=(0,t),\qquad
		\delta_2(t)=(t,0).
	\]
	于是
	\[
		\delta_0^*\omega=(1-t)\,\mathrm{d}t,\qquad
		\delta_1^*\omega=0,\qquad
		\delta_2^*\omega=0.
	\]
	将各面的系数 \(+1,-1,+1\) 保留下来，得到
	\[
		\int_{\partial\Delta^2}\omega
		=\int_0^1(1-t)\,\mathrm{d}t-\int_0^1 0\,\mathrm{d}t+\int_0^1 0\,\mathrm{d}t
		=\frac12.
	\]
	第二条边虽然在本例中积分为零，它的负号仍不可删去：负号把 \(v_0\to v_2\) 改成实际边界上的 \(v_2\to v_0\)，三条边才形成逆时针绕行。

	对一条有向边，边界总是终点减起点。因此三条边再取边界，给出
	\[
		([v_2]-[v_1])-([v_2]-[v_0])+([v_1]-[v_0])=0.
	\]
	每个顶点各以正、负号出现一次。这既核实了本例的端点符号，也解释了 \(\partial^2=0\) 在二维单形上的具体抵消。
\end{solution}

\begin{exercise}\label{g1:ex:two-puncture-independent-periods}
	在 \(P=\mathbb R^2\setminus\{(-1,0),(1,0)\}\) 上定义
	\[
		\alpha_-=\frac{-y\,\mathrm{d}x+(x+1)\,\mathrm{d}y}{(x+1)^2+y^2},\qquad
		\alpha_+=\frac{-y\,\mathrm{d}x+(x-1)\,\mathrm{d}y}{(x-1)^2+y^2}.
	\]
	证明对任意 \(a,b\in\mathbb R\)，形式 \(a\alpha_-+b\alpha_+\) 闭；它恰当当且仅当 \(a=b=0\)。证明时分别沿两个缺失点周围的小圆积分，并说明在哪一步可以使用圆盘上的 Stokes 定理。
\end{exercise}
\begin{solution}
	每个 \(\alpha_\pm\) 都是\cref{g1:exm:angular-period}中角形式平移后的表达式。外微分与平移拉回相容，故 \(\mathrm{d}\alpha_-=\mathrm{d}\alpha_+=0\)；线性性说明任意所述线性组合都闭。

	固定 \(0<\varepsilon<1\)，取正向小圆
	\[
		\gamma_-(t)=(-1+\varepsilon\cos t,\varepsilon\sin t),\qquad
		\gamma_+(t)=(1+\varepsilon\cos t,\varepsilon\sin t),\qquad 0\le t\le2\pi.
	\]
	分别将这两条参数曲线代入各自中心的角形式，可得
	\[
		\gamma_-^*\alpha_-=\mathrm{d}t,\qquad
		\gamma_+^*\alpha_+=\mathrm{d}t,
	\]
	所以对应积分均为 \(2\pi\)。

	以 \((-1,0)\) 为圆心、\(\varepsilon\) 为半径的闭圆盘 \(D_-\) 不含 \((1,0)\)。虽然 \(D_-\) 并不包含于 \(P\)，但 \(\alpha_+\) 的同一公式在 \(D_-\) 的一个邻域上光滑且闭。因此可以在圆盘上对这个延伸后的一形式使用\cref{g1:thm:stokes}，得到
	\[
		\int_{\gamma_-}\alpha_+=\int_{D_-}\mathrm{d}\alpha_+=0.
	\]
	同理，以 \((1,0)\) 为中心的小圆盘上可以使用光滑的 \(\alpha_-\)，从而 \(\int_{\gamma_+}\alpha_-=0\)。不能对 \(\alpha_-\) 在 \(D_-\) 上直接使用 Stokes，因为该形式恰在圆盘中心没有定义；另一个中心处的情形相同。

	若 \(a\alpha_-+b\alpha_+=\mathrm{d}f\)，\cref{g1:lem:exact-path-period}说明它沿两条闭曲线的积分都为零。于是
	\[
		0=\int_{\gamma_-}(a\alpha_-+b\alpha_+)=2\pi a,\qquad
		0=\int_{\gamma_+}(a\alpha_-+b\alpha_+)=2\pi b,
	\]
	故 \(a=b=0\)。反向若两个系数都为零，该形式就是 \(\mathrm{d}0\)，当然恰当。两个小圆分别检测两个系数，说明这两个闭形式代表的 de Rham 上同调类线性无关；此处没有断言它们已张成全部上同调。
\end{solution}

\begin{exercise}\label{g1:ex:cylinder-first-de-rham}
	令 \(C=S^1\times\mathbb R\)，\(\pi:C\to S^1\) 为第一因子投影，\(j:S^1\to C\) 为 \(j(u)=(u,0)\)。令 \(\sigma\) 是标准正向单位圆上的角形式，满足对 \(q(t)=(\cos t,\sin t)\) 有 \(q^*\sigma=\mathrm{d}t\)。只用本章已经证明的同伦公式和圆周周期判据，证明
	\[
		H_{\mathrm{dR}}^1(C)\cong\mathbb R,
	\]
	并证明 \([\pi^*\sigma]\) 是一个基。明确写出一个由周期给出的线性同构，不使用 de Rham 比较定理。
\end{exercise}
\begin{solution}
	定义圆柱收缩到中间圆周的光滑同伦
	\[
		H:[0,1]\times C\longrightarrow C,\qquad H(t,(u,s))=(u,(1-t)s).
	\]
	它满足 \(H_0=\operatorname{id}_C\)、\(H_1=j\circ\pi\)。对任意闭一形式 \(\eta\)，\cref{g1:thm:homotopy-formula}给出
	\[
		\pi^*j^*\eta-\eta=\mathrm{d}(K_H\eta).
	\]
	圆周上的一形式 \(j^*\eta\) 根据\cref{g1:prop:circle-exactness-period}可唯一写成
	\[
		j^*\eta=c\,\sigma+\mathrm{d}h,\qquad
		c=\frac1{2\pi}\int_{S^1}j^*\eta,
	\]
	其中 \(h\) 在相差常数的意义下唯一。拉回并整理，得到
	\[
		\eta=c\,\pi^*\sigma+\mathrm{d}\bigl(h\circ\pi-K_H\eta\bigr).
	\]
	因此每个闭一形式都与 \(c\,\pi^*\sigma\) 相差一个恰当形式，证明了这个类的张成性。

	具体定义
	\[
		\Phi:H_{\mathrm{dR}}^1(C)\longrightarrow\mathbb R,\qquad
		\Phi([\eta])=\frac1{2\pi}\int_{S^1}j^*\eta.
	\]
	若把 \(\eta\) 换成 \(\eta+\mathrm{d}f\)，积分增加的是 \(\int_{S^1}\mathrm{d}(f\circ j)=0\)，所以 \(\Phi\) 与代表元无关；积分的线性性说明它线性。上面的分解表明 \(\Phi([\eta])=0\) 时 \(\eta\) 恰当，故 \(\Phi\) 单射。又因为 \(j^*\pi^*\sigma=\sigma\) 且 \(\int_{S^1}\sigma=2\pi\)，有 \(\Phi([\pi^*\sigma])=1\)，故 \(\Phi\) 满射，逆映射就是 \(c\mapsto c[\pi^*\sigma]\)。

	这同时说明 \([\pi^*\sigma]\ne0\)，所以它确实构成一个基。整个计算只使用闭形式模恰当形式的定义、显式同伦算子及圆周的一元周期积分；无需先知道圆柱的奇异上同调。
\end{solution}

\begin{exercise}\label{g1:ex:no-smooth-ball-retraction}
	设 \(n\ge2\)，\(\overline B^n=\{x\in\mathbb R^n:\lVert x\rVert\le1\}\)。证明不存在光滑映射
	\[
		r:\overline B^n\longrightarrow S^{n-1}
	\]
	使 \(r(x)=x\) 对所有 \(x\in S^{n-1}=\partial\overline B^n\) 成立。这里光滑性包括边界处的局部光滑延拓。证明应使用球面的一个正体积形式及 Stokes 定理。
\end{exercise}
\begin{solution}
	给闭球以标准定向 \(\mathrm{d}x^1\wedge\cdots\wedge\mathrm{d}x^n\)，给球面以外法向优先的边界定向。记 \(i:S^{n-1}\hookrightarrow\overline B^n\) 为包含映射。在球面上定义
	\[
		\sigma=\left.\left(\sum_{a=1}^n(-1)^{a-1}x^a\,
		\mathrm{d}x^1\wedge\cdots\wedge\widehat{\mathrm{d}x^a}\wedge\cdots\wedge\mathrm{d}x^n\right)\right|_{S^{n-1}}.
	\]
	这里的限制指把环境形式拉回球面。若 \(v_1,\ldots,v_{n-1}\) 是球面上一组正向切向基，则外向单位法向为 \(x\)，从而
	\[
		\sigma_x(v_1,\ldots,v_{n-1})
		=(\mathrm{d}x^1\wedge\cdots\wedge\mathrm{d}x^n)_x(x,v_1,\ldots,v_{n-1})>0.
	\]
	所以 \(\sigma\) 是光滑正体积形式；球面紧且非空，积分的正性给出 \(\int_{S^{n-1}}\sigma>0\)。同时 \(\mathrm{d}\sigma=0\)，因为这是 \((n-1)\) 维流形上的 \(n\) 形式，必为零。

	假设存在题述映射 \(r\)。则 \(r\circ i=\operatorname{id}_{S^{n-1}}\)，所以
	\[
		i^*(r^*\sigma)=\sigma,\qquad
		\mathrm{d}(r^*\sigma)=r^*(\mathrm{d}\sigma)=0.
	\]
	形式 \(r^*\sigma\) 在整个闭球上光滑；闭球紧，满足\cref{g1:thm:stokes}的支集条件。于是
	\[
		0<\int_{S^{n-1}}\sigma
		=\int_{\partial\overline B^n}i^*(r^*\sigma)
		=\int_{\overline B^n}\mathrm{d}(r^*\sigma)=0,
	\]
	矛盾。因此这样的光滑映射不存在。

	这里 \(\mathrm{d}\sigma=0\) 是球面上按自身维数计算的结论。上面括号中的环境 \((n-1)\) 形式在 \(\mathbb R^n\) 上的外微分其实为 \(n\,\mathrm{d}x^1\wedge\cdots\wedge\mathrm{d}x^n\)，并不为零；把它限制到球面后才变为零。这一区别保证证明没有把球面上的闭性误用成整个球内的闭性。
\end{solution}


% !TeX root = ../../Book1.tex
% 纲要标题：## 第15章　联络、曲率与平行移动
% 纲要位置：几何学计划纲要.md，第 690 行。
% 正文状态：初稿；尚未正式审核。

% 纲要内容开始
% 
% > 本章目标与证明范围：从“不同点的向量如何比较”构造联络，再以曲率测量路径依赖。完整证明局部公式、变换律、Bianchi 和平坦局部模型；不把曲率只写成一个矩阵公式。
% 
% 纲要内容结束

\chapter{联络、曲率与平行移动}
\label{g1:ch:15}
\BookChapterNumber{g1:ch:15}

\begin{chapterabstract}
	不同点的纤维没有默认的比较方式，联络为这种比较提供一阶规则。我们先从 Leibniz 法则构造局部联络矩阵，证明整体存在、换标架规律与仿射空间结构，再把联络化为沿路径的线性常微分方程。平行移动反过来恢复协变导数，并把度量相容解释成长度保持。曲率由两次协变微分的不交换性产生，其结构方程和小环路展开连接局部计算与实际运输；Bianchi 恒等式表达这些变化之间的相容关系。对切丛，还能用 Lie 括号定义挠率，区分平行标架与坐标标架。最后由平坦方程逐方向构造平行标架，用连续同伦的网格论证证明路径无关性，并由基本群表示记录仍可能存在的整体单值。
\end{chapterabstract}

\begin{learninggoals}
	\begin{itemize}
		\item 理解联络两个输入的不同线性要求，证明局部矩阵公式、整体存在及两联络之差的张量性，构造对偶与张量上的诱导联络。
		\item 完整建立沿路径的平行移动，掌握拼接、反向、参数依赖和恢复联络的公式，并证明度量相容与运输保长等价。
		\item 证明曲率的张量性、结构方程及换标架共轭规律，用小环路计算解释曲率，并推导 Bianchi 恒等式。
		\item 构造切丛联络的挠率，推导第一结构方程，区分无挠性、平坦性和度量相容性，判断平行标架何时可由坐标实现。
		\item 从曲率为零构造局部平行标架，证明平坦运输的同伦不变性与单连通区域上的路径无关性，用单值表示解释整体障碍。
	\end{itemize}
\end{learninggoals}

在 Euclid 空间中，我们习惯把两个不同点的向量平移到一起相减。第八、九章已经提醒我们：曲面的切向量属于不同切平面，必须先说明如何比较，才谈得上“方向没有改变”。第十二章将这种情形推广到一般向量丛，第十四章又建立了不依赖坐标的微分形式运算。现在要把这两条线连接起来：既允许纤维随点变化，又保留微分的乘法法则。

\ProseParagraphBreak

局部标架能暂时把邻近纤维认作同一个向量空间，但这种认法随标架选择而变。联络矩阵补偿标架的变化；沿路径解相应方程，便得到由联络指定的比较。两条路径的比较结果不一致时，局部的小环路由曲率控制，较大的绕行还可能涉及基本群。圆周上的例子将展示曲率处处为零、绕行运输仍非平凡的联络，从而把局部曲率与整体单值分开。

\ProseParagraphBreak

本章默认底流形有限维、Hausdorff、第二可数，向量丛为有限秩实光滑丛；除明确限于内点的坐标实现及小方环计算外，允许流形有边界。需要的前置工具是\chapref{g1:ch:12}的局部标架、张量构造和单位分解，\secref{g1:sec:1301}的常微分方程存在唯一性与参数依赖，以及\chapref{g1:ch:14}的外微分和局部恰当性。\secref{g1:sec:1505}还会回顾\chapref{g1:ch:05}的路径同伦、基本群与覆盖。调用时均重申具体条件，不要求读者记住前章每个技术细节。

\ProseParagraphBreak

本章的联络构造、局部公式、运输、曲率恒等式与平坦局部模型均在正文证明。延伸阅读可查 Tu 的第一版：第6.2节 的 Torsion and Curvature，第10--11节 的向量丛联络、存在性及联络、曲率、挠率形式，第13.1节 的 Covariant Differentiation Along a Curve，以及 第14.5--14.7节 的 Parallel Translation 及其存在性、度量情形\cite[第6.2、10--11、13.1、14.5--14.7节]{Tu2017Connections}。本书继续使用行标架 \(e\)、列分量 \(u\)、\(\nabla(eu)=e(\mathrm{d}u+Au)\) 的约定；对读其他文献时，应先核对矩阵放在哪一侧，再比较符号。

% !TeX root = ../../Book1.tex
% 纲要标题：### 15.1 丛联络及其空间
% 纲要位置：几何学计划纲要.md，第 694 行。
% 正文状态：初稿；尚未正式审核。

% 纲要内容开始
% 
% * Leibniz 规则与局部联络矩阵
% * 联络的整体存在
% * 两联络之差与仿射空间结构
% * 对偶、张量及外幂上的诱导联络
% 
% 纲要内容结束

\section{丛联络及其空间}
\label{g1:sec:1501}

一个截面 \(s\) 在不同点取值于不同纤维，所以表达式 \(s(q)-s(p)\) 通常没有意义。局部标架能把这些值写成同一个 \(\mathbb R^r\) 中的列向量，于是可以对分量求导；然而标架自己也可能随基点变化，只对分量求导会依赖这次选择。曲面上的协变导数曾用“先在环境空间求导，再投影回切空间”处理这一问题。一般向量丛未指定环境空间，我们需要直接规定一种微分运算，使它在换标架后仍表示同一个几何对象。

\ProseParagraphBreak

本节讨论有限秩实光滑向量丛 \(E\to M\)，其中 \(M\) 有限维、Hausdorff、第二可数，允许有边界。局部平凡化和截面沿用第12章的定义：标架写成行 \(e=(e_1,\ldots,e_r)\)，截面写成 \(s=e u\)，其中 \(u\) 是列分量。边界处的光滑函数与截面均按局部延拓定义；联络作用于完整的切空间，不要求求导方向切于边界。

\subsection{Leibniz 规则与局部联络矩阵}
\label{g1:subsec:1501:01}

对函数倍乘的截面 \(fs\) 求导时，一部分变化来自 \(f\)，另一部分来自 \(s\)。前一部分应当仍由普通方向导数 \(Xf\) 描述。这个要求给出联络的乘积法则。方向 \(X\) 则只指定在哪个切向量上求导，不应再对 \(X\) 的系数求导；因此两个输入承担不同的线性条件。

\begin{definition}[丛联络]\label{g1:def:bundle-connection}
	设 \(M\) 为有限维、Hausdorff、第二可数光滑流形，允许有边界，\(E\to M\) 为有限秩实光滑向量丛。一个实双线性映射
	\[
		\nabla:\mathfrak X(M)\times\Gamma(E)\longrightarrow\Gamma(E),\qquad (X,s)\longmapsto\nabla_Xs,
	\]
	若对每个 \(f\in C^\infty(M)\) 满足
	\[
		\nabla_{fX}s=f\nabla_Xs,\qquad
		\nabla_X(fs)=X(f)s+f\nabla_Xs,
	\]
	就称为 \(E\) 上的\term[conglianluo]{丛联络}[vector bundle connection]，简称联络。\(\nabla_Xs\) 称为 \(s\) 沿 \(X\) 的协变导数。这里的丛联络都是线性联络，即遵守所写的实线性条件。
\end{definition}

第一个公式是对方向输入的 \(C^\infty(M)\) 线性，第二个公式是对截面输入的 Leibniz 规则，不能交换二者的职责。例如在平凡实线丛 \(M\times\mathbb R\) 的标准标架 \(e\) 中，\(\nabla_X(fe)=X(f)e\) 就满足定义。若把第二个输入也要求成 \(C^\infty(M)\) 线性，便会删掉函数本身的变化，无法包含这个最基本的例子。

\begin{proposition}[联络的局部性]\label{g1:prop:connection-locality}
	设 \(M\) 为有限维、Hausdorff、第二可数光滑流形，允许有边界，\(E\to M\) 为有限秩实光滑向量丛，\(\nabla\) 为其联络。则有以下结论。
	\begin{enumerate}
		\item \(\nabla_Xs\) 在开集 \(U\) 上的限制，只依赖 \(X,s\) 在 \(U\) 上的限制。因此 \(\nabla\) 唯一地作用于任意开集上的局部向量场和局部截面，并满足相同的线性与乘积规则。
		\item 对 \(p\in M\)，\((\nabla_Xs)(p)\) 对方向场只依赖 \(X_p\)。给定任意 \(v\in \mathrm{T}_pM\) 及 \(p\) 附近的局部截面 \(s\)，可以良定义
		\[
			\nabla_v s\coloneqq(\nabla_Xs)(p),\qquad X_p=v.
		\]
		它属于 \(E_p\)，且对 \(v\) 实线性。
	\end{enumerate}
\end{proposition}
\begin{proof}
	先设整体截面 \(s\) 在 \(p\) 的开邻域 \(U\) 中为零。由\cref{g1:cor:smooth-cutoff-closed}，可取在 \(p\) 附近等于 \(1\)、支集紧且包含于 \(U\) 的光滑函数 \(\chi\)。因为 \(\chi s=0\)，Leibniz 规则给出
	\[
		0=\nabla_X(\chi s)=X(\chi)s+\chi\nabla_Xs.
	\]
	在 \(\chi=1\) 的邻域中，\(X(\chi)=0\)，所以 \(\nabla_Xs=0\)。若整体向量场 \(X\) 在 \(U\) 上为零，则 \(\chi X=0\)，由方向输入的函数线性得到 \(0=\nabla_{\chi X}s=\chi\nabla_Xs\)，结论相同。作差便知局部相同的输入给出局部相同的输出。

	现设 \(X,s\) 只在 \(U\) 上定义。对每个 \(p\in U\) 选上述 \(\chi\)，将 \(\chi X,\chi s\) 在 \(U\) 外补零。所得整体场和截面光滑：在 \(U\) 内按乘法光滑，在 \(U\) 外每点附近因支集包含于 \(U\) 而恒为零。用整体联络求导后，在 \(\chi=1\) 的邻域内作为局部定义。不同截断或延拓在较小邻域中相同，刚证明的局部性保证结果也相同，故各点的结果粘合为 \(U\) 上的光滑截面。对局部函数同样用截断延拓，便能从整体规则得到局部规则。唯一性来自每点附近都可以这样替换成整体输入。

	最后，在 \(p\) 附近取坐标 \(x^1,\ldots,x^n\)，写 \(X=\sum_iX^i\partial_i\)。已建立的局部规则给出
	\[
		(\nabla_Xs)(p)=\sum_iX^i(p)(\nabla_{\partial_i}s)(p).
	\]
	因此 \(X_p=0\) 时输出为零，两个方向场在 \(p\) 相同时也给出相同值。任意 \(v=\sum_iv^i\partial_i|_p\) 可以由局部常分量场 \(\sum_iv^i\partial_i\) 实现，再乘紧支截断延成整体场，故所有切向量都能作为输入。边界图的坐标向量仍构成完整的 \(n\) 维切空间基，这一构造没有使用保持 \(M\) 的流，也没有限制 \(v\) 的法向分量。线性直接由公式得到。
\end{proof}

点处求值的两个输入仍不对称。\(\nabla_v s\) 可以只记方向向量 \(v\)，却通常不能只记纤维向量 \(s(p)\)。例如在 \(\mathbb R\) 上的平凡线丛中，标准联络使截面 \(s(x)=xe\) 满足 \(s(0)=0\)、\(\nabla_{\partial_x|_0}s=e\)，而零截面有相同的点值，导数却为零。因此联络并没有给每个 \(v\) 一个可直接作用于 \(E_p\) 的线性自同态；它还需要知道截面在附近如何变化。

\begin{definition}\label{g1:def:bundle-valued-form}
	设 \(M\) 为有限维、Hausdorff、第二可数光滑流形，允许有边界，\(E\to M\) 为有限秩实光滑向量丛。记
	\[
		\Omega^k(M;E)\coloneqq\Gamma(\Lambda^k\mathrm{T}^*M\otimes E),\qquad k\geq0,
	\]
	其元素称为 \(E\) 值微分形式，也称\term[congzhiweifenxingshi]{丛值微分形式}[bundle-valued differential form]，\(\Omega^0(M;E)=\Gamma(E)\)。\symidx{\(\Omega^k(M;E)\)：丛值微分形式空间}对 \(k=1\)，它在每个点是从 \(\mathrm{T}_pM\) 到 \(E_p\) 的线性映射，随基点光滑变化。
\end{definition}

这里的丛已由\cref{g1:prop:tensor-bundle-constructions}构造，\(\mathrm{T}_p^*M\otimes E_p\simeq\operatorname{Hom}(\mathrm{T}_pM,E_p)\) 的识别来自\cref{g1:prop:tensor-hom-duality}及因子交换。固定截面 \(s\)，上一命题说明 \(v\mapsto\nabla_v s\) 给出一个 \(E\) 值一形式：在坐标中，它写成 \(\sum_i\mathrm{d}x^i\otimes\nabla_{\partial_i}s\)，所以光滑。记此一形式为 \(\nabla s\)，联络也就可以写成
\[
	\nabla:\Gamma(E)\longrightarrow\Omega^1(M;E),\qquad
	\nabla(fs)=\mathrm{d}f\otimes s+f\nabla s.
\]
反过来，一个满足此式的实线性映射通过 \(\nabla_Xs=(\nabla s)(X)\) 恢复定义中的两条规则。因此两种写法等价：前者突出求导方向，后者把各方向的结果放在同一个一形式中。

\ProseParagraphBreak

在局部标架 \(e=(e_1,\ldots,e_r)\) 中，各 \(\nabla e_b\) 是 \(E\) 值一形式，因而唯一写成 \(\nabla e_b=\sum_a A^a{}_b\otimes e_a\)。把一形式 \(A^a{}_b\) 排成矩阵 \(A\)，称为该标架中的\term[lianluojuzhen]{联络矩阵}[connection matrix]。为使行标架计算简洁，约定对一形式列 \(\eta\) 写 \(e\eta=\sum_a\eta^a\otimes e_a\)；于是 \(\nabla e=eA\) 表示刚才的逐列等式。它并没有改变 \(E\) 值形式中两个张量因子的类型。

\begin{proposition}[联络的局部表示]\label{g1:prop:connection-local-data}
	设 \(M\) 为有限维、Hausdorff、第二可数光滑流形，允许有边界，\(E\to M\) 为秩 \(r\) 的实光滑向量丛。联络 \(\nabla\) 在局部标架 \(e\) 中由一个 \(r\times r\) 的光滑一形式矩阵 \(A\) 唯一决定，具体为
	\begin{equation}
		\label{g1:eq:connection-local-form}
		\nabla e=eA,\qquad
		\nabla(eu)=e(\mathrm{d}u+A u),\qquad
		\nabla_X(eu)=e\bigl(X(u)+A(X)u\bigr).
	\end{equation}
	这里 \(\mathrm{d}u\) 与 \(X(u)\) 均逐分量计算。若换标架为 \(e'=eg\)，其中 \(g:U\to\operatorname{GL}(r,\mathbb R)\) 光滑，则新矩阵满足
	\begin{equation}
		\label{g1:eq:connection-frame-change}
		A'=g^{-1}Ag+g^{-1}\mathrm{d}g.
	\end{equation}
	反过来，若一个标架覆盖上的光滑一形式矩阵在交叠处满足此规律，则\eqref{g1:eq:connection-local-form}唯一地定义一个整体联络。
\end{proposition}
\begin{proof}
	把 \(s=\sum_bu^b e_b\) 代入 Leibniz 规则，得到
	\[
		\nabla s=\sum_b\mathrm{d}u^b\otimes e_b+\sum_{a,b}u^bA^a{}_b\otimes e_a
		=e(\mathrm{d}u+A u).
	\]
	对 \(X\) 求值便得到方向公式。反向，在一张有标架的开集上，任意光滑矩阵一形式 \(A\) 由此定义光滑输出；\(X(fu)=X(f)u+fX(u)\) 和 \(A(fX)=fA(X)\) 直接保证联络的两条规则。

	换标架时，逐列使用乘积法则可得
	\[
		\nabla e'=\nabla(eg)=e(Ag+\mathrm{d}g)=e'g^{-1}(Ag+\mathrm{d}g).
	\]
	由于 \(\nabla e'=e'A'\)，这就是\eqref{g1:eq:connection-frame-change}。同一截面的分量满足 \(u=gu'\)，于是该规律还给出
	\[
		e'(\mathrm{d}u'+A'u')
		=e\bigl(g\,\mathrm{d}u'+\mathrm{d}g\,u'+Ag u'\bigr)
		=e\bigl(\mathrm{d}(gu')+A(gu')\bigr).
	\]
	因此满足变换规律的局部矩阵对同一截面给出相同的一形式，能够逐点粘合。光滑性和联络规则都是局部性质，故粘合结果是整体联络；各张图中的公式又完全确定它，所以唯一。边界图中的微分按局部光滑延拓计算，\cref{g1:prop:boundary-smooth-derivative}保证链式法则及所得导数不依赖延拓，上述论证同样成立。
\end{proof}

若原来采用第12章的记号 \(e_j=e_i g_{ij}\)，这里的公式就是 \(A_j=g_{ij}^{-1}A_i g_{ij}+g_{ij}^{-1}\mathrm{d}g_{ij}\)。注意它含有 \(g\) 的导数。联络矩阵不是一个单纯按共轭变换的张量，\(A=0\) 也不是换标架后保持不变的断言。

\begin{example}
	在平凡实线丛中，令标准标架 \(e\) 下的联络为 \(\nabla(fe)=\mathrm{d}f\otimes e\)，于是 \(A=0\)。取另一标架 \(e'=e e^h\)，其中 \(h\) 是任意光滑实函数，则 \(g=e^h\)，从而
	\[
		A'=e^{-h}\mathrm{d}(e^h)=\mathrm{d}h.
	\]
	同一截面 \(fe\) 在新标架中的分量为 \(u'=e^{-h}f\)。因此
	\[
		e'\bigl(\mathrm{d}(e^{-h}f)+(\mathrm{d}h)e^{-h}f\bigr)=e\,\mathrm{d}f.
	\]
	对分量求导时出现的额外项，恰被联络矩阵抵消。新标架中的矩阵非零，并不表示我们更换了联络。
\end{example}

\subsection{联络的整体存在}
\label{g1:subsec:1501:02}

在一张平凡化中，把联络矩阵取为零就得到联络。困难在于不同平凡化的“零矩阵”一般不满足变换公式，所以不能直接宣称它们粘成整体对象。单位分解提供另一种做法：先把各个局部联络的输出乘以权重，再相加。Leibniz 规则中需要保留的普通导数项，正好由权重之和等于 \(1\) 来保证。

\begin{theorem}[联络的存在]\label{g1:thm:connection-existence}
	设 \(M\) 为有限维、Hausdorff、第二可数光滑流形，允许有边界，\(E\to M\) 为有限秩实光滑向量丛。则 \(E\) 上存在光滑联络。
\end{theorem}
\begin{proof}
	取局部标架覆盖 \(\{U_i,e_i\}\)，在每个 \(U_i\) 上令
	\[
		\nabla^{(i)}(e_i u)=e_i\,\mathrm{d}u.
	\]
	这是局部联络。由\cref{g1:thm:smooth-partition-unity}，存在从属于 \(\{U_i\}\) 的光滑单位分解 \(\{\rho_i\}\)，满足支集族局部有限、\(\operatorname{supp}\rho_i\subseteq U_i\) 及 \(\sum_i\rho_i=1\)。对整体截面 \(s\)，定义
	\[
		\nabla s=\sum_i\rho_i\nabla^{(i)}(s|_{U_i}),
	\]
	其中每项在 \(U_i\) 外补零。这样的补零光滑：任何 \(U_i\) 外的点都有一个避开闭集 \(\operatorname{supp}\rho_i\) 的邻域，在此邻域该项恒为零。每点附近只有有限个非零项，所以和是光滑的 \(E\) 值一形式，且对 \(s\) 实线性。

	对 \(f\in C^\infty(M)\)，在每个点的有限求和邻域内，有
	\begin{align*}
		\nabla(fs)
		&=\sum_i\rho_i\bigl(\mathrm{d}f\otimes s+f\nabla^{(i)}s\bigr)\\
		&=\left(\sum_i\rho_i\right)\mathrm{d}f\otimes s+f\sum_i\rho_i\nabla^{(i)}s
		=\mathrm{d}f\otimes s+f\nabla s.
	\end{align*}
	于是按一形式对方向求值，就得到所需联络。证明使用的局部标架与单位分解在带边界流形上也存在，故包含边界情形。秩零时所有截面和输出均为零，结论也成立。
\end{proof}

支集包含于 \(U_i\) 与局部有限性各自有明确作用：前者使局部输出能够光滑补零，后者使求和、求导都成为邻域内的有限运算。权重在此乘在联络的输出上，验证乘积法则时不对 \(\rho_i\) 求导。构造依赖所选局部标架与单位分解，存在定理没有给出一个由向量丛本身唯一指定的联络。

\begin{example}
	Möbius 实线丛的局部过渡函数可以取为 \(1\) 或 \(-1\)，在每个交叠连通分支上为常数，见\cref{g1:exm:mobius-line-bundle}。因此 \(\mathrm{d}g_{ij}=0\)，在这些局部标架中同时取 \(A_i=0\) 就满足联络变换公式，直接给出整体联络。这个丛没有全局处处非零截面，却仍有联络。由此可见，给出一种相容的求导规则，并不等于给出一个全局标架。
\end{example}

\subsection{两联络之差与仿射空间结构}
\label{g1:subsec:1501:03}

单个联络对截面求导时包含函数的一阶变化。若对同一个截面使用两个联络，再把结果相减，普通微分项就会消掉。留下的运算只观察该截面在一点的值，因而成为真正的张量。这也解释了为什么联络的矩阵变换带一个附加项，而两个矩阵的差没有。

\begin{proposition}[联络的仿射空间]\label{g1:prop:connection-affine-space}
	设 \(M\) 为有限维、Hausdorff、第二可数光滑流形，允许有边界，\(E\to M\) 为有限秩实光滑向量丛。两个联络 \(\nabla^0,\nabla^1\) 唯一确定
	\[
		B\in\Omega^1(M;\operatorname{End}E)
	\]
	使
	\[
		\nabla^1_Xs-\nabla^0_Xs=B(X)s
	\]
	对所有 \(X,s\) 成立。反过来，给定联络 \(\nabla\) 和任意这样的 \(B\)，公式 \((\nabla+B)_Xs=\nabla_Xs+B(X)s\) 定义一个联络。因此全部联络的集合 \(\operatorname{Conn}(E)\) 是以 \(\Omega^1(M;\operatorname{End}E)\) 为方向空间的非空仿射空间。\symidx{\(\operatorname{Conn}(E)\)：联络的仿射空间}
\end{proposition}
\begin{proof}
	令 \(D(X,s)=\nabla^1_Xs-\nabla^0_Xs\)。方向输入的函数线性由两个联络的相同性质得到；对截面输入，Leibniz 项相消，故
	\[
		D(fX,s)=fD(X,s),\qquad D(X,fs)=fD(X,s).
	\]
	先前的局部性保证 \(D\) 可作用于局部输入。在 \(p\) 附近取坐标及标架，写 \(X=\sum_iX^i\partial_i\)、\(s=\sum_a u^a e_a\)，便有
	\[
		D(X,s)(p)=\sum_{i,a}X^i(p)u^a(p)D(\partial_i,e_a)(p).
	\]
	右端只依赖 \(X_p\) 与 \(s(p)\)，并且对二者实双线性。任何 \(v\in \mathrm{T}_pM\)、\(w\in E_p\) 都能由局部常分量场及截面实现，再用紧支截断延成整体输入。因此定义 \(B_p(v)w=D(X,s)(p)\) 与选择无关，也唯一。

	各 \(D(\partial_i,e_a)\) 光滑，所以上式同时给出 \(B\) 的光滑局部系数，得到 \(B\in\Gamma(\mathrm{T}^*M\otimes\operatorname{End}E)\)。在一组标架中若两个联络矩阵为 \(A^0,A^1\)，则 \(B\) 的矩阵一形式就是 \(A^1-A^0\)。由\eqref{g1:eq:connection-frame-change}，换标架后
	\[
		B'=g^{-1}Bg,
	\]
	因为两个 \(g^{-1}\mathrm{d}g\) 项消掉了。这正是自同态值一形式所应有的变换规律。

	反向，\(B(X)s\) 对两个输入都为函数线性，因此将它加到联络上，方向线性不变，对截面的乘积法则仍只有原来的 \(X(f)s\) 项。故 \(\nabla+B\) 是联络。存在定理保证联络集合非空；上面的差与加法互相对应，满足
	\[
		(\nabla+B)+C=\nabla+(B+C),\qquad \nabla+0=\nabla,
	\]
	而任意两个联络之间有且只有一个 \(B\) 把前者移到后者。这就是仿射空间的加法公理，不要求方向空间有限维。
\end{proof}

选择一个参考联络 \(\nabla^0\) 后，任何联络都唯一写成 \(\nabla^0+B\)，从而可用向量空间 \(\Omega^1(M;\operatorname{End}E)\) 参数化。但参考联络是一个选择；更换它会平移全部参数。尤其不能把矩阵 \(A=0\) 当成与标架无关的“零联络”。在正维流形上的正秩丛中，零算子一般连 Leibniz 规则都不满足：它会把 \(\nabla_X(fs)\) 置为零，却遗漏可能非零的 \(X(f)s\)。

\ProseParagraphBreak

仿射组合则没有这一问题。如果 \(\nabla^0,\ldots,\nabla^m\) 是联络，光滑函数 \(f_0,\ldots,f_m\) 满足 \(\sum_j f_j=1\)，那么 \(\sum_j f_j\nabla^j\) 仍是联络，证明与存在定理中的乘积法则计算相同；此处不要求 \(f_j\) 非负。特别地，\((1-t)\nabla^0+t\nabla^1\) 对每个实数 \(t\) 都是联络。若 \(\dim M=0\) 或 \(\operatorname{rank}E=0\)，方向空间为零，联络只有一个，这与仿射描述也相容。

\subsection{对偶、张量及外幂上的诱导联络}
\label{g1:subsec:1501:04}

给 \(E\) 选择求导规则之后，它的对偶与张量构造不应再独立任选一套规则。例如 \(\lambda\in\Gamma(E^*)\) 与 \(s\in\Gamma(E)\) 配对得到普通函数 \(\lambda(s)\)；若对它求导，结果应与先分别求导再配对一致。这个要求决定了对偶联络的负号。张量积和外幂则应沿每个因子分别求导，再把结果相加。

\begin{proposition}[诱导联络]\label{g1:prop:induced-connections}
	设 \(M\) 为有限维、Hausdorff、第二可数光滑流形，允许有边界，\(E,F\to M\) 为有限秩实光滑向量丛，分别给定联络 \(\nabla^E,\nabla^F\)。存在以下唯一的诱导联络。
	\begin{enumerate}
		\item 对偶丛 \(E^*\) 上的联络满足
		\[
			X\bigl(\lambda(s)\bigr)=(\nabla^{E^*}_X\lambda)(s)+\lambda(\nabla^E_Xs),
			\qquad \lambda\in\Gamma(E^*),\ s\in\Gamma(E).
		\]
		\item 张量积丛 \(E\otimes F\) 上的联络满足
		\[
			\nabla^{E\otimes F}_X(s\otimes t)
			=(\nabla^E_Xs)\otimes t+s\otimes\nabla^F_Xt.
		\]
		重复此构造得到任意有限个张量因子上的联络；它与因子的排列及对偶求值产生的收缩相容。
		\item 对每个整数 \(k\geq1\)，外幂 \(\Lambda^kE\) 上的联络满足
		\[
			\nabla^{\Lambda^kE}_X(s_1\wedge\cdots\wedge s_k)
			=\sum_{j=1}^k s_1\wedge\cdots\wedge\nabla^E_Xs_j\wedge\cdots\wedge s_k.
		\]
		零次外幂是平凡实线丛，其诱导联络为 \(\nabla_Xf=Xf\)。对称幂上有把 \(\wedge\) 换成 \(\odot\) 的同类公式。
		\item 同态丛 \(\operatorname{Hom}(E,F)\) 上的联络满足
		\[
			(\nabla^{\operatorname{Hom}}_X\mathcal B)(s)
			=\nabla^F_X(\mathcal B s)-\mathcal B(\nabla^E_Xs),
			\qquad \mathcal B\in\Gamma(\operatorname{Hom}(E,F)).
		\]
	\end{enumerate}
\end{proposition}
\begin{proof}
	先构造对偶联络。给定 \(X,\lambda\)，令
	\[
		Q(s)=X\bigl(\lambda(s)\bigr)-\lambda(\nabla^E_Xs).
	\]
	若把 \(s\) 换成 \(fs\)，第一项中的 \(X(f)\lambda(s)\) 与第二项中的相同项相消，故 \(Q(fs)=fQ(s)\)。所以 \(Q\) 在一点处只依赖 \(s\) 在该点的值：可以先用紧支截断把局部截面延拓，再展开为局部标架的函数组合，得到 \(Q(s)(p)=\sum_a u^a(p)Q(e_a)(p)\)。这与\cref{g1:prop:connection-locality}中的局部化论证相同，所需的函数线性刚刚已经验证。于是 \(Q\) 唯一确定一个对偶截面，规定为 \(\nabla^{E^*}_X\lambda\)。

	它的光滑性和具体符号也可直接看出。令 \(\lambda=\sum_a\ell_a e^a\)，其中 \((e^a)\) 为 \((e_a)\) 的对偶标架，\(\ell\) 为列分量。若 \(\nabla^E\) 的矩阵为 \(A\)，对 \(s=e u\) 展开上式得到
	\[
		Q(s)=X(\ell^{\mathsf T}u)-\ell^{\mathsf T}\bigl(X(u)+A(X)u\bigr)
		=\bigl(X(\ell)-A(X)^{\mathsf T}\ell\bigr)^{\mathsf T}u.
	\]
	所以对偶联络的列分量为 \(\mathrm{d}\ell-A^{\mathsf T}\ell\)，对偶标架中的联络矩阵为 \(-A^{\mathsf T}\)。它对方向输入函数线性；把 \(\lambda\) 换成 \(f\lambda\)，公式给出 \(\nabla^{E^*}_X(f\lambda)=X(f)\lambda+f\nabla^{E^*}_X\lambda\)。故它确实为联络。求值恒等式不依赖标架，所以各图中的构造相容，且该恒等式决定了它在每个纤维向量上的值，保证唯一。

	再构造张量积联络。在局部标架 \((e_a)\)、\((h_b)\) 中，\(e_a\otimes h_b\) 构成 \(E\otimes F\) 的标架，见\cref{g1:thm:tensor-product-universal,g1:prop:tensor-bundle-constructions}。先规定
	\[
		\nabla_X(e_a\otimes h_b)
		=(\nabla^E_Xe_a)\otimes h_b+e_a\otimes\nabla^F_Xh_b,
	\]
	再对一般截面 \(w=\sum_{a,b}w^{ab}e_a\otimes h_b\) 按 \(X(w^{ab})\) 和 Leibniz 规则定义。基展开唯一，所以得到一个良定义的局部联络。把 \(s=\sum_a s^a e_a\)、\(t=\sum_b t^b h_b\) 代入，普通乘积法则 \(X(s^at^b)=X(s^a)t^b+s^aX(t^b)\) 正好给出命题所述纯张量公式。特别地，在 \((fs)\otimes t=s\otimes(ft)\) 的两种写法中，额外的 \(X(f)s\otimes t\) 项相同，不会产生不同结果。

	任取另一组标架，它的每个张量基元素仍是两个局部截面的纯张量，故刚构造的联络也满足该组基上的同一规定。再用对系数的 Leibniz 规则，得到两组构造在任意截面上相同，因而能够整体粘合。纯张量基上的规定也证明唯一性。多个张量因子时，重复同一构造得到“对每个位置求导后相加”的公式。改变括号或排列因子，只会重新排列这组求和项，所以与这些自然识别相容。

	收缩的相容性来自对偶公式。对纯张量中一对因子 \(\lambda,s\) 作收缩时，它们被替换成函数 \(\lambda(s)\)。对这一函数求导产生的两个项，正好是先对 \(\lambda\) 或 \(s\) 求导再收缩的结果；其余因子上的求导项保持不变。纯张量局部张成全部张量，故一般收缩也与诱导联络相容。

	接下来说明外幂公式确实尊重外幂的关系，而不仅是在可分解元素上写出一个表达。对 \(E^{\otimes k}\)，记 \(P_\sigma\) 为按置换 \(\sigma\) 重新排列因子的光滑丛映射，定义交替化算子
	\[
		\operatorname{Alt}=\frac1{k!}\sum_{\sigma\in S_k}\operatorname{sgn}(\sigma)P_\sigma.
	\]
	\(\operatorname{Alt}^2=\operatorname{Alt}\)：在两次求和中，每个合成置换恰出现 \(k!\) 次，符号相乘等于合成的符号。若纯张量有两个相同因子，按交换这两个位置把求和项配对，便知其交替化为零。因此 \(\operatorname{Alt}\) 消去\cref{g1:def:symmetric-exterior-powers}中定义外幂的关系子空间，诱导出 \(\Lambda^kE\to\operatorname{im}(\operatorname{Alt})\)。

	设 \(q:E^{\otimes k}\to\Lambda^kE\) 为自然商映射。交换输入改变楔积符号，所以 \(qP_\sigma=\operatorname{sgn}(\sigma)q\)，进而 \(q\operatorname{Alt}=q\)。因此刚才的诱导映射与 \(q\) 限制在 \(\operatorname{im}(\operatorname{Alt})\) 上的映射互逆：一个方向由 \(q\operatorname{Alt}=q\) 给出，另一方向由 \(\operatorname{Alt}^2=\operatorname{Alt}\) 给出。这些映射在张量标架中都是固定的有限线性组合，故为光滑识别。

	张量积联络与每个 \(P_\sigma\) 相容，因而与 \(\operatorname{Alt}\) 相容。若 \(w=\operatorname{Alt}w\)，则 \(\nabla_Xw=\operatorname{Alt}(\nabla_Xw)\)，所以它保持交替化像中的截面。通过上述识别，将此联络传到 \(\Lambda^kE\)，就得到外幂联络。对 \(q(s_1\otimes\cdots\otimes s_k)\) 使用张量求导公式，恰得到所述楔积公式。递增指标楔积构成局部标架，见\cref{g1:prop:symmetric-exterior-bases}，故该公式也保证唯一性。\(k>\operatorname{rank}E\) 时外幂为零丛，结论仍成立。\(k=0\) 按普通函数求导，\(k=1\) 则恢复原联络。去掉交替化中的符号得到对称化平均；它消去排列之差，且与张量联络相容，同一论证便给出对称幂上的联络。

	最后，\cref{g1:prop:tensor-hom-duality}给出 \(\operatorname{Hom}(E,F)\simeq F\otimes E^*\)，故已有一个张量诱导联络。对纯张量 \(\mathcal B=t\otimes\lambda\)，其作用为 \(\mathcal B s=\lambda(s)t\)，于是
	\begin{align*}
		\nabla^F_X(\mathcal B s)-\mathcal B(\nabla^E_Xs)
		&=X(\lambda(s))t+\lambda(s)\nabla^F_Xt-\lambda(\nabla^E_Xs)t\\
		&=(\nabla^{E^*}_X\lambda)(s)t+\lambda(s)\nabla^F_Xt\\
		&=\bigl(\nabla_X^{F\otimes E^*}(t\otimes\lambda)\bigr)(s).
	\end{align*}
	一般同态截面在每张图中是有限个这样的纯张量之和，所以上式对它也成立。求值于每个纤维向量能确定一个线性映射，故公式保证唯一性。以上所有构造只用了光滑系数、函数微分与有限次代数运算，边界处按局部延拓仍满足同样的公式。
\end{proof}

外幂联络在求导时保持外幂次数，因此
\[
	\nabla_X(\alpha\wedge\beta)=(\nabla_X\alpha)\wedge\beta+\alpha\wedge\nabla_X\beta
\]
没有外微分乘积法则中的 \((-1)^{\deg\alpha}\)。两种运算的类型不同：这里 \(\alpha,\beta\) 可以是 \(\Lambda^pE,\Lambda^qE\) 的截面，\(\nabla_X\) 在每个纤维的外代数上保持次数；第14章的 \(\mathrm{d}\) 作用于底流形的微分形式并使次数增加 \(1\)。即使取 \(E=\mathrm{T}^*M\)，也应保留这种区别。

\ProseParagraphBreak

诱导联络还给出了以后常用的求值公式。例如对光滑双线性形式场 \(h\in\Gamma(E^*\otimes E^*)\)，求导与收缩相容意味着
\begin{equation}
	\label{g1:eq:connection-bilinear-form}
	(\nabla_Xh)(s,t)=X\bigl(h(s,t)\bigr)-h(\nabla_Xs,t)-h(s,\nabla_Xt).
\end{equation}
这是先把 \(h,s,t\) 全部看作张量因子、对它们分别求导，再把 \(h\) 的导数移到等式一侧得到的。它同时记录了系数的变化与输入截面的变化，不能只保留第一项。

\ProseParagraphBreak

把同态联络写成矩阵，更能看清这种修正。在 \(E,F\) 的标架 \(e,h\) 中，设原联络矩阵分别为 \(A,C\)，同态截面 \(\mathcal B\) 由矩阵函数 \(B\) 表示，即 \(\mathcal B(eu)=hBu\)。对上面同态公式的右侧计算，有
\begin{align*}
	\nabla^F_X(\mathcal B(eu))&=h\bigl(X(B)u+B X(u)+C(X)Bu\bigr),\\
	\mathcal B(\nabla^E_X(eu))&=hB\bigl(X(u)+A(X)u\bigr).
\end{align*}
相减后含 \(X(u)\) 的项消去，说明结果确实只对输入向量的点值线性。同态值一形式 \(\nabla^{\operatorname{Hom}}\mathcal B\) 的矩阵因此为 \(\mathrm{d}B+CB-BA\)。当 \(F=E\)、两边使用同一个联络时，这成为
\begin{equation}
	\label{g1:eq:endomorphism-connection}
	(\nabla^{\operatorname{End}}\mathcal B)_e
	=\mathrm{d}B+[A,B],\qquad [A,B]=AB-BA.
\end{equation}
下标 \(e\) 表示取该自同态值一形式在标架 \(e\) 中的矩阵；\(B\) 是零次矩阵函数，故这里是通常的矩阵交换子，没有分次符号。

\ProseParagraphBreak

例如恒等自同态的矩阵为 \(I\)，所以 \(\nabla^{\operatorname{End}}\operatorname{id}_E=0\)。对两条自同态截面的复合，矩阵乘积法则给出
\[
	\mathrm{d}(BC)+A(BC)-(BC)A
	=(\mathrm{d}B+AB-BA)C+B(\mathrm{d}C+AC-CA).
\]
因此 \(\nabla_X^{\operatorname{End}}(\mathcal B\mathcal C)=(\nabla_X^{\operatorname{End}}\mathcal B)\mathcal C+\mathcal B\nabla_X^{\operatorname{End}}\mathcal C\)。对偶、张量和自同态上的这些规则使一个原始联络能够同时处理向量、线性观测与几何张量，而无需为每一种对象重新选择求导方式。
% 纲要条目排为编号小节标题，并进入目录与书签；说明段仍保留于纲要注释中。

% !TeX root = ../../Book1.tex
% 纲要标题：### 15.2 沿路径的平行移动
% 纲要位置：几何学计划纲要.md，第 701 行。
% 正文状态：初稿；尚未正式审核。

% 纲要内容开始
% 
% * 协变常微分方程及参数依赖
% * 路径拼接、反向和平行移动
% * 由平行移动恢复联络
% * 度量联络和长度保持
% 
% 纲要内容结束

\section{沿路径的平行移动}
\label{g1:sec:1502}

联络规定了截面怎样沿一个瞬时方向变化，但两个不同点处的向量仍属于不同纤维，不能直接相减。如果沿一条路径要求协变导数始终为零，就可以由初始向量确定沿途的“不变”向量。这将给出两个端点纤维之间的线性同构。比较一般截面时，可以先把终点向量运回起点再作差；这个差商又会恢复最初的联络。

\ProseParagraphBreak

本节的底流形 \(M\) 为有限维、Hausdorff、第二可数的光滑流形，允许有边界；\(E\to M\) 是实秩 \(r\) 光滑向量丛，已给定\cref{g1:def:bundle-connection}中的联络 \(\nabla\)。局部标架写成行 \(e=(e_1,\ldots,e_r)\)，截面分量写成列，约定 \(\nabla e=eA\)。存在性所需的分析输入是\cref{g1:thm:ode-local-existence-uniqueness,g1:thm:ode-smooth-dependence}：光滑常微分方程有唯一局部解，且在共同存在区间上光滑依赖初值和有限维参数。我们还要补明，当前的线性方程为何能走完整个给定的紧时间区间。

\subsection{协变常微分方程及参数依赖}
\label{g1:subsec:1502:01}

设 \(\gamma:I\to M\) 是光滑曲线。沿 \(\gamma\) 的一个向量 \(s(t)\) 位于 \(E_{\gamma(t)}\)，所以应当把它看成拉回丛 \(\gamma^*E\to I\) 的截面，参见\cref{g1:prop:pullback-bundle}。具体说，\(s:I\to E\) 光滑且 \(\pi\circ s=\gamma\)。即使 \(\gamma(t_1)=\gamma(t_2)\)，也不要求 \(s(t_1)=s(t_2)\)。因此不能要求每个沿曲线的截面都由底流形上的同一个截面限制而来；自交与常值路径已经会使这种要求失败。

\begin{definition}[沿曲线的协变导数]\label{g1:def:along-curve-covariant-derivative}
	设 \(E\to M\) 为实有限秩光滑向量丛，\(\nabla\) 为联络，\(I\subseteq\mathbb R\) 为区间，\(\gamma:I\to M\) 为光滑曲线，\(s\) 为拉回丛 \(\gamma^*E\) 的光滑截面，\(M\) 可带边界。在曲线落于局部标架 \(e\) 的定义域内的参数子区间上，令 \(\nabla e=eA\)，并写 \(s(t)=e(\gamma(t))u(t)\)，定义
	\begin{equation}\label{g1:eq:along-curve-covariant-derivative}
		\mathrm{D}_t s=e(\gamma(t))\left(\frac{\mathrm{d}u}{\mathrm{d}t}+A_{\gamma(t)}(\gamma'(t))u(t)\right).
	\end{equation}
	这称为 \(s\) 沿 \(\gamma\) 的协变导数。满足 \(\mathrm{D}_t s=0\) 的截面称为沿 \(\gamma\) 的\term[pingxingjiemian]{平行截面}[parallel section]。
\end{definition}

先检验这个定义不依赖标架。若 \(\widetilde e=e g\)、\(u=g(\gamma(t))\widetilde u\)，由\eqref{g1:eq:connection-frame-change}，
\[
	\widetilde A=g^{-1}Ag+g^{-1}\mathrm{d}g.
\]
记 \(B(t)=A_{\gamma(t)}(\gamma'(t))\)、\(g(t)=g(\gamma(t))\)，则 \(\widetilde B=g^{-1}Bg+g^{-1}g'\)，从而
\[
	u'+Bu=g'\widetilde u+g\widetilde u'+Bg\widetilde u
	=g(\widetilde u'+\widetilde B\widetilde u).
\]
乘上 \(e(\gamma(t))\) 后，两套标架给出同一个纤维向量。因此局部公式可以粘合。它也直接给出实线性及乘积法则 \(\mathrm{D}_t(fs)=f's+f\mathrm{D}_ts\)。若 \(S\) 是底流形上的局部截面，则链式法则与\eqref{g1:eq:connection-local-form}说明
\[
	\mathrm{D}_t(S\circ\gamma)=(\nabla_{\gamma'(t)}S)_{\gamma(t)}.
\]
这个相容关系适用于确实来自局部截面的情形，但不是沿曲线导数的定义前提。

\ProseParagraphBreak

现在平行条件成为线性初值问题
\begin{equation}\label{g1:eq:parallel-transport-ode}
	\frac{\mathrm{d}u}{\mathrm{d}t}=-B(t)u(t),\qquad
	B(t)=A_{\gamma(t)}(\gamma'(t)),\qquad u(t_0)=u_0.
\end{equation}
负号来自把联络项移到等式右侧。一般情况下，\(B(t_1)\) 与 \(B(t_2)\) 不交换，因而不能未经验证就把解写成 \(\exp(-\int B)\)。存在唯一性本身并不需要这种显式表达。

\begin{lemma}[线性方程的全区间解]\label{g1:lem:linear-transport-ode}
	设 \(a<b\)，\(B:[a,b]\to\operatorname{Mat}_{r\times r}(\mathbb R)\) 光滑，包括端点处的局部光滑延拓。对任意 \(t_0\in[a,b]\) 与 \(u_0\in\mathbb R^r\)，方程 \(u'=-Bu\)、\(u(t_0)=u_0\) 在整个 \([a,b]\) 上有唯一光滑解。若 \(\|B(t)\|\le L\)，则
	\begin{equation}\label{g1:eq:linear-transport-bound}
		|u(t)|\le e^{L|t-t_0|}|u_0|.
	\end{equation}
	解线性依赖 \(u_0\)。记 \(U(t,s)\) 为从时刻 \(s\) 的分量到时刻 \(t\) 的分量的解算子，则 \(U(t,s)\) 可逆，其逆为 \(U(s,t)\)。若 \(B=B(t,\lambda)\) 光滑依赖有限维参数，解在每个参数点附近对 \(t,t_0,u_0,\lambda\) 联合光滑。
\end{lemma}
\begin{proof}
	局部存在唯一性适用于 \(F(t,u)=-B(t)u\)，因为其右端在时间的局部延拓区间与整个 \(\mathbb R^r\) 上光滑。积分方程给出
	\[
		|u(t)|\le |u_0|+L\int_{t_0}^t|u(v)|\,\mathrm{d}v\qquad(t\ge t_0).
	\]
	\cref{g1:lem:gronwall-estimate}给出\eqref{g1:eq:linear-transport-bound}；向过去求解时先把时间反向，得到相同估计。

	为排除有限时间中断，假设解向前只能存在到 \(c\le b\)。上述界使解在 \([t_0,c)\) 上有共同界 \(C\)，故 \(|u'|\le LC\)。于是 \(|u(t)-u(v)|\le LC|t-v|\)，当 \(t\uparrow c\) 时存在有限极限 \(u_c\)。在 \((c,u_c)\) 处重新使用局部初值定理。旧解的积分方程在 \(c\) 处取极限后仍成立，接上新解便得到穿过 \(c\) 的解，与不能延拓矛盾。若 \(c=b\)，同样得到端点处及其局部延拓中的光滑解。向后完全相同，故解存在于整个紧区间。这里起作用的是线性增长估计，而不是一般非线性方程都能在有限区间内求解。

	同一方程的两解作任意线性组合，仍满足方程，初值也作相同组合；唯一性因此证明初值依赖线性。从时刻 \(s\) 到 \(t\) 后，再以所得值从 \(t\) 解回 \(s\)，唯一性使返回值等于原初值。交换两个时刻也成立，故 \(U(s,t)\) 与 \(U(t,s)\) 互逆。

	最后考虑参数。对每个固定参数，前面的证明已经给出完整时间区间上的解。系数在 \([a,b]\) 与一个较小参数邻域上有共同局部光滑延拓；必要时沿紧时间区间取有限覆盖并用时间变量的光滑权重拼接这些矩阵系数的延拓。\cref{g1:thm:ode-smooth-dependence}适用于光滑右端 \(-B(t,\lambda)u\)：该定理关于已有紧解区间的结论，给出附近初值、初始时刻及参数的共同解，并保证全部变量的联合光滑性。参数若落在半空间中，先局部延拓矩阵系数，再把所得光滑解限制回来，结论也成立。
\end{proof}

\begin{theorem}[平行移动]\label{g1:thm:parallel-transport}
	设 \(M\) 为有限维、Hausdorff、第二可数的光滑流形，允许有边界，\(E\to M\) 为实有限秩光滑向量丛，带联络 \(\nabla\)。设 \(\gamma:[a,b]\to M\) 连续且分片光滑，即存在有限分割，使每个闭子区间上的限制光滑。对任意 \(t_0\in[a,b]\) 和 \(v\in E_{\gamma(t_0)}\)，存在唯一连续、分片光滑的沿曲线截面 \(s_v\)，在每个光滑片上平行且 \(s_v(t_0)=v\)。映射
	\[
		P_{\gamma,t\leftarrow t_0}:E_{\gamma(t_0)}\longrightarrow E_{\gamma(t)},\qquad
		v\longmapsto s_v(t)
	\]
	为线性同构，称为沿该路径段的\term[pingxingyidong]{平行移动}[parallel transport]。记 \(P_\gamma=P_{\gamma,b\leftarrow a}\)。

	若 \(\gamma=\gamma_\lambda\) 是一个有限维光滑参数族，或者各片具有共同有限分割且对时间与参数联合光滑，则在参数点附近，端点局部标架中的 \(P_{\gamma_\lambda}\) 矩阵光滑依赖参数；光滑片内的到达与起始时刻也具有联合光滑依赖。
\end{theorem}
\begin{proof}
	先设路径光滑。其紧像被有限个标架域覆盖。在时间区间的这个开覆盖中取足够细的有限分割，使每个闭子区间的像都落在某个标架域内；这可由紧区间的 Lebesgue 数性质得到。对包含初始时刻的区间使用\cref{g1:lem:linear-transport-ode}，再把端点向量作为下一段的初值，依次向两个方向求解。换标架时，前面已验证的变换公式使同一个沿曲线截面满足相同的平行条件。局部唯一性因而保证相邻解在重叠处一致，所得截面光滑。

	若另有一个同初值的平行截面，两者先在初始段相同，再由共同端点初值逐段相同。因此整体唯一。这也证明选择不同标架域或细分区间不改变结果：两种构造都是同一初值问题的解。对分片光滑路径，先把给定的有限分割加入上述分割，然后逐片连接；拐角只要求向量连续，不要求底路径在拐角有一个共同速度。唯一性与选择无关性仍由逐段论证得到。线性及可逆性来自每段的线性解算子与反向求解，有限次复合保留这些性质。

	对参数族，在一个固定参数 \(\lambda_0\) 处选择上述有限分割，并使每一段的像紧含在所需的标架域中。连续性与时间区间的紧性保证：缩小 \(\lambda_0\) 的邻域后，同一批分割和标架仍适用于附近全部路径。在每段，矩阵 \(A_{\gamma_\lambda(t)}(\partial_t\gamma_\lambda(t))\) 对时间和参数光滑，故线性方程引理给出光滑解矩阵。段与段之间的标架转换矩阵，也是在光滑变化的路径端点求值的光滑矩阵。端点运输矩阵是这些矩阵的有限次乘积，因此光滑。片内可变时刻的结论使用同一引理关于初始和到达时刻的光滑依赖。
\end{proof}

底流形有边界并不破坏上述论证。我们已经给定一条留在 \(M\) 中的路径，只解纤维分量所满足的 Euclid 线性方程；不需要让底流形上的某个向量场在边界处产生双侧流。时间端点处只延拓局部坐标中的矩阵系数，也没有要求把底路径延伸到 \(M\) 之外。

\begin{example}\label{g1:exm:line-connection-transport}
	在平凡实线丛 \(M\times\mathbb R\) 的标准标架中，联络矩阵就是一个实一形式 \(\alpha\)。方程 \(u'=-\alpha(\gamma')u\) 的解为
	\[
		u(t)=\exp\left(-\int_{t_0}^t \alpha_{\gamma(v)}(\gamma'(v))\,\mathrm{d}v\right)u_0.
	\]
	直接求导并检验初值，唯一性就确认了这个公式。特别地，若 \(\alpha=\mathrm{d}f\)，则 \(P_\gamma\) 在标准纤维坐标中是乘以 \(e^{f(\gamma(a))-f(\gamma(b))}\)，只依赖两个端点。这里是秩一标量方程，所以积分后的指数表达确实成立。
\end{example}

\subsection{路径拼接、反向和平行移动}
\label{g1:subsec:1502:02}

一个端点纤维的向量，可以先沿一条路径运到中间点，再沿下一条路径继续。为避免书写次序与实际运动相反，沿用前文道路约定：\(\gamma*\eta\) 表示先走 \(\gamma\)，后走 \(\eta\)。因线性映射的复合从右向左作用，最终公式的次序应当是 \(P_\eta P_\gamma\)。

\begin{proposition}[平行移动的运算规律]\label{g1:prop:parallel-transport-laws}
	设 \(E\to M\) 为带联络的实有限秩光滑向量丛，\(M\) 可带边界。对连续且有限分片光滑的路径 \(\gamma,\eta:[0,1]\to M\)，在 \(\gamma(1)=\eta(0)\) 时，以 \(\gamma*\eta\) 表示先走 \(\gamma\) 后走 \(\eta\) 的拼接。平行移动满足：
	\begin{enumerate}
		\item 常值路径的平行移动为该纤维的恒等映射。同一路径上有
		\[
			P_{\gamma,t\leftarrow s}P_{\gamma,s\leftarrow r}=P_{\gamma,t\leftarrow r}.
		\]
		\item 若 \(\gamma(1)=\eta(0)\)，则
		\[
			P_{\gamma*\eta}=P_\eta P_\gamma.
		\]
		\item 对反向路径 \(\overline\gamma(t)=\gamma(1-t)\)，有 \(P_{\overline\gamma}=P_\gamma^{-1}\)。
		\item 另取连续且有限分片光滑的路径 \(\gamma:[a,b]\to M\)。若 \(h:[c,d]\to[a,b]\) 是光滑非减映射，\(h(c)=a\)、\(h(d)=b\)，则 \(P_{\gamma\circ h}=P_\gamma\)。允许 \(h'\) 为零，特别允许路径在中间停顿。
	\end{enumerate}
\end{proposition}
\begin{proof}
	常值路径的速度为零，局部方程变成 \(u'=0\)，所以纤维向量不变。第一项中的复合公式则由唯一性得到：从时刻 \(r\) 的初值先解到 \(s\)，再以该值继续，所得截面仍是同一条初值解；任意时刻次序都可以用全区间的唯一解作这一比较。

	对重参数化，若 \(s(t)=e(\gamma(t))u(t)\) 沿 \(\gamma\) 平行，链式法则给出
	\[
		\mathrm{D}_\tau(s\circ h)
		=e(\gamma(h(\tau)))\,h'(\tau)\bigl(u'(h(\tau))+B(h(\tau))u(h(\tau))\bigr)=0.
	\]
	因此 \(s\circ h\) 就是重参数路径上具有相同起点向量的平行截面。这个计算没有除以 \(h'\)，所以在停顿区间仍有效。分片路径可在各分段点的原像端点处细分，再逐片应用同一等式。

	把 \(\gamma,\eta\) 各自压缩到 \([0,\tfrac12]\)、\([\tfrac12,1]\)，前段终点向量为 \(P_\gamma v\)，后段终点为 \(P_\eta(P_\gamma v)\)。上述重参数化规律和连接初值的唯一性证明拼接公式。反向参数化也满足同一个链式法则，只是导数为负。因而沿反向路径运输一条原平行截面的终点值，恰好回到其起点值，给出 \(P_{\overline\gamma}=P_\gamma^{-1}\)。
\end{proof}

在两个端点更换标架时，运输矩阵也有清楚的变换律。若原矩阵为 \(U\)，新标架为 \(\widetilde e=e g\)，则起点旧分量为 \(g(\gamma(a))\widetilde u_a\)，终点再换成新分量，故
\begin{equation}\label{g1:eq:transport-frame-change}
	\widetilde U=g(\gamma(b))^{-1}U\,g(\gamma(a)).
\end{equation}
矩阵变了，所表示的纤维线性同构没有改变。取初始纤维的一组基并逐个平行移动，还得到 \(\gamma^*E\) 在整个区间上的平行标架。这不意味着它能成为底空间上一个单值标架：若路径闭合，末端基可能不等于初始基。

\begin{example}\label{g1:exm:path-dependent-line-transport}
	在 \(\mathbb R^2\times\mathbb R\) 上取 \(\alpha=x\,\mathrm{d}y\)。从 \(p=(0,0)\) 到 \(q=(1,1)\)，路径 \(\gamma\) 先沿水平边到 \((1,0)\)，再沿竖直边到 \(q\)；路径 \(\eta\) 先到 \((0,1)\)，再到 \(q\)。沿 \(\gamma\)，第一段 \(\alpha(\gamma')=0\)，第二段 \(\int\alpha=1\)；沿 \(\eta\)，两段的积分都为零。因此
	\[
		P_\gamma(v)=e^{-1}v,\qquad P_\eta(v)=v.
	\]
	同样的起点、终点给出不同运输，尽管这两条路径位于可收缩的平面内。可见仅有底空间的拓扑条件还不够；联络自身也必须满足额外条件，才能获得路径无关性。
	\input{Figures/Chapter15-PathTransport.tex}
\end{example}

\subsection{由平行移动恢复联络}
\label{g1:subsec:1502:03}

给定局部截面 \(S\)，式 \(S(\gamma(t))-S(\gamma(0))\) 一般没有意义，因为两个向量不在同一纤维。先用反向运输 \(P_{\gamma,0\leftarrow t}\) 把第一个向量送回起点，就能在固定向量空间中取差商。这个操作说明平行移动保存了联络的全部一阶信息。

\begin{theorem}[由运输恢复联络]\label{g1:thm:connection-from-transport}
	设 \(E\to M\) 为实有限秩光滑向量丛，\(M\) 可带边界，\(\nabla\) 为给定联络，\(P\) 是由该联络构造的平行移动。对光滑曲线 \(\gamma:I\to M\)、任意沿曲线的光滑截面 \(s\) 与固定 \(t_0\in I\)，在固定纤维 \(E_{\gamma(t_0)}\) 中有
	\begin{equation}\label{g1:eq:transported-section-derivative}
		\frac{\mathrm{d}}{\mathrm{d}t}\bigl(P_{\gamma,t_0\leftarrow t}s(t)\bigr)
		=P_{\gamma,t_0\leftarrow t}\mathrm{D}_t s.
	\end{equation}
	特别地，若 \(\gamma(0)=p\)、\(\gamma'(0)=v\)，\(S\) 是 \(p\) 附近的光滑截面，则
	\begin{equation}\label{g1:eq:connection-recovery-limit}
		(\nabla_vS)_p
		=\lim_{t\to0}\frac{P_{\gamma,0\leftarrow t}S(\gamma(t))-S(p)}{t}.
	\end{equation}
	若曲线只在零点的一侧有定义，则使用相应的单侧极限。因此同一向量丛上，若两个光滑联络沿所有分片光滑路径产生相同的平行移动，则两个联络相同。
\end{theorem}
\begin{proof}
	先在包含所考察时刻的一个标架域内计算。设从该局部初始时刻到 \(t\) 的平行移动矩阵为 \(U(t)\)，则 \(U'=-BU\)，且 \(U\) 可逆。对 \(U^{-1}U=I\) 求导，得到
	\[
		(U^{-1})'=-U^{-1}U'U^{-1}=U^{-1}B.
	\]
	若 \(s(t)=e(\gamma(t))u(t)\)，则把它运回局部初始纤维后的分量为 \(U(t)^{-1}u(t)\)，而
	\[
		\frac{\mathrm{d}}{\mathrm{d}t}(U^{-1}u)=U^{-1}(u'+Bu).
	\]
	右端恰是 \(\mathrm{D}_ts\) 反向运输后的分量。若固定的 \(t_0\) 不在这一标架域内，先运回局部初始时刻，再接上从该时刻到 \(t_0\) 的固定线性映射；固定线性映射与求导相容，故\eqref{g1:eq:transported-section-derivative}仍成立。

	现在取 \(s=S\circ\gamma\)、\(t_0=0\)，在零点求导。由 \(P_{\gamma,0\leftarrow0}=\operatorname{id}\) 及沿曲线导数与 \(\nabla_{\gamma'}S\) 的相容关系，得到\eqref{g1:eq:connection-recovery-limit}。同一局部计算也证明单侧极限的版本。

	在内部点，每个 \(v\in \mathrm{T}_pM\) 都能由坐标直线给出一条双侧小曲线的初速度，故相同的运输决定所有 \(\nabla_vS\)。在边界点取 \(x^n\ge0\)、\(x(p)=0\) 的图。若 \(v^n\ge0\)，坐标曲线 \(t\mapsto tv\)、\(t\ge0\)，留在半空间中，单侧公式决定 \(\nabla_vS\)。若 \(v^n<0\)，则 \(-v\) 是向内方向，用已经得到的值及 \(\nabla_vS=-\nabla_{-v}S\) 决定它。故两个给定联络若有相同运输，在边界点也相同。零维底流形没有非零切向方向，唯一性同样成立。
\end{proof}

边界处需要保留上述区别。一条以边界点为内部参数时刻、始终留在半空间中的双侧曲线，其法向坐标在该时刻取得局部最小值，所以法向速度必为零。不能用这种曲线代表所有切向量；向内方向由单侧曲线实现，向外方向再通过联络关于方向的线性性决定。

\ProseParagraphBreak

定理恢复的是一个已经给定的联络所产生的运输。若只向每条路径任意指定一个线性同构，即使其中部分映射满足拼接规律，也还不能直接把差商宣布为联络：差商的存在、对路径一阶信息的依赖及所需光滑性都要另外验证。本节得到的是联络与其运输之间的确定关系和唯一性。

\subsection{度量联络和长度保持}
\label{g1:subsec:1502:04}

丛度量 \(g\) 为每个纤维提供内积，存在性已由\cref{g1:thm:bundle-metric-existence}证明。但任意联络不一定在运输时保长；前面 \(x\,\mathrm{d}y\) 的例子，在标准纤维内积下就可能缩短向量。要求“沿途不变的向量保持内积”，应当对应一条关于 \(g\) 与 \(\nabla\) 的微分条件。

\begin{definition}[度量联络]\label{g1:def:metric-connection}
	设 \(E\to M\) 为实有限秩光滑向量丛，\(g\) 是光滑正定丛度量。若联络 \(\nabla\) 对所有局部向量场 \(X\) 和局部截面 \(S,T\) 满足
	\[
		X\bigl(g(S,T)\bigr)=g(\nabla_XS,T)+g(S,\nabla_XT),
	\]
	则称其为关于 \(g\) 的\term[dulianglianluo]{度量联络}[metric connection]，也称与 \(g\) 相容。这就是诱导张量联络下的 \(\nabla g=0\)。
\end{definition}

\begin{proposition}[度量相容与平行移动]\label{g1:prop:metric-parallel-transport}
	设 \(M\) 为有限维光滑流形，允许有边界，\(E\to M\) 为实有限秩光滑向量丛，带光滑正定度量 \(g\) 和联络 \(\nabla\)。以下条件等价：
	\begin{enumerate}
		\item \(\nabla g=0\)。
		\item 在每个局部标架 \(e\) 中，令 \(G=(g(e_i,e_j))\)、\(\nabla e=eA\)，则
		\begin{equation}\label{g1:eq:metric-connection-matrix}
			\mathrm{d}G=A^{\mathsf T}G+GA.
		\end{equation}
		\item 对每条光滑曲线及任意两个沿曲线的光滑截面 \(s,t\)，有
		\[
			\frac{\mathrm{d}}{\mathrm{d}\tau}g(s(\tau),t(\tau))
			=g(\mathrm{D}_\tau s,t)+g(s,\mathrm{D}_\tau t).
		\]
		\item 每条有限分片光滑路径的平行移动都是端点纤维之间的内积等距同构。
	\end{enumerate}
	因而也可以把最后一项等价地写为所有平行移动保持每个向量的长度。
\end{proposition}
\begin{proof}
	将度量相容条件应用于局部基截面 \(e_i,e_j\)。\(\nabla e_j=\sum_k e_kA^k{}_j\)，所以右边的两个矩阵分别为 \(A^{\mathsf T}G\) 和 \(GA\)，得到第二项。反过来，把任意截面写成 \(eu,ev\)，对 \(u^{\mathsf T}Gv\) 使用乘积法则，第二项就给出第一项。

	沿曲线记 \(B=A(\gamma')\)。第二项拉回曲线成为 \(G'=B^{\mathsf T}G+GB\)。于是
	\[
		\frac{\mathrm{d}}{\mathrm{d}\tau}(u^{\mathsf T}Gv)
		=(u'+Bu)^{\mathsf T}Gv+u^{\mathsf T}G(v'+Bv),
	\]
	恰为第三项。两个截面都平行时，右端为零，故内积沿每个光滑片恒定；连续性使这些常值在分段端点吻合。这证明第四项。

	反过来假设第四项。对 \(p\) 附近的局部截面 \(S,T\) 及过 \(p\) 的曲线，令
	\[
		\widehat S(\tau)=P_{\gamma,0\leftarrow\tau}S(\gamma(\tau)),\qquad
		\widehat T(\tau)=P_{\gamma,0\leftarrow\tau}T(\gamma(\tau)).
	\]
	运输保内积，故 \(g_{\gamma(\tau)}(S,T)=g_p(\widehat S(\tau),\widehat T(\tau))\)。在零点求导，再用\cref{g1:thm:connection-from-transport}，得到方向 \(v=\gamma'(0)\) 上的度量相容等式。内部点可取任意 \(v\)；边界处先取向内的单侧曲线，再利用等式两边关于 \(v\) 的线性性处理向外方向。因此第一项成立。

	最后，保内积必然保长度；反过来，实内积由
	\[
		g(v,w)=\frac12\bigl(|v+w|_g^2-|v|_g^2-|w|_g^2\bigr)
	\]
	恢复。平行移动已经线性，故保全部向量的长度也保证保内积。
\end{proof}

在正交归一标架中 \(G=I_r\)，矩阵条件简化为 \(A^{\mathsf T}+A=0\)，即每个 \(A(X)\) 为斜对称矩阵。在一般标架中却不能直接要求 \(A\) 斜对称，必须保留 \(\mathrm{d}G\) 这一项。度量可能随基点变化，而非正交标架自身的长度也会变化。

\begin{proposition}[度量联络的存在]\label{g1:prop:metric-connection-existence}
	设 \(M\) 为有限维、Hausdorff、第二可数的光滑流形，允许有边界，\(E\to M\) 为实有限秩光滑向量丛，\(g\) 为光滑正定丛度量。则 \(E\) 上存在与 \(g\) 相容的联络。给定一个度量联络后，所有其他度量联络恰为 \(\nabla+K\)，其中 \(K\in\Omega^1(M;\operatorname{End}E)\) 满足每个 \(K(X)\) 关于 \(g\) 斜伴随。
\end{proposition}
\begin{proof}
	由\cref{g1:thm:connection-existence}取任意联络 \(\nabla^0\)，再修正它不保持度量的部分。定义
	\[
		Q_X(S,T)=Xg(S,T)-g(\nabla^0_XS,T)-g(S,\nabla^0_XT).
	\]
	Leibniz 法则使这里对 \(S,T\) 的导数项在乘以函数时抵消，故 \(Q\) 对 \(X,S,T\) 都是逐点线性的；它还是光滑的，且对 \(S,T\) 对称。正定性使每个纤维的 \(g\) 非退化，因此存在唯一的 \(C_X\in\operatorname{End}(E)\)，满足
	\[
		g(C_XS,T)=\frac12Q_X(S,T).
	\]
	在局部标架中，若 \(Q_X\) 的对称矩阵为 \(Q(X)\)，这个算子为 \(C_X=\tfrac12G^{-1}Q(X)\)。矩阵逆光滑，故 \(C\) 是光滑的 \(\operatorname{End}E\) 值一形式。

	由\cref{g1:prop:connection-affine-space}，\(\nabla=\nabla^0+C\) 仍是联络。对称性给出
	\begin{align*}
		(\nabla_Xg)(S,T)&=Q_X(S,T)-g(C_XS,T)-g(S,C_XT)\\
		&=Q_X(S,T)-\tfrac12Q_X(S,T)-\tfrac12Q_X(S,T)=0.
	\end{align*}
	因此修正后的联络与度量相容。在矩阵语言中，这一修正是
	\[
		A=A^0+\frac12G^{-1}\bigl(\mathrm{d}G-(A^0)^{\mathsf T}G-GA^0\bigr).
	\]

	若 \(\widetilde\nabla=\nabla+K\)，其中 \(\nabla\) 已度量相容，则
	\[
		(\widetilde\nabla_Xg)(S,T)=-g(K(X)S,T)-g(S,K(X)T).
	\]
	因而 \(\widetilde\nabla\) 也相容，当且仅当每个 \(K(X)\) 斜伴随。这证明了最后的全部刻画。
\end{proof}

\begin{example}\label{g1:exm:rotating-metric-transport}
	在平凡二维丛中取标准内积，令
	\[
		J=\begin{pmatrix}0&-1\\1&0\end{pmatrix},\qquad A=J\,\alpha
	\]
	其中 \(\alpha\) 是任意实一形式。因为 \(J^{\mathsf T}=-J\)，这是度量联络。沿路径记 \(\theta(t)=\int_{t_0}^t \alpha(\gamma'(v))\,\mathrm{d}v\)，则
	\[
		u(t)=\exp(-\theta(t)J)u_0
		=\begin{pmatrix}\cos\theta(t)&\sin\theta(t)\\-\sin\theta(t)&\cos\theta(t)\end{pmatrix}u_0.
	\]
	矩阵求导给出 \(u'=-J\,\alpha(\gamma')u\)，因此这就是平行解。它保持长度，转角却是 \(-\theta\)，再次显示平行方程中的负号。保长不表示每条环路运输都为恒等；它只把可能的运输限制在纤维的正交变换之中。
\end{example}

% 纲要条目排为编号小节标题，并进入目录与书签；说明段仍保留于纲要注释中。

% !TeX root = ../../Book1.tex
% 纲要标题：### 15.3 曲率与结构方程
% 纲要位置：几何学计划纲要.md，第 708 行。
% 正文状态：初稿；尚未正式审核。

% 纲要内容开始
% 
% * 二阶协变微分及张量性
% * 外协变微分与曲率矩阵
% * 曲率的换标架规律与小环路运输
% * Bianchi 恒等式的结构证明
% 
% 纲要内容结束

\section{曲率与结构方程}
\label{g1:sec:1503}

联络使我们能够沿路径比较纤维中的向量，但两条同起点、同终点的路径未必给出同一种比较。在第九章的曲面例子中，绕小圈后的转角由曲率控制。一般向量丛中的返回变换可以伸缩、剪切，也可以旋转，因而需要一个取值于纤维自同态的对象来记录这种变化。本节把它从二阶协变微分中构造出来，再说明它怎样出现在小环路的运输中。

\ProseParagraphBreak

以下 \(E\to M\) 是有限秩实光滑向量丛，\(M\) 有限维、Hausdorff、第二可数，允许边界，\(\nabla\) 是光滑联络。沿用行标架 \(e\)、列分量 \(u\) 和 \(\nabla s=e(\mathrm{d}u+Au)\) 的约定。外微分及分次乘积法则使用\cref{g1:thm:exterior-derivative}；向量场括号的顺序仍为 \([X,Y]=XY-YX\)。边界处使用局部光滑延拓定义括号；\eqref{g1:eq:lie-bracket-product}中的乘法法则由同一坐标计算继续成立。

\subsection{二阶协变微分及张量性}
\label{g1:subsec:1503:01}

单独的 \(\nabla_X\nabla_Ys\) 会依赖 \(Y\) 及 \(s\) 在邻近点的变化。交换 \(X,Y\) 可以消去一部分二阶项，但若方向场自身不交换，还须减去它们的 Lie 括号所贡献的一阶变化。这正是第八章曲面曲率定义中出现第三项的原因，现在同一构造适用于任意向量丛。

\begin{definition}[丛联络的曲率]\label{g1:def:bundle-curvature}
	设 \(E\to M\) 为实光滑向量丛，\(\nabla\) 为其联络，允许 \(M\) 带边界。定义
	\[
		R^\nabla(X,Y)s=\nabla_X\nabla_Ys-\nabla_Y\nabla_Xs-\nabla_{[X,Y]}s,
		\qquad X,Y\in\mathfrak X(M),\quad s\in\Gamma(E).
	\]
	称 \(R^\nabla\) 为该联络的曲率；不致混淆时记作 \(R\)。
\end{definition}

\begin{proposition}[曲率的张量性]\label{g1:prop:bundle-curvature-tensorial}
	设 \(E\to M\) 为实光滑向量丛，\(M\) 可带边界，\(\nabla\) 为光滑联络。对光滑向量场 \(X,Y\) 和光滑截面 \(s\)，令
	\[
		R(X,Y)s=\nabla_X\nabla_Ys-\nabla_Y\nabla_Xs-\nabla_{[X,Y]}s.
	\]
	此运算对 \(X,Y,s\) 均为 \(C^\infty(M)\) 线性，并对前两个变量反对称。因此它唯一确定一个光滑截面
	\[
		R\in\Omega^2(M;\operatorname{End}E)
		=\Gamma(\Lambda^2\mathrm{T}^*M\otimes\operatorname{End}E).
	\]
	在点 \(p\)，\(R_p(v,w)\) 是 \(E_p\) 上的线性自同态，只依赖 \(v,w\in \mathrm{T}_pM\) 和被作用的纤维向量。
\end{proposition}
\begin{proof}
	实多线性和反对称性由定义得到。对光滑函数 \(f\)，先回忆\eqref{g1:eq:lie-bracket-product}的 \([fX,Y]=f[X,Y]-Y(f)X\)。该等式在边界处仍由局部延拓成立。于是
	\begin{align*}
		R(fX,Y)s
		&=f\nabla_X\nabla_Ys-\nabla_Y(f\nabla_Xs)-\nabla_{f[X,Y]-Y(f)X}s\\
		&=fR(X,Y)s.
	\end{align*}
	两个含 \(Y(f)\nabla_Xs\) 的项相消；第二变量的函数线性再由反对称性得到。

	对截面 \(fs\)，展开
	\[
		\nabla_X\nabla_Y(fs)
		=X(Yf)s+(Yf)\nabla_Xs+(Xf)\nabla_Ys+f\nabla_X\nabla_Ys.
	\]
	交换 \(X,Y\) 后，一阶交叉项抵消。再减去 \(\nabla_{[X,Y]}(fs)\)，额外剩余项为 \((X(Yf)-Y(Xf)-[X,Y]f)s=0\)，故 \(R(X,Y)(fs)=fR(X,Y)s\)。在局部坐标标架和丛标架中展开三个输入，函数线性把所有分量移到外面，点处结果便只依赖这些分量在该点的值。各基向量上的结果光滑，故得到所述光滑丛值二形式。
\end{proof}

这里的“二形式”只计算两个切向输入；它的输出不是实数，而是一个还可以作用于纤维向量的线性映射。若 \(E=\mathrm{T}M\)，这就是\cref{g1:def:surface-curvature-tensor}所用曲率的高维版本；若 \(E\) 是一条线丛，则每个纤维自同态只是一个实数倍。

\begin{definition}\label{g1:def:flat-connection}
	若光滑向量丛联络的曲率处处为零，即 \(R^\nabla=0\)，称该联络为\term[pingtanlianluo]{平坦联络}[flat connection]。这个条件针对给定联络；说一个丛具有平坦联络，还须指定或证明这样的联络存在。
\end{definition}

平坦的定义目前只表达一个局部微分条件。它是否足以得到平行标架，以及是否使所有同端点路径的运输相同，需要\secref{g1:sec:1505}的构造来回答。

\subsection{外协变微分与曲率矩阵}
\label{g1:subsec:1503:02}

普通外微分作用于实值形式，联络作用于丛截面。要把它们组织成同一个运算，回顾\cref{g1:def:bundle-valued-form}的记号
\[
	\Omega^k(M;E)=\Gamma(\Lambda^k\mathrm{T}^*M\otimes E).
\]
其元素在一点接受 \(k\) 个切向量，交替多线性地输出一个纤维向量。在标架 \(e\) 中写为 \(\eta=\sum_a\eta^a\otimes e_a\)，其中 \(\eta^a\) 是普通 \(k\) 形式；简记成 \(e\eta_e\)，\(\eta_e=(\eta^a)\) 为列。换标架 \(e'=eg\) 时，列分量变为 \(\eta_{e'}=g^{-1}\eta_e\)。

\ProseParagraphBreak

矩阵值形式相乘时，同时作矩阵乘法和形式楔积。若 \(B\) 是矩阵值 \(p\) 形式、\(C\) 是矩阵值 \(q\) 形式，规定
\[
	(B\wedge C)^a{}_b=\sum_c B^a{}_c\wedge C^c{}_b.
\]
它结合，但一般不满足普通标量形式的分次交换律。特别地，\(A\wedge A\) 未必为零，因为不同矩阵系数未必交换。矩阵乘列形式也按同样规则定义。

\begin{proposition}[外协变微分]\label{g1:prop:exterior-covariant-derivative}
	设 \(\nabla\) 为实光滑向量丛 \(E\to M\) 上的联络。存在唯一的实线性局部算子
	\[
		\mathrm{d}_\nabla:\Omega^k(M;E)\longrightarrow\Omega^{k+1}(M;E)
	\]
	在零次截面上等于 \(\nabla\)，并满足对标量 \(p\) 形式 \(\alpha\) 与丛值形式 \(\eta\)，
	\[
		\mathrm{d}_\nabla(\alpha\wedge\eta)
		=\mathrm{d}\alpha\wedge\eta+(-1)^p\alpha\wedge\mathrm{d}_\nabla\eta.
	\]
	在联络矩阵为 \(A\) 的标架中，它由
	\begin{equation}\label{g1:eq:exterior-covariant-local}
		\mathrm{d}_\nabla(e\eta_e)=e(\mathrm{d}\eta_e+A\wedge\eta_e)
	\end{equation}
	给出，称为\term[waixiebianweifen]{外协变微分}[exterior covariant derivative]。
\end{proposition}
\begin{proof}
	先按局部公式定义。令 \(e'=eg\)，则\eqref{g1:eq:connection-frame-change}给出 \(A'=g^{-1}Ag+g^{-1}\mathrm{d}g\)，而微分 \(g^{-1}g=I\) 得
	\[
		\mathrm{d}(g^{-1})=-g^{-1}(\mathrm{d}g)g^{-1}.
	\]
	于是
	\begin{align*}
		\mathrm{d}(g^{-1}\eta_e)+A'\wedge(g^{-1}\eta_e)
		&=-g^{-1}(\mathrm{d}g)g^{-1}\wedge\eta_e+g^{-1}\mathrm{d}\eta_e\\
		&\quad+g^{-1}A\wedge\eta_e+g^{-1}(\mathrm{d}g)g^{-1}\wedge\eta_e\\
		&=g^{-1}(\mathrm{d}\eta_e+A\wedge\eta_e).
	\end{align*}
	结果按丛值形式的分量变换，因而粘成整体算子。它在零次时就是联络，分次乘积法则则由普通外微分的规则及 \(A\wedge\alpha=(-1)^p\alpha\wedge A\) 得到；这里 \(\alpha\) 是标量，允许这样换序。

	反过来，任意局部丛值形式都是 \(\sum_a\eta^a\otimes e_a\)。对每项使用规定的乘积法则，结果只能为 \(\mathrm{d}\eta^a\otimes e_a+(-1)^k\eta^a\wedge\nabla e_a\)，移过一次形式后正是\eqref{g1:eq:exterior-covariant-local}，所以唯一。
\end{proof}

后文有时将 \(\mathrm{d}_\nabla\) 简记为 \(\mathrm{D}\)，将自同态丛诱导联络的外协变微分简记为 \(\mathrm{D}_{\operatorname{End}}\)。普通外微分平方为零，外协变微分却不一定如此。差别正好是曲率。

\begin{theorem}[曲率的结构方程]\label{g1:thm:curvature-structure-equation}
	设 \(\nabla\) 为实光滑向量丛 \(E\to M\) 上的联络，\(A\) 是某局部标架中的联络矩阵。曲率 \(R\) 在该标架中的矩阵值二形式为
	\begin{equation}\label{g1:eq:curvature-matrix}
		F=\mathrm{d}A+A\wedge A.
	\end{equation}
	对每个 \(\eta\in\Omega^k(M;E)\)，有 \(\mathrm{d}_\nabla^2\eta=R\wedge\eta\)，局部分量即 \(F\wedge\eta_e\)。
\end{theorem}
\begin{proof}
	由局部公式及 \(\mathrm{d}^2=0\)，逐项计算
	\begin{align*}
		(\mathrm{d}+A\wedge)^2\eta_e
		&=\mathrm{d}A\wedge\eta_e-A\wedge\mathrm{d}\eta_e
		+A\wedge\mathrm{d}\eta_e+A\wedge A\wedge\eta_e\\
		&=(\mathrm{d}A+A\wedge A)\wedge\eta_e.
	\end{align*}
	还须确认这确实是先前定义的曲率。对零次截面 \(s=e u\)，把两次协变导数展开并扣除括号项，导数作用于 \(u\) 的项抵消，余下
	\[
		e^{-1}R(X,Y)s
		=\bigl(X(A(Y))-Y(A(X))-A([X,Y])
		+A(X)A(Y)-A(Y)A(X)\bigr)u.
	\]
	由\eqref{g1:eq:exterior-one-form-evaluation}，前三项是 \(\mathrm{d}A(X,Y)\)，后两项是 \((A\wedge A)(X,Y)\)，故所列 \(F\) 表示 \(R\)。一般次数的结论随第一段计算得到。
\end{proof}

在坐标 \(x^i\) 中写 \(A=\sum_i A_i\,\mathrm{d}x^i\)，则
\[
	F=\sum_{i<j}F_{ij}\,\mathrm{d}x^i\wedge\mathrm{d}x^j,
	\qquad F_{ij}=\partial_iA_j-\partial_jA_i+[A_i,A_j].
\]
这里对 \(i<j\) 求和，没有额外的 \(1/2\)。最后一项是矩阵交换子，它记录纤维方向的两次线性变化不交换，不能用标量一形式自楔为零的规则删掉。

\begin{example}[常系数也可能有曲率]\label{g1:exm:constant-noncommuting-connection}
	在 \(\mathbb R^2\times\mathbb R^2\) 的标准标架中取
	\[
		A=B\,\mathrm{d}x+C\,\mathrm{d}y,\qquad
		B=\begin{pmatrix}0&1\\0&0\end{pmatrix},\quad
		C=\begin{pmatrix}0&0\\1&0\end{pmatrix}.
	\]
	虽然 \(\mathrm{d}A=0\)，但
	\[
		F=[B,C]\,\mathrm{d}x\wedge\mathrm{d}y
		=\begin{pmatrix}1&0\\0&-1\end{pmatrix}\mathrm{d}x\wedge\mathrm{d}y\ne0.
	\]
	因此系数恒定、底空间单连通，都不能替代曲率为零。沿一条路径的方程只涉及一个时间变量；同时要求两个方向上的运输相容，才会发现这里的交换子。
\end{example}

在线丛的局部标架中，\(A\) 是标量一形式，\(A\wedge A=0\)，所以 \(F=\mathrm{d}A\)。在第九章的正向正交曲面标架中，\eqref{g1:eq:surface-frame-connection}给出
\[
	A=\begin{pmatrix}0&-\omega\\\omega&0\end{pmatrix},\qquad
	F=\begin{pmatrix}0&-\mathrm{d}\omega\\\mathrm{d}\omega&0\end{pmatrix}.
\]
该例中 \(A\wedge A=0\)，并由\eqref{g1:eq:surface-curvature-form}得到 \(\mathrm{d}\omega=-K\,\mathrm{d}A\)。这里最后的 \(\mathrm{d}A\) 沿用曲面面积形式的记号，与前面“对矩阵 \(A\) 求外微分”的 \(\mathrm{d}A\) 对象不同；在一般丛公式中，\(A\) 始终表示联络矩阵。

\subsection{曲率的换标架规律与小环路运输}
\label{g1:subsec:1503:03}

联络矩阵的变换律含有 \(\mathrm{d}g\)，说明它不按张量变换。曲率已经证明是张量，因而应当只发生纤维基变换，不再出现新的导数项。

\begin{proposition}[曲率的换标架公式]\label{g1:prop:curvature-frame-change}
	若同一实向量丛联络在标架 \(e\) 与 \(e'=eg\) 中的曲率矩阵分别为 \(F,F'\)，则
	\[
		F'=g^{-1}Fg.
	\]
	因此曲率是否为零与标架无关；对固定的切向输入，所得自同态的迹、行列式等共轭不变量也与标架无关。
\end{proposition}
\begin{proof}
	对任意列向量 \(u'\)，同一截面为 \(s=e'u'=e(gu')\)。曲率的纤维线性给出 \(R(X,Y)s=eF(X,Y)g u'\)，另一方面它等于 \(e'F'(X,Y)u'=egF'(X,Y)u'\)。因此 \(Fg=gF'\)，得到结论。这个证明也说明变换律中没有导数项的原因：\(R\) 对截面是逐点线性的。
\end{proof}

\begin{proposition}[小方环的运输]\label{g1:prop:small-loop-curvature}
	设 \(p\) 是光滑流形 \(M\) 的内点，\(\dim M\geq2\)，\(E\to M\) 为秩 \(r\) 实光滑向量丛，\(\nabla\) 为其光滑联络。取以 \(p\) 为原点的坐标及丛标架，记联络矩阵为 \(A\)，曲率矩阵为 \(F=\mathrm{d}A+A\wedge A=\sum_{i<j}F_{ij}\,\mathrm{d}x^i\wedge\mathrm{d}x^j\)。令 \(\gamma_\varepsilon\) 依次沿坐标方形的 \(+x^1,+x^2,-x^1,-x^2\) 四边行进，每边坐标长度为 \(\varepsilon>0\)，其余坐标为零。则在 \(E_p\) 的该标架中，
	\begin{equation}\label{g1:eq:small-loop-curvature}
		P_{\gamma_\varepsilon}
		=I-\varepsilon^2F_{12}(p)+O(\varepsilon^3).
	\end{equation}
	余项指选定矩阵范数下存在与充分小 \(\varepsilon\) 无关的常数 \(C\)，使其范数不超过 \(C\varepsilon^3\)。
\end{proposition}
\begin{proof}
	先用换标架消去一点处的联络系数。令 \(g(x)=I-\sum_i x^iA_i(p)\)，缩小邻域可使 \(g\) 可逆。由变换律，\(g(p)=I\) 且 \(A'_i(p)=A_i(p)+\partial_i g(p)=0\)。新旧标架在基点相同，故返回映射的基点矩阵不变，曲率矩阵在基点也不变。以下在这组新标架中计算并省去撇号。

	四边上 \(A_i=O(\varepsilon)\)。把一条边用坐标长度参数化，平行移动的矩阵满足 \(U'=-B(t)U\)、\(U(0)=I\)，其中时间长度为 \(\varepsilon\)、\(\|B\|\leq C\varepsilon\)。积分方程及\cref{g1:lem:gronwall-estimate}给出 \(\|U\|\leq e^{C\varepsilon^2}\) 和 \(U-I=O(\varepsilon^2)\)，于是再次代回积分方程得
	\[
		U=I-\int_{\text{该边}}A+O(\varepsilon^4).
	\]
	把四段按实际行进顺序相乘，两个积分项的乘积也是 \(O(\varepsilon^4)\)，所以
	\[
		P_{\gamma_\varepsilon}=I-\oint_{\gamma_\varepsilon}A+O(\varepsilon^4).
	\]
	对每个矩阵条目用\cref{g1:thm:stokes-corners}的矩形情形，按\eqref{g1:eq:rectangle-boundary-signs}取正向边界，再对系数作 Taylor 展开，有
	\[
		\oint_{\gamma_\varepsilon}A
		=\int_0^\varepsilon\int_0^\varepsilon
		(\partial_1A_2-\partial_2A_1)(s,t,0,\ldots,0)\,\mathrm{d}s\,\mathrm{d}t
		=\varepsilon^2F_{12}(p)+O(\varepsilon^3).
	\]
	最后一步使用光滑性和 \(A(p)=0\)，此时基点的交换子项为零。结合两式得到结论。矩形缩在内点坐标域内，故没有额外的边界条件。
\end{proof}

\input{Figures/Chapter15-CurvatureLoop.tex}

\cref{g1:fig:curvature-loop}中的两条折线路径有相同端点，但沿第一条出发、沿第二条反向返回，会形成这个方环。\eqref{g1:eq:small-loop-curvature}说明它们的差首先在面积阶出现，系数就是曲率。负号来自平行方程 \(u'=-A(\dot\gamma)u\) 和所选绕行方向；换成反向环路便反号。一般矩阵系数不交换，不能把小环路的一阶面积近似误写成任意大环路上的 \(P_\gamma=\exp(-\int F)\)。

\subsection{Bianchi 恒等式的结构证明}
\label{g1:subsec:1503:04}

曲率是联络的二阶变化，它自身并不是任意的矩阵值二形式：结构方程强制它满足一个协变微分恒等式。为了准确写出这个式子，先将\cref{g1:prop:induced-connections}给出的 \(\operatorname{End}E\) 联络延拓到自同态值形式。对矩阵值 \(k\) 形式 \(B\)，其外协变微分为
\begin{equation}\label{g1:eq:endomorphism-exterior-derivative}
	\mathrm{d}_{\nabla^{\operatorname{End}}}B
	=\mathrm{d}B+A\wedge B-(-1)^k B\wedge A.
\end{equation}
这里把整体自同态值形式与其局部矩阵暂用同一字母。对 \(B=\beta C\)，其中 \(\beta\) 为标量 \(k\) 形式、\(C\) 为自同态矩阵函数，先用 \(\nabla^{\operatorname{End}}C=\mathrm{d}C+AC-CA\)，再用外协变微分的分次乘积法则，就得到该式；一般 \(B\) 是这类项之和。因而它来自已定义的诱导联络，而不是任意补入的修正项。

下面的\term[bianchihengdengshi]{Bianchi 恒等式}[Bianchi identity]\footnote{Bianchi，Luigi Bianchi（1856--1928），意大利数学家，研究曲面、Riemann 几何及连续变换群。}表达曲率的这种相容性。

\begin{theorem}[Bianchi 恒等式]\label{g1:thm:bianchi-identity}
	设 \(E\to M\) 为实光滑向量丛，\(\nabla\) 为光滑联络，\(R\) 为其曲率，\(\nabla^{\operatorname{End}}\) 为在 \(\operatorname{End}E\) 上诱导的联络。在任一光滑局部标架中，记 \(A,F\) 分别为联络及曲率的矩阵形式。则
	\[
		\mathrm{d}_{\nabla^{\operatorname{End}}}R=0.
	\]
	局部矩阵形式为
	\begin{equation}\label{g1:eq:bianchi-matrix}
		\mathrm{d}F+A\wedge F-F\wedge A=0.
	\end{equation}
\end{theorem}
\begin{proof}
	代入结构方程 \(F=\mathrm{d}A+A\wedge A\)。由于 \(A\) 是一次形式，
	\[
		\mathrm{d}F=\mathrm{d}A\wedge A-A\wedge\mathrm{d}A.
	\]
	另外
	\[
		A\wedge F-F\wedge A
		=A\wedge\mathrm{d}A+A\wedge A\wedge A
		-\mathrm{d}A\wedge A-A\wedge A\wedge A.
	\]
	两式相加，各项成对抵消。上述局部表达是整体 \(\operatorname{End}E\) 联络的外协变微分，所以每张图中为零就意味着整体为零。整个论证只用外微分平方为零、分次乘积法则及矩阵乘法结合性，不要求度量或无挠条件。
\end{proof}

这个恒等式通常称为第二 Bianchi 恒等式；切丛上另有涉及挠率的第一 Bianchi 恒等式，将在下一节得到。

\ProseParagraphBreak

在坐标向量场上，它的可计算形式是
\[
	\partial_iF_{jk}+\partial_jF_{ki}+\partial_kF_{ij}
	+[A_i,F_{jk}]+[A_j,F_{ki}]+[A_k,F_{ij}]=0.
\]
因此一般不能删掉交换子而断言 \(\mathrm{d}F=0\)。在线丛上矩阵系数交换，式子才退化为普通的 \(\mathrm{d}F=0\)。在二维底流形上，恒等式因三形式为零而自动成立；它既不迫使二维曲率为零，也不说明联络可以取成零矩阵。

% 纲要条目排为编号小节标题，并进入目录与书签；说明段仍保留于纲要注释中。

% !TeX root = ../../Book1.tex
% 纲要标题：### 15.4 切丛联络与挠率
% 纲要位置：几何学计划纲要.md，第 715 行。
% 正文状态：初稿；尚未正式审核。

% 纲要内容开始
% 
% * 挠率的坐标无关定义
% * 无挠性和坐标标架
% * 第一结构方程
% * 切丛与一般向量丛的区别
% 
% 纲要内容结束

\section{切丛联络与挠率}
\label{g1:sec:1504}

一般向量丛的联络接受一个切向量方向和一个丛截面，输出同一丛的截面。对切丛，方向场和被微分的截面恰好都是向量场，因而可以交换二者，还可以把结果与 Lie 括号比较。这种额外的比较产生挠率。曲率考察同一向量连续经过两次协变微分后的差别；挠率则从一次协变微分出发，比较方向与被微分向量的交换。

\ProseParagraphBreak

本节 \(\nabla\) 是 \(\mathrm{T}M\) 上的光滑联络，\(M\) 有限维、Hausdorff、第二可数，允许边界。涉及坐标向量场的局部可实现性时，将明确限制在内点。度量不是挠率定义所需的数据，以下也不默认联络保持任何度量。

\subsection{挠率的坐标无关定义}
\label{g1:subsec:1504:01}

在 Euclid 空间的通常导数下，\(\partial_XY-\partial_YX=[X,Y]\)。因此，直接把 \(\nabla_XY-\nabla_YX\) 当成偏离 Euclid 情形的量并不合适：即使使用通常导数，只要所选方向场不交换，它也可能不为零。应当先减去由方向场本身造成的 Lie 括号。

\begin{definition}[挠率]\label{g1:def:torsion}
	切丛联络 \(\nabla\) 的\term[naolv]{挠率}[torsion]定义为
	\begin{equation}\label{g1:eq:torsion-definition}
		T^\nabla(X,Y)=\nabla_XY-\nabla_YX-[X,Y].
	\end{equation}
	若 \(T^\nabla=0\)，称 \(\nabla\) 为\term[wunaolianluo]{无挠联络}[torsion-free connection]。下文不致混淆时简记为 \(T\)。
\end{definition}

第七章也用“挠率”表示空间曲线的标量 \(\tau\)，它描述 Frenet 标架的转动。这里的 \(T\) 是切丛联络的二输入张量，二者的定义和所需数据不同，不能因中文名称相同便代入同一个公式。

\begin{proposition}[挠率的张量性]\label{g1:prop:torsion-tensorial}
	设 \(M\) 为可带边界的光滑流形，\(\nabla\) 为 \(\mathrm{T}M\) 上的光滑线性联络，\(T\) 为其挠率。则 \(T\) 是一个光滑的切丛值二形式，即 \(T\in\Omega^2(M;\mathrm{T}M)\)。在任意局部标架 \((e_i)\) 中，若
	\[
		\nabla_{e_i}e_j=\sum_k\Gamma^k_{ij}e_k,\qquad
		[e_i,e_j]=\sum_k c^k_{ij}e_k,
	\]
	则
	\begin{equation}\label{g1:eq:torsion-frame-components}
		T(e_i,e_j)=\sum_k\bigl(\Gamma^k_{ij}-\Gamma^k_{ji}-c^k_{ij}\bigr)e_k.
	\end{equation}
\end{proposition}
\begin{proof}
	由定义，\(T\) 实双线性且反对称。用\eqref{g1:eq:lie-bracket-product}中的 \([fX,Y]=f[X,Y]-Y(f)X\)，得到
	\begin{align*}
		T(fX,Y)
		&=f\nabla_XY-\nabla_Y(fX)-[fX,Y]\\
		&=f\nabla_XY-Y(f)X-f\nabla_YX-f[X,Y]+Y(f)X\\
		&=fT(X,Y).
	\end{align*}
	反对称性随后给出第二个输入的函数线性。因此在局部标架中展开 \(X,Y\)，结果只依赖它们在求值点的分量；各分量又光滑，故为光滑张量。把两个输入取成 \(e_i,e_j\)，直接代入定义便得分量公式。边界处同样的局部延拓计算成立。
\end{proof}

在\secref{g1:sec:0803}中，曲面的诱导联络已经满足 \(\nabla_XY-\nabla_YX=[X,Y]\)，见\cref{g1:prop:surface-connection-properties}。现在可以把该性质简洁地表述为“曲面的诱导联络无挠”。这不是所有切丛联络自动具有的性质。

\subsection{无挠性和坐标标架}
\label{g1:subsec:1504:02}

坐标标架满足 \([\partial_i,\partial_j]=0\)，所以\eqref{g1:eq:torsion-frame-components}化为
\[
	T(\partial_i,\partial_j)=\sum_k(\Gamma^k_{ij}-\Gamma^k_{ji})\partial_k.
\]
于是无挠性等价于每个坐标图中 \(\Gamma^k_{ij}=\Gamma^k_{ji}\)。这里两个下指标分别表示微分方向和被微分向量；只有在坐标标架中，才能不加修正地使用这条对称性判据。一般标架中的正确判据是 \(\Gamma^k_{ij}-\Gamma^k_{ji}=c^k_{ij}\)。

\begin{proposition}[无挠联络的存在与差]\label{g1:prop:torsion-free-affine-space}
	设 \(M\) 为有限维、Hausdorff、第二可数的光滑流形，允许有边界。则 \(\mathrm{T}M\) 上存在无挠光滑联络。对任意两个光滑切丛联络 \(\nabla,\nabla'\)，若 \(\nabla' =\nabla+B\)，其中 \(B\in\Omega^1(M;\operatorname{End}(\mathrm{T}M))\)，则
	\begin{equation}\label{g1:eq:torsion-difference}
		T^{\nabla'}(X,Y)=T^\nabla(X,Y)+B(X)Y-B(Y)X.
	\end{equation}
	无挠联络组成一个仿射空间，其平移向量空间为 \(\Gamma(\operatorname{Sym}^2\mathrm{T}^*M\otimes \mathrm{T}M)\)。
\end{proposition}
\begin{proof}
	将 \(\nabla'\) 代入挠率定义，Lie 括号项相消，就得到差公式。由\cref{g1:thm:connection-existence}，先任选一个切丛联络 \(\nabla\)。其挠率已经证明为张量，所以
	\[
		\nabla^0_XY=\nabla_XY-\tfrac12T^\nabla(X,Y)
	\]
	仍满足联络的两条函数乘法规则。利用 \(T^\nabla(Y,X)=-T^\nabla(X,Y)\)，差公式给出 \(T^{\nabla^0}=T^\nabla-\tfrac12T^\nabla-\tfrac12T^\nabla=0\)。若 \(\nabla\) 已无挠，则 \(\nabla+B\) 无挠当且仅当 \(B(X)Y=B(Y)X\)，这正是 \(\mathrm{T}M\) 值二阶协变张量的对称性。联络之差的唯一性和仿射性质已由\cref{g1:prop:connection-affine-space}给出，限制到该对称子空间即得结论。
\end{proof}

这个构造只解决无挠性。即使原联络保持某个度量，减去一半挠率后也未必继续保持它；章末习题将给出直接计算。我们应当把无挠性、度量相容性和平坦性分别检查。

\begin{example}\label{g1:exm:torsion-free-curved}
	在 \(\mathbb R^2\) 的坐标标架 \((\partial_x,\partial_y)\) 中规定
	\[
		\nabla_{\partial_y}\partial_y=x\partial_x,
		\qquad\nabla_{\partial_x}\partial_x
		=\nabla_{\partial_x}\partial_y
		=\nabla_{\partial_y}\partial_x=0.
	\]
	各 \(\Gamma^k_{ij}\) 对 \(i,j\) 对称，所以 \(T=0\)。然而
	\[
		A=\begin{pmatrix}0&x\,\mathrm{d}y\\0&0\end{pmatrix},\qquad
		F=\mathrm{d}A+A\wedge A=\begin{pmatrix}0&\mathrm{d}x\wedge\mathrm{d}y\\0&0\end{pmatrix}.
	\]
	特别地，\(R(\partial_x,\partial_y)\partial_y=\partial_x\ne0\)。无挠并不保证平坦。
\end{example}

\begin{example}\label{g1:exm:flat-with-torsion}
	仍在 \(\mathbb R^2\)，取处处线性无关的向量场 \(e_1=\partial_x\)、\(e_2=e^x\partial_y\)，规定 \(\nabla_Xe_i=0\) 对所有 \(X\) 成立。对 \(Y=y^1e_1+y^2e_2\)，这具体表示 \(\nabla_XY=X(y^1)e_1+X(y^2)e_2\)，因而确实定义联络。该标架中 \(A=0\)，故 \(F=0\)。但
	\[
		[e_1,e_2]=e_2,\qquad T(e_1,e_2)=-e_2\ne0.
	\]
	若取 \(g=\mathrm{d}x^2+e^{-2x}\mathrm{d}y^2\)，则 \(e_1,e_2\) 为正交规范标架，平行移动保持其分量，因而保持 \(g\)。这同时给出平坦、保持度量而有挠率的联络。
\end{example}

后一个例子还能说明“把联络矩阵变成零”和“找到使所有坐标向量场平行的坐标”不同。它已有整体平行标架，但该标架的 Lie 括号不为零，不能成为坐标标架。

\subsection{第一结构方程}
\label{g1:subsec:1504:03}

设 \(e=(e_1,\ldots,e_n)\) 是切丛局部行标架，\(\theta=(\theta^1,\ldots,\theta^n)^{\mathsf T}\) 是对偶余标架的列，满足 \(\theta^i(e_j)=\delta^i_j\)。这样每个向量场都可写成 \(X=e\theta(X)\)。令 \(\nabla e=eA\)，并将挠率写成 \(T(X,Y)=e\Theta(X,Y)\)，其中 \(\Theta\) 是二形式列。

\begin{proposition}[第一结构方程]\label{g1:prop:first-structure-equation}
	设 \(M\) 为可带边界的光滑 \(n\) 维流形，\(\nabla\) 为其光滑切丛联络，\(T\) 为挠率。在开集 \(U\) 上取光滑行标架 \(e=(e_1,\ldots,e_n)\) 及对偶余标架列 \(\theta\)，令 \(\nabla e=eA\)，并以 \(T(X,Y)=e\Theta(X,Y)\) 定义挠率分量列 \(\Theta\)。则
	\begin{equation}\label{g1:eq:first-structure-equation}
		\Theta=\mathrm{d}\theta+A\wedge\theta.
	\end{equation}
	若 \(g:U\to\operatorname{GL}(n,\mathbb R)\) 光滑，且换标架为 \(e'=eg\)，则 \(\theta'=g^{-1}\theta\)、\(\Theta'=g^{-1}\Theta\)。因而右侧虽由局部形式组成，却表示同一个整体切丛值二形式。
\end{proposition}
\begin{proof}
	联络的局部表达给出
	\[
		\nabla_XY=e\bigl(X(\theta(Y))+A(X)\theta(Y)\bigr).
	\]
	交换 \(X,Y\) 并减去 \([X,Y]=e\theta([X,Y])\)，分量成为
	\[
		X(\theta(Y))-Y(\theta(X))-\theta([X,Y])
		+A(X)\theta(Y)-A(Y)\theta(X).
	\]
	由\eqref{g1:eq:exterior-one-form-evaluation}，前三项正是 \(\mathrm{d}\theta(X,Y)\)；后两项正是 \((A\wedge\theta)(X,Y)\)。这证明第一结构方程。

	由 \(X=e\theta(X)=eg\theta'(X)\) 得 \(\theta'=g^{-1}\theta\)，由 \(T=e\Theta=eg\Theta'\) 得 \(\Theta'=g^{-1}\Theta\)。也可以使用\eqref{g1:eq:connection-frame-change}直接验算：\(\mathrm{d}(g^{-1})\wedge\theta\) 与联络变换中的 \(g^{-1}\mathrm{d}g\wedge g^{-1}\theta\) 恰好抵消，留下 \(g^{-1}(\mathrm{d}\theta+A\wedge\theta)\)。
\end{proof}

现在有两条结构方程：第一条 \(\Theta=\mathrm{d}\theta+A\wedge\theta\) 描述挠率，第二条 \(F=\mathrm{d}A+A\wedge A\) 描述曲率。第一条需要切丛自带的余标架；第二条对任意向量丛成立。在坐标标架中 \(\theta^i=\mathrm{d}x^i\)，所以 \(\mathrm{d}\theta=0\)，无挠性变成 \(A\wedge\theta=0\)。一般标架下不能漏掉 \(\mathrm{d}\theta\)。

\begin{example}\label{g1:exm:polar-torsion-cancellation}
	在极坐标局部区域 \(r>0\) 中，Euclid 平面的正交规范标架为 \(e_1=\partial_r\)、\(e_2=r^{-1}\partial_\varphi\)，其对偶列为 \(\theta=(\mathrm{d}r,r\,\mathrm{d}\varphi)^{\mathsf T}\)。把二者写成笛卡儿分量，\(e_1=(\cos\varphi,\sin\varphi)\)、\(e_2=(-\sin\varphi,\cos\varphi)\)。直接微分得到通常 Euclid 联络的矩阵
	\[
		A=\begin{pmatrix}0&-\mathrm{d}\varphi\\\mathrm{d}\varphi&0\end{pmatrix}.
	\]
	于是
	\[
		\mathrm{d}\theta=\begin{pmatrix}0\\\mathrm{d}r\wedge\mathrm{d}\varphi\end{pmatrix},\qquad
		A\wedge\theta=\begin{pmatrix}0\\-\mathrm{d}r\wedge\mathrm{d}\varphi\end{pmatrix}.
	\]
	两项抵消，\(\Theta=0\)。虽然 \([e_1,e_2]=-r^{-1}e_2\ne0\)，联络仍无挠；标架的不交换性已被 \(\nabla_{e_1}e_2-\nabla_{e_2}e_1\) 准确补偿。这里 \(\mathrm{d}A=0\)、\(A\wedge A=0\)，故曲率也为零。
\end{example}

\begin{proposition}[第一 Bianchi 恒等式]\label{g1:prop:first-bianchi-identity}
	设 \(M\) 为可带边界的光滑流形，\(\nabla\) 为其光滑切丛联络，\(T,R\) 分别为挠率与曲率。任取光滑局部行标架 \(e\) 及对偶余标架列 \(\theta\)，以 \(\nabla e=eA\)、\(T(X,Y)=e\Theta(X,Y)\)、\(R(X,Y)(eu)=eF(X,Y)u\) 定义 \(A,\Theta,F\)。则
	\begin{equation}\label{g1:eq:first-bianchi-identity}
		\mathrm{d}\Theta+A\wedge\Theta=F\wedge\theta.
	\end{equation}
	特别地，若联络无挠，则对任意光滑向量场 \(X,Y,Z\) 有
	\begin{equation}\label{g1:eq:torsion-free-cyclic-curvature}
		R(X,Y)Z+R(Y,Z)X+R(Z,X)Y=0.
	\end{equation}
\end{proposition}
\begin{proof}
	恒等映射 \(\operatorname{Id}_{\mathrm{T}M}\) 可看成一个切丛值一形式，它在 \(X\) 上取值 \(X\)，在标架 \(e\) 中的分量正是 \(\theta\)。第一结构方程说明 \(\mathrm{D}\operatorname{Id}_{\mathrm{T}M}=T\)。对它再作一次外协变微分，\cref{g1:thm:curvature-structure-equation}中的 \(\mathrm{D}^2\eta=F\wedge\eta\) 给出 \(\mathrm{D}T=R\wedge\operatorname{Id}_{\mathrm{T}M}\)，其分量就是\eqref{g1:eq:first-bianchi-identity}。

	若 \(T=0\)，左侧为零。右侧在 \(X,Y,Z\) 上的值为
	\[
		R(X,Y)Z-R(X,Z)Y+R(Y,Z)X.
	\]
	这是二形式与一形式楔积的三个保序项；用 \(R(X,Z)=-R(Z,X)\)，便得到循环恒等式。
\end{proof}

这里的“第一”恒等式把挠率与曲率联系起来；上一节证明的 \(\mathrm{D}_{\operatorname{End}}R=0\) 通常称为第二 Bianchi 恒等式，它对任意向量丛联络成立，并不要求无挠。对有挠率的切丛联络，不能直接套用\eqref{g1:eq:torsion-free-cyclic-curvature}，必须保留\eqref{g1:eq:first-bianchi-identity}的左侧。

\subsection{切丛与一般向量丛的区别}
\label{g1:subsec:1504:04}

挠率构造中的特殊数据，是 \(\mathrm{T}M\) 的纤维正好就是底空间的切向量空间。因此存在无需选择的 \(\mathrm{T}M\) 值一形式 \(X\mapsto X\)，也能对两个切丛截面取 Lie 括号。一般向量丛 \(E\) 的纤维可能是余切空间、张量空间，或者与切向量没有自然关系的向量空间；既没有默认的映射 \(\mathrm{T}M\to E\)，也没有默认的截面 Lie 括号。仅给 \(E\) 一个联络，便没有相应的规范挠率。

\ProseParagraphBreak

如果另行给出 \(E\) 值一形式 \(\eta\)，当然可以构造 \(\mathrm{D}\eta\)。但这是依赖 \(\eta\) 的额外对象，不能只由联络决定。切丛的情形恰好有一个天然选择 \(\eta=\operatorname{Id}_{\mathrm{T}M}\)，才使 \(\mathrm{D}\eta\) 成为联络自身的挠率。

\begin{proposition}[平行标架的坐标实现]\label{g1:prop:parallel-coordinate-frame}
	设 \(M\) 为光滑 \(n\) 维流形，\(U\subseteq\operatorname{Int}M\) 为开集，\(\nabla\) 为 \(\mathrm{T}U\) 上的光滑线性联络，\(e_1,\ldots,e_n\) 是相对于 \(\nabla\) 的光滑平行标架，即 \(\nabla_Xe_i=0\) 对所有向量场 \(X\) 成立。则下列条件等价：联络在 \(U\) 上无挠；\([e_i,e_j]=0\) 对所有 \(i,j\) 成立；每点附近都有坐标 \((y^1,\ldots,y^n)\) 使 \(e_i=\partial/\partial y^i\)。
\end{proposition}
\begin{proof}
	平行性使 \(T(e_i,e_j)=-[e_i,e_j]\)，而挠率是张量，因此前两个条件等价。坐标向量场两两交换，所以第三个条件蕴含第二个。

	现设括号全部为零，取对偶余标架 \(\theta^i\)。由于 \(\theta^i(e_j)=\delta^i_j\) 是常数，外微分的一形式公式给出
	\[
		\mathrm{d}\theta^i(e_j,e_k)=-\theta^i([e_j,e_k])=0.
	\]
	这些向量场张成切空间，故 \(\mathrm{d}\theta^i=0\)。由\cref{g1:cor:local-exactness}，在每个内点的足够小共同邻域上存在光滑函数 \(y^i\)，使 \(\theta^i=\mathrm{d}y^i\)。提醒这里调用的是正次闭形式的局部恰当性，无须这些一形式在整个 \(U\) 上恰当。映射 \(y=(y^1,\ldots,y^n)\) 的微分由余标架组成，处处可逆。\cref{g1:thm:inverse-function}于是给出局部坐标，且 \(\mathrm{d}y^i(e_j)=\delta^i_j\) 强制 \(e_j=\partial/\partial y^j\)。
\end{proof}

下一节将证明，曲率为零恰好允许在足够小的邻域中构造平行标架。与此命题结合，便得到流形内部的结论：平坦且无挠的切丛联络局部可以用联络系数全为零的坐标表示。曲率消失负责产生平行标架，挠率消失负责使这个标架来自坐标。\cref{g1:exm:flat-with-torsion}说明第二步确有独立条件。

\ProseParagraphBreak

边界处须保留区别。闭余标架仍有局部原函数，它们也给出到 Euclid 空间的局部光滑嵌入，但像的边界未必是一个坐标超平面，不能因此断言它们构成标准半空间型边界坐标。例如，Euclid 平面中闭圆盘的通常联络平坦且无挠；若新坐标也使联络系数为零，则逆坐标映射的所有二阶导数都为零，因而它必须局部为仿射映射，而仿射变换不能把圆弧变成直线。因而上面的坐标实现命题特意陈述在内点，边界邻域的平行标架则不受这个限制。

% !TeX root = ../../Book1.tex
% 纲要标题：### 15.5 平坦性、单值与局部实现
% 纲要位置：几何学计划纲要.md，第 722 行。
% 正文状态：初稿；尚未正式审核。

% 纲要内容开始
% 
% * 平坦联络的平行标架构造
% * 单连通区域上的路径无关性
% * 非单连通情形的单值表示
% * 局部曲率与整体平坦丛的区别
% 
% ---
% 
% 纲要内容结束

\section{平坦性、单值与局部实现}
\label{g1:sec:1505}

沿一条给定路径，平行移动总能由常微分方程求出。但如果希望在一个区域内同时比较所有纤维，就要问：改变路径会不会改变比较结果？曲率记录局部绕行造成的差异；当曲率处处为零时，还须判断一条较大的环路能否在区域内收缩。前者是局部微分条件，后者涉及区域的整体拓扑。

\ProseParagraphBreak

本节的底空间是有限维、Hausdorff、第二可数光滑流形，允许有边界，向量丛为有限秩实光滑丛。沿用行标架 \(e=(e_1,\ldots,e_r)\) 与列分量 \(u\) 的约定：
\[
	\nabla(eu)=e(\mathrm{d}u+A u),\qquad
	A'=g^{-1}Ag+g^{-1}\mathrm{d}g,\qquad
	F=\mathrm{d}A+A\wedge A.
\]
其中 \(e'=eg\)。\cref{g1:def:flat-connection}称曲率恒为零的联络为平坦联络；由\cref{g1:eq:curvature-matrix,g1:prop:curvature-frame-change}，这等价于在任意局部标架中 \(F=0\)。下面把这一矩阵条件转成实际可用的标架与运输。

\subsection{平坦联络的平行标架构造}
\label{g1:subsec:1505:01}

若能换到一个联络矩阵为零的标架，平行方程便成为 \(\mathrm{d}u/\mathrm{d}t=0\)：向量沿路径移动时，其分量保持不变。这正是我们希望由 \(F=0\) 得到的局部模型。困难在于同时消去所有坐标方向的联络系数；只沿一条路径求解，并没有完成这件事。

\begin{definition}
	\label{g1:def:parallel-frame}
	设 \(M\) 为可带边界的光滑流形，\(E\to M\) 为实有限秩光滑向量丛，\(\nabla\) 为其光滑联络，\(U\subset M\) 开。若 \(E|_U\) 的光滑局部标架 \(f=(f_1,\ldots,f_r)\) 对每个 \(a\) 都满足 \(\nabla f_a=0\)，则称 \(f\) 为\term[pingxingbiaojia]{平行标架}[parallel frame]。等价地，对任意光滑列函数 \(u\)，有 \(\nabla(fu)=f\,\mathrm{d}u\)。
\end{definition}

\begin{theorem}[平坦局部模型]
	\label{g1:thm:flat-parallel-frame}
	设 \(E\to M\) 是有限秩实光滑向量丛，\(M\) 是有限维、Hausdorff、第二可数光滑流形，允许有边界，\(\nabla\) 是光滑联络。则 \(\nabla\) 平坦当且仅当每点附近存在平行标架。平坦时，给定 \(p\in M\) 及 \(E_p\) 的一个有序基，可以在某个连通邻域 \(U\ni p\) 上选平行标架，使其在 \(p\) 处等于这个基；一旦 \(U\) 固定，这样的平行标架唯一。
\end{theorem}
\begin{proof}
	先从 \(F=0\) 构造标架。秩为零或底空间维数为零时结论直接成立，以下设秩 \(r\geq1\)、维数 \(n\geq1\)。在 \(p\) 附近取坐标及丛标架 \(e\)，使 \(p\) 对应原点，且 \(e(p)\) 等于指定基。缩小后，坐标域可以取为盒
	\[
		B=(-\varepsilon,\varepsilon)^n
		\quad\text{或}\quad
		B=(-\varepsilon,\varepsilon)^{n-1}\times[0,\varepsilon).
	\]
	后一种用于边界点，最后一个坐标是向内的坐标。写 \(A=\sum_{i=1}^n A_i\,\mathrm{d}x^i\)。变换公式说明，把第一个系数变成零，需要解
	\begin{equation}
		\label{g1:eq:flat-first-coordinate-ode}
		\partial_1g(x^1,x')=-A_1(x^1,x')g(x^1,x'),\qquad
		g(0,x')=I_r,
	\end{equation}
	其中 \(x'=(x^2,\ldots,x^n)\) 暂作参数。

	\cref{g1:thm:ode-local-existence-uniqueness,g1:thm:ode-smooth-dependence}给出共同小时间内的解及对 \(x'\) 的光滑依赖，故适当缩小盒后得到光滑矩阵 \(g\)。这里把矩阵的 \(r^2\) 个条目作为普通 Euclid 未知量即可。为核实 \(g\) 始终可逆，同样解 \(\partial_1k=kA_1\)、\(k(0,x')=I_r\)。乘积法则给出 \(\partial_1(kg)=0\)，所以 \(kg=I_r\)，方阵 \(g\) 因而可逆。边界情形中，光滑系数按定义可延拓到半盒周围的开集；在延拓后的开集应用同一个参数定理，再限制回来，便得到光滑到边界的解。最后剩下一维半区间时，只保留从零开始的非负时间。延拓只用于获得解和参数正则性，不要求延拓后的联络仍然平坦。

	在标架 \(e'=eg\) 中记联络矩阵为 \(B=\sum_iB_i\,\mathrm{d}x^i\)。由\eqref{g1:eq:flat-first-coordinate-ode}，\(B_1=0\)。曲率按共轭变换仍为零，而其 \((1,j)\) 分量为
	\[
		0=F'_{1j}
		=\partial_1B_j-\partial_jB_1+B_1B_j-B_jB_1
		=\partial_1B_j,\qquad j\geq2.
	\]
	固定 \(x'\) 后沿第一坐标积分，便得到 \(B_j(x^1,x')=B_j(0,x')\)。因此余下的联络系数实际上只依赖 \(n-1\) 个变量。

	现在对维数归纳。一维时已经全部消去。若 \(n>1\)，截面 \(x^1=0\) 上的矩阵形式 \(\sum_{j=2}^n B_j(0,x')\,\mathrm{d}x^j\) 仍然平坦，因为它的曲率分量就是原来为零的 \(F'_{ij}\)，其中 \(i,j\geq2\)。归纳假设给出 \(h(x')\in\operatorname{GL}(r,\mathbb R)\)，在原点取 \(I_r\)，并消去这个低一维联络的全部系数。让 \(h\) 沿第一坐标方向保持不变，标架 \(f=e' h\) 的第一系数仍为零，其他系数也为零。归纳中的每一步只需缩小坐标盒，并且对半盒采用刚才的延拓与限制方法，所以边界处同样成立。所得 \(f(p)=e(p)\)，完成存在性。

	反过来，若每点附近有平行标架，则其中 \(A=0\)，所以 \(F=\mathrm{d}A+A\wedge A=0\)。曲率是整体定义的张量，局部为零便处处为零。

	最后设 \(f,\widehat f\) 是连通 \(U\) 上的两个平行标架，写 \(\widehat f=f c\)。由 \(\nabla\widehat f=f\,\mathrm{d}c=0\)，得到 \(\mathrm{d}c=0\)。在每个坐标小盒内，沿线段积分各个条目的导数，可知 \(c\) 局部常值；半盒边界也由同一结论或连续性得到。连通性于是使 \(c\) 为一个常矩阵。若两标架在 \(p\) 相同，则 \(c(p)=I_r\)，所以处处相同。
\end{proof}

证明中 \(F_{1j}=0\) 的作用可以具体看见：消去第一方向以后，它保证其余系数不会沿第一方向重新变化。因此逐方向构造可以继续，而不会破坏已经完成的一步。对一般联络，仍能沿第一方向解同一个方程，但剩下的 \(\partial_1B_j\) 可能不为零，归纳便不能如此进行。

\ProseParagraphBreak

若两个平行标架在某个开集上都有定义，换标架公式中的两套联络矩阵都为零，转移矩阵的微分也就为零。因此它在交叠的每个连通分支上恒定，不同分支上可以取不同的常矩阵。这些常矩阵将用于把一条路径上的局部运输接合起来。

\subsection{单连通区域上的路径无关性}
\label{g1:subsec:1505:02}

\cref{g1:thm:parallel-transport}已经给出分片光滑路径上的平行移动。为了用\cref{g1:def:path-homotopy,g1:def:simply-connected}中的连续同伦比较运输，需要先把运输延伸到连续路径。平行标架提供了这个延伸：在每张图中保持分量，再利用交叠处的常转移矩阵接合，路径本身便无须可微。

\begin{lemma}[连续路径上的平坦运输]
	\label{g1:lem:flat-continuous-transport}
	设 \(E\to M\) 是有限秩实光滑向量丛，底流形有限维、Hausdorff、第二可数，允许有边界，\(\nabla\) 平坦。对每条连续路径 \(\gamma:[0,1]\to M\)，可以定义线性同构 \(P_\gamma:E_{\gamma(0)}\to E_{\gamma(1)}\)，使路径位于一个平行标架域内时，运输保持该标架中的分量。此定义与局部标架及细分无关；在分片光滑路径上等于原来的平行移动，并满足下列规律：对另一连续路径 \(\eta:[0,1]\to M\)，当 \(\gamma(1)=\eta(0)\) 时，令 \(\gamma*\eta\) 先走 \(\gamma\) 后走 \(\eta\)；令 \(\bar\gamma(t)=\gamma(1-t)\)，\(e_p(t)=p\) 为常路径。则
	\[
		P_{\gamma*\eta}=P_\eta\circ P_\gamma,\qquad
		P_{\bar\gamma}=P_\gamma^{-1},\qquad
		P_{e_p}=\operatorname{id}_{E_p}.
	\]
	其中 \(p\in M\)。
\end{lemma}
\begin{proof}
	把平行标架 \(f\) 在每点看成同构 \(f(x):\mathbb R^r\to E_x\)。对其定义域 \(U\) 中的两点 \(x,y\)，令
	\[
		T^U_{y\leftarrow x}=f(y)f(x)^{-1}.
	\]
	若 \(\widehat f=fc\) 是另一个平行标架，而 \(x,y\) 位于交叠的同一个连通分支内，则 \(c(y)=c(x)\)，所以两个标架给出的 \(T\) 相同。特别地，若一段连续路径同时落在两个标架域内，其像是连通集，以上比较便适用。

	平行标架域的原像覆盖紧区间。由\cref{g1:lem:lebesgue-number}，存在有限分点 \(0=t_0<\cdots<t_N=1\)，使每段 \(\gamma([t_{j-1},t_j])\) 都落在某个标架域 \(U_j\)。定义
	\[
		P_\gamma
		=T^{U_N}_{\gamma(t_N)\leftarrow\gamma(t_{N-1})}\circ\cdots\circ
		T^{U_1}_{\gamma(t_1)\leftarrow\gamma(t_0)}.
	\]
	在同一标架域内增加分点，不改变这个复合，因为
	\(T^U_{z\leftarrow y}T^U_{y\leftarrow x}=T^U_{z\leftarrow x}\)。
	两种任意细分可以先取共同细分；每个共同小段同时落在两种选择的标架域内，由前段比较得到相同的因子。所以定义不依赖选择。

	分片光滑路径在平行标架中满足的方程是 \(u'=0\)，故上述局部运输就是\cref{g1:thm:parallel-transport}的解；有限段拼接后仍相同。对于连续路径，拼接式由把两组分点接起来得到，反向式由反转因子的顺序并逐个取逆得到，常路径的全部因子都是恒等。因而这些规律与\cref{g1:prop:parallel-transport-laws}相容。
\end{proof}

\begin{theorem}[平坦运输与路径无关性]
	\label{g1:thm:flat-path-independence}
	设 \(E\to M\) 是有限秩实光滑向量丛，\(M\) 有限维、Hausdorff、第二可数，允许有边界，\(\nabla\) 平坦。若两条连续路径 \(\gamma_0,\gamma_1:[0,1]\to M\) 定端点连续同伦，则 \(P_{\gamma_0}=P_{\gamma_1}\)。特别地，若开集 \(U\subset M\) 非空且单连通，则 \(E|_U\) 上的运输只依赖起点与终点；给定 \(p\in U\) 的任意纤维基，存在唯一的整体平行标架在 \(p\) 取该基。
\end{theorem}
\begin{proof}
	取定端点同伦 \(H:[0,1]^2\to M\)，其中横坐标 \(t\) 沿路径行进，纵坐标 \(s\) 改变路径。将平行标架域的原像作为紧方形的开覆盖。\cref{g1:lem:lebesgue-number}给出一个正数，使直径足够小的每个网格闭矩形都落在某个原像中。于是可以选有限网格
	\[
		0=t_0<\cdots<t_m=1,\qquad
		0=s_0<\cdots<s_\ell=1,
	\]
	使每个小矩形的整个像都位于一个平行标架域内。这里直接细分同伦的参数方形，不需要把 \(H\) 改成光滑映射。

	令 \(h_{ij}\) 表示第 \(j\) 行上从 \((t_i,s_j)\) 到 \((t_{i+1},s_j)\) 的像路径运输，令 \(v_{ij}\) 表示从 \((t_i,s_j)\) 到 \((t_i,s_{j+1})\) 的像路径运输。每个小矩形内，先向右再向上与先向上再向右，只是在同一平行标架中保持相同分量，因此
	\begin{equation}
		\label{g1:eq:flat-homotopy-cell}
		v_{i+1,j}\circ h_{ij}=h_{i,j+1}\circ v_{ij}.
	\end{equation}
	记整行运输为 \(H_j=h_{m-1,j}\circ\cdots\circ h_{0j}\)。从右端开始依次使用\eqref{g1:eq:flat-homotopy-cell}，得到
	\[
		v_{mj}\circ H_j=H_{j+1}\circ v_{0j}.
	\]
	这一步移动的是相邻的复合因子，并未交换不相邻矩阵的次序。由于同伦固定端点，两侧竖边的像都是常路径，故 \(v_{0j}\) 与 \(v_{mj}\) 都是恒等映射。因此 \(H_j=H_{j+1}\)，逐行得到 \(P_{\gamma_0}=P_{\gamma_1}\)。

	现在设 \(U\) 单连通。按照\cref{g1:def:simply-connected,g1:prop:path-classes}，这包括道路连通性，并保证同起终点的任意两条连续路径定端点同伦。固定基 \(b_1,\ldots,b_r\in E_p\)，对每个 \(q\in U\)，任选 \(p\) 到 \(q\) 的路径 \(\gamma\)，定义 \(f_a(q)=P_\gamma b_a\)。前段结论保证定义与路径无关，各点又都是线性同构作用于一个基，所以 \(f_1(q),\ldots,f_r(q)\) 构成基。

	还须证明这些截面光滑且平行。取 \(q\) 附近连通的平行标架域 \(V\subset U\)，标架记为 \(e\)，先固定一条从 \(p\) 到 \(q\) 的路径。在它后面追加 \(V\) 中从 \(q\) 到 \(y\) 的任意路径，便得到
	\[
		(f_1(y),\ldots,f_r(y))
		=e(y)e(q)^{-1}(f_1(q),\ldots,f_r(q)),\qquad y\in V.
	\]
	右边是平行标架乘一个固定可逆矩阵，因此在 \(V\) 上光滑且平行。这样的 \(V\) 覆盖 \(U\)，便得到整体结论。若另一平行标架在 \(p\) 有相同值，沿每条路径其分量也必须保持不变，所以它在所有 \(q\) 处与已构造标架相同。
\end{proof}

\cref{g1:fig:flat-homotopy-grid}把同伦中的比较画在参数方形上；真正需要的只是每个小格有一个共同的平行标架。

\input{Figures/Chapter15-FlatHomotopyGrid.tex}

单连通性在这个证明中只用于最后把所有同端点路径归入同一个同伦类。平坦性本身已经保证连续同伦不改变运输；即使区域并不单连通，这一部分仍然成立。由此剩下的路径依赖可以精确地交给基本群记录。

\subsection{非单连通情形的单值表示}
\label{g1:subsec:1505:03}

固定基点以后，每条环路都给出同一纤维到自身的线性同构。要把它们组织成矩阵表示，必须先核对乘法方向：本书的道路乘积先走左边，函数复合却先作用右边。这与第五章覆盖单值采用右作用的原因相同，见\cref{g1:prop:monodromy-action}。

\begin{proposition}[平坦联络的单值]
	\label{g1:prop:flat-monodromy}
	设 \(E\to M\) 是秩 \(r\) 实光滑向量丛，\(M\) 非空、连通、有限维、Hausdorff、第二可数，允许有边界，\(\nabla\) 平坦。固定 \(p\in M\) 和基同构 \(b:\mathbb R^r\to E_p\)。约定 \(\pi_1(M,p)\) 的乘积 \([\gamma][\eta]=[\gamma*\eta]\) 由先走 \(\gamma\) 后走 \(\eta\) 的拼接表示。对基于 \(p\) 的连续环路 \(\gamma\)，定义
	\[
		C_\gamma=b^{-1}P_\gamma b,\qquad
		\rho([\gamma])\coloneqq C_\gamma^{-1}.
	\]
	则 \(\rho:\pi_1(M,p)\to\operatorname{GL}(r,\mathbb R)\) 是群同态，称为所选基下的\term[danzhibiaoshi]{单值表示}[monodromy representation]。将基改成 \(b'=bB\) 时，表示变成 \(\rho'=B^{-1}\rho B\)。联络存在整体平行标架，当且仅当 \(\rho\) 为平凡表示。
\end{proposition}
\begin{proof}
	\cref{g1:thm:flat-path-independence}使 \(C_\gamma\) 及其逆只依赖环路类。由于先走 \(\gamma\) 再走 \(\eta\)，
	\[
		C_{\gamma*\eta}=C_\eta C_\gamma,\qquad
		\rho([\gamma][\eta])
		=(C_\eta C_\gamma)^{-1}
		=C_\gamma^{-1}C_\eta^{-1}
		=\rho([\gamma])\rho([\eta]).
	\]
	常路径给单位矩阵，所以这确是同态。换基后 \(C'_\gamma=B^{-1}C_\gamma B\)，取逆便得到所写共轭公式。

	若有整体平行标架，则所有环路都保持该标架中的分量，故 \(P_\gamma=\operatorname{id}\)，从而 \(\rho\) 平凡。反之，若 \(\rho\) 平凡，则任意从 \(p\) 到 \(q\) 的两条路径 \(\gamma,\eta\) 满足
	\[
		P_\eta^{-1}P_\gamma=P_{\gamma*\bar\eta}=\operatorname{id}_{E_p},
	\]
	所以它们的运输相同。连通流形局部道路连通，因而道路连通：一个道路分支在坐标小盒中包含每个点的邻域，故各分支均开，连通性迫使只有一个分支。于是可以用 \(P_\gamma b\) 在每点定义标架，再按上一证明的局部公式验证光滑性与平行性。
\end{proof}

若直接令向量沿环路移动，则 \(v\cdot[\gamma]=P_\gamma v\) 是 \(E_p\) 上的右作用，因为 \((v\cdot[\gamma])\cdot[\eta]=v\cdot([\gamma][\eta])\)。本节对这组运输取逆矩阵，得到通常意义的左矩阵表示。选基只会共轭表示，所以“是否平凡”以及不动向量空间的维数不依赖所选基。

\begin{example}
	\label{g1:exm:flat-circle-holonomy}
	在 \(S^1\subset\mathbb R^2\) 上，记
	\[
		\vartheta=(x\,\mathrm{d}y-y\,\mathrm{d}x)|_{S^1}.
	\]
	局部取角度 \(\theta\) 时有 \(\vartheta=\mathrm{d}\theta\)；不同角度分支相差常数 \(2\pi k\)，所以这些局部微分确实给出一个整体一形式。符号 \(\mathrm{d}\theta\) 在这里不表示存在整体实值角度函数。

	在平凡实线丛 \(S^1\times\mathbb R\) 上取整体标架 \(e=1\)，并令
	\[
		\nabla(eu)=e(\mathrm{d}u+a\vartheta\,u),\qquad a\in\mathbb R.
	\]
	一维底空间上的二形式为零，故此联络平坦。沿正向一周路径 \(\gamma(t)=e^{2\pi i t}\)，平行方程为 \(u'+2\pi a u=0\)，于是
	\[
		P_\gamma=e^{-2\pi a},\qquad
		\rho(1)=e^{2\pi a}.
	\]
	这里用\cref{g1:thm:circle-fundamental}将 \(\pi_1(S^1)\) 按绕数识别为 \(\mathbb Z\)，故 \(\rho(k)=e^{2\pi ak}\)。当 \(a\ne0\) 时，没有整体平行标架，尽管底层线丛已经有整体普通标架 \(e\)。在每个角度区间内，\(e\,e^{-a\theta}\) 都是平行标架；绕一周后出现的因子 \(e^{-2\pi a}\) 正是其不能成为整体标架的障碍。
\end{example}

也可以从一个给定表示反向造出平坦丛。在万有覆盖上先取普通微分，再用表示识别不同覆盖层；这些识别在局部是常矩阵，因而能将各片的零联络粘合起来，并把指定的单值保留下来。

\begin{proposition}[由表示构造平坦丛]
	\label{g1:prop:flat-bundle-representation}
	设 \(M\) 是非空、连通、有限维、Hausdorff、第二可数光滑流形，允许有边界，固定 \(p\in M\)，令 \(G=\pi_1(M,p)\)。对任一表示 \(\rho:G\to\operatorname{GL}(r,\mathbb R)\)，存在秩 \(r\) 实光滑向量丛及平坦联络，其在一个指定基下的单值表示正好为 \(\rho\)。
\end{proposition}
\begin{proof}
	流形的坐标小盒或半盒可缩，因此局部道路连通且半局部单连通；连通性又给出道路连通性。\cref{g1:thm:path-class-universal-cover}于是提供道路类万有覆盖 \(\widetilde M\to M\)：点 \(q=[\alpha]\) 记录从 \(p\) 出发的路径类，投影为其终点。记 \(q_0\) 为常路径类。按照\cref{g1:prop:universal-deck-discontinuous}，群 \(G\) 的左覆盖变换作用是
	\[
		L_g[\alpha]=[\gamma*\alpha],\qquad g=[\gamma],
	\]
	并且在每根覆盖纤维上自由、传递。定义
	\[
		g\cdot(q,v)=(L_gq,\rho(g)v),\qquad
		E_\rho=(\widetilde M\times\mathbb R^r)/G.
	\]
	这是左作用，因为 \(L_gL_h=L_{gh}\) 且 \(\rho(g)\rho(h)=\rho(gh)\)。商投影 \([q,v]\mapsto p(q)\) 良定义。

	取连通的均匀覆盖坐标邻域 \(U\)，并在其上选一张覆盖片的局部逆 \(\sigma_U:U\to\widetilde M\)。映射
	\[
		\Phi_U:U\times\mathbb R^r\longrightarrow E_\rho|_U,\qquad
		(x,v)\longmapsto[\sigma_U(x),v]
	\]
	为双射：任一覆盖点都唯一写成 \(L_g\sigma_U(x)\)，而
	\[
		[L_g\sigma_U(x),v]=[\sigma_U(x),\rho(g)^{-1}v].
	\]
	它也是同胚。事实上，乘积中的轨道商映射是开映射，因为开集的轨道饱和集是其全部群平移的并；将商映射限制到一张覆盖片与 \(\mathbb R^r\) 的乘积上，便得到开连续双射。

	若另选 \(\sigma_V\)，在交叠附近有唯一 \(g(x)\in G\) 使 \(\sigma_V(x)=L_{g(x)}\sigma_U(x)\)，而 \(g(x)\) 局部常值。为看出这一点，在交叠点附近再选连通的均匀覆盖小邻域，两条局部逆都取值于各自固定的覆盖片；使它们在一点相同的覆盖变换，在这整个小邻域上都使它们相同。因此两套纤维坐标由局部常矩阵 \(\rho(g(x))^{-1}\) 过渡，尤其光滑且线性。这些图定义了光滑向量丛。不同底点可由底空间的 Hausdorff 性分离，同一纤维的不同点可在同一丛图内分离；底空间第二可数，可从这些平凡化域中选可数子覆盖，再取各 \(U\times\mathbb R^r\) 的可数基，得到总空间第二可数。

	在每张丛图中令联络矩阵为零。由于转移矩阵局部常值，换标架公式中的 \(\mathrm{d}g\) 为零，故这些局部联络相容并粘成整体联络，其曲率为零。以 \(v\mapsto[q_0,v]\) 识别基点纤维。环路 \(\gamma\) 的覆盖提升从 \(q_0\) 到 \(L_{[\gamma]}q_0\)，沿途保持纤维分量 \(v\)；因此
	\[
		P_\gamma[q_0,v]=[L_{[\gamma]}q_0,v]
		=[q_0,\rho([\gamma])^{-1}v].
	\]
	所以实际运输矩阵为 \(\rho([\gamma])^{-1}\)，按本节定义取逆后正好恢复指定表示 \(\rho\)。
\end{proof}

\subsection{局部曲率与整体平坦丛的区别}
\label{g1:subsec:1505:04}

一个向量丛连同指定的平坦联络，称为\term[pingtancong]{平坦丛}[flat bundle]。这一说法包含了联络数据：同一个光滑向量丛可以带不同联络，它们的曲率或单值未必相同。平坦丛的局部平行标架有局部常数转移矩阵，反过来，一套局部常数转移矩阵也能如上证明那样，将各片上的零联络粘成平坦联络。

\begin{example}
	\label{g1:exm:flat-mobius-line}
	考虑商空间
	\[
		L=(\mathbb R\times\mathbb R)/\bigl((t,v)\sim(t+2\pi,-v)\bigr)
		\longrightarrow S^1,\qquad
		[t,v]\longmapsto e^{it}.
	\]
	在长度小于 \(2\pi\) 的角度区间上选择角度分支，就得到实线丛图；分支相差 \(2\pi k\) 时，纤维坐标相差因子 \((-1)^k\)。因此这是光滑线丛，即 Möbius 实线丛。转移系数为常数 \(\pm1\)，所以局部取 \(\nabla=\mathrm{d}\) 给出整体平坦联络，且 \(v^2\) 给出与之相容的整体纤维度量。

	沿正向圆周一周，从 \([0,v]\) 开始保持 \(v\) 不变，终点是
	\[
		[2\pi,v]=[0,-v].
	\]
	故 \(P_\gamma=-1\)，并且 \(\rho(1)=-1\)。这里甚至无从通过更换普通整体标架来讨论单值，因为该线丛根本没有整体标架：若有处处非零连续截面，拉回到 \(\mathbb R\) 后可写成连续函数 \(u(t)\)，满足 \(u(t+2\pi)=-u(t)\)。于是 \(u(0)\) 与 \(u(2\pi)\) 符号相反，中值定理迫使它在中间为零，矛盾。
\end{example}

\cref{g1:exm:flat-circle-holonomy,g1:exm:flat-mobius-line}区分了两种不同现象。第一例底层线丛平凡，障碍在指定联络的环路运输；第二例底层线丛本身已经扭转。两例的曲率却都为零，因为它们在每一个足够小的区间内都与普通微分完全相同。曲率是局部量，不能单独区分这些整体情形。

\ProseParagraphBreak

平凡丛上的联络也可以有非零曲率：在\cref{g1:exm:path-dependent-line-transport}中，\(\mathbb R^2\times\mathbb R\) 的联络矩阵是 \(A=x\,\mathrm{d}y\)，其曲率为 \(F=\mathrm{d}x\wedge\mathrm{d}y\ne0\)。同一底层丛也可以选 \(A=0\)。因此“丛平凡”“联络平坦”“存在整体平行标架”是三个不同断言，其中最后一个同时消除了局部曲率与整体单值两个障碍。

\ProseParagraphBreak

本节的“局部实现”是把平坦联络实现成普通微分：在合适标架中，\(\nabla(fu)=f\,\mathrm{d}u\)。若丛恰为切丛，平行标架也未必是某组坐标的坐标向量场；还须考虑上一节的挠率以及标架之间的括号。一般向量丛的纤维向量甚至不属于底流形的切空间，因而不能把平行标架自动解释为底空间坐标。

% 纲要条目排为编号小节标题，并进入目录与书签；说明段仍保留于纲要注释中。


\begin{chaptersummary}
	联络对方向输入函数线性，对截面输入遵守 Leibniz 法则。行标架中 \(\nabla(eu)=e(\mathrm{d}u+Au)\)，换为 \(e'=eg\) 时有 \(A'=g^{-1}Ag+g^{-1}\mathrm{d}g\)。单位分解把局部联络合成整体联络；两个联络之差是 \(\operatorname{End}E\) 值一形式，故联络集合是非空仿射空间。配对与乘法法则进一步决定对偶、张量、外幂和同态上的诱导联络。

	\ProseParagraphBreak

	沿路径的平行方程为 \(u'+A(\gamma')u=0\)。线性方程在紧参数区间上的延拓、唯一性和参数依赖，使运输成为端点纤维之间的线性同构。先走 \(\gamma\) 后走 \(\eta\) 时，\(P_{\gamma*\eta}=P_\eta P_\gamma\)；反向路径给逆映射。先把邻点截面值运回固定纤维，再取差商，可以恢复联络。度量联络恰是使全部平行移动保持内积的联络。

	\ProseParagraphBreak

	曲率 \(R(X,Y)=\nabla_X\nabla_Y-\nabla_Y\nabla_X-\nabla_{[X,Y]}\) 对所有输入逐点线性，局部矩阵为 \(F=\mathrm{d}A+A\wedge A\)。外协变微分满足 \(\mathrm{D}^2\eta=R\wedge\eta\)，自同态值形式上的 Bianchi 恒等式为 \(\mathrm{D}_{\operatorname{End}}R=0\)。换标架使 \(F'=g^{-1}Fg\)；正向小方环的返回变换为 \(I-\varepsilon^2F_{12}(p)+O(\varepsilon^3)\)，给出曲率作为局部路径依赖的具体含义。

	\ProseParagraphBreak

	切丛的特殊结构使 \(T(X,Y)=\nabla_XY-\nabla_YX-[X,Y]\) 成为有意义的张量。局部余标架 \(\theta\) 与挠率列 \(\Theta\) 满足 \(\Theta=\mathrm{d}\theta+A\wedge\theta\)，再求外协变微分得到 \(\mathrm{D}\Theta=F\wedge\theta\)。无挠性在坐标标架中等价于联络系数两个下指标对称；任意平行标架能在内点附近成为坐标标架，当且仅当挠率为零。无挠、平坦和保持度量是不同条件。

	\ProseParagraphBreak

	曲率为零时，逐坐标方向解线性方程，可以依次消去全部联络系数，得到局部平行标架；平坦性保证已消去的方向不会重新出现。平行标架的过渡矩阵局部常值，使运输延伸到连续路径，并由同伦网格证明其只依赖定端点同伦类。单连通区域上存在由基点纤维基唯一决定的整体平行标架。一般连通底空间上，以环路运输的逆矩阵定义同态 \(\rho:\pi_1(M,p)\to\operatorname{GL}(r,\mathbb R)\)；整体平行标架存在当且仅当 \(\rho\) 平凡。平凡丛、平坦联络与整体平行标架因此必须分别判断。
\end{chaptersummary}

\BookExercises

\begin{exercise}\label{g1:ex:rotating-frame-connection}
	在 \(\mathbb R^2\) 上的平凡丛 \(E=\mathbb R^2\times\mathbb R^2\) 中，取标准标架 \(e=(e_1,e_2)\) 和逐分量联络 \(\nabla(eu)=e\,\mathrm{d}u\)。令 \(\theta:\mathbb R^2\to\mathbb R\) 光滑，取
	\[
		g=\begin{pmatrix}\cos\theta&-\sin\theta\\\sin\theta&\cos\theta\end{pmatrix},\qquad e'=eg.
	\]
	求新标架中的联络矩阵，并分别写出 \(\nabla e'_1,\nabla e'_2\)。对截面 \(s=xe_1+ye_2\)，求它在 \(e'\) 中的列分量，使用新联络矩阵计算 \(\nabla s\)，核对它与旧标架中的结果一致。
\end{exercise}
\begin{solution}
	令 \(J=\begin{pmatrix}0&-1\\1&0\end{pmatrix}\)。对 \(g\) 逐项求导，得到 \(\mathrm{d}g=gJ\,\mathrm{d}\theta\)，故
	\[
		A'=g^{-1}\mathrm{d}g=J\,\mathrm{d}\theta
		=\begin{pmatrix}0&-\mathrm{d}\theta\\\mathrm{d}\theta&0\end{pmatrix}.
	\]
	\(\nabla e'=e'A'\) 的第一、第二列分别给出
	\[
		\nabla e'_1=\mathrm{d}\theta\otimes e'_2,\qquad
		\nabla e'_2=-\mathrm{d}\theta\otimes e'_1.
	\]
	这些符号也符合几何意义：把 \(e'_1\) 的转角增加一点时，它沿 \(e'_2\) 方向变化。

	同一截面的新分量为
	\[
		u'=g^{-1}\begin{pmatrix}x\\y\end{pmatrix}
		=\begin{pmatrix}x\cos\theta+y\sin\theta\\-x\sin\theta+y\cos\theta\end{pmatrix}.
	\]
	逐项求导可得
	\begin{align*}
		\mathrm{d}u'^1&=\cos\theta\,\mathrm{d}x+\sin\theta\,\mathrm{d}y+u'^2\,\mathrm{d}\theta,\\
		\mathrm{d}u'^2&=-\sin\theta\,\mathrm{d}x+\cos\theta\,\mathrm{d}y-u'^1\,\mathrm{d}\theta.
	\end{align*}
	而 \(A'u'=(-u'^2\,\mathrm{d}\theta,u'^1\,\mathrm{d}\theta)^{\mathsf T}\)，恰好消掉这两行中来自旋转标架的附加项。于是
	\[
		\nabla s=e'(\mathrm{d}u'+A'u')
		=e'g^{-1}\begin{pmatrix}\mathrm{d}x\\\mathrm{d}y\end{pmatrix}
		=\mathrm{d}x\otimes e_1+\mathrm{d}y\otimes e_2.
	\]
	这与原标架中直接对 \((x,y)^{\mathsf T}\) 求导的结果相同。当 \(\mathrm{d}\theta\ne0\) 时，新矩阵虽然非零，联络本身却没有改变。
\end{solution}

\begin{exercise}\label{g1:ex:projection-induced-connection}
	设 \(M\) 为有限维、Hausdorff、第二可数光滑流形，允许有边界，\(P:M\to\operatorname{Mat}_{N\times N}(\mathbb R)\) 光滑，满足 \(P^2=P\) 且秩恒为 \(r\)。令 \(E_p=\operatorname{im}P(p)\)、\(F_p=\ker P(p)\)。
	\begin{enumerate}
		\item 说明 \(E,F\) 是平凡丛 \(M\times\mathbb R^N\) 的光滑子丛，且 \(E\oplus F=M\times\mathbb R^N\)。证明 \(\nabla^Es=P\,\mathrm{d}s\) 为 \(E\) 上的联络，\(\nabla^Ft=(I-P)\,\mathrm{d}t\) 为 \(F\) 上的联络，这里 \(\mathrm{d}\) 对环境空间分量逐项求导。
		\item 将两联络的直和看成平凡丛上的联络，证明它与标准联络的差为
		\[
			\nabla^\oplus=\mathrm{d}+(2P-I)\,\mathrm{d}P.
		\]
		\item 证明 \(P(\mathrm{d}P)P=0\)，并直接用自同态联络公式验证 \((\nabla^\oplus)^{\operatorname{End}}P=0\)。
	\end{enumerate}
\end{exercise}
\begin{solution}
	1. 固定 \(p\)，从矩阵 \(P(p)\) 的列中选出 \(r\) 个线性无关者。相应的非零子式在一个邻域内仍非零，故这些光滑列向量在那里独立；秩恒为 \(r\)，使它们恰好张成每个 \(E_q\)。这给出 \(E\) 的局部光滑标架。同样，\(I-P\) 的像等于 \(\ker P\)，秩为 \(N-r\)，对它的列作同一选择即得到 \(F\) 的局部标架。任意向量唯一分解为 \(v=Pv+(I-P)v\)，两部分分别属于 \(E,F\)，所以这两个子丛互补。

	对 \(E\) 的截面 \(s\)，\(P\,\mathrm{d}s\) 是取值于 \(E\) 的光滑一形式，并对 \(s\) 实线性。若 \(f\) 为光滑函数，由 \(Ps=s\)，
	\[
		P\,\mathrm{d}(fs)=P(\mathrm{d}f\otimes s+f\,\mathrm{d}s)
		=\mathrm{d}f\otimes s+fP\,\mathrm{d}s.
	\]
	因此它是联络。把 \(P\) 换成 \(I-P\) 得到 \(F\) 上的相同结论；此处无须假设投影正交。

	2. 平凡丛的任意截面 \(v\) 对应直和分解 \(Pv+(I-P)v\)，所以直和联络为
	\begin{align*}
		\nabla^\oplus v
		&=P\,\mathrm{d}(Pv)+(I-P)\,\mathrm{d}((I-P)v)\\
		&=P(\mathrm{d}P)v+P^2\mathrm{d}v-(I-P)(\mathrm{d}P)v+(I-P)^2\mathrm{d}v\\
		&=\mathrm{d}v+(2P-I)(\mathrm{d}P)v.
	\end{align*}
	故其相对于标准联络的差确为自同态值一形式 \(A=(2P-I)\mathrm{d}P\)，也与\cref{g1:prop:connection-affine-space}相符。

	3. 对 \(P^2=P\) 求导得到
	\[
		(\mathrm{d}P)P+P\,\mathrm{d}P=\mathrm{d}P.
	\]
	等式左右同时乘 \(P\)，便有 \(2P(\mathrm{d}P)P=P(\mathrm{d}P)P\)，所以 \(P(\mathrm{d}P)P=0\)。再计算
	\[
		AP=(2P-I)(\mathrm{d}P)P=-(\mathrm{d}P)P,\qquad
		PA=P(2P-I)\mathrm{d}P=P\,\mathrm{d}P.
	\]
	于是
	\[
		(\nabla^\oplus)^{\operatorname{End}}P
		=\mathrm{d}P+AP-PA
		=\mathrm{d}P-(\mathrm{d}P)P-P\,\mathrm{d}P=0.
	\]
	这也表达了直和联络的构造目的：求导与投影 \(P\) 相容，因此可以分别在两个子丛内求导。
\end{solution}

\begin{exercise}\label{g1:ex:endomorphism-connection-commutant}
	在 \(\mathbb R\) 上的平凡秩二实丛中，给定联络矩阵
	\[
		A=N\,\mathrm{d}x,\qquad N=\begin{pmatrix}0&1\\0&0\end{pmatrix}.
	\]
	求所有在标准标架中为常矩阵、且满足 \(\nabla^{\operatorname{End}}B=0\) 的自同态截面。它们之中有哪些是投影，即满足 \(B^2=B\)？再检验矩阵函数
	\[
		P(x)=\begin{pmatrix}1&x\\0&0\end{pmatrix}
	\]
	是否给出满足 \(\nabla^{\operatorname{End}}P=0\) 的秩一投影，并说明这与前面的常矩阵结论为何不矛盾。
\end{exercise}
\begin{solution}
	若 \(B=\begin{pmatrix}a&b\\c&d\end{pmatrix}\) 为常矩阵，\(\mathrm{d}B=0\)，所以\eqref{g1:eq:endomorphism-connection}给出
	\[
		\nabla^{\operatorname{End}}B=[N,B]\,\mathrm{d}x
		=\begin{pmatrix}c&d-a\\0&-c\end{pmatrix}\mathrm{d}x.
	\]
	因此所求条件为 \(c=0\)、\(d=a\)，全体解恰为
	\[
		B=aI+bN=\begin{pmatrix}a&b\\0&a\end{pmatrix},\qquad a,b\in\mathbb R.
	\]
	因 \(N^2=0\)，有 \(B^2=a^2I+2abN\)。投影条件变成 \(a^2=a\)、\((2a-1)b=0\)。第一式使 \(a=0\) 或 \(a=1\)，两种情形第二式都使 \(b=0\)，故其中只有零投影与恒等投影。

	对题给矩阵函数，直接相乘得 \(P(x)^2=P(x)\)，其第一列非零、第二行全零，故秩恒为 \(1\)。此外
	\[
		\mathrm{d}P=N\,\mathrm{d}x,\qquad [N,P]=-N,
	\]
	所以 \(\nabla^{\operatorname{End}}P=\mathrm{d}P+[N,P]\,\mathrm{d}x=0\)。这里普通导数与联络修正项恰好抵消。它不是常矩阵，不属于前半题所分类的集合。

	具体地，\(\operatorname{im}P\) 由 \(e_1\) 张成，\(\ker P\) 由 \(-xe_1+e_2\) 张成。两生成截面的协变导数都为零：\(Ne_1=0\)，而后一截面的普通导数为 \(-e_1\,\mathrm{d}x\)，联络项为 \(e_1\,\mathrm{d}x\)。所以两个子丛确实都被联络保持。常矩阵分类不能排除随基点变化的投影，也不能把“矩阵系数恒定”与“协变导数为零”混为一谈。
\end{solution}

\begin{exercise}\label{g1:ex:determinant-line-connection}
	设 \(M\) 为有限维、Hausdorff、第二可数光滑流形，允许有边界，\(E\to M\) 为秩 \(r\geq1\) 的实光滑向量丛，给定联络 \(\nabla\)。在局部标架 \(e=(e_1,\ldots,e_r)\) 中记联络矩阵为 \(A\)，并令 \(\delta=e_1\wedge\cdots\wedge e_r\)。
	\begin{enumerate}
		\item 证明行列式线丛 \(\det E=\Lambda^rE\) 的诱导联络满足 \(\nabla^{\det E}\delta=(\operatorname{tr}A)\otimes\delta\)。若 \(e'=eg\)，推导
		\[
			\operatorname{tr}A'=\operatorname{tr}A+\mathrm{d}\log|\det g|.
		\]
		\item 若给定处处非零的截面 \(\nu\in\Gamma(\det E)\)，证明存在 \(E\) 上的联络 \(\widehat\nabla\)，使 \(\widehat\nabla^{\det E}\nu=0\)。要求从任意已给联络 \(\nabla\) 出发，写出一个具体修正公式。
	\end{enumerate}
\end{exercise}
\begin{solution}
	1. 外幂联络沿各因子求导，故对任意方向 \(X\)，
	\[
		\nabla_X^{\det E}\delta
		=\sum_{j=1}^r\sum_{a=1}^r A^a{}_j(X)\,
		e_1\wedge\cdots\wedge e_a\wedge\cdots\wedge e_r,
	\]
	其中 \(e_a\) 位于原来第 \(j\) 个位置。若 \(a\ne j\)，楔积中出现两个 \(e_a\)，该项为零；只有 \(a=j\) 的项留下，系数为 \(\sum_jA^j{}_j(X)\)。所以 \(\nabla^{\det E}\delta=(\operatorname{tr}A)\otimes\delta\)。

	换标架使 \(\delta'=(\det g)\delta\)。把它视为线丛标架的变化，再用 Leibniz 规则，得到
	\[
		\nabla^{\det E}\delta'
		=\bigl(\mathrm{d}(\det g)+(\det g)\operatorname{tr}A\bigr)\otimes\delta
		=\left((\det g)^{-1}\mathrm{d}(\det g)+\operatorname{tr}A\right)\otimes\delta'.
	\]
	左侧在新标架中也等于 \((\operatorname{tr}A')\otimes\delta'\)，故
	\[
		\operatorname{tr}A'=\operatorname{tr}A+(\det g)^{-1}\mathrm{d}(\det g)
		=\operatorname{tr}A+\mathrm{d}\log|\det g|.
	\]
	由于 \(\det g\) 处处非零，它在每个交叠连通分支上的符号固定，\(\log|\det g|\) 光滑，最后的普通函数链式法则因此合法。

	2. 非零截面 \(\nu\) 是行列式线丛的全局标架，所以存在唯一光滑一形式 \(\theta\)，使
	\[
		\nabla^{\det E}\nu=\theta\otimes\nu.
	\]
	局部写 \(\nu=f\delta\)，其中 \(f\ne0\)，就有 \(\theta=f^{-1}\mathrm{d}f+\operatorname{tr}A\)；这个式子也直接说明 \(\theta\) 光滑。定义
	\[
		\widehat\nabla=\nabla-\frac1r\theta\,\operatorname{id}_E.
	\]
	被减去的项是整体的自同态值一形式，故\cref{g1:prop:connection-affine-space}保证 \(\widehat\nabla\) 仍是联络。它在 \(e\) 中的矩阵为 \(\widehat A=A-\tfrac1r\theta I_r\)，于是 \(\operatorname{tr}\widehat A=\operatorname{tr}A-\theta\)。对 \(\nu=f\delta\) 计算，得到
	\[
		\widehat\nabla^{\det E}\nu
		=\bigl(\mathrm{d}f+f(\operatorname{tr}A-\theta)\bigr)\otimes\delta=0.
	\]
	这给出了所需修正。条件 \(r\geq1\) 使除以 \(r\) 有意义，给定非零 \(\nu\) 则使整体一形式 \(\theta\) 能够定义；并未假设原来的联络已经保持 \(\nu\)。
\end{solution}

\begin{exercise}\label{g1:ex:transport-variation-of-constants}
	设 \(E\to M\) 为带联络的实有限秩光滑向量丛，\(\gamma:[a,b]\to M\) 光滑，\(f\) 为一个沿 \(\gamma\) 的光滑截面。给定 \(v\in E_{\gamma(a)}\)，证明非齐次初值问题
	\[
		\mathrm{D}_ts(t)=f(t),\qquad s(a)=v
	\]
	有唯一光滑解，且
	\[
		s(t)=P_{\gamma,t\leftarrow a}v+
		\int_a^t P_{\gamma,t\leftarrow\tau}f(\tau)\,\mathrm{d}\tau.
	\]
	说明积分在哪个向量空间中进行，为什么这个写法不依赖局部标架。
\end{exercise}
\begin{solution}
	先把所有向量运回固定纤维 \(E_{\gamma(a)}\)。若 \(s\) 为解，令 \(q(t)=P_{\gamma,a\leftarrow t}s(t)\)，则\eqref{g1:eq:transported-section-derivative}给出
	\[
		q'(t)=P_{\gamma,a\leftarrow t}f(t),\qquad q(a)=v.
	\]
	这是固定有限维向量空间中的普通微积分问题，所以唯一可能的解为
	\[
		q(t)=v+\int_a^tP_{\gamma,a\leftarrow\tau}f(\tau)\,\mathrm{d}\tau.
	\]
	被积向量总属于 \(E_{\gamma(a)}\)，可在任意固定基中逐分量积分。换一个固定基只是施加固定线性变换，而固定线性变换与积分相容，因此结果与该基无关。

	反过来用此式定义 \(q\)，再令 \(s(t)=P_{\gamma,t\leftarrow a}q(t)\)。运输的光滑时刻依赖保证 \(s\) 光滑。由同一求导公式，
	\[
		P_{\gamma,a\leftarrow t}\mathrm{D}_ts=q'(t)=P_{\gamma,a\leftarrow t}f(t).
	\]
	运输可逆，故 \(\mathrm{D}_ts=f\)，且初值为 \(v\)。将 \(P_{\gamma,t\leftarrow a}\) 乘入 \(q\) 的表达式，再用复合规律
	\[
		P_{\gamma,t\leftarrow a}P_{\gamma,a\leftarrow\tau}
		=P_{\gamma,t\leftarrow\tau},
	\]
	便得到题述公式。对每个固定 \(t\)，末式中的积分在固定纤维 \(E_{\gamma(t)}\) 内进行；它是上一个固定纤维积分经线性同构传送后的结果，所以同样不依赖标架。

	若两条解具有相同初值，反向运输后的差满足 \(q'=0\)、\(q(a)=0\)，故差恒为零。这也直接证明了唯一性。
\end{solution}

\begin{exercise}\label{g1:ex:parallel-invariant-subbundle}
	设 \(E\to M\) 为带联络 \(\nabla\) 的实有限秩光滑向量丛，\(F\subset E\) 为光滑子丛，\(M\) 允许有边界。证明下列条件等价：
	\begin{enumerate}
		\item 对每个局部向量场 \(X\) 及每个取值于 \(F\) 的局部光滑截面 \(S\)，\(\nabla_XS\) 也取值于 \(F\)。
		\item 对每条分片光滑路径 \(\gamma\)，有 \(P_\gamma(F_{\gamma(a)})=F_{\gamma(b)}\)。
	\end{enumerate}
\end{exercise}
\begin{solution}
	先假设第一项。取一个适应 \(F\) 的局部标架，使前 \(k=\operatorname{rank}F\) 个基向量张成 \(F\)，并把联络矩阵写成分块形式
	\[
		A=\begin{pmatrix}A_{11}&A_{12}\\A_{21}&A_{22}\end{pmatrix}.
	\]
	前 \(k\) 个基截面的协变导数取值于 \(F\)，恰表示 \(A_{21}=0\)。沿位于该标架域中的路径，把分量写成 \(u=(u_1,u_2)^{\mathsf T}\)，则平行方程的下半部分为
	\[
		u_2'=-A_{22}(\gamma')u_2.
	\]
	若初值位于 \(F\)，则 \(u_2(a)=0\)，唯一性使 \(u_2\equiv0\)。分割路径并逐段使用这个结论，得到 \(P_\gamma(F_{\gamma(a)})\subset F_{\gamma(b)}\)。再对反向路径作同样推理，结合 \(P_{\overline\gamma}=P_\gamma^{-1}\)，便得到相反包含，从而是等号。

	再假设第二项。固定 \(p\in M\)，取一个局部截面 \(S\) 使其取值于 \(F\)，并令 \(\gamma(0)=p\)。对足够小的允许时刻 \(t\)，有
	\[
		P_{\gamma,0\leftarrow t}S(\gamma(t))\in F_p,\qquad S(p)\in F_p.
	\]
	所以两者的差商属于同一个有限维线性子空间 \(F_p\subset E_p\)。这个子空间闭，故由\eqref{g1:eq:connection-recovery-limit}，极限 \((\nabla_{\gamma'(0)}S)_p\) 也在 \(F_p\) 中。内部点可用双侧曲线实现任意方向；边界点先用向内的单侧曲线实现对应方向，再对向外方向使用线性性。因此所有 \(\nabla_XS\) 都取值于 \(F\)，证明第一项。

	当 \(k=0\) 或 \(k=\operatorname{rank}E\) 时，相应的矩阵块为空，上述论证分别表示零子丛和整个丛都被保持，也包括在结论中。
\end{solution}

\begin{exercise}\label{g1:ex:metric-line-connection-unique}
	在平凡实线丛 \(E=M\times\mathbb R\) 上，给定光滑正函数 \(h:M\to(0,\infty)\)，以
	\[
		g_p(u,v)=h(p)uv
	\]
	定义丛度量。证明关于此度量的联络唯一，在标准标架中其联络一形式为 \(\tfrac12\mathrm{d}\log h\)。求出任意路径的平行移动，并证明另一个正函数 \(\widetilde h\) 给出同一度量联络，当且仅当 \(\widetilde h/h\) 在每个连通分支上为常数。
\end{exercise}
\begin{solution}
	在标准标架中，一个实线丛联络由实一形式 \(a\) 给出。此时度量矩阵只有一个条目 \(G=h\)，所以\eqref{g1:eq:metric-connection-matrix}成为
	\[
		\mathrm{d}h=2h\,a.
	\]
	因 \(h>0\)，唯一可能的形式为
	\[
		a=\frac{\mathrm{d}h}{2h}=\frac12\mathrm{d}\log h.
	\]
	它光滑，并确实满足度量相容条件，因此存在且唯一。具体地，对截面 \(u\) 和向量场 \(X\)，有 \(\nabla_Xu=Xu+\tfrac12X(\log h)\,u\)。

	沿 \(\gamma:[a,b]\to M\)，标量平行方程给出
	\[
		P_\gamma(v)
		=\exp\left(-\frac12\int_a^b\frac{\mathrm{d}}{\mathrm{d}t}\log h(\gamma(t))\,\mathrm{d}t\right)v
		=\sqrt{\frac{h(\gamma(a))}{h(\gamma(b))}}\,v.
	\]
	对分片光滑路径逐段求积，内部分点的因子相消，仍是同一公式。标准分量可以改变，但长度保持，因为
	\[
		h(\gamma(b))\bigl(P_\gamma(v)\bigr)^2=h(\gamma(a))v^2.
	\]
	这个例子说明非正交归一标架中的非零联络系数，并不意味着运输改变了给定度量的长度。

	最后，两个度量联络相同当且仅当
	\[
		\mathrm{d}\log\widetilde h=\mathrm{d}\log h,
		\quad\text{即}\quad
		\mathrm{d}\log(\widetilde h/h)=0.
	\]
	由\cref{g1:cor:local-exactness}关于零次形式的结论，导数为零的光滑函数在每个连通分支上为常数；指数函数把这一结论转为 \(\widetilde h/h\) 在每个连通分支上为正常数。反向若这些比值为常数，两个对数的微分相同，故联络相同。
\end{solution}

\begin{exercise}\label{g1:ex:noncommuting-path-transport}
	在 \(\mathbb R^2\times\mathbb R^2\) 的标准标架中，取联络矩阵
	\[
		A=B\,\mathrm{d}x+C\,\mathrm{d}y,\qquad
		B=\begin{pmatrix}0&1\\0&0\end{pmatrix},\qquad
		C=\begin{pmatrix}0&0\\1&0\end{pmatrix}.
	\]
	令 \(\gamma\) 从 \((0,0)\) 先到 \((1,0)\) 再到 \((1,1)\)，\(\eta\) 从 \((0,0)\) 先到 \((0,1)\) 再到 \((1,1)\)，每段均为直线。计算 \(P_\gamma\)、\(P_\eta\) 的矩阵，证明两条路径的矩阵值积分 \(\int_\gamma A\)、\(\int_\eta A\) 相同，运输却不同。由此说明一般运输不能仅凭 \(\exp(-\int_\gamma A)\) 给出。
\end{exercise}
\begin{solution}
	沿正向单位水平线段，以 \(x\) 为参数，方程为 \(u'=-Bu\)。因为 \(B^2=0\)，矩阵 \(I-xB\) 的导数为 \(-B\)，且
	\[
		-B(I-xB)=-B,
	\]
	所以长度为一的水平段运输矩阵为 \(I-B\)。同理，由 \(C^2=0\)，正向单位竖直段的运输矩阵为 \(I-C\)。这里直接检验方程即可，不需要借用一般矩阵指数公式。

	拼接时后走的路径所对应矩阵乘在左边。因此
	\[
		P_\gamma=(I-C)(I-B)
		=\begin{pmatrix}1&-1\\-1&2\end{pmatrix},\qquad
		P_\eta=(I-B)(I-C)
		=\begin{pmatrix}2&-1\\-1&1\end{pmatrix}.
	\]
	例如它们把第一基向量分别送到 \((1,-1)^{\mathsf T}\)、\((2,-1)^{\mathsf T}\)，故确实是不同运输。差别来自
	\[
		CB=\begin{pmatrix}0&0\\0&1\end{pmatrix},\qquad
		BC=\begin{pmatrix}1&0\\0&0\end{pmatrix}.
	\]

	另一方面，每条路径都有一个正向单位水平段和一个正向单位竖直段，逐条目积分均得
	\[
		\int_\gamma A=B+C=\int_\eta A.
	\]
	所以只依赖总积分矩阵的表达式，特别是 \(\exp(-\int A)\)，会给这两条路径相同的结果，与上述解方程所得运输不符。把不同时刻的矩阵直接相加再取指数，会遗漏实际运动中不可交换的先后次序。
\end{solution}

\begin{exercise}\label{g1:ex:bianchi-nonclosed-matrix}
	在 \(\mathbb R^3\) 上的平凡秩二实向量丛中，取
	\[
		B=\begin{pmatrix}0&1\\0&0\end{pmatrix},\qquad
		C=\begin{pmatrix}0&0\\1&0\end{pmatrix},\qquad
		A=xB\,\mathrm{d}y+C\,\mathrm{d}z.
	\]
	计算 \(F\)、\(\mathrm{d}F\) 及 \(A\wedge F-F\wedge A\)，直接核对 Bianchi 恒等式。说明为什么不能把 Bianchi 恒等式简写成 \(\mathrm{d}F=0\)。
\end{exercise}
\begin{solution}
	令 \(H=[B,C]=\operatorname{diag}(1,-1)\)。先有 \(\mathrm{d}A=B\,\mathrm{d}x\wedge\mathrm{d}y\)，而 \(A\wedge A\) 中只有交叉项留下，得到
	\[
		F=B\,\mathrm{d}x\wedge\mathrm{d}y+xH\,\mathrm{d}y\wedge\mathrm{d}z,
		\qquad\mathrm{d}F=H\,\mathrm{d}x\wedge\mathrm{d}y\wedge\mathrm{d}z\ne0.
	\]
	含 \(xB\,\mathrm{d}y\) 的乘积均因 \(\mathrm{d}y\) 重复而消失，\(C\,\mathrm{d}z\) 与曲率第二项的乘积也为零。因此
	\begin{align*}
		A\wedge F-F\wedge A
		&=CB\,\mathrm{d}z\wedge\mathrm{d}x\wedge\mathrm{d}y-BC\,\mathrm{d}x\wedge\mathrm{d}y\wedge\mathrm{d}z\\
		&=(CB-BC)\,\mathrm{d}x\wedge\mathrm{d}y\wedge\mathrm{d}z\\
		&=-H\,\mathrm{d}x\wedge\mathrm{d}y\wedge\mathrm{d}z.
	\end{align*}
	它与 \(\mathrm{d}F\) 恰好抵消。曲率取值于随点变化的纤维自同态空间，协变微分中的两个联络项负责比较这些自同态；普通外微分不包含这种比较。在秩一情形矩阵交换子项消失，才可以在局部写成 \(\mathrm{d}F=0\)。
\end{solution}

\begin{exercise}\label{g1:ex:curvature-connection-difference}
	设 \(\nabla' =\nabla+B\)，其中 \(B\in\Omega^1(M;\operatorname{End}(E))\)。证明
	\[
		R^{\nabla'}=R^\nabla+\mathrm{D}_{\operatorname{End}}B+B\wedge B.
	\]
	再在 \(\mathbb R^2\) 上的平凡秩二丛中取两个常矩阵 \(P,Q\)，满足 \([P,Q]\ne0\)。说明矩阵分别为 \(P\,\mathrm{d}x\) 与 \(Q\,\mathrm{d}y\) 的联络都平坦，而它们的非平凡凸组合未必平坦。
\end{exercise}
\begin{solution}
	在同一局部标架中，设两个联络的矩阵分别为 \(A\) 与 \(A+B\)。逐项展开给出
	\begin{align*}
		F'&=\mathrm{d}(A+B)+(A+B)\wedge(A+B)\\
		&=F+\mathrm{d}B+A\wedge B+B\wedge A+B\wedge B.
	\end{align*}
	由于 \(B\) 是一次形式，\eqref{g1:eq:endomorphism-exterior-derivative}中的符号为 \(-(-1)^1=+1\)，故中间三项是 \(\mathrm{D}_{\operatorname{End}}B\)。各项都是整体 \(\operatorname{End}(E)\) 值二形式，局部等式因此给出所需整体公式。也可在 \(X,Y\) 上求值得到最后一项 \(B(X)B(Y)-B(Y)B(X)\)。

	两个给定矩阵的系数均常值，且各自只有一个坐标微分，故曲率都为零。参数 \(0<t<1\) 的凸组合为
	\[
		A_t=(1-t)P\,\mathrm{d}x+tQ\,\mathrm{d}y.
	\]
	于是
	\[
		F_t=t(1-t)[P,Q]\,\mathrm{d}x\wedge\mathrm{d}y\ne0.
	\]
	例如可取上一题的 \(P=B,Q=C\)。联络构成仿射空间，平坦条件却包含二次项；这正是平坦联络不必构成仿射子空间的原因。
\end{solution}

\begin{exercise}\label{g1:ex:determinant-connection-curvature}
	设 \(E\) 为秩 \(r\ge1\) 的实光滑向量丛，底流形有限维、Hausdorff、第二可数，允许边界，给定联络 \(\nabla\)。\cref{g1:ex:determinant-line-connection}已算出行列式线丛 \(\det E=\Lambda^rE\) 的诱导联络一形式为 \(\operatorname{tr}A\)。证明其曲率二形式为 \(\operatorname{tr}F\)，并说明它是全局闭形式。再通过选取一个参考度量联络，证明对实向量丛而言 \(\operatorname{tr}F\) 实际上总是恰当形式；不要求 \(E\) 或 \(\det E\) 平凡。
\end{exercise}
\begin{solution}
	行列式线丛的联络一形式为 \(\operatorname{tr}A\)。标量一形式与自身楔积为零，因而其曲率为 \(\mathrm{d}\operatorname{tr}A\)。另一方面，
	\[
		\operatorname{tr}(A\wedge A)=\sum_{i,j}A^i{}_j\wedge A^j{}_i=0,
	\]
	因为对角项为零，而 \(i\ne j\) 的项成对反号。所以 \(\operatorname{tr}F=\mathrm{d}\operatorname{tr}A\)，恰是行列式线丛的曲率。

	换标架时 \(F'=g^{-1}Fg\)。在每点对两个切向量求值后，这是通常矩阵的共轭，故 \(\operatorname{tr}F'=\operatorname{tr}F\)，局部形式拼成全局二形式。在每个标架邻域，\(\mathrm{d}\operatorname{tr}F=\mathrm{d}^2\operatorname{tr}A=0\)，因此它全局闭。至此只有局部原形式，不能直接把各片的 \(\operatorname{tr}A\) 拼在一起。

	为得到整体原形式，由\cref{g1:thm:bundle-metric-existence,g1:prop:metric-connection-existence}选择正定丛度量及其度量联络 \(\nabla^0\)。在正交规范标架中，\(A^0\) 为反对称实矩阵一形式，所以 \(\operatorname{tr}A^0=0\)，从而 \(\operatorname{tr}F^0=\mathrm{d}\operatorname{tr}A^0=0\)。这又是标架无关的结论。

	联络之差 \(B=\nabla-\nabla^0\) 为整体自同态值一形式，故 \(\beta=\operatorname{tr}B\) 是全局实一形式。在任意共同标架中，它写成 \(\operatorname{tr}(A-A^0)\)，因此
	\[
		\mathrm{d}\beta=\mathrm{d}\operatorname{tr}A-\mathrm{d}\operatorname{tr}A^0
		=\operatorname{tr}F-\operatorname{tr}F^0=\operatorname{tr}F.
	\]
	这给出了实际的整体原形式。我们只要求参考联络存在，并没有要求存在整体正交规范标架；正交标架只用于每片上验证 \(\operatorname{tr}F^0=0\)。
\end{solution}

\begin{exercise}\label{g1:ex:metric-curvature-skew}
	设实向量丛 \(E\) 配有正定度量 \(g\) 和度量联络 \(\nabla\)。证明每个 \(R(X,Y)\) 关于 \(g\) 反对称，即
	\[
		g(R(X,Y)s,t)+g(s,R(X,Y)t)=0.
	\]
	由此说明实线丛的每个度量联络都平坦，并在平凡秩二丛上给出保持标准度量而曲率非零的联络。
\end{exercise}
\begin{solution}
	局部取光滑正交规范标架。\eqref{g1:eq:metric-connection-matrix}在此成为 \(A(X)^{\mathsf T}=-A(X)\)。直接用二形式求值公式，
	\[
		F(X,Y)=X(A(Y))-Y(A(X))-A([X,Y])+[A(X),A(Y)].
	\]
	前三项仍为反对称矩阵；若 \(P,Q\) 反对称，则
	\[
		[P,Q]^{\mathsf T}=Q^{\mathsf T}P^{\mathsf T}-P^{\mathsf T}Q^{\mathsf T}
		=QP-PQ=-[P,Q].
	\]
	因此 \(F(X,Y)\) 反对称，翻译回纤维内积就是所需恒等式。该等式与标架无关，故整体成立。

	一维实内积空间的反对称自同态只有零，故实线丛的度量联络曲率为零。这是局部曲率的结论，不能据此断言每个环路的运输都为恒等；\secref{g1:sec:1505}中的 Möbius 线丛会显示区别。秩二时取
	\[
		J=\begin{pmatrix}0&-1\\1&0\end{pmatrix},\qquad A=xJ\,\mathrm{d}y
	\]
	作为 \(\mathbb R^2\) 上平凡丛标准标架中的矩阵。由于 \(A\) 反对称，它定义度量联络；但 \(F=J\,\mathrm{d}x\wedge\mathrm{d}y\ne0\)。保持长度并不排除旋转型的曲率。
\end{solution}

\begin{exercise}\label{g1:ex:heisenberg-parallel-frame-torsion}
	在 \(\mathbb R^3\) 上取全局标架
	\[
		e_1=\partial_x,\qquad e_2=\partial_y+x\partial_z,\qquad e_3=\partial_z,
	\]
	并规定 \(\nabla_Xe_i=0\)。求其对偶余标架 \(\theta\)、联络矩阵 \(A\)、挠率二形式列 \(\Theta\) 和曲率矩阵 \(F\)，直接验证第一 Bianchi 恒等式。判断这组标架是否可能在某个非空开集上成为坐标标架，以及是否可能有其他坐标使该联络的全部坐标系数为零。
\end{exercise}
\begin{solution}
	三个向量场在标准坐标标架中的列矩阵为
	\[
		\begin{pmatrix}1&0&0\\0&1&0\\0&x&1\end{pmatrix},
	\]
	行列式为一，故确实处处构成标架。逐个求值可得其对偶余标架
	\[
		\theta^1=\mathrm{d}x,\qquad
		\theta^2=\mathrm{d}y,\qquad
		\theta^3=\mathrm{d}z-x\,\mathrm{d}y.
	\]
	例如 \(\theta^3(e_2)=x-x=0\)，而 \(\theta^3(e_3)=1\)。规定每个 \(e_i\) 平行，具体就是对 \(Y=\sum_i u^i e_i\) 定义 \(\nabla_XY=\sum_iX(u^i)e_i\)，所以 \(A=0\)，并且
	\[
		F=\mathrm{d}A+A\wedge A=0,\qquad
		\Theta=\mathrm{d}\theta
		=\begin{pmatrix}0\\0\\-\mathrm{d}x\wedge\mathrm{d}y\end{pmatrix}.
	\]
	因此
	\[
		T(e_1,e_2)=-e_3,\qquad T(e_1,e_3)=T(e_2,e_3)=0.
	\]
	直接从 Lie 括号也可核对：\([e_1,e_2]=\partial_z=e_3\)，其余两个括号为零；由于标架平行，挠率就是这些括号的负值。

	在当前标架中，第一 Bianchi 恒等式的左端为
	\[
		\mathrm{d}\Theta+A\wedge\Theta=\mathrm{d}\Theta=0,
	\]
	因为 \(\mathrm{d}(\mathrm{d}x\wedge\mathrm{d}y)=0\)；右端 \(F\wedge\theta\) 也为零。这个三维例子中，二形式 \(\Theta^3\) 非零而其外微分为零，恒等式约束的并不是挠率本身必须消失。

	坐标向量场必须两两交换，而 \([e_1,e_2]=e_3\) 在每一点都非零，所以给定标架在任何非空开集上都不是坐标标架。也不存在另一组坐标使联络矩阵为零：若有，该坐标标架既平行又两两交换，挠率在该标架上便全为零。挠率的张量性将推出同一开集上 \(T=0\)，与刚算出的 \(T(e_1,e_2)=-e_3\) 矛盾。
\end{solution}

\begin{exercise}\label{g1:ex:torsion-correction-breaks-metric}
	在 \(\mathbb R^2\) 上令
	\[
		e_1=\partial_x,\qquad e_2=e^x\partial_y,\qquad
		g=\mathrm{d}x^2+e^{-2x}\mathrm{d}y^2,
	\]
	并规定 \(\nabla_Xe_i=0\)。验证 \(\nabla\) 保持 \(g\)，求其挠率 \(T\)。令
	\[
		\widehat\nabla_XY=\nabla_XY-\frac12T(X,Y).
	\]
	计算四个 \(\widehat\nabla_{e_i}e_j\)，验证 \(\widehat\nabla\) 无挠，并用一个明确的分量说明它不保持 \(g\)。
\end{exercise}
\begin{solution}
	因为 \(g(e_i,e_j)=\delta_{ij}\)，标架 \(e_1,e_2\) 正交归一。在该标架中 \(\nabla\) 的矩阵为零，所以\eqref{g1:eq:metric-connection-matrix}给出 \(\nabla g=0\)。也可以对 \(Y=u^1e_1+u^2e_2\)、\(Z=v^1e_1+v^2e_2\) 直接使用普通乘积法则，得到
	\[
		Xg(Y,Z)=X\left(\sum_i u^iv^i\right)
		=g(\nabla_XY,Z)+g(Y,\nabla_XZ).
	\]

	由 \([e_1,e_2]=e_2\) 及标架平行性，
	\[
		T(e_1,e_2)=-e_2,\qquad T(e_2,e_1)=e_2,\qquad
		T(e_1,e_1)=T(e_2,e_2)=0.
	\]
	代入修正公式便得到
	\[
		\widehat\nabla_{e_1}e_1=0,\qquad
		\widehat\nabla_{e_1}e_2=\frac12e_2,\qquad
		\widehat\nabla_{e_2}e_1=-\frac12e_2,\qquad
		\widehat\nabla_{e_2}e_2=0.
	\]
	因此
	\[
		T^{\widehat\nabla}(e_1,e_2)
		=\frac12e_2-\left(-\frac12e_2\right)-[e_1,e_2]=0.
	\]
	挠率是反对称张量，两个标架向量张成每个切空间，故这已经证明 \(T^{\widehat\nabla}=0\)。

	但度量的协变导数满足
	\[
		(\widehat\nabla_{e_1}g)(e_2,e_2)
		=e_1\bigl(g(e_2,e_2)\bigr)
		-2g(\widehat\nabla_{e_1}e_2,e_2)
		=0-2\cdot\frac12=-1.
	\]
	所以 \(\widehat\nabla g\ne0\)。原标架中的 \(g(e_2,e_2)\) 恒为一，而修正后的导数让 \(e_2\) 沿 \(e_1\) 方向出现自身的非零分量；这正是度量乘积法则失效的具体位置。
\end{solution}

\begin{exercise}\label{g1:ex:autoparallel-curves-torsion-invisible}
	设 \(M\) 为无边界光滑流形，\(\nabla,\widetilde\nabla\) 为两个切丛联络，写 \(\widetilde\nabla_XY=\nabla_XY+B(X)Y\)。本题称满足
	\[
		\mathrm{D}_t^\nabla\gamma'(t)=0
	\]
	的光滑曲线为 \(\nabla\) 的带参数自平行曲线；比较两条曲线时保留原来的参数，不允许另行重参数化。证明以下条件等价：
	\begin{enumerate}
		\item 两个联络具有完全相同的带参数自平行曲线。
		\item 对每个 \(p\in M\)、\(v\in \mathrm{T}_pM\)，有 \(B_p(v)v=0\)。
		\item 对任意 \(X,Y\)，有 \(B(X)Y=-B(Y)X\)。
	\end{enumerate}
	证明中应说明每个内点初值 \((p,v)\) 都有局部自平行曲线。由此解释，联络与其无挠修正 \(\nabla^0_XY=\nabla_XY-\tfrac12T^\nabla(X,Y)\) 为什么具有相同的带参数自平行曲线。
\end{exercise}
\begin{solution}
	先说明初值存在性。在坐标域中写 \(\nabla_{\partial_i}\partial_j=\sum_k\Gamma^k_{ij}\partial_k\)，令曲线的坐标为 \(x^k(t)\)、速度分量为 \(v^k(t)\)。沿曲线导数公式把自平行方程写成
	\[
		\frac{\mathrm{d}x^k}{\mathrm{d}t}=v^k,\qquad
		\frac{\mathrm{d}v^k}{\mathrm{d}t}
		=-\sum_{i,j}\Gamma^k_{ij}(x)v^iv^j.
	\]
	右端是坐标开集与 \(\mathbb R^n\) 上的光滑函数。\cref{g1:thm:ode-local-existence-uniqueness}因此对任意 \((x(p),v)\) 给出唯一局部解，缩短时间后其位置留在原坐标域。这里没有边界，故任意切向初速度都可以实现为双侧局部曲线。

	两个沿曲线导数的差只取决于速度与当前向量，所以
	\[
		\mathrm{D}_t^{\widetilde\nabla}\gamma'
		=\mathrm{D}_t^\nabla\gamma'+B_{\gamma(t)}(\gamma'(t))\gamma'(t).
	\]
	若第一项成立，取任意初值 \((p,v)\) 的一条 \(\nabla\) 自平行曲线；它也是 \(\widetilde\nabla\) 自平行曲线。在初始时刻代入上式，得到 \(B_p(v)v=0\)，即第二项。反过来，若第二项成立，上式中的差对每条光滑曲线都为零，所以一条曲线满足一个自平行方程当且仅当满足另一个，证明第一项。

	第二项蕴含第三项可逐点作极化。对同一纤维中的 \(u,v\)，双线性给出
	\[
		0=B(u+v)(u+v)-B(u)u-B(v)v=B(u)v+B(v)u.
	\]
	第三项反过来令 \(u=v\)，得到 \(2B(v)v=0\)，在实数域上便是第二项。

	最后，修正项 \(B(X)Y=-\tfrac12T^\nabla(X,Y)\) 对 \(X,Y\) 反对称，故满足第三项。于是 \(\nabla\) 与 \(\nabla^0\) 的带参数自平行曲线相同，尽管两者的挠率可能不同。更一般地，添加任意反对称的 \(\mathrm{T}M\) 值二阶张量都不改变这些曲线，却使挠率增加 \(2B\)。因此仅由这族曲线不能恢复联络的挠率信息。
\end{solution}

\begin{exercise}\label{g1:ex:metric-connection-determined-by-torsion}
	设 \(g\) 为实切丛 \(\mathrm{T}M\) 上的光滑正定度量，\(\nabla,\widetilde\nabla\) 都保持 \(g\)，并令
	\[
		B(X)Y=\widetilde\nabla_XY-\nabla_XY,\qquad
		\tau(X,Y,Z)=g\bigl(T^{\widetilde\nabla}(X,Y)-T^\nabla(X,Y),Z\bigr).
	\]
	证明
	\[
		2g(B(X)Y,Z)
		=\tau(X,Y,Z)-\tau(Y,Z,X)+\tau(Z,X,Y).
	\]
	由此证明两个保持同一度量的联络若具有相同挠率，则必相同。说明为什么这一结论不能去掉“保持同一度量”的假设。
\end{exercise}
\begin{solution}
	记 \(b(X,Y,Z)=g(B(X)Y,Z)\)。两个联络的度量相容公式作差，函数导数项相消，得到
	\[
		g(B(X)Y,Z)+g(Y,B(X)Z)=0,
	\]
	也就是 \(b(X,Y,Z)=-b(X,Z,Y)\)。因此 \(b\) 的后两个位置反对称；这里没有假定前两个位置具有任何对称性。

	挠率差公式给出
	\begin{align*}
		\tau(X,Y,Z)&=g(B(X)Y-B(Y)X,Z)\\
		&=b(X,Y,Z)-b(Y,X,Z)\\
		&=b(X,Y,Z)+b(Y,Z,X).
	\end{align*}
	循环替换三个输入，还得到
	\[
		\begin{aligned}
			\tau(Y,Z,X)&=b(Y,Z,X)+b(Z,X,Y),\\
			\tau(Z,X,Y)&=b(Z,X,Y)+b(X,Y,Z).
		\end{aligned}
	\]
	将第一条等式减去第二条、再加第三条，两个中间项相消，结果正是
	\[
		\tau(X,Y,Z)-\tau(Y,Z,X)+\tau(Z,X,Y)=2b(X,Y,Z).
	\]
	这证明了所求公式。若两个挠率相同，则 \(\tau=0\)，故 \(g(B(X)Y,Z)=0\) 对所有 \(Z\) 成立。内积非退化，于是 \(B(X)Y=0\) 对所有 \(X,Y\) 成立，即 \(\widetilde\nabla=\nabla\)。

	若不要求保持同一度量，联络差可以含有非零对称部分而不改变挠率。例如在 \(\mathbb R\) 的标准坐标中，取
	\[
		\nabla_{\partial_x}\partial_x=0,\qquad
		\widetilde\nabla_{\partial_x}\partial_x=\partial_x.
	\]
	二者不同，但一维空间上所有反对称二阶张量均为零，所以两者挠率都为零。关于同一标准度量 \(\mathrm{d}x^2\)，第一个联络相容，第二个却有
	\[
		(\widetilde\nabla_{\partial_x}g)(\partial_x,\partial_x)=-2.
	\]
	这说明固定度量并要求两个联络都与它相容，正是上述唯一性论证中后两个输入反对称的来源。
\end{solution}

% 第15.5节配套习题与解答；供章末习题区整合。

\begin{exercise}\label{g1:ex:flat-matrix-local-model}
	在平凡丛 \(\mathbb R^2\times\mathbb R^2\) 的标准行标架 \(e\) 中，令
	\[
		A=
		\begin{pmatrix}1&2y\\0&-1\end{pmatrix}\mathrm{d}x+
		\begin{pmatrix}0&1\\0&0\end{pmatrix}\mathrm{d}y,\qquad
		g(x,y)=
		\begin{pmatrix}e^x&ye^x\\0&e^{-x}\end{pmatrix}.
	\]
	直接计算 \(F=\mathrm{d}A+A\wedge A\)，并验证 \(A=g^{-1}\mathrm{d}g\)。求一个整体平行标架以及从 \(p\) 到 \(q\) 的平行移动矩阵。说明为什么本例不能通过假定 \(A_x,A_y\) 交换来解释平坦性。
\end{exercise}
\begin{solution}
	令 \(H(y)=\begin{psmallmatrix}1&2y\\0&-1\end{psmallmatrix}\)、\(N=\begin{psmallmatrix}0&1\\0&0\end{psmallmatrix}\)。则
	\[
		\mathrm{d}A=(\partial_xA_y-\partial_yA_x)\,\mathrm{d}x\wedge\mathrm{d}y
		=-2N\,\mathrm{d}x\wedge\mathrm{d}y.
	\]
	另一方面，\(HN=N\)、\(NH=-N\)，所以
	\[
		A\wedge A=[H,N]\,\mathrm{d}x\wedge\mathrm{d}y
		=2N\,\mathrm{d}x\wedge\mathrm{d}y.
	\]
	两项恰好抵消，故 \(F=0\)，而 \([A_x,A_y]=2N\ne0\)。

	为计算 \(g^{-1}\mathrm{d}g\)，写
	\[
		g=D(x)U(y),\qquad
		D(x)=\begin{pmatrix}e^x&0\\0&e^{-x}\end{pmatrix},\qquad
		U(y)=\begin{pmatrix}1&y\\0&1\end{pmatrix}.
	\]
	乘积法则给出
	\[
		g^{-1}\mathrm{d}g
		=U^{-1}\begin{pmatrix}1&0\\0&-1\end{pmatrix}U\,\mathrm{d}x
		+U^{-1}\begin{pmatrix}0&1\\0&0\end{pmatrix}\mathrm{d}y
		=H(y)\,\mathrm{d}x+N\,\mathrm{d}y=A.
	\]
	因此换标架矩阵应取 \(g^{-1}\)，而不是 \(g\)。事实上，\(\mathrm{d}(g^{-1})=-g^{-1}(\mathrm{d}g)g^{-1}\)，所以标架 \(f=e g^{-1}\) 中的联络矩阵为
	\[
		gAg^{-1}+g\,\mathrm{d}(g^{-1})
		=(\mathrm{d}g)g^{-1}-(\mathrm{d}g)g^{-1}=0.
	\]
	它就是整体平行标架。

	也可以直接看平行方程。若 \(u\) 是标准标架分量，则沿路径有
	\[
		\frac{\mathrm{d}}{\mathrm{d}t}\bigl(g(\gamma(t))u(t)\bigr)
		=\bigl(\mathrm{d}g(\gamma')-gA(\gamma')\bigr)u=0.
	\]
	故从 \(p\) 到 \(q\) 的运输矩阵为 \(g(q)^{-1}g(p)\)，仅依赖端点。本例的平坦性来自系数导数项与矩阵交换子项的抵消，不能把非交换项直接略去。
\end{solution}

\begin{exercise}\label{g1:ex:flat-periodic-line-gauge}
	设 \(b:\mathbb R\to\mathbb R\) 光滑且 \(2\pi\) 周期。在 \(S^1\times\mathbb R\) 的标准标架 \(e\) 中定义联络 \(A=b(\theta)\,\mathrm{d}\theta\)，其中角度分支按\cref{g1:exm:flat-circle-holonomy}理解。令
	\[
		\overline b=\frac1{2\pi}\int_0^{2\pi}b(t)\,\mathrm{d}t.
	\]
	计算正向一周的平行移动，并证明存在整体平行标架当且仅当 \(\overline b=0\)。进一步构造正的周期函数 \(g\)，使标架 \(eg\) 中的联络矩阵为 \(\overline b\,\mathrm{d}\theta\)；说明周期换标架为何不能改变这个平均值。
\end{exercise}
\begin{solution}
	把一周路径提升到角度区间 \([0,2\pi]\)，平行方程为 \(u'+bu=0\)，解是
	\[
		u(\theta)=\exp\left(-\int_0^\theta b(t)\,\mathrm{d}t\right)u(0).
	\]
	所以正向一周的运输乘以 \(e^{-2\pi\overline b}\)。实指数函数只有在指数为零时等于 \(1\)，故由\cref{g1:prop:flat-monodromy}，整体平行标架存在当且仅当 \(\overline b=0\)。也可直接从解式看出：非零初值产生处处非零周期解，当且仅当这个因子为 \(1\)。

	令
	\[
		g(\theta)=\exp\left(-\int_0^\theta(b(t)-\overline b)\,\mathrm{d}t\right).
	\]
	积分中的函数周期且一周期积分为零，所以其原函数在相差 \(2\pi\) 的两点取值相同；因此 \(g\) 是正的光滑周期函数，定义了圆周上的合法换标架。由秩一变换公式，
	\[
		A'=A+g^{-1}\mathrm{d}g
		=\left(b+\frac{g'}g\right)\mathrm{d}\theta
		=\overline b\,\mathrm{d}\theta.
	\]
	任意其他合法换标架函数 \(h:S^1\to\mathbb R^\times\) 在连通圆周上符号恒定，因此 \(\log|h|\) 是整体光滑函数，且 \(h^{-1}\mathrm{d}h=\mathrm{d}\log|h|\)。沿一周积分为零，故变换前后的联络系数平均值相同。当 \(\overline b\ne0\) 时，局部仍能消去联络系数，但消去全部系数所需的指数函数不能满足周期条件。
\end{solution}

\begin{exercise}\label{g1:ex:flat-parallel-section-invariants}
	设 \(E\to M\) 是连通光滑流形上的平坦实向量丛，允许底流形有边界。固定基点 \(p\) 与基 \(b:\mathbb R^r\to E_p\)，单值表示记为 \(\rho\)。证明求值映射给出线性同构
	\[
		\{s\in\Gamma(E):\nabla s=0\}
		\longrightarrow
		\{v\in\mathbb R^r:\rho(g)v=v\text{ 对所有 }g\in\pi_1(M,p)\},
		\qquad s\longmapsto b^{-1}s(p).
	\]
	再对 \(T^2=S^1\times S^1\) 上由表示
	\[
		\rho(m,n)=\operatorname{diag}(1,(-1)^m,e^n),\qquad(m,n)\in\mathbb Z^2,
	\]
	构造的平坦丛，求整体平行截面空间的维数，并判断是否有整体平行标架。
\end{exercise}
\begin{solution}
	若 \(\nabla s=0\)，则在每个平行标架域中，\(s\) 的分量导数为零，故其分量在连通小域中恒定。因此沿任意连续路径 \(\gamma\) 都有 \(s(\gamma(1))=P_\gamma s(\gamma(0))\)。对于基点环路，这说明 \(s(p)\) 被 \(P_\gamma\) 固定；于是 \(v=b^{-1}s(p)\) 被 \(C_\gamma=b^{-1}P_\gamma b\) 固定，也就被其逆 \(\rho([\gamma])\) 固定。因为连通流形道路连通，\(s(p)\) 还通过运输决定所有点的取值，所以求值映射单射。

	反过来，设 \(v\) 被所有 \(\rho(g)\) 固定，令 \(v_p=bv\)。对 \(q\in M\)，任选从 \(p\) 到 \(q\) 的路径 \(\gamma\)，定义 \(s(q)=P_\gamma v_p\)。若另选 \(\eta\)，则
	\[
		P_\eta^{-1}P_\gamma v_p=P_{\gamma*\bar\eta}v_p=v_p,
	\]
	最后一个等式来自 \(v\) 同时被 \(\rho([\gamma*\bar\eta])\) 及其逆固定。因此定义与路径无关。在连通平行标架域中，\(s\) 是该标架的一个常系数线性组合，所以它光滑且平行。构造关于 \(v\) 线性，得到所需同构。

	由\cref{g1:thm:circle-fundamental,g1:prop:product-fundamental}，\(\pi_1(T^2)=\mathbb Z^2\)，两个整数记录两个圆周方向的绕数。题给矩阵满足 \(\rho(m,n)\rho(m',n')=\rho(m+m',n+n')\)，确为表示。若 \(v=(v_1,v_2,v_3)\) 是不动向量，取 \((m,n)=(1,0)\) 得 \(v_2=-v_2\)，所以 \(v_2=0\)；再取 \((0,1)\) 得 \(v_3=e\,v_3\)，所以 \(v_3=0\)。第一分量没有限制，不动空间恰为第一坐标轴。因此整体平行截面空间为一维。表示并不平凡，所以没有整体平行标架；存在一个处处非零平行截面，并不足以构成秩三丛的平行标架。
\end{solution}

\begin{exercise}\label{g1:ex:curvature-zero-at-one-point}
	在 \(\mathbb R^2\times\mathbb R\) 的标准标架中取联络 \(A=x^2\,\mathrm{d}y\)。计算曲率。令 \(\gamma_h\) 为正向绕行正方形 \([0,h]^2\) 的分片光滑边界路径，从原点出发，\(h>0\)。求 \(P_{\gamma_h}\)，并据此判断原点附近是否存在平行标架。解释为什么这不与原点处曲率为零矛盾。
\end{exercise}
\begin{solution}
	秩一时 \(A\wedge A=0\)，因此
	\[
		F=\mathrm{d}A=2x\,\mathrm{d}x\wedge\mathrm{d}y.
	\]
	它在原点乃至整条直线 \(x=0\) 上为零，但在原点的任一开邻域内都不恒为零。

	沿正方形边界，底边与顶边有 \(\mathrm{d}y=0\)，左边有 \(x=0\)，只有右边贡献积分。正向绕行时右边从 \((h,0)\) 到 \((h,h)\)，所以
	\[
		\int_{\gamma_h}A=h^2\int_0^h\mathrm{d}y=h^3.
	\]
	由秩一平行方程的解式，\(P_{\gamma_h}\) 是乘以 \(e^{-h^3}\)，当 \(h>0\) 时不等于恒等。任一原点邻域都包含充分小的这种正方形；若其中有平行标架，沿这条环路应保持全部分量，矛盾。

	平坦性要求曲率在一个区域内处处为零，\cref{g1:thm:flat-parallel-frame}使用的正是这个条件。单点的曲率为零只消去该点处环路运输的最低阶曲率贡献，并不保证附近的曲率也消失；本例的运输差异 \(e^{-h^3}-1=-h^3+O(h^6)\) 仍然存在。
\end{solution}


% !TeX root = ../Book1.tex
% 纲要标题：# 第五编　李群、齐性空间与规范几何
% 纲要位置：几何学计划纲要.md，第 731 行。
% 正文状态：编导读初稿；各章状态见相应文件。

% 纲要内容开始
% 
% 纲要内容结束

\BookPart{李群、齐性空间与规范几何}{g1:part:05}
\BookPartGuide
前四编已从具体空间走到流形、向量场与联络。第五编重新把对称性放到中心：连续变换自身组成怎样的流形，单位元附近的变化如何决定群运算，这些群又怎样作用于几何空间？第十六章从李群和李代数出发，建立指数、闭子群及同态积分，为后面的群作用与商空间提供必要条件。

\ProseParagraphBreak

第十七章将讨论轨道、稳定子与齐性空间，区分一个集合商与具有良好光滑结构的商。第十八章再用主丛组织随位置变化的群对称，将前一编的联络语言接到规范变换和平坦结构上。第十九、二十章为选读部分，分别展开紧李群的极大环面与根几何，以及表示和几何实现；阅读这些内容时，以相应章首列明的条件为准。

% !TeX root = ../../Book1.tex
% 纲要标题：## 第16章　李群的局部结构与整体积分
% 纲要位置：几何学计划纲要.md，第 733 行。
% 正文状态：初稿；尚未正式审核。

% 纲要内容开始
% 
% > 本章目标与证明范围：在有限维实李群范围内建立指数、闭子群和同态积分。李代数到李群的局部与整体结论分开，证明单连通性消除的具体障碍。
% 
% 纲要内容结束

\chapter{李群的局部结构与整体积分}
\label{g1:ch:16}
\BookChapterNumber{g1:ch:16}

\begin{chapterabstract}
	李群把光滑变化与群乘法放在同一对象上，单位元切空间则记录这些变化的一阶方向。左平移将切向量延伸为向量场，其括号给出李代数。我们先构造指数映射，证明局部可逆性，并以 BCH 的有限阶展开解释非交换乘法怎样出现在切空间中。随后用 Frobenius 定理把李子代数积分成浸入子群，再证明闭子群必为嵌入子群。最后把李代数同态先积分到单位元附近，再沿路径延拓；连续同伦的网格论证说明单连通性恰好消除了绕行带来的歧义。万有覆盖群和离散中心商进一步说明，相同李代数可以对应不同的整体李群。
\end{chapterabstract}

\begin{learninggoals}
	\begin{itemize}
		\item 从矩阵例子与一般定义认识李群，构造单位元切空间上的李括号，证明伴随作用和同态微分的相容性。
		\item 证明左不变场完备，建立指数的局部逆、自然性及微分公式，推导 BCH 三阶展开和乘积极限。
		\item 利用左不变积分分布构造连通浸入子群，证明闭子群定理，并用无理斜率子群辨别叶拓扑与子空间拓扑。
		\item 从局部同态构造路径延拓，证明单连通源上的整体积分，构造万有覆盖群及离散中心商，并判断同态何时能下降到商群。
	\end{itemize}
\end{learninggoals}

旋转既可以连续改变角度，也可以前后复合；可逆线性变换也同时具有光滑参数和群运算。前面我们分别研究了变换群、流形及向量场，现在可以把这些结构放在一起。群乘法使每一点都能与单位元比较，因此看似遍布整个群的微分信息，有希望由一个切空间组织起来。

\ProseParagraphBreak

不过，一个切空间究竟能记住多少整体信息？实数加法群与圆周群的单位切空间都是一维、括号都是零，但沿同一个局部方向运动，前者不会绕回起点，后者会。反过来，环面上的一条无理斜率运动虽由一个实参数描述，其像却可以稠密。本章一边建立局部与整体之间的联系，一边保留周期、闭性和拓扑所造成的差别。

\ProseParagraphBreak

本章的李群均为有限维实光滑李群，流形无边界、Hausdorff、第二可数；复矩阵群在这里也先按实流形处理。读者需要\chapref{g1:ch:11}的逆函数、常秩及嵌入概念，\chapref{g1:ch:13}的常微分方程、李括号与 Frobenius 定理，以及\chapref{g1:ch:05}的路径同伦和万有覆盖。指数微分的计算还会回顾\secref{g1:sec:1401}的一形式外微分。每次调用时将指出实际条件，尤其区分浸入与嵌入、连通与单连通。

\ProseParagraphBreak

本章在正文证明指数、闭子群与同态积分。这里整体同态的存在定理以已给定的李群 \(G,H\) 为对象，并要求源群 \(G\) 单连通；任意抽象有限维李代数能否先找到一个李群来实现，是另一项存在问题，不包含在本章的证明范围内。我们也不预设抽象李群都可以忠实地表示为矩阵群。

\ProseParagraphBreak

延伸阅读可查 Lee 第二版第19章的 Lie Subalgebras and Lie Subgroups（第505--506页），以及第20章的 One-Parameter Subgroups and the Exponential Map（第516--521页）、The Closed Subgroup Theorem（第522--524页）和 The Lie Correspondence（第530--532页）\cite[第505--506、516--524、530--532页]{Lee2012Smooth}。矩阵计算可对照 Hall 第二版 §§2.1--2.4、3.3--3.5；指数微分、BCH、同态及覆盖的阅读位置为 §§5.3--5.9\cite[\S\S2.1--2.4, 3.3--3.5, 5.3--5.9]{Hall2015}。不同文献可能从矩阵群或抽象流形出发，对照时应核清对象范围及左、右不变场的约定。

% !TeX root = ../../Book1.tex
% 纲要标题：### 16.1 Lie 群与 Lie 代数
% 纲要位置：几何学计划纲要.md，第 737 行。
% 正文状态：初稿；尚未正式审核。

% 纲要内容开始
% 
% * 矩阵群的光滑结构
% * 左不变向量场和单位元切空间
% * 伴随作用、伴随表示与括号
% * Lie 群同态的微分
% 
% 纲要内容结束

\section{Lie 群与 Lie 代数}
\label{g1:sec:1601}

旋转、平移和可逆线性变换既能复合，又能随参数光滑变化。群结构记录复合的规则，流形结构允许我们对这些变化求导。两者相容之后，单位元附近的变化便能通过平移传播到整个群；原本分散在各点的切向量，可以统一放回单位元的切空间研究。不过，一个向量空间本身还没有记录变换复合的次序。我们将看到，左不变向量场的 Lie 括号恰好给这个切空间补上一种运算，使它保留局部复合的部分信息。

\ProseParagraphBreak

本节中的 Lie 群均为有限维实 Lie 群，底流形 Hausdorff、第二可数且无边界，不要求连通。即使元素写成复矩阵，其流形维数也先按实数计算。矩阵群提供可计算的模型，后面的定义和证明则适用于一般 Lie 群，不要求预先把群嵌入某个矩阵空间。

\subsection{矩阵群的光滑结构}
\label{g1:subsec:1601:01}

要对群运算求导，先要规定乘法与求逆怎样和流形结构相容。矩阵群提供了直接的检验方法：先由矩阵方程建立流形结构，再把乘法与求逆限制到这个流形上。

\begin{definition}[Lie 群]\label{g1:def:lie-group}
	一个群 \(G\) 若同时是有限维、Hausdorff、第二可数的无边界实光滑流形，并且乘法与求逆映射
	\[
		m:G\times G\longrightarrow G,\qquad m(g,h)=gh,
		\qquad
		\iota:G\longrightarrow G,\qquad \iota(g)=g^{-1}
	\]
	均光滑，就称为\term[liqun]{Lie 群}[Lie group]。群的单位元记为 \(e\)。两个 Lie 群之间的光滑群同态称为\term[liquntongtai]{Lie 群同态}[Lie group homomorphism]；若它还是微分同胚，就称为 Lie 群同构。
\end{definition}

例如 \((\mathbb R^n,+)\) 是 Lie 群，单位元为 \(0\)。有限个 Lie 群的直积取乘积流形结构后仍是 Lie 群，因为乘法和求逆逐分量进行。可数离散群取零维流形结构也是 Lie 群：所有点都是坐标邻域，所有映射都光滑，可数性则保证第二可数。这里允许零维和不连通的例子，是因为群的离散信息并不会自动消失，只是稍后的切空间无法记录全部这类信息。

\ProseParagraphBreak

对 \(g\in G\)，记左平移和右平移为
\[
	L_g(h)=gh,\qquad R_g(h)=hg.
\]
它们都是微分同胚，逆分别为 \(L_{g^{-1}},R_{g^{-1}}\)：光滑性来自固定乘法的一个变量，逆的公式来自群律。特别地，各点切空间都与 \(\mathrm{T}_eG\) 线性同构。乘法在单位元处的微分还有一个直接的表达式：由 \(m(g,e)=g\)、\(m(e,h)=h\) 及微分的线性，
\begin{equation}\label{g1:eq:lie-group-unit-differentials}
	\mathrm{d}m_{(e,e)}(X,Y)=X+Y,\qquad \mathrm{d}\iota_e(X)=-X.
\end{equation}
第二式是对恒等式 \(m(g,\iota(g))=e\) 求导的结果，所用链式法则已在流形微分的定义中建立。

\begin{proposition}[经典矩阵群]\label{g1:prop:matrix-classical-lie-groups}
	设 \(n\geq1\)。下表中的矩阵集合具有其所在实矩阵空间的自然开子流形或嵌入子流形结构，配上矩阵乘法都是实 Lie 群。这里 \(A^*=\overline A^{\mathsf T}\)，复矩阵空间按实向量空间理解；表中的切空间均在单位矩阵 \(I\) 处。
	\[
		\begin{array}{c|c|c}
			G&\dim_{\mathbb R}G&\mathrm{T}_IG\\ \hline
			\operatorname{GL}(n,\mathbb R)&n^2&\operatorname{Mat}_{n\times n}(\mathbb R)\\
			\operatorname{GL}(n,\mathbb C)&2n^2&\operatorname{Mat}_{n\times n}(\mathbb C)\\
			\operatorname{SL}(n,\mathbb R)&n^2-1&\{X:\operatorname{tr}X=0\}\\
			\operatorname{SL}(n,\mathbb C)&2n^2-2&\{X:\operatorname{tr}X=0\}\\
			\operatorname{O}(n),\ \operatorname{SO}(n)&n(n-1)/2&\{X:X^{\mathsf T}+X=0\}\\
			\operatorname{U}(n)&n^2&\{X:X^*+X=0\}\\
			\operatorname{SU}(n)&n^2-1&\{X:X^*+X=0,\ \operatorname{tr}X=0\}.
		\end{array}
	\]
	其中 \(\operatorname{GL}\) 由可逆矩阵组成，\(\operatorname{SL}\) 再要求行列式为 \(1\)；\(\operatorname{O}(n)=\{A\in\operatorname{GL}(n,\mathbb R):A^{\mathsf T}A=I\}\)，\(\operatorname{SO}(n)=\operatorname{O}(n)\cap\operatorname{SL}(n,\mathbb R)\)；\(\operatorname{U}(n)=\{A\in\operatorname{GL}(n,\mathbb C):A^*A=I\}\)，\(\operatorname{SU}(n)=\operatorname{U}(n)\cap\operatorname{SL}(n,\mathbb C)\)。表中每行的矩阵域与该行的群一致。
\end{proposition}
\begin{proof}
	行列式是矩阵分量的多项式，故非零行列式的集合 \(\operatorname{GL}(n,\mathbb R)\)、\(\operatorname{GL}(n,\mathbb C)\) 是开集，实维数分别为 \(n^2,2n^2\)。矩阵乘法的分量是多项式，求逆的公式为 \(A^{-1}=\operatorname{adj}(A)/\det A\)，因此两种运算都光滑。复矩阵情形也如此，因为复数的加法、乘法、共轭以及非零复数的求逆都是实光滑映射。

	先计算行列式的微分。按列的多线性展开 \(\det(I+tB)\)，常数项为 \(1\)，一次项只能从一列取 \(tB\)、其余列取单位矩阵的列，故一次项为 \(t\sum_jB_{jj}=t\operatorname{tr}B\)。于是对可逆 \(A\)，
	\begin{equation}\label{g1:eq:determinant-differential}
		\mathrm{d}(\det)_A(H)=\det A\,\operatorname{tr}(A^{-1}H).
	\end{equation}
	在实矩阵的 \(\det A=1\) 处，取 \(H=(c/n)A\) 就使微分等于任意 \(c\in\mathbb R\)；在复矩阵情形，取 \(H=(z/n)A\) 则得到任意 \(z\in\mathbb C\)。因此 \(1\) 分别是到 \(\mathbb R\) 和到 \(\mathbb C\simeq\mathbb R^2\) 的行列式映射的正则值。由\cref{g1:thm:regular-level-set}，两种 \(\operatorname{SL}\) 都是嵌入子流形，其余维分别为 \(1,2\)，单位元切空间均由 \(\operatorname{tr}X=0\) 给出。

	对正交群，令 \(\operatorname{Sym}_n(\mathbb R)\) 为实对称矩阵空间，其维数为 \(n(n+1)/2\)，考虑
	\[
		Q:\operatorname{GL}(n,\mathbb R)\longrightarrow\operatorname{Sym}_n(\mathbb R),\qquad Q(A)=A^{\mathsf T}A.
	\]
	这里把靶空间取成对称矩阵空间很重要：若仍取全体矩阵，微分便不可能满射。在 \(Q(A)=I\) 处，
	\[
		\mathrm{d}Q_A(H)=A^{\mathsf T}H+H^{\mathsf T}A.
	\]
	给定对称矩阵 \(S\)，取 \(H=AS/2\)，就有 \(\mathrm{d}Q_A(H)=S\)。故 \(I\) 是正则值，\(\operatorname{O}(n)\) 是维数 \(n^2-n(n+1)/2=n(n-1)/2\) 的嵌入子流形，其单位元切空间为反对称矩阵。正交矩阵满足 \((\det A)^2=1\)，所以 \(\operatorname{SO}(n)\) 正是 \(\operatorname{O}(n)\) 中 \(\det A>0\) 的开子集，具有相同的局部维数和单位元切空间。

	复矩阵情形令 \(\operatorname{Herm}_n(\mathbb C)=\{S:S^*=S\}\)。它的对角项有 \(n\) 个实自由度，上三角有 \(n(n-1)/2\) 个复自由度，实维数为 \(n^2\)。映射 \(Q(A)=A^*A\) 的微分为 \(A^*H+H^*A\)。在 \(A^*A=I\) 处，对任意 Hermitian 矩阵 \(S\)，\(H=AS/2\) 仍给出 \(\mathrm{d}Q_A(H)=S\)。故 \(\operatorname{U}(n)\) 是实维数 \(2n^2-n^2=n^2\) 的嵌入子流形，\(\mathrm{T}_I\operatorname{U}(n)\) 由 \(X^*+X=0\) 给出。

	最后讨论 \(\operatorname{SU}(n)\)。圆周 \(S^1=\{z\in\mathbb C:|z|^2=1\}\) 是一维嵌入子流形，因为 \(z\mapsto|z|^2\) 在圆周上微分非零，其切空间为 \(\mathrm{T}_zS^1=iz\mathbb R\)。酉矩阵的行列式属于 \(S^1\)，由\cref{g1:prop:embedded-initial-property}，行列式限制给出光滑映射 \(\det:\operatorname{U}(n)\to S^1\)。上述 \(Q\) 的微分还说明
	\[
		\mathrm{T}_A\operatorname{U}(n)=\{AB:B^*+B=0\}.
	\]
	由\eqref{g1:eq:determinant-differential}，\(\mathrm{d}(\det)_A(AB)=\det A\,\operatorname{tr}B\)。取 \(B=(it/n)I\)，可得到 \(i\det A\,t\)，遍历 \(\mathrm{T}_{\det A}S^1\)。于是 \(1\) 是这次行列式映射的正则值，\(\operatorname{SU}(n)\) 的实余维为 \(1\)，单位元切空间再加上 \(\operatorname{tr}X=0\) 这一条件。反 Hermitian 矩阵的迹本来就是纯虚数，因此这里的迹零只有一个实条件，与 \(\operatorname{SL}(n,\mathbb C)\) 中的两个实条件并不矛盾。

	各集合由定义直接对乘法和求逆封闭。它们在上述构造中都具有嵌入子流形结构；环境矩阵群的乘法和求逆限制后仍光滑地取值于相应集合，理由再次是\cref{g1:prop:embedded-initial-property}。这些子流形的拓扑是子空间拓扑，因此 Hausdorff、第二可数，无边界性由正则值定理或开子流形构造保证。
\end{proof}

以上矩阵群的流形结构来自方程微分的秩。例如，正交群的证明把对称矩阵作为靶空间，使单位矩阵成为正则值。稍后的一般闭子群定理将从子群的闭性统一构造局部坐标，给出另一种建立流形结构的方法。

\subsection{左不变向量场和单位元切空间}
\label{g1:subsec:1601:02}

现在设 \(G\) 是任意有限维实 Lie 群，记 \(\mathfrak g=\mathrm{T}_eG\)。一个方向 \(X\in\mathfrak g\) 可以被左平移送到每个点：在 \(g\) 处取 \(\mathrm{d}(L_g)_eX\)。这个向量场精确表达了“各点都朝同一个群方向前进”。它一般不是某个坐标系中的常向量场；保持不变的是左平移给出的比较方式。

\begin{proposition}[左不变场]\label{g1:prop:left-invariant-fields}
	设 \(G\) 为有限维实 Lie 群。若 \(Z\in\mathfrak X(G)\) 对任意 \(g,h\in G\) 都满足
	\[
		Z_{gh}=\mathrm{d}(L_g)_hZ_h.
	\]
	则称它为\term[zuobubianxiangliangchang]{左不变向量场}[left-invariant vector field]。在单位元处求值，得到左不变向量场空间到 \(\mathfrak g=\mathrm{T}_eG\) 的实线性同构，其逆为
	\begin{equation}\label{g1:eq:left-invariant-extension}
		X\longmapsto X^L,\qquad X^L_g=\mathrm{d}(L_g)_eX.
	\end{equation}
	两个左不变场的 Lie 括号仍左不变。\(\mathfrak g\) 的任意一组基经上述延拓，得到 \(\mathrm{T}G\) 的一个整体光滑标架。
\end{proposition}
\begin{proof}
	乘法 \(m(g,h)=gh\) 光滑，对第二个变量在 \(h=e\) 处求微分，得到随 \(g\) 光滑变化的线性映射 \(\mathrm{d}(L_g)_e\)。因而 \(X^L\) 确为光滑向量场。由 \(L_{gh}=L_g\circ L_h\) 及链式法则，
	\[
		X^L_{gh}=\mathrm{d}(L_g)_h\,\mathrm{d}(L_h)_eX=\mathrm{d}(L_g)_hX^L_h,
	\]
	所以它左不变。反过来，若 \(Z\) 左不变，在定义中令 \(h=e\)，立即得到 \(Z_g=\mathrm{d}(L_g)_eZ_e\)，故 \(Z\) 完全由 \(Z_e\) 确定。两种操作互逆且都线性。

	左不变等价于 \((L_g)_*Z=Z\)。由\cref{g1:prop:lie-derivative-natural}中微分同胚保持 Lie 括号的结论，若 \(Z,W\) 左不变，则
	\[
		(L_g)_*[Z,W]=[(L_g)_*Z,(L_g)_*W]=[Z,W].
	\]
	故其括号也左不变。若 \(X_1,\ldots,X_d\) 为 \(\mathfrak g\) 的基，\(\mathrm{d}(L_g)_e\) 是线性同构，故 \(X_1^L(g),\ldots,X_d^L(g)\) 在每点构成基，给出整体标架。具体的切丛平凡化为 \((g,X)\mapsto\mathrm{d}(L_g)_eX\)，其逆在 \(\mathrm{T}_gG\) 上为 \(v\mapsto(g,\mathrm{d}(L_{g^{-1}})_gv)\)；乘法、求逆及其微分的光滑性说明正逆映射都光滑。
\end{proof}

\begin{definition}[Lie 代数]\label{g1:def:lie-algebra-group}
	实向量空间 \(V\) 上一个实双线性运算 \([\ ,\ ]:V\times V\to V\)，若满足反对称性 \([X,Y]=-[Y,X]\) 及 Jacobi 恒等式
	\[
		[X,[Y,Z]]+[Y,[Z,X]]+[Z,[X,Y]]=0,
	\]
	就使 \(V\) 成为一个\term[lidaishu]{Lie 代数}[Lie algebra]。对有限维实 Lie 群 \(G\)，在 \(\mathfrak g=\mathrm{T}_eG\) 上定义
	\[
		[X,Y]\coloneqq[X^L,Y^L]_e.
	\]
	所得 Lie 代数称为 \(G\) 的 Lie 代数，记为 \(\operatorname{Lie}(G)=\mathfrak g\)。\symidx{\(\mathfrak g=\mathrm{T}_eG\)：Lie 群的 Lie 代数}由\cref{g1:prop:left-invariant-fields}，其左不变延拓满足 \(([X,Y])^L=[X^L,Y^L]\)；实双线性、反对称性和 Jacobi 恒等式随即由\cref{g1:prop:lie-bracket-properties}的向量场版本传下来。
\end{definition}

一般向量场的 Lie 括号在一点的值还依赖附近的变化，不能只由该点的两个切向量计算。这里能够定义 \([X,Y]\in \mathrm{T}_eG\)，并非删掉了这种依赖，而是由左不变性为每个单位元切向量选定了唯一的整体延拓。因此群乘法参与了括号的定义；只给同一个流形选一个点，通常没有这样的运算。

\ProseParagraphBreak

以后还会用到两个代数概念。Lie 代数之间保持括号的线性映射称为\term[lidaishutongtai]{Lie 代数同态}[Lie algebra homomorphism]；线性子空间 \(\mathfrak h\subset\mathfrak g\) 若满足 \([\mathfrak h,\mathfrak h]\subset\mathfrak h\)，称为\term[lizidaishu]{Lie 子代数}[Lie subalgebra]。若更强地满足 \([\mathfrak g,\mathfrak h]\subset\mathfrak h\)，则称 \(\mathfrak h\) 为理想。前者保证子空间内部可以作括号，后者还控制它与整个 Lie 代数的相互作用。

\begin{proposition}[矩阵括号]\label{g1:prop:matrix-lie-bracket}
	设 \(n\ge1\)，\(\mathbb F=\mathbb R\) 或 \(\mathbb C\)，\(G\) 为 \(\operatorname{GL}(n,\mathbb F)\)、\(\operatorname{SL}(n,\mathbb F)\)、\(\operatorname{O}(n)\)、\(\operatorname{SO}(n)\)、\(\operatorname{U}(n)\)、\(\operatorname{SU}(n)\) 之一。这些矩阵群的流形结构见\cref{g1:prop:matrix-classical-lie-groups}。在各群相应的矩阵域中，将 \(\mathfrak g=\mathrm{T}_IG\) 视为矩阵空间的实线性子空间。则 Lie 代数的括号为
	\[
		[X,Y]=XY-YX.
	\]
\end{proposition}
\begin{proof}
	先在 \(\operatorname{GL}(n,\mathbb R)\) 或 \(\operatorname{GL}(n,\mathbb C)\) 中计算。左平移 \(L_A(B)=AB\) 的微分为 \(H\mapsto AH\)，故 \(X^L_A=AX\)、\(Y^L_A=AY\)。把它们看成开矩阵空间上的向量值函数，\(\mathrm{d}Y^L_A(H)=HY\)、\(\mathrm{d}X^L_A(H)=HX\)。依\eqref{g1:eq:lie-bracket-coordinates}的符号约定，
	\[
		[X^L,Y^L]_A
		=\mathrm{d}Y^L_A(AX)-\mathrm{d}X^L_A(AY)
		=A(XY-YX).
	\]
	在 \(A=I\) 处即得所需公式。复矩阵情况下以上乘法仍是实线性微分，计算不变。

	其余矩阵群 \(H\) 的左平移仍由矩阵乘法给出。若 \(X\in \mathrm{T}_IH\)，\(A\mapsto AX\) 在 \(H\) 上取值于 \(\mathrm{T}H\)，且通过嵌入 \(H\hookrightarrow\operatorname{GL}\) 与环境的同名左不变场相关。再次用\cref{g1:prop:lie-derivative-natural}，它们的括号也与环境括号相关，所以单位元处的公式仍为 \(XY-YX\)。这也从几何上证明了表中各切空间对矩阵交换子封闭。
\end{proof}

这些 Lie 代数通常记作 \(\mathfrak{gl}(n,\mathbb R)\)、\(\mathfrak{sl}(n,\mathbb R)\)、\(\mathfrak{so}(n)\)、\(\mathfrak u(n)\)、\(\mathfrak{su}(n)\) 等；\(\operatorname{O}(n)\) 和 \(\operatorname{SO}(n)\) 具有同一个 \(\mathfrak{so}(n)\)。\(\mathfrak{gl}(n,\mathbb C)\)、\(\mathfrak{sl}(n,\mathbb C)\) 此处仍可视为实 Lie 代数，只是它们还额外具有与括号相容的复向量空间结构。

\ProseParagraphBreak

Lie 代数会丢失一些整体信息。例如 \((\mathbb R,+)\) 和 \(S^1=\operatorname{U}(1)\) 的 Lie 代数都是一维实向量空间，其括号因反对称性而恒为零；但两群不可能同构，因为圆周紧而实直线不紧。任意离散 Lie 群的 Lie 代数甚至都是零空间。因此本章后面既要研究局部积分，也要另行考虑覆盖、单连通性和离散子群。

\subsection{伴随作用、伴随表示与括号}
\label{g1:subsec:1601:03}

固定 \(g\in G\)，共轭 \(C_g(h)=ghg^{-1}\) 把单位元固定，因而其微分成为 \(\mathfrak g\) 的一个线性变换。左平移把单位元处的方向搬到 \(g\)；共轭则将改变后的方向仍放在单位元处比较，由此得到伴随表示。

\begin{proposition}[伴随表示]\label{g1:prop:adjoint-representation}
	设 \(G\) 为有限维实 Lie 群，\(\mathfrak g=\mathrm{T}_eG\)。对 \(g\in G\)，令共轭映射 \(C_g:G\to G\) 为 \(C_g(h)=ghg^{-1}\)，并定义
	\[
		\operatorname{Ad}_g\coloneqq\mathrm{d}(C_g)_e:\mathfrak g\longrightarrow\mathfrak g.
	\]
	则 \(\operatorname{Ad}:G\to\operatorname{GL}(\mathfrak g)\) 是光滑群同态，且每个 \(\operatorname{Ad}_g\) 保持 Lie 括号。它称为 \(G\) 的\term[bansuibiaoshi]{伴随表示}[adjoint representation]。这里一般称一个到有限维向量空间可逆线性变换群的光滑群同态为该 Lie 群的表示。
\end{proposition}
\begin{proof}
	共轭 \(C_g\) 是微分同胚，逆为 \(C_{g^{-1}}\)，因此 \(\operatorname{Ad}_g\) 可逆。由 \(C_{gh}=C_g\circ C_h\) 及它们都固定 \(e\)，链式法则给出
	\[
		\operatorname{Ad}_{gh}=\operatorname{Ad}_g\operatorname{Ad}_h,\qquad
		\operatorname{Ad}_e=\operatorname{id}_{\mathfrak g}.
	\]
	映射 \((g,h)\mapsto ghg^{-1}\) 光滑，对第二个变量在 \(h=e\) 处求微分，得到的矩阵分量随 \(g\) 光滑。这说明 \(\operatorname{Ad}\) 是到 \(\operatorname{End}(\mathfrak g)\) 的光滑映射，而其值落在开子集 \(\operatorname{GL}(\mathfrak g)\)，所以也是到后者的光滑映射。

	恒等式 \(C_g\circ L_h=L_{C_g(h)}\circ C_g\) 在单位元处求导，给出
	\[
		(C_g)_*X^L=(\operatorname{Ad}_gX)^L.
	\]
	微分同胚保持向量场的 Lie 括号，故
	\[
		(C_g)_*[X^L,Y^L]
		=[(\operatorname{Ad}_gX)^L,(\operatorname{Ad}_gY)^L].
	\]
	两边在 \(e\) 处求值，即得 \(\operatorname{Ad}_g[X,Y]=[\operatorname{Ad}_gX,\operatorname{Ad}_gY]\)。
\end{proof}

再对 \(\operatorname{Ad}\) 在单位元处求微分。利用 \(\mathrm{T}_I\operatorname{GL}(\mathfrak g)=\operatorname{End}(\mathfrak g)\)，记
\[
	\operatorname{ad}\coloneqq\mathrm{d}(\operatorname{Ad})_e:
	\mathfrak g\longrightarrow\operatorname{End}(\mathfrak g),\qquad
	X\longmapsto\operatorname{ad}_X.
\]
\symidx{\(\operatorname{Ad}_g,\operatorname{ad}_X\)：伴随表示及其微分}大写 \(\operatorname{Ad}\) 的输入是群元素，小写 \(\operatorname{ad}\) 的输入是单位元切向量。下面的公式把这次求导与此前定义的 Lie 括号联系起来。

\begin{proposition}[伴随微分]\label{g1:prop:adjoint-differential-bracket}
	设 \(G\) 为有限维实 Lie 群，\(\mathfrak g=\mathrm{T}_eG\)，\(\operatorname{Ad}:G\to\operatorname{GL}(\mathfrak g)\) 为伴随表示。在 \(\mathrm{T}_I\operatorname{GL}(\mathfrak g)=\operatorname{End}(\mathfrak g)\) 的识别下，令 \(\operatorname{ad}=(\mathrm{d}\operatorname{Ad})_e\)。则对任意 \(X,Y\in\mathfrak g\)，
	\begin{equation}\label{g1:eq:ad-bracket}
		\operatorname{ad}_XY=[X,Y].
	\end{equation}
	此外 \(\operatorname{ad}_X\) 是括号的导子，而且
	\[
		[\operatorname{ad}_X,\operatorname{ad}_Y]=\operatorname{ad}_{[X,Y]},
	\]
	左边是线性自同态的交换子。
\end{proposition}
\begin{proof}
	由\cref{g1:thm:maximal-integral-curve}，左不变场 \(X^L\) 在单位元有局部积分曲线 \(\gamma(t)\)，满足 \(\gamma(0)=e\)、\(\gamma'(0)=X\)。只需在 \(0\) 附近取它，不需要知道它是否能延拓到所有时间。对固定的 \(g\)，左不变性给出
	\[
		\frac{\mathrm{d}}{\mathrm{d}t}\bigl(g\gamma(t)\bigr)
		=\mathrm{d}(L_g)_{\gamma(t)}X^L_{\gamma(t)}
		=X^L_{g\gamma(t)}.
	\]
	它在 \(t=0\) 经过 \(g\)，所以由积分曲线唯一性，\(X^L\) 的局部流为
	\begin{equation}\label{g1:eq:left-invariant-local-flow}
		\Phi_t(g)=g\gamma(t)=R_{\gamma(t)}(g)
	\end{equation}
	在两边定义的局部时间范围内成立。因而流的逆局部映射是 \(R_{\gamma(t)^{-1}}\)。注意，左不变场的流在这里以右乘写出；这是 \(g\) 固定而单位元曲线被左平移的结果。

	按照\cref{g1:prop:flow-pullback-derivative}的向量场拉回定义，在单位元有
	\begin{align*}
		(\Phi_t^*Y^L)_e
		&=\mathrm{d}(R_{\gamma(t)^{-1}})_{\gamma(t)}Y^L_{\gamma(t)}\\
		&=\mathrm{d}(R_{\gamma(t)^{-1}})_{\gamma(t)}\,\mathrm{d}(L_{\gamma(t)})_eY\\
		&=\mathrm{d}(C_{\gamma(t)})_eY
		=\operatorname{Ad}_{\gamma(t)}Y.
	\end{align*}
	对 \(t\) 在 \(0\) 处求导。左边由流拉回的导数公式等于 \([X^L,Y^L]_e=[X,Y]\)，右边由 \(\gamma'(0)=X\) 等于 \(\operatorname{ad}_XY\)。这证明\eqref{g1:eq:ad-bracket}，也固定了它与第13章括号约定的符号。

	最后，Jacobi 恒等式改写为
	\[
		[X,[Y,Z]]=[[X,Y],Z]+[Y,[X,Z]],
	\]
	恰好说明 \(\operatorname{ad}_X[Y,Z]=[\operatorname{ad}_XY,Z]+[Y,\operatorname{ad}_XZ]\)，即导子性质。将第二项移到左边，又得到 \((\operatorname{ad}_X\operatorname{ad}_Y-\operatorname{ad}_Y\operatorname{ad}_X)Z=\operatorname{ad}_{[X,Y]}Z\)。对所有 \(Z\) 成立便是所需的自同态恒等式。
\end{proof}

对矩阵群，共轭的微分直接给出 \(\operatorname{Ad}_AY=AYA^{-1}\)。若 \(A(t)\) 是经过 \(I\)、速度为 \(X\) 的群内曲线，对 \(A(t)A(t)^{-1}=I\) 求导可得 \((A^{-1})'(0)=-X\)，从而 \(\frac{\mathrm{d}}{\mathrm{d}t}|_0 A(t)YA(t)^{-1}=XY-YX\)，与抽象证明一致。

\subsection{Lie 群同态的微分}
\label{g1:subsec:1601:04}

群同态保持乘法，因而应当保持由乘法定义的无穷小运算。左不变场把这个问题转为向量场的比较：同态的微分先把它们变成相关向量场，再利用第13章已经证明的括号自然性。

\begin{proposition}[同态的微分]\label{g1:prop:lie-homomorphism-differential}
	设 \(F:G\to H\) 是有限维实 Lie 群之间的 Lie 群同态，\(\mathfrak g=\mathrm{T}_{e_G}G\)、\(\mathfrak h=\mathrm{T}_{e_H}H\)。则
	\[
		\varphi\coloneqq\mathrm{d}F_{e_G}:\mathfrak g\longrightarrow\mathfrak h
	\]
	是 Lie 代数同态。具体地，对任意 \(g\in G\)、\(X,Y\in\mathfrak g\)，
	\begin{align*}
		\mathrm{d}F_g(X^L_g)&=(\varphi X)^L_{F(g)},\\
		\varphi[X,Y]&=[\varphi X,\varphi Y],\\
		\varphi\circ\operatorname{Ad}_g&=\operatorname{Ad}_{F(g)}\circ\varphi.
	\end{align*}
	同态复合的单位元微分为单位元微分的复合，恒等同态的微分为恒等映射。特别地，\(\operatorname{im}\varphi\) 是 \(\mathfrak h\) 的 Lie 子代数，\(\ker\varphi\) 是 \(\mathfrak g\) 的理想。
\end{proposition}
\begin{proof}
	群同态满足 \(F(e_G)=e_H\) 及 \(F\circ L_g=L_{F(g)}\circ F\)。后一式在 \(e_G\) 处求导，得到
	\[
		\mathrm{d}F_g\,\mathrm{d}(L_g)_{e_G}X
		=\mathrm{d}(L_{F(g)})_{e_H}\,\mathrm{d}F_{e_G}X,
	\]
	即第一式。因此 \(X^L\) 与 \((\varphi X)^L\) 是 \(F\)-相关场，\(Y^L\) 与 \((\varphi Y)^L\) 也如此。\cref{g1:prop:lie-derivative-natural}说明它们的括号仍相关。这一结论无需 \(F\) 可逆或满射：对 \(H\) 上任意局部光滑函数 \(f\)，只要把两次求导展开，就有
	\[
		[X^L,Y^L](f\circ F)
		=\bigl([(\varphi X)^L,(\varphi Y)^L]f\bigr)\circ F.
	\]
	在 \(e_G\) 处求值，恰得 \(\varphi[X,Y]=[\varphi X,\varphi Y]\)。同态还满足 \(F\circ C_g=C_{F(g)}\circ F\)，在单位元求导便给出伴随表示的交织公式。

	复合与恒等映射的断言是链式法则。若 \(U=\varphi X\)、\(V=\varphi Y\)，则 \([U,V]=\varphi[X,Y]\) 仍在像中；若 \(Y\in\ker\varphi\)，则对任意 \(X\in\mathfrak g\)，\(\varphi[X,Y]=[\varphi X,0]=0\)。这分别证明像是 Lie 子代数、核是理想。
\end{proof}

\begin{proposition}[同态的秩与核]\label{g1:prop:lie-homomorphism-constant-rank}
	设 \(F:G\to H\) 是有限维实 Lie 群之间的 Lie 群同态，令 \(r=\operatorname{rank}\mathrm{d}F_{e_G}\)。则 \(F\) 在所有点的秩都等于 \(r\)。核
	\[
		K=\ker F=F^{-1}(e_H)
	\]
	是 \(G\) 的闭嵌入子流形，维数为 \(\dim G-r\)。它以限制的群运算成为 Lie 群，且
	\[
		\mathrm{T}_{e_G}K=\ker\mathrm{d}F_{e_G};
	\]
	这里用包含映射的微分将左边视为 \(\mathrm{T}_{e_G}G\) 的子空间。
\end{proposition}
\begin{proof}
	前一命题的第一式可以写成
	\[
		\mathrm{d}F_g
		=\mathrm{d}(L_{F(g)})_{e_H}\circ\mathrm{d}F_{e_G}\circ\mathrm{d}(L_{g^{-1}})_g.
	\]
	两端的平移微分都是同构，故各点的秩都等于 \(r\)。固定 \(k\in K\)，由\cref{g1:cor:manifold-constant-rank}，在 \(k\) 和 \(e_H\) 附近可取分别以这两点为原点的坐标，使 \(F\) 写成
	\[
		(u^1,\ldots,u^m)\longmapsto(u^1,\ldots,u^r,0,\ldots,0),\qquad m=\dim G.
	\]
	必要时缩小源坐标邻域，使其像包含于选定的靶坐标邻域。于是核在这个源邻域中恰由 \(u^1=\cdots=u^r=0\) 给出，而 \((u^{r+1},\ldots,u^m)\) 是核的局部坐标。坐标转换由环境转换限制到这些切片得到，光滑且互逆。这些局部描述与子空间拓扑一致，给出 \(K\) 的嵌入子流形结构；它从 \(G\) 继承 Hausdorff 性和第二可数性。其切空间在每个 \(k\in K\) 都是 \(\ker\mathrm{d}F_k\)，因为坐标中两者均由前 \(r\) 个分量为零的向量组成。

	单点 \(\{e_H\}\) 在 Hausdorff 空间 \(H\) 中闭，\(F\) 连续，故 \(K\) 闭。群同态的核是子群，因此 \(G\) 的乘法和求逆限制后取值于 \(K\)。这些限制作为到 \(G\) 的映射光滑，由\cref{g1:prop:embedded-initial-property}，作为到 \(K\) 的映射也光滑。由此 \(K\) 满足 Lie 群的全部条件，单位元切空间的公式已在局部坐标计算中得到。
\end{proof}

恒秩定理把核写成坐标切片，给出它的嵌入子流形结构；如果 \(\mathrm{d}F_{e_G}\) 不满射，\(e_H\) 并非正则值，这时正是恒秩标准形发挥作用。像的整体拓扑还受远处源点的累积影响，\(F(G)\) 在 \(H\) 中是否闭或嵌入，需要下一节对子群作进一步分析。

\ProseParagraphBreak

例如实矩阵群的行列式同态 \(\det:\operatorname{GL}(n,\mathbb R)\to\mathbb R^\times\) 在单位元的微分是迹，核是 \(\operatorname{SL}(n,\mathbb R)\)，与此前的直接计算一致。若一般同态的单位元微分单射，它在各点都是浸入；但它仍可能把不同群元素送到同一点。下一节的指数映射将进一步把单位元切向量与群内曲线联系起来，之后才能系统讨论一个 Lie 代数同态能够恢复多少群同态的信息。

% 纲要条目排为编号小节标题，并进入目录与书签；说明段仍保留于纲要注释中。

% !TeX root = ../../Book1.tex
% 纲要标题：### 16.2 指数映射与局部乘法
% 纲要位置：几何学计划纲要.md，第 744 行。
% 正文状态：初稿；尚未正式审核。

% 纲要内容开始
% 
% * 单参数子群及完备性
% * 指数映射的微分和局部逆
% * 共轭与指数的相容
% * 局部乘法与 BCH 的三阶展开
% 
% 纲要内容结束

\section{指数映射与局部乘法}
\label{g1:sec:1602}

一个单位元处的切向量描述瞬时运动。左平移把同一个瞬时方向放到每一点，形成左不变向量场；沿它运动，就有望把切向量变成群元素。这条路线给出指数映射。它还回答另一个问题：把两个很小的运动相乘，结果与将两个切向量直接相加有多大差别？Lie 括号将成为最先出现的修正项。

\ProseParagraphBreak

本节的 Lie 群均为有限维实 Lie 群，流形无边界、Hausdorff、第二可数。记单位元为 \(e\)，Lie 代数为 \(\mathfrak g=\mathrm{T}_eG\)，\(X\in\mathfrak g\) 对应的左不变向量场记为 \(X^L_g=(\mathrm{d}L_g)_eX\)。这一对应及其括号相容性已在\secref{g1:sec:1601}建立。以下所有群乘积都有顺序，不能先把它们当作向量加法。

\subsection{单参数子群及完备性}
\label{g1:subsec:1602:01}

\begin{definition}\label{g1:def:one-parameter-subgroup}
	设 \(G\) 为有限维实 Lie 群。光滑同态 \(c:(\mathbb R,+)\to G\) 称为\term[dancanshuziqun]{单参数子群}[one-parameter subgroup]，即
	\[
		c(s+t)=c(s)c(t),\qquad c(0)=e.
	\]
	这里名称指带参数的同态，不要求 \(c\) 单射，也不预设其像是嵌入子流形。
\end{definition}

若 \(X=c'(0)\)，对 \(c(t+h)=c(t)c(h)\) 在 \(h=0\) 求导，便得 \(c'(t)=(\mathrm{d}L_{c(t)})_eX\)。所以单参数子群必是从单位元出发的左不变场积分曲线。反过来，局部初值定理只给出一小段积分曲线；要获得整个 \(\mathbb R\) 上的同态，还须证明它不会在有限时间终止。

\begin{theorem}[指数映射]\label{g1:thm:lie-exponential}
	设 \(G\) 为有限维实 Lie 群。对每个 \(X\in\mathfrak g\)，左不变场 \(X^L\) 完备，且存在唯一单参数子群 \(c_X\) 满足 \(c_X'(0)=X\)。定义
	\[
		\exp_G:\mathfrak g\longrightarrow G,\qquad \exp_GX=c_X(1).
	\]
	则 \(\exp_G\) 光滑，并且
	\begin{equation}\label{g1:eq:one-parameter-exponential}
		c_X(t)=\exp_G(tX),\qquad
		\exp_G((s+t)X)=\exp_G(sX)\exp_G(tX).
	\end{equation}
	\(X^L\) 的整体流为 \(\Phi_t(g)=g\exp_G(tX)\)。
\end{theorem}
\begin{proof}
	由\cref{g1:thm:maximal-integral-curve}先取过 \(e\) 的局部积分曲线 \(c\)，将其限制在 \((-2a,2a)\)，其中 \(a>0\)。左不变性说明，对于固定 \(s\)，\(t\mapsto c(s)c(t)\) 与 \(t\mapsto c(s+t)\) 都是 \(X^L\) 的积分曲线，并在 \(t=0\) 取相同值。因此只要相关参数留在这段局部区间中，唯一性给出局部乘法规律 \(c(s)c(t)=c(s+t)\)。

	现在把同一小段反复左平移。对每个整数 \(n\)，在区间 \((na-a,na+a)\) 上令
	\[
		\widehat c(t)=c(a)^n c(t-na).
	\]
	相邻区间交叠时，令 \(u=t-(n+1)a\in(-a,0)\)，则 \(c(a)c(u)=c(a+u)\)，所以两式相同。每段是光滑积分曲线，因而粘成整个实轴上的光滑积分曲线，且在零点取 \(e\)。任意 \(g\in G\) 出发的曲线 \(g\widehat c(t)\) 也在整个实轴有定义，所以 \(X^L\) 完备。对全局曲线再次使用初值唯一性，得到 \(\widehat c(s+t)=\widehat c(s)\widehat c(t)\)。前段对任意单参数子群的求导说明，另一个具有相同初速度的单参数子群必须是同一初值解，故唯一。记它为 \(c_X\)。

	曲线 \(s\mapsto c_X(ts)\) 满足场 \((tX)^L\) 的方程，并在零点取单位元，包括 \(t=0\) 和负 \(t\) 的情形。因此唯一性给出 \(c_{tX}(s)=c_X(ts)\)，令 \(s=1\) 便得 \(\exp_G(tX)=c_X(t)\)。对 \(X\) 的光滑性调用\cref{g1:thm:ode-smooth-dependence}：场 \((g,X)\mapsto(\mathrm{d}L_g)_eX\) 联合光滑，任意固定 \(X_0\) 的轨道已在紧时间区间 \([0,1]\) 上存在，所以附近 \(X\) 的解在共同区间上光滑依赖参数。于 \(t=1\) 求值，得到 \(\exp_G\) 在每个 \(X_0\) 附近光滑。最后，从 \(g\) 出发的积分曲线已经写成 \(g c_X(t)\)，即所列流公式。
\end{proof}

这个映射称为 Lie 群的\term[zhishuyingshe]{指数映射}[exponential map]。\symidx{\(\exp_G\)：Lie 群的指数映射}它的定义依赖群结构和左不变场，与第九章由度量和测地线定义的指数映射具有不同的输入数据。两者都把切向量送到某条指定曲线的终点，但不能未经说明互换。

\begin{example}\label{g1:exm:matrix-exponential-series}
	在 \(\operatorname{GL}(n,\mathbb R)\) 中，\(X^L_A=AX\)，因此积分曲线满足 \(C'=CX\)、\(C(0)=I\)。矩阵级数
	\[
		e^{tX}=\sum_{k=0}^{\infty}\frac{t^kX^k}{k!}
	\]
	在每个有界时间区间绝对一致收敛，其逐项导数级数也如此，因为可分别用标量指数级数控制矩阵范数。故它可逐项求导，满足所需方程及初值。\(e^{tX}e^{-tX}=I\) 由乘积求导及零点值验证，所以该曲线始终可逆；唯一性给出 \(\exp_{\operatorname{GL}(n,\mathbb R)}X=e^X\)。复矩阵按实变量处理时完全相同。对抽象 Lie 群，\(X\) 未必是矩阵，定义仍是前面的积分曲线，不能预先写一个不存在的矩阵幂级数。
\end{example}

在加法群 \(\mathbb R^n\) 中，指数就是 \(X\mapsto X\)；在正实数乘法群中，\(X^L_x=xX\)，指数就是通常的 \(X\mapsto e^X\)。在圆周群中，把 \(\mathrm{T}_1S^1\) 写成 \(i\mathbb R\)，便有 \(\exp_{S^1}(it)=e^{it}\)。这三个例子分别显示指数可能是恒等映射、一般微分同胚，也可能因周期而失去单射性。

\subsection{指数映射的微分和局部逆}
\label{g1:subsec:1602:02}

\begin{proposition}[指数坐标]\label{g1:prop:exponential-local-diffeomorphism}
	设 \(G\) 为有限维实 Lie 群。用向量空间的平移识别 \(\mathrm{T}_0\mathfrak g=\mathfrak g\)，则
	\[
		(\mathrm{d}\exp_G)_0=\operatorname{id}_{\mathfrak g}.
	\]
	因而零点有一个开邻域 \(V\subset\mathfrak g\)，单位元有一个开邻域 \(U\subset G\)，使 \(\exp_G:V\to U\) 为微分同胚。其局部逆记为 \(\log:U\to V\)，称由此得到的坐标为指数坐标。
\end{proposition}
\begin{proof}
	对 \(Y\in\mathfrak g\)，曲线 \(t\mapsto tY\) 在零点速度为 \(Y\)，而
	\[
		(\mathrm{d}\exp_G)_0Y
		=\left.\frac{\mathrm{d}}{\mathrm{d}t}\right|_{t=0}\exp_G(tY)
		=c_Y'(0)=Y.
	\]
	于是微分可逆。\cref{g1:thm:inverse-function}适用于这两个同维无边界流形间的光滑映射，给出所需邻域及光滑局部逆。必要时可继续缩小 \(V\)，使其关于原点对称且为星形。
\end{proof}

结论只保证单位元附近的可逆性。圆周的指数已有 \(2\pi\) 周期；一般连通群的指数甚至未必满射。不过，单位元的小邻域已足以生成整个单位连通分支。

\begin{proposition}[由指数生成单位分支]\label{g1:prop:exponentials-generate-identity-component}
	设 \(G\) 为有限维实 Lie 群，\(G^0\) 为包含单位元的连通分支。则 \(G^0\) 是开正规子群，并且由任意足够小的对称指数邻域 \(U=\exp_G(V)\) 生成。特别地，连通 Lie 群的每个元素都是有限个指数的乘积。
\end{proposition}
\begin{proof}
	流形局部道路连通，故道路分支均开，且连通分支就是道路分支。若 \(g,h\) 都能用路径从 \(e\) 到达，路径的逐点乘积和逐点逆给出到 \(gh\)、\(g^{-1}\) 的路径，所以 \(G^0\) 为子群。共轭把单位元保持不变，又把路径变成路径，故 \(G^0\) 正规且开。

	取包含于 \(G^0\) 的对称指数邻域 \(U\)。令 \(K=\bigcup_{m\ge1}U^m\)。由于 \(e\in U=U^{-1}\)，这确是由 \(U\) 生成的子群。它开，因为每个有限乘积集是开集的平移之并；它在 \(G^0\) 中的每个陪集也开，因此补集开。于是 \(K\) 在连通空间 \(G^0\) 中既开又闭且非空，只能等于 \(G^0\)。
\end{proof}

为了计算局部乘法的高阶项，还需要知道非零点处的指数微分。这里应先把输出切向量移回单位元，才能与输入 \(Y\in\mathfrak g\) 比较。定义一个 \(\mathfrak g\) 值一形式
\[
	\theta_g=(\mathrm{d}L_{g^{-1}})_g:\mathrm{T}_gG\longrightarrow\mathfrak g.
\]
它满足 \(\theta(X^L)=X\)。在矩阵群中，\(\theta_g(v)=g^{-1}v\)；一般 Lie 群没有矩阵乘法，也仍由左平移微分定义这个对象。

\begin{proposition}[指数微分的积分公式]\label{g1:prop:lie-exponential-differential}
	设 \(G\) 为有限维实 Lie 群，\(\mathfrak g=\mathrm{T}_eG\)，\(X,Y\in\mathfrak g\)。令 \(\theta_g=(\mathrm{d}L_{g^{-1}})_g:\mathrm{T}_gG\to\mathfrak g\)，并记 \(\operatorname{ad}_X(Y)=[X,Y]\)，则
	\begin{equation}\label{g1:eq:lie-dexp-integral}
		\theta_{\exp_GX}\bigl((\mathrm{d}\exp_G)_XY\bigr)
		=J_XY,\qquad
		J_X=\int_0^1 e^{-s\operatorname{ad}_X}\,\mathrm{d}s.
	\end{equation}
	右侧的指数是有限维向量空间 \(\mathfrak g\) 上线性算子的指数，因而有收敛展开
	\begin{equation}\label{g1:eq:lie-dexp-series}
		J_X=\sum_{k=0}^{\infty}\frac{(-1)^k}{(k+1)!}\operatorname{ad}_X^k
		=I-\tfrac12\operatorname{ad}_X+\tfrac16\operatorname{ad}_X^2+O(\|X\|^3).
	\end{equation}
\end{proposition}
\begin{proof}
	先计算 \(\theta\) 的外微分。将\eqref{g1:eq:exterior-one-form-evaluation}逐分量应用于 \(\mathfrak g\) 值形式，对左不变场有
	\[
		\mathrm{d}\theta(X^L,Y^L)=-\theta([X^L,Y^L])=-[X,Y].
	\]
	两边都是对切向输入逐点双线性的二形式，而左不变场的点值张成每个切空间，故一般输入 \(V,W\) 满足 \(\mathrm{d}\theta(V,W)=-[\theta(V),\theta(W)]\)。

	取二参数映射 \(a(s,t)=\exp_G(s(X+tY))\)，并记
	\[
		v(s,t)=\theta(\partial_sa),\qquad u(s,t)=\theta(\partial_ta).
	\]
	对 \(a^*\theta=v\,\mathrm{d}s+u\,\mathrm{d}t\) 作外微分。外微分与拉回交换，见\cref{g1:prop:exterior-pullback}，所以刚才的恒等式给出
	\[
		\partial_su-\partial_tv=-[v,u].
	\]
	固定 \(t\) 时，\(s\mapsto a(s,t)\) 是初速度 \(X+tY\) 的单参数子群，因此 \(v=X+tY\)。在 \(t=0\) 处得到有限维线性方程
	\[
		\partial_su(s,0)=Y-\operatorname{ad}_Xu(s,0),\qquad u(0,0)=0.
	\]
	其解为 \(u(s,0)=\int_0^s e^{-(s-r)\operatorname{ad}_X}Y\,\mathrm{d}r\)：直接求导检查方程与初值，再用唯一性即可确认。令 \(s=1\)，左端是所需的左平移后微分，换元得到\eqref{g1:eq:lie-dexp-integral}。算子指数级数在 \(s\in[0,1]\) 上绝对一致收敛，可以逐项积分；双线性括号在有限维空间中有界，故 \(\|\operatorname{ad}_X\|\le C\|X\|\)，也就得到余项估计。
\end{proof}

公式中分母 \((k+1)!\) 使级数对每个 \(X\) 收敛，但 \(J_X\) 未必处处可逆。指数的局部可逆性在零点附近来自 \(J_0=I\)；离开这个邻域以后，不能仅凭级数存在就保证没有奇点。

\subsection{共轭与指数的相容}
\label{g1:subsec:1602:03}

\begin{proposition}[指数的自然性]\label{g1:prop:exponential-naturality}
	设 \(F:G\to H\) 为有限维实 Lie 群的光滑同态，\(\varphi=(\mathrm{d}F)_{e_G}\)。则对每个 \(X\in\mathfrak g\)，有
	\begin{equation}\label{g1:eq:exponential-homomorphism}
		F(\exp_GX)=\exp_H(\varphi X).
	\end{equation}
	特别地，若 \(g\in G\)，则
	\begin{equation}\label{g1:eq:exponential-conjugation}
		g\exp_G(X)g^{-1}=\exp_G(\operatorname{Ad}_gX),\qquad
		\operatorname{Ad}_{\exp_GX}=e^{\operatorname{ad}_X}.
	\end{equation}
\end{proposition}
\begin{proof}
	复合曲线 \(t\mapsto F(\exp_G(tX))\) 是 \(H\) 中的单参数子群，零点速度为 \(\varphi X\)。单参数子群的唯一性因此使它等于 \(t\mapsto\exp_H(t\varphi X)\)，令 \(t=1\) 得第一式。取同态 \(F=C_g\)、\(C_g(h)=ghg^{-1}\)，其单位元微分是 \(\operatorname{Ad}_g\)，便得共轭公式。

	再取 \(F=\operatorname{Ad}:G\to\operatorname{GL}(\mathfrak g)\)。\cref{g1:prop:adjoint-representation}及\eqref{g1:eq:ad-bracket}说明它是光滑同态，单位元微分为 \(X\mapsto\operatorname{ad}_X\)。刚证明的自然性给出 \(\operatorname{Ad}_{\exp_GX}=\exp_{\operatorname{GL}(\mathfrak g)}(\operatorname{ad}_X)\)，而\cref{g1:exm:matrix-exponential-series}把右端识别为算子指数。
\end{proof}

两个公式分别联系群上的共轭、单位切空间上的线性作用及 Lie 代数的括号。在矩阵群中，它们成为 \(g e^Xg^{-1}=e^{gXg^{-1}}\) 和 \(e^XYe^{-X}=e^{\operatorname{ad}_X}Y\)。后一式的前几项是 \(Y+[X,Y]+\tfrac12[X,[X,Y]]+\cdots\)，因此括号描述共轭造成的最初变化。

\begin{corollary}[交换方向的指数]\label{g1:cor:commuting-exponentials}
	设 \(G\) 为有限维实 Lie 群，\(X,Y\in\mathfrak g\)，且 \([X,Y]=0\)。则对任意实数 \(s,t\)，\(\exp_G(sX)\) 与 \(\exp_G(tY)\) 交换，并且 \(\exp_G(X+Y)=\exp_G(X)\exp_G(Y)\)。
\end{corollary}
\begin{proof}
	由于 \(\operatorname{ad}_XY=0\)，有 \(e^{s\operatorname{ad}_X}Y=Y\)。将其代入共轭公式，得到两条单参数子群的所有元素交换。令 \(c(t)=\exp_G(tX)\exp_G(tY)\)。因这些元素交换，\(c(s+t)=c(s)c(t)\)，故 \(c\) 是单参数子群。群乘法在 \((e,e)\) 处的微分为 \((X,Y)\mapsto X+Y\)：分别固定一个因子时都是恒等映射，微分又线性。因而 \(c'(0)=X+Y\)，唯一性使 \(c(t)=\exp_G(t(X+Y))\)，令 \(t=1\) 即得结论。
\end{proof}

\begin{corollary}[连通源上的唯一性]\label{g1:cor:lie-homomorphism-uniqueness}
	设 \(G,H\) 为有限维实 Lie 群，且 \(G\) 连通。若两个光滑同态 \(F_1,F_2:G\to H\) 在单位元处具有相同微分，则 \(F_1=F_2\)。
\end{corollary}
\begin{proof}
	自然性说明两同态在全部指数上相同。由\cref{g1:prop:exponentials-generate-identity-component}，\(G\) 的每个元素都是有限个指数的乘积；同态保持乘积，所以在每个群元素上都相同。
\end{proof}

这个结论提供唯一性，没有提供存在性。给定一个 Lie 代数同态能否成为群同态的微分，还可能受到环路的限制；\secref{g1:sec:1604}将具体构造延拓并解释单连通性的作用。若源群不连通，单位元微分甚至不记录其他连通分支上的离散选择，例如离散群的单位切空间总是零。

\subsection{局部乘法与 BCH 的三阶展开}
\label{g1:subsec:1602:04}

缩小指数坐标域，使 \(X,Y\) 充分小时 \(\exp_GX\exp_GY\) 仍在 \(\log\) 的定义域，记
\[
	\mu(X,Y)=\log(\exp_GX\exp_GY).
\]
这个局部运算把原来的群乘法写回向量空间。若 \(X,Y\) 交换，上一推论说明它就是 \(X+Y\)；不交换时，括号给出修正。下面证明\term[BCHgongshi]{Baker--Campbell--Hausdorff 公式}[Baker--Campbell--Hausdorff formula]的三阶形式，简称 BCH 公式。\footnote{Baker，Henry Frederick Baker（1866--1956），英国数学家，研究几何与分析；Campbell，John Edward Campbell（1862--1924），生于爱尔兰、在英国任教的数学家，研究连续变换群与微分几何。Hausdorff 的人物介绍见\chapref{g1:ch:04}。}

\begin{proposition}[BCH 的三阶展开]\label{g1:prop:bch-finite-order}
	设 \(G\) 为有限维实 Lie 群，在 \(\mathfrak g\) 上任选一个范数。对充分小的 \(X,Y\)，有
	\begin{equation}\label{g1:eq:bch-third-order}
		\begin{aligned}
			\log(\exp_GX\exp_GY)
			={}&X+Y+\tfrac12[X,Y]\\
			&+\tfrac1{12}[X,[X,Y]]+\tfrac1{12}[Y,[Y,X]]
			+O((\|X\|+\|Y\|)^4).
		\end{aligned}
	\end{equation}
	余项是在某个共同邻域内、相对于所选范数的统一估计。本命题只陈述有限阶 Taylor 展开。
\end{proposition}
\begin{proof}
	我们把第二个因子逐渐乘进去，比较这条运动的左平移后速度。对 \(0\le t\le1\)，令
	\[
		Z(t)=\log(\exp_GX\exp_G(tY)).
	\]
	只要 \(X,Y\) 在足够小的共同邻域中，整条路径都在 \(\log\) 的定义域。其左平移后的速度是 \(Y\)：乘上固定的左因子 \(\exp_GX\) 不改变 \(\theta\) 对速度的值，而 \(\exp_G(tY)\) 的该值本来就是 \(Y\)。由\eqref{g1:eq:lie-dexp-integral}，得到
	\[
		J_{Z(t)}Z'(t)=Y,\qquad Z(0)=X.
	\]
	因为 \(J_0=I\)，小邻域内 \(J_Z\) 可逆。将\eqref{g1:eq:lie-dexp-series}乘以待定的二次多项式，逐次比较可得
	\[
		J_Z^{-1}=I+\tfrac12\operatorname{ad}_Z+\tfrac1{12}\operatorname{ad}_Z^2+O(\|Z\|^3).
	\]
	确切地说，右侧多项式与 \(J_Z\) 的乘积是 \(I+O(\|Z\|^3)\)；再乘以在小邻域上一致有界的 \(J_Z^{-1}\)，便证明所列余项。因此
	\begin{equation}\label{g1:eq:bch-evolution}
		Z'=Y+\tfrac12[Z,Y]+\tfrac1{12}[Z,[Z,Y]]+O(\delta^4),
		\qquad \delta=\|X\|+\|Y\|.
	\end{equation}
	这里 \(Z(t)=O(\delta)\) 对 \(0\le t\le1\) 一致成立：映射 \((t,X,Y)\mapsto Z(t)\) 光滑，在 \((X,Y)=(0,0)\) 为零，沿连接零点的线段积分其一阶微分，并利用紧时间区间上的共同界即可。故逆算子余项作用于 \(Y\) 后确为 \(O(\delta^4)\)。

	先从上式积分，得到 \(Z=X+tY+O(\delta^2)\)。再代入一次，利用 \([Y,Y]=0\)，得到
	\[
		Z=X+tY+\tfrac t2[X,Y]+O(\delta^3).
	\]
	将这个精度代入一次括号项，将 \(X+tY+O(\delta^2)\) 代入二重括号项，便有
	\begin{align*}
		Z'
		&=Y+\tfrac12[X,Y]+\tfrac t4[[X,Y],Y]\\
		&\quad+\tfrac1{12}[X+tY,[X,Y]]+O(\delta^4)\\
		&=Y+\tfrac12[X,Y]+\tfrac1{12}[X,[X,Y]]
		-\tfrac t6[Y,[X,Y]]+O(\delta^4).
	\end{align*}
	所有余项在同一紧时间区间上一致，积分保持其阶数。从 \(t=0\) 积分到 \(1\)，并使用 \(-[Y,[X,Y]]=[Y,[Y,X]]\)，即得\eqref{g1:eq:bch-third-order}。
\end{proof}

三阶公式的推导不要求把 \(G\) 嵌入矩阵群，也没有用一个未证明收敛的无穷括号级数替代群乘法。对后面闭子群定理的局部论证，二阶及其余项控制已经够用；保留三阶则能看清两个方向在高一阶怎样再次相互影响。

\begin{corollary}[乘积极限]\label{g1:cor:lie-product-limit}
	设 \(G\) 为有限维实 Lie 群，\(X,Y\in\mathfrak g\)、\(t\in\mathbb R\)。则
	\[
		\left(\exp_G\left(\frac{tX}{n}\right)
		\exp_G\left(\frac{tY}{n}\right)\right)^n
		\longrightarrow\exp_G(t(X+Y))\qquad(n\to\infty).
	\]
\end{corollary}
\begin{proof}
	当 \(n\) 充分大时，小乘积在对数邻域中。令其对数为 \(Z_n\)，BCH 的二阶精度给出
	\[
		Z_n=\frac t n(X+Y)+O(n^{-2}),\qquad nZ_n\longrightarrow t(X+Y).
	\]
	每个小乘积恰等于 \(\exp_GZ_n\)，而单参数子群的乘法规律给出 \((\exp_GZ_n)^n=\exp_G(nZ_n)\)。最后使用指数的连续性即可。这里取幂的是同一个群元素，无须交换两个不同方向的指数。
\end{proof}

\begin{example}\label{g1:exm:heisenberg-bch}
	三阶上三角幺幂矩阵是对角条目均为 \(1\) 的上三角矩阵；三个上三角条目给出 \(\mathbb R^3\) 坐标，乘法和求逆为多项式，因此它们构成 Lie 群。在其 Lie 代数中取 \(X=aE_{12}+bE_{23}+cE_{13}\)，其中 \(E_{ij}\) 表示只有第 \((i,j)\) 项为 \(1\) 的矩阵。因 \(X^2=abE_{13}\)、\(X^3=0\)，有
	\[
		\exp X=I+aE_{12}+bE_{23}+\left(c+\tfrac12ab\right)E_{13}.
	\]
	对 \(Y=a'E_{12}+b'E_{23}+c'E_{13}\)，有 \([X,Y]=(ab'-a'b)E_{13}\)，且所有二重括号都为零。直接将上面的两个指数矩阵相乘，再用同一公式读回 \((a+a',b+b',c+c'+\tfrac12(ab'-a'b))\)，得到精确等式
	\[
		\log(\exp X\exp Y)=X+Y+\tfrac12[X,Y].
	\]
	精确性来自本例的实际矩阵计算；一般 Lie 群不能只因有限阶展开里暂时没有写出高阶项，就把余项删去。
\end{example}

% !TeX root = ../../Book1.tex
% 纲要标题：### 16.3 Lie 子代数的积分与闭子群定理
% 纲要位置：几何学计划纲要.md，第 751 行。
% 正文状态：初稿；尚未正式审核。

% 纲要内容开始
% 
% * Lie 子代数的积分分布
% * 浸入子群的构造
% * 闭子群定理
% * 闭性与嵌入性的关系
% 
% 纲要内容结束

\section{Lie 子代数的积分与闭子群定理}
\label{g1:sec:1603}

单位元处的一组允许方向，怎样长成一个子群？线性子空间还不够，因为沿两种方向交替运动会产生括号方向；要求这些新方向仍被允许，就得到 Lie 子代数。Frobenius 定理将把它积分成经过单位元的最大叶，但这片叶首先只有自己的流形拓扑，未必采用原群的子空间拓扑。闭子群定理将进一步说明，若一个代数子群在原群中闭，单凭这个拓扑条件就足以恢复嵌入子流形结构。

\ProseParagraphBreak

以下 Lie 群均为有限维实、Hausdorff、第二可数且无边界。固定 Lie 群 \(G\)，单位元为 \(e\)，Lie 代数为 \(\mathfrak g=\mathrm{T}_eG\)。我们使用\cref{g1:thm:frobenius,g1:thm:maximal-connected-leaf}中已经构造的最大积分叶，并使用上一节的指数映射及局部逆。本节只讨论子群自身的流形结构；陪集空间的流形结构将在下一章构造。

\subsection{Lie 子代数的积分分布}
\label{g1:subsec:1603:01}

回顾前面的定义：Lie 子代数是满足 \([\mathfrak h,\mathfrak h]\subseteq\mathfrak h\) 的线性子空间 \(\mathfrak h\subseteq\mathfrak g\)。这个条件可以在整个群上展开：把 \(\mathfrak h\) 从单位元左移到每个点，得到
\begin{equation}\label{g1:eq:subalgebra-left-distribution}
	D_g=(\mathrm{d}L_g)_e\mathfrak h,\qquad g\in G.
\end{equation}
这里 \(L_g(a)=ga\)，所以每个 \(D_g\) 都有相同维数。我们将证明括号封闭性正好保证这个分布可积。

\begin{proposition}[左不变积分分布]\label{g1:prop:lie-subalgebra-left-distribution}
	设 \(G\) 为有限维实 Lie 群，\(\mathfrak g=\mathrm{T}_eG\)，\(\mathfrak h\subseteq\mathfrak g\) 为 \(r\) 维 Lie 子代数。令 \(D_g=(\mathrm{d}L_g)_e\mathfrak h\)，即\eqref{g1:eq:subalgebra-left-distribution}中的分布。则 \(D\) 是秩 \(r\) 的光滑对合分布，所有左平移都保持此分布。对 \(X\in\mathfrak h\)，还有
	\[
		\operatorname{Ad}_{\exp(tX)}\mathfrak h=\mathfrak h\qquad(t\in\mathbb R).
	\]
\end{proposition}
\begin{proof}
	选 \(\mathfrak h\) 的基 \(X_1,\ldots,X_r\)，将其延成左不变向量场 \(X_1^L,\ldots,X_r^L\)。这些场光滑，且在每点线性无关并张成 \(D_g\)，故给出光滑分布。Lie 代数括号的定义及 \(\mathfrak h\) 的封闭性说明，每个 \([X_i^L,X_j^L]\) 都是这些基场的常系数线性组合。对一般局部截面 \(U=\sum_i f_iX_i^L\)、\(V=\sum_j h_jX_j^L\)，有
	\[
		[U,V]=\sum_{i,j}\left(f_ih_j[X_i^L,X_j^L]
		+f_iX_i^L(h_j)X_j^L-h_jX_j^L(f_i)X_i^L\right),
	\]
	仍取值于 \(D\)，所以分布对合。又因为 \(L_aL_g=L_{ag}\)，有 \(\mathrm{d}L_a(D_g)=D_{ag}\)，证明左平移不变性。

	最后固定 \(X,Y\in\mathfrak h\)。由\cref{g1:prop:exponential-naturality}，伴随作用与算子指数满足
	\[
		\operatorname{Ad}_{\exp(tX)}Y
		=e^{t\operatorname{ad}_X}Y
		=\sum_{j=0}^{\infty}\frac{t^j}{j!}\operatorname{ad}_X^jY.
	\]
	括号封闭性说明 \(\operatorname{ad}_X=[X,\cdot]\) 保持 \(\mathfrak h\)，所以级数的每个部分和均在 \(\mathfrak h\) 中。有限维线性子空间闭，级数极限也在其中。因此 \(\operatorname{Ad}_{\exp(tX)}\mathfrak h\subseteq\mathfrak h\)；把 \(t\) 换成 \(-t\)，再用两伴随算子互逆，得到反向包含。
\end{proof}

Frobenius 定理现在给出叶状坐标盒，在其中 \(D\) 由前 \(r\) 个坐标方向张成。后面的横向坐标在每张叶片上恒定。我们会反复使用这个局部事实来判断一个环境光滑映射何时也对叶自身的拓扑光滑。

\begin{lemma}[取值于叶的光滑映射]\label{g1:lem:leaf-valued-map}
	设 \(M,N\) 为有限维、Hausdorff、第二可数的无边界光滑流形，\(D\) 为 \(M\) 上的常秩光滑对合分布。若 \(N\) 连通，\(F:N\to M\) 光滑，且
	\[
		\mathrm{d}F_q(\mathrm{T}_qN)\subseteq D_{F(q)}
	\]
	对每个 \(q\in N\) 成立，则 \(F(N)\) 包含在一片最大叶中。将值域限制为该叶后，\(F\) 对叶的光滑结构仍然光滑。
\end{lemma}
\begin{proof}
	固定 \(q\in N\)，取连通的小坐标邻域 \(V\ni q\)，使 \(F(V)\) 落在一个叶状坐标盒中，盒内坐标记为 \((u,v)\)，其中 \(v\) 是横向坐标。由于 \(\mathrm{d}v\) 在 \(D\) 上为零，链式法则给出 \(\mathrm{d}(v\circ F)=0\)。在连通坐标小球内沿线段积分导数，便知 \(v\circ F\) 恒为一个常值。因此 \(F(V)\) 位于一张叶片中。

	在 \(N\) 中，把像位于同一最大叶的点归成一类。上面的局部结论说明这些类都是开集，它们两两不交并覆盖 \(N\)，故连通性使只有一个非空类。这证明全像位于一片叶。每个 \(F(V)\) 又位于该叶的一张开叶片中，叶片坐标下的 \(F\) 正是光滑函数 \(u\circ F\)，所以到叶的映射局部光滑，进而整体光滑。这里没有要求 \(F\) 是浸入，也没有把叶拓扑替换成环境子空间拓扑。
\end{proof}

\subsection{浸入子群的构造}
\label{g1:subsec:1603:02}

先说清要构造哪一种子群。一个子群 \(H\subseteq G\) 可以带有自己的 Lie 群结构；只要求包含映射光滑且微分单射时，得到的是浸入子群。如果还要求这份拓扑等于 \(G\) 在 \(H\) 上诱导的子空间拓扑，才得到嵌入子群。两种情形中的群乘法都必须对所用的流形结构光滑，不能从集合上的子群性质直接略过这一步。

\begin{definition}[浸入 Lie 子群]\label{g1:def:immersed-lie-subgroup}
	设 \(G\) 为有限维实 Lie 群。子群 \(H\subseteq G\) 若带有有限维、Hausdorff、第二可数的无边界 Lie 群结构，且包含映射 \(i:H\hookrightarrow G\) 为光滑浸入，则称为\term[jinruliziqun]{浸入 Lie 子群}[immersed Lie subgroup]。若包含映射还是光滑嵌入，则称为嵌入 Lie 子群。通过 \(\mathrm{d}i_e\) 将 \(\mathrm{T}_eH\) 识别为 \(\mathfrak g\) 的子空间。
\end{definition}

\begin{theorem}[Lie 子代数的积分]\label{g1:thm:lie-subalgebra-integration}
	设 \(G\) 为有限维实、Hausdorff、第二可数的无边界 Lie 群，\(\mathfrak h\subseteq\mathfrak g\) 为 Lie 子代数。存在唯一的连通浸入 Lie 子群 \(H\subseteq G\)，满足 \(\operatorname{Lie}(H)=\mathfrak h\)。其底集合及流形结构就是左不变分布 \(D_g=(\mathrm{d}L_g)_e\mathfrak h\) 经过 \(e\) 的最大叶，并且
	\begin{equation}\label{g1:eq:integral-subgroup-generated-exponentials}
		H=\langle\exp_G(\mathfrak h)\rangle.
	\end{equation}
	这里右侧表示由全部 \(\exp_G X\)、\(X\in\mathfrak h\) 生成的代数子群。
\end{theorem}
\begin{proof}
	由\cref{g1:prop:lie-subalgebra-left-distribution,g1:thm:maximal-connected-leaf}，经过 \(e\) 的最大叶 \(H\) 是连通 \(r\) 维流形，包含映射为浸入，且 \(\mathrm{T}_gH\) 在 \(\mathrm{T}_gG\) 中的像是 \(D_g\)。先提醒叶拓扑所需的两个可数与分离条件。它使每张局部叶片为开坐标片，包含映射连续；环境中的两个不交开邻域限制到叶上仍然开，所以 \(H\) Hausdorff。第二可数性则来自叶片构造：在 \(G\) 的可数个叶状坐标盒中，一张叶片与任意另一盒的交只有可数个连通分支，每个分支只落在那一盒的一张叶片中。沿有限相交链逐步加入叶片，每一步至多可数，故经过 \(e\) 的叶只需可数张叶片覆盖。取各片可数坐标基之并，得到叶的可数基。这正是最大叶定理中用于叶拓扑的论证，并非从环境子空间拓扑直接继承。

	左平移 \(L_a\) 保持 \(D\)，所以它把经过 \(e\) 的最大叶送到经过 \(a\) 的最大叶：若后者还有更大的连通积分流形，左移回来就会违背前者的最大性。若 \(a\in H\)，经过 \(a\) 的叶仍是 \(H\)，故 \(aH=H\)。于是 \(a,b\in H\) 蕴含 \(ab\in H\)，而 \(e\in aH\) 蕴含 \(a^{-1}\in H\)。这先证明 \(H\) 在集合上是子群。各 \(L_a|_H\)、\(a\in H\)，由\cref{g1:lem:leaf-valued-map}光滑，逆 \(L_{a^{-1}}|_H\) 同样光滑，所以已经是叶流形上的微分同胚。

	若 \(X\in\mathfrak h\)，曲线 \(t\mapsto\exp(tX)\) 的速度为 \(X^L\)，始终切于 \(D\)。它经过 \(e\)，故由取值于叶的映射引理，整条曲线在 \(H\) 中且对叶拓扑光滑。令 \(H'\) 是由这些曲线的值生成的子群，则 \(H'\subseteq H\)。为证明相等，选 \(\mathfrak h\) 的基 \(X_1,\ldots,X_r\)，考虑
	\[
		f(t^1,\ldots,t^r)=\exp(t^1X_1)\cdots\exp(t^rX_r).
	\]
	其值在 \(H'\) 中。记第 \(i\) 个因子之后的乘积为 \(b_i\)。对 \(t^i\) 求导，再用左平移送回单位元，所得方向为 \(\operatorname{Ad}_{b_i^{-1}}X_i\)。每个 \(b_i\) 都是 \(\exp(\mathfrak h)\) 中元素的乘积，上一命题说明它的伴随作用保持 \(\mathfrak h\)，故 \(\mathrm{d}f\) 的像包含于 \(D\)。取值于叶的映射引理使 \(f:\mathbb R^r\to H\) 光滑。在零点，它的微分把标准基送到 \(X_1,\ldots,X_r\)，所以到 \(\mathrm{T}_eH\) 为同构。逆函数定理说明 \(H'\) 包含 \(H\) 中单位元的一个开邻域。

	由 \(H\) 中的左平移已是微分同胚，\(H'\) 是 \(H\) 中的开子群；它的其他左陪集也开，故补集开。连通性迫使 \(H'=H\)，证明\eqref{g1:eq:integral-subgroup-generated-exponentials}。同时，生成元的伴随作用保持 \(\mathfrak h\)，于是
	\[
		\operatorname{Ad}_a\mathfrak h=\mathfrak h\qquad(a\in H).
	\]

	现在验证群运算在叶拓扑下光滑。对 \(a,b\in H\)、\(X,Y\in\mathfrak h\)，把切向量写成左平移的形式，乘法 \(\mu\) 的微分为
	\begin{equation}\label{g1:eq:subgroup-multiplication-differential}
		\mathrm{d}\mu_{(a,b)}\bigl(\mathrm{d}L_aX,\mathrm{d}L_bY\bigr)
		=\mathrm{d}L_{ab}\bigl(\operatorname{Ad}_{b^{-1}}X+Y\bigr)\in D_{ab}.
	\end{equation}
	这个公式可对曲线
	\[
		a\exp(tX)b\exp(tY)
		=ab\,\exp\bigl(t\operatorname{Ad}_{b^{-1}}X\bigr)\exp(tY)
	\]
	在零点求导得到。环境乘法组成的映射 \(H\times H\to G\) 光滑，微分又取值于 \(D\)，所以取值于叶的映射引理表明它作为 \(H\times H\to H\) 也光滑。对逆映射 \(\iota(a)=a^{-1}\)，同理有
	\[
		\mathrm{d}\iota_a(\mathrm{d}L_aX)
		=\mathrm{d}L_{a^{-1}}(-\operatorname{Ad}_aX)\in D_{a^{-1}},
	\]
	这来自 \((a\exp(tX))^{-1}=a^{-1}\exp(-t\operatorname{Ad}_aX)\)。因此逆映射也对叶结构光滑，\(H\) 确实成为 Lie 群。

	包含映射是光滑群同态，单位元处微分像为 \(\mathfrak h\)，所以其 Lie 代数恰为给定子代数。若 \(K\subseteq G\) 是另一个连通浸入 Lie 子群，且单位元切空间同样为 \(\mathfrak h\)，同态与左平移相容使它在每点的切空间像为 \(D\)。因此 \(K\) 是经过 \(e\) 的连通积分流形。最大叶定理给出包含 \(K\to H\)，并保证它为开嵌入。其像又是 \(H\) 的子群，故是开子群；由 \(H\) 连通，像等于 \(H\)。这个开嵌入遂为微分同胚，证明底集合和 Lie 群结构都唯一。若 \(r=0\)，最大叶只是 \(\{e\}\)，上述结论同样成立。
\end{proof}

浸入子群使用叶拓扑，群运算的光滑性由上面的三个性质保证：微分切于分布、横向坐标恒定，以及叶片本身是叶中的开集。对嵌入子群，则可直接调用\cref{g1:prop:embedded-initial-property}中的限制映射性质。

\subsection{闭子群定理}
\label{g1:subsec:1603:03}

现在反向开始：只给 \(G\) 的一个代数子群 \(K\)，暂时不赋予任何流形结构。若它闭，我们可以用全部留在 \(K\) 内的单参数子群识别候选切空间，再证明附近没有额外的横向部分。闭性将在两次取极限时起决定作用。

\begin{lemma}[小元素的极限方向]\label{g1:lem:closed-subgroup-limit-direction}
	设 \(K\) 是有限维实 Lie 群 \(G\) 的闭子群，给 \(\mathfrak g\) 任选 Euclid 范数。若 \(Z_j\ne0\)、\(Z_j\to0\)、\(\exp Z_j\in K\)，且 \(Z_j/\lVert Z_j\rVert\to Z\)，则
	\[
		\exp(tZ)\in K\qquad\text{对每个 }t\in\mathbb R.
	\]
\end{lemma}
\begin{proof}
	固定 \(t\in\mathbb R\)，令 \(\varepsilon_j=\lVert Z_j\rVert\)，选整数 \(m_j\) 使 \(|m_j\varepsilon_j-t|\le\varepsilon_j\)。因为 \(\varepsilon_j\to0\)，有
	\[
		m_jZ_j=(m_j\varepsilon_j)\frac{Z_j}{\varepsilon_j}\longrightarrow tZ.
	\]
	单参数子群规律给出
	\[
		(\exp Z_j)^{m_j}=\exp(m_jZ_j)\in K.
	\]
	指数映射连续，而 \(K\) 闭，故极限 \(\exp(tZ)\) 也属于 \(K\)。整数 \(m_j\) 可以为负，所以论证同时包含正、负时间；\(t\) 任意，结论成立。
\end{proof}

\begin{theorem}[闭子群定理]\label{g1:thm:closed-subgroup}
	设 \(G\) 为有限维实、Hausdorff、第二可数的无边界 Lie 群，\(K\subseteq G\) 为闭的代数子群。则 \(K\) 具有唯一的嵌入 Lie 子群结构，其 Lie 代数为
	\begin{equation}\label{g1:eq:closed-subgroup-lie-algebra}
		\mathfrak k=\{X\in\mathfrak g:\exp(tX)\in K\text{ 对所有 }t\in\mathbb R\}.
	\end{equation}
	对任意线性补空间 \(\mathfrak g=\mathfrak k\oplus\mathfrak m\)，可以选择零点的邻域 \(V\subseteq\mathfrak k\)、\(W\subseteq\mathfrak m\)，使
	\[
		\Psi:V\times W\longrightarrow U,\qquad \Psi(X,Y)=\exp X\exp Y
	\]
	为到 \(G\) 中单位元开邻域的微分同胚，并且
	\begin{equation}\label{g1:eq:closed-subgroup-local-slice}
		K\cap U=\Psi(V\times\{0\})=\exp V.
	\end{equation}
\end{theorem}
\begin{proof}
	先证明\eqref{g1:eq:closed-subgroup-lie-algebra}确实给出 Lie 子代数。对每个固定 \(t\)，\(X\mapsto\exp(tX)\) 连续，所以 \(\mathfrak k\) 是闭集 \(K\) 的这些原像之交，因而在 \(\mathfrak g\) 中闭。它包含零，并且对任意实数倍封闭，因为 \(t(cX)=(tc)X\)。若 \(X,Y\in\mathfrak k\)，由\cref{g1:cor:lie-product-limit}，对每个实数 \(t\) 有
	\[
		\exp(t(X+Y))
		=\lim_{n\to\infty}\bigl(\exp(tX/n)\exp(tY/n)\bigr)^n.
	\]
	右侧每一项都在子群 \(K\) 中，闭性使极限也在 \(K\) 中。因此 \(X+Y\in\mathfrak k\)，\(\mathfrak k\) 是闭线性子空间。

	对 \(k\in K\)，\cref{g1:prop:exponential-naturality}说明指数与共轭相容，故
	\[
		\exp(t\operatorname{Ad}_kX)=k\exp(tX)k^{-1}\in K
		\qquad(X\in\mathfrak k).
	\]
	于是 \(\operatorname{Ad}_k\mathfrak k=\mathfrak k\)，反向包含使用 \(k^{-1}\)。若 \(X,Y\in\mathfrak k\)，则 \(\exp(sX)\in K\)，所以 \(\operatorname{Ad}_{\exp(sX)}Y\) 始终位于闭线性空间 \(\mathfrak k\)。在零点取差商极限，得到
	\[
		[X,Y]=\left.\frac{\mathrm{d}}{\mathrm{d}s}\right|_{s=0}
		\operatorname{Ad}_{\exp(sX)}Y\in\mathfrak k.
	\]
	这证明其括号封闭性。

	取补空间 \(\mathfrak m\)。我们先证明存在 \(\delta>0\)，使
	\[
		Y\in\mathfrak m,\quad \lVert Y\rVert<\delta,\quad \exp Y\in K
		\quad\Longrightarrow\quad Y=0.
	\]
	若不存在，就可取非零 \(Y_j\in\mathfrak m\)，满足 \(Y_j\to0\)、\(\exp Y_j\in K\)。\(\mathfrak m\) 的单位球面紧，故取子列后 \(Y_j/\lVert Y_j\rVert\to Y\in\mathfrak m\)，且 \(\lVert Y\rVert=1\)。由\cref{g1:lem:closed-subgroup-limit-direction}，\(\exp(tY)\in K\) 对所有 \(t\) 成立，即 \(Y\in\mathfrak k\)。这与 \(\mathfrak k\cap\mathfrak m=\{0\}\) 矛盾。若 \(\mathfrak m=\{0\}\)，这一断言直接成立。

	考虑光滑映射 \(\Psi(X,Y)=\exp X\exp Y\)。\cref{g1:prop:exponential-local-diffeomorphism}给出 \(\mathrm{d}(\exp)_0=\operatorname{id}_{\mathfrak g}\)，群乘法在 \((e,e)\) 的微分是向量相加，故
	\[
		\mathrm{d}\Psi_{(0,0)}(X,Y)=X+Y.
	\]
	直和分解使它可逆。逆函数定理给出零点附近的局部微分同胚；再把定义域缩成乘积 \(V\times W\)，并使 \(W\) 包含于上面的 \(\delta\) 球中。这个乘积的像 \(U\) 仍在 \(G\) 中开。

	若 \(k\in K\cap U\)，唯一写成 \(k=\exp X\exp Y\)，其中 \(X\in V\)、\(Y\in W\)。由于 \(X\in\mathfrak k\)，有 \(\exp X\in K\)，所以
	\[
		\exp Y=(\exp X)^{-1}k\in K.
	\]
	横向排除结论强制 \(Y=0\)。反过来，\(X\in V\) 时 \(\exp X\in K\cap U\)。这就证明了完整的切片等式\eqref{g1:eq:closed-subgroup-local-slice}；它排除了附近所有额外分支，而不仅仅给出一个包含于 \(K\) 的参数片。

	将这张切片用 \(L_k\)、\(k\in K\)，移到每个子群点附近，得到 \(K\) 的嵌入子流形结构。它取环境子空间拓扑，Hausdorff 性和第二可数性从 \(G\) 继承；切片坐标转换由环境坐标转换限制得到，故光滑。群乘法与逆映射在 \(G\) 上光滑，其限制值域在 \(K\) 中，\cref{g1:prop:embedded-initial-property}因而保证它们在这个子流形结构下光滑。所以 \(K\) 是嵌入 Lie 子群。单位元的切片为 \(X\mapsto\exp X\)、\(X\in V\)，其微分为包含 \(\mathfrak k\hookrightarrow\mathfrak g\)，故 Lie 代数恰为 \(\mathfrak k\)。最后，嵌入子流形的光滑结构由包含映射的初始性质唯一决定，证明唯一性。
\end{proof}

当 \(\mathfrak k=\{0\}\) 时，局部切片只是单位元，左移以后每个点都孤立，所以 \(K\) 是离散子群。当 \(\mathfrak k=\mathfrak g\) 时，切片占满单位元的一个开邻域，所以 \(K\) 是开子群。两种极端都包含在同一证明中。

\ProseParagraphBreak

Lie 代数只描述单位元所在的连通部分，不能确定全部离散分支。例如加法群 \(\mathbb R\) 中，\(\{0\}\) 与 \(\mathbb Z\) 都是闭子群，二者 Lie 代数都为零，但它们不是同一个子群。因此前面的积分唯一性明确要求连通。对一般闭子群 \(K\)，单位元连通分支 \(K^0\) 是开连通子群：流形局部路径连通使连通分支开，乘法与逆映射又把含单位元的连通集送入同一分支。它具有与 \(K\) 相同的单位元切空间，所以正是 \(\mathfrak k\) 的唯一连通积分子群。

\subsection{闭性与嵌入性的关系}
\label{g1:subsec:1603:04}

一般嵌入子流形不必闭，例如平面中的开线段。然而子群的平移结构会把单位元附近的局部闭性传播到整个群。另一方面，如果一个已有自身拓扑的浸入 Lie 子群碰巧具有闭像，还须说明原拓扑与闭子群定理给出的拓扑一致，才能说原包含映射是嵌入。

\begin{proposition}[闭性与嵌入性]\label{g1:prop:immersed-subgroup-closed-embedded}
	设 \(H\) 为有限维实 Lie 群 \(G\) 的浸入 Lie 子群，两个 Lie 群都为 Hausdorff、第二可数。则包含映射 \(H\hookrightarrow G\) 是嵌入，当且仅当子集 \(H\) 在 \(G\) 中闭。
\end{proposition}
\begin{proof}
	先设包含映射是嵌入。由\cref{g1:prop:embedding-slice-characterization}，存在 \(G\) 中单位元的开邻域 \(U\)，使 \(H\cap U\) 是一张坐标切片，因而在 \(U\) 中闭。令 \(K=\overline H\)。乘法与逆映射连续，故 \(K\) 仍是子群：逼近 \(a,b\in K\) 的子群点列，其乘积及逆分别逼近 \(ab,a^{-1}\)，这些极限仍在闭包中。因为 \(U\) 开，局部闭性给出 \(K\cap U=H\cap U\)。于是 \(H\) 包含 \(K\) 中单位元的开邻域，故是 \(K\) 的开子群。每个左陪集也开，所以 \(H\) 的补集在 \(K\) 中开，即 \(H\) 又在 \(K\) 中闭。它在 \(K=\overline H\) 中稠密，只能有 \(H=K\)，证明闭性。

	反过来，假设像 \(H\) 在 \(G\) 中闭，令 \(K\) 表示同一个集合上由闭子群定理给出的嵌入 Lie 群结构。原包含映射通过 \(K\) 分解为一个双射光滑群同态 \(j:H\to K\)，光滑性使用嵌入子流形的初始性质。需要证明 \(j^{-1}\) 连续，而不是把这一点当作双射的自动结论。

	给定 \(H\) 中单位元的任意开邻域 \(O\)，可取较小开邻域 \(V\)，使 \(C=\overline V\) 紧且 \(C^{-1}C\subseteq O\)。具体可先由乘法与逆的连续性取 \(W\) 满足 \(W^{-1}W\subseteq O\)，再在坐标图中取闭包包含于 \(W\) 的小球作为 \(V\)。第二可数性使覆盖 \(\{hV:h\in H\}\) 有可数子覆盖，故对某列 \(h_\ell\) 有
	\[
		K=\bigcup_{\ell=1}^{\infty}j(h_\ell)\,j(C).
	\]
	每个 \(j(C)\) 都紧，在 Hausdorff 流形 \(K\) 中闭。它不能没有内点。否则右侧是可数个闭的无内点集；在 \(K\) 的一张坐标图中先取一个具有非空内部的紧闭球，再逐次在前一闭球的内部选一个更小的非空闭球，避开第 \(\ell\) 个闭集。这些嵌套紧集有公共点，而该点不属于任何一个所列闭集，与覆盖矛盾。零维时闭的无内点集只能为空，结论直接成立。

	因此可取 \(K\) 中非空开集 \(W_0\subseteq j(C)\) 及 \(k_0\in W_0\)。则
	\[
		k_0^{-1}W_0\subseteq j(C)^{-1}j(C)
		=j(C^{-1}C)\subseteq j(O).
	\]
	左侧是单位元的开邻域，所以 \(j(O)\) 包含单位元邻域。在 \(O\) 的每个点先平移一个更小的单位元邻域，并应用相同论证，可知 \(j(O)\) 本身开。再用群平移，\(j\) 把任意开集送到开集。于是双射 \(j\) 为同胚，原 \(H\) 的拓扑等于 \(G\) 的子空间拓扑。原包含映射又已知是浸入，故按定义为光滑嵌入。
\end{proof}

这个证明中的第二可数性用于取得可数紧集覆盖；紧性和环境的 Hausdorff 性用于把它们的像看成闭集。它们共同排除了闭像还带有一套更细 Lie 群拓扑的可能。对于一般的单射浸入映射，没有群平移结构时，不能照搬这个结论。

\begin{example}[无理斜率子群]\label{g1:exm:irrational-torus-subgroup}
	在环面 Lie 群 \(T^2=\mathbb R^2/\mathbb Z^2\) 中，固定无理数 \(\alpha\)，定义
	\[
		f_\alpha:\mathbb R\longrightarrow T^2,\qquad
		f_\alpha(t)=[t,\alpha t].
	\]
	它是光滑群同态，微分像为直线 \(\mathfrak h_\alpha=\mathbb R(1,\alpha)\)。若 \(f_\alpha(t)=0\)，则 \(t\in\mathbb Z\) 且 \(\alpha t\in\mathbb Z\)，无理性迫使 \(t=0\)，所以它单射；微分处处非零，故为浸入。将像 \(H_\alpha\) 赋予从 \(\mathbb R\) 传来的 Lie 群结构，便得到 \(\mathfrak h_\alpha\) 的连通积分子群。

	\cref{g1:prop:linear-torus-leaves}已经证明这条轨道在环面中稠密。其关键是整数倍 \(\alpha\) 模 \(1\) 稠密：对任意指定的第一坐标 \(u\)，时刻 \(u+n\) 的第二坐标为 \(\alpha u+n\alpha\)，因此能够进入任意给定小矩形。这个像却不是整个环面，因为它与第一坐标为零的圆周只相交于可数个点 \([0,n\alpha]\)。于是 \(H_\alpha\) 是真稠密子群，既不闭也不嵌入。

	拓扑差异还可直接检测。前述整数逼近给出 \(|n_j|\to\infty\)，但 \(n_j\alpha\bmod1\to0\)，所以 \(f_\alpha(n_j)\to f_\alpha(0)\) 在环面中成立；在 \(H_\alpha\) 自身的参数坐标里，\(n_j\) 却不趋于零。这正是从子空间像返回叶参数的逆映射不连续。子群自身始终是一维 Lie 群，闭包却是二维环面，其 Lie 代数也从一条直线增大为整个 \(\mathbb R^2\)。
\end{example}

因此，Lie 子代数的积分定理提供连通浸入子群，闭子群定理提供采用子空间拓扑的嵌入子群，两者处理的条件不同。后面构造齐性空间时，将用闭性把子群附近的横向坐标组织为陪集空间的光滑坐标。

% 纲要条目排为编号小节标题，并进入目录与书签；说明段仍保留于纲要注释中。

% !TeX root = ../../Book1.tex
% 纲要标题：### 16.4 同态积分及覆盖群
% 纲要位置：几何学计划纲要.md，第 758 行。
% 正文状态：初稿；尚未正式审核。

% 纲要内容开始
% 
% * 局部同态的积分构造
% * 沿路径延拓和同伦无关性
% * 单连通源群上的整体同态
% * 连通 Lie 群的万有覆盖群
% * 离散中心商与同 Lie 代数的不同 Lie 群
% 
% ---
% 
% 纲要内容结束

\section{同态积分及覆盖群}
\label{g1:sec:1604}

Lie 群同态在单位元的微分必为 Lie 代数同态，见\cref{g1:prop:lie-homomorphism-differential}。反过来，一个保持括号的线性映射，能否来自某个群同态？局部问题由微分方程的相容性解决；整体问题还要检查，沿不同路径把这个局部映射继续下去，是否会得到相同结果。圆周上“角度乘以 \(\lambda\)”的例子已经提示：局部允许任意实数 \(\lambda\)，绕完一周后的闭合要求却可能排除其中大部分。

\ProseParagraphBreak

本节始终讨论有限维实 Lie 群，底流形无边界、Hausdorff 且第二可数。记 \(e_G,e_H\) 为单位元，\(\mathfrak g=\mathrm{T}_{e_G}G\)、\(\mathfrak h=\mathrm{T}_{e_H}H\) 为相应 Lie 代数。需要的积分工具是\cref{g1:thm:frobenius,g1:lem:local-integral-plaque}：对合分布局部有叶片，且位于同一叶状坐标盒内的连通积分流形只能处于一张叶片中。整体延拓则使用第五章的连续路径与定端点同伦，拼接 \(\alpha*\beta\) 仍然表示先走 \(\alpha\)，后走 \(\beta\)。

\subsection{局部同态的积分构造}
\label{g1:subsec:1604:01}

\begin{definition}
	\label{g1:def:local-lie-homomorphism}
	设 \(G,H\) 为有限维实 Lie 群，\(U\subset G\) 是单位元的开邻域。光滑映射 \(f:U\to H\) 若满足 \(f(e_G)=e_H\)，并且
	\[
		f(xy)=f(x)f(y)\qquad\text{只要 }x,y,xy\in U,
	\]
	则称为从 \(G\) 到 \(H\) 的\term[jubuliquntongtai]{局部 Lie 群同态}[local Lie group homomorphism]。两个这样的映射局部相同，是指它们在某个共同的单位元邻域上相同。
\end{definition}

要从线性映射 \(\varphi:\mathfrak g\to\mathfrak h\) 找到 \(f\)，可以先问它的图像应有什么切空间。若 \(f\) 是同态，固定 \(g\) 后对 \(f(gx)=f(g)f(x)\) 在 \(x=e_G\) 求导，便得到
\[
	\mathrm{d}f_g\circ\mathrm{d}(L_g)_{e_G}
	=\mathrm{d}(L_{f(g)})_{e_H}\circ\varphi.
\]
因此，图像在 \((g,f(g))\) 处的切向量应当把同一个 \(X\in\mathfrak g\) 送成一对左平移向量。这个要求直接给出待积分的分布。

\begin{theorem}[局部同态积分]
	\label{g1:thm:local-lie-homomorphism}
	设 \(G,H\) 为有限维实 Lie 群，\(\varphi:\mathfrak g\to\mathfrak h\) 为 Lie 代数同态。则存在单位元的连通对称开邻域 \(U\subset G\)，以及局部 Lie 群同态 \(f:U\to H\)，满足 \(\mathrm{d}f_{e_G}=\varphi\)。这样的局部同态在单位元附近唯一。
\end{theorem}
\begin{proof}
	在 \(G\times H\) 上定义秩 \(\dim G\) 的光滑分布
	\begin{equation}
		\label{g1:eq:homomorphism-graph-distribution}
		D_{(g,h)}
		=\left\{\left(\mathrm{d}(L_g)_{e_G}X,\,
		\mathrm{d}(L_h)_{e_H}\varphi X\right):X\in\mathfrak g\right\}.
	\end{equation}
	取 \(\mathfrak g\) 的基 \(X_1,\ldots,X_m\)。对应向量场
	\(Z_i=(X_i^L,(\varphi X_i)^L)\)
	处处线性无关，并张成 \(D\)，因为其第一分量已经线性无关。由\cref{g1:prop:left-invariant-fields,g1:def:lie-algebra-group}以及 \(\varphi\) 保括号，
	\[
		[Z_i,Z_j]
		=\left([X_i,X_j]^L,[\varphi X_i,\varphi X_j]^L\right)
		=\left([X_i,X_j]^L,(\varphi[X_i,X_j])^L\right)
	\]
	仍取值于 \(D\)。一般局部截面是这些 \(Z_i\) 的光滑函数线性组合；括号乘法法则将其括号展开为上述括号与 \(Z_i\) 的函数线性组合，所以 \(D\) 对合。

	由 Frobenius 定理，\((e_G,e_H)\) 附近有叶状坐标盒 \(\Omega\)，并有经过该点的中央叶片 \(P\)。投影 \(\pi_G:G\times H\to G\) 的微分在 \(D\) 上为同构。因此由\cref{g1:thm:inverse-function}，缩小中央叶片后，\(\pi_G|_P\) 是到单位元某开邻域 \(U_0\) 的微分同胚。可以同时缩小叶状坐标盒，使它的整个中央叶片就是这一小片。于是
	\[
		P=\{(g,f(g)):g\in U_0\}
	\]
	定义光滑映射 \(f:U_0\to H\)，且 \(f(e_G)=e_H\)。图像切空间等于 \(D\)，给出
	\begin{equation}
		\label{g1:eq:local-homomorphism-differential}
		\mathrm{d}f_g\circ\mathrm{d}(L_g)_{e_G}
		=\mathrm{d}(L_{f(g)})_{e_H}\circ\varphi.
	\end{equation}
	特别地，在单位元得到 \(\mathrm{d}f_{e_G}=\varphi\)。

	还须证明图像尊重乘法。分布 \(D\) 在 \(G\times H\) 的每个左平移下不变。由乘法及 \(f\) 的连续性，可以选连通对称单位元邻域 \(U\subset U_0\)，使 \(U^2\subset U_0\)，并使
	\[
		(ab,f(a)f(b))\in\Omega\qquad(a,b\in U).
	\]
	这里是对全部 \(a,b\in U\) 的共同缩小：映射 \((a,b)\mapsto(ab,f(a)f(b))\) 在 \((e_G,e_G)\) 取值于开集 \(\Omega\)，故其原像包含一个乘积邻域，再在两个因子的交中取小的连通对称邻域即可。

	固定 \(a\in U\)。映射 \(b\mapsto(ab,f(a)f(b))\)、\(b\in U\)，是中央叶片一部分的左平移，因而是连通积分流形，且全部位于 \(\Omega\) 内。它在 \(b=e_G\) 时经过 \((a,f(a))\in P\)。由\cref{g1:lem:local-integral-plaque}，它只能位于同一张中央叶片 \(P\)。投影到第一分量，便得到
	\[
		f(ab)=f(a)f(b)\qquad(a,b\in U).
	\]
	左边在 \(U_0\) 中有定义；将 \(f\) 限制到 \(U\) 后，尤其得到定义中要求的、当 \(a,b,ab\in U\) 时的局部乘法律。

	最后设 \(\widehat f\) 是另一个微分为 \(\varphi\) 的局部同态。对固定 \(g\) 及充分接近单位元的 \(x\)，对 \(\widehat f(gx)=\widehat f(g)\widehat f(x)\) 求导，得到与\eqref{g1:eq:local-homomorphism-differential}相同的等式。因此 \(\widehat f\) 的图像也是 \(D\) 的积分流形。取足够小的连通定义域，使图像落在 \(\Omega\) 中；它经过 \((e_G,e_H)\)，故也位于中央叶片 \(P\)。中央叶片已经是 \(f\) 的图像，所以两映射在这个共同邻域内相同。
\end{proof}

保持括号的条件在这里承担了明确的相容性责任：它使图分布对合，因而存在所需局部图像。局部唯一性则意味着，后面沿路径延拓时不会在每个小邻域重新引入任意选择。此处得到的仍是局部同态；要把定义域扩展到整个群，必须进一步比较路径。

\subsection{沿路径延拓和同伦无关性}
\label{g1:subsec:1604:02}

设 \(f:U\to H\) 是上面得到的局部同态，\(U=U^{-1}\) 为单位元邻域。从点 \(g\) 走到附近的点 \(g'\)，其左增量是 \(g^{-1}g'\)。只要增量足够小，就可以用 \(f(g^{-1}g')\) 记录目标群中的相应变化。沿整条路径把这些小变化相乘，便是局部同态继续到远处的候选值。

\begin{lemma}[局部同态的路径延拓]
	\label{g1:lem:local-hom-path-extension}
	设 \(G,H\) 为有限维实 Lie 群，\(U=U^{-1}\subset G\) 是单位元开邻域，\(f:U\to H\) 为局部同态。对任意连续路径 \(\gamma:[0,1]\to G\)，可以选有限分割 \(0=t_0<\cdots<t_N=1\)，使
	\[
		\gamma(s)^{-1}\gamma(t)\in U
		\quad\text{对同一小区间内的所有 }s,t\text{ 成立}.
	\]
	定义
	\begin{equation}
		\label{g1:eq:local-hom-path-product}
		J_f(\gamma)
		=f(\gamma(t_0)^{-1}\gamma(t_1))\cdots
		f(\gamma(t_{N-1})^{-1}\gamma(t_N)).
	\end{equation}
	则此值与这样的分割无关，并在连续定端点同伦下不变。对另一连续路径 \(\eta:[0,1]\to G\)，在 \(\gamma(1)=\eta(0)\) 时令 \(\gamma*\eta\) 表示先走 \(\gamma\) 后走 \(\eta\) 的拼接，并令 \(\bar\gamma(t)=\gamma(1-t)\)。它还满足
	\[
		J_f(\gamma*\eta)=J_f(\gamma)J_f(\eta),\qquad
		J_f(\bar\gamma)=J_f(\gamma)^{-1},\qquad
		J_f(a\gamma)=J_f(\gamma),
	\]
	其中 \((a\gamma)(t)=a\,\gamma(t)\) 表示固定群元素 \(a\) 对路径的左平移。常路径的值为 \(e_H\)。若两个局部同态在单位元附近相同，则它们给出相同的路径值。
\end{lemma}
\begin{proof}
	先取对称单位元邻域 \(W\)，使 \(W^{-1}W\subset U\)。各左平移 \(aW\) 覆盖 \(G\)，同一个 \(aW\) 内任意两点的相对增量属于 \(W^{-1}W\)。将这个覆盖沿 \(\gamma\) 拉回紧区间，由\cref{g1:lem:lebesgue-number}选择有限细分，使每个小区间的像都位于某个 \(aW\) 中，便得到所需分割。

	在一个已有小区间中增加分点 \(s\)，原增量分解为
	\[
		\gamma(t_{i-1})^{-1}\gamma(t_i)
		=\bigl(\gamma(t_{i-1})^{-1}\gamma(s)\bigr)
		\bigl(\gamma(s)^{-1}\gamma(t_i)\bigr).
	\]
	两个因子及其乘积都在 \(U\) 中，所以局部乘法律使原来的一个 \(f\) 因子等于细分后的两个因子。反复加入分点不改变总乘积。任意两种分割有共同细分，而且细分仍满足整段增量条件，故定义无关。拼接对应接上两组分割，反向对应逆序并逐项取逆；这里 \(f(x^{-1})=f(x)^{-1}\) 来自 \(xx^{-1}=e_G\)。左平移不改变增量，因为 \((a\gamma(s))^{-1}(a\gamma(t))=\gamma(s)^{-1}\gamma(t)\)。这些事实给出所有运算规律。

	现在取连续定端点同伦 \(K:[0,1]^2\to G\)。仍用左平移 \(aW\) 的原像覆盖紧方形，取足够细的有限矩形网格，使每格的像落在某个 \(aW\) 中。若一格四个顶点的像依次为左下 \(a\)、右下 \(b\)、左上 \(c\)、右上 \(d\)，其边路径的 \(J_f\) 值分别是
	\[
		h=f(a^{-1}b),\quad v'=f(b^{-1}d),\quad
		v=f(a^{-1}c),\quad h'=f(c^{-1}d).
	\]
	因为 \(a^{-1}b,b^{-1}d,a^{-1}d\) 都在 \(U\)，以及另一条边路的三个增量也都在 \(U\)，局部乘法律给出
	\begin{equation}
		\label{g1:eq:local-hom-homotopy-cell}
		hv'=f(a^{-1}d)=vh'.
	\end{equation}
	这个等式允许按原来的顺序把经过一格下边、右边的乘积换成左边、上边的乘积。

	具体地，记网格第 \(j\) 行的水平边值为 \(h_{0j},\ldots,h_{m-1,j}\)，竖边值为 \(v_{0j},\ldots,v_{mj}\)，整行乘积为 \(r_j=h_{0j}\cdots h_{m-1,j}\)。逐格应用\eqref{g1:eq:local-hom-homotopy-cell}得到
	\[
		r_jv_{mj}=v_{0j}r_{j+1}.
	\]
	同伦固定端点，所以左右外边的值 \(v_{0j},v_{mj}\) 都为 \(e_H\)，故 \(r_j=r_{j+1}\)。逐行比较便得到两端路径具有相同的 \(J_f\) 值。整个论证没有交换不相邻因子，也没有要求 \(K\) 光滑。

	若 \(f,\widehat f\) 在更小的单位元邻域 \(U'\) 上相同，用 \(U'\) 再作一次上述细分，每个因子都相同。由已经证明的细分独立性，得到 \(J_f=J_{\widehat f}\)。
\end{proof}

这个证明与\cref{g1:thm:flat-path-independence}使用同一种有限网格方法。此处的 \(J_f\) 是按行程顺序相乘的群元素，故拼接公式也是按这个顺序书写；上一章的 \(P_\gamma\) 是作用于向量的线性映射，其函数复合按另一种书写次序进行。记住所比较的对象，便能核对每个公式中的顺序。

\subsection{单连通源群上的整体同态}
\label{g1:subsec:1604:03}

\begin{theorem}[Lie 群同态的积分]
	\label{g1:thm:lie-homomorphism-integration}
	设 \(G,H\) 为有限维实 Lie 群，\(G\) 连通且单连通。每个 Lie 代数同态 \(\varphi:\mathfrak g\to\mathfrak h\) 都唯一地来自一个光滑群同态 \(F:G\to H\)，即 \(\mathrm{d}F_{e_G}=\varphi\)。唯一性只需要源群连通：两个从同一连通 Lie 群出发的光滑群同态，若在单位元有相同微分，则处处相同。
\end{theorem}
\begin{proof}
	由\cref{g1:thm:local-lie-homomorphism}取得局部同态 \(f:U\to H\)。对 \(g\in G\)，取从 \(e_G\) 到 \(g\) 的连续路径 \(\gamma\)，定义 \(F(g)=J_f(\gamma)\)。单连通性包括道路连通性，所以路径存在；同端点的路径定端点同伦，\cref{g1:lem:local-hom-path-extension}因而保证定义与路径无关。常路径给出 \(F(e_G)=e_H\)。

	先证明 \(F\) 光滑。取足够小的单位元坐标球 \(V\subset U\)，使 \(V^{-1}V\subset U\)，并使坐标把 \(e_G\) 送到原点。对 \(x\in V\)，用坐标中的径向线段给出路径 \(\kappa_x:[0,1]\to V\)，从 \(e_G\) 走到 \(x\)。这一整段满足增量条件，所以 \(J_f(\kappa_x)=f(x)\)。固定 \(g\) 及一条从 \(e_G\) 到 \(g\) 的路径，在其后接上 \(g\kappa_x\)，便有
	\[
		F(gx)=F(g)f(x),\qquad x\in V.
	\]
	右边对 \(x\) 光滑，故 \(F\) 在每个 \(gV\) 上光滑。在 \(g=e_G\) 时得到 \(F|_V=f|_V\)，所以 \(\mathrm{d}F_{e_G}=\varphi\)。

	对 \(g,h\in G\)，取路径 \(\alpha:e_G\to g\)、\(\beta:e_G\to h\)。路径 \(\alpha*(g\beta)\) 从 \(e_G\) 到 \(gh\)，由路径值的拼接与左平移规律，
	\[
		F(gh)=J_f(\alpha*(g\beta))
		=J_f(\alpha)J_f(g\beta)=F(g)F(h).
	\]
	因此 \(F\) 为群同态。

	唯一性由\cref{g1:cor:lie-homomorphism-uniqueness}得到。具体地，单位元微分相同使两同态在所有指数上一致，而连通群的每个元素都是有限个指数的乘积，故两同态处处相同。这部分只需要源群连通，不需要再次使用单连通性。
\end{proof}

\begin{corollary}[环路的积分障碍]
	\label{g1:cor:lie-integration-loop-criterion}
	设 \(G,H\) 为有限维实 Lie 群，\(G\) 连通，\(\varphi:\operatorname{Lie}(G)\to\operatorname{Lie}(H)\) 为 Lie 代数同态。取单位元对称开邻域 \(U\subseteq G\) 及光滑局部同态 \(f:U\to H\)，满足 \((\mathrm{d}f)_{e_G}=\varphi\)。对连续路径 \(\gamma:[0,1]\to G\)，取使同段增量 \(\gamma(s)^{-1}\gamma(t)\) 均在 \(U\) 内的有限分割 \(0=t_0<\cdots<t_N=1\)，按时间顺序定义
	\[
		J_f(\gamma)=\prod_{j=1}^N f\bigl(\gamma(t_{j-1})^{-1}\gamma(t_j)\bigr).
	\]
	约定 \(\pi_1(G,e_G)\) 的乘法先走左边的环路，再走右边的环路。则 \(\varphi\) 能积分成 \(G\to H\) 的整体同态，当且仅当
	\[
		J_f(\gamma)=e_H
		\qquad\text{对每条基于 }e_G\text{ 的连续环路 }\gamma\text{ 成立}.
	\]
	映射 \([\gamma]\mapsto J_f(\gamma)\) 本身是 \(\pi_1(G,e_G)\to H\) 的群同态，并只依赖 \(\varphi\)。
\end{corollary}
\begin{proof}
	同伦不变性保证映射在环路类上良定义，拼接式保证它是群同态，局部积分的唯一性则保证构造不依赖局部积分的选择。若已有整体同态 \(F\)，对任意细分使用局部相同及同态性质，路径乘积逐项消去后等于
	\(F(\gamma(0))^{-1}F(\gamma(1))\)，所以环路值为单位元。

	反之，若全部环路值平凡，而 \(\gamma,\eta\) 都从 \(e_G\) 到 \(g\)，则
	\[
		e_H=J_f(\gamma*\bar\eta)=J_f(\gamma)J_f(\eta)^{-1},
	\]
	故两条路径给出相同值。上一证明中 \(F(g)=J_f(\gamma)\) 的构造及其光滑性、乘法性质便仍然适用。
\end{proof}

\begin{example}
	\label{g1:exm:circle-lie-integration-obstruction}
	将 \(S^1=\{e^{it}:t\in\mathbb R\}\) 的 Lie 代数按单位角速度识别为 \(\mathbb R\)。因为括号为零，每个线性映射 \(X\mapsto\lambda X\)、\(\lambda\in\mathbb R\)，都是 Lie 代数同态。其局部积分是在一个小角度区间内定义的
	\[
		f(e^{it})=e^{i\lambda t}.
	\]
	沿正向一周环路细分，小角度增量之和为 \(2\pi\)，所以路径乘积为 \(e^{2\pi i\lambda}\)。由\cref{g1:thm:circle-fundamental}，任意环路由其整数绕数确定；绕数 \(k\) 的路径值为 \(e^{2\pi i\lambda k}\)。因此整体积分存在当且仅当 \(\lambda\in\mathbb Z\)，此时就是 \(F(z)=z^\lambda\)。

	若把源群改成加法群 \(\mathbb R\)，则 \(t\mapsto e^{i\lambda t}\) 对每个实数 \(\lambda\) 都是整体同态。两种源群有同一个一维 Lie 代数；差别是圆周把相差 \(2\pi\) 的角度识别了，目标值也必须尊重这个识别。
\end{example}

\subsection{连通 Lie 群的万有覆盖群}
\label{g1:subsec:1604:04}

对于不是单连通的 \(G\)，可以先把路径的同伦历史保留下来。第五章的道路类万有覆盖已经这样做了；现在还须证明它能够承载群乘法，并且所得结构确实是通常意义下的 Lie 群。

\begin{theorem}[万有覆盖群]
	\label{g1:thm:universal-cover-lie-group}
	设 \(G\) 为连通有限维实 Lie 群。以 \(e_G\) 为基点的道路类万有覆盖
	\(p:(\widetilde G,\widetilde e)\to(G,e_G)\)
	具有以 \(\widetilde e\) 为单位元的唯一光滑 Lie 群结构，使 \(p\) 为光滑覆盖同态，并在每点为局部微分同胚。其核 \(\ker p\) 为离散中心子群，而且自然同构于 \(\pi_1(G,e_G)\)。特别地，连通 Lie 群的基本群为交换群。这个带群结构的覆盖称为\term[wanyoufugaiqun]{万有覆盖群}[universal covering group]。
\end{theorem}
\begin{proof}
	连通 Lie 群道路连通，而坐标小球可缩，所以第五章\cref{g1:thm:path-class-universal-cover}适用。其点为从 \(e_G\) 出发的连续路径的定端点同伦类，投影取终点，总空间已经知道连通且单连通。对路径 \(\alpha,\beta\) 定义逐点乘积 \((\alpha\cdot\beta)(t)=\alpha(t)\beta(t)\)，并令 \(\alpha^\iota(t)=\alpha(t)^{-1}\)。置
	\begin{equation}
		\label{g1:eq:universal-cover-group-product}
		[\alpha][\beta]=[\alpha\cdot\beta],\qquad
		[\alpha]^{-1}=[\alpha^\iota],\qquad
		\widetilde e=[e_{e_G}].
	\end{equation}
	若改变路径代表，逐点相乘或取逆相应的同伦，就得到新路径之间的定端点同伦，所以运算良定义。结合律、单位律和逆元律由逐点的群公理得到，终点映射 \(p\) 也立即保持乘法。此处 \(\alpha^\iota\) 是逐点取群逆，不是把路径参数反向；它仍从单位元出发。

	为了比较第五章的路径拼接，记 \(a=\alpha(1)\)。有
	\begin{equation}
		\label{g1:eq:pointwise-product-concatenation}
		[\alpha\cdot\beta]=[\alpha*(a\beta)].
	\end{equation}
	事实上，映射 \(Q(s,t)=\alpha(s)\beta(t)\) 将方形对角线送成逐点乘积路径，将先沿底边、再沿右边的折线路径送成 \(\alpha*(a\beta)\)。这两条方形内路径的端点都是 \((0,0),(1,1)\)；按坐标作线性插值给出定端点同伦，再复合 \(Q\) 即得所需同伦。\cref{g1:fig:cover-path-product}显示了这两种行进方式。

	接着核对局部光滑结构。取单位元的一个坐标小球 \(U\)，使其坐标像为以零为中心的凸球。记 \(\kappa_x(t)\) 为坐标径向线段从 \(e_G\) 到 \(x\) 的路径。对 \(q=[\alpha]\)、\(a=p(q)\)，
	\[
		x\longmapsto q_x=[\alpha*(a\kappa_x)],\qquad x\in U,
	\]
	正是 \(aU\) 上经过 \(q\) 的一张覆盖片的局部逆，因为在可缩邻域 \(aU\) 内，追加路径的终点决定其路径类。它与\cref{g1:lem:path-class-basis}的拓扑一致。把 \(G\) 的坐标通过 \(p\) 提升到每张片，就得到相容的局部 Euclid 坐标；交叠处的转换由底空间坐标转换给出。

	乘法在这些坐标中光滑，可以直接核算。设 \(r=[\beta]\)、\(b=p(r)\)，并使用 \(q_x,r_y\) 的上述表示。两条拼接路径逐点相乘后，前半段是 \(\alpha\cdot\beta\)，后半段为
	\[
		a\kappa_x(t)b\kappa_y(t)
		=ab\bigl(b^{-1}\kappa_x(t)b\kappa_y(t)\bigr).
	\]
	固定一个足够小的可缩单位元坐标球 \(W\)。缩小 \(x,y\) 的范围，可以使括号中的整条路径都留在 \(W\) 中：关于 \((x,y,t)\) 的表达连续，在 \(x=y=e_G\) 时对全部 \(t\in[0,1]\) 都为单位元，紧区间给出共同的缩小范围。其终点为 \(w=b^{-1}xby\)。在 \(W\) 中，这条路径与径向路径 \(\kappa_w\) 定端点同伦，故
	\[
		q_xr_y=[(\alpha\cdot\beta)*(ab\kappa_w)],\qquad w=b^{-1}xby.
	\]
	因此乘法的局部坐标表达就是 \(w=b^{-1}xby\)，为光滑映射。类似地，对 \(q_x\) 的代表逐点取逆，其末段为
	\[
		(a\kappa_x(t))^{-1}
		=a^{-1}\bigl(a\kappa_x(t)^{-1}a^{-1}\bigr).
	\]
	缩小后括号内路径留在一个可缩小邻域中，取逆的局部坐标表达为 \(x\mapsto ax^{-1}a^{-1}\)，也光滑。这也先行证明了运算关于覆盖拓扑的连续性。

	还要核对 Hausdorff 性和第二可数性。不同投影点可由底空间的不交开集分离；投影相同的不同覆盖点，则位于某均匀覆盖邻域上方的不交片内，所以 \(\widetilde G\) Hausdorff。第二可数性不能仅由底空间第二可数推出。取一个对称开单位元邻域 \(V\subset\widetilde G\)，包含在一张覆盖坐标片内；它本身第二可数。由于运算已经连续，\(\bigcup_{n\geq1}V^n\) 是开子群，其各陪集也开；\(\widetilde G\) 连通，所以这个开子群等于整个 \(\widetilde G\)。每个乘法映射
	\[
		V^{\times n}\longrightarrow V^n,\qquad
		(v_1,\ldots,v_n)\longmapsto v_1\cdots v_n
	\]
	都是开满射：乘积空间中的开矩形映成若干群开集的乘积，而该乘积是开集，因为可固定前面因子再左平移最后一个开集，随后取并。因此 \(V^{\times n}\) 的可数基的像给出 \(V^n\) 的可数基。各 \(V^n\) 在 \(\widetilde G\) 中开，且可数个这样的集合覆盖 \(\widetilde G\)，所有这些基的并便是总空间的可数基。至此，提升的坐标使 \(\widetilde G\) 成为有限维实 Lie 群，\(p\) 是局部微分同胚。

	其核是覆盖的单位元纤维，因而离散。固定 \(k\in\ker p\)，连续映射 \(q\mapsto qkq^{-1}\) 的像仍在 \(\ker p\) 中。连通空间到离散空间的连续映射为常值，在 \(q=\widetilde e\) 处取值为 \(k\)，故 \(k\) 与每个 \(q\) 交换。因此核位于中心。核的点恰是基于 \(e_G\) 的环路类；对两条环路，\eqref{g1:eq:pointwise-product-concatenation}中的 \(a=e_G\)，所以核乘法正是第五章基本群的拼接乘法。于是 \(\ker p\cong\pi_1(G,e_G)\)，而中心的子群交换，得到最后结论。

	最后说明唯一性。使 \(p\) 成为局部微分同胚的光滑结构由各覆盖片唯一确定。若还有另一套使 \(p\) 为同态的连续群乘法，那么两种乘法都是映射 \((q,r)\mapsto p(q)p(r)\) 到覆盖空间的提升，并在 \((\widetilde e,\widetilde e)\) 取相同值。空间 \(\widetilde G\times\widetilde G\) 连通，故\cref{g1:lem:lift-uniqueness}使两种乘法相同；单位元和逆元也随之确定。
\end{proof}

\input{Figures/Chapter16-CoverPathProduct.tex}

在这个群结构下，基本群的左覆盖变换作用也有直接解释。若 \(k=[\gamma]\in\ker p\)，则对任意 \(q=[\alpha]\)，\eqref{g1:eq:pointwise-product-concatenation}给出
\[
	kq=[\gamma*\alpha].
\]
这恰是第五章前置环路的左作用。因此“左覆盖变换”与“由核元素作左乘”在这里一致；核的中心性又说明它同时等于右乘。

\subsection{离散中心商与同 Lie 代数的不同 Lie 群}
\label{g1:subsec:1604:05}

万有覆盖同态把离散中心核中的元素识别为单位元，却不改变任何点附近的微分结构。反向地，从一个 Lie 群出发作同样的识别，也应得到具有相同 Lie 代数的 Lie 群。此处的离散性让商空间可以由不相交的小邻域直接构造。

\begin{proposition}[离散中心商]
	\label{g1:prop:discrete-central-quotient}
	设 \(G\) 是连通有限维实 Lie 群，\(\Gamma\subset Z(G)\) 为离散子群，其中离散指其子空间拓扑离散。则 \(\Gamma\) 闭，商群 \(G/\Gamma\) 具有唯一的光滑 Lie 群结构，使自然投影 \(q:G\to G/\Gamma\) 为光滑覆盖同态及局部微分同胚。特别地，\(\mathrm{d}q_{e_G}\) 是 Lie 代数同构。
\end{proposition}
\begin{proof}
	先证明闭性。由离散性取单位元开邻域 \(O\)，使 \(O\cap\Gamma=\{e_G\}\)，再取单位元邻域 \(V\)，使 \(V^{-1}V\subset O\)。每个左平移 \(xV\) 至多含一个 \(\Gamma\) 中的点，因为若 \(\gamma_1=xv_1,\gamma_2=xv_2\)，则 \(\gamma_1^{-1}\gamma_2=v_1^{-1}v_2\in O\cap\Gamma\)。若 \(x\in\overline\Gamma\)，开邻域 \(xV\) 必含一个 \(\gamma\in\Gamma\)，并且只含这一个。于是 \(x\) 属于单点 \(\{\gamma\}\) 的闭包；Hausdorff 性使单点闭，所以 \(x=\gamma\in\Gamma\)。故 \(\Gamma\) 闭。

	\(\Gamma\) 位于中心，因而正规，商集上有群乘法。给它商拓扑。投影 \(q\) 开，因为 \(q^{-1}(q(A))=\bigcup_{\gamma\in\Gamma}A\gamma\) 对每个开集 \(A\subset G\) 都开。继续缩小 \(V\)，使它成为坐标邻域且 \(V^{-1}V\cap\Gamma=\{e_G\}\)。则各 \(V\gamma\) 两两不交，且
	\[
		q^{-1}(q(V))=\coprod_{\gamma\in\Gamma}V\gamma.
	\]
	投影在每片上为开连续双射，故为同胚；左平移后，所有 \(q(gV)\) 都有同样的局部形式。因此 \(q\) 是覆盖。

	商空间也是 Hausdorff。若 \(x\Gamma\ne y\Gamma\)，则 \(x^{-1}y\notin\Gamma\)。由 \(\Gamma\) 闭及 \((u,v)\mapsto u^{-1}v\) 连续，可取 \(x,y\) 的开邻域 \(A,B\)，使 \(A^{-1}B\cap\Gamma=\varnothing\)。于是开集 \(q(A),q(B)\) 不交。又因为 \(q\) 开且满，\(G\) 的一个可数基的像给出商空间的可数基。

	把每个 \(q(gV)\) 的坐标从 \(gV\) 传递过去。两套局部逆在交叠点附近相差一个固定的右平移 \(R_\gamma\)：在足够小的连通均匀覆盖邻域上，选定的两张片不能突然改变所对应的 \(\gamma\)。右平移光滑，所以这些坐标相容，并使 \(q\) 局部微分同胚。商乘法的光滑性可在两个因子附近分别选择这样的光滑局部逆，再作 \(G\) 的乘法并投影；商取逆同理。它们与陪集定义一致，因而得到 Lie 群结构。任何使 \(q\) 为局部微分同胚的结构都必须使用同样的局部坐标，故唯一。最后，\(q\) 为光滑同态，微分保括号；局部微分同胚又使其微分可逆，因此得到 Lie 代数同构。
\end{proof}

还需要一个把“微分同构”转换成“覆盖同态”的结论。仅有局部逆函数定理尚未说明映射满射，也未说明所有原像能被同一个小邻域同时控制。

\begin{proposition}[微分同构与覆盖同态]
	\label{g1:prop:lie-homomorphism-local-isomorphism}
	设 \(F:G\to H\) 为连通有限维实 Lie 群之间的光滑同态。若 \(\mathrm{d}F_{e_G}\) 为同构，则 \(F\) 满射且为光滑覆盖映射，\(\ker F\) 是离散中心子群，并诱导 Lie 群同构 \(G/\ker F\cong H\)。
\end{proposition}
\begin{proof}
	逆函数定理给出单位元开邻域 \(U\subset G\)、\(V\subset H\)，使 \(F|_U:U\to V\) 为微分同胚。像 \(F(G)\) 是含 \(V\) 的子群，所以为开子群：对每个像点，左平移 \(V\) 仍留在像中。其余陪集也开，因此像的补集开。由 \(H\) 连通，得到 \(F(G)=H\)。

	记 \(K=\ker F\)。单位元在 \(K\) 中孤立，平移后每个核元素也孤立，所以 \(K\) 离散。它正规；对固定 \(k\in K\)，映射 \(g\mapsto gkg^{-1}\) 连续且取值于离散 \(K\)，源群连通使其恒等于 \(k\)，故 \(K\subset Z(G)\)。

	现在精确写出覆盖片。任意 \(x\in F^{-1}(V)\) 都有唯一 \(u\in U\) 满足 \(F(u)=F(x)\)，于是 \(x=ku\)，其中 \(k=xu^{-1}\in K\)。各 \(kU\) 两两不交：若 \(ku=k'u'\)，则 \(F(u)=F(u')\)，由 \(F|_U\) 单射先得 \(u=u'\)，再得 \(k=k'\)。因此
	\[
		F^{-1}(V)=\coprod_{k\in K}kU,
	\]
	且每片经 \(F\) 微分同胚于 \(V\)。对任意 \(h\in H\)，选 \(g\in G\) 满足 \(F(g)=h\)，左平移这个分解便覆盖 \(hV\)。故 \(F\) 为覆盖。最后，\(F\) 的纤维正是 \(K\) 陪集，因此诱导群双射 \(G/K\to H\)；它在刚构造的小片坐标中是微分同胚，故整体为 Lie 群同构。
\end{proof}

\begin{corollary}[同 Lie 代数的连通 Lie 群]
	\label{g1:cor:connected-groups-same-lie-algebra}
	设 \(G,H\) 是连通有限维实 Lie 群，且 \(\mathfrak g\cong\mathfrak h\)。则存在一个单连通 Lie 群 \(S\) 及其离散中心子群 \(\Gamma_G,\Gamma_H\)，使
	\[
		G\cong S/\Gamma_G,\qquad H\cong S/\Gamma_H.
	\]
	这里可取 \(S\) 为 \(G\) 的万有覆盖群；它同时也是 \(H\) 的一个万有覆盖群。
\end{corollary}
\begin{proof}
	取 \(p:S\to G\) 为\cref{g1:thm:universal-cover-lie-group}构造的万有覆盖同态，并固定 Lie 代数同构 \(\varphi:\mathfrak g\to\mathfrak h\)。复合 \(\varphi\circ\mathrm{d}p_e:\operatorname{Lie}(S)\to\mathfrak h\) 是 Lie 代数同构。因为 \(S\) 单连通，\cref{g1:thm:lie-homomorphism-integration}将它积分为同态 \(F:S\to H\)。由上一命题，\(F\) 是满射覆盖，核 \(\Gamma_H\) 离散且中心。令 \(\Gamma_G=\ker p\)，分别应用上一命题得到两个商同构。\(S\) 单连通，所以 \(F\) 又是 \(H\) 的万有覆盖。
\end{proof}

例如加法群 \(\mathbb R\) 和圆周群 \(S^1\) 都有一维交换 Lie 代数，其共同单连通覆盖群是 \(\mathbb R\)：前者对应核 \(\{0\}\)，后者通过 \(t\mapsto e^{it}\) 对应核 \(2\pi\mathbb Z\)。它们一个非紧、一个紧，因此并不同构。另一方面，\(\mathbb R/(a\mathbb Z)\) 在 \(a>0\) 改变时都同构于圆周；不同的离散子群未必给出不同的群同构类型。

\ProseParagraphBreak

这些结果从已经给定的 Lie 群及其 Lie 代数出发：括号确定局部同态，单连通源群允许唯一的整体积分，离散中心核则记录可以怎样作整体识别。局部微分数据与整体群的关系由此分成了可分别检查的两部分。

% 纲要条目排为编号小节标题，并进入目录与书签；说明段仍保留于纲要注释中。


\begin{chaptersummary}
	单位元切空间经左平移成为左不变向量场，场的括号给出 \(\mathfrak g=\mathrm{T}_eG\) 的李代数结构。共轭微分形成伴随表示 \(\operatorname{Ad}\)，再在单位元求微分便得到 \(\operatorname{ad}_XY=[X,Y]\)。群同态的微分保持括号；矩阵群中，这个括号是 \(XY-YX\)。

	\ProseParagraphBreak

	左不变性让局部积分曲线能够重复平移，得到完备的单参数子群。指数满足 \((\mathrm{d}\exp_G)_0=I\)，故提供单位元附近的坐标；连通群由这样的指数邻域生成，但单个指数未必能覆盖整个群。指数的自然性把同态和共轭连接到李代数，BCH 展开则以 \(\tfrac12[X,Y]\) 为局部乘法偏离加法的首个修正项。

	\ProseParagraphBreak

	李子代数确定左不变的对合分布，过单位元的最大叶形成唯一连通浸入子群。这里必须使用叶自身的流形拓扑。若子群在环境群中闭，指数横截方向不能无限靠近单位元，因而得到局部嵌入结构；无理斜率子群展示了缺少闭性时可能出现的稠密浸入。

	\ProseParagraphBreak

	李代数同态先确定局部群同态。沿路径把小增量依次相乘可延拓其值，同伦网格保证端点固定的同伦不改变结果。源群单连通时得到唯一整体同态。一般连通李群的万有覆盖带有唯一相容李群结构，覆盖核离散且位于中心；同一个已实现李代数的连通李群因而可由其单连通群的不同离散中心商得到。把同态下降到商群时，还必须检查相应离散核是否被保持。
\end{chaptersummary}

\BookExercises

\begin{exercise}\label{g1:ex:affine-group-lie-data}
	在 \(G=(0,\infty)\times\mathbb R\) 上定义
	\[
		(a,b)(c,d)=(ac,b+ad).
	\]
	证明这是一个李群。以单位元处的坐标向量 \(H=\partial_a|_{(1,0)}\)、\(E=\partial_b|_{(1,0)}\) 为基，求它们的左不变延拓与括号，并求 \(\operatorname{Ad}_{(a,b)}(xH+yE)\)。直接对所得伴随表示求导，核对 \(\operatorname{ad}_XY=[X,Y]\)。
\end{exercise}
\begin{solution}
	把 \((a,b)\) 对应到矩阵 \(\begin{pmatrix}a&b\\0&1\end{pmatrix}\)，矩阵相乘恰好给出题设运算，且对应单射，因此运算满足结合律。单位元是 \((1,0)\)，逆为 \((a^{-1},-b/a)\)。\(G\) 是 \(\mathbb R^2\) 的开集，乘法由多项式给出，求逆在 \(a>0\) 时光滑，所以它是二维实李群。这也是实直线上变换 \(t\mapsto at+b\) 的复合群。

	固定 \((a,b)\) 后，左平移为 \(L_{(a,b)}(c,d)=(ac,b+ad)\)。它在单位元的微分把 \((x,y)\) 送到 \((ax,ay)\)，故
	\[
		H^L=a\partial_a,\qquad E^L=a\partial_b.
	\]
	按坐标括号公式，\(H^L\) 对 \(E^L\) 的 \(\partial_b\) 系数 \(a\) 求导得到 \(a\)，而 \(E^L\) 对 \(H^L\) 的 \(\partial_a\) 系数求导为零。因此 \([H^L,E^L]=a\partial_b=E^L\)，即 \([H,E]=E\)。实双线性和反对称性进一步给出
	\[
		[xH+yE,x'H+y'E]=(xy'-yx')E.
	\]

	利用上面的逆元公式，逐项计算共轭得
	\[
		(a,b)(c,d)(a,b)^{-1}=(c,ad+(1-c)b).
	\]
	在 \((c,d)=(1,0)\) 求导，因此
	\[
		\operatorname{Ad}_{(a,b)}(xH+yE)=xH+(ay-bx)E.
	\]
	任取经过单位元的曲线 \((a(t),b(t))\)，若 \(a'(0)=x\)、\(b'(0)=y\)，则对 \(Y=x'H+y'E\)，
	\[
		\left.\frac{\mathrm{d}}{\mathrm{d}t}\right|_{t=0}\operatorname{Ad}_{(a(t),b(t))}Y
		=(xy'-yx')E=[xH+yE,Y].
	\]
	这正是 \(\operatorname{ad}_XY=[X,Y]\)，同时显示了缩放方向与平移方向的括号为何非零：改变缩放会改变平移的大小。
\end{solution}

\begin{exercise}\label{g1:ex:right-invariant-bracket-sign}
	设 \(G\) 为有限维实李群，\(X,Y\in\mathfrak g=\mathrm{T}_eG\)。定义右不变场 \(X^R_g=\mathrm{d}(R_g)_eX\)，左不变延拓仍记为 \(X^L\)。证明求逆微分同胚 \(\iota(g)=g^{-1}\) 满足 \(\iota_*X^L=-X^R\)，并推出
	\[
		[X^R,Y^R]=-[X,Y]^R,\qquad [X^L,Y^R]=0.
	\]
	说明为什么用右不变场定义单位元切空间的括号时，需要调整一个负号才能沿用本书的李代数约定。
\end{exercise}
\begin{solution}
	由乘法光滑，\(g\mapsto\mathrm{d}(R_g)_eX\) 光滑；又由 \(R_a\circ R_g=R_{ga}\)，
	\[
		\mathrm{d}(R_a)_gX^R_g=\mathrm{d}(R_{ga})_eX=X^R_{ga},
	\]
	所以 \(X^R\) 确为右不变场。恒等式 \(\iota\circ L_g=R_{g^{-1}}\circ\iota\) 在 \(e\) 处求导，并用\eqref{g1:eq:lie-group-unit-differentials}的 \(\mathrm{d}\iota_e=-\operatorname{id}\)，得到
	\[
		\mathrm{d}\iota_gX^L_g
		=\mathrm{d}(R_{g^{-1}})_e\mathrm{d}\iota_eX
		=-\mathrm{d}(R_{g^{-1}})_eX=-X^R_{g^{-1}}.
	\]
	这就是 \(\iota_*X^L=-X^R\)。微分同胚保持向量场的括号，故
	\[
		-[X,Y]^R
		=\iota_*([X,Y])^L
		=\iota_*[X^L,Y^L]
		=[-X^R,-Y^R]=[X^R,Y^R].
	\]

	为证明混合括号为零，在单位元取 \(X^L\) 的局部积分曲线 \(\gamma(t)\)。由\eqref{g1:eq:left-invariant-local-flow}，\(X^L\) 的局部流是 \(\Phi_t=R_{\gamma(t)}\)。右不变性意味着 \((R_{\gamma(t)})_*Y^R=Y^R\)，等价地 \(\Phi_t^*Y^R=Y^R\)，在相应的局部定义域内成立。用\cref{g1:prop:flow-pullback-derivative}在 \(t=0\) 求导，得到 \([X^L,Y^R]=0\)。这个论证只需小时间流，不需要先证明左不变场完备。

	最后，第一式在 \(e\) 处给出 \([X^R,Y^R]_e=-[X,Y]\)。因此若希望仍使用由左不变场定义的括号，右不变场版本应写成 \([X,Y]=-[X^R,Y^R]_e\)。负号来自求逆把单位元速度变成相反方向，不能在两种不变性之间直接省去。
\end{solution}

\begin{exercise}\label{g1:ex:so-three-cross-product}
	对 \(v=(v_1,v_2,v_3)\in\mathbb R^3\)，令
	\[
		\widehat v=
		\begin{pmatrix}
			0&-v_3&v_2\\
			v_3&0&-v_1\\
			-v_2&v_1&0
		\end{pmatrix}.
	\]
	证明 \(v\mapsto\widehat v\) 是 \(\mathbb R^3\) 到 \(\mathfrak{so}(3)\) 的线性同构，并满足 \([\widehat v,\widehat w]=\widehat{v\times w}\)。对 \(R\in\operatorname{SO}(3)\)，证明 \(\operatorname{Ad}_R\widehat v=\widehat{Rv}\)。若把 \(R\) 换成任意 \(A\in\operatorname{O}(3)\)，求正确的公式。
\end{exercise}
\begin{solution}
	由\cref{g1:prop:matrix-classical-lie-groups}，\(\mathfrak{so}(3)\) 是全体实反对称 \(3\times3\) 矩阵。任意这样的矩阵恰由三个上三角元素决定，且唯一写成题设形式，所以 \(v\mapsto\widehat v\) 是线性同构。直接乘矩阵可见 \(\widehat vu=v\times u\)。对任意 \(u\in\mathbb R^3\)，展开叉积的三个坐标可得
	\[
		v\times(w\times u)=w(v\cdot u)-u(v\cdot w).
	\]
	交换 \(v,w\) 后相减，便有
	\begin{align*}
		[\widehat v,\widehat w]u
		&=v\times(w\times u)-w\times(v\times u)\\
		&=w(v\cdot u)-v(w\cdot u)=(v\times w)\times u.
	\end{align*}
	因此 \([\widehat v,\widehat w]=\widehat{v\times w}\)。这里左边采用的确实是本书的李代数括号，因为\cref{g1:prop:matrix-lie-bracket}已证明它等于矩阵交换子。

	对正交矩阵 \(A\)，利用叉积的刻画 \((v\times w)\cdot u=\det(v,w,u)\)，可得
	\[
		\bigl((Av)\times(Aw)\bigr)\cdot(Au)
		=\det A\,\det(v,w,u)
		=\bigl(\det A\,A(v\times w)\bigr)\cdot(Au).
	\]
	由于 \(Au\) 遍历所有向量，
	\[
		(Av)\times(Aw)=\det A\,A(v\times w).
	\]
	令 \(w=A^{-1}u\)，并用 \((\det A)^2=1\)，得到
	\[
		(A\widehat v A^{-1})u=A(v\times A^{-1}u)
		=\det A\,(Av)\times u.
	\]
	因而
	\[
		\operatorname{Ad}_A\widehat v=\det A\,\widehat{Av}.
	\]
	在 \(A=R\in\operatorname{SO}(3)\) 时 \(\det R=1\)，便化为题目要求的 \(\operatorname{Ad}_R\widehat v=\widehat{Rv}\)。行列式为 \(-1\) 的正交变换反转空间定向，正是它带来一般公式中的额外负号。
\end{solution}

\begin{exercise}\label{g1:ex:circle-power-homomorphism-kernel}
	设 \(m\) 为非零整数。证明 \(P_m:S^1\to S^1\)、\(P_m(z)=z^m\) 是李群同态，求其单位元微分与核，证明它在每一点都是局部微分同胚，并判断何时为李群同构。再考虑加法李群之间的同态 \(F:\mathbb R\to\mathbb R^2\)、\(F(t)=(t,0)\)：判断 \((0,0)\) 是否为正则值，并直接说明 \(\ker F\) 的光滑结构。这两个例子分别说明了从同态的微分推断整体性质时的什么限制？
\end{exercise}
\begin{solution}
	圆周上的复数乘法交换，故 \((zw)^m=z^mw^m\)，包括负整数 \(m\) 在内。非负整数次幂由复数乘法组成，负整数次幂还要复合光滑求逆，因此 \(P_m\) 光滑，是李群同态。由 \(\mathrm{T}_1S^1=i\mathbb R\)，取经过 \(1\) 的曲线 \(z(t)=\cos t+i\sin t\)，它的初速度为 \(i\)，而 \(P_m(z(t))=\cos(mt)+i\sin(mt)\)，故
	\[
		\mathrm{d}(P_m)_1(it)=imt.
	\]
	\(m\neq0\) 时这个微分可逆，由\cref{g1:prop:lie-homomorphism-constant-rank}，其他点的微分也可逆；逆函数定理说明 \(P_m\) 在每点都是局部微分同胚。

	令 \(q=|m|\)。每个圆周点可写成 \(\cos\theta+i\sin\theta\)，\(z^m=1\) 等价于 \(m\theta\in2\pi\mathbb Z\)。因此
	\[
		\ker P_m=\left\{\cos\frac{2\pi k}{q}+i\sin\frac{2\pi k}{q}:k=0,1,\ldots,q-1\right\}.
	\]
	它是 \(q\) 个点组成的离散李群，其切空间为零，符合 \(\ker\mathrm{d}(P_m)_1=0\)。给定 \(w=\cos\theta+i\sin\theta\)，取 \(z=\cos(\theta/m)+i\sin(\theta/m)\) 即有 \(P_m(z)=w\)，所以 \(P_m\) 总是满射。当 \(|m|>1\) 时核不止一个点，它不是单射；当 \(m=1\) 时是恒等映射，当 \(m=-1\) 时是求逆，二者都有光滑逆。因此恰在 \(m=\pm1\) 时是李群同构。

	对第二个同态，\(\mathrm{d}F_t(v)=(v,0)\) 在每点的秩均为 \(1\)，没有满射到二维靶切空间。唯一原像点为 \(t=0\)，所以 \((0,0)\) 不是正则值。不过 \(\ker F=\{0\}\) 显然是零维嵌入子流形，以唯一的点为零维坐标图即可；其切空间为 \(0=\ker\mathrm{d}F_0\)，也符合恒秩核的结论。

	第一个例子表明微分处处可逆仍不能排除离散的多对一现象，局部微分同胚未必是整体同构。第二个例子表明同态的核是嵌入子流形，并不意味着靶单位元是正则值；在微分不满射时，保证核光滑的是恒秩标准形给出的纤维切片。
\end{solution}

\begin{exercise}\label{g1:ex:affine-exponential}
	令 \(G=\{\left(\begin{smallmatrix}a&b\\0&1\end{smallmatrix}\right):a>0,\ b\in\mathbb R\}\)，采用参数 \((a,b)\) 给出的光滑结构。求其李代数与指数映射，并证明 \(\exp_G:\mathfrak g\to G\) 是整体微分同胚。
\end{exercise}
\begin{solution}
	矩阵乘法和求逆在半平面 \(a>0\) 上光滑，因此 \(G\) 是二维实李群。单位元为 \((1,0)\)，切向矩阵恰为
	\[
		X=\begin{pmatrix}x&y\\0&0\end{pmatrix},\qquad x,y\in\mathbb R.
	\]
	两个这样的矩阵有括号 \([X(x,y),X(u,v)]=X(0,xv-uy)\)，从而给出这二维空间的完整李代数结构。对 \(k\geq1\)，\(X^k=\left(\begin{smallmatrix}x^k&x^{k-1}y\\0&0\end{smallmatrix}\right)\)。由矩阵指数级数及包含同态的自然性，
	\[
		\exp_GX=\begin{pmatrix}e^x&y\psi(x)\\0&1\end{pmatrix},\qquad
		\psi(x)=\int_0^1e^{tx}\,\mathrm{d}t
		=\begin{cases}(e^x-1)/x,&x\ne0,\\1,&x=0.\end{cases}
	\]
	积分表达式说明 \(\psi\) 在整个实轴光滑且严格为正，因而分段写法在零点没有奇性。给定 \((a,b)\) 且 \(a>0\)，唯一原像为
	\[
		x=\log a,\qquad y=\frac{b}{\psi(\log a)}.
	\]
	这也光滑，所以指数是整体微分同胚。这里全局可逆是明确计算的结果，不能从任意李群的单位元附近局部可逆性直接推出。
\end{solution}

\begin{exercise}\label{g1:ex:non-surjective-exponential}
	在 \(\operatorname{SL}(2,\mathbb R)\) 中令 \(A=\operatorname{diag}(-2,-1/2)\)。证明 \(A\) 属于单位连通分支，且是两个指数的乘积，但不存在实矩阵 \(X\) 满足 \(e^X=A\)。这说明连通李群的指数不一定满射。
\end{exercise}
\begin{solution}
	令
	\[
		J=\begin{pmatrix}0&-1\\1&0\end{pmatrix},\qquad D=\begin{pmatrix}1&0\\0&-1\end{pmatrix}.
	\]
	二者迹为零，属于 \(\mathfrak{sl}(2,\mathbb R)\)。直接求指数得 \(e^{\pi J}=-I\)、\(e^{(\log2)D}=\operatorname{diag}(2,1/2)\)，故
	\[
		A=e^{\pi J}e^{(\log2)D}.
	\]
	曲线 \(t\mapsto e^{t\pi J}e^{t(\log2)D}\)，\(0\leq t\leq1\)，留在 \(\operatorname{SL}(2,\mathbb R)\) 中，并从 \(I\) 走到 \(A\)。所以 \(A\) 属于单位连通分支 \(G^0\)，两个因子的整条单参数子群也都在 \(G^0\) 中。

	若实矩阵 \(X\) 满足 \(e^X=A\)，则 \(X\) 与 \(A\) 交换，因为 \(X\) 与自身的每个幂交换，进而与收敛的指数级数交换。写出 \(XA=AX\) 的两个非对角元，由 \(-2\ne-1/2\) 得 \(X_{12}=X_{21}=0\)。于是 \(X\) 为实对角矩阵，其指数的两个对角元都严格为正，不可能等于 \(A\)。特别地，连通李群 \(G^0\) 的指数不能到达 \(A\)，但有限个指数的乘积仍能到达它。
\end{solution}

\begin{exercise}\label{g1:ex:commutator-exponential-limit}
	设 \(G\) 为有限维实李群，\(X,Y\in\mathfrak g\)。对充分小的实数 \(s,t\)，令
	\[
		K(s,t)=\exp_G(sX)\exp_G(tY)\exp_G(-sX)\exp_G(-tY).
	\]
	证明 \(\log K(s,t)=st[X,Y]+O((|s|+|t|)^3)\)，并由此证明
	\[
		K(1/n,1/n)^{n^2}\longrightarrow\exp_G([X,Y]).
	\]
	说明改变交换子的次序为何会改变首项的符号。
\end{exercise}
\begin{solution}
	由共轭与指数的相容性，前面三个因子的乘积为
	\[
		\exp_G\bigl(t e^{s\operatorname{ad}_X}Y\bigr).
	\]
	固定 \(X,Y\)，记 \(\delta=|s|+|t|\)，并令 \(U=t e^{s\operatorname{ad}_X}Y\)、\(V=-tY\)。算子指数展开给出
	\[
		U=tY+st[X,Y]+O(\delta^3),\qquad U+V=st[X,Y]+O(\delta^3).
	\]
	\([tY,-tY]=0\)，故 \([U,V]=O(\delta^3)\)。将它们代入 BCH 的二阶展开，得到所需对数公式；此处 \(K\) 对充分小的参数确在单位元对数邻域内。

	令 \(Z_n=\log K(1/n,1/n)\)。则 \(Z_n=n^{-2}[X,Y]+O(n^{-3})\)，所以 \(n^2Z_n\to[X,Y]\)。单参数子群规律给出 \(K(1/n,1/n)^{n^2}=\exp_G(n^2Z_n)\)，连续性即得极限。若把 \(X,Y\) 两方向的角色交换，首项变为 \(st[Y,X]=-st[X,Y]\)；这里的符号由写出的乘积顺序决定。
\end{solution}

\begin{exercise}\label{g1:ex:adjoint-kernel-center}
	设 \(G\) 为连通有限维实李群，\(Z(G)=\{g\in G:gh=hg\text{ 对所有 }h\in G\}\)。证明 \(\ker\operatorname{Ad}=Z(G)\)。再证明若 \([X,Y]=0\) 对所有 \(Y\in\mathfrak g\) 成立，则 \(\exp_GX\in Z(G)\)，并给出逆命题不成立的例子。
\end{exercise}
\begin{solution}
	若 \(g\in Z(G)\)，共轭 \(C_g\) 是恒等映射，故其单位元微分为恒等，即 \(g\in\ker\operatorname{Ad}\)。反之，若 \(\operatorname{Ad}_g=I\)，指数自然性说明
	\[
		g\exp_G(Y)g^{-1}=\exp_G(Y)\qquad(Y\in\mathfrak g).
	\]
	因 \(G\) 连通，每个群元素都是有限个指数的乘积，故 \(g\) 与所有群元素交换。这一步正是连通性的用途。

	若 \(\operatorname{ad}_X=0\)，则 \(\operatorname{Ad}_{\exp_GX}=e^{\operatorname{ad}_X}=I\)，于是 \(\exp_GX\) 为中心元素。反向可能因指数的周期而失败。取 \(G=\operatorname{SL}(2,\mathbb R)^0\)，并令
	\[
		X=\pi\begin{pmatrix}0&-1\\1&0\end{pmatrix},\qquad
		Y=\begin{pmatrix}1&0\\0&-1\end{pmatrix}.
	\]
	\(X,Y\in\operatorname{Lie}(G)=\mathfrak{sl}(2,\mathbb R)\)，其中等号来自单位分支为开子群。\(\exp_GX=-I\) 与所有矩阵交换，但
	\[
		[X,Y]=2\pi\begin{pmatrix}0&1\\1&0\end{pmatrix}\ne0.
	\]
	因此“某一个指数落入中心”不能替代“整个方向与所有方向的括号为零”。
\end{solution}

\begin{exercise}\label{g1:ex:closed-subgroups-real-line}
	证明加法李群 \((\mathbb R,+)\) 的闭子群恰为 \(\{0\}\)、\(\mathbb R\) 以及 \(a\mathbb Z\)、\(a>0\)。求它们的李代数与单位元连通分支，并说明为什么李代数不能区分所有闭子群。
\end{exercise}
\begin{solution}
	设 \(K\subseteq\mathbb R\) 为闭子群。若 \(K\) 没有正元素，由取负封闭可知它也没有负元素，故 \(K=\{0\}\)。以下设正元素集非空，并令
	\[
		a=\inf\{x\in K:x>0\}\ge0.
	\]

	若 \(a>0\)，由下确界的定义可取 \(x_j\in K\)、\(x_j>0\)，使 \(x_j\to a\)。闭性给出 \(a\in K\)。任取 \(x\in K\)，选择整数 \(m\) 使 \(ma\le x<(m+1)a\)，则 \(r=x-ma\in K\) 且 \(0\le r<a\)。若 \(r>0\)，就与 \(a\) 是正元素下确界矛盾，所以 \(r=0\)，即 \(x\in a\mathbb Z\)。反向包含由 \(a\in K\) 和子群性质得到，故 \(K=a\mathbb Z\)。

	若 \(a=0\)，可取 \(0<x_j\in K\) 且 \(x_j\to0\)。对任意 \(y\in\mathbb R\)，选整数 \(m_j\) 使
	\[
		|m_jx_j-y|\le x_j.
	\]
	于是 \(m_jx_j\in K\) 且趋于 \(y\)，闭性迫使 \(y\in K\)。因此 \(K=\mathbb R\)。反过来，\(\{0\}\) 和 \(\mathbb R\) 显然闭；当 \(a>0\) 时，\(a\mathbb Z\) 的每个点与其余点距离至少为 \(a\)，而不在该集合中的点位于两个相邻整数倍之间，存在避开它的开区间，所以 \(a\mathbb Z\) 也闭。

	加法群的指数为恒等映射，故闭子群定理给出的李代数为
	\[
		\mathfrak k=\{X\in\mathbb R:tX\in K\text{ 对所有 }t\in\mathbb R\}.
	\]
	对 \(K=\mathbb R\)，它等于 \(\mathbb R\)；对 \(K=\{0\}\) 或 \(K=a\mathbb Z\)，它等于 \(\{0\}\)，因为非零 \(X\) 的全部实数倍充满 \(\mathbb R\)。相应的单位元连通分支分别是 \(\mathbb R\)、\(\{0\}\) 和 \(\{0\}\)。所以所有这些非平凡离散子群都具有同一个零李代数；不同的离散分支不会出现在单位元切空间中。
\end{solution}

\begin{exercise}\label{g1:ex:triangular-two-dimensional-subgroups}
	考虑上三角矩阵李群
	\[
		N=\left\{
		n(x,y,z)=
		\begin{pmatrix}
			1&x&z\\
			0&1&y\\
			0&0&1
		\end{pmatrix}:x,y,z\in\mathbb R
		\right\}.
	\]
	记其李代数的基为 \(X=E_{12}\)、\(Y=E_{23}\)、\(Z=E_{13}\)。证明每个二维李子代数都恰为
	\[
		\mathfrak h_{a,b}=\operatorname{span}\{aX+bY,Z\},
		\qquad(a,b)\ne(0,0).
	\]
	求出相应的连通积分子群，并判断它是否闭。说明为什么 \(\operatorname{span}\{X,Y\}\) 不能成为二维浸入李子群的李代数。
\end{exercise}
\begin{solution}
	矩阵乘法与交换子分别给出
	\[
		n(x,y,z)n(x',y',z')
		=n(x+x',y+y',z+z'+xy'),
	\]
	以及
	\[
		[X,Y]=Z,\qquad[X,Z]=[Y,Z]=0.
	\]
	所以对 \(U=aX+bY+cZ\)、\(V=a'X+b'Y+c'Z\)，有
	\[
		[U,V]=(ab'-ba')Z.
	\]

	设 \(\mathfrak h\) 为二维李子代数。若 \(Z\notin\mathfrak h\)，则投影
	\[
		\pi:\mathfrak h\longrightarrow\operatorname{span}\{X,Y\},
		\qquad aX+bY+cZ\longmapsto aX+bY
	\]
	的核为零，所以它在两个二维空间之间是同构。取 \(\mathfrak h\) 的一组基 \(U,V\)，其投影独立，因而 \(ab'-ba'\ne0\)。括号封闭性随即给出 \(Z\in\mathfrak h\)，矛盾。因此 \(Z\in\mathfrak h\)。再选一个不与 \(Z\) 共线的元素 \(aX+bY+cZ\)，减去 \(cZ\) 即得所述形式，其中 \((a,b)\ne(0,0)\)。反过来，任意这样的平面内两个元素的括号都为零，故确为李子代数。

	令
	\[
		H_{a,b}=\{n(x,y,z):bx-ay=0\}
		=\{n(as,bs,z):s,z\in\mathbb R\}.
	\]
	乘法公式说明它对乘积封闭。逆矩阵为
	\[
		n(x,y,z)^{-1}=n(-x,-y,xy-z),
	\]
	所以它也对取逆封闭。约束 \(bx-ay=0\) 是闭的线性条件，故 \(H_{a,b}\) 在 \(N\) 中闭。参数 \((s,z)\mapsto n(as,bs,z)\) 是从 \(\mathbb R^2\) 到这个平面的微分同胚，说明它连通；在单位元求导得到切空间 \(\operatorname{span}\{aX+bY,Z\}\)。由\cref{g1:thm:lie-subalgebra-integration}的唯一性，它就是所求的连通积分子群。这里的参数 \(z\) 是矩阵右上角的实际条目，无须先把它换成指数坐标。

	最后，\([X,Y]=Z\notin\operatorname{span}\{X,Y\}\)，所以这个平面不是李子代数。浸入李子群的单位切空间必须对括号封闭，故不可能由任何二维浸入李子群实现。两个切向方向逐次运动产生的第三方向，正是本例的障碍。
\end{solution}

\begin{exercise}\label{g1:ex:vector-line-stabilizer-subgroups}
	设 \(n\ge2\)，\(v\in\mathbb R^n\setminus\{0\}\)。在 \(\operatorname{GL}(n,\mathbb R)\) 中定义
	\[
		K_v=\{g:gv=v\},\qquad
		K_{\mathbb Rv}=\{g:g(\mathbb Rv)=\mathbb Rv\}.
	\]
	证明这两个集合都是闭子群，分别求其李代数和维数。这里的闭性相对于 \(\operatorname{GL}(n,\mathbb R)\)，不要求它们在所有 \(n\times n\) 矩阵组成的向量空间中闭。
\end{exercise}
\begin{solution}
	保持一个向量或一条直线的可逆线性映射，其复合与逆仍保持相应对象，所以两集合都是子群。选一组以 \(v\) 为首向量的基。在这组基中，它们分别具有块矩阵形式
	\[
		K_v=
		\left\{
		\begin{pmatrix}
			1&a\\
			0&B
		\end{pmatrix}:
		\begin{aligned}
			&a\in\operatorname{Mat}_{1,n-1}(\mathbb R),\\
			&B\in\operatorname{GL}(n-1,\mathbb R)
		\end{aligned}
		\right\},
	\]
	\[
		K_{\mathbb Rv}=
		\left\{
		\begin{pmatrix}
			c&a\\
			0&B
		\end{pmatrix}:
		\begin{aligned}
			&c\in\mathbb R\setminus\{0\},\\
			&a\in\operatorname{Mat}_{1,n-1}(\mathbb R),\\
			&B\in\operatorname{GL}(n-1,\mathbb R)
		\end{aligned}
		\right\}.
	\]
	前者由第一列等于 \((1,0,\ldots,0)^{\mathsf T}\) 的闭条件描述；后者由第一列的后 \(n-1\) 项为零的闭条件描述。后一情形中，\(c\ne0\) 与 \(B\) 可逆已由整个矩阵可逆保证。因此两者在 \(\operatorname{GL}(n,\mathbb R)\) 中闭，可以使用\cref{g1:thm:closed-subgroup}。

	若 \(\exp(tA)v=v\) 对所有 \(t\) 成立，在 \(t=0\) 求导得到 \(Av=0\)。反之，若 \(Av=0\)，则 \(A^kv=0\) 对所有 \(k\ge1\) 成立，指数级数给出 \(\exp(tA)v=v\)。于是
	\[
		\operatorname{Lie}(K_v)=\{A:Av=0\}
		=
		\left\{
		\begin{pmatrix}
			0&a\\
			0&C
		\end{pmatrix}
		\right\}.
	\]

	若 \(\exp(tA)v\in\mathbb Rv\) 对所有 \(t\) 成立，对商空间投影 \(\mathbb R^n\to\mathbb R^n/\mathbb Rv\) 求导可得 \(Av\in\mathbb Rv\)。反之，若 \(Av=\lambda v\)，则指数级数给出 \(\exp(tA)v=e^{t\lambda}v\)，它始终张成原来的直线。故
	\[
		\operatorname{Lie}(K_{\mathbb Rv})=\{A:Av\in\mathbb Rv\}
		=
		\left\{
		\begin{pmatrix}
			\lambda&a\\
			0&C
		\end{pmatrix}
		\right\}.
	\]
	这里 \(C\) 是任意 \((n-1)\times(n-1)\) 实矩阵。计数自由条目，得到
	\[
		\dim K_v=(n-1)+(n-1)^2=n(n-1),\qquad
		\dim K_{\mathbb Rv}=n^2-n+1.
	\]
	相差的一个维数来自沿指定直线的伸缩：固定向量禁止这一变化，保持直线则允许它。
\end{solution}

\begin{exercise}\label{g1:ex:connected-subgroup-normal-in-closure}
	设 \(H\) 为有限维实李群 \(G\) 的连通浸入李子群，\(\mathfrak h=\operatorname{Lie}(H)\)，令 \(K=\overline H\)，并赋予它闭子群定理给出的李群结构，记 \(\mathfrak k=\operatorname{Lie}(K)\)。证明
	\[
		\operatorname{Ad}_k\mathfrak h=\mathfrak h\quad(k\in K),
		\qquad[\mathfrak k,\mathfrak h]\subseteq\mathfrak h,
	\]
	且 \(H\) 是 \(K\) 的正规子群。第二个包含式称 \(\mathfrak h\) 为 \(\mathfrak k\) 的理想。用无理斜率环面子群说明，结论并不意味着 \(\mathfrak h=\mathfrak k\)。
\end{exercise}
\begin{solution}
	包含映射 \(H\to G\) 的像落在嵌入子流形 \(K\) 中，由\cref{g1:prop:embedded-initial-property}，它作为到 \(K\) 的映射仍光滑。因而单位元微分给出 \(\mathfrak h\subseteq\mathfrak k\)。对 \(h\in H\)，内部共轭是 \(H\) 的光滑自同构，包含到 \(G\) 后与环境共轭相容，所以
	\[
		\operatorname{Ad}_h\mathfrak h=\mathfrak h.
	\]

	任取 \(k\in K\)。流形在每点有可数邻域基，由闭包的定义可以选 \(h_j\in H\) 使 \(h_j\to k\)，这里的收敛在环境群中理解。对 \(Y\in\mathfrak h\)，伴随作用连续，所以
	\[
		\operatorname{Ad}_{h_j}Y\longrightarrow\operatorname{Ad}_kY.
	\]
	有限维向量空间的线性子空间是闭集，而左侧每一项均属于 \(\mathfrak h\)，故 \(\operatorname{Ad}_kY\in\mathfrak h\)。再对 \(k^{-1}\) 作同一论证，得到反向包含，证明 \(\operatorname{Ad}_k\mathfrak h=\mathfrak h\)。

	若 \(X\in\mathfrak k\)，则闭子群定理保证 \(\exp_G(tX)\in K\) 对所有实数 \(t\) 成立。因此对 \(Y\in\mathfrak h\)，上一结论给出
	\[
		\operatorname{Ad}_{\exp_G(tX)}Y\in\mathfrak h.
	\]
	在零点取差商，使用\eqref{g1:eq:ad-bracket}以及 \(\mathfrak h\) 的闭线性性质，得到 \([X,Y]\in\mathfrak h\)。这就是理想条件。

	最后固定 \(k\in K\)。共轭 \(C_k\) 是 \(G\) 的微分同胚，把 \(H\) 的流形结构传到 \(kHk^{-1}\)，便得到一个连通浸入李子群，其李代数为 \(\operatorname{Ad}_k\mathfrak h=\mathfrak h\)。由\cref{g1:thm:lie-subalgebra-integration}，同一李子代数的连通积分子群唯一，所以 \(kHk^{-1}=H\)。这对每个 \(k\in K\) 成立，即 \(H\) 在 \(K\) 中正规；这里使用的是 \(H\) 自身的连通性。

	在\cref{g1:exm:irrational-torus-subgroup}中，\(H=\{[t,\alpha t]:t\in\mathbb R\}\)，\(\alpha\) 无理，且 \(K=T^2\)。于是 \(\mathfrak h=\mathbb R(1,\alpha)\) 真包含于 \(\mathfrak k=\mathbb R^2\)。环面群交换，两者的括号都为零，任何线性子空间都是理想。这同时满足本题所有结论，并显示闭包确实可以增大李代数。
\end{solution}

\begin{exercise}\label{g1:ex:torus-homomorphism-integer-matrix}
	记 \(T^m=\mathbb R^m/\mathbb Z^m\)，用商映射的单位元微分把其李代数识别为 \(\mathbb R^m\)，\(T^n\) 同理。设 \(A:\mathbb R^m\to\mathbb R^n\) 为实线性映射。证明 \(A\) 能积分为李群同态 \(T^m\to T^n\)，当且仅当 \(A\) 在标准基下的全部矩阵条目都是整数；存在时写出这个同态。对 \(m=n=1\)，说明非零整数 \(k\) 所给同态的每个纤维有多少点，并说明 \(k=0\) 的情形。
\end{exercise}
\begin{solution}
	设 \(q_m:\mathbb R^m\to T^m\)、\(q_n:\mathbb R^n\to T^n\) 为商同态。若 \(F:T^m\to T^n\) 的单位元微分为 \(A\)，则 \(F\circ q_m\) 与 \(q_n\circ A\) 都是从连通群 \(\mathbb R^m\) 到 \(T^n\) 的同态，且微分相同。由\cref{g1:thm:lie-homomorphism-integration}的唯一性，
	\[
		F(q_m(x))=q_n(Ax).
	\]
	对每个 \(z\in\mathbb Z^m\)，左边在 \(x=z\) 时为单位元，所以 \(Az\in\mathbb Z^n\)。特别地，对标准基向量逐一使用这个条件，得到 \(A\) 的每一列都是整数向量，即所有矩阵条目为整数。

	反过来，若 \(A\mathbb Z^m\subset\mathbb Z^n\)，定义 \(F([x])=[Ax]\)。改变代表 \(x\) 为 \(x+z\) 时，目标改变 \(Az\in\mathbb Z^n\)，所以映射良定义。它保持加法，并在商覆盖的局部坐标中由线性映射 \(A\) 表示，因而光滑，微分也正是 \(A\)。

	一维时映射为 \([x]\mapsto[kx]\)。若 \(k\ne0\)，其核由 \(kx\in\mathbb Z\) 决定，共有 \(|k|\) 个模 \(\mathbb Z\) 的不同解。映射满射，因为对任意代表 \(y\in\mathbb R\) 可取 \(x=y/k\)；每个纤维是核的一个陪集，故也有 \(|k|\) 个点。\(k=0\) 时则为常值同态，单位元的纤维是整个圆周，其他点的纤维为空，因而不是到 \(T^1\) 的覆盖。
\end{solution}

\begin{exercise}\label{g1:ex:connected-abelian-lie-group-quotient}
	设 \(G\) 是 \(n\) 维连通实李群，且其李代数括号恒为零。利用同态积分证明 \(G\) 为交换群，并证明存在加法群 \(\mathbb R^n\) 的离散子群 \(\Gamma\)，使 \(G\cong\mathbb R^n/\Gamma\)。这里已经给定 \(G\)，不要求使用交换李群的分类定理。
\end{exercise}
\begin{solution}
	加法群 \(\mathbb R^n\) 的李代数也是 \(\mathbb R^n\)，括号恒为零。选线性同构 \(L:\mathbb R^n\to\mathfrak g\)，因为两边括号都为零，它是李代数同构。源群 \(\mathbb R^n\) 连通且单连通，所以\cref{g1:thm:lie-homomorphism-integration}将 \(L\) 积分成同态 \(F:\mathbb R^n\to G\)。

	\(\mathrm{d}F_0=L\) 为同构，目标 \(G\) 连通，因此\cref{g1:prop:lie-homomorphism-local-isomorphism}使 \(F\) 成为满射覆盖，核 \(\Gamma\) 离散，并给出 \(G\cong\mathbb R^n/\Gamma\)。加法群 \(\mathbb R^n\) 交换，其任意同态像都交换；由满射性，\(G\) 交换。也可在商的表示中直接读出这一点。\(n=0\) 时，连通零维群只有一个点，同一结论取 \(\Gamma=\{0\}\)。
\end{solution}

\begin{exercise}\label{g1:ex:discrete-normal-subgroup-central}
	证明连通李群的任意离散正规子群都位于中心。再在 \(\operatorname{O}(2)\) 中取旋转 \(r=R_{2\pi/3}\)，证明 \(\Gamma=\{I,r,r^2\}\) 是离散正规子群，但不位于中心，由此说明连通条件的作用。
\end{exercise}
\begin{solution}
	设 \(G\) 连通、\(\Gamma\) 离散且正规。固定 \(\gamma\in\Gamma\)。正规性保证连续映射
	\[
		G\longrightarrow\Gamma,\qquad g\longmapsto g\gamma g^{-1}
	\]
	确实取值于 \(\Gamma\)。连通空间的连续像连通，而离散空间中连通子集至多含一点，所以该映射恒等于它在单位元的值 \(\gamma\)。因此 \(g\gamma=\gamma g\) 对所有 \(g\) 成立，即 \(\Gamma\subset Z(G)\)。

	现在令
	\[
		R_\theta=\begin{pmatrix}\cos\theta&-\sin\theta\\
		\sin\theta&\cos\theta\end{pmatrix},\qquad
		S=\begin{pmatrix}1&0\\0&-1\end{pmatrix}.
	\]
	每个二维正交矩阵都是某个 \(R_\theta\) 或 \(R_\theta S\)：第一列为单位向量 \((\cos\theta,\sin\theta)\)，第二列因正交性及单位长度只能是 \(\pm(-\sin\theta,\cos\theta)\)。旋转彼此交换，直接相乘又得 \(SR_\theta S^{-1}=R_{-\theta}\)。所以任意正交矩阵共轭 \(\Gamma\) 时，都将它保持不变，\(\Gamma\) 正规。它有限，故在 Hausdorff 矩阵空间中离散。

	但 \(SrS^{-1}=r^{-1}=r^2\ne r\)，所以 \(r\) 不在中心。共轭映射在 \(\operatorname{O}(2)\) 的旋转分支与反射分支上可以取不同值；源空间不连通，便不能用前段的常值论证把两者连在一起。
\end{solution}

\begin{exercise}\label{g1:ex:lie-homomorphism-descent}
	设 \(S\) 为单连通实李群，\(\Gamma\subset Z(S)\) 为离散子群，\(q:S\to S/\Gamma\) 为商同态。给定有限维实李群 \(H\) 及李代数同态 \(\varphi:\operatorname{Lie}(S)\to\mathfrak h\)，令 \(\widetilde F:S\to H\) 为其整体积分。证明存在同态 \(F:S/\Gamma\to H\) 满足 \(F\circ q=\widetilde F\)，当且仅当 \(\Gamma\subset\ker\widetilde F\)，并说明存在时为何光滑且唯一。将结论用于 \(S=(\mathbb R,+)\)、\(\Gamma=2\pi\mathbb Z\)、\(H=(\mathbb R,+)\) 及 \(\varphi(t)=\lambda t\)，判断哪些 \(\lambda\) 可以下降。
\end{exercise}
\begin{solution}
	若已有 \(F\)，则对 \(\gamma\in\Gamma\)，有
	\[
		\widetilde F(\gamma)=F(q(\gamma))=F(e_{S/\Gamma})=e_H,
	\]
	所以 \(\Gamma\subset\ker\widetilde F\)。

	反过来，若这个包含成立，定义 \(F(s\Gamma)=\widetilde F(s)\)。当 \(s'=s\gamma\) 时，
	\(\widetilde F(s')=\widetilde F(s)\widetilde F(\gamma)=\widetilde F(s)\)，故定义无关于陪集代表；同态性质由 \(\widetilde F\) 传下。由\cref{g1:prop:discrete-central-quotient}，\(q\) 在每个点附近有光滑局部逆 \(\sigma\)。在相应商邻域上，
	\[
		F=\widetilde F\circ\sigma,
	\]
	所以 \(F\) 光滑。投影 \(q\) 满射，等式 \(Fq=\widetilde F\) 又唯一决定 \(F\)。

	在最后给定的加法群例子中，\(\widetilde F(t)=\lambda t\)。下降要求 \(\widetilde F(2\pi)=2\pi\lambda=0\)，所以只有 \(\lambda=0\) 可以下降。于是任意光滑同态 \(S^1\cong\mathbb R/(2\pi\mathbb Z)\to(\mathbb R,+)\) 都是零同态。与圆周作为目标时的整数条件相比，这里目标没有可容纳非零一周增量的周期识别。
\end{solution}

\begin{exercise}\label{g1:ex:local-isomorphism-simply-connected-target}
	设 \(F:G\to H\) 是连通有限维实李群之间的光滑同态，且 \(\mathrm{d}F_{e_G}\) 为同构。证明：若目标群 \(H\) 单连通，则 \(F\) 为李群同构。再给出目标群不是单连通时结论失败的例子，并核对其单位元微分。
\end{exercise}
\begin{solution}
	令 \(\varphi=(\mathrm{d}F_{e_G})^{-1}:\mathfrak h\to\mathfrak g\)。它是李代数同态，因为李代数同构的逆保持括号。目标群 \(H\) 现在作为源群，具有单连通性，所以\cref{g1:thm:lie-homomorphism-integration}给出 \(K:H\to G\)，满足 \(\mathrm{d}K_{e_H}=\varphi\)。

	链式法则说明 \(K\circ F:G\to G\) 的单位元微分为恒等，而恒等同态的微分也是恒等。源群 \(G\) 连通，同态唯一性给出 \(KF=\operatorname{id}_G\)。同样，\(FK=\operatorname{id}_H\)。因此 \(F\) 有光滑同态逆 \(K\)，为李群同构。这个证明没有预先假定 \(G\) 单连通；它作为最终与 \(H\) 微分同胚的结果才随之成立。

	反例为
	\[
		F:(\mathbb R,+)\longrightarrow S^1,\qquad F(t)=e^{it}.
	\]
	单位元处的切向量 \(1\in \mathrm{T}_0\mathbb R\) 被送成 \(i\in \mathrm{T}_1S^1\)，按单位角速度识别两边的李代数后，微分就是 \(\mathbb R\) 上的恒等映射。但 \(\ker F=2\pi\mathbb Z\)，故 \(F\) 不单射。它是覆盖同态；非平凡的离散核正是局部同构没有变成整体同构的原因。
\end{solution}

% !TeX root = ../../Book1.tex
% 纲要标题：## 第17章　群作用、齐性空间与商结构
% 纲要位置：几何学计划纲要.md，第 768 行。
% 正文状态：初稿；尚未正式审核。

% 纲要内容开始
% 
% > 本章目标与证明范围：先证明轨道和商流形，再处理非自由作用。球面、射影空间和 Grassmann 流形承担切空间及不变结构的实际计算；有限稳定子商保留到局部模型层次。
% 
% 纲要内容结束

\chapter{群作用、齐性空间与商结构}
\label{g1:ch:17}
\BookChapterNumber{g1:ch:17}

\begin{chapterabstract}
	Lie 群作用把群中的连续运动变成流形上的向量场，稳定子则记录哪些运动在指定点上无法看见。本章先由闭子群构造齐性空间，把轨道的切空间写成 Lie 代数之商，再证明自由固有作用的轨道空间具有唯一相容的光滑结构。紧群上的归一化积分将直接由左不变密度构造，用于平均度量和固定点线性化。球面、射影空间与 Grassmann 流形提供切表示、图坐标和度量的具体计算。最后以固有作用的切片定理描述非自由作用，说明有限稳定子产生的局部商图为什么还须保存群数据。
\end{chapterabstract}

\begin{learninggoals}
	\begin{itemize}
		\item 构造基本向量场，核对左作用下的括号符号，证明稳定子、轨道秩与齐性空间切表示的关系。
		\item 从横截面和固有性证明自由固有作用的局部乘积及商流形定理，理解商映射的光滑泛性质，并辨别条件缺失的反例。
		\item 构造紧 Lie 群的归一化左右不变积分，用平均得到不变度量，并以稳定子表示组织齐性空间上的不变张量。
		\item 计算球面、实复射影空间及 Grassmann 流形的切空间、维数和不变度量，明确实维数与复维数的区别。
		\item 证明固定点线性化与固有作用切片定理，辨别底层商空间、有限稳定子数据、反射和非有效局部图。
	\end{itemize}
\end{learninggoals}

把平面中的点绕原点旋转，同一个圆周上的点就成为同一类。把球面上所有方向一起旋转，又可以把一个点送到任何另一个点。这两种观察分别引出轨道空间与齐性空间：\(M/G\) 用一点代表一整条轨道，而 \(G/G_p\) 用群参数描述经过 \(p\) 的那一条轨道。两个商的角色不同，前者保留轨道之间的方向，后者保留轨道自身的方向。

\ProseParagraphBreak

集合上能够分类，尚不足以得到光滑坐标。轨道可能稠密，稳定子可能改变大小，商中不同点也可能无法用不交邻域分开。本章的证明会分别说明自由性、固有性和稳定子紧性解决什么问题。特别是出现固定点时，商的底拓扑即使仍像一个圆盘，也未必保留全部几何数据。

\ProseParagraphBreak

本章的 Lie 群均为有限维实 Lie 群，光滑流形默认有限维、无边界、Hausdorff、第二可数；需要允许边界的商模型会另作说明。所有群作用默认左作用，基本向量场采用 \(\exp(tX)\cdot p\) 的正号初速度。复矩阵群及复参数空间先作为实流形使用，复维数在相应位置另外标明。

\ProseParagraphBreak

主要准备是\chapref{g1:ch:16}的指数映射、伴随表示与闭子群定理，\chapref{g1:ch:11}的逆函数及常秩定理，以及\chapref{g1:ch:12}的张量与度量。紧群积分使用\chapref{g1:ch:14}已经建立的密度积分和换元公式；这里会重新说明紧性怎样保证归一化与参数求导。切片证明采用紧稳定子的平均坐标。

\ProseParagraphBreak

延伸阅读可查 Lee 第二版第20章的 Infinitesimal Generators of Group Actions（第525--529页），以及第21章的 Quotients of Manifolds by Group Actions（第541--547页）和 Homogeneous Spaces（第550--554页）\cite[第525--529、541--547、550--554页]{Lee2012Smooth}。比较基本场公式时须先核对左、右作用及指数中的正负号；比较有限群商时则须核对文献是否允许反射和非有效局部群。

% !TeX root = ../../Book1.tex
% 纲要标题：### 17.1 作用和轨道的微分几何
% 纲要位置：几何学计划纲要.md，第 772 行。
% 正文状态：初稿；尚未正式审核。

% 纲要内容开始
% 
% * 无穷小作用与基本向量场
% * 稳定子、轨道切空间与秩
% * 齐性空间的光滑结构
% * 等变映射及其微分
% 
% 纲要内容结束

\section{作用和轨道的微分几何}
\label{g1:sec:1701}

在\secref{g1:sec:0104}中，群作用把集合分解为轨道，轨道又与稳定子群的陪集空间对应。现在群和被作用的空间都是流形，我们要把这份集合论描述变成微分几何：群中的一个切向量怎样推动点运动，轨道有多少个独立方向，陪集空间又怎样取得光滑坐标？这里必须保留一个区别：轨道可以有自己的流形拓扑，而这个拓扑未必就是环境流形给它的子空间拓扑。

\subsection{无穷小作用与基本向量场}
\label{g1:subsec:1701:01}

设 \(G\) 为 Lie 群，\(M\) 为光滑流形。\term[guanghuaqunzuoyong]{光滑群作用}[smooth group action]是光滑映射 \(a:G\times M\to M\)，满足 \(e\cdot p=p\) 和 \(g\cdot(h\cdot p)=(gh)\cdot p\)。本章默认左作用，记 \(a_g(p)=g\cdot p\)；\(a_g\) 是微分同胚，逆为 \(a_{g^{-1}}\)。右作用也可用同样方法处理，但变换次序与某些括号符号需要随之调整。

\ProseParagraphBreak

记 \(\mathfrak g=\mathrm{T}_eG\)。沿群中的单参数子群 \(\exp(tX)\) 运动，会使每个点 \(p\) 描出一条曲线。把其初速度收集起来，便得到 \(M\) 上的向量场。

\begin{definition}\label{g1:def:fundamental-vector-field}
	设 \(G\) 为有限维实 Lie 群，\(M\) 为光滑流形，\(a:G\times M\to M\) 为光滑左作用，记 \(g\cdot p=a(g,p)\)，\(\mathfrak g=\operatorname{Lie}(G)\)。对 \(X\in\mathfrak g\)，定义对应的\term[jibenxiangliangchang]{基本向量场}[fundamental vector field]为
	\begin{equation}\label{g1:eq:fundamental-field-convention}
		X_M(p)=\left.\frac{\mathrm{d}}{\mathrm{d}t}\right|_{t=0}\exp(tX)\cdot p.
	\end{equation}
	映射 \(X\mapsto X_M\) 称为该作用的\term[wuqiongxiaozuoyong]{无穷小作用}[infinitesimal action]，本章始终采用式中的正号约定。\symidx{\(X_M\)：正号约定下的基本向量场}
\end{definition}

\begin{proposition}[基本场的流与括号]\label{g1:prop:fundamental-fields}
	设有限维实 Lie 群 \(G\) 光滑左作用于无边界光滑流形 \(M\)。基本场 \(X_M\) 光滑且完备，其流为 \(\Phi_t(p)=\exp(tX)\cdot p\)。映射 \(X\mapsto X_M\) 线性，并满足
	\begin{equation}\label{g1:eq:fundamental-field-pushforward-bracket}
		(a_g)_*X_M=(\operatorname{Ad}_gX)_M,
		\qquad [X_M,Y_M]=-[X,Y]_M.
	\end{equation}
\end{proposition}
\begin{proof}
	对固定 \(p\)，定义轨道映射 \(\alpha_p:G\to M\)，\(\alpha_p(g)=g\cdot p\)。由于 \((\mathrm{d}\exp)_0=I\)，有 \(X_M(p)=(\mathrm{d}\alpha_p)_eX\)。微分对向量线性，并随 \(p\) 光滑变化，所以基本场光滑，且对 \(X\) 线性。

	单参数子群的乘法给出 \(\Phi_{t+s}=\Phi_s\circ\Phi_t\)。在 \(s=0\) 求导得 \(\partial_t\Phi_t(p)=X_M(\Phi_t(p))\)，初值为 \(p\)。此式对一切实数 \(t\) 都有定义，故它是全局流。

	由\cref{g1:prop:exponential-naturality}，
	\[
		g\exp(tX)\cdot p=\exp(t\operatorname{Ad}_gX)\cdot(g\cdot p).
	\]
	在 \(t=0\) 求导便得推前公式。为核对括号，回忆\cref{g1:prop:flow-pullback-derivative}：\([X_M,Y_M]=\left.\frac{\mathrm{d}}{\mathrm{d}t}\right|_0\Phi_t^*Y_M\)，其中场的拉回用逆向流的微分。因此
	\[
		\Phi_t^*Y_M=(a_{\exp(-tX)})_*Y_M
		=(\operatorname{Ad}_{\exp(-tX)}Y)_M.
	\]
	再用 \(\operatorname{Ad}_{\exp(tX)}=e^{t\operatorname{ad}_X}\)，在零点求导得到 \(-[X,Y]_M\)。
\end{proof}

负号来自左作用与\eqref{g1:eq:fundamental-field-convention}的正号约定。若改用 \(X\mapsto-X_M\)，就得到保持括号的映射。具体地，\(G\) 左乘自身时，基本场在 \(g\) 处是曲线 \(\exp(tX)g\) 的速度，因而是右不变场；前章的左不变场则沿 \(g\exp(tX)\) 运动。

\begin{example}\label{g1:exm:linear-action-fundamental}
	\(\operatorname{GL}(n,\mathbb R)\) 自然作用于 \(\mathbb R^n\) 时，\(X_M(v)=Xv\)。按向量场的坐标括号，
	\[
		[X_M,Y_M](v)=YXv-XYv=-(XY-YX)v,
	\]
	直接验证了负号。对平面旋转，取 \(J=\begin{pmatrix}0&-1\\1&0\end{pmatrix}\)，得到 \(J_M=-y\partial_x+x\partial_y\)。原点处这个场为零，其余点沿以原点为圆心的圆周运动。场在某点为零，意味着这一个单参数运动固定该点，并不意味着整个群都固定该点。
\end{example}

\subsection{稳定子、轨道切空间与秩}
\label{g1:subsec:1701:02}

固定 \(p\in M\)，记 \(G_p=\{g:g\cdot p=p\}\)、\(G\cdot p=\{g\cdot p:g\in G\}\)。这仍是\cref{g1:prop:orbit-stabilizer}中的稳定子群与轨道。回顾第10章的自由性：若每个 \(G_p=\{e\}\)，则称该作用为自由作用；若 \(\bigcap_{p\in M}G_p=\{e\}\)，则称该作用为\term[youxiaozuoyong]{有效作用}[effective action]。前者逐点排除非平凡稳定子，后者只排除处处不起作用的元素，二者不能混同。

\begin{proposition}[轨道映射的秩]\label{g1:prop:orbit-map-rank}
	设有限维实 Lie 群 \(G\) 光滑左作用于无边界光滑流形 \(M\)，\(p\in M\)。稳定子 \(H=G_p\) 是闭的嵌入 Lie 子群，其 Lie 代数为
	\begin{equation}\label{g1:eq:stabilizer-lie-algebra}
		\mathfrak h=\ker(\mathrm{d}\alpha_p)_e
		=\{X\in\mathfrak g:X_M(p)=0\}.
	\end{equation}
	轨道映射 \(\alpha_p\) 的秩恒等于 \(\dim G-\dim H\)，其在单位元处的微分像为 \(\{X_M(p):X\in\mathfrak g\}\)。此外，\(G_{g\cdot p}=gG_pg^{-1}\)。
\end{proposition}
\begin{proof}
	\(H=\alpha_p^{-1}(\{p\})\)，单点集在 Hausdorff 流形中闭，所以 \(H\) 闭。\cref{g1:thm:closed-subgroup}赋予它嵌入 Lie 子群结构，并且说明 \(X\in\mathfrak h\) 当且仅当 \(\exp(tX)\in H\) 对所有 \(t\) 成立。这一条件显然推出 \(X_M(p)=0\)。反之，若 \(X_M(p)=0\)，常值曲线 \(p\) 和曲线 \(\exp(tX)\cdot p\) 都是过 \(p\) 的积分曲线。由\secref{g1:sec:1301}的积分曲线唯一性，它们在整个共同定义域 \(\mathbb R\) 上相等，故 \(\exp(tX)\in H\)。这证明核的描述。

	恒等式 \(\alpha_p\circ L_g=a_g\circ\alpha_p\) 在单位元的微分给出
	\[
		(\mathrm{d}\alpha_p)_g\circ(\mathrm{d}L_g)_e
		=(\mathrm{d}a_g)_p\circ(\mathrm{d}\alpha_p)_e.
	\]
	两端的平移微分均可逆，故秩处处相同，等于 \(\dim\mathfrak g-\dim\mathfrak h\)。最后，\(k\) 固定 \(g\cdot p\) 当且仅当 \(g^{-1}kg\) 固定 \(p\)，这给出稳定子共轭公式。
\end{proof}

这个命题给出了轨道映射的秩及其核。核中的方向来自稳定子，因此可以先除去稳定子，把轨道参数化为一个陪集空间。常秩定理控制每个群元素附近的像；群中远处的元素仍可能把轨道送回当前邻域，所以陪集空间的流形拓扑与轨道像的子空间拓扑需要分别考察。

\subsection{齐性空间的光滑结构}
\label{g1:subsec:1701:03}

设 \(H\subseteq G\) 为闭子群。记 \(G/H\) 为左陪集 \(gH\) 的集合，赋予商映射 \(q(g)=gH\) 所决定的商拓扑。这里把同一个右 \(H\) 轨道压成一点；如果 \(H\) 不正规，不能把 \(G/H\) 当作商群。

\begin{theorem}[闭子群的齐性空间]\label{g1:thm:homogeneous-space}
	设 \(G\) 为有限维实 Lie 群，\(H\subseteq G\) 为闭子群。\(G/H\) 具有唯一的无边界、Hausdorff、第二可数光滑流形结构，使 \(q:G\to G/H\) 为光滑淹没。这个结构具有原来的商拓扑，维数为 \(\dim G-\dim H\)，左作用 \(G\times G/H\to G/H\) 光滑。\(q\) 在每点附近有光滑局部截面；对任意流形 \(N\) 和集合映射 \(f:G/H\to N\)，
	\[
		f\text{ 光滑}\quad\Longleftrightarrow\quad f\circ q\text{ 光滑}.
	\]
\end{theorem}
\begin{proof}
	先处理拓扑。若 \(A\subseteq G\) 开，则 \(q^{-1}(q(A))=AH=\bigcup_{h\in H}Ah\) 开，所以 \(q\) 是开映射。不同陪集 \(xH,yH\) 满足 \(x^{-1}y\notin H\)。由 \(H\) 闭及乘法连续，可取 \(x,y\) 的开邻域 \(A,B\)，使 \(A^{-1}B\cap H=\varnothing\)。于是 \(q(A)\) 与 \(q(B)\) 不交，证明商空间 Hausdorff。\(G\) 的一个可数开基经开映射 \(q\) 取像，就是 \(G/H\) 的可数开基。

	记 \(\mathfrak h=\mathrm{T}_eH\)，取线性补空间 \(\mathfrak g=\mathfrak m\oplus\mathfrak h\)。考虑
	\[
		F:\mathfrak m\times H\longrightarrow G,
		\qquad F(X,h)=\exp X\,h.
	\]
	它在 \((0,e)\) 的微分是 \((X,Y)\mapsto X+Y\)，为同构。逆函数定理给出 \(0,e\) 的邻域 \(W_0\subseteq\mathfrak m\)、\(V\subseteq H\)，使 \(F\) 在 \(W_0\times V\) 上是到开集的微分同胚。取更小的 \(W\subseteq W_0\)，使 \((\mathrm{d}F)_{(X,e)}\) 对 \(X\in W\) 都可逆，并且
	\[
		X,Y\in W,\quad \exp(-Y)\exp X\in H
		\quad\Longrightarrow\quad \exp(-Y)\exp X\in V.
	\]
	这一缩小可行，因为 \(H\) 取子空间拓扑，存在 \(G\) 中单位元邻域 \(D\) 满足 \(H\cap D\subseteq V\)，而 \(\exp(-Y)\exp X\) 在 \(X,Y\to0\) 时趋向 \(e\)。

	我们断言 \(F:W\times H\to (\exp W)H\) 是微分同胚。若 \(\exp X\,h=\exp Y\,k\)，则 \(u=\exp(-Y)\exp X=kh^{-1}\in V\)，从而 \(F(X,e)=F(Y,u)\)。局部单射性给出 \(X=Y\)、\(u=e\)，继而 \(h=k\)。故 \(F\) 单射。它在 \((X,e)\) 的微分可逆，利用第二因子和目标上的右平移，便知在 \((X,h)\) 处也可逆。因此它处处为局部微分同胚，像 \((\exp W)H\) 开；结合单射性，得到断言。

	现在 \(X\mapsto q(\exp X)\) 是 \(W\) 到开集 \(q((\exp W)H)\) 的同胚：在上述乘积坐标中，\(q\) 正是忘掉 \(H\) 因子的投影，故其逆坐标连续。把这张图用左平移移到任意 \(gH\) 附近，得到局部参数 \(X\mapsto q(g\exp X)\)。在两张图的交叠处，坐标转换由
	\[
		X\longmapsto k^{-1}g\exp X
		\longmapsto F_k^{-1}(k^{-1}g\exp X)
		\longmapsto\text{第一分量}
	\]
	给出，其中 \(F_k\) 是第二张图所用的横截乘积。各箭头均光滑，因而得到光滑图册。它的维数为 \(\dim\mathfrak m\)，且 \(q\) 在乘积坐标下为投影，所以是淹没。

	局部截面为 \(q(g\exp X)\mapsto g\exp X\)，记为 \(\sigma\)。若 \(f\circ q\) 光滑，则 \(f=(f\circ q)\circ\sigma\) 在每个这样的邻域光滑；反向由复合直接得到。左作用在局部写成 \((g,x)\mapsto q(g\sigma(x))\)，故光滑。

	最后，若另一光滑结构也使 \(q\) 为淹没，则由\cref{g1:thm:immersion-submersion-normal-form}，\(q\) 对该结构同样有光滑局部截面，并且是开映射。其拓扑因此仍为商拓扑。用两种结构各自的局部截面，恒等映射在两个方向都可写成一个局部截面与另一结构下的 \(q\) 的复合，故双向光滑，证明唯一性。
\end{proof}

定理中的最后一个性质说明：若先在 \(G\) 上定义光滑映射 \(F\)，并且证明它在每个陪集上常值，那么下降后的映射不必重新逐张坐标检查。下面的交换图同时表示 \(F=f\circ q\) 与 \(f\) 的唯一性：
\[
	\begin{tikzcd}
		G \arrow[r,"F"] \arrow[d,"q"'] & N \\
		G/H \arrow[ur,dashed,"f"'] &
	\end{tikzcd}
\]
局部截面把 \(F\) 的光滑性传递给 \(f\)。轨道映射恰好在稳定子的陪集上常值，因而可以用这个下降性质把它分解为陪集投影与轨道参数化。

\begin{proposition}[轨道的浸入结构]\label{g1:prop:orbit-immersion}
	设有限维实 Lie 群 \(G\) 光滑左作用于无边界光滑流形 \(M\)，\(p\in M\)。映射
	\[
		\overline\alpha_p:G/G_p\longrightarrow M,
		\qquad gG_p\longmapsto g\cdot p
	\]
	是单射浸入，像为 \(G\cdot p\)。把其定义域的光滑结构传给轨道后，
	\begin{equation}\label{g1:eq:orbit-tangent-space}
		\mathrm{T}_p(G\cdot p)=\{X_M(p):X\in\mathfrak g\}
		\cong\mathfrak g/\mathfrak g_p.
	\end{equation}
\end{proposition}
\begin{proof}
	集合上的单射性和像的描述正是轨道与陪集的对应。\(\alpha_p=\overline\alpha_p\circ q\)，所以前一定理的下降性质使 \(\overline\alpha_p\) 光滑。\((\mathrm{d}q)_e\) 满射，其核为 \(\mathfrak g_p\)；\cref{g1:prop:orbit-map-rank}给出 \((\mathrm{d}\alpha_p)_e\) 的同一个核。因此诱导的 \((\mathrm{d}\overline\alpha_p)_{eG_p}\) 单射。左平移和作用微分把此结论传到每点，证明浸入，并得到切空间公式。
\end{proof}

浸入轨道可能具有不同于像的子空间拓扑。例如 \(\mathbb R\) 按 \(t\cdot(x,y)=(x+t,y+\lambda t)\) 作用于 \(\mathbb R^2/\mathbb Z^2\)，其中 \(\lambda\) 无理。稳定子为零，轨道是\secref{g1:sec:1603}讨论过的稠密无理斜率线：其自身拓扑是一条实直线的拓扑，而远处参数会反复进入像中的任一小邻域。下一节的固有性将控制这种返回现象。

\ProseParagraphBreak

若 \(M\ne\varnothing\) 且只有一个轨道，则称该作用为\term[chuandizuoyong]{传递作用}[transitive action]；带有传递光滑 \(G\) 作用的流形称为 \(G\) 的\term[qixingkongjian]{齐性空间}[homogeneous space]。此时 \(G/G_p\) 与 \(M\) 集合上一一对应。传递性还使轨道映射满射；结合常秩定理与维数比较，可以进一步证明这个对应是微分同胚。

\begin{proposition}[传递作用的模型]\label{g1:prop:transitive-action-model}
	设有限维实 Lie 群 \(G\) 在非空无边界光滑流形 \(M\) 上传递光滑左作用，\(p\in M\)，\(G_p=\{g\in G:g\cdot p=p\}\)。则映射 \(\overline\alpha_p:G/G_p\to M\)、\(\overline\alpha_p(gG_p)=g\cdot p\) 为微分同胚。
\end{proposition}
\begin{proof}
	先说明满射的常秩映射 \(F:P\to M^n\) 必有秩 \(n\)，其中 \(P,M\) 为本章约定的流形。若秩 \(r<n\)，常秩定理给出定义域和像空间的局部坐标，使每个局部像包含在一个 \(r\) 维坐标平面中。用第二可数性及局部紧性，可取可数个紧坐标小球 \(K_j\) 覆盖 \(P\)，每个 \(K_j\) 都落在一张这样的定义域图内。于是 \(F(K_j)\) 在 \(M\) 中紧且闭，并因包含在低维坐标平面中而内部为空。

	非空流形不能被这些闭且内部为空的集合覆盖。这里可在 \(M\) 的一张坐标图内直接证明：取一个有内部的紧 Euclid 球，在其内部选择避开 \(F(K_1)\) 的更小闭球；再在后者内部选择避开 \(F(K_2)\) 的闭球，依次进行，并令半径趋于零。每一步可行，因为这些闭集内部为空。嵌套紧球有公共点，该点避开所有 \(F(K_j)\)，与覆盖矛盾。因此 \(F\) 的秩必须为 \(n\)。

	传递性使 \(\alpha_p:G\to M\) 满射，前面已证明其常秩，故它是淹没。于是 \(\dim(G/G_p)=\dim M\)。\(\overline\alpha_p\) 是同维流形之间的双射浸入，微分处处可逆，逆函数定理使其逆在每点附近光滑，从而它是微分同胚。
\end{proof}

为了以后计算不变度量，还需要知道稳定子怎样作用在一个切空间上。若 \(h\in G_p\)，则 \(a_h(p)=p\)，故 \((\mathrm{d}a_h)_p\) 是 \(\mathrm{T}_pM\) 的线性自同构。这给出\term[wendingzibiaoshi]{稳定子表示}[isotropy representation]
\[
	\rho_p:G_p\longrightarrow\operatorname{GL}(\mathrm{T}_pM),\qquad
	\rho_p(h)=(\mathrm{d}a_h)_p.
\]
链式法则保证 \(\rho_p(hk)=\rho_p(h)\rho_p(k)\)。它也保持轨道切空间，因此在法商空间 \(\mathrm{T}_pM/\mathrm{T}_p(G\cdot p)\) 上诱导线性作用。

\begin{proposition}[齐性空间的切表示]\label{g1:prop:homogeneous-tangent-isotropy}
	设 \(G\) 为有限维实 Lie 群，\(H\subseteq G\) 为闭子群，\(o=eH\in G/H\)，\(\mathfrak h=\mathrm{T}_eH\)。商映射的微分诱导同构 \(\mathrm{T}_o(G/H)\cong\mathfrak g/\mathfrak h\)。在此同构下，\(h\in H\) 的稳定子表示为
	\[
		\rho_o(h)(X+\mathfrak h)=\operatorname{Ad}_hX+\mathfrak h.
	\]
\end{proposition}
\begin{proof}
	\(q\) 是淹没，且局部乘积表明 \(\ker(\mathrm{d}q)_e=\mathfrak h\)，所以先得切空间同构。共轭保持 \(H\)，故 \(\operatorname{Ad}_h\) 保持 \(\mathfrak h\)，所写商空间算子有定义。再由
	\[
		h\exp(tX)H=\exp(t\operatorname{Ad}_hX)H
	\]
	在 \(t=0\) 求导，即得表示公式。
\end{proof}

\subsection{等变映射及其微分}
\label{g1:subsec:1701:04}

两个带有群作用的空间之间，普通光滑映射可以任意改变轨道。若希望映射保留对称性，就需要它与作用相容。设 \(G\) 作用于 \(M\)、\(K\) 作用于 \(N\)，\(\varphi:G\to K\) 为光滑同态。若光滑映射 \(F:M\to N\) 满足
\[
	F(g\cdot p)=\varphi(g)\cdot F(p).
\]
则称它为关于 \(\varphi\) 的\term[dengbianyingshe]{等变映射}[equivariant map]。若 \(G=K\)、\(\varphi=\operatorname{id}\)，就简称 \(G\) 等变。等变映射由一点的取值控制整条轨道，因此稳定子也会对这个取值施加限制。

\begin{proposition}[等变映射的微分]\label{g1:prop:equivariant-map-differential}
	设有限维实 Lie 群 \(G,K\) 分别光滑左作用于无边界光滑流形 \(M,N\)，作用映射记为 \(a_g(p)=g\cdot p\)、\(b_k(q)=k\cdot q\)。设 \(\varphi:G\to K\) 为光滑同态，\(F:M\to N\) 为满足 \(F(g\cdot p)=\varphi(g)\cdot F(p)\) 的光滑映射。对 \(p\in M\)，记两作用的稳定子为 \(G_p,K_{F(p)}\)，切表示为 \(\rho_p(h)=(\mathrm{d}a_h)_p\)、\(\rho_{F(p)}(k)=(\mathrm{d}b_k)_{F(p)}\)。则 \(\varphi(G_p)\subseteq K_{F(p)}\)，且对 \(X\in\operatorname{Lie}(G)\) 有
	\begin{equation}\label{g1:eq:equivariant-infinitesimal}
		(\mathrm{d}F)_pX_M(p)=((\mathrm{d}\varphi)_eX)_N(F(p)).
	\end{equation}
	若 \(h\in G_p\)，则其切表示满足
	\[
		(\mathrm{d}F)_p\circ\rho_p(h)
		=\rho_{F(p)}(\varphi(h))\circ(\mathrm{d}F)_p.
	\]
\end{proposition}
\begin{proof}
	若 \(h\cdot p=p\)，等变式给出 \(\varphi(h)\cdot F(p)=F(p)\)，故有稳定子包含。把 \(g=\exp(tX)\) 代入，并用指数自然性 \(\varphi(\exp(tX))=\exp(t(\mathrm{d}\varphi)_eX)\)，对 \(t\) 求导得\eqref{g1:eq:equivariant-infinitesimal}。最后对 \(F\circ a_h=a_{\varphi(h)}\circ F\) 在 \(p\) 求微分，得到切表示的相容式。
\end{proof}

在齐性空间上，等变映射由一个基点的像决定：其余点都由群作用得到。要使这个规定与基点的不同群参数表示相容，基点的稳定子必须固定所选像点。这样，构造映射的问题就转成了稳定子之间的包含关系。

\begin{proposition}[齐性空间之间的等变映射]\label{g1:prop:equivariant-homogeneous-maps}
	设 \(G\) 为有限维实 Lie 群，\(H,K\subseteq G\) 为闭子群。\(G/H,G/K\) 取使自然商投影为光滑淹没的商流形结构，并取 \(G\) 的标准左作用。\(G\) 等变光滑映射 \(G/H\to G/K\) 与满足 \(a^{-1}Ha\subseteq K\) 的陪集 \(aK\) 一一对应，映射为
	\[
		F_a(gH)=gaK.
	\]
	特别地，若 \(H\subseteq K\)，则 \(F_e:gH\mapsto gK\) 为光滑淹没，在基点的微分是自然满射 \(\mathfrak g/\mathfrak h\to\mathfrak g/\mathfrak k\)。
\end{proposition}
\begin{proof}
	若 \(F\) 等变且 \(F(eH)=aK\)，则 \(F(gH)=gaK\)，所以取值已被确定。\(H\) 固定 \(eH\)，故必须固定 \(aK\)，即 \(a^{-1}Ha\subseteq K\)。反之，在此包含条件下，\(g'=gh\) 给出 \(g'aK=ga(a^{-1}ha)K=gaK\)，所以公式与代表元无关。\(g\mapsto gaK\) 光滑，在 \(H\) 陪集上常值，由\cref{g1:thm:homogeneous-space}下降为光滑映射；等变性从公式直接得到。两个这样的映射相同，当且仅当它们在 \(eH\) 处的取值 \(aK\) 相同。

	若 \(H\subseteq K\)，关系 \(q_K=F_e\circ q_H\) 在单位元求微分，得所写商向量空间满射。左作用将满射性传到各点，故 \(F_e\) 是淹没。
\end{proof}

上述微分有清楚的几何含义：从 \(G/H\) 再压到 \(G/K\)，就是把原来仍能看见的 \(\mathfrak k/\mathfrak h\) 方向也视为等价。它的核给出被新等价关系消去的方向；保留下来的方向才是目标空间的切向量。下一节将把这种局部投影结构推广到不必传递的群作用。

% !TeX root = ../../Book1.tex
% 纲要标题：### 17.2 自由固有作用的商定理
% 纲要位置：几何学计划纲要.md，第 779 行。
% 正文状态：初稿；尚未正式审核。

% 纲要内容开始
% 
% * 固有性的局部和序列判据
% * 局部横截面与局部乘积
% * 商的 Hausdorff 性和光滑图册
% * 商映射的泛性质
% * 缺少自由或固有条件的反例
% 
% 纲要内容结束

\section{自由固有作用的商定理}
\label{g1:sec:1702}

把同一轨道上的点合并以后，横向方向应当留下来，沿轨道运动的方向则应被忘掉。自由性保证各条轨道具有同一个维数，但它不能阻止一条轨道反复回到很小的区域。为了让横向坐标真正成为商空间的坐标，还需要控制这种返回。本节的固有性正是一个同时限制出发点、到达点和群参数的条件。

\subsection{固有性的局部和序列判据}
\label{g1:subsec:1702:01}

\begin{definition}[固有作用]\label{g1:def:proper-lie-group-action}
	设有限维实 Lie 群 \(G\) 光滑左作用于有限维无边界光滑流形 \(M\)，两者均为 Hausdorff、第二可数。若映射
	\[
		\Theta:G\times M\longrightarrow M\times M,\qquad
		\Theta(g,p)=(gp,p)
	\]
	为固有映射，即每个紧集的原像都紧，则称这个作用为\term[guyouzuoyong]{固有作用}[proper action]。
\end{definition}

这里采用的是\cref{g1:def:proper-map-manifolds}的固有映射概念。需要约束的映射是 \(\Theta\)，而非单独的作用映射 \((g,p)\mapsto gp\)：后者忘掉出发点，即使直线上的平移也不满足同样的紧集原像条件。对两个子集 \(A,B\subseteq M\)，记
\[
	T(A,B)=\{g\in G:gA\cap B\ne\varnothing\}.
\]
它收集所有能够把 \(A\) 中某点送入 \(B\) 的群元素，下面称为从 \(A\) 到 \(B\) 的搬运集。这个记号中的先后次序与 \(\Theta\) 的两个输出次序不同：\(\Theta^{-1}(B\times A)\) 投影到群因子，才得到 \(T(A,B)\)。

\begin{proposition}[固有性的判据]\label{g1:prop:proper-action-criteria}
	设有限维实 Lie 群 \(G\) 光滑作用于有限维无边界光滑流形 \(M\)，且两者均为 Hausdorff、第二可数。对 \(A,B\subseteq M\)，令 \(T(A,B)=\{g\in G:gA\cap B\ne\varnothing\}\)。下列条件等价。
	\begin{equivalences}
		\item 作用固有。
		\item 对任意紧集 \(K,L\subseteq M\)，\(T(K,L)\) 在 \(G\) 中紧。
		\item 对任意 \(p,q\in M\)，存在开邻域 \(U\ni p\)、\(V\ni q\)，使 \(T(U,V)\) 在 \(G\) 中的闭包紧。
		\item 若 \(p_j\to p\)、\(g_jp_j\to q\) 在 \(M\) 中收敛，则 \((g_j)\) 有收敛子列。
	\end{equivalences}
	第二个条件也可只对 \(K=L\) 检查。
\end{proposition}
\begin{proof}
	先说明本证明使用序列的依据。流形有可数局部基，紧子集中的任意序列都有收敛子列：若无无限重复值，紧性给出一个聚点，再沿其递减局部基逐次选出子列。反过来，在第二可数空间的子集中，若每个序列都有收敛子列，则该子集紧。否则一个没有有限子覆盖的开覆盖可缩成可数子覆盖 \(O_1,O_2,\ldots\)，逐次选取 \(x_j\notin O_1\cup\cdots\cup O_j\)。任何收敛子列的极限都落在某个 \(O_k\)，从而子列最终也落在 \(O_k\)，矛盾。这使以下序列论证适用于本命题，而无须把它误用到任意拓扑空间。

	I 蕴含 II，因为 \(\Theta^{-1}(L\times K)\) 紧，而它在 \(G\) 中的投影就是 \(T(K,L)\)。反过来，设 II 成立，取紧集 \(C\subseteq M\times M\)，其两个投影分别记为 \(L,K\)。那么
	\[
		\Theta^{-1}(C)\subseteq T(K,L)\times K.
	\]
	右侧紧；\(C\) 在 Hausdorff 空间中闭，故左侧是右侧的闭子集，也紧，证明 I。若仅已知每个 \(T(K,K)\) 紧，则 \(T(K,L)\subseteq T(K\cup L,K\cup L)\)。此外 \(T(K,L)\) 总是闭的：若 \(g_j\to g\) 且 \(g_jx_j\in L\)、\(x_j\in K\)，紧性允许取 \(x_j\to x\in K\) 的子列，连续性给出 \(gx\in L\)，从而 \(g\in T(K,L)\)。流形中闭性可由序列检验，所以紧集的这个闭子集仍紧。

	II 蕴含 III，只须在 \(p,q\) 的坐标图中各取一个闭包紧的小开球。若反过来已知 III，对紧集 \(K,L\)，用相应乘积邻域 \(U\times V\) 覆盖 \(K\times L\)，再取有限子覆盖 \(U_i\times V_i\)。每个 \(g\in T(K,L)\) 都有见证点 \(x\in K\)、\(gx\in L\)；这对点落在某个 \(U_i\times V_i\)，所以
	\[
		T(K,L)\subseteq\bigcup_i\overline{T(U_i,V_i)}.
	\]
	右侧是有限个紧集之并。刚才已证左侧闭，故它紧。这也解释了为什么局部条件要允许任意一对 \(p,q\)，而不只检查同一点附近。

	I 蕴含 IV：序列 \((g_jp_j,p_j)\) 连同其极限组成紧集，固有性使全部 \((g_j,p_j)\) 落在一个紧集内，再取收敛子列即可。反之，若 IV 成立，对任意紧集 \(C\subseteq M\times M\) 及任意序列 \((g_j,p_j)\in\Theta^{-1}(C)\)，先由 \(C\) 紧取子列，使 \(p_j\to p\)、\(g_jp_j\to q\)。再由 IV 取 \(g_j\to g\) 的子列，便有 \((g_j,p_j)\to(g,p)\)。由于 \(C\) 闭且 \(\Theta\) 连续，这个极限仍在 \(\Theta^{-1}(C)\) 中。因此原像序列紧，按开头的说明即为紧集。
\end{proof}

若 \(G\) 离散，紧子集只有有限个点，于是第二个判据成为“每个紧集只被有限个群元素搬回与自身相交的位置”。这恢复了\secref{g1:sec:1003}讨论离散等距群时采用的固有性，但此时仍允许有限而非平凡的稳定子。

\begin{proposition}[固有作用的稳定子与轨道]\label{g1:prop:proper-action-orbits}
	设有限维实 Lie 群 \(G\) 固有地光滑作用于有限维无边界光滑流形 \(M\)，两者均为 Hausdorff、第二可数。则每个稳定子 \(G_p\) 紧，每条轨道 \(Gp\) 为 \(M\) 中的闭嵌入子流形，而且轨道等价关系
	\[
		R=\{(q,p)\in M\times M:q\in Gp\}
	\]
	在 \(M\times M\) 中闭。
\end{proposition}
\begin{proof}
	稳定子恰为 \(T(\{p\},\{p\})\)，因此紧。若 \((q_j,p_j)\in R\) 收敛到 \((q,p)\)，选取 \(g_j\) 使 \(q_j=g_jp_j\)。由序列判据，取子列后 \(g_j\to g\)，于是 \(q=gp\)，所以 \((q,p)\in R\)。这证明 \(R\) 闭。固定第二个变量为 \(p\)，\(Gp\) 就是连续映射 \(q\mapsto(q,p)\) 下 \(R\) 的原像，因而闭。

	为说明嵌入性，对任意紧集 \(C\subseteq M\)，轨道映射的原像 \(\alpha_p^{-1}(C)=T(\{p\},C)\) 紧。令 \(q:G\to G/G_p\)，则 \(\overline\alpha_p^{-1}(C)=q(\alpha_p^{-1}(C))\) 也紧。因此\cref{g1:prop:orbit-immersion}的单射浸入 \(\overline\alpha_p:G/G_p\to M\) 固有，由\cref{g1:prop:proper-injective-immersion}，它是闭像嵌入。这也说明在固有作用下，轨道的自然流形拓扑就是环境子空间拓扑。
\end{proof}

\begin{proposition}[紧群作用]\label{g1:prop:compact-group-action-proper}
	设紧的有限维实 Lie 群 \(G\) 光滑作用于有限维、Hausdorff、第二可数的无边界光滑流形 \(M\)。则作用固有。
\end{proposition}
\begin{proof}
	对紧集 \(C\subseteq M\times M\)，记第二个投影像为 \(K\)。\(\Theta^{-1}(C)\) 是紧集 \(G\times K\) 的闭子集，所以紧。
\end{proof}

\subsection{局部横截面与局部乘积}
\label{g1:subsec:1702:02}

在自由作用下，\cref{g1:prop:orbit-map-rank}给出单射 \(\mathfrak g\to \mathrm{T}_pM\)、\(X\mapsto X_M(p)\)，其像是轨道方向。选一个线性补空间，再用流形坐标把它实现成经过 \(p\) 的小子流形 \(S\)。这份\term[jubuhengjiemian]{局部横截面}[local transversal]的任务是为附近每条轨道提供一个代表。逆函数定理首先只保证小群参数内的代表唯一；固有性还将排除其他群元素把横截面送回自身。

\begin{lemma}[自由固有作用的局部乘积]\label{g1:lem:free-proper-local-product}
	设有限维实 Lie 群 \(G\) 自由且固有地光滑作用于 \(m\) 维无边界光滑流形 \(M\)，两者均为 Hausdorff、第二可数。令 \(r=\dim G\)。对每个 \(p\in M\)，存在经过 \(p\) 的 \(m-r\) 维嵌入子流形 \(S\)，使
	\begin{equation}\label{g1:eq:free-action-local-product}
		\Phi:G\times S\longrightarrow GS,\qquad \Phi(g,s)=gs
	\end{equation}
	为微分同胚，且饱和集 \(GS=\{gs:g\in G,\ s\in S\}\) 在 \(M\) 中开。每条与 \(S\) 相交的轨道只与 \(S\) 相交一次。
\end{lemma}
\begin{proof}
	先取 \(\mathrm{T}_pM\) 的直和分解
	\[
		\mathrm{T}_pM=\{X_M(p):X\in\mathfrak g\}\oplus E.
	\]
	在以 \(p\) 为原点的坐标图中，对坐标再作可逆线性变换，就可令最后 \(m-r\) 个坐标方向对应 \(E\)。前 \(r\) 个坐标置零，得到经过 \(p\) 的嵌入子流形 \(S_0\)，且 \(\mathrm{T}_pS_0=E\)。作用映射限制为 \(\Phi:G\times S_0\to M\)，其在 \((e,p)\) 的微分是
	\[
		\mathrm{d}\Phi_{(e,p)}(X,v)=X_M(p)+v.
	\]
	自由性和上述直和使它为同构。由\cref{g1:thm:inverse-function}，可以选取单位元开邻域 \(U\subseteq G\) 与 \(p\) 在 \(S_0\) 中的开邻域 \(S_1\)，使 \(\Phi|_{U\times S_1}\) 为到 \(M\) 中一个开集的微分同胚。此时 \((e,s)\)、\(s\in S_1\)，也都具有可逆微分。

	下一步是控制 \(U\) 外的群元素。我们断言，存在 \(p\) 在 \(M\) 中的开邻域 \(V\)，使
	\begin{equation}\label{g1:eq:proper-action-no-return}
		gV\cap V\ne\varnothing\quad\Longrightarrow\quad g\in U.
	\end{equation}
	若否，在 \(p\) 的一列递减坐标邻域 \(V_j\) 中，可选 \(s_j,t_j\to p\) 及 \(g_j\notin U\)，使 \(g_js_j=t_j\)。固有性的序列判据给出收敛子列 \(g_j\to g\)。取极限得 \(gp=p\)，自由性于是迫使 \(g=e\)。由于 \(U\) 是 \(e\) 的开邻域，这与 \(g_j\notin U\) 矛盾，证明断言。

	取 \(S=S_1\cap V\)，必要时再缩小为其中经过 \(p\) 的坐标邻域。如果 \(gs=t\)，其中 \(s,t\in S\)，则\eqref{g1:eq:proper-action-no-return}说明 \(g\in U\)。现在
	\[
		\Phi(g,s)=t=\Phi(e,t),
	\]
	而 \(\Phi\) 在 \(U\times S_1\) 上单射，故 \(g=e\)、\(s=t\)。若一般地 \(gs=ht\)，则 \(h^{-1}gs=t\)，由同一结论得 \(g=h\)、\(s=t\)。因此整个 \(G\times S\) 上的 \(\Phi\) 单射，而不只是小群邻域上的单射。

	还须说明这个全局于群因子的映射具有光滑逆。固定 \(g\in G\)，由
	\[
		\Phi(gk,s)=g\Phi(k,s)
	\]
	可知 \(\mathrm{d}\Phi_{(g,s)}\) 是 \(\mathrm{d}\Phi_{(e,s)}\) 与两侧平移微分同构的复合，所以处处可逆。逆函数定理使 \(\Phi\) 为局部微分同胚，因而为开映射，其像 \(GS\) 开。它到像又是双射，各局部光滑逆在交叠处由唯一性相同，便拼成全局光滑逆。
\end{proof}

\cref{g1:fig:free-action-product}显示这里的两个参数：群参数沿轨道变化，\(S\) 上的参数在不同轨道之间变化。图中画出的只是 \(G\) 的一小块；定理中的乘积却包含整个 \(G\)。因此它并不声称轨道本身像直线，也不要求 \(G\) 连通或紧。

\begin{figure}[htbp]
	\centering
	\input{Figures/Chapter17-FreeActionProduct}
	\caption[自由固有作用的局部乘积]{局部乘积 \(G\times S\cong GS\) 中的轨道方向与横向方向。每条水平线表示一条轨道的一部分，竖线表示横截面 \(\{e\}\times S\)；商映射忘掉群参数，只保留横向坐标。图中直线仅表示乘积坐标，不表示原流形中的轨道形状。}
	\label{g1:fig:free-action-product}
\end{figure}

\subsection{商的 Hausdorff 性和光滑图册}
\label{g1:subsec:1702:03}

给轨道集 \(M/G\) 赋予商拓扑，记 \(\pi(p)=[p]=Gp\)。对任意开集 \(O\subseteq M\)，有
\[
	\pi^{-1}(\pi(O))=GO=\bigcup_{g\in G}gO,
\]
它是开集，故 \(\pi\) 为开映射。这一点对任意连续群作用都成立；要得到流形，还须证明不同轨道能够分开，并为这些轨道配上相容的光滑坐标。

\begin{theorem}[自由固有作用的商流形]\label{g1:thm:free-proper-quotient}
	设 \(G\) 为 \(r\) 维实 Lie 群，\(M\) 为 \(m\) 维无边界光滑流形，两者均为 Hausdorff、第二可数。若 \(G\) 自由且固有地光滑作用于 \(M\)，则商拓扑下的 \(M/G\) 具有唯一的 \(m-r\) 维无边界光滑流形结构，使投影
	\[
		\pi:M\longrightarrow M/G
	\]
	成为光滑满射淹没。在\cref{g1:lem:free-proper-local-product}的每个局部乘积 \(G\times S\cong GS\) 中，\(\pi\) 对应第二因子投影，且 \(\pi|_S:S\to\pi(GS)\) 是微分同胚。
\end{theorem}
\begin{proof}
	先核验商空间的分离性与可数性。若 \([p]\ne[q]\)，则 \((q,p)\) 不在闭关系集 \(R\) 中。由\cref{g1:prop:proper-action-orbits}，可取开邻域 \(U\ni p\)、\(V\ni q\)，使 \((V\times U)\cap R=\varnothing\)。这就是说没有轨道同时经过 \(U\) 和 \(V\)，所以 \(\pi(U)\)、\(\pi(V)\) 是两个不交的开邻域，分别包含 \([p]\)、\([q]\)。故商空间 Hausdorff。取 \(M\) 的可数开基 \((B_j)\)，开集 \(\pi(B_j)\) 构成商空间的可数基：对 \([p]\in W\) 且 \(W\) 开，选 \(p\) 的基邻域 \(B_j\subseteq\pi^{-1}(W)\)，即可得到 \([p]\in\pi(B_j)\subseteq W\)。

	对局部横截面 \(S\)，记 \(W=\pi(GS)\)。\(GS\) 开且饱和，故 \(W\) 开；每条轨道与 \(S\) 交于唯一点，使 \(\pi|_S:S\to W\) 成为连续双射。它还开：若 \(A\subseteq S\) 开，则 \(G\times A\) 在 \(G\times S\) 中开，局部乘积微分同胚把它送到 \(GA\)，这是 \(M\) 的开集，所以 \(\pi(A)=\pi(GA)\) 在商中开。因此 \(\pi|_S\) 是同胚。取 \(S\) 上的光滑坐标 \(u:S\to\mathbb R^{m-r}\) 到一个开集，便得到 \(W\) 上的坐标 \(u\circ(\pi|_S)^{-1}\)。

	检验交叠时，设另一个横截面为 \(S'\)，局部乘积逆映射为 \(\Phi'^{-1}:GS'\to G\times S'\)。若 \(s\in S\cap GS'\)，唯一写成 \(s=g s'\)，其中 \(s'\in S'\)。用两张横截面表示同一轨道时，相应的转换正是
	\[
		S\cap GS'\longrightarrow S',\qquad
		s\longmapsto\operatorname{pr}_{S'}\bigl(\Phi'^{-1}(s)\bigr).
	\]
	定义域在 \(S\) 中开，右侧是光滑映射的复合，因此转换光滑；交换 \(S,S'\) 后得到光滑逆。于是上述坐标构成相容光滑图册，维数为 \(m-r\)，且没有边界。

	在乘积坐标中，\(\pi(gs)=\pi(s)\)，所以 \(\pi\) 就是 \((g,s)\mapsto s\)。它光滑且微分满射，因而为淹没；定义又使它满射。这些坐标还给出每个商点附近的光滑截面 \(\sigma:W\to M\)，即 \(\sigma([s])=s\)。

	最后证明唯一性。设同一商拓扑上还有另一份光滑结构，使同一个投影 \(\pi'\) 为光滑满射淹没。从刚构造的商到另一份商的恒等映射，在 \(W\) 上可写为 \(\pi'\circ\sigma\)，故光滑。反方向中，\cref{g1:thm:immersion-submersion-normal-form}说明 \(\pi'\) 的局部表达为投影；把被忘掉的坐标取常值，就得到目标点附近的光滑截面 \(\tau\)。从另一份商到本商的恒等映射在这片邻域上为 \(\pi\circ\tau\)，也光滑。两个恒等映射互逆，所以两份光滑结构相同。
\end{proof}

自由性使轨道方向始终具有维数 \(\dim G\)，固有性则既排除横截面的远处返回，也保证商的 Hausdorff 性。证明中第二可数性同样保留下来，不能只凭“每点附近看起来像 Euclid 空间”就称所得空间为本书意义的流形。

\begin{example}[正比例缩放的商]\label{g1:exm:positive-scaling-quotient}
	乘法 Lie 群 \((0,\infty)\) 作用于 \(\mathbb R^n\setminus\{0\}\)、\(n\ge1\)，作用为 \(a\cdot x=ax\)。它自由。若 \(x_j\to x\ne0\)、\(a_jx_j\to y\ne0\)，则
	\[
		a_j=\frac{\lVert a_jx_j\rVert}{\lVert x_j\rVert}
		\longrightarrow\frac{\lVert y\rVert}{\lVert x\rVert}>0,
	\]
	故作用固有。每条轨道恰与单位球面 \(S^{n-1}\) 相交一次，并且
	\[
		(0,\infty)\times S^{n-1}\longrightarrow\mathbb R^n\setminus\{0\},
		\qquad(a,u)\longmapsto au
	\]
	的光滑逆为 \(x\mapsto(\lVert x\rVert,x/\lVert x\rVert)\)。所以这里甚至有全局乘积，商就是 \(S^{n-1}\)，商投影对应 \(x\mapsto x/\lVert x\rVert\)。在 \(n=1\) 时，球面 \(S^0\) 为两个点；正负半轴分别是一条轨道，这也与零维商一致。
\end{example}

\subsection{商映射的泛性质}
\label{g1:subsec:1702:04}

拓扑商的泛性质见\cref{g1:prop:quotient-descent}：在每条轨道上取相同值的连续映射，才能成为商上的连续映射。现在局部横截面使同一结论提升到光滑范畴。实际使用时，这通常比直接计算商坐标更方便。

\begin{proposition}[光滑映射的下降]\label{g1:prop:quotient-smooth-universal-property}
	设有限维实 Lie 群 \(G\) 自由且固有地光滑作用于有限维无边界光滑流形 \(M\)，并设 \(N\) 为有限维无边界光滑流形；所有流形均为 Hausdorff、第二可数。给 \(M/G\) 以上一定理的光滑结构。则集合映射 \(h:M/G\to N\) 光滑，当且仅当 \(h\circ\pi:M\to N\) 光滑。

	特别地，光滑映射 \(f:M\to N\) 满足 \(f(gp)=f(p)\) 对所有 \(g,p\) 成立，当且仅当存在唯一光滑映射 \(\bar f:M/G\to N\)，使 \(f=\bar f\circ\pi\)。
\end{proposition}
\begin{proof}
	若 \(h\) 光滑，复合 \(h\circ\pi\) 光滑。反方向中，取商空间任意点附近的光滑截面 \(\sigma:W\to M\)，则
	\[
		h|_W=(h\circ\pi)\circ\sigma,
	\]
	所以它光滑。这样的 \(W\) 覆盖整个商，故 \(h\) 光滑；证明甚至不必预先假定 \(h\) 连续。

	若 \(f\) 在轨道上常值，定义 \(\bar f([p])=f(p)\)。轨道常值性保证定义不依赖代表元，满射性保证存在且唯一，刚证的光滑性判据保证 \(\bar f\) 光滑。反过来，任何复合 \(\bar f\circ\pi\) 在轨道上都取常值。
\end{proof}

这一泛性质可以画成
\[
	\begin{tikzcd}[column sep=large,row sep=large]
		M \arrow[r,"f\text{ 光滑且不变}"] \arrow[d,two heads,"\pi\text{ 淹没}"'] & N \\
		M/G \arrow[ur,dashed,"\bar f\text{ 唯一光滑}"'] &
	\end{tikzcd}
\]
图中的 \(N\) 是任意测试流形，不承担给定的 \(G\) 作用。若目标也带作用，等变映射未必沿源轨道取常值；要下降到轨道商，通常应同时把目标送到它的轨道商，不能直接套用这里的不变性条件。

\subsection{缺少自由或固有条件的反例}
\label{g1:subsec:1702:05}

\begin{example}[自由作用的非 Hausdorff 商]\label{g1:exm:irrational-flow-quotient}
	固定无理数 \(\alpha\)，令加法群 \(\mathbb R\) 作用于 \(T^2=\mathbb R^2/\mathbb Z^2\)：
	\[
		t\cdot[x,y]=[x+t,y+\alpha t].
	\]
	若某个 \(t\) 固定一点，则 \(t\in\mathbb Z\)、\(\alpha t\in\mathbb Z\)，所以 \(t=0\)，作用自由。\cref{g1:exm:irrational-torus-subgroup}已证明单位轨道是环面的真稠密子群；其他轨道是它的平移，也都稠密。取非空开集 \(O\subseteq T^2\)，每条轨道都与 \(O\) 相交，所以其饱和集 \(\mathbb R O\) 为整个环面。于是轨道商只有空集和全空间两个开集，而又有不止一条轨道，因此不是 Hausdorff 空间。

	这里固有性明确失败：紧集 \(K=T^2\) 的搬运集 \(T(K,K)=\mathbb R\) 非紧。每条轨道本身仍有一维浸入流形结构；失败的是把所有轨道合并以后，无法给它们彼此分离的邻域。
\end{example}

\begin{example}[旋转商的端点]\label{g1:exm:rotation-half-line-quotient}
	圆周群 \(S^1\) 通过复数乘法作用于 \(\mathbb C\cong\mathbb R^2\)。群紧，所以作用固有；原点的稳定子是整个 \(S^1\)，其他点的稳定子平凡，因此作用不自由。轨道由半径确定，\([z]\mapsto|z|\) 给出
	\[
		\mathbb C/S^1\cong[0,\infty).
	\]
	这是同胚：模长连续且在轨道上常值，故诱导连续映射；其逆 \(r\mapsto[r]\) 由正实半轴的连续包含给出。端点附近删去端点以后仍连通，而开区间中删去内点会出现两个分支，所以该商不是一维无边界流形。

	即使允许边界，商投影也不会满足前一定理的淹没结论。常用的光滑不变量 \(u(z)=|z|^2\) 在原点的微分为零。更一般地，若实值光滑函数 \(f\) 对旋转不变，则 \(f(-z)=f(z)\)，在原点求导得到 \(\mathrm{d}f_0=-\mathrm{d}f_0\)，因此 \(\mathrm{d}f_0=0\)。任何通过一维商坐标描述的光滑投影，在原点都不可能有秩 \(1\)。
\end{example}

非自由性也不意味着商一定出现坏点。例如 \(S^1\) 在任意流形 \(M\) 上的平凡作用是固有而不自由的，轨道商仍是 \(M\)，投影就是恒等映射。改变的是“轨道维数等于 \(\dim G\)”这一前提，故不能使用 \(\dim M-\dim G\) 的维数公式。下一步研究非自由作用时，需要保留稳定子及其横向作用，才能区分这种无害的冗余对称与旋转例中的端点。

% 纲要条目排为编号小节标题，并进入目录与书签；说明段仍保留于纲要注释中。

% !TeX root = ../../Book1.tex
% 纲要标题：### 17.3 不变结构与齐性模型
% 纲要位置：几何学计划纲要.md，第 787 行。
% 正文状态：初稿；尚未正式审核。

% 纲要内容开始
% 
% * 紧 Lie 群上的归一化不变积分
% 
%   以左不变密度及紧性构造积分，并用模函数的紧像证明右不变性；随即用于平均。本卷第19章进一步发展双不变度量，不反向承担此处引理。
% 
% * 紧群平均及不变度量
% * 稳定子表示与不变张量
% * 球面和射影空间的齐性表达
% * Grassmann 流形的切空间与维数
% 
% 纲要内容结束

\section{不变结构与齐性模型}
\label{g1:sec:1703}

球面的旋转既移动点，也保持各点的长度和夹角。一般群作用只规定怎样移动点，并不自动保持我们任意选取的度量。若群紧，可以把一个度量的所有变换平均起来；若作用传递，还可以只在一个点选取数据，再用群作用搬到全空间。这两种办法都需要解决选择无关性：平均要有不偏向某个群元素的积分，而沿群搬运要消除稳定子的影响。

\subsection{紧 Lie 群上的归一化不变积分}
\label{g1:subsec:1703:01}

有限群的平均是 \(|K|^{-1}\sum_{k\in K}\)。对圆周群，相应公式应为 \((2\pi)^{-1}\int_0^{2\pi}f(e^{it})\,\mathrm{d}t\)。两者都保持平移：重新选择群中的起点不改变平均值。一般紧 Lie 群没有一组整体坐标供我们直接积分，不过\cref{g1:prop:left-invariant-fields}已经说明，左平移能把单位元处的线性数据光滑地搬到全群；这次搬运的对象是一条正密度。

\begin{theorem}[紧群上的不变积分]\label{g1:thm:compact-group-integration}
	设 \(K\) 是有限维实紧 Lie 群，不要求连通。则 \(K\) 上存在唯一左不变的光滑正密度 \(\mu_K\)，满足 \(\int_K\mu_K=1\)。它还满足
	\[
		L_a^*\mu_K=\mu_K,\qquad R_a^*\mu_K=\mu_K,\qquad
		\iota^*\mu_K=\mu_K\quad(a\in K),
	\]
	其中 \(\iota(k)=k^{-1}\)。因此对任意连续实值函数 \(f\)，
	\begin{equation}\label{g1:eq:compact-group-integral-invariance}
		\begin{aligned}
			\int_K f(ak)\,\mathrm{d}\mu_K(k)
			&=\int_K f(ka)\,\mathrm{d}\mu_K(k)\\
			&=\int_K f(k^{-1})\,\mathrm{d}\mu_K(k)
			=\int_K f(k)\,\mathrm{d}\mu_K(k).
		\end{aligned}
	\end{equation}
	该积分线性、保持非负性，且 \(f\geq0\)、\(f\not\equiv0\) 时积分严格为正。
\end{theorem}
\begin{proof}
	先在 \(\mathfrak k=\mathrm{T}_eK\) 上任取一个正密度 \(\nu_e\)。它可以由任一正定内积的体积密度给出，见\cref{g1:prop:gram-volume-density}。对 \(v_1,\ldots,v_d\in \mathrm{T}_kK\) 定义
	\[
		\nu_k(v_1,\ldots,v_d)
		=\nu_e\bigl(\mathrm{d}(L_{k^{-1}})_kv_1,\ldots,
		\mathrm{d}(L_{k^{-1}})_kv_d\bigr),\qquad d=\dim K.
	\]
	左平移的微分及其逆随 \(k\) 光滑，所以 \(\nu\) 是光滑正密度。由 \(L_{(ak)^{-1}}\circ L_a=L_{k^{-1}}\)，有 \(L_a^*\nu=\nu\)。任意左不变密度都由单位元处的值唯一决定，而单位元处的密度空间是一维，故任意另一条左不变正密度都是 \(\nu\) 的正数倍。

	紧性使 \(\nu\) 本身具有紧支集。\cref{g1:thm:integral-well-defined}保证 \(c=\int_K\nu\) 有限且严格为正，故 \(\mu_K=\nu/c\) 是唯一满足归一化条件的左不变正密度。右不变性也可由左不变密度的唯一性得到：左、右平移交换，故 \(R_a^*\mu_K\) 也是左不变正密度。因而存在唯一正数 \(\Delta(a)\) 使
	\[
		R_a^*\mu_K=\Delta(a)\mu_K.
	\]
	在单位元的一组固定基上求值可知 \(\Delta\) 光滑。右平移复合及密度拉回的复合律给出 \(\Delta(ab)=\Delta(a)\Delta(b)\)。这是本处采用的\term[mohanshu]{模函数}[modular function]约定：记录右拉回对左不变密度的倍率。

	连续同态 \(\log\Delta:K\to(\mathbb R,+)\) 的像是紧加法子群。若其中有非零实数 \(r\)，它就含有所有整数倍 \(nr\)，与有界性矛盾。因此像只有零，\(\Delta\equiv1\)，得到右不变性。这一步同样适用于不连通的 \(K\)，没有把 \(K\) 换成单位元分支。

	由 \(\iota\circ L_a=R_{a^{-1}}\circ\iota\)，\(\iota^*\mu_K\) 是左不变正密度。\cref{g1:thm:oriented-change-variables}的密度换元公式表明 \(\int_K\iota^*\mu_K=\int_K\mu_K=1\)，故归一化的唯一性给出 \(\iota^*\mu_K=\mu_K\)。对三类微分同胚应用同一换元公式，就得到光滑 \(f\) 的\eqref{g1:eq:compact-group-integral-invariance}。

	连续 \(f\) 的积分也可直接用有限坐标覆盖和单位分解定义：逐片对连续紧支系数作 Euclid 积分。第14章证明中的坐标换元、共同细化与有限求和仍适用，所以积分和换元公式不变。非负性逐片成立；若 \(f(k_0)>0\)，连续性使它在一个坐标小邻域中有正下界，正密度便使这一块贡献正积分。零维时坐标积分就是点上的求和，论证也成立。
\end{proof}

以后写 \(\int_K f(k)\,\mathrm{d}\mu_K(k)\)，就是 \(\int_K f\mu_K\)，称为\term[guiyihuabubianjifen]{归一化不变积分}[normalized invariant integral]。这里积分的是密度 \(f\mu_K\)，无须选择定向。

\begin{example}\label{g1:exm:compact-integral-circle-finite}
	零维紧 Lie 群离散且紧，所以是有限群。左平移使各点的密度权重相同，归一化以后每点权重为 \(1/|K|\)，恰好恢复有限平均。对 \(S^1\)，角坐标中的密度 \((2\pi)^{-1}|\mathrm{d}t|\) 在换角度分支时相容，且被旋转保持；它的总积分为一，故就是上述 \(\mu_{S^1}\)。例如 \(m\in\mathbb Z\) 时
	\[
		\int_{S^1} z^m\,\mathrm{d}\mu_{S^1}(z)
		=\frac1{2\pi}\int_0^{2\pi}e^{imt}\,\mathrm{d}t
		=\begin{cases}1,&m=0,\\0,&m\ne0.\end{cases}
	\]
	复值积分按实部、虚部分别积分理解。紧性在一般定理中首先保证有限总量；在加法群 \(\mathbb R\) 上，左不变正密度必为 \(c|\mathrm{d}x|\)，其中 \(c>0\)，总积分无限，无法作同样的归一化。
\end{example}

\subsection{紧群平均及不变度量}
\label{g1:subsec:1703:02}

把光滑对象逐个平均，还须证明所得对象仍光滑。平均所用的群变量范围固定且紧，这正好提供了把参数微分移入积分所需的统一控制。以下引理是\subsecref{g1:subsec:1402:01}中紧区间参数积分论证的流形版本。

\begin{lemma}[紧参数积分]\label{g1:lem:compact-parameter-integral}
	设 \(K\) 是紧无边界光滑流形，\(\mu\) 为光滑密度，\(U\subset\mathbb R^m\) 开，\(F:K\times U\to\mathbb R\) 光滑。则
	\[
		I(x)=\int_K F(k,x)\,\mathrm{d}\mu(k)
	\]
	在 \(U\) 上光滑，且对每个多重指标 \(\alpha\)，有
	\[
		\partial_x^\alpha I(x)
		=\int_K\partial_x^\alpha F(k,x)\,\mathrm{d}\mu(k).
	\]
\end{lemma}
\begin{proof}
	给 \(K\) 取有限坐标覆盖和从属光滑单位分解。每一项都是对固定紧支集上的 Euclid 坐标变量积分，密度系数与单位分解均不依赖 \(x\)。固定 \(x_0\in U\)，取闭包紧且包含于 \(U\) 的小坐标盒 \(B\)。对第 \(j\) 个参数，微积分基本定理给出
	\[
		\frac{F(k,x+he_j)-F(k,x)}{h}
		=\int_0^1\partial_jF(k,x+the_j)\,\mathrm{d}t.
	\]
	当 \(x,x+he_j\in B\) 且 \(h\to0\) 时，右侧在紧参数集 \(K\times B\) 上局部一致趋于 \(\partial_jF(k,x)\)，因为该导数连续且在紧集上一致连续。各坐标片的密度系数绝对值可积，所以积分差的绝对值受“一致差上界乘有限积分”控制，差商可以移入积分。相同的一致控制还给出所得一阶导数的连续性。对任意 \(\partial_x^\alpha F\) 重复此论证，得到所有阶导数及其连续性。
\end{proof}

\begin{theorem}[平均度量]\label{g1:thm:compact-group-average}
	设紧 Lie 群 \(K\) 光滑地左作用于无边界光滑流形 \(M\)，记 \(a_k(p)=k\cdot p\)，\(\mu_K\) 为 \(K\) 上总质量为 \(1\) 的左右不变光滑正密度。给定 \(M\) 上任意光滑 Riemann 度量 \(g\)，公式
	\begin{equation}\label{g1:eq:averaged-metric}
		\overline g_p(v,w)
		=\int_K g_{k\cdot p}\bigl(\mathrm{d}(a_k)_pv,
		\mathrm{d}(a_k)_pw\bigr)\,\mathrm{d}\mu_K(k)
	\end{equation}
	定义一个 \(K\) 不变的光滑 Riemann 度量，即 \(a_h^*\overline g=\overline g\) 对每个 \(h\in K\) 成立。因此这样的作用总有不变 Riemann 度量，且不要求 \(M\) 紧。
\end{theorem}
\begin{proof}
	固定 \(p\) 时，各项都已拉回到 \(\mathrm{T}_pM\) 上的双线性型，故可以逐分量积分，而不是把不同切空间中的向量直接相加。双线性和对称性由积分的线性得到。若 \(v\ne0\)，微分同胚 \(a_k\) 的微分可逆，故 \(g_{k\cdot p}(\mathrm{d}a_kv,\mathrm{d}a_kv)>0\)。被积函数连续且处处正，\cref{g1:thm:compact-group-integration}给出 \(\overline g_p(v,v)>0\)。

	在 \(M\) 的一张坐标图内，把 \(v,w\) 依次取为坐标向量场。作用映射的微分随 \((k,p)\) 光滑，所得标量被积函数在 \(K\times U\) 上联合光滑。\cref{g1:lem:compact-parameter-integral}说明 \(\overline g\) 的每个局部分量光滑。最后，由 \(a_k\circ a_h=a_{kh}\) 得
	\[
		a_h^*\overline g
		=\int_K a_h^*a_k^*g\,\mathrm{d}\mu_K(k)
		=\int_K a_{kh}^*g\,\mathrm{d}\mu_K(k)
		=\overline g,
	\]
	最后一步用右不变积分。任意初始 \(g\) 的存在由\cref{g1:thm:bundle-metric-existence}用于 \(\mathrm{T}M\) 保证。
\end{proof}

同样可以平均协变张量、函数和微分形式，因为它们拉回后都属于固定纤维中的向量空间。对一般张量，平均未必保留非线性约束：例如一个非零向量场的平均可能为零，一个非退化二形式的平均也可能退化。度量的正定性能够保留，原因是正定双线性型构成凸集。

\begin{example}\label{g1:exm:rotation-average-metric}
	在 \(\mathbb R^2\) 上取常系数度量 \(g=a\,\mathrm{d}x^2+2b\,\mathrm{d}x\,\mathrm{d}y+c\,\mathrm{d}y^2\)，其中 \(a>0\)、\(ac-b^2>0\)。记其矩阵为 \(Q\)，平面旋转群的平均矩阵为
	\[
		\frac1{2\pi}\int_0^{2\pi}R_t^{\mathsf T}QR_t\,\mathrm{d}t
		=\frac{a+c}{2}I.
	\]
	这是因为 \(\cos^2t\)、\(\sin^2t\) 的平均均为 \(1/2\)，\(\sin t\cos t\) 的平均为零。平均消除了原来的优先方向，但保留了正定性。这里的底空间 \(\mathbb R^2\) 非紧，并不妨碍对紧的旋转参数积分。
\end{example}

\subsection{稳定子表示与不变张量}
\label{g1:subsec:1703:03}

在齐性空间 \(G/H\) 上，任何点都可写成 \(gH\)。若在基点 \(o=eH\) 选定一个内积，似乎只要用 \(g\) 把它搬到 \(gH\) 即可；但是同一点也等于 \(ghH\)。这两次搬运相差 \(h\) 在 \(\mathrm{T}_o(G/H)\) 上的作用，因此内积必须被整个稳定子保持。\cref{g1:prop:homogeneous-tangent-isotropy}已把这一作用写成
\[
	\mathrm{T}_o(G/H)\simeq\mathfrak g/\mathfrak h,\qquad
	\rho(h)(X+\mathfrak h)=\operatorname{Ad}_hX+\mathfrak h.
\]
此处 \(H\) 是闭 Lie 子群，可以不连通；只检查 \(\mathfrak h\) 的无穷小作用未必控制它的其余分支。

\begin{proposition}[不变张量的基点数据]\label{g1:prop:invariant-tensor-isotropy}
	设 \(G\) 为有限维实 Lie 群，\(H\subset G\) 为闭子群，\(o=eH\)，并给 \(G/H\) 取\cref{g1:thm:homogeneous-space}的光滑结构。对任意非负整数 \(r,s\)，取基点值给出如下两个实向量空间的线性同构：\(G/H\) 上的光滑 \(G\) 不变 \((r,s)\) 型张量场，以及 \(T^r_s(\mathrm{T}_o(G/H))\) 中被稳定子表示 \(\rho(H)\) 固定的张量。

	特别地，\(G\) 不变 Riemann 度量恰好对应 \(\mathfrak g/\mathfrak h\) 上的 \(\operatorname{Ad}(H)\) 不变正定内积。
\end{proposition}
\begin{proof}
	记 \(\ell_g(xH)=gxH\)。若 \(A:V\to W\) 是线性同构，张量的搬运在向量因子上用 \(A\)，在对偶因子上用 \((A^{-1})^*\)；记这个同构为 \(A_\#\)。它保持张量积、收缩和复合。对协变二阶张量，公式具体为 \((A_\#b)(v,w)=b(A^{-1}v,A^{-1}w)\)。

	若 \(T\) 是不变张量场，则 \(T_{gH}=(\mathrm{d}\ell_g)_\#T_o\)，所以它由 \(T_o\) 唯一决定。令 \(g=h\in H\)，因 \(hH=o\)，便知 \(T_o\) 被 \(\rho(h)=\mathrm{d}(\ell_h)_o\) 固定。反过来，给定这样的固定张量 \(t\)，规定
	\[
		T_{gH}=(\mathrm{d}(\ell_g)_o)_\#t.
	\]
	若改用代表 \(gh\)，则 \(\ell_{gh}=\ell_g\circ\ell_h\)，所得值为 \((\mathrm{d}\ell_g)_\#\rho(h)_\#t=(\mathrm{d}\ell_g)_\#t\)，因此良定义。

	还需证明它光滑。商投影 \(q:G\to G/H\) 在每点附近有光滑局部截面 \(s:U\to G\)，这是\cref{g1:thm:homogeneous-space}的局部乘积坐标所给出的结论。在 \(U\) 上，刚才的定义写成 \(T_x=(\mathrm{d}(\ell_{s(x)})_o)_\#t\)。作用映射及其微分光滑，逆微分也光滑，所以这些局部分量光滑。群乘法直接给出 \((\mathrm{d}\ell_a)_\#T_x=T_{ax}\)，即不变性。两种构造互逆，且线性。最后，同构搬运保持对称性和正定性，故正定内积恰好给出不变 Riemann 度量。
\end{proof}

\begin{corollary}[紧稳定子的度量]\label{g1:cor:compact-isotropy-metric}
	设 \(G\) 为有限维实 Lie 群，\(H\subset G\) 为紧子群，则 \(G/H\) 存在光滑 \(G\) 不变 Riemann 度量。这里不要求 \(G\) 紧。
\end{corollary}
\begin{proof}
	紧子群在 Hausdorff 的 \(G\) 中闭，故商流形和稳定子表示都有定义。在 \(\mathfrak g/\mathfrak h\) 上任取正定内积 \(b\)，令
	\[
		\overline b(u,v)
		=\int_H b\bigl(\rho(h)u,\rho(h)v\bigr)\,\mathrm{d}\mu_H(h).
	\]
	与\cref{g1:thm:compact-group-average}相同，线性、正性和右不变积分保证它为 \(H\) 不变正定内积。前一命题把它唯一延拓到 \(G/H\)。
\end{proof}

紧性在这里是一个充分条件。比如 \(G=\mathbb R^n\) 平移作用于自身，稳定子只有单位元，任何常系数正定内积都给出不变度量，尽管 \(G\) 非紧。相反，\(\mathbb R_{>0}\) 以伸缩作用于 \(\mathbb R\) 时不可能有不变 Riemann 度量：原点被固定，若其内积值为 \(b>0\)，伸缩 \(x\mapsto2x\) 的不变性便要求 \(4b=b\)。这说明在固定点处，稳定子的线性作用已经可能阻止不变度量的存在。

\subsection{球面和射影空间的齐性表达}
\label{g1:subsec:1703:04}

以下例子把商空间写法变成切空间和度量的具体计算。\cref{g1:prop:matrix-classical-lie-groups}已经构造了正交群、特殊正交群和酉群的 Lie 群结构。它们还是紧群：矩阵方程定义闭集，而每一列长度为一使所有分量有界。稳定子的块对角写法则来自“保持指定向量”或“保持指定子空间”这两个不同条件。

\begin{proposition}[球面模型]\label{g1:prop:sphere-homogeneous-model}
	对 \(n\geq1\)，映射
	\[
		\operatorname{SO}(n+1)/\operatorname{SO}(n)\longrightarrow S^n,
		\qquad AH\longmapsto Ae_0,
	\]
	是与 \(\operatorname{SO}(n+1)\) 作用相容的微分同胚。此处 \(\operatorname{SO}(n)\) 嵌入为 \(\operatorname{diag}(1,B)\)。在基点 \(e_0\)，商切空间通过
	\[
		\begin{pmatrix}0&-v^{\mathsf T}\\v&B\end{pmatrix}
		+\mathfrak{so}(n)\longmapsto v\in\mathbb R^n
	\]
	识别为 \(\mathrm{T}_{e_0}S^n=e_0^\perp\)，稳定子作用为 \(v\mapsto Bv\)。其标准内积延拓得到单位球面的诱导度量。
\end{proposition}
\begin{proof}
	任意单位向量 \(p\) 都能补成正交单位基。因 \(n\geq1\)，必要时反转另一列，可把该基的行列式调为 \(1\) 而保持第一列为 \(p\)，故作用传递。固定 \(e_0\) 的正交矩阵必保持 \(e_0^\perp\)，因此恰为 \(\operatorname{diag}(1,B)\)、\(B\in\operatorname{SO}(n)\)。矩阵作用在球面坐标中光滑，由\cref{g1:prop:transitive-action-model}得到所述微分同胚。

	轨道映射 \(A\mapsto Ae_0\) 的单位元微分是 \(X\mapsto Xe_0\)。反对称矩阵按上述块式写出后，微分取出 \(v\)，核正是右下角的 \(\mathfrak{so}(n)\)。稳定子在 \(e_0^\perp\) 上的作用是其右下角矩阵，因而保持标准内积。所有 \(A\in\operatorname{SO}(n+1)\) 都保持环境 Euclid 内积，所以球面诱导度量正是由基点内积延拓得到的那个不变度量。
\end{proof}

当 \(n=1\) 时，\(\operatorname{SO}(1)=\{1\}\)，上式就是 \(S^1\simeq\operatorname{SO}(2)\)。当 \(n=0\) 时，\(S^0\) 有两个点，\(\operatorname{SO}(1)\) 不能传递作用，不能原样套用该式；改用 \(\operatorname{O}(1)\) 便有 \(S^0\simeq\operatorname{O}(1)/\{1\}\)。

\ProseParagraphBreak

实射影空间的光滑图已在\cref{g1:exm:projective-smooth-atlas}中建立。复射影空间 \(\mathbb{CP}^n\) 同样由 \(\mathbb C^{n+1}\) 中的一维复子空间组成；这里把它视为实流形。给它单位球面按复相位 \(z\sim\lambda z\)、\(|\lambda|=1\) 取商的拓扑，在 \(z_i\ne0\) 处取坐标 \((z_j/z_i)_{j\ne i}\in\mathbb C^n\)。坐标逆把 \(1\) 插入第 \(i\) 位，再归一化到球面并取商，因而连续。换图公式仍是各分量除以一个非零分量，是实光滑映射；\(\mathbb C^n\simeq\mathbb R^{2n}\) 给出实维数 \(2n\)。映射 \([z]\mapsto zz^*/(z^*z)\) 是到 Hermitian 矩阵空间的连续单射，不同投影可用不交开球分离，故商 Hausdorff；有限张坐标图保证第二可数。于是这些图确实给出一个紧光滑流形。

\begin{proposition}[射影空间模型]\label{g1:prop:projective-homogeneous-models}
	对 \(n\geq0\)，有等变微分同胚
	\begin{align*}
		\mathbb{RP}^n&\simeq\operatorname{O}(n+1)/(\operatorname{O}(1)\times\operatorname{O}(n)),\\
		\mathbb{CP}^n&\simeq\operatorname{U}(n+1)/(\operatorname{U}(1)\times\operatorname{U}(n)).
	\end{align*}
	分母按相应的块对角矩阵嵌入，\(\operatorname{O}(0)\)、\(\operatorname{U}(0)\) 均理解为平凡群。统一记 \(\mathbb F=\mathbb R\) 或 \(\mathbb C\)。在一条直线 \(L\) 处，有自然识别
	\[
		\mathrm{T}_L\mathbb{FP}^n\simeq\operatorname{Hom}_{\mathbb F}(L,\mathbb F^{n+1}/L)
		\simeq\operatorname{Hom}_{\mathbb F}(L,L^\perp),
	\]
	第二个识别使用标准内积。对 \(L=\mathbb Fe_0\)，将 \(A\) 识别为 \(v=Ae_0\)，稳定子 \(\operatorname{diag}(\lambda,B)\) 作用为 \(v\mapsto Bv\lambda^{-1}\)。实内积 \(\operatorname{Re}(v^*w)\) 因而给出射影空间上的不变 Riemann 度量。
\end{proposition}
\begin{proof}
	正交或酉变换可以把任意单位向量送到另一单位向量，所以在直线上传递。固定 \(\mathbb Fe_0\) 的变换也保持其正交补，恰为所列块对角矩阵。矩阵作用在射影图中由线性分量相除写出，分母在选定图域非零，故光滑。由\cref{g1:prop:transitive-action-model}，两种集合上的轨道对应都提升为所述微分同胚。

	先在 \(L\) 的某个补空间 \(Q\) 中描述邻近直线。凡投影到 \(L\) 非零的直线都唯一写成 \(\{w+A(w):w\in L\}\)，其中 \(A\in\operatorname{Hom}_{\mathbb F}(L,Q)\)。这正是以 \(L\) 为坐标原点的射影图。因此过 \(L\) 的曲线 \(A(t)\)、\(A(0)=0\) 的切向量为 \(A'(0)\)。把 \(A'(0)w\) 取模 \(L\)，便消除了补空间选择，得到第一个自然识别；取 \(Q=L^\perp\) 得第二个识别。

	稳定子将图中向量 \(w+A(w)\) 送到 \(\lambda w+B A(w)\)，改以 \(w'=\lambda w\) 作参数，新的线性映射为 \(B A\lambda^{-1}\)。于是基点切向量变为 \(Bv\lambda^{-1}\)。因为 \(B\) 正交或酉且 \(|\lambda|=1\)，它保持 \(\operatorname{Re}(v^*w)\)。\cref{g1:prop:invariant-tensor-isotropy}给出不变度量。
\end{proof}

这个度量还能直接算到坐标中。取射影图 \([1:z]\)，\(z\in\mathbb F^n\)，以 \(u,v\in\mathbb F^n\) 表示坐标切向量，令 \(r=1+z^*z\)。曲线 \([1:z+tu]\) 的代表向量为 \(w(t)=(1,z+tu)\)，速度为 \((0,u)\)。先把速度投影到 \(w(0)^\perp\)，再除以 \(|w(0)|\)，便得到定义在单位线向量上的切映射。故
\begin{equation}\label{g1:eq:projective-invariant-metric}
	g_z(u,v)=\operatorname{Re}\frac{r\,u^*v-(u^*z)(z^*v)}{r^2}.
\end{equation}
分母中的两次归一化分别来自正交投影和线向量的长度。由 Cauchy--Schwarz 不等式，分子在 \(u=v\ne0\) 时至少为 \(|u|^2>0\)，所以公式也直接显示正定性。实情形的 \(S^n\to\mathbb{RP}^n\) 是局部等距映射：单位代表 \(p\) 处的速度 \(u\perp p\) 对应 \(A(p)=u\)，内积没有改变。复情形还要去掉相位方向 \(ip\)，因此不能把奇维球面上的所有切向量直接当作射影切向量。

\subsection{Grassmann 流形的切空间与维数}
\label{g1:subsec:1703:05}

\cref{g1:def:grassmannian}已经把 \(\operatorname{Gr}(k,V)\) 定义为 \(V\) 中所有 \(k\) 维线性子空间的集合，\subsecref{g1:subsec:0206:02}还用矩阵的非零子式给出了局部参数。现在要说明这些参数怎样组成光滑流形，以及一个子空间的“无穷小移动”是什么。我们继续允许 \(\mathbb F=\mathbb R\) 或 \(\mathbb C\)，并把 \(k=0\) 的情形补入约定：它只有零子空间这一个点。

\begin{theorem}[Grassmann 流形]\label{g1:thm:grassmann-homogeneous-model}
	设 \(\mathbb F=\mathbb R\) 或 \(\mathbb C\)，\(n\ge0\) 及 \(0\leq k\leq n\) 均为整数，\(V=\mathbb F^n\)，分别赋予标准 Euclid 内积或标准 Hermitian 内积。则 \(\operatorname{Gr}(k,V)\) 有自然的紧光滑流形结构，并有等变微分同胚
	\begin{align*}
		\operatorname{Gr}(k,\mathbb R^n)&\simeq\operatorname{O}(n)/(\operatorname{O}(k)\times\operatorname{O}(n-k)),\\
		\operatorname{Gr}(k,\mathbb C^n)&\simeq\operatorname{U}(n)/(\operatorname{U}(k)\times\operatorname{U}(n-k)).
	\end{align*}
	它们的实维数分别为 \(k(n-k)\)、\(2k(n-k)\)。在任意 \(W\in\operatorname{Gr}(k,V)\) 处，有自然同构
	\begin{equation}\label{g1:eq:grassmann-tangent}
		\mathrm{T}_W\operatorname{Gr}(k,V)
		\simeq\operatorname{Hom}_{\mathbb F}(W,V/W).
	\end{equation}
	复情形右边的复向量空间在此先作为实向量空间理解；其复维数为 \(k(n-k)\)。
\end{theorem}
\begin{proof}
	记 \(E=\mathbb F^k\oplus0\)、\(Q=0\oplus\mathbb F^{n-k}\)。给任意 \(k\) 维子空间选正交单位基并补成 \(V\) 的正交单位基，可见正交群或酉群在所有这样的子空间上传递。保持 \(E\) 的等距线性映射也保持 \(Q=E^\perp\)，稳定子恰好为块对角群 \(H\)。它是闭子群，所以\cref{g1:thm:homogeneous-space}先在 \(G/H\) 上给出光滑结构。集合双射 \(gH\mapsto gE\) 将这个结构赋予 Grassmann 空间；\(G\) 紧，故它也紧。

	为了辨认这套结构，考虑所有投影到 \(E\) 为同构的子空间。它们恰好是线性映射 \(A:E\to Q\) 的图：
	\begin{equation}\label{g1:eq:grassmann-graph-chart}
		W_A=\{x+Ax:x\in E\}.
	\end{equation}
	存在性是先用到 \(E\) 的投影把 \(W\) 参数化，唯一性来自每个 \(x\) 只有一个相应图向量。若 \(g\in G\) 的前 \(k\) 列按 \(E\oplus Q\) 写成 \(\binom{a}{c}\)，则此投影可逆恰好等价于 \(a\) 可逆，此时 \(A=ca^{-1}\)。该条件在 \(G\) 中开，且对右乘 \(H\) 不变；商投影是开映射，所以这些子空间组成 \(G/H\) 的开集。用商的光滑局部截面代入 \(ca^{-1}\)，可知图坐标光滑。

	反向，对块矩阵 \(\left(\begin{smallmatrix}I&0\\A&I\end{smallmatrix}\right)\) 的列依次作正交归一化，得到 \(g(A)\in G\)，其前 \(k\) 列仍张成 \(W_A\)。这个矩阵处处可逆，故每一步减去前面投影后的向量不为零；各步只用光滑加减、内积与正数平方根，因此 \(g(A)\) 光滑依赖 \(A\)。于是 \(A\mapsto g(A)H\) 光滑，确为上述坐标的逆。更换 \(E\oplus Q\) 的坐标后，若同一个向量的分量按固定可逆块矩阵 \(\left(\begin{smallmatrix}a&b\\c&d\end{smallmatrix}\right)\) 变换，则
	\begin{equation}\label{g1:eq:grassmann-chart-transition}
		A'=(c+dA)(a+bA)^{-1}.
	\end{equation}
	交叠条件恰是 \(a+bA\) 可逆，因此转换光滑；复情形还由复有理函数给出。各坐标 \(k\) 平面对应的这类图覆盖全空间，因为任意秩为 \(k\) 的基矩阵总有一个非零 \(k\) 阶子式。这也恢复了第2章的矩阵参数，说明上述结构与那些图坐标相容。图空间 \(\operatorname{Hom}_{\mathbb F}(E,Q)\) 给出所述维数。

	最后在任意 \(W\) 处选补空间并取以 \(W\) 为原点的图坐标。若 \(W(t)\) 是过 \(W\) 的光滑曲线，对 \(w\in W\) 取光滑向量曲线 \(w(t)\in W(t)\)、\(w(0)=w\)，规定其切向量对应的映射为
	\[
		w\longmapsto w'(0)+W\in V/W.
	\]
	这与向量曲线的选择无关：两种选择之差 \(v(t)\in W(t)\) 满足 \(v(0)=0\)，用 \(W(t)\) 的一组光滑基写为 \(\sum_i b_i(t)e_i(t)\)，有 \(b_i(0)=0\)，故 \(v'(0)=\sum_i b_i'(0)e_i(0)\in W\)。若图坐标为 \(A(t)\)、\(A(0)=0\)，可取 \(w(t)=w+A(t)w\)，所得映射就是 \(w\mapsto A'(0)w+W\)。因此它线性、仅依赖切向量，而且给出同构；模 \(W\) 的描述同时表明不依赖补空间。\(k=0\) 或 \(k=n\) 时只有一个子空间，图空间及切空间均为零，结论仍成立。
\end{proof}

\input{Figures/Chapter17-GrassmannGraph}

\cref{g1:fig:grassmann-graph}表示图坐标的几何意义。我们固定横向的 \(E\)，让子空间 \(W_A\) 倾斜；坐标 \(A\) 记录每个 \(x\in E\) 需要补上的竖向分量。图中只有一个横向和一个竖向维度，但同一构造适用于任意 \(k,n\)。切向量记录的是斜率映射的变化率，而非对子空间内的每个向量任意指定一个独立移动。

\begin{proposition}[Grassmann 流形的度量]\label{g1:prop:grassmann-invariant-metric}
	设 \(\mathbb F=\mathbb R\) 或 \(\mathbb C\)，\(n\ge0\) 及 \(0\leq k\leq n\) 均为整数，\(V=\mathbb F^n\) 分别带标准 Euclid 内积或标准 Hermitian 内积。在 \(\operatorname{Gr}(k,V)\) 上取\cref{g1:thm:grassmann-homogeneous-model}的自然光滑结构，实、复情形分别考虑 \(\operatorname{O}(n)\)、\(\operatorname{U}(n)\) 的标准作用，并用正交投影将 \(\mathrm{T}_W\operatorname{Gr}(k,V)\simeq\operatorname{Hom}_{\mathbb F}(W,V/W)\) 识别为 \(\operatorname{Hom}_{\mathbb F}(W,W^\perp)\)。则
	\begin{equation}\label{g1:eq:grassmann-hilbert-schmidt}
		g_W(S,T)=\operatorname{Re}\operatorname{tr}(S^*T),
		\qquad S,T\in\operatorname{Hom}_{\mathbb F}(W,W^\perp),
	\end{equation}
	定义光滑不变 Riemann 度量。稳定子 \(\operatorname{diag}(C,D)\) 的切表示为 \(T\mapsto DTC^{-1}\)。在 \(\dim_{\mathbb F}E=k\) 的正交分解 \(V=E\oplus Q\) 所给的图坐标 \(A:E\to Q\) 中，对坐标切向量 \(U_1,U_2:E\to Q\)，该度量为
	\begin{equation}\label{g1:eq:grassmann-chart-metric}
		g_A(U_1,U_2)=\operatorname{Re}\operatorname{tr}
		\bigl((I+A^*A)^{-1}U_1^*(I+AA^*)^{-1}U_2\bigr).
	\end{equation}
\end{proposition}
\begin{proof}
	稳定子将 \(x+Tx\) 送到 \(Cx+DTx\)，用 \(x'=Cx\) 改写，图映射变为 \(DTC^{-1}\)，在 \(T=0\) 处也就是所述切表示。由于 \(C,D\) 正交或酉，
	\[
		\operatorname{tr}\bigl((DSC^{-1})^*(DTC^{-1})\bigr)
		=\operatorname{tr}(CS^*TC^{-1})=\operatorname{tr}(S^*T).
	\]
	在正交单位基中，\(\operatorname{tr}(T^*T)\) 是 \(T\) 的所有矩阵分量的模平方之和，故为正定实内积。这是线性映射空间上的平方和内积。\cref{g1:prop:invariant-tensor-isotropy}把它延拓为光滑不变度量；等距变换把正交补送到正交补，所以延拓后的每点公式仍为\eqref{g1:eq:grassmann-hilbert-schmidt}。

	为计算图坐标表达，设 \(S_0=I+A^*A\)、\(T_0=I+AA^*\)。图映射 \(x\mapsto(x,Ax)\) 将 \(W_A\) 的内积拉回为 \(x^*S_0y\)。坐标曲线 \(A+tU\) 中，图向量 \((x,(A+tU)x)\) 的速度为 \((0,Ux)\)。而
	\[
		W_A^\perp=\{(-A^*b,b):b\in Q\}.
	\]
	若 \((0,Ux)\) 的正交投影为 \((-A^*b,b)\)，与任意 \((-A^*c,c)\) 作内积便得 \(c^*T_0b=c^*Ux\)，故 \(b=T_0^{-1}Ux\)。因此两个坐标速度 \(U_1,U_2\) 在输出空间的内积矩阵为 \(U_1^*T_0^{-1}U_2\)。输入空间的度量矩阵是 \(S_0\)，要取正交单位输入基的平方和，就需取 \(\operatorname{tr}(S_0^{-1}U_1^*T_0^{-1}U_2)\)；这可以把一组满足 \(B^*S_0B=I\) 的输入基代入后，用 \(BB^*=S_0^{-1}\) 直接核对。取实部即得\eqref{g1:eq:grassmann-chart-metric}。两个正定矩阵始终可逆，公式中的各项光滑，亦可直接看出整个图上的光滑性。
\end{proof}

当 \(k=1\) 时，Grassmann 流形就是射影空间，\eqref{g1:eq:grassmann-chart-metric}化为\eqref{g1:eq:projective-invariant-metric}：恒等式 \((I+zz^*)^{-1}=I-zz^*/(1+z^*z)\) 可直接相乘验证。更一般地，\(\operatorname{Gr}(2,\mathbb R^4)\) 的一张图由一个 \(2\times2\) 实矩阵的四个分量参数化，实维数为四；\(\operatorname{Gr}(2,\mathbb C^4)\) 使用四个复参数，实维数为八。这里的四维或八维指的是“子空间组成的参数空间”的维数，不能与每个被参数化子空间本身的维数二混为一谈。

% 纲要条目排为编号小节标题，并进入目录与书签；说明段仍保留于纲要注释中。

% !TeX root = ../../Book1.tex
% 纲要标题：### 17.4 非自由作用的局部模型
% 纲要位置：几何学计划纲要.md，第 798 行。
% 正文状态：初稿；尚未正式审核。

% 纲要内容开始
% 
% * 紧 Lie 群固定点附近的线性化
% * 有限稳定子与局部有限群商图
% * 底层商空间和稳定子数据
% * 固有作用的切片定理
% 
% ---
% 
% 纲要内容结束

\section{非自由作用的局部模型}
\label{g1:sec:1704}

自由性使轨道方向能被完全消去，留下普通流形坐标。若一个点仍被某些群元素固定，横截轨道以后，这些元素还会作用在横截面上；商中保留下来的局部识别正由它们产生。本节先把固定点附近的作用写成线性作用，再证明固有性允许将这种局部描述推广到整条轨道的邻域。有限稳定子给出有限群商图，而正维稳定子也能用同一切片构造处理。

\subsection{紧 Lie 群固定点附近的线性化}
\label{g1:subsec:1704:01}

设有限群 \(H\) 光滑作用于 \(M\)，并固定 \(p\)。每个 \(h\in H\) 在 \(p\) 的微分组成线性表示 \(\rho(h)=\mathrm{d}(a_h)_p\)，其中 \(a_h(x)=hx\)。任意坐标 \(c\) 通常不满足 \(c(hx)=\rho(h)c(x)\)。不过，\(\rho(h)^{-1}c(hx)\) 是一份把坐标先移过去、再按线性作用移回来的修正；对所有 \(h\) 平均，就能让每个群元素享有同样地位。

\begin{proposition}[固定点线性化]\label{g1:prop:fixed-point-linearization}
	设有限维实紧 Lie 群 \(H\) 光滑作用于有限维、无边界、Hausdorff、第二可数实光滑流形 \(M\)，且 \(p\) 被 \(H\) 的每个元素固定。令 \(\rho(h)=\mathrm{d}(a_h)_p\)。则存在 \(H\) 不变的开邻域 \(U\ni p\)、\(\mathrm{T}_pM\) 上的 \(H\) 不变内积及关于此内积的零点开球 \(B\)，使某个微分同胚
	\[
		F:U\longrightarrow B\subset \mathrm{T}_pM
	\]
	满足 \(F(p)=0\)、\(\mathrm{d}F_p=\operatorname{id}_{\mathrm{T}_pM}\) 及 \(F(hx)=\rho(h)F(x)\)。有限群作为零维紧 Lie 群包含在此结论中。
\end{proposition}
\begin{proof}
	先取一张以 \(\mathrm{T}_pM\) 为坐标值空间的坐标图 \(c:U_0\to \mathrm{T}_pM\)，使 \(c(p)=0\)、\(\mathrm{d}c_p=I\)；任意坐标图经一个可逆线性变换即可满足后一个条件。我们希望对同一个 \(x\) 的全部 \(hx\) 求坐标，故将定义域缩到
	\[
		U_1=\{x\in M:Hx\subset U_0\}.
	\]
	它含 \(p\) 且 \(H\) 不变。它也是开集：若 \(Hx\subset U_0\)，作用的连续性给每个 \(h\in H\) 一对邻域 \(V_h\ni h\)、\(W_h\ni x\)，使 \(V_hW_h\subset U_0\)。从 \(\{V_h\}\) 取有限子覆盖，再交相应的 \(W_h\)，所得开邻域 \(W\ni x\) 满足 \(HW\subset U_0\)。有限群时，\(U_1\) 也就是有限个开集 \(h^{-1}U_0\) 的交。

	有限群的平均坐标为
	\[
		F(x)=\frac1{|H|}\sum_{h\in H}\rho(h)^{-1}c(hx).
	\]
	一般紧 Lie 群则使用\cref{g1:thm:compact-group-integration}的归一化左右不变积分，定义
	\[
		F(x)=\int_H\rho(h)^{-1}c(hx)\,\mathrm{d}\mu(h).
	\]
	这里逐坐标积分。被积函数光滑，\cref{g1:lem:compact-parameter-integral}允许逐阶在积分下求导，故 \(F\) 光滑：所需的一致控制来自紧参数集 \(H\) 及充分小坐标邻域的紧闭包。有限群时，这个积分直接化为有限个光滑映射之和。对固定 \(k\in H\)，作右平移换元 \(h'=hk\)，得到
	\[
		F(kx)=\int_H\rho(k)\rho(h')^{-1}c(h'x)\,\mathrm{d}\mu(h')
		=\rho(k)F(x).
	\]
	而 \(F(p)=0\)，每一项在 \(p\) 的微分均为 \(\rho(h)^{-1}I\rho(h)=I\)，故 \(\mathrm{d}F_p=I\)。

	由\cref{g1:thm:inverse-function}，\(F\) 在 \(p\) 的一个开邻域 \(V\subset U_1\) 上为到开集的微分同胚。不能直接假定 \(V\) 不变，但可再次使用上面的紧性论证，取开不变邻域 \(V_1=\{x:Hx\subset V\}\)。因为 \(V_1\subset V\)，限制 \(F|_{V_1}\) 仍为微分同胚，且其像开、含零并在 \(\rho(H)\) 下不变。

	最后，将任意正定内积按 \(H\) 平均，得到
	\[
		\langle u,v\rangle_H=\int_H\langle\rho(h)u,\rho(h)v\rangle_0\,\mathrm{d}\mu(h).
	\]
	其正定性与不变性和\cref{g1:thm:compact-group-average}中的论证相同。在开集 \(F(V_1)\) 中取关于此内积的充分小零点球 \(B\)，令 \(U=F^{-1}(B)\)。\(B\) 与 \(U\) 都不变，所需性质随限制保留。
\end{proof}

这说明在固定点附近，我们不仅能得到作用的一阶线性近似，还能选好坐标，使作用本身就等于这个线性表示。若某个 \(h\) 的微分 \(\rho(h)\) 为恒等，它便在这个共同邻域中逐点固定所有点。反过来，\(\rho(h)\ne I\) 时，总有任意靠近 \(p\) 的点被它移动。后面保留稳定子数据时，这一区别会很重要。

\subsection{有限稳定子与局部有限群商图}
\label{g1:subsec:1704:02}

现在 \(p\) 不必被整个 \(G\) 固定，只被稳定子 \(H=G_p\) 固定。\(H\) 在 \(\mathrm{T}_pM\) 上的微分表示保持轨道切空间 \(\mathrm{T}_p(Gp)\)：群元素把轨道送回自身。若作用固有，\cref{g1:prop:proper-action-orbits}保证 \(H\) 紧，因而平均内积给出一个 \(H\) 不变的正交补 \(N\)，满足
\[
	\mathrm{T}_pM=\mathrm{T}_p(Gp)\oplus N.
\]
\(N\) 记录离开轨道的方向。我们需要把它变成真正的横截面，并排除很远的群元素将这张小横截面送回自身的情形。

\begin{lemma}[横截面与返回]\label{g1:lem:proper-slice-construction}
	设有限维实 Lie 群 \(G\) 光滑且固有地作用于有限维实光滑流形 \(M\)，二者无边界、Hausdorff、第二可数。固定 \(p\in M\)，令 \(H=G_p\)、\(\mathfrak g=\mathrm{T}_eG\)、\(\mathfrak h=\mathrm{T}_eH\)。取 \(H\) 不变补空间 \(\mathrm{T}_pM=\mathrm{T}_p(Gp)\oplus N\)，并取线性补空间 \(\mathfrak g=\mathfrak h\oplus\mathfrak m\)。存在 \(H\) 不变嵌入子流形 \(S\ni p\)，使 \(S\) 与 \(N\) 中的一个不变零点开球 \(B\) 等变微分同胚，且具有以下性质：
	\begin{enumerate}
		\item \(\mathrm{T}_pS=N\)，并存在 \(\mathfrak m\) 中的零点邻域 \(V\)，使 \((X,s)\mapsto\exp_G(X)s\) 是 \(V\times S\) 到 \(M\) 中某个开集的微分同胚。
		\item 若 \(s,s'\in S\) 且 \(gs=s'\)，则 \(g\in H\)。因此 \(G_s=H_s\) 对所有 \(s\in S\) 成立。
		\item \(GS\) 在 \(M\) 中开，包含整条轨道 \(Gp\)。包含 \(S\hookrightarrow GS\) 诱导同胚 \(S/H\simeq GS/G\)。
	\end{enumerate}
\end{lemma}
\begin{proof}
	由\cref{g1:thm:closed-subgroup}及\cref{g1:prop:proper-action-orbits}，\(H\) 是紧嵌入 Lie 子群。将\cref{g1:prop:fixed-point-linearization}应用于 \(H\) 在 \(p\) 附近的作用。在得到的等变坐标中取 \(N\) 的小球，再取其逆像 \(S_0\)。这是一张 \(H\) 不变嵌入横截面，\(\mathrm{T}_pS_0=N\)。

	考虑
	\[
		\Psi:\mathfrak m\times S_0\longrightarrow M,\qquad
		\Psi(X,s)=\exp_G(X)s.
	\]
	它在 \((0,p)\) 的微分为 \((X,v)\mapsto X_M(p)+v\)。由\cref{g1:prop:orbit-map-rank}，\(X\mapsto X_M(p)\) 的核是 \(\mathfrak h\)，像是 \(\mathrm{T}_p(Gp)\)，所以限制到 \(\mathfrak m\) 后是到轨道切空间的线性同构。再用直和分解，\(\mathrm{d}\Psi_{(0,p)}\) 可逆。逆函数定理给出局部微分同胚。将其定义域缩成乘积 \(V\times S_1\)，其中 \(V\subset\mathfrak m\) 开且含零，\(S_1\subset S_0\) 仍由一个不变零点小球给出，便有 \(\Psi|_{V\times S_1}\) 单射且微分处处可逆。

	先控制靠近 \(H\) 的群元素。映射
	\[
		H\times\mathfrak m\longrightarrow G,\qquad (h,X)\longmapsto h\exp_G X
	\]
	在每个 \((h,0)\) 处微分可逆，因为左平移后就是 \(\mathfrak h\oplus\mathfrak m\to\mathfrak g\) 的相加映射。因此 \(H\exp_G(V)\) 包含 \(G\) 中一个开邻域 \(W\supset H\)。若 \(g\in W\)、\(s,s'\in S_1\) 且 \(gs=s'\)，可写 \(g=h\exp_G X\)，其中 \(h\in H\)、\(X\in V\)。由 \(S_1\) 不变，\(h^{-1}s'\in S_1\)，于是
	\[
		\Psi(X,s)=h^{-1}s'=\Psi(0,h^{-1}s').
	\]
	\(\Psi\) 的单射性迫使 \(X=0\)，因而 \(g=h\in H\)。这里 \(S_1\) 的不变性让同一个局部逆同时适用于所有 \(h\in H\)，无需逐个群元素选择横截面。

	接着控制不属于 \(W\) 的群元素。若任意小的 \(H\) 不变球形横截面仍有这样的返回，就能取 \(s_j,s'_j\in S_1\)，使 \(s_j,s'_j\to p\)，并取 \(g_j\notin W\) 满足 \(g_js_j=s'_j\)。由\cref{g1:prop:proper-action-criteria}的序列判据，\(g_j\) 有收敛子列，设其极限为 \(g\)。作用连续性给出 \(gp=p\)，即 \(g\in H\subset W\)。因为 \(W\) 开，该子列最终位于 \(W\)，矛盾。因此可将 \(S_1\) 缩成同类不变横截面 \(S\)，使一切满足 \(gs=s'\)、\(s,s'\in S\) 的 \(g\) 都属于 \(W\)，再由上一段得到 \(g\in H\)。这就是所需的返回结论，取 \(s'=s\) 便得 \(G_s=H_s\)。

	\(\Psi(V\times S)\) 开且包含 \(S\)，又包含于 \(GS\)，故其所有 \(G\) 平移的并恰为 \(GS\)，从而 \(GS\) 开。每条其中的轨道都与 \(S\) 相交，返回结论又说明两点 \(s,s'\in S\) 在同一 \(G\) 轨道上，当且仅当它们在同一 \(H\) 轨道上。因此得到连续双射 \(S/H\to GS/G\)。为证明逆连续，任取 \(S/H\) 中开集，设其在 \(S\) 中的原像为开集 \(A\)。\(\Psi(V\times A)\) 在 \(M\) 中开，故其所有 \(G\) 平移的并 \(GA\) 开。\(GA\) 恰是对应商集在 \(GS\) 中的原像，因而该商集在 \(GS/G\) 中开。双射于是为同胚。
\end{proof}

固有性在证明中承担了一个具体任务：如果横截面上的点与其像都逼近 \(p\)，相关群元素不能逃离每个紧集。没有这一步，逆函数定理只能阻止靠近单位元的返回，不能阻止来自群的远处的返回。

\begin{definition}[局部有限群商图]\label{g1:def:orbifold-chart-local}
	设 \(X\) 为拓扑空间，\(U\subset X\) 开。一个\term[orbifoldtu]{局部有限群商图}[orbifold chart]是三元组 \((B,H,\chi)\)：\(B\) 为有限维实内积空间 \(N\) 中的零点开球，有限群 \(H\) 线性地保持 \(B\)，而 \(\chi:B/H\to U\) 是同胚。本节允许表示 \(H\to\operatorname{GL}(N)\) 有非平凡核，也允许其中有反射；表示核非平凡的图称为非有效图，表示忠实的图称为有效图。
\end{definition}

图中的点 \([v]\) 还带有稳定子 \(H_v=\{h\in H:hv=v\}\)；更换同一轨道的代表只把这个子群变为共轭子群，所以其同构类型不变。这正是普通商坐标之外需要保留的数据。一般 局部有限群商图册还要规定不同图之间的相容映射；本章只建立来自给定群作用的局部模型，不展开抽象图册的全部理论。

\begin{corollary}[有限稳定子商的局部模型]\label{g1:cor:finite-stabilizer-local-quotient}
	设有限维实 Lie 群 \(G\) 光滑且固有地作用于有限维无边界实光滑流形 \(M\)，二者 Hausdorff、第二可数，且所有稳定子有限。则每个商点 \([p]\in M/G\) 都有上述 局部有限群商图，其群为 \(G_p\)，其线性作用为 \(G_p\) 在横向空间 \(N\simeq \mathrm{T}_pM/\mathrm{T}_p(Gp)\) 上的表示。局部图的维数为 \(\dim M-\dim G\)，并且图中每点的稳定子与原作用在相应点的稳定子同构。
\end{corollary}
\begin{proof}
	对 \(p\) 应用\cref{g1:lem:proper-slice-construction}。\(S\) 等变微分同胚于 \(N\) 中的球 \(B\)，而 \(S/H\simeq GS/G\)，后者在 \(M/G\) 中开，因此给出所需图。有限群 \(H\) 的 Lie 代数为零，故轨道维数为 \(\dim G\)，横向维数为 \(\dim M-\dim G\)。最后，切片内 \(G_s=H_s\)；在 \(gs\) 处稳定子为 \(gH_sg^{-1}\)，所以图保存了原稳定子的同构类型。
\end{proof}

有限稳定子本身不够。例如 \(\mathbb R\) 沿无理斜率作用于 \(T^2\) 时稳定子平凡，轨道却稠密；商甚至不是 Hausdorff 空间。每个有限群的 Euclid 球商都是 Hausdorff 的，故这样的稠密轨道商不能由上述图描述。固有性与稳定子的有限性分别控制返回现象及残余群的大小，二者不能互相替代。

\subsection{底层商空间和稳定子数据}
\label{g1:subsec:1704:03}

同一个拓扑空间可以有不同的局部群数据。最简单的例子已经出现在二维：一个锥点在底层拓扑上未必是特殊点，但围绕它的有限旋转仍是商结构的一部分。

\begin{example}[有限旋转商]\label{g1:exm:rotation-orbifold-chart}
	令 \(m\geq2\)，\(\mathbb Z/m\mathbb Z\) 在 \(\mathbb C\simeq\mathbb R^2\) 上按 \(z\mapsto\exp(2\pi\mathrm{i}/m)z\) 作用。映射 \(z\mapsto z^m\) 在每条轨道上常值，而且两个点具有相同的 \(m\) 次幂，当且仅当它们在同一轨道上。因此它诱导连续双射
	\[
		\mathbb C/(\mathbb Z/m\mathbb Z)\longrightarrow\mathbb C,\qquad [z]\longmapsto z^m.
	\]
	这是同胚。具体地，若闭集 \(A\subset\mathbb C\) 的像中有收敛序列 \(z_j^m\to w\)，则 \(|z_j|=|z_j^m|^{1/m}\) 有界，存在收敛子列，其极限 \(z\in A\) 满足 \(z^m=w\)。所以幂映射是闭映射，由商拓扑得到上述双射也闭，因而逆连续。

	底空间虽然是一张平面，原点稳定子却是整个 \(\mathbb Z/m\mathbb Z\)，其余点的稳定子平凡。不同的 \(m\)，以及没有群作用的普通平面，具有相同底拓扑而不同稳定子。甚至在坐标 \(w=z^m\) 中，光滑不变函数 \(|z|^2\) 对应 \(|w|^{2/m}\)，后者在原点不光滑。这说明这张拓扑坐标不能直接代替有限群商的光滑数据。
\end{example}

\begin{example}[反射商]\label{g1:exm:reflection-orbifold-chart}
	令二阶群按 \((x,y)\mapsto(x,-y)\) 作用于 \(\mathbb R^2\)。映射 \((x,y)\mapsto(x,|y|)\) 将商同胚地识别为上半平面：每条轨道有唯一的上半平面代表，取该代表的映射连续；商映射限制在上半平面上给出逆映射。边界 \(y=0\) 上的稳定子为整个二阶群，内部点的稳定子平凡。因此，即使被作用的流形无边界，允许反射的有限群商也可能有底拓扑的边界。

	本节的 局部有限群商图允许这种反射。若采用排除反射的文献约定，就必须另加条件，不能将两种定义下的“无边界”直接等同。\cref{g1:fig:finite-quotient-data}将有限旋转商与反射商并列，图中标注的是稳定子，而非额外删去的点。
\end{example}

\begin{figure}[htbp]
	\centering
	\input{Figures/Chapter17-FiniteQuotients}
	\caption[有限群商中的稳定子]{有限旋转将若干点识别成一个点，商仍同胚于圆盘，但中心保留 \(\mathbb Z/m\mathbb Z\)；反射商是半圆盘，直径上的点保留二阶稳定子。左图以 \(m=3\) 示意。}
	\label{g1:fig:finite-quotient-data}
\end{figure}

还有一种不同的信息损失：某些群元素在整张图中什么也不移动。若 \(K=\ker(H\to\operatorname{GL}(N))\)，则 \(K\) 正规，且 \(H/K\) 在 \(B\) 上有效作用。两者轨道完全相同，因而 \(B/H=B/(H/K)\) 作为拓扑空间自然一致；但是
\[
	(H/K)_v=H_v/K.
\]
将 \(H\) 换为 \(H/K\) 称为这张图的有效化，它同时除掉了每一点稳定子中的 \(K\)，并不保存原来的稳定子数据。

\begin{example}[相同轨道与不同稳定子]\label{g1:exm:ineffective-local-quotient}
	在前面的旋转例子中改用群 \(H=(\mathbb Z/m\mathbb Z)\times(\mathbb Z/k\mathbb Z)\)，其中 \(k\geq2\)，让第二因子恒等作用。商空间及每条轨道均未改变，但原点稳定子变为 \(H\)，非零点稳定子变为 \(\{0\}\times\mathbb Z/k\mathbb Z\)。有效化回到原来的 \(m\) 阶旋转图，恰好删去这些共同的稳定子。这也是本节不自动要求局部群表示忠实的原因。
\end{example}

\subsection{固有作用的切片定理}
\label{g1:subsec:1704:04}

前面的横截面构造并未用到 \(H\) 有限，只用到稳定子紧及作用固有。有限性只是把 \(S/H\) 限制为有限群商。对一般紧稳定子，我们还可以说明整个轨道邻域怎样由一张横截面组装起来：每个 \(g\) 带来一份 \(gS\)，而 \((g,s)\) 与 \((gh^{-1},hs)\) 表示同一个点。这一识别由一个自由作用产生，因此可以使用已经证明的商流形定理。

\begin{theorem}[固有作用的切片定理]\label{g1:thm:proper-action-slice}
	设有限维实 Lie 群 \(G\) 光滑且固有地作用于有限维无边界实光滑流形 \(M\)，二者 Hausdorff、第二可数。对任意 \(p\in M\)，令 \(H=G_p\)，令 \(N\subset \mathrm{T}_pM\) 为某个 \(H\) 不变的轨道切空间补空间，并在 \(N\) 上取 \(H\) 不变内积。存在零点开球 \(B\subset N\)、\(H\) 等变嵌入 \(j:B\to M\) 以及 \(G\) 不变开邻域 \(U\supset Gp\)，满足 \(j(0)=p\)、\(\mathrm{d}j_0\) 为 \(N\hookrightarrow \mathrm{T}_pM\) 的包含，且
	\begin{equation}\label{g1:eq:proper-slice-diffeomorphism}
		G\times_H B\longrightarrow U,\qquad [g,v]\longmapsto g j(v)
	\end{equation}
	为 \(G\) 等变微分同胚。这里
	\[
		G\times_H B=(G\times B)/H,\qquad
		h\cdot(g,v)=(gh^{-1},hv),
	\]
	\(G\) 在其上按 \(a[g,v]=[ag,v]\) 作用。特别地，\(U/G\simeq B/H\)，并且
	\[
		G_{g j(v)}=gH_vg^{-1}.
	\]
\end{theorem}
\begin{proof}
	用\cref{g1:lem:proper-slice-construction}选取 \(S=j(B)\)，令 \(U=GS\)。若事先在 \(N\) 上选定了另一个不变内积，只需在已有线性化坐标域中取关于该内积的更小零点球；引理的局部逆及返回性质在缩小后仍成立。

	先核实 \(G\times_H B\) 是已建立的商流形。上式确实定义左 \(H\) 作用，因为连续作用两次给出
	\[
		h_1\cdot\bigl(h_2\cdot(g,v)\bigr)
		=(gh_2^{-1}h_1^{-1},h_1h_2v)
		=(h_1h_2)\cdot(g,v).
	\]
	它自由，因为 \(gh^{-1}=g\) 强制 \(h=e\)。它固有，因为 \(H\) 紧：对紧集 \(K\subset(G\times B)^2\)，作用映射的原像是 \(H\times\operatorname{pr}_2K\) 中的闭集，故紧。\cref{g1:thm:free-proper-quotient}于是赋予商一个无边界、Hausdorff、第二可数的光滑流形结构，商映射为淹没。\(G\) 的左乘作用与 \(H\) 作用交换，故下降为商上的光滑作用；其光滑性可用商映射的局部截面，在 \(G\times(G\times_HB)\) 的局部坐标中检验。

	光滑映射 \((g,v)\mapsto gj(v)\) 在 \(H\) 轨道上常值，由\cref{g1:prop:quotient-smooth-universal-property}得到光滑映射 \(\overline\Phi:G\times_HB\to U\)。它满射。若 \(g j(v)=g'j(v')\)，则 \(g'^{-1}g\) 将 \(S\) 中一点送到另一点，引理的返回性质给出 \(h=g'^{-1}g\in H\)，且 \(v'=hv\)。于是 \((g',v')=(gh^{-1},hv)\)，两者属于同一个 \(H\) 轨道，证明单射。

	还须检验逆光滑。在 \((e,v)\) 处，原映射沿 \(\mathfrak m\) 方向和 \(\mathrm{T}_vB\) 方向的微分，正是引理中 \((X,s)\mapsto\exp_G(X)s\) 的可逆微分与 \(\mathrm{d}j_v\) 的组合，因而满射。左平移将这点传到每个 \((g,v)\)。原映射通过商映射分解，故 \(\mathrm{d}\overline\Phi\) 处处满射；而
	\[
		\dim(G\times_HB)=\dim G+\dim N-\dim H=\dim M,
	\]
	所以它处处可逆。逆函数定理使 \(\overline\Phi\) 为局部微分同胚，双射性再将各处的局部光滑逆拼成全局光滑逆。这就证明\eqref{g1:eq:proper-slice-diffeomorphism}，等变性由公式直接得到。最后，商空间同胚已由引理证明；稳定子公式由 \(G_{j(v)}=H_{j(v)}=H_v\) 及共轭变换得到。
\end{proof}

这样的 \(S=j(B)\) 称为过 \(p\) 的\term[qiepian]{切片}[slice]。其中 \(N\) 可同构地换成商空间 \(\mathrm{T}_pM/\mathrm{T}_p(Gp)\)，后者的 \(H\) 表示不依赖补空间的选择。不变内积只帮助选择具体代表和小球，并未给 \(M\) 预设一个全局 \(G\) 不变度量。定理也不要求 \(G\) 紧；非紧群的固有性已经足以提供紧稳定子并排除远处返回。

\ProseParagraphBreak

当 \(H=\{e\}\) 时，模型退化为 \(G\times B\)，恢复自由固有作用的局部乘积；当 \(H\) 有限时，商局部为上述有限群商图；当 \(H\) 正维时，\(B/H\) 可以呈现不同维数轨道相邻的现象。例如 \(\operatorname{SO}(2)\) 旋转作用于 \(\mathbb R^2\)，原点的切片就是一个圆盘，商为半开半径区间 \([0,\varepsilon)\)，非零点的轨道是一维圆，原点轨道却是零维。它的底空间恰巧也像反射商的一维截面，但原点稳定子是整条圆周群，不能用有限稳定子图保留同样的数据。

\ProseParagraphBreak

本节已证明的是光滑固有作用在一条轨道附近的切片模型。关于不同稳定子类型如何组成分层、不同切片如何组织为完整 局部有限群商图册，以及带边界或其他范畴中的切片版本，还需要相应的定义与证明。后续使用本节结果时，可以直接调用\eqref{g1:eq:proper-slice-diffeomorphism}及其商 \(B/H\)，而不能由此把所有非自由商都当作普通流形。


\begin{chaptersummary}
	沿 \(\exp(tX)\cdot p\) 取初速度得到完备基本场 \(X_M\)，左作用和正号约定使 \([X_M,Y_M]=-[X,Y]_M\)。稳定子 \(G_p\) 闭，其 Lie 代数是轨道微分的核；轨道自然具有 \(G/G_p\) 的浸入流形结构，切空间为 \(\mathfrak g/\mathfrak g_p\)。传递作用使这一模型成为整个流形的微分同胚描述。

	\ProseParagraphBreak

	固有性控制两端留在紧区域内的全部群参数，使稳定子紧、轨道闭嵌入，并排除横截面的远处返回。与自由性合用，得到 \(G\times S\cong GS\) 的局部乘积，因而 \(M/G\) 是唯一使投影成为淹没的光滑商。轨道不变光滑映射可唯一光滑下降；只要求投影光滑不足以保证这一性质。

	\ProseParagraphBreak

	左平移延拓单位元处的正密度，紧性给出有限总量；模函数的紧像只能平凡，于是归一化积分同时左右不变。平均度量保留正定性，紧参数积分保证光滑性。齐性空间上的不变张量则由基点处的稳定子不变张量唯一决定，不必对整个群都作积分。

	\ProseParagraphBreak

	球面和射影空间由“固定向量”与“固定直线”的不同稳定子产生。Grassmann 流形以线性映射的图作局部坐标，\(\mathrm{T}_W\operatorname{Gr}(k,V)\cong\operatorname{Hom}_{\mathbb F}(W,V/W)\)。选定内积后可用 \(W^\perp\) 代表商空间，从而直接计算稳定子作用及平方和度量；复情形的实维数须乘二。

	\ProseParagraphBreak

	紧群固定点附近可通过平均坐标把作用完全线性化。对固有作用，这给出轨道邻域 \(G\times_{G_p}B\)，商邻域为 \(B/G_p\)。有限稳定子产生局部有限群商图，其中反射可能带来底拓扑边界，非有效核则保留额外稳定子。底空间的同胚类型不能代替这些群数据；一般正维稳定子的商也不自动是普通流形或有限群商。
\end{chaptersummary}

\BookExercises


\begin{exercise}\label{g1:ex:conjugation-orbit-dimension}
	令 \(G=\operatorname{GL}(n,\mathbb R)\) 在 \(M_n(\mathbb R)\) 上按共轭作用。求 \(X\in\mathfrak{gl}(n,\mathbb R)\) 对应的基本场、矩阵 \(A\) 的稳定子 Lie 代数及共轭轨道切空间。若 \(A\) 可实对角化，互异特征值的重数为 \(m_1,\ldots,m_s\)，求轨道维数。
\end{exercise}
\begin{solution}
	对 \(\exp(tX)A\exp(-tX)\) 求导，得到
	\[
		X_M(A)=XA-AX.
	\]
	稳定子是可逆中心化子 \(G_A=\{g:gA=Ag\}\)，由\cref{g1:prop:orbit-map-rank}，其 Lie 代数为 \(\mathfrak g_A=\{X:XA=AX\}\)。\cref{g1:prop:orbit-immersion}进一步给出
	\[
		\mathrm{T}_A(G\cdot A)=\{XA-AX:X\in M_n(\mathbb R)\}.
	\]
	这里切空间对应轨道的自然浸入结构，不需要先假定轨道闭。

	取特征空间直和 \(\mathbb R^n=E_1\oplus\cdots\oplus E_s\)，\(A|_{E_i}=\lambda_iI\)。将 \(X\) 写为分块矩阵 \((X_{ij})\)，交换关系在第 \((i,j)\) 块成为 \((\lambda_j-\lambda_i)X_{ij}=0\)。不同特征值使非对角块都为零；对角块 \(X_{ii}\in\operatorname{End}(E_i)\) 可以任取。因此
	\[
		\dim\mathfrak g_A=\sum_{i=1}^s m_i^2,
		\qquad \dim(G\cdot A)=n^2-\sum_{i=1}^s m_i^2.
	\]
	特别地，互异实特征值给出维数 \(n^2-n\)，标量矩阵给出零维单点轨道。后一个例子也说明群维数大，并不意味着它在每个对象上都产生很多运动方向。
\end{solution}

\begin{exercise}\label{g1:ex:affine-action-fundamental-fields}
	令 \(G=(0,\infty)\times\mathbb R\)，乘法为 \((a,b)(c,d)=(ac,b+ad)\)，作用于 \(\mathbb R\) 的公式为 \((a,b)\cdot x=ax+b\)。取单位元处的基 \(H=\partial_a\)、\(E=\partial_b\)。求基本场、各点的稳定子及稳定子 Lie 代数，并把 \(\mathbb R\) 写成 \(G\) 的齐性空间。说明该作用是否自由、是否有效。
\end{exercise}
\begin{solution}
	单参数子群分别为 \(\exp(tH)=(e^t,0)\)、\(\exp(tE)=(1,t)\)，所以
	\[
		H_{\mathbb R}=x\partial_x,\qquad E_{\mathbb R}=\partial_x,
		\qquad [H_{\mathbb R},E_{\mathbb R}]=-\partial_x.
	\]
	前章算得 \([H,E]=E\)，因此这个例子核对了正号基本场约定下的反同态公式。固定 \(x\) 的条件是 \(ax+b=x\)，故
	\[
		G_x=\{(a,(1-a)x):a>0\},\qquad
		\mathfrak g_x=\mathbb R(H-xE).
	\]
	后一个等式也可以从 \((uH+vE)_{\mathbb R}(x)=(ux+v)\partial_x\) 的核直接得到。

	平移可把 \(0\) 送到任意 \(x\)，所以作用传递。令 \(H_0=G_0=\{(a,0):a>0\}\)，则\cref{g1:prop:transitive-action-model}给出 \(G/H_0\cong\mathbb R\)，对应为 \((a,b)H_0\mapsto b\)。也可直接核对：右乘 \((c,0)\) 只改变第一坐标，\(b\) 不变，且 \(x\mapsto(1,x)H_0\) 是光滑逆。每点稳定子都有一维，故作用不自由。如果 \(ax+b=x\) 对所有 \(x\) 成立，先令 \(x=0\) 得 \(b=0\)，再令 \(x=1\) 得 \(a=1\)，所以作用有效。
\end{solution}

\begin{exercise}\label{g1:ex:circle-equivariant-maps}
	设 \(a,b>0\)。加法群 \(\mathbb R\) 按平移作用于 \(\mathbb R/a\mathbb Z\) 和 \(\mathbb R/b\mathbb Z\)。求两者之间的所有 \(\mathbb R\) 等变光滑映射，并求每个映射的纤维基数及微分。它们何时为微分同胚？
\end{exercise}
\begin{solution}
	设 \(F(0+a\mathbb Z)=c+b\mathbb Z\)。等变性强制
	\[
		F(t+a\mathbb Z)=t+c+b\mathbb Z.
	\]
	代表元 \(t\) 与 \(t+a\) 给出同一个目标值，当且仅当 \(a\in b\mathbb Z\)，即 \(a=kb\)，其中 \(k\) 为正整数。这与\cref{g1:prop:equivariant-homogeneous-maps}中的 \(a\mathbb Z\subseteq b\mathbb Z\) 完全一致。若此条件不成立，就没有等变映射；若成立，参数 \(c\) 按模 \(b\mathbb Z\) 取值，给出全部映射。

	固定目标点 \(s+b\mathbb Z\)，其原像由 \(s-c+jb\)、\(j\in\mathbb Z\)，在模 \(a\mathbb Z\) 的意义下给出。因为 \(a=kb\)，恰有 \(k\) 个不同类，可取 \(0\le j<k\)。在局部提升到实直线的坐标中，\(F\) 为 \(t\mapsto t+c\)，故微分是恒等线性映射，处处为局部微分同胚。它为微分同胚恰好当 \(k=1\)，也即 \(a=b\)。注意局部微分可逆并不能排除整体上出现多个原像。
\end{solution}

\begin{exercise}\label{g1:ex:equivariant-automorphisms-normalizer}
	设 \(H\) 为 Lie 群 \(G\) 的闭子群，记正规化子 \(N_G(H)=\{a\in G:aHa^{-1}=H\}\)。证明 \(G/H\) 的全部 \(G\) 等变微分同胚与 \(N_G(H)/H\) 一一对应。核对复合次序，并说明如何选取公式使这一对应保持群乘法。
\end{exercise}
\begin{solution}
	由\cref{g1:prop:equivariant-homogeneous-maps}，等变映射必为 \(F_a(gH)=gaH\)。若它还是双射，则对任意 \(k\in G\)，
	\[
		k\cdot F_a(eH)=F_a(eH)
		\quad\Longleftrightarrow\quad F_a(kH)=F_a(eH)
		\quad\Longleftrightarrow\quad k\in H.
	\]
	另一方面，\(aH\) 的稳定子是 \(aHa^{-1}\)，所以必须有 \(aHa^{-1}=H\)。反之，若 \(a\in N_G(H)\)，则 \(F_a\) 与 \(F_{a^{-1}}\) 都有定义且光滑，互为逆。\(F_a=F_b\) 当且仅当 \(aH=bH\)，得到所需集合对应。

	\(N_G(H)\) 是子群，且 \(H\) 在其中正规，故 \(N_G(H)/H\) 是群。按普通函数复合顺序，
	\[
		(F_a\circ F_b)(gH)=gbaH=F_{ba}(gH).
	\]
	因此 \(aH\mapsto F_a\) 反转乘法。改定义 \(R_a(gH)=ga^{-1}H\)，则
	\[
		(R_a\circ R_b)(gH)=gb^{-1}a^{-1}H
		=g(ab)^{-1}H=R_{ab}(gH),
	\]
	从而 \(aH\mapsto R_a\) 是到等变微分同胚群的群同构。公式中逆元的位置由右乘与函数复合的次序决定。
\end{solution}


\begin{exercise}\label{g1:ex:screw-action-quotient}
	令 \(R_t\) 表示平面上转角 \(t\) 的旋转矩阵。加法群 \(\mathbb R\) 在 \(\mathbb R^2\times\mathbb R\) 上作用为
	\[
		t\cdot(v,z)=(R_tv,z+t).
	\]
	证明此作用自由且固有，求其基本向量场，并用一个明确的光滑不变量把商识别为 \(\mathbb R^2\)。
\end{exercise}
\begin{solution}
	旋转矩阵满足 \(R_{s+t}=R_sR_t\)，因此给出的公式确为光滑作用。若 \(t\cdot(v,z)=(v,z)\)，第三坐标给出 \(t=0\)，故自由。若 \((v_j,z_j)\to(v,z)\)，且 \(t_j\cdot(v_j,z_j)\to(w,u)\)，则由第三坐标 \(z_j+t_j\to u\)，得到 \(t_j\to u-z\)。因此\cref{g1:prop:proper-action-criteria}的序列条件成立，作用固有。

	以 \(1\in\mathbb R\) 为 Lie 代数生成元，在 \(t=0\) 求导，得到基本场
	\[
		X=-y\partial_x+x\partial_y+\partial_z.
	\]
	考虑 \(F(v,z)=R_{-z}v\)。由于 \(R_{-(z+t)}R_t=R_{-z}\)，它在轨道上常值。每条轨道恰与平面 \(z=0\) 相交一次，交点为 \((F(v,z),0)\)。因此下降映射 \(\bar F:(\mathbb R^2\times\mathbb R)/\mathbb R\to\mathbb R^2\) 双射且光滑，其逆 \(w\mapsto[(w,0)]\) 也是光滑映射。更直接地，
	\[
		(t,w)\longmapsto(R_tw,t),\qquad
		(v,z)\longmapsto(z,R_{-z}v)
	\]
	是互逆的全局光滑乘积坐标。轨道的螺旋形状在这些坐标中变为固定 \(w\)、只改变 \(t\)。
\end{solution}

\begin{exercise}\label{g1:ex:diagonal-proper-action}
	设有限维实 Lie 群 \(G\) 自由且固有地光滑作用于无边界光滑流形 \(M\)，并任意光滑作用于无边界光滑流形 \(N\)；所有流形均为有限维、Hausdorff、第二可数。证明对角作用
	\[
		g\cdot(p,n)=(gp,gn)
	\]
	在 \(M\times N\) 上自由且固有。特别地，取 \(M=G\) 且作用为左平移，证明 \((G\times N)/G\) 微分同胚于 \(N\)。
\end{exercise}
\begin{solution}
	若 \(g\) 固定 \((p,n)\)，它固定 \(p\)，所以 \(g=e\)。若 \((p_j,n_j)\) 收敛且 \((g_jp_j,g_jn_j)\) 收敛，则 \(p_j\) 与 \(g_jp_j\) 都收敛。原作用的固有性给出 \(g_j\) 的收敛子列，因此对角作用也满足序列判据。

	左平移在 \(G\) 上自由且固有：如果 \(a_j\to a\)、\(g_ja_j\to b\)，就有 \(g_j=(g_ja_j)a_j^{-1}\to ba^{-1}\)。在 \(G\times N\) 上令
	\[
		F(a,n)=a^{-1}n.
	\]
	则 \(F(ga,gn)=a^{-1}n\)，故它下降为光滑映射 \(\bar F:(G\times N)/G\to N\)。每条轨道包含唯一一点 \((e,a^{-1}n)\)，所以 \(\bar F\) 为双射，且逆映射 \(n\mapsto[(e,n)]\) 光滑。这证明所需微分同胚。注意不需要 \(G\) 在 \(N\) 上的作用自由或固有；第一因子的平移已控制了群参数。
\end{solution}

\begin{exercise}\label{g1:ex:closed-subgroup-action-quotients}
	设有限维实 Lie 群 \(G\) 自由且固有地光滑作用于有限维、Hausdorff、第二可数的无边界光滑流形 \(M\)，\(H\subseteq G\) 为闭子群。证明限制的 \(H\) 作用也自由且固有，并证明
	\[
		F:M/H\longrightarrow M/G,\qquad F(Hp)=Gp
	\]
	是光滑满射淹没。说明若把 \(H\) 换成带自身 Lie 群拓扑的非闭浸入子群，固有性不再能由同一论证保证，并给出反例。
\end{exercise}
\begin{solution}
	闭子群定理使 \(H\) 成为具有子空间拓扑的嵌入 Lie 子群。自由性直接继承。记 \(\Theta_G(g,p)=(gp,p)\)，则对每个紧集 \(C\subseteq M\times M\)，有
	\[
		\Theta_H^{-1}(C)=\Theta_G^{-1}(C)\cap(H\times M).
	\]
	右侧是紧集的闭子集，因而紧。这里 \(H\) 的闭性与子空间拓扑同时被使用。

	记商投影为 \(\pi_H\)、\(\pi_G\)。映射 \(\pi_G\) 光滑且在 \(H\) 轨道上常值，所以光滑下降性质给出唯一光滑 \(F\)，满足 \(F\circ\pi_H=\pi_G\)。由 \(\pi_G\) 满射可知 \(F\) 满射。对 \(p\in M\) 求微分，得到
	\[
		\mathrm{d}F_{Hp}\circ\mathrm{d}(\pi_H)_p=\mathrm{d}(\pi_G)_p.
	\]
	右侧满射，因此 \(\mathrm{d}F_{Hp}\) 满射，证明 \(F\) 是淹没。定义 \(F\) 无须假定 \(H\) 正规：这里比较的是轨道空间，尚未给 \(M/H\) 定义 \(G/H\) 的群作用。

	反例取 \(G=M=T^2\)，作用为左平移；它自由且因 \(G\) 紧而固有。取\cref{g1:exm:irrational-torus-subgroup}的无理斜率浸入子群 \(H_\alpha\cong\mathbb R\)。限制作用正是本章的无理流。其底空间 \(M\) 紧，但搬运集 \(T(M,M)=H_\alpha\cong\mathbb R\) 非紧，故作用不固有。
\end{solution}

\begin{exercise}\label{g1:ex:smooth-quotient-without-submersion}
	令 \(\mathbb R\) 以 \(t\cdot(x,y)=(x+t,y)\) 作用于 \(\mathbb R^2\)，轨道集记为 \(Q\)。通常用 \(y\) 给 \(Q\) 坐标。现改用
	\[
		u:Q\longrightarrow\mathbb R,\qquad u([x,y])=y^3
	\]
	作为全局坐标，赋予 \(Q\) 一份光滑结构。证明这份结构与原结构具有相同拓扑，商投影仍光滑，但不是淹没，并举出一个不再能光滑下降的不变光滑函数。这说明商流形定理的唯一性为什么要求投影为淹没。
\end{exercise}
\begin{solution}
	通常的商坐标 \([x,y]\mapsto y\) 是到 \(\mathbb R\) 的同胚；\(s\mapsto s^3\) 又是 \(\mathbb R\) 的同胚，逆为连续函数 \(u\mapsto\sqrt[3]{u}\)。因此 \(u\) 仍是商拓扑的一张全局坐标，作为单张图可以规定光滑结构。

	在新坐标下，投影 \(\pi:\mathbb R^2\to Q\) 的表达为 \(u\circ\pi(x,y)=y^3\)，故光滑。但其微分为 \(3y^2\,\mathrm{d}y\)，在直线 \(y=0\) 上为零，所以不是到一维流形的淹没。原来的函数 \(f(x,y)=y\) 光滑且平移不变，其唯一集合意义的下降在新坐标下为 \(\bar f(u)=\sqrt[3]{u}\)，在零点不可微。因此它不能光滑下降。两份结构之间，原商到新商的恒等映射表达为 \(s\mapsto s^3\)，反向为立方根，故恒等双射不是微分同胚。
\end{solution}

\begin{exercise}\label{g1:ex:closed-orbits-nonproper-action}
	加法群 \(\mathbb R\) 在 \(M=\mathbb R^2\setminus\{(0,0)\}\) 上作用为
	\[
		t\cdot(x,y)=(e^t x,e^{-t}y).
	\]
	证明作用自由，且每条轨道在 \(M\) 中闭；再证明作用不固有，轨道商也不是 Hausdorff 空间。由此说明“每条轨道闭”不能替代固有性。
\end{exercise}
\begin{solution}
	因为 \((x,y)\ne(0,0)\)，至少有一个坐标非零。若 \(e^t x=x\)、\(e^{-t}y=y\)，由该非零坐标得到 \(t=0\)，所以自由。

	乘积 \(xy\) 与每个非零坐标的符号在作用下不变。如果 \(xy=c\ne0\)，固定 \(x\) 的符号以后，曲线 \(xy=c\) 的相应分支就是一条轨道，因为两个同号的 \(x\) 值之比可写成 \(e^t\)。这条分支在 \(M\) 中闭：任何有限极限仍满足 \(xy=c\ne0\)，且 \(x\) 不会从正变负或从负变正。若 \(xy=0\)，轨道是四条坐标半轴之一；它们在去掉原点的 \(M\) 中也闭。因此所有轨道都闭。

	取 \(p_j=(e^{-j},1)\)、\(t_j=j\)，则
	\[
		p_j\longrightarrow p=(0,1),\qquad
		t_jp_j=(1,e^{-j})\longrightarrow q=(1,0).
	\]
	群参数 \(t_j\) 没有收敛子列，故作用不固有。另一方面，\(p,q\) 位于不同轨道。若商中分别包含 \([p]\)、\([q]\) 的开邻域为 \(U,V\)，其原像是分别包含 \(p,q\) 的开集。充分大的 \(j\) 满足 \(p_j\in\pi^{-1}(U)\)、\(t_jp_j\in\pi^{-1}(V)\)，而这两点属于同一轨道，故 \(U\cap V\ne\varnothing\)。所以这两个商点不能由不交开邻域分开。
\end{solution}


\begin{exercise}\label{g1:ex:finite-nonorthogonal-average}
	令 \(C=\left(\begin{smallmatrix}0&-1\\1&-1\end{smallmatrix}\right)\)。验证 \(C^3=I\)，并把 \(\mathbb R^2\) 上的标准内积对群 \(\{I,C,C^2\}\) 平均。求所得矩阵 \(Q\)，直接验证 \(C^{\mathsf T}QC=Q\)，并说明为什么 \(C\) 不保持标准内积，却能保持一个正定内积。
\end{exercise}
\begin{solution}
	相乘得到
	\[
		C^2=\begin{pmatrix}-1&1\\-1&0\end{pmatrix},\qquad C^3=I.
	\]
	有限群的归一化积分为各元素取权重 \(1/3\)，故
	\[
		Q=\frac13\bigl(I+C^{\mathsf T}C+(C^2)^{\mathsf T}C^2\bigr)
		=\frac13\begin{pmatrix}4&-2\\-2&4\end{pmatrix}.
	\]
	它的特征值为 \(2/3\)、\(2\)，所以正定。直接相乘有
	\[
		QC=\frac13\begin{pmatrix}-2&-2\\4&-2\end{pmatrix},\qquad
		C^{\mathsf T}QC=\frac13\begin{pmatrix}4&-2\\-2&4\end{pmatrix}=Q.
	\]
	另一方面，\(C^{\mathsf T}C=\left(\begin{smallmatrix}1&-1\\-1&2\end{smallmatrix}\right)\ne I\)，故标准内积不被保持。正交性依赖选定内积；有限群平均构造了一个与该群作用相容的新内积。它的单位圆在原坐标中是一条椭圆，\(C\) 保持的是这条椭圆。
\end{solution}

\begin{exercise}\label{g1:ex:average-projection-fixed-vectors}
	设 \(K\) 为紧 Lie 群，\(\rho:K\to\operatorname{GL}(V)\) 为有限维实光滑表示。逐矩阵分量定义 \(P=\int_K\rho(k)\,\mathrm{d}\mu_K(k)\)。证明 \(P^2=P\)，且 \(\operatorname{im}P=V^K\)，其中 \(V^K=\{v:\rho(k)v=v\text{ 对所有 }k\in K\}\)。进一步证明
	\[
		\dim V^K=\int_K\operatorname{tr}\rho(k)\,\mathrm{d}\mu_K(k),
	\]
	并求圆周群在 \(\mathbb R^2\) 上的标准旋转表示的 \(P\)。
\end{exercise}
\begin{solution}
	矩阵分量是连续函数，紧性保证所有积分有限。对 \(h\in K\)，利用左、右不变积分分别得到
	\[
		\rho(h)P=\int_K\rho(hk)\,\mathrm{d}\mu_K(k)=P,
		\qquad P\rho(h)=P.
	\]
	所以 \(Pv\) 被每个 \(\rho(h)\) 固定，\(\operatorname{im}P\subseteq V^K\)。若 \(v\in V^K\)，则 \(Pv=\int_Kv\,\mathrm{d}\mu_K=v\)，反向包含也成立，而且 \(P^2v=Pv\)。这证明 \(P\) 是到固定向量空间的投影，并非任意向量的平均都非零。

	投影给出直和 \(V=\operatorname{im}P\oplus\ker P\)，因为 \(v=Pv+(v-Pv)\)，两子空间相交为零。在适应该直和的基中，\(P\) 的矩阵为 \(\operatorname{diag}(I,0)\)，故 \(\operatorname{tr}P=\dim V^K\)。迹是对角分量的有限和，可以与积分交换，得到题中公式。旋转表示的矩阵为 \(R_t=\left(\begin{smallmatrix}\cos t&-\sin t\\\sin t&\cos t\end{smallmatrix}\right)\)，各分量在一周上的平均均为零，所以 \(P=0\)，也与没有非零向量被全部旋转固定一致。
\end{solution}

\begin{exercise}\label{g1:ex:sphere-invariant-tensors}
	对 \(n\geq1\)，证明 \(S^n\) 上每个 \(\operatorname{SO}(n+1)\) 不变 Riemann 度量都是单位球面标准度量的正常数倍。再求所有 \(\operatorname{SO}(n+1)\) 不变向量场，并说明 \(n=1\) 与 \(n\geq2\) 的区别。
\end{exercise}
\begin{solution}
	由\cref{g1:prop:sphere-homogeneous-model,g1:prop:invariant-tensor-isotropy}，问题等价于求 \(\mathbb R^n\) 上的 \(\operatorname{SO}(n)\) 不变正定内积及固定向量。若 \(n\geq2\)，\(\operatorname{SO}(n)\) 在单位球面 \(S^{n-1}\) 上传递：把任一单位向量补成正向正交单位基，再将两组这样的基对应，即得所需矩阵。因而一个不变内积 \(b\) 的二次型在标准单位向量上取同一值 \(c>0\)，按齐次性有 \(b(v,v)=c|v|^2\)。极化公式
	\[
		2b(v,w)=b(v+w,v+w)-b(v,v)-b(w,w)
	\]
	于是给出 \(b(v,w)=c\langle v,w\rangle\)。当 \(n=1\) 时，一维空间上的任意正定内积本来就是该形式。基点数据的唯一延拓证明了全部 \(n\geq1\) 的度量结论。

	若 \(n\geq2\)，任意非零 \(v\in\mathbb R^n\) 都能与一个垂直向量张成二维平面；在该平面转动 \(\pi\)，在正交补上取恒等，得到 \(\operatorname{SO}(n)\) 中把 \(v\) 变为 \(-v\) 的元素。因此固定向量只有零，不变向量场也只有零。若 \(n=1\)，稳定子 \(\operatorname{SO}(1)\) 平凡，基点切空间的每个向量都能延拓。将 \(S^1\) 写为 \(p=(\cos t,\sin t)\)，所有不变向量场正是
	\[
		X_p=c(-\sin t,\cos t),\qquad c\in\mathbb R.
	\]
	右侧也可写成 \(cJp\)，所以与角度分支的选择无关。
\end{solution}

\begin{exercise}\label{g1:ex:projective-line-radius}
	给 \(\mathbb{CP}^1\) 配备\eqref{g1:eq:projective-invariant-metric}的度量。证明仿射坐标 \(z=x+iy\) 中
	\[
		g=\frac{\mathrm{d}x^2+\mathrm{d}y^2}{(1+x^2+y^2)^2},
	\]
	并证明映射
	\[
		F([1:z])=\frac1{1+|z|^2}
		\left(x,y,\frac{|z|^2-1}{2}\right)
	\]
	延拓为 \(\mathbb{CP}^1\) 到半径 \(1/2\) 的 Euclid 球面的等距微分同胚。
\end{exercise}
\begin{solution}
	一维复坐标中，\((1+|z|^2)\overline uv-(\overline uz)(\overline zv)=\overline uv\)，故\eqref{g1:eq:projective-invariant-metric}立即给出题中的实度量。记 \(s=x^2+y^2\)、\(r=1+s\)。直接计算
	\[
		|F|^2=\frac{s+(s-1)^2/4}{(1+s)^2}=\frac14,
	\]
	所以像落在指定球面。有限 \(z\) 的像不含北极 \((0,0,1/2)\)，其余点 \((X,Y,Z)\) 的逆为 \(z=(X+iY)/(1/2-Z)\)。这证明该坐标片到去掉北极的球面的双射。

	在另一张射影图中令 \(w=1/z=a+ib\)。对 \(w\ne0\)，同一公式变为
	\[
		F([w:1])=\frac1{1+|w|^2}
		\left(a,-b,\frac{1-|w|^2}{2}\right),
	\]
	右侧在 \(w=0\) 光滑，取值恰为北极。由两张图及球面的相应坐标逆，可得整体微分同胚。

	最后核对度量。对 \(\xi=(u,v)\in\mathbb R^2\) 令 \(q=xu+yv\)。\(F\) 的前两个分量的微分为 \(\xi/r-2q(x,y)/r^2\)，第三个分量的微分为 \(2q/r^2\)。故
	\begin{align*}
		|\mathrm{d}F(\xi)|^2
		&=\frac{|\xi|^2}{r^2}-\frac{4q^2}{r^3}
		+\frac{4q^2s}{r^4}+\frac{4q^2}{r^4}\\
		&=\frac{|\xi|^2}{r^2}.
	\end{align*}
	极化后得到两切向量的内积也一致；北极处由另一张光滑图中的连续性成立。因此它是等距微分同胚。球半径为 \(1/2\)，具体反映本节射影度量的归一化，不能在保持同一个度量的同时把它写成单位球面。
\end{solution}

\begin{exercise}\label{g1:ex:grassmann-complement-isometry}
	设 \(V=\mathbb F^n\) 带标准内积，\(\mathbb F=\mathbb R\) 或 \(\mathbb C\)。证明正交补映射
	\[
		c:\operatorname{Gr}(k,V)\longrightarrow\operatorname{Gr}(n-k,V),
		\qquad W\longmapsto W^\perp
	\]
	是光滑微分同胚，且对\eqref{g1:eq:grassmann-hilbert-schmidt}定义的两边度量是等距映射。明确写出 \(\mathrm{d}c_W\) 在切空间 \(\operatorname{Hom}_{\mathbb F}(W,W^\perp)\) 上的表达。
\end{exercise}
\begin{solution}
	固定正交分解 \(V=W\oplus W^\perp\)，邻近的子空间写成 \(W_A=\{x+Ax:x\in W\}\)。向量 \(u+v\)、\(u\in W\)、\(v\in W^\perp\) 垂直于全部 \(x+Ax\)，当且仅当 \(u=-A^*v\)，故
	\[
		W_A^\perp=\{-A^*v+v:v\in W^\perp\}.
	\]
	在目标使用交换后的分解 \(V=W^\perp\oplus W\)，正交补映射的坐标表达就是 \(A\mapsto-A^*\)。它在实情形线性，在复情形实线性，因而都是实光滑映射；不能把复共轭误认为复线性。又因 \((W^\perp)^\perp=W\)，其逆也是同一正交补构造，故为微分同胚。

	在 \(A=0\) 处求导得到
	\[
		\mathrm{d}c_W(T)=-T^*:
		W^\perp\longrightarrow W.
	\]
	因此对 \(S,T:W\to W^\perp\)，目标度量给出
	\[
		\operatorname{Re}\operatorname{tr}\bigl((-S^*)^*(-T^*)\bigr)
		=\operatorname{Re}\operatorname{tr}(ST^*)
		=\operatorname{Re}\operatorname{tr}(S^*T).
	\]
	最后一步可在正交单位基中展开为 \(\sum_{i,j}\operatorname{Re}(s_{ij}\overline{t_{ij}})\)，与 \(\sum_{i,j}\operatorname{Re}(\overline{s_{ij}}t_{ij})\) 相等；这也适用于两个迹分别在不同维数空间上计算的情形。于是每点的内积都被保持，得到等距性。\(k=0\) 或 \(k=n\) 时两空间均为单点，结论按零切空间解释成立。
\end{solution}


\begin{exercise}\label{g1:ex:compact-action-fixed-set}
	设紧 Lie 群 \(K\) 光滑作用于无边界、Hausdorff、第二可数光滑流形 \(M\)，定义共同固定点集
	\[
		M^K=\{p\in M:kp=p\text{ 对所有 }k\in K\}.
	\]
	证明 \(M^K\) 闭，并且每个固定点附近的固定点集都是嵌入子流形，切空间为
	\[
		\mathrm{T}_p(M^K)=\{v\in \mathrm{T}_pM:\mathrm{d}(a_k)_pv=v\text{ 对所有 }k\in K\}.
	\]
	说明各连通分支的维数可以不同，并用 \(\mathbb{RP}^2\) 上的反射给出一个例子。
\end{exercise}
\begin{solution}
	对每个 \(k\)，集合 \(\{p:kp=p\}\) 是映射 \(p\mapsto(kp,p)\) 对闭对角线的原像，故闭；这些闭集的交 \(M^K\) 也闭。固定 \(p\in M^K\)，由\cref{g1:prop:fixed-point-linearization}，在 \(p\) 的某个不变邻域中有等变坐标 \(F:U\to B\subset \mathrm{T}_pM\)，且 \(\mathrm{d}F_p=I\)。令 \(V\) 为题目右侧的固定向量子空间。等变性与 \(F\) 的单射性给出
	\[
		F(U\cap M^K)=B\cap V.
	\]
	右侧是线性子空间中的开球，因而是一张嵌入坐标切片；其零点切空间是 \(V\)，经 \(\mathrm{d}F_p=I\) 送回即得所求公式。固定点集的维数因此在每个这样的局部片内恒定，从而在每个连通分支上恒定，但不同分支未必有相同维数。

	具体地，在 \(\mathbb{RP}^2\) 上令二阶群的生成元作用为
	\[
		[x:y:z]\longmapsto[x:y:-z].
	\]
	这是由线性反射诱导的光滑作用。一个射影点固定，当且仅当存在非零实数 \(\lambda\)，使 \((x,y,-z)=\lambda(x,y,z)\)。若 \((x,y)\ne(0,0)\)，必有 \(\lambda=1\)、\(z=0\)，得到射影直线 \(\{z=0\}\simeq\mathbb{RP}^1\)。若 \((x,y)=(0,0)\)，得到单点 \([0:0:1]\)。二者不相交，固定点集因此有一个一维连通分支和一个零维连通分支。故未经补充不能把整个固定点集说成同一固定维数的子流形。
\end{solution}

\begin{exercise}\label{g1:ex:weighted-circle-local-quotients}
	设正整数 \(a,b\geq1\)，令 \(S^1\) 光滑作用于 \(S^3\subset\mathbb C^2\)：
	\[
		\lambda\cdot(z_1,z_2)=(\lambda^a z_1,\lambda^b z_2).
	\]
	求两坐标均非零处及两条坐标圆上的稳定子。计算 \(p=(1,0)\)、\(q=(0,1)\) 处的横向稳定子表示，说明这些局部图何时有效；再求整个作用的核以及有效化后两处的旋转阶数。
\end{exercise}
\begin{solution}
	作用保持 \(|z_1|^2+|z_2|^2\)，并且是光滑作用。由于 \(S^1\) 紧，作用固有，因此可用本节的切片结果。记 \(d=\gcd(a,b)\)，并以 \(\mu_r=\{\lambda\in S^1:\lambda^r=1\}\) 表示 \(r\) 次单位根群。如果 \(z_1z_2\ne0\)，稳定条件为 \(\lambda^a=\lambda^b=1\)，恰好等价于 \(\lambda^d=1\)：一向由 \(d\mid a,b\)，另一向由整数 \(u,v\) 使 \(ua+vb=d\)。故一般点稳定子为 \(\mu_d\)。当 \(z_2=0\) 时稳定子是 \(\mu_a\)，当 \(z_1=0\) 时稳定子是 \(\mu_b\)。

	在 \(p=(1,0)\)，切空间为
	\[
		\mathrm{T}_pS^3=\{(\mathrm{i}t,w):t\in\mathbb R,\ w\in\mathbb C\}.
	\]
	轨道方向由 \((\mathrm{i}a,0)\) 张成，故可取横向空间 \(N_p=\{(0,w):w\in\mathbb C\}\)。\(\lambda\in\mu_a\) 在这个复直线上的作用为 \(w\mapsto\lambda^b w\)。还可以直接取小圆盘上的横截面
	\[
		j_p(w)=\bigl(\sqrt{1-|w|^2},w\bigr),
	\]
	它对 \(\mu_a\) 等变；这也在实际坐标中显示了上述表示。其核为 \(\mu_a\cap\mu_b=\mu_d\)，像是阶数 \(a/d\) 的旋转群。交换两坐标，在 \(q\) 得到 \(\mu_b\) 的表示 \(w\mapsto\lambda^a w\)，核仍为 \(\mu_d\)，像的阶数为 \(b/d\)。两张图都有效当且仅当 \(d=1\)。

	一个 \(\lambda\) 若在整个 \(S^3\) 上恒等作用，检验 \(p,q\) 即得 \(\lambda^a=\lambda^b=1\)，所以全局作用核也是 \(\mu_d\)。有效化后，一般点稳定子变平凡，两处特殊图的旋转阶数分别为 \(a/d\)、\(b/d\)。阶数为 \(1\) 时该处就是普通圆盘图。有效化前则不能删除共同的 \(\mu_d\)，否则会改变每个点的稳定子。
\end{solution}

\begin{exercise}\label{g1:ex:finite-action-faithful-linearization}
	设有限群 \(K\) 有效地光滑作用于连通无边界光滑流形 \(M\)。证明对任意 \(p\in M\)，稳定子的表示 \(K_p\to\operatorname{GL}(\mathrm{T}_pM)\) 忠实。这里“有效”指只有单位群元素在整个 \(M\) 上恒等作用。若去掉连通性，给出反例。
\end{exercise}
\begin{solution}
	设 \(h\in K_p\)，且 \(\mathrm{d}(a_h)_p=I\)。对有限循环群 \(\langle h\rangle\) 应用固定点线性化，可知 \(h\) 在 \(p\) 的某个邻域上恒等。我们还要把这条局部结论传播到整个连通流形；一般光滑映射在开集上恒等并不足够，\(h\) 的有限阶性质在下一步仍要使用。

	令
	\[
		A=\{x\in M:hx=x,\ \mathrm{d}(a_h)_x=I\}.
	\]
	这个集合非空。它闭：若 \(x_j\in A\) 趋于 \(x\)，连续性先给 \(hx=x\)；在 \(x\) 附近的一张共同坐标图中，\(a_h\) 的导数矩阵连续，所以其在 \(x\) 的值也是单位矩阵。它也开：对每个 \(x\in A\)，再次对 \(\langle h\rangle\) 在固定点 \(x\) 应用线性化，得到 \(h\) 在 \(x\) 的整个小邻域上恒等，因而该邻域均在 \(A\) 中。连通性迫使 \(A=M\)，所以 \(h\) 恒等作用。全局有效性给出 \(h=e\)，证明忠实性。

	若 \(M=\mathbb R\sqcup\mathbb R\)，让二阶群在第一份直线上恒等作用，在第二份上按 \(t\mapsto-t\) 作用。这是有效光滑作用，但第一份上每个点都被整个群固定，且稳定子的微分表示全为恒等。反例表明，连通性负责把局部不动传播到全局，不能省略。
\end{solution}

\begin{exercise}\label{g1:ex:effective-transitive-ineffective-slice}
	令 \(G=\operatorname{SO}(3)\)，\(r=\operatorname{diag}(-1,-1,1)\)，\(H=\{I,r\}\)。考虑 \(G\) 在齐性空间 \(M=G/H\) 上的左作用。证明这个作用有效、固有，所有稳定子均为二阶群，但每个点的横向切片都是零维且带平凡的稳定子表示。说明它为什么不与上一题矛盾，又为什么不能将该商只按一个普通点理解而保留所有稳定子信息。
\end{exercise}
\begin{solution}
	\(H\) 是有限闭子群，因此\cref{g1:thm:homogeneous-space}给出 \(G/H\) 的光滑流形结构，维数为 \(3\)。左作用传递；在 \(gH\) 处，稳定子由 \(agH=gH\) 算得为 \(gHg^{-1}\)，所以均为二阶群。\(G\) 紧，故作用固有。

	整个作用的核是 \(\bigcap_{g\in G}gHg^{-1}\)，特别地包含于 \(H\)。非单位元素 \(r\) 是绕第三坐标轴的半周旋转；取把第三坐标轴转到第一坐标轴的 \(g\)，则 \(grg^{-1}=\operatorname{diag}(1,-1,-1)\ne r\)。两个二阶子群 \(H\) 与 \(gHg^{-1}\) 的交只有单位元，故全局核平凡，作用有效。

	传递性使每个点的轨道就是整个 \(M\)，所以 \(\mathrm{T}_p(Gp)=\mathrm{T}_pM\)，横向空间 \(N\) 为零空间。其零点球只有一个点，整个二阶稳定子都在这一点上平凡作用；故切片图非有效。上一题讨论的是有限群作用在整个连通流形上的微分表示，这里作用群 \(G\) 正维，而且我们只保留商掉轨道方向以后的横向表示，假设与表示对象都不同。

	底层商 \(M/G\) 是一个点，而切片图还记录 \(H\simeq\mathbb Z/2\mathbb Z\)。把该图有效化后成为带平凡群的一个点，底拓扑不变，稳定子却丢失。这一例子说明全局作用有效，并不保证横向切片表示有效。
\end{solution}


% !TeX root = ../../Book1.tex
% 纲要标题：## 第18章　主丛、规范变换与平坦结构
% 纲要位置：几何学计划纲要.md，第 807 行。
% 正文状态：初稿；尚未正式审核。

% 纲要内容开始
% 
% > 本章目标与证明范围：将过渡函数、水平分布和局部规范场统一，证明主丛联络与伴随丛联络的关系。规范等价和基空间坐标变化必须各有清楚作用对象。
% 
% 纲要内容结束

\chapter{主丛、规范变换与平坦结构}
\label{g1:ch:18}
\BookChapterNumber{g1:ch:18}

\begin{chapterabstract}
	主丛把同一纤维中的标架选择组织为群的自由传递作用。过渡函数描述这些选择如何粘合，主联络则规定沿底空间移动时如何比较它们。本章构造主丛、关联丛与结构群约化，证明水平分布、联络形式和局部联络数据的等价性，并建立关联向量丛上的协变微分。水平提升的全区间存在、曲率结构方程、Bianchi 恒等式及运输变分说明局部微分怎样控制路径比较。最后区分规范变换、换局部截面与底空间坐标变化，以单值表示分类平坦主丛，并用圆周、环面和 Hopf 丛检验这份分类的范围。
\end{chapterabstract}

\begin{learninggoals}
	\begin{itemize}
		\item 从局部截面求过渡函数，由余循环重建主丛，构造关联丛与标架丛，并证明结构群约化的截面判据。
		\item 在水平分布、主联络形式与局部形式之间转换，证明主联络存在，计算联络之差及关联丛上的诱导协变微分。
		\item 证明紧参数区间上路径的水平提升存在且唯一，建立曲率结构方程与 Bianchi 恒等式，用曲率计算运输变分和小环路的首项。
		\item 区分改变局部截面、施行规范变换与改变底空间坐标，计算规范轨道，并用环路运输的中心化子确定规范稳定子。
		\item 由单值表示构造平坦主丛，证明表示模共轭的分类及固定底层丛时的限制，计算圆周、环面的相位和 Hopf 丛的曲率积分。
	\end{itemize}
\end{learninggoals}

给一个向量丛选取局部标架之后，纤维中的向量就有了列坐标。但我们也可以把标架本身当作研究对象：同一点的两组标架之间有唯一可逆矩阵，却没有哪一组天然应当充当单位标架。若要求标架正交，允许的变换便缩为正交群；若所记录的只是一个单位复数相位，允许的变换则是圆周群。主丛用同一种结构容纳这些选择，联络在这个选择空间中规定如何随底点移动。

\ProseParagraphBreak

局部表达与整体比较提出不同的问题。一个变化的标架能使原本为零的联络矩阵出现非零项，但没有改变联络；一条环路的水平运输也可能不平凡，即使曲率处处为零。我们需要同时保留总空间上的几何对象、截面中的表达和环路留下的群元素。平坦结构的分类会把最后一类信息组织成基本群表示，同时说明有些底层主丛根本不容许平坦联络。

\ProseParagraphBreak

本章的 Lie 群均为有限维实 Lie 群，流形默认有限维、无边界、Hausdorff、第二可数；复线性表示另作说明。主丛统一使用右作用 \(p\mapsto pg\)，局部截面满足 \(s_j=s_i g_{ij}\)。因此其基本场取 \(X_p^\#=\left.\frac{\mathrm{d}}{\mathrm{d}t}\right|_0p\exp(tX)\)，括号符号须与前章默认的左作用区分。局部联络公式延续\chapref{g1:ch:15}的行标架与列分量约定；Lie 群及光滑商使用\chapref{g1:ch:16}和\chapref{g1:ch:17}。水平提升、平坦性和单值构造分别调用前面已经建立的常微分方程、Frobenius 定理与万有覆盖，使用时会恢复所需条件。

\ProseParagraphBreak

对照阅读可查 Tu 的主丛、主联络、水平提升、曲率及关联丛各节\cite[第27--31节]{Tu2017Connections}。比较局部公式时，先核对主作用在左还是在右，以及过渡函数从哪一张截面换到哪一张；这些选择决定逆元与共轭的位置。

% !TeX root = ../../Book1.tex
% 纲要标题：### 18.1 主丛和关联构造
% 纲要位置：几何学计划纲要.md，第 811 行。
% 正文状态：初稿；尚未正式审核。

% 纲要内容开始
% 
% * 主作用、局部截面与过渡函数
% * 由余循环重建主丛
% * 关联向量丛和标架丛
% * 结构群约化及截面判据
% 
% 纲要内容结束

\section{主丛和关联构造}
\label{g1:sec:1801}

在向量丛的同一条纤维中，两组标架之间有唯一的可逆换基矩阵，但没有一组标架天然应当充当单位元。主丛保留的正是这种“任意两点可以唯一比较，却未指定起点”的结构。\cref{g1:def:frame-bundle}已经给出了一个实例：标架 \(u:\mathbb R^r\to E_x\) 可以右乘矩阵，\(u\cdot A=u\circ A\)。下面把矩阵群换成一般 Lie 群，先确定局部选择如何相容，再由这些选择构造向量丛及额外几何结构。本节的流形均为有限维、无边界、Hausdorff、第二可数的光滑流形，Lie 群也采用\chapref{g1:ch:16}的约定。

\subsection{主作用、局部截面与过渡函数}
\label{g1:subsec:1801:01}

在主丛的一条纤维中，选定一个参考点后，每个点都由唯一的群元素得到；参考点随底点光滑变化时，便形成局部乘积。局部截面负责选取这些参考点，过渡函数记录两套选择之间的群元素。

\begin{definition}[主丛]\label{g1:def:principal-bundle}
	设 \(G\) 为有限维实 Lie 群，\(P,M\) 为无边界、Hausdorff、第二可数的光滑流形。一个光滑\term[zhucong]{主丛}[principal bundle]由光滑满射 \(\pi:P\to M\) 和光滑右作用 \(P\times G\to P\)、\((p,g)\mapsto pg\) 组成，满足 \(\pi(pg)=\pi(p)\)，且每个 \(x\in M\) 都有开邻域 \(U\) 及微分同胚
	\[
		\Phi_U:U\times G\longrightarrow\pi^{-1}(U),
		\qquad \pi(\Phi_U(x,a))=x,\qquad
		\Phi_U(x,ag)=\Phi_U(x,a)g.
	\]
	称 \(G\) 为结构群，\(P\) 为总空间，\(M\) 为底空间。覆盖恒等映射的主丛同构是保持右作用的微分同胚；覆盖一般光滑映射的主丛映射也要求保持右作用与投影。
\end{definition}

局部模型说明每条纤维 \(P_x\) 是一条自由的 \(G\) 轨道：给定 \(p,q\in P_x\)，存在唯一 \(g\in G\) 使 \(q=pg\)。因此，在一个纤维中选定 \(p\) 后，\(g\mapsto pg\) 才把它识别为群 \(G\)；没有这个选择时，不能在 \(P_x\) 中直接把两个点相乘。局部平凡化提供光滑截面 \(s_U(x)=\Phi_U(x,e)\)，并满足 \(\Phi_U(x,g)=s_U(x)g\)。反过来，任意光滑局部截面都给出这样的平凡化：在已有平凡化中写 \(s_U(x)=\Phi_U^0(x,h(x))\)，则相应映射为 \((x,g)\mapsto\Phi_U^0(x,h(x)g)\)，其逆由左乘 \(h(x)^{-1}\) 给出。

\begin{proposition}[主丛与自由固有作用]\label{g1:prop:principal-free-proper}
	设 \(G\) 为有限维实 Lie 群，所有流形均无边界、Hausdorff、第二可数。主 \(G\) 丛 \(\pi:P\to M\) 的右作用自由且固有，纤维恰为轨道，并且 \(M\) 与轨道商 \(P/G\) 自然微分同胚。反之，\(G\) 在 \(P\) 上的任意光滑自由固有右作用都使商投影 \(P\to P/G\) 成为主 \(G\) 丛。
\end{proposition}
\begin{proof}
	自由性与纤维等于轨道已由局部模型得到。为证明固有性，设 \(p_j\to p\)、\(p_jg_j\to q\)。投影连续给出 \(\pi(p)=\pi(q)=x\)。选取 \(x\) 附近的主丛平凡化，把充分大的 \(j\) 对应的两点写成 \((x_j,a_j)\)、\((x_j,b_j)\)，则 \(a_j\to a\)、\(b_j\to b\)，并且 \(b_j=a_jg_j\)。故 \(g_j=a_j^{-1}b_j\to a^{-1}b\)。\cref{g1:prop:proper-action-criteria}的序列判据适用于右作用，只须把它改写为左作用 \(g*p=pg^{-1}\)；群元素取逆不会改变子列收敛条件。因此作用固有。

	\cref{g1:thm:free-proper-quotient}赋予 \(P/G\) 光滑结构。诱导的集合双射 \(P/G\to M\) 在任一局部乘积中都是 \([s_U(x)g]\mapsto x\)，它及其逆均光滑，因而为微分同胚。

	反过来，对左作用 \(g*p=pg^{-1}\) 应用\cref{g1:lem:free-proper-local-product}，得到横截面 \(S\) 与 \(G\times S\to SG\)、\((g,s)\mapsto sg^{-1}\) 的微分同胚。交换两个因子并取逆，便得到右作用所需的 \(S\times G\to SG\)、\((s,g)\mapsto sg\)。商定理又把 \(S\) 微分同胚地识别为底空间的开集，所以这正是主丛平凡化。
\end{proof}

\begin{example}[齐性空间上的主丛]\label{g1:exm:homogeneous-principal-bundle}
	设 \(H\) 为有限维实 Lie 群 \(G\) 的闭子群。\cref{g1:thm:homogeneous-space}不仅给出 \(G/H\) 的光滑结构，还给出局部乘积 \((\exp W)\times H\to(\exp W)H\) 及其左平移。因此 \(G\to G/H\) 是主 \(H\) 丛，右作用为群内的右乘。这里结构群是 \(H\)，不是 \(G\)；底空间也不必是一个群。
\end{example}

现在取一族局部截面 \(s_i:U_i\to P\)。交叠处，每条纤维内的自由传递性给出唯一函数 \(g_{ij}:U_i\cap U_j\to G\)，使
\begin{equation}\label{g1:eq:principal-transition}
	s_j(x)=s_i(x)g_{ij}(x),\qquad
	g_{ii}=e,\qquad g_{ij}g_{jk}=g_{ik}.
\end{equation}
在 \(s_i\) 给出的平凡化中，\(g_{ij}\) 就是 \(s_j\) 的群坐标，所以光滑。三重交叠上，\(s_k=s_jg_{jk}=s_i g_{ij}g_{jk}\)，再由自由性得到余循环等式；特别地 \(g_{ji}=g_{ij}^{-1}\)。同一点 \(p\in P_x\) 若写成 \(s_i(x)a_i=s_j(x)a_j\)，则
\begin{equation}\label{g1:eq:principal-coordinate-transition}
	a_i=g_{ij}(x)a_j.
\end{equation}
截面从 \(i\) 换到 \(j\) 时右乘 \(g_{ij}\)；表示同一点的群坐标从 \(j\) 换到 \(i\) 时，按上式左乘它。

\ProseParagraphBreak

如果换用 \(s_i'=s_i h_i\)，其中 \(h_i:U_i\to G\) 光滑，则
\begin{equation}\label{g1:eq:principal-transition-change}
	g_{ij}'=h_i^{-1}g_{ij}h_j.
\end{equation}
这是在同一个主丛上改选局部截面，底点 \(x\) 保持不变。局部联络形式也随截面改变，其变换律将同时涉及这里的群元素及其导数。

\subsection{由余循环重建主丛}
\label{g1:subsec:1801:02}

反过来，能否只给出 \(g_{ij}\)，就得到一个总空间？\cref{g1:thm:bundle-cocycle-gluing}已经对向量丛证明过这种粘合。这里用群中的左乘代替矩阵作用；需要保留的不只是局部公式，还包括不同纤维能够分离、总空间具有可数坐标基这些整体条件。

\begin{theorem}[主丛的余循环构造]\label{g1:thm:principal-cocycle-gluing}
	设 \(M\) 为有限维、无边界、Hausdorff、第二可数光滑流形，\(G\) 为有限维实 Lie 群，\(\{U_i\}_{i\in I}\) 为 \(M\) 的开覆盖。设光滑映射 \(g_{ij}:U_i\cap U_j\to G\) 满足 \(g_{ii}=e\)、\(g_{ij}g_{jk}=g_{ik}\)。则存在主 \(G\) 丛 \(\pi:P\to M\) 及局部截面 \(s_i\)，满足 \(s_j=s_i g_{ij}\)。它在保持这些局部截面的主丛同构意义下唯一。
\end{theorem}
\begin{proof}
	取不交并 \(\mathcal P=\coprod_i(U_i\times G)\)，规定
	\[
		(i,x,a)\sim(j,y,b)
		\quad\Longleftrightarrow\quad
		x=y\quad\text{且}\quad a=g_{ij}(x)b.
	\]
	自反性来自 \(g_{ii}=e\)，对称性来自 \(g_{ji}=g_{ij}^{-1}\)，传递性来自 \(g_{ij}g_{jk}=g_{ik}\)。令 \(P=\mathcal P/\!\sim\)，赋予商拓扑，定义 \(\pi([i,x,a])=x\)。关系没有改变 \(x\)，故 \(\pi\) 良定义且连续。

	商映射 \(Q:\mathcal P\to P\) 为开映射。若 \(O\subset\mathcal P\) 开，其饱和集在第 \(j\) 个副本中的部分为
	\[
		\bigcup_i\{(x,b):x\in U_i\cap U_j,\ (i,x,g_{ij}(x)b)\in O\}.
	\]
	每项是光滑同胚 \((x,b)\mapsto(x,g_{ij}(x)b)\) 下的开原像，所以饱和集开。于是
	\[
		\Phi_i:U_i\times G\longrightarrow\pi^{-1}(U_i),\qquad
		\Phi_i(x,a)=[i,x,a]
	\]
	是连续开双射，因而为同胚。双射性来自每个相应等价类可以唯一换算到第 \(i\) 个副本。它们的转换为
	\[
		\Phi_i^{-1}\Phi_j(x,b)=(x,g_{ij}(x)b),
	\]
	光滑且具有光滑逆。把 \(U_i\) 和 \(G\) 再限制到坐标域，便得到维数 \(\dim M+\dim G\) 的相容图册。

	若两点投影不同，用 \(M\) 中不交开邻域的逆像分离；若投影相同，把它们放进同一个开集 \(\pi^{-1}(U_i)\)，利用 \(U_i\times G\) 的 Hausdorff 性分离。故 \(P\) 为 Hausdorff 空间。第二可数性也要单独说明：从 \(M\) 的一个可数基中，对每个包含于某个 \(U_i\) 的基元素选择这样的 \(U_i\)，所选 \(U_i\) 构成可数子覆盖。各 \(U_i\times G\) 都有可数基，其像的可数并就给出 \(P\) 的可数基。因而任意开覆盖参与粘合也不会破坏总空间的第二可数性。

	定义 \([i,x,a]g=[i,x,ag]\)。若 \(a=g_{ij}b\)，则 \(ag=g_{ij}(bg)\)，故右作用良定义；在每个 \(\Phi_i\) 中就是群的右乘，因而光滑。\(\pi\) 局部为乘积投影，\(\Phi_i\) 为等变平凡化，故得到主丛。取 \(s_i(x)=[i,x,e]\)，则
	\[
		s_j(x)=[j,x,e]=[i,x,g_{ij}(x)]=s_i(x)g_{ij}(x).
	\]
	最后，若另一个主丛有同样关系的截面 \(t_i\)，公式 \([i,x,a]\mapsto t_i(x)a\) 良定义，在各局部乘积中是恒等坐标映射，故为主丛同构。保持指定截面与右作用的任何同构都必须使用这一公式，因此唯一。
\end{proof}

余循环本身仍依赖截面。如果两组余循环满足\eqref{g1:eq:principal-transition-change}，把旧余循环构造出的点 \([i,x,a]\) 送到新余循环构造中的 \([i,x,h_i(x)^{-1}a]\)，便得到主丛同构：由 \(a=g_{ij}b\) 可算出 \(h_i^{-1}a=g_{ij}'h_j^{-1}b\)，所以公式与代表无关，且局部光滑可逆。覆盖不同时，先取两组开集的交集作为共同细化即可比较。这里保持局部截面的唯一性，与更换截面后允许的主丛同构，是两种不同的陈述。

\begin{example}[有全局截面的主丛]\label{g1:exm:principal-global-section}
	主 \(G\) 丛 \(P\to M\) 与 \(M\times G\) 同构，当且仅当它有一个光滑全局截面 \(s\)。一方面，乘积中的 \(x\mapsto(x,e)\) 是截面；另一方面，\((x,g)\mapsto s(x)g\) 在每条纤维上双射，其在任何局部平凡化中的表达为 \((x,g)\mapsto(x,h(x)g)\)，故它与逆都光滑。局部截面在定义中总存在，全局截面的存在则是额外结论。向量丛即使非平凡也有零截面，主丛纤维没有相应的零点，因此不能用零截面的直觉推断这个条件。
\end{example}

\subsection{关联向量丛和标架丛}
\label{g1:subsec:1801:03}

主丛的一点可以理解为一种局部参照。若一个空间 \(F\) 上也规定了 \(G\) 怎样改变参照，就应能把 \(F\) 随着主丛一起粘到各个底点上。关键是：主丛中的参照乘上 \(g\) 后，\(F\) 中的分量必须按 \(g^{-1}\) 改变，才能表示同一个对象。

\begin{proposition}[关联丛]\label{g1:prop:associated-fiber-bundle}
	设 \(\pi:P\to M\) 为有限维实 Lie 群 \(G\) 的光滑主丛，\(F\) 为具有光滑左 \(G\) 作用的有限维光滑流形；所有流形均无边界、Hausdorff、第二可数。规定
	\[
		(p,y)\cdot g=(pg,g^{-1}\cdot y),\qquad
		P\times_G F\coloneqq(P\times F)/G.
	\]
	该商具有自然的光滑流形结构，投影 \([p,y]\mapsto\pi(p)\) 局部平凡，纤维为 \(F\)。若主丛截面满足 \(s_j=s_i g_{ij}\)，则局部表示 \((x,y_i)\mapsto[s_i(x),y_i]\) 的转换为 \(y_i=g_{ij}(x)\cdot y_j\)。称此丛为 \(P\) 的\term[guanliancong]{关联丛}[associated bundle]。
\end{proposition}
\begin{proof}
	先核对所写公式确为右作用：连续作用 \(g\) 再作用 \(h\)，第二因子变为 \(h^{-1}\cdot(g^{-1}\cdot y)=(gh)^{-1}\cdot y\)。它自由，因为第一因子的作用自由。它也固有：若 \((p_j,y_j)\to(p,y)\)，且 \((p_jg_j,g_j^{-1}\cdot y_j)\) 收敛，则主丛作用的固有性给出 \(g_j\) 的收敛子列，正好满足\cref{g1:prop:proper-action-criteria}的判据。于是\cref{g1:thm:free-proper-quotient}使商为无边界、Hausdorff、第二可数光滑流形，商映射为淹没。

	在一个主丛平凡化中，每个 \(p\) 唯一写成 \(s_i(x)a\)。定义
	\[
		[p,y]\longmapsto(x,a\cdot y),\qquad
		(x,z)\longmapsto[s_i(x),z].
	\]
	若 \((p,y)\) 换成 \((pg,g^{-1}\cdot y)\)，则第一式的群坐标由 \(a\) 变成 \(ag\)，输出仍为 \((x,a\cdot y)\)，所以良定义。两式互逆；第二式显然光滑，第一式与商映射的复合光滑，故由\cref{g1:prop:quotient-smooth-universal-property}也光滑。它们给出所需局部平凡化。最后
	\[
		[s_j(x),y_j]=[s_i(x)g_{ij}(x),y_j]
		=[s_i(x),g_{ij}(x)\cdot y_j],
	\]
	得到转换公式。
\end{proof}

\begin{theorem}[关联向量丛]\label{g1:thm:associated-vector-bundle}
	设 \(\pi:P\to M\) 为有限维实 Lie 群 \(G\) 的光滑主右丛，\(M,P\) 无边界、Hausdorff、第二可数，\(V\) 为有限维实或复向量空间，\(\rho:G\to\operatorname{GL}(V)\) 为光滑线性表示。取开覆盖 \(\{U_i\}\) 上的光滑局部截面 \(s_i:U_i\to P\)，并以 \(s_j=s_i g_{ij}\) 定义过渡函数 \(g_{ij}:U_i\cap U_j\to G\)。则
	\[
		E=P\times_\rho V\coloneqq(P\times V)/G,
		\qquad(p,v)\cdot g=(pg,\rho(g)^{-1}v)
	\]
	具有自然的秩 \(\dim_{\mathbb K}V\) 的光滑 \(\mathbb K\) 向量丛结构，其中 \(\mathbb K=\mathbb R\) 或 \(\mathbb C\)。局部坐标 \([s_i(x),v_i]\) 满足
	\begin{equation}\label{g1:eq:associated-vector-transition}
		v_i=\rho(g_{ij})v_j.
	\end{equation}
	对固定 \(p\in P_x\)，映射 \(v\mapsto[p,v]\) 是 \(V\to E_x\) 的线性同构。
\end{theorem}
\begin{proof}
	应用\cref{g1:prop:associated-fiber-bundle}，所得局部转换为 \((x,v_j)\mapsto(x,\rho(g_{ij}(x))v_j)\)。因为 \(\rho\) 是同态，这些可逆线性转换满足同样的余循环等式。由此，局部乘积中的加法、数乘在交叠处相容，给出各纤维与参照无关的向量空间结构，局部平凡化线性且光滑。复情形只把所有纤维运算和转换理解为复线性的，底空间及总空间的光滑结构仍按实流形理解。固定 \(p\) 后，每个类都可唯一写成 \([p,v]\)，故得到所述同构。
\end{proof}

这类构造也常称“伴随向量丛”。本书用“关联向量丛”表示一般 \(\rho\) 的情形，把伴随表示 \(\operatorname{Ad}\) 给出的特例记为 \(\operatorname{ad}(P)\)。表示不必忠实；关联构造可能丢掉主丛的一部分信息，不能从某一个关联向量丛自动恢复原主丛。

\begin{proposition}[截面与等变函数]\label{g1:prop:associated-sections}
	设 \(P\to M\) 为有限维实 Lie 群 \(G\) 的光滑主丛，\(\rho:G\to\operatorname{GL}(V)\) 为有限维实或复光滑表示，所有流形无边界、Hausdorff、第二可数。关联向量丛 \(E=P\times_\rho V\) 的光滑截面，与满足
	\begin{equation}\label{g1:eq:associated-equivariant-function}
		f(pg)=\rho(g^{-1})f(p)
	\end{equation}
	的光滑函数 \(f:P\to V\) 一一对应，对应式为 \(\sigma(x)=[p,f(p)]\)、\(p\in P_x\)。取开覆盖 \(\{U_i\}\) 上的光滑主丛截面 \(s_i\)，并以 \(s_j=s_i g_{ij}\) 定义过渡函数；其局部分量为 \(f_i=f\circ s_i\)，并满足 \(f_i=\rho(g_{ij})f_j\)。
\end{proposition}
\begin{proof}
	给定 \(\sigma\)，固定 \(p\in P_x\) 后，上一结果给出的线性同构使 \(\sigma(x)=[p,f(p)]\) 唯一确定 \(f(p)\)。换用 \(pg\) 时，\([pg,f(pg)]=[p,\rho(g)f(pg)]\)，故等于原截面值当且仅当\eqref{g1:eq:associated-equivariant-function}成立。在 \(p=s_i(x)a\) 的局部坐标中，\(f(p)=\rho(a^{-1})f_i(x)\)，因此截面光滑蕴含 \(f\) 光滑。

	反过来，等变关系使 \([p,f(p)]\) 与 \(p\in P_x\) 的选择无关，所以给出一个截面。它在 \(s_i\) 中的分量是光滑函数 \(f_i\)，因而光滑。两种构造互逆；而 \(s_j=s_i g_{ij}\) 给出 \(f_j=\rho(g_{ij}^{-1})f_i\)，等价于所写分量转换。
\end{proof}

\begin{proposition}[标架丛与向量丛]\label{g1:prop:frame-associated-equivalence}
	设 \(M\) 为有限维、无边界、Hausdorff、第二可数光滑流形，\(E\to M\) 为秩 \(r\geq1\) 的实光滑向量丛。\(\operatorname{Fr}(E)\to M\) 是主 \(\operatorname{GL}(r,\mathbb R)\) 丛，且标准表示给出自然向量丛同构
	\[
		\operatorname{Fr}(E)\times_{\operatorname{GL}(r,\mathbb R)}\mathbb R^r
		\longrightarrow E,\qquad [u,v]\longmapsto u(v).
	\]
	反之，若 \(P\to M\) 是主 \(\operatorname{GL}(r,\mathbb R)\) 丛，令 \(E=P\times_{\operatorname{GL}(r,\mathbb R)}\mathbb R^r\)，则
	\[
		P\longrightarrow\operatorname{Fr}(E),\qquad
		p\longmapsto u_p,\qquad u_p(v)=[p,v]
	\]
	是自然的主丛同构。复向量丛与 \(\operatorname{GL}(r,\mathbb C)\) 满足相同结论。
\end{proposition}
\begin{proof}
	\cref{g1:def:frame-bundle}已给出 \(\operatorname{Fr}(E)\) 的局部坐标 \(u=e_i(x)A\)，右作用在其中是 \((x,A)B=(x,AB)\)，所以满足主丛定义。第一式在等价代表下不变，因为 \((uA)(A^{-1}v)=u(v)\)。在参考标架 \(e_i\) 下，它把 \([e_i(x),v]\) 送到 \(e_i(x)v\)，局部就是恒等向量坐标，故光滑可逆并且逐纤维线性。

	对于第二式，\(u_p\) 是线性同构，且
	\[
		u_{pA}(v)=[pA,v]=[p,Av]=(u_p\circ A)(v),
	\]
	所以映射保持右作用。固定 \(p\in P_x\)，任何标架 \(u:\mathbb R^r\to E_x\) 都唯一写成 \(u_p\circ A\)，从而此映射在每条纤维上双射。在局部截面 \(s_i\) 给出的主丛坐标及标架 \(u_{s_i}\) 给出的标架丛坐标中，它为 \((x,A)\mapsto(x,A)\)，故为微分同胚。复情形逐步使用复线性同构，证明相同。
\end{proof}

\begin{definition}[伴随丛]\label{g1:def:adjoint-bundle}
	设 \(P\to M\) 为有限维实 Lie 群 \(G\) 的光滑主丛，\(\mathfrak g\) 为 \(G\) 的 Lie 代数。由伴随表示构造的实向量丛
	\[
		\operatorname{ad}(P)\coloneqq P\times_{\operatorname{Ad}}\mathfrak g
	\]
	称为\term[banshuicong]{伴随丛}[adjoint bundle]。其等价关系为 \((p,X)\sim(pg,\operatorname{Ad}_{g^{-1}}X)\)，取光滑局部截面 \(s_i:U_i\to P\)，以 \(s_j=s_i g_{ij}\) 定义过渡函数，并将局部元素写为 \([s_i(x),X_i(x)]\)。则局部分量满足 \(X_i=\operatorname{Ad}_{g_{ij}}X_j\)。
\end{definition}

每条纤维还带有 Lie 括号：用同一个 \(p\in P_x\) 表示两个元素，定义 \(\bigl[[p,X],[p,Y]\bigr]=[p,[X,Y]]\)，这里左侧外层括号表示纤维中的 Lie 括号。由于 \(\operatorname{Ad}_g[X,Y]=[\operatorname{Ad}_gX,\operatorname{Ad}_gY]\)，改变 \(p\) 后所得类不变；在局部坐标中就是 \(\mathfrak g\) 的固定双线性括号，故随底点光滑。若 \(G\) 交换，则所有 \(\operatorname{Ad}_g\) 都是恒等映射，于是 \([p,X]\mapsto(\pi(p),X)\) 给出典范同构 \(\operatorname{ad}(P)\cong M\times\mathfrak g\)，即使 \(P\) 本身可能并不平凡。后文的曲率会取值于这个伴随丛。它的纤维是带括号的向量空间 \(\mathfrak g\)，而原主丛的纤维是一条自由传递的群轨道。

\subsection{结构群约化及截面判据}
\label{g1:subsec:1801:04}

给向量丛指定内积后，正交单位标架之间只需正交矩阵换算。\cref{g1:def:frame-structure-reduction}已把这称为结构群约化。对于一般主 \(G\) 丛，问题变为：能否在每条 \(G\) 纤维中，光滑地选出一条固定闭子群 \(H\subseteq G\) 的轨道？选择单个点会使主丛平凡，而选择一条 \(H\) 轨道只要求减少允许的参照。

\begin{definition}[主丛的结构群约化]\label{g1:def:principal-reduction}
	设 \(P\to M\) 为有限维实 Lie 群 \(G\) 的光滑主丛，\(H\subseteq G\) 为闭子群，所有流形无边界、Hausdorff、第二可数。一个 \(H\) 结构群约化是 \(P\) 的光滑嵌入子流形 \(Q\)，它在右 \(H\) 作用下不变，且限制投影 \(Q\to M\) 与该作用组成主 \(H\) 丛。
\end{definition}

这里要求 \(H\) 闭，是为了使用\cref{g1:thm:homogeneous-space}中 \(G/H\) 的 Hausdorff 光滑结构及局部截面。我们将 \(P/H\) 定义为主丛总空间对右 \(H\) 作用的轨道商；它的一个点是一条 \(H\) 轨道，并非整条 \(G\) 纤维。仍有投影 \(r:P/H\to M\)、\(pH\mapsto\pi(p)\)。

\begin{theorem}[约化的截面判据]\label{g1:thm:reduction-section-criterion}
	设 \(\pi:P\to M\) 为有限维实 Lie 群 \(G\) 的光滑主丛，\(H\subseteq G\) 为闭子群，所有流形无边界、Hausdorff、第二可数。令 \(P/H\) 为 \(P\) 对右 \(H\) 作用的轨道商，赋予自然商流形结构，并定义 \(r:P/H\to M\)、\(r(pH)=\pi(p)\)。则 \(r\) 是纤维为 \(G/H\) 的光滑局部平凡丛，并有自然同构
	\[
		P\times_G(G/H)\longrightarrow P/H,\qquad
		[p,gH]\longmapsto pgH.
	\]
	主丛 \(P\) 的 \(H\) 约化与 \(r\) 的光滑截面一一对应，具体为
	\[
		Q\longmapsto\sigma_Q(x)=Q_x\in P/H,
		\qquad
		\sigma\longmapsto Q_\sigma=\{p\in P:pH=\sigma(\pi(p))\}.
	\]
	若 \(Q\) 为约化，则由 \(H\) 在 \(G\) 上的左乘构造的丛 \(Q\times_HG\) 具有自然主 \(G\) 丛结构，并有主丛同构
	\begin{equation}\label{g1:eq:extension-structure-group}
		Q\times_HG\longrightarrow P,\qquad[q,g]\longmapsto qg.
	\end{equation}
\end{theorem}
\begin{proof}
	先构造 \(P/H\)。右 \(H\) 作用自由。若 \(p_j\to p\)、\(p_jh_j\to q\)，主 \(G\) 作用的固有性给出 \(h_j\) 在 \(G\) 中的收敛子列；因为 \(H\) 闭，极限仍在 \(H\)，故 \(H\) 作用固有。商定理给出 \(P/H\) 的光滑结构。在主丛局部坐标中，商映射为
	\[
		U_i\times G\longrightarrow U_i\times(G/H),\qquad
		(x,g)\longmapsto(x,gH).
	\]
	\cref{g1:thm:homogeneous-space}保证 \(G\to G/H\) 有光滑局部截面，故上述商坐标和其逆光滑。交叠时 \(g_jH\) 变成 \(g_{ij}(x)g_jH\)，于是 \(r\) 局部平凡。公式 \([p,gH]\mapsto pgH\) 不随代表变化，因为 \((pa)(a^{-1}g)H=pgH\)；它在刚才的局部坐标中为恒等映射，所以给出所述自然同构。

	若已有 \(H\) 约化 \(Q\)，每个 \(Q_x\) 是唯一的一条 \(H\) 轨道，故 \(\sigma_Q\) 良定义。在 \(Q\) 的局部截面 \(t:U\to Q\) 中，\(\sigma_Q(x)=t(x)H\)，是光滑映射，且 \(r\sigma_Q=\operatorname{id}_M\)。

	反过来，给定光滑截面 \(\sigma\)，在 \(P\) 的局部截面 \(s_i\) 下，把它写成 \(\sigma(x)=(x,k(x)H)\)。这里首先确定的是一个光滑映射 \(U_i\to G/H\)，不能预先假定已有全局的 \(G\) 值函数 \(k\)。在任一点附近，使用 \(G\to G/H\) 的光滑局部截面，并适当缩小 \(U_i\)，才得到光滑提升 \(k:U_i\to G\)。于是
	\[
		Q_\sigma|_{U_i}=\{s_i(x)k(x)h:x\in U_i,\ h\in H\}.
	\]
	在 \(P|_{U_i}\cong U_i\times G\) 中，光滑坐标改变 \((x,g)\mapsto(x,k(x)^{-1}g)\) 把它送到 \(U_i\times H\)。闭子群 \(H\) 是嵌入 Lie 子群，故这些局部表达证明 \(Q_\sigma\) 是嵌入子流形，且 \((x,h)\mapsto s_i(x)k(x)h\) 给出主 \(H\) 丛的局部平凡化。它继承 \(P\) 的 Hausdorff 性与第二可数性。两种构造逐纤维互逆，故得到一一对应。

	最后，\(Q\times_HG\) 的关系为 \((q,g)\sim(qh,h^{-1}g)\)。右 \(G\) 作用定义为 \([q,g]a=[q,ga]\)，与等价关系相容；取 \(Q\) 的局部截面 \(t\)，则 \((x,g)\mapsto[t(x),g]\) 给出主 \(G\) 丛平凡化。\eqref{g1:eq:extension-structure-group}与代表无关，因为 \((qh)(h^{-1}g)=qg\)，并保持右 \(G\) 作用。在上述平凡化及 \(t\) 作为 \(P\) 的局部截面所给出的平凡化中，它为恒等坐标映射，故为主丛同构。
\end{proof}

由此还能直接读出过渡函数判据：\(P\) 有一个 \(H\) 约化，当且仅当经过细化覆盖，可以选取局部截面使全部 \(g_{ij}\) 取值于 \(H\)。若已有 \(Q\)，取它的局部截面即可；反过来，\(Q|_{U_i}=s_i(U_i)H\) 因 \(H\) 值过渡而相容，并由局部 \(U_i\times H\) 成为约化。这恢复了\chapref{g1:ch:12}以标架定义的结构群约化。

\begin{example}[度量选择的是一类标架]\label{g1:exm:metric-as-principal-reduction}
	设 \(E\to M\) 为秩 \(r\geq1\) 的实光滑向量丛。取 \(G=\operatorname{GL}(r,\mathbb R)\)、\(H=\operatorname{O}(r)\)。给定丛度量 \(h\)，全部正交单位标架组成 \(Q\subset\operatorname{Fr}(E)\)。\cref{g1:prop:metric-orthogonal-reduction}的光滑 Gram--Schmidt 构造给出局部正交单位标架，所以它是主 \(H\) 约化。对应的截面 \(M\to\operatorname{Fr}(E)/\operatorname{O}(r)\) 在每点选择的是“所有对 \(h_x\) 正交的标架”这一类，而不是某一组标架。

	同一事实可直接逐纤维核验：一个标架 \(u:\mathbb R^r\to E_x\) 确定内积 \(h_u(v,w)=\langle u^{-1}v,u^{-1}w\rangle\)，把 \(u\) 右乘正交矩阵不会改变它；反之，给出相同内积的两个标架恰相差一个正交矩阵。于是该商纤维就是 \(E_x\) 上所有正定内积的空间。度量总存在，并不意味着可以在整个底空间上选出一组光滑正交标架；后者还要求这个主 \(\operatorname{O}(r)\) 丛本身平凡。
\end{example}
% 纲要条目排为编号小节标题，并进入目录与书签；说明段仍保留于纲要注释中。

% !TeX root = ../../Book1.tex
% 纲要标题：### 18.2 主联络的等价描述
% 纲要位置：几何学计划纲要.md，第 818 行。
% 正文状态：初稿；尚未正式审核。

% 纲要内容开始
% 
% * 水平分布和等变性
% * 联络形式及再现条件
% * 局部联络形式的粘合
% * 主联络和关联丛上的协变微分
% 
% 纲要内容结束

\section{主联络的等价描述}
\label{g1:sec:1802}

主丛的一条纤维记录同一点上所有允许的标架或参考状态。沿底空间移动时，我们希望把一个状态连续地带到邻近纤维，但局部平凡化提供的“保持群坐标不变”会随所选截面改变。联络要做的是在总空间中规定哪些速度算作纯粹的底空间运动，再把这个规定写成可以计算、可以换局部截面的形式。本节固定有限维实 Lie 群 \(G\) 及光滑主右 \(G\) 丛 \(\pi:P\to M\)；所有流形无边界、Hausdorff 且第二可数。

\subsection{水平分布和等变性}
\label{g1:subsec:1802:01}

先分清两种速度。在 \(p\in P\) 处，沿纤维运动的速度组成
\[
	V_pP\coloneqq\ker(\mathrm{d}\pi)_p.
\]
局部平凡化把 \(\pi\) 写成 \(U\times G\to U\) 的投影，故这些核构成秩为 \(\dim G\) 的光滑子丛 \(VP\subset \mathrm{T}P\)，称为\term[chuizhifenbu]{垂直分布}[vertical distribution]。仅知道底空间速度 \(v\in \mathrm{T}_xM\)，还不能确定总空间速度：若 \(w\) 投影到 \(v\)，加上任意垂直向量仍投影到 \(v\)。

\begin{definition}\label{g1:def:principal-horizontal-connection}
	设 \(\pi:P\to M\) 为光滑主右 \(G\) 丛。一个光滑子丛 \(H\subset \mathrm{T}P\)，若对任意 \(p\in P\)、\(g\in G\) 满足
	\[
		\mathrm{T}_pP=H_p\oplus V_pP,\qquad
		(\mathrm{d}R_g)_pH_p=H_{pg},
	\]
	就称为一个\term[zhulianluo]{主联络}[principal connection]，其中 \(R_g(p)=pg\)。\(H\) 称为它的\term[shuipingfenbu]{水平分布}[horizontal distribution]，\(H_p\) 中的向量称为水平向量。
\end{definition}

直和条件使 \((\mathrm{d}\pi)_p|_{H_p}:H_p\to \mathrm{T}_{\pi(p)}M\) 成为线性同构。因此每个 \(v\in \mathrm{T}_{\pi(p)}M\) 有唯一水平向量 \(v_p^H\) 投影到它；这个向量称为 \(v\) 在 \(p\) 处的水平提升。其对 \(p,v\) 的依赖光滑，因为在局部丛标架中只是在求一个可逆矩阵的逆。等变条件进一步给出
\[
	(\mathrm{d}R_g)_p v_p^H=v_{pg}^H.
\]
它要求先改变纤维中的参考状态再提升，与先提升再作同一个右平移得到相同结果。任意垂直补空间未必满足这一条件；主联络要求这组选择与群作用相容。

\ProseParagraphBreak

垂直空间本身有一个来自群的自然坐标。对 \(X\in\mathfrak g=\mathrm{T}_eG\)，沿用基本向量场的构造，但本章使用右作用，记
\begin{equation}\label{g1:eq:principal-fundamental-field}
	X_p^\#\coloneqq\left.\frac{\mathrm{d}}{\mathrm{d}t}\right|_{t=0}p\exp(tX).
\end{equation}
这里使用右作用，基本场的括号因而带正号；第17章的左作用约定则带负号。

\begin{proposition}[主作用的垂直方向]\label{g1:prop:principal-fundamental-fields}
	设 \(G\) 为有限维实 Lie 群，\(\pi:P\to M\) 为光滑主右 \(G\) 丛。记 \(\mathfrak g=\operatorname{Lie}(G)\)，对 \(X\in\mathfrak g\) 令 \(X_p^\#=\left.\frac{\mathrm{d}}{\mathrm{d}t}\right|_{t=0}p\exp(tX)\)，并记 \(V_pP=\ker(\mathrm{d}\pi)_p\)、\(R_g(p)=pg\)。则映射 \(X\mapsto X_p^\#\) 是 \(\mathfrak g\to V_pP\) 的线性同构，并随 \(p\) 光滑变化。其基本场满足
	\begin{equation}\label{g1:eq:principal-fundamental-identities}
		(\mathrm{d}R_g)_pX_p^\#=(\operatorname{Ad}_{g^{-1}}X)_{pg}^\#,
		\qquad [X^\#,Y^\#]=[X,Y]^\#.
	\end{equation}
\end{proposition}
\begin{proof}
	取局部截面 \(s:U\to P\) 及平凡化 \(\Phi(x,h)=s(x)h\)。在这些坐标中，\(X^\#\) 对应 \((0,X_h^L)\)，其中 \(X_h^L=(\mathrm{d}L_h)_eX\) 是第16章的左不变场。因此它逐点给出所有垂直方向，而且 \(X\mapsto X_h^L\) 是随 \(h\) 光滑变化的同构。左不变场的括号定义了 \(\mathfrak g\) 的括号，故本章右作用的正号基本场保持括号。

	由\cref{g1:prop:exponential-naturality}的共轭公式，
	\[
		p\exp(tX)g=pg\exp(t\operatorname{Ad}_{g^{-1}}X).
	\]
	在零点求导便得到右平移公式。若要与\cref{g1:prop:fundamental-fields}比较，可把当前右作用改写为左作用 \(g\cdot p=pg^{-1}\)；那个左作用的基本场是 \(-X^\#\)，代入其反同态公式也得到此处的正号。
\end{proof}

\begin{example}\label{g1:exm:horizontal-additive-bundle}
	取加法群 \(G=\mathbb R\)、主丛 \(P=\mathbb R_x\times\mathbb R_t\to\mathbb R_x\)，右作用为 \((x,t)a=(x,t+a)\)。垂直方向由 \(\partial_t\) 张成。给定光滑函数 \(a(x)\)，令
	\[
		H_{(x,t)}=\operatorname{span}\{\partial_x-a(x)\partial_t\}.
	\]
	它与垂直方向互补，且右平移不改变 \(a(x)\)、\(\partial_x\) 和 \(\partial_t\)，所以是主联络。底空间单位速度的水平提升带有纤维分量 \(-a(x)\partial_t\)。水平不是图面上预先指定的“横着走”，而是由联络规定如何同时调整纤维状态。若把系数改为依赖 \(t\) 的一般函数 \(a(x,t)\)，仍有垂直补空间，却通常不再右平移不变。
\end{example}

\subsection{联络形式及再现条件}
\label{g1:subsec:1802:02}

有了水平与垂直的分解，任意速度 \(w\in \mathrm{T}_pP\) 都能写成“水平部分加 \(X_p^\#\)”。把这个 \(X\) 记下来，就得到一个 \(\mathfrak g\) 值一形式。这里 \(\Omega^k(P;\mathfrak g)\) 表示取值于固定向量空间 \(\mathfrak g\) 的交替 \(k\) 形式，等价于\cref{g1:def:bundle-valued-form}中平凡丛 \(P\times\mathfrak g\) 值的形式。取一组 \(\mathfrak g\) 的基后可逐分量拉回、求外微分；换一组常数基只作线性组合，故所得运算与基无关。

\begin{definition}\label{g1:def:principal-connection-form}
	设 \(G\) 为有限维实 Lie 群，\(\pi:P\to M\) 为光滑主右 \(G\) 丛，\(\mathfrak g=\operatorname{Lie}(G)\)。对 \(X\in\mathfrak g\) 令 \(X_p^\#=\left.\frac{\mathrm{d}}{\mathrm{d}t}\right|_{t=0}p\exp(tX)\)，并令 \(R_g(p)=pg\)。若光滑一形式 \(\omega\in\Omega^1(P;\mathfrak g)\) 对每个 \(X\in\mathfrak g\)、\(g\in G\) 都满足
	\begin{equation}\label{g1:eq:principal-connection-axioms}
		\omega_p(X_p^\#)=X,\qquad
		R_g^*\omega=\operatorname{Ad}_{g^{-1}}\omega
	\end{equation}
	则称其为\term[zhulianluoxingshi]{主联络形式}[principal connection form]。第一式称为再现条件；第二式的意思是 \(\omega_{pg}((\mathrm{d}R_g)_pw)=\operatorname{Ad}_{g^{-1}}(\omega_p(w))\)。
\end{definition}

再现条件固定了联络在垂直方向上的值，等变条件说明这些垂直坐标在右移后如何变化。\(\omega\) 不要求右平移不变：基本场自身按 \(\operatorname{Ad}_{g^{-1}}\) 变换，记录它的形式也必须作同样变换。只有在伴随作用平凡的情形，第二式才成为通常的不变性。

\begin{theorem}[主联络的两种描述]\label{g1:thm:principal-connection-equivalence}
	设 \(G\) 为有限维实 Lie 群，\(\pi:P\to M\) 为光滑主右 \(G\) 丛。对 \(X\in\operatorname{Lie}(G)\) 记 \(X_p^\#=\left.\frac{\mathrm{d}}{\mathrm{d}t}\right|_{t=0}p\exp(tX)\)。则右等变垂直补空间意义的主联络与主联络形式一一对应。对应由
	\[
		H_p=\ker\omega_p,\qquad
		w=w^H+\bigl(\omega_p(w)\bigr)_p^\#
	\]
	给出，其中 \(w^H\) 是 \(w\) 的水平部分。
\end{theorem}
\begin{proof}
	从 \(H\) 出发，沿光滑直和 \(\mathrm{T}P=H\oplus VP\) 投影到 \(VP\)，再用\cref{g1:prop:principal-fundamental-fields}的垂直同构取逆，得到光滑丛映射 \(\omega:\mathrm{T}P\to\mathfrak g\)。它逐纤维线性，因此是一形式。对 \(X_p^\#\)，垂直投影不作改变，故 \(\omega_p(X_p^\#)=X\)。若 \(w=w^H+X_p^\#\)，右平移把第一项送到水平空间，把第二项送到 \((\operatorname{Ad}_{g^{-1}}X)_{pg}^\#\)。于是右平移后的垂直坐标是 \(\operatorname{Ad}_{g^{-1}}X\)，正是等变条件。

	反向给定 \(\omega\)。再现条件使 \(\omega_p:\mathrm{T}_pP\to\mathfrak g\) 处处满射，故其核为光滑的秩 \(\dim M\) 子丛；这也可在局部矩阵中选可逆子块直接求核的光滑基。任意 \(w\) 都有分解
	\[
		w=\bigl(w-(\omega_p(w))_p^\#\bigr)+(\omega_p(w))_p^\#.
	\]
	第一项被 \(\omega_p\) 消去，第二项垂直。若一个垂直向量 \(X_p^\#\) 也在核中，再现条件给出 \(X=0\)，所以这个和是直和。等变条件将核送入核，应用 \(R_{g^{-1}}\) 又得到反向包含，故核右平移不变。两种构造都通过同一个唯一的水平—垂直分解恢复原对象，因此互为逆。
\end{proof}

\cref{g1:exm:horizontal-additive-bundle}的形式就是 \(\omega=\mathrm{d}t+a(x)\,\mathrm{d}x\)。它在 \(\partial_t\) 上取值为 \(1\)，在 \(\partial_x-a(x)\partial_t\) 上取值为零，因而把两种描述直接联系起来。

\ProseParagraphBreak

计算一般主丛的局部公式前，要把群方向上的微分写准确。回顾\cref{g1:prop:lie-exponential-differential}之前构造的形式
\[
	\theta_g^L=(\mathrm{d}L_{g^{-1}})_g:\mathrm{T}_gG\longrightarrow\mathfrak g.
\]
它称为左\term[maurercartanxingshi]{Maurer--Cartan 形式}[Maurer--Cartan form]。\footnote{Maurer，Ludwig Maurer（1859--1927），德国数学家，研究连续变换群与早期代数群理论。}上标强调采用左平移识别切空间；矩阵群中可以简写为 \(\theta^L=g^{-1}\mathrm{d}g\)，而一般 Lie 群始终用平移微分解释。

\begin{proposition}[群乘法的微分形式]\label{g1:prop:maurer-cartan-product}
	设 \(G\) 为有限维实 Lie 群，\(\theta^L_g=(\mathrm{d}L_{g^{-1}})_g:\mathrm{T}_gG\to\operatorname{Lie}(G)\) 为左 Maurer--Cartan 形式，\(m(g,h)=gh\)。对 \(v\in \mathrm{T}_gG\)、\(w\in \mathrm{T}_hG\)，有
	\begin{equation}\label{g1:eq:maurer-cartan-product}
		\theta^L_{gh}\bigl(\mathrm{d}m_{(g,h)}(v,w)\bigr)
		=\operatorname{Ad}_{h^{-1}}\theta^L_g(v)+\theta^L_h(w).
	\end{equation}
	特别地，\(L_a^*\theta^L=\theta^L\)、\(R_a^*\theta^L=\operatorname{Ad}_{a^{-1}}\theta^L\)。若 \(\rho:G\to\operatorname{GL}(V)\) 是有限维实或复表示，\(\rho_*=(\mathrm{d}\rho)_e\)，则
	\begin{equation}\label{g1:eq:maurer-cartan-representation}
		\rho_*\theta^L=\rho^{-1}\mathrm{d}\rho.
	\end{equation}
	右侧是 \(\operatorname{End}(V)\) 值一形式，表示先求 \(\rho\) 的微分，再左乘矩阵 \(\rho(g)^{-1}\)。
\end{proposition}
\begin{proof}
	乘法的微分为 \(\mathrm{d}R_h(v)+\mathrm{d}L_g(w)\)。写 \(v=(\mathrm{d}L_g)_eX\)，第一项是曲线 \(g\exp(tX)h=gh\exp(t\operatorname{Ad}_{h^{-1}}X)\) 的零点速度，应用 \(\theta^L_{gh}\) 后得到 \(\operatorname{Ad}_{h^{-1}}X\)。写 \(w=(\mathrm{d}L_h)_eY\)，第二项是 \(gh\exp(tY)\) 的速度，得到 \(Y\)。线性相加便得乘法公式。分别固定第一个或第二个因子，即得左右平移公式。

	同态满足 \(\rho\circ L_g=L_{\rho(g)}\circ\rho\)。由\cref{g1:prop:lie-homomorphism-differential}中这一恒等式的微分，\(\mathrm{d}\rho_g((\mathrm{d}L_g)_eX)=\rho(g)\rho_*(X)\)。左乘 \(\rho(g)^{-1}\)，便得最后一式。
\end{proof}

沿任何一条主纤维取轨道参数 \(\alpha_p(g)=pg\)，再现条件给出 \(\alpha_p^*\omega=\theta^L\)：因为 \(\mathrm{d}\alpha_p(X_g^L)=X_{pg}^\#\)。因此 \(\omega\) 的纤维方向部分已经完全固定；可以自由选择的，是底空间方向怎样与这些垂直方向混合。

\subsection{局部联络形式的粘合}
\label{g1:subsec:1802:03}

取主丛的局部截面 \(s_i:U_i\to P\)，令 \(\Phi_i(x,g)=s_i(x)g\)。把联络形式限制到截面上，得到 \(\mathfrak g\) 值一形式 \(A_i=s_i^*\omega\)。它只在底空间 \(U_i\) 上，通常比定义在整个 \(P\) 上的 \(\omega\) 更便于计算。若 \(G\) 是矩阵群，\(A_i\) 是矩阵值一形式，因此也称局部联络矩阵；一般 \(G\) 的 \(A_i\) 仍是 \(\mathfrak g\) 值形式，不必先选择一个矩阵实现。

\begin{theorem}[主联络的局部数据]\label{g1:thm:principal-connection-local-data}
	设 \(G\) 为有限维实 Lie 群，\(\pi:P\to M\) 为光滑主右 \(G\) 丛，\(\mathfrak g=\operatorname{Lie}(G)\)。取开覆盖 \(\{U_i\}\) 上的光滑局部截面 \(s_i:U_i\to P\)，以 \(s_j=s_i g_{ij}\) 定义过渡函数。令 \(\Phi_i:U_i\times G\to P|_{U_i}\)、\(\Phi_i(x,g)=s_i(x)g\)，并令 \(\theta^L_g=(\mathrm{d}L_{g^{-1}})_g\) 为左 Maurer--Cartan 形式。任意主联络形式 \(\omega\) 由 \(A_i=s_i^*\omega\) 唯一确定，且
	\begin{equation}\label{g1:eq:principal-connection-trivialization}
		(\Phi_i^*\omega)_{(x,g)}(v,\xi)
		=\operatorname{Ad}_{g^{-1}}\bigl((A_i)_x(v)\bigr)+\theta_g^L(\xi).
	\end{equation}
	在交叠域上，这些局部形式满足
	\begin{equation}\label{g1:eq:principal-connection-change}
		A_j=\operatorname{Ad}_{g_{ij}^{-1}}A_i+g_{ij}^*\theta^L.
	\end{equation}
	反过来，满足此规律的一族 \(A_i\in\Omega^1(U_i;\mathfrak g)\) 通过\eqref{g1:eq:principal-connection-trivialization}唯一地粘成主联络形式。
\end{theorem}
\begin{proof}
	\(\mathrm{d}\Phi_i(v,0)\) 是 \((\mathrm{d}s_i)v\) 经 \(R_g\) 推前的向量，等变条件使其 \(\omega\) 值等于 \(\operatorname{Ad}_{g^{-1}}A_i(v)\)。写 \(\xi=(\mathrm{d}L_g)_eX\)，则 \(\mathrm{d}\Phi_i(0,\xi)=X_{s_i(x)g}^\#\)，再现条件给出 \(X=\theta_g^L(\xi)\)。相加得到局部表示，因此它必然唯一。

	截面 \(s_j\) 在第 \(i\) 个平凡化中的表达为 \(x\mapsto(x,g_{ij}(x))\)。把这个映射代入局部表示，便得到\eqref{g1:eq:principal-connection-change}。

	反向给定 \(A_i\)，先在每个 \(U_i\times G\) 上定义右端的一形式。基本场为 \((0,X_g^L)\)，所以再现条件成立。右乘固定的 \(h\) 后，第一项变为 \(\operatorname{Ad}_{h^{-1}}\operatorname{Ad}_{g^{-1}}A_i\)，第二项由\cref{g1:prop:maurer-cartan-product}也变为 \(\operatorname{Ad}_{h^{-1}}\theta^L\)，因而等变条件成立。

	交叠处的坐标变换是 \(\Psi_{ij}(x,h)=(x,g_{ij}(x)h)\)，满足 \(\Phi_j=\Phi_i\circ\Psi_{ij}\)。记 \(\widetilde\omega_i=\operatorname{Ad}_{g^{-1}}A_i+\theta^L\)，乘法公式给出
	\begin{align*}
		\Psi_{ij}^*\widetilde\omega_i
		&=\operatorname{Ad}_{h^{-1}}\bigl(\operatorname{Ad}_{g_{ij}^{-1}}A_i+g_{ij}^*\theta^L\bigr)+\theta^L\\
		&=\operatorname{Ad}_{h^{-1}}A_j+\theta^L.
	\end{align*}
	这里各项均已拉回到交叠域与 \(G\) 的乘积，省略了底空间投影的记号。结果正是第 \(j\) 张图的形式，所以它们在总空间重叠部分相等，粘合为光滑 \(\omega\)。两条联络条件是局部可验证的，故粘合后仍成立。
\end{proof}

在矩阵群中，以上两式分别是
\[
	\Phi_i^*\omega=g^{-1}A_i g+g^{-1}\mathrm{d}g,
	\qquad A_j=g_{ij}^{-1}A_i g_{ij}+g_{ij}^{-1}\mathrm{d}g_{ij}.
\]
第二式恰是\cref{g1:prop:connection-local-data}中的换标架公式。现在可以看见其中两项的不同来源：共轭来自纤维参考状态的改变，附加微分项来自新截面本身随基点移动。局部截面 \(s_i\) 的全部切向量水平当且仅当 \(A_i=0\)。这样的截面比局部平凡化多满足一个微分条件；下一节将由曲率判断它何时存在。

\begin{example}\label{g1:exm:principal-circle-local-change}
	在平凡主圆周丛 \(U\times\operatorname{U}(1)\) 上，\(\mathfrak u(1)=i\mathbb R\)。选一个实一形式 \(a\)，令 \(A=i a\)，则
	\[
		\omega=i a+z^{-1}\mathrm{d}z.
	\]
	因为圆周群交换，伴随项没有共轭变化。若换截面 \(s'=s e^{i\varphi(x)}\)，其中 \(\varphi\) 为光滑实函数，则 \(A'=i(a+\mathrm{d}\varphi)\)。这只是同一个 \(\omega\) 的新表达。一个一般的 \(U\to\operatorname{U}(1)\) 映射未必有全局实值角度 \(\varphi\)，但 \(g^{-1}\mathrm{d}g\) 总有定义，所以用群值过渡函数写的公式比角度公式具有更大的适用范围。
\end{example}

变换律中的附加微分项说明，各个截面下的零形式一般不相容。整体存在性可以沿用单位分解的方法；这次把局部联络形式拉到总空间后加权，权重之和为一使其保留竖直方向的再现条件。

\begin{theorem}[主联络的存在]\label{g1:thm:principal-connection-existence}
	设 \(G\) 为有限维实 Lie 群，\(M\) 为无边界、Hausdorff、第二可数光滑流形。每个光滑主右 \(G\) 丛 \(P\to M\) 都有光滑主联络，无需假定 \(G\) 或 \(M\) 紧。
\end{theorem}
\begin{proof}
	取局部截面覆盖 \(\{U_i\}\)，在 \(\pi^{-1}(U_i)\) 上令 \(A_i=0\)，用\eqref{g1:eq:principal-connection-trivialization}得到局部主联络 \(\omega_i\)。这些 \(\omega_i\) 不要求在交叠处相同。由\cref{g1:thm:smooth-partition-unity}取从属单位分解 \(\{\chi_i\}\)，支集包含于 \(U_i\)、支集族局部有限、总和为 \(1\)。定义
	\[
		\omega=\sum_i(\chi_i\circ\pi)\omega_i,
	\]
	每项在 \(\pi^{-1}(U_i)\) 外补零。因为 \(\operatorname{supp}\chi_i\subset U_i\)，补零后在该开集外每点附近恒为零，故光滑；底空间支集族局部有限，拉回后的求和也局部有限。

	每个 \(\chi_i\circ\pi\) 右平移不变，故等变条件可逐项相加。再现条件为
	\[
		\omega(X^\#)=\sum_i(\chi_i\circ\pi)X=X.
	\]
	因此 \(\omega\) 是主联络形式。这里没有对非紧群作平均；起作用的是底空间的单位分解以及联络公理对权重和为 \(1\) 的加法稳定。
\end{proof}

\begin{proposition}[主联络的仿射空间]\label{g1:prop:principal-connection-affine}
	设 \(P\to M\) 为上述光滑主右 \(G\) 丛。任意两个主联络形式之差 \(\beta=\omega_1-\omega_0\) 在垂直方向为零，并满足 \(R_g^*\beta=\operatorname{Ad}_{g^{-1}}\beta\)。这些形式与 \(\Omega^1(M;\operatorname{ad}P)\) 自然一一对应。因此主联络的集合是以 \(\Omega^1(M;\operatorname{ad}P)\) 为方向空间的非空仿射空间。
\end{proposition}
\begin{proof}
	两个再现条件相减给出 \(\beta(X^\#)=0\)，两个等变条件相减给出所述变换式。回顾\cref{g1:def:adjoint-bundle}，伴随丛为 \(\operatorname{ad}P=P\times_{\operatorname{Ad}}\mathfrak g\)，满足 \([pg,Y]=[p,\operatorname{Ad}_gY]\)。对 \(v\in \mathrm{T}_xM\)，选 \(p\in P_x\) 及任意 \(w\in \mathrm{T}_pP\) 使 \(\mathrm{d}\pi(w)=v\)，令
	\[
		B_x(v)=[p,\beta_p(w)].
	\]
	同一点 \(p\) 处改变提升 \(w\) 只增加垂直向量，不影响值；把 \(p\) 改为 \(pg\)，可取 \(\mathrm{d}R_g(w)\)，此时
	\[
		[pg,\beta_{pg}(\mathrm{d}R_gw)]
		=[pg,\operatorname{Ad}_{g^{-1}}\beta_p(w)]
		=[p,\beta_p(w)].
	\]
	所以 \(B\) 与选择无关。在局部截面中它由 \(s_i^*\beta\) 表示，因而光滑且逐点对 \(v\) 线性。

	反向给定 \(B\in\Omega^1(M;\operatorname{ad}P)\)，对每个 \(p\in P_x\)、\(w\in \mathrm{T}_pP\)，由线性同构 \(Y\mapsto[p,Y]\) 唯一解出 \(\beta_p(w)\)，使 \([p,\beta_p(w)]=B_x(\mathrm{d}\pi(w))\)。局部表达保证光滑，定义直接保证消去垂直方向以及所需等变性。故 \(\omega_0+\beta\) 保留两条联络公理，也是主联络。两种对应互为逆，且任何两个联络之间有唯一的差；存在定理使集合非空，于是得到所称的仿射空间。
\end{proof}

这与\cref{g1:prop:connection-affine-space}的向量丛结论相呼应：单个联络包含一个规定好的微分部分，两联络之差才是底空间上的丛值一形式。\(\omega\) 本身不能直接下降成 \(\operatorname{ad}P\) 值一形式，因为它在垂直方向上的再现值通常非零。

\subsection{主联络和关联丛上的协变微分}
\label{g1:subsec:1802:04}

设 \(\rho:G\to\operatorname{GL}(V)\) 为有限维实或复表示，\(E=P\times_\rho V\)。\cref{g1:prop:associated-sections}把 \(E\) 的截面 \(\sigma\) 与等变函数 \(f:P\to V\) 对应：
\[
	\sigma(x)=[p,f(p)],\qquad
	f(pg)=\rho(g)^{-1}f(p).
\]
普通微分 \(\mathrm{d}f\) 还会记录沿纤维改变参考状态所造成的分量变化。协变微分只沿水平提升求导，正好排除这部分变化。

\begin{theorem}[关联丛上的诱导联络]\label{g1:thm:associated-covariant-derivative}
	设 \(G\) 为有限维实 Lie 群，\(\pi:P\to M\) 为光滑主右 \(G\) 丛，\(\omega\) 为主联络，\(\rho:G\to\operatorname{GL}(V)\) 为有限维实或复光滑表示。令 \(E=P\times_\rho V\)，采用 \((p,v)\sim(pg,\rho(g)^{-1}v)\) 的商约定，并记 \(\rho_*=(\mathrm{d}\rho)_e\)。对 \(\sigma\in\Gamma(E)\)，令 \(f:P\to V\) 为满足 \(f(pg)=\rho(g)^{-1}f(p)\)、\(\sigma(x)=[p,f(p)]\) 的光滑函数。对 \(v\in\mathrm{T}_xM\)、\(p\in P_x\)，以 \(v_p^H\) 表示满足 \((\mathrm{d}\pi)_p(v_p^H)=v\)、\(\omega_p(v_p^H)=0\) 的水平提升。则公式
	\begin{equation}\label{g1:eq:associated-horizontal-derivative}
		(\nabla\sigma)_x(v)=[p,(\mathrm{d}f)_p(v_p^H)],
		\qquad p\in P_x,
	\end{equation}
	定义与 \(p\) 的选择无关的丛联络。在局部截面 \(s_i\) 下，令 \(f_i=f\circ s_i\)、\(A_i=s_i^*\omega\)，则
	\begin{equation}\label{g1:eq:associated-local-derivative}
		\nabla\sigma=[s_i,\mathrm{d}f_i+\rho_*(A_i)f_i].
	\end{equation}
	这里的方括号对每个切向输入求值；实情形得到实联络，复情形得到对截面复线性的复联络。
\end{theorem}
\begin{proof}
	水平提升满足 \(v_{pg}^H=\mathrm{d}R_g(v_p^H)\)，等变函数满足 \(f\circ R_g=\rho(g)^{-1}f\)。对后一式求导，得到
	\[
		(\mathrm{d}f)_{pg}(v_{pg}^H)=\rho(g)^{-1}(\mathrm{d}f)_p(v_p^H).
	\]
	按关联丛的等价关系，两边与 \(pg,p\) 组合后表示同一纤维向量，证明选择无关。光滑性可在任意局部截面下检验。对底空间函数 \(h\)，截面 \(h\sigma\) 对应 \((h\circ\pi)f\)；在水平向量上使用普通乘积法则，得到 \(\nabla(h\sigma)=\mathrm{d}h\otimes\sigma+h\nabla\sigma\)。方向输入的函数线性及截面输入的实或复线性由微分直接继承，故确为联络。

	为求局部表达，先沿垂直基本场微分等变函数。指数自然性给出
	\[
		f(p\exp(tX))=\rho(\exp(tX))^{-1}f(p),
		\qquad (\mathrm{d}f)_p(X_p^\#)=-\rho_*(X)f(p).
	\]
	把 \(w=w^H+(\omega_p(w))_p^\#\) 代入，便得
	\begin{equation}\label{g1:eq:equivariant-covariant-derivative}
		(\mathrm{d}f)_p(w^H)=(\mathrm{d}f)_p(w)+\rho_*(\omega_p(w))f(p).
	\end{equation}
	因此总空间上的表达为 \(\mathrm{d}f+\rho_*(\omega)f\)，它消去垂直方向。拉回到 \(s_i\) 就得到\eqref{g1:eq:associated-local-derivative}。

	还可直接核对局部粘合。令 \(r_{ij}=\rho(g_{ij})\)。由\cref{g1:prop:lie-homomorphism-differential}的伴随交织公式及\eqref{g1:eq:maurer-cartan-representation}，
	\[
		\rho_*(A_j)=r_{ij}^{-1}\rho_*(A_i)r_{ij}+r_{ij}^{-1}\mathrm{d}r_{ij}.
	\]
	而 \(f_i=r_{ij}f_j\)，故乘积法则给出
	\[
		\mathrm{d}f_i+\rho_*(A_i)f_i
		=r_{ij}\bigl(\mathrm{d}f_j+\rho_*(A_j)f_j\bigr).
	\]
	这正是关联丛值一形式的分量变换，也说明局部公式与第15章的联络定义一致。
\end{proof}

\begin{corollary}[伴随丛的协变微分]\label{g1:cor:adjoint-connection-local}
	设 \(G\) 为有限维实 Lie 群，\(\pi:P\to M\) 为光滑主右 \(G\) 丛，\(\omega\) 为主联络，\(\mathfrak g=\operatorname{Lie}(G)\)，\(\operatorname{ad}P=P\times_{\operatorname{Ad}}\mathfrak g\)。取光滑局部截面 \(s_i:U_i\to P\)，令 \(A_i=s_i^*\omega\)，并把 \(\xi\in\Gamma(\operatorname{ad}P)\) 写为 \(\xi|_{U_i}=[s_i,\xi_i]\)。则 \(\omega\) 在 \(\operatorname{ad}P\) 上诱导联络，其局部表达为
	\[
		\nabla^{\operatorname{ad}}\xi=[s_i,\mathrm{d}\xi_i+[A_i,\xi_i]],
		\qquad [A_i,\xi_i](v)=[A_i(v),\xi_i].
	\]
\end{corollary}
\begin{proof}
	在上一定理中取 \(\rho=\operatorname{Ad}\)。\cref{g1:prop:adjoint-differential-bracket}证明了 \(\rho_*(X)=\operatorname{ad}_X\)，且 \(\operatorname{ad}_X(Y)=[X,Y]\)，代入即得公式。
\end{proof}

从主联络到关联丛联络的方向对所有表示都成立，反方向则取决于表示能看见多少信息。标架丛给出一个确切的双向对应。

\begin{proposition}[标架丛与线性联络]\label{g1:prop:frame-connection-equivalence}
	设 \(E\to M\) 为秩 \(r\geq1\) 的实光滑向量丛，其标架丛为 \(\operatorname{Fr}(E)\)，结构群为 \(\operatorname{GL}(r,\mathbb R)\)。\(\operatorname{Fr}(E)\) 上的主联络与 \(E\) 上的线性联络自然一一对应，正向对应为标准表示下的诱导联络。正有限秩复向量丛与 \(\operatorname{GL}(r,\mathbb C)\) 主标架丛有同样结论，其中线性联络取复线性。
\end{proposition}
\begin{proof}
	\cref{g1:prop:frame-associated-equivalence}给出 \(\operatorname{Fr}(E)\times_{\operatorname{GL}(r)}\mathbb F^r\cong E\)，映射为 \([u,v]\mapsto u(v)\)，其中 \(\mathbb F=\mathbb R\) 或 \(\mathbb C\)。因此主联络通过标准表示诱导 \(E\) 上的联络。

	反向，给定 \(E\) 上的联络 \(\nabla\)，在局部行标架 \(e_i\) 中取其联络矩阵 \(A_i\)。\(e_i\) 同时是 \(\operatorname{Fr}(E)\) 的局部截面，且 \(e_j=e_i g_{ij}\)。\cref{g1:prop:connection-local-data}已证明
	\[
		A_j=g_{ij}^{-1}A_i g_{ij}+g_{ij}^{-1}\mathrm{d}g_{ij}.
	\]
	复情形逐分量应用同一乘积法则，所得公式不变。\(\mathfrak{gl}(r,\mathbb F)\) 就是所有 \(r\times r\) 的 \(\mathbb F\) 矩阵，所以这些 \(A_i\) 均是合格的 Lie 代数值局部数据。由\cref{g1:thm:principal-connection-local-data}，它们唯一地构造出主联络。标准表示的微分是矩阵上的恒等映射，故再诱导回 \(E\) 后，局部矩阵仍为原来的 \(A_i\)，恢复 \(\nabla\)。反过来，从主联络得到 \(A_i\) 后再按局部公式重建，也由唯一性恢复原主联络。
\end{proof}

\begin{example}\label{g1:exm:adjoint-loses-central-connection}
	取任意主 \(\operatorname{U}(1)\) 丛。圆周群交换，所以 \(\operatorname{Ad}\) 是平凡表示，\(\operatorname{ad}P\) 自然等于 \(M\times i\mathbb R\)。无论原主联络是什么，诱导在伴随丛上的联络都是普通微分 \(\mathrm{d}\)，因为 \([A_i,\xi_i]=0\)。例如在 \(\mathbb R^2\times\operatorname{U}(1)\) 上，局部形式 \(A_0=0\) 与 \(A_1=i x\,\mathrm{d}y\) 确定不同主联络，却有完全相同的伴随丛联络。因此“主联络诱导伴随丛联络”并不意味着仅凭后者总能恢复前者。
\end{example}

一般表示中，两个主联络的差 \(B\in\Omega^1(M;\operatorname{ad}P)\) 在关联丛上变成 \(\rho_*(B)\)。若 \(\rho_*\) 单射，这些差仍可被区分；若 \(\rho_*\) 有核，就可能丢失方向。即使 \(\rho_*\) 单射，也不能说关联丛上的每个联络都来自原结构群：其局部矩阵还必须落在 \(\rho_*(\mathfrak g)\) 中。标架丛的标准表示之所以给出全体联络的双向对应，正是因为这个像等于所有矩阵。

% !TeX root = ../../Book1.tex
% 纲要标题：### 18.3 水平提升与曲率
% 纲要位置：几何学计划纲要.md，第 825 行。
% 正文状态：初稿；尚未正式审核。

% 纲要内容开始
% 
% * 路径提升及完备性范围
% * 曲率结构方程
% * Bianchi 恒等式的主丛形式
% * 小环路和平行移动的变化
% 
% 纲要内容结束

\section{水平提升与曲率}
\label{g1:sec:1803}

水平空间规定了怎样离开一个纤维而不额外作群方向的运动。沿底空间的一条路径不断采用这个方向，就应当得到两个端点纤维之间的比较。这里有两个需要分别回答的问题：所给路径能否完整提升，以及沿不同路径得到的比较为何会不同。前者由主作用的等变性控制，后者由曲率控制。

\ProseParagraphBreak

本节设 \(\pi:P\to M\) 为有限维实 Lie 群 \(G\) 的光滑右主丛，流形均有限维、Hausdorff、第二可数且无边界，\(\omega\in\Omega^1(P;\mathfrak g)\) 为主联络形式，\(H=\ker\omega\)。沿用上一节的约定：\(X^\#_p=\left.\frac{\mathrm{d}}{\mathrm{d}t}\right|_0p\exp(tX)\)、\(\omega(X^\#)=X\)，且 \(R_g^*\omega=\operatorname{Ad}_{g^{-1}}\omega\)。若 \(s:U\to P\) 为局部截面，则 \(A=s^*\omega\)，并以 \(p=s(x)g\) 表示局部纤维坐标。

\subsection{路径提升及完备性范围}
\label{g1:subsec:1803:01}

设 \(\gamma:[a,b]\to U\) 光滑，将它的一条提升写成 \(q(t)=s(\gamma(t))g(t)\)。由\eqref{g1:eq:principal-connection-trivialization}，水平条件正好是
\[
	0=\omega(q')=\operatorname{Ad}_{g^{-1}}A(\gamma')+\theta^L_g(g'),
\]
其中 \(\theta^L_g=(\mathrm{d}L_{g^{-1}})_g\)。把这个等式左平移回 \(\mathrm{T}_gG\)，得到
\begin{equation}\label{g1:eq:principal-horizontal-ode}
	g'(t)=-(\mathrm{d}R_{g(t)})_eB(t),\qquad B(t)=A_{\gamma(t)}(\gamma'(t)).
\end{equation}
矩阵群中它写成 \(g'=-Bg\)。右边是随时间变化的右不变向量场；“右不变”指对每个固定时间作右平移后不变，并不是说系数 \(B(t)\) 与时间无关。这个特征使我们能够在非紧 Lie 群上排除有限时间爆破。

\begin{lemma}[右不变方程的全区间解]\label{g1:lem:right-invariant-time-ode}
	设 \(G\) 为有限维实 Lie 群，\(\mathfrak g=\operatorname{Lie}(G)\)，\(B:[a,b]\to\mathfrak g\) 光滑，端点处按局部光滑延拓理解。对任意 \(t_0\in[a,b]\) 和 \(g_0\in G\)，方程
	\[
		g'(t)=-(\mathrm{d}R_{g(t)})_eB(t),\qquad g(t_0)=g_0
	\]
	（即\eqref{g1:eq:principal-horizontal-ode}）在整个 \([a,b]\) 上存在唯一解。解光滑依赖初始时刻、初值及 \(B\) 所含的有限维光滑参数。
\end{lemma}
\begin{proof}
	\cref{g1:thm:ode-local-existence-uniqueness}在单位元的一张坐标图中给出局部解。对每个 \(\tau\in[a,b]\)，该定理还给出初始时刻的一个邻域，以及从单位元出发都可使用的一段共同时间。用有限多个这样的邻域覆盖紧区间 \([a,b]\)，取各共同时间的一个更小正下界 \(\delta\)。于是从任意 \(\tau\in[a,b]\) 和单位元出发，都能至少求解到 \(|t-\tau|<\delta\) 与给定时间区间的交；\(\delta\) 不依赖 \(\tau\)。

	若 \(u(t)\) 是从单位元出发的这条短时间解，则 \(u(t)g_0\) 满足同一方程，因为
	\[
		\frac{\mathrm{d}}{\mathrm{d}t}(u g_0)
		=-(\mathrm{d}R_{g_0})_u(\mathrm{d}R_u)_eB
		=-(\mathrm{d}R_{u g_0})_eB.
	\]
	因此相同的时间 \(\delta\) 适用于任意 \(g_0\)，即使这些初值不在同一个紧集中。把 \([a,b]\) 分成长度小于 \(\delta/2\) 的有限小段，用前一段的终值作为后一段的初值，便可向前、向后走完整个区间。局部唯一性使重叠部分一致，也证明所得全区间解唯一。

	参数依赖由\cref{g1:thm:ode-smooth-dependence}及上述有限次拼接得到：固定一组参数后，每一步在一个共同局部解邻域内光滑依赖起始数据；缩小参数邻域，有限步的中间值仍落在相应邻域中。有限次复合保持光滑性。初始时刻和初值也包括在该定理的光滑变量内。
\end{proof}

每一步都能把初值右平移回同一个单位元附近，因此同一步长可以重复使用，直至走完整个紧参数区间。这个论证适用于一般的 \(G\) 与 \(M\)。

\begin{theorem}[主丛的水平提升]\label{g1:thm:principal-horizontal-lift}
	设 \(P\to M\) 为有限维实 Lie 群的光滑右主丛，\(M\) 无边界，\(\omega\) 为主联络。对任意连续、有限分片光滑路径 \(\gamma:[a,b]\to M\)、任意 \(t_0\in[a,b]\) 和 \(p_0\in P_{\gamma(t_0)}\)，存在唯一连续、分片光滑提升 \(q:[a,b]\to P\)，使下列等式在每个光滑片上成立：
	\[
		\pi\circ q=\gamma,\qquad q(t_0)=p_0,\qquad \omega(q')=0
	\]
	称这条提升为具有该初值的\term[shuipingtisheng]{水平提升}[horizontal lift]。若路径及初值光滑依赖有限维参数，并具有共同有限光滑分割，则提升在各片内都光滑依赖这些参数，端点运输在端点丛图中的表达也光滑依赖这些参数。
\end{theorem}
\begin{proof}
	先将光滑路径的紧像用有限个主丛平凡化域覆盖。时间区间的 Lebesgue 数性质给出有限分割，使每段闭区间的像包含在一个平凡化域内。在每段解\eqref{g1:eq:principal-horizontal-ode}，前一引理保证这一段可以完整走完。接续时将实际终点写入下一张丛图，作为新的初值。

	局部公式是内在条件 \(\omega(q')=0\) 的等价表达，因此换丛图不会改变所求对象。若两条提升具有同一初值，局部方程的唯一性先使它们在初始段相同，再依次使全部相邻段相同。这同时证明细分与丛图的选择无关。分片光滑路径只需把它的有限分段点加入上述分割；拐角处保留连续性，各片分别满足水平条件。

	对参数族，先在一个固定参数处选择有限分割，并让各段紧像留在所选丛图的内部。连续性与时间区间的紧性保证附近参数可以沿用同一批丛图和分割。局部方程的系数、各段末端的坐标变换与初值都光滑依赖参数，前一引理的参数结论及有限次复合遂给出所述光滑依赖。
\end{proof}

\begin{proposition}[主丛运输的运算规律]\label{g1:prop:principal-transport-laws}
	设 \(G\) 为有限维实 Lie 群，\(\pi:P\to M\) 为光滑主右 \(G\) 丛，\(M\) 无边界，\(\omega\) 为主联络。设 \(\gamma:[a,b]\to M\)、\(\eta:[c,d]\to M\) 均连续且有限分片光滑，并满足 \(\gamma(b)=\eta(c)\)。按\cref{g1:thm:principal-horizontal-lift}，以 \(\omega\) 水平提升的终点定义
	\[
		T_\gamma:P_{\gamma(a)}\longrightarrow P_{\gamma(b)}.
	\]
	它是 \(G\) 等变微分同胚，满足 \(T_\gamma(p g)=T_\gamma(p)g\)。以 \(\gamma*\eta\) 表示先走 \(\gamma\)、后走 \(\eta\) 的拼接，以 \(\bar\gamma(t)=\gamma(a+b-t)\) 表示反向路径。常值路径给出恒等映射，且有 \(T_{\gamma*\eta}=T_\eta\circ T_\gamma\)，而 \(T_{\bar\gamma}=T_\gamma^{-1}\)。若 \(h:[a',b']\to[a,b]\) 光滑非减，且 \(h(a')=a\)、\(h(b')=b\)，则 \(T_{\gamma\circ h}=T_\gamma\)，允许速度为零的停顿。
\end{proposition}
\begin{proof}
	若 \(q\) 水平，则 \(R_g\circ q\) 也水平，因为 \(R_g^*\omega=\operatorname{Ad}_{g^{-1}}\omega\)。其初值是 \(q(a)g\)，唯一性给出等变性。端点映射对初值光滑，由反向路径得到的逆映射也光滑。常值、拼接和反向的提升分别是常值提升、逐段连接的提升与反向提升；检验初值后使用唯一性即得前三条运算规律。最后，若时间换元为 \(t=h(\tau)\)，则 \(\omega((q\circ h)')=h'\omega(q')=0\)，没有除以 \(h'\)，所以停顿也不影响结论。
\end{proof}

对表示 \(\rho:G\to\operatorname{GL}(V)\)，关联丛中的向量 \([p,v]\) 沿 \(\gamma\) 变成 \([T_\gamma(p),v]\)。这正是\cref{g1:thm:associated-covariant-derivative}所构造联络的平行移动：在丛图中，\(q=s(\gamma)g\) 给出的向量分量为 \(\rho(g)v\)，对\eqref{g1:eq:principal-horizontal-ode}应用 \(\rho\) 后，其导数为 \(-\rho_*(A(\gamma'))\rho(g)v\)，恰好满足\eqref{g1:eq:parallel-transport-ode}。

\ProseParagraphBreak

若 \(\gamma\) 是以 \(x\) 为基点的环路，选定 \(p\in P_x\) 后唯一写成 \(T_\gamma(p)=p h_\gamma\)。运输与右作用等变，所以
\begin{equation}\label{g1:eq:principal-loop-composition}
	h_{\gamma*\eta}=h_\eta h_\gamma.
\end{equation}
在下一节将运输写成基本群表示时，这个次序正是需要取逆的原因。

\begin{example}[一般水平分布的提升可能中断]\label{g1:exm:incomplete-ehresmann-lift}
	在投影 \(\mathbb R^2\to\mathbb R\)、\((x,y)\mapsto x\) 上，令
	\[
		H_{(x,y)}=\operatorname{span}\{\partial_x+y^2\partial_y\}.
	\]
	它光滑且与竖直空间 \(\operatorname{span}\{\partial_y\}\) 互补，因而是一般纤维丛意义的水平分布。路径 \(x(t)=t\)、\(0\le t\le2\)，从 \((0,1)\) 出发的水平提升却须满足 \(y'=y^2\)，故 \(y(t)=(1-t)^{-1}\)，在 \(t=1\) 中断。纤维上的平移 \(y\mapsto y+c\) 不保持这个分布，所以它不是加法群主作用的主联络。由此可见，局部互补条件本身不能推出任意给定路径的完整提升；上述定理额外使用了主作用等变性。
\end{example}

\subsection{曲率结构方程}
\label{g1:subsec:1803:02}

若先沿一个水平方向、再沿另一个水平方向运动，与交换次序后的结果不同，差别可能含竖直方向。微分形式会把这种竖直偏差记录下来。对两个水平向量场 \(U,V\)，\eqref{g1:eq:exterior-one-form-evaluation}已经给出 \(\mathrm{d}\omega(U,V)=-\omega([U,V])\)。不过 \(\mathrm{d}\omega\) 在任意切向输入上未必只记录水平变化，还须补上一个 Lie 括号项。

\ProseParagraphBreak

为明确这个项的系数，选 \(\mathfrak g\) 的任一基 \(E_a\)，对 \(\mathfrak g\) 值形式 \(\alpha=\sum_a\alpha^aE_a\)、\(\beta=\sum_b\beta^bE_b\) 定义
\begin{equation}\label{g1:eq:lie-bracket-wedge}
	[\alpha\wedge\beta]\coloneqq\sum_{a,b}\alpha^a\wedge\beta^b[E_a,E_b].
\end{equation}
右边只是把双线性 Lie 括号作用于系数，因此与所选基无关。特别地，对一形式 \(\omega\)，有 \([\omega\wedge\omega](U,V)=2[\omega(U),\omega(V)]\)。这解释了下一定义中的 \(1/2\)。

\begin{definition}[主联络的曲率]\label{g1:def:principal-curvature}
	设 \(\omega\) 为光滑右主 \(G\) 丛 \(P\to M\) 上的联络形式。其\term[zhulianluoqulv]{曲率形式}[curvature form]定义为
	\begin{equation}\label{g1:eq:principal-curvature-structure}
		\Omega\coloneqq\mathrm{d}\omega+\tfrac12[\omega\wedge\omega]\in\Omega^2(P;\mathfrak g).
	\end{equation}
	若 \(\Omega=0\)，则称该主联络平坦。若 \(G\) 是矩阵 Lie 群，括号是矩阵交换子，上式也写成 \(\Omega=\mathrm{d}\omega+\omega\wedge\omega\)，此处后一个楔积同时使用矩阵乘法。
\end{definition}

\begin{proposition}[曲率的水平性与局部表达]\label{g1:prop:principal-curvature-properties}
	设 \(P\to M\) 为有限维实 Lie 群的光滑右主丛，\(\omega\) 为主联络，曲率为 \(\Omega\)。则
	\begin{enumerate}
		\item \(\Omega\) 是水平的，即任意一个输入为竖直向量时其值为零；并且 \(R_g^*\Omega=\operatorname{Ad}_{g^{-1}}\Omega\)。
		\item 对水平向量场 \(U,V\)，有 \(\Omega(U,V)=-\omega([U,V])\)。
		\item 若 \(s_i:U_i\to P\) 为局部截面，\(A_i=s_i^*\omega\)，则
		\begin{equation}\label{g1:eq:principal-local-curvature}
			F_i\coloneqq s_i^*\Omega=\mathrm{d}A_i+\tfrac12[A_i\wedge A_i],\qquad
			F_j=\operatorname{Ad}_{g_{ij}^{-1}}F_i\quad(s_j=s_i g_{ij}).
		\end{equation}
		\item 对光滑有限维线性表示 \(\rho:G\to\operatorname{GL}(V)\)，关联丛联络的局部曲率矩阵为 \(\rho_*(F_i)\)。
	\end{enumerate}
\end{proposition}
\begin{proof}
	对固定 \(X\in\mathfrak g\)，基本场 \(X^\#\) 的流为 \(R_{\exp(tX)}\)。对联络的等变公式求导，得到
	\[
		\mathcal L_{X^\#}\omega=-[X,\omega].
	\]
	\cref{g1:thm:cartan-formula}逐分量适用，而 \(\iota_{X^\#}\omega=X\) 为常值，故 \(\iota_{X^\#}\mathrm{d}\omega=-[X,\omega]\)。另一方面，\(\tfrac12\iota_{X^\#}[\omega\wedge\omega]=[X,\omega]\)。两项相消，得到 \(\iota_{X^\#}\Omega=0\)。每个竖直向量都是某个 \(X^\#\) 的点值，因此 \(\Omega\) 水平。固定群元素的伴随作用与外微分交换，又保持 Lie 括号；对结构方程取 \(R_g^*\) 便得到等变性。

	对水平 \(U,V\)，括号楔项为零，外微分的一形式公式化为 \(\Omega(U,V)=-\omega([U,V])\)，说明曲率由水平场括号的竖直分量确定。

	拉回与外微分交换，见\cref{g1:prop:exterior-pullback}，并且拉回保持\eqref{g1:eq:lie-bracket-wedge}，故 \(F_i\) 满足所列结构方程。计算 \(s_j=s_i g_{ij}\) 的微分时，一项是 \(\mathrm{d}R_{g_{ij}}\mathrm{d}s_i\)，另一项来自 \(g_{ij}\) 的变化，属于竖直空间。曲率的水平性消去后一项，等变性则将前一项变成 \(\operatorname{Ad}_{g_{ij}^{-1}}F_i\)。这正是曲率变换式中没有 \(\mathrm{d}g_{ij}\) 的原因。

	最后，\cref{g1:thm:associated-covariant-derivative}给出关联联络矩阵 \(\rho_*(A_i)\)。由\cref{g1:thm:curvature-structure-equation}，其曲率矩阵为
	\[
		\mathrm{d}(\rho_*A_i)+(\rho_*A_i)\wedge(\rho_*A_i)
		=\rho_*\bigl(\mathrm{d}A_i+\tfrac12[A_i\wedge A_i]\bigr),
	\]
	其中用了\cref{g1:prop:lie-homomorphism-differential}中 Lie 群同态的微分保持 Lie 括号这一结论。因此它等于 \(\rho_*(F_i)\)。
\end{proof}

曲率因而也可看作底空间上取值于伴随丛 \(\operatorname{ad}P=P\times_{\operatorname{Ad}}\mathfrak g\) 的二形式：在 \(x\) 处取 \(p\in P_x\)，把底空间的输入水平提升至 \(p\)，再取 \([p,\Omega_p(\cdot,\cdot)]\)。改取 \(pg\) 时，输入由 \(\mathrm{d}R_g\) 搬运，曲率值变为 \(\operatorname{Ad}_{g^{-1}}\Omega_p\)，恰好被关联丛的等价关系抵消。因此这个底空间对象无需选择局部截面；\(F_i\) 是它在所选截面中的表示。

\begin{example}[曲率与水平括号]\label{g1:exm:principal-additive-curvature}
	在 \(P=\mathbb R^2\times\mathbb R\to\mathbb R^2\) 上，让加法群 \(\mathbb R\) 作用于最后一个坐标 \(z\)，取 \(\omega=\mathrm{d}z+x\,\mathrm{d}y\)。这是主联络：它对 \(\partial_z\) 取值为 \(1\)，且在 \(z\) 平移下不变。水平分布由
	\[
		U=\partial_x,\qquad V=\partial_y-x\partial_z
	\]
	张成，而 \([U,V]=-\partial_z\)、\(\Omega=\mathrm{d}x\wedge\mathrm{d}y\)。于是 \(\Omega(U,V)=1=-\omega([U,V])\)。沿正向矩形边界提升时，\(z'=-x y'\)，走完后 \(z\) 减少矩形面积。这一竖直位移就是此处的环路返回，未必是角度或旋转。
\end{example}

\begin{theorem}[平坦联络的局部水平截面]\label{g1:thm:flat-principal-local-section}
	设 \(\pi:P\to M\) 为有限维实 Lie 群 \(G\) 的光滑主右丛，流形无边界，\(\omega\) 为主联络，\(\Omega\) 为其曲率；对任一局部截面 \(s\) 记 \(A_s=s^*\omega\)。以下条件等价：
	\begin{enumerate}
		\item \(\Omega=0\)。
		\item 水平分布 \(H=\ker\omega\) 对合。
		\item 每个 \(p\in P\) 都位于某个局部水平截面 \(s:U\to P\) 上，即 \(s(\pi(p))=p\) 且 \(s^*\omega=0\)。
	\end{enumerate}
	因此平坦主联络在足够小的截面中具有 \(A_s=0\) 的局部形式；这样的截面之间的过渡函数局部常值。
\end{theorem}
\begin{proof}
	曲率水平，故要判断它是否为零，只须代入两个水平输入。前一命题给出 \(\Omega(U,V)=-\omega([U,V])\)，所以它处处为零恰好等价于水平场的括号仍水平。这证明前两项等价。

	若 \(H\) 对合，\cref{g1:thm:frobenius}适用于总空间 \(P\) 上秩为 \(\dim M\) 的光滑分布，给出通过 \(p\) 的局部嵌入积分片 \(L\)。投影微分 \(\mathrm{d}\pi:H_p\to \mathrm{T}_{\pi(p)}M\) 是同构，所以\cref{g1:thm:inverse-function}使 \(\pi|_L\) 在 \(p\) 附近成为微分同胚。其局部逆 \(s\) 满足 \(\pi\circ s=\operatorname{id}\)，并且 \(\mathrm{d}s(\mathrm{T}U)\subset H\)，即为所求水平截面。反过来，若每个 \(p\) 都有这样的截面，那么 \(s^*\Omega=\mathrm{d}(s^*\omega)+\tfrac12[s^*\omega\wedge s^*\omega]=0\)。\(\mathrm{d}s\) 在该点张成全部水平空间，曲率的水平性遂说明 \(\Omega_p=0\)。

	若两个水平截面满足 \(s_j=s_i g_{ij}\)，联络变换式给出 \(0=\theta^L(g_{ij})\)，即 \((\mathrm{d}L_{g_{ij}^{-1}})\mathrm{d}g_{ij}=0\)。左平移微分可逆，所以 \(\mathrm{d}g_{ij}=0\)。光滑映射微分为零时，在一个充分小的连通坐标邻域中沿坐标线积分可知其值不变，故 \(g_{ij}\) 局部常值。
\end{proof}

这个结果是局部陈述。水平截面在绕行之后能否重新拼成同一个截面，还要看环路运输；下一节将保留这部分整体信息。

\subsection{Bianchi 恒等式的主丛形式}
\label{g1:subsec:1803:03}

曲率的局部表达在更换截面时按伴随表示变化。因此对曲率求导时，需要补上这一表示的联络项，才能将各局部结果粘成同一个对象。对水平且满足 \(R_g^*\eta=\operatorname{Ad}_{g^{-1}}\eta\) 的 \(\mathfrak g\) 值 \(k\) 形式，定义
\begin{equation}\label{g1:eq:principal-exterior-covariant-derivative}
	\mathrm{D}\eta\coloneqq\mathrm{d}\eta+[\omega\wedge\eta].
\end{equation}
这仍是水平等变形式。事实上，\(\mathcal L_{X^\#}\eta=-[X,\eta]\)，Cartan 公式与 \(\iota_{X^\#}\eta=0\) 给出 \(\iota_{X^\#}\mathrm{d}\eta=-[X,\eta]\)，而 \(\iota_{X^\#}[\omega\wedge\eta]=[X,\eta]\)，两项抵消。等变性由外微分及 Lie 括号与固定伴随变换相容得到。零次形式时插入项本来为零，同一计算也适用。局部截面中的公式为 \(s^*\mathrm{D}\eta=\mathrm{d}(s^*\eta)+[A\wedge s^*\eta]\)，正是伴随 Lie 代数丛联络的外协变微分。

\begin{theorem}[主联络的 Bianchi 恒等式]\label{g1:thm:principal-bianchi}
	设 \(\pi:P\to M\) 为有限维实 Lie 群 \(G\) 的光滑主右丛，\(\omega\) 为主联络，\(\Omega\) 为其曲率。对光滑局部截面 \(s:U\to P\)，记 \(A=s^*\omega\)、\(F=s^*\Omega\)。则
	\begin{equation}\label{g1:eq:principal-bianchi}
		\mathrm{D}\Omega=\mathrm{d}\Omega+[\omega\wedge\Omega]=0.
	\end{equation}
	在局部截面中，这成为 \(\mathrm{d}F+[A\wedge F]=0\)。矩阵群情形即 \(\mathrm{d}F+A\wedge F-F\wedge A=0\)。
\end{theorem}
\begin{proof}
	Lie 括号的 Jacobi 恒等式提供了所需的抵消。对次数为 \(p,q\) 的形式，由定义逐项交换标量楔积与 Lie 括号，得到
	\[
		[\alpha\wedge\beta]=-(-1)^{pq}[\beta\wedge\alpha],\qquad
		\mathrm{d}[\alpha\wedge\beta]=[\mathrm{d}\alpha\wedge\beta]+(-1)^p[\alpha\wedge\mathrm{d}\beta].
	\]
	特别地，\(\mathrm{d}[\omega\wedge\omega]=-2[\omega\wedge\mathrm{d}\omega]\)。另外，直接在三个切向输入 \(U,V,W\) 上展开有
	\begin{align*}
		[\omega\wedge[\omega\wedge\omega]](U,V,W)
		={}&2[\omega(U),[\omega(V),\omega(W)]]\\
		&+2[\omega(V),[\omega(W),\omega(U)]]\\
		&+2[\omega(W),[\omega(U),\omega(V)]]=0,
	\end{align*}
	最后一步就是 \(\mathfrak g\) 的 Jacobi 恒等式。因此结构方程给出
	\[
		\mathrm{d}\Omega=-[\omega\wedge\mathrm{d}\omega],\qquad
		[\omega\wedge\Omega]=[\omega\wedge\mathrm{d}\omega],
	\]
	相加即为零。拉回到局部截面得到局部公式。矩阵情形的一次形式与二次形式之间，括号楔积为 \([A\wedge F]=A\wedge F-F\wedge A\)，故与\eqref{g1:eq:bianchi-matrix}完全一致。
\end{proof}

这个恒等式适用于任意主丛联络。若 \(G\) 为交换群，括号项消失，恒等式退化为 \(\mathrm{d}\Omega=0\)；一般情形删去括号项就改变了所比较的对象。

\subsection{小环路和平行移动的变化}
\label{g1:subsec:1803:04}

在\cref{g1:prop:principal-curvature-properties}中，曲率记录了水平场括号的竖直部分。现在把它与实际运输联系起来：先计算一族路径变化时终点怎样变化，再说明一个小正向方环为何产生负的曲率面积项。两种计算都保持上一节的右主作用约定。

\begin{theorem}[水平运输的变分]\label{g1:thm:principal-transport-variation}
	设 \(G\) 为有限维实 Lie 群，\(\pi:P\to M\) 为带主联络 \(\omega\) 的光滑主右 \(G\) 丛，\(\Omega\) 为其曲率，\(\theta^L_g=(\mathrm{d}L_{g^{-1}})_g\) 为 \(G\) 的左 Maurer--Cartan 形式。设 \(\Gamma:I\times[a,b]\to M\) 为光滑路径族，\(I\subset\mathbb R\) 为开区间。设 \(c:I\to P\) 光滑，\(\pi(c(s))=\Gamma(s,a)\)。以 \(c(s)\) 为初值，沿 \(t\mapsto\Gamma(s,t)\) 的水平提升记为 \(q(s,t)\)。令
	\[
		W(s,t)=\omega_{q(s,t)}(\partial_s q).
	\]
	则
	\begin{equation}\label{g1:eq:principal-variation-intrinsic}
		W(s,b)=\omega_{c(s)}(c'(s))+\int_a^b\Omega_{q(s,t)}(\partial_tq,\partial_sq)\,\mathrm{d}t.
	\end{equation}
	若 \(\Gamma\) 的像在同一个截面 \(\sigma:U\to P\) 的定义域内，取 \(c(s)=\sigma(\Gamma(s,a))\)、\(q(s,t)=\sigma(\Gamma(s,t))u(s,t)\)，则 \(u(s,a)=e\)。记 \(A=\sigma^*\omega\)、\(F=\sigma^*\Omega\)、\(A_s(s,t)=A(\partial_s\Gamma)\)，有
	\begin{equation}\label{g1:eq:principal-variation-endpoints}
		\begin{split}
			\theta^L_{u(s,b)}(\partial_su(s,b))
			={}&A_s(s,a)-\operatorname{Ad}_{u(s,b)^{-1}}A_s(s,b)\\
			&+\int_a^b\operatorname{Ad}_{u(s,t)^{-1}}F_{\Gamma(s,t)}(\partial_t\Gamma,\partial_s\Gamma)\,\mathrm{d}t.
		\end{split}
	\end{equation}
\end{theorem}
\begin{proof}
	水平提升的参数依赖保证 \(q\) 光滑，因而可以对拉回形式求外微分。沿 \(t\) 方向水平意味着
	\[
		q^*\omega=W\,\mathrm{d}s.
	\]
	它与自身的括号楔积为零，所以曲率结构方程拉回后给出
	\[
		q^*\Omega=\mathrm{d}(W\,\mathrm{d}s)=\partial_tW\,\mathrm{d}t\wedge\mathrm{d}s.
	\]
	代入 \((\partial_t,\partial_s)\)，得到 \(\partial_tW=\Omega(q_t,q_s)\)。从 \(a\) 到 \(b\) 积分，并注意 \(q(s,a)=c(s)\)，便得第一式。积分中的值都在同一个有限维向量空间 \(\mathfrak g\) 内，按任一基逐分量积分即可。

	在所给局部截面中，联络平凡化公式给出
	\[
		W(s,t)=\operatorname{Ad}_{u(s,t)^{-1}}A_s(s,t)+\theta^L_{u(s,t)}(\partial_su(s,t)).
	\]
	曲率水平且等变，故 \(\Omega(q_t,q_s)=\operatorname{Ad}_{u^{-1}}F(\Gamma_t,\Gamma_s)\)。初端 \(u(s,a)=e\) 对 \(s\) 常值，因此 \(W(s,a)=A_s(s,a)\)。将这三项代入第一式，再把终端的 \(A_s\) 项移到右边，得到第二式。
\end{proof}

\input{Figures/Chapter18-TransportVariation.tex}

\cref{g1:fig:principal-transport-variation}把两种变化放在同一张参数图里：\(t\) 沿着路径前进，\(s\) 改变整条路径。公式中曲率的输入次序是 \((\Gamma_t,\Gamma_s)\)；改成 \((\Gamma_s,\Gamma_t)\) 就应补一个负号。若起点、终点随 \(s\) 移动，\eqref{g1:eq:principal-variation-endpoints}中的两个端点项不能省略。若两端固定，端点项消失，而积分中的 \(\operatorname{Ad}_{u^{-1}}\) 将不同位置的曲率放到提升所选的同一比较方式中；一般非交换群不能删去这一项。

\ProseParagraphBreak

特别地，对一族基点固定的环路，从同一个 \(p\in P_x\) 提升，写 \(q(s,b)=p h(s)\)，则
\[
	\theta^L_{h(s)}(h'(s))=\int_a^b\Omega(q_t,q_s)\,\mathrm{d}t.
\]
平坦时右侧为零，因而 \(h'(s)=0\)。这已经说明光滑定端点同伦不改变平坦运输；下一节将用局部水平截面把运输延伸到连续路径，并证明基本群所采用的连续同伦也保持它。

\begin{proposition}[小方环的曲率项]\label{g1:prop:principal-small-loop}
	设 \(\pi:P\to M\) 为光滑主右 \(G\) 丛，\(G\) 为有限维实 Lie 群，\(\mathfrak g=\operatorname{Lie}(G)\)，\(\omega\) 为给定主联络，\(\Omega\) 为其曲率。以 \(T_\gamma\) 表示 \(\omega\) 的水平运输。取内点 \(x_0\) 附近的坐标，使 \(x(x_0)=0\)，并取局部截面 \(\sigma\)，令 \(p=\sigma(x_0)\)、\(F=\sigma^*\Omega\)。假设 \(\dim M\ge2\)，并记 \(F_{12}(x_0)=F_{x_0}(\partial/\partial x^1,\partial/\partial x^2)\)。令 \(\gamma_\varepsilon\) 从原点沿第一坐标正向、第二坐标正向、第一坐标反向、第二坐标反向绕行边长 \(\varepsilon>0\) 的方环，其余坐标为零。写
	\[
		T_{\gamma_\varepsilon}(p)=p h_\varepsilon.
	\]
	当 \(\varepsilon\) 充分小时，\(h_\varepsilon\) 位于指数映射的局部可逆邻域，且
	\begin{equation}\label{g1:eq:principal-small-loop}
		\log h_\varepsilon=-\varepsilon^2F_{12}(x_0)+O(\varepsilon^3).
	\end{equation}
	余项用 \(\mathfrak g\) 上任一固定范数衡量。对矩阵 Lie 群，等价地有 \(h_\varepsilon=I-\varepsilon^2F_{12}(x_0)+O(\varepsilon^3)\)。
\end{proposition}
\begin{proof}
	先选一个在基点处一阶变化恰好抵消 \(A=\sigma^*\omega\) 的截面变换。令
	\[
		k(x)=\exp\left(-\sum_i x^iA_i(x_0)\right),\qquad \widetilde\sigma=\sigma k.
	\]
	由\cref{g1:prop:exponential-local-diffeomorphism}，\((\mathrm{d}\exp)_0=\operatorname{id}\)，且指数映射在零点附近有光滑逆。因为 \(k(x_0)=e\)，联络的局部变换公式给出 \(\widetilde A_i(x_0)=0\)。新旧截面在基点相同，故 \(p,h_\varepsilon\) 不变；曲率的等变变换又保证 \(\widetilde F(x_0)=F(x_0)\)。以下在这个新截面中计算，省去波浪号。

	把四条边都用坐标长度参数化，总时间为 \(4\varepsilon\)。沿全部边，系数 \(B(t)=A(\gamma_\varepsilon'(t))\) 分片光滑且 \(\|B(t)\|\le C\varepsilon\)，因为 \(A(x_0)=0\)。选单位元附近的指数坐标 \(\xi=\log g\)。在这张图中，右不变方程成为
	\[
		\xi'=-K(\xi)B(t),\qquad
		K(\xi)=(\mathrm{d}\log)_{\exp\xi}\circ(\mathrm{d}R_{\exp\xi})_e,
	\]
	其中 \(K\) 光滑且 \(K(0)=I\)。在固定小闭球中，\(K\) 有界且 \(K(\xi)=I+O(\|\xi\|)\)。从 \(\xi(0)=0\) 出发，在解尚未离开这个球时，积分估计给出 \(\|\xi(t)\|\le C'\varepsilon t\le4C'\varepsilon^2\)。若 \(\varepsilon\) 足够小，该界严格小于球半径，所以解不可能首次碰到球边界。于是全程 \(\xi=O(\varepsilon^2)\)，特别 \(h_\varepsilon\) 确实在 \(\log\) 的定义域内。

	将这个界代回微分方程，\(K(\xi)-I\) 对积分的贡献至多为 \(O(\varepsilon^2)O(\varepsilon)O(\varepsilon)=O(\varepsilon^4)\)，因此
	\[
		\log h_\varepsilon=-\oint_{\gamma_\varepsilon}A+O(\varepsilon^4).
	\]
	对 \(\mathfrak g\) 的各分量应用矩形 Stokes 公式\cref{g1:thm:stokes-corners}，再作一阶余项估计，得到
	\begin{align*}
		\oint_{\gamma_\varepsilon}A
		&=\int_0^\varepsilon\int_0^\varepsilon(\partial_1A_2-\partial_2A_1)(s,t,0,\ldots,0)\,\mathrm{d}s\,\mathrm{d}t\\
		&=\varepsilon^2F_{12}(x_0)+O(\varepsilon^3).
	\end{align*}
	最后一步用到了 \(A(x_0)=0\)，所以基点的交换子项消失。两式合并证明一般 Lie 群的结论。矩阵情形再用 \(\exp Z=I+Z+O(\|Z\|^2)\) 即得矩阵表达；前面的论证本身没有要求 \(G\) 具有忠实矩阵表示。
\end{proof}

这与\eqref{g1:eq:small-loop-curvature}的向量丛公式同号：负号来自水平方程 \(g'=-A(\gamma')g\) 及规定的正向绕行。曲率为零意味着所有这些面积首项消失，但单个环路恰好返回原点不意味着邻近曲率为零；而一般结构群非交换的情形也不能把面积首项直接改写成任意大小曲面的精确指数积分。

% 纲要条目排为编号小节标题，并进入目录与书签；说明段仍保留于纲要注释中。

% !TeX root = ../../Book1.tex
% 纲要标题：### 18.4 规范群与平坦模问题
% 纲要位置：几何学计划纲要.md，第 832 行。
% 正文状态：初稿；尚未正式审核。

% 纲要内容开始
% 
% * 规范变换及局部公式
% * 规范轨道和稳定子
% * 平坦丛的单值构造
% * 表示空间模共轭与平坦结构的对应范围
% * 圆周、环面和 Hopf 丛的对照
% 
% ---
% 
% 纲要内容结束

\section{规范群与平坦模问题}
\label{g1:sec:1804}

局部联络形式会随着截面变化，但联络还可以被一个真正作用在总空间上的自同构改变。怎样判断两份局部公式只是同一几何对象的不同表示，还是经由自同构成为等价的联络？平坦情形中，曲率已经全部消失，这个问题又留下什么可比较的数据？本节先分清这些操作，再用环路运输组织平坦结构的整体分类。

\subsection{规范变换及局部公式}
\label{g1:subsec:1804:01}

固定右主 \(G\) 丛 \(\pi:P\to M\)。这里 \(G\) 为有限维实 Lie 群，\(P,M\) 为有限维、无边界、Hausdorff、第二可数光滑流形。底点表示实际位置，纤维则容纳同一位置的不同标架或内部选择。保持底点、同时与右群作用相容的自同构，是本节要研究的变换。

\begin{definition}\label{g1:def:gauge-group}
	设 \(G\) 为有限维实 Lie 群，\(\pi:P\to M\) 为光滑主右 \(G\) 丛，\(P,M\) 无边界、Hausdorff、第二可数。若光滑微分同胚 \(\Phi:P\to P\) 满足 \(\pi\circ\Phi=\pi\) 和 \(\Phi(pg)=\Phi(p)g\)，则称其为 \(P\) 的一个\term[guifanbianhuan]{规范变换}[gauge transformation]。所有规范变换按函数复合组成\term[guifanqun]{规范群}[gauge group]，记为 \(\mathcal G(P)\)。\symidx{\(\mathcal G(P)\)：主丛的规范群}
\end{definition}

\begin{proposition}[规范变换的表示]\label{g1:prop:gauge-equivariant-functions}
	设 \(P\to M\) 为光滑右主 \(G\) 丛，\(G\) 为有限维实 Lie 群。规范变换与光滑映射 \(u:P\to G\) 一一对应，其对应式及条件为
	\begin{equation}\label{g1:eq:gauge-equivariant-function}
		\Phi_u(p)=p\,u(p),\qquad u(pg)=g^{-1}u(p)g.
	\end{equation}
	若局部截面满足 \(s_j=s_i g_{ij}\)，记 \(u_i=u\circ s_i\)，则
	\[
		u_j=g_{ij}^{-1}u_i g_{ij},\qquad
		\Phi_u(s_i(x)g)=s_i(x)u_i(x)g.
	\]
	在这些函数下，复合对应逐点乘法：\(\Phi_u\circ\Phi_v=\Phi_{uv}\)。
\end{proposition}
\begin{proof}
	每根纤维上的右作用自由且传递，所以 \(\Phi(p)\) 唯一写成 \(p u(p)\)。在局部乘积坐标 \(p=s_i(x)g\) 中，\(\Phi(s_i(x))=s_i(x)u_i(x)\) 给出光滑 \(u_i\)，等变性进一步给出 \(u(s_i(x)g)=g^{-1}u_i(x)g\)。故 \(u\) 光滑且满足所写条件。反之，满足该条件的光滑 \(u\) 给出覆盖恒等的等变光滑映射，其局部表达 \((x,g)\mapsto(x,u_i(x)g)\) 有光滑逆 \((x,g)\mapsto(x,u_i(x)^{-1}g)\)，因而为规范变换。

	令 \(p=s_i(x)\)、\(g=g_{ij}(x)\)，立即得到交叠公式。最后，
	\[
		\Phi_u(\Phi_v(p))=p\,v(p)\,u(pv(p))
		=p\,v(p)\bigl(v(p)^{-1}u(p)v(p)\bigr)=p\,u(p)v(p),
	\]
	证明复合规律。注意中间必须使用 \(u\) 的等变性，不能把 \(u(pv(p))\) 直接写成 \(u(p)\)。
\end{proof}

对共轭左作用 \(g\cdot a=gag^{-1}\)，\cref{g1:prop:associated-fiber-bundle}给出群纤维丛 \(\operatorname{Ad}(P)=P\times_GG\)。式\eqref{g1:eq:gauge-equivariant-function}说明规范群也可识别为这个丛的光滑截面群。大写 \(\operatorname{Ad}(P)\) 的纤维是群，小写 \(\operatorname{ad}(P)=P\times_G\mathfrak g\) 的纤维是 Lie 代数；联络之差取值于后者，不要把两个对象混为一谈。

\begin{proposition}[联络的规范变换]\label{g1:prop:gauge-connection-formula}
	设 \(G\) 为有限维实 Lie 群，\(\pi:P\to M\) 为光滑主右 \(G\) 丛，\(\omega\) 为主联络，\(\Omega\) 为其曲率。令 \(u:P\to G\) 光滑且满足 \(u(pg)=g^{-1}u(p)g\)，并定义规范变换 \(\Phi_u(p)=pu(p)\)。则 \(\Phi_u^*\omega\) 仍为主联络，并有
	\begin{equation}\label{g1:eq:gauge-global-connection}
		\Phi_u^*\omega=\operatorname{Ad}_{u^{-1}}\omega+u^*\theta^L,
	\end{equation}
	其中 \(\theta^L\) 是 \(G\) 的左 Maurer--Cartan 形式。取固定局部截面 \(s_i\)，令 \(u_i=u\circ s_i\)、\(A_i=s_i^*\omega\)、\(F_i=s_i^*\Omega\)，则变换后的局部形式和曲率满足
	\begin{equation}\label{g1:eq:gauge-local-connection}
		A_i^u=\operatorname{Ad}_{u_i^{-1}}A_i+u_i^*\theta^L,
		\qquad F_i^u=\operatorname{Ad}_{u_i^{-1}}F_i.
	\end{equation}
	矩阵群时，第一式为 \(A_i^u=u_i^{-1}A_i u_i+u_i^{-1}\mathrm{d}u_i\)。
\end{proposition}
\begin{proof}
	\(\Phi_u\) 与右作用交换，并把每个基本场 \(X^\#\) 推到同一个基本场，所以拉回保持联络形式的等变性和再现条件。由\cref{g1:thm:principal-connection-equivalence}，拉回确实是主联络。

	对过 \(p\) 的曲线 \(p(t)\)，微分曲线 \(p(t)u(p(t))\)。由乘积法则，其速度分成把 \(p'(0)\) 右移的部分，以及由 \(u(p(t))\) 的变化产生的竖直部分。联络在前者上取值 \(\operatorname{Ad}_{u(p)^{-1}}\omega_p(p'(0))\)，在后者上取值 \(\theta^L_{u(p)}((\mathrm{d}u)_p p'(0))\)。这正是\eqref{g1:eq:gauge-global-connection}；它也就是式\eqref{g1:eq:principal-connection-trivialization}中允许截面和群参数同时变化的计算。再沿 \(s_i\) 拉回，得到局部联络公式。

	曲率的定义与拉回交换，故新曲率为 \(\Phi_u^*\Omega\)。原曲率水平，刚才速度分解中的竖直部分对曲率求值没有贡献。保留下来的右移部分由等变性给出 \(\Phi_u^*\Omega=\operatorname{Ad}_{u^{-1}}\Omega\)，再沿截面拉回即得曲率公式。
\end{proof}

现在可以比较三种容易混淆的操作。被动地把截面 \(s_i\) 换成 \(s_i k_i\)，联络 \(\omega\) 不变，局部表达却按同一个公式变化；主动地施行 \(\Phi_u\)，则保持所选 \(s_i\) 不变，把联络改为 \(\Phi_u^*\omega\)。两者公式相同，是因为 \(\Phi_u\circ s_i=s_i u_i\)，但各片任取的换截面函数 \(k_i\) 不必满足 \(u_j=g_{ij}^{-1}u_i g_{ij}\)，所以并不自动拼成一个整体规范变换。

\ProseParagraphBreak

底空间坐标变化又是第三件事。若 \(x=x(y)\)，同一个一形式 \(A=\sum_iA_i\,\mathrm{d}x^i\) 在新坐标下写成 \(A'_a=\sum_iA_i\,\partial x^i/\partial y^a\)，只使用协变指标的 Jacobian 变换。若真正将底点按微分同胚 \(f:M\to M\) 移动，首先得到的是拉回丛 \(f^*P\) 及其拉回联络；还须给定 \(P\) 与 \(f^*P\) 的丛同构才能在同一个总空间上比较。这不属于本节覆盖 \(\operatorname{id}_M\) 的规范群。

\subsection{规范轨道和稳定子}
\label{g1:subsec:1804:02}

以 \(\mathcal A(P)\) 记 \(P\) 上全部光滑主联络。它由\cref{g1:prop:principal-connection-affine}成为以 \(\Omega^1(M;\operatorname{ad}(P))\) 为平移空间的仿射空间。令 \(\omega\cdot\Phi=\Phi^*\omega\)，这给出右规范作用，因为
\[
	(\omega\cdot\Phi)\cdot\Psi=\Psi^*\Phi^*\omega
	=(\Phi\circ\Psi)^*\omega=\omega\cdot(\Phi\circ\Psi).
\]
同一轨道中的联络称为规范等价。我们此时只定义等价类集合 \(\mathcal A(P)/\mathcal G(P)\)，尚未给这个无穷维商建立流形结构。

\ProseParagraphBreak

给定基点 \(x_0\) 和纤维点 \(p_0\in P_{x_0}\)，把沿分段光滑路径 \(\gamma\) 的水平运输记为 \(T_\gamma\)。\cref{g1:thm:principal-horizontal-lift,g1:prop:principal-transport-laws}保证它在整个紧参数区间存在，且与右作用交换。对基于 \(x_0\) 的环路，有唯一 \(h_\gamma\in G\) 满足 \(T_\gamma(p_0)=p_0h_\gamma\)。所有这样的 \(h_\gamma\) 构成 \(G\) 的子群，记为 \(\operatorname{Hol}_{p_0}(\omega)\)，称为该点的\term[zhulianluoholonomyqun]{环路平行移动群}[holonomy group of a principal connection]。它未必闭，这里不预设其具有某种嵌入 Lie 群结构。

\begin{theorem}[规范稳定子与环路运输]\label{g1:thm:gauge-stabilizer-centralizer}
	设 \(G\) 为有限维实 Lie 群，\(\pi:P\to M\) 为光滑主右 \(G\) 丛，\(M\) 非空、连通、有限维、无边界、Hausdorff、第二可数，\(\omega\) 为光滑主联络。固定 \(x_0\in M\) 及 \(p_0\in P_{x_0}\)。把规范变换写成 \(\Phi_u(p)=pu(p)\)，其中 \(u:P\to G\) 光滑且 \(u(pg)=g^{-1}u(p)g\)；令 \(\mathcal A(P)\) 为 \(P\) 上全部光滑主联络组成的空间。对基于 \(x_0\) 的连续有限分片光滑环路 \(\gamma\)，以 \(\omega\) 的水平运输 \(T_\gamma(p_0)=p_0h_\gamma\) 定义返回元素，令 \(\operatorname{Hol}_{p_0}(\omega)=\{h_\gamma\}\)。则取值映射 \(\Phi_u\mapsto u(p_0)\) 给出群同构
	\[
		\operatorname{Stab}_{\mathcal G(P)}(\omega)
		\cong Z_G(\operatorname{Hol}_{p_0}(\omega)),
	\]
	右侧为与全部环路运输元素交换的群元素所成的中心化子。特别地，固定 \(p_0\) 的规范子群在 \(\mathcal A(P)\) 上自由作用。
\end{theorem}
\begin{proof}
	先核对环路平行移动群确为子群。若先走 \(\gamma\) 再走 \(\eta\)，等变性给出
	\[
		T_{\gamma*\eta}(p_0)=T_\eta(p_0h_\gamma)
		=p_0h_\eta h_\gamma,
		\qquad h_{\gamma*\eta}=h_\eta h_\gamma.
	\]
	常路径给出 \(e\)，反向路径给出逆元；虽然所写乘法次序反转，全部取值仍对群乘法封闭。

	若 \(\Phi^*\omega=\omega\)，则 \(\Phi\) 保持水平分布，因此由提升唯一性与所有 \(T_\gamma\) 交换。记 \(\Phi(p_0)=p_0a\)。对环路，一方面 \(\Phi(T_\gamma p_0)=p_0a h_\gamma\)，另一方面 \(T_\gamma(\Phi p_0)=p_0h_\gamma a\)，故 \(a\) 与全部 \(h_\gamma\) 交换。对从 \(x_0\) 到 \(x\) 的任意分段光滑路径 \(\gamma\)，等变性又迫使
	\begin{equation}\label{g1:eq:gauge-stabilizer-transport}
		\Phi(T_\gamma(p_0)g)=T_\gamma(p_0)a g.
	\end{equation}
	连通光滑流形可用有限段坐标线连接任意两点，所以此式确定 \(\Phi\) 在整个 \(P\) 上的取值，证明取值映射单射。

	反过来，给定中心化子中的 \(a\)，用\eqref{g1:eq:gauge-stabilizer-transport}定义 \(\Phi\)。若另一路径 \(\beta\) 到达同一点，写 \(T_\beta(p_0)=T_\gamma(p_0)k\)。沿 \(\beta*\bar\gamma\) 运输可知 \(k\in\operatorname{Hol}_{p_0}(\omega)\)。同一点的两种纤维表示于是满足 \(g=kg'\)，而
	\[
		T_\beta(p_0)a g'=T_\gamma(p_0)ka g'
		=T_\gamma(p_0)a k g'=T_\gamma(p_0)a g.
	\]
	所以定义与路径无关。

	为核对光滑性，在任一点附近取坐标球，先沿一条固定路径到球心，再沿坐标径向线到各点。水平提升对参数光滑依赖，故 \(s(x)=T_{\gamma_x}(p_0)\) 是该球上的光滑截面。在此平凡化中 \(\Phi(s(x)g)=s(x)a g\)，所以 \(\Phi\) 光滑；把 \(a\) 换成 \(a^{-1}\) 给出光滑逆。定义也表明它与沿任意路径的水平运输交换，故保持水平分布，因而保持联络。最后，复合在 \(p_0\) 的取值是 \(ab\)，所以取值映射保持群运算。若 \(\Phi\) 还固定 \(p_0\)，则 \(a=e\)，上式给出 \(\Phi=\operatorname{id}\)，得到自由性结论。
\end{proof}

稳定子的出现是具体的对称现象。例如 \(G\) 的中心总包含在上述中心化子中；对交换群，只要底空间连通，保持联络的规范变换就由一个常群元素确定。使用带基点的规范群可以消去这些稳定子，但仅凭作用自由仍不能断言商是流形：还需要选择函数空间、拓扑和适当的局部分析。

\subsection{平坦丛的单值构造}
\label{g1:subsec:1804:03}

主联络 \(\omega\) 平坦，指 \(\Omega=0\)。\cref{g1:thm:flat-principal-local-section}已把这个条件转成局部存在水平截面 \(s\)，即 \(s^*\omega=0\)。如果两个这样的截面满足 \(s'=s k\)，局部变换公式给出 \(0=k^*\theta^L\)。\(\theta^L_g:\mathrm{T}_gG\to\mathfrak g\) 可逆，故 \(\mathrm{d}k=0\)，\(k\) 在每个交叠连通分支上为常数。这个局部常数性质使第十五章的平坦运输论证可以不依赖向量空间结构而继续成立。

\begin{proposition}[平坦主丛的同伦运输]\label{g1:prop:flat-principal-homotopy}
	设有限维实 Lie 群 \(G\) 的光滑右主丛 \(P\to M\) 带平坦主联络，\(M\) 为无边界、Hausdorff、第二可数光滑流形。水平运输唯一延拓到连续路径，且只依赖定端点连续同伦类。若 \(M\) 非空且单连通，则给定 \(p_0\in P_{x_0}\)，存在唯一整体水平截面 \(s\)，满足 \(s(x_0)=p_0\)；从而带联络的丛同构于具有乘积水平分布的 \(M\times G\)。
\end{proposition}
\begin{proof}
	在水平截面域 \(U\) 内定义
	\[
		T^U_{y\leftarrow x}(s(x)g)=s(y)g.
	\]
	若改用 \(s'=s k\)，而 \(x,y\) 在交叠的同一连通分支内，则 \(k(x)=k(y)\)，所以这个映射不变。连续路径的像可以细分成有限段，每段落在一个这样的截面域内；将局部运输依路径次序复合。增加分点不改变结果，因为 \(T^U_{z\leftarrow y}T^U_{y\leftarrow x}=T^U_{z\leftarrow x}\)。两种选择取共同细分后逐段比较，便得到定义无关性。对于分段光滑路径，水平截面坐标中的群分量恒定，故这里的运输与前节水平提升一致。

	对一个定端点连续同伦，将参数方形细分，使每个闭小格的像落在一个水平截面域内。这样的有限网格由紧性和\cref{g1:lem:lebesgue-number}得到。每个小格上，先横后竖和先竖后横都保持同一个群分量，因此边运输满足 \(v_{i+1,j}\circ h_{ij}=h_{i,j+1}\circ v_{ij}\)。将一行的这些等式按次序复合，得到 \(v_{mj}\circ H_j=H_{j+1}\circ v_{0j}\)。固定端点使左右两边的竖边运输都是恒等，故相邻行运输相同，最后得到同伦两端的运输相同。这就是\cref{g1:thm:flat-path-independence}证明中网格消去的同一个论证，未交换非交换群中的因子次序，也不需要先把连续同伦光滑化。

	若 \(M\) 单连通，从 \(x_0\) 到 \(x\) 的任意两条路径定端点同伦，故可定义 \(s(x)=T_\gamma(p_0)\)。在任一连通水平截面域 \(U\) 内，固定到其中一点的路径，再在 \(U\) 内追加路径，便知 \(s\) 等于一个局部水平截面右乘固定群元素。因此 \(s\) 光滑且水平。运输唯一性还保证通过 \(p_0\) 的水平截面唯一。映射 \((x,g)\mapsto s(x)g\) 是主丛微分同胚，并把保持 \(g\) 不变的乘积水平分布送到原水平分布。
\end{proof}

\begin{proposition}[主丛的单值表示]\label{g1:prop:flat-principal-monodromy}
	设 \(\pi:P\to M\) 为有限维实 Lie 群 \(G\) 的光滑主右丛，指定平坦光滑主联络 \(\omega\)，\(M\) 非空、连通、无边界、Hausdorff、第二可数。固定 \(x_0\in M\)、\(p_0\in P_{x_0}\)，对基于 \(x_0\) 的连续环路 \(\gamma\)，以按局部水平截面延伸到连续路径的平行移动 \(T_\gamma\) 定义 \(T_\gamma(p_0)=p_0h_\gamma\)。约定 \(\pi_1(M,x_0)\) 的乘积 \([\gamma][\eta]=[\gamma*\eta]\) 先走 \(\gamma\) 后走 \(\eta\)。则
	\begin{equation}\label{g1:eq:principal-monodromy-convention}
		\rho:\pi_1(M,x_0)\longrightarrow G,\qquad
		\rho([\gamma])=h_\gamma^{-1}
	\end{equation}
	是群同态。将 \(p_0\) 换成 \(p_0a\) 时，\(\rho\) 换成 \(a^{-1}\rho a\)。存在整体水平截面，当且仅当 \(\rho\) 平凡。
\end{proposition}
\begin{proof}
	前一命题使 \(h_\gamma\) 只依赖环路类。由于 \(h_{\gamma*\eta}=h_\eta h_\gamma\)，取逆得到 \(\rho([\gamma][\eta])=\rho([\gamma])\rho([\eta])\)。换基点后，\(T_\gamma(p_0a)=p_0h_\gamma a=(p_0a)(a^{-1}h_\gamma a)\)，故有共轭公式。

	整体水平截面使全部环路运输恒等，所以 \(\rho\) 平凡。反之，若所有环路运输恒等，则两条从 \(x_0\) 到 \(x\) 的路径 \(\gamma,\beta\) 满足 \(T_\beta^{-1}T_\gamma=T_{\gamma*\bar\beta}=\operatorname{id}\)。于是可沿路径定义 \(s(x)=T_\gamma(p_0)\)，并用前一命题中的局部公式证明它光滑且水平。
\end{proof}

这里取逆的原因与\cref{g1:prop:flat-monodromy}完全相同：道路乘积先走左边，运输复合却先作用右边。若 \(\tau:G\to\operatorname{GL}(V)\) 是表示，则关联向量丛的实际环路运输为 \(\tau(h_\gamma)\)，按第十五章约定得到的左单值表示恰是 \(\tau\circ\rho\)。

\begin{theorem}[由单值构造平坦主丛]\label{g1:thm:principal-monodromy-construction}
	设 \(M\) 为非空、连通、有限维、无边界、Hausdorff、第二可数光滑流形，\(G\) 为有限维实 Lie 群，\(\Gamma=\pi_1(M,x_0)\)。给定群同态 \(\rho:\Gamma\to G\)，令 \(\widetilde M\to M\) 为道路类万有覆盖，并取前置环路的左覆盖变换作用 \(L_\gamma\)。则
	\begin{equation}\label{g1:eq:flat-principal-representation-quotient}
		P_\rho=(\widetilde M\times G)/\Gamma,
		\qquad \gamma\cdot(q,g)=(L_\gamma q,\rho(\gamma)g)
	\end{equation}
	带有自然光滑右主 \(G\) 丛结构。局部保持第二因子恒定给出平坦主联络 \(\omega_\rho\)；相对于 \(p_0=[\widetilde x_0,e]\)，其单值恰为 \(\rho\)。
\end{theorem}
\begin{proof}
	流形的坐标球可缩，故\cref{g1:thm:path-class-universal-cover}适用。\cref{g1:prop:universal-deck-discontinuous}给出覆盖纤维上自由传递的左覆盖变换作用，满足 \(L_\gamma L_\delta=L_{\gamma\delta}\)。\(\rho\) 是同态，因此题中公式确为左作用。它与第二因子上的右 \(G\) 作用交换，所以 \([q,g]a=[q,ga]\) 良定义。

	取连通的均匀覆盖坐标球 \(U\)，选一张覆盖片的局部逆 \(\sigma_U:U\to\widetilde M\)。映射
	\[
		U\times G\longrightarrow P_\rho|_U,
		\qquad (x,g)\longmapsto[\sigma_U(x),g]
	\]
	是双射：覆盖纤维的任意点唯一写成 \(L_\gamma\sigma_U(x)\)，而 \([L_\gamma\sigma_U(x),g]=[\sigma_U(x),\rho(\gamma)^{-1}g]\)。它也是同胚，因为轨道商映射开，限制在一张覆盖片乘 \(G\) 后仍是开连续双射。

	若 \(\sigma_V=L_{c(x)}\sigma_U\)，则 \(c(x)\in\Gamma\) 局部常值：在一个更小的连通均匀覆盖邻域中，两局部逆各自落在固定覆盖片内，使它们在一点相等的覆盖变换便在整片上使它们相等。因而截面 \(s_U=[\sigma_U,e]\)、\(s_V=[\sigma_V,e]\) 满足 \(s_V=s_U\rho(c(x))^{-1}\)。过渡函数局部常值，尤其光滑。由这些乘积图得到光滑主丛；不同底点的分离由 \(M\) 的 Hausdorff 性给出，同一纤维的不同点在同一乘积图内分开。底空间有可数平凡化子覆盖，各乘积图均有可数基，故总空间第二可数。

	在每个局部截面中取 \(A_U=0\)。由于过渡函数局部常值，\cref{g1:thm:principal-connection-local-data}的粘合公式成立，故得到主联络，其曲率为零。沿环路 \(\gamma\) 的覆盖提升，从 \(\widetilde x_0\) 到 \(L_{[\gamma]}\widetilde x_0\)，保持 \(g=e\) 给出水平提升。因此
	\[
		T_\gamma(p_0)=[L_{[\gamma]}\widetilde x_0,e]
		=[\widetilde x_0,\rho([\gamma])^{-1}]
		=p_0\rho([\gamma])^{-1},
	\]
	按照\eqref{g1:eq:principal-monodromy-convention}取逆，恢复给定的 \(\rho\)。
\end{proof}

\subsection{表示空间模共轭与平坦结构的对应范围}
\label{g1:subsec:1804:04}

上述构造表明每个表示都能实现。为了得到分类，还必须证明任意平坦主丛都由自己的单值重建，并准确说明什么时候将两个对象视为相同。这里始终固定底流形及其恒等映射；允许底流形微分同胚的分类是另一问题。

\begin{theorem}[平坦主丛的分类]\label{g1:thm:flat-principal-classification}
	设 \(M\) 为非空、连通、有限维、无边界、Hausdorff、第二可数光滑流形，\(G\) 为有限维实 Lie 群，固定 \(x_0\in M\)。光滑右主 \(G\) 丛连同指定平坦主联络，在覆盖 \(\operatorname{id}_M\) 且保持联络的主丛同构下的等价类，与
	\[
		\operatorname{Hom}(\pi_1(M,x_0),G)/G
	\]
	一一对应，其中 \(G\) 按同时共轭作用。若还指定 \(p_0\in P_{x_0}\)，并要求同构保持该点，则对应未经共轭的表示集合 \(\operatorname{Hom}(\pi_1(M,x_0),G)\)。
\end{theorem}
\begin{proof}
	取带点平坦主丛 \((P,\omega,p_0)\)，令 \(\rho\) 为其单值。对 \(q=[\alpha]\in\widetilde M\)，定义
	\[
		\widehat F(q,g)=T_\alpha(p_0)g.
	\]
	同伦不变性保证它与 \(\alpha\) 的代表无关。若 \(\gamma\in\Gamma\)，则
	\[
		T_{\gamma*\alpha}(p_0)=T_\alpha(p_0)\rho(\gamma)^{-1},
		\qquad \widehat F(L_\gamma q,\rho(\gamma)g)=\widehat F(q,g).
	\]
	所以它下降为 \(F:P_\rho\to P\)。

	在一个足够小的连通水平截面域 \(U\) 内，先固定到其中一点的路径，再在 \(U\) 中追加路径，所得 \(T_\alpha(p_0)\) 给出一张水平截面。上述 \(F\) 在对应的局部乘积图中就是 \((x,g)\mapsto s(x)g\)。因此它在每根纤维上双射，有光滑局部逆，整体为主丛微分同胚；它还把乘积的水平曲线送到水平曲线，所以保持联络，并将 \([\widetilde x_0,e]\) 送到 \(p_0\)。这证明重建性。

	若两个带点平坦主丛同构，等变性及水平提升唯一性使它们的 \(h_\gamma\) 完全相同，故表示相同；反过来，相同表示使两者都等价于同一个 \(P_\rho\)，故带点等价类由表示分类。

	忘掉纤维点以后，可以把 \(p_0\) 换成任意 \(p_0a\)，表示正好变成 \(a^{-1}\rho a\)。每个这样的共轭都能由换点实现，且不改变底层带联络的主丛，故得到未指定纤维点的分类。也可直接核对：若 \(\rho'=a^{-1}\rho a\)，则 \([q,g]\mapsto[q,a^{-1}g]\) 给出 \(P_\rho\to P_{\rho'}\) 的保持联络的主丛同构。
\end{proof}

若一开始已固定底层光滑主丛 \(P\)，令 \(\mathcal A_{\mathrm{flat}}(P)\) 表示它的平坦联络集合，则正确的对应是
\begin{equation}\label{g1:eq:fixed-bundle-flat-moduli}
	\begin{split}
		&\mathcal A_{\mathrm{flat}}(P)/\mathcal G(P)\\
		&\quad\cong\{[\rho]\in\operatorname{Hom}(\pi_1(M,x_0),G)/G:
		P_\rho\cong P\text{ 作为光滑主丛}\}.
	\end{split}
\end{equation}
因为选取一个 \(P_\rho\to P\) 的光滑主丛同构可搬运联络，而换一个这样的同构恰好相差规范变换。右侧可能为空，也可能只占全部表示类的一部分。定理分类的是带联络对象，不是把全部表示类不加条件地塞进任意固定主丛。

\ProseParagraphBreak

“表示空间模共轭”在这里首先是一个集合商。即使给表示集合赋予逐点收敛拓扑，共轭商也未必为 Hausdorff 空间。例如取 \(M=S^1\)、\(G=\operatorname{GL}(2,\mathbb R)\)，表示由生成元的像决定。矩阵
\[
	U_t=\begin{pmatrix}1&t\\0&1\end{pmatrix}\qquad(t>0)
\]
都与 \(U_1\) 共轭，因为 \(U_t=\operatorname{diag}(t,1)U_1\operatorname{diag}(t,1)^{-1}\)，但 \(U_t\to I\)，而 \(I\) 不与 \(U_1\) 共轭。因此 \(U_1\) 的共轭轨道不是闭集，商中的对应单点也不闭。若希望得到光滑模空间，必须进一步讨论对象范围、稳定子与商拓扑，不能仅靠这份集合分类。

\subsection{圆周、环面和 Hopf 丛的对照}
\label{g1:subsec:1804:05}

先看圆周。\cref{g1:thm:circle-fundamental}给出 \(\pi_1(S^1)=\mathbb Z\)，所以一般 \(G\) 的平坦主丛由一个群元素的共轭类描述。一个元素描述的不是曲率，而是绕一周以后仍可能留下的整体运输。

\begin{example}[圆周上的相位]\label{g1:exm:circle-principal-flat-moduli}
	取 \(G=\operatorname{U}(1)\)，其 Lie 代数是 \(\mathrm{i}\mathbb R\)。在平凡主丛 \(S^1\times\operatorname{U}(1)\) 的标准截面中，取
	\[
		A_a=\mathrm{i}a\,\vartheta,\qquad a\in\mathbb R,
		\qquad \vartheta=(x\,\mathrm{d}y-y\,\mathrm{d}x)|_{S^1}.
	\]
	局部角度 \(\theta\) 中 \(\vartheta=\mathrm{d}\theta\)，但 \(\theta\) 不是圆周上的整体实值函数。一维底空间上曲率恒为零。沿正向一周，水平方程给出 \(h=\exp(-2\pi\mathrm{i}a)\)，所以 \(\rho(1)=\exp(2\pi\mathrm{i}a)\)。规范变换 \(u(e^{\mathrm{i}\theta})=e^{\mathrm{i}m\theta}\)、\(m\in\mathbb Z\)，使 \(A_a\) 变成 \(A_{a+m}\)。

	这些公式给出全部规范类。事实上，任意局部形式在整体标准截面中写成 \(A=\mathrm{i}f(\theta)\,\mathrm{d}\theta\)，其中 \(f\) 光滑且 \(2\pi\) 周期。令 \(a=(2\pi)^{-1}\int_0^{2\pi}f(t)\,\mathrm{d}t\)，则 \(\chi(\theta)=\int_0^\theta(a-f(t))\,\mathrm{d}t\) 周期，故 \(u=e^{\mathrm{i}\chi}\) 是整体规范变换，并将 \(A\) 变为 \(A_a\)。一般 \(u:S^1\to\operatorname{U}(1)\) 在实线提升为 \(e^{\mathrm{i}\chi(\theta)}\)，满足 \(\chi(\theta+2\pi)-\chi(\theta)=2\pi m\)；这是第五章圆周覆盖提升和整数绕数的结论，提升局部光滑。故 \(\int_{S^1}u^{-1}\mathrm{d}u=2\pi\mathrm{i}m\)，说明 \(A_a,A_b\) 规范等价当且仅当 \(b-a\in\mathbb Z\)。

	所以模集合为 \(\mathbb R/\mathbb Z\simeq\operatorname{U}(1)\)。每个表示都由上述平凡丛上的联络实现，分类定理于是还说明，圆周上每个平坦主 \(\operatorname{U}(1)\) 丛的底层光滑丛都平凡。然而只有 \(a\in\mathbb Z\) 的类具有整体水平截面。
\end{example}

\begin{example}[环面上的两个相位]\label{g1:exm:torus-principal-flat-moduli}
	由\cref{g1:prop:product-fundamental}，\(\pi_1(T^2)\cong\mathbb Z^2\)。一般 \(G\) 的表示由一对交换元素 \((a,b)\) 决定，等价关系为同时共轭。交换条件来自环面的基本群关系，不能省略成任意两个群元素。

	对 \(\operatorname{U}(1)\)，在 \(T^2\times\operatorname{U}(1)\) 上取
	\[
		A=\mathrm{i}(a\,\mathrm{d}\theta+b\,\mathrm{d}\varphi),\qquad a,b\in\mathbb R,
	\]
	其中两角度微分分别是两个圆周因子的整体角度一形式。曲率为零，两个生成元的单值为 \((e^{2\pi\mathrm{i}a},e^{2\pi\mathrm{i}b})\)。每一对相位都能这样得到，所以分类定理表明全部平坦主 \(\operatorname{U}(1)\) 丛均在这份平凡底层丛上实现；其规范类为 \((\mathbb R/\mathbb Z)^2\)。具体的 \(u=e^{\mathrm{i}(m\theta+n\varphi)}\) 给出整数平移 \((a,b)\mapsto(a+m,b+n)\)，而不同相位不能规范等价，因为交换群中的共轭没有作用。规范稳定子在每一类上均为常数 \(\operatorname{U}(1)\)。
\end{example}

前两例中，局部平坦性与非平凡整体运输可以并存。下面的圆周主丛则连平坦联络也不存在；我们用一个直接积分证明这一点，不把障碍留给后卷的示性类理论。

\begin{example}[Hopf 丛与曲率通量]\label{g1:exm:hopf-principal-curvature}
	令 \(\operatorname{U}(1)\) 在 \(S^3\subset\mathbb C^2\) 上右作用为 \((z_0,z_1)\lambda=(z_0\lambda,z_1\lambda)\)。作用自由且固有，商为 \(\mathbb{CP}^1\)，得到\term[hopfcong]{Hopf 丛}[Hopf bundle]
	\[
		\operatorname{U}(1)\longrightarrow S^3\longrightarrow\mathbb{CP}^1.
	\]
	这里 \(\mathbb{CP}^1\) 微分同胚于 \(S^2\)：在图 \([1:w]\)、\(w=x+\mathrm{i}y\) 中，映射为
	\[
		[1:w]\longmapsto\frac{(2x,2y,1-|w|^2)}{1+|w|^2}.
	\]
	其像是去掉南极的单位球面，逆由 \(w=(X+\mathrm{i}Y)/(1+Z)\) 给出；在另一图 \([v:1]\) 中代入 \(w=1/v\)，便在南极光滑延拓，且给出相应光滑局部逆。这也说明可用两张射影图描述球面。下面取复坐标 \((x,y)\) 给出的定向。

	总空间上的 \(\mathrm{i}\mathbb R\) 值一形式
	\[
		\omega=\overline z_0\,\mathrm{d}z_0+\overline z_1\,\mathrm{d}z_1
	\]
	是主联络。它在 \(S^3\) 上纯虚，因为实部为 \(\tfrac12\mathrm{d}(|z_0|^2+|z_1|^2)=0\)；它在基本向量 \((z_0\mathrm{i}t,z_1\mathrm{i}t)\) 上取值 \(\mathrm{i}t\)，且在右乘固定 \(\lambda\in\operatorname{U}(1)\) 下不变。再现条件与等变性因而都成立。

	在 \(U_0=\{z_0\ne0\}\)、\(U_1=\{z_1\ne0\}\) 上分别取截面
	\[
		s_0(w)=\frac{(1,w)}{\sqrt{1+|w|^2}},\qquad
		s_1(v)=\frac{(v,1)}{\sqrt{1+|v|^2}}.
	\]
	交叠处 \(v=1/w\)，直接相乘得到 \(s_1(1/w)=s_0(w)g_{01}(w)\)，其中
	\begin{equation}\label{g1:eq:hopf-transition}
		g_{01}(w)=\frac{|w|}{w}.
	\end{equation}
	因此局部截面不是任意两份独立坐标：在重叠环带上，它们相差一个沿圆周绕行的群元素。

	记 \(r=1+|w|^2\)。将 \(s_0\) 代入 \(\omega\)，微分归一化因子后得到
	\[
		A_0=\frac{\overline w\,\mathrm{d}w-w\,\mathrm{d}\overline w}{2r}
		=\frac{\mathrm{i}(x\,\mathrm{d}y-y\,\mathrm{d}x)}{1+x^2+y^2}.
	\]
	因为群为交换群，曲率在底空间上为整体二形式 \(F\)，其此图表达为
	\begin{equation}\label{g1:eq:hopf-curvature-flux}
		\begin{gathered}
			F|_{U_0}=\mathrm{d}A_0
			=\frac{\mathrm{d}\overline w\wedge\mathrm{d}w}{(1+|w|^2)^2}
			=\frac{2\mathrm{i}\,\mathrm{d}x\wedge\mathrm{d}y}{(1+x^2+y^2)^2},\\
			\frac1{2\pi\mathrm{i}}\int_{\mathbb{CP}^1}F=1.
		\end{gathered}
	\end{equation}
	最后一个积分可在此图中用极坐标计算：去掉的单点对光滑二形式积分没有贡献，而 \(\int_0^\infty 2R(1+R^2)^{-2}\,\mathrm{d}R=1\)，角变量积分为 \(2\pi\)。曲率在另一图光滑，亦保证这一非紧图积分就是整个紧曲面的积分。

	若该主丛有整体截面 \(s\)，则整体一形式 \(A=s^*\omega\) 满足 \(F=\mathrm{d}A\)。由闭曲面上的 Stokes 定理，\(\int F=0\)，与\eqref{g1:eq:hopf-curvature-flux}矛盾，所以丛不平凡。更强地，任何另一主联络与 \(\omega\) 的差都为 \(\pi^*\beta\)，其中 \(\beta\) 是底空间上的 \(\mathrm{i}\mathbb R\) 值一形式；这是\cref{g1:prop:principal-connection-affine}在交换群、平凡伴随表示下的结论。其曲率为 \(F+\mathrm{d}\beta\)，积分仍为 \(2\pi\mathrm{i}\)，所以不可能恒为零。Hopf 丛没有平坦主联络。
\end{example}

\begin{center}
	\small
	\begin{tabular}{p{.19\textwidth}p{.32\textwidth}p{.37\textwidth}}
		\toprule
		底层主丛 & 平坦联络与规范类 & 整体信息 \\
		\midrule
		\(S^1\times\operatorname{U}(1)\) & 相位 \(e^{2\pi\mathrm{i}a}\) & 曲率为零仍可有非平凡环路运输。 \\
		\(T^2\times\operatorname{U}(1)\) & 两个相位 & 两个基本圆周各自记录单值。 \\
		Hopf 丛 \(S^3\to S^2\) & 不存在平坦联络 & 曲率积分非零，且不随联络改变。 \\
		\bottomrule
	\end{tabular}
\end{center}

表中的区别同时涉及联络和底层丛。圆周与环面的例子表明，平凡丛不要求所选平坦联络有整体水平截面；Hopf 例子则说明，有些底层丛连平坦联络都不容许。另一方面，\(S^2\) 单连通，\cref{g1:prop:sphere-simply-connected}与分类定理也说明：其上所有平坦主丛都必须平凡。这与刚才的直接曲率积分相互核对，但不能用“局部可以令联络为零”来跳过整体障碍。


\begin{chaptersummary}
	主丛的局部截面满足 \(s_j=s_i g_{ij}\)，三重交叠给出 \(g_{ij}g_{jk}=g_{ik}\)；反过来，光滑余循环通过局部乘积的粘合重建主丛。表示把主丛变成关联向量丛，其截面对应总空间上的等变函数。对闭子群 \(H\subseteq G\)，结构群约化等价于在 \(P/H\to M\) 中选取光滑截面。

	\ProseParagraphBreak

	主联络既是右平移不变的垂直补空间，也是满足再现与伴随等变条件的一形式 \(\omega\)。其局部形式 \(A_i=s_i^*\omega\) 按 \(A_j=\operatorname{Ad}_{g_{ij}^{-1}}A_i+g_{ij}^*\theta^L\) 变换。单位分解保证整体存在，两联络之差属于 \(\Omega^1(M;\operatorname{ad}(P))\)。表示 \(\rho\) 诱导局部协变微分 \(\mathrm{d}+\rho_*(A_i)\)；标架丛给出与线性联络的双向对应，一般关联表示则可能丢失主联络信息。

	\ProseParagraphBreak

	主作用的等变性使水平提升在给定路径的整个紧参数区间存在，无需群紧。曲率 \(\Omega=\mathrm{d}\omega+\tfrac12[\omega\wedge\omega]\) 是水平等变二形式，记录水平场括号的垂直部分，并满足 \(\mathrm{D}\Omega=0\)。它的积分控制运输随路径的变化；正向小方环满足 \(\log h_\varepsilon=-\varepsilon^2F_{12}+O(\varepsilon^3)\)。平坦性等价于局部存在水平截面，其交叠过渡函数局部常值。

	\ProseParagraphBreak

	规范变换是覆盖底空间恒等映射的主丛自同构。它可以改变联络，而换局部截面只改变表达，底空间坐标变化则作用于微分形式的坐标指标。连通底空间上，保持联络的规范变换由一个纤维点处的取值确定，这些取值组成全部环路运输的中心化子。规范商因而可能有稳定子，不能仅从等价关系就认定它具有流形结构。

	\ProseParagraphBreak

	平坦运输只依赖定端点同伦类。若 \(T_\gamma(p_0)=p_0h_\gamma\)，则 \(\rho([\gamma])=h_\gamma^{-1}\) 按本书道路乘法约定成为同态。在非空连通底空间上，将平坦主丛连同其联络按保持联络且覆盖恒等映射的同构视为等价，则其等价类由 \(\operatorname{Hom}(\pi_1(M),G)/G\) 分类；预先固定底层主丛时，只保留能够重建该丛的表示类。圆周和环面的平坦圆周主丛分别记录一个与两个相位，曲率为零仍可没有整体水平截面。Hopf 丛的曲率积分在本章定向下为 \(2\pi\mathrm{i}\)，不随主联络改变，因而没有平坦联络。
\end{chaptersummary}

\BookExercises

以下习题依次对应主丛构造、联络的等价描述、水平运输与曲率、规范等价与平坦结构。计算时固定右作用及局部截面的变换方向；涉及环路时，分别写出实际运输元素与按本书约定取逆得到的单值表示。

\begin{exercise}\label{g1:ex:principal-pullback}
	设 \(\pi:P\to M\) 为有限维实 Lie 群 \(G\) 的光滑主丛，\(f:N\to M\) 为光滑映射，所有流形均有限维、无边界、Hausdorff、第二可数。
	\begin{enumerate}
		\item 证明 \(f^*P=\{(x,p)\in N\times P:f(x)=\pi(p)\}\) 在子空间拓扑下自然成为主 \(G\) 丛，投影为 \((x,p)\mapsto x\)，作用为 \((x,p)g=(x,pg)\)。求其过渡函数。
		\item 若 \(Q\to N\) 为主 \(G\) 丛，\(F:Q\to P\) 是覆盖 \(f\) 的光滑 \(G\) 等变映射，证明存在唯一覆盖 \(\operatorname{id}_N\) 的主丛同构 \(\widehat F:Q\to f^*P\)，使第二因子投影与 \(\widehat F\) 的复合为 \(F\)。
	\end{enumerate}
\end{exercise}
\begin{solution}
	1. 设 \(P|_{U_i}\) 的局部截面为 \(s_i\)。在 \(f^{-1}(U_i)\) 上定义
	\[
		(x,g)\longmapsto(x,s_i(f(x))g).
	\]
	它到相应部分的 \(f^*P\) 双射。在 \(N\times P|_{U_i}\) 的乘积坐标 \((x,y,g)\) 中，\(f^*P\) 是由 \(y=f(x)\) 描述的图；\((x,g)\) 为光滑子流形坐标，并与子空间拓扑相同。因此 \(f^*P\) 为嵌入子流形，继承 Hausdorff 性和第二可数性。上述图中作用是 \((x,g)a=(x,ga)\)，故它是主丛。若 \(s_j=s_i g_{ij}\)，拉回截面 \(\widetilde s_i(x)=(x,s_i(f(x)))\) 满足
	\[
		\widetilde s_j(x)=\widetilde s_i(x)g_{ij}(f(x)),
	\]
	所以过渡函数是 \(g_{ij}\circ f\)，定义在 \(f^{-1}(U_i\cap U_j)\) 上。

	2. 设 \(\pi_Q:Q\to N\)，唯一可能的公式是
	\[
		\widehat F(q)=(\pi_Q(q),F(q)).
	\]
	覆盖条件 \(\pi(F(q))=f(\pi_Q(q))\) 保证其值属于拉回丛。映射光滑、等变，并覆盖恒等映射。对每个 \(x\in N\)，两边纤维都是自由传递的 \(G\) 集；固定 \(q_0\in Q_x\) 后，\(\widehat F(q_0g)=\widehat F(q_0)g\)，所以逐纤维双射。在 \(Q\) 的局部截面 \(t\) 下，\(F(t(x))\) 给出拉回丛的局部截面，\(\widehat F\) 在这两组平凡化中是 \((x,g)\mapsto(x,g)\)。故逆也光滑，从而为主丛同构。第一分量由覆盖恒等映射确定，第二分量由复合条件确定，所以唯一。
\end{solution}

\begin{exercise}\label{g1:ex:mobius-principal-double-cover}
	在 \(P=S^1\subset\mathbb C\) 上令 \(\operatorname{O}(1)=\{1,-1\}\) 右作用为乘法，投影 \(\pi:P\to S^1\) 为 \(\pi(z)=z^2\)。
	\begin{enumerate}
		\item 证明这是主 \(\operatorname{O}(1)\) 丛，且没有连续全局截面。
		\item 对符号表示 \(\rho(\varepsilon)v=\varepsilon v\)（\(v\in\mathbb R\)），证明 \(P\times_\rho\mathbb R\) 同构于 Möbius 线丛 \(L=(\mathbb R\times\mathbb R)/((t+2\pi,v)\sim(t,-v))\)。
		\item 对平凡一维表示 \(\rho_0(\varepsilon)=1\)，证明关联线丛平凡。说明这一比较为何不能仅从一个关联丛恢复主丛。
	\end{enumerate}
\end{exercise}
\begin{solution}
	1. 每个短圆弧上可选择光滑角度 \(t(x)\)，则 \(s(x)=e^{\mathrm{i}t(x)/2}\) 是 \(\pi\) 的局部截面。每个平方根唯一写成 \(s(x)\varepsilon\)，其中 \(\varepsilon=\pm1\)，所以得到等变平凡化。若存在连续全局截面 \(s\)，映射 \((x,\varepsilon)\mapsto s(x)\varepsilon\) 就是从 \(S^1\times\{1,-1\}\) 到 \(S^1\) 的连续双射。源紧、目标 Hausdorff，故它是同胚；但源有两个非空的开闭分支，目标连通，矛盾。因此连连续截面也不存在。

	2. 定义
	\[
		[e^{\mathrm{i}t/2},v]\longmapsto[t,v]\in L.
	\]
	同一个 \(z=e^{\mathrm{i}t/2}\) 的参数可改变 \(4\pi k\)，而 \([t+4\pi k,v]=[t,v]\)，故与角度选择无关。关联关系把 \((z,v)\) 与 \((-z,-v)\) 视为相同，后者可用 \((t+2\pi,-v)\) 表示；在 \(L\) 中 \([t+2\pi,-v]=[t,v]\)，所以公式与关联商关系相容。反向公式为 \([t,v]\mapsto[e^{\mathrm{i}t/2},v]\)，同样良定义。两式互逆，且在任意短圆弧的角度坐标中都为恒等纤维坐标，故是光滑线丛同构。\cref{g1:exm:mobius-line-bundle}证明的反周期条件 \(a(t+2\pi)=-a(t)\) 还说明，该关联丛没有处处非零截面。

	3. 平凡表示下，\((z,v)\sim(-z,v)\)，所以 \([z,v]\mapsto(z^2,v)\) 是到 \(S^1\times\mathbb R\) 的线丛同构；局部平方根给出光滑逆。主丛 \(P\) 仍是第 1 问中非平凡的双覆盖，而任意主 \(\operatorname{O}(1)\) 丛用平凡表示关联得到的线丛都具有同一形式 \(M\times\mathbb R\)。因此表示可以把主丛的区别全部消去；关联构造的输出必须连同所用表示一起理解。
\end{solution}

\begin{exercise}\label{g1:ex:associated-tensor-operations}
	设 \(P\to M\) 为有限维实 Lie 群 \(G\) 的光滑主丛，\(\rho:G\to\operatorname{GL}(V)\)、\(\tau:G\to\operatorname{GL}(W)\) 为同一基域 \(\mathbb K=\mathbb R\) 或 \(\mathbb C\) 上的有限维光滑表示。记 \(E=P\times_\rho V\)、\(F=P\times_\tau W\)。证明存在自然向量丛同构
	\[
		P\times_{\rho\otimes\tau}(V\otimes W)\cong E\otimes F,
		\qquad P\times_{\rho^*}V^*\cong E^*,
	\]
	其中 \(\rho^*(g)\lambda=\lambda\circ\rho(g^{-1})\)。写出同构公式，并说明两条不同关联丛的向量为什么要用同一个 \(p\in P_x\) 表示后，才能直接对分量作张量积。
\end{exercise}
\begin{solution}
	固定 \(p\in P_x\)，分别有线性同构 \(V\to E_x\)、\(W\to F_x\)。它们的张量积给出
	\[
		[p,v\otimes w]\longmapsto[p,v]\otimes[p,w],
	\]
	再按线性延拓到任意张量。将 \(p\) 换为 \(pg\) 时，同一个左侧元素的分量变成 \(\rho(g^{-1})v\otimes\tau(g^{-1})w\)，右侧两个类分别不变，故公式良定义。固定 \(p\) 时它是两个线性同构的张量积，因此逐纤维可逆；局部取 \(p=s_i(x)\)，公式是固定向量空间中的张量积识别，故它及逆光滑。

	对偶情形，把 \([p,\lambda]\) 送到 \(E_x\) 上的线性泛函
	\[
		[p,v]\longmapsto\lambda(v).
	\]
	若使用 \(pg\)，同一个对偶类由 \(\rho^*(g^{-1})\lambda=\lambda\circ\rho(g)\) 表示，同一个向量由 \(\rho(g^{-1})v\) 表示，所以求值仍为 \(\lambda(\rho(g)\rho(g^{-1})v)=\lambda(v)\)。因此公式独立于参照；它逐纤维为对偶空间的线性同构，在局部截面中是固定的求值配对，故光滑可逆。

	若两个向量最初表示为 \([p,v]\)、\([q,w]\)，其中 \(q=pg\)，则第二个向量应先改写为 \([p,\tau(g)w]\)。因此它们的张量积对应 \([p,v\otimes\tau(g)w]\)，不能直接使用 \(v\otimes w\)。两列数字所属的参照不同；共同参照正是把两个纤维对象的张量积正确翻译为分量运算所需的数据。
\end{solution}

\begin{exercise}\label{g1:ex:orientation-reduction-cover}
	设 \(E\to M\) 为秩 \(r\geq1\) 的实光滑向量丛，\(M\) 有限维、无边界、Hausdorff、第二可数。令 \(P=\operatorname{Fr}(E)\)、\(H=\operatorname{GL}^+(r,\mathbb R)\)。证明 \(P/H\to M\) 是纤维有两个点的光滑覆盖，并将其纤维解释为 \(E_x\) 的两种定向。证明 \(E\) 可定向当且仅当该覆盖有光滑截面，等价地，当且仅当该覆盖同构于 \(M\times\{+1,-1\}\)。若 \(M\) 非空且连通，证明可定向时恰有两种定向。
\end{exercise}
\begin{solution}
	行列式符号给出 \(\operatorname{GL}(r,\mathbb R)/H\cong\{+1,-1\}\)。\(H\) 既开又闭，因而是闭 Lie 子群；\cref{g1:thm:reduction-section-criterion}给出 \(P/H\) 的局部乘积 \(U_i\times\{+1,-1\}\)。每个 \(U_i\) 的逆像分成两个开集，各自微分同胚地投影到 \(U_i\)，所以这是光滑双覆盖。同一纤维中两个标架相差正行列式矩阵，当且仅当它们给出相同定向，因此两个陪集就是两种定向。过渡作用由 \(\operatorname{sgn}\det g_{ij}\) 决定。

	光滑截面在每点连续地选择一个定向，正是向量丛可定向的含义。在局部双覆盖中，连续截面的第二分量取值于离散两点集，因而局部常值，所以它自动光滑。这也表明连续定向与光滑定向在这里没有额外差别。若给定一个截面 \(\sigma\)，把 \((x,+1)\) 送到 \(\sigma(x)\)，把 \((x,-1)\) 送到同一纤维中另一点，局部是两个片的排列，故给出覆盖同构。反过来，乘积覆盖的常值正号截面经同构给出截面。

	当 \(M\) 连通时，任何另一个定向与指定定向的比较符号是连续函数 \(M\to\{+1,-1\}\)，只能为常值。两个常值都能实现，所以恰有两种定向。非空条件排除空底空间只有一个空选择的情形；连通条件则排除不同连通分支可以独立反转定向的情形。
\end{solution}

\begin{exercise}\label{g1:ex:principal-determinant-connection}
	设 \(E\to M\) 为秩 \(r\geq1\) 的实光滑向量丛，\(\nabla\) 在局部标架 \(e=(e_1,\ldots,e_r)\) 中的联络矩阵为 \(A\)。把行列式线丛 \(\det E=\Lambda^rE\) 看成标架丛经表示 \(\det:\operatorname{GL}(r,\mathbb R)\to\mathbb R^\times\) 得到的关联线丛。
	\begin{enumerate}
		\item 求其诱导联络在 \(e_1\wedge\cdots\wedge e_r\) 中的矩阵，并核对与逐因子求协变导数的结果一致。
		\item 两个 \(E\) 上的联络在 \(\det E\) 上诱导相同联络，当且仅当它们的差满足什么条件？在秩 \(2\) 的平凡丛上给出不同联络诱导相同联络的例子。
	\end{enumerate}
\end{exercise}
\begin{solution}
	1. 在单位矩阵处，行列式的微分为 \(X\mapsto\operatorname{tr}X\)。这可由多重线性直接计算：\(\det(I+tX)\) 的一次项必须从恰好一列取 \(tX\)，其余列取标准基，非零贡献恰为 \(t\sum_aX^a{}_a\)。因此由\cref{g1:thm:associated-covariant-derivative}，行列式线的局部联络形式为 \(\operatorname{tr}A\)。逐因子求导也给出
	\[
		\nabla(e_1\wedge\cdots\wedge e_r)
		=\sum_{b=1}^r e_1\wedge\cdots\wedge(\nabla e_b)\wedge\cdots\wedge e_r
		=(\operatorname{tr}A)(e_1\wedge\cdots\wedge e_r),
	\]
	因为 \(\nabla e_b=\sum_a A^a{}_b e_a\)，而 \(a\ne b\) 的项有重复因子，楔积为零。

	2. 记两联络之差为 \(B\in\Omega^1(M;\operatorname{End}E)\)。所得行列式线联络之差为 \(\operatorname{tr}B\)，所以条件是 \(\operatorname{tr}B(v)=0\) 对所有底空间切向量 \(v\) 成立。这是整体条件，因为换标架时 \(B\) 共轭变换，而迹不变。在 \(\mathbb R\times\mathbb R^2\) 的标准标架中取
	\[
		A_0=0,\qquad A_1=\begin{pmatrix}1&0\\0&-1\end{pmatrix}\mathrm{d}x.
	\]
	两者对常截面 \(e_1\) 的协变导数分别为零和 \(e_1\mathrm{d}x\)，故不同；但迹均为零，所以在行列式线上都诱导标准联络。行列式记录总体积变化，不能记录这类互相抵消的两个方向的变化。
\end{solution}

\begin{exercise}\label{g1:ex:principal-connection-relative-extension}
	设 \(P\to M\) 为光滑主右 \(G\) 丛，\(K\subset M\) 闭，\(U\supset K\) 开，且 \(\pi^{-1}(U)\) 上已给定主联络 \(\omega_U\)。证明存在 \(P\) 上的主联络，在 \(K\) 的某个开邻域上等于 \(\omega_U\)。若两整体主联络均在某个固定开集 \(W\) 上等于指定联络，证明它们可由一条在 \(W\) 上保持该联络不变的光滑联络族连接。
\end{exercise}
\begin{solution}
	由\cref{g1:thm:principal-connection-existence}取任意整体主联络 \(\omega_0\)。由\cref{g1:cor:smooth-cutoff-closed}取 \(\chi:M\to[0,1]\)，使 \(\chi=1\) 于 \(K\) 的某开邻域 \(W_0\subset U\)，且 \(\operatorname{supp}\chi\subset U\)。在 \(\pi^{-1}(U)\) 上，\(\omega_U-\omega_0\) 消去垂直方向并按伴随表示等变。令
	\[
		\omega=\omega_0+(\chi\circ\pi)(\omega_U-\omega_0),
	\]
	第二项在 \(\pi^{-1}(U)\) 外补零。支集包含条件保证补零光滑：在 \(U\) 外每点附近，\(\chi\) 恒为零。权重来自底空间，所以这项仍消去垂直方向并等变；\cref{g1:prop:principal-connection-affine}便保证 \(\omega\) 是主联络。\(\chi=1\) 的地方有 \(\omega=\omega_U\)，满足要求。

	若给定两整体联络 \(\omega_0,\omega_1\)，令 \(\omega_t=(1-t)\omega_0+t\omega_1\)，其中 \(0\leq t\leq1\)。等变条件可直接相加，再现条件给出 \(\omega_t(X^\#)=((1-t)+t)X=X\)，故每个 \(\omega_t\) 都是主联络。系数随 \(t\) 光滑变化；在 \(W\) 上二者原本相同，故整族也保持指定联络不变。
\end{solution}

\begin{exercise}\label{g1:ex:pullback-principal-connection}
	设 \(F:N\to M\) 为无边界光滑流形之间的光滑映射，\(P\to M\) 为主右 \(G\) 丛且带主联络 \(\omega\)。令
	\[
		F^*P=\{(y,p)\in N\times P:F(y)=\pi(p)\},\qquad
		\widehat F(y,p)=p.
	\]
	证明 \(F^*P\to N\) 是主右 \(G\) 丛，\(\widehat F^*\omega\) 是其主联络；不需要 \(F\) 是浸入或淹没。写出拉回联络的局部形式，并说明这一构造与复合光滑映射相容。
\end{exercise}
\begin{solution}
	取 \(P\) 的局部截面 \(s_i:U_i\to P\)。在 \(F^{-1}(U_i)\) 上，
	\[
		(y,g)\longmapsto(y,s_i(F(y))g)
	\]
	给出到 \(F^*P\) 的局部平凡化，逆由 \(P|_{U_i}\) 的平凡化求群坐标得到。其过渡函数为 \(g_{ij}\circ F\)，故光滑；不同底点可由 \(N\) 的开集分离，同一纤维的不同点可在一张乘积图中分离，可数局部图又保证第二可数性。右作用为 \((y,p)g=(y,pg)\)，在这些图中就是群坐标右乘，所以给出主丛。

	记 \(\omega'=\widehat F^*\omega\)。基本场满足 \(\mathrm{d}\widehat F(X^\#_{(y,p)})=X^\#_p\)，因此 \(\omega'(X^\#)=X\)。又因 \(\widehat F\circ R_g=R_g\circ\widehat F\)，有
	\[
		R_g^*\omega'=\widehat F^*(R_g^*\omega)
		=\widehat F^*(\operatorname{Ad}_{g^{-1}}\omega)
		=\operatorname{Ad}_{g^{-1}}\omega'.
	\]
	所以它是主联络。新局部截面为 \(s_i'(y)=(y,s_i(F(y)))\)，故局部形式为
	\[
		(s_i')^*\omega'=F^*(s_i^*\omega)=F^*A_i.
	\]
	以上构造只用局部乘积图与微分的拉回，没有对 \(\mathrm{d}F\) 的秩提出条件。

	若还有 \(H:L\to N\)，则自然映射 \(H^*(F^*P)\to(F\circ H)^*P\) 把 \((z,(H(z),p))\) 送到 \((z,p)\)，在局部群坐标中为恒等，故为主丛同构。两边联络的局部形式分别是 \(H^*(F^*A_i)\) 与 \((F\circ H)^*A_i\)，由链式法则相等，因此自然同构保持联络。
\end{solution}

\begin{exercise}\label{g1:ex:circle-weight-connection-detection}
	设 \(P\to M\) 为主 \(\operatorname{U}(1)\) 丛，\(m\in\mathbb Z\)，取表示 \(\rho_m(z)v=z^m v\) 及关联复线丛 \(L_m=P\times_{\rho_m}\mathbb C\)。
	\begin{enumerate}
		\item 设主联络在局部截面中的形式为 \(A_i\in\Omega^1(U_i;i\mathbb R)\)，求 \(L_m\) 的局部联络形式。
		\item 证明当 \(m\ne0\) 时，两个主联络若在 \(L_m\) 上诱导相同联络，就已经相同；当 \(m=0\) 时，每个主联络都诱导标准平凡线丛上的普通微分。说明这与 \(\rho_m\) 在 \(|m|>1\) 时不忠实并不矛盾。
	\end{enumerate}
\end{exercise}
\begin{solution}
	1. 对 \(X=i a\in i\mathbb R\)，有 \(\rho_m(e^{tX})=e^{tmX}\)，故 \((\rho_m)_*(X)=mX\)。由\eqref{g1:eq:associated-local-derivative}，局部形式为 \(mA_i\)，而对局部分量 \(f_i\)，协变微分为 \(\mathrm{d}f_i+mA_i f_i\)。

	2. 设两主联络的局部形式为 \(A_i,B_i\)。诱导联络相同意味着 \(mA_i=mB_i\)。若 \(m\ne0\)，得到 \(A_i=B_i\)，再由\cref{g1:thm:principal-connection-local-data}的唯一性得到两主联络相同。若 \(m=0\)，表示平凡，\([p,v]\mapsto(\pi(p),v)\) 给出 \(L_0\cong M\times\mathbb C\) 的自然平凡化，局部联络形式处处为零，因此诱导联络就是 \(\mathrm{d}\)。

	当 \(|m|>1\) 时，\(\rho_m\) 的核是 \(|m|\) 个单位根，故它不是忠实群表示；但这个核离散，其 Lie 代数为零，微分 \(X\mapsto mX\) 仍单射。比较同一个主丛上的两主联络，只须判断其 Lie 代数值一形式之差，因此这里仍能全部检出。群表示有非平凡离散核，与微分是否丢失联络方向，是不同的问题。
\end{solution}

\begin{exercise}\label{g1:ex:heisenberg-principal-rectangle}
	令 \(G\) 为实三阶上三角、对角线全为 \(1\) 的矩阵 Lie 群，其 Lie 代数包含 \(X=E_{12}\)、\(Y=E_{23}\)、\(Z=E_{13}\)。在 \(\mathbb R^2\times G\) 上以 \(A=X\,\mathrm{d}x+Y\,\mathrm{d}y\) 定义主联络。计算曲率，并求从 \((0,0,e)\) 出发、沿矩形 \([0,s]\times[0,t]\) 正向边界提升的返回群元素，其中 \(s,t>0\)。说明这里的精确面积公式为何不与一般结构群非交换的情形的限制矛盾。
\end{exercise}
\begin{solution}
	矩阵单位满足 \(XY=Z\)、\(YX=0\)，而 \(Z\) 与 \(X,Y\) 的乘积都为零。因此 \([X,Y]=Z\)，并且 \(Z\) 在这个 Lie 代数中居中。\(A\) 的系数为常值，所以
	\[
		F=\mathrm{d}A+A\wedge A=Z\,\mathrm{d}x\wedge\mathrm{d}y.
	\]
	四条边分别给出左乘 \(e^{-sX}\)、\(e^{-tY}\)、\(e^{sX}\)、\(e^{tY}\)。按实际作用次序，返回元素为
	\[
		h=e^{tY}e^{sX}e^{-tY}e^{-sX}.
	\]
	因为 \(X^2=Y^2=Z^2=0\)，这些指数只有常数项与一次项。直接相乘得到
	\[
		(I+tY)(I+sX)(I-tY)(I-sX)=I-stZ=\exp(-stZ).
	\]
	所以该公式对任意正的 \(s,t\) 精确成立。这里原联络系数不交换，但它们的交换子 \(Z\) 与全部系数交换，更高交换子为零；同时曲率为常值。这些特殊条件消去了更高阶修正，不能由这个例子推出任意结构群非交换的联络都满足同一精确公式。
\end{solution}

\begin{exercise}\label{g1:ex:circle-area-principal-holonomy}
	在 \(\mathbb R^2\times\operatorname{U}(1)\) 上取主联络的局部形式
	\[
		A=\frac{\mathrm i\lambda}{2}(x\,\mathrm{d}y-y\,\mathrm{d}x),\qquad \lambda\in\mathbb R.
	\]
	计算曲率，以及沿半径 \(r>0\) 的正向圆周从 \((r,0,1)\) 出发的返回元素。判断哪些 \(\lambda\) 使联络平坦，并说明返回元素为 \(1\) 的单个圆周能否检测平坦性。
\end{exercise}
\begin{solution}
	\(\operatorname{U}(1)\) 为交换群，故括号项为零，而
	\[
		F=\mathrm{d}A=\mathrm i\lambda\,\mathrm{d}x\wedge\mathrm{d}y.
	\]
	以 \(x=r\cos\theta\)、\(y=r\sin\theta\)、\(0\le\theta\le2\pi\) 参数化，有 \(x\,\mathrm{d}y-y\,\mathrm{d}x=r^2\,\mathrm{d}\theta\)。水平方程为 \(g'=-(\mathrm i\lambda r^2/2)g\)，初值 \(g(0)=1\)，因而
	\[
		g(\theta)=\exp(-\mathrm i\lambda r^2\theta/2),\qquad
		h=\exp(-\mathrm i\lambda\pi r^2).
	\]
	联络平坦当且仅当 \(\lambda=0\)。但只要 \(\lambda r^2\in2\mathbb Z\)，该圆周的返回元素就是 \(1\)；例如 \(\lambda=2/r^2\) 时曲率仍处处非零。单个有限大小的环路只给出一个群元素，指数映射的周期可以使非零面积贡献返回单位元。与此不同，所有充分小圆周的返回都为单位元，会由 \(h=1-\mathrm i\lambda\pi r^2+O(r^4)\) 迫使 \(\lambda=0\)。
\end{solution}

\begin{exercise}\label{g1:ex:associated-curvature-detection}
	设 \(\rho:G\to\operatorname{GL}(V)\) 为有限维实 Lie 群的光滑表示。证明：若 \(\rho_*:\mathfrak g\to\operatorname{End}(V)\) 单射，则关联丛联络平坦当且仅当原主联络平坦。用 \(G=(\mathbb R^2,+)\)、\(\rho(a,b)=e^a\in\operatorname{GL}(1,\mathbb R)\) 构造反例，说明不加这个条件时，关联丛可以看不见非零主曲率。再说明检测曲率并不要求 \(\rho\) 本身单射。
\end{exercise}
\begin{solution}
	在任意局部截面中，关联曲率为 \(\rho_*(F)\)。若 \(F=0\)，显然关联曲率也为零；若 \(\rho_*\) 单射，\(\rho_*(F(v,w))=0\) 又推出每个 \(F(v,w)=0\)，所以反向也成立。各局部截面覆盖底空间，故这个结论是整体的。

	对所给 \(\mathbb R^2\) 表示，有 \(\rho_*(a,b)=a\)。在底空间 \(\mathbb R^2\) 的平凡主丛上取 \(A=(0,x\,\mathrm{d}y)\)，其中群坐标与底空间坐标是不同的两组变量。群为交换群，所以
	\[
		F=(0,\mathrm{d}x\wedge\mathrm{d}y)\ne0,
	\]
	但关联联络形式 \(\rho_*(A)=0\)，故关联丛联络平坦。这是表示在第二个 Lie 代数方向上完全消失所造成的。

	最后，取 \(G=\operatorname{U}(1)\)，以 \(\rho(z)v=z^m v\) 作用于 \(V=\mathbb C\cong\mathbb R^2\)，其中整数 \(m\ge2\)。它的核包含全部 \(m\) 次单位根，所以群表示不单射；但微分把 \(\mathrm i a\) 送到复乘法 \(\mathrm i m a\)，因而单射。离散核不影响 Lie 代数层面的曲率检测，尽管它可能影响整体环路返回的检测。
\end{solution}

\begin{exercise}\label{g1:ex:thin-homotopy-principal-transport}
	设 \(P\to M\) 为带主联络的光滑右主 \(G\) 丛，不假设联络平坦。设 \(\Gamma:[0,1]^2\to M\) 光滑到边界，满足 \(\Gamma(s,0)=x\)、\(\Gamma(s,1)=y\)，并且 \(\mathrm{d}\Gamma\) 在每点的秩至多为 \(1\)。证明所有路径 \(t\mapsto\Gamma(s,t)\) 的主丛运输相同。解释这里为何不需要曲率为零，以及为何这并未证明一般定端点同伦都保持运输。
\end{exercise}
\begin{solution}
	固定 \(p\in P_x\)，令 \(q(s,t)\) 是每条路径从 \(p\) 出发的水平提升。初值曲线常值，所以变分公式给出
	\[
		\omega(q_s(s,1))=\int_0^1\Omega(q_t,q_s)\,\mathrm{d}t.
	\]
	在每个点，\(\Gamma_t,\Gamma_s\) 线性相关。曲率水平，而 \(\mathrm{d}\pi:H\to \mathrm{T}M\) 为同构，故它在 \((q_t,q_s)\) 上的取值只取决于这两个底空间输入的水平提升。两个输入线性相关时，交替二形式取值为零。因此积分右边为零。

	终点 \(\Gamma(s,1)=y\) 固定，故 \(q_s(s,1)\) 是竖直向量；主联络在每个竖直空间上与 \(\mathfrak g\) 同构，所以 \(\omega(q_s)=0\) 迫使 \(q_s(s,1)=0\)。终点曲线因而常值，所有运输把 \(p\) 送到同一点。任意其他初值都唯一写成 \(p g\)，运输的等变性进一步说明全部纤维映射相同。

	这个条件通常称为薄同伦：路径族没有扫出二维切向面积，因而任意曲率二形式在该参数族上的拉回都为零。一般同伦允许微分秩为 \(2\)，这时曲率积分可以非零；本题只消除了当前参数族能够测到的曲率，并没有假设或推出底空间上的曲率消失。
\end{solution}

\begin{exercise}\label{g1:ex:infinitesimal-principal-gauge}
	设 \(P\to M\) 为有限维实 Lie 群 \(G\) 的光滑右主丛，\(\xi\) 为 \(\operatorname{ad}(P)\) 的光滑截面。把它写成等变函数 \(\xi:P\to\mathfrak g\)，令 \(u_t(p)=\exp(t\xi(p))\)。证明 \(u_t\) 给出一族规范变换，并求固定局部截面中联络及曲率的一阶变化。说明这些局部变化为何能拼成整体对象。
\end{exercise}
\begin{solution}
	伴随丛截面的等变条件为 \(\xi(pg)=\operatorname{Ad}_{g^{-1}}\xi(p)\)。指数映射与共轭相容，所以
	\[
		u_t(pg)=\exp(t\operatorname{Ad}_{g^{-1}}\xi(p))
		=g^{-1}\exp(t\xi(p))g.
	\]
	由\cref{g1:prop:gauge-equivariant-functions}，\(\Phi_{u_t}\) 是规范变换。取固定局部截面 \(s_i\)，记 \(\xi_i=\xi\circ s_i\)、\(A_i=s_i^*\omega\)。局部规范公式为
	\[
		A_i(t)=\operatorname{Ad}_{\exp(-t\xi_i)}A_i
		+(\exp(t\xi_i))^*\theta^L.
	\]
	第一项在 \(t=0\) 的导数为 \(-[\xi_i,A_i]=[A_i,\xi_i]\)。第二项的导数为 \(\mathrm{d}\xi_i\)：在群单位元附近取坐标，\(\exp(t\xi_i(x))=e+t\xi_i(x)+O(t^2)\)，而 \(\theta^L_e=\operatorname{id}\)，故对任意底空间切向量 \(v\)，该项为 \(t\,\mathrm{d}\xi_i(v)+O(t^2)\)。这里是坐标中的一阶展开，不要求 \(G\) 为矩阵群。因此
	\[
		\left.\frac{\mathrm{d}}{\mathrm{d}t}\right|_{0}A_i(t)
		=\mathrm{d}\xi_i+[A_i,\xi_i]=\mathrm{D}_{A_i}\xi_i.
	\]
	同理，由 \(F_i(t)=\operatorname{Ad}_{\exp(-t\xi_i)}F_i\)，有
	\[
		\left.\frac{\mathrm{d}}{\mathrm{d}t}\right|_{0}F_i(t)=[F_i,\xi_i].
	\]
	在交叠处，\(\xi_j=\operatorname{Ad}_{g_{ij}^{-1}}\xi_i\)。\cref{g1:cor:adjoint-connection-local}说明 \(\mathrm{D}_A\xi\) 是伴随丛值一形式，局部表达按同一伴随规则变化；曲率与括号的相容性也使 \([F_i,\xi_i]\) 拼成伴随丛值二形式。因而这两项分别是整体联络空间中的切向变化与曲率的变化。特别地，平坦联络沿规范轨道仍平坦。
\end{solution}

\begin{exercise}\label{g1:ex:punctured-plane-flat-principal}
	设 \(p_j=(a_j,b_j)\)、\(1\le j\le k\)，是平面上互异的点，\(k\ge1\)，令 \(M=\mathbb R^2\setminus\{p_1,\ldots,p_k\}\)。定义
	\[
		\alpha_j=\frac{(x-a_j)\,\mathrm{d}y-(y-b_j)\,\mathrm{d}x}
		{(x-a_j)^2+(y-b_j)^2}.
	\]
	在平凡主 \(\operatorname{U}(1)\) 丛上取 \(A=\mathrm{i}\sum_j c_j\alpha_j\)，\(c_j\in\mathbb R\)。求曲率和绕各缺失点正向一周的单值。证明这族联络实现 \(M\) 上全部平坦主 \(\operatorname{U}(1)\) 丛的规范类，并给出参数之间的等价关系。
\end{exercise}
\begin{solution}
	在任何可以选取相对于 \(p_j\) 的角度 \(\theta_j\) 的小邻域内，\(\alpha_j=\mathrm{d}\theta_j\)，所以 \(\mathrm{d}\alpha_j=0\)。\(\operatorname{U}(1)\) 为交换群，故 \(F=\mathrm{d}A=0\)。取从同一基点到 \(p_j\) 附近的小圆周的道路，绕该圆周正向一周后原路返回，记所得环路为 \(\gamma_j\)。来回连接道路的积分相消，小圆周上的积分满足
	\[
		\int_{\gamma_j}\alpha_\ell=2\pi\delta_{j\ell}.
	\]
	其中 \(j=\ell\) 时用圆周角度直接计算；\(j\ne\ell\) 时，\(\alpha_\ell\) 在该小圆周围成的整个圆盘上光滑且闭，由 Stokes 定理积分为零。因此实际返回元素为 \(h_{\gamma_j}=e^{-2\pi\mathrm{i}c_j}\)，本章的左单值为
	\[
		\rho([\gamma_j])=e^{2\pi\mathrm{i}c_j}.
	\]
	\cref{g1:prop:punctured-euclidean}证明 \(\pi_1(M)\) 是以这些环路为生成元的自由群 \(F_k\)。向 \(\operatorname{U}(1)\) 的同态可在自由生成元上任意指定 \(k\) 个相位，而交换目标群中的共轭不改变同态。上述公式实现全部相位，故\cref{g1:thm:flat-principal-classification}说明全部平坦主丛都同构于这族联络中的一员，特别其底层光滑主丛均平凡。规范等价的充要条件是
	\[
		(c_1',\ldots,c_k')-(c_1,\ldots,c_k)\in\mathbb Z^k,
	\]
	所以规范类集合为 \((\mathbb R/\mathbb Z)^k\)。也可以直接实现整数平移：函数
	\[
		u(x,y)=\prod_{j=1}^k
		\left(\frac{x-a_j+\mathrm{i}(y-b_j)}{\sqrt{(x-a_j)^2+(y-b_j)^2}}\right)^{m_j}
	\]
	在 \(M\) 上光滑，且 \(u^{-1}\mathrm{d}u=\mathrm{i}\sum_jm_j\alpha_j\)。这核对了局部规范公式与单值分类中的等价关系。自由群中的非交换信息在这个以交换群为目标的表示下被消去，但缺失点补空间的基本群本身仍是 \(F_k\)，不是预先被改成 \(\mathbb Z^k\)。
\end{solution}

\begin{exercise}\label{g1:ex:so3-gauge-stabilizers}
	设圆周上的平坦主 \(\operatorname{SO}(3)\) 丛的单值生成元是一绕固定轴旋转 \(\theta\) 的矩阵，\(0\le\theta\le\pi\)。求其规范稳定子，分别处理 \(\theta=0\)、\(0<\theta<\pi\)、\(\theta=\pi\)。说明为什么这些稳定子可以发生跳变，而固定纤维基点的规范子群仍自由作用。
\end{exercise}
\begin{solution}
	由\cref{g1:thm:gauge-stabilizer-centralizer}，稳定子同构于单值像的中心化子。实际环路返回是该旋转的逆，但一个元素与其逆具有同一中心化子。选定轴为第三坐标轴，旋转写成 \(R=\operatorname{diag}(R_\theta,1)\)。若 \(\theta=0\)，则 \(R=I\)，其中心化子为整个 \(\operatorname{SO}(3)\)。

	若 \(0<\theta\le\pi\)，\(R\) 的实特征值 \(1\) 的特征空间恰为旋转轴。任何与 \(R\) 交换的正交矩阵都保持这条直线，也保持其正交平面，故写成 \(Q=\operatorname{diag}(B,\varepsilon)\)，其中 \(B\in\operatorname{O}(2)\)、\(\varepsilon=\pm1\)，并有 \(\varepsilon=\det B\)。交换条件归结为 \(BR_\theta=R_\theta B\)。

	当 \(0<\theta<\pi\) 时，写 \(R_\theta=\cos\theta\,I+\sin\theta\,J\)，其中 \(J=\left(\begin{smallmatrix}0&-1\\1&0\end{smallmatrix}\right)\)。因为 \(\sin\theta\ne0\)，条件等价于 \(BJ=JB\)。设 \(B=\left(\begin{smallmatrix}a&b\\c&d\end{smallmatrix}\right)\)，直接比较两边得到 \(d=a\)、\(c=-b\)；再由正交性得 \(a^2+b^2=1\)，所以 \(B\in\operatorname{SO}(2)\)。此时 \(\varepsilon=1\)，中心化子是绕原轴的全部旋转，同构于 \(\operatorname{SO}(2)\)。

	当 \(\theta=\pi\) 时，\(R_\pi=-I\) 在垂直平面上为标量，故任意 \(B\in\operatorname{O}(2)\) 都与它交换。中心化子为
	\[
		\{\operatorname{diag}(B,\det B):B\in\operatorname{O}(2)\}
		\cong\operatorname{O}(2).
	\]
	其中行列式为 \(-1\) 的平面变换还反转旋转轴，这部分在一般角度时不再交换。因此稳定子不仅在恒等单值处增大，在半周旋转处也多出一个连通分支。固定 \(p_0\) 的规范子群则要求稳定子对应元素 \(u(p_0)=e\)，由同一定理只能是恒等规范变换，所以其作用一直自由；这不与全规范群的稳定子跳变矛盾。
\end{solution}

\begin{exercise}\label{g1:ex:hopf-powers-flat-obstruction}
	在 Hopf 丛的两张射影图上，令 \(g_{01}(w)=|w|/w\)，并以 \(g_{01}^m\) 为过渡函数构造主 \(\operatorname{U}(1)\) 丛 \(P_m\to\mathbb{CP}^1\)，其中 \(m\in\mathbb Z\)。证明原 Hopf 联络的局部形式 \(A_0,A_1\) 乘以 \(m\) 后给出 \(P_m\) 上的联络，求其曲率积分，并判断 \(P_m\) 何时平凡、何时容许平坦联络。
\end{exercise}
\begin{solution}
	连同 \(g_{10}^m=(g_{01}^m)^{-1}\)、\(g_{ii}^m=1\)，这组函数满足余循环条件，故\cref{g1:thm:principal-cocycle-gluing}给出 \(P_m\)。原来的局部联络满足 \(A_1=A_0+g_{01}^{-1}\mathrm{d}g_{01}\)。对整数 \(m\)，包括负整数，由乘积求导及交换性可得
	\[
		(g_{01}^m)^{-1}\mathrm{d}(g_{01}^m)
		=m\,g_{01}^{-1}\mathrm{d}g_{01}.
	\]
	于是 \(mA_1=mA_0+(g_{01}^m)^{-1}\mathrm{d}(g_{01}^m)\)，满足主联络的粘合公式。其曲率是 \(F_m=mF\)，按正文的复定向有
	\[
		\frac1{2\pi\mathrm{i}}\int_{\mathbb{CP}^1}F_m=m.
	\]
	若 \(m\ne0\)，全局截面会使 \(F_m\) 成为整体恰当二形式，Stokes 定理迫使积分为零，矛盾，故 \(P_m\) 不平凡。对任意另一联络，\(\operatorname{U}(1)\) 的伴随表示平凡，联络之差下降为底空间上的整体 \(\mathrm{i}\mathbb R\) 值一形式 \(\beta\)，曲率之差为 \(\mathrm{d}\beta\)。积分因而仍是 \(2\pi\mathrm{i}m\)，所以没有平坦联络。

	若 \(m=0\)，过渡函数全部为 \(1\)，主丛为乘积，局部形式 \(mA_i=0\) 粘成乘积平坦联络。因此平凡性和容许平坦联络在这族中都恰好等价于 \(m=0\)。这里断言的是存在平坦联络，不是说平凡丛上的每个联络都平坦；后一说法在二维底空间上并不成立。
\end{solution}

% !TeX root = ../../Book1.tex
% 纲要标题：## 第19章*　紧 Lie 群的极大环面与根几何
% 纲要位置：几何学计划纲要.md，第 842 行。
% 正文状态：初稿；尚未正式审核。

% 纲要内容开始
% 
% > 本章目标与证明范围：保留紧连通 Lie 群极大环面存在、元素共轭入环面及环面共轭性的完整证明。证明所需的紧 Riemann 引理在本章自备；完整根系分类不作为此前流形理论的前提。
% 
% 纲要内容结束

\OptionalChapter{紧 Lie 群的极大环面与根几何}{g1:ch:19}
\BookChapterNumber{g1:ch:19}

\begin{chapterabstract}
	紧连通 Lie 群中的单个元素虽然未必与其他元素交换，却总能放进一个连通交换的闭子群。这样的极大环面彼此共轭，使群内的共轭问题归结为环面上的有限对称。本章先由归一化积分构造双不变度量，补足紧群上最短路径与指数满射所需的论证，再证明极大环面的存在、元素入环面及环面共轭定理。伴随作用的权分解给出正则元素、根方向和反射，整数格则补充无穷小结构不能记录的整体信息。最后在酉群、旋转群和紧辛群中计算环面、根、Weyl 群、中心与基本群，并用低维覆盖比较相同 Lie 代数的不同群。
\end{chapterabstract}

\begin{learninggoals}
	\begin{itemize}
		\item 从紧群不变积分构造伴随不变内积，说明双不变度量中的一参数子群为何是测地线，并证明指数映射满射。
		\item 由一参数子群的闭包得到环面，证明紧连通交换群的格商结构，并把任一元素放进极大环面。
		\item 用环面作用的无穷小分解辨认正则元素，证明极大环面共轭，区分群元素共轭与 Lie 代数元素共轭。
		\item 解释正规化子、中心化子、根空间与根反射的作用，辨认根格、特征标及指数核分别保存的信息。
		\item 计算经典紧群中的具体环面与 Weyl 群，用中心、基本群和二重覆盖检验局部 Lie 代数与整体群的区别。
	\end{itemize}
\end{learninggoals}

一个酉矩阵可以在某组酉基下写成对角矩阵，但两个酉矩阵通常不能在同一组基下同时对角化。这里已经隐含两个层次的问题：单个元素可以移到怎样的交换子群中，以及不同选择的交换子群之间有什么关系。三维旋转给出另一个图像：每个旋转都有一条轴，绕固定轴的全部旋转组成圆周群；改变轴则通过整体旋转实现。极大环面定理把这两种经验推广到一般紧连通 Lie 群。

\ProseParagraphBreak

从矩阵模型猜测结论并不困难，困难在于一般 Lie 群没有事先给定的特征向量。我们将先用紧性取得伴随不变内积，再沿一条最短路径把任一群元素写成一个指数。该一参数子群的闭包是紧连通交换群，因此成为环面。环面的共轭性则来自另一个极值问题：把伴随轨道上的一个线性函数取到最大，极值条件会迫使两个 Lie 代数元素交换。读者可以把这两个极值步骤分开记忆：前者寻找通往给定群元素的路径，后者改变无穷小方向的基准位置。

\ProseParagraphBreak

本章的 Lie 群均为有限维实 Lie 群，默认的流形条件仍为无边界、Hausdorff、第二可数；涉及主结论时明确列出紧性与连通性。主要准备是\chapref{g1:ch:16}的指数、闭子群与同态积分，\chapref{g1:ch:17}的齐性空间与紧群平均，以及\chapref{g1:ch:18}的主丛水平提升。最短路径所需的度量论证在本章建立。经典群基本群的计算还使用第五章的覆盖与 van Kampen 定理；相应论证会恢复实际需要的版本。

\ProseParagraphBreak

环面记为 \(T\)，其 Lie 代数记为 \(\mathfrak t\)，与切丛的正体记号 \(\mathrm{T}M\) 区分。根采用实线性泛函 \(\alpha\in\mathfrak t^*\)：根向量 \(Z\) 满足 \([H,Z]=\mathrm{i}\alpha(H)Z\)，所以 \(\operatorname{Ad}_{\exp H}Z=e^{\mathrm{i}\alpha(H)}Z\)。由此，指数核 \(\Lambda=\ker(\exp|_{\mathfrak t})\) 与根之间的整数条件是 \(\alpha(\Lambda)\subseteq2\pi\mathbb Z\)。保留这项归一化，在比较圆周角度、根和特征标时就不会遗漏 \(2\pi\)。

\ProseParagraphBreak

延伸阅读可查 Hall 的《Lie Groups, Lie Algebras, and Representations》中题为“Compact Lie Groups and Maximal Tori”的一章\cite{Hall2015}。本章以极大环面定理及具体根模型为范围；一般根系的分类及最高权表示理论需要进一步的代数工具。

% !TeX root = ../../Book1.tex
% 纲要标题：### 19.1 不变积分与双不变度量
% 纲要位置：几何学计划纲要.md，第 846 行。
% 正文状态：初稿；尚未正式审核。

% 纲要内容开始
% 
% * 紧 Lie 群体积的归一化
% * 左右不变性与平均
% * 伴随不变内积和双不变度量
% * 紧连通 Lie 群的最短测地线与指数满射
% 
% 纲要内容结束

\section{不变积分与双不变度量}
\label{g1:sec:1901}

在一般连通 Lie 群中，一个元素可以写成有限个指数的乘积，却未必是单个指数。紧群上的长度提供了进一步的控制：若能找到从单位元到给定元素的最短曲线，并证明这条曲线来自单参数子群，就能把有限乘积改为一个指数。为此，我们需要一个与左右平移同时相容的度量。本节先把第17章的紧群积分用于伴随表示，再从指数微分公式出发建立所需的最短曲线论证。

\subsection{紧 Lie 群体积的归一化}
\label{g1:subsec:1901:01}
设 \(G\) 为有限维实紧 Lie 群，暂不要求连通，\(\mathfrak g=\mathrm{T}_eG\)。\cref{g1:thm:compact-group-integration}给出了 \(G\) 上唯一的归一化左不变光滑正密度 \(\mu_G\)，而且它同时右不变、取逆不变。这里回顾归一化中的两个要点。先把 \(\mathfrak g\) 上任一正密度 \(\nu_e\) 用左平移延拓，得到
\[
	\nu_g(v_1,\ldots,v_n)
	=\nu_e\bigl(\mathrm{d}(L_{g^{-1}})_gv_1,\ldots,
	\mathrm{d}(L_{g^{-1}})_gv_n\bigr),\qquad n=\dim G.
\]
紧性保证 \(0<\int_G\nu<\infty\)，因而可以令 \(\mu_G=\nu/\int_G\nu\)。单位元处的密度空间是一维，所有左不变正密度只相差一个正常数；除以总量后便与最初选择无关。

\ProseParagraphBreak

所谓体积归一化，归一化的是整个群的总量，而非每个坐标片的 Euclid 体积。它也不依赖定向：正密度在换坐标时使用 Jacobian 行列式的绝对值。即使 \(G\) 不连通，上式仍沿左平移规定每个分支上的密度。紧 Lie 群只有有限个连通分支，因为分支是开集，而所有分支构成一个开覆盖；左平移互换这些分支，所以每个分支的归一化体积相同。

\begin{example}\label{g1:exm:normalized-volume-torus}
	在 \(T^r=(S^1)^r\) 上，用角坐标 \(z_j=e^{\mathrm{i}t_j}\) 表示积分，则
	\[
		\int_{T^r}f\,\mathrm{d}\mu_{T^r}
		=\frac1{(2\pi)^r}\int_{[0,2\pi]^r}
		f(e^{\mathrm{i}t_1},\ldots,e^{\mathrm{i}t_r})\,\mathrm{d}t_1\cdots\mathrm{d}t_r.
	\]
	任一常系数正定内积给出的体积密度只是角坐标密度的正常数倍，故归一化后都得到这一积分。相反，加法群 \(\mathbb R^r\) 在 \(r>0\) 时的左不变正密度总量无限，不能作此归一化。归一化积分的存在与任意长度单位的选择是不同问题。
\end{example}

\subsection{左右不变性与平均}
\label{g1:subsec:1901:02}
右不变性是平均伴随表示时实际使用的条件。回顾\cref{g1:thm:compact-group-integration}的证明：左右平移交换，使 \(R_a^*\mu_G\) 仍左不变，故为 \(\Delta(a)\mu_G\)，其中 \(\Delta:G\to\mathbb R_{>0}\) 是连续同态。紧像 \(\log\Delta(G)\) 是 \(\mathbb R\) 的紧加法子群；若含非零数，就含其所有整数倍而无界。因此 \(\Delta\equiv1\)。这一步说明为何我们可以在同一个积分中同时使用左换元和右换元。

\begin{proposition}[表示的平均内积]\label{g1:prop:compact-representation-inner-product}
	设 \(G\) 为有限维实紧 Lie 群，\(\rho:G\to\operatorname{GL}(V)\) 为有限维实光滑表示，\(b\) 为 \(V\) 上的正定内积，\(\mu_G\) 为 \(G\) 上满足 \(\mu_G(G)=1\) 的左右不变 Haar 测度。则
	\[
		\langle u,v\rangle
		=\int_G b(\rho(g)u,\rho(g)v)\,\mathrm{d}\mu_G(g)
	\]
	是 \(\rho(G)\) 不变的正定内积。若 \(V\) 为复向量空间、\(\rho\) 为复线性表示，则平均任一正定 Hermitian 内积也得到不变正定 Hermitian 内积。
\end{proposition}
\begin{proof}
	在固定基下逐分量积分，所得形式仍保留被积式的双线性和对称性。对非零 \(u\)，每个 \(\rho(g)\) 可逆，所以 \(b(\rho(g)u,\rho(g)u)>0\)。\cref{g1:thm:compact-group-integration}的正性给出 \(\langle u,u\rangle>0\)。对 \(a\in G\)，右换元 \(g\mapsto ga\) 给出
	\[
		\langle\rho(a)u,\rho(a)v\rangle
		=\int_G b(\rho(ga)u,\rho(ga)v)\,\mathrm{d}\mu_G(g)
		=\langle u,v\rangle.
	\]
	复情形逐实部、虚部积分，半双线性、Hermitian 对称性和正性由相同论证得到。
\end{proof}

特别地，取 \(V=\mathfrak g\)、\(\rho=\operatorname{Ad}\)，就得到对所有 \(a\in G\) 满足
\[
	\langle\operatorname{Ad}_aX,\operatorname{Ad}_aY\rangle
	=\langle X,Y\rangle
\]
的\term[bansui-bubian-neiji]{伴随不变内积}[Ad-invariant inner product]。若只在 \(G^0\) 上平均，得到的内积未必被其余分支保持；对整个紧群积分消除了这个问题。

\subsection{伴随不变内积和双不变度量}
\label{g1:subsec:1901:03}
左平移能把任意 \(\mathfrak g\) 上的内积延拓成左不变度量，但右平移会在单位元数据上引入共轭。伴随不变性恰好使这两种要求相容。

\begin{proposition}[双不变度量]\label{g1:prop:bi-invariant-metric}
	设 \(G\) 为有限维实 Lie 群，\(\mathfrak g=\mathrm{T}_eG\)，\(\langle\ ,\ \rangle\) 为 \(\mathfrak g\) 上的正定内积。定义
	\[
		h_g(v,w)=\left\langle\mathrm{d}(L_{g^{-1}})_gv,
		\mathrm{d}(L_{g^{-1}})_gw\right\rangle.
	\]
	则 \(h\) 是光滑左不变 Riemann 度量；它同时右不变，当且仅当该内积 \(\operatorname{Ad}(G)\) 不变。每个紧 Lie 群因而具有\term[shuangbubian-duliang]{双不变度量}[bi-invariant metric]。若 \(G\) 连通，则伴随不变性又等价于
	\begin{equation}\label{g1:eq:ad-skew-inner-product}
		\langle[Z,X],Y\rangle+\langle X,[Z,Y]\rangle=0
		\qquad(X,Y,Z\in\mathfrak g).
	\end{equation}
\end{proposition}
\begin{proof}
	平移微分光滑且可逆，故上式光滑、正定。\(L_{(ag)^{-1}}\circ L_a=L_{g^{-1}}\) 说明左不变性。另一方面，在 \(g\) 处求微分有
	\[
		\mathrm{d}(L_{(ga)^{-1}})_{ga}\circ\mathrm{d}(R_a)_g
		=\operatorname{Ad}_{a^{-1}}\circ\mathrm{d}(L_{g^{-1}})_g.
	\]
	因此右平移保持 \(h\) 恰好等价于所有 \(\operatorname{Ad}_{a^{-1}}\) 保持单位元内积；由于取逆遍历 \(G\)，这就是所述条件。紧群情形应用前一命题。

	若内积伴随不变，对 \(a=\exp_G(tZ)\) 在 \(t=0\) 求导，并用\cref{g1:prop:adjoint-differential-bracket}中已证的 \(\operatorname{ad}_ZX=[Z,X]\)，得到\eqref{g1:eq:ad-skew-inner-product}。反之，该式说明每个 \(\operatorname{ad}_Z\) 是斜伴随算子，于是
	\[
		\frac{\mathrm{d}}{\mathrm{d}t}
		\langle e^{t\operatorname{ad}_Z}X,e^{t\operatorname{ad}_Z}Y\rangle=0.
	\]
	\cref{g1:prop:exponential-naturality}给出 \(\operatorname{Ad}_{\exp_G(tZ)}=e^{t\operatorname{ad}_Z}\)，所以每个指数的伴随作用保持内积。\cref{g1:prop:exponentials-generate-identity-component}说明连通 \(G\) 的每个元素是有限个指数的乘积；正交变换的乘积仍正交，故得到整个 \(G\) 的伴随不变性。
\end{proof}

\begin{example}\label{g1:exm:unitary-trace-metric}
	在 \(\mathfrak u(n)=\{X:X^*=-X\}\) 上，取
	\[
		\langle X,Y\rangle=-\operatorname{Re}\operatorname{tr}(XY).
	\]
	它对称，且 \(\langle X,X\rangle=\operatorname{tr}(X^*X)>0\) 对 \(X\ne0\) 成立。由迹的循环性质，酉共轭保持它，故它给出 \(\operatorname{U}(n)\) 的双不变度量，也可限制到其 Lie 子群。并非每个左不变度量都双不变。例如在满足 \([H,E]=E\) 的二维仿射 Lie 代数上，\eqref{g1:eq:ad-skew-inner-product}若成立，就有 \(2\langle E,E\rangle=0\)，与正定性矛盾。
\end{example}

为了把一参数运动与“测地线”联系起来，还须指出这里用哪个联络。第15章已定义度量相容与挠率；双不变内积允许直接写出同时满足度量相容与无挠性的联络，再用它验证一参数子群的测地线方程。

\begin{proposition}[双不变度量的测地线]\label{g1:prop:bi-invariant-geodesics}
	设 \(G\) 为有限维实 Lie 群，配备双不变 Riemann 度量 \(h\)。则存在唯一保持 \(h\) 的无挠联络 \(\nabla\)，在左不变向量场上满足
	\begin{equation}\label{g1:eq:bi-invariant-connection}
		\nabla_{X^L}Y^L=\tfrac12[X,Y]^L.
	\end{equation}
	其测地线，即满足 \(\nabla_{\dot\gamma}\dot\gamma=0\) 的曲线，恰为
	\[
		\gamma(t)=g\exp_G(tX),\qquad g\in G,\ X\in\mathfrak g,
	\]
	或这些曲线的区间限制，参数起点可任意平移。
\end{proposition}
\begin{proof}
	取 \(\mathfrak g\) 的一组基，左平移给出整体标架 \(E_i\)。规定 \(\nabla_{E_i}E_j=\tfrac12[E_i,E_j]\)，再按第一输入的函数线性和第二输入的 Leibniz 法则延拓，便得到光滑联络；这是\cref{g1:prop:connection-local-data}中按标架指定联络系数的构造，此处标架覆盖全群，因而无需另行粘合。左不变场上的挠率为 \(\tfrac12[X,Y]-\tfrac12[Y,X]-[X,Y]=0\)，由\cref{g1:prop:torsion-tensorial}，挠率处处为零。左不变场的内积为常数，而\eqref{g1:eq:ad-skew-inner-product}给出
	\[
		h(\nabla_{Z^L}X^L,Y^L)+h(X^L,\nabla_{Z^L}Y^L)=0.
	\]
	用 Leibniz 法则展开一般场的分量，就得到度量相容条件。

	唯一性也可直接检查。两个这样的联络之差记为 \(B(U,V)\)；联络法则说明它是双线性张量。无挠性使 \(B(U,V)=B(V,U)\)，度量相容性使 \(b(U,V,W)=h(B(U,V),W)\) 关于后两个输入反对称。因此
	\[
		\begin{aligned}
			b(U,V,W)&=b(V,U,W)=-b(V,W,U)\\
			&=-b(W,V,U)=b(W,U,V)\\
			&=b(U,W,V)=-b(U,V,W).
		\end{aligned}
	\]
	故 \(b=0\)，正定性推出 \(B=0\)。这个唯一的无挠度量联络称为 Levi--Civita 联络。\footnote{Tullio Levi-Civita（1873--1941），意大利数学家，发展了绝对微分学与张量分析，对微分几何有重要贡献。}

	沿曲线写 \(\dot\gamma=\sum_i u_iE_i\)，并记左平移回单位元的速度为 \(u(t)=\sum_i u_i(t)E_i(e)\)。\cref{g1:def:along-curve-covariant-derivative}的沿曲线公式及\eqref{g1:eq:bi-invariant-connection}给出
	\[
		\nabla_{\dot\gamma}\dot\gamma
		=\sum_i u_i'E_i+\tfrac12\sum_{i,j}u_iu_j[E_i,E_j]
		=\sum_i u_i'E_i.
	\]
	最后一项因系数对称、括号反对称而抵消。所以测地线条件恰为 \(u'=0\)。设常值为 \(X\)，则曲线满足左不变场 \(X^L\) 的方程，\cref{g1:thm:lie-exponential}的存在唯一性给出 \(g\exp_G(tX)\)，并保证它对所有实参数有定义。
\end{proof}

\subsection{紧连通 Lie 群的最短测地线与指数满射}
\label{g1:subsec:1901:04}
测地线方程给出了每个初速度对应的曲线，却还没有说明它能否到达指定终点，也没有说明它是否最短。我们先证明短指数弧最短，再用紧性选择一条没有角点的最短折线。以下 \(G\) 配备双不变度量，\(|X|=\sqrt{\langle X,X\rangle}\) 表示单位元处的范数，\(|\dot c(t)|=\sqrt{h_{c(t)}(\dot c(t),\dot c(t))}\) 表示曲线的速度。分段 \(C^1\) 曲线 \(c:[a,b]\to G\) 的长度定义为
\[
	L(c)=\int_a^b|\dot c(t)|\,\mathrm{d}t,
\]
积分在各光滑段计算后相加。

\begin{lemma}[指数坐标中的径向估计]\label{g1:lem:compact-group-radial-estimate}
	设 \(G\) 为有限维实 Lie 群，具有双不变度量。在 \(\mathfrak g\) 中可取 \(r>0\)，使 \(\exp_G\) 在包含闭球 \(\overline B(0,2r)\) 的邻域上为到其像的微分同胚。记其局部逆为 \(\log\)。若 \(|X|<r\)，则从 \(e\) 到 \(\exp_GX\) 的任意分段 \(C^1\) 曲线长度至少为 \(|X|\)，而 \(t\mapsto\exp_G(tX)\)、\(0\le t\le1\)，的长度等于 \(|X|\)。从 \(e\) 到 \(G\setminus\exp_GB(0,r)\) 的任意此类曲线长度至少为 \(r\)。
\end{lemma}
\begin{proof}
	\cref{g1:prop:exponential-local-diffeomorphism}保证指数在零点附近可逆；取足够小的 \(r\) 即满足球的要求。\cref{g1:prop:lie-exponential-differential}说明，\((\mathrm{d}\exp_G)_ZY\) 左平移回单位元后等于
	\[
		J_ZY=\int_0^1e^{-s\operatorname{ad}_Z}Y\,\mathrm{d}s.
	\]
	伴随不变性使被积算子正交，而 \(e^{s\operatorname{ad}_Z}Z=Z\)，所以
	\begin{equation}\label{g1:eq:compact-group-radial-pairing}
		J_ZZ=Z,\qquad \langle J_ZY,Z\rangle=\langle Y,Z\rangle.
	\end{equation}
	这就是说，指数的微分保持径向速度，并把垂直于径向的向量仍送到垂直方向；它不声称所有方向的长度都保持。

	若曲线留在该指数坐标域中，写成 \(c(t)=\exp_GZ(t)\)，则速度范数为 \(|J_{Z(t)}Z'(t)|\)。为了允许曲线经过 \(e\)，对 \(\varepsilon>0\) 考虑光滑径向函数 \(r_\varepsilon(t)=\sqrt{|Z(t)|^2+\varepsilon^2}\)。由\eqref{g1:eq:compact-group-radial-pairing}与 Cauchy--Schwarz 不等式，
	\[
		|r_\varepsilon'(t)|
		=\frac{|\langle Z(t),J_{Z(t)}Z'(t)\rangle|}
		{\sqrt{|Z(t)|^2+\varepsilon^2}}
		\le |\dot c(t)|.
	\]
	在各光滑段积分、相加后令 \(\varepsilon\to0\)，得到
	\[
		\bigl||Z(b)|-|Z(a)|\bigr|\le L(c|_{[a,b]}).
	\]
	因此，从 \(e\) 出发而留在指数球内的曲线，其长度至少是终点的径向距离。

	还要排除先离开坐标域再绕回来的竞争曲线。若曲线从 \(e\) 首次离开 \(\exp_GB(0,r)\)，其首次到达的边界点形如 \(\exp_GZ\)、\(|Z|=r\)：这是因为闭球在较大的指数坐标域内，指数把闭球的边界准确映到该邻域的边界。前段曲线留在闭指数球内，上述估计给出其长度至少为 \(r\)。故当终点为 \(\exp_GX\)、\(|X|<r\) 时，任何离开指数球的竞争曲线都比 \(|X|\) 长；未离开的曲线则受径向估计控制。径向曲线的左平凡化速度恒为 \(X\)，长度正好 \(|X|\)，于是达到下界。
\end{proof}

这个局部估计也让我们能够严格使用全群的距离。设 \(G\) 连通，定义
\[
	d(p,q)=\inf_c L(c),
\]
其中 \(c:[a,b]\to G\) 遍历从 \(p\) 到 \(q\) 的分段 \(C^1\) 曲线。任意两点可由此类曲线连接：从一点可达的点集是开集，其补集也开，因为坐标小球内的直线像给出光滑连接；连通性排除了两个非空开集的分离。因此下确界有限。反向参数化与拼接分别给出对称性及三角不等式，左右平移保持长度，因而也保持 \(d\)。径向估计进一步给出
\begin{equation}\label{g1:eq:compact-group-local-distance}
	d(e,\exp_GX)=|X|\quad(|X|<r),\qquad
	B_d(e,r)=\exp_GB(0,r).
\end{equation}
其中第二式的反向包含用到了离开指数球至少花费长度 \(r\)。在更小半径重复同一估计，说明 \(d\) 正定且其开球给出原流形拓扑；左平移将结论搬到每一点。所以 \(d\) 确实是与原拓扑相容的距离。三角不等式还给出
\[
	|d(p,q)-d(p',q')|\le d(p,p')+d(q,q'),
\]
故 \(d\) 在 \(G\times G\) 上连续。

\begin{theorem}[紧群的最短测地线与指数满射]\label{g1:thm:compact-exp-surjective}
	设 \(G\) 为有限维实紧连通 Lie 群，并配备任一双不变 Riemann 度量。对任意 \(p,q\in G\)，存在 \(X\in\mathfrak g\)，使
	\[
		\gamma(t)=p\exp_G(tX),\qquad 0\le t\le1,
	\]
	连接 \(p,q\)，并且 \(L(\gamma)=|X|=d(p,q)\)。特别地，\(\exp_G:\mathfrak g\to G\) 满射。
\end{theorem}
\begin{proof}
	若 \(G\) 为单点，结论取 \(X=0\) 即得。以下设维数正。证明的难点是：一般的极小化曲线列可能没有足够好的导数收敛。我们只选择有限个中间点，把极小化化成紧空间上的连续函数问题，再利用局部指数弧消去所有角点。

	因为 \(d\) 连续而 \(G\times G\) 紧，其最大值 \(D\) 有限。取上一个引理中的 \(r>0\)，再取整数 \(N\ge2\) 满足 \(D/\sqrt N<r\)。固定 \(x_0=p\)、\(x_N=q\)，在紧空间 \(G^{N-1}\) 上定义
	\[
		E_N(x_1,\ldots,x_{N-1})
		=N\sum_{i=0}^{N-1}d(x_i,x_{i+1})^2.
	\]
	它连续，故有最小值 \(m\)。系数 \(N\) 对应把每一段安排在长度为 \(1/N\) 的参数区间上；最小化平方距离和，会同时惩罚过长的单段。

	我们先证 \(m\le d(p,q)^2\)。任取长度 \(L>0\) 的连接曲线，沿其累计弧长选取 \(N+1\) 个点，使相邻两点之间的弧长均为 \(L/N\)。累计弧长是从 \(0\) 到 \(L\) 的连续非减函数，因此即使原曲线有静止段，这些点也能按参数顺序选出。相邻两点的距离至多为 \(L/N\)，于是这个点组满足 \(E_N\le L^2\)。取长度趋于 \(d(p,q)\) 的曲线列，便有 \(m\le d(p,q)^2\)；若 \(p=q\)，也可直接取所有点等于 \(p\)。

	选定一个极小点组 \((x_i)\)。每段距离都满足
	\[
		d(x_i,x_{i+1})\le\sqrt{m/N}\le D/\sqrt N<r.
	\]
	所以\eqref{g1:eq:compact-group-local-distance}允许定义唯一的短向量
	\[
		X_i=\log(x_i^{-1}x_{i+1}),\qquad
		|X_i|=d(x_i,x_{i+1}).
	\]
	把短指数弧 \(x_i\exp_G(sX_i)\)、\(0\le s\le1\)，依次拼接，其长度为 \(\sum_i|X_i|\)。Cauchy--Schwarz 不等式及距离的定义给出
	\[
		d(p,q)\le\sum_i|X_i|
		\le\sqrt{N\sum_i|X_i|^2}=\sqrt m\le d(p,q).
	\]
	因此所有不等式都取等，折线已达到全局最短长度。但若只停在这一步，仍未说明折线在接点处光滑。

	为消去角点，先计算局部平方距离的一阶变化。在 \(\exp_GB(0,r)\) 上令 \(F(y)=\tfrac12|\log y|^2\)。若 \(y(t)=\exp_GZ(t)\)，其左平凡化速度为 \(v=\mathrm{d}(L_{y^{-1}})_y\dot y=J_ZZ'\)，由\eqref{g1:eq:compact-group-radial-pairing}有
	\begin{equation}\label{g1:eq:compact-group-distance-variation}
		\frac{\mathrm{d}}{\mathrm{d}t}F(y(t))
		=\langle Z,Z'\rangle=\langle Z,v\rangle.
	\end{equation}
	固定 \(1\le i\le N-1\)，把中间点改成 \(x_i(t)=x_i\exp_G(tV)\)，其中 \(V\in\mathfrak g\) 任意。其余点不动。由于相邻距离严格小于 \(r\)，小 \(t\) 时相关平方距离仍有光滑的上述表达。对前一段，左平移后为
	\[
		x_{i-1}^{-1}x_i(t)=\exp_G(X_{i-1})\exp_G(tV),
	\]
	在 \(t=0\) 的左平凡化速度为 \(V\)，所以半平方距离的导数为 \(\langle X_{i-1},V\rangle\)。对后一段先用距离的对称性，再左平移成
	\[
		x_{i+1}^{-1}x_i(t)=\exp_G(-X_i)\exp_G(tV).
	\]
	其半平方距离导数为 \(-\langle X_i,V\rangle\)。极小点的一阶导数为零，因此
	\[
		0=\left.\frac{\mathrm{d}}{\mathrm{d}t}\right|_{t=0}E_N
		=2N\langle X_{i-1}-X_i,V\rangle
		\qquad(V\in\mathfrak g).
	\]
	内积非退化，故 \(X_{i-1}=X_i\)。所有短向量于是等于同一个 \(Y\)，从而 \(x_i=p\exp_G(iY)\)。此前的折线以每段相同的时间 \(1/N\) 参数化，恰好成为一条光滑曲线 \(p\exp_G(tNY)\)。令 \(X=NY\)，由\cref{g1:prop:bi-invariant-geodesics}它是测地线，而已经证明的长度等式给出 \(|X|=d(p,q)\)。最后取 \(p=e\)、\(q\) 任意，便有 \(q=\exp_GX\)。
\end{proof}

\begin{remark}\label{g1:rem:compact-exponential-boundaries}
	紧性在证明中提供有限直径和有限点组极小值，双不变性则提供径向配对及相接速度的统一表达。结论不包括指数单射，也不说每条测地线段都最短：在标准圆周群上，\(t\mapsto e^{\mathrm{i}t\theta}\)、\(0\le t\le1\)，若 \(|\theta|>\pi\)，便可换一条更短的圆弧。若紧群不连通，指数的像只能落在 \(G^0\)，因为 \(t\mapsto\exp_G(tX)\) 把每个指数连到 \(e\)；上述定理应用于紧连通的 \(G^0\)，说明指数的像恰为 \(G^0\)。
\end{remark}
% 纲要条目排为编号小节标题，并进入目录与书签；说明段仍保留于纲要注释中。

% !TeX root = ../../Book1.tex
% 纲要标题：### 19.2 极大环面的存在
% 纲要位置：几何学计划纲要.md，第 853 行。
% 正文状态：初稿；尚未正式审核。

% 纲要内容开始
% 
% * 一参数子群的闭包
% * 紧连通交换 Lie 群的结构
% * 极大环面的构造
% * 每个元素包含在极大环面中的证明
% 
% 纲要内容结束

\section{极大环面的存在}
\label{g1:sec:1902}

圆周群中的运动可以反复绕行，两个圆周的乘积还允许以互不成比例的速度同时转动。后一种运动虽然只有一个实参数，其像的闭包却可能是整个二维环面。因此，在紧群中寻找交换方向时，合适的对象不是单参数子群的像本身，而是它的闭包。闭子群定理赋予这个闭包 Lie 群结构；紧连通交换 Lie 群的结构定理则将说明，它正是若干圆周的乘积。

\subsection{一参数子群的闭包}
\label{g1:subsec:1902:01}
设 \(G\) 为紧 Lie 群，\(X\in\mathfrak g\)。记
\[
	A_X\coloneqq\overline{\{\exp_G(tX):t\in\mathbb R\}},
\]
其中闭包取在 \(G\) 中。参数化的像可能不是嵌入子流形；其闭包却是闭子群，可以使用闭子群定理。连通性与交换性也传到闭包，因而将得到一个环面。

\begin{lemma}[交换子群的闭包]\label{g1:lem:abelian-subgroup-closure}
	设 \(G\) 为有限维实 Lie 群，\(A\subseteq G\) 为交换子群。则 \(\overline A\) 仍是交换子群。若 \(A\) 连通，则 \(\overline A\) 连通；若 \(G\) 紧，则 \(\overline A\) 是紧 Lie 子群。特别地，对 \(X\in\operatorname{Lie}(G)\)，令 \(A_X=\overline{\{\exp_G(tX):t\in\mathbb R\}}\)；若 \(G\) 紧，则 \(A_X\) 是紧连通交换 Lie 群。
\end{lemma}
\begin{proof}
	Lie 群为局部 Euclid 且第二可数的空间，点的闭包关系可以用序列检测。取 \(a,b\in\overline A\)，选 \(a_j,b_j\in A\) 分别趋于它们。连续性给出 \(a_jb_j^{-1}\to ab^{-1}\)，所以 \(ab^{-1}\in\overline A\)，这证明闭包为子群。又因 \(a_jb_j=b_ja_j\)，取极限得 \(ab=ba\)。

	连通集的闭包连通：若闭包被两个非空相对开闭集分开，稠密子集 \(A\) 必须与两者都相交，于是也被分开，矛盾。紧群的闭子集紧，而\cref{g1:thm:closed-subgroup}使闭子群成为采用子空间拓扑的嵌入 Lie 子群。最后，\(t\mapsto\exp_G(tX)\) 是连通空间 \(\mathbb R\) 的连续像，且单参数乘法规律使像为交换子群，因此适用于上述结论。
\end{proof}

\cref{g1:exm:irrational-torus-subgroup}中的 \(t\mapsto(e^{\mathrm{i}t},e^{\mathrm{i}\alpha t})\)、\(\alpha\notin\mathbb Q\)，在 \(S^1\times S^1\) 中稠密。这说明 \(\operatorname{Lie}(A_X)\) 可能严格包含 \(\mathbb RX\)：闭包加入了极限方向，不能按最初参数的个数计算闭包的维数。相反，若这条单参数子群有非零周期，它的像已是一个圆周，因而无需再添极限点。\(X=0\) 时闭包只是单位元。

\subsection{紧连通交换 Lie 群的结构}
\label{g1:subsec:1902:02}
现在从一个抽象的紧连通交换 Lie 群 \(A\) 出发。其 Lie 代数 \(\mathfrak a\) 的括号为零，故\cref{g1:cor:commuting-exponentials}表明指数映射保持加法：
\[
	\exp_A(X+Y)=\exp_A X\,\exp_A Y.
\]
指数在零点附近是微分同胚，所以像是包含单位元邻域的子群，因而是开子群。\(A\) 连通，开子群的各陪集又都是开集，故像只能为整个 \(A\)。因此分类问题变成：指数映射的核在向量空间里是什么样的子群？

\begin{lemma}[离散子群的格基]\label{g1:lem:discrete-subgroup-lattice-basis}
	设 \(V\) 为 \(r\) 维实向量空间，\(\Lambda\subset V\) 为离散加法子群。则存在实线性无关的向量 \(v_1,\ldots,v_k\)，其中 \(0\le k\le r\)，使
	\[
		\Lambda=\mathbb Zv_1\oplus\cdots\oplus\mathbb Zv_k.
	\]
	商空间 \(V/\Lambda\) 紧当且仅当 \(k=r\)。当 \(k=r\) 时，称 \(\Lambda\) 为 \(V\) 中的\term[manzhige]{满秩格}[full lattice]。
\end{lemma}
\begin{proof}
	在 \(V\) 上选 Euclid 内积。离散性给出 \(\varepsilon>0\)，使 \(\Lambda\cap B_\varepsilon(0)=\{0\}\)。不同格点之差仍在 \(\Lambda\) 中，所以不同格点的距离至少为 \(\varepsilon\)。这也说明每个有界集只包含有限多个格点，并且 \(\Lambda\) 闭：收敛的格点序列为 Cauchy 序列，最终必须恒定。

	对 \(r\) 归纳。\(\Lambda=\{0\}\) 时取空基。否则选长度最小的非零 \(v_1\in\Lambda\)；最小值存在，因为先取一个非零格点，其长度以内的格点只有有限个。若 \(sv_1\in\Lambda\)，减去 \(\lfloor s\rfloor v_1\) 后得到长度小于 \(\lVert v_1\rVert\) 的格点，故余项为零。这证明
	\[
		\Lambda\cap\mathbb Rv_1=\mathbb Zv_1.
	\]
	令 \(p:V\to v_1^\perp\) 为正交投影。我们需要先证明 \(p(\Lambda)\) 仍离散；任意离散集合的投影并没有这一性质。若存在非零 \(p(\lambda_j)\to0\)，将 \(\lambda_j\) 减去适当整数倍的 \(v_1\)，可使其沿 \(v_1\) 的分量都位于线段 \([0,v_1]\) 内。这些调整后的格点有界，因而只有有限种；其非零投影便不可能趋于零，矛盾。

	由归纳假设，\(p(\Lambda)\) 有格基 \(w_2,\ldots,w_k\)。各取一个提升 \(v_j\in\Lambda\)，满足 \(p(v_j)=w_j\)。任意 \(\lambda\in\Lambda\) 的投影都是这些 \(w_j\) 的整数线性组合，减去相应 \(v_j\) 后，余项位于 \(\Lambda\cap\mathbb Rv_1=\mathbb Zv_1\)。因此 \(v_1,\ldots,v_k\) 生成 \(\Lambda\)。若它们有实线性关系，先投影便得到 \(w_2,\ldots,w_k\) 的实线性关系，迫使相应系数为零，再由 \(v_1\ne0\) 得剩余系数为零。

	记 \(W=\operatorname{span}_{\mathbb R}\Lambda\)。映射 \(V/\Lambda\to V/W\)、\([v]\mapsto v+W\) 连续且满射。若 \(V/\Lambda\) 紧，则向量空间 \(V/W\) 紧，只能有 \(W=V\)，即 \(k=r\)。反之，\(k=r\) 时每个陪集都有代表元位于紧平行多面体
	\[
		\left\{\sum_{j=1}^r t_jv_j:0\le t_j\le1\right\},
	\]
	因为可将各坐标减去整数部分。商空间是这个紧集的连续像，所以紧。
\end{proof}

\begin{theorem}[紧连通交换 Lie 群的结构]\label{g1:thm:compact-connected-abelian-group}
	设 \(A\) 为 \(r\) 维紧连通交换实 Lie 群，\(\mathfrak a=\operatorname{Lie}(A)\)。则
	\[
		\Lambda_A\coloneqq\ker(\exp_A:\mathfrak a\to A)
	\]
	是 \(\mathfrak a\) 中的满秩格，指数诱导 Lie 群同构
	\[
		\mathfrak a/\Lambda_A\ \cong\ A\ \cong\ (S^1)^r.
	\]
	\(r=0\) 时右侧解释为平凡群。
\end{theorem}
\begin{proof}
	前面的论证已给出指数为满的光滑同态。由\cref{g1:prop:exponential-local-diffeomorphism}，零点有一个邻域只包含核的零元，故 \(\Lambda_A\) 离散；它也是闭集 \(\{e\}\) 的原像。\cref{g1:thm:homogeneous-space}应用于加法 Lie 群 \(\mathfrak a\) 及其闭子群 \(\Lambda_A\)，赋予商空间光滑结构。指数在陪集上常值，因此下降成光滑双射 \(\overline{\exp}_A:\mathfrak a/\Lambda_A\to A\)。

	这不仅是一个集合双射。商映射与指数映射在零点附近都是局部微分同胚，故 \(\overline{\exp}_A\) 在单位元附近具有光滑局部逆。它又是同态，用群平移即可将局部逆移到每一点。因此它是 Lie 群同构，商空间也由 \(A\) 的紧性而紧。\cref{g1:lem:discrete-subgroup-lattice-basis}于是给出满秩格基 \(v_1,\ldots,v_r\)。线性坐标 \(\sum_jt_jv_j\mapsto(t_1,\ldots,t_r)\) 将格化为 \(\mathbb Z^r\)，再用
	\[
		(t_1,\ldots,t_r)+\mathbb Z^r\longmapsto
		(e^{2\pi\mathrm{i}t_1},\ldots,e^{2\pi\mathrm{i}t_r})
	\]
	得到到 \((S^1)^r\) 的 Lie 群同构。
\end{proof}

选择格基才得到与标准圆周积的具体同构；指数核本身不依赖这项选择。以后比较具有相同 Lie 代数的紧群时，核格的差别将记录部分整体信息。这里的 \(2\pi\) 来自圆周的角度坐标：若用 \(\mathrm{i}\mathbb R\) 表示 \(\operatorname{Lie}(S^1)\)，指数核为 \(2\pi\mathrm{i}\mathbb Z\)，而不是 \(\mathrm{i}\mathbb Z\)。

\subsection{极大环面的构造}
\label{g1:subsec:1902:03}
\begin{definition}\label{g1:def:maximal-torus}
	设 \(G\) 为有限维实 Lie 群。\(G\) 中的紧连通交换 Lie 子群称为\term[huanmian]{环面}[torus]。若一个环面 \(T\subseteq G\) 不真包含于任何更大的环面，则称 \(T\) 为\term[jidahuanmian]{极大环面}[maximal torus]。平凡群也允许作为零维环面。\symidx{\(T,\mathfrak t\)：极大环面及其 Lie 代数}
\end{definition}

这个定义中的“极大”按包含关系判断。它并不表示存在一个包含全部环面的“最大环面”。例如绕不同轴旋转得到 \(\operatorname{SO}(3)\) 中不同的圆周子群，它们彼此不包含。随后证明的共轭定理将说明，这些极大环面仍具有共同的维数。

\begin{proposition}[环面的极大扩张]\label{g1:prop:maximal-torus-existence}
	设 \(G\) 为有限维紧实 Lie 群。任意环面 \(T_0\subseteq G\) 都包含于一个极大环面。特别地，\(G\) 存在极大环面。
\end{proposition}
\begin{proof}
	考虑全部包含 \(T_0\) 的环面，其维数是 \(0,1,\ldots,\dim G\) 中的整数，且这个集合非空。取其中维数最大的一个 \(T\)。若 \(T\subseteq S\) 而 \(S\) 也是环面，维数最大性给出 \(\dim T=\dim S\)。包含映射在单位元处的微分因维数相同而可逆，故 \(T\) 包含 \(S\) 中单位元的一个开邻域。于是 \(T\) 是连通群 \(S\) 的开子群，只能有 \(T=S\)。这证明 \(T\) 极大。取 \(T_0=\{e\}\) 即得存在性。
\end{proof}

相同论证说明，环面之间的真包含必使维数严格增加；因此此处的极大扩张只需有限维数论证。\cref{g1:lem:abelian-subgroup-closure,g1:thm:compact-connected-abelian-group}又表明，每条单参数子群的闭包 \(A_X\) 是一个环面，所以任意 \(X\in\mathfrak g\) 都属于某个极大环面的 Lie 代数：包含 \(A_X\) 的环面包含整条 \(\exp_G(tX)\)，在 \(t=0\) 求导即可。

\subsection{每个元素包含在极大环面中的证明}
\label{g1:subsec:1902:04}
上一步已经处理所有无穷小方向。要把任意群元素也放入环面，关键是它能否写成一个指数。\cref{g1:thm:compact-exp-surjective}对紧连通群给出了肯定答案；连通性在这里真正进入。

\begin{theorem}[元素属于极大环面]\label{g1:thm:element-in-maximal-torus}
	设 \(G\) 为紧连通实 Lie 群。对每个 \(g\in G\)，存在极大环面 \(T\subseteq G\)，使 \(g\in T\)。
\end{theorem}
\begin{proof}
	由紧连通群的指数满射性，取 \(X\in\mathfrak g\) 满足 \(g=\exp_GX\)。单参数子群的闭包 \(A_X\) 是一个环面，并且包含 \(g\)。由\cref{g1:prop:maximal-torus-existence}，它可扩张到极大环面 \(T\)，于是 \(g\in A_X\subseteq T\)。
\end{proof}

若 \(G\) 不连通，任何环面由于连通且包含单位元，都位于 \(G^0\)，所以别的连通分支中的元素不可能包含在环面中。例如 \(\operatorname{O}(2)\) 中的反射不在任何环面中。对紧连通群，定理则给出
\[
	G=\bigcup_{T\text{ 为极大环面}}T.
\]
这个等式暂时允许不同环面具有不同位置。下一节将证明，它们全部通过群内共轭相互得到，于是可以用一个固定环面研究所有群元素的共轭类。

% !TeX root = ../../Book1.tex
% 纲要标题：### 19.3 共轭性与中心化子
% 纲要位置：几何学计划纲要.md，第 860 行。
% 正文状态：初稿；尚未正式审核。

% 纲要内容开始
% 
% * 环面作用及无穷小分解
% * 正则元素和中心化子
% * 极大环面的共轭定理
% * 环面内共轭与正规化子
% 
% 纲要内容结束

\section{共轭性与中心化子}
\label{g1:sec:1903}

极大环面已经存在，但目前它可能随群元素改变。为了比较两个环面，先研究一个固定环面在 Lie 代数上的伴随作用。交换性使这些线性变换可以同时分解，紧性使它们成为旋转；由此选出一个向量，使与它交换的方向尽量少。将任意伴随轨道与这个向量作内积，再利用紧性取极值，就能把整条轨道带到固定环面的 Lie 代数中。

\ProseParagraphBreak

以下 \(G\) 为紧连通实 Lie 群，\(T\subseteq G\) 为一个极大环面，\(\mathfrak t=\operatorname{Lie}(T)\)。由\cref{g1:prop:bi-invariant-metric}，\(\mathfrak g\) 上有正定的 \(\operatorname{Ad}(G)\) 不变内积 \(\langle\cdot,\cdot\rangle\)。本节每次取正交补或内积极值，都使用这同一个内积。

\subsection{环面作用及无穷小分解}
\label{g1:subsec:1903:01}
对 \(H\in\mathfrak t\)，算子 \(\operatorname{ad}_H:X\mapsto[H,X]\) 是斜伴随的。事实上，\(\operatorname{Ad}_{\exp(sH)}\) 保持内积，对
\[
	\langle\operatorname{Ad}_{\exp(sH)}X,
	\operatorname{Ad}_{\exp(sH)}Y\rangle=\langle X,Y\rangle
\]
在 \(s=0\) 求导，得到
\begin{equation}\label{g1:eq:invariant-inner-product-bracket}
	\langle[H,X],Y\rangle+\langle X,[H,Y]\rangle=0.
\end{equation}
由于 \(\mathfrak t\) 交换，Jacobi 恒等式还给出 \([\operatorname{ad}_H,\operatorname{ad}_K]=\operatorname{ad}_{[H,K]}=0\)。因此我们面对的是一族彼此交换的斜伴随算子。

\begin{lemma}[交换斜 Hermitian 算子的同时分解]\label{g1:lem:commuting-skew-hermitian}
	设 \(V\) 为有限维复 Hermitian 空间，\(A_1,\ldots,A_r\) 为彼此交换的斜 Hermitian 线性算子，即 \(A_j^*=-A_j\)。则 \(V\) 有一组共同特征向量组成的标准正交基；各 \(A_j\) 的特征值均为纯虚数。
\end{lemma}
\begin{proof}
	先看一个算子 \(A\)。复数域上的特征多项式有根，故有单位特征向量 \(v\)。若 \(Av=\lambda v\)，斜 Hermitian 性给出 \(\lambda=-\overline\lambda\)，所以 \(\lambda\) 纯虚。\(v^\perp\) 在 \(A\) 下不变，因为对 \(w\perp v\)，有 \(\langle Aw,v\rangle=\langle w,-Av\rangle=0\)。在 \(v^\perp\) 上归纳，得到标准正交特征基。

	现在先分解 \(A_1\) 的各特征空间。由 \(A_jA_1=A_1A_j\)，每个 \(A_j\) 保持这些空间，其限制仍斜 Hermitian。在每个空间内再按 \(A_2\) 分解，继续至 \(A_r\)。每一步只细分已有的正交分解，所以最后得到共同特征向量的标准正交基。
\end{proof}

为使用此引理，把 \(\mathfrak g\) 复化为 \(\mathfrak g_{\mathbb C}=\mathfrak g\otimes_{\mathbb R}\mathbb C\)。原实内积复双线性延拓为 \(B_{\mathbb C}\)，再定义 Hermitian 内积 \((Z,W)=B_{\mathbb C}(Z,\overline W)\)，其中共轭固定实空间 \(\mathfrak g\)，本式对第一变量线性。\(\operatorname{ad}_H\) 复线性延拓后是斜 Hermitian 算子。取 \(\mathfrak t\) 的一组基同时对角化，每个共同特征空间上的特征值随 \(H\) 实线性变化，可以唯一写成 \(\mathrm{i}\alpha(H)\)，其中 \(\alpha\in\mathfrak t^*\) 为实线性泛函。

\begin{proposition}[环面伴随作用的权分解]\label{g1:prop:torus-weight-decomposition}
	设 \(G\) 为紧连通实 Lie 群，\(T\subseteq G\) 为环面，\(\mathfrak g=\operatorname{Lie}(G)\)、\(\mathfrak t=\operatorname{Lie}(T)\)。固定 \(\mathfrak g\) 上的 \(\operatorname{Ad}(G)\) 不变正定实内积 \(B\)，在 \(\mathfrak g_{\mathbb C}\) 上取其 Hermitian 延拓，约定第一变量线性、第二变量共轭线性。存在有限集 \(R\subseteq\mathfrak t^*\setminus\{0\}\)，使
	\begin{equation}\label{g1:eq:torus-adjoint-weight-decomposition}
		\mathfrak g_{\mathbb C}=\mathfrak g_0\oplus
		\bigoplus_{\alpha\in R}\mathfrak g_\alpha,
		\qquad
		\mathfrak g_\alpha=\{Z:[H,Z]=\mathrm{i}\alpha(H)Z
		\;\text{对所有 }H\in\mathfrak t\}.
	\end{equation}
	这些空间相对于复化后所得的 Hermitian 内积彼此正交，且 \(\overline{\mathfrak g_\alpha}=\mathfrak g_{-\alpha}\)。零权空间为
	\[
		\mathfrak g_0=\{X\in\mathfrak g:[X,H]=0
		\;\text{对所有 }H\in\mathfrak t\}_{\mathbb C}.
	\]
	若 \(\Lambda_T=\ker(\exp_G|_{\mathfrak t})\)，则对 \(\alpha\in R\) 有 \(\alpha(\Lambda_T)\subseteq2\pi\mathbb Z\)，并且
	\begin{equation}\label{g1:eq:torus-weight-character}
		\operatorname{Ad}_{\exp_GH}Z=e^{\mathrm{i}\alpha(H)}Z
		\qquad(Z\in\mathfrak g_\alpha).
	\end{equation}
\end{proposition}
\begin{proof}
	前面的同时分解给出有限个实线性泛函。合并相同泛函的共同特征空间，并将非零者组成 \(R\)，就得到直和分解。若两个泛函不同，可取某个 \(H\) 使其取值不同；斜 Hermitian 算子 \(\operatorname{ad}_H\) 的不同特征空间正交，所以整个分解正交。零权条件就是与 \(\mathfrak t\) 的所有元素交换；将 \(Z=X+\mathrm{i}Y\) 分成实虚部可知它是所列实空间的复化。

	由于 \(H\) 为实向量，对 \([H,Z]=\mathrm{i}\alpha(H)Z\) 取共轭得到 \([H,\overline Z]=-\mathrm{i}\alpha(H)\overline Z\)，所以共轭交换 \(\mathfrak g_\alpha\) 与 \(\mathfrak g_{-\alpha}\)。由\cref{g1:prop:exponential-naturality}，\(\operatorname{Ad}_{\exp_GH}=e^{\operatorname{ad}_H}\)。将算子指数作用在特征空间上，便得\eqref{g1:eq:torus-weight-character}。若 \(H\in\Lambda_T\)，左侧为恒等作用；\(\mathfrak g_\alpha\ne0\) 迫使 \(e^{\mathrm{i}\alpha(H)}=1\)，即 \(\alpha(H)\in2\pi\mathbb Z\)。
\end{proof}

我们把 \(\alpha\) 称为这一作用的\term[wuqiongxiaoquan]{无穷小权}[infinitesimal weight]。其群上版本是\term[tezhengbiao]{特征标}[character]，即光滑群同态 \(\chi_\alpha:T\to S^1\)，这里由 \(\chi_\alpha(\exp H)=e^{\mathrm{i}\alpha(H)}\) 给出。\cref{g1:thm:compact-connected-abelian-group}保证每个 \(t\in T\) 都有这样的指数表示，而刚证明的整性保证换一个相差格向量的 \(H\) 不改变取值。复化将实旋转写成标量乘法；共轭配对说明这些复特征空间仍来自实几何。

\subsection{正则元素和中心化子}
\label{g1:subsec:1903:02}
对一个群子集 \(A\subseteq G\)，与它的每个元素都交换的群元素组成\term[zhongxinhua]{中心化子}[centralizer]
\[
	Z_G(A)\coloneqq\{g\in G:ga=ag\;\text{对所有 }a\in A\}.
\]
它是闭子群，因为每个交换条件都是连续映射的闭水平集。由\cref{g1:thm:closed-subgroup}，其 Lie 代数可以用全部留在中心化子中的单参数子群检测；指数自然性给出
\[
	\operatorname{Lie}(Z_G(A))
	=\{X\in\mathfrak g:\operatorname{Ad}_aX=X\;\text{对所有 }a\in A\}.
\]
这里两个方向都可直接核对：若 \(\exp(sX)\) 与 \(a\) 交换，微分得 \(\operatorname{Ad}_aX=X\)；反之此等式通过指数自然性恢复全部 \(s\) 的交换关系。对于 Lie 代数子集 \(E\subseteq\mathfrak g\)，相应记号为
\[
	\mathfrak z_{\mathfrak g}(E)\coloneqq\{X:[X,Y]=0\;\text{对所有 }Y\in E\}.
\]

\begin{proposition}[极大环面的无穷小中心化子]\label{g1:prop:maximal-torus-lie-centralizer}
	设 \(G\) 为紧连通实 Lie 群，\(T\subseteq G\) 为极大环面，\(\mathfrak t=\operatorname{Lie}(T)\)。则
	\[
		\operatorname{Lie}(Z_G(T))=\mathfrak z_{\mathfrak g}(\mathfrak t)=\mathfrak t.
	\]
	特别地，\(\mathfrak t\) 是 \(\mathfrak g\) 的极大交换 Lie 子代数，且权分解的零权空间为 \(\mathfrak g_0=\mathfrak t_{\mathbb C}\)。
\end{proposition}
\begin{proof}
	因为 \(T=\exp_G(\mathfrak t)\)，与全部 \(\operatorname{Ad}_t\)、\(t\in T\)，相容的不动条件等价于 \([H,X]=0\) 对所有 \(H\in\mathfrak t\) 成立：一个方向对 \(t=\exp(sH)\) 求导，另一个方向用 \(\operatorname{Ad}_{\exp H}=e^{\operatorname{ad}_H}\)。这证明前两个空间相等。

	若 \(X\in\mathfrak z_{\mathfrak g}(\mathfrak t)\)，则 \(\mathfrak a=\mathfrak t+\mathbb RX\) 是交换 Lie 子代数。由\cref{g1:cor:commuting-exponentials}，\(\exp_G(\mathfrak a)\) 是连通交换子群。它的闭包由\cref{g1:lem:abelian-subgroup-closure}是紧连通交换 Lie 群，因而是包含 \(T\) 的环面。\(T\) 的极大性迫使这个闭包等于 \(T\)。于是 \(\exp_G(sX)\in T\) 对全部 \(s\) 成立，在零点求导得 \(X\in\mathfrak t\)。反向包含来自 \(\mathfrak t\) 本身交换。任何包含 \(\mathfrak t\) 的交换 Lie 子代数都在其中心化子中，所以只能等于 \(\mathfrak t\)。零权空间的结论随即由\eqref{g1:eq:torus-adjoint-weight-decomposition}得到。
\end{proof}

对 \(H\in\mathfrak t\)，权分解现在给出更精确的式子
\begin{equation}\label{g1:eq:regular-element-centralizer}
	\mathfrak z_{\mathfrak g}(H)_{\mathbb C}
	=\mathfrak t_{\mathbb C}\oplus
	\bigoplus_{\substack{\alpha\in R\\\alpha(H)=0}}\mathfrak g_\alpha.
\end{equation}
所以，想使与 \(H\) 交换的方向尽量少，只需让 \(H\) 避开有限多个超平面 \(\ker\alpha\)。

\begin{definition}\label{g1:def:regular-toral-element}
	设 \(T\) 为紧连通 Lie 群 \(G\) 的极大环面。满足 \(\mathfrak z_{\mathfrak g}(H)=\mathfrak t\) 的 \(H\in\mathfrak t\) 称为\term[zhengzeyuansu]{正则元素}[regular element]。它们组成的集合记为 \(\mathfrak t_{\mathrm{reg}}\)。
\end{definition}

\begin{lemma}[正则元素的存在]\label{g1:lem:regular-toral-elements}
	设 \(G\) 为紧连通实 Lie 群，\(T\) 为其极大环面，\(\mathfrak g=\operatorname{Lie}(G)\)、\(\mathfrak t=\operatorname{Lie}(T)\)。令 \(R\) 为所有满足下列条件的非零实泛函 \(\alpha\in\mathfrak t^*\) 的集合：存在非零 \(Z\in\mathfrak g_{\mathbb C}\)，使 \([H,Z]=\mathrm{i}\alpha(H)Z\) 对每个 \(H\in\mathfrak t\) 成立。以 \(\mathfrak t_{\mathrm{reg}}\) 记在 \(\mathfrak g\) 中中心化子恰为 \(\mathfrak t\) 的 \(H\in\mathfrak t\) 所成的集合。则
	\[
		\mathfrak t_{\mathrm{reg}}=
		\mathfrak t\setminus\bigcup_{\alpha\in R}\ker\alpha
	\]
	是 \(\mathfrak t\) 中非空的稠密开集。若 \(R=\varnothing\)，则 \(\mathfrak t_{\mathrm{reg}}=\mathfrak t\)。
\end{lemma}
\begin{proof}
	所列等式来自\eqref{g1:eq:regular-element-centralizer}。每个 \(\alpha\ne0\) 的核都是闭的真超平面，其补集开且稠密：在任意开球中沿一个 \(\alpha\) 不为零的方向稍作移动，就能离开超平面。有限多个稠密开集的交仍稠密，可依次在给定非空开集中选一个更小的非空开球避开下一个超平面。这个交也是开集，故有结论。空根集时没有需要避开的超平面。
\end{proof}

这里的正则性首先属于 Lie 代数。对群元素 \(\exp H\)，中心化子的 Lie 代数却由 \(\operatorname{Ad}_{\exp H}\) 的不动向量给出，条件是 \(\alpha(H)\in2\pi\mathbb Z\)，而非 \(\alpha(H)=0\)。例如在 \(\operatorname{SU}(2)\) 中，\(H=\operatorname{diag}(\mathrm{i}\pi,-\mathrm{i}\pi)\) 的 Lie 代数中心化子是对角子代数，但 \(\exp H=-I\) 与全群交换。这一差别将在处理群元素的共轭时再次出现。

\subsection{极大环面的共轭定理}
\label{g1:subsec:1903:03}
固定一个正则 \(H\in\mathfrak t\)。若某个向量 \(Z\) 与 \(H\) 交换，正则性就迫使 \(Z\in\mathfrak t\)。因此，把任意 \(X\in\mathfrak g\) 共轭入 \(\mathfrak t\)，可以转化为寻找一个共轭代表，使它与 \(H\) 交换。内积函数在极值点的一阶条件恰好给出这个括号方程。

\begin{proposition}[伴随轨道与固定环面相交]\label{g1:prop:adjoint-orbit-meets-torus}
	设 \(G\) 为紧连通实 Lie 群，\(T\subseteq G\) 为一个极大环面。对任意 \(X\in\mathfrak g\)，存在 \(g\in G\) 使 \(\operatorname{Ad}_gX\in\mathfrak t\)。
\end{proposition}
\begin{proof}
	选定 \(\operatorname{Ad}(G)\) 不变正定内积，以及\cref{g1:lem:regular-toral-elements}给出的 \(H\in\mathfrak t_{\mathrm{reg}}\)。连续函数
	\[
		f:G\to\mathbb R,\qquad f(g)=\langle\operatorname{Ad}_gX,H\rangle
	\]
	在紧空间 \(G\) 上达到最大值。设最大点为 \(g_0\)，记 \(Z=\operatorname{Ad}_{g_0}X\)。对任意 \(Y\in\mathfrak g\)，曲线 \(s\mapsto\exp(sY)g_0\) 经过最大点，故
	\[
		0=\left.\frac{\mathrm{d}}{\mathrm{d}s}\right|_{s=0}
		f(\exp(sY)g_0)
		=\langle[Y,Z],H\rangle
		=\langle Y,[Z,H]\rangle.
	\]
	最后一等式由\eqref{g1:eq:invariant-inner-product-bracket}及括号反对称性得到。它对所有 \(Y\) 成立，特别取 \(Y=[Z,H]\)，正定性给出 \([Z,H]=0\)。正则性说明 \(Z\in\mathfrak z_{\mathfrak g}(H)=\mathfrak t\)，所以 \(g=g_0\) 满足要求。
\end{proof}

\begin{theorem}[极大环面的共轭定理]\label{g1:thm:maximal-tori-conjugate}
	设 \(G\) 为紧连通实 Lie 群。任意两个极大环面 \(T,S\subseteq G\) 共轭，即存在 \(g\in G\)，使 \(gSg^{-1}=T\)。因此所有极大环面的维数相等。
\end{theorem}
\begin{proof}
	记 \(\mathfrak t=\operatorname{Lie}(T)\)、\(\mathfrak s=\operatorname{Lie}(S)\)。在 \(\mathfrak s\) 中选择正则元素 \(X\)，使 \(\mathfrak z_{\mathfrak g}(X)=\mathfrak s\)。由\cref{g1:prop:adjoint-orbit-meets-torus}，有 \(g\in G\) 使 \(\operatorname{Ad}_gX\in\mathfrak t\)。由于 \(\mathfrak t\) 交换，得到
	\[
		\mathfrak t\subseteq\mathfrak z_{\mathfrak g}(\operatorname{Ad}_gX)
		=\operatorname{Ad}_g\mathfrak z_{\mathfrak g}(X)
		=\operatorname{Ad}_g\mathfrak s.
	\]
	中间的等式来自共轭保持括号。右端是交换 Lie 子代数，而\cref{g1:prop:maximal-torus-lie-centralizer}说明 \(\mathfrak t\) 已经是极大交换 Lie 子代数，所以这个包含只能是等号。

	两个环面都是连通 Lie 子群，且 \(gSg^{-1}\) 与 \(T\) 的 Lie 代数相同；\cref{g1:thm:lie-subalgebra-integration}中的连通积分子群唯一性给出 \(gSg^{-1}=T\)。也可直接用环面指数的满射性以及指数与共轭相容，写成 \(gSg^{-1}=\exp_G(\operatorname{Ad}_g\mathfrak s)=\exp_G\mathfrak t=T\)。共轭是微分同胚，故维数相同。
\end{proof}

\begin{definition}\label{g1:def:compact-group-rank}
	紧连通 Lie 群 \(G\) 的任一极大环面的维数称为 \(G\) 的\term[jinqunzhizhishu]{秩}[rank]，记为 \(\operatorname{rank}G\)。
\end{definition}

秩的选择无关性现在来自共轭定理。由\cref{g1:prop:adjoint-orbit-meets-torus}及\eqref{g1:eq:regular-element-centralizer}，对任意 \(X\in\mathfrak g\) 都有 \(\dim\mathfrak z_{\mathfrak g}(X)\ge\operatorname{rank}G\)，而正则元素恰使这个维数取最小值。因此可以对一般 \(X\in\mathfrak g\) 用最小中心化子维数定义正则性，与刚才在 \(\mathfrak t\) 中的定义一致。

\begin{proposition}[极大环面的中心化子]\label{g1:prop:maximal-torus-centralizer}
	设 \(G\) 为紧连通实 Lie 群，\(T\subseteq G\) 为极大环面。则 \(Z_G(T)=T\)。
\end{proposition}
\begin{proof}
	\(T\subseteq Z_G(T)\) 来自交换性。反过来，取 \(a\in Z_G(T)\)。由\cref{g1:thm:element-in-maximal-torus}，取包含 \(a\) 的极大环面 \(S\)。令 \(K=Z_G(a)\)，并取其单位连通分支 \(K^0\)。它是紧连通 Lie 群：\(K\) 是 \(G\) 的闭子群，而连通分支也闭。

	由于 \(a\) 与 \(T\) 中所有元素交换，\(T\subseteq K\)；由于 \(a\in S\) 且 \(S\) 交换，\(S\subseteq K\)。两环面连通且包含单位元，所以其实都包含于 \(K^0\)。它们在 \(K^0\) 中仍极大，因为任何更大的环面也会是 \(G\) 中更大的环面。将\cref{g1:thm:maximal-tori-conjugate}应用于紧连通群 \(K^0\)，得到 \(k\in K^0\) 使 \(kSk^{-1}=T\)。但 \(k\in Z_G(a)\)，故 \(a=kak^{-1}\in kSk^{-1}=T\)。
\end{proof}

等式 \(\operatorname{Lie}(Z_G(T))=\mathfrak t\) 确定的是中心化子的单位连通分支。证明进一步在 \(Z_G(a)^0\) 中使用共轭定理，排除了环面中心化子的额外分支。对单个元素的中心化子，额外分支仍可能真实存在。

\subsection{环面内共轭与正规化子}
\label{g1:subsec:1903:04}
共轭定理讨论整个子群能否相互移到一起。另一个问题是：已经放在同一个环面里的两个元素，什么时候仍在全群中共轭？环面交换，所以只用环面自身共轭不会改变任何元素，必须考虑那些把环面整体保持住、但不一定逐点固定它的变换。

\begin{definition}\label{g1:def:torus-normalizer}
	设 \(T\) 为 Lie 群 \(G\) 的环面。其\term[zhengguihuazi]{正规化子}[normalizer]定义为
	\[
		N_G(T)\coloneqq\{g\in G:gTg^{-1}=T\}.
	\]
	它通过共轭作用于 \(T\)，通过伴随作用于 \(\mathfrak t\)。中心化子 \(Z_G(T)\) 是这个共轭作用的核。\symidx{\(N_G(T),Z_G(T)\)：环面的正规化子与中心化子}
\end{definition}

\begin{proposition}[环面内共轭的判据]\label{g1:prop:torus-conjugacy-normalizer}
	设 \(G\) 为紧连通实 Lie 群，\(T\subseteq G\) 为极大环面。
	\begin{enumerate}
		\item 对 \(a,b\in T\)，它们在 \(G\) 中共轭，当且仅当存在 \(n\in N_G(T)\) 使 \(nan^{-1}=b\)。
		\item 对 \(X,Y\in\mathfrak t\)，它们在 \(\operatorname{Ad}(G)\) 下属于同一轨道，当且仅当存在 \(n\in N_G(T)\) 使 \(\operatorname{Ad}_nX=Y\)。
	\end{enumerate}
\end{proposition}
\begin{proof}
	1. 若 \(b=gag^{-1}\)，则 \(T\) 与 \(gTg^{-1}\) 都包含于 \(Z_G(b)\)：前者与 \(b\in T\) 交换，后者与 \(b\in gTg^{-1}\) 交换。这两个环面连通，故都在 \(K=Z_G(b)^0\) 中；它们仍是 \(K\) 的极大环面。\(K\) 紧且连通，所以共轭定理给出 \(k\in K\)，使 \(k(gTg^{-1})k^{-1}=T\)。令 \(n=kg\)，便有 \(n\in N_G(T)\)，并且
	\[
		nan^{-1}=kbk^{-1}=b.
	\]
	反向因 \(N_G(T)\subseteq G\) 直接成立。

	2. 若 \(Y=\operatorname{Ad}_gX\)，令
	\[
		K=\{h\in G:\operatorname{Ad}_hY=Y\}^0.
	\]
	括号中的集合是光滑作用的闭稳定子群，所以 \(K\) 是紧连通 Lie 群。由于 \(Y\in\mathfrak t\)，\(T\) 的伴随作用固定 \(Y\)；由于 \(X\in\mathfrak t\)，\(gTg^{-1}\) 的伴随作用也固定 \(Y\)。两环面均连通，所以都包含于 \(K\)，且仍极大。再选 \(k\in K\)，使 \(k(gTg^{-1})k^{-1}=T\)，并令 \(n=kg\)。则 \(n\in N_G(T)\) 且 \(\operatorname{Ad}_nX=\operatorname{Ad}_kY=Y\)。反向同样直接成立。
\end{proof}

因此，每个群元素都能先通过共轭放进固定环面 \(T\)，放进去以后剩余的共轭关系由 \(N_G(T)\) 决定；每个 Lie 代数向量也有对应的陈述。两种分类仍不相同，因为群元素还要经过指数核格的识别。例如在圆周群中，\(0\) 与 \(2\pi\mathrm{i}\) 是不同的 Lie 代数共轭类，却指数到同一个群元素。下一节将把正规化子中逐点固定 \(T\) 的部分除去，得到真正起作用的有限群。

% !TeX root = ../../Book1.tex
% 纲要标题：### 19.4 Weyl 群与根空间
% 纲要位置：几何学计划纲要.md，第 867 行。
% 正文状态：初稿；尚未正式审核。

% 纲要内容开始
% 
% * 正规化子模中心化子
% * 根空间与根的整性
% * 根平面、反射及 Weyl 室
% * 整数格、根格和整体群的区别
% 
% 纲要内容结束

\section{Weyl 群与根空间}
\label{g1:sec:1904}

把元素共轭入一个固定极大环面以后，仍可能有若干环面内的代表。例如，对角酉矩阵可以重新排列对角元；环面本身的共轭作用却是平凡的，不能记录这种排列。本节先把这些剩余对称组成有限群，再研究它为何与一组实超平面的反射有关。始终设 \(G\) 为紧连通实 Lie 群，\(T\) 为极大环面，\(\mathfrak t=\operatorname{Lie}(T)\)，并在 \(\mathfrak g\) 上固定一个伴随不变正定内积 \(B\)。

\subsection{正规化子模中心化子}
\label{g1:subsec:1904:01}
\begin{definition}\label{g1:def:weyl-group}
	设 \(G\) 为紧连通实 Lie 群，\(T\) 为其极大环面。其\term[weylqun]{Weyl 群}[Weyl group]是下列商群。\footnote{Hermann Weyl（1885--1955），德国数学家，在 Lie 群表示论与微分几何方面作出重要贡献。}
	\[
		W=W(G,T)\coloneqq N_G(T)/Z_G(T)=N_G(T)/T.
	\]
	最后一个等号来自\cref{g1:prop:maximal-torus-centralizer}。它通过共轭作用于 \(T\)，通过伴随作用于 \(\mathfrak t\)。\symidx{\(W(G,T)\)：极大环面的 Weyl 群}
\end{definition}

模掉中心化子的理由是它恰好逐点固定环面。若 \(n\) 与 \(nt\) 代表同一个商类，\(t\in T\) 的共轭作用在 \(T\) 上为恒等，故得到的作用不依赖代表。\(T=\exp\mathfrak t\) 又说明：一个元素逐点固定 \(\mathfrak t\)，当且仅当逐点固定 \(T\)。因此这两个作用都忠实。

\begin{proposition}[Weyl 群的有限性]\label{g1:prop:weyl-group-finite}
	设 \(G\) 为紧连通实 Lie 群，\(T\) 为极大环面，\(\mathfrak g=\operatorname{Lie}(G)\)、\(\mathfrak t=\operatorname{Lie}(T)\)。固定 \(\mathfrak g\) 上的 \(\operatorname{Ad}(G)\) 不变正定内积 \(B\)，并在 \(\mathfrak t\) 上取其限制。令 \(W(G,T)=N_G(T)/Z_G(T)\)。则 \(W(G,T)\) 是有限群，在 \(\mathfrak t\) 上忠实地作用为正交变换，并保持整数格 \(\Lambda_T=\ker(\exp_G|_{\mathfrak t})\)。
\end{proposition}
\begin{proof}
	正规化子是闭集：若 \(n_j\to n\) 且 \(n_jTn_j^{-1}=T\)，对每个 \(t\in T\) 取极限得到 \(ntn^{-1}\in T\)；再对 \(n_j^{-1}\) 使用同一论证，便得反向包含。故它是紧 Lie 子群。

	若 \(X\in\operatorname{Lie}(N_G(T))\)，则 \(\operatorname{Ad}_{\exp(sX)}\) 保持 \(\mathfrak t\)，微分得 \([X,H]\in\mathfrak t\) 对所有 \(H\in\mathfrak t\) 成立。另一方面，对于 \(K\in\mathfrak t\)，内积不变性给出
	\[
		B([X,H],K)=B(X,[H,K])=0.
	\]
	所以 \([X,H]\) 同时属于 \(\mathfrak t\) 及其正交补，只能为零。由\cref{g1:prop:maximal-torus-lie-centralizer}，\(X\in\mathfrak t\)。反向包含来自 \(T\subseteq N_G(T)\)，故商 \(N_G(T)/T\) 的维数为零。它既紧又离散，因而有限。

	忠实性已由作用的核说明；正交性来自 \(B\) 的伴随不变性。若 \(H\in\Lambda_T\)，则指数自然性给出
	\[
		\exp_G(\operatorname{Ad}_nH)=n\exp_G(H)n^{-1}=e.
	\]
	故 \(\operatorname{Ad}_n\) 保持 \(\Lambda_T\)，对 \(n^{-1}\) 再用一次得到集合相等。
\end{proof}

\cref{g1:prop:torus-conjugacy-normalizer}现在可以简洁地表述为：\(G\) 的共轭类由 \(T/W\) 参数化，\(\mathfrak g\) 的伴随轨道由 \(\mathfrak t/W\) 参数化。这里先指集合的一一对应；每个轨道与环面相交，而交点恰是同一 \(W\) 轨道。两种参数空间仍有指数核格的差别，不能把 \(T/W\) 与 \(\mathfrak t/W\) 混同。

\subsection{根空间与根的整性}
\label{g1:subsec:1904:02}
\cref{g1:prop:torus-weight-decomposition,g1:prop:maximal-torus-lie-centralizer}已经把互相交换的斜对称算子同时分解为
\begin{equation}\label{g1:eq:root-space-decomposition}
	\mathfrak g_{\mathbb C}=\mathfrak t_{\mathbb C}\oplus
	\bigoplus_{\alpha\in R}\mathfrak g_\alpha,
	\qquad
	[H,Z]=\mathrm{i}\alpha(H)Z\quad(Z\in\mathfrak g_\alpha).
\end{equation}
这里 \(R\subseteq\mathfrak t^*\setminus\{0\}\) 有限。伴随表示的这些非零无穷小权称为\term[gen]{根}[root]，\(\mathfrak g_\alpha\) 称为\term[genkongjian]{根空间}[root space]。我们使用实值的根，因而特征值中显式写出 \(\mathrm{i}\)。取共轭将 \(\mathfrak g_\alpha\) 变成 \(\mathfrak g_{-\alpha}\)；根从一开始就成正负对出现。\symidx{\(R,\mathfrak g_\alpha\)：根集与复根空间}

对于 \(Z\in\mathfrak g_\alpha\)、\(U\in\mathfrak g_\beta\)，Jacobi 恒等式给出
\[
	[H,[Z,U]]=[[H,Z],U]+[Z,[H,U]]
	=\mathrm{i}(\alpha+\beta)(H)[Z,U].
\]
因此 \([\mathfrak g_\alpha,\mathfrak g_\beta]\subseteq\mathfrak g_{\alpha+\beta}\)，其中约定 \(\mathfrak g_0=\mathfrak t_{\mathbb C}\)，不存在的权空间为零。权分解与 Lie 括号相容，正是根比一般表示的权更受约束的原因。下一步要证明每个非零根空间实际上只有一维；一般环面表示的权空间没有这个性质。

\ProseParagraphBreak
反复取括号会沿一个根方向升降特征值。引入相互配合的升、降与权算子，把它们的作用组织成有限链，就能控制根空间的维数，并从链的长度读出整数条件。

\begin{lemma}[升降算子的有限链]\label{g1:lem:sl2-weight-chain}
	设有限维复 Hermitian 空间 \(V\) 上的算子 \(E,F,L\) 满足
	\[
		[L,E]=2E,\qquad [L,F]=-2F,\qquad [E,F]=L,
		\qquad E^*=F,\quad L^*=L.
	\]
	则 \(V\) 是若干相互正交的不变子空间的直和。每个非零子空间有基
	\(v,Fv,\ldots,F^nv\)，其中 \(n\) 是非负整数，且
	\[
		Ev=0,\quad Lv=nv,\quad
		L(F^kv)=(n-2k)F^kv,\quad
		E(F^kv)=k(n-k+1)F^{k-1}v.
	\]
	特别地，\(L\) 的全部特征值都是整数。
\end{lemma}
\begin{proof}
	自伴算子 \(L\) 有正交特征分解。取最大特征值 \(\lambda\) 的非零特征向量 \(v\)。关系 \([L,E]=2E\) 表明 \(Ev\) 若非零就具有特征值 \(\lambda+2\)，所以 \(Ev=0\)。由交换关系归纳得到
	\[
		L(F^kv)=(\lambda-2k)F^kv,\qquad
		E(F^kv)=k(\lambda-k+1)F^{k-1}v.
	\]
	非零的 \(F^kv\) 具有互不相同的特征值，故有限维性保证存在最大整数 \(n\ge0\) 使 \(F^nv\ne0\)，而 \(F^{n+1}v=0\)。在第二式中取 \(k=n+1\)，得到
	\(0=(n+1)(\lambda-n)F^nv\)，故 \(\lambda=n\)。这些向量张成的空间在三个算子下都不变；由于 \(E^*=F\)、\(L^*=L\)，其正交补也不变。例如 \(u\) 在正交补内、\(w\) 在链空间内时，\(\langle Eu,w\rangle=\langle u,Fw\rangle=0\)。在正交补中重复，维数每次严格下降，最终得到所述分解。
\end{proof}

\begin{theorem}[根空间与根的整性]\label{g1:thm:compact-root-spaces}
	设 \(G\) 为紧连通实 Lie 群，\(T\) 为极大环面，\(\mathfrak g=\operatorname{Lie}(G)\)、\(\mathfrak t=\operatorname{Lie}(T)\)。固定 \(\mathfrak g\) 上的 \(\operatorname{Ad}(G)\) 不变正定实内积 \(B\)。采用\eqref{g1:eq:root-space-decomposition}的约定，对 \(\alpha\in\mathfrak t^*\) 令
	\[
		\mathfrak g_\alpha=\{Z\in\mathfrak g_{\mathbb C}:[H,Z]=\mathrm{i}\alpha(H)Z\text{ 对每个 }H\in\mathfrak t\},
		\qquad R=\{\alpha\ne0:\mathfrak g_\alpha\ne0\}.
	\]
	则每个 \(\mathfrak g_\alpha\) 的复维数为一。以 \(B\) 识别 \(\mathfrak t\) 与其对偶，并以 \((\alpha,\beta)\) 记诱导的对偶内积，则
	\[
		\frac{2(\beta,\alpha)}{(\alpha,\alpha)}\in\mathbb Z
		\qquad(\alpha,\beta\in R).
	\]
	若两根 \(\alpha,\beta\) 共线，则 \(\beta=\alpha\) 或 \(\beta=-\alpha\)。
\end{theorem}
\begin{proof}
	把 \(B\) 复双线性延拓为 \(B_{\mathbb C}\)，令 Hermitian 内积为 \(\langle Z,U\rangle=B_{\mathbb C}(Z,\overline U)\)。固定根 \(\alpha\)，取范数一的 \(Z\in\mathfrak g_\alpha\)，并记 \(A_\alpha\in\mathfrak t\) 为满足 \(B(A_\alpha,H)=\alpha(H)\) 的向量，\(a^2=(\alpha,\alpha)=B(A_\alpha,A_\alpha)>0\)。由括号的权相加性质，\([Z,\overline Z]\in\mathfrak t_{\mathbb C}\)；对 \(H\in\mathfrak t\)，计算
	\[
		B_{\mathbb C}([Z,\overline Z],H)
		=B_{\mathbb C}(Z,[\overline Z,H])
		=\mathrm{i}\alpha(H).
	\]
	所以 \([Z,\overline Z]=\mathrm{i}A_\alpha\)。令
	\begin{equation}\label{g1:eq:root-sl2-triple}
		e=\frac{\sqrt2}{a}Z,\qquad
		f=-\frac{\sqrt2}{a}\overline Z,\qquad
		h=-\frac{2\mathrm{i}}{a^2}A_\alpha.
	\end{equation}
	直接代入得 \([e,f]=h\)、\([h,e]=2e\)、\([h,f]=-2f\)。又有 \(\overline e=-f\)、\(\overline h=-h\)。内积不变性说明 \((\operatorname{ad}_u)^*=-\operatorname{ad}_{\overline u}\)，故 \(E=\operatorname{ad}_e\)、\(F=\operatorname{ad}_f\)、\(L=\operatorname{ad}_h\) 满足前一引理。在 \(\mathfrak g_\beta\) 上，\(L\) 的特征值是
	\[
		\frac{2\beta(A_\alpha)}{a^2}
		=\frac{2(\beta,\alpha)}{(\alpha,\alpha)},
	\]
	所以它是整数。

	为控制 \(\mathfrak g_\alpha\) 的维数，只看不变子空间
	\[
		V_\alpha=\mathfrak t_{\mathbb C}\oplus
		\bigoplus_{\substack{k\in\mathbb Z\setminus\{0\}\\k\alpha\in R}}
		\mathfrak g_{k\alpha}.
	\]
	它的 \(L\) 特征值都是偶数，零特征空间恰为 \(\mathfrak t_{\mathbb C}\)。升算子在这个零特征空间上满足 \(E(H)=-\mathrm{i}\alpha(H)e\)，所以其秩恰为一。引理中的每条非平凡偶数链都经过权零与权二，且从权零到权二的映射非零；因此这种链只能有一条。另一方面，\(\operatorname{span}_{\mathbb C}\{e,h,f\}\) 已是一条最高权为二的链，其正交补不变。于是余下的链全为平凡链。由此得到 \(\dim_{\mathbb C}\mathfrak g_\alpha=1\)，并且 \(k\alpha\notin R\) 对所有整数 \(|k|\ge2\) 成立。

	最后若 \(\beta=c\alpha\)、\(c>0\)，上面的整性分别用于 \((\beta,\alpha)\) 与 \((\alpha,\beta)\)，得到 \(2c\) 和 \(2/c\) 都是正整数；它们的乘积为四，故 \(c\in\{1/2,1,2\}\)。\(c=2\) 已被排除，\(c=1/2\) 则对根 \(\beta\) 使用同一结论排除。负倍数的情形先将 \(\beta\) 换为 \(-\beta\)，所以只剩 \(\pm\alpha\)。
\end{proof}

这里还可以把根与群元素的正则性联系起来。对 \(g\in G\)，若 \(\dim Z_G(g)=\operatorname{rank}G\)，则称该元素\term[qunyuansuzhengzexing]{正则}[regular]。对 \(t\in T\)，权分解给出
\[
	\operatorname{Lie}(Z_G(t))_{\mathbb C}
	=\mathfrak t_{\mathbb C}\oplus
	\bigoplus_{\substack{\alpha\in R\\\chi_\alpha(t)=1}}\mathfrak g_\alpha,
	\qquad \chi_\alpha(\exp H)=e^{\mathrm{i}\alpha(H)}.
\]
因此正则性等价于每个 \(\chi_\alpha(t)\ne1\)。这只控制中心化子的维数，即其单位分支；中心化子仍可能有其他分支。与 Lie 代数条件 \(\alpha(H)\ne0\) 相比，群元素 \(\exp H\) 要避开的是全部 \(\alpha(H)\in2\pi\mathbb Z\)。

\subsection{根平面、反射及 Weyl 室}
\label{g1:subsec:1904:03}
复根空间为一维，但原来的 Lie 代数是实空间。每一对 \(\{\alpha,-\alpha\}\) 恢复一个实二维平面
\[
	P_\alpha=(\mathfrak g_\alpha\oplus\mathfrak g_{-\alpha})\cap\mathfrak g,
	\qquad P_\alpha=P_{-\alpha},
\]
称为相应的\term[genpingmian]{根平面}[root plane]。取上一证明中范数一的 \(Z\in\mathfrak g_\alpha\)，则
\[
	u=\frac{Z+\overline Z}{\sqrt2},\qquad
	v=\frac{\mathrm{i}(Z-\overline Z)}{\sqrt2}
\]
构成 \(P_\alpha\) 的实正交单位基。这里 \(B_{\mathbb C}(Z,Z)=0\)：对取值 \(\alpha(H)\ne0\) 的 \(H\) 应用内积不变性，得到 \(2\mathrm{i}\alpha(H)B_{\mathbb C}(Z,Z)=0\)，再与 \(B_{\mathbb C}(Z,\overline Z)=1\) 合用便可核对正交与单位长度。括号为
\[
	[H,u]=\alpha(H)v,\qquad [H,v]=-\alpha(H)u.
\]
所以 \(\operatorname{Ad}_{\exp H}\) 在这个平面上旋转角度 \(\alpha(H)\)。选取每一对根中的一个组成集合 \(R_+\)，就有实正交分解
\[
	\mathfrak g=\mathfrak t\oplus\bigoplus_{\alpha\in R_+}P_\alpha,
	\qquad \dim_{\mathbb R}G=\operatorname{rank}G+2|R_+|.
\]
根平面在 \(\mathfrak g\) 内，而方程 \(\alpha(H)=0\) 的超平面在 \(\mathfrak t\) 内；它们是两个不同空间里的对象。前者表示被旋转的方向，后者表示旋转的无穷小速度为零的参数。

\ProseParagraphBreak
令
\begin{equation}\label{g1:eq:coroot-reflection}
	\alpha^\vee\coloneqq\frac{2A_\alpha}{(\alpha,\alpha)},
	\qquad
	s_\alpha(H)\coloneqq H-\alpha(H)\alpha^\vee.
\end{equation}
向量 \(\alpha^\vee\) 称为\term[yugen]{余根}[coroot]。\symidx{\(\alpha^\vee,s_\alpha\)：余根与根反射}由于 \(\alpha(\alpha^\vee)=2\)，\(s_\alpha\) 固定 \(\ker\alpha\)，并把它的正交方向 \(\alpha^\vee\) 变号，所以它就是关于 \(\ker\alpha\) 的正交反射。定义一个线性反射尚不足以证明它来自群的共轭作用；下面给出实现它的具体群元素。

\begin{proposition}[根反射的群内实现]\label{g1:prop:root-reflection-realized}
	设 \(G\) 为紧连通实 Lie 群，\(T\) 为极大环面，\(\mathfrak g=\operatorname{Lie}(G)\)、\(\mathfrak t=\operatorname{Lie}(T)\)。固定 \(\mathfrak g\) 上的 \(\operatorname{Ad}(G)\) 不变正定内积 \(B\)，以 \((\alpha,\beta)\) 记 \(B|_{\mathfrak t}\) 诱导的对偶内积。令 \(R\) 为所有满足下列条件的非零实泛函 \(\alpha\in\mathfrak t^*\) 的集合：存在 \(0\ne Z\in\mathfrak g_{\mathbb C}\)，使 \([H,Z]=\mathrm{i}\alpha(H)Z\) 对每个 \(H\in\mathfrak t\) 成立。对 \(\alpha\in R\)，以 \(B(H_\alpha,H)=\alpha(H)\)（\(H\in\mathfrak t\)）定义 \(H_\alpha\)，并令 \(\alpha^\vee=2H_\alpha/B(H_\alpha,H_\alpha)\)。令 \(s_\alpha(H)=H-\alpha(H)\alpha^\vee\)。则每个根 \(\alpha\) 的反射 \(s_\alpha\) 都属于 \(W(G,T)\) 在 \(\mathfrak t\) 上的作用。整个 \(W(G,T)\) 保持根集；其对根的作用为 \((w\alpha)(H)=\alpha(w^{-1}H)\)。特别地，
	\[
		s_\alpha\beta
		=\beta-\frac{2(\beta,\alpha)}{(\alpha,\alpha)}\alpha
		\in R\qquad(\alpha,\beta\in R).
	\]
\end{proposition}
\begin{proof}
	沿用\eqref{g1:eq:root-sl2-triple}中的 \(e,f,h\)，令
	\[
		x=e-f,\qquad y=\mathrm{i}(e+f),\qquad k=\mathrm{i}h=\alpha^\vee.
	\]
	它们在复共轭下不变，故属于实 Lie 代数，且括号满足
	\[
		[x,y]=2k,\qquad [y,k]=2x,\qquad[k,x]=2y.
	\]
	因此 \(\operatorname{ad}_x\) 在 \(\operatorname{span}_{\mathbb R}\{y,k\}\) 上的矩阵为
	\(\begin{psmallmatrix}0&-2\\2&0\end{psmallmatrix}\)。由\cref{g1:prop:exponential-naturality}，
	\[
		n_\alpha=\exp_G(\tfrac\pi2x)
		\quad\Longrightarrow\quad
		\operatorname{Ad}_{n_\alpha}k=-k.
	\]
	若 \(H\in\ker\alpha\)，则它与 \(e,f\) 都交换，因而与 \(x\) 交换，得到 \(\operatorname{Ad}_{n_\alpha}H=H\)。将一般 \(H\in\mathfrak t\) 写成
	\[
		H=\left(H-\frac{\alpha(H)}2k\right)+\frac{\alpha(H)}2k,
	\]
	括号内属于 \(\ker\alpha\)，从而 \(\operatorname{Ad}_{n_\alpha}H=H-\alpha(H)k=s_\alpha(H)\)。该作用保持 \(\mathfrak t\)，故指数自然性和 \(T=\exp\mathfrak t\) 给出 \(n_\alpha Tn_\alpha^{-1}=T\)。所以它确实代表一个 Weyl 群元素。

	对任意 \(n\in N_G(T)\)，令 \(w=\operatorname{Ad}_n|_{\mathfrak t}\)。若 \(Z\in\mathfrak g_\beta\)，则
	\[
		[H,\operatorname{Ad}_nZ]
		=\operatorname{Ad}_n[w^{-1}H,Z]
		=\mathrm{i}\beta(w^{-1}H)\operatorname{Ad}_nZ.
	\]
	因 \(\operatorname{Ad}_nZ\ne0\)，\(w\beta\) 仍是根。取 \(w=s_\alpha=w^{-1}\)，代入\eqref{g1:eq:coroot-reflection}，便得最后的公式。
\end{proof}

至此，根集已具有有限性、成对性、整性和反射不变性，且一条直线上只有一对根。在 \(\operatorname{span}_{\mathbb R}R\) 中，这些性质正组成一个约化的晶体根系。根也未必张成全部 \(\mathfrak t^*\)：其公共零空间恰为 \(\mathfrak g\) 的中心。事实上，\(H\in\mathfrak t\) 与每个根空间交换当且仅当全部 \(\alpha(H)=0\)，而零权空间本来就与 \(H\) 交换；反过来，中心元素先与 \(\mathfrak t\) 交换，故属于 \(\mathfrak t\)。

\begin{definition}\label{g1:def:weyl-chamber}
	设 \(T\) 为紧连通 Lie 群的极大环面，\(R\) 为其根集。开集
	\[
		\mathfrak t_{\mathrm{reg}}
		=\mathfrak t\setminus\bigcup_{\alpha\in R}\ker\alpha
	\]
	的各连通分支称为\term[weylshi]{Weyl 室}[Weyl chamber]。室的闭包中余维一的边界部分称为它的墙。
\end{definition}

选定 \(H_0\in\mathfrak t_{\mathrm{reg}}\)，令 \(R_+=\{\alpha:\alpha(H_0)>0\}\)。包含 \(H_0\) 的室就是
\[
	C=\{H:\alpha(H)>0\;\text{对全部 }\alpha\in R_+\}.
\]
一方面，一个连通分支内每个根的符号不能变化；另一方面，右侧是有限个开半空间的交，凸且连通，所以恰为这个分支。室是开凸多面锥；若存在中心方向，它沿这些方向不受限制。若 \(G\) 本身是环面，根集为空，唯一的室就是整个 \(\mathfrak t\)。

\begin{proposition}[相邻室与反射]\label{g1:prop:chamber-reflection-transitivity}
	设 \(G\) 为紧连通实 Lie 群，\(T\) 为极大环面。\(W(G,T)\) 置换 Weyl 室。由所有根反射生成的子群在这些室上作用传递；隔着同一墙相邻的两个室由关于该墙的反射交换。
\end{proposition}
\begin{proof}
	Weyl 群保持根集，所以保持根超平面的并，进而置换其补集的连通分支。设两个室沿 \(\ker\alpha\) 的一片墙相邻。在这片墙上取不属于其他不同根超平面的点 \(q\)，再取充分小的球，使球中只经过这一个根超平面。反射 \(s_\alpha\) 固定 \(q\)，交换球内的两个半球，又保持整个超平面排列；因此它把这两个半球所在的室相互交换。

	任取两个室 \(C,D\)，选 \(p\in C\)，再在 \(D\) 中选 \(q\)，使线段 \([p,q]\) 不经过两个不同根超平面的交。这样的 \(q\) 存在：对每个余维至少为二的交空间 \(L\)，所有可能使线段碰到 \(L\) 的端点都落在包含 \(p\) 与 \(L\) 的一个真仿射子空间内；有限个真仿射子空间不能覆盖开集 \(D\)。这条线段只在有限个互不相同的时刻逐个穿过根超平面，依次经过一列相邻室。每一次相邻转移由刚证明的一个根反射实现，将这些反射复合，就把 \(C\) 送到 \(D\)。根集为空时只有一个室，结论同样成立。
\end{proof}

\input{Figures/Chapter19-RootPlanes}

\cref{g1:fig:root-arrangements}画出秩二根系的两个模型。图中箭头把根借内积画成向量，虚线才是它们的核；位于根方向的箭头与同根对应的反射墙互相垂直。一般群的根反射与室之间的关系由上面的命题保证；下一节还会从矩阵直接算出经典群的整个 Weyl 群。

\subsection{整数格、根格和整体群的区别}
\label{g1:subsec:1904:04}
根来自伴随表示的微分，主要描述 Lie 代数。要恢复环面中的周期，还必须保留指数核
\(\Lambda_T\subseteq\mathfrak t\)。为了把这些不同层次的数据放在一起，先把环面的全部特征标写成线性泛函。

\begin{proposition}[环面的特征标格]\label{g1:prop:torus-character-lattice}
	设 \(T=\mathfrak t/\Lambda_T\) 为有限维紧连通交换 Lie 群。光滑群同态 \(\chi:T\to S^1\) 与集合
	\[
		X^*(T)\coloneqq\{\lambda\in\mathfrak t^*:
		\lambda(\Lambda_T)\subseteq2\pi\mathbb Z\}
	\]
	中的泛函一一对应，对应公式为 \(\chi_\lambda(\exp H)=e^{\mathrm{i}\lambda(H)}\)。该集合在加法下是 \(\mathfrak t^*\) 中的满秩格。\symidx{\(\Lambda_T,X^*(T)\)：环面的指数核格与特征标格}
\end{proposition}
\begin{proof}
	同态 \(\chi\) 的单位元微分是实线性映射 \(\mathfrak t\to\mathrm{i}\mathbb R\)，可唯一写成 \(H\mapsto\mathrm{i}\lambda(H)\)。指数自然性给出所述公式。对于 \(H\in\Lambda_T\)，必须有 \(e^{\mathrm{i}\lambda(H)}=1\)，这就是整性条件。反过来，满足该条件时，公式在 \(H\) 换成 \(H+L\)、\(L\in\Lambda_T\) 后不变，所以能下降到商环面，且是光滑群同态。指数的满射性保证对应的唯一性。

	取 \(\Lambda_T\) 的格基 \(L_1,\ldots,L_r\)，令 \(\lambda_j(L_k)=2\pi\delta_{jk}\)。上述条件恰好要求 \(\lambda\) 是 \(\lambda_1,\ldots,\lambda_r\) 的整数线性组合，故它们组成 \(X^*(T)\) 的格基。
\end{proof}

这里的特征标是一维酉表示；对于一般有限维表示，常把表示矩阵的迹也称为特征标，一维情形两者恰好一致。上一节已经得到每个根都满足整性，故有\term[genge]{根格}[root lattice]
\[
	Q\coloneqq\sum_{\alpha\in R}\mathbb Z\alpha
	\subseteq X^*(T).
\]
\(Q\) 位于对偶空间，而 \(\Lambda_T\) 位于原空间，不能直接比较包含关系。两者通过取值配对联系：\(\lambda(L)/(2\pi)\in\mathbb Z\)。根格在 \(\operatorname{span}_{\mathbb R}R\) 中是满秩格，因为它是满秩格 \(X^*(T)\) 的离散子群，并张成这个子空间；若有连续中心方向，根格的秩可能小于 \(\dim T\)。

\ProseParagraphBreak
余根属于原空间，因而可以与周期格直接比较；正确的系数是 \(2\pi\)。

\begin{proposition}[余根给出的周期]\label{g1:prop:coroot-period}
	设 \(G\) 为紧连通实 Lie 群，\(T\) 为极大环面，\(\mathfrak g=\operatorname{Lie}(G)\)、\(\mathfrak t=\operatorname{Lie}(T)\)。固定 \(\mathfrak g\) 上的 \(\operatorname{Ad}(G)\) 不变正定内积 \(B\)，以 \((\alpha,\beta)\) 记 \(B|_{\mathfrak t}\) 诱导的对偶内积。令 \(R\) 为所有满足下列条件的非零实泛函 \(\alpha\in\mathfrak t^*\) 的集合：存在 \(0\ne Z\in\mathfrak g_{\mathbb C}\)，使 \([H,Z]=\mathrm{i}\alpha(H)Z\) 对每个 \(H\in\mathfrak t\) 成立。对 \(\alpha\in R\)，以 \(B(H_\alpha,H)=\alpha(H)\)（\(H\in\mathfrak t\)）定义 \(H_\alpha\)，并令 \(\alpha^\vee=2H_\alpha/B(H_\alpha,H_\alpha)\)。根与余根采用\eqref{g1:eq:root-space-decomposition}及\eqref{g1:eq:coroot-reflection}的归一化，并记 \(\Lambda_T=\ker(\exp_G|_{\mathfrak t})\)。则
	\[
		\exp_G(2\pi\alpha^\vee)=e\qquad(\alpha\in R).
	\]
	因此若 \(Q^\vee=\sum_{\alpha\in R}\mathbb Z\alpha^\vee\)，则 \(2\pi Q^\vee\subseteq\Lambda_T\)。
\end{proposition}
\begin{proof}
	在\cref{g1:prop:root-reflection-realized}的证明中，实向量 \(x,y,k=\alpha^\vee\) 满足与下列矩阵相同的括号关系：
	\[
		x_0=\begin{pmatrix}0&1\\-1&0\end{pmatrix},\qquad
		y_0=\begin{pmatrix}0&\mathrm{i}\\\mathrm{i}&0\end{pmatrix},\qquad
		k_0=\begin{pmatrix}\mathrm{i}&0\\0&-\mathrm{i}\end{pmatrix}.
	\]
	它们是 \(\mathfrak{su}(2)\) 的实基，所以 \(x_0\mapsto x\)、\(y_0\mapsto y\)、\(k_0\mapsto k\) 定义了 Lie 代数同态 \(\varphi:\mathfrak{su}(2)\to\mathfrak g\)。这里需要一个整体同态来保持完整的周期，而不仅是局部指数公式。

	每个 \(\operatorname{SU}(2)\) 矩阵唯一写成
	\[
		\begin{pmatrix}z&w\\-\overline w&\overline z\end{pmatrix},
		\qquad |z|^2+|w|^2=1.
	\]
	这由两行酉正交及行列式等于一直接得到；取出 \(z,w\) 给出光滑逆映射，所以 \(\operatorname{SU}(2)\) 微分同胚于 \(S^3\)。由\cref{g1:prop:sphere-simply-connected}，它连通且单连通。\cref{g1:thm:lie-homomorphism-integration}遂把 \(\varphi\) 积分为 \(\Phi:\operatorname{SU}(2)\to G\)。再用指数自然性和对角矩阵指数，得到
	\[
		\exp_G(2\pi\alpha^\vee)
		=\Phi\bigl(\exp_{\operatorname{SU}(2)}(2\pi k_0)\bigr)
		=\Phi(I)=e.
	\]
	余根都在交换子代数 \(\mathfrak t\) 内，其整数线性组合的指数等于相应指数的乘积，故得到所列格包含。
\end{proof}

若根张成 \(\mathfrak t^*\)，还可定义\term[quange]{权格}[weight lattice]及\term[yuquange]{余权格}[coweight lattice]
\[
	P=\{\lambda\in\mathfrak t^*:\lambda(\alpha^\vee)\in\mathbb Z
	\;\text{对所有 }\alpha\in R\},\qquad
	P^\vee=\{H\in\mathfrak t:\alpha(H)\in\mathbb Z
	\;\text{对所有 }\alpha\in R\}.
\]
根张成假设保证这两者都是满秩格。的确，\(Q\) 已是满秩格，而 \(P^\vee\) 是它的整数对偶格；\(Q^\vee\subseteq\Lambda_T/(2\pi)\) 是离散且张成 \(\mathfrak t\) 的子群，故也是满秩格，\(P\) 是其整数对偶格。已有结论合成
\begin{equation}\label{g1:eq:root-lattice-inclusions}
	Q\subseteq X^*(T)\subseteq P,
	\qquad
	Q^\vee\subseteq\frac1{2\pi}\Lambda_T\subseteq P^\vee.
\end{equation}
\symidx{\(Q,P,Q^\vee,P^\vee\)：根格、权格、余根格与余权格}前一组右侧的包含是把 \(\lambda(\Lambda_T)\subseteq2\pi\mathbb Z\) 用于 \(2\pi\alpha^\vee\)；后一组右侧的包含则是把根的整性用于任意周期。若有中心方向，\(Q\)、\(Q^\vee\) 只张成相应的非中心子空间，不能把这些定义不加条件地称为整个空间的满秩格。

\ProseParagraphBreak
根还直接检测群中心，但需要在环面上取指数以后的值。

\begin{proposition}[由根检测中心]\label{g1:prop:center-from-roots}
	设 \(G\) 为紧连通实 Lie 群，\(T\) 为极大环面，\(\mathfrak g=\operatorname{Lie}(G)\)、\(\mathfrak t=\operatorname{Lie}(T)\)。令 \(R\) 为所有满足下列条件的非零实泛函 \(\alpha\in\mathfrak t^*\) 的集合：存在 \(0\ne Z\in\mathfrak g_{\mathbb C}\)，使 \([H,Z]=\mathrm{i}\alpha(H)Z\) 对每个 \(H\in\mathfrak t\) 成立。对 \(\alpha\in R\)，令根特征标 \(\chi_\alpha:T\to S^1\) 满足 \(\chi_\alpha(\exp_G H)=e^{\mathrm{i}\alpha(H)}\)，并记
	\[
		\Lambda_T=\ker(\exp_G|_{\mathfrak t}),\qquad
		P^\vee=\{H\in\mathfrak t:\alpha(H)\in\mathbb Z\text{ 对所有 }\alpha\in R\}.
	\]
	则
	\[
		Z(G)=\{t\in T:\chi_\alpha(t)=1\;\text{对所有 }\alpha\in R\}.
	\]
	若根张成 \(\mathfrak t^*\)，则还有自然群同构
	\[
		Z(G)\cong (2\pi P^\vee)/\Lambda_T.
	\]
\end{proposition}
\begin{proof}
	中心元素与 \(T\) 交换，故由 \(Z_G(T)=T\) 必在 \(T\) 中。对于 \(t\in T\)，其伴随作用在 \(\mathfrak t_{\mathbb C}\) 上恒等，在 \(\mathfrak g_\alpha\) 上乘以 \(\chi_\alpha(t)\)；所以所列根条件等价于 \(\operatorname{Ad}_t\) 在整个 \(\mathfrak g\) 上恒等。指数自然性于是说明 \(t\) 与所有指数交换，而连通群由指数生成，故 \(t\in Z(G)\)。这也证明反向包含。

	写 \(t=\exp H\)，根条件为 \(\alpha(H)\in2\pi\mathbb Z\)，即 \(H\in2\pi P^\vee\)。指数在此子群上的核仍为 \(\Lambda_T\)，所以得到最后的商群描述。
\end{proof}

例如，将 \(G\) 除以有限中心子群 \(C\)，不会改变 Lie 代数、根空间和根的配对，却会把周期格变成
\[
	\Lambda_{T/C}=\{H\in\mathfrak t:\exp_GH\in C\}\supseteq\Lambda_T.
\]
这里 \(C\subseteq T\) 已由中心判据保证。一个原来的特征标能下降到 \(T/C\)，当且仅当它在 \(C\) 上恒为一，故商群的特征标格反而缩小。下一节的 \(\operatorname{SU}(2)\) 与 \(\operatorname{SO}(3)\) 将把这一现象写成 \(2\pi\mathbb Z\) 与 \(\pi\mathbb Z\) 的具体区别。根记录括号和无穷小共轭，周期格记录哪些一参数运动已经闭合；分清这两类数据，才能讨论整体群的中心和覆盖。
% 纲要条目排为编号小节标题，并进入目录与书签；说明段仍保留于纲要注释中。

% !TeX root = ../../Book1.tex
% 纲要标题：### 19.5 经典紧群的完整算例
% 纲要位置：几何学计划纲要.md，第 874 行。
% 正文状态：初稿；尚未正式审核。

% 纲要内容开始
% 
% * 酉群和特殊酉群的环面与根
% * 正交群与辛群的根模型
% * 低秩同构、中心及基本群
% 
% ---
% 
% 纲要内容结束

\section{经典紧群的完整算例}
\label{g1:sec:1905}

抽象理论把共轭问题集中到一个环面，再把环面内剩余的识别交给 Weyl 群。对矩阵群，这两步分别对应选择正交或酉基，以及重新排列特征方向。本节把环面、根、整数格、中心和基本群放在同一组模型中计算。根描述伴随作用的无穷小旋转，整数格则记录哪些旋转已经回到单位元；两者同时出现，才不会把具有相同 Lie 代数的群误认为同一个群。

\subsection{酉群和特殊酉群的环面与根}
\label{g1:subsec:1905:01}
先看最熟悉的酉对角化。记 \(\mathfrak u(n)=\{X:X^*=-X\}\)，并令 \(E_{jk}\) 为第 \((j,k)\) 项等于一、其余项为零的矩阵。

\begin{proposition}[酉群的环面与根]\label{g1:prop:unitary-torus-roots}
	设 \(n\ge2\)。在 \(\operatorname{U}(n)\) 中，
	\[
		T=\{\operatorname{diag}(e^{\mathrm{i}\theta_1},\ldots,e^{\mathrm{i}\theta_n}):\theta_j\in\mathbb R\}
	\]
	是极大环面，其 Lie 代数用 \(H(\theta)=\mathrm{i}\operatorname{diag}(\theta_1,\ldots,\theta_n)\) 识别为 \(\mathbb R^n\)。若 \(\varepsilon_j(H(\theta))=\theta_j\)，则
	\[
		\Lambda=2\pi\mathbb Z^n,\qquad
		R=\{\varepsilon_j-\varepsilon_k:j\ne k\},\qquad
		W\cong S_n.
	\]
	在 \(\operatorname{SU}(n)\) 中，相应极大环面为 \(T_0=T\cap\operatorname{SU}(n)\)，Lie 代数及整数格分别是
	\[
		\mathfrak t_0=\{H(\theta):\textstyle\sum_j\theta_j=0\},\qquad
		\Lambda_0=\{2\pi(m_1,\ldots,m_n):m_j\in\mathbb Z,\ \textstyle\sum_jm_j=0\}.
	\]
	其根是上述根在 \(\mathfrak t_0\) 上的限制，Weyl 群仍为 \(S_n\)。
\end{proposition}
\begin{proof}
	酉矩阵有酉特征基。这里所需的谱论可以从一个单位特征向量开始归纳：复特征值存在，酉性使其模为一；该向量的正交补在矩阵及其逆下都不变，因而可继续选基。于是每个酉矩阵酉共轭于 \(T\)。改变对角元的角度给出从单位矩阵到它的路径，再作同一共轭，说明 \(\operatorname{U}(n)\) 连通。对于特殊酉矩阵，特征角之和属于 \(2\pi\mathbb Z\)；给其中一个角减去这个整数倍后，可使角度和恰为零，路径 \(\operatorname{diag}(e^{\mathrm{i}t\theta_j})\) 就始终在 \(\operatorname{SU}(n)\) 中。对角化所用的酉矩阵也可右乘一个对角酉矩阵，将其行列式调成一而不改变对角化。因此 \(\operatorname{SU}(n)\) 连通，且每个元素都在该群内共轭入 \(T_0\)。两群由矩阵方程定义为有界闭集，所以紧。

	\(T_0\) 可由前 \(n-1\) 个相位自由参数化，最后一个取其乘积的逆，故是连通环面。在两种情形，与所有这些对角矩阵交换的矩阵必须是对角矩阵：对第 \((j,k)\) 项选取第 \(j,k\) 个相位不同的环面元即可。于是两群中的中心化子分别就是 \(T,T_0\)，任何包含它们的环面也不能更大。这直接验证了极大性。

	\(\mathfrak u(n)\) 的复化是 \(\mathfrak{gl}(n,\mathbb C)\)，而 \(\mathfrak{su}(n)\) 的复化是 \(\mathfrak{sl}(n,\mathbb C)\)。逐项计算得
	\[
		[H(\theta),E_{jk}]=\mathrm{i}(\theta_j-\theta_k)E_{jk}.
	\]
	故每个非对角矩阵单位给出所列根，零权空间是对角子代数；取迹零只删去中心的一维零权方向。指数是逐项取指数，所以整数格正是所写的集合。

	正规化环面的酉矩阵必须置换其共同特征直线 \(\mathbb Ce_j\)，因而是一个置换矩阵与对角酉矩阵的乘积；反之这样的矩阵确实正规化环面。模掉对角部分只剩置换。在特殊酉群中，任意置换矩阵都能乘以对角酉矩阵将行列式调成一，仍实现相同的置换。因此两个商群都是 \(S_n\)。
\end{proof}

上述根系在迹零超平面内通常记作 \(A_{n-1}\)。一个 Weyl 室可取 \(\theta_1>\theta_2>\cdots>\theta_n\)，在特殊酉情形还要求和为零。相邻两角相等的墙记录两个特征方向不能再由伴随作用区分。\(\operatorname{U}(n)\) 还保留所有角同时增加的中心方向，所以它的根不张成整个 \(\mathfrak t^*\)。

\begin{example}\label{g1:exm:su2-so3-lattices}
	在 \(\operatorname{SU}(2)\) 中写 \(H(t)=\operatorname{diag}(\mathrm{i}t,-\mathrm{i}t)\)。正根为 \(\alpha(H(t))=2t\)，而 \(\exp H(t)=I\) 当且仅当 \(t\in2\pi\mathbb Z\)。商群 \(\operatorname{SU}(2)/\{\pm I\}\) 使用同一个 Lie 代数与同一条根，却在 \(t\in\pi\mathbb Z\) 时就回到单位元。因此根空间完全相同，整数格却扩大一倍。前一个环面的特征标为 \(e^{\mathrm{i}mt}\)，\(m\in\mathbb Z\)；能下降到商环面的特征标还必须在 \(t=\pi\) 取值一，即 \(m\) 必须为偶数。这个商群与 \(\operatorname{SO}(3)\) 的同构将在本节后面直接构造。
\end{example}

\subsection{正交群与辛群的根模型}
\label{g1:subsec:1905:02}
正交群的实特征方向未必是一维。非实特征值成共轭对，对应实二维旋转平面；这使它的环面参数成对出现。令
\[
	R(\theta)=\begin{pmatrix}\cos\theta&-\sin\theta\\\sin\theta&\cos\theta\end{pmatrix},\qquad
	J=\begin{pmatrix}0&-1\\1&0\end{pmatrix}.
\]

\begin{proposition}[旋转群的根模型]\label{g1:prop:orthogonal-torus-roots}
	令 \(J=\begin{pmatrix}0&-1\\1&0\end{pmatrix}\)。在 \(\operatorname{SO}(2r)\)（\(r\ge2\)）中取由旋转块 \(\exp(\theta_jJ)\) 组成的分块对角环面，在 \(\operatorname{SO}(2r+1)\)（\(r\ge1\)）中另加一个末尾的 \(1\) 块。以 \(H(\theta)=\operatorname{diag}(\theta_1J,\ldots,\theta_rJ)\) 表示其 Lie 代数，奇数维时另加零块，并置 \(\varepsilon_j(H(\theta))=\theta_j\)。这两个环面都是极大的，整数格均为 \(2\pi\mathbb Z^r\)，其根与 Weyl 群为
	\[
		\begin{array}{c|c|c}
		G&R&W\\ \hline
		\operatorname{SO}(2r+1)&\{\pm\varepsilon_j\pm\varepsilon_k:j<k\}\cup\{\pm\varepsilon_j\}&(\mathbb Z/2)^r\rtimes S_r\\
		\operatorname{SO}(2r)&\{\pm\varepsilon_j\pm\varepsilon_k:j<k\}&(\mathbb Z/2)^{r-1}\rtimes S_r.
		\end{array}
	\]
	偶数维一行中的符号变换要求改变偶数个坐标的符号。
\end{proposition}
\begin{proof}
	\cref{g1:thm:orthogonal-blocks}的实正交标准分块给出旋转块及 \(1,-1\) 块。行列式为一迫使 \(-1\) 块的数目为偶数，可将它们两两写成 \(R(\pi)\)；\(1\) 块也两两写成 \(R(0)\)，奇数维只留一个 \(1\)。若对角化基的定向不合要求，反转其中一根基向量即可，至多把一个旋转角变号。于是可在特殊正交群内共轭入所列环面。沿各角的线性路径回到零，说明特殊正交群连通；它由正交方程定义为有界闭集，所以紧。

	为计算伴随作用，把标准实表示复化。第 \(j\) 个旋转平面有权 \(\varepsilon_j,-\varepsilon_j\)，奇数维还有一个零权方向。记复双线性标准配对为 \(B\)。映射
	\[
		u\wedge v\longmapsto\bigl(w\longmapsto B(u,w)v-B(v,w)u\bigr)
	\]
	给出 \(\Lambda^2\mathbb C^m\cong\mathfrak{so}(m,\mathbb C)\)，其中 \(m=2r\) 或 \(2r+1\)。在标准基上，它把 \(e_j\wedge e_k\) 送到基本反对称矩阵，因此是线性同构；保持 \(B\) 的矩阵与此映射交换，故还是表示的同构。外积中两权相加，于是不同旋转平面贡献 \(\pm\varepsilon_j\pm\varepsilon_k\)，奇数维的零权方向再贡献 \(\pm\varepsilon_j\)。同一平面的正负权相加为零，共有 \(r\) 个零权方向，恰好构成 \(\mathfrak t_{\mathbb C}\)。因此与 \(\mathfrak t\) 交换的 Lie 代数就是 \(\mathfrak t\)，任何更大的环面都会给出更大的交换子代数，矛盾。

	正规化子置换这些实旋转平面，并可在各平面上反转定向，分别对应置换 \(\theta_j\) 及改变其符号。置换两个二维平面的行列式为一，反转一个平面的行列式为负。在奇数维，可同时反转剩余的一维轴以恢复总行列式一，所以任意符号变化都可实现；偶数维没有这根轴，必须反转偶数个平面。对平面排列及其定向作用平凡的部分就是各平面内的旋转，模掉环面后不再有其他元素。这证明了两个 Weyl 群公式。
\end{proof}

奇数维与偶数维的上述根模型分别记作 \(B_r,D_r\)。\(\operatorname{SO}(2)\cong S^1\) 则是无根的交换群，不能把它当作有非零根的普通 \(D_r\) 例子。\(D_2\) 的四个根分成两组互相正交的正负根，这将在 \(\operatorname{SO}(4)\) 的双覆盖中得到解释。

\ProseParagraphBreak

本节的\term[jinxinqun]{紧辛群}[compact symplectic group]用 \(\operatorname{Sp}(r)\) 表示，它是保持四元数 Hermitian 内积的群，与保持实反对称双线性型的非紧群 \(\operatorname{Sp}(2r,\mathbb R)\) 不同。四元数代数 \(\mathbb H\) 有实基 \(1,\mathbf i,\mathbf j,\mathbf k\)，满足 \(\mathbf i^2=\mathbf j^2=\mathbf k^2=-1\)、\(\mathbf i\mathbf j=\mathbf k=-\mathbf j\mathbf i\)。共轭把虚部取负，且 \(\overline{ab}=\bar b\bar a\)、\(a\bar a=|a|^2\)，故非零元的逆是 \(\bar a/|a|^2\)。把向量写成右 \(\mathbb H\) 向量空间中的列，定义
\[
	\operatorname{Sp}(r)=\{A\in\operatorname{Mat}_{r\times r}(\mathbb H):A^*A=I\},\qquad A^*=\overline A^{\mathsf T}.
\]

\begin{proposition}[紧辛群的根模型]\label{g1:prop:symplectic-torus-roots}
	设 \(r\ge1\)，\(\mathbb H\) 为四元数体，固定四元数单位 \(\mathbf i\)，并令 \(\operatorname{Sp}(r)=\{A\in\operatorname{Mat}_{r\times r}(\mathbb H):A^*A=I\}\)，其中 \(A^*\) 表示四元数共轭转置。则 \(\operatorname{Sp}(r)\) 是实维数 \(r(2r+1)\) 的紧连通 Lie 群。其极大环面可取 \(\operatorname{diag}(e^{\mathbf i\theta_1},\ldots,e^{\mathbf i\theta_r})\)，整数格为 \(2\pi\mathbb Z^r\)。令 \(H(\theta)=\operatorname{diag}(\mathbf i\theta_1,\ldots,\mathbf i\theta_r)\)，用 \(\varepsilon_j(H(\theta))=\theta_j\) 表示坐标权，则
	\[
		R=\{\pm\varepsilon_j\pm\varepsilon_k:j<k\}\cup\{\pm2\varepsilon_j\},\qquad
		W\cong(\mathbb Z/2)^r\rtimes S_r.
	\]
\end{proposition}
\begin{proof}
	映射 \(A\mapsto A^*A\) 的值在四元数 Hermitian 矩阵空间中。在 \(A^*A=I\) 处，微分是 \(V\mapsto A^*V+V^*A\)；对任意 Hermitian 矩阵 \(S\)，取 \(V=AS/2\) 就得到 \(S\)。因而正则值定理给出嵌入子流形结构。Hermitian 矩阵有 \(r\) 个实对角元及 \(r(r-1)/2\) 个任意四元数上三角元，实维数为 \(2r^2-r\)，故群的维数为 \(4r^2-(2r^2-r)=r(2r+1)\)。乘法和求逆光滑，紧性来自闭性及每列长度为一。

	四元数 Gram--Schmidt 与实、复情形采用同一投影步骤：从 \(v\) 中减去 \(\sum_j u_j\langle u_j,v\rangle\)，再除以正实范数。右乘系数的次序不能交换。这说明群传递地作用于 \(S^{4r-1}\)，第一列为指定单位向量的稳定子同构于 \(\operatorname{Sp}(r-1)\)。\cref{g1:thm:homogeneous-space,g1:prop:transitive-action-model}把此球面识别为稳定子的商，并给出局部截面，故取第一列是局部平凡丛。\(\operatorname{Sp}(1)=S^3\) 连通；归纳假定纤维连通，沿球面路径用有限个局部平凡化提升，再在终点纤维内连接，就能把任一矩阵连到单位矩阵。因此各 \(\operatorname{Sp}(r)\) 都连通。

	写四元数矩阵为 \(Z+W\mathbf j\)，其中 \(Z,W\) 是复矩阵。代入 \(\mathbf jz=\bar z\mathbf j\) 可验证
	\[
		Z+W\mathbf j\longmapsto
		\begin{pmatrix}Z&W\\-\bar W&\bar Z\end{pmatrix}
	\]
	保持乘法及共轭转置。故 Lie 代数可写成
	\[
		\begin{pmatrix}A&B\\-\bar B&\bar A\end{pmatrix},\qquad A^*=-A,\quad B^{\mathsf T}=B.
	\]
	环面在这个复矩阵模型中为 \(\operatorname{diag}(e^{\mathrm{i}\theta_1},\ldots,e^{\mathrm{i}\theta_r},e^{-\mathrm{i}\theta_1},\ldots,e^{-\mathrm{i}\theta_r})\)。\(A\) 的非对角项给出 \(\varepsilon_j-\varepsilon_k\)，\(B\) 的非对角对称项给出 \(\varepsilon_j+\varepsilon_k\)，其对角项给出 \(2\varepsilon_j\)；共轭块给出各负根。零权空间只有 \(A\) 的纯虚对角元，故所列环面极大。

	置换四元数坐标实现任意 \(S_r\) 作用，在某一坐标左乘单位四元数 \(\mathbf j\)，则由 \(\mathbf j e^{\mathbf i\theta}\mathbf j^{-1}=e^{-\mathbf i\theta}\) 实现单个符号反转。这些矩阵属于紧辛群。反过来，正规化子在标准复表示中只能置换成对的权线 \(\varepsilon_j,-\varepsilon_j\)，所以作用必是带符号的置换；作用平凡的部分是环面。这就确定了 Weyl 群。
\end{proof}

这组根记作 \(C_r\)。\(B_r\) 与 \(C_r\) 有相同的反射超平面与 Weyl 群，但轴向根分别是 \(\varepsilon_j\) 和 \(2\varepsilon_j\)，它们与斜向根的相对长度不同。只画反射墙会丢掉这一差别；比较根模型还必须保留根向量与所用坐标尺度。

\subsection{低秩同构、中心及基本群}
\label{g1:subsec:1905:03}
根与 Weyl 群已经能控制许多共轭计算，但它们尚未决定中心与覆盖关系。我们先构造低维覆盖，再用球面上的标架丛计算基本群，避免把覆盖次数从根图中猜出来。

\begin{proposition}[四元数给出的低维群映射]\label{g1:prop:quaternion-low-rank}
	有光滑群同构 \(\operatorname{Sp}(1)\cong\operatorname{SU}(2)\)，其共同底流形为 \(S^3\)。另有满射 Lie 群同态
	\[
		\operatorname{Sp}(1)\longrightarrow\operatorname{SO}(3),\qquad
		\operatorname{Sp}(1)\times\operatorname{Sp}(1)\longrightarrow\operatorname{SO}(4),\qquad
		\operatorname{Sp}(2)\longrightarrow\operatorname{SO}(5),
	\]
	其核依次为 \(\{\pm1\}\)、\(\{(1,1),(-1,-1)\}\)、\(\{\pm I\}\)，所以均为二重覆盖。
\end{proposition}
\begin{proof}
	单位四元数 \(z+w\mathbf j\) 对应
	\[
		\begin{pmatrix}z&w\\-\bar w&\bar z\end{pmatrix},\qquad |z|^2+|w|^2=1.
	\]
	其列正交且行列式为一；反过来，特殊酉二阶矩阵的第二行由第一行及行列式一唯一确定为这一形式。上面已经核对乘法相容，所以得到光滑群同构及 \(S^3\) 模型。

	单位四元数在三维虚四元数空间上作用为 \(v\mapsto qvq^{-1}\)。共轭保持虚部与范数，且从单位元连续变化，故像在 \(\operatorname{SO}(3)\)。若作用恒等，\(q\) 与 \(\mathbf i,\mathbf j,\mathbf k\) 都交换，逐项比较便知 \(q\) 是实数，故核为 \(\{\pm1\}\)。源、靶同为三维，离散核使单位元微分单射，因而为同构。像是开子群，又因源紧而闭；靶连通，所以像为整个群。\cref{g1:prop:lie-homomorphism-local-isomorphism}于是给出二重覆盖。也可直接令 \(q=\cos t+u\sin t\)，其中 \(u\) 是单位虚四元数：共轭固定 \(u\)，在 \(u^\perp\) 上旋转 \(2t\)。这解释了\cref{g1:exm:su2-so3-lattices}中的两倍角。

	在四维实空间 \(\mathbb H\) 上令 \((a,b)\) 作用为 \(x\mapsto axb^{-1}\)。它保持范数，且源连通，像在 \(\operatorname{SO}(4)\)。若作用恒等，代入 \(x=1\) 得 \(a=b\)，再代入所有四元数得 \(a=b=\pm1\)。两群维数均为六，刚才的开闭像论证再次适用，得到第二个覆盖。

	最后令 \(V\) 为迹零的二阶四元数 Hermitian 矩阵空间：
	\[
		V=\left\{\begin{pmatrix}r&q\\\bar q&-r\end{pmatrix}:r\in\mathbb R,\ q\in\mathbb H\right\},\qquad
		\|X\|^2=\tfrac12\operatorname{Re}\operatorname{tr}(X^2)=r^2+|q|^2.
	\]
	它是五维实 Euclid 空间。\(A\in\operatorname{Sp}(2)\) 的作用为 \(X\mapsto AXA^*\)。四元数满足 \(\operatorname{Re}(ab)=\operatorname{Re}(ba)\)，按分量求和便得到实部迹的循环性；因而这个作用保持迹零条件及所写范数。源连通，像落在 \(\operatorname{SO}(V)\)。核中的 \(A\) 必须与 \(\operatorname{diag}(1,-1)\) 交换，所以 \(A=\operatorname{diag}(a,d)\)；再与所有非对角 \(X\) 交换，得 \(aq=qd\) 对所有 \(q\in\mathbb H\) 成立。取 \(q=1,\mathbf i,\mathbf j\) 得 \(a=d=\pm1\)。两群维数都是十，同样得到满射二重覆盖。
\end{proof}

\begin{lemma}[球面上的标架丛与基本群]\label{g1:lem:sphere-bundle-pi1}
	设 \(K\to E\xrightarrow{p}S^m\) 是光滑主 \(K\) 丛，\(K\) 为连通 Lie 群，\(m\ge3\)。则纤维包含 \(K\hookrightarrow E\) 诱导基本群同构。
\end{lemma}
\begin{proof}
	取分别去掉南、北极的开集 \(U,V\)，二者都微分同胚于 \(\mathbb R^m\)，交集同伦等价于 \(S^{m-1}\)。丛在 \(U,V\) 上各有全局截面。为说明这一点，可按\cref{g1:thm:principal-connection-existence}取主联络，在 Euclid 坐标中从原点向每点作径向路径，再从固定纤维点水平提升；\cref{g1:thm:principal-horizontal-lift}保证提升在整个区间存在并光滑依赖终点，因此给出所需截面。这只用主联络沿紧区间的提升，不要求底空间完备。

	于是 \(p^{-1}(U)\) 与 \(p^{-1}(V)\) 都是可缩开集与 \(K\) 的乘积；交集用任一平凡化写成 \((U\cap V)\times K\)。因 \(m-1\ge2\)，\cref{g1:prop:sphere-simply-connected}说明 \(U\cap V\) 单连通。乘积中的路径与同伦逐分量进行，所以这三个空间的基本群都通过同一基点纤维与 \(\pi_1(K)\) 同构。两条包含映射在该纤维上就是恒等，故在这些识别下也都是恒等。\cref{g1:thm:van-kampen}将两个 \(\pi_1(K)\) 沿共同的同一 \(\pi_1(K)\) 粘合，所得群仍为 \(\pi_1(K)\)，且同构正由纤维包含诱导。
\end{proof}

\begin{proposition}[经典紧群的中心与基本群]\label{g1:prop:classical-centers-pi1}
	下表列出相应连通经典紧群的中心与基本群。其中 \(\mu_n=\{z\in S^1:z^n=1\}\)，\(r\ge1\)，并单独保留低维旋转群。
	\[
		\begin{array}{c|c|c}
		G&Z(G)&\pi_1(G)\\ \hline
		\operatorname{U}(n),\ n\ge1&S^1I&\mathbb Z\\
		\operatorname{SU}(n),\ n\ge2&\mu_nI&0\\
		\operatorname{Sp}(r)&\{\pm I\}&0\\
		\operatorname{SO}(2)&\operatorname{SO}(2)&\mathbb Z\\
		\operatorname{SO}(2r+1),\ r\ge1&\{I\}&\mathbb Z/2\\
		\operatorname{SO}(2r),\ r\ge2&\{\pm I\}&\mathbb Z/2.
		\end{array}
	\]
\end{proposition}
\begin{proof}
	酉或特殊酉群的中心元必须与其对角环面交换，所以为对角矩阵；再与混合任意两个坐标的特殊酉二阶块交换，迫使所有对角元相同。行列式条件分别给出 \(S^1I\) 与 \(\mu_nI\)。紧辛群中，先与复对角环面交换得到复对角矩阵，再与各坐标上的 \(\mathbf j\) 交换，迫使对角元为实数；置换坐标使它们相同，故只有 \(\pm I\)。对于 \(m\ge3\)，若实线性变换与所有坐标平面旋转交换，则对每个单位向量 \(v\)，它必须保持固定 \(v\) 的旋转子群的共同不动空间 \(\mathbb Rv\)。因此每条直线不变，由线性性是标量变换。正交性及行列式一分别给出奇数维的 \(I\) 和偶数维的 \(\pm I\)。\(\operatorname{SO}(2)\) 本来交换，须另行处理。

	取第一列给出主丛
	\[
		\operatorname{SU}(n-1)\longrightarrow\operatorname{SU}(n)\longrightarrow S^{2n-1},\qquad
		\operatorname{Sp}(r-1)\longrightarrow\operatorname{Sp}(r)\longrightarrow S^{4r-1}.
	\]
	第一条中的传递性由酉 Gram--Schmidt 得到：选定第一列后，用其余列之一调整行列式而不改变第一列。稳定子恰为低一阶的特殊酉群。第二条的传递性和稳定子已经在\cref{g1:prop:symplectic-torus-roots}中证明。它们均由齐性空间的局部截面构成主丛，符合\cref{g1:exm:homogeneous-principal-bundle}的条件。以 \(\operatorname{SU}(1)=\operatorname{Sp}(0)=\{1\}\) 为起点，反复应用上一引理，得到两类群的基本群都为零。\(\operatorname{SU}(2)\cong S^3\) 还给出独立的起始核验。

	映射
	\[
		S^1\times\operatorname{SU}(n)\longrightarrow\operatorname{U}(n),\qquad
		(z,A)\longmapsto\operatorname{diag}(z,1,\ldots,1)A
	\]
	是微分同胚：逆映射先取 \(z=\det U\)，再取 \(A=\operatorname{diag}(z^{-1},1,\ldots,1)U\)。它通常不是直积群同构，但已足够说明 \(\pi_1(\operatorname{U}(n))\cong\pi_1(S^1)=\mathbb Z\)，且行列式诱导此同构。

	\(S^3\to\operatorname{SO}(3)\) 是前一命题给出的二重覆盖；\(S^3\) 单连通，所以\cref{g1:thm:normal-covering-deck}将基本群识别为二元覆盖变换群，给出 \(\pi_1(\operatorname{SO}(3))\cong\mathbb Z/2\)。对 \(n\ge4\)，第一列又给出主丛
	\[
		\operatorname{SO}(n-1)\longrightarrow\operatorname{SO}(n)\longrightarrow S^{n-1}.
	\]
	正交标架可在不改变第一列的条件下调整定向，故作用传递，稳定子为 \(\operatorname{SO}(n-1)\)。球面维数至少三，上一引理说明包含映射诱导基本群同构，归纳便得全部 \(n\ge3\) 的结论。\(\operatorname{SO}(2)\) 的一整周旋转生成 \(\mathbb Z\)，进入 \(\operatorname{SO}(3)\) 后只保留其奇偶性。
\end{proof}

至此，\(\operatorname{Sp}(1)\times\operatorname{Sp}(1)\to\operatorname{SO}(4)\) 与 \(\operatorname{Sp}(2)\to\operatorname{SO}(5)\) 的源群也已证明单连通，因而是相应的万有覆盖。微分给出 \(\mathfrak{so}(4)\cong\mathfrak{su}(2)\oplus\mathfrak{su}(2)\) 及 \(\mathfrak{so}(5)\cong\mathfrak{sp}(2)\)，但群本身仍须除以所列二元中心核。根模型中的 \(D_2\) 分裂及 \(B_2,C_2\) 的低秩对应因此各有具体的矩阵实现。

\ProseParagraphBreak

这些算例显示三种资料各自回答不同问题。根给出共轭轨道的无穷小方向，Weyl 群给出同一环面内的有限识别，整数格与中心商则决定哪些环路真正闭合以及覆盖会怎样改变群。后一类信息不能从 Lie 代数同构中自动恢复。


\begin{chaptersummary}
	紧群上的归一化不变积分把任意内积平均为伴随不变内积，左平移后得到双不变度量。径向长度比较与紧性共同保证任意两点之间存在最短测地线；在双不变度量中，从单位元出发的测地线是一参数子群。因此紧连通 Lie 群的每个元素都可以写成 \(\exp X\)。

	\ProseParagraphBreak

	一参数子群的闭包是紧连通交换 Lie 群，而这样的群具有 \(\mathfrak t/\Lambda\) 的形式，其中 \(\Lambda\) 是满秩格。维数的有界性使每个环面包含在极大环面中。于是每个群元素属于某个极大环面；伴随轨道上的极值论证进一步证明所有极大环面共轭。固定一个极大环面后，群内的共轭问题便集中到这个环面。

	\ProseParagraphBreak

	中心化子说明哪些方向仍与给定对象交换，正规化子还允许重新安排环面内部的方向。有限商 \(W=N_G(T)/Z_G(T)\) 描述后一种作用。伴随表示的权分解把非交换方向组织为根空间，根超平面分隔正则方向，根反射由群内变换实现。群元素的正则条件涉及 \(e^{\mathrm{i}\alpha(H)}\)，Lie 代数元素的正则条件则涉及 \(\alpha(H)\)，两者不能直接混同。

	\ProseParagraphBreak

	酉群给出差根 \(\varepsilon_j-\varepsilon_k\)，旋转群给出 \(B_r,D_r\) 模型，紧辛群给出 \(C_r\) 模型。矩阵计算同时确定整数格、Weyl 群和中心。\(\operatorname{SU}(2)\to\operatorname{SO}(3)\) 等低维覆盖说明：微分同构保留 Lie 代数，却可能丢失离散中心核。特征标能否下降、环路绕行多少次及群是否单连通，都需要进一步保留这些整体信息。
\end{chaptersummary}

\BookExercises

以下习题从平均与最短路径开始，逐步转到环面的闭包、共轭、根与经典群计算。涉及指数时分别检查 Lie 代数中的等式和模整数格后的群等式；涉及共轭时明确允许的共轭群。

\begin{exercise}\label{g1:ex:unitary-conjugation-average}
	对任意复矩阵 \(A\in M_n(\mathbb C)\)，令
	\[
		P(A)=\int_{\operatorname{U}(n)}UAU^{-1}\,\mathrm{d}\mu(U).
	\]
	证明 \(P(A)=\frac{\operatorname{tr}A}{n}I\)。进而求 \(\int_{\operatorname{U}(n)}|U_{11}|^2\,\mathrm{d}\mu(U)\)，并在 \(n>1\) 时举例说明，共轭平均不保持矩阵的秩。
\end{exercise}
\begin{solution}
	对任意 \(V\in\operatorname{U}(n)\)，左不变换元 \(U\mapsto VU\) 给出 \(VP(A)V^{-1}=P(A)\)。若一个矩阵与所有对角酉矩阵交换，则其每个非对角元均为零：独立改变第 \(j\) 与第 \(k\) 个对角相位，会把该位置乘以任意单位复数，而不变性要求原值不变。故 \(P(A)\) 对角。再与所有置换矩阵交换，得到各对角元相同，所以 \(P(A)=cI\)。迹与有限分量积分交换，且 \(\operatorname{tr}(UAU^{-1})=\operatorname{tr}A\)，于是 \(nc=\operatorname{tr}A\)，证明公式。

	令 \(A=E_{11}\)。因 \(U^{-1}=U^*\)，矩阵 \(UE_{11}U^*\) 的第 \((1,1)\) 项为 \(|U_{11}|^2\)，而平均为 \(I/n\)，故所求积分是 \(1/n\)。\(E_{11}\) 的秩为 \(1\)，但 \(P(E_{11})=I/n\) 的秩为 \(n\)；当 \(n>1\) 时二者不同。
\end{solution}

\begin{exercise}\label{g1:ex:bi-invariant-curvature}
	设有限维实 Lie 群 \(G\) 配备双不变度量。采用本书的曲率约定，证明对单位元切向量 \(X,Y,Z\)，有
	\[
		R(X,Y)Z=-\tfrac14[[X,Y],Z].
	\]
	当 \(X,Y\) 线性无关时，由它们张成平面的\term[jiemianqulv]{截面曲率}[sectional curvature]定义为
	\[
		K(X,Y)=\frac{\langle R(X,Y)Y,X\rangle}{|X|^2|Y|^2-\langle X,Y\rangle^2}.
	\]
	求此处的截面曲率，证明其非负，并证明该度量平坦当且仅当 \(\mathfrak g\) 交换。最后说明没有连通性时不能由平坦推出 \(G\) 交换。
\end{exercise}
\begin{solution}
	用左不变场延拓 \(X,Y,Z\)。由\eqref{g1:eq:bi-invariant-connection}和曲率定义，
	\[
		R(X,Y)Z
		=\tfrac14[X,[Y,Z]]-\tfrac14[Y,[X,Z]]-\tfrac12[[X,Y],Z].
	\]
	Jacobi 恒等式使前两项之和为 \(\tfrac14[[X,Y],Z]\)，故得到所求公式。记 \(A=[X,Y]\)。伴随不变内积的斜伴随恒等式给出
	\[
		\langle[A,Y],X\rangle
		=-\langle[Y,A],X\rangle
		=\langle A,[Y,X]\rangle=-|A|^2.
	\]
	所以
	\[
		K(X,Y)=\frac{|[X,Y]|^2}
		{4\bigl(|X|^2|Y|^2-\langle X,Y\rangle^2\bigr)}\ge0.
	\]
	若 \(\mathfrak g\) 交换，曲率公式直接给出 \(R=0\)。反之，若 \(R=0\)，对任意线性无关的 \(X,Y\) 有 \([X,Y]=0\)；线性相关时括号也为零，故整个 Lie 代数交换。取非交换有限群 \(F\)，并令 \(G=S^1\times F\)。它的每个分支都是标准圆周，可由 \(\mathfrak g=\mathbb R\) 上的标准内积得到双不变平坦度量，但群乘法包含非交换的 \(F\)，所以 \(G\) 不交换。
\end{solution}

\begin{exercise}\label{g1:ex:compact-group-power-map}
	设 \(G\) 为有限维实紧连通 Lie 群，\(m\ge1\) 为整数。证明幂映射 \(P_m:G\to G\)、\(P_m(g)=g^m\)，满射。用 \(\operatorname{SU}(2)\) 中的具体矩阵说明 \(P_2\) 未必是群同态，并用 \(\operatorname{O}(2)\) 说明去掉连通性后 \(P_2\) 未必满射。
\end{exercise}
\begin{solution}
	给定 \(g\in G\)，\cref{g1:thm:compact-exp-surjective}给出 \(X\in\mathfrak g\) 使 \(g=\exp_GX\)。令 \(h=\exp_G(X/m)\)，单参数子群规律给出 \(h^m=\exp_GX=g\)，因此满射。

	在 \(\operatorname{SU}(2)\) 中取
	\[
		A=\begin{pmatrix}\mathrm{i}&0\\0&-\mathrm{i}\end{pmatrix},
		\qquad B=\begin{pmatrix}0&1\\-1&0\end{pmatrix}.
	\]
	二者均酉且行列式为 \(1\)，并满足 \(A^2=B^2=-I\)、\(AB=-BA\)。所以 \((AB)^2=-A^2B^2=-I\)，但 \(A^2B^2=I\)，幂映射不保持乘法。在 \(\operatorname{O}(2)\) 中，任意 \(Q\) 均有 \(\det(Q^2)=1\)，故行列式为 \(-1\) 的反射没有平方根。连通性在满射命题中不能省略。
\end{solution}

\begin{exercise}\label{g1:ex:flat-torus-shortest-logarithms}
	在 \(T^r=\mathbb R^r/\mathbb Z^r\) 上取由正定对称矩阵 \(Q\) 给出的常系数度量，记 \(|v|_Q^2=v^{\mathsf T}Qv\)。证明
	\[
		d(0,x+\mathbb Z^r)=\min_{z\in\mathbb Z^r}|x-z|_Q.
	\]
	当 \(r=2\)、\(Q=\left(\begin{smallmatrix}1&9/10\\9/10&1\end{smallmatrix}\right)\)、\(x=(49/100,49/100)\) 时，求所有最短测地线的初速度，并说明为什么逐坐标选取 \([-1/2,1/2]\) 中的代表元未必得到最短对数。
\end{exercise}
\begin{solution}
	商映射是局部等距映射，且指数映射为 \(v\mapsto v+\mathbb Z^r\)。\cref{g1:thm:compact-exp-surjective}给出一条初速度 \(v\) 的最短测地线；其终点条件为 \(v=x-z\)，其中 \(z\in\mathbb Z^r\)，故 \(d(0,x+\mathbb Z^r)=|x-z|_Q\)。反过来，每个这样的 \(x-z\) 都给出一条连接两点的直线投影，其长度为 \(|x-z|_Q\)，所以距离不大于任何一个候选值，证明最小值公式。也可直接从 \(Q\) 正定看出候选值在 \(|z|\to\infty\) 时趋于无穷，因此最小值确能取得。

	写 \(z=(m,n)\)，并令 \(u=49/100-m\)、\(v=49/100-n\)。则
	\[
		|(u,v)|_Q^2
		=\frac{19}{20}(u+v)^2+\frac1{20}(u-v)^2
		=\frac{19}{20}\bigl(49/50-m-n\bigr)^2
		+\frac1{20}(n-m)^2.
	\]
	若 \(m+n\) 为偶数，则 \(|49/50-m-n|\ge49/50\)，故平方长度至少为 \(45619/50000\)，在 \(m=n=0\) 时取得。若 \(m+n\) 为奇数，则 \(|49/50-m-n|\ge1/50\) 且 \(|n-m|\ge1\)，故平方长度至少为 \(2519/50000\)。等号要求 \(m+n=1\)、\(|n-m|=1\)，恰好给出 \((m,n)=(1,0)\) 或 \((0,1)\)。因此最短测地线的初速度恰有两个：
	\[
		(-51/100,49/100),\qquad(49/100,-51/100).
	\]
	原代表元 \((49/100,49/100)\) 的每个坐标都在指定半开或闭区间内，却具有较大长度。原因是 \(Q\) 的交叉项使两个坐标方向并不正交，逐坐标最小化不能代替按完整二次型最小化。
\end{solution}

\begin{exercise}\label{g1:ex:torus-integer-homomorphism}
	把 \(n\) 维环面写为 \(T^n=\mathbb R^n/\mathbb Z^n\)。证明每个光滑群同态 \(F:T^n\to T^m\) 都唯一写成
	\[
		F(x+\mathbb Z^n)=Ax+\mathbb Z^m,
	\]
	其中 \(A\) 是一个 \(m\times n\) 整数矩阵。进而证明 \(T^n\) 的 Lie 群自同构恰对应于 \(\operatorname{GL}(n,\mathbb Z)\)，并说明一个可逆实矩阵未必诱导环面自同构。
\end{exercise}
\begin{solution}
	商环面的 Lie 代数为 \(\mathbb R^n\)，其指数映射为 \(x\mapsto x+\mathbb Z^n\)。令 \(A=(\mathrm{d}F)_0\)，指数自然性给出
	\[
		F(x+\mathbb Z^n)=Ax+\mathbb Z^m.
	\]
	若 \(z\in\mathbb Z^n\)，左侧为单位元，故 \(Az\in\mathbb Z^m\)。对每个标准基向量应用此条件，可知 \(A\) 的每一列为整数列。反之，整数矩阵满足 \(A\mathbb Z^n\subseteq\mathbb Z^m\)，所以公式与代表元无关，且由线性映射下降成光滑群同态。其微分就是 \(A\)，因此矩阵唯一。

	若 \(F\) 为 Lie 群自同构，\(F^{-1}\) 同样由整数矩阵 \(B\) 给出，链式法则使 \(AB=BA=I\)。因而 \(\det A\)、\(\det B\) 都是整数且乘积为 \(1\)，故 \(\det A=\pm1\)。反之，整数矩阵若行列式为 \(\pm1\)，伴随矩阵公式给出整数逆矩阵，遂诱导逆同态。故自同构群恰为 \(\operatorname{GL}(n,\mathbb Z)\)。例如一维的 \(A=\sqrt2\) 不保持整数格，连良定义的商映射也不给出；\(A=2\) 虽然可逆且保持整数格，但诱导的圆周映射有两个原像，因而不是自同构。
\end{solution}

\begin{exercise}\label{g1:ex:three-torus-flow-closure}
	在 \(T^3=(S^1)^3\) 中考虑单参数子群
	\[
		c(t)=(e^{\mathrm{i}t},e^{\mathrm{i}\sqrt2t},e^{\mathrm{i}(1+\sqrt2)t}).
	\]
	求它的像的闭包及闭包的 Lie 代数，并比较该 Lie 代数与 \(c'(0)\) 张成的直线。
\end{exercise}
\begin{solution}
	每个 \(c(t)\) 都满足 \(z_3=z_1z_2\)，所以闭包包含于闭子群
	\[
		A=\{(z_1,z_2,z_1z_2):z_1,z_2\in S^1\}.
	\]
	映射 \(j:T^2\to A\)、\(j(z_1,z_2)=(z_1,z_2,z_1z_2)\) 是 Lie 群同构，逆为前两个坐标的投影。由\cref{g1:exm:irrational-torus-subgroup}，\(t\mapsto(e^{\mathrm{i}t},e^{\mathrm{i}\sqrt2t})\) 在 \(T^2\) 中稠密，因为 \(\sqrt2\) 无理。将其送入同胚 \(j\)，可知 \(c(\mathbb R)\) 在 \(A\) 中稠密。因此所求闭包恰为 \(A\)。

	在单位元处微分 \(j\)，得到
	\[
		\operatorname{Lie}(A)=\{(\mathrm{i}u,\mathrm{i}v,\mathrm{i}(u+v)):u,v\in\mathbb R\}.
	\]
	这是二维实平面，而 \(c'(0)=\mathrm{i}(1,\sqrt2,1+\sqrt2)\) 只张成其中一条直线。闭包增加了一维的无穷小方向，尽管所有原像点都由同一个实参数产生。
\end{solution}

\begin{exercise}\label{g1:ex:commuting-elements-no-common-torus}
	在 \(\operatorname{SO}(3)\) 中令
	\[
		A=\operatorname{diag}(1,-1,-1),\qquad
		B=\operatorname{diag}(-1,1,-1).
	\]
	证明 \(A\) 与 \(B\) 交换，但不存在同时包含二者的环面。说明这与“每个元素属于某个极大环面”并不矛盾。
\end{exercise}
\begin{solution}
	两矩阵都正交且行列式为 \(1\)，作为对角矩阵彼此交换。它们满足 \(A^2=B^2=I\)，且 \(I,A,B,AB\) 是四个不同元素。

	先看绕第三轴旋转的标准圆周 \(T\subseteq\operatorname{SO}(3)\)。若 \(C\) 与 \(T\) 的所有元素交换，便保持 \(T\) 的公共不动轴 \(\mathbb Re_3\) 及其正交平面。它在平面上的限制与所有平面旋转交换。写一个实 \(2\times2\) 矩阵并令它与 \(\begin{psmallmatrix}0&-1\\1&0\end{psmallmatrix}\) 交换，可得它形如 \(\begin{psmallmatrix}a&-b\\b&a\end{psmallmatrix}\)；正交性给出 \(a^2+b^2=1\)，故这是行列式为 \(1\) 的平面旋转。\(C\) 在第三轴上的因子由 \(\det C=1\) 也等于 \(1\)。所以 \(Z_{\operatorname{SO}(3)}(T)=T\)，任何包含 \(T\) 的环面都只能等于 \(T\)。这证明 \(T\) 极大。

	\(\operatorname{SO}(3)\) 连通：每个三维正交行列式为 \(1\) 的矩阵都有特征值 \(1\)，并在其正交平面上是旋转，连续缩小旋转角便给出到单位元的道路。因此极大环面共轭定理适用，所有极大环面都是圆周，任意环面又包含于一个极大环面。如果一个环面同时包含 \(A,B\)，相应圆周便至少有四个满足 \(x^2=e\) 的元素。然而 \(S^1\) 中方程 \(z^2=1\) 只有 \(1,-1\) 两解，矛盾。逐个元素的存在定理允许为不同元素选择不同环面，并没有给出有限个交换元素的共同环面。
\end{solution}

\begin{exercise}\label{g1:ex:center-intersection-maximal-tori}
	设 \(G\) 为紧连通 Lie 群。证明其中心恰为全部极大环面的交：
	\[
		Z(G)=\bigcap_{T\text{ 为极大环面}}T.
	\]
	由此说明，即使 \(G\) 连通，这个交也未必连通，并给出具体例子。
\end{exercise}
\begin{solution}
	若 \(z\in Z(G)\)，则 \(z\) 与任一极大环面 \(T\) 的所有元素交换。由\cref{g1:prop:maximal-torus-centralizer}，\(z\in Z_G(T)=T\)，所以中心包含于所列交集。

	反过来，若 \(z\) 属于每个极大环面，任取 \(g\in G\)。由\cref{g1:thm:element-in-maximal-torus}，可选包含 \(g\) 的极大环面 \(T_g\)。假设也给出 \(z\in T_g\)，由交换性得 \(zg=gz\)。\(g\) 任意，故 \(z\in Z(G)\)。

	例如 \(\operatorname{SU}(2)\) 的中心为 \(\{I,-I\}\)。可直接核对：与全部对角矩阵 \(\operatorname{diag}(e^{\mathrm{i}t},e^{-\mathrm{i}t})\) 交换的矩阵必须对角；这样的矩阵若再与 \(\begin{psmallmatrix}0&1\\-1&0\end{psmallmatrix}\) 交换，两对角元必须相等，行列式为 \(1\) 便给出 \(\pm I\)。反向这两个标量矩阵确实与全群交换。\(\operatorname{SU}(2)\) 同胚于 \(S^3\)，是连通群，而两个不同点组成的中心不连通。
\end{solution}

\begin{exercise}\label{g1:ex:compact-power-map}
	设 \(G\) 为紧连通 Lie 群，\(m\) 为正整数。证明一个 \(r\) 维环面中每个元素恰有 \(m^r\) 个 \(m\) 次方根。一般紧连通群中的根数是否必然有限？用 \(\operatorname{SO}(3)\) 中的平方根说明。
\end{exercise}
\begin{solution}
	在 \(T\cong\mathbb R^r/\mathbb Z^r\) 中，求 \(x\) 的 \(m\) 次方根就是求 \(my=x\pmod{\mathbb Z^r}\)。固定一个代表元 \(x\in\mathbb R^r\)，全部解为
	\[
		y=\frac{x+k}{m}+\mathbb Z^r,
		\qquad k\in\{0,1,\ldots,m-1\}^r.
	\]
	任意整数向量可模 \(m\) 化成上述唯一代表；不同代表的差除以 \(m\) 不在 \(\mathbb Z^r\)，故这些解彼此不同，恰有 \(m^r\) 个。

	全群中的根数却可以无限。\(\operatorname{SO}(3)\) 中，绕任意一条过原点的轴旋转 \(\pi\) 都满足 \(R^2=I\)，不同的无向轴给出不同矩阵。事实上若 \(R^2=I\)，则 \(R^{-1}=R\) 且正交性给出 \(R^{\mathsf T}=R\)。实对称谱分解使特征值全为 \(\pm1\)，再由行列式为 \(1\)，可知 \(R=I\)，或 \(R\) 在一维空间上为 \(1\)、其正交平面上为 \(-1\)。后一种恰是绕该直线的半周旋转。因此 \(I\) 的平方根集合由孤立点 \(I\) 及一个以 \(\mathbb{RP}^2\) 为参数的无限族组成；不同环面中的根不能按单个环面的有限计数合并。
\end{solution}

\begin{exercise}\label{g1:ex:root-height-a2}
	在 \(\operatorname{SU}(3)\) 的对角子代数中写 \(H=\mathrm{i}\operatorname{diag}(\theta_1,\theta_2,\theta_3)\)，其中 \(\theta_1+\theta_2+\theta_3=0\)。设 \(\alpha=\varepsilon_1-\varepsilon_2\)、\(\beta=\varepsilon_2-\varepsilon_3\)。用 \(B(X,Y)=-\operatorname{tr}(XY)\) 计算两个余根和对应反射；描述含 \(H_0=\mathrm{i}\operatorname{diag}(2,0,-2)\) 的 Weyl 室，并把 \(\mathrm{i}\operatorname{diag}(-3,2,1)\) 用反射移到该室。
\end{exercise}
\begin{solution}
	在角度坐标中，内积是 \(\mathbb R^3\) 的普通内积在和为零平面上的限制。因此
	\[
		A_\alpha=(1,-1,0),\quad A_\beta=(0,1,-1),\quad
		(\alpha,\alpha)=(\beta,\beta)=2,
	\]
	所以 \(\alpha^\vee=(1,-1,0)\)、\(\beta^\vee=(0,1,-1)\)。反射公式给出
	\[
		s_\alpha(\theta_1,\theta_2,\theta_3)=(\theta_2,\theta_1,\theta_3),\qquad
		s_\beta(\theta_1,\theta_2,\theta_3)=(\theta_1,\theta_3,\theta_2).
	\]
	六根为 \(\pm\alpha,\pm\beta,\pm(\alpha+\beta)\)。在 \(H_0\) 上三个所列正根都取正值，故相应室由 \(\theta_1>\theta_2>\theta_3\)、\(\sum\theta_j=0\) 定义。角度向量 \((-3,2,1)\) 先经 \(s_\alpha\) 变为 \((2,-3,1)\)，再经 \(s_\beta\) 变为 \((2,1,-3)\)，后者属于所求室。整个过程保持角度和与 Euclid 范数，并对应群内共轭。
\end{solution}

\begin{exercise}\label{g1:ex:b2-coroots}
	在 Euclid 平面 \(\mathbb R^2\) 中考虑根集
	\[
		R=\{\pm\varepsilon_1,\pm\varepsilon_2,
		\pm(\varepsilon_1+\varepsilon_2),\pm(\varepsilon_1-\varepsilon_2)\}.
	\]
	计算 \(Q,Q^\vee,P,P^\vee\)，并求 \(P/Q\) 与 \(P^\vee/Q^\vee\)。说明短根与长根取余根以后长度次序如何改变。
\end{exercise}
\begin{solution}
	通过标准内积识别向量与协向量。短根 \(\pm\varepsilon_j\) 的余根为 \(\pm2e_j\)；长根 \(\pm(\varepsilon_1\pm\varepsilon_2)\) 的余根为 \(\pm(e_1\pm e_2)\)。因此
	\[
		Q=\mathbb Z^2,\qquad
		Q^\vee=\{(m,n)\in\mathbb Z^2:m+n\text{ 为偶数}\}.
	\]
	后一个等式可由生成元 \((1,1),(1,-1)\) 看出：它们的整数线性组合恰好使两坐标同奇偶，而 \((2,0),(0,2)\) 已在其中。余权格由根的整数取值定义，由于 \(\varepsilon_1,\varepsilon_2\) 本身都是根，得到 \(P^\vee=\mathbb Z^2\)。

	权格要求 \(2a,2b,a+b,a-b\) 都是整数，等价于
	\[
		P=\mathbb Z^2\ \cup\ \bigl((\tfrac12,\tfrac12)+\mathbb Z^2\bigr).
	\]
	所以 \(P/Q\) 由半整数坐标的类生成，阶为二；\(P^\vee/Q^\vee\) 由坐标和的奇偶分类，也同构于 \(\mathbb Z/2\mathbb Z\)。原短根长为一、原长根长为 \(\sqrt2\)，余根相应长为二、\(\sqrt2\)，故长短次序反转。它解释了同一反射超平面排列仍能配上不同的根长数据。
\end{solution}

\begin{exercise}\label{g1:ex:regular-so3-disconnected-centralizer}
	令 \(g=\operatorname{diag}(-1,-1,1)\in\operatorname{SO}(3)\)，即绕第三坐标轴旋转 \(\pi\)。求 \(Z_{\operatorname{SO}(3)}(g)\) 及其单位分支，证明 \(g\) 是正则群元素，但其中心化子不是极大环面。
\end{exercise}
\begin{solution}
	与 \(g\) 交换的矩阵必须保持它的两个实特征空间 \(E_-\) 与 \(E_+\)，所以具有分块形式 \(\operatorname{diag}(A,\epsilon)\)，其中 \(A\in\operatorname{O}(2)\)、\(\epsilon\in\{1,-1\}\)。行列式等于一恰好要求 \(\epsilon=\det A\)，故
	\[
		Z_{\operatorname{SO}(3)}(g)
		=\{\operatorname{diag}(A,\det A):A\in\operatorname{O}(2)\}
		\cong\operatorname{O}(2).
	\]
	其单位分支为 \(\operatorname{diag}(\operatorname{SO}(2),1)\)，就是绕第三坐标轴的旋转环面 \(T\)。这个环面极大：与全部平面旋转交换的三维特殊正交矩阵必须保留轴线；在平面上只有旋转与全部旋转交换，行列式约束随后使轴线方向取正。故 \(Z_{\operatorname{SO}(3)}(T)=T\)，任何包含 \(T\) 的环面都只能是 \(T\)。于是群的秩为一，\(g\) 的中心化子也为一维，满足正则性定义。但是 \(\operatorname{O}(2)\) 有两个连通分支，例如 \(\operatorname{diag}(1,-1,-1)\) 在非单位分支中。因此整个中心化子严格大于 \(T\)。
\end{solution}

\begin{exercise}\label{g1:ex:su3-central-quotient-lattice}
	取 \(G=\operatorname{SU}(3)\) 及其中心子群 \(C=\{\zeta I:\zeta^3=1\}\)。用角度和为零的坐标描述 \(G/C\) 的周期格，并判定 \(\operatorname{SU}(3)\) 环面特征标
	\(\chi_m(\theta)=\exp(\mathrm{i}\sum_{j=1}^3m_j\theta_j)\)、\(m_j\in\mathbb Z\)，何时下降到商环面。由此求商环面的特征标格在原特征标格中的指数。
\end{exercise}
\begin{solution}
	指数落在 \(C\) 中，当且仅当三个相位相同且为三次单位根。因此周期格是
	\[
		\Lambda'=
		\{2\pi(m_1+k/3,m_2+k/3,m_3+k/3):
		m_j,k\in\mathbb Z,\ \sum_jm_j=-k\}.
	\]
	其中 \(k\) 改变三的倍数不改变相应陪集，所以原来的格 \(\Lambda=2\pi\{m\in\mathbb Z^3:\sum m_j=0\}\) 在 \(\Lambda'\) 中的指数为三。直接地，指数映射给出 \(\Lambda'/\Lambda\cong C\)。

	在和为零的平面上，两个整数三元组 \(m,m'\) 给出同一线性泛函，当且仅当 \(m-m'\) 是 \((1,1,1)\) 的整数倍。每个环面特征标都可如此表示：格基 \(2\pi(1,0,-1)\)、\(2\pi(0,1,-1)\) 上的取值除以 \(2\pi\) 是任意两个整数，可取 \(m_3=0\) 实现。对 \(\zeta I\)，有 \(\chi_m(\zeta I)=\zeta^{m_1+m_2+m_3}\)。故下降条件是 \(m_1+m_2+m_3\equiv0\pmod3\)；换一个表示三元组会给这个和加上三的倍数，条件不变。坐标和模三给出原特征标格到 \(\mathbb Z/3\mathbb Z\) 的满同态，其核就是可下降的格，故指数为三。
\end{solution}

\begin{exercise}\label{g1:ex:root-data-central-torus}
	设 \(K\) 为紧连通 Lie 群，\(T_K\) 为其极大环面，\(r\ge1\)。求 \(G=K\times(S^1)^r\) 的一个极大环面、Weyl 群、根集、周期格及特征标格，并描述它的 Weyl 室。说明单凭根集为何不能记录新增的连续中心方向。
\end{exercise}
\begin{solution}
	令 \(T=T_K\times(S^1)^r\)。若一个环面包含 \(T\)，它在 \(K\) 中的投影是包含 \(T_K\) 的紧连通交换子群，故由极大性等于 \(T_K\)；于是原环面已包含于 \(T\)，证明 \(T\) 极大。第二因子为中心，故
	\[
		N_G(T)=N_K(T_K)\times(S^1)^r,\qquad
		W(G,T)\cong W(K,T_K).
	\]
	Lie 代数为 \(\mathfrak k\oplus\mathbb R^r\)，第二个直和项的括号恒零，因此根恰为 \((\alpha,0)\)、\(\alpha\in R_K\)，各根空间来自原来的 \(\mathfrak k_\alpha\)。取标准角度坐标后，
	\[
		\Lambda_T=\Lambda_{T_K}\oplus2\pi\mathbb Z^r,\qquad
		X^*(T)=X^*(T_K)\oplus\mathbb Z^r.
	\]
	全部根超平面为 \(\ker\alpha\times\mathbb R^r\)，所以每个室是 \(C_K\times\mathbb R^r\)。根格始终为 \(Q_K\oplus\{0\}\)，新增维数全部是零权方向。若只保留根张成空间及其中的根集，各个 \(r\) 得到同一份数据；要区分这些群，至少还须记录完整的环面空间及其周期格。
\end{solution}

\begin{exercise}\label{g1:ex:unitary-multiplicity-orbit}
	设 \(\lambda,\mu,\nu\in S^1\) 两两不同，且 \(g=\operatorname{diag}(\lambda,\lambda,\mu,\nu)\in\operatorname{U}(4)\)。求 \(Z_{\operatorname{U}(4)}(g)\) 及 \(g\) 的共轭轨道维数。若再假设 \(\lambda^2\mu\nu=1\)，求特殊酉群中的中心化子及共轭轨道维数。说明为什么两种轨道的维数相同，而稳定子并不相同。
\end{exercise}
\begin{solution}
	与 \(g\) 交换的线性映射必须保持三个特征子空间 \(\mathbb C^2\)、\(\mathbb C\)、\(\mathbb C\)，在各子空间上又可任取酉变换。因此
	\[
		Z_{\operatorname{U}(4)}(g)=\operatorname{U}(2)\times\operatorname{U}(1)\times\operatorname{U}(1),
	\]
	实维数为 \(4+1+1=6\)。由轨道的齐性空间模型，共轭轨道维数为 \(16-6=10\)。若 \(g\in\operatorname{SU}(4)\)，中心化子还需满足 \(\det A\,z\,w=1\)，故为
	\[
		\{(A,z,w)\in\operatorname{U}(2)\times S^1\times S^1:\det A\,z\,w=1\}.
	\]
	用 \(w=(\det A\,z)^{-1}\) 可把它参数化为 \(\operatorname{U}(2)\times S^1\)，维数为五。于是特殊酉共轭轨道维数为 \(15-5=10\)。从酉群改为特殊酉群，同时删去了作用群与稳定子中同一个一维中心方向，因而轨道维数不变。实际上两轨道作为矩阵集合也相同：任意酉共轭矩阵可乘一个适当的标量矩阵调成行列式一，而标量不会改变共轭作用。
\end{solution}

\begin{exercise}\label{g1:ex:so4-oriented-conjugacy}
	取 \(0<a<b<\pi\)，令
	\[
		g=\operatorname{diag}(R(a),R(b)),\qquad h=\operatorname{diag}(R(a),R(-b))\in\operatorname{SO}(4).
	\]
	证明 \(g,h\) 有相同的复特征值和特征多项式，并在 \(\operatorname{O}(4)\) 中共轭，但在 \(\operatorname{SO}(4)\) 中不共轭。再判断 \(g\) 与 \(g^{-1}\) 在 \(\operatorname{SO}(4)\) 中是否共轭。
\end{exercise}
\begin{solution}
	两个旋转块的复特征值分别为 \(e^{\pm\mathrm{i}a},e^{\pm\mathrm{i}b}\)，改变第二个角的符号只交换它的两特征值。因此两矩阵的特征值多重集相同。令 \(C=\operatorname{diag}(1,-1)\)，则 \(CR(b)C^{-1}=R(-b)\)，故 \(\operatorname{diag}(I_2,C)\) 给出正交共轭，其行列式为负。

	若在特殊正交群内共轭，环面内的两个元素必须在同一 Weyl 群轨道中。\(\operatorname{SO}(4)\) 的 Weyl 群只允许交换两角或同时改变两角的符号。由于 \(a,b\) 严格位于 \((0,\pi)\) 且不同，模 \(2\pi\) 后，这些操作都不能把 \((a,b)\) 变成 \((a,-b)\)。因此不存在特殊正交共轭。另一方面，\(g^{-1}\) 对应 \((-a,-b)\)，同时在两个平面上反转定向由 \(\operatorname{diag}(C,C)\in\operatorname{SO}(4)\) 实现，所以 \(g,g^{-1}\) 在特殊正交群内共轭。特征多项式没有记录实特征平面组合的这项定向信息。
\end{solution}

\begin{exercise}\label{g1:ex:sp2-centralizers}
	在 \(\operatorname{Sp}(2)\) 中，分别求下列元素的中心化子及共轭轨道实维数：
	\[
		g_1=\operatorname{diag}(e^{\mathbf i a},e^{\mathbf i b}),\quad
		g_2=e^{\mathbf i a}I_2,\quad g_3=-I_2,
	\]
	其中 \(0<a<b<\pi\)。特别说明 \(g_2\) 虽然写成标量矩阵，为何不在群的中心。
\end{exercise}
\begin{solution}
	在复矩阵模型中，\(g_1\) 的四个特征值 \(e^{\pm\mathrm{i}a},e^{\pm\mathrm{i}b}\) 两两不同。与它交换的酉矩阵为对角矩阵，再要求具有紧辛群的共轭配对形式，便只剩 \(\operatorname{diag}(e^{\mathbf i t_1},e^{\mathbf i t_2})\)。所以中心化子为二维极大环面，轨道维数为 \(10-2=8\)。

	对于 \(g_2\)，交换条件逐个四元数分量写成 \(q e^{\mathbf i a}=e^{\mathbf i a}q\)。因 \(\sin a\ne0\)，这等价于 \(q\mathbf i=\mathbf i q\)，故所有矩阵分量必须落在 \(\mathbb C=\mathbb R+\mathbb R\mathbf i\) 中。余下的方程 \(A^*A=I\) 正是 \(\operatorname{U}(2)\)，因此中心化子维数四，轨道维数六。含 \(\mathbf j\) 的矩阵一般不与 \(g_2\) 交换，所以这里的四元数标量矩阵不是中心元；只有实标量才对所有四元数分量交换。\(g_3=-I_2\) 属于中心，中心化子为整个十维群，轨道是单点。
\end{solution}

\begin{exercise}\label{g1:ex:unitary-central-cover}
	设 \(n\ge2\)，定义 \(F:S^1\times\operatorname{SU}(n)\to\operatorname{U}(n)\) 为 \(F(z,A)=zA\)。证明这是一个 \(n\) 重覆盖同态，确定其核，并说明在由圆周第一因子及行列式确定的基本群生成元下，\(F_*:\mathbb Z\to\mathbb Z\) 是乘以 \(n\)。与正文中 \(S^1\times\operatorname{SU}(n)\cong\operatorname{U}(n)\) 的微分同胚比较，解释两个结论为何没有矛盾。
\end{exercise}
\begin{solution}
	标量矩阵与所有矩阵交换，所以 \(F\) 是群同态。给定 \(U\)，选取 \(z\in S^1\) 满足 \(z^n=\det U\)，则 \(A=z^{-1}U\in\operatorname{SU}(n)\)，因此满射。核为
	\[
		\ker F=\{(z,z^{-1}I):z^n=1\}\cong\mu_n.
	\]
	单位元微分为 \((\mathrm{i}t,X)\mapsto\mathrm{i}tI+X\)。取迹可知其核为零，且两边维数都为 \(n^2\)，所以它是局部微分同胚。由\cref{g1:prop:lie-homomorphism-local-isomorphism}，\(F\) 为覆盖，纤维含 \(n\) 个点。

	源群的基本群由 \(s\mapsto(e^{2\pi\mathrm{i}s},I)\) 生成。它在 \(F\) 下变为 \(e^{2\pi\mathrm{i}s}I\)，行列式绕圆周 \(n\) 次。正文已证明行列式在基本群上为同构，故 \(F_*\) 是乘以 \(n\)。正文的微分同胚则把 \((z,A)\) 送到 \(\operatorname{diag}(z,1,\ldots,1)A\)，它是另一张映射，通常不保持直积群乘法；其圆周生成环路的行列式只绕一次。微分同胚描述底流形，中心覆盖同态描述群乘法及离散核，所保留的结构不同。
\end{solution}

% !TeX root = ../../Book1.tex
% 纲要标题：## 第20章*　紧群表示与几何实现
% 纲要位置：几何学计划纲要.md，第 882 行。
% 正文状态：初稿；尚未正式审核。

% 纲要内容开始
% 
% > 本章目标与证明范围：不把 Peter–Weyl、最高权和线丛实现压成两个提要。正文选定紧 Lie 群的 Peter–Weyl 证明；一般最高权分类和 Borel–Weil 的完整理论仍在后续表示专题中，不作为本章前提。
% 
% 纲要内容结束

\OptionalChapter{紧群表示与几何实现}{g1:ch:20}
\BookChapterNumber{g1:ch:20}

\begin{chapterabstract}
	圆周上的 Fourier 展开把函数分解为各个整数频率；紧群的矩阵系数把这一方法推广到具有非交换对称的空间。本章从平均内积与 Schur 正交关系出发，以平滑卷积构造紧自伴算子，证明所需的谱分解，再由其有限维特征空间证明 Peter--Weyl 定理。随后将稳定子表示构造成齐性向量丛，以等变函数描述截面，并证明光滑截面的 Frobenius 互反。最后完整计算 \(\operatorname{SU}(2)\) 的不可约表示及其权，把齐次多项式实现为球面上线丛的截面，比较有限维多项式空间与全部光滑截面的表示分解。
\end{chapterabstract}

\begin{learninggoals}
	\begin{itemize}
		\item 用紧群平均证明有限维复表示可以酉化且完全可约，通过平均交织映射推出矩阵系数与特征标的正交关系。
		\item 构造平滑逼近单位，证明连续核算子的紧性及紧自伴谱分解，解释非零特征空间怎样产生有限维光滑表示。
		\item 证明矩阵系数的一致稠密性与正交 Hilbert 分解，区分正交展开、专门选择的一致逼近及光滑逼近。
		\item 从闭稳定子的表示构造齐性向量丛，核对截面的右等变条件与左群作用，并用 Frobenius 互反计算表示的出现重数。
		\item 分类 \(\operatorname{SU}(2)\) 的有限维复不可约表示，计算特征标与下降到 \(\operatorname{SO}(3)\) 的条件，将其多项式模型与球面线丛联系起来。
	\end{itemize}
\end{learninggoals}

在圆周上，旋转不会把一个纯频率变成另一个频率：若 \(f_m(e^{\mathrm{i}t})=e^{\mathrm{i}mt}\)，则平移角度 \(\theta\) 只把它乘以 \(e^{\mathrm{i}m\theta}\)。因此每个频率张成一个在群作用下不变的一维空间。球面上的旋转却会混合三个坐标函数；需要一起保留它们张成的三维空间，才能描述旋转的作用。由此自然产生一个问题：在一般紧群或紧齐性空间上，能否找到足够多的有限维不变函数空间，使任意函数都由这些部分逼近？

\ProseParagraphBreak

有限维表示给出候选函数。若 \(\pi(g)\) 是表示矩阵，它的每个条目都是群上的光滑函数，而矩阵乘法说明平移会把这些条目线性组合起来。平均内积使表示酉化，平均交织映射则使不同不可约表示的条目相互正交。但正交性本身并不保证没有遗漏函数。为证明完备性，我们先把函数与集中在单位元附近的平滑核卷积，再利用紧自伴算子的非零特征空间产生表示。对有限谱和再卷积一次，就能把原有的平方范数误差控制改进为一致误差。这是从分析走向表示的关键步骤。

\ProseParagraphBreak

表示还可以规定几何对象的转换规律。在齐性空间 \(G/H\) 上，先给定稳定子 \(H\) 在一个向量空间上的作用，就能把这个空间搬到各个底点，组成齐性向量丛。其截面对应群上的等变函数，因此前半章的平移和分解方法可以继续使用。对 \(\operatorname{SU}(2)/T\cong S^2\)，一个整数权决定一条复线丛；齐次多项式又提供其中可明确写出的有限维截面空间。我们将核对这两个整数的符号，而不只比较维数。

\ProseParagraphBreak

本章的群为有限维实 Lie 群，有限维表示空间取复数域，Hermitian 内积对第一变量线性。前半章使用\cref{g1:thm:compact-group-integration}的归一化左右不变积分、\cref{g1:thm:smooth-partition-unity}的单位分解及有限维线性代数；Hilbert 完备化与紧自伴谱引理在本章说明。\footnote{David Hilbert（1862--1943），德国数学家，在不变量理论、几何公理化及积分方程理论中作出重要贡献。}后半章使用第十七、十八章的齐性空间与关联丛，\(\operatorname{SU}(2)\) 的分类调用\cref{g1:lem:sl2-weight-chain}的有限权链及第十九章的极大环面。Peter--Weyl 部分不要求群连通，最后的球面模型则明确固定 \(\operatorname{SU}(2)\)。

\ProseParagraphBreak

延伸阅读可查 Hall 的《Lie Groups, Lie Algebras, and Representations》中题为“The Compact Group Approach to Representation Theory”的一章\cite{Hall2015}。这里证明紧 Lie 群的 Peter--Weyl 定理及 \(\operatorname{SU}(2)\) 的具体分类；一般最高权分类和 Borel--Weil 理论还需要根系、复结构及全纯截面的进一步工具，留在后续表示专题中讨论。

% !TeX root = ../../Book1.tex
% 纲要标题：### 20.1 有限维酉表示
% 纲要位置：几何学计划纲要.md，第 886 行。
% 正文状态：初稿；尚未正式审核。

% 纲要内容开始
% 
% * 平均内积和完全可约性
% * Schur 引理及矩阵系数正交
% * 正则表示的几何意义
% 
% 纲要内容结束

\section{有限维酉表示}
\label{g1:sec:2001}

圆周上的函数 \(z\mapsto z^m\)、\(m\in\mathbb Z\)，既是群同态，也是在 Fourier 展开中出现的基本函数。一般紧群通常不能用一维同态描述：例如三维旋转会混合向量的不同坐标，必须同时保留若干个矩阵分量。本节先研究这些有限维矩阵作用，证明它们的分量在群上积分时满足正交关系。随后把群的平移作用放到函数空间中，说明有限维表示怎样成为分析一般函数的材料。

\ProseParagraphBreak

本节中的 \(G\) 是有限维实紧 Lie 群，不要求连通；表示空间均为有限维复向量空间。归一化积分 \(\mu_G\) 取自\cref{g1:thm:compact-group-integration}，满足左右平移及取逆不变性。所有 Hermitian 内积对第一变量线性、对第二变量共轭线性，因此 \(\langle av,bw\rangle=a\overline b\langle v,w\rangle\)。固定这一约定后，矩阵分量和正交公式的指标才不会在取共轭时互换。

\subsection{平均内积和完全可约性}
\label{g1:subsec:2001:01}
\begin{definition}\label{g1:def:finite-dimensional-representation}
	设 \(G\) 为实 Lie 群，\(V\) 为有限维复向量空间。光滑群同态
	\[
		\pi:G\longrightarrow\operatorname{GL}(V)
	\]
	称为 \(G\) 在 \(V\) 上的\term[youxianweibiaoshi]{有限维表示}[finite-dimensional representation]，其中 \(\operatorname{GL}(V)\) 取实 Lie 群结构，\(\pi(g)\) 对 \(V\) 复线性。若 \(V\) 已赋予正定 Hermitian 内积，且每个 \(\pi(g)\) 保持它，则称 \(\pi\) 为\term[youbiaoshi]{酉表示}[unitary representation]。
\end{definition}

光滑性是这里的定义条件。以后从一个函数空间的有限维子空间构造表示时，也须验证参数 \(g\) 对应的矩阵随 \(g\) 光滑，不能仅凭每个 \(g\) 都给出线性变换就略去这一要求。群同态条件意味着 \(\pi(e)=I\)、\(\pi(g^{-1})=\pi(g)^{-1}\)，所以所有群元素都作用为可逆线性映射。

\begin{definition}\label{g1:def:irreducible-intertwiner}
	设 \(\pi:G\to\operatorname{GL}(V)\) 为有限维复表示。若复线性子空间 \(W\subseteq V\) 对每个 \(g\in G\) 都满足 \(\pi(g)W\subseteq W\)，则称 \(W\) 为不变子空间。若 \(V\ne0\)，且仅有 \(0,V\) 两个不变子空间，则称 \(\pi\)\term[bukeyuebiaoshi]{不可约}[irreducible]。若 \(V\) 可分解为有限个不可约不变子空间的直和，则称 \(\pi\)\term[wanquankeyue]{完全可约}[completely reducible]；零空间对应空直和。

	给定另一个表示 \(\sigma:G\to\operatorname{GL}(W)\)。若复线性映射 \(A:V\to W\) 对全部 \(g\) 都满足 \(A\pi(g)=\sigma(g)A\)，则称其为\term[jiaozhiyingshe]{交织映射}[intertwining map]，其空间记为 \(\operatorname{Hom}_G(V,W)\)。若存在可逆交织映射，称两表示等价。
\end{definition}

不变子空间条件实际给出 \(\pi(g)W=W\)，因为也可以使用 \(g^{-1}\)。在 \(W\) 上限制矩阵作用，就得到一个表示；在商空间 \(V/W\) 上也有诱导表示。仅有这些构造还不能把 \(V\) 拆成两部分，因为商空间未必能作为另一个不变子空间放回 \(V\)。紧性将通过平均内积解决这个问题。

\begin{proposition}[酉化与完全可约性]\label{g1:prop:compact-unitarization-complete-reducibility}
	设 \(G\) 为有限维实紧 Lie 群，\(\pi:G\to\operatorname{GL}(V)\) 为有限维复光滑表示。则 \(V\) 上存在 \(\pi(G)\) 不变的正定 Hermitian 内积。相对于任一这样的内积，每个不变子空间的正交补仍不变，且 \(V\) 可以写成不可约不变子空间的正交直和。
\end{proposition}
\begin{proof}
	任取正定 Hermitian 内积 \(h_0\)。\cref{g1:prop:compact-representation-inner-product}适用于整个紧群 \(G\) 的复线性表示，给出
	\[
		h(v,w)=\int_G h_0(\pi(g)v,\pi(g)w)\,\mathrm{d}\mu_G(g).
	\]
	回顾它在此处的用途：\(v\ne0\) 时各 \(\pi(g)v\ne0\)，所以平均后仍正定；右换元 \(g\mapsto ga\) 给出 \(h(\pi(a)v,\pi(a)w)=h(v,w)\)。由此，换用 \(h\) 的标准正交基以后，所有表示矩阵都为酉矩阵。

	若 \(W\) 不变、\(v\in W^\perp\)、\(w\in W\)，则
	\[
		h(\pi(g)v,w)=h(v,\pi(g)^{-1}w)=0,
	\]
	因为 \(\pi(g)^{-1}w\in W\)。所以 \(W^\perp\) 也不变，\(V=W\oplus W^\perp\) 是不变正交直和。最后对 \(\dim_{\mathbb C}V\) 归纳：零空间无需分解，非零不可约空间已经符合要求，否则取一个非零真不变子空间 \(W\)，分别分解两个维数更小的空间 \(W,W^\perp\)，再合并即可。
\end{proof}

这里保证分解存在，不保证各个具体不可约子空间唯一。例如在二维平凡表示中，每条直线都不变，任选一组标准正交基就有一种分解。稍后，特征标正交性将说明每一种不可约等价类出现的次数却与分解选择无关。

\begin{example}\label{g1:exm:noncompact-nonsemisimple-representation}
	加法群 \(\mathbb R\) 在 \(\mathbb C^2\) 上的表示
	\[
		\pi(t)=\begin{pmatrix}1&t\\0&1\end{pmatrix}
	\]
	不是完全可约的。直线 \(\mathbb Ce_1\) 不变；若它有一维不变补空间，可将该补空间的生成元写成 \((a,1)\)。不变性要求 \((a+t,1)\) 是 \((a,1)\) 的倍数，比较第二个分量使倍数为 \(1\)，再比较第一个分量便要求 \(t=0\)，与任意 \(t\) 矛盾。因此它也不可能有不变正定 Hermitian 内积，否则正交补会给出刚才不存在的不变补空间。紧性在酉化中的作用不能删除。
\end{example}

\subsection{Schur 引理及矩阵系数正交}
\label{g1:subsec:2001:02}

完全可约性把有限维表示拆成不可约部分。要比较这些部分，还需要知道它们之间有哪些保持群作用的线性映射。下面的\term[Schuryinli]{Schur 引理}[Schur's lemma]\footnote{Schur，Issai Schur（1875--1941），生于俄国、主要在德国工作的数学家，研究群表示与特征标。}说明，不可约性把这样的映射限制得很强；复数域则使不可约表示的自交织映射全部退化为标量。

\begin{lemma}[复 Schur 引理]\label{g1:lem:complex-schur}
	设 \(G\) 为实 Lie 群，\(\pi:G\to\operatorname{GL}(V)\)、\(\sigma:G\to\operatorname{GL}(W)\) 为非零有限维复不可约光滑表示。则任何非零交织映射 \(A:V\to W\) 都是同构；若两表示不等价，则 \(\operatorname{Hom}_G(V,W)=0\)。此外，\(\operatorname{End}_G(V)=\mathbb C I_V\)。
\end{lemma}
\begin{proof}
	交织条件给出 \(\pi(g)\ker A\subseteq\ker A\)、\(\sigma(g)\operatorname{im}A\subseteq\operatorname{im}A\)，所以核与像是不变子空间。若 \(A\ne0\)，不可约性迫使 \(\ker A=0\)、\(\operatorname{im}A=W\)，故 \(A\) 是同构。不存在同构时，所有交织映射便只能为零。

	再设 \(A\in\operatorname{End}_G(V)\)。复数域上特征多项式有根，故 \(A\) 有特征值 \(\lambda\)。\(A-\lambda I\) 仍是交织映射，却有非零核，所以由刚证的结论，它不可能是非零映射。因此 \(A=\lambda I\)。
\end{proof}

此引理不要求 \(G\) 紧。紧性要在接下来的平均中使用。复数域的条件则不能省略：圆周在 \(\mathbb R^2\) 上的旋转表示没有不变实直线，因而实不可约，但旋转 \(\pi/2\) 的矩阵与所有旋转交换，并非实标量矩阵。转到复数域以后，这个二维表示会分成两个一维表示，与刚才的结论相容。

\ProseParagraphBreak

对任意线性映射 \(A:W\to V\)，通常 \(\pi(g)A\ne A\sigma(g)\)。将它在群作用下的所有变换平均，恰好能得到一个交织映射。这里平均的不是向量的长度，而是映射本身，所得零或非零的结果将产生矩阵系数的正交关系。

\begin{proposition}[平均交织映射]\label{g1:prop:averaged-intertwiner}
	设 \(G\) 为有限维实紧 Lie 群，\(\mu_G\) 为 \(G\) 上满足 \(\mu_G(G)=1\) 的左右不变 Haar 测度（简记为 \(\mu\)），\(\pi\) 与 \(\sigma\) 为有限维复酉表示，空间分别为 \(V,W\)。对 \(A\in\operatorname{Hom}_{\mathbb C}(W,V)\)，定义
	\[
		\mathcal P_{\pi,\sigma}(A)
		\coloneqq\int_G\pi(g)A\sigma(g)^{-1}\,\mathrm{d}\mu_G(g).
	\]
	则 \(\mathcal P_{\pi,\sigma}(A)\) 从 \(\sigma\) 交织到 \(\pi\)，且 \(\mathcal P_{\pi,\sigma}\) 是到交织映射空间的线性投影。若 \(\pi,\sigma\) 不可约且不等价，则该投影为零；若 \(\sigma=\pi\)，且 \(\pi\) 不可约、\(d_\pi=\dim_{\mathbb C}V\)，则
	\begin{equation}\label{g1:eq:averaged-endomorphism-trace}
		\mathcal P_{\pi,\pi}(A)=\frac{\operatorname{tr}A}{d_\pi}I_V.
	\end{equation}
\end{proposition}
\begin{proof}
	积分在有限维映射空间中逐矩阵分量进行。对 \(h\in G\)，左换元 \(g\mapsto hg\) 给出
	\[
		\begin{aligned}
			\pi(h)\mathcal P_{\pi,\sigma}(A)\sigma(h)^{-1}
			&=\int_G\pi(hg)A\sigma(hg)^{-1}\,\mathrm{d}\mu_G(g)\\
			&=\mathcal P_{\pi,\sigma}(A).
		\end{aligned}
	\]
	这正是交织条件。若 \(A\) 原本已经交织，则 \(\pi(g)A\sigma(g)^{-1}=A\) 对每个 \(g\) 成立，归一化积分使 \(\mathcal P_{\pi,\sigma}(A)=A\)。因而它在其像上是恒等，确实为投影。

	不可约不等价情形的结论由\cref{g1:lem:complex-schur}得到。\(\sigma=\pi\) 时，该引理给出 \(\mathcal P_{\pi,\pi}(A)=cI_V\)。取迹并利用相似变换不改变迹，便有
	\[
		c\,d_\pi=\int_G\operatorname{tr}(\pi(g)A\pi(g)^{-1})\,\mathrm{d}\mu_G(g)
		=\operatorname{tr}A,
	\]
	从而得到\eqref{g1:eq:averaged-endomorphism-trace}。
\end{proof}

对表示 \(\pi\)、向量 \(v\in V\) 及复线性泛函 \(\ell\in V^*\)，函数 \(g\mapsto\ell(\pi(g)v)\) 称为\term[juzhenxishu]{矩阵系数}[matrix coefficient]。在酉表示的标准正交基 \(e_1,\ldots,e_{d_\pi}\) 中，我们固定
\begin{equation}\label{g1:eq:matrix-coefficient-convention}
	\pi_{ij}(g)\coloneqq\langle\pi(g)e_j,e_i\rangle,
	\qquad
	\pi(g)e_j=\sum_i\pi_{ij}(g)e_i.
\end{equation}
因此 \(i\) 是行指标，\(j\) 是列指标。一般矩阵系数正是这些函数的线性组合，而且都光滑。

\begin{theorem}[Schur 正交关系]\label{g1:thm:schur-orthogonality}
	设 \(G\) 为有限维实紧 Lie 群，\(\mu_G\) 为 \(G\) 上满足 \(\mu_G(G)=1\) 的左右不变 Haar 测度（简记为 \(\mu\)）。从每个有限维复不可约表示的等价类中选一个酉代表，并在其空间上固定标准正交基。对两个所选代表 \(\pi,\sigma\)，有
	\begin{equation}\label{g1:eq:schur-matrix-orthogonality}
		\int_G\pi_{ij}(g)\overline{\sigma_{kl}(g)}\,\mathrm{d}\mu_G(g)
		=\begin{cases}
			\dfrac{\delta_{ik}\delta_{jl}}{d_\pi},&\pi=\sigma,\\[4pt]
			0,&\pi\not\cong\sigma.
		\end{cases}
	\end{equation}
	其中 \(d_\pi=\dim_{\mathbb C}V_\pi\)，\(\delta_{ab}\) 在 \(a=b\) 时为 \(1\)，否则为 \(0\)。
\end{theorem}
\begin{proof}
	令 \(e_j\) 为 \(V_\pi\) 的基，\(f_l\) 为 \(V_\sigma\) 的基。取秩一映射 \(A:V_\sigma\to V_\pi\)，将 \(f_l\) 送到 \(e_j\)，将其余基向量送到零。平均映射的第 \((i,k)\) 个矩阵分量为
	\[
		\bigl(\mathcal P_{\pi,\sigma}(A)\bigr)_{ik}
		=\int_G\pi_{ij}(g)\bigl(\sigma(g)^{-1}\bigr)_{lk}\,\mathrm{d}\mu_G(g)
		=\int_G\pi_{ij}(g)\overline{\sigma_{kl}(g)}\,\mathrm{d}\mu_G(g),
	\]
	最后一步用了酉性 \(\sigma(g)^{-1}=\sigma(g)^*\)。若 \(\pi\not\cong\sigma\)，\cref{g1:prop:averaged-intertwiner}使平均映射为零，得到第二种情形。若 \(\pi=\sigma\)，两组基就是同一组，此时 \(\operatorname{tr}A=\delta_{jl}\)。\eqref{g1:eq:averaged-endomorphism-trace}使平均映射等于 \(\delta_{jl}I/d_\pi\)，其第 \((i,k)\) 个分量为 \(\delta_{jl}\delta_{ik}/d_\pi\)，得到第一种情形。
\end{proof}

“从每个等价类只取一个代表”消除了公式中的选择歧义。两个等价表示若分别使用任意基，它们的矩阵系数应先通过交织同构作基变换，不能仅凭等价就直接令不同表示空间上的基指标相等。酉不可约表示之间的同构总可取为酉同构：若 \(A\) 交织，则 \(A^*A\) 是正定自交织算子，Schur 引理使它等于 \(cI\)、\(c>0\)，所以 \(A/\sqrt c\) 酉。

\begin{definition}\label{g1:def:representation-character}
	有限维复表示 \(\pi\) 的\term[biaoshitezhengbiao]{表示特征标}[character of a representation]为
	\[
		\chi_\pi(g)\coloneqq\operatorname{tr}\pi(g).
	\]
	它是与基选择无关的光滑复值函数，并在每个共轭类上常值。
\end{definition}

当 \(\dim V=1\) 时，\(\chi_\pi=\pi\)，它就是\chapref{g1:ch:19}中的一维特征标；一般维数的迹却不是到 \(S^1\) 的群同态。共轭不变性来自 \(\pi(hgh^{-1})=\pi(h)\pi(g)\pi(h)^{-1}\) 及迹的循环性质。取酉模型时，还能读出 \(\chi_\pi(e)=d_\pi\)、\(\chi_\pi(g^{-1})=\overline{\chi_\pi(g)}\)。

\begin{corollary}[特征标正交与重数]\label{g1:cor:character-orthogonality}
	设 \(G\) 为有限维实紧 Lie 群，\(\mu_G\) 为 \(G\) 上满足 \(\mu_G(G)=1\) 的左右不变 Haar 测度（简记为 \(\mu\)），\(\pi,\sigma\) 为所选的有限维复不可约酉表示代表。则
	\[
		\int_G\chi_\pi(g)\overline{\chi_\sigma(g)}\,\mathrm{d}\mu_G(g)
		=\begin{cases}1,&\pi=\sigma,\\0,&\pi\not\cong\sigma.\end{cases}
	\]
	对任意有限维复光滑表示 \(\rho\)，其完全可约分解中 \(\pi\) 出现的次数为
	\begin{equation}\label{g1:eq:character-multiplicity-integral}
		m_\pi(\rho)=\int_G\chi_\rho(g)\overline{\chi_\pi(g)}\,\mathrm{d}\mu_G(g).
	\end{equation}
	因此各重数与分解选择无关；两个有限维复表示等价，当且仅当它们的特征标相同。
\end{corollary}
\begin{proof}
	将 \(\chi_\pi=\sum_i\pi_{ii}\)、\(\chi_\sigma=\sum_k\sigma_{kk}\) 代入积分。Schur 正交关系在不等价时使每项为零；在同一代表时，只有 \(i=k\) 的 \(d_\pi\) 项各贡献 \(1/d_\pi\)，总和为 \(1\)。

	由完全可约性，\(\rho\) 是有限个不可约表示的直和。直和矩阵的迹相加，故 \(\chi_\rho=\sum_\sigma m_\sigma\chi_\sigma\)，这里只含有限个非零重数。与 \(\chi_\pi\) 作上述积分便抽出 \(m_\pi\)，证明\eqref{g1:eq:character-multiplicity-integral}及选择无关性。等价表示的矩阵彼此相似，所以特征标相同；反之，特征标相同给出所有重数相同，将各不可约直和因子逐个对应，即得表示同构。
\end{proof}

\begin{example}\label{g1:exm:torus-irreducible-representations}
	若 \(G=T\) 为环面，则每个有限维复不可约表示都是一维的。确实，\(T\) 交换，所以每个 \(\pi(t)\) 都与全部 \(\pi(T)\) 交换；Schur 引理使它们都是标量。若表示空间维数大于一，任一复直线便不变，违背不可约性。因此环面的不可约表示正是\cref{g1:prop:torus-character-lattice}给出的格上特征标。对 \(S^1\)，它们为 \(\chi_m(z)=z^m\)、\(m\in\mathbb Z\)，正交关系化为
	\[
		\frac1{2\pi}\int_0^{2\pi}e^{\mathrm{i}(m-n)t}\,\mathrm{d}t=\delta_{mn}.
	\]
	一般紧群中，单个不可约表示维数可以大于一；相应的正交对象就从一个指数函数扩展为 \(d_\pi^2\) 个矩阵分量。
\end{example}

\subsection{正则表示的几何意义}
\label{g1:subsec:2001:03}

每个表示给出群上的一些函数。反过来，群本身的平移也会移动函数，由此得到一个同时容纳许多有限维表示的空间。先在 \(C(G)\)、即 \(G\) 上连续复值函数的空间中工作，定义
\begin{equation}\label{g1:eq:regular-actions-convention}
	(L_gf)(x)\coloneqq f(g^{-1}x),\qquad
	(R_gf)(x)\coloneqq f(xg).
\end{equation}
它们分别称为\term[zuozhengzebiaoshi]{左正则表示}[left regular representation]和\term[youzhengzebiaoshi]{右正则表示}[right regular representation]。\symidx{\(L_g,R_g\)：函数上的左右正则作用}左式中的逆元保证两个算子按群乘法的同一顺序复合；右式已有参数写在 \(x\) 的右边，故不再取逆。这些是无限维函数空间上的线性表示，前面定义中的有限维光滑性条件尚不直接适用于整个空间。

\begin{proposition}[左右正则作用]\label{g1:prop:regular-representations}
	设 \(G\) 为有限维实紧 Lie 群，\(\mu_G\) 为 \(G\) 上满足 \(\mu_G(G)=1\) 的左右不变 Haar 测度（简记为 \(\mu\)）。令 \(C(G)=C(G,\mathbb C)\)，按\eqref{g1:eq:regular-actions-convention}的约定定义
	\[
		L_gf(x)=f(g^{-1}x),\qquad R_gf(x)=f(xg).
	\]
	它们是 \(C(G)\) 上两组可逆复线性算子，满足
	\[
		L_gL_h=L_{gh},\qquad R_gR_h=R_{gh},\qquad L_gR_h=R_hL_g.
	\]
	它们保持上确界范数以及积分内积
	\begin{equation}\label{g1:eq:compact-group-function-inner-product}
		\langle f,u\rangle_2
		\coloneqq\int_G f(x)\overline{u(x)}\,\mathrm{d}\mu_G(x).
	\end{equation}
	它们也保持子空间 \(C^\infty(G)\)。
\end{proposition}
\begin{proof}
	逐点代入可得
	\[
		\begin{aligned}
			(L_gL_hf)(x)&=f(h^{-1}g^{-1}x)=f((gh)^{-1}x),\\
			(R_gR_hf)(x)&=f((xg)h)=f(x(gh)),\\
			(L_gR_hf)(x)&=f(g^{-1}xh)=(R_hL_gf)(x).
		\end{aligned}
	\]
	单位元作用是恒等，\(g^{-1}\) 给出逆算子。左、右平移都是 \(G\) 上的双射，所以保持上确界范数；它们又是光滑映射，故复合保持光滑性。

	积分内积的线性及 Hermitian 对称性来自积分。若连续函数 \(f\ne0\)，则 \(|f|^2\) 是非零非负连续函数，\cref{g1:thm:compact-group-integration}给出 \(\langle f,f\rangle_2>0\)，因此它确为内积。该定理的左、右换元公式分别给出
	\[
		\langle L_gf,L_gu\rangle_2=\langle f,u\rangle_2,
		\qquad
		\langle R_gf,R_gu\rangle_2=\langle f,u\rangle_2.
	\]
\end{proof}

\(C(G)\) 对此积分范数未必完备；下一节将对它作 Hilbert 完备化，得到 \(L^2(G)\)。目前的矩阵系数已经是光滑函数，所有点值、平移和积分都可以直接计算。

\begin{definition}\label{g1:def:representative-functions}
	设 \(G\) 为有限维实紧 Lie 群。记 \(\widehat G\) 为 \(G\) 的有限维复不可约表示的等价类集合，并为每一类选定一个酉代表 \(\pi:G\to\operatorname{U}(V_\pi)\) 及标准正交基 \(e_1,\ldots,e_{d_\pi}\)，其中 \(d_\pi=\dim_{\mathbb C}V_\pi\)，\(\pi_{ij}(g)=\langle\pi(g)e_j,e_i\rangle\)，内积对第一变量线性。\symidx{\(\widehat G\)：紧群的不可约酉表示等价类集合}记
	\[
		\mathcal H_\pi\coloneqq\operatorname{span}_{\mathbb C}\{\pi_{ij}:1\le i,j\le d_\pi\}
	\]
	为该表示的矩阵系数空间。所有有限维复光滑表示的矩阵系数的有限线性组合，组成\term[daibiaohanshu]{代表函数}[representative function]空间，记为 \(\mathcal R(G)\)。\symidx{\(\mathcal H_\pi\)：不可约表示的矩阵系数空间}\symidx{\(\mathcal R(G)\)：紧群的代表函数代数}
\end{definition}

\begin{proposition}[矩阵系数空间的两侧作用]\label{g1:prop:matrix-coefficient-biregular-model}
	设 \(G\) 为有限维实紧 Lie 群，\(\mu_G\) 为 \(G\) 上满足 \(\mu_G(G)=1\) 的左右不变 Haar 测度（简记为 \(\mu\)），\(\widehat G\) 为有限维复不可约光滑表示的等价类集合，每类选一个酉代表。对一个所选代表 \(\pi:G\to\operatorname{U}(V_\pi)\)，固定标准正交基 \(e_i\)，令 \(d_\pi=\dim_{\mathbb C}V_\pi\)、\(\pi_{ij}(g)=\langle\pi(g)e_j,e_i\rangle\)，内积对第一变量线性。记 \(\mathcal H_\pi=\operatorname{span}_{\mathbb C}\{\pi_{ij}\}\)，\(\mathcal R(G)\) 为所有有限维复光滑表示的矩阵系数的线性张成空间，并令 \(L_gf(x)=f(g^{-1}x)\)、\(R_hf(x)=f(xh)\)。则 \(\mathcal H_\pi\) 的复维数为 \(d_\pi^2\)，且 \(\{\sqrt{d_\pi}\,\pi_{ij}\}\) 是积分内积下的标准正交基。不同 \(\pi\) 的空间彼此正交，且
	\[
		\mathcal R(G)=\bigoplus_{\pi\in\widehat G}^{\mathrm{alg}}\mathcal H_\pi,
	\]
	其中上标 \(\mathrm{alg}\) 表示每个元素只涉及有限个分量。每个 \(\mathcal H_\pi\) 同时在左右正则作用下不变。若 \(\pi^\vee(g)\ell=\ell\circ\pi(g^{-1})\) 为 \(V_\pi^*\) 上的对偶表示，则
	\begin{equation}\label{g1:eq:coefficient-tensor-model}
		V_\pi^*\otimes V_\pi\longrightarrow\mathcal H_\pi,
		\qquad\ell\otimes v\longmapsto\bigl[x\mapsto\ell(\pi(x)v)\bigr]
	\end{equation}
	是线性同构，将 \(\pi^\vee(g)\otimes\pi(h)\) 对应到 \(L_gR_h\)。
\end{proposition}
\begin{proof}
	Schur 正交关系说明同一代表的 \(d_\pi^2\) 个系数彼此正交且范数均为 \(1/\sqrt{d_\pi}\)，故线性无关；同时也说明不同代表的系数空间正交。任意有限维表示完全可约，把它按不变直和写成分块对角矩阵，其矩阵系数便是有限个不可约矩阵系数的线性组合。再作任意基变换也只改变这些线性组合。因此代表函数空间恰是所列代数直和。

	由群同态性质，有
	\begin{equation}\label{g1:eq:left-right-matrix-coefficients}
		\begin{aligned}
			L_g\pi_{ij}&=\sum_k\pi_{ik}(g^{-1})\pi_{kj},\\
			R_h\pi_{ij}&=\sum_k\pi_{kj}(h)\pi_{ik}.
		\end{aligned}
	\end{equation}
	所以 \(\mathcal H_\pi\) 在两种作用下不变。\eqref{g1:eq:coefficient-tensor-model}把对偶基与原基的张量送到全部矩阵系数；两边都有维数 \(d_\pi^2\)，且这些系数线性无关，故为同构。最后，直接计算
	\[
		\begin{aligned}
			(L_gR_h[\ell(\pi(\cdot)v)])(x)
			&=\ell(\pi(g^{-1}xh)v)\\
			&=(\pi^\vee(g)\ell)(\pi(x)\pi(h)v),
		\end{aligned}
	\]
	证明其与两侧作用相容。
\end{proof}

特别地，固定行指标 \(i\) 后，\(\operatorname{span}\{\pi_{i1},\ldots,\pi_{id_\pi}\}\) 在右正则作用下是一个与 \(\pi\) 等价的副本；\(\mathcal H_\pi\) 因而包含 \(d_\pi\) 个这样的副本。这里的重复次数来自另一个未被右平移改变的指标。这比“每个不可约表示出现一次”更准确，也解释了为什么非交换群的 Fourier 系数应当是矩阵。

\begin{proposition}[代表函数的代数性质]\label{g1:prop:representative-function-algebra}
	设 \(G\) 为有限维实紧 Lie 群。令 \(\mathcal R(G)\) 为函数 \(g\mapsto\ell(\pi(g)v)\) 的有限复线性组合组成的空间，其中 \(\pi:G\to\operatorname{GL}(V)\) 遍历有限维复光滑表示，\(v\in V\)、\(\ell\in V^*\)。则 \(\mathcal R(G)\subseteq C^\infty(G)\) 包含常数函数，对逐点乘法及复共轭封闭，并在左右平移下不变。
\end{proposition}
\begin{proof}
	平凡一维表示给出常数 \(1\)，其标量倍给出所有常数。若 \(f(g)=\ell(\pi(g)v)\)、\(u(g)=m(\sigma(g)w)\)，则张量积表示 \((\pi\otimes\sigma)(g)=\pi(g)\otimes\sigma(g)\) 有矩阵系数
	\[
		(fu)(g)=(\ell\otimes m)((\pi\otimes\sigma)(g)(v\otimes w)).
	\]
	张量积矩阵的各分量是原来分量的乘积，故它仍是光滑有限维表示。对复共轭向量空间 \(\overline V\)，共轭矩阵 \(\overline{\pi(g)}\) 同样给出光滑表示，其矩阵系数就是原系数的复共轭。因此有限线性组合也对乘法和共轭封闭。平移不变性由\eqref{g1:eq:left-right-matrix-coefficients}及代数直和分解得到。
\end{proof}

还需要反向的结论：如果在函数空间中找到一个有限维平移不变子空间，其中的函数是否自动成为某个表示的矩阵系数？这个问题正好连接下一节的算子特征空间与本节的代数结果。

\begin{proposition}[有限维正则子空间的矩阵系数]\label{g1:prop:finite-regular-subspace-coefficients}
	设 \(G\) 为有限维实紧 Lie 群，\(\mu_G\) 为 \(G\) 上满足 \(\mu_G(G)=1\) 的左右不变 Haar 测度（简记为 \(\mu\)），\(E\subseteq C^\infty(G)\) 为有限维复线性子空间，令 \(R_gf(x)=f(xg)\)，并假设 \(R_gE\subseteq E\) 对每个 \(g\in G\) 成立。则 \(\rho(g)=R_g|_E\) 是 \(E\) 上的有限维复光滑酉表示，内积取积分内积的限制。对每个 \(f\in E\)，有
	\begin{equation}\label{g1:eq:evaluation-recovers-coefficient}
		f(g)=\operatorname{ev}_e(\rho(g)f),\qquad
		\operatorname{ev}_e(u)=u(e).
	\end{equation}
	因此 \(E\subseteq\mathcal R(G)\)。
\end{proposition}
\begin{proof}
	\(E=0\) 时结论成立。设 \(\dim_{\mathbb C}E=d>0\)。\cref{g1:prop:regular-representations}已经给出群乘法规律与积分内积不变性。又因 \(R_{g^{-1}}E\subseteq E\)，有 \(R_gE=E\)，所以限制算子可逆。剩下的关键是验证 \(g\mapsto\rho(g)\) 光滑。

	所有点评价 \(\operatorname{ev}_x|_E\)、\(x\in G\)，张成 \(E^*\)。若其线性张成不为整个对偶空间，有限维线性代数便给出非零 \(f\in E\)，被全部这些泛函消去，也就是 \(f(x)=0\) 对全部 \(x\) 成立，矛盾。因此可以从中选出 \(d\) 个组成 \(E^*\) 的基，设对应点为 \(x_1,\ldots,x_d\)。求值映射
	\[
		J:E\longrightarrow\mathbb C^d,
		\qquad J(f)=(f(x_1),\ldots,f(x_d))
	\]
	为线性同构。固定 \(f\in E\)，\(J(\rho(g)f)=(f(x_1g),\ldots,f(x_dg))\) 的每个分量都随 \(g\) 光滑，因为 \(f\) 及群乘法光滑。将 \(f\) 依次取为 \(E\) 的基，再乘固定的线性映射 \(J^{-1}\)，就得到 \(\rho(g)\) 的全部矩阵分量光滑。故它确为本节定义意义下的表示。

	最后，\(\operatorname{ev}_e(\rho(g)f)=(R_gf)(e)=f(g)\)。右端由这个有限维表示、向量 \(f\) 及线性泛函 \(\operatorname{ev}_e|_E\) 给出，正是矩阵系数的定义，因此 \(f\in\mathcal R(G)\)。
\end{proof}

至此，我们已经知道矩阵系数相互正交，也知道有限维平移不变的光滑函数空间会产生矩阵系数。还没有证明这些函数足以逼近任意连续函数。下一节将构造与右平移交换的紧算子：它们的非零特征空间有限维，经过平滑性核验后可用上面的命题，再由逼近单位使这些有限维空间合起来逼近整个函数空间。

% !TeX root = ../../Book1.tex
% 纲要标题：### 20.2 Peter–Weyl 定理的分析证明
% 纲要位置：几何学计划纲要.md，第 892 行。
% 正文状态：初稿；尚未正式审核。

% 纲要内容开始
% 
% * 紧群上的逼近单位
% * 平滑卷积核和正则表示
% * 卷积算子与紧自伴谱分解
% 
%   就地建立所用紧自伴算子的谱分解引理，再把其非零特征空间用于表示分解；本章不要求先学第三卷全部椭圆谱理论。
% 
% * 非零特征空间的有限维不变性
% * 有限谱逼近与矩阵系数分离
% * 矩阵系数的稠密性
% * Peter–Weyl 分解与收敛
% 
% 纲要内容结束

\section{Peter–Weyl 定理的分析证明}
\label{g1:sec:2002}

圆周上的三角多项式既来自一维表示，也能逼近连续函数。一般紧群没有统一的角坐标，表示还可能高于一维；我们要证明，所有有限维表示的矩阵系数仍足以逼近任意连续函数。困难在于不能预先假定已经找到足够多的表示。这里反过来从函数空间出发：用平滑卷积构造紧自伴算子，其非零特征空间自动有限维并被群作用保持，由此产生所需表示。

\ProseParagraphBreak

本节 \(G\) 是有限维实紧 Lie 群，不要求连通；\(\mu_G\) 是\cref{g1:thm:compact-group-integration}的归一化左右不变密度，以下简写为 \(\mu\)。函数和表示空间均取复数值，内积对第一变量线性。沿用\secref{g1:sec:2001}的记号 \(\mathcal R(G)\)，表示有限维光滑表示的矩阵系数所张成的空间；有限维完全可约性使它也等于不可约酉表示系数的线性张成。

\subsection{紧群上的逼近单位}
\label{g1:subsec:2002:01}
先解释证明中要使用的函数空间。沿用\eqref{g1:eq:compact-group-function-inner-product}，在 \(C(G)\) 上取积分内积及其范数
\[
	\langle f,h\rangle_2=\int_G f(x)\overline{h(x)}\,\mathrm{d}\mu(x),
	\qquad \|f\|_2=\sqrt{\langle f,f\rangle_2}.
\]
连续非零函数在某个开集上有非零下界，正密度使该开集有正积分，所以 \(\|f\|_2=0\) 仅当 \(f=0\)。本章把这个内积空间的完备化记为 \(L^2(G)\)：元素由 \(C(G)\) 中的 \(\|\cdot\|_2\)-Cauchy 列表示，两列之差趋于零时表示同一个元素，内积取代表列内积的极限。Cauchy--Schwarz 不等式保证这个极限存在且与代表无关。逐项选取误差小于 \(2^{-j}\) 的连续近似，再作对角选择，便使完备化中的每个 Cauchy 列收敛；因此所得空间是 Hilbert 空间，即完备的内积空间。这个模型也与通常按密度测度定义的平方可积函数空间一致，但以下论证只使用完备化构造。\symidx{\(L^2(G)\)：紧群上的平方可积函数空间}

\ProseParagraphBreak

由 \(\mu(G)=1\) 有 \(\|f\|_2\le\|f\|_\infty\)。连续函数在 \(L^2(G)\) 中稠密是上述定义的一部分，而 \(L^2\) 元素一般没有指定的逐点值。以后对它写卷积积分时，先在连续函数上定义，再按范数连续延拓；卷积产生的连续函数则具有真实的逐点值。

\begin{lemma}[平滑逼近单位]\label{g1:lem:compact-group-approximate-identity}
	设 \(G\) 为有限维实紧 Lie 群，\(\mu_G\) 为 \(G\) 上满足 \(\mu_G(G)=1\) 的左右不变 Haar 测度（简记为 \(\mu\)）。存在实值光滑函数列 \(k_n\)，满足
	\[
		k_n\ge0,\qquad \int_Gk_n\,\mathrm{d}\mu=1,
		\qquad k_n(x^{-1})=k_n(x),
	\]
	且其支集收缩到单位元：每个 \(e\) 的邻域都包含充分大 \(n\) 时的 \(\operatorname{supp}k_n\)。对任意 \(f\in C(G)\)，
	\[
		T_nf(x)=\int_G k_n(y)f(y^{-1}x)\,\mathrm{d}\mu(y)
	\]
	一致趋于 \(f\)，并且 \(\|T_nf\|_\infty\le\|f\|_\infty\)、\(\|T_nf\|_2\le\|f\|_2\)。因此 \(T_n\) 延拓为 \(L^2(G)\) 上的压缩算子，且 \(T_nf\to f\) 对每个 \(f\in L^2(G)\) 按 \(L^2\) 范数成立。
\end{lemma}
\begin{proof}
	取一列收缩到 \(e\) 的对称开邻域 \(U_n=U_n^{-1}\)，例如在局部坐标小邻域与其逆像中逐次缩小。由\cref{g1:lem:smooth-bump-function}的坐标截断及零延拓，存在 \(b_n\ge0\)，在 \(e\) 附近等于 \(1\)，且支集包含于 \(U_n\)。令
	\[
		k_n(x)=\frac{b_n(x)+b_n(x^{-1})}{2\int_Gb_n\,\mathrm{d}\mu}.
	\]
	分母严格为正，取逆不变换元保证积分为 \(1\)，其余要求由定义成立。零维紧群为有限群，这时也可取 \(U_n=\{e\}\)，构造仍适用。

	给定 \(\varepsilon>0\)，紧性与 \((y,x)\mapsto f(y^{-1}x)-f(x)\) 的连续性给出 \(e\) 的邻域 \(U\)，使该差对 \(y\in U\)、\(x\in G\) 一律小于 \(\varepsilon\)。确实，每个 \((e,x)\) 有这样的积邻域，有限个 \(x\) 邻域覆盖 \(G\)，再取对应 \(e\) 邻域之交即可。充分大 \(n\) 时 \(k_n\) 支于 \(U\)，所以
	\[
		|T_nf(x)-f(x)|
		\le\int_G k_n(y)|f(y^{-1}x)-f(x)|\,\mathrm{d}\mu(y)
		\le\varepsilon.
	\]
	同理得到一致范数的压缩估计。将 Cauchy--Schwarz 不等式用于权重为 \(k_n\) 的积分，再对 \(x\) 积分，有
	\[
		\|T_nf\|_2^2
		\le\int_G\int_G k_n(y)|f(y^{-1}x)|^2\,
		\mathrm{d}\mu(y)\,\mathrm{d}\mu(x)=\|f\|_2^2.
	\]
	这里被积函数连续，紧参数下的逐次积分可以交换；左不变换元使内层 \(x\) 积分等于 \(\|f\|_2^2\)。因此 \(T_n\) 唯一延拓到完备化，算子范数不超过 \(1\)。若 \(f\in L^2(G)\)，选 \(h\in C(G)\) 使 \(\|f-h\|_2<\varepsilon\)，便有
	\[
		\|T_nf-f\|_2\le2\varepsilon+\|T_nh-h\|_2.
	\]
	末项由一致收敛趋于零，随后令 \(\varepsilon\to0\)，得到所述强收敛。
\end{proof}

这些 \(k_n\) 称为\term[bijindanwei]{逼近单位}[approximate identity]。名称描述的是卷积趋近恒等算子；函数 \(k_n\) 本身并不趋于常数函数 \(1\)，而是把总量 \(1\) 集中到越来越小的单位元邻域。

\subsection{平滑卷积核和正则表示}
\label{g1:subsec:2002:02}
要把逼近与表示联系起来，需要卷积既能改善函数的正则性，又能与一侧正则作用交换。以下始终采用左卷积
\begin{equation}\label{g1:eq:left-convolution-kernel}
	T_kf(x)=\int_Gk(y)f(y^{-1}x)\,\mathrm{d}\mu(y)
	=\int_Gk(xz^{-1})f(z)\,\mathrm{d}\mu(z).
\end{equation}
第二式由 \(z=y^{-1}x\) 换元得到，使用左右不变性及取逆不变性。它把作用于 \(f\) 的平移改写为作用于核的平移，便于对输出求导。

\begin{proposition}[卷积的平滑性与交换性]\label{g1:prop:smoothing-convolution}
	设 \(G\) 为有限维实紧 Lie 群，\(\mu_G\) 为 \(G\) 上满足 \(\mu_G(G)=1\) 的左右不变 Haar 测度（简记为 \(\mu\)）。对 \(k\in C(G)\)，按\eqref{g1:eq:left-convolution-kernel}的约定，对 \(f\in C(G)\) 定义
	\[
		T_kf(x)=\int_G k(y)f(y^{-1}x)\,\mathrm{d}\mu(y)
		=\int_G k(xz^{-1})f(z)\,\mathrm{d}\mu(z).
	\]
	则 \(T_k\) 唯一延拓为有界映射 \(T_k:L^2(G)\to C(G)\)，并满足
	\begin{equation}\label{g1:eq:convolution-l2-sup-bound}
		\|T_kf\|_\infty\le\|k\|_2\,\|f\|_2.
	\end{equation}
	若 \(k\) 光滑，则输出光滑；在任一闭包紧含于坐标域的坐标片上，每个多重指标 \(\alpha\) 都有常数 \(C_{k,\alpha}\)，使 \(\|\partial^\alpha T_kf\|_\infty\le C_{k,\alpha}\|f\|_2\)。此外，右正则表示 \(R_gf(x)=f(xg)\) 在 \(L^2(G)\) 上酉且强连续，并满足
	\[
		T_kR_g=R_gT_k\qquad(g\in G).
	\]
\end{proposition}
\begin{proof}
	对连续 \(f\)，在第二个积分中用 Cauchy--Schwarz 不等式。函数 \(z\mapsto k(xz^{-1})\) 的 \(L^2\) 范数等于 \(\|k\|_2\)，由此得到\eqref{g1:eq:convolution-l2-sup-bound}。随 \(x\) 变化，该核在一致范数、因而也在 \(L^2\) 范数中连续，所以输出连续。若连续函数列在 \(L^2\) 中趋于 \(f\)，其卷积由此估计一致 Cauchy；\(C(G)\) 对一致范数完备，故得到唯一的连续输出及同一估计。

	若 \(k\) 光滑，在 \(x\) 的坐标小片上，核 \(k(xz^{-1})\) 的每阶坐标导数随 \((x,z)\) 连续，在闭坐标片乘紧群上有界。微积分基本定理使核的差商在 \(z\) 上一致趋于相应导数，因而也在 \(L^2\) 中趋于该导数。与固定 \(f\in L^2\) 配对是连续线性运算，所以可以反复对核求导，得到
	\[
		\partial_x^\alpha T_kf(x)
		=\int_G\partial_x^\alpha\bigl(k(xz^{-1})\bigr)f(z)\,\mathrm{d}\mu(z).
	\]
	对一般 \(f\)，右侧仍按连续延拓理解。Cauchy--Schwarz 不等式及核导数的紧集界给出所需估计，也给出每阶导数的连续性。

	\cref{g1:prop:regular-representations}已在 \(C(G)\) 上验证 \(R_gR_h=R_{gh}\) 及积分内积不变，所以 \(R_g\) 唯一酉延拓到 \(L^2\)。对连续 \(f\)，紧性给出 \(\|R_gf-f\|_\infty\to0\)（\(g\to e\)）；对一般 \(f\) 用连续近似和酉性，得到 \(\|R_gf-f\|_2\to0\)。群律把单位元处的结论移到任意 \(g\)，这正是强连续性。最后直接计算
	\[
		T_kR_gf(x)=\int_Gk(y)f(y^{-1}xg)\,\mathrm{d}\mu(y)
		=R_gT_kf(x).
	\]
	连续函数的稠密性将等式延拓到 \(L^2\)。
\end{proof}

\cref{g1:lem:compact-group-approximate-identity}与\cref{g1:prop:smoothing-convolution}合起来说明 \(C^\infty(G)\) 在 \(C(G)\) 中一致稠密，也在 \(L^2(G)\) 中稠密：对连续 \(f\)，\(T_nf\) 就是所需的光滑近似。\(T_k\) 与右正则作用交换；一般核并不使它与左正则作用交换。

\subsection{卷积算子与紧自伴谱分解}
\label{g1:subsec:2002:03}
若一个有界线性算子把任意有界序列映为具有收敛子列的序列，则称它为\term[jinsuanzi]{紧算子}[compact operator]。有限秩算子具有这个性质，因为有界像留在有限维空间内。与有限维矩阵相比，紧算子仍可能有无限多个非零特征值，但这些特征值只能趋向零。我们先把平滑卷积放入这个范围，再证明准确需要的谱结论。

\begin{lemma}[连续核的紧性]\label{g1:lem:continuous-kernel-compact}
	设 \(G\) 为有限维实紧 Lie 群，\(\mu_G\) 为 \(G\) 上满足 \(\mu_G(G)=1\) 的左右不变 Haar 测度（简记为 \(\mu\)），\(K\in C(G\times G)\)。积分算子
	\[
		Af(x)=\int_GK(x,z)f(z)\,\mathrm{d}\mu(z)
	\]
	延拓为 \(L^2(G)\) 上的紧算子。若 \(K(x,z)=\overline{K(z,x)}\)，则它自伴，即 \(\langle Af,h\rangle_2=\langle f,Ah\rangle_2\)。特别地，对实值光滑函数 \(k_n\) 且 \(k_n(x^{-1})=k_n(x)\)，令 \(T_nf(x)=\int_G k_n(y)f(y^{-1}x)\,\mathrm{d}\mu(y)\)，则 \(T_n\) 紧、自伴，并与所有右平移 \(R_gf(x)=f(xg)\) 交换。
\end{lemma}
\begin{proof}
	Cauchy--Schwarz 不等式给出 \(\|Af\|_2\le\|K\|_\infty\|f\|_2\)，所以算子可唯一延拓。固定 \(\varepsilon>0\)。核的紧集一致连续性给出有限个开集 \(U_j\) 和点 \(x_j\)，使它们覆盖 \(G\)，且
	\[
		|K(x,z)-K(x_j,z)|<\varepsilon
		\qquad(x\in U_j,\ z\in G).
	\]
	这里的控制对全部 \(z\) 一致；可逐点取积邻域后利用第二个因子的紧性，再用第一个因子的紧性取得有限覆盖。取从属单位分解 \(\phi_j\)，令
	\[
		K_\varepsilon(x,z)=\sum_j\phi_j(x)K(x_j,z).
	\]
	\cref{g1:thm:smooth-partition-unity}保证可以选光滑的非负 \(\phi_j\)，且和为 \(1\)。因此 \(\|K-K_\varepsilon\|_\infty\le\varepsilon\)。对应算子
	\[
		A_\varepsilon f=\sum_j\phi_j
		\int_GK(x_j,z)f(z)\,\mathrm{d}\mu(z)
	\]
	的像包含于 \(\operatorname{span}\{\phi_j\}\)，是有限秩算子，且 \(\|A-A_\varepsilon\|\le\varepsilon\)。这就给出了连续核的有限秩逼近。

	为说明算子范数极限仍紧，取有界序列 \(f_\ell\)、\(\|f_\ell\|_2\le M\)，并取 \(\varepsilon_j\to0\)。依次选子列使每个 \(A_{\varepsilon_j}f_\ell\) 收敛，再取对角子列。对该子列中的两项，
	\[
		\|Af_\ell-Af_m\|_2
		\le2M\varepsilon_j+
		\|A_{\varepsilon_j}f_\ell-A_{\varepsilon_j}f_m\|_2.
	\]
	先取大 \(j\)，再取大 \(\ell,m\)，可见其像 Cauchy；Hilbert 空间完备性给出收敛。因此 \(A\) 紧。

	若核满足所述对称性，对连续 \(f,h\) 交换紧空间上的两次积分，就有
	\[
		\langle Af,h\rangle_2
		=\int_G\int_GK(x,z)f(z)\overline{h(x)}\,
		\mathrm{d}\mu(z)\,\mathrm{d}\mu(x)
		=\langle f,Ah\rangle_2.
	\]
	由有界性及稠密性，这对一般 \(L^2\) 元素也成立。\(T_n\) 的核为 \(k_n(xz^{-1})\)；它实值、取逆不变，故满足对称性。交换性已经在前一命题证明。
\end{proof}

以下引理只使用 Hilbert 空间中的正交与紧性。它也解释为何必须保留零特征空间：紧算子的核可以无限维，只有非零特征空间必定有限维。

\begin{lemma}[紧自伴谱分解]\label{g1:lem:compact-selfadjoint-spectral}
	设 \(H\) 为复 Hilbert 空间，\(A:H\to H\) 为紧自伴算子。它的非零特征值均为实数，每个对应特征空间 \(E_\lambda\) 有限维，不同特征空间正交。非零特征值至多可数；若有无限多个，把它们依绝对值非增次序列出为 \(\lambda_1,\lambda_2,\ldots\)，则 \(\lambda_j\to0\)。令 \(P_j\) 为到 \(E_{\lambda_j}\) 的正交投影，则
	\[
		H=\ker A\mathbin{\widehat\oplus}
		\mathop{\widehat\bigoplus}_{j}E_{\lambda_j},
		\qquad Af=\sum_j\lambda_jP_jf.
	\]
	直和与级数按 Hilbert 范数理解，而且 \(\sum_{j\le N}\lambda_jP_j\to A\) 按算子范数成立。若非零特征值只有有限多个，取有限和。
\end{lemma}
\begin{proof}
	证明先寻找绝对值最大的特征值，然后在它的正交补中重复。关键是让极值序列真正收敛；这一点由紧性承担。

	先证明任意有界自伴算子都有
	\[
		\|A\|=\sup_{\|u\|=1}|\langle Au,u\rangle|.
	\]
	记右端为 \(M\)，则 \(M\le\|A\|\)。反向，对单位向量 \(u,v\)，把 \(v\) 乘一个单位复数，使 \(\langle Au,v\rangle\) 为非负实数。记 \(q(w)=\langle Aw,w\rangle\)，自伴性使它为实数，而极化给出
	\[
		|\langle Au,v\rangle|
		=\tfrac14\bigl(q(u+v)-q(u-v)\bigr)
		\le\tfrac M4(\|u+v\|^2+\|u-v\|^2)=M.
	\]
	取 \(v=Au/\|Au\|\)（\(Au\ne0\) 时），得到 \(\|Au\|\le M\)，故等式成立。

	现在假设 \(A\ne0\)。存在 \(\lambda\in\{\|A\|,-\|A\|\}\) 及单位向量列 \(u_n\)，使 \(\langle Au_n,u_n\rangle\to\lambda\)。于是
	\[
		\|(A-\lambda I)u_n\|^2
		=\|Au_n\|^2-2\lambda\langle Au_n,u_n\rangle+\lambda^2
		\le2\lambda^2-2\lambda\langle Au_n,u_n\rangle\longrightarrow0.
	\]
	紧性使 \(Au_n\) 有收敛子列；\(\lambda\ne0\) 及上式使相应 \(u_n\) 也收敛到单位向量 \(u\)，且 \(Au=\lambda u\)。因此每个非零紧自伴算子都具有绝对值等于算子范数的特征值。

	若 \(Au=\lambda u\)、\(u\ne0\)，则 \(\lambda\|u\|^2=\langle Au,u\rangle\) 为实数，所以 \(\lambda\) 实。若另有 \(Av=\eta v\)，自伴性给出 \((\lambda-\eta)\langle u,v\rangle=0\)，所以不同特征空间正交。非零特征空间不能无限维，否则逐次作有限维正交投影就能在其中构造一列正交单位向量 \(u_n\)；其像 \(\lambda u_n\) 两两相距 \(\sqrt2|\lambda|\)，违背紧性。同理，对每个 \(\delta>0\)，绝对值至少为 \(\delta\) 的特征值对应空间之和必为有限维；否则取一列来自这些空间的正交单位特征向量，其像也没有 Cauchy 子列。

	取刚找到的极值特征值 \(\lambda_1\)，移去整个有限维空间 \(E_{\lambda_1}\)。其正交补闭且被 \(A\) 保持，因为对 \(v\perp E_{\lambda_1}\)、\(u\in E_{\lambda_1}\)，有 \(\langle Av,u\rangle=\langle v,Au\rangle=0\)。限制算子仍紧、自伴。若限制非零，便在其中选 \(\lambda_2\)，继续此过程。它或者有限步后留下零算子，或者产生无限个特征空间。后一情形中 \(|\lambda_j|\) 非增；前段的 \(\delta\) 论证说明 \(\lambda_j\to0\)。在前 \(N\) 个空间的正交补上，算子范数恰为下一步的 \(|\lambda_{N+1}|\)。故
	\begin{equation}\label{g1:eq:compact-spectrum-tail}
		\left\|A-\sum_{j=1}^N\lambda_jP_j\right\|
		=|\lambda_{N+1}|\longrightarrow0.
	\end{equation}
	这还说明没有遗漏非零特征值，因为遗漏的非零特征向量会留在每一个这样的正交补中，与剩余算子范数趋于零矛盾。

	最后，对任意 \(f\in H\)，有限维正交投影的勾股等式给出 \(\sum_{j\le N}\|P_jf\|^2\le\|f\|^2\)。因此正交级数 \(\sum_jP_jf\) Cauchy，由完备性收敛到某个 \(v\)。余项 \(w=f-v\) 与全部 \(E_{\lambda_j}\) 正交，\eqref{g1:eq:compact-spectrum-tail}遂给出 \(Aw=0\)。反过来，若 \(Aw=0\)，则 \(\lambda_j\langle w,u\rangle=\langle w,Au\rangle=\langle Aw,u\rangle=0\)，所以 \(w\) 与全部非零特征空间正交。这证明所述正交分解及唯一性。算子级数的范数收敛已经由尾项估计给出。
\end{proof}

逼近单位中 \(k_n\ge0\) 表示对函数值作非负加权平均；它并不意味着自伴算子 \(T_n\) 的每个特征值非负。刚才的谱引理允许正、负特征值，后面的论证只需排除零值，不需要附加算子正性。

\subsection{非零特征空间的有限维不变性}
\label{g1:subsec:2002:04}
卷积与右平移交换，因此每个特征空间都在右作用下不变。对非零特征值，特征函数又由平滑卷积得到；选取有限个求值点，就能把右平移写成随群参数光滑变化的矩阵。这样，谱引理给出的有限维空间便承载一个 Lie 群表示。

\begin{proposition}[卷积特征空间上的表示]\label{g1:prop:convolution-eigenspace-representation}
	设 \(G\) 为有限维实紧 Lie 群，\(\mu_G\) 为 \(G\) 上满足 \(\mu_G(G)=1\) 的左右不变 Haar 测度（简记为 \(\mu\)），\(k\in C^\infty(G)\) 满足 \(k(x^{-1})=\overline{k(x)}\)。令 \(T_kf(x)=\int_G k(y)f(y^{-1}x)\,\mathrm{d}\mu(y)\)、\(R_gf(x)=f(xg)\)，并以 \(\mathcal R(G)\) 记有限维复光滑表示的矩阵系数的线性张成空间。若 \(\lambda\ne0\) 是 \(T_k\) 在 \(L^2(G)\) 上的特征值，则
	\[
		E_\lambda=\{f\in L^2(G):T_kf=\lambda f\}
	\]
	为有限维右正则不变空间，每个元素都有唯一的光滑函数代表。限制作用 \(R|_{E_\lambda}\) 是光滑酉表示；\(E_\lambda\) 中的每个函数属于 \(\mathcal R(G)\)，并能写成不可约酉表示的矩阵系数的有限线性组合。
\end{proposition}
\begin{proof}
	连续核引理使 \(T_k\) 紧、自伴，谱引理给出 \(E_\lambda\) 有限维。若 \(f\in E_\lambda\)，则 \(f=\lambda^{-1}T_kf\) 在 \(L^2\) 中成立，右侧由\cref{g1:prop:smoothing-convolution}为光滑函数。连续函数到 \(L^2\) 的嵌入是单射，故这个光滑代表唯一。

	因 \(T_kR_g=R_gT_k\)，有 \(T_k(R_gf)=\lambda R_gf\)，所以空间被 \(R_g\) 保持；\(R_{g^{-1}}\) 给出逆作用。积分内积被 \(R_g\) 保持，故限制表示酉。\cref{g1:prop:finite-regular-subspace-coefficients}适用于这个有限维的 \(R\) 不变光滑函数空间：该命题用有限组点评价给空间取坐标，从而证明 \(g\mapsto R_g|_{E_\lambda}\) 光滑，并且
	\[
		f(g)=(R_gf)(e)=\operatorname{ev}_e(R_gf).
	\]
	评价泛函 \(\operatorname{ev}_e\) 在有限维 \(E_\lambda\) 上是普通线性泛函，所以这正是一个矩阵系数。最后，\cref{g1:prop:compact-unitarization-complete-reducibility}把该有限维酉表示正交分解为有限个不可约子表示；将向量 \(f\) 和评价泛函限制到各分量，便得到所述有限线性组合。
\end{proof}

在无限维 \(L^2\) 上，点评价一般不是连续线性泛函，因此不能一开始就把整个正则表示的“系数”写成 \(f(g)\)。先平滑，再进入有限维特征空间，正好消除了这个障碍。零特征空间没有 \(f=\lambda^{-1}T_kf\) 的表达式，也就不能由同一论证获得光滑性或有限维性。

\subsection{有限谱逼近与矩阵系数分离}
\label{g1:subsec:2002:05}
一个固定卷积算子的零空间可能很大，所产生的有限维表示也可能漏掉群中的某些运动。让核收缩到单位元，卷积就趋近恒等；这样不同点最终不能被所有产生的表示同时忽略。要把这个直觉落实到点评价，须先把谱引理中的 \(L^2\) 逼近改进为一致逼近。

\begin{lemma}[再卷积后的有限谱逼近]\label{g1:lem:smoothed-spectral-approximation}
	设 \(G\) 为有限维实紧 Lie 群，\(\mu_G\) 为 \(G\) 上满足 \(\mu_G(G)=1\) 的左右不变 Haar 测度（简记为 \(\mu\)），固定一个光滑实值函数 \(k_n\ge0\)，满足 \(\int_Gk_n\,\mathrm{d}\mu=1\) 及 \(k_n(x^{-1})=k_n(x)\)，以
	\[
		Tf(x)=T_nf(x)=\int_G k_n(y)f(y^{-1}x)\,\mathrm{d}\mu(y)
	\]
	记其卷积算子。将 \(T\) 的各个不同非零特征值列为 \(\lambda_j\)，相应特征空间的正交投影记为 \(P_j\)，并以 \(\mathcal R(G)\) 记有限维复光滑表示的矩阵系数的线性张成空间。对任意 \(f\in L^2(G)\)，有限和
	\[
		h_N=\sum_{j=1}^N\lambda_j^2P_jf
	\]
	属于 \(\mathcal R(G)\)，并一致收敛于 \(T^2f\)。若特征值只有有限多个，取全部项即得等式。事实上，对固定的任意阶数，\(h_N\to T^2f\) 的各阶导数也一致收敛。
\end{lemma}
\begin{proof}
	由\cref{g1:lem:compact-selfadjoint-spectral}，
	\[
		S_Nf=\sum_{j=1}^N\lambda_jP_jf\longrightarrow Tf
	\]
	按 \(L^2\) 范数成立。再作用一次 \(T\)，由\eqref{g1:eq:convolution-l2-sup-bound}得到
	\begin{equation}\label{g1:eq:smoothed-spectral-sup-error}
		\|TS_Nf-T^2f\|_\infty
		\le\|k_n\|_2\,\|S_Nf-Tf\|_2\longrightarrow0.
	\end{equation}
	每个 \(P_jf\in E_{\lambda_j}\)，故 \(TS_Nf=h_N\)；前一命题又保证每项为代表函数。因此这是一列有限矩阵系数和的一致逼近。对各阶导数，使用\cref{g1:prop:smoothing-convolution}中的相应估计代替一致范数估计即可。此处 \(n\) 固定，允许估计常数依赖该核。
\end{proof}

\begin{corollary}[矩阵系数分离点]\label{g1:cor:matrix-coefficients-separate-points}
	设 \(G\) 为有限维实紧 Lie 群。若 \(p,q\in G\)、\(p\ne q\)，则存在有限维不可约光滑酉表示 \(\pi\)，使 \(\pi(p)\ne\pi(q)\)。
\end{corollary}
\begin{proof}
	取在 \(p\) 附近等于 \(1\)、支集不含 \(q\) 的光滑截断函数 \(f\)，则 \(f(p)=1\)、\(f(q)=0\)。逼近单位的压缩估计给出
	\[
		\|T_n^2f-f\|_\infty
		\le\|T_n(T_nf-f)\|_\infty+\|T_nf-f\|_\infty
		\le2\|T_nf-f\|_\infty\longrightarrow0.
	\]
	先选 \(n\) 使这个误差小于 \(1/8\)，再用前一引理选 \(h\in\mathcal R(G)\)，使 \(\|h-T_n^2f\|_\infty<1/8\)。于是 \(|h(p)-h(q)|>1/2\)。将 \(h\) 写成不可约系数的有限线性组合，其中至少一项在 \(p,q\) 的值不同；相应表示矩阵也不同。
\end{proof}

该推论给出了紧群的有限维表示足以检测群元素的结论。它尚未指定一张同时区分所有元素的有限维矩阵表示；这里的量词是对每一对不同点可以选择一个表示。

\subsection{矩阵系数的稠密性}
\label{g1:subsec:2002:06}
点分离的证明其实提供了一个更强的逼近过程：先选择核使卷积靠近目标函数，再选择有限谱截断使再卷积靠近有限矩阵系数和。这个过程直接给出 Peter--Weyl 定理的一致稠密性。\footnote{Fritz Peter（1899--1949），德国数学家，与 Weyl 共同证明了紧群矩阵系数的完备性定理。}

\begin{theorem}[Peter--Weyl 一致稠密性]\label{g1:thm:peter-weyl-density}
	设 \(G\) 为有限维实紧 Lie 群，\(\mu_G\) 为 \(G\) 上满足 \(\mu_G(G)=1\) 的左右不变 Haar 测度（简记为 \(\mu\)）。对任意 \(f\in C(G)\) 和任意 \(\varepsilon>0\)，存在有限个不可约光滑酉表示 \(\pi_1,\ldots,\pi_s\) 及其矩阵系数的线性组合 \(h\)，使
	\[
		\|f-h\|_\infty<\varepsilon.
	\]
	等价地，\(\mathcal R(G)\) 在 \(C(G)\) 中对一致范数稠密。它也在 \(L^2(G)\) 中稠密。
\end{theorem}
\begin{proof}
	由逼近单位引理，选 \(n\) 使 \(2\|T_nf-f\|_\infty<\varepsilon/2\)，于是 \(\|T_n^2f-f\|_\infty<\varepsilon/2\)。固定这个 \(n\)，\cref{g1:lem:smoothed-spectral-approximation}给出 \(h\in\mathcal R(G)\)，使 \(\|h-T_n^2f\|_\infty<\varepsilon/2\)。三角不等式给出结论。\(\mathcal R(G)\) 中每个函数均为有限维表示的系数有限和，再用完全可约性就能使出现的表示不可约。

	对 \(f\in L^2(G)\)，先取 \(u\in C(G)\) 使 \(\|f-u\|_2<\varepsilon/2\)，再取 \(h\in\mathcal R(G)\) 使 \(\|u-h\|_\infty<\varepsilon/2\)。因为 \(\|u-h\|_2\le\|u-h\|_\infty\)，得到 \(\|f-h\|_2<\varepsilon\)，证明 \(L^2\) 稠密性。
\end{proof}

两次选择的次序不可颠倒：\(\|k_n\|_2\) 随核支集收缩可能变大，\eqref{g1:eq:smoothed-spectral-sup-error}中的常数没有统一上界。先固定足够集中的核，再把该核的谱截断误差压小，才能确保最终的一致误差。再卷积给出的有界映射 \(L^2\to C\) 将这里的 \(L^2\) 误差转成一致误差。

\subsection{Peter–Weyl 分解与收敛}
\label{g1:subsec:2002:07}
现在可以把上一节的正交性与本节的稠密性合在一起。正交性告诉我们不同不可约表示的系数互不干扰；稠密性则排除了整个函数空间中仍有一块与所有系数正交的部分。

\begin{theorem}[Peter--Weyl 正交分解]\label{g1:thm:peter-weyl-hilbert}
	设 \(G\) 为有限维实紧 Lie 群，\(\mu_G\) 为 \(G\) 上满足 \(\mu_G(G)=1\) 的左右不变 Haar 测度（简记为 \(\mu\)），\(\widehat G\) 表示其有限维不可约光滑酉表示的等价类，每类选一个代表 \(\pi:G\to\operatorname{U}(V_\pi)\) 及标准正交基。记 \(d_\pi=\dim_{\mathbb C}V_\pi\)、\(\pi_{ij}(g)=\langle\pi(g)e_j,e_i\rangle\)，并令
	\[
		\mathcal H_\pi=\operatorname{span}_{\mathbb C}
		\{\pi_{ij}:1\le i,j\le d_\pi\}.
	\]
	则 \(\widehat G\) 至多可数，\(\dim \mathcal H_\pi=d_\pi^2\)，并且
	\begin{equation}\label{g1:eq:peter-weyl-hilbert-sum}
		L^2(G)=\mathop{\widehat\bigoplus}_{[\pi]\in\widehat G}\mathcal H_\pi.
	\end{equation}
	函数族 \(\{\sqrt{d_\pi}\,\pi_{ij}\}_{\pi,i,j}\) 是 \(L^2(G)\) 的标准正交基。对 \(f\in L^2(G)\)，令 \(a_{\pi,ij}=\langle f,\pi_{ij}\rangle_2\)，则
	\[
		f=\sum_{[\pi]\in\widehat G}d_\pi
		\sum_{i,j=1}^{d_\pi}a_{\pi,ij}\pi_{ij}
	\]
	按 \(L^2\) 范数收敛，并满足 Parseval 等式：\footnote{Parseval，Marc-Antoine Parseval des Chênes（1755--1836），法国数学家，研究无穷级数与偏微分方程；以其名字命名的恒等式现用于 Fourier 分析。}
	\begin{equation}\label{g1:eq:peter-weyl-parseval}
		\|f\|_2^2=\sum_{[\pi]\in\widehat G}d_\pi
		\sum_{i,j=1}^{d_\pi}|a_{\pi,ij}|^2.
	\end{equation}
	在右正则表示中，每个 \(\pi\) 出现的重数为 \(d_\pi\)。
\end{theorem}
\begin{proof}
	\cref{g1:thm:schur-orthogonality}给出不同代表的系数正交，同一表示内部有
	\[
		\langle\pi_{ij},\pi_{k\ell}\rangle_2
		=\frac{\delta_{ik}\delta_{j\ell}}{d_\pi}.
	\]
	故各 \(\mathcal H_\pi\) 的维数为 \(d_\pi^2\)，缩放后的系数构成标准正交族。\cref{g1:thm:peter-weyl-density}说明它们的线性张成在 \(L^2\) 中稠密，因此这是完备的标准正交族。

	先由 \(L^2(G)\) 的可分性证明展开项的可数性。紧流形 \(G\) 正则、第二可数，\cref{g1:thm:urysohn-metrization}给出与拓扑相容的距离 \(\delta\)。对每个正整数 \(n\)，取有限个直径小于 \(1/n\) 的开集覆盖 \(G\)，并取从属光滑单位分解 \(\phi_{n,j}\)。所有有限和 \(\sum_jq_j\phi_{n,j}\)，其中 \(q_j\in\mathbb Q+\mathrm{i}\mathbb Q\)，随 \(n\) 变化构成可数集合。它在 \(C(G)\) 中一致稠密：给定连续 \(f\)，用紧性取得一致连续性，再在各开集选一点 \(x_{n,j}\)，取 \(q_j\) 接近 \(f(x_{n,j})\)；单位分解的非负性和总和 \(1\) 便把误差控制为各开集内的振幅加有理近似误差。于是 \(L^2(G)\) 也可分，即有可数稠密子集。

	可分 Hilbert 空间的标准正交族至多可数：在可数稠密集内为每个单位向量选一个距离小于 \(1/3\) 的点，不同正交单位向量相距 \(\sqrt2\)，所以不能选到同一个点。这个单射证明系数族可数，也就证明 \(\widehat G\) 至多可数。

	将规范化系数枚举为 \(u_1,u_2,\ldots\)。有限正交投影给出
	\[
		\left\|f-\sum_{j=1}^N\langle f,u_j\rangle_2u_j\right\|_2^2
		=\|f\|_2^2-\sum_{j=1}^N|\langle f,u_j\rangle_2|^2.
	\]
	右侧非负，故系数平方和有界，正交级数按范数收敛。其余项与所有 \(u_j\) 正交；由线性张成稠密，它也与自身正交，因而为零。于是得到展开和 Parseval 等式。正交级数尾项的范数平方就是尾部系数平方和，故任何重排以及按有限表示组逐步穷尽的截断，都给出同一个 \(L^2\) 极限。\eqref{g1:eq:peter-weyl-hilbert-sum}中的帽号表示允许这种平方可和的无限正交和；不带帽的代数直和只允许有限和。

	最后固定 \(\pi\) 及行指标 \(i\)。矩阵乘法给出
	\[
		R_g\pi_{ij}(x)=\pi_{ij}(xg)
		=\sum_k\pi_{ik}(x)\pi_{kj}(g).
	\]
	因此映射 \(e_j\mapsto\sqrt{d_\pi}\,\pi_{ij}\) 把 \(V_\pi\) 酉等变地送到固定第 \(i\) 行的系数空间。不同的 \(i\) 对应正交子空间，合起来恰为 \(\mathcal H_\pi\)。所以 \(\mathcal H_\pi\) 是 \(d_\pi\) 个 \(\pi\) 的正交和；不同等价类的系数空间互相正交且已穷尽 \(L^2\)，给出所述重数。
\end{proof}

若希望把标量系数整理成矩阵，可以定义
\[
	\widehat f(\pi)=\int_G f(g)\pi(g^{-1})\,\mathrm{d}\mu(g).
\]
\symidx{\(\widehat f(\pi)\)：紧群函数的 Fourier 矩阵}对 \(L^2\) 元素，这只是有限个连续内积泛函组成的矩阵，其第 \((j,i)\) 项为 \(a_{\pi,ij}\)。于是展开和范数公式写成
\[
	f=\sum_{[\pi]}d_\pi\operatorname{tr}\bigl(\widehat f(\pi)\pi(\,\cdot\,)\bigr),
	\qquad \|f\|_2^2=\sum_{[\pi]}d_\pi\|\widehat f(\pi)\|_{\mathrm{HS}}^2.
\]
这里 Hilbert--Schmidt 范数的平方就是矩阵所有条目的绝对值平方和。转置的指标顺序来自 \(\pi(g^{-1})=\pi(g)^*\)，不能把它无说明地换成 \(\pi(g)\)。

\begin{example}[圆周上的展开]\label{g1:ex:peter-weyl-circle}
	对 \(S^1\)，不可约酉表示为 \(\pi_m(e^{\mathrm{i}t})=e^{\mathrm{i}mt}\)、\(m\in\mathbb Z\)。交换群的不可约复表示是一维的：每个表示矩阵都与整个表示交换，Schur 引理使它为标量，不可约性再使维数为 \(1\)。一维光滑同态拉回到 \(\mathbb R\) 后是 \(e^{\mathrm{i}at}\)，周期 \(2\pi\) 要求 \(a\in\mathbb Z\)；这里用了单参数子群的唯一性。于是上一定理成为
	\[
		f=\sum_{m\in\mathbb Z}\widehat f(m)e^{\mathrm{i}mt},
		\qquad \|f\|_2^2=\sum_{m\in\mathbb Z}|\widehat f(m)|^2,
	\]
	其中 \(\widehat f(m)=(2\pi)^{-1}\int_0^{2\pi}f(e^{\mathrm{i}t})e^{-\mathrm{i}mt}\,\mathrm{d}t\)。第一个等式在本定理中指 \(L^2\) 收敛；它并未给一般 \(L^2\) 函数指定逐点值。
\end{example}

对光滑函数，前面的有核逼近还能保持任意有限阶导数。取 \(\mathfrak g\) 的一组基，对应左不变场 \(X_1^L,\ldots,X_d^L\)，并令
\[
	p_m(f)=\max_{0\le r\le m}\ \max_{i_1,\ldots,i_r}
	\|X_{i_1}^L\cdots X_{i_r}^Lf\|_\infty.
\]
\(r=0\) 项为 \(\|f\|_\infty\)；若 \(G\) 零维，则统一取 \(p_m(f)=\|f\|_\infty\)。这些半范数定义 \(C^\infty\) 收敛：每个固定阶数的导数都一致收敛。左不变场给出整体标架，坐标导数与这些场之间通过光滑系数互相表示；在有限个较小坐标片的闭包上系数有界，故这与通常坐标中的定义一致。

\begin{proposition}[光滑函数的系数逼近]\label{g1:prop:peter-weyl-smooth-density}
	设 \(G\) 为有限维实紧 Lie 群，固定 \(\mathfrak g=\operatorname{Lie}(G)\) 的一组基及对应的左不变向量场 \(X_1^L,\ldots,X_d^L\)。对整数 \(r\ge0\)，令
	\[
		p_r(f)=\max_{\substack{0\le k\le r\\1\le i_1,\ldots,i_k\le d}}
		\|X_{i_1}^L\cdots X_{i_k}^Lf\|_\infty,
	\]
	其中 \(k=0\) 项为 \(\|f\|_\infty\)，\(d=0\) 时只取此项。以 \(\mathcal R(G)\) 记有限维复光滑表示的矩阵系数的线性张成空间。则对任意 \(f\in C^\infty(G)\)，存在 \(h_m\in\mathcal R(G)\)，使对每个固定 \(r\)，均有 \(p_r(h_m-f)\to0\)。
\end{proposition}
\begin{proof}
	左不变场的流为 \(x\mapsto x\exp_G(tX)\)。对光滑 \(f\)，在左卷积公式中沿这个流求导，紧参数积分引理允许把求导移入积分，得到 \(X^LT_nf=T_nX^Lf\)。对任意有限次复合重复此等式，逼近单位的一致收敛给出 \(p_r(T_nf-f)\to0\)。而每个导数分量上 \(T_n\) 都为一致范数压缩，故 \(p_r(T_n^2f-f)\to0\) 也成立。

	给定 \(m\)，先选 \(n\) 使 \(p_m(T_n^2f-f)<1/(2m)\)。固定此核后，\cref{g1:lem:smoothed-spectral-approximation}的各阶导数估计允许选有限谱和 \(h_m\in\mathcal R(G)\)，使 \(p_m(h_m-T_n^2f)<1/(2m)\)。于是 \(p_m(h_m-f)<1/m\)。对每个固定 \(r\)，当 \(m\ge r\) 时 \(p_r\le p_m\)，所以得到所需收敛。
\end{proof}

上述三种结论承担不同任务。正交展开直接由积分系数确定，其基本保证是 \(L^2\) 收敛；一致或光滑逼近则使用先选择平滑核、再选择谱截断的过程。\cref{g1:tab:peter-weyl-convergence}将逼近对象、构造方式与所用的范数对应起来。
\begin{table}[htbp]
	\centering
	\caption{Peter--Weyl 逼近的三种收敛意义}
	\label{g1:tab:peter-weyl-convergence}
	\begin{tabularx}{\textwidth}{@{}lXX@{}}
		\toprule
		目标函数 & 近似的选择 & 本节得到的收敛\\
		\midrule
		\(f\in L^2(G)\) & 正交系数展开，有限项逐步穷尽 & \(L^2\) 范数收敛及 Parseval 等式\\
		\(f\in C(G)\) & 先固定逼近单位核，再截断其非零谱并再卷积 & 函数值一致收敛\\
		\(f\in C^\infty(G)\) & 对逐渐增大的导数阶数作上述两次选择 & 每个固定阶数的导数一致收敛\\
		\bottomrule
	\end{tabularx}
\end{table}
一致稠密性本身不能推出任意顺序的正交部分和一致收敛，光滑逼近命题也没有对任意谱排列给出导数估计。对这类具体的求和方式，需要相应的核界或系数衰减来控制误差及其导数。

% !TeX root = ../../Book1.tex
% 纲要标题：### 20.3 从表示到齐性向量丛
% 纲要位置：几何学计划纲要.md，第 905 行。
% 正文状态：初稿；尚未正式审核。

% 纲要内容开始
% 
% * 稳定子表示与关联丛
% * 截面的等变函数模型
% * 截面表示与 Frobenius 互反
% * 复射影直线上的齐性线丛
% 
% 纲要内容结束

\section{从表示到齐性向量丛}
\label{g1:sec:2003}

一个群的表示在固定向量空间中改变向量；一个齐性空间上的向量丛则让向量空间随底点移动。两者的联系来自稳定子：先规定稳定子怎样作用于一个基点上的纤维，再用全群把这条纤维搬到各处。所得截面可以重新写成群上的函数，群表示又表现为函数的平移。本节使用\cref{g1:exm:homogeneous-principal-bundle,g1:thm:associated-vector-bundle}中的右主丛与关联丛约定，重点说明齐性空间额外提供的左作用以及它怎样控制截面。

\subsection{稳定子表示与关联丛}
\label{g1:subsec:2003:01}
设 \(G\) 为有限维实 Lie 群，\(H\subseteq G\) 为闭子群，\(o=eH\in G/H\)。左作用下，\(o\) 的稳定子就是 \(H\)。若向量丛 \(E\to G/H\) 带有覆盖左作用的光滑 \(G\) 作用，并且每个群元素在纤维之间作线性映射，那么 \(H\) 保持 \(E_o\)，从而给出一个纤维上的表示。这种带作用的向量丛称为\term[qixingxiangliangcong]{齐性向量丛}[homogeneous vector bundle]。下面的构造说明，这一个稳定子表示已经决定整个丛及其 \(G\) 作用。

\begin{proposition}[稳定子表示构造齐性向量丛]\label{g1:prop:homogeneous-associated-bundle}
	设 \(G\) 为有限维实 Lie 群，\(H\) 为闭子群，\(W\) 为有限维复向量空间，\(\rho:H\to\operatorname{GL}(W)\) 为光滑表示。则
	\[
		E_\rho=G\times_H W\longrightarrow G/H,\qquad
		[g,w]\longmapsto gH,
		\qquad [gh,w]=[g,\rho(h)w]
	\]
	是光滑复向量丛，且 \(a\cdot[g,w]=[ag,w]\) 定义了齐性向量丛结构。反之，每个有限秩光滑复齐性向量丛都与其基点纤维 \(W=E_o\) 的稳定子表示所构造的 \(G\times_H W\) 自然地 \(G\) 等变同构。实向量丛有相同结论。
\end{proposition}
\begin{proof}
	\(H\) 闭保证 \(G/H\) 是 Hausdorff 光滑流形，而且 \(G\to G/H\) 有光滑局部截面，是右主 \(H\) 丛。将\cref{g1:thm:associated-vector-bundle}用于这个主丛，得到所列关联丛；其等价关系来自右作用 \((g,w)\cdot h=(gh,\rho(h)^{-1}w)\)。左乘与这个右作用交换，因为 \(a(gh)=(ag)h\)，所以 \([g,w]\mapsto[ag,w]\) 良定义。在商映射之前它是光滑映射，局部平凡化又说明下降后的作用光滑，并且逐纤维复线性。

	反过来，设已有这样的 \(E\)，令 \(\rho(h)w=h\cdot w\)、\(w\in E_o\)。作用的光滑性及纤维线性性使 \(\rho\) 成为光滑线性表示。定义
	\[
		\Psi:G\times_H E_o\longrightarrow E,\qquad
		\Psi([g,w])=g\cdot w.
	\]
	它不随代表改变，因为 \((gh)\cdot w=g\cdot\rho(h)w\)。在 \(gH\) 上，群元素 \(g\) 将 \(E_o\) 线性同构地送到 \(E_{gH}\)，所以 \(\Psi\) 逐纤维双射。若 \(s:U\to G\) 是主丛局部截面，则 \(\Psi\) 的局部表达为 \((x,w)\mapsto s(x)\cdot w\)，其逆为 \(v\mapsto(p(v),s(p(v))^{-1}\cdot v)\)。这两式都光滑，故 \(\Psi\) 为丛同构。最后 \(\Psi(a\cdot[g,w])=ag\cdot w=a\cdot\Psi([g,w])\)，证明等变性。
\end{proof}

这里分类的是带指定 \(G\) 作用的丛。忘掉作用以后，不同的稳定子表示可能给出同构的普通向量丛。若 \(A:W\to W'\) 满足 \(A\rho(h)=\rho'(h)A\)，则 \([g,w]\mapsto[g,Aw]\) 定义一个等变丛映射；反过来，任何覆盖 \(G/H\) 恒等映射的等变丛映射都由其在 \(o\) 处的限制确定。这是因为它必须把 \(g\cdot w\) 送到 \(g\cdot A(w)\)。

\ProseParagraphBreak
一般表示给出的构造仍称关联丛；只有结构群的伴随表示给出的特例称为伴随丛。对于当前主 \(H\) 丛 \(G\to G/H\)，\cref{g1:def:adjoint-bundle}中的记号为
\[
	\operatorname{ad}(G\to G/H)=G\times_{\operatorname{Ad}|_H}\mathfrak h.
\]
这条丛的纤维是 \(\mathfrak h\)，不是 \(\mathfrak g\)。表示 \(\rho\) 决定了纤维中的对象及其转换规律，不能把一般关联构造中的表示都换成 \(\operatorname{Ad}\)。

\begin{example}\label{g1:exm:homogeneous-tangent-associated}
	设 \(G\) 为有限维实 Lie 群，\(H\) 为闭子群。\cref{g1:prop:homogeneous-tangent-isotropy}将 \(\mathrm{T}_o(G/H)\) 识别为 \(\mathfrak g/\mathfrak h\)，稳定子作用为 \(X+\mathfrak h\mapsto\operatorname{Ad}_hX+\mathfrak h\)。因此齐性切丛为
	\[
		\mathrm{T}(G/H)\cong G\times_H(\mathfrak g/\mathfrak h),\qquad
		[g,X+\mathfrak h]\longmapsto
		\left.\frac{\mathrm{d}}{\mathrm{d}t}\right|_{t=0}g\exp(tX)H.
	\]
	在基点，商空间同构由商映射的微分给出；在其他点用左平移搬运，正是前一命题的构造。若 \(G\) 紧，取伴随不变内积后，\(\mathfrak m=\mathfrak h^\perp\) 在 \(\operatorname{Ad}(H)\) 下不变，且 \(\mathfrak m\to\mathfrak g/\mathfrak h\) 是 \(H\) 等变线性同构，故也可用 \(G\times_H\mathfrak m\) 表示切丛。稳定子表示作用于商空间或正交补；它一般不同于作用在 \(\mathfrak h\) 上的伴随表示。
\end{example}

\subsection{截面的等变函数模型}
\label{g1:subsec:2003:02}
要在 \(gH\) 上指定截面值，可以先写成 \([g,f(g)]\)。同一个底点也能写成 \(ghH\)，所以函数值不能任意变化；它必须恰好补偿代表 \(g\) 的改变。这恢复了第十八章的等变函数关系，并且让截面变成可以直接平移、求导和积分的对象。

\begin{proposition}[齐性丛截面的函数模型]\label{g1:prop:equivariant-functions-sections}
	设 \(G\) 为有限维实 Lie 群，\(H\) 为闭子群，\(\rho:H\to\operatorname{GL}(W)\) 为有限维复光滑表示。令 \(E_\rho=G\times_H W\to G/H\)，采用 \([gh,w]=[g,\rho(h)w]\) 的商约定，并令
	\[
		C^\infty(G,W)_\rho
		=\{f\in C^\infty(G,W):f(gh)=\rho(h)^{-1}f(g)
		\;\text{对所有 }g\in G,h\in H\}.
	\]
	则线性映射
	\[
		f\longmapsto\sigma_f,\qquad
		\sigma_f(gH)=[g,f(g)]
	\]
	给出 \(C^\infty(G,W)_\rho\cong\Gamma^\infty(G/H,E_\rho)\)。若主丛局部截面满足 \(s_j=s_i h_{ij}\)，则相应局部分量 \(f_i=f\circ s_i\) 满足 \(f_i=\rho(h_{ij})f_j\)。
\end{proposition}
\begin{proof}
	这属于\cref{g1:prop:associated-sections}的情形：主丛为 \(G\to G/H\)，其右作用是 \(g\mapsto gh\)。具体地，
	\[
		[gh,f(gh)]=[g,\rho(h)f(gh)]=[g,f(g)],
	\]
	所以所给公式不依赖余类代表。在局部截面 \(s_i\) 下，\(\sigma_f(x)=[s_i(x),f(s_i(x))]\)，故截面光滑。

	反过来，对每个 \(g\)，映射 \(w\mapsto[g,w]\) 是到纤维 \((E_\rho)_{gH}\) 的线性同构，所以 \(\sigma(gH)=[g,f(g)]\) 唯一确定 \(f(g)\)。比较 \(g\) 与 \(gh\) 得到所述等变性；在 \(g=s_i(x)h\) 的局部坐标中，有 \(f(g)=\rho(h)^{-1}f_i(x)\)，所以它光滑。两种构造逐点互逆。最后，\(f_j(x)=f(s_i(x)h_{ij}(x))=\rho(h_{ij}(x))^{-1}f_i(x)\)，即所写的转换式。
\end{proof}

若 \(\rho\) 是平凡的一维表示，等变条件退化为 \(f(gh)=f(g)\)，也就是普通函数从 \(G\) 下降到 \(G/H\) 的条件。非平凡表示允许函数沿同一条纤维变化，但变化已经由 \(\rho\) 完全规定。\cref{g1:fig:associated-section-descent}把这个条件写成两个商映射之间的相容关系：函数与群元素配在一起后，才能下降为截面。

\input{Figures/Chapter20-AssociatedSection}

\subsection{截面表示与 Frobenius 互反}
\label{g1:subsec:2003:03}
群元素 \(a\) 把 \(x\) 处的纤维送到 \(ax\) 处。要得到一条新截面在 \(x\) 处的值，应先去 \(a^{-1}x\) 处读取原截面，再由 \(a\) 将这个值送回来。这给出公式
\[
	(U_a\sigma)(x)=a\cdot\sigma(a^{-1}x).
\]
逆元出现于底点，正是为了使结果仍落在 \(E_x\)，同时使 \(U_aU_b=U_{ab}\)。

\begin{proposition}[截面上的左表示]\label{g1:prop:section-left-representation}
	设 \(G\) 为有限维实 Lie 群，\(H\) 为闭子群，\(\rho:H\to\operatorname{GL}(W)\) 为有限维复光滑表示。令 \(E_\rho=G\times_H W\)，采用 \([gh,w]=[g,\rho(h)w]\) 的商约定。公式 \((U_a\sigma)(x)=a\cdot\sigma(a^{-1}x)\) 在 \(\Gamma^\infty(G/H,E_\rho)\) 上给出一个线性表示，其中总空间的作用为 \(a\cdot[g,w]=[ag,w]\)。以 \(\sigma_f(gH)=[g,f(g)]\) 识别满足 \(f(gh)=\rho(h)^{-1}f(g)\) 的光滑函数，即\cref{g1:prop:equivariant-functions-sections}的函数模型。在此模型中，它是
	\begin{equation}\label{g1:eq:homogeneous-section-left-action}
		(U_af)(g)=f(a^{-1}g).
	\end{equation}
	若 \(G\) 紧，以 \(\mu_G\) 为 \(G\) 上满足 \(\mu_G(G)=1\) 的左右不变 Haar 测度，并在 \(W\) 上选取 \(H\) 不变 Hermitian 内积，则
	\[
		\langle\sigma_f,\sigma_k\rangle_{L^2}
		=\int_G\langle f(g),k(g)\rangle_W\,\mathrm{d}\mu_G(g)
	\]
	给出截面空间上的正定内积，各 \(U_a\) 保持此内积，因而延拓为其 Hilbert 完备化上的酉算子。
\end{proposition}
\begin{proof}
	截面公式的纤维位置已核对；光滑性来自群作用及截面的光滑性。直接计算
	\[
		(U_aU_b\sigma)(x)
		=a\cdot\bigl(b\cdot\sigma(b^{-1}a^{-1}x)\bigr)
		=(ab)\cdot\sigma((ab)^{-1}x).
	\]
	单位元作用为恒等，所以确为表示。若 \(x=gH\)，则
	\[
		(U_a\sigma_f)(gH)
		=a\cdot[a^{-1}g,f(a^{-1}g)]
		=[g,f(a^{-1}g)],
	\]
	这就得到函数模型。左平移保持右等变条件，因为
	\(f(a^{-1}gh)=\rho(h)^{-1}f(a^{-1}g)\)。

	紧性使 \(H\) 也紧，所以可以平均一个 Hermitian 内积使其 \(H\) 不变。于是
	\[
		\langle f(gh),k(gh)\rangle_W=\langle f(g),k(g)\rangle_W,
	\]
	被积函数下降为底空间上的函数；积分也等于对 \(G/H\) 上由 \(\mu_G\) 推前的概率测度积分。若 \(f\ne0\)，则它在某个开集上范数有正下界，而归一化光滑体积在非空开集上为正，故所得内积正定。左不变性及代换 \(g=au\) 给出
	\[
		\int_G\langle f(a^{-1}g),k(a^{-1}g)\rangle_W\,\mathrm{d}\mu_G(g)
		=\int_G\langle f(u),k(u)\rangle_W\,\mathrm{d}\mu_G(u).
	\]
	因此每个 \(U_a\) 是等距线性算子，唯一延拓到完备化，延拓后的逆为 \(U_{a^{-1}}\)，故为酉算子。
\end{proof}

截面空间通常是无限维的。若想在其中寻找一个指定有限维表示 \(V\)，直接构造每一条截面会显得繁重；等变性允许只观察基点 \(o\) 上的值。以下互反关系将问题化为两个有限维稳定子表示之间的线性映射。

\begin{theorem}[光滑截面的 Frobenius 互反]\label{g1:thm:smooth-frobenius-reciprocity}
	设 \(G\) 为有限维实 Lie 群，\(H\) 为闭子群，\(\pi:G\to\operatorname{GL}(V)\)、\(\rho:H\to\operatorname{GL}(W)\) 为有限维复光滑表示。令 \(E_\rho=G\times_H W\)，采用 \([gh,w]=[g,\rho(h)w]\) 的商约定；对光滑截面 \(\sigma\)，以 \(\sigma(gH)=[g,f_\sigma(g)]\) 定义其等变函数。截面上取 \(G\) 作用 \((U_a\sigma)(x)=a\cdot\sigma(a^{-1}x)\)，其中 \(a\cdot[g,w]=[ag,w]\)。则求值于单位元给出自然线性同构
	\begin{equation}\label{g1:eq:smooth-frobenius-reciprocity}
		\operatorname{Hom}_G\bigl(V,\Gamma^\infty(G/H,E_\rho)\bigr)
		\cong\operatorname{Hom}_H(V|_H,W).
	\end{equation}
	这里两侧的 \(\operatorname{Hom}\) 都指复线性的交织映射。若左侧映射为 \(\Phi\)，则对应的 \(A\) 是
	\(A(v)=f_{\Phi(v)}(e)\)；反过来，\(A\) 对应的截面由
	\begin{equation}\label{g1:eq:frobenius-section-formula}
		f_v(g)=A\bigl(\pi(g^{-1})v\bigr),\qquad
		\Phi_A(v)(gH)=[g,f_v(g)]
	\end{equation}
	给出。
\end{theorem}
\begin{proof}
	设 \(\Phi\) 为 \(G\) 交织映射。对于 \(h\in H\)，先用 \(\Phi\) 的交织性，再用函数模型及右等变性，得到
	\[
		A(\pi(h)v)
		=f_{U_h\Phi(v)}(e)
		=f_{\Phi(v)}(h^{-1})
		=\rho(h)f_{\Phi(v)}(e)=\rho(h)A(v).
	\]
	所以 \(A\) 确为 \(H\) 交织映射。

	反过来，给定 \(A\in\operatorname{Hom}_H(V|_H,W)\)，\eqref{g1:eq:frobenius-section-formula}中的函数光滑，且
	\[
		f_v(gh)=A\bigl(\pi(h^{-1})\pi(g^{-1})v\bigr)
		=\rho(h)^{-1}f_v(g).
	\]
	所以它给出光滑截面，且线性依赖于 \(v\)。对 \(a\in G\)，
	\[
		f_{\pi(a)v}(g)=A\bigl(\pi(g^{-1}a)v\bigr)
		=f_v(a^{-1}g)=(U_af_v)(g),
	\]
	因此 \(\Phi_A\) 为 \(G\) 交织映射。其在单位元的值为 \(A(v)\)，证明一个方向的复合为恒等。若从 \(\Phi\) 出发再构造，则对任意 \(g\)，
	\[
		f_{\Phi(v)}(g)
		=f_{U_{g^{-1}}\Phi(v)}(e)
		=f_{\Phi(\pi(g^{-1})v)}(e)
		=A\bigl(\pi(g^{-1})v\bigr).
	\]
	故重建的截面就是原截面，另一方向的复合也为恒等。
\end{proof}

例如，取 \(V=\mathbb C\) 为平凡表示，得到
\[
	\Gamma^\infty(G/H,E_\rho)^G\cong W^H.
\]
群不变截面对应函数模型中的常函数，而这样的常函数必须取值于稳定子的不动子空间。若 \(V\) 不可约，任一非零交织映射 \(V\to\Gamma^\infty(E_\rho)\) 的核都是不变子空间，故它单射；于是 \(\operatorname{Hom}_H(V|_H,W)\ne0\) 正是这个不可约表示能够出现为截面子空间的判据。它不意味着全部光滑截面都是有限多个不可约分量的和；无限展开仍须另行说明收敛方式。

\subsection{复射影直线上的齐性线丛}
\label{g1:subsec:2003:04}
现在取 \(G=\operatorname{SU}(2)\) 及极大环面
\[
	T=\{t_\theta=\operatorname{diag}(e^{\mathrm{i}\theta},e^{-\mathrm{i}\theta}):\theta\in\mathbb R\}.
\]
对整数 \(m\)，令 \(\chi_m(t_\theta)=e^{\mathrm{i}m\theta}\)，并记
\[
	E_m=\operatorname{SU}(2)\times_{\chi_m}\mathbb C.
\]
这些线丛把上一章环面的整数权转变为球面上的具体几何对象。\symidx{\(E_m\)：球面上权为 \(m\) 的关联线丛}我们将直接核对符号：在本书的右主丛约定下，权一对应射影直线自身组成的线丛，而其对偶对应权负一。

\begin{proposition}[球面上的齐性线丛]\label{g1:prop:su2-hopf-associated-lines}
	设 \(T=\{t(\theta)=\operatorname{diag}(e^{\mathrm{i}\theta},e^{-\mathrm{i}\theta}):\theta\in\mathbb R\}\subseteq\operatorname{SU}(2)\)，\(m\in\mathbb Z\)，\(\chi_m(t(\theta))=e^{\mathrm{i}m\theta}\)，\(E_m=\operatorname{SU}(2)\times_{\chi_m}\mathbb C\)，采用 \([gt,z]=[g,\chi_m(t)z]\) 的商约定。由 \(gT\mapsto\mathbb C\,g(1,0)^{\mathsf T}\) 得到自然微分同胚 \(\operatorname{SU}(2)/T\cong\mathbb{CP}^1\cong S^2\)。在这个识别下，\(E_1\) 与线丛
	\[
		L=\{(\ell,v):\ell\in\mathbb{CP}^1,\ v\in\ell\}\longrightarrow\mathbb{CP}^1
	\]
	等变同构。记 \(L=\mathcal O(-1)\)，以对偶及张量幂定义 \(\mathcal O(k)\)、\(k\in\mathbb Z\)，则有光滑等变同构 \(E_m\cong\mathcal O(-m)\)。在两张射影图
	\[
		U_0:\ z=z_1/z_0,\qquad U_1:\ w=z_0/z_1
	\]
	上，可取局部标架 \(\epsilon_0^{(m)},\epsilon_1^{(m)}\) 使交叠处
	\begin{equation}\label{g1:eq:projective-line-associated-transition}
		w=1/z,\qquad
		\epsilon_1^{(m)}=w^m\epsilon_0^{(m)},\qquad
		a_0(z)=z^{-m}a_1(1/z)
	\end{equation}
	成立，其中最后一式是同一截面的局部分量转换。
\end{proposition}
\begin{proof}
	令 \(e_0=(1,0)^{\mathsf T}\)。映射 \(g\mapsto g\mathbb Ce_0\) 将 \(\operatorname{SU}(2)\) 送到 \(\mathbb{CP}^1\)。对任意单位向量 \(u=(u_0,u_1)^{\mathsf T}\)，矩阵
	\[
		g_u=\begin{pmatrix}u_0&-\overline{u_1}\\u_1&\overline{u_0}\end{pmatrix}
	\]
	属于 \(\operatorname{SU}(2)\)，并且 \(g_ue_0=u\)，所以作用传递。保持 \(\mathbb Ce_0\) 的特殊酉矩阵必须对角，恰为 \(T\)。由\cref{g1:prop:transitive-action-model}，轨道对应给出所述光滑商模型。

	取第一列还把 \(\operatorname{SU}(2)\) 微分同胚地识别为 \(S^3\)，逆映射就是 \(u\mapsto g_u\)；右乘 \(t_\theta\) 使第一列乘以 \(e^{\mathrm{i}\theta}\)。因此当前主丛正是\cref{g1:exm:hopf-principal-curvature}的 Hopf 丛。该例中的球极坐标同胚在这里写成
	\[
		[1:z]\longmapsto
		\frac{(2\operatorname{Re}z,2\operatorname{Im}z,1-|z|^2)}{1+|z|^2}.
	\]
	它在另一图 \([w:1]\)、\(w=1/z\) 上光滑延拓到南极，逆映射也由两张球极坐标给出，故 \(\mathbb{CP}^1\cong S^2\)。

	定义
	\[
		E_1\longrightarrow L,\qquad
		[g,c]\longmapsto(g\mathbb Ce_0,cge_0).
	\]
	若把 \([g,c]\) 换成 \([gt_\theta,c]\)，输出的向量变成 \(ce^{\mathrm{i}\theta}ge_0\)，与 \([g,e^{\mathrm{i}\theta}c]\) 的输出相同，所以公式良定义。它逐纤维是线性同构，在局部截面中及其逆都光滑，并与 \(\operatorname{SU}(2)\) 左作用相容。因此 \(E_1\cong L\)。对于任意整数 \(m,n\)，使用同一个 \(g\) 表示同一纤维中的元素，可定义
	\[
		E_m\otimes E_n\longrightarrow E_{m+n},\qquad
		[g,c]\otimes[g,d]\longmapsto[g,cd].
	\]
	改变代表 \(g\mapsto gt_\theta\) 时，两因子的分量分别乘以 \(e^{-\mathrm{i}m\theta}\)、\(e^{-\mathrm{i}n\theta}\)，乘积正好按权 \(m+n\) 变化，所以该映射良定义且逐纤维可逆。又有 \(E_0\cong\mathbb{CP}^1\times\mathbb C\)，于是 \(E_{-m}\cong E_m^*\)。从 \(m=1\) 出发便得全部 \(E_m\cong\mathcal O(-m)\)。

	最后计算局部转换。两张图上的主丛截面可取
	\[
		s_0(z)=\frac1{\sqrt{1+|z|^2}}
		\begin{pmatrix}1&-\overline z\\z&1\end{pmatrix},\qquad
		s_1(w)=\frac1{\sqrt{1+|w|^2}}
		\begin{pmatrix}w&-1\\1&\overline w\end{pmatrix}.
	\]
	它们的第一列分别代表 \([1:z]\)、\([w:1]\)。在交叠处，直接比较第一列得到
	\[
		s_1(1/z)=s_0(z)
		\operatorname{diag}\!\left(\frac{|z|}{z},\frac{z}{|z|}\right).
	\]
	第二列由第一列及特殊酉条件唯一确定，所以这也是完整矩阵的等式。故 \(E_m\) 的单位局部标架 \(b_i^{(m)}=[s_i,1]\) 满足
	\(b_1^{(m)}=(|z|/z)^m b_0^{(m)}\)。为了得到更简单的有理转换，取
	\[
		\epsilon_0^{(m)}=(1+|z|^2)^{m/2}b_0^{(m)},\qquad
		\epsilon_1^{(m)}=(1+|w|^2)^{m/2}b_1^{(m)}.
	\]
	由于 \(1+|1/z|^2=(1+|z|^2)/|z|^2\)，得到 \(\epsilon_1^{(m)}=z^{-m}\epsilon_0^{(m)}=w^m\epsilon_0^{(m)}\)。对于 \(m=1\)，它们在线丛 \(L\) 中正是 \(e_0(z)=(1,z)\)、\(e_1(w)=(w,1)\)。同一截面写成 \(\epsilon_0^{(m)}a_0=\epsilon_1^{(m)}a_1\)，便得最后的分量关系。
\end{proof}

局部标架的转换与截面分量的转换方向相反。比如 \(m=-n\)、\(n\ge0\) 时，\(E_{-n}=\mathcal O(n)\) 的分量必须满足 \(a_0(z)=z^n a_1(1/z)\)。这个条件将使下一节的二元齐次多项式恰好粘成全局截面；对于 \(E_n=\mathcal O(-n)\)，相同的多项式规则则不成立。

\begin{proposition}[这些线丛的非平凡性]\label{g1:prop:projective-line-bundle-winding}
	设 \(m,n\in\mathbb Z\)，\(T=\{t(\theta)=\operatorname{diag}(e^{\mathrm{i}\theta},e^{-\mathrm{i}\theta}):\theta\in\mathbb R\}\subseteq\operatorname{SU}(2)\)。对 \(j=m,n\)，令 \(\chi_j(t(\theta))=e^{\mathrm{i}j\theta}\)，\(E_j=\operatorname{SU}(2)\times_{\chi_j}\mathbb C\to\operatorname{SU}(2)/T\cong\mathbb{CP}^1\)，采用 \([gt,z]=[g,\chi_j(t)z]\) 的商约定。则在 \(\mathbb{CP}^1\) 上，\(E_m\) 作为光滑复线丛平凡，当且仅当 \(m=0\)。更一般地，\(E_m\) 与 \(E_n\) 作为覆盖底空间恒等映射的光滑复线丛同构，当且仅当 \(m=n\)。
\end{proposition}
\begin{proof}
	\(m=0\) 的平凡化已经构造。若 \(E_m\) 平凡，就有处处非零的光滑截面，其在 \(U_0,U_1\) 的分量 \(a_0,a_1\) 处处非零。在赤道 \(z=e^{\mathrm{i}\theta}\)、\(w=e^{-\mathrm{i}\theta}\) 上，转换关系给出
	\[
		a_0(e^{\mathrm{i}\theta})
		=e^{-\mathrm{i}m\theta}a_1(e^{-\mathrm{i}\theta}).
	\]
	任意非零复数环路先除以模长，就成为圆周环路；\cref{g1:thm:circle-fundamental}用连续角度的总增量定义其绕数。环路相乘时角度相加，所以绕数相加；反向参数化使绕数变号。若环路能在圆盘上保持非零地延拓，则径向收缩给出零同伦，绕数为零。

	这里 \(a_0\) 与 \(a_1\) 分别在 \(|z|\le1\)、\(|w|\le1\) 的圆盘上非零延拓，所以上式左边的绕数为零，右边第二因子的绕数也为零。第一因子的绕数为 \(-m\)，故 \(0=-m\)，只能有 \(m=0\)。最后，若 \(E_m\cong E_n\)，则 \(E_m^*\otimes E_n\cong E_{n-m}\) 有由该同构给出的处处非零截面，因而平凡；刚才的结论迫使 \(m=n\)。反向取恒等同构即可。
\end{proof}

这些线丛即使非平凡，也有大量光滑截面：在一张图内取紧支集光滑函数乘以局部标架，再在图外延拓为零即可。非平凡性排除的是处处非零的截面。另一方面，\(E_m\) 在基点的稳定子表示为 \(\chi_m\)，所以 \(m\ne0\) 时没有非零的 \(\operatorname{SU}(2)\) 不变截面。普通光滑截面的存在、不变截面的存在和全局标架的存在，是三个不同的问题。
% 纲要条目排为编号小节标题，并进入目录与书签；说明段仍保留于纲要注释中。

% !TeX root = ../../Book1.tex
% 纲要标题：### 20.4 SU(2) 的最高权分类与球面线丛实现
% 纲要位置：几何学计划纲要.md，第 912 行。
% 正文状态：初稿；尚未正式审核。

% 纲要内容开始
% 
% 以 SU(2) 的权分解和球面线丛说明后续最高权及全纯实现的问题；明确哪些结论已由本章证明，哪些需要第四卷复几何。
% 
% ---
% 
% 纲要内容结束

\section{SU(2) 的最高权分类与球面线丛实现}
\label{g1:sec:2004}

Peter--Weyl 定理告诉我们所有不可约表示的矩阵系数合在一起足够丰富，却没有逐个列出这些表示。上一节又把表示与齐性丛截面联系起来，但一个光滑截面空间通常无限维。现在要分别解决两个更具体的问题：怎样用权辨认不可约表示，以及怎样在一个截面空间中找出自然的有限维表示。我们在 \(\operatorname{SU}(2)\) 上把这两步都算出来，再说明一般紧群所需的新结构。

\ProseParagraphBreak

固定
\[
	T=\{t(\theta)=\operatorname{diag}(e^{\mathrm{i}\theta},e^{-\mathrm{i}\theta}):\theta\in\mathbb R\},
	\qquad k=\begin{pmatrix}\mathrm{i}&0\\0&-\mathrm{i}\end{pmatrix}.
\]
环面的周期为 \(2\pi\)，所以它的特征标恰为 \(t(\theta)\mapsto e^{\mathrm{i}m\theta}\)、\(m\in\mathbb Z\)。若 \(V\) 是有限维复酉表示，将 \(\mathrm{d}\pi(k)\) 对角化，就得到
\[
	V=\bigoplus_{m\in\mathbb Z}V[m],\qquad
	V[m]=\{v:\pi(t(\theta))v=e^{\mathrm{i}m\theta}v\text{ 对所有 }\theta\}.
\]
这里只出现有限多个非零空间。整数性也可直接核验：斜 Hermitian 算子的特征值写成 \(\mathrm{i}a\)，而 \(t(2\pi)=I\) 迫使 \(e^{2\pi\mathrm{i}a}=1\)。本节将整数 \(m\) 简称为相应的权；它代表 Lie 代数上的实泛函 \(\theta k\mapsto m\theta\)。第十九章的正根为 \(\alpha(\theta k)=2\theta\)，所以沿根方向改变权时，整数标签改变二。

\ProseParagraphBreak

先构造候选表示。对整数 \(n\ge0\)，记
\[
	V_n=\{p(z_0,z_1):p\text{ 是总次数 }n\text{ 的复齐次多项式}\},
	\qquad (\pi_n(g)p)(z)=p(g^{-1}z).
\]
\symidx{\(V_n,\pi_n\)：\(\operatorname{SU}(2)\) 的 \(n\) 次齐次多项式表示}常数构成 \(V_0\)。逆矩阵保证这是左表示：依次代入 \(g^{-1}\)、\(h^{-1}\) 得到 \((gh)^{-1}\)，故 \(\pi_n(g)\pi_n(h)=\pi_n(gh)\)。系数光滑依赖 \(g\)，所以它是光滑表示。其维数为 \(n+1\)，由平均内积可以酉化。

\begin{theorem}[特殊酉二阶群的不可约表示]\label{g1:thm:su2-polynomial-classification}
	取 \(T=\{t(\theta)=\operatorname{diag}(e^{\mathrm{i}\theta},e^{-\mathrm{i}\theta}):\theta\in\mathbb R\}\subseteq\operatorname{SU}(2)\)。在表示 \(\pi\) 中，约定非零向量 \(v\) 满足 \(\pi(t(\theta))v=e^{\mathrm{i}m\theta}v\)（所有 \(\theta\in\mathbb R\)）时称其为权 \(m\) 向量，其中 \(m\in\mathbb Z\)。对每个整数 \(n\ge0\)，二元 \(n\) 次复齐次多项式空间 \(V_n\) 配备作用 \((\pi_n(g)p)(z)=p(g^{-1}z)\)，是 \(\operatorname{SU}(2)\) 的不可约复表示，权为
	\[
		n,n-2,\ldots,-n,
	\]
	每个权空间均为一维。每个有限维不可约复光滑表示都与唯一一个 \(V_n\) 等价。\(n\) 是该表示的\term[zuigaoquan]{最高权}[highest weight]。
\end{theorem}
\begin{proof}
	先在候选表示中实际计算升降算子。\(\mathfrak{su}(2)\) 的复化为 \(\mathfrak{sl}(2,\mathbb C)\)，取
	\[
		e=\begin{pmatrix}0&1\\0&0\end{pmatrix},\quad
		f=\begin{pmatrix}0&0\\1&0\end{pmatrix},\quad
		h=\begin{pmatrix}1&0\\0&-1\end{pmatrix}.
	\]
	它们满足 \([h,e]=2e\)、\([h,f]=-2f\)、\([e,f]=h\)，且 \(k=\mathrm{i}h\)。对 \(p(g^{-1}z)\) 微分，得到
	\[
		E=-z_1\frac{\partial}{\partial z_0},\qquad
		F=-z_0\frac{\partial}{\partial z_1},\qquad
		L=-z_0\frac{\partial}{\partial z_0}+z_1\frac{\partial}{\partial z_1}.
	\]
	因此 \(z_0^jz_1^{n-j}\) 的权为 \(n-2j\)，且 \(F\) 将它送到 \(-(n-j)z_0^{j+1}z_1^{n-j-1}\)，\(E\) 则向相反方向移动，系数为 \(-j\)。两端以外的这些系数都非零。

	设 \(W\subseteq V_n\) 是非零群不变子空间。微分后，它在 \(E,F,L\) 下也不变。\(L\) 有互不相同的特征值，每个特征空间投影都可写成 \(L\) 的插值多项式，所以 \(W\) 包含某个单项式。反复使用 \(E,F\) 就得到全部单项式，故 \(W=V_n\)。这证明不可约性。

	反过来，设 \(\pi\) 是任一不可约表示，先用\cref{g1:prop:compact-unitarization-complete-reducibility}选不变 Hermitian 内积。把微分复线性延拓，记 \(E=\mathrm{d}\pi(e)\)、\(F=\mathrm{d}\pi(f)\)、\(L=\mathrm{d}\pi(h)\)。实 Lie 代数中的
	\[
		x=e-f,\qquad y=\mathrm{i}(e+f),\qquad k=\mathrm{i}h
	\]
	在酉表示中微分为斜 Hermitian 算子。由 \(e=(x-\mathrm{i}y)/2\)、\(f=(-x-\mathrm{i}y)/2\)、\(h=-\mathrm{i}k\)，可得 \(E^*=F\)、\(L^*=L\)。三个交换关系也由微分保持括号而成立。因此\cref{g1:lem:sl2-weight-chain}的条件全部满足，空间分成有限条升降链。

	每条链都在 \(E,F,L\) 下不变，因而在实 Lie 代数下不变，再由指数自然性在每个 \(\pi(\exp X)\) 下不变。\(\operatorname{SU}(2)\) 连通，\cref{g1:prop:exponentials-generate-identity-component}说明指数生成整个群，所以每条链也是群不变子空间。不可约性迫使只有一条链。于是存在非负整数 \(n\) 及非零向量 \(v\)，使
	\[
		Ev=0,\quad Lv=nv,\quad
		v,Fv,\ldots,F^nv
	\]
	为一组基，且 \(E(F^jv)=j(n-j+1)F^{j-1}v\)。这正是候选表示中从 \(z_1^n\) 开始的链所满足的公式。把 \(F^jv\) 送到 \(F^jz_1^n\)，得到与三个算子交换的线性同构。它因而与所有指数的表示交换，再与整个群交换，给出 \(\pi\cong\pi_n\)。不同 \(n\) 的维数不同，唯一性随之成立。
\end{proof}

最高权来自升降链的终止。选定正根以后，升算子不断把权加二；有限维性保证这个过程终止，链顶被升算子杀掉。交换关系又强迫链顶是非负整数，并准确决定向下还能走多少步。一般紧群有多个根方向，这才出现如何让不同升降方向相容的新问题。

\begin{corollary}[特征标与旋转群表示]\label{g1:cor:su2-character-so3-descent}
	设 \(n\ge0\) 为整数，\(V_n\) 为二元 \(n\) 次复齐次多项式空间，取 \(\operatorname{SU}(2)\) 作用 \((\pi_n(g)p)(z)=p(g^{-1}z)\)，并令 \(t(\theta)=\operatorname{diag}(e^{\mathrm{i}\theta},e^{-\mathrm{i}\theta})\)。则 \(V_n\) 的特征标在该对角环面上为
	\[
		\chi_n(t(\theta))=\sum_{j=0}^n e^{\mathrm{i}(n-2j)\theta}
		=\frac{\sin((n+1)\theta)}{\sin\theta}\qquad(\sin\theta\ne0).
	\]
	在 \(\theta=0\) 时其值为 \(n+1\)，在 \(\theta=\pi\) 时为 \((-1)^n(n+1)\)。该表示下降到 \(\operatorname{SO}(3)\cong\operatorname{SU}(2)/\{\pm I\}\) 当且仅当 \(n\) 为偶数。因此 \(\operatorname{SO}(3)\) 的不可约复光滑表示恰由非负整数 \(\ell\) 标记，维数为 \(2\ell+1\)。
\end{corollary}
\begin{proof}
	迹等于各一维权空间上的特征值之和，有限等比级数求和给出正弦商；分母为零时直接使用有限和，不把商式当作定义。中心元 \(-I\) 作用为 \(p(z)\mapsto p(-z)=(-1)^np(z)\)。由\cref{g1:prop:quaternion-low-rank}的二重覆盖，下降恰要求这个核作用平凡，即 \(n\) 为偶数。下降后的表示光滑，因为覆盖有光滑局部逆。反过来，将旋转群的不可约表示沿该满射拉回，仍不可约，且核作用平凡；分类定理于是给出唯一的 \(n=2\ell\)。
\end{proof}

\input{Figures/Chapter20-SU2Weights}

\cref{g1:fig:su2-weight-chains}显示了权的奇偶性如何记录整体群。相同 Lie 代数上的升降公式不能单独保证表示下降；必须检查周期格，也即这里的中心元 \(-I\)。这正是第十九章根数据与整体群区别在表示论中的一个可计算后果。

现在回到球面上的线丛。沿用\cref{g1:prop:su2-hopf-associated-lines}的记号，\(E_m=\operatorname{SU}(2)\times_{\chi_m}\mathbb C\to\mathbb{CP}^1\)，其中 \(\chi_m(t(\theta))=e^{\mathrm{i}m\theta}\)。\(E_1\) 是重言线丛，纤维为点 \(\ell\in\mathbb{CP}^1\) 所代表的直线本身；\(E_{-n}\) 是它的 \(n\) 次对偶张量幂。为了说明多项式怎样变成截面，先把每条直线上的齐次性解释为该纤维中的一个元素。

\begin{proposition}[球面线丛中的多项式表示]\label{g1:prop:polynomial-sections-sphere-line}
	设 \(n\ge0\) 为整数，\(V_n\) 为二元 \(n\) 次复齐次多项式空间，取 \(\operatorname{SU}(2)\) 作用 \((\pi_n(g)p)(z)=p(g^{-1}z)\)。令 \(T=\{t(\theta)=\operatorname{diag}(e^{\mathrm{i}\theta},e^{-\mathrm{i}\theta}):\theta\in\mathbb R\}\subseteq\operatorname{SU}(2)\)；对 \(j\in\mathbb Z\)，置 \(\chi_j(t(\theta))=e^{\mathrm{i}j\theta}\)、\(E_j=\operatorname{SU}(2)\times_{\chi_j}\mathbb C\)，采用 \([gt,c]=[g,\chi_j(t)c]\) 的商约定。每个 \(p\in V_n\) 确定 \(E_{-n}\) 的光滑截面，等变函数为
	\[
		f_p(g)=p(ge_0),\qquad e_0=(1,0)^{\mathsf T}.
	\]
	这一映射 \(V_n\to\Gamma^\infty(E_{-n})\) 单射且 \(\operatorname{SU}(2)\) 等变。在 \(U_0=\{z_0\ne0\}\)、\(U_1=\{z_1\ne0\}\) 的坐标 \(z=z_1/z_0\)、\(w=z_0/z_1\) 中，相对于 \((1,z)\)、\((w,1)\) 的对偶张量标架，截面分量分别为
	\[
		a_0(z)=p(1,z),\qquad a_1(w)=p(w,1),\qquad
		a_0(z)=z^n a_1(1/z).
	\]
	它的像恰为这两张图中分量都为多项式的截面。若将 \(n\) 换成负整数，则满足相同多项式过渡条件的截面只有零截面。
\end{proposition}
\begin{proof}
	齐次性使 \(p\) 在一条直线 \(\ell\) 上的限制成为 \((\ell^*)^{\otimes n}\) 的一个元素：若选 \(v\ne0\) 作基，只需规定它在 \(v^{\otimes n}\) 上取值 \(p(v)\)；把基改成 \(cv\) 时，输入和取值都乘 \(c^n\)，故与选择无关。把这个构造放在各条直线上，就得到所述截面。在两张图中写出的分量都是多项式，因此光滑。

	用单位标架 \(ge_0\) 表示同一个截面，所得函数正是 \(f_p\)。右乘 \(t(\theta)\) 使第一列乘以 \(e^{\mathrm{i}\theta}\)，故
	\[
		f_p(gt(\theta))=e^{\mathrm{i}n\theta}f_p(g)
		=\chi_{-n}(t(\theta))^{-1}f_p(g),
	\]
	符合\cref{g1:prop:equivariant-functions-sections}的条件。若截面为零，则 \(p\) 在所有单位向量上为零；齐次性使它在全部非零向量上为零，也在零点为零，故 \(p=0\)。对 \(a\in\operatorname{SU}(2)\)，
	\[
		f_{\pi_n(a)p}(g)=p(a^{-1}ge_0)=f_p(a^{-1}g),
	\]
	正是截面上的左作用，所以映射等变。

	在交集上，\((w,1)=w(1,z)\)、\(w=1/z\)，故对偶的 \(n\) 次张量幂给出 \(a_0(z)=z^n a_1(1/z)\)。若 \(a_1(w)=\sum_{j=0}^N b_jw^j\)，右侧为 \(\sum_j b_jz^{n-j}\)。它与多项式 \(a_0\) 在 \(\mathbb C^*\) 上相等，说明所有负次幂的系数都为零。对 \(n\ge0\)，只有 \(0\le j\le n\) 可以出现，而
	\[
		p(z_0,z_1)=\sum_{j=0}^n b_j z_0^jz_1^{n-j}
	\]
	恰好产生这两个分量。对 \(n<0\)，所有可能的次数 \(n-j\) 都为负，只能全部系数为零。这里的比较只是有限 Laurent 多项式系数的唯一性，无须对一般解析函数作展开。
\end{proof}

\(E_{-n}\) 的光滑截面远不止这些多项式截面。实际上，每一种奇偶性合适、足够大的不可约表示都会出现。以下结果精确说明两种空间的差别。

\begin{proposition}[球面线丛的光滑表示重数]\label{g1:prop:sphere-line-section-multiplicity}
	设 \(m\in\mathbb Z\)，\(V_n\) 为二元 \(n\) 次复齐次多项式表示，\(n\ge0\) 为整数，作用为 \((\pi_n(g)p)(z)=p(g^{-1}z)\)。令 \(T=\{t(\theta)=\operatorname{diag}(e^{\mathrm{i}\theta},e^{-\mathrm{i}\theta}):\theta\in\mathbb R\}\subseteq\operatorname{SU}(2)\)，\(\chi_m(t(\theta))=e^{\mathrm{i}m\theta}\)，\(E_m=\operatorname{SU}(2)\times_{\chi_m}\mathbb C\)，采用 \([gt,c]=[g,\chi_m(t)c]\) 的商约定。截面上取左作用 \((U_a\sigma)(x)=a\cdot\sigma(a^{-1}x)\)，其中 \(a\cdot[g,c]=[ag,c]\)。则
	\[
		\dim\operatorname{Hom}_{\operatorname{SU}(2)}
		\bigl(V_n,\Gamma^\infty(E_m)\bigr)
		=\begin{cases}
			1,& n\ge|m|\text{ 且 }n\equiv m\pmod2,\\
			0,&\text{其他情形}.
		\end{cases}
	\]
	把光滑截面按不变 Hermitian 度量和归一化商体积作 \(L^2\) 完备化，所得 Hilbert 空间 \(\mathcal H_m\) 在左群作用下分解为
	\[
		\mathcal H_m\cong
		\widehat{\bigoplus}_{\substack{n\ge|m|\\n\equiv m\ (2)}}V_n,
	\]
	每个表示恰出现一次。这是 Hilbert 空间中的收敛陈述，不表示任意光滑截面的任意排序级数都一致或逐阶可微收敛。
\end{proposition}
\begin{proof}
	\cref{g1:thm:smooth-frobenius-reciprocity}把所求同态空间识别为 \(\operatorname{Hom}_T(V_n,\mathbb C_m)\)，其中 \(T\) 在 \(\mathbb C_m\) 上按 \(\chi_m\) 作用。一个这样的线性映射在权 \(q\ne m\) 的空间上必为零：取使 \(e^{\mathrm{i}q\theta}\ne e^{\mathrm{i}m\theta}\) 的 \(\theta\)，等变性就迫使该空间的像为零。在权 \(m\) 空间上则可任取线性泛函。\(V_n\) 的权是 \(n,n-2,\ldots,-n\)，各重数一，故得到第一式。

	再说明为何这份重数表穷尽整个 \(L^2\) 空间。令 \(\mathrm{d}t\) 为 \(T\) 的归一化积分，在 \(L^2(\operatorname{SU}(2))\) 上定义
	\[
		(P_m f)(g)=\int_T\chi_m(t)f(gt)\,\mathrm{d}t.
	\]
	积分可先对连续函数定义，再由 \(\|P_mf\|_2\le\|f\|_2\) 延拓；这个估计来自 Cauchy--Schwarz 不等式与右不变换元。对连续函数换元 \(u=st\) 得 \((P_mf)(gs)=\chi_m(s)^{-1}(P_mf)(g)\)，而已经满足该条件的函数被 \(P_m\) 固定。因此 \(P_m^2=P_m\)。酉性、取逆不变性及 \(\overline{\chi_m(t)}=\chi_m(t^{-1})\) 给出 \(P_m^*=P_m\)，故它是正交投影，并与所有左平移交换。

	\cref{g1:thm:peter-weyl-hilbert}给出矩阵系数有限和在群上 \(L^2\) 稠密。\(P_m\) 把每个有限维系数空间仍送回该空间，所得函数光滑且满足上述等变性，所以其像正是 \(\mathcal H_m\) 中的一个稠密子空间。等变函数的模沿右 \(T\) 陪集不变，群积分下推到商上的积分，因而该识别保持范数。不同不可约类型正交，由第一部分给出的重数，每个允许的 \(V_n\) 恰留下一份，其他类型消失，得到 Hilbert 直和。这里也可注意 \(V_n\) 的对偶仍不可约，权集合取负后不变，因此按分类仍等价于 \(V_n\)；使用左正则表示或其对偶的标记不会改变此处的重数表。
\end{proof}

\begin{example}\label{g1:exm:sphere-sections-smooth-polynomial}
	\(m=0\) 时，线丛平凡，光滑截面就是球面上的光滑函数。上式给出维数依次为 \(1,3,5,\ldots\) 的不可约部分。常数构成 \(V_0\)；在单位球面 \(S^2\subseteq\mathbb R^3\) 上，三个坐标函数的复线性张成空间构成三维的 \(V_2\)。后一断言可由第十九章 \(\operatorname{SU}(2)\) 在虚四元数上的共轭作用核验：其复化是 \(\mathfrak{sl}(2,\mathbb C)\) 的三维升降链，因而对应 \(n=2\)。

	对 \(m=1\)，重言线丛虽然没有处处非零的光滑截面，却有一个二维不变截面空间。给定 \(a\in\mathbb C^2\)，令 \(\sigma_a(\ell)\) 为 \(a\) 在直线 \(\ell\) 上的正交投影。它光滑依赖 \(\ell\)，且酉变换保持正交投影，故 \(a\mapsto\sigma_a\) 等变。取 \(\ell=\mathbb Ca\) 可知它单射；非零 \(a\) 对应的截面恰在 \(\ell=a^\perp\) 处为零。这与丛的非平凡性相容。与此同时，前一命题在 \(n=-1\) 的结论说明 \(E_1\) 没有非零的两图多项式截面。光滑条件与多项式条件已经选出了不同的表示空间。
\end{example}

由此可以看清一般理论下一步要增加什么。权是极大环面的特征标，根描述不同权空间之间可由哪些 Lie 代数方向连接；最高权还依赖正根的选择。\(\operatorname{SU}(2)\) 只有一个正根方向，有限升降链已经给出完整分类。对于一般紧连通 Lie 群，需要研究多个根方向的相容性，并同时保留特征标格对整体群的限制；一般最高权分类在这一层面组织不可约表示。

\ProseParagraphBreak

几何方面，多项式截面给出了 \(V_n\) 的显式实现，但选取更一般的截面类别还需要底空间及线丛的复结构。Borel--Weil 理论\footnote{Armand Borel（1923--2003），瑞士数学家，研究代数拓扑与线性代数群。André Weil（1906--1998），法国数学家，对代数几何与数论有重要贡献；其姓氏 Weil 与本章 Weyl 是不同姓名。}研究在旗流形的全纯线丛截面中实现表示。第四卷将建立全纯坐标、全纯丛与相应的解析工具，这些概念使“哪些截面应当保留”成为内在的几何条件。在已经算出的秩一模型中，整数权决定线丛的过渡函数，多项式条件则选出了相应的有限维表示。
% 本节沿用纲要的连续叙述安排，不另造小节标题。


\begin{chaptersummary}
	归一化不变积分把有限维表示酉化，使每个不变子空间都有不变正交补。Schur 引理再把平均交织映射化为零或标量，从而给出矩阵系数与特征标的正交关系。固定一个维数为 \(d_\pi\) 的不可约表示，其全部矩阵系数张成 \(d_\pi^2\) 维空间，在右正则作用下包含 \(d_\pi\) 个该表示的副本。

	\ProseParagraphBreak

	平滑对称逼近单位产生与右正则作用交换的紧自伴卷积算子。其非零特征空间有限维、由光滑函数组成，并在群作用下不变。紧自伴谱分解先给出 \(L^2\) 逼近，再卷积的估计将它提升为一致逼近，由此证明代表函数在 \(C(G)\) 与 \(L^2(G)\) 中稠密。正交性与稠密性合起来得到 Peter--Weyl 分解及 Parseval 等式；一般正交展开的直接结论是 \(L^2\) 收敛，光滑逼近则使用另行选择的平滑核与有限谱截断。

	\ProseParagraphBreak

	稳定子表示 \(\rho:H\to\operatorname{GL}(W)\) 决定齐性向量丛 \(G\times_H W\)。截面对应满足 \(f(gh)=\rho(h)^{-1}f(g)\) 的函数，群在截面上的作用对应左平移。求值于单位元给出 Frobenius 互反，把截面空间中一个有限维表示的出现问题化为稳定子上的交织映射问题。对球面模型，权 \(m\) 对应 \(E_m=\mathcal O(-m)\)，局部标架与截面分量按互逆方向转换。

	\ProseParagraphBreak

	\(\operatorname{SU}(2)\) 的不可约复表示由非负整数 \(n\) 分类：\(V_n\) 是二元 \(n\) 次齐次多项式空间，权为 \(n,n-2,\ldots,-n\)，维数为 \(n+1\)。中心元素 \(-I\) 作用为 \((-1)^n\)，故恰在 \(n\) 为偶数时下降到 \(\operatorname{SO}(3)\)。\(V_n\) 实现为 \(E_{-n}\) 的多项式截面；而 \(E_m\) 的全部平方可积截面还包含所有满足 \(n\ge|m|\)、\(n\equiv m\pmod2\) 的 \(V_n\)，各出现一次。有限维几何实现与无限维截面分解由此在同一个例子中得到比较。
\end{chaptersummary}

\BookExercises

以下习题依次练习平均、正交关系、卷积逼近、齐性丛与低秩表示。处理无限展开时注明收敛范数；处理线丛时先写出所用的代表关系和局部分量转换，再判断符号与重数。

\begin{exercise}\label{g1:ex:representation-average-fixed-vectors}
	设 \(G\) 为紧 Lie 群，\(\pi:G\to\operatorname{U}(V)\) 为有限维复光滑酉表示，令
	\[
		P=\int_G\pi(g)\,\mathrm{d}\mu_G(g),\qquad
		V^G=\{v:\pi(g)v=v\;\text{对全部 }g\in G\}.
	\]
	证明 \(P\) 是到 \(V^G\) 的正交投影，并推出 \(\dim_{\mathbb C}V^G=\int_G\chi_\pi\,\mathrm{d}\mu_G\)。对圆周表示 \(\pi(z)=\operatorname{diag}(z^{m_1},\ldots,z^{m_N})\)，其中 \(m_j\in\mathbb Z\)，求 \(P\)。
\end{exercise}
\begin{solution}
	对 \(h\in G\)，左不变性给出
	\[
		\pi(h)P=\int_G\pi(hg)\,\mathrm{d}\mu_G(g)=P.
	\]
	因此 \(\operatorname{im}P\subseteq V^G\)。若 \(v\in V^G\)，被积向量恒为 \(v\)，归一化给出 \(Pv=v\)。所以像恰为 \(V^G\)，且 \(P^2=P\)。又因表示酉且积分取逆不变，
	\[
		P^*=\int_G\pi(g)^*\,\mathrm{d}\mu_G(g)
		=\int_G\pi(g^{-1})\,\mathrm{d}\mu_G(g)=P.
	\]
	自伴投影的核与像正交：若 \(Pv=0\)、\(u=Pw\)，则 \(\langle v,u\rangle=\langle Pv,w\rangle=0\)。结合 \(V=\operatorname{im}P\oplus\ker P\)，可知 \(P\) 就是到 \(V^G\) 的正交投影。

	在适应这一直和的基中，\(P\) 在像上为 \(1\)、核上为 \(0\)，故 \(\operatorname{tr}P=\dim V^G\)。取迹与有限维积分交换，得到所需特征标积分公式。圆周情形逐对角分量积分，而 \(\int_{S^1}z^m\,\mathrm{d}\mu=0\) 对 \(m\ne0\) 成立，\(m=0\) 时为 \(1\)，因此
	\[
		P=\operatorname{diag}(\varepsilon_1,\ldots,\varepsilon_N),
		\qquad \varepsilon_j=\begin{cases}1,&m_j=0,\\0,&m_j\ne0.\end{cases}
	\]
	平均恰好保留零权坐标。
\end{solution}

\begin{exercise}\label{g1:ex:intertwiner-character-dimension}
	设 \(\rho,\sigma\) 为紧 Lie 群 \(G\) 的有限维复光滑表示。证明
	\[
		\dim_{\mathbb C}\operatorname{Hom}_G(V_\rho,V_\sigma)
		=\int_G\chi_\sigma(g)\overline{\chi_\rho(g)}\,\mathrm{d}\mu_G(g).
	\]
	进而证明非零表示 \(\rho\) 不可约，当且仅当 \(\int_G|\chi_\rho|^2\,\mathrm{d}\mu_G=1\)。对 \(G=S^1\)、\(\rho(z)=\operatorname{diag}(z,z,z^{-1})\)、\(\sigma(z)=\operatorname{diag}(z,z^{-1},z^{-1},z^2)\)，计算交织映射空间及自交织映射空间 \(\operatorname{End}_G(V_\rho)\) 的维数。
\end{exercise}
\begin{solution}
	由完全可约性，把两个表示按所选不可约代表写成
	\[
		V_\rho\cong\bigoplus_{\pi\in\widehat G}V_\pi^{\oplus a_\pi},
		\qquad
		V_\sigma\cong\bigoplus_{\pi\in\widehat G}V_\pi^{\oplus b_\pi},
	\]
	其中每个直和只有有限个非零项。任意线性映射可按这些因子写成矩阵分块；交织条件逐块成立。Schur 引理使不等价因子之间的块为零，相同代表之间的每块是一个任意复标量，所以
	\[
		\operatorname{Hom}_G(V_\rho,V_\sigma)
		\cong\bigoplus_\pi\operatorname{Mat}_{b_\pi\times a_\pi}(\mathbb C),
		\qquad \dim\operatorname{Hom}_G(V_\rho,V_\sigma)=\sum_\pi a_\pi b_\pi.
	\]
	另一方面，\(\chi_\rho=\sum a_\pi\chi_\pi\)、\(\chi_\sigma=\sum b_\pi\chi_\pi\)，由特征标正交性，题中积分也等于 \(\sum a_\pi b_\pi\)。

	取 \(\sigma=\rho\)，得到 \(\int_G|\chi_\rho|^2\,\mathrm{d}\mu_G=\sum a_\pi^2\)。\(\rho\ne0\) 意味着至少一个 \(a_\pi\) 为正整数，因此右侧等于 \(1\) 当且仅当恰有一个重数为 \(1\)、其余为零，也就是 \(\rho\) 不可约。

	在所给圆周例中，\(\rho\) 的权 \(1,-1\) 的重数分别为 \(2,1\)，\(\sigma\) 的权 \(1,-1,2\) 的重数分别为 \(1,2,1\)。所以 \(\dim\operatorname{Hom}_G(V_\rho,V_\sigma)=2\cdot1+1\cdot2=4\)，而 \(\dim\operatorname{End}_G(V_\rho)=2^2+1^2=5\)。
\end{solution}

\begin{exercise}\label{g1:ex:finite-cyclic-fourier}
	令 \(C_m=\{1,\zeta,\ldots,\zeta^{m-1}\}\)，其中 \(\zeta=e^{2\pi\mathrm{i}/m}\)、\(m\ge1\)。求 \(C_m\) 的全部复不可约表示，并对任意函数 \(f:C_m\to\mathbb C\) 证明展开公式及平方范数恒等式
	\[
		f(\zeta^r)=\sum_{k=0}^{m-1}\widehat f(k)\zeta^{kr},
		\qquad
		\frac1m\sum_{r=0}^{m-1}|f(\zeta^r)|^2
		=\sum_{k=0}^{m-1}|\widehat f(k)|^2,
	\]
	其中 \(\widehat f(k)=m^{-1}\sum_{r=0}^{m-1}f(\zeta^r)\zeta^{-kr}\)。说明此处无需使用无限维函数空间的完备性。
\end{exercise}
\begin{solution}
	\(C_m\) 交换。和环面的情形一样，每个表示矩阵都与全体表示矩阵交换，复 Schur 引理使不可约表示中的矩阵全为标量，不可约性再使表示空间一维。一个一维表示由 \(\zeta\) 的像 \(\lambda\) 决定，而 \(\zeta^m=1\) 要求 \(\lambda^m=1\)。因此全部不可约表示为 \(\chi_k(\zeta^r)=\zeta^{kr}\)，\(0\le k<m\)。有限群是零维 Lie 群，这些映射自动光滑。

	对 \(k,l\) 取值于上述范围，有限等比数列给出
	\[
		\frac1m\sum_{r=0}^{m-1}\chi_k(\zeta^r)\overline{\chi_l(\zeta^r)}
		=\frac1m\sum_{r=0}^{m-1}(\zeta^{k-l})^r=\delta_{kl}.
	\]
	\(k=l\) 时每项为 \(1\)；\(k\ne l\) 时分子 \(1-(\zeta^{k-l})^m\) 为零、分母 \(1-\zeta^{k-l}\) 非零，所以和为零。全部函数组成的复向量空间维数为 \(m\)，这 \(m\) 个正交非零函数因而已是一组标准正交基。有限维正交展开使 \(f=\sum_k\langle f,\chi_k\rangle\chi_k\)，系数正是所列 \(\widehat f(k)\)，取范数即得第二个恒等式。这里的展开为有限和，完备性来自维数计数，没有涉及无限级数的收敛问题。
\end{solution}

\begin{exercise}\label{g1:ex:su2-regular-row-subspace}
	将 \(\operatorname{SU}(2)\) 的元素写为
	\[
		g=\begin{pmatrix}a(g)&b(g)\\-\overline{b(g)}&\overline{a(g)}\end{pmatrix},
		\qquad |a(g)|^2+|b(g)|^2=1.
	\]
	证明其在 \(\mathbb C^2\) 上的定义表示不可约，并用 Schur 正交关系求四个坐标函数 \(a,b,\overline a,\overline b\) 的范数及两两内积。证明 \(E=\operatorname{span}_{\mathbb C}\{a,b\}\) 在右正则作用下不变，却不在左正则作用下不变；给出一个具体的左平移作为例子。
\end{exercise}
\begin{solution}
	若定义表示有非零真不变子空间，它必为复直线。取 \(\operatorname{diag}(e^{\mathrm{i}t},e^{-\mathrm{i}t})\)，其中 \(t\notin\pi\mathbb Z\)，两个特征值不同，该矩阵的不变复直线只能是两条坐标轴。可是 \(J=\begin{psmallmatrix}0&1\\-1&0\end{psmallmatrix}\in\operatorname{SU}(2)\) 交换两条坐标轴，故没有一条对全群不变。表示因此不可约。

	四个矩阵系数依次是 \(a,b,-\overline b,\overline a\)，维数 \(d=2\)。Schur 正交关系说明它们两两正交，且每个平方范数为 \(1/2\)。去掉一个函数前的负号并改变排列不影响此结论，所以 \(a,b,\overline a,\overline b\) 两两正交，范数都为 \(1/\sqrt2\)，特别地这四个函数线性无关。

	写 \(h=\begin{psmallmatrix}c&d\\-\overline d&\overline c\end{psmallmatrix}\)。右乘直接给出
	\[
		R_ha=ca-\overline d\,b,\qquad
		R_hb=da+\overline c\,b,
	\]
	所以 \(E\) 右不变，且在基 \((a,b)\) 中，限制表示的矩阵就是 \(h\)。另一方面，\(J^{-1}=\begin{psmallmatrix}0&-1\\1&0\end{psmallmatrix}\)，计算 \(J^{-1}g\) 的左上角可得 \(L_Ja=\overline b\)。四个函数的线性无关性说明 \(\overline b\notin E\)，所以 \(E\) 不左不变。左右正则作用互相交换，并不意味着一个任意选定的右不变子空间也必须左不变；在完整的四维系数空间中，两种作用才同时保持该空间。
\end{solution}

\begin{exercise}\label{g1:ex:nonnegative-kernel-negative-spectrum}
	在 \(S^1\) 上取归一化角度积分，令 \(k(e^{\mathrm{i}t})=1-a\cos t\)，其中 \(0<a\le1\)。证明 \(k\) 非负、积分为 \(1\)、取逆不变，但卷积算子 \(T_k\) 有负特征值。求其全部特征值和非零特征空间，并说明“非负核”与“非负自伴算子”的区别。
\end{exercise}
\begin{solution}
	由 \(-1\le\cos t\le1\) 有 \(k\ge1-a\ge0\)，余弦的平均为零，故积分为 \(1\)，且 \(\cos(-t)=\cos t\) 给出取逆不变性。对 \(e_m(e^{\mathrm{i}t})=e^{\mathrm{i}mt}\)，直接代入左卷积，得
	\[
		T_ke_m(e^{\mathrm{i}t})
		=e^{\mathrm{i}mt}\frac1{2\pi}\int_0^{2\pi}
		(1-a\cos s)e^{-\mathrm{i}ms}\,\mathrm{d}s.
	\]
	所以 \(e_0\) 的特征值为 \(1\)，\(e_1,e_{-1}\) 的特征值为 \(-a/2\)，其余 \(e_m\) 的特征值为 \(0\)。\cref{g1:ex:peter-weyl-circle}说明这些指数函数为完备正交系，因此没有其他特征值；非零特征空间分别是 \(\mathbb C e_0\) 与 \(\operatorname{span}\{e_1,e_{-1}\}\)，算子秩为 \(3\)。

	非负核保证 \(f\ge0\) 时 \(T_kf\ge0\)，这是函数值的次序性质。非负自伴算子则要求 \(\langle T_kf,f\rangle_2\ge0\) 对所有复函数成立。取 \(f=e_1\) 即得 \(\langle T_kf,f\rangle_2=-a/2<0\)，所以后一性质不成立。
\end{solution}

\begin{exercise}\label{g1:ex:approximate-identity-not-operator-norm}
	设 \(T_n\) 是本节在 \(S^1\) 上构造的平滑对称逼近单位算子。证明对每个 \(n\)，都有 \(\|I-T_n\|_{L^2\to L^2}\ge1\)。因此 \(T_nf\to f\) 对每个 \(f\in L^2\) 成立，并不意味着 \(T_n\to I\) 按算子范数成立。
\end{exercise}
\begin{solution}
	固定 \(n\)，记 \(T=T_n\)，并取正交单位函数列 \(e_m(t)=e^{\mathrm{i}mt}\)、\(m=1,2,\ldots\)。有限正交投影的勾股等式说明，对每个 \(h\in L^2\)，有 \(\sum_m|\langle e_m,h\rangle_2|^2\le\|h\|_2^2\)，故 \(\langle e_m,h\rangle_2\to0\)。自伴性遂给出 \(\langle Te_m,h\rangle_2=\langle e_m,Th\rangle_2\to0\)。

	紧性使 \((Te_m)\) 的每个子列还有范数收敛子列，而任何这样的极限 \(v\) 均满足 \(\langle v,h\rangle_2=0\) 对所有 \(h\) 成立，所以只能是 \(0\)。若 \(\|Te_m\|_2\) 不趋于零，便存在范数有正下界的子列；取它的收敛子列会与极限为零矛盾。因此 \(\|Te_m\|_2\to0\)。由
	\[
		\|(I-T)e_m\|_2\ge1-\|Te_m\|_2
	\]
	得到 \(\|I-T\|\ge1\)。这对每个固定 \(n\) 成立，而逼近单位的强收敛已经由正文证明。两者的量词不同：强收敛先固定目标函数；算子范数收敛则要求对整个单位球同时控制误差。
\end{solution}

\begin{exercise}\label{g1:ex:convolution-fourier-order}
	设 \(G\) 为有限维实紧 Lie 群，\(k,f\in C(G)\)，取本节约定的左卷积及 Fourier 矩阵。证明
	\[
		\widehat{k*f}(\pi)=\widehat f(\pi)\widehat k(\pi).
	\]
	若 \(k\) 共轭不变，即 \(k(hxh^{-1})=k(x)\)，证明对每个不可约酉表示 \(\pi\)，\(\widehat k(\pi)=c_\pi I\)，求 \(c_\pi\)，并说明 \(T_k\) 在 \(\mathcal H_\pi\) 上的作用。
\end{exercise}
\begin{solution}
	全部被积函数连续，紧性允许交换积分。将 \(x=yz\) 代入 Fourier 矩阵的定义，有
	\[
		\begin{aligned}
		\widehat{k*f}(\pi)
		&=\int_G\int_Gk(y)f(z)\pi((yz)^{-1})\,
		\mathrm{d}\mu(z)\,\mathrm{d}\mu(y)\\
		&=\left(\int_Gf(z)\pi(z^{-1})\,\mathrm{d}\mu(z)\right)
		\left(\int_Gk(y)\pi(y^{-1})\,\mathrm{d}\mu(y)\right).
		\end{aligned}
	\]
	这正是所求次序；逆元把乘法顺序反转，所以不能按标量 Fourier 变换的习惯随意交换两个矩阵。

	若 \(k\) 共轭不变，用共轭换元得到
	\[
		\pi(h)\widehat k(\pi)\pi(h)^{-1}
		=\int_Gk(y)\pi(hy^{-1}h^{-1})\,\mathrm{d}\mu(y)
		=\widehat k(\pi).
	\]
	Schur 引理使它为标量矩阵。取迹得
	\[
		c_\pi=\frac1{d_\pi}\int_Gk(y)\overline{\chi_\pi(y)}\,\mathrm{d}\mu(y).
	\]
	直接把 \(\pi(y^{-1}x)=\pi(y^{-1})\pi(x)\) 代入卷积，得到
	\[
		T_k\pi_{ij}(x)
		=\sum_a\widehat k(\pi)_{ia}\pi_{aj}(x)=c_\pi\pi_{ij}(x).
	\]
	因此 \(T_k\) 在整个 \(\mathcal H_\pi\) 上都是乘以 \(c_\pi\)。
\end{solution}

\begin{exercise}\label{g1:ex:character-convolution-projection}
	设 \(G\) 为有限维实紧 Lie 群，\(\pi\) 为不可约酉表示。证明从 \(L^2(G)\) 到 \(\mathcal H_\pi\) 的正交投影为
	\[
		Q_\pi f(x)=d_\pi\int_G\chi_\pi(y)f(y^{-1}x)\,\mathrm{d}\mu(y).
	\]
	据此求 \(\chi_\pi*\chi_\sigma\)，其中 \(\pi,\sigma\) 是已选定的不可约代表，并说明投影的秩为什么是 \(d_\pi^2\) 而不是 \(d_\pi\)。
\end{exercise}
\begin{solution}
	特征标函数共轭不变。由上一题的公式以及\cref{g1:cor:character-orthogonality}，核 \(k=d_\pi\chi_\pi\) 在 \(\mathcal H_\sigma\) 上的卷积特征值为
	\[
		\frac{d_\pi}{d_\sigma}
		\int_G\chi_\pi(y)\overline{\chi_\sigma(y)}\,\mathrm{d}\mu(y)
		=\begin{cases}1,&\pi=\sigma,\\0,&\pi\ne\sigma.\end{cases}
	\]
	因此该卷积在 \(\mathcal H_\pi\) 上为恒等，在其他系数空间上为零。它是有界算子，Peter--Weyl 的正交 Hilbert 直和将这个结论延拓到全部 \(L^2\)，恰好得到正交投影。把 \(f=\chi_\sigma\) 代入，得到
	\[
		\chi_\pi*\chi_\sigma
		=\begin{cases}\chi_\pi/d_\pi,&\pi=\sigma,\\0,&\pi\ne\sigma.\end{cases}
	\]
	像空间由 \(d_\pi^2\) 个独立矩阵条目张成，故秩为 \(d_\pi^2\)。\(d_\pi\) 是原表示空间的维数；正则表示中相应系数空间含有 \(d_\pi\) 个该表示的副本，所以两种维数不同。
\end{solution}

\begin{exercise}\label{g1:ex:characters-class-hilbert-basis}
	设 \(G\) 为有限维实紧 Lie 群，在 \(L^2(G)\) 上令 \(C_hf(x)=f(h^{-1}xh)\)。把共轭不变子空间定义为
	\[
		L^2(G)^{\mathrm{cl}}=\{f:C_hf=f\text{ 对每个 }h\in G\}.
	\]
	证明全部不可约特征标 \(\chi_\pi\) 构成这个闭子空间的标准正交基，并写出相应的展开和 Parseval 等式。对一般 \(L^2\) 元素，应如何理解题中的不变性？
\end{exercise}
\begin{solution}
	\(C_h=L_hR_h\) 为酉算子，所以每个 \(\ker(C_h-I)\) 闭，它们的交 \(L^2(G)^{\mathrm{cl}}\) 也闭。左右正则作用保持每个 \(\mathcal H_\pi\)，因而 \(C_h\) 保持它。它还保持各 \(\mathcal H_\pi\) 的正交补，所以相应正交投影与 \(C_h\) 交换。

	每个 \(\mathcal H_\pi\) 中的函数可唯一写为 \(f_A(x)=\operatorname{tr}(A\pi(x))\)。唯一性来自矩阵系数线性无关。计算得
	\[
		C_hf_A(x)=\operatorname{tr}\bigl(\pi(h)A\pi(h)^{-1}\pi(x)\bigr).
	\]
	因此 \(f_A\) 被所有 \(C_h\) 固定，当且仅当 \(A\) 与所有 \(\pi(h)\) 交换。Schur 引理使 \(A=cI\)，所以 \(\mathcal H_\pi\) 内的固定子空间恰为 \(\mathbb C\chi_\pi\)。投影与共轭作用交换，Peter--Weyl 分解遂给出
	\[
		L^2(G)^{\mathrm{cl}}
		=\mathop{\widehat\bigoplus}_{[\pi]\in\widehat G}\mathbb C\chi_\pi.
	\]
	特征标正交关系给出 \(\|\chi_\pi\|_2=1\)，于是对 \(f\in L^2(G)^{\mathrm{cl}}\)，
	\[
		f=\sum_{[\pi]}\langle f,\chi_\pi\rangle_2\chi_\pi,
		\qquad\|f\|_2^2=\sum_{[\pi]}|\langle f,\chi_\pi\rangle_2|^2.
	\]
	展开按 \(L^2\) 范数成立。不变性是在 Hilbert 空间中对每个固定 \(h\) 的等式；若用平方可积函数模零的模型，它表示相应等式几乎处处成立，不能擅自为未选定的代表断言每一点都共轭不变。
\end{solution}

\begin{exercise}\label{g1:ex:peter-weyl-convergence-counterexamples}
	在 \(S^1\) 上令
	\[
		f_N(e^{\mathrm{i}t})=\frac1N\sum_{m=0}^{N-1}e^{\mathrm{i}mt},
		\qquad h_N(e^{\mathrm{i}t})=\frac{e^{\mathrm{i}Nt}}N.
	\]
	证明 \(f_N\to0\) 按 \(L^2\) 范数成立但不一致收敛；证明 \(h_N\to0\) 一致成立但不按 \(C^1\) 收敛。说明即使所有近似函数都是代表函数，仍必须注明所用收敛意义。
\end{exercise}
\begin{solution}
	圆周指数函数标准正交，故 \(\|f_N\|_2^2=N/N^2=1/N\to0\)。另一方面，三角不等式给出 \(|f_N|\le1\)，而 \(f_N(1)=1\)，所以 \(\|f_N\|_\infty=1\)，不可能一致趋于零。

	对第二列，\(\|h_N\|_\infty=1/N\to0\)，但
	\[
		\frac{\mathrm{d}}{\mathrm{d}t}h_N(e^{\mathrm{i}t})
		=\mathrm{i}e^{\mathrm{i}Nt},
	\]
	其一致范数恒为 \(1\)，故导数不能一致趋于零，不是 \(C^1\) 收敛。两列都是有限 Fourier 和，因而都是代表函数；属于同一个函数代数并不会使三种拓扑的收敛条件相同。
\end{solution}

\begin{exercise}\label{g1:ex:extendible-isotropy-trivial-bundle}
	设 \(G\) 为实 Lie 群，\(H\subseteq G\) 为闭子群，且 \(H\) 的有限维复表示 \(\rho\) 可以延拓为 \(G\) 的表示 \(\pi:G\to\operatorname{GL}(W)\)。构造 \(E_\rho\) 到 \((G/H)\times W\) 的光滑丛同构，并写出这个平凡化中 \(G\) 的作用。将结论应用于 \(\operatorname{SU}(2)\) 的标准二维表示，证明 \(E_1\oplus E_{-1}\) 平凡，并解释这个齐性丛仍没有非零的不变截面。
\end{exercise}
\begin{solution}
	定义
	\[
		[g,w]\longmapsto(gH,\pi(g)w).
	\]
	由于 \(\pi(gh)w=\pi(g)\rho(h)w\)，它与关联丛的等价关系相容。在一个余类代表 \(g\) 下，逆映射为 \((gH,v)\mapsto[g,\pi(g)^{-1}v]\)；换用 \(gh\) 会使括号内分量乘以 \(\rho(h)^{-1}\)，所以逆也与代表无关。两式在主丛局部截面中光滑，且逐纤维线性互逆。左乘 \(a\) 后，第二坐标从 \(\pi(g)w\) 变成 \(\pi(ag)w=\pi(a)\pi(g)w\)，故平凡化中的群作用为
	\[
		a\cdot(x,v)=(ax,\pi(a)v).
	\]
	它同时移动底点与纤维，不能把它理解为只在底点上作用。

	标准二维表示限制到对角环面时分解为 \(\chi_1\oplus\chi_{-1}\)，关联构造与直和相容，所以所得丛正是 \(E_1\oplus E_{-1}\)。上述公式给出其全局平凡化。对于不变截面，Frobenius 互反的平凡表示特例将其识别为 \((\mathbb C^2)^T\)。若 \((v_0,v_1)\) 被全部 \(\operatorname{diag}(e^{\mathrm{i}\theta},e^{-\mathrm{i}\theta})\) 固定，则两个分量都必须为零。因此不变截面空间为零。普通平凡化并不要求所选常值标架为群不变标架。
\end{solution}

\begin{exercise}\label{g1:ex:circle-finite-isotropy-sections}
	取 \(G=S^1\)、\(H=\mu_q=\{\zeta:\zeta^q=1\}\)，其中 \(q\ge2\)。对整数 \(r\) 令 \(\rho_r(\zeta)=\zeta^r\)，并以 \(x=z^q\) 识别 \(G/H\cong S^1\)。证明关联线丛 \(E_{\rho_r}\) 光滑平凡，写出一个处处非零截面。再求特征标 \(\pi_n(a)=a^n\) 在其光滑截面空间中出现的条件及交织映射空间的维数。
\end{exercise}
\begin{solution}
	函数 \(f_r(z)=z^{-r}\) 满足 \(f_r(z\zeta)=\zeta^{-r}f_r(z)\)，所以
	\[
		s_r(z^q)=[z,z^{-r}]
	\]
	给出光滑截面，并且处处非零。映射 \((x,c)\mapsto c\,s_r(x)\) 因而是全局平凡化。

	由 Frobenius 互反，所求交织空间同构于 \(\operatorname{Hom}_{\mu_q}(\mathbb C_n,\mathbb C_r)\)。一维空间之间的非零线性映射交织，当且仅当 \(\zeta^n=\zeta^r\) 对全部 \(q\) 次单位根成立，也就是 \(n\equiv r\pmod q\)。条件成立时交织空间维数为一，否则为零。对应的函数可以直接取 \(f_n(z)=z^{-n}\)，因为
	\[
		(U_af_n)(z)=f_n(a^{-1}z)=a^n f_n(z).
	\]
	在标架 \(s_r\) 中，截面写成 \(s_r(x)b(x)\)，函数模型为 \(z^{-r}b(z^q)\)，左作用变成
	\[
		(U_ab)(x)=a^r b(a^{-q}x).
	\]
	特别地，\(b(x)=x^k\) 的特征标为 \(a^{r-qk}\)，与模 \(q\) 的条件一致。不同的 \(r\bmod q\) 给出不同的齐性作用，尽管底层光滑线丛都平凡。
\end{solution}

\begin{exercise}\label{g1:ex:projective-tangent-weight}
	把 \(\mathbb{CP}^1\) 的实切丛用复坐标中的乘 \(\mathrm{i}\) 看成复线丛。求基点 \([1:0]\) 的 \(\operatorname{SU}(2)\) 稳定子在这条复切线上的特征标，并将整个切丛识别为某个 \(E_m\)。用坐标 \(w=1/z\) 的导数独立核对所得局部转换。
\end{exercise}
\begin{solution}
	在图 \([1:z]\) 中，\(t_\theta=\operatorname{diag}(e^{\mathrm{i}\theta},e^{-\mathrm{i}\theta})\) 将点送到
	\[
		[e^{\mathrm{i}\theta}:e^{-\mathrm{i}\theta}z]
		=[1:e^{-2\mathrm{i}\theta}z].
	\]
	因此在 \(z=0\) 处的微分为复数乘法 \(v\mapsto e^{-2\mathrm{i}\theta}v\)，稳定子特征标是 \(\chi_{-2}\)。坐标转换 \(w=1/z\) 的导数 \(-z^{-2}\) 是复线性映射，所以两图中的乘 \(\mathrm{i}\) 相容，确实给出复线丛结构。群作用在射影图中为复分式线性映射，其微分也复线性；由齐性丛的稳定子分类，
	\[
		\mathrm{T}\mathbb{CP}^1\cong E_{-2}\cong\mathcal O(2)
	\]
	作为光滑复线丛等变同构。

	独立看转换：若向量在 \(z\) 图中的复速度为 \(v_z\)，则在 \(w\) 图中为 \(v_w=-z^{-2}v_z\)，即 \(v_z=-z^2v_w\)。因此坐标切标架满足 \(e_w=-z^2e_z\)。在第二图将标架换为 \(-e_w\)，得到转换 \((-e_w)=z^2e_z=w^{-2}e_z\)，正是 \(E_{-2}\) 的标架转换。常数负号只改变一张图上的标架，不改变丛。
\end{solution}

\begin{exercise}\label{g1:ex:tautological-projection-sections}
	沿用\cref{g1:exm:sphere-sections-smooth-polynomial}的投影截面 \(\sigma_v(\ell)=\operatorname{proj}_\ell v\)。用归一化群积分证明
	\[
		\langle\sigma_v,\sigma_w\rangle_{L^2}
		=\tfrac12\langle v,w\rangle_{\mathbb C^2}.
	\]
	若 \(e_0,e_1\) 为 \(\mathbb C^2\) 的标准正交基，证明 \(\|\sigma_{e_0}(\ell)\|^2+\|\sigma_{e_1}(\ell)\|^2=1\) 对每条直线成立，并写出从全部平方可积截面到这两个截面张成空间的正交投影。
\end{exercise}
\begin{solution}
	用单位向量 \(u=ge_0\) 表示 \(\ell\)，正交投影矩阵为 \(P_\ell=uu^*\)，故 \(\sigma_v(\ell)=uu^*v\)。\(uu^*\) 不随相位改变，且在射影图中光滑，所以这是光滑截面。函数模型为 \(f_v(g)=u^*v\)，它满足 \(f_v(gt_\theta)=e^{-\mathrm{i}\theta}f_v(g)\)，符合 \(E_1\) 的等变条件。

	令 \(M=\int_G uu^*\,\mathrm{d}\mu_G(g)\)。左不变性给出 \(aMa^{-1}=M\) 对全部 \(a\in\operatorname{SU}(2)\) 成立。与所有对角相位矩阵交换迫使 \(M\) 对角，再与矩阵 \(\begin{psmallmatrix}0&-1\\1&0\end{psmallmatrix}\) 交换迫使两个对角元相同。又因 \(\operatorname{tr}(uu^*)=1\)，有 \(\operatorname{tr}M=1\)，故 \(M=I/2\)。本章内积对第一变量线性，因而
	\[
		\int_G(u^*v)\overline{(u^*w)}\,\mathrm{d}\mu_G(g)
		=w^*Mv=\tfrac12w^*v
		=\tfrac12\langle v,w\rangle.
	\]
	因为 \(\|\sigma_{e_j}(\ell)\|^2=|u^*e_j|^2=|u_j|^2\)，两项之和为 \(1\)。积分公式又说明 \(\sqrt2\sigma_{e_0},\sqrt2\sigma_{e_1}\) 是标准正交基，所以所求正交投影为
	\[
		Q\sigma=2\sum_{j=0}^1\langle\sigma,\sigma_{e_j}\rangle_{L^2}\sigma_{e_j}.
	\]
	对一般平方可积截面，这两个内积按完备化连续延拓。此处点上的平方和等于 \(1\)，每条截面的积分平方范数却为 \(1/2\)，说明逐点归一化与整体正交归一化是不同要求。
\end{solution}

\begin{exercise}\label{g1:ex:su2-tensor-decomposition}
	设 \(a,b\) 为非负整数。利用特征标证明
	\[
		V_a\otimes V_b\cong\bigoplus_{j=0}^{\min(a,b)}V_{a+b-2j}.
	\]
	据此分解 \(V_2\otimes V_3\)，并求 \((V_a\otimes V_b)^{\operatorname{SU}(2)}\) 的维数。
\end{exercise}
\begin{solution}
	张量积的特征标为两个迹的乘积。设 \(a\ge b\)，并把环面特征标暂写成 Laurent 多项式 \(\chi_n(u)=u^n+u^{n-2}+\cdots+u^{-n}\)。乘积中 \(u^{a+b-2k}\) 的系数是满足 \(i+j=k\)、\(0\le i\le a\)、\(0\le j\le b\) 的整数对数目，等于
	\[
		\min(b,k,a+b-k)+1\qquad(0\le k\le a+b).
	\]
	另一方面，它在 \(\sum_{j=0}^b\chi_{a+b-2j}(u)\) 中的系数是满足 \(0\le j\le b\)、\(j\le k\le a+b-j\) 的整数 \(j\) 的数目，恰是同一个数。因此两个 Laurent 多项式相同。代入 \(u=e^{\mathrm{i}\theta}\) 得到环面上的特征标恒等式；每个群元素共轭入环面，迹又在共轭下不变，故恒等式在全群成立。

	有限维表示完全可约，不可约特征标正交，所以用各 \(\chi_n\) 作内积即可唯一恢复每个不可约重数。于是特征标恒等式给出所述直和同构。具体有
	\[
		V_2\otimes V_3\cong V_5\oplus V_3\oplus V_1,
	\]
	两边维数均为 \(12\)。不变向量来自平凡表示 \(V_0\)；它出现当且仅当 \(a+b-2j=0\) 对某个 \(j\le\min(a,b)\) 成立，也即 \(a=b\)。此时恰出现一次，所以所求维数在 \(a=b\) 时为一，否则为零。
\end{solution}

\begin{exercise}\label{g1:ex:su2-polynomial-kernel}
	对 \(n\ge0\)，确定 \(\pi_n:\operatorname{SU}(2)\to\operatorname{GL}(V_n)\) 的核，并判断哪些 \(V_n\) 为忠实表示。对偶数 \(n>0\)，判断下降到 \(\operatorname{SO}(3)\) 的表示是否忠实。
\end{exercise}
\begin{solution}
	\(n=0\) 是平凡表示，核为全群。以下设 \(n>0\)。取核中的 \(g\)，由极大环面定理把它共轭为 \(t(\theta)\)。核为正规子群，所以 \(t(\theta)\) 仍在核中。这等价于全部权满足
	\[
		e^{\mathrm{i}(n-2j)\theta}=1\qquad(0\le j\le n).
	\]
	若 \(n\) 为奇数，权集合包含 \(1\)，故 \(\theta\in2\pi\mathbb Z\)，于是 \(g\) 共轭于单位矩阵，只能为单位矩阵。若 \(n\) 为正偶数，权集合包含 \(2\)，故 \(\theta\in\pi\mathbb Z\)，对应 \(\pm I\)。二者都是中心元，共轭不改变它们；所有权均为偶数又保证这两个元素确在核内。所以
	\[
		\ker\pi_n=\begin{cases}
			\operatorname{SU}(2),&n=0,\\
			\{I\},&n\text{ 为奇数},\\
			\{I,-I\},&n>0\text{ 为偶数}.
		\end{cases}
	\]
	忠实者恰为奇数 \(n\) 的表示。正偶数情形下降时恰好除去全部核，所以所得旋转群表示忠实。这里没有只检查 Lie 代数的核；环面上的周期条件同时确定了可能残留的离散核。
\end{solution}

\begin{exercise}\label{g1:ex:polynomial-product-exact-sequence}
	设 \(a,b\ge1\)。把 \(V_a\otimes V_b\) 识别为两个向量变量 \(z,w\in\mathbb C^2\) 的双齐次多项式，次数分别为 \(a,b\)，群同时作用于两变量。令
	\[
		D(z,w)=z_0w_1-z_1w_0,\qquad
		M(P)(z)=P(z,z).
	\]
	证明有 \(\operatorname{SU}(2)\) 等变短正合列
	\[
		0\longrightarrow V_{a-1}\otimes V_{b-1}
		\xrightarrow{\;P\mapsto DP\;}V_a\otimes V_b
		\xrightarrow{\;M\;}V_{a+b}\longrightarrow0.
	\]
	解释为什么它存在等变线性分裂，以及这个结论并不要求 \(M\) 保持多项式的全部乘法结构。
\end{exercise}
\begin{solution}
	矩阵的行列式为一，所以 \(D(g^{-1}z,g^{-1}w)=D(z,w)\)。乘以 \(D\) 因而等变，且使两个次数各增加一。多项式环无零因子，\(D\ne0\)，故第一张非零映射单射。对角代入也与群作用交换，所以 \(M\) 等变。它满射：次数 \(a+b\) 的每个单项式都能把其 \(a+b\) 个一次因子分成大小为 \(a\)、\(b\) 的两组，写成一个 \(a\) 次单项式与一个 \(b\) 次单项式之积，因此在 \(M\) 的像中。

	由于 \(D(z,z)=0\)，第一张映射的像包含于 \(\ker M\)。两者维数相同：前者维数为 \(ab\)，后者维数为
	\[
		(a+1)(b+1)-(a+b+1)=ab.
	\]
	所以像等于核，证明正合性。平均内积使中间表示酉化，\(\ker M\) 的正交补为群不变子空间，且 \(M\) 在该补空间上的限制是到 \(V_{a+b}\) 的等变线性同构；取逆即得分裂。它所保持的是表示的线性与群作用结构。固定次数的空间本身并非对多项式乘法封闭，因此这里并没有声称一个分次代数的分裂。
\end{solution}


% !TeX root = ../Book1.tex
% 纲要标题：# 第一卷附录
% 纲要位置：几何学计划纲要.md，第 918 行。
% 正文状态：附录 A—D 初稿完成，均含完整习题解答；尚未正式审核。

% 纲要内容开始
% 
% 纲要内容结束

\appendix
\renewcommand{\thechapter}{\Alph{chapter}}
\ctexset{chapter={number=\Alph{chapter}}}
\BookAppendixPart{g1:part:appendices}

% !TeX root = ../../Book1.tex
% 纲要标题：## 附录 A　公理化几何与几何史的方法
% 纲要位置：几何学计划纲要.md，第 920 行。
% 正文状态：初稿；尚未正式审核。

% 纲要内容开始
% 
% > 选读；保留公理、模型、坐标化之间的逻辑辨析，不充当正文证明的替代。
% 
% * **公理与模型**
%   * 关联、顺序、合同及连续性
%   * 平行公设和非欧模型
%   * 独立性与相对一致性
% 
% * **坐标化及历史论证**
%   * 几何对象与代数模型的双向建立
%   * 历史证明的原始语境与现代重述
% 
% ---
% 
% 纲要内容结束

\chapter{公理化几何与几何史的方法}
\label{g1:app:A}
\BookChapterNumber{g1:app:A}

\begin{chapterabstract}
	几何证明既要说明从哪些假设出发，也要说明“点”“直线”“相等”等词在所用模型中指什么。本附录以有限组关联公理、欧氏与双曲平面、仿射坐标和等边三角形构造为例，辨别公理、模型与计算之间的关系。我们在明确列出的公理系统中比较平行命题的不同真值，解释模型如何支持独立性与相对一致性的判断；再把已有的仿射及欧氏结构写成坐标，并核对坐标变换保留什么。历史部分回到具体原始论证，区分原文的构造条件、现代补充的假设和另一种证明。
\end{chapterabstract}

\begin{learninggoals}
	\begin{itemize}
		\item 区分点线关联、介于关系、合同关系与连续性分别规定的内容，并在具体模型中核验所列条件。
		\item 在欧氏与双曲平面中判断给定点线之间的平行关系，说明模型中使用的直线与图纸上的形状有何区别。
		\item 用同一公理系统的两个模型解释独立性，区分模型存在、形式推导与依赖背景理论的相对一致性。
		\item 由实仿射标架建立坐标，恢复直线、介于及内积数据，核对换标架时对象和矩阵怎样变化。
		\item 阅读一个历史几何构造时，分别检查目标、已有命题、对象存在性和现代重述中增加的条件。
	\end{itemize}
\end{learninggoals}

在一条线段的两个端点各画一个圆，半径都取这条线段的长度；若两圆相交，交点与两个端点便组成等边三角形。最后一句的长度比较很直接，真正需要停下来说明的是前一句中的“相交”。图纸能够帮助找到候选位置，却不能替代所选几何体系中的存在性依据。若改用坐标求解，方程给出候选点后，仍须检查这些坐标属于允许使用的数系。一个熟悉的作图问题就这样同时涉及构造、代数和公理。

\ProseParagraphBreak

本卷正文主要在已经指定结构的空间中工作：实仿射空间给出点差和直线，内积给出距离，流形与度量给出局部坐标和测地线。公理化研究则允许先把点、线及其关系视为尚未解释的对象，问哪些假设足以推出一项结论。模型再为这些词指定具体含义，使每条假设成为能够核验的数学陈述。两种做法解决的问题不同，但可以相互校验：模型检查假设是否相容，也提醒我们一个证明究竟用了哪些额外结构。

\ProseParagraphBreak

这里使用\chapref{g1:ch:01}的实仿射空间、内积与标架，以及\secref{g1:sec:1001}已经构造并证明性质的双曲平面。逻辑部分只介绍当前模型论证所需的满足关系与推理保真性，不预设数理逻辑课程。坐标部分从已有实仿射结构出发，不把选取标架等同于从任意抽象点线系统重建一个数域。历史部分采用可定位的具体文本，依其所陈述的条件比较证明，而不从现代计算反推作者当时的想法。

\section{公理与模型}
\label{g1:sec:A01}

“两点确定一条直线”可以在纸上画出的直线中理解，也可以在双曲平面的测地线中理解。换一种解释，并不改变这句话的量词；改变的是点、线以及“位于其上”的含义。要用这种方法判断平行公设是否来自其他假设，必须先把那些假设写清楚。否则，一个图形即使没有唯一平行线，也可能同时破坏了原本希望保留的条件。

\subsection{关联、顺序、合同及连续性}
\label{g1:subsec:A01:01}

先只考虑点和线的相互归属，不讨论长度、角度或介于关系。这允许我们把模型核验压缩到一组明确的条件，并看见后来加入的结构分别做了什么。

\begin{definition}[点线关联结构]
	\label{g1:def:incidence-model}
	设 \(X\) 是点集，\(\mathcal L\) 是线集，\(I\subset X\times\mathcal L\) 是一个关系。称 \((X,\mathcal L,I)\) 为一个\term[guanlianjiegou]{关联结构}[incidence structure]；\(I(A,\ell)\) 表示点 \(A\) 位于线 \(\ell\) 上。三个点若同在某条线上，就称它们共线。本节以 \(\mathsf I\) 表示以下三个公理组成的系统：
	\symidx{\(\mathsf I\)：附录 A 的三个关联公理}
	\begin{enumerate}
		\item 任意两个不同点恰好同在一条线上。
		\item 每条线上至少有两个不同点。
		\item 存在三个不共线的点。
	\end{enumerate}
	为点、线及关联关系指定具体集合与关系，并使这三个条件全部成立，就得到 \(\mathsf I\) 的一个\term[moxing]{模型}[model]。
\end{definition}

第一条中的“任意”“存在”和“唯一”不能在核验中互相替代。一个图中画出的两点可能有连线，仍不足以说明所有两点都有连线；若点集或线集受到限制，熟悉的直线也可能根本不属于所选模型。第二、三条排除只有空线、单点线或所有点都挤在同一条线上的退化情况。

\begin{proposition}[两线的公共点]
	\label{g1:prop:incidence-intersection}
	在满足 \(\mathsf I\) 的关联结构中，两条不同的线至多有一个公共点；对每条线，至少有一个点不在该线上。
\end{proposition}
\begin{proof}
	若两条线有不同的公共点 \(A,B\)，则它们都经过 \(A,B\)，由第一条的唯一性只能是同一条线。再固定一条线 \(\ell\)，取第三条保证的三个不共线点。它们不能全在 \(\ell\) 上，故其中至少一个点不在 \(\ell\) 上。
\end{proof}

这些公理尚未规定一条线上的点怎样排列。\term[jieyuguanxi]{介于关系}[betweenness relation]用三点关系 \(B(A,B,C)\) 表示“\(B\) 严格介于 \(A,C\) 之间”；严格意味着三点互不相同。\term[xianduanhetong]{线段合同}[segment congruence]则用四点关系 \(AB\equiv CD\) 表示两段可以作为同一种长度来比较。在公理化处理时，可以把这些关系也作为原始关系，另列它们应满足的条件。关联只说明哪些点共线，顺序说明谁在谁之间，合同说明哪些段相等，三者并不是同一信息。

\ProseParagraphBreak

本卷正文已经使用距离，因此可以在模型中直接解释这些新关系。设点集 \(X\) 另带距离 \(d\)，约定
\[
	\begin{aligned}
		B(A,B,C)\iff{}& A,B,C\text{ 互异且共线，}\\
		&\text{且 }d(A,C)=d(A,B)+d(B,C),\\
		AB\equiv CD\iff{}& d(A,B)=d(C,D).
	\end{aligned}
\]
共线条件仍须保留；在一般距离空间里，三角不等式取等号并不能单独指定几何直线。下面的性质说明，把每条线参数化为通常的实数轴，能够保证哪些顺序、合同与连续性条件。

\begin{proposition}[实参数给出的线内结构]
	\label{g1:prop:axiom-line-parameter}
	设关联结构 \((X,\mathcal L,I)\) 满足 \(\mathsf I\)，点集带距离 \(d\)。假设对每条线 \(\ell\)，都有一个从 \(\mathbb R\) 到 \(\ell\) 上点集的双射 \(\gamma_\ell\)，满足
	\[
		d\bigl(\gamma_\ell(s),\gamma_\ell(t)\bigr)=|s-t|.
	\]
	按上述方式定义介于与线段合同，则有以下性质：
	\begin{enumerate}
		\item \(B(A,B,C)\) 与 \(B(C,B,A)\) 等价；三个互异共线点恰有一个介于另外两个之间。
		\item 对不同点 \(A,B\)，存在 \(C\) 使 \(B(A,B,C)\) 成立。
		\item 固定起点 \(A\) 和一条从 \(A\) 出发的射线，对任意 \(r>0\)，射线上恰有一个点 \(D\) 满足 \(d(A,D)=r\)。特别地，任意非退化线段都可以在该射线上唯一复制。
		\item 线段合同是点对上的等价关系；若 \(B(A,B,C)\)、\(B(A',B',C')\)，且 \(AB\equiv A'B'\)、\(BC\equiv B'C'\)，则 \(AC\equiv A'C'\)。
		\item 每条线选定一个方向后，其中每个非空、有上界的点集都有最小上界。
	\end{enumerate}
\end{proposition}
\begin{proof}
	固定一条线并用 \(\gamma_\ell\) 的实参数表示其点。对三个互异参数 \(a,b,c\)，等式
	\[
		|a-c|=|a-b|+|b-c|
	\]
	成立，当且仅当 \(a<b<c\) 或 \(c<b<a\)。这给出第一项；取 \(c=2b-a\) 则给出第二项。

	射线可以表示为 \(\gamma_\ell([a,\infty))\) 或 \(\gamma_\ell(({-\infty},a])\)。相应的唯一参数为 \(a+r\) 或 \(a-r\)，从而得到第三项。合同的自反、对称与传递性来自实数相等；第四项最后的断言来自两组距离的相加。

	最后按所选方向令参数递增；若原参数方向相反，改用 \(t\mapsto\gamma_\ell(-t)\)。给定非空、有上界的点集 \(S\)，其参数集 \(\gamma_\ell^{-1}(S)\) 是非空、有上界的实数集，故有最小上界 \(s_0\)。点 \(\gamma_\ell(s_0)\) 是 \(S\) 的一个上界，而任何其他上界的参数都不小于 \(s_0\)，所以它就是所需的最小上界。
\end{proof}

这里的射线、方向和顺序没有依赖某个特选参数。事实上，若 \(\gamma,\widetilde\gamma\) 是同一条线的两种上述参数，令 \(h=\gamma^{-1}\circ\widetilde\gamma\)，则 \(|h(s)-h(t)|=|s-t|\)。记 \(b=h(0)\)，由到 \(0\) 的距离知 \(h(t)-b\) 的绝对值为 \(|t|\)，再与 \(h(1)-b=1\) 或 \(-1\) 比较到 \(1\) 的距离，分别得到 \(h(t)=b+t\) 或 \(h(t)=b-t\)。故换参数至多改变原点和方向，不改变介于关系，也不改变一个方向上最小上界是否存在。

\ProseParagraphBreak

第五项是一种\term[xuwanbeixing]{序完备性}[order completeness]要求。它与前四项的区别在于，需要控制无限点集留下的空隙；证明实际调用了实数的最小上界性质。若只允许变量遍历点和线，“任意点集”就不能直接写成这两类变量的一次全称量化。要把这条要求纳入形式公理系统，必须另行说明采用集合变量，还是采用某种较弱的一阶公理方案。下文讨论一阶独立性时，所用公理系统仍是明确列出的 \(\mathsf I\)。

\ProseParagraphBreak

也不能把上面的五项性质当成整个平面几何的公理表。它们没有规定一条线穿过三角形一边之后如何与其他边相交，也没有给出角合同以及边角边条件。Hilbert 的《几何基础》把这些不同任务分组处理；本附录所据 Townsend 英译重印本的关联、顺序、平行、合同、连续性依次见\cite[Chapter I, \S\S 2--8]{Hilbert1950Foundations}。其中 \S 8 以等段反复截取的 Archimedes 条件开头，随后另论完备性，不能把不同连续性要求视为同一个句子。这里选择较小的系统，是为了让每一个模型判断都有可以逐项检查的假设。

\subsection{平行公设和非欧模型}
\label{g1:subsec:A01:02}

在关联结构中，可以只用“有没有公共点”定义不相交，而无须先定义角度。对点线组成的整个结构，考虑如下\term[pingxinggongse]{平行公设}[parallel postulate]：
\[
	\begin{gathered}
		\mathsf P:\quad\text{对每条线 \(\ell\) 和每个不在 \(\ell\) 上的点 \(A\)，}\\
		\text{恰有一条过 \(A\) 的线与 \(\ell\) 不相交。}
	\end{gathered}
\]
这里采用的是唯一平行线的表述，与上述英译本\cite[Chapter I, \S 5]{Hilbert1950Foundations}一致。“不相交”指在模型的点集中没有公共点。命题 \(\mathsf P\) 同时包含存在性和唯一性；否定它，可以是某组点线根本没有这样的线，也可以是至少有两条。

\begin{proposition}[欧氏平面的模型核验]
	\label{g1:prop:euclidean-incidence-model}
	取点集 \(X=\mathbb R^2\)，线集为全部仿射直线 \(a+\mathbb R u\)，其中 \(u\ne0\)，关联解释为点属于直线的关系。这个结构满足 \(\mathsf I\) 和 \(\mathsf P\)。配以欧氏距离后，每条线满足\cref{g1:prop:axiom-line-parameter}的参数条件。
\end{proposition}
\begin{proof}
	两不同点 \(a,b\) 确定直线 \(a+\mathbb R(b-a)\)。若另一条直线包含它们，其方向空间必包含非零向量 \(b-a\)，因方向空间是一维的，它只能等于 \(\mathbb R(b-a)\)，故两线相同。每条线包含 \(a\) 和 \(a+u\)；点 \((0,0),(1,0),(0,1)\) 不共线。因此三条关联公理成立。

	固定 \(\ell=a+\mathbb R u\) 和 \(p\notin\ell\)。线 \(m=p+\mathbb R u\) 与 \(\ell\) 不交，否则 \(p-a\in\mathbb R u\)，与 \(p\notin\ell\) 矛盾。若过 \(p\) 的另一条线方向为 \(v\)，且 \(v\) 与 \(u\) 线性无关，则 \(u,v\) 是 \(\mathbb R^2\) 的一组基，所以方程 \(a+su=p+tv\) 有解，两线相交。因而不相交的线只能取方向 \(\mathbb R u\)，也就只能是 \(m\)。最后将 \(u\) 归一化为单位向量，\(\gamma(t)=a+tu\) 就是所需参数。
\end{proof}

同样的关联公理也允许另一种解释。取上半平面
\[
	\mathbb H=\{x+\mathrm i y\in\mathbb C:y>0\},
\]
把垂直于实轴的直线，以及圆心在实轴上的开上半圆，全部称为模型的“线”。\cref{g1:thm:hyperbolic-distance-geodesics}已证明，这些正是度量 \((\mathrm dx^2+\mathrm dy^2)/y^2\) 的完整测地线，每对点间的测地线段唯一且实现距离。现在先核验点线关联；这一步甚至还不需要距离公式。

\begin{proposition}[双曲平面的关联与平行]
	\label{g1:prop:hyperbolic-incidence-model}
	以上半平面 \(\mathbb H\) 为点集，以竖直线和实轴为圆心的开上半圆为线集，并以点属于曲线解释关联，得到 \(\mathsf I\) 的一个模型。该模型不满足 \(\mathsf P\)。
\end{proposition}
\begin{proof}
	取不同点 \(z_j=x_j+\mathrm i y_j\)，其中 \(y_j>0\)。若 \(x_1=x_2\)，竖直线 \(x=x_1\) 经过两点；一个圆心为 \((c,0)\) 的圆若也经过两点，则
	\[
		(x_1-c)^2+y_1^2=(x_2-c)^2+y_2^2
	\]
	迫使 \(y_1=y_2\)，与两点不同矛盾。因此竖直线是唯一选择。若 \(x_1\ne x_2\)，竖直线不可能同时经过两点，上式唯一确定圆心
	\[
		c=\frac{x_1^2+y_1^2-x_2^2-y_2^2}{2(x_1-x_2)}.
	\]
	半径 \(r=\sqrt{(x_1-c)^2+y_1^2}>0\) 随之确定，并且所得上半圆确实经过两点。这证明了第一条关联公理。每条竖直线和每个开上半圆都包含无穷多个点；而 \(\mathrm i,2\mathrm i,1+\mathrm i\) 不共线，因为前两点确定竖直线 \(x=0\)，第三点不在其上。

	为否定 \(\mathsf P\)，只需找到一组点线。令
	\[
		\ell=\{\mathrm i y:y>0\},\qquad p=1+\mathrm i.
	\]
	过 \(p\) 的竖直线 \(m_1=\{1+\mathrm i y:y>0\}\) 与 \(\ell\) 不交。开上半圆
	\[
		m_2=\{x+\mathrm i y:(x-1)^2+y^2=1,\ y>0\}
	\]
	也过 \(p\)，而且同样与 \(\ell\) 不交：令 \(x=0\) 后圆方程只允许 \(y=0\)，这不属于 \(\mathbb H\)。由于 \(m_1\ne m_2\)，唯一性失败。
\end{proof}

\begin{figure}[htbp]
	\centering
	% !TeX root = ../Book1.tex
\begin{tikzpicture}[>=stealth,font=\small,x=1.35cm,y=1.35cm]
	\begin{scope}
		\draw[thick,<->] (0,-.3)--(0,2.25) node[above] {\(\ell\)};
		\draw[blue!70!black,thick,<->] (1,-.3)--(1,2.25) node[above] {\(m\)};
		\fill (1,1) circle (1.5pt) node[right] {\(p\)};
		\node at (.5,-.65) {欧氏平面};
	\end{scope}
	\begin{scope}[xshift=4.5cm]
		\draw[dashed,gray,->] (-.35,0)--(2.5,0) node[right] {\(\mathbb R\)};
		\draw[thick,->] (0,0)--(0,2.25) node[above] {\(\ell\)};
		\draw[blue!70!black,thick,->] (1,0)--(1,2.25) node[above] {\(m_1\)};
		\draw[red!70!black,thick] (0,0) arc[start angle=180,end angle=0,radius=1];
		\node[red!70!black,right] at (1.8,.65) {\(m_2\)};
		\fill (1,1) circle (1.5pt) node[above right] {\(p=(1,1)\)};
		\foreach \x in {0,1,2} {
			\draw[gray,fill=white] (\x,0) circle (1.6pt);
		}
		\node[below] at (0,0) {\(0\)};
		\node[below] at (2,0) {\(2\)};
		\node at (1,-.65) {上半平面模型};
	\end{scope}
\end{tikzpicture}

	\caption[平行关系依赖所选模型]{左图在欧氏平面中，过 \(p\) 与 \(\ell\) 不交的直线只有 \(m\)。右图在上半平面模型中，\(m_1\) 与 \(m_2\) 是两条不同的模型直线，都经过 \(p\) 且不与 \(\ell\) 相交；空心端点位于被排除的实轴上。图只显示线的一部分。}
	\label{g1:fig:axiom-parallel-models}
\end{figure}

\cref{g1:fig:axiom-parallel-models}中，右侧半圆在图纸上弯曲，在模型中却是一条直线。把图纸上的欧氏直线含义强加给它，会混淆解释模型所用的环境与模型自身的关系。同样，\((0,0)\) 只是 \(\ell\) 和 \(m_2\) 的共同理想端点；它不是模型中的点，因而不能充当交点。

\begin{proposition}[双曲直线的实参数]
	\label{g1:prop:hyperbolic-line-parameter}
	以上半平面的双曲距离度量点集，并以竖直线及实轴为圆心的开上半圆为线，则每条线都有\cref{g1:prop:axiom-line-parameter}要求的等距参数 \(\mathbb R\rightarrow\ell\)。因此该模型具有该命题列出的介于、线段复制、合同相加及序完备性性质。
\end{proposition}
\begin{proof}
	对竖直线 \(x=c\)，取 \(\gamma(t)=c+\mathrm i e^t\)。它是从 \(\mathbb R\) 到该线的双射。由\eqref{g1:eq:half-plane-distance}，
	\[
		\cosh d\bigl(\gamma(s),\gamma(t)\bigr)
		=1+\frac{(e^s-e^t)^2}{2e^{s+t}}
		=\cosh(s-t).
	\]
	由于距离非负，且 \(\cosh\) 在 \([0,\infty)\) 上单调递增，距离就是 \(|s-t|\)。

	再取理想端点为实数 \(a<b\) 的上半圆。变换
	\[
		F(z)=\frac{z-a}{b-z}
	\]
	对应的实矩阵为 \(\begin{psmallmatrix}1&-a\\-1&b\end{psmallmatrix}\)，其行列式为 \(b-a>0\)；乘以 \((b-a)^{-1/2}\) 后行列式为 \(1\)，而分式不变。由\cref{g1:thm:hyperbolic-isometry-group}中已直接验证的度量保持公式，\(F\) 是上半平面的双曲等距。具体地，
	\[
		\operatorname{Im}F(z)=\frac{(b-a)y}{|b-z|^2},\qquad
		|F'(z)|=\frac{b-a}{|b-z|^2},
	\]
	所以长度元 \(|\mathrm dz|/y\) 保持不变；逆变换是 \(F^{-1}(w)=(bw+a)/(1+w)\)。圆方程为 \(y^2=(x-a)(b-x)\)，而
	\[
		\operatorname{Re}F(z)=\frac{(x-a)(b-x)-y^2}{|b-z|^2}.
	\]
	因此 \(F\) 将该开上半圆双射到正虚轴。用 \(\gamma(t)=F^{-1}(\mathrm i e^t)\) 即得所需参数，距离保持性将竖直线的结论传回原来的半圆。
\end{proof}

这个核验说明，双曲模型中并非连线段延长、长度复制和线内连续性也一并失效。两种模型在线内满足上述共同性质，却有不同的平行关系。但是，这仍只证明了已经列出的性质；如果把公理系统加强到包含所有平面顺序与角合同公理，就需要相应增加模型核验。

\subsection{独立性与相对一致性}
\label{g1:subsec:A01:03}

现在把模型比较转成关于证明的结论。设 \(T\) 是一组公理，\(\varphi\) 是用同一种语言写出的命题。记
\[
	\mathcal M\models T
\]
表示结构 \(\mathcal M\) 中 \(T\) 的每条公理都成立。例如，在本节的语言中，“任意点”遍历 \(\mathcal M\) 的点集，“存在一条线”遍历它的线集，原始关系 \(I\) 由模型指定的关联关系解释。这称为\term[manzuguanxi]{满足关系}[satisfaction relation]。\symidx{\(\mathcal M\models T\)：模型满足公理系统}记 \(T\vdash\varphi\) 表示从 \(T\) 出发，按照固定的一阶逻辑推演规则，经过有限步得到 \(\varphi\)；它说的是存在一份形式证明。\symidx{\(T\vdash\varphi\)：从公理系统形式推导命题}

\ProseParagraphBreak

两种记号之间，当前需要的是“证明不会把模型中的真命题变成假命题”这一方向。为看清它的依据，可考虑以逻辑公理、等式规则、分离规则和全称化组成的通常一阶推演：分离规则由 \(\alpha\) 及 \(\alpha\Rightarrow\beta\) 推出 \(\beta\)；全称化由对任意变量赋值均成立的 \(\alpha(x)\) 推出 \(\forall x\,\alpha(x)\)，且该变量不能在尚未解除的额外假设中自由出现。等式规则只允许用相等对象替换同一表达式中的对象。逻辑公理包括命题逻辑恒真式以及合法代入和量词分配等公理，它们的真值由逻辑连接词、量词与相等的解释确定，与选哪一个几何模型无关。

\begin{proposition}[推导的保真性]
	\label{g1:prop:geometric-soundness}
	设 \(T\) 是由一阶句子组成的公理系统，推导使用上述通常的一阶逻辑规则。若 \(\mathcal M\models T\) 且 \(T\vdash\varphi\)，则 \(\mathcal M\models\varphi\)。
\end{proposition}
\begin{proof}
	按一份有限证明中各行的次序归纳，证明每行在 \(\mathcal M\) 的任意变量赋值下成立。若该行是 \(T\) 的公理，它由模型假设成立；若是逻辑公理，它由逻辑语义成立。例如 \(\forall x\,\alpha(x)\Rightarrow\alpha(t)\) 在代入不使变量被错误绑定时成立，是因为全称命题包含了由项 \(t\) 指定的那个对象。量词分配
	\[
		\forall x\,(\alpha\Rightarrow\beta)
		\Rightarrow\bigl(\alpha\Rightarrow\forall x\,\beta\bigr)
	\]
	在 \(x\) 不自由出现在 \(\alpha\) 中时成立，是因为改变 \(x\) 的赋值不改变 \(\alpha\) 的真值。等式的自反性与替换规则来自相同对象具有相同解释。

	若该行由分离规则得到，归纳假设已知 \(\alpha\) 和 \(\alpha\Rightarrow\beta\) 对同一赋值为真，故 \(\beta\) 为真。若由全称化得到，归纳假设说明前一公式对 \(x\) 的每个允许取值都为真，因而其全称量化为真。这里 \(T\) 由句子组成，没有自由变量妨碍全称化；若在带临时假设的证明中使用该规则，就必须保留前述自由变量限制。由归纳，最后一行 \(\varphi\) 在模型中成立。
\end{proof}

\begin{definition}[独立性与一致性]
	\label{g1:def:axiom-independence}
	给定公理系统 \(T\) 与固定的推演规则，若 \(T\nvdash\varphi\) 且 \(T\nvdash\neg\varphi\)，则称命题 \(\varphi\) 对 \(T\) 是\term[dulixing]{独立的}[independent]。若不存在一个句子 \(\psi\) 使 \(T\vdash\psi\) 与 \(T\vdash\neg\psi\) 同时成立，则称 \(T\) 是\term[yizhixing]{一致的}[consistent]。
\end{definition}

\begin{theorem}[本节关联公理中的平行独立性]
	\label{g1:thm:incidence-parallel-independence}
	对\cref{g1:def:incidence-model}的三个关联公理 \(\mathsf I\)，唯一平行线命题 \(\mathsf P\) 是独立的。系统 \(\mathsf I\cup\{\mathsf P\}\) 和 \(\mathsf I\cup\{\neg\mathsf P\}\) 都有本节构造的模型。
\end{theorem}
\begin{proof}
	\cref{g1:prop:euclidean-incidence-model}给出一个满足 \(\mathsf I\) 与 \(\mathsf P\) 的模型，而\cref{g1:prop:hyperbolic-incidence-model}给出一个满足 \(\mathsf I\) 与 \(\neg\mathsf P\) 的模型。若 \(\mathsf I\vdash\mathsf P\)，\cref{g1:prop:geometric-soundness}便迫使双曲模型满足 \(\mathsf P\)，与已核验的两条不交线矛盾。若 \(\mathsf I\vdash\neg\mathsf P\)，同一保真性又迫使欧氏模型满足 \(\neg\mathsf P\)，与欧氏模型中存在且唯一的平行线矛盾。这排除了两个方向的形式证明。
\end{proof}

同样的论证解释了模型存在为何能支持一致性。若某个有模型的公理系统同时推出 \(\psi\) 与 \(\neg\psi\)，保真性便使同一模型中的同一句子既真又假，这是不可能的。这里没有使用“一切在所有模型中成立的命题都可证明”的反方向，也就没有调用逻辑完备性定理。

\ProseParagraphBreak

不过，我们用来构造模型的实数、集合、函数及有限推导本身，也处在一个背景理论中。上述证明说明：只要背景理论中的这些构造和推理有效，\(\mathsf I\) 的两个扩充就各有模型，因而不可能产生形式矛盾。这就是当前的\term[xiangduiyizhixing]{相对一致性}[relative consistency]结论所依赖的背景。它没有脱离一切数学前提证明一个绝对的“永不矛盾”；它把待研究的公理系统与已经承担构造任务的背景数学联系起来。

\ProseParagraphBreak

最后，独立性总是相对于写明的公理系统而言。从 \(\mathsf I\) 推不出 \(\mathsf P\)，并不自动说明向 \(\mathsf I\) 加入任意新公理后仍推不出它。经典几何基础中更强的任务，是保留选定体系的其他全部公理，再说明平行公设不能推出；这要求非欧模型也满足那些顺序、合同及连续性条件。Hilbert 在所据版本\cite[Chapter II, \S\S 9--10]{Hilbert1950Foundations}中先讨论公理的相容性，再讨论平行公理的独立性。阅读这种论证时，应把所保留的公理表与模型解释逐项对照；本节三个关联公理的定理则提供了一次可以直接检查的较小规模实例。

\section{坐标化及历史论证}
\label{g1:sec:A02}

把几何对象写成数对，可以将构造问题变成方程；但数对与原对象之间的对应仍须建立。先从已经给定的实仿射空间检验这项对应，再回到一项古典作图，就能分清两件事：现代坐标证明是否成立，以及它与原始证明实际保留了多少相同的假设和步骤。

\subsection{几何对象与代数模型的双向建立}
\label{g1:subsec:A02:01}

给平面上的点写上两个实数，并不会自动证明一个公理系统中的每个平面都能这样表示。坐标化至少包含两个不同的问题：已经有向量空间和平移运算时，怎样用坐标忠实地记录几何关系；只给抽象的点、线与若干关系时，又能否从这些关系中构造数域及向量运算。\chapref{g1:ch:01}解决的是前一个问题，其出发点在\cref{g1:def:affine-space}中已经写明：方向空间是有限维实向量空间，且任意两点之间有唯一的平移向量。本小节把这个来回构造及其条件展开，以便准确判断“化为代数问题”究竟保留了什么。

\ProseParagraphBreak

设 \(A\) 为 \(n\) 维实仿射空间，方向空间为 \(V\)，其中 \(n\ge1\)。选一个原点 \(o\in A\) 和 \(V\) 的基 \(e_1,\ldots,e_n\)，记这组选择为 \(\mathcal F\)。\cref{g1:prop:affine-coordinates}中的标架是 \((o,o+e_1,\ldots,o+e_n)\)。该命题的存在与唯一性给出互逆映射
\[
	\begin{aligned}
		\Phi_{\mathcal F}:A&\longrightarrow\mathbb R^n,
		& p=o+\sum_i x_ie_i&\longmapsto x,\\
		\Psi_{\mathcal F}:\mathbb R^n&\longrightarrow A,
		& x&\longmapsto o+\sum_i x_ie_i.
	\end{aligned}
\]
这是点集之间的双射。要说它们是几何结构的同构，还必须检查直线、介于关系和所要保留的其他关系怎样对应。

\begin{proposition}[坐标中的几何关系]\label{g1:prop:affine-model-relations}
	设 \(A\) 为 \(n\) 维实仿射空间，\(\mathcal F=(o;e_1,\ldots,e_n)\) 为一组仿射标架数据。若 \(p,q,r\) 的坐标分别为 \(x,y,z\)，则有以下对应。
	\begin{enumerate}
		\item 对 \(p\ne q\)，过 \(p,q\) 的直线在坐标中恰为 \(x+\mathbb R(y-x)\)。因此点在线上的关联关系被 \(\Phi_{\mathcal F}\) 及其逆保持。
		\item 对 \(p\ne q\)，\(r\) 严格介于 \(p,q\) 之间，当且仅当存在 \(0<t<1\) 使 \(z=(1-t)x+ty\)。这里严格介于的定义是 \(p\ne q\) 且 \(r=p+t(q-p)\)、\(0<t<1\)。
		\item 若 \(V\) 另有给定的正定内积，记 \(G=(\langle e_i,e_j\rangle)\)，则 \(G\) 对称正定，且
		\begin{equation}\label{g1:eq:coordinate-model-distance}
			d(p,q)^2=(y-x)^{\mathsf T}G(y-x).
		\end{equation}
		所以线段 \(pq\) 与 \(p'q'\) 合同，等价于它们的坐标差对二次型 \(G\) 具有相同取值。对以 \(p\) 为顶点、方向非零的两条射线 \(pq,pr\)，夹角 \(\theta\in[0,\pi]\) 满足
		\[
			\cos\theta=
			\frac{(y-x)^{\mathsf T}G(z-x)}
			{\sqrt{(y-x)^{\mathsf T}G(y-x)}\,
			\sqrt{(z-x)^{\mathsf T}G(z-x)}}.
		\]
	\end{enumerate}
\end{proposition}
\begin{proof}
	令 \(E:\mathbb R^n\to V\) 为 \(Ex=\sum_i x_ie_i\)。它是线性同构，且点差满足 \(q-p=E(y-x)\)。因而
	\[
		p+t(q-p)=o+E\bigl(x+t(y-x)\bigr).
	\]
	令 \(t\) 遍历全部实数，得到直线的精确对应；再将参数限制为 \(0<t<1\)，得到严格介于关系的对应。由于坐标映射双射，这些关系在两个方向都被保持。

	若有给定内积，双线性使 \(\langle Eu,Ev\rangle=u^{\mathsf T}Gv\)。矩阵 \(G\) 的对称性来自内积对称性；对非零 \(u\)，\(Eu\ne0\)，所以 \(u^{\mathsf T}Gu=\langle Eu,Eu\rangle>0\)。把 \(u=y-x\) 代入就得到距离公式。线段合同指两段距离相等，而距离非负，故等价于平方相等。夹角公式则是\cref{g1:def:inner-product}之后给出的内积夹角定义在此坐标中的表达；两射线方向非零使分母严格为正。\(\cos\) 在 \([0,\pi]\) 上单射，所以比较夹角也可等价地比较右侧数值。
\end{proof}

反方向从代数数据开始同样具体。给定 \(\mathbb R^n\) 和对称正定实矩阵 \(G\)，以坐标加法作平移，令直线为 \(x+\mathbb Rv\)、\(v\ne0\)，令严格介于由 \(0<t<1\) 的仿射组合定义，并令 \(\langle u,v\rangle_G=u^{\mathsf T}Gv\)。这满足\cref{g1:def:affine-space,g1:def:inner-product}的条件，因而给出一个实欧氏仿射空间。对原来的 \(A\) 按此办法重建后，\(\Psi_{\mathcal F}\) 恢复原点集、点差、直线、介于及内积，而不只给出一幅外形相似的图。若没有指定 \(G\)，相同的坐标数据只给出仿射结构及其基于实数次序的介于关系。

\ProseParagraphBreak

各项假设在这里承担不同任务。实向量空间给出加法与实数倍，所以能用一组基坐标化；实数的次序使 \(0<t<1\) 有意义；正定内积才给出欧氏长度和夹角。仅就仿射组合及介于的公式而言，可以把 \(\mathbb R\) 换成一个已经给定的有序域，但相应平方根是否存在、直线是否具备实数的连续性，还需另外讨论。不能因为写出了同样的式子，就把这些额外性质一起带入。

\ProseParagraphBreak

换一组标架会怎样？设新原点的旧坐标为 \(a\)，新基在旧基中的列矩阵为 \(S\in\operatorname{GL}(n,\mathbb R)\)。\cref{g1:prop:affine-coordinates}给出同一点的新旧坐标 \(x=a+Sx'\)，故
\begin{equation}\label{g1:eq:coordinate-model-change}
	x'=S^{-1}(x-a),\qquad G'=S^{\mathsf T}GS.
\end{equation}
第二式由新基向量的内积直接得到。下面的交换图表示先取旧坐标再换坐标，与直接取新坐标相同；左边始终是同一个几何点。
\[
	\begin{tikzcd}[column sep=large,row sep=large]
		A \arrow[r,"{\Phi_{\mathcal F}}"] \arrow[d,equal]
		& (\mathbb R^n,G) \arrow[d,"{x\mapsto S^{-1}(x-a)}"]\\
		A \arrow[r,"{\Phi_{\mathcal F'}}"']
		& (\mathbb R^n,G')
	\end{tikzcd}
\]
对点差 \(u=Su'\)，有 \(u^{\mathsf T}Gu=(u')^{\mathsf T}G'u'\)，所以两套坐标计算同一个距离。\(G\) 与 \(G'\) 的数值可以不同，几何长度并没有改变。特别地，非正交标架中的坐标差一般不能用各分量的平方和计算长度。

\ProseParagraphBreak

同一计算也解释了变换群怎样随保留结构而改变。在固定坐标模型中，\(x\mapsto Mx+b\)、\(M\in\operatorname{GL}(n,\mathbb R)\)，保持直线、介于及全部仿射组合；它保持 \(G\) 给出的数值距离，当且仅当 \(M^{\mathsf T}GM=G\)。必要性可先比较每个向量的长度平方，再用\cref{g1:prop:polarization}恢复两个向量的配对；充分性则由距离公式直接得到。标准正交坐标中的条件就是 \(M^{\mathsf T}M=I\)。这里只在描述一组已经证明为仿射的映射；若一开始只给一个集合映射，就不能未经论证把它写成矩阵加平移。

\begin{proposition}[连续的中点保持映射]\label{g1:prop:continuous-midpoint-affine}
	设 \(A,B\) 为有限维实仿射空间，赋予任一仿射坐标确定的通常拓扑。若映射 \(f:A\to B\) 在至少一个点连续，且对任意 \(p,q\in A\) 满足
	\[
		f\bigl(\tfrac12p+\tfrac12q\bigr)
		=\tfrac12 f(p)+\tfrac12 f(q),
	\]
	则 \(f\) 为实仿射映射。反之，每个实仿射映射都连续且保持中点。
\end{proposition}
\begin{proof}
	取 \(A\) 的一个原点 \(o\)，以 \(f(o)\) 为 \(B\) 的原点，并选两边方向基。于是只需考察一个保持中点、满足 \(F(0)=0\) 的映射 \(F:\mathbb R^n\to\mathbb R^m\)。中点条件对 \(0,2x\) 给出 \(F(2x)=2F(x)\)，所以
	\[
		F(x+y)=2F\bigl((x+y)/2\bigr)=F(x)+F(y).
	\]
	由加法性可依次得到整数倍以及有理数倍相容性，即 \(F(rx)=rF(x)\)、\(r\in\mathbb Q\)。若连续点为 \(x_0\)，则当 \(h\to0\) 时，\(F(h)=F(x_0+h)-F(x_0)\to0\)，所以 \(F\) 在零点连续；再由 \(F(x+h)-F(x)=F(h)\)，它在每一点都连续。

	对任意 \(t\in\mathbb R\)，取有理数列 \(r_j\to t\)。连续性使
	\[
		F(tx)=\lim_jF(r_jx)=\lim_jr_jF(x)=tF(x).
	\]
	因此 \(F\) 实线性，恢复原点便得 \(f\) 仿射。反之，仿射映射在坐标中为有限维线性映射加常向量，因而连续，并由\cref{g1:prop:affine-preservation}保持所有仿射组合，特别保持中点。
\end{proof}

这个命题把\cref{g1:thm:isometry-affine}证明中用到的恢复步骤分离出来，明确了连续性的作用：中点等式先恢复有理数倍，实数倍还须由极限恢复。对欧氏空间的距离保持映射，连续性来自距离等式，中点保持来自距离对中点的唯一刻画，所以那一定理能够进一步推出仿射性；它没有预先把仿射性放进“等距”的定义。

\begin{example}[关联和介于不足以确定合同]\label{g1:exm:incidence-order-not-congruence}
	在同一个点集 \(\mathbb R^2\) 上，使用同一族通常直线及同一介于关系，分别取
	\[
		G_1=\begin{pmatrix}1&0\\0&1\end{pmatrix},
		\qquad G_2=\begin{pmatrix}4&0\\0&1\end{pmatrix}.
	\]
	两种结构的点线关联及介于关系完全相同，但线段 \((0,0)(1,0)\) 与 \((0,0)(0,1)\) 在第一种结构中合同，在第二种结构中长度分别为 \(2,1\)，并不合同。因此关联和介于数据没有唯一指定合同关系。两模型作为未标记的欧氏空间仍然等距：\((x_1,x_2)\mapsto(2x_1,x_2)\) 将第二种距离变成第一种距离。这里不同的是同一批指定点之间保留了哪些关系，不能把这一差别误说成存在两种非同构的二维实欧氏空间。

	若只把 \(G_1\) 改为 \(cG_1\)、\(c>0\)，则所有长度同时乘 \(\sqrt c\)，线段之间是否合同和夹角大小都不变。因此即使给出了全部合同关系，数值长度仍需要选定单位。把“保持合同关系”与“逐个保持给定的数值距离”混为一谈，也会错误地删去统一伸缩这一自由度。
\end{example}

\chapref{g1:ch:02}的射影模型有同样需要说明的出发点。\cref{g1:def:projective-space}先给定实或复向量空间 \(V\)，再把一维子空间当作点；\cref{g1:prop:projective-charts}在一个非零齐次坐标处归一化，得到仿射图。这建立了向量空间与所构造射影空间之间的对应，却没有从任意抽象点线关联体系中推出底层数域。对于后一问题，域运算怎样定义、是否满足域公理、几何关系是否全部恢复，都需要专门的条件与证明。分清这两种坐标化任务，才能知道代数模型是在实现给定公理、描述一个已经构造的空间，还是承担某个几何对象的分类。

\subsection{历史证明的原始语境与现代重述}
\label{g1:subsec:A02:02}

一幅图中两条圆周相交，与一套公理能够保证它们相交，是两个需要分别说明的事实。阅读历史证明时，还须追问所读的文字属于哪个版本：是原著的译文，后人的编订，还是加入新公理后的重写。下面以一个等边三角形作图和两种平行条目为例，把文本中实际写出的论证，与本书在实坐标中补出的论证并排考察。

\ProseParagraphBreak

这里通过后世编订本考察 Euclid 的《几何原本》。\footnote{欧几里得（\OriginalName{Εὐκλείδης}，拉丁转写 Eukleidēs，约公元前300年活动，确切生卒年不详），古希腊数学家，活动于亚历山大里亚，以《几何原本》对几何与数论的系统论述著称。}具体使用的是 Casey 于1885年出版的第三版。\footnote{John Casey（约翰·凯西，1820--1891），爱尔兰数学家，研究初等几何与圆的几何，编写过《Euclid 几何原本》及其续编。}它在题名中明确列出注释与习题，并保留了大量编者说明，因此以下只把它当作十九世纪的一个可核查文本，不把它逐字归给古希腊原著。该版 Book I 的 Prop. I 要求在给定的有限直线上构造等边三角形；其中的 \emph{finite right line} 指两端已给定的直线段，\emph{right line} 在这里是直线，并非带有直角条件的线。该版在 Postulates, I 后也明确解释了这个用语。\cite[Book I, Postulates, I; Prop. I]{Casey1885Euclid}

\ProseParagraphBreak

把已给线段记为 \(AB\)，该文本的构造分为三个动作：以 \(A\) 为圆心、\(AB\) 为半径作圆；以 \(B\) 为圆心、\(BA\) 为半径作第二个圆，并取交点 \(C\)；连接 \(CA,CB\)。随后利用同圆半径相等得到 \(AC=AB\) 与 \(BC=BA\)，再用“等于同一量的量彼此相等”得到三边相等。文本对作圆和连接两点分别援引第三、第一公设，却没有在取交点处另给一个保证交点存在的论证。它证明了交点一旦取得便给出所求三角形；若要把这段论证改写成明确的现代公理演绎，交点的存在及其不在直线 \(AB\) 上，还须从所选公理推出。

\ProseParagraphBreak

在本书已经构造的实欧氏平面中，可以直接补上这一步。先选择一个定向，令 \(J\) 表示向该正向旋转 \(\pi/2\) 的线性映射。若 \(v=B-A\ne0\)，则 \(Jv\perp v\)、\(\lVert Jv\rVert=\lVert v\rVert\)。为了找同时到 \(A,B\) 等距的点，先取中点，再沿垂直方向移动。所需移动量由勾股关系确定为边长的 \(\sqrt3/2\) 倍。

\begin{proposition}[两个圆的交点与等边三角形]\label{g1:prop:equilateral-coordinate-reconstruction}
	设 \(A,B\) 是实欧氏平面中两个不同的点，固定该平面的一个定向，并令 \(J\) 为正向直角旋转。以 \(A,B\) 为圆心、以 \(\lVert B-A\rVert\) 为共同半径的两个圆周，恰好交于
	\begin{equation}\label{g1:eq:equilateral-coordinate-intersections}
		C_\pm=\frac{A+B}{2}\pm\frac{\sqrt3}{2}J(B-A).
	\end{equation}
	两个交点都不在直线 \(AB\) 上，且 \(ABC_+\)、\(ABC_-\) 都是等边三角形。
\end{proposition}
\begin{proof}
	记 \(v=B-A\)、\(r=\lVert v\rVert>0\)、\(M=(A+B)/2\)。因 \(v,Jv\) 为一组正交基，任意点 \(C\) 唯一写成 \(M+sv+tJv\)。到两个圆心的距离平方分别为
	\[
		\begin{aligned}
			\lVert C-A\rVert^2&=r^2\bigl((s+\tfrac12)^2+t^2\bigr),\\
			\lVert C-B\rVert^2&=r^2\bigl((s-\tfrac12)^2+t^2\bigr).
		\end{aligned}
	\]
	两式都等于 \(r^2\) 时，相减得 \(2s=0\)，代回得 \(t^2=3/4\)。实数中恰有 \(t=\pm\sqrt3/2\) 两解，所以交点恰为所列两点。它们的 \(Jv\) 分量非零，因而不在直线 \(AB\) 上；两个圆周的方程又给出 \(AC_\pm=BC_\pm=AB\)，即得等边性。
\end{proof}

% !TeX root = ../Book1.tex
\begin{figure}[htbp]
	\centering
	\begin{tikzpicture}[scale=2,font=\small]
		\draw[GeometryGray] (0,0) circle[radius=1];
		\draw[GeometryGray] (1,0) circle[radius=1];
		\fill[GeometryBlue!8] (0,0)--(1,0)--(0.5,{sqrt(3)/2})--cycle;
		\draw[thick,GeometryBlue] (0,0)--(1,0)--(0.5,{sqrt(3)/2})--cycle;
		\draw[GeometryGray,dashed] (0.5,{-sqrt(3)/2})--(0.5,{sqrt(3)/2});
		\draw[GeometryGray,dashed] (0,0)--(0.5,{-sqrt(3)/2})--(1,0);
		\foreach \x/\y in {0/0,1/0,0.5/0}{\fill (\x,\y) circle[radius=0.025];}
		\fill[GeometryBlue] (0.5,{sqrt(3)/2}) circle[radius=0.025];
		\fill[GeometryBlue] (0.5,{-sqrt(3)/2}) circle[radius=0.025];
		\node[below left] at (0,0) {\(A\)};
		\node[below right] at (1,0) {\(B\)};
		\node[below right] at (0.5,0) {\(M\)};
		\node[above] at (0.5,{sqrt(3)/2}) {\(C_+\)};
		\node[below] at (0.5,{-sqrt(3)/2}) {\(C_-\)};
		\node[right] at (0.5,0.44) {\(h\)};
	\end{tikzpicture}
	\caption[等边三角形作图中的两个交点]{在实欧氏平面中，两圆的共同点位于 \(AB\) 的中垂线上，距中点 \(M\) 的长度为 \(h=\sqrt3\lVert B-A\rVert/2\)。图画出的是实数模型；若改变点集，图中可见的位置未必仍是该模型中的点。}
	\label{g1:fig:equilateral-historical-construction}
\end{figure}


\cref{g1:fig:equilateral-historical-construction}把这两个交点画在实平面中。命题的证明使用了实数域、向量运算、内积和正平方根，故它是在本书所用模型中对作图的现代重述。它也说明相同的公式在适当的有序子域上仍可能有效，但并没有把这些代数与度量结构从一套更少的抽象公理中重新构造出来。若研究的问题是某组关联、顺序及合同公理是否足够，必须继续追查这些运算与交点存在性的公理依据。

\begin{example}[两个圆周分别存在而没有共同点]\label{g1:exm:rational-circles-no-intersection}
	把点集改为 \(\mathbb Q^2\)，仅考察由有理系数方程定义的两个集合
	\[
		\begin{aligned}
			S_A&=\{(x,y)\in\mathbb Q^2:x^2+y^2=1\},\\
			S_B&=\{(x,y)\in\mathbb Q^2:(x-1)^2+y^2=1\}.
		\end{aligned}
	\]
	它们都非空，例如 \((1,0)\in S_A\)、\((0,0)\in S_B\)。但若某点同时属于两者，相减得 \(x=1/2\)，随后得到 \(y^2=3/4\)，所以 \(2y\) 必须是 \(3\) 的有理平方根。这不可能：若互素整数 \(m,n\)、\(n\ne0\) 满足 \(m^2=3n^2\)，则 \(3\mid m\)，代回又得 \(3\mid n\)，矛盾。因此这两个集合不相交。

	这个计算只反驳“定义了两个圆周，就已保证它们具有所画出的交点”这一推断，并未证明某套完整几何公理中的独立性。事实上，\(\mathbb Q^2\) 配通常的距离相等关系就不满足任意射线上的线段搬移：从 \((0,0)\) 沿正 \(x\) 轴，无法取得与线段 \((0,0)(1,1)\) 等长的有理点，因为该点横坐标必须为 \(\sqrt2\)。所以不能把本例直接称作满足全部关联、顺序和合同公理的平面。
\end{example}

不同版本对平行条件的组织也不能只凭熟悉的名称判断。Casey 第三版把作图要求列为三个 Postulates，而把一个平行条件列在 Axioms, XII：若两条直线与第三条直线相遇，所成的两个同侧内角之和小于两直角，则延长这两条直线后，它们在有限距离处相遇。\cite[Book I, Axioms, XII]{Casey1885Euclid}这里的假设涉及第三条直线、内角以及角和比较；不能把所读条目的编号改写成“该版第五公设”。另一方面，Hilbert 的 Townsend 英译本以“过直线外一点恰有一条与该直线不相交的直线”表述平行公理，并随即把它拆为存在与唯一两个断言。\cite[Chapter I, §5]{Hilbert1950Foundations}

\ProseParagraphBreak

这两种表述提出的判定方式不同。前者由角的条件推出相交，后者直接限制不相交直线的条数。要在一个抽象体系内把它们互相推出，须先确定共同的其他公理，并检查角的作出、角与顺序的关系以及所用三角形结论怎样得到；把一句话换成另一句话本身不能代替这些推导。在实欧氏模型中，两种条件均可用已有坐标和角度计算检验，这也只说明它们在该模型中成立，不能单凭这个事实得到任意公理模型中的等价性。

\ProseParagraphBreak

因此，重述历史证明时应保留两条可查的联系：一条从具体版本通向它实际采用的词义、假设与推理；另一条从现代陈述通向所选公理或模型中的证明。前者帮助我们避免把编者补充写成原著文字，后者帮助我们说明补入的平方根、连续性或唯一性究竟解决了哪一个问题。等边三角形的例子所揭示的关键，正是“构造得到的点满足要求”与“所需的点确实存在”这两项证明任务。

\begin{chaptersummary}
	公理规定对象之间应满足的关系，模型为对象与关系提供解释。关联条件不决定长度，直线上的顺序也不自动提供连续性；增加合同与连续性条件时，需要明确新增了哪些可用的构造和结论。在坐标模型中，实数的代数、顺序与完备性分别承担不同的存在性依据。

	\ProseParagraphBreak

	欧氏平面的直线和上半平面的双曲直线都满足本附录列出的关联公理，而唯一平行命题在两者中取不同真值。推理的保真性由此排除了从这组公理推出该命题或其否定的可能。独立性总是相对于指定公理系统的；模型存在所依赖的背景理论也必须保留，因而相对一致性并不是一个无条件的绝对保证。

	\ProseParagraphBreak

	给定实仿射空间以后，标架把点对应到实数坐标，直线和介于关系也随之得到代数表达；增加内积后，Gram 矩阵再记录长度与角度。改变标架会改变分量，却不改变原几何对象。仅保留关联的变换可以改变合同关系，这说明选择坐标与选择几何结构是两个不同步骤。

	\ProseParagraphBreak

	阅读历史证明时，应先恢复原文中的对象、构造要求和已用命题，再比较现代重述。等边三角形的坐标构造给出了明确交点，但它使用的实数运算并非原文的逐步翻译。原始证明、补足存在性依据后的重述和另一种现代证明，各有明确的内容，不能只凭结论相同便视为同一个论证。
\end{chaptersummary}

\BookExercises
以下练习分别考察有限模型、结构之间的区别、坐标数据的恢复及构造所需的数域。每题都应先明确正在保留哪些关系，再判断结论能否推出。

\begin{exercise}
	\label{g1:ex:four-point-incidence}
	令 \(X=\{a,b,c,d\}\)，把 \(X\) 的每个二元子集都当作一条线，关联为点属于子集。
	\begin{enumerate}
		\item 核验三个关联公理 \(\mathsf I\)，并证明该模型满足唯一平行线命题 \(\mathsf P\)。
		\item 证明：在这个模型中，无法定义一种严格介于关系，同时使“介于的三点互异且共线”和“任意线段都能向一端延长”成立。
	\end{enumerate}
\end{exercise}
\begin{solution}
	(1) 两个不同点恰好确定由它们组成的二元子集，所以连线存在且唯一；每条线正有两个点；任取三个不同点，它们不可能同在一个二元子集中，故不共线。三条关联公理都成立。固定线 \(\ell\) 与其外一点 \(p\)，补集 \(X\setminus\ell\) 恰有两个点并包含 \(p\)，所以它是一条过 \(p\) 且与 \(\ell\) 不交的线。任何与 \(\ell\) 不交的线都是这个二元补集，故唯一。

	(2) 取任意两个不同点 \(a,b\)。若线段能延长，就应有点 \(e\) 使 \(B(a,b,e)\) 成立。严格介于要求 \(a,b,e\) 互异且共线，但经过 \(a,b\) 的唯一线只有 \(a,b\) 两点，不可能含有第三点。因此关联公理加上唯一平行线，并没有保证通常的线内顺序与延长性质。
\end{solution}

\begin{exercise}
	\label{g1:ex:seven-point-incidence}
	在二元域 \(\mathbb F_2=\{0,1\}\) 上取向量空间 \(V=\mathbb F_2^3\)。以 \(V\setminus\{0\}\) 为点集；对每个非零线性函数 \(f:V\rightarrow\mathbb F_2\)，取 \(\ker f\setminus\{0\}\) 为一条线。
	\begin{enumerate}
		\item 证明共有七点、七线，每条线有三个点，且满足 \(\mathsf I\)。
		\item 证明任意两条不同的线恰好有一个公共点。说明 \(\mathsf P\) 在此模型中的失败与双曲模型中的失败有什么不同。
	\end{enumerate}
\end{exercise}
\begin{solution}
	(1) \(V\) 有 \(2^3=8\) 个向量，去掉零向量后有七点。对偶空间也有八个元素，故有七个非零线性函数。每个非零 \(f\) 的核是二维空间，含四个向量，所以对应的线有三个点。不同的非零函数给出不同的核：若两个函数核相同，它们都在核上取零，并在核外取唯一的非零标量 \(1\)，于是函数相同。这也证明共有七条线。

	两个不同的非零向量 \(u,v\) 线性无关，因为 \(\mathbb F_2\) 的非零标量只有 \(1\)。经过两点的线必须由一个在 \(\operatorname{span}(u,v)\) 上为零的非零线性函数给出。这些函数构成对偶空间中一个一维子空间，因而恰有一个非零元素，证明连线存在且唯一。每线已有三个点；而标准基 \(e_1,e_2,e_3\) 不共线，因为非零函数不能同时在一组基上取零。故 \(\mathsf I\) 成立。

	(2) 若 \(f,g\) 非零且不同，则它们在对偶空间中线性无关。线性映射 \(v\mapsto(f(v),g(v))\) 的秩为 \(2\)，其核 \(\ker f\cap\ker g\) 维数为 \(1\)，因此恰含一个非零向量。该向量就是两线唯一的公共点。

	任取一条线及其外一点，所有过该点的线都与原线相交，没有不交线。因此这里破坏的是平行线的存在性。双曲模型的既有反例则给出了同一外点处至少两条不交线，破坏的是唯一性。两者都否定 \(\mathsf P\)，却不是同一种几何情形。
\end{solution}

\begin{exercise}
	\label{g1:ex:taxicab-betweenness}
	在 \(\mathbb R^2\) 上保留通常的仿射直线，改用距离
	\[
		d_1\bigl((x_1,y_1),(x_2,y_2)\bigr)=|x_1-x_2|+|y_1-y_2|.
	\]
	\begin{enumerate}
		\item 核验 \(d_1\) 是距离，并为每条仿射直线构造一个保持距离的实参数。
		\item 取 \(A=(0,0),B=(1,0),C=(1,1)\)，证明距离可加，但三点不共线。
		\item 以距离相等定义线段合同，找出两条 \(d_1\)-合同却并非欧氏等长的线段。说明点线关联相同为何不能决定合同关系。
	\end{enumerate}
\end{exercise}
\begin{solution}
	(1) 非负性、对称性以及距离为零恰在两点相同时，均来自实数绝对值。把两个坐标各自的三角不等式相加，得到 \(d_1(P,R)\le d_1(P,Q)+d_1(Q,R)\)。对直线 \(a+\mathbb R u\)，其中 \(u=(u_1,u_2)\ne0\)，令 \(c=|u_1|+|u_2|>0\)，取
	\[
		\gamma(t)=a+\frac{t}{c}u.
	\]
	它是从 \(\mathbb R\) 到该线的双射，且 \(d_1(\gamma(s),\gamma(t))=|s-t|\)。因此正文列出的线内顺序、唯一复制长度与序完备性仍成立。

	(2) 三段距离为 \(d_1(A,B)=1\)、\(d_1(B,C)=1\)、\(d_1(A,C)=2\)，满足可加等式。然而 \(B-A=(1,0)\) 与 \(C-A=(1,1)\) 线性无关，所以三点不共线。若删除正文定义中独立列出的共线条件，就会错误地把 \(B\) 宣布为介于 \(A,C\) 的点。

	(3) 取从原点到 \((1,0)\) 及从原点到 \((1/2,1/2)\) 的两段。它们的 \(d_1\)-长度都是 \(1\)，而欧氏长度分别为 \(1\) 与 \(1/\sqrt2\)。两种距离是在同一仿射点线结构上增加的不同数据，所以保留全部关联关系仍可改变合同判断。
\end{solution}

\begin{exercise}
	\label{g1:ex:model-inference-limits}
	设 \(T\) 是使用可靠的一阶推演规则的公理系统，\(\varphi\) 是同语言中的句子，模型 \(\mathcal M,\mathcal N\) 满足
	\[
		\mathcal M\models T\cup\{\varphi\},\qquad
		\mathcal N\models T\cup\{\neg\varphi\}.
	\]
	\begin{enumerate}
		\item 证明 \(\varphi\) 对 \(T\) 独立，并说明两个扩充系统的一致性依据。
		\item 有人说：“因为 \(T\) 有模型，所以 \(T\) 的所有模型都同构。”以上材料能否支持这一结论？
		\item 设 \(T'=T\cup\{\varphi\}\)。解释为何 \(\varphi\) 对 \(T\) 独立，却不对 \(T'\) 独立；指出原来的模型反证在哪一步不能用于 \(T'\)。
	\end{enumerate}
\end{exercise}
\begin{solution}
	(1) 若 \(T\vdash\varphi\)，保真性使 \(\mathcal N\models\varphi\)，与 \(\mathcal N\models\neg\varphi\) 矛盾；若 \(T\vdash\neg\varphi\)，同理与 \(\mathcal M\models\varphi\) 矛盾，所以两个方向都不能证明。每个扩充系统都有题设给出的模型；若其中一个推出某句子及其否定，该模型就会同时满足二者。因此在构造模型并验证推演保真性的背景理论中，两个系统均一致。

	(2) 不能。对同一种语言，同构在每个基本关系上保持真假，因而对由这些关系经逻辑连接和量化组成的句子也保持真假。具体地，连接词的步骤由真值表保证；存在量词的见证可由同构映到另一模型，全称量词则由同构的满射性保证。对句子的构成归纳，便知同构模型必须在 \(\varphi\) 上同真同假。题设的 \(\mathcal M,\mathcal N\) 却在 \(\varphi\) 上取相反真值，所以它们恰好给出了不可能同构的两个模型。

	(3) \(\varphi\) 已是 \(T'\) 的一条公理，故有一行证明 \(T'\vdash\varphi\)，独立性失效。\(\mathcal N\) 虽然满足 \(T\)，却不满足新加入的 \(\varphi\)，因此不是 \(T'\) 的模型。把它用于排除从 \(T'\) 出发的证明，违反了保真性命题中的模型假设。这说明独立性结论必须连同基准公理系统一起陈述。
\end{solution}

\begin{exercise}\label{g1:ex:oblique-frame-model}
	在通常欧氏平面中取原点 \(o=(1,-1)\)，方向基 \(e_1=(2,0)\)、\(e_2=(1,1)\)，记一个点的新坐标为 \((u,v)\)。求此标架中的内积矩阵、点 \(p=(5,1)\) 的坐标及 \(d(o,p)\)，并用 \((u,v)\) 写出线段 \(op\) 的中垂线方程。若方向基改为 \(e'_1=e_1\)、\(e'_2=e_1+e_2\)，原点不变，求新内积矩阵及 \(p\) 的新坐标，并核对长度不变。
\end{exercise}
\begin{solution}
	原来的直角坐标与 \((u,v)\) 的关系为
	\[
		(x,y)=(1,-1)+u(2,0)+v(1,1)=(1+2u+v,-1+v).
	\]
	所以 \(p\) 的坐标为 \((1,2)\)。基向量的配对给出
	\[
		G=\begin{pmatrix}4&2\\2&2\end{pmatrix},\qquad
		d(o,p)^2=\begin{pmatrix}1&2\end{pmatrix}G\begin{pmatrix}1\\2\end{pmatrix}=20.
	\]
	因此 \(d(o,p)=2\sqrt5\)。若把这套非正交坐标误当成标准正交坐标，会错误地得到 \(\sqrt{1^2+2^2}=\sqrt5\)。

	令 \(x=(u,v)^{\mathsf T}\)、\(a=(1,2)^{\mathsf T}\)。与 \(o,p\) 等距的条件为 \(x^{\mathsf T}Gx=(x-a)^{\mathsf T}G(x-a)\)，消去二次项得到 \(2x^{\mathsf T}Ga=a^{\mathsf T}Ga\)。由于 \(Ga=(8,6)^{\mathsf T}\)，中垂线为 \(4u+3v=5\)。此处的线性方程来自度量比较，而不是把纸面坐标轴预设为垂直。

	换基矩阵为 \(S=\begin{psmallmatrix}1&1\\0&1\end{psmallmatrix}\)，满足旧坐标列等于 \(S\) 乘新坐标列。因此
	\[
		G'=S^{\mathsf T}GS=\begin{pmatrix}4&6\\6&10\end{pmatrix},
		\qquad p'=S^{-1}\begin{pmatrix}1\\2\end{pmatrix}=\begin{pmatrix}-1\\2\end{pmatrix}.
	\]
	直接计算 \((p')^{\mathsf T}G'p'=4-24+40=20\)，恢复同一个长度平方。
\end{solution}

\begin{exercise}\label{g1:ex:congruence-determines-scale}
	设 \(b_1,b_2\) 是同一个非零有限维实向量空间 \(V\) 上的两个正定内积，令 \(q_j(v)=b_j(v,v)\)。假设它们给出的等长关系相同，即对全部 \(u,v\in V\)，有
	\[
		q_1(u)=q_1(v)\quad\Longleftrightarrow\quad q_2(u)=q_2(v).
	\]
	证明存在常数 \(c>0\) 使 \(b_2=cb_1\)。若再指定一个非零向量在两个度量下长度相同，证明两个内积必相同。解释这与“合同关系还没有指定长度单位”的联系。
\end{exercise}
\begin{solution}
	选 \(q_1(e)=1\) 的向量 \(e\)，令 \(c=q_2(e)>0\)。任意满足 \(q_1(u)=1\) 的向量都与 \(e\) 对 \(q_1\) 等长，故由假设也有 \(q_2(u)=c\)。对非零 \(v\)，将其写成 \(v=\sqrt{q_1(v)}\,u\)，其中 \(q_1(u)=1\)，利用二次齐次性得到
	\[
		q_2(v)=q_1(v)q_2(u)=c\,q_1(v).
	\]
	零向量也满足此式。由极化恒等式，
	\[
		b_2(u,v)=\frac{q_2(u+v)-q_2(u)-q_2(v)}2
		=c\,b_1(u,v).
	\]
	若某个 \(v_0\ne0\) 在两种度量下长度相同，则 \(q_2(v_0)=q_1(v_0)>0\)，代入上式得 \(c=1\)。在已有实仿射和内积结构的范围内，全部线段合同关系因而确定内积到一个正比例因子；再用一条非零线段校准单位，才确定具体的数值长度。
\end{solution}

\begin{exercise}\label{g1:ex:frame-distances-rigidity}
	设 \(f:A\to B\) 是两个 \(n\) 维实欧氏仿射空间之间的连续中点保持映射，其中 \(n\ge1\)。取 \(A\) 的仿射标架 \((p_0,\ldots,p_n)\)，并假设
	\[
		d_B(f(p_i),f(p_j))=d_A(p_i,p_j)\qquad(0\le i,j\le n).
	\]
	证明 \(f\) 保持任意两点的距离。若只要求 \(p_0p_i\)、\(1\le i\le n\)，这 \(n\) 条线段的长度保持，结论是否仍成立？给出二维反例。
\end{exercise}
\begin{solution}
	由\cref{g1:prop:continuous-midpoint-affine}，\(f\) 有实线性部分 \(L\)。令 \(e_i=p_i-p_0\)，这些向量组成 \(A\) 的方向空间的基。给定距离等式先给出 \(\lVert Le_i\rVert=\lVert e_i\rVert\)，又给出 \(\lVert Le_i-Le_j\rVert=\lVert e_i-e_j\rVert\)。用极化恒等式得到
	\[
		\langle Le_i,Le_j\rangle_B
		=\frac{\lVert Le_i\rVert_B^2+\lVert Le_j\rVert_B^2-\lVert Le_i-Le_j\rVert_B^2}{2}
		=\langle e_i,e_j\rangle_A.
	\]
	对基展开并用双线性，可知 \(L\) 保持所有向量的内积，特别保持长度。又因 \(f(q)-f(p)=L(q-p)\)，任意两点距离均被保持。正定性还说明 \(L\) 单射，同维有限维性使它可逆，所以 \(f\) 实际为等距同构。

	只保留从 \(p_0\) 出发的边长不足以控制不同基方向之间的夹角。在通常欧氏平面中取标准标架 \(p_0=(0,0),p_1=(1,0),p_2=(0,1)\)，并令
	\[
		f(x,y)=\bigl(x+\tfrac12y,\tfrac{\sqrt3}{2}y\bigr).
	\]
	它线性，因而连续且保持中点。两向量 \(f(p_1)=(1,0)\)、\(f(p_2)=(1/2,\sqrt3/2)\) 都长为 \(1\)，满足所述 \(p_0p_i\) 的边长条件；但它们之间的距离为 \(1\)，而 \(d(p_1,p_2)=\sqrt2\)。所以 \(f\) 不是等距映射，遗漏的标架边恰好承载了夹角信息。
\end{solution}

\begin{exercise}\label{g1:ex:one-circle-construction-not-all}
	令 \(K=\mathbb Q(\sqrt3)\subset\mathbb R\)，以 \(K^2\) 为点集，距离平方使用通常的 \((x_1-y_1)^2+(x_2-y_2)^2\)。证明每一条端点在 \(K^2\) 中的非零线段都能在 \(K^2\) 中作为某个等边三角形的一边。再证明以 \((0,0)\)、\((1,0)\) 为圆心、共同半径为 \(\sqrt3\) 的两个圆周，在 \(K^2\) 中各自非空却没有共同点。解释为什么完成第一项构造仍不足以保证所有实平面中相交的圆在这个点集中相交。
\end{exercise}
\begin{solution}
	映射 \(J(x,y)=(-y,x)\) 保持 \(K^2\)，而 \(1/2\)、\(\sqrt3/2\) 都属于 \(K\)。因此对 \(A,B\in K^2\)、\(A\ne B\)，式\eqref{g1:eq:equilateral-coordinate-intersections}给出的两个点仍在 \(K^2\) 内，且同样的距离平方计算给出三边相等。非共线性也由垂直分量非零得到。

	第二组圆周各含点 \((\sqrt3,0)\) 与 \((1+\sqrt3,0)\)，故分别非空。共同点必须满足
	\[
		x^2+y^2=3,\qquad (x-1)^2+y^2=3.
	\]
	两式相减给出 \(x=1/2\)，代回得 \(y^2=11/4\)。如果 \(y\in K\)，则 \(\sqrt{11}\in K\)。设 \(\sqrt{11}=a+b\sqrt3\)、\(a,b\in\mathbb Q\)，平方后由 \(1,\sqrt3\) 在 \(\mathbb Q\) 上线性无关得到 \(2ab=0\) 和 \(a^2+3b^2=11\)。若 \(b=0\)，则 \(a^2=11\)；若 \(a=0\)，则 \(b^2=11/3\)。这两种情形都不可能，因为一个非零有理数的平方在约分后的分子、分母中，每个素因子的指数都为偶数，而 \(11\) 和 \(11/3\) 不具备该性质。因此两圆周没有 \(K^2\) 中的共同点。

	在 \(\mathbb R^2\) 中，这两个圆确实交于 \((1/2,\pm\sqrt{11}/2)\)。第一项构造只需要 \(\sqrt3\)，第二项交点却需要额外的平方根 \(\sqrt{11}\)。一种作图能够完成，并不意味着其他作图所需的代数运算也已经在同一个点集中封闭。本题没有断言 \(K^2\) 满足一套完整的中性几何公理。
\end{solution}


% !TeX root = ../../Book1.tex
% 纲要标题：## 附录 B　离散与凸几何
% 纲要位置：几何学计划纲要.md，第 935 行。
% 正文状态：初稿；尚未正式审核。

% 纲要内容开始
% 
% > 从正文凸方法进入格与离散结构；给出可核验定理和模型，不扩为组合优化全书。
% 
% * **格点与多面体**
%   * 格、基本平行多面体及体积
%   * 二维 Pick 公式及三角剖分证明
%   * Ehrhart 多项式的代表性陈述与例子
% 
% * **Voronoi 与 Delaunay**
%   * 距离平方与分区
%   * 空球条件和对偶分解
%   * 退化位置及一般位置的区别
% 
% * **支撑函数与离散逼近**
%   * 凸体的多面体逼近
%   * Hausdorff 距离与支撑函数
%   * 数值表示的证书
% 
% ---
% 
% 纲要内容结束

\chapter{离散与凸几何}
\label{g1:app:B}
\BookChapterNumber{g1:app:B}

\begin{chapterabstract}
	凸集的形状可以由顶点、半空间或支撑函数描述；加入整数格与有限采样点之后，还要辨别计数、分解和逼近各自保留了哪些信息。本附录从格的基本区域出发，用格三角剖分证明二维 Pick 公式，计算若干 Ehrhart 多项式；再由距离平方的线性差与抛物面提升建立 Voronoi 分区和 Delaunay 分解，说明共球位置为何产生非单纯形单元。最后用有限方向上的支撑值构造内外多面体，并把 Hausdorff 误差转成可检查的函数界与有限证书。
\end{chapterabstract}

\begin{learninggoals}
	\begin{itemize}
		\item 从格基计算基本区域的体积，辨别保持格的坐标变换与一般实线性变换。
		\item 用三角剖分证明凸格多边形的 Pick 公式，计算膨胀后多边形、箱体与单形的格点数。
		\item 把最近站点的距离比较写成线性不等式，并用空球条件识别相应的 Delaunay 单元。
		\item 区分一般位置的单纯形剖分与共球退化的多胞形分解，说明额外选择对角线改变了什么。
		\item 用支撑函数计算紧凸集的 Hausdorff 距离，为有限采样构造带误差界的内外逼近。
		\item 检验一个有限数值表示所附的包含关系与误差证书，分清近似结果与已获证明的结论。
	\end{itemize}
\end{learninggoals}

边长为 \(m\) 的正方形面积是 \(m^2\)，但当 \(m\) 为非负整数时，其中连同边界共有 \((m+1)^2\) 个整数点。两者的差为 \(2m+1\)：面积刻画了计数的主要增长，边界与顶点还留下了较低阶的信息。若把正方形剪成三角形，公共边上的格点就会被重复计入；要从局部计数恢复整体，必须把这种重复明确扣除。Pick 公式与本附录的 Ehrhart 计数正是从这一问题展开。

\ProseParagraphBreak

有限数据还可能承担另外两种任务。给定若干站点，可以按“离谁最近”划分整个空间；给定一个凸体，可以只保留有限个支撑方向，尝试用多面体近似它。前一个问题要求公共边界与退化单元的定义准确，后一个问题要求未采样方向上的误差仍受控制。画出一个看似合理的分区，或计算出一批精确小数，都不能独自回答这些问题。

\ProseParagraphBreak

本附录使用\chapref{g1:ch:03}的有限维凸性。\cref{g1:thm:polytope-hv}把有界半空间交与有限点凸包联系起来，\cref{g1:prop:support-function-representation}用全部方向的支撑值恢复闭凸集，\cref{g1:prop:distance-support-formula}则由最近点给出距离的支撑表示。下文会在使用时恢复所需条件；另需的线性代数、整数最大公因数、有限求和与欧氏体积均在具体计算中说明。格点计数的补充阅读可查\cite[Sections 2.2--2.3, 2.6, 3.1--3.5]{BeckRobins2015}；这里完整证明所选模型，不以算例代替所有格多胞形的一般理论。

\section{格点与多面体}
\label{g1:sec:B01}

把一个三角形放大两倍，面积变成原来的四倍，里面的整数点数却未必恰好乘以四：边界也会经过新的整数点。要把面积和点数联系起来，先要说明哪些点算作整数点，以及一份基本体积有多大。本节从格的基本平行多面体出发，在平面中把多边形分成面积最小的格三角形，再用同一计数解释整数倍放大时出现的多项式。

\subsection{格、基本平行多面体及体积}
\label{g1:subsec:B01:01}

\secref{g1:sec:0304}已经把格引入为一组实基的整数线性组合。这里保持这个定义，并把计数时所需的换基与体积性质展开。设 \(E\) 是 \(d\) 维实内积空间，\(b_1,\ldots,b_d\) 是其一组实基，则
\[
	\Lambda=\mathbb Zb_1+\cdots+\mathbb Zb_d
	=\left\{\sum_{i=1}^d m_i b_i:m_i\in\mathbb Z\right\}
\]
是 \(E\) 中的一个格，\(b_1,\ldots,b_d\) 称为它的\term[geji]{格基}[lattice basis]。每个格点在这组基下有唯一的整数坐标。以下只讨论张成整个环境空间的格；若格只张成一个子空间，就把该子空间作为环境，并采用它的诱导内积。

\begin{proposition}[格基变换]
	\label{g1:prop:lattice-basis-change}
	设 \(b=(b_1,\ldots,b_d)\) 是格 \(\Lambda\subset E\) 的一组格基。另一组实基 \(b'=bU\) 是同一格的格基，当且仅当 \(U\) 是整数矩阵且 \(\det U=\pm1\)。此外，\(E\) 的每个有界子集只包含有限个格点。
\end{proposition}
\begin{proof}
	若 \(b'\) 也是 \(\Lambda\) 的格基，则每个 \(b'_j\) 在 \(b\) 下有整数坐标，故 \(U\) 的元素为整数。同样，\(b=b'V\) 对某个整数矩阵 \(V\) 成立，于是 \(UV=I\)，\(V=U^{-1}\)。两个行列式都是整数且乘积为 \(1\)，所以 \(\det U=\pm1\)。反之，若 \(U\) 为整数矩阵且行列式为 \(\pm1\)，伴随矩阵公式说明 \(U^{-1}\) 也是整数矩阵。因此 \(U\mathbb Z^d=\mathbb Z^d\)，两组基生成同一格。

	记线性同构 \(B:\mathbb R^d\rightarrow E\) 为 \(B(t)=\sum_i t_i b_i\)。若 \(K\subset E\) 有界，则 \(B^{-1}(K)\) 有界，因而其各坐标都落在某个有界实区间中。每个区间只含有限个整数，故 \(B^{-1}(K)\cap\mathbb Z^d\)，也就是 \(K\cap\Lambda\)，是有限集。
\end{proof}

行列式为 \(\pm1\) 的整数矩阵称为\term[yaomojuzhen]{幺模矩阵}[unimodular matrix]，它们构成 \(\operatorname{GL}(d,\mathbb Z)\)。例如 \((1,0),(1,1)\) 是 \(\mathbb Z^2\) 的一组格基，而 \((2,0),(0,1)\) 只生成第一坐标为偶数的子格。两组向量都是实基，但后一组漏掉了 \((1,0)\)，所以“线性无关且属于格”仍不足以保证它们是格基。

\begin{definition}[基本平行多面体与余体积]
	\label{g1:def:lattice-covolume}
	设 \(b_1,\ldots,b_d\) 是 \(E\) 中格 \(\Lambda\) 的一组格基。相应的半开\term[jibenpingxingduomianti]{基本平行多面体}[fundamental parallelepiped]为
	\[
		F_b=\left\{\sum_{i=1}^d t_i b_i:0\le t_i<1\right\}.
	\]
	它的 \(d\) 维欧氏体积称为格的\term[yutiji]{余体积}[covolume]，记为 \(\operatorname{covol}(\Lambda)\)。对可测集合 \(P\subset E\)，记
	\[
		\operatorname{Vol}_{\Lambda}(P)
		=\frac{\operatorname{Vol}_d(P)}{\operatorname{covol}(\Lambda)},
	\]
	称为相对于格的体积；于是一个基本平行多面体的这种体积为 \(1\)。
\end{definition}

\symidx{\(\operatorname{covol}(\Lambda)\)：格的余体积}
\symidx{\(\operatorname{Vol}_{\Lambda}(P)\)：相对于格的体积}
这一定义中的余体积不应依赖所选格基，否则面积与格点数的比较也会随记法变化。半开条件则承担另一个任务：每个点只归入一份基本区域，不在相邻两份区域的边界上重复计数。

\begin{proposition}[平铺与余体积公式]
	\label{g1:prop:lattice-fundamental-domain}
	对 \(d\) 维实内积空间中的格 \(\Lambda\) 和任一格基 \(b\)，各集合 \(F_b+\lambda\)，\(\lambda\in\Lambda\)，两两不交且并成 \(E\)。每个加法商类 \(x+\Lambda\) 在 \(F_b\) 中恰有一个代表。余体积满足
	\[
		\operatorname{covol}(\Lambda)
		=\sqrt{\det\bigl(\langle b_i,b_j\rangle\bigr)_{i,j=1}^d},
	\]
	并且与格基无关。若 \(b_i\) 在一组标准正交基下组成列矩阵 \(B\)，则余体积也等于 \(|\det B|\)。
\end{proposition}
\begin{proof}
	对任意 \(x\in E\)，唯一写成 \(x=\sum_i s_i b_i\)，再把每个实数分成整数与小数部分：
	\[
		s_i=\lfloor s_i\rfloor+t_i,\qquad 0\le t_i<1.
	\]
	于是 \(x=\lambda+f\)，其中 \(\lambda=\sum_i\lfloor s_i\rfloor b_i\in\Lambda\)，\(f=\sum_i t_i b_i\in F_b\)。若还有 \(x=\lambda'+f'\)，则 \(f-f'\in\Lambda\)，但其各坐标都在 \((-1,1)\) 中且为整数，故只能全为零。因此 \(f=f'\)、\(\lambda=\lambda'\)，既证明了不交平铺，也证明了商类代表的唯一性。

	半开平行多面体与其闭包只相差有限个低维面，不改变 \(d\) 维体积。\cref{g1:prop:gram-volume}给出闭包体积的平方为 Gram 行列式，所以得到所列公式；在标准正交坐标下，Gram 矩阵为 \(B^{\mathsf T}B\)，其行列式为 \((\det B)^2\)。若换格基为 \(B'=BU\)，\cref{g1:prop:lattice-basis-change}给出 \(|\det U|=1\)，从而 \(|\det B'|=|\det B|\)，证明余体积与换基无关。
\end{proof}

选择格基后，\(B^{-1}\) 把 \(\Lambda\) 变为 \(\mathbb Z^d\)，也把相对于格的体积变为通常坐标体积。它一般不保持长度和角度，却同时保持我们即将比较的格点数与相对体积。例如 \(\Lambda=\mathbb Z(2,0)+\mathbb Z(1,1)\) 的余体积为 \(2\)；同样欧氏面积的区域，相对于这个格所占的基本面积份数是相对于 \(\mathbb Z^2\) 的一半。

\ProseParagraphBreak

本节沿用第3章的格多胞形概念，即有限个格点的凸包，也称格多面体。由\cref{g1:thm:polytope-hv}，它是有界多面体，全部顶点都是格点。以下的格多边形特指平面中内部非空的凸格多胞形；顶点按边界次序排列，边界上一条边中间的格点不另外称为多边形的顶点。后面构造三角剖分时，这些点才会成为剖分的顶点。

\subsection{二维 Pick 公式及三角剖分证明}
\label{g1:subsec:B01:02}

先在标准格 \(\mathbb Z^2\) 中工作。对格多边形 \(P\)，记内部格点数为 \(I(P)\)，边界格点数为 \(B(P)\)。面积与这两个整数的联系来自一个二维特征：没有其他格点的格三角形，其面积被迫为 \(1/2\)。证明这个事实之后，才可以用这样的三角形拼回多边形。

\begin{lemma}[不含额外格点的三角形]
	\label{g1:lem:empty-lattice-triangle}
	设 \(a,b,c\in\mathbb Z^2\) 不共线，\(T=\operatorname{conv}\{a,b,c\}\)。以下条件等价：
	\begin{enumerate}
		\item \(T\cap\mathbb Z^2=\{a,b,c\}\)，即边上及内部都没有其他格点；
		\item \(b-a,c-a\) 是 \(\mathbb Z^2\) 的一组格基；
		\item \(\operatorname{Area}(T)=1/2\)。
	\end{enumerate}
\end{lemma}
\begin{proof}
	平移后可设 \(a=0\)，记 \(u=b-a\)、\(v=c-a\)，令 \(A\) 为以 \(u,v\) 为列的整数矩阵。其行列式为非零整数，而三角形面积为 \(|\det A|/2\)。因此第二项与第三项由\cref{g1:prop:lattice-basis-change}等价。

	若 \(|\det A|=1\)，则 \(A^{-1}\) 是整数矩阵。任取 \(x\in T\cap\mathbb Z^2\)，其坐标 \(A^{-1}x=(s,t)\) 是整数，又满足 \(s,t\ge0\)、\(s+t\le1\)，只能是 \((0,0),(1,0),(0,1)\)。故 \(T\) 没有其他格点。

	反过来，若 \(|\det A|>1\)，则 \(A^{-1}\) 不可能是整数矩阵，否则其行列式 \(1/\det A\) 应为整数。于是某个标准基向量 \(e\) 的坐标 \(A^{-1}e\) 不是整数向量。逐坐标取小数部分，得到 \(s,t\in[0,1)\)，不同时为零，并且
	\[
		w=su+tv=e-A\lfloor A^{-1}e\rfloor\in\mathbb Z^2.
	\]
	若 \(s+t\le1\)，则 \(w\in T\)，并且 \(w\) 不是三个顶点中的任何一个。若 \(s+t>1\)，改取
	\[
		w'=u+v-w=(1-s)u+(1-t)v.
	\]
	此时两个系数都严格为正，和为 \(2-s-t<1\)，所以 \(w'\) 是 \(T\) 的一个内部格点。两种情形都与第一项矛盾，故第一项也迫使 \(|\det A|=1\)。
\end{proof}

第二项解释了面积 \(1/2\) 的来源：这样的三角形正好是一个基本平行四边形的一半。条件中“边上也没有其他格点”不能删去，例如顶点为 \((0,0),(m,0),(0,1)\) 的三角形没有内部格点，面积却可以随正整数 \(m\) 任意增大。

\begin{lemma}[包含全部格点的三角剖分]
	\label{g1:lem:lattice-polygon-triangulation}
	每个内部非空的凸格多边形 \(P\subset\mathbb R^2\) 都有一个有限直线三角剖分，其顶点集恰好是 \(P\cap\mathbb Z^2\)，每个三角形的面积为 \(1/2\)。若剖分的顶点、边、三角形数为 \(V,E,F\)，则 \(V-E+F=1\)。
\end{lemma}
\begin{proof}
	设 \(P\) 的顶点按边界次序为 \(p_1,\ldots,p_r\)，其中 \(r\ge3\)。由凸性，从 \(p_1\) 向 \(p_3,\ldots,p_{r-1}\) 作对角线，得到三角形
	\[
		\operatorname{conv}\{p_1,p_i,p_{i+1}\},\qquad 2\le i\le r-1.
	\]
	这些三角形覆盖 \(P\)，内部分离，且相交时只沿公共边或顶点相交：从 \(p_1\) 指向多边形中一点的射线，离开 \(P\) 时落在某条边 \([p_i,p_{i+1}]\) 上，给出覆盖；相邻顶点方向之间的扇形内部互不重叠，给出交集性质。因此这是一个直线三角剖分，初始计数为 \(V=r\)、\(E=2r-3\)、\(F=r-2\)，交替和等于 \(1\)。

	由\cref{g1:prop:lattice-basis-change}，\(P\cap\mathbb Z^2\) 有限。依次插入尚未成为剖分顶点的格点 \(q\)。若 \(q\) 位于某个三角形内部，就把它与该三角形的三个顶点相连，把一面分成三面；此时 \((V,E,F)\) 的增量为 \((1,3,2)\)。若 \(q\) 位于边界上一条边的内部，就把该边分成两段，并将 \(q\) 与相邻三角形的对顶点相连，增量为 \((1,2,1)\)。若 \(q\) 位于内部边的内部，就同时对这条边两侧的三角形作这个操作，增量为 \((1,3,2)\)。每次新线段都留在原三角形内，两侧对共有边的划分一致，所以仍是一个直线三角剖分。

	这些操作都保持 \(V-E+F\)。有限次后，全部格点都已成为剖分顶点；若还有一个小三角形在边上或内部含有非顶点格点，该点就应是剖分的另一顶点，违反单形相交只能沿公共面的要求。因此各小三角形都满足\cref{g1:lem:empty-lattice-triangle}，面积为 \(1/2\)，而交替和仍为 \(1\)。
\end{proof}

这里直接追踪了构造过程中的计数，也与\cref{g1:prop:euler-triangulation-independent}对闭圆盘给出的 Euler 数 \(1\) 一致。直接构造保留了直线边和格点顶点；一般拓扑三角剖分的存在性本身并不保证这两个额外条件。下面的\term[Pickgongshi]{Pick 公式}[Pick's theorem]由这些小三角形的面积和边界计数得到。

把这些三角形的面积与剖分中的顶点、边、面计数相比较，得到\term[Pickgongshi]{Pick 公式}[Pick's theorem]。\footnote{Pick（皮克，Georg Alexander Pick，1859--1942），奥地利数学家，研究函数论与几何，以格多边形的面积公式著称。}

\begin{theorem}[Pick 公式]
	\label{g1:thm:pick-formula}
	设 \(P\subset\mathbb R^2\) 是内部非空的凸格多边形，\(I(P)\) 与 \(B(P)\) 分别是它的内部、边界格点数，则
	\begin{equation}
		\label{g1:eq:pick-formula}
		\operatorname{Area}(P)=I(P)+\frac{B(P)}2-1.
	\end{equation}
	更一般地，若顶点属于平面格 \(\Lambda\)，则同一公式成立，只须将左侧改为 \(\operatorname{Area}(P)/\operatorname{covol}(\Lambda)\)，右侧按 \(\Lambda\) 计算格点。
\end{theorem}
\begin{proof}
	先用\cref{g1:lem:lattice-polygon-triangulation}取得包含全部格点的剖分，简记 \(I=I(P)\)、\(B=B(P)\)。其顶点数为 \(V=I+B\)。边界上的格点按一圈排列，相邻两点间恰有一条剖分边，所以边界边数也是 \(B\)。记内部边数为 \(E_{\mathrm{int}}\)，则
	\[
		E=E_{\mathrm{int}}+B,\qquad
		3F=2E_{\mathrm{int}}+B=2E-B.
	\]
	第二式逐三角形数边：内部边属于两面，边界边只属于一面。再用 \(V-E+F=1\)，消去 \(V,E\)，得到
	\[
		F=2I+B-2.
	\]
	每个小三角形面积为 \(1/2\)，有限面积相加给出 \(\operatorname{Area}(P)=F/2\)，即所需公式。对任意平面格，选择格基并施以相应的逆线性坐标变换；格点与内外部关系不变，面积则除以 \(\operatorname{covol}(\Lambda)\)，归结到已经证明的标准格情形。
\end{proof}

\begin{figure}[htbp]
	\centering
	% !TeX root = ../Book1.tex
\begin{tikzpicture}[>=stealth,font=\small,x=.95cm,y=.95cm]
	\begin{scope}
		\fill[blue!10] (0,0)--(2,0)--(3,1)--(1,1)--cycle;
		\foreach \p in {(0,0),(2,0),(1,1),(3,1),(0,2),(2,2)} {
			\fill[gray] \p circle (1.2pt);
		}
		\draw[blue!70!black,thick] (2,0)--(0,0)--(1,1);
		\draw[blue!70!black,thick,dashed] (2,0)--(3,1)--(1,1);
		\draw[->] (0,0)--(2,0) node[midway,below] {\(b_1\)};
		\draw[->] (0,0)--(1,1) node[midway,left] {\(b_2\)};
		\foreach \p in {(2,0),(3,1),(1,1)} {
			\draw[blue!70!black,fill=white] \p circle (1.8pt);
		}
		\fill[blue!70!black] (0,0) circle (1.8pt);
		\node at (1.5,1.55) {\(F_b\)};
		\node at (1.5,-.75) {余体积 \(2\)};
	\end{scope}
	\begin{scope}[xshift=5.5cm]
		\foreach \x in {0,1,2,3} {
			\foreach \y in {0,1,2} {
				\fill[gray!60] (\x,\y) circle (1pt);
			}
		}
		\fill[blue!10] (0,0)--(3,0)--(0,2)--cycle;
		\draw[thick] (0,0)--(3,0)--(0,2)--cycle;
		\foreach \p in {(0,0),(1,0),(2,0),(3,0),(0,2),(0,1)} {
			\draw[blue!70!black] (1,1)--\p;
			\fill \p circle (1.6pt);
		}
		\fill[red!70!black] (1,1) circle (2pt);
		\node[above right] at (1,1) {\((1,1)\)};
		\node[below left] at (0,0) {\((0,0)\)};
		\node[below right] at (3,0) {\((3,0)\)};
		\node[above] at (0,2) {\((0,2)\)};
		\node at (1.5,-.75) {\(I=1,\ B=6,\ F=6\)};
	\end{scope}
\end{tikzpicture}

	\caption[基本区域与格三角剖分]{左图为 \(\Lambda=\mathbb Z(2,0)+\mathbb Z(1,1)\) 的基本平行四边形，虚线边及空心顶点不属于半开区域，余体积为 \(2\)。右图在标准格中，把顶点为 \((0,0),(3,0),(0,2)\) 的三角形分成六个面积 \(1/2\) 的小三角形；内部有一个格点，边界有六个格点，所以面积为 \(1+6/2-1=3\)。}
	\label{g1:fig:lattice-fundamental-pick}
\end{figure}

公式把面积恢复为整数与半整数的组合，也说明格多边形的面积必为半整数。\cref{g1:fig:lattice-fundamental-pick}右侧三角形是一个具体核验。上述证明以凸性取得第一份扇形三角剖分；对非凸简单格多边形，要先另证存在保留原顶点的直线三角剖分，才能继续作同样的格点插入。这里的定理范围保持为凸格多边形，不借一般曲面的拓扑剖分省略这一几何步骤。

\subsection{Ehrhart 多项式的代表性陈述与例子}
\label{g1:subsec:B01:03}

面积随整数倍放大按二次律增长，边界上的格点则按一次律增长。Pick 公式因而预示：凸格多边形的全部格点数应由一个二次多项式给出。先把边界计数中需要的整除条件说清楚。

\begin{lemma}[格线段上的点数]
	\label{g1:lem:lattice-segment-count}
	设 \(a,b\in\mathbb Z^2\) 不同，\(b-a=(u,v)\)，记 \(g=\gcd(|u|,|v|)>0\)，其中 \(\gcd(k,0)=k\)。则闭线段 \([a,b]\) 上的格点恰为
	\[
		a+\frac{k}{g}(b-a),\qquad k=0,1,\ldots,g.
	\]
	因此它有 \(g+1\) 个格点；若凸格多边形各边向量为 \((u_i,v_i)\)，则
	\[
		B(P)=\sum_i\gcd(|u_i|,|v_i|),\qquad
		B(nP)=nB(P)\quad(n\ge1).
	\]
\end{lemma}
\begin{proof}
	写 \((u,v)=g(u',v')\)，则 \(u',v'\) 互素。若 \(a+t(b-a)\) 为格点，\(0\le t\le1\)，则 \(tgu'\) 与 \(tgv'\) 都为整数。取整数 \(r,s\) 使 \(ru'+sv'=1\)，可得 \(tg=r(tgu')+s(tgv')\in\mathbb Z\)。所以 \(t=k/g\)，且 \(0\le k\le g\)。反过来，这些参数确实给出格点。

	把各边的点数 \(g_i+1\) 相加时，每个多边形顶点重复计数一次，而顶点数等于边数，所以 \(B(P)=\sum_i g_i\)。放大 \(n\) 倍把每个边向量乘以 \(n\)，其两个坐标的最大公因数也乘以 \(n\)，因而得到最后一式。
\end{proof}

对非空格多胞形 \(P\subset\mathbb R^d\)，定义其整数倍放大的计数函数
\[
	L_P(n)=\#(nP\cap\mathbb Z^d),\qquad n\in\mathbb Z_{\ge0},
\]
其中 \(nP=\{nx:x\in P\}\)，特别地 \(0P=\{0\}\)，所以 \(L_P(0)=1\)。当这些值由一个多项式给出时，称该多项式为 \(P\) 的\term[Ehrhartduoxiangshi]{Ehrhart 多项式}[Ehrhart polynomial]。\footnote{Ehrhart（埃尔哈特，Eugène Ehrhart，1906--2000），法国数学家，研究多面体中的整数点计数，建立了格多胞形膨胀计数的多项式理论。}下面分别证明二维凸格多边形、整数箱体和任意维格单形的多项式公式。

\symidx{\(L_P(n)\)：整数倍多胞形的格点计数}
\begin{proposition}[二维 Ehrhart 公式]
	\label{g1:prop:polygon-ehrhart}
	设 \(P\subset\mathbb R^2\) 是内部非空的凸格多边形，面积为 \(A\)，边界格点数为 \(B\)。则对每个整数 \(n\ge0\)，
	\begin{equation}
		\label{g1:eq:polygon-ehrhart}
		L_P(n)=An^2+\frac B2n+1.
	\end{equation}
	对 \(n\ge1\)，内部格点数为
	\[
		\#\bigl(\operatorname{int}(nP)\cap\mathbb Z^2\bigr)
		=An^2-\frac B2n+1=L_P(-n),
	\]
	其中 \(L_P(-n)\) 表示把负数代入已得到的多项式。
\end{proposition}
\begin{proof}
	当 \(n\ge1\) 时，\(nP\) 仍为凸格多边形，面积是 \(An^2\)，边界格点数由\cref{g1:lem:lattice-segment-count}等于 \(nB\)。将这些数代入 Pick 公式，先得内部点数 \(An^2-nB/2+1\)，再加上边界点数即得 \(L_P(n)=An^2+nB/2+1\)。在 \(n=0\) 时，该多项式的值为 \(1\)，与单点 \(0P\) 一致。最后直接代入 \(-n\)，得到所写的内部点数公式。
\end{proof}

对\cref{g1:fig:lattice-fundamental-pick}右图的三角形，\(A=3\)、\(B=6\)，所以 \(L_P(n)=3n^2+3n+1\)，在 \(n=1,2\) 时分别给出 \(7,19\) 个格点。最高次系数记录面积，一次系数记录边界的格长度信息，常数项反映零倍收缩后的一个点。负参数只在已经确定的多项式中求值；不能把 \(L_P(-n)\) 在这里理解为反射后区域 \((-n)P\) 的全部格点数。

\begin{proposition}[整数箱体与标准单形]
	\label{g1:prop:box-simplex-ehrhart}
	设 \(a_1,\ldots,a_d\) 为正整数，\(Q=\prod_{i=1}^d[0,a_i]\)。再令
	\[
		\Delta_d=\operatorname{conv}\{0,e_1,\ldots,e_d\}
		=\left\{x\in\mathbb R^d:x_i\ge0,\ \sum_{i=1}^d x_i\le1\right\}.
	\]
	则对每个整数 \(n\ge0\)，
	\[
		L_Q(n)=\prod_{i=1}^d(a_i n+1),\qquad
		L_{\Delta_d}(n)=\binom{n+d}{d}
		=\frac{(n+1)\cdots(n+d)}{d!}.
	\]
\end{proposition}
\begin{proof}
	箱体中第 \(i\) 个坐标可以独立选取 \(0,1,\ldots,a_i n\)，所以由乘法原理得到第一式。对标准单形，格点就是非负整数解 \(m_1+\cdots+m_d\le n\)。引入余数坐标 \(m_0=n-\sum_{i=1}^d m_i\)，得到 \(d+1\) 个非负整数之和为 \(n\) 的表示。把 \(n\) 个相同记号排成一行，再放置 \(d\) 个隔板，允许隔板相邻或位于两端；相邻隔板之间的记号数依次给出 \(m_0,\ldots,m_d\)，反之也唯一决定这种排列。因此共有 \(\binom{n+d}{d}\) 种，得到第二式及其多项式表达。
\end{proof}

一般格单形不一定能经幺模变换化为标准单形。仍然可以计数，但须先把所有倍数 \(nS\) 放进同一个高一维空间：把顶点 \(v_i\) 提升为 \((v_i,1)\)，则高度为 \(n\) 的截面恰好是 \(nS\)。对提升后的基方向取整数与小数部分，可以把无限锥中的格点分成有限种余数类型。

\begin{proposition}[格单形的半开锥计数]
	\label{g1:prop:lattice-simplex-ehrhart}
	设 \(v_0,\ldots,v_d\in\mathbb Z^d\) 仿射无关，\(S=\operatorname{conv}\{v_0,\ldots,v_d\}\)，记 \(w_i=(v_i,1)\in\mathbb Z^{d+1}\)。定义有界半开区域
	\[
		\Pi=\left\{\sum_{i=0}^d t_iw_i:0\le t_i<1\right\},
	\]
	并令 \(h_j\) 为 \(\Pi\cap\mathbb Z^{d+1}\) 中末坐标等于 \(j\) 的点数。则 \(0\le j\le d\)、\(h_0=1\)，且作为形式幂级数有
	\[
		\sum_{n\ge0}L_S(n)z^n
		=\frac{h_0+h_1z+\cdots+h_dz^d}{(1-z)^{d+1}}.
	\]
	特别地，对每个整数 \(n\ge0\)，
	\[
		L_S(n)=\sum_{j=0}^d h_j\binom{n+d-j}{d},
	\]
	右侧是一个次数为 \(d\) 的多项式。此处二项式按多项式 \(\binom{t}{d}=t(t-1)\cdots(t-d+1)/d!\) 解释。
\end{proposition}
\begin{proof}
	仿射无关性意味着 \(w_0,\ldots,w_d\) 线性无关：若 \(\sum_i c_iw_i=0\)，最后一坐标给出 \(\sum_i c_i=0\)，其余坐标再给出 \(\sum_i c_iv_i=0\)，故所有 \(c_i=0\)。因此锥
	\[
		C=\left\{\sum_{i=0}^d s_iw_i:s_i\ge0\right\}
	\]
	中每个点都有唯一的非负系数表示。在高度 \(n>0\) 时，\(\sum_i s_i=n\)，其前 \(d\) 个坐标为 \(n\sum_i(s_i/n)v_i\)，正好遍历 \(nS\)。高度 \(0\) 时只剩原点，也与 \(0S\) 一致。

	对锥中任一格点 \(x=\sum_i s_iw_i\)，令 \(k_i=\lfloor s_i\rfloor\)、\(t_i=s_i-k_i\)，则
	\[
		x=p+\sum_{i=0}^d k_iw_i,
		\qquad p=\sum_{i=0}^d t_iw_i\in\Pi\cap\mathbb Z^{d+1}.
	\]
	由于 \(w_i\) 是整数向量，\(p\) 确为格点；由系数表示及整数小数分解的唯一性，这种分解也唯一。反之，任一 \(p\in\Pi\cap\mathbb Z^{d+1}\) 加上非负整数倍 \(w_i\)，都给出锥中格点。

	区域 \(\Pi\) 有界，因而它的格点数有限。\(p\) 的高度是 \(\sum_i t_i\)，为 \([0,d+1)\) 中的整数，故只可能是 \(0,1,\ldots,d\)；高度为零迫使全部 \(t_i=0\)，所以 \(h_0=1\)。每加一个 \(w_i\)，高度增加 \(1\)，于是上述唯一分解给出形式生成函数
	\[
		\sum_{x\in C\cap\mathbb Z^{d+1}}z^{x_{d+1}}
		=\left(\sum_{j=0}^d h_jz^j\right)
		\prod_{i=0}^d\left(\sum_{k_i\ge0}z^{k_i}\right).
	\]
	各次系数只涉及有限项，因此该等式不需要解析收敛；代入几何级数的形式恒等式便得所述分式。高度为 \(n\) 的项，对每个高度为 \(j\) 的余数点，计数 \(k_0+\cdots+k_d=n-j\) 的非负整数解。由上一命题的隔板计数，在 \(n\ge j\) 时有 \(\binom{n-j+d}{d}\) 个；在 \(n<j\le d\) 时没有解，而多项式中的上参数 \(n+d-j\) 落在 \(0,\ldots,d-1\)，其值也为零。因此多项式公式对全部 \(n\ge0\) 成立。最高次系数为 \((h_0+\cdots+h_d)/d!>0\)，所以次数恰为 \(d\)。
\end{proof}

标准单形的提升向量构成 \(\mathbb Z^{d+1}\) 的格基，所以其 \(\Pi\) 中只有格点 \(0\)，上述公式退化为 \(\binom{n+d}{d}\)。对其他格单形，有限个余数点保存了标准单形中没有的整数结构。一般格多胞形可以涉及多个单形及其公共面，不能把闭单形的计数简单相加；本节已经证明的公式分别适用于上述三类对象，并未用这些例子替代任意维格多胞形的统一定理。

\section{Voronoi 与 Delaunay}
\label{g1:sec:B02}

在平面上给定几个服务点，每个位置应当归属于离它最近的哪一个点？若两个服务点同样近，分界线在哪里？反过来，哪些服务点应当连起来，才能记录这些区域怎样相邻？这两个问题分别产生 Voronoi 分区和 Delaunay 分解。\footnote{Voronoi（沃罗诺伊，\OriginalName{Георгій Феодосійович Вороний}，1868--1908），乌克兰数学家，研究数论与二次型。Delaunay（德劳内，\OriginalName{Борис Николаевич Делоне}，1890--1980），俄罗斯数学家，研究数的几何与离散几何；其姓氏也常转写为 Delone。}距离本来含有平方根，但比较两个距离的平方时，关于位置的二次项恰好消去；再把服务点提升到一个抛物面上，空圆或空球条件也会成为支撑超平面的条件。本节使用\cref{g1:thm:polytope-hv,g1:thm:polytope-faces}中已经证明的有限半空间表示与面结构，不要求预先掌握计算几何算法。

\subsection{距离平方与分区}
\label{g1:subsec:B02:01}

设 \(S\subset\mathbb R^d\) 为非空有限互异点集，\(d\ge1\)，使用标准欧氏内积和距离。把 \(S\) 中的点称为站点。固定站点 \(p\)，位置 \(x\) 离 \(p\) 不比离另一站点 \(q\) 远，等价于
\begin{equation}\label{g1:eq:voronoi-halfspace}
	\lVert x-p\rVert^2\le\lVert x-q\rVert^2
	\quad\Longleftrightarrow\quad
	2\langle x,q-p\rangle\le\lVert q\rVert^2-\lVert p\rVert^2.
\end{equation}
当 \(p\ne q\) 时，右侧是一个闭半空间，其边界是线段 \(pq\) 的垂直平分超平面。共同的 \(\lVert x\rVert^2\) 消去了，所以分界是平的。

\begin{definition}[最近站点与 Voronoi 胞]\label{g1:def:voronoi-cell}
	对非空有限互异点集 \(S\subset\mathbb R^d\)，定义位置 \(x\) 的最近站点集为
	\[
		N(x)=\{p\in S:\lVert x-p\rVert=\min_{q\in S}\lVert x-q\rVert\}.
	\]
	站点 \(p\in S\) 的\term[voronoibao]{Voronoi 胞}[Voronoi cell]为
	\[
		V_p=\{x\in\mathbb R^d:p\in N(x)\}
		=\bigcap_{q\in S\setminus\{p\}}
		\{x:2\langle x,q-p\rangle\le\lVert q\rVert^2-\lVert p\rVert^2\}.
	\]
	这些胞连同它们的非空面组成\term[voronoitu]{Voronoi 图}[Voronoi diagram]。这里的“图”包括各维多面体，并非只有一张边与顶点组成的图。
\end{definition}

\begin{proposition}[Voronoi 胞的覆盖与交界]\label{g1:prop:voronoi-cover}
	设 \(S\subset\mathbb R^d\) 为非空有限互异点集。每个 \(V_p\) 都是含 \(p\) 为内点的闭凸多面体，而且
	\[
		\bigcup_{p\in S}V_p=\mathbb R^d,
		\qquad
		\operatorname{int}V_p
		=\{x:\lVert x-p\rVert<\lVert x-q\rVert\text{ 对全部 }q\ne p\}.
	\]
	不同胞的内部不相交；它们可以共享边界。对非空 \(T\subseteq S\)，若 \(V_T=\bigcap_{p\in T}V_p\ne\varnothing\)，则 \(V_T\) 是每个 \(V_p\)、\(p\in T\)，的面。所有非空 Voronoi 面都可写成这种交集。
\end{proposition}
\begin{proof}
	半空间公式给出闭性、凸性和有限半空间表示，正是\cref{g1:def:polyhedron-polytope}中的多面体。在 \(x=p\) 处，对每个 \(q\ne p\) 都有严格不等式；由有限性，它们在 \(p\) 的一个小球内同时严格成立。若 \(S=\{p\}\)，则 \(V_p=\mathbb R^d\)，结论也成立。每个位置到有限站点集的最短距离都能取到，故这些胞覆盖全空间。

	全部距离比较严格时，仍由有限性得到包含于 \(V_p\) 的邻域。若某个比较取等号，则 \(x\) 位于一个真正的限制超平面上；向被该不等式排除的一侧作任意小移动就会离开 \(V_p\)，所以 \(x\) 不是内点。这证明内部公式，也证明不同内部不交。

	固定 \(p_0\in T\)。在 \(V_{p_0}\) 内，属于 \(V_T\) 恰好要求与 \(T\) 中各站点的比较都取等号。每条这样的等式给出支撑面，交集仍满足\cref{g1:def:polytope-face}中的线段性质；同一定义和论证也适用于无界多面体。

	最后说明这些交集没有遗漏别的面。设 \(F\) 为 \(V_{p_0}\) 的非空面，取 \(x\in\operatorname{ri}F\)，令 \(G\) 为 \(V_{p_0}\) 中满足 \(x\) 处全部活动等式的集合。每条活动等式都在整个 \(F\) 上成立：否则由 \(x\) 的相对内部性，沿 \(F\) 内一条线段向相反方向稍作延长，就会违反该不等式。因此 \(F\subseteq G\)。对 \(y\in G\)，点 \(z=x+\varepsilon(x-y)\) 保持全部活动等式；其他有限条不等式有严格余量，所以足够小的 \(\varepsilon>0\) 使 \(z\in V_{p_0}\)。把 \(x\) 写成 \(y,z\) 的严格凸组合，面的线段性质给出 \(y\in F\)。于是 \(F=G=V_{N(x)}\)。这也是\cref{g1:thm:polytope-faces}中活动等式论证在此处的具体使用。
\end{proof}

闭胞覆盖并不表示每个位置只属于一个闭胞。在两站点等距且最近的位置，两胞都包含该点；在多个站点同时最近的位置，共享的胞更多。只有删去这些边界、只看各胞的内部时，归属才唯一。Voronoi 胞也不必有界：有限站点不能从所有方向围住每一个站点。

\begin{example}[三个站点的完整分区]\label{g1:exm:three-site-voronoi}
	取 \(a=(0,0)\)、\(b=(2,0)\)、\(c=(0,2)\)。逐对消去距离平方中的二次项，得到
	\[
		\begin{aligned}
			V_a&=\{(x,y):x\le1,\ y\le1\},\\
			V_b&=\{(x,y):x\ge1,\ y\le x\},\\
			V_c&=\{(x,y):y\ge1,\ x\le y\}.
		\end{aligned}
	\]
	两两交界分别为 \(x=1,y\le1\)、\(y=1,x\le1\) 和 \(x=y\ge1\)，三者在 \((1,1)\) 相遇。这个公共点到三站点都距离 \(\sqrt2\)。三条完整垂直平分线并不是 Voronoi 图的全部边界：每条线上只有仍然不比第三站点远的部分被保留。
\end{example}

\subsection{空球条件和对偶分解}
\label{g1:subsec:B02:02}

若位置 \(c\) 同时以若干站点为最近站点，则以 \(c\) 为中心、最短距离为半径的开球不含任何站点，而全部最近站点都在球面上。这给出从 Voronoi 面寻找站点组合的办法。不过，把每一组拥有共同最近位置的站点随意连成单纯形，尚不能保证这些图形按公共面拼接。为证明真正的分解性质，需要把整个构造放在同一个凸包上。

\begin{definition}[抛物面提升与下部面]\label{g1:def:delaunay-lift}
	设 \(S\subset\mathbb R^d\) 为非空有限互异点集。令
	\[
		\widehat p=(p,\lVert p\rVert^2),\qquad
		Q=\operatorname{conv}\{\widehat p:p\in S\}\subset\mathbb R^{d+1},
		\qquad \pi(y,z)=y.
	\]
	对 \(c\in\mathbb R^d\)，令
	\[
		b(c)=\min_{p\in S}(\lVert p\rVert^2-2\langle c,p\rangle),
		\qquad
		F(c)=\{(y,z)\in Q:z-2\langle c,y\rangle=b(c)\}.
	\]
	称这样的非空暴露面为 \(Q\) 的下部面，其投影
	\[
		D(c)=\pi F(c)=\operatorname{conv}N(c)
	\]
	称为一个\term[delaunaydanyuan]{Delaunay 单元}[Delaunay cell]。不同中心可能给出同一个单元，收集单元时只保留不同的集合。
\end{definition}

这里的下部面由法向量 \((2c,-1)\) 暴露，最后一坐标严格为负；竖直支撑超平面的最后一法向坐标为零，不在这个定义中。等式 \(D(c)=\operatorname{conv}N(c)\) 来自
\[
	\lVert p\rVert^2-2\langle c,p\rangle
	=\lVert p-c\rVert^2-\lVert c\rVert^2.
\]
线性函数在有限凸包上的最小值由生成点决定，达到最小值的点组成对应的面，因而正好选出了最近站点。

\begin{proposition}[空球与支撑超平面]\label{g1:prop:empty-sphere-lift}
	设 \(S\subset\mathbb R^d\) 为非空有限互异点集，\(c\in\mathbb R^d\)、\(r\ge0\)。开球 \(\{x:\lVert x-c\rVert<r\}\) 不含 \(S\) 中的点且球面包含至少一个站点，当且仅当
	\[
		z=2\langle c,y\rangle+r^2-\lVert c\rVert^2
	\]
	从下方支撑提升凸包 \(Q\)。接触面恰由球面上的全部站点的提升生成，其投影为一个 Delaunay 单元。这个单元的顶点恰是球面上的站点，且单元中不含其他站点。
\end{proposition}
\begin{proof}
	对任意 \(p\in S\)，有
	\[
		\lVert p\rVert^2-\bigl(2\langle c,p\rangle+r^2-\lVert c\rVert^2\bigr)
		=\lVert p-c\rVert^2-r^2.
	\]
	右侧全部非负且至少一次为零，恰好等价于所给平面在每个生成点下方并与凸包接触。这些非负差在凸组合上仍非负；一个凸组合的差为零，当且仅当其中正系数项的差都为零。因此接触面正由球面站点生成。

	设这些站点组成 \(T\)。对 \(x=\sum_{p\in T}\lambda_pp\)、\(\lambda_p\ge0\)、\(\sum_p\lambda_p=1\)，有
	\[
		r^2-\lVert x-c\rVert^2
		=\sum_{p\in T}\lambda_p\lVert p-x\rVert^2.
	\]
	如果组合实际使用至少两个不同点，右侧严格为正，故 \(x\) 在开球内，不可能是站点。特别地，\(T\) 中任何点都不能由其余点非平凡地凸组合得到，所以它们全是顶点；\(\operatorname{conv}T\) 中的站点也只能是 \(T\)。半径为零时 \(T\) 只有一个点，结论仍成立。
\end{proof}

空球在这里允许球面上有额外站点。“开球中没有站点”与“球面上恰有 \(d+1\) 个站点”是两件事；后一条件决定所得单元是否一定是单纯形。

\begin{theorem}[Delaunay 多胞形剖分]\label{g1:thm:delaunay-subdivision}
	设 \(S\subset\mathbb R^d\) 为有限互异点集，\(d\ge1\)，且 \(\operatorname{aff}S=\mathbb R^d\)。全部不同的 Delaunay 单元组成有限个多胞形，具有以下性质：
	\begin{enumerate}
		\item 它们的并为 \(K=\operatorname{conv}S\)。
		\item 每个单元的非空面仍是一个单元；两个单元的交为空集或它们的公共面。
		\item 每个按包含关系极大的单元都为 \(d\) 维；全部顶点恰为 \(S\)。
	\end{enumerate}
	这称为 \(S\) 的\term[delaunaypoufen]{Delaunay 多胞形剖分}[Delaunay subdivision]。
\end{theorem}
\begin{proof}
	提升把问题变成同一个凸包的下表面：只要投影不把不同高度的面叠在一起，面的拼接就能保留下来。由\cref{g1:thm:polytope-hv,g1:thm:polytope-faces}，\(Q\) 有有限半空间表示、有限多个面。每个下部面位于一个图形超平面 \(z=2\langle c,y\rangle+b(c)\) 中，故 \(\pi\) 在该面上是仿射双射；投影不会降低其维数。

	先证明覆盖。由 \(\pi Q=K\)，对任意 \(y\in K\)，竖直纤维中存在最低点 \((y,h(y))\in Q\)，因为 \(Q\) 紧。满维假设使 \(\dim Q\) 只能是 \(d\) 或 \(d+1\)。若 \(\dim Q=d\)，则 \(\pi:\operatorname{aff}Q\to\mathbb R^d\) 是仿射同构，因此 \(\operatorname{aff}Q\) 是某个 \(z=a\cdot y+b\) 的图形。取 \(c=a/2\) 就有 \(F(c)=Q\)、\(D(c)=K\)，覆盖已经成立。

	若 \(\dim Q=d+1\)，给 \(Q\) 取一个没有冗余不等式的有限半空间表示。每条限制的等号集在 \(Q\) 中都是 \(d\) 维棱面：删去该限制后可找到违反它而满足其余限制的点，将此点与 \(Q\) 的内点连线，首次碰到该等号超平面的位置仍严格满足其余限制，故在该超平面内有一个相对邻域属于 \(Q\)。在最低点 \((y,h(y))\) 处，必有一条活动限制的外法向量最后一坐标为负。否则把该点稍微向下移动，最后一法向坐标非负的活动限制仍成立，而非活动的有限条限制由严格余量也仍成立，违背最低性。将这条负坐标归一化为 \(-1\)，得到一个下部棱面。因此每个 \(y\in K\) 都在一个 \(d\) 维下部面的投影内。

	再证明面相容。若 \((y,z)\in F(c)\)，整个 \(Q\) 都在它所在的图形超平面上方，所以 \(z=h(y)\)。于是若两个下部面的投影具有公共点，它们在该点上方的提升必为同一个最低点，从而
	\[
		\pi F(c)\cap\pi F(c')=\pi\bigl(F(c)\cap F(c')\bigr).
	\]
	两个面的交仍是各自的面，仿射双射保持面，故此交为空集或投影单元的公共面。

	还须证明下部面的面仍可由最后一坐标为负的法向量暴露。设 \(E\) 是 \(F(c)\) 的非空面，由\cref{g1:thm:polytope-faces}选一个线性函数 \(v\) 在 \(F(c)\) 上恰好以 \(E\) 为最大值集。原法向量为 \(u=(2c,-1)\)。在 \(Q\) 的有限顶点上，\(u\) 对 \(F(c)\) 外的顶点有严格差距；故对充分小的 \(\varepsilon>0\)，\(u+\varepsilon v\) 的最大值仍只在 \(F(c)\) 的顶点中取得，并在其中恰选出 \(E\)。其最后一坐标仍负，归一化后说明 \(E\) 是下部面。这里 \(v\) 从面的仿射包络延拓到环境空间即可，常数项不影响最大值集。因而每个单元的面也属于这份名单。

	前面的覆盖实际上由 \(d\) 维单元完成。任取一个单元的相对内点，它属于某个 \(d\) 维单元；两者的交是前者的面，并包含其相对内点，所以包含前者全体。这说明极大单元都是 \(d\) 维。最后，对任意 \(p\in S\) 取中心 \(c=p\)，有 \(N(c)=\{p\}\)，故每个站点都给出一个零维单元；\cref{g1:prop:empty-sphere-lift}又说明全部单元的顶点都在 \(S\) 中。
\end{proof}

定理特意包括全部提升点共面的情形。例如所有站点共球时，提升凸包可能只有 \(d\) 维，整个 \(K\) 本身就是一个 Delaunay 单元。若只把下部面理解成一个满 \((d+1)\) 维凸包的部分棱面，就会遗漏这种最重要的退化情形。

\begin{theorem}[Voronoi 与 Delaunay 的反序对应]\label{g1:thm:voronoi-delaunay-duality}
	设 \(S\subset\mathbb R^d\) 为有限互异点集，\(\operatorname{aff}S=\mathbb R^d\)。对 Delaunay 单元 \(\sigma\)，记它的顶点集为 \(T_\sigma=S\cap\sigma\)，并定义
	\[
		V(\sigma)=\bigcap_{p\in T_\sigma}V_p.
	\]
	这个对应是全部非空 Delaunay 单元与全部非空 Voronoi 面之间的双射，而且
	\[
		\sigma\subseteq\tau\quad\Longleftrightarrow\quad
		V(\tau)\subseteq V(\sigma),
		\qquad \dim\sigma+\dim V(\sigma)=d.
	\]
	若 \(c\in\operatorname{ri}V(\sigma)\)，则 \(N(c)=T_\sigma\)。
\end{theorem}
\begin{proof}
	先写清 \(V(\sigma)\) 的方程。取 \(p_0\in T_\sigma\)，它由
	\[
		\begin{aligned}
			2\langle c,p-p_0\rangle&=\lVert p\rVert^2-\lVert p_0\rVert^2
			&& (p\in T_\sigma),\\
			2\langle c,q-p_0\rangle&\le\lVert q\rVert^2-\lVert p_0\rVert^2
			&& (q\in S\setminus T_\sigma)
		\end{aligned}
	\]
	给出。因为 \(\sigma=D(c_0)\) 对某个 \(c_0\) 成立，\cref{g1:prop:empty-sphere-lift}保证 \(N(c_0)=T_\sigma\)，所以第二组不等式在 \(c_0\) 处全严格。若 \(\dim\sigma=k\)，第一组方程的法向量张成 \(k\) 维空间，解集是 \(d-k\) 维仿射空间；有限个严格余量使 \(V(\sigma)\) 包含 \(c_0\) 在此空间中的一个邻域。故其仿射包络正是该方程解集，维数为 \(d-k\)。第二组中任一等号都只能出现在相对边界：若在相对内点取等号，沿仿射包络稍作反向移动便会违反它，因为它并非在整个包络上恒等。这同时说明相对内点的最近站点恰为 \(T_\sigma\)。

	任取一个非空 Voronoi 面 \(W=V_T\)，选 \(c\in\operatorname{ri}W\)，令 \(U=N(c)\)。显然 \(T\subseteq U\)，所以 \(V_U\subseteq W\)。反过来，对 \(q\in U\) 及固定 \(p_0\in T\)，距离平方差在 \(W\) 上非负、在相对内点 \(c\) 为零；它是仿射函数，故必须在整个 \(W\) 上为零，否则沿 \(W\) 的相对邻域就会取得负值。因此 \(W\subseteq V_U\)，从而 \(W=V_U=V(D(c))\)。这证明满射，也说明如何从一个 Voronoi 面恢复单元。前一段证明相对内点恢复的是完整顶点集，故该恢复唯一。

	若 \(\sigma\subseteq\tau\)，则 \(T_\sigma\subseteq T_\tau\)，对更多站点要求同时最近，只会使交集缩小，得到 \(V(\tau)\subseteq V(\sigma)\)。若后一个包含成立，取 \(c\in\operatorname{ri}V(\tau)\)，则 \(T_\sigma\subseteq N(c)=T_\tau\)，取凸包得到 \(\sigma\subseteq\tau\)。
\end{proof}

例如\cref{g1:exm:three-site-voronoi}中的三角形对应一个 Voronoi 顶点，三条三角形边分别对应三条无界射线，三个站点分别对应二维胞。维数和等于二，但并不要求对应的两个图形有相同大小或都有限长。

\ProseParagraphBreak

边界的准确含义也可由不等式读出。若 \(p_0\in T_\sigma\)，从 \(c\in V(\sigma)\) 出发，整条射线 \(c+tv\)、\(t\ge0\)，仍在该面内，当且仅当
\begin{equation}\label{g1:eq:voronoi-recession}
	\langle v,q-p_0\rangle\le0\quad(q\in S),
	\qquad
	\langle v,p-p_0\rangle=0\quad(p\in T_\sigma).
\end{equation}
这就是说线性函数 \(\langle v,-\rangle\) 在整个 \(K\) 上于 \(\sigma\) 取最大值。如果 \(v\ne0\)，满维性保证其最大值面是 \(K\) 的真支撑面。因此 \(V(\sigma)\) 无界，当且仅当 \(\sigma\) 包含于 \(K\) 的一个真支撑面。此处“无界必有射线”也可以直接检查：取 \(c_j\in V(\sigma)\)、\(\lVert c_j-c\rVert\to\infty\)，对单位方向 \((c_j-c)/\lVert c_j-c\rVert\) 取收敛子列，在各线性限制中除以长度再取极限，就得到满足\eqref{g1:eq:voronoi-recession}的单位向量。反方向按该式作射线即可。Voronoi 图覆盖整个 \(\mathbb R^d\)，Delaunay 剖分只覆盖 \(K\)；凸包边界通过无界 Voronoi 面参与上述对偶，无须另添一个假定的“无穷远站点”。

\subsection{退化位置及一般位置的区别}
\label{g1:subsec:B02:03}

多胞形剖分并不一定是单纯形剖分。若一个空球的球面上恰好有 \(d+1\) 个仿射无关站点，它们生成一个 \(d\) 维单纯形；若还有更多站点同时在球面上，就必须先保留完整接触面，再决定是否另作细分。

\begin{proposition}[一般位置下的单纯形剖分]\label{g1:prop:delaunay-general-position}
	设有限互异点集 \(S\subset\mathbb R^d\) 满足 \(\operatorname{aff}S=\mathbb R^d\)。假定任意 \(d+1\) 个站点仿射无关，且没有 \(d+2\) 个站点同在一个欧氏球面上，则 Delaunay 剖分的每个极大单元都是 \(d\) 维单纯形，其全部面组成一个唯一确定的几何单纯形剖分。二维时，这两个条件就是没有三点共线、没有四点共圆。
\end{proposition}
\begin{proof}
	由\cref{g1:thm:delaunay-subdivision}，极大单元为 \(d\) 维，故至少有 \(d+1\) 个顶点。它的全部顶点同在一个空球的球面上；无 \(d+2\) 点共球的条件使顶点数至多为 \(d+1\)。于是恰有 \(d+1\) 个，且由维数可知它们仿射无关。所有其他单元都是这些单纯形的面，公共面性质已经由剖分定理证明。唯一性来自下部面的定义：给定 \(S\) 就确定 \(Q\)，并确定全部下部面，此时没有需要另外选择的非单纯形单元。
\end{proof}

这些条件是本节采用的一组明确的一般位置条件，并非每个具体结论所需的最弱假设。即使某几个站点共球，只要那个球内部还有其他站点，它也不会形成下部接触面；是否产生 Delaunay 退化，仍要核对空球条件。

\begin{example}[共圆四点与对角线的改变]\label{g1:exm:delaunay-square-perturbation}
	取
	\[
		A=(0,0),\quad B=(1,0),\quad C=(1,1),\quad D_\varepsilon=(0,1+\varepsilon),
		\qquad |\varepsilon|<\tfrac12.
	\]
	当 \(\varepsilon=0\) 时，四点同在中心 \((1/2,1/2)\)、半径 \(1/\sqrt2\) 的空圆上，四个 Voronoi 胞共享这个顶点。唯一的二维 Delaunay 单元是整个正方形，其棱只有四条边，任何一条对角线都不是该多胞形剖分的棱。可以另选任一对角线把它细分成两个三角形；两种细分都满足“外接开圆无站点”，因为第四点仍在圆周上。这说明允许共圆时，空圆条件本身没有选出唯一的三角剖分。

	现在把 \(D_0\) 沿竖直方向稍作移动。圆 \(ABC\) 的方程是 \(x^2+y^2-x-y=0\)，将 \(D_\varepsilon\) 代入左边，得到
	\[
		\varepsilon(1+\varepsilon).
	\]
	若 \(\varepsilon>0\)，该值为正，所以 \(D_\varepsilon\) 在圆外；圆 \(ACD_\varepsilon\) 的方程为 \(x^2+y^2-(1-\varepsilon)x-(1+\varepsilon)y=0\)，站点 \(B\) 的代入值为 \(\varepsilon>0\)。因此两个三角形 \(ABC\)、\(ACD_\varepsilon\) 都有严格空圆，Delaunay 对角线为 \(AC\)。

	若 \(\varepsilon<0\)，则 \(D_\varepsilon\) 落入圆 \(ABC\) 内，原来的三角形不再是 Delaunay 单元。圆 \(ABD_\varepsilon\) 的方程为 \(x^2+y^2-x-(1+\varepsilon)y=0\)，站点 \(C\) 的代入值为 \(-\varepsilon>0\)。圆 \(BCD_\varepsilon\) 的方程为
	\[
		x^2+y^2-(1-\varepsilon-\varepsilon^2)x-y-\varepsilon(1+\varepsilon)=0,
	\]
	站点 \(A\) 的代入值为 \(-\varepsilon(1+\varepsilon)>0\)。所以这次的两个三角形为 \(ABD_\varepsilon\)、\(BCD_\varepsilon\)，对角线改为 \(BD_\varepsilon\)。每一种情形中，这两个三角形都已覆盖四边形，由剖分的公共面性质便没有另一条可同时穿过内部的对角线。
\end{example}

\begin{figure}[htbp]
	\centering
	% !TeX root = ../Book1.tex
\begin{tikzpicture}[scale=2.45]
	\begin{scope}
		\draw[GeometryBlue,thick] (.5,-.35) -- (.5,1.65);
		\draw[GeometryBlue,thick] (-.45,.5) -- (1.55,.5);
		\draw[thick] (0,0) -- (1,0) -- (1,1) -- (0,1) -- cycle;
		\fill (0,0) circle (.6pt) node[below left] {\(A\)};
		\fill (1,0) circle (.6pt) node[below right] {\(B\)};
		\fill (1,1) circle (.6pt) node[above right] {\(C\)};
		\fill (0,1) circle (.6pt) node[above left] {\(D_0\)};
		\fill[GeometryBlue] (.5,.5) circle (.7pt);
		\node at (.5,-.64) {\(\varepsilon=0\)};
	\end{scope}
	\begin{scope}[xshift=2.75cm]
		\draw[GeometryBlue,thick] (.5,-.35) -- (.5,.5) -- (1.55,.5);
		\draw[GeometryBlue,thick] (-.45,.7) -- (.3,.7) -- (.68,1.65);
		\draw[GeometryBlue,thick] (.3,.7) -- (.5,.5);
		\draw[thick] (0,0) -- (1,0) -- (1,1) -- (0,1.4) -- cycle;
		\draw[thick] (0,0) -- (1,1);
		\fill (0,0) circle (.6pt) node[below left] {\(A\)};
		\fill (1,0) circle (.6pt) node[below right] {\(B\)};
		\fill (1,1) circle (.6pt) node[above right] {\(C\)};
		\fill (0,1.4) circle (.6pt) node[above left] {\(D_\varepsilon\)};
		\fill[GeometryBlue] (.5,.5) circle (.7pt);
		\fill[GeometryBlue] (.3,.7) circle (.7pt);
		\node at (.5,-.64) {\(\varepsilon=2/5\)};
	\end{scope}
\end{tikzpicture}

	\caption[共圆退化及扰动后的分解]{共圆退化及扰动后的分解。黑线为 Delaunay 单元边界，蓝线为 Voronoi 边界。左图四胞共用一个顶点，正方形保持为一个单元；右图取 \(\varepsilon=2/5\)，出现对角线 \(AC\)，对应两个 Voronoi 顶点之间的一条短边。}
	\label{g1:fig:voronoi-degeneracy}
\end{figure}

共圆例子还指出一种容易混淆的“对偶”。若只记录哪些子族 \(\{V_p:p\in T\}\) 具有非空公共交，就得到一个抽象单纯复形，即该族集合的神经。在正方形例子中，四个闭胞具有公共交，因而神经包含由四个抽象顶点组成的三维单纯形；将这些顶点放回平面的原站点，并不能得到一个几何三维单纯形。甚至一对相对顶点的 Voronoi 胞也在中心相交，但这对顶点的连线只是正方形的对角线，并非 Delaunay 多胞形剖分的面。\cref{g1:thm:voronoi-delaunay-duality}恢复单元时使用的是一个 Voronoi 面的相对内点处的\emph{全部}最近站点，不是任取一个具有公共交的站点子集。

\ProseParagraphBreak

严格退化与计算误差也要分开。给定仿射无关的 \(p_0,\ldots,p_d\)，外接球心由线性方程
\[
	2\langle c,p_i-p_0\rangle=\lVert p_i\rVert^2-\lVert p_0\rVert^2,
	\qquad 1\le i\le d,
\]
唯一确定，令 \(r^2=\lVert p_0-c\rVert^2\)。另一站点 \(q\) 在球外、球面上或球内，分别由
\[
	\delta(q)=\lVert q-c\rVert^2-r^2
\]
的正、零、负判定。等价地，可以考察以 \((p_i,\lVert p_i\rVert^2,1)\) 及 \((q,\lVert q\rVert^2,1)\) 为行的 \((d+2)\) 阶行列式；零值表示提升点共面，其符号还要结合前 \(d+1\) 点的定向来解释。有限个严格非零的判定在足够小扰动下保持符号，而\cref{g1:exm:delaunay-square-perturbation}中的零值两侧产生不同剖分。浮点计算得到接近零的数，不能据此宣布原数据严格共球：有理输入可用精确算术核对，近似输入则需有误差范围才能确认符号；不能确认时应保留几何关系未定这一事实。

\ProseParagraphBreak

若 \(S\) 没有张成整个环境空间，应先转到 \(L=\operatorname{aff}S\) 中讨论相应维数的 Delaunay 剖分。设 \(W\) 为 \(L\) 的方向空间，任意位置唯一写成 \(x=y+v\)、\(y\in L\)、\(v\in W^\perp\)，则对每个 \(p\in S\)，
\[
	\lVert x-p\rVert^2=\lVert y-p\rVert^2+\lVert v\rVert^2.
\]
共同增加的第二项不改变最近站点，所以环境空间中的 Voronoi 胞只是 \(L\) 内的胞沿 \(W^\perp\) 方向延伸。低维站点集与共球退化是两种不同情形：前者改变有效环境维数，后者可以在满维凸包中保留非单纯形单元。

\section{支撑函数与离散逼近}
\label{g1:sec:B03}

给出有限个边界点，可以把凸集从内部围起来；给出有限个方向上的支撑半空间，可以把它从外部围起来。两种近似分别遗漏了未采样的点和未采样的方向。本节要估计这些遗漏究竟造成多大误差，并说明怎样把一次计算的结果变成可核验的不等式。所需工具是\cref{g1:thm:convex-projection}的最近点、\cref{g1:prop:support-function-representation}的支撑表示以及\cref{g1:thm:polytope-hv}的有界多面体两种表示。

\subsection{凸体的多面体逼近}
\label{g1:subsec:B03:01}

以下环境为 \(\mathbb R^n\)、\(n\ge1\)，使用标准内积，记闭单位球为 \(\mathbb B\)、单位球面为 \(S^{n-1}\)。本节把具有非空内部的紧凸集称为\term[tuti]{凸体}[convex body]；下面的逼近结论还适用于单点、线段等非空紧凸集。对集合 \(A\) 及 \(r\ge0\)，记
\[
	A+r\mathbb B=\{a+rb:a\in A,\ b\in\mathbb B\}.
\]
若 \(A\) 非空且紧，这正是距离 \(A\) 不超过 \(r\) 的点集：最近点存在，所以距离上界可以写成上述分解。我们先用这样的包含关系表达误差。

\ProseParagraphBreak

给定 \(\delta>0\)，单位球面上的有限集合 \(U=\{u_1,\ldots,u_m\}\) 称为球面的一个\term[deltawang]{\(\delta\)-网}[\(\delta\)-net]，如果每个 \(v\in S^{n-1}\) 都满足某个 \(\lVert v-u_i\rVert\le\delta\)。这里用的是环境欧氏距离，而不是球面上的弧长。这样的有限集合总能选到：以球面上全部点为中心的半径 \(\delta\) 开球覆盖球面，由紧性取有限子覆盖即可。

\begin{lemma}[有界集合的方向误差]\label{g1:lem:support-lipschitz-bound}
	设 \(K\subseteq\mathbb R^n\) 非空且紧，并有 \(K\subseteq R\mathbb B\)、\(R\ge0\)。则对任意 \(u,v\in\mathbb R^n\)，有
	\begin{equation}\label{g1:eq:support-lipschitz-bound}
		|h_K(u)-h_K(v)|\le R\lVert u-v\rVert.
	\end{equation}
\end{lemma}
\begin{proof}
	对每个 \(x\in K\)，有
	\[
		\langle u,x\rangle
		=\langle v,x\rangle+\langle u-v,x\rangle
		\le h_K(v)+R\lVert u-v\rVert.
	\]
	对 \(x\) 取最大值得到一个方向的不等式；交换 \(u,v\) 得到另一个方向。紧性保证两个支撑值均为有限实数，所以可以相减。
\end{proof}

同一个方向网给出两种近似。在方向 \(u_i\) 上取一个达到最大值的点 \(p_i\in K\)，用这些点作凸包；同时保留不等式 \(\langle u_i,x\rangle\le h_K(u_i)\)，取它们的交。前者肯定在 \(K\) 内，后者肯定包含 \(K\)。误差估计还要利用方向网没有留下太大的空隙。

\begin{theorem}[内外多面体逼近]\label{g1:thm:inner-outer-polytope-approximation}
	设 \(K\subseteq\mathbb R^n\) 为非空紧凸集，且 \(K\subseteq R\mathbb B\)、\(R>0\)。取 \(0<\delta<1\) 及 \(S^{n-1}\) 的有限 \(\delta\)-网 \(U=\{u_1,\ldots,u_m\}\)。为每个 \(i\) 选择 \(p_i\in K\) 满足 \(\langle u_i,p_i\rangle=h_K(u_i)\)，并令
	\[
		Q=\operatorname{conv}\{p_1,\ldots,p_m\},\qquad
		P=\bigcap_{i=1}^m\{x:\langle u_i,x\rangle\le h_K(u_i)\},
		\qquad c_\delta=1-\frac{\delta^2}{2}.
	\]
	则 \(Q,P\) 均为非空多胞形，并且
	\begin{equation}\label{g1:eq:inner-outer-polytope-error}
		Q\subseteq K\subseteq P,\qquad
		K\subseteq Q+2R\delta\mathbb B,\qquad
		P\subseteq K+R\delta(1+c_\delta^{-1})\mathbb B.
	\end{equation}
	此外 \(P\subseteq(R/c_\delta)\mathbb B\)。特别地，任意给定的 \(\varepsilon>0\) 都允许选取有限的内外多胞形，使式中两项误差不超过 \(\varepsilon\)。
\end{theorem}
\begin{proof}
	紧性保证各 \(p_i\) 存在。由凸性和支撑函数定义，\(Q\subseteq K\subseteq P\)。\(Q\) 是有限点集的凸包，因而紧。先估计内侧误差。任取单位向量 \(v\)，选择 \(\lVert v-u_i\rVert\le\delta\) 的一个方向，由前一引理得到
	\begin{align*}
		h_K(v)-h_Q(v)
		&\le h_K(v)-\langle v,p_i\rangle\\
		&=h_K(v)-h_K(u_i)+\langle u_i-v,p_i\rangle
		\le2R\delta.
	\end{align*}
	若 \(x\in K\setminus Q\)，取其到 \(Q\) 的最近点 \(q\)，令 \(v=(x-q)/\lVert x-q\rVert\)。\cref{g1:prop:distance-support-formula}的证明说明，此方向满足
	\[
		d(x,Q)=\langle v,x\rangle-h_Q(v)
		\le h_K(v)-h_Q(v)\le2R\delta.
	\]
	对 \(x\in Q\) 距离为零，故全部 \(K\) 都包含于 \(Q+2R\delta\mathbb B\)。

	外侧先须排除无界性。任取 \(x\in P\)、\(x\ne0\)，对单位向量 \(x/\lVert x\rVert\) 选取邻近的 \(u_i\)。因为二者单位长，
	\[
		\left\langle u_i,\frac{x}{\lVert x\rVert}\right\rangle
		=1-\frac12\left\lVert u_i-\frac{x}{\lVert x\rVert}\right\rVert^2
		\ge c_\delta>0.
	\]
	由 \(x\in P\) 得 \(c_\delta\lVert x\rVert\le h_K(u_i)\le R\)。零点也满足所得范数界，因此 \(P\) 有界。它是非空闭半空间交，故为紧多面体；\cref{g1:thm:polytope-hv}说明它也具有有限点的凸包表示。

	现取 \(x\in P\setminus K\)，令 \(p\) 为其到 \(K\) 的最近点，\(v=(x-p)/\lVert x-p\rVert\)。再次选择 \(\lVert v-u_i\rVert\le\delta\)，有
	\begin{align*}
		d(x,K)
		&=\langle v,x\rangle-h_K(v)\\
		&=\langle u_i,x\rangle-h_K(u_i)
		+\langle v-u_i,x\rangle+h_K(u_i)-h_K(v)\\
		&\le\delta\lVert x\rVert+R\delta
		\le R\delta(1+c_\delta^{-1}).
	\end{align*}
	第一项非正是外侧半空间约束的作用。\(x\in K\) 的情形距离为零，由此得到外侧包含。最后 \(0<\delta<1\) 时 \(c_\delta>1/2\)，两项误差分别小于或等于 \(2R\delta\)、\(3R\delta\)。选取 \(0<\delta<\min\{1,\varepsilon/(3R)\}\) 即可。
\end{proof}

定理不要求支撑点唯一。一个方向若支撑在整条边或整个面上，从其中任取一个点仍满足证明所用的等式；不应把任意选择误作几何上的唯一性。方向过稀时，内多胞形还可能降维。若另知 \(a+r\mathbb B\subseteq K\)、\(r>0\)，则证明中的支撑差值给出
\[
	h_Q(v)\ge\langle v,a\rangle+r-2R\delta\qquad(\lVert v\rVert=1).
\]
当 \(2R\delta<r\) 时，由\cref{g1:prop:support-function-representation}可得 \(a+(r-2R\delta)\mathbb B\subseteq Q\)，所以这时内近似也是凸体。这里的 \(R\) 只需要是包含 \(K\) 的一个已知球半径；把坐标原点移到靠近 \(K\) 的位置，往往能减小这一误差常数。

% !TeX root = ../Book1.tex
\begin{figure}[htbp]
	\centering
	\begin{tikzpicture}[scale=2.2,font=\small]
		\draw[thick,GeometryGray]
			(22.5:{1/cos(22.5)})--(67.5:{1/cos(22.5)})--(112.5:{1/cos(22.5)})--(157.5:{1/cos(22.5)})--(202.5:{1/cos(22.5)})--(247.5:{1/cos(22.5)})--(292.5:{1/cos(22.5)})--(337.5:{1/cos(22.5)})--cycle;
		\filldraw[fill=GeometryBlue!10,draw=GeometryBlue,thick]
			(0:1)--(45:1)--(90:1)--(135:1)--(180:1)--(225:1)--(270:1)--(315:1)--cycle;
		\draw[black,thick] (0,0) circle[radius=1];
		\foreach \a in {0,45,90,135,180,225,270,315}{
			\fill[GeometryBlue] (\a:1) circle[radius=0.018];
		}
		\draw[->,GeometryGray] (0,0)--(0:0.78) node[midway,below] {\(u_0\)};
		\node[GeometryBlue] at (-0.3,-0.1) {\(Q\)};
		\node[right] at (1.16,0.85) {\(P\)：支撑半平面的交};
		\draw[GeometryGray] (1.12,0.83)--(1,{tan(22.5)});
		\node[right] at (1.16,0.35) {\(\mathbb B\)：单位圆盘};
		\draw (1.12,0.33)--(0.95,0.31);
		\node[right,GeometryBlue] at (1.16,-0.18) {\(Q\)：支撑点的凸包};
	\end{tikzpicture}
	\caption[有限方向给出的内外近似]{八个等角单位方向给出 \(Q\subseteq\mathbb B\subseteq P\)。内多边形连接取到支撑值的点，外多边形保留相应支撑半平面；未采样方向上的两侧间隙由误差估计控制。}
	\label{g1:fig:support-polytope-approximation}
\end{figure}


\subsection{Hausdorff 距离与支撑函数}
\label{g1:subsec:B03:02}

上面的两项包含分别控制了“\(K\) 的点离近似多胞形多远”与“近似多胞形的点离 \(K\) 多远”。如果只保留其中一项，任意一个集合的很小子集都可能给出零误差。集合之间的距离因此需要同时检查两个方向。

\begin{definition}[Hausdorff 距离]\label{g1:def:hausdorff-distance}
	对非空紧集 \(K,L\subseteq\mathbb R^n\)，定义它们的\term[hausdorffjuli]{Hausdorff 距离}[Hausdorff distance]为
	\begin{equation}\label{g1:eq:hausdorff-distance}
		d_H(K,L)=\max\left\{\max_{x\in K}d(x,L),\ \max_{y\in L}d(y,K)\right\}.
	\end{equation}
	这里 \(d(x,L)=\min_{y\in L}\lVert x-y\rVert\)。
	\symidx{\(d_H(K,L)\)：非空紧集的 Hausdorff 距离}
\end{definition}

定义中的最小值由紧性取得。对任意 \(x,x'\)，三角不等式给出 \(|d(x,L)-d(x',L)|\le\lVert x-x'\rVert\)，所以距离函数连续，外层的两个最大值也能取得。由最近点存在，对每个 \(\varepsilon\ge0\) 有
\begin{equation}\label{g1:eq:hausdorff-inclusions}
	d_H(K,L)\le\varepsilon
	\quad\Longleftrightarrow\quad
	K\subseteq L+\varepsilon\mathbb B\ \text{且}\ L\subseteq K+\varepsilon\mathbb B.
\end{equation}
这也核验了距离公理：对称性和非负性直接成立；零距离使两个紧集互相包含；若再有非空紧集 \(M\)，沿两次包含相加得到
\[
	d_H(K,M)\le d_H(K,L)+d_H(L,M).
\]
这里用到 \(a\mathbb B+b\mathbb B=(a+b)\mathbb B\)、\(a,b\ge0\)：一个方向来自三角不等式，反方向在 \(a+b>0\) 时把同一个单位球中向量分别乘 \(a,b\)，零值情形单独成立。

\begin{theorem}[Hausdorff 距离的支撑表示]\label{g1:thm:hausdorff-support-isometry}
	设 \(K,L\subseteq\mathbb R^n\) 为非空紧凸集。则
	\begin{equation}\label{g1:eq:hausdorff-support-isometry}
		d_H(K,L)=\max_{u\in S^{n-1}}|h_K(u)-h_L(u)|.
	\end{equation}
	因此，把一个非空紧凸集对应到它在单位球面上的支撑函数，保持距离；右侧使用连续函数的一致范数。
\end{theorem}
\begin{proof}
	先把距离中的外扩集合写成支撑值。\(L+\varepsilon\mathbb B\) 为紧凸集：它是两个紧集之积在连续加法映射下的像，凸性则逐项取凸组合验证。对任意 \(u\)，两个分量可以独立选择，所以
	\[
		h_{L+\varepsilon\mathbb B}(u)
		=\max_{y\in L,\,b\in\mathbb B}\langle u,y+\varepsilon b\rangle
		=h_L(u)+\varepsilon\lVert u\rVert.
	\]
	由\cref{g1:prop:support-function-representation}，\(K\subseteq L+\varepsilon\mathbb B\) 当且仅当 \(h_K(u)\le h_L(u)+\varepsilon\lVert u\rVert\) 对全部 \(u\) 成立。支撑函数的正齐次性说明，只检查单位方向就足够，零方向的条件自动成立。交换 \(K,L\)，再用式\eqref{g1:eq:hausdorff-inclusions}，便得到
	\[
		d_H(K,L)\le\varepsilon
		\quad\Longleftrightarrow\quad
		|h_K(u)-h_L(u)|\le\varepsilon\quad(\lVert u\rVert=1).
	\]
	由\cref{g1:lem:support-lipschitz-bound}，两个支撑函数均连续，故右侧差值在单位球面上有最大值。取这个最大值为 \(\varepsilon\)，以及反过来取 \(\varepsilon=d_H(K,L)\)，即得等式。
\end{proof}

\cref{g1:thm:inner-outer-polytope-approximation}现在可统一写成
\[
	d_H(K,Q)\le2R\delta,\qquad
	d_H(K,P)\le R\delta(1+c_\delta^{-1}).
\]
这些估计对任意紧凸集有效；特定形状往往给出更好的误差。图中的圆盘便可以精确计算。

\begin{example}[圆盘的正多边形逼近]\label{g1:exm:regular-polygon-support-error}
	在 \(\mathbb R^2\) 中令 \(K=\mathbb B\)，取 \(m\ge3\) 个单位方向
	\[
		u_j=\bigl(\cos(2\pi j/m),\sin(2\pi j/m)\bigr),\qquad 0\le j<m,
	\]
	并令 \(Q_m=\operatorname{conv}\{u_j\}\)、\(P_m=\{x:\langle u_j,x\rangle\le1\ \text{对全部 }j\}\)。对任意单位方向 \(v\)，它与最近 \(u_j\) 的夹角 \(\theta\) 介于 \(0\) 和 \(\pi/m\) 之间，因而
	\[
		h_{Q_m}(v)=\cos\theta,\qquad
		d_H(Q_m,\mathbb B)=1-\cos(\pi/m).
	\]
	最大误差在相邻两个采样方向的角平分线方向取得。对外近似，射线上的点 \(tv\)、\(t\ge0\)，属于 \(P_m\) 当且仅当
	\[
		t\max_j\langle u_j,v\rangle\le1.
	\]
	上式中的最大值为 \(\cos\theta\ge\cos(\pi/m)>0\)，所以 \(P_m\) 中最大的范数为 \(1/\cos(\pi/m)\)，也在角平分线方向取得。点到单位圆盘的距离为 \(\max\{0,\lVert x\rVert-1\}\)，故
	\[
		d_H(P_m,\mathbb B)=\frac1{\cos(\pi/m)}-1.
	\]
	例如 \(m=4\) 时，内近似为菱形，误差为 \(1-1/\sqrt2\)；外近似为正方形，误差为 \(\sqrt2-1\)。两者取样方向相同，误差仍然不同。这些公式直接来自圆盘的形状，不能把其中更快的收敛速度无条件推广给所有紧凸集。
\end{example}

凸性在式\eqref{g1:eq:hausdorff-support-isometry}中不可删除。单位圆周 \(S^1\) 与闭单位圆盘 \(\mathbb B\subset\mathbb R^2\) 的支撑函数相同，但 \(d_H(S^1,\mathbb B)=1\)：圆盘中心到圆周的距离为 \(1\)，其他圆盘点的距离为 \(1-\lVert x\rVert\)。支撑函数只恢复闭凸包，这与\cref{g1:prop:support-function-basic}一致；它不会记录集合内部的空洞。

\subsection{数值表示的证书}
\label{g1:subsec:B03:03}

计算给出一张顶点表、一组半空间或若干近似支撑值后，需要说明哪些有限数据能证明所声称的包含关系及误差。本节把能够通过明确等式、不等式核验某项结论的有限数据称为该结论的\term[zhengshu]{证书}[certificate]。证书可以由程序寻找，但核验所需的推理应能从数据本身读出。例如，声称两份多面体描述相同，需要两个方向的包含证书，不能只验证列出的点都满足列出的不等式。

\begin{proposition}[有限包含与支撑界的证书]\label{g1:prop:finite-polytope-certificates}
	设 \(v_1,\ldots,v_N\in\mathbb R^n\)、\(N\ge1\)，\(Q=\operatorname{conv}\{v_1,\ldots,v_N\}\)，并设 \(A\) 为实 \(s\times n\) 矩阵、\(b\in\mathbb R^s\)、\(P=\{x:Ax\le b\}\)。下列有限数据具有所述作用。
	\begin{enumerate}
		\item \(Q\subseteq P\) 当且仅当对所有 \(j\) 有 \(Av_j\le b\)。
		\item 若 \(P\) 非空有界，并已列出其全部顶点 \(w_1,\ldots,w_M\)，则对每个 \(k\) 给出
		\[
			w_k=\sum_j\theta_{kj}v_j,\qquad
			\theta_{kj}\ge0,\qquad \sum_j\theta_{kj}=1,
		\]
		便可证明 \(P\subseteq Q\)。
		\item 对 \(c\in\mathbb R^n\)，若 \(\lambda\in\mathbb R^s\) 满足 \(\lambda\ge0\)、\(A^{\mathsf T}\lambda=c\)，则 \(\langle c,x\rangle\le\langle\lambda,b\rangle\) 对所有 \(x\in P\) 成立。若再给出 \(p\in P\)，就得到
		\[
			\langle c,p\rangle\le h_P(c)\le\langle\lambda,b\rangle.
		\]
	\end{enumerate}
\end{proposition}
\begin{proof}
	第一项的必要性来自 \(v_j\in Q\)；充分性来自每条线性不等式在凸组合下保持。第二项由\cref{g1:thm:polytope-hv}知 \(P\) 是其全部顶点的凸包；给出的每个凸组合使 \(w_k\in Q\)，再用 \(Q\) 的凸性便得到全部 \(P\subseteq Q\)。第三项对每个 \(x\in P\) 按非负系数相加：
	\[
		\langle c,x\rangle=\langle A^{\mathsf T}\lambda,x\rangle
		=\langle\lambda,Ax\rangle\le\langle\lambda,b\rangle.
	\]
	取上确界得到上界；可行点 \(p\) 的函数值是上确界的下界。
\end{proof}

“全部顶点”也有有限核验办法。先给一个满足 \(Ax\le b\) 的点证明非空，再证明有界，例如分别为 \(\pm x_k\le R_k\) 给出第三项中的非负乘子。于是 \(P\subseteq\prod_k[-R_k,R_k]\)。接着枚举全部 \(n\) 条线性无关约束的组合，解相应等式组并保留满足全部不等式的解，便得到全部顶点。其原因与第三章一致：若某点的活动法向量不张成 \(\mathbb R^n\)，就能沿一个同时垂直的非零方向作两侧小扰动，有限条非活动约束仍保持，故该点不是顶点；若活动法向量张成全空间，其中 \(n\) 条就唯一固定该点。这个办法允许冗余约束和低维多面体，重复出现的顶点只需合并；顶点数很多时枚举可能昂贵，但所得证书的含义没有改变。

\ProseParagraphBreak

例如，对三角形
\[
	P=\{(x,y):x\ge0,\ y\ge0,\ x+y\le1\},\qquad c=(3/5,4/5),
\]
按 \(-x\le0,-y\le0,x+y\le1\) 排列约束。取 \(p=(0,1)\) 和 \(\lambda=(1/5,0,4/5)\)，可以直接验算 \(A^{\mathsf T}\lambda=c\) 及
\[
	\langle c,p\rangle=4/5=\langle\lambda,b\rangle.
\]
上下界相等，故这份有限数据证明了 \(h_P(c)=4/5\)，无需把某个优化程序报告的“最优”二字当作证明。

\ProseParagraphBreak

要把有限个支撑值的误差扩展到全部方向，还需要方向覆盖的证据。除了平面中的等角网，一般维数可以使用整数网格。取整数 \(M>\sqrt n\)，令
\begin{equation}\label{g1:eq:integer-direction-net}
	U_M=\left\{\frac{q}{\lVert q\rVert}:q\in\mathbb Z^n,\ 0<\max_k|q_k|\le M\right\}.
\end{equation}
这是有限集合，并且是半径 \(\delta_M=\sqrt n/M\) 的方向网。事实上，对任意单位向量 \(u\)，把 \(Mu\) 的各坐标舍入到最近整数，得到上述范围内的 \(q\)，并有 \(z=q/M\) 满足 \(\lVert z-u\rVert\le\sqrt n/(2M)<1\)，所以 \(q\ne0\)。再由反三角不等式，
\[
	\left\lVert\frac{q}{\lVert q\rVert}-u\right\rVert
	\le\left\lVert\frac z{\lVert z\rVert}-z\right\rVert+\lVert z-u\rVert
	=|1-\lVert z\rVert|+\lVert z-u\rVert
	\le\frac{\sqrt n}{M}.
\]
因此网格的有限生成规则本身就提供了覆盖半径的证明。

\begin{proposition}[支撑区间给出的距离证书]\label{g1:prop:finite-support-distance-certificate}
	设 \(K,L\subseteq\mathbb R^n\) 为非空紧凸集，\(K\subseteq R_K\mathbb B\)、\(L\subseteq R_L\mathbb B\)，其中 \(R_K,R_L\ge0\)。取单位球面的一个非空有限 \(\delta\)-网 \(\{u_i\}_{i=1}^m\)，并对每个 \(i\) 给出已经证明的区间
	\[
		a_i\le h_K(u_i)\le b_i,\qquad
		c_i\le h_L(u_i)\le d_i.
	\]
	令
	\[
		D_-=\max\bigl(\{0\}\cup\{a_i-d_i,c_i-b_i:1\le i\le m\}\bigr),
	\]
	\[
		D_+=\max_{1\le i\le m}\max\{b_i-c_i,d_i-a_i\}.
	\]
	则有可核验的双侧误差界
	\begin{equation}\label{g1:eq:finite-support-distance-certificate}
		D_-\le d_H(K,L)\le D_++(R_K+R_L)\delta.
	\end{equation}
\end{proposition}
\begin{proof}
	在每个采样方向上，\(h_K(u_i)-h_L(u_i)\ge a_i-d_i\)，而反向差值不小于 \(c_i-b_i\)。\cref{g1:thm:hausdorff-support-isometry}因此给出下界。另一方面，两份区间使
	\[
		|h_K(u_i)-h_L(u_i)|\le\max\{b_i-c_i,d_i-a_i\}\le D_+.
	\]
	任取单位 \(v\)，选择 \(\lVert v-u_i\rVert\le\delta\)，并对两个支撑函数使用\cref{g1:lem:support-lipschitz-bound}，得到
	\[
		|h_K(v)-h_L(v)|
		\le D_++R_K\delta+R_L\delta.
	\]
	取全部单位方向的最大值，再用支撑表示即得上界。
\end{proof}

特别地，支撑值精确已知时，有限最大误差 \(E=\max_i|h_K(u_i)-h_L(u_i)|\) 满足
\[
	E\le d_H(K,L)\le E+(R_K+R_L)\delta.
\]
末项控制未采样方向，不能因为已测误差为零就删去。例如单位圆盘与正方形 \([-1,1]^2\) 在四个坐标轴方向的支撑值全为 \(1\)，Hausdorff 距离却为 \(\sqrt2-1\)。对一般 \(K\)，支撑值的下界可由已经证明属于 \(K\) 的点提供，上界则须由有效的半空间包含或其他已证估计提供；对半空间表示的多胞形，前一命题的可行点与非负乘子恰好给出这两种数据。

\ProseParagraphBreak

浮点误差也必须落在明确的界内。例如，已知 \(P=\{x:Ax\le b\}\subseteq R\mathbb B\) 非空，给出已确认非负的乘子 \(\lambda\ge0\)，但等式 \(A^{\mathsf T}\lambda=c\) 只近似成立。若已证明残差满足 \(\lVert c-A^{\mathsf T}\lambda\rVert\le\rho\)，则仍有
\begin{equation}\label{g1:eq:dual-residual-support-bound}
	h_P(c)\le\langle\lambda,b\rangle+R\rho.
\end{equation}
确实，对每个 \(x\in P\)，把 \(c\) 分成 \(A^{\mathsf T}\lambda\) 和残差，配对后分别使用 \(Ax\le b\) 与 Cauchy--Schwarz 不等式即可。若内积本身也由区间运算给出，就取其已证明的上端点；\(\rho\) 同样须是范数的上界。乘子的非负性必须确实成立，不能把一个小负数按显示精度忽略。下界也需要可行性：若只有一个近似点 \(z\)，但已证明存在 \(p\in P\) 使 \(\lVert z-p\rVert\le\eta\)，才能推出 \(h_P(c)\ge\langle c,z\rangle-\lVert c\rVert\eta\)。

\begin{example}[小约束残差未必意味着小距离]\label{g1:exm:small-residual-large-distance}
	给定 \(a>0\)，令
	\[
		T_a=\{(x,y):x\ge0,\ ax+y\le a,\ ax-y\le a\}
		=\operatorname{conv}\{(0,a),(0,-a),(1,0)\}.
	\]
	最后两条不等式除以 \(\sqrt{1+a^2}\) 后，法向量均为单位向量。点 \(z=(2,0)\) 对第一条约束可行，对其余每条归一化约束的超出量都只有 \(a/\sqrt{1+a^2}\)。然而 \(T_a\) 内所有点的横坐标不超过 \(1\)，且 \((1,0)\in T_a\)，故 \(d(z,T_a)=1\)。当 \(a\to0\) 时，两条支撑边越来越接近平行，约束残差趋于零，几何距离仍为 \(1\)。因此“每条约束的误差小于容差”本身不能代替点到集合的距离证书，即使各法向量已经归一化。
\end{example}

这些证书分别控制数据的三种误差：有限点与半空间的包含关系、有限方向以外的遗漏，以及数值运算的误差。只要半径界、方向覆盖和支撑区间各有依据，式\eqref{g1:eq:finite-support-distance-certificate}就把它们合成一个对整个凸集有效的距离界。

\begin{chaptersummary}
	格给欧氏空间增加了整数结构。换格基的行列式只能为 \(\pm1\)，因此基本区域的体积与格基选择无关。二维中，把凸格多边形分成不含额外格点的三角形，既能按面积相加，也能按内部点与边界点计数；两种计算相合便得到 Pick 公式。膨胀后的面积按平方增长，边界格点按一次增长，由此得到二维 Ehrhart 多项式。

	\ProseParagraphBreak

	距离平方的共同二次项相消，使最近站点的比较成为半空间不等式。抛物面提升又把空球条件化为支撑超平面条件，连接 Voronoi 面与 Delaunay 单元。共球点可以共同构成一个非单纯形单元；再作三角剖分需要额外选择，不能把该选择当成距离数据已经唯一决定的结果。

	\ProseParagraphBreak

	对于非空紧凸集，支撑函数不仅记录包含关系，也精确记录 Hausdorff 距离。有限方向网配合统一半径界，可以控制没有采样的方向，从而为内外多面体逼近给出误差。数值表示是否可信，取决于是否附有可以核算的约束、权重和误差界，而不取决于图像是否平滑或输出了多少位小数。
\end{chaptersummary}

\BookExercises
以下习题分别训练格基与计数、空球判定、退化分析和凸体误差估计。计算题中的每个数值都须连同适用条件解释，证书题则须同时给出所声称上下界的依据。

\begin{exercise}
	\label{g1:ex:lattice-basis-cell}
	在标准欧氏平面中，令 \(b_1=(2,0)\)、\(b_2=(1,1)\)，\(\Lambda=\mathbb Zb_1+\mathbb Zb_2\)。
	\begin{enumerate}
		\item 给出整数点 \((x,y)\) 属于 \(\Lambda\) 的充要条件，并计算 \(\operatorname{covol}(\Lambda)\)。
		\item 判断 \(c_1=(3,1)\)、\(c_2=(4,2)\) 是否构成同一格的格基。
		\item 把 \(q=(11/2,9/4)\) 唯一写为 \(q=\lambda+f\)，其中 \(\lambda\in\Lambda\)、\(f\in F_b\)。
	\end{enumerate}
\end{exercise}
\begin{solution}
	(1) 格点具有形式 \(m(2,0)+n(1,1)=(2m+n,n)\)，所以整数点 \((x,y)\) 属于 \(\Lambda\) 当且仅当 \(x-y\) 是偶数；满足此条件时，取 \(n=y\)、\(m=(x-y)/2\) 即可。基矩阵行列式为 \(2\)，所以余体积为 \(2\)。

	(2) \(c_1=b_1+b_2\)、\(c_2=b_1+2b_2\)，变换矩阵为
	\[
		U=\begin{pmatrix}1&1\\1&2\end{pmatrix},\qquad \det U=1.
	\]
	它是幺模矩阵，因此两向量是同一格的格基。也可由 \(b_1=2c_1-c_2\)、\(b_2=c_2-c_1\) 直接看出反向整数生成。

	(3) \(q\) 的 \(b\)-坐标为 \((13/8,9/4)\)，整数部分为 \((1,2)\)，小数部分为 \((5/8,1/4)\)。因此
	\[
		\lambda=b_1+2b_2=(4,2),\qquad
		f=\frac58b_1+\frac14b_2=\left(\frac32,\frac14\right).
	\]
	两个小数坐标都在 \([0,1)\) 中，故 \(f\in F_b\)。若还有另一分解，两余数点的差既有整数格坐标，又在各坐标的 \((-1,1)\) 内，只能为零，所以分解唯一。
\end{solution}

\begin{exercise}
	\label{g1:ex:pick-quadrilateral}
	令 \(P\) 为依次以 \((0,0),(4,0),(3,2),(0,3)\) 为顶点的凸四边形。计算面积、边界格点数和内部格点数，并求 \(L_P(n)\)。再分别求 \(2P\) 的内部格点数与全部格点数。
\end{exercise}
\begin{solution}
	沿从 \((0,0)\) 到 \((3,2)\) 的对角线分成两个三角形，面积分别为
	\[
		\frac12\left|\det\begin{pmatrix}4&3\\0&2\end{pmatrix}\right|=4,
		\qquad
		\frac12\left|\det\begin{pmatrix}3&0\\2&3\end{pmatrix}\right|=\frac92.
	\]
	所以 \(A=17/2\)。四条边向量的坐标最大公因数依次为 \(4,1,1,3\)，于是 \(B=9\)。Pick 公式给出 \(I=17/2-9/2+1=5\)。可用实际内部点 \((1,1),(2,1),(3,1),(1,2),(2,2)\) 再核对这个数。

	因此
	\[
		L_P(n)=\frac{17}{2}n^2+\frac92n+1.
	\]
	当 \(n=2\) 时，面积为 \(34\)，边界格点数为 \(18\)，内部格点数为 \(34-9+1=26\)，全部格点数为 \(26+18=44\)。
\end{solution}

\begin{exercise}
	\label{g1:ex:rational-count-parity}
	取有理顶点矩形 \(Q=[0,1/2]\times[0,1]\)。计算 \(\#(nQ\cap\mathbb Z^2)\)，分别写出 \(n\) 为偶数与奇数时的多项式表达，并证明不存在一个多项式在所有非负整数 \(n\) 上给出该计数。
\end{exercise}
\begin{solution}
	第一坐标可以是 \(0,\ldots,\lfloor n/2\rfloor\)，第二坐标可以是 \(0,\ldots,n\)，所以计数为
	\[
		\left(\left\lfloor\frac n2\right\rfloor+1\right)(n+1)
		=\begin{cases}
			\dfrac12n^2+\dfrac32n+1,&n\text{ 为偶数},\\[3pt]
			\dfrac12n^2+n+\dfrac12,&n\text{ 为奇数}.
		\end{cases}
	\]
	若一个多项式 \(f\) 给出全部计数，那么它与第一行的多项式在无穷多个非负偶数上相等。非零多项式只有有限个根，所以 \(f\) 必恒等于第一行。但两行多项式之差为 \((n+1)/2\)，在正奇数上不为零，矛盾。这个例子显示，把整顶点放宽为有理顶点后，点数可能随倍数的奇偶而改变公式。
\end{solution}

\begin{exercise}
	\label{g1:ex:simplex-cone-remainders}
	对格三角形 \(S=\operatorname{conv}\{(0,0),(2,0),(0,1)\}\)，写出\cref{g1:prop:lattice-simplex-ehrhart}的半开区域 \(\Pi\) 中全部格点，求 \(h_0,h_1,h_2\)，并据此求 \(L_S(n)\)。再按水平整数坐标逐行计数，独立核验结果。
\end{exercise}
\begin{solution}
	提升向量为 \(w_0=(0,0,1)\)、\(w_1=(2,0,1)\)、\(w_2=(0,1,1)\)，故 \(\Pi\) 中的点具有形式
	\[
		(2t_1,t_2,t_0+t_1+t_2),\qquad 0\le t_0,t_1,t_2<1.
	\]
	若它是整数点，第二坐标迫使 \(t_2=0\)，第一坐标迫使 \(t_1=0\) 或 \(1/2\)。当 \(t_1=0\) 时，末坐标为整数迫使 \(t_0=0\)；当 \(t_1=1/2\) 时，末坐标为整数迫使 \(t_0=1/2\)。因此仅有格点 \((0,0,0)\)、\((1,0,1)\)，从而 \(h_0=1,h_1=1,h_2=0\)，得到
	\[
		L_S(n)=\binom{n+2}{2}+\binom{n+1}{2}=(n+1)^2.
	\]
	另一方面，\(nS=\{(x,y):x,y\ge0,\ x+2y\le2n\}\)。固定整数 \(y=k\in\{0,\ldots,n\}\)，可取 \(2(n-k)+1\) 个整数横坐标。因此总数为
	\[
		\sum_{k=0}^n\bigl(2(n-k)+1\bigr)
		=2\frac{n(n+1)}2+(n+1)=(n+1)^2,
	\]
	与半开锥计数相符。
\end{solution}

\begin{exercise}
	\label{g1:ex:empty-lattice-tetrahedron}
	对整数 \(m\ge2\)，令
	\[
		T_m=\operatorname{conv}\{(0,0,0),(1,0,0),(0,1,0),(1,1,m)\}.
	\]
	证明 \(T_m\) 除四个顶点外没有其他格点，并计算其体积。说明这一族如何阻止把二维“不含额外格点的三角形面积固定”直接推广到三维。
\end{exercise}
\begin{solution}
	任一点均可写为
	\[
		(x,y,z)=\alpha(1,0,0)+\beta(0,1,0)+\gamma(1,1,m),
		\qquad \alpha,\beta,\gamma\ge0,\quad\alpha+\beta+\gamma\le1.
	\]
	所以 \(z=m\gamma\in[0,m]\)。若它是格点且 \(0<z<m\)，则 \(0<\gamma<1\)，又有 \(x=\alpha+\gamma>0\)、\(y=\beta+\gamma>0\)。两个正整数都至少为 \(1\)，于是
	\[
		\alpha+\beta+\gamma=x+y-\gamma\ge2-\gamma>1,
	\]
	与凸组合条件矛盾。若 \(z=0\)，则 \(\gamma=0\)，非负整数 \(x,y\) 满足 \(x+y\le1\)，只得到底面的三个顶点。若 \(z=m\)，则 \(\gamma=1\)，迫使 \(\alpha=\beta=0\)，只得到第四个顶点。

	在高度 \(z\in[0,m]\) 处，截面是底面标准三角形缩放 \(1-z/m\) 倍后再平移所得，面积为 \(\tfrac12(1-z/m)^2\)。沿高度积分得
	\[
		\operatorname{Vol}_3(T_m)
		=\frac12\int_0^m\left(1-\frac zm\right)^2\mathrm dz
		=\frac m6.
	\]
	因此这些没有额外格点的四面体，其体积仍可任意增大。二维引理证明中，把平行四边形中的格点 \(w\) 换成 \(u+v-w\) 就能落入三角形；更多系数时，反射后的系数和未必小于 \(1\)，该步骤不再适用。
\end{solution}

\begin{exercise}\label{g1:ex:voronoi-interior-site}
	在通常欧氏平面中取站点集 \(S=\{O,E,N,W,H\}\)，其中
	\[
		O=(0,0),\quad E=(2,0),\quad N=(0,2),\quad W=(-2,0),\quad H=(0,-2).
	\]
	求全部 Voronoi 胞和二维 Delaunay 单元。四个外侧站点同在一个圆上，这是否使 Delaunay 剖分出现一个四边形单元？说明理由。
\end{exercise}
\begin{solution}
	与四个外侧站点比较距离，得到
	\[
		V_O=\{(x,y):|x|\le1,\ |y|\le1\}.
	\]
	对 \(E\)，与 \(O\) 比较给出 \(x\ge1\)，与 \(N,H\) 比较给出 \(x\ge y\)、\(x\ge-y\)；与 \(W\) 比较仅要求 \(x\ge0\)，已被前面的条件包含。因此全部外侧胞为
	\[
		\begin{aligned}
			V_E&=\{(x,y):x\ge1,\ x\ge|y|\},&
			V_W&=\{(x,y):-x\ge1,\ -x\ge|y|\},\\
			V_N&=\{(x,y):y\ge1,\ y\ge|x|\},&
			V_H&=\{(x,y):-y\ge1,\ -y\ge|x|\}.
		\end{aligned}
	\]
	它们和中央正方形覆盖平面，中央胞有界，四个外侧胞无界。

	Voronoi 顶点为 \((1,1),(-1,1),(-1,-1),(1,-1)\)。以 \((1,1)\) 为中心、半径 \(\sqrt2\) 的圆通过 \(O,E,N\)，而到 \(W,H\) 的距离平方都为 \(10>2\)。其他三个中心由旋转得到。因此二维 Delaunay 单元恰为
	\[
		\operatorname{conv}\{O,E,N\},\quad
		\operatorname{conv}\{O,N,W\},\quad
		\operatorname{conv}\{O,W,H\},\quad
		\operatorname{conv}\{O,H,E\}.
	\]
	这四个三角形已经覆盖 \(\operatorname{conv}S=\{(x,y):|x|+|y|\le2\}\)，并按公共边拼接，故没有遗漏其他二维单元。外侧四点所在的圆中心为 \(O\)、半径为 \(2\)，其开圆内部包含站点 \(O\)，所以不是空圆。共圆只有同时满足空圆条件时，才会产生这里所说的非单纯形接触单元。
\end{solution}

\begin{exercise}\label{g1:ex:largest-empty-disk-constrained}
	设有限互异点集 \(S\subset\mathbb R^2\) 张成平面，记 \(K=\operatorname{conv}S\)，并把空圆的中心限制在 \(K\) 中。证明最大的空开圆半径
	\[
		R=\max_{x\in K}\min_{p\in S}\lVert x-p\rVert
	\]
	可在某个多胞形 \(K\cap V_p\) 的顶点处取得。对 \(S=\{A=(0,0),B=(4,0),C=(1,1)\}\)，求 \(R\) 及一个取到最大值的中心，并说明只检查位于 \(K\) 内的 Voronoi 顶点为什么会失败。
\end{exercise}
\begin{solution}
	函数 \(\rho(x)=\min_{p\in S}\lVert x-p\rVert\) 是有限个连续函数的最小值，因而连续；\(K\) 紧，所以最大值存在。在 \(P_p=K\cap V_p\) 上，有 \(\rho(x)^2=\lVert x-p\rVert^2\)。\(P_p\) 是有界多面体，由\cref{g1:thm:polytope-hv}等于其有限顶点的凸包。若 \(x=\sum_i\lambda_i v_i\) 为顶点凸组合，则
	\[
		\lVert x-p\rVert^2
		\le\sum_i\lambda_i\lVert v_i-p\rVert^2
		\le\max_i\lVert v_i-p\rVert^2.
	\]
	第一个不等式可由展开平方直接得到，其差为 \(\sum_i\lambda_i\lVert v_i-x\rVert^2\ge0\)。故在每个 \(P_p\) 上都能到顶点取最大值；这有限个多胞形覆盖 \(K\)，结论成立。

	在给定三角形中，三个 Voronoi 胞为
	\[
		\begin{aligned}
			V_A&=\{(x,y):x\le2,\ x+y\le1\},\\
			V_B&=\{(x,y):x\ge2,\ 3x-y\ge7\},\\
			V_C&=\{(x,y):x+y\ge1,\ 3x-y\le7\}.
		\end{aligned}
	\]
	三胞唯一的公共点是 \((2,-1)\)，在 \(K\) 外。限制到 \(K\) 后，三个多胞形除原站点外的顶点为
	\[
		(1,0),\quad (1/2,1/2),\quad (7/3,0),\quad (5/2,1/2).
	\]
	其中前两点由直线 \(x+y=1\) 与边 \(AB,AC\) 相交得到，后两点由 \(3x-y=7\) 与边 \(AB,BC\) 相交得到。它们的最近站点距离依次为
	\[
		1,\qquad \frac1{\sqrt2},\qquad \frac53,\qquad \sqrt{\frac52}.
	\]
	三个原站点的最近距离为零。比较这些候选值得到 \(R=5/3\)，在中心 \((7/3,0)\) 取得。该圆通过 \(B,C\)，而到 \(A\) 的距离为 \(7/3>5/3\)。最大值中心是 Voronoi 边与凸包边界的交点，平面 Voronoi 图的唯一顶点却在约束区域外，因此只查后者会漏掉答案。
\end{solution}

\begin{exercise}\label{g1:ex:incircle-determinant}
	设平面中的 \(A=(a_1,a_2)\)、\(B=(b_1,b_2)\)、\(C=(c_1,c_2)\) 不共线，且
	\[
		\Delta=\det\begin{pmatrix}a_1&a_2&1\\b_1&b_2&1\\c_1&c_2&1\end{pmatrix}>0.
	\]
	记其外接圆心为 \(o\)、半径为 \(r\)。对 \(P=(p_1,p_2)\)，证明
	\[
		\det\begin{pmatrix}
			a_1&a_2&\lVert A\rVert^2&1\\
			b_1&b_2&\lVert B\rVert^2&1\\
			c_1&c_2&\lVert C\rVert^2&1\\
			p_1&p_2&\lVert P\rVert^2&1
		\end{pmatrix}
		=-\Delta\bigl(\lVert P-o\rVert^2-r^2\bigr).
	\]
	据此判定 \(P\) 在圆内、圆上、圆外时行列式的符号。再取 \(A=(0,0),B=(1,0),C=(0,1)\)、\(P=(1,1+10^{-m})\)，其中 \(m\) 为正整数，精确求此行列式，并解释为什么不能把数值很小的结果直接视为共圆。
\end{exercise}
\begin{solution}
	记 \(o=(o_1,o_2)\)。把行列式的第三列减去第一列的 \(2o_1\) 倍、第二列的 \(2o_2\) 倍，再加上第四列的 \(\lVert o\rVert^2-r^2\) 倍，行列式不变。新第三列的各项为对应点 \(X\) 的 \(\lVert X-o\rVert^2-r^2\)。前三项为零，最后一项为 \(\delta=\lVert P-o\rVert^2-r^2\)。沿该列展开，唯一非零项的代数余子式符号为 \((-1)^{4+3}=-1\)，得到 \(-\Delta\delta\)。因 \(\Delta>0\)，圆内、圆上、圆外分别给出正、零、负值。若交换 \(A,B\) 的顺序，\(\Delta\) 和大行列式同时变号，判定时必须保留这个定向约定。

	具体三角形满足 \(\Delta=1\)、\(o=(1/2,1/2)\)、\(r^2=1/2\)。置 \(t=10^{-m}>0\)，有
	\[
		\delta=(1-1/2)^2+(1+t-1/2)^2-1/2=t+t^2,
	\]
	所以行列式精确等于 \(-10^{-m}-10^{-2m}<0\)。无论 \(m\) 多大，\(P\) 都严格在圆外。固定数值阈值会把足够大的 \(m\) 误判成共圆；若输入本身只是近似值，还必须计入输入误差，才能判断真实符号是否已被确定。
\end{solution}

\begin{exercise}\label{g1:ex:voronoi-similarity}
	设 \(S\subset\mathbb R^d\) 为非空有限互异点集，\(f(x)=aUx+b\)，其中 \(a>0\)、\(U^{\mathsf T}U=I\)。证明 \(f\) 把 \(S\) 的 Voronoi 胞和 Delaunay 单元分别映为 \(f(S)\) 的对应胞与单元。这个结论能否把相似变换扩大为任意可逆仿射变换？用正方形顶点集及剪切 \(g(x,y)=(x+y/2,y)\) 说明。
\end{exercise}
\begin{solution}
	正交性给出 \(\lVert f(x)-f(p)\rVert=a\lVert x-p\rVert\)。由于所有距离乘以同一正数，最短距离的大小顺序和等号都保持，故
	\[
		N_{f(S)}(f(x))=f(N_S(x)),\qquad
		V^{f(S)}_{f(p)}=f(V^S_p).
	\]
	上标在此标明所用站点集。仿射映射保持凸组合，所以
	\[
		\operatorname{conv}N_{f(S)}(f(x))
		=f\bigl(\operatorname{conv}N_S(x)\bigr).
	\]
	又因 \(f\) 双射，右边随 \(x\) 遍历得到全部对应单元，证明所需结论。

	取正方形顶点 \(A=(0,0),B=(1,0),C=(1,1),D=(0,1)\)。原 Delaunay 剖分只有一个二维单元，即正方形。剪切后四点成为
	\[
		A'=(0,0),\quad B'=(1,0),\quad C'=(3/2,1),\quad D'=(1/2,1).
	\]
	圆 \(A'B'D'\) 的方程为 \(x^2+y^2-x-3y/4=0\)，\(C'\) 代入左边为 \(1>0\)，所以三角形 \(A'B'D'\) 有严格空圆。平行四边形关于中心的对称把它映为三角形 \(B'C'D'\)，并保持距离，所以后者也有严格空圆。两个三角形沿 \(B'D'\) 拼成整个平行四边形，构成其 Delaunay 剖分。原来的一个四边形单元因剪切变成两个三角形，故任意仿射变换并不保持这种欧氏构造。剪切保持直线和凸组合，却一般改变距离比较，这是结论失效的原因。
\end{solution}

\begin{exercise}\label{g1:ex:hausdorff-affine-stability}
	设 \(K,L\subseteq\mathbb R^n\) 为非空紧凸集，\(f(x)=Tx+a\) 为到 \(\mathbb R^m\) 的仿射映射。令 \(\lVert T\rVert_{\mathrm{op}}=\max_{\lVert v\rVert=1}\lVert Tv\rVert\)。证明
	\[
		h_{f(K)}(u)=\langle u,a\rangle+h_K(T^{\mathsf T}u),\qquad
		d_H(f(K),f(L))\le\lVert T\rVert_{\mathrm{op}}d_H(K,L).
	\]
	当 \(m=n\) 且 \(T\) 可逆时，再证明
	\[
		\frac{d_H(K,L)}{\lVert T^{-1}\rVert_{\mathrm{op}}}
		\le d_H(f(K),f(L)).
	\]
	说明当 \(T\) 不单射时，一般不存在只依赖 \(T\) 的正常数作为这种下界的系数。
\end{exercise}
\begin{solution}
	对 \(x\in K\) 展开 \(\langle u,Tx+a\rangle=\langle T^{\mathsf T}u,x\rangle+\langle u,a\rangle\)，取最大值便得支撑函数等式。令 \(\varepsilon=d_H(K,L)\)。对每个 \(x\in K\)，取 \(y\in L\) 满足 \(\lVert x-y\rVert\le\varepsilon\)，则
	\[
		\lVert f(x)-f(y)\rVert=\lVert T(x-y)\rVert
		\le\lVert T\rVert_{\mathrm{op}}\varepsilon.
	\]
	反向从 \(L\) 的点到 \(K\) 同理，因此两个方向的最大距离都受同一个界控制，得到所求上界。这里允许 \(T=0\)，此时两个像均为同一个单点。

	若 \(T\) 可逆，把刚证得的不等式应用于 \(f^{-1}(z)=T^{-1}(z-a)\) 及 \(f(K),f(L)\)，得到
	\[
		d_H(K,L)\le\lVert T^{-1}\rVert_{\mathrm{op}}d_H(f(K),f(L)),
	\]
	移项即得下界。若 \(T\) 不单射，取非零 \(v\in\ker T\)，令 \(K=\{0\}\)、\(L=\{v\}\)。它们的原距离为 \(\lVert v\rVert>0\)，两个像却均为 \(\{a\}\)，像的距离为零。因此不存在对所有非空紧凸集都成立的正比例下界。
\end{solution}

\begin{exercise}\label{g1:ex:perturbed-vertex-hulls}
	给定两组点 \(v_1,\ldots,v_N\) 与 \(w_1,\ldots,w_N\) 在 \(\mathbb R^n\) 中，\(N\ge1\)。设 \(\varepsilon\ge0\)，且 \(\lVert v_j-w_j\rVert\le\varepsilon\) 对所有 \(j\) 成立。证明两个凸包的 Hausdorff 距离不超过 \(\varepsilon\)，并证明常数 \(1\) 不能普遍减小。若两组点的数目不相同，但每个第一组点到某个第二组点的距离不超过 \(\varepsilon\)，反向也成立，结论是否仍然成立？最后说明两个凸包的距离为零，并不保证给定编号下的对应点相同。
\end{exercise}
\begin{solution}
	记两个凸包为 \(K,L\)。对任意 \(x=\sum_j\lambda_jv_j\in K\)，其中 \(\lambda_j\ge0\)、\(\sum_j\lambda_j=1\)，取 \(y=\sum_j\lambda_jw_j\in L\)。三角不等式给出
	\[
		\lVert x-y\rVert
		\le\sum_j\lambda_j\lVert v_j-w_j\rVert\le\varepsilon.
	\]
	所以 \(K\subseteq L+\varepsilon\mathbb B\)，反向包含完全相同，由式\eqref{g1:eq:hausdorff-inclusions}得到结论。取各只有一个点的两组数据 \(v_1=0\)、\(w_1=\varepsilon e_1\)，其中 \(\varepsilon>0\)，两凸包的距离恰为 \(\varepsilon\)，因此不能减小常数。

	数目不相同的情形也成立。对第一组的每个点 \(v_j\)，选一个距它不超过 \(\varepsilon\) 的第二组点 \(w_{k(j)}\)。使用同一凸组合系数，\(\sum_j\lambda_jw_{k(j)}\) 仍在第二组的凸包中，即使其中有重复点也没有影响。于是得到一个方向的包含；反向重新选择邻近点，再得到另一个方向。

	Hausdorff 距离比较的是集合，没有记录输入名单的编号。比如 \(v_1=-1,v_2=1\) 与 \(w_1=1,w_2=-1\) 在 \(\mathbb R\) 中给出相同线段 \([-1,1]\)，凸包距离为零，两个指定编号的点差却都为 \(2\)。点数据的对应误差是一个方便的充分证书，不是集合距离的必要刻画。
\end{solution}

\begin{exercise}\label{g1:ex:planar-polytope-distance-certificate}
	在 \(\mathbb R^2\) 中取
	\[
		P=\{(x,y):x\ge0,\ y\ge0,\ 2x+y\le2,\ x+2y\le2\},
		\qquad T=\operatorname{conv}\{(0,0),(1,0),(0,1)\}.
	\]
	列出并核验 \(P\) 的全部顶点，证明 \(T\subseteq P\)。对单位方向 \(u=(1,1)/\sqrt2\)，给出计算 \(h_P(u)\) 的可行点和非负乘子证书。进而求 \(d_H(P,T)\)，并给出与所得数值相等的上下界证据。
\end{exercise}
\begin{solution}
	零点满足全部约束，且由非负性有 \(0\le x\le1\)、\(0\le y\le1\)，故 \(P\) 非空有界。枚举两条独立活动边界的交点并核验余下约束，得到
	\[
		(0,0),\quad (1,0),\quad (2/3,2/3),\quad (0,1).
	\]
	另两组候选交点 \((0,2)\)、\((2,0)\) 分别违反另一条斜边约束；四个法向方向没有平行对，因而这六种组合已包括全部候选。活动法向量判据与\cref{g1:thm:polytope-hv}说明上面四点恰为全部顶点，且其凸包就是 \(P\)。\(T\) 的三个生成点都满足 \(P\) 的不等式，故 \(T\subseteq P\)。

	按 \(-x\le0,-y\le0,2x+y\le2,x+2y\le2\) 排列矩阵 \(A\) 和右端 \(b\)。取
	\[
		p=(2/3,2/3),\qquad
		\lambda=\left(0,0,\frac1{3\sqrt2},\frac1{3\sqrt2}\right).
	\]
	它们满足 \(p\in P\)、\(\lambda\ge0\)、\(A^{\mathsf T}\lambda=u\)，且
	\[
		\langle u,p\rangle=\frac4{3\sqrt2}=\langle\lambda,b\rangle.
	\]
	因此 \(h_P(u)=4/(3\sqrt2)\)。而 \(h_T(u)=1/\sqrt2\)，支撑差值给出下界 \(d_H(P,T)\ge1/(3\sqrt2)\)。

	为得到上界，三个属于 \(T\) 的顶点无需移动；唯一额外顶点 \(p\) 距点 \(q=(1/2,1/2)\in T\) 为
	\[
		\lVert p-q\rVert=\sqrt{(1/6)^2+(1/6)^2}=\frac1{3\sqrt2}.
	\]
	令 \(\varepsilon=1/(3\sqrt2)\)，全部四个顶点都在凸集 \(T+\varepsilon\mathbb B\) 中，所以它们的凸包 \(P\) 也在其中。结合 \(T\subseteq P\)，有 \(d_H(P,T)\le\varepsilon\)。两界相等，故所求距离恰为 \(1/(3\sqrt2)\)。这份核验同时给出了全部顶点、支撑值上下界及全局距离上界的有限数据。
\end{solution}

\begin{exercise}\label{g1:ex:inexact-outer-support-data}
	设 \(K\subseteq R\mathbb B\subseteq\mathbb R^n\) 为非空紧凸集，\(R>0\)。取 \(0<\delta<1\) 的有限单位方向网 \(\{u_i\}_{i=1}^m\)，并设给定数值满足
	\[
		h_K(u_i)\le s_i\le h_K(u_i)+\tau\qquad(1\le i\le m),\qquad \tau\ge0.
	\]
	令 \(P_\tau=\{x:\langle u_i,x\rangle\le s_i\ \text{对全部 }i\}\)、\(c_\delta=1-\delta^2/2\)。证明
	\[
		K\subseteq P_\tau\subseteq\frac{R+\tau}{c_\delta}\mathbb B,
		\qquad
		d_H(K,P_\tau)\le
		\tau+\delta\left(R+\frac{R+\tau}{c_\delta}\right).
	\]
	若只知道 \(|s_i-h_K(u_i)|\le\tau\)，直接使用 \(s_i\) 能否保证外包含？说明如何修正右端以恢复外包含，并给出修正后的误差界。
\end{exercise}
\begin{solution}
	有效的支撑上界给出 \(K\subseteq P_\tau\)，所以 \(P_\tau\) 非空。对 \(x\in P_\tau\)、\(x\ne0\)，取距 \(x/\lVert x\rVert\) 不超过 \(\delta\) 的 \(u_i\)，则
	\[
		c_\delta\lVert x\rVert\le\langle u_i,x\rangle
		\le s_i\le h_K(u_i)+\tau\le R+\tau.
	\]
	因此 \(P_\tau\) 有界，也是闭凸集。令 \(M=(R+\tau)/c_\delta\)。对其外于 \(K\) 的一点 \(x\)，取到 \(K\) 的最近点所给的单位方向 \(v\)，再取邻近的 \(u_i\)，有
	\begin{align*}
		d(x,K)
		&=\langle v,x\rangle-h_K(v)\\
		&\le\langle u_i,x\rangle-h_K(u_i)
		+\delta\lVert x\rVert+R\delta\\
		&\le\tau+\delta(M+R).
	\end{align*}
	\(x\in K\) 时距离为零，故全部 \(P_\tau\) 到 \(K\) 的距离受此界控制；反向由 \(K\subseteq P_\tau\) 为零，得到所求Hausdorff界。

	只有双侧数值误差时，某个 \(s_i\) 可能小于真实支撑值，因而切掉 \(K\) 的一部分。例如单位圆盘在方向 \((1,0)\) 的支撑值为 \(1\)，使用 \(s_i=1-\tau\)、\(0<\tau<1\)，就排除了圆盘点 \((1,0)\)。应改用 \(\widetilde s_i=s_i+\tau\)。此时
	\[
		h_K(u_i)\le\widetilde s_i\le h_K(u_i)+2\tau,
	\]
	所以对修正后的半空间交 \(\widetilde P\)，前面的结论给出
	\[
		K\subseteq\widetilde P,\qquad
		d_H(K,\widetilde P)\le2\tau+
		\delta\left(R+\frac{R+2\tau}{c_\delta}\right).
	\]
	这个修正使用的是误差的已证上界；仅凭最后几位小数看起来稳定，不能据此断言这些不等式。
\end{solution}


% !TeX root = ../../Book1.tex
% 纲要标题：## 附录 C　几何计算与图示
% 纲要位置：几何学计划纲要.md，第 966 行。
% 正文状态：初稿；尚未正式审核。

% 纲要内容开始
% 
% > 图示服务于局部图、变换和不变量的核验，不以软件输出代替证明。
% 
% * **坐标、符号与可验证计算**
%   * 图册覆盖范围和坐标奇点
%   * 联络、曲率及 Bianchi 核验
%   * 参数模型、精确恒等式与误差界
% 
% * **数值误差与几何解释**
% 
%   区分几何退化、坐标退化和离散化误差；记录度量、法向、取向、步长与精度，使计算能够独立复核。
% 
% ---
% 
% 纲要内容结束

\chapter{几何计算与图示}
\label{g1:app:C}
\BookChapterNumber{g1:app:C}

\begin{chapterabstract}
	同一个几何对象可以在一张坐标图中表现为规则的公式，在另一张图中出现发散系数或退化矩阵。计算的任务首先是辨认这些表示的定义域与变换规律，再判断数值结果究竟说明了什么。本附录从球面的经纬参数和立体投影图出发，核对覆盖、取向与度量；随后计算联络、曲率及 Bianchi 恒等式，并用旋转椭球展示参数条件、精确恒等式和法向误差界。最后通过中心差分的误差分析与球面曲率实验，区分几何退化、坐标失效、截断误差及输入误差，使图示与数值能够服务于可复核的几何论证。
\end{chapterabstract}

\begin{learninggoals}
	\begin{itemize}
		\item 写出球面坐标图的准确覆盖范围、逆映射和过渡映射，识别经纬参数在极点的失效。
		\item 在坐标或标架变化后核对度量、法向、联络及曲率的变换，而不直接比较不具不变意义的系数。
		\item 展开一个非交换矩阵值联络的曲率与 Bianchi 核验，说明恒等式核验的适用范围。
		\item 对旋转椭球计算基本形式和曲率，给出参数条件及有限代数证书。
		\item 从输入误差和正则性下界推导法向与差分误差，说明步长和精度为何需要共同选择。
		\item 记录数值实验的对象、度量、取向、离散公式及误差依据，区分有限观测和已证结论。
	\end{itemize}
\end{learninggoals}

单位球面在经纬参数中的度量为 \(\mathrm{d}\theta^2+\sin^2\theta\,\mathrm{d}\phi^2\)。当 \(\theta\) 趋近零时，一个系数趋于零，逆矩阵中的相应系数趋于无穷；若只绘制这些系数，很容易把北极误认为曲面的奇点。然而北极附近可以写成光滑图像 \(z=\sqrt{1-x^2-y^2}\)，球面的曲率在那里仍然等于 \(1\)。两个公式描述的是同一个几何对象，出现差别的原因必须在坐标及其适用范围中寻找。

\ProseParagraphBreak

计算中还会出现更隐蔽的问题：一个曲率表达式在许多测试点上都给出正确数值，仍可能在遗漏的参数范围内失效；将一个恒等式的两边相减后得到零，也可能只是因为受维数限制，其中的高次微分形式都为零。因而需要把三种论据分清。精确推导说明公式为何成立，数值实验揭示特定输入下的表现，误差界则把有限精度的数据与连续几何量联系起来。图示呈现覆盖范围和局部模型。

\ProseParagraphBreak

本附录使用\secref{g1:sec:0801}的正则参数曲面与第一基本形式、\secref{g1:sec:0802}的形算子及曲率公式，以及\chapref{g1:ch:11}的坐标图与光滑过渡。联络计算另需\chapref{g1:ch:14}的外微分、楔积和\chapref{g1:ch:15}的局部联络矩阵；下文会重述所用约定。尚未学习联络的读者可以先读图册、参数模型和数值实验，再回到中间的小节。所选例子都给出可直接重算的公式，不要求使用某一种计算软件。

\section{坐标、符号与可验证计算}
\label{g1:sec:C01}

要比较两次计算，先要知道两次计算是否表示同一个点、同一个切向量及同一个法向。坐标图给出定义域和点的对应，度量与标架决定分量的意义，参数条件则保证求逆和归一化合法。本节把这些问题放进具体模型中：先计算球面图册，再处理联络中的非交换乘法，最后把一个曲面族的精确公式转化为带条件的误差估计。

\subsection{图册覆盖范围和坐标奇点}
\label{g1:subsec:C01:01}

球面的经线在两极汇合。把经纬网格画得更密，并不能在极点区分经度方向；但极点附近的球面仍是一块光滑曲面。计算中看到分母趋零、度量矩阵退化或许多参数点具有同一个像时，首先要查的是当前参数表示的定义域与逆映射。只有明确了这些范围，公式才能用于判断几何对象。

\ProseParagraphBreak

回顾\cref{g1:def:smooth-manifold}：坐标图是流形中一个开集 \(U\) 到Euclid 开集 \(D\) 的同胚 \(\varphi:U\rightarrow D\)，而计算曲面时写出的参数通常是相反方向的 \(\Phi=\varphi^{-1}:D\rightarrow U\)。因此，给出 \(\Phi\) 的表达式后，仍须确定 \(D\)、像 \(U\)、单射性及逆映射的光滑性。若要用若干张图描述整个对象，还须核对这些 \(U\) 的并是否覆盖它。

\begin{proposition}[由参数取得坐标图]
	\label{g1:prop:parameter-chart-test}
	设 \(M\) 是已给定光滑结构的二维无边界流形，\(D\subset\mathbb R^2\) 是开集，\(\Phi:D\rightarrow M\) 光滑。若 \(\Phi\) 单射，且对每个 \(a\in D\)，\(\mathrm{d}\Phi_a:\mathbb R^2\rightarrow\mathrm{T}_{\Phi(a)}M\) 的秩均为 \(2\)，则 \(U=\Phi(D)\) 是 \(M\) 的开集，\(\Phi:D\rightarrow U\) 是微分同胚，\((U,\Phi^{-1})\) 是一张光滑坐标图。
\end{proposition}
\begin{proof}
	源、目标的维数都为 \(2\)，所以满秩微分是线性同构。由\cref{g1:thm:inverse-function}，每个 \(a\in D\) 都有一个开邻域 \(D_a\)，使 \(\Phi\) 把它微分同胚地映到 \(M\) 的一个开集 \(U_a\)。因此 \(U=\bigcup_{a\in D}U_a\) 开。全局单射性保证各局部逆在交叠处给出同一个原像，因而粘成 \(U\) 上的逆映射。该逆映射在每个 \(U_a\) 上光滑，所以在 \(U\) 上光滑。
\end{proof}

满秩保证局部可逆，单射性进一步排除不同参数分支映到同一点；两者结合，便在像 \(U\) 上得到一张图。它覆盖的范围由参数像决定，可能小于整个 \(M\)。同一球面参数可以同时展示这三项条件的作用。

\ProseParagraphBreak

令 \(S^2=\{(x,y,z):x^2+y^2+z^2=1\}\)，赋予\cref{g1:exm:sphere-smooth-atlas}中的标准光滑结构。取余纬 \(\theta\) 和经度 \(\phi\)，定义
\begin{equation}
	\label{g1:eq:spherical-chart-parameters}
	\Phi(\theta,\phi)
	=(\sin\theta\cos\phi,\sin\theta\sin\phi,\cos\theta),
	\qquad 0<\theta<\pi,\quad-\pi<\phi<\pi.
\end{equation}
这里纬度等于 \(\pi/2-\theta\)；采用余纬是为了使北极对应 \(\theta=0\)。经度区间的两端不包括在内，是为了避免 \(-\pi\) 与 \(\pi\) 表示同一条经线。

\begin{proposition}[经纬图的精确范围]
	\label{g1:prop:spherical-chart-domain}
	设 \(S^2\subset\mathbb R^3\) 为带标准光滑结构及诱导度量的单位球面，定义
	\[
		\Phi(\theta,\phi)=(\sin\theta\cos\phi,\sin\theta\sin\phi,\cos\theta),
		\qquad (\theta,\phi)\in(0,\pi)\times(-\pi,\pi).
	\]
	令
	\[
		\Gamma=\{(x,y,z)\in S^2:y=0,\ x\le0\},\qquad
		U=S^2\setminus\Gamma.
	\]
	则 \(\Phi\) 把开矩形 \((0,\pi)\times(-\pi,\pi)\) 微分同胚地映到 \(U\)。其诱导度量为
	\begin{equation}
		\label{g1:eq:spherical-chart-metric}
		g=\mathrm{d}\theta^2+\sin^2\theta\,\mathrm{d}\phi^2.
	\end{equation}
	若把同一三角函数表达式延伸到 \(\theta=0\) 或 \(\theta=\pi\)，其微分在这些参数处的秩降为 \(1\)。
\end{proposition}
\begin{proof}
	集合 \(\Gamma\) 是一条包含两极的闭半经线，所以 \(U\) 在球面中开。对 \(p=(x,y,z)\in U\)，令 \(\rho=\sqrt{x^2+y^2}\)。此时 \(\rho>0\)，并且 \(\rho+x>0\)：等号只能发生在 \(y=0,x\le0\) 的被删去部分。于是可定义
	\[
		\theta=\arccos z,\qquad
		\phi=2\arctan\frac{y}{\rho+x}.
	\]
	第一项落在 \((0,\pi)\)，第二项落在 \((-\pi,\pi)\)。半角公式给出 \(\cos\phi=x/\rho\)、\(\sin\phi=y/\rho\)，而 \(\sin\theta=\rho\)，所以代回得到 \(\Phi(\theta,\phi)=p\)。反过来，原参数区间中 \(\cos\theta\) 唯一确定 \(\theta\)，经度在这个长度为 \(2\pi\) 的开区间内也唯一，故所给反解确为逆映射。\(\arccos\) 在 \((-1,1)\) 上光滑，平方根和分式在所列范围内光滑，因此逆映射光滑，原映射也光滑。

	直接求导得
	\[
		\begin{aligned}
			\Phi_\theta&=(\cos\theta\cos\phi,\cos\theta\sin\phi,-\sin\theta),\\
			\Phi_\phi&=(-\sin\theta\sin\phi,\sin\theta\cos\phi,0).
		\end{aligned}
	\]
	两个向量正交，长度分别为 \(1\) 与 \(\sin\theta\)，故在开矩形内线性无关，第一基本形式就是\eqref{g1:eq:spherical-chart-metric}。在 \(\theta=0,\pi\) 处，\(\Phi_\phi=0\)，而 \(\Phi_\theta\) 仍为单位向量，所以延伸后的参数映射秩为 \(1\)。
\end{proof}

这张图漏掉的不只是两个极点，还漏掉整条选作经度割线的半经线。在 \(0<\theta<\pi\)、\(\phi=\pm\pi\) 处，三角函数表达的微分仍满秩；问题是两端重复表示同一点，因而不能把包含两端的矩形直接当作一张Euclid 开集坐标图。把割线移到另一条经线，可以覆盖这些非极点。两极则不同：所有经度在同一极点汇合，必须改用另一种参数。

\ProseParagraphBreak

度量矩阵 \(\operatorname{diag}(1,\sin^2\theta)\) 在 \(\theta=0\) 的退化，恰好来自 \(\Phi_\phi=0\)。它是在一个已经不再浸入的映射下拉回度量所得的退化形式，不能当作极点处的有效坐标矩阵。我们把这类由表示方式造成的失效称为\term[zuobiaoqidian]{坐标奇点}[coordinate singularity]；下面换图后，极点处的度量仍是正定的。

\ProseParagraphBreak

取北极 \(N=(0,0,1)\)、南极 \(S=(0,0,-1)\)。\cref{g1:ex:stereographic-smooth-charts}已经给出一般球面的立体投影；这里重述二维公式，并把覆盖范围和度量计算连在一起。以下按投影中心命名两张图，记
\[
	U_N=S^2\setminus\{N\},\qquad U_S=S^2\setminus\{S\},
\]
并定义
\[
	\sigma_N(x,y,z)=\frac{(x,y)}{1-z},\qquad
	\sigma_S(x,y,z)=\frac{(x,y)}{1+z}.
\]
例如，从 \(N\) 到 \(p\ne N\) 的直线与平面 \(z=0\) 相交，交点的前两坐标就是 \(\sigma_N(p)\)。这张图删除投影中心 \(N\)，却完整包含 \(S\)。

\begin{proposition}[两张立体投影图及其度量]
	\label{g1:prop:stereographic-chart-metric}
	设 \(S^2\subset\mathbb R^3\) 为标准光滑单位球面，\(N=(0,0,1)\)、\(S=(0,0,-1)\)，\(U_N=S^2\setminus\{N\}\)、\(U_S=S^2\setminus\{S\}\)。定义 \(\sigma_N(x,y,z)=(x,y)/(1-z)\)、\(\sigma_S(x,y,z)=(x,y)/(1+z)\)。则 \(\sigma_N:U_N\to\mathbb R^2\)、\(\sigma_S:U_S\to\mathbb R^2\) 是光滑坐标图，并且 \(U_N\cup U_S=S^2\)。对 \(u=(u_1,u_2)\)，它们的逆映射分别为
	\begin{align}
		\Phi_N(u)&=\left(\frac{2u_1}{1+|u|^2},\frac{2u_2}{1+|u|^2},\frac{|u|^2-1}{1+|u|^2}\right),
		\label{g1:eq:stereographic-north-inverse}\\
		\Phi_S(u)&=\left(\frac{2u_1}{1+|u|^2},\frac{2u_2}{1+|u|^2},\frac{1-|u|^2}{1+|u|^2}\right).
		\label{g1:eq:stereographic-south-inverse}
	\end{align}
	在交叠区域中，过渡映射为
	\begin{equation}
		\label{g1:eq:stereographic-real-transition}
		\tau(u)=(\sigma_S\circ\Phi_N)(u)=\frac{u}{|u|^2},
		\qquad u\in\mathbb R^2\setminus\{0\},
	\end{equation}
	且 \(\tau^{-1}=\tau\)。单位球面的诱导度量在两张图中均为
	\begin{equation}
		\label{g1:eq:stereographic-sphere-metric}
		g=\frac{4}{(1+u_1^2+u_2^2)^2}
		\bigl(\mathrm{d}u_1^2+\mathrm{d}u_2^2\bigr).
	\end{equation}
\end{proposition}
\begin{proof}
	先解 \(u=(x,y)/(1-z)\)。球面方程给出
	\[
		|u|^2(1-z)^2=1-z^2=(1-z)(1+z).
	\]
	在 \(U_N\) 上 \(1-z\ne0\)，约去后得到 \(z=(|u|^2-1)/(1+|u|^2)\)，再代回 \((x,y)=(1-z)u\)，便得\eqref{g1:eq:stereographic-north-inverse}。该式对所有 \(u\in\mathbb R^2\) 有定义，像的第三坐标严格小于 \(1\)，直接代回两式得到两个复合恒等映射。将 \(z\) 换成 \(-z\) 得到另一张图及其逆。两图的坐标函数都是在分母非零处定义的光滑函数；逆式是分母始终为正的有理式。与第11章半球图比较时，只需在这些表达式中删除一个坐标，或把半球图的光滑平方根反解代入其中，所以它们与标准球面图册光滑相容。

	两个删去的点不同，故两图覆盖球面。在 \(\Phi_N\) 中，\(u=0\) 对应 \(S\)，恰是另一张图所删除的点。因此交叠的参数域是 \(\mathbb R^2\setminus\{0\}\)。将 \(\Phi_N\) 代入 \(\sigma_S\)，得到 \(u/|u|^2\)；再次应用同一个式子恢复 \(u\)，所以正逆过渡均光滑。

	最后计算度量。令 \(q=1+|u|^2\)，对参数向量 \(h\in\mathbb R^2\)，记 \(a=\langle u,h\rangle\)。微分为
	\[
		\mathrm{d}\Phi_N(u)[h]
		=\left(\frac{2h}{q}-\frac{4au}{q^2},\frac{4a}{q^2}\right).
	\]
	因此
	\[
		\left|\mathrm{d}\Phi_N(u)[h]\right|^2
		=\frac{4|h|^2}{q^2}-\frac{16a^2}{q^3}
		+\frac{16a^2(|u|^2+1)}{q^4}
		=\frac{4|h|^2}{q^2}.
	\]
	这给出\eqref{g1:eq:stereographic-sphere-metric}。\(\Phi_S\) 的最后一个微分分量只差负号，平方范数不变，所以它给出相同形式的度量矩阵。
\end{proof}

\begin{figure}[htbp]
	\centering
	% !TeX root = ../Book1.tex
\begin{tikzpicture}[>=stealth,font=\small,x=2.1cm,y=2.1cm]
	\draw[gray,->] (-1.25,0)--(2.55,0) node[right] {\(z=0\)};
	\draw[gray,dashed] (0,-1.2)--(0,1.2);
	\draw[thick] (0,0) circle (1);
	\draw[blue!70!black,thick,->] (0,1)--(2,0);
	\draw[orange!80!black,thick] (0,-1)--(.8,.6);
	\draw[orange!80!black,->] (0,-1)--(.5,0);
	\fill (0,1) circle (1.5pt) node[above left] {\(N\)};
	\fill (0,-1) circle (1.5pt) node[below left] {\(S\)};
	\fill (.8,.6) circle (1.8pt) node[above right] {\(p\)};
	\fill[blue!70!black] (2,0) circle (1.6pt) node[below] {\(2\)};
	\fill[orange!80!black] (.5,0) circle (1.6pt) node[below right] {\(1/2\)};
\end{tikzpicture}

	\caption[两张立体投影图的覆盖]{在 \(y=0\) 的子午截面中，点 \(p=(4/5,0,3/5)\) 从北极、南极向赤道平面投影，所得横坐标分别为 \(2\) 与 \(1/2\)。两条直线显示投影中心与交点的关系；靠近北极时，南极投影仍给出接近零的有限坐标。}
	\label{g1:fig:stereographic-chart-coverage}
\end{figure}

两张图中的零坐标代表不同的点：\(\Phi_N(0)=S\)，\(\Phi_S(0)=N\)。在这两个点，度量矩阵都是 \(4I_2\)，正定而无退化。过渡式在 \(u=0\) 未定义，原因是对应点不属于两图的交集。

\ProseParagraphBreak

还可以把经纬表达与这两张图直接对照。在 \(0<\theta<\pi\) 内，半角公式给出
\[
	\begin{aligned}
		\sigma_N(\Phi(\theta,\phi))
		&=\cot\frac\theta2\,(\cos\phi,\sin\phi),\\
		\sigma_S(\Phi(\theta,\phi))
		&=\tan\frac\theta2\,(\cos\phi,\sin\phi).
	\end{aligned}
\]
当 \(\theta\to0\) 时，第一组坐标的模趋于无穷，而第二组趋于 \(0\)。球面上的点只是趋向北极，并没有在球面上趋向无穷远；发散的是一张删去北极的图中的坐标。与此同时，第二组是平面中的极坐标形式，\(\tan(\theta/2)=0\) 时所有经度给出同一个平面原点，这又解释了经纬参数的秩为何下降。

\ProseParagraphBreak

实际需要有限范围内的坐标时，可以在 \(z\ge0\) 采用 \(\sigma_S\)，在 \(z\le0\) 采用 \(\sigma_N\)。前一情形有
\[
	|\sigma_S(p)|^2=\frac{x^2+y^2}{(1+z)^2}
	=\frac{1-z}{1+z}\le1,
\]
后一情形同理有 \(|\sigma_N(p)|\le1\)。这些半球规定选图的规则，各张图仍使用原来的开图域。在选出的坐标值处，\eqref{g1:eq:stereographic-sphere-metric}满足
\[
	|h|^2\le g_u(h,h)\le4|h|^2.
\]
因此两张图足以把整个球面的计算安排在坐标模不超过 \(1\) 的区域内；在赤道附近二者都可使用，并由过渡式比较结果。

\begin{example}[参数秩下降而像集光滑]
	\label{g1:exm:cubic-plane-parametrization}
	映射
	\[
		F:\mathbb R^2\longrightarrow\mathbb R^3,\qquad
		F(s,t)=(s^3,t,0)
	\]
	是到整个平面 \(P=\{z=0\}\) 的连续双射，逆映射为 \((x,y,0)\mapsto(\sqrt[3]{x},y)\)，所以它还是同胚。但其微分的两列为 \((3s^2,0,0)\)、\((0,1,0)\)，在 \(s=0\) 时秩为 \(1\)，拉回Euclid 度量为
	\[
		F^*g=9s^4\,\mathrm{d}s^2+\mathrm{d}t^2.
	\]
	参数映射的微分将这个退化方向映成了零向量。平面本身仍有处处正则的参数 \((x,y)\mapsto(x,y,0)\)；问题在于原来的逆含有在零点不可微的立方根。因而 \((s,t)\) 不能作为标准光滑平面在该直线附近的光滑坐标。这把\cref{g1:exm:cubic-rank-change}中的一维现象放到了一个像集完全明确的曲面上。
\end{example}

要判断参数公式给出的对象，可以先确定其像，再寻找该像附近有效的坐标或正则定义方程。对球面，方程 \(H(x,y,z)=x^2+y^2+z^2-1=0\) 的微分在每个球面点都非零，\cref{g1:thm:regular-level-set}给出二维光滑水平集；对上例平面，定义方程 \(z=0\) 同样处处正则。因此，可以分别检查参数映射是否降秩，以及像集本身是否光滑。

\subsection{联络、曲率及 Bianchi 核验}
\label{g1:subsec:C01:02}

一份联络计算中，矩阵的行列次序、微分形式的楔积次序和坐标标架的选择会同时影响结果。例如，平直Euclid 平面在极坐标正交标架中已有非零联络矩阵；另一方面，一组常值矩阵也可能由于不交换而产生非零曲率。因此，核验应从“这组矩阵表示哪个标架中的哪一个联络”开始，再比较曲率定义、结构方程和恒等式。本小节把这些关系放进可逐项重算的模型中。

\ProseParagraphBreak

先固定\cref{g1:prop:connection-local-data}的约定。对秩 \(m\) 的实光滑向量丛，局部标架写为行 \(e=(e_1,\ldots,e_m)\)，截面写为 \(s=eu\)，其中 \(u\) 是列分量。联络矩阵 \(A\) 的第 \(j\) 列记录 \(\nabla e_j\)，即
\[
	\nabla e=eA,\qquad \nabla(eu)=e(\mathrm{d}u+Au).
\]
若 \(e'=eg\)，则同一截面的分量满足 \(u=gu'\)。对 \(eg\) 使用乘积法则，
\[
	\nabla(eg)=e(Ag+\mathrm{d}g)
	=e'g^{-1}(Ag+\mathrm{d}g),
\]
所以新矩阵为
\[
	A'=g^{-1}Ag+g^{-1}\mathrm{d}g.
\]
这里的附加项来自标架本身随位置变化。只有比较同一个联络时，才可以把两组局部矩阵放进这个公式。

\ProseParagraphBreak

矩阵值形式的乘积同时保留矩阵乘法与楔积。例如
\[
	(A\wedge A)^a{}_b=\sum_c A^a{}_c\wedge A^c{}_b.
\]
若在坐标中写 \(A=\sum_i A_i\,\mathrm{d}x^i\)，则把 \((i,j)\) 与 \((j,i)\) 两项合并，得到
\[
	A_iA_j\,\mathrm{d}x^i\wedge\mathrm{d}x^j
	+A_jA_i\,\mathrm{d}x^j\wedge\mathrm{d}x^i
	=[A_i,A_j]\,\mathrm{d}x^i\wedge\mathrm{d}x^j.
\]
因此\cref{g1:thm:curvature-structure-equation}给出的曲率具体为
\begin{equation}\label{g1:eq:connection-check-curvature-components}
	F=\mathrm{d}A+A\wedge A
	=\sum_{i<j}F_{ij}\,\mathrm{d}x^i\wedge\mathrm{d}x^j,
	\qquad F_{ij}=\partial_iA_j-\partial_jA_i+[A_i,A_j].
\end{equation}
这里对 \(i<j\) 求和，不另乘 \(1/2\)。并且 \(R(X,Y)(eu)=eF(X,Y)u\)，其中
\[
	R(X,Y)s=\nabla_X\nabla_Ys-\nabla_Y\nabla_Xs-\nabla_{[X,Y]}s.
\]
这个作用公式使矩阵中的正负号最终可以用一个实际的纤维向量来校准。

\ProseParagraphBreak

在曲率二形式上，本小节将自同态丛的外协变微分简记为
\[
	\mathrm{D}F=\mathrm{d}F+A\wedge F-F\wedge A.
\]
最后一项的负号来自\eqref{g1:eq:endomorphism-exterior-derivative}中二次形式的偶数次数。\cref{g1:thm:bianchi-identity}已证明，对由 \(A\) 计算出的曲率总有 \(\mathrm{D}F=0\)。在三维坐标 \((x,y,z)\) 中，记 \(\nu=\mathrm{d}x\wedge\mathrm{d}y\wedge\mathrm{d}z\)，则
\begin{equation}\label{g1:eq:connection-check-bianchi-coefficient}
	\begin{aligned}
		\mathrm{D}F
		={}&\bigl(\partial_xF_{yz}-\partial_yF_{xz}+\partial_zF_{xy}\\
		&\quad+[A_x,F_{yz}]-[A_y,F_{xz}]+[A_z,F_{xy}]\bigr)\nu.
	\end{aligned}
\end{equation}
例如 \(\mathrm{d}y\wedge\mathrm{d}x\wedge\mathrm{d}z=-\nu\)，解释了中间导数项的负号。将 \(F_{zx}=-F_{xz}\) 代入，就恢复第十五章中的循环求和形式。若使用的基底是一般余标架 \(\theta^a\)，则 \(\mathrm{d}\theta^a\) 未必为零，不能直接把它们当成坐标微分代入这些偏导公式。

\begin{example}[极坐标平面的联络与无挠性]\label{g1:exm:connection-check-polar}
	在Euclid 平面的一张极坐标图中取 \(r>0\)，角度 \(\varphi\) 只在该局部图上作为实值函数使用。笛卡儿平行标架记为 \(e^0=(\partial_x,\partial_y)\)，正向正交单位标架为
	\[
		e_1=\partial_r=\cos\varphi\,\partial_x+\sin\varphi\,\partial_y,\qquad
		e_2=r^{-1}\partial_\varphi=-\sin\varphi\,\partial_x+\cos\varphi\,\partial_y.
	\]
	令
	\[
		Q(\varphi)=
		\begin{pmatrix}\cos\varphi&-\sin\varphi\\
		\sin\varphi&\cos\varphi\end{pmatrix},
		\qquad
		J=\begin{pmatrix}0&-1\\1&0\end{pmatrix}.
	\]
	于是 \(e=e^0Q\)。通常Euclid 联络在 \(e^0\) 中的矩阵为零，在 \(e\) 中则为
	\[
		A=Q^{-1}\mathrm{d}Q=J\,\mathrm{d}\varphi
		=\begin{pmatrix}0&-\mathrm{d}\varphi\\
		\mathrm{d}\varphi&0\end{pmatrix}.
	\]
	这一等式也可以逐列检查：\(\mathrm{d}e_1=e_2\,\mathrm{d}\varphi\)、\(\mathrm{d}e_2=-e_1\,\mathrm{d}\varphi\)。矩阵不是零，但它反对称，与正交标架中的度量相容性一致。

	对偶余标架列为 \(\theta=(\mathrm{d}r,r\,\mathrm{d}\varphi)^{\mathsf T}\)。由\cref{g1:prop:first-structure-equation}，挠率二形式列为 \(\Theta=\mathrm{d}\theta+A\wedge\theta\)，而
	\[
		\mathrm{d}\theta=
		\begin{pmatrix}0\\\mathrm{d}r\wedge\mathrm{d}\varphi\end{pmatrix},
		\qquad
		A\wedge\theta=
		\begin{pmatrix}
			-\mathrm{d}\varphi\wedge(r\,\mathrm{d}\varphi)\\
			\mathrm{d}\varphi\wedge\mathrm{d}r
		\end{pmatrix}
		=\begin{pmatrix}0\\-\mathrm{d}r\wedge\mathrm{d}\varphi\end{pmatrix}.
	\]
	所以 \(\Theta=0\)。若只看 \(A\wedge\theta\)，就会把余标架的变化漏掉，错误地得出联络有挠的结论。

	还可以直接对向量场核验。由 \(A(e_1)=0\)、\(A(e_2)=J/r\)，有
	\[
		\nabla_{e_1}e_2=0,\qquad
		\nabla_{e_2}e_1=r^{-1}e_2,\qquad
		[e_1,e_2]=-r^{-1}e_2.
	\]
	故 \(T(e_1,e_2)=0-r^{-1}e_2-(-r^{-1}e_2)=0\)。这一计算与前面的二形式计算核验的是同一张量。最后，\(\mathrm{d}A=0\)，且 \(A\wedge A=J^2\,\mathrm{d}\varphi\wedge\mathrm{d}\varphi=0\)，所以 \(F=0\)。标架在旋转、标架向量不交换，都不意味着Euclid 联络有曲率或挠率。以上计算只在 \(r>0\) 的图中进行，不能把极坐标原点处的表达式当作光滑标架。
\end{example}

平面模型检验换标架项和无挠方程；它的曲率为零，对整体反号不敏感。用同一方法再计算正、负曲率模型，就能确定曲率的符号约定。

\begin{proposition}[一类曲面度量的曲率校准]\label{g1:prop:connection-check-warped-metric}
	设 \(I\subset\mathbb R\) 为开区间，\(f:I\to(0,\infty)\) 光滑。在坐标区域 \(I\times J_0\) 上取度量
	\[
		g_f=\mathrm{d}r^2+f(r)^2\,\mathrm{d}\varphi^2,
	\]
	其中 \(J_0\subset\mathbb R\) 为开区间。对正向正交单位标架
	\[
		e_1=\partial_r,\qquad e_2=f(r)^{-1}\partial_\varphi,
		\qquad \theta^1=\mathrm{d}r,\quad \theta^2=f(r)\,\mathrm{d}\varphi,
	\]
	记 \(J=\begin{psmallmatrix}0&-1\\1&0\end{psmallmatrix}\)。采用行标架 \(e=(e_1,e_2)\) 与列分量 \(u\) 的约定 \(\nabla(eu)=e(\mathrm du+Au)\)，并令 \(F=\mathrm dA+A\wedge A\)。唯一保持该度量且无挠的联络具有联络矩阵和曲率矩阵
	\[
		A=J f'(r)\,\mathrm{d}\varphi,\qquad
		F=J f''(r)\,\mathrm{d}r\wedge\mathrm{d}\varphi.
	\]
	按本书的 \(R(X,Y)=\nabla_X\nabla_Y-\nabla_Y\nabla_X-\nabla_{[X,Y]}\) 约定，Gauss 曲率为
	\[
		K=-\frac{f''(r)}{f(r)}.
	\]
\end{proposition}
\begin{proof}
	先由所需几何性质确定联络，而不预先猜曲率。正交单位标架中的度量矩阵为 \(I\)。对标架内积求导，度量相容要求 \(A^{\mathsf T}+A=0\)，因此 \(A=J\omega\)，其中 \(\omega=a\,\mathrm{d}r+b\,\mathrm{d}\varphi\) 是某个实一形式。这里的反向也成立：若 \(A\) 反对称，对任意两截面分量使用乘积法则，就得到内积的协变求导公式。

	第一结构方程的两行分别为
	\[
		\Theta^1=-\omega\wedge(f\,\mathrm{d}\varphi)
		=-fa\,\mathrm{d}r\wedge\mathrm{d}\varphi,\qquad
		\Theta^2=f'\,\mathrm{d}r\wedge\mathrm{d}\varphi+\omega\wedge\mathrm{d}r
		=(f'-b)\,\mathrm{d}r\wedge\mathrm{d}\varphi.
	\]
	由于 \(f>0\)，无挠性恰好给出 \(a=0\)、\(b=f'\)。这既给出矩阵，也证明了所述唯一性。任意光滑矩阵一形式都按\cref{g1:prop:connection-local-data}定义联络，因此所得矩阵确实实现这些条件。

	现在 \(\omega=f'\,\mathrm{d}\varphi\)，故 \(\mathrm{d}\omega=f''\,\mathrm{d}r\wedge\mathrm{d}\varphi\)，而 \(A\wedge A=J^2\omega\wedge\omega=0\)，得到所列 \(F\)。注意
	\[
		(\mathrm{d}r\wedge\mathrm{d}\varphi)(e_1,e_2)=\frac1f.
	\]
	于是 \(F(e_1,e_2)=(f''/f)J\)。矩阵的第二列为 \((-f''/f,0)^{\mathsf T}\)，所以
	\[
		R(e_1,e_2)e_2=-\frac{f''}{f}e_1,\qquad
		K=g_f(R(e_1,e_2)e_2,e_1)=-\frac{f''}{f}.
	\]
	这一步同时校准了矩阵各列的含义和曲率符号。
\end{proof}

取 \(f(r)=r\)、\(r>0\)，得到平面上的 \(K=0\)。取 \(f(r)=R\sin(r/R)\)、\(0<r<\pi R\)，其中 \(R>0\)，得到半径 \(R\) 球面的 \(K=1/R^2\)。取 \(f(r)=R\sinh(r/R)\)、\(r>0\)，则得到 \(K=-1/R^2\)。三次计算使用同一个公式和同一标架约定。又因正向面积形式为 \(\theta^1\wedge\theta^2=f\,\mathrm{d}r\wedge\mathrm{d}\varphi\)，它们都满足
\[
	\mathrm{d}\omega=-K\,\theta^1\wedge\theta^2,
\]
与第九章的\eqref{g1:eq:surface-curvature-form}一致。这里把面积写成 \(\theta^1\wedge\theta^2\)，以免与联络矩阵 \(A\) 的外微分 \(\mathrm{d}A\) 混淆。

\ProseParagraphBreak

这些曲面模型中的第二 Bianchi 恒等式不能提供一次非平凡的三形式检查。二维底空间上只有两个独立坐标微分，每个三重楔积都包含重复因子，故任何矩阵值三形式都为零。即使误写了曲率二形式的某些系数，\(\mathrm{d}F+A\wedge F-F\wedge A\) 仍可能因次数而自动消失。要实际检验导数项与交换子项的消去，需要至少三个独立底空间方向。

\begin{example}[三维底空间上的非交换消去]\label{g1:exm:connection-check-nonabelian}
	在平凡实秩二丛 \(\mathbb R^3\times\mathbb R^2\) 的标准标架中，以 \((x,y,z)\) 为底空间坐标，取常矩阵
	\[
		B=\begin{pmatrix}0&1\\0&0\end{pmatrix},\qquad
		C=\begin{pmatrix}0&0\\1&0\end{pmatrix},\qquad
		H=\begin{pmatrix}1&0\\0&-1\end{pmatrix}.
	\]
	直接乘法给出 \(BC=\operatorname{diag}(1,0)\)、\(CB=\operatorname{diag}(0,1)\)，并有
	\[
		[B,C]=H,\qquad [H,B]=2B,\qquad [H,C]=-2C.
	\]
	定义联络矩阵
	\[
		A=B\,\mathrm{d}x+xC\,\mathrm{d}y+yH\,\mathrm{d}z.
	\]
	这是一般向量丛的联络，不把这个秩二丛认作三维底空间的切丛，也不预设它保持标准内积。

	先分开计算外微分和二次项：
	\[
		\begin{aligned}
			\mathrm{d}A&=C\,\mathrm{d}x\wedge\mathrm{d}y
			+H\,\mathrm{d}y\wedge\mathrm{d}z,\\
			A\wedge A&=xH\,\mathrm{d}x\wedge\mathrm{d}y
			-2yB\,\mathrm{d}x\wedge\mathrm{d}z
			+2xyC\,\mathrm{d}y\wedge\mathrm{d}z.
		\end{aligned}
	\]
	所以三个有序系数为
	\begin{equation}\label{g1:eq:connection-check-three-curvatures}
		F_{xy}=C+xH,\qquad F_{xz}=-2yB,\qquad F_{yz}=H+2xyC.
	\end{equation}
	例如 \(F_{xy}\) 中的 \(C\) 来自 \(\partial_x(xC)\)，\(xH\) 来自 \([B,xC]\)；若把矩阵当成交换标量，就会漏掉后一项以及另外两个交换子项。

	现在按 \(\nu=\mathrm{d}x\wedge\mathrm{d}y\wedge\mathrm{d}z\) 收集三形式。三个导数项为
	\[
		\partial_xF_{yz}=2yC,\qquad
		-\partial_yF_{xz}=2B,\qquad
		\partial_zF_{xy}=0,
	\]
	因此 \(\mathrm{d}F=2(B+yC)\nu\ne0\)。再分别计算\eqref{g1:eq:connection-check-bianchi-coefficient}中的三个交换子：
	\[
		\begin{aligned}
			[A_x,F_{yz}]&=[B,H+2xyC]=-2B+2xyH,\\
			-[A_y,F_{xz}]&=-[xC,-2yB]=-2xyH,\\
			[A_z,F_{xy}]&=[yH,C+xH]=-2yC.
		\end{aligned}
	\]
	它们的和为 \(-2(B+yC)\)，恰好抵消普通外微分的系数。于是
	\[
		\mathrm{D}F
		=\bigl(2yC+2B-2B+2xyH-2xyH-2yC\bigr)\nu=0.
	\]
	这次等式的两部分都不恒为零，确实使用了三维方向、矩阵不交换性及楔积次序。
\end{example}

这个模型还可以说明，不同的核验方法各能发现哪些问题。若误把曲率写成 \(\widetilde F=\mathrm{d}A\)，则 \(\mathrm{d}\widetilde F=0\)，但对上面的联络，
\[
	\mathrm{D}\widetilde F=-2(B+yC)\nu\ne0,
\]
所以 Bianchi 核验会发现这一次遗漏。然而，由外协变微分的线性可知，对任何固定常数 \(\lambda\)，都有 \(\mathrm{D}(\lambda F)=0\)；特别地，错误候选 \(-F\) 和零形式同样能通过这一项检查。恒等式为零只能说明候选形式满足一条必要关系，不能证明它就是给定 \(A\) 的曲率。

\ProseParagraphBreak

对切丛上的度量 \(h\)，还须核对输入数据与几何性质。若标架的度量矩阵为 \(G=(h(e_i,e_j))\)，度量相容条件是 \(\mathrm{d}G=A^{\mathsf T}G+GA\)，只有在正交单位标架中才简化为 \(A^{\mathsf T}+A=0\)。无挠性则由 \(\mathrm{d}\theta+A\wedge\theta=0\) 检查。换标架时，除了 \(A'\) 的附加微分项，还必须同时变换 \(\theta'=g^{-1}\theta\) 与 \(F'=g^{-1}Fg\)，后一个公式由\cref{g1:prop:curvature-frame-change}保证。这些关系分别确定联络的几何含义、候选曲率是否按定义生成，以及不同标架是否在描述同一个对象。

\subsection{参数模型、精确恒等式与误差界}
\label{g1:subsec:C01:03}

同一段曲面计算通常含有三类判断：参数是否仍给出正则曲面，所得公式是否在整个参数域内成立，以及输入有误差时输出能偏离多少。三者需要不同的依据。下面用一个带形状参数的椭球族贯穿这些判断；计算中始终保留参数范围和法向选择，使每一步都能从原始模型重新得到。

\begin{example}[轴对称椭球的精确计算]\label{g1:exm:ellipsoid-exact-certificate}
	固定 \(a>0\)，考虑
	\[
		X_a(\theta,\phi)
		=(\sin\theta\cos\phi,\sin\theta\sin\phi,a\cos\theta),
		\qquad 0<\theta<\pi.
	\]
	经度 \(\phi\) 可取任意开区间；若要求这张参数片单射，可将其限制在长度不超过 \(2\pi\) 的开区间内。求 \(\theta,\phi\) 的偏导时，\(a\) 是固定的形状参数，不参与求导。若另行令 \(a\) 随 \(\theta\) 变化，得到的已经是另一张曲面，以下公式须重新推导。

	记
	\[
		s=\sin\theta>0,\qquad c=\cos\theta,\qquad
		W=\sqrt{c^2+a^2s^2}>0.
	\]
	直接求导得
	\[
		X_\theta=(c\cos\phi,c\sin\phi,-as),\qquad
		X_\phi=(-s\sin\phi,s\cos\phi,0),
	\]
	从而
	\begin{equation}\label{g1:eq:ellipsoid-area-normal}
		\begin{aligned}
			X_\theta\times X_\phi&=s(as\cos\phi,as\sin\phi,c),\\
			\lVert X_\theta\times X_\phi\rVert&=sW,\\
			n&=\frac{(as\cos\phi,as\sin\phi,c)}{W}.
		\end{aligned}
	\end{equation}
	因为 \(sW>0\)，\cref{g1:def:surface-immersion}中的秩条件在这一区域内成立。这个叉积法向恰为外向法向：对椭球的定义函数
	\[
		F_a(x,y,z)=x^2+y^2+z^2/a^2-1
	\]
	有 \(n=\frac{a}{2W}\nabla F_a(X_a)\)，而系数为正，\(F_a\) 向外增大。这里的取向判断不能只用 \(\lVert n\rVert=1\) 代替，因为 \(-n\) 也满足该等式。

	第一基本形式由切向量的内积给出：
	\[
		E=W^2,\qquad F=0,\qquad G=s^2,\qquad
		I=W^2\,\mathrm d\theta^2+s^2\,\mathrm d\phi^2.
	\]
	特别地，\(D=EG-F^2=s^2W^2\)，与叉积的平方范数一致。按\cref{g1:def:shape-operator}取 \(S=-\mathrm dn\)；\cref{g1:prop:shape-self-adjoint}说明第二基本形式的系数为 \(h_{ij}=\langle X_{ij},n\rangle\)，且在基 \(X_\theta,X_\phi\) 下有 \([S]=g^{-1}h\)。本例所需二阶导数是
	\[
		\begin{aligned}
			X_{\theta\theta}&=(-s\cos\phi,-s\sin\phi,-ac),\\
			X_{\theta\phi}&=(-c\sin\phi,c\cos\phi,0),\\
			X_{\phi\phi}&=(-s\cos\phi,-s\sin\phi,0).
		\end{aligned}
	\]
	分别与 \(n\) 作内积，得到
	\[
		\begin{aligned}
			h_{\theta\theta}&=-\frac aW,\qquad h_{\theta\phi}=0,
			\qquad h_{\phi\phi}=-\frac{as^2}{W},\\
			\mathrm{II}&=-\frac aW\,\mathrm d\theta^2-\frac{as^2}{W}\,\mathrm d\phi^2.
		\end{aligned}
	\]
	例如第一项的分子是 \(-as^2(\cos^2\phi+\sin^2\phi)-ac^2=-a\)，所以它的符号并非来自对凸性的猜测。由于两个基本形式在这组坐标下都为对角形式，形算子与两项主曲率为
	\[
		[S]=\begin{pmatrix}-a/W^3&0\\0&-a/W\end{pmatrix},
		\qquad k_\theta=-\frac a{W^3},\qquad k_\phi=-\frac aW.
	\]
	下标表示对应的坐标方向，没有约定两项按大小排列。由\cref{g1:prop:surface-curvature-formulas}，
	\begin{equation}\label{g1:eq:ellipsoid-exact-curvatures}
		K=\frac{a^2}{W^4},\qquad
		H=-\frac{a(1+W^2)}{2W^3}.
	\end{equation}
	当 \(a=1\) 时，\(W=1\)，得到外向单位球面的 \(S=-\operatorname{Id}\)、\(H=-1\)、\(K=1\)。在赤道 \(\theta=\pi/2\) 上，则有 \(K=1/a^2\)、\(H=-(1+a^2)/(2a^2)\)。这两次代入可帮助发现符号或幂次错误；整个参数域上的结论则由上面的求导和内积运算得到。
\end{example}

两极处 \(s=0\)，原参数的 \(X_\phi\) 为零，不能把曲率公式中约去 \(s^2\) 的步骤直接用于这些参数点。但椭球本身在那里仍然正则。其北极、南极附近分别可用
\[
	Z_a^\pm(x,y)=\bigl(x,y,\pm a\sqrt{1-x^2-y^2}\bigr),
	\qquad x^2+y^2<1
\]
参数化；前两个坐标保证导数的秩为二，且平方根在这个开圆盘上取正支。也可在 \(F_a=0\) 中使用\cref{g1:thm:implicit-function-parameters}：形状参数对应 \(a>0\)，待解变量为 \(z\)，在两极有 \(\partial_zF_a=\pm2/a\ne0\)，故所得图像对 \(a,x,y\) 联合光滑。外向法向 \(\nabla F_a/\lVert\nabla F_a\rVert\) 在整个椭球上光滑，因为 \(\nabla F_a\) 在 \(F_a=0\) 上不可能为零。由这些图像计算的 \(H,K\) 因而连续，式\eqref{g1:eq:ellipsoid-exact-curvatures}在两极的极限给出 \(H=-a\)、\(K=a^2\)。这里先确认了对象及其曲率在极点有定义，再解释原公式的极限。

\ProseParagraphBreak

精确计算若要交给另一人或另一套程序复核，保留最终简化式还不够。例如 \(W\) 中有平方根，\(s,c\) 也不是相互独立的变量。可以暂将 \(a,s,c,p,q,W\) 看作形式变量，其中实际代入时 \(p=\cos\phi\)、\(q=\sin\phi\)，并记录三个关系
\[
	r_1=s^2+c^2-1=0,\qquad r_2=p^2+q^2-1=0,
	\qquad r_3=W^2-c^2-a^2s^2=0.
\]
以下在多项式环 \(\mathbb Q[a,s,c,p,q,W]\) 中验证等式。若待核验等式的两边相减、清去分母后得到多项式 \(N\)，给出具体多项式 \(A_1,A_2,A_3\) 使
\begin{equation}\label{g1:eq:polynomial-identity-certificate}
	N=A_1r_1+A_2r_2+A_3r_3
\end{equation}
作为形式多项式恒等式成立，就给出了 \(N=0\) 在所列关系下成立的充分条件。这是一份有限的代数证据：展开并逐项比较系数即可复核，再代入 \(r_i=0\) 得到结论。

\begin{example}[法向与曲率恒等式的代数证据]\label{g1:exm:ellipsoid-polynomial-witnesses}
	对式\eqref{g1:eq:ellipsoid-area-normal}中的 \(n\)，将 \(\lVert n\rVert^2-1\) 乘以 \(W^2\)，其分子满足
	\[
		a^2s^2(p^2+q^2)+c^2-W^2=a^2s^2r_2-r_3.
	\]
	右边在实际参数上为零，且 \(W>0\)，所以 \(n\) 确实具有单位长度。类似地，计算 \(\langle X_{\theta\theta},n\rangle+a/W\) 的分子时，所需的恒等式是
	\[
		a\bigl(1-s^2(p^2+q^2)-c^2\bigr)=-ar_1-as^2r_2.
	\]
	它把第二基本形式的一项计算化成了可逐项展开的多项式等式。

	还可以核验平均曲率与 Gauss 曲率之差：将式\eqref{g1:eq:ellipsoid-exact-curvatures}代入，有
	\[
		4W^6(H^2-K)=a^2\bigl((1+W^2)^2-4W^2\bigr)=a^2(W^2-1)^2.
	\]
	因为
	\[
		\begin{aligned}
			&(W^2-1)^2-(a^2-1)^2s^4\\
			&\quad=(r_1+r_3)\bigl(W^2-1+(a^2-1)s^2\bigr),
		\end{aligned}
	\]
	得到精确恒等式
	\begin{equation}\label{g1:eq:ellipsoid-umbilic-witness}
		H^2-K=\frac{a^2(a^2-1)^2s^4}{4W^6}.
	\end{equation}
	分母严格为正，分子是平方的乘积；因此右边非负，并且当 \(a\ne1\) 时，在 \(0<\theta<\pi\) 内严格为正。这与两个主曲率在此区域内不相等的结论一致。两极的 \(H=-a\)、\(K=a^2\) 则满足 \(H^2=K\)。
\end{example}

清分母与平方都带有条件。由 \(N=0\) 恢复 \(N/Q=0\)，必须确认 \(Q\ne0\)；例如约去 \(s^2\) 后的表达式不能自行赋予极点一个正则参数。关系 \(W^2=c^2+a^2s^2\) 同时容许正、负两个 \(W\)，真正的归一化使用的是 \(W>0\)。同样，\(\sqrt{s^2W^2}=sW\) 还用了 \(s>0\)。多项式等式不会自动保存这些符号信息，也不会区分 \(n\) 与 \(-n\) 所决定的取向。将外向法向改为内向法向时，\(\mathrm{II},S,H\) 变号，\(K\) 不变；这也是为什么只核验 \(K\) 可能漏掉取向错误。

\ProseParagraphBreak

有限个数值点可以暴露一条恒等式的错误，却一般不能证明它。即使没有舍入误差，给定有限个参数点 \((\theta_j,\phi_j)\)，非零多项式
\[
	P(\theta,\phi)=\prod_j\bigl((\theta-\theta_j)^2+(\phi-\phi_j)^2\bigr)
\]
也在这些点全部为零，而在其余点为正。因而两个公式可以在全部采样点吻合，仍在别处不同。式\eqref{g1:eq:polynomial-identity-certificate}提供的则是对任意满足已列关系与非零条件的参数都成立的推导；若只得到很小的浮点残差，还需要把该残差同输入误差及正则性条件联系起来。

\ProseParagraphBreak

在椭球计算中，法向需要将叉积除以自身长度。两个切向量近乎共线或其中一个很短时，即使各自的绝对误差小，叉积的方向仍可能明显改变。下面用带有明确常数的估计说明这个现象；所有范数均为 \(\mathbb R^3\) 的Euclid 范数。

\begin{proposition}[切向量扰动下的单位法向]\label{g1:prop:cross-product-normal-stability}
	设 \(v,w,\widetilde v,\widetilde w\in\mathbb R^3\)，并给定 \(m>0\)、\(\alpha,\beta\ge0\)，满足
	\[
		\lVert v\times w\rVert\ge m,\qquad
		\lVert\widetilde v-v\rVert\le\alpha,\qquad
		\lVert\widetilde w-w\rVert\le\beta.
	\]
	令 \(\eta=\beta\lVert v\rVert+\alpha\lVert w\rVert+\alpha\beta\)。若 \(\eta<m\)，则 \(\widetilde v,\widetilde w\) 线性无关。按相同的向量次序取单位法向
	\[
		n=\frac{v\times w}{\lVert v\times w\rVert},\qquad
		\widetilde n=\frac{\widetilde v\times\widetilde w}{\lVert\widetilde v\times\widetilde w\rVert},
	\]
	有
	\begin{equation}\label{g1:eq:cross-product-normal-error}
		\begin{aligned}
			\bigl|\lVert\widetilde v\times\widetilde w\rVert-\lVert v\times w\rVert\bigr|&\le\eta,\\
			\lVert\widetilde n-n\rVert&\le\frac{2\eta}{\lVert v\times w\rVert}\le\frac{2\eta}{m}.
		\end{aligned}
	\end{equation}
	而且 \(\langle\widetilde n,n\rangle>0\)，所以这两个有向法向之间的夹角小于 \(\pi/2\)。
\end{proposition}
\begin{proof}
	记 \(e=\widetilde v-v\)、\(f=\widetilde w-w\)，以及 \(z=v\times w\)、\(\widetilde z=\widetilde v\times\widetilde w\)。叉积的双线性给出
	\[
		\widetilde z-z=v\times f+e\times w+e\times f.
	\]
	利用 \(\lVert x\times y\rVert\le\lVert x\rVert\lVert y\rVert\)，逐项估计便得 \(\lVert\widetilde z-z\rVert\le\eta\)。范数的反三角不等式于是给出
	\[
		\bigl|\lVert\widetilde z\rVert-\lVert z\rVert\bigr|\le\eta,
		\qquad\lVert\widetilde z\rVert\ge m-\eta>0.
	\]
	因此新法向有定义。为估计归一化造成的变化，按同一个分母 \(\lVert z\rVert\) 分拆：
	\[
		\frac{\widetilde z}{\lVert\widetilde z\rVert}-\frac z{\lVert z\rVert}
		=\frac{\widetilde z-z}{\lVert z\rVert}
		+\widetilde z\left(\frac1{\lVert\widetilde z\rVert}-\frac1{\lVert z\rVert}\right).
	\]
	第一项的范数不超过 \(\eta/\lVert z\rVert\)，第二项的范数等于
	\[
		\frac{\bigl|\lVert z\rVert-\lVert\widetilde z\rVert\bigr|}{\lVert z\rVert}
		\le\frac\eta{\lVert z\rVert}.
	\]
	相加即得法向误差界。最后，
	\[
		\langle z,\widetilde z\rangle
		=\lVert z\rVert^2+\langle z,\widetilde z-z\rangle
		\ge\lVert z\rVert(\lVert z\rVert-\eta)>0,
	\]
	除以两个正范数即得夹角的断言。
\end{proof}

这个估计也控制切向量 Gram 矩阵的行列式。记
\[
	g=\begin{pmatrix}\langle v,v\rangle&\langle v,w\rangle\\
	\langle w,v\rangle&\langle w,w\rangle\end{pmatrix},
	\qquad \widetilde g=\begin{pmatrix}\langle\widetilde v,\widetilde v\rangle&\langle\widetilde v,\widetilde w\rangle\\
	\langle\widetilde w,\widetilde v\rangle&\langle\widetilde w,\widetilde w\rangle\end{pmatrix}.
\]
由 \(\det g=\lVert v\times w\rVert^2\) 及前述范数界，立刻有
\[
	\det\widetilde g\ge(m-\eta)^2>0,\qquad
	|\det\widetilde g-\det g|
	\le\eta\bigl(2\lVert v\times w\rVert+\eta\bigr).
\]
后一式只是两平方之差的因式分解；前一式则说明误差仍不足以破坏正则性。两种结论分别控制面积密度的变化和参数是否失去秩。

\ProseParagraphBreak

若要对一整族椭球给出统一误差界，应同时限制形状参数与坐标范围。设
\[
	0<a_0\le a\le a_1,\qquad
	0<\theta_0<\pi/2,\qquad\theta_0\le\theta\le\pi-\theta_0.
\]
由于 \(s\ge\sin\theta_0\) 及 \(c^2+s^2=1\)，有
\[
	W\ge\min\{1,a_0\},\qquad
	\lVert X_\theta\rVert=W\le M\coloneqq\max\{1,a_1\},\qquad
	\lVert X_\phi\rVert=s\le1.
\]
所以可取统一的正则性下界
\[
	m=\sin\theta_0\min\{1,a_0\}>0.
\]
若两项近似切向量的误差都不超过 \(\varepsilon\ge0\)，令 \(\eta_0=(M+1)\varepsilon+\varepsilon^2\)。只要 \(\eta_0<m\)，\cref{g1:prop:cross-product-normal-stability}就对这整个参数区域给出
\begin{equation}\label{g1:eq:ellipsoid-uniform-normal-bound}
	\lVert\widetilde n-n\rVert
	\le\frac{2\bigl((M+1)\varepsilon+\varepsilon^2\bigr)}
	{\sin\theta_0\min\{1,a_0\}}.
\end{equation}
这里 \(\widetilde n\) 指近似切向量的叉积经精确归一化所得的单位向量；若归一化本身还带来已证明不超过 \(\zeta\) 的误差，由三角不等式在右端再加 \(\zeta\) 即可。

\ProseParagraphBreak

当 \(\theta_0\) 接近零时，经纬参数在极点附近变差，分母中的 \(\sin\theta_0\) 将这种影响明确显示出来；这不表示固定椭球的极点不光滑。当 \(a_0\) 接近零时，则整族参数已经接近压扁的极限，在赤道有 \(\lVert X_\theta\times X_\phi\rVert=a\)。准确给出 \(a_0,\theta_0\) 与输入误差，才能判断一次很小的计算残差是否足以控制几何误差。精确恒等式负责确认理想模型，带有正则性下界的估计负责说明近似数据能够支持多强的结论。

\section{数值误差与几何解释}
\label{g1:sec:C02}

曲率涉及曲面的二阶导数，图像通常只使用有限个位置值；参数变换又可能把一块平整的曲面压缩到很窄的参数区域。因而一个异常数值至少有三种可能来源：所研究的几何对象确实在退化，坐标对这个对象的描述正在失效，或者离散计算尚未准确反映连续量。本节用明确的误差模型把这些可能性分开。

\begin{example}[同一现象的不同来源]\label{g1:exm:three-degeneracy-sources}
	单位球面的经纬参数在极点满足 \(\det g=0\)，但\subsecref{g1:subsec:C01:01}的立体投影坐标在极点仍有正定度量。这是参数失效。比较两个矩阵的行列式之前必须记住\cref{g1:prop:surface-metric-change}：换元导数为 \(J\) 时，\(\widetilde g=J^{\mathsf T}gJ\)，行列式多出因子 \((\det J)^2\)。行列式的数值大小本身不是几何不变量。

	另一种情形是半径趋于零的圆柱族
	\[
		X_\varepsilon(\theta,z)=(\varepsilon\cos\theta,\varepsilon\sin\theta,z),
		\qquad \varepsilon>0.
	\]
	在任一长度小于 \(2\pi\) 的开角区间上，它是正则参数片。取外向法向 \(n=(\cos\theta,\sin\theta,0)\)，有
	\[
		g=\begin{pmatrix}\varepsilon^2&0\\0&1\end{pmatrix},\qquad
		h=\begin{pmatrix}-\varepsilon&0\\0&0\end{pmatrix},\qquad
		K=0,\quad H=-\frac1{2\varepsilon}.
	\]
	小的度量行列式还可以局部用 \(s=\varepsilon\theta\) 改成 \(g=I\)；但横截圆的长度 \(2\pi\varepsilon\) 趋于零，\(H\) 的绝对值趋于无穷，均不会因换坐标消失。对固定高度区间，这些圆柱在空间中趋近一条线段，极限不再是一张二维正则曲面。这是这个几何族的退化；\(K\) 恒为零，说明只看一个不变量仍可能遗漏它。

	最后，对固定单位球面使用差分，粗步长下得到 \(K_h\ne1\) 完全可能。若随着步长缩小，差值按已证误差界衰减，它描述的是离散化误差；若步长过小而误差重新增大，还须检查数据和算术误差。下面将把这个说法写成可以定量核验的式子。
\end{example}

为讨论步长，先把曲面的一个参数方向固定。设 \(f\) 是取值于 \(\mathbb R^m\) 的函数，\(t\) 是待求导的点，\(h>0\) 是步长；范数均为Euclid 范数。只使用对称位置上的函数值，得到一阶和二阶\term[zhongxinchafen]{中心差分}[central difference]
\[
	\delta_h f(t)=\frac{f(t+h)-f(t-h)}{2h},\qquad
	\delta_h^2 f(t)=\frac{f(t+h)-2f(t)+f(t-h)}{h^2}.
\]
这里上标 \(2\) 表示这一二阶差分公式，并不是把步长为 \(h\) 的一阶差分算子复合两遍。实际得到的可能是近似数据 \(\widetilde f\)，因此应把公式对精确数据的偏差与输入误差分别估计。

\begin{proposition}[中心差分的误差界]\label{g1:prop:central-difference-error}
	固定 \(t\in\mathbb R\)、\(h>0\)、\(\eta\ge0\)。设 \(f\) 在包含 \([t-h,t+h]\) 的开区间上取值于 \(\mathbb R^m\)，给定三个采样值 \(\widetilde f(s)\)，并对每个 \(s\in\{t-h,t,t+h\}\) 要求 \(\|\widetilde f(s)-f(s)\|\le\eta\)。对函数或采样数据 \(v\)，记
	\[
		\delta_h v(t)=\frac{v(t+h)-v(t-h)}{2h},\qquad
		\delta_h^2v(t)=\frac{v(t+h)-2v(t)+v(t-h)}{h^2}.
	\]
	若 \(f\in C^3\)，且区间上 \(\|f^{(3)}\|\le M_3\)，则
	\begin{equation}\label{g1:eq:first-difference-error}
		\|\delta_h\widetilde f(t)-f'(t)\|
		\le\frac{M_3h^2}{6}+\frac\eta h.
	\end{equation}
	若 \(f\in C^4\)，且区间上 \(\|f^{(4)}\|\le M_4\)，则
	\begin{equation}\label{g1:eq:second-difference-error}
		\|\delta_h^2\widetilde f(t)-f''(t)\|
		\le\frac{M_4h^2}{12}+\frac{4\eta}{h^2}.
	\end{equation}
	两式中的差分均按精确算术作用于给定数据；若最终计算值另有范数不超过 \(\rho\) 的算术误差，右端再加 \(\rho\)。
\end{proposition}
\begin{proof}
	由带积分余项的 Taylor 公式，
	\[
		f(t\pm h)=f(t)\pm hf'(t)+\frac{h^2}{2}f''(t)+r_\pm,
		\qquad \|r_\pm\|\le\frac{M_3h^3}{6}.
	\]
	这个余项界由三次积分直接得到，对向量值函数同样成立：积分范数不超过范数的积分。两式相减再除以 \(2h\)，精确数据的偏差不超过 \(M_3h^2/6\)。两个输入误差相减后的范数至多为 \(2\eta\)，除以 \(2h\) 给出 \(\eta/h\)，于是得到第一式。

	二阶差分则展开到三阶项：
	\[
		f(t\pm h)=f(t)\pm hf'(t)+\frac{h^2}{2}f''(t)
		\pm\frac{h^3}{6}f^{(3)}(t)+q_\pm,
		\qquad \|q_\pm\|\le\frac{M_4h^4}{24}.
	\]
	相加并减去 \(2f(t)\)，奇数阶项消去，余项除以 \(h^2\) 后至多为 \(M_4h^2/12\)。三个输入误差的系数绝对值之和为 \(1+2+1=4\)，所以数据误差至多为 \(4\eta/h^2\)。最后一项算术误差直接用三角不等式加入。
\end{proof}

右端的 \(h^2\) 项称为\term[jieduanwucha]{截断误差}[truncation error]的上界；它来自以有限差分代替导数。另两项分别来自采样数据与执行运算。这里的 \(\eta\) 必须是输入误差的已知上界，可以来自测量精度或函数求值的误差分析；仅写“使用双精度”并没有给出这个上界。尤其两个很接近的数相减时，即使各自的相对误差小，差的相对误差仍可能大。

\ProseParagraphBreak

如果固定 \(\eta>0\)，减小 \(h\) 只能降低截断项，却会放大输入误差项。暂时忽略额外算术误差，且假定 \(M_3,M_4>0\)，将右端看作 \(h\) 的函数并求导，可得相应的最小值点
\[
	h_1=\left(\frac{3\eta}{M_3}\right)^{1/3},\qquad
	h_2=\left(\frac{48\eta}{M_4}\right)^{1/4}.
\]
它们只是这个误差上界模型中的选择，还须保证采样点留在定义域、导数界在相应区间成立。固定精度下让 \(h\) 无限趋于零，并不能由这两个估计推出导数近似收敛。若误差上界依赖于步长，则使 \(\eta(h)/h^2\to0\) 并令额外算术误差趋于零，是二阶差分收敛的一个充分条件。

\begin{example}[球面北极的离散曲率]\label{g1:exm:sphere-difference-curvature}
	在单位球面北半球采用图像参数
	\[
		X(u,v)=(u,v,\sqrt{1-u^2-v^2}),\qquad u^2+v^2<1,
	\]
	并取外向法向。计算点为 \((u,v)=(0,0)\)，此时 \(n=(0,0,1)\)、\(g=I\)、\(h_{ij}=-\delta_{ij}\)，所以 \(K=1\)、\(H=-1\)。这里的 \(h_{ij}\) 是第二基本形式的系数，以下单独的 \(h\) 是采样步长。

	取 \(0<h\le1/2\)，分别在两个参数方向使用中心差分，并用
	\[
		\delta_{uv,h}X(0,0)=
		\frac{X(h,h)-X(h,-h)-X(-h,h)+X(-h,-h)}{4h^2}
	\]
	近似混合二阶导数。全部采样点都在参数域内。对精确的采样值，一阶差分恰为 \((1,0,0)\)、\((0,1,0)\)，混合差分恰为零；两个纯二阶差分的法向分量相同，等于
	\begin{equation}\label{g1:eq:sphere-second-difference}
		b_h=\frac{2(\sqrt{1-h^2}-1)}{h^2}
		=-\frac{2}{1+\sqrt{1-h^2}}.
	\end{equation}
	由近似切向量的叉积得到的法向仍是 \((0,0,1)\)。把这些离散数据代入\cref{g1:prop:surface-curvature-formulas}所给的代数表达式，定义本例的近似曲率
	\begin{equation}\label{g1:eq:sphere-discrete-curvatures}
		H_h=b_h,\qquad K_h=b_h^2
		=\frac4{(1+\sqrt{1-h^2})^2}.
	\end{equation}
	这个定义说明了如何从离散数据产生数值；它并不把差分数据自动认作另一张曲面的基本形式。

	令 \(s=\sqrt{1-h^2}\)，利用 \((1-s)(1+s)=h^2\) 可得
	\begin{equation}\label{g1:eq:sphere-discrete-curvature-error}
		K_h-1=h^2\frac{3+s}{(1+s)^3},\qquad
		\lim_{h\to0}\frac{K_h-1}{h^2}=\frac12.
	\end{equation}
	因此这里的二阶收敛由精确恒等式证明。还可直接取得统一界：\(s\ge\sqrt3/2\)，于是
	\[
		0\le K_h-1\le
		\frac4{(1+\sqrt3/2)^3}\,h^2\qquad(0<h\le1/2).
	\]
	\cref{g1:tab:sphere-discrete-curvature}用有理步长代入约化后的式\eqref{g1:eq:sphere-discrete-curvatures}，以十进制 \(60\) 位工作精度计算，再将 \(K_h\) 舍入到小数点后 \(12\) 位、误差比舍入到 \(6\) 位。误差列在舍入 \(K_h\) 之前计算；这些显示位数用于复算，统一误差界仍由上面的等式和不等式提供。
\end{example}

\begin{table}[htbp]
	\centering
	\caption{球面北极的中心差分曲率}
	\label{g1:tab:sphere-discrete-curvature}
	\begin{tabular}{cccc}
		\toprule
		步长 \(h\) & \(K_h\) & \(e_h=K_h-1\) & \(e_{2h}/e_h\)\\
		\midrule
		\(0.2\) & \(1.020514433644\) & \(2.051443364\times10^{-2}\) & —\\
		\(0.1\) & \(1.005031470404\) & \(5.031470404\times10^{-3}\) & \(4.077224\)\\
		\(0.05\) & \(1.001251956549\) & \(1.251956549\times10^{-3}\) & \(4.018886\)\\
		\(0.025\) & \(1.000312622124\) & \(3.126221237\times10^{-4}\) & \(4.004696\)\\
		\bottomrule
	\end{tabular}
\end{table}

这个表比较的是同一个离散公式在精确数据下的截断效应，不包含任意测量噪声。式\eqref{g1:eq:sphere-second-difference}左边先做两个接近 \(1\) 的数的相减；当 \(h\) 很小时，求值误差随后还要除以 \(h^2\)。右边的有理化形式避开了这一步相减，因此更适于本模型的求值。但这种代数化简只有在已知球面公式时才可用，不能替代对一般采样数据的误差估计。换用向内法向时，所有 \(H_h\) 变号，\(K_h\) 不变；只核对 \(K_h\) 的数值仍不足以检查法向约定。

\ProseParagraphBreak

没有已知真值时，常比较 \(Q_h,Q_{h/2},Q_{h/4}\) 三个近似量。若已知存在展开 \(Q_h=Q+ch^p+o(h^p)\)，其中 \(c\ne0\)、\(p>0\)，相减得到
\[
	\frac{Q_h-Q_{h/2}}{Q_{h/2}-Q_{h/4}}\longrightarrow2^p.
\]
在上述误差展开下，差值比反映误差随步长变化的阶数，却对固定偏差不敏感。例如 \(Q_h=Q+b+ch^2\)，其中固定偏差 \(b\ne0\)，差值比仍恒等于 \(4\)，极限却是 \(Q+b\)。坐标、单位或曲率符号若从一开始就设错，网格加密也可能表现得十分规律。

\begin{proposition}[图像接近不控制曲率]\label{g1:prop:position-does-not-control-curvature}
	在固定开正方形 \(U=(-1,1)^2\) 上，令
	\[
		X_\varepsilon(u,v)=(u,v,f_\varepsilon(u,v)),\qquad
		f_\varepsilon(u,v)=\varepsilon^2
		\left(\cos\frac u\varepsilon+\cos\frac v\varepsilon\right),
		\qquad \varepsilon>0.
	\]
	当 \(\varepsilon\downarrow0\) 时，\(X_\varepsilon\) 及其一阶导数在 \(U\) 上分别一致趋于平面参数 \(X_0(u,v)=(u,v,0)\) 及其一阶导数。但选择向上的法向后，每张图像在 \((0,0)\) 都有 \(K_\varepsilon=1\)、\(H_\varepsilon=-1\)，不趋于平面的曲率。
\end{proposition}
\begin{proof}
	图像参数的前两个分量保证其处处正则。由
	\[
		|f_\varepsilon|\le2\varepsilon^2,\qquad
		|(f_\varepsilon)_u|,|(f_\varepsilon)_v|\le\varepsilon
	\]
	得到所述一致收敛。在原点，一阶偏导为零，向上法向是 \((0,0,1)\)，因而 \(g=I\)。又
	\[
		(f_\varepsilon)_{uu}(0,0)=(f_\varepsilon)_{vv}(0,0)=-1,
		\qquad(f_\varepsilon)_{uv}(0,0)=0.
	\]
	由 \(h_{ij}=\langle X_{ij},n\rangle\) 及 \([S]=g^{-1}h\)，原点的形算子为 \(-I\)，故 \(K=1,H=-1\)。平面的二阶导数则全部为零。
\end{proof}

这组图像的振幅越来越小，波长也越来越短。一张固定分辨率的图片可能看不出起伏，一阶数据也确实变得接近平面，但二阶数据没有随之变小。要由数值位置恢复曲率，需要与目标导数阶数匹配的正则性和采样条件。即使位置拟合的误差已有可靠上界，它仍不自动成为曲率误差界。

\ProseParagraphBreak

一次能够独立复核的几何实验，应让读者知道究竟比较了什么。\cref{g1:tab:geometry-experiment-data}列出本节算例需要的记录及其数学作用。若只展示一张图，仍应给出所用参数域、坐标轴尺度和法向方向；透视投影会改变图上的长度和夹角，均匀参数网格通常也不是曲面上等面积的网格。

\begin{table}[htbp]
	\centering
	\caption{几何数值实验的复核数据}
	\label{g1:tab:geometry-experiment-data}
	\begin{tabularx}{\textwidth}{@{}p{0.22\textwidth}XX@{}}
		\toprule
		记录内容 & 本节球面实验 & 对结论的作用\\
		\midrule
		对象与坐标 & 单位球面北半球图像，计算点为北极 & 确认坐标覆盖及所有采样点属于同一参数片\\
		度量与单位 & 环境标准Euclid 度量，球半径为长度单位 & 确定长度、曲率量纲及参数步长的含义\\
		法向与符号 & 向外法向，\(S=-\mathrm{d}n\)，\(H=\operatorname{tr}S/2\) & 区分符号差异、平均值差异与计算错误\\
		离散公式 & 对称一阶、二阶及四点混合差分 & 明确哪些数据生成近似量，以及误差阶数的来源\\
		步长与精度 & 表中四个步长，约化公式以十进制 \(60\) 位计算 & 使表值可以复算，并区分截断和算术效应\\
		误差依据 & 式\eqref{g1:eq:sphere-discrete-curvature-error}及其统一界 & 将有限个表值与已证明的全步长结论区分\\
		\bottomrule
	\end{tabularx}
\end{table}

对一般曲面，具体诊断可以从保持对象不变的比较开始：在重叠区域换一张正则坐标图，保持同一度量和几何法向，再计算同一不变量；随后固定坐标，改变步长和工作精度，检查变化是否符合误差估计。前一种比较检查坐标及其变换的实现是否正确，后一种比较检查离散化与算术运算的影响。只有参数族中的对象本身发生改变时，才进一步比较曲率、长度、面积等几何量随参数的变化。

\begin{chaptersummary}
	坐标公式必须连同定义域、逆映射和重叠区域一起使用。球面经纬参数在极点失效，立体投影图仍能给出正定度量；同一像集也可能由秩退化的映射描述。判断奇点属于参数还是对象，需要换图或采用与坐标无关的量。

	\ProseParagraphBreak

	联络矩阵随标架变化出现非齐次项，曲率则按共轭变换。矩阵值微分形式的乘法同时涉及矩阵次序和楔积符号。极坐标、旋转型度量及三维底空间上的显式计算，分别检验无挠条件、曲率号约定与 Bianchi 消去。

	\ProseParagraphBreak

	旋转椭球的计算把参数范围、正则性、外向法向和曲率公式联系起来。精确恒等式需要在合法的定义域上化简，法向误差界则需要切向量外积的正下界。中心差分还表明：步长减小时截断误差下降，输入误差却可能被求导放大。表值、图片与加密趋势应和误差依据一起解释，几何结论才具有可复核的范围。
\end{chaptersummary}

\BookExercises

\begin{exercise}
	\label{g1:ex:longitude-cut-transition}
	把\eqref{g1:eq:spherical-chart-parameters}中的经度区间改为 \(0<\widetilde\phi<2\pi\)，余纬仍满足 \(0<\theta<\pi\)。
	\begin{enumerate}
		\item 写出新参数图的精确像，并给出光滑逆映射。
		\item 求它与原经度区间 \((-\pi,\pi)\) 所给图之间的过渡映射。说明过渡映射为何允许在交叠区域的不同连通分支上采用不同公式。
	\end{enumerate}
\end{exercise}
\begin{solution}
	(1) 新图删除的半经线是
	\[
		\widetilde\Gamma=\{(x,y,z)\in S^2:y=0,\ x\ge0\},
	\]
	包含两个极点。对补集中的点令 \(\rho=\sqrt{x^2+y^2}\)，则 \(\rho-x>0\)，逆映射为
	\[
		\theta=\arccos z,\qquad
		\widetilde\phi=\pi+2\arctan\frac{-y}{\rho-x}.
	\]
	后式取值在 \((0,2\pi)\)，半角公式给出 \(\cos\widetilde\phi=x/\rho\)、\(\sin\widetilde\phi=y/\rho\)，因此代回后得到原来的球面点。原参数区间内经纬取值唯一，所以这是双侧逆。分母严格为正且 \(-1<z<1\)，也保证逆映射光滑。

	(2) 两图交叠处恰有 \(y\ne0\)，分成 \(y>0\) 和 \(y<0\) 两个开分支。保留余纬不变，经度关系为
	\[
		\widetilde\phi=
		\begin{cases}
			\phi,&0<\phi<\pi,\\
			\phi+2\pi,&-\pi<\phi<0.
		\end{cases}
	\]
	相应逆式在 \(0<\widetilde\phi<\pi\) 时为 \(\phi=\widetilde\phi\)，在 \(\pi<\widetilde\phi<2\pi\) 时为 \(\phi=\widetilde\phi-2\pi\)。每个公式都是平移，其导数矩阵为单位阵。分段的接缝不属于交叠域，所以不需要在接缝处比较左右极限；光滑性只要求在实际定义域的每个点附近成立。这两张经纬图的并覆盖所有非极点，仍没有覆盖两极。
\end{solution}

\begin{exercise}
	\label{g1:ex:curve-through-stereographic-pole}
	考虑单位球面上的曲线 \(\gamma(t)=(\sin t,0,\cos t)\)，其中 \(|t|<\pi/2\)。分别写出它在 \(\sigma_N\)、\(\sigma_S\) 中的坐标，说明两种表达在 \(t=0\) 附近的适用范围，并用\eqref{g1:eq:stereographic-sphere-metric}核验其速度恒为 \(1\)。
\end{exercise}
\begin{solution}
	半角公式给出
	\[
		\sigma_S(\gamma(t))=\left(\tan\frac t2,0\right)
		\quad(|t|<\pi/2),
	\]
	而
	\[
		\sigma_N(\gamma(t))=\left(\cot\frac t2,0\right)
		\quad(0<|t|<\pi/2).
	\]
	第一式在 \(t=0\) 的值为 \((0,0)\)，光滑通过北极。第二式在 \(t=0\) 处没有定义，其第一坐标从两侧分别趋向不同符号的无穷；原因是 \(\gamma(0)=N\) 不属于 \(U_N\)。这不影响原曲线在 \(t=0\) 的光滑性。

	令 \(u(t)=\tan(t/2)\)，则 \(u'(t)=\tfrac12\sec^2(t/2)\)。在南极投影图中，球面速度平方为
	\[
		\frac{4}{(1+u(t)^2)^2}\,u'(t)^2
		=\frac{4}{\sec^4(t/2)}\cdot\frac14\sec^4(t/2)=1.
	\]
	这个等式也在 \(t=0\) 成立。对 \(t\ne0\)，另一图给出 \(v(t)=\cot(t/2)\)、\(v'(t)=-\tfrac12\csc^2(t/2)\)，代入同一个度量公式仍得到 \(1\)。坐标导数的大小可以很不相同，但与各自的度量系数合在一起，给出的球面速度相同。直接计算 \(\gamma'(t)=(\cos t,0,-\sin t)\) 的Euclid 长度，也得到 \(1\)。
\end{solution}

\begin{exercise}
	\label{g1:ex:stereographic-transition-orientation}
	考虑同号实立体投影图之间的过渡 \(\tau(u)=u/|u|^2\)，\(u\in\mathbb R^2\setminus\{0\}\)。
	\begin{enumerate}
		\item 计算 \(\mathrm{D}\tau(u)\)，求它在径向与垂直径向上的特征值，并判断该过渡是否保持取向。
		\item 用导数直接核对两图中的度量\eqref{g1:eq:stereographic-sphere-metric}相互变换一致。
		\item 若将第二张图改为 \(\widetilde\sigma_S=(\sigma_S^1,-\sigma_S^2)\)，以 \(z=u_1+\mathrm i u_2\)、\(w=\widetilde\sigma_S^1+\mathrm i\widetilde\sigma_S^2\) 表示坐标，求新的复数过渡式，并比较新旧图册的取向关系。
	\end{enumerate}
\end{exercise}
\begin{solution}
	(1) 记 \(r=|u|>0\)，微分为
	\[
		\mathrm{D}\tau(u)=\frac1{r^2}\left(I_2-2\frac{uu^{\mathsf T}}{r^2}\right).
	\]
	括号内的矩阵把径向向量 \(u\) 变为 \(-u\)，把与 \(u\) 正交的向量保持不变。因此两个特征值是 \(-r^{-2}\) 与 \(r^{-2}\)，行列式为 \(-r^{-4}<0\)。所以原两张实图的过渡反转取向，不能把它们按同样的坐标轴取向直接称作一个定向图册。

	(2) 括号内矩阵是正交反射，所以对任意向量 \(h\)，有 \(|\mathrm{D}\tau(u)h|^2=r^{-4}|h|^2\)。又 \(|\tau(u)|^2=r^{-2}\)，于是第二张图的度量拉回到 \(u\) 处给出
	\[
		\frac{4}{(1+|\tau(u)|^2)^2}|\mathrm{D}\tau(u)h|^2
		=\frac{4r^{-4}}{(1+r^{-2})^2}|h|^2
		=\frac{4}{(1+r^2)^2}|h|^2.
	\]
	这正是第一张图的度量。度量保持性不要求坐标变换保持取向。

	(3) 改变第二坐标的符号后，
	\[
		w=\frac{u_1-\mathrm i u_2}{u_1^2+u_2^2}=\frac1z.
	\]
	新过渡比 \(\tau\) 多复合了一次平面反射，因此导数行列式为 \(r^{-4}>0\)，两张图现在具有相容的取向；反射又保持Euclid 平方长度，所以度量公式不变。若另外指定球面的外法向取向，仍须核对符号：在 \(\Phi_N(0)=S\) 处，两个参数导数为 \((2,0,0)\)、\((0,2,0)\)，叉积指向正 \(z\) 方向，与南极的外法向相反。因此这一配对给出与外法向相反的全局取向；同时在两图各反转一个坐标，才得到外法向取向。
\end{solution}

\begin{exercise}\label{g1:ex:connection-check-polar-coordinate}
	在Euclid 平面的极坐标区域 \(r>0\) 中，正交单位标架 \(e=(\partial_r,r^{-1}\partial_\varphi)\) 的联络矩阵为
	\[
		A=\begin{pmatrix}0&-\mathrm{d}\varphi\\
		\mathrm{d}\varphi&0\end{pmatrix}.
	\]
	换成坐标标架 \(e'=(\partial_r,\partial_\varphi)\)，求变换矩阵 \(g\)、联络矩阵 \(A'\)、度量矩阵 \(G'\) 和余标架列 \(\theta'\)。直接展开验证曲率和挠率均为零，并核对度量相容性。解释为什么 \(A'\) 不反对称并不意味着联络不保持度量。
\end{exercise}
\begin{solution}
	有 \(e'=eg\)，其中 \(g=\operatorname{diag}(1,r)\)。于是
	\[
		A'=g^{-1}Ag+g^{-1}\mathrm{d}g
		=\begin{pmatrix}
			0&-r\,\mathrm{d}\varphi\\
			r^{-1}\mathrm{d}\varphi&r^{-1}\mathrm{d}r
		\end{pmatrix},
		\quad
		G'=\begin{pmatrix}1&0\\0&r^2\end{pmatrix},
		\quad
		\theta'=\begin{pmatrix}\mathrm{d}r\\\mathrm{d}\varphi\end{pmatrix}.
	\]
	特别地，\(g^{-1}\mathrm{d}g\) 的右下角为 \(r^{-1}\mathrm{d}r\)，不能省略。

	令 \(\eta=\mathrm{d}r\wedge\mathrm{d}\varphi\)。逐项外微分与矩阵楔积给出
	\[
		\mathrm{d}A'
		=\begin{pmatrix}0&-\eta\\-r^{-2}\eta&0\end{pmatrix},
		\qquad
		A'\wedge A'
		=\begin{pmatrix}0&\eta\\r^{-2}\eta&0\end{pmatrix}.
	\]
	两者相消，故 \(F'=0\)。在这个坐标余标架中，\(\mathrm{d}\theta'=0\)，而
	\[
		A'\wedge\theta'
		=\begin{pmatrix}
			-r\,\mathrm{d}\varphi\wedge\mathrm{d}\varphi\\
			r^{-1}\mathrm{d}\varphi\wedge\mathrm{d}r+
			r^{-1}\mathrm{d}r\wedge\mathrm{d}\varphi
		\end{pmatrix}=0,
	\]
	所以挠率也为零。

	度量相容性由矩阵乘法核验为
	\[
		(A')^{\mathsf T}G'+G'A'
		=\begin{pmatrix}0&0\\0&2r\,\mathrm{d}r\end{pmatrix}
		=\mathrm{d}G'.
	\]
	\(G'\) 随 \(r\) 变化，所以正确条件包含 \(\mathrm{d}G'\)，也包含左右两处 \(G'\)。只有在度量矩阵恒为 \(I\) 的正交单位标架中，度量相容条件才要求联络矩阵反对称。
\end{solution}

\begin{exercise}\label{g1:ex:connection-check-same-curvature}
	在 \(r>0\) 的坐标区域中取度量
	\[
		g_f=\mathrm{d}r^2+\sinh^2r\,\mathrm{d}\varphi^2,
		\qquad e=(\partial_r,(\sinh r)^{-1}\partial_\varphi),
		\qquad J=\begin{pmatrix}0&-1\\1&0\end{pmatrix}.
	\]
	在这同一个固定标架中，分别用
	\[
		A=J\cosh r\,\mathrm{d}\varphi,\qquad
		\widetilde A=J(\cosh r\,\mathrm{d}\varphi+\mathrm{d}r)
	\]
	定义两个切丛联络。证明它们都保持 \(g_f\)，具有相同曲率矩阵，但只有第一个无挠。两者是否都通过第二 Bianchi 恒等式核验？这些事实是否足以把第二个联络也认作该度量的无挠度量联络？
\end{exercise}
\begin{solution}
	两矩阵在给定正交单位标架中都反对称，所以都保持度量。每个矩阵都形如 \(J\omega\)，故自楔为 \(J^2\omega\wedge\omega=0\)。又因为 \(\mathrm{d}(\mathrm{d}r)=0\)，两曲率相同：
	\[
		F=\widetilde F=J\sinh r\,\mathrm{d}r\wedge\mathrm{d}\varphi.
	\]
	对 \(A\)，\cref{g1:prop:connection-check-warped-metric}给出挠率为零，并给出该度量的 Gauss 曲率 \(-1\)。

	余标架列是 \(\theta=(\mathrm{d}r,\sinh r\,\mathrm{d}\varphi)^{\mathsf T}\)。由于第一个联络的挠率为零，可得第二个联络的挠率为
	\[
		\widetilde\Theta
		=(\widetilde A-A)\wedge\theta
		=J\,\mathrm{d}r\wedge
		\begin{pmatrix}\mathrm{d}r\\\sinh r\,\mathrm{d}\varphi\end{pmatrix}
		=\begin{pmatrix}
			-\sinh r\,\mathrm{d}r\wedge\mathrm{d}\varphi\\0
		\end{pmatrix}\ne0.
	\]
	具体在正交单位基上求值，有 \(\widetilde T(e_1,e_2)=-e_1\)。

	两者都通过第二 Bianchi 恒等式；在这个二维底空间上，相应三形式本来就恒为零。即使已经分别核对度量相容和曲率相同，也没有核对无挠条件。第二个联络确实是另一个度量联络，不能作为该度量的无挠度量联络。这里两个矩阵是在同一个固定标架和同一余标架中比较；若把矩阵改变解释成换标架，还须相应改变余标架，不能只替换 \(A\)。
\end{solution}

\begin{exercise}\label{g1:ex:connection-check-gauge-three-dimensional}
	在\cref{g1:exm:connection-check-nonabelian}的三维底空间模型中，保持 \(B,C,H\) 及
	\[
		A=B\,\mathrm{d}x+xC\,\mathrm{d}y+yH\,\mathrm{d}z
	\]
	不变，令 \(e'=eg\)，其中 \(g=\operatorname{diag}(e^z,e^{-z})\)。求 \(A'\)，再从 \(A'\) 的坐标系数直接计算 \(F'\)，核对 \(F'=g^{-1}Fg\)。若遗漏换标架公式中的 \(g^{-1}\mathrm{d}g\)，仅取 \(\widetilde A=g^{-1}Ag\)，其曲率是否仍等于 \(g^{-1}Fg\)？对这组错误变换后的数据，单独验证“由 \(\widetilde A\) 算出的曲率满足 Bianchi 恒等式”能否识别错误？
\end{exercise}
\begin{solution}
	直接共轭和微分得到
	\[
		g^{-1}Bg=e^{-2z}B,\qquad
		g^{-1}Cg=e^{2z}C,\qquad
		g^{-1}Hg=H,\qquad
		g^{-1}\mathrm{d}g=H\,\mathrm{d}z.
	\]
	所以
	\[
		A'=e^{-2z}B\,\mathrm{d}x+xe^{2z}C\,\mathrm{d}y+(y+1)H\,\mathrm{d}z.
	\]
	使用\eqref{g1:eq:connection-check-curvature-components}逐对计算：
	\[
		\begin{aligned}
			F'_{xy}
			&=e^{2z}C+[e^{-2z}B,xe^{2z}C]=e^{2z}C+xH,\\
			F'_{xz}
			&=2e^{-2z}B+[e^{-2z}B,(y+1)H]
			=-2ye^{-2z}B,\\
			F'_{yz}
			&=H-2xe^{2z}C+[xe^{2z}C,(y+1)H]
			=H+2xye^{2z}C.
		\end{aligned}
	\]
	分别对\eqref{g1:eq:connection-check-three-curvatures}中的三个矩阵作 \(g^{-1}(\cdot)g\)，恰好得到这些结果。含 \(z\) 的偏导项与附加联络项的交换子在这里发生了必要抵消。

	若遗漏 \(H\,\mathrm{d}z\)，则 \(\widetilde A_z=yH\)，其他两个系数不变。同样计算给出
	\[
		\widetilde F_{xy}=e^{2z}C+xH,\qquad
		\widetilde F_{xz}=2(1-y)e^{-2z}B,\qquad
		\widetilde F_{yz}=H+2x(y-1)e^{2z}C.
	\]
	后两个系数与正确共轭结果不同；特别地，在 \(x=y=0\) 时，正确的 \(F'_{xz}\) 为零，错误结果却为 \(2e^{-2z}B\ne0\)。所以它们没有表示原联络在新标架中的曲率。

	然而，\(\widetilde A\) 本身仍是一组光滑矩阵一形式，会在标架 \(e'\) 中定义另一个联络。按 \(\widetilde F=\mathrm{d}\widetilde A+\widetilde A\wedge\widetilde A\) 计算，\cref{g1:thm:bianchi-identity}仍保证
	\[
		\mathrm{d}\widetilde F+\widetilde A\wedge\widetilde F-\widetilde F\wedge\widetilde A=0.
	\]
	因此这项恒等式能与错误的换标架操作同时成立；识别此处错误需要比较联络变换律或曲率共轭律，不能只检查各组数据内部的 Bianchi 恒等式。
\end{solution}

\begin{exercise}\label{g1:ex:ellipsoid-polar-graph}
	设 \(a>0\)。对椭球 \(x^2+y^2+z^2/a^2=1\)，在北极附近采用图像参数
	\[
		Z_a(x,y)=\bigl(x,y,a\sqrt{1-x^2-y^2}\bigr),\qquad x^2+y^2<1,
	\]
	并取外向单位法向，以 \(S=-\mathrm{d}n\)、\(H=\tfrac12\operatorname{tr}S\) 为约定。直接计算 \((x,y)=(0,0)\) 处的第一、第二基本形式、形算子、平均曲率和 Gauss 曲率。说明该点是否为脐点，以及经纬参数在那里失去秩为何不表示椭球有奇点。若在同一图像上改取相反的法向，哪些计算结果会改变？
\end{exercise}
\begin{solution}
	令 \(f(x,y)=a\sqrt{1-x^2-y^2}\)。在原点有
	\[
		f_x=f_y=0,\qquad f_{xx}=f_{yy}=-a,\qquad f_{xy}=0.
	\]
	例如 \(f_x=-ax/\sqrt{1-x^2-y^2}\)，再对 \(x\) 求导并在原点代入即得 \(f_{xx}=-a\)；其余项同理。因此 \(Z_x(0)=(1,0,0)\)、\(Z_y(0)=(0,1,0)\)，二者线性无关，第一基本形式的矩阵为单位矩阵。

	向上的图像法向在北极为 \(n(0)=(0,0,1)\)，也正是外向法向。由 \(h_{ij}=\langle Z_{ij},n\rangle\)，第二基本形式的矩阵是 \(-aI_2\)。于是，在这两个切向量组成的基下，
	\[
		I=\mathrm dx^2+\mathrm dy^2,\qquad
		\mathrm{II}=-a(\mathrm dx^2+\mathrm dy^2),\qquad
		[S]=-aI_2,\qquad H=-a,\qquad K=a^2.
	\]
	两项主曲率相同，所以北极为脐点。经纬参数的经度在极点无法区分点的移动方向，表现为 \(X_\phi=0\)；图像参数的两个导数仍独立，直接证明这里是坐标退化而非曲面奇点。

	若取 \(-n\)，内积 \(\langle Z_{ij},n\rangle\) 全部变号，故 \(\mathrm{II},S,H\) 变号，两项主曲率变为 \(a\)。第一基本形式由切向量的内积定义，不随法向改变；\(K\) 是两项主曲率的乘积，仍等于 \(a^2\)，脐点性质也不变。
\end{solution}

\begin{exercise}\label{g1:ex:finite-exact-polynomial-samples}
	设 \(J\subset\mathbb R\) 为非空开区间，\(d\ge0\) 为整数。两个有理函数 \(R_1=p_1/q_1\)、\(R_2=p_2/q_2\) 的系数均为实数，且 \(q_1,q_2\) 在 \(J\) 上没有零点。假设多项式 \(N=p_1q_2-p_2q_1\) 的次数至多为 \(d\)，并在 \(J\) 内 \(d+1\) 个互不相同的点精确得到 \(R_1=R_2\)。证明两函数在整个 \(J\) 上相等。说明若没有次数界，仅给定任意有限组采样点为何不足以推出同一结论；若只知道采样值之差小于某个正容差，又为何不足以证明精确相等。
\end{exercise}
\begin{solution}
	记采样点为 \(t_0,\ldots,t_d\)。由于两分母在这些点非零，精确相等给出 \(N(t_j)=0\)。若多项式 \(N\) 不恒为零，由因式定理，\(N(t_0)=0\) 蕴含 \(N(t)=(t-t_0)N_1(t)\)。对 \(j\ne0\)，\(t_j-t_0\ne0\)，所以 \(N_1(t_j)=0\)。逐次重复，得到
	\[
		N(t)=\prod_{j=0}^d(t-t_j)\,Q(t)
	\]
	且 \(Q\) 是非零多项式，迫使 \(\deg N\ge d+1\)，与假设矛盾。因此 \(N=0\) 为多项式恒等式。对每个 \(t\in J\)，\(q_1(t)q_2(t)\ne0\)，故 \(R_1(t)-R_2(t)=N(t)/(q_1(t)q_2(t))=0\)。

	若没有次数界，对任意有限个不同采样点 \(t_1,\ldots,t_m\)，取 \(R_1(t)=0\)、\(R_2(t)=\prod_{j=1}^m(t-t_j)\)。它们在全部采样点相等，却在 \(J\) 的其余点不同。若采样组为空，非零常数同样构成反例。

	最后，即使有次数界 \(d=0\)，也可取 \(R_1=0\)、\(R_2=\delta/2\)，其中 \(\delta>0\) 是容差。每个点的差都严格小于 \(\delta\)，两函数仍不相等。因此，要以有限采样作为证明的依据，必须同时具备已证次数界、足够多个互异点、非零分母与精确等式；浮点近似值本身不满足最后一个条件。
\end{solution}

\begin{exercise}\label{g1:ex:posterior-normal-certificate}
	设 \(v,w\in\mathbb R^3\) 是未知的真实切向量，\(\widetilde v,\widetilde w\) 是给定近似值，并已证明
	\[
		\lVert v-\widetilde v\rVert\le\alpha,\qquad
		\lVert w-\widetilde w\rVert\le\beta,
		\qquad\alpha,\beta\ge0.
	\]
	令
	\[
		\widehat\eta=\beta\lVert\widetilde v\rVert
		+\alpha\lVert\widetilde w\rVert+\alpha\beta,
		\qquad r=\lVert\widetilde v\times\widetilde w\rVert.
	\]
	证明，若 \(r>\widehat\eta\)，则真实切向量线性无关，按相同次序定义的真实单位法向满足
	\[
		\left\lVert\frac{v\times w}{\lVert v\times w\rVert}
		-\frac{\widetilde v\times\widetilde w}{r}\right\rVert
		\le\frac{2\widehat\eta}{r}.
	\]
	这个判据只含给定近似值及误差上界，因而可在不知道真实叉积下界时使用。再取 \(t>0\)，考察
	\[
		v_t=\widetilde v_t=(1,0,0),\qquad
		w_t=(1,t,0),\qquad\widetilde w_t=(1,t,t/2).
	\]
	计算其切向量误差与单位法向误差，并说明若不保留正则性的统一下界，为何不能仅从切向量误差趋于零推出法向误差趋于零。
\end{exercise}
\begin{solution}
	写 \(v=\widetilde v+e\)、\(w=\widetilde w+f\)。展开叉积得到
	\[
		v\times w-\widetilde v\times\widetilde w
		=\widetilde v\times f+e\times\widetilde w+e\times f,
	\]
	故两叉积的距离不超过 \(\widehat\eta\)。由范数的反三角不等式，真实叉积的长度至少为 \(r-\widehat\eta>0\)，真实切向量因此独立。对两个非零叉积，采用\cref{g1:prop:cross-product-normal-stability}证明中的归一化分拆，并以已知叉积的长度 \(r\) 为分母，得到
	\[
		\left\lVert\frac{v\times w}{\lVert v\times w\rVert}
		-\frac{\widetilde v\times\widetilde w}{r}\right\rVert
		\le\frac{\lVert v\times w-\widetilde v\times\widetilde w\rVert}{r}
		+\frac{|\lVert v\times w\rVert-r|}{r}
		\le\frac{2\widehat\eta}{r}.
	\]
	若实际计算的 \(r\) 和 \(\widehat\eta\) 也仅有区间界，应用判据时须用 \(r\) 的已证下界与 \(\widehat\eta\) 的已证上界来比较。

	对题中的向量族，第一项误差为零，第二项误差为 \(t/2\)，趋于零。两个叉积则为
	\[
		v_t\times w_t=(0,0,t),\qquad
		\widetilde v_t\times\widetilde w_t=(0,-t/2,t).
	\]
	因为 \(t>0\)，相应单位法向分别为
	\[
		n_t=(0,0,1),\qquad
		\widetilde n_t=(0,-1/\sqrt5,2/\sqrt5).
	\]
	它们的距离恒为
	\[
		\lVert\widetilde n_t-n_t\rVert
		=\sqrt{2-\frac4{\sqrt5}}>0.
	\]
	每一对真实切向量都独立，但 \(\lVert v_t\times w_t\rVert=t\to0\)，没有统一的正下界。数据误差与真实叉积长度按相同比例趋于零，归一化后的方向误差便不必趋于零。这也说明逐点正则与对一族数据具有统一稳定性是不同的要求。
\end{solution}

\begin{exercise}\label{g1:ex:difference-noise-sharpness}
	在 \([-h,h]\) 上令真实函数 \(f\equiv0\)，每个采样值的绝对误差不超过 \(\eta>0\)。分别给出使 \(|\delta_h\widetilde f(0)|=\eta/h\) 和使 \(|\delta_h^2\widetilde f(0)|=4\eta/h^2\) 的数据。若 \(\eta(h)=h^\alpha\)，求保证任意这类数据的二阶差分都趋于零的 \(\alpha\) 范围。进一步，在截断误差恰为零时，解释为何仍不能只靠缩小步长消除固定噪声。
\end{exercise}
\begin{solution}
	对一阶差分取 \(\widetilde f(h)=\eta\)、\(\widetilde f(-h)=-\eta\)，便得 \(\delta_h\widetilde f(0)=\eta/h\)。对二阶差分取 \(\widetilde f(\pm h)=\eta\)、\(\widetilde f(0)=-\eta\)，便得 \(\delta_h^2\widetilde f(0)=4\eta/h^2\)。各项数据均满足给定误差上界，因此命题中的这两个数据误差常数无法在同样假设下降低。

	当 \(\eta(h)=h^\alpha\) 时，任意数据的二阶差分绝对值至多为 \(4h^{\alpha-2}\)，所以 \(\alpha>2\) 足以保证收敛到零。若 \(\alpha=2\)，刚才的极端数据始终给出 \(4\)；若 \(\alpha<2\)，则该值趋于无穷。因此保证所有允许数据收敛的充要条件是 \(\alpha>2\)。这些例子的真实函数为常数，截断误差已经为零；固定 \(\eta\) 时仍有允许数据产生 \(h^{-2}\) 量级的误差，问题来自求导放大了输入误差。
\end{solution}

\begin{exercise}\label{g1:ex:discrete-sphere-extrapolation}
	对\cref{g1:exm:sphere-difference-curvature}中的 \(K_h\)，令 \(E_h=(4K_{h/2}-K_h)/3\)。证明当 \(h\to0\) 时
	\[
		K_h=1+\frac12h^2+\frac5{16}h^4+O(h^6),\qquad
		E_h=1-\frac5{64}h^4+O(h^6).
	\]
	若实际取得的 \(K_h,K_{h/2}\) 各有绝对值不超过 \(\delta\) 的误差，估计由此给 \(E_h\) 造成的误差。若两项另有相同的固定偏差 \(b\)，这种组合能否消除它？
\end{exercise}
\begin{solution}
	在零附近，\(\sqrt{1-h^2}=1-h^2/2-h^4/8+O(h^6)\)，其中余项由平方根函数在 \(1\) 附近的有界三阶导数给出。因此
	\[
		K_h=\left(1-\frac{h^2}4-\frac{h^4}{16}+O(h^6)\right)^{-2}.
	\]
	用 \((1-z)^{-2}=1+2z+3z^2+O(z^3)\) 展开，得到 \(h^2\) 的系数为 \(1/2\)，\(h^4\) 的系数为 \(1/8+3/16=5/16\)。将同一展开中的 \(h\) 换成 \(h/2\)，有
	\[
		\frac{4K_{h/2}-K_h}{3}
		=1+\frac13\left(\frac{4}{16}-1\right)\frac5{16}h^4+O(h^6)
		=1-\frac5{64}h^4+O(h^6).
	\]
	两项独立误差经系数 \(4/3\)、\(-1/3\) 组合，绝对值至多为 \(5\delta/3\)。相同的固定偏差则变成 \((4b-b)/3=b\)，完全保留。这种组合消去了已知形式的主要截断项，并不同时消除数据误差或系统偏差。
\end{solution}

\begin{exercise}\label{g1:ex:curvature-under-scaling}
	设 \(X:U\to\mathbb R^3\) 为光滑浸入，具有选定单位法向 \(n\)，沿用 \(S=-\mathrm{d}n\)、\(h(v,w)=g(Sv,w)\)、\(H=\tfrac12\operatorname{tr}S\) 的约定。对常数 \(\lambda>0\)，令 \(\widehat X=\lambda X\)，并在对应点保留同方向法向。求 \(\widehat g,\widehat h,\widehat H,\widehat K\) 与原数据的关系。再对球面北极图像
	\[
		X_R(u,v)=(u,v,\sqrt{R^2-u^2-v^2}),\qquad R>0,
	\]
	用本节的差分方法求 \(K_{R,h}\)，并证明其相对误差只依赖于无量纲比 \(h/R\)。据此说明为什么比较不同半径球面的误差时应报告步长与半径的比值。
\end{exercise}
\begin{solution}
	由 \(\widehat X_i=\lambda X_i\)、\(\widehat X_{ij}=\lambda X_{ij}\) 得到 \(\widehat g=\lambda^2g\)、\(\widehat h=\lambda h\)，故
	\[
		[\widehat S]=\widehat g^{-1}\widehat h=\lambda^{-1}[S],\qquad
		\widehat H=\lambda^{-1}H,\qquad \widehat K=\lambda^{-2}K.
	\]
	这也说明平均曲率和 Gauss 曲率分别具有长度的负一次和负二次量纲。

	在 \((0,0)\) 使用 \(0<h\le R/2\) 时，近似度量仍是单位阵，混合差分仍为零，两个纯二阶差分的法向分量为
	\[
		\frac{2(\sqrt{R^2-h^2}-R)}{h^2}
		=-\frac2{R+\sqrt{R^2-h^2}}.
	\]
	因此
	\[
		K_{R,h}=\frac4{(R+\sqrt{R^2-h^2})^2}
		=\frac1{R^2}\frac4{(1+\sqrt{1-(h/R)^2})^2}.
	\]
	真实曲率为 \(R^{-2}\)，相对误差是
	\[
		\frac{K_{R,h}-R^{-2}}{R^{-2}}
		=\frac4{(1+\sqrt{1-(h/R)^2})^2}-1.
	\]
	它只依赖于 \(h/R\)。同一个带长度单位的步长，对小球和大球代表不同的相对采样尺度；报告 \(h/R\) 才能在本例中作直接比较。
\end{solution}

% !TeX root = ../../Book1.tex
% 纲要标题：## 附录 D　G-结构与 Cartan 几何
% 纲要位置：几何学计划纲要.md，第 981 行。
% 正文状态：初稿完成，含完整证明、模型计算及习题解答；尚未正式审核。

% 纲要内容开始
% 
% > 沿标架丛约化发展结构方程；选定低阶模型给出完整计算，一般等价算法另属专题。
% 
% * **结构群约化与内蕴挠率**
%   * G-结构及适应标架
%   * 联络选择与内蕴挠率
%   * 正交、几乎复及接触模型
% 
% * **Cartan 联络与齐性模型**
%   * Cartan 联络的公理
%   * 曲率、平坦性及路径展开
%   * Cartan 联络与主联络
% 
% * **标架归一化与几何等价**
%   * 结构方程中的规范自由
%   * 曲率梯度与曲面标架的归一化
%   * 局部等价与整体单值
% 
% ---
% 
% 纲要内容结束

\chapter{G-结构与 Cartan 几何}
\label{g1:app:D}
\BookChapterNumber{g1:app:D}

\begin{chapterabstract}
	几何结构可以通过允许哪些标架来记录，而这些标架随位置变化时，又会产生挠率与曲率。这个附录把结构群约化推进到微分层面：先用联络的差说明内蕴挠率为何与联络选择无关，再通过正交、几乎复及接触模型比较不同约化的障碍。随后定义以齐性空间为模型的 Cartan 联络，证明其平坦性与局部模型实现的关系，并解释它和主丛水平联络的区别。最后以曲面曲率梯度固定标架，计算剩余结构函数，说明局部等价的检验及整体单值所记录的信息。
\end{chapterabstract}

\begin{learninggoals}
	\begin{itemize}
		\item 用切丛标架的约化描述几何结构，写出适应标架上的典范形式及其变换规律。
		\item 计算改变相容联络时挠率的变化，并在商空间中识别不能消去的部分。
		\item 比较正交、几乎复和接触模型，区分无挠、平坦与具有局部标准形。
		\item 核对 Cartan 联络的三项公理，计算其曲率，并说明平坦模型的局部构造。
		\item 通过展开方程比较路径与齐性模型，辨别 Cartan 联络和 Ehresmann 主联络所提供的数据。
		\item 在曲率梯度非零处归一化曲面标架，用结构方程检验局部等价，并辨认整体周期障碍。
	\end{itemize}
\end{learninggoals}

在平面开集上取两条处处独立的一形式
\[
	\theta^1=\mathrm{d}x,\qquad \theta^2=e^{x^2}\,\mathrm{d}y.
\]
它们构成一个余标架，因而可以用来记录切向量的两个分量。但若希望它们分别是两个新坐标的微分，就遇到了
\[
	\mathrm{d}\theta^2=2x\,\theta^1\wedge\theta^2.
\]
在 \(x\ne0\) 处，第二条形式不是闭形式，不能写成局部坐标函数的微分。把它乘以 \(e^{-x^2}\) 虽然恢复了 \(\mathrm{d}y\)，却改变了度量 \(g=(\theta^1)^2+(\theta^2)^2\)。如果希望保留这个度量，就只能用正交矩阵改变正交余标架，不能任意伸缩其中一条形式。

\ProseParagraphBreak

这说明“选择方便的标架”有一个先于计算的问题：哪些变化仍然保留所研究的结构？\chapref{g1:ch:12}和\chapref{g1:ch:18}用结构群约化回答了逐点的选择问题。现在还要比较这些标架在邻近点的变化。相容联络提供一种比较方法，但联络往往不唯一；改变联络能够消去哪些挠率分量，又有哪些分量属于几何结构本身，是本附录的第一个问题。

\ProseParagraphBreak

即使可以把挠率全部消去，几何也未必成为Euclid 几何。曲面上的无挠度量联络仍可有曲率。若把“允许的转动”和“沿曲面前进的位移”放入同一个Lie 代数值形式，便得到一种同时记录标架及切向运动的描述。Cartan 联络把这种构造推广到一般齐性模型；其平坦性表达局部与模型一致，整体是否能用同一份模型坐标描述，还取决于绕行和延拓。

\ProseParagraphBreak

本附录使用\chapref{g1:ch:12}的标架丛、关联丛与光滑单位分解，\chapref{g1:ch:13}的Lie 括号和 Frobenius 定理，以及\chapref{g1:ch:15}的挠率、曲率与结构方程。主丛、局部截面、联络形式和群值微分方程采用\chapref{g1:ch:18}的右作用约定；齐性空间的光滑结构接续\chapref{g1:ch:17}。曲面模型只需恢复\chapref{g1:ch:09}的正向正交标架及 \(\mathrm{d}\alpha=-K\,\mathrm{d}A\) 约定。各节将从所用对象的定义出发，把结构等价的问题化成适应标架与结构函数的比较。

\section{结构群约化与内蕴挠率}
\label{g1:sec:D01}

第十二章把度量、几乎复结构等逐点数据组织成标架丛的结构群约化，第十八章又说明主联络怎样在不同纤维之间作比较。把两件事放在一起，会遇到一个新的问题：既要求联络保持指定结构，又要求它无挠，这两项能否同时做到？第十五章对任意切丛联络作修正可以消去挠率，但修正后的联络未必还保持原来的几何结构。本节把允许的修正写成一个有限维线性映射，再在目标空间中对该映射的像取商，提取不随联络改变的部分。

\subsection{G-结构及适应标架}
\label{g1:subsec:D01:01}

本节中，\(M\) 是有限维、无边界、Hausdorff、第二可数的光滑流形，\(\dim M=n\geq1\)，\(V=\mathbb R^n\)，\(G\subset\operatorname{GL}(V)\) 是闭Lie 子群，Lie 代数为 \(\mathfrak g\)。沿用\cref{g1:def:frame-bundle}的约定，\(x\) 处的标架是线性同构 \(u:V\to\mathrm{T}_xM\)，右作用为 \(u g=u\circ g\)。指定标准基以后，\(u\) 也就是一个有序切向量基。

\begin{definition}
	\label{g1:def:g-structure}
	设 \(n\ge1\)，\(M\) 为 \(n\) 维、无边界、Hausdorff、第二可数光滑流形，\(V=\mathbb R^n\)，\(G\subset\operatorname{GL}(V)\) 为闭 Lie 子群。标架 \(u:V\to\mathrm T_xM\) 的右作用为 \(ug=u\circ g\)。\(M\) 上的一个 \term[Gjiegou]{\(G\)-结构}[\(G\)-structure]是标架丛中的光滑主 \(G\) 子丛
	\[
		\pi:P\subset\operatorname{Fr}(\mathrm{T}M)\longrightarrow M.
	\]
	它的局部截面称为\term[shiyingbiaojia]{适应标架}[adapted frame]。具体说，存在开覆盖 \(\{U_i\}\) 及光滑标架 \(s_i:U_i\to P\)，使
	\[
		P|_{U_i}=\{s_i(x)g:x\in U_i,\ g\in G\},
		\qquad s_j=s_i g_{ij},\quad g_{ij}:U_i\cap U_j\to G.
	\]
	这里要求局部表达 \((x,g)\mapsto s_i(x)g\) 给出主丛平凡化。
\end{definition}

这正是\cref{g1:def:frame-structure-reduction}在切丛上的主丛表述。若从局部适应标架及其 \(G\) 值过渡函数出发，上述集合在交叠处一致；用 \(U_i\times G\) 的坐标粘合，就得到所需主子丛。反过来，主子丛的局部截面给出这样的标架。因此，局部适应标架及其光滑的过渡关系共同确定一个约化。

\ProseParagraphBreak

标架还提供了仅由主丛结构无法得到的信息：它把固定空间 \(V\) 与底空间的切空间联系起来。这个联系可以记在 \(P\) 上的一形式中。对 \(u\in P_x\)、\(\xi\in\mathrm{T}_uP\)，定义
\begin{equation}
	\label{g1:eq:tautological-form}
	\theta_u(\xi)=u^{-1}\bigl(\mathrm{d}\pi_u\xi\bigr)\in V.
\end{equation}
称 \(\theta\in\Omega^1(P;V)\) 为标架约化上的\term[dianfanyixingshi]{典范一形式}[tautological form]。\symidx{\(\theta\)：标架丛及其约化上的典范一形式}它先将总空间切向量投到底空间，再用当前标架读取分量；这个定义没有选择联络。

\begin{proposition}[典范形式与余标架]
	\label{g1:prop:tautological-form-properties}
	设 \(M\) 为 \(n\) 维光滑流形，\(V=\mathbb R^n\)，\(G\subset\operatorname{GL}(V)\) 为闭 Lie 子群，\(\pi:P\subset\operatorname{Fr}(\mathrm{T}M)\to M\) 为 \(G\)-结构。对标架 \(u:V\to\mathrm T_xM\)，定义典范形式 \(\theta_u(\xi)=u^{-1}(\mathrm d\pi_u\xi)\)，并以 \(R_g(u)=ug\) 记右作用。则 \(\theta\) 光滑，每个 \(\theta_u:\mathrm{T}_uP\to V\) 满射，且
	\[
		\ker\theta_u=\ker\mathrm{d}\pi_u,\qquad
		R_g^*\theta=g^{-1}\theta\quad(g\in G).
	\]
	若 \(s=(e_1,\ldots,e_n)\) 是局部适应标架，则 \(s^*\theta\) 是其对偶余标架组成的 \(V\) 值一形式。
\end{proposition}
\begin{proof}
	在局部表达 \(u=s(x)g\) 下，\eqref{g1:eq:tautological-form}成为 \(g^{-1}s(x)^{-1}\mathrm{d}\pi\)，其矩阵系数光滑，所以 \(\theta\) 光滑。\(\pi\) 为淹没且 \(u^{-1}\) 可逆，故 \(\theta_u\) 满射，核恰为垂直空间。又因 \(\pi\circ R_g=\pi\)，有
	\[
		\theta_{ug}(\mathrm{d}R_g\xi)
		=(ug)^{-1}\mathrm{d}\pi_u\xi
		=g^{-1}\theta_u(\xi).
	\]
	最后，\(\pi\circ s=\operatorname{id}\) 给出 \((s^*\theta)_x(X)=s(x)^{-1}X\)。把 \(X=\sum_iX^ie_i\) 代入，结果正是列向量 \((X^1,\ldots,X^n)^{\mathsf T}\)，所以各分量为对偶余标架。
\end{proof}

在底空间上，\(s^*\theta\) 的 \(n\) 个分量确是余标架；在 \(P\) 上，\(\theta\) 却有垂直核，通常不是总空间的余标架，因为 \(\dim P=n+\dim G\)。它也不是主联络形式：典范形式取值于 \(V\)，并在所有垂直向量上为零；主联络形式取值于 \(\mathfrak g\)，需要读出垂直方向中的群运动。下一节的 Cartan 联络将同时处理这两种信息。

\subsection{联络选择与内蕴挠率}
\label{g1:subsec:D01:02}

\cref{g1:thm:principal-connection-existence}保证 \(P\) 上存在主联络。固定一个主联络形式 \(\omega\in\Omega^1(P;\mathfrak g)\)，其公理为
\[
	\omega(A^\#)=A,\qquad R_g^*\omega=\operatorname{Ad}_{g^{-1}}\omega.
\]
这里 \(A^\#_u=\left.\frac{\mathrm{d}}{\mathrm{d}t}\right|_{0}u\exp(tA)\) 是 \(A\in\mathfrak g\) 产生的垂直基本向量。标准表示给出
\[
	P\times_G V\longrightarrow\mathrm{T}M,\qquad [u,v]\longmapsto u(v).
\]
它与\cref{g1:prop:frame-associated-equivalence}中的映射相同：\([ug,g^{-1}v]\) 的像仍为 \(u(v)\)，在适应标架中则是逐纤维恒等坐标，故为光滑向量丛同构。由\cref{g1:thm:associated-covariant-derivative}，\(\omega\) 因而诱导切丛联络。若 \(a=s^*\omega\)，则其局部公式为 \(\nabla(s f)=s(\mathrm{d}f+a f)\)。这样的联络称为适应该结构的 \(G\)-联络；局部联络矩阵取值于 \(\mathfrak g\)，正是允许的标架变换所要求的限制。

\begin{definition}
	\label{g1:def:g-structure-torsion-form}
	设 \(M\) 为 \(n\) 维光滑流形，\(V=\mathbb R^n\)，\(G\subset\operatorname{GL}(V)\) 为闭 Lie 子群，Lie 代数为 \(\mathfrak g\)，\(\pi:P\subset\operatorname{Fr}(\mathrm TM)\to M\) 为 \(G\)-结构。令 \(\omega\in\Omega^1(P;\mathfrak g)\) 为主 \(G\) 联络，\(\theta_u(\xi)=u^{-1}(\mathrm d\pi_u\xi)\) 为典范 \(V\) 值形式。定义 \(\omega\) 的挠率形式为
	\begin{equation}
		\label{g1:eq:g-structure-torsion-form}
		T^\omega=\mathrm{d}\theta+\omega\wedge\theta\in\Omega^2(P;V),
	\end{equation}
	其中矩阵对向量的作用约定为
	\[
		(\omega\wedge\theta)(\xi,\eta)
		=\omega(\xi)\theta(\eta)-\omega(\eta)\theta(\xi).
	\]
\end{definition}

\begin{proposition}[挠率的水平性]
	\label{g1:prop:g-torsion-horizontal}
	设 \(M\) 为 \(n\) 维光滑流形，\(V=\mathbb R^n\)，\(G\subset\operatorname{GL}(V)\) 为闭 Lie 子群，\(\pi:P\subset\operatorname{Fr}(\mathrm TM)\to M\) 为 \(G\)-结构，\(\omega\) 为其主 \(G\) 联络。令 \(\theta_u=u^{-1}\mathrm d\pi_u\)、\(T^\omega=\mathrm d\theta+\omega\wedge\theta\)，并通过 \([u,v]\mapsto u(v)\) 将标准表示诱导的关联丛联络视为切丛联络 \(\nabla\)。则 \(T^\omega\) 是水平且等变的，即只要一个输入为垂直向量，其值便为零，并且 \(R_g^*T^\omega=g^{-1}T^\omega\)。对光滑向量场 \(X,Y\)，它对应的切丛值二形式正是诱导联络的挠率
	\[
		T^\nabla(X,Y)=\nabla_XY-\nabla_YX-[X,Y].
	\]
\end{proposition}
\begin{proof}
	由典范形式的等变性，沿 \(A^\#\) 的流求导得到 \(\mathcal L_{A^\#}\theta=-A\theta\)。又有 \(\iota_{A^\#}\theta=0\)，所以对每个分量使用 Cartan 公式，得
	\[
		\iota_{A^\#}\mathrm{d}\theta=-A\theta,
		\qquad
		\iota_{A^\#}(\omega\wedge\theta)=A\theta.
	\]
	两项相消，而各点的垂直空间由基本向量填满，故 \(T^\omega\) 水平。对于固定 \(g\)，外微分与拉回交换，且
	\[
		(g^{-1}\omega g)\wedge(g^{-1}\theta)
		=g^{-1}(\omega\wedge\theta),
	\]
	所以得出等变性。

	再令 \(\vartheta=s^*\theta\)、\(a=s^*\omega\)。拉回后
	\[
		s^*T^\omega=\mathrm{d}\vartheta+a\wedge\vartheta.
	\]
	对底空间向量场 \(X,Y\)，其右边等于
	\[
		X(\vartheta(Y))-Y(\vartheta(X))-\vartheta([X,Y])
		+a(X)\vartheta(Y)-a(Y)\vartheta(X).
	\]
	利用局部协变导数公式，这就是 \(\vartheta(\nabla_XY-\nabla_YX-[X,Y])\)。它同时复现了\cref{g1:prop:first-structure-equation}中的第一结构方程。
\end{proof}

水平性使挠率可以完全由底空间方向的标架分量描述。令
\[
	W=\Lambda^2V^*\otimes V.
\]
每个 \(u\in P\) 唯一确定 \(t^\omega(u)\in W\)，使
\begin{equation}
	\label{g1:eq:g-torsion-function}
	T^\omega_u(\xi,\eta)
	=t^\omega(u)\bigl(\theta_u\xi,\theta_u\eta\bigr).
\end{equation}
确实，\(\theta_u\) 满射；对给定 \(v,w\in V\) 任选提升，两个不同提升之差垂直，故 \(T^\omega\) 的值相同。在局部适应标架下，此定义就是把光滑挠率张量取分量，因而 \(t^\omega\) 光滑。称它为相对于联络 \(\omega\) 的挠率函数。

\ProseParagraphBreak

现在改变联络。由\cref{g1:prop:principal-connection-affine}，两个主联络之差 \(\beta=\omega'-\omega\) 是水平的伴随等变一形式。因此唯一存在光滑函数
\[
	A:P\longrightarrow C\coloneqq V^*\otimes\mathfrak g,
	\qquad \beta_u=A(u)\circ\theta_u.
\]
这里 \(A(u)\) 把一个底空间方向的标架分量映成允许的无穷小标架变换。把这个变换依次作用在另一个方向上，再作交替化，就得到线性映射
\begin{equation}
	\label{g1:eq:intrinsic-torsion-delta}
	\delta:C\longrightarrow W,\qquad
	(\delta A)(v,w)=A(v)w-A(w)v.
\end{equation}

\begin{proposition}[改变联络时的挠率变化]
	\label{g1:prop:g-torsion-change}
	设 \(M\) 为 \(n\) 维光滑流形，\(V=\mathbb R^n\)，\(G\subset\operatorname{GL}(V)\) 为闭 Lie 子群，Lie 代数为 \(\mathfrak g\)，\(\pi:P\subset\operatorname{Fr}(\mathrm TM)\to M\) 为 \(G\)-结构，\(\theta_u=u^{-1}\mathrm d\pi_u\) 为典范形式。对 \(P\) 上的两主联络 \(\omega,\omega'\)，设 \(A:P\to V^*\otimes\mathfrak g\) 满足 \((\omega'-\omega)_u=A(u)\circ\theta_u\)。对任一主联络 \(\eta\)，令 \(T^\eta=\mathrm d\theta+\eta\wedge\theta\)，并由
	\[
		T^\eta_u(\xi,\zeta)=t^\eta(u)(\theta_u\xi,\theta_u\zeta)
	\]
	定义 \(t^\eta:P\to\Lambda^2V^*\otimes V\)。定义交替化映射 \((\delta B)(v,w)=B(v)w-B(w)v\)，其中 \(B\in V^*\otimes\mathfrak g\)。则
	\begin{equation}
		\label{g1:eq:g-torsion-change}
		t^{\omega'}=t^\omega+\delta A.
	\end{equation}
	在表示
	\[
		(g\cdot t)(v,w)=g\,t(g^{-1}v,g^{-1}w),
		\qquad
		(g\cdot A)(v)=g A(g^{-1}v)g^{-1}
	\]
	下，\(\delta\) 是 \(G\) 等变映射，而 \(t^\omega(ug)=g^{-1}\cdot t^\omega(u)\)、\(A(ug)=g^{-1}\cdot A(u)\)。
\end{proposition}
\begin{proof}
	将两个挠率公式相减，得 \(T^{\omega'}-T^\omega=\beta\wedge\theta\)。在 \(\theta\xi=v\)、\(\theta\eta=w\) 的两个输入上，右边为 \(A(v)w-A(w)v\)，得到\eqref{g1:eq:g-torsion-change}。

	等变性可直接核对。例如
	\[
		\bigl(\delta(g\cdot A)\bigr)(v,w)
		=g\bigl(A(g^{-1}v)g^{-1}w-A(g^{-1}w)g^{-1}v\bigr)
		=(g\cdot\delta A)(v,w).
	\]
	\(T^\omega\) 与 \(\theta\) 的等变性给出
	\[
		t^\omega(ug)(v,w)=g^{-1}t^\omega(u)(gv,gw),
	\]
	即所述公式。对 \(\beta=A\theta\) 使用 \(R_g^*\beta=g^{-1}\beta g\) 同样得到 \(A\) 的变换式。
\end{proof}

\begin{definition}
	\label{g1:def:intrinsic-torsion}
	设 \(n\ge1\)，\(M\) 为 \(n\) 维、无边界、Hausdorff、第二可数光滑流形，\(V=\mathbb R^n\)，\(G\subset\operatorname{GL}(V)\) 为闭 Lie 子群，Lie 代数为 \(\mathfrak g\)，\(\pi:P\subset\operatorname{Fr}(\mathrm TM)\to M\) 为 \(G\)-结构。任取主 \(G\) 联络 \(\omega\)，以 \(\theta_u=u^{-1}\mathrm d\pi_u\) 记典范形式，并由
	\[
		(\mathrm d\theta+\omega\wedge\theta)_u(\xi,\zeta)
		=t^\omega(u)(\theta_u\xi,\theta_u\zeta)
	\]
	定义挠率函数 \(t^\omega:P\to\Lambda^2V^*\otimes V\)。令 \(\delta:V^*\otimes\mathfrak g\to\Lambda^2V^*\otimes V\) 满足 \((\delta A)(v,w)=A(v)w-A(w)v\)，并记
	\[
		W_G=\frac{\Lambda^2V^*\otimes V}{\delta(V^*\otimes\mathfrak g)}.
	\]
	\(G\)-结构 \(P\) 的\term[neiyunnaolv]{内蕴挠率}[intrinsic torsion]是等变函数
	\[
		\tau_P:P\longrightarrow W_G,\qquad
		\tau_P(u)=[t^\omega(u)],
	\]
	此商类与主联络 \(\omega\) 的选择无关，也可视为关联丛 \(P\times_G W_G\to M\) 的光滑截面。\symidx{\(\tau_P\)：\(G\)-结构的内蕴挠率}
\end{definition}

由\cref{g1:prop:g-torsion-change}，更换联络只会增加 \(\delta\) 的像中的项，因此这个商类与联络无关。\(\delta\) 等变又保证其像为不变子空间，商空间具有良定义的 \(G\) 作用；换适应标架时，各局部商类正按这个表示变换。因此内蕴挠率是该约化本身的几何数据。它是商空间中的类；若要用 \(W\) 中某个补空间的张量来代表它，需再选取与换标架相容的补空间。

\begin{theorem}[无挠适应联络的存在]
	\label{g1:thm:intrinsic-torsion-zero}
	设 \(n\ge1\)，\(M\) 为 \(n\) 维、无边界、Hausdorff、第二可数光滑流形，\(G\subset\operatorname{GL}(n,\mathbb R)\) 为闭Lie 子群，\(\pi:P\subset\operatorname{Fr}(\mathrm{T}M)\to M\) 为 \(G\)-结构。记 \(V=\mathbb R^n\)、\(\mathfrak g=\operatorname{Lie}(G)\)，令 \(\delta:V^*\otimes\mathfrak g\to\Lambda^2V^*\otimes V\) 满足 \((\delta A)(v,w)=A(v)w-A(w)v\)。对任一主联络 \(\omega\)，以 \(\theta_u=u^{-1}\mathrm d\pi_u\) 记典范形式，由 \((\mathrm d\theta+\omega\wedge\theta)_u(\xi,\zeta)=t^\omega(u)(\theta_u\xi,\theta_u\zeta)\) 定义挠率函数，并令 \(\tau_P(u)=[t^\omega(u)]\in\operatorname{coker}\delta\)。则 \(\tau_P=0\) 当且仅当 \(P\) 存在无挠 \(G\)-联络。若存在，则所有这样的联络组成仿射空间，其差由 \(P\times_G\ker\delta\) 的光滑截面给出。
\end{theorem}
\begin{proof}
	若存在无挠联络，其挠率函数为零，故内蕴挠率为零。反向，先由\cref{g1:thm:principal-connection-existence}取一个全局主联络 \(\omega\)。若 \(\tau_P=0\)，则 \(t^\omega(u)\in\operatorname{im}\delta\) 对每个 \(u\) 成立。

	取局部适应标架 \(s_i:U_i\to P\)。有限维满射 \(\delta:C\to\operatorname{im}\delta\) 有线性右逆 \(R:\operatorname{im}\delta\to C\)：在 \(C\) 中选取 \(\ker\delta\) 的一个线性补空间，将 \(\delta\) 在该补空间上的限制求逆即可。定义
	\[
		A_i(s_i(x))=-R\bigl(t^\omega(s_i(x))\bigr),
		\qquad
		A_i(s_i(x)g)=g^{-1}\cdot A_i(s_i(x)).
	\]
	平凡化中的 \((x,g)\) 唯一，所以后一式给出光滑、等变的延拓。由 \(\delta\) 等变性，\(\delta A_i=-t^\omega\) 在整个 \(P|_{U_i}\) 成立。于是 \(\beta_i=A_i\theta\) 水平且伴随等变，\(\omega_i=\omega+\beta_i\) 是 \(P|_{U_i}\) 上的主联络，且由\eqref{g1:eq:g-torsion-change}无挠。这里右逆 \(R\) 只用于一个局部平凡化，并没有要求它对 \(G\) 等变。

	取从属于 \(\{U_i\}\) 的局部有限光滑单位分解 \(\{\chi_i\}\)，其支集包含于相应 \(U_i\)。将 \((\chi_i\circ\pi)\beta_i\) 在 \(P|_{U_i}\) 之外补零；支集条件保证补零后仍光滑。定义
	\[
		\omega_0=\omega+\sum_i(\chi_i\circ\pi)\beta_i.
	\]
	该和局部有限，各项均水平且伴随等变，故 \(\omega_0\) 是全局主联络。在每一点附近也可写成仿射组合 \(\sum_i(\chi_i\circ\pi)\omega_i\)，权重和为一，因而
	\[
		T^{\omega_0}
		=\sum_i(\chi_i\circ\pi)
		\bigl(\mathrm{d}\theta+\omega_i\wedge\theta\bigr)=0.
	\]
	这一步没有权重的导数项，因为挠率公式中对 \(\theta\) 求外微分，对联络形式只作代数运算。于是得到全局无挠适应联络。

	最后，固定一个无挠联络后，另一个联络由 \(\beta=A\theta\) 给出。它也无挠恰好等价于 \(\delta A=0\)。等变函数 \(A:P\to\ker\delta\) 与关联丛 \(P\times_G\ker\delta\) 的截面一一对应，这证明最后一项。
\end{proof}

这个结论同时区分了存在性和唯一性：商空间中的类控制能否消去挠率，\(\ker\delta\) 则控制消去以后还有多少联络可选。即使无挠联络存在，它的曲率也未必为零；即使某类几何具有统一的局部标准模型，该模型的内蕴挠率也未必为零。下面的例子将分别说明这些区别。

\subsection{正交、几乎复及接触模型}
\label{g1:subsec:D01:03}

先设 \(g\) 是 \(M\) 上的Riemann 度量。由\cref{g1:prop:metric-orthogonal-reduction}，正交单位标架组成 \(\operatorname{O}(n)\)-结构。在这样的标架中，度量矩阵为 \(I\)，\cref{g1:eq:metric-connection-matrix}化为 \(a^{\mathsf T}+a=0\)。因此 \(\operatorname{O}(n)\)-联络恰为保持 \(g\) 的度量联络。

\begin{proposition}[正交约化的挠率修正]
	\label{g1:prop:orthogonal-delta-isomorphism}
	设 \(V\) 是 \(n\) 维实内积空间。线性映射
	\[
		\delta:V^*\otimes\mathfrak{so}(V)\longrightarrow\Lambda^2V^*\otimes V
	\]
	由 \((\delta A)(v,w)=A(v)w-A(w)v\) 定义，则它是同构。因此，任意 \(n\) 维光滑流形上的 Riemann 度量所确定的正交标架约化具有恒为零的内蕴挠率，并且有唯一的无挠适应联络。
\end{proposition}
\begin{proof}
	先证单射。设 \(\delta A=0\)，令
	\[
		a(v,w,z)=\langle A(v)w,z\rangle.
	\]
	\(\delta A=0\) 说明 \(a\) 的前两个位置对称；\(A(v)\in\mathfrak{so}(V)\) 说明后两个位置反对称。依次使用这两种关系，得到
	\[
	\begin{aligned}
		a(v,w,z)&=a(w,v,z)=-a(w,z,v)=-a(z,w,v)\\
		&=a(z,v,w)=a(v,z,w)=-a(v,w,z).
	\end{aligned}
	\]
	故 \(a=0\)，由内积非退化性得 \(A=0\)。两边的维数均为 \(n^2(n-1)/2\)，所以单射即为同构。于是 \(W_{\operatorname{O}(n)}=0\) 且 \(\ker\delta=0\)，由\cref{g1:thm:intrinsic-torsion-zero}分别得到存在性与唯一性。
\end{proof}

这也从结构群约化的角度解释了\chapref{g1:ch:19}所用的 Levi--Civita 联络：斜对称矩阵的自由度恰能唯一地修正全部挠率。内蕴挠率在所有Riemann 度量上都为零，所以它不负责区分平面、球面等不同曲率的几何；那些差别还须由这个唯一联络的曲率来测量。

\ProseParagraphBreak

设 \(J:\mathrm{T}M\to\mathrm{T}M\) 光滑且 \(J^2=-I\)，\(\dim M=2m\)。\cref{g1:prop:complex-frame-reduction}给出 \(\operatorname{GL}(m,\mathbb C)\) 约化，相应的实适应标架可以写成 \(e_1,Je_1,\ldots,e_m,Je_m\)。在适应标架中 \(J\) 是常矩阵 \(J_0\)，所以
\[
	\nabla J=0\quad\Longleftrightarrow\quad aJ_0=J_0a.
\]
右边恰好要求 \(a\) 取值于 \(\mathfrak{gl}(m,\mathbb C)\) 的实表示。因此这里的适应联络就是保持 \(J\) 的联络。

\ProseParagraphBreak

复化后的 Lie 括号直接给出一种在改变适应联络时保持不变的挠率分量。把切丛和 \(J\) 复化，令
\[
	\mathrm{T}^{1,0}M=\ker(J-\mathrm iI),\qquad
	\mathrm{T}^{0,1}M=\ker(J+\mathrm iI).
\]
投影 \(\operatorname{pr}^{1,0}=\tfrac12(I-\mathrm iJ)\)、\(\operatorname{pr}^{0,1}=\tfrac12(I+\mathrm iJ)\) 光滑且互补，两个像均为复秩 \(m\) 的子丛。对局部 \(\mathrm{T}^{1,0}M\) 截面 \(Z,W\)，定义
\begin{equation}
	\label{g1:eq:almost-complex-bracket-obstruction}
	\mathcal K_J(Z,W)=\operatorname{pr}^{0,1}[Z,W].
\end{equation}
这里括号按复双线性延拓。若在输入中乘以光滑复函数，括号的额外导数项仍在 \(\mathrm{T}^{1,0}M\)，投影后消失；例如 \(\operatorname{pr}^{0,1}[fZ,W]=f\operatorname{pr}^{0,1}[Z,W]\)。所以 \(\mathcal K_J\) 是良定义的光滑张量，测量两个 \((1,0)\) 方向的括号落出该子丛的部分。

\begin{proposition}[几乎复结构的括号障碍]
	\label{g1:prop:almost-complex-intrinsic-component}
	设 \(M\) 为实 \(2m\) 维光滑流形，\(J\in\Gamma(\operatorname{End}(\mathrm TM))\) 光滑且 \(J^2=-I\)，\(\nabla\) 为 \(\mathrm TM\) 上满足 \(\nabla J=0\) 的线性联络，挠率为 \(T^\nabla(X,Y)=\nabla_XY-\nabla_YX-[X,Y]\)。复化切丛后，令 \(\mathrm T^{1,0}M=\ker(J-\mathrm iI)\)、\(\operatorname{pr}^{0,1}=(I+\mathrm iJ)/2\)，并对光滑截面 \(Z,W\) 定义 \(\mathcal K_J(Z,W)=\operatorname{pr}^{0,1}[Z,W]\)。将联络及其挠率复化，则对 \(Z,W\in\Gamma(\mathrm{T}^{1,0}M)\)，有
	\[
		\operatorname{pr}^{0,1}T^\nabla(Z,W)=-\mathcal K_J(Z,W).
	\]
	因此右边是其 \(\operatorname{GL}(m,\mathbb C)\)-结构的内蕴挠率中不随适应联络改变的一个分量；若 \(\mathcal K_J\neq0\)，则不存在无挠适应联络。
\end{proposition}
\begin{proof}
	因为 \(\nabla J=0\)，由 \(JW=\mathrm iW\) 求协变导数可得 \(J\nabla_ZW=\mathrm i\nabla_ZW\)。所以 \(\nabla_ZW\) 与 \(\nabla_WZ\) 都在 \(\mathrm{T}^{1,0}M\) 中。将挠率恒等式投影到 \(\mathrm{T}^{0,1}M\)，这两项消失，只留下 \(-\operatorname{pr}^{0,1}[Z,W]\)。这也说明改变适应联络所产生的 \(\delta A\) 在这个投影下为零，故该分量确由内蕴商类决定。
\end{proof}

若 \(J\) 来自一组复坐标，那么 \(\mathrm{T}^{1,0}M\) 局部由相互对易的复坐标向量场张成，其括号仍落在本子丛，故 \(\mathcal K_J=0\)。反过来从括号条件构造复坐标，需要额外的可积性定理；非零的括号障碍则已经排除了复坐标的存在，下面就计算一个这样的例子。

\begin{example}[不能消去的几乎复挠率分量]
	\label{g1:exm:almost-complex-nonzero-intrinsic}
	在 \(\mathbb R^4\) 的坐标 \((x,y,z,t)\) 中取全局标架
	\[
		e_1=\partial_x,\quad e_2=\partial_y,\quad
		e_3=\partial_z,\quad e_4=\partial_t+x\partial_y.
	\]
	规定 \(Je_1=e_2\)、\(Je_2=-e_1\)、\(Je_3=e_4\)、\(Je_4=-e_3\)，得到光滑几乎复结构。唯一非零的基本括号为 \([e_1,e_4]=e_2\) 及其反向。令
	\[
		Z=e_1-\mathrm i e_2,\qquad W=e_3-\mathrm i e_4,
	\]
	它们属于 \(\mathrm{T}^{1,0}M\)，且
	\[
		[Z,W]=-\mathrm i e_2=\frac12(Z-\overline Z),
		\qquad \mathcal K_J(Z,W)=-\frac12\overline Z\neq0.
	\]
	因此此结构没有无挠适应联络。适应联络本身并不难构造：规定 \(\nabla e_i=0\) 就保持了 \(J\)，但其挠率满足 \(T^\nabla(e_1,e_4)=-e_2\)。换用其他适应联络能够改变某些分量，却不能消去上面算出的 \((0,1)\) 投影。
\end{example}

最后回到\subsecref{g1:subsec:1305:04}的三维接触分布。设 \(D\subset\mathrm{T}M\) 为光滑秩二分布，令 \(V_0=\operatorname{span}\{\varepsilon_1,\varepsilon_2\}\subset\mathbb R^3\)。保留此标准平面的群为
\[
	G_D=\left\{\begin{pmatrix}B&b\\0&c\end{pmatrix}:
	B\in\operatorname{GL}(2,\mathbb R),\ b\in\mathbb R^2,\ c\neq0\right\}.
\]
局部取 \(D\) 的两个光滑基截面，再补一个横向向量，得到适应标架。因此
\[
	P_D=\{u\in\operatorname{Fr}(\mathrm{T}M):u(V_0)=D_{\pi(u)}\}
\]
是光滑 \(G_D\)-结构。其适应联络恰为保持分布的联络，即 \(Y\in\Gamma(D)\) 时，对任意 \(X\) 都有 \(\nabla_XY\in\Gamma(D)\)；这可由局部联络矩阵的左下块为零直接读出。

\ProseParagraphBreak

前文已经定义括号张量
\[
	\mathcal B_x(v,w)=[X,Y]_x\pmod{D_x},
	\qquad v,w\in D_x,\quad X_x=v,\ Y_x=w,
\]
其中 \(X,Y\) 为局部 \(D\) 截面。其值与延拓无关，因为 \([fX,Y]=f[X,Y]-Y(f)X\) 中的导数项仍属于 \(D\)。它是 \(\Lambda^2D^*\otimes(\mathrm{T}M/D)\) 的光滑截面，接触条件要求它在每点非零。对保持 \(D\) 的任意联络，挠率公式给出
\begin{equation}
	\label{g1:eq:distribution-intrinsic-torsion}
	T^\nabla(X,Y)\pmod D=-\mathcal B(X,Y),\qquad X,Y\in\Gamma(D),
\end{equation}
因为两项协变导数都在 \(D\) 中。从 \(\delta\) 也能看出同一事实：若 \(v,w\in V_0\)，而 \(A\) 取值于保持 \(V_0\) 的Lie 代数，则 \(\delta A(v,w)\in V_0\)，其在 \(V/V_0\) 中的分量为零。因此 \(-\mathcal B\) 是这个约化内蕴挠率的一个固定分量；接触分布的适应联络必定有挠。

\begin{example}[标准接触模型的挠率]
	\label{g1:exm:contact-intrinsic-torsion}
	沿用第十三章的
	\[
		\alpha=\mathrm{d}z-x\,\mathrm{d}y,\qquad
		A=\partial_x,\quad B=\partial_y+x\partial_z,
	\]
	并记 \(C=\partial_z\)。\((A,B,C)\) 是 \(D=\ker\alpha\) 的适应标架，对偶余标架为 \((\mathrm{d}x,\mathrm{d}y,\alpha)\)。\cref{g1:prop:contact-model-nonintegrable}已经证明 \([A,B]=C\)；故 \(\mathcal B(A,B)=[C]\neq0\)。令这个标架平行，便得到保持 \(D\) 的联络，其局部联络矩阵为零。因此第一结构方程给出挠率的分量列
	\[
		\begin{pmatrix}0\\0\\-\mathrm{d}x\wedge\mathrm{d}y\end{pmatrix},
		\qquad T^\nabla(A,B)=-C,
	\]
	与\eqref{g1:eq:distribution-intrinsic-torsion}完全相符。同时
	\[
		\mathrm{d}\alpha=-\mathrm{d}x\wedge\mathrm{d}y,
		\qquad
		\alpha\wedge\mathrm{d}\alpha
		=-\mathrm{d}x\wedge\mathrm{d}y\wedge\mathrm{d}z\neq0.
	\]
	这个已经写成标准形式的接触模型仍具有非零内蕴挠率。它的平面分布不能误认为坐标水平面 \(\ker\mathrm{d}z\)，后者的括号张量为零。
\end{example}

还要区分固定分布与固定定义形式。如果只要求 \(\ker\alpha=D\)，就允许将 \(\alpha\) 换成任意处处非零光滑倍数 \(f\alpha\)，上述群中的 \(c\) 可以取任意非零值。若把一个具体的 \(\alpha\) 也作为结构的一部分，并要求 \(u^*\alpha=\varepsilon^3\)，则适应标架间的变换必须满足 \(c=1\)。两者允许的结构群不同，因而允许的联络修正也不同。判断内蕴挠率之前，必须先明确几何结构究竟保留哪些数据。

\section{Cartan 联络与齐性模型}
\label{g1:sec:D02}

Euclid 平面上一副正向正交标架既有位置，也有方向。沿曲线移动这副标架时，需要同时记录位移与转动；只记录转动的普通主联络，没有把两类信息合在一起。Cartan 联络所做的事情，是把标架总空间在每一点的全部速度，与一个模型变换群的Lie 代数逐点对应。模型可以是Euclid 平面，也可以是球面或其他齐性空间；曲率衡量这种逐点对应能否在邻域中积分成与模型的实际等价。

\subsection{Cartan 联络的公理}
\label{g1:subsec:D02:01}

设 \(L\) 是有限维实Lie 群，\(H\subseteq L\) 是闭子群，Lie 代数分别为 \(\mathfrak l,\mathfrak h\)。由\cref{g1:thm:homogeneous-space}，商 \(L/H\) 是光滑流形，\(q:L\to L/H\) 是主右 \(H\) 丛，并在每个点附近有光滑局部截面。这里 \(L\) 描述模型的运动，\(H\) 描述固定模型基点 \(o=eH\) 的运动；二者的作用不同。以下流形均无边界、Hausdorff 且第二可数。

\ProseParagraphBreak

先看模型总空间 \(L\) 自身。它的左 Maurer--Cartan 形式为
\[
	(\theta_L)_a=(\mathrm{d}L_{a^{-1}})_a:
	\mathrm{T}_aL\longrightarrow\mathfrak l.
\]
它用左平移把每个速度搬回单位元，是逐点线性同构。若只把 \(L\) 看成主 \(H\) 丛，垂直速度为
\[
	X_a^\#=\left.\frac{\mathrm{d}}{\mathrm{d}t}\right|_{t=0}
	a\exp(tX),\qquad X\in\mathfrak h,
\]
所以 \(\theta_L(X^\#)=X\)。\cref{g1:prop:maurer-cartan-product}给出的乘法公式又说明 \(R_h^*\theta_L=\operatorname{Ad}_{h^{-1}}\theta_L\)。把这三项性质保留下来，就能在一个未必齐性的空间上规定相同类型的无穷小模型。

\begin{definition}[Cartan 联络]\label{g1:def:cartan-connection}
	设 \(L\) 为有限维实Lie 群，\(H\subseteq L\) 为闭子群，\(\pi:P\to M\) 为无边界光滑流形上的光滑主右 \(H\) 丛。若 \(\mathfrak l\) 值光滑一形式 \(\omega\in\Omega^1(P;\mathfrak l)\) 满足下列条件，就称它为以 \(L/H\) 为模型的\term[cartanlianluo]{Cartan 联络}[Cartan connection]：
	\begin{enumerate}
		\item 对每个 \(p\in P\)，\(\omega_p:\mathrm{T}_pP\to\mathfrak l\) 是线性同构；
		\item 对每个 \(h\in H\)，有 \(R_h^*\omega=\operatorname{Ad}_{h^{-1}}\omega\)；
		\item 对每个 \(X\in\mathfrak h\)，有 \(\omega(X^\#)=X\)，其中 \(X^\#_p=\left.\frac{\mathrm{d}}{\mathrm{d}t}\right|_0p\exp(tX)\)。
	\end{enumerate}
	二元组 \((P,\omega)\) 称为一个模型为 \(L/H\) 的\term[cartanjihe]{Cartan 几何}[Cartan geometry]。两个 Cartan 几何之间的同构是覆盖底空间微分同胚的主 \(H\) 丛同构；沿它拉回目标的 Cartan 联络，应当得到源的 Cartan 联络。
\end{definition}

第一条公理称为\term[jueduipingxing]{绝对平行}[absolute parallelism]条件：固定 \(\mathfrak l\) 的一组基，便得到 \(P\) 上一组处处线性无关的向量场。这个名称不表示这些向量场交换，也不表示某个已经给定的协变导数使它们平行。第二条说明改变标架后的速度记录如何变化，第三条则把纯粹的标架变化与 \(\mathfrak h\) 中的元素对应。三条公理同时满足时，沿底空间运动的信息也已经包含在 \(\omega\) 中。

\begin{proposition}[Cartan 联络的焊接同构]\label{g1:prop:cartan-soldering}
	设 \((P\to M,\omega)\) 为模型 \(L/H\) 的光滑 Cartan 几何，\(L\) 为有限维实Lie 群，\(H\) 为闭子群。则
	\[
		\dim M=\dim L-\dim H=\dim(L/H),\qquad
		\omega_p^{-1}(\mathfrak h)=V_pP.
	\]
	商表示 \(h\cdot(X+\mathfrak h)=\operatorname{Ad}_hX+\mathfrak h\) 给出自然的向量丛同构
	\begin{equation}\label{g1:eq:cartan-soldering}
		\mathrm{T}M\cong P\times_H(\mathfrak l/\mathfrak h).
	\end{equation}
\end{proposition}
\begin{proof}
	绝对平行条件给出 \(\dim P=\dim\mathfrak l\)，主丛局部乘积给出 \(\dim P=\dim M+\dim H\)，从而得到维数等式。由\cref{g1:prop:principal-fundamental-fields}，\(X\mapsto X_p^\#\) 是 \(\mathfrak h\to V_pP\) 的同构；再现条件说明 \(\omega_p\) 在这两个子空间之间正是它的逆。由于 \(\omega_p\) 在整个切空间上也单射，\(\mathfrak h\) 的原像恰为 \(V_pP\)。

	给定 \(v\in\mathrm{T}_xM\) 和 \(p\in P_x\)，选取 \(\widetilde v\in\mathrm{T}_pP\) 满足 \((\mathrm{d}\pi)_p\widetilde v=v\)，定义
	\[
		\overline\omega_p(v)=\omega_p(\widetilde v)+\mathfrak h.
	\]
	两个提升之差垂直，故其 \(\omega\) 值之差在 \(\mathfrak h\) 中，这一定义与提升无关。前段的原像等式说明 \(\overline\omega_p\) 单射；维数相同，故它是同构。

	因为 \(\operatorname{Ad}_h\mathfrak h=\mathfrak h\)，商表示有定义。等变性给出
	\[
		\overline\omega_{ph}(v)
		=\operatorname{Ad}_{h^{-1}}\overline\omega_p(v).
	\]
	关联丛约定为 \([ph,\xi]=[p,h\cdot\xi]\)，见\cref{g1:thm:associated-vector-bundle}，所以
	\[
		v\longmapsto[p,\overline\omega_p(v)]
	\]
	与 \(p\) 的选择无关。它在局部截面中由光滑可逆矩阵给出，因此是所求光滑向量丛同构。
\end{proof}

这里的商值一形式把抽象模型方向 \(\mathfrak l/\mathfrak h\) 与真实切向量联系起来，因而称为焊接形式。它没有要求先选一个补空间 \(\mathfrak l=\mathfrak h\oplus\mathfrak m\)；商空间本身已经足够。特别地，Cartan 联络不是任意主 \(H\) 丛上都能附加的数据，底空间维数和切丛的关联表示必须相容。

\ProseParagraphBreak

为了计算，取局部截面 \(s:U\to P\)，记 \(B=s^*\omega\)，并用 \(\Psi(x,h)=s(x)h\) 平凡化。与\cref{g1:thm:principal-connection-local-data}中的推导一样，把 \(\mathrm{d}\Psi(v,\eta)\) 拆成 \((\mathrm{d}R_h)\mathrm{d}s(v)\) 和垂直部分；分别用等变性与再现条件，便得到
\begin{equation}\label{g1:eq:cartan-local-form}
	\Psi^*\omega
	=\operatorname{Ad}_{h^{-1}}B+\theta_H,
	\qquad
	B\bmod\mathfrak h:
	\mathrm{T}_xU\longrightarrow\mathfrak l/\mathfrak h
	\text{ 为同构}.
\end{equation}
这里 \(\theta_H\) 是 \(H\) 的左 Maurer--Cartan 形式，通过包含映射 \(\mathfrak h\subseteq\mathfrak l\)，将其视为 \(\mathfrak l\) 值形式；第一项作用于 \(U\) 方向，第二项作用于 \(H\) 方向。若换成 \(s'=s\,k\)，其中 \(k:U\to H\) 光滑，则
\[
	B'=\operatorname{Ad}_{k^{-1}}B+k^*\theta_H.
\]
这个变换律与普通主联络相似，但 \(B\) 还含有不能丢掉的 \(\mathfrak l/\mathfrak h\) 分量。

\subsection{曲率、平坦性及路径展开}
\label{g1:subsec:D02:02}

模型上的 \(\theta_L\) 不只是逐点可逆，它还满足一个微分相容条件。对左不变场 \(X^L,Y^L\)，有 \(\theta_L(X^L)=X\)、\(\theta_L(Y^L)=Y\)，二者为常值，因此
\[
	\mathrm{d}\theta_L(X^L,Y^L)
	=-\theta_L([X^L,Y^L])=-[X,Y].
\]
左不变场逐点张成切空间，故
\begin{equation}\label{g1:eq:cartan-model-mc}
	\mathrm{d}\theta_L+\tfrac12[\theta_L\wedge\theta_L]=0.
\end{equation}
这里沿用\eqref{g1:eq:lie-bracket-wedge}的约定，即
\[
	[\omega\wedge\omega](U,V)=2[\omega(U),\omega(V)].
\]
因子 \(1/2\) 恰好抵消交替求和产生的重复。

\begin{definition}\label{g1:def:cartan-curvature}
	设 \((P\to M,\omega)\) 为模型 \(L/H\) 的光滑 Cartan 几何。定义它的\term[cartanqulv]{Cartan 曲率}[Cartan curvature]为
	\begin{equation}\label{g1:eq:cartan-curvature}
		\Omega=\mathrm{d}\omega+\tfrac12[\omega\wedge\omega]
		\in\Omega^2(P;\mathfrak l).
	\end{equation}
	若 \(\Omega=0\)，称这个 Cartan 几何平坦。
\end{definition}

\begin{proposition}[曲率的水平性、等变性与 Bianchi 恒等式]\label{g1:prop:cartan-curvature-identities}
	设 \((P\to M,\omega)\) 为模型 \(L/H\) 的光滑 Cartan 几何，其中 \(L\) 为有限维实Lie 群、\(H\subseteq L\) 闭。其曲率满足
	\begin{equation}\label{g1:eq:cartan-curvature-properties}
		\iota_{X^\#}\Omega=0\quad(X\in\mathfrak h),\qquad
		R_h^*\Omega=\operatorname{Ad}_{h^{-1}}\Omega\quad(h\in H),
	\end{equation}
	以及
	\begin{equation}\label{g1:eq:cartan-bianchi}
		\mathrm{d}\Omega+[\omega\wedge\Omega]=0.
	\end{equation}
\end{proposition}
\begin{proof}
	将 \(R_{\exp(tX)}^*\omega=\operatorname{Ad}_{\exp(-tX)}\omega\) 在零点求导，得
	\[
		\mathcal L_{X^\#}\omega=-[X,\omega].
	\]
	由于 \(\iota_{X^\#}\omega=X\) 是常函数，Cartan 的 Lie 导数公式给出
	\[
		\iota_{X^\#}\mathrm{d}\omega=-[X,\omega].
	\]
	另一方面，插入括号楔积给出
	\(\frac12\iota_{X^\#}[\omega\wedge\omega]=[X,\omega]\)，两项抵消，证明水平性。固定 \(h\in H\) 时，\(\operatorname{Ad}_{h^{-1}}\) 是常系数的Lie 代数自同构；它与外微分相交换并保持括号。因此直接拉回\eqref{g1:eq:cartan-curvature}即得等变性。

	Bianchi 恒等式只使用Lie 代数括号的代数性质，不要求 \(\omega\) 是 \(\mathfrak h\) 值。具体地，\cref{g1:thm:principal-bianchi}中已由分量计算得到
	\[
		\mathrm{d}[\omega\wedge\omega]
		=-2[\omega\wedge\mathrm{d}\omega],
		\qquad
		[\omega\wedge[\omega\wedge\omega]]=0.
	\]
	第二式是对三个输入交替求和后的 Jacobi 恒等式。于是
	\[
		\mathrm{d}\Omega=-[\omega\wedge\mathrm{d}\omega],
		\qquad
		[\omega\wedge\Omega]=[\omega\wedge\mathrm{d}\omega],
	\]
	相加便得结论。
\end{proof}

曲率的“水平性”表示有一个垂直输入时曲率为零，因此它通过底空间的切向量取值。结合焊接同构，可以把 \(\Omega_p\) 写成
\[
	\kappa(p)\in
	\operatorname{Hom}\bigl(\Lambda^2(\mathfrak l/\mathfrak h),\mathfrak l\bigr),
	\qquad
	\kappa(p)(\overline X,\overline Y)
	=\Omega_p(\omega_p^{-1}X,\omega_p^{-1}Y).
\]
改变 \(X,Y\) 的代表只会给输入增加垂直向量，所以这个式子有定义。它说明 Cartan 曲率记录的是两个底空间方向之间的相容障碍。

\begin{example}[定向曲面的Euclid  Cartan 几何]\label{g1:exm:cartan-euclidean-surface}
	设 \(M\) 为定向的二维光滑Riemann 流形，\(P\to M\) 为其正向正交标架丛。把标架 \(p=(e_1,e_2)\) 看作等距同构 \(p:\mathbb R^2\to\mathrm{T}_{\pi(p)}M\)，定义焊接形式
	\[
		\theta_p(V)=p^{-1}(\mathrm{d}\pi)_pV
		=\begin{pmatrix}\theta^1_p(V)\\\theta^2_p(V)\end{pmatrix}.
	\]
	它在垂直向量上为零，且 \(R_h^*\theta=h^{-1}\theta\)。令 \(A=J_0\alpha\) 为该度量的无挠度量联络在 \(P\) 上的主联络形式，其中
	\[
		J_0=\begin{pmatrix}0&-1\\1&0\end{pmatrix},
		\qquad
		\alpha\in\Omega^1(P).
	\]
	局部截面拉回的 \(\alpha\) 正是\cref{g1:def:surface-connection-form}中满足 \(\nabla e_1=\alpha e_2\)、\(\nabla e_2=-\alpha e_1\) 的形式；其换架公式把这些局部数据粘成上述主联络。

	取 \(L=\operatorname{SE}(2)=\mathbb R^2\rtimes\operatorname{SO}(2)\)，\(H=\operatorname{SO}(2)\)，于是 \(L/H\) 是定向Euclid 平面。\(\mathfrak l=\mathfrak{so}(2)\ltimes\mathbb R^2\) 的括号为
	\[
		[(B,u),(C,v)]=([B,C],Bv-Cu).
	\]
	在 \(P\) 上令
	\[
		\omega=(J_0\alpha,\theta).
	\]
	若 \(\omega_p(V)=0\)，则 \(\theta_p(V)=0\) 使 \(V\) 垂直，再现条件又使 \(A_p(V)=0\) 只能发生于 \(V=0\)。两边维数均为 \(3\)，故 \(\omega_p\) 是同构。主联络与焊接形式的等变律，恰好分别是 \(\operatorname{Ad}_{h^{-1}}\) 在旋转分量与平移分量上的作用；垂直再现也成立。因此 \(\omega\) 是 Cartan 联络。

	按半直积括号展开曲率，得到
	\begin{equation}\label{g1:eq:cartan-surface-curvature}
		\Omega=
		\left(
		J_0\,\mathrm{d}\alpha,\ 
		\mathrm{d}\theta+J_0\alpha\wedge\theta
		\right).
	\end{equation}
	第二分量是切丛联络的挠率形式。此处联络无挠，因而它为零；分量式是
	\[
		\mathrm{d}\theta^1=\alpha\wedge\theta^2,\qquad
		\mathrm{d}\theta^2=-\alpha\wedge\theta^1.
	\]
	由\cref{g1:prop:connection-curvature-form}，第一分量满足
	\[
		\mathrm{d}\alpha=-(K\circ\pi)\theta^1\wedge\theta^2.
	\]
	该式在局部截面上就是 \(\mathrm{d}\alpha=-K\,\mathrm{d}A\)；两边都水平且等变，所以它在整个标架丛上成立。因此这个以Euclid 平面为模型的 Cartan 几何平坦，当且仅当 \(K=0\)。单位球面虽然无挠，仍有非零的Euclid  Cartan 曲率。
\end{example}

为什么要把旋转曲率与挠率放进同一个 \(\Omega\)？局部Euclid 坐标不仅要求标架在各方向上的转动相容，也要求记录的位移真正来自坐标。前者对应旋转分量，后者对应平移分量。对一般齐性模型，下面的定理把这两种相容要求统一为 Maurer--Cartan 方程。

\begin{theorem}[平坦 Cartan 几何的局部模型]\label{g1:thm:flat-cartan-local-model}
	设 \(L\) 为有限维实Lie 群，\(H\subseteq L\) 为闭子群，\((P\to M,\omega)\) 为无边界光滑流形上的 Cartan 几何。记 \(q:L\to L/H\) 为商投影，\((\theta_L)_a=(\mathrm dL_{a^{-1}})_a\) 为左 Maurer--Cartan 形式。其曲率为零，当且仅当每个 \(x_0\in M\) 都有开邻域 \(U\)，以及 \(L/H\) 中的开集 \(V\)，使 \((P|_U,\omega)\) 与 \((q^{-1}(V),\theta_L)\) 作为主 \(H\) 丛及 Cartan 联络同构。
\end{theorem}
\begin{proof}
	先设 \(\Omega=0\)。我们先在总空间的一点附近积分 \(\omega\)，再借助主丛局部截面将所得映射扩展到整条 \(H\) 纤维。第二步是从总空间的局部坐标得到底空间局部等价所必需的。

	对 \(X\in\mathfrak l\)，令 \(E_X\) 为由 \(\omega(E_X)=X\) 定义的光滑向量场。由曲率为零，且等式两边均为常值，可得
	\[
		0=\Omega(E_X,E_Y)
		=-\omega([E_X,E_Y])+[X,Y],
		\qquad [E_X,E_Y]=E_{[X,Y]}.
	\]
	在 \(P\times L\) 上考虑分布
	\[
		\mathcal D_{(p,a)}
		=\{(v,w):\omega_p(v)=(\theta_L)_a(w)\}.
	\]
	两形式都逐点可逆，故此分布秩恒为 \(\dim L\)。若 \(X_i\) 是 \(\mathfrak l\) 的一组基，则 \((E_{X_i},X_i^L)\) 给出它的光滑标架，且
	\[
		[(E_X,X^L),(E_Y,Y^L)]
		=(E_{[X,Y]},[X,Y]^L).
	\]
	标架的括号仍在分布中；再用带函数系数的括号公式，得 \(\mathcal D\) 对合。

	固定 \(p_0\in P_{x_0}\)。\cref{g1:thm:frobenius}适用于无边界流形 \(P\times L\) 上的这个恒秩光滑分布，给出通过 \((p_0,e)\) 的局部嵌入积分片。第一投影在 \(\mathcal D\) 上的微分为同构，因为每个 \(v\) 唯一决定 \(w\)。由\cref{g1:thm:inverse-function}，缩小积分片后，它是某个光滑映射 \(f:W\to L\) 的图像，其中 \(W\) 为 \(p_0\) 的开邻域、\(f(p_0)=e\)。图像的切空间在 \(\mathcal D\) 中，正好说明
	\[
		f^*\theta_L=\omega|_W.
	\]
	两个一形式逐点可逆，故 \(\mathrm{d}f\) 逐点可逆。

	取经过 \(p_0\) 的局部截面 \(s:U\to P\)，缩小 \(U\) 使 \(s(U)\subset W\)，并记 \(a=f\circ s\)。定义
	\[
		\Phi:P|_U\longrightarrow L,\qquad
		\Phi(s(x)h)=a(x)h,\quad h\in H.
	\]
	主丛平凡化给出 \(s(x)h\) 的唯一表示，故 \(\Phi\) 光滑、有定义，并对所有 \(h\in H\) 等变；这里不要求 \(H\) 连通。沿截面有 \(a^*\theta_L=s^*\omega\)。由 Maurer--Cartan 乘法公式及\eqref{g1:eq:cartan-local-form}，在 \(U\times H\) 上
	\[
		\Phi^*\theta_L
		=\operatorname{Ad}_{h^{-1}}a^*\theta_L+\theta_H
		=\operatorname{Ad}_{h^{-1}}s^*\omega+\theta_H
		=\omega.
	\]
	此式是在主丛平凡化中书写的，因而给出整个 \(P|_U\) 上的等式。

	定义底映射 \(\phi(x)=a(x)H\)。它的微分可由 \(s^*\omega\bmod\mathfrak h\) 计算：左平移 \(a(x)^{-1}\) 将 \(\mathrm{T}_{\phi(x)}(L/H)\) 识别为 \(\mathfrak l/\mathfrak h\)，在此识别下，\(\mathrm{d}\phi_x\) 就是 \(\overline\omega_{s(x)}\)。\cref{g1:prop:cartan-soldering}说明它为同构。再缩小 \(U\)，逆函数定理使 \(\phi:U\to V\) 成为到某个开集 \(V\subseteq L/H\) 的微分同胚。

	对每个 \(x\in U\)，\(\Phi\) 将整条纤维 \(s(x)H\) 双射到 \(a(x)H=q^{-1}(\phi(x))\)。因此它是 \(P|_U\to q^{-1}(V)\) 的双射。目标截面 \(y\mapsto a(\phi^{-1}(y))\) 给出局部平凡化，在两边平凡化中 \(\Phi\) 写成 \((x,h)\mapsto(\phi(x),h)\)，故其逆光滑。这就得到保持 Cartan 联络的主丛同构。

	反之，若每一点附近存在这样的同构，拉回模型上的\eqref{g1:eq:cartan-model-mc}，便得 \(\Omega|_{P|_U}=0\)。这些 \(U\) 覆盖 \(M\)，故 \(\Omega=0\)。
\end{proof}

证明中的积分片把两个空间里由同一个 \(X\in\mathfrak l\) 指定的运动配对；曲率为零恰好保证这些配对在取括号后仍相容。局部截面随后把积分片产生的映射延伸到整条 \(H\) 纤维，才得到底空间开片上的几何等价。

\ProseParagraphBreak

曲率不为零时，仍然可以沿一条指定路径，把实际空间中的运动记录到模型中。给定有限分片光滑路径 \(\gamma:[0,1]\to M\) 和初始标架 \(p_0\in P_{\gamma(0)}\)，先选取同样正则的提升 \(p(t)\)，满足 \(p(0)=p_0\)。这种提升可在有限个局部平凡化中逐段构造：每段保持选定的群坐标不变，在接点调节后一段的群坐标使两段接合。

\begin{proposition}[路径的展开]\label{g1:prop:cartan-path-development}
	设 \((P\to M,\omega)\) 为模型 \(L/H\) 的光滑 Cartan 几何，\(L\) 为有限维实Lie 群，\(H\subseteq L\) 闭。给定有限分片光滑路径 \(\gamma:[0,1]\to M\) 及 \(p_0\in P_{\gamma(0)}\)，选择有限分片光滑提升 \(p:[0,1]\to P\)，使 \(\pi\circ p=\gamma\)、\(p(0)=p_0\)。记 \((\theta_L)_a=(\mathrm dL_{a^{-1}})_a\) 为 \(L\) 上的左 Maurer--Cartan 形式。则未知路径 \(g:[0,1]\to L\) 的初值问题
	\begin{equation}\label{g1:eq:cartan-development-ode}
		\theta_L(g'(t))=\omega(p'(t)),\qquad g(0)=e
	\end{equation}
	在整个区间上有唯一有限分片光滑解。路径
	\[
		\operatorname{dev}_{p_0}(\gamma)(t)=g(t)H
	\]
	与保持 \(p(0)=p_0\) 的提升选择无关，称为 \(\gamma\) 相对于 \(p_0\) 的\term[lujingzhankai]{路径展开}[development]。对任意 \(h_0\in H\)，若把初始标架改成 \(p_0h_0\)，则展开路径变为 \(h_0^{-1}g(t)H\)。
\end{proposition}
\begin{proof}
	记 \(b(t)=\omega(p'(t))\)。方程等价于 \(g'=(\mathrm{d}L_g)_e b(t)\)。令 \(u=g^{-1}\)，对 \(g(t)u(t)=e\) 求导可得
	\[
		u'=-(\mathrm{d}R_u)_e b(t).
	\]
	\cref{g1:lem:right-invariant-time-ode}说明，光滑Lie 代数值系数在紧时间区间上定义的这类右平移方程有全区间唯一解；其证明用单位元附近的统一短时存在性及右平移延拓，不要求 \(L\) 紧。对每个光滑子区间应用该结论并接合初值，即得此处的解。

	另一个提升可唯一写成 \(\widehat p(t)=p(t)h(t)\)，其中 \(h(t)\in H\) 有限分片光滑。对变化的 \(h(t)\)，等变性和再现条件给出
	\[
		\omega(\widehat p')
		=\operatorname{Ad}_{h^{-1}}\omega(p')+\theta_H(h').
	\]
	另一方面，\cref{g1:prop:maurer-cartan-product}给出
	\[
		\theta_L((gh)')
		=\operatorname{Ad}_{h^{-1}}\theta_L(g')+\theta_H(h').
	\]
	若两提升初值均为 \(p_0\)，则 \(h(0)=e\)，所以 \(gh\) 满足新方程及单位元初值。唯一性使新解正是 \(gh\)，投影后 \((gh)H=gH\)，从而展开路径不变。

	若初始标架改为 \(p_0h_0\)，可选择提升 \(ph_0\)。新系数为 \(\operatorname{Ad}_{h_0^{-1}}b\)，新解为 \(h_0^{-1}gh_0\)，因为它在零时刻为单位元且满足相应方程。其陪集为 \(h_0^{-1}gH\)，得到最后的变换式。
\end{proof}

在\cref{g1:thm:flat-cartan-local-model}的平坦邻域中，若局部同构归一化为 \(\Phi(p_0)=e\)，则 \(\Phi(p(t))\) 满足展开方程，故展开就是底空间局部映射 \(\phi(\gamma(t))\)。一般情形下，展开可以依赖整条路径；不能仅由终点 \(x\) 写出一个单值的 \(\operatorname{dev}(x)\)。即使曲率为零，把不同邻域的映射接成全局单值映射也还涉及绕环路的延拓，\secref{g1:sec:D03}将对此作具体计算。

\ProseParagraphBreak

Cartan 几何的完备性提出反向的延拓问题：给定模型中的运动，实际空间中是否始终存在相应运动？这与沿实际空间中已给路径求解展开方程使用的数据不同。例如开区间 \((0,1)\) 上的 \(\omega=\mathrm{d}x\) 以加法群 \(\mathbb R\) 为模型，曲率为零；但从 \(x=1/2\) 按单位速度前进，只能在到达区间边界之前留在原空间中。局部平坦、全局单值和运动能否无限延拓，是三个不同问题。

\subsection{Cartan 联络与主联络}
\label{g1:subsec:D02:03}

第18章通过一个与右作用相容的水平补空间定义主联络，这也是主丛上的 Ehresmann 联络。\footnote{Charles Ehresmann（1905--1979），法国数学家，研究纤维丛、联络与微分拓扑。}若用一形式表示，它是 \(\mathfrak h\) 值的 \(\eta\)，满足等变性和垂直再现，并由 \(\ker\eta\) 给出水平分布。Cartan 联络则是更大Lie 代数 \(\mathfrak l\) 值的逐点同构，因此
\[
	\ker\omega_p=0,\qquad
	\omega_p^{-1}(\mathfrak h)=V_pP.
\]
把 \(\ker\omega\) 当成水平分布，会使所有非零底空间速度都无从提升。

\ProseParagraphBreak

两类形式在垂直方向上的行为相同，但在底空间方向上承担的任务不同。\cref{g1:tab:cartan-ehresmann}把这一区别列在同一个主 \(H\) 丛上；其中 \(\mathcal H\) 表示水平分布，避免与结构群 \(H\) 混用。
\begin{table}[htbp]
	\centering
	\caption{主联络与 Cartan 联络所记录的数据}
	\label{g1:tab:cartan-ehresmann}
	\begin{tabularx}{\linewidth}{@{}lXX@{}}
		\toprule
		比较项 & 主丛 Ehresmann 联络 \(\eta\) & 模型 \(L/H\) 的 Cartan 联络 \(\omega\)\\
		\midrule
		取值空间 & \(\mathfrak h\) & \(\mathfrak l\)\\
		逐点作用 & 记录垂直分量，\(\eta|_{VP}\) 可逆 & 记录全部速度，\(\omega_p\) 可逆\\
		核 & \(\mathcal H_p\)，维数为 \(\dim M\) & 零空间\\
		底空间方向 & 投影 \(\mathcal H_p\to\mathrm{T}_xM\) 为同构 & 选定 \(p\in P_x\) 后给出 \(\mathrm{T}_xM\cong\mathfrak l/\mathfrak h\)\\
		维数限制 & 与某个外部模型群无关 & \(\dim M=\dim(L/H)\)\\
		平坦性的局部含义 & 每个总空间点均有局部水平截面通过 & 与 \(L/H\) 的开片局部等价\\
		\bottomrule
	\end{tabularx}
\end{table}

若模型Lie 代数存在与 \(H\) 作用相容的补空间，两类联络仍有直接联系。这里必须要求补空间在整个 \(\operatorname{Ad}H\) 下不变；任取一个线性补空间虽能分解向量，却未必能分解出满足主丛等变律的形式。

\begin{proposition}[不变补空间下的分解]\label{g1:prop:reductive-cartan-splitting}
	设 \(L\) 为有限维实Lie 群，\(H\subseteq L\) 为闭子群，且有 \(\operatorname{Ad}H\) 不变的直和
	\[
		\mathfrak l=\mathfrak h\oplus\mathfrak m.
	\]
	对任意以 \(L/H\) 为模型的光滑 Cartan 几何 \((P\to M,\omega)\)，记
	\[
		\eta=\operatorname{pr}_{\mathfrak h}\omega,\qquad
		\vartheta=\operatorname{pr}_{\mathfrak m}\omega.
	\]
	则 \(\eta\) 是主 \(H\) 联络，\(\vartheta\) 为水平、等变的 \(\mathfrak m\) 值形式，且诱导 \(\mathrm{T}M\cong P\times_H\mathfrak m\)。此时主联络的水平空间为 \(\omega^{-1}(\mathfrak m)\)。
\end{proposition}
\begin{proof}
	两直和分量都在 \(\operatorname{Ad}H\) 下不变，故两投影与 \(\operatorname{Ad}H\) 交换。\(\omega\) 的等变性因此分别下降为 \(\eta,\vartheta\) 的等变性。对 \(X\in\mathfrak h\)，再现条件给出 \(\eta(X^\#)=X\)、\(\vartheta(X^\#)=0\)。由\cref{g1:thm:principal-connection-equivalence}，前两项性质使 \(\eta\) 成为主联络，而后一等式给出 \(\vartheta\) 的水平性。投影 \(\mathfrak m\to\mathfrak l/\mathfrak h\) 是 \(H\) 等变同构，故焊接同构直接成为 \(\mathrm{T}M\cong P\times_H\mathfrak m\)。最后，\(\eta_p(v)=0\) 等价于 \(\omega_p(v)\in\mathfrak m\)，所以得到所述水平空间。
\end{proof}

以上讨论的是具有不变补空间的模型（reductive model）。即使不变补空间存在，也可能不唯一；由它投影得到的普通主联络随所选分解而定。如果不存在不变补空间，就不能用一个任意线性投影，从 \(\omega\) 自然得到普通主联络。主 \(H\) 丛本身仍可另选主联络，这是另一项选择。

\ProseParagraphBreak

在这个分解下，Cartan 曲率的两分量写成
\[
	\begin{aligned}
		\Omega_{\mathfrak h}
		&=\mathrm{d}\eta+\tfrac12[\eta\wedge\eta]
		+\tfrac12[\vartheta\wedge\vartheta]_{\mathfrak h},\\
		\Omega_{\mathfrak m}
		&=\mathrm{d}\vartheta+[\eta\wedge\vartheta]
		+\tfrac12[\vartheta\wedge\vartheta]_{\mathfrak m}.
	\end{aligned}
\]
这里 \([\mathfrak h,\mathfrak m]\subseteq\mathfrak m\) 来自不变性，因此混合项只出现在第二行。Euclid 模型的平移部分满足 \([\mathfrak m,\mathfrak m]=0\)，于是两行分别是普通联络的曲率与挠率；\cref{g1:exm:cartan-euclidean-surface}正是这种情形。一般模型中还要保留 \([\vartheta\wedge\vartheta]\) 的贡献，所以“Cartan 平坦”必须连同所选模型一起理解。

\section{标架归一化与几何等价}
\label{g1:sec:D03}

同一张平面可以用直角标架或旋转标架描述，联络系数随之改变；周长不同的平圆柱却可以具有相同的零曲率，而不具有相同的整体几何。等价问题必须同时回答两件事：哪些差异只是表示方式的选择，哪些数据在改变表示后仍然存在。本节先在结构方程中辨认这些自由，再用曲率不恒定的度量曲面给出一次具体的标架归一化，最后用显式展开映射检验局部结论能否成为整体结论。

\subsection{结构方程中的规范自由}
\label{g1:subsec:D03:01}

设 \(G\subseteq\operatorname{GL}(n,\mathbb R)\) 是矩阵Lie 群，在一个带 \(G\)-结构的 \(n\) 维流形的开集 \(U\) 上选适应标架 \(e=(e_1,\ldots,e_n)\)。对偶余标架写成列 \(\theta=(\theta^1,\ldots,\theta^n)^{\mathsf T}\)，即 \(V=e\,\theta(V)\)。选定一个相容联络后，沿用 \(\nabla e=eA\)，其中 \(A\) 是取值于 \(\mathfrak g\) 的局部联络矩阵。

\begin{proposition}[换标架与结构方程]\label{g1:prop:equivalence-frame-change}
	设 \(G\subseteq\operatorname{GL}(n,\mathbb R)\) 为矩阵Lie 群，\(e\) 是某个 \(n\) 维光滑流形的开集 \(U\) 上一个 \(G\)-结构的光滑适应标架，\(\nabla\) 为相容联络，\(\theta,A\) 分别为此标架中的余标架与联络矩阵。将标架改为 \(e'=eg\)，其中 \(g:U\to G\) 光滑，联络本身保持不变。则
	\begin{equation}\label{g1:eq:equivalence-frame-change}
		\theta'=g^{-1}\theta,\qquad
		A'=g^{-1}Ag+g^{-1}\mathrm{d}g.
	\end{equation}
	记挠率形式及曲率矩阵为 \(T=\mathrm{d}\theta+A\wedge\theta\)、\(F_A=\mathrm{d}A+A\wedge A\)，则
	\[
		T'=g^{-1}T,\qquad F_{A'}=g^{-1}F_Ag.
	\]
	这里矩阵值形式的楔积同时使用矩阵乘法，向量值形式的楔积按矩阵作用在列向量上解释。
\end{proposition}
\begin{proof}
	恒等式 \(e\theta(V)=eg\theta'(V)\) 对每个切向量成立，所以 \(\theta'=g^{-1}\theta\)。再对 \(e'=eg\) 使用联络的乘法法则，得到
	\[
		\nabla e'=e(Ag+\mathrm{d}g)=e'\bigl(g^{-1}Ag+g^{-1}\mathrm{d}g\bigr),
	\]
	即第二式。这也恢复了\cref{g1:thm:principal-connection-local-data}在标架丛上的局部公式：固定的主联络在不同截面上的表达不同。

	由 \(\mathrm{d}(g^{-1})=-g^{-1}(\mathrm{d}g)g^{-1}\)，
	\[
		\begin{aligned}
			\mathrm{d}\theta'&=-g^{-1}\mathrm{d}g\wedge g^{-1}\theta+g^{-1}\mathrm{d}\theta,\\
			A'\wedge\theta'&=g^{-1}A\wedge\theta+g^{-1}\mathrm{d}g\wedge g^{-1}\theta.
		\end{aligned}
	\]
	相加时含 \(\mathrm{d}g\) 的两项抵消，得到挠率公式。对联络的算子 \(\mathrm{D}=\mathrm{d}+A\)，新分量中的算子为 \(\mathrm{D}'=g^{-1}\mathrm{D}g\)；连续作用两次得到 \((\mathrm{D}')^2=g^{-1}\mathrm{D}^2g\)。由于 \(\mathrm{D}^2\) 在列向量值函数上等于左乘 \(F_A\)，曲率公式随之成立。
\end{proof}

这里的换标架没有移动底空间的点，也没有改变联络。它与\cref{g1:prop:gauge-connection-formula}中的主动规范变换具有同样的局部公式，但主动规范变换将联络拉回到另一个联络，所选截面可以保持不动。另一方面，两张 \(G\)-结构之间的几何等价还要有底空间的局部微分同胚，并要求其微分把适应标架送到适应标架。这是对底点和切向量的共同要求，不能由任取一族矩阵 \(g(x)\) 代替。

\ProseParagraphBreak

例如，设 \(\Phi:U\to\widehat U\) 是待检验的局部微分同胚，两边的标架分别为 \(e,\widehat e\)。若写
\[
	\Phi_*e=\widehat e\,k,\qquad k:U\to G,
\]
则余标架必须满足 \(\Phi^*\widehat\theta=k\theta\)。若还要求保持所选联络，则必须同时满足
\[
	A=k^{-1}(\Phi^*\widehat A)k+k^{-1}\mathrm{d}k.
\]
第二式由对 \(\Phi_*e=\widehat e k\) 求协变导数得到，第一式则直接来自切向量的分量。两个条件中的同一个 \(k\) 必须与 \(\mathrm{d}\Phi\) 相容。只把联络矩阵变成某种相似形式，尚未得到这样的 \(\Phi\)。

\ProseParagraphBreak

另一个不同的选择是固定标架，改变相容联络。若 \(A\) 改为 \(A+B\)，其中 \(B\in\Omega^1(U;\mathfrak g)\)，则
\[
	T_{A+B}=T_A+B\wedge\theta.
\]
用 \(\theta\) 将切空间识别为 \(V=\mathbb R^n\)，在一点把 \(B\) 看作 \(V^*\otimes\mathfrak g\) 的元素，挠率的变化就是
\[
	(\delta B)(v,w)=B(v)w-B(w)v.
\]
这正是\subsecref{g1:subsec:D01:02}中讨论的线性映射。给定挠率中落在 \(\operatorname{im}\delta\) 的一项，才可能选 \(B\) 将它消去；商空间中的内蕴挠率不受这一操作影响。如果有多个 \(B\) 能消去同一项，它们之差落在 \(\ker\delta\)。因此，选联络消去挠率与选标架减少结构群是两种操作，各自留下的自由也不同。

\begin{example}[二维正交结构中的消挠率]\label{g1:exm:surface-torsion-absorption}
	在定向Riemann 度量曲面上取正向正交余标架 \(\theta^1,\theta^2\)，记 \(\Theta=\theta^1\wedge\theta^2\)。由于二维二形式只有一个独立分量，存在唯一的一对光滑函数 \(c_1,c_2\) 使
	\[
		\mathrm{d}\theta^1=c_1\Theta,\qquad
		\mathrm{d}\theta^2=c_2\Theta.
	\]
	所有度量相容的局部联络矩阵都形如
	\[
		A=J_0\alpha,\qquad
		J_0=\begin{pmatrix}0&-1\\1&0\end{pmatrix},\qquad
		\alpha=a\theta^1+b\theta^2.
	\]
	把它代入挠率公式，逐项得到
	\[
		T^1=\mathrm{d}\theta^1-\alpha\wedge\theta^2=(c_1-a)\Theta,
		\qquad
		T^2=\mathrm{d}\theta^2+\alpha\wedge\theta^1=(c_2-b)\Theta.
	\]
	因此无挠条件唯一确定 \(a=c_1,b=c_2\)。这一维数下的 \(\delta\) 既能消去所有挠率，又没有剩余的联络差；所得正是\cref{g1:prop:surface-connection-properties}的唯一度量相容无挠联络。

	但是正交标架仍可旋转。若
	\[
		g=\begin{pmatrix}\cos\varphi&-\sin\varphi\\
		\sin\varphi&\cos\varphi\end{pmatrix},
	\]
	则 \(g^{-1}J_0g=J_0\)、\(g^{-1}\mathrm{d}g=J_0\mathrm{d}\varphi\)，故 \(\alpha'=\alpha+\mathrm{d}\varphi\)。这与\cref{g1:prop:connection-frame-change}的符号一致。消去挠率并未消去旋转自由，也不能任意使 \(\alpha'=0\)：若在一个开集上能如此选择，便有 \(\mathrm{d}\alpha=0\)，由\cref{g1:prop:connection-curvature-form}必有 \(K=0\)。
\end{example}

更一般的结构可能留下非零内蕴挠率、非唯一相容联络或无法一次固定的结构群。下面只对刚才的曲面模型继续归一化：我们将使用一个由度量决定且非零的切向量来选定旋转角。这个具体构造既说明归一化怎样工作，也明确显示它在哪些点会失效。

\subsection{曲率梯度与曲面标架的归一化}
\label{g1:subsec:D03:02}

设 \((U,h)\) 为定向光滑Riemann 度量曲面的一块开区域，\(K\) 为 Gauss 曲率，并假设 \(\mathrm{d}K\) 在 \(U\) 上处处非零。这里讨论的是带度量的实二维曲面。梯度由
\[
	h(\operatorname{grad}_h K,V)=\mathrm{d}K(V)
\]
对全部切向量 \(V\) 定义；度量正定使这一向量唯一，且 \(\mathrm{d}K\ne0\) 保证它不为零。令
\begin{equation}\label{g1:eq:curvature-gradient-frame}
	q=|\operatorname{grad}_h K|_h>0,\qquad
	e_1=\frac{\operatorname{grad}_h K}{q},\qquad e_2=Je_1,
\end{equation}
其中 \(J\) 为由 \(h\) 与定向确定的正向直角旋转。于是 \(e_1\) 指向曲率增加最快的方向，\(e_2\) 的方向由定向确定；\(\operatorname{SO}(2)\) 中保持这组标架的元素只有单位元。

\ProseParagraphBreak

记对偶余标架为 \(\theta^i=h(e_i,-)\)。由梯度定义，
\[
	\mathrm{d}K=q\theta^1,\qquad \theta^1=\frac{\mathrm{d}K}{q},
	\qquad h=(\theta^1)^2+(\theta^2)^2.
\]
相应联络一形式仍取第九章的约定
\[
	\alpha(V)=h(\nabla_Ve_1,e_2),\qquad
	\nabla e_1=\alpha e_2,\qquad \nabla e_2=-\alpha e_1.
\]
因为标架已经由度量选定，\(\alpha\) 在这一步也不再含任意旋转角。

\begin{proposition}[归一化曲面标架的结构式]\label{g1:prop:normalized-surface-structure}
	设 \((U,h)\) 为定向光滑 Riemann 度量曲面，其 Gauss 曲率 \(K\) 满足 \(\mathrm dK\) 处处非零。令 \(J\) 为各切平面内的正向直角旋转，\(q=|\operatorname{grad}_h K|_h\)、\(e_1=\operatorname{grad}_h K/q\)、\(e_2=Je_1\)、\(\theta^i=h(e_i,-)\)。取 \(h\) 的 Levi--Civita 联络 \(\nabla\)，定义 \(\alpha(V)=h(\nabla_Ve_1,e_2)\)，并写 \(\alpha=a\theta^1+b\theta^2\)。则
	\begin{equation}\label{g1:eq:normalized-surface-structure}
		\begin{aligned}
			\mathrm{d}\theta^1&=\alpha\wedge\theta^2=a\theta^1\wedge\theta^2,\\
			\mathrm{d}\theta^2&=-\alpha\wedge\theta^1=b\theta^1\wedge\theta^2,\\
			\mathrm{d}\alpha&=-K\theta^1\wedge\theta^2.
		\end{aligned}
	\end{equation}
	此外，
	\begin{equation}\label{g1:eq:normalized-surface-coefficients}
		a=\frac{e_2(q)}q,\qquad
		K=e_2(a)-e_1(b)-a^2-b^2,\qquad
		[e_1,e_2]=-ae_1-be_2.
	\end{equation}
\end{proposition}
\begin{proof}
	度量相容和无挠性给出
	\[
		\begin{aligned}
			\mathrm{d}\theta^1(V,W)
			&=Vh(e_1,W)-Wh(e_1,V)-h(e_1,[V,W])\\
			&=h(\nabla_Ve_1,W)-h(\nabla_We_1,V)\\
			&=\alpha(V)\theta^2(W)-\alpha(W)\theta^2(V).
		\end{aligned}
	\]
	第二行消去了 \(h(e_1,\nabla_VW-\nabla_WV-[V,W])=0\)。所以 \(\mathrm{d}\theta^1=\alpha\wedge\theta^2\)；对 \(e_2\) 作同样计算，利用 \(\nabla e_2=-\alpha e_1\)，得到第二式。

	曲率式的符号也可直接核对。对 \(e_1\) 应用 \(R(V,W)=\nabla_V\nabla_W-\nabla_W\nabla_V-\nabla_{[V,W]}\)，沿 \(e_1\) 的乘积项抵消，余下
	\[
		R(V,W)e_1=\mathrm{d}\alpha(V,W)e_2.
	\]
	\cref{g1:prop:surface-curvature-two-dim}在正交单位基上给出 \(R(e_1,e_2)e_1=-Ke_2\)，故 \(\mathrm{d}\alpha(e_1,e_2)=-K\)，也就得到第三式。

	再对 \(\mathrm{d}K=q\theta^1\) 取外微分。因为 \(\mathrm{d}q=e_1(q)\theta^1+e_2(q)\theta^2\)，有
	\[
		0=\mathrm{d}q\wedge\theta^1+q\mathrm{d}\theta^1
		=\bigl(-e_2(q)+qa\bigr)\theta^1\wedge\theta^2,
	\]
	从而 \(a=e_2(q)/q\)。将 \(\alpha=a\theta^1+b\theta^2\) 求导，得到
	\[
		\mathrm{d}\alpha
		=\bigl(-e_2(a)+a^2+e_1(b)+b^2\bigr)\theta^1\wedge\theta^2.
	\]
	与曲率式比较便得 \(K\) 的表达。最后，\(\theta^i(e_j)\) 为常数，故 \(\mathrm{d}\theta^i(e_1,e_2)=-\theta^i([e_1,e_2])\)。前两条结构式分别给出括号的两个分量 \(-a,-b\)。
\end{proof}

这些式子中的 \(K,q,a,b\) 都由带定向的度量确定。继续沿已由度量选定的 \(e_1,e_2\) 求导，所得量也由同一几何数据决定；任意坐标偏导数则还依赖坐标选择。特别地，比较两个曲面的局部几何时，归一化余标架给出了一个明确的充分必要条件。

\begin{proposition}[归一化余标架的局部等价判据]\label{g1:prop:normalized-coframe-equivalence}
	设 \((U,h)\)、\((\widehat U,\widehat h)\) 为定向光滑Riemann 度量曲面，两个曲率函数的微分均处处非零。对第一曲面，以 \(K\) 记 Gauss 曲率，\(J\) 记切平面内的正向直角旋转，令 \(q=|\operatorname{grad}_h K|_h\)、\(e_1=\operatorname{grad}_h K/q\)、\(e_2=Je_1\)、\(\theta^i=h(e_i,-)\)；以 \(\nabla\) 记其 Levi--Civita 联络，令 \(\alpha(V)=h(\nabla_Ve_1,e_2)\)。在第二曲面上以同样的定义得到带帽记号 \(\widehat K,\widehat q,\widehat e_i,\widehat\theta^i,\widehat\alpha\)。对光滑映射 \(\Phi:U\to\widehat U\)，下列条件等价：
	\begin{equivalences}
		\item \(\Phi\) 是保持定向的局部等距。
		\item \(\Phi^*\widehat\theta^1=\theta^1\) 且 \(\Phi^*\widehat\theta^2=\theta^2\)。
	\end{equivalences}
	满足这些条件时，\(\widehat K\circ\Phi=K\)、\(\widehat q\circ\Phi=q\)，并有 \(\Phi^*\widehat\alpha=\alpha\)。
\end{proposition}
\begin{proof}
	先设 \(\Phi\) 是保持定向的局部等距。沿 \(\Phi\) 拉回目标上的度量相容无挠联络，仍满足度量相容和无挠性；\cref{g1:prop:surface-connection-properties}的唯一性说明它就是源上的联络。因此 \(\Phi\) 也保持曲率，二维曲率表达给出 \(\widehat K\circ\Phi=K\)。

	对每个切向量 \(V\)，使用梯度的定义和等距性，
	\[
		\widehat h\bigl(\Phi_*\operatorname{grad}_hK,\Phi_*V\bigr)
		=h(\operatorname{grad}_hK,V)
		=\mathrm{d}\widehat K(\Phi_*V).
	\]
	\(\Phi_*\) 在每点是同构，故 \(\Phi_*\operatorname{grad}_hK=\operatorname{grad}_{\widehat h}\widehat K\)，并保持其长度。于是 \(\Phi_*e_1=\widehat e_1\)；保持定向与度量还保证 \(\Phi_*J=\widehat J\Phi_*\)，所以 \(\Phi_*e_2=\widehat e_2\)。对偶后得到两条余标架等式。联络被保持且归一化标架对应，便有 \(\Phi^*\widehat\alpha=\alpha\)。

	反过来，两条余标架等式使 \(\mathrm{d}\Phi\) 的核为零：若 \(\Phi_*V=0\)，则 \(\theta^1(V)=\theta^2(V)=0\)，故 \(V=0\)。两侧维数同为二，所以逆函数定理使 \(\Phi\) 成为局部微分同胚。又
	\[
		\Phi^*\widehat h
		=\Phi^*\bigl((\widehat\theta^1)^2+(\widehat\theta^2)^2\bigr)
		=(\theta^1)^2+(\theta^2)^2=h,
	\]
	并且 \(\Phi^*(\widehat\theta^1\wedge\widehat\theta^2)=\theta^1\wedge\theta^2\) 保持正向面积形式。因此 \(\Phi\) 是保持定向的局部等距。
\end{proof}

判据中的等式是整个邻域上的一形式等式，不是只比较某个点的几个数。它也给出待求映射应满足的一阶方程：若 \(\theta^i=C^i_j\mathrm{d}x^j\)、\(\widehat\theta^i=\widehat C^i_a\mathrm{d}y^a\)，则
\[
	\widehat C^i_a(\Phi(x))\,\partial_j\Phi^a(x)=C^i_j(x).
\]
找到满足这些式子的局部映射，才得到等价。有限组标量相等给出必要限制，邻域上的等价还要求解这些一形式方程。

\begin{example}[一个曲率梯度标架]\label{g1:exm:normalized-warped-surface}
	在 \(x>0\) 的平面区域上取定向 \(\mathrm{d}x\wedge\mathrm{d}y>0\) 及度量
	\[
		h_f=\mathrm{d}x^2+f(x)^2\mathrm{d}y^2,\qquad f(x)>0.
	\]
	先用 \(E_1=\partial_x\)、\(E_2=f^{-1}\partial_y\) 及对偶 \(\eta^1=\mathrm{d}x\)、\(\eta^2=f\mathrm{d}y\)。由于
	\[
		\mathrm{d}\eta^1=0,\qquad
		\mathrm{d}\eta^2=\frac{f'}f\eta^1\wedge\eta^2,
	\]
	\cref{g1:exm:surface-torsion-absorption}确定联络一形式为 \(\beta=(f'/f)\eta^2=f'\mathrm{d}y\)。再算
	\[
		\mathrm{d}\beta=f''\mathrm{d}x\wedge\mathrm{d}y
		=\frac{f''}f\eta^1\wedge\eta^2,
	\]
	与曲率式比较得到 \(K=-f''/f\)。

	现在取 \(f=e^{x^2}\)。因 \(f''=(2+4x^2)f\)，有
	\[
		K=-2-4x^2,\qquad\operatorname{grad}_{h_f}K=-8x\partial_x,
		\qquad q=8x.
	\]
	所以归一化标架及余标架为
	\[
		e_1=-\partial_x,\qquad e_2=-f^{-1}\partial_y,
		\qquad \theta^1=-\mathrm{d}x,\quad\theta^2=-f\mathrm{d}y.
	\]
	这里两个向量都改变了符号，是旋转 \(\pi\)，仍保持定向。旋转角为常数，故 \(\alpha=\beta=f'\mathrm{d}y=-2x\theta^2\)，从而 \(a=0,b=-2x\)。直接核验，
	\[
		\mathrm{d}\theta^1=0,\qquad
		\mathrm{d}\theta^2=-2x\theta^1\wedge\theta^2,
		\qquad \mathrm{d}\alpha=(2+4x^2)\theta^1\wedge\theta^2.
	\]
	同时 \(e_2(q)=0\)、\(e_1(b)=2\)，故式\eqref{g1:eq:normalized-surface-coefficients}给出的曲率为 \(-2-4x^2\)，与直接计算一致。\(x=0\) 处度量仍然光滑正定，但 \(\mathrm{d}K=0\)，不能再用曲率梯度来选 \(e_1\)。
\end{example}

\begin{example}[相同的曲率和梯度长度仍不足以确定度量]\label{g1:exm:same-curvature-different-metrics}
	在同一个 \(x>0\) 区域取
	\[
		f_1(x)=e^{x^2},\qquad
		I(x)=\int_0^x e^{-2t^2}\,\mathrm{d}t,\qquad
		f_2(x)=e^{x^2}I(x).
	\]
	两函数均为正。因为 \(I'=e^{-2x^2}\)、\(I''=-4xI'\)，
	\[
		f_2''=f_1''I+2f_1'I'+f_1I''
		=(2+4x^2)f_2.
	\]
	中间两项因 \(f_1'=2xf_1\) 而相消。因此两度量 \(h_i=\mathrm{d}x^2+f_i(x)^2\mathrm{d}y^2\) 具有完全相同的函数 \(K=-2-4x^2\)，而且 \(q=8x\)。归一化标架下却有
	\[
		a_1=a_2=0,\qquad b_1=-2x,\qquad
		b_2=-2x-\frac{I'}I.
	\]
	由于 \(I'/I>0\)，两项 \(b\) 不相同。

	还可以直接排除任何局部等距，连反向等距也包括在内。若局部等距写成 \(\Phi(x,y)=(X(x,y),Y(x,y))\)，曲率必须相等。函数 \(-2-4x^2\) 在 \(x>0\) 上严格单调，所以 \(X=x\)。将此代入 \(\Phi^*h_2=h_1\)，比较 \(\mathrm{d}x^2\) 的系数得 \(1+f_2(x)^2Y_x^2=1\)，故 \(Y_x=0\)。缩小到一个开矩形后，可写 \(Y=\psi(y)\)。再比较 \(\mathrm{d}y^2\) 的系数，得到
	\[
		|\psi'(y)|=\frac{f_1(x)}{f_2(x)}=\frac1{I(x)}.
	\]
	左边只依赖 \(y\)，右边随 \(x\) 严格变化，不可能在一个开矩形上相等。因此没有这样的局部等距。相同的 \(K,q\) 留下了未被测量的横向伸缩，而 \(b\) 或完整余标架能够识别这一区别。
\end{example}

在 \(\mathrm{d}K=0\) 的点，式\eqref{g1:eq:curvature-gradient-frame}的分母为零；若 \(K\) 在开集上恒定，则整个开集都不能用这种归一化。此时可以另找非零的几何数据，或保留适当的标架自由。是否有这样的数据、剩余群是否还能约化，都须针对所选结构继续分析。

\subsection{局部等价与整体单值}
\label{g1:subsec:D03:03}

上面的标架计算比较邻域中的几何。即使每个邻域都等价于同一个齐性模型，沿环路接续这些等价时，也可能回到模型中的另一个位置或另一个标架。第十八章对平坦主联络用环路运输记录单值，\cref{g1:prop:flat-principal-monodromy}的表示取值于结构群。本节的Euclid 几何还记录位置，因此使用保向Euclid 运动群 \(\operatorname{SE}(2)=\mathbb R^2\rtimes\operatorname{SO}(2)\)，其中一个元素作用为 \(z\mapsto Qz+t\)。

\ProseParagraphBreak

下面在平圆柱、平环面与Euclid 开子域中具体构造这种接续。设 \(p:\widetilde M\to M\) 为相应模型的万有覆盖，\(\Gamma\) 为其左覆盖变换群。接续后的局部Euclid 坐标给出局部等距 \(\operatorname{dev}:\widetilde M\to\mathbb R^2\)，以及同态 \(\rho:\Gamma\to\operatorname{SE}(2)\)，满足
\begin{equation}\label{g1:eq:euclidean-developing-equivariance}
	\operatorname{dev}(\gamma\cdot\widetilde x)
	=\rho(\gamma)\operatorname{dev}(\widetilde x).
\end{equation}
这个式子比较绕路后的位置与标架变化：左边先改变覆盖片，右边在模型中施加Euclid 运动。下面在每个模型中写出 \(p,\operatorname{dev},\rho\)，从局部坐标的接续直接得到此式。更换模型中的初始Euclid 坐标 \(\operatorname{dev}'=a\circ\operatorname{dev}\)，其中 \(a\in\operatorname{SE}(2)\)，会将 \(\rho\) 改为 \(a\rho a^{-1}\)，这是将两边同时复合 \(a\) 的直接结果。

\begin{example}[平圆柱的平移单值]\label{g1:exm:flat-cylinder-developing}
	给定 \(\ell>0\)，令
	\[
		C_\ell=\mathbb R^2/\langle T\rangle,
		\qquad T(x,y)=(x+\ell,y).
	\]
	平移保持Euclid 度量，故 \(\mathrm{d}x^2+\mathrm{d}y^2\) 下降为 \(C_\ell\) 上的光滑度量。投影 \(p(x,y)=[x,y]\) 在宽度小于 \(\ell\) 的足够小开矩形上单射，给出均匀覆盖邻域；\(\mathbb R^2\) 单连通，所以它是万有覆盖。覆盖变换为 \(T^n(x,y)=(x+n\ell,y)\)，\(n\in\mathbb Z\)。取
	\[
		\operatorname{dev}(x,y)=(x,y),\qquad
		\rho(n)(u,v)=(u+n\ell,v).
	\]
	\(\rho(n+m)=\rho(n)\rho(m)\)，且 \(\operatorname{dev}\circ T^n=\rho(n)\circ\operatorname{dev}\)，故式\eqref{g1:eq:euclidean-developing-equivariance}成立。

	各小矩形都等距于平面，因而 \(K=0\)。向量场 \(\partial_x,\partial_y\) 也在平移下不变，下降成整个圆柱上的平行正交标架。因此切丛的环路平行移动全部为恒等。可是 \(\rho(1)\) 是长度为 \(\ell\) 的平移，并不平凡：只看切向量的环路平行移动，没有记录这一位置变化。局部一形式 \(\mathrm{d}x\) 同样下降为全局闭一形式 \(\xi\)，沿一周圆柱有 \(\int\xi=\ell\)，故它不是一个全局实值函数的微分。
\end{example}

\cref{g1:fig:euclidean-cylinder-strip}将一个基本带的两条竖边按相同高度对应；一次绕行在覆盖平面上留下长度为 \(\ell\) 的端点差。圆柱的周期长度可以仅用度量和环路恢复，因此不能由局部坐标选择消去。

% !TeX root = ../Book1.tex
\begin{figure}[htbp]
	\centering
	\begin{tikzpicture}[x=1.25cm,y=0.9cm,font=\small]
		\fill[GeometryBlue!7] (0,-1.5) rectangle (3,1.5);
		\draw[dashed,GeometryGray] (0,-1.5)--(3,-1.5);
		\draw[dashed,GeometryGray] (0,1.5)--(3,1.5);
		\draw[thick,GeometryBlue] (0,-1.5)--(0,1.5);
		\draw[thick,GeometryBlue] (3,-1.5)--(3,1.5);
		\draw[->,very thick,GeometryBlue] (0,-0.8)--(0,-0.15);
		\draw[->,very thick,GeometryBlue] (3,-0.8)--(3,-0.15);
		\draw[->,thick] (0,0.45)--(3,0.45);
		\fill (0,0.45) circle (1.3pt);
		\fill (3,0.45) circle (1.3pt);
		\node[above] at (1.5,0.45) {\(\ell\)};
		\node[left] at (0,0.45) {\((0,t)\)};
		\node[right] at (3,0.45) {\((\ell,t)\)};
		\draw[->,GeometryGray] (0,1.85)--(3,1.85);
		\node[above] at (1.5,1.85) {\(T(s,t)=(s+\ell,t)\)};
		\node[below] at (0,-1.5) {\(s=0\)};
		\node[below] at (3,-1.5) {\(s=\ell\)};
	\end{tikzpicture}
	\caption[平圆柱的基本带与平移周期]{平圆柱由基本带的两条竖边按同高点粘合得到，两边箭头同向。横线段投影为一次闭路，在覆盖平面上的端点相差平移 \(\ell\)。上下虚线只表示绘图截断，实际基本带沿 \(t\) 方向无限延伸。}
	\label{g1:fig:euclidean-cylinder-strip}
\end{figure}


\begin{proposition}[平圆柱的最短非平凡环路]\label{g1:prop:cylinder-period-length}
	设 \(\ell>0\)，以平移 \((x,y)\mapsto(x+\ell,y)\) 生成的群作用定义圆柱 \(C_\ell=\mathbb R^2/(\ell\mathbb Z\times\{0\})\)，并赋予标准平度量 \(\mathrm dx^2+\mathrm dy^2\) 下降的商度量。在 \(C_\ell\) 上，所有不可缩的分段 \(C^1\) 环路的长度至少为 \(\ell\)，且存在长度恰为 \(\ell\) 的不可缩环路。对任意 \(\ell'>0\)，以同样方式定义带商度量的 \(C_{\ell'}\)。因此若 \(C_\ell\) 与 \(C_{\ell'}\) 整体等距，则 \(\ell=\ell'\)。
\end{proposition}
\begin{proof}
	将一条环路从 \((x_0,y_0)\) 提升到 \(\mathbb R^2\)，终点必为 \((x_0+n\ell,y_0)\)，其中 \(n\in\mathbb Z\)。由\cref{g1:cor:closed-lift-criterion}，万有覆盖上的提升闭合当且仅当原环路可缩；这里 \(\mathbb R^2\) 的基本群平凡，所以不可缩恰等价于 \(n\ne0\)。投影是局部等距，故原曲线与提升曲线长度相同。提升曲线的Euclid 长度不小于端点距离，因而
	\[
		L(\gamma)\ge |n|\ell\ge\ell.
	\]
	水平线段 \(t\mapsto(x_0+t\ell,y_0)\)、\(0\le t\le1\)，投影后是一条环路，提升的端点差对应 \(n=1\)，所以它不可缩且长度恰为 \(\ell\)。整体等距是微分同胚，保持环路的可缩性与曲线长度，故必须保持这个最小值。
\end{proof}

所以所有 \(C_\ell\) 的局部曲率都相同，切向量的平行移动单值也都平凡，仍有连续变化的整体周期数据。环面上的现象相似：若 \(\Lambda\subset\mathbb R^2\) 是满格，则平环面 \(\mathbb R^2/\Lambda\) 的覆盖空间仍为 \(\mathbb R^2\)，可取 \(\operatorname{dev}(z)=z\)、\(\rho(\lambda)(z)=z+\lambda\)。格向量记录两个方向的平移周期，不能从 \(K=0\) 中恢复。

\ProseParagraphBreak

即使完整Euclid 单值 \(\rho\) 平凡，也只说明展开映射在各覆盖片上的取值相同。由式\eqref{g1:eq:euclidean-developing-equivariance}，此时可定义
\[
	\overline{\operatorname{dev}}(p(\widetilde x))
	=\operatorname{dev}(\widetilde x).
\]
同一纤维的点相差覆盖变换，故定义无关；在均匀覆盖邻域上使用一个局部逆即可看出它光滑且为局部等距。这个结论没有保证 \(\overline{\operatorname{dev}}\) 单射，也没有保证它覆盖整个 \(\mathbb R^2\)。

\begin{example}[平凡单值不保证整体等价于平面]\label{g1:exm:trivial-monodromy-incomplete-model}
	单位开圆盘 \(B\subset\mathbb R^2\) 带Euclid 度量时，万有覆盖就是自身，\(\operatorname{dev}:B\hookrightarrow\mathbb R^2\) 为包含映射，\(\rho\) 平凡。它是单值且单射的局部等距，但不满射。圆盘凸，任意两点之间的Euclid 线段仍在圆盘内，所以内在距离等于Euclid 距离，直径为 \(2\)；整个Euclid 平面的直径无界。因此圆盘也不可能通过另一个映射整体等距于平面。

	单射性也可能失败。取 \(N=(0,\infty)\times\mathbb R\)，赋予度量 \(\mathrm{d}r^2+r^2\mathrm{d}t^2\)，并令
	\[
		D(r,t)=(r\cos t,r\sin t).
	\]
	其两个偏导线性无关，且 \(D^*(\mathrm{d}x^2+\mathrm{d}y^2)=\mathrm{d}r^2+r^2\mathrm{d}t^2\)，所以 \(D\) 是局部等距。\(N\) 是凸开集，故单连通；把它作为底流形时，万有覆盖可取恒等映射，单值仍平凡。然而 \(D(r,t+2\pi)=D(r,t)\)，所以展开映射并不单射，像也只是穿孔平面。
\end{example}

在这些例子中，局部结构方程确定邻域中的模型，单值确定绕行后的识别。当单值平凡、展开映射下降到底流形后，所得局部等距若再是到整个模型的双射，便给出整体同构。因而还需分别考察它的单射性与像。

\begin{chaptersummary}
	一个 \(G\)-结构规定了每点允许的标架，典范形式把标架坐标同真实切向量联系起来。相容联络的差是张量，它通过交替化改变挠率；在相应商空间中留下的类就是内蕴挠率。正交结构的内蕴挠率为零且无挠相容联络唯一，但其曲率可以非零。几乎复和接触模型则显示，适应标架的括号可能产生不能由联络选择消去的分量。

	\ProseParagraphBreak

	Cartan 联络在主 \(H\) 丛的每个切空间上给出到模型Lie 代数 \(\mathfrak l\) 的同构，既再现竖直方向，又记录底空间方向；它的核为零。Ehresmann 主联络则用 \(\mathfrak h\) 值形式的核选择水平空间。只有在给出相容的分解后，才能把 Cartan 形式拆成主联络部分与切向部分。曲率为零时，结构方程可以局部积分成齐性模型，路径展开进一步描述了这种比较怎样沿路延伸。

	\ProseParagraphBreak

	等价问题需要同时处理标架自由和结构函数。在曲面上，非零的曲率梯度能够选出一组规范的正交标架，随后由联络系数及其变化检验度量；曲率函数相同本身仍不足够。这种归一化有明确的适用区域，不能直接跨过梯度为零的点。平圆柱的周期及模型中的开子域则表明，局部模型、整体单值和全局同构是不同层次的结论。
\end{chaptersummary}

\BookExercises

\begin{exercise}
	\label{g1:ex:intrinsic-torsion-global-frame}
	在 \(\mathbb R^2\) 上取标架 \(e_1=\partial_x\)、\(e_2=e^x\partial_y\)。将这个指定标架视为 \(G=\{I\}\) 的结构群约化。
	\begin{enumerate}
		\item 求对偶余标架以及唯一适应联络的挠率。
		\item 证明不存在任何非空开集上的坐标 \((u,v)\)，使同时有 \(e_1=\partial_u\)、\(e_2=\partial_v\)。
		\item 将第二个标架向量改为 \(\widetilde e_2=e^{-x}e_2\)，所得标架的内蕴挠率为零。这是否与内蕴挠率不随适应标架改变相矛盾？
	\end{enumerate}
\end{exercise}
\begin{solution}
	(1) 对偶余标架为 \(\vartheta^1=\mathrm{d}x\)、\(\vartheta^2=e^{-x}\mathrm{d}y\)。由于 \(\mathfrak g=0\)，唯一适应联络满足 \(\nabla e_1=\nabla e_2=0\)。其挠率分量为
	\[
		\mathrm{d}\vartheta^1=0,\qquad
		\mathrm{d}\vartheta^2=-e^{-x}\mathrm{d}x\wedge\mathrm{d}y
		=-\vartheta^1\wedge\vartheta^2.
	\]
	因此 \(T(e_1,e_2)=-e_2\)。这里 \(C=V^*\otimes\mathfrak g=0\)，\(\operatorname{im}\delta=0\)，所以这一挠率本身就是内蕴挠率，处处非零。

	(2) 直接计算 \([e_1,e_2]=e^x\partial_y=e_2\)，处处不为零。而同一坐标系中的两个坐标向量场必定对易，因此不可能在任何非空开集同时满足题设两式。

	(3) 新标架为 \((\partial_x,\partial_y)\)，对偶形式均闭，其唯一适应联络无挠。但从旧标架到新标架的矩阵为 \(\operatorname{diag}(1,e^{-x})\)，并非恒等群 \(\{I\}\) 值函数，因此不是同一约化内允许的适应标架变换。题目改变了 \(G\)-结构本身，内蕴挠率可以随之改变。
\end{solution}

\begin{exercise}
	\label{g1:ex:conformal-coframe-torsion-correction}
	设 \(U\subset\mathbb R^2\) 为开集，\(f\in C^\infty(U)\)，在 \(U\) 上取度量
	\[
		g=e^{2f}(\mathrm{d}x^2+\mathrm{d}y^2)
	\]
	及正交单位余标架 \(\vartheta^1=e^f\mathrm{d}x\)、\(\vartheta^2=e^f\mathrm{d}y\)。在其对偶标架中，把度量联络矩阵写为
	\[
		a=\begin{pmatrix}0&-\eta\\\eta&0\end{pmatrix}.
	\]
	从无挠结构方程求出一形式 \(\eta\)，并在 \(f(x,y)=x^2+y\) 时写出结果。说明无挠为何不要求这两个余标架分量分别闭。
\end{exercise}
\begin{solution}
	写 \(\eta=p\,\mathrm{d}x+q\,\mathrm{d}y\)。外微分为
	\[
		\mathrm{d}\vartheta^1=-e^ff_y\,\mathrm{d}x\wedge\mathrm{d}y,
		\qquad
		\mathrm{d}\vartheta^2=e^ff_x\,\mathrm{d}x\wedge\mathrm{d}y.
	\]
	方程 \(\mathrm{d}\vartheta+a\wedge\vartheta=0\) 的两个分量分别为
	\[
		\mathrm{d}\vartheta^1-\eta\wedge\vartheta^2
		=-e^f(f_y+p)\,\mathrm{d}x\wedge\mathrm{d}y=0,
	\]
	\[
		\mathrm{d}\vartheta^2+\eta\wedge\vartheta^1
		=e^f(f_x-q)\,\mathrm{d}x\wedge\mathrm{d}y=0.
	\]
	故唯一解为 \(p=-f_y\)、\(q=f_x\)，即
	\[
		\eta=-f_y\,\mathrm{d}x+f_x\,\mathrm{d}y.
	\]
	当 \(f=x^2+y\) 时，\(\eta=-\mathrm{d}x+2x\,\mathrm{d}y\)。此时 \(\mathrm{d}\vartheta^1\) 处处非零，但它恰被 \(-\eta\wedge\vartheta^2\) 抵消。无挠要求 \(\mathrm{d}\vartheta+a\wedge\vartheta=0\)，所以余标架变化与联络作用要合在一起；只有联络矩阵也为零时，无挠才要求各余标架分量闭。
\end{solution}

\begin{exercise}
	\label{g1:ex:almost-complex-variable-obstruction}
	设 \(I\subset\mathbb R\) 是非空开区间，\(f\in C^\infty(I)\)。在 \(I\times\mathbb R^3\) 上取
	\[
		e_1=\partial_x,\quad e_2=\partial_y,\quad
		e_3=\partial_z,\quad e_4=\partial_t+f(x)\partial_y,
	\]
	并按 \(Je_1=e_2\)、\(Je_3=e_4\)、\(J^2=-I\) 定义几乎复结构。
	\begin{enumerate}
		\item 计算 \(\mathcal K_J(e_1-\mathrm i e_2,e_3-\mathrm i e_4)\)。
		\item 证明这个结构在整个 \(I\times\mathbb R^3\) 上有无挠适应联络，当且仅当 \(f\) 为常数。
		\item 当 \(f=c\) 为常数时，给出实坐标，使 \(J\) 在坐标标架中的矩阵为标准常矩阵。
	\end{enumerate}
\end{exercise}
\begin{solution}
	(1) 唯一可能非零的基本括号为 \([e_1,e_4]=f'(x)e_2\)。令 \(Z=e_1-\mathrm i e_2\)、\(W=e_3-\mathrm i e_4\)，则
	\[
		[Z,W]=-\mathrm i f'(x)e_2
		=\frac{f'(x)}2(Z-\overline Z),
	\]
	故 \(\mathcal K_J(Z,W)=-\tfrac12 f'(x)\overline Z\)。

	(2) 若有无挠适应联络，由\cref{g1:prop:almost-complex-intrinsic-component}，上一式处处为零，所以在 \(I\) 上有 \(f'=0\)。区间连通，故 \(f\) 为常数。反之，若 \(f=c\)，所有基本括号都为零。规定 \(\nabla e_i=0\)，得到保持 \(J\) 的联络，且 \(T(e_i,e_j)=-[e_i,e_j]=0\)。由于挠率是张量，这就说明其在所有输入上为零。

	(3) 取新坐标 \((u,v,w,s)=(x,y-ct,z,t)\)，逆变换为 \((x,y,z,t)=(u,v+cs,w,s)\)。因此
	\[
		\partial_u=e_1,\quad \partial_v=e_2,\quad
		\partial_w=e_3,\quad \partial_s=e_4.
	\]
	在此坐标中，\(J\partial_u=\partial_v\)、\(J\partial_w=\partial_s\)，故确为标准常矩阵。这个结论来自本模型的显式换坐标，不需要调用一般的几乎复可积性定理。
\end{solution}

\begin{exercise}
	\label{g1:ex:contact-form-scaling-intrinsic}
	在 \(\mathbb R^3\) 上令 \(\alpha=\mathrm{d}z-x\,\mathrm{d}y\)、\(D=\ker\alpha\)，并令 \(h\) 为处处非零光滑实函数，\(\widetilde\alpha=h\alpha\)。
	\begin{enumerate}
		\item 证明两一形式定义同一个 \(G_D\)-结构，并计算 \(\widetilde\alpha\wedge\mathrm{d}\widetilde\alpha\)。
		\item 对 \(X,Y\in\Gamma(D)\)，证明
		\[
			\alpha([X,Y])=-\mathrm{d}\alpha(X,Y),\qquad
			\widetilde\alpha([X,Y])=h\alpha([X,Y]).
		\]
		解释为什么右边的倍数变化没有改变商丛值括号张量 \(\mathcal B\)。
		\item 设 \((A,B,C)\) 为正文中的标架。求将它改成满足 \(u^*\widetilde\alpha=\varepsilon^3\) 的标架的一种方法，并说明该变换是否保持严格条件 \(u^*\alpha=\varepsilon^3\)。
	\end{enumerate}
\end{exercise}
\begin{solution}
	(1) \(h\) 处处非零，故 \(\ker\widetilde\alpha=\ker\alpha=D\)。\(P_D\) 只由标架的前两列是否张成 \(D\) 决定，所以约化相同。又
	\[
		\mathrm{d}\widetilde\alpha=\mathrm{d}h\wedge\alpha+h\,\mathrm{d}\alpha,
	\]
	从而
	\[
		\widetilde\alpha\wedge\mathrm{d}\widetilde\alpha
		=h^2\alpha\wedge\mathrm{d}\alpha
		=-h^2\mathrm{d}x\wedge\mathrm{d}y\wedge\mathrm{d}z\neq0.
	\]
	其中含两次 \(\alpha\) 的项为零。

	(2) 由一形式外微分公式及 \(\alpha(X)=\alpha(Y)=0\)，有 \(\mathrm{d}\alpha(X,Y)=-\alpha([X,Y])\)。第二式直接由 \(\widetilde\alpha=h\alpha\) 得出。这里 \(\mathcal B(X,Y)\) 是 \([X,Y]\) 在 \(\mathrm{T}M/D\) 中的类，从未用 \(\alpha\) 选择数值坐标。\(\alpha\) 诱导同构 \(\mathrm{T}M/D\to\mathbb R\)，\(\widetilde\alpha\) 则把这个同构乘以 \(h\)。变的是同一个商丛元素的标量表示，商丛值张量本身不变。

	(3) 取 \((A,B,C/h)\)。前两列仍属于 \(D\)，且 \(\widetilde\alpha(C/h)=1\)，所以满足新形式的严格条件。相对于原标架，变换矩阵为 \(\operatorname{diag}(1,1,h^{-1})\)，属于 \(G_D\)。但 \(\alpha(C/h)=h^{-1}\)，故它只有在 \(h=1\) 的点才保持原来的严格条件；要在整个区域保持，则必须有 \(h\equiv1\)。这正体现了保留分布与保留一个指定一形式的区别。
\end{solution}

\begin{exercise}\label{g1:ex:cartan-trivial-bundle}
	设 \(L\) 为有限维实Lie 群，\(H\subseteq L\) 为闭子群，\(U\) 为无边界光滑流形。给定 \(B\in\Omega^1(U;\mathfrak l)\)，在主右 \(H\) 丛 \(P=U\times H\) 上定义
	\[
		\omega_{(x,h)}(v,\eta)
		=\operatorname{Ad}_{h^{-1}}B_x(v)+(\theta_H)_h(\eta).
	\]
	证明：\(\omega\) 为模型 \(L/H\) 的 Cartan 联络，当且仅当每个 \(x\in U\) 处
	\[
		B_x\bmod\mathfrak h:\mathrm{T}_xU\longrightarrow\mathfrak l/\mathfrak h
	\]
	为同构。并证明，在这一条件下，\(\omega\) 平坦当且仅当
	\[
		\mathrm{d}B+\tfrac12[B\wedge B]=0.
	\]
	这是否要求 \(B_x:\mathrm{T}_xU\to\mathfrak l\) 本身为同构？
\end{exercise}
\begin{solution}
	固定 \(k\in H\)。右作用为 \((x,h)k=(x,hk)\)，Maurer--Cartan 形式的右平移律给出
	\[
		\begin{aligned}
			(R_k^*\omega)_{(x,h)}(v,\eta)
			&=\operatorname{Ad}_{(hk)^{-1}}B_x(v)
			+(\theta_H)_{hk}((\mathrm{d}R_k)_h\eta)\\
			&=\operatorname{Ad}_{k^{-1}}
			\bigl(\operatorname{Ad}_{h^{-1}}B_x(v)+(\theta_H)_h(\eta)\bigr).
		\end{aligned}
	\]
	故等变性对任意 \(B\) 成立。基本向量为 \((0,X_h^L)\)，代入即得 \(\omega(0,X_h^L)=X\)，故垂直再现也总成立。

	先假定商映射为同构。若 \(\omega(v,\eta)=0\)，取模 \(\mathfrak h\)，得
	\[
		\operatorname{Ad}_{h^{-1}}(B_x(v)+\mathfrak h)=0,
	\]
	从而 \(v=0\)。随后 \(\theta_H(\eta)=0\) 给出 \(\eta=0\)。商同构还给出 \(\dim U+\dim H=\dim L\)，所以 \(\omega\) 为同维空间之间的单射，因而是同构。

	反之，若 \(\omega\) 为 Cartan 联络，在 \((x,e)\) 处它是
	\[
		(v,X)\longmapsto B_x(v)+X
	\]
	这个线性同构。任意 \(Y\in\mathfrak l\) 可唯一如此表示，取商说明 \(B_x\bmod\mathfrak h\) 满射。若 \(B_x(v)\in\mathfrak h\)，则 \((v,-B_x(v))\) 被映到零，同构性使 \(v=0\)，故商映射也单射。

	令 \(s(x)=(x,e)\)，则 \(s^*\omega=B\)，从而
	\[
		s^*\Omega=\mathrm{d}B+\tfrac12[B\wedge B].
	\]
	若 \(\Omega=0\)，右端当然为零。反过来，若 \(s^*\Omega=0\)，则两个输入都沿截面切向时，曲率的取值为零；水平性使插入任何垂直分量的项也为零。由于截面方向与垂直方向张成 \(\mathrm{T}_{(x,e)}P\)，曲率在 \(s(U)\) 上为零；右等变性再将零传到 \((x,h)\) 的每一点。

	\(B_x\) 本身不必、通常也不可能满射到 \(\mathfrak l\)：若 \(\dim H>0\)，源维数比目标小 \(\dim H\)。它只需在取商后为同构；缺少的 \(\mathfrak h\) 方向由 \(\theta_H\) 提供。
\end{solution}

\begin{exercise}\label{g1:ex:cartan-spherical-model}
	沿用\cref{g1:exm:cartan-euclidean-surface}的定向二维Riemann 流形及正向正交标架丛 \(P\)，保留 \(\alpha,\theta^1,\theta^2,J_0\) 的约定。现在取
	\[
		L=\operatorname{SO}(3),\qquad
		H=\{\operatorname{diag}(h,1):h\in\operatorname{SO}(2)\},
	\]
	并在同一个 \(P\) 上定义 \(\mathfrak{so}(3)\) 值形式
	\[
		\widehat\omega=
		\begin{pmatrix}
			0&-\alpha&\theta^1\\
			\alpha&0&\theta^2\\
			-\theta^1&-\theta^2&0
		\end{pmatrix}.
	\]
	证明它是模型为单位球面 \(\operatorname{SO}(3)/\operatorname{SO}(2)\) 的 Cartan 联络，直接计算其曲率，并判定它何时平坦。解释为什么单位球面的这个 Cartan 几何平坦，与该球面的Euclid  Cartan 几何不平坦并不矛盾。
\end{exercise}
\begin{solution}
	先将 \(P\) 的结构群通过 \(h\mapsto\operatorname{diag}(h,1)\) 识别为给定的 \(H\)。记 \(A=J_0\alpha\) 和 \(\theta=(\theta^1,\theta^2)^{\mathsf T}\)，则
	\[
		\widehat\omega=
		\begin{pmatrix}A&\theta\\-\theta^{\mathsf T}&0\end{pmatrix}.
	\]
	\(\widehat\omega\) 反对称，因而取值于 \(\mathfrak{so}(3)\)。主联络满足 \(R_h^*A=h^{-1}Ah\)，焊接形式满足 \(R_h^*\theta=h^{-1}\theta\)。由于 \(h^{-1}=h^{\mathsf T}\)，整个分块矩阵满足
	\[
		R_h^*\widehat\omega
		=\operatorname{diag}(h,1)^{-1}\widehat\omega
		\operatorname{diag}(h,1).
	\]
	垂直输入时 \(\theta=0\)，而 \(A\) 再现 \(\mathfrak{so}(2)\)，故再现公理成立。若 \(\widehat\omega(V)=0\)，先由 \(\theta(V)=0\) 知 \(V\) 垂直，再由 \(A(V)=0\) 得 \(V=0\)。\(\dim P=\dim\mathfrak{so}(3)=3\)，所以逐点为同构。

	\(\operatorname{SO}(3)\) 在单位球面上传递，固定北极的稳定子恰为上述 \(H\)。由\cref{g1:prop:transitive-action-model}，商模型是单位球面。矩阵Lie 代数中的曲率为
	\(\widehat\Omega=\mathrm{d}\widehat\omega+\widehat\omega\wedge\widehat\omega\)。逐项相乘得
	\[
		\widehat\Omega=
		\begin{pmatrix}
			0&-\mathrm{d}\alpha-\theta^1\wedge\theta^2&
			\mathrm{d}\theta^1-\alpha\wedge\theta^2\\
			\mathrm{d}\alpha+\theta^1\wedge\theta^2&0&
			\mathrm{d}\theta^2+\alpha\wedge\theta^1\\
			-\mathrm{d}\theta^1+\alpha\wedge\theta^2&
			-\mathrm{d}\theta^2-\alpha\wedge\theta^1&0
		\end{pmatrix}.
	\]
	无挠结构方程使最后一列和最后一行的非对角项为零。又有 \(\mathrm{d}\alpha=-(K\circ\pi)\theta^1\wedge\theta^2\)，所以
	\[
		\widehat\Omega=
		\begin{pmatrix}
			0&(K\circ\pi-1)\theta^1\wedge\theta^2&0\\
			(1-K\circ\pi)\theta^1\wedge\theta^2&0&0\\
			0&0&0
		\end{pmatrix}.
	\]
	焊接形式在底空间方向给出余标架，故 \(\theta^1\wedge\theta^2\) 在这些方向上非零。因此球面模型的 Cartan 曲率为零，当且仅当 \(K\equiv1\)。

	同一底空间和同一主 \(\operatorname{SO}(2)\) 丛可以承载不同模型的 Cartan 联络。Euclid 模型中两个平移方向的括号为零，曲率的旋转部分是 \(J_0\,\mathrm{d}\alpha\)；球面模型中对应方向的括号产生附加项 \(J_0\theta^1\wedge\theta^2\)。单位球面与球面模型局部相同，所以后者平坦；它与Euclid 平面局部不等距，所以前者不平坦。平坦性所比较的模型已经改变。
\end{solution}

\begin{exercise}\label{g1:ex:cartan-development-path-dependence}
	取加法Lie 群 \(L=\mathbb R^2\)、平凡子群 \(H=\{0\}\)，于是主丛总空间就是 \(P=M=\mathbb R^2\)。定义
	\[
		\omega=\begin{pmatrix}\mathrm{d}x\\e^x\,\mathrm{d}y\end{pmatrix}.
	\]
	证明这给出一个非平坦 Cartan 几何。比较从 \((0,0)\) 到 \((1,1)\) 的两条折线路径：第一条先沿 \(x\) 轴到 \((1,0)\)，再沿竖直方向到达终点；第二条先沿 \(y\) 轴到 \((0,1)\)，再沿水平方向到达终点。求各自展开路径的终点，说明在单连通底空间上，曲率不为零时为何仍不能只由路径终点定义展开映射。
\end{exercise}
\begin{solution}
	在坐标基下，\(\omega\) 的矩阵为 \(\operatorname{diag}(1,e^x)\)，行列式始终为正，所以逐点可逆。结构群平凡，等变性和再现条件没有额外要求，故 \(\omega\) 是 Cartan 联络。加法Lie 代数的括号为零，因而
	\[
		\Omega=\mathrm{d}\omega
		=\begin{pmatrix}0\\e^x\,\mathrm{d}x\wedge\mathrm{d}y\end{pmatrix}\ne0.
	\]
	模型上的 Maurer--Cartan 形式是恒等余标架。展开 \(g=(u,v)\) 满足
	\[
		u'=x',\qquad v'=e^{x(t)}y',\qquad
		(u(0),v(0))=(0,0).
	\]
	故 \(u(1)=1\)，而 \(v(1)=\int_\gamma e^x\,\mathrm{d}y\)。

	沿第一条路径，水平段贡献为零，竖直段 \(x=1\)，所以 \(v(1)=e\)，展开终点是 \((1,e)\)。沿第二条路径，竖直段 \(x=0\)，随后水平段没有贡献，所以展开终点是 \((1,1)\)。两条路径在 \(\mathbb R^2\) 中定端点同伦，却有不同的展开终点；单连通性不能消除曲率造成的路径依赖。

	差值还可以直接由曲率看出。将第一条路径与第二条的反向拼接，得到单位正方形的正向边界；Stokes 公式给出
	\[
		e-1
		=\int_{\partial[0,1]^2}e^x\,\mathrm{d}y
		=\int_{[0,1]^2}e^x\,\mathrm{d}x\wedge\mathrm{d}y.
	\]
	在这个交换模型中，路径展开的差恰好由曲率的面积积分表示。
\end{solution}

\begin{exercise}\label{g1:ex:cartan-one-dimensional-completeness}
	取 \(L=\mathbb R\)、\(H=\{0\}\)，并在 \(M=\mathbb R\) 上给定 \(\omega=f(x)\,\mathrm{d}x\)，其中 \(f:\mathbb R\to(0,\infty)\) 光滑。证明这总是平坦 Cartan 联络。令
	\[
		F(x)=\int_0^x f(s)\,\mathrm{d}s.
	\]
	求从 \(x_0\) 出发的路径展开，刻画 \(F\) 的像，并证明由 \(\omega(E)=1\) 定义的向量场 \(E\) 完备，当且仅当
	\[
		\int_0^\infty f(s)\,\mathrm{d}s=\infty,
		\qquad
		\int_{-\infty}^0 f(s)\,\mathrm{d}s=\infty.
	\]
	分别讨论 \(f=1\) 与 \(f=e^x\)，说明即使底空间是单连通的整条实直线，平坦性也不保证所有模型运动都能无限延拓。
\end{exercise}
\begin{solution}
	\(f>0\) 保证 \(\omega\) 逐点可逆，另外两条 Cartan 公理在平凡结构群下自动成立。加法Lie 代数的括号为零，而一维流形上二形式为零，故 \(\Omega=\mathrm{d}(f\,\mathrm{d}x)=0\)。

	若 \(\gamma(t)=x(t)\) 且 \(x(0)=x_0\)，展开方程为 \(g'=f(x)x'\)、\(g(0)=0\)。由链式法则，
	\[
		g(t)=F(x(t))-F(x_0).
	\]
	因此此处展开确实只依赖当前终点。又因为 \(F'=f>0\)，\(F\) 严格递增且处处为局部微分同胚。记扩展实数
	\[
		a=-\int_{-\infty}^0 f(s)\,\mathrm{d}s,\qquad
		b=\int_0^\infty f(s)\,\mathrm{d}s.
	\]
	单调性及积分定义给出 \(\lim_{x\to-\infty}F(x)=a\)、\(\lim_{x\to+\infty}F(x)=b\)。由介值定理，\(F\) 将 \(\mathbb R\) 双射到 \((a,b)\)；局部光滑逆拼接成光滑全局逆，故它是微分同胚。

	由 \(\omega(E)=1\) 得 \(E=f^{-1}\partial_x\)。其过 \(x_0\) 的积分曲线满足
	\[
		F(x(t))=F(x_0)+t,\qquad
		x(t)=F^{-1}(F(x_0)+t).
	\]
	该式恰在
	\[
		a-F(x_0)<t<b-F(x_0)
	\]
	上有定义。若积分曲线能在 \(\mathbb R\) 中延拓到某个有限端点，连续性会使该端点的 \(F\) 值属于 \((a,b)\)，与它等于 \(a\) 或 \(b\) 矛盾，所以这是最大定义区间。于是 \(E\) 的所有积分曲线在整个实时间轴上存在，当且仅当 \(a=-\infty\)、\(b=\infty\)，即题述两个积分均发散。

	当 \(f=1\) 时，\(F(x)=x\)，\(E=\partial_x\) 完备。当 \(f=e^x\) 时，
	\[
		F(x)=e^x-1,\qquad
		F(\mathbb R)=(-1,\infty),\qquad
		x(t)=\log(e^{x_0}+t).
	\]
	这条积分曲线只在 \(t>-e^{x_0}\) 上存在，故 \(E=e^{-x}\partial_x\) 不完备。底空间的拓扑与平坦性均没有改变，但模型中的匀速负向运动在有限时间后离开了 \(F(\mathbb R)\)。
\end{solution}

\begin{exercise}\label{g1:ex:rotating-plane-frame-gauge}
	在定向Euclid 平面上从标架 \((\partial_x,\partial_y)\) 出发，将标架在点 \((x,y)\) 处旋转角度 \(x\)。求新的余标架 \(\theta^1,\theta^2\) 和联络一形式 \(\alpha\)，并直接核验
	\[
		\mathrm{d}\theta^1=\alpha\wedge\theta^2,\qquad
		\mathrm{d}\theta^2=-\alpha\wedge\theta^1,\qquad
		\mathrm{d}\alpha=0.
	\]
	说明非零的 \(\alpha\) 为何不表示度量发生改变，也不表示曲率非零。
\end{exercise}
\begin{solution}
	旋转后的标架为
	\[
		e_1=\cos x\,\partial_x+\sin x\,\partial_y,\qquad
		e_2=-\sin x\,\partial_x+\cos x\,\partial_y.
	\]
	对偶余标架按旋转矩阵的逆变换，得到
	\[
		\theta^1=\cos x\,\mathrm{d}x+\sin x\,\mathrm{d}y,\qquad
		\theta^2=-\sin x\,\mathrm{d}x+\cos x\,\mathrm{d}y.
	\]
	原标架的联络一形式为零，换架公式给出 \(\alpha=\mathrm{d}x\)。也可直接计算 \(\nabla_Ve_1=\mathrm{d}x(V)e_2\) 得到相同结果。求外微分，
	\[
		\begin{aligned}
			\mathrm{d}\theta^1&=\cos x\,\mathrm{d}x\wedge\mathrm{d}y
			=\mathrm{d}x\wedge\theta^2,\\
			\mathrm{d}\theta^2&=-\sin x\,\mathrm{d}x\wedge\mathrm{d}y
			=-\mathrm{d}x\wedge\theta^1,
		\end{aligned}
	\]
	而 \(\mathrm{d}\alpha=\mathrm{d}(\mathrm{d}x)=0\)。因此挠率与曲率仍为零。

	将两个余标架平方相加，混合项抵消，得到 \((\theta^1)^2+(\theta^2)^2=\mathrm{d}x^2+\mathrm{d}y^2\)。度量没有改变；非零的 \(\alpha\) 记录所选标架随位置旋转的速率。曲率比较不同方向上的变化能否相容，这里由 \(\mathrm{d}\alpha=0\) 表示，不能仅由 \(\alpha\) 是否为零判断。
\end{solution}

\begin{exercise}\label{g1:ex:normalized-curvature-hessian}
	设 \((U,h)\) 为定向光滑Riemann 度量曲面，\(\mathrm{d}K\) 处处非零。使用式\eqref{g1:eq:curvature-gradient-frame}的归一化标架，并写 \(\alpha=a\theta^1+b\theta^2\)。对
	\[
		\operatorname{Hess}K(V,W)=h(\nabla_V\operatorname{grad}_hK,W),
	\]
	证明其在 \(e_1,e_2\) 下的矩阵为
	\[
		\begin{pmatrix}e_1(q)&e_2(q)\\e_2(q)&qb\end{pmatrix}.
	\]
	再证明曲率水平集沿单位切向量 \(e_2\) 定向时，其有向测地曲率等于 \(b\)。这里采用第九章的约定 \(k_g=h(\nabla_TT,JT)\)。
\end{exercise}
\begin{solution}
	由 \(\operatorname{grad}_hK=qe_1\) 与联络公式，
	\[
		\nabla_{e_1}\operatorname{grad}_hK=e_1(q)e_1+qa e_2,
		\qquad
		\nabla_{e_2}\operatorname{grad}_hK=e_2(q)e_1+qb e_2.
	\]
	分别与 \(e_1,e_2\) 作内积，先得到矩阵
	\[
		\begin{pmatrix}e_1(q)&qa\\e_2(q)&qb\end{pmatrix}.
	\]
	由式\eqref{g1:eq:normalized-surface-coefficients}中的 \(qa=e_2(q)\)，即得所求对称矩阵。也可以由无挠性直接检查对称性：
	\[
		\operatorname{Hess}K(V,W)-\operatorname{Hess}K(W,V)
		=VW(K)-WV(K)-[V,W](K)=0.
	\]
	这里展开 Hessian 时用了 \(\operatorname{Hess}K(V,W)=V(WK)-(\nabla_VW)K\)。

	因为 \(\mathrm{d}K\ne0\)，每个曲率水平集在局部是一条光滑曲线；\(e_2(K)=0\) 且 \(|e_2|_h=1\)，故可按 \(e_2\) 定向并用弧长参数化。其切向量 \(T=e_2\) 满足
	\[
		\nabla_TT=\nabla_{e_2}e_2=-b e_1,\qquad JT=Je_2=-e_1.
	\]
	于是 \(k_g=h(-be_1,-e_1)=b\)。因此归一化后的系数 \(b\) 同时记录 Hessian 的横向分量和曲率水平线的有向测地弯曲。
\end{solution}

\begin{exercise}\label{g1:ex:warped-surface-local-symmetries}
	在半平面 \(x>0\) 上赋予度量 \(h=\mathrm{d}x^2+e^{2x^2}\mathrm{d}y^2\) 及定向 \(\mathrm{d}x\wedge\mathrm{d}y>0\)。证明任意保持定向的局部自等距，在每个足够小的连通邻域内都形如
	\[
		\Phi(x,y)=(x,y+c)
	\]
	其中 \(c\) 为常数；反之，这些映射确实是等距。再求反转定向的局部自等距。
\end{exercise}
\begin{solution}
	\cref{g1:exm:normalized-warped-surface}已计算 \(K=-2-4x^2\)，在 \(x>0\) 上严格单调。若 \(\Phi=(X,Y)\) 是局部自等距，曲率保持性给出 \(-2-4X^2=-2-4x^2\)。由于源、目标的第一坐标都为正，只能有 \(X=x\)。

	将其代入度量拉回式，得到
	\[
		\mathrm{d}x^2+e^{2x^2}(Y_x\mathrm{d}x+Y_y\mathrm{d}y)^2
		=\mathrm{d}x^2+e^{2x^2}\mathrm{d}y^2.
	\]
	比较 \(\mathrm{d}x^2\) 系数，\(Y_x=0\)；比较 \(\mathrm{d}y^2\) 系数，\(Y_y^2=1\)。在一个小连通邻域内，连续且不为零的 \(Y_y\) 符号恒定。Jacobian 行列式恰为 \(Y_y\)，所以保持定向时有 \(Y_y=1\)，从而 \(\mathrm{d}(Y-y)=0\)，得到 \(Y=y+c\)。反转定向时有 \(Y_y=-1\)，同理得到 \(Y=-y+c\)。

	反向代入即可核验：这两类映射都保持 \(x\) 和 \(\mathrm{d}y^2\)，故保持度量，Jacobian 的符号分别为正、负。局部计算因而确定了全部局部自等距，而非只找到其中一族。
\end{solution}

\begin{exercise}\label{g1:ex:cylinder-local-isometry-covering}
	给定 \(\ell>0\) 与整数 \(m\ge2\)。对平圆柱
	\[
		C_{m\ell}=\mathbb R^2/(m\ell\mathbb Z\times\{0\}),\qquad
		C_\ell=\mathbb R^2/(\ell\mathbb Z\times\{0\}),
	\]
	考察映射 \(F([x,y]_{m\ell})=[x,y]_\ell\)。证明它良定义，是保持定向的局部等距及 \(m\) 叶覆盖，但不是整体等距同构。计算源圆柱上绕行一周的环路映到目标后的绕行次数，并写出两侧展开单值的对应关系。说明即使两边的切丛平行移动单值都平凡，这个覆盖仍不可能是单射。
\end{exercise}
\begin{solution}
	源中的两个代表元相差 \((km\ell,0)\)，它们在目标中也相差一个允许的平移，所以 \(F\) 良定义。两个商投影的局部逆将 \(F\) 表示为Euclid 平面中的恒等映射，故 \(F\) 光滑、保持定向且局部等距。

	对目标中的 \([x,y]_\ell\)，所有原像恰为
	\[
		[x+j\ell,y]_{m\ell},\qquad j=0,1,\ldots,m-1.
	\]
	它们两两不同，因为 \((j-k)\ell\) 不是非零的 \(m\ell\) 整数倍。取目标中由宽度小于 \(\ell\) 的小矩形给出的坐标域，它在源中的逆像分成这 \(m\) 个互不相交的小矩形，每一片由 \(F\) 同胚地映到原坐标域。因此 \(F\) 是 \(m\) 叶覆盖，且 \(m\ge2\) 说明它不单射。

	源中一周环路可写为 \(t\mapsto[x_0+tm\ell,y_0]_{m\ell}\)、\(0\le t\le1\)。在目标中它的平面提升仍从 \((x_0,y_0)\) 走到 \((x_0+m\ell,y_0)\)，故绕行次数为 \(m\)。相应地，基本群上的同态为 \(n\mapsto mn\)。若两边都取恒等的平面展开映射，则
	\[
		\rho_{m\ell}(n)(u,v)=(u+nm\ell,v)
		=\rho_\ell(mn)(u,v).
	\]
	两侧都有由 \(\partial_x,\partial_y\) 下降的整体平行标架，所以切丛平行移动单值平凡；展开单值却分别记录周期 \(m\ell\) 与 \(\ell\)。它们的最短不可缩环路长度不同，由\cref{g1:prop:cylinder-period-length}，甚至不存在其他整体等距把这两个圆柱识别。即使存在这样的局部等距与有限叶覆盖，两个圆柱的整体周期仍有差异。
\end{solution}



\backmatter
\cleardoublepage
\pagenumbering{roman}
\bookmarksetup{startatroot}
% !TeX root = ../Book1.tex
\cleardoublepage
% !TeX root = ../Book1.tex
\chapter*{后记}
\addcontentsline{toc}{chapter}{后记}
回看本卷，最初的坐标变换已经有了几种不同的作用。在仿射与射影几何中，它帮助我们把对象化成标准形；在曲面和流形上，它把同一对象的局部表达联系起来；在向量丛与主丛中，它记录不同纤维标架怎样粘合。相似的矩阵公式由此承担了不同的几何任务。能够辨认公式正在改变什么、保留什么，是读完这些内容后值得保留的一种习惯。

\ProseParagraphBreak

局部计算与整体构造之间的距离，也在各编中逐渐显露。曲线的曲率、挠率等微分数据可以用来解出局部标架，曲线是否闭合还要检查走完一周后的结果；曲面的基本形式满足相容条件后，可以恢复局部浸入，沿不同道路延拓时又会遇到整体问题。到了平坦联络，曲率为零与环路平行移动非平凡可以同时发生。基本群和覆盖因此再次出现：早先用来辨别空间的工具，现在也记录几何数据的延拓。

\ProseParagraphBreak

本卷保留了古典几何、曲面和一般流形各自的篇幅。前者提供能够画出、计算和分类的对象，后者让这些经验适用于更广的空间。两者之间的联系有时要到第二次阅读才更清楚。例如，重新计算球面的曲率，可以同时观察局部坐标、移动标架与整体积分；重新看射影空间，可以把齐次坐标、商构造和群作用放在一起。读者也可以选定一个熟悉的空间，逐章积累它的不同描述，再尝试说明这些描述为何相容。

\ProseParagraphBreak

后续学习将沿本卷留下的具体问题展开。覆盖与曲面分类引出了更高维空间的不变量，第二卷用同调、上同调和同伦研究它们，并用横截性与临界点恢复光滑几何中的信息。度量、测地线与曲率引出了距离的整体比较和几何方程，第三卷继续处理这些问题。几乎复结构、微分形式和主丛中的相容条件，在第四卷的复几何、辛几何中有了新的内容；射影构造和由方程定义的几何对象，则在第五卷进入代数几何的讨论。这些方向各有进一步的基础要求，可以按所关心的问题接续阅读。

\ProseParagraphBreak

若希望继续研究本卷的专题，可以沿各章已给出的文献寻找更系统的处理：由曲面走向整体微分几何，由紧群走向表示论，由主丛与联络走向几何结构与规范理论。附录则保留了几条较短的延伸，例如用反射生成对称群、用有限构型研究多面体、用计算检验坐标公式。再次阅读时，选择一项能够独立完成的工作很有用：重建一个标准模型，算出一组过渡函数，或把一个局部解延拓到整个空间。遇到不能继续的地方，检查阻碍来自条件、拓扑还是计算，便能把下一步要学的问题说得更清楚。

\cleardoublepage
% !TeX root = ../Book1.tex
\PrintBookBibliography
% 仅在核读并引用真实资料后补入书目。

\cleardoublepage
% !TeX root = ../Book1.tex
\PrintBookIndex{chinese}{中文索引}

\cleardoublepage
% !TeX root = ../Book1.tex
\PrintBookIndex{foreign}{外文索引}

\cleardoublepage
% !TeX root = ../Book1.tex
\PrintBookIndex{symbols}{符号索引}


\end{document}
